\documentclass[11pt]{article}

\usepackage[T1]{fontenc}
\usepackage{lmodern}
\usepackage[margin=1in]{geometry}
\usepackage{amsmath,amssymb,amsthm,mathtools}
\usepackage{enumitem}
\usepackage{booktabs,longtable,array}
\usepackage{xcolor}
\usepackage[colorlinks=true,linkcolor=blue!55!black,urlcolor=blue!55!black]{hyperref}
\usepackage[nameinlink,noabbrev]{cleveref}

% Compatibility dispatcher for inherited comma-separated equation-reference lists.
\makeatletter
\let\renewal@singleeqref\eqref
\newcommand{\renewal@eqrefdispatch}[1]{%
  \in@{,}{#1}%
  \ifin@\labelcref{#1}\else\renewal@singleeqref{#1}\fi}
\renewcommand{\eqref}[1]{%
  \begingroup\renewal@eqrefdispatch{#1}\endgroup}
\makeatother

\newtheorem{theorem}{Theorem}[section]
\newtheorem{lemma}[theorem]{Lemma}
\newtheorem{proposition}[theorem]{Proposition}
\newtheorem{corollary}[theorem]{Corollary}
\newtheorem{counterexample}[theorem]{Exact separator}
\newtheorem{openproblem}[theorem]{Open problem}
\theoremstyle{definition}
\newtheorem{definition}[theorem]{Definition}
\newtheorem{assessment}[theorem]{Assessment}
% Compatibility declarations for environments first used in the appended V8 material.
\newtheorem{example}[theorem]{Example}
\newtheorem{firewall}[theorem]{Firewall}
\theoremstyle{remark}
\newtheorem{remark}[theorem]{Remark}

\newcommand{\C}{\mathbb C}
\newcommand{\R}{\mathbb R}
\newcommand{\Ran}{\operatorname{Ran}}
\newcommand{\Ker}{\operatorname{Ker}}
\newcommand{\Tr}{\operatorname{Tr}}
\newcommand{\rank}{\operatorname{rank}}
\newcommand{\HS}{\mathrm{HS}}
\newcommand{\op}{\mathrm{op}}
\newcommand{\dd}{\mathrm d}
\newcommand{\cF}{\mathcal F}
\newcommand{\cH}{\mathcal H}
\newcommand{\cS}{\mathcal S}
\newcommand{\cY}{\mathcal Y}
\newcommand{\id}{\mathrm{id}}

\title{Renewal Standard--Model--Einstein Continuum Research Notes\\
Version V1: Curvature of Moving Relation Supports}
\author{Append-only research continuation}
\date{September 4, 2026}

\begin{document}
\maketitle

\begin{abstract}
This is the first note in a separate append-only series whose long-term target
is the interacting Standard--Model--Einstein continuum described in the
supplied Renewal manuscripts.  It does not claim that continuum.  The latest
attached SM--Einstein continuation stops at a shrinking, common-action,
relation-covariant RPESM command square.  This note audits that stopping point
and proves that a second plaquette is necessary: the scalar likelihood
plaquette may be exactly flat while the moving returned-word support has
nonzero Kato curvature.  We derive the curvature, its finite polar-chord
holonomy, an exact rank-one separator, and the resulting area-scaled cutoff
requirement.  These results sharpen the next finite acquisition without
repeating the one-parameter support transport proved in the attached V30 note.
\end{abstract}

\tableofcontents

\section{Context, imported results, and scope}
\label{sec:context}

The sources audited for this note are the following five supplied files:
\begin{enumerate}[leftmargin=2.2em]
\item \texttt{Renewal\_Grand\_Tensor\_Monograph\_V067.tex};
\item \texttt{Renewal\_Einstein\_Standard\_Model\_Physical\_Source\_Metric\_Duhamel\_}
      \texttt{Ancestry\_and\_Quantum\_Continuum\_Programme\_V30.tex};
\item \texttt{Renewal\_Yang\_Mills\_Boundary\_Sector\_Cofinal\_Contraction\_Programme\_V29.tex};
\item \texttt{Renewal\_Yang\_Mills\_Reduced\_Fibre\_Wilson\_Shell\_Symbol\_and\_Local\_}
      \texttt{Vacuum\_Programme\_V30.tex};
\item \texttt{Renewal\_NS\_Terminal\_Source\_Concentration\_and\_Singular\_Thread\_}
      \texttt{Closure\_Programme\_V29.tex}.
\end{enumerate}
The repository manuscripts and their proof-triage ledgers were also checked.
Their existing finite-dimensional, probability, Schur-complement, continuum
handoff, Yang--Mills, and Navier--Stokes results are treated as imported and are
not reproved here.

The immediate input is the Version V30 terminal acquisition in the second
file.  It asks for a two-parameter shrinking RPESM command square carrying one
common unnormalized cylinder, returned-word support Grams at all corners,
direct responses, polar support chords, physical supported transport,
determinant/Pfaffian sewing, held-out reads, and derivative-scaled cutoff
rectangles.  Version V30 proves the one-parameter Kato identity
\[
 P_tU_t=U_tP_0,
 \qquad
 \left.\frac{\dd}{\dd t}\right|_{t=0}X_tU_t
 =\dot X_0P_0
\]
for a descended response $X_t=X_tP_t$, and proves scalar common-action
reconstruction from multiplicative likelihood plaquettes.

Those two results do not settle the compatibility of support transport in two
command directions.  The missing object is the curvature of the bundle
$\Ran P_\theta$.  This is distinct from the determinant-line Berry curvature
already retained in the V21 phase ledger: the present curvature belongs to the
returned-word relation quotient itself.  Searches of all five supplied files
found no theorem computing this Kato support curvature or its finite
polar-chord plaquette.

The attached sources also state that they do not print the selected RPESM
corner matrices or the numerical three-arm and command-square values.  Hence a
numerical population of the requested physical square cannot be reconstructed
from the attachments alone.  The rigorous step available upstream is to
identify every logically necessary row before new regulator data are acquired.

\paragraph{Append-only convention.}
Future versions of this notebook will retain this entire V1 text and append new
material before the sole final \verb|\end{document}|.  A proved finite lemma is
kept separate from a conditional continuum handoff, and a separator is not
reported as a construction of the physical continuum.

\section{The projection bundle and its Kato curvature}
\label{sec:curvature}

Let $\Theta\subset\R^d$ be open, let $E$ be a finite-dimensional complex
Hilbert space, and let
\[
 P:\Theta\longrightarrow\operatorname{End}(E)
\]
be a $C^2$ family of orthogonal projections of constant rank.  Write
$P_i=\partial_iP$ and $\cS_\theta=\Ran P_\theta$.  Constant rank gives a
$C^2$ Hermitian vector bundle $\cS\to\Theta$.

\begin{lemma}[Off-diagonal differential of a projection]
\label{lem:off-diagonal}
For every coordinate direction $i$,
\begin{equation}
 PP_iP=0,
 \qquad
 (I-P)P_i(I-P)=0,
 \qquad
 P_i=PP_i(I-P)+(I-P)P_iP.
\label{eq:off-diagonal}
\end{equation}
Thus $P_i$ is self-adjoint and exchanges $\Ran P$ and $\Ker P$.
\end{lemma}

\begin{proof}
Differentiate $P^2=P$ to obtain $P_iP+PP_i=P_i$.  Multiplication on the
left and right by $P$ gives $2PP_iP=PP_iP$, hence $PP_iP=0$.  Applying the
same argument to $I-P$ gives the second identity.  Expanding $P_i$ in the
orthogonal decomposition $E=\Ran P\oplus\Ker P$ gives the third.  Finally,
$P_i^*=P_i$ follows by differentiating $P^*=P$.
\end{proof}

\begin{definition}[Kato connection]
\label{def:Kato-connection}
For a $C^1$ section $s$ of $\cS$, define
\begin{equation}
 \nabla_i^{\rm K}s:=P\partial_i s.
\label{eq:Kato-connection}
\end{equation}
Its curvature is
\begin{equation}
 \cF_{ij}^{\rm K}
 :=[\nabla_i^{\rm K},\nabla_j^{\rm K}].
\label{eq:Kato-curvature-def}
\end{equation}
\end{definition}

\begin{theorem}[Exact curvature of a moving relation support]
\label{thm:Kato-curvature}
The Kato connection is Hermitian and its curvature is the supported
commutator
\begin{equation}
 \boxed{
 \cF_{ij}^{\rm K}
 =P[P_i,P_j]P\big|_{\Ran P}.}
\label{eq:Kato-curvature}
\end{equation}
In particular $(\cF_{ij}^{\rm K})^*=-\cF_{ij}^{\rm K}$.  It vanishes in
direction $(i,j)$ exactly when
\begin{equation}
 PP_i(I-P)P_jP=PP_j(I-P)P_iP.
\label{eq:flatness-criterion}
\end{equation}
\end{theorem}

\begin{proof}
For two sections $s,t$ of $\cS$, the identities $Ps=s$ and $Pt=t$ give
\[
 \partial_i\langle s,t\rangle
 =\langle P\partial_i s,t\rangle
  +\langle s,P\partial_i t\rangle,
\]
so the connection is Hermitian.  Next,
\begin{align*}
 [\nabla_i^{\rm K},\nabla_j^{\rm K}]s
 &=P\partial_i(P\partial_j s)-P\partial_j(P\partial_i s)\\
 &=PP_i\partial_j s-PP_j\partial_i s.
\end{align*}
Differentiate $s=Ps$ to write
$\partial_j s=P_j s+P\partial_j s$.  By
\cref{lem:off-diagonal}, $PP_iP=0$; hence
\[
 PP_i\partial_j s=PP_iP_jPs.
\]
Interchanging $i$ and $j$ proves \eqref{eq:Kato-curvature}.  The adjoint
claim follows because $P_i$ and $P_j$ are self-adjoint.  Inserting the
off-diagonal decomposition \eqref{eq:off-diagonal} into the commutator gives
\eqref{eq:flatness-criterion}.
\end{proof}

The physical responses in the V30 acquisition are operators which descend
through the current relation quotient.  Their induced connection has the
opposite curvature.

\begin{theorem}[Curvature seen by a relation-descended response]
\label{thm:response-curvature}
Let $\cY$ be a fixed finite-dimensional Hilbert space and let
$X_\theta:E\to\cY$ be $C^2$ with
\begin{equation}
 X_\theta=X_\theta P_\theta.
\label{eq:response-descent}
\end{equation}
Regard $X$ as a section of $\operatorname{Hom}(\cS,\Theta\times\cY)$ and
equip the target with its flat connection.  Then
\begin{equation}
 (\nabla_i^{\rm Hom}X)P=(\partial_iX)P
\label{eq:response-connection}
\end{equation}
and
\begin{equation}
 \boxed{
 [\nabla_i^{\rm Hom},\nabla_j^{\rm Hom}]X
 =-X\cF_{ij}^{\rm K}.}
\label{eq:response-curvature}
\end{equation}
Consequently, scalar action flatness does not make mixed aligned response
derivatives commute unless $X\cF_{ij}^{\rm K}=0$.
\end{theorem}

\begin{proof}
For a section $s$ of $\cS$, the induced connection is defined by
\[
 (\nabla_i^{\rm Hom}X)s
 =\partial_i(Xs)-X(\nabla_i^{\rm K}s).
\]
Using $X=XP$ gives
\[
 (\nabla_i^{\rm Hom}X)s
 = (\partial_iX)s+X(I-P)\partial_i s
 = (\partial_iX)s,
\]
which proves \eqref{eq:response-connection}.  The general curvature identity
for a Hom bundle now reduces to the source curvature with a minus sign.  It can
also be checked directly: apply the commutator to $Xs$ and use the flatness of
the target to get
\[
 ([\nabla_i^{\rm Hom},\nabla_j^{\rm Hom}]X)s
 =-X[\nabla_i^{\rm K},\nabla_j^{\rm K}]s.
\]
Then use \cref{thm:Kato-curvature}.
\end{proof}

\section{Finite polar chords and the support plaquette}
\label{sec:finite-holonomy}

For equal-rank orthogonal projections $P,Q$ with $\|P-Q\|_{\op}<1$, define
the polar chord already used in the attached V30 note by
\begin{equation}
 U_{Q\leftarrow P}:=QP(PQP)^{-1/2}\quad\hbox{on }\Ran P.
\label{eq:polar-chord}
\end{equation}
It is the canonical isometry from $\Ran P$ onto $\Ran Q$ with positive old
support overlap.

Fix $\theta\in\Theta$, two coordinate directions $i,j$, and a sufficiently
small $h>0$.  Put
\begin{align}
 T_{ij}(\theta,h)
 &:={U}_{P_{\theta+h(e_i+e_j)}\leftarrow P_{\theta+he_i}}
    {U}_{P_{\theta+he_i}\leftarrow P_\theta},
\label{eq:path-ij}\\
 T_{ji}(\theta,h)
 &:={U}_{P_{\theta+h(e_i+e_j)}\leftarrow P_{\theta+he_j}}
    {U}_{P_{\theta+he_j}\leftarrow P_\theta}.
\label{eq:path-ji}
\end{align}
Both are isometries from $\cS_\theta$ to
$\cS_{\theta+h(e_i+e_j)}$.

\begin{definition}[Support holonomy and debit]
\label{def:support-debit}
The finite polar support holonomy and its positive debit are
\begin{equation}
 \cH_{ij,h}:=T_{ji}^*T_{ij},
 \qquad
 \Delta_{ij,h}^{\rm supp}
 :=\|T_{ij}-T_{ji}\|_{\HS}^2.
\label{eq:support-debit}
\end{equation}
They are computed after the scalar common-action plaquette, because they live
on the relation-support bundle rather than on the positive density line.
\end{definition}

\begin{theorem}[Polar-chord holonomy converges to Kato curvature]
\label{thm:finite-holonomy}
If $P$ is $C^3$, then, as $h\downarrow0$,
\begin{align}
 \boxed{
 \cH_{ij,h}
 =P_\theta-h^2\cF_{ij}^{\rm K}(\theta)+O(h^3)}
 &\quad\text{on }\cS_\theta,
\label{eq:holonomy-expansion}\\
 \boxed{
 \frac{\Delta_{ij,h}^{\rm supp}}{h^4}
 \longrightarrow
 \|\cF_{ij}^{\rm K}(\theta)\|_{\HS}^2.}
\label{eq:debit-limit}
\end{align}
The $O(h^3)$ remainder is uniform on compact subsets on which the first three
derivatives of $P$ are bounded and the chord angle stays below one.
\end{theorem}

\begin{proof}
Choose a $C^3$ local orthonormal frame
$e(\theta):\C^r\to E$ for $\cS$, so $e^*e=I_r$ and $P=ee^*$.  Let
\begin{equation}
 A_i=e^*\partial_i e.
\label{eq:frame-connection}
\end{equation}
Differentiating $e^*e=I_r$ shows $A_i^*=-A_i$.  In these frames the polar
chord from $\theta$ to $\theta+he_i$ is the unitary polar factor of
$e(\theta+he_i)^*e(\theta)$.  Taylor expansion and smoothness of the polar
map near the identity give
\begin{equation}
 u_i(\theta,h)=I_r-hA_i(\theta)+O(h^2).
\label{eq:edge-expansion}
\end{equation}
More precisely, smoothness gives matrices $B_i(\theta)$ such that
\begin{equation}
 u_i(\theta,h)
 =I_r-hA_i(\theta)+h^2B_i(\theta)+O(h^3).
\label{eq:edge-second-order}
\end{equation}
Expanding the connection at the shifted starting point and multiplying in path
order yields
\begin{align*}
 t_{ij}
 &=I_r-h(A_i+A_j)
   +h^2\bigl(B_i+B_j+A_jA_i-\partial_iA_j\bigr)+O(h^3),\\
 t_{ji}
 &=I_r-h(A_i+A_j)
   +h^2\bigl(B_i+B_j+A_iA_j-\partial_jA_i\bigr)+O(h^3).
\end{align*}
Thus the second-order contribution of each individual edge cancels, and
\begin{align*}
 t_{ij}-t_{ji}
 =-h^2\bigl(
 \partial_iA_j-\partial_jA_i+[A_i,A_j]
 \bigr)+O(h^3).
\end{align*}
The expression in parentheses is the frame matrix of the curvature:
\begin{equation}
 \partial_iA_j-\partial_jA_i+[A_i,A_j]
 =e^*P[P_i,P_j]Pe.
\label{eq:frame-curvature}
\end{equation}
Since $t_{ij}$ and $t_{ji}$ are unitary and have the same first-order term,
\[
 t_{ji}^*t_{ij}
 =I_r-h^2e^*\cF_{ij}^{\rm K}e+O(h^3),
\]
which proves \eqref{eq:holonomy-expansion}.  The same path-difference
expansion, unitary invariance of the Hilbert--Schmidt norm, and division by
$h^4$ prove \eqref{eq:debit-limit}.  Taylor remainders and smoothness of the
inverse square root are uniform under the stated compactness and angle
conditions.
\end{proof}

\begin{corollary}[Route dependence of a corner response]
\label{cor:response-route}
Under the hypotheses of \cref{thm:finite-holonomy}, let $X=X P$ be a $C^2$
descended response.  At the opposite corner define the two base-aligned
responses
\[
 F_{ij,h}=X_{\theta+h(e_i+e_j)}T_{ij},
 \qquad
 F_{ji,h}=X_{\theta+h(e_i+e_j)}T_{ji}.
\]
After identifying the corner support with the base along either path,
\begin{equation}
 \boxed{
 \frac{F_{ij,h}-F_{ji,h}}{h^2}
 \longrightarrow -X_\theta\cF_{ij}^{\rm K}(\theta).}
\label{eq:response-route-limit}
\end{equation}
Thus a held-out command superposition or a mixed command derivative needs a
support-holonomy row even when both edgewise tangents exist.
\end{corollary}

\begin{proof}
The proof of \cref{thm:finite-holonomy} gives, in a common corner frame,
$T_{ij}-T_{ji}=-h^2T_{ji}\cF_{ij}^{\rm K}+O(h^3)$.  Multiply on the left by
$X_{\theta+h(e_i+e_j)}$, use $X_{\theta+h(e_i+e_j)}\to X_\theta$, and pull
back by the unitary path, which tends to the identity.  This yields
\eqref{eq:response-route-limit}.
\end{proof}

\section{An exact flat-action/nonflat-support separator}
\label{sec:separator}

The next calculation is finite, exact, and independent of any numerical
spectral approximation.

\begin{counterexample}[Flat likelihood square with curved relation support]
\label{ex:rank-one-separator}
Let $E=\C^2$ and, near $(0,0)\in\R^2$, set
\begin{equation}
 s(x,y)=\frac{1}{\sqrt{1+x^2+y^2}}
 \binom{1}{x+iy},
 \qquad
 P(x,y)=s(x,y)s(x,y)^*.
\label{eq:rank-one-family}
\end{equation}
Let the history space consist of one point, take the unnormalized density
$a_{x,y}=1$, and put $X(x,y)=s(x,y)^*$.  Then:
\begin{enumerate}[label=\textup{(S\arabic*)},leftmargin=3.2em]
\item every likelihood edge ratio is one, so every multiplicative action
      plaquette vanishes exactly;
\item $X=XP$, every support has rank one, and all sufficiently short polar
      chords exist;
\item at the origin,
      \begin{equation}
       \boxed{\cF_{xy}^{\rm K}(0,0)=2iP(0,0)\ne0;}
       \label{eq:rank-one-curvature}
      \end{equation}
\item on the square with vertices $(0,0),(h,0),(0,h),(h,h)$,
      \begin{equation}
       \boxed{
       \cH_{xy,h}
       =\frac{(1+h^2-ih^2)^2}{(1+h^2)^2+h^4}P(0,0),}
       \label{eq:exact-holonomy}
      \end{equation}
      and
      \begin{equation}
       \boxed{
       \Delta_{xy,h}^{\rm supp}
       =\frac{4h^4}{(1+h^2)^2+h^4}>0
       \quad(h>0).}
       \label{eq:exact-support-debit}
      \end{equation}
\end{enumerate}
The determinant and Pfaffian line rows may be chosen constant and nonzero, so
this obstruction is neither scalar action curvature nor line holonomy.
\end{counterexample}

\begin{proof}
Items (S1) and (S2) are immediate.  At the origin,
\begin{equation}
 P_0=\begin{pmatrix}1&0\\0&0\end{pmatrix},
 \quad
 P_x=\begin{pmatrix}0&1\\1&0\end{pmatrix},
 \quad
 P_y=\begin{pmatrix}0&-i\\i&0\end{pmatrix}.
\label{eq:Pauli-derivatives}
\end{equation}
Therefore
\[
 P_0[P_x,P_y]P_0=2iP_0,
\]
which proves (S3).

For a rank-one projection $P_u=uu^*$ with $\|u\|=1$, the polar chord sends
$u$ to
\begin{equation}
 U_{P_v\leftarrow P_u}u
 =v\frac{v^*u}{|v^*u|}.
\label{eq:rank-one-chord}
\end{equation}
The first edge on either route has positive overlap.  At the opposite-corner
edges the two phases are
\begin{equation}
 c_{xy}
 =\frac{1+h^2-ih^2}{\sqrt{(1+h^2)^2+h^4}},
 \qquad
 c_{yx}=\overline{c_{xy}}.
\label{eq:two-phases}
\end{equation}
Thus $T_{xy}e_1=s(h,h)c_{xy}$ and
$T_{yx}e_1=s(h,h)c_{yx}$.  Their relative holonomy is
$\overline{c_{yx}}c_{xy}=c_{xy}^2$, which is
\eqref{eq:exact-holonomy}.  Their squared distance is
\[
 |c_{xy}-c_{yx}|^2
 =\frac{4h^4}{(1+h^2)^2+h^4},
\]
proving (S4).  Notice that
$\cH_{xy,h}=P_0-2ih^2P_0+O(h^4)$ and
$h^{-4}\Delta_{xy,h}^{\rm supp}\to4
=\|2iP_0\|_{\HS}^2$, in agreement with
\cref{thm:finite-holonomy}.
\end{proof}

\begin{assessment}[What the separator decides]
\label{ass:separator}
The example proves the strict logical separation
\begin{equation}
 \boxed{
 \text{common positive action}
 \not\!\Longrightarrow
 \text{flat returned-word support transport}.}
\label{eq:logical-separation}
\end{equation}
It does not show that the selected RPESM support is curved.  That is a finite
question for its actual corner Grams.  It shows that the answer cannot be
inferred from the action plaquette, stable rank, endpoint descent, or trivial
determinant-line sewing.
\end{assessment}

\section{Area-scaled cutoff transport for a command square}
\label{sec:area-scaled}

Version V30 proves that first central secants amplify endpoint cutoff errors by
one inverse command radius.  A two-command plaquette or mixed derivative has a
stronger, inverse-area scaling.

\begin{theorem}[Exact cutoff bound for a mixed command secant]
\label{thm:area-bound}
At cutoff $n$, let $F_{n,ab}(h)$, $(a,b)\in\{0,1\}^2$, be four response
amplitudes transported to a common relation, physical-connection, and output
frame.  Define
\begin{equation}
 C_{n,h}:=\frac{F_{n,11}-F_{n,10}-F_{n,01}+F_{n,00}}{h^2}.
\label{eq:mixed-secant}
\end{equation}
Let $J_n$ and $V_n$ identify the input and output carriers at adjacent
cutoffs, and suppose
\begin{equation}
 \|F_{n+1,ab}J_n-V_nF_{n,ab}\|
 \le\delta_{n,ab}(h).
\label{eq:corner-errors}
\end{equation}
Then
\begin{equation}
 \boxed{
 \|C_{n+1,h}J_n-V_nC_{n,h}\|
 \le
 \frac{\delta_{n,11}(h)+\delta_{n,10}(h)
       +\delta_{n,01}(h)+\delta_{n,00}(h)}{h^2}.}
\label{eq:area-bound}
\end{equation}
Therefore endpoint convergence, and even the $o(h)$ control sufficient for a
first derivative, do not transport a mixed command row.  An area-normalized
row requires $o(h^2)$ control, or a summable version of it along a cofinal
diagonal.
\end{theorem}

\begin{proof}
Insert \eqref{eq:mixed-secant} at cutoffs $n+1$ and $n$, distribute the
transport maps, and group the four corner discrepancies with signs
$+,-,-,+$.  The triangle inequality gives \eqref{eq:area-bound}.
\end{proof}

\begin{counterexample}[First-derivative control need not control the square]
\label{ex:area-sharp}
On one-dimensional carriers, let the four coarse amplitudes be zero and the
fine amplitudes be
\begin{equation}
 F_{1,11}(h)=h^2,
 \qquad
 F_{1,10}(h)=F_{1,01}(h)=F_{1,00}(h)=0.
\label{eq:area-sharp-data}
\end{equation}
With identity transports every endpoint error is $O(h^2)=o(h)$, but
\begin{equation}
 \boxed{C_{1,h}-C_{0,h}=1}
\label{eq:area-order-one}
\end{equation}
for every $h>0$.  The exponent $h^{-2}$ in
\eqref{eq:area-bound} is therefore sharp.
\end{counterexample}

\begin{proof}
Substitute \eqref{eq:area-sharp-data} into \eqref{eq:mixed-secant}.
\end{proof}

\begin{corollary}[Cofinal transport of support curvature]
\label{cor:cofinal-curvature}
Let $h_n\downarrow0$.  Suppose at each cutoff the four support corners and
both polar paths are defined, and suppose the adjacent cutoff defects of every
route-resolved corner amplitude satisfy
\begin{equation}
 \sum_n h_n^{-2}
 \sum_{a,b\in\{0,1\}}
 \delta_{n,ab}(h_n)<\infty.
\label{eq:area-summability}
\end{equation}
Assume also a summable defect for changing $h_n$ to $h_{n+1}$ on the common
fine card.  Fix either route convention $ij$ or $ji$ at every cutoff.  The
corresponding mixed secants are then Cauchy on every fixed transported screen.
If their within-card Taylor remainders tend to zero, they converge to the
corresponding route-ordered covariant second derivative.  When both route
conventions satisfy the hypotheses, the difference of their limits is exactly
$-X\cF_{ij}^{\rm K}$, with the orientation fixed by
\eqref{eq:response-route-limit}.
\end{corollary}

\begin{proof}
Apply \cref{thm:area-bound} at $h_n$.  The first summability assumption makes
the adjacent-cutoff telescoping series convergent; the second joins the
diagonal radii.  The Taylor identification follows from
\cref{cor:response-route} and the ordinary two-variable mixed-difference
remainder.
\end{proof}

\section{Revised finite acquisition and continuum status}
\label{sec:revised-acquisition}

The V30 command square should be augmented, rather than replaced, by the
following finite rows.  At every elementary two-command square one records:
\begin{enumerate}[label=\textup{(K\arabic*)},leftmargin=3.2em]
\item the scalar likelihood plaquette $\Delta_{ij,h}^{\rm act}$ already
      required in V30;
\item all four canonical returned-word supports and the four edge polar
      chords;
\item both corner transports $T_{ij}$ and $T_{ji}$, the unitary support
      holonomy $\cH_{ij,h}$, and the positive debit
      $\Delta_{ij,h}^{\rm supp}$;
\item the response-visible curvature debit
      $\|X_{11}(T_{ij}-T_{ji})\|_{\HS}^2$;
\item the corresponding loop for the independently supplied supported
      physical reflection, flavour, and line connection;
\item the determinant/Pfaffian line holonomy as a separate row, so it is not
      identified with returned-word support curvature;
\item adjacent-cutoff errors divided by $h$ for first tangents and by $h^2$
      for mixed or plaquette-normalized rows.
\end{enumerate}

These rows lead to a terminating finite classification:
\begin{enumerate}[label=\textup{(B\arabic*)},leftmargin=3.2em]
\item a nonzero scalar plaquette is a common-action obstruction;
\item a zero scalar plaquette with nonzero support debit is a relation-support
      curvature event, not a source birth;
\item vanishing support debit with nontrivial physical or line holonomy is an
      internal-connection event;
\item after all applicable loop rows are retained, a positive quotient
      response reserve may be passed to the existing primitive and
      semigroup-cyclic shorts;
\item failure of the $h^{-2}$ cutoff row blocks a continuum mixed derivative
      even if every fixed cutoff square passes.
\end{enumerate}

\begin{theorem}[V1 support-plaquette theorem]
\label{thm:V1-terminal}
For a $C^3$ constant-rank returned-word support family, the first-order Kato
transport extends to a two-command square with curvature
$P[P_i,P_j]P$.  The finite polar-chord loop detects that curvature to second
order and its squared path debit, divided by $h^4$, converges to the
Hilbert--Schmidt curvature norm.  Scalar common-action flatness is logically
independent of this support flatness.  Transport of a mixed command row across
cutoffs requires errors controlled in units of command area $h^2$.
\end{theorem}

\begin{proof}
The curvature formula is \cref{thm:Kato-curvature}.  Its response action is
\cref{thm:response-curvature}.  Finite detection and the debit limit are
\cref{thm:finite-holonomy,cor:response-route}.  Logical independence follows
from the exact separator \cref{ex:rank-one-separator}.  The area-scaled cutoff
claim is \cref{thm:area-bound,ex:area-sharp,cor:cofinal-curvature}.
\end{proof}

\begin{assessment}[Programme boundary after V1]
\label{ass:V1-boundary}
This note closes one missing finite geometric implication: the relation
support has its own command plaquette, and its curvature must be transported
at area scale.  It does not populate the selected RPESM square because the
attachments do not supply its corner matrices or executable certificates.
It also does not prove regulator removal, uniform reflection positivity,
four-dimensional interacting tightness, renormalized Standard--Model
matching, Lorentzian reconstruction, the Einstein equation in the quantum
limit, or scattering.  Those remain separate analytic and physical proof
obligations.

The next iteration should first seek or reconstruct one explicit selected
RPESM boundary record with its finite corner action and returned-word matrices.
If those data remain unavailable, the next upstream step is to derive the
support-curvature matrix directly from a declared finite RPESM word synthesis
$W(\theta)$, with certified gap bounds ensuring constant rank, and then to
identify which physical command pair gives a nonzero or vanishing row.
\end{assessment}

\section*{V1 append-only record}
\addcontentsline{toc}{section}{V1 append-only record}
This V1 note imports the supplied manuscripts, adds the Kato support-curvature
and finite-holonomy layer, and ends with one active document-closing command.
The next notebook version must copy this file, remove only that final command,
append the V2 material, restore one final command, and use a filename ending in
\texttt{\_V2.tex}.



% =============================================================================
% BEGIN APPENDED VERSION V2 MATERIAL
% =============================================================================
\clearpage
\hypersetup{
 pdftitle={Renewal Standard--Model--Einstein Continuum Research Notes V2},
 pdfsubject={Exact score-support curvature and the two-command RPESM reduction}
}

\part*{Version V2 continuation: exact score-support curvature and a two-command RPESM subcard}
\addcontentsline{toc}{part}{Version V2 continuation: exact score-support curvature and a two-command RPESM subcard}

\section{V2 direction and the concrete finite target}
\label{sec:V2-direction}

Version V1 proved that a moving relation or source support has a Kato curvature
which is independent of scalar common-action flatness.  The next task is to
compute that curvature from quantities present in a physical command
experiment, without differentiating an $81\times81$ projector numerically.

The attached SM--Einstein V7 material declares on its $N=2$, $81$-record
finite representative the two local commands
\begin{equation}
 v_g=-\operatorname{Re}u_0,
 \qquad
 v_4=\phi_0^4,
 \qquad
 u_0\in\mathbb Z_3,
 \quad \phi_0\in\{-1,0,1\}.
\label{eq:V2-physical-commands}
\end{equation}
The same material proves that every record has strictly positive base weight.
It reports the gauge variance, but the cited complete exact weight table and
certificate output are not among the supplied files.  We therefore separate
what can be proved from the printed model from the one short numerical read
which remains unavailable.

For the declared exponential deformation, this continuation proves:
\begin{enumerate}[label=\textup{(D\arabic*)},leftmargin=3.2em]
\item the unnormalized two-command square and every larger rectangle have
      exactly zero scalar action plaquette;
\item the normalized score carrier has constant rank two on the entire finite
      command plane;
\item its Kato curvature is an explicit expression involving only the Fisher
      covariance and the centered third-moment tensor;
\item for two commands, flatness is decided by one scalar;
\item for the physical value patterns in \eqref{eq:V2-physical-commands}, all
      required moments are determined by four aggregated cell masses.
\end{enumerate}
This is a source-score subcard of the V30 acquisition.  It does not pretend to
be the full returned-word, overlap-sign, determinant/Pfaffian, or Duhamel card.

\section{Finite exponential commands and the score carrier}
\label{sec:V2-score-carrier}

Let $\Omega$ be a finite set, let $q(\omega)>0$, and let
$v_1,\ldots,v_p:\Omega\to\R$ be fixed real command writers.  For
$\theta\in\R^p$ define
\begin{equation}
 \Lambda_\theta(\omega)
 :=q(\omega)e^{-\theta\cdot v(\omega)},
 \qquad
 Z_\theta:=\sum_{\omega\in\Omega}\Lambda_\theta(\omega),
 \qquad
 \mu_\theta:=Z_\theta^{-1}\Lambda_\theta.
\label{eq:V2-exponential-family}
\end{equation}
Write $\mathbb E_\theta$ for expectation under $\mu_\theta$ and set
\begin{equation}
 m_a(\theta)=\mathbb E_\theta v_a,
 \qquad
 w_a(\theta)=v_a-m_a(\theta),
 \qquad
 s_\theta(\omega)=\sqrt{\mu_\theta(\omega)}.
\label{eq:V2-centered-writers}
\end{equation}
On $\mathcal H=\ell^2(\Omega)$ with counting inner product, define the score
synthesis
\begin{equation}
 B_\theta:\C^p\longrightarrow\mathcal H,
 \qquad
 B_\theta e_a=w_a(\theta)s_\theta.
\label{eq:V2-score-synthesis}
\end{equation}
Its Gram and centered third moments are
\begin{align}
 M_{ab}(\theta)
 &:=B_\theta^*B_\theta
 =\mathbb E_\theta(w_aw_b),
\label{eq:V2-Fisher}\\
 T_{abi}(\theta)
 &:=\mathbb E_\theta(w_aw_bw_i),
 \qquad
 (T_i)_{ab}:=T_{abi}.
\label{eq:V2-third-tensor}
\end{align}
The tensor $T_{abi}$ is symmetric in all three indices.

\begin{lemma}[Exact derivatives of the root, scores, and Fisher Gram]
\label{lem:V2-score-derivatives}
For every $a,b,i$,
\begin{align}
 \partial_i\log\mu_\theta&=-w_i,
 &\partial_i s_\theta&=-\tfrac12w_i s_\theta,
\label{eq:V2-root-derivative}\\
 \partial_iw_a&=M_{ai},
 &\partial_iM_{ab}&=-T_{abi},
\label{eq:V2-moment-derivatives}\\
 (\partial_iB)e_a
 &=\left(M_{ai}-\tfrac12w_aw_i\right)s_\theta.
\label{eq:V2-B-derivative}
\end{align}
\end{lemma}

\begin{proof}
Differentiate \eqref{eq:V2-exponential-family}.  Since
$\partial_i\log Z_\theta=-\mathbb E_\theta v_i$, one gets
$\partial_i\log\mu_\theta=-v_i+\mathbb E_\theta v_i=-w_i$ and then the root
formula.  For any fixed writer $f$,
\begin{equation}
 \partial_i\mathbb E_\theta f
 =-\mathbb E_\theta\bigl((f-\mathbb E_\theta f)w_i\bigr).
\label{eq:V2-expectation-derivative}
\end{equation}
Taking $f=v_a$ gives $\partial_i m_a=-M_{ai}$, hence
$\partial_iw_a=M_{ai}$.  Differentiate
$M_{ab}=\mathbb E_\theta(w_aw_b)$.  The two derivatives of the centered
writers are constants and their expectations are multiplied by
$\mathbb E_\theta w_a=\mathbb E_\theta w_b=0$.  Only the derivative of the
law remains, giving $\partial_iM_{ab}=-T_{abi}$.  Finally, apply the product
rule to $B e_a=w_as_\theta$.
\end{proof}

\begin{lemma}[Constant rank and canonical score support]
\label{lem:V2-score-support}
Assume that no nonzero $c\in\R^p$ makes $\sum_ac_av_a$ constant on $\Omega$.
Then $M_\theta\succ0$ for every $\theta\in\R^p$, the score synthesis has rank
$p$, and its range projection is
\begin{equation}
 \boxed{
 Q_\theta=B_\theta M_\theta^{-1}B_\theta^*.}
\label{eq:V2-score-projection}
\end{equation}
On every compact command set there is a positive Fisher floor
$M_\theta\succeq m_*I_p$.
\end{lemma}

\begin{proof}
For $c\in\R^p$,
\[
 c^*M_\theta c
 =\operatorname{Var}_{\mu_\theta}\!\left(\sum_ac_av_a\right).
\]
Because $\mu_\theta$ has full support, this variance vanishes exactly when the
displayed linear combination is constant on all of $\Omega$.  The hypothesis
therefore gives $M_\theta\succ0$.  The standard full-column-rank formula makes
$BM^{-1}B^*$ the orthogonal projection onto $\Ran B$.  Continuity of the
least eigenvalue and compactness give the uniform floor.
\end{proof}

The projection $Q_\theta$ lives on the history half-density carrier.
Equivalently, the descended read
$X_\theta=B_\theta^*:\mathcal H\to\C^p$ satisfies
$X_\theta=X_\theta Q_\theta$.  Thus the V1 support-transport theorem applies
to this physical score subcard.


\section{Curvature from third moments}
\label{sec:V2-third-curvature}

\begin{theorem}[Exact score-support curvature formula]
\label{thm:V2-score-curvature}
Under the hypotheses of \cref{lem:V2-score-support}, let
\begin{equation}
 E_\theta:=B_\theta M_\theta^{-1/2}:\C^p\longrightarrow\mathcal H.
\label{eq:V2-orthonormal-score-frame}
\end{equation}
Then $E_\theta^*E_\theta=I_p$ and $Q_\theta=E_\theta E_\theta^*$.  In this
orthonormal frame the Kato curvature of the score range is
\begin{equation}
 \boxed{
 E^*\cF_{ij}^{\rm K}E
 =-\frac14M^{-1/2}
 \left(T_iM^{-1}T_j-T_jM^{-1}T_i\right)
 M^{-1/2}.}
\label{eq:V2-score-curvature}
\end{equation}
In particular the fourth centered moments cancel from the curvature.  The
score support is flat in directions $(i,j)$ exactly when
\begin{equation}
 \boxed{T_iM^{-1}T_j=T_jM^{-1}T_i.}
\label{eq:V2-flatness}
\end{equation}
\end{theorem}

\begin{proof}
For any orthonormal frame $E$ of a projection $Q=EE^*$, put
\begin{equation}
 C_i=(I-Q)\partial_iE.
\label{eq:V2-second-fundamental}
\end{equation}
The projection identities from V1 give the finite-dimensional Gauss formula
\begin{equation}
 E^*\cF_{ij}^{\rm K}E=C_i^*C_j-C_j^*C_i.
\label{eq:V2-Gauss}
\end{equation}
Indeed $(\partial_iQ)E=(I-Q)\partial_iE=C_i$, so compression of
$Q[\partial_iQ,\partial_jQ]Q$ by $E$ is exactly the right-hand side.

Differentiate $E=BM^{-1/2}$.  The term containing
$\partial_iM^{-1/2}$ has range in $\Ran B$ and is killed by $I-Q$.  Hence
\begin{equation}
 C_i=(I-Q)(\partial_iB)M^{-1/2}.
\label{eq:V2-Ci}
\end{equation}
The derivative formula \eqref{eq:V2-B-derivative} gives, entrywise,
\begin{align}
 \bigl((\partial_iB)^*(\partial_jB)\bigr)_{ab}
 &=\frac14\mathbb E_\theta(w_aw_bw_iw_j),
\label{eq:V2-fourth-cancellation-a}\\
 \bigl(B^*\partial_iB\bigr)_{ca}
 &=-\frac12T_{cai}=-\frac12(T_i)_{ca}.
\label{eq:V2-fourth-cancellation-b}
\end{align}
For the first identity, the constant--quadratic cross terms cancel the product
$M_{ai}M_{bj}$; for the second, the constant term has zero inner product with
every centered score.  Since $Q=BM^{-1}B^*$,
\begin{equation}
 (\partial_iB)^*Q(\partial_jB)
 =\frac14T_iM^{-1}T_j.
\label{eq:V2-projected-derivatives}
\end{equation}
The fourth-moment term in \eqref{eq:V2-fourth-cancellation-a} is symmetric
under $i\leftrightarrow j$ and therefore cancels in
$C_i^*C_j-C_j^*C_i$.  Substitution into \eqref{eq:V2-Gauss} proves
\eqref{eq:V2-score-curvature}.  Invertibility of $M^{-1/2}$ gives the exact
flatness criterion.
\end{proof}

\begin{corollary}[Coordinate covariance]
\label{cor:V2-coordinate-covariance}
Replace the command bank by
$v'_\alpha=\sum_aA_{a\alpha}v_a+c_\alpha$ with
$A\in\operatorname{GL}_p(\R)$.  The score ranges and their projection
$Q_\theta$ are unchanged after the corresponding inverse command-coordinate
change.  Formula \eqref{eq:V2-score-curvature} changes by the tensorial command-coordinate
rule and by the unitary transformation between the two orthonormal score
frames.  Hence flatness, rank, and polar holonomy are intrinsic.  For fixed
tangent vectors $\xi,\eta$, the norm $\|\cF^{\rm K}(\xi,\eta)\|_{\HS}$
is coordinate independent; a displayed coordinate component such as
$\|\cF_{12}^{\rm K}\|_{\HS}$ rescales when the command axes are rescaled.
\end{corollary}

\begin{proof}
Constants disappear on centering and the two score syntheses satisfy
$B'=BA$.  Therefore
\[
 B'(B'^*B')^{-1}B'^*
 =BA(A^*MA)^{-1}A^*B^*
 =BM^{-1}B^*=Q.
\]
The curvature is the operator-valued two-form
$Q[\dd Q,\dd Q]Q$, so it obeys the usual covariant rule under a command
coordinate change.  Two orthonormal frames of $\Ran Q$ differ by a unitary
matrix, and their compressed curvature matrices are unitarily conjugate.
Evaluation on fixed tangent vectors therefore has invariant Hilbert--Schmidt
norm, while coordinate components transform tensorially.
\end{proof}

\begin{corollary}[One-scalar test for two whitened commands]
\label{cor:V2-two-command}
At a fixed command point, whiten the command bank so that $M=I_2$ and write
the symmetric third tensor as
\begin{equation}
 T_{111}=A,\qquad T_{112}=B,\qquad
 T_{122}=C,\qquad T_{222}=D.
\label{eq:V2-whitened-third}
\end{equation}
Put
\begin{equation}
 \kappa:=AC+BD-B^2-C^2.
\label{eq:V2-kappa-white}
\end{equation}
Then
\begin{equation}
 \boxed{
 E^*\cF_{12}^{\rm K}E
 =-\frac14
 \begin{pmatrix}0&\kappa\\-\kappa&0\end{pmatrix},
 \qquad
 \|\cF_{12}^{\rm K}\|_{\HS}^2=\frac{\kappa^2}{8}.}
\label{eq:V2-whitened-curvature}
\end{equation}
Thus the two-command score support is flat at that point if and only if
$\kappa=0$.
\end{corollary}

\begin{proof}
In whitened coordinates,
\[
 T_1=\begin{pmatrix}A&B\\B&C\end{pmatrix},
 \qquad
 T_2=\begin{pmatrix}B&C\\C&D\end{pmatrix}.
\]
Direct multiplication gives
\[
 [T_1,T_2]
 =\begin{pmatrix}
 0&AC+BD-B^2-C^2\\
 -(AC+BD-B^2-C^2)&0
 \end{pmatrix}.
\]
Apply \cref{thm:V2-score-curvature}.  The displayed real skew matrix has
squared Hilbert--Schmidt norm $2\kappa^2$ before multiplication by $1/4$.
\end{proof}

For a calculation in the original, unwhitened command coordinates, write
\begin{equation}
 M=\begin{pmatrix}a&c\\c&b\end{pmatrix},
 \qquad
 T_1=\begin{pmatrix}A&B\\B&C\end{pmatrix},
 \qquad
 T_2=\begin{pmatrix}B&C\\C&D\end{pmatrix},
 \qquad
 \Delta=ab-c^2>0.
\label{eq:V2-unwhitened-data}
\end{equation}

\begin{corollary}[Unwhitened scalar certificate]
\label{cor:V2-unwhitened}
Define
\begin{equation}
 k:=
 \frac{
 b(AC-B^2)+a(BD-C^2)+c(BC-AD)}
 {\Delta}.
\label{eq:V2-unwhitened-k}
\end{equation}
Then
\begin{equation}
 T_1M^{-1}T_2-T_2M^{-1}T_1
 =\begin{pmatrix}0&k\\-k&0\end{pmatrix},
\label{eq:V2-unwhitened-commutator}
\end{equation}
and
\begin{equation}
 \boxed{
 \cF_{12}^{\rm K}=0\Longleftrightarrow k=0,
 \qquad
 \|\cF_{12}^{\rm K}\|_{\HS}^2
 =\frac{k^2}{8\Delta}.}
\label{eq:V2-unwhitened-norm}
\end{equation}
Thus three second moments and four third moments give a complete scalar
support-curvature certificate for two commands.
\end{corollary}

\begin{proof}
Use
$M^{-1}=\Delta^{-1}\begin{psmallmatrix}b&-c\\-c&a\end{psmallmatrix}$
and multiply.  The $(1,2)$ entry is precisely
\eqref{eq:V2-unwhitened-k}, and skew-adjointness supplies the other
off-diagonal entry.  Formula \eqref{eq:V2-score-curvature} proves flatness.
For the norm, let $J=\begin{psmallmatrix}0&1\\-1&0\end{psmallmatrix}$.
For every real symmetric $2\times2$ matrix $S$,
$SJS=(\det S)J$.  Taking $S=M^{-1/2}$ gives
\[
 M^{-1/2}(kJ)M^{-1/2}=\frac{k}{\sqrt\Delta}J.
\]
Its squared Hilbert--Schmidt norm is $2k^2/\Delta$; the factor $1/16$ from
\eqref{eq:V2-score-curvature} yields \eqref{eq:V2-unwhitened-norm}.
\end{proof}


\section{Exact finite checks}
\label{sec:V2-exact-checks}

\begin{counterexample}[A three-record exponential family has curved score support]
\label{ex:V2-three-record}
Let $\Omega=\{0,1,2\}$, let $\mu_0$ be uniform, and take
\[
 v_1=\mathbf1_{\{1\}},
 \qquad
 v_2=\mathbf1_{\{2\}}.
\]
At $\theta=0$,
\begin{equation}
 M=\frac19\begin{pmatrix}2&-1\\-1&2\end{pmatrix},
 \qquad
 M^{-1}=\begin{pmatrix}6&3\\3&6\end{pmatrix},
\label{eq:V2-three-M}
\end{equation}
and
\begin{equation}
 T_1=\frac1{27}\begin{pmatrix}2&-1\\-1&-1\end{pmatrix},
 \qquad
 T_2=\frac1{27}\begin{pmatrix}-1&-1\\-1&2\end{pmatrix}.
\label{eq:V2-three-T}
\end{equation}
Consequently
\begin{equation}
 T_1M^{-1}T_2-T_2M^{-1}T_1
 =\frac1{27}\begin{pmatrix}0&-1\\1&0\end{pmatrix}.
\label{eq:V2-three-commutator}
\end{equation}
The score range is the full centered plane
$s_0^\perp\subset\R^3$, and its curvature as an operator on $\R^3$ is
\begin{equation}
 \boxed{
 \cF_{12}^{\rm K}
 =\frac1{36}
 \begin{pmatrix}
 0&1&-1\\
 -1&0&1\\
 1&-1&0
 \end{pmatrix},
 \qquad
 \|\cF_{12}^{\rm K}\|_{\HS}^2=\frac1{216}.}
\label{eq:V2-three-curvature}
\end{equation}
\end{counterexample}

\begin{proof}
The centered writer vectors on the three equally weighted atoms are
\[
 w_1=(-1/3,2/3,-1/3),
 \qquad
 w_2=(-1/3,-1/3,2/3).
\]
Direct averaging gives \eqref{eq:V2-three-M} and
\eqref{eq:V2-three-T}; multiplication gives
\eqref{eq:V2-three-commutator}.  Since the two centered writers are
independent and $\dim s_0^\perp=2$, the score projection is
$Q=I-s_0s_0^*$.  From \eqref{eq:V2-root-derivative},
\[
 \partial_iQ
 =\frac12(M_{w_i}s_0s_0^*+s_0s_0^*M_{w_i}).
\]
Multiplying $Q[\partial_1Q,\partial_2Q]Q$ gives the matrix in
\eqref{eq:V2-three-curvature}.  It has six off-diagonal entries of magnitude
$1/36$, so its squared Hilbert--Schmidt norm is
$6/36^2=1/216$.  Independently,
\cref{cor:V2-unwhitened} has
$\Delta=\det M=1/27$ and $k=-1/27$, giving
$k^2/(8\Delta)=1/216$.
\end{proof}

The preceding calculation is a score-support example, rather than a reuse of
the V1 arbitrary rank-one projection separator.  It proves that the curvature
can occur inside the normalized half-density geometry of an ordinary positive
exponential family.

\section{Four-cell compression of the physical command pair}
\label{sec:V2-cell-compression}

\begin{lemma}[Sufficient-cell compression]
\label{lem:V2-cell-compression}
Let $z:\Omega\to\mathcal Z$ be finite and suppose every command writer is a
function of $z$.  At a base point define
\begin{equation}
 p_\zeta=\sum_{\omega:z(\omega)=\zeta}\mu_0(\omega),
 \qquad
 e_\zeta(\omega)
 =\frac{\mathbf1_{\{z(\omega)=\zeta\}}s_0(\omega)}
 {\sqrt{p_\zeta}}.
\label{eq:V2-cell-frame}
\end{equation}
Then the vectors $(e_\zeta)_{\zeta\in\mathcal Z}$ are orthonormal and remain
fixed under every exponential command tilt.  Moreover,
\begin{equation}
 s_\theta=\sum_{\zeta\in\mathcal Z}
 \sqrt{p_\zeta(\theta)}\,e_\zeta,
 \qquad
 B_\theta e_a=\sum_{\zeta\in\mathcal Z}
 w_a(\zeta;\theta)\sqrt{p_\zeta(\theta)}\,e_\zeta,
\label{eq:V2-cell-score}
\end{equation}
where
\begin{equation}
 p_\zeta(\theta)
 =\frac{p_\zeta e^{-\theta\cdot v(\zeta)}}
 {\sum_{\xi\in\mathcal Z}p_\xi e^{-\theta\cdot v(\xi)}}.
\label{eq:V2-cell-law}
\end{equation}
Thus the score projection, its polar chords, its Kato curvature, and every
moment in \eqref{eq:V2-score-curvature} are unitarily equivalent to the
corresponding objects on the compressed cell law $(p_\zeta(\theta))$.
\end{lemma}

\begin{proof}
Disjoint cell supports make the vectors in \eqref{eq:V2-cell-frame}
orthonormal.  Within a cell, the tilt multiplier is constant, so
\[
 \mu_\theta(\omega\mid z=\zeta)
 =\mu_0(\omega\mid z=\zeta).
\]
The conditional square-root vector $e_\zeta$ is therefore independent of
$\theta$.  Summing the tilted weights gives \eqref{eq:V2-cell-law}, and
expanding the root and centered-score columns in the fixed cell frame gives
\eqref{eq:V2-cell-score}.  Every listed operator has range in the fixed span
of the $e_\zeta$, and its matrix there is exactly the compressed-cell matrix.
\end{proof}

For the physical commands \eqref{eq:V2-physical-commands}, put
\begin{equation}
 X=\mathbf1_{\{u_0\ne1\}},
 \qquad
 Y=\phi_0^4=\mathbf1_{\{\phi_0\ne0\}}.
\label{eq:V2-binary-commands}
\end{equation}
Then
\begin{equation}
 v_g=-1+\frac32X,
 \qquad
 v_4=Y.
\label{eq:V2-affine-binary}
\end{equation}
Constants disappear on centering and the nonzero factor $3/2$ is an
invertible command-coordinate change.  If $\eta$ is the command coordinate
conjugate to $X$ and $\alpha$ is conjugate to $v_g$, then
$\eta=3\alpha/2$ after normalization.  Consequently
\begin{equation}
 \cF_{g4}^{\rm K}=\frac32\cF_{XY}^{\rm K},
 \qquad
 \|\cF_{g4}^{\rm K}\|_{\HS}^2
 =\frac94\|\cF_{XY}^{\rm K}\|_{\HS}^2.
\label{eq:V2-physical-curvature-scale}
\end{equation}
The score support itself is determined by the four masses
\begin{equation}
 p_{xy}:=\mu_0(X=x,Y=y),
 \qquad (x,y)\in\{0,1\}^2.
\label{eq:V2-four-masses}
\end{equation}

\begin{lemma}[Binary moment reduction]
\label{lem:V2-binary-moments}
Let
\begin{equation}
 r=\mathbb E X=p_{10}+p_{11},
 \qquad
 s=\mathbb E Y=p_{01}+p_{11},
 \qquad
 c=p_{11}-rs,
\label{eq:V2-binary-means}
\end{equation}
and put $a=r(1-r)$ and $b=s(1-s)$.  For the command pair $(X,Y)$,
\begin{equation}
 M=\begin{pmatrix}a&c\\c&b\end{pmatrix}
\label{eq:V2-binary-M}
\end{equation}
and the four independent third moments are
\begin{equation}
 \boxed{
 A=a(1-2r),\qquad
 B=c(1-2r),\qquad
 C=c(1-2s),\qquad
 D=b(1-2s).}
\label{eq:V2-binary-T}
\end{equation}
If all four $p_{xy}$ are positive, then $M\succ0$.
\end{lemma}

\begin{proof}
The covariance formula is immediate.  A Bernoulli variable satisfies
\begin{equation}
 (X-r)^2=(1-2r)(X-r)+r(1-r),
\label{eq:V2-Bernoulli-square}
\end{equation}
and the analogous identity holds for $Y$.  Taking expectations after
multiplication by $X-r$ or $Y-s$ gives
\eqref{eq:V2-binary-T}.  If $M$ were singular, some nonzero affine
combination of $X$ and $Y$ would be constant almost surely.  Positive mass on
all four corners rules this out.
\end{proof}

\begin{theorem}[Exact physical two-command source-score subcard]
\label{thm:V2-RPESM-subcard}
On the $N=2$ finite representative, freeze the base unnormalized weight
$q(\omega)$, including its bosonic, determinant-modulus, and Pfaffian-modulus
factors, and deform only by the two declared local action writers:
\begin{equation}
 \Lambda_{\alpha,\beta}(\omega)
 =q(\omega)\exp[-\alpha v_g(\omega)-\beta v_4(\omega)].
\label{eq:V2-RPESM-cylinder}
\end{equation}
For every step $h$ the same-history positive edge ratios are
\begin{equation}
 r_g=e^{-hv_g},
 \qquad
 r_4=e^{-hv_4}.
\label{eq:V2-RPESM-edges}
\end{equation}
They obey every elementary and held-out rectangular common-action identity
exactly.  The score synthesis for $(v_g,v_4)$ has rank two for all
$(\alpha,\beta)$, has a positive Fisher floor on every compact command
rectangle, and its read $B_{\alpha,\beta}^*$ descends exactly through the
score-range projection.

At the base point, the support curvature and its Hilbert--Schmidt norm are
determined by \eqref{eq:V2-unwhitened-k}--\eqref{eq:V2-unwhitened-norm}, with
the moments obtained from the four masses \eqref{eq:V2-four-masses} using
\eqref{eq:V2-binary-M}--\eqref{eq:V2-binary-T} and the affine coordinate
change \eqref{eq:V2-affine-binary}.  No $81\times81$ projector derivative is
needed.
\end{theorem}

\begin{proof}
The edge ratios follow by division of \eqref{eq:V2-RPESM-cylinder}.  Around
any rectangle, the two products are
\[
 e^{-hv_g}e^{-kv_4}=e^{-kv_4}e^{-hv_g},
\]
so the scalar action plaquette is identically zero, history by history.
The attached V7 finite proof gives strictly positive weight to all $81$
records.  Each of the four pairs $(X,Y)\in\{0,1\}^2$ occurs among those
records, so its aggregate mass is positive.  By
\cref{lem:V2-binary-moments}, the Fisher matrix is positive definite.
\Cref{lem:V2-score-support} gives constant rank, compact Fisher floors, and
the exact descended read.  The cell and moment reductions are
\cref{lem:V2-cell-compression,lem:V2-binary-moments}, and the curvature
certificate is \cref{thm:V2-score-curvature,cor:V2-unwhitened}.
\end{proof}

\begin{remark}[What is frozen in this subcard]
The determinant and Pfaffian sections are part of $q$ and are held fixed as
functions on the record set while the two action coefficients are tilted.
Their values still affect every cell mass.  This is the two local-action
experiment declared in the V7 exponential family; it is not a determinant
command, a Wilson-sign command, or a claim that the physical line connection
has been measured.
\end{remark}


\begin{lemma}[A direct Fisher floor from four positive cell masses]
\label{lem:V2-cell-floor}
If $p_{xy}\ge\eta>0$ for all four binary cells, then the covariance matrix of
$(X,Y)$ satisfies
\begin{equation}
 \boxed{M_{X,Y}\succeq2\eta^2I_2.}
\label{eq:V2-cell-floor}
\end{equation}
The same lower bound holds for $(v_g,v_4)$ in
\eqref{eq:V2-affine-binary}.
\end{lemma}

\begin{proof}
For $Z=\alpha X+\beta Y$, the pairwise variance identity is
\[
 \operatorname{Var}Z
 =\sum_{\zeta<\xi}p_\zeta p_\xi
  \bigl(Z(\zeta)-Z(\xi)\bigr)^2.
\]
Keep only the two pairs which differ in $X$ at fixed $Y$ and the two pairs
which differ in $Y$ at fixed $X$.  This gives
\[
 \operatorname{Var}Z
 \ge(p_{00}p_{10}+p_{01}p_{11})\alpha^2
   +(p_{00}p_{01}+p_{10}p_{11})\beta^2
 \ge2\eta^2(\alpha^2+\beta^2).
\]
The change from $X$ to $v_g=-1+3X/2$ multiplies the first centered column by
$3/2$, whose least singular value relative to the unchanged second column is
one.
\end{proof}

\begin{counterexample}[Exact nonflat four-cell control]
\label{ex:V2-four-cell}
On the four cells $(00),(10),(01),(11)$ take
\begin{equation}
 (p_{00},p_{10},p_{01},p_{11})
 =\frac15(1,1,1,2).
\label{eq:V2-four-cell-law}
\end{equation}
For the binary commands $(X,Y)$,
\begin{equation}
 r=s=\frac35,\qquad
 M=\frac1{25}\begin{pmatrix}6&1\\1&6\end{pmatrix},
 \qquad
 \Delta=\frac7{125},
\label{eq:V2-four-cell-M}
\end{equation}
and
\begin{equation}
 (A,B,C,D)
 =-\frac1{125}(6,1,1,6).
\label{eq:V2-four-cell-T}
\end{equation}
The scalar certificate and curvature norm are
\begin{equation}
 \boxed{
 k=\frac1{875},
 \qquad
 \|\cF_{12}^{\rm K}\|_{\HS}^2=\frac1{343000}>0.}
\label{eq:V2-four-cell-curvature}
\end{equation}
By contrast, every product cell law $p_{xy}=p_xp_y$ has $c=B=C=0$ and
therefore $k=0$.  Strict positivity and the command value pattern alone do
not decide the physical curvature.
\end{counterexample}

\begin{proof}
Equations \eqref{eq:V2-four-cell-M} and
\eqref{eq:V2-four-cell-T} follow from
\cref{lem:V2-binary-moments}.  Substitution in
\eqref{eq:V2-unwhitened-k} gives $k=1/875$, and
\[
 \frac{k^2}{8\Delta}
 =\frac{1/875^2}{8(7/125)}
 =\frac1{343000}.
\]
For a product law, $\operatorname{Cov}(X,Y)=0$.  The binary third-moment
formulas then give $B=C=0$, and
\eqref{eq:V2-unwhitened-k} gives $k=0$.
\end{proof}

\begin{assessment}[Status of the actual $81$-record curvature read]
\label{ass:V2-actual-read}
The exact control \cref{ex:V2-four-cell} is not a measurement of the printed
RPESM weights.  For those weights, the missing numerical input has now been
reduced to
\[
 \boxed{(p_{00},p_{10},p_{01},p_{11}),\qquad \sum_{x,y}p_{xy}=1.}
\]
These masses are obtained by summing the cited exact $81$-record weight table
over the four cells \eqref{eq:V2-binary-commands}.  Once supplied, they give
the Fisher floor, the complete third tensor, the curvature scalar $k$, and
the infinitesimal polar-holonomy debit by closed formulas.

The V7 gauge variance supplies only
\[
 M_{gg}=\frac94r(1-r).
\]
It does not determine $s$ or $c=p_{11}-rs$, and therefore cannot decide
$k$.  Assigning either the flat product value or the nonflat value from the
variance alone would be an unsupported inference.
\end{assessment}

\section{The revised finite and cofinal acquisition}
\label{sec:V2-acquisition}

For the gauge/quartic source-score subcard, the next exact finite calculation
is now:
\begin{enumerate}[label=\textup{(A\arabic*)},leftmargin=3.2em]
\item sum the exact unnormalized $81$-record weights into the four cells
      $(X,Y)=(0,0),(1,0),(0,1),(1,1)$;
\item normalize once, report all four outward intervals, and certify a positive
      lower cell mass;
\item compute $r,s,c$, then $M$ and $(A,B,C,D)$ from
      \eqref{eq:V2-binary-means}--\eqref{eq:V2-binary-T};
\item evaluate the single rational expression
      \eqref{eq:V2-unwhitened-k} in the binary coordinates and apply the
      physical scaling \eqref{eq:V2-physical-curvature-scale};
\item compare the predicted $h^4\|\cF_{g4}^{\rm K}\|_{\HS}^2$ with the
      two polar-chord route debit on at least one shrinking command square;
\item retain the scalar action plaquette, which is exactly zero by
      \cref{thm:V2-RPESM-subcard}, as a separate row;
\item at an adjacent regulator cutoff, transport both routes with the
      area-scaled error budget proved in V1.
\end{enumerate}
Items (A1)--(A4) replace an $81\times81$ projector differentiation by four
positive sums and one $2\times2$ calculation.  Item (A5) is a held-out
geometric check, rather than a fitted definition of curvature.

For a cofinal sequence, a uniform lower cell mass gives a uniform Fisher floor
by \cref{lem:V2-cell-floor}.  If a cell mass tends to zero, the score carrier
may approach a rank transition and the correct output is a support-boundary
event.  If the cell masses remain positive but the area-scaled route errors
are not summable, each finite curvature read remains valid while its continuum
transport is unresolved.

\begin{theorem}[Version V2 two-command score-support theorem]
\label{thm:V2-terminal}
Preserving all V1 results, Version V2 proves:
\begin{enumerate}[label=\textup{(V2.\arabic*)},leftmargin=5.2em]
\item the Kato curvature of a finite exponential-family score range is given
      exactly by the Fisher matrix and centered third moments through
      \eqref{eq:V2-score-curvature};
\item for two commands, one scalar $k$ is necessary and sufficient for
      support flatness, and $k^2/(8\det M)$ is the squared curvature norm;
\item a positive three-record exponential family and a positive four-cell
      family have explicitly nonzero score-support curvature;
\item the declared RPESM gauge/quartic exponential deformation has an exactly
      flat scalar action square, a globally rank-two score carrier, and a
      compact-command Fisher floor;
\item its score-support curvature is determined entirely by the four joint
      cell masses of $u_0=1$ versus $u_0\ne1$ and $\phi_0=0$ versus
      $\phi_0\ne0$;
\item the supplied gauge variance alone cannot determine that curvature.
\end{enumerate}
\end{theorem}

\begin{proof}
Item (V2.1) is \cref{thm:V2-score-curvature}.  Item (V2.2) is
\cref{cor:V2-two-command,cor:V2-unwhitened}.  Item (V2.3) is
\cref{ex:V2-three-record,ex:V2-four-cell}.  Items (V2.4)--(V2.5) are
\cref{lem:V2-cell-compression,lem:V2-binary-moments,lem:V2-cell-floor,thm:V2-RPESM-subcard}.  The last item is
\cref{ass:V2-actual-read}.
\end{proof}

\begin{assessment}[Programme boundary after V2]
\label{ass:V2-boundary}
This version populates the common-action, full-support, constant-rank, and
finite formula rows for one physically declared two-command source-score
subcard.  It identifies exactly three independent scalar cell probabilities
as the missing data for its numerical curvature verdict.  The selected
$81$-record weight table or certificate output is still absent, so no sign or
nonzero lower bound for the actual RPESM value of $k$ is claimed.

The full V30 acquisition also requires the returned-word and overlap-sign
supports, the independently physical reflection/flavour/line transport,
determinant/Pfaffian sewing under commands which move those coefficients,
direct Duhamel responses, held-out words, and adjacent-cutoff data.  Beyond
that finite square, regulator removal still requires uniform reflection
positivity, tightness, locality, renormalized matching, and the separate
Einstein--matter continuum hypotheses.  None is discharged by a finite
score-support calculation.
\end{assessment}

\section*{V2 exact checks and append-only record}
\addcontentsline{toc}{section}{V2 exact checks and append-only record}
The exact rational checks used above are:
\[
 \det\!\left(\frac1{25}
 \begin{pmatrix}6&1\\1&6\end{pmatrix}\right)=\frac7{125},
 \qquad
 k=\frac1{875},
 \qquad
 \frac{k^2}{8\det M}=\frac1{343000}.
\]
A separate exact-arithmetic replay also verifies
\eqref{eq:V2-three-commutator} and
\eqref{eq:V2-three-curvature}.  These controls test the proofs; they are not
substitutes for the missing RPESM cell sums.

The V1 source before its final active document-closing command is retained
verbatim.  Its complete-file SHA--256 digest is
\begin{center}\small\ttfamily
39bb974e2db7352111b16ffb897ef5d652ac79513059301316224fee9ea651b7.
\end{center}
A Version V3 must retain the complete V2 source, remove only the final active
document-closing command, append new material, restore one final command, and
use a filename ending in \texttt{\_V3.tex}.

% =============================================================================
% END APPENDED VERSION V2 MATERIAL
% =============================================================================

% =============================================================================
% BEGIN APPENDED VERSION V3 MATERIAL
% =============================================================================

\clearpage
\section{Version V3: global refresh, spatial jump locality, and Green transport}
\label{sec:V3-refresh-locality}

Version V2 reduced one missing source-score curvature read to four cell masses.
Those masses are not present in the supplied sources, so assigning them would
manufacture a physical conclusion.  This version instead moves upstream to a
different continuum interface that is already populated symbolically in the
Grand-Tensor construction.  At each cutoff the declared parent generator has
the form
\begin{equation}
 \mathcal L_X^{\rm par}
 =\rho_X(\mathsf R_X-I)+\mathcal L_X^{\rm loc},
 \qquad
 \mathsf R_Xh=\mu_X[h]\mathbf 1.
 \label{eq:V3-parent-generator}
\end{equation}
The supplied monograph and its Lean rendering already prove the exact
fixed-cutoff semigroup identities which remove this refresh, the corresponding
split of the finite-time carr\'e-du-champ, and the scalar rescaling of selected
lag moments.  They also prove that the gap created by $\rho_X$ is not evidence
for a Yang--Mills gap.  None of those results is repeated here.

The new question is spatial and cofinal.  A global reset is one jump which can
change two arbitrarily separated records at once.  We first give an exact
operator test for that event, then determine what must replace the refresh
bound when one transports zero-mass Green responses to the local generator.
The outcome is a three-level de-refreshing criterion: a uniform local gap
transports arbitrary sources, one local-generator factor transports Green
quadratic forms, and a fractional extra generator factor transports response
vectors even without a uniform gap.

\begin{assessment}[V3 nonduplication boundary]
\label{ass:V3-nonduplication}
The fixed-cutoff identities
\[
 P_X^{\rm par}(t)
 =R_X+e^{-\rho_Xt}Q_XP_X^{\rm loc}(t)Q_X,
 \qquad
 P_X^{\rm loc}(t)
 =R_X+e^{\rho_Xt}Q_XP_X^{\rm par}(t)Q_X
\]
are imported.  V3 studies instead the double multiplication commutator of the
\emph{generator}, the resulting separated-region jump test, and uniform
comparison of $(H_X+\rho_XI)^{-1}$ with $H_X^\dagger$ on physical source
ranges.  These statements are absent from the attached de-refresh theorem.
\end{assessment}

% -----------------------------------------------------------------------------
\subsection{An exact two-region diagnostic for a finite jump generator}
\label{subsec:V3-double-commutator}
% -----------------------------------------------------------------------------

Let $(\Omega,\mathcal F,\mu)$ be a probability space.  We use bounded
measurable functions and bounded operators on $L^2(\mu)$.  This formulation is
deliberate: a spacetime cutoff has finitely many sites, but its Higgs
configuration space is still continuous.  We write $M_f$ for multiplication
by $f$ and $\mu[f]=\int_\Omega f\,\dd\mu$.

\begin{definition}[Two-region multiplication commutator]
\label{def:V3-double-commutator}
For a conservative operator $L$ on functions on $\Omega$, define
\begin{equation}
 \boxed{
 \mathfrak C_L(f,g):=[[L,M_f],M_g].}
 \label{eq:V3-double-commutator}
\end{equation}
This operator measures an instantaneous event which changes both $f$ and $g$.
It is bilinear in $(f,g)$, symmetric in $f,g$, and linear in $L$.
\end{definition}

\begin{theorem}[Exact jump formula and short-time meaning]
\label{thm:V3-jump-formula}
Suppose $q(x,\dd y)$ is a finite positive jump kernel and the bounded
generator $L$ on $L^2(\mu)$ is
\begin{equation}
 (Lh)(x)=\int_\Omega\bigl(h(y)-h(x)\bigr)q(x,\dd y).
 \label{eq:V3-jump-generator}
\end{equation}
Then
\begin{equation}
 \boxed{
 (\mathfrak C_L(f,g)h)(x)
 =\int_\Omega
   \bigl(f(y)-f(x)\bigr)\bigl(g(y)-g(x)\bigr)h(y)\,q(x,\dd y).}
 \label{eq:V3-jump-commutator}
\end{equation}
If $(X_t)$ is the continuous-time chain with generator $L$, then
\begin{equation}
 \boxed{
 \left.\frac{\dd}{\dd t}\right|_{t=0+}
 \mathbb E_x\!\left[
 (f(X_t)-f(x))(g(X_t)-g(x))
 \right]
 =(\mathfrak C_L(f,g)\mathbf1)(x).}
 \label{eq:V3-increment-derivative}
\end{equation}
\end{theorem}

\begin{proof}
Since multiplication operators commute,
\[
 \mathfrak C_L(f,g)h
 =L(fgh)-fL(gh)-gL(fh)+fgLh.
\]
Insert \eqref{eq:V3-jump-generator}.  All terms multiplying $h(x)$ cancel.
The coefficient of $h(y)$ is
\[
 f(y)g(y)-f(x)g(y)-g(x)f(y)+f(x)g(x),
\]
which factors into the two increments in
\eqref{eq:V3-jump-commutator}.  For the last assertion, put
$F_x(y)=(f(y)-f(x))(g(y)-g(x))$.  Because $F_x(x)=0$, the defining derivative
of the Markov semigroup gives
\[
 \left.\frac{\dd}{\dd t}\right|_{0+}\mathbb E_xF_x(X_t)
 =(LF_x)(x)
 =\int_\Omega F_x(y)\,q(x,\dd y).
\]
The jump formula with $h=\mathbf1$ identifies the last expression.
\end{proof}

We next state the locality consequence in a form which does not assume that
the local updates are single-site heat baths.

\begin{definition}[Finite update range]
\label{def:V3-update-range}
Let $\Lambda$ be a finite metric set and
$\Omega=\prod_{z\in\Lambda}\Omega_z$.  A function is $S$-local if it depends
only on the coordinates in $S\subseteq\Lambda$.  An update generator $L_A$ is
supported in $A\subseteq\Lambda$ when its jump rate $q_A(x,y)$ vanishes unless
$x$ and $y$ agree on $A^c$.  A generator is carried by an update family
$\mathcal A$ when
\[
 L^{\rm loc}=\sum_{A\in\mathcal A}L_A.
\]
\end{definition}

\begin{proposition}[Separated local updates have zero two-region commutator]
\label{prop:V3-local-zero}
Let $f$ be $S$-local and $g$ be $T$-local.  If no update set
$A\in\mathcal A$ meets both $S$ and $T$, then
\begin{equation}
 \boxed{\mathfrak C_{L^{\rm loc}}(f,g)=0.}
 \label{eq:V3-local-zero}
\end{equation}
In particular, if every $A$ has diameter at most $r$ and
$\operatorname{dist}(S,T)>r$, then \eqref{eq:V3-local-zero} holds.
\end{proposition}

\begin{proof}
Fix $A$ and a transition with $q_A(x,y)>0$.  If $A\cap S=\varnothing$, then
$f(y)=f(x)$; if $A\cap T=\varnothing$, then $g(y)=g(x)$.  At least one of
these alternatives holds by hypothesis.  Every summand in
\eqref{eq:V3-jump-commutator} therefore vanishes for $L_A$.  Sum over $A$.
The metric statement follows because an update meeting both $S$ and $T$ would
contain two points at distance greater than $r$.
\end{proof}

% -----------------------------------------------------------------------------
\subsection{The exact global-refresh defect}
\label{subsec:V3-refresh-defect}
% -----------------------------------------------------------------------------

\begin{theorem}[Global refresh has an explicit all-distance commutator]
\label{thm:V3-refresh-commutator}
Let
\[
 L^{\rm ref}=\rho(R-I),
 \qquad Rh=\mu[h]\mathbf1,
 \qquad \rho\ge0.
\]
Then for all $f,g,h$,
\begin{equation}
 \boxed{
 (\mathfrak C_{L^{\rm ref}}(f,g)h)(x)
 =\rho\int_\Omega
 (f(y)-f(x))(g(y)-g(x))h(y)\,\dd\mu(y).}
 \label{eq:V3-refresh-commutator}
\end{equation}
If $f$ and $g$ are real and centred, then
\begin{equation}
 \boxed{
 \mathfrak C_{L^{\rm ref}}(f,g)\mathbf1
 =\rho\bigl(fg+\mu[fg]\mathbf1\bigr),}
 \label{eq:V3-refresh-on-one}
\end{equation}
and hence
\begin{equation}
 \boxed{
 \lVert\mathfrak C_{L^{\rm ref}}(f,g)\rVert_{2\to2}
 \ge
 \rho\sqrt{\mu[f^2g^2]+3\mu[fg]^2}.}
 \label{eq:V3-refresh-lower}
\end{equation}
\end{theorem}

\begin{proof}
The refresh is the jump generator with $q(x,\dd y)=\rho\mu(\dd y)$, so
\eqref{eq:V3-refresh-commutator} is
\cref{thm:V3-jump-formula}.  With $h=\mathbf1$, expand the product of
increments.  The two linear terms vanish because $f$ and $g$ are centred,
which proves \eqref{eq:V3-refresh-on-one}.  Since
$\lVert\mathbf1\rVert_2=1$,
\begin{align*}
 \lVert\mathfrak C_{L^{\rm ref}}(f,g)\rVert_{2\to2}^2
 &\ge \rho^2\mu\!\left[(fg+\mu[fg])^2\right]\\
 &=\rho^2\bigl(\mu[f^2g^2]+3\mu[fg]^2\bigr).
\end{align*}
Taking square roots gives the claim.
\end{proof}

\begin{corollary}[Exact locality residual of the parent generator]
\label{cor:V3-parent-residual}
In \eqref{eq:V3-parent-generator}, suppose $f$ and $g$ obey the separated-update
hypothesis of \cref{prop:V3-local-zero}.  Then
\begin{equation}
 \boxed{
 \mathfrak C_{\mathcal L_X^{\rm par}}(f,g)
 =\mathfrak C_{\rho_X(\mathsf R_X-I)}(f,g).}
 \label{eq:V3-parent-residual}
\end{equation}
For real centred $f,g$, the right side has the lower bound
\eqref{eq:V3-refresh-lower} with $(\rho,\mu)=(\rho_X,\mu_X)$.
\end{corollary}

\begin{proof}
Use linearity in the generator and
\cref{prop:V3-local-zero,thm:V3-refresh-commutator}.
\end{proof}

\begin{theorem}[Cofinal refresh-rate necessity for this locality test]
\label{thm:V3-rho-must-vanish}
Let $X_n$ be a cutoff sequence.  Suppose there are real centred separated
local tests $f_n,g_n$ for which
\begin{equation}
 \mu_n[f_n^2g_n^2]\ge\chi>0
 \label{eq:V3-nondegenerate-separated-tests}
\end{equation}
and the local update commutator vanishes.  Then
\begin{equation}
 \lVert\mathfrak C_{\mathcal L_n^{\rm par}}(f_n,g_n)\rVert_{2\to2}
 \ge \rho_n\sqrt\chi.
 \label{eq:V3-cofinal-locality-lower}
\end{equation}
Consequently, vanishing of this normalized two-region locality residual forces
$\rho_n\to0$.  A fixed lower refresh rate $\rho_n\ge\rho_*>0$ is incompatible
with that locality requirement.
\end{theorem}

\begin{proof}
Apply \cref{cor:V3-parent-residual} and discard the nonnegative term
$3\mu_n[f_ng_n]^2$ in \eqref{eq:V3-refresh-lower}.
\end{proof}

\begin{remark}[What the test does and does not say]
\label{rem:V3-locality-scope}
For a local $f$, $L^{\rm ref}f=\rho(\mu[f]-f)$ is again a local function plus a
constant.  Thus \cref{thm:V3-rho-must-vanish} is not a claim that the refresh
necessarily destroys invariance of every local observable algebra.  It detects
the more specific fact that one elementary jump has simultaneous increments in
two separated regions.  Static projectivity, the stationary measure, and an
already proved Osterwalder--Schrader square are untouched.  The consequence is
for interpreting the parent path as finite-range physical noise or transporting
its refresh-created coercivity into a local continuum dynamics.
\end{remark}

\begin{counterexample}[Two Rademacher sites separate local heat baths from refresh]
\label{ex:V3-rademacher}
Let $\Omega=\{-1,1\}^2$ with uniform measure.  Let $E_1$ resample the first
coordinate uniformly while retaining the second, and let $E_2$ resample the
second while retaining the first.  For $a,b>0$ put
\[
 L^{\rm loc}=a(E_1-I)+b(E_2-I),
 \qquad f(x)=x_1,\qquad g(x)=x_2.
\]
Then
\begin{equation}
 \mathfrak C_{L^{\rm loc}}(f,g)=0,
 \qquad
 \lVert\mathfrak C_{\rho(R-I)}(f,g)\rVert_{2\to2}=\rho.
 \label{eq:V3-rademacher-separation}
\end{equation}
At the same time, $H^{\rm loc}=-L^{\rm loc}$ has centred eigenvalues
$a,b,a+b$, and therefore gap $\min(a,b)>0$.
\end{counterexample}

\begin{proof}
Each local update changes only one of the two coordinates, so the first equality
is \cref{prop:V3-local-zero}.  Expand
$h=c_0+c_1x_1+c_2x_2+c_{12}x_1x_2$ in the orthonormal Walsh basis.  Formula
\eqref{eq:V3-refresh-commutator} gives
\[
 \rho^{-1}\mathfrak C_{\rho(R-I)}(f,g)h
 =c_{12}-c_2x_1-c_1x_2+c_0x_1x_2.
\]
The coefficient map
$(c_0,c_1,c_2,c_{12})\mapsto(c_{12},-c_2,-c_1,c_0)$ is orthogonal, proving the
operator-norm identity.  Direct application of $E_1,E_2$ to
$x_1,x_2,x_1x_2$ gives the three stated eigenvalues.
\end{proof}

% -----------------------------------------------------------------------------
\subsection{The coercivity--locality fork}
\label{subsec:V3-coercivity-fork}
% -----------------------------------------------------------------------------

Let $\mathcal H_0=L^2_0(\mu)$ and put
\[
 H=-L^{\rm loc}|_{\mathcal H_0}\succeq0,
 \qquad H_\rho=H+\rho I.
\]
Assume $H$ is bounded and self-adjoint; this covers a finite sum of bounded
reversible conditional-expectation updates even when the cutoff field space is
continuous.  The symbol $H^\dagger$ denotes the spectral pseudoinverse on
$\overline{\Ran H}$, possibly unbounded.  Every occurrence below is either
under a positive gap or on a stated source range where the displayed form is
bounded.

\begin{proposition}[Exact source of the parent Green bound]
\label{prop:V3-gap-shift}
For $\rho>0$, if $\lambda_{\rm loc}=\inf\sigma(H)\ge0$, then
\begin{equation}
 \boxed{
 \operatorname{gap}(H_\rho)=\lambda_{\rm loc}+\rho,
 \qquad
 \lVert H_\rho^{-1}\rVert
 =\frac1{\lambda_{\rm loc}+\rho}.}
 \label{eq:V3-gap-shift}
\end{equation}
Thus the universal estimate $\lVert H_\rho^{-1}\rVert\le\rho^{-1}$ contains
no lower bound on $\lambda_{\rm loc}$.  If the locality test forces
$\rho_n\to0$, this estimate alone becomes nonuniform.
\end{proposition}

\begin{proof}
Adding $\rho I$ translates the spectrum of $H$ by $\rho$.  The spectral
infimum and the norm of the positive inverse are therefore exactly as
displayed.
\end{proof}

\begin{theorem}[Uniform-gap de-refreshing of arbitrary responses]
\label{thm:V3-uniform-gap-derefresh}
Let $H_n\succeq\lambda_*I$ on $\mathcal H_{0,n}$ for one $\lambda_*>0$, and let
$\rho_n\ge0$.  Then
\begin{equation}
 \boxed{
 \lVert H_n^{-1}-(H_n+\rho_nI)^{-1}\rVert
 \le
 \frac{\rho_n}{\lambda_*(\lambda_*+\rho_n)}
 \le\frac{\rho_n}{\lambda_*^2}.}
 \label{eq:V3-gap-resolvent-bound}
\end{equation}
Consequently, for every source synthesis $B_n:U_n\to\mathcal H_{0,n}$,
\begin{equation}
 \lVert H_n^{-1}B_n-(H_n+\rho_nI)^{-1}B_n\rVert
 \le \frac{\rho_n}{\lambda_*^2}\lVert B_n\rVert.
 \label{eq:V3-gap-response-bound}
\end{equation}
In particular, $\rho_n\to0$ and $\sup_n\lVert B_n\rVert<\infty$ give uniform
response de-refreshing.
\end{theorem}

\begin{proof}
The spectral multiplier of the difference is
\[
 \frac1\lambda-\frac1{\lambda+\rho_n}
 =\frac{\rho_n}{\lambda(\lambda+\rho_n)}.
\]
It decreases for $\lambda>0$, so its supremum on
$[\lambda_*,\infty)$ is the first right side in
\eqref{eq:V3-gap-resolvent-bound}.  Composition with $B_n$ proves the response
bound.
\end{proof}

The uniform-gap hypothesis is strong.  The next result isolates what can still
be transported when only the selected source range is regular relative to the
local generator.

\begin{theorem}[Source-relative Green de-refreshing without a gap]
\label{thm:V3-source-form-derefresh}
Let $H\succeq0$ and let $B:U\to\mathcal H_0$ factor as
\begin{equation}
 B=HC
 \label{eq:V3-one-factor-source}
\end{equation}
for a bounded $C:U\to\mathcal H_0$.  Define the local and refreshed source
Green forms
\[
 G_0(B)=B^*H^\dagger B,
 \qquad
 G_\rho(B)=B^*(H+\rho I)^{-1}B
 \quad(\rho>0).
\]
Then
\begin{equation}
 \boxed{
 0\preceq G_0(B)-G_\rho(B)
 \preceq\rho C^*C
 \preceq\rho\lVert C\rVert^2I_U.}
 \label{eq:V3-source-form-bound}
\end{equation}
Hence, along a cutoff sequence, $\rho_n\to0$ and
$\sup_n\lVert C_n\rVert<\infty$ transport these selected Green quadratic forms
even when the global local gap tends to zero.
\end{theorem}

\begin{proof}
By \eqref{eq:V3-one-factor-source} and the spectral theorem,
\begin{align*}
 G_0(B)-G_\rho(B)
 &=C^*\left(H-H(H+\rho I)^{-1}H\right)C\\
 &=C^*\left(\rho H(H+\rho I)^{-1}\right)C.
\end{align*}
The scalar multiplier $\rho\lambda/(\lambda+\rho)$ lies in $[0,\rho]$ for
$\lambda\ge0$.  Functional calculus gives all three operator inequalities.
\end{proof}

One generator factor is enough for the quadratic Green read but is not enough
for uniform convergence of the response vector.  A fractional additional
factor supplies the missing infrared control.

\begin{theorem}[Gap-free response de-refreshing from source regularity]
\label{thm:V3-fractional-source-derefresh}
Let $0<\alpha\le1$ and let $\rho>0$.  Suppose
\begin{equation}
 B=H^{1+\alpha}D
 \label{eq:V3-fractional-source}
\end{equation}
for a bounded $D:U\to\mathcal H_0$.  Then
\begin{equation}
 \boxed{
 \lVert H^\dagger B-(H+\rho I)^{-1}B\rVert
 \le c_\alpha\rho^\alpha\lVert D\rVert,}
 \label{eq:V3-fractional-response-bound}
\end{equation}
where
\begin{equation}
 c_\alpha=
 \begin{cases}
  \alpha^\alpha(1-\alpha)^{1-\alpha},&0<\alpha<1,\\
  1,&\alpha=1.
 \end{cases}
 \label{eq:V3-fractional-constant}
\end{equation}
This estimate is independent of the least positive eigenvalue of $H$.
\end{theorem}

\begin{proof}
On a spectral value $\lambda>0$, the multiplier of the response difference is
\[
 \lambda^\alpha-\frac{\lambda^{1+\alpha}}{\lambda+\rho}
 =\frac{\rho\lambda^\alpha}{\lambda+\rho}
 =\rho^\alpha\frac{(\lambda/\rho)^\alpha}{1+\lambda/\rho}.
\]
It is zero at $\lambda=0$.  For $0<\alpha<1$, elementary differentiation
shows
\[
 \sup_{z\ge0}\frac{z^\alpha}{1+z}
 =\alpha^\alpha(1-\alpha)^{1-\alpha};
\]
for $\alpha=1$ the supremum is $1$.  Functional calculus and composition with
$D$ prove the result.
\end{proof}

\begin{counterexample}[One source factor does not control response vectors]
\label{ex:V3-one-factor-insufficient}
Let $\rho_n\downarrow0$, let $\mathcal H_{0,n}=\R$, and let
\[
 H_n=\rho_n^2I,
 \qquad C_n=I,
 \qquad B_n=H_nC_n.
\]
Then
\begin{equation}
 H_n^\dagger B_n=I,
 \qquad
 (H_n+\rho_nI)^{-1}B_n=\frac{\rho_n}{1+\rho_n}I,
 \label{eq:V3-response-counterexample}
\end{equation}
so the response-vector difference tends to $I$.  Nevertheless,
\begin{equation}
 0\le G_0(B_n)-G_{\rho_n}(B_n)
 =\frac{\rho_n^2}{1+\rho_n}I
 \le\rho_n I,
 \label{eq:V3-form-counterexample}
\end{equation}
in agreement with \cref{thm:V3-source-form-derefresh}.
\end{counterexample}

\begin{proof}
All assertions follow by scalar substitution.  This one-dimensional family is
also a sharp warning that convergence of selected Green costs is weaker than
convergence of the corresponding response vectors.
\end{proof}

% -----------------------------------------------------------------------------
\subsection{Consequences for the RPESM continuum route}
\label{subsec:V3-RPESM-consequences}
% -----------------------------------------------------------------------------

\begin{theorem}[Refresh-locality and Green-transport trichotomy]
\label{thm:V3-trichotomy}
Consider the RPESM parent generators
\eqref{eq:V3-parent-generator} on a cofinal sequence and suppose the declared
local part is carried by uniformly finite-range update families.  For any
nondegenerate separated tests satisfying
\eqref{eq:V3-nondegenerate-separated-tests}, the following alternatives are
rigorous.
\begin{enumerate}[label=\textup{(R\arabic*)},leftmargin=3.2em]
\item If $\inf_n\rho_n>0$ and the parent generator itself is retained as the
      proposed physical path generator, its simultaneous-jump locality residual
      stays bounded away from zero.
\item If $\rho_n\to0$ and the local centred generators have one uniform gap,
      arbitrary uniformly bounded source responses de-refresh by
      \eqref{eq:V3-gap-response-bound}.
\item Without a uniform gap, selected Green forms still de-refresh under the
      source factorization $B_n=H_nC_n$ with uniformly bounded $C_n$; selected
      response vectors de-refresh under
      $B_n=H_n^{1+\alpha}D_n$ with one $\alpha\in(0,1]$ and uniformly bounded $D_n$.
\end{enumerate}
If none of (R2)--(R3) is populated, the parent Green bound cannot by itself be
transported to a local zero-mass response.
\end{theorem}

\begin{proof}
Item (R1) is \cref{thm:V3-rho-must-vanish}.  Item (R2) is
\cref{thm:V3-uniform-gap-derefresh}.  The two claims in (R3) are
\cref{thm:V3-source-form-derefresh,thm:V3-fractional-source-derefresh}.
\Cref{prop:V3-gap-shift} proves the last sentence.
\end{proof}

\begin{assessment}[How this changes the continuum audit]
\label{ass:V3-continuum-audit}
The finite RPESM law, its exact projectivity, and its reflected square may use
the global refresh as an auxiliary population device.  The existing exact
semigroup identity then reconstructs the local finite-cutoff dynamics.  The
remaining continuum statement must be made on that reconstructed local
generator, with estimates uniform after the refresh contribution is removed.

There are now explicit source-native rows to acquire.  If the Einstein--matter
handoff uses only a Green quadratic form $B_X^*H_X^\dagger B_X$, it is enough
for this interface to construct a uniformly bounded solution writer
$C_X$ to $B_X=H_XC_X$.  If it requires the response vector itself, a uniform
local gap or a stronger fractional range estimate is needed.  The latter can
be tested on the same physical source bank; it is not a claim about the entire
spectrum.  The parent estimate $\lVert(H_X+\rho_XI)^{-1}\rVert\le\rho_X^{-1}$
does not populate either row.
\end{assessment}

\begin{openproblem}[Executable V4 acquisition: local generator source regularity]
\label{op:V3-next-acquisition}
For one named RPESM source bank $B_X$, after applying the already proved exact
de-refreshing identity, acquire the following data on two adjacent cutoffs:
\begin{enumerate}[label=\textup{(D\arabic*)},leftmargin=3.2em]
\item the literal local update blocks and their maximal physical diameter;
\item two normalized separated local record functions with a certified lower
      bound for $\mu_X[f_X^2g_X^2]$;
\item the actual refresh rates $\rho_X$ in physical clock units;
\item either a local centred gap lower bound, a factorization
      $B_X=H_XC_X$, or a fractional factorization
      $B_X=H_X^{1+\alpha}D_X$;
\item a cross-cutoff bound for the chosen factor writer, not merely for the
      refreshed inverse.
\end{enumerate}
The outputs are: a local-response transport bound; a source-Green-only bound;
or a typed infrared/source-range obstruction.  Before this card is populated,
the correct status is that the finite parent construction is exact while
local zero-mass Green transport remains conditional.
\end{openproblem}

\begin{theorem}[Version V3 terminal theorem]
\label{thm:V3-terminal}
Preserving V1 and V2 verbatim before this appended block, V3 proves:
\begin{enumerate}[label=\textup{(V3.\arabic*)},leftmargin=5.2em]
\item the double multiplication commutator of a bounded pure-jump generator
      is the
      exact simultaneous-increment kernel
      \eqref{eq:V3-jump-commutator};
\item finite-range updates give zero commutator on sufficiently separated
      records, while global refresh gives the explicit nonzero operator
      \eqref{eq:V3-refresh-commutator};
\item a nondegenerate cofinal separated-record test forces the refresh rate to
      vanish if this instantaneous locality residual is to vanish;
\item the refresh Green bound is exactly a spectral shift and supplies no
      uniform estimate for the de-refreshed local inverse;
\item a uniform local gap transports arbitrary responses, one local-generator
      source factor transports Green forms, and one fractional additional
      factor transports response vectors without a uniform gap;
\item exact finite examples separate all three claims.
\end{enumerate}
\end{theorem}

\begin{proof}
Items (V3.1)--(V3.3) are
\cref{thm:V3-jump-formula,prop:V3-local-zero,thm:V3-refresh-commutator,thm:V3-rho-must-vanish}.
Item (V3.4) is \cref{prop:V3-gap-shift}.  Item (V3.5) is
\cref{thm:V3-uniform-gap-derefresh,thm:V3-source-form-derefresh,thm:V3-fractional-source-derefresh}.
The exact controls are
\cref{ex:V3-rademacher,ex:V3-one-factor-insufficient}.
\end{proof}

\begin{assessment}[Programme boundary after V3]
\label{ass:V3-boundary}
V3 does not prove or disprove the existence of the full interacting
Standard-Model--Einstein continuum.  It removes one ambiguity in the proposed
route: auxiliary global refresh may populate a finite reflected path, but its
simultaneous-jump locality and its inverse bound cannot be promoted as local
continuum physics.  Exact finite de-refreshing is already available; the new
remaining burden is a uniform local-generator or source-range estimate.

The four-cell physical curvature read from V2 remains unpopulated because the
exact $81$-record weights are absent.  The broader continuum still requires the
declared uniform reflection, tightness, quasilocal product, determinant and
Pfaffian, anomaly, stress-tensor, renormalized coupling, Lorentzian, and
Einstein closure rows.  No assertion in V3 substitutes the refresh diagnostic
for those independent gates.
\end{assessment}

\section*{V3 exact checks and append-only record}
\addcontentsline{toc}{section}{V3 exact checks and append-only record}
The Rademacher coefficient map in \cref{ex:V3-rademacher} is an orthogonal
signed permutation.  The scalar family in
\cref{ex:V3-one-factor-insufficient} gives exactly
\[
 \frac{\rho_n^2}{\rho_n^2+\rho_n}
 =\frac{\rho_n}{1+\rho_n},
 \qquad
 \rho_n^2-\frac{\rho_n^4}{\rho_n^2+\rho_n}
 =\frac{\rho_n^2}{1+\rho_n}.
\]
These identities, the maximum in
\eqref{eq:V3-fractional-constant}, and the operator inequalities above are
independently replayed in the structural verification of this source.

The complete V2 source before removal of its final active document-closing
command is retained verbatim.  Its complete-file SHA--256 digest is
\begin{center}\small\ttfamily
9652d8352d2f9166d6c1aafac12ac3203e7f20debe70e8a3b13283a78fcacb3f.
\end{center}
A Version V4 must retain the complete V3 source, remove only the final active
document-closing command, append new material, restore one final command, and
use a filename ending in \texttt{\_V4.tex}.

% =============================================================================
% END APPENDED VERSION V3 MATERIAL
% =============================================================================


% =============================================================================
% BEGIN APPENDED VERSION V4 MATERIAL
% =============================================================================

\clearpage
\section{Version V4: projective de-refreshing and a refresh-free ultraviolet gap}
\label{sec:V4-projective-derefresh}

Version V3 showed that a positive global reset rate leaves an all-distance
simultaneous-jump signature, even when every other update is local.  The
authoritative RPESM definition makes the issue concrete: equation \textup{(RP.8)}
uses one fixed $\rho>0$, while \textup{(RP.14)} requires exact transport of the
old Markov dynamics.  The fixed rate is therefore part of the projective
construction, not a disposable cutoff parameter.

This version resolves the apparent conflict at finite time.  With a common
rate, exact projectivity of the parent semigroups descends exactly to the
de-refreshed local semigroups.  Thus the parent path can remain an auxiliary
population device while the local dynamics retain exact projectivity.  A
rate-rigidity theorem shows why varying $\rho_n\to0$ inside the same exact
intertwining diagram is not an equivalent repair.

The remaining difficulty is the zero-mass Green integral.  De-refreshing
multiplies the centred parent kernel by $e^{\rho t}$ and cancels every unit of
decay supplied by the refresh.  We derive the exact spectral and time-tail
criterion for the local Green form.  Finally, a new bounded-density comparison
for conditional heat baths applies to the manuscript's uniformly equivalent
ultraviolet tail law and gives an explicit local, refresh-free gap.  This closes
one selected ultraviolet row; it does not manufacture the interacting
infrared or full Standard-Model--Einstein continuum.

\begin{assessment}[Imported facts and new work in V4]
\label{ass:V4-nonduplication}
The exact fixed-cutoff de-refresh identity, the statement that a refresh gap is
not Yang--Mills coercivity, the RPESM projective kernel \textup{(RP.13)}, the
bounded ultraviolet likelihood \textup{(RS.11)}, and the Higgs scaling
trichotomy \textup{(RS.13)--(RS.14)} are imported.  V4 proves the commuting
projective de-refresh square, refresh-rate rigidity, the refresh-weighted
zero-mass Green criterion, the heat-bath density-comparison theorem, and its
explicit application to the selected ultraviolet tail.  None of those five
bridges occurs in the supplied sources.
\end{assessment}

% -----------------------------------------------------------------------------
\subsection{The projective de-refresh square}
\label{subsec:V4-projective-square}
% -----------------------------------------------------------------------------

Let $(\Omega_n,\mu_n)$ and $(\Omega_m,\mu_m)$ be two cutoff probability spaces,
$m\succeq n$, with a measurable projection
$\pi_{m/n}:\Omega_m\to\Omega_n$ satisfying
$(\pi_{m/n})_*\mu_m=\mu_n$.  Its pullback
\[
 J_{m/n}:L^2(\mu_n)\longrightarrow L^2(\mu_m),
 \qquad J_{m/n}f=f\circ\pi_{m/n},
\]
is an isometry.  Put
\[
 R_kf=\mu_k[f]\mathbf1,\qquad Q_k=I-R_k
 \quad(k=n,m).
\]

\begin{lemma}[Constants and centred subspaces transport exactly]
\label{lem:V4-RQ-intertwine}
One has
\begin{equation}
 \boxed{
 R_mJ_{m/n}=J_{m/n}R_n,
 \qquad
 Q_mJ_{m/n}=J_{m/n}Q_n.}
 \label{eq:V4-RQ-intertwine}
\end{equation}
\end{lemma}

\begin{proof}
For $f\in L^2(\mu_n)$,
\[
 \mu_m[J_{m/n}f]
 =\int_{\Omega_m}f\circ\pi_{m/n}\,\dd\mu_m
 =\int_{\Omega_n}f\,\dd\mu_n.
\]
Both sides of the first identity are the corresponding constant function.
Subtracting it from $J_{m/n}$ proves the second identity.
\end{proof}

For each cutoff let $S_k^{\rm par}(t)$ and $S_k^{\rm loc}(t)$ be the parent
and local Markov semigroups, related at one common rate $\rho>0$ by
\begin{equation}
 \begin{aligned}
 S_k^{\rm par}(t)
 &=R_k+e^{-\rho t}Q_kS_k^{\rm loc}(t)Q_k,\\
 S_k^{\rm loc}(t)
 &=R_k+e^{\rho t}Q_kS_k^{\rm par}(t)Q_k.
 \end{aligned}
 \label{eq:V4-imported-derefresh}
\end{equation}

\begin{theorem}[Exact projective descent after de-refreshing]
\label{thm:V4-projective-derefresh}
Fix $t\ge0$.  If the parent kernel transports exactly,
\begin{equation}
 S_m^{\rm par}(t)J_{m/n}
 =J_{m/n}S_n^{\rm par}(t),
 \label{eq:V4-parent-intertwine}
\end{equation}
then the de-refreshed local kernel also transports exactly:
\begin{equation}
 \boxed{
 S_m^{\rm loc}(t)J_{m/n}
 =J_{m/n}S_n^{\rm loc}(t).}
 \label{eq:V4-local-intertwine}
\end{equation}
If \eqref{eq:V4-parent-intertwine} holds for all $t\ge0$, then so does
\eqref{eq:V4-local-intertwine}.  For bounded generators, differentiation at
$t=0$ gives
\begin{equation}
 L_m^{\rm loc}J_{m/n}=J_{m/n}L_n^{\rm loc}.
 \label{eq:V4-local-generator-intertwine}
\end{equation}
\end{theorem}

\begin{proof}
Use the inverse identity in \eqref{eq:V4-imported-derefresh} and
\cref{lem:V4-RQ-intertwine}:
\begin{align*}
 S_m^{\rm loc}(t)J_{m/n}
 &=R_mJ_{m/n}
   +e^{\rho t}Q_mS_m^{\rm par}(t)Q_mJ_{m/n}\\
 &=J_{m/n}R_n
   +e^{\rho t}Q_mS_m^{\rm par}(t)J_{m/n}Q_n\\
 &=J_{m/n}R_n
   +e^{\rho t}J_{m/n}Q_nS_n^{\rm par}(t)Q_n\\
 &=J_{m/n}S_n^{\rm loc}(t).
\end{align*}
The all-time statement follows pointwise in $t$.  Bounded semigroups are norm
differentiable at zero, so differentiating the intertwining identity proves
\eqref{eq:V4-local-generator-intertwine}.
\end{proof}

\begin{corollary}[Application to the RPESM shell kernel]
\label{cor:V4-RPESM-projective-local}
Suppose the one-step kernels in \textup{(RP.13)--(RP.14)} are parent kernels
with the common rate fixed in \textup{(RP.8)} and satisfy the imported
de-refresh identity.  Then their local de-refreshed one-step kernels satisfy
the same exact old-observable intertwining \textup{(RP.14)}.  If the complete
parent semigroups intertwine, the complete local semigroups and their bounded
generators intertwine.
\end{corollary}

\begin{proof}
Take $J_{m/n}$ to be pullback by the RPESM field projection and apply
\cref{thm:V4-projective-derefresh}.  Static projectivity supplies the isometry
and \cref{lem:V4-RQ-intertwine}.
\end{proof}

This result identifies a viable logical order: construct the reflected parent
path, use it to populate the finite event law, de-refresh at every cutoff, and
take the projective continuum limit on the local kernels.

% -----------------------------------------------------------------------------
\subsection{Refresh-rate rigidity inside an exact projective diagram}
\label{subsec:V4-rate-rigidity}
% -----------------------------------------------------------------------------

The use of one common rate in the previous theorem is necessary if both the
parent and local generators are required to intertwine.

\begin{theorem}[Exact adjacent-cutoff defect identity]
\label{thm:V4-rate-defect}
Let
\[
 L_k^{\rm par}=\rho_k(R_k-I)+L_k^{\rm loc}
 =-\rho_kQ_k+L_k^{\rm loc},
 \qquad k=n,m.
\]
Then
\begin{equation}
 \boxed{
 L_m^{\rm par}J_{m/n}-J_{m/n}L_n^{\rm par}
 =
 L_m^{\rm loc}J_{m/n}-J_{m/n}L_n^{\rm loc}
 -(\rho_m-\rho_n)J_{m/n}Q_n.}
 \label{eq:V4-rate-defect}
\end{equation}
Consequently, if the parent generators intertwine, their local defect is
forced to be
\begin{equation}
 \boxed{
 L_m^{\rm loc}J_{m/n}-J_{m/n}L_n^{\rm loc}
 =(\rho_m-\rho_n)J_{m/n}Q_n.}
 \label{eq:V4-compensating-defect}
\end{equation}
\end{theorem}

\begin{proof}
Expand both parent generators and use
$Q_mJ_{m/n}=J_{m/n}Q_n$ from
\cref{lem:V4-RQ-intertwine}.  Rearrangement gives both formulas.
\end{proof}

\begin{corollary}[Refresh-rate rigidity]
\label{cor:V4-rate-rigidity}
Assume the parent and local generators both intertwine and the old centred
space $Q_nL^2(\mu_n)$ is nonzero.  Then
\begin{equation}
 \boxed{\rho_m=\rho_n.}
 \label{eq:V4-rate-rigidity}
\end{equation}
Along a connected nested cutoff sequence, all refresh rates are equal.
\end{corollary}

\begin{proof}
Both defects on the left of \eqref{eq:V4-rate-defect} vanish.  Choose a nonzero
$f\in Q_nL^2(\mu_n)$.  The remaining identity is
$(\rho_m-\rho_n)J_{m/n}f=0$.  The pullback is an isometry, so
$\rho_m-\rho_n=0$.
\end{proof}

\begin{theorem}[Projectivity--parent-locality incompatibility]
\label{thm:V4-projectivity-locality-incompatibility}
Suppose along a nontrivial nested sequence that:
\begin{enumerate}[label=\textup{(\roman*)},leftmargin=2.8em]
\item both parent and local bounded generators intertwine exactly;
\item every $\rho_n>0$;
\item the separated local tests of \cref{thm:V3-rho-must-vanish} exist with one
      nondegeneracy constant $\chi>0$;
\item the parent generators themselves are required to have vanishing
      simultaneous-jump locality residual.
\end{enumerate}
Then the four requirements are inconsistent.
\end{theorem}

\begin{proof}
By \cref{cor:V4-rate-rigidity}, $\rho_n=\rho_1>0$ for every $n$.  By
\cref{thm:V3-rho-must-vanish}, requirement (iv) forces $\rho_n\to0$, a
contradiction.
\end{proof}

\begin{assessment}[The correct branch after rate rigidity]
\label{ass:V4-auxiliary-branch}
The contradiction concerns interpreting the \emph{parent} generator as the
local physical dynamics.  It does not obstruct the auxiliary-parent branch.
On that branch the common positive rate is retained to preserve the exact
parent projective packet, and \cref{thm:V4-projective-derefresh} removes it
before spatial locality or physical local Green response is asserted.

If one instead changes $\rho_n$ across cutoffs while retaining exact parent
projectivity, \eqref{eq:V4-compensating-defect} is the mandatory local-generator
defect.  Its global centred-projection form must be acquired and justified; it
cannot be hidden in an unspecified local correction.
\end{assessment}

% -----------------------------------------------------------------------------
\subsection{The exact refresh-weighted Green criterion}
\label{subsec:V4-weighted-Green}
% -----------------------------------------------------------------------------

Let $H\succeq0$ be the nonnegative self-adjoint local generator on the centred
Hilbert space $\mathcal H_0$, and let
\[
 T^{\rm loc}(t)=e^{-tH},
 \qquad
 T^{\rm par}(t)=e^{-t(H+\rho I)}
 =e^{-\rho t}T^{\rm loc}(t).
\]
Let $U$ be finite dimensional and let $B:U\to\mathcal H_0$ be a bounded source
synthesis.  Define the source kernels
\[
 K_B^{\rm loc}(t)=B^*T^{\rm loc}(t)B,
 \qquad
 K_B^{\rm par}(t)=B^*T^{\rm par}(t)B.
\]

\begin{theorem}[Weighted parent kernel reconstructs the local Green form]
\label{thm:V4-weighted-Green}
For every $t,T\ge0$,
\begin{equation}
 \boxed{
 K_B^{\rm loc}(t)=e^{\rho t}K_B^{\rm par}(t),}
 \label{eq:V4-kernel-derefresh}
\end{equation}
\begin{equation}
 \boxed{
 G_{B,T}^{\rm loc}
 :=\int_0^T K_B^{\rm loc}(t)\,\dd t
 =\int_0^T e^{\rho t}K_B^{\rm par}(t)\,\dd t.}
 \label{eq:V4-truncated-Green}
\end{equation}
Let $E_H$ be the spectral measure of $H$, let $P_0=E_H(\{0\})$, and for
$u\in U$ put
\[
 \nu_u(A)=\lVert E_H(A)Bu\rVert^2.
\]
Then
\begin{equation}
 \boxed{
 \begin{aligned}
 \langle u,G_{B,T}^{\rm loc}u\rangle
 ={}&T\lVert P_0Bu\rVert^2\\
 &+\int_{(0,\infty)}
   \frac{1-e^{-T\lambda}}{\lambda}\,\dd\nu_u(\lambda).
 \end{aligned}}
 \label{eq:V4-spectral-truncated-Green}
\end{equation}
The increasing limit $G_B^{\rm loc}=\lim_{T\to\infty}G_{B,T}^{\rm loc}$ is a
finite operator on $U$ if and only if, for every $u\in U$,
\begin{equation}
 P_0Bu=0,
 \qquad
 \int_{(0,\infty)}\lambda^{-1}\,\dd\nu_u(\lambda)<\infty.
 \label{eq:V4-zero-mass-domain}
\end{equation}
On this domain, let
$H^{\dagger/2}=\int_{(0,\infty)}\lambda^{-1/2}\,\dd E_H(\lambda)$.
Then $H^{\dagger/2}B$ is bounded and
\begin{equation}
 \boxed{
 G_B^{\rm loc}
 =(H^{\dagger/2}B)^*(H^{\dagger/2}B)
 =:B^*H^\dagger B
 =\int_0^\infty e^{\rho t}K_B^{\rm par}(t)\,\dd t.}
 \label{eq:V4-local-Green-reconstruction}
\end{equation}
Thus $B^*H^\dagger B$ is a quadratic-form compression; the stronger condition
$Bu\in\operatorname{Dom}H^\dagger$ is not assumed.
By contrast,
\begin{equation}
 \boxed{
 G_B^{\rm par}=B^*(H+\rho I)^{-1}B
 =\int_0^\infty K_B^{\rm par}(t)\,\dd t
 \preceq\rho^{-1}B^*B}
 \label{eq:V4-parent-Green}
\end{equation}
is always finite and does not imply \eqref{eq:V4-zero-mass-domain}.
\end{theorem}

\begin{proof}
The first identity is the scalar commutation
$e^{-t(H+\rho I)}=e^{-\rho t}e^{-tH}$; integration gives the second.
For a fixed $u$, the spectral theorem and Tonelli's theorem give
\[
 \int_0^T\langle Bu,e^{-tH}Bu\rangle\,\dd t
 =\int_{[0,\infty)}
   \left(\int_0^Te^{-t\lambda}\,\dd t\right)\dd\nu_u(\lambda).
\]
The inner integral is $T$ at $\lambda=0$ and
$(1-e^{-T\lambda})/\lambda$ otherwise, which proves
\eqref{eq:V4-spectral-truncated-Green}.  Monotone convergence gives exactly
the two conditions in \eqref{eq:V4-zero-mass-domain}.  Since $U$ is finite
dimensional, finiteness of every quadratic form makes
$H^{\dagger/2}B$ bounded.  Polarization identifies the limit with
$(H^{\dagger/2}B)^*(H^{\dagger/2}B)$.  Applying the spectral
theorem to $H+\rho I\succeq\rho I$ proves
\eqref{eq:V4-parent-Green}.
\end{proof}

The theorem identifies the exact datum lost by the parent bound: decay strictly
faster than the refresh factor.

\begin{corollary}[Uniform excess decay is sufficient]
\label{cor:V4-excess-decay}
Let $B_n:U\to\mathcal H_{0,n}$ be a cofinal source bank and suppose, for fixed
$C,\gamma>0$ and one positive matrix $M$ on $U$,
\begin{equation}
 0\preceq K_{B_n}^{\rm par}(t)
 \preceq C e^{-(\rho+\gamma)t}M
 \quad(t\ge0,\ n\ge1).
 \label{eq:V4-excess-decay}
\end{equation}
Then every local Green form exists and
\begin{equation}
 \boxed{
 0\preceq G_{B_n}^{\rm loc}\preceq\frac C\gamma M,
 \qquad
 0\preceq
 \int_T^\infty K_{B_n}^{\rm loc}(t)\,\dd t
 \preceq\frac C\gamma e^{-\gamma T}M.}
 \label{eq:V4-excess-tail}
\end{equation}
\end{corollary}

\begin{proof}
Multiply \eqref{eq:V4-excess-decay} by $e^{\rho t}$ and integrate over
$[0,\infty)$ and $[T,\infty)$.
\end{proof}

\begin{theorem}[Cofinal weighted-tail criterion]
\label{thm:V4-cofinal-weighted-tail}
Suppose $K_{B_n}^{\rm par}(t)$ converges in norm locally uniformly in $t$ after
the declared source and cutoff identifications.  If
\begin{equation}
 \boxed{
 \lim_{T\to\infty}\sup_n
 \left\|
 \int_T^\infty e^{\rho t}K_{B_n}^{\rm par}(t)\,\dd t
 \right\|=0,}
 \label{eq:V4-uniform-weighted-tail}
\end{equation}
then the local zero-mass Green matrices converge in norm.  Conversely, a
uniform Green-tail conclusion obtained through the parent kernels is exactly
a bound of the form \eqref{eq:V4-uniform-weighted-tail}; the unweighted parent
bound is insufficient.
\end{theorem}

\begin{proof}
Given $\varepsilon>0$, choose $T$ so that both weighted tails are below
$\varepsilon$.  On $[0,T]$, local uniform convergence and
\eqref{eq:V4-truncated-Green} give convergence of the integrals.  The two tails
then bound the full difference by $2\varepsilon$ plus the finite-window
difference.  The converse is the definition of uniform tail exhaustion after
using \eqref{eq:V4-kernel-derefresh}.
\end{proof}

\begin{counterexample}[A parent Green bound can hide a local zero mode]
\label{ex:V4-refresh-hides-zero}
On the centred line of the uniform two-point space take $H=0$ and let
$B:\R\to\mathcal H_0$ synthesize the unit spin.  Then
\[
 K_B^{\rm par}(t)=e^{-\rho t},
 \qquad
 G_B^{\rm par}=\rho^{-1},
\]
while
\[
 e^{\rho t}K_B^{\rm par}(t)=1,
 \qquad
 G_{B,T}^{\rm loc}=T\longrightarrow\infty.
\]
Thus the finite parent Green value neither removes the local zero mode nor
defines a local zero-mass Green response.
\end{counterexample}

\begin{counterexample}[No zero mode, but no uniform local Green bound]
\label{ex:V4-collapsing-local-gap}
Let $\lambda_n\downarrow0$.  On the uniform two-point space let
\[
 L_n^{\rm loc}=\lambda_n(R-I)
\]
and let $B_n$ synthesize the unit centred spin.  Each local generator has no
centred zero mode, but
\begin{equation}
 \boxed{
 G_{B_n}^{\rm par}=\frac1{\rho+\lambda_n}\le\frac1\rho,
 \qquad
 G_{B_n}^{\rm loc}=\frac1{\lambda_n}\longrightarrow\infty.}
 \label{eq:V4-collapsing-gap}
\end{equation}
The parent kernel has excess decay $\gamma_n=\lambda_n$, and the collapse of
that excess is exactly the local Green obstruction.
\end{counterexample}

\begin{proof}
The centred spin is an eigenvector of $-L_n^{\rm loc}$ with eigenvalue
$\lambda_n$.  The displayed formulas are scalar spectral calculus.
\end{proof}

% -----------------------------------------------------------------------------
\subsection{A bounded-density comparison for conditional heat baths}
\label{subsec:V4-density-comparison}
% -----------------------------------------------------------------------------

We now construct a same-support route to a local gap which does not use global
refresh.  Let
$\Omega=\prod_{i\in I}\Omega_i$ be a finite product, or a countable product
with all statements initially made on cylinder functions.  Let $\nu$ and
$\mu$ be probability measures on the same measurable carrier, with
\begin{equation}
 \dd\mu=w\,\dd\nu,
 \qquad
 0<m\le w\le M<\infty.
 \label{eq:V4-density-window}
\end{equation}
Let $E_i^\nu$ and $E_i^\mu$ be conditional expectation with all coordinates
outside block $i$ fixed.  For rates $r_i\ge0$, define the heat-bath forms
\begin{align}
 \mathcal E_\nu(f)
 &=\sum_i r_i\,\mathbb E_\nu
   \operatorname{Var}_\nu(f\mid\mathcal F_{i^c}),\\
 \mathcal E_\mu(f)
 &=\sum_i r_i\,\mathbb E_\mu
   \operatorname{Var}_\mu(f\mid\mathcal F_{i^c}).
 \label{eq:V4-heatbath-forms}
\end{align}

\begin{theorem}[Same-support heat-bath gap comparison]
\label{thm:V4-heatbath-comparison}
If
\begin{equation}
 \mathcal E_\nu(f)\ge\lambda_\nu\operatorname{Var}_\nu(f)
 \label{eq:V4-reference-gap}
\end{equation}
for every form-core function, then
\begin{equation}
 \boxed{
 \mathcal E_\mu(f)
 \ge
 \left(\frac mM\right)^2
 \lambda_\nu\operatorname{Var}_\mu(f).}
 \label{eq:V4-compared-gap}
\end{equation}
The inequality extends to the closed heat-bath form whenever the cylinder core
is form dense.
\end{theorem}

\begin{proof}
Fix a block $i$ and write
\[
 \overline w_i
 =\mathbb E_\nu[w\mid\mathcal F_{i^c}].
\]
Then $m\le\overline w_i\le M$.  By disintegration, the conditional density is
\[
 \frac{\dd\mu_i(\,\cdot\mid x_{i^c})}
      {\dd\nu_i(\,\cdot\mid x_{i^c})}
 =\frac{w}{\overline w_i}\ge\frac mM.
\]
Using conditional variance as an infimum over constants gives
\[
 \operatorname{Var}_\mu(f\mid\mathcal F_{i^c})
 \ge\frac mM
 \operatorname{Var}_\nu(f\mid\mathcal F_{i^c}).
\]
The outside marginal of $\mu$ has density $\overline w_i$ relative to the
outside marginal of $\nu$, hence
\[
 \mathbb E_\mu\operatorname{Var}_\mu(f\mid\mathcal F_{i^c})
 \ge\frac{m^2}{M}
 \mathbb E_\nu\operatorname{Var}_\nu(f\mid\mathcal F_{i^c}).
\]
Multiply by $r_i$ and sum:
\[
 \mathcal E_\mu(f)\ge\frac{m^2}{M}\mathcal E_\nu(f).
\]
Finally,
\[
 \operatorname{Var}_\mu(f)
 =\inf_c\int|f-c|^2w\,\dd\nu
 \le M\operatorname{Var}_\nu(f),
\]
so $\operatorname{Var}_\nu(f)\ge M^{-1}\operatorname{Var}_\mu(f)$.
Combining the last three inequalities proves
\eqref{eq:V4-compared-gap}.  Closure follows by lower semicontinuity of the
forms.
\end{proof}

\begin{corollary}[Product reference and bounded interaction]
\label{cor:V4-product-comparison}
Suppose $\nu=\bigotimes_i\nu_i$, every $r_i\ge r_*>0$, and $f$ is a cylinder
function.  Then
\begin{equation}
 \mathcal E_\nu(f)\ge r_*\operatorname{Var}_\nu(f).
 \label{eq:V4-product-gap}
\end{equation}
Consequently \eqref{eq:V4-density-window} implies
\begin{equation}
 \boxed{
 \mathcal E_\mu(f)
 \ge r_*\left(\frac mM\right)^2
 \operatorname{Var}_\mu(f).}
 \label{eq:V4-product-perturbed-gap}
\end{equation}
If $w=e^{-V}/\nu[e^{-V}]$ and
$\operatorname*{ess\,osc}V\le A$, the lower constant is
$r_*e^{-2A}$.
\end{corollary}

\begin{proof}
The Efron--Stein inequality for the product law is
\[
 \operatorname{Var}_\nu(f)
 \le\sum_i\mathbb E_\nu
   \operatorname{Var}_\nu(f\mid\mathcal F_{i^c}).
\]
For completeness, expose the finitely many coordinates on which $f$ depends,
write its Doob martingale, and decompose its variance into the orthogonal
martingale differences.  Replacing one revealed coordinate by an independent
copy bounds each difference by the corresponding conditional variance; summing
gives the displayed inequality.  Multiplication by $r_*$ proves
\eqref{eq:V4-product-gap}, and
\cref{thm:V4-heatbath-comparison} proves
\eqref{eq:V4-product-perturbed-gap}.  For an exponential density,
$M/m\le e^A$.
\end{proof}

% -----------------------------------------------------------------------------
\subsection{An explicit refresh-free gap on the selected RPESM ultraviolet tail}
\label{subsec:V4-UV-gap}
% -----------------------------------------------------------------------------

The manuscript's ultraviolet equivalence theorem supplies a product-tail
reference $\nu_\infty$ and a physical tail law $\mu_\infty$ with
\begin{equation}
 e^{-5/64}\le
 L_\infty=\frac{\dd\mu_\infty}{\dd\nu_\infty}
 \le e^{5/64}.
 \label{eq:V4-imported-UV-window}
\end{equation}
On the cylinder core, equip every declared tail block with the
single-block conditional heat-bath form at rate $\rho_F>0$.

\begin{theorem}[Selected ultraviolet tail has a local gap without refresh]
\label{thm:V4-UV-local-gap}
Let $H_{\rm UV}$ be the Friedrichs operator of the physical-tail heat-bath
form just defined.  Then on centred tail cylinder functions,
\begin{equation}
 \boxed{
 \langle f,H_{\rm UV}f\rangle_{\mu_\infty}
 \ge
 \rho_F e^{-5/16}\lVert f\rVert_{L^2(\mu_\infty)}^2.}
 \label{eq:V4-UV-gap}
\end{equation}
Consequently,
\begin{equation}
 \boxed{
 \lVert H_{\rm UV}^{-1}\rVert_{L^2_0(\mu_\infty)}
 \le\frac{e^{5/16}}{\rho_F}.}
 \label{eq:V4-UV-Green-bound}
\end{equation}
For every bounded centred source synthesis $B:U\to L^2_0(\mu_\infty)$,
\begin{equation}
 \boxed{
 0\preceq B^*H_{\rm UV}^{-1}B
 \preceq\frac{e^{5/16}}{\rho_F}B^*B.}
 \label{eq:V4-UV-source-Green}
\end{equation}
These estimates use no global reset term.
\end{theorem}

\begin{proof}
The product reference heat bath has gap $\rho_F$ by
\cref{cor:V4-product-comparison}.  In
\eqref{eq:V4-imported-UV-window},
\[
 m=e^{-5/64},\qquad M=e^{5/64},
 \qquad
 \left(\frac mM\right)^2=e^{-20/64}=e^{-5/16}.
\]
Apply \cref{thm:V4-heatbath-comparison} on the cylinder core and close the
form.  The inverse and source bounds follow from the spectral theorem.
\end{proof}

\begin{corollary}[The ultraviolet tail has uniform excess decay]
\label{cor:V4-UV-excess-decay}
If an auxiliary global refresh of rate $\rho>0$ is added to
$H_{\rm UV}$, then every bounded centred ultraviolet source satisfies
\begin{equation}
 \boxed{
 K_B^{\rm par}(t)
 \preceq
 e^{-(\rho+\rho_Fe^{-5/16})t}B^*B.}
 \label{eq:V4-UV-excess-decay}
\end{equation}
Thus its refresh-weighted local Green tail obeys
\begin{equation}
 \boxed{
 \int_T^\infty e^{\rho t}K_B^{\rm par}(t)\,\dd t
 \preceq
 \frac{e^{5/16}}{\rho_F}
 e^{-\rho_Fe^{-5/16}T}B^*B.}
 \label{eq:V4-UV-weighted-tail}
\end{equation}
\end{corollary}

\begin{proof}
Equation \eqref{eq:V4-UV-gap} gives
$e^{-tH_{\rm UV}}\preceq
e^{-\rho_Fe^{-5/16}t}I$ on the centred space.  Add the scalar refresh decay
and apply \cref{cor:V4-excess-decay}.
\end{proof}

\begin{assessment}[What the ultraviolet gap closes]
\label{ass:V4-UV-scope}
The selected ultraviolet tail now has a same-support, conditional-heat-bath
Poincare inequality and a zero-mass source Green bound which survive complete
removal of global refresh.  This strengthens the existing mutual absolute
continuity and zero--one conclusion by adding a quantitative dynamical row.

The result is intentionally limited.  The tail blocks must still be identified
with physically local spacetime update regions before
\eqref{eq:V4-UV-gap} is called a spacetime-local clock.  More decisively, the
same selected tail is the branch on which the elementary Higgs coordinate
collapses or becomes Gaussian white noise under the imported scaling
trichotomy.  Its new Markov gap therefore controls an ultraviolet auxiliary
sector; it does not provide an interacting macroscopic Higgs field.
\end{assessment}

% -----------------------------------------------------------------------------
\subsection{Why the compact control gap is not yet the unbounded RPESM gap}
\label{subsec:V4-support-firewall}
% -----------------------------------------------------------------------------

The earlier field-shell theorem has a uniform constant
\[
 g_{48}=\min\{m_*,\rho_F(1-\alpha_*)\}>0
\]
on the compact Higgs carrier and the product field--particle law
$\mu_n^{\rm com}$.  The later RPESM measure in \textup{(RP.7)} instead has an
unbounded Higgs carrier and additional full-coupling, line-modulus, route,
scheduler, gauge-fixing, and counterterm factors.

\begin{proposition}[Support mismatch blocks direct density transfer]
\label{prop:V4-support-mismatch}
Let $\nu$ be a control law supported on a fixed compact Higgs ball
$\mathscr K_H^{V_n}$.  Let $\mu$ be an unbounded quartically controlled Higgs
law whose density is positive on a set of positive Lebesgue measure outside
that ball.  Then $\mu\not\ll\nu$.  In particular, no constants
$0<m\le M<\infty$ can satisfy
\[
 \dd\mu=w\,\dd\nu,\qquad m\le w\le M,
\]
and \cref{thm:V4-heatbath-comparison} cannot transport the compact-control gap
to $\mu$.
\end{proposition}

\begin{proof}
Let $A$ be the set of configurations with at least one Higgs coordinate
outside $\mathscr K_H$.  Then $\nu(A)=0$ by support, while $\mu(A)>0$ by the
assumed positive unbounded density.  This contradicts $\mu\ll\nu$, which is
necessary for a Radon--Nikodym window.
\end{proof}

\begin{assessment}[The actual local-gap ledger after V4]
\label{ass:V4-gap-ledger}
There are now three distinct statements.
\begin{enumerate}[label=\textup{(G\arabic*)},leftmargin=3em]
\item The compact protected field--reaction control slab has the imported gap
      $g_{48}$.
\item The selected ultraviolet product-tail perturbation has the new
      refresh-free gap $\rho_Fe^{-5/16}$.
\item The complete unbounded full-coupling RPESM local generator has no proved
      identification with either heat-bath form and no proved uniform local
      gap in the supplied corpus.
\end{enumerate}
The positive first two rows cannot be substituted for the third.  Likewise,
the covariance-walk lower kinetic coefficient in the selected Wilsonian
benchmark is a spatial quadratic-symbol statement, not a spectral gap of the
RPESM Markov generator.
\end{assessment}

% -----------------------------------------------------------------------------
\subsection{Revised continuum acquisition}
\label{subsec:V4-acquisition}
% -----------------------------------------------------------------------------

\begin{openproblem}[Executable V5 acquisition: the unbounded local clock]
\label{op:V4-next-acquisition}
Keep the common positive refresh rate only on the auxiliary parent packet and
use \cref{thm:V4-projective-derefresh} to work with the exact projective local
semigroups.  On one named unbounded RPESM source bank:
\begin{enumerate}[label=\textup{(A\arabic*)},leftmargin=3.2em]
\item specify the actual local conditional-update blocks, their physical
      diameters, rates, and common invariant law;
\item prove their incidence with the unbounded RPESM measure
      \textup{(RP.7)}, rather than the compact control law or a selected scalar
      benchmark;
\item acquire either a same-support density/Dobrushin window, a direct
      source-range factorization, or the refresh-weighted tail
      \eqref{eq:V4-uniform-weighted-tail};
\item test a current or stress source whose cutoff transport is already part
      of the RPESM event record;
\item retain separately the interacting infrared window, fixed-source
      non-Gaussian margin, determinant/Pfaffian line, and Palatini residual.
\end{enumerate}
A success gives a projective, refresh-free local Green response for that
physical source.  Failure returns a block-incidence defect, density-window
collapse, local zero mode, inverse-energy tail, or nonlocal update witness.
\end{openproblem}

\begin{theorem}[Version V4 terminal theorem]
\label{thm:V4-terminal}
Preserving V1--V3 verbatim before this appended block, V4 proves:
\begin{enumerate}[label=\textup{(V4.\arabic*)},leftmargin=5.2em]
\item common-rate parent projectivity descends exactly to the de-refreshed
      local semigroups and bounded generators;
\item simultaneous exact parent and local generator projectivity forces the
      refresh rate to be constant across cutoffs;
\item parent-as-physical jump locality, positive refresh, and both exact
      intertwinements are therefore incompatible on the V3 separated tests;
\item the local zero-mass Green form is exactly the exponentially weighted
      parent kernel integral and exists precisely under
      \eqref{eq:V4-zero-mass-domain};
\item a parent Green bound alone can hide both a local zero mode and a
      collapsing positive local gap;
\item a bounded same-support density perturbation transports a conditional
      heat-bath gap with explicit factor $(m/M)^2$;
\item the selected RPESM ultraviolet tail consequently has the refresh-free
      gap $\rho_Fe^{-5/16}$ and Green bound
      $e^{5/16}/\rho_F$;
\item support mismatch prevents direct promotion of the compact control gap to
      the complete unbounded full-coupling RPESM law.
\end{enumerate}
\end{theorem}

\begin{proof}
Items (V4.1)--(V4.3) are
\cref{thm:V4-projective-derefresh,thm:V4-rate-defect,cor:V4-rate-rigidity,thm:V4-projectivity-locality-incompatibility}.
Items (V4.4)--(V4.5) are
\cref{thm:V4-weighted-Green,thm:V4-cofinal-weighted-tail,ex:V4-refresh-hides-zero,ex:V4-collapsing-local-gap}.
Item (V4.6) is
\cref{thm:V4-heatbath-comparison,cor:V4-product-comparison}.
Item (V4.7) is
\cref{thm:V4-UV-local-gap,cor:V4-UV-excess-decay}.
The final item is \cref{prop:V4-support-mismatch}.
\end{proof}

\begin{assessment}[Programme boundary after V4]
\label{ass:V4-boundary}
V4 closes finite-time projective de-refreshing and one quantitative
refresh-free ultraviolet Green sector.  It also removes a false option:
one cannot preserve both exact generator intertwinements, start from a
positive refresh, and make the parent generator pass the V3 simultaneous-jump
locality test by silently sending its rate to zero.

The full Standard-Model--Einstein continuum remains unproved.  Its active
analytic frontier is narrower: identify and control the actual unbounded
local RPESM clock on physical current/stress sources, while separately closing
the interacting infrared, non-Gaussian, line, renormalization, Lorentzian, and
Palatini rows.  The ultraviolet gap proved here is evidence for that selected
tail only.
\end{assessment}

\section*{V4 exact checks and append-only record}
\addcontentsline{toc}{section}{V4 exact checks and append-only record}
The density constant used in \cref{thm:V4-UV-local-gap} is exactly
\[
 \left(\frac{e^{-5/64}}{e^{5/64}}\right)^2
 =e^{-20/64}=e^{-5/16}.
\]
The two-state controls give
\[
 \int_0^\infty e^{-\rho t}\,\dd t=\rho^{-1},
 \qquad
 \int_0^T e^{\rho t}e^{-\rho t}\,\dd t=T,
\]
and
\[
 \int_0^\infty e^{-(\rho+\lambda_n)t}\,\dd t
 =\frac1{\rho+\lambda_n},
 \qquad
 \int_0^\infty e^{\rho t}e^{-(\rho+\lambda_n)t}\,\dd t
 =\frac1{\lambda_n}.
\]
These scalar identities, the projective matrix identities, and the
conditional-density comparison are independently replayed in the structural
verification of this source.

The complete V3 source before removal of its final active document-closing
command is retained verbatim.  Its complete-file SHA--256 digest is
\begin{center}\small\ttfamily
3c76d8f3e538895d22b4bcc4b0ec75f3939d78147b9f1413ae78d4a53b1350e9.
\end{center}
A Version V5 must retain the complete V4 source, remove only the final active
document-closing command, append new material, restore one final command, and
use a filename ending in \texttt{\_V5.tex}.

% =============================================================================
% END APPENDED VERSION V4 MATERIAL
% =============================================================================


% =============================================================================
% BEGIN APPENDED VERSION V5 MATERIAL
% =============================================================================

\clearpage
\section{Version V5: refresh-free kinetic incidence and the invariant-score deficit}
\label{sec:V5-local-incidence}

Version V4 isolated the unbounded local-clock bottleneck: the global refresh
can be removed projectively, but the supplied regulator does not yet identify
a quantitatively controlled physical local generator.  The present version
moves one step upstream.  It does not assume a local spectral gap.  Instead it
starts from the actual law-dependent conditional heat baths and computes the
source created when the invariant law is varied.

The attached monograph already proves, on its finite representative, the
refreshed identity
\[
 \mathcal J_\rho=\tfrac12H_\rho W,
 \qquad
 4\mathcal J_\rho^*H_\rho^{-2}\mathcal J_\rho=W^*W
 \tag{HIT.10 revisited}
\]
for the score synthesis $W$.  That identity uses
$H_\rho=H_0+\rho I$ on the centred space.  It therefore cannot by itself
decide whether the refresh-free heat baths see the same score directions.
The new point below is an exact local-only decomposition:
\[
 W^*W
 =
 4(H_0^\dagger\mathcal J_0)^*
   (H_0^\dagger\mathcal J_0)
 +W^*P_{\rm inv}W.
\]
The second term is the Fisher mass of the score component invariant under
every declared local update.  It is zero for a full-coordinate Gibbs sampler
under a measure equivalent to a product reference, but need not be zero for
the unspecified local generator in \textup{(RP.8)}.

% -----------------------------------------------------------------------------
\subsection{Square-root heat baths and the local tangent column}
\label{subsec:V5-heatbath-setup}
% -----------------------------------------------------------------------------

Let $(\Omega,\mathcal F,\lambda)$ be a $\sigma$-finite measured space and let
$\mu_t$ be a probability law with strictly positive density with respect to
$\lambda$ for $t$ near zero.  Write
\[
 s_t=\sqrt{\frac{\dd\mu_t}{\dd\lambda}},
 \qquad
 R_t=|s_t\rangle\langle s_t|,
 \qquad
 Q_t=I-R_t.
\]
For one command direction assume Hellinger differentiability
\begin{equation}
 \frac{s_t-s}{t}\longrightarrow-\frac12b
 \quad\hbox{in }L^2(\lambda),
 \qquad
 b=M_ws,\qquad
 w\in L^2_0(\mu).
 \label{eq:V5-Hellinger-tangent}
\end{equation}
An exponential deformation
$\dd\mu_t=Z_t^{-1}e^{-tv}\dd\mu$ has this form with
$w=v-\mu(v)$ whenever its square-root path is differentiable in
$L^2(\lambda)$.

Let $E_{j,t}$ be conditional expectation for a finite family of update blocks,
and let
\[
 \mathcal U_t:L^2(\mu_t)\longrightarrow L^2(\lambda),
 \qquad
 \mathcal U_tf=s_tf,
 \qquad
 P_{j,t}=\mathcal U_tE_{j,t}\mathcal U_t^*.
\]
Each $P_{j,t}$ is an orthogonal projection and $P_{j,t}s_t=s_t$.  Fix rates
$c_j>0$ and put
\[
 H_{0,t}:=\sum_{j=1}^Jc_j(I-P_{j,t}),
 \qquad
 H_{\rho,t}:=H_{0,t}+\rho(I-R_t).
\]
Assume only that $H_{0,t}\to H_0$ strongly at zero.  This is automatic on a
finite carrier and is the precise continuity row needed below on a general
carrier.  All these operators are bounded by $2\sum_jc_j+\rho$.

For a finite-dimensional real command space $U$ with basis $(e_a)$, use one
path \eqref{eq:V5-Hellinger-tangent} per basis vector and define
\[
 We_a=b_a=M_{w_a}s,
 \qquad
 M:=W^*W,
 \qquad
 M_{ab}=\mathbb E_\mu(w_aw_b).
 \label{eq:V5-score-bank}
\]
Thus $W:U\to\mathcal K:=s^\perp$ is the square-root score synthesis and $M$ is
its Fisher Gram.

\begin{theorem}[Actual refresh-free heat-bath incidence]
\label{thm:V5-local-tangent-factor}
The local and refreshed tangent columns, evaluated on the fixed stationary
vector,
\[
 \mathcal J_0e_a
 :=\lim_{t\to0}\frac{(H_{0,t}-H_0)s}{t},
 \qquad
 \mathcal J_\rho e_a
 :=\lim_{t\to0}\frac{(H_{\rho,t}-H_{\rho,0})s}{t},
\]
exist and satisfy
\begin{equation}
 \boxed{
 \mathcal J_0=\frac12H_0W,
 \qquad
 \mathcal J_\rho=\frac12(H_0+\rho I_{\mathcal K})W,
 \qquad
 \mathcal J_\rho-\mathcal J_0=\frac\rho2W.}
 \label{eq:V5-local-parent-incidence}
\end{equation}
In particular, the actual local source has the one-generator factorization
required by \cref{thm:V3-source-form-derefresh}, with $C=W/2$.  No inverse and
no spectral gap are used to obtain this factorization.
\end{theorem}

\begin{proof}
The stationarity identities
$H_{0,t}s_t=0$ and $H_0s=0$ give
\[
 \frac{(H_{0,t}-H_0)s}{t}
 =-H_{0,t}\frac{s_t-s}{t}.
\]
The operators $H_{0,t}$ are uniformly bounded, strong convergence holds by
hypothesis, and \eqref{eq:V5-Hellinger-tangent} converges in norm.  The right
side therefore converges to $H_0b/2$.  This proves the first identity for each
command column.

Likewise $H_{\rho,t}s_t=0$.  On $s^\perp$ the operator
$H_{\rho,0}$ is $H_0+\rho I_{\mathcal K}$, so the same stationary-line
argument gives the second identity.  Subtraction proves the third.
Equivalently, differentiating the refresh directly gives
$-\rho\dot R\,s=\rho b/2$ because $b\perp s$.
\end{proof}

\begin{remark}[Why the strong formulation matters]
\label{rem:V5-strong-formulation}
The proof differentiates the generator only on its stationary vector.  It
does not require an operator-norm derivative of every conditional projection.
Consequently an unbounded polynomial score is admissible as soon as its
Hellinger tangent lies in $L^2$ and the heat baths converge strongly.  This is
the natural level for the unbounded Higgs carrier in \textup{(RP.7)}.
\end{remark}
% -----------------------------------------------------------------------------
\subsection{The invariant-score Fisher decomposition}
\label{subsec:V5-Fisher-split}
% -----------------------------------------------------------------------------

Restrict $H_0$ to the centred Hilbert space $\mathcal K=s^\perp$.  It is
bounded, nonnegative, and self-adjoint.  Let
\[
 P_{\rm inv}:=P_{\Ker H_0},
 \qquad
 P_{\rm mov}:=I_{\mathcal K}-P_{\rm inv},
\]
and let $H_0^\dagger$ be the supported Moore--Penrose inverse.  Its domain
contains $\Ran H_0$, and spectral calculus gives
\[
 H_0^\dagger H_0=P_{\rm mov}.
 \label{eq:V5-supported-projector}
\]

\begin{theorem}[Exact local Fisher Pythagoras]
\label{thm:V5-local-Fisher-Pythagoras}
For the source in \cref{thm:V5-local-tangent-factor}, define
\begin{align}
 M_{\rm mov}
 &:=4(H_0^\dagger\mathcal J_0)^*
          (H_0^\dagger\mathcal J_0),\\
 M_{\rm inv}
 &:=W^*P_{\rm inv}W.
\end{align}
Then
\begin{equation}
 \boxed{
 H_0^\dagger\mathcal J_0=\frac12P_{\rm mov}W,
 \qquad
 M_{\rm mov}=W^*P_{\rm mov}W,
 \qquad
 M=M_{\rm mov}+M_{\rm inv}.}
 \label{eq:V5-Fisher-Pythagoras}
\end{equation}
Both summands are positive semidefinite.  Moreover
\begin{equation}
 \boxed{
 \mathcal J_0^*H_0^\dagger\mathcal J_0
 =\frac14W^*H_0W.}
 \label{eq:V5-local-Green-Dirichlet}
\end{equation}
For the refreshed tangent,
\begin{equation}
 \boxed{
 (H_0+\rho I)^{-1}\mathcal J_\rho=\frac12W,
 \qquad
 4\bigl((H_0+\rho I)^{-1}\mathcal J_\rho\bigr)^*
   \bigl((H_0+\rho I)^{-1}\mathcal J_\rho\bigr)=M.}
 \label{eq:V5-parent-full-Fisher}
\end{equation}
Thus the refreshed identity always returns the full Fisher matrix, while the
local identity returns it if and only if
\begin{equation}
 \boxed{P_{\rm inv}W=0.}
 \label{eq:V5-invariant-score-test}
\end{equation}
On an invariant score component the contrast is exact:
\[
 P_{\rm inv}\mathcal J_0=0,
 \qquad
 P_{\rm inv}\mathcal J_\rho=\frac\rho2P_{\rm inv}W.
 \label{eq:V5-refresh-created-incidence}
\]
\end{theorem}

\begin{proof}
Insert $\mathcal J_0=H_0W/2$ into
\eqref{eq:V5-supported-projector}.  This gives
$H_0^\dagger\mathcal J_0=P_{\rm mov}W/2$.  Taking its Gram gives
$M_{\rm mov}=W^*P_{\rm mov}W$.  Since
$P_{\rm mov}+P_{\rm inv}=I_{\mathcal K}$, the Fisher Pythagoras follows.
Likewise,
\[
 \mathcal J_0^*H_0^\dagger\mathcal J_0
 =\frac14W^*H_0P_{\rm mov}W
 =\frac14W^*H_0W.
\]
The refreshed operator is strictly positive on $\mathcal K$, so
\eqref{eq:V5-parent-full-Fisher} follows immediately from the second identity
in \eqref{eq:V5-local-parent-incidence}.  Finally $H_0P_{\rm inv}=0$, and
projecting \eqref{eq:V5-local-parent-incidence} onto $\Ker H_0$ gives
\eqref{eq:V5-refresh-created-incidence}.
\end{proof}

\begin{corollary}[Rank and identifiability]
\label{cor:V5-local-rank}
One has
\[
 \Ker\mathcal J_0=\Ker M_{\rm mov}
 =\{u\in U:Wu\in\Ker H_0\}.
\]
Consequently
\[
 \Ker\mathcal J_0=\Ker M
 \quad\Longleftrightarrow\quad
 \Ran W\cap\Ker H_0=\{0\}.
\]
The stronger condition \eqref{eq:V5-invariant-score-test} identifies the
entire local response with $W/2$, not merely its rank.
\end{corollary}

\begin{proof}
By \eqref{eq:V5-local-parent-incidence},
$\mathcal J_0u=0$ precisely when $Wu\in\Ker H_0$.  A Gram and its synthesis
have the same kernel.  Since $\Ker M=\Ker W$, the second statement follows.
The final statement is the first identity in
\eqref{eq:V5-Fisher-Pythagoras}.
\end{proof}

% -----------------------------------------------------------------------------
\subsection{Conditional variances and the exact invariant field}
\label{subsec:V5-conditional-form}
% -----------------------------------------------------------------------------

Return to the heat-bath representation.  At $t=0$ let $E_j$ be conditional
expectation under $\mu$ and let
$\mathcal U:L^2(\mu)\to L^2(\lambda)$ be multiplication by $s$.  Then
$P_j=\mathcal UE_j\mathcal U^*$ and
\begin{equation}
 (I-P_j)M_ws=\mathcal U(w-E_jw).
 \label{eq:V5-score-residual}
\end{equation}
Define the closed common-invariant subspace
\[
 \mathcal I
 :=\bigcap_{j=1}^J\Ran E_j
 =\{f\in L^2(\mu):E_jf=f\text{ for every }j\},
\]
and let $\Pi_{\mathcal I}$ be its orthogonal projection.

\begin{theorem}[Heat-bath kernel and conditional-score formula]
\label{thm:V5-heatbath-kernel}
For positive rates $c_j$,
\begin{equation}
 \boxed{
 \Ker H_0
 =\mathcal U(\mathcal I\cap1^\perp)
 \quad\hbox{on }\mathcal K.}
 \label{eq:V5-kernel-intersection}
\end{equation}
For the centred command scores,
\begin{align}
 (M_{\rm inv})_{ab}
 &=\left\langle\Pi_{\mathcal I}w_a,
                    \Pi_{\mathcal I}w_b\right\rangle_{L^2(\mu)},
 \label{eq:V5-invariant-Fisher}\\
 \bigl(\mathcal J_0^*H_0^\dagger\mathcal J_0\bigr)_{ab}
 &=\frac14\sum_{j=1}^Jc_j\,
   \mathbb E_\mu\!\left[
    (w_a-E_jw_a)(w_b-E_jw_b)\right].
 \label{eq:V5-conditional-Green}
\end{align}
Hence the local Fisher deficit is exactly the covariance retained by every
update, while the local Green form is exactly one quarter of the
rate-weighted conditional-score Dirichlet matrix.
\end{theorem}

\begin{proof}
For $f\in\mathcal K$,
\[
 \langle f,H_0f\rangle
 =\sum_jc_j\|(I-P_j)f\|^2.
\]
Every rate is positive, so this vanishes precisely when
$f\in\bigcap_j\Ran P_j=\mathcal U\mathcal I$.  Restriction to $\mathcal K$
removes the stationary line and proves \eqref{eq:V5-kernel-intersection}.
Since every $w_a$ is centred, its projection onto $\mathcal I$ is also
centred.  Conjugating the projection in
\eqref{eq:V5-Fisher-Pythagoras} proves
\eqref{eq:V5-invariant-Fisher}.  Finally
\eqref{eq:V5-score-residual} and
\eqref{eq:V5-local-Green-Dirichlet} give
\[
 \langle b_a,H_0b_b\rangle
 =\sum_jc_j
 \langle w_a-E_jw_a,w_b-E_jw_b\rangle_{L^2(\mu)},
\]
which is \eqref{eq:V5-conditional-Green}.
\end{proof}

\begin{theorem}[Full-coordinate positivity removes exact hidden scores]
\label{thm:V5-product-ergodicity}
Suppose
\[
 \Omega=\prod_{j=1}^J\Omega_j,
 \qquad
 \lambda=\bigotimes_{j=1}^J\lambda_j,
 \qquad
 \mu\sim\lambda,
 \label{eq:V5-product-equivalence}
\]
and the update family contains, with positive rate, the heat bath of every
single coordinate $j$ conditional on all other coordinates.  Then
$\mathcal I=\mathbb R1$ and
\begin{equation}
 \boxed{
 \Ker H_0=\{0\}\text{ on }\mathcal K,
 \qquad
 M_{\rm inv}=0,
 \qquad
 M_{\rm mov}=M.}
 \label{eq:V5-full-coordinate-conclusion}
\end{equation}
Only equivalence in \eqref{eq:V5-product-equivalence} is used; no bounded
density ratio, Dobrushin estimate, or spectral-gap lower bound is required.
\end{theorem}

\begin{proof}
Let $f\in\mathcal I$.  For each $j$, the equality $E_jf=f$ says that $f$ has
a version independent of coordinate $j$.  Equivalence of $\mu$ and $\lambda$
identifies their null sets, so these coordinate-independence statements also
hold $\lambda$-almost everywhere.

Take independent product points $x,y$.  Replace the coordinates of $x$ by
those of $y$ one at a time.  At the $j$th replacement the value of $f$ is
unchanged almost surely because $f$ is independent of that coordinate.
Fubini's theorem makes the finite chain simultaneous, and hence
$f(x)=f(y)$ for $\lambda\otimes\lambda$-almost every $(x,y)$.  Therefore $f$
is constant $\lambda$-almost everywhere and, by equivalence, $\mu$-almost
everywhere.  Thus $\mathcal I=\mathbb R1$.  The remaining conclusions follow
from \cref{thm:V5-heatbath-kernel,thm:V5-local-Fisher-Pythagoras}.
\end{proof}

\begin{remark}[What this theorem does not quantify]
\label{rem:V5-no-gap}
At each finite cutoff, exact positivity and a complete block list can eliminate
extra zero modes even when the smallest positive eigenvalue is arbitrarily
small.  Therefore \cref{thm:V5-product-ergodicity} is an identifiability
theorem, not a uniform mixing theorem.  Approximate invariants can still
develop along a cofinal sequence.
\end{remark}
% -----------------------------------------------------------------------------
\subsection{Gapless Green integration and quantitative de-refreshing}
\label{subsec:V5-gapless-Green}
% -----------------------------------------------------------------------------

The local incidence factor provides more than a kernel test.  It cancels one
power of every small local eigenvalue before the Green inverse is taken.

\begin{theorem}[Exact gapless Green primitive for the actual local source]
\label{thm:V5-gapless-Green}
For every $T>0$,
\begin{align}
 \int_0^T e^{-tH_0}\mathcal J_0\,\dd t
 &=\frac12(I-e^{-TH_0})W,
 \label{eq:V5-vector-primitive}\\
 \int_0^T\mathcal J_0^*e^{-tH_0}\mathcal J_0\,\dd t
 &=\frac14W^*H_0(I-e^{-TH_0})W.
 \label{eq:V5-form-primitive}
\end{align}
As $T\to\infty$, the first identity converges strongly and the second in
operator norm on the finite-dimensional command space:
\begin{align}
 \boxed{
 \lim_{T\to\infty}\int_0^T e^{-tH_0}\mathcal J_0\,\dd t
 =\frac12P_{\rm mov}W,}
 \label{eq:V5-vector-limit}\\
 \boxed{
 \int_0^\infty
 \mathcal J_0^*e^{-tH_0}\mathcal J_0\,\dd t
 =\frac14W^*H_0W.}
 \label{eq:V5-form-limit}
\end{align}
These improper integrals exist without a positive gap.
\end{theorem}

\begin{proof}
Because $H_0$ is bounded and commutes with its semigroup,
\[
 \int_0^T e^{-tH_0}H_0\,\dd t=I-e^{-TH_0}.
\]
Substitute $\mathcal J_0=H_0W/2$ to obtain
\eqref{eq:V5-vector-primitive}.  Substitution on both sides of the semigroup
gives
\[
 \frac14W^*\!\left(\int_0^TH_0^2e^{-tH_0}\,\dd t\right)W
 =\frac14W^*H_0(I-e^{-TH_0})W,
\]
which is \eqref{eq:V5-form-primitive}.  Spectral calculus gives
$e^{-TH_0}\to P_{\rm inv}$ strongly.  This proves the vector limit.  The
command space is finite dimensional, so entrywise convergence of the bounded
positive matrices in \eqref{eq:V5-form-primitive} is norm convergence and
gives \eqref{eq:V5-form-limit}.
\end{proof}

For comparison with the auxiliary refreshed clock, keep the \emph{same local
source} $\mathcal J_0$ and define
\[
 G_0:=\mathcal J_0^*H_0^\dagger\mathcal J_0,
 \qquad
 G_\rho^{[0]}:=
 \mathcal J_0^*(H_0+\rho I)^{-1}\mathcal J_0.
\]
The bracket in $G_\rho^{[0]}$ records that the source is $\mathcal J_0$, not
the different refreshed tangent $\mathcal J_\rho$.

\begin{theorem}[Fisher-controlled de-refreshing error]
\label{thm:V5-Fisher-derefresh}
For every $\rho>0$,
\begin{align}
 G_\rho^{[0]}
 &=\frac14W^*H_0(H_0+\rho I)^{-1}H_0W,
 \label{eq:V5-refreshed-local-source}\\
 \boxed{
 0\preceq G_0-G_\rho^{[0]}
 =\frac\rho4W^*H_0(H_0+\rho I)^{-1}W
 \preceq\frac\rho4M.}
 \label{eq:V5-Fisher-error}
\end{align}
The vector-response difference is
\begin{equation}
 \boxed{
 H_0^\dagger\mathcal J_0
 -(H_0+\rho I)^{-1}\mathcal J_0
 =\frac\rho2(H_0+\rho I)^{-1}P_{\rm mov}W.}
 \label{eq:V5-vector-error}
\end{equation}
It tends strongly to zero for a fixed cutoff as $\rho\downarrow0$, but no
cutoff-uniform operator-norm rate follows without further low-spectrum
control.
\end{theorem}

\begin{proof}
The first formula follows by inserting $\mathcal J_0=H_0W/2$.  Spectral
calculus gives
\[
 H_0-H_0(H_0+\rho I)^{-1}H_0
 =\rho H_0(H_0+\rho I)^{-1},
\]
and the middle factor lies between zero and the identity.  This proves
\eqref{eq:V5-Fisher-error}.  On the positive spectral subspace,
\[
 1-\frac{\lambda}{\lambda+\rho}=\frac{\rho}{\lambda+\rho};
\]
on the kernel both sides of \eqref{eq:V5-vector-error} vanish.  The displayed
identity and its strong limit follow.
\end{proof}

\begin{corollary}[Cofinal quadratic de-refreshing without local gaps]
\label{cor:V5-cofinal-derefresh}
Let $(H_{0,n},W_n)$ be a cutoff family on a fixed finite command space and
suppose
\[
 W_n^*W_n\preceq C_MI
 \quad\hbox{for every }n.
 \label{eq:V5-uniform-Fisher}
\]
For any $\rho_n\downarrow0$,
\[
 \boxed{
 \|G_{0,n}-G_{\rho_n,n}^{[0]}\|
 \le\frac{C_M}{4}\rho_n\longrightarrow0.}
 \label{eq:V5-cofinal-Green-error}
\]
No lower bound on the positive spectrum of $H_{0,n}$ is needed.
\end{corollary}

\begin{proof}
Apply \eqref{eq:V5-Fisher-error} at each cutoff and take the operator norm on
the finite command space.
\end{proof}

\begin{assessment}[Quadratic response closes before vector response]
\label{ass:V5-form-vector-fork}
The cofinal Green form in \eqref{eq:V5-cofinal-Green-error} is controlled by
the Fisher Gram because the actual local tangent carries one factor of
$H_{0,n}$.  The vector response in \eqref{eq:V5-vector-error} has no analogous
uniform conclusion when positive eigenvalues collapse.  Any continuum claim
that needs the full response vector, rather than its source Gram, must still
supply a source-weighted low-spectrum condition or a local gap.
\end{assessment}

% -----------------------------------------------------------------------------
\subsection{Sharp missing-block separator}
\label{subsec:V5-missing-block}
% -----------------------------------------------------------------------------

\begin{counterexample}[Refresh manufactures an unseen command direction]
\label{ex:V5-missing-coordinate}
Let $\Omega=\{-1,1\}^2$ with the uniform law, use the orthonormal centred
scores
\[
 w_1(x)=x_1,\qquad w_2(x)=x_2,
\]
and update only coordinate $1$ at rate $c>0$.  Thus
$E_1f=\mathbb E[f\mid x_2]$ and $H_0=c(I-E_1)$ on the centred space.  In the
command basis $(e_1,e_2)$,
\begin{align}
 M&=I_2,
 &
 \mathcal J_0&=\frac12(cw_1,0),
 \label{eq:V5-two-bit-source}\\
 M_{\rm mov}&=
 \begin{pmatrix}1&0\\0&0\end{pmatrix},
 &
 M_{\rm inv}&=
 \begin{pmatrix}0&0\\0&1\end{pmatrix},
 &
 G_0&=\frac14
 \begin{pmatrix}c&0\\0&0\end{pmatrix}.
 \label{eq:V5-two-bit-split}
\end{align}
By contrast,
\[
 \mathcal J_\rho
 =\frac12\bigl((c+\rho)w_1,\rho w_2\bigr)
\]
and its refreshed inverse reconstructs $M=I_2$.  If the coordinate-$2$ heat
bath is added at rate $d>0$, then the centred kernel becomes zero and
\[
 M_{\rm mov}=I_2,\qquad
 G_0=\frac14\begin{pmatrix}c&0\\0&d\end{pmatrix}.
 \label{eq:V5-two-bit-complete}
\]
The limit $d\downarrow0$ retains exact finite-$d$ identifiability while losing
a uniform gap.
\end{counterexample}

\begin{proof}
Conditional averaging gives
$E_1w_1=0$ and $E_1w_2=w_2$, hence
$H_0w_1=cw_1$ and $H_0w_2=0$.  All matrices in
\eqref{eq:V5-two-bit-source}--\eqref{eq:V5-two-bit-split} follow from
\cref{thm:V5-local-Fisher-Pythagoras}.  With refresh,
$H_0+\rho I$ has eigenvalues $c+\rho$ and $\rho$ on these two score
directions, proving the displayed parent source and its exact reconstruction.
The second heat bath satisfies $H_0w_2=dw_2$ and leaves the first eigenvalue
$c$, which proves \eqref{eq:V5-two-bit-complete}.
\end{proof}

This four-state separator shows why the refreshed Fisher identity cannot be
used as evidence that the chosen local blocks see every physical command.  It
also gives a complete finite test: compute $M_{\rm inv}$ or, equivalently,
the common-invariant conditional expectation of every score.

% -----------------------------------------------------------------------------
\subsection{Projective source transport has two independent defects}
\label{subsec:V5-projective-source}
% -----------------------------------------------------------------------------

The factorization at each cutoff does not by itself make the source
projective.  The score bank and the local generator must both transport.

\begin{proposition}[Exact projective local-incidence defect]
\label{prop:V5-incidence-defect}
Let $\iota:\mathcal K_n\to\mathcal K_m$ be an isometry between two centred
cutoff spaces, let $H_n,H_m$ be their local nonnegative generators, and let
$W_n,W_m:U\to\mathcal K_n,\mathcal K_m$ be score syntheses on the same command
space.  Put $\mathcal J_k=H_kW_k/2$.  Then
\begin{equation}
 \boxed{
 2(\mathcal J_m-\iota\mathcal J_n)
 =H_m(W_m-\iota W_n)
 +(H_m\iota-\iota H_n)W_n.}
 \label{eq:V5-projective-source-defect}
\end{equation}
Consequently,
\begin{equation}
 \|\mathcal J_m-\iota\mathcal J_n\|
 \le\frac12\|H_m\|\,\|W_m-\iota W_n\|
 +\frac12\|(H_m\iota-\iota H_n)W_n\|.
 \label{eq:V5-projective-source-bound}
\end{equation}
Exact score transport and exact local-generator intertwining imply exact
transport of the actual local tangent.
\end{proposition}

\begin{proof}
Add and subtract $H_m\iota W_n/2$ in
$H_mW_m/2-\iota H_nW_n/2$.  The norm bound is the triangle inequality.
\end{proof}

\begin{remark}[No converse from source transport alone]
\label{rem:V5-no-converse}
The two terms on the right of
\eqref{eq:V5-projective-source-defect} can cancel on a selected command bank.
Therefore exact transport of $\mathcal J_n$ alone does not certify either
score transport or generator projectivity.  Both defects must be retained in
a cofinal audit.
\end{remark}
% -----------------------------------------------------------------------------
\subsection{Application to the unbounded RPESM regulator}
\label{subsec:V5-RPESM-application}
% -----------------------------------------------------------------------------

The measure in \textup{(RP.7)} has the form
\[
 \dd\mu_n
 =Z_n^{-1}e^{-S_{B,n}}\|\sigma_n^+\|^2|p_n^+|^2\,\dd\lambda_n
\]
on a finite-cutoff carrier with unbounded Higgs coordinates.  Its moment
theorem states that fixed-support polynomial action and score writers are
integrable beyond first order.  Applied to the square of such a polynomial,
that assertion supplies the $L^2(\mu_n)$ score condition used in
\eqref{eq:V5-score-bank}.  Along any coefficient interval on which the
confining quartic coefficient remains positive, the corresponding
exponential action deformation supplies the Hellinger paths required above.

\begin{proposition}[Conditional RPESM local-score closure]
\label{prop:V5-RPESM-score-closure}
Fix a cutoff $n$ and a finite bank of actual action-deformation scores for
\textup{(RP.7)}.  Suppose:
\begin{enumerate}[label=\textup{(R\arabic*)},leftmargin=3.2em]
\item the reference $\lambda_n$ is written as a product over all active
      finite-cutoff coordinates and $\mu_n\sim\lambda_n$;
\item $L_n^{\rm loc}$ is chosen to contain every singleton-coordinate
      $\mu_n$-conditional heat bath with positive rate;
\item the conditional projections vary strongly along each declared action
      deformation.
\end{enumerate}
Let $H_{0,n}=-L_n^{\rm loc}$ on centred square-root coordinates and let $W_n$
be the score synthesis.  Then
\begin{equation}
 \boxed{
 \mathcal J_{0,n}=\frac12H_{0,n}W_n,\qquad
 H_{0,n}^\dagger\mathcal J_{0,n}=\frac12W_n,}
 \label{eq:V5-RPESM-local-response}
\end{equation}
and
\begin{equation}
 \boxed{
 4(H_{0,n}^\dagger\mathcal J_{0,n})^*
   (H_{0,n}^\dagger\mathcal J_{0,n})=M_n.}
 \label{eq:V5-RPESM-local-Fisher}
\end{equation}
The local Green matrix is the explicit conditional-score form
\begin{equation}
 \boxed{
 (G_{0,n})_{ab}
 =\frac14\sum_jc_{j,n}\,
  \mathbb E_{\mu_n}\!\left[
   (w_{a,n}-E_{j,n}w_{a,n})
   (w_{b,n}-E_{j,n}w_{b,n})\right].}
 \label{eq:V5-RPESM-conditional-Green}
\end{equation}
If the Fisher matrices are uniformly bounded, the quadratic replacement of
this local Green matrix by the same-source refreshed Green matrix has the
uniform error \eqref{eq:V5-cofinal-Green-error}.
\end{proposition}

\begin{proof}
The moment and deformation hypotheses make
\cref{thm:V5-local-tangent-factor} applicable.  Assumptions \textup{(R1)} and
\textup{(R2)} give \eqref{eq:V5-full-coordinate-conclusion}; hence
$P_{\rm inv}W_n=0$.  The response and Fisher identities follow from
\cref{thm:V5-local-Fisher-Pythagoras}.  The conditional formula is
\eqref{eq:V5-conditional-Green}, and the final assertion is
\cref{cor:V5-cofinal-derefresh}.
\end{proof}

\begin{assessment}[What is now proved and what \textup{(RP.8)} still omits]
\label{ass:V5-RPESM-ledger}
The proposition is constructive at one cutoff: choose all coordinate heat
baths and compute the score residuals in
\eqref{eq:V5-RPESM-conditional-Green}.  It applies to an unbounded Higgs
carrier and needs no compact-support density comparison.  It proves that the
actual local action tangent has a finite, refresh-free Green form even if the
local gap is small.

The supplied statement \textup{(RP.8)}, however, says only that
$L_n^{\rm loc}$ is a finite sum of local conditional-expectation generators.
It does not list the blocks and rates, assert product equivalence of the
positive-modulus cylinder, or prove that the full-coordinate choices
intertwine under the declared cutoff maps.  Therefore
\cref{prop:V5-RPESM-score-closure} is an exact acquisition theorem, not an
unconditional certification of the generator already named in
\textup{(RP.8)}.

A current or stress writer enters this theorem only when it is the score of an
actual differentiable gauge, coframe, or metric deformation of the same
cylinder.  Merely recording a current or stress observable does not establish
that incidence.  Its cofinal passage must additionally drive both terms in
\eqref{eq:V5-projective-source-defect} to zero.
\end{assessment}

% -----------------------------------------------------------------------------
\subsection{Revised acquisition after the local-score result}
\label{subsec:V5-acquisition}
% -----------------------------------------------------------------------------

\begin{openproblem}[Executable V6 acquisition: physical block and stress incidence]
\label{op:V5-next-acquisition}
For each cutoff, enumerate the actual coordinates in \textup{(RP.7)} and
declare a positive-rate conditional block cover.  Then:
\begin{enumerate}[label=\textup{(A\arabic*)},leftmargin=3.2em]
\item verify pointwise positivity, or record the zero set, of
      $\|\sigma_n^+\|^2|p_n^+|^2$ relative to the product reference;
\item compute the invariant-score matrix $M_{{\rm inv},n}$ for the gauge,
      Higgs-quartic, hopping, current, and stress command bank;
\item realize current and stress as derivatives of explicit gauge/coframe
      deformations of the same unnormalized cylinder;
\item evaluate the conditional-score Green matrix
      \eqref{eq:V5-RPESM-conditional-Green} and obtain a uniform Fisher bound;
\item evaluate separately the score-transport and generator-intertwining
      defects in \eqref{eq:V5-projective-source-defect};
\item retain the vector low-spectrum tail whenever the continuum observable
      requires more than the quadratic Green form.
\end{enumerate}
A zero invariant-score matrix and vanishing projective defects give a
refresh-free cofinal Green form for that physical command bank without a
uniform local gap.  Failure produces a named zero-set, missing-block,
incidence, score-transport, generator-transport, or low-spectrum witness.
\end{openproblem}

\begin{theorem}[Version V5 terminal theorem]
\label{thm:V5-terminal}
Preserving V1--V4 verbatim before this appended block, V5 proves:
\begin{enumerate}[label=\textup{(V5.\arabic*)},leftmargin=5.2em]
\item the actual law-dependent local heat-bath tangent factors exactly as
      $\mathcal J_0=H_0W/2$ under Hellinger and strong heat-bath continuity;
\item the Fisher matrix splits orthogonally into locally movable and
      common-invariant score Grams;
\item the refreshed tangent fills every invariant score direction by the
      explicit source $\rho P_{\rm inv}W/2$;
\item complete singleton heat baths under a law equivalent to a product
      reference have no nonconstant exact invariant;
\item the actual local source has exact vector and quadratic Green primitives
      without a local spectral gap;
\item the same-source refreshed Green error is bounded by $\rho M/4$, uniformly
      along any cutoff family with a uniform Fisher bound;
\item the four-state missing-coordinate model sharply separates refreshed
      Fisher reconstruction from local identifiability;
\item projective transport of the local tangent splits exactly into a score
      defect and a local-generator defect;
\item the unbounded RPESM action-score application is reduced to three
      checkable rows: product equivalence, a complete physical block cover,
      and strong deformation continuity.
\end{enumerate}
\end{theorem}

\begin{proof}
Items \textup{(V5.1)}--\textup{(V5.3)} are
\cref{thm:V5-local-tangent-factor,thm:V5-local-Fisher-Pythagoras}.
Item \textup{(V5.4)} is \cref{thm:V5-product-ergodicity}.
Items \textup{(V5.5)}--\textup{(V5.6)} are
\cref{thm:V5-gapless-Green,thm:V5-Fisher-derefresh,cor:V5-cofinal-derefresh}.
Item \textup{(V5.7)} is \cref{ex:V5-missing-coordinate}.
Item \textup{(V5.8)} is \cref{prop:V5-incidence-defect}.
The final item is \cref{prop:V5-RPESM-score-closure,ass:V5-RPESM-ledger}.
\end{proof}

\begin{assessment}[Programme boundary after V5]
\label{ass:V5-boundary}
V5 closes a source-relative part of the unbounded local-clock problem.  For
actual action deformations, the local tangent automatically has the factor
needed to remove refresh from its quadratic Green form, and the only exact
Fisher loss is the explicitly computable common-invariant score component.
This avoids substituting the compact control gap for the unbounded law and
avoids treating the refreshed Fisher identity as evidence of local
identifiability.

The full Standard--Model--Einstein continuum remains unproved.  The next
finite acquisition must instantiate the block cover and the physical
current/stress deformations on \textup{(RP.7)}, then prove their projective
transport.  Interacting infrared control, fixed-source non-Gaussianity,
determinant/Pfaffian transport, renormalization, Lorentzian reconstruction, and
the Palatini residual remain separate continuum rows.
\end{assessment}

\section*{V5 exact checks and append-only record}
\addcontentsline{toc}{section}{V5 exact checks and append-only record}
For the missing-coordinate separator, the centred Walsh basis
$(w_1,w_2,w_1w_2)$ gives
\[
 H_0=\operatorname{diag}(c,0,c),
 \qquad
 H_0+\rho I=\operatorname{diag}(c+\rho,\rho,c+\rho).
\]
The two score columns therefore reproduce
\eqref{eq:V5-two-bit-source}--\eqref{eq:V5-two-bit-complete} by direct matrix
multiplication.  For a scalar spectral value $\lambda\ge0$, the two identities
behind the gapless and de-refreshing calculations are
\[
 \int_0^T\lambda e^{-t\lambda}\,\dd t
 =1-e^{-T\lambda},
 \qquad
 \lambda-\frac{\lambda^2}{\lambda+\rho}
 =\frac{\rho\lambda}{\lambda+\rho}.
\]
These scalar identities and the four-state matrices are independently
replayed in the structural verification of this source.

The complete V4 source before removal of its final active document-closing
command is retained verbatim.  Its complete-file SHA--256 digest is
\begin{center}\small\ttfamily
b7b42803684fbfa0fe9bf429225ab7de4964a246faddf7c0acc3f68099ad2869.
\end{center}
A Version V6 must retain the complete V5 source, remove only the final active
document-closing command, append new material, restore one final command, and
use a filename ending in \texttt{\_V6.tex}.

% =============================================================================
% END APPENDED VERSION V5 MATERIAL
% =============================================================================


% =============================================================================
% BEGIN APPENDED VERSION V6 MATERIAL
% =============================================================================

\clearpage
\section{Version V6: projective local shell heat baths and total-cylinder stress incidence}
\label{sec:V6-projective-local-shell}

Version V5 proved that an actual local heat-bath tangent has the exact
factorization $\mathcal J_0=H_0W/2$.  Its remaining cutoff problem was not the
Green inverse: it was the simultaneous choice of physical local blocks and
projective transport.  The supplied construction \textup{(RP.13)} obtains
exact transport by drawing the entire fine-detail field after every old
transition.  That kernel is a valid auxiliary projective skeleton, but an
all-detail draw can change separated detail coordinates in one jump.
Conversely, a fine singleton heat bath is jump-local but generally conditions
on the retained detail and therefore need not reproduce the old heat bath.

This version proves the exact compatibility criterion and gives a constructive
middle route.  A bounded-range factorization of the shell kernel determines a
Markov blanket for each old coordinate.  Updating that old coordinate together
with precisely its incident new details gives a reversible, exactly projective,
bounded-diameter heat bath.  The construction also exposes the precise
obstruction when the factorization is absent.

A second upstream issue concerns current and stress incidence.  The probability
score of the full cylinder \textup{(RP.7)} is not generally the centred
derivative of $S_{B,n}$ alone.  The positive determinant/Pfaffian control
modulus and a moving reference measure contribute to the same score.  Their
exact contribution is derived below before the V5 local-tangent theorem is
applied.

% -----------------------------------------------------------------------------
\subsection{The exact support gate in the positive-modulus cylinder}
\label{subsec:V6-support-gate}
% -----------------------------------------------------------------------------

At one cutoff suppress $n$ and write
\[
 a:=\|\sigma^+\|^2|p^+|^2,
 \qquad
 \dd\mu=Z^{-1}e^{-S_B}a\,\dd\lambda,
 \qquad 0<Z<\infty.
 \label{eq:V6-cylinder-density}
\]

\begin{proposition}[Equivalence is exactly a pointwise modulus question]
\label{prop:V6-equivalence-gate}
Assume $S_B$ is finite $\lambda$-almost everywhere.  Then
\begin{equation}
 \boxed{
 \mu\sim\lambda
 \quad\Longleftrightarrow\quad
 a>0\quad\lambda\text{-almost everywhere}.}
 \label{eq:V6-equivalence-criterion}
\end{equation}
If $N=\{a=0\}$ has positive $\lambda$-measure, then $\mu(N)=0$ and
$\lambda(N)>0$, so the product-equivalence hypothesis of
\cref{thm:V5-product-ergodicity} fails.  The statement that a section is
nonzero as a global vector does not by itself imply the almost-everywhere
condition in \eqref{eq:V6-equivalence-criterion}.
\end{proposition}

\begin{proof}
Since $S_B$ is finite, $e^{-S_B}$ is strictly positive
$\lambda$-almost everywhere.  The Radon--Nikodym density in
\eqref{eq:V6-cylinder-density} is therefore positive almost everywhere
exactly when $a$ is.  A probability law with positive finite density is
equivalent to its reference measure.  If $a$ vanishes on $N$, the displayed
density vanishes there, giving the stated obstruction.
\end{proof}

\begin{remark}[No bounded-density conclusion]
\label{rem:V6-no-density-window}
The criterion proves equivalence only.  It supplies neither a positive
essential infimum nor a finite essential supremum for $\dd\mu/\dd\lambda$.
Thus it can remove exact invariant functions through
\cref{thm:V5-product-ergodicity}, but it cannot invoke the quantitative
density-window gap theorem of V4.
\end{remark}

% -----------------------------------------------------------------------------
\subsection{One-shell notation and the Markov-blanket criterion}
\label{subsec:V6-blanket}
% -----------------------------------------------------------------------------

Let $X=(X_i,X_{-i})$ be an old finite-cutoff configuration with law $\mu$ on a
standard Borel product space.  Let $Y=(Y_a)_{a\in\mathcal A}$ be the new detail
and let
\[
 \nu(\dd x,\dd y)=\mu(\dd x)K(\dd y\mid x).
 \label{eq:V6-fine-law}
\]
The lift
\[
 J:L^2(\mu)\longrightarrow L^2(\nu),
 \qquad (Jf)(x,y)=f(x)
 \label{eq:V6-coarse-lift}
\]
is an isometry.  Let
\[
 E_i f=\mathbb E_\mu[f(X)\mid X_{-i}]
\]
be the old coordinate heat bath.  Choose a set $A_i\subseteq\mathcal A$ of
detail variables to update together with $X_i$, and define the fine block
heat bath
\[
 \widetilde E_iF
 :=\mathbb E_\nu[
 F\mid X_{-i},Y_{A_i^c}].
 \label{eq:V6-fine-block}
\]
It changes only the block
\[
 B_i:=\{X_i\}\cup\{Y_a:a\in A_i\}.
 \label{eq:V6-update-block}
\]

\begin{theorem}[Exact projectivity criterion for a local fine heat bath]
\label{thm:V6-blanket-criterion}
The following are equivalent:
\begin{enumerate}[label=\textup{(\roman*)},leftmargin=3em]
\item $\widetilde E_iJ=JE_i$ on $L^2(\mu)$;
\item for every bounded measurable $f$,
      \[
       \mathbb E_\nu[f(X)\mid X_{-i},Y_{A_i^c}]
       =\mathbb E_\mu[f(X)\mid X_{-i}];
       \]
\item
      \begin{equation}
       \boxed{X_i\ \perp\!\!\!\perp\ Y_{A_i^c}\mid X_{-i}}
       \label{eq:V6-blanket-independence}
      \end{equation}
      under $\nu$.
\end{enumerate}
Thus $A_i$ must contain a Markov blanket of $X_i$ in the new detail.  For
\[
 D_i:=\widetilde E_iJ-JE_i
 \label{eq:V6-block-defect}
\]
one has, for every $f\in L^2(\mu)$,
\begin{equation}
 \boxed{
 \|D_if\|_{L^2(\nu)}^2
 =
 \|\mathbb E_\nu[f(X)\mid X_{-i},Y_{A_i^c}]\|_2^2
 -\|\mathbb E_\mu[f(X)\mid X_{-i}]\|_2^2}
 \label{eq:V6-defect-Pythagoras}
\end{equation}
and equivalently
\begin{equation}
 \boxed{
 \|D_if\|_2^2
 =
 \mathbb E_\nu\operatorname{Var}\!\left(
 \mathbb E_\nu[f(X)\mid X_{-i},Y_{A_i^c}]
 \,\middle|\,X_{-i}\right).}
 \label{eq:V6-defect-variance}
\end{equation}
At a finite cutoff, vanishing on the point-indicator basis is a finite
necessary and sufficient test for \eqref{eq:V6-blanket-independence}.
\end{theorem}

\begin{proof}
By definition,
\[
 \widetilde E_iJf
 =\mathbb E_\nu[f(X)\mid X_{-i},Y_{A_i^c}].
\]
This is $JE_if$ for all bounded $f$ exactly when the conditional law of $X_i$
given $(X_{-i},Y_{A_i^c})$ equals its conditional law given $X_{-i}$.
That equality of regular conditional laws is equivalent to
\eqref{eq:V6-blanket-independence}.

For the quantitative identities, the tower property gives
\[
 \mathbb E_\nu[\widetilde E_iJf\mid X_{-i}]=JE_if.
\]
Hence $\widetilde E_iJf-JE_if$ is orthogonal to $JE_if$, proving
\eqref{eq:V6-defect-Pythagoras}.  Conditional variance gives
\eqref{eq:V6-defect-variance}.  On a finite carrier, conditional laws are
determined by their values on point indicators.
\end{proof}

\begin{corollary}[Exact generator defect]
\label{cor:V6-generator-defect}
Let
\[
 L=\sum_i c_i(E_i-I)
\]
be the old heat-bath generator.  On the fine carrier let
\[
 \widetilde L
 =\sum_i c_i(\widetilde E_i-I)
  +\sum_{a\in\mathcal A}d_a(F_a-I),
 \qquad d_a\ge0,
 \label{eq:V6-fine-generator}
\]
where $F_a$ is the heat bath of detail coordinate $Y_a$ given all other
coordinates.  Then
\begin{equation}
 \boxed{
 \widetilde LJ-JL=\sum_i c_iD_i.}
 \label{eq:V6-generator-defect}
\end{equation}
In particular, the Markov-blanket condition for every $i$ implies exact
generator and semigroup projectivity,
\[
 \widetilde LJ=JL,
 \qquad
 e^{t\widetilde L}J=Je^{tL}
 \quad(t\ge0).
 \label{eq:V6-semigroup-projectivity}
\]
\end{corollary}

\begin{proof}
Every detail-only heat bath fixes a coarse lift: $F_aJ=J$.  Expanding
\eqref{eq:V6-fine-generator} on $Jf$ gives
\eqref{eq:V6-generator-defect}.  If every $D_i$ vanishes, bounded-generator
intertwining passes to every power and hence to the exponential series.
\end{proof}
% -----------------------------------------------------------------------------
\subsection{A factorized shell gives a local projective dynamics}
\label{subsec:V6-factorized-shell}
% -----------------------------------------------------------------------------

Assume the detail kernel factorizes over a finite incidence hypergraph:
\begin{equation}
 \boxed{
 K(\dd y\mid x)
 =\bigotimes_{a\in\mathcal A}
   K_a(\dd y_a\mid x_{N(a)}),}
 \label{eq:V6-factor-kernel}
\end{equation}
where $N(a)$ is the set of old coordinates on which the $a$th detail factor
depends.  Define the incident-detail set
\[
 A_i:=\{a\in\mathcal A:i\in N(a)\}.
 \label{eq:V6-incident-set}
\]

\begin{theorem}[Markov-blanket block lift]
\label{thm:V6-block-lift}
For \eqref{eq:V6-factor-kernel} and \eqref{eq:V6-incident-set},
\[
 X_i\ \perp\!\!\!\perp\ Y_{A_i^c}\mid X_{-i}.
 \label{eq:V6-factor-independence}
\]
Consequently the generator \eqref{eq:V6-fine-generator}, with the block
$B_i=\{X_i\}\cup Y_{A_i}$ for every old coordinate and arbitrary detail-only
rates, is reversible with respect to $\nu$ and satisfies
\[
 \boxed{\widetilde LJ=JL,\qquad
 e^{t\widetilde L}J=Je^{tL}.}
 \label{eq:V6-local-projective-dynamics}
\]
Suppose the coordinates carry a metric and
\begin{equation}
 \sup_i\operatorname{diam}B_i\le R.
 \label{eq:V6-block-radius}
\end{equation}
If bounded multiplication writers $f,g$ have supports at distance greater
than $R$, then
\begin{equation}
 \boxed{\mathfrak C_{\widetilde L}(f,g)=0.}
 \label{eq:V6-local-commutator-zero}
\end{equation}
Thus exact projectivity and bounded-range jump locality coexist.

If, in addition, $\mu$ is equivalent to a product reference on the old
coordinates, every $K_a(\cdot\mid x_{N(a)})$ has an almost-everywhere positive
density relative to a product detail reference, and every $c_i,d_a$ is
positive, then
\begin{equation}
 \boxed{\Ker(-\widetilde L)=\mathbb R1.}
 \label{eq:V6-fine-ergodicity}
\end{equation}
There is no nonconstant exact invariant score on the fine carrier.
\end{theorem}

\begin{proof}
Integrating the factors with indices in $A_i$ gives one.  The conditional
marginal of $Y_{A_i^c}$ given $x$ is
\[
 \bigotimes_{a\notin A_i}K_a(\dd y_a\mid x_{N(a)}).
\]
None of the sets $N(a)$ in this product contains $i$, so the marginal depends
only on $x_{-i}$.  This proves
\eqref{eq:V6-factor-independence}.  Exact projectivity follows from
\cref{thm:V6-blanket-criterion,cor:V6-generator-defect}.  Conditional
expectations are self-adjoint orthogonal projections in $L^2(\nu)$, hence
their positive-rate sum is reversible.

Every jump of a block heat bath changes only that block.  Under
\eqref{eq:V6-block-radius}, no block meets both separated supports.  The
increment product in \eqref{eq:V3-jump-commutator} is therefore zero for every
possible jump, proving \eqref{eq:V6-local-commutator-zero}.

For the last assertion, let $F$ lie in the kernel of $-\widetilde L$.  Its
Dirichlet form is a positive sum of squared conditional residuals, so $F$ is
fixed by every detail-only heat bath and every old-coordinate block heat bath.
Product equivalence and the detail heat baths make $F$ independent of every
$Y_a$, hence $F=Jf$ for some old function $f$.  The block projectivity identity
then gives $JE_if=Jf$ for every $i$.  Product equivalence on the old carrier
and the coordinate-replacement proof of
\cref{thm:V5-product-ergodicity} make $f$ constant.
\end{proof}

\begin{corollary}[Uniform geometric range from local incidence]
\label{cor:V6-uniform-range}
Suppose every old coordinate and every detail coordinate has a spatial
location, and for some $r<\infty$,
\[
 i\in N(a)\quad\Longrightarrow\quad
 \operatorname{dist}(i,a)\le r.
 \label{eq:V6-incidence-range}
\]
Then $\operatorname{diam}B_i\le2r$ for every $i$.  A cofinal family with the
same $r$ therefore has a cutoff-independent simultaneous-jump range, even
though the number of coordinates in a block may vary.
\end{corollary}

\begin{proof}
Every point of $B_i$ is within distance $r$ of $i$.  The triangle inequality
bounds the distance between any two points of the block by $2r$.
\end{proof}

\begin{proposition}[A nonempty factorized subclass of the RPESM shell]
\label{prop:V6-RPESM-factor-shell}
In \textup{(RP.12)}, choose
\[
 \kappa(\dd y\mid x)
 =\bigotimes_{a\in\mathcal A}\kappa_a(\dd y_a\mid x_{N(a)}),
 \qquad
 V(x,y)=\sum_{a\in\mathcal A}V_a(x_{N(a)},y_a),
 \label{eq:V6-local-reference-action}
\]
with finite positive local partition functions
\[
 Z_a(x_{N(a)})
 =\int e^{-V_a(x_{N(a)},y_a)}
       \kappa_a(\dd y_a\mid x_{N(a)}).
\]
Then
\begin{equation}
 Z(x)=\prod_aZ_a(x_{N(a)}),
 \qquad
 K(\dd y\mid x)
 =\bigotimes_a
  \frac{e^{-V_a(x_{N(a)},y_a)}}{Z_a(x_{N(a)})}
  \kappa_a(\dd y_a\mid x_{N(a)}).
 \label{eq:V6-RPESM-factorization}
\end{equation}
If the local factors are chosen in reflection pairs and at least one
$V_a$ depends nontrivially on both old and detail variables, this is a
reflection-covariant, genuinely coupled instance of \textup{(RP.12)}.
With bounded-range $N(a)$, \cref{thm:V6-block-lift} supplies a reversible
local dynamics on the same static fine law which is exactly projective onto
the old heat-bath dynamics.
\end{proposition}

\begin{proof}
Both the exponential weight and the reference kernel factor over $a$.
Tonelli's theorem gives the product partition function and normalized kernel
in \eqref{eq:V6-RPESM-factorization}.  Reflection covariance follows by
pairing reflected factors.  A mixed nonconstant $V_a$ makes its conditional
detail law depend on the old variable, so the fine law is not a product of its
old and detail marginals.  The final assertion is
\cref{thm:V6-block-lift,cor:V6-uniform-range}.
\end{proof}

\begin{assessment}[What the construction changes]
\label{ass:V6-local-lift-scope}
The local block lift uses the same projective static law \textup{(RP.12)} but
need not equal the all-detail one-step skeleton \textup{(RP.13)}.  It supplies
a different reversible dynamics whose restriction to old observables is
exactly the old heat-bath dynamics.  This is sufficient for a projective
local-clock construction once the old generator itself is chosen from the
same recursive class.  Rates, physical durations, and block cardinalities
remain data to be normalized along the cofinal sequence.
\end{assessment}

% -----------------------------------------------------------------------------
\subsection{Why the two naive lifts fail}
\label{subsec:V6-naive-failures}
% -----------------------------------------------------------------------------

\begin{counterexample}[A retained detail biases a singleton old update]
\label{ex:V6-noisy-copy}
Let $X,Y\in\{-1,1\}$, let $X$ be uniform, and for
$0<\varepsilon<1/2$ set
\[
 \mathbb P(Y=X\mid X)=1-\varepsilon,
 \qquad
 \mathbb P(Y=-X\mid X)=\varepsilon.
\]
The old full heat bath sends the centred writer $f(X)=X$ to zero.  The fine
singleton heat bath of $X$ retains $Y$ and gives
\begin{equation}
 \boxed{
 \mathbb E[X\mid Y]=(1-2\varepsilon)Y,
 \qquad
 \|D_Xf\|_2^2=(1-2\varepsilon)^2>0.}
 \label{eq:V6-noisy-copy-defect}
\end{equation}
It is therefore not projective.  Updating the Markov-blanket block $(X,Y)$
instead gives $\mathbb E[X]=0$ and restores exact projectivity.
\end{counterexample}

\begin{proof}
Bayes' rule and the uniform prior give
$\mathbb P(X=Y\mid Y)=1-\varepsilon$.  Hence
$\mathbb E[X\mid Y]=(1-\varepsilon)Y-\varepsilon Y$.
The old conditional mean is zero, so the defect norm is the variance of the
displayed fine conditional mean.  With no retained variable, the joint block
conditional expectation is the unconditional mean.
\end{proof}

\begin{proposition}[The all-detail projective skeleton can be nonlocal]
\label{prop:V6-global-detail-nonlocal}
Let
\[
 P^\uparrow((x,y),\dd x'\dd y')
 =P(x,\dd x')K(\dd y'\mid x')
 \label{eq:V6-global-lift}
\]
be the kernel in \textup{(RP.13)}, and let
$L^\uparrow=P^\uparrow-I$.  Suppose $a,b$ are separated detail coordinates and
there exist a state $(x,y)$ and measurable sets $A,B$, with
$y_a\notin A$, $y_b\notin B$, such that
\begin{equation}
 \int P(x,\dd x')\,
 K(Y_a\in A,Y_b\in B\mid x')>0.
 \label{eq:V6-joint-detail-change}
\end{equation}
For $f=\mathbf1_{\{Y_a\in A\}}$ and
$g=\mathbf1_{\{Y_b\in B\}}$,
\begin{equation}
 \boxed{
 (\mathfrak C_{L^\uparrow}(f,g)\mathbf1)(x,y)
 =
 \int P(x,\dd x')K(Y_a\in A,Y_b\in B\mid x')>0.}
 \label{eq:V6-global-lift-residual}
\end{equation}
Thus the exact projectivity of \textup{(RP.14)} does not make the
all-detail skeleton a bounded-range physical update.
\end{proposition}

\begin{proof}
At the selected initial state both writers vanish.  Their two jump increments
are therefore the two terminal indicators, whose product is the indicator of
the joint event in \eqref{eq:V6-joint-detail-change}.  Apply the exact jump
formula \eqref{eq:V3-jump-commutator}.
\end{proof}

The two failures are complementary.  Retaining all detail keeps a singleton
update local but leaks information backward and breaks projectivity.  Drawing
all detail after an old transition preserves projectivity but can create an
all-distance simultaneous jump.  The Markov-blanket block lift removes both
defects when the shell incidence graph has bounded range.
% -----------------------------------------------------------------------------
\subsection{The total-cylinder score for current and stress deformations}
\label{subsec:V6-total-score}
% -----------------------------------------------------------------------------

Let $\theta$ be a gauge, coframe, metric, or coupling parameter.  To compare
all laws on one carrier, choose a fixed master measure $\lambda$ and write
\[
 \dd\lambda_\theta=j_\theta\,\dd\lambda,
 \qquad
 a_\theta=\|\sigma_\theta^+\|^2|p_\theta^+|^2,
\]
\[
 q_\theta
 =e^{-S_{B,\theta}}a_\theta j_\theta,
 \qquad
 \dd\mu_\theta=Z_\theta^{-1}q_\theta\,\dd\lambda.
 \label{eq:V6-parametric-cylinder}
\]
Work on a parameter interval where $q_\theta>0$ almost everywhere.  Assume
the logarithmic derivatives below lie in $L^2(\mu_0)$ and that differentiation
under the normalization integral is justified.

For a tangent direction $\xi$, put
\begin{equation}
 r_\xi
 :=\left.\partial_\xi\right|_0\log(a_\theta j_\theta),
 \qquad
 \tau_\xi
 :=\dot S_{B,\xi}-r_\xi,
 \qquad
 w_\xi:=\tau_\xi-\mathbb E_\mu\tau_\xi.
 \label{eq:V6-total-potential}
\end{equation}

\begin{theorem}[Exact total-cylinder Hellinger score]
\label{thm:V6-total-cylinder-score}
The normalized law and its square root satisfy
\begin{equation}
 \boxed{
 \left.\partial_\xi\right|_0
 \log\frac{\dd\mu_\theta}{\dd\lambda}
 =-w_\xi,
 \qquad
 \dot s_\xi=-\frac12M_{w_\xi}s.}
 \label{eq:V6-total-Hellinger}
\end{equation}
On a nonvanishing line chart with metric-compatible derivatives,
\begin{equation}
 \boxed{
 r_\xi
 =
 2\operatorname{Re}
 \frac{\langle\sigma^+,\nabla_\xi\sigma^+\rangle}
      {\|\sigma^+\|^2}
 +2\operatorname{Re}\frac{\nabla_\xi p^+}{p^+}
 +\dot{\log j}_\xi.}
 \label{eq:V6-line-reference-score}
\end{equation}
This expression is invariant under unitary frame changes of the determinant
line and sign-preserving frame changes of the real Pfaffian line.

Let
\[
 w_\xi^{B}:=\dot S_{B,\xi}-\mathbb E_\mu\dot S_{B,\xi},
 \qquad
 r_\xi^\circ:=r_\xi-\mathbb E_\mu r_\xi.
\]
Then
\begin{equation}
 \boxed{w_\xi=w_\xi^B-r_\xi^\circ.}
 \label{eq:V6-bare-total-split}
\end{equation}
For a finite tangent bank, let $W_B,W_{\rm tot},R$ synthesize respectively
$M_{w_\xi^B}s$, $M_{w_\xi}s$, and $M_{r_\xi^\circ}s$.  The exact incidence
defect is
\begin{equation}
 \boxed{
 W_B-W_{\rm tot}=R,
 \qquad
 R^*R=\bigl(\operatorname{Cov}_\mu(r_\xi,r_\zeta)\bigr)_{\xi,\zeta}.}
 \label{eq:V6-amplitude-incidence-defect}
\end{equation}
Thus a bare common-action derivative is the probability score precisely when
the centred line-modulus/reference derivative vanishes on the selected
tangent bank.
\end{theorem}

\begin{proof}
From \eqref{eq:V6-parametric-cylinder},
\[
 \partial_\xi\log q_\theta\big|_0
 =-\dot S_{B,\xi}+r_\xi=-\tau_\xi.
\]
Differentiating $Z_\theta=\int q_\theta\,\dd\lambda$ gives
$\partial_\xi\log Z_\theta|_0=-\mathbb E_\mu\tau_\xi$.  Subtracting the
normalization derivative proves the first identity in
\eqref{eq:V6-total-Hellinger}; differentiating the square root proves the
second.

Metric compatibility gives
\[
 \partial_\xi\log\|\sigma_\theta^+\|^2
 =2\operatorname{Re}
   \frac{\langle\sigma^+,\nabla_\xi\sigma^+\rangle}{\|\sigma^+\|^2},
\]
and the scalar real-line calculation gives the second term in
\eqref{eq:V6-line-reference-score}.  Logarithmic derivatives of norms do not
depend on a unitary frame, and a sign-preserving real frame has zero
logarithmic sign derivative.  Centering
$\tau_\xi=\dot S_{B,\xi}-r_\xi$ proves
\eqref{eq:V6-bare-total-split}.  Multiplication by $s$ converts that identity
to the synthesis identity; taking the Gram of $R$ gives the covariance matrix.
\end{proof}

\begin{corollary}[The actual local current/stress tangent]
\label{cor:V6-total-local-tangent}
Under the heat-bath continuity hypotheses of
\cref{thm:V5-local-tangent-factor}, the local kinetic source created by the
deformation \eqref{eq:V6-parametric-cylinder} is
\begin{equation}
 \boxed{
 \mathcal J_{0,\xi}
 =\frac12H_0M_{w_\xi}s
 =\frac12H_0M_{w_\xi^B-r_\xi^\circ}s.}
 \label{eq:V6-total-local-incidence}
\end{equation}
Using $w_\xi^B$ alone changes the tangent by $H_0M_{r_\xi^\circ}s/2$.
Its vanishing is an incidence statement and must be checked before the
probability tangent is identified with a current or stress insertion.
\end{corollary}

\begin{proof}
Apply \cref{thm:V5-local-tangent-factor} to
\eqref{eq:V6-total-Hellinger} and then use
\eqref{eq:V6-bare-total-split}.
\end{proof}

\begin{assessment}[Stress normalization versus represented stress]
\label{ass:V6-stress-scope}
Equation \eqref{eq:V6-total-local-incidence} identifies the derivative of the
stationary heat-bath Hamiltonian.  Conventional factors converting a metric
variation into $T_{\mu\nu}$ can be applied afterward.  This probability-level
incidence does not by itself prove the OS-represented Ward, contact,
symmetry, common-domain, or energy--Hamiltonian rows already required by the
attached stress/current programme.  It supplies the correct source to which
those downstream tests must be applied.
\end{assessment}

% -----------------------------------------------------------------------------
\subsection{Projective compression of the actual local tangent}
\label{subsec:V6-score-compression}
% -----------------------------------------------------------------------------

Consider a differentiable projective pair
\[
 \nu_\theta(\dd x,\dd y)
 =\mu_\theta(\dd x)K_\theta(\dd y\mid x),
 \qquad
 \pi_*\nu_\theta=\mu_\theta.
 \label{eq:V6-parametric-projective-pair}
\]
Use the sign convention
\[
 \partial_\theta\log\mu_\theta|_0=-w_c,
 \qquad
 \partial_\theta\log\nu_\theta|_0=-w_f.
\]
Differentiating the marginal identity gives the one-shell score martingale
\begin{equation}
 \boxed{\mathbb E_\nu[w_f\mid X]=w_c.}
 \label{eq:V6-score-martingale}
\end{equation}
This is the one-shell specialization of the canonical record-likelihood and
Fisher tower already proved in the supplied monograph; it is imported here.

Let
\[
 \iota:L^2(\mu)\longrightarrow L^2(\nu)
\]
be the square-root lift, so that $\iota(fs_c)=(Jf)s_f$.  For a command bank,
let $W_c,W_f$ be the corresponding square-root score syntheses.  Equation
\eqref{eq:V6-score-martingale} is
\begin{equation}
 \iota^*W_f=W_c.
 \label{eq:V6-score-compression}
\end{equation}
Put
\[
 W_{\rm sh}:=(I-\iota\iota^*)W_f.
 \label{eq:V6-shell-score}
\]
The imported Fisher Pythagoras is
\[
 W_f^*W_f=W_c^*W_c+W_{\rm sh}^*W_{\rm sh}.
 \label{eq:V6-shell-Fisher}
\]

\begin{theorem}[Exact compression and shell split of local kinetic incidence]
\label{thm:V6-local-tangent-compression}
Let $H_c,H_f$ be the coarse and fine nonnegative local heat-bath Hamiltonians.
Assume
\begin{equation}
 \boxed{H_f\iota=\iota H_c.}
 \label{eq:V6-Hamiltonian-projectivity}
\end{equation}
Define the actual local tangent columns
\[
 \mathcal J_c=\frac12H_cW_c,
 \qquad
 \mathcal J_f=\frac12H_fW_f.
\]
Then
\begin{align}
 \boxed{\iota^*\mathcal J_f=\mathcal J_c,}
 \label{eq:V6-tangent-compression}\\
 \boxed{
 \mathcal J_f
 =\iota\mathcal J_c+\mathcal J_{\rm sh},
 \qquad
 \mathcal J_{\rm sh}
 =\frac12H_fW_{\rm sh},
 \qquad
 \iota^*\mathcal J_{\rm sh}=0.}
 \label{eq:V6-tangent-shell-split}
\end{align}
The two subspaces reduce $H_f$, and the supported Green forms split exactly:
\begin{equation}
 \boxed{
 \mathcal J_f^*H_f^\dagger\mathcal J_f
 =
 \mathcal J_c^*H_c^\dagger\mathcal J_c
 +\frac14W_{\rm sh}^*H_fW_{\rm sh}.}
 \label{eq:V6-Green-shell-Pythagoras}
\end{equation}
The shell term is finite without a fine local gap and is the Green price of
the newly visible conditional score, rather than a duplicated price for the
old score.
\end{theorem}

\begin{proof}
Self-adjointness and \eqref{eq:V6-Hamiltonian-projectivity} imply
$\iota^*H_f=H_c\iota^*$.  With \eqref{eq:V6-score-compression},
\[
 \iota^*\mathcal J_f
 =\frac12H_c\iota^*W_f
 =\frac12H_cW_c=\mathcal J_c.
\]
Also $\iota W_c=\iota\iota^*W_f$, and hence
\[
 \mathcal J_f-\iota\mathcal J_c
 =\frac12H_f(W_f-\iota W_c)
 =\frac12H_fW_{\rm sh}.
\]
The orthogonal complement of $\Ran\iota$ is invariant under $H_f$ because
$\iota^*H_f=H_c\iota^*$; therefore
$\iota^*\mathcal J_{\rm sh}=0$ and the decomposition reduces $H_f$.
Spectral calculus on the two reducing subspaces gives
$H_f^\dagger\iota=\iota H_c^\dagger$ on the supported domains.  The Green
cross terms vanish.  The coarse term becomes the first term on the right of
\eqref{eq:V6-Green-shell-Pythagoras}, while
\cref{thm:V5-local-Fisher-Pythagoras} applied to
$\mathcal J_{\rm sh}=H_fW_{\rm sh}/2$ gives the second.
\end{proof}

\begin{corollary}[Factorized shells transport the correct stress tangent]
\label{cor:V6-factor-stress-transport}
For the factorized local block lift of
\cref{thm:V6-block-lift,prop:V6-RPESM-factor-shell}, any differentiable
total-cylinder command bank satisfying
\eqref{eq:V6-total-Hellinger} obeys
\eqref{eq:V6-tangent-compression} and
\eqref{eq:V6-Green-shell-Pythagoras}.  A parameter dependence of the shell
kernel appears only through $W_{\rm sh}$; it does not alter the compressed old
tangent.
\end{corollary}

\begin{proof}
The local block construction gives
\eqref{eq:V6-Hamiltonian-projectivity}.  Projectivity of the parameterized
static laws gives the score martingale.  Apply
\cref{thm:V6-local-tangent-compression}.
\end{proof}
% -----------------------------------------------------------------------------
\subsection{A cofinal local-Green telescope}
\label{subsec:V6-Green-telescope}
% -----------------------------------------------------------------------------

The supplied likelihood tower makes Fisher increments positive.  The preceding
local-incidence result gives a different positive increment: the Green price
of each shell conditional score after the local clock has acted.

\begin{theorem}[Exact shell telescope for projective local Green forms]
\label{thm:V6-Green-telescope}
Let
\[
 (\mathcal K_n,H_n,W_n,\mathcal J_n),
 \qquad
 \mathcal J_n=\frac12H_nW_n,
\]
be a cofinal finite-cutoff family on one finite command space $U$.  Suppose
each adjacent isometry $\iota_n:\mathcal K_n\to\mathcal K_{n+1}$ satisfies
\[
 H_{n+1}\iota_n=\iota_nH_n,
 \qquad
 \iota_n^*W_{n+1}=W_n.
 \label{eq:V6-adjacent-score-H}
\]
Put
\[
 W_n^{\rm sh}
 :=(I-\iota_n\iota_n^*)W_{n+1},
 \qquad
 \Delta_n^G
 :=\frac14(W_n^{\rm sh})^*H_{n+1}W_n^{\rm sh}.
 \label{eq:V6-shell-Green-increment}
\]
Then $\Delta_n^G\succeq0$ and
\begin{equation}
 \boxed{
 G_{n+1}=G_n+\Delta_n^G,
 \qquad
 G_n:=\mathcal J_n^*H_n^\dagger\mathcal J_n.}
 \label{eq:V6-adjacent-Green}
\end{equation}
Consequently
\begin{equation}
 \boxed{
 G_N=G_0+\sum_{n=0}^{N-1}\Delta_n^G.}
 \label{eq:V6-Green-telescope}
\end{equation}
The following are equivalent:
\begin{enumerate}[label=\textup{(\roman*)},leftmargin=3em]
\item $(G_N)$ converges in operator norm;
\item $\sup_N\operatorname{Tr}G_N<\infty$;
\item $\sum_n\operatorname{Tr}\Delta_n^G<\infty$.
\end{enumerate}
When they hold, the limit is the positive form
\[
 G_\infty=G_0+\sum_{n\ge0}\Delta_n^G.
 \label{eq:V6-Green-limit}
\]
No cutoffwise or uniform local spectral gap appears in this criterion.
\end{theorem}

\begin{proof}
Apply \cref{thm:V6-local-tangent-compression} to every adjacent pair.  This
gives \eqref{eq:V6-adjacent-Green}, and induction gives
\eqref{eq:V6-Green-telescope}.  The sequence is increasing in the Loewner
order.  On the finite-dimensional command space, an increasing sequence of
positive matrices converges in norm exactly when its traces are bounded.
Taking traces in the telescope makes that condition equivalent to summability
of the nonnegative shell traces.
\end{proof}

\begin{corollary}[Joint refresh removal and shell convergence]
\label{cor:V6-refresh-shell-limit}
In addition to the hypotheses of \cref{thm:V6-Green-telescope}, suppose
\[
 \sup_n\|W_n^*W_n\|\le C_M
\]
and choose $\rho_n\downarrow0$.  Let
$G_{\rho_n,n}^{[0]}$ be the same-local-source refreshed form of V5.  If
$\sum_n\operatorname{Tr}\Delta_n^G<\infty$, then
\begin{equation}
 \boxed{
 G_{\rho_n,n}^{[0]}\longrightarrow G_\infty,
 \qquad
 \|G_{n}-G_{\rho_n,n}^{[0]}\|
 \le\frac{C_M}{4}\rho_n.}
 \label{eq:V6-joint-limit}
\end{equation}
\end{corollary}

\begin{proof}
The Green telescope gives $G_n\to G_\infty$.  The quantitative difference is
\eqref{eq:V5-cofinal-Green-error}; the triangle inequality proves the limit.
\end{proof}

\begin{remark}[The new summability target]
\label{rem:V6-shell-summability}
A uniform Fisher bound controls removal of the auxiliary refresh but does not
by itself sum the physical shell contributions.  The exact remaining
quadratic target is
\[
 \sum_n\operatorname{Tr}
 \bigl((W_n^{\rm sh})^*H_{n+1}W_n^{\rm sh}\bigr)<\infty.
 \label{eq:V6-physical-shell-target}
\]
This target is source relative and can hold even when the ambient local gap
collapses.
\end{remark}

% -----------------------------------------------------------------------------
\subsection{RPESM consequence and remaining physical rows}
\label{subsec:V6-RPESM-consequence}
% -----------------------------------------------------------------------------

\begin{theorem}[Conditional projective local-clock packet for RPESM]
\label{thm:V6-RPESM-local-clock}
Suppose a nested RPESM cutoff family is chosen so that:
\begin{enumerate}[label=\textup{(P\arabic*)},leftmargin=3.2em]
\item the initial positive-modulus law passes
      \eqref{eq:V6-equivalence-criterion};
\item every adjacent shell has the reflection-covariant factorization
      \eqref{eq:V6-RPESM-factorization}, with a uniform incidence range $r$ and
      almost-everywhere positive local densities;
\item the local generators are defined recursively by the Markov-blanket
      blocks and positive detail heat baths in
      \eqref{eq:V6-fine-generator};
\item a gauge, coframe, metric, or coupling command bank uses the
      total-cylinder scores \eqref{eq:V6-total-potential};
\item the Hellinger and strong heat-bath differentiability hypotheses hold.
\end{enumerate}
Then at every adjacent pair:
\begin{enumerate}[label=\textup{(\alph*)},leftmargin=3em]
\item the local semigroups are exactly projective;
\item every jump block has diameter at most $2r$;
\item the centred local generator has no exact zero mode;
\item the actual local current/stress tangent compresses exactly to its old
      tangent;
\item its Green form obeys the positive shell identity
      \eqref{eq:V6-adjacent-Green}.
\end{enumerate}
If also the shell trace series
\eqref{eq:V6-physical-shell-target} converges and the Fisher matrices are
uniformly bounded, the quadratic local response has the refresh-free limit
\eqref{eq:V6-joint-limit}.
\end{theorem}

\begin{proof}
The support and positive-factor hypotheses propagate equivalence to the
product reference.  Items \textup{(a)}--\textup{(c)} follow from
\cref{thm:V6-block-lift,cor:V6-uniform-range}.  The total-cylinder theorem
identifies the correct score, and
\cref{cor:V6-factor-stress-transport} gives items \textup{(d)}--\textup{(e)}.
The final assertion is \cref{cor:V6-refresh-shell-limit}.
\end{proof}

\begin{assessment}[Status of the supplied regulator after V6]
\label{ass:V6-RPESM-status}
The theorem constructs a nonempty local and projective subclass within the
freedom of the shell definition \textup{(RP.12)}.  It replaces the false
choice between singleton locality and all-detail projectivity by a
Markov-blanket update.  It also turns cofinal local Green convergence into the
explicit positive series \eqref{eq:V6-physical-shell-target}.

The supplied manuscripts do not print a factorization of their selected
$K_{n+1/n}$ into the local factors of
\eqref{eq:V6-RPESM-factorization}, do not state that the control sections are
almost everywhere nonvanishing, and do not give the line-modulus/reference
score $r_\xi$ for the stress/current bank.  Therefore the hypotheses of
\cref{thm:V6-RPESM-local-clock} are a rigorous construction target, not a
claim that the already named full-coupling sequence satisfies them.

Even after those rows pass, the probability tangent must still pass the
represented stress/current residual, Ward-contact, energy--Hamiltonian,
common-core, and OS-null tests in the attached programme.  The local-clock
packet does not replace them.
\end{assessment}

% -----------------------------------------------------------------------------
\subsection{Revised continuum acquisition}
\label{subsec:V6-acquisition}
% -----------------------------------------------------------------------------

\begin{openproblem}[Executable V7 acquisition: populate one physical factor shell]
\label{op:V6-next-acquisition}
Choose one adjacent pair in the declared RPESM sequence and print the following
data on the same history:
\begin{enumerate}[label=\textup{(A\arabic*)},leftmargin=3.2em]
\item the product reference coordinates and the zero set of
      $\|\sigma_n^+\|^2|p_n^+|^2$;
\item the local detail factors $K_a$, their dependency sets $N(a)$, and the
      largest physical block diameter;
\item the point-indicator defect matrices from
      \eqref{eq:V6-defect-Pythagoras} for every retained detail outside the
      proposed blankets;
\item the complete gauge/coframe score
      $\dot S_B-r_\xi$ with the determinant, Pfaffian, and reference-Jacobian
      pieces of \eqref{eq:V6-line-reference-score};
\item the conditional shell score $W_n^{\rm sh}$ and its positive Green
      increment \eqref{eq:V6-shell-Green-increment};
\item the direct represented Ward/contact and energy--Hamiltonian residuals
      for the same current/stress direction.
\end{enumerate}
Repeat on a cofinal family and test the nonnegative series
\eqref{eq:V6-physical-shell-target}.  A passing packet supplies a genuinely
local, projective, refresh-free quadratic current/stress response.  Failure
returns an explicit modulus zero set, nonlocal dependency edge, blanket
defect, line/reference incidence term, divergent shell Green series, or
represented Ward residual.
\end{openproblem}

\begin{theorem}[Version V6 terminal theorem]
\label{thm:V6-terminal}
Preserving V1--V5 verbatim before this appended block, V6 proves:
\begin{enumerate}[label=\textup{(V6.\arabic*)},leftmargin=5.2em]
\item equivalence of the RPESM positive-modulus law to its reference is exactly
      the almost-everywhere nonvanishing of its control modulus;
\item a fine block heat bath projects to the old coordinate heat bath exactly
      when its retained detail is conditionally independent of the updated old
      coordinate;
\item the projectivity defect has the exact conditional-variance
      Pythagoras \eqref{eq:V6-defect-variance};
\item factorized finite-range shells admit reversible Markov-blanket heat
      baths that are both exactly projective and uniformly jump-local;
\item positive full-coordinate coverage removes exact centred zero modes for
      those block lifts;
\item the all-detail skeleton \textup{(RP.13)} has a sharp separated-jump
      locality residual whenever two distant details can change together;
\item the total cylinder current/stress score contains the line-modulus and
      reference-measure derivatives, with exact bare-action incidence defect
      \eqref{eq:V6-amplitude-incidence-defect};
\item projective local kinetic tangents compress exactly and their Green forms
      split into old and positive shell contributions;
\item the local Green forms telescope monotonically, with finite cofinal limit
      exactly under the shell trace summability criterion;
\item under the displayed conditional RPESM hypotheses, the quadratic
      current/stress response has a projective, jump-local, refresh-free
      limit without a uniform local gap.
\end{enumerate}
\end{theorem}

\begin{proof}
Items \textup{(V6.1)}--\textup{(V6.3)} are
\cref{prop:V6-equivalence-gate,thm:V6-blanket-criterion}.
Items \textup{(V6.4)}--\textup{(V6.5)} are
\cref{thm:V6-block-lift,prop:V6-RPESM-factor-shell}.
Item \textup{(V6.6)} is \cref{prop:V6-global-detail-nonlocal}.
Item \textup{(V6.7)} is
\cref{thm:V6-total-cylinder-score,cor:V6-total-local-tangent}.
Item \textup{(V6.8)} is
\cref{thm:V6-local-tangent-compression,cor:V6-factor-stress-transport}.
Item \textup{(V6.9)} is \cref{thm:V6-Green-telescope}.
The final item is
\cref{cor:V6-refresh-shell-limit,thm:V6-RPESM-local-clock}.
\end{proof}

\begin{assessment}[Programme boundary after V6]
\label{ass:V6-boundary}
V6 supplies the first explicit architecture in these notes that makes the
RPESM shell dynamics simultaneously local and exactly projective.  It also
identifies the correct total-cylinder probability score and gives a positive
shell-by-shell convergence criterion for its local Green form.  These results
remove the local-clock incompatibility as an abstract obstruction on the
factorized-shell branch.

The full Standard--Model--Einstein continuum remains unproved.  The next
required evidence is a populated physical factor shell with its control
zero-set audit, total stress/current score, Green increment, and represented
Ward residual.  Cofinal shell summability, interacting infrared control,
determinant/Pfaffian transport on the selected physical line, renormalized
composite domains, Lorentzian reconstruction, and the Palatini residual
remain open rows.
\end{assessment}

\section*{V6 exact checks and append-only record}
\addcontentsline{toc}{section}{V6 exact checks and append-only record}
For the noisy-copy separator,
\[
 \mathbb E[X\mid Y]=(1-2\varepsilon)Y,
 \qquad
 \mathbb E\bigl[\mathbb E[X\mid Y]^2\bigr]
 =(1-2\varepsilon)^2.
\]
For a factor kernel
$K(y_1,y_2\mid x_1,x_2)=K_1(y_1\mid x_1)K_2(y_2\mid x_2)$,
the $(x_1,y_1)$ block conditional expectation of every coarse $f(x_1,x_2)$
is exactly $\mathbb E[f\mid x_2]$, while the singleton $x_1$ update generally
conditions on $y_1$.  These finite identities, the separated-jump
double-commutator, the total-score normalization derivative, and the Green
telescope are independently replayed in the structural verification of this
source.

The complete V5 source before removal of its final active document-closing
command is retained verbatim.  Its complete-file SHA--256 digest is
\begin{center}\small\ttfamily
d5600aeb45cdb17ae22d153e1782f823c727d32a9edd1db68e401b290c2c6ab5.
\end{center}
A Version V7 must retain the complete V6 source, remove only the final active
document-closing command, append new material, restore one final command, and
use a filename ending in \texttt{\_V7.tex}.

% =============================================================================
% END APPENDED VERSION V6 MATERIAL
% =============================================================================


% =============================================================================
% BEGIN APPENDED VERSION V7 MATERIAL
% =============================================================================

\clearpage
\section{Version V7: Fisher-to-Green shell summability and an exact \texorpdfstring{$\mathbb Z_3$}{Z3} control tower}
\label{sec:V7-Fisher-Green}

Version V6 reduced refresh-free local response convergence to the positive
series
\[
 \sum_n\operatorname{Tr}
 \bigl((W_n^{\rm sh})^*H_{n+1}W_n^{\rm sh}\bigr).
\]
A bound by the ambient norm $\|H_{n+1}\|$ is poorly adapted to a growing
lattice: the generator contains more update blocks at every cutoff even when
the shell score is supported near one physical cell.  The correct estimate
counts only update blocks that act nontrivially on the shell-score range.

This version proves the source-relative estimate
\[
 \Delta_n^G\preceq\frac{C_n^{\rm act}}4\,\Delta_n^M,
\]
where $\Delta_n^M$ is the already positive Fisher increment and
$C_n^{\rm act}$ is the sum of rates of the source-active blocks.  A uniform
active-rate bound converts a bounded Fisher tower into a convergent local
Green tower, without an ambient generator bound or a spectral gap.

The final part populates every formula on an exact compact
$\mathbb Z_3$ factor-shell tower.  It is a rigorous gauge-like control of the
new mechanism, not a replacement for the unbounded Higgs, line, represented
stress, or Einstein rows.

% -----------------------------------------------------------------------------
\subsection{Conditional factor scores of a projective shell}
\label{subsec:V7-factor-scores}
% -----------------------------------------------------------------------------

Let $\theta\in\mathbb R^p$ and suppose the parameterized fine law has the
factorized projective form
\[
 \nu_\theta(\dd x,\dd y)
 =\mu_\theta(\dd x)
  \bigotimes_{a\in\mathcal A}
  K_{a,\theta}(\dd y_a\mid x_{N(a)}).
 \label{eq:V7-parametric-factor-law}
\]
Use the V6 sign convention
\[
 \partial_\xi\log\mu_\theta|_0=-w_{c,\xi},
 \qquad
 \partial_\xi\log\nu_\theta|_0=-w_{f,\xi}.
\]
For each detail factor define
\begin{equation}
 \eta_{a,\xi}
 :=-\left.\partial_\xi\right|_0
       \log K_{a,\theta}(Y_a\mid X_{N(a)}).
 \label{eq:V7-local-shell-score}
\end{equation}
Assume these derivatives belong to $L^2(\nu)$.

\begin{theorem}[Orthogonal factor resolution of the shell Fisher increment]
\label{thm:V7-factor-Fisher}
For every $a$ and command direction $\xi$,
\begin{equation}
 \boxed{\mathbb E_\nu[\eta_{a,\xi}\mid X]=0.}
 \label{eq:V7-factor-centering}
\end{equation}
The fine score decomposes as
\begin{equation}
 \boxed{
 w_{f,\xi}=w_{c,\xi}+\sum_{a\in\mathcal A}\eta_{a,\xi}.}
 \label{eq:V7-score-factor-split}
\end{equation}
Distinct factors are conditionally orthogonal:
\begin{equation}
 \boxed{
 \mathbb E_\nu[\eta_{a,\xi}\eta_{b,\zeta}\mid X]=0
 \quad(a\ne b).}
 \label{eq:V7-factor-orthogonality}
\end{equation}
Consequently the shell Fisher increment in
\eqref{eq:V6-shell-Fisher} has the exact local sum
\begin{equation}
 \boxed{
 (\Delta^M)_{\xi\zeta}
 :=(W_{\rm sh}^*W_{\rm sh})_{\xi\zeta}
 =\sum_{a\in\mathcal A}
  \mathbb E_\nu[\eta_{a,\xi}\eta_{a,\zeta}].}
 \label{eq:V7-factor-Fisher-sum}
\end{equation}

If
\[
 K_{a,\theta}(\dd y_a\mid x_{N(a)})
 =Z_{a,\theta}(x_{N(a)})^{-1}
  e^{-V_{a,\theta}(x_{N(a)},y_a)}
  \kappa_a(\dd y_a\mid x_{N(a)})
 \label{eq:V7-Gibbs-factor}
\]
with parameter-independent $\kappa_a$, then
\begin{equation}
 \boxed{
 \eta_{a,\xi}
 =\dot V_{a,\xi}
  -\mathbb E_{K_a}[\dot V_{a,\xi}\mid X].}
 \label{eq:V7-centered-local-action}
\end{equation}
Thus the shell Fisher increment is a sum of conditional covariances of local
action derivatives.
\end{theorem}

\begin{proof}
Differentiate
$\int K_{a,\theta}(\dd y_a\mid x_{N(a)})=1$.  The result is
\[
 \int \partial_\xi\log K_{a,\theta}\big|_0\,
       K_a(\dd y_a\mid x_{N(a)})=0,
\]
which proves \eqref{eq:V7-factor-centering}.  Differentiating the logarithm of
\eqref{eq:V7-parametric-factor-law} and reversing the sign gives
\eqref{eq:V7-score-factor-split}.  Conditional on $X$, the detail variables
are independent; combining that independence with
\eqref{eq:V7-factor-centering} proves
\eqref{eq:V7-factor-orthogonality}.  The coarse score is also orthogonal to
every $\eta_a$ by conditional centering.  Taking the Gram of the shell sum
therefore gives \eqref{eq:V7-factor-Fisher-sum}.

For \eqref{eq:V7-Gibbs-factor},
\[
 -\partial_\xi\log K_{a,\theta}|_0
 =\dot V_{a,\xi}+\partial_\xi\log Z_{a,\theta}|_0.
\]
Differentiating the partition function gives
$\partial_\xi\log Z_{a,\theta}|_0
=-\mathbb E_{K_a}[\dot V_{a,\xi}\mid X]$, proving
\eqref{eq:V7-centered-local-action}.
\end{proof}

\begin{remark}[Imported versus new content]
\label{rem:V7-nonduplication}
The coarse/shell Fisher Pythagoras itself is the one-shell specialization of
the supplied record-likelihood tower and was imported in V6.  The new content
of \cref{thm:V7-factor-Fisher} is the resolution of that shell term into
conditionally orthogonal local factor scores, which makes its update incidence
and continuum summability computable.
\end{remark}

% -----------------------------------------------------------------------------
\subsection{Only source-active heat baths contribute}
\label{subsec:V7-active-blocks}
% -----------------------------------------------------------------------------

Let
\[
 H=\sum_{b\in\mathcal B}r_b(I-P_b),
 \qquad r_b\ge0,
 \label{eq:V7-heatbath-H}
\]
where the $P_b$ are square-root conjugates of conditional expectations.  For a
bounded source synthesis $W:U\to\mathcal H$, define
\[
 \mathcal B_{\rm act}(W)
 :=\{b\in\mathcal B:(I-P_b)W\ne0\},
 \qquad
 C^{\rm act}(W)
 :=\sum_{b\in\mathcal B_{\rm act}(W)}r_b.
 \label{eq:V7-active-rate}
\]

\begin{theorem}[Source-active Dirichlet domination]
\label{thm:V7-active-Dirichlet}
For every finite-dimensional command space $U$,
\begin{equation}
 \boxed{
 0\preceq W^*HW
 \preceq C^{\rm act}(W)\,W^*W.}
 \label{eq:V7-active-domination}
\end{equation}
Suppose the columns of $W$ are multiplication scores supported on a common
coordinate set $S$.  If a block disjoint from $S$ fixes every such score, then
\[
 \mathcal B_{\rm act}(W)
 \subseteq\{b:B_b\cap S\ne\varnothing\}.
 \label{eq:V7-active-support}
\]
If every coordinate lies in at most $D$ update blocks, $|S|\le m$, and
$r_b\le r_*$, then
\begin{equation}
 \boxed{C^{\rm act}(W)\le r_*Dm.}
 \label{eq:V7-incidence-bound}
\end{equation}
This estimate is independent of the total number of lattice blocks.
\end{theorem}

\begin{proof}
Since $I-P_b$ is an orthogonal projection,
\[
 0\preceq W^*(I-P_b)W\preceq W^*W.
\]
Terms outside $\mathcal B_{\rm act}(W)$ vanish.  Multiplying by $r_b$ and
summing proves \eqref{eq:V7-active-domination}.  A conditional update in a
block disjoint from the coordinate support leaves the score unchanged, giving
\eqref{eq:V7-active-support}.  Count incidences between the at most $m$
coordinates of $S$ and their at most $D$ blocks, then multiply by the largest
rate to obtain \eqref{eq:V7-incidence-bound}.
\end{proof}

\begin{corollary}[Local factor score bound]
\label{cor:V7-local-factor-energy}
In the factorized shell of \cref{thm:V7-factor-Fisher}, suppose the union of
the supports of the nonzero $\eta_{a,\xi}$ meets at most $m$ coordinates,
uniformly over unit command directions.  Under the update-degree and rate
bounds of \cref{thm:V7-active-Dirichlet},
\begin{equation}
 \boxed{
 W_{\rm sh}^*H_fW_{\rm sh}
 \preceq r_*Dm\,\Delta^M.}
 \label{eq:V7-local-factor-energy}
\end{equation}
\end{corollary}

\begin{proof}
Use $W=W_{\rm sh}$ in
\eqref{eq:V7-active-domination} and
$W_{\rm sh}^*W_{\rm sh}=\Delta^M$.
\end{proof}
% -----------------------------------------------------------------------------
\subsection{Fisher domination closes the Green shell series}
\label{subsec:V7-Fisher-dominates-Green}
% -----------------------------------------------------------------------------

For the projective local-tangent family of V6, write
\[
 \Delta_n^M
 :=(W_n^{\rm sh})^*W_n^{\rm sh}
 =M_{n+1}-M_n,
 \qquad
 \Delta_n^G
 :=\frac14(W_n^{\rm sh})^*H_{n+1}W_n^{\rm sh}.
 \label{eq:V7-two-increments}
\]
Let
\[
 C_n^{\rm act}:=C^{\rm act}(W_n^{\rm sh})
 \label{eq:V7-shell-active-rate}
\]
for the fine generator $H_{n+1}$.

\begin{theorem}[Fisher-to-Green shell domination]
\label{thm:V7-Fisher-Green-domination}
At every shell,
\begin{equation}
 \boxed{
 0\preceq\Delta_n^G
 \preceq\frac{C_n^{\rm act}}4\,\Delta_n^M.}
 \label{eq:V7-Fisher-Green-domination}
\end{equation}
Hence, for $N>M$,
\begin{equation}
 \boxed{
 0\preceq G_N-G_M
 \preceq\frac14
 \sum_{n=M}^{N-1}C_n^{\rm act}\Delta_n^M.}
 \label{eq:V7-Green-tail-general}
\end{equation}
If
\[
 C_n^{\rm act}\le C_*
 \quad\text{for every }n,
 \label{eq:V7-uniform-active-rate}
\]
then
\begin{equation}
 \boxed{
 0\preceq G_N-G_M
 \preceq\frac{C_*}{4}(M_N-M_M).}
 \label{eq:V7-Green-Fisher-tail}
\end{equation}
In particular, a uniformly bounded Fisher tower
\begin{equation}
 M_n\preceq C_MI
 \label{eq:V7-bounded-Fisher-tower}
\end{equation}
has limits $M_\infty$ and $G_\infty$ in operator norm, and
\begin{equation}
 \boxed{
 0\preceq G_\infty-G_N
 \preceq\frac{C_*}{4}(M_\infty-M_N).}
 \label{eq:V7-limit-tail}
\end{equation}
Thus V6's shell trace summability follows from two local rows: bounded Fisher
information and bounded source-active update rate.
\end{theorem}

\begin{proof}
Apply \cref{thm:V7-active-Dirichlet} with
$W=W_n^{\rm sh}$ and divide by four.  This proves
\eqref{eq:V7-Fisher-Green-domination}.  Sum it and use the exact Green and
Fisher telescopes to obtain \eqref{eq:V7-Green-tail-general}.  Under
\eqref{eq:V7-uniform-active-rate}, the Fisher increments telescope to
$M_N-M_M$, proving \eqref{eq:V7-Green-Fisher-tail}.

The Fisher matrices form an increasing sequence of positive matrices on the
fixed finite-dimensional command space.  The uniform bound makes them
norm-convergent.  Equation \eqref{eq:V7-Green-Fisher-tail} then makes the
Green sequence Cauchy and gives \eqref{eq:V7-limit-tail}.
\end{proof}

\begin{corollary}[Weighted alternative when incidence grows]
\label{cor:V7-weighted-incidence}
Without \eqref{eq:V7-uniform-active-rate}, it is enough that
\begin{equation}
 \sum_{n\ge0}C_n^{\rm act}\operatorname{Tr}\Delta_n^M<\infty.
 \label{eq:V7-weighted-Fisher-tail}
\end{equation}
Then $G_n$ converges and
\[
 \operatorname{Tr}(G_\infty-G_N)
 \le\frac14\sum_{n\ge N}
 C_n^{\rm act}\operatorname{Tr}\Delta_n^M.
 \label{eq:V7-weighted-Green-tail}
\]
\end{corollary}

\begin{proof}
Take traces in \eqref{eq:V7-Green-tail-general}.  The right side is the tail
of a convergent nonnegative series.  Trace convergence of increasing positive
matrices on a finite-dimensional space implies norm convergence.
\end{proof}

\begin{assessment}[Why the ambient norm is unnecessary]
\label{ass:V7-ambient-norm}
The norm of a sum of heat-bath generators may grow like the number of sites.
That growth is irrelevant to a shell score fixed by all but finitely many
blocks.  The quantity $C_n^{\rm act}$ records exactly the rates that survive
after restriction to the score range.  A continuum proof should therefore
audit block incidence before imposing an ambient generator bound.
\end{assessment}

% -----------------------------------------------------------------------------
\subsection{Uniform Fisher bounds for total-cylinder commands}
\label{subsec:V7-total-Fisher}
% -----------------------------------------------------------------------------

Recall the V6 synthesis identity
\[
 W_{\rm tot}=W_B-R,
\]
where $W_B$ is the centred bare-action derivative and $R$ is the centred
line-modulus/reference derivative.

\begin{theorem}[Two-row control of the total Fisher matrix]
\label{thm:V7-total-Fisher-bound}
Let
\[
 M_B=W_B^*W_B,\qquad M_r=R^*R,\qquad
 M_{\rm tot}=W_{\rm tot}^*W_{\rm tot}.
\]
Then
\begin{equation}
 \boxed{M_{\rm tot}\preceq2M_B+2M_r.}
 \label{eq:V7-total-Fisher-bound}
\end{equation}
If the bare command bank $v:U\to L^2(\mu)$ satisfies
\begin{equation}
 \operatorname{osc}(v(u))
 \le A\|u\|_U
 \quad(u\in U),
 \label{eq:V7-command-oscillation}
\end{equation}
then
\begin{equation}
 \boxed{M_B\preceq\frac{A^2}{4}I_U.}
 \label{eq:V7-oscillation-Fisher}
\end{equation}
Consequently, if $M_r\preceq R_*^2I_U$,
\begin{equation}
 \boxed{
 M_{\rm tot}
 \preceq\left(\frac{A^2}{2}+2R_*^2\right)I_U.}
 \label{eq:V7-total-uniform-bound}
\end{equation}
\end{theorem}

\begin{proof}
For every $u\in U$,
\[
 \|(W_B-R)u\|^2
 \le2\|W_Bu\|^2+2\|Ru\|^2,
\]
which is \eqref{eq:V7-total-Fisher-bound} in quadratic-form order.  The scalar
variance bound
$\operatorname{Var}Z\le(\operatorname{ess\,sup}Z-
\operatorname{ess\,inf}Z)^2/4$ applied to $v(u)$ proves
\eqref{eq:V7-oscillation-Fisher}.  Combine the two estimates.
\end{proof}

\begin{corollary}[Smeared compact gauge writers]
\label{cor:V7-smeared-gauge-Fisher}
Suppose
\[
 v_n(u)=\sum_{p\in\mathcal P_n}\alpha_{p,n}(u)\psi_{p,n},
 \qquad
 |\psi_{p,n}|\le1,
 \]
and
\begin{equation}
 \sup_n\sup_{\|u\|=1}
 \sum_{p\in\mathcal P_n}|\alpha_{p,n}(u)|
 \le A_*.
 \label{eq:V7-l1-smearing}
\end{equation}
Then $\operatorname{osc}(v_n(u))\le2A_*\|u\|$ and
\[
 M_{B,n}\preceq A_*^2I_U.
 \label{eq:V7-gauge-Fisher-bound}
\]
If the line/reference correction has a uniform Fisher bound, the
total-cylinder gauge command bank satisfies
\eqref{eq:V7-bounded-Fisher-tower}.  This applies to fixed-support or
volume-normalized Wilson writers with uniformly bounded coefficient
$\ell^1$ norm.
\end{corollary}

\begin{proof}
The absolute value of the sum is at most
$A_*\|u\|$, so its range has length at most $2A_*\|u\|$.  Apply
\cref{thm:V7-total-Fisher-bound}.
\end{proof}

\begin{remark}[Unbounded Higgs and stress commands]
\label{rem:V7-unbounded-Fisher}
An unbounded polynomial Higgs or coframe command is not covered by an
oscillation estimate.  For it, \eqref{eq:V7-bounded-Fisher-tower} is exactly a
uniform local second-moment requirement for the \emph{total} score, including
the line and reference terms.  Cutoffwise $L^{1+\eta}$ integrability in
\textup{(RP.8)} does not by itself give this uniform bound.
\end{remark}
% -----------------------------------------------------------------------------
\subsection{An exact populated \texorpdfstring{$\mathbb Z_3$}{Z3} factor-shell tower}
\label{subsec:V7-Z3-tower}
% -----------------------------------------------------------------------------

Let
\[
 \mathbb Z_3=\{1,\omega,\omega^2\},
 \qquad \omega=e^{2\pi i/3},
\]
and put
\[
 r(x,y):=\operatorname{Re}(\overline{x}y)
 \in\left\{1,-\frac12\right\}.
 \label{eq:V7-Z3-relative-writer}
\]
Let $X$ be uniform on $\mathbb Z_3$.  At shell $k$, conditionally on $X=x$,
draw $Y_k$ with
\begin{equation}
 K_{\beta_k}(y\mid x)
 =\frac{\exp(\beta_kr(x,y))}{Z(\beta_k)},
 \qquad
 Z(\beta)=e^\beta+2e^{-\beta/2}.
 \label{eq:V7-Z3-kernel}
\end{equation}
For $N\ge1$ define the projective law
\begin{equation}
 \nu_N(x,y_1,\ldots,y_N)
 =\frac13\prod_{k=1}^NK_{\beta_k}(y_k\mid x).
 \label{eq:V7-Z3-law}
\end{equation}
Forgetting $Y_N$ sends $\nu_N$ to $\nu_{N-1}$.

Let $E_{0,N}$ be the heat bath of the one spatial-cell block
$(X,Y_1,\ldots,Y_N)$, and let $F_{k,N}$ update only $Y_k$ conditional on all
other variables.  With $c>0$ and $d_k>0$, set
\begin{equation}
 L_N
 =c(E_{0,N}-I)+\sum_{k=1}^Nd_k(F_{k,N}-I),
 \qquad H_N=-L_N.
 \label{eq:V7-Z3-generator}
\end{equation}
All variables are assigned to the same physical cell.  The first block may
grow in cardinality, but its spatial diameter is zero.

Now choose real coefficients $\alpha_k$ and deform
$\beta_k(t)=\beta_k+t\alpha_k$.  Define
\begin{align}
 m(\beta)
 &:=\mathbb E_{K_\beta}[r(X,Y)\mid X]
 =\frac{e^\beta-e^{-\beta/2}}
        {e^\beta+2e^{-\beta/2}},
 \label{eq:V7-Z3-mean}\\
 v(\beta)
 &:=\operatorname{Var}_{K_\beta}(r(X,Y)\mid X)
 =\frac{9}{2}
   \frac{e^{3\beta/2}}{(e^{3\beta/2}+2)^2}.
 \label{eq:V7-Z3-variance}
\end{align}

\begin{theorem}[Exact projective \texorpdfstring{$\mathbb Z_3$}{Z3} Green tower]
\label{thm:V7-Z3-Green-tower}
The laws \eqref{eq:V7-Z3-law} have strictly positive density relative to
product counting measure.  The generators are reversible, have only constants
in their kernels, and intertwine exactly under every forgetting map.

For the deformation $\beta_k(t)=\beta_k+t\alpha_k$, put
\[
 \eta_k
 :=\alpha_k\bigl(m(\beta_k)-r(X,Y_k)\bigr).
 \label{eq:V7-Z3-score}
\]
Then
\begin{equation}
 \boxed{
 w_N=\sum_{k=1}^N\eta_k,
 \qquad
 \langle\eta_j,\eta_k\rangle_{\nu_N}=0
 \quad(j\ne k),}
 \label{eq:V7-Z3-score-sum}
\end{equation}
and
\begin{equation}
 \boxed{
 H_N\eta_k=(c+d_k)\eta_k.}
 \label{eq:V7-Z3-score-eigen}
\end{equation}
The Fisher and actual local Green forms are the exact scalar sums
\begin{align}
 \boxed{
 M_N=\sum_{k=1}^N\alpha_k^2v(\beta_k),}
 \label{eq:V7-Z3-Fisher}\\
 \boxed{
 G_N
 =\frac14\sum_{k=1}^N
  (c+d_k)\alpha_k^2v(\beta_k).}
 \label{eq:V7-Z3-Green}
\end{align}
Moreover
\begin{equation}
 0<v(\beta)\le\frac9{16}
 \quad(\beta\in\mathbb R).
 \label{eq:V7-Z3-variance-bound}
\end{equation}
If $(\alpha_k)\in\ell^2$ and
$\sup_k(c+d_k)\le C$, then
\begin{equation}
 \boxed{
 M_\infty\le\frac9{16}\sum_{k\ge1}\alpha_k^2,
 \qquad
 G_\infty\le\frac{9C}{64}\sum_{k\ge1}\alpha_k^2.}
 \label{eq:V7-Z3-limits}
\end{equation}
Thus this compact interacting factor-shell tower has a nonzero, exactly
projective, spatially local, refresh-free continuum command response.
\end{theorem}

\begin{proof}
Every weight in \eqref{eq:V7-Z3-kernel} is positive.  The factorization makes
the Markov-blanket block for $X$ equal to the single spatial cell containing
$X$ and all incident details.  The construction and kernel statement follow
from \cref{thm:V6-block-lift}.  Exact generator projectivity can also be seen
directly: $E_{0,N+1}$ applied to a function of the first $N$ variables is its
$\nu_N$ mean, the old $F_{k,N+1}$ reduce to $F_{k,N}$ by conditional
independence, and the new detail heat bath fixes every old function.

The logarithmic derivative of \eqref{eq:V7-Z3-kernel} is
\[
 \alpha_k\bigl(r(X,Y_k)-m(\beta_k)\bigr).
\]
Reversing the sign gives \eqref{eq:V7-Z3-score}.  Conditional centering and
conditional independence prove \eqref{eq:V7-Z3-score-sum}.  Both
$E_{0,N}\eta_k$ and $F_{k,N}\eta_k$ vanish.  For $j\ne k$, the writer
$\eta_k$ is fixed by $F_{j,N}$.  Expanding \eqref{eq:V7-Z3-generator} gives
\eqref{eq:V7-Z3-score-eigen}.  The Fisher formula is the orthogonal score
sum, and the V5 identity $G_N=W_N^*H_NW_N/4$ gives
\eqref{eq:V7-Z3-Green}.

Under $K_\beta$, the value $1$ has probability
\[
 p=\frac{e^{3\beta/2}}{e^{3\beta/2}+2}
\]
and the value $-1/2$ has probability $1-p$.  A two-point variance is the
product of the two probabilities times the squared separation, so
\[
 v(\beta)=p(1-p)\left(\frac32\right)^2,
\]
which is \eqref{eq:V7-Z3-variance} and is at most $9/16$.  Summing this bound,
and then the bound weighted by $c+d_k\le C$, proves
\eqref{eq:V7-Z3-limits}.
\end{proof}

\begin{corollary}[The active-rate inequality is sharp]
\label{cor:V7-active-sharp}
For the $k$th score in the preceding tower, exactly two heat-bath terms act
nontrivially: the cell block at rate $c$ and its detail block at rate $d_k$.
Thus
\[
 C^{\rm act}(\eta_k)=c+d_k,
 \qquad
 \frac14\langle\eta_k,H_N\eta_k\rangle
 =\frac{C^{\rm act}(\eta_k)}4\|\eta_k\|^2.
\]
Equality holds in \eqref{eq:V7-Fisher-Green-domination}.
\end{corollary}

\begin{proof}
This is exactly \eqref{eq:V7-Z3-score-eigen}.
\end{proof}

\begin{assessment}[Control-tower firewall]
\label{ass:V7-Z3-scope}
The tower is interacting: every conditional detail law depends on the old
$\mathbb Z_3$ variable when $\beta_k\ne0$.  It also gives an actual infinite
projective response, rather than only a finite separator.  Its growing
one-cell block, however, has unbounded scale cardinality, and its carrier
contains neither the full Standard-Model gauge group nor the unbounded Higgs,
fermion/Pfaffian line, coframe, represented stress, or Einstein residual.
It proves the sharpness and consistency of the V7 convergence mechanism, not
the physical continuum.
\end{assessment}
% -----------------------------------------------------------------------------
\subsection{Conditional RPESM consequence}
\label{subsec:V7-RPESM-consequence}
% -----------------------------------------------------------------------------

\begin{theorem}[Bounded-incidence closure of the RPESM local Green row]
\label{thm:V7-RPESM-Green-closure}
Assume the projective factor-shell and total-cylinder hypotheses of
\cref{thm:V6-RPESM-local-clock}.  For one fixed finite command bank, suppose
in addition that
\begin{enumerate}[label=\textup{(B\arabic*)},leftmargin=3.2em]
\item the factor scores have the conditional resolution
      \eqref{eq:V7-centered-local-action};
\item their source-active rates satisfy
      $C_n^{\rm act}\le C_*$;
\item the total-cylinder Fisher matrices satisfy
      $M_n\preceq C_MI$.
\end{enumerate}
Then the refresh-free local Green matrices converge in operator norm and
\begin{equation}
 \boxed{
 0\preceq G_\infty-G_N
 \preceq\frac{C_*}{4}(M_\infty-M_N).}
 \label{eq:V7-RPESM-Green-tail}
\end{equation}
For any auxiliary refresh sequence $\rho_N\downarrow0$, the same-source parent
forms converge to this limit with
\begin{equation}
 \boxed{
 \|G_{\rho_N,N}^{[0]}-G_\infty\|
 \le\frac{C_M}{4}\rho_N+\|G_N-G_\infty\|.}
 \label{eq:V7-RPESM-parent-tail}
\end{equation}
If the command writers have uniformly bounded coordinate support, update
degree, and rates, item \textup{(B2)} follows from
\eqref{eq:V7-incidence-bound}.  If they are compact gauge writers satisfying
\eqref{eq:V7-l1-smearing} and the line/reference Fisher row is uniformly
bounded, item \textup{(B3)} follows from
\eqref{eq:V7-total-uniform-bound}.
\end{theorem}

\begin{proof}
Items \textup{(B2)}--\textup{(B3)} are precisely the hypotheses of
\cref{thm:V7-Fisher-Green-domination}, which gives
\eqref{eq:V7-RPESM-Green-tail}.  The parent estimate is the triangle inequality
combined with \eqref{eq:V5-cofinal-Green-error}.  The two final implications
are \cref{thm:V7-active-Dirichlet,cor:V7-smeared-gauge-Fisher}.
\end{proof}

\begin{assessment}[What V7 closes]
\label{ass:V7-RPESM-status}
On any bounded-incidence factor-shell branch, the separate V6 demand for
summability of Green increments is no longer an independent analytic
hypothesis: it follows from a bounded total Fisher tower.  For compact,
properly normalized gauge commands, the bare-action part of that Fisher bound
is automatic.  The remaining gauge rows are the line/reference score and the
actual population of the factor shell.

For unbounded Higgs and coframe/stress commands, a uniform total Fisher bound
remains substantive.  The cutoffwise moment statement in \textup{(RP.8)} has
no printed uniform constant.  Promoting quartic coercivity to a uniform local
second-moment estimate is therefore the next upstream analytic task.
\end{assessment}

% -----------------------------------------------------------------------------
\subsection{Revised continuum acquisition}
\label{subsec:V7-acquisition}
% -----------------------------------------------------------------------------

\begin{openproblem}[Executable V8 acquisition: uniform quartic local Fisher control]
\label{op:V7-next-acquisition}
On one factorized unbounded RPESM shell family:
\begin{enumerate}[label=\textup{(A\arabic*)},leftmargin=3.2em]
\item write the conditional Higgs density at one updated block and isolate a
      cutoff-independent positive quartic coefficient after all negative
      local terms are absorbed by Young inequalities;
\item prove a uniform conditional moment or exponential-tail estimate strong
      enough for the Higgs action, hopping, and coframe derivatives in the
      total score \eqref{eq:V6-total-potential};
\item include the determinant/Pfaffian modulus and reference-Jacobian
      derivatives rather than bounding the bare bosonic polynomial alone;
\item verify a uniform update-degree/rate bound for the support of each
      conditional factor score;
\item insert the resulting Fisher constant into
      \eqref{eq:V7-RPESM-Green-tail};
\item retain separately the represented Ward/contact,
      energy--Hamiltonian, line transport, and Palatini residuals.
\end{enumerate}
A successful estimate closes the refresh-free quadratic local Green row for
that unbounded Higgs/stress bank.  Failure returns a named loss of quartic
coercivity, a line-score singularity, an unbounded reference Jacobian, growing
update incidence, or a divergent Fisher tail.
\end{openproblem}

\begin{theorem}[Version V7 terminal theorem]
\label{thm:V7-terminal}
Preserving V1--V6 verbatim before this appended block, V7 proves:
\begin{enumerate}[label=\textup{(V7.\arabic*)},leftmargin=5.2em]
\item a factorized parameterized shell score is the sum of conditionally
      centred local action derivatives;
\item distinct detail-factor scores are conditionally orthogonal, so the
      shell Fisher increment is the exact sum of local conditional
      covariances;
\item the Dirichlet energy of a source is bounded by the sum of rates of only
      those heat baths acting nontrivially on its range;
\item bounded coordinate support, update degree, and rates make that
      source-active bound independent of the lattice volume;
\item every Green shell increment is dominated by one quarter of its Fisher
      increment times the source-active rate;
\item a bounded Fisher tower and a uniform active-rate bound imply a
      convergent refresh-free Green tower with the tail
      \eqref{eq:V7-limit-tail};
\item the full cylinder Fisher matrix is controlled by separate bare-action
      and line/reference Fisher rows;
\item normalized compact gauge writers with uniformly bounded coefficient
      $\ell^1$ norm have a uniform bare Fisher bound;
\item the explicit interacting $\mathbb Z_3$ factor-shell tower has exact
      projective Fisher and Green limits for square-summable shell commands;
\item on a conditional RPESM factor-shell branch, these estimates close the
      quadratic local Green row without an ambient generator bound or a local
      spectral gap.
\end{enumerate}
\end{theorem}

\begin{proof}
Items \textup{(V7.1)}--\textup{(V7.2)} are
\cref{thm:V7-factor-Fisher}.
Items \textup{(V7.3)}--\textup{(V7.4)} are
\cref{thm:V7-active-Dirichlet,cor:V7-local-factor-energy}.
Items \textup{(V7.5)}--\textup{(V7.6)} are
\cref{thm:V7-Fisher-Green-domination,cor:V7-weighted-incidence}.
Items \textup{(V7.7)}--\textup{(V7.8)} are
\cref{thm:V7-total-Fisher-bound,cor:V7-smeared-gauge-Fisher}.
Item \textup{(V7.9)} is
\cref{thm:V7-Z3-Green-tower,cor:V7-active-sharp}.
The final item is \cref{thm:V7-RPESM-Green-closure}.
\end{proof}

\begin{assessment}[Programme boundary after V7]
\label{ass:V7-boundary}
V7 removes the Green-shell summability condition as an independent bottleneck
for bounded-incidence, bounded-Fisher command banks.  It also supplies a fully
populated infinite factor-shell control in which local projectivity, Fisher
increments, local kinetic incidence, Green convergence, and refresh removal
all coexist exactly.

The full Standard--Model--Einstein continuum remains unproved.  The principal
next analytic row is uniform total-score control on the unbounded quartic Higgs
carrier.  The selected physical shell factors and control-section zero set
are still unprinted, and the represented current/stress, interacting
infrared, determinant/Pfaffian, renormalized-domain, Lorentzian, and Palatini
rows remain open.
\end{assessment}

\section*{V7 exact checks and append-only record}
\addcontentsline{toc}{section}{V7 exact checks and append-only record}
For the $\mathbb Z_3$ factor,
\[
 \mathbb P(r=1)
 =\frac{e^{3\beta/2}}{e^{3\beta/2}+2},
 \qquad
 \mathbb P(r=-1/2)
 =\frac{2}{e^{3\beta/2}+2},
\]
so direct two-point arithmetic gives
\[
 \operatorname{Var}(r)
 =\frac94\,\mathbb P(r=1)\mathbb P(r=-1/2)
 =\frac92\frac{e^{3\beta/2}}{(e^{3\beta/2}+2)^2}.
\]
The maximum $9/16$ occurs when the two displayed probabilities are equal.
Finite enumerations independently verify conditional centering, orthogonality,
generator projectivity, the eigenvalue $c+d_k$, and the Fisher and Green sums.
Random positive-matrix tests independently verify the Loewner
Fisher-to-Green bound and its telescoping tail.

The complete V6 source before removal of its final active document-closing
command is retained verbatim.  Its complete-file SHA--256 digest is
\begin{center}\small\ttfamily
9c87c1c0673e4675e20d59b495e10987aaf17a2c0d760d622156b2f465cf406f.
\end{center}
A Version V8 must retain the complete V7 source, remove only the final active
document-closing command, append new material, restore one final command, and
use a filename ending in \texttt{\_V8.tex}.

% =============================================================================
% END APPENDED VERSION V7 MATERIAL
% =============================================================================


% =============================================================================
% BEGIN APPENDED VERSION V8 MATERIAL
% =============================================================================

\clearpage
\section{Version V8: quartic-tail export to total Fisher and local Green control}
\label{sec:V8-quartic-Fisher}

Version V7 reduced convergence of the refresh-free local Green response to two
uniform estimates: a bound on the total-cylinder Fisher matrix and a bound on
the rates of the heat baths which actually move the score.  Compact gauge
writers already satisfy the first estimate.  The next case is the unbounded
Higgs carrier.

The attached monograph contains two relevant statements which must not be
identified.  Its two-shell hard action has the printed volume-averaged fourth
moment bound
\[
 \frac1{|\mathcal H|}\sum_{x\in\mathcal H}
 \mathbb E\|\eta_x\|^4\le10+2\sqrt{10}.
\]
Later, the proof of the root-subsequence theorem invokes a
\emph{uniform local quartic-exponential estimate}, but no separately stated
estimate, constant, or proof accompanies that invocation.  A fourth moment by
itself does not bound the Fisher information of a quartic-coupling score.  This
version fills that precise analytic gap for bounded-degree quartic lattice
actions, exports the result to local polynomial scores, and treats the
determinant/Pfaffian line score by two independent sufficient criteria.

The result closes the V7 Fisher and Green rows for the pinned hard Higgs
family and for bounded-density perturbations having a controlled line score.
It does not establish those hypotheses for the complete interacting RPESM
density or construct the Standard--Model--Einstein continuum.

% -----------------------------------------------------------------------------
\subsection{Why the printed fourth moment is insufficient}
\label{subsec:V8-fourth-not-eighth}
% -----------------------------------------------------------------------------

\begin{proposition}[A fourth moment does not control a quartic score]
\label{prop:V8-fourth-not-Fisher}
There are probability laws $\mu_N$ on $\mathbb R$ such that
\[
 \sup_N\mathbb E_{\mu_N}|X|^4=1
\]
while the Fisher variance of the centred quartic writer diverges:
\begin{equation}
 \operatorname{Var}_{\mu_N}(|X|^4)=N^4-1\longrightarrow\infty.
 \label{eq:V8-fourth-separator}
\end{equation}
Consequently, a uniform fourth-moment estimate cannot replace a uniform
eighth moment or a suitable quartic-exponential estimate in the V7
total-Fisher hypothesis.
\end{proposition}

\begin{proof}
Give $X$ the value $N$ with probability $N^{-4}$ and the value zero otherwise.
Then
\[
 \mathbb E|X|^4=1,\qquad
 \mathbb E|X|^8=N^4.
\]
The variance identity proves \eqref{eq:V8-fourth-separator}.  The quartic
coupling deformation has raw derivative $|X|^4$ and centred score
$|X|^4-\mathbb E|X|^4$, so the displayed variance is exactly its scalar
Fisher information.
\end{proof}

% -----------------------------------------------------------------------------
\subsection{A uniform local exponential estimate from weighted coercivity}
\label{subsec:V8-weighted-coercivity}
% -----------------------------------------------------------------------------

Let $G=(V,E)$ be a finite graph of maximum degree at most $\Delta\ge1$.
At every $i\in V$ let $x_i\in\mathbb R^d$.  For an unoriented edge
$\{i,j\}$ let $U_{ij}$ be a linear contraction and let $J_{ij}\in\mathbb R$.
Consider
\begin{equation}
 S_G(x)
 =\sum_{i\in V}\bigl(\lambda_i\|x_i\|^4+q_i\|x_i\|^2\bigr)
 -\sum_{\{i,j\}\in E}J_{ij}\langle x_i,U_{ij}x_j\rangle
 \label{eq:V8-quartic-action}
\end{equation}
and the normalized law
\[
 \mu_G(\dd x)=Z_G^{-1}e^{-S_G(x)}\prod_{i\in V}\dd x_i.
\]
Assume
\begin{equation}
 \lambda_i\ge\lambda_0>0,\qquad
 q_i\ge-Q_0,\qquad
 \sup_i\sum_{j:\{i,j\}\in E}|J_{ij}|\le J_0.
 \label{eq:V8-coercive-data}
\end{equation}

\begin{theorem}[Uniform local quartic-exponential bound]
\label{thm:V8-local-quartic-exponential}
Fix a nonempty finite set $S\subset V$ and choose
$0<\alpha<\Delta^{-1}$.  Put
\begin{equation}
 B_\alpha
 =\frac{J_0(1+\alpha^{-1})+3Q_0}{\lambda_0}.
 \label{eq:V8-Balpha}
\end{equation}
For every
\begin{equation}
 0<\delta\le\frac{\lambda_0}{8|S|}
 \label{eq:V8-delta-range}
\end{equation}
one has the volume-independent estimate
\begin{equation}
 \boxed{
 \mathbb E_{\mu_G}
 \exp\!\left(\delta\sum_{o\in S}\|x_o\|^4\right)
 \le
 \exp\!\left[
  \frac{\delta|S|}{1-\Delta\alpha}
  \left(\frac d{\lambda_0}+\frac{B_\alpha^2}{4}\right)
 \right].}
 \label{eq:V8-local-exponential}
\end{equation}
The right side depends on the local support size and the uniform coercivity
and incidence constants, but not on $|V|$.
\end{theorem}

\begin{proof}
For $i\in V$ define the summable graph weight
\[
 a_i=\sum_{o\in S}\alpha^{\operatorname{dist}(o,i)},
\]
with $\alpha^\infty=0$ on components disjoint from $S$.  The elementary
growth bound on graph spheres gives
\begin{equation}
 \sum_i a_i\le\frac{|S|}{1-\Delta\alpha},
 \qquad
 0\le a_i\le |S|.
 \label{eq:V8-weight-sum}
\end{equation}
If $\{i,j\}\in E$ and $a_i+a_j>0$, the triangle inequality for graph distance
gives
\begin{equation}
 \alpha a_i\le a_j\le\alpha^{-1}a_i.
 \label{eq:V8-neighbour-weight}
\end{equation}
Also $a_o\ge1$ for $o\in S$, so it is enough to bound the exponential moment
of $\sum_i a_i\|x_i\|^4$.

Write
\[
 u_i=\frac{\delta a_i}{\lambda_i},
 \qquad
 s_i=(1-2u_i)^{-1/4}.
\]
By \eqref{eq:V8-delta-range}, $0\le2u_i\le1/4$.  In the numerator of
\[
 \mathbb E_{\mu_G}e^{\delta\sum_i a_i\|x_i\|^4}
\]
make the diagonal change of variables $x_i=s_i y_i$.  Its quartic
coefficient satisfies
\begin{equation}
 (\lambda_i-\delta a_i)s_i^4
 =\lambda_i+\frac{\delta a_i}{1-2u_i}
 \ge\lambda_i+\delta a_i.
 \label{eq:V8-quartic-buffer}
\end{equation}
For $0\le2u\le1/4$, elementary one-variable calculus gives
\begin{equation}
 s_i s_j-1
 \le\frac{2\delta}{\lambda_0}(a_i+a_j),
 \qquad
 s_i^2-1\le\frac{3\delta}{\lambda_0}a_i.
 \label{eq:V8-scaling-bounds}
\end{equation}
The part of $q_i$ which is nonnegative only decreases the transformed
integrand.  The negative part and the change of the edge term are, using
\eqref{eq:V8-neighbour-weight} and
$2\|y_i\|\|y_j\|\le\|y_i\|^2+\|y_j\|^2$, bounded above by
\begin{equation}
 \delta B_\alpha\sum_i a_i\|y_i\|^2.
 \label{eq:V8-quadratic-error}
\end{equation}
The quartic buffer absorbs this error pointwise because
\begin{equation}
 B_\alpha r^2-r^4\le\frac{B_\alpha^2}{4}
 \qquad(r\ge0).
 \label{eq:V8-Young-buffer}
\end{equation}
Finally, the Jacobian satisfies
\begin{equation}
 \log\prod_i s_i^d
 =-\frac d4\sum_i\log(1-2u_i)
 \le\frac{d\delta}{\lambda_0}\sum_i a_i.
 \label{eq:V8-Jacobian}
\end{equation}
After \eqref{eq:V8-quartic-buffer}--\eqref{eq:V8-Jacobian}, the transformed
numerator is at most
\[
 \exp\!\left[
  \delta\left(\frac d{\lambda_0}+\frac{B_\alpha^2}{4}\right)
  \sum_i a_i\right]Z_G.
\]
Divide by $Z_G$ and use \eqref{eq:V8-weight-sum}.
\end{proof}
\begin{corollary}[The missing hard-root local tail estimate]
\label{cor:V8-hard-root-tail}
Rewrite the symmetric quadratic term in the monograph's pinned hard action
\[
 \mathcal S_{\mathcal H,q}(\eta)
 =\sum_{x\in\mathcal H}
   \bigl(\|\eta_x\|^4+q\|\eta_x\|^2\bigr)
  -\langle\eta,\mathsf A_{\mathcal H}\eta\rangle
\]
as an edge sum, and let $\Delta_{\mathcal H}$ and $J_{\mathcal H}$ be uniform
bounds on its graph degree and absolute interaction row sum.  For any fixed
finite set $S\subset\mathcal H$, any
$0<\alpha<\Delta_{\mathcal H}^{-1}$, and
$0<\delta\le(8|S|)^{-1}$,
\begin{equation}
 \boxed{
 \sup_{\mathcal H}\sup_{q\ge0}
 \mathbb E_q
 e^{\delta\sum_{x\in S}\|\eta_x\|^4}
 \le
 \exp\!\left[
 \frac{\delta|S|}{1-\Delta_{\mathcal H}\alpha}
 \left(4+\frac{J_{\mathcal H}^2
 (1+\alpha^{-1})^2}{4}\right)\right].}
 \label{eq:V8-hard-root-tail}
\end{equation}
The first supremum ranges over the compatible finite periodic volumes.
In particular, the estimate invoked in the monograph's root-subsequence proof
holds with a printed constant and uniformly throughout its root bracket.
\end{corollary}

\begin{proof}
Apply \cref{thm:V8-local-quartic-exponential} with
$d=4$, $\lambda_0=1$, and $Q_0=0$.  The hard adjacency has uniformly bounded
degree and row sum, independently of the periodic volume.  No lower bound on
$q$ beyond $q\ge0$ and no Dobrushin contraction enter the proof.
\end{proof}

\begin{assessment}[The estimate does not select the Gaussian branch]
\label{ass:V8-tail-not-mixing}
The proof of \cref{cor:V8-hard-root-tail} uses quartic coercivity and bounded
incidence only.  It supplies local tightness and uniform integrability at the
retuned roots even when an influence margin approaches one.  It therefore
does not invoke the fixed strict-mixing hypothesis which the monograph proves
leads to a Gaussian-or-zero physical scaling limit.
\end{assessment}

% -----------------------------------------------------------------------------
\subsection{Exponential tails give explicit Fisher constants}
\label{subsec:V8-tail-to-Fisher}
% -----------------------------------------------------------------------------

\begin{lemma}[Exponential-moment export]
\label{lem:V8-exponential-export}
Let $R\ge0$ and suppose
\begin{equation}
 \mathbb E e^{\delta R}\le K
 \qquad(\delta,K>0).
 \label{eq:V8-QET}
\end{equation}
For every $p>0$,
\begin{equation}
 \boxed{
 \mathbb E R^p
 \le K\left(\frac{p}{e\delta}\right)^p.}
 \label{eq:V8-moment-export}
\end{equation}
In particular,
\begin{equation}
 \mathbb E R^2\le\frac{4K}{e^2\delta^2}.
 \label{eq:V8-second-R}
\end{equation}
\end{lemma}

\begin{proof}
The maximum of $r^pe^{-\delta r}$ over $r\ge0$ is attained at
$r=p/\delta$ and equals $(p/(e\delta))^p$.  Hence
\[
 R^p\le\left(\frac{p}{e\delta}\right)^pe^{\delta R}.
\]
Take expectations.
\end{proof}

Let $E$ be a finite-dimensional real command space.  At cutoff $n$, write the
raw derivative of the bosonic action as $b_n(u)$ and the derivative of the
line modulus and reference Jacobian as $r_n(u)$.  With the sign convention of
V6, the centred total-cylinder score is
\begin{equation}
 w_n(u)
 =b_n(u)-r_n(u)
  -\mathbb E_n\bigl[b_n(u)-r_n(u)\bigr].
 \label{eq:V8-total-score}
\end{equation}
Its Fisher operator is the positive operator $M_n$ determined by
\[
 \langle u,M_nu\rangle_E=\mathbb E_n|w_n(u)|^2.
\]

\begin{theorem}[Quartic-tail to total-Fisher bound]
\label{thm:V8-tail-to-total-Fisher}
Suppose that, for one fixed command support and all cutoffs $n$, there is a
nonnegative random variable $R_n$ such that
\begin{align}
 \mathbb E_ne^{\delta R_n}&\le K,
 \label{eq:V8-uniform-QET}\\
 |b_n(u)|&\le\|u\|_E(A+BR_n),
 \label{eq:V8-bosonic-envelope}\\
 \mathbb E_n|r_n(u)|^2&\le L^2\|u\|_E^2
 \label{eq:V8-line-L2}
\end{align}
for fixed $A,B,L,\delta,K$.  Then
\begin{equation}
 \boxed{
 0\preceq M_n\preceq C_FI_E,\qquad
 C_F=3A^2+\frac{12B^2K}{e^2\delta^2}+3L^2.}
 \label{eq:V8-Fisher-constant}
\end{equation}
No mixing, Poincar\'e, or spectral-gap estimate is required.
\end{theorem}

\begin{proof}
Centering is orthogonal projection away from the constants in $L^2(\mu_n)$,
so
\[
 \mathbb E_n|w_n(u)|^2
 \le\mathbb E_n|b_n(u)-r_n(u)|^2.
\]
The scalar inequality $(a+b+c)^2\le3(a^2+b^2+c^2)$ and
\eqref{eq:V8-bosonic-envelope} give
\[
 |b_n(u)-r_n(u)|^2
 \le3\|u\|_E^2(A^2+B^2R_n^2)+3|r_n(u)|^2.
\]
Apply \eqref{eq:V8-second-R} and \eqref{eq:V8-line-L2}.  The resulting
quadratic-form bound for every $u$ is exactly
\eqref{eq:V8-Fisher-constant}.
\end{proof}

\begin{lemma}[Every fixed quartic polynomial has the required envelope]
\label{lem:V8-polynomial-envelope}
Let $S$ be a fixed finite site set and put
\[
 R_S(h)=\sum_{x\in S}\|h_x\|^4.
\]
Let $P_n(u;h)$ be linear in $u\in E$, polynomial of total degree at most four
in the real coordinates of $(h_x)_{x\in S}$, and suppose the sum of the
absolute values of its coefficients is at most $C_P\|u\|_E$, uniformly in
$n$.  Then
\begin{equation}
 |P_n(u;h)|\le C_P\|u\|_E(1+R_S(h)).
 \label{eq:V8-polynomial-envelope}
\end{equation}
Thus \cref{thm:V8-tail-to-total-Fisher} applies with $A=B=C_P$ whenever
$R_S$ satisfies \eqref{eq:V8-uniform-QET}.
\end{lemma}

\begin{proof}
For a monomial of positive total degree $r\le4$, weighted arithmetic--geometric
mean bounds its absolute value by a convex combination of $r$th powers of the
coordinates which occur.  Since $|t|^r\le1+|t|^4$ and
$\sum_a|h_{x,a}|^4\le\|h_x\|^4$, every such monomial is at most
$1+R_S(h)$.  The same is true of the constant monomial.  Sum with the
absolute coefficient weights.
\end{proof}

\begin{corollary}[Local Higgs-action Fisher at the hard roots]
\label{cor:V8-hard-root-Fisher}
Every fixed finite bank of local Higgs potential, mass, and hopping
deformations of the pinned hard law whose coefficient norms remain bounded
has a Fisher matrix bounded uniformly in the periodic volume and uniformly
for $q$ in the full root bracket.  The conclusion includes the local
quartic-coupling score, for which the printed fourth-moment estimate alone was
insufficient.
\end{corollary}

\begin{proof}
Potential and mass derivatives have degrees four and two.  A hopping
derivative is bilinear in two neighboring fields.  On a fixed support they
satisfy \eqref{eq:V8-polynomial-envelope}.  Use
\cref{cor:V8-hard-root-tail,thm:V8-tail-to-total-Fisher} with $r_n=0$.
\end{proof}

\begin{lemma}[Bounded-density transfer]
\label{lem:V8-bounded-density-transfer}
Let $\mu_n=L_n\nu_n$, where $L_n\ge0$,
$\mathbb E_{\nu_n}L_n=1$, and $\|L_n\|_\infty\le C_{\rm den}$.  If
\[
 \mathbb E_{\nu_n}e^{\delta R_n}\le K,
\]
then
\begin{equation}
 \mathbb E_{\mu_n}e^{\delta R_n}\le C_{\rm den}K.
 \label{eq:V8-density-QET}
\end{equation}
\end{lemma}

\begin{proof}
Integrate $L_ne^{\delta R_n}\le
C_{\rm den}e^{\delta R_n}$ against $\nu_n$.
\end{proof}

% -----------------------------------------------------------------------------
\subsection{The selected quartic tail as an exact benchmark}
\label{subsec:V8-quartic-benchmark}
% -----------------------------------------------------------------------------

\begin{lemma}[Exact radial law of the selected quartic proposal]
\label{lem:V8-quartic-radial-law}
For
\[
 q_H(\dd h)=\frac2{\pi^2}e^{-\|h\|^4}\dd h
 \qquad(h\in\mathbb C^2\simeq\mathbb R^4),
\]
the random variable $T=\|H\|^4$ has the exponential law of rate one.  Hence,
for independent sites and $R_S=\sum_{x\in S}\|H_x\|^4$ with $m=|S|$,
\begin{align}
 \mathbb E e^{\delta R_S}&=(1-\delta)^{-m}
 &&(0\le\delta<1),
 \label{eq:V8-gamma-mgf}\\
 \mathbb E R_S^2&=m(m+1).
 \label{eq:V8-gamma-second}
\end{align}
\end{lemma}

\begin{proof}
The area of the unit sphere in $\mathbb R^4$ is $2\pi^2$.  In polar
coordinates,
\[
 \frac2{\pi^2}(2\pi^2)r^3e^{-r^4}\dd r
 =e^{-t}\dd t,\qquad t=r^4.
\]
Thus $T$ is exponential.  The sum of $m$ independent copies is gamma of shape
$m$ and rate one, which gives both formulas.
\end{proof}

\begin{corollary}[Explicit selected-tail Fisher constant]
\label{cor:V8-selected-tail-Fisher}
For the selected ultraviolet tail in the monograph, whose physical law obeys
\[
 \dd\mu_\infty=L_\infty\dd\nu_\infty,\qquad
 e^{-5/64}\le L_\infty\le e^{5/64},
\]
every fixed $m$-site set satisfies
\begin{equation}
 \mathbb E_{\mu_\infty}e^{\delta R_S}
 \le e^{5/64}(1-\delta)^{-m}
 \qquad(0\le\delta<1).
 \label{eq:V8-selected-tail-QET}
\end{equation}
If the bosonic derivative satisfies
$|b(u)|\le\|u\|_E(A+BR_S)$ and the line/reference row satisfies
\eqref{eq:V8-line-L2}, then its total Fisher operator obeys the sharper bound
\begin{equation}
 \boxed{
 M_\infty\preceq
 \bigl(3A^2+3B^2e^{5/64}m(m+1)+3L^2\bigr)I_E.}
 \label{eq:V8-selected-tail-Fisher}
\end{equation}
\end{corollary}

\begin{proof}
The first claim is
\cref{lem:V8-quartic-radial-law,lem:V8-bounded-density-transfer}.
For the second, repeat the proof of
\cref{thm:V8-tail-to-total-Fisher} and use
$\mathbb E_{\mu_\infty}R_S^2\le
e^{5/64}\mathbb E_{\nu_\infty}R_S^2$ together with
\eqref{eq:V8-gamma-second}.
\end{proof}

\begin{assessment}[Benchmark scope]
\label{ass:V8-selected-tail-scope}
The selected tail gives an exact noncompact test of the tail-to-Fisher
mechanism and closes its local bosonic row.  The monograph separately proves
that this particular ultraviolet tail is boundedly equivalent to a product
quartic tail and that its elementary field has a collapse/white-noise scaling
alternative.  The estimate is therefore a control benchmark, not evidence
for the desired interacting Higgs continuum.
\end{assessment}
% -----------------------------------------------------------------------------
\subsection{Line-modulus Fisher control without hiding the zero set}
\label{subsec:V8-line-score}
% -----------------------------------------------------------------------------

The bosonic tail estimate does not by itself control a logarithmic derivative
of a determinant or Pfaffian.  There are two different sufficient routes.
The first keeps the logarithmic derivative and controls an inverse.  The
second works with the section before division and lets the modulus weight
cancel the apparent singularity.

\begin{proposition}[Inverse-moment determinant and Pfaffian bounds]
\label{prop:V8-inverse-line-bound}
Let $D_\theta(x)$ be a differentiable family of invertible complex
$r\times r$ matrices and let
$\dot D(u)=\partial_uD_\theta|_{\theta=0}$.  Under any probability law $\mu$,
suppose
\begin{equation}
 \|\dot D(u)\|_{\rm HS}\le K_D\|u\|_E,
 \qquad
 \mathbb E_\mu\|D_0^{-1}\|_{\rm HS}^2\le\Xi_D.
 \label{eq:V8-inverse-moment}
\end{equation}
Then the determinant-modulus derivative satisfies
\begin{equation}
 \mathbb E_\mu
 \left|\partial_u\log|\det D_\theta|\big|_{\theta=0}\right|^2
 \le K_D^2\Xi_D\|u\|_E^2.
 \label{eq:V8-det-line-bound}
\end{equation}
If $A_\theta$ is an invertible skew matrix and the line section is
$\operatorname{Pf}A_\theta$, the right side is divided by four with
$D$ replaced by $A$.

In particular, the pointwise bounds
\[
 \sigma_{\min}(D_0)\ge m_D>0,\qquad
 \|\dot D(u)\|_{\rm HS}\le K_D\|u\|_E
\]
give $\Xi_D\le r/m_D^2$.
\end{proposition}

\begin{proof}
Jacobi's formula and the Pfaffian derivative formula give
\begin{align*}
 \partial_u\log|\det D_\theta|\big|_0
  &=\operatorname{Re}\operatorname{Tr}(D_0^{-1}\dot D(u)),\\
 \partial_u\log|\operatorname{Pf}A_\theta|\big|_0
  &=\frac12\operatorname{Re}
    \operatorname{Tr}(A_0^{-1}\dot A(u)).
\end{align*}
The Hilbert--Schmidt Cauchy--Schwarz inequality bounds the absolute trace by
the product of the two Hilbert--Schmidt norms.  Square, integrate, and apply
\eqref{eq:V8-inverse-moment}.  A singular-value floor gives
$\|D_0^{-1}\|_{\rm HS}^2\le r/m_D^2$.
\end{proof}

\begin{proposition}[Zero-safe section criterion]
\label{prop:V8-zero-safe-line}
Let $\rho$ be a positive measure, let $a\ge0$, and let
$s_\theta(x)\in\mathbb C$ be differentiable in the command direction
$u\in E$.  For $\kappa>0$ define
\begin{equation}
 \mu_\theta(\dd x)
 =Z_\theta^{-1}|s_\theta(x)|^\kappa a(x)\rho(\dd x).
 \label{eq:V8-section-law}
\end{equation}
At $\theta=0$, define the logarithmic line derivative on $\{s_0\ne0\}$ by
\[
 r(u)=\kappa\operatorname{Re}\frac{\dot s(u)}{s_0}.
\]
The zero set has $\mu_0$-measure zero after removal of any part on which the
whole density vanishes identically, and
\begin{equation}
 \boxed{
 \mathbb E_{\mu_0}|r(u)|^2
 \le\frac{\kappa^2}{Z_0}
 \int |\dot s(u)|^2|s_0|^{\kappa-2}a\,\dd\rho.}
 \label{eq:V8-zero-safe-bound}
\end{equation}
For $\kappa=2$ this becomes
\begin{equation}
 \mathbb E_{\mu_0}|r(u)|^2
 \le\frac4{Z_0}\int|\dot s(u)|^2a\,\dd\rho,
 \label{eq:V8-square-modulus-bound}
\end{equation}
with no inverse section and no spectral floor.  Therefore uniform bounds
$Z_0\ge z_*>0$ and
\[
 \int|\dot s(u)|^2a\,\dd\rho\le K_s^2\|u\|_E^2
\]
give the line constant $L^2=4K_s^2/z_*$ required in
\eqref{eq:V8-line-L2}.
\end{proposition}

\begin{proof}
The formula for $r(u)$ is the derivative of
$\log|s_\theta|^\kappa$.  Multiply its square by the density in
\eqref{eq:V8-section-law} and use
$(\operatorname{Re}z)^2\le|z|^2$:
\[
 |s_0|^\kappa |r(u)|^2
 \le\kappa^2|\dot s(u)|^2|s_0|^{\kappa-2}.
\]
Integration and division by $Z_0$ prove
\eqref{eq:V8-zero-safe-bound}; setting $\kappa=2$ proves the remaining
claims.
\end{proof}

\begin{assessment}[What the two line criteria mean]
\label{ass:V8-line-dichotomy}
A uniform inverse floor is sufficient but stronger than necessary.  A
squared-modulus determinant or Pfaffian section can have zeros and still have
finite score Fisher information because the density cancels the logarithmic
singularity exactly as in \eqref{eq:V8-square-modulus-bound}.  For
$0<\kappa<2$, the negative power in
\eqref{eq:V8-zero-safe-bound} is a real zero-set integrability condition.
For a growing global fermion matrix, rank growth, loss of a normalizer floor,
or growth of the section derivative can still destroy uniformity.  Neither
criterion orients the Pfaffian line; orientation transport remains a separate
cofinal requirement.
\end{assessment}

% -----------------------------------------------------------------------------
\subsection{Closure of the V7 Green row under checkable quartic data}
\label{subsec:V8-Green-closure}
% -----------------------------------------------------------------------------

\begin{theorem}[Coercive local-score closure theorem]
\label{thm:V8-coercive-Green-closure}
Assume the projective factor-shell, total-cylinder, and local heat-bath
hypotheses of \cref{thm:V6-RPESM-local-clock,thm:V7-RPESM-Green-closure}.
Fix a finite command bank $E$ with a Higgs support of at most $m$ sites.
Suppose, uniformly in the cutoff:
\begin{enumerate}[label=\textup{(C\arabic*)},leftmargin=3.2em]
\item a reference Higgs law has the form
      \eqref{eq:V8-quartic-action} with fixed
      $\lambda_0,Q_0,J_0,\Delta$ as in
      \eqref{eq:V8-coercive-data};
\item the physical Higgs marginal has normalized density at most
      $C_{\rm den}$ relative to that reference law;
\item the bosonic raw derivatives obey
      \[
       |b_n(u)|\le C_P\|u\|_E(1+R_{S,n});
      \]
\item the combined line/reference derivative obeys
      \eqref{eq:V8-line-L2}, for example by
      \cref{prop:V8-inverse-line-bound} or
      \cref{prop:V8-zero-safe-line};
\item the source-active heat-bath rates obey $C_n^{\rm act}\le C_*$.
\end{enumerate}
Choose $0<\alpha<\Delta^{-1}$ and
$0<\delta\le\lambda_0/(8m)$, and put
\begin{equation}
 K_{\rm phys}
 =C_{\rm den}\exp\!\left[
  \frac{\delta m}{1-\Delta\alpha}
  \left(\frac d{\lambda_0}+\frac{B_\alpha^2}{4}\right)
 \right].
 \label{eq:V8-Kphys}
\end{equation}
Then the total Fisher tower is uniformly bounded:
\begin{equation}
 \boxed{
 M_n\preceq
 \left[
 3C_P^2+\frac{12C_P^2K_{\rm phys}}{e^2\delta^2}
 +3L^2\right]I_E.}
 \label{eq:V8-full-Fisher-bound}
\end{equation}
Consequently it has an operator limit $M_\infty$, the refresh-free local
Green matrices have an operator limit $G_\infty$, and
\begin{equation}
 \boxed{
 0\preceq G_\infty-G_N
 \preceq\frac{C_*}{4}(M_\infty-M_N).}
 \label{eq:V8-Green-tail}
\end{equation}
For an auxiliary global refresh $\rho_N\downarrow0$,
\begin{equation}
 \|G_{\rho_N,N}^{[0]}-G_\infty\|
 \le\frac{\rho_N}{4}
 \left[
 3C_P^2+\frac{12C_P^2K_{\rm phys}}{e^2\delta^2}
 +3L^2\right]
 +\|G_N-G_\infty\|.
 \label{eq:V8-refreshed-tail}
\end{equation}
\end{theorem}

\begin{proof}
By \cref{thm:V8-local-quartic-exponential} the reference law has exponential
constant equal to the exponential factor in \eqref{eq:V8-Kphys}.
\Cref{lem:V8-bounded-density-transfer} supplies the factor
$C_{\rm den}$.  Apply \cref{thm:V8-tail-to-total-Fisher} with
$A=B=C_P$ to obtain \eqref{eq:V8-full-Fisher-bound}.

Projective factorization makes the Fisher increments positive, so a bounded
finite-dimensional increasing Fisher tower converges in operator norm.
Item \textup{(C5)} and
\cref{thm:V7-Fisher-Green-domination} give
\eqref{eq:V8-Green-tail}.  The V5 de-refreshing estimate with the uniform
Fisher constant gives \eqref{eq:V8-refreshed-tail}.
\end{proof}

\begin{corollary}[Two rigorous populated branches]
\label{cor:V8-populated-branches}
The hypotheses controlling the bosonic score in
\cref{thm:V8-coercive-Green-closure} are populated in each of the following
cases.
\begin{enumerate}[label=\textup{(\roman*)},leftmargin=2.8em]
\item For the pinned hard-root Higgs law,
      \cref{cor:V8-hard-root-tail,cor:V8-hard-root-Fisher} give the local
      exponential and Fisher bounds directly, uniformly through the root
      bracket.
\item For the selected product-quartic ultraviolet tail and its
      bounded-density physical perturbation,
      \cref{cor:V8-selected-tail-Fisher} gives the bound with explicit
      constants.
\end{enumerate}
On either branch the Green conclusion follows if, in addition, its actual
shells satisfy the V6 projective factorization, the line/reference row meets
one of the V8 line criteria, and the source-active rates are uniformly
bounded.
\end{corollary}

\begin{proof}
Only the cited estimates are used.  The final sentence lists the remaining
hypotheses of \cref{thm:V8-coercive-Green-closure}; it does not infer them
from quartic coercivity.
\end{proof}
\begin{firewall}[Quartic tightness is not continuum interaction]
\label{fire:V8-quartic-scope}
A local quartic-exponential bound supplies all fixed local polynomial moments
and closes a finite command's Fisher row.  It does not supply a positive
macroscopic correlation length, a transported non-Gaussian source margin, or
a renormalized product.  In the complete SM--Einstein programme it must be
combined with the interacting infrared packet, determinant/Pfaffian line
transport, represented Ward and stress identities, reflection-positive
source convergence, coframe nondegeneracy, and the Einstein residual.
\end{firewall}

% -----------------------------------------------------------------------------
\subsection{Revised upstream acquisition after the quartic estimate}
\label{subsec:V8-acquisition}
% -----------------------------------------------------------------------------

\begin{openproblem}[Executable V9 acquisition: physical density and zero-safe line]
\label{op:V8-next-acquisition}
For one printed factorized RPESM shell family, perform the following steps on
the same cutoff sequence and command bank.
\begin{enumerate}[label=\textup{(A\arabic*)},leftmargin=3.2em]
\item Write the complete normalized Higgs marginal, including compact gauge
      hopping, control sections, determinant/Pfaffian modulus, and reference
      Jacobian.
\item Extract uniform constants
      $\lambda_0,Q_0,J_0,\Delta$ for its quartic and quadratic bosonic part.
      If the remaining normalized density is not uniformly bounded, extend
      the weighted rescaling proof of
      \cref{thm:V8-local-quartic-exponential} through that precise density
      rather than replacing it by a cutoffwise $L^{1+\varepsilon}$ statement.
\item For each line factor, print its modulus exponent $\kappa$ and either
      verify the weighted section integral
      \eqref{eq:V8-zero-safe-bound} with a uniform normalizer floor or prove
      the inverse-square estimate \eqref{eq:V8-inverse-moment}.
\item Show that the actual Higgs, hopping, and represented stress derivatives
      have a fixed support and a uniform polynomial envelope.  Treat inverse
      coframe factors separately; they are not covered by
      \cref{lem:V8-polynomial-envelope}.
\item Count the actual shell-update incidence and obtain
      $C_n^{\rm act}\le C_*$.
\item Insert the resulting constants in
      \eqref{eq:V8-full-Fisher-bound} and transport one predeclared
      non-Gaussian physical source through the same projective family.
\end{enumerate}
Success closes the unbounded local total-score and Green rows on an actual
RPESM shell.  Failure returns a named obstruction: loss of quartic
coercivity, an unbounded density ratio, a nonintegrable line zero, vanishing
line normalizer, growing section derivative or rank, an inverse-coframe
singularity, growing update incidence, or loss of the interacting source.
\end{openproblem}

\begin{theorem}[Version V8 terminal theorem]
\label{thm:V8-terminal}
Preserving V1--V7 verbatim before this appended block, V8 proves:
\begin{enumerate}[label=\textup{(V8.\arabic*)},leftmargin=5.2em]
\item a uniform fourth moment alone cannot control the Fisher information of
      a quartic-coupling score;
\item bounded-degree quartic lattice actions with a uniform positive quartic
      coefficient, a uniformly bounded negative mass, and bounded quadratic
      incidence have volume-independent local quartic-exponential moments;
\item this weighted-coercivity estimate supplies the previously invoked but
      unprinted local exponential estimate for the pinned hard-root family
      throughout its critical root bracket, without a mixing hypothesis;
\item a quartic-exponential bound gives explicit moments of every order and
      an explicit total-Fisher constant for every fixed degree-four local
      command;
\item the selected law
      $q_H=(2/\pi^2)e^{-\|h\|^4}\dd h$ has
      $\|H\|^4\sim\operatorname{Exp}(1)$, giving exact local tail and moment
      constants and their bounded-density transfer;
\item a determinant/Pfaffian line score is uniformly square integrable under
      a Hilbert--Schmidt inverse-square estimate;
\item a modulus-squared line section admits a zero-safe Fisher estimate in
      terms of the undivided section derivative, so a spectral floor is
      sufficient but not necessary;
\item the resulting explicit Fisher constant, together with bounded
      source-active rates, closes the V7 refresh-free local Green telescope
      and its de-refreshing estimate.
\end{enumerate}
\end{theorem}

\begin{proof}
Item \textup{(V8.1)} is \cref{prop:V8-fourth-not-Fisher}.
Items \textup{(V8.2)}--\textup{(V8.3)} are
\cref{thm:V8-local-quartic-exponential,cor:V8-hard-root-tail}.
Item \textup{(V8.4)} is
\cref{lem:V8-exponential-export,thm:V8-tail-to-total-Fisher,
lem:V8-polynomial-envelope}.
Item \textup{(V8.5)} is
\cref{lem:V8-quartic-radial-law,cor:V8-selected-tail-Fisher}.
Items \textup{(V8.6)}--\textup{(V8.7)} are
\cref{prop:V8-inverse-line-bound,prop:V8-zero-safe-line}.
The final item is \cref{thm:V8-coercive-Green-closure}.
\end{proof}

\begin{assessment}[Programme boundary after V8]
\label{ass:V8-boundary}
V8 removes the local quartic-exponential estimate as an unproved step for the
bounded-degree pinned hard Higgs action and turns that estimate into the
uniform local Higgs Fisher bound required by V7.  It also identifies two
precise, non-equivalent ways to close the determinant/Pfaffian modulus row.

The full Standard--Model--Einstein continuum remains unproved.  For the full
RPESM law, a uniform normalized-density or direct weighted-coercivity estimate
including the actual line and reference factors is not printed.  Uniform
zero-safe line control, inverse-coframe and represented stress bounds,
critical interacting source survival, OS reconstruction, renormalized local
products, Lorentzian continuation, and the Palatini--Einstein limit remain
open.
\end{assessment}

\section*{V8 exact checks and append-only record}
\addcontentsline{toc}{section}{V8 exact checks and append-only record}
The weighted proof uses only the scalar inequalities, valid for
$0\le t\le1/4$,
\[
 (1-t)^{-1/4}-1\le\frac t2,\qquad
 (1-t)^{-1/2}-1\le\frac{t}{1-t},\qquad
 -\log(1-t)\le\frac{t}{1-t},
\]
together with graph-sphere summability and
$Br^2-r^4\le B^2/4$.  These inequalities and the constants propagated from
them were checked directly on the full interval.  Independent radial
integration gives
\[
 \int_{\mathbb C^2}\frac2{\pi^2}e^{-\|h\|^4}\dd h=1,\quad
 \mathbb E\|H\|^4=1,\quad
 \mathbb E\|H\|^8=2,
\]
confirming the exponential radial law used above.

The complete V7 source before removal of its final active document-closing
command is retained verbatim.  Its complete-file SHA--256 digest is
\begin{center}\small\ttfamily
5afc13ae1b5ad224a66429e3b2507c80740722d99e58528a36799d1fe1be2a58.
\end{center}
A Version V9 must retain the complete V8 source, remove only the final active
document-closing command, append new material, restore one final command, and
use a filename ending in \texttt{\_V9.tex}.

% =============================================================================
% END APPENDED VERSION V8 MATERIAL
% =============================================================================


% =============================================================================
% BEGIN APPENDED VERSION V9 MATERIAL
% =============================================================================

\clearpage
\section{Version V9: polynomial line weights, zero-safe shell scores, and local rank control}
\label{sec:V9-line-shell}

Version V8 proved a local quartic-exponential estimate for bounded-degree
Higgs actions and showed how such an estimate controls every fixed local
degree-four bosonic score.  Its application to a physical cylinder still used
either a bounded density ratio or a separately assumed line-score bound.

The actual RPESM control cylinder has the more specific form
\[
 \dd\Lambda_n^+
 =e^{-S_{B,n}}\|\sigma_n^+\|^2|p_n^+|^2\,\dd\lambda_n.
\]
The weight is a product of squared section moduli.  At each finite cutoff the
sections are polynomially bounded and the normalizer is positive, but the
source supplies no cutoff-independent degree, coefficient, normalizer, or
section-derivative estimate.  Positivity at every cutoff is not enough for a
uniform Fisher bound.

This version proves the missing quantitative bridge.  A polynomially growing
positive weight preserves a quartic-exponential moment if its normalized
mass has a uniform positive floor.  For a squared line section, the Fisher
integrand can be written before division by the section, and a finite
polynomial feature Gram gives both the normalizer floor and the score bound.
On a multiplicative projective line extension, the old global determinant
cancels from a detail-shell score.  Only the fixed-size detail section enters,
so growth of the total fermion matrix rank is not by itself a local Fisher
obstruction.

% -----------------------------------------------------------------------------
\subsection{Finite nonvanishing does not give a uniform line metric}
\label{subsec:V9-nonvanishing-no-Fisher}
% -----------------------------------------------------------------------------

\begin{proposition}[A nonzero section with divergent squared-modulus Fisher]
\label{prop:V9-nonzero-no-uniform}
Let $\nu$ be a symmetric probability law on $\mathbb R$ with
$0<v=\mathbb E_\nu X^2<\infty$.  For $\varepsilon>0$ define
\[
 s_{\varepsilon,\theta}(x)=\varepsilon+\theta x,\qquad
 \mu_{\varepsilon,\theta}(\dd x)
 =Z_{\varepsilon,\theta}^{-1}
  |s_{\varepsilon,\theta}(x)|^2\nu(\dd x).
\]
At $\theta=0$, the section is everywhere nonzero, $Z_{\varepsilon,0}>0$,
and $\mu_{\varepsilon,0}=\nu$, but
\begin{equation}
 \boxed{
 \operatorname{Var}_{\mu_{\varepsilon,0}}
 \left(\left.\partial_\theta
 \log|s_{\varepsilon,\theta}|^2\right|_{\theta=0}\right)
 =\frac{4v}{\varepsilon^2}.}
 \label{eq:V9-vanishing-normalizer-Fisher}
\end{equation}
Thus finite-cutoff nonvanishing and fixed polynomial degree do not imply a
uniform line-score Fisher bound.
\end{proposition}

\begin{proof}
At $\theta=0$ the weight is the constant $\varepsilon^2$, so
$Z_{\varepsilon,0}=\varepsilon^2$ and the normalized law is $\nu$.  The raw
logarithmic derivative is $2X/\varepsilon$.  Symmetry centers it, and its
variance is \eqref{eq:V9-vanishing-normalizer-Fisher}.
\end{proof}

\begin{assessment}[The missing datum in the RPESM existence statement]
\label{ass:V9-RPESM-normalizer-gap}
The RPESM statement $0<Z_n<\infty$ at each cutoff proves existence of every
finite law.  It does not exclude $Z_n\downarrow0$ in a chosen line
trivialization, nor does the phrase \emph{nonzero massive control section}
supply a printed uniform lower bound for the normalized local section mass.
The counterexample shows why that quantitative row cannot be inferred.
\end{assessment}

% -----------------------------------------------------------------------------
\subsection{Polynomial weights preserve local quartic tails}
\label{subsec:V9-polynomial-reweighting}
% -----------------------------------------------------------------------------

For $p\ge0$ and $\epsilon>0$ put
\begin{equation}
 \Psi_p(\epsilon)=
 \begin{cases}
  1,&p=0,\\[0.2em]
  2^{(p-1)_+}
  \left[1+\left(\dfrac{p}{e\epsilon}\right)^p\right],
     &p>0.
 \end{cases}
 \label{eq:V9-Psi}
\end{equation}

\begin{theorem}[Polynomial reweighting of a quartic-exponential law]
\label{thm:V9-polynomial-QET}
Let $\nu_n$ be probability laws, let $R_n\ge0$, and suppose
\begin{equation}
 \sup_n\mathbb E_{\nu_n}e^{\delta_0R_n}\le K_0.
 \label{eq:V9-reference-QET}
\end{equation}
Let $w_n\ge0$ satisfy
\begin{equation}
 w_n\le C_w(1+R_n)^p,\qquad
 z_n:=\mathbb E_{\nu_n}w_n\ge z_*>0,
 \label{eq:V9-weight-data}
\end{equation}
and define $\mu_n=z_n^{-1}w_n\nu_n$.  Then, for every
$0\le\delta<\delta_0$,
\begin{equation}
 \boxed{
 \sup_n\mathbb E_{\mu_n}e^{\delta R_n}
 \le
 \frac{C_wK_0}{z_*}
 \Psi_p(\delta_0-\delta).}
 \label{eq:V9-polynomial-QET-bound}
\end{equation}
\end{theorem}

\begin{proof}
For $p>0$,
\[
 (1+r)^p\le2^{(p-1)_+}(1+r^p),\qquad
 r^pe^{-(\delta_0-\delta)r}
 \le\left(\frac{p}{e(\delta_0-\delta)}\right)^p.
\]
Therefore
\[
 w_ne^{\delta R_n}
 \le C_w\Psi_p(\delta_0-\delta)e^{\delta_0R_n}.
\]
Integrate and divide by $z_n\ge z_*$.  For $p=0$, use
$w_n\le C_w$ directly.
\end{proof}

\begin{lemma}[Polynomial section growth in quartic radius]
\label{lem:V9-section-growth}
Let $S$ be a fixed finite site set,
$R_S(h)=\sum_{x\in S}\|h_x\|^4$, and let $s(h)$ be a scalar polynomial of
total degree at most $D$ in the real field coordinates on $S$.  If the sum of
the absolute values of its coefficients is at most $C_s$, then
\begin{equation}
 |s(h)|^2\le C_s^2(1+R_S(h))^{D/2}.
 \label{eq:V9-section-growth}
\end{equation}
Consequently a squared polynomial section has
\eqref{eq:V9-weight-data} with $p=D/2$ once its normalized section mass has a
uniform floor.
\end{lemma}

\begin{proof}
Every coordinate obeys
$|h_{x,a}|\le R_S(h)^{1/4}$ when $R_S\ge1$, while every coordinate is at most
one when $R_S\le1$.  Hence every monomial of degree at most $D$ is bounded by
$(1+R_S)^{D/4}$.  Sum with the coefficient weights and square.
\end{proof}

\begin{corollary}[Unbounded polynomial control factors fit the V8 tail theorem]
\label{cor:V9-polynomial-physical-tail}
Suppose the V8 reference Higgs law satisfies
\eqref{eq:V8-local-exponential} on $S$ at exponent $\delta_0$ with constant
$K_0$.  If the physical local factor is a squared polynomial section of
degree at most $D$, coefficient norm at most $C_s$, and section mass at least
$z_*>0$, then its normalized physical law satisfies
\eqref{eq:V9-polynomial-QET-bound} with
\[
 C_w=C_s^2,\qquad p=D/2.
\]
Thus bounded density is not required.
\end{corollary}

\begin{proof}
Combine
\cref{thm:V9-polynomial-QET,lem:V9-section-growth}.
\end{proof}
% -----------------------------------------------------------------------------
\subsection{A finite feature Gram closes the squared-section line score}
\label{subsec:V9-feature-Gram}
% -----------------------------------------------------------------------------

Let $\mathfrak m_x$ be a finite positive measure on one detail block $Y$,
depending measurably on an old boundary variable $x$.  Let
$\Phi_x:Y\to\mathbb C^r$ be a fixed finite feature bank and define
\begin{equation}
 \mathcal G_x
 =\int_Y\Phi_x(y)\Phi_x(y)^*\,\mathfrak m_x(\dd y).
 \label{eq:V9-feature-Gram}
\end{equation}
For a coefficient vector $c_x$, put
\[
 s_x(y)=c_x^*\Phi_x(y),\qquad
 z_x=\int_Y|s_x|^2\,\dd\mathfrak m_x
     =c_x^*\mathcal G_xc_x.
\]
A command direction $u\in E$ has section derivative
$\dot s_x(u;y)=d_x(u)^*\Phi_x(y)$.

\begin{theorem}[Conditional zero-safe feature-Gram bound]
\label{thm:V9-feature-line-Fisher}
Assume $z_x>0$ and define the squared-section conditional law
\begin{equation}
 K_x(\dd y)=z_x^{-1}|s_x(y)|^2\mathfrak m_x(\dd y).
 \label{eq:V9-line-kernel}
\end{equation}
On $\{s_x\ne0\}$ let
\[
 r_x(u;y)=2\operatorname{Re}\frac{\dot s_x(u;y)}{s_x(y)}.
\]
Then
\begin{equation}
 \boxed{
 \operatorname{Var}_{K_x}(r_x(u;\cdot))
 \le\mathbb E_{K_x}|r_x(u;\cdot)|^2
 \le
 4\,\frac{d_x(u)^*\mathcal G_xd_x(u)}
         {c_x^*\mathcal G_xc_x}.}
 \label{eq:V9-feature-line-bound}
\end{equation}
If, uniformly in $x$ and the cutoff,
\begin{equation}
 g_0I\preceq\mathcal G_x\preceq g_1I,\qquad
 \|c_x\|\ge c_0>0,\qquad
 \|d_x(u)\|\le c_1\|u\|_E,
 \label{eq:V9-Gram-window}
\end{equation}
then
\begin{equation}
 \boxed{
 \operatorname{Var}_{K_x}(r_x(u;\cdot))
 \le L_{\rm line}^2\|u\|_E^2,\qquad
 L_{\rm line}^2=\frac{4g_1c_1^2}{g_0c_0^2}.}
 \label{eq:V9-line-constant}
\end{equation}
No inverse of the determinant or Pfaffian coefficient occurs.
\end{theorem}

\begin{proof}
Centering decreases the second moment.  The squared modulus in
\eqref{eq:V9-line-kernel} cancels the denominator of the logarithmic
derivative:
\begin{align*}
 \mathbb E_{K_x}|r_x(u)|^2
 &\le\frac4{z_x}
 \int_Y|\dot s_x(u;y)|^2\,\mathfrak m_x(\dd y)\\
 &=4\,\frac{d_x(u)^*\mathcal G_xd_x(u)}
          {c_x^*\mathcal G_xc_x}.
\end{align*}
This is \eqref{eq:V9-feature-line-bound}.  The spectral and coefficient
bounds in \eqref{eq:V9-Gram-window} give
\[
 c_x^*\mathcal G_xc_x\ge g_0c_0^2,\qquad
 d_x(u)^*\mathcal G_xd_x(u)\le g_1c_1^2\|u\|_E^2,
\]
which proves \eqref{eq:V9-line-constant}.
\end{proof}

\begin{corollary}[Integrated zero-safe alternative]
\label{cor:V9-integrated-line}
Let $\rho_n(\dd x)$ be the old marginal and suppose only that
$z_{n,x}>0$ for $\rho_n$-almost every $x$.  The centred conditional line
score obeys
\begin{equation}
 \boxed{
 \mathbb E_{\rho_n}
 \operatorname{Var}_{K_{n,x}}(r_{n,x}(u))
 \le4\mathbb E_{\rho_n}
 \frac{d_{n,x}(u)^*\mathcal G_{n,x}d_{n,x}(u)}
      {c_{n,x}^*\mathcal G_{n,x}c_{n,x}}.}
 \label{eq:V9-integrated-line-bound}
\end{equation}
Thus a uniform bound on the expectation of the displayed quotient is
sufficient even when no pointwise normalizer floor exists.
\end{corollary}

\begin{proof}
Integrate \eqref{eq:V9-feature-line-bound} against the old marginal and use
Tonelli's theorem.
\end{proof}

\begin{lemma}[Finite positivity versus uniform positivity]
\label{lem:V9-feature-positive}
Suppose $\mathfrak m_x$ has a density which is strictly positive on a
nonempty open subset of $\mathbb R^d$, and the components of $\Phi_x$ are
linearly independent polynomials there.  Then $\mathcal G_x\succ0$ at that
fixed $x$.  This conclusion alone supplies no cutoff-independent lower bound
on $\lambda_{\min}(\mathcal G_x)$.
\end{lemma}

\begin{proof}
If $a^*\mathcal G_xa=0$, then the polynomial $a^*\Phi_x$ vanishes
$\mathfrak m_x$-almost everywhere on the open set.  A polynomial which
vanishes on an open set is identically zero, so polynomial independence gives
$a=0$.  Hence the Gram is positive definite.  A sequence of positive
definite finite matrices may have least eigenvalue tending to zero; the
one-dimensional family $\mathcal G_n=[1/n]$ is the elementary separator.
\end{proof}

\begin{assessment}[What must be acquired]
\label{ass:V9-Gram-acquisition}
Finite polynomial independence explains why every cutoff can have a positive
line normalizer.  Uniform Fisher control requires the stronger Gram window
\eqref{eq:V9-Gram-window} or the integrated quotient
\eqref{eq:V9-integrated-line-bound}.  This is the exact quantitative datum
missing from a statement of finite nonvanishing.
\end{assessment}

\begin{example}[Exact one-site quartic control section]
\label{ex:V9-linear-quartic-section}
Under
\[
 q_H(\dd h)=\frac2{\pi^2}e^{-\|h\|^4}\dd h
 \quad\text{on }\mathbb C^2,
\]
take $\Phi(h)=(1,h_1)^{\mathsf T}$ and
$s(h)=m+ah_1$.  Rotational invariance and the radial calculation in V8 give
\begin{equation}
 \mathcal G
 =\begin{pmatrix}1&0\\0&\sqrt\pi/4\end{pmatrix},
 \qquad
 z=|m|^2+\frac{\sqrt\pi}{4}|a|^2.
 \label{eq:V9-linear-Gram}
\end{equation}
For a coefficient direction $(\dot m,\dot a)$,
\begin{equation}
 \operatorname{Var}_{|s|^2q_H/z}
 \left(2\operatorname{Re}
 \frac{\dot m+\dot a h_1}{m+ah_1}\right)
 \le
 \frac{4|\dot m|^2+\sqrt\pi|\dot a|^2}
 {|m|^2+(\sqrt\pi/4)|a|^2}.
 \label{eq:V9-linear-line-bound}
\end{equation}
The section may have zeros when $a\ne0$, but the bound is finite whenever
$(m,a)\ne(0,0)$.  A uniform Euclidean lower bound on $(m,a)$ gives a uniform
line Fisher constant without a singular-value floor.
\end{example}

\begin{proof}
The off-diagonal moments vanish by phase rotation, and
$\mathbb E_{q_H}|h_1|^2=\sqrt\pi/4$.  Insert these values in
\eqref{eq:V9-feature-line-bound}.
\end{proof}
% -----------------------------------------------------------------------------
\subsection{Projective multiplication localizes the line score}
\label{subsec:V9-projective-line}
% -----------------------------------------------------------------------------

Consider one adjacent shell.  Let $x$ be the old field, $y$ the detail, and
let $\mathfrak m_{x,\theta}(\dd y)$ be a differentiable unnormalized positive
reference measure.  Let $s_{x,\theta}(y)$ be the local positive-control line
section and put
\begin{equation}
 z_{x,\theta}
 =\int|s_{x,\theta}(y)|^2\mathfrak m_{x,\theta}(\dd y),
 \qquad
 K_{x,\theta}(\dd y)
 =z_{x,\theta}^{-1}|s_{x,\theta}(y)|^2
   \mathfrak m_{x,\theta}(\dd y).
 \label{eq:V9-projective-line-kernel}
\end{equation}
Suppose the full control line is multiplicative in the chosen adjacent
trivialization:
\begin{equation}
 \Sigma_{n+1,\theta}(x,y)=\Sigma_n(x)s_{x,\theta}(y).
 \label{eq:V9-line-factorization}
\end{equation}
The detail command changes only
$\mathfrak m_{x,\theta}$ and $s_{x,\theta}$.

\begin{theorem}[Exact cancellation of the old determinant line]
\label{thm:V9-projective-line-localization}
Let
\[
 b_x(u;y)
 =\left.\partial_u
  \log\frac{\dd\mathfrak m_{x,\theta}}
                 {\dd\mathfrak m_{x,0}}(y)\right|_{\theta=0},
 \qquad
 \ell_x(u;y)
 =2\operatorname{Re}\frac{\dot s_x(u;y)}{s_x(y)}.
\]
Then the score of the normalized detail factor is exactly
\begin{equation}
 \boxed{
 \omega_x(u;y)
 =b_x(u;y)+\ell_x(u;y)
  -\mathbb E_{K_x}[b_x(u;\cdot)+\ell_x(u;\cdot)].}
 \label{eq:V9-local-line-score}
\end{equation}
No derivative of $\Sigma_n$, no inverse of the old Dirac matrix, and no old
matrix-rank factor occurs.

Assume, uniformly in the old boundary and cutoff, that for a local quartic
radius $R$,
\begin{align}
 \mathbb E_{\nu_x}e^{\delta_0R}&\le K_0,
 \label{eq:V9-shell-reference-tail}\\
 |s_x(y)|^2&\le C_s^2(1+R(y))^{D/2},
 \qquad
 \mathbb E_{\nu_x}|s_x|^2\ge z_*,
 \label{eq:V9-shell-section-data}\\
 |b_x(u;y)|&\le C_P\|u\|_E(1+R(y)),
 \label{eq:V9-shell-action-envelope}\\
 \mathbb E_{K_x}|\ell_x(u;\cdot)|^2
 &\le L_{\rm line}^2\|u\|_E^2,
 \label{eq:V9-shell-line-L2}
\end{align}
where $\nu_x$ is the normalization of $\mathfrak m_{x,0}$ and the section
mass in \eqref{eq:V9-shell-section-data} is taken under $\nu_x$.
For $0<\delta<\delta_0$, set
\begin{equation}
 K_{\rm sh}
 =\frac{C_s^2K_0}{z_*}
   \Psi_{D/2}(\delta_0-\delta).
 \label{eq:V9-shell-K}
\end{equation}
Then
\begin{equation}
 \boxed{
 \mathbb E_{K_x}e^{\delta R}\le K_{\rm sh},}
 \label{eq:V9-shell-physical-tail}
\end{equation}
and the conditional shell Fisher operator satisfies
\begin{equation}
 \boxed{
 M_x^{\rm sh}
 \preceq
 \left(
 3C_P^2+\frac{12C_P^2K_{\rm sh}}{e^2\delta^2}
 +3L_{\rm line}^2
 \right)I_E.}
 \label{eq:V9-shell-Fisher}
\end{equation}
The line estimate \eqref{eq:V9-shell-line-L2} follows from
\eqref{eq:V9-line-constant} under the feature-Gram window.
\end{theorem}

\begin{proof}
Differentiate the logarithm of
\eqref{eq:V9-projective-line-kernel}.  Differentiation under the finite
coercive integral gives
\[
 \partial_u\log z_{x,\theta}|_0
 =\mathbb E_{K_x}[b_x(u;\cdot)+\ell_x(u;\cdot)],
\]
which proves \eqref{eq:V9-local-line-score}.  The old factor
$\Sigma_n(x)$ in \eqref{eq:V9-line-factorization} is independent of the
detail command and cancels before this calculation.

Apply \cref{thm:V9-polynomial-QET} conditionally in $x$ to obtain
\eqref{eq:V9-shell-physical-tail}.  The score in
\eqref{eq:V9-local-line-score} is the orthogonal centering of
$b_x+\ell_x$.  The proof of
\cref{thm:V8-tail-to-total-Fisher}, with
$A=B=C_P$, $K=K_{\rm sh}$, and $L=L_{\rm line}$, therefore gives
\eqref{eq:V9-shell-Fisher}.  The final assertion is
\cref{thm:V9-feature-line-Fisher}.
\end{proof}

\begin{corollary}[The RPESM block factors remove global rank from a local shell score]
\label{cor:V9-RPESM-line-localization}
For the block-triangular and congruent extensions in the attached RPESM
construction,
\[
 D_{\chi,n+1}=\mathsf L_n(D_{\chi,n}\oplus E_n)\mathsf R_n,
 \qquad
 A_{n+1}^{M}=\mathsf Q_n^{\mathsf T}
             (A_n^M\oplus B_n)\mathsf Q_n,
\]
the marked sections obey
\[
 \det D_{\chi,n+1}
 =\det D_{\chi,n}\det E_n,\qquad
 \operatorname{Pf}A_{n+1}^{M}
 =\operatorname{Pf}A_n^M\operatorname{Pf}B_n.
\]
If the positive control sections are chosen in this same multiplicative
trivialization, their local combined detail section is
\begin{equation}
 s_n(x,y)=\det E_n(x,y)\operatorname{Pf}B_n(x,y).
 \label{eq:V9-RPESM-detail-section}
\end{equation}
A command acting on the new detail therefore has the localized score
\eqref{eq:V9-local-line-score}; the ranks of $D_{\chi,n}$ and $A_n^M$ do not
enter its line bound.

If the detail blocks have uniformly bounded size and their section and
command-derivative polynomial coefficient norms are uniformly bounded, and
their feature Grams satisfy
\eqref{eq:V9-Gram-window} or the integrated quotient
\eqref{eq:V9-integrated-line-bound}, the line part of every fixed local shell
score has a cutoff-independent Fisher bound.
\end{corollary}

\begin{proof}
Determinant-one triangular multiplication and Pfaffian congruence give the two
displayed product identities.  Under the stated compatible choice of positive
control sections, their product is
\eqref{eq:V9-line-factorization} with the detail factor
\eqref{eq:V9-RPESM-detail-section}.  Apply
\cref{thm:V9-projective-line-localization}.  Fixed block size and the two coefficient bounds put the section and its
derivative in a fixed polynomial feature bank;
then \cref{thm:V9-feature-line-Fisher,cor:V9-integrated-line} gives the last
claim.
\end{proof}

\begin{assessment}[Imported transport versus the new estimate]
\label{ass:V9-imported-new}
The determinant/Pfaffian product identities, divisor transport, and exact
projective kernel are imported from the attached RPESM construction.  V9 does
not reprove them.  The new content is the cancellation
\eqref{eq:V9-local-line-score}, the polynomial-tail constant
\eqref{eq:V9-shell-K}, and the zero-safe feature-Gram Fisher bound.  The source
does not print the uniform fixed-block, coefficient, feature-Gram, or
integrated-quotient data needed to apply the last sentence of
\cref{cor:V9-RPESM-line-localization} to its full selected sequence.
\end{assessment}
% -----------------------------------------------------------------------------
\subsection{A summable local-line criterion for the Green telescope}
\label{subsec:V9-summable-line}
% -----------------------------------------------------------------------------

\begin{theorem}[Square-summable shell coefficients close total Fisher and Green]
\label{thm:V9-summable-line-Green}
Assume the projective factor-shell hypotheses of V6--V7.  Suppose every shell
satisfies \cref{thm:V9-projective-line-localization} with common
$\delta,K_{\rm sh},g_0,g_1,c_0$, but allow shell-dependent action and section
derivative constants
\begin{equation}
 |b_{n,x}(u;y)|
 \le C_{P,n}\|u\|_E(1+R(y)),\qquad
 \|d_{n,x}(u)\|\le c_{1,n}\|u\|_E.
 \label{eq:V9-shell-amplitudes}
\end{equation}
If
\begin{equation}
 \sum_{n\ge0}C_{P,n}^2<\infty,\qquad
 \sum_{n\ge0}c_{1,n}^2<\infty,
 \label{eq:V9-square-summability}
\end{equation}
then the total Fisher tower is bounded and converges.  More precisely, with
\begin{equation}
 A_{\rm Q}=1+\frac{4K_{\rm sh}}{e^2\delta^2},
 \qquad
 A_{\rm L}=\frac{4g_1}{g_0c_0^2},
 \label{eq:V9-tail-constants}
\end{equation}
its tail satisfies
\begin{equation}
 \boxed{
 0\preceq M_\infty-M_N
 \preceq
 3\left(
 A_{\rm Q}\sum_{n\ge N}C_{P,n}^2
 +A_{\rm L}\sum_{n\ge N}c_{1,n}^2
 \right)I_E.}
 \label{eq:V9-Fisher-tail}
\end{equation}
If also $C_n^{\rm act}\le C_*$, then
\begin{equation}
 \boxed{
 0\preceq G_\infty-G_N
 \preceq\frac{3C_*}{4}\left(
 A_{\rm Q}\sum_{n\ge N}C_{P,n}^2
 +A_{\rm L}\sum_{n\ge N}c_{1,n}^2
 \right)I_E.}
 \label{eq:V9-Green-tail}
\end{equation}
\end{theorem}

\begin{proof}
The feature-Gram estimate gives
$L_{{\rm line},n}^2\le A_{\rm L}c_{1,n}^2$.  Equation
\eqref{eq:V9-shell-Fisher} therefore gives
\[
 \Delta M_n
 \preceq3\bigl(A_{\rm Q}C_{P,n}^2
              +A_{\rm L}c_{1,n}^2\bigr)I_E.
\]
The V7 conditional orthogonality theorem makes the total Fisher operator the
positive sum of its shell increments.  Sum the displayed estimate to prove
\eqref{eq:V9-Fisher-tail}.  The V7 source-active domination
$\Delta G_n\preceq C_n^{\rm act}\Delta M_n/4$ and
$C_n^{\rm act}\le C_*$ give \eqref{eq:V9-Green-tail}.
\end{proof}

\begin{assessment}[The remaining quantitative RPESM test]
\label{ass:V9-RPESM-test}
For the RPESM line extension, fixed detail-block size removes the growing
global rank, but it does not imply
\eqref{eq:V9-square-summability}.  The actual command normalization must make
the local bosonic and section derivatives square summable, or an independent
bounded total-Fisher estimate must replace that condition.  No such
coefficient schedule or feature-Gram window is printed for the complete
physical cylinder.
\end{assessment}

% -----------------------------------------------------------------------------
\subsection{An exact noncompact zero-safe projective line tower}
\label{subsec:V9-line-control-tower}
% -----------------------------------------------------------------------------

Let $q_H$ be the quartic law of
\cref{lem:V8-quartic-radial-law}, set $c_H=\sqrt\pi/4$, and fix
$m,a\in\mathbb C$ with $m\ne0$.  For a real command parameter $\theta$ and a
real sequence $(\gamma_k)$ define
\begin{equation}
 s_{k,\theta}(h)
 =m+(a+\theta\gamma_k)h_1,\qquad
 \nu_{k,\theta}(\dd h)
 =z_{k,\theta}^{-1}|s_{k,\theta}(h)|^2q_H(\dd h),
 \label{eq:V9-control-factor}
\end{equation}
where
\begin{equation}
 z_{k,\theta}
 =|m|^2+c_H|a+\theta\gamma_k|^2.
 \label{eq:V9-control-normalizer}
\end{equation}
Put $\nu_{N,\theta}=\bigotimes_{k=1}^N\nu_{k,\theta}$.  Let $E_{k,N}$ be
conditional expectation over the $k$th coordinate and take the positive
heat-bath Hamiltonian
\begin{equation}
 H_N=\sum_{k=1}^Nd_k(I-E_{k,N}),\qquad d_k>0.
 \label{eq:V9-control-H}
\end{equation}

\begin{theorem}[Zero-safe line tower with finite Fisher and Green limits]
\label{thm:V9-zero-safe-line-tower}
Let
\begin{equation}
 \eta_k(h_k)
 =\left.\partial_\theta\log\nu_{k,\theta}(h_k)
  \right|_{\theta=0}.
 \label{eq:V9-control-score}
\end{equation}
Then the finite laws and heat baths are exactly projective, the score is the
orthogonal sum
\begin{equation}
 w_N=\sum_{k=1}^N\eta_k,\qquad
 \langle\eta_j,\eta_k\rangle_{\nu_N}=0\quad(j\ne k),
 \label{eq:V9-control-orthogonal}
\end{equation}
and
\begin{equation}
 H_N\eta_k=d_k\eta_k.
 \label{eq:V9-control-eigen}
\end{equation}
Each line-score increment obeys
\begin{equation}
 0<\|\eta_k\|_{L^2(\nu_k)}^2
 \le\frac{\sqrt\pi}{|m|^2}\gamma_k^2
 \quad\text{when }\gamma_k\ne0.
 \label{eq:V9-control-Fisher-bound}
\end{equation}
Consequently, if $(\gamma_k)\in\ell^2$ and
$\sup_kd_k\le d_*$, the exact Fisher and local Green limits satisfy
\begin{align}
 M_\infty
 &=\sum_{k\ge1}\|\eta_k\|^2
 \le\frac{\sqrt\pi}{|m|^2}\sum_{k\ge1}\gamma_k^2,
 \label{eq:V9-control-M}\\
 G_\infty
 &=\frac14\sum_{k\ge1}d_k\|\eta_k\|^2
 \le\frac{d_*\sqrt\pi}{4|m|^2}
       \sum_{k\ge1}\gamma_k^2.
 \label{eq:V9-control-G}
\end{align}
If $a\ne0$, every base section has a zero at $h_1=-m/a$, so no pointwise inverse
floor exists.  The Fisher and Green limits nevertheless remain finite.
\end{theorem}

\begin{proof}
Equation \eqref{eq:V9-control-normalizer} follows from
$\mathbb E_{q_H}h_1=0$ and
$\mathbb E_{q_H}|h_1|^2=c_H$.  Product marginalization and coordinate
heat-bath intertwining prove exact projectivity.  Differentiating the
normalized factor gives a centered function of $h_k$ alone, so independence
proves \eqref{eq:V9-control-orthogonal}.  Conditional expectation over its
own coordinate kills $\eta_k$, while every other heat bath fixes it; this is
\eqref{eq:V9-control-eigen}.

Before centering, the logarithmic line derivative is
\[
 \ell_k(h)=2\operatorname{Re}
 \frac{\gamma_kh_1}{m+ah_1}.
\]
The zero-safe cancellation gives
\begin{align*}
 \mathbb E_{\nu_k}|\ell_k|^2
 &\le\frac{4\gamma_k^2}{z_{k,0}}
       \mathbb E_{q_H}|h_1|^2\\
 &\le\frac{4c_H}{|m|^2}\gamma_k^2
 =\frac{\sqrt\pi}{|m|^2}\gamma_k^2.
\end{align*}
Centering decreases the second moment.  If $\gamma_k\ne0$, the logarithmic
derivative is not almost surely constant, so its variance is positive.
The Fisher equality is orthogonal Pythagoras.  The V5 identity
$G_N=W_N^*H_NW_N/4$, together with
\eqref{eq:V9-control-eigen}, gives the Green equality.  Sum the bounds.
Finally, the section zero for $a\ne0$ is immediate, and it has
$q_H$-measure zero.
\end{proof}

\begin{assessment}[Control-tower scope]
\label{ass:V9-control-scope}
This tower is noncompact, line weighted, projective, and genuinely zero safe.
It proves that a pointwise inverse floor is not necessary for a cofinal
line-score response.  Its factors are independent and contain neither
gauge--Higgs interaction, chiral phase transport, a real Pfaffian
orientation, represented stress, nor Einstein dynamics.  It is a sharp
control for the analytic mechanism, not the physical continuum.
\end{assessment}
% -----------------------------------------------------------------------------
\subsection{Revised continuum acquisition}
\label{subsec:V9-acquisition}
% -----------------------------------------------------------------------------

\begin{openproblem}[Executable V10 acquisition: integrated feature-Gram control]
\label{op:V9-next-acquisition}
For the actual RPESM detail blocks in one selected projective sequence:
\begin{enumerate}[label=\textup{(A\arabic*)},leftmargin=3.2em]
\item Print the conditional reference measure
      $\mathfrak m_{n,x}$ and the multiplicative positive detail section
      $s_{n,x}=\det E_n\,\operatorname{Pf}B_n$ in a fixed polynomial feature
      bank.
\item Verify a cutoff-independent detail-block dimension and polynomial
      degree, and bound the section coefficient norm.
\item Prove either the pointwise feature-Gram window
      \eqref{eq:V9-Gram-window} or the weaker integrated quotient
      \eqref{eq:V9-integrated-line-bound}.  If the old boundary produces an
      unbounded linear field, derive the dependence of the inverse Gram on
      that field and integrate it using the V8 quartic-exponential estimate.
\item Obtain a conditional section-mass floor strong enough for the
      polynomial tail transfer \eqref{eq:V9-shell-physical-tail}, or prove a
      direct weighted-tail estimate which includes the squared section.
\item Compute the physical shell normalization of the bosonic and section
      derivatives and test \eqref{eq:V9-square-summability}.
\item Retain the complex phase connection, real Pfaffian orientation, divisor
      orders, and crossing signs as separate line data; the positive
      squared-modulus estimate does not decide them.
\item Insert the resulting Fisher tail in \eqref{eq:V9-Green-tail}, then test
      the same sequence for a transported non-Gaussian infrared source and
      the represented stress/coframe envelope.
\end{enumerate}
A pass closes the positive-modulus Higgs--fermion local Fisher and Green rows
on the actual projective cylinder.  A failure returns a precise fixed-block,
feature-Gram, section-mass, coefficient-summability, phase/orientation,
infrared-source, or stress/coframe obstruction.
\end{openproblem}

\begin{theorem}[Version V9 terminal theorem]
\label{thm:V9-terminal}
Preserving V1--V8 verbatim before this appended block, V9 proves:
\begin{enumerate}[label=\textup{(V9.\arabic*)},leftmargin=5.2em]
\item finite-cutoff nonvanishing of a squared line section does not imply a
      uniform line-score Fisher bound;
\item a positive polynomial weight with a uniform normalized-mass floor
      preserves a local quartic-exponential moment, with the explicit
      constant \eqref{eq:V9-polynomial-QET-bound};
\item a degree-$D$ squared polynomial section has growth exponent $D/2$ in
      the quartic radius;
\item the squared-section logarithmic score has the zero-safe feature-Gram
      bound \eqref{eq:V9-feature-line-bound}, and either a uniform Gram window
      or an integrated Gram quotient closes its Fisher row without an inverse
      matrix;
\item the exact linear control section over $q_H$ has an explicit two-feature
      Gram and remains Fisher controlled even when the section has zeros;
\item multiplicative projective line transport cancels the entire old
      determinant/Pfaffian section from a detail-shell score;
\item under fixed local block data, the RPESM determinant/Pfaffian product
      identities reduce the shell line bound to
      $\det E_n\,\operatorname{Pf}B_n$, independently of global matrix rank;
\item square-summable local action and section-derivative coefficients give
      explicit total-Fisher and Green tails;
\item an exact noncompact projective squared-line tower has finite Fisher and
      Green limits despite having no pointwise inverse floor.
\end{enumerate}
\end{theorem}

\begin{proof}
Item \textup{(V9.1)} is \cref{prop:V9-nonzero-no-uniform}.
Items \textup{(V9.2)}--\textup{(V9.3)} are
\cref{thm:V9-polynomial-QET,lem:V9-section-growth}.
Items \textup{(V9.4)}--\textup{(V9.5)} are
\cref{thm:V9-feature-line-Fisher,cor:V9-integrated-line,
ex:V9-linear-quartic-section}.
Items \textup{(V9.6)}--\textup{(V9.7)} are
\cref{thm:V9-projective-line-localization,
cor:V9-RPESM-line-localization}.
Item \textup{(V9.8)} is \cref{thm:V9-summable-line-Green}.
The final item is \cref{thm:V9-zero-safe-line-tower}.
\end{proof}

\begin{assessment}[Programme boundary after V9]
\label{ass:V9-boundary}
V9 removes two false bottlenecks.  An unbounded polynomial line weight need
not be replaced by a bounded density, and growth of the full determinant
matrix rank need not enter a local projective shell score.  The correct
positive-modulus data are a local polynomial degree and coefficient bound, a
section-mass or integrated feature-Gram estimate, and square-summable
physical shell derivatives.

The full Standard--Model--Einstein continuum remains unproved.  The attached
RPESM construction asserts finite polynomial growth, finite normalizers,
exact line products, and finite population, but does not print the uniform
local feature-Gram, section-mass, or derivative-summability constants required
above.  The chiral phase, Pfaffian orientation, interacting infrared source,
inverse-coframe and represented-stress control, OS continuum, local products,
Lorentzian reconstruction, and Palatini--Einstein limit also remain open.
\end{assessment}

\section*{V9 exact checks and append-only record}
\addcontentsline{toc}{section}{V9 exact checks and append-only record}
For $p>0$ and $\epsilon>0$, direct maximization gives
\[
 \sup_{r\ge0}r^pe^{-\epsilon r}
 =\left(\frac{p}{e\epsilon}\right)^p,
\]
which verifies the only nonalgebraic constant in
\eqref{eq:V9-polynomial-QET-bound}.  The feature-Gram identity was checked by
direct expansion:
\[
 \int|c^*\Phi|^2\,\dd\mathfrak m=c^*\mathcal Gc,\qquad
 \int|d^*\Phi|^2\,\dd\mathfrak m=d^*\mathcal Gd.
\]
For the exact control tower, radial integration gives
$\mathbb E h_1=0$ and
$\mathbb E|h_1|^2=\sqrt\pi/4$, hence
\[
 \int|m+ah_1|^2q_H(\dd h)
 =|m|^2+\frac{\sqrt\pi}{4}|a|^2.
\]
Finite product enumerations of the conditional-expectation algebra verify
score centering, cross-factor orthogonality, generator projectivity, and
$H_N\eta_k=d_k\eta_k$.  Direct numerical sampling was used only as a
consistency check; none of the proofs depends on it.

The complete V8 source before removal of its final active document-closing
command is retained verbatim.  Its complete-file SHA--256 digest is
\begin{center}\small\ttfamily
3d843d5341e95cd62c91f6bc6995d7ed89a7c076db455be4d6cd8fb4bd7c78ad.
\end{center}
A Version V10 must retain the complete V9 source, remove only the final active
document-closing command, append new material, restore one final command, and
use a filename ending in \texttt{\_V10.tex}.

% =============================================================================
% END APPENDED VERSION V9 MATERIAL
% =============================================================================


% =============================================================================
% BEGIN APPENDED VERSION V10 MATERIAL
% =============================================================================

\clearpage
\section{Version V10: boundary-dependent feature Grams and integrated zero-safe closure}
\label{sec:V10-integrated-Gram}

Version V9 reduced the positive determinant/Pfaffian modulus to a fixed local
detail section and identified a finite feature-Gram quotient as the remaining
line-score estimate.  Requiring that Gram to have a boundary-independent
pointwise lower eigenvalue is still too strong for an unbounded Higgs
boundary.  A large linear boundary field can concentrate a one-block quartic
law and make its polynomial Gram poorly conditioned.

The correct loss has the scale $\exp(C|b|^{4/3})$, the Legendre scale dual to
a quartic potential.  This version proves explicit upper and lower feature-
Gram bounds at that scale.  It then proves that the V8 local quartic-
exponential estimate integrates the loss whenever the boundary field is a
bounded-incidence linear function of neighboring Higgs variables.  The same
argument transfers a local quartic tail through the squared polynomial line
section without a pointwise conditional normalizer floor.

This closes the analytic feature-Gram step for any printed fixed-size detail
block satisfying the stated uniform polynomial nondegeneracy and coefficient
conditions.  Whether the abstract RPESM population selects such blocks and a
square-summable physical command schedule remains a separate acquisition.

% -----------------------------------------------------------------------------
\subsection{Quartic conditional laws with an unbounded linear boundary}
\label{subsec:V10-conditional-quartic}
% -----------------------------------------------------------------------------

Let $\mathscr K$ be a compact parameter set.  For
$\xi\in\mathscr K$, $b\in\mathbb R^d$, and $y\in\mathbb R^d$, consider
\begin{equation}
 \nu_{\xi,b}(\dd y)
 =Z_{\xi,b}^{-1}
  e^{-V_\xi(y)+\langle b,y\rangle}\dd y.
 \label{eq:V10-conditional-law}
\end{equation}
Assume that for fixed $\lambda_0>0$, $Q_0,C_0\ge0$,
\begin{equation}
 V_\xi(y)\ge
 \lambda_0\|y\|^4-Q_0\|y\|^2-C_0
 \quad(\xi\in\mathscr K,\ y\in\mathbb R^d),
 \label{eq:V10-potential-lower}
\end{equation}
and, for some fixed $R>0$,
\begin{equation}
 V_R:=\sup_{\xi\in\mathscr K,\ \|y\|\le R}V_\xi(y)<\infty.
 \label{eq:V10-ball-upper}
\end{equation}
Let $\Phi_\xi:\mathbb R^d\to\mathbb C^r$ be a polynomial feature bank with
\begin{align}
 \|\Phi_\xi(y)\|^2
 &\le C_\Phi(1+\|y\|)^{2D},
 \label{eq:V10-feature-growth}\\
 \int_{\|y\|\le R}
 \Phi_\xi(y)\Phi_\xi(y)^*\,\dd y
 &\succeq\kappa_RI_r
 \quad(\kappa_R>0).
 \label{eq:V10-ball-Gram}
\end{align}
Write
\begin{equation}
 \mathcal G_{\xi,b}
 =\mathbb E_{\nu_{\xi,b}}
   [\Phi_\xi\Phi_\xi^*].
 \label{eq:V10-conditional-Gram}
\end{equation}
Set
\begin{align}
 I_0&=\int_{\mathbb R^d}
       e^{-(\lambda_0/2)\|y\|^4}\dd y,
 \label{eq:V10-I0}\\
 I_D&=\int_{\mathbb R^d}
       (1+\|y\|)^{2D}
       e^{-(\lambda_0/2)\|y\|^4}\dd y,
 \label{eq:V10-ID}\\
 \mathcal A(B)
 &=C_0+\frac{Q_0^2}{\lambda_0}
   +\frac34\lambda_0^{-1/3}B^{4/3}
   +V_R+BR.
 \label{eq:V10-A}
\end{align}

\begin{theorem}[Explicit boundary-dependent feature-Gram window]
\label{thm:V10-boundary-Gram}
For every $\xi\in\mathscr K$ and $b\in\mathbb R^d$,
\begin{equation}
 \boxed{
 \frac{\kappa_R}{I_0}e^{-\mathcal A(\|b\|)}I_r
 \preceq\mathcal G_{\xi,b}
 \preceq
 \frac{C_\Phi I_D}{|B_R|}
 e^{\mathcal A(\|b\|)}I_r.}
 \label{eq:V10-Gram-window}
\end{equation}
Consequently,
\begin{equation}
 \boxed{
 \frac{\lambda_{\max}(\mathcal G_{\xi,b})}
      {\lambda_{\min}(\mathcal G_{\xi,b})}
 \le
 \frac{C_\Phi I_DI_0}{|B_R|\kappa_R}
 e^{2\mathcal A(\|b\|)}.}
 \label{eq:V10-Gram-condition}
\end{equation}
Here $|B_R|$ is the Euclidean volume of the radius-$R$ ball.  All constants
are independent of the finite lattice volume.
\end{theorem}

\begin{proof}
Put $B=\|b\|$ and $r=\|y\|$.  The two scalar Young bounds
\begin{align}
 Q_0r^2&\le\frac{\lambda_0}{4}r^4+\frac{Q_0^2}{\lambda_0},
 \label{eq:V10-Young-quadratic}\\
 Br&\le\frac{\lambda_0}{4}r^4
      +\frac34\lambda_0^{-1/3}B^{4/3}
 \label{eq:V10-Young-linear}
\end{align}
and \eqref{eq:V10-potential-lower} give
\begin{equation}
 Z_{\xi,b}
 \le I_0\exp\!\left(
 C_0+\frac{Q_0^2}{\lambda_0}
 +\frac34\lambda_0^{-1/3}B^{4/3}\right).
 \label{eq:V10-Z-upper}
\end{equation}
On the ball $\|y\|\le R$,
$-V_\xi(y)+\langle b,y\rangle\ge-V_R-BR$, hence
\begin{equation}
 Z_{\xi,b}\ge |B_R|e^{-V_R-BR}.
 \label{eq:V10-Z-lower}
\end{equation}

For any $a\in\mathbb C^r$, restrict the numerator of
$a^*\mathcal G_{\xi,b}a$ to the radius-$R$ ball.  Equations
\eqref{eq:V10-ball-Gram,eq:V10-Z-upper} give
\[
 a^*\mathcal G_{\xi,b}a
 \ge\frac{\kappa_R}{I_0}e^{-\mathcal A(B)}\|a\|^2.
\]
For the upper bound, use
\eqref{eq:V10-feature-growth} and the same Young estimates in the numerator,
then divide by \eqref{eq:V10-Z-lower}:
\[
 a^*\mathcal G_{\xi,b}a
 \le\frac{C_\Phi I_D}{|B_R|}
      e^{\mathcal A(B)}\|a\|^2.
\]
This proves \eqref{eq:V10-Gram-window}; division of the two scalar bounds gives
\eqref{eq:V10-Gram-condition}.
\end{proof}

\begin{lemma}[Compact polynomial independence supplies the ball Gram]
\label{lem:V10-compact-ball-Gram}
Suppose the coefficients of $\Phi_\xi$ depend continuously on
$\xi\in\mathscr K$ and its components are linearly independent as
polynomials for every $\xi$.  Then
\[
 \inf_{\xi\in\mathscr K}
 \lambda_{\min}\!\left(
 \int_{\|y\|\le R}\Phi_\xi\Phi_\xi^*\,\dd y\right)>0.
\]
Thus \eqref{eq:V10-ball-Gram} follows from a finite algebraic independence
test plus compactness.
\end{lemma}

\begin{proof}
At fixed $\xi$, a nonzero linear combination of the features is a nonzero
polynomial and cannot vanish on the ball.  Its squared integral is therefore
positive.  The Gram entries and their least eigenvalue depend continuously on
$\xi$.  A positive continuous function on compact $\mathscr K$ has a positive
minimum.
\end{proof}

\begin{lemma}[Quartic tails integrate every quartic-dual boundary loss]
\label{lem:V10-boundary-integrability}
Let $X_1,\ldots,X_m$ be random vectors and suppose
\begin{equation}
 \mathbb E\exp\!\left(
 \delta\sum_{j=1}^m\|X_j\|^4\right)\le K
 \qquad(\delta,K>0).
 \label{eq:V10-neighbour-QET}
\end{equation}
If
\begin{equation}
 B\le J\sum_{j=1}^m\|X_j\|,
 \label{eq:V10-boundary-incidence}
\end{equation}
then, for every $a,p\ge0$,
\begin{equation}
 \boxed{
 \mathbb E(1+B)^p e^{aB^{4/3}}<\infty,}
 \label{eq:V10-dual-integrability}
\end{equation}
with a bound depending only on
$\delta,K,J,m,a,p$.
\end{lemma}

\begin{proof}
First,
\[
 B^{4/3}
 \le J^{4/3}m^{1/3}
     \sum_{j=1}^m\|X_j\|^{4/3}.
\]
For every $c,\epsilon>0$, direct maximization gives
\begin{equation}
 cr^{4/3}
 \le\epsilon r^4+
 \frac{2c^{3/2}}{3\sqrt{3\epsilon}}
 \qquad(r\ge0).
 \label{eq:V10-dual-Young}
\end{equation}
Also $(1+B)^p\le C_{p,\eta}e^{\eta B^{4/3}}$ for every $\eta>0$, because
the quotient has finite supremum on $[0,\infty)$.  Apply
\eqref{eq:V10-dual-Young} to every neighbor with $\epsilon=\delta$ after
combining the coefficients $a+\eta$, then use
\eqref{eq:V10-neighbour-QET}.
\end{proof}
% -----------------------------------------------------------------------------
\subsection{Integrated line Fisher without a pointwise section-mass floor}
\label{subsec:V10-integrated-line}
% -----------------------------------------------------------------------------

Let a local section and its command derivative have the feature
representations
\begin{equation}
 s_{\xi,b}(y)=c_{\xi,b}^*\Phi_\xi(y),\qquad
 \dot s_{\xi,b}(u;y)=d_{\xi,b}(u)^*\Phi_\xi(y).
 \label{eq:V10-section-features}
\end{equation}

\begin{theorem}[Boundary-integrated zero-safe line estimate]
\label{thm:V10-integrated-line}
Assume the hypotheses of \cref{thm:V10-boundary-Gram}.  Suppose, for constants
$c_0,c_1>0$ and exponents $q,m\ge0$,
\begin{equation}
 \|c_{\xi,b}\|\ge c_0(1+\|b\|)^{-q},\qquad
 \|d_{\xi,b}(u)\|
 \le c_1(1+\|b\|)^m\|u\|_E.
 \label{eq:V10-coefficient-window}
\end{equation}
Let
\[
 K_{\xi,b}(\dd y)
 =\frac{|s_{\xi,b}(y)|^2\nu_{\xi,b}(\dd y)}
        {\mathbb E_{\nu_{\xi,b}}|s_{\xi,b}|^2}
\]
and let $\ell_{\xi,b}(u;y)
=2\operatorname{Re}(\dot s_{\xi,b}(u;y)/s_{\xi,b}(y))$.
Then
\begin{equation}
 \boxed{
 \mathbb E_{K_{\xi,b}}|\ell_{\xi,b}(u)|^2
 \le
 C_{\rm line}
 (1+\|b\|)^{2(m+q)}
 e^{2\mathcal A(\|b\|)}\|u\|_E^2,}
 \label{eq:V10-pointwise-line}
\end{equation}
where
\begin{equation}
 C_{\rm line}
 =\frac{4c_1^2C_\Phi I_DI_0}
        {c_0^2|B_R|\kappa_R}.
 \label{eq:V10-Cline}
\end{equation}
If the boundary is random, obeys \eqref{eq:V10-boundary-incidence}, and its
neighbor fields satisfy \eqref{eq:V10-neighbour-QET} uniformly in the cutoff,
then
\begin{equation}
 \boxed{
 \sup_n\mathbb E
 |\ell_{\xi_n,B_n}(u)|^2
 \le\overline L_{\rm line}^2\|u\|_E^2<\infty.}
 \label{eq:V10-integrated-line-final}
\end{equation}
Thus the integrated V9 feature-Gram quotient is finite even though its
pointwise denominator may tend to zero as $\|b\|\to\infty$.
\end{theorem}

\begin{proof}
The zero-safe feature estimate \eqref{eq:V9-feature-line-bound} and
\eqref{eq:V10-coefficient-window} give
\[
 \mathbb E_{K_{\xi,b}}|\ell_{\xi,b}(u)|^2
 \le\frac{4c_1^2}{c_0^2}(1+\|b\|)^{2(m+q)}
 \frac{\lambda_{\max}(\mathcal G_{\xi,b})}
      {\lambda_{\min}(\mathcal G_{\xi,b})}\|u\|_E^2.
\]
Insert \eqref{eq:V10-Gram-condition} to obtain
\eqref{eq:V10-pointwise-line}.  Since
\[
 \mathcal A(B)\le C_{\mathcal A}(1+B^{4/3})
\]
for a constant depending only on the structural data,
\cref{lem:V10-boundary-integrability} makes the right side integrable with a
cutoff-independent bound.  This proves
\eqref{eq:V10-integrated-line-final}.
\end{proof}

% -----------------------------------------------------------------------------
\subsection{The squared section also preserves the detail-field tail}
\label{subsec:V10-section-tail}
% -----------------------------------------------------------------------------

For $0<\delta_1<\lambda_0$, put
\begin{equation}
 \lambda_1=\lambda_0-\delta_1,\qquad
 I_D(\lambda_1)
 =\int_{\mathbb R^d}(1+\|y\|)^{2D}
   e^{-(\lambda_1/2)\|y\|^4}\dd y.
 \label{eq:V10-ID-lambda}
\end{equation}

\begin{theorem}[Conditional squared-section quartic tail]
\label{thm:V10-section-tail}
Assume \eqref{eq:V10-potential-lower},
\eqref{eq:V10-ball-upper}, and
\eqref{eq:V10-ball-Gram}.  In addition to the lower coefficient estimate in
\eqref{eq:V10-coefficient-window}, suppose
\begin{equation}
 \|c_{\xi,b}\|\le c_2(1+\|b\|)^{q_2}
 \label{eq:V10-coefficient-upper}
\end{equation}
for $c_2>0$ and $q_2\ge0$.  Then the squared-section law satisfies
\begin{equation}
 \boxed{
 \mathbb E_{K_{\xi,b}}e^{\delta_1\|Y\|^4}
 \le C_{\rm tail}(1+\|b\|)^{2(q+q_2)}
 \exp\!\left[
  C_0+\frac{Q_0^2}{\lambda_1}
  +\frac34\lambda_1^{-1/3}\|b\|^{4/3}
  +V_R+R\|b\|\right],}
 \label{eq:V10-conditional-tail}
\end{equation}
where
\begin{equation}
 C_{\rm tail}
 =\frac{c_2^2C_\Phi I_D(\lambda_1)}
        {c_0^2\kappa_R}.
 \label{eq:V10-Ctail}
\end{equation}
If $b$ is a bounded-incidence linear boundary field whose neighbors satisfy
\eqref{eq:V10-neighbour-QET}, then
\begin{equation}
 \boxed{
 \sup_n\mathbb E
 e^{\delta_1\|Y_n\|^4}<\infty}
 \label{eq:V10-joint-physical-tail}
\end{equation}
under the joint old--detail squared-section law.  No pointwise lower bound on
$\mathbb E_{\nu_{\xi,b}}|s_{\xi,b}|^2$ is assumed.
\end{theorem}

\begin{proof}
The unnormalized denominator of $K_{\xi,b}$ is
\[
 \int|c_{\xi,b}^*\Phi_\xi(y)|^2
 e^{-V_\xi(y)+\langle b,y\rangle}\dd y.
\]
Restricting to $\|y\|\le R$ and using
\eqref{eq:V10-ball-upper,eq:V10-ball-Gram} and the lower coefficient bound
gives
\begin{equation}
 \int|s_{\xi,b}|^2e^{-V_\xi+\langle b,y\rangle}\dd y
 \ge c_0^2(1+\|b\|)^{-2q}
     \kappa_Re^{-V_R-R\|b\|}.
 \label{eq:V10-section-denominator}
\end{equation}
For the numerator, the feature growth and upper coefficient bound give
\[
 |s_{\xi,b}(y)|^2
 \le c_2^2C_\Phi(1+\|b\|)^{2q_2}
     (1+\|y\|)^{2D}.
\]
Apply \eqref{eq:V10-Young-quadratic,eq:V10-Young-linear} with
$\lambda_0$ replaced by $\lambda_1=\lambda_0-\delta_1$.  The numerator is at
most
\[
 c_2^2C_\Phi(1+\|b\|)^{2q_2}I_D(\lambda_1)
 \exp\!\left[
 C_0+\frac{Q_0^2}{\lambda_1}
 +\frac34\lambda_1^{-1/3}\|b\|^{4/3}\right].
\]
Divide by \eqref{eq:V10-section-denominator} to obtain
\eqref{eq:V10-conditional-tail}.  Integrate the boundary-dependent right side
using \cref{lem:V10-boundary-integrability}.
\end{proof}

\begin{corollary}[Fixed-size Higgs--line blocks meet the V8 total-score hypotheses]
\label{cor:V10-fixed-block-total-score}
Consider a projective shell whose new Higgs block has a conditional potential
of the form \eqref{eq:V10-conditional-law}, with compact gauge/geometric
parameter $\xi$ and boundary field $b$ satisfying
\eqref{eq:V10-boundary-incidence}.  Suppose:
\begin{enumerate}[label=\textup{(\roman*)},leftmargin=2.8em]
\item the potential and polynomial feature bank obey the uniform data
      \eqref{eq:V10-potential-lower}--\eqref{eq:V10-ball-Gram};
\item the positive detail section and its command derivative obey
      \eqref{eq:V10-coefficient-window,eq:V10-coefficient-upper};
\item the neighboring old Higgs fields have the V8 local
      quartic-exponential bound;
\item the bosonic raw command derivative is a fixed local polynomial of
      degree at most four with bounded coefficient norm.
\end{enumerate}
Then the physical squared-section detail law has a uniform local
quartic-exponential moment, its line derivative has a uniform integrated
second moment, and its complete local total-cylinder score has a
cutoff-independent Fisher bound.
\end{corollary}

\begin{proof}
Item \textup{(iii)} and bounded incidence give
\eqref{eq:V10-neighbour-QET}.  Apply
\cref{thm:V10-integrated-line,thm:V10-section-tail}.  Item \textup{(iv)} and
\cref{lem:V8-polynomial-envelope,thm:V8-tail-to-total-Fisher} then give the
total-score Fisher bound.
\end{proof}
% -----------------------------------------------------------------------------
\subsection{Projective and Green consequences}
\label{subsec:V10-projective-consequence}
% -----------------------------------------------------------------------------

\begin{theorem}[Integrated fixed-block shell criterion]
\label{thm:V10-fixed-block-shell}
Assume the projective local line factorization of
\cref{thm:V9-projective-line-localization}.  At shell $n$, let the bosonic
command envelope have coefficient $C_{P,n}$ and let the section-derivative
coefficient in \eqref{eq:V10-coefficient-window} be $c_{1,n}$, with all other
structural constants fixed.  Suppose the neighboring physical Higgs fields
satisfy one common estimate \eqref{eq:V10-neighbour-QET}.  Then there are
finite constants $A_{\rm B},A_{\rm L}$, independent of $n$, such that
\begin{equation}
 \boxed{
 \Delta M_n
 \preceq
 (A_{\rm B}C_{P,n}^2+A_{\rm L}c_{1,n}^2)I_E.}
 \label{eq:V10-shell-Fisher}
\end{equation}
If
\begin{equation}
 \sum_nC_{P,n}^2+\sum_nc_{1,n}^2<\infty
 \label{eq:V10-shell-sum}
\end{equation}
and $C_n^{\rm act}\le C_*$, then the total Fisher and refresh-free Green
towers converge and
\begin{align}
 0\preceq M_\infty-M_N
 &\preceq
 \left(A_{\rm B}\sum_{n\ge N}C_{P,n}^2
      +A_{\rm L}\sum_{n\ge N}c_{1,n}^2\right)I_E,
 \label{eq:V10-total-Fisher-tail}\\
 0\preceq G_\infty-G_N
 &\preceq\frac{C_*}{4}
 \left(A_{\rm B}\sum_{n\ge N}C_{P,n}^2
      +A_{\rm L}\sum_{n\ge N}c_{1,n}^2\right)I_E.
 \label{eq:V10-total-Green-tail}
\end{align}
\end{theorem}

\begin{proof}
Theorems
\cref{thm:V10-integrated-line,thm:V10-section-tail} give, respectively, a
line second-moment constant proportional to $c_{1,n}^2$ and a common physical
quartic-exponential constant.  The bosonic polynomial envelope and the proof
of \cref{thm:V8-tail-to-total-Fisher} give the term proportional to
$C_{P,n}^2$.  This proves \eqref{eq:V10-shell-Fisher}.  Conditional
orthogonality of distinct factor scores makes the total Fisher tail their
positive sum.  Sum \eqref{eq:V10-shell-Fisher}; then use the V7 inequality
$\Delta G_n\preceq C_n^{\rm act}\Delta M_n/4$.
\end{proof}

\begin{assessment}[One-step closure and cofinal propagation are distinct]
\label{ass:V10-propagation-scope}
The hypotheses of \cref{thm:V10-fixed-block-shell} include a common quartic
tail for the neighboring \emph{physical} fields.  V10 proves that one further
squared-section detail block inherits such a tail, but the displayed constant
need not be a contraction under repeated refinement.  A uniform cofinal
physical-tail theorem therefore requires either a stable Lyapunov recursion,
a direct weighted estimate for the whole line-weighted specification, or
shell coefficients whose decay absorbs the propagated constants.  Finite
iteration alone is not used as a uniform conclusion.
\end{assessment}

% -----------------------------------------------------------------------------
\subsection{A populated interacting one-shell control}
\label{subsec:V10-interacting-control}
% -----------------------------------------------------------------------------

Let
\begin{equation}
 q_4(\dd x)=\frac{2}{\Gamma(1/4)}e^{-x^4}\dd x
 \qquad(x\in\mathbb R).
 \label{eq:V10-q4}
\end{equation}
For $J\in\mathbb R$ define
\begin{equation}
 \nu_x(\dd y)
 =Z_x^{-1}e^{-y^4+Jxy}\dd y.
 \label{eq:V10-control-reference}
\end{equation}
Fix $(m,a)\ne(0,0)$ and set
\begin{equation}
 K_x(\dd y)
 =\frac{|m+ay|^2\nu_x(\dd y)}
        {\mathbb E_{\nu_x}|m+aY|^2},
 \qquad
 \mu(\dd x,\dd y)=q_4(\dd x)K_x(\dd y).
 \label{eq:V10-control-joint}
\end{equation}

\begin{proposition}[Interacting quartic detail with a zero-safe line]
\label{prop:V10-interacting-control}
The law \eqref{eq:V10-control-joint} is exactly projective onto $q_4$ and is
nonproduct when $J\ne0$.  For every $0<\delta<1$,
\begin{equation}
 \mathbb E_{q_4}e^{\delta X^4}=(1-\delta)^{-1/4}.
 \label{eq:V10-q4-mgf}
\end{equation}
With $\Phi(y)=(1,y)^{\mathsf T}$, its radius-$R$ ball Gram is
\begin{equation}
 \int_{-R}^{R}\Phi(y)\Phi(y)^{\mathsf T}\dd y
 =\begin{pmatrix}2R&0\\0&2R^3/3\end{pmatrix}.
 \label{eq:V10-control-ball-Gram}
\end{equation}
Consequently:
\begin{enumerate}[label=\textup{(\roman*)},leftmargin=2.8em]
\item $\mathbb E_\mu e^{\delta_1Y^4}<\infty$ for every
      $0<\delta_1<1$;
\item every coefficient command
      $(m,a)\mapsto(m,a)+\theta(\dot m,\dot a)$ has finite squared-modulus
      line Fisher information;
\item if $a\ne0$, the section vanishes at $y=-m/a$ when this value is real,
      yet the preceding Fisher conclusion still holds.
\end{enumerate}
Thus the V10 integrated mechanism is populated on an interacting unbounded
old--detail pair without a Dobrushin or inverse-floor hypothesis.
\end{proposition}

\begin{proof}
Normalization of $K_x$ gives exact projectivity.  When $J\ne0$, its reference
conditional density depends nontrivially on $x$, and multiplication by the
fixed nonzero polynomial cannot remove that dependence, so the joint law is
not a product.  Scaling the integral of $e^{-(1-\delta)x^4}$ gives
\eqref{eq:V10-q4-mgf}; direct integration gives
\eqref{eq:V10-control-ball-Gram}.

Apply \cref{thm:V10-section-tail,thm:V10-integrated-line} with
$d=1$, $V(y)=y^4$, $b=Jx$, and the constant feature coefficient vectors
$c=(\overline m,\overline a)^{\mathsf T}$ and
$d=(\overline{\dot m},\overline{\dot a})^{\mathsf T}$.  Equation
\eqref{eq:V10-q4-mgf} supplies the boundary quartic tail.  The zero statement
is elementary and has $\nu_x$-measure zero.
\end{proof}

\begin{assessment}[Control scope]
\label{ass:V10-control-scope}
The control is genuinely interacting across its old--detail edge and its
positive line factor may vanish.  It verifies the new integrated estimates
on an exact unbounded model.  It has only one real scalar field and one shell;
it does not supply a chiral Standard-Model coefficient, Pfaffian orientation,
cofinal interacting source, represented stress, or Einstein dynamics.
\end{assessment}
% -----------------------------------------------------------------------------
\subsection{A nonlinear Lyapunov recursion gives cofinal tail stability}
\label{subsec:V10-Lyapunov}
% -----------------------------------------------------------------------------

\begin{lemma}[Sublinear exponential recursion]
\label{lem:V10-sublinear-recursion}
Let $R_n\ge0$ be adapted to a projective chain and suppose, for fixed
$\delta>0$, $A\ge1$, and $0\le\varepsilon<1$,
\begin{equation}
 \mathbb E[e^{\delta R_{n+1}}\mid R_0,\ldots,R_n]
 \le A e^{\varepsilon\delta R_n}.
 \label{eq:V10-Lyapunov-drift}
\end{equation}
Then
\begin{equation}
 \boxed{
 \sup_n\mathbb E e^{\delta R_n}
 \le
 \max\left\{\mathbb E e^{\delta R_0},
             A^{1/(1-\varepsilon)}\right\}.}
 \label{eq:V10-Lyapunov-fixed-point}
\end{equation}
\end{lemma}

\begin{proof}
Set $K_n=\mathbb E e^{\delta R_n}\ge1$.  Conditional expectation and
concavity of $t\mapsto t^\varepsilon$ give
\[
 K_{n+1}\le A\mathbb E(e^{\delta R_n})^\varepsilon
 \le AK_n^\varepsilon.
\]
If
$K_*=\max\{K_0,A^{1/(1-\varepsilon)}\}$ and $K_n\le K_*$, then
$K_{n+1}\le AK_*^\varepsilon\le K_*$.  Induction proves the claim.
\end{proof}

\begin{theorem}[Cofinal physical quartic tail on a bounded dependency cone]
\label{thm:V10-cofinal-tail}
Consider a projective sequence with a selected ancestry chain of
fixed-dimensional local blocks.  Let $X_n$ be the bounded old blanket used to
generate a new detail block $Y_{n+1}$, and let $R_n$ be the sum of fourth
powers on $X_n$.  Assume the blanket passed to the next step is a function of
the new block and satisfies
\[
 R_{n+1}\le a_0+a_1\sum_j\|Y_{n+1,j}\|^4
\]
for fixed $a_0,a_1$.  Assume uniformly:
\begin{enumerate}[label=\textup{(T\arabic*)},leftmargin=3.2em]
\item the conditional squared-section kernels satisfy the structural
      hypotheses of \cref{thm:V10-section-tail};
\item every boundary field has the bounded-incidence form
      \[
       \|b_n\|\le J\sum_{j=1}^{m}\|X_{n,j}\|;
      \]
\item the coefficient lower and upper bounds have the fixed polynomial
      exponents in
      \eqref{eq:V10-coefficient-window,eq:V10-coefficient-upper};
\item the initial blanket has a quartic-exponential moment.
\end{enumerate}
Then there is a $\delta>0$ such that
\begin{equation}
 \boxed{\sup_n\mathbb E e^{\delta R_n}<\infty.}
 \label{eq:V10-cofinal-physical-QET}
\end{equation}
Consequently, if the shell command coefficients obey
\eqref{eq:V10-shell-sum} and the source-active rates are uniformly bounded,
the total local Fisher and refresh-free Green towers converge with the tails
\eqref{eq:V10-total-Fisher-tail,eq:V10-total-Green-tail}.
\end{theorem}

\begin{proof}
Choose $\delta>0$ smaller than the uniform quartic coefficient in
\cref{thm:V10-section-tail}.  After decreasing $\delta$ to account for $a_1$ and absorbing $a_0$, its
conditional estimate has the form
\[
 \mathbb E[e^{\delta R_{n+1}}\mid\mathcal F_n]
 \le C(1+\|b_n\|)^p
       e^{C(1+\|b_n\|^{4/3})}.
\]  Fix any $0<\varepsilon<1$.  The incidence
bound, the polynomial absorption used in
\cref{lem:V10-boundary-integrability}, and
\eqref{eq:V10-dual-Young} give a finite $A_\varepsilon$, independent of $n$,
such that
\[
 C(1+\|b_n\|)^p e^{C(1+\|b_n\|^{4/3})}
 \le A_\varepsilon e^{\varepsilon\delta R_n}.
\]
Apply \cref{lem:V10-sublinear-recursion} to prove
\eqref{eq:V10-cofinal-physical-QET}.  This supplies the common neighboring
physical tail required by \cref{thm:V10-fixed-block-shell}; its Fisher and
Green conclusions then follow under the two stated summability hypotheses.
\end{proof}

\begin{assessment}[Dependency-cone hypothesis]
\label{ass:V10-dependency-cone}
The theorem is cofinal, but it is local.  It applies when the Markov blanket
needed at the next scale can be carried by a uniformly bounded local state or
by a fixed finite coloring of such states.  A refinement whose dependency
cone accumulates unboundedly many old coordinates must first be compressed,
or its growing blanket must enter the incidence and Lyapunov estimates
explicitly.
\end{assessment}
% -----------------------------------------------------------------------------
\subsection{Revised continuum acquisition}
\label{subsec:V10-acquisition}
% -----------------------------------------------------------------------------

\begin{openproblem}[Executable V11 acquisition: the actual RPESM dependency cone]
\label{op:V10-next-acquisition}
On one selected RPESM projective sequence:
\begin{enumerate}[label=\textup{(A\arabic*)},leftmargin=3.2em]
\item Print the fixed-size detail variables, their compact gauge/geometric
      parameter $\xi_n$, the unbounded linear Higgs boundary $b_n$, and the
      complete potential $V_{\xi_n}$.
\item Verify the uniform quartic lower envelope
      \eqref{eq:V10-potential-lower}, one ball upper bound, and the
      bounded-incidence representation \eqref{eq:V10-boundary-incidence}.
\item Express
      $\det E_n\,\operatorname{Pf}B_n$ and every command derivative in one
      polynomial feature bank.  Test polynomial independence on a fixed ball
      and the coefficient window
      \eqref{eq:V10-coefficient-window,eq:V10-coefficient-upper}.
\item Determine whether the next-scale Markov blanket is replaced by the new
      detail block, as required by \cref{thm:V10-cofinal-tail}, or retains old
      coordinates.  In the latter case, construct a decaying weighted
      dependency-cone radius and prove its sublinear drift instead of invoking
      the replacement theorem.
\item Compute $C_{P,n}$ and $c_{1,n}$ in physical normalization and verify
      \eqref{eq:V10-shell-sum}; otherwise retain their divergent weighted sum
      as the explicit Fisher obstruction.
\item Apply \eqref{eq:V10-total-Green-tail} on the same projective line
      trivialization.
\item Separately transport the chiral phase connection, divisor, real
      Pfaffian orientation, and crossing sign.  The squared-modulus estimates
      do not determine them.
\item Test one predeclared non-Gaussian infrared source and the actual
      inverse-coframe/represented-stress writer on the same cofinal sequence.
\end{enumerate}
A pass closes the full positive-modulus local Higgs--line Fisher and Green
rows on the selected regulator.  Failure returns a named potential, feature,
coefficient, dependency-cone, shell-summability, phase/orientation, infrared,
or stress obstruction.
\end{openproblem}

\begin{theorem}[Version V10 terminal theorem]
\label{thm:V10-terminal}
Preserving V1--V9 verbatim before this appended block, V10 proves:
\begin{enumerate}[label=\textup{(V10.\arabic*)},leftmargin=5.6em]
\item a quartic conditional law with an arbitrary linear boundary has
      explicit upper and lower polynomial feature-Gram bounds at the dual
      scale $\exp(C\|b\|^{4/3})$;
\item compact polynomial independence reduces the required local Gram datum
      to a finite algebraic test on one ball;
\item a local quartic-exponential moment integrates every polynomial times
      $\exp(C\|b\|^{4/3})$ for a bounded-incidence boundary field;
\item the squared-section line score has a uniform integrated second moment
      under polynomial coefficient windows, even when its pointwise
      conditional normalizer tends to zero;
\item the squared polynomial section preserves a detail-field
      quartic-exponential moment without a pointwise section-mass floor;
\item a fixed-size Higgs--line detail block satisfying the printed
      hypotheses has a cutoff-independent total-score Fisher bound;
\item square-summable local action and section derivatives yield explicit
      total-Fisher and refresh-free Green tails;
\item a sublinear exponential Lyapunov recursion propagates the physical
      quartic tail cofinally along a bounded replacement dependency cone;
\item an exact interacting quartic old--detail law with a possibly vanishing
      line section populates the integrated mechanism.
\end{enumerate}
\end{theorem}

\begin{proof}
Items \textup{(V10.1)}--\textup{(V10.2)} are
\cref{thm:V10-boundary-Gram,lem:V10-compact-ball-Gram}.
Item \textup{(V10.3)} is \cref{lem:V10-boundary-integrability}.
Items \textup{(V10.4)}--\textup{(V10.5)} are
\cref{thm:V10-integrated-line,thm:V10-section-tail}.
Items \textup{(V10.6)}--\textup{(V10.7)} are
\cref{cor:V10-fixed-block-total-score,thm:V10-fixed-block-shell}.
Item \textup{(V10.8)} is
\cref{lem:V10-sublinear-recursion,thm:V10-cofinal-tail}.
The final item is \cref{prop:V10-interacting-control}.
\end{proof}

\begin{assessment}[Programme boundary after V10]
\label{ass:V10-boundary}
V10 removes the pointwise conditional section-mass floor from the analytic
requirements.  Large unbounded Higgs boundary fields may make a detail
feature Gram ill conditioned, but the loss occurs at the quartic-dual scale
and is integrable under the V8 tail estimate.  V10 also gives a genuine
cofinal tail theorem when the local scale ancestry has a bounded replacement
state.

The full Standard--Model--Einstein continuum remains unproved.  The abstract
RPESM construction does not print the detail potentials, polynomial control
sections, uniform coefficient windows, physical shell normalization, or
dependency-cone geometry needed to apply these results.  A cumulative
unbounded dependency cone requires a new weighted drift estimate.  The
chiral phase, Pfaffian orientation, interacting infrared source,
inverse-coframe and represented-stress bounds, OS continuum, renormalized
products, Lorentzian reconstruction, and Palatini--Einstein limit remain
open.
\end{assessment}

\section*{V10 exact checks and append-only record}
\addcontentsline{toc}{section}{V10 exact checks and append-only record}
The constants in \eqref{eq:V10-Young-quadratic} and
\eqref{eq:V10-Young-linear} are the exact maxima
\[
 \sup_{r\ge0}\left(Qr^2-\frac{\lambda}{4}r^4\right)
 =\frac{Q^2}{\lambda},
 \qquad
 \sup_{r\ge0}\left(Br-\frac{\lambda}{4}r^4\right)
 =\frac34\lambda^{-1/3}B^{4/3}.
\]
Likewise,
\[
 \sup_{r\ge0}(cr^{4/3}-\epsilon r^4)
 =\frac{2c^{3/2}}{3\sqrt{3\epsilon}},
\]
which checks \eqref{eq:V10-dual-Young}.  The radial quartic integrals
$I_0$, $I_D$, and $I_D(\lambda_1)$ are finite by direct polar comparison.
The recursion fixed point was checked algebraically from
$AK_*^\varepsilon\le K_*$.

For the interacting control, direct quadrature verifies the normalizer and
feature-Gram formulas over a range of positive and negative boundary fields,
and sampling verifies the predicted finite quartic and zero-safe line
moments.  These computations are consistency checks only; the stated results
follow from the displayed inequalities.

The complete V9 source before removal of its final active document-closing
command is retained verbatim.  Its complete-file SHA--256 digest is
\begin{center}\small\ttfamily
a9f61c95047f26a0ffe2bfbecb7f18776929c11635675925760350cf911d8fac.
\end{center}
A Version V11 must retain the complete V10 source, remove only the final
active document-closing command, append new material, restore one final
command, and use a filename ending in \texttt{\_V11.tex}.

% =============================================================================
% END APPENDED VERSION V10 MATERIAL
% =============================================================================


% =============================================================================
% BEGIN APPENDED VERSION V11 MATERIAL
% =============================================================================

\clearpage
\section{Version V11: decaying ancestry cones and cofinal line-weighted tails}
\label{sec:V11-ancestry-tail}

Version V10 proved cofinal quartic-tail stability when the local Markov
blanket passed to the next scale is controlled by the newly sampled detail
block.  A projective refinement may instead retain older coordinates.  Then
the unweighted sum of all ancestral fourth powers grows with the cutoff and
the replacement recursion does not apply.

This version introduces a geometrically weighted ancestry radius.  Old
coordinates remain in the record, but their Lyapunov weight decays with
scale age.  A weighted H\"older inequality identifies the exact incidence
quantity
\[
 \sum_j\|A_{n,j}\|^{4/3}\beta^{-(n-j)/3}.
\]
If this quantity is uniformly bounded, the quartic-dual boundary loss of V10
is sublinear in the weighted radius.  The resulting nonlinear recurrence has
a cutoff-independent fixed point.  Geometrically decaying ancestry
coefficients pass precisely when their fourth-power decay dominates the
chosen memory weight.

The final part populates the result on an infinite, interacting, noncompact,
projective squared-line tower with dependence on every preceding scale.  Its
factor scores have a finite cofinal Fisher limit without a pointwise inverse
floor.  The model is an analytic control and is not the RPESM continuum.

% -----------------------------------------------------------------------------
\subsection{The weighted incidence norm}
\label{subsec:V11-weighted-incidence}
% -----------------------------------------------------------------------------

Let $X_0,\ldots,X_n$ be finite-dimensional real vectors, let
$0<\beta<1$, and define
\begin{equation}
 \mathcal R_n
 =\sum_{j=0}^n\beta^{\,n-j}\|X_j\|^4.
 \label{eq:V11-weighted-radius}
\end{equation}
Let the boundary field driving the next conditional block be
\begin{equation}
 b_n=\sum_{j=0}^nA_{n,j}X_j
 \label{eq:V11-boundary-field}
\end{equation}
and put
\begin{equation}
 \mathcal I_{\beta,n}
 =\sum_{j=0}^n
   \|A_{n,j}\|^{4/3}\beta^{-(n-j)/3}.
 \label{eq:V11-weighted-incidence}
\end{equation}

\begin{lemma}[Weighted ancestry H\"older bound]
\label{lem:V11-weighted-Holder}
For every $n$,
\begin{equation}
 \boxed{
 \|b_n\|^{4/3}
 \le\mathcal I_{\beta,n}\mathcal R_n^{1/3}.}
 \label{eq:V11-boundary-radius}
\end{equation}
Thus a uniform bound on $\mathcal I_{\beta,n}$ converts the quartic-dual
boundary scale into a sublinear function of the retained ancestry radius.
\end{lemma}

\begin{proof}
Apply H\"older with exponents $4$ and $4/3$:
\begin{align*}
 \|b_n\|
 &\le\sum_{j=0}^n
   \bigl(\beta^{(n-j)/4}\|X_j\|\bigr)
   \bigl(\|A_{n,j}\|\beta^{-(n-j)/4}\bigr)\\
 &\le\mathcal R_n^{1/4}
 \left(\sum_{j=0}^n
 \|A_{n,j}\|^{4/3}\beta^{-(n-j)/3}\right)^{3/4}.
\end{align*}
Raise both sides to the power $4/3$.
\end{proof}

\begin{lemma}[Sublinear absorption at the weighted radius]
\label{lem:V11-sublinear-absorption}
For every $a,\eta>0$ and $r\ge0$,
\begin{equation}
 ar^{1/3}
 \le\eta r+\frac{2a^{3/2}}{3\sqrt{3\eta}}.
 \label{eq:V11-radius-Young}
\end{equation}
Consequently, if
$\sup_n\mathcal I_{\beta,n}\le I_*$, then for every
$C,p,\eta>0$ there is $K_{C,p,\eta}<\infty$ such that
\begin{equation}
 (1+\|b_n\|)^p e^{C\|b_n\|^{4/3}}
 \le K_{C,p,\eta}e^{\eta\mathcal R_n}
 \quad\text{for all }n.
 \label{eq:V11-boundary-absorption}
\end{equation}
\end{lemma}

\begin{proof}
The first inequality follows by maximizing
$ar^{1/3}-\eta r$.  A polynomial is bounded by an arbitrarily small
exponential in $\|b_n\|^{4/3}$.  Combine that fact with
\eqref{eq:V11-boundary-radius}, replace $a$ by the resulting constant times
$I_*$, and apply \eqref{eq:V11-radius-Young}.
\end{proof}

% -----------------------------------------------------------------------------
\subsection{Cofinal stability with retained ancestry}
\label{subsec:V11-cofinal-stability}
% -----------------------------------------------------------------------------

\begin{theorem}[Weighted-ancestry quartic-tail theorem]
\label{thm:V11-weighted-tail}
Let $(X_n)_{n\ge0}$ be generated by a projective sequence of conditional
kernels.  Suppose there are fixed $\delta_0,C_0,C_1>0$ and $p\ge0$ such that
\begin{equation}
 \mathbb E\!\left[
 e^{\delta_0\|X_{n+1}\|^4}\mid\mathcal F_n\right]
 \le C_0(1+\|b_n\|)^p
       e^{C_1\|b_n\|^{4/3}},
 \label{eq:V11-conditional-tail}
\end{equation}
where $b_n$ is \eqref{eq:V11-boundary-field}.  Assume
\begin{equation}
 \sup_n\mathcal I_{\beta,n}\le I_*<\infty
 \label{eq:V11-incidence-window}
\end{equation}
for some $0<\beta<1$, and
$\mathbb E e^{\delta\|X_0\|^4}<\infty$ for some
$0<\delta\le\delta_0$.

Then, after decreasing $\delta$ if necessary,
\begin{equation}
 \boxed{
 \sup_{n\ge0}\mathbb E e^{\delta\mathcal R_n}<\infty.}
 \label{eq:V11-uniform-weighted-tail}
\end{equation}
In particular,
\begin{equation}
 \boxed{
 \sup_{n\ge0}\mathbb E e^{\delta\|X_n\|^4}<\infty.}
 \label{eq:V11-uniform-current-tail}
\end{equation}
\end{theorem}

\begin{proof}
The identity
\begin{equation}
 \mathcal R_{n+1}
 =\beta\mathcal R_n+\|X_{n+1}\|^4
 \label{eq:V11-radius-recursion}
\end{equation}
is exact.  Choose $0<\varepsilon<1-\beta$.  By monotonicity, the left side
of \eqref{eq:V11-conditional-tail} remains bounded by the same right side
when $\delta_0$ is replaced by $\delta\le\delta_0$.  Apply
\eqref{eq:V11-boundary-absorption} with
$\eta=\varepsilon\delta$ to obtain a constant $A_\varepsilon$ such that
\[
 \mathbb E[e^{\delta\|X_{n+1}\|^4}\mid\mathcal F_n]
 \le A_\varepsilon e^{\varepsilon\delta\mathcal R_n}.
\]
Multiplication by $e^{\beta\delta\mathcal R_n}$ and
\eqref{eq:V11-radius-recursion} give
\[
 \mathbb E[e^{\delta\mathcal R_{n+1}}\mid\mathcal F_n]
 \le A_\varepsilon
 e^{(\beta+\varepsilon)\delta\mathcal R_n}.
\]
Since $\beta+\varepsilon<1$,
\cref{lem:V10-sublinear-recursion} yields
\eqref{eq:V11-uniform-weighted-tail}.  The newest coordinate has weight one
in \eqref{eq:V11-weighted-radius}, which proves
\eqref{eq:V11-uniform-current-tail}.
\end{proof}

\begin{corollary}[Geometric ancestry threshold]
\label{cor:V11-geometric-threshold}
If, for fixed $J_0>0$ and $0\le\vartheta<1$,
\begin{equation}
 \|A_{n,j}\|\le J_0\vartheta^{\,n-j},
 \label{eq:V11-geometric-incidence}
\end{equation}
then \eqref{eq:V11-incidence-window} holds whenever
\begin{equation}
 \boxed{\vartheta^4<\beta<1,}
 \label{eq:V11-decay-threshold}
\end{equation}
and
\begin{equation}
 \mathcal I_{\beta,n}
 \le\frac{J_0^{4/3}}
 {1-\vartheta^{4/3}\beta^{-1/3}}.
 \label{eq:V11-incidence-sum}
\end{equation}
For a memory of at most $r$ preceding scales, the incidence window holds for
every $0<\beta<1$, with a constant depending on $r$.
\end{corollary}

\begin{proof}
Insert \eqref{eq:V11-geometric-incidence} in
\eqref{eq:V11-weighted-incidence} and sum the geometric series in
$k=n-j$.  Its ratio is
$\vartheta^{4/3}\beta^{-1/3}<1$ exactly when
$\vartheta^4<\beta$.  A finite-memory sum has only $r$ terms.
\end{proof}
\begin{corollary}[A finite persistent core may be conditioned rather than forgotten]
\label{cor:V11-persistent-core}
Let $C$ be a fixed finite-dimensional core variable and replace
\eqref{eq:V11-boundary-field} by
\begin{equation}
 b_n=B_nC+\sum_{j=0}^nA_{n,j}X_j,
 \qquad \sup_n\|B_n\|\le J_C.
 \label{eq:V11-core-boundary}
\end{equation}
Assume the tail incidence window \eqref{eq:V11-incidence-window}, and suppose
there are $\delta_C>0$ and finite constants such that
\begin{align}
 \mathbb E e^{\delta_C\|C\|^4}&<\infty,
 \label{eq:V11-core-tail}\\
 \mathbb E[e^{\delta\mathcal R_0}\mid C]
 &\le A_0(1+\|C\|)^{p_0}e^{a_0\|C\|^{4/3}}.
 \label{eq:V11-core-initial}
\end{align}
Then, after decreasing $\delta>0$,
\begin{equation}
 \boxed{\sup_n\mathbb E e^{\delta\mathcal R_n}<\infty.}
 \label{eq:V11-core-conclusion}
\end{equation}
Thus finitely many nondecaying ancestral coordinates need not be assigned
artificially decaying weights.
\end{corollary}

\begin{proof}
Condition on $C$.  The contribution of $B_nC$ changes the constant in
\eqref{eq:V11-boundary-absorption} to
\[
 A(C)\le A_1(1+\|C\|)^{p_1}e^{a_1\|C\|^{4/3}}.
\]
The proof of \cref{thm:V11-weighted-tail} gives, conditionally on $C$,
\[
 K_{n+1}(C)\le A(C)K_n(C)^\rho,
 \qquad \rho=\beta+\varepsilon<1.
\]
Together with \eqref{eq:V11-core-initial},
\[
 \sup_nK_n(C)
 \le\max\{K_0(C),A(C)^{1/(1-\rho)}\}.
\]
Every polynomial times $e^{a\|C\|^{4/3}}$ is integrable under
\eqref{eq:V11-core-tail}, by the scalar Young estimate
\eqref{eq:V10-dual-Young}.  Integrate the conditional bound.
\end{proof}

\begin{corollary}[V10 line-weighted blocks satisfy the ancestry hypothesis]
\label{cor:V11-V10-handoff}
A fixed-dimensional squared-section block satisfying the uniform hypotheses
of \cref{thm:V10-section-tail} obeys
\eqref{eq:V11-conditional-tail}.  Hence
\cref{thm:V11-weighted-tail} supplies a uniform physical local quartic tail
for any retained ancestry satisfying \eqref{eq:V11-incidence-window}, with
the persistent-core extension of \cref{cor:V11-persistent-core} when needed.

If, on the same projective sequence, the local bosonic and section-derivative
coefficients are square summable and the source-active rates are uniformly
bounded, the Fisher and Green conclusions of
\cref{thm:V10-fixed-block-shell} follow.
\end{corollary}

\begin{proof}
The right side of \eqref{eq:V10-conditional-tail} is a fixed polynomial in
$1+\|b_n\|$ times the exponential of a constant multiple of
$1+\|b_n\|^{4/3}$.  Absorb the constant into $C_0$ in
\eqref{eq:V11-conditional-tail}.  Apply
\cref{thm:V11-weighted-tail,cor:V11-persistent-core}, then
\cref{thm:V10-fixed-block-shell}.
\end{proof}

\begin{assessment}[What the ancestry estimate asks of the actual shell]
\label{ass:V11-RPESM-incidence}
The abstract RPESM kernel states only that the old--detail interaction is
finite range.  Spatial finite range does not by itself give the scale-age
bound \eqref{eq:V11-incidence-window}: a detail may couple to a fixed spatial
cell containing an increasing number of retained scales.  The required
audit is now numerical and local.  Its output is the weighted row sum
\eqref{eq:V11-weighted-incidence}, together with a finite list of persistent
core coordinates.  No global lattice cardinality enters.
\end{assessment}

% -----------------------------------------------------------------------------
\subsection{An infinite interacting all-ancestor line tower}
\label{subsec:V11-control-tower}
% -----------------------------------------------------------------------------

Let $q_4$ be \eqref{eq:V10-q4}.  Fix real $m\ne0$, $a\ne0$, $J\ne0$,
and $0<\vartheta<1$.  Given $X_0,\ldots,X_n$, put
\begin{equation}
 b_n=J\sum_{j=0}^n\vartheta^{\,n-j}X_j.
 \label{eq:V11-control-boundary}
\end{equation}
For a real command parameter $\tau$ and coefficients $\gamma_n$, define
\begin{equation}
 K_{n,\tau}(\dd y\mid X_0,\ldots,X_n)
 =
 \frac{|m+(a+\tau\gamma_n)y|^2e^{-y^4+b_ny}\dd y}
 {\displaystyle\int_{\mathbb R}
  |m+(a+\tau\gamma_n)z|^2e^{-z^4+b_nz}\dd z}.
 \label{eq:V11-control-kernel}
\end{equation}
Starting with $X_0\sim q_4$, let
\begin{equation}
 \mu_{N,\tau}(\dd x_0\cdots\dd x_N)
 =q_4(\dd x_0)
  \prod_{n=0}^{N-1}
  K_{n,\tau}(\dd x_{n+1}\mid x_0,\ldots,x_n).
 \label{eq:V11-control-law}
\end{equation}

\begin{theorem}[Cofinal interacting zero-safe line Fisher tower]
\label{thm:V11-control-tower}
The laws \eqref{eq:V11-control-law} are exactly projective and nonproduct.
For every $\beta$ with $\vartheta^4<\beta<1$, there is $\delta>0$ such that
\begin{equation}
 \sup_N\mathbb E_{\mu_{N,0}}
 \exp\!\left(
 \delta\sum_{j=0}^N\beta^{N-j}|X_j|^4\right)<\infty,
 \qquad
 \sup_N\mathbb E_{\mu_{N,0}}e^{\delta|X_N|^4}<\infty.
 \label{eq:V11-control-tail}
\end{equation}
Let
\begin{equation}
 \eta_n
 =\left.\partial_\tau
 \log K_{n,\tau}(X_{n+1}\mid X_0,\ldots,X_n)
 \right|_{\tau=0}.
 \label{eq:V11-control-score}
\end{equation}
Then $(\eta_n)$ is a martingale-difference family and
\begin{equation}
 \mathbb E_{\mu_{N,0}}\eta_j\eta_k=0
 \quad(j\ne k).
 \label{eq:V11-control-orthogonality}
\end{equation}
There is a constant $C_{\rm ctl}<\infty$, independent of $n$, such that
\begin{equation}
 \mathbb E_{\mu_{N,0}}\eta_n^2
 \le C_{\rm ctl}\gamma_n^2.
 \label{eq:V11-control-increment}
\end{equation}
Consequently, if $(\gamma_n)\in\ell^2$, the total score
$\sum_{n<N}\eta_n$ has the finite cofinal Fisher limit
\begin{equation}
 \boxed{
 M_\infty
 =\sum_{n\ge0}\mathbb E\eta_n^2
 \le C_{\rm ctl}\sum_{n\ge0}\gamma_n^2<\infty.}
 \label{eq:V11-control-Fisher}
\end{equation}
Every conditional section vanishes at $y=-m/a$, so the construction has no
pointwise line inverse floor.
\end{theorem}

\begin{proof}
Each factor in \eqref{eq:V11-control-kernel} is a normalized probability
kernel, which proves projectivity.  Since $J\ne0$, the conditional law depends
on the retained ancestry through $b_n$ and is not a product.

The potential is $V(y)=y^4$, the feature bank is $(1,y)$, and the coefficient
vector $(m,a)$ is fixed and nonzero.  Thus
\cref{thm:V10-section-tail} gives \eqref{eq:V11-conditional-tail}.  Here
$A_{n,j}=J\vartheta^{n-j}$, so
\cref{cor:V11-geometric-threshold,thm:V11-weighted-tail} prove
\eqref{eq:V11-control-tail}.

Normalized-kernel differentiation gives
$\mathbb E[\eta_n\mid\mathcal F_n]=0$.  If $j<k$, then $\eta_j$ is
$\mathcal F_k$-measurable, and conditional expectation proves
\eqref{eq:V11-control-orthogonality}.  The derivative of the local section
has feature coefficient proportional to $\gamma_n$.  The integrated
zero-safe estimate \cref{thm:V10-integrated-line}, together with
\eqref{eq:V11-control-tail}, gives
\eqref{eq:V11-control-increment}.  Orthogonal Pythagoras and monotone
convergence give \eqref{eq:V11-control-Fisher}.  The zero of the real affine
section is immediate and has conditional Lebesgue measure zero.
\end{proof}

\begin{assessment}[Control-tower scope]
\label{ass:V11-control-scope}
The tower is infinite, noncompact, interacting with every preceding scale,
projective, zero safe, and has a finite total line Fisher response.  It does
not carry a reversible local Green generator, the Standard-Model gauge and
chiral line, Pfaffian orientation, represented stress, or Einstein dynamics.
It proves consistency of the retained-ancestry analytic mechanism rather than
the physical continuum.
\end{assessment}
% -----------------------------------------------------------------------------
\subsection{Revised continuum acquisition}
\label{subsec:V11-acquisition}
% -----------------------------------------------------------------------------

\begin{openproblem}[Executable V12 acquisition: physical ancestry and coframe stress]
\label{op:V11-next-acquisition}
For one selected RPESM shell family:
\begin{enumerate}[label=\textup{(A\arabic*)},leftmargin=3.2em]
\item Expand the old-variable dependence of each local detail potential and
      record the matrices $A_{n,j}$ in
      \eqref{eq:V11-boundary-field}.
\item Choose a memory weight $\beta$ and compute the exact row sum
      $\mathcal I_{\beta,n}$.  Either prove
      \eqref{eq:V11-incidence-window}, isolate finitely many nondecaying
      coordinates for \cref{cor:V11-persistent-core}, or return the divergent
      ancestry row as the obstruction.
\item On the same factorization, verify the V10 polynomial feature and
      coefficient hypotheses for
      $\det E_n\,\operatorname{Pf}B_n$.
\item Compute the physically normalized action and section derivatives and
      test their square-summability.  Insert a passing result in
      \eqref{eq:V10-total-Fisher-tail,eq:V10-total-Green-tail}.
\item With the positive-modulus Higgs--line row fixed, write the complete
      coframe-dependent total score.  Separate polynomial volume/coframe
      terms from every inverse coframe or inverse metric factor.
\item Prove a uniform nondegeneracy or integrable inverse-determinant estimate
      for those geometric factors, or return a named coframe-collapse
      concentration witness.
\item Retain the chiral phase connection, Pfaffian orientation, crossing sign,
      non-Gaussian infrared source, and Palatini residual as independent
      rows.
\end{enumerate}
A pass through items \textup{(A1)}--\textup{(A4)} closes the cumulative
positive-modulus local Fisher and Green problem.  Items
\textup{(A5)}--\textup{(A6)} then open the actual represented stress and
Einstein source row.
\end{openproblem}

\begin{theorem}[Version V11 terminal theorem]
\label{thm:V11-terminal}
Preserving V1--V10 verbatim before this appended block, V11 proves:
\begin{enumerate}[label=\textup{(V11.\arabic*)},leftmargin=5.8em]
\item the quartic-dual boundary loss on a retained ancestry is controlled by
      the weighted incidence norm
      \eqref{eq:V11-weighted-incidence};
\item a uniform weighted incidence bound converts every polynomial times
      $\exp(C\|b_n\|^{4/3})$ into an arbitrarily small exponential of the
      ancestry radius;
\item the decaying-radius identity
      $\mathcal R_{n+1}=\beta\mathcal R_n+\|X_{n+1}\|^4$ gives a sublinear
      Lyapunov recursion and a uniform cofinal physical quartic tail;
\item geometric scale-age couplings
      $\|A_{n,j}\|\le J_0\vartheta^{n-j}$ pass exactly when one can choose
      $\vartheta^4<\beta<1$;
\item finitely many nondecaying ancestral coordinates can be retained as a
      persistent core under a quartic-exponential core law;
\item fixed-size V10 squared-section blocks inherit the weighted-ancestry
      theorem and hence the V10 Fisher/Green conclusion under
      square-summable command derivatives;
\item the only additional shell datum is the local weighted row sum, rather
      than the global number of lattice or fermion coordinates;
\item an exact infinite all-ancestor quartic line tower is projective,
      interacting, zero safe, uniformly locally tight, and has a finite
      cofinal Fisher response for square-summable line commands.
\end{enumerate}
\end{theorem}

\begin{proof}
Items \textup{(V11.1)}--\textup{(V11.2)} are
\cref{lem:V11-weighted-Holder,lem:V11-sublinear-absorption}.
Items \textup{(V11.3)}--\textup{(V11.4)} are
\cref{thm:V11-weighted-tail,cor:V11-geometric-threshold}.
Item \textup{(V11.5)} is \cref{cor:V11-persistent-core}.
Items \textup{(V11.6)}--\textup{(V11.7)} are
\cref{cor:V11-V10-handoff,ass:V11-RPESM-incidence}.
The final item is \cref{thm:V11-control-tower}.
\end{proof}

\begin{assessment}[Programme boundary after V11]
\label{ass:V11-boundary}
V11 removes the bounded replacement-cone restriction from V10.  Retained
scale ancestry is compatible with a uniform physical quartic tail whenever
its weighted $4/3$ incidence row is bounded; geometric decay has the explicit
threshold \eqref{eq:V11-decay-threshold}.  The exact control tower shows that
cumulative interaction, projectivity, section zeros, local tightness, and
finite cofinal line Fisher information can coexist.

The full Standard--Model--Einstein continuum remains unproved.  The attached
RPESM source does not print the matrices $A_{n,j}$, their scale decay,
persistent core, physical command coefficients, or the local detail section
features needed to apply V11.  Uniform inverse-coframe and represented-stress
control, the chiral phase and Pfaffian orientation, interacting infrared
source survival, OS continuum, renormalized products, Lorentzian
reconstruction, and the Palatini--Einstein limit remain open.
\end{assessment}

\section*{V11 exact checks and append-only record}
\addcontentsline{toc}{section}{V11 exact checks and append-only record}
The weighted H\"older calculation was checked directly by raising
\[
 \sum_j
 (\beta^{(n-j)/4}\|X_j\|)
 (\|A_{n,j}\|\beta^{-(n-j)/4})
 \le\mathcal R_n^{1/4}\mathcal I_{\beta,n}^{3/4}
\]
to the power $4/3$.  The geometric series has ratio
$\vartheta^{4/3}\beta^{-1/3}$, so convergence is equivalent to
$\vartheta^4<\beta$.  Direct random positive-array tests verify
\eqref{eq:V11-boundary-radius}, the geometric sum
\eqref{eq:V11-incidence-sum}, and the sublinear fixed-point recursion.

For the control tower, numerical conditional quadrature and sequential
sampling verify projectivity, nontrivial dependence on the retained boundary,
the uniform recent-coordinate tail, martingale-difference centering, and
cross-shell score orthogonality on finite truncations.  The computations are
consistency checks; the theorem follows from the analytic estimates.

The complete V10 source before removal of its final active document-closing
command is retained verbatim.  Its complete-file SHA--256 digest is
\begin{center}\small\ttfamily
56901a0e4bda0577179eca4f9f7318b0bce53957d0e514c791b9d2178527836b.
\end{center}
A Version V12 must retain the complete V11 source, remove only the final
active document-closing command, append new material, restore one final
command, and use a filename ending in \texttt{\_V12.tex}.

% =============================================================================
% END APPENDED VERSION V11 MATERIAL
% =============================================================================


% =============================================================================
% BEGIN APPENDED VERSION V12 MATERIAL
% =============================================================================

\clearpage
\section{Version V12: coframe-collapse tails and represented stress Fisher control}
\label{sec:V12-coframe-stress}

Versions V8--V11 control the unbounded Higgs field and the positive
determinant/Pfaffian modulus on projective shells satisfying explicit local
incidence conditions.  The next total-score term is geometric.  The attached
monograph already proves two finite algebraic facts: for the spatial frame
source $E$, the Gram $q=E^*E$ has
$\det q=\|\bigwedge^3E\|^2$, and a uniform positive source-frame floor gives
a uniform determinant-frame floor.  It also states the normalized DeWitt
response in terms of $q^{-1}$.  Those results are imported here.

The unbounded projective RPESM law supplies neither a hard cofinal frame floor
nor a uniform negative determinant moment.  Finite-cutoff positivity is not
enough: normalized volume, inverse metric, unit-normal, and represented
stress writers can concentrate near frame collapse.  This version separates
the polynomial, densitized Palatini core from normalized inverse-metric
readouts and proves a probabilistic alternative to the hard frame floor.
A joint positive-field/collapse-barrier Orlicz estimate gives explicit
negative determinant moments, inverse-frame moments, and total-score Fisher
bounds.

No volume-collapse barrier is asserted to occur in the attached physical
cylinder.  The result identifies exactly what would have to be acquired, or
what concentration witness is returned if it fails.

% -----------------------------------------------------------------------------
\subsection{Finite positivity is not cofinal inverse control}
\label{subsec:V12-finite-positive}
% -----------------------------------------------------------------------------

\begin{proposition}[Pointwise frame positivity with divergent normalized score]
\label{prop:V12-positive-no-inverse}
Let $\mu_n$ give probability $1/2$ to each of
\[
 q_n^{(1)}=\operatorname{diag}(n^{-1},1,1),
 \qquad
 q_n^{(2)}=\operatorname{diag}(2n^{-1},1,1),
\]
and let $h=\operatorname{diag}(1,0,0)$.  Every $q$ is positive definite and
$\|q\|\le1$, but
\begin{align}
 \mathbb E_{\mu_n}\|q^{-1}\|^p&\longrightarrow\infty
 &&(p>0),
 \label{eq:V12-inverse-divergence}\\
 \operatorname{Var}_{\mu_n}\Tr(q^{-1}h)
 &=\frac{n^2}{16}\longrightarrow\infty.
 \label{eq:V12-logdet-Fisher-divergence}
\end{align}
The second quantity is the Fisher information at zero of the normalized
determinant command
\begin{equation}
 \mu_{n,\theta}(\dd q)
 =
 \frac{\det(q+\theta h)/\det q}
      {\mathbb E_{\mu_n}[\det(q+\theta h)/\det q]}
 \,\mu_n(\dd q)
 \label{eq:V12-det-command}
\end{equation}
for sufficiently small $\theta\ge0$.
\end{proposition}

\begin{proof}
The inverse norms are $n$ and $n/2$.  Jacobi's formula gives
\[
 \left.\partial_\theta
 \log\frac{\det(q+\theta h)}{\det q}\right|_{\theta=0}
 =\Tr(q^{-1}h),
\]
whose two values are $n$ and $n/2$.  Centering gives deviations
$\pm n/4$, proving \eqref{eq:V12-logdet-Fisher-divergence}.  The command
\eqref{eq:V12-det-command} is a positive probability law for
$\theta\ge0$ and its score is exactly that centered derivative.
\end{proof}

\begin{assessment}[Meaning for the RPESM chart]
\label{ass:V12-finite-positive-scope}
A positive rate/speed chart and a faithful frame at every finite cutoff do not
supply a cofinal inverse estimate.  One needs the deterministic frame floor
already isolated in the monograph, a negative-volume moment estimate, or an
equivalent probabilistic exclusion of the collapsing stratum.
\end{assessment}

% -----------------------------------------------------------------------------
\subsection{Polynomial Palatini data versus normalized inverse data}
\label{subsec:V12-polynomial-normalized}
% -----------------------------------------------------------------------------

\begin{proposition}[The inverse enters only after normalization or metric readout]
\label{prop:V12-polynomial-split}
For a $3\times3$ symmetric matrix $q$, $\det q$ is cubic and
$\operatorname{adj}q$ is quadratic.  On $q\succ0$,
\begin{align}
 D(\det)_q[h]
 &=\Tr((\operatorname{adj}q)h)
  =(\det q)\Tr(q^{-1}h),
 \label{eq:V12-det-first}\\
 -\frac{D^2(\det)_q[h,k]}{\det q}
 &=\Tr(q^{-1}h\,q^{-1}k)
   -\Tr(q^{-1}h)\Tr(q^{-1}k),
 \label{eq:V12-det-normalized-second}\\
 -D^2(\log\det)_q[h,k]
 &=\Tr(q^{-1}h\,q^{-1}k).
 \label{eq:V12-logdet-second}
\end{align}
The unnormalized first and second derivatives of $\det q$ extend
polynomially across $\det q=0$.  The normalized DeWitt and logarithmic
responses do not.

Likewise, the Palatini--Holst bulk density $e\wedge e\wedge F$ and the
cosmological density $e^{\wedge4}$ are polynomial in the coframe and
curvature before conversion to inverse-metric, unit-normal, or normalized
volume variables.
\end{proposition}

\begin{proof}
The adjugate identity gives \eqref{eq:V12-det-first}.  Differentiate Jacobi's
formula and use
$D(q^{-1})[k]=-q^{-1}kq^{-1}$ to obtain
\eqref{eq:V12-det-normalized-second,eq:V12-logdet-second}.  Since the
determinant is a cubic polynomial, its unnormalized derivatives are
polynomial everywhere.  The two exterior-form statements follow directly
from multilinearity of the wedge product.
\end{proof}

\begin{firewall}[Densitized closure is not normalized metric closure]
\label{fire:V12-densitized-scope}
A polynomial Palatini variation can have finite moments while a normalized
inverse-metric or unit-normal writer diverges.  Passing the densitized weak
equation therefore does not license the DeWitt inverse, a normalized stress
tensor, Brown--York unit normals, or an Einstein tensor on a collapsing
coframe sequence.
\end{firewall}
% -----------------------------------------------------------------------------
\subsection{A two-sided Orlicz bridge for inverse geometry}
\label{subsec:V12-orlicz-bridge}
% -----------------------------------------------------------------------------

For the rest of this version write
\[
 \Delta(q):=\det q,\qquad L(q):=\|q\|_{\mathrm{op}}.
\]
The collapse variable is $\Delta^{-1}$; it must not be confused with the
positive large-field radius.

\begin{lemma}[Spectral conversion from volume collapse to inverse writers]
\label{lem:V12-spectral-conversion}
Let $q=E^*E\succ0$ on a three-dimensional Euclidean fibre.  Then
\begin{align}
 \|\operatorname{adj}q\|_{\mathrm{op}}
 &\le L(q)^2,
 \label{eq:V12-adj-bound}\\
 \|q^{-1}\|_{\mathrm{op}}
 &\le \frac{L(q)^2}{\Delta(q)},
 \label{eq:V12-qinverse-bound}\\
 \|E^{-1}\|_{\mathrm{op}}
 =\|Eq^{-1}\|_{\mathrm{op}}
 &\le \frac{L(q)}{\Delta(q)^{1/2}}.
 \label{eq:V12-Einverse-bound}
\end{align}
Consequently every monomial bounded by
$C L(q)^a\|q^{-1}\|_{\mathrm{op}}^b
 \|E^{-1}\|_{\mathrm{op}}^c$
is bounded by
\begin{equation}
 C L(q)^{a+2b+c}\Delta(q)^{-b-c/2}.
 \label{eq:V12-rational-monomial}
\end{equation}
\end{lemma}

\begin{proof}
Let $0<\lambda_1\le\lambda_2\le\lambda_3=L(q)$ be the eigenvalues of
$q$.  The eigenvalues of $\operatorname{adj}q$ are the pairwise products,
so their maximum is $\lambda_2\lambda_3\le L(q)^2$.  Moreover,
\[
 \|q^{-1}\|_{\mathrm{op}}
 =\lambda_1^{-1}
 =\frac{\lambda_2\lambda_3}{\lambda_1\lambda_2\lambda_3}
 \le\frac{L(q)^2}{\Delta(q)}.
\]
The smallest singular value of $E$ is $\lambda_1^{1/2}$.  Hence
\[
 \|E^{-1}\|_{\mathrm{op}}
 =\|Eq^{-1}\|_{\mathrm{op}}
 =\lambda_1^{-1/2}
 =\left(\frac{\lambda_2\lambda_3}{\Delta(q)}\right)^{1/2}
 \le\frac{L(q)}{\Delta(q)^{1/2}}.
\]
Multiplication of these inequalities proves
\eqref{eq:V12-rational-monomial}.
\end{proof}

\begin{theorem}[Joint large-field/collapse Orlicz export]
\label{thm:V12-joint-orlicz-export}
Let $(R,\Delta)$ be random variables with $R\ge0$ and $\Delta>0$ almost
surely.  Suppose that for some $\delta,\kappa,s>0$ and $K\ge1$,
\begin{equation}
 \mathbb E\exp\{\delta R+\kappa\Delta^{-s}\}\le K.
 \label{eq:V12-joint-orlicz}
\end{equation}
For $a,b\ge0$ define
\begin{align}
 A_a(\delta)&:=
 \begin{cases}
  1,&a=0,\\
  \bigl(a/(e\delta)\bigr)^a,&a>0,
 \end{cases}
 &
 B_b(\kappa,s)&:=
 \begin{cases}
  1,&b=0,\\
  \bigl(b/(e\kappa s)\bigr)^{b/s},&b>0.
 \end{cases}
 \label{eq:V12-ABconstants}
\end{align}
Then
\begin{equation}
 \mathbb E[R^a\Delta^{-b}]
 \le K A_a(\delta)B_b(\kappa,s).
 \label{eq:V12-mixed-moment}
\end{equation}
In particular, for every $p>0$,
\begin{equation}
 \mathbb E\Delta^{-p}
 \le K\left(\frac{p}{e\kappa s}\right)^{p/s}.
 \label{eq:V12-negative-det-moment}
\end{equation}
If $R\ge L(q)$ and $\Delta=\det q$, all finite powers of $q^{-1}$ and
$E^{-1}$ have finite moments, with explicit bounds obtained from
Lemma~\ref{lem:V12-spectral-conversion}.
\end{theorem}

\begin{proof}
For $x\ge0$ and $\alpha,\eta>0$, elementary differentiation gives
\[
 \sup_{x\ge0}x^\alpha e^{-\eta x}
 =\left(\frac{\alpha}{e\eta}\right)^\alpha.
\]
Put $Y=\Delta^{-s}$.  Since $\Delta^{-b}=Y^{b/s}$, pointwise
\[
 R^a\Delta^{-b}
 \le A_a(\delta)B_b(\kappa,s)
       e^{\delta R+\kappa Y}.
\]
Taking expectations proves \eqref{eq:V12-mixed-moment}; setting $a=0$
gives \eqref{eq:V12-negative-det-moment}.  The last assertion follows
from \eqref{eq:V12-rational-monomial}.
\end{proof}

\begin{corollary}[Separate positive and collapse exponentials may be joined]
\label{cor:V12-separate-to-joint}
If
\begin{equation}
 \mathbb E e^{2\delta R}\le K_+,
 \qquad
 \mathbb E e^{2\kappa\Delta^{-s}}\le K_-,
 \label{eq:V12-separate-orlicz}
\end{equation}
then
\begin{equation}
 \mathbb E e^{\delta R+\kappa\Delta^{-s}}
 \le(K_+K_-)^{1/2}.
 \label{eq:V12-holder-join}
\end{equation}
No independence is required.
\end{corollary}

\begin{proof}
Apply Cauchy--Schwarz to
$e^{\delta R}e^{\kappa\Delta^{-s}}$.
\end{proof}

\begin{assessment}[The exact new input]
\label{ass:V12-exact-new-input}
The positive-field estimate of Versions V10--V11 can provide the first
bound in \eqref{eq:V12-separate-orlicz}, after enlarging $R$ to include
the coframe operator norm if that field is not already compact.  What is
still absent from the attached RPESM construction is the second bound:
a cofinal estimate of
$\mathbb E\exp\{2\kappa(\det q)^{-s}\}$, or a deterministic frame floor
which makes it automatic.
\end{assessment}
% -----------------------------------------------------------------------------
\subsection{Represented geometric score and Fisher control}
\label{subsec:V12-geometric-score}
% -----------------------------------------------------------------------------

\begin{theorem}[Fisher bound for rational coframe and stress writers]
\label{thm:V12-geometric-Fisher}
Assume \eqref{eq:V12-joint-orlicz}, let $q=E^*E$, set
$\Delta=\det q$, and suppose $R\ge L(q)$.  Let $u$ be a vector in a
finite-dimensional command space and let $z(u)$ be an uncentred
directional log-density derivative satisfying
\begin{equation}
 |z(u)|
 \le C\|u\|(1+R^m)
       (1+\|q^{-1}\|_{\mathrm{op}}^r)
 \qquad(m,r\ge0).
 \label{eq:V12-score-envelope}
\end{equation}
Define the probability score
$s(u):=z(u)-\mathbb E z(u)$.  Then
\begin{equation}
 \mathbb E|s(u)|^2
 \le \mathfrak C_{m,r}\|u\|^2,
 \label{eq:V12-score-Fisher}
\end{equation}
where the completely explicit constant
\begin{align}
 \mathfrak C_{m,r}
 :=4C^2\bigl\{&
 1+K A_{2m}(\delta)
 +K A_{4r}(\delta)B_{2r}(\kappa,s)
 \notag\\
 &+K A_{2m+4r}(\delta)B_{2r}(\kappa,s)
 \bigr\}.
 \label{eq:V12-geometric-constant}
\end{align}
Thus the Fisher quadratic form
$\mathcal I(u,v):=\mathbb E[s(u)s(v)]$ is finite and
$|\mathcal I(u,v)|\le\mathfrak C_{m,r}\|u\|\|v\|$.

The same conclusion holds for a finite sum of rational frame,
inverse-metric, normalized-volume, unit-normal, and represented stress
components whenever their tensor contractions imply envelopes of the
form \eqref{eq:V12-score-envelope}.
\end{theorem}

\begin{proof}
Lemma~\ref{lem:V12-spectral-conversion} and $R\ge L(q)$ give
\[
 \|q^{-1}\|_{\mathrm{op}}^r
 \le R^{2r}\Delta^{-r}.
\]
Twice using $(1+x)^2\le2(1+x^2)$ yields
\begin{align*}
 \frac{|z(u)|^2}{C^2\|u\|^2}
 &\le
 4(1+R^{2m})(1+R^{4r}\Delta^{-2r})\\
 &=4\{1+R^{2m}
       +R^{4r}\Delta^{-2r}
       +R^{2m+4r}\Delta^{-2r}\}.
\end{align*}
Theorem~\ref{thm:V12-joint-orlicz-export} bounds the last three
expectations by the three corresponding terms in
\eqref{eq:V12-geometric-constant}.  Finally,
\[
 \mathbb E|s(u)|^2=\operatorname{Var}z(u)
 \le\mathbb E|z(u)|^2.
\]
Cauchy--Schwarz gives the bilinear bound.  For finitely many tensor
components, sum their scalar envelopes and apply
$(\sum_{j=1}^Jx_j)^2\le J\sum_{j=1}^Jx_j^2$.
\end{proof}

\begin{corollary}[Joining the V11 positive sector to represented geometry]
\label{cor:V12-total-score-join}
Let $z_{\mathrm{old}}(u)$ denote any finite collection of Higgs,
line-determinant/Pfaffian, and Green-response score components already
controlled by the Version V11 positive-radius exponential estimate.
Let $z_{\mathrm{geo}}(u)$ satisfy \eqref{eq:V12-score-envelope}.  If the
coframe positive norm is included in the V11 radius and the collapse
estimate
\begin{equation}
 \sup_n\mathbb E_n
 \exp\{2\kappa(\det q_n)^{-s}\}\le K_-
 \label{eq:V12-cofinal-collapse-input}
\end{equation}
holds, then the centered total score associated with
$z_{\mathrm{old}}+z_{\mathrm{geo}}$ has a cutoff-uniform Fisher bound.

This conclusion requires no probabilistic independence between the
Standard Model fields and the coframe.
\end{corollary}

\begin{proof}
Use Corollary~\ref{cor:V12-separate-to-joint} to join the V11
positive-radius exponential with
\eqref{eq:V12-cofinal-collapse-input}.  Theorem
\ref{thm:V12-geometric-Fisher} controls the geometric term, and the V11
estimate controls the old term.  The inequality
$(x+y)^2\le2x^2+2y^2$ completes the proof.
\end{proof}

\begin{corollary}[The hard-floor route]
\label{cor:V12-hard-floor-route}
If $q\succeq cI$ almost surely with $c>0$, then
$\Delta(q)\ge c^3$ and $\|q^{-1}\|_{\mathrm{op}}\le c^{-1}$.
Accordingly, every score satisfying \eqref{eq:V12-score-envelope} is
controlled by a positive moment of $R$ alone.  In particular, the
uniform source-frame floor theorem of the attached monograph would
close the inverse-geometric part of
Corollary~\ref{cor:V12-total-score-join} if its hypothesis were proved
uniformly for the cofinal RPESM laws.
\end{corollary}

\begin{proof}
The eigenvalues of $q$ are at least $c$.  Their product is at least
$c^3$ and the largest eigenvalue of $q^{-1}$ is at most $c^{-1}$.
Substitute in \eqref{eq:V12-score-envelope}.
\end{proof}

% -----------------------------------------------------------------------------
\subsection{Failure certificate: a coframe-collapse concentration witness}
\label{subsec:V12-collapse-witness}
% -----------------------------------------------------------------------------

\begin{theorem}[Negative-volume divergence produces a collapse witness]
\label{thm:V12-collapse-witness}
Let $\mu_n$ be probability laws on positive-definite $3\times3$ matrices,
let $\Delta(q)=\det q$, and fix $p>0$.  Suppose
\[
 0<M_n:=\int\Delta(q)^{-p}\,\mu_n(\dd q)<\infty,
 \qquad M_n\longrightarrow\infty.
\]
Define the inverse-volume tilted probability law
\begin{equation}
 \nu_n(\dd q):=\frac{\Delta(q)^{-p}}{M_n}\,\mu_n(\dd q).
 \label{eq:V12-collapse-tilt}
\end{equation}
Then for every $\varepsilon>0$,
\begin{equation}
 \nu_n\{\Delta\ge\varepsilon\}
 \le\frac{\varepsilon^{-p}}{M_n}
 \longrightarrow0.
 \label{eq:V12-collapse-concentration}
\end{equation}
Thus $\nu_n$ concentrates on every neighbourhood of the degenerate-frame
stratum.  The pair $(\mu_n,\nu_n)$ is an explicit, normalized obstruction
certificate for any proposed uniform $p$th inverse-volume estimate.
\end{theorem}

\begin{proof}
On $\{\Delta\ge\varepsilon\}$ one has
$\Delta^{-p}\le\varepsilon^{-p}$.  Therefore
\[
 \nu_n\{\Delta\ge\varepsilon\}
 =M_n^{-1}\int_{\{\Delta\ge\varepsilon\}}
       \Delta^{-p}\,\dd\mu_n
 \le M_n^{-1}\varepsilon^{-p}.
\]
\end{proof}

\begin{corollary}[Infinite moments admit truncated witnesses]
\label{cor:V12-truncated-collapse-witness}
If $\int\Delta^{-p}\,\dd\mu=\infty$, put
$w_N=\min\{\Delta^{-p},N\}$ and
$\nu_N=w_N\mu/\int w_N\,\dd\mu$.  Then some sequence $N_j\to\infty$
satisfies
\[
 \nu_{N_j}\{\Delta\ge\varepsilon\}\longrightarrow0
 \qquad(\varepsilon>0).
\]
\end{corollary}

\begin{proof}
Monotone convergence gives $\int w_N\,\dd\mu\to\infty$, while
$\int_{\{\Delta\ge\varepsilon\}}w_N\,\dd\mu\le\varepsilon^{-p}$.
Apply the proof of Theorem~\ref{thm:V12-collapse-witness}; in fact the
conclusion holds along the full limit $N\to\infty$.
\end{proof}

\begin{assessment}[Operational value of the witness]
\label{ass:V12-witness-value}
A failed inverse estimate need not be recorded merely as divergence of
a number.  The tilted laws \eqref{eq:V12-collapse-tilt} identify the
ancestry fibres and shells carrying the failure.  They can therefore be
fed upstream to the scheduler/source-frame analysis: one tests whether
the mass comes from loss of source rank, from a rate/speed boundary
chart, or from a normalization introduced only at readout.
\end{assessment}
% -----------------------------------------------------------------------------
\subsection{Two populated acquisition mechanisms}
\label{subsec:V12-acquisition}
% -----------------------------------------------------------------------------

The preceding estimates are useful only if the collapse hypothesis can be
proved from a finite-cylinder law.  The next theorem gives a local sufficient
condition which is stable under arbitrary boundary data as long as two
potential bounds are uniform.

\begin{theorem}[Uniform singular-coframe barrier criterion]
\label{thm:V12-barrier-criterion}
Let
\[
 \mathsf G^+
 :=\{E\in\mathbb R^{3\times3}:\det E>0\},
 \qquad
 \Delta(E):=\det(E^*E)=(\det E)^2.
\]
Let $V_\omega:\mathsf G^+\to\mathbb R$ be measurable finite-dimensional
potentials indexed by boundary data $\omega$.  Suppose there are
$\alpha,\kappa_0,s>0$, constants $C,C_Q<\infty$, and a fixed measurable
set $Q\Subset\mathsf G^+$ of positive Lebesgue measure such that, for every
$\omega$,
\begin{align}
 V_\omega(E)
 &\ge \alpha\|E\|_{\mathrm F}^4
       +\kappa_0\Delta(E)^{-s}-C
 &&\text{for almost every }E,
 \label{eq:V12-barrier-lower}\\
 V_\omega(E)&\le C_Q
 &&\text{for almost every }E\in Q.
 \label{eq:V12-barrier-well}
\end{align}
Then
\[
 K_\omega(\dd E)
 :=Z_\omega^{-1}e^{-V_\omega(E)}
     \mathbf1_{\mathsf G^+}(E)\,\dd E
\]
is a probability kernel.  For every $0\le\delta<\alpha$ and
$0\le\kappa<\kappa_0$,
\begin{equation}
 \sup_\omega\int
 \exp\{\delta(1+\|E\|_{\mathrm F}^4)
              +\kappa\Delta(E)^{-s}\}\,
 K_\omega(\dd E)
 \le \overline K_{\delta,\kappa}<\infty,
 \label{eq:V12-uniform-barrier-Orlicz}
\end{equation}
where one may take
\begin{equation}
 \overline K_{\delta,\kappa}
 =
 \frac{e^{\delta+C+C_Q}}{|Q|}
 \int_{\mathbb R^9}
 e^{-(\alpha-\delta)\|E\|_{\mathrm F}^4}\,\dd E.
 \label{eq:V12-barrier-constant}
\end{equation}
Since
$1+\|E\|_{\mathrm F}^4\ge\|E^*E\|_{\mathrm{op}}$,
Theorems~\ref{thm:V12-joint-orlicz-export}
and \ref{thm:V12-geometric-Fisher} apply uniformly to this kernel.
\end{theorem}

\begin{proof}
The well condition gives
$Z_\omega\ge |Q|e^{-C_Q}>0$.  The lower bound gives
\[
 Z_\omega
 \le e^C\int_{\mathbb R^9}
       e^{-\alpha\|E\|_{\mathrm F}^4}\,\dd E<\infty.
\]
For the numerator in \eqref{eq:V12-uniform-barrier-Orlicz}, use
\eqref{eq:V12-barrier-lower} to obtain
\begin{align*}
 &\int_{\mathsf G^+}
 e^{\delta(1+\|E\|_{\mathrm F}^4)+\kappa\Delta^{-s}}
 e^{-V_\omega(E)}\,\dd E\\
 &\quad\le
 e^{\delta+C}\int_{\mathsf G^+}
 e^{-(\alpha-\delta)\|E\|_{\mathrm F}^4
    -(\kappa_0-\kappa)\Delta^{-s}}\,\dd E\\
 &\quad\le
 e^{\delta+C}\int_{\mathbb R^9}
 e^{-(\alpha-\delta)\|E\|_{\mathrm F}^4}\,\dd E.
\end{align*}
Divide by the lower bound for $Z_\omega$.  Finally, if
$x=\|E\|_{\mathrm F}^2$, then
$\|E^*E\|_{\mathrm{op}}\le x\le1+x^2
=1+\|E\|_{\mathrm F}^4$.
\end{proof}

\begin{example}[The collapse hypothesis is populated]
\label{ex:V12-populated-barrier}
For fixed $\alpha,\kappa_0,s>0$, the law
\begin{equation}
 \mu(\dd E)
 =
 Z^{-1}\exp\{-\alpha\|E\|_{\mathrm F}^4
              -\kappa_0(\det(E^*E))^{-s}\}
 \mathbf1_{\{\det E>0\}}\,\dd E
 \label{eq:V12-populated-law}
\end{equation}
satisfies \eqref{eq:V12-uniform-barrier-Orlicz} for every
$\delta<\alpha$ and $\kappa<\kappa_0$.  Hence all polynomial moments of
$E$, $E^{-1}$, $q$, and $q^{-1}$ are finite.  This example proves
non-vacuity of the analytic criterion; it does not assert that the
barrier term occurs in the physical RPESM cylinder.
\end{example}

A uniform one-step barrier is stronger than necessary.  A contracting
collapse drift also propagates a cofinal bound.

\begin{theorem}[Cofinal collapse Lyapunov recursion]
\label{thm:V12-collapse-Lyapunov}
Let $(\mathcal F_n)$ be a filtration and let
$Y_n=\Delta_n^{-s}\ge0$.  Suppose that for some
$\kappa>0$, $a,b\ge0$, and $0\le\rho<1$,
\begin{equation}
 \mathbb E\!\left[e^{\kappa Y_{n+1}}\mid\mathcal F_n\right]
 \le a e^{\kappa\rho Y_n}+b
 \quad\text{almost surely}.
 \label{eq:V12-collapse-drift}
\end{equation}
If $M_0:=\mathbb E e^{\kappa Y_0}<\infty$, then
\begin{equation}
 \sup_{n\ge0}\mathbb E e^{\kappa Y_n}\le B,
 \qquad
 B:=\max\!\left\{
 M_0,1,(2a)^{1/(1-\rho)},2b
 \right\}.
 \label{eq:V12-collapse-Lyapunov-bound}
\end{equation}
Consequently the uniform negative determinant moments
\begin{equation}
 \sup_n\mathbb E\Delta_n^{-p}
 \le B\left(\frac{p}{e\kappa s}\right)^{p/s}
 \qquad(p>0)
 \label{eq:V12-cofinal-negative-moments}
\end{equation}
hold.
\end{theorem}

\begin{proof}
Write $M_n=\mathbb E e^{\kappa Y_n}$.  Taking expectations in
\eqref{eq:V12-collapse-drift} and using concavity of $x^\rho$ gives
\[
 M_{n+1}\le a\,\mathbb E(e^{\kappa Y_n})^\rho+b
 \le aM_n^\rho+b.
\]
If $M_n\le B$, then the definition of $B$ gives
$aB^\rho\le B/2$ and $b\le B/2$, hence $M_{n+1}\le B$.
Induction proves \eqref{eq:V12-collapse-Lyapunov-bound}.
The moment estimate follows from
$\sup_{y\ge0}y^{p/s}e^{-\kappa y}
=(p/(e\kappa s))^{p/s}$.
\end{proof}

\begin{assessment}[Upstream form of the V12 bottleneck]
\label{ass:V12-upstream-bottleneck}
For the actual RPESM projective family, it is enough to prove one of the
following statements uniformly in the retained boundary:
\begin{enumerate}
 \item the existing source-frame hypothesis yields $q\succeq cI$;
 \item the conditional coframe potential satisfies
       \eqref{eq:V12-barrier-lower,eq:V12-barrier-well}; or
 \item the shell ancestry satisfies the weaker drift
       \eqref{eq:V12-collapse-drift}.
\end{enumerate}
If all three routes fail, Theorem~\ref{thm:V12-collapse-witness} returns
a normalized family concentrating on the precise collapse stratum which
must be traced into the source/scheduler map.  This is the upstream
replacement for assuming that finite-cutoff positivity persists
cofinally.
\end{assessment}
% -----------------------------------------------------------------------------
\subsection{Passage of noncollapse and rational writers to a weak limit}
\label{subsec:V12-limit-passage}
% -----------------------------------------------------------------------------

\begin{theorem}[The collapse barrier survives continuum extraction]
\label{thm:V12-limit-passage}
Let
\[
 \mathcal X=[0,\infty)\times\operatorname{Sym}^+_{\ge0}(3)
\]
with coordinates $(R,q)$, and let $\mu_n$ be probability laws supported
on $\det q>0$ and $R\ge\|q\|_{\mathrm{op}}$.  Suppose
\begin{equation}
 \sup_n\int
 e^{\delta R+\kappa(\det q)^{-s}}\,\dd\mu_n\le K
 \label{eq:V12-uniform-limit-Orlicz}
\end{equation}
for some $\delta,\kappa,s>0$.  If a subsequence converges weakly,
$\mu_{n_j}\Rightarrow\mu$, then
\begin{align}
 \mu\{\det q=0\}&=0,
 \label{eq:V12-limit-nondegenerate}\\
 \int e^{\delta R+\kappa(\det q)^{-s}}\,\dd\mu&\le K.
 \label{eq:V12-limit-Orlicz}
\end{align}
Moreover, let $W$ be continuous on $\{\det q>0\}$ and satisfy
\begin{equation}
 |W(R,q)|\le C(1+R^a(\det q)^{-b})
 \qquad(a,b\ge0).
 \label{eq:V12-rational-writer-growth}
\end{equation}
Then $\{W(R,q)\}_{\mu_n}$ is uniformly integrable and
\begin{equation}
 \int W\,\dd\mu_{n_j}\longrightarrow\int W\,\dd\mu.
 \label{eq:V12-writer-expectation-limit}
\end{equation}
The same holds for second moments and covariance bilinear forms of any
finite family of such writers.
\end{theorem}

\begin{proof}
Extend
\[
 F(R,q)=e^{\delta R+\kappa(\det q)^{-s}}
\]
to the singular stratum by setting $F=+\infty$ when $\det q=0$.
This is a nonnegative lower-semicontinuous function on $\mathcal X$.
The Portmanteau theorem and \eqref{eq:V12-uniform-limit-Orlicz} give
\[
 \int F\,\dd\mu
 \le\liminf_{j\to\infty}\int F\,\dd\mu_{n_j}\le K.
\]
A finite integral of $F$ forces
\eqref{eq:V12-limit-nondegenerate}, and also proves
\eqref{eq:V12-limit-Orlicz}.

Choose any $\eta>0$.  Theorem
\ref{thm:V12-joint-orlicz-export}, applied uniformly in $n$, bounds
\[
 \sup_n\int
 R^{a(1+\eta)}(\det q)^{-b(1+\eta)}\,\dd\mu_n<\infty.
\]
Together with \eqref{eq:V12-rational-writer-growth}, this gives a
uniform $(1+\eta)$ moment of $W$, hence uniform integrability.  The
limit law charges only the open set $\{\det q>0\}$, where $W$ is
continuous.  The continuous mapping theorem gives convergence in law
of the writers, and uniform integrability upgrades it to
\eqref{eq:V12-writer-expectation-limit}.  Applying the same argument
with an exponent strictly larger than $2$ proves convergence of second
moments.  Polarization gives covariance convergence.
\end{proof}

\begin{corollary}[Conditional represented-stress limit]
\label{cor:V12-stress-limit}
Under the hypotheses of Theorem~\ref{thm:V12-limit-passage}, every
finite-dimensional represented stress or geometric score writer whose
frame contractions obey \eqref{eq:V12-rational-writer-growth} has a
finite, subsequence-stable mean and covariance.  In particular, the
Fisher forms controlled in Theorem~\ref{thm:V12-geometric-Fisher}
survive weak continuum extraction whenever the score writers themselves
converge on the nondegenerate stratum.
\end{corollary}

\begin{proof}
Apply Theorem~\ref{thm:V12-limit-passage} to each writer and each
pairwise product.  Centering is a finite linear combination of their
first and second moments.
\end{proof}

\begin{firewall}[What V12 does not prove]
\label{fire:V12-no-full-continuum}
Theorem~\ref{thm:V12-limit-passage} is a conditional compactness
interface.  It does not construct the weakly convergent projective
subsequence, identify an RPESM coframe barrier, prove the Einstein
equation, or establish gauge/reflection-positive reconstruction.
Its role is narrower and indispensable: once a cofinal collapse
estimate is acquired, inverse geometry and represented stress cannot
silently disappear at the limiting step.
\end{firewall}

% -----------------------------------------------------------------------------
\subsection{V12 theorem ledger and next acquisition target}
\label{subsec:V12-ledger}
% -----------------------------------------------------------------------------

\paragraph{New conclusions proved in V12.}
Finite-cutoff frame positivity was separated from cofinal inverse control by
Proposition~\ref{prop:V12-positive-no-inverse}.  The polynomial
densitized sector was separated from normalized inverse readouts in
Proposition~\ref{prop:V12-polynomial-split}.  Lemma
\ref{lem:V12-spectral-conversion} reduces every rational coframe monomial
to positive coframe powers and negative determinant powers.  Theorem
\ref{thm:V12-joint-orlicz-export} exports all such moments from one
two-sided Orlicz estimate, and Theorem
\ref{thm:V12-geometric-Fisher} converts them into a represented-stress
Fisher bound.  Theorem~\ref{thm:V12-collapse-witness} supplies a
normalized collapse witness on failure.  Theorems
\ref{thm:V12-barrier-criterion} and
\ref{thm:V12-collapse-Lyapunov} give strong local and weaker cofinal
routes for acquiring the missing estimate.  Theorem
\ref{thm:V12-limit-passage} proves that noncollapse and rational writers
survive weak extraction.

\paragraph{Imported conclusions not reproved.}
The Gram/exterior determinant identity, the normalized DeWitt formula, and
the conditional hard source-frame floor are used from the attached grand
tensor manuscript.  Positive-field, determinant/Pfaffian-line, Green, and
ancestry estimates are used from Versions V8--V11.

\paragraph{Unresolved physical acquisition.}
The actual finite-cylinder RPESM coframe kernel has not yet been shown to
satisfy a hard floor, the barrier criterion, or the collapse Lyapunov drift.
Therefore the full Standard Model--Einstein continuum is not claimed.
The next note must inspect the scheduler-to-coframe/source map upstream and
derive one of these three statements from its physical finite data.  If that
derivation fails, the tilt \eqref{eq:V12-collapse-tilt} should be pulled back
through the source map to determine whether rank loss, a chart boundary, or
readout normalization is responsible.  Only after that acquisition should
the programme assemble the cofinal Einstein stress identity and its Ward
compatibility with the already controlled Standard Model sectors.

% =============================================================================
% END APPENDED VERSION V12 MATERIAL
% =============================================================================


% =============================================================================
% BEGIN APPENDED VERSION V13 MATERIAL
% =============================================================================

\clearpage
\section{Version V13: scheduler scale--shape noncollapse and full ADM inverse control}
\label{sec:V13-scheduler-noncollapse}

Version V12 reduced normalized geometric and represented-stress control to a
cofinal exclusion of coframe collapse.  The appropriate upstream variables are
already present in the attached grand tensor manuscript.  Its scheduler theorem
defines a positive predictable bracket $\mathcal B_h$, a positive speed density
$\varrho_h$, and
\[
 q_h=g_h
 =\varrho_h^{2/3}(\det\mathcal B_h)^{1/3}\mathcal B_h^{-1},
 \qquad
 N_h=\varrho_h^{1/3}(\det\mathcal B_h)^{1/6}.
\]
The manuscript proves that this is an analytic ADM parametrization and gives
explicit flat and de~Sitter vacuum schedulers.  It does not separate the four
different ways in which the interacting projective law can leave every compact
nondegenerate ADM set.

This version derives that separation.  Spatial volume, bracket scale, bracket
shape, and shift are distinct control rows.  The reference scheduler lies in
the interior of all four, but compactness of the scheduler domain at each
finite cutoff does not by itself give a common compact subset of the open
positive chart.  The results below turn the missing cofinal assertion into
quantitative scheduler inequalities and normalized obstruction laws.  The
existing vacuum realizations and the existing global Einstein-or-obstruction
theorem are imported rather than reproved.

% -----------------------------------------------------------------------------
\subsection{Exact scheduler scale--shape coordinates}
\label{subsec:V13-scale-shape}
% -----------------------------------------------------------------------------

\begin{definition}[Bracket volume and unimodular shape]
\label{def:V13-bracket-shape}
For $\mathcal B_h\in\operatorname{Sym}_3^+$ set
\begin{equation}
 \upsilon_h:=(\det\mathcal B_h)^{1/3},
 \qquad
 \mathcal S_h:=\upsilon_h^{-1}\mathcal B_h.
 \label{eq:V13-shape-definition}
\end{equation}
Then $\det\mathcal S_h=1$.  We call $\upsilon_h$ the bracket-volume scale and
$\mathcal S_h$ the bracket shape.  Let $q_h$ denote the spatial metric, let
$E_h=q_h^{1/2}$ be its positive spatial coframe, and define the time-gauge
coframe density
\begin{equation}
 \Omega_h:=N_h\det E_h=N_h\sqrt{\det q_h}.
 \label{eq:V13-Omega-definition}
\end{equation}
\end{definition}

\begin{theorem}[Exact scheduler factorization of spatial and spacetime volume]
\label{thm:V13-exact-factorization}
The scheduler ADM formula is equivalent to
\begin{align}
 q_h&=\varrho_h^{2/3}\mathcal S_h^{-1},
 &
 q_h^{-1}&=\varrho_h^{-2/3}\mathcal S_h,
 \label{eq:V13-q-shape}\\
 N_h&=\varrho_h^{1/3}\upsilon_h^{1/2},
 &
 \det q_h&=\varrho_h^2,
 \label{eq:V13-lapse-volume}\\
 \det E_h&=\varrho_h,
 &
 \Omega_h&=\varrho_h^{4/3}\upsilon_h^{1/2}.
 \label{eq:V13-coframe-volume}
\end{align}
For the ADM Lorentzian metric
\begin{equation}
 \mathbf g_h=
 \begin{pmatrix}
 -N_h^2+\beta_h^{\mathsf T}q_h\beta_h&
 \beta_h^{\mathsf T}q_h\\
 q_h\beta_h&q_h
 \end{pmatrix},
 \label{eq:V13-ADM-matrix}
\end{equation}
one has
\begin{equation}
 \det\mathbf g_h
 =-N_h^2\det q_h
 =-\varrho_h^{8/3}\upsilon_h,
 \qquad
 |\det\mathbf g_h|=\Omega_h^2.
 \label{eq:V13-full-determinant}
\end{equation}
Consequently:
\begin{enumerate}[label=\textup{(\roman*)},leftmargin=2.8em]
\item spatial determinant collapse is exactly $\varrho_h\downarrow0$;
\item lapse collapse is controlled jointly by
      $\varrho_h\downarrow0$ and $\upsilon_h\downarrow0$;
\item spatial anisotropy is entirely carried by $\mathcal S_h$;
\item the shift changes neither determinant.
\end{enumerate}
\end{theorem}

\begin{proof}
Write $\mathcal B_h=\upsilon_h\mathcal S_h$.  Since
$\det\mathcal S_h=1$, substitution into the imported scheduler formula gives
\[
 q_h
 =\varrho_h^{2/3}\upsilon_h
   (\upsilon_h\mathcal S_h)^{-1}
 =\varrho_h^{2/3}\mathcal S_h^{-1},
\]
and inversion gives the second identity in
\eqref{eq:V13-q-shape}.  Taking a determinant yields
\[
 \det q_h=(\varrho_h^{2/3})^3\det(\mathcal S_h^{-1})
 =\varrho_h^2.
\]
The lapse formula is
$N_h=\varrho_h^{1/3}\upsilon_h^{1/2}$.
The positive square root $E_h=q_h^{1/2}$ has
$\det E_h=\sqrt{\det q_h}=\varrho_h$, proving
\eqref{eq:V13-coframe-volume}.

Let
\[
 L_{\beta_h}:=
 \begin{pmatrix}1&0\\ \beta_h&I_3\end{pmatrix},
 \qquad
 D_h:=\operatorname{diag}(-N_h^2,q_h).
\]
Then \eqref{eq:V13-ADM-matrix} is
$\mathbf g_h=L_{\beta_h}^{\mathsf T}D_hL_{\beta_h}$ and
$\det L_{\beta_h}=1$.  Hence
$\det\mathbf g_h=-N_h^2\det q_h$; substitution gives
\eqref{eq:V13-full-determinant}.  The four conclusions now follow directly.
\end{proof}

\begin{theorem}[Quantitative full inverse and inverse-coframe bounds]
\label{thm:V13-full-inverse}
With the notation of Theorem~\ref{thm:V13-exact-factorization},
\begin{align}
 \|\mathbf g_h^{-1}\|_{\mathrm{op}}
 &\le
 (1+\|\beta_h\|)^2\varrho_h^{-2/3}
 \max\{\upsilon_h^{-1},\|\mathcal S_h\|_{\mathrm{op}}\},
 \label{eq:V13-full-inverse-bound}\\
 \|\mathsf e_h^{-1}\|_{\mathrm{op}}
 &\le
 (1+\|\beta_h\|)\varrho_h^{-1/3}
 \max\{\upsilon_h^{-1/2},
          \|\mathcal S_h\|_{\mathrm{op}}^{1/2}\},
 \label{eq:V13-full-coframe-inverse-bound}
\end{align}
where
\begin{equation}
 \mathsf e_h:=
 \operatorname{diag}(N_h,E_h)L_{\beta_h}
 \label{eq:V13-ADM-coframe}
\end{equation}
is the time-gauge coframe matrix.  Moreover
\begin{align}
 \|q_h\|_{\mathrm{op}}
 &=\varrho_h^{2/3}\|\mathcal S_h^{-1}\|_{\mathrm{op}}
 \le\varrho_h^{2/3}\|\mathcal S_h\|_{\mathrm{op}}^2,
 \label{eq:V13-q-positive-bound}\\
 \|q_h^{-1}\|_{\mathrm{op}}
 &=\varrho_h^{-2/3}\|\mathcal S_h\|_{\mathrm{op}}.
 \label{eq:V13-q-inverse-exact}
\end{align}
\end{theorem}

\begin{proof}
The factorization in the preceding proof gives
\[
 \mathbf g_h^{-1}
 =L_{\beta_h}^{-1}
  \operatorname{diag}(-N_h^{-2},q_h^{-1})
  L_{\beta_h}^{-\mathsf T}.
\]
Since $L_{\beta_h}^{-1}=I+K_{\beta_h}$ with
$\|K_{\beta_h}\|_{\mathrm{op}}=\|\beta_h\|$, its norm is at most
$1+\|\beta_h\|$.  Equations
\eqref{eq:V13-q-shape,eq:V13-lapse-volume} then prove
\eqref{eq:V13-full-inverse-bound}.

Likewise,
$\mathsf e_h^{-1}=L_{\beta_h}^{-1}
 \operatorname{diag}(N_h^{-1},E_h^{-1})$.
Because
\[
 \|E_h^{-1}\|_{\mathrm{op}}
 =\|q_h^{-1}\|_{\mathrm{op}}^{1/2}
 =\varrho_h^{-1/3}\|\mathcal S_h\|_{\mathrm{op}}^{1/2},
\]
the coframe estimate follows.  The equalities in
\eqref{eq:V13-q-positive-bound,eq:V13-q-inverse-exact} follow from
\eqref{eq:V13-q-shape}.  Finally, if the eigenvalues of
$\mathcal S_h$ are $0<s_1\le s_2\le s_3$ with product one, then
$s_1^{-1}=s_2s_3\le s_3^2$, which proves the last inequality in
\eqref{eq:V13-q-positive-bound}.
\end{proof}

\begin{assessment}[What the V12 determinant variable sees]
\label{ass:V13-V12-translation}
For the spatial Gram used in V12,
\[
 \Delta_h=\det q_h=\varrho_h^2.
\]
Thus the V12 collapse input
$\mathbb E\exp\{\kappa\Delta_h^{-s}\}<\infty$ is exactly
$\mathbb E\exp\{\kappa\varrho_h^{-2s}\}<\infty$.
Full spacetime writers additionally require negative
$\upsilon_h$ moments, while anisotropic inverse writers require positive
$\|\mathcal S_h\|$ moments and coordinate components require shift moments.
A spatial determinant estimate alone cannot control those three rows.
\end{assessment}
% -----------------------------------------------------------------------------
\subsection{Four independent scheduler degenerations}
\label{subsec:V13-four-degenerations}
% -----------------------------------------------------------------------------

\begin{proposition}[Exact separators for scale, speed, shape, and shift]
\label{prop:V13-four-separators}
Every item below consists of strictly positive finite scheduler data.
Nevertheless the indicated normalized writer diverges.
\begin{enumerate}[label=\textup{(\alph*)},leftmargin=2.8em]
\item Set
      $\varrho_n=1$, $\mathcal S_n=I_3$,
      $\upsilon_n=n^{-2}$, and $\beta_n=0$.  Then
      \[
       q_n=I_3,\qquad \det q_n=1,\qquad
       N_n=n^{-1},\qquad |\det\mathbf g_n|=n^{-2}.
       \tag{V13.S1}
      \]
      Thus the spatial geometry stays fixed while the lapse and spacetime
      coframe collapse.
\item Set
      \[
       \varrho_n=\upsilon_n=1,\qquad
       \mathcal S_n=\operatorname{diag}(n^{-2},n,n),\qquad
       \beta_n=0.
      \]
      Then
      \[
       \det q_n=|\det\mathbf g_n|=1,\qquad
       q_n=\operatorname{diag}(n^2,n^{-1},n^{-1}),
       \qquad
       \|q_n^{-1}\|_{\mathrm{op}}=n.
       \tag{V13.S2}
      \]
      Fixed spatial and spacetime volume therefore do not exclude
      anisotropic coframe collapse.
\item Set
      $\varrho_n=n^{-3}$, $\upsilon_n=1$,
      $\mathcal S_n=I_3$, and $\beta_n=0$.  Then
      \[
       q_n=n^{-2}I_3,\quad N_n=n^{-1},\quad
       \det q_n=n^{-6},\quad
       |\det\mathbf g_n|=n^{-8}.
       \tag{V13.S3}
      \]
      This is pure speed-density collapse.
\item Set
      $\varrho_n=\upsilon_n=1$, $\mathcal S_n=I_3$, and
      $\beta_n=ne_1$.  Then
      \[
       \det q_n=|\det\mathbf g_n|=1,
       \qquad
       \|\mathbf g_n^{-1}\|_{\mathrm{op}}\ge n^2-1.
       \tag{V13.S4}
      \]
      The coordinate inverse diverges through shift escape even though the
      coframe determinant stays fixed.
\end{enumerate}
No one of the four control rows can therefore be silently replaced by the
other three.
\end{proposition}

\begin{proof}
The first three displays follow by direct substitution in
\eqref{eq:V13-q-shape,eq:V13-lapse-volume,eq:V13-full-determinant}.
For the fourth, the inverse factorization gives
\[
 \mathbf g_n^{-1}
 =
 \begin{pmatrix}
 -1&\beta_n^{\mathsf T}\\
 \beta_n&I_3-\beta_n\beta_n^{\mathsf T}
 \end{pmatrix}.
\]
Its spatial compression has eigenvalue $1-n^2$ in the $e_1$ direction, so
the operator norm is at least $n^2-1$.
\end{proof}

\begin{firewall}[A faithful ADM chart is not a uniform noncollapse estimate]
\label{fire:V13-open-chart}
Analytic invertibility on the open cone
$\{\mathcal B\succ0,\varrho>0\}$ is a pointwise statement.  Each sequence in
Proposition~\ref{prop:V13-four-separators} remains in that cone.  Likewise,
compact support at each cutoff allows the supporting compact sets to approach
the cone boundary or escape in shape and shift.  A cofinal hard floor requires
one common interior tube or the probabilistic Orlicz alternative below.
\end{firewall}

% -----------------------------------------------------------------------------
\subsection{A quantitative tube around the RPESM reference scheduler}
\label{subsec:V13-reference-tube}
% -----------------------------------------------------------------------------

\begin{theorem}[Reference-tube hard noncollapse]
\label{thm:V13-reference-tube}
Fix $0\le\eta<1$, $0\le r<1$, and $b<\infty$.  Suppose
\begin{equation}
 \|\mathcal B_h-I_3\|_{\mathrm{op}}\le\eta,
 \qquad
 |\varrho_h-1|\le r,
 \qquad
 \|\beta_h\|\le b
 \label{eq:V13-reference-tube}
\end{equation}
throughout a family of scheduler configurations.  Put
\begin{equation}
 a_\eta:=\frac{1-\eta}{1+\eta}.
 \label{eq:V13-aeta}
\end{equation}
Then
\begin{align}
 1-\eta\le\upsilon_h&\le1+\eta,
 &
 a_\eta I_3\preceq\mathcal S_h&\preceq a_\eta^{-1}I_3,
 \label{eq:V13-tube-shape}\\
 (1-r)^{2/3}a_\eta I_3
 \preceq q_h
 &\preceq
 (1+r)^{2/3}a_\eta^{-1}I_3,
 \label{eq:V13-tube-q}\\
 (1-r)^{1/3}(1-\eta)^{1/2}
 \le N_h
 &\le
 (1+r)^{1/3}(1+\eta)^{1/2},
 \label{eq:V13-tube-lapse}\\
 |\det\mathbf g_h|
 &\ge(1-r)^{8/3}(1-\eta).
 \label{eq:V13-tube-fourdet}
\end{align}
In addition,
\begin{align}
 \|\mathbf g_h^{-1}\|_{\mathrm{op}}
 &\le
 (1+b)^2(1-r)^{-2/3}a_\eta^{-1},
 \label{eq:V13-tube-full-inverse}\\
 \|\mathsf e_h^{-1}\|_{\mathrm{op}}
 &\le
 (1+b)(1-r)^{-1/3}a_\eta^{-1/2}.
 \label{eq:V13-tube-coframe-inverse}
\end{align}
Thus a cutoff-independent tube of the form
\eqref{eq:V13-reference-tube} supplies the deterministic hard-floor route
required in V12 for all spatial and full ADM inverse writers.
\end{theorem}

\begin{proof}
The eigenvalues of $\mathcal B_h$ lie in
$[1-\eta,1+\eta]$.  Their geometric mean
$\upsilon_h$ lies in the same interval.  Dividing an eigenvalue by that
geometric mean proves \eqref{eq:V13-tube-shape}.  Equation
\eqref{eq:V13-q-shape} then gives \eqref{eq:V13-tube-q}, and the lapse formula
gives \eqref{eq:V13-tube-lapse}.  The determinant bound follows from
\eqref{eq:V13-full-determinant}.

For the inverse metric, Theorem~\ref{thm:V13-full-inverse} and
\eqref{eq:V13-tube-shape} give
\[
 \|\mathbf g_h^{-1}\|_{\mathrm{op}}
 \le(1+b)^2(1-r)^{-2/3}
 \max\{(1-\eta)^{-1},a_\eta^{-1}\}.
\]
Since $a_\eta^{-1}=(1+\eta)/(1-\eta)\ge(1-\eta)^{-1}$, this is
\eqref{eq:V13-tube-full-inverse}.  The same argument with square roots proves
\eqref{eq:V13-tube-coframe-inverse}.
\end{proof}

\begin{corollary}[The displayed reference and vacuum schedulers pass locally]
\label{cor:V13-reference-passes}
The RPESM reference values
$\mathcal B_h=I_3$, $\varrho_h=1$, $\beta_h=0$ pass
Theorem~\ref{thm:V13-reference-tube} with $\eta=r=b=0$.
The attached flat-vacuum family passes globally, and its homogeneous
de~Sitter family passes on every compact time slab with constants depending
only on the slab and $H$.
\end{corollary}

\begin{proof}
The reference and flat statements are immediate.  On a compact de~Sitter
slab, the imported formulas
$\mathcal B_h(t)=e^{-2Ht}I_3$,
$\varrho_h(t)=e^{3Ht}$, and $\beta_h(t)=0$ have positive upper and lower
bounds independent of $h$.  Applying
Theorems~\ref{thm:V13-exact-factorization}
and \ref{thm:V13-full-inverse} gives the claim.  No interacting-field
conclusion is used.
\end{proof}
% -----------------------------------------------------------------------------
\subsection{Root-rate conditions which imply scheduler noncollapse}
\label{subsec:V13-root-rates}
% -----------------------------------------------------------------------------

Let $\mathscr R=\{v_i\}_{i=1}^m\subset\mathbb R^3$ be the directed root list,
including repeated opposite orientations when they use the same rank-one
matrix.  Set
\begin{equation}
 \mathcal G_{\mathscr R}:=\sum_{i=1}^m v_iv_i^{\mathsf T},
 \qquad
 \gamma_{\mathscr R}:=\lambda_{\min}(\mathcal G_{\mathscr R})>0,
 \qquad
 M_{\mathscr R}:=\max_i\|v_i\|^2.
 \label{eq:V13-root-constants}
\end{equation}
For rates $k_{i,h}>0$, define
\begin{equation}
 K_h:=\sum_i k_{i,h},\qquad
 p_{i,h}:=\frac{k_{i,h}}{K_h},\qquad
 t_h:=h^2K_h,\qquad
 \mathcal A_h:=\sum_i p_{i,h}v_iv_i^{\mathsf T}.
 \label{eq:V13-rate-coordinates}
\end{equation}
Then $\mathcal B_h=t_h\mathcal A_h$.

\begin{theorem}[Rate-share, scale, and shape certificate]
\label{thm:V13-rate-certificate}
Suppose
\begin{equation}
 p_{i,h}\ge\varepsilon>0
 \qquad(1\le i\le m)
 \label{eq:V13-rate-share-floor}
\end{equation}
uniformly in $h$, and put
$m_\varepsilon:=\varepsilon\gamma_{\mathscr R}$.
Then
\begin{align}
 m_\varepsilon I_3
 \preceq\mathcal A_h
 &\preceq M_{\mathscr R}I_3,
 \label{eq:V13-A-bounds}\\
 t_hm_\varepsilon
 \le\upsilon_h
 &\le t_hM_{\mathscr R},
 \label{eq:V13-upsilon-rate}\\
 \frac{m_\varepsilon}{M_{\mathscr R}}I_3
 \preceq\mathcal S_h
 &\preceq
 \frac{M_{\mathscr R}}{m_\varepsilon}I_3.
 \label{eq:V13-shape-rate}
\end{align}
The winner--waiting log-rate Fisher matrix imported from the scheduler theorem
is $\operatorname{diag}(p_{1,h},\ldots,p_{m,h})$; hence
\eqref{eq:V13-rate-share-floor} is simultaneously a uniform score-frame floor
and a sufficient bracket-shape certificate.

If the roots are paired as $r_a$ with directed probabilities $p_{a,h}^\pm$,
define
\begin{equation}
 d_h:=\sum_a(p_{a,h}^+-p_{a,h}^-)r_a.
 \label{eq:V13-directed-imbalance}
\end{equation}
Then the ADM shift satisfies the exact scaling law
\begin{equation}
 \beta_h=\frac{t_h}{h}d_h.
 \label{eq:V13-shift-scaling}
\end{equation}
Thus, when $t_h$ stays bounded above and below, bounded shift requires and is
equivalent to $d_h=O(h)$.
\end{theorem}

\begin{proof}
Since $p_{i,h}\ge\varepsilon$,
\[
 \mathcal A_h
 \succeq\varepsilon\sum_i v_iv_i^{\mathsf T}
 \succeq\varepsilon\gamma_{\mathscr R}I_3.
\]
Also
$v_iv_i^{\mathsf T}\preceq\|v_i\|^2I_3
 \preceq M_{\mathscr R}I_3$, and the $p_{i,h}$ sum to one.
This proves \eqref{eq:V13-A-bounds}.  The geometric mean of the eigenvalues of
$\mathcal A_h$ lies in $[m_\varepsilon,M_{\mathscr R}]$.  Since
\[
 \upsilon_h
 =t_h(\det\mathcal A_h)^{1/3},
 \qquad
 \mathcal S_h
 =\frac{\mathcal A_h}{(\det\mathcal A_h)^{1/3}},
\]
the next two bounds follow.

The Fisher identity is part of the imported scheduler calculation; its least
eigenvalue is $\min_i p_{i,h}$.  Finally,
\[
 \beta_h
 =h\sum_a(k_{a,h}^+-k_{a,h}^-)r_a
 =hK_hd_h=\frac{t_h}{h}d_h,
\]
which proves \eqref{eq:V13-shift-scaling}.
\end{proof}

\begin{corollary}[A rate-level hard ADM floor]
\label{cor:V13-rate-hard-floor}
In addition to \eqref{eq:V13-rate-share-floor}, suppose
\begin{equation}
 0<t_-\le t_h\le t_+<\infty,\qquad
 0<\varrho_-\le\varrho_h\le\varrho_+<\infty,\qquad
 \|d_h\|\le \frac{b h}{t_h}.
 \label{eq:V13-rate-hard-inputs}
\end{equation}
Then
\begin{align}
 q_h&\succeq
 \varrho_-^{2/3}\frac{m_\varepsilon}{M_{\mathscr R}}I_3,
 \label{eq:V13-rate-q-floor}\\
 N_h&\ge
 \varrho_-^{1/3}(t_-m_\varepsilon)^{1/2},
 \label{eq:V13-rate-lapse-floor}
\end{align}
and $\|\mathbf g_h^{-1}\|_{\mathrm{op}}$ and
$\|\mathsf e_h^{-1}\|_{\mathrm{op}}$ have cutoff-independent upper bounds.
\end{corollary}

\begin{proof}
Equation \eqref{eq:V13-shift-scaling} gives $\|\beta_h\|\le b$.
The lower metric bound follows from
$q_h=\varrho_h^{2/3}\mathcal S_h^{-1}$ and the upper shape bound in
\eqref{eq:V13-shape-rate}.  Equations
\eqref{eq:V13-upsilon-rate,eq:V13-lapse-volume} give the lapse floor.
Theorem~\ref{thm:V13-full-inverse} supplies the inverse bounds.
\end{proof}

\begin{remark}[The certificate is sufficient rather than necessary]
\label{rem:V13-share-stronger}
A bracket can remain positive when a redundant directed root probability tends
to zero, so a spatial shape floor need not imply the full directed-rate Fisher
floor.  Conversely, a uniform directed-rate Fisher floor controls shape but
does not control $t_h$, $\varrho_h$, or the $O(h)$ cancellation in
\eqref{eq:V13-directed-imbalance}.  These remain independent physical rows.
\end{remark}

\begin{example}[The monograph's reference normalization]
\label{ex:V13-A3-reference}
For the twelve standard oriented $A_3$ roots,
\[
 \sum_{i=1}^{12}v_iv_i^{\mathsf T}=8I_3,
 \qquad M_{\mathscr R}=2.
\]
At $k_{i,h}=1/(8h^2)$,
\[
 p_{i,h}=\frac1{12},\qquad
 t_h=\frac32,\qquad
 \mathcal A_h=\frac23I_3,\qquad
 \mathcal B_h=I_3,\qquad d_h=0.
\]
This recovers the exact reference ADM values.  What the interacting
finite-cylinder definition does not yet state is that its scheduler law obeys
\eqref{eq:V13-rate-hard-inputs}, or a probabilistic replacement, uniformly
over the projective family.
\end{example}
% -----------------------------------------------------------------------------
\subsection{A full scheduler Orlicz packet}
\label{subsec:V13-Orlicz-packet}
% -----------------------------------------------------------------------------

\begin{lemma}[Finite multivariable Orlicz export]
\label{lem:V13-multi-Orlicz}
Let $X_1,\ldots,X_J\ge0$ and suppose
\begin{equation}
 \mathbb E\exp\!\left\{\sum_{j=1}^Jc_jX_j^{\alpha_j}\right\}
 \le K,
 \qquad c_j,\alpha_j>0.
 \label{eq:V13-multi-Orlicz}
\end{equation}
For $p=(p_1,\ldots,p_J)\in[0,\infty)^J$, define
\begin{equation}
 \mathfrak M(p):=
 K\prod_{j:p_j>0}
 \left(\frac{p_j}{e c_j\alpha_j}\right)^{p_j/\alpha_j}.
 \label{eq:V13-multi-moment-constant}
\end{equation}
Then
\begin{equation}
 \mathbb E\prod_{j=1}^JX_j^{p_j}\le\mathfrak M(p).
 \label{eq:V13-multi-moment}
\end{equation}
\end{lemma}

\begin{proof}
For $x\ge0$,
\[
 \sup_{x\ge0}x^{p_j}e^{-c_jx^{\alpha_j}}
 =
 \left(\frac{p_j}{e c_j\alpha_j}\right)^{p_j/\alpha_j}
\]
when $p_j>0$, while the factor is one for $p_j=0$.  Multiply these pointwise
bounds by the exponential in \eqref{eq:V13-multi-Orlicz} and take
expectations.
\end{proof}

For a scheduler-coupled Standard Model cylinder, set
\begin{equation}
 \begin{aligned}
 X_{1,h}&=R_{{\rm SM},h},&
 X_{2,h}&=\varrho_h,&
 X_{3,h}&=\varrho_h^{-1},\\
 X_{4,h}&=\upsilon_h,&
 X_{5,h}&=\upsilon_h^{-1},&
 X_{6,h}&=\|\mathcal S_h\|_{\mathrm{op}},&
 X_{7,h}&=1+\|\beta_h\|.
 \end{aligned}
 \label{eq:V13-packet-variables}
\end{equation}
Here $R_{{\rm SM},h}$ is the positive field/line/Green radius already used in
Versions V10--V11.

\begin{definition}[Full scheduler Orlicz packet]
\label{def:V13-full-packet}
A projective family has a full scheduler Orlicz packet when there are
$c_j,\alpha_j>0$ and $K<\infty$, independent of $h$, such that
\begin{equation}
 \sup_h\mathbb E_h
 \exp\!\left\{\sum_{j=1}^7c_jX_{j,h}^{\alpha_j}\right\}
 \le K.
 \label{eq:V13-full-packet}
\end{equation}
\end{definition}

\begin{theorem}[The packet controls every finite rational ADM writer]
\label{thm:V13-packet-controls-ADM}
Assume \eqref{eq:V13-full-packet}.  Every finite tensor contraction built
polynomially from the already controlled Standard Model fields and from
\[
 q_h,\ q_h^{-1},\ E_h,\ E_h^{-1},\
 N_h,\ N_h^{-1},\ \beta_h,\ \mathsf e_h,\ \mathsf e_h^{-1}
\]
has moments of every finite order, uniformly in $h$.

More explicitly, suppose an uncentred directional score has an algebraic
envelope
\begin{equation}
 |z_h(u)|
 \le\|u\|\sum_{\ell=1}^LC_\ell
       \prod_{j=1}^7X_{j,h}^{p_{\ell j}},
 \qquad p_{\ell j}\ge0.
 \label{eq:V13-score-monomial-envelope}
\end{equation}
Then its centered Fisher form satisfies
\begin{equation}
 \mathbb E_h|z_h(u)-\mathbb E_hz_h(u)|^2
 \le\mathfrak C_{\rm ADM}\|u\|^2
 \label{eq:V13-ADM-Fisher}
\end{equation}
uniformly, where
\begin{equation}
 \mathfrak C_{\rm ADM}
 :=
 L\sum_{\ell=1}^LC_\ell^2\,
 \mathfrak M(2p_{\ell1},\ldots,2p_{\ell7}).
 \label{eq:V13-ADM-Fisher-constant}
\end{equation}
The associated covariance bilinear form is bounded by
$\mathfrak C_{\rm ADM}\|u\|\|v\|$.
\end{theorem}

\begin{proof}
Theorems~\ref{thm:V13-exact-factorization}
and \ref{thm:V13-full-inverse} bound every listed geometric norm by a finite
sum of monomials in $X_{2,h},\ldots,X_{7,h}$.  For example,
\begin{align*}
 \|q_h\|&\le X_{2,h}^{2/3}X_{6,h}^2,&
 \|q_h^{-1}\|&=X_{3,h}^{2/3}X_{6,h},\\
 N_h&=X_{2,h}^{1/3}X_{4,h}^{1/2},&
 N_h^{-1}&=X_{3,h}^{1/3}X_{5,h}^{1/2},\\
 \|\mathsf e_h^{-1}\|
 &\le X_{7,h}X_{3,h}^{1/3}
       (X_{5,h}^{1/2}+X_{6,h}^{1/2}).
\end{align*}
The positive coframe norm is controlled in the same way by square roots of
$q_h$ and by $N_h$.  A finite tensor contraction is therefore bounded by a
finite monomial sum.  Lemma~\ref{lem:V13-multi-Orlicz} controls every power of
every such monomial.

For \eqref{eq:V13-score-monomial-envelope}, use
$(\sum_{\ell=1}^Ly_\ell)^2\le L\sum_\ell y_\ell^2$ and the lemma to obtain
\[
 \mathbb E_h|z_h(u)|^2
 \le\mathfrak C_{\rm ADM}\|u\|^2.
\]
Variance is at most the uncentred second moment.  Cauchy--Schwarz proves the
bilinear estimate.
\end{proof}

\begin{corollary}[Explicit spatial and spacetime negative moments]
\label{cor:V13-volume-moments}
Under \eqref{eq:V13-full-packet}, for every $p>0$,
\begin{align}
 \sup_h\mathbb E_h(\det q_h)^{-p}
 &=\sup_h\mathbb E_h\varrho_h^{-2p}<\infty,
 \label{eq:V13-spatial-negative}\\
 \sup_h\mathbb E_h|\det\mathbf g_h|^{-p}
 &=\sup_h\mathbb E_h
   \varrho_h^{-8p/3}\upsilon_h^{-p}<\infty,
 \label{eq:V13-spacetime-negative}\\
 \sup_h\mathbb E_h\Omega_h^{-p}
 &=\sup_h\mathbb E_h
   \varrho_h^{-4p/3}\upsilon_h^{-p/2}<\infty.
 \label{eq:V13-coframe-negative}
\end{align}
\end{corollary}

\begin{proof}
Use \eqref{eq:V13-lapse-volume,eq:V13-full-determinant} and apply
Lemma~\ref{lem:V13-multi-Orlicz} to the indicated powers of
$X_{3,h}$ and $X_{5,h}$.
\end{proof}

\begin{corollary}[Seven marginal shell drifts suffice]
\label{cor:V13-marginal-drifts}
Independence among the seven variables in
\eqref{eq:V13-packet-variables} is unnecessary.  It is enough that, for every
$j$, the cofinal shell ancestry satisfies the V12 collapse-Lyapunov
conditional drift for
\[
 Y_{j,n}=X_{j,n}^{\alpha_j}
\]
at exponential coefficient $7c_j$, with cutoff-independent drift constants
and finite initial exponential moment.  These seven scalar drifts imply
\eqref{eq:V13-full-packet}.
\end{corollary}

\begin{proof}
The V12 collapse-Lyapunov theorem gives constants $K_j$ such that
\[
 \sup_n\mathbb E e^{7c_jX_{j,n}^{\alpha_j}}\le K_j.
\]
Generalized H\"older with seven factors yields
\[
 \mathbb E\exp\!\left\{\sum_{j=1}^7c_jX_{j,n}^{\alpha_j}\right\}
 \le\prod_{j=1}^7K_j^{1/7}.
\]
\end{proof}

\begin{assessment}[Probabilistic replacement of a fixed tube]
\label{ass:V13-packet-meaning}
The deterministic tube in Theorem~\ref{thm:V13-reference-tube} is sufficient
but stronger than needed.  Definition~\ref{def:V13-full-packet} allows rare
large fields, small speed density, small bracket scale, anisotropic shape, and
large shift, provided their joint tails are exponentially controlled.  It is
therefore the natural finite-cylinder target when the interacting scheduler
must remain unbounded.
\end{assessment}
% -----------------------------------------------------------------------------
\subsection{Non-vacuity and scheduler-native obstruction laws}
\label{subsec:V13-populated-witness}
% -----------------------------------------------------------------------------

\begin{proposition}[A populated full scheduler packet]
\label{prop:V13-populated-packet}
Let
$\mathfrak p=\{A\in\operatorname{Sym}_3:\Tr A=0\}$ and use coordinates
\[
 x,y\in\mathbb R,\qquad A\in\mathfrak p,\qquad\beta\in\mathbb R^3
\]
with
\[
 \varrho=e^x,\qquad
 \upsilon=e^y,\qquad
 \mathcal S=e^A.
\]
Fix positive coefficients and exponents for the six scheduler variables
\[
 \varrho,\quad\varrho^{-1},\quad
 \upsilon,\quad\upsilon^{-1},\quad
 \|\mathcal S\|_{\mathrm{op}},\quad1+\|\beta\|.
\]
If $\Phi$ is their weighted power sum, then
\begin{equation}
 Z^{-1}e^{-2\Phi(x,y,A,\beta)}
 \,\dd x\,\dd y\,\dd A\,\dd\beta
 \label{eq:V13-populated-packet-law}
\end{equation}
is a probability law satisfying the corresponding six-variable version of
\eqref{eq:V13-full-packet}.  If an independent finite Standard Model law has
the required doubled exponential moment of $R_{\rm SM}$, their product
satisfies the seven-variable packet.  Tilting this product by
$e^{-W}$ for any bounded nonconstant mixed interaction $W$ preserves the
packet and gives a nonproduct law.
\end{proposition}

\begin{proof}
The positive and negative exponential powers of $e^x$ and $e^y$ make the
$x$ and $y$ integrals finite at both ends.  The shift term makes the
$\beta$ integral finite.  For $A\in\mathfrak p$,
\[
 \lambda_{\max}(A)\ge\frac{\|A\|_{\mathrm F}}{\sqrt6},
 \qquad
 \|e^A\|_{\mathrm{op}}=e^{\lambda_{\max}(A)}.
\]
Indeed, if $m=\lambda_{\max}(A)$, tracelessness implies that every eigenvalue
has absolute value at most $2m$, and
$\|A\|_{\mathrm F}^2\le6m^2$.  The shape factor in $e^{-2\Phi}$ therefore
decays superexponentially in $\|A\|_{\mathrm F}$, proving $0<Z<\infty$.

Multiplication by $e^\Phi$ leaves the integrable density $e^{-\Phi}/Z$, so
the six-variable packet holds.  The product claim follows from factorization.
If $|W|\le C$, the tilted density and its normalizing constant are bounded
above and below by fixed multiples of those of the product law; all packet
bounds survive.  A nonseparable $W$ which is not almost surely a sum of two
marginal functions gives a nonproduct tilt.
\end{proof}

\begin{assessment}[Scope of the populated law]
\label{ass:V13-populated-scope}
Proposition~\ref{prop:V13-populated-packet} proves that simultaneous full-ADM
noncollapse, unbounded fields, and nonzero mixed coupling are analytically
compatible.  It does not identify the scheduler part of the declared RPESM
action or reference measure.  That identification remains a physical
population step.
\end{assessment}

For a normalized failure certificate, define the scheduler exhaustion
\begin{equation}
 \begin{split}
 \mathcal K_M:=\{&
 R_{\rm SM}\le M,\quad
 M^{-1}\le\varrho\le M,\quad
 M^{-1}\le\upsilon\le M,\\
 &\|\mathcal S\|_{\mathrm{op}}\le M,\quad
 \|\beta\|\le M\},
 \qquad M>1.
 \end{split}
 \label{eq:V13-scheduler-exhaustion}
\end{equation}
Because $\det\mathcal S=1$, the shape condition also gives
$\|\mathcal S^{-1}\|_{\mathrm{op}}\le M^2$.

\begin{theorem}[Pullback of an inverse-writer failure to scheduler escape]
\label{thm:V13-scheduler-witness}
Let $\mu_n$ be scheduler-coupled cylinder laws and let $W_n\ge0$ be measurable
writers such that
\begin{equation}
 0<Z_n^W:=\int W_n\,\dd\mu_n<\infty,
 \qquad Z_n^W\longrightarrow\infty.
 \label{eq:V13-writer-divergence}
\end{equation}
Assume that for every fixed $M>1$,
\begin{equation}
 \sup_n\ \operatorname*{ess\,sup}_{\mathcal K_M}W_n
 \le C_M<\infty.
 \label{eq:V13-local-writer-bound}
\end{equation}
Define
\begin{equation}
 \nu_n^W(\dd\omega):=\frac{W_n(\omega)}{Z_n^W}\,\mu_n(\dd\omega).
 \label{eq:V13-writer-tilt}
\end{equation}
Then
\begin{equation}
 \nu_n^W(\mathcal K_M)\le\frac{C_M}{Z_n^W}\longrightarrow0
 \qquad(M>1).
 \label{eq:V13-witness-escape}
\end{equation}
Thus the normalized writer energy escapes through at least one of:
large Standard Model radius, speed-density zero or infinity, bracket-scale
zero or infinity, shape anisotropy, or shift infinity.

For the determinant writers this becomes more precise:
\begin{enumerate}[label=\textup{(\roman*)},leftmargin=2.8em]
\item if
      $\mathbb E_n(\det q_n)^{-p}
       =\mathbb E_n\varrho_n^{-2p}\to\infty$,
      its normalized tilt concentrates on $\{\varrho_n<\varepsilon\}$ for
      every $\varepsilon>0$;
\item if
      $\mathbb E_n|\det\mathbf g_n|^{-p}\to\infty$,
      its normalized tilt concentrates on
      \[
       \{\varrho_n<\varepsilon\}\cup\{\upsilon_n<\varepsilon\}
       \tag{V13.W1}
      \]
      for every $\varepsilon>0$.
\end{enumerate}
\end{theorem}

\begin{proof}
On $\mathcal K_M$, the numerator of
\eqref{eq:V13-writer-tilt} is at most $C_M$.  Division by $Z_n^W$ proves
\eqref{eq:V13-witness-escape}.

For item~\textup{(i)}, on $\{\varrho\ge\varepsilon\}$ the weight
$\varrho^{-2p}$ is at most $\varepsilon^{-2p}$, so the normalized mass there
tends to zero.  For item~\textup{(ii)}, outside the union in
\textup{(V13.W1)} one has
\[
 |\det\mathbf g|^{-p}
 =\varrho^{-8p/3}\upsilon^{-p}
 \le\varepsilon^{-11p/3}.
\]
The same normalization argument proves the claim.
\end{proof}

\begin{corollary}[The obstruction is expressed in physical scheduler rows]
\label{cor:V13-witness-physical}
For every rational ADM writer covered by
Theorem~\ref{thm:V13-packet-controls-ADM},
condition \eqref{eq:V13-local-writer-bound} holds whenever its fixed tensor
coefficients are cutoff-uniform on $\mathcal K_M$.  Therefore failure of its
uniform first or second moment yields the explicit probability laws
\eqref{eq:V13-writer-tilt}; the failure cannot remain an unnamed
\emph{geometric singularity}.
\end{corollary}

\begin{proof}
All factors in the norm bounds of
Theorems~\ref{thm:V13-exact-factorization}
and \ref{thm:V13-full-inverse} are bounded on $\mathcal K_M$.  A finite
rational tensor contraction is consequently bounded there.  Apply
Theorem~\ref{thm:V13-scheduler-witness}.
\end{proof}
% -----------------------------------------------------------------------------
\subsection{Full Lorentzian noncollapse under weak extraction}
\label{subsec:V13-limit}
% -----------------------------------------------------------------------------

\begin{theorem}[The scheduler packet survives a weak continuum limit]
\label{thm:V13-scheduler-limit}
Let $\mu_n$ be tight laws of the seven variables in
\eqref{eq:V13-packet-variables}, with
$\mathcal S_n\succ0$ and $\det\mathcal S_n=1$, and suppose the full packet
\eqref{eq:V13-full-packet} holds.  Let a subsequence converge weakly in the
ambient finite-coordinate state space, allowing a priori
$\varrho=0$ or $\upsilon=0$.  Then its limit law $\mu$ satisfies
\begin{equation}
 \mu\{\varrho>0,\ \upsilon>0,\ \mathcal S\succ0\}=1.
 \label{eq:V13-limit-interior}
\end{equation}
The derived $q,N,\mathsf e$, and $\mathbf g$ are therefore nondegenerate
almost surely, $\mathbf g$ has Lorentzian inertia $(1,3)$, and the limiting
law satisfies the same packet bound.

Every continuous rational ADM/Standard Model writer with an envelope of the
form \eqref{eq:V13-score-monomial-envelope} is uniformly integrable along the
subsequence.  Its expectation and every finite covariance matrix converge to
those of the limiting writer.
\end{theorem}

\begin{proof}
Extend $e^{c_3\varrho^{-\alpha_3}}$ and
$e^{c_5\upsilon^{-\alpha_5}}$ by $+\infty$ at zero.  These are nonnegative
lower-semicontinuous functions.  The entire packet exponential is likewise
lower semicontinuous.  Portmanteau gives
\[
 \int\exp\!\left\{\sum_{j=1}^7c_jX_j^{\alpha_j}\right\}\dd\mu
 \le\liminf_n
 \int\exp\!\left\{\sum_{j=1}^7c_jX_{j,n}^{\alpha_j}\right\}\dd\mu_n
 \le K.
\]
A finite integral excludes $\varrho=0$ and $\upsilon=0$.
Continuity of determinant gives $\det\mathcal S=1$ in the limit; a
positive-semidefinite finite matrix with determinant one is positive definite.
This proves \eqref{eq:V13-limit-interior}.

The maps in
\eqref{eq:V13-q-shape,eq:V13-lapse-volume,eq:V13-ADM-matrix}
are continuous on this interior.  Their factorization by
$L_\beta^{\mathsf T}\operatorname{diag}(-N^2,q)L_\beta$ gives Lorentzian
inertia.  Lemma~\ref{lem:V13-multi-Orlicz} gives a uniform moment of order
strictly larger than any prescribed writer power.  Hence the writer and its
square are uniformly integrable.  The continuous mapping theorem followed by
uniform integrability proves convergence of first and second moments;
polarization proves covariance convergence.
\end{proof}

\begin{corollary}[Conditional full represented-stress passage]
\label{cor:V13-full-stress-passage}
If the hypotheses of Theorem~\ref{thm:V13-scheduler-limit} are combined with
the V11 line/Green controls and the V12 rational stress envelopes, then every
finite represented stress and geometric-score covariance covered by those
envelopes passes to the weak limit.  No independence between the scheduler
and Standard Model variables is needed.
\end{corollary}

\begin{proof}
All variables occur in the single joint packet
\eqref{eq:V13-full-packet}.  Apply
Theorem~\ref{thm:V13-scheduler-limit} to the finitely many stress and score
writers.
\end{proof}

\begin{firewall}[Moment passage is not the Einstein equation]
\label{fire:V13-limit-scope}
Nondegenerate Lorentzian support and convergent represented-stress moments do
not prove Palatini stationarity, identify a renormalized stress distribution,
or pass the complete Ward and boundary identities.  Those conclusions still
require the vanishing residuals, compact-screen estimates, and action
comparison in the imported global Einstein-or-obstruction theorem.
\end{firewall}

% -----------------------------------------------------------------------------
\subsection{V13 source audit, theorem ledger, and next acquisition}
\label{subsec:V13-ledger}
% -----------------------------------------------------------------------------

\paragraph{What the source audit establishes.}
The attached RPESM definition fixes the reference rates
$k_{\alpha,h}=1/(8h^2)$, hence
$\mathcal B_h=I_3$, and fixes reference
$\varrho_h=1$, $\beta_h=0$.  Its finite normalizability proof states that the
scheduler variables have finite volume at each cutoff.  The later population
theorem places the scheduler-derived ADM packet and stress rows on one event
law.  Neither theorem specifies a cutoff-independent scheduler support tube,
a coercive scheduler potential in the cylinder action, or conditional shell
drifts for the seven variables in \eqref{eq:V13-packet-variables}.  Finite
population and finite volume therefore do not imply
Theorem~\ref{thm:V13-reference-tube} or
Definition~\ref{def:V13-full-packet}.

\paragraph{New conclusions proved in V13.}
Theorem~\ref{thm:V13-exact-factorization} gives the exact scale--shape
decomposition of the imported scheduler formula, including spatial and full
spacetime determinants.  Theorem~\ref{thm:V13-full-inverse} controls the full
inverse metric and coframe.  Proposition~\ref{prop:V13-four-separators}
proves that speed, bracket scale, shape, and shift are independent cofinal
rows.  Theorems~\ref{thm:V13-reference-tube}
and \ref{thm:V13-rate-certificate} give deterministic scheduler- and
rate-level hard-floor certificates, including the necessary $O(h)$ directed
rate cancellation for bounded shift.  Theorem
\ref{thm:V13-packet-controls-ADM} gives the unbounded probabilistic
alternative and an explicit Fisher constant.  Proposition
\ref{prop:V13-populated-packet} proves compatibility with a bounded nonproduct
mixed interaction.  Theorem~\ref{thm:V13-scheduler-witness} pulls a failed
inverse writer back to normalized physical scheduler escape, and
Theorem~\ref{thm:V13-scheduler-limit} passes the successful packet to a
Lorentzian weak limit.

\paragraph{Imported conclusions not reproved.}
The root-rate bracket formula, its analytic ADM reconstruction, the
winner--waiting Fisher identity, the reference scheduler, the flat and
de~Sitter vacuum branches, finite RPESM population, and the global
Einstein-or-obstruction alternative remain imported from the attached
monograph.  The V11 Standard Model radius and line/Green controls and the V12
collapse/stress interface remain imported from the earlier notes.

\paragraph{Next physical acquisition.}
The full interacting continuum is not yet proved.  The next iteration must
freeze the actual scheduler component of the RPESM base measure
$\dd\lambda_n$ and gravitational action $S_{g,n}$, in the coordinates
\[
 (\log\varrho_n,\log\upsilon_n,\log\mathcal S_n,\beta_n).
\]
It must then prove one of:
\begin{enumerate}[label=\textup{(\arabic*)},leftmargin=2.8em]
\item a single support tube satisfying
      \eqref{eq:V13-reference-tube};
\item a uniform potential lower bound and reference well which imply
      \eqref{eq:V13-full-packet}; or
\item the seven conditional shell drifts of
      Corollary~\ref{cor:V13-marginal-drifts}.
\end{enumerate}
If no route passes, the first divergent rational stress or inverse writer must
be tilted as in \eqref{eq:V13-writer-tilt}; its mass then identifies the
specific scheduler boundary to be repaired.  Only after this scheduler card
passes can the programme safely assemble the cofinal Palatini/Einstein
first-variation identity with the controlled Standard Model stress.

% =============================================================================
% END APPENDED VERSION V13 MATERIAL
% =============================================================================


% =============================================================================
% BEGIN APPENDED VERSION V14 MATERIAL
% =============================================================================

\clearpage
\section{Version V14: a projective logarithmic scheduler cylinder}
\label{sec:V14-projective-scheduler}

Version V13 identified the exact missing scheduler rows but did not construct
their joint projective law.  This version supplies a compatible finite-cylinder
completion.  The construction uses all ten ADM scheduler directions, realizes
them by positive $A_3$ root rates at every sufficiently fine cutoff, permits
bounded nonseparable interactions with the Standard Model fields, preserves
exact old marginals, and converges strongly inside one nondegenerate ADM tube.

This is an existential completion of the scheduler component left unspecified
in the RPESM base measure and gravitational action.  It is not asserted to be
the historical or physically selected scheduler law.  Its purpose is to prove
that exact projectivity, full local ADM variability, mixed coupling, and
cofinal coframe noncollapse can hold simultaneously.  After this construction,
the first unclosed positive-branch row is no longer scheduler existence; it is
the vanishing common Palatini--matter first variation required by the imported
Einstein handoff.

% -----------------------------------------------------------------------------
\subsection{The global ten-dimensional logarithmic ADM chart}
\label{subsec:V14-log-chart}
% -----------------------------------------------------------------------------

Let
\begin{equation}
 \mathcal Z
 :=
 \mathbb R_x\oplus\mathbb R_y
 \oplus\operatorname{Sym}_0(3)_A\oplus\mathbb R^3_b,
 \qquad
 \dim\mathcal Z=1+1+5+3=10,
 \label{eq:V14-Z-space}
\end{equation}
where $\operatorname{Sym}_0(3)$ denotes real traceless symmetric matrices.
Equip it with
\begin{equation}
 \|z\|_{\mathcal Z}
 :=\max\{|x|,|y|,\|A\|_{\mathrm{op}},\|b\|\},
 \qquad z=(x,y,A,b).
 \label{eq:V14-Z-norm}
\end{equation}

\begin{theorem}[Global logarithmic scheduler chart]
\label{thm:V14-global-chart}
The map
\begin{equation}
 \begin{aligned}
 \Psi:\mathcal Z&\longrightarrow
 (0,\infty)\times\operatorname{Sym}_3^+\times\mathbb R^3,\\
 (x,y,A,b)&\longmapsto
 \bigl(\varrho=e^x,\ \mathcal B=e^y e^A,\ \beta=b\bigr)
 \end{aligned}
 \label{eq:V14-Psi}
\end{equation}
is a real-analytic diffeomorphism.  In these coordinates,
\begin{align}
 \upsilon&=e^y,&
 \mathcal S&=e^A,
 \label{eq:V14-log-shape}\\
 q&=e^{2x/3}e^{-A},&
 N&=e^{x/3+y/2},
 \label{eq:V14-log-ADM}\\
 \det q&=e^{2x},&
 |\det\mathbf g|&=e^{8x/3+y}.
 \label{eq:V14-log-determinants}
\end{align}
The reference scheduler
$(\mathcal B,\varrho,\beta)=(I_3,1,0)$ is $z=0$.
\end{theorem}

\begin{proof}
The forward map is analytic and has positive scalar and matrix coordinates.
Given $(\varrho,\mathcal B,\beta)$, its inverse is
\begin{equation}
 x=\log\varrho,\qquad
 y=\frac13\log\det\mathcal B,\qquad
 A=\log\mathcal B-yI_3,\qquad b=\beta.
 \label{eq:V14-Psi-inverse}
\end{equation}
The positive matrix logarithm is real analytic on
$\operatorname{Sym}_3^+$.  Moreover
$\Tr A=\Tr\log\mathcal B-3y=\log\det\mathcal B-\log\det\mathcal B=0$,
and exponentiation recovers $\mathcal B$.  This proves bijectivity and
analyticity of the inverse.  Equations
\eqref{eq:V14-log-shape,eq:V14-log-ADM,eq:V14-log-determinants} are the V13
scale--shape identities after substitution.
\end{proof}

\begin{theorem}[Uniform hard floor on a logarithmic ball]
\label{thm:V14-log-ball-floor}
If $\|z\|_{\mathcal Z}\le R$, then
\begin{align}
 e^{-R}\le\varrho,\upsilon&\le e^R,
 &
 e^{-R}I_3\preceq\mathcal S&\preceq e^RI_3,
 \label{eq:V14-log-basic-bounds}\\
 e^{-5R/3}I_3\preceq q&\preceq e^{5R/3}I_3,
 &
 e^{-5R/6}\le N&\le e^{5R/6},
 \label{eq:V14-log-qN-bounds}\\
 e^{-11R/3}\le|\det\mathbf g|&\le e^{11R/3},
 &
 \|\mathbf g^{-1}\|_{\mathrm{op}}
 &\le(1+R)^2e^{5R/3},
 \label{eq:V14-log-full-bounds}\\
 \|\mathsf e^{-1}\|_{\mathrm{op}}
 &\le(1+R)e^{5R/6}.
 \label{eq:V14-log-coframe-bound}
\end{align}
Every derivative of the maps
$z\mapsto(q,N,\mathsf e,\mathbf g)$ and their inverses is bounded on the
closed ball $\{\|z\|_{\mathcal Z}\le R\}$.
\end{theorem}

\begin{proof}
The spectral theorem gives
$e^{-R}I_3\preceq e^A\preceq e^RI_3$.  The first line follows, and
$q=e^{2x/3}e^{-A}$ gives the metric bounds.  The lapse and determinant
bounds follow from \eqref{eq:V14-log-ADM,eq:V14-log-determinants}.
The V13 full inverse estimates give
\[
 \|\mathbf g^{-1}\|
 \le(1+R)^2e^{2R/3}\max\{e^R,e^R\}
 =(1+R)^2e^{5R/3}
\]
and similarly
$\|\mathsf e^{-1}\|\le(1+R)e^{R/3}e^{R/2}$.
Analyticity and compactness of the finite-dimensional closed ball give the
derivative bounds.
\end{proof}

\begin{corollary}[All rational ADM writers are bounded on the ball]
\label{cor:V14-ball-writers}
For fixed tensor degree and fixed coefficient norms, every polynomial
contraction in the ADM variables and their inverses has a finite upper bound
depending only on $R$ and those degrees and coefficients.
\end{corollary}

\begin{proof}
Apply Theorem~\ref{thm:V14-log-ball-floor} to each factor and use
submultiplicativity of tensor and operator norms.
\end{proof}
% -----------------------------------------------------------------------------
\subsection{Positive root-rate realization of the logarithmic chart}
\label{subsec:V14-rate-realization}
% -----------------------------------------------------------------------------

Use the six unoriented $A_3$ root classes $r_a$ of the imported scheduler
theorem and define
\begin{equation}
 \mathcal L:\mathbb R^6\to\operatorname{Sym}_3,\qquad
 \mathcal L(u)=\sum_{a=1}^6u_ar_ar_a^{\mathsf T},
 \label{eq:V14-L-map}
\end{equation}
and
\begin{equation}
 \mathcal D:\mathbb R^6\to\mathbb R^3,\qquad
 \mathcal D(d)=\sum_{a=1}^6d_ar_a.
 \label{eq:V14-D-map}
\end{equation}
The imported root calculation says that $\mathcal L$ is an isomorphism and
$\mathcal D$ is onto.  Fix a right inverse
$\mathcal J:\mathbb R^3\to\mathbb R^6$ of $\mathcal D$, and put
\begin{equation}
 C_{\mathcal L}:=\|\mathcal L^{-1}\|_{\mathrm{op}\to\ell^\infty},
 \qquad
 C_{\mathcal J}:=\|\mathcal J\|_{\ell^2\to\ell^\infty}.
 \label{eq:V14-LJ-constants}
\end{equation}

\begin{theorem}[Uniformly positive directed-rate realization]
\label{thm:V14-positive-rates}
Let
\begin{equation}
 \delta_{\mathcal L}:=\frac1{8C_{\mathcal L}}.
 \label{eq:V14-delta-L}
\end{equation}
Suppose
\begin{equation}
 \|\mathcal B-I_3\|_{\mathrm{op}}\le\delta_{\mathcal L},
 \qquad
 hC_{\mathcal J}\|\beta\|\le\frac1{16}.
 \label{eq:V14-rate-domain}
\end{equation}
Define
\begin{align}
 u&:=\mathcal L^{-1}(\mathcal B),
 &
 s_{a,h}&:=h^{-2}u_a,
 \label{eq:V14-rate-sums}\\
 d_h&:=h^{-1}\mathcal J\beta,
 &
 k_{a,h}^{\pm}&:=\frac12(s_{a,h}\pm d_{a,h}).
 \label{eq:V14-directed-rates}
\end{align}
Then every directed rate is strictly positive and
\begin{equation}
 \frac1{32h^2}\le k_{a,h}^{\pm}\le\frac7{32h^2}.
 \label{eq:V14-rate-bounds}
\end{equation}
They reconstruct the prescribed bracket and shift exactly:
\begin{equation}
 h^2\sum_a(k_{a,h}^++k_{a,h}^-)r_ar_a^{\mathsf T}
 =\mathcal B,
 \qquad
 h\sum_a(k_{a,h}^+-k_{a,h}^-)r_a=\beta.
 \label{eq:V14-rate-reconstruction}
\end{equation}
Furthermore, with $K_h=\sum_a(k_{a,h}^++k_{a,h}^-)$,
\begin{equation}
 \frac34\le h^2K_h\le\frac94,
 \qquad
 \frac{k_{a,h}^{\pm}}{K_h}\ge\frac1{72}.
 \label{eq:V14-rate-scale-share}
\end{equation}
Thus the same construction supplies the V13 rescaled-rate, rate-share, shape,
and bounded-shift certificates.
\end{theorem}

\begin{proof}
At the reference scheduler,
\[
 \mathcal L(1/4,\ldots,1/4)=I_3.
\]
The first inequality in \eqref{eq:V14-rate-domain} gives
\[
 \|u-(1/4,\ldots,1/4)\|_{\ell^\infty}
 \le C_{\mathcal L}\|\mathcal B-I_3\|_{\mathrm{op}}\le\frac18.
\]
Hence $1/8\le u_a\le3/8$.  The shift condition gives
\[
 |d_{a,h}|
 \le h^{-1}C_{\mathcal J}\|\beta\|
 \le\frac1{16h^2}.
\]
Substitution in \eqref{eq:V14-directed-rates} proves
\eqref{eq:V14-rate-bounds}.

The first reconstruction identity is
$\mathcal L(u)=\mathcal B$.  The second follows from
$\mathcal D\mathcal J=I_3$.  Summing the six rate sums gives
$3/(4h^2)\le K_h\le9/(4h^2)$.  Dividing the lower individual-rate bound by
the upper total-rate bound gives $1/72$.
\end{proof}

\begin{corollary}[One fixed logarithmic ball has positive rates at all fine cutoffs]
\label{cor:V14-log-ball-rates}
Choose $R_0>0$ so that
\begin{equation}
 e^{2R_0}-1\le\delta_{\mathcal L}.
 \label{eq:V14-R0-choice}
\end{equation}
For every $\|z\|_{\mathcal Z}\le R_0$,
\[
 \|\mathcal B-I_3\|_{\mathrm{op}}\le e^{2R_0}-1.
\]
Consequently the positive rates
\eqref{eq:V14-directed-rates} realize the entire logarithmic ball whenever
\begin{equation}
 h\le h_0:=\frac1{16C_{\mathcal J}R_0}.
 \label{eq:V14-h0}
\end{equation}
The six bracket coordinates, three selected difference coordinates, and the
speed coordinate give all ten ADM tangent directions.  The unused
three-dimensional difference kernel remains the circulation sector.
\end{corollary}

\begin{proof}
The eigenvalues of $e^ye^A$ lie in $[e^{-2R_0},e^{2R_0}]$, so its
operator-norm distance from the identity is at most $e^{2R_0}-1$.
Also $\|\beta\|\le R_0$, and \eqref{eq:V14-h0} implies the second condition
in \eqref{eq:V14-rate-domain}.  The dimension statement follows from
invertibility of $\mathcal L$, surjectivity of $\mathcal D$, and the
independent speed coordinate.
\end{proof}

\begin{corollary}[Scheduler reflection covariance]
\label{cor:V14-rate-reflection}
Under the involution
\begin{equation}
 \Theta_{\mathcal Z}(x,y,A,b)=(x,y,A,-b),
 \label{eq:V14-Z-reflection}
\end{equation}
the rate sums stay fixed and
$k_{a,h}^+\leftrightarrow k_{a,h}^-$.  Hence a
$\Theta_{\mathcal Z}$-covariant law in the logarithmic chart induces the
required directed-root reflection covariance.
\end{corollary}

\begin{proof}
Only $d_h=h^{-1}\mathcal Jb$ changes sign.  The definition
\eqref{eq:V14-directed-rates} then exchanges the two directed rates.
\end{proof}
% -----------------------------------------------------------------------------
\subsection{Exact projective cylinder with bounded mixed shell interactions}
\label{subsec:V14-projective-cylinder}
% -----------------------------------------------------------------------------

Let
\[
 \mathbb B_{\mathcal Z}:=\{\xi\in\mathcal Z:\|\xi\|_{\mathcal Z}\le1\}.
\]
Choose a level-zero scheduler $z_0\in\mathcal Z$ and $a_j\ge0$ with
\begin{equation}
 A_*:=\sum_{j=1}^{\infty}a_j,
 \qquad \|z_0\|_{\mathcal Z}+A_*\le R_0,
 \label{eq:V14-summable-amplitudes}
\end{equation}
and define
\begin{equation}
 z_n:=z_0+\sum_{j=1}^na_j\xi_j,
 \qquad \xi_j\in\mathbb B_{\mathcal Z}.
 \label{eq:V14-scheduler-series}
\end{equation}
Let $\omega_n$ denote the complete old event record through level $n$,
including its Standard Model fields.  For every $j$, let $\kappa_j$ be a
probability law on $\mathbb B_{\mathcal Z}$ and let
$V_j(\omega_{j-1},\xi)$ be a measurable shell interaction satisfying
\begin{equation}
 |V_j(\omega_{j-1},\xi)|\le M_j<\infty.
 \label{eq:V14-bounded-shell-interaction}
\end{equation}
Set
\begin{align}
 Z_j(\omega_{j-1})
 &:=
 \int_{\mathbb B_{\mathcal Z}}
 e^{-V_j(\omega_{j-1},\xi)}\,\kappa_j(\dd\xi),
 \label{eq:V14-shell-Z}\\
 K_j(\dd\xi\mid\omega_{j-1})
 &:=
 Z_j(\omega_{j-1})^{-1}
 e^{-V_j(\omega_{j-1},\xi)}\,\kappa_j(\dd\xi).
 \label{eq:V14-shell-kernel}
\end{align}

\begin{theorem}[Nontrivial projective scheduler cylinder]
\label{thm:V14-projective-cylinder}
Starting from any finite old-cylinder probability law $\mu_0$, define
recursively
\begin{equation}
 \mu_j(\dd\omega_{j-1},\dd\xi_j)
 =
 \mu_{j-1}(\dd\omega_{j-1})
 K_j(\dd\xi_j\mid\omega_{j-1}).
 \label{eq:V14-projective-law}
\end{equation}
Then:
\begin{enumerate}[label=\textup{(\roman*)},leftmargin=2.8em]
\item dropping the last shell gives the exact old marginal
      \[
       (\pi_{j/j-1})_*\mu_j=\mu_{j-1};
       \tag{V14.P1}
      \]
\item every scheduler value satisfies
      $\|z_j\|_{\mathcal Z}\le A_*\le R_0$, hence the hard floors of
      Theorem~\ref{thm:V14-log-ball-floor} uniformly;
\item at every mesh $h_j\le h_0$, the positive directed rates of
      Theorem~\ref{thm:V14-positive-rates} realize the scheduler exactly and
      obey \eqref{eq:V14-rate-bounds,eq:V14-rate-scale-share};
\item the conditional density is uniformly equivalent to its reference:
      \begin{equation}
       e^{-2M_j}
       \le\frac{\dd K_j(\cdot\mid\omega_{j-1})}{\dd\kappa_j}
       \le e^{2M_j};
       \label{eq:V14-density-equivalence}
      \end{equation}
\item if $\kappa_j$ has full support in a neighbourhood of zero and
      $a_j>0$, the new scheduler detail has all ten local ADM directions;
\item if the measure-valued map
      $\omega\mapsto K_j(\,\cdot\mid\omega)$ is not almost surely constant,
      then the new shell is not a product of the old event and its scheduler
      detail.
\end{enumerate}
If $\kappa_j$ is invariant under \eqref{eq:V14-Z-reflection} and
\begin{equation}
 V_j(\Theta\omega,\Theta_{\mathcal Z}\xi)=V_j(\omega,\xi),
 \label{eq:V14-interaction-reflection}
\end{equation}
then the projective cylinder is reflection covariant.
\end{theorem}

\begin{proof}
Equation \eqref{eq:V14-shell-kernel} is normalized, so integration over
$\xi_j$ proves exact marginality.  The triangle inequality and
\eqref{eq:V14-summable-amplitudes} give
\[
 \|z_j\|_{\mathcal Z}
 \le\|z_0\|_{\mathcal Z}
    +\sum_{\ell=1}^ja_\ell\|\xi_\ell\|_{\mathcal Z}
 \le\|z_0\|_{\mathcal Z}+A_*\le R_0.
\]
The hard-floor and rate conclusions follow from
Theorems~\ref{thm:V14-log-ball-floor}
and \ref{thm:V14-positive-rates}.

The potential bound gives
$e^{-M_j}\le Z_j(\omega)\le e^{M_j}$ and therefore
\eqref{eq:V14-density-equivalence}.  Equivalence of measures preserves
support.  Since $\mathcal Z$ has dimension ten, a reference density positive
near zero supplies all ten local directions.  A joint law is a product exactly when its conditional scheduler law is
almost surely constant, proving item~\textup{(vi)}.  Finally,
\eqref{eq:V14-interaction-reflection} and invariance of $\kappa_j$ imply
$Z_j(\Theta\omega)=Z_j(\omega)$ and covariance of $K_j$.
\end{proof}

\begin{corollary}[Explicit bounded nonseparable coupling]
\label{cor:V14-explicit-mixed}
Let $F_j(\omega_{j-1})$ be a bounded, not almost surely constant Standard
Model writer and let $\ell_j$ be bounded and not
$\kappa_j$-almost surely constant on $\mathbb B_{\mathcal Z}$.
For $\epsilon_j\ne0$ set
\begin{equation}
 V_j(\omega_{j-1},\xi)
 =\epsilon_jF_j(\omega_{j-1})\ell_j(\xi).
 \label{eq:V14-explicit-interaction}
\end{equation}
Then \eqref{eq:V14-bounded-shell-interaction} holds and the conditional
scheduler law varies with the Standard Model event.  Choosing $F_j$ and
$\ell_j$ with matching reflection parity gives
\eqref{eq:V14-interaction-reflection}.
\end{corollary}

\begin{proof}
Boundedness is immediate.  Write $K_t$ for the exponential tilt with
parameter $t=\epsilon_jF_j(\omega)$.  If $K_{t_1}=K_{t_2}$ for
$t_1\ne t_2$, equality of their densities would make
$e^{-(t_1-t_2)\ell_j}$ constant $\kappa_j$-almost surely, contradicting the
hypothesis on $\ell_j$.  Thus $t\mapsto K_t$ is injective, and the
nonconstant $F_j$ makes the conditional law depend on the old event.  The
parity statement follows by multiplication.
\end{proof}

\begin{assessment}[An explicit choice of base measure and scheduler action]
\label{ass:V14-base-action}
Relative to the product reference
$\prod_j\kappa_j$, equations
\eqref{eq:V14-shell-Z,eq:V14-shell-kernel} specify the scheduler base measure
and the normalized shell action
\[
 V_j(\omega_{j-1},\xi_j)+\log Z_j(\omega_{j-1}).
\]
This is exactly the conditional normalization mechanism of the imported
RPESM projective transport theorem.  It fills the previously unspecified
scheduler-measure row while retaining the old marginal.  It does not identify
this bounded shell action with the physical Palatini action.
\end{assessment}
% -----------------------------------------------------------------------------
\subsection{Pathwise projective limit and strong ADM convergence}
\label{subsec:V14-strong-limit}
% -----------------------------------------------------------------------------

Put
\begin{equation}
 r_n:=\sum_{j>n}a_j\longrightarrow0.
 \label{eq:V14-tail-radius}
\end{equation}

\begin{theorem}[Pathwise scheduler limit without subsequence extraction]
\label{thm:V14-pathwise-limit}
The consistent cylinder laws of
Theorem~\ref{thm:V14-projective-cylinder} have an infinite projective law
$\mu_\infty$.  On every path,
\begin{equation}
 z_\infty:=z_0+\sum_{j=1}^{\infty}a_j\xi_j
 \label{eq:V14-z-infinity}
\end{equation}
exists in $\mathcal Z$ and
\begin{equation}
 \|z_\infty-z_n\|_{\mathcal Z}\le r_n.
 \label{eq:V14-z-tail}
\end{equation}
For every ADM variable or inverse variable
\[
 F\in\{\varrho,\upsilon,\mathcal S,\mathcal B,q,N,
       \mathsf e,\mathbf g,
       q^{-1},N^{-1},\mathsf e^{-1},\mathbf g^{-1}\},
\]
there is $L_{F,R_0}<\infty$ such that
\begin{equation}
 \|F(z_\infty)-F(z_n)\|
 \le L_{F,R_0}r_n
 \label{eq:V14-F-tail}
\end{equation}
pathwise.  The limiting coframe is uniformly nondegenerate and every fixed
rational ADM writer converges in every $L^p(\mu_\infty)$, $p<\infty$.
\end{theorem}

\begin{proof}
All spaces are standard Borel and
\eqref{eq:V14-shell-kernel} is a probability kernel.  The
Ionescu--Tulcea extension theorem therefore gives one infinite law with the
finite marginals \eqref{eq:V14-projective-law}.  Absolute summability gives
\eqref{eq:V14-z-infinity,eq:V14-z-tail} on every path.

Every displayed map $F$ has bounded derivative on the logarithmic ball by
Theorem~\ref{thm:V14-log-ball-floor}.  The finite-dimensional mean-value
inequality proves \eqref{eq:V14-F-tail}.  The hard floor holds at the limit
because the closed ball is preserved.  A fixed rational writer is bounded on
that ball by Corollary~\ref{cor:V14-ball-writers}; bounded pointwise
convergence then gives convergence in every finite $L^p$ by dominated
convergence.
\end{proof}

\begin{corollary}[Normalized rates converge while finite shift survives]
\label{cor:V14-normalized-rate-limit}
Let $h_n\downarrow0$ and use the rates of
\eqref{eq:V14-directed-rates}.  Then
\begin{equation}
 h_n^2k_{a,n}^{\pm}
 =
 \frac12u_a(\mathcal B_n)
 \pm\frac{h_n}{2}(\mathcal J\beta_n)_a
 \longrightarrow
 \frac12u_a(\mathcal B_\infty).
 \label{eq:V14-normalized-rate-limit}
\end{equation}
The plus/minus normalized rates have the same limit, while
\[
 h_n\sum_a(k_{a,n}^+-k_{a,n}^-)r_a=\beta_n\longrightarrow\beta_\infty.
\]
Thus a finite continuum shift is encoded in the first antisymmetric
$h^{-1}$ correction to the symmetric $h^{-2}$ scheduler, rather than in its
leading normalized rate profile.
\end{corollary}

\begin{proof}
Multiply \eqref{eq:V14-directed-rates} by $h_n^2$ and use
Theorem~\ref{thm:V14-pathwise-limit}.  The shift identity is
\eqref{eq:V14-rate-reconstruction}.
\end{proof}

The construction also yields the strong local convergence required for
quadratic continuum insertions when the details have spatial regularity.

\begin{theorem}[Summable regular details give strong local ADM convergence]
\label{thm:V14-strong-local}
Let $K$ be a compact regulator patch and $0<\alpha\le1$.  Suppose now that
\[
 \xi_j\in C^{0,\alpha}(K;\mathcal Z),
 \qquad
 \|\xi_j\|_{C^{0,\alpha}}\le1,
\]
and retain \eqref{eq:V14-summable-amplitudes}.  Then
$z_n\to z_\infty$ in $C^{0,\alpha}(K;\mathcal Z)$ with
\begin{equation}
 \|z_\infty-z_n\|_{C^{0,\alpha}}\le r_n.
 \label{eq:V14-Calpha-tail}
\end{equation}
Every ADM variable and inverse listed in
Theorem~\ref{thm:V14-pathwise-limit} converges strongly in
$C^{0,\alpha}(K)$.

Let $v_n$ be a sequence of scheduler time, command, or weak-variation
velocities such that
\begin{equation}
 v_n\longrightarrow v
 \quad\hbox{strongly in }L^2(K;\mathcal Z).
 \label{eq:V14-velocity-convergence}
\end{equation}
For every listed analytic writer $F$,
\begin{equation}
 DF(z_n)v_n\longrightarrow DF(z_\infty)v
 \quad\hbox{strongly in }L^2(K).
 \label{eq:V14-derived-velocity}
\end{equation}
Consequently finite quadratic products of these derived velocities converge
in $L^1(K)$.
\end{theorem}

\begin{proof}
The Banach-space triangle inequality gives
\eqref{eq:V14-Calpha-tail}.  On the common bounded range, the first two
derivatives of every analytic finite-dimensional map $F$ are bounded.  The
standard composition estimate in the H\"older Banach algebra gives
\[
 \|F(z_n)-F(z_\infty)\|_{C^{0,\alpha}}
 \le C_{F,R_0,K}\|z_n-z_\infty\|_{C^{0,\alpha}}.
\]
This proves the first claim, including inverse writers because the hard floor
keeps the range inside their analytic domain.

For the differentiated writers, write
\[
 DF(z_n)v_n-DF(z_\infty)v
 =DF(z_n)(v_n-v)+[DF(z_n)-DF(z_\infty)]v.
\]
The first term tends to zero in $L^2$ by the uniform derivative bound.  The
second is bounded by
$C\|z_n-z_\infty\|_{L^\infty}\|v\|_{L^2}$ and also tends to zero.
Finally, strong $L^2$ convergence of each factor implies strong $L^1$
convergence of pairwise products by Cauchy--Schwarz.
\end{proof}

\begin{assessment}[Which Einstein-handoff row is advanced]
\label{ass:V14-handoff-row}
Theorem~\ref{thm:V14-strong-local} supplies a direct compact-screen route for
the scheduler-derived metric/coframe variables and their quadratic insertions,
provided the physical scheduler details and their variation velocities obey
the displayed summability.  It does not control curvature or prove that the
common Palatini--matter first-variation residual vanishes.  Those are the next
action-level rows of the imported Einstein handoff.
\end{assessment}
% -----------------------------------------------------------------------------
\subsection{Action-selected scheduler details}
\label{subsec:V14-action-selected}
% -----------------------------------------------------------------------------

The preceding scheduler series allows arbitrary bounded details.  The next
result gives a sufficient condition under which the details are instead forced
by the cutoff common-action equation.  It uses invertibility after the
physical gauge and relation quotient; positivity of the Einstein Hessian is
not assumed.

\begin{theorem}[Uniform inverse-function transport of cutoff residuals]
\label{thm:V14-residual-transport}
Give $\mathcal Z$ a Euclidean norm, let
$B_R=\{z:\|z\|_{\mathcal Z}\le R\}$, and let
\[
 F_n:B_R\longrightarrow\mathcal Z
\]
be continuously differentiable transported first-variation residuals.
Let $H:\mathcal Z\to\mathcal Z$ be invertible.  Suppose that for some
$0\le\theta<1$,
\begin{align}
 \sup_n\sup_{z\in B_R}
 \|I-H^{-1}DF_n(z)\|_{\mathrm{op}}
 &\le\theta,
 \label{eq:V14-uniform-linearization}\\
 \sup_n\|H^{-1}F_n(0)\|
 &\le(1-\theta)R.
 \label{eq:V14-center-residual}
\end{align}
Then every $F_n$ has a unique zero $z_n\in B_R$.

Define the adjacent transported action defect
\begin{equation}
 \epsilon_n:=
 \sup_{z\in B_R}
 \|H^{-1}(F_{n+1}(z)-F_n(z))\|.
 \label{eq:V14-adjacent-action-defect}
\end{equation}
The selected zeros obey
\begin{equation}
 \|z_{n+1}-z_n\|
 \le\frac{\epsilon_n}{1-\theta}.
 \label{eq:V14-root-transport}
\end{equation}
If $\sum_n\epsilon_n<\infty$, then $z_n$ converges in $\mathcal Z$ and
\begin{equation}
 \|z_\infty-z_n\|
 \le\frac1{1-\theta}\sum_{j\ge n}\epsilon_j.
 \label{eq:V14-root-tail}
\end{equation}
Moreover $H^{-1}F_n$ converges uniformly on $B_R$ to a continuous map
$H^{-1}F_\infty$, and
\begin{equation}
 F_\infty(z_\infty)=0.
 \label{eq:V14-limit-residual-zero}
\end{equation}
\end{theorem}

\begin{proof}
Put $T_n(z)=z-H^{-1}F_n(z)$.  The derivative bound gives
\[
 \|T_n(z)-T_n(w)\|\le\theta\|z-w\|.
\]
Also
\[
 \|T_n(z)\|
 \le\theta\|z\|+\|H^{-1}F_n(0)\|
 \le\theta R+(1-\theta)R=R.
\]
Thus $T_n$ is a contraction of the complete ball $B_R$ into itself.
Banach's fixed-point theorem gives a unique $z_n\in B_R$ with
$F_n(z_n)=0$.

For $z,w\in B_R$, the line segment lies in $B_R$ and
\[
 H^{-1}(F_n(z)-F_n(w))
 =
 \int_0^1H^{-1}DF_n(w+t(z-w))(z-w)\,\dd t.
\]
The operator in the integral differs from the identity by norm at most
$\theta$, so
\begin{equation}
 \|H^{-1}(F_n(z)-F_n(w))\|
 \ge(1-\theta)\|z-w\|.
 \label{eq:V14-lower-inverse-bound}
\end{equation}
Take $z=z_{n+1}$ and $w=z_n$.  Since both are their respective zeros,
\[
 H^{-1}F_n(z_{n+1})
 =
 H^{-1}(F_n(z_{n+1})-F_{n+1}(z_{n+1})),
\]
whose norm is at most $\epsilon_n$.  Equation
\eqref{eq:V14-lower-inverse-bound} proves
\eqref{eq:V14-root-transport}.  Summability gives the Cauchy property and
\eqref{eq:V14-root-tail}.  The same telescoping bound makes
$H^{-1}F_n$ uniformly Cauchy.  Passing to the limit in
$F_n(z_n)=0$ proves \eqref{eq:V14-limit-residual-zero}.
\end{proof}

\begin{corollary}[Common-action population of the summable scheduler series]
\label{cor:V14-action-populates-series}
Suppose $F_n(\omega,z)$ is jointly measurable in $\omega,z$, continuous in
$z$, and is the Riesz representative, on one transported
finite determining scheduler screen, of the complete gravitational plus
Standard Model first variation at cutoff $n$.  Assume
\eqref{eq:V14-uniform-linearization,eq:V14-center-residual} hold uniformly in
the event $\omega$, and that the deterministic bounds
$\epsilon_n$ in \eqref{eq:V14-adjacent-action-defect} are summable.  Then the
unique roots $z_n(\omega)$ are measurable, solve the total scheduler
first-variation equation at every cutoff, and converge eventwise to a root of
the limiting residual.

If
\begin{equation}
 \|z_0\|+\frac1{1-\theta}\sum_{n\ge0}\epsilon_n\le R_0,
 \label{eq:V14-action-tube-budget}
\end{equation}
their successive differences can be written
$z_{n+1}-z_n=a_{n+1}\xi_{n+1}$ with
$\|\xi_{n+1}\|\le1$ and $\sum a_n\le R_0$.
Hence they populate the hard-floor scheduler series of
Theorem~\ref{thm:V14-projective-cylinder}.
\end{corollary}

\begin{proof}
Contraction iteration starting at zero is measurable at every step, and its
pointwise limit is the measurable fixed point $z_n(\omega)$.
Theorem~\ref{thm:V14-residual-transport} gives all other assertions.  Take
$a_{n+1}=\epsilon_n/(1-\theta)$ when this is positive and normalize the
difference; take $\xi_{n+1}=0$ when $a_{n+1}=0$.
\end{proof}

\begin{firewall}[The invertible screen must be physically obtained]
\label{fire:V14-Hessian-screen}
The map $H$ cannot be declared invertible before diffeomorphism, local Lorentz,
gauge, relation, and boundary-null directions have been removed on the
physical determining quotient.  Failure of
\eqref{eq:V14-uniform-linearization} is itself a supported Hessian singular
direction or nonlinear radius obstruction.  The theorem neither deletes the
DeWitt conformal direction nor turns an indefinite Einstein Hessian positive;
it only requires an inverse on the final transported finite screen.
\end{firewall}
\begin{proposition}[Finite residual pass-or-witness card]
\label{prop:V14-residual-witness}
On a fixed finite scheduler screen and compact ball $B_R$, every hypothesis of
Theorem~\ref{thm:V14-residual-transport} has a finite obstruction witness:
\begin{enumerate}[label=\textup{(\roman*)},leftmargin=2.8em]
\item failure of invertibility of $H$ gives a unit vector in $\ker H$;
\item failure of \eqref{eq:V14-uniform-linearization} gives
      $z\in B_R$ and a unit $v$ with
      \[
       \|(I-H^{-1}DF_n(z))v\|>\theta;
       \tag{V14.R1}
      \]
\item failure of \eqref{eq:V14-center-residual} gives the explicit center
      residual vector $H^{-1}F_n(0)$;
\item for a proposed adjacent tolerance $\overline\epsilon_n$, failure of
      $\epsilon_n\le\overline\epsilon_n$ gives $z\in B_R$ with
      \[
       \|H^{-1}(F_{n+1}(z)-F_n(z))\|>\overline\epsilon_n.
       \tag{V14.R2}
      \]
\end{enumerate}
For continuous $DF_n$ these witnesses may be chosen at maximizers.  Thus the
action-selected scheduler construction either passes on the declared screen
or returns a concrete Hessian, radius, or cutoff-transport obstruction.
\end{proposition}

\begin{proof}
Singular-value decomposition gives item~\textup{(i)}.  The closed ball and
unit sphere are compact in the finite screen, and the displayed norm functions
are continuous.  Their suprema are attained, proving
items~\textup{(ii)} and \textup{(iv)}.  Item~\textup{(iii)} is already a
finite vector.
\end{proof}

% -----------------------------------------------------------------------------
\subsection{V14 theorem ledger and next action-level acquisition}
\label{subsec:V14-ledger}
% -----------------------------------------------------------------------------

\paragraph{New conclusions proved in V14.}
Theorem~\ref{thm:V14-global-chart} gives a global analytic ten-dimensional
logarithmic chart for the scheduler ADM variables.  Theorem
\ref{thm:V14-log-ball-floor} turns a single summable-coordinate budget into
uniform spatial, lapse, full-volume, inverse-metric, and inverse-coframe
bounds.  Theorem~\ref{thm:V14-positive-rates} lifts one fixed chart ball to
strictly positive $A_3$ directed rates with uniform $h^{-2}$ bounds, rate
shares, and exact finite shift.  Theorem
\ref{thm:V14-projective-cylinder} supplies normalized, reflection-covariant,
nonproduct shell kernels while preserving every old marginal.
Theorem~\ref{thm:V14-pathwise-limit} constructs the infinite projective law
and a pathwise nondegenerate ADM limit without subsequence extraction.
Theorem~\ref{thm:V14-strong-local} upgrades summable regular details and
velocities to strong local ADM and quadratic-insertion convergence.
Finally, Theorem~\ref{thm:V14-residual-transport} shows how a uniformly
invertible transported common-action residual selects the scheduler details
and transports exact cutoff stationarity to the limit.

\paragraph{What is imported rather than repeated.}
The $A_3$ bracket and shift maps, their rank statements, the scheduler-to-ADM
formula, the abstract normalized projective shell kernel, the reference and
vacuum schedulers, and the conditional Einstein handoff are imported from the
grand tensor manuscript.  V14 specializes those interfaces by supplying a
single compatible scheduler measure, a positive-rate lift, quantitative
noncollapse, and an action-root transport theorem.

\paragraph{Exact remaining gap.}
The V14 cylinder is an explicit admissible completion of the scheduler row,
but its bounded mixed potential has not been identified with the physical
gravitational plus Standard Model common action.  The full continuum is
therefore not claimed.  On a selected RPESM determining screen one must still
construct the actual transported residual $F_n$ and verify:
\begin{enumerate}[label=\textup{(\arabic*)},leftmargin=2.8em]
\item removal of all physical null directions before choosing $H$;
\item the uniform inverse-radius bound
      \eqref{eq:V14-uniform-linearization};
\item the tube budget \eqref{eq:V14-action-tube-budget};
\item summability of the adjacent total-action defects
      \eqref{eq:V14-adjacent-action-defect};
\item the connection/curvature, Ward, counterterm, and boundary rows of the
      imported Einstein handoff.
\end{enumerate}
The next iteration should attack item~\textup{(1)} by forming the supported
gauge/relation short of the finite common-action Hessian.  A positive supported
singular-value floor gives the inverse $H^{-1}$ used here; failure gives its
source-minimal kernel vector.  Only after that finite Hessian card is
physically populated should the Palatini--matter residual be promoted to a
continuum Einstein identity.

\begin{firewall}[V14 closes scheduler existence, not physical selection]
\label{fire:V14-final-scope}
The explicit logarithmic cylinder proves compatibility of projectivity,
ten-dimensional ADM variation, positive root rates, mixed coupling,
noncollapse, and strong limit passage.  It does not prove that the selected
microscopic RPESM action chooses this cylinder, that its limiting residual is
the Einstein--Standard Model residual, or that the remaining Ward and boundary
certificates pass.
\end{firewall}

% =============================================================================
% END APPENDED VERSION V14 MATERIAL
% =============================================================================


% =============================================================================
% BEGIN APPENDED VERSION V15 MATERIAL
% =============================================================================

\clearpage
\section{Version V15: the indefinite gauge--relation Hessian screen}
\label{sec:V15-Hessian-screen}

Version V14 constructed a nondegenerate projective scheduler and showed that a
uniformly invertible transported common-action derivative selects a unique
summable scheduler response.  The next task is to obtain that inverse without
mistaking Fisher positivity for invertibility of the physical Einstein
Hessian.

The attached monograph contains three ingredients which are imported here.
Its command-quotient theorem gives exact descent conditions and a harmonic
lift when a relevant action block is positive.  Its common-action Schur theorem
also assumes a nonnegative Hessian.  Its full same-cylinder identity, however,
states that the partition Hessian is the sum of a positive score covariance
and an independent direct second-action term of either sign.  The DeWitt
metric itself has a physical negative conformal direction.  Therefore the
Einstein screen must be reduced by verified null directions and tested by
singular values, not made positive by deleting negative directions.

This version constructs that signed finite screen.  Gauge and relation nulls
are removed by an exact quotient or, equivalently, a bordered system.  The
remaining positive normal operator measures invertibility without changing
the Hessian inertia.  A transported singular-value floor then supplies the
inverse-function hypothesis of V14.  Failure returns an explicit gauge
leakage, physical kernel, cokernel, or radius witness.

% -----------------------------------------------------------------------------
\subsection{Why Fisher positivity does not supply the action inverse}
\label{subsec:V15-Fisher-separator}
% -----------------------------------------------------------------------------

\begin{proposition}[Same score Fisher matrix, arbitrary signed action Hessian]
\label{prop:V15-Fisher-not-Hessian}
Let $X$ take values $\pm1$ with probability $1/2$, and for $c\in\mathbb R$
define the positive unnormalized deformation
\begin{equation}
 W_{c,\theta}(X)
 :=\exp\{\theta X+\tfrac12c\theta^2\}.
 \label{eq:V15-Wc}
\end{equation}
At $\theta=0$, every member has the same centered score $X$ and hence Fisher
information one.  Its log-partition Hessian is
\begin{equation}
 \left.\partial_\theta^2\log
 \mathbb E W_{c,\theta}\right|_{\theta=0}
 =1+c.
 \label{eq:V15-signed-Hessian}
\end{equation}
Thus the same positive Fisher block accompanies a positive, zero, or negative
total Hessian according as $c>-1$, $c=-1$, or $c<-1$.
\end{proposition}

\begin{proof}
The first log-weight derivative at zero is $X$ and the second is the constant
$c$.  Moreover
\[
 \log\mathbb E W_{c,\theta}
 =\frac c2\theta^2+\log\cosh\theta.
\]
Twice differentiating at zero gives $c+1$.  This is exactly the
direct-second-writer plus score-covariance decomposition.
\end{proof}

\begin{firewall}[The covariance normal operator must be formed after assembly]
\label{fire:V15-Fisher-scope}
The positive covariance of first scores may be used to measure a source frame,
but it cannot replace the signed derivative of the complete first-variation
map.  Gauge descent, counterterms, contact terms, determinant/Pfaffian
normalization, and the direct second action must be assembled before the
physical Hessian kernel and inertia are tested.
\end{firewall}

% -----------------------------------------------------------------------------
\subsection{Exact gauge--relation descent for a signed Hessian}
\label{subsec:V15-descent}
% -----------------------------------------------------------------------------

Let $\mathcal X$ be a finite real Hilbert command screen.  Let
\begin{equation}
 C:\mathcal G\longrightarrow\mathcal X
 \label{eq:V15-null-synthesis}
\end{equation}
be an injective, source-minimal synthesis of the declared gauge and
future-null relation directions.  Put
\begin{equation}
 \mathcal K:=\operatorname{Ran}C,\qquad
 P_{\mathcal K}:=C(C^*C)^{-1}C^*,\qquad
 P:=I-P_{\mathcal K}.
 \label{eq:V15-physical-projector}
\end{equation}
Let $\mathsf H=\mathsf H^*$ be the signed Hessian of the completely assembled
common action on this screen.

\begin{theorem}[Positive leakage certificate for signed quotient descent]
\label{thm:V15-leakage}
Define
\begin{equation}
 \mathbb L_C
 :=C^*\mathsf H^2C
 =(\mathsf HC)^*(\mathsf HC)\succeq0.
 \label{eq:V15-leakage-Gram}
\end{equation}
Then the following are equivalent:
\begin{equation}
 \mathbb L_C=0
 \quad\Longleftrightarrow\quad
 \mathsf HC=0
 \quad\Longleftrightarrow\quad
 \mathsf H=P\mathsf HP.
 \label{eq:V15-descent-equivalence}
\end{equation}
On this branch the bilinear Hessian form descends uniquely to
$\mathcal X/\mathcal K$, represented by
\begin{equation}
 \mathsf H_{\rm ph}:=
 P\mathsf HP\big|_{P\mathcal X}:P\mathcal X\to P\mathcal X.
 \label{eq:V15-physical-Hessian}
\end{equation}
If $\mathbb L_C\ne0$, the polar range of $\mathsf HC$ is the
source-minimal physical response generated by a purported null direction.
\end{theorem}

\begin{proof}
The first equivalence is the zero-Gram criterion.  If $\mathsf HC=0$, then
$\mathsf HP_{\mathcal K}=0$.  Self-adjointness gives
$P_{\mathcal K}\mathsf H=0$, hence $\mathsf H=P\mathsf HP$.
The converse is immediate because $PC=0$.

When these conditions hold, changing a representative by $Cg$ changes neither
slot of the bilinear form, so it defines a unique form on the quotient.
The orthogonal representative identifies that quotient with $P\mathcal X$.
For nonzero leakage, factorization through the polar range is the ordinary
source-minimal Gram realization.
\end{proof}

\begin{assessment}[Relation to the imported command quotient]
\label{ass:V15-command-quotient}
The imported command-quotient theorem tests first and mixed response descent
and uses a positive harmonic short when descent fails on a source-free
operation fibre.  Theorem~\ref{thm:V15-leakage} addresses the later signed
common-action Hessian: it gives a positive test for exact null descent without
requiring $\mathsf H\succeq0$ and does not minimize away any negative physical
direction.
\end{assessment}
% -----------------------------------------------------------------------------
\subsection{Supported singular floor and bordered gauge system}
\label{subsec:V15-bordered}
% -----------------------------------------------------------------------------

\begin{theorem}[The physical normal operator tests an indefinite Hessian]
\label{thm:V15-normal-floor}
Assume the descent conditions
\eqref{eq:V15-descent-equivalence} and define
\begin{equation}
 \mathsf N_{\rm ph}:=
 \mathsf H_{\rm ph}^*\mathsf H_{\rm ph}
 =\mathsf H_{\rm ph}^2\succeq0.
 \label{eq:V15-normal-operator}
\end{equation}
For $\sigma>0$, the following are equivalent:
\begin{align}
 \mathsf N_{\rm ph}&\succeq\sigma^2I_{P\mathcal X},
 \label{eq:V15-normal-floor}\\
 \|\mathsf H_{\rm ph}x\|&\ge\sigma\|x\|
 \quad(x\in P\mathcal X),
 \label{eq:V15-H-lower}\\
 \mathsf H_{\rm ph}\ \hbox{is invertible and}\quad
 \|\mathsf H_{\rm ph}^{-1}\|&\le\sigma^{-1}.
 \label{eq:V15-H-inverse}
\end{align}
The optimal $\sigma$ is the least absolute eigenvalue of
$\mathsf H_{\rm ph}$.  The test retains every positive and negative
eigenvalue of $\mathsf H_{\rm ph}$; it removes only the verified null space
$\mathcal K$.

If the floor fails at zero, a unit eigenvector
$v\in P\mathcal X$ with $\mathsf H_{\rm ph}v=0$ is an explicit physical soft
direction.  If a proposed positive floor $\sigma$ fails, a unit right singular
vector gives
$\|\mathsf H_{\rm ph}v\|<\sigma$.
\end{theorem}

\begin{proof}
The first two conditions are identical after taking the quadratic form of
$\mathsf H_{\rm ph}^*\mathsf H_{\rm ph}$.  In finite equal-dimensional domain
and codomain, a lower bound makes $\mathsf H_{\rm ph}$ injective and hence
bijective; the inverse bound is immediate.  Conversely an inverse bound gives
\eqref{eq:V15-H-lower}.  The spectral theorem for the self-adjoint
$\mathsf H_{\rm ph}$ identifies the optimal constant and the witnesses.
Squaring changes eigenvalues $\lambda$ to $\lambda^2$ but does not alter the
signed Hessian used in the action equation.
\end{proof}

\begin{theorem}[Bordered system removes nulls without changing physical inertia]
\label{thm:V15-bordered-system}
Assume \eqref{eq:V15-descent-equivalence},
\eqref{eq:V15-normal-floor}, and
\begin{equation}
 C^*C\succeq\gamma^2I_{\mathcal G}
 \qquad(\gamma>0).
 \label{eq:V15-C-floor}
\end{equation}
Define the symmetric bordered operator on $\mathcal X\oplus\mathcal G$,
\begin{equation}
 \mathbb B_C:=
 \begin{pmatrix}
 \mathsf H&C\\
 C^*&0
 \end{pmatrix}.
 \label{eq:V15-bordered-operator}
\end{equation}
Then $\mathbb B_C$ is invertible and
\begin{equation}
 s_{\min}(\mathbb B_C)
 =\min\{s_{\min}(\mathsf H_{\rm ph}),s_{\min}(C)\}
 \ge\min\{\sigma,\gamma\}.
 \label{eq:V15-bordered-floor}
\end{equation}
If $\operatorname{In}(\mathsf H_{\rm ph})=(n_+,n_-,0)$ and
$g=\dim\mathcal G$, then
\begin{equation}
 \operatorname{In}(\mathbb B_C)=(n_++g,n_-+g,0).
 \label{eq:V15-bordered-inertia}
\end{equation}
In particular, a physical negative conformal direction remains negative.

For every $f\in\mathcal X$, the constrained system
\begin{equation}
 \mathsf Hx+C\lambda=f,\qquad C^*x=0
 \label{eq:V15-KKT-system}
\end{equation}
has the unique solution
\begin{equation}
 x=\mathsf H_{\rm ph}^{-1}Pf,\qquad
 \lambda=(C^*C)^{-1}C^*f.
 \label{eq:V15-KKT-solution}
\end{equation}
\end{theorem}

\begin{proof}
Use the orthogonal decomposition
$\mathcal X=P\mathcal X\oplus\mathcal K$.  Exact descent gives
$\mathsf H=\mathsf H_{\rm ph}\oplus0$.  Since $C$ is a bijection from
$\mathcal G$ onto $\mathcal K$, the bordered operator is orthogonally
equivalent to the direct sum of $\mathsf H_{\rm ph}$ and
\[
 \begin{pmatrix}0&C\\C^*&0\end{pmatrix}
 \quad\hbox{on }\mathcal K\oplus\mathcal G.
\]
The latter has eigenvalues $\pm s_j(C)$.  This proves the singular-value and
inertia statements.

The constraint gives $x\in P\mathcal X$.  Projecting the first equation by
$P$ gives $\mathsf H_{\rm ph}x=Pf$.  Projecting onto $\mathcal K$ gives
$C\lambda=P_{\mathcal K}f$, whose source-minimal solution is the displayed
formula.
\end{proof}

\begin{corollary}[Stability under a retained descent defect]
\label{cor:V15-bordered-perturbation}
Let $\mathbb B_{C,0}$ satisfy
Theorem~\ref{thm:V15-bordered-system} with
$m_0=\min\{\sigma,\gamma\}$.  If a completely assembled bordered operator is
\begin{equation}
 \mathbb B_C=\mathbb B_{C,0}+\mathbb E,
 \qquad \|\mathbb E\|<m_0,
 \label{eq:V15-bordered-perturbation}
\end{equation}
then
\begin{equation}
 s_{\min}(\mathbb B_C)\ge m_0-\|\mathbb E\|,
 \qquad
 \|\mathbb B_C^{-1}\|
 \le(m_0-\|\mathbb E\|)^{-1}.
 \label{eq:V15-bordered-perturbed-floor}
\end{equation}
The error $\mathbb E$ must remain in the residual ledger; this estimate does
not declare approximate gauge leakage to be zero.
\end{corollary}

\begin{proof}
For every unit vector $u$,
\[
 \|\mathbb B_Cu\|
 \ge\|\mathbb B_{C,0}u\|-\|\mathbb Eu\|
 \ge m_0-\|\mathbb E\|.
\]
Finite dimensionality gives invertibility and the inverse bound.
\end{proof}

% -----------------------------------------------------------------------------
\subsection{Two-sided quotient for a represented residual map}
\label{subsec:V15-two-sided}
% -----------------------------------------------------------------------------

Let $\mathcal Y$ be a finite represented-residual screen and let
$R:\mathcal N\to\mathcal Y$ be an injective source-minimal synthesis of
declared target nuisance rows.  Put
\begin{equation}
 Q:=I-R(R^*R)^{-1}R^*.
 \label{eq:V15-target-projector}
\end{equation}

\begin{theorem}[Domain descent, determining coverage, and singular witnesses]
\label{thm:V15-two-sided-quotient}
Let $J:\mathcal X\to\mathcal Y$ be the derivative of a represented
first-variation map and define
\begin{equation}
 A:=QJP:P\mathcal X\longrightarrow Q\mathcal Y,
 \qquad
 \mathbb L_{J,C}:=C^*J^*QJC\succeq0.
 \label{eq:V15-reduced-J}
\end{equation}
Then $QJ$ descends through the domain quotient
$\mathcal X/\mathcal K$ exactly when
\begin{equation}
 \mathbb L_{J,C}=0
 \quad\Longleftrightarrow\quad QJC=0.
 \label{eq:V15-J-descent}
\end{equation}
On that branch:
\begin{enumerate}[label=\textup{(\roman*)},leftmargin=2.8em]
\item $A$ is injective with lower constant $\sigma_X>0$ exactly when
      \[
       A^*A\succeq\sigma_X^2I_{P\mathcal X};
       \tag{V15.Q1}
      \]
\item $A$ covers the determining target with lower constant $\sigma_Y>0$
      exactly when
      \[
       AA^*\succeq\sigma_Y^2I_{Q\mathcal Y};
       \tag{V15.Q2}
      \]
\item $A$ is bijective exactly when both conditions hold; in equal dimensions
      either one implies the other, and
      \[
       \|A^{-1}\|=s_{\min}(A)^{-1}.
       \tag{V15.Q3}
      \]
\end{enumerate}
Failure of domain descent gives a right singular vector of $QJC$.
Failure of \textup{(V15.Q1)} gives a physical kernel or soft right singular
vector.  Failure of \textup{(V15.Q2)} gives a target cokernel or soft left
singular vector.
\end{theorem}

\begin{proof}
The zero-Gram criterion gives \eqref{eq:V15-J-descent}.  If it holds,
$QJ(x+Cg)=QJx$, so $A$ is the orthogonal-representative form of the quotient
map.  The two normal-operator inequalities are respectively the lower
singular-value tests for $A$ and $A^*$.  In finite dimensions they characterize
injectivity and surjectivity.  Singular-value decomposition supplies all
claimed witnesses and the inverse norm.
\end{proof}
% -----------------------------------------------------------------------------
\subsection{Transported floor implies the V14 inverse radius}
\label{subsec:V15-transport}
% -----------------------------------------------------------------------------

Let $\mathcal X_*$ and $\mathcal Y_*$ be fixed source-minimal physical screens.
At cutoff $n$, suppose the relation/gauge and determining-target projectors
have already passed their Kato transport tests, giving isometries
\begin{equation}
 U_n:\mathcal X_*\to P_n\mathcal X_n,
 \qquad
 V_n:\mathcal Y_*\to Q_n\mathcal Y_n.
 \label{eq:V15-screen-isometries}
\end{equation}
These are the transported quotient frames; their construction is imported
from the earlier support-transport notes.

Let $\mathcal F_n$ be the completely assembled represented first-variation
map on the cutoff chart and define
\begin{align}
 \widehat F_n(z)
 &:=
 V_n^*Q_n\mathcal F_n(U_nz),
 \label{eq:V15-aligned-residual}\\
 \widehat J_n(z)
 &:=
 D\widehat F_n(z)
 =V_n^*Q_nD\mathcal F_n(U_nz)U_n.
 \label{eq:V15-aligned-Jacobian}
\end{align}

\begin{theorem}[Supported Hessian floor gives the nonlinear inverse radius]
\label{thm:V15-floor-to-radius}
Assume $\dim\mathcal X_*=\dim\mathcal Y_*$ and let
$H_*:\mathcal X_*\to\mathcal Y_*$ be invertible with
\begin{equation}
 s_{\min}(H_*)=\sigma_*>0.
 \label{eq:V15-base-floor}
\end{equation}
Suppose exact domain descent and target coverage hold at every cutoff and
\begin{equation}
 \sup_n\sup_{\|z\|\le R}
 \|\widehat J_n(z)-H_*\|\le\delta<\sigma_*.
 \label{eq:V15-uniform-J-defect}
\end{equation}
Then
\begin{align}
 s_{\min}(\widehat J_n(z))
 &\ge\sigma_*-\delta,
 \label{eq:V15-J-floor}\\
 \|I-H_*^{-1}\widehat J_n(z)\|
 &\le\theta:=\frac{\delta}{\sigma_*}<1.
 \label{eq:V15-V14-theta}
\end{align}
Thus \eqref{eq:V15-V14-theta} is precisely the uniform linearization
hypothesis of Theorem~\ref{thm:V14-residual-transport} with $H=H_*$.

A sufficient check for \eqref{eq:V15-uniform-J-defect} is
\begin{equation}
 \sup_n\|\widehat J_n(0)-H_*\|\le d,
 \qquad
 \sup_n\sup_{\|z\|\le R}
 \|D\widehat F_n(z)-D\widehat F_n(0)\|\le L\|z\|,
 \label{eq:V15-center-Lipschitz}
\end{equation}
with
\begin{equation}
 d+LR<\sigma_*.
 \label{eq:V15-radius-budget}
\end{equation}
\end{theorem}

\begin{proof}
For every unit $x$,
\[
 \|\widehat J_n(z)x\|
 \ge\|H_*x\|-\|(\widehat J_n(z)-H_*)x\|
 \ge\sigma_*-\delta.
\]
Also $\|H_*^{-1}\|=\sigma_*^{-1}$, so
\[
 \|I-H_*^{-1}\widehat J_n(z)\|
 \le\|H_*^{-1}\|\,\|H_*-\widehat J_n(z)\|
 \le\delta/\sigma_*.
\]
The bounds in \eqref{eq:V15-center-Lipschitz} give
$\|\widehat J_n(z)-H_*\|\le d+L\|z\|\le d+LR$.
\end{proof}

\begin{corollary}[Action-selected noncollapsing Einstein scheduler branch]
\label{cor:V15-action-selected-branch}
In addition to Theorem~\ref{thm:V15-floor-to-radius}, suppose
\begin{align}
 \sup_n\|H_*^{-1}\widehat F_n(0)\|
 &\le(1-\theta)R,
 \label{eq:V15-center-force}\\
 \sum_n\sup_{\|z\|\le R}
 \|H_*^{-1}(\widehat F_{n+1}(z)-\widehat F_n(z))\|
 &<\infty,
 \label{eq:V15-action-defect-sum}
\end{align}
and the resulting root-motion budget fits the V14 logarithmic scheduler ball.
Then:
\begin{enumerate}[label=\textup{(\roman*)},leftmargin=2.8em]
\item every cutoff has a unique physical scheduler root
      $\widehat F_n(z_n)=0$ in the ball;
\item the roots converge with the explicit V14 tail bound;
\item their positive $A_3$ rate lifts have the uniform bounds of V14;
\item the coframe, inverse metric, and represented rational stress writers
      converge strongly as in V14;
\item the limiting represented scheduler residual vanishes.
\end{enumerate}
In the self-adjoint action case
$\mathcal Y_n=\mathcal X_n$, $Q_n=P_n$, and $V_n=U_n$, if the nonlinear
Ward identities make the complete first variation orthogonal to
$\mathcal K_n$ throughout the ball, the quotient zero lifts to the complete
constrained zero.  The bordered system then supplies its unique gauge
multiplier.
\end{corollary}

\begin{proof}
Equations \eqref{eq:V15-V14-theta,eq:V15-center-force} and
\eqref{eq:V15-action-defect-sum} are the hypotheses of
Theorem~\ref{thm:V14-residual-transport}.  Its action-root conclusion and the
V14 scheduler construction give items~\textup{(i)--(v)}.  On the stated self-adjoint branch, exact nonlinear Ward descent puts the
residual in $P_n\mathcal X_n$; its projected zero is therefore its complete
zero.  Theorem~\ref{thm:V15-bordered-system} uniquely enforces the gauge
constraint.
\end{proof}

\begin{corollary}[A fixed supported floor gives an explicit radius]
\label{cor:V15-explicit-radius}
Under \eqref{eq:V15-center-Lipschitz}, if $L>0$ and $d<\sigma_*$, every
\begin{equation}
 0<R<\frac{\sigma_*-d}{L}
 \label{eq:V15-explicit-radius}
\end{equation}
passes the nonlinear inverse test.  If $L=0$, any radius contained in the
physical chart passes provided $d<\sigma_*$.
\end{corollary}

\begin{proof}
This is exactly \eqref{eq:V15-radius-budget}.
\end{proof}

% -----------------------------------------------------------------------------
\subsection{The executable same-cylinder Hessian card}
\label{subsec:V15-card}
% -----------------------------------------------------------------------------

\begin{definition}[Signed physical Hessian card]
\label{def:V15-Hessian-card}
On one predeclared cutoff command screen, a signed physical Hessian card
contains:
\begin{enumerate}[label=\textup{(H\arabic*)},leftmargin=3em]
\item the complete first log-weight writers $a_i$ and direct second writers
      $b_{ij}$ on one unnormalized cylinder;
\item the signed cylinder-response matrix
      \begin{equation}
       (\mathsf H_n^{\rm cyl})_{ij}
       =\mathbb E_n b_{ij}
        +\operatorname{Cov}_n(a_i,a_j);
       \label{eq:V15-same-cylinder-H}
      \end{equation}
\item the declared first-variation residual $\mathcal F_n$, its derivative
      $D\mathcal F_n(0)$ on the same commands, and the incidence Gram
      \begin{equation}
       \mathbb I_n^{\rm Hess}
       :=
       (D\mathcal F_n(0)-\mathsf H_n^{\rm cyl})^*
       (D\mathcal F_n(0)-\mathsf H_n^{\rm cyl});
       \label{eq:V15-Hessian-incidence}
      \end{equation}
\item source-minimal gauge/relation and target-nuisance syntheses $C_n,R_n$,
      their supported floors, and the leakage Grams
      \eqref{eq:V15-leakage-Gram,eq:V15-reduced-J};
\item the quotient Hessian or represented Jacobian, its complete singular
      values, and its signed inertia when self-adjoint;
\item the Kato screen transports $U_n,V_n$, the center mismatch $d_n$, a
      verified nonlinear Lipschitz constant $L_n$, the center force, and the
      adjacent action defect;
\item one held-out palindromic second command and one held-out Ward direction
      testing the fitted Hessian and the declared null space.
\end{enumerate}
\end{definition}

\begin{theorem}[The card is a terminating finite pass-or-witness test]
\label{thm:V15-card-alternative}
For every finite populated card in
Definition~\ref{def:V15-Hessian-card}, the pass branch below occurs exactly
when all tests pass; otherwise at least one of the displayed witness branches
is returned:
\begin{enumerate}[label=\textup{(\alph*)},leftmargin=2.8em]
\item the Hessian-incidence, descent, coverage, singular-floor, held-out,
      inverse-radius, and adjacent-defect tests all pass, producing the
      constants required by
      Corollary~\ref{cor:V15-action-selected-branch};
\item a nonzero gauge/relation leakage vector is obtained from
      \eqref{eq:V15-leakage-Gram} or \eqref{eq:V15-J-descent};
\item a physical right soft mode or determining-target left soft mode is
      obtained from the quotient singular-value decomposition;
\item the fitted second derivative or Ward kernel fails its held-out command;
\item an attained center, nonlinear-radius, support-transport, or adjacent
      action-defect witness is returned.
\end{enumerate}
The cylinder matrix is used as the action-residual derivative only when
$\mathbb I_n^{\rm Hess}=0$.  No positive Fisher or Schur block is substituted
for the signed derivative.
\end{theorem}

\begin{proof}
All spaces and matrices are finite.  The Gram
\eqref{eq:V15-Hessian-incidence} decides Hessian incidence; the other Gram
zero tests decide descent, singular value decomposition decides domain and
target floors, and signed diagonalization gives inertia.  Continuity on the compact command ball makes
the nonlinear and adjacent suprema attained.  The two held-out rows are direct
evaluations.  These alternatives exhaust the finitely listed tests.
\end{proof}
\begin{example}[Populated indefinite bordered branch]
\label{ex:V15-indefinite-populated}
Let $\mathcal X=\mathbb R^4$, $\mathcal G=\mathbb R$,
\[
 \mathsf H=\operatorname{diag}(-2,3,5,0),
 \qquad
 C\lambda=(0,0,0,\lambda).
\]
Then $\mathbb L_C=0$,
$\mathsf H_{\rm ph}=\operatorname{diag}(-2,3,5)$,
$\sigma=2$, and $\gamma=1$.  The bordered operator has eigenvalues
\[
 -2,\ 3,\ 5,\ -1,\ 1
\]
and inertia $(3,2,0)$.  Thus the negative physical direction is retained while
the gauge zero is replaced by one positive and one negative multiplier
direction.

For $f\in\mathbb R^3$ with
$\|\mathsf H_{\rm ph}^{-1}f\|\le R/2$, put
$f_n=(1-2^{-n})f$ and
\[
 \mathscr A_n(x)
 =\frac12\langle x,\mathsf Hx\rangle
  -\langle(f_n,0),x\rangle.
\]
On the constraint $C^*x=0$, the unique stationary point is
\[
 x_n=(\mathsf H_{\rm ph}^{-1}f_n,0),
 \qquad
 \|x_\infty-x_n\|
 =2^{-n}\|\mathsf H_{\rm ph}^{-1}f\|.
\]
This is an exact summable action-selected branch with an indefinite physical
Hessian.
\end{example}

\begin{proof}
All assertions follow by diagonalization.  The stationary equation on the
physical subspace is
$\mathsf H_{\rm ph}x_{\rm ph}=f_n$, and the displayed tail follows by
subtraction.
\end{proof}

% -----------------------------------------------------------------------------
\subsection{From finite supported floors to an infinite determining screen}
\label{subsec:V15-cofinal-floor}
% -----------------------------------------------------------------------------

\begin{theorem}[Cofinal physical Hessian floor or soft sequence]
\label{thm:V15-cofinal-floor}
Let $\mathcal X_\infty$ be a real Hilbert determining space, let
$P_m\uparrow I$ be increasing finite-rank orthogonal projections, and let
$\mathsf H=\mathsf H^*\in\mathcal B(\mathcal X_\infty)$.  Set
\begin{equation}
 \mathsf H_m:=P_m\mathsf HP_m\big|_{P_m\mathcal X_\infty}.
 \label{eq:V15-finite-compressions}
\end{equation}
If
\begin{equation}
 \mathsf H_m^2\succeq\sigma^2I_{P_m\mathcal X_\infty}
 \qquad\hbox{for every }m
 \label{eq:V15-uniform-cofinal-floor}
\end{equation}
with one $\sigma>0$, then
\begin{equation}
 \|\mathsf Hx\|\ge\sigma\|x\|
 \quad(x\in\mathcal X_\infty),
 \qquad
 \mathsf H^{-1}\in\mathcal B(\mathcal X_\infty),
 \qquad
 \|\mathsf H^{-1}\|\le\sigma^{-1}.
 \label{eq:V15-infinite-inverse}
\end{equation}

Conversely, suppose
$s_{\min}(\mathsf H_m)\to0$ along a subsequence and the exterior leakage
satisfies
\begin{equation}
 \ell_m:=\|(I-P_m)\mathsf HP_m\|\longrightarrow0.
 \label{eq:V15-exterior-leakage}
\end{equation}
Then there are unit vectors $v_m\in P_m\mathcal X_\infty$ such that
\begin{equation}
 \|\mathsf Hv_m\|
 \le s_{\min}(\mathsf H_m)+\ell_m\longrightarrow0.
 \label{eq:V15-Weyl-soft}
\end{equation}
After passage to a weakly convergent subsequence
$v_m\rightharpoonup v$, either $v\ne0$ and $v$ is an actual kernel vector, or
$v=0$ and the $v_m$ form an explicit escaping cofinal soft-mode sequence.
\end{theorem}

\begin{proof}
Let $x$ belong to the dense union $\bigcup_mP_m\mathcal X_\infty$, and choose
$m_0$ with $x=P_{m_0}x$.  For $m\ge m_0$,
\[
 \|P_m\mathsf Hx\|
 =\|\mathsf H_mx\|\ge\sigma\|x\|.
\]
Strong convergence $P_m\to I$ gives
$\|\mathsf Hx\|\ge\sigma\|x\|$.  Density and boundedness extend this estimate
to every $x$.  Hence $\mathsf H$ is injective with closed range.  Since it is
self-adjoint,
$(\operatorname{Ran}\mathsf H)^\perp=\ker\mathsf H=\{0\}$, so the closed range
is all of $\mathcal X_\infty$.  This proves \eqref{eq:V15-infinite-inverse}.

For the converse, take a unit right singular vector $v_m$ of
$\mathsf H_m$ at its least singular value.  Then
\[
 \|\mathsf Hv_m\|
 \le\|P_m\mathsf Hv_m\|
    +\|(I-P_m)\mathsf HP_mv_m\|
 \le s_{\min}(\mathsf H_m)+\ell_m.
\]
The unit ball is weakly compact, so pass to
$v_m\rightharpoonup v$.  Since $\mathsf Hv_m\to0$ strongly and $\mathsf H$ is
bounded, weak continuity gives $\mathsf Hv=0$.  A nonzero $v$ is therefore an
actual kernel vector; $v=0$ is precisely the weakly escaping soft-mode case.
\end{proof}

\begin{corollary}[Cofinal bordered inverse]
\label{cor:V15-cofinal-bordered}
Suppose the transported finite gauge/relation syntheses have one source floor
$\gamma>0$, exact leakage zero, and their physical quotient screens identify
with the nested $P_m\mathcal X_\infty$ above.  If
\eqref{eq:V15-uniform-cofinal-floor} holds, the finite bordered inverses have
the uniform bound
\begin{equation}
 \sup_m\|\mathbb B_{C_m}^{-1}\|
 \le\max\{\sigma^{-1},\gamma^{-1}\}.
 \label{eq:V15-uniform-bordered-inverse}
\end{equation}
If the physical floor collapses and
\eqref{eq:V15-exterior-leakage} holds, the vectors
\eqref{eq:V15-Weyl-soft} are physical rather than gauge modes.
\end{corollary}

\begin{proof}
Apply Theorem~\ref{thm:V15-bordered-system} at every finite screen and
Theorem~\ref{thm:V15-cofinal-floor} to the physical blocks.  Exact leakage
places the soft vectors in the physical quotient.
\end{proof}
\begin{firewall}[The cofinal operator is Riesz normalized]
\label{fire:V15-bounded-normalization}
Theorem~\ref{thm:V15-cofinal-floor} concerns the bounded Riesz representative
of the Hessian in a declared physical action or graph norm.  A raw
differential Einstein Hessian on $L^2$ is generally unbounded.  Passing from
the normalized operator back to that differential realization requires the
separate common-domain, closed-form, and Mosco or resolvent convergence rows
already isolated in the monograph.
\end{firewall}

% -----------------------------------------------------------------------------
\subsection{V15 theorem ledger and next Palatini acquisition}
\label{subsec:V15-ledger}
% -----------------------------------------------------------------------------

\paragraph{New conclusions proved in V15.}
Proposition~\ref{prop:V15-Fisher-not-Hessian} gives an exact positive-cylinder
separator showing that one Fisher matrix is compatible with every Hessian
sign and with a zero Hessian.  Theorem~\ref{thm:V15-leakage} constructs a
positive leakage Gram which tests exact gauge/relation descent of a signed
common-action Hessian.  Theorem~\ref{thm:V15-normal-floor} uses the physical
normal operator to test invertibility while retaining signed inertia.
Theorem~\ref{thm:V15-bordered-system} proves that the bordered gauge system
is invertible with an explicit floor and preserves every negative physical
direction.  Theorem~\ref{thm:V15-two-sided-quotient} adds target coverage and
separate right/left singular witnesses.  Theorem
\ref{thm:V15-floor-to-radius} converts the supported finite Hessian floor into
the exact nonlinear radius required by V14.  Theorem
\ref{thm:V15-card-alternative} packages these quantities into a terminating
same-cylinder test.  Finally, Theorem~\ref{thm:V15-cofinal-floor} promotes a
uniform nested finite-screen floor to an inverse on the infinite determining
space, or returns an actual kernel or escaping soft sequence.

\paragraph{Imported results not repeated.}
The command-operation descent test, positive harmonic lift, positive
common-action Schur residual, full same-cylinder second-derivative identity,
palindromic Hessian command, Kato support transport, DeWitt inertia, and
conditional Einstein handoff are imported from the attached monograph and
earlier notes.  V15 supplies the missing signed quotient, bordered system,
Hessian-incidence row, and cofinal inverse alternative.

\paragraph{Authoritative status of the physical card.}
The attached RPESM population theorem says that action, Ward, stress, and
second-action rows occur on one finite event law.  It does not print the
selected command matrices
\[
 \mathsf H_n^{\rm cyl},\quad
 D\mathcal F_n(0),\quad C_n,\quad R_n,\quad U_n,\quad V_n
\]
or cutoff-uniform singular values and nonlinear radii for an active
interacting Einstein--Standard Model history.  Therefore none of
\eqref{eq:V15-Hessian-incidence,eq:V15-normal-floor} is asserted for
that physical sequence.

\paragraph{Next acquisition.}
On one selected projective history and increasing compact determining screens,
the next note must:
\begin{enumerate}[label=\textup{(\arabic*)},leftmargin=2.8em]
\item select the physical action/graph norm, form its bounded Riesz
      representative, and evaluate the direct second-action and covariance
      pieces of \eqref{eq:V15-same-cylinder-H} on the same commands;
\item verify $\mathbb I_n^{\rm Hess}=0$ and both gauge/relation leakage Grams;
\item diagonalize the signed quotient Hessian and either obtain a common
      $\sigma_*>0$ or retain its soft sequence;
\item bound the nonlinear Hessian oscillation $L_n$, center force, Kato support
      debit, exterior leakage, and adjacent total-action defect;
\item on the pass branch, insert the resulting action-selected scheduler roots
      into the localized Palatini first variation and prove convergence on a
      dense compactly supported test core.
\end{enumerate}
Only item~\textup{(5)}, together with the remaining Ward, counterterm, memory,
and boundary rows, can promote the finite stationary roots to the
distributional Einstein--Standard Model equation.

\begin{firewall}[No Hessian card has been numerically populated]
\label{fire:V15-no-population}
V15 proves a finite and cofinal certification theorem.  It does not supply the
missing same-cylinder numerical Hessian of the selected interacting RPESM
history.  The full continuum is therefore not claimed.  If the supported
floor collapses, the correct output is the physical soft sequence of
Theorem~\ref{thm:V15-cofinal-floor}, not deletion of the negative or soft
direction.
\end{firewall}

% =============================================================================
% END APPENDED VERSION V15 MATERIAL
% =============================================================================


% =============================================================================
% BEGIN APPENDED VERSION V16 MATERIAL
% =============================================================================

\clearpage
\section{Version V16: curvature lift and same-core Einstein closure}
\label{sec:V16-Einstein-closure}

Version V15 reduced physical action selection to an indefinite
gauge--relation quotient with a supported singular-value floor.  Its pass
branch produces summably transported scheduler roots, but the V14 regularity
theorem only controlled pointwise ADM writers and quadratic velocity
insertions.  Curvature uses two spacetime derivatives, and the stress which is
represented on the Osterwalder--Schrodinger quotient must still be identified
with the metric derivative of the same action whose root was selected.

Several stronger statements are already present in the attached monograph and
are imported rather than repeated.  The localized renewal--Palatini theorem
passes a strong-coframe/weak-curvature product.  The assembled spectral-tail
theorem promotes a finite Palatini residual to an exhausting variation bank.
The finite and cofinal stress/current theorems construct represented Ward
stress on a common graph core.  The conditional Einstein handoff then closes
once all those objects are known to be variations of one common action on one
determining test core.

This note supplies the synchronization between those interfaces.  First it
raises the summable logarithmic scheduler from $C^{0,\alpha}$ to
$C^{2,\alpha}$ and obtains an explicit tail for the Levi--Civita connection,
curvature, and densitized Einstein writer.  Second it gives an $H^1$ incidence
test which identifies an independently occurring chronological connection
with that scheduler lift.  Third it places action stationarity, Palatini
incidence, and stress incidence in one Riesz-normalized test space.  Their
signed sum is exactly the Einstein residual.  This isolates the remaining
physical acquisitions without declaring any of them from kinematics.

% -----------------------------------------------------------------------------
\subsection{A summable two-derivative scheduler lift}
\label{subsec:V16-curvature-lift}
% -----------------------------------------------------------------------------

Let $K$ be the closure of a bounded coordinate cylinder with Lipschitz
boundary, compactly contained in one spacetime chart, and fix
$0<\alpha<1$.  Retain the ten-dimensional logarithmic scheduler space
$\mathcal Z$ and analytic ADM map of V14.  Thus
\[
 z=(x,y,A,b),\qquad
 \varrho=e^x,\quad \upsilon=e^y,\quad
 \mathcal S=e^A,\quad \beta=b,
\]
where $A$ is symmetric and trace free.  Write
\[
 \Phi(z)=(\mathbf g(z),\mathsf e(z))
\]
for the full Lorentz metric and the fixed-gauge coframe constructed by that
map.

Suppose the projective scheduler has the pathwise expansion
\begin{equation}
 z_n=z_0+\sum_{j=1}^na_j\xi_j
 \label{eq:V16-C2-series}
\end{equation}
and define its two-derivative tail
\begin{equation}
 \rho_n^{(2,\alpha)}
 :=
 \sum_{j>n}a_j\|\xi_j\|_{C^{2,\alpha}(K;\mathcal Z)}.
 \label{eq:V16-C2-tail}
\end{equation}

\begin{theorem}[Quantitative Levi--Civita and curvature lift]
\label{thm:V16-curvature-lift}
Assume
\begin{equation}
 \|z_0\|_{C^{2,\alpha}}
 +\sum_{j\ge1}a_j
   \|\xi_j\|_{C^{2,\alpha}}
 \le R_0,
 \qquad
 \rho_n^{(2,\alpha)}\longrightarrow0,
 \label{eq:V16-C2-budget}
\end{equation}
where $R_0$ is contained in the V14 noncollapse ball.  Then
$z_n\to z_\infty$ in $C^{2,\alpha}(K;\mathcal Z)$ and there is a constant
$C_K$, independent of $n$, such that
\begin{equation}
 \|\mathbf g_n-\mathbf g_\infty\|_{C^{2,\alpha}}
 +\|\mathbf g_n^{-1}-\mathbf g_\infty^{-1}\|_{C^{2,\alpha}}
 +\|\mathsf e_n-\mathsf e_\infty\|_{C^{2,\alpha}}
 +\|\mathsf e_n^{-1}-\mathsf e_\infty^{-1}\|_{C^{2,\alpha}}
 \le C_K\rho_n^{(2,\alpha)}.
 \label{eq:V16-metric-tail}
\end{equation}
All metrics have one common Lorentz signature and a uniform determinant and
inverse-metric floor.

Let $\Gamma_n$ be the Levi--Civita coordinate connection of $\mathbf g_n$,
let $\omega_n^{\rm LC}$ be its spin connection in the coframe
$\mathsf e_n$, and let
\[
 \operatorname{Riem}_n,\qquad F_n^{\rm LC},\qquad G_n
\]
be their Riemann, Lorentz curvature, and Einstein tensors.  Then
\begin{align}
 \|\Gamma_n-\Gamma_\infty\|_{C^{1,\alpha}}
 +\|\omega_n^{\rm LC}-\omega_\infty^{\rm LC}\|_{C^{1,\alpha}}
 &\le C_K\rho_n^{(2,\alpha)},
 \label{eq:V16-connection-tail}\\
 \|\operatorname{Riem}_n-\operatorname{Riem}_\infty\|_{C^{0,\alpha}}
 +\|F_n^{\rm LC}-F_\infty^{\rm LC}\|_{C^{0,\alpha}}
 +\|G_n-G_\infty\|_{C^{0,\alpha}}
 &\le C_K\rho_n^{(2,\alpha)}.
 \label{eq:V16-curvature-tail}
\end{align}
For fixed $\Lambda\in\mathbb R$, the densitized Einstein writer
\begin{equation}
 \mathfrak G_n
 :=
 \sqrt{-\det\mathbf g_n}\,
 (G_n+\Lambda\mathbf g_n)
 \label{eq:V16-densitized-Einstein}
\end{equation}
obeys
\begin{equation}
 \|\mathfrak G_n-\mathfrak G_\infty\|_{C^{0,\alpha}}
 \le C_{K,\Lambda}\rho_n^{(2,\alpha)}.
 \label{eq:V16-Einstein-tail}
\end{equation}
The coframe $\mathsf e_n$ is exactly torsion free with respect to
$\omega_n^{\rm LC}$ at every $n$.
\end{theorem}

\begin{proof}
The triangle inequality in the Banach space
$C^{2,\alpha}(K;\mathcal Z)$ gives
\[
 \|z_\infty-z_n\|_{C^{2,\alpha}}
 \le\rho_n^{(2,\alpha)}.
\]
In particular every pointwise value remains in the compact logarithmic ball
of V14.  On a neighborhood of that compact set, all derivatives of the
finite-dimensional analytic maps
\[
 z\longmapsto
 \mathbf g,\ \mathbf g^{-1},\ \mathsf e,\ \mathsf e^{-1},
 \ \sqrt{-\det\mathbf g}
\]
are bounded.  The composition theorem and product estimate in the
$C^{2,\alpha}$ Banach algebra therefore give
\eqref{eq:V16-metric-tail}.  The pointwise V14 floors persist on $K$ and give
the asserted determinant and inverse bounds.

In coordinates,
\[
 \Gamma^\lambda_{\mu\nu}(\mathbf g)
 =
 \frac12\mathbf g^{\lambda\rho}
 \bigl(
  \partial_\mu\mathbf g_{\rho\nu}
 +\partial_\nu\mathbf g_{\rho\mu}
 -\partial_\rho\mathbf g_{\mu\nu}
 \bigr).
\]
Subtracting the formulas at $n$ and at infinity, using
\eqref{eq:V16-metric-tail}, and applying the
$C^{1,\alpha}$ product estimate proves the first term of
\eqref{eq:V16-connection-tail}.  The spin connection is an algebraic sum of
products of $\mathsf e$, $\mathsf e^{-1}$, $\partial\mathsf e$, and
$\Gamma$.  The same estimate proves the second term.

The coordinate and spin curvatures have the forms
\[
 \operatorname{Riem}=\partial\Gamma+\Gamma\Gamma,
 \qquad
 F^{\rm LC}=\dd\omega^{\rm LC}
             +\omega^{\rm LC}\wedge\omega^{\rm LC}.
\]
Their differences are controlled in $C^{0,\alpha}$ by
\eqref{eq:V16-connection-tail} and the uniform $C^{1,\alpha}$ bounds.
Ricci contraction and scalar contraction use only the uniformly bounded
inverse metric.  This proves \eqref{eq:V16-curvature-tail}.  Multiplication
by the determinant density and addition of the cosmological term prove
\eqref{eq:V16-Einstein-tail}.  Finally, $\omega_n^{\rm LC}$ is defined as
the unique metric-compatible connection satisfying
$\dd\mathsf e_n+\omega_n^{\rm LC}\wedge\mathsf e_n=0$, so its torsion
vanishes identically.
\end{proof}

\begin{corollary}[The lifted scheduler meets the analytic Palatini compactness row]
\label{cor:V16-Palatini-compactness}
Under Theorem~\ref{thm:V16-curvature-lift},
\[
 \mathsf e_n\longrightarrow\mathsf e_\infty
 \quad\hbox{strongly in }L^q(K)\quad(1\le q\le\infty),
\]
and
\[
 F_n^{\rm LC}\longrightarrow F_\infty^{\rm LC}
 \quad\hbox{strongly in }L^p(K)\quad(1\le p\le\infty).
\]
Thus, on the Levi--Civita branch, the strong $L^2$ coframe, uniform $L^6$
coframe, and weak $L^p$ curvature assumptions of the imported localized
Palatini theorem hold, in particular with $p=2$.  Moreover, for every
compactly supported inverse-metric test $k$,
\begin{equation}
 \left|
 \chi\int_K
  \langle\mathfrak G_n-\mathfrak G_\infty,k\rangle
 \right|
 \le
 |\chi|C_{K,\Lambda}\rho_n^{(2,\alpha)}
 \|k\|_{L^1(K)}.
 \label{eq:V16-weak-Einstein-tail}
\end{equation}
\end{corollary}

\begin{proof}
The strong Lebesgue conclusions follow from the uniform estimates on the
finite-volume cylinder.  Equation \eqref{eq:V16-weak-Einstein-tail} is
Holder's inequality applied to \eqref{eq:V16-Einstein-tail}.
\end{proof}

\begin{firewall}[Two derivatives are an acquired row]
\label{fire:V16-C2-not-automatic}
Absolute summability of the values of the scheduler details does not imply
\eqref{eq:V16-C2-budget}.  A sequence can have summable amplitudes and
unbounded second derivatives.  The $C^{2,\alpha}$ detail budget must come from
the microscopic action, a regulator consistency estimate, or a compact
higher-regularity screen.  Kinematic noncollapse alone supplies none of these.
\end{firewall}

% -----------------------------------------------------------------------------
\subsection{Chronological connection incidence with the scheduler lift}
\label{subsec:V16-connection-incidence}
% -----------------------------------------------------------------------------

The scheduler determines a Levi--Civita connection, whereas the renewal
history supplies a chronological Cartan connection through its own face
writer.  These connections live in different constructions until an incidence
row identifies them.  Fix the coframe gauge and transport both connections to
one trivialized $\mathfrak{so}(1,3)$ bundle over $K$.  Let
$\omega_n^{\rm chr}$ be the chronological metric-compatible connection and
put
\begin{equation}
 a_n:=\omega_n^{\rm chr}-\omega_n^{\rm LC},
 \qquad
 \kappa_n:=\|a_n\|_{H^1(K)}.
 \label{eq:V16-connection-defect}
\end{equation}
Because the difference of two connections is a tensor, this definition is
independent of the chosen coordinate chart after the common Lorentz gauge has
been fixed.

\begin{theorem}[Connection incidence controls chronological curvature]
\label{thm:V16-connection-incidence}
In four spacetime dimensions, suppose
\begin{equation}
 \kappa_n\longrightarrow0.
 \label{eq:V16-connection-incidence-zero}
\end{equation}
Under Theorem~\ref{thm:V16-curvature-lift}, there is a constant $C_K$ such
that
\begin{align}
 \|\omega_n^{\rm chr}-\omega_\infty^{\rm LC}\|_{L^4(K)}
 &\le
 C_K\bigl(\kappa_n+\rho_n^{(2,\alpha)}\bigr),
 \label{eq:V16-chronological-connection-tail}\\
 \|F(\omega_n^{\rm chr})-F_\infty^{\rm LC}\|_{L^2(K)}
 &\le
 C_K\bigl(
 \kappa_n+\kappa_n^2+\rho_n^{(2,\alpha)}
 \bigr),
 \label{eq:V16-chronological-curvature-tail}\\
 \|T(\mathsf e_n,\omega_n^{\rm chr})\|_{L^2(K)}
 &\le C_K\kappa_n.
 \label{eq:V16-chronological-torsion-tail}
\end{align}
Consequently the actual chronological coframe and curvature satisfy the
strong-coframe/weak-$L^p$ curvature hypothesis of the imported localized
Palatini theorem with $p=2$, and in fact the chronological curvature
converges strongly.

The positive connection-incidence number
\begin{equation}
 \mathbb I_n^\omega
 :=
 \langle a_n,a_n\rangle_{H^1(K)}
 =\kappa_n^2
 \ge0
 \label{eq:V16-connection-Gram}
\end{equation}
vanishes exactly when the two connections coincide.  If
$\limsup_n\kappa_n>0$, a subsequence and the normalized tensors
$a_{n_j}/\kappa_{n_j}$ give unit $H^1$ connection-mismatch witnesses.
\end{theorem}

\begin{proof}
On a bounded Lipschitz domain in dimension four, Sobolev embedding gives
\[
 \|a_n\|_{L^4}\le C_K\|a_n\|_{H^1}=C_K\kappa_n.
\]
Together with \eqref{eq:V16-connection-tail}, this proves
\eqref{eq:V16-chronological-connection-tail}.  The exact curvature-difference
identity is
\begin{equation}
 F(\omega_n^{\rm chr})-F(\omega_n^{\rm LC})
 =
 D_{\omega_n^{\rm LC}}a_n+a_n\wedge a_n.
 \label{eq:V16-curvature-difference}
\end{equation}
The $C^{1,\alpha}$ bound on $\omega_n^{\rm LC}$ and the Sobolev embedding
give
\[
 \|D_{\omega_n^{\rm LC}}a_n\|_{L^2}\le C_K\kappa_n,
 \qquad
 \|a_n\wedge a_n\|_{L^2}
 \le C_K\|a_n\|_{L^4}^2
 \le C_K\kappa_n^2.
\]
Add the Levi--Civita curvature tail
\eqref{eq:V16-curvature-tail} to obtain
\eqref{eq:V16-chronological-curvature-tail}.  Since the Levi--Civita torsion
is zero,
\[
 T(\mathsf e_n,\omega_n^{\rm chr})
 =a_n\wedge\mathsf e_n.
\]
The uniform $L^\infty$ coframe bound proves
\eqref{eq:V16-chronological-torsion-tail}.  The zero and normalized-witness
claims follow directly from the Hilbert norm.
\end{proof}

\begin{assessment}[Relation to connection stationarity]
\label{ass:V16-connection-stationarity}
The imported Cartan theorem proves that exact spinless connection stationarity
and a nondegenerate coframe force zero torsion, hence the Levi--Civita
connection.  Equation \eqref{eq:V16-connection-Gram} is the cutoff transport
version needed here.  A small algebraic torsion residual controls the
connection difference in $L^2$ under the imported Cartan floor, but it does
not by itself control its first derivative.  The $H^1$ incidence in
\eqref{eq:V16-connection-incidence-zero}, or an equivalent curvature
compactness estimate, is therefore an additional acquisition rather than a
consequence of scheduler noncollapse.

If the chronological connection is exactly metric compatible and torsion free
at every cutoff, uniqueness makes $a_n=0$, so the entire incidence row passes
without a separate estimate.
\end{assessment}

% -----------------------------------------------------------------------------
\subsection{One determining core for action, Palatini, and represented stress}
\label{subsec:V16-one-core}
% -----------------------------------------------------------------------------

Fix $s>2$ and let
\begin{equation}
 \mathscr T_K
 :=
 H_0^s(K;\operatorname{Sym}^2 TK)
 \label{eq:V16-test-space}
\end{equation}
be a Hilbert space of compactly supported symmetric inverse-metric tests.
The inclusion $\mathscr T_K\hookrightarrow L^1(K)$ is continuous.  Every
continuous weak metric variation is therefore an element of
$\mathscr T_K^*$.  For $L\in\mathscr T_K^*$, write
$\mathcal R_KL\in\mathscr T_K$ for its Riesz representative:
\begin{equation}
 L(k)=\langle\mathcal R_KL,k\rangle_{\mathscr T_K}.
 \label{eq:V16-Riesz}
\end{equation}

At cutoff $n$, transport all the following functionals to this same test
space:
\begin{enumerate}[label=\textup{(\roman*)},leftmargin=2.8em]
\item $\mathcal P_n$, the chronological, boundary-assembled
      Palatini--Cartan gravitational metric variation;
\item $\mathcal V_n$, the Standard Model, gauge, potential, counterterm,
      improvement, contact, and boundary metric variation of the same common
      action;
\item $\mathcal A_n:=\mathcal P_n+\mathcal V_n$, the complete common-action
      metric first variation;
\item the scheduler geometric functional
      \begin{equation}
       \mathcal G_n(k)
       :=
       \chi\int_K\langle\mathfrak G_n,k\rangle;
       \label{eq:V16-geometric-functional}
      \end{equation}
\item the represented stress functional
      \begin{equation}
       \mathcal S_n(k)
       :=
       -\frac12\langle\mathbf T_n^{\rm OS},k\rangle,
       \label{eq:V16-stress-functional}
      \end{equation}
      where $\mathbf T_n^{\rm OS}$ is the densitized stress distribution
      reconstructed on the common OS graph core.
\end{enumerate}
The finite stress/current criterion imported from the monograph establishes
Ward relations for $\mathbf T_n^{\rm OS}$.  It does not identify
\eqref{eq:V16-stress-functional} with $\mathcal V_n$; that is a separate
same-action incidence statement.

Define four Riesz vectors
\begin{align}
 r_n^{\rm act}
 &:=
 \mathcal R_K\mathcal A_n,
 \label{eq:V16-action-vector}\\
 r_n^{\rm Pal}
 &:=
 \mathcal R_K(\mathcal G_n-\mathcal P_n),
 \label{eq:V16-Palatini-incidence-vector}\\
 r_n^{\rm str}
 &:=
 \mathcal R_K(\mathcal S_n-\mathcal V_n),
 \label{eq:V16-stress-incidence-vector}\\
 r_n^{E}
 &:=
 \mathcal R_K(\mathcal G_n+\mathcal S_n).
 \label{eq:V16-Einstein-vector}
\end{align}

\begin{lemma}[Exact three-defect identity]
\label{lem:V16-three-defect}
On $\mathscr T_K$,
\begin{equation}
 \boxed{
 r_n^E
 =
 r_n^{\rm act}
 +r_n^{\rm Pal}
 +r_n^{\rm str}.}
 \label{eq:V16-three-defect}
\end{equation}
Consequently the Einstein residual can vanish only after action stationarity,
Palatini incidence, and stress incidence have been assembled with their
displayed signs.
\end{lemma}

\begin{proof}
Using $\mathcal A_n=\mathcal P_n+\mathcal V_n$,
\[
 \mathcal A_n+(\mathcal G_n-\mathcal P_n)
             +(\mathcal S_n-\mathcal V_n)
 =\mathcal G_n+\mathcal S_n.
\]
The Riesz map is linear, which gives \eqref{eq:V16-three-defect}.
\end{proof}

Let $\mathscr L_K\ge0$ be a self-adjoint compact-resolvent determining
operator on $\mathscr T_K$.  Choose $\Lambda_m\uparrow\infty$ and set
\begin{equation}
 P_m:=\mathbf1_{[0,\Lambda_m]}(\mathscr L_K),
 \qquad P_m\uparrow I.
 \label{eq:V16-test-screens}
\end{equation}
For $u\in\mathscr T_K$, use
$(u\otimes u)v=\langle u,v\rangle u$.  The finite Palatini and stress
incidence Grams are
\begin{align}
 \mathbb I_{n,m}^{\rm Pal}
 &:=
 (P_mr_n^{\rm Pal})\otimes(P_mr_n^{\rm Pal})\succeq0,
 \label{eq:V16-Palatini-Gram}\\
 \mathbb I_{n,m}^{\rm str}
 &:=
 (P_mr_n^{\rm str})\otimes(P_mr_n^{\rm str})\succeq0.
 \label{eq:V16-stress-Gram}
\end{align}
Their vanishing is respectively equivalent to
$\mathcal G_n=\mathcal P_n$ and
$\mathcal S_n=\mathcal V_n$ on $P_m\mathscr T_K$.

\begin{theorem}[Quantitative same-core Einstein closure]
\label{thm:V16-Einstein-closure}
For $q\in\{\mathrm{act},\mathrm{Pal},\mathrm{str}\}$, suppose there are
$\sigma_q>0$ and $M_q<\infty$ such that
\begin{equation}
 \sup_n
 \|(I+\mathscr L_K)^{\sigma_q/2}r_n^q\|_{\mathscr T_K}
 \le M_q.
 \label{eq:V16-three-regularity-bounds}
\end{equation}
Put
\begin{equation}
 \varepsilon_{n,m}^{\rm act}
 :=\|P_mr_n^{\rm act}\|_{\mathscr T_K}
 \label{eq:V16-action-screen-residual}
\end{equation}
and define
\begin{align}
 B_{n,m}:={}&
 \left[
  (\varepsilon_{n,m}^{\rm act})^2
  +(1+\Lambda_m)^{-\sigma_{\rm act}}M_{\rm act}^2
 \right]^{1/2}
 \nonumber\\
 &+
 \left[
  \|\mathbb I_{n,m}^{\rm Pal}\|_{\rm op}
  +(1+\Lambda_m)^{-\sigma_{\rm Pal}}M_{\rm Pal}^2
 \right]^{1/2}
 \nonumber\\
 &+
 \left[
  \|\mathbb I_{n,m}^{\rm str}\|_{\rm op}
  +(1+\Lambda_m)^{-\sigma_{\rm str}}M_{\rm str}^2
 \right]^{1/2}.
 \label{eq:V16-master-bound}
\end{align}
Then the complete finite-cutoff Einstein residual satisfies
\begin{equation}
 \boxed{
 \|r_n^E\|_{\mathscr T_K}\le B_{n,m}.}
 \label{eq:V16-finite-Einstein-bound}
\end{equation}

Assume in addition the $C^{2,\alpha}$ scheduler hypotheses of
Theorem~\ref{thm:V16-curvature-lift} and a summable represented-stress
transport bound
\begin{equation}
 \|\mathcal S_{j+1}-\mathcal S_j\|_{\mathscr T_K^*}
 \le d_j^{\rm str},
 \qquad
 \sum_{j\ge0}d_j^{\rm str}<\infty.
 \label{eq:V16-stress-transport}
\end{equation}
Then $\mathcal S_n$ converges in $\mathscr T_K^*$ to a functional
$\mathcal S_\infty$, and
\begin{equation}
 \boxed{
 \|\mathcal R_K(\mathcal G_\infty+\mathcal S_\infty)\|_{\mathscr T_K}
 \le
 B_{n,m}
 +C_{\mathscr T,K}|\chi|C_{K,\Lambda}
  \rho_n^{(2,\alpha)}
 +\sum_{j\ge n}d_j^{\rm str}.}
 \label{eq:V16-limit-Einstein-bound}
\end{equation}
Here $C_{\mathscr T,K}$ is the norm of
$\mathscr T_K\hookrightarrow L^1(K)$.

If there is a cofinal choice $m=m(n)\to\infty$ for which
$B_{n,m(n)}\to0$, then
\begin{equation}
 \mathcal G_\infty+\mathcal S_\infty=0
 \quad\hbox{in }\mathscr T_K^*.
 \label{eq:V16-limit-Einstein-zero}
\end{equation}
If the limiting represented stress has the tensor-distribution
normalization
\begin{equation}
 \langle\mathbf T_\infty^{\rm OS},k\rangle
 =
 \int_K\sqrt{-\det\mathbf g_\infty}\,
 T_{\mu\nu}^{\rm ren}k^{\mu\nu},
 \label{eq:V16-stress-normalization}
\end{equation}
then
\begin{equation}
 \boxed{
 2\chi(G_{\mu\nu}(\mathbf g_\infty)
       +\Lambda\mathbf g_{\infty,\mu\nu})
 =
 T_{\mu\nu}^{\rm ren}}
 \label{eq:V16-distributional-Einstein}
\end{equation}
holds distributionally on $K$.
\end{theorem}

\begin{proof}
For $u$ in the domain of
$(I+\mathscr L_K)^{\sigma/2}$, the spectral theorem gives
\begin{equation}
 \|(I-P_m)u\|^2
 \le
 (1+\Lambda_m)^{-\sigma}
 \|(I+\mathscr L_K)^{\sigma/2}u\|^2.
 \label{eq:V16-spectral-inequality}
\end{equation}
The two spectral pieces are orthogonal.  Hence
\[
 \|r_n^{\rm act}\|
 \le
 \left[
  (\varepsilon_{n,m}^{\rm act})^2
  +(1+\Lambda_m)^{-\sigma_{\rm act}}M_{\rm act}^2
 \right]^{1/2}.
\]
For a rank-one positive operator,
$\|u\otimes u\|_{\rm op}=\|u\|^2$.  Applying
\eqref{eq:V16-spectral-inequality} to the Palatini and stress incidence
vectors gives the other two summands in
\eqref{eq:V16-master-bound}.  The triangle inequality in the exact identity
\eqref{eq:V16-three-defect} proves
\eqref{eq:V16-finite-Einstein-bound}.  This is the imported assembled
spectral-tail mechanism applied to the two incidence rows which that theorem
does not identify.

Summability in \eqref{eq:V16-stress-transport} makes
$(\mathcal S_n)$ Cauchy in the complete dual space and gives
\[
 \|\mathcal S_\infty-\mathcal S_n\|_{\mathscr T_K^*}
 \le\sum_{j\ge n}d_j^{\rm str}.
\]
For a unit test $k\in\mathscr T_K$, the curvature estimate gives
\[
 |(\mathcal G_\infty-\mathcal G_n)(k)|
 \le
 C_{\mathscr T,K}|\chi|C_{K,\Lambda}
 \rho_n^{(2,\alpha)}.
\]
The Riesz map is an isometry.  Add these two bounds to
\eqref{eq:V16-finite-Einstein-bound} to obtain
\eqref{eq:V16-limit-Einstein-bound}.  All three terms on its right tend to
zero on the stated cofinal choice, proving
\eqref{eq:V16-limit-Einstein-zero}.

Finally, insert \eqref{eq:V16-geometric-functional},
\eqref{eq:V16-stress-functional}, and
\eqref{eq:V16-stress-normalization} into
\eqref{eq:V16-limit-Einstein-zero}.  The determinant density is everywhere
nonzero by Theorem~\ref{thm:V16-curvature-lift}.  Since
$C_c^\infty(K;\operatorname{Sym}^2TK)$ is contained in the test core, the
resulting equality is precisely
\eqref{eq:V16-distributional-Einstein} in the sense of distributions.
\end{proof}

\begin{firewall}[A scheduler root is only one variation block]
\label{fire:V16-root-not-full-residual}
The V15 inverse theorem can set the transported scheduler component of the
common-action first variation to zero.  It does not imply
$\varepsilon_{n,m}^{\rm act}=0$ on an enlarged determining metric--matter
test screen.  Constraint, connection, matter, counterterm, contact, and
boundary variations must either be included in the invertible quotient map or
be tested separately before the first summand of
\eqref{eq:V16-master-bound} is declared small.
\end{firewall}

% -----------------------------------------------------------------------------
\subsection{Reducing Palatini incidence to face and connection incidence}
\label{subsec:V16-Palatini-reduction}
% -----------------------------------------------------------------------------

Let $\mathcal P^{\rm can}(\mathsf e,\omega)$ denote the symmetric coframe
variation of the canonically normalized Palatini--Holst--volume functional on
$K$, including the same boundary convention used in $\mathcal P_n$.  Normalize
its Palatini and volume coefficients so that
\begin{equation}
 \mathcal P^{\rm can}(\mathsf e_n,\omega_n^{\rm LC})
 =\mathcal G_n.
 \label{eq:V16-canonical-Palatini-Einstein}
\end{equation}
This equality is the ordinary torsion-free Palatini identity; on symmetric
metric variations the Levi--Civita Holst term vanishes by the first Bianchi
identity.

Define the face-writer consistency error
\begin{equation}
 \eta_n^{\rm face}
 :=
 \|\mathcal P_n
 -\mathcal P^{\rm can}(\mathsf e_n,\omega_n^{\rm chr})
 \|_{\mathscr T_K^*}.
 \label{eq:V16-face-consistency}
\end{equation}

\begin{proposition}[Palatini incidence bound from two physical rows]
\label{prop:V16-Palatini-incidence-bound}
Under Theorems~\ref{thm:V16-curvature-lift} and
\ref{thm:V16-connection-incidence},
\begin{equation}
 \boxed{
 \|r_n^{\rm Pal}\|_{\mathscr T_K}
 \le
 \eta_n^{\rm face}
 +C_K(\kappa_n+\kappa_n^2).}
 \label{eq:V16-Palatini-incidence-bound}
\end{equation}
Hence the Palatini-incidence summand of
\eqref{eq:V16-master-bound} tends to zero if the chronological connection
incidence and the common-face consistency error tend to zero.
\end{proposition}

\begin{proof}
The coframe variation of the Palatini and Holst terms is linear in curvature
and polynomial of fixed degree in $\mathsf e_n$ and its inverse.  The
cosmological term is independent of the connection.  The V14 hard floor and
Theorem~\ref{thm:V16-curvature-lift} uniformly bound every coframe
coefficient.  Therefore
\[
 \|\mathcal P^{\rm can}(\mathsf e_n,\omega_n^{\rm chr})
   -\mathcal P^{\rm can}(\mathsf e_n,\omega_n^{\rm LC})
 \|_{\mathscr T_K^*}
 \le
 C_K\|F(\omega_n^{\rm chr})-F(\omega_n^{\rm LC})\|_{L^2}.
\]
Equation \eqref{eq:V16-curvature-difference} bounds the right side by
$C_K(\kappa_n+\kappa_n^2)$.  Add and subtract the canonical chronological
functional, use \eqref{eq:V16-canonical-Palatini-Einstein}, and apply the
triangle inequality.  The Riesz map is an isometry, proving
\eqref{eq:V16-Palatini-incidence-bound}.
\end{proof}

\begin{firewall}[The face error contains physical normalization]
\label{fire:V16-face-error}
The number $\eta_n^{\rm face}$ is not a generic finite-difference error.  It
tests that the chronological face score, after memory shorting, determinant
orientation, counterterms, boundary assembly, and physical clock
normalization, is the same canonical Palatini functional used in
\eqref{eq:V16-canonical-Palatini-Einstein}.  A small connection mismatch
cannot repair a wrong coefficient, a missing source, or an unassembled
boundary term.
\end{firewall}

% -----------------------------------------------------------------------------
\subsection{Cofinal synchronization or a supported obstruction}
\label{subsec:V16-diagonal}
% -----------------------------------------------------------------------------

For the three residual vectors define
\begin{align}
 c_{n,m}
 &:=
 \|P_mr_n^{\rm act}\|
 +\|P_mr_n^{\rm Pal}\|
 +\|P_mr_n^{\rm str}\|,
 \label{eq:V16-screen-sum}\\
 q_m
 &:=
 \sup_n\left(
 \|(I-P_m)r_n^{\rm act}\|
 +\|(I-P_m)r_n^{\rm Pal}\|
 +\|(I-P_m)r_n^{\rm str}\|
 \right).
 \label{eq:V16-uniform-tail}
\end{align}
The regularity hypotheses
\eqref{eq:V16-three-regularity-bounds} imply $q_m\to0$.

\begin{theorem}[Diagonal same-core pass or physical witness]
\label{thm:V16-diagonal-alternative}
Suppose
\begin{equation}
 q_m\longrightarrow0,
 \qquad
 \liminf_{n\to\infty}c_{n,m}=0
 \quad\hbox{for every fixed }m.
 \label{eq:V16-diagonal-hypotheses}
\end{equation}
Then there is a strictly increasing sequence $n_m$ such that
\begin{equation}
 \|r_{n_m}^E\|
 \le 2^{-m}+q_m\longrightarrow0.
 \label{eq:V16-diagonal-bound}
\end{equation}
If the scheduler and represented-stress limits of
Theorem~\ref{thm:V16-Einstein-closure} exist along the full sequence, this
subsequence proves the same distributional Einstein identity
\eqref{eq:V16-distributional-Einstein}.

If \eqref{eq:V16-diagonal-hypotheses} fails, one of the following explicit
obstructions occurs:
\begin{enumerate}[label=\textup{(\roman*)},leftmargin=2.8em]
\item there are a fixed $m_0$, $\epsilon_*>0$, and arbitrarily large $n$
      for which one of the three vectors has a finite-screen component of
      norm at least $\epsilon_*/3$; normalizing that component gives a unit
      test in $P_{m_0}\mathscr T_K$ detecting the corresponding action,
      Palatini, or stress defect;
\item there are $m_j\uparrow\infty$, cutoffs $n_j$, one fixed defect type
      after taking a subsequence, and unit tests
      $v_j\in(I-P_{m_j})\mathscr T_K$ such that
      \begin{equation}
       |\langle r_{n_j}^q,v_j\rangle|
       \ge\epsilon_*>0.
       \label{eq:V16-escaping-witness}
      \end{equation}
      The $v_j$ converge weakly to zero, so this is a genuine escaping
      determining variation.
\end{enumerate}
\end{theorem}

\begin{proof}
For each $m$, the second condition in
\eqref{eq:V16-diagonal-hypotheses} permits an $n_m>n_{m-1}$ with
$c_{n_m,m}\le2^{-m}$.  By \eqref{eq:V16-three-defect},
\[
 \|r_{n_m}^E\|
 \le c_{n_m,m}
 +\sum_q\|(I-P_m)r_{n_m}^q\|
 \le2^{-m}+q_m.
\]
This proves \eqref{eq:V16-diagonal-bound}.  Since both limiting functionals
are fixed by the full-sequence limits, convergence of their sum along one
cofinal subsequence forces their limiting sum to vanish.

If the fixed-screen condition fails, then for some $m_0$ and
$\epsilon_*>0$, $c_{n,m_0}\ge\epsilon_*$ eventually.  At least one of its
three nonnegative summands is at least $\epsilon_*/3$ along a subsequence.
Normalize that projected vector.

If $q_m$ does not tend to zero, choose $\epsilon_*>0$, $m_j\uparrow\infty$,
and $n_j$ for which the sum in \eqref{eq:V16-uniform-tail} is at least
$3\epsilon_*$.  One of the three tail components is at least
$\epsilon_*$ along a subsequence of a fixed type $q$.  Normalize it to obtain
$v_j$ and \eqref{eq:V16-escaping-witness}.  For a fixed $w\in\mathscr T_K$,
choose $M$ with $\|(I-P_M)w\|$ small.  When $m_j\ge M$,
$\langle v_j,P_Mw\rangle=0$, hence
$|\langle v_j,w\rangle|\le\|(I-P_M)w\|$.  Thus $v_j\rightharpoonup0$.
\end{proof}

% -----------------------------------------------------------------------------
\subsection{Executable V16 acquisition card}
\label{subsec:V16-card}
% -----------------------------------------------------------------------------

\begin{definition}[Curvature--stress closure card]
\label{def:V16-card}
On each selected compact cylinder and cutoff, record:
\begin{enumerate}[label=\textup{(C\arabic*)},leftmargin=3em]
\item the V15 transported physical action root and quotient Hessian floor;
\item the scheduler details through two derivatives, their
      $\rho_n^{(2,\alpha)}$ tail, and the V14 determinant and inverse floors;
\item the chronological connection, the scheduler Levi--Civita connection,
      and the incidence number $\mathbb I_n^\omega$;
\item the completely memory-shorted and boundary-assembled chronological
      Palatini writer, its physical coefficients, and
      $\eta_n^{\rm face}$;
\item one common transported test inner product, determining operator
      $\mathscr L_K$, and finite screens $P_m$;
\item the complete action vector $r_n^{\rm act}$, including every metric,
      connection, constraint, matter, counterterm, contact, and boundary
      variation assigned to the screen;
\item the represented OS stress and the common-action matter metric
      variation on the same continued test core, together with
      $\mathbb I_{n,m}^{\rm str}$;
\item the three high-screen tail norms, the adjacent represented-stress
      transport debit, and one held-out variation outside the fitting screen.
\end{enumerate}
All source, boundary, and nuisance lists are frozen before the terminal
residuals are evaluated.
\end{definition}

\begin{theorem}[The V16 card terminates with a bound or a witness]
\label{thm:V16-card-alternative}
A populated finite V16 card has exactly the following certification output.
If its finite screen rows and declared tail margins pass, it returns the
number $B_{n,m}$ in \eqref{eq:V16-master-bound} and the verified estimate
\eqref{eq:V16-finite-Einstein-bound}.  A cofinal sequence of passing cards
which satisfies the scheduler and stress transport tails returns the
distributional identity \eqref{eq:V16-distributional-Einstein}.

Otherwise the card returns at least one concrete obstruction:
\begin{enumerate}[label=\textup{(\alph*)},leftmargin=2.8em]
\item a V15 gauge leakage, right soft mode, target cokernel, nonlinear-radius,
      or adjacent-action witness;
\item a two-derivative scheduler-tail or determinant-floor failure;
\item the normalized chronological connection mismatch from
      \eqref{eq:V16-connection-Gram};
\item the signed face-writer difference in
      \eqref{eq:V16-face-consistency};
\item a nonzero complete-action, Palatini-incidence, or stress-incidence
      Riesz vector on the finite determining screen;
\item a held-out variation on which a fitted zero fails;
\item the weakly escaping variation
      \eqref{eq:V16-escaping-witness};
\item a nonsummable represented-stress transport debit.
\end{enumerate}
No zero vacuum expectation, Fisher block, Ward relation, or scheduler root is
substituted for any missing row.
\end{theorem}

\begin{proof}
Every cutoff screen is finite dimensional.  The V15 singular-value and Gram
tests decide item~\textup{(a)}.  Direct norms decide the scheduler,
connection, face, Riesz, held-out, and transport rows.  The positive rank-one
operators \eqref{eq:V16-Palatini-Gram,eq:V16-stress-Gram} vanish exactly when
their incidence vectors vanish on the screen.  Passing rows give
Theorem~\ref{thm:V16-Einstein-closure}; failure of uniform exhaustion gives
Theorem~\ref{thm:V16-diagonal-alternative}.  These tests exhaust the entries
of Definition~\ref{def:V16-card}.
\end{proof}

\begin{firewall}[The OS-to-Lorentz test transport is physical data]
\label{fire:V16-Wick-transport}
The represented stress in \eqref{eq:V16-stress-functional} begins on an
OS-reconstructed Euclidean graph core, while
\eqref{eq:V16-geometric-functional} is written on the Lorentzian scheduler
cylinder.  Placing them in one $\mathscr T_K$ presupposes the determinant-line,
time-orientation, analytic-continuation, domain, and source transports already
isolated in the monograph.  The stress-incidence Gram tests that transported
object.  Formal replacement of Euclidean time by imaginary Lorentz time does
not populate this row.
\end{firewall}

% -----------------------------------------------------------------------------
\subsection{V16 theorem ledger and next physical acquisition}
\label{subsec:V16-ledger}
% -----------------------------------------------------------------------------

\paragraph{New conclusions proved in V16.}
Theorem~\ref{thm:V16-curvature-lift} upgrades a summable
two-derivative logarithmic scheduler to quantitative convergence of the
metric, coframe, Levi--Civita connection, curvature, and densitized Einstein
writer.  Theorem~\ref{thm:V16-connection-incidence} converts one $H^1$
chronological-connection incidence row into strong $L^2$ curvature
convergence.  Lemma~\ref{lem:V16-three-defect} gives the exact signed
decomposition of the Einstein residual into action, Palatini, and stress
defects.  Theorem~\ref{thm:V16-Einstein-closure} supplies a computable
spectral bound for those three defects and proves the distributional equation
on their cofinal pass branch.  Proposition
\ref{prop:V16-Palatini-incidence-bound} reduces the Palatini incidence to a
face-writer error and the connection incidence.  Theorem
\ref{thm:V16-diagonal-alternative} synchronizes increasing finite screens or
returns a supported or weakly escaping variation.

\paragraph{Imported results not repeated.}
The analytic ADM chart and its hard floors, the localized Palatini product
limit, connection-stationarity torsion theorem, assembled Palatini spectral
tail, finite and cofinal OS stress/current criteria, common-action exactness,
boundary response, and conditional Einstein handoff are imported.  V16 adds
the scheduler curvature estimate, chronological/scheduler connection
incidence, represented-stress/common-action incidence, and their exact
same-core sum.

\paragraph{Authoritative physical status.}
The attached files do not provide, for one selected interacting RPESM
history, the $C^{2,\alpha}$ scheduler detail bounds, the chronological
connection matrices in the transported scheduler gauge, the values of
$\eta_n^{\rm face}$, or the same-core Riesz vectors
\[
 r_n^{\rm act},\qquad r_n^{\rm Pal},\qquad r_n^{\rm str}.
\]
They also do not print a summable OS-to-Lorentz stress transport debit.
Consequently the hypotheses of
Theorem~\ref{thm:V16-Einstein-closure} have not been physically populated,
and the full Standard Model--Einstein continuum is not claimed.

\paragraph{Next acquisition.}
The next iteration should go upstream to the actual face data and populate one
minimal cylinder card in this order:
\begin{enumerate}[label=\textup{(\arabic*)},leftmargin=2.8em]
\item choose one compact cylinder, one physical graph norm, and one increasing
      eigenbasis of $\mathscr L_K$;
\item differentiate the selected V15 action roots through two spacetime
      derivatives and measure $\rho_n^{(2,\alpha)}$;
\item transport the occurring chronological connection into the same coframe
      gauge and evaluate $\mathbb I_n^\omega$;
\item evaluate the memory-shorted Palatini face writer on the determining
      metric tests and compute $\eta_n^{\rm face}$;
\item continue the represented OS stress to those same tests and directly
      compare it with the common-action matter metric derivative;
\item report $B_{n,m}$ at adjacent cutoffs, then enlarge $m$ only after the
      three high-screen tails have been measured.
\end{enumerate}
On the pass branch these values feed directly into
\eqref{eq:V16-limit-Einstein-bound}.  On failure the first nonzero vector is
retained as the physical obstruction and the programme moves upstream to its
source, gauge, memory, boundary, or continuation row.

\begin{firewall}[V16 is a closure theorem, not a populated continuum]
\label{fire:V16-final-scope}
The estimates in this version close the metric equation once a single
physical sequence supplies the displayed data.  They do not manufacture the
two-derivative action bound, chronological face normalization, OS-to-Lorentz
stress transport, remaining matter Euler variations, Ward conservation,
counterterm limits, or boundary exhaustion.  Those are occurrence and
transport questions for the attached RPESM programme.
\end{firewall}

% =============================================================================
% END APPENDED VERSION V16 MATERIAL
% =============================================================================



% =============================================================================
% BEGIN APPENDED VERSION V17 MATERIAL
% =============================================================================

\clearpage
\section{Version V17: graph-norm action roots and metric-secanted stress incidence}
\label{sec:V17-graph-roots}

Version V16 isolated a quantitative closure theorem, but two of its hypotheses
were still presented as external data: the action-selected scheduler had to
converge through two spacetime derivatives, and the represented OS stress had
to agree with the matter metric derivative of the same common action.  The
latest attached source-metric programme supplies relation-covariant shrinking
command secants and the required reciprocal-radius cutoff scaling, but only on
fixed finite coefficient screens.  It does not turn those secants into a
spacetime graph-norm inverse or a stress/action identity.

This version builds that bridge.  A uniformly invertible gauge-reduced action
linearization between Sobolev graph spaces upgrades the V14--V15 root
transport to $H^s$.  For $s>4+\alpha$ in four dimensions, Sobolev--Morrey
embedding supplies exactly the $C^{2,\alpha}$ scheduler tail used in V16.
Finite spectral shells and relation-covariant command secants give testable
sufficient conditions for the graph inverse.  A second central secant,
formed from a physical metric command on the common unnormalized cylinder,
bounds the represented-stress/common-action incidence on a source-complete
finite atlas.

\paragraph{Inherited token repair.}
The V1 display defining its mixed command secant contained one form-feed
control byte where the backslash of \verb|\frac| belongs.  V17 replaces that
single byte by the two characters \verb|\f|, so the source now contains the
intended \verb|\frac|.  No theorem text or mathematical content is otherwise
changed.  The append-only audit for this version compares against V16 after
reversing precisely this one-token repair.

% -----------------------------------------------------------------------------
\subsection{Gauge-reduced action equations in a spacetime graph norm}
\label{subsec:V17-graph-action}
% -----------------------------------------------------------------------------

Let $K$ be the four-dimensional compact coordinate cylinder of V16, let
$0<\alpha<1$, and choose
\begin{equation}
 s>4+\alpha.
 \label{eq:V17-Sobolev-index}
\end{equation}
Fix boundary or Cauchy data and all diffeomorphism, local Lorentz, internal
gauge, and relation quotients.  Let
\begin{equation}
 \mathscr X_K:=H_B^s(K;\mathcal Z_{\rm tot}),
 \qquad
 \mathscr Y_K:=H_B^{s-\nu}(K;\mathcal W),
 \qquad 0\le\nu\le s,
 \label{eq:V17-graph-spaces}
\end{equation}
be the resulting real Hilbert graph spaces.  The complete field fibre
$\mathcal Z_{\rm tot}$ has a bounded projection
$\pi_{\rm sch}:\mathcal Z_{\rm tot}\to\mathcal Z$ onto the established
ten-dimensional scheduler fibre; $\mathcal W$ is the complete determining
response fibre.  We take $\|\pi_{\rm sch}\|\le1$ by rescaling the product
norm.  The subscript $B$ records the fixed boundary or Cauchy conditions.
Sobolev--Morrey embedding
gives a constant $C_{\rm emb}$ such that
\begin{equation}
 \|u\|_{C^{2,\alpha}(K)}
 \le C_{\rm emb}\|u\|_{\mathscr X_K}.
 \label{eq:V17-Sobolev-embedding}
\end{equation}

Choose a reference field $z_{\rm ref}$ and the affine ball
\begin{equation}
 \mathscr U_R
 :=
 z_{\rm ref}+\{u\in\mathscr X_K:\|u\|_{\mathscr X_K}\le R\}.
 \label{eq:V17-affine-ball}
\end{equation}
Let
\[
 \mathcal F_n:\mathscr U_R\longrightarrow\mathscr Y_K
\]
be the completely assembled, transported, gauge-reduced common-action first
variation at cutoff $n$.  Let
\begin{equation}
 H_*:\mathscr X_K\longrightarrow\mathscr Y_K
 \label{eq:V17-base-isomorphism}
\end{equation}
be a bounded isomorphism.  It is the graph-space realization of the signed
physical Hessian, rather than a positive Fisher surrogate.

\begin{theorem}[Graph-space action roots populate the V16 curvature tail]
\label{thm:V17-graph-root}
Suppose every $\mathcal F_n$ is continuously differentiable and, for one
$0\le\theta<1$,
\begin{align}
 \sup_n\sup_{z\in\mathscr U_R}
 \|I-H_*^{-1}D\mathcal F_n(z)\|_{\mathscr X_K\to\mathscr X_K}
 &\le\theta,
 \label{eq:V17-graph-linearization}\\
 \sup_n
 \|H_*^{-1}\mathcal F_n(z_{\rm ref})\|_{\mathscr X_K}
 &\le(1-\theta)R.
 \label{eq:V17-graph-center}
\end{align}
Define the adjacent preconditioned action defect
\begin{equation}
 \epsilon_n^{(s)}
 :=
 \sup_{z\in\mathscr U_R}
 \|H_*^{-1}
   (\mathcal F_{n+1}(z)-\mathcal F_n(z))
 \|_{\mathscr X_K}.
 \label{eq:V17-graph-adjacent-defect}
\end{equation}
If
\begin{equation}
 \sum_{n\ge0}\epsilon_n^{(s)}<\infty,
 \label{eq:V17-graph-summability}
\end{equation}
then:
\begin{enumerate}[label=\textup{(\roman*)},leftmargin=2.8em]
\item every cutoff has a unique root $z_n\in\mathscr U_R$;
\item the roots converge in $\mathscr X_K$ to $z_\infty$ and
      \begin{equation}
       \|z_\infty-z_n\|_{\mathscr X_K}
       \le
       \frac1{1-\theta}
       \sum_{j\ge n}\epsilon_j^{(s)};
       \label{eq:V17-graph-root-tail}
      \end{equation}
\item they converge in $C^{2,\alpha}(K)$ with
      \begin{equation}
       \|z_\infty-z_n\|_{C^{2,\alpha}}
       \le
       \frac{C_{\rm emb}}{1-\theta}
       \sum_{j\ge n}\epsilon_j^{(s)}.
       \label{eq:V17-C2-root-tail}
      \end{equation}
\end{enumerate}
If in addition
\begin{equation}
 \|\pi_{\rm sch}z_0\|_{C^{2,\alpha}}
 +\frac{C_{\rm emb}}{1-\theta}
  \sum_{j\ge0}\epsilon_j^{(s)}
 \le R_0,
 \label{eq:V17-noncollapse-budget}
\end{equation}
where $R_0$ is the V14 logarithmic noncollapse radius, then the projected
increments $\pi_{\rm sch}(z_{n+1}-z_n)$ form a scheduler series satisfying the V16
two-derivative budget.  Consequently all V16 Levi--Civita connection,
curvature, and densitized Einstein tails hold with
\begin{equation}
 \rho_n^{(2,\alpha)}
 \le
 \frac{C_{\rm emb}}{1-\theta}
 \sum_{j\ge n}\epsilon_j^{(s)}.
 \label{eq:V17-to-V16-tail}
\end{equation}
\end{theorem}

\begin{proof}
For $u$ in the closed radius-$R$ ball of $\mathscr X_K$, set
\[
 T_n(u)
 :=
 u-H_*^{-1}\mathcal F_n(z_{\rm ref}+u).
\]
The derivative estimate \eqref{eq:V17-graph-linearization} gives
\[
 \|T_n(u)-T_n(v)\|_{\mathscr X_K}
 \le\theta\|u-v\|_{\mathscr X_K}.
\]
Moreover,
\[
 \|T_n(u)\|
 \le\theta\|u\|
    +\|H_*^{-1}\mathcal F_n(z_{\rm ref})\|
 \le\theta R+(1-\theta)R=R.
\]
Thus $T_n$ is a contraction of a complete closed ball into itself.  Its
unique fixed point $u_n$ gives $z_n=z_{\rm ref}+u_n$ and
$\mathcal F_n(z_n)=0$.

For $z,w\in\mathscr U_R$, the segment joining them lies in the ball, and the
Banach-space fundamental theorem of calculus gives
\[
 H_*^{-1}(\mathcal F_n(z)-\mathcal F_n(w))
 =
 \int_0^1
 H_*^{-1}D\mathcal F_n(w+t(z-w))(z-w)\,\dd t.
\]
Every integrand differs from $z-w$ by norm at most
$\theta\|z-w\|$, so
\begin{equation}
 \|H_*^{-1}(\mathcal F_n(z)-\mathcal F_n(w))\|
 \ge(1-\theta)\|z-w\|.
 \label{eq:V17-graph-lower-bound}
\end{equation}
Take $z=z_{n+1}$ and $w=z_n$.  Since both are roots at their respective
cutoffs,
\[
 \|H_*^{-1}\mathcal F_n(z_{n+1})\|
 =
 \|H_*^{-1}
   (\mathcal F_n(z_{n+1})-\mathcal F_{n+1}(z_{n+1}))
 \|
 \le\epsilon_n^{(s)}.
\]
Equation \eqref{eq:V17-graph-lower-bound} yields
\begin{equation}
 \|z_{n+1}-z_n\|_{\mathscr X_K}
 \le\frac{\epsilon_n^{(s)}}{1-\theta}.
 \label{eq:V17-adjacent-root-bound}
\end{equation}
Summation proves the Cauchy property and
\eqref{eq:V17-graph-root-tail}.  Apply
\eqref{eq:V17-Sobolev-embedding} to obtain
\eqref{eq:V17-C2-root-tail}.

Set
\[
 a_{n+1}:=\|\pi_{\rm sch}(z_{n+1}-z_n)\|_{C^{2,\alpha}}.
\]
When $a_{n+1}>0$, put
$\xi_{n+1}=\pi_{\rm sch}(z_{n+1}-z_n)/a_{n+1}$; otherwise put
$\xi_{n+1}=0$.  Then
\[
 \pi_{\rm sch}z_n=\pi_{\rm sch}z_0+\sum_{j=1}^na_j\xi_j,
 \qquad
 \|\xi_j\|_{C^{2,\alpha}}\le1,
\]
and \eqref{eq:V17-adjacent-root-bound} gives
\[
 \sum_{j>n}a_j
 \le
 \frac{C_{\rm emb}}{1-\theta}
 \sum_{j\ge n}\epsilon_j^{(s)}.
\]
The budget \eqref{eq:V17-noncollapse-budget} therefore places every pointwise
scheduler value in the V14 ball and proves
\eqref{eq:V17-to-V16-tail}.  Apply the V16 curvature-lift theorem.
\end{proof}

\begin{corollary}[The limiting graph-space action equation]
\label{cor:V17-limit-action}
Under Theorem~\ref{thm:V17-graph-root}, the maps
$H_*^{-1}\mathcal F_n$ converge uniformly on $\mathscr U_R$ to a continuous
map $H_*^{-1}\mathcal F_\infty$, and
\begin{equation}
 \mathcal F_\infty(z_\infty)=0.
 \label{eq:V17-limit-action-zero}
\end{equation}
If the derivatives also converge uniformly, then
$D\mathcal F_\infty(z)$ retains the lower graph-space estimate
\begin{equation}
 \|H_*^{-1}D\mathcal F_\infty(z)u\|
 \ge(1-\theta)\|u\|.
 \label{eq:V17-limit-derivative-floor}
\end{equation}
\end{corollary}

\begin{proof}
The definition \eqref{eq:V17-graph-adjacent-defect} and summability give
\[
 \sup_{z\in\mathscr U_R}
 \|H_*^{-1}(\mathcal F_m(z)-\mathcal F_n(z))\|
 \le\sum_{j=n}^{m-1}\epsilon_j^{(s)}.
\]
Hence the preconditioned maps are uniformly Cauchy.  Since $z_n\to z_\infty$
and $\mathcal F_n(z_n)=0$, uniform convergence and the common Lipschitz bound
give \eqref{eq:V17-limit-action-zero}.  The derivative statement follows by
passing to the limit in the integral-free pointwise estimate implied by
\eqref{eq:V17-graph-linearization}.
\end{proof}

\begin{firewall}[The graph inverse is a dynamical certificate]
\label{fire:V17-graph-inverse}
A uniform singular-value floor on each finite coefficient matrix does not by
itself construct the isomorphism \eqref{eq:V17-base-isomorphism}.  The latter
also requires a common domain, compatible boundary or Cauchy data, control of
high-frequency leakage, and a uniform graph estimate.  On a Euclidean
gauge-fixed branch it may come from an elliptic boundary estimate.  On a
Lorentzian branch it must come from a well-posed causal energy estimate.
Theorem~\ref{thm:V17-graph-root} assumes neither positivity nor ellipticity;
it uses the actual bounded inverse in the norm appropriate to the selected
branch.
\end{firewall}

% -----------------------------------------------------------------------------
\subsection{A spectral shell criterion for the two-derivative budget}
\label{subsec:V17-spectral-shell}
% -----------------------------------------------------------------------------

The graph norm can be measured by finite spectral cards.  Let
$\mathscr L\ge0$ be a positive self-adjoint elliptic operator of order two on
the compact cylinder with the fixed boundary conditions.  Assume the
equivalent norm
\begin{equation}
 c_s\|u\|_{H^s}
 \le
 \|(I+\mathscr L)^{s/2}u\|_{L^2}
 \le C_s\|u\|_{H^s}.
 \label{eq:V17-spectral-norm-equivalence}
\end{equation}
For $\Lambda\ge0$, write
$P_\Lambda=\mathbf1_{[0,\Lambda]}(\mathscr L)$.

\begin{proposition}[Band and leakage decomposition of a scheduler shell]
\label{prop:V17-spectral-shell}
For every $u\in H^s$ and every $\Lambda\ge0$,
\begin{equation}
 \boxed{
 \|u\|_{C^{2,\alpha}}
 \le
 \frac{C_{\rm emb}}{c_s}
 \left[
  (1+\Lambda)^{s/2}\|P_\Lambda u\|_{L^2}
  +\|(I+\mathscr L)^{s/2}
       (I-P_\Lambda)u\|_{L^2}
 \right].}
 \label{eq:V17-band-leakage-bound}
\end{equation}
Consequently, for projected scheduler increments
$u_j=\pi_{\rm sch}(z_{j+1}-z_j)$ and chosen $\Lambda_j$, the sufficient finite-shell condition
\begin{equation}
 \sum_j
 \left[
  (1+\Lambda_j)^{s/2}\|P_{\Lambda_j}u_j\|_{L^2}
  +\|(I+\mathscr L)^{s/2}
       (I-P_{\Lambda_j})u_j\|_{L^2}
 \right]
 <\infty
 \label{eq:V17-weighted-shell-sum}
\end{equation}
implies the V16 $C^{2,\alpha}$ detail budget.  If every increment is
spectrally supported below $\Lambda_j$, the second term vanishes.
\end{proposition}

\begin{proof}
Spectral calculus gives
\[
 \|(I+\mathscr L)^{s/2}P_\Lambda u\|_{L^2}
 \le(1+\Lambda)^{s/2}\|P_\Lambda u\|_{L^2}.
\]
Split $u=P_\Lambda u+(I-P_\Lambda)u$, apply the triangle inequality,
the left inequality in \eqref{eq:V17-spectral-norm-equivalence}, and
\eqref{eq:V17-Sobolev-embedding}.  Summing
\eqref{eq:V17-band-leakage-bound} over the increments proves the last claim.
\end{proof}

\begin{assessment}[What a divergent shell upper bound means]
\label{ass:V17-shell-upper}
Condition \eqref{eq:V17-weighted-shell-sum} is sufficient for absolute
graph-norm transport.  Divergence of its right side is not proof that the
scheduler fails to converge, because different increments may cancel.
A nonvanishing high-frequency term is nevertheless a literal failure of the
declared spectrally localized shell model.  If
\[
 \|(I+\mathscr L)^{s/2}
 (I-P_{\Lambda_j})u_j\|_{L^2}\ge\varepsilon_*
\]
along a subsequence, its normalized high component is an explicit
graph-normalized leakage witness.  It must be retained rather than hidden
inside the low-band amplitude.
\end{assessment}

% -----------------------------------------------------------------------------
\subsection{Finite relation-covariant jet commands certify the graph inverse}
\label{subsec:V17-jet-card}
% -----------------------------------------------------------------------------

Let $P_m$ be a finite-rank orthogonal projection in $\mathscr X_K$ and
$Q_m$ one in $\mathscr Y_K$.  Assume the screens are compatible with the base
isomorphism:
\begin{equation}
 H_*P_m=Q_mH_*P_m.
 \label{eq:V17-compatible-screens}
\end{equation}
Fix $0<\eta<1$ and let $\mathcal N_{m,\eta}$ be a finite $\eta$-net of the
unit sphere of $P_m\mathscr X_K$.

\begin{lemma}[Finite-net operator bound]
\label{lem:V17-net-bound}
If $A:P_m\mathscr X_K\to Q_m\mathscr Y_K$ is linear, then
\begin{equation}
 \boxed{
 \|A\|_{\rm op}
 \le
 \frac1{1-\eta}
 \max_{v\in\mathcal N_{m,\eta}}\|Av\|.}
 \label{eq:V17-net-bound}
\end{equation}
\end{lemma}

\begin{proof}
For a unit vector $x$, choose $v$ in the net with $\|x-v\|\le\eta$.  Then
\[
 \|Ax\|\le\|Av\|+\eta\|A\|_{\rm op}.
\]
Take the supremum over $x$ and rearrange.
\end{proof}

Fix $z\in\mathscr U_R$ and a command radius $\varepsilon>0$ for which
$z\pm\varepsilon v\in\mathscr U_R$ on the net.  At each endpoint perform the
attached V30 common-action plaquette, endpoint relation, Kato polar-chord,
physical supported-connection, and determinant/Pfaffian line transports.
For command superpositions, also retain the V1 support-holonomy row.  Denote
the resulting base-aligned endpoint values in $Q_m\mathscr Y_K$ by
\[
 \widetilde{\mathcal F}_{n,z}(\pm\varepsilon v).
\]
Define the aligned central jet
\begin{equation}
 S_{n,m,\varepsilon}(z;v)
 :=
 \frac{
 \widetilde{\mathcal F}_{n,z}(\varepsilon v)
 -\widetilde{\mathcal F}_{n,z}(-\varepsilon v)}
 {2\varepsilon}.
 \label{eq:V17-aligned-central-jet}
\end{equation}
On a passing relation/line card this is the central difference of
$Q_m\mathcal F_n$ in one fixed physical frame.

Assume the derivative has chordwise Lipschitz constant $L_{n,m}$:
\begin{equation}
 \|D\mathcal F_n(z+tv)-D\mathcal F_n(z)\|_{\mathscr X_K\to\mathscr Y_K}
 \le L_{n,m}|t|
 \quad
 (|t|\le\varepsilon,\ v\in\mathcal N_{m,\eta}).
 \label{eq:V17-derivative-Lipschitz}
\end{equation}
Measure the finite net discrepancy
\begin{equation}
 d_{n,m,\varepsilon}(z)
 :=
 \max_{v\in\mathcal N_{m,\eta}}
 \|S_{n,m,\varepsilon}(z;v)-Q_mH_*v\|_{\mathscr Y_K}.
 \label{eq:V17-net-jet-discrepancy}
\end{equation}
Also record the target leakage and unresolved domain tail
\begin{align}
 \ell_{n,m}^{\rm out}(z)
 &:=
 \|(I-Q_m)D\mathcal F_n(z)P_m\|_{\rm op},
 \label{eq:V17-target-leakage}\\
 \tau_{n,m}^{\rm high}(z)
 &:=
 \|(D\mathcal F_n(z)-H_*)(I-P_m)\|_{\rm op}.
 \label{eq:V17-domain-tail}
\end{align}

\begin{theorem}[Finite jet card gives the infinite graph-radius bound]
\label{thm:V17-jet-inverse}
For the preceding card,
\begin{equation}
 \|S_{n,m,\varepsilon}(z;v)
   -Q_mD\mathcal F_n(z)v\|
 \le\frac12L_{n,m}\varepsilon
 \quad(v\in\mathcal N_{m,\eta}).
 \label{eq:V17-central-remainder}
\end{equation}
Consequently
\begin{equation}
 \boxed{
 \|D\mathcal F_n(z)-H_*\|_{\mathscr X_K\to\mathscr Y_K}
 \le
 \Delta_{n,m,\varepsilon}(z),}
 \label{eq:V17-full-derivative-defect}
\end{equation}
where
\begin{equation}
 \Delta_{n,m,\varepsilon}(z)
 :=
 \frac{
 d_{n,m,\varepsilon}(z)+L_{n,m}\varepsilon/2}
 {1-\eta}
 +\ell_{n,m}^{\rm out}(z)
 +\tau_{n,m}^{\rm high}(z).
 \label{eq:V17-jet-master-defect}
\end{equation}
In particular, if
\begin{equation}
 \sup_n\sup_{z\in\mathscr U_R}
 \|H_*^{-1}\|\,\Delta_{n,m(n),\varepsilon_n}(z)
 \le\theta<1,
 \label{eq:V17-jet-radius-pass}
\end{equation}
then the graph-space linearization hypothesis
\eqref{eq:V17-graph-linearization} holds.

Failure of the finite portion of this test returns one of the measured
endpoint relation, action plaquette, support-holonomy, line-sewing, net-jet,
target-leakage, or held-out command vectors.  Failure of the high-domain term
returns a unit vector in $(I-P_m)\mathscr X_K$ on which the derivative differs
from $H_*$ by the reported amount.
\end{theorem}

\begin{proof}
On a passing alignment card, the fundamental theorem of calculus gives
\[
 S_{n,m,\varepsilon}(z;v)
 =
 \frac1{2\varepsilon}
 \int_{-\varepsilon}^{\varepsilon}
 Q_mD\mathcal F_n(z+tv)v\,\dd t.
\]
Subtract $Q_mD\mathcal F_n(z)v$ and use
\eqref{eq:V17-derivative-Lipschitz}.  Since
\[
 \frac1{2\varepsilon}
 \int_{-\varepsilon}^{\varepsilon}|t|\,\dd t
 =\frac{\varepsilon}{2},
\]
this proves \eqref{eq:V17-central-remainder}.

Put
\[
 A:=Q_m(D\mathcal F_n(z)-H_*)P_m.
\]
On the net,
\[
 \|Av\|
 \le d_{n,m,\varepsilon}(z)+L_{n,m}\varepsilon/2.
\]
Lemma~\ref{lem:V17-net-bound} bounds $\|A\|$ by the first term of
\eqref{eq:V17-jet-master-defect}.  Compatibility
\eqref{eq:V17-compatible-screens} makes
\[
 (I-Q_m)(D\mathcal F_n-H_*)P_m
 =(I-Q_m)D\mathcal F_nP_m,
\]
which is bounded by \eqref{eq:V17-target-leakage}.  The action on the
remaining domain is \eqref{eq:V17-domain-tail}.  Splitting an arbitrary input
with $P_m+(I-P_m)$ and applying the triangle inequality proves
\eqref{eq:V17-full-derivative-defect}.  Finally,
\[
 \|I-H_*^{-1}D\mathcal F_n(z)\|
 \le\|H_*^{-1}\|\,
     \|D\mathcal F_n(z)-H_*\|,
\]
so \eqref{eq:V17-jet-radius-pass} implies
\eqref{eq:V17-graph-linearization}.  All stated failures are direct norm or
finite Gram failures; norm attainment holds on the finite screens, and a
near-maximizer gives the high-domain witness in the infinite complement.
\end{proof}

\begin{corollary}[Derivative-scaled cutoff transport of the jet card]
\label{cor:V17-jet-cutoff}
Let the two aligned endpoints in
\eqref{eq:V17-aligned-central-jet} have adjacent-cutoff route errors
$\delta_{n,+}(v,\varepsilon_n)$ and
$\delta_{n,-}(v,\varepsilon_n)$ after all input and output transports.
If
\begin{equation}
 \sum_n
 \sup_{v\in\mathcal N_{m,\eta}}
 \frac{
 \delta_{n,+}(v,\varepsilon_n)
 +\delta_{n,-}(v,\varepsilon_n)}
 {\varepsilon_n}
 <\infty
 \label{eq:V17-scaled-jet-transport}
\end{equation}
and the common-fine-card radius-change errors are summable, then the aligned
central jets are Cauchy on that fixed screen.  If in addition
$L_{n,m}\varepsilon_n\to0$, their limit is the transported physical
derivative on the screen.
\end{corollary}

\begin{proof}
The attached V30 derivative-scaled endpoint bound, applied uniformly to the
finite net, makes the successive central jets summably Cauchy.  The separate
radius-change row joins the diagonal radii.  Equation
\eqref{eq:V17-central-remainder} identifies the limit with the derivative.
\end{proof}

% -----------------------------------------------------------------------------
\subsection{A metric command secant for represented stress incidence}
\label{subsec:V17-stress-secant}
% -----------------------------------------------------------------------------

Fix a finite determining metric screen $P_m\mathscr T_K$ from V16 and a
source-complete OS graph-core atlas.  Let $\mathscr W_{n,m}$ be the finite
Hilbert response carrier containing all matrix elements on that atlas.  Let
\begin{equation}
 S_{n,m}:P_m\mathscr T_K\longrightarrow\mathscr W_{n,m}
 \label{eq:V17-direct-stress-map}
\end{equation}
be the direct represented stress synthesis with the normalization of
V16.  For a unit metric test $k$, let
\[
 t\longmapsto
 \mathscr A_{n,m}(t;k)\in\mathscr W_{n,m}
\]
be the completely aligned response of the same unnormalized common action
under the physical metric command $tk$.  All matter, gauge, potential,
counterterm, improvement, contact, vacuum, and boundary terms are included
before this response is differentiated.  Assume
\begin{equation}
 V_{n,m}k
 :=
 \left.\frac{\dd}{\dd t}\right|_{t=0}
 \mathscr A_{n,m}(t;k)
 \label{eq:V17-action-metric-derivative}
\end{equation}
is linear in $k$.  This is the common-action matter metric variation on the
represented source atlas.

For a radius $\varepsilon>0$, define the central metric-command tangent
\begin{equation}
 C_{n,m,\varepsilon}^{g}(k)
 :=
 \frac{
 \mathscr A_{n,m}(\varepsilon;k)
 -\mathscr A_{n,m}(-\varepsilon;k)}
 {2\varepsilon}.
 \label{eq:V17-metric-central-secant}
\end{equation}
Use a finite $\eta$-net
$\mathcal M_{m,\eta}$ of the unit sphere of $P_m\mathscr T_K$ and record
\begin{equation}
 d_{n,m,\varepsilon}^{\rm str}
 :=
 \max_{k\in\mathcal M_{m,\eta}}
 \|S_{n,m}k-C_{n,m,\varepsilon}^{g}(k)\|_{\mathscr W_{n,m}}.
 \label{eq:V17-direct-secant-stress-defect}
\end{equation}

\begin{theorem}[Metric secants bound stress/action incidence]
\label{thm:V17-stress-secant}
Suppose $t\mapsto\mathscr A_{n,m}(t;k)$ is three times continuously
differentiable and, for every $k\in\mathcal M_{m,\eta}$,
\begin{equation}
 \sup_{|t|\le\varepsilon}
 \left\|
 \frac{\dd^3}{\dd t^3}\mathscr A_{n,m}(t;k)
 \right\|_{\mathscr W_{n,m}}
 \le M_{n,m}^{(3)}.
 \label{eq:V17-third-metric-command}
\end{equation}
Then
\begin{equation}
 \boxed{
 \|S_{n,m}-V_{n,m}\|_{\rm op}
 \le
 D_{n,m,\varepsilon}^{\rm str}
 :=
 \frac{
 d_{n,m,\varepsilon}^{\rm str}
 +M_{n,m}^{(3)}\varepsilon^2/6}
 {1-\eta}.}
 \label{eq:V17-stress-incidence-bound}
\end{equation}
In particular, if the V16 scalar stress and action functionals are obtained
from this response carrier by a readout of norm at most
$C_{\rm src}$, then
\begin{equation}
 \|P_mr_n^{\rm str}\|_{\mathscr T_K}
 \le C_{\rm src}D_{n,m,\varepsilon}^{\rm str},
 \qquad
 \|\mathbb I_{n,m}^{\rm str}\|_{\rm op}
 \le C_{\rm src}^2
      (D_{n,m,\varepsilon}^{\rm str})^2.
 \label{eq:V17-to-V16-stress-Gram}
\end{equation}
If the atlas is source complete, zero incidence on an exhausting sequence of
screens identifies the represented stress with the common-action metric
variation on the limiting common core.
\end{theorem}

\begin{proof}
Taylor's formula with integral remainder, valid for Hilbert-valued maps,
gives
\[
 \left\|
 C_{n,m,\varepsilon}^{g}(k)-V_{n,m}k
 \right\|
 \le\frac16M_{n,m}^{(3)}\varepsilon^2.
\]
Indeed, each endpoint third-order remainder has norm at most
$M_{n,m}^{(3)}\varepsilon^3/6$; subtracting the endpoint expansions and
dividing by $2\varepsilon$ gives the stated constant.  Therefore, on the
finite net,
\[
 \|(S_{n,m}-V_{n,m})k\|
 \le
 d_{n,m,\varepsilon}^{\rm str}
 +M_{n,m}^{(3)}\varepsilon^2/6.
\]
The difference $S_{n,m}-V_{n,m}$ is linear, so
Lemma~\ref{lem:V17-net-bound} proves
\eqref{eq:V17-stress-incidence-bound}.

Composition with a readout of norm $C_{\rm src}$ gives the first inequality
in \eqref{eq:V17-to-V16-stress-Gram}.  The V16 incidence Gram is the rank-one
square of this Riesz vector, so its operator norm is the square of the vector
norm.  Finally, source completeness turns equality of every retained graph
matrix element into equality of the represented operators on the common
core; exhaustion and the separately required graph transport pass that
identity to the limit.
\end{proof}

\begin{corollary}[Cofinal metric-secant stress landing]
\label{cor:V17-stress-landing}
Choose $m_n\uparrow\infty$ and $\varepsilon_n\downarrow0$.  Suppose:
\begin{enumerate}[label=\textup{(\roman*)},leftmargin=2.8em]
\item the action endpoints and direct stress maps use one common source,
      relation, determinant-line, OS graph, and Lorentzian test transport;
\item $C_{\rm src}$ is uniform;
\item $D_{n,m_n,\varepsilon_n}^{\rm str}\to0$;
\item the high-screen stress-incidence tail tends to zero by the V16 spectral
      estimate;
\item the adjacent stress transports are summable.
\end{enumerate}
Then $r_n^{\rm str}\to0$ in $\mathscr T_K$, and the limiting represented
stress is the metric derivative of the limiting common action on the
determining core.
\end{corollary}

\begin{proof}
Equation \eqref{eq:V17-to-V16-stress-Gram} controls the finite screen.
The V16 high-screen estimate controls its orthogonal complement, so the full
Riesz vector tends to zero.  Summable graph and stress transport identifies
both limits on the same core.
\end{proof}

\begin{firewall}[A finite difference of normalized expectations is not stress]
\label{fire:V17-normalized-stress}
The response $\mathscr A_{n,m}$ must retain the unnormalized common-action
and absolute vacuum/pressure component.  Differentiating only normalized
expectations removes the identity direction and can change the stress by a
central vacuum term.  Likewise, a Ward-conserved direct insertion need not be
the action derivative.  Equation \eqref{eq:V17-direct-secant-stress-defect}
compares the two occurring objects before either inference is made.
\end{firewall}

% -----------------------------------------------------------------------------
\subsection{Combined graph-root to Einstein estimate}
\label{subsec:V17-combined}
% -----------------------------------------------------------------------------

\begin{theorem}[Graph-root and metric-secant quantitative handoff]
\label{thm:V17-combined-handoff}
Assume the graph-space root in Theorem~\ref{thm:V17-graph-root} is the
complete common-action root and its metric component is exactly
$\mathcal A_n$ from V16.  Hence
\begin{equation}
 r_n^{\rm act}=0.
 \label{eq:V17-action-component-zero}
\end{equation}
Assume the V16 connection and face hypotheses, and define
\begin{equation}
 E_n^{\rm Pal}
 :=
 \eta_n^{\rm face}+C_K(\kappa_n+\kappa_n^2).
 \label{eq:V17-Palatini-error}
\end{equation}
Assume also the V16 stress-incidence regularity bound with exponent
$\sigma_{\rm str}$ and constant $M_{\rm str}$.  For a metric screen
$P_m=\mathbf1_{[0,\Lambda_m]}(\mathscr L_K)$, put
\begin{equation}
 E_{n,m}^{\rm str}
 :=
 \left[
 C_{\rm src}^2
 (D_{n,m,\varepsilon}^{\rm str})^2
 +(1+\Lambda_m)^{-\sigma_{\rm str}}M_{\rm str}^2
 \right]^{1/2}.
 \label{eq:V17-stress-total-error}
\end{equation}
Then
\begin{equation}
 \boxed{
 \|r_n^E\|_{\mathscr T_K}
 \le E_n^{\rm Pal}+E_{n,m}^{\rm str}.}
 \label{eq:V17-finite-combined-bound}
\end{equation}

If the represented stresses obey the summable transport
\eqref{eq:V16-stress-transport}, then
\begin{align}
 &\|\mathcal R_K(\mathcal G_\infty+\mathcal S_\infty)\|_{\mathscr T_K}
 \nonumber\\
 &\qquad\le
 E_n^{\rm Pal}+E_{n,m}^{\rm str}
 +\frac{
 C_{\mathscr T,K}|\chi|C_{K,\Lambda}C_{\rm emb}}
 {(1-\theta)}
 \sum_{j\ge n}\epsilon_j^{(s)}
 +\sum_{j\ge n}d_j^{\rm str}.
 \label{eq:V17-limit-combined-bound}
\end{align}
Therefore any cofinal choice with
\begin{equation}
 \eta_n^{\rm face}\to0,\qquad
 \kappa_n\to0,\qquad
 E_{n,m(n)}^{\rm str}\to0
 \label{eq:V17-combined-pass}
\end{equation}
satisfies the distributional Einstein equation
\[
 2\chi(G_{\mu\nu}+\Lambda g_{\mu\nu})
 =T_{\mu\nu}^{\rm ren}
\]
on $K$, with the stress normalization and continuation hypotheses of V16.
\end{theorem}

\begin{proof}
Equation \eqref{eq:V17-action-component-zero} removes the first term in the
exact V16 three-defect identity.  Proposition
\ref{prop:V16-Palatini-incidence-bound} gives
$\|r_n^{\rm Pal}\|\le E_n^{\rm Pal}$.  The finite metric-secant estimate
\eqref{eq:V17-to-V16-stress-Gram}, combined with the V16 spectral tail,
gives $\|r_n^{\rm str}\|\le E_{n,m}^{\rm str}$.  The triangle inequality
proves \eqref{eq:V17-finite-combined-bound}.

The V16 limit estimate adds the scheduler curvature tail and the represented
stress transport tail.  Substitute \eqref{eq:V17-to-V16-tail} for the former
to obtain \eqref{eq:V17-limit-combined-bound}.  Every term tends to zero
under \eqref{eq:V17-graph-summability,eq:V17-combined-pass} and the
stress assumptions.  The final distributional identity is the V16
same-core conclusion.
\end{proof}

\begin{remark}[Matter and Ward equations remain separate components]
\label{rem:V17-matter-components}
If the graph-space residual $\mathcal F_n$ includes source-complete matter
variations and the transported Ward rows, its limiting zero also supplies
their Euler and conservation identities after their own nonlinear product and
domain limits pass.  Theorem~\ref{thm:V17-combined-handoff} proves only the
metric equation because the attached files do not populate those additional
graph-space components on the selected sequence.
\end{remark}

% -----------------------------------------------------------------------------
\subsection{Executable V17 card and obstruction order}
\label{subsec:V17-card}
% -----------------------------------------------------------------------------

\begin{definition}[Graph-root and metric-stress card]
\label{def:V17-card}
A V17 card on one compact cylinder contains:
\begin{enumerate}[label=\textup{(G\arabic*)},leftmargin=3em]
\item the physical graph spaces, boundary or Cauchy conditions, quotient
      maps, reference field, and proposed isomorphism $H_*$;
\item compatible finite domain and target screens, a declared finite
      $\eta$-net, and every common-action, relation-support, physical
      connection, line, and held-out row required to align its endpoints;
\item the central action jets, their net discrepancy, chordwise derivative
      Lipschitz bound, target leakage, high-domain tail, and the master defect
      \eqref{eq:V17-jet-master-defect};
\item the center residual and the adjacent preconditioned action defects
      $\epsilon_n^{(s)}$;
\item the scheduler shell bands, low-mode $L^2$ lengths, and graph-normalized
      spectral leakage in \eqref{eq:V17-weighted-shell-sum};
\item a source-complete represented stress atlas, the direct stress map,
      the two unnormalized metric-command endpoints, their third-command
      remainder, and $D_{n,m,\varepsilon}^{\rm str}$;
\item the V16 chronological connection, face-writer, high-screen stress,
      and OS-to-Lorentz transport rows.
\end{enumerate}
\end{definition}

\begin{theorem}[V17 finite and cofinal pass-or-witness alternative]
\label{thm:V17-card-alternative}
Every populated finite V17 card either:
\begin{enumerate}[label=\textup{(\roman*)},leftmargin=2.8em]
\item passes \eqref{eq:V17-jet-radius-pass}, the center condition, the finite
      action and stress incidences, and all declared margins; or
\item returns the first nonzero action plaquette, support holonomy, endpoint
      relation, physical-connection, determinant-line, net-jet,
      target-leakage, high-domain, center, adjacent-action, scheduler-shell,
      direct-stress, held-out-source, or third-remainder witness.
\end{enumerate}
A cofinal sequence of passing cards with summable adjacent action and stress
defects, vanishing face and connection errors, and exhausting spectral tails
satisfies \eqref{eq:V17-limit-combined-bound}.  Otherwise a failed finite row
or a normalized escaping graph mode remains as the physical output.
\end{theorem}

\begin{proof}
All aligned screen and net quantities are finite matrices or finite response
vectors, so their norms and Grams decide their listed rows.  Theorem
\ref{thm:V17-jet-inverse} converts a passing jet card into the graph-radius
bound.  Theorem~\ref{thm:V17-graph-root} converts that bound and the adjacent
defects into the $C^{2,\alpha}$ scheduler tail.  Theorem
\ref{thm:V17-stress-secant} decides the represented-stress/action incidence.
The V16 connection and spectral-tail alternatives decide the remaining
Palatini and high-screen rows.  Their composition is
Theorem~\ref{thm:V17-combined-handoff}.  Each failure listed in
item~\textup{(ii)} is the norm-attaining finite vector or the normalized
near-maximizer already constructed in those component theorems.
\end{proof}

% -----------------------------------------------------------------------------
\subsection{V17 theorem ledger and next graph-inverse acquisition}
\label{subsec:V17-ledger}
% -----------------------------------------------------------------------------

\paragraph{New conclusions proved in V17.}
Theorem~\ref{thm:V17-graph-root} lifts the action-selected root construction
from a finite coefficient norm to a spacetime Sobolev graph norm and derives
the exact $C^{2,\alpha}$ tail required by V16.  Proposition
\ref{prop:V17-spectral-shell} makes that regularity testable through low-band
shell length plus graph-normalized leakage.  Theorem
\ref{thm:V17-jet-inverse} converts finitely many relation-covariant central
command secants, a sphere net, and two leakage norms into the full
graph-linearization bound.  Theorem~\ref{thm:V17-stress-secant} gives a
source-complete, third-remainder-controlled finite test of represented stress
against the metric derivative of the same unnormalized action.  Theorem
\ref{thm:V17-combined-handoff} inserts these two acquisitions into the V16
three-defect identity and prints the resulting Einstein-residual bound.

\paragraph{Imported results not repeated.}
The V14 pointwise scheduler and finite contraction proof, V15 signed quotient
Hessian, V16 curvature and same-core closure, the attached V30
relation-covariant central-secant scaling, the V1 returned-word support
holonomy, the monograph's OS stress reconstruction, and its Palatini
spectral-tail theorem are imported.  V17 supplies the missing graph-space,
finite-net, spectral-shell, and metric-stress incidence links between them.

\paragraph{Authoritative physical status.}
No attached file supplies a gauge-fixed differential operator
$H_*:\mathscr X_K\to\mathscr Y_K$ for the selected interacting RPESM cylinder,
its boundary or Cauchy realization, a uniform inverse norm, or the high-domain
tails \eqref{eq:V17-domain-tail}.  The attachments also do not contain the
unnormalized metric-command endpoint responses and direct represented stress
matrices needed for
\eqref{eq:V17-direct-secant-stress-defect}.  Therefore neither
\eqref{eq:V17-jet-radius-pass} nor
\eqref{eq:V17-combined-pass} is asserted for a physical regulator sequence.
The full Standard Model--Einstein continuum remains unproved.

\paragraph{Next acquisition.}
The next iteration should attack the base graph inverse rather than add
another finite response quotient.  On one selected compact cylinder it should:
\begin{enumerate}[label=\textup{(\arabic*)},leftmargin=2.8em]
\item write the principal and lower-order blocks of the completely gauge-fixed
      Palatini--Standard-Model linearization, including constraint and
      boundary or Cauchy conditions;
\item decide whether the selected branch is an elliptic Euclidean boundary
      problem or a Lorentzian hyperbolic initial-boundary problem;
\item prove the corresponding Fredholm/Schauder or causal energy estimate in
      the physical quotient graph norm;
\item identify every finite-dimensional kernel and cokernel with a declared
      gauge, relation, boundary, or genuine physical mode;
\item construct $H_*^{-1}$ or return a normalized kernel, cokernel, glancing,
      boundary, or high-frequency quasimode;
\item only then populate the shrinking jet and metric-stress endpoint cards of
      V17.
\end{enumerate}
This is upstream of the remaining curvature and stress values: without a
graph inverse, a sequence of passing low-mode Hessian cards can still lose its
margin into unresolved spacetime frequency.

\begin{firewall}[V17 proves an implication, not the graph estimate]
\label{fire:V17-final-scope}
V17 proves that a physical graph inverse plus measurable adjacent defects
generates the V16 curvature branch, and that an unnormalized metric-command
secant tests stress/action incidence.  It does not infer the graph inverse
from finite Fisher information, infer a causal estimate from an elliptic
symbol, supply the missing RPESM metric-command endpoints, or establish the
remaining matter, Ward, counterterm, memory, and global boundary limits.
\end{firewall}

% =============================================================================
% END APPENDED VERSION V17 MATERIAL
% =============================================================================



% =============================================================================
% BEGIN APPENDED VERSION V18 MATERIAL
% =============================================================================

\clearpage
\section{Version V18: the gauge-fixed graph-inverse alternative}
\label{sec:V18-graph-inverse}

Version V17 showed that a bounded inverse for the completely reduced
common-action linearization converts summable adjacent action defects into the
two-derivative scheduler convergence needed by V16.  A finite Hessian floor is
not enough: the differential realization must include its domain, boundary or
Cauchy data, gauge propagation, high-frequency estimate, and the couplings
between gravity and the Standard Model.

This version constructs the missing analytic interface.  It first computes
the reduced Einstein principal symbol and proves why the Euclidean and
Lorentzian graph problems cannot share one inverse theorem.  It then gives a
compact-remainder Fredholm criterion for the Euclidean branch and a
symmetric-hyperbolic energy criterion for the Lorentzian branch.  Finally, an
exact block Schur theorem assembles either gravitational inverse with the
matter, gauge, ghost, and constraint blocks without assuming positivity of the
coupled Hessian.  Every positive conclusion remains conditional on the
physical RPESM linearization satisfying the displayed coefficient, domain,
and incidence hypotheses.

% -----------------------------------------------------------------------------
\subsection{The reduced Einstein symbol and the two analytic branches}
\label{subsec:V18-reduced-symbol}
% -----------------------------------------------------------------------------

Let $g$ be a smooth nondegenerate background metric and let
$h\in\Gamma(\operatorname{Sym}^2T^*M)$.  Define the de Donder gauge one-form
\begin{equation}
 \mathcal C_g(h)_\nu
 :=
 \nabla^\mu h_{\mu\nu}
 -\frac12\nabla_\nu\operatorname{tr}_g h.
 \label{eq:V18-de-Donder}
\end{equation}
Parentheses around tensor indices include the factor $1/2$.  Define the
reduced linearized Ricci operator
\begin{equation}
 \mathcal L_g^{\rm red}h
 :=
 D\operatorname{Ric}_g[h]
 -\nabla_{(\mu}\mathcal C_g(h)_{\nu)}.
 \label{eq:V18-reduced-Ricci}
\end{equation}

\begin{theorem}[Exact reduced Einstein principal symbol]
\label{thm:V18-reduced-symbol}
Modulo curvature and other zeroth-order terms,
\begin{equation}
 \boxed{
 (\mathcal L_g^{\rm red}h)_{\mu\nu}
 =
 -\frac12\nabla^\alpha\nabla_\alpha h_{\mu\nu}
 +\operatorname{l.o.t.}}
 \label{eq:V18-reduced-principal-part}
\end{equation}
Consequently, with the symbol convention $\partial_\alpha\mapsto i\xi_\alpha$,
\begin{equation}
 \boxed{
 \sigma_2(\mathcal L_g^{\rm red})(x,\xi)h
 =
 \frac12g^{\alpha\beta}(x)\xi_\alpha\xi_\beta\,h.}
 \label{eq:V18-reduced-symbol}
\end{equation}
Algebraic trace reversal and the cosmological term do not change this
principal symbol.

If $g$ is Riemannian, the operator is strongly elliptic on symmetric tensors.
If $g$ is Lorentzian, its characteristic set is the null cone and the operator
is normally hyperbolic.  Hence a Euclidean elliptic boundary inverse and a
Lorentzian causal inverse are different graph realizations.

For a vector field $X$, let
$(K_gX)_{\mu\nu}=2\nabla_{(\mu}X_{\nu)}$.  Then
\begin{equation}
 \boxed{
 \mathcal C_g(K_gX)_\nu
 =
 \nabla^\alpha\nabla_\alpha X_\nu
 +\operatorname{Ric}_{\nu\mu}X^\mu.}
 \label{eq:V18-gauge-vector-operator}
\end{equation}
Thus gauge fixing replaces infinitesimal diffeomorphism zero directions by a
vector Laplace or wave operator; it does not remove their boundary, initial,
or global zero modes automatically.
\end{theorem}

\begin{proof}
The second-derivative part of the Ricci variation is
\[
 D\operatorname{Ric}_g[h]_{\mu\nu}
 =
 \frac12\left(
 \nabla^\alpha\nabla_\mu h_{\alpha\nu}
 +\nabla^\alpha\nabla_\nu h_{\alpha\mu}
 -\nabla^\alpha\nabla_\alpha h_{\mu\nu}
 -\nabla_\mu\nabla_\nu\operatorname{tr}_g h
 \right)
 +\operatorname{l.o.t.}
\]
Commuting derivatives changes only curvature terms.  The two mixed terms and
the trace term combine to
$\nabla_{(\mu}\mathcal C_g(h)_{\nu)}$ at principal order.  Subtraction gives
\eqref{eq:V18-reduced-principal-part}, and replacing derivatives by
$i\xi$ gives \eqref{eq:V18-reduced-symbol}.  Trace reversal is an invertible
pointwise algebraic map in dimension four, while $\Lambda h$ has order zero.

For Riemannian $g$, $g^{\alpha\beta}\xi_\alpha\xi_\beta>0$ whenever
$\xi\ne0$, proving strong ellipticity.  For Lorentzian $g$, that scalar
vanishes precisely on nonzero null covectors, which is the normally
hyperbolic symbol.

Finally,
\[
 \begin{aligned}
 \mathcal C_g(K_gX)_\nu
 &=
 \nabla^\mu(\nabla_\mu X_\nu+\nabla_\nu X_\mu)
 -\nabla_\nu\nabla_\mu X^\mu\\
 &=
 \nabla^\mu\nabla_\mu X_\nu
 +[\nabla^\mu,\nabla_\nu]X_\mu.
 \end{aligned}
\]
The commutator contraction is
$\operatorname{Ric}_{\nu\mu}X^\mu$, proving
\eqref{eq:V18-gauge-vector-operator}.
\end{proof}

\begin{corollary}[Gauge propagation is a separate homogeneous equation]
\label{cor:V18-gauge-propagation}
Let a reduced Einstein equation be written
\begin{equation}
 G_{\mu\nu}+\Lambda g_{\mu\nu}-T_{\mu\nu}
 -\nabla_{(\mu}C_{\nu)}
 +\frac12g_{\mu\nu}\nabla_\alpha C^\alpha=0.
 \label{eq:V18-nonlinear-reduced-equation}
\end{equation}
If $\nabla^\mu T_{\mu\nu}=0$, then
\begin{equation}
 \boxed{
 \nabla^\alpha\nabla_\alpha C_\nu
 +\operatorname{Ric}_{\nu\mu}C^\mu=0.}
 \label{eq:V18-gauge-propagation}
\end{equation}
On a Lorentzian Cauchy branch, $C$ and its normal derivative vanishing on the
initial surface imply $C=0$ wherever uniqueness holds.  On a Euclidean
boundary branch, the same conclusion requires a kernel-free realization of
the vector operator in \eqref{eq:V18-gauge-propagation}.
\end{corollary}

\begin{proof}
Take the divergence of
\eqref{eq:V18-nonlinear-reduced-equation}.  The contracted Bianchi identity,
metric compatibility, and stress conservation remove the first three terms.
For the gauge terms,
\[
 \nabla^\mu\nabla_{(\mu}C_{\nu)}
 -\frac12\nabla_\nu\nabla_\alpha C^\alpha
 =
 \frac12\left(
 \nabla^\alpha\nabla_\alpha C_\nu
 +\operatorname{Ric}_{\nu\mu}C^\mu
 \right).
\]
This proves the homogeneous equation.  The final statements are uniqueness
for the corresponding causal or elliptic vector realization.
\end{proof}

% -----------------------------------------------------------------------------
\subsection{The Euclidean compact-remainder inverse}
\label{subsec:V18-Euclidean}
% -----------------------------------------------------------------------------

The reduced symbol alone does not imply invertibility.  Boundary
complementarity, kernel removal, and target coverage remain necessary.  The
following functional-analytic theorem isolates exactly those rows.

\begin{theorem}[Compact-remainder graph inverse or a Fredholm witness]
\label{thm:V18-compact-remainder}
Let $\mathscr X,\mathscr Y,\mathscr Z$ be Hilbert spaces,
$L\in\mathcal B(\mathscr X,\mathscr Y)$, and
$K\in\mathcal K(\mathscr X,\mathscr Z)$ compact.  Assume the
compact-remainder estimate
\begin{equation}
 \|u\|_{\mathscr X}
 \le C_E\bigl(
 \|Lu\|_{\mathscr Y}+\|Ku\|_{\mathscr Z}
 \bigr)
 \quad(u\in\mathscr X).
 \label{eq:V18-compact-remainder-estimate}
\end{equation}
Then:
\begin{enumerate}[label=\textup{(\roman*)},leftmargin=2.8em]
\item $\ker L$ is finite dimensional;
\item if $\ker L=\{0\}$, there is $C_L<\infty$ such that
      \begin{equation}
       \|u\|_{\mathscr X}\le C_L\|Lu\|_{\mathscr Y},
       \label{eq:V18-kernel-free-estimate}
      \end{equation}
      and $\operatorname{Ran}L$ is closed;
\item if also $\ker L^*=\{0\}$, then $L$ is a bounded isomorphism and
      $\|L^{-1}\|\le C_L$;
\item a nonzero element of $\ker L$ is a domain obstruction, while a nonzero
      $y\in\ker L^*$ is a target-cokernel obstruction satisfying
      \begin{equation}
       \langle Lu,y\rangle_{\mathscr Y}=0
       \quad(u\in\mathscr X).
       \label{eq:V18-cokernel-witness}
      \end{equation}
\end{enumerate}
If no constant $C_E$ makes
\eqref{eq:V18-compact-remainder-estimate} true for the declared compact
remainder, there are unit vectors $u_j\in\mathscr X$ with
\begin{equation}
 \|Lu_j\|_{\mathscr Y}+\|Ku_j\|_{\mathscr Z}
 \longrightarrow0.
 \label{eq:V18-Euclidean-quasimode}
\end{equation}
They are explicit high-frequency or boundary quasimode witnesses once a
violating sequence is extracted.
\end{theorem}

\begin{proof}
On $\ker L$, estimate \eqref{eq:V18-compact-remainder-estimate} says
$\|u\|_{\mathscr X}\le C_E\|Ku\|_{\mathscr Z}$.  If the kernel were infinite
dimensional, choose an orthonormal sequence $u_j$ in it.  Compactness gives a
Cauchy subsequence of $Ku_j$, and the estimate applied to differences would
make the corresponding $u_j$ Cauchy, a contradiction.  This proves
item~\textup{(i)}.

Suppose $\ker L=0$ and \eqref{eq:V18-kernel-free-estimate} fails.  There are
$u_j$ with $\|u_j\|_{\mathscr X}=1$ and $\|Lu_j\|_{\mathscr Y}\to0$.
After taking a subsequence, $Ku_j$ is Cauchy.  Apply
\eqref{eq:V18-compact-remainder-estimate} to $u_j-u_k$; the subsequence is
Cauchy in $\mathscr X$ and converges to a unit vector $u$.  Continuity gives
$Lu=0$, contradicting the trivial kernel.  Thus
\eqref{eq:V18-kernel-free-estimate} holds.  It immediately implies that
$\operatorname{Ran}L$ is closed.

For bounded operators between Hilbert spaces,
$(\operatorname{Ran}L)^\perp=\ker L^*$.  A zero adjoint kernel therefore
makes the already closed range dense and hence equal to $\mathscr Y$.
The inverse bound is \eqref{eq:V18-kernel-free-estimate}.  Item~\textup{(iv)}
is the kernel definition and the same orthogonal-range identity.

If the compact-remainder estimate fails for every positive integer constant,
choose $v_j$ with
\[
 \|v_j\|_{\mathscr X}
 >
 j\bigl(\|Lv_j\|_{\mathscr Y}+\|Kv_j\|_{\mathscr Z}\bigr)
\]
and normalize.  This gives \eqref{eq:V18-Euclidean-quasimode}.
\end{proof}

\begin{corollary}[Explicit absorption of the compact remainder]
\label{cor:V18-compact-absorption}
In Theorem~\ref{thm:V18-compact-remainder}, suppose
\begin{equation}
 \|Ku\|_{\mathscr Z}
 \le a\|Lu\|_{\mathscr Y}+b\|u\|_{\mathscr X},
 \qquad C_Eb<1.
 \label{eq:V18-compact-absorption}
\end{equation}
Then
\begin{equation}
 \boxed{
 \|u\|_{\mathscr X}
 \le
 \frac{C_E(1+a)}{1-C_Eb}
 \|Lu\|_{\mathscr Y}.}
 \label{eq:V18-explicit-Euclidean-inverse}
\end{equation}
If $\ker L^*=0$, this is an explicit inverse bound.
\end{corollary}

\begin{proof}
Insert \eqref{eq:V18-compact-absorption} in
\eqref{eq:V18-compact-remainder-estimate} and move
$C_Eb\|u\|_{\mathscr X}$ to the left.
\end{proof}

\begin{corollary}[Elliptic gauge-fixed Euclidean branch]
\label{cor:V18-Euclidean-branch}
Let $L_E$ be the completely assembled Euclidean
Palatini--Standard-Model linearization after diffeomorphism, local Lorentz,
internal gauge, relation, and boundary-null reduction.  Give each field its
natural order in a common Douglis--Nirenberg graph norm.  Suppose:
\begin{enumerate}[label=\textup{(E\arabic*)},leftmargin=3em]
\item the interior principal symbol is strongly elliptic;
\item the chosen boundary operators satisfy the complementing boundary
      condition and give the estimate
      \eqref{eq:V18-compact-remainder-estimate}, with $K$ the compact
      one-order-lower embedding;
\item every domain kernel and adjoint kernel has been computed after the
      physical quotient and both are zero.
\end{enumerate}
Then $L_E$ is the graph isomorphism $H_*$ required by V17.  A failure of
item~\textup{(E1)} is a principal-symbol vector, failure of
item~\textup{(E2)} is a boundary-symbol or quasimode sequence, and failure of
item~\textup{(E3)} is the kernel or cokernel vector in
Theorem~\ref{thm:V18-compact-remainder}.
\end{corollary}

\begin{proof}
Strong ellipticity and the complementing boundary condition supply the stated
compact-remainder estimate in the declared graph norms.  Apply
Theorem~\ref{thm:V18-compact-remainder} with $L=L_E$.
\end{proof}

\begin{counterexample}[A Euclidean inverse does not give a two-ended Lorentzian inverse]
\label{cth:V18-Euclidean-not-Lorentzian}
Let
\[
 K=[0,\pi]_t\times\mathbb S^1_x
\]
with Dirichlet conditions at $t=0,\pi$ and periodicity in $x$.  The Euclidean
operator
\begin{equation}
 L_E=-\partial_t^2-\partial_x^2
 \label{eq:V18-Euclidean-control}
\end{equation}
has eigenvalues $n^2+k^2>0$ on
$\sin(nt)e^{ikx}$, $n\ge1$, $k\in\mathbb Z$, and is invertible.  The
Lorentzian operator
\begin{equation}
 L_L=\partial_t^2-\partial_x^2
 \label{eq:V18-Lorentzian-control}
\end{equation}
with the same two-ended boundary conditions has the nonzero kernel
\begin{equation}
 L_L(\sin t\,e^{ix})=0.
 \label{eq:V18-Lorentzian-kernel}
\end{equation}
Thus Wick-signature change does not transport an elliptic boundary inverse
into a Lorentzian boundary inverse.  The Lorentzian problem must use Cauchy or
a separately justified initial-boundary data set.
\end{counterexample}

\begin{proof}
Direct differentiation gives the displayed eigenvalues and the zero in
\eqref{eq:V18-Lorentzian-kernel}.
\end{proof}

% -----------------------------------------------------------------------------
\subsection{The Lorentzian causal graph inverse}
\label{subsec:V18-Lorentzian}
% -----------------------------------------------------------------------------

Let $\Sigma$ be a compact three-manifold without boundary and
$K_T=[0,T]\times\Sigma$.  A physical timelike boundary can be added only with
a separately verified maximally dissipative or constraint-preserving boundary
condition.  Consider the real first-order system
\begin{equation}
 \mathcal P U
 :=
 A^0(t,x)\partial_tU
 +\sum_{a=1}^3A^a(t,x)\nabla_aU
 +B(t,x)U
 =F.
 \label{eq:V18-symmetric-hyperbolic-system}
\end{equation}

\begin{theorem}[Symmetric-hyperbolic energy graph estimate]
\label{thm:V18-hyperbolic-energy}
Assume:
\begin{enumerate}[label=\textup{(H\arabic*)},leftmargin=3em]
\item $A^0$ and every $A^a$ are symmetric;
\item for constants $0<a_0\le a_1<\infty$,
      \begin{equation}
       a_0I\preceq A^0(t,x)\preceq a_1I;
       \label{eq:V18-symmetrizer-floor}
      \end{equation}
\item the coefficients and the derivatives needed for the $H^r$ commutators
      are bounded, for an integer $r\ge0$.
\end{enumerate}
Then every smooth solution satisfies
\begin{equation}
 \boxed{
 \|U\|_{C([0,T];H^r(\Sigma))}
 \le
 C_{T,r}\left(
 \|U(0)\|_{H^r(\Sigma)}
 +\|F\|_{L^1([0,T];H^r(\Sigma))}
 \right).}
 \label{eq:V18-hyperbolic-energy}
\end{equation}
The constant depends only on $T$, the symmetrizer floor and ceiling, the
background spatial geometry, and the displayed coefficient bounds.

After completion in the energy graph norm, the map
\begin{equation}
 H_LU:=(\mathcal PU,U(0))
 \label{eq:V18-causal-graph-map}
\end{equation}
is injective with closed range.  Standard Galerkin convergence for the stated
linear symmetric-hyperbolic system makes the range all of
\[
 L^1([0,T];H^r(\Sigma))\oplus H^r(\Sigma).
\]
Hence $H_L$ is a bounded isomorphism onto the source-plus-Cauchy-data space,
with inverse norm bounded by $C_{T,r}$.
\end{theorem}

\begin{proof}
For $r=0$, define
\[
 E(t):=\frac12\int_\Sigma
 \langle A^0U,U\rangle\,\dd\operatorname{vol}_\Sigma.
\]
Multiply \eqref{eq:V18-symmetric-hyperbolic-system} by $U$, integrate over
the closed spatial slice, and integrate the symmetric spatial terms by
parts.  There is no boundary flux.  Derivatives of $A^0,A^a$, the connection
coefficients, and $B$ are lower-order terms, so
\[
 \frac{\dd}{\dd t}E(t)
 \le C E(t)+C E(t)^{1/2}\|F(t)\|_{L^2}.
\]
The floor and ceiling in \eqref{eq:V18-symmetrizer-floor}, followed by
Gronwall's inequality, give
\[
 \|U(t)\|_{L^2}
 \le C_T\left(
 \|U(0)\|_{L^2}
 +\int_0^t\|F(\tau)\|_{L^2}\,\dd\tau
 \right).
\]

For integer $r>0$, commute covariant spatial derivatives of order at most $r$
through the equation.  The leading matrices remain symmetric and have the
same $A^0$ floor.  The commutators contain at most $r$ derivatives of $U$
with bounded coefficients.  Summing the preceding energy estimate over the
commuted system proves \eqref{eq:V18-hyperbolic-energy}.

The estimate with zero source and initial data proves injectivity.  Applied to
differences, it also proves that a graph-Cauchy sequence with convergent image
is Cauchy in $C H^r$, so the range is closed.  For existence, project the
system onto the first $N$ eigenfunctions of a positive elliptic operator on
$\Sigma$.  The projected equations are finite-dimensional linear ODEs.
The same estimate is uniform in $N$; weak compactness gives an energy
solution, and the estimate gives uniqueness and strong convergence at lower
order.  Approximation of the data and the equation upgrades to the completed
$H^r$ graph space.  Thus every source and initial datum is attained, and the
same estimate bounds the inverse.
\end{proof}

\begin{corollary}[Constraint and gauge propagation with a measured debit]
\label{cor:V18-constraint-propagation}
Let the gauge, Hamiltonian, momentum, Gauss, or Cartan constraint vector
$c=\mathcal CU$ satisfy a subsidiary symmetric-hyperbolic system
\begin{equation}
 \mathcal P_Cc=R_C
 \label{eq:V18-subsidiary-system}
\end{equation}
with a symmetrizer floor and coefficient bounds as above.  Then
\begin{equation}
 \boxed{
 \|c\|_{C H^{r_C}}
 \le C_{C,T}\left(
 \|c(0)\|_{H^{r_C}}
 +\|R_C\|_{L^1H^{r_C}}
 \right).}
 \label{eq:V18-constraint-debit}
\end{equation}
In particular compatible zero initial constraints and $R_C=0$ propagate
exactly.  At finite cutoff the right side is the quantitative constraint
transport debit and cannot be discarded.
\end{corollary}

\begin{proof}
Apply Theorem~\ref{thm:V18-hyperbolic-energy} to the subsidiary system.
\end{proof}

\begin{corollary}[Lorentzian reduced Einstein--Standard-Model branch]
\label{cor:V18-Lorentzian-branch}
Suppose the completely gauge-fixed linearized Palatini--Standard-Model
evolution on $K_T$ admits a first-order reduction of the form
\eqref{eq:V18-symmetric-hyperbolic-system}, with:
\begin{enumerate}[label=\textup{(L\arabic*)},leftmargin=3em]
\item a uniformly positive symmetrizer on the physical evolution variables;
\item compatible initial data and a common coefficient/domain identification
      across cutoffs;
\item homogeneous subsidiary equations for every declared constraint, up to
      the measured residual $R_C$;
\item no unaccounted timelike-boundary flux.
\end{enumerate}
Then the source-plus-Cauchy-data map $H_L$ is a valid base graph isomorphism
$H_*$ for V17.  The V17 root theorem applies to the nonlinear
preconditioned residual if its derivative stays within the declared
$\theta<1$ radius.  Constraint error is bounded independently by
\eqref{eq:V18-constraint-debit}.
\end{corollary}

\begin{proof}
The graph inverse is Theorem~\ref{thm:V18-hyperbolic-energy}.  Constraint
propagation is Corollary~\ref{cor:V18-constraint-propagation}.  Substitution
of this inverse into the V17 hypotheses gives the last assertion.
\end{proof}

\begin{firewall}[A positive action Hessian is not the hyperbolic symmetrizer]
\label{fire:V18-symmetrizer}
The symmetrizer $A^0$ controls a first-order causal evolution after gauge and
constraint reduction.  It is not the signed second variation of the
Einstein action and does not erase the DeWitt conformal direction.  Conversely,
an invertible indefinite action Hessian does not supply
\eqref{eq:V18-symmetrizer-floor}.  The action-incidence and causal-energy rows
must both be checked on the same physical variables.
\end{firewall}

% -----------------------------------------------------------------------------
\subsection{Exact coupling of the gravitational and Standard-Model graph blocks}
\label{subsec:V18-coupled-Schur}
% -----------------------------------------------------------------------------

Let $\mathscr X_g,\mathscr Y_g$ be the reduced gravitational graph spaces and
$\mathscr X_m,\mathscr Y_m$ the direct sum of the scalar, gauge, fermion,
ghost, source, and retained constraint spaces.  Suppose the diagonal graph
operators
\[
 L_g:\mathscr X_g\to\mathscr Y_g,
 \qquad
 L_m:\mathscr X_m\to\mathscr Y_m
\]
are bounded isomorphisms on one of the two branches above.  Let
\[
 B:\mathscr X_m\to\mathscr Y_g,
 \qquad
 C:\mathscr X_g\to\mathscr Y_m
\]
be the completely assembled off-diagonal variations from the same common
action.  Define
\begin{equation}
 \mathbb H
 :=
 \begin{pmatrix}
  L_g&B\\ C&L_m
 \end{pmatrix},
 \qquad
 A:=L_g^{-1}B,
 \qquad
 D:=L_m^{-1}C.
 \label{eq:V18-coupled-operator}
\end{equation}

\begin{theorem}[Exact graph-space Schur alternative]
\label{thm:V18-graph-Schur}
The coupled operator $\mathbb H$ is invertible if and only if
\begin{equation}
 S_m:=I_{\mathscr X_m}-DA
 \label{eq:V18-matter-Schur}
\end{equation}
is invertible.  Equivalently, the gravitational Schur operator
$S_g=I_{\mathscr X_g}-AD$ is invertible.  The exact factorization is
\begin{equation}
 \begin{pmatrix}L_g^{-1}&0\\0&L_m^{-1}\end{pmatrix}
 \mathbb H
 =
 \begin{pmatrix}I&0\\D&I\end{pmatrix}
 \begin{pmatrix}I&A\\0&S_m\end{pmatrix}.
 \label{eq:V18-Schur-factorization}
\end{equation}
For preconditioned data $(f,g)$, the solution is
\begin{align}
 y&=S_m^{-1}(g-Df),
 \label{eq:V18-Schur-y}\\
 x&=f-Ay.
 \label{eq:V18-Schur-x}
\end{align}

The sufficient small-gain condition
\begin{equation}
 q:=\|D\|\,\|A\|<1
 \label{eq:V18-small-gain}
\end{equation}
gives
\begin{equation}
 \|S_m^{-1}\|\le\frac1{1-q}.
 \label{eq:V18-Schur-inverse-bound}
\end{equation}
Using sum norms on the product spaces,
\begin{equation}
 \boxed{
 \|\mathbb H^{-1}\|
 \le
 \frac{(1+\|A\|)(1+\|D\|)}{1-q}
 \max\{\|L_g^{-1}\|,\|L_m^{-1}\|\}.}
 \label{eq:V18-coupled-inverse-bound}
\end{equation}

Condition \eqref{eq:V18-small-gain} is not necessary.  On any branch, a
nonzero $y\in\ker S_m$ gives the coupled kernel vector
\begin{equation}
 \binom{-Ay}{y}\in\ker\mathbb H.
 \label{eq:V18-coupled-kernel}
\end{equation}
A nonzero $z\in\ker S_m^*$ gives the preconditioned left-cokernel vector
\begin{equation}
 \binom{-D^*z}{z}
 \label{eq:V18-coupled-cokernel}
\end{equation}
and hence, after applying the inverse adjoint of the diagonal preconditioner, a
target obstruction for $\mathbb H$.
\end{theorem}

\begin{proof}
Left multiplication by the diagonal inverse gives the block operator
\[
 M=\begin{pmatrix}I&A\\D&I\end{pmatrix}.
\]
Direct multiplication proves \eqref{eq:V18-Schur-factorization}.  The first
triangular factor is always invertible, while the second is invertible exactly
when $S_m$ is.  Forward substitution gives
\eqref{eq:V18-Schur-y,eq:V18-Schur-x}.  Eliminating $y$ instead proves the
equivalent $S_g$ criterion.

Under \eqref{eq:V18-small-gain}, $\|DA\|\le q<1$, so the Neumann series gives
\eqref{eq:V18-Schur-inverse-bound}.  In the sum norm, the inverse of the first
triangular factor has norm at most $1+\|D\|$.  The inverse of the second has
norm at most $(1+\|A\|)\|S_m^{-1}\|$.  Multiplication by the diagonal inverse
proves \eqref{eq:V18-coupled-inverse-bound}.

If $S_my=0$, substitution shows
\[
 M\binom{-Ay}{y}
 =
 \binom0{(I-DA)y}=0.
\]
Also
\[
 M^*\binom{-D^*z}{z}
 =
 \binom0{(I-A^*D^*)z}
 =
 \binom0{S_m^*z}.
\]
This proves the kernel and cokernel assertions, with the diagonal
preconditioner restored for the original target pairing.
\end{proof}

\begin{corollary}[Compact Euclidean coupling gives a finite mixed obstruction]
\label{cor:V18-compact-coupling}
On the Euclidean branch, suppose $A$ and $D$ are compact; this occurs, for
example, when the off-diagonal action variations are lower order relative to
the diagonal elliptic graph norms on a compact cylinder.  Then $DA$ and $AD$
are compact, and $S_m=I-DA$ is Fredholm of index zero.  Therefore either
$\mathbb H$ is invertible or
\eqref{eq:V18-coupled-kernel} and
\eqref{eq:V18-coupled-cokernel} contain nonzero finite-dimensional mixed
gravity--matter obstructions.
\end{corollary}

\begin{proof}
A composition containing a compact operator is compact.  The
Riesz--Schauder theorem makes identity minus compact Fredholm of index zero.
Apply Theorem~\ref{thm:V18-graph-Schur}.
\end{proof}

\begin{corollary}[Short-time Lorentzian coupled inverse]
\label{cor:V18-short-time-coupling}
On a Lorentzian causal slab, let $L_g^{-1}$ and $L_m^{-1}$ be the diagonal
retarded solution operators including the Cauchy data.  If the causal
off-diagonal gains obey
\begin{equation}
 \|L_m^{-1}C\|\,\|L_g^{-1}B\|
 \le q_T<1,
 \label{eq:V18-causal-small-gain}
\end{equation}
then the complete coupled source-plus-data map is invertible with
\eqref{eq:V18-coupled-inverse-bound}.  If $q_T<1$ holds successively on a
finite cover of a longer time interval and the terminal data of each slab are
the initial data of the next, the inverse propagates across the cover.
Constraint debits accumulate according to
\eqref{eq:V18-constraint-debit}.
\end{corollary}

\begin{proof}
Apply Theorem~\ref{thm:V18-graph-Schur} on each slab.  Uniqueness makes the
successive solutions agree at shared Cauchy surfaces, so their concatenation
solves the full interval problem.  The constraint estimate is applied and
telescoped on the same slab decomposition.
\end{proof}

\begin{corollary}[The coupled inverse supplies the V17 preconditioner]
\label{cor:V18-to-V17}
Assume either Corollary~\ref{cor:V18-Euclidean-branch} with the exact Schur
test, or Corollary~\ref{cor:V18-Lorentzian-branch} with
\eqref{eq:V18-causal-small-gain}.  If the diagonal and off-diagonal writers
are incident with the same finite common-action Hessian after all quotient
and boundary rows pass, then
\begin{equation}
 H_*:=\mathbb H
 \label{eq:V18-V17-preconditioner}
\end{equation}
is the graph preconditioner required in V17.  Its inverse bound is either the
kernel-free Euclidean bound or
\eqref{eq:V18-coupled-inverse-bound}.  Substitution in the V17 jet-radius,
center, and adjacent-defect estimates produces the $C^{2,\alpha}$ scheduler
and V16 Einstein-residual bounds.
\end{corollary}

\begin{proof}
The branch theorems give the two diagonal graph inverses.
Theorem~\ref{thm:V18-graph-Schur} gives the coupled inverse.  Same-action
incidence identifies it with the derivative used by the V17 finite jet card.
The remaining claims are direct substitution into the V17 estimates.
\end{proof}

\begin{firewall}[Diagonal solvability does not imply coupled solvability]
\label{fire:V18-diagonal-not-coupled}
Even when $L_g=L_m=I$, taking $A=D=I$ gives
$S_m=0$ and the mixed kernel $\{(-y,y)\}$.  Separate gravitational and matter
floors therefore cannot replace the mixed Schur row.  Conversely, failure of
the small-gain inequality is inconclusive when $S_m$ remains invertible; the
exact Schur spectrum, rather than the upper product bound, decides that
branch.
\end{firewall}

% -----------------------------------------------------------------------------
\subsection{Finite certificates for symbols, symmetrizers, and the mixed Schur row}
\label{subsec:V18-finite-certificates}
% -----------------------------------------------------------------------------

\begin{lemma}[Lipschitz finite-net singular floor]
\label{lem:V18-symbol-net}
Let $\mathcal K$ be a compact metric space, let
$p:\mathcal K\to\mathcal B(E,F)$ be operator-Lipschitz with constant $L_p$,
and let $\{y_1,\ldots,y_N\}$ be a $\delta$-net.  If
\begin{equation}
 s_{\min}(p(y_j))\ge\mu
 \quad(1\le j\le N),
 \label{eq:V18-net-sampled-floor}
\end{equation}
then
\begin{equation}
 \boxed{
 s_{\min}(p(y))\ge\mu-L_p\delta
 \quad(y\in\mathcal K).}
 \label{eq:V18-uniform-symbol-floor}
\end{equation}
Thus $\mu>L_p\delta$ certifies a uniform floor.  Failure at a sampled point
returns its least right singular vector.
\end{lemma}

\begin{proof}
Choose $y_j$ with $d(y,y_j)\le\delta$.  For every unit $v$,
\[
 \|p(y)v\|
 \ge\|p(y_j)v\|-\|(p(y)-p(y_j))v\|
 \ge\mu-L_p\delta.
\]
Take the infimum over $v$.
\end{proof}

\begin{corollary}[Finite Euclidean symbol and Lorentzian symmetrizer cards]
\label{cor:V18-two-net-cards}
On the Euclidean branch, apply Lemma~\ref{lem:V18-symbol-net} to the normalized
principal symbol on the compact unit cosphere of the cylinder.  A positive
margin $\mu-L_p\delta$ proves uniform interior ellipticity.

On the Lorentzian first-order branch, apply it to
$(A^0)^{1/2}$, or use the Hermitian eigenvalue perturbation inequality directly.
If the sampled least eigenvalues of $A^0$ are at least $\mu$ and
$A^0$ is operator-Lipschitz with constant $L_0$, then
\begin{equation}
 A^0(t,x)\succeq(\mu-L_0\delta)I.
 \label{eq:V18-finite-symmetrizer-floor}
\end{equation}
These interior cards do not decide a complementing Euclidean boundary symbol
or a Lorentzian boundary flux; those require their own boundary-covector
tests.
\end{corollary}

\begin{proof}
The Euclidean assertion is the lemma.  For the symmetrizer,
Weyl's eigenvalue inequality gives
\[
 \lambda_{\min}(A^0(y))
 \ge\lambda_{\min}(A^0(y_j))
    -\|A^0(y)-A^0(y_j)\|
 \ge\mu-L_0\delta.
\]
\end{proof}

Let $\widehat A_N:\mathscr X_m\to\mathscr X_g$ and
$\widehat D_N:\mathscr X_g\to\mathscr X_m$ be finite-rank approximations to
the preconditioned off-diagonal maps $A,D$.  Put
\begin{equation}
 \widehat S_N:=I-\widehat D_N\widehat A_N,
 \qquad
 \sigma_N:=s_{\min}(\widehat S_N).
 \label{eq:V18-finite-Schur}
\end{equation}
The identity on the orthogonal complement makes $\sigma_N$ a finite matrix
calculation together with the value one.

\begin{theorem}[Finite-to-infinite Schur certificate]
\label{thm:V18-finite-Schur}
Assume $\sigma_N>0$ and define
\begin{equation}
 e_N
 :=
 \|D-\widehat D_N\|\,\|A\|
 +\|\widehat D_N\|\,\|A-\widehat A_N\|.
 \label{eq:V18-Schur-tail}
\end{equation}
Then
\begin{equation}
 \|S_m-\widehat S_N\|\le e_N.
 \label{eq:V18-Schur-perturbation}
\end{equation}
If $e_N<\sigma_N$, the exact infinite Schur operator is invertible and
\begin{equation}
 \boxed{
 \|S_m^{-1}\|
 \le\frac1{\sigma_N-e_N}.}
 \label{eq:V18-infinite-Schur-bound}
\end{equation}
If $\sigma_N$ is small, its least singular vectors give finite mixed
gravity--matter soft modes.  If $e_N$ is not controlled below the floor, the
finite card is inconclusive and the two off-diagonal tail operators in
\eqref{eq:V18-Schur-tail} are retained as the missing physical rows.
\end{theorem}

\begin{proof}
Expand
\[
 DA-\widehat D_N\widehat A_N
 =(D-\widehat D_N)A
  +\widehat D_N(A-\widehat A_N)
\]
and take norms to obtain
\eqref{eq:V18-Schur-perturbation}.  Since
$\|\widehat S_N^{-1}\|=\sigma_N^{-1}$,
\[
 \|\widehat S_N^{-1}(S_m-\widehat S_N)\|
 \le e_N/\sigma_N<1.
\]
The Neumann perturbation lemma gives invertibility and
\[
 \|S_m^{-1}\|
 \le
 \frac{\|\widehat S_N^{-1}\|}
 {1-\|\widehat S_N^{-1}\|\|S_m-\widehat S_N\|}
 \le\frac1{\sigma_N-e_N}.
\]
The witness statements are the singular-value decomposition and the explicit
two terms in the perturbation expansion.
\end{proof}

\begin{assessment}[Order of a physical V18 graph card]
\label{ass:V18-card-order}
A physical graph card must be assembled in the following order:
\begin{enumerate}[label=\textup{(\arabic*)},leftmargin=2.8em]
\item freeze the common action, field variables, physical quotient, and
      boundary or Cauchy data;
\item compute the reduced gravitational and matter principal symbols and
      their same-action incidence;
\item choose exactly one Euclidean elliptic or Lorentzian causal realization;
\item certify its interior symbol or symmetrizer, then its boundary or initial
      data and constraint propagation;
\item invert the diagonal graph blocks or return their kernel, cokernel, or
      quasimode;
\item form the same-action off-diagonal maps and test the exact Schur operator;
\item transport the resulting inverse and its tails into the V17 shrinking
      jet card.
\end{enumerate}
Changing signature, boundary conditions, quotient, or field variables between
these steps changes the operator and invalidates the inverse certificate.
\end{assessment}

% -----------------------------------------------------------------------------
\subsection{Removing the gauge--Ward circularity}
\label{subsec:V18-gauge-Ward}
% -----------------------------------------------------------------------------

At finite cutoff neither the reduced Einstein residual nor the represented
Ward divergence is identically zero a priori.  They must enter the subsidiary
gauge equation as sources.

\begin{proposition}[Forced gauge propagation identity]
\label{prop:V18-forced-gauge}
Define
\begin{align}
 \mathscr E_{\mu\nu}
 &:=
 G_{\mu\nu}+\Lambda g_{\mu\nu}-T_{\mu\nu}
 -\nabla_{(\mu}C_{\nu)}
 +\frac12g_{\mu\nu}\nabla_\alpha C^\alpha,
 \label{eq:V18-reduced-residual}\\
 J_\nu&:=\nabla^\mu T_{\mu\nu}.
 \label{eq:V18-Ward-defect}
\end{align}
Then the exact identity
\begin{equation}
 \boxed{
 \nabla^\alpha\nabla_\alpha C_\nu
 +\operatorname{Ric}_{\nu\mu}C^\mu
 =
 -2\left(
 J_\nu+\nabla^\mu\mathscr E_{\mu\nu}
 \right)}
 \label{eq:V18-forced-gauge}
\end{equation}
holds.  On a Lorentzian causal branch its first-order reduction and
Theorem~\ref{thm:V18-hyperbolic-energy} give
\begin{align}
 \|C\|_{\rm gauge}
 \le C_T\bigl(
 &\|C|_{\Sigma_0}\|+\|\nabla_nC|_{\Sigma_0}\|
 \nonumber\\
 &+\|J+\operatorname{div}\mathscr E\|_{\rm source}
 \bigr).
 \label{eq:V18-gauge-Ward-debit}
\end{align}
Thus gauge closure follows from vanishing initial gauge data, the Ward defect,
and the differentiated reduced-equation defect.  None may be assumed from
either of the other two.
\end{proposition}

\begin{proof}
Take the divergence of \eqref{eq:V18-reduced-residual}.  The contracted
Bianchi identity and metric compatibility give
\[
 \nabla^\mu\mathscr E_{\mu\nu}
 =
 -J_\nu
 -\frac12\left(
 \nabla^\alpha\nabla_\alpha C_\nu
 +\operatorname{Ric}_{\nu\mu}C^\mu
 \right),
\]
using the gauge-term calculation in
Corollary~\ref{cor:V18-gauge-propagation}.  Rearrangement proves
\eqref{eq:V18-forced-gauge}.  Apply the causal energy estimate to the
first-order reduction of this vector wave equation to obtain
\eqref{eq:V18-gauge-Ward-debit}.
\end{proof}

\begin{firewall}[Differentiated residuals need their own graph norm]
\label{fire:V18-divergence-residual}
Smallness of $\mathscr E$ in a weak dual norm does not imply smallness of
$\operatorname{div}\mathscr E$ in the source norm of
\eqref{eq:V18-gauge-Ward-debit}.  The V17 action card must either control one
additional derivative of the reduced residual or formulate the subsidiary
system distributionally in a matching negative Sobolev norm.  This is a
domain issue, not an algebraic consequence of the Einstein residual.
\end{firewall}

% -----------------------------------------------------------------------------
\subsection{Executable V18 graph-inverse card}
\label{subsec:V18-card}
% -----------------------------------------------------------------------------

\begin{definition}[Branch-typed graph-inverse card]
\label{def:V18-card}
A V18 card contains:
\begin{enumerate}[label=\textup{(I\arabic*)},leftmargin=3em]
\item the common-action gravitational, matter, and mixed linearization blocks
      after every physical quotient;
\item the direct high-frequency principal-symbol response and its incidence
      with the declared reduced differential operator;
\item a branch tag selecting either the Euclidean elliptic boundary
      realization or the Lorentzian source-plus-Cauchy-data realization;
\item on the Euclidean branch, the interior symbol net, complementing boundary
      symbol, compact-remainder constant, kernel, and adjoint kernel;
\item on the Lorentzian branch, the symmetrizer net, coefficient bounds,
      Cauchy data, boundary flux if present, and the forced gauge--Ward debit;
\item the two diagonal inverse bounds, the off-diagonal maps $A,D$, the exact
      finite Schur singular values, and both tail terms in
      \eqref{eq:V18-Schur-tail};
\item adjacent-cutoff domain, coefficient, inverse, constraint, and boundary
      transport debits required by the V17 graph card;
\item one held-out high-frequency polarization, one held-out boundary or
      initial datum, and one held-out mixed gravity--matter command.
\end{enumerate}
\end{definition}

\begin{theorem}[V18 inverse or obstruction alternative]
\label{thm:V18-card-alternative}
For a populated card on a declared branch:
\begin{enumerate}[label=\textup{(\roman*)},leftmargin=2.8em]
\item if its symbol or symmetrizer, domain, boundary or Cauchy, kernel,
      cokernel, constraint, same-action incidence, exact Schur, and tail rows
      pass with positive margins, the complete coupled operator is the
      bounded graph isomorphism $H_*$ required by V17;
\item otherwise the first failed row returns a principal-symbol vector,
      boundary-symbol vector, Euclidean quasimode, causal energy or boundary
      flux defect, initial constraint, Ward-forced gauge source, diagonal
      kernel or cokernel, mixed Schur singular vector, off-diagonal tail, or
      held-out command witness.
\end{enumerate}
On the pass branch, the inverse norm is bounded by
\eqref{eq:V18-explicit-Euclidean-inverse} and
\eqref{eq:V18-coupled-inverse-bound} in the Euclidean small-remainder branch,
or by \eqref{eq:V18-hyperbolic-energy} and
\eqref{eq:V18-coupled-inverse-bound} in the causal small-gain branch.  The
exact Schur branch may use \eqref{eq:V18-infinite-Schur-bound} instead of the
small-gain estimate.
\end{theorem}

\begin{proof}
The reduced principal symbol is
Theorem~\ref{thm:V18-reduced-symbol}.  The Euclidean graph alternative is
Theorem~\ref{thm:V18-compact-remainder} and
Corollary~\ref{cor:V18-Euclidean-branch}.  The Lorentzian graph alternative is
Theorem~\ref{thm:V18-hyperbolic-energy}, with constraint and gauge errors
controlled by
Corollary~\ref{cor:V18-constraint-propagation} and
Proposition~\ref{prop:V18-forced-gauge}.  The coupled inverse and its exact
mixed witnesses are Theorem~\ref{thm:V18-graph-Schur}; finite tail promotion
is Theorem~\ref{thm:V18-finite-Schur}.  These rows exhaust
Definition~\ref{def:V18-card}.  Corollary~\ref{cor:V18-to-V17} supplies the
handoff to V17.
\end{proof}

% -----------------------------------------------------------------------------
\subsection{V18 theorem ledger and next symbol-incidence acquisition}
\label{subsec:V18-ledger}
% -----------------------------------------------------------------------------

\paragraph{New conclusions proved in V18.}
Theorem~\ref{thm:V18-reduced-symbol} computes the gauge-reduced Einstein
principal symbol and separates the Euclidean elliptic and Lorentzian normally
hyperbolic realizations.  Corollary~\ref{cor:V18-gauge-propagation} identifies
the vector operator which must remove residual diffeomorphism gauge.
Theorem~\ref{thm:V18-compact-remainder} proves a kernel/cokernel-complete
Euclidean graph inverse from a compact-remainder estimate.
Theorem~\ref{thm:V18-hyperbolic-energy} proves the causal $H^r$ graph estimate
and Corollary~\ref{cor:V18-constraint-propagation} retains every constraint
debit.  Theorem~\ref{thm:V18-graph-Schur} gives the exact coupled
gravity--matter inverse, its small-gain bound, and its mixed kernel and
cokernel.  Lemma~\ref{lem:V18-symbol-net} and
Theorem~\ref{thm:V18-finite-Schur} turn the interior and mixed rows into
finite-net plus tail certificates.  Proposition
\ref{prop:V18-forced-gauge} removes the circular use of Ward conservation in
gauge propagation.

\paragraph{Imported results not repeated.}
The common-action and relation-support command squares, signed finite Hessian
quotient, graph-root transport, metric-stress secant, Palatini limit, OS Ward
stress reconstruction, and conditional Einstein handoff are imported from
the attached sources and V1--V17.  V18 supplies the missing branch-typed
differential inverse and coupled graph assembly.

\paragraph{Authoritative physical status.}
The supplied manuscripts do not print the principal and lower-order
coefficient matrices of one completely gauge-fixed interacting RPESM
Palatini--Standard-Model linearization.  They provide neither a complementing
Euclidean boundary symbol nor a Lorentzian symmetrizer and constraint-
preserving boundary/Cauchy card for that selected history.  They also do not
populate the off-diagonal graph maps $A,D$ or the tails
\eqref{eq:V18-Schur-tail}.  Therefore V18 does not assert that either branch
passes for the physical regulator, and the full Standard Model--Einstein
continuum is not claimed.

\paragraph{Next acquisition.}
The next iteration should extract the differential symbol from the occurring
common-action response rather than declare a continuum operator.  A suitable
test uses localized oscillatory jet commands
\[
 u_{\lambda,x,\xi,v}
 =\chi_x\,e^{i\lambda\phi_\xi}v
\]
and compares the rescaled action response
$\lambda^{-\nu}e^{-i\lambda\phi_\xi}D\mathcal F_nu_{\lambda,x,\xi,v}$ with the
declared principal symbol.  The required theorem must:
\begin{enumerate}[label=\textup{(\arabic*)},leftmargin=2.8em]
\item give an explicit remainder in terms of coefficient derivatives, cutoff
      scale, localization radius, and $\lambda$;
\item preserve the returned-word, physical-connection, and determinant-line
      transports of the V1/V30 command square;
\item reconstruct the full symbol matrix by finite polarization on a
      source-complete field and response frame;
\item separate a genuine symbol zero from boundary leakage and an unresolved
      high-frequency tail;
\item feed the resulting uniform symbol or symmetrizer margin into the V18
      branch card.
\end{enumerate}
This moves upstream of the unavailable coefficient table while still testing
the action which actually occurs on the regulator cylinder.

\begin{firewall}[V18 does not choose the analytic branch]
\label{fire:V18-final-scope}
The Euclidean and Lorentzian theorems are rigorous conditional alternatives.
Analytic continuation of an OS theory does not by itself select compatible
Lorentzian Cauchy data, and a Lorentzian scheduler does not supply an elliptic
OS boundary problem.  The physical line, orientation, common action, and
domain data must identify which branch is being used before any inverse is
inserted into V17.
\end{firewall}

% =============================================================================
% END APPENDED VERSION V18 MATERIAL
% =============================================================================



% =============================================================================
% BEGIN APPENDED VERSION V19 MATERIAL
% =============================================================================

\clearpage
\section{Version V19: oscillatory action commands and principal-symbol incidence}
\label{sec:V19-symbol-incidence}

Version V18 proved that a graph inverse can be certified only after the
reduced differential symbol, analytic branch, domain, diagonal blocks, and
mixed Schur row are tied to the same physical common action.  The attached
programmes do not print that differential coefficient table.  Declaring the
standard continuum Einstein--Standard-Model symbol would therefore skip the
occurrence question.

This version extracts the symbol from the occurring linearized response.
Localized oscillatory field commands separate the highest differential order
from all lower orders.  Their matched, demodulated response converges to the
principal symbol with an explicit localization/frequency error.  A centered
lattice calculation then supplies the additional resolvability condition
$h\lambda\to0$.  Finite field, response, base-point, and covector frames turn
the construction into a positive incidence card for the V18 interior symbol.
Relation-support, physical-connection, determinant-line, and nonlinear
command-radius errors remain separate inputs from V1, V17, and the attached
V30 programme.

% -----------------------------------------------------------------------------
\subsection{Localized packets recover a differential principal symbol}
\label{subsec:V19-packet}
% -----------------------------------------------------------------------------

Work in a coordinate ball $U\Subset\mathbb R^d$ and trivialize finite
dimensional Hermitian bundles $E,F$.  Put $D_j=-i\partial_j$ and consider
\begin{equation}
 L=\sum_{|\gamma|\le m}A_\gamma(x)D^\gamma,
 \qquad
 p_m(x,\xi)=\sum_{|\gamma|=m}A_\gamma(x)\xi^\gamma.
 \label{eq:V19-differential-operator}
\end{equation}
Assume the principal coefficients are uniformly continuous and all lower
coefficients are bounded.  Define
\begin{equation}
 \omega_m(r)
 :=
 \sup_{\substack{x,y\in U\\|x-y|\le r}}
 \|p_m(x,\cdot)-p_m(y,\cdot)\|_
 {C(\mathbb S^{d-1};\mathcal B(E,F))}.
 \label{eq:V19-symbol-modulus}
\end{equation}

Choose a real $\chi\in C_c^\infty(B_1(0))$ with
$\|\chi\|_{L^2}=1$.  If $B_r(x_0)\Subset U$, set
\begin{equation}
 \chi_{r,x_0}(x)
 :=r^{-d/2}\chi\left(\frac{x-x_0}{r}\right)
 \label{eq:V19-rescaled-cutoff}
\end{equation}
and, for $|\xi|=1$ and $v\in E$, define
\begin{equation}
 u_{\lambda,r,x_0,\xi,v}(x)
 :=
 \chi_{r,x_0}(x)
 e^{i\lambda\xi\cdot(x-x_0)}v.
 \label{eq:V19-oscillatory-packet}
\end{equation}
It has $L^2$ norm $\|v\|$.

The matched symbol read is the vector in $F$
\begin{align}
 \mathfrak S_{\lambda,r}^{L}(x_0,\xi)v
 &:=
 \int_U
 \chi_{r,x_0}(x)e^{-i\lambda\xi\cdot(x-x_0)}
 \lambda^{-m}
 L u_{\lambda,r,x_0,\xi,v}(x)\,\dd x.
 \label{eq:V19-matched-symbol-read}
\end{align}

\begin{theorem}[Quantitative oscillatory extraction]
\label{thm:V19-symbol-extraction}
There is a constant $C_{L,\chi,m}$, independent of
$x_0,\xi,v,r,\lambda$, such that whenever
\[
 0<r\le1,\qquad \lambda\ge1,\qquad\lambda r\ge1,
\]
one has
\begin{equation}
 \boxed{
 \|\mathfrak S_{\lambda,r}^{L}(x_0,\xi)v
    -p_m(x_0,\xi)v\|_F
 \le
 \left[
 \omega_m(r)
 +C_{L,\chi,m}
  \left(\frac1{\lambda r}+\frac1\lambda\right)
 \right]\|v\|_E.}
 \label{eq:V19-symbol-extraction}
\end{equation}
If the principal coefficients are Lipschitz with constant $L_p$, then
$\omega_m(r)\le L_pr$.  The choice $r=\lambda^{-1/2}$ gives an
$O(\lambda^{-1/2})$ symbol error.

The same estimate holds with a smooth positive density in place of $\dd x$
after normalizing $\chi_{r,x_0}$ in that density; its Lipschitz variation is
absorbed into the $O(r)$ term.
\end{theorem}

\begin{proof}
The multi-index Leibniz formula gives
\begin{align}
 D^\gamma\left(
 e^{i\lambda\xi\cdot(x-x_0)}\chi_{r,x_0}v
 \right)
 =
 e^{i\lambda\xi\cdot(x-x_0)}
 \sum_{\beta\le\gamma}
 {\gamma\choose\beta}
 (\lambda\xi)^{\gamma-\beta}
 D^\beta\chi_{r,x_0}\,v.
 \label{eq:V19-Leibniz}
\end{align}
For $|\gamma|=m$ and $\beta=0$, division by $\lambda^m$ gives
$A_\gamma(x)\xi^\gamma\chi_{r,x_0}v$.  Its matched integral differs from
$p_m(x_0,\xi)v$ by at most
$\omega_m(r)\|v\|$, because
$\int\chi_{r,x_0}^2=1$ and the packet is supported in $B_r(x_0)$.

For a principal term with $|\beta|\ge1$,
\[
 \lambda^{-m}
 \|(\lambda\xi)^{\gamma-\beta}
   D^\beta\chi_{r,x_0}\|_{L^2}
 \le
 C_{\gamma,\beta,\chi}(\lambda r)^{-|\beta|}.
\]
Since $\lambda r\ge1$, the finite sum of these terms is bounded by
$C(\lambda r)^{-1}$.  For $|\gamma|\le m-1$, the term with $\beta=0$ is
bounded after division by $\lambda^m$ by $C\lambda^{-1}$.  Terms with
$\beta\ne0$ have the additional factor $(\lambda r)^{-|\beta|}$ and obey the
same or a smaller bound.  Multiplication by the matched cutoff has
$L^2$ norm one, so Cauchy--Schwarz gives
\eqref{eq:V19-symbol-extraction}.  Lipschitz continuity gives
$\omega_m(r)\le L_pr$.

For a smooth positive density, use its normalized square-root cutoff.  On the
radius-$r$ ball the density and its normalization differ from their central
values by $O(r)$, and derivatives of that factor enter the already estimated
cutoff terms.
\end{proof}

\begin{corollary}[Real command implementation]
\label{cor:V19-real-commands}
If $L$ has real coefficients, the complex packet
\eqref{eq:V19-oscillatory-packet} is reconstructed from the two real commands
\[
 \chi_{r,x_0}\cos(\lambda\xi\cdot(x-x_0))v,
 \qquad
 \chi_{r,x_0}\sin(\lambda\xi\cdot(x-x_0))v.
\]
Applying the real response to both and combining them as real and imaginary
parts gives exactly \eqref{eq:V19-matched-symbol-read}.  Thus complexification
introduces no unphysical command.
\end{corollary}

\begin{proof}
The complex packet is the first real command plus $i$ times the second.
The complex-linear extension of a real linear response has the same
decomposition, and demodulation is a known post-processing multiplication.
\end{proof}

% -----------------------------------------------------------------------------
\subsection{A resolved lattice diagonal}
\label{subsec:V19-lattice}
% -----------------------------------------------------------------------------

On a Cartesian mesh of spacing $h$, define the centered realization of
$D_j=-i\partial_j$ by
\begin{equation}
 D_{h,j}f(x)
 :=
 \frac{f(x+he_j)-f(x-he_j)}{2ih}.
 \label{eq:V19-centered-difference}
\end{equation}

\begin{lemma}[Exact centered-difference dispersion]
\label{lem:V19-centered-symbol}
For a plane wave,
\begin{equation}
 D_{h,j}e^{i\lambda\xi\cdot x}
 =
 \frac{\sin(h\lambda\xi_j)}{h}
 e^{i\lambda\xi\cdot x}.
 \label{eq:V19-centered-plane-wave}
\end{equation}
If $|h\lambda\xi_j|\le1$, then
\begin{equation}
 \left|
 \frac{\sin(h\lambda\xi_j)}{h\lambda}-\xi_j
 \right|
 \le\frac16(h\lambda)^2|\xi_j|^3.
 \label{eq:V19-sine-error}
\end{equation}
For a homogeneous matrix polynomial $p_m(x,\xi)$ of degree $m$, uniformly
bounded on the unit base/covector set,
\begin{equation}
 \left\|
 p_m\left(x,
 \frac{\sin(h\lambda\xi)}{h\lambda}\right)
 -p_m(x,\xi)
 \right\|
 \le C_{p,m}(h\lambda)^2.
 \label{eq:V19-polynomial-dispersion}
\end{equation}
\end{lemma}

\begin{proof}
Substitution gives \eqref{eq:V19-centered-plane-wave}.
The inequality $|\sin z-z|\le|z|^3/6$ gives
\eqref{eq:V19-sine-error}.  The normalized discrete covector stays in a
fixed compact ball when $|h\lambda|\le1$.  A degree-$m$ polynomial is
Lipschitz there, so the componentwise estimate proves
\eqref{eq:V19-polynomial-dispersion}.
\end{proof}

Let $L_h$ be obtained from \eqref{eq:V19-differential-operator} by replacing
each $D_j$ with $D_{h,j}$, sampling coefficients on the lattice, and using
the normalized lattice $L^2$ pairing in
\eqref{eq:V19-matched-symbol-read}.  Restrict the packet to a ball whose
two-mesh collar is contained in $U$.

\begin{theorem}[Localized centered-lattice symbol read]
\label{thm:V19-lattice-symbol}
There is a constant $C$ depending on the coefficient bounds, $m$, and
finitely many derivatives of $\chi$ such that, when
\[
 h\lambda\le1,\qquad
 \lambda r\ge1,\qquad
 h/r\le c_0<1,
\]
the discrete matched read satisfies
\begin{align}
 &\|\mathfrak S_{\lambda,r}^{L_h}(x_0,\xi)v
       -p_m(x_0,\xi)v\|
 \nonumber\\
 &\quad\le
 \left[
 \omega_m(r)
 +C\left(
 \frac1{\lambda r}+\frac1\lambda
 +(h\lambda)^2+\frac hr
 \right)
 \right]\|v\|.
 \label{eq:V19-lattice-symbol-bound}
\end{align}
Hence the scale conditions
\begin{equation}
 r_n\longrightarrow0,\qquad
 \lambda_nr_n\longrightarrow\infty,\qquad
 h_n\lambda_n\longrightarrow0,\qquad
 h_n/r_n\longrightarrow0
 \label{eq:V19-four-scale-conditions}
\end{equation}
recover the principal symbol on every interior compact set.

If the principal coefficients are Lipschitz, the explicit choice
\begin{equation}
 \boxed{
 \lambda_n=h_n^{-4/5},
 \qquad
 r_n=h_n^{2/5}}
 \label{eq:V19-balanced-scales}
\end{equation}
makes the right side of
\eqref{eq:V19-lattice-symbol-bound} $O(h_n^{2/5})$.
\end{theorem}

\begin{proof}
For a smooth scalar function, Taylor's theorem gives
\[
 \|D_{h,j}f-D_jf\|_{L^2}
 \le C h^2\|\partial_j^3f\|_{L^2}
\]
on the interior mesh collar.  Derivatives of the packet satisfy
\[
 \|\partial^\ell u_{\lambda,r}\|_{L^2}
 \le C_\ell(\lambda+r^{-1})^\ell\|v\|.
\]
Replace the centered factors in each order-$m$ monomial one at a time.
After division by $\lambda^m$, the sum of the resulting third-derivative
remainders is bounded by $C(h\lambda)^2\|v\|$, because
$\lambda r\ge1$.  Equivalently, its leading plane-wave term is exactly
\eqref{eq:V19-polynomial-dispersion}; derivatives falling on the cutoff have
the same or a smaller bound.  The continuous cutoff and lower-order errors
are those in Theorem~\ref{thm:V19-symbol-extraction}.

The normalized lattice pairing is a Riemann sum on a function varying at
envelope scale $r$ after demodulation.  Its quadrature and normalization error
is bounded by $C(h/r)\|v\|$.  Adding the four contributions proves
\eqref{eq:V19-lattice-symbol-bound}.

Conditions \eqref{eq:V19-four-scale-conditions} make every term vanish.
For \eqref{eq:V19-balanced-scales},
\[
 r_n=h_n^{2/5},\quad
 (\lambda_nr_n)^{-1}=h_n^{2/5},\quad
 (h_n\lambda_n)^2=h_n^{2/5},\quad
 h_n/r_n=h_n^{3/5},\quad
 \lambda_n^{-1}=h_n^{4/5}.
\]
This proves the asserted rate.
\end{proof}

\begin{firewall}[A root scheduler is not automatically a centered stencil]
\label{fire:V19-stencil}
Theorem~\ref{thm:V19-lattice-symbol} is an exact consistency result for the
declared centered local realization.  The RPESM scheduler uses directed root
rates, chronological faces, and transported relation supports.  Its actual
high-frequency response must first be shown incident with this stencil, or
with another explicitly analyzed stencil.  Positive rates and the correct
low-frequency bracket do not determine the full lattice dispersion at
$h\lambda>0$.
\end{firewall}


\subsection{The actual common-action Hessian versus a declared local symbol}
\label{subsec:V19-action-symbol-incidence}

Fix transported orthonormal frames for the input bundle $E$ and output
bundle $F$ over an interior chart.  Let
\[
 J_n=D\mathcal F_n(z_n):X_n\longrightarrow Y_n
\]
be the aligned derivative of the \emph{actual} finite common-action residual
at an action root $z_n$, as constructed in V17.  Let $L_{n,h_n}$ be the
centered-lattice realization of a declared order-$m$ local operator $L_n$,
with principal symbol $p_{n,m}$.  No equality between $J_n$ and
$L_{n,h_n}$ is assumed.

For a unit fibre vector $v$, write $u_n=u_{\lambda_n,r_n,x,\xi,v}$ and
define the normalized packet-incidence debit
\begin{equation}
 \eta_n^{\mathrm{loc}}(x,\xi,v)
 :=
 \lambda_n^{-m}
 \|(J_n-L_{n,h_n})u_n\|_{2,h_n}.
 \label{eq:V19-locality-debit}
\end{equation}
The actual experiment is normally a symmetric action secant rather than
direct access to $J_n$:
\begin{equation}
 \mathcal D_{n,\varepsilon_n}u
 :=
 \frac{\mathcal F_n(z_n+\varepsilon_nu)
             -\mathcal F_n(z_n-\varepsilon_nu)}
      {2\varepsilon_n}.
 \label{eq:V19-action-secant}
\end{equation}
After all relation, support, chronological, and line transports have been
applied, denote its response by
$\widetilde{\mathcal D}_{n,\varepsilon_n}u$.  We keep the two additional
debits
\begin{align}
 \eta_n^{\mathrm{jet}}(u)
 &:=
 \lambda_n^{-m}
 \|\mathcal D_{n,\varepsilon_n}u-J_nu\|_{2,h_n},
 \label{eq:V19-jet-debit}\\
 \eta_n^{\mathrm{route}}(u)
 &:=
 \lambda_n^{-m}
 \|\widetilde{\mathcal D}_{n,\varepsilon_n}u
          -\mathcal D_{n,\varepsilon_n}u\|_{2,h_n}.
 \label{eq:V19-route-debit}
\end{align}
Thus locality, nonlinear differentiation, and alignment are separate
incidences and cannot cancel in the bookkeeping.

Let $\mathcal M_{n,x,\xi}:Y_n\to F_x$ be the normalized cutoff,
demodulation, and fibre-transport pairing used in
\eqref{eq:V19-matched-symbol-read}.  Its normalization is chosen so that
$\|\mathcal M_{n,x,\xi}\|\le1$.  The observed symbol column is
\begin{equation}
 \widehat p_n(x,\xi)v
 :=
 \mathcal M_{n,x,\xi}
 \bigl(\lambda_n^{-m}
       \widetilde{\mathcal D}_{n,\varepsilon_n}u_n\bigr).
 \label{eq:V19-observed-symbol-column}
\end{equation}

\begin{theorem}[Common-action packet-incidence estimate]
\label{thm:V19-action-symbol-incidence}
Assume the hypotheses of Theorem~\ref{thm:V19-lattice-symbol}, uniformly in
$n$, and put
\[
 \eta_n^{\mathrm{coef}}
 :=
 \sup_{(x,\xi)\in K\times S^{d-1}}
 \|p_{n,m}(x,\xi)-p_m(x,\xi)\|.
\]
For every $x\in K$, $|\xi|=1$, and unit $v\in E_x$,
\begin{align}
 \|\widehat p_n(x,\xi)v-p_m(x,\xi)v\|
 \le{}&
 \omega_{n,m}(r_n)
 +C\left((\lambda_nr_n)^{-1}+\lambda_n^{-1}
              +(h_n\lambda_n)^2+h_n/r_n\right)
 \nonumber\\
 &+\eta_n^{\mathrm{coef}}
  +\eta_n^{\mathrm{loc}}(x,\xi,v)
  +\eta_n^{\mathrm{jet}}(u_n)
  +\eta_n^{\mathrm{route}}(u_n).
 \label{eq:V19-total-symbol-debit}
\end{align}
Consequently, uniform convergence of every term on the right proves
uniform convergence of the observed common-action symbol to $p_m$ on the
compact unit cosphere.
\end{theorem}

\begin{proof}
Insert, in order,
$\mathcal D_{n,\varepsilon_n}$, $J_n$, $L_{n,h_n}$,
$p_{n,m}(x,\xi)$, and $p_m(x,\xi)$ into
\eqref{eq:V19-observed-symbol-column}.  The first three differences are
bounded by \eqref{eq:V19-route-debit},
\eqref{eq:V19-jet-debit}, and \eqref{eq:V19-locality-debit}, because
$\mathcal M_{n,x,\xi}$ is a contraction.  The lattice packet term is
bounded by Theorem~\ref{thm:V19-lattice-symbol}; the last term is
$\eta_n^{\mathrm{coef}}$.  The triangle inequality proves
\eqref{eq:V19-total-symbol-debit}.  Uniform convergence is immediate.
\end{proof}

\begin{lemma}[A sufficient amplitude law for the central jet]
\label{lem:V19-amplitude-law}
Suppose that, on the relevant ball in $X_n$, the third derivative obeys
\[
 \|D^3\mathcal F_n(w)[a,b,c]\|_{Y_n}
 \le M_3\|a\|_{X_n}\|b\|_{X_n}\|c\|_{X_n}
\]
uniformly in $n$, and that a unit-$L^2$ packet satisfies
$\|u_n\|_{X_n}\le C\lambda_n^s$.  Then
\begin{equation}
 \eta_n^{\mathrm{jet}}(u_n)
 \le \frac{M_3C^3}{6}\,
       \varepsilon_n^2\lambda_n^{3s-m}.
 \label{eq:V19-amplitude-law}
\end{equation}
In particular,
$\varepsilon_n^2\lambda_n^{3s-m}\to0$ is a sufficient noncircular
amplitude condition.
\end{lemma}

\begin{proof}
Taylor's formula with integral remainder, applied at
$z_n\pm\varepsilon_nu_n$, cancels the even terms in the symmetric
difference and gives
\[
 \|\mathcal D_{n,\varepsilon_n}u_n-J_nu_n\|_{Y_n}
 \le \frac{M_3}{6}\varepsilon_n^2\|u_n\|_{X_n}^3.
\]
Multiply by $\lambda_n^{-m}$ and use the packet bound.  This is the
high-frequency specialization of the central-jet control in
Theorem~\ref{thm:V17-jet-inverse}.
\end{proof}


For clarity, the modulus used in
\eqref{eq:V19-total-symbol-debit} is
\[
 \omega_{n,m}(r)
 :=
 \max_{|\gamma|=m}\ 
 \sup_{\substack{x,y\in K_r\\ |x-y|\le r}}
 \|A_{n,\gamma}(x)-A_{n,\gamma}(y)\|,
\]
where $K_r$ is the $r$-neighbourhood of $K$ in the chart.

\subsection{Exterior leakage and exact finite-stencil reconstruction}
\label{subsec:V19-exterior-leakage}

Assume that $\chi$ is supported in the unit ball.  If the longest offset
appearing in the order-$m$ centered stencil has length at most $c_mh_n$,
let $Q^{\mathrm{out}}_{n,x,r}$ be multiplication by the complement of
$B(x,r+c_mh_n)$.  Define
\begin{equation}
 \ell_n^{\mathrm{ext}}(x,\xi,v)
 :=
 \lambda_n^{-m}
 \|Q^{\mathrm{out}}_{n,x,r_n}
       \widetilde{\mathcal D}_{n,\varepsilon_n}u_n\|_{2,h_n}.
 \label{eq:V19-exterior-leakage}
\end{equation}

\begin{proposition}[Exterior leakage is a locality obstruction]
\label{prop:V19-exterior-obstruction}
For the preceding packet and local comparator,
\begin{equation}
 \ell_n^{\mathrm{ext}}(x,\xi,v)
 \le
 \eta_n^{\mathrm{loc}}(x,\xi,v)
 +\eta_n^{\mathrm{jet}}(u_n)
 +\eta_n^{\mathrm{route}}(u_n).
 \label{eq:V19-exterior-obstruction}
\end{equation}
Therefore a positive lower limit of the left side, after the jet and route
debits vanish, disproves incidence with every comparator having the stated
stencil radius.
\end{proposition}

\begin{proof}
A finite-difference operator only samples points in its stencil, so
\[
 \operatorname{supp}(L_{n,h_n}u_n)
 \subset B(x,r_n+c_mh_n).
\]
Consequently $Q^{\mathrm{out}}_{n,x,r_n}L_{n,h_n}u_n=0$.
Insert $L_{n,h_n}u_n$, $J_nu_n$, and
$\mathcal D_{n,\varepsilon_n}u_n$ in
\eqref{eq:V19-exterior-leakage}; the exterior projection is a contraction.
The triangle inequality gives \eqref{eq:V19-exterior-obstruction}.
\end{proof}

The converse requires more than a few packets.  The following exact
finite-dimensional statement records the required audit.

\begin{lemma}[Impulse locality reconstructs the actual stencil]
\label{lem:V19-impulse-stencil}
Let $\Lambda_h$ be a finite lattice and let
$J_h:\ell^2(\Lambda_h;E)\to\ell^2(\Lambda_h;F)$ be linear.  Write its block
kernel as
\[
 (J_hu)(x)=\sum_{y\in\Lambda_h}K_h(x,y)u(y).
\]
The following are equivalent for an integer $c\ge0$:
\begin{enumerate}
 \item $K_h(x,y)=0$ whenever $|x-y|>ch$;
 \item for every $S\subset\Lambda_h$,
 \[
  \operatorname{supp}(J_hu)\subset
  N_{ch}(S)\quad\hbox{whenever}\quad
  \operatorname{supp}u\subset S;
 \]
 \item the support statement in item\/ {\rm(2)} holds for every
 fibre-basis impulse supported at a single lattice point.
\end{enumerate}
When these conditions hold, the actual frozen lattice symbol is the finite
matrix sum
\begin{equation}
 \sigma_{J_h}(x,\theta)
 =
 \sum_{\substack{z\in\mathbb Z^d\\ |z|\le c}}
 K_h(x,x+hz)e^{i\theta\cdot z},
 \label{eq:V19-actual-stencil-symbol}
\end{equation}
with absent lattice points omitted.
\end{lemma}

\begin{proof}
If item\/ {\rm(1)} holds, every nonzero summand has
$x\in N_{ch}(\operatorname{supp}u)$, proving item\/ {\rm(2)}.
Item\/ {\rm(2)} trivially implies item\/ {\rm(3)}.  Conversely, apply
item\/ {\rm(3)} to each basis vector $e_a$ at $y$.  Its value at $x$ is
$K_h(x,y)e_a$, so that value vanishes for $|x-y|>ch$.  Varying $a$ proves
item\/ {\rm(1)}.  Substitution of
$u(x+hz)=e^{i\theta\cdot z}u(x)$ gives
\eqref{eq:V19-actual-stencil-symbol}.
\end{proof}

\begin{firewall}[Packet leakage is necessary, not sufficient]
\label{fire:V19-packet-locality}
Small values of \eqref{eq:V19-exterior-leakage} for a finite packet family
do not prove that $J_n$ is differential or even local.  Exact locality on
the finite level requires the impulse audit in
Lemma~\ref{lem:V19-impulse-stencil}, or an equivalent uniform kernel
estimate.  Passage from a local stencil to an order-$m$ continuum operator
also requires the moment conditions proved next.
\end{firewall}


\begin{theorem}[Finite-stencil moment certificate]
\label{thm:V19-stencil-moments}
Suppose the impulse audit gives, at an interior lattice point,
\begin{equation}
 (J_hu)(x)
 =
 h^{-m}\sum_{z\in Z_c}B_{h,z}(x)u(x+hz),
 \qquad
 Z_c:=\{z\in\mathbb Z^d:|z|\le c\}.
 \label{eq:V19-stencil-form}
\end{equation}
Set
\[
 M_{h,\alpha}(x):=\sum_{z\in Z_c}B_{h,z}(x)z^\alpha,
 \qquad
 p^{\mathrm{mom}}_{h,m}(x,\xi)
 :=
 \sum_{|\alpha|=m}\frac{i^m}{\alpha!}
 M_{h,\alpha}(x)\xi^\alpha .
\]
For $\theta=h\lambda\le1$ and $|\xi|=1$, the frozen normalized response
\[
 q_{h,\lambda}(x,\xi)
 :=
 \lambda^{-m}h^{-m}
 \sum_{z\in Z_c}B_{h,z}(x)e^{ih\lambda\xi\cdot z}
\]
satisfies
\begin{align}
 \|q_{h,\lambda}(x,\xi)-p^{\mathrm{mom}}_{h,m}(x,\xi)\|
 \le{}&
 \rho_h^{<m}(x,\theta)+C_m\theta\,\mathcal B_{h,m+1}(x),
 \label{eq:V19-moment-bound}\\
 \rho_h^{<m}(x,\theta)
 :={}&
 \sum_{k=0}^{m-1}\theta^{k-m}
 \sum_{|\alpha|=k}\frac{\|M_{h,\alpha}(x)\|}{\alpha!},
 \label{eq:V19-lower-moment-debit}\\
 \mathcal B_{h,m+1}(x)
 :={}&
 \sum_{z\in Z_c}\|B_{h,z}(x)\|\,|z|^{m+1}.
 \label{eq:V19-stencil-size}
\end{align}
In particular, the exact cancellations
\begin{equation}
 M_{h,\alpha}(x)=0
 \quad(0\le|\alpha|<m)
 \label{eq:V19-moment-cancellations}
\end{equation}
remove every divergent lower-order term.  If
$\mathcal B_{h,m+1}$ is uniformly bounded and
$p^{\mathrm{mom}}_{h,m}\to p_m$ uniformly, then any cofinal scale with
$h\lambda\to0$ recovers $p_m$ from the frozen actual stencil.
\end{theorem}

\begin{proof}
Put $\theta=h\lambda$ and expand
\[
 e^{i\theta\xi\cdot z}
 =
 \sum_{k=0}^{m}\frac{i^k\theta^k}{k!}(\xi\cdot z)^k
 +R_{m+1}(\theta,\xi,z).
\]
The multinomial identity gives
\[
 \frac{(\xi\cdot z)^k}{k!}
 =
 \sum_{|\alpha|=k}\frac{\xi^\alpha z^\alpha}{\alpha!}.
\]
After multiplication by $\theta^{-m}B_{h,z}$ and summation in $z$, the
terms with $k<m$ are bounded by
\eqref{eq:V19-lower-moment-debit}, while the $k=m$ term is exactly
$p^{\mathrm{mom}}_{h,m}$.  Taylor's remainder obeys
\[
 |R_{m+1}(\theta,\xi,z)|
 \le C_m\theta^{m+1}|z|^{m+1}
\]
because $Z_c$ is finite and $\theta\le1$.  Division by $\theta^m$ gives
the last term in \eqref{eq:V19-moment-bound}.  The final assertion follows
from \eqref{eq:V19-moment-cancellations} and the triangle inequality.
\end{proof}

\begin{corollary}[Centered first and second differences pass the moments]
\label{cor:V19-centered-moments}
For
\[
 D_{h,j}u(x)=\frac{u(x+he_j)-u(x-he_j)}{2ih},
\]
the zeroth moment vanishes and the first moment gives $\xi_j$.  For
\[
 \delta^{[2]}_{h,j}u(x)
 =\frac{2u(x)-u(x+he_j)-u(x-he_j)}{h^2},
\]
the moments of orders zero and one vanish and the second moment gives
$\xi_j^2$.  Thus these stencils have the $D=-i\partial$ principal symbols
asserted in Theorem~\ref{thm:V19-lattice-symbol}.
\end{corollary}

\begin{proof}
For the first stencil, $B_{e_j}=1/(2i)$,
$B_{-e_j}=-1/(2i)$, so $M_0=0$ and $iM_{e_j}=1$.
For the second, $B_0=2$, $B_{e_j}=B_{-e_j}=-1$; hence
$M_0=M_{e_j}=0$ and $(i^2/2)M_{2e_j}=1$.
\end{proof}

\begin{firewall}[Low-frequency brackets do not supply the moment row]
\label{fire:V19-moment-row}
The renewal bracket can identify a first formal derivative while
\eqref{eq:V19-lower-moment-debit} still diverges for the actual
common-action Hessian.  The impulse kernel, the cancellations
\eqref{eq:V19-moment-cancellations}, and the top moment are therefore
physical incidence data.  They cannot be imported from the declared
continuum operator.
\end{firewall}


\subsection{Finite matrix reads and a certified symbol net}
\label{subsec:V19-symbol-matrix-net}

Let $q=\operatorname{rank}E$ and let $(e_a)_{a=1}^q$ be an orthonormal
input frame.  Running \eqref{eq:V19-observed-symbol-column} for every
$e_a$ reconstructs the matrix $\widehat p_n(x,\xi)$.  If a target reduced
symbol $p^{\mathrm{tar}}(x,\xi)$ has been declared, define the observable
positive incidence matrix
\begin{equation}
 \mathbb I_n^{\mathrm{sym}}(x,\xi)
 :=
 \bigl(\widehat p_n-p^{\mathrm{tar}}\bigr)^*
 \bigl(\widehat p_n-p^{\mathrm{tar}}\bigr)
 \succeq0.
 \label{eq:V19-symbol-incidence-Gram}
\end{equation}

\begin{lemma}[Columns, Gram incidence, and operator error]
\label{lem:V19-column-Gram}
If
\[
 \|(\widehat p-p)e_a\|\le\epsilon_{\mathrm{col}}
 \quad(1\le a\le q),
\]
then
\begin{equation}
 \|\widehat p-p\|_{\mathrm{op}}
 \le \sqrt q\,\epsilon_{\mathrm{col}}.
 \label{eq:V19-column-operator}
\end{equation}
Moreover,
\begin{equation}
 \|\widehat p-p^{\mathrm{tar}}\|_{\mathrm{op}}
 =
 \sqrt{\lambda_{\max}(\mathbb I^{\mathrm{sym}})}.
 \label{eq:V19-Gram-operator}
\end{equation}
Both conclusions remain valid for rectangular symbols $E\to F$.
\end{lemma}

\begin{proof}
For $v=\sum_av_ae_a$ with $\|v\|=1$,
\[
 \|(\widehat p-p)v\|
 \le\epsilon_{\mathrm{col}}\sum_a|v_a|
 \le\sqrt q\,\epsilon_{\mathrm{col}}
\]
by Cauchy--Schwarz.  The second identity is the standard equality
$\|A\|_{\mathrm{op}}^2=\lambda_{\max}(A^*A)$.
\end{proof}

\begin{theorem}[Observed finite-net singular floor]
\label{thm:V19-observed-net-floor}
For a linear map $A:E\to F$, write
\[
 \sigma_{\min}(A):=\inf_{\|v\|=1}\|Av\|.
\]
Let $\mathcal N_x$ and $\mathcal N_\xi$ be respectively a
$\delta_x$-net of a compact base set $K$ and a $\delta_\xi$-net of the
unit covectors.  Suppose the actual principal symbol $p$ satisfies
\begin{equation}
 \|p(x,\xi)-p(y,\zeta)\|_{\mathrm{op}}
 \le L_x|x-y|+L_\xi|\xi-\zeta|.
 \label{eq:V19-symbol-Lipschitz}
\end{equation}
Assume at every net pair that
\[
 \sigma_{\min}(\widehat p_n(x_\nu,\xi_\mu))
 \ge\widehat\sigma,
 \qquad
 \|\widehat p_n(x_\nu,\xi_\mu)-p(x_\nu,\xi_\mu)\|
 \le\epsilon_{\mathrm{op}}.
\]
Then
\begin{equation}
 \boxed{
 \inf_{\substack{x\in K\\|\xi|=1}}
 \sigma_{\min}(p(x,\xi))
 \ge
 \widehat\sigma-\epsilon_{\mathrm{op}}
 -L_x\delta_x-L_\xi\delta_\xi.}
 \label{eq:V19-observed-net-floor}
\end{equation}
For columnwise packet bounds one may take
$\epsilon_{\mathrm{op}}=\sqrt q\,\epsilon_{\mathrm{col}}$.
A positive right side is therefore a finite certificate of injective
interior symbol.  For a rectangular symbol, surjective coverage requires
the corresponding lower-singular-value test for the adjoint, unless the
dimensions already force equivalence.
\end{theorem}

\begin{proof}
Choose a nearest net pair.  The smallest singular value is
$1$-Lipschitz in operator norm, so
\[
 \sigma_{\min}(p(x,\xi))
 \ge\sigma_{\min}(\widehat p_n(x_\nu,\xi_\mu))
 -\|\widehat p_n-p\|
 -\|p(x,\xi)-p(x_\nu,\xi_\mu)\|.
\]
Apply \eqref{eq:V19-symbol-Lipschitz}.  The columnwise assertion is
Lemma~\ref{lem:V19-column-Gram}.  Applying the same argument to $p^*$
gives the stated coverage test.
\end{proof}

If the Gram incidence data instead obey
$\lambda_{\max}(\mathbb I_n^{\mathrm{sym}})\le\tau_n^2$ on the net, then
Lemma~\ref{lem:V19-column-Gram} gives
$\|\widehat p_n-p^{\mathrm{tar}}\|\le\tau_n$ there.  Combining this
observable target error with \eqref{eq:V19-total-symbol-debit} and the two
Lipschitz tails gives a uniform incidence certificate.  This is the
action-Hessian input needed by Lemma~\ref{lem:V18-symbol-net}; it does not
supply the boundary complementing or causal-energy rows of V18.


\subsection{The Euclidean and Lorentzian reads must remain separate}
\label{subsec:V19-branch-symbol-reads}

Let $N=\operatorname{rank}S^2T^*M$ and recall from
Theorem~\ref{thm:V18-reduced-symbol} that the gauge-reduced gravitational
target is
\begin{equation}
 p_g(x,\xi)
 =
 \frac12G^{\alpha\beta}(x)\xi_\alpha\xi_\beta I_N,
 \qquad G^{\alpha\beta}=g^{\alpha\beta}.
 \label{eq:V19-Einstein-target-symbol}
\end{equation}

\begin{proposition}[Branch-typed interpretation of the packet matrix]
\label{prop:V19-branch-interpretation}
The following are distinct certificates.
\begin{enumerate}[label=\textup{(\alph*)},leftmargin=3em]
\item If $G$ is Riemannian and
$\lambda_{\min}(G)\ge\gamma>0$, while
\[
 \sup_{x,|\xi|=1}
 \|\widehat p_n(x,\xi)-p_g(x,\xi)\|
 \le e_n,
\]
then
\[
 \inf_{x,|\xi|=1}\sigma_{\min}(\widehat p_n(x,\xi))
 \ge\frac{\gamma}{2}-e_n.
\]
A positive finite-net version of this bound certifies the interior
strong-ellipticity row used in Corollary~\ref{cor:V18-Euclidean-branch}.

\item If $G$ is Lorentzian, no positive singular floor can hold on the full
unit covector sphere, because \eqref{eq:V19-Einstein-target-symbol} vanishes
on every nonzero null covector.  Instead, suppose the symbol is already
certified to be a homogeneous quadratic matrix polynomial.  Evaluate it at
\[
 \mathcal Q
 :=
 \{e_\alpha\}_{\alpha=1}^d
 \cup
 \{(e_\alpha+e_\beta)/\sqrt2:\alpha<\beta\}.
 \label{eq:V19-polarization-set}
\]
If all these values are scalar matrices, write them as
$q_\alpha I_N$ and $q_{\alpha\beta}I_N$.  Then
\begin{align}
 G^{\alpha\alpha}&=2q_\alpha,
 \label{eq:V19-diagonal-metric-read}\\
 G^{\alpha\beta}
 &=2q_{\alpha\beta}-q_\alpha-q_\beta
 \qquad(\alpha<\beta)
 \label{eq:V19-cross-metric-read}
\end{align}
and these values determine \eqref{eq:V19-Einstein-target-symbol} for every
$\xi$.  If each scalar read has error at most $\epsilon$, the reconstructed
matrix has spectral error at most $4d\epsilon$.

If the reconstructed symmetric matrix $\widehat G$ has one negative and
$d-1$ positive eigenvalues and
\[
 \min_j|\lambda_j(\widehat G)|>4d\epsilon,
 \label{eq:V19-signature-margin}
\]
then the true $G$ has the same Lorentzian signature.
This identifies the normally hyperbolic second-order symbol.  The positive
first-order symmetrizer, its coefficient bounds, the Cauchy data, and the
constraint residual must still satisfy
Theorem~\ref{thm:V18-hyperbolic-energy} and
Corollary~\ref{cor:V18-constraint-propagation}.
\end{enumerate}
\end{proposition}

\begin{proof}
For part (a),
$\sigma_{\min}(p_g(x,\xi))
=\frac12G^{\alpha\beta}\xi_\alpha\xi_\beta\ge\gamma/2$.
The singular-value perturbation inequality gives the result.

For part (b), every Lorentzian quadratic form has nonzero null covectors,
so the asserted obstruction is immediate.  A symmetric matrix-valued
quadratic polynomial is determined by its diagonal values and its
polarizations.  Substitution of \eqref{eq:V19-Einstein-target-symbol} at
the points in \eqref{eq:V19-polarization-set} gives
\eqref{eq:V19-diagonal-metric-read} and
\eqref{eq:V19-cross-metric-read}.  A diagonal coefficient has error at
most $2\epsilon$ and an off-diagonal coefficient at most $4\epsilon$.
Thus every entry error is at most $4\epsilon$, and the spectral norm is at
most $d$ times the largest entry.  Weyl's eigenvalue inequality and
\eqref{eq:V19-signature-margin} preserve the number of positive and negative
eigenvalues.  The last statement follows because normal hyperbolicity and
symmetric-hyperbolic energy positivity are different hypotheses, as recorded
in Firewall~\ref{fire:V18-symmetrizer}.
\end{proof}

\begin{firewall}[A refined Lorentzian net approaches the null cone]
\label{fire:V19-null-cone}
A coarse covector net can display a spurious positive minimum by missing the
null cone.  Refining such a singular-value test drives the minimum to zero.
The polarization and signature test in
Proposition~\ref{prop:V19-branch-interpretation}, followed by the independent
symmetrizer test \eqref{eq:V18-symmetrizer-floor}, is the branch-correct
certificate.
\end{firewall}


\subsection{Cofinal interior-symbol handoff}
\label{subsec:V19-cofinal-handoff}

\begin{theorem}[Cofinal common-action principal-symbol certificate]
\label{thm:V19-cofinal-symbol}
Let $h_n\downarrow0$, take the balanced scales
\eqref{eq:V19-balanced-scales}, and choose the symmetric command amplitudes
so that
\begin{equation}
 \varepsilon_n^2\lambda_n^{3s-m}\longrightarrow0.
 \label{eq:V19-cofinal-amplitude}
\end{equation}
Assume the V17 action roots and common transports exist and their route
debits vanish uniformly.  Suppose either of the following two incidence
routes is completed on every compact interior chart.

\begin{enumerate}[label=\textup{(R\arabic*)},leftmargin=3em]
\item There are local comparators $L_{n,h_n}$ for which
$\omega_{n,m}(r_n)$, $\eta_n^{\mathrm{coef}}$, and
$\eta_n^{\mathrm{loc}}$ in
\eqref{eq:V19-total-symbol-debit} tend uniformly to zero.

\item The impulse audit of Lemma~\ref{lem:V19-impulse-stencil} reconstructs
the actual Hessian stencil; its stencil sizes
\eqref{eq:V19-stencil-size} are uniformly bounded, its lower-moment debits
\eqref{eq:V19-lower-moment-debit} vanish at
$\theta_n=h_n\lambda_n$, and its top-moment polynomials converge uniformly
to $p_m$.
\end{enumerate}

Let the base and covector net radii tend to zero and retain all column
matrices and Gram incidences.  Then the actual aligned common-action
linearization has principal symbol $p_m$ in the packet sense, uniformly on
compact unit cospheres.

If the target is \eqref{eq:V19-Einstein-target-symbol}, the following
branch conclusions are valid.

\begin{enumerate}[label=\textup{(B\arabic*)},leftmargin=3em]
\item On the Riemannian branch, a positive right side in
\eqref{eq:V19-observed-net-floor} proves the interior strong-ellipticity
hypothesis of the V18 Euclidean graph construction.

\item On the Lorentzian branch, the polarization residuals tending to zero
and a uniform signature margin in \eqref{eq:V19-signature-margin} prove the
normally hyperbolic reduced Einstein principal part.  A causal graph inverse
follows only after the independent first-order symmetrizer, Cauchy or
boundary, coefficient, and constraint hypotheses of
Corollary~\ref{cor:V18-Lorentzian-branch} are verified.
\end{enumerate}
\end{theorem}

\begin{proof}
Under route {\rm(R1)}, Lemma~\ref{lem:V19-amplitude-law},
the balanced-scale calculation in
Theorem~\ref{thm:V19-lattice-symbol}, and the assumed route and coefficient
limits make every term in
\eqref{eq:V19-total-symbol-debit} tend uniformly to zero.  Under route
{\rm(R2)}, Theorem~\ref{thm:V19-stencil-moments} gives the same conclusion
for the frozen actual stencil; localization contributes the cutoff and
coefficient moduli already estimated in
Theorem~\ref{thm:V19-symbol-extraction}.  The central-jet and route estimates
then transfer that conclusion to the measured common-action response.

Lemma~\ref{lem:V19-column-Gram} promotes convergence of all frame columns to
operator-norm convergence.  The shrinking nets and the Lipschitz tail in
Theorem~\ref{thm:V19-observed-net-floor} promote the finite reads to the
compact unit cosphere.  The two branch conclusions are precisely
Proposition~\ref{prop:V19-branch-interpretation}, followed respectively by
Corollaries~\ref{cor:V18-Euclidean-branch} and
\ref{cor:V18-Lorentzian-branch}.  The latter citation is conditional because
an interior second-order symbol does not prove an energy estimate.
\end{proof}

\begin{corollary}[What V19 supplies to the V18 Schur row]
\label{cor:V19-Schur-handoff}
When Theorem~\ref{thm:V19-cofinal-symbol} is applied to every diagonal and
mixed field block of the same common-action Hessian, it certifies the
interior high-frequency part of those blocks.  The coupled graph inverse of
Theorem~\ref{thm:V18-graph-Schur} still requires:
\begin{enumerate}
 \item the selected Euclidean boundary or Lorentzian causal graph domains;
 \item control of the compact or finite-frequency remainders;
 \item the finite-to-infinite Schur tail
       \eqref{eq:V18-Schur-tail}; and
 \item the forced gauge--Ward debit
       \eqref{eq:V18-gauge-Ward-debit}.
\end{enumerate}
Thus an interior symbol pass is a genuine advance but is not a full
Standard-Model--Einstein continuum proof.
\end{corollary}

\begin{proof}
Theorem~\ref{thm:V19-cofinal-symbol} identifies the principal action
incidence.  The hypotheses listed are exactly the remaining domain,
remainder, coupling, and gauge inputs in
Theorems~\ref{thm:V18-compact-remainder},
\ref{thm:V18-graph-Schur}, and
Proposition~\ref{prop:V18-forced-gauge}.
\end{proof}


\subsection{Executable V19 acquisition card}
\label{subsec:V19-card}

For each cutoff $n$, chart, base-net point, covector-net point, field-frame
column, and retained branch, the finite card is the following.

\begin{enumerate}[label=\textup{(\arabic*)},leftmargin=3em]
\item Retain the V17 action root $z_n$, the exact residual
$\mathcal F_n(z_n)$, the relation/support/chronological transports, and the
source and target norms.

\item Issue both real commands
\[
 \chi_{r_n,x}\cos(\lambda_n\xi\cdot(\,\cdot-x\,))e_a,
 \qquad
 \chi_{r_n,x}\sin(\lambda_n\xi\cdot(\,\cdot-x\,))e_a
\]
at $\lambda_n=h_n^{-4/5}$, $r_n=h_n^{2/5}$, and an amplitude satisfying
\eqref{eq:V19-cofinal-amplitude}.  Retain the two endpoint residuals used in
every symmetric secant.

\item Recombine the real responses into
\eqref{eq:V19-observed-symbol-column}; retain every column, the full matrix,
the Gram incidence \eqref{eq:V19-symbol-incidence-Gram}, the exterior
leakage \eqref{eq:V19-exterior-leakage}, and the jet and route debits.

\item Complete either the comparator residual
\eqref{eq:V19-locality-debit} or the fibre-basis impulse audit of
Lemma~\ref{lem:V19-impulse-stencil}.  On the impulse route retain the block
kernel, every lower moment in
\eqref{eq:V19-lower-moment-debit}, the top moment polynomial, and the stencil
size \eqref{eq:V19-stencil-size}.

\item Run disjoint certification and held-out base/covector nets.  Record
the measured Lipschitz constants or verified upper bounds used for the
finite-to-continuum tail.  A held-out failure enlarges the net and invalidates
the current pass.

\item On the Euclidean branch, report the margin in
\eqref{eq:V19-observed-net-floor}.  On the Lorentzian branch, report the
polarized coefficient matrix, its signature margin, the scalar-matrix
residual, and the separate first-order $A^0$ floor and ceiling from
\eqref{eq:V18-symmetrizer-floor}.
\end{enumerate}

A machine-checkable packet-pass number is
\begin{align}
 E_n^{\mathrm{pkt}}
 :={}&
 \sup_{x,\xi,a}\bigg[
 \omega_{n,m}(r_n)
 +C\left((\lambda_nr_n)^{-1}+\lambda_n^{-1}
 +(h_n\lambda_n)^2+h_n/r_n\right)
 \nonumber\\
 &\hspace{7em}
 +\eta_n^{\mathrm{coef}}
 +\eta_n^{\mathrm{loc}}(x,\xi,e_a)
 +\eta_n^{\mathrm{jet}}(u_{n,a})
 +\eta_n^{\mathrm{route}}(u_{n,a})
 \bigg].
 \label{eq:V19-machine-packet-bound}
\end{align}
The comparator route passes only if $E_n^{\mathrm{pkt}}\to0$.  On the
impulse route, replace the coefficient and locality terms by the
top-moment mismatch, \eqref{eq:V19-lower-moment-debit}, and
$C_mh_n\lambda_n\mathcal B_{h_n,m+1}$.

Every failure has a retained witness:
\begin{align}
 W_n^{\mathrm{ext}}&=(x,\xi,e_a,
   Q^{\mathrm{out}}\widetilde{\mathcal D}u),\\
 W_n^{\mathrm{mom}}&=(x,\alpha,M_{h_n,\alpha}),\\
 W_n^{\mathrm{sym}}&=(x,\xi,e_a,
   (\widehat p_n-p^{\mathrm{tar}})e_a),\\
 W_n^{\mathrm{floor}}&=(x,\xi,v,
   \|\widehat p_n v\|),\\
 W_n^{\mathrm{sig}}&=(x,\widehat G,
   \operatorname{spec}\widehat G),\\
 W_n^{\mathrm{hyp}}&=(t,x,w,
   \langle A^0w,w\rangle).
 \label{eq:V19-witness-family}
\end{align}
The first nonvanishing witness decides where to go upstream: transport,
nonlinear amplitude, locality, stencil consistency, target incidence,
signature, or causal reduction.

\subsection{V19 status and next upstream acquisition}
\label{subsec:V19-status}

V19 proves the following new analytic statements.

\begin{enumerate}
\item A localized oscillatory command recovers the principal symbol with an
explicit cutoff and coefficient error.

\item A centered lattice command has the additional dispersion and quadrature
errors in \eqref{eq:V19-lattice-symbol-bound}; the balanced scale
\eqref{eq:V19-balanced-scales} makes all analytic terms vanish at rate
$O(h_n^{2/5})$ for Lipschitz principal coefficients.

\item The action-secant, route, local-comparator, and coefficient errors enter
the exact additive bound \eqref{eq:V19-total-symbol-debit}.

\item An impulse support audit reconstructs the actual finite stencil, and
Theorem~\ref{thm:V19-stencil-moments} states the precise lower-moment
cancellations needed to prevent divergent false continuum terms.

\item Matrix columns, positive Gram incidences, and finite nets give a
quantitative global symbol certificate.  The Lorentzian test uses
polarization, signature, and an independent energy symmetrizer rather than
a false positive floor across the null cone.
\end{enumerate}

The attached grand tensor manuscript and the four programme notes do not
contain the oscillatory common-action responses, impulse kernels, moment
tables, or held-out net matrices required by the V19 card.  Therefore the
physical principal-symbol incidence remains an acquisition row; V19 does
not mark it passed and does not claim the full Standard-Model--Einstein
continuum.

After this interior row, the next upstream bottleneck is the graph domain at
a boundary.  The Euclidean acquisition must use tangential packets
\[
 e^{i\lambda x'\cdot\xi'}e^{-\lambda\kappa x_d}v
\]
to reconstruct the principal boundary symbol and test the complementing
condition uniformly in $(x',\xi')$.  The Lorentzian acquisition must
reconstruct the boundary flux $U^*A^\nu U$, certify a maximal
nonpositive incoming subspace of the correct dimension, and verify the
constraint-preserving boundary map.  Those boundary-layer incidences, plus
the compact-remainder and Schur-tail rows retained in
Corollary~\ref{cor:V19-Schur-handoff}, are the mathematically prior target
for the next appended version.

\begin{firewall}[V19 closes an analytic implication, not a physical pass]
\label{fire:V19-status}
Theorems~\ref{thm:V19-action-symbol-incidence},
\ref{thm:V19-stencil-moments}, and
\ref{thm:V19-cofinal-symbol} say exactly which finite common-action data
would prove the interior continuum symbol and what error margin they must
satisfy.  They do not manufacture those data.  Until the card is populated,
neither the Euclidean Fredholm branch nor the Lorentzian causal branch can
be promoted to the complete coupled continuum.
\end{firewall}

% END APPENDED VERSION V19 MATERIAL



% =============================================================================
% BEGIN APPENDED VERSION V20 MATERIAL
% =============================================================================

\clearpage
\section{Version V20: boundary symbols, complementing modes, and causal flux}
\label{sec:V20-boundary-incidence}

V19 identified the occurring interior principal symbol by oscillatory
common-action commands.  That symbol does not choose an operator domain.
V18 therefore left two genuinely different boundary rows: a complementing
elliptic boundary map on the Euclidean branch, or a maximally dissipative
and constraint-preserving flux subspace on a Lorentzian initial-boundary
branch.

The attached programmes mention boundary fluxes and Friedrichs--Dirichlet
realizations, but they do not print the stable normal-mode map, its smallest
singular value, the causal flux matrix, or the induced incoming constraint
map for the assembled Einstein--Standard-Model Hessian.  V20 constructs
finite tests for exactly those objects.  It first proves that localized
tangential oscillations with a rescaled decaying normal profile recover an
elliptic boundary symbol.  It then proves the algebraic and energy estimates
for a Lorentzian maximal nonnegative flux graph.  The two tests are retained
as different cards.

% -----------------------------------------------------------------------------
\subsection{The stable normal space and the complementing map}
\label{subsec:V20-stable-space}
% -----------------------------------------------------------------------------

Use boundary coordinates $(y,x_d)\in\mathbb R^{d-1}\times[0,\infty)$, with
the domain on the side $x_d>0$, and put $D=-i\partial$.  Let
\begin{align}
 P&=\sum_{|\alpha|+k\le2m}
 A_{\alpha k}(y,x_d)D_y^\alpha D_d^k,
 \label{eq:V20-elliptic-operator}\\
 B_j&=\sum_{|\alpha|+k\le m_j}
 C_{j,\alpha k}(y)D_y^\alpha D_d^k,
 \qquad 0\le m_j\le2m-1.
 \label{eq:V20-boundary-operators}
\end{align}
At $x_0=(y_0,0)$ their principal frozen polynomials are
\begin{align}
 p^0_{x_0}(\xi',\zeta)
 &:=
 \sum_{|\alpha|+k=2m}
 A_{\alpha k}(x_0)(\xi')^\alpha\zeta^k,
 \label{eq:V20-interior-normal-pencil}\\
 b^0_{j,x_0}(\xi',\zeta)
 &:=
 \sum_{|\alpha|+k=m_j}
 C_{j,\alpha k}(y_0)(\xi')^\alpha\zeta^k.
 \label{eq:V20-boundary-pencil}
\end{align}

Assume proper ellipticity and a uniform normal-root separation on the
compact unit tangential cosphere under consideration.  Define the stable
normal space
\begin{equation}
 \mathcal E^s(x_0,\xi')
 :=
 \left\{
 V:\ p^0_{x_0}(\xi',D_t)V=0,\quad
 \sum_{k=0}^{2m}\|e^{ct}D_t^kV\|_{L^2(\mathbb R_+)}<\infty
 \right\}
 \label{eq:V20-stable-normal-space}
\end{equation}
for a fixed uniform $c>0$.  Repeated roots are included through their
polynomial-exponential generalized modes.  Give this finite-dimensional
space the initial-jet norm
\begin{equation}
 \|V\|_s^2
 :=
 \sum_{k=0}^{2m-1}\|D_t^kV(0)\|^2.
 \label{eq:V20-stable-jet-norm}
\end{equation}
The principal boundary map is
\begin{equation}
 \mathcal B^s(x_0,\xi')V
 :=
 \left(
 b^0_{j,x_0}(\xi',D_t)V(0)
 \right)_{j=1}^J
 \in\mathcal G:=\bigoplus_{j=1}^J G_j .
 \label{eq:V20-complementing-map}
\end{equation}

\begin{definition}[Quantitative complementing margin]
\label{def:V20-complementing-margin}
When $\dim\mathcal E^s=\dim\mathcal G$, put
\begin{equation}
 \beta_\partial
 :=
 \inf_{\substack{x_0\in\partial M,\ |\xi'|=1\\
                  V\in\mathcal E^s(x_0,\xi'),\ \|V\|_s=1}}
 \|\mathcal B^s(x_0,\xi')V\|_{\mathcal G}.
 \label{eq:V20-Lopatinski-floor}
\end{equation}
The boundary operators satisfy the uniform complementing condition exactly
when $\beta_\partial>0$.
\end{definition}

\begin{lemma}[The floor is the inverse boundary-mode bound]
\label{lem:V20-floor-inverse}
At a fixed $(x_0,\xi')$, assume
$\dim\mathcal E^s=\dim\mathcal G$.  Then
$\mathcal B^s$ is bijective if and only if its pointwise version of
\eqref{eq:V20-Lopatinski-floor} is positive.  In that case
\begin{equation}
 \|(\mathcal B^s)^{-1}\|_{\mathcal G\to\mathcal E^s}
 \le\beta_\partial^{-1}.
 \label{eq:V20-boundary-inverse}
\end{equation}
If the stable spaces and maps vary continuously on the compact unit
tangential cosphere, failure of the uniform floor produces an attained
nonzero decaying witness $V_*$ with
\begin{equation}
 p^0_{x_*}(\xi'_*,D_t)V_*=0,\qquad
 \mathcal B^s(x_*,\xi'_*)V_*=0.
 \label{eq:V20-complementing-witness}
\end{equation}
\end{lemma}

\begin{proof}
In equal finite dimensions, a linear map is bijective exactly when it is
bounded below.  The definition of the least singular value then gives
\eqref{eq:V20-boundary-inverse}.  For the final assertion, use local
continuous frames of the stable bundle, pass to a convergent subsequence of
points and unit coefficient vectors, and use compactness.  Continuity
preserves both the stable equation and the vanishing boundary value.
\end{proof}


% -----------------------------------------------------------------------------
\subsection{Boundary-layer commands recover the principal boundary map}
\label{subsec:V20-boundary-packets}
% -----------------------------------------------------------------------------

Choose $\chi\in C_c^\infty(B_1^{d-1})$ with $\|\chi\|_2=1$ and a normal
cutoff $\psi\in C_c^\infty([0,\rho_0))$ equal to one on
$[0,\rho_0/2]$.  For $V\in\mathcal E^s(x_0,\xi')$, define
\begin{equation}
 U_{\lambda,r,V}(y,x_d)
 :=
 \lambda^{1/2}\chi_{r,y_0}(y)
 e^{i\lambda\xi'\cdot(y-y_0)}
 \psi(x_d)V(\lambda x_d).
 \label{eq:V20-boundary-layer-packet}
\end{equation}
Apart from an exponentially small normal-cutoff tail, its $L^2$ norm is
$\|V\|_{L^2(\mathbb R_+)}$.

The matched interior profile and matched boundary reads are
\begin{align}
 \mathfrak P_{\lambda,r}V(t)
 &:=
 \lambda^{-2m-1/2}
 \int
 \chi_{r,y_0}(y)e^{-i\lambda\xi'\cdot(y-y_0)}
 (PU_{\lambda,r,V})(y,t/\lambda)\,\dd y,
 \label{eq:V20-matched-interior-profile}\\
 \mathfrak B_{j,\lambda,r}V
 &:=
 \lambda^{-m_j-1/2}
 \int
 \chi_{r,y_0}(y)e^{-i\lambda\xi'\cdot(y-y_0)}
 (B_jU_{\lambda,r,V})(y,0)\,\dd y.
 \label{eq:V20-matched-boundary-read}
\end{align}

\begin{theorem}[Quantitative boundary-symbol extraction]
\label{thm:V20-boundary-extraction}
Assume the principal coefficients of $P$ and $B_j$ are Lipschitz near the
boundary, the lower-order coefficients are bounded, and the stable normal
spaces have the uniform exponential and initial-jet bounds stated in
\eqref{eq:V20-stable-normal-space}--\eqref{eq:V20-stable-jet-norm}.  There is
a constant $C$, uniform on the compact unit tangential cosphere, such that
for $\lambda\ge1$, $0<r\le1$, and $\lambda r\ge1$,
\begin{align}
 &\left\|
 \mathfrak P_{\lambda,r}V
 -p^0_{x_0}(\xi',D_t)V
 \right\|_{L^2(\mathbb R_+)}
 \le
 C\left(r+\frac1{\lambda r}+\frac1\lambda\right)\|V\|_s,
 \label{eq:V20-interior-profile-error}\\
 &\left(
 \sum_{j=1}^J
 \left\|
 \mathfrak B_{j,\lambda,r}V
 -b^0_{j,x_0}(\xi',D_t)V(0)
 \right\|^2
 \right)^{1/2}
 \le
 C\left(r+\frac1{\lambda r}+\frac1\lambda\right)\|V\|_s.
 \label{eq:V20-boundary-map-error}
\end{align}
For a stable profile the comparison term in
\eqref{eq:V20-interior-profile-error} is zero.  The choice
$r=\lambda^{-1/2}$ therefore makes both errors $O(\lambda^{-1/2})$.
\end{theorem}

\begin{proof}
For a principal monomial with $|\alpha|+k=2m$, tangential Leibniz expansion
gives
\begin{align*}
 D_y^\alpha D_d^kU_{\lambda,r,V}
 =
 \lambda^{k+1/2}e^{i\lambda\xi'\cdot(y-y_0)}
 \sum_{\beta\le\alpha}{\alpha\choose\beta}
 (\lambda\xi')^{\alpha-\beta}
 D_y^\beta\chi_{r,y_0}\,
 D_t^kV(\lambda x_d)
\end{align*}
where $\psi=1$.  After division by $\lambda^{2m+1/2}$, the term
$\beta=0$ is the corresponding monomial in
$p^0_{(y,x_d)}(\xi',D_t)V$.  The matched tangential integral averages it
with weight $\chi_{r,y_0}^2$.

At $x_d=t/\lambda$, Lipschitz continuity bounds the difference from the
frozen coefficient by $C(r+t/\lambda)$.  The uniform exponential bounds on
$V$ imply
\[
 \|(1+t)D_t^kV\|_{L^2_t}\le C\|V\|_s,
\]
so coefficient freezing costs $C(r+\lambda^{-1})\|V\|_s$.  Every term with
$|\beta|\ge1$ contains $(\lambda r)^{-|\beta|}$ and their finite sum is
bounded by $C(\lambda r)^{-1}\|V\|_s$.  A monomial below order $2m$ has at
least one additional factor $\lambda^{-1}$; tangential cutoff derivatives
only decrease it further.  Derivatives of $\psi$ occur at
$t\ge\lambda\rho_0/2$, where exponential decay makes them $O(\lambda^{-N})$
for every fixed $N$.  This proves
\eqref{eq:V20-interior-profile-error}.

At $x_d=0$, the same computation for an order-$m_j$ boundary monomial has
leading scale $\lambda^{m_j+1/2}$.  Coefficient freezing now costs $Cr$,
tangential cutoff derivatives cost $C(\lambda r)^{-1}$, and lower-order
terms cost $C\lambda^{-1}$.  Initial normal derivatives are controlled by
\eqref{eq:V20-stable-jet-norm}.  Summing the finite boundary blocks proves
\eqref{eq:V20-boundary-map-error}.  The stable equation and the stated
choice of $r$ give the last assertions.
\end{proof}

\begin{corollary}[Real boundary commands]
\label{cor:V20-real-boundary-commands}
For real coefficients, the complex tangential carrier in
\eqref{eq:V20-boundary-layer-packet} is exactly reconstructed from cosine
and sine commands with the same real normal profile components.  Hence the
matrix of \eqref{eq:V20-complementing-map} requires only real common-action
secants and complex post-processing.
\end{corollary}

\begin{proof}
Apply the real response separately to the real and imaginary parts, exactly
as in Corollary~\ref{cor:V19-real-commands}.
\end{proof}


% -----------------------------------------------------------------------------
\subsection{The occurring stable space and a finite Lopatinski card}
\label{subsec:V20-finite-Lopatinski}
% -----------------------------------------------------------------------------

The stable profiles must come from the occurring interior pencil, not merely
from the declared continuum pencil.  When the leading normal coefficient is
invertible, rewrite the order-$2m$ normal equation as a first-order jet
system
\begin{equation}
 \partial_tW=\mathcal G(x_0,\xi')W.
 \label{eq:V20-normal-companion}
\end{equation}
Choose a positively oriented contour $\Gamma$ enclosing all eigenvalues
with negative real part and no other spectrum.

\begin{lemma}[Certified transport of the stable normal space]
\label{lem:V20-stable-Riesz}
Let
\[
 M_\Gamma:=
 \sup_{z\in\Gamma}\|(zI-\mathcal G)^{-1}\|,
 \qquad
 \widetilde{\mathcal G}=\mathcal G+E,
 \qquad
 M_\Gamma\|E\|<1.
\]
Then $\Gamma$ remains in the resolvent of
$\widetilde{\mathcal G}$, the two Riesz projections satisfy
\begin{equation}
 \|\widetilde\Pi^s-\Pi^s\|
 \le
 \delta_\Pi
 :=
 \frac{|\Gamma|}{2\pi}
 \frac{M_\Gamma^2\|E\|}{1-M_\Gamma\|E\|},
 \label{eq:V20-Riesz-projection-bound}
\end{equation}
and their ranks agree.  If $\delta_\Pi<1$, then
\[
 T:=\widetilde\Pi^s|_{\operatorname{Ran}\Pi^s}
 :
 \operatorname{Ran}\Pi^s\longrightarrow
 \operatorname{Ran}\widetilde\Pi^s
\]
is an isomorphism with
\begin{equation}
 \|T-I\|\le\delta_\Pi,
 \qquad
 \|T^{-1}\|\le(1-\delta_\Pi)^{-1}.
 \label{eq:V20-stable-identification}
\end{equation}
\end{lemma}

\begin{proof}
The factorization
\[
 zI-\widetilde{\mathcal G}
 =
 (zI-\mathcal G)
 \left[I-(zI-\mathcal G)^{-1}E\right]
\]
and the Neumann series give
\[
 \|(zI-\widetilde{\mathcal G})^{-1}\|
 \le\frac{M_\Gamma}{1-M_\Gamma\|E\|}.
\]
The resolvent identity therefore gives
\[
 \|(zI-\widetilde{\mathcal G})^{-1}
 -(zI-\mathcal G)^{-1}\|
 \le\frac{M_\Gamma^2\|E\|}{1-M_\Gamma\|E\|}.
\]
Integrating around $\Gamma$ proves
\eqref{eq:V20-Riesz-projection-bound}.  The rank of a continuous family of
finite-dimensional projections is locally constant, so the ranks agree.
For $w\in\operatorname{Ran}\Pi^s$,
$\|Tw-w\|\le\delta_\Pi\|w\|$, hence
$\|Tw\|\ge(1-\delta_\Pi)\|w\|$.  Thus $T$ is injective; equal rank makes it
onto, and gives \eqref{eq:V20-stable-identification}.
\end{proof}

Let $q_s=\dim\mathcal E^s$ and choose an orthonormal stable-jet frame
$(V_a)_{a=1}^{q_s}$.  Apply the actual aligned common-action boundary
response to \eqref{eq:V20-boundary-layer-packet} and place the matched reads
in the columns of $\widehat C_n(x_\nu,\xi'_\mu)$.  Write
$\epsilon_{n,\partial}^{\mathrm{col}}$ for a certified upper bound on each
column error after adding:
\begin{enumerate}
 \item the analytic extraction error
       $C(r_n+(\lambda_nr_n)^{-1}+\lambda_n^{-1})$;
 \item the central-action-jet and route errors from
       \eqref{eq:V19-jet-debit}--\eqref{eq:V19-route-debit}, with the
       boundary normalization $\lambda_n^{-m_j-1/2}$;
 \item the actual-versus-local boundary-response discrepancy; and
 \item the stable-space transport error obtained from
       \eqref{eq:V20-Riesz-projection-bound}.
\end{enumerate}
Set
\begin{equation}
 \epsilon_{n,\partial}^{\mathrm{op}}
 :=\sqrt{q_s}\,\epsilon_{n,\partial}^{\mathrm{col}}.
 \label{eq:V20-boundary-column-error}
\end{equation}

\begin{theorem}[Finite observed complementing certificate]
\label{thm:V20-finite-complementing}
Let the boundary base points and unit tangential covectors form
$\delta_x$- and $\delta_\xi$-nets.  Suppose the occurring stable boundary
map, after the Riesz identification, has Lipschitz constant
$L_\partial$ in these variables.  If
\begin{equation}
 \min_{\nu,\mu}
 \sigma_{\min}\bigl(\widehat C_n(x_\nu,\xi'_\mu)\bigr)
 \ge\widehat\beta_n,
 \label{eq:V20-observed-boundary-floor}
\end{equation}
then its uniform complementing margin obeys
\begin{equation}
 \boxed{
 \beta_{\partial,n}
 \ge
 \widehat\beta_n
 -\epsilon_{n,\partial}^{\mathrm{op}}
 -L_\partial(\delta_x+\delta_\xi).}
 \label{eq:V20-certified-boundary-floor}
\end{equation}
A positive right side, together with the exact stable-mode/target dimension
count, certifies the complementing boundary row for that cutoff.
\end{theorem}

\begin{proof}
Lemma~\ref{lem:V19-column-Gram} converts column errors to
\eqref{eq:V20-boundary-column-error}.  At an arbitrary point choose a nearest
net pair.  The least singular value is $1$-Lipschitz in operator norm.
Subtract first the observed-matrix error and then the Lipschitz net tail from
\eqref{eq:V20-observed-boundary-floor}.  Taking the infimum proves
\eqref{eq:V20-certified-boundary-floor}.  Lemma~\ref{lem:V20-floor-inverse}
gives the final assertion.
\end{proof}

\begin{firewall}[The boundary basis cannot be frozen before the interior read]
\label{fire:V20-stable-basis}
A positive boundary matrix on stable modes of the declared pencil says
nothing about the occurring action Hessian unless the two stable spaces are
identified.  The interior packet matrix of V19, a normal-root resolvent
margin, and Lemma~\ref{lem:V20-stable-Riesz} must precede
\eqref{eq:V20-certified-boundary-floor}.  A root reaching the separating
contour is a glancing or loss-of-ellipticity witness, not a small boundary
tail.
\end{firewall}


% -----------------------------------------------------------------------------
\subsection{A passed model block and the boundary-quasimode converse}
\label{subsec:V20-boundary-converse}
% -----------------------------------------------------------------------------

\begin{proposition}[Laplace-type mixed boundary blocks]
\label{prop:V20-Laplace-mixed}
At a Riemannian boundary point choose orthonormal boundary-normal
coordinates.  For the reduced Einstein Laplace-type principal symbol
\[
 p_g^0(\xi',D_t)
 =
 \frac12\left(|\xi'|^2+D_t^2\right)I_N
\]
and $|\xi'|=1$, every stable profile is
\[
 V(t)=e^{-t}a,\qquad a\in\mathbb C^N.
\]
Let $Q_D,Q_N$ be orthogonal complementary projections and impose the
principal mixed conditions
\begin{equation}
 Q_DV(0)=g_D,\qquad Q_ND_tV(0)=g_N.
 \label{eq:V20-mixed-Laplace-boundary}
\end{equation}
Then the stable boundary map is
\begin{equation}
 a\longmapsto(Q_Da,iQ_Na)
 \label{eq:V20-mixed-stable-map}
\end{equation}
and has complementing margin $2^{-1/2}$ in the stable-jet norm
\eqref{eq:V20-stable-jet-norm}.  Dirichlet, Neumann, and a Robin condition
whose added term has order zero are included.
\end{proposition}

\begin{proof}
Since $D_t^2e^{-t}=-e^{-t}$, the displayed profiles solve the normal
equation and exhaust its decaying roots.  Also
$D_tV(0)=ia$.  Orthogonality and complementarity give
\[
 \|(Q_Da,iQ_Na)\|^2=\|a\|^2,
 \qquad
 \|V\|_s^2=\|a\|^2+\|ia\|^2=2\|a\|^2.
\]
This proves the floor.  A zeroth-order Robin term is absent from the
principal boundary polynomial and does not change the calculation.
\end{proof}

\begin{firewall}[The model block is not the assembled boundary condition]
\label{fire:V20-model-block}
Proposition~\ref{prop:V20-Laplace-mixed} covers a scalar-principal
second-order block.  It does not select compatible projections for the
metric, tetrad, gauge, Higgs, ghost, and fermion fields; it does not prove
the gauge boundary map; and it does not treat the first-order Dirac stable
space.  Those choices must be read from the same assembled common action and
tested by \eqref{eq:V20-certified-boundary-floor}.
\end{firewall}

\begin{theorem}[Failure of complementarity produces a boundary quasimode]
\label{thm:V20-boundary-quasimode}
Assume $d\ge2$ and that at $(x_*,\xi'_*)$ there is a stable witness $V_*$
as in \eqref{eq:V20-complementing-witness}.  With
$r_\lambda=\lambda^{-1/2}$, form
\[
 W_\lambda
 :=
 \lambda^{-2m}U_{\lambda,r_\lambda,V_*}.
\]
After multiplication by a fixed constant,
\begin{align}
 \|W_\lambda\|_{H^{2m}}&\longrightarrow1,
 \label{eq:V20-quasimode-normalization}\\
 \|PW_\lambda\|_{L^2}
 +\sum_{j=1}^J
 \|B_jW_\lambda\|_
 {H^{\,2m-m_j-1/2}(\partial M)}
 &\longrightarrow0,
 \label{eq:V20-boundary-quasimode}\\
 W_\lambda&\rightharpoonup0
 \quad\hbox{in }H^{2m}.
 \label{eq:V20-quasimode-weak}
\end{align}
Consequently, for every compact
$K:H^{2m}\to\mathscr Z$,
$\|KW_\lambda\|_{\mathscr Z}\to0$.  No estimate of the form
\begin{equation}
 \|u\|_{H^{2m}}
 \le C\left(
 \|Pu\|_{L^2}
 +\sum_j\|B_ju\|_{H^{2m-m_j-1/2}}
 +\|Ku\|_{\mathscr Z}\right)
 \label{eq:V20-forbidden-boundary-estimate}
\end{equation}
can therefore hold.

The same conclusion follows from a sequence whose complementing margins
tend to zero, by choosing a diagonal sequence of frequencies.
\end{theorem}

\begin{proof}
Every derivative of total order $2m$ falling on the carrier or normal
profile contributes $\lambda^{2m}$, whereas a tangential derivative falling
on the envelope contributes only $r_\lambda^{-1}=\lambda^{1/2}$.
Since $|\xi'_*|=1$ and $V_*\ne0$, the order-$2m$ derivative norm of
$U_{\lambda,r_\lambda,V_*}$ is comparable to $\lambda^{2m}$.
All lower derivatives are smaller.  This proves
\eqref{eq:V20-quasimode-normalization} after one fixed normalization.

The leading order-$2m$ interior term is
$p^0_{x_*}(\xi'_*,D_t)V_*=0$.  Repeating the full-norm estimates in the proof
of Theorem~\ref{thm:V20-boundary-extraction} shows that the remaining
interior response, after the factor $\lambda^{-2m}$, is
$O(\lambda^{-1/2})$ in $L^2$.  At the boundary, the unscaled leading
$B_j$ trace has $L^2$ size $\lambda^{m_j+1/2}$.
The tangential Sobolev multiplier of order
$2m-m_j-1/2$ contributes
$\lambda^{2m-m_j-1/2}$, exactly cancelling the factor
$\lambda^{m_j+1/2-2m}$.  Its leading coefficient is
$b^0_{j,x_*}(\xi'_*,D_t)V_*(0)=0$; coefficient freezing and an envelope
derivative again gain $\lambda^{-1/2}$.  This proves
\eqref{eq:V20-boundary-quasimode}.

The tangential Fourier support is centered at
$\lambda\xi'_*$ with width $O(\lambda^{1/2})$.  It escapes every fixed
frequency ball.  Boundedness in $H^{2m}$ and Fourier truncation therefore
give \eqref{eq:V20-quasimode-weak}.  Compact operators send weakly convergent
bounded sequences to strongly convergent sequences.  Substitution in
\eqref{eq:V20-forbidden-boundary-estimate} contradicts
\eqref{eq:V20-quasimode-normalization}.  If only a vanishing sequence of
margins is available, first choose a unit stable profile with boundary value
tending to zero and then choose $\lambda$ fast enough that the analytic
packet errors are smaller than that boundary value.
\end{proof}


% -----------------------------------------------------------------------------
\subsection{Lorentzian maximal flux graphs}
\label{subsec:V20-Lorentzian-flux}
% -----------------------------------------------------------------------------

Let $\Omega$ be a compact spatial domain with outward conormal $\nu$ and
consider the symmetric-hyperbolic system
\eqref{eq:V18-symmetric-hyperbolic-system}.  Its boundary matrix is
\begin{equation}
 A^\nu:=\sum_{a=1}^3\nu_aA^a.
 \label{eq:V20-boundary-flux-matrix}
\end{equation}
With the outward-normal convention, integration by parts places
$+\frac12\int_{\partial\Omega}\langle A^\nu U,U\rangle$ on the left of the
energy identity.  Thus a homogeneous allowed trace space must be
nonnegative for
\[
 q_\nu(u):=\langle A^\nu u,u\rangle.
\]

Let
$\mathcal H=\mathcal H_+\oplus\mathcal H_0\oplus\mathcal H_-$ be the
positive, zero, and negative spectral splitting of $A^\nu$, and define
\[
 \|u_+\|_+^2:=\langle A^\nu u_+,u_+\rangle,
 \qquad
 \|u_-\|_-^2:=-\langle A^\nu u_-,u_-\rangle.
\]

\begin{theorem}[Algebraic form of a maximal nonnegative boundary space]
\label{thm:V20-maximal-flux-graph}
A subspace $\mathcal M\subset\mathcal H$ is maximal among subspaces on which
$q_\nu\ge0$ if and only if
\begin{equation}
 \mathcal M
 =
 \{u_++Ku_++u_0:
   u_+\in\mathcal H_+,\ u_0\in\mathcal H_0\},
 \label{eq:V20-maximal-flux-space}
\end{equation}
where $K:\mathcal H_+\to\mathcal H_-$ is a contraction in the weighted
norms:
\begin{equation}
 \|Ku_+\|_-\le\|u_+\|_+.
 \label{eq:V20-flux-contraction}
\end{equation}
Every such space has dimension
$\dim\mathcal H_++\dim\mathcal H_0$.  If the contraction norm is
$\kappa<1$, then
\begin{equation}
 q_\nu(u_++Ku_++u_0)
 \ge(1-\kappa^2)\|u_+\|_+^2.
 \label{eq:V20-strict-flux-margin}
\end{equation}
\end{theorem}

\begin{proof}
A nonnegative subspace can contain no nonzero vector of
$\mathcal H_-$.  Modulo its intersection with $\mathcal H_0$, projection
onto $\mathcal H_+$ is therefore injective, so its dimension is at most
$\dim\mathcal H_++\dim\mathcal H_0$.  Maximality forces inclusion of
$\mathcal H_0$, since adding a zero mode preserves $q_\nu\ge0$, and it
forces the projection onto $\mathcal H_+$ to be onto; otherwise one can add
a positive vector orthogonal to that projection.  The remaining negative
component is therefore a linear map $K$ of $u_+$.  On its graph,
\[
 q_\nu(u_++Ku_++u_0)=\|u_+\|_+^2-\|Ku_+\|_-^2.
\]
Nonnegativity is equivalent to \eqref{eq:V20-flux-contraction}, and the
strict bound gives \eqref{eq:V20-strict-flux-margin}.  Conversely, that
calculation makes the graph nonnegative and its maximal possible dimension
makes it maximal.
\end{proof}

\begin{theorem}[Energy estimate with an inhomogeneous incoming trace]
\label{thm:V20-boundary-energy}
Assume the hypotheses of Theorem~\ref{thm:V18-hyperbolic-energy} on
$[0,T]\times\Omega$, allow a constant-rank zero-speed subspace, and impose
\begin{equation}
 U_-=KU_++g,\qquad \|K\|_{+\to-}\le\kappa<1,
 \label{eq:V20-incoming-boundary-condition}
\end{equation}
with the zero-speed variables left to their separately compatible
tangential equations.  Then there are $c_\partial>0$ and $C_T<\infty$ such
that every smooth compatible solution satisfies
\begin{align}
 &\|U(t)\|_{L^2(\Omega)}^2
 +c_\partial\int_0^t
   \|U_+(\tau)\|_{+,L^2(\partial\Omega)}^2\,\dd\tau
 \nonumber\\
 &\qquad\le
 C_T\left(
 \|U(0)\|_{L^2(\Omega)}^2
 +\|F\|_{L^2([0,t]\times\Omega)}^2
 +\|g\|_{-,L^2([0,t]\times\partial\Omega)}^2
 \right).
 \label{eq:V20-boundary-energy}
\end{align}
For $\kappa=1$ and homogeneous data $g=0$, the same argument gives an
energy estimate without a positive boundary-trace term.  Higher $H^r$
estimates follow when the differentiated
system has compatible boundary conditions and the required coefficient
bounds.
\end{theorem}

\begin{proof}
Put
$E(t)=\frac12\int_\Omega\langle A^0U,U\rangle$.  Symmetry and integration
by parts give
\[
 E'(t)+\frac12\int_{\partial\Omega}q_\nu(U)
 \le CE(t)+C\|F(t)\|_{L^2(\Omega)}^2.
\]
Writing $x=\|U_+\|_+$ and $y=\|g\|_-$ pointwise,
\[
 q_\nu(U)
 =x^2-\|KU_++g\|_-^2
 \ge(1-\kappa^2)x^2-2\kappa xy-y^2.
\]
Young's inequality yields
\begin{equation}
 q_\nu(U)
 \ge
 \frac{1-\kappa^2}{2}x^2
 -\frac{1+\kappa^2}{1-\kappa^2}y^2.
 \label{eq:V20-inhomogeneous-flux-bound}
\end{equation}
Insert this bound in the energy inequality, integrate, use the floor and
ceiling of $A^0$, and apply Gronwall.  This proves
\eqref{eq:V20-boundary-energy}.  If $\kappa=1$ and $g=0$, retain only
$q_\nu\ge0$.  Commuting compatible tangential and time derivatives
gives the final assertion; normal derivatives are recovered from the
equation on the nonzero-speed subspace.
\end{proof}

\begin{firewall}[The sign is fixed by the outward energy identity]
\label{fire:V20-flux-sign}
Some conventions call the same trace space maximal nonpositive after
reversing the normal.  In V20 the normal is outward and
$+\frac12q_\nu$ occurs on the left of the energy identity, so the allowed
homogeneous space is maximal nonnegative.  The negative spectral variables
are the incoming variables prescribed by
\eqref{eq:V20-incoming-boundary-condition}.
\end{firewall}


% -----------------------------------------------------------------------------
\subsection{Finite flux reconstruction and constraint preservation}
\label{subsec:V20-flux-card}
% -----------------------------------------------------------------------------

Work on the realification of every complex field so that $A^\nu$ is a real
symmetric matrix.  For an orthonormal frame $(e_a)_{a=1}^q$, the quadratic
flux reads
\[
 Q_a=q_\nu(e_a),\qquad
 Q_{ab}=q_\nu(e_a+e_b)
\]
reconstruct it by
\begin{equation}
 (A^\nu)_{aa}=Q_a,\qquad
 (A^\nu)_{ab}
 =\frac12(Q_{ab}-Q_a-Q_b)\quad(a<b).
 \label{eq:V20-flux-polarization}
\end{equation}

\begin{lemma}[Finite flux and inertia certificate]
\label{lem:V20-flux-reconstruction}
If every scalar quadratic read in
\eqref{eq:V20-flux-polarization} has absolute error at most
$\epsilon_Q$, then
\begin{equation}
 \|\widehat A^\nu-A^\nu\|_{\mathrm{op}}
 \le\frac{3q}{2}\epsilon_Q.
 \label{eq:V20-flux-matrix-error}
\end{equation}
If
\begin{equation}
 \min_{\lambda\in\operatorname{spec}\widehat A^\nu}|\lambda|
 >
 \frac{3q}{2}\epsilon_Q,
 \label{eq:V20-flux-gap}
\end{equation}
then $A^\nu$ and $\widehat A^\nu$ have the same positive and negative
inertias and no zero mode.  With a separated zero cluster, the same
conclusion holds for the positive, zero, and negative spectral clusters
using a Riesz projection bound as in
Lemma~\ref{lem:V20-stable-Riesz}.
\end{lemma}

\begin{proof}
A diagonal entry has error at most $\epsilon_Q$ and an off-diagonal entry
at most $3\epsilon_Q/2$.  The spectral norm is at most $q$ times the
largest entry, proving \eqref{eq:V20-flux-matrix-error}.  Weyl's eigenvalue
inequality proves inertia stability under \eqref{eq:V20-flux-gap}.  For a
zero cluster, use disjoint contours and the resolvent proof of
Lemma~\ref{lem:V20-stable-Riesz}.
\end{proof}

Let a homogeneous principal boundary condition be
\[
 R_-U_-+R_+U_+=0
\]
on the nonzero-speed subspaces.  In flux-normalized coordinates set
\begin{equation}
 \overline R_-:=R_-(-A^\nu_-)^{-1/2},
 \qquad
 \overline R_+:=R_+(A^\nu_+)^{-1/2}.
 \label{eq:V20-weighted-boundary-blocks}
\end{equation}

\begin{proposition}[Finite contraction-margin certificate]
\label{prop:V20-contraction-certificate}
Assume $\overline R_-$ is square.  Suppose the observed weighted blocks
have common operator error at most
$\epsilon_R$ and
\[
 \widehat\rho_-:=
 \sigma_{\min}(\widehat{\overline R}_-)>\epsilon_R.
\]
Then $R_-$ is invertible and the allowed trace is the graph
$U_-=KU_+$ with
\begin{equation}
 \boxed{
 \|K\|_{+\to-}
 \le
 \frac{\|\widehat{\overline R}_+\|+\epsilon_R}
      {\widehat\rho_- -\epsilon_R}.}
 \label{eq:V20-certified-contraction}
\end{equation}
If the displayed right side is below one and the zero modes have the
compatible treatment in Theorem~\ref{thm:V20-maximal-flux-graph}, the
boundary space is maximal nonnegative with a certified strict flux margin.
\end{proposition}

\begin{proof}
The singular-value perturbation inequality gives
$\sigma_{\min}(\overline R_-)
\ge\widehat\rho_- -\epsilon_R>0$.
The exact incoming graph in normalized variables is
\[
 (-A^\nu_-)^{1/2}U_-
 =
 -\overline R_-^{-1}\overline R_+
 (A^\nu_+)^{1/2}U_+.
\]
Hence
\[
 \|K\|_{+\to-}
 \le\|\overline R_-^{-1}\|\,\|\overline R_+\|
 \le
 \frac{\|\widehat{\overline R}_+\|+\epsilon_R}
      {\widehat\rho_- -\epsilon_R}.
\]
Apply Theorem~\ref{thm:V20-maximal-flux-graph}.
\end{proof}

Let the subsidiary constraint variables $c=\mathcal CU$ obey a
symmetric-hyperbolic system with boundary splitting
$P_C^++P_C^0+P_C^-=I$ and desired contraction $K_C$.  Let
$T_{\mathcal M}$ parametrize the allowed homogeneous main-system traces,
including any boundary equations needed to replace normal derivatives in
$\mathcal C$.  Define the induced constraint-boundary defect
\begin{equation}
 \mathfrak D_C
 :=
 (P_C^- -K_CP_C^+)\mathcal C_\partial T_{\mathcal M}.
 \label{eq:V20-constraint-boundary-defect}
\end{equation}

\begin{theorem}[Constraint-preserving boundary handoff]
\label{thm:V20-constraint-boundary}
Assume $\|K_C\|_{C,+\to C,-}\le\kappa_C<1$ and all traces in
\eqref{eq:V20-constraint-boundary-defect} are defined in the energy class.
If $w$ denotes the free main-system boundary trace, then
\begin{align}
 \|c(t)\|_{L^2}^2
 \le C_T\bigl(
 &\|c(0)\|_{L^2}^2
 +\|R_C\|_{L^2([0,t]\times\Omega)}^2
 \nonumber\\
 &+\|\mathfrak D_Cw\|_
 {L^2([0,t]\times\partial\Omega)}^2
 \bigr),
 \label{eq:V20-constraint-boundary-estimate}
\end{align}
where $R_C$ is the interior subsidiary residual.  In particular,
\begin{equation}
 c(0)=0,\qquad R_C=0,\qquad\mathfrak D_C=0
 \quad\Longrightarrow\quad c=0.
 \label{eq:V20-exact-constraint-preservation}
\end{equation}
\end{theorem}

\begin{proof}
On an allowed main trace,
\[
 P_C^-c=K_CP_C^+c+\mathfrak D_Cw.
\]
This is exactly the inhomogeneous incoming condition
\eqref{eq:V20-incoming-boundary-condition} for the subsidiary system.
Apply Theorem~\ref{thm:V20-boundary-energy} with source $R_C$ and boundary
datum $\mathfrak D_Cw$.  The zero-data implication follows from the
resulting estimate.
\end{proof}

\begin{firewall}[Main-system dissipation does not imply constraint preservation]
\label{fire:V20-main-versus-constraint}
The contraction bound \eqref{eq:V20-certified-contraction} controls the
energy of $U$.  It does not force
$\mathfrak D_C=0$.  Conversely, a constraint-preserving boundary rule can
still inject main-system energy if its flux graph is not nonnegative.
Both matrices must be read on the same action variables and transports.
\end{firewall}


% -----------------------------------------------------------------------------
\subsection{One-sided lattice boundary closures}
\label{subsec:V20-one-sided}
% -----------------------------------------------------------------------------

The centered stencil of V19 cannot be used at $x_d=0$ without values outside
the domain.  Let an actual order-$m_j$ boundary closure have the finite form
\begin{equation}
 (B_{j,h}u)(y,0)
 =
 h^{-m_j}\sum_{z=(z',z_d)\in Z_j}
 C_{j,h,z}(y)u(y+hz',hz_d),
 \qquad z_d\ge0.
 \label{eq:V20-one-sided-stencil}
\end{equation}
For $\theta=h\lambda$, define its combined tangential-normal moment
operators
\begin{equation}
 \mathscr M_{j,h,\ell}(\xi',D_t)
 :=
 \frac{i^\ell}{\ell!}
 \sum_{z\in Z_j}C_{j,h,z}
 \bigl(\xi'\cdot z'+z_dD_t\bigr)^\ell.
 \label{eq:V20-boundary-moment-operator}
\end{equation}

\begin{theorem}[One-sided boundary-stencil moment certificate]
\label{thm:V20-boundary-stencil-moments}
For a smooth normal profile $V$ and a frozen tangential carrier, the
normalized boundary response is
\begin{align}
 &\lambda^{-m_j}h^{-m_j}
 \sum_{z\in Z_j}C_{j,h,z}
 e^{ih\lambda\xi'\cdot z'}V(h\lambda z_d)
 \nonumber\\
 &\quad=
 \sum_{\ell=0}^{m_j-1}
 \theta^{\ell-m_j}
 \mathscr M_{j,h,\ell}(\xi',D_t)V(0)
 +\mathscr M_{j,h,m_j}(\xi',D_t)V(0)
 +R_{j,h,\theta}V,
 \label{eq:V20-boundary-moment-expansion}
\end{align}
where, for $\theta\le1$ and a fixed finite stencil,
\begin{equation}
 \|R_{j,h,\theta}V\|
 \le
 C_j\theta
 \max_{\substack{0\le t\le c_j\theta\\0\le k\le m_j+1}}
 \|D_t^kV(t)\|.
 \label{eq:V20-boundary-moment-remainder}
\end{equation}
Thus a sufficient differential-incidence audit is
\begin{align}
 \mathscr M_{j,h,\ell}&=0
 &&(0\le\ell<m_j),
 \label{eq:V20-boundary-lower-moments}\\
 \mathscr M_{j,h,m_j}(\xi',D_t)
 &\longrightarrow b^0_{j,x_0}(\xi',D_t)
 &&\text{uniformly on the stable jet space}.
 \label{eq:V20-boundary-top-moment}
\end{align}
Without the lower cancellations, a term in
\eqref{eq:V20-boundary-moment-expansion} can diverge as
$\theta^{\ell-m_j}$.
\end{theorem}

\begin{proof}
Tangential translation of the carrier and normal translation of the profile
combine as
\[
 e^{i\theta\xi'\cdot z'}V(\theta z_d)
 =
 e^{i\theta(\xi'\cdot z'+z_dD_t)}V(0),
\]
because $\partial_t=iD_t$.  Taylor-expand this exponential through order
$m_j$.  After division by $\theta^{m_j}$, its terms are exactly
\eqref{eq:V20-boundary-moment-operator}.  Taylor's remainder and finiteness
of $Z_j$ give \eqref{eq:V20-boundary-moment-remainder}.  The last assertions
follow directly from the expansion.
\end{proof}

For example, the forward realization of $D_d$ is
\begin{equation}
 D_{h,d}^+f(0):=\frac{f(h)-f(0)}{ih},
 \qquad
 \lambda^{-1}D_{h,d}^+V(\lambda\,\cdot)(0)
 =
 \frac{V(\theta)-V(0)}{i\theta}
 =
 D_tV(0)+O(\theta).
 \label{eq:V20-forward-normal-difference}
\end{equation}
Its generic error is first order in $h\lambda$, unlike the second-order
centered dispersion error in V19.

\begin{corollary}[Resolved boundary-layer scale]
\label{cor:V20-boundary-scale}
After the moment audit, a first-order-consistent one-sided boundary closure
and normalized lattice quadrature give the total analytic scale error
\begin{equation}
 C\left(
 r+\frac1{\lambda r}+\frac1\lambda
 +h\lambda+\frac hr
 \right).
 \label{eq:V20-boundary-lattice-error}
\end{equation}
The choice
\begin{equation}
 \boxed{
 \lambda_n=h_n^{-2/3},
 \qquad
 r_n=h_n^{1/3}}
 \label{eq:V20-balanced-boundary-scales}
\end{equation}
makes \eqref{eq:V20-boundary-lattice-error} $O(h_n^{1/3})$.
\end{corollary}

\begin{proof}
The first three terms are
Theorem~\ref{thm:V20-boundary-extraction}.  The one-sided Taylor remainder
is $O(h\lambda)$ by
Theorem~\ref{thm:V20-boundary-stencil-moments}; normalized tangential
quadrature and envelope sampling cost $O(h/r)$.  At the stated scales,
\[
 r=(\lambda r)^{-1}=h\lambda=h^{1/3},
 \qquad
 \lambda^{-1}=h^{2/3},
 \qquad
 h/r=h^{2/3}.
\]
\end{proof}

\begin{firewall}[A boundary closure is physical incidence data]
\label{fire:V20-boundary-closure}
Ghost cells, reflected values, a one-sided closure, and a boundary action
term can share the same interior stencil while producing different
operators in \eqref{eq:V20-boundary-moment-expansion}.  The actual
common-action closure and all of its lower moments must be retained; it
cannot be inferred from the V19 interior symbol.
\end{firewall}


% -----------------------------------------------------------------------------
\subsection{The first-order fermion boundary block}
\label{subsec:V20-first-order}
% -----------------------------------------------------------------------------

The even-order notation in
\eqref{eq:V20-elliptic-operator} is appropriate for the Laplace-type metric,
gauge, ghost, and Higgs blocks after gauge fixing.  The fermion Hessian has a
first-order principal block and must have its own stable-space count.

\begin{proposition}[First-order boundary-symbol certificate]
\label{prop:V20-first-order-boundary}
Let $Q$ have constant rank and let
\begin{equation}
 \mathscr D
 =
 A^d(x)D_d+\sum_{a<d}A^a(x)D_a+\operatorname{l.o.t.},
 \qquad A^d(x_0)\ \hbox{invertible},
 \label{eq:V20-first-order-operator}
\end{equation}
and let $Q(x,D_y)$ be an order-zero principal boundary map.  At
$(x_0,\xi')$ define
\begin{equation}
 \mathcal E_{\mathscr D}^s(x_0,\xi')
 :=
 \{V:\ (A^dD_t+A^a\xi'_a)V=0,\quad
          V(t)\text{ decays exponentially}\}.
 \label{eq:V20-Dirac-stable-space}
\end{equation}
Then the first-order complementing condition is exactly
\begin{equation}
 Q^0(x_0,\xi')|_{\mathcal E_{\mathscr D}^s}
 :
 \mathcal E_{\mathscr D}^s
 \longrightarrow\operatorname{Ran}Q^0
 \quad\hbox{is an isomorphism}.
 \label{eq:V20-first-order-complementing}
\end{equation}
Its quantitative floor is certified by the same column, Riesz-projection,
and finite-net bounds as
\eqref{eq:V20-certified-boundary-floor}.

For the packet \eqref{eq:V20-boundary-layer-packet},
\begin{align}
 \left\|
 \lambda^{-3/2}\mathfrak M_y(\mathscr DU)
 -(A^dD_t+A^a\xi'_a)V
 \right\|_{L^2_t}
 &\le C\left(r+\frac1{\lambda r}+\frac1\lambda\right)\|V\|_s,
 \label{eq:V20-first-order-interior-read}\\
 \left\|
 \lambda^{-1/2}\mathfrak M_y(QU)
 -Q^0V(0)
 \right\|
 &\le C\left(r+\frac1{\lambda r}+\frac1\lambda\right)\|V\|_s,
 \label{eq:V20-first-order-boundary-read}
\end{align}
where $\mathfrak M_y$ denotes tangential cutoff matching and demodulation.
If \eqref{eq:V20-first-order-complementing} fails, then
$\lambda^{-1}U_{\lambda,\lambda^{-1/2},V_*}$ is a bounded $H^1$ boundary
quasimode with vanishing interior and natural $H^{1/2}$ boundary residual.
\end{proposition}

\begin{proof}
The normal equation is first order, so its initial value identifies a stable
profile with a vector in the stable spectral subspace of the corresponding
normal generator.  The boundary data are uniquely solvable precisely when
the restriction in \eqref{eq:V20-first-order-complementing} is bijective.
The estimates follow from the proof of
Theorem~\ref{thm:V20-boundary-extraction} with total interior order one and
boundary order zero: the packet amplitude is $\lambda^{1/2}$, so the leading
interior and boundary scales are respectively $\lambda^{3/2}$ and
$\lambda^{1/2}$.  If a nonzero stable vector lies in the boundary kernel,
the proof of Theorem~\ref{thm:V20-boundary-quasimode}, with $2m$ replaced by
one, gives the stated quasimode.
\end{proof}

\begin{firewall}[A squared Dirac inverse does not choose the Dirac domain]
\label{fire:V20-Dirac-domain}
An elliptic boundary condition for $\mathscr D^*\mathscr D$ does not by
itself specify the domain or boundary eta data of $\mathscr D$.  The
first-order stable space, its target dimension, and the map
\eqref{eq:V20-first-order-complementing} must be retained for the occurring
fermion block.
\end{firewall}


% -----------------------------------------------------------------------------
\subsection{Cofinal boundary-domain handoff}
\label{subsec:V20-cofinal}
% -----------------------------------------------------------------------------

\begin{theorem}[Branch-typed boundary certificate for the common action]
\label{thm:V20-cofinal-boundary}
Assume the V19 interior symbol certificate has passed on the same transported
field variables and action roots.

On the Euclidean branch, suppose:
\begin{enumerate}[label=\textup{(E\arabic*)},leftmargin=3em]
 \item the occurring normal companion has a uniform separating contour and
       satisfies the Riesz perturbation bound
       \eqref{eq:V20-Riesz-projection-bound};
 \item the actual one-sided boundary closure satisfies
       \eqref{eq:V20-boundary-lower-moments} and
       \eqref{eq:V20-boundary-top-moment};
 \item the central-jet, route, local-response, extraction, and net errors in
       $\epsilon_{n,\partial}^{\mathrm{op}}$ tend uniformly to zero at the
       boundary scales \eqref{eq:V20-balanced-boundary-scales};
 \item the stable-mode/target dimensions agree and the certified right side
       of \eqref{eq:V20-certified-boundary-floor} has a positive uniform
       lower bound, for every even-order and first-order block.
\end{enumerate}
Then the occurring Euclidean boundary operators satisfy the uniform
complementing condition.  With the coefficient regularity and global graph
hypotheses of Corollary~\ref{cor:V18-Euclidean-branch}, this supplies its
boundary-principal row.  Failure of that row yields either a normal-root
contour witness, a lower-moment witness, a dimension mismatch, or the
boundary quasimode of Theorem~\ref{thm:V20-boundary-quasimode}.

On the Lorentzian branch, suppose:
\begin{enumerate}[label=\textup{(L\arabic*)},leftmargin=3em]
 \item the occurring first-order reduction has the $A^0$ floor and
       coefficient bounds of Theorem~\ref{thm:V18-hyperbolic-energy};
 \item the flux polarization error obeys
       \eqref{eq:V20-flux-matrix-error}, with uniformly separated spectral
       clusters and constant inertia;
 \item the incoming block is square, its certified contraction bound
       \eqref{eq:V20-certified-contraction} is uniformly below one, and all
       zero-speed equations are compatible;
 \item the subsidiary system has a contraction below one and its measured
       defect \eqref{eq:V20-constraint-boundary-defect}, initial constraint,
       and interior residual converge to zero in
       \eqref{eq:V20-constraint-boundary-estimate}.
\end{enumerate}
Then the main system has the boundary energy estimate
\eqref{eq:V20-boundary-energy}, and its constraints converge to zero in the
same energy class.  Exact zero data give exact constraint propagation.
Existence in the claimed causal graph additionally requires a construction
with the selected boundary and corner compatibility conditions; the energy
estimate alone proves uniqueness and closed range, not surjectivity.
\end{theorem}

\begin{proof}
On the Euclidean branch, item {\rm(E1)} identifies the occurring stable
bundle with the measured one.  Items {\rm(E2)} and {\rm(E3)}, together with
Theorems~\ref{thm:V20-boundary-stencil-moments} and
\ref{thm:V20-boundary-extraction}, identify the measured boundary matrix
with the occurring principal boundary map.  Item {\rm(E4)} and
Theorem~\ref{thm:V20-finite-complementing} give a uniform positive
Lopatinski floor.  Proposition~\ref{prop:V20-first-order-boundary} covers the
fermion block.  The failure alternatives are the contrapositives of the
four steps and Theorem~\ref{thm:V20-boundary-quasimode}.

On the Lorentzian branch, Lemma~\ref{lem:V20-flux-reconstruction} fixes the
inertia and flux splitting.  Proposition~\ref{prop:V20-contraction-certificate}
and Theorem~\ref{thm:V20-maximal-flux-graph} give a maximal nonnegative
boundary space with strict margin.  Theorem~\ref{thm:V20-boundary-energy}
gives the main energy estimate, and
Theorem~\ref{thm:V20-constraint-boundary} gives the constraint estimate.
An a priori estimate implies uniqueness and closed range by the difference
and graph-Cauchy arguments used in
Theorem~\ref{thm:V18-hyperbolic-energy}; it does not construct a solution,
which proves the final limitation.
\end{proof}

\begin{corollary}[Exact boundary debit in the coupled graph programme]
\label{cor:V20-coupled-boundary-debit}
After Theorem~\ref{thm:V20-cofinal-boundary}, the V18 coupled graph card may
replace an unspecified suitable-boundary assertion by the explicit tuple
\begin{equation}
 \mathfrak R_{\partial,n}
 =
 \left(
 \delta_{\Pi,n},\
 \rho^{<m}_{\partial,n},\
 \beta_{\partial,n},\
 \operatorname{Inertia}(A^\nu_n),\
 \kappa_n,\
 \|\mathfrak D_{C,n}\|
 \right).
 \label{eq:V20-boundary-debit-tuple}
\end{equation}
The Euclidean branch uses the first three entries; the Lorentzian branch
uses the last three, together with the $A^0$ floor.  Passing one branch does
not populate the other.
\end{corollary}

\begin{proof}
Each entry is respectively certified by
Lemma~\ref{lem:V20-stable-Riesz},
Theorem~\ref{thm:V20-boundary-stencil-moments},
Theorem~\ref{thm:V20-finite-complementing},
Lemma~\ref{lem:V20-flux-reconstruction},
Proposition~\ref{prop:V20-contraction-certificate}, and
Theorem~\ref{thm:V20-constraint-boundary}.
\end{proof}


% -----------------------------------------------------------------------------
\subsection{Executable V20 acquisition card}
\label{subsec:V20-card}
% -----------------------------------------------------------------------------

The Euclidean boundary card is populated as follows.

\begin{enumerate}[label=\textup{(E\arabic*)},leftmargin=3em]
 \item Retain the V19 interior symbol matrices and construct the occurring
 normal companion \eqref{eq:V20-normal-companion}.  Store the contour,
 resolvent maximum, Riesz projection, stable rank, and every root approaching
 the contour.

 \item At the boundary scales \eqref{eq:V20-balanced-boundary-scales}, issue
 cosine and sine packets for an orthonormal frame of every occurring stable
 space.  Retain both endpoints of each symmetric action secant, all
 relation/support/line transports, the matched interior residual, and every
 matched boundary column.

 \item Reconstruct the actual one-sided closure
 \eqref{eq:V20-one-sided-stencil}.  Store its offsets, block coefficients,
 lower moment operators, top moment, and the discrepancy from the declared
 principal boundary polynomial.

 \item Assemble $\widehat C_n$ for the even-order blocks and separately for
 the first-order fermion blocks.  Store the stable and target dimensions,
 singular values, certification net, held-out net, and Lipschitz tail.

 \item Define the scaled lower-moment debit
 \begin{equation}
  \rho_{\partial,n}^{<m}
  :=
  \max_{\substack{j,\ 0\le\ell<m_j\\x,\xi',\ \|V\|_s=1}}
  (h_n\lambda_n)^{\ell-m_j}
  \|\mathscr M_{j,h_n,\ell}(\xi',D_t)V(0)\|.
  \label{eq:V20-scaled-boundary-moment-debit}
 \end{equation}
 The certified complementing pass requires
 \begin{equation}
  \delta_{\Pi,n}\to0,\qquad
  \rho_{\partial,n}^{<m}\to0,\qquad
  \epsilon_{n,\partial}^{\mathrm{op}}\to0,\qquad
  \underline\beta_{\partial,n}
  :=
  \widehat\beta_n-\epsilon_{n,\partial}^{\mathrm{op}}
  -L_\partial(\delta_x+\delta_\xi)
  \ge\beta_0>0.
  \label{eq:V20-Euclidean-pass}
 \end{equation}
\end{enumerate}

The Lorentzian boundary card is different.

\begin{enumerate}[label=\textup{(L\arabic*)},leftmargin=3em]
 \item Retain the actual first-order reduction, its $A^0$ symmetrizer floor
 and ceiling, and its spatial coefficient matrices in the same field frame
 used by the action and constraint derivatives.

 \item At every boundary net point, run the real polarization reads
 \eqref{eq:V20-flux-polarization}.  Store $\widehat A^\nu$, its certified
 error, all spectral projectors, inertia, zero-speed cluster, and gap.

 \item Store the principal boundary matrix $R$, its weighted blocks
 \eqref{eq:V20-weighted-boundary-blocks}, the incoming block floor, and the
 right side of \eqref{eq:V20-certified-contraction}.  The main energy pass
 requires a uniform number $\kappa_0<1$ bounding that side.

 \item Derive the subsidiary system from the occurring gauge, Gauss,
 Hamiltonian, momentum, and Cartan constraints.  Store its flux splitting,
 contraction $K_C$, trace eliminations, and the full matrix
 $\mathfrak D_C$ in
 \eqref{eq:V20-constraint-boundary-defect}.  Exact preservation requires a
 zero matrix; convergence requires its norm and the interior forced
 gauge--Ward debit to vanish in the energy topology.

 \item Retain initial-boundary corner compatibility jets and run held-out
 boundary sources through \eqref{eq:V20-boundary-energy}.  An a priori pass
 is not recorded as a range or existence pass until a solution construction
 reaches every declared source, initial datum, and admissible boundary datum.
\end{enumerate}

Every failed row carries a finite witness:
\begin{align}
 W_n^{\mathrm{root}}
 &:=(x,\xi',z,\|(zI-\mathcal G_n)^{-1}\|),\\
 W_n^{\mathrm{bmom}}
 &:=(x,\xi',j,\ell,V,
     \mathscr M_{j,h_n,\ell}V(0)),\\
 W_n^{\mathrm{Lop}}
 &:=(x,\xi',V,\|\mathcal B_n^sV\|),\\
 W_n^{\mathrm{flux}}
 &:=(x,u,\langle A_n^\nu u,u\rangle),\\
 W_n^{\mathrm{in}}
 &:=(x,u_+,\|K_nu_+\|_-/\|u_+\|_+),\\
 W_n^{\mathrm{con}}
 &:=(x,w,\mathfrak D_{C,n}w).
 \label{eq:V20-boundary-witnesses}
\end{align}
A root witness sends the programme back to the interior pencil; a moment
witness to the lattice closure; a Lopatinski witness to the Euclidean
boundary operator; and the last three respectively to the causal reduction,
incoming rule, or constraint trace.


\subsection{V20 status and next upstream row}
\label{subsec:V20-status}

V20 adds the following results without repeating the V18 graph alternative or
the V19 interior packet calculation.

\begin{enumerate}
 \item The stable normal-mode space and its boundary map are converted into
 a quantitative Lopatinski floor.  Boundary-layer packets recover both the
 interior normal equation and boundary matrix with an explicit error.

 \item A Riesz resolvent estimate transports the declared stable space to
 the occurring common-action stable space.  This removes the circular step
 of testing a boundary operator on modes that have not themselves been
 observed.

 \item Failure of complementarity produces the normalized boundary
 quasimode \eqref{eq:V20-boundary-quasimode}, which directly violates the
 V18 compact-remainder estimate.

 \item The actual one-sided lattice closure has its own moment expansion and
 the resolved scale \eqref{eq:V20-balanced-boundary-scales}.  It is not
 replaced by the centered interior stencil.

 \item The first-order fermion domain is tested on its own stable space
 rather than inferred from a squared operator.

 \item On the causal branch, the finite flux polarization determines
 $A^\nu$, the incoming rule is certified as a strict contraction, and the
 induced subsidiary boundary map gives the exact constraint-preservation
 debit.
\end{enumerate}

The attached files do not contain the boundary-layer action responses,
one-sided closure kernels, stable companion matrices, Lopatinski singular
values, causal flux polarizations, incoming boundary blocks, or induced
constraint matrices required by \eqref{eq:V20-boundary-debit-tuple}.
Accordingly, V20 proves the implication and supplies obstruction witnesses,
but it does not record a physical Euclidean or Lorentzian boundary pass.

The next mathematically prior row is the \emph{adjoint domain and range}.
On the Euclidean branch, the Green boundary form must be extracted from the
same action, the annihilator boundary condition for the adjoint must be
computed, and both the domain kernel and adjoint kernel in
Theorem~\ref{thm:V18-compact-remainder} must be decided.  On the Lorentzian
branch, the selected maximal flux graph must be used in an actual
initial-boundary solution construction with all corner compatibility jets;
only then is surjectivity of the source--Cauchy--boundary graph map proved.
Those two tasks precede the final finite-to-infinite Schur and forced
gauge--Ward closure.

\begin{firewall}[V20 supplies boundary logic, not missing measurements]
\label{fire:V20-status}
A positive model calculation such as
Proposition~\ref{prop:V20-Laplace-mixed} is not the assembled
Einstein--Standard-Model boundary matrix.  The full boundary row passes only
through the actual common-action data and margins in
Theorem~\ref{thm:V20-cofinal-boundary}.  Until then, the full
Standard-Model--Einstein continuum remains open.
\end{firewall}

% END APPENDED VERSION V20 MATERIAL



% =============================================================================
% BEGIN APPENDED VERSION V21 MATERIAL
% =============================================================================

\clearpage
\section{Version V21: Green boundary concomitants, adjoint domains, and range}
\label{sec:V21-adjoint-domain}

V20 certified which boundary traces can define an elliptic or causal graph
domain, but V18 also requires the adjoint kernel to vanish.  A formal
coefficient adjoint does not determine that kernel: the adjoint operator has
a boundary domain fixed by the Green boundary concomitant and by the actual
allowed trace subspace.  The renewal Green forms earlier in the manuscript
are covariance and resolvent forms, and the finite Hilbert--Schmidt adjoints
act on process-history spaces.  Neither object is the differential boundary
concomitant used here.

V21 derives the adjoint boundary domain from finite trace matrices, gives a
stable numerical certificate for its nullspace, and separates self-adjoint
from merely formally self-adjoint realizations.  It then reduces the
Euclidean cokernel row to the already available V17 finite-screen and
high-tail tests.  On the Lorentzian branch it proves the causal range theorem
from a maximal-dissipativity resolvent condition, so that a boundary energy
estimate is no longer mistaken for existence.

% -----------------------------------------------------------------------------
\subsection{Abstract Green identity and the exact adjoint trace}
\label{subsec:V21-Green-adjoint}
% -----------------------------------------------------------------------------

Let
\[
 L_{\max}:\mathcal D_{\max}\subset H_0\to H_1,
 \qquad
 L^\sharp_{\max}:\mathcal D^\sharp_{\max}\subset H_1\to H_0
\]
be a formal differential adjoint pair with maximal graph domains.  Assume
continuous trace maps
\[
 \Gamma:\mathcal D_{\max}\to\mathcal T,
 \qquad
 \Gamma^\sharp:\mathcal D^\sharp_{\max}\to\mathcal T^\sharp
\]
and a matrix $\mathbb J:\mathcal T^\sharp\to\mathcal T$ for which
\begin{equation}
 \langle L_{\max}u,v\rangle_{H_1}
 -\langle u,L^\sharp_{\max}v\rangle_{H_0}
 =
 (\Gamma u)^*\mathbb J(\Gamma^\sharp v).
 \label{eq:V21-Green-identity}
\end{equation}
The trace spaces may be finite fibre spaces at principal-symbol level or
Sobolev boundary spaces at graph level.  Assume that compactly supported
smooth fields are trace-zero and dense in the minimal domains, and that the
allowed trace values below have continuous right inverses.

Let $\mathcal M\subset\mathcal T$ be the actual homogeneous allowed trace
space and let $T_{\mathcal M}:\mathbb C^r\to\mathcal T$ have range
$\mathcal M$.  Define
\begin{equation}
 \mathcal D(L_{\mathcal M})
 :=
 \{u\in\mathcal D_{\max}:\Gamma u\in\mathcal M\}.
 \label{eq:V21-primal-domain}
\end{equation}

\begin{theorem}[Adjoint-domain annihilator]
\label{thm:V21-adjoint-domain}
Under the preceding trace hypotheses,
\begin{equation}
 \boxed{
 \mathcal D(L_{\mathcal M}^*)
 =
 \left\{
 v\in\mathcal D^\sharp_{\max}:
 T_{\mathcal M}^*\mathbb J\,\Gamma^\sharp v=0
 \right\},}
 \label{eq:V21-adjoint-domain}
\end{equation}
and on this domain
\begin{equation}
 L_{\mathcal M}^*v=L^\sharp_{\max}v.
 \label{eq:V21-adjoint-action}
\end{equation}
Thus the adjoint principal boundary matrix is
\begin{equation}
 \boxed{B^\sharp_{\partial}:=T_{\mathcal M}^*\mathbb J.}
 \label{eq:V21-adjoint-boundary-matrix}
\end{equation}
It depends on the allowed primal trace space and the Green matrix, not only
on the formal coefficient adjoint.
\end{theorem}

\begin{proof}
Suppose first that $v\in\mathcal D(L_{\mathcal M}^*)$.  There is
$w\in H_0$ such that
\[
 \langle L_{\mathcal M}u,v\rangle_{H_1}
 =\langle u,w\rangle_{H_0}
 \qquad(u\in\mathcal D(L_{\mathcal M})).
\]
Testing on compactly supported smooth $u$ shows distributionally that
$w=L^\sharp v$, hence $v\in\mathcal D^\sharp_{\max}$ and
$w=L^\sharp_{\max}v$.  Green's identity now says
\[
 (\Gamma u)^*\mathbb J\Gamma^\sharp v=0
 \qquad(u\in\mathcal D(L_{\mathcal M})).
\]
The right-inverse hypothesis realizes every trace
$\Gamma u=T_{\mathcal M}a$.  Varying $a$ gives
$T_{\mathcal M}^*\mathbb J\Gamma^\sharp v=0$.

Conversely, if $v\in\mathcal D^\sharp_{\max}$ satisfies the displayed
boundary condition, the right side of \eqref{eq:V21-Green-identity} vanishes
for every $u\in\mathcal D(L_{\mathcal M})$.  Therefore
\[
 \langle L_{\mathcal M}u,v\rangle
 =\langle u,L^\sharp_{\max}v\rangle,
\]
which is exactly the definition of membership in
$\mathcal D(L_{\mathcal M}^*)$ and proves
\eqref{eq:V21-adjoint-action}.
\end{proof}

\begin{firewall}[Formal adjoint is not domain adjoint]
\label{fire:V21-formal-domain-adjoint}
Transport of a finite Hilbert--Schmidt adjoint or equality
$L^\sharp=L$ in the interior does not imply
$L_{\mathcal M}^*=L_{\mathcal M}$.  Equality of realizations additionally
requires equality of the two trace domains in
\eqref{eq:V21-adjoint-domain}.
\end{firewall}


% -----------------------------------------------------------------------------
\subsection{Finite acquisition of the Green matrix}
\label{subsec:V21-Green-acquisition}
% -----------------------------------------------------------------------------

At cutoff $n$, let $L_n^{\mathrm{bulk}}$ be the aligned bulk differential
linearization and let $(L_n^{\mathrm{bulk}})^\sharp$ be its coefficient
formal adjoint in the retained physical pairing.  Boundary action terms are
not dropped: they determine the conormal trace variables and the allowed
trace space $\mathcal M_n$.

Choose trace frames $(e_a)_{a=1}^{p}$ and
$(f_b)_{b=1}^{q}$ and extensions
$E_a\in\mathcal D_{\max}$ and
$E_b^\sharp\in\mathcal D^\sharp_{\max}$ such that
\[
 \Gamma E_a=e_a,\qquad
 \Gamma^\sharp E_b^\sharp=f_b.
\]
The Green matrix entries are the finite bilinear reads
\begin{equation}
 (\mathbb J_n)_{ab}
 :=
 \langle L_n^{\mathrm{bulk}}E_a,E_b^\sharp\rangle
 -
 \langle E_a,
        (L_n^{\mathrm{bulk}})^\sharp E_b^\sharp\rangle.
 \label{eq:V21-Green-matrix-read}
\end{equation}

\begin{lemma}[Extension independence and entrywise certification]
\label{lem:V21-Green-read}
The value in \eqref{eq:V21-Green-matrix-read} depends only on the two trace
vectors, not on the selected extensions.  If every observed entry has error
at most $\epsilon_{\mathrm{ent}}$, then
\begin{equation}
 \|\widehat{\mathbb J}_n-\mathbb J_n\|_{\mathrm{op}}
 \le\sqrt{pq}\,\epsilon_{\mathrm{ent}}.
 \label{eq:V21-Green-matrix-error}
\end{equation}
The entry error may be decomposed into the V17 central-jet, relation/line
route, bulk-locality, quadrature, and trace-extension residuals; these debits
must be retained separately.
\end{lemma}

\begin{proof}
Two extensions of the same first trace differ by a trace-zero field.
Applying \eqref{eq:V21-Green-identity} to that difference makes the
bilinear read vanish.  The same argument applies to the second extension.
For the error matrix $E$, the Frobenius estimate gives
\[
 \|E\|_{\mathrm{op}}\le\|E\|_{\mathrm F}
 \le\sqrt{pq}\,\epsilon_{\mathrm{ent}}.
\]
\end{proof}

Suppose an observed matrix $\widehat T_{\mathcal M}$ spans the allowed
primal traces, with
\[
 \|\widehat T_{\mathcal M}-T_{\mathcal M}\|\le\delta_T,
 \qquad
 \|\widehat{\mathbb J}-\mathbb J\|\le\delta_J.
\]
Put
\[
 C:=T_{\mathcal M}^*\mathbb J,\qquad
 \widehat C:=\widehat T_{\mathcal M}^*\widehat{\mathbb J},
\]
and define the computable error
\begin{equation}
 \delta_C
 :=
 \delta_T(\|\widehat{\mathbb J}\|+\delta_J)
 +\|\widehat T_{\mathcal M}\|\delta_J.
 \label{eq:V21-adjoint-matrix-error}
\end{equation}
Then $\|\widehat C-C\|\le\delta_C$.

\begin{lemma}[Stable adjoint-trace nullspace]
\label{lem:V21-adjoint-nullspace}
Assume $\widehat C$ has full row rank and let
$\widehat\rho$ be its smallest nonzero singular value.  If
$\delta_C<\widehat\rho$, then $C$ has the same full row rank,
\begin{equation}
 \rho_{\min}^+(C)\ge\widehat\rho-\delta_C,
 \label{eq:V21-adjoint-row-gap}
\end{equation}
and the orthogonal projectors onto the observed and actual adjoint trace
spaces obey
\begin{equation}
 \boxed{
 \|\Pi_{\ker\widehat C}-\Pi_{\ker C}\|
 \le
 \frac{\delta_C}{\widehat\rho}
 +\frac{\delta_C}{\widehat\rho-\delta_C}.}
 \label{eq:V21-adjoint-projector-error}
\end{equation}
\end{lemma}

\begin{proof}
Singular-value perturbation gives
\eqref{eq:V21-adjoint-row-gap}, so both row ranks and hence both kernel
dimensions agree.  If $x\in\ker C$ is unit and
$x=y+z$ with $y\in\ker\widehat C$ and
$z\perp\ker\widehat C$, then
\[
 \widehat\rho\|z\|
 \le\|\widehat Cz\|
 =\|\widehat Cx\|
 \le\delta_C.
\]
Thus
$\|(I-\Pi_{\ker\widehat C})\Pi_{\ker C}\|
\le\delta_C/\widehat\rho$.
Interchanging $C$ and $\widehat C$ and using
\eqref{eq:V21-adjoint-row-gap} gives the other directed gap.  The projector
difference is bounded by the sum of the two directed gaps, proving
\eqref{eq:V21-adjoint-projector-error}.
\end{proof}

\begin{firewall}[A scalar action value cannot determine the adjoint trace]
\label{fire:V21-scalar-Green}
The adjoint boundary matrix needs the complete bilinear panel
\eqref{eq:V21-Green-matrix-read}.  Diagonal action values alone do not
determine nonsymmetric or complex off-diagonal Green entries.  Polarization
is valid only after the physical pairing and common trace frame have been
fixed.
\end{firewall}


% -----------------------------------------------------------------------------
\subsection{Self-adjointness is a boundary Lagrangian condition}
\label{subsec:V21-Lagrangian}
% -----------------------------------------------------------------------------

Assume now $H_0=H_1$, $L^\sharp=L$, the two trace maps agree, and
$\mathbb J$ is nondegenerate and skew-Hermitian.  Define
\begin{equation}
 \mathcal M^{\perp_{\mathbb J}}
 :=
 \{z\in\mathcal T:m^*\mathbb Jz=0
                  \text{ for every }m\in\mathcal M\}.
 \label{eq:V21-J-annihilator}
\end{equation}

\begin{proposition}[Lagrangian criterion for a self-adjoint realization]
\label{prop:V21-Lagrangian-domain}
The realization $L_{\mathcal M}$ is symmetric if and only if
\[
 \mathcal M\subset\mathcal M^{\perp_{\mathbb J}},
\]
and it is self-adjoint if and only if
\begin{equation}
 \boxed{\mathcal M=\mathcal M^{\perp_{\mathbb J}}.}
 \label{eq:V21-Lagrangian-domain}
\end{equation}
In a $2N$-dimensional nondegenerate trace space, an isotropic
$N$-dimensional allowed trace space is therefore self-adjoint.
\end{proposition}

\begin{proof}
For $u,v$ in the primal domain, Green's identity shows that symmetry is
equivalent to $m^*\mathbb Jn=0$ for every $m,n\in\mathcal M$, which is the
displayed inclusion.  Theorem~\ref{thm:V21-adjoint-domain} says that the
adjoint trace space is exactly
$\mathcal M^{\perp_{\mathbb J}}$.  Equality of operator domains is therefore
equivalent to \eqref{eq:V21-Lagrangian-domain}.  Nondegeneracy gives
$\dim\mathcal M^{\perp_{\mathbb J}}=2N-\dim\mathcal M$; an isotropic
half-dimensional space must equal its annihilator.
\end{proof}

\begin{proposition}[Exact Laplace Dirichlet, Neumann, and Robin adjoints]
\label{prop:V21-Laplace-adjoint}
For a real Laplace-type bulk operator with trace
\[
 \Gamma u=(u|_{\partial M},\partial_\nu u|_{\partial M}),
\]
Green's formula has
\begin{equation}
 \mathbb J
 =
 \begin{pmatrix}0&I\\-I&0\end{pmatrix},
 \qquad
 (\Gamma u)^*\mathbb J\Gamma v
 =
 u^*\partial_\nu v-(\partial_\nu u)^*v.
 \label{eq:V21-Laplace-Green-matrix}
\end{equation}
Consequently:
\begin{enumerate}
 \item the adjoint of the Dirichlet realization is Dirichlet;
 \item the adjoint of the Neumann realization is Neumann;
 \item the adjoint of
       $\partial_\nu u+Su=0$ is
       $\partial_\nu v+S^*v=0$.
\end{enumerate}
The Robin realization is self-adjoint exactly when $S=S^*$ in the boundary
pairing.  Orthogonal mixed Dirichlet--Robin projections have the same
conclusion on their respective summands.
\end{proposition}

\begin{proof}
For Dirichlet, the allowed traces are $(0,b)$.  Pairing with
$(v,\partial_\nu v)$ in \eqref{eq:V21-Laplace-Green-matrix} gives
$-b^*v$, which vanishes for all $b$ exactly when $v=0$ at the boundary.
For Neumann, the traces are $(a,0)$ and the pairing is
$a^*\partial_\nu v$, giving Neumann again.  For Robin, the traces are
$(a,-Sa)$ and the pairing is
\[
 a^*\partial_\nu v+(-Sa)^*(-v)
 =a^*(\partial_\nu v+S^*v).
\]
Theorem~\ref{thm:V21-adjoint-domain} gives all three adjoint domains.
The last statements follow from
Proposition~\ref{prop:V21-Lagrangian-domain}.
\end{proof}

For a first-order Dirac-type block, the trace contains only the boundary
spinor and the Green matrix is the normal principal matrix, up to the fixed
symbol and inner-product convention.  Thus a projector that complements the
stable modes in Proposition~\ref{prop:V20-first-order-boundary} is
self-adjoint only if its kernel is also its Green annihilator.  Squaring the
Dirac operator tests neither equality.

\begin{counterexample}[A formally self-adjoint bulk expression can have a
non-self-adjoint domain]
\label{cth:V21-formal-not-selfadjoint}
On $[0,1]$, take $L=-\partial_x^2$ with
\[
 u(0)=0,\qquad u'(1)+\alpha u(1)=0,
\]
where $\alpha\in\mathbb C\setminus\mathbb R$.  The bulk expression is
formally self-adjoint, but the adjoint domain has
\[
 v(0)=0,\qquad v'(1)+\overline\alpha\,v(1)=0.
\]
The domains differ, so the realization is not self-adjoint.
\end{counterexample}

\begin{proof}
Apply Proposition~\ref{prop:V21-Laplace-adjoint} at the two endpoints.
Complex conjugation in the boundary pairing replaces $\alpha$ by
$\overline\alpha$.
\end{proof}


% -----------------------------------------------------------------------------
\subsection{Two-sided kernel and cokernel certification}
\label{subsec:V21-kernel-cokernel}
% -----------------------------------------------------------------------------

The compact-remainder estimate in
Theorem~\ref{thm:V18-compact-remainder} makes the primal kernel finite
dimensional; the adjoint requires its own domain and tail estimate.  A
finite matrix alone does not prove that either kernel is zero.  The missing
implication is supplied by a domain-tail estimate.

Let $P_N:\mathscr X\to\mathscr X$ and
$Q_N:\mathscr Y\to\mathscr Y$ be finite-rank orthogonal projections for a
closed realization $L:\mathscr X\to\mathscr Y$.  Define
\begin{align}
 \sigma_N
 &:=
 \inf_{\substack{x\in P_N\mathscr X\\\|x\|=1}}
 \|Q_NLP_Nx\|,
 \label{eq:V21-low-floor}\\
 c_N
 &:=
 \|Q_NL(I-P_N)\|_{\mathscr X\to\mathscr Y}.
 \label{eq:V21-cross-tail}
\end{align}
Assume a certified high-domain estimate
\begin{equation}
 \|(I-P_N)u\|_{\mathscr X}
 \le
 a_N\|Lu\|_{\mathscr Y}
 +b_N\|P_Nu\|_{\mathscr X}
 \qquad(u\in\mathscr X).
 \label{eq:V21-high-domain-estimate}
\end{equation}

\begin{theorem}[Finite low floor plus high tail removes the kernel]
\label{thm:V21-kernel-certificate}
If
\begin{equation}
 \boxed{\sigma_N>c_Nb_N,}
 \label{eq:V21-kernel-margin}
\end{equation}
then $\ker L=\{0\}$.
\end{theorem}

\begin{proof}
Let $Lu=0$ and split $u=u_L+u_H$ with
$u_L=P_Nu$.  Then
\[
 Q_NLu_L=-Q_NLu_H,
\]
so
\[
 \sigma_N\|u_L\|
 \le c_N\|u_H\|
 \le c_Nb_N\|u_L\|
\]
by \eqref{eq:V21-high-domain-estimate}.  The strict margin forces
$u_L=0$, and the same high-domain estimate then gives $u_H=0$.
\end{proof}

Construct the adjoint realization using
Theorem~\ref{thm:V21-adjoint-domain}, choose finite projections
$P_N^\sharp,Q_N^\sharp$ on its graph and target spaces, and define
$\sigma_N^\sharp,c_N^\sharp,a_N^\sharp,b_N^\sharp$ by the same formulas
with $L^*$.

\begin{corollary}[Two-sided finite certificate gives the Euclidean graph
isomorphism]
\label{cor:V21-Euclidean-isomorphism}
Assume the V20 Euclidean interior and complementing rows, the
compact-remainder estimate \eqref{eq:V18-compact-remainder-estimate}, and
the two strict margins
\begin{equation}
 \sigma_N>c_Nb_N,\qquad
 \sigma_N^\sharp>c_N^\sharp b_N^\sharp.
 \label{eq:V21-two-sided-kernel-margin}
\end{equation}
Then
\[
 \ker L=\ker L^*=\{0\},
\]
and $L$ is a bounded graph isomorphism.  If
Proposition~\ref{prop:V21-Lagrangian-domain} has already proved that the
realization is self-adjoint, the two tests coincide and only one is needed.
\end{corollary}

\begin{proof}
Theorem~\ref{thm:V21-kernel-certificate} applied to $L$ and $L^*$ removes
both kernels.  Theorem~\ref{thm:V18-compact-remainder} then gives closed
range, target coverage, and the bounded inverse.  For a self-adjoint
realization the operators and graph domains are identical.
\end{proof}

The finite matrices in \eqref{eq:V21-low-floor} are acquired by the V17
central-jet commands and their singular values are exact finite
computations.  The cross-tail and high-domain constants must be bounded by
the V17 target-leakage and unresolved-domain panels, using the primal
boundary basis for $L$ and the Green-derived adjoint boundary basis for
$L^*$.  Proposition~\ref{prop:V17-spectral-shell} supplies the common
spectral norm bookkeeping; it does not make $b_N$ small without measured
tail data.

\begin{assessment}[Meaning of a failed finite margin]
\label{ass:V21-failed-kernel-margin}
Failure of \eqref{eq:V21-kernel-margin} is not by itself a kernel.  It returns
a finite low singular vector and the high-tail block responsible for the
unclosed estimate.  A genuine kernel is a vector satisfying the full domain
and equation; a failed compact-remainder estimate instead returns the
quasimode of Theorem~\ref{thm:V18-compact-remainder} or the boundary
quasimode of Theorem~\ref{thm:V20-boundary-quasimode}.
\end{assessment}


% -----------------------------------------------------------------------------
\subsection{The Lorentzian range row is a resolvent row}
\label{subsec:V21-maximal-accretive}
% -----------------------------------------------------------------------------

Let $\mathcal A:\mathcal D(\mathcal A)\subset\mathcal H\to\mathcal H$ be the
spatial graph operator obtained from the first-order Lorentzian system after
multiplication by $(A^0)^{-1}$, with the V20 incoming boundary condition in
its domain.  The $A^0$ energy is the Hilbert norm.  The boundary flux
estimate and bounded lower-order terms give a number $\omega$ such that
\begin{equation}
 \operatorname{Re}\langle(\mathcal A+\omega I)u,u\rangle_{\mathcal H}
 \ge0.
 \label{eq:V21-quasi-accretive}
\end{equation}
Put $\mathcal A_\omega=\mathcal A+\omega I$.

\begin{lemma}[Adjoint-kernel criterion for maximal accretivity]
\label{lem:V21-maximal-accretive}
Let $A$ be densely defined, closed, and accretive.  For every $\lambda>0$,
\begin{equation}
 \|(I+\lambda A)u\|\ge\|u\|,
 \label{eq:V21-resolvent-lower-bound}
\end{equation}
so $\operatorname{Ran}(I+\lambda A)$ is closed.  Moreover,
\begin{equation}
 \operatorname{Ran}(I+\lambda A)=\mathcal H
 \quad\Longleftrightarrow\quad
 \ker(I+\lambda A^*)=\{0\}.
 \label{eq:V21-range-adjoint-kernel}
\end{equation}
Thus a single adjoint-kernel pass proves that $A$ is maximal accretive.
The domain of $A^*$ is the Green-annihilator domain in
Theorem~\ref{thm:V21-adjoint-domain}.
\end{lemma}

\begin{proof}
Accretivity gives
\[
 \|(I+\lambda A)u\|\,\|u\|
 \ge
 \operatorname{Re}\langle(I+\lambda A)u,u\rangle
 \ge\|u\|^2,
\]
which proves \eqref{eq:V21-resolvent-lower-bound}.  The same inequality
makes the range closed by the graph-Cauchy argument.  For every densely
defined operator,
\[
 \operatorname{Ran}(I+\lambda A)^\perp
 =\ker(I+\lambda A^*).
\]
A zero adjoint kernel therefore makes the closed range dense and hence all
of $\mathcal H$.  Surjectivity of one positive resolvent is the range
definition of maximal accretivity.
\end{proof}

\begin{theorem}[Causal range from the resolvent certificate]
\label{thm:V21-causal-range}
Let $A$ be maximal accretive on $\mathcal H$.  Then $-A$ generates a
strongly continuous contraction semigroup $S(t)$, and for every
$u_0\in\mathcal H$ and $f\in L^1(0,T;\mathcal H)$,
\begin{equation}
 u(t)=S(t)u_0+\int_0^tS(t-s)f(s)\,\dd s
 \label{eq:V21-Duhamel-solution}
\end{equation}
is the unique mild solution of
\[
 u'(t)+Au(t)=f(t),\qquad u(0)=u_0,
\]
with
\begin{equation}
 \boxed{
 \|u(t)\|
 \le\|u_0\|+\int_0^t\|f(s)\|\,\dd s.}
 \label{eq:V21-contraction-estimate}
\end{equation}
For $u_0\in\mathcal D(A)$ and sufficiently regular $f$, this is the unique
strong solution.  For the quasi-accretive operator $\mathcal A$ in
\eqref{eq:V21-quasi-accretive}, conjugation by $e^{-\omega t}$ gives the
same statements with the factor $e^{\omega t}$.
\end{theorem}

\begin{proof}
For $\tau>0$, maximal accretivity makes
$J_\tau=(I+\tau A)^{-1}$ an everywhere-defined contraction by
\eqref{eq:V21-resolvent-lower-bound}.  Put
$A_\tau=AJ_\tau=(I-J_\tau)/\tau$.  This is a bounded accretive operator, so
$S_\tau(t)=e^{-tA_\tau}$ is a contraction.  If $x\in\mathcal D(A)$, then
$A_\tau x=J_\tau Ax$ and hence
\begin{equation}
 \|A_\tau S_\tau(t)x\|
 \le\|Ax\|.
 \label{eq:V21-Yosida-derivative-bound}
\end{equation}

For $\tau,\sigma>0$, write
$u_\tau=S_\tau(t)x$, $p_\tau=J_\tau u_\tau$, and
$a_\tau=Ap_\tau=A_\tau u_\tau$, and similarly for $\sigma$.  Accretivity
of $A$ gives
\begin{align*}
 \frac12\frac{\dd}{\dd t}\|u_\tau-u_\sigma\|^2
 &=-\operatorname{Re}
   \langle a_\tau-a_\sigma,u_\tau-u_\sigma\rangle\\
 &=-\operatorname{Re}
   \langle a_\tau-a_\sigma,
           p_\tau-p_\sigma+\tau a_\tau-\sigma a_\sigma\rangle\\
 &\le(\tau+\sigma)\|a_\tau\|\,\|a_\sigma\|.
\end{align*}
Together with \eqref{eq:V21-Yosida-derivative-bound}, this yields
\begin{equation}
 \|S_\tau(t)x-S_\sigma(t)x\|^2
 \le2t(\tau+\sigma)\|Ax\|^2.
 \label{eq:V21-Yosida-Cauchy}
\end{equation}
Thus $S_\tau(t)x$ converges uniformly on compact time intervals for
$x\in\mathcal D(A)$.  Density and contractivity extend the limit $S(t)$ to
all $x\in\mathcal H$.  Passing to the limit in the semigroup law proves
$S(t+s)=S(t)S(s)$ and strong continuity.  The identities
$p_\tau=u_\tau-\tau a_\tau$ and
$u_\tau(t)-x=-\int_0^ta_\tau(s)\,\dd s$, together with closedness of the
graph of $A$, identify the generator as $-A$.

For step functions $f$, variation of constants gives
\eqref{eq:V21-Duhamel-solution}; contraction gives
\eqref{eq:V21-contraction-estimate}.  Approximation in $L^1$ defines the
Bochner integral and extends both statements to arbitrary $f$.  Applying
the estimate to the difference gives uniqueness.  Difference quotients and
closedness of $A$ give the strong-solution assertion for regular data.
Finally, $v=e^{-\omega t}u$ satisfies an equation with
$\mathcal A_\omega$, proving the quasi-accretive version.
\end{proof}

Combining Lemma~\ref{lem:V21-maximal-accretive} with
Theorem~\ref{thm:V21-kernel-certificate} yields a finite causal range card:
apply the low-floor/high-tail test to
\[
 L_\lambda^\sharp:=I+\lambda\mathcal A_\omega^*
\]
on the Green-derived adjoint boundary domain.  A strict kernel margin makes
$L_\lambda^\sharp$ injective, hence
$I+\lambda\mathcal A_\omega$ surjective.  Theorem~\ref{thm:V21-causal-range}
then constructs every source and initial datum rather than merely bounding
solutions that already exist.


\begin{corollary}[Bounded time-dependent perturbations]
\label{cor:V21-bounded-evolution}
Let $A$ be maximal accretive and let
$K\in L^\infty(0,T;\mathcal B(\mathcal H))$.  Then
\begin{equation}
 u'(t)+Au(t)+K(t)u(t)=f(t),\qquad u(0)=u_0
 \label{eq:V21-bounded-nonautonomous}
\end{equation}
has a unique mild solution for every
$u_0\in\mathcal H$ and $f\in L^1(0,T;\mathcal H)$, and
\begin{equation}
 \|u(t)\|
 \le
 e^{Mt}\left(\|u_0\|+\int_0^t\|f(s)\|\,\dd s\right),
 \qquad
 M:=\operatorname*{ess\,sup}_t\|K(t)\|.
 \label{eq:V21-bounded-perturbation-estimate}
\end{equation}
\end{corollary}

\begin{proof}
On an interval of length $\delta$ with $M\delta<1$, the map
\[
 \Phi(u)(t)
 =
 S(t)u_0+\int_0^tS(t-s)(f(s)-K(s)u(s))\,\dd s
\]
is a contraction on $C([0,\delta];\mathcal H)$.  Its fixed point solves
\eqref{eq:V21-bounded-nonautonomous}.  Iterating adjacent intervals reaches
$T$ and uniqueness patches the solutions.  Contractivity of $S$ gives
\[
 \|u(t)\|\le\|u_0\|+\int_0^t\|f(s)\|\,\dd s
             +M\int_0^t\|u(s)\|\,\dd s,
\]
and Gronwall proves \eqref{eq:V21-bounded-perturbation-estimate}.
\end{proof}

This corollary permits time-dependent lower-order couplings after one
maximal-accretive principal domain has been fixed.  It does not cover an
unbounded change of principal part or a changing boundary domain; those
require a common-domain evolution theorem and a measured graph-Lipschitz
bound.

% -----------------------------------------------------------------------------
\subsection{Initial-boundary corner compatibility}
\label{subsec:V21-corners}
% -----------------------------------------------------------------------------

For a time-dependent equation written as
\begin{equation}
 U'(t)+\mathcal A(t)U(t)=F(t),
 \qquad
 R(t)\Gamma U(t)=g(t),
 \label{eq:V21-time-boundary-system}
\end{equation}
define the formal initial time jets recursively by
\begin{align}
 U^{[0]}&:=U_0,
 \label{eq:V21-initial-jet-zero}\\
 U^{[k+1]}
 &:=
 F^{(k)}(0)
 -\sum_{\ell=0}^{k}{k\choose\ell}
   \mathcal A^{(k-\ell)}(0)U^{[\ell]}.
 \label{eq:V21-initial-jet-recursion}
\end{align}

\begin{proposition}[Exact corner-compatibility ledger]
\label{prop:V21-corner-compatibility}
Every $C^r$ strong solution of
\eqref{eq:V21-time-boundary-system} necessarily satisfies, for
$0\le k\le r-1$,
\begin{equation}
 \boxed{
 \sum_{\ell=0}^{k}{k\choose\ell}
 R^{(k-\ell)}(0)\Gamma U^{[\ell]}
 =
 g^{(k)}(0).}
 \label{eq:V21-corner-compatibility}
\end{equation}
Conversely, these identities remove the order-$r-1$ algebraic corner
obstruction; existence still requires the maximal-accretive or
nonautonomous range construction in the chosen graph class.

The finite residual
\begin{equation}
 \chi_{n,k}^{\mathrm{corner}}
 :=
 \sum_{\ell=0}^{k}{k\choose\ell}
 R_n^{(k-\ell)}(0)\Gamma_n U_n^{[\ell]}
 -g_n^{(k)}(0)
 \label{eq:V21-corner-residual}
\end{equation}
is the exact cutoffwise compatibility debit.
\end{proposition}

\begin{proof}
Differentiate the evolution equation $k$ times at $t=0$ and use Leibniz's
rule.  This gives \eqref{eq:V21-initial-jet-recursion} inductively.
Differentiate the boundary equation $k$ times and apply the same rule,
giving \eqref{eq:V21-corner-compatibility}.  A nonzero left-minus-right
value is incompatible with the boundary trace of any $C^r$ solution, while
its vanishing is only an algebraic compatibility statement and therefore
does not assert analytic existence.
\end{proof}

The constraint boundary condition has its own copy of
\eqref{eq:V21-corner-compatibility}.  It must be evaluated after the main
boundary rule is inserted into
$\mathfrak D_C$ from \eqref{eq:V20-constraint-boundary-defect}; otherwise a
main-system-compatible corner can still inject an incoming constraint.


% -----------------------------------------------------------------------------
\subsection{Cofinal adjoint-domain and range handoff}
\label{subsec:V21-cofinal}
% -----------------------------------------------------------------------------

\begin{theorem}[Branch-typed V21 graph-range certificate]
\label{thm:V21-cofinal-range}
Assume the V19 interior and V20 boundary-principal incidence rows are
evaluated on the same common-action roots and transported physical frames.

On the Euclidean branch, suppose:
\begin{enumerate}[label=\textup{(E\arabic*)},leftmargin=3em]
 \item the complete Green matrix and allowed trace basis are acquired with
       errors satisfying
       $\delta_C<\widehat\rho$ in
       Lemma~\ref{lem:V21-adjoint-nullspace};
 \item the resulting adjoint boundary space has the required dimension and
       passes its V20 stable-mode complementing test;
 \item the primal compact-remainder estimate holds and the primal and
       adjoint low-floor/high-tail margins in
       \eqref{eq:V21-two-sided-kernel-margin} are positive.
\end{enumerate}
Then the occurring Euclidean graph linearization is a bounded isomorphism.
If the Green matrix is nondegenerate and the allowed trace is certified
Lagrangian, the realization is self-adjoint and one kernel test suffices.

On an autonomous Lorentzian slab, suppose:
\begin{enumerate}[label=\textup{(L\arabic*)},leftmargin=3em]
 \item the V20 flux, strict incoming contraction, $A^0$, and constraint
       boundary rows pass;
 \item the spatial graph operator is closed, densely defined, and
       quasi-accretive as in \eqref{eq:V21-quasi-accretive};
 \item for some $\lambda>0$, the Green-derived adjoint realization
       $I+\lambda\mathcal A_\omega^*$ has a positive
       low-floor/high-tail kernel margin;
 \item the required corner residuals
       \eqref{eq:V21-corner-residual} vanish.
\end{enumerate}
Then $\mathcal A_\omega$ is maximal accretive,
\eqref{eq:V21-Duhamel-solution} reaches every source and initial datum, and
the V20 energy and constraint estimates hold for the constructed solution.
Bounded time-dependent lower-order Standard-Model couplings are admitted by
Corollary~\ref{cor:V21-bounded-evolution}.

A time-dependent principal part or boundary domain is not covered unless it
is transported to this fixed domain with a separately proved common-domain
graph estimate.
\end{theorem}

\begin{proof}
In the Euclidean case, Theorem~\ref{thm:V21-adjoint-domain} and
Lemma~\ref{lem:V21-adjoint-nullspace} identify the actual adjoint trace
domain.  The adjoint complementing and high-tail hypotheses make its finite
kernel test legitimate.  Corollary~\ref{cor:V21-Euclidean-isomorphism} gives
the bounded inverse.  Proposition~\ref{prop:V21-Lagrangian-domain} gives the
self-adjoint shortcut.

In the Lorentzian case, item {\rm(L1)} and
Theorem~\ref{thm:V20-boundary-energy} give quasi-accretivity.  Theorem
\ref{thm:V21-kernel-certificate} removes the adjoint resolvent kernel in
item {\rm(L3)}.  Lemma~\ref{lem:V21-maximal-accretive} then proves the range
condition, and Theorem~\ref{thm:V21-causal-range} constructs the solution.
Theorem~\ref{thm:V20-constraint-boundary} applies to that solution.
Proposition~\ref{prop:V21-corner-compatibility} gives the smooth corner
conditions.  The bounded perturbation statement is
Corollary~\ref{cor:V21-bounded-evolution}; its hypothesis explains the final
limitation.
\end{proof}

\begin{firewall}[Uniqueness, range, and smooth compatibility are three rows]
\label{fire:V21-three-range-rows}
The V20 energy estimate proves uniqueness for an existing causal solution.
The adjoint-resolvent test proves range and hence existence in the energy
class.  The corner ledger is additionally required for the claimed strong
regularity.  None of the three statements implies the other two without
the hypotheses in Theorem~\ref{thm:V21-cofinal-range}.
\end{firewall}

% -----------------------------------------------------------------------------
\subsection{Executable V21 acquisition card}
\label{subsec:V21-card}
% -----------------------------------------------------------------------------

For each cutoff and analytic branch, retain the following objects.

\begin{enumerate}[label=\textup{(\arabic*)},leftmargin=3em]
 \item Fix the bulk physical pairing, the primal and formal-adjoint trace
 frames, and trace right inverses.  Run every bilinear panel in
 \eqref{eq:V21-Green-matrix-read} through aligned central action jets and
 store the individual jet, route, locality, and quadrature errors.

 \item Store the V20 allowed-trace basis $T_{\mathcal M}$, the observed
 Green matrix, and
 \[
  \widehat C=\widehat T_{\mathcal M}^*
                    \widehat{\mathbb J}.
 \]
 Report $\delta_T,\delta_J,\delta_C,\widehat\rho$ and the projector error
 \eqref{eq:V21-adjoint-projector-error}.  Use the resulting nullspace basis,
 not the coefficient formal adjoint alone, in every adjoint command.

 \item On the Euclidean branch, rerun the V20 complementing matrix on the
 adjoint stable modes.  Then record for both $L$ and $L^*$ the finite low
 matrices, $\sigma_N,c_N,a_N,b_N$, and the margins
 \begin{equation}
  \mu_N:=\sigma_N-c_Nb_N,\qquad
  \mu_N^\sharp:=\sigma_N^\sharp-c_N^\sharp b_N^\sharp.
  \label{eq:V21-Euclidean-margin-card}
 \end{equation}
 A graph inverse is recorded only when the compact remainder and both
 margins pass, or when a proved Lagrangian equality identifies the two.

 \item On the Lorentzian branch, form
 $I+\lambda\mathcal A_\omega^*$ on the Green-derived adjoint incoming
 domain.  Record its finite low matrix and high tail, and the range margin
 \begin{equation}
  \mu_{\lambda,N}^{\mathrm{range}}
  :=
  \sigma_{\lambda,N}^\sharp
  -c_{\lambda,N}^\sharp b_{\lambda,N}^\sharp.
  \label{eq:V21-Lorentzian-range-margin}
 \end{equation}
 A positive number certifies the adjoint kernel is zero; the accretive lower
 bound then certifies closed target coverage.

 \item Generate the time jets
 \eqref{eq:V21-initial-jet-recursion} and store every main and constraint
 corner residual through the declared regularity order.  Also store whether
 time dependence is bounded on the fixed energy space, bounded only in the
 graph norm, or changes the domain; only the first case is covered by
 Corollary~\ref{cor:V21-bounded-evolution}.
\end{enumerate}

The failure payload is
\begin{align}
 W_n^{\mathrm G}
 &:=(a,b,\widehat{\mathbb J}_{ab},
     \epsilon_{ab}^{\mathrm{jet}},
     \epsilon_{ab}^{\mathrm{route}}),\\
 W_n^{\mathrm{ad}}
 &:=(\widehat C,\widehat\rho,\delta_C,
     \widehat\Pi_{\ker\widehat C}),\\
 W_n^{\ker}
 &:=(u_{N,\min},\sigma_N,c_N,b_N),\\
 W_n^{\operatorname{coker}}
 &:=(v_{N,\min}^\sharp,\sigma_N^\sharp,
     c_N^\sharp,b_N^\sharp),\\
 W_n^{\mathrm{range}}
 &:=(v_{\lambda,N}^\sharp,
     \mu_{\lambda,N}^{\mathrm{range}}),\\
 W_{n,k}^{\mathrm{corner}}
 &:=(k,\chi_{n,k}^{\mathrm{corner}}).
 \label{eq:V21-witness-family}
\end{align}
These witnesses distinguish a trace-frame failure, an unstable adjoint
domain, a primal kernel candidate, a cokernel candidate, a causal range
failure, and a corner incompatibility.


\subsection{V21 status and next upstream row}
\label{subsec:V21-status}

V21 proves the following new implications.

\begin{enumerate}
 \item The Green identity determines the exact adjoint trace domain by the
 annihilator matrix $T_{\mathcal M}^*\mathbb J$.  A full-row singular gap
 gives an explicit perturbation bound for that domain.

 \item Formal self-adjointness is separated from domain self-adjointness.
 The latter is exactly the Lagrangian trace condition
 \eqref{eq:V21-Lagrangian-domain}; Dirichlet, Neumann, and Robin adjoints are
 computed as checks.

 \item A finite low singular floor removes a continuum kernel only when the
 cross coupling and high-domain tail satisfy
 \eqref{eq:V21-kernel-margin}.  Running the test on the Green-derived
 adjoint domain decides the cokernel row needed by V18.

 \item On the causal branch, a single adjoint-resolvent kernel certificate
 promotes the V20 accretive energy estimate to maximal accretivity and hence
 to the constructive Duhamel range theorem
 \eqref{eq:V21-Duhamel-solution}.

 \item The exact corner recursion
 \eqref{eq:V21-corner-compatibility} separates energy-class existence from
 the higher regularity claimed for a smooth continuum.
\end{enumerate}

The attached manuscripts do not print the differential Green matrix in the
assembled physical trace frame, an allowed-trace basis for all
Einstein--Standard-Model fields, the primal and adjoint high-domain
constants, an adjoint resolvent matrix, or the corner jets.  Their
Hilbert--Schmidt formal adjoints and renewal Green forms do not fill these
rows.  V21 therefore supplies rigorous acquisition and implication
theorems, but it does not declare either kernel, cokernel, or causal range
physically passed.

The next upstream obstruction is a moving principal operator or moving
boundary domain.  A time-dependent Einstein--Standard-Model background
generally changes the flux spectral projectors and allowed trace space.
The next version must construct a bounded transport $W(t)$ of those
projectors, conjugate the problem to one fixed domain, and prove a
graph-Lipschitz nonautonomous evolution estimate.  Only after that range row
and the actual V21 cards pass can the V18 finite-to-infinite Schur and
forced gauge--Ward rows be assembled into the nonlinear continuum inverse.

\begin{firewall}[No graded or indefinite shortcut]
\label{fire:V21-graded-shortcut}
Ghost, fermion, gauge, and gravitational blocks need the pairing appropriate
to their realified or graded variables.  A positive Euclidean scalar block,
a symmetric action Hessian as a map to its dual, or a transported
Hilbert--Schmidt adjoint does not make the full gauge-fixed graph realization
self-adjoint.  The Green matrix and trace annihilator must be computed block
by block and then assembled in the common physical pairing.
\end{firewall}

% END APPENDED VERSION V21 MATERIAL



% =============================================================================
% BEGIN APPENDED VERSION V22 MATERIAL
% =============================================================================

\clearpage
\section{Version V22: moving boundary domains and nonautonomous causal range}
\label{sec:V22-moving-domain}

V21 proved the causal range theorem after one maximal-accretive spatial
domain had been fixed.  On an evolving Einstein--Standard-Model background,
the $A^0$ metric, flux eigenspaces, and allowed incoming trace graph can all
move with time.  Applying the autonomous theorem separately at each time
does not construct one evolution, because its operator domains are
different.

The earlier Kato transports in the attached programmes align finite
relation or overlap-projector supports.  Their ranges are not Sobolev PDE
domains and they do not lift a causal boundary trace into the bulk field
space.  Likewise, the attached Navier--Stokes evolution family acts on a
fixed Fourier domain without a physical boundary.  V22 supplies the missing
chain: an exact projector ODE, a boundary-collar lift, metric normalization,
conjugation to a fixed graph domain, and a frozen-time product theorem with
an explicit graph-Lipschitz residual.

% -----------------------------------------------------------------------------
\subsection{Exact unitary transport of an allowed trace space}
\label{subsec:V22-projector-transport}
% -----------------------------------------------------------------------------

Let $\Pi(t)$ be a strongly $C^1$ family of orthogonal projections on a
finite-dimensional or Hilbert trace space $\mathcal T$.  Define
\begin{equation}
 K_\Pi(t):=[\dot\Pi(t),\Pi(t)]
 =\dot\Pi(t)\Pi(t)-\Pi(t)\dot\Pi(t)
 \label{eq:V22-Kato-generator}
\end{equation}
and let $W_\partial(t)$ solve
\begin{equation}
 \dot W_\partial(t)=K_\Pi(t)W_\partial(t),
 \qquad W_\partial(0)=I.
 \label{eq:V22-boundary-transport-ODE}
\end{equation}

\begin{theorem}[Exact moving-trace transport]
\label{thm:V22-projector-transport}
The solution of \eqref{eq:V22-boundary-transport-ODE} is unitary and
\begin{equation}
 \boxed{
 \Pi(t)W_\partial(t)=W_\partial(t)\Pi(0).}
 \label{eq:V22-projector-intertwining}
\end{equation}
Consequently
$W_\partial(t)\operatorname{Ran}\Pi(0)=\operatorname{Ran}\Pi(t)$ and
\[
 W_\partial(t)\ker\Pi(0)=\ker\Pi(t).
\]
Moreover,
\begin{equation}
 \|W_\partial(t)-I\|
 \le
 2\int_0^t\|\dot\Pi(s)\|\,\dd s.
 \label{eq:V22-transport-length}
\end{equation}
\end{theorem}

\begin{proof}
Because $\Pi$ and $\dot\Pi$ are self-adjoint,
$K_\Pi^*=-K_\Pi$.  Hence
\[
 \frac{\dd}{\dd t}(W_\partial^*W_\partial)
 =W_\partial^*(K_\Pi^*+K_\Pi)W_\partial=0,
\]
so $W_\partial$ is unitary.  Differentiate $\Pi^2=\Pi$ to obtain
\[
 \Pi\dot\Pi\Pi=0,\qquad
 \dot\Pi=\dot\Pi\Pi+\Pi\dot\Pi.
\]
These identities give
\[
 [\Pi,K_\Pi]=-\dot\Pi.
\]
Therefore
\[
 \frac{\dd}{\dd t}
 (W_\partial^*\Pi W_\partial)
 =
 W_\partial^*(\dot\Pi+[\Pi,K_\Pi])W_\partial=0.
\]
Its value is $\Pi(0)$, proving
\eqref{eq:V22-projector-intertwining}.  Finally
$\|K_\Pi\|\le2\|\dot\Pi\|$ and integration of the ODE, using unitarity,
gives \eqref{eq:V22-transport-length}.
\end{proof}

\begin{lemma}[Measured-projector transport error]
\label{lem:V22-projector-transport-error}
Let $\widehat\Pi(t)$ be another $C^1$ orthogonal projector family with
transport $\widehat W_\partial$, and set
\[
 \delta_0(t)=\|\widehat\Pi(t)-\Pi(t)\|,
 \qquad
 \delta_1(t)=\|\dot{\widehat\Pi}(t)-\dot\Pi(t)\|.
\]
Then
\begin{align}
 \|\widehat K_\Pi-K_\Pi\|
 &\le
 2\delta_1
 +\delta_0\bigl(\|\dot\Pi\|+\|\dot{\widehat\Pi}\|\bigr),
 \label{eq:V22-generator-error}\\
 \|\widehat W_\partial(t)-W_\partial(t)\|
 &\le
 \int_0^t
 \left[
 2\delta_1(s)
 +\delta_0(s)
  \bigl(\|\dot\Pi(s)\|+\|\dot{\widehat\Pi}(s)\|\bigr)
 \right]\dd s.
 \label{eq:V22-transport-error}
\end{align}
\end{lemma}

\begin{proof}
Insert and subtract cross terms:
\begin{align*}
 \widehat K_\Pi-K_\Pi
 &=(\dot{\widehat\Pi}-\dot\Pi)\widehat\Pi
   +\dot\Pi(\widehat\Pi-\Pi)\\
 &\quad-(\widehat\Pi-\Pi)\dot{\widehat\Pi}
   -\Pi(\dot{\widehat\Pi}-\dot\Pi).
\end{align*}
Taking norms proves \eqref{eq:V22-generator-error}.  Variation of constants
for the two skew-adjoint ODEs gives
\[
 \widehat W_\partial(t)-W_\partial(t)
 =
 \int_0^t
 \widehat W_\partial(t,s)
 (\widehat K_\Pi(s)-K_\Pi(s))W_\partial(s)\,\dd s.
\]
Both propagators are unitary, so
\eqref{eq:V22-transport-error} follows.
\end{proof}


\begin{lemma}[Derivative of a separated spectral projector]
\label{lem:V22-spectral-projector-derivative}
Let $A(t)$ be a $C^1$ matrix family and let a fixed contour $\Gamma$
separate a spectral cluster for all $t$.  With
\[
 M_\Gamma=\sup_{t,z\in\Gamma}\|(zI-A(t))^{-1}\|,
\]
its Riesz projector satisfies
\begin{align}
 \Pi_A(t)
 &=
 \frac1{2\pi i}\int_\Gamma(zI-A(t))^{-1}\,\dd z,
 \label{eq:V22-Riesz-projector}\\
 \dot\Pi_A(t)
 &=
 \frac1{2\pi i}\int_\Gamma
 (zI-A(t))^{-1}\dot A(t)(zI-A(t))^{-1}\,\dd z,
 \label{eq:V22-Riesz-derivative}\\
 \|\dot\Pi_A(t)\|
 &\le
 \frac{|\Gamma|}{2\pi}M_\Gamma^2\|\dot A(t)\|.
 \label{eq:V22-Riesz-derivative-bound}
\end{align}
\end{lemma}

\begin{proof}
Differentiate the resolvent identity
$(zI-A)(zI-A)^{-1}=I$ to obtain
\[
 \frac{\dd}{\dd t}(zI-A)^{-1}
 =(zI-A)^{-1}\dot A(zI-A)^{-1}.
\]
Uniform separation permits differentiation under the finite contour
integral.  Taking norms proves the bound.
\end{proof}

The allowed trace projector need not be a spectral projector.  If a
full-column-rank matrix $T(t)$ parametrizes the trace space, then
\begin{equation}
 \Pi_{\mathcal M}(t)
 =
 T(t)(T(t)^*T(t))^{-1}T(t)^*.
 \label{eq:V22-trace-graph-projector}
\end{equation}

\begin{lemma}[Conditioned derivative of a trace-graph projector]
\label{lem:V22-trace-projector-derivative}
If $\sigma_{\min}(T(t))\ge s_0>0$, then
\begin{equation}
 \|\dot\Pi_{\mathcal M}(t)\|
 \le
 2\|\dot T\|\|T\|s_0^{-2}
 +2\|\dot T\|\|T\|^3s_0^{-4}.
 \label{eq:V22-trace-projector-derivative}
\end{equation}
Thus the V20 incoming graph and zero-speed trace columns give a computable
projector-speed bound once their column floor is retained.
\end{lemma}

\begin{proof}
Put $G=T^*T$.  Differentiating
\eqref{eq:V22-trace-graph-projector} gives
\[
 \dot\Pi
 =\dot TG^{-1}T^*+TG^{-1}\dot T^*
 -TG^{-1}(\dot T^*T+T^*\dot T)G^{-1}T^*.
\]
Use $\|G^{-1}\|\le s_0^{-2}$ and
$\|\dot T^*T+T^*\dot T\|\le2\|T\|\|\dot T\|$.
\end{proof}

% -----------------------------------------------------------------------------
\subsection{A boundary-collar lift into the PDE domain}
\label{subsec:V22-collar-lift}
% -----------------------------------------------------------------------------

Let the first-order field bundle be trivialized in a boundary collar
$\partial\Omega\times[0,\rho_0)$, and assume
$K_\Pi(t,y)$ has the bounded spatial derivatives required below.  Choose
$\chi\in C_c^\infty([0,\rho_0))$ with $\chi(0)=1$ and solve pointwise
\begin{equation}
 \partial_t\mathcal W(t,y,x_d)
 =
 \chi(x_d)K_\Pi(t,y)\mathcal W(t,y,x_d),
 \qquad
 \mathcal W(0,y,x_d)=I.
 \label{eq:V22-collar-transport}
\end{equation}
Extend $\mathcal W$ by the identity outside the collar and let the same
letter denote multiplication by this matrix on bulk fields.

\begin{theorem}[Unitary lift of the moving causal domain]
\label{thm:V22-collar-lift}
For every $(t,y,x_d)$, $\mathcal W$ is unitary.  At the boundary it equals
$W_\partial(t,y)$ and, for each integer $s\ge0$ for which the coefficient
derivatives are bounded,
\begin{equation}
 \|\mathcal W(t)\|_{H^s\to H^s}
 +\|\mathcal W(t)^{-1}\|_{H^s\to H^s}
 +\|\partial_t\mathcal W(t)\|_{H^s\to H^s}
 \le C_{s,T}.
 \label{eq:V22-collar-Sobolev-bound}
\end{equation}
If
\begin{equation}
 \mathcal D_t
 :=
 \{u\in H^1(\Omega;E):
       \gamma u\in\operatorname{Ran}\Pi(t)\},
 \label{eq:V22-moving-domain}
\end{equation}
then
\begin{equation}
 \boxed{\mathcal W(t)\mathcal D_0=\mathcal D_t.}
 \label{eq:V22-domain-intertwining}
\end{equation}
\end{theorem}

\begin{proof}
The generator in \eqref{eq:V22-collar-transport} is skew-adjoint, so the
pointwise solution is unitary.  At $x_d=0$ the ODE is
\eqref{eq:V22-boundary-transport-ODE}, while outside the support of $\chi$
it has the constant solution $I$.  Differentiating the ODE in $t$, $y$, and
$x_d$, followed by induction and Gronwall, bounds all coefficient
derivatives through order $s$.  The Sobolev multiplication theorem then
gives \eqref{eq:V22-collar-Sobolev-bound}.

The trace multiplication law and
\eqref{eq:V22-projector-intertwining} give
\[
 \gamma(\mathcal Wu)=W_\partial\gamma u,\qquad
 W_\partial\operatorname{Ran}\Pi(0)
 =\operatorname{Ran}\Pi(t).
\]
Thus $\mathcal W\mathcal D_0\subset\mathcal D_t$.  Apply the same argument
to $\mathcal W^{-1}$ for the reverse inclusion.
\end{proof}

\begin{firewall}[A trace unitary needs a bulk lift]
\label{fire:V22-trace-versus-domain}
The matrix $W_\partial$ alone does not act on $L^2(\Omega)$ and cannot
conjugate a PDE.  The collar lift, its Sobolev bounds, and its compatibility
with every transported field and constraint trace are separate incidence
rows.  For higher-order elliptic jet traces, an arbitrary jet-space unitary
need not be the trace of a bounded zeroth-order bulk multiplier; V22's lift
is asserted for the first-order causal value trace.
\end{firewall}


% -----------------------------------------------------------------------------
\subsection{Metric normalization and fixed-domain conjugation}
\label{subsec:V22-metric-conjugation}
% -----------------------------------------------------------------------------

Let $M(t,x)=A^0(t,x)$ satisfy
\begin{equation}
 0<a_0I\preceq M(t,x)\preceq a_1I
 \label{eq:V22-moving-metric-floor}
\end{equation}
and have the bounded time and spatial derivatives required by the Sobolev
order.  Define the pointwise metric normalizer
\begin{equation}
 R(t,x):=M(t,x)^{-1/2}M(0,x)^{1/2}.
 \label{eq:V22-metric-normalizer}
\end{equation}
Then
\begin{equation}
 R(t)^*M(t)R(t)=M(0).
 \label{eq:V22-metric-isometry}
\end{equation}
Pull the physical allowed trace space back by $R(t)|_{\partial\Omega}$
and take its orthogonal projector in the fixed $M(0)$ metric.  Conjugate
pointwise by $S_0(x)=M(0,x)^{1/2}$ so this metric becomes the standard one,
apply Theorem~\ref{thm:V22-projector-transport}, extend the resulting
standard skew generator through the collar as in
\eqref{eq:V22-collar-transport}, and conjugate the lift back by $S_0$.
Denote this $M(0)$-unitary collar lift by $\mathcal W(t)$ and put
\begin{equation}
 \mathcal U(t):=R(t)\mathcal W(t).
 \label{eq:V22-total-domain-transport}
\end{equation}

\begin{theorem}[Metric-unitary moving-domain identification]
\label{thm:V22-metric-domain-transport}
The multiplier $\mathcal U(t)$ maps the fixed physical domain
$\mathcal D_0$ bijectively onto $\mathcal D_t$ and satisfies
\begin{equation}
 \boxed{
 \mathcal U(t)^*M(t)\mathcal U(t)=M(0).}
 \label{eq:V22-total-metric-isometry}
\end{equation}
Its Sobolev operator norms, inverse norms, and time-derivative norms are
bounded in terms of $a_0,a_1$, the derivatives of $M$, and the bounds in
\eqref{eq:V22-collar-Sobolev-bound}.

If
\begin{equation}
 u'(t)+\mathcal A(t)u(t)=f(t),
 \qquad u(t)\in\mathcal D_t,
 \label{eq:V22-original-moving-equation}
\end{equation}
then $u=\mathcal Uv$ is equivalent to the fixed-domain equation
\begin{align}
 v'(t)+\widetilde{\mathcal A}(t)v(t)
 &=\widetilde f(t),\qquad v(t)\in\mathcal D_0,
 \label{eq:V22-fixed-domain-equation}\\
 \widetilde{\mathcal A}(t)
 &:=
 \mathcal U(t)^{-1}\mathcal A(t)\mathcal U(t)
 +\mathcal U(t)^{-1}\dot{\mathcal U}(t),
 \label{eq:V22-conjugated-generator}\\
 \widetilde f(t)&:=\mathcal U(t)^{-1}f(t).
 \label{eq:V22-conjugated-source}
\end{align}
\end{theorem}

\begin{proof}
Equation \eqref{eq:V22-metric-isometry} is direct multiplication.  The
pulled-back boundary projector and
Theorem~\ref{thm:V22-collar-lift} show that $\mathcal W$ maps the initial
pulled-back trace space onto the time-$t$ one; multiplication by $R$ returns
it to the physical trace space.  This proves the domain assertion.
The $S_0$ conjugation makes the extended collar generator skew in the
standard metric; after conjugating back, $\mathcal W^*M(0)\mathcal W=M(0)$
throughout the collar.  Therefore
\[
 \mathcal U^*M(t)\mathcal U
 =\mathcal W^*R^*M(t)R\mathcal W
 =\mathcal W^*M(0)\mathcal W=M(0).
\]
Uniform positive spectral bounds make the matrix square root and inverse
square root smooth functions of $M$.  Differentiating their functional
calculus, followed by the Sobolev multiplication estimate and
Theorem~\ref{thm:V22-collar-lift}, gives the asserted bounds.  Finally,
substitute $u=\mathcal Uv$ into
\eqref{eq:V22-original-moving-equation}, expand
$u'=\dot{\mathcal U}v+\mathcal Uv'$, and multiply by
$\mathcal U^{-1}$.
\end{proof}

Let
\[
 C_{\mathcal U}(t):=\mathcal U(t)^{-1}\dot{\mathcal U}(t).
\]
Differentiating \eqref{eq:V22-total-metric-isometry} gives the exact
connection identity
\begin{equation}
 2\operatorname{Re}
 \langle C_{\mathcal U}v,v\rangle_{M(0)}
 =
 -\langle \dot M\,\mathcal Uv,\mathcal Uv\rangle.
 \label{eq:V22-metric-connection-identity}
\end{equation}
Thus, if
\[
 \kappa_M
 :=
 \sup_t\|M(t)^{-1/2}\dot M(t)M(t)^{-1/2}\|,
\]
then
\begin{equation}
 \operatorname{Re}
 \langle C_{\mathcal U}v,v\rangle_{M(0)}
 \ge-\frac{\kappa_M}{2}\|v\|_{M(0)}^2.
 \label{eq:V22-connection-lower-bound}
\end{equation}

\begin{lemma}[Maximal accretivity survives the fixed-domain transport]
\label{lem:V22-transported-maximality}
Suppose $\mathcal A(t)+\omega I$ is maximal accretive in the $M(t)$ metric
on $\mathcal D_t$.  Then
$\widetilde{\mathcal A}(t)+\Omega I$ is maximal accretive in the fixed
$M(0)$ metric for every uniform
\[
 \Omega\ge\omega+\sup_t\|C_{\mathcal U}(t)\|.
\]
\end{lemma}

\begin{proof}
Metric-unitary conjugation preserves maximal accretivity.  The bounded
operator $C_{\mathcal U}+\|C_{\mathcal U}\|I$ is accretive.  The sum of a
maximal accretive operator and a bounded accretive operator is maximal
accretive: accretivity is immediate, and for sufficiently small $\lambda$,
\[
 I+\lambda(A+B)
 =
 \left[I+\lambda B(I+\lambda A)^{-1}\right](I+\lambda A)
\]
is onto by the Neumann series.  Lemma~\ref{lem:V21-maximal-accretive} then
gives maximal accretivity.  Adding any further nonnegative scalar preserves
the conclusion.
\end{proof}

\begin{firewall}[Metric normalization changes the connection]
\label{fire:V22-metric-connection}
Replacing $A^0(t)$ by the identity is not a free notation change.  Its time
derivative reappears exactly as
\eqref{eq:V22-metric-connection-identity}.  Omitting that bounded connection
can reverse the claimed energy sign.
\end{firewall}


% -----------------------------------------------------------------------------
\subsection{A common-domain nonautonomous product theorem}
\label{subsec:V22-common-domain}
% -----------------------------------------------------------------------------

Let $\mathcal D$ be a common dense domain in a Hilbert space $\mathcal H$.
For $A(t):\mathcal D\to\mathcal H$, use the graph norm
\begin{equation}
 \|u\|_t:=\|u\|+\|A(t)u\|.
 \label{eq:V22-time-graph-norm}
\end{equation}

\begin{theorem}[Frozen-time construction on a common graph domain]
\label{thm:V22-common-domain-evolution}
Assume:
\begin{enumerate}[label=\textup{(N\arabic*)},leftmargin=3em]
 \item every $A(t)$ is maximal accretive on $\mathcal H$ with domain
       $\mathcal D$;
 \item for a constant $L$,
       \begin{equation}
        \|(A(t)-A(s))u\|
        \le L|t-s|\|u\|_s
        \qquad(u\in\mathcal D).
        \label{eq:V22-graph-Lipschitz}
       \end{equation}
\end{enumerate}
Then there is a unique contraction evolution family
$\mathscr U(t,s)$ on $\mathcal H$, $0\le s\le t\le T$, such that for
$u_0\in\mathcal H$ and $f\in L^1(0,T;\mathcal H)$,
\begin{equation}
 u(t)=\mathscr U(t,0)u_0+
 \int_0^t\mathscr U(t,s)f(s)\,\dd s
 \label{eq:V22-nonautonomous-Duhamel}
\end{equation}
is the unique mild solution of
\begin{equation}
 u'(t)+A(t)u(t)=f(t),\qquad u(0)=u_0.
 \label{eq:V22-nonautonomous-equation}
\end{equation}
It satisfies
\begin{equation}
 \boxed{
 \|u(t)\|
 \le\|u_0\|+\int_0^t\|f(s)\|\,\dd s.}
 \label{eq:V22-nonautonomous-energy}
\end{equation}
For $u_0\in\mathcal D$ and graph-continuous
$f:[0,T]\to\mathcal D$, the solution is strong away from the usual
data-regularity endpoints and is the limit of the frozen-time products
below.  A uniform quasi-accretive shift $\Omega$ changes the right side by
the factor $e^{\Omega t}$.
\end{theorem}

\begin{proof}
First take $u_0\in\mathcal D$ and graph-continuous
$f:[0,T]\to\mathcal D$.  For a partition
$\pi=\{0=t_0<\cdots<t_N=T\}$, define $u_\pi$ on
$[t_k,t_{k+1}]$ by
\begin{equation}
 u_\pi'(t)+A(t_k)u_\pi(t)=f(t),
 \qquad
 u_\pi(t_k+)=u_\pi(t_k-).
 \label{eq:V22-frozen-equation}
\end{equation}
Theorem~\ref{thm:V21-causal-range} constructs each interval.  Its semigroup
commutes with $A(t_k)$ on $\mathcal D$, so it is a contraction in
$\|\cdot\|_{t_k}$.  Condition
\eqref{eq:V22-graph-Lipschitz} also gives
\begin{equation}
 \|u\|_{t_{k+1}}
 \le(1+L|t_{k+1}-t_k|)\|u\|_{t_k}.
 \label{eq:V22-adjacent-graph-norms}
\end{equation}
Iteration and variation of constants therefore give a partition-independent
bound
\begin{equation}
 \sup_{t\in[0,T]}\|u_\pi(t)\|_t
 \le
 C_{L,T}\left(
 \|u_0\|_0+
 \int_0^T\|f(s)\|_s\,\dd s\right)
 =:C_*.
 \label{eq:V22-frozen-graph-bound}
\end{equation}

Viewed as an approximate solution of the time-dependent equation,
$u_\pi$ has residual
\begin{equation}
 r_\pi(t)
 :=
 (A(t)-A(t_k))u_\pi(t),
 \qquad t_k\le t<t_{k+1}.
 \label{eq:V22-frozen-residual}
\end{equation}
Equations \eqref{eq:V22-graph-Lipschitz} and
\eqref{eq:V22-frozen-graph-bound} imply
\begin{equation}
 \|r_\pi\|_{L^1(0,T;\mathcal H)}
 \le LTC_*|\pi|,
 \qquad
 |\pi|:=\max_k(t_{k+1}-t_k).
 \label{eq:V22-frozen-residual-bound}
\end{equation}

For two partitions, $w=u_\pi-u_{\pi'}$ is continuous and piecewise strong,
and
\[
 w'+A(t)w=r_\pi-r_{\pi'}
\]
away from their combined nodes.  Accretivity gives
\[
 \frac12\frac{\dd}{\dd t}\|w\|^2
 \le\|r_{\pi'}-r_\pi\|\,\|w\|.
\]
Integrating across all subintervals yields
\begin{equation}
 \sup_{t\le T}\|u_\pi(t)-u_{\pi'}(t)\|
 \le
 \|r_\pi\|_{L^1}+\|r_{\pi'}\|_{L^1}.
 \label{eq:V22-partition-Cauchy}
\end{equation}
Thus $u_\pi$ converges in $C([0,T];\mathcal H)$ as $|\pi|\to0$.

The graph norms are uniformly equivalent by
\eqref{eq:V22-adjacent-graph-norms}, and
\eqref{eq:V22-frozen-graph-bound} makes the approximants weakly bounded in
the common graph space.  Weak closedness of the linear graphs and
\eqref{eq:V22-frozen-residual-bound} identify the limit with a strong
solution of \eqref{eq:V22-nonautonomous-equation} for the present data.
For two such solutions, the same accretive energy calculation with zero
residual proves
\[
 \|u(t)-v(t)\|
 \le\|u(0)-v(0)\|+\int_0^t\|f-g\|\,\dd s.
 \label{eq:V22-solution-stability}
\]
Density of $\mathcal D$ in $\mathcal H$ and of graph-continuous functions
in $L^1(0,T;\mathcal H)$ therefore extends the solution uniquely to all
stated data and proves \eqref{eq:V22-nonautonomous-energy}.

Solving from an arbitrary entrance time and using uniqueness defines
$\mathscr U(t,s)$ and its composition law.  Applying the construction to
step sources and passing by $L^1$ approximation gives
\eqref{eq:V22-nonautonomous-Duhamel}.  For quasi-accretive $A(t)$, set
$v(t)=e^{-\Omega t}u(t)$ and apply the proved result to $A(t)+\Omega I$.
\end{proof}


% -----------------------------------------------------------------------------
\subsection{Finite certification of the graph-Lipschitz row}
\label{subsec:V22-graph-card}
% -----------------------------------------------------------------------------

For maximal accretive $A(s)$, define the bounded relative variation
\begin{equation}
 \mathfrak G(t,s)
 :=
 (A(t)-A(s))(I+A(s))^{-1}.
 \label{eq:V22-relative-graph-variation}
\end{equation}

\begin{lemma}[Resolvent form of graph Lipschitz continuity]
\label{lem:V22-resolvent-graph-Lipschitz}
The estimate
\begin{equation}
 \|\mathfrak G(t,s)\|\le L_R|t-s|
 \label{eq:V22-relative-variation-bound}
\end{equation}
implies \eqref{eq:V22-graph-Lipschitz} with $L=L_R$.  Conversely,
\eqref{eq:V22-graph-Lipschitz} implies
\begin{equation}
 \|\mathfrak G(t,s)\|\le3L|t-s|.
 \label{eq:V22-graph-to-resolvent}
\end{equation}
\end{lemma}

\begin{proof}
For $u\in\mathcal D$, put $x=(I+A(s))u$.  Then
\[
 (A(t)-A(s))u=\mathfrak G(t,s)x,
 \qquad
 \|x\|\le\|u\|_s,
\]
which proves the first implication.  Conversely,
\eqref{eq:V21-resolvent-lower-bound} gives
$\|u\|\le\|x\|$, while $A(s)u=x-u$ gives
$\|A(s)u\|\le2\|x\|$.  Hence $\|u\|_s\le3\|x\|$, and
\eqref{eq:V22-graph-Lipschitz} proves
\eqref{eq:V22-graph-to-resolvent}.
\end{proof}

Let $P_N,Q_N$ be finite screens in the fixed energy space and let
$\mathcal N_{N,\eta}$ be an $\eta$-net of the unit sphere of
$P_N\mathcal H$.  Retain
\begin{align}
 d_{N,\eta}(t,s)
 &:=
 \max_{v\in\mathcal N_{N,\eta}}
 \|Q_N\mathfrak G(t,s)P_Nv\|,
 \label{eq:V22-relative-net}\\
 \ell_N^{\mathrm{out}}(t,s)
 &:=
 \|(I-Q_N)\mathfrak G(t,s)P_N\|,
 \label{eq:V22-relative-target-tail}\\
 \tau_N^{\mathrm{high}}(t,s)
 &:=
 \|\mathfrak G(t,s)(I-P_N)\|.
 \label{eq:V22-relative-domain-tail}
\end{align}

\begin{proposition}[Finite relative-graph certificate]
\label{prop:V22-finite-graph-Lipschitz}
For every $t\ne s$,
\begin{equation}
 \frac{\|\mathfrak G(t,s)\|}{|t-s|}
 \le
 \frac1{|t-s|}
 \left[
 \frac{d_{N,\eta}(t,s)}{1-\eta}
 +\ell_N^{\mathrm{out}}(t,s)
 +\tau_N^{\mathrm{high}}(t,s)
 \right].
 \label{eq:V22-finite-relative-bound}
\end{equation}
A uniform bound on the right over a time-pair net, together with a measured
modulus for the time derivative of $\mathfrak G$, certifies
\eqref{eq:V22-relative-variation-bound} on the full time square.
\end{proposition}

\begin{proof}
Apply Lemma~\ref{lem:V17-net-bound} to
$Q_N\mathfrak G P_N$, then split an arbitrary input by
$P_N+(I-P_N)$ and the output by $Q_N+(I-Q_N)$.  This gives the bracket in
\eqref{eq:V22-finite-relative-bound}.  A nearest time-pair net point and
the derivative modulus bound the unsampled tail by the net radius.
\end{proof}

The resolvent in \eqref{eq:V22-relative-graph-variation} is not a formal
inverse.  It must be supplied by the V21 adjoint-resolvent range card on the
transported fixed domain.  The time derivative of the allowed projector is
likewise measured by symmetric time secants:
\begin{equation}
 \widehat{\dot\Pi}_\delta(t)
 :=
 \frac{\widehat\Pi(t+\delta)-\widehat\Pi(t-\delta)}{2\delta}.
 \label{eq:V22-projector-time-secant}
\end{equation}
If $\|\Pi^{(3)}\|\le M_3$, its analytic error is at most
$M_3\delta^2/6$; endpoint route errors are divided by $2\delta$ and must be
added before using \eqref{eq:V22-transport-error}.

\begin{firewall}[Pointwise maximality does not give an evolution]
\label{fire:V22-pointwise-maximality}
Even if every frozen operator is maximal accretive, arbitrarily fast graph
variation can prevent the frozen products from converging.  The relative
resolvent matrix \eqref{eq:V22-relative-graph-variation}, its high-domain
tail, and a uniform time modulus are independent acquisition rows.
\end{firewall}


% -----------------------------------------------------------------------------
\subsection{Inhomogeneous boundary data are a lifting row}
\label{subsec:V22-boundary-lifting}
% -----------------------------------------------------------------------------

The fixed operator domain is homogeneous.  To reach prescribed incoming
data, let $\mathcal G_\partial$ be the boundary data space and assume a
time-dependent right inverse
\begin{equation}
 \mathcal E_\partial(t):\mathcal G_\partial\to\mathcal D_{\max},
 \qquad
 R(t)\Gamma\mathcal E_\partial(t)g=g.
 \label{eq:V22-boundary-right-inverse}
\end{equation}

\begin{proposition}[Exact reduction of inhomogeneous boundary data]
\label{prop:V22-boundary-lifting}
Let $g$ be such that
\[
 G(t):=\mathcal E_\partial(t)g(t)
 \in W^{1,1}(0,T;\mathcal H),
 \qquad
 A(t)G(t)\in L^1(0,T;\mathcal H).
\]
Then $u$ solves
\begin{equation}
 u'+A(t)u=f,\qquad
 R(t)\Gamma u=g,\qquad u(0)=u_0
 \label{eq:V22-inhomogeneous-boundary-problem}
\end{equation}
if and only if $w=u-G$ solves the homogeneous-domain problem
\begin{equation}
 w'+A(t)w
 =
 f-G'-A(t)G,
 \qquad
 w(0)=u_0-G(0).
 \label{eq:V22-lifted-source}
\end{equation}
Hence Theorem~\ref{thm:V22-common-domain-evolution} reaches every boundary
datum for which the displayed lifting source is in $L^1$.  Different right
inverses give the same $u$ by uniqueness.
\end{proposition}

\begin{proof}
The right-inverse identity gives
$R\Gamma(u-G)=0$.  Substitution of $u=w+G$ gives
\eqref{eq:V22-lifted-source}, including its initial value.  The converse is
the same calculation in reverse.  Two lift choices produce two solutions
of \eqref{eq:V22-inhomogeneous-boundary-problem}; their difference has zero
source, initial value, and homogeneous boundary trace, so the energy
estimate makes it zero.
\end{proof}

The regularity of $G$ is not optional.  A boundary datum in the V20 energy
trace class may yield only a transposition solution, whereas V17's nonlinear
graph inverse may require a strong solution.  The corner residuals in
\eqref{eq:V21-corner-residual} test compatibility of $u_0$, $g$, and the
chosen lifting at their common edge.


% -----------------------------------------------------------------------------
\subsection{Cofinal moving-domain causal graph theorem}
\label{subsec:V22-cofinal}
% -----------------------------------------------------------------------------

\begin{theorem}[Moving-domain Einstein--Standard-Model causal handoff]
\label{thm:V22-cofinal-moving-domain}
Assume on a compact time interval:
\begin{enumerate}[label=\textup{(M\arabic*)},leftmargin=3em]
 \item the V19 interior symbol, V20 flux/incoming/constraint, and V21
       Green-adjoint resolvent rows pass in common physical frames, uniformly
       in time;
 \item $A^0=M$ satisfies \eqref{eq:V22-moving-metric-floor}, the allowed
       causal trace has constant rank, and its projector and time derivative
       satisfy the bounds in
       Theorem~\ref{thm:V22-projector-transport} and
       Lemmas~\ref{lem:V22-spectral-projector-derivative}--%
       \ref{lem:V22-trace-projector-derivative};
 \item the metric and collar transport
       $\mathcal U(t)$ has the bounds in
       Theorem~\ref{thm:V22-metric-domain-transport};
 \item for one uniform $\Omega$, every fixed-domain operator
       $\widetilde{\mathcal A}(t)+\Omega I$ is maximal accretive and the
       finite relative-graph card proves
       \eqref{eq:V22-relative-variation-bound};
 \item the inhomogeneous boundary right inverse satisfies
       Proposition~\ref{prop:V22-boundary-lifting}, all required main and
       constraint corner residuals vanish, and the forced gauge--Ward and
       induced constraint-boundary debits converge to zero in their energy
       norms.
\end{enumerate}
Then the physical moving-domain graph map
\begin{equation}
 \mathcal H_{\mathrm{mov}}u
 :=
 \left(
 u'+\mathcal A(t)u,\
 u(0),\
 R(t)\Gamma u
 \right)
 \label{eq:V22-moving-graph-map}
\end{equation}
is one-to-one and reaches every datum in the declared source, initial, and
liftable boundary-data class satisfying the corner conditions.  Its inverse
is bounded by the metric, transport, graph-Lipschitz, quasi-accretive, and
lifting constants displayed above.

For homogeneous boundary data, the solution is explicitly
\begin{equation}
 u(t)
 =
 \mathcal U(t)\left[
 \widetilde{\mathscr U}(t,0)u_0+
 \int_0^t
 \widetilde{\mathscr U}(t,s)
 \mathcal U(s)^{-1}f(s)\,\dd s
 \right],
 \label{eq:V22-physical-evolution}
\end{equation}
where $\widetilde{\mathscr U}$ is the fixed-domain evolution family from
Theorem~\ref{thm:V22-common-domain-evolution}.  Inhomogeneous boundary data
are added by Proposition~\ref{prop:V22-boundary-lifting}.

Thus $\mathcal H_{\mathrm{mov}}$ is a valid causal base graph isomorphism
$H_*$ for the V17 action-root theorem.  The nonlinear conclusion still
requires the actual derivative-radius, mixed Schur-tail, and forced
gauge--Ward margins of V17--V18.
\end{theorem}

\begin{proof}
Items {\rm(M2)} and {\rm(M3)} give
$\mathcal U(t)\mathcal D_0=\mathcal D_t$ and transform the moving-metric
problem into \eqref{eq:V22-fixed-domain-equation}.  Item {\rm(M4)} and
Theorem~\ref{thm:V22-common-domain-evolution} construct the unique
fixed-domain evolution family and its Duhamel solution.  Conjugating back
gives \eqref{eq:V22-physical-evolution}; metric isometry and the bounds on
$\mathcal U$ transfer the energy estimate.  Proposition
\ref{prop:V22-boundary-lifting} adds every admitted boundary datum and
preserves uniqueness.  Item {\rm(M5)} supplies the strong corner and
constraint conditions.

The constructed solution proves range, while the energy estimate for a
difference proves injectivity and bounds the inverse.  This is exactly the
source--initial--boundary graph map needed as $H_*$.  The last statement
follows because the V17 nonlinear root theorem and V18 Schur theorem have
additional hypotheses beyond invertibility of the diagonal causal graph.
\end{proof}

\begin{firewall}[A fixed-domain formula is not a fixed background]
\label{fire:V22-fixed-domain-not-fixed-coefficients}
Conjugation makes the operator domain constant; it does not freeze the
Einstein--Standard-Model coefficients.  Their remaining time variation is
the relative graph operator \eqref{eq:V22-relative-graph-variation}.
Omitting it would replace the physical evolution by a different autonomous
problem.
\end{firewall}


% -----------------------------------------------------------------------------
\subsection{Executable V22 acquisition card}
\label{subsec:V22-card}
% -----------------------------------------------------------------------------

At every cutoff retain the following time-resolved matrices and errors.

\begin{enumerate}[label=\textup{(\arabic*)},leftmargin=3em]
 \item In the V20 realified field frame, store $A^0_n(t)$,
 $A^\nu_n(t)$, the incoming boundary blocks, zero-speed cluster, allowed
 trace basis $T_n(t)$, and their symmetric time secants.  Every eigencluster
 must retain the resolvent contour and gap used in
 \eqref{eq:V22-Riesz-derivative-bound}.

 \item Form the orthogonal allowed-trace projector by
 \eqref{eq:V22-trace-graph-projector}.  Store its column floor, derivative
 bound, central-secant error, endpoint route error, Kato generator, and the
 integrated transport debit
 \begin{equation}
  \Delta_{W,n}(T)
  :=
  \int_0^T
  \left[
   2\delta_{1,n}
   +\delta_{0,n}
    \bigl(\|\dot\Pi_n\|+\|\dot{\widehat\Pi}_n\|\bigr)
  \right]\dd t.
  \label{eq:V22-transport-debit}
 \end{equation}

 \item Store $M_n^{-1/2}$, the pulled-back trace projector, the collar
 multiplier $\mathcal W_n$, and the total transport $\mathcal U_n$.
 Directly test
 \[
  \mathcal U_n^*M_n\mathcal U_n=M_n(0),\qquad
  \Gamma_n\mathcal U_n\mathcal D_{n,0}
  \subset\operatorname{Ran}\Pi_n(t),
 \]
 and retain the metric-connection matrix and
 $\kappa_{M,n}$.

 \item On the fixed domain, rerun the V21 adjoint-resolvent range test for
 $\widetilde{\mathcal A}_n(t)+\Omega I$ at every certification time.
 Held-out times must preserve the resolvent kernel margin and the common
 domain.

 \item For every time-pair net entry, issue the resolved commands for
 $\mathfrak G_n(t,s)$ in
 \eqref{eq:V22-relative-graph-variation}.  Store the finite net response,
 target leakage, and high-domain tail.  Report
 \begin{equation}
  L_{n,N}^{\mathrm{cert}}
  :=
  \sup_{t\ne s}
  \frac1{|t-s|}
  \left[
   \frac{d_{N,\eta}(t,s)}{1-\eta}
   +\ell_N^{\mathrm{out}}(t,s)
   +\tau_N^{\mathrm{high}}(t,s)
  \right].
  \label{eq:V22-certified-graph-speed}
 \end{equation}

 \item Store the boundary right inverse, the three terms
 $G'$, $\widetilde{\mathcal A}G$, and the transformed source in
 \eqref{eq:V22-lifted-source}.  Run every main and subsidiary corner
 residual from V21 in the transported frame.

 \item Refine frozen-time partitions and retain the observable residual
 bound
 \begin{equation}
  \epsilon_{\pi,n}^{\mathrm{freeze}}
  :=
  L_{n,N}^{\mathrm{cert}}T C_{*,n}|\pi|.
  \label{eq:V22-freezing-debit}
 \end{equation}
 A causal range pass requires this number and every preceding transport,
 tail, constraint, and corner debit to converge to zero with a uniform
 inverse bound.
\end{enumerate}

The failure witnesses are
\begin{align}
 W_n^{\dot\Pi}
 &:=(t,\widehat{\dot\Pi}_n,\delta_{0,n},\delta_{1,n}),\\
 W_n^{\mathrm{lift}}
 &:=(t,\Gamma_n\mathcal U_n,
       \Pi_n(t),\Delta_{W,n}),\\
 W_n^{M}
 &:=(t,\mathcal U_n^*M_n\mathcal U_n-M_n(0)),\\
 W_n^{\mathrm{dom}}
 &:=(t,u,\Gamma_n\mathcal U_nu,
       (I-\Pi_n(t))\Gamma_n\mathcal U_nu),\\
 W_n^{\mathrm{graph}}
 &:=(t,s,v,\mathfrak G_n(t,s)v),\\
 W_n^{\mathrm{freeze}}
 &:=(\pi,r_{\pi,n},\epsilon_{\pi,n}^{\mathrm{freeze}}),\\
 W_n^{\partial g}
 &:=(g,G',\widetilde{\mathcal A}_nG),\\
 W_{n,k}^{\mathrm{corner}}
 &:=(k,\chi_{n,k}^{\mathrm{corner}}).
 \label{eq:V22-witness-family}
\end{align}
These distinguish projector speed, bulk lifting, metric normalization,
domain transport, graph variation, time freezing, boundary-data lifting,
and corner failures.


\subsection{V22 status and next upstream row}
\label{subsec:V22-status}

V22 proves the following new analytic statements.

\begin{enumerate}
 \item The commutator ODE
 $W_\partial'=[\dot\Pi,\Pi]W_\partial$ transports the allowed trace space
 exactly, and \eqref{eq:V22-transport-error} propagates measured projector
 and derivative errors into the transport.

 \item Spectral flux projectors and incoming-graph projectors have explicit
 derivative bounds.  A pointwise unitary collar lift converts the trace
 transport into a bounded Sobolev-domain transport for the causal
 first-order system.

 \item The moving $A^0$ metric is normalized without discarding its time
 derivative.  The exact connection identity
 \eqref{eq:V22-metric-connection-identity} supplies the added
 quasi-accretive debit.

 \item The frozen-time theorem constructs a nonautonomous evolution family
 on a common domain.  Its proof includes the $O(|\pi|)$ residual
 \eqref{eq:V22-frozen-residual-bound} and the inter-partition Cauchy estimate
 \eqref{eq:V22-partition-Cauchy}.

 \item Relative graph continuity is reduced to the bounded resolvent panel
 \eqref{eq:V22-relative-graph-variation}, with finite-screen, target-tail,
 and high-domain terms.  Inhomogeneous boundary data are included by an
 exact lifting formula.
\end{enumerate}

The attached manuscripts do not print a time-resolved physical allowed-trace
projector for the assembled Standard Model--Einstein system, its derivative,
a boundary-collar lift, the metric connection, the fixed-domain adjoint
resolvents, or the relative graph high tail.  Their finite Kato support
transports and fixed-domain Navier--Stokes evolution family do not populate
these data.  V22 therefore closes the moving-domain implication but does not
claim that the physical causal graph card passes.

With V19--V22, the branch-typed interior symbol, boundary domain, adjoint
range, and moving-domain causal implications are explicit.  The next
upstream task is the \emph{coupled nonlinear assembly}: place the verified
gravity, gauge, Higgs, fermion, ghost, source, and constraint graph inverses
in one common norm; propagate every finite-to-infinite mixed Schur tail and
forced gauge--Ward debit through the moving-domain inverse; and state one
cofinal contraction margin for the full action derivative.  That master
margin, together with actual populated acquisition cards, is the next
necessary step toward the full continuum.

\begin{firewall}[V22 is conditional on the measured graph speed]
\label{fire:V22-status}
Exact projector transport does not control
$(\widetilde{\mathcal A}(t)-\widetilde{\mathcal A}(s))
(I+\widetilde{\mathcal A}(s))^{-1}$.  If the high-domain term in
\eqref{eq:V22-certified-graph-speed} is unbounded, the common-domain
evolution theorem does not apply even when every frozen boundary problem
passes.  That failure must be retained as a time-dependent physical
witness.
\end{firewall}

% END APPENDED VERSION V22 MATERIAL



% =============================================================================
\section{Version V23: full-sector nonlinear assembly in one transported graph norm}
\label{sec:V23-full-sector-assembly}
% =============================================================================

\subsection{Purpose, upstream position, and nonduplication audit}
\label{subsec:V23-purpose}

Versions V17 and V18 leave a precise gap between two correct statements.
The V17 root theorem applies once one bounded common-action graph
isomorphism \(H_*\) and one contraction number are known.  V18 constructs
an exact two-block gravity--matter Schur alternative once the diagonal
operators and the two aggregate mixed maps are already operators on fixed
graph spaces.  Versions V19--V22 then resolve the interior symbol, boundary,
adjoint-range, and moving-time-domain implications.  They do not place the
gravity, Yang--Mills, Higgs, fermion, ghost/BRST, retained source, and
subsidiary constraint blocks in one norm, nor do they propagate all
finite-screen, cofinal, and field-dependent transport errors into the V17
number.

This version supplies that missing implication.  It first trivializes the
field-dependent residual bundle and displays the derivative of that
trivialization.  It then gives a weighted block majorant, a full-sector
finite-to-infinite inverse theorem, and one nonlinear master margin.  Finally
it passes the forced gauge--Ward and constraint sources through the same
moving-domain inverse.  These are analytic implications.  The physical
matrices which would make their hypotheses true are not present in the
attached files, so no physical pass is asserted.

The full-primitive Schur shorts in the attached SM--Einstein note concern
finite positive source Grams.  The V18 Schur theorem concerns two graph
blocks.  Neither is the field-bundle, seven-sector graph estimate below.
The Perron--Doob constructions in the grand manuscript concern positive
path generators and are likewise distinct from the norm weights used here.


% -----------------------------------------------------------------------------
\subsection{The nonlinear residual is a section of a moving graph bundle}
\label{subsec:V23-residual-bundle}
% -----------------------------------------------------------------------------

At cutoff \(n\), let \(\mathscr C_n\) be the reduced configuration manifold
after the declared gauge, relation, line, and boundary choices.  The
common-action residual is naturally a section
\[
 \mathfrak F_n:\mathscr C_n\longrightarrow\mathbb Y_n,
 \qquad
 \mathfrak F_n(q)\in\mathbb Y_{n,q},
\]
rather than a map into one fixed vector space.  On an affine ball
\(\mathscr U_R\subset\mathscr X\), choose a \(C^1\) chart
\[
 \chi_n:\mathscr U_R\longrightarrow\mathscr C_n,
 \qquad
 J_n(z):=D\chi_n(z):\mathscr X\longrightarrow T_{\chi_n(z)}\mathscr C_n,
\]
and a \(C^1\) target trivialization
\[
 V_n(z):\mathbb Y_{n,\chi_n(z)}\longrightarrow\mathscr Y.
\]
The V22 metric and collar maps, the source--initial--boundary graph
coordinates, and the adjacent-cutoff relation and determinant-line
transports are included in \(J_n\) and \(V_n\).  Define the fixed-frame
residual
\begin{equation}
 \mathcal F_n(z)
 :=
 V_n(z)\mathfrak F_n(\chi_n(z)).
 \label{eq:V23-pulled-residual}
\end{equation}

\begin{lemma}[Exact connection term in a residual trivialization]
\label{lem:V23-trivialized-derivative}
For \(h\in\mathscr X\),
\begin{equation}
 \boxed{
 D\mathcal F_n(z)h
 =
 \bigl(DV_n(z)[h]\bigr)\mathfrak F_n(\chi_n(z))
 +
 V_n(z)D\mathfrak F_n(\chi_n(z))J_n(z)h.}
 \label{eq:V23-trivialized-derivative}
\end{equation}
If a second target frame is
\(\overline V_n(z)=S_n(z)V_n(z)\), where \(S_n(z)\) is boundedly
invertible, then
\begin{align}
 \overline{\mathcal F}_n(z)&=S_n(z)\mathcal F_n(z),
 \label{eq:V23-frame-residual}\\
 D\overline{\mathcal F}_n(z)h
 &=DS_n(z)[h]\mathcal F_n(z)
   +S_n(z)D\mathcal F_n(z)h.
 \label{eq:V23-frame-derivative}
\end{align}
Thus the zero set is frame independent, and at a root \(z_*\),
\begin{equation}
 D\overline{\mathcal F}_n(z_*)
 =S_n(z_*)D\mathcal F_n(z_*).
 \label{eq:V23-root-frame-conjugacy}
\end{equation}
Away from a root, the first term in
\eqref{eq:V23-frame-derivative} is a genuine nonlinear-radius debit.
\end{lemma}

\begin{proof}
Apply the Banach-space product and chain rules to
\eqref{eq:V23-pulled-residual}.  The same product rule applied to
\(S_n\mathcal F_n\) gives
\eqref{eq:V23-frame-derivative}.  Bounded invertibility of \(S_n(z)\)
makes \(\overline{\mathcal F}_n(z)=0\) equivalent to
\(\mathcal F_n(z)=0\).  At such a zero the term containing \(DS_n\)
vanishes, which proves \eqref{eq:V23-root-frame-conjugacy}.
\end{proof}

\begin{firewall}[A transported Hessian without the connection term is not the pulled derivative]
\label{fire:V23-transport-connection}
The second term of \eqref{eq:V23-trivialized-derivative} is the transported
physical linearization.  It is not the whole derivative unless the target
frame is constant or the residual is already zero.  Central jets used on
the V17 ball must be jets of the complete map
\eqref{eq:V23-pulled-residual}.  In particular, derivatives of the V22
metric square root, allowed-trace projector, collar lift, and target graph
coordinates must either be included in the endpoint transport or charged
as the first term of \eqref{eq:V23-trivialized-derivative}.
\end{firewall}


% -----------------------------------------------------------------------------
\subsection{Sector normalization and a weighted block majorant}
\label{subsec:V23-weighted-blocks}
% -----------------------------------------------------------------------------

Let
\begin{equation}
 \mathfrak S
 =
 \{\mathrm g,\mathrm{YM},\mathrm H,\mathrm f,
   \mathrm{gh},\mathrm{src},\mathrm c\}
 \label{eq:V23-sector-set}
\end{equation}
denote the gravitational, Yang--Mills, Higgs, realified fermion,
ghost/BRST, retained source/response, and subsidiary gauge/constraint
sectors.  If a source is externally prescribed, its domain factor is
\(\{0\}\) and it remains only as a target row.  Put
\[
 \mathscr X=\bigoplus_{\alpha\in\mathfrak S}\mathscr X_\alpha,
 \qquad
 \mathscr Y=\bigoplus_{\alpha\in\mathfrak S}\mathscr Y_\alpha.
\]
All sums are finite.  The spaces are real Hilbert graph spaces; complex
bosonic and fermionic variables are realified before this decomposition.

For every sector choose a verified diagonal graph isomorphism
\[
 L_\alpha:\mathscr X_\alpha\longrightarrow\mathscr Y_\alpha,
 \qquad
 \mathscr L:=\operatorname{diag}_{\alpha\in\mathfrak S}L_\alpha.
\]
On the causal branch, \(L_\alpha^{-1}\) is the physical source--initial--
boundary inverse obtained after V22 domain transport.  On an elliptic
branch it is the corresponding V18--V21 boundary inverse.  Choose weights
\(w_\alpha>0\) and define
\begin{align}
 \|x\|_w^2
 &:=
 \sum_{\alpha\in\mathfrak S}
 w_\alpha^{-2}\|x_\alpha\|_{\mathscr X_\alpha}^2,
 \label{eq:V23-domain-weighted-norm}\\
 \|y\|_{\mathscr L,w}^2
 &:=
 \sum_{\alpha\in\mathfrak S}
 w_\alpha^{-2}
 \|L_\alpha^{-1}y_\alpha\|_{\mathscr X_\alpha}^2.
 \label{eq:V23-target-weighted-norm}
\end{align}
Then \(\mathscr L:(\mathscr X,\|\cdot\|_w)\to
(\mathscr Y,\|\cdot\|_{\mathscr L,w})\) is an isometry.  These are honest
common norms.  If \(w_{\min}\) and \(w_{\max}\) are the extreme weights,
then
\begin{align}
 w_{\max}^{-1}\|x\|_\oplus
 &\le\|x\|_w\le w_{\min}^{-1}\|x\|_\oplus,
 \label{eq:V23-domain-norm-equivalence}\\
 \frac{\|y\|_\oplus}{w_{\max}\|\mathscr L\|}
 &\le\|y\|_{\mathscr L,w}
 \le\frac{\|\mathscr L^{-1}\|}{w_{\min}}\|y\|_\oplus.
 \label{eq:V23-target-norm-equivalence}
\end{align}
Thus no sector is removed by weighting.

Let \(W=\operatorname{diag}(w_\alpha I_{\mathscr X_\alpha})\).  For a
bounded block operator \(T=(T_{\alpha\beta})\) on \(\mathscr X\), define
the nonnegative scalar matrix
\begin{equation}
 G_w(T)_{\alpha\beta}
 :=
 \frac{w_\beta}{w_\alpha}\,
 \|T_{\alpha\beta}\|_{\mathscr X_\beta\to\mathscr X_\alpha}.
 \label{eq:V23-block-majorant}
\end{equation}

\begin{lemma}[Weighted Hilbert block majorant]
\label{lem:V23-block-majorant}
Every bounded block operator satisfies
\begin{equation}
 \boxed{
 \|T\|_{w\to w}
 =
 \|W^{-1}TW\|_{\oplus\to\oplus}
 \le
 \|G_w(T)\|_{\ell^2(\mathfrak S)\to\ell^2(\mathfrak S)}.}
 \label{eq:V23-block-majorant-bound}
\end{equation}
The same conclusion holds when the entries in
\eqref{eq:V23-block-majorant} are certified upper bounds rather than exact
block norms.
\end{lemma}

\begin{proof}
For \(x\in\mathscr X\), put
\(u_\beta=\|x_\beta\|/w_\beta\).  The triangle inequality gives
\[
 \frac{\|(Tx)_\alpha\|}{w_\alpha}
 \le
 \sum_\beta
 \frac{w_\beta}{w_\alpha}\|T_{\alpha\beta}\|u_\beta
 =
 (G_w(T)u)_\alpha.
\]
Taking the Euclidean norm over \(\alpha\) proves
\[
 \|Tx\|_w
 \le\|G_w(T)u\|_{\ell^2}
 \le\|G_w(T)\|_{2\to2}\|u\|_{\ell^2}
 =\|G_w(T)\|_{2\to2}\|x\|_w.
\]
Replacing exact entries by larger nonnegative entries preserves the
componentwise inequality.
\end{proof}

\begin{corollary}[Constructive weights from a sector gain matrix]
\label{cor:V23-Perron-weights}
Let \(G\) be a finite nonnegative matrix of unscaled sector gains.  If
\(\rho(G)<1\), then for every \(\rho(G)<q<1\) there are explicit positive
weights satisfying
\begin{equation}
 Gw<qw.
 \label{eq:V23-weighted-row-contraction}
\end{equation}
Consequently a block operator \(T\) with
\(\|T_{\alpha\beta}\|\le G_{\alpha\beta}\) has norm less than \(q\) in
the weighted maximum norm
\(\|x\|_{w,\infty}=\max_\alpha\|x_\alpha\|/w_\alpha\).
In particular, \(I+T\) is invertible.
\end{corollary}

\begin{proof}
Since \(\rho(q^{-1}G)<1\), the finite-dimensional Neumann series
\[
 w:=(qI-G)^{-1}{\bf1}
   =q^{-1}\sum_{k\ge0}q^{-k}G^k{\bf1}
\]
converges and has strictly positive entries.  The identity
\((qI-G)w={\bf1}\) gives \(Gw=qw-{\bf1}<qw\).
The weighted row estimate yields
\[
 \|T\|_{w,\infty}
 \le\max_\alpha\frac{(Gw)_\alpha}{w_\alpha}<q.
\]
The Neumann series for \(I+T\) then converges.  This is merely a
constructive sufficient branch; the exact singular-floor theorem below
also covers invertible assemblies for which this gain test fails.
\end{proof}


% -----------------------------------------------------------------------------
\subsection{A full-sector finite-to-infinite graph inverse}
\label{subsec:V23-full-inverse}
% -----------------------------------------------------------------------------

Fix a reference coordinate \(z_{\rm ref}\in\mathscr U_R\).  For the complete
pulled residual \(\mathcal F_n\), define
\begin{equation}
 M_n(z)
 :=
 \mathscr L^{-1}D\mathcal F_n(z)
 =
 \bigl(M_{n,\alpha\beta}(z)\bigr)_{\alpha,\beta\in\mathfrak S}.
 \label{eq:V23-preconditioned-derivative}
\end{equation}
Its reference correction blocks are
\begin{equation}
 K_{n,\alpha\beta}
 :=
 M_{n,\alpha\beta}(z_{\rm ref})
 -\delta_{\alpha\beta}I_{\mathscr X_\beta}.
 \label{eq:V23-reference-correction}
\end{equation}
Thus diagonal drift, every off-diagonal common-action variation, and the
transport-connection term in
\eqref{eq:V23-trivialized-derivative} are treated by the same operator
calculus.

Assume the adjacent transported derivative debits
\begin{equation}
 d^{D}_{j,\alpha\beta}
 :=
 \|K_{j+1,\alpha\beta}-K_{j,\alpha\beta}\|
 \label{eq:V23-adjacent-derivative-debit}
\end{equation}
are summable for every pair.  Then
\(K_{n,\alpha\beta}\) has an operator-norm limit
\(K_{\infty,\alpha\beta}\), and
\begin{equation}
 \|K_{\infty,\alpha\beta}-K_{n,\alpha\beta}\|
 \le
 d^{D,\infty}_{n,\alpha\beta}
 :=
 \sum_{j\ge n}d^{D}_{j,\alpha\beta}.
 \label{eq:V23-derivative-tail}
\end{equation}
Put \(M_\infty=I+K_\infty\).

Let \(P_{\alpha,N}\) be finite-rank orthogonal graph projections and
\(P_N=\bigoplus_\alpha P_{\alpha,N}\).  A measured finite block
\(\widehat K_{n,N,\alpha\beta}\) maps
\(P_{\beta,N}\mathscr X_\beta\) to
\(P_{\alpha,N}\mathscr X_\alpha\) and is extended by zero.  Retain the
three distinct errors
\begin{align}
 r^{\rm read}_{n,N,\alpha\beta}
 &:=
 \|P_{\alpha,N}K_{n,\alpha\beta}P_{\beta,N}
       -\widehat K_{n,N,\alpha\beta}\|,
 \label{eq:V23-block-read-error}\\
 \ell^{\rm out}_{n,N,\alpha\beta}
 &:=
 \|(I-P_{\alpha,N})K_{n,\alpha\beta}P_{\beta,N}\|,
 \label{eq:V23-block-target-leakage}\\
 \tau^{\rm high}_{n,N,\alpha\beta}
 &:=
 \|K_{n,\alpha\beta}(I-P_{\beta,N})\|.
 \label{eq:V23-block-domain-tail}
\end{align}
The read error includes endpoint alignment, V22 moving-domain transport,
and finite numerical error.  The other two quantities are operator tails,
not sampled matrix residuals.  Set
\begin{align}
 E_{n,N,\alpha\beta}
 &:=
 r^{\rm read}_{n,N,\alpha\beta}
 +\ell^{\rm out}_{n,N,\alpha\beta}
 +\tau^{\rm high}_{n,N,\alpha\beta}
 +d^{D,\infty}_{n,\alpha\beta},
 \label{eq:V23-block-total-error}\\
 \mathbf E_{n,N}^{w}(\alpha,\beta)
 &:=
 \frac{w_\beta}{w_\alpha}E_{n,N,\alpha\beta},
 \qquad
 e_{n,N}^{w}:=\|\mathbf E_{n,N}^{w}\|_{2\to2}.
 \label{eq:V23-full-tail-majorant}
\end{align}
Finally define
\begin{equation}
 \widehat M_{n,N}:=I+\widehat K_{n,N},
 \qquad
 \sigma_{n,N}^{w}
 :=
 s_{\min}\!\left(W^{-1}\widehat M_{n,N}W\right).
 \label{eq:V23-finite-full-floor}
\end{equation}
Because \(\widehat K_{n,N}=P_N\widehat K_{n,N}P_N\), this singular
floor is the minimum of one and the least singular value of one finite
matrix.

\begin{theorem}[Full-sector finite-to-infinite inverse certificate]
\label{thm:V23-full-sector-inverse}
For every populated card above,
\begin{equation}
 \boxed{
 \|M_\infty-\widehat M_{n,N}\|_{w\to w}
 \le e_{n,N}^{w}.}
 \label{eq:V23-full-operator-error}
\end{equation}
If
\begin{equation}
 \gamma_{n,N}^{w}
 :=
 \sigma_{n,N}^{w}-e_{n,N}^{w}>0,
 \label{eq:V23-inverse-margin}
\end{equation}
then \(M_\infty\) and
\begin{equation}
 H_*:=\mathscr L M_\infty:
 (\mathscr X,\|\cdot\|_w)
 \longrightarrow
 (\mathscr Y,\|\cdot\|_{\mathscr L,w})
 \label{eq:V23-master-preconditioner}
\end{equation}
are bounded isomorphisms, with
\begin{align}
 \|M_\infty^{-1}\|_{w\to w}
 &\le(\gamma_{n,N}^{w})^{-1},
 \label{eq:V23-normalized-inverse-bound}\\
 \|H_*^{-1}\|_{\mathscr L,w\to w}
 &\le(\gamma_{n,N}^{w})^{-1}.
 \label{eq:V23-master-inverse-bound}
\end{align}
In the original direct-sum norms this implies
\begin{equation}
 \|H_*^{-1}\|_{\oplus\to\oplus}
 \le
 \frac{w_{\max}}{w_{\min}}\,
 \frac{\|\mathscr L^{-1}\|}{\gamma_{n,N}^{w}}.
 \label{eq:V23-original-norm-inverse}
\end{equation}
A small \(\sigma_{n,N}^{w}\) returns a finite mixed sector vector.  An
uncontrolled \(e_{n,N}^{w}\) returns the first named read, target-leakage,
high-domain, or adjacent-derivative block; it does not prove a kernel.
\end{theorem}

\begin{proof}
For each block, the exact identity
\[
 K_{n,\alpha\beta}
 -P_{\alpha,N}K_{n,\alpha\beta}P_{\beta,N}
 =
 (I-P_{\alpha,N})K_{n,\alpha\beta}P_{\beta,N}
 +K_{n,\alpha\beta}(I-P_{\beta,N})
\]
gives
\[
 \|K_{\infty,\alpha\beta}
      -\widehat K_{n,N,\alpha\beta}\|
 \le E_{n,N,\alpha\beta}
\]
after adding \eqref{eq:V23-derivative-tail} and the read error.
Lemma~\ref{lem:V23-block-majorant}, applied to these upper bounds, proves
\eqref{eq:V23-full-operator-error}.

The operator \(\widehat M_{n,N}\) is the identity on
\((I-P_N)\mathscr X\) and a finite matrix on \(P_N\mathscr X\).
Consequently \(\sigma_{n,N}^{w}>0\) makes it invertible and
\[
 \|\widehat M_{n,N}^{-1}\|_{w\to w}
 =(\sigma_{n,N}^{w})^{-1}.
\]
By \eqref{eq:V23-full-operator-error},
\[
 \|\widehat M_{n,N}^{-1}
      (M_\infty-\widehat M_{n,N})\|_{w\to w}
 \le
 \frac{e_{n,N}^{w}}{\sigma_{n,N}^{w}}<1.
\]
The Neumann perturbation lemma gives
\[
 \|M_\infty^{-1}\|_{w\to w}
 \le
 \frac{(\sigma_{n,N}^{w})^{-1}}
 {1-e_{n,N}^{w}/\sigma_{n,N}^{w}}
 =
 \frac1{\gamma_{n,N}^{w}}.
\]
Since \(\mathscr L\) is an isometry between the normalized spaces,
\(H_*^{-1}=M_\infty^{-1}\mathscr L^{-1}\) proves
\eqref{eq:V23-master-inverse-bound}.  The two norm equivalences
\eqref{eq:V23-domain-norm-equivalence}--%
\eqref{eq:V23-target-norm-equivalence} give
\eqref{eq:V23-original-norm-inverse}.  The witness statements follow
from the finite singular-value decomposition and the four summands in
\eqref{eq:V23-block-total-error}.
\end{proof}

\begin{remark}[Relation with the V18 two-block Schur test]
\label{rem:V23-Schur-relation}
For two sectors, exact blocks, and no tail,
Theorem~\ref{thm:V23-full-sector-inverse} tests the least singular value
of the complete normalized block matrix.  The V18 Schur factorization is
an algebraically sharper way to evaluate the same invertibility question
when its diagonal blocks have already been inverted.  V23 adds the
remaining sectors, the common weights, the field and cutoff transports,
and every blockwise finite-to-infinite debit.  Either finite inverse
calculation may be used, but its error must still be bounded by
\eqref{eq:V23-full-tail-majorant}.
\end{remark}

\begin{firewall}[A finite coupled matrix does not control its high graph complement]
\label{fire:V23-finite-not-infinite}
The number \(\sigma_{n,N}^{w}\) contains no information about
\(\tau^{\rm high}_{n,N,\alpha\beta}\).  Setting those terms to zero
because a command was not issued would replace
\eqref{eq:V23-full-operator-error} by a false finite-dimensional
inference.  When a high-domain tail is not bounded, the correct result is
an unresolved block and its near-maximizing graph vector.
\end{firewall}


% -----------------------------------------------------------------------------
\subsection{One nonlinear radius for the complete action derivative}
\label{subsec:V23-nonlinear-radius}
% -----------------------------------------------------------------------------

Fix one passing inverse card \((n_0,N,w)\) and abbreviate
\[
 \sigma:=\sigma_{n_0,N}^{w},
 \qquad
 e:=e_{n_0,N}^{w},
 \qquad
 \gamma:=\sigma-e>0.
\]
For \(n\ge n_0\), define the block oscillation on the declared affine ball
\begin{equation}
 \nu_{n,\alpha\beta}(R)
 :=
 \sup_{z\in\mathscr U_R}
 \|M_{n,\alpha\beta}(z)
      -M_{n,\alpha\beta}(z_{\rm ref})\|.
 \label{eq:V23-block-nonlinear-oscillation}
\end{equation}
Form the nonnegative matrix
\begin{equation}
 \mathbf A_{n,R}^{w}(\alpha,\beta)
 :=
 \frac{w_\beta}{w_\alpha}
 \left(
  d^{D,\infty}_{n,\alpha\beta}
  +\nu_{n,\alpha\beta}(R)
 \right)
 \label{eq:V23-nonlinear-majorant-matrix}
\end{equation}
and the scalar quantities
\begin{align}
 \Theta_R
 &:=
 \frac1{\gamma}
 \sup_{n\ge n_0}\|\mathbf A_{n,R}^{w}\|_{2\to2},
 \label{eq:V23-master-theta}\\
 \mathfrak b
 &:=
 \frac1{\gamma}
 \sup_{n\ge n_0}
 \|\mathscr L^{-1}\mathcal F_n(z_{\rm ref})\|_w,
 \label{eq:V23-master-center}\\
 a_n
 &:=
 \sup_{z\in\mathscr U_R}
 \|\mathscr L^{-1}
   (\mathcal F_{n+1}(z)-\mathcal F_n(z))\|_w.
 \label{eq:V23-adjacent-action-debit}
\end{align}
When \(\sigma>0\), the dimensionless master margin is
\begin{equation}
 \boxed{
 \mathfrak m_{23}(R;n_0,N,w)
 :=
 \min\left\{
  1-\frac e\sigma,\
  1-\Theta_R-\frac{\mathfrak b}{R}
 \right\}.}
 \label{eq:V23-master-margin}
\end{equation}
The first entry is the finite-to-infinite inverse reserve.  The second is
simultaneously the contraction and ball-invariance reserve.

\begin{lemma}[Second-derivative majorant for the nonlinear blocks]
\label{lem:V23-Hessian-majorant}
Suppose the complete pulled derivative is differentiable and, for
\(z\in\mathscr U_R\),
\begin{equation}
 \|D_\delta M_{n,\alpha\beta}(z)h_\delta\|_{\op}
 \le
 h_{n,\alpha\beta\delta}\|h_\delta\|_{\mathscr X_\delta},
 \qquad h_{n,\alpha\beta\delta}\ge0.
 \label{eq:V23-third-block-bound}
\end{equation}
Then
\begin{equation}
 \boxed{
 \nu_{n,\alpha\beta}(R)
 \le
 R\sum_{\delta\in\mathfrak S}
 h_{n,\alpha\beta\delta}w_\delta.}
 \label{eq:V23-nonlinear-oscillation-bound}
\end{equation}
The constants in \eqref{eq:V23-third-block-bound} must be taken after the
field-dependent transports in
\eqref{eq:V23-trivialized-derivative}; raw physical Hessian bounds alone
do not control this quantity.
\end{lemma}

\begin{proof}
For \(z=z_{\rm ref}+u\in\mathscr U_R\), the weighted ball gives
\(\|u_\delta\|\le w_\delta R\).  Integrating the derivative along
\(z_{\rm ref}+tu\) yields
\[
 M_{n,\alpha\beta}(z)-M_{n,\alpha\beta}(z_{\rm ref})
 =
 \int_0^1
 \sum_\delta
 D_\delta M_{n,\alpha\beta}(z_{\rm ref}+tu)u_\delta\,\dd t.
\]
Use \eqref{eq:V23-third-block-bound}, the triangle inequality, and the
component bounds.  Taking the supremum proves
\eqref{eq:V23-nonlinear-oscillation-bound}.
\end{proof}

\begin{theorem}[Full-sector cofinal contraction theorem]
\label{thm:V23-master-contraction}
Assume:
\begin{enumerate}[label=\textup{(C\arabic*)},leftmargin=3em]
 \item the residuals \(\mathcal F_n:\mathscr U_R\to\mathscr Y\) are
       continuously differentiable in the fixed frames of
       \eqref{eq:V23-pulled-residual};
 \item the finite full-sector card has \(\gamma>0\);
 \item \(\mathfrak m_{23}(R;n_0,N,w)>0\);
 \item
       \begin{equation}
        \sum_{n\ge n_0}a_n<\infty.
        \label{eq:V23-action-summability}
       \end{equation}
\end{enumerate}
Then for every \(n\ge n_0\) there is a unique
\(z_n\in\mathscr U_R\) satisfying
\begin{equation}
 \mathcal F_n(z_n)=0.
 \label{eq:V23-finite-action-root}
\end{equation}
The roots converge in the full weighted graph norm to a limit \(z_\infty\),
and
\begin{equation}
 \boxed{
 \|z_\infty-z_n\|_w
 \le
 \frac1{(1-\Theta_R)\gamma}
 \sum_{j\ge n}a_j.}
 \label{eq:V23-root-tail}
\end{equation}
Moreover, \(\mathscr L^{-1}\mathcal F_n\) converges uniformly on
\(\mathscr U_R\) to a continuous map
\(\mathscr L^{-1}\mathcal F_\infty\), and
\begin{equation}
 \mathcal F_\infty(z_\infty)=0.
 \label{eq:V23-limit-full-action}
\end{equation}
Every sector component, including the mixed and subsidiary rows, is part
of this same zero.

The V17 preconditioned hypotheses hold with the explicit bounds
\begin{align}
 \sup_{n\ge n_0}\sup_{z\in\mathscr U_R}
 \|I-H_*^{-1}D\mathcal F_n(z)\|_{w\to w}
 &\le\Theta_R,
 \label{eq:V23-to-V17-derivative}\\
 \sup_{n\ge n_0}
 \|H_*^{-1}\mathcal F_n(z_{\rm ref})\|_w
 &\le\mathfrak b,
 \label{eq:V23-to-V17-center}\\
 \sup_{z\in\mathscr U_R}
 \|H_*^{-1}(\mathcal F_{n+1}(z)-\mathcal F_n(z))\|_w
 &\le\frac{a_n}{\gamma}.
 \label{eq:V23-to-V17-adjacent}
\end{align}
\end{theorem}

\begin{proof}
Theorem~\ref{thm:V23-full-sector-inverse} gives
\(\|M_\infty^{-1}\|_{w\to w}\le\gamma^{-1}\) and
\(H_*=\mathscr L M_\infty\).  Hence
\begin{align*}
 I-H_*^{-1}D\mathcal F_n(z)
 &=
 I-M_\infty^{-1}M_n(z)\\
 &=
 M_\infty^{-1}\bigl(M_\infty-M_n(z)\bigr).
\end{align*}
The \((\alpha,\beta)\) block of the last difference has norm at most
\[
 d^{D,\infty}_{n,\alpha\beta}
 +\nu_{n,\alpha\beta}(R).
\]
Lemma~\ref{lem:V23-block-majorant} and
\eqref{eq:V23-master-theta} prove
\eqref{eq:V23-to-V17-derivative}.  The inverse bound applied to the
center and adjacent residuals proves
\eqref{eq:V23-to-V17-center} and
\eqref{eq:V23-to-V17-adjacent}.

For \(u\) in the closed weighted radius-\(R\) ball, define
\[
 T_n(u)
 :=
 u-H_*^{-1}\mathcal F_n(z_{\rm ref}+u).
\]
Equation \eqref{eq:V23-to-V17-derivative} gives
\(\|T_n(u)-T_n(v)\|_w\le\Theta_R\|u-v\|_w\).
Also
\[
 \|T_n(u)\|_w
 \le
 \Theta_R\|u\|_w
 +\|H_*^{-1}\mathcal F_n(z_{\rm ref})\|_w
 \le
 \Theta_R R+\mathfrak b<R,
\]
where the final strict inequality is the second entry of
\eqref{eq:V23-master-margin}.  Banach's theorem gives a unique fixed
point in the ball, which is \eqref{eq:V23-finite-action-root}.

For \(z,v\in\mathscr U_R\), integrate the derivative along their segment:
\[
 \|H_*^{-1}(\mathcal F_n(z)-\mathcal F_n(v))\|_w
 \ge
 (1-\Theta_R)\|z-v\|_w.
\]
Use \(z=z_{n+1}\), \(v=z_n\), the two root identities, and
\eqref{eq:V23-to-V17-adjacent} to obtain
\[
 \|z_{n+1}-z_n\|_w
 \le
 \frac{a_n}{(1-\Theta_R)\gamma}.
\]
Summation proves convergence and \eqref{eq:V23-root-tail}.

Finally, \eqref{eq:V23-action-summability} makes
\(\mathscr L^{-1}\mathcal F_n\) uniformly Cauchy on the ball.  Let its
uniform limit be \(\mathscr L^{-1}\mathcal F_\infty\).  Uniform
convergence, continuity of the limit, \(z_n\to z_\infty\), and
\(\mathcal F_n(z_n)=0\) give
\(\mathcal F_\infty(z_\infty)=0\).
\end{proof}

\begin{corollary}[Common \(C^{2,\alpha}\) scheduler and curvature tail]
\label{cor:V23-scheduler-tail}
If the weighted graph space embeds into the V17 local
\(C^{2,\alpha}\) scheduler space with constant \(C_{w,\rm emb}\), then
\begin{equation}
 \|\pi_{\rm sch}(z_\infty-z_n)\|_{C^{2,\alpha}}
 \le
 \frac{C_{w,\rm emb}}{(1-\Theta_R)\gamma}
 \sum_{j\ge n}a_j.
 \label{eq:V23-scheduler-tail}
\end{equation}
If the right side at \(n_0\), plus the reference scheduler radius, is
within the V14 noncollapse ball, all V16 connection, curvature, and
densitized Einstein tails follow on that same cofinal sequence.
\end{corollary}

\begin{proof}
Apply the stated embedding to \eqref{eq:V23-root-tail}, then invoke the
V14--V16 curvature lift with this explicit detail budget.
\end{proof}

\begin{firewall}[Separate sector radii do not imply the master margin]
\label{fire:V23-separate-radii}
Even if every diagonal nonlinear block is a contraction, the matrix
\(\mathbf A_{n,R}^{w}\) can have a large mixed singular direction.
Conversely, a large entrywise upper bound may coexist with an invertible
full finite matrix.  The positive decision is therefore
\eqref{eq:V23-master-margin}, using either the exact full singular floor
or a proved weighted gain contraction.  Independent sector fits cannot
be spliced after the calculation.
\end{firewall}


% -----------------------------------------------------------------------------
\subsection{Forced gauge--Ward and constraint rows in the same norm}
\label{subsec:V23-subsidiary}
% -----------------------------------------------------------------------------

The subsidiary equations must be driven by the residuals which actually
occur, rather than by equations assumed to vanish in advance.  Let
\(\mathscr P_{n}^{\rm gw}(z)\) be the first-order graph operator for the
gauge vector \(C_n\), with the V22 moving boundary and initial data included.
Assume
\begin{equation}
 \|(\mathscr P_{n}^{\rm gw}(z))^{-1}\|
 \le\kappa_{\rm gw}
 \quad
 (n\ge n_0,\ z\in\mathscr U_R).
 \label{eq:V23-gauge-inverse-bound}
\end{equation}
The exact V18 identity, in the fixed frame, has the form
\begin{equation}
 \mathscr P_{n}^{\rm gw}(z)C_n(z)
 =
 -2\left(\mathscr J_n(z)+
          \mathscr D_n\mathscr E_n(z)\right)
 +b_n^{\rm gw}(z),
 \label{eq:V23-forced-gauge-fixed-frame}
\end{equation}
where \(\mathscr J_n\) is the Ward defect,
\(\mathscr E_n\) is the reduced Einstein residual,
\(\mathscr D_n\) is its covariant divergence in the declared source norm,
and \(b_n^{\rm gw}\) contains initial, boundary, collar, and frame-transport
debits.

Assume that bounded target readouts \(Q_n^J,Q_n^E\) and incidence errors
\(r_n^J,r_n^E\) satisfy
\begin{align}
 \mathscr J_n(z)
 &=Q_n^J\mathcal F_n(z)+r_n^J(z),
 \label{eq:V23-Ward-readout}\\
 \mathscr D_n\mathscr E_n(z)
 &=Q_n^E\mathcal F_n(z)+r_n^E(z).
 \label{eq:V23-divEinstein-readout}
\end{align}
These identities require the target norm of \(\mathcal F_n\) to control
the derivative used by \(\mathscr D_n\).

\begin{proposition}[Moving-domain gauge--Ward debit at the nonlinear root]
\label{prop:V23-gauge-Ward-landing}
Under \eqref{eq:V23-gauge-inverse-bound}--%
\eqref{eq:V23-divEinstein-readout},
\begin{align}
 \|C_n(z)\|_{\rm gw}
 \le\kappa_{\rm gw}\Bigl[
 &2(\|Q_n^J\|+\|Q_n^E\|)
     \|\mathcal F_n(z)\|_{\mathscr L,w}
 \nonumber\\
 &+2\|r_n^J(z)\|+2\|r_n^E(z)\|
   +\|b_n^{\rm gw}(z)\|
 \Bigr].
 \label{eq:V23-gauge-Ward-master-debit}
\end{align}
At the roots from Theorem~\ref{thm:V23-master-contraction},
\begin{equation}
 \boxed{
 \|C_n(z_n)\|_{\rm gw}
 \le
 \kappa_{\rm gw}
 \left[
  2\|r_n^J(z_n)\|+2\|r_n^E(z_n)\|
  +\|b_n^{\rm gw}(z_n)\|
 \right].}
 \label{eq:V23-gauge-Ward-root-debit}
\end{equation}
Therefore vanishing incidence, initial, boundary, and transport debits imply
\(C_n(z_n)\to0\).  No Ward equation or reduced Einstein equation is used
before its corresponding row of \(\mathcal F_n\) is controlled.
\end{proposition}

\begin{proof}
Apply \((\mathscr P_n^{\rm gw}(z))^{-1}\) to
\eqref{eq:V23-forced-gauge-fixed-frame}, substitute
\eqref{eq:V23-Ward-readout}--%
\eqref{eq:V23-divEinstein-readout}, and use the triangle inequality and
\eqref{eq:V23-gauge-inverse-bound}.  Equation
\eqref{eq:V23-finite-action-root} removes the two readout terms at \(z_n\).
The limiting statement follows from
\eqref{eq:V23-gauge-Ward-root-debit}.
\end{proof}

The other subsidiary constraints have the same noncircular form.  For a
finite family \(\mathfrak Q\) containing the Yang--Mills Gauss, Einstein
Hamiltonian and momentum, BRST, and any retained boundary constraints,
write
\begin{equation}
 \mathscr P_n^q(z)c_n^q(z)
 =
 S_n^q\mathcal F_n(z)+r_n^q(z)+b_n^q(z),
 \qquad q\in\mathfrak Q,
 \label{eq:V23-general-constraint-equation}
\end{equation}
with
\(\|(\mathscr P_n^q(z))^{-1}\|\le\kappa_q\).

\begin{corollary}[Simultaneous subsidiary landing]
\label{cor:V23-constraint-landing}
At a V23 root,
\begin{equation}
 \|c_n^q(z_n)\|_q
 \le
 \kappa_q\bigl(\|r_n^q(z_n)\|+\|b_n^q(z_n)\|\bigr).
 \label{eq:V23-constraint-root-debit}
\end{equation}
If the right sides tend to zero for every \(q\in\mathfrak Q\), all
subsidiary constraints vanish in their declared limit norms.
\end{corollary}

\begin{proof}
Set \(\mathcal F_n(z_n)=0\) in
\eqref{eq:V23-general-constraint-equation} and apply the inverse bound.
The family is finite, so componentwise convergence is simultaneous.
\end{proof}

\begin{firewall}[The divergence row fixes the target regularity]
\label{fire:V23-divergence-regularity}
If \(\mathscr D_n:\mathscr Y_E\to\mathscr Z_{\rm gw}\) is unbounded in
the target norm used for \(\mathcal F_n\), then
\eqref{eq:V23-divEinstein-readout} is not a bounded readout and
\eqref{eq:V23-gauge-Ward-master-debit} is unavailable.  The valid remedies
are to strengthen the Einstein residual graph norm by one derivative or
to place the subsidiary identity in a matching negative Sobolev space.
Smallness of the undifferentiated residual cannot be substituted.
\end{firewall}


% -----------------------------------------------------------------------------
\subsection{Conditional full Standard--Model--Einstein continuum handoff}
\label{subsec:V23-full-handoff}
% -----------------------------------------------------------------------------

\begin{theorem}[Cofinal full-sector continuum implication on a compact cylinder]
\label{thm:V23-full-continuum-implication}
Let \(K\) be a compact Lorentzian cylinder, or a compact Euclidean region
with one fixed complementing boundary realization.  Suppose:
\begin{enumerate}[label=\textup{(F\arabic*)},leftmargin=3em]
 \item the common-action residual contains source-complete variations in the
       gravity, Yang--Mills, Higgs, realified fermion, ghost/BRST, retained
       source/response, and subsidiary sectors, all in the fixed
       configuration and target frames of
       \eqref{eq:V23-pulled-residual};
 \item the V19 interior symbol, V20 boundary/flux/constraint, V21
       adjoint-range, and V22 metric/domain/evolution cards pass for every
       diagonal sector, giving the \(L_\alpha\) used in
       \eqref{eq:V23-target-weighted-norm};
 \item one cofinal family of full-sector cards has a uniform positive
       \(\mathfrak m_{23}\), satisfies
       \eqref{eq:V23-action-summability}, and its root tail is within the V14
       noncollapse radius;
 \item all V16--V17 common-action metric secants, Palatini incidence errors,
       stress high tails, line orientations, renormalized product limits,
       and OS-to-Lorentz transports are taken on this same cofinal family
       and have the stated vanishing or summable bounds;
 \item the bounded readout identities
       \eqref{eq:V23-Ward-readout}--%
       \eqref{eq:V23-general-constraint-equation} hold, and every incidence,
       initial, boundary, collar, and transport debit on their right sides
       tends to zero.
\end{enumerate}
Then the roots \(z_n\) converge in the common local graph norm to a
noncollapsed field \(z_\infty\).  The limiting source-complete
common-action equation is
\begin{equation}
 \mathcal F_\infty(z_\infty)=0.
 \label{eq:V23-complete-limit-equation}
\end{equation}
Its sector components give the limiting Yang--Mills, Higgs, fermion,
ghost/BRST, source/response, and reduced gravitational Euler equations
in the distribution spaces declared by their graph rows.  The subsidiary
gauge vector and every constraint in \(\mathfrak Q\) vanish.  The reduced
gravitational equation therefore descends to the unreduced equation.
The metric component also satisfies
\begin{equation}
 \boxed{
 2\chi\bigl(G_{\mu\nu}(g_\infty)+\Lambda g_{\infty,\mu\nu}\bigr)
 =
 T_{\mu\nu}^{\rm ren}(z_\infty)}
 \quad\hbox{in }\mathcal D'(K).
 \label{eq:V23-full-Einstein-equation}
\end{equation}
Thus assumptions \textup{(F1)--(F5)} are a sufficient, noncircular
completion criterion for the interacting Standard--Model--Einstein
continuum on \(K\).
\end{theorem}

\begin{proof}
Theorem~\ref{thm:V23-master-contraction} gives the cofinal roots, the graph
tail, and \eqref{eq:V23-complete-limit-equation}.  The noncollapse
condition and Corollary~\ref{cor:V23-scheduler-tail} put the metric and
coframe in the V14--V16 regularity class.  Source completeness and the
product and domain limits in \textup{(F4)} permit passage of every
component of \eqref{eq:V23-complete-limit-equation} to its declared
distributional Euler equation.

Proposition~\ref{prop:V23-gauge-Ward-landing} and
Corollary~\ref{cor:V23-constraint-landing} make the gauge and constraint
fields vanish without using the desired limiting equations as premises.
The V18 gauge reduction then identifies the reduced gravitational
solution with the unreduced one.  On the metric row, the action residual
is zero.  The V16 exact three-defect identity, the vanishing Palatini
incidence error, and the V17 common-action stress landing give
\eqref{eq:V23-full-Einstein-equation}.  The line and OS hypotheses in
\textup{(F4)} identify the same real Lorentzian stress and fermion
orientation throughout.  Every conclusion is therefore on one cofinal
family and in one transported physical frame.
\end{proof}

\begin{corollary}[Compatible compact exhaustion gives a global local solution]
\label{cor:V23-compact-exhaustion}
Let \(K_1\Subset K_2\Subset\cdots\) exhaust the selected spacetime.
Suppose Theorem~\ref{thm:V23-full-continuum-implication} holds on each
\(K_r\), and either the finite roots are restrictions of one cutoff
configuration or the limiting initial-boundary problem is unique on
overlaps.  Then the fields \(z_\infty^{(r)}\) agree on
\(K_r\cap K_s\) and define
\[
 z_\infty\in\mathscr X_{\rm loc}.
\]
All sector equations, subsidiary constraints, and
\eqref{eq:V23-full-Einstein-equation} hold distributionally on the
exhausted spacetime.
\end{corollary}

\begin{proof}
In the restriction branch, convergence in the graph norm on the smaller
compact set identifies both limits with the limit of the same restricted
finite sequence.  In the uniqueness branch, both limits solve the same
limiting initial-boundary problem on the overlap, so uniqueness identifies
them.  The compatible family therefore glues.  Every compactly supported
test section lies in some \(K_r\), where the asserted distributional
identity already holds.
\end{proof}

\begin{assessment}[Logical force of the V23 completion theorem]
\label{ass:V23-logical-force}
Theorem~\ref{thm:V23-full-continuum-implication} closes the implication
from populated physical cards to the full coupled continuum.  It does not
populate those cards.  In particular, the theorem cannot be invoked from
sector names, formal continuum Lagrangians, or finite positive source
Grams.  It requires the occurring pulled common-action jets, diagonal
moving-domain inverses, full mixed matrix, high graph tails, nonlinear
radius, and subsidiary readouts on one cofinal sequence.
\end{assessment}


% -----------------------------------------------------------------------------
\subsection{Executable V23 card and failure witnesses}
\label{subsec:V23-card}
% -----------------------------------------------------------------------------

A V23 acquisition must preserve the following order.

\begin{enumerate}[label=\textup{(\arabic*)},leftmargin=3em]
 \item Freeze one action, signature, physical quotient, line orientation,
       configuration chart, boundary or Cauchy realization, and compact
       cylinder.  Declare the seven sector source and target graph spaces.
 \item Store the field-dependent chart derivative \(J_n(z)\), target
       trivialization \(V_n(z)\), and their first derivatives.  At every
       command endpoint retain both terms of
       \eqref{eq:V23-trivialized-derivative}.
 \item Populate the V19--V22 diagonal graph cards and store
       \(L_\alpha^{-1}\), its physical norm bound, and the input and output
       transports which put it in the fixed frame.
 \item Choose positive weights \(w_\alpha\), print
       \eqref{eq:V23-domain-weighted-norm}--%
       \eqref{eq:V23-target-norm-equivalence}, and retain either the
       constructive gain calculation
       \eqref{eq:V23-weighted-row-contraction} or the exact scaled finite
       singular calculation.
 \item On source-complete sector screens issue central commands for the
       complete pulled residual.  Assemble every block of
       \(K_{n,\alpha\beta}\), including diagonal drift, mixed action
       derivatives, and transport-connection derivatives.  Hold out at least
       one command in every nonzero block.
 \item Store \(r^{\rm read}_{n,N,\alpha\beta}\),
       \(\ell^{\rm out}_{n,N,\alpha\beta}\),
       \(\tau^{\rm high}_{n,N,\alpha\beta}\), and
       \(d^{D}_{j,\alpha\beta}\) separately.  Report the matrix
       \(\mathbf E_{n,N}^{w}\), its norm \(e_{n,N}^{w}\), the finite floor
       \(\sigma_{n,N}^{w}\), and \(\gamma_{n,N}^{w}\).
 \item Bound the complete pulled second derivative on the nonlinear ball.
       The finite command remainder, target leakage, high-domain tail, and
       field-transport derivative are all included in
       \(h_{n,\alpha\beta\delta}\).  Form
       \(\mathbf A_{n,R}^{w}\) and \(\Theta_R\).
 \item Read the unnormalized center residual and adjacent common-action
       differences in the same target norm.  Report \(\mathfrak b\), every
       \(a_n\), their cofinal sum, and the root-tail bound
       \eqref{eq:V23-root-tail}.
 \item Build the differentiated Einstein and Ward readouts in the declared
       source norm.  Rerun their V21--V22 adjoint-range, boundary, corner,
       and transport cards, then report
       \eqref{eq:V23-gauge-Ward-root-debit} and every constraint debit
       \eqref{eq:V23-constraint-root-debit}.
 \item On the identical cofinal roots run the V16--V17 Palatini, metric
       secant, stress tail, product, determinant-line, and OS-to-Lorentz
       cards.  No independently fitted matter or stress history may be
       inserted.
\end{enumerate}

The scalar decision printed by the nonlinear assembly is
\begin{equation}
 \left(
 \mathfrak m_{23}(R;n_0,N,w),\
 \sum_{n\ge n_0}a_n,\
 \frac1{(1-\Theta_R)\gamma}
 \sum_{n\ge n_0}a_n
 \right).
 \label{eq:V23-executable-output}
\end{equation}
A positive first entry, finite second entry, and an admissible third entry
are respectively the inverse/radius pass, cofinal action transport, and
noncollapse detail budget.  The subsidiary and stress rows remain
separate required outputs.

Canonical failure witnesses include
\begin{align}
 W_n^{\rm conn}
 &:=
 \left(z,h,\,
 (DV_n(z)[h])\mathfrak F_n(\chi_n(z))\right),
 \label{eq:V23-connection-witness}\\
 W_{n,N}^{\rm mix}
 &:=
 \left(W^{-1}\widehat M_{n,N}W,\,
 \widehat x_{\min},\,
 \sigma_{n,N}^{w}\right),
 \label{eq:V23-mixed-witness}\\
 W_{n,N,\alpha\beta}^{\rm tail}
 &:=
 \left(\alpha,\beta,\,
 r^{\rm read},\ell^{\rm out},\tau^{\rm high},
 d^{D,\infty}\right),
 \label{eq:V23-tail-witness}\\
 W_{n,\alpha\beta\delta}^{\rm nl}
 &:=
 \left(z,h_\delta,\,
 D_\delta M_{n,\alpha\beta}(z)h_\delta\right),
 \label{eq:V23-nonlinear-witness}\\
 W_n^{\rm ctr}
 &:=
 \left(\mathscr L^{-1}\mathcal F_n(z_{\rm ref}),
       \mathfrak b,R,\Theta_R\right),
 \label{eq:V23-center-witness}\\
 W_n^{\rm gw}
 &:=
 \left(\mathscr J_n-Q_n^J\mathcal F_n,\,
       \mathscr D_n\mathscr E_n-Q_n^E\mathcal F_n,\,
       b_n^{\rm gw}\right).
 \label{eq:V23-gauge-Ward-witness}
\end{align}
Here \(\widehat x_{\min}\) is a least right singular vector.  For a tail
norm which is not attained, a graph-normalized near-maximizer with its
achieved response is retained.  These witnesses distinguish a bad
trivialization, a genuine mixed soft mode, missing graph exhaustion,
nonlinear curvature of the residual, an excessive center force, and a
forced subsidiary source.


% -----------------------------------------------------------------------------
\subsection{V23 theorem ledger and next upstream row}
\label{subsec:V23-status}
% -----------------------------------------------------------------------------

\paragraph{New conclusions proved in V23.}
Lemma~\ref{lem:V23-trivialized-derivative} gives the exact derivative of
the field-dependent residual trivialization and isolates the connection
term which disappears only at a root.  Lemma~\ref{lem:V23-block-majorant}
places every sector in one weighted Hilbert graph norm, while
Corollary~\ref{cor:V23-Perron-weights} gives constructive weights on the
small-gain branch.  Theorem~\ref{thm:V23-full-sector-inverse} promotes one
full finite sector matrix to an infinite graph isomorphism after charging
read, target, domain, and cofinal derivative tails.
Theorem~\ref{thm:V23-master-contraction} compresses the complete inverse,
nonlinear oscillation, center, and adjacent-action data into
\(\mathfrak m_{23}\) and proves the explicit cofinal root tail.
Proposition~\ref{prop:V23-gauge-Ward-landing} and
Corollary~\ref{cor:V23-constraint-landing} pass every forced subsidiary
source through the moving-domain inverse without circular use of the
desired equations.  Theorem~\ref{thm:V23-full-continuum-implication}
then states and proves the single-sequence implication to all
Standard--Model Euler rows and the distributional Einstein equation.

\paragraph{Results imported without repetition.}
The finite common-action relation squares, source-Gram Schur shorts,
support and line transports, scheduler noncollapse, curvature lift,
Palatini incidence, stress reconstruction, graph root theorem, two-block
Schur factorization, symbol extraction, boundary complementing/flux
tests, adjoint range, and time-dependent domain transport remain the
results of the attached sources and Versions V1--V22.  V23 uses their
outputs and proves only the missing field-bundle and full-sector
assembly.

\paragraph{Authoritative physical status.}
The attached manuscripts do not print a single source-complete,
field-dependent matrix for
\(D\mathcal F_n\) in the seven fixed graph sectors.  They do not provide
the derivatives of the assembled metric/trace/collar trivialization over
a nonlinear ball, the forty-nine block read and tail bounds, a scaled
full-sector singular floor, or the differentiated Einstein and Ward
readout incidence in one source norm.  Consequently
\(\mathfrak m_{23}\) cannot be evaluated from the supplied tables.
The hypotheses of
Theorem~\ref{thm:V23-full-continuum-implication} are not physically
populated, and the full Standard--Model--Einstein continuum is not yet
claimed.

\paragraph{Next upstream acquisition.}
The sharp next row is the parameter-dependent domain calculus.  Starting
from the occurring common-action coefficients, one must differentiate the
V22 Riesz allowed-trace projector, the metric square root, the Kato
collar lift, and the source--initial--boundary graph coordinates with
respect to each of the seven field sectors.  Those derivatives must have
tame Sobolev bounds on the V23 ball and must be compared with central
jets of the complete pulled action.  Only then can the forty-nine blocks
of \(K_n\), their high graph tails, and
\(\mathfrak m_{23}\) be populated.  If differentiability loses a
derivative or becomes circular, the next version must go upstream to a
Nash--Moser scale or to a first-order enlarged system rather than report
a Banach contraction.

\begin{firewall}[A conditional completion criterion is not a populated continuum]
\label{fire:V23-status}
The conclusion of
Theorem~\ref{thm:V23-full-continuum-implication} is rigorous under its
listed hypotheses.  The absence of a counterexample to those hypotheses
is not evidence that they hold.  Until the physical card prints a
positive \(\mathfrak m_{23}\), summable adjacent actions, vanishing
subsidiary debits, and the V16--V17 stress/Palatini passage on the same
roots, the continuum remains an open programme.
\end{firewall}

% END APPENDED VERSION V23 MATERIAL



% =============================================================================
\section{Version V24: field-parameter domain calculus and tame collar charts}
\label{sec:V24-parameter-domain}
% =============================================================================

\subsection{Purpose and exact relation to the earlier transport results}
\label{subsec:V24-purpose}

The V23 derivative is a derivative on configuration space, whereas V22
controls motion along one prescribed time path.  Smooth dependence on time
does not by itself give a \(C^2\) chart over the full nonlinear ball.  The
missing data are directional derivatives, in every gravity, gauge, Higgs,
fermion, ghost, source, and constraint direction, of the metric
normalization, allowed-trace projector, and collar lift.

V1 already proves that polar chords around a field-space plaquette retain
the Kato curvature.  That curvature must not be set to zero.  V24 instead
uses the direct polar chord from one fixed base projector to each nearby
projector, together with the complementary chord, to obtain a full unitary
chart.  Mixed derivatives of this chart remain order dependent through the
V1 curvature identities.  V22's pathwise commutator ODE and nonautonomous
evolution theorem are imported, not repeated.

This version proves quantitative first- and second-derivative formulae for
positive inverse square roots, Riesz projectors, trace-graph projectors, the
fixed-base unitary, and its unitary collar extension.  It then identifies
the precise Sobolev regularity at which these operations populate the V23
transport-connection blocks.  A Fourier separator shows when a same-space
Banach contraction is impossible because the coefficients lose
derivatives.


% -----------------------------------------------------------------------------
\subsection{Positive metric normalization with two parameter derivatives}
\label{subsec:V24-metric-calculus}
% -----------------------------------------------------------------------------

Let \(\mathcal B\) be an open subset of a real Banach space \(\mathscr Z\),
let \(\mathscr H\) be a Hilbert space, and suppose
\[
 M:\mathcal B\longrightarrow\mathcal B(\mathscr H)
\]
is \(C^2\), self-adjoint, and satisfies the uniform floor
\begin{equation}
 M(z)\succeq m_*I,\qquad m_*>0.
 \label{eq:V24-metric-floor}
\end{equation}
For directions \(h,k\in\mathscr Z\), write
\(M_h=DM(z)[h]\) and \(M_{hk}=D^2M(z)[h,k]\).

\begin{theorem}[Frechet calculus for the inverse metric square root]
\label{thm:V24-inverse-square-root}
The map \(z\mapsto M(z)^{-1/2}\) is \(C^2\), and with
\(R_t=(M(z)+tI)^{-1}\),
\begin{align}
 M^{-1/2}
 &=
 \frac1\pi\int_0^\infty t^{-1/2}R_t\,\dd t,
 \label{eq:V24-inverse-root-integral}\\
 D(M^{-1/2})[h]
 &=
 -\frac1\pi\int_0^\infty
 t^{-1/2}R_tM_hR_t\,\dd t,
 \label{eq:V24-inverse-root-first}\\
 D^2(M^{-1/2})[h,k]
 &=
 \frac1\pi\int_0^\infty t^{-1/2}
 \Bigl(
 R_tM_kR_tM_hR_t
 +R_tM_hR_tM_kR_t
 -R_tM_{hk}R_t
 \Bigr)\,\dd t.
 \label{eq:V24-inverse-root-second}
\end{align}
In particular,
\begin{align}
 \|D(M^{-1/2})[h]\|
 &\le
 \frac12m_*^{-3/2}\|M_h\|,
 \label{eq:V24-inverse-root-first-bound}\\
 \|D^2(M^{-1/2})[h,k]\|
 &\le
 \frac34m_*^{-5/2}\|M_h\|\|M_k\|
 +\frac12m_*^{-3/2}\|M_{hk}\|.
 \label{eq:V24-inverse-root-second-bound}
\end{align}
For a fixed base metric \(M_0\), the normalization
\begin{equation}
 R_M(z):=M(z)^{-1/2}M_0^{1/2}
 \label{eq:V24-metric-normalizer}
\end{equation}
satisfies
\begin{equation}
 R_M(z)^*M(z)R_M(z)=M_0,
 \label{eq:V24-metric-normalizer-identity}
\end{equation}
and its first two derivative bounds are
\eqref{eq:V24-inverse-root-first-bound}--%
\eqref{eq:V24-inverse-root-second-bound} multiplied by
\(\|M_0^{1/2}\|\).
\end{theorem}

\begin{proof}
For a positive scalar \(\lambda\),
\[
 \lambda^{-1/2}
 =
 \frac1\pi\int_0^\infty
 \frac{t^{-1/2}}{\lambda+t}\,\dd t.
\]
The spectral theorem gives \eqref{eq:V24-inverse-root-integral}.  The
resolvent bound
\(\|R_t\|\le(m_*+t)^{-1}\) supplies an integrable uniform majorant for
the resolvent and its first two parameter derivatives.  Differentiation
under the integral is therefore valid.  The inverse rule
\[
 DR_t[h]=-R_tM_hR_t
\]
gives \eqref{eq:V24-inverse-root-first}; differentiating it in direction
\(k\) gives \eqref{eq:V24-inverse-root-second}.

The beta integrals
\begin{align*}
 \frac1\pi\int_0^\infty
 \frac{t^{-1/2}}{(m_*+t)^2}\,\dd t
 &=\frac12m_*^{-3/2},\\
 \frac2\pi\int_0^\infty
 \frac{t^{-1/2}}{(m_*+t)^3}\,\dd t
 &=\frac34m_*^{-5/2}
\end{align*}
prove the two norm estimates.  Equation
\eqref{eq:V24-metric-normalizer-identity} follows by direct
multiplication, and \(M_0^{1/2}\) is parameter independent.
\end{proof}

If \(R_M^{-1}\) is also needed, no new functional calculus is required:
\begin{align}
 D(R_M^{-1})[h]
 &=-R_M^{-1}(DR_M[h])R_M^{-1},
 \label{eq:V24-normalizer-inverse-first}\\
 D^2(R_M^{-1})[h,k]
 &=
 R_M^{-1}(DR_M[k])R_M^{-1}(DR_M[h])R_M^{-1}
 \nonumber\\
 &\quad+
 R_M^{-1}(DR_M[h])R_M^{-1}(DR_M[k])R_M^{-1}
 -R_M^{-1}(D^2R_M[h,k])R_M^{-1}.
 \label{eq:V24-normalizer-inverse-second}
\end{align}
These identities give explicit bounds from a uniform inverse floor and the
preceding theorem.


% -----------------------------------------------------------------------------
\subsection{Two parameter derivatives of the allowed-trace projector}
\label{subsec:V24-projector-calculus}
% -----------------------------------------------------------------------------

There are two physical constructions of an allowed-trace projector.  A
flux cluster is obtained by a Riesz contour; a Robin or graph boundary
space is given by a full-column trace frame.  Both admit quantitative
field derivatives.

\begin{lemma}[Second-order Riesz projector calculus]
\label{lem:V24-Riesz-C2}
Let \(A:\mathcal B\to\mathcal B(\mathscr E)\) be \(C^2\).  Suppose one
fixed positively oriented contour \(\mathfrak C\) lies in the resolvent
set of every \(A(z)\) and separates the same finite spectral cluster.
Put
\[
 R_\lambda(z):=(\lambda I-A(z))^{-1},
 \qquad
 r_{\mathfrak C}:=
 \sup_{z\in\mathcal B,\,\lambda\in\mathfrak C}
 \|R_\lambda(z)\|.
\]
Then
\begin{equation}
 P(z):=
 \frac1{2\pi i}\int_{\mathfrak C}R_\lambda(z)\,\dd\lambda
 \label{eq:V24-Riesz-projector}
\end{equation}
is \(C^2\), with
\begin{align}
 DP(z)[h]
 &=
 \frac1{2\pi i}\int_{\mathfrak C}
 R_\lambda A_hR_\lambda\,\dd\lambda,
 \label{eq:V24-Riesz-first}\\
 D^2P(z)[h,k]
 &=
 \frac1{2\pi i}\int_{\mathfrak C}
 \Bigl(
 R_\lambda A_kR_\lambda A_hR_\lambda
 +R_\lambda A_hR_\lambda A_kR_\lambda
 +R_\lambda A_{hk}R_\lambda
 \Bigr)\,\dd\lambda.
 \label{eq:V24-Riesz-second}
\end{align}
If \(|\mathfrak C|\) is the contour length, then
\begin{align}
 \|DP[h]\|
 &\le
 \frac{|\mathfrak C|}{2\pi}
 r_{\mathfrak C}^2\|A_h\|,
 \label{eq:V24-Riesz-first-bound}\\
 \|D^2P[h,k]\|
 &\le
 \frac{|\mathfrak C|}{2\pi}
 \left(
 2r_{\mathfrak C}^3\|A_h\|\|A_k\|
 +r_{\mathfrak C}^2\|A_{hk}\|
 \right).
 \label{eq:V24-Riesz-second-bound}
\end{align}
\end{lemma}

\begin{proof}
The fixed-contour resolvent identity makes the resolvent \(C^2\) uniformly
on \(\mathfrak C\).  Since
\[
 D R_\lambda[h]=R_\lambda A_hR_\lambda,
\]
one differentiation of \eqref{eq:V24-Riesz-projector} gives
\eqref{eq:V24-Riesz-first}; a second differentiation gives the three
terms in \eqref{eq:V24-Riesz-second}.  Estimating the contour integrals
by length times the supremum proves the bounds.
\end{proof}

For the graph construction, let
\(T:\mathcal B\to\mathcal B(\mathscr E_-,\mathscr E_\partial)\) be \(C^2\)
and suppose
\begin{equation}
 T(z)^*T(z)\succeq s_*^2I_{\mathscr E_-}.
 \label{eq:V24-trace-column-floor}
\end{equation}
Put \(T^\dagger=(T^*T)^{-1}T^*\) and \(P_T=TT^\dagger\).

\begin{lemma}[Second-order trace-graph projector bound]
\label{lem:V24-trace-projector-C2}
The orthogonal range projector \(P_T\) is \(C^2\), and
\begin{align}
 DP_T[h]
 &=
 (I-P_T)T_hT^\dagger
 +(T^\dagger)^*T_h^*(I-P_T),
 \label{eq:V24-trace-projector-first}\\
 D(T^\dagger)[h]
 &=
 -T^\dagger T_hT^\dagger
 +(T^*T)^{-1}T_h^*(I-P_T).
 \label{eq:V24-pseudoinverse-first}
\end{align}
Consequently, if for unit directions
\[
 \|T_h\|\le t_1,\qquad
 \|T_{hk}\|\le t_2,
\]
then
\begin{align}
 \|DP_T[h]\|
 &\le\frac{2t_1}{s_*},
 \label{eq:V24-trace-projector-first-bound}\\
 \|D^2P_T[h,k]\|
 &\le
 \frac{2t_2}{s_*}
 +\frac{8t_1^2}{s_*^2}.
 \label{eq:V24-trace-projector-second-bound}
\end{align}
\end{lemma}

\begin{proof}
Differentiate \(P_T=T(T^*T)^{-1}T^*\) and use
\(T^\dagger T=I\) to collect the two off-range terms, giving
\eqref{eq:V24-trace-projector-first}.  Differentiating
\(T^\dagger=(T^*T)^{-1}T^*\) and collecting the same range projection
gives \eqref{eq:V24-pseudoinverse-first}.  Since
\(\|T^\dagger\|\le s_*^{-1}\) and
\(\|(T^*T)^{-1}\|\le s_*^{-2}\),
\[
 \|D(T^\dagger)[h]\|\le\frac{2t_1}{s_*^2}.
\]
Equation \eqref{eq:V24-trace-projector-first} immediately gives
\eqref{eq:V24-trace-projector-first-bound}.  Differentiate its first
summand in direction \(k\):
\[
 -DP_T[k]\,T_hT^\dagger
 +(I-P_T)T_{hk}T^\dagger
 +(I-P_T)T_hD(T^\dagger)[k].
\]
Its norm is at most
\(t_2/s_*+4t_1^2/s_*^2\).  The differentiated adjoint summand has the
same bound.  Their sum proves
\eqref{eq:V24-trace-projector-second-bound}.
\end{proof}

\begin{firewall}[A moving spectral cluster requires one uniform contour]
\label{fire:V24-fixed-contour}
Pointwise eigenvalue separation does not give
\eqref{eq:V24-Riesz-first-bound} unless one contour remains in the
resolvent set over the whole nonlinear ball.  A gap closing, rank change,
or contour crossing stops the \(C^2\) domain chart.  Its physical output is
the crossing parameter, the colliding spectral vectors, and the lost
resolvent margin.
\end{firewall}


% -----------------------------------------------------------------------------
\subsection{A fixed-base full unitary and its quantitative derivatives}
\label{subsec:V24-direct-rotation}
% -----------------------------------------------------------------------------

Let \(P_0=P(z_{\rm ref})\) and suppose the orthogonal projectors above
satisfy
\begin{equation}
 \delta:=
 \sup_{z\in\mathcal B}\|P(z)-P_0\|
 \le\frac12.
 \label{eq:V24-projector-ball}
\end{equation}
Write
\begin{equation}
 E(z):=P(z)-P_0,\qquad
 H(z):=I-E(z)^2,
 \label{eq:V24-direct-H}
\end{equation}
and
\begin{equation}
 S(z):=P(z)P_0+(I-P(z))(I-P_0),
 \qquad
 U_P(z):=S(z)H(z)^{-1/2}.
 \label{eq:V24-direct-unitary}
\end{equation}

\begin{theorem}[Fixed-base direct rotation with explicit \(C^2\) bounds]
\label{thm:V24-direct-rotation}
For every \(z\in\mathcal B\), \(U_P(z)\) is unitary and
\begin{equation}
 U_P(z)P_0=P(z)U_P(z).
 \label{eq:V24-direct-intertwining}
\end{equation}
Its restriction to \(\operatorname{Ran}P_0\) is exactly the V1 polar chord
\begin{equation}
 U_P(z)|_{\operatorname{Ran}P_0}
 =
 P(z)P_0\bigl(P_0P(z)P_0\bigr)^{-1/2},
 \label{eq:V24-direct-old-chord}
\end{equation}
and the complementary restriction is the analogous chord from
\(\ker P_0\) to \(\ker P(z)\).

Suppose for unit parameter directions
\[
 \|DP(z)[h]\|\le p_1,\qquad
 \|D^2P(z)[h,k]\|\le p_2.
\]
Put
\begin{equation}
 a:=1-\delta^2.
 \label{eq:V24-direct-floor}
\end{equation}
Then
\begin{align}
 \|DU_P(z)[h]\|
 &\le
 u_1:=
 p_1\left(a^{-1/2}+\delta a^{-3/2}\right),
 \label{eq:V24-direct-first-bound}\\
 \|D^2U_P(z)[h,k]\|
 &\le
 u_2:=
 p_2a^{-1/2}
 +\bigl((1+2\delta)p_1^2+\delta p_2\bigr)a^{-3/2}
 +3\delta^2p_1^2a^{-5/2}.
 \label{eq:V24-direct-second-bound}
\end{align}
Moreover,
\begin{equation}
 \|U_P(z)-I\|
 \le
 q_\delta:=
 \frac{1+\delta}{\sqrt{1-\delta^2}}-1<1.
 \label{eq:V24-direct-distance}
\end{equation}
\end{theorem}

\begin{proof}
The projection algebra gives
\begin{equation}
 S^*S=I-(P-P_0)^2=H,
 \qquad
 SP_0=PS,
 \label{eq:V24-direct-algebra}
\end{equation}
and \(H\) commutes with both \(P_0\) and \(P\).  Since
\(H\succeq aI\), the polar factor \(U_P=SH^{-1/2}\) is defined.
Also
\[
 S-I=(P-P_0)(2P_0-I),
\]
so \(\|S-I\|\le\delta<1\); hence \(S\) is invertible.  Therefore its
isometric polar factor is onto and is unitary.  Commutation of
\(H^{-1/2}\) with the two projectors and the second identity in
\eqref{eq:V24-direct-algebra} prove
\eqref{eq:V24-direct-intertwining}.  Restricted to
\(\operatorname{Ran}P_0\), \(S=P P_0\) and
\(H=P_0PP_0\), which proves \eqref{eq:V24-direct-old-chord}.  The
complementary formula is identical with \(P\) replaced by \(I-P\).

For the derivative bounds,
\[
 DS[h]=DP[h](2P_0-I),\qquad
 D^2S[h,k]=D^2P[h,k](2P_0-I),
\]
while
\begin{align*}
 DH[h]&=-DP[h]E-E\,DP[h],\\
 D^2H[h,k]
 &=-DP[h]DP[k]-DP[k]DP[h]\\
 &\quad-D^2P[h,k]E-E\,D^2P[h,k].
\end{align*}
Consequently
\[
 \|DH[h]\|\le2\delta p_1,\qquad
 \|D^2H[h,k]\|\le2p_1^2+2\delta p_2.
\]
Apply Theorem~\ref{thm:V24-inverse-square-root} to \(H\):
\begin{align*}
 \|D(H^{-1/2})[h]\|
 &\le\delta p_1a^{-3/2},\\
 \|D^2(H^{-1/2})[h,k]\|
 &\le
 3\delta^2p_1^2a^{-5/2}
 +(p_1^2+\delta p_2)a^{-3/2}.
\end{align*}
Differentiate \(U_P=SH^{-1/2}\).  Using
\(\|S\|\le1\) and \(\|H^{-1/2}\|\le a^{-1/2}\) gives
\eqref{eq:V24-direct-first-bound}.  The four terms
\[
 D^2S\,H^{-1/2}
 +DS[h]D(H^{-1/2})[k]
 +DS[k]D(H^{-1/2})[h]
 +S D^2(H^{-1/2})[h,k]
\]
give \eqref{eq:V24-direct-second-bound}.

Finally, the spectrum of \(H\) lies in \([a,1]\), so
\[
 \|H^{-1/2}-I\|\le a^{-1/2}-1.
\]
The identity
\(U_P-I=(S-I)H^{-1/2}+(H^{-1/2}-I)\) gives
\[
 \|U_P-I\|
 \le\delta a^{-1/2}+a^{-1/2}-1=q_\delta.
\]
At \(\delta=1/2\), the right side is \(\sqrt3-1<1\), and it is increasing
on the stated interval.
\end{proof}

Because \(q_\delta<1\), define the norm-convergent principal logarithm
\begin{equation}
 B_P(z):=\log U_P(z)
 =
 \sum_{j\ge1}\frac{(-1)^{j+1}}j(U_P(z)-I)^j.
 \label{eq:V24-unitary-log}
\end{equation}

\begin{corollary}[Skew logarithm and its first two parameter bounds]
\label{cor:V24-log-bounds}
The operator \(B_P(z)\) is skew-adjoint and
\begin{align}
 \|DB_P(z)[h]\|
 &\le
 b_1:=\frac{u_1}{1-q_\delta},
 \label{eq:V24-log-first-bound}\\
 \|D^2B_P(z)[h,k]\|
 &\le
 b_2:=
 \frac{u_2}{1-q_\delta}
 +\frac{u_1^2}{(1-q_\delta)^2}.
 \label{eq:V24-log-second-bound}
\end{align}
\end{corollary}

\begin{proof}
The spectrum of \(U_P\) lies in the disk
\(\{|\lambda-1|\le q_\delta\}\), which avoids the nonpositive real axis.
Functional calculus for the principal logarithm of a unitary gives
\(B_P^*=-B_P\).  Termwise differentiation of
\eqref{eq:V24-unitary-log} is justified uniformly by \(q_\delta<1\).
For the first derivative, the \(j\)-th power contributes at most
\(j q_\delta^{j-1}u_1\); division by \(j\) and summation gives
\eqref{eq:V24-log-first-bound}.  The part of the second derivative
containing \(D^2U_P\) has the same geometric factor.  The two first
derivatives have \(j(j-1)\) ordered insertion positions, so after
division by \(j\) their sum is bounded by
\[
 u_1^2\sum_{j\ge2}(j-1)q_\delta^{j-2}
 =\frac{u_1^2}{(1-q_\delta)^2}.
\]
This proves \eqref{eq:V24-log-second-bound}.
\end{proof}


% -----------------------------------------------------------------------------
\subsection{A unitary collar chart over the nonlinear ball}
\label{subsec:V24-collar-chart}
% -----------------------------------------------------------------------------

Choose a fixed collar coordinate \(0\le\rho<\rho_0\) and a smooth cutoff
\(\zeta\) with \(\zeta(0)=1\) and \(\zeta=0\) near \(\rho_0\).  In the
standard boundary metric define
\begin{equation}
 \widetilde{\mathcal W}_P(z;\rho,x)
 :=
 \exp\!\bigl(\zeta(\rho)B_P(z,x)\bigr),
 \label{eq:V24-collar-chart}
\end{equation}
and extend it by the identity outside the collar.

\begin{theorem}[Field-parameter unitary collar chart]
\label{thm:V24-collar-chart}
The multiplier \(\widetilde{\mathcal W}_P(z)\) is pointwise unitary,
equals \(U_P(z)\) on the boundary, and equals the identity off the collar.
It maps the fixed allowed trace space to the current one:
\begin{equation}
 P(z)\Gamma\widetilde{\mathcal W}_P(z)u
 =
 \Gamma\widetilde{\mathcal W}_P(z)u
 \quad\Longleftrightarrow\quad
 P_0\Gamma u=\Gamma u.
 \label{eq:V24-collar-domain-map}
\end{equation}
For unit parameter directions, its pointwise operator derivatives obey
\begin{align}
 \|D\widetilde{\mathcal W}_P(z)[h]\|
 &\le b_1,
 \label{eq:V24-collar-first-bound}\\
 \|D^2\widetilde{\mathcal W}_P(z)[h,k]\|
 &\le b_2+b_1^2.
 \label{eq:V24-collar-second-bound}
\end{align}
\end{theorem}

\begin{proof}
Skew-adjointness of \(B_P\) makes the exponential in
\eqref{eq:V24-collar-chart} unitary.  At \(\rho=0\) it is
\(\exp B_P=U_P\); where \(\zeta=0\) it is the identity.  The trace of a
smooth multiplier is its boundary value, so
\eqref{eq:V24-direct-intertwining} proves
\eqref{eq:V24-collar-domain-map}.

For a skew-adjoint \(X\), the derivative formula
\[
 D\exp_X[H]
 =
 \int_0^1e^{(1-s)X}He^{sX}\,\dd s
\]
has norm at most \(\|H\|\).  Here \(X=\zeta B_P\), and
\(|\zeta|\le1\), so \eqref{eq:V24-log-first-bound} gives
\eqref{eq:V24-collar-first-bound}.  Differentiating once more produces
one term containing \(\zeta D^2B_P\) and the two ordered simplex terms
containing \(\zeta DB_P[h]\) and \(\zeta DB_P[k]\).  The total volume of
the two ordered simplices is one.  Unitarity of all exponential factors
therefore gives
\[
 \|D^2\widetilde{\mathcal W}_P[h,k]\|
 \le\|D^2B_P[h,k]\|+\|DB_P[h]\|\|DB_P[k]\|,
\]
which is \eqref{eq:V24-collar-second-bound}.
\end{proof}

\begin{remark}[The fixed-base chart does not flatten support curvature]
\label{rem:V24-curvature-retained}
On \(\operatorname{Ran}P_0\), the boundary value of
\(\widetilde{\mathcal W}_P\) is the V1 polar chord.  In this local frame,
the Kato curvature is the conjugate
\[
 P_0U_P^*\,P[DP[h],DP[k]]P\,U_PP_0.
\]
It need not vanish.  A mixed command square must still retain the V1
support-holonomy debit.  The direct chart gives a common domain for
differentiation; it does not make covariant mixed derivatives commute.
\end{remark}


% -----------------------------------------------------------------------------
\subsection{Metric and boundary data give one \(C^2\) domain chart}
\label{subsec:V24-total-chart}
% -----------------------------------------------------------------------------

Return to the physical boundary fibre.  Let \(M_0=M(z_{\rm ref})\) and
\(S_0=M_0^{1/2}\).  Suppose \(T_{\rm phys}(z)\) is a full-column frame for
the allowed physical trace.  After metric normalization its standard-frame
columns are
\begin{equation}
 \overline T(z)
 :=
 S_0\,R_M(z)|_{\partial}^{-1}T_{\rm phys}(z).
 \label{eq:V24-pulled-trace-frame}
\end{equation}
Assume the column floor
\begin{equation}
 \overline T(z)^*\overline T(z)\succeq s_*^2I
 \label{eq:V24-pulled-column-floor}
\end{equation}
and form the standard orthogonal projector
\begin{equation}
 \overline P(z)
 :=
 \overline T(z)
 \bigl(\overline T(z)^*\overline T(z)\bigr)^{-1}
 \overline T(z)^*.
 \label{eq:V24-pulled-projector}
\end{equation}
A spectral incoming projector may be used instead, provided the fixed
contour hypotheses of Lemma~\ref{lem:V24-Riesz-C2} hold.

Apply Theorem~\ref{thm:V24-direct-rotation} and
Theorem~\ref{thm:V24-collar-chart} to \(\overline P\).  Return to the
fixed \(M_0\) metric by
\begin{equation}
 \mathcal W(z):=
 S_0^{-1}\widetilde{\mathcal W}_{\overline P}(z)S_0,
 \qquad
 \mathcal U(z):=R_M(z)\mathcal W(z).
 \label{eq:V24-total-parameter-chart}
\end{equation}

\begin{theorem}[Field-parameter metric-unitary domain chart]
\label{thm:V24-total-domain-chart}
Suppose \(M\) and \(T_{\rm phys}\) are \(C^2\), the metric and column
floors \eqref{eq:V24-metric-floor} and
\eqref{eq:V24-pulled-column-floor} hold, and
\[
 \sup_{z\in\mathcal B}
 \|\overline P(z)-\overline P(z_{\rm ref})\|\le\frac12.
\]
Then \(\mathcal U(z)\) is a \(C^2\) family of boundedly invertible
multipliers satisfying
\begin{align}
 \mathcal U(z)^*M(z)\mathcal U(z)&=M_0,
 \label{eq:V24-total-metric-identity}\\
 \mathcal U(z)\mathcal D_{\rm ref}&=\mathcal D_z.
 \label{eq:V24-total-domain-identity}
\end{align}
Here \(\mathcal D_z\) is the physical graph domain with trace in
\(\operatorname{Ran}T_{\rm phys}(z)\).

The parameter derivatives are exactly
\begin{align}
 D\mathcal U[h]
 &=
 DR_M[h]\,\mathcal W+R_M\,D\mathcal W[h],
 \label{eq:V24-total-chart-first}\\
 D^2\mathcal U[h,k]
 &=
 D^2R_M[h,k]\,\mathcal W
 +DR_M[h]\,D\mathcal W[k]
 +DR_M[k]\,D\mathcal W[h]
 +R_M\,D^2\mathcal W[h,k].
 \label{eq:V24-total-chart-second}
\end{align}
Consequently every first and second chart norm is bounded explicitly by
Theorem~\ref{thm:V24-inverse-square-root},
Lemma~\ref{lem:V24-Riesz-C2} or
Lemma~\ref{lem:V24-trace-projector-C2}, and
\eqref{eq:V24-direct-first-bound}--%
\eqref{eq:V24-collar-second-bound}.
\end{theorem}

\begin{proof}
Equation \eqref{eq:V24-pulled-trace-frame} says precisely that a
fixed-metric trace \(v\) is allowed after multiplication by \(R_M\) if
and only if \(S_0v\in\operatorname{Ran}\overline P(z)\).
The collar identity \eqref{eq:V24-collar-domain-map}, conjugated by
\(S_0\), therefore makes \(\mathcal W(z)\) carry the reference pulled
trace space onto the current pulled trace space.  Multiplication by
\(R_M(z)\) returns it to the physical trace space.  This proves
\eqref{eq:V24-total-domain-identity} and also bijectivity.

Since \(\widetilde{\mathcal W}_{\overline P}\) is unitary,
\(\mathcal W^*M_0\mathcal W=M_0\).  Combine this with
\eqref{eq:V24-metric-normalizer-identity}:
\[
 \mathcal U^*M\mathcal U
 =
 \mathcal W^*R_M^*MR_M\mathcal W
 =
 \mathcal W^*M_0\mathcal W=M_0.
\]
All constituent maps are \(C^2\) by the cited results and the inverse
rules \eqref{eq:V24-normalizer-inverse-first}--%
\eqref{eq:V24-normalizer-inverse-second}.  The two product rules give
\eqref{eq:V24-total-chart-first}--%
\eqref{eq:V24-total-chart-second}; their triangle inequalities give the
announced explicit bounds.
\end{proof}

The natural source--initial--boundary target trivialization is a block
diagonal product of \(\mathcal U(z)^{-1}\), its initial-time restriction,
and its boundary trace map.  The inverse differentiation rules show that
this target map is also \(C^2\).  It may therefore be used as \(V_n(z)\)
in \eqref{eq:V23-pulled-residual}; its derivative is no longer an
unrecorded object.


% -----------------------------------------------------------------------------
\subsection{Sobolev realization and the exact derivative-loss separator}
\label{subsec:V24-Sobolev}
% -----------------------------------------------------------------------------

The preceding operator formulae are pointwise until spatial multiplier
regularity is specified.  Let \(X\) be a compact \(D\)-dimensional
cylinder or compact boundary chart.  Fix an integer \(q>D/2\).  Matrix
valued \(H^q(X)\) is a Banach algebra, and multiplication satisfies
\begin{equation}
 \|ab\|_{H^q}\le C_q\|a\|_{H^q}\|b\|_{H^q}.
 \label{eq:V24-Sobolev-algebra}
\end{equation}
Choose \(q\) at least as large as the largest Sobolev exponent on which
the domain multiplier must act.

Let \(j^r z=(z,\partial z,\ldots,\partial^r z)\), and suppose a coefficient
field has the local form
\begin{equation}
 C(z)(x)=c\bigl(x,j^r z(x)\bigr),
 \label{eq:V24-jet-coefficient}
\end{equation}
where \(c\) is smooth on a compact range containing the declared ball.

\begin{proposition}[Tame \(C^2\) jet substitution at the correct levels]
\label{prop:V24-tame-jet}
If \(z\) ranges in a bounded subset of \(H^{q+r}(X)\), then
\[
 C:H^{q+r}(X)\longrightarrow H^q(X)
\]
is \(C^2\).  There is a ball-dependent constant \(C_{q,r,R}\) such that
\begin{align}
 \|DC(z)[h]\|_{H^q}
 &\le C_{q,r,R}\|h\|_{H^{q+r}},
 \label{eq:V24-tame-first}\\
 \|D^2C(z)[h,k]\|_{H^q}
 &\le C_{q,r,R}
 \|h\|_{H^{q+r}}\|k\|_{H^{q+r}}.
 \label{eq:V24-tame-second}
\end{align}
If the pointwise metric, spectral, column, and chord floors above are
uniform, the maps
\[
 M^{-1/2},\quad P,\quad U_P,\quad B_P,\quad
 \mathcal W,\quad\mathcal U
\]
are also \(C^2\) from the same \(H^{q+r}\) ball into \(H^q\), with tame
bounds obtained by substituting
\eqref{eq:V24-tame-first}--%
\eqref{eq:V24-tame-second} into the explicit V24 formulae.
Their multiplication operators and first two parameter derivatives are
bounded on \(H^q\).
\end{proposition}

\begin{proof}
The jet map is bounded
\(H^{q+r}\to H^q\).  The Sobolev composition theorem on the bounded
range of \(j^rz\), together with
\eqref{eq:V24-Sobolev-algebra}, gives
\[
 DC(z)[h]
 =
 D_{\rm jet}c(x,j^rz)\,j^rh
\]
and
\[
 D^2C(z)[h,k]
 =
 D_{\rm jet}^2c(x,j^rz)[j^rh,j^rk],
\]
with the two displayed bounds.  Inversion is smooth in the open group of
invertible elements of the unital matrix \(H^q\) algebra.  The resolvent
integrals, fixed contour, inverse square root, logarithm series, and
exponential series converge uniformly on the stated positive margins.
Differentiating them in the Banach algebra reproduces the pointwise
formulae with the algebra constant absorbed into
\(C_{q,r,R}\).  Finally, \eqref{eq:V24-Sobolev-algebra} makes every
\(H^q\) coefficient a bounded multiplier on \(H^q\), including its first
two parameter derivatives.
\end{proof}

\begin{theorem}[Fourier separator for an uncharged derivative loss]
\label{thm:V24-derivative-loss}
For every \(r\ge1\), the map
\[
 \Phi_r:z\longmapsto\partial_1^rz
\]
is not a bounded map \(H^s(\mathbb T^D)\to H^s(\mathbb T^D)\) for any
finite \(s\).  In particular, a same-\(H^s\) \(C^1\) estimate for a
field-dependent domain chart cannot be inferred when its coefficient
linearization contains an uncancelled \(r\)-th derivative of the field.
\end{theorem}

\begin{proof}
Let
\[
 h_N(x):=(1+N^2)^{-s/2}e^{iNx_1}.
\]
Then \(\|h_N\|_{H^s}=1\), up to the fixed Fourier normalization, while
\[
 \|\partial_1^rh_N\|_{H^s}=N^r.
\]
Thus the operator norm of \(\partial_1^r:H^s\to H^s\) is infinite.
The derivative of \(\Phi_r\) is this same linear operator.  Any proposed
same-space chart estimate whose linearization contains this term without
a proved cancellation would imply the false boundedness just disproved.
\end{proof}

\begin{assessment}[The two valid branches after a loss is detected]
\label{ass:V24-loss-branches}
Proposition~\ref{prop:V24-tame-jet} gives a Banach branch by placing the
configuration in \(H^{q+r}\) and the coefficient/domain chart in \(H^q\).
A first-order enlargement is another possible Banach branch: introduce
the required jets as independent variables and add their differential and
integrability constraints to the V23 subsidiary sector.  That branch is
valid only after those new constraint graph inverses and mixed tails pass.
If neither construction preserves the nonlinear inverse estimates, the
appropriate remaining route is a tame scale with smoothing, such as a
Nash--Moser theorem.  V24 does not claim that the physical regulator
satisfies the hypotheses of that stronger route.
\end{assessment}


% -----------------------------------------------------------------------------
\subsection{Quantitative landing in the V23 transport-connection blocks}
\label{subsec:V24-to-V23}
% -----------------------------------------------------------------------------

The chart estimates enter the V23 nonlinear matrix through an elementary
but indispensable product rule.  In any fixed local raw target frame,
write
\[
 G_n(z):=\mathfrak F_n(\chi_n(z)),
 \qquad
 \mathcal F_n(z)=V_n(z)G_n(z).
\]

\begin{lemma}[Transport contribution to the complete second derivative]
\label{lem:V24-residual-second-derivative}
For directions \(h,k\),
\begin{align}
 D^2\mathcal F_n(z)[h,k]
 &=
 D^2V_n(z)[h,k]G_n(z)
 +DV_n(z)[h]DG_n(z)[k]
 \nonumber\\
 &\quad+
 DV_n(z)[k]DG_n(z)[h]
 +V_n(z)D^2G_n(z)[h,k].
 \label{eq:V24-residual-second-derivative}
\end{align}
If on the declared ball
\[
 \|V_n\|\le v_0,\quad
 \|DV_n\|\le v_1,\quad
 \|D^2V_n\|\le v_2,
 \qquad
 \|G_n\|\le g_0,\quad
 \|DG_n\|\le g_1,\quad
 \|D^2G_n\|\le g_2,
\]
in compatible multilinear norms, then
\begin{equation}
 \boxed{
 \|D^2\mathcal F_n\|
 \le
 v_2g_0+2v_1g_1+v_0g_2.}
 \label{eq:V24-complete-Hessian-bound}
\end{equation}
Sectorwise versions of the four terms give the constants
\(h_{n,\alpha\beta\delta}\) in
\eqref{eq:V23-third-block-bound}.
\end{lemma}

\begin{proof}
Differentiate
\(D\mathcal F_n[h]=DV_n[h]G_n+V_nDG_n[h]\) in direction \(k\).
This gives the four terms in
\eqref{eq:V24-residual-second-derivative}.  Their norms and the triangle
inequality prove \eqref{eq:V24-complete-Hessian-bound}.
\end{proof}

\begin{corollary}[V24 domain chart populates the analytic V23 connection row]
\label{cor:V24-populates-V23}
Suppose the physical coefficient maps satisfy the no-loss Sobolev branch
of Proposition~\ref{prop:V24-tame-jet}, all V24 floors are uniform, and
the raw common-action residual \(G_n\) has the displayed \(C^2\) bounds.
Then the pulled residual \(\mathcal F_n\) is \(C^2\) on the V23 affine
ball.  Its complete Hessian, including the metric, trace-projector,
collar, and target-frame derivatives, is bounded by
\eqref{eq:V24-complete-Hessian-bound}.  Hence
Lemma~\ref{lem:V23-Hessian-majorant} yields a finite
\(\nu_{n,\alpha\beta}(R)\) without omitting the V23 connection term.
\end{corollary}

\begin{proof}
Theorem~\ref{thm:V24-total-domain-chart} and
Proposition~\ref{prop:V24-tame-jet} give the \(C^2\) source and target
trivializations.  Apply Lemma~\ref{lem:V24-residual-second-derivative},
then Lemma~\ref{lem:V23-Hessian-majorant}.
\end{proof}


\begin{proposition}[Adjacent-cutoff derivative debit with moving frames]
\label{prop:V24-adjacent-pulled-derivative}
Let
\(\mathcal F_n=V_nG_n\) in fixed source coordinates.  With
\(\Delta X_n=X_{n+1}-X_n\),
\begin{align}
 \Delta(D\mathcal F)_n
 &=
 \Delta(DV)_n\,G_{n+1}
 +DV_n\,\Delta G_n
 \nonumber\\
 &\quad+
 \Delta V_n\,DG_{n+1}
 +V_n\,\Delta(DG)_n.
 \label{eq:V24-adjacent-derivative-identity}
\end{align}
Suppose
\[
 \|V_n\|\le v_0,\quad\|DV_n\|\le v_1,\quad
 \|G_n\|\le g_0,\quad\|DG_n\|\le g_1
\]
and define the adjacent norms
\[
 \eta_n^V=\|\Delta V_n\|,\quad
 \eta_n^{DV}=\|\Delta(DV)_n\|,\quad
 \eta_n^G=\|\Delta G_n\|,\quad
 \eta_n^{DG}=\|\Delta(DG)_n\|.
\]
Then
\begin{equation}
 \boxed{
 \|\Delta(D\mathcal F)_n\|
 \le
 g_0\eta_n^{DV}
 +v_1\eta_n^G
 +g_1\eta_n^V
 +v_0\eta_n^{DG}.}
 \label{eq:V24-adjacent-derivative-bound}
\end{equation}
After left multiplication by the fixed diagonal inverse
\(\mathscr L^{-1}\), sectorwise versions of this estimate populate the
V23 debits \(d^D_{n,\alpha\beta}\).
\end{proposition}

\begin{proof}
The first derivative is
\(D\mathcal F_n=DV_n\,G_n+V_nDG_n\).  For the first product use
\[
 DV_{n+1}G_{n+1}-DV_nG_n
 =
 \Delta(DV)_nG_{n+1}+DV_n\Delta G_n,
\]
and use the analogous identity for the second product.  This proves
\eqref{eq:V24-adjacent-derivative-identity}.  The triangle inequality
gives \eqref{eq:V24-adjacent-derivative-bound}.
\end{proof}

\begin{corollary}[Analytic discharge of the V23 bundle hypothesis]
\label{cor:V24-discharges-bundle}
On the no-loss Sobolev branch, suppose the V24 metric, contour or column,
chord, and Sobolev margins are uniform in \(n\); the raw residuals have
uniform \(C^2\) bounds; and the four right-hand sequences in
\eqref{eq:V24-adjacent-derivative-bound} are summable sectorwise.
Then the field-dependent source and target domain charts required by V23
exist, the V23 nonlinear oscillation constants are finite, and the
adjacent pulled-derivative debits are summable.  Thus these three analytic
requirements of the V23 full-sector inverse and contraction theorems need
not be assumed separately.
\end{corollary}

\begin{proof}
Theorem~\ref{thm:V24-total-domain-chart} and
Proposition~\ref{prop:V24-tame-jet} give the uniform \(C^2\) charts.
Corollary~\ref{cor:V24-populates-V23} gives the nonlinear oscillation
bounds.  Sum \eqref{eq:V24-adjacent-derivative-bound} and use the fixed
bounded diagonal inverse.
\end{proof}


% -----------------------------------------------------------------------------
\subsection{Executable V24 parameter-domain card}
\label{subsec:V24-card}
% -----------------------------------------------------------------------------

A physical V24 card is evaluated before the V23 full-sector singular
matrix.

\begin{enumerate}[label=\textup{(\arabic*)},leftmargin=3em]
 \item Declare the configuration ball, Sobolev levels \(H^{q+r}\to H^q\),
       the largest jet order \(r\), the boundary collar, and the fixed
       reference configuration.  Record whether a first-order enlarged
       state has already replaced the jets.
 \item For every field sector retain \(M(z)\), \(DM(z)[h]\), and
       \(D^2M(z)[h,k]\), together with the pointwise metric floor.  Evaluate
       \eqref{eq:V24-inverse-root-first-bound}--%
       \eqref{eq:V24-inverse-root-second-bound}; do not differentiate
       eigenvectors.
 \item For a flux projector retain one fixed contour, its resolvent maximum,
       \(DA\), and \(D^2A\).  For a trace-frame projector retain
       \(T_{\rm phys}\), its first two derivatives, and the pulled column
       floor.  Evaluate the corresponding V24 projector bounds.
 \item Compute the maximum chord angle \(\delta\), the direct rotation
       \(U_P\), the exact intertwining residual, and the constants
       \(a,u_1,u_2,q_\delta,b_1,b_2\).  Independently retain every V1 mixed
       support-holonomy row.
 \item Build the logarithmic collar multiplier and test its boundary value,
       off-collar identity, metric unitarity, domain mapping, and first two
       parameter derivatives.  Store the spatial \(H^q\) multiplier norms,
       not merely pointwise matrix norms.
 \item Form \(R_M,\mathcal W,\mathcal U\) and verify
       \eqref{eq:V24-total-metric-identity} and
       \eqref{eq:V24-total-domain-identity} on held-out trace vectors.
       Retain the exact product-rule terms in
       \eqref{eq:V24-total-chart-first}--%
       \eqref{eq:V24-total-chart-second}.
 \item For the raw residual and its target frame store
       \(v_0,v_1,v_2,g_0,g_1,g_2\), evaluate
       \eqref{eq:V24-complete-Hessian-bound}, and compare it with held-out
       central action jets.
 \item Across adjacent cutoffs retain all four frame/residual increments in
       \eqref{eq:V24-adjacent-derivative-bound}.  Their sectorwise sums must
       be finite before a V23 derivative tail is declared.
 \item Run the Fourier or spectral high-frequency panel at the claimed
       same-space regularity.  If the normalized derivative response grows,
       move to the declared higher-domain, first-order enlarged, or tame-scale
       branch before forming \(\mathfrak m_{23}\).
\end{enumerate}

The compact analytic output is
\begin{equation}
 \mathfrak C_{24}
 :=
 \left(
 m_*,s_*,
 r_{\mathfrak C}^{-1},
 \frac12-\delta,
 1-q_\delta,
 u_1,u_2,b_1,b_2,
 C_{q,r,R},
 v_2g_0+2v_1g_1+v_0g_2
 \right),
 \label{eq:V24-card-output}
\end{equation}
with the unused contour or column entry omitted.  Every margin preceding
the upper-bound constants must be positive.

Canonical failures are
\begin{align}
 W^{M}_{z,h}
 &:=
 \left(z,h,M(z),DM(z)[h],\lambda_{\min}M(z)\right),
 \label{eq:V24-metric-witness}\\
 W^{\mathfrak C}_{z,\lambda}
 &:=
 \left(z,\lambda,R_\lambda(z),r_{\mathfrak C}\right),
 \label{eq:V24-contour-witness}\\
 W^{T}_{z,c}
 &:=
 \left(z,c,\overline T(z)c,\|\overline T(z)c\|\right),
 \label{eq:V24-column-witness}\\
 W^{P}_{z}
 &:=
 \left(z,P(z)-P_0,\delta,U_P(z)\right),
 \label{eq:V24-chord-witness}\\
 W^{\rm loss}_{N}
 &:=
 \left(h_N,\partial_1^rh_N,N^r\right),
 \label{eq:V24-loss-witness}\\
 W^{\rm frame}_{n}
 &:=
 \left(\Delta V_n,\Delta(DV)_n,
       \Delta G_n,\Delta(DG)_n\right).
 \label{eq:V24-frame-witness}
\end{align}
They separate a metric-floor loss, spectral crossing, rank loss, large
chord, Sobolev derivative loss, and nonsummable cofinal frame motion.


% -----------------------------------------------------------------------------
\subsection{V24 theorem ledger and next upstream row}
\label{subsec:V24-status}
% -----------------------------------------------------------------------------

\paragraph{New conclusions proved in V24.}
Theorem~\ref{thm:V24-inverse-square-root} gives exact resolvent formulae
and constants for two field derivatives of the metric normalization.
Lemma~\ref{lem:V24-Riesz-C2} and
Lemma~\ref{lem:V24-trace-projector-C2} do the same for the two physical
allowed-trace constructions.  Theorem~\ref{thm:V24-direct-rotation}
extends the V1 allowed-space chord to a full fixed-base unitary and gives
explicit \(C^2\) bounds.  Corollary~\ref{cor:V24-log-bounds} and
Theorem~\ref{thm:V24-collar-chart} turn it into a unitary collar chart.
Theorem~\ref{thm:V24-total-domain-chart} combines this chart with the
metric normalizer.

Proposition~\ref{prop:V24-tame-jet} proves the no-loss Sobolev realization
at the correct domain and coefficient levels.
Theorem~\ref{thm:V24-derivative-loss} supplies an exact Fourier obstruction
to an unjustified same-space derivative estimate.
Lemma~\ref{lem:V24-residual-second-derivative} and
Proposition~\ref{prop:V24-adjacent-pulled-derivative} propagate the chart
through the nonlinear Hessian and adjacent-cutoff derivative rows required
by V23.

\paragraph{Imported conclusions not repeated.}
V1's Kato curvature and polar-chord plaquette, the attached moving-wall
material derivative, V19 symbol extraction, V20 boundary flux, V21
adjoint range, V22 time-path transport/evolution, and V23 full-sector
inverse and contraction remain imported.  V24 neither flattens their
support curvature nor reconstructs their missing physical coefficient
tables.

\paragraph{Authoritative physical status.}
The attached files assume smooth parameter pullbacks in selected finite
or material calculations, but they do not provide the complete first and
second field derivatives of the assembled SM--Einstein symmetrizer,
incoming projector or trace frame, collar lift, and source--initial--
boundary target chart.  They also do not specify one Sobolev pair
\(H^{q+r}\to H^q\) for all seven sectors or populate the cofinal increments
in \eqref{eq:V24-adjacent-derivative-bound}.  Therefore the V24 analytic
card is not physically populated, V23's number
\(\mathfrak m_{23}\) remains unevaluated, and the full continuum is not
claimed.

\paragraph{Next upstream row.}
The derivative-loss separator points to a concrete construction rather
than an impasse.  The next version should build one first-order enlarged
Einstein--Standard-Model state: introduce the metric/coframe, connection,
Yang--Mills curvature, Higgs derivative, and fermion derivative variables
which occur in the coefficient jets; write their defining and curl
constraints; prove a graph inverse and propagation estimate for that
jet-constraint complex; and compute how its triangular augmentation changes
the V23 full-sector singular margin.  If the augmentation is not triangular
after gauge reduction, its new mixed blocks and high tails must enter the
full matrix.

\begin{firewall}[Smooth finite matrices do not establish a tame field chart]
\label{fire:V24-status}
A \(C^2\) matrix formula at each fixed spacetime point does not imply a
\(C^2\) multiplier on the graph space.  The positive floors, spatial
Sobolev norms, jet order, and cofinal frame increments in
\(\mathfrak C_{24}\) are required data.  Until they pass on the occurring
common-action history, V24 is an analytic bridge rather than a physical
continuum construction.
\end{firewall}

% END APPENDED VERSION V24 MATERIAL



% =============================================================================
\section{Version V25: first-order jet enlargement and compatible constraint complex}
\label{sec:V25-jet-enlargement}
% =============================================================================

\subsection{Purpose and nonduplication boundary}
\label{subsec:V25-purpose}

V18 assumes that a gauge-fixed first-order Einstein--Standard-Model
reduction and its subsidiary system have already been formed.  V20 then
tests the boundary flux of such a system.  Neither version constructs the
jet variables which remove the V24 derivative loss, identifies the range
of their redundant curl rows, or computes how their addition changes the
V23 inverse margin.  The attached manuscripts likewise contain
moving-material derivatives and finite constraint cards, but no common
Sobolev graph factorization of the complete enlarged field system.

V25 supplies that factorization.  The defining constraint is made an
actual equation, while curl and Bianchi rows are kept as compatible
outputs rather than freely prescribable data.  Exact Cartan identities
then show how the coframe, spin connection, Yang--Mills curvature, Higgs
derivative, fermion derivative, and ghost derivative rows fit one
constraint complex.  The construction is analytic and does not assert
that the missing physical common-action matrices pass its estimates.


% -----------------------------------------------------------------------------
\subsection{Nonlinear jet enlargement and exact root equivalence}
\label{subsec:V25-abstract-enlargement}
% -----------------------------------------------------------------------------

Let \(\mathscr X\) be the primary-field graph space,
\(\mathscr Z\) a lower-order jet space, and \(\mathscr Y\) the main
equation target.  The mixed Sobolev levels are part of these definitions;
for example,
\[
 \mathscr X=H^{q+1},\qquad \mathscr Z=H^q,
\]
makes one spatial derivative bounded from \(\mathscr X\) to
\(\mathscr Z\).  Let
\[
 \mathcal J:\mathscr U\subset\mathscr X\longrightarrow\mathscr Z,
 \qquad
 \mathcal E:\mathscr U\times\mathscr Z\longrightarrow\mathscr Y
\]
be \(C^2\).  Here \(\mathcal J(x)\) is the geometric first jet, including
the selected connection, and \(\mathcal E(x,p)\) is the prolonged main
equation with every coefficient depending algebraically on the
independent jet \(p\).  Define
\begin{align}
 \mathcal R(x)&:=\mathcal E(x,\mathcal J(x)),
 \label{eq:V25-reduced-residual}\\
 \mathbb F(x,p)
 &:=
 \binom{\mathcal E(x,p)}{p-\mathcal J(x)}.
 \label{eq:V25-augmented-residual}
\end{align}

\begin{theorem}[Exact nonlinear jet-root equivalence]
\label{thm:V25-root-equivalence}
The projection \((x,p)\mapsto x\) is a bijection between zeros of
\(\mathbb F\) and zeros of \(\mathcal R\), with inverse
\(x\mapsto(x,\mathcal J(x))\).  No solution is introduced or removed by
the defining-jet equation.
\end{theorem}

\begin{proof}
If \(\mathbb F(x,p)=0\), its second component gives
\(p=\mathcal J(x)\); its first component is then
\(\mathcal R(x)=0\).  Conversely, if \(\mathcal R(x)=0\), substitution
of \(p=\mathcal J(x)\) makes both components of
\eqref{eq:V25-augmented-residual} vanish.  The two operations are inverse.
\end{proof}

Fix a constraint-satisfying point \((x,\mathcal J(x))\), and set
\begin{equation}
 A:=D_x\mathcal E(x,\mathcal J(x)),\qquad
 B:=D_p\mathcal E(x,\mathcal J(x)),\qquad
 J:=D\mathcal J(x).
 \label{eq:V25-linear-blocks}
\end{equation}
Then
\begin{equation}
 H:=D\mathcal R(x)=A+BJ.
 \label{eq:V25-reduced-linearization}
\end{equation}

\begin{theorem}[Triangular jet factorization and inverse]
\label{thm:V25-triangular-factorization}
The augmented derivative and the jet-coordinate shear are
\begin{equation}
 \mathbb H
 =
 D\mathbb F(x,\mathcal J(x))
 =
 \begin{pmatrix}A&B\\-J&I\end{pmatrix},
 \qquad
 T_J:=
 \begin{pmatrix}I&0\\J&I\end{pmatrix}.
 \label{eq:V25-augmented-derivative}
\end{equation}
They satisfy the exact factorization
\begin{equation}
 \boxed{
 \mathbb H T_J
 =
 \begin{pmatrix}H&B\\0&I\end{pmatrix}.}
 \label{eq:V25-triangular-factorization}
\end{equation}
Consequently \(\mathbb H\) is a bounded isomorphism if and only if \(H\)
is.  For data \((f,c)\), its inverse is
\begin{align}
 h&=H^{-1}(f-Bc),
 \label{eq:V25-augmented-inverse-h}\\
 q&=Jh+c.
 \label{eq:V25-augmented-inverse-q}
\end{align}

In Hilbert product norms, if
\begin{equation}
 \|H^{-1}\|\le\kappa,\qquad
 \|B\|\le b,\qquad
 \|J\|\le j,
 \label{eq:V25-block-bounds}
\end{equation}
then
\begin{equation}
 \boxed{
 \|\mathbb H^{-1}\|
 \le
 \kappa_{\rm aug}
 :=
 (1+j)
 \left\|
 \begin{pmatrix}
  \kappa&\kappa b\\
  0&1
 \end{pmatrix}
 \right\|_{2\to2}.}
 \label{eq:V25-augmented-inverse-bound}
\end{equation}
Thus the augmented singular floor is at least
\begin{equation}
 \gamma_{\rm aug}:=\kappa_{\rm aug}^{-1}.
 \label{eq:V25-augmented-floor}
\end{equation}
\end{theorem}

\begin{proof}
The chain rule gives \eqref{eq:V25-reduced-linearization}.  Direct block
multiplication gives
\[
 \begin{pmatrix}A&B\\-J&I\end{pmatrix}
 \begin{pmatrix}I&0\\J&I\end{pmatrix}
 =
 \begin{pmatrix}A+BJ&B\\0&I\end{pmatrix},
\]
which is \eqref{eq:V25-triangular-factorization}.  The shear and its
inverse are bounded.  The upper triangular factor is invertible exactly
when \(H\) is, and forward substitution gives
\eqref{eq:V25-augmented-inverse-h}--%
\eqref{eq:V25-augmented-inverse-q}.

The inverse of the upper triangular factor is
\[
 \begin{pmatrix}
 H^{-1}&-H^{-1}B\\0&I
 \end{pmatrix}.
\]
Lemma~\ref{lem:V23-block-majorant} bounds its norm by the scalar matrix in
\eqref{eq:V25-augmented-inverse-bound}.  Since
\(T_J=I+\left(\begin{smallmatrix}0&0\\J&0\end{smallmatrix}\right)\),
\(\|T_J\|\le1+j\).  The inverse factorization
\[
 \mathbb H^{-1}
 =
 T_J
 \begin{pmatrix}
 H^{-1}&-H^{-1}B\\0&I
 \end{pmatrix}
\]
proves the bound and hence \eqref{eq:V25-augmented-floor}.
\end{proof}

\begin{corollary}[Finite read errors in the reduced jet block]
\label{cor:V25-reduced-read-error}
Let finite or transported approximations
\(\widehat A,\widehat B,\widehat J\) satisfy
\[
 \|A-\widehat A\|\le e_A,\quad
 \|B-\widehat B\|\le e_B,\quad
 \|J-\widehat J\|\le e_J,
 \quad
 \|J\|\le j,\quad\|\widehat B\|\le\widehat b.
\]
For
\(\widehat H=\widehat A+\widehat B\widehat J\),
\begin{equation}
 \boxed{
 \|H-\widehat H\|
 \le e_H:=e_A+e_Bj+\widehat b\,e_J.}
 \label{eq:V25-reduced-read-error}
\end{equation}
If \(s_{\min}(\widehat H)=\sigma_H>e_H\), then
\[
 \kappa\le(\sigma_H-e_H)^{-1}
\]
in \eqref{eq:V25-augmented-inverse-bound}.
\end{corollary}

\begin{proof}
Use
\[
 H-\widehat H
 =
 (A-\widehat A)+(B-\widehat B)J
 +\widehat B(J-\widehat J)
\]
and take norms.  The last assertion is the Neumann perturbation estimate.
\end{proof}

\begin{firewall}[A formal jet variable does not remove an unbounded defining row]
\label{fire:V25-mixed-levels}
The factorization is a bounded-operator statement only when
\(J:\mathscr X\to\mathscr Z\) is bounded.  Taking both spaces to be the
same \(H^q\) while \(J\) contains a derivative would reproduce the V24
Fourier obstruction.  The correct mixed graph levels, or an independently
proved tame-scale inverse, must be declared before
\eqref{eq:V25-augmented-floor} is evaluated.
\end{firewall}


% -----------------------------------------------------------------------------
\subsection{Curl rows have a compatible target, not free data}
\label{subsec:V25-compatible-range}
% -----------------------------------------------------------------------------

Let \(K:\mathscr Z\to\mathscr W\) and
\(R:\mathscr X\to\mathscr W\) be bounded operators satisfying the complex
identity
\begin{equation}
 KJ=R.
 \label{eq:V25-complex-identity}
\end{equation}
The linear defining and curl defects are
\[
 c:=q-Jh,\qquad d:=Kq-Rh.
\]
Then
\begin{equation}
 d=Kc.
 \label{eq:V25-curl-from-defining}
\end{equation}

\begin{theorem}[Range and inverse of the augmented constraint complex]
\label{thm:V25-compatible-range}
Define
\begin{equation}
 \mathbb H_{\rm cx}(h,q)
 :=
 \left(
 Ah+Bq,\ q-Jh,\ Kq-Rh
 \right).
 \label{eq:V25-complex-graph-map}
\end{equation}
If \(H=A+BJ\) is an isomorphism, then
\(\mathbb H_{\rm cx}\) is one-to-one and has range exactly
\begin{equation}
 \mathscr Y_{\rm cx}
 :=
 \{(f,c,d)\in\mathscr Y\oplus\mathscr Z\oplus\mathscr W:
       d=Kc\}.
 \label{eq:V25-compatible-target}
\end{equation}
For \((f,c,Kc)\) in this target, the inverse is still
\eqref{eq:V25-augmented-inverse-h}--%
\eqref{eq:V25-augmented-inverse-q}.  Thus the redundant curl row neither
creates a kernel nor permits an independently prescribed curl datum.
\end{theorem}

\begin{proof}
Equation \eqref{eq:V25-complex-identity} gives
\[
 Kq-Rh=K(q-Jh),
\]
so every image lies in \(\mathscr Y_{\rm cx}\).  Given
\((f,c,Kc)\), equations
\eqref{eq:V25-augmented-inverse-h}--%
\eqref{eq:V25-augmented-inverse-q} solve the first two components, and
the identity makes the third component \(Kc\).  This proves surjectivity
onto the compatible target.  Zero compatible data give \(h=q=0\), so the
map is one-to-one.
\end{proof}

More generally, if a nonlinear curl map is affine in its jet argument,
\[
 \mathcal K(x,p)=K(x)p-r(x),
 \qquad
 \mathcal K(x,\mathcal J(x))=0,
\]
then the exact identity
\begin{equation}
 \mathcal K(x,p)
 =
 K(x)\bigl(p-\mathcal J(x)\bigr)
 \label{eq:V25-nonlinear-compatible-row}
\end{equation}
holds.  At a constraint-satisfying point its derivative is therefore the
linear compatible row in Theorem~\ref{thm:V25-compatible-range}.
\begin{firewall}[Overdetermined curl data are not a range failure]
\label{fire:V25-overdetermined-data}
A triple with \(d\ne Kc\) is outside the physical target
\eqref{eq:V25-compatible-target}.  Failure to solve it is not a cokernel
of the prolonged field equations; it is incompatible prescribed data.
A genuine adjoint obstruction is tested only after the target is reduced
to this compatibility graph.
\end{firewall}


% -----------------------------------------------------------------------------
\subsection{Covariant jet identities and finite residual propagation}
\label{subsec:V25-covariant-jet}
% -----------------------------------------------------------------------------

The abstract complex has a direct geometric realization.  Fix the convention
\begin{equation}
 [D_\mu,D_\nu]u=\mathcal R_{\mu\nu}u
 \label{eq:V25-curvature-convention}
\end{equation}
in a coordinate frame.  Introduce independent covariant jet variables
\(p_\mu\) and define
\begin{align}
 C_\mu&:=p_\mu-D_\mu u,
 \label{eq:V25-defining-constraint}\\
 Z_{\mu\nu}
 &:=
 D_\mu p_\nu-D_\nu p_\mu-\mathcal R_{\mu\nu}u.
 \label{eq:V25-curl-constraint}
\end{align}

\begin{lemma}[Exact covariant curl identity]
\label{lem:V25-covariant-curl}
For every \(u,p\),
\begin{equation}
 \boxed{
 Z_{\mu\nu}
 =
 D_\mu C_\nu-D_\nu C_\mu.}
 \label{eq:V25-curl-identity}
\end{equation}
Thus the curl datum is determined by the defining-constraint datum.
\end{lemma}

\begin{proof}
Substitute \(p_\mu=D_\mu u+C_\mu\) in
\eqref{eq:V25-curl-constraint}.  The terms containing \(u\) cancel by
\eqref{eq:V25-curvature-convention}; the remaining terms are the right
side of \eqref{eq:V25-curl-identity}.
\end{proof}

Let the primary evolution and its differentiated prolongation have residuals
\begin{align}
 E_u
 &:=
 D_0u-\mathcal G(u,p_i),
 \label{eq:V25-primary-evolution-residual}\\
 E_i
 &:=
 D_0p_i-D_i\mathcal G(u,p_j)-\mathcal R_{0i}u.
 \label{eq:V25-prolonged-evolution-residual}
\end{align}

\begin{proposition}[Exact source of the defining-jet constraint]
\label{prop:V25-defining-propagation}
The spatial defining constraints satisfy
\begin{equation}
 \boxed{
 D_0C_i=E_i-D_iE_u.}
 \label{eq:V25-defining-propagation}
\end{equation}
Suppose \(D_0=\partial_t+B_0(t,x)\) on the constraint bundle and
\(\|B_0(t)\|_{H^q\to H^q}\le b_0(t)\).  Then
\begin{align}
 \|C(t)\|_{H^q}
 \le
 \exp\!\left(\int_0^tb_0(\tau)\,\dd\tau\right)
 \Bigl[
 &\|C(0)\|_{H^q}
 \nonumber\\
 &+\int_0^t
 \bigl(\|E_i(\tau)\|_{H^q}
       +\|D_iE_u(\tau)\|_{H^q}\bigr)\,\dd\tau
 \Bigr].
 \label{eq:V25-defining-energy}
\end{align}
In particular, exact prolonged equations and exact initial definitions
give \(C_i=0\), hence \(Z_{ij}=0\), for the whole interval.
\end{proposition}

\begin{proof}
Using \eqref{eq:V25-curvature-convention},
\[
 D_0D_i u=D_iD_0u+\mathcal R_{0i}u.
\]
Subtract this identity from
\(D_0p_i=E_i+D_i\mathcal G+\mathcal R_{0i}u\).  The result is
\[
 D_0(p_i-D_i u)
 =
 E_i+D_i(\mathcal G-D_0u)
 =
 E_i-D_iE_u,
\]
which proves \eqref{eq:V25-defining-propagation}.  Take the \(H^q\)
energy pairing, bound the zeroth-order connection term by \(b_0(t)\),
and apply Gronwall's inequality to obtain
\eqref{eq:V25-defining-energy}.  The exact conclusion follows from zero
right side, and Lemma~\ref{lem:V25-covariant-curl} then gives \(Z=0\).
\end{proof}

\begin{firewall}[The differentiated main residual has its own Sobolev cost]
\label{fire:V25-prolongation-source}
The finite debit in \eqref{eq:V25-defining-energy} contains
\(D_iE_u\).  Controlling only \(E_u\) in \(H^q\) does not control this
term in \(H^q\).  A strong prolongation card must measure \(E_u\) in
\(H^{q+1}\), or formulate \eqref{eq:V25-defining-propagation} and the
constraint energy at one lower or dual level.  This is the jet analogue
of the V23 differentiated Einstein-residual firewall.
\end{firewall}


% -----------------------------------------------------------------------------
\subsection{The Einstein--Standard-Model Cartan constraint complex}
\label{subsec:V25-SM-complex}
% -----------------------------------------------------------------------------

The following identities specify one common enlarged variable list.  They
are differential identities, independent of the field equations.  All
products use the represented bundle action and one fixed sign convention.

For a coframe \(e\), spin connection \(\omega\), independent torsion
\(\mathsf T\), and independent curvature \(\mathsf R\), put
\begin{align}
 C_{\mathsf R}
 &:=
 \mathsf R-(\dd\omega+\omega\wedge\omega),
 \label{eq:V25-curvature-definition-defect}\\
 C_{\mathsf T}
 &:=
 \mathsf T-D_\omega e,
 \label{eq:V25-torsion-definition-defect}\\
 B_{\mathsf R}
 &:=
 D_\omega\mathsf R,
 \label{eq:V25-curvature-Bianchi-defect}\\
 B_{\mathsf T}
 &:=
 D_\omega\mathsf T-\mathsf R\wedge e.
 \label{eq:V25-torsion-Bianchi-defect}
\end{align}
For a Yang--Mills connection \(A\) and independent curvature \(F\), put
\begin{align}
 C_F&:=F-(\dd A+A\wedge A),
 \label{eq:V25-YM-curvature-defect}\\
 B_F&:=D_AF.
 \label{eq:V25-YM-Bianchi-defect}
\end{align}
For a Higgs field \(\Phi\) and independent covariant derivative \(\Pi\),
put
\begin{align}
 C_\Pi&:=\Pi-D_A\Phi,
 \label{eq:V25-Higgs-jet-defect}\\
 B_\Pi&:=D_A\Pi-F\mathbin{\cdot}\Phi.
 \label{eq:V25-Higgs-integrability-defect}
\end{align}
Let \(D_{\rm tot}\) be the spin--gauge connection on the fermion bundle,
let \(\mathsf R_{\rm tot}\) be the curvature assembled from the independent
\(\mathsf R,F\), and let
\[
 C_{\mathsf R_{\rm tot}}
 :=
 \mathsf R_{\rm tot}-\operatorname{curv}(D_{\rm tot}).
\]
For an independent fermion derivative \(\Xi\), put
\begin{align}
 C_\Xi&:=\Xi-D_{\rm tot}\psi,
 \label{eq:V25-fermion-jet-defect}\\
 B_\Xi&:=D_{\rm tot}\Xi-\mathsf R_{\rm tot}\mathbin{\cdot}\psi.
 \label{eq:V25-fermion-integrability-defect}
\end{align}
Ghost and antighost jets use the same last pair in their respective
representations.

\begin{theorem}[Exact Cartan--SM compatibility identities]
\label{thm:V25-Cartan-identities}
The defects above satisfy
\begin{align}
 B_{\mathsf R}
 &=D_\omega C_{\mathsf R},
 \label{eq:V25-curvature-Bianchi-identity}\\
 B_{\mathsf T}
 &=D_\omega C_{\mathsf T}-C_{\mathsf R}\wedge e,
 \label{eq:V25-torsion-Bianchi-identity}\\
 B_F&=D_AC_F,
 \label{eq:V25-YM-Bianchi-identity}\\
 B_\Pi&=D_AC_\Pi-C_F\mathbin{\cdot}\Phi,
 \label{eq:V25-Higgs-integrability-identity}\\
 B_\Xi
 &=D_{\rm tot}C_\Xi
   -C_{\mathsf R_{\rm tot}}\mathbin{\cdot}\psi.
 \label{eq:V25-fermion-integrability-identity}
\end{align}
Hence every Bianchi and matter-integrability row vanishes when the
corresponding defining curvature and jet rows vanish.  Their target data
are constrained by these identities.
\end{theorem}

\begin{proof}
The connection curvatures obey
\[
 D_\omega(\dd\omega+\omega\wedge\omega)=0,
 \qquad
 D_A(\dd A+A\wedge A)=0.
\]
Substitution of
\(\mathsf R=\operatorname{curv}(\omega)+C_{\mathsf R}\) and
\(\mathsf T=D_\omega e+C_{\mathsf T}\), together with
\(D_\omega^2e=\operatorname{curv}(\omega)\wedge e\), gives the first two
identities.  The Yang--Mills Bianchi identity gives the third.
Next use
\[
 D_A^2\Phi=(\dd A+A\wedge A)\mathbin{\cdot}\Phi
\]
and substitute \(F=\operatorname{curv}(A)+C_F\) and
\(\Pi=D_A\Phi+C_\Pi\); this gives
\eqref{eq:V25-Higgs-integrability-identity}.  The identical calculation
with
\(D_{\rm tot}^2\psi=\operatorname{curv}(D_{\rm tot})\mathbin{\cdot}\psi\)
gives \eqref{eq:V25-fermion-integrability-identity}.
\end{proof}

For a coordinate first-order reduction of the reduced Einstein equation,
one may instead introduce
\[
 Q_{\alpha\mu\nu}:=\partial_\alpha g_{\mu\nu}
\]
as an independent variable.  Its defining and curl defects are
\begin{align}
 C^g_{\alpha\mu\nu}
 &:=
 Q_{\alpha\mu\nu}-\partial_\alpha g_{\mu\nu},
 \label{eq:V25-metric-jet-defect}\\
 Z^g_{\alpha\beta\mu\nu}
 &:=
 \partial_\alpha Q_{\beta\mu\nu}
 -\partial_\beta Q_{\alpha\mu\nu}
 =
 \partial_\alpha C^g_{\beta\mu\nu}
 -\partial_\beta C^g_{\alpha\mu\nu}.
 \label{eq:V25-metric-curl-identity}
\end{align}
A physical card must choose this metric reduction or the coframe--connection
reduction before forming the principal matrix.  Mixing them after the
fact changes the unknown and constraint spaces.

\begin{corollary}[No independent Bianchi row in the V23 inverse]
\label{cor:V25-no-independent-Bianchi}
On an exact enlarged root, all defining defects
\[
 C^g,\ C_{\mathsf R},\ C_{\mathsf T},\
 C_F,\ C_\Pi,\ C_\Xi
\]
vanish.  Theorem~\ref{thm:V25-Cartan-identities} then makes every displayed
curl, Bianchi, and integrability defect vanish.  These rows belong to the
compatible target of Theorem~\ref{thm:V25-compatible-range}; adding them
as independent target coordinates would create a false cokernel.
\end{corollary}

\begin{proof}
Apply the identities after setting their defining defects to zero.
The target statement is exactly
\eqref{eq:V25-compatible-target}, with \(K\) given by the appropriate
covariant exterior derivative and algebraic curvature action.
\end{proof}


% -----------------------------------------------------------------------------
\subsection{The augmented inverse and nonlinear V23 margin}
\label{subsec:V25-to-V23}
% -----------------------------------------------------------------------------

Apply the preceding construction to the direct sum of all selected
first-order variables.  The map \(\mathcal J\) contains the metric or
coframe jet, spin connection/curvature definitions, Yang--Mills curvature,
Higgs derivative, fermion derivative, and ghost derivatives.  Its
derivative \(J\) is generally not sector diagonal: varying the gauge or
spin connection changes the covariant jets of matter.  All such blocks
are included in the norm \(j=\|J\|\) and in the V23 full matrix.

Suppose the reduced full-sector preconditioner obtained in V23 has inverse
margin
\begin{equation}
 \|H_*^{-1}\|_{w\to w}\le\gamma_{23}^{-1}.
 \label{eq:V25-imported-V23-floor}
\end{equation}
At the same reference field let \(B_*\) be the derivative of the prolonged
main equations with respect to the independent jets and \(J_*\) the
defining-jet derivative.  Define
\begin{align}
 b_*&:=\|B_*\|,\qquad j_*:=\|J_*\|,
 \label{eq:V25-reference-jet-gains}\\
 \kappa_{25}
 &:=
 (1+j_*)
 \left\|
 \begin{pmatrix}
  \gamma_{23}^{-1}&\gamma_{23}^{-1}b_*\\
  0&1
 \end{pmatrix}
 \right\|_{2\to2},
 \qquad
 \gamma_{25}:=\kappa_{25}^{-1}.
 \label{eq:V25-reference-augmented-margin}
\end{align}

For the pulled cutoff residuals, let
\[
 \mathbb F_n:\mathbb U_R
 \subset\mathscr X\oplus\mathscr Z
 \longrightarrow\mathscr Y\oplus\mathscr Z
\]
include the defining constraints but not duplicate compatible curl
coordinates.  Let \(\mathbb H_*\) be the reference augmented derivative.
After all V24 transports, define
\begin{align}
 \Delta_{25}(R)
 &:=
 \sup_{n\ge n_0}\sup_{(x,p)\in\mathbb U_R}
 \|D\mathbb F_n(x,p)-\mathbb H_*\|,
 \label{eq:V25-augmented-derivative-radius}\\
 \Theta_{25}(R)
 &:=
 \kappa_{25}\Delta_{25}(R),
 \label{eq:V25-augmented-theta}\\
 \mathfrak b_{25}
 &:=
 \kappa_{25}
 \sup_{n\ge n_0}
 \|\mathbb F_n(x_{\rm ref},p_{\rm ref})\|,
 \label{eq:V25-augmented-center}\\
 \mathfrak m_{25}(R)
 &:=
 1-\Theta_{25}(R)-\frac{\mathfrak b_{25}}R.
 \label{eq:V25-augmented-master-margin}
\end{align}
Here \(p_{\rm ref}\in\mathscr Z\) is one fixed reference jet for every
cutoff, for example \(\mathcal J_{n_0}(x_{\rm ref})\).  The defining
component
\[
 p_{\rm ref}-\mathcal J_n(x_{\rm ref})
\]
is part of \(\mathfrak b_{25}\); it vanishes only when the transported
cutoff jet agrees with the fixed anchor.

A directly measurable upper bound for
\eqref{eq:V25-augmented-derivative-radius} is obtained as follows.  If
the transported blocks \(A,B,J\) differ from their reference blocks by
at most \(a_R,b_R,j_R\), respectively, then
\begin{equation}
 \boxed{
 \Delta_{25}(R)
 \le
 \left\|
 \begin{pmatrix}
  a_R&b_R\\
  j_R&0
 \end{pmatrix}
 \right\|_{2\to2}.}
 \label{eq:V25-augmented-radius-majorant}
\end{equation}
This is Lemma~\ref{lem:V23-block-majorant} applied to
\[
 D\mathbb F_n-\mathbb H_*
 =
 \begin{pmatrix}
  A_n-A_*&B_n-B_*\\
  -(J_n-J_*)&0
 \end{pmatrix}.
\]

Define the adjacent augmented action debit
\begin{equation}
 \mathfrak a_n^{25}
 :=
 \sup_{(x,p)\in\mathbb U_R}
 \|\mathbb F_{n+1}(x,p)-\mathbb F_n(x,p)\|.
 \label{eq:V25-adjacent-augmented-action}
\end{equation}
Its constraint component contains
\(\mathcal J_{n+1}(x)-\mathcal J_n(x)\), so adjacent motion of the
connection and jet definition cannot be omitted.

\begin{theorem}[Cofinal contraction for the enlarged first-order state]
\label{thm:V25-augmented-contraction}
Assume:
\begin{enumerate}[label=\textup{(A\arabic*)},leftmargin=3em]
 \item the V23 reduced full-sector inverse has
       \(\gamma_{23}>0\);
 \item \(B_*,J_*\) are bounded on the declared mixed graph levels;
 \item \(\mathfrak m_{25}(R)>0\);
 \item
       \begin{equation}
        \sum_{n\ge n_0}\mathfrak a_n^{25}<\infty.
        \label{eq:V25-augmented-action-sum}
       \end{equation}
\end{enumerate}
Then every \(n\ge n_0\) has a unique enlarged root
\((x_n,p_n)\in\mathbb U_R\), the roots converge in the product graph norm,
and
\begin{equation}
 \boxed{
 \|(x_\infty,p_\infty)-(x_n,p_n)\|
 \le
 \frac{\kappa_{25}}{1-\Theta_{25}(R)}
 \sum_{j\ge n}\mathfrak a_j^{25}.}
 \label{eq:V25-augmented-root-tail}
\end{equation}
For every cutoff and at the limit,
\begin{equation}
 p_n=\mathcal J_n(x_n),\qquad
 p_\infty=\mathcal J_\infty(x_\infty),
 \label{eq:V25-jets-land}
\end{equation}
and the primary fields are exactly the roots of the corresponding reduced
equations.  All compatible curl and Bianchi rows vanish at these roots.
\end{theorem}

\begin{proof}
Theorem~\ref{thm:V25-triangular-factorization} and
\eqref{eq:V25-imported-V23-floor} give the reference augmented inverse
bound \(\kappa_{25}\).  Equations
\eqref{eq:V25-augmented-derivative-radius}--%
\eqref{eq:V25-augmented-master-margin} are precisely the derivative and
center estimates used in the V17 contraction proof, now on the product
graph space.  The positive margin makes
\[
 (x,p)\longmapsto
 (x,p)-\mathbb H_*^{-1}\mathbb F_n(x,p)
\]
a strict self-map contraction of the closed radius-\(R\) ball.
The adjacent-root argument from
Theorem~\ref{thm:V23-master-contraction}, with inverse norm
\(\kappa_{25}\), gives \eqref{eq:V25-augmented-root-tail}.

The second component of every root is the defining equation, so
Theorem~\ref{thm:V25-root-equivalence} gives the reduced roots.
Uniform convergence of the adjacent maps and of their roots passes this
identity to the limit.  Lemma~\ref{lem:V25-covariant-curl} and
Theorem~\ref{thm:V25-Cartan-identities} then annihilate the compatible
curl and Bianchi rows.
\end{proof}

\begin{corollary}[First-order branch discharges the V24 derivative-loss alternative]
\label{cor:V25-discharges-loss}
Suppose every coefficient used by the V24 metric, projector, and collar
chart is a smooth algebraic function of the enlarged state \((x,p)\) at
the \(H^q\) level, the defining map
\(\mathcal J:H^{q+1}\to H^q\) is \(C^2\), and the V25 margin is positive.
Then the field-domain chart is \(C^2\) on the mixed product graph ball
without asserting a false \(H^q\to H^q\) derivative bound for
\(x\mapsto\partial x\).  The V23 full-sector theorem may be applied to
the enlarged residual, with base inverse margin \(\gamma_{25}\).
\end{corollary}

\begin{proof}
The independent variable \(p\) enters the coefficient Nemytskii maps
without differentiation, so Proposition~\ref{prop:V24-tame-jet} applies
with jet order zero to the enlarged coefficient coordinates.  The only
derivative of \(x\) occurs in the defining constraint, where it is bounded
from \(H^{q+1}\) to \(H^q\).  Theorem
\ref{thm:V25-augmented-contraction} supplies the required inverse and
root equivalence.
\end{proof}


% -----------------------------------------------------------------------------
\subsection{Cofinal handoff to the full continuum criterion}
\label{subsec:V25-continuum-handoff}
% -----------------------------------------------------------------------------

\begin{theorem}[Jet-enlarged SM--Einstein continuum implication]
\label{thm:V25-continuum-handoff}
Assume on each compact cylinder:
\begin{enumerate}[label=\textup{(J\arabic*)},leftmargin=3em]
 \item the selected metric or coframe first-order state contains all
       coefficient jets, with defining rows chosen from
       \eqref{eq:V25-defining-constraint} and
       \eqref{eq:V25-curvature-definition-defect}--%
       \eqref{eq:V25-fermion-jet-defect}, and its common-action realization
       is the multiplier action \eqref{eq:V25-KKT-action};
 \item the mixed Sobolev graph levels make every defining map bounded, and
       the V24 metric, projector, collar, and target-frame cards pass on the
       enlarged ball;
 \item the reduced seven-sector Hessian has a positive V23 inverse margin,
       the KKT floor \(\gamma_{25}^{\rm act}\) is positive, and the KKT
       action gradient satisfies the V23 nonlinear radius, fixed-center,
       and summable adjacent-action conditions;
 \item the same-frame incidence between the KKT Euler operator and the
       prolonged causal graph passes, and the V18--V22 main and subsidiary
       symbol, boundary, adjoint-range, and moving-domain cards pass for
       that enlarged first-order system;
 \item the V16--V17 Palatini, stress, product, determinant-line, and
       OS-to-Lorentz rows pass on the same enlarged roots.
\end{enumerate}
Then Theorem~\ref{thm:V23-full-continuum-implication} applies to the
enlarged residual.  Its limit has genuine geometric jets by
\eqref{eq:V25-jets-land}; all curl, Cartan--Bianchi, Gauss, gauge, and
integrability constraints vanish; and its primary fields satisfy all
conclusions of
\eqref{eq:V23-complete-limit-equation}--%
\eqref{eq:V23-full-Einstein-equation}.  No extra physical solution is
created by the enlargement.
\end{theorem}

\begin{proof}
Corollary~\ref{cor:V25-discharges-loss} supplies the \(C^2\) fixed-domain
coordinates required by V23.  Theorems~\ref{thm:V25-KKT-critical} and
\ref{thm:V25-KKT-congruence} supply the common-action residual, its
signed Hessian inverse, and exact projection to the reduced critical
point.  The V23 contraction theorem therefore supplies the KKT roots and
their cofinal graph tail.  The same-frame incidence in \textup{(J4)}
identifies their Euler equations with the prolonged causal graph.
The defining components give genuine jets, and the exact compatibility
identities annihilate the redundant differential rows.  Assumptions
\textup{(J4)--(J5)} are the remaining hypotheses of
Theorem~\ref{thm:V23-full-continuum-implication}; applying it proves the
sector equations and Einstein identity.  The critical-point bijection
proves that no extra primary solution is introduced.
\end{proof}

\begin{assessment}[What V25 resolves and what it leaves physical]
\label{ass:V25-logical-status}
Theorem~\ref{thm:V25-triangular-factorization} proves that a correctly
formed defining-jet enlargement has no new kernel and quantifies its
inverse cost.  Theorem~\ref{thm:V25-compatible-range} prevents a redundant
curl row from being misread as a cokernel.  These statements resolve the
analytic derivative-loss branch conditionally.  They do not prove that
the occurring regulator action differentiates to the selected enlarged
equations, that its augmented symmetrizer is positive, or that its
boundary data preserve every defining constraint.
\end{assessment}


% -----------------------------------------------------------------------------
\subsection{Common-action enlargement requires a multiplier}
\label{subsec:V25-KKT-action}
% -----------------------------------------------------------------------------

The two-variable residual \eqref{eq:V25-augmented-residual} is sufficient
for a prolonged causal PDE graph, but it is not automatically the gradient
of an action.  The common-action V23 card therefore uses the following
multiplier construction.

Let \(\mathcal S_1:\mathscr X\times\mathscr Z\to\mathbb R\) be the
first-order action before imposing the jet definition, and define
\begin{align}
 \mathcal S_{\rm red}(x)
 &:=
 \mathcal S_1(x,\mathcal J(x)),
 \label{eq:V25-reduced-action}\\
 \mathcal S_{\rm KKT}(x,p,\lambda)
 &:=
 \mathcal S_1(x,p)
 +\langle\lambda,p-\mathcal J(x)\rangle_{\mathscr Z}.
 \label{eq:V25-KKT-action}
\end{align}
Real Hilbert identifications are used for the multiplier and dual spaces;
the fermion sector is realified first.

\begin{theorem}[Critical-point equivalence for the first-order common action]
\label{thm:V25-KKT-critical}
Critical points of \(\mathcal S_{\rm red}\) are in bijection with critical
points of \(\mathcal S_{\rm KKT}\).  The lift is
\begin{equation}
 x\longmapsto
 \left(
 x,\mathcal J(x),
 -D_p\mathcal S_1(x,\mathcal J(x))
 \right).
 \label{eq:V25-KKT-critical-lift}
\end{equation}
\end{theorem}

\begin{proof}
The three KKT equations are
\begin{align*}
 D_x\mathcal S_1-J(x)^*\lambda&=0,\\
 D_p\mathcal S_1+\lambda&=0,\\
 p-\mathcal J(x)&=0,
\end{align*}
where \(J(x)=D\mathcal J(x)\).  The last two equations give
\(p=\mathcal J(x)\) and
\(\lambda=-D_p\mathcal S_1(x,\mathcal J(x))\).  Substitution in the first
gives
\[
 D_x\mathcal S_1(x,\mathcal J(x))
 +J(x)^*D_p\mathcal S_1(x,\mathcal J(x))
 =
 D\mathcal S_{\rm red}(x),
\]
which proves both directions and uniqueness of the lifted variables.
\end{proof}

At a KKT critical point, let
\[
 J=D\mathcal J,\qquad
 B=D_{xp}^2\mathcal S_1,\qquad
 C=D_{pp}^2\mathcal S_1,
\]
and let \(A\) denote the actual \(xx\) block of
\(D^2\mathcal S_{\rm KKT}\), including the
\(-\langle\lambda,D^2\mathcal J\rangle\) term.  Then
\begin{equation}
 \mathbb K
 :=
 D^2\mathcal S_{\rm KKT}
 =
 \begin{pmatrix}
  A&B&-J^*\\
  B^*&C&I\\
  -J&I&0
 \end{pmatrix}.
 \label{eq:V25-KKT-Hessian}
\end{equation}
The reduced Hessian is
\begin{equation}
 H_{\rm red}
 =
 A+BJ+J^*B^*+J^*CJ.
 \label{eq:V25-reduced-Hessian}
\end{equation}
Put
\begin{equation}
 D:=B+J^*C
 \label{eq:V25-KKT-D}
\end{equation}
and define
\begin{equation}
 T_{\rm KKT}
 :=
 \begin{pmatrix}
  I&0&0\\
  J&I&0\\
  -D^*&-\tfrac12C&I
 \end{pmatrix}.
 \label{eq:V25-KKT-shear}
\end{equation}

\begin{theorem}[Exact KKT congruence and common-action inverse floor]
\label{thm:V25-KKT-congruence}
The Hessian has the exact congruence
\begin{equation}
 \boxed{
 T_{\rm KKT}^*\mathbb K T_{\rm KKT}
 =
 \begin{pmatrix}
  H_{\rm red}&0&0\\
  0&0&I\\
  0&I&0
 \end{pmatrix}.}
 \label{eq:V25-KKT-congruence}
\end{equation}
Consequently \(\mathbb K\) is invertible if and only if
\(H_{\rm red}\) is invertible.  If
\[
 \|H_{\rm red}^{-1}\|\le\kappa,\qquad
 \|J\|\le j,\qquad
 \|C\|\le c,\qquad
 \|D\|\le d,
\]
and
\begin{equation}
 \tau_{\rm KKT}
 :=
 \left\|
 \begin{pmatrix}
  1&0&0\\
  j&1&0\\
  d&c/2&1
 \end{pmatrix}
 \right\|_{2\to2},
 \label{eq:V25-KKT-shear-bound}
\end{equation}
then
\begin{align}
 \|\mathbb K^{-1}\|
 &\le
 \tau_{\rm KKT}^2\max\{\kappa,1\},
 \label{eq:V25-KKT-inverse-bound}\\
 s_{\min}(\mathbb K)
 &\ge
 \gamma_{25}^{\rm act}
 :=
 \frac1{\tau_{\rm KKT}^2\max\{\kappa,1\}}.
 \label{eq:V25-KKT-action-floor}
\end{align}
Thus a multiplier adds an indefinite unit saddle block but no kernel.
\end{theorem}

\begin{proof}
Use the coordinates
\begin{equation}
 q=Jh+c_0,\qquad
 \mu=\nu-D^*h-\tfrac12Cc_0.
 \label{eq:V25-KKT-coordinate-change}
\end{equation}
The quadratic form of \(\mathbb K\) is
\[
 \langle h,Ah\rangle
 +2\operatorname{Re}\langle h,Bq\rangle
 +\langle q,Cq\rangle
 -2\operatorname{Re}\langle Jh,\mu\rangle
 +2\operatorname{Re}\langle q,\mu\rangle.
\]
Substitute \eqref{eq:V25-KKT-coordinate-change}.  The terms containing
\(\mu\), the \(h\)--\(c_0\) cross terms, and the \(C\) quadratic terms
cancel pairwise, leaving
\[
 \langle h,H_{\rm red}h\rangle
 +2\operatorname{Re}\langle c_0,\nu\rangle.
\]
Polarization proves the operator congruence
\eqref{eq:V25-KKT-congruence}.  Its right side is invertible exactly when
\(H_{\rm red}\) is, because
\(\left(\begin{smallmatrix}0&I\\I&0\end{smallmatrix}\right)\) is a
self-adjoint unitary.

Lemma~\ref{lem:V23-block-majorant} applied to
\eqref{eq:V25-KKT-shear} gives
\(\|T_{\rm KKT}\|\le\tau_{\rm KKT}\).  Inverting the congruence gives
\[
 \mathbb K^{-1}
 =
 T_{\rm KKT}
 \begin{pmatrix}
  H_{\rm red}^{-1}&0&0\\
  0&0&I\\
  0&I&0
 \end{pmatrix}
 T_{\rm KKT}^*.
\]
This proves \eqref{eq:V25-KKT-inverse-bound}; the reciprocal gives
\eqref{eq:V25-KKT-action-floor}.
\end{proof}

\begin{corollary}[Action-compatible replacement of the V25 base margin]
\label{cor:V25-action-margin}
For the common-action V23 card, replace the two-variable base floor
\(\gamma_{25}\) by \(\gamma_{25}^{\rm act}\) from
\eqref{eq:V25-KKT-action-floor}, and issue nonlinear commands in all
\((x,p,\lambda)\) directions.  The V23 contraction proof then applies to
the KKT action gradient.  At every root, Theorem
\ref{thm:V25-KKT-critical} removes \(p,\lambda\) and returns exactly a
critical point of the reduced common action.
\end{corollary}

\begin{proof}
Theorem~\ref{thm:V25-KKT-congruence} supplies the signed physical Hessian
isomorphism required by V17--V23.  Its inverse bound replaces the base
inverse constant in their derivative and center estimates.  The critical
point equivalence supplies the final projection.
\end{proof}

\begin{firewall}[The causal prolongation and the action Hessian are distinct rows]
\label{fire:V25-PDE-versus-action}
The matrix \(\mathbb H\) in
\eqref{eq:V25-augmented-derivative} certifies a prolonged equation graph.
The symmetric KKT matrix \(\mathbb K\) certifies the first variation of
one constrained common action.  They may be related by equation
multipliers, but one cannot be substituted for the other without a
same-frame incidence calculation.  The causal symmetrizer, KKT Hessian,
and their boundary domains remain three separately tested structures.
\end{firewall}


% -----------------------------------------------------------------------------
\subsection{Executable V25 jet and common-action card}
\label{subsec:V25-card}
% -----------------------------------------------------------------------------

The first-order card is assembled before the V23 finite full-sector matrix.

\begin{enumerate}[label=\textup{(\arabic*)},leftmargin=3em]
 \item Choose exactly one metric-coordinate or coframe--connection reduction.
       Print every primary field, independent jet, curvature, torsion, and
       multiplier, with its Sobolev domain and target level.
 \item Construct \(\mathcal J\), the prolonged main residual
       \(\mathcal E\), the reduced residual \(\mathcal R\), and the defining
       residual.  On held-out fields verify the two directions of
       Theorem~\ref{thm:V25-root-equivalence}.
 \item Read every block \(A,B,J\) in one V24 transported frame.  Store their
       finite read, target-leakage, high-domain, and adjacent-cutoff errors.
       Evaluate \(e_H\), the reduced floor, \(\kappa_{\rm aug}\), and
       \(\gamma_{\rm aug}\).
 \item Form each curl or Bianchi operator \(K\) and curvature term \(R\).
       Verify \(KJ=R\) on the finite screen and retain its high graph tail.
       Reduce the target to \(d=Kc\) before running any adjoint-kernel test.
 \item Evaluate all identities
       \eqref{eq:V25-curl-identity} and
       \eqref{eq:V25-curvature-Bianchi-identity}--%
       \eqref{eq:V25-metric-curl-identity} on common-action commands.
       A separately fitted Bianchi packet is inadmissible.
 \item Store the primary and prolonged evolution residuals
       \(E_u,E_i\), the initial defining constraints, and the differentiated
       source norm in \eqref{eq:V25-defining-energy}.  At a timelike
       boundary also retain the V20 incoming constraint flux.
 \item Construct the first-order action \(\mathcal S_1\) and multiplier
       action \(\mathcal S_{\rm KKT}\).  Directly compare its first variation
       with the prolonged causal equations after all equation multipliers,
       gauge reductions, and boundary traces.
 \item Read the KKT blocks \(A,B,C,J\), including
       \(-\langle\lambda,D^2\mathcal J\rangle\), and test the symmetry
       incidences \(B^*,J^*,C=C^*\).  Compute \(D\),
       \(\tau_{\rm KKT}\), and \(\gamma_{25}^{\rm act}\), and verify the
       congruence \eqref{eq:V25-KKT-congruence} on held-out vectors.
 \item Issue nonlinear commands in every primary, jet, and multiplier
       direction.  Form the V23 nonlinear, center, adjacent-action, and
       full high-tail margins using the KKT inverse floor.
 \item Only after these rows pass, project the KKT roots by
       \eqref{eq:V25-KKT-critical-lift} and run the V16--V17 stress,
       Palatini, line, product, and OS-to-Lorentz cards on the same roots.
\end{enumerate}

A concise card output is
\begin{equation}
 \mathfrak C_{25}
 :=
 \left(
 \|J\|,\|B\|,\|C\|,
 \sigma_H-e_H,
 \gamma_{\rm aug},
 \tau_{\rm KKT},
 \gamma_{25}^{\rm act},
 \mathfrak m_{25},
 \sum_n\mathfrak a_n^{25},
 \sup_n\|KJ-R\|,
 \sup_t\|C(t)\|_{H^q}
 \right).
 \label{eq:V25-card-output}
\end{equation}
The KKT nonlinear margin used by the final action card is the V23 margin
with base inverse constant
\((\gamma_{25}^{\rm act})^{-1}\); it is printed separately from the
causal PDE value \(\mathfrak m_{25}\).

Canonical failures include
\begin{align}
 W^{\rm def}_{x,p}
 &:=
 \bigl(x,p,p-\mathcal J(x)\bigr),
 \label{eq:V25-defining-witness}\\
 W^{\rm cx}_{h}
 &:=
 \bigl(h,KJh-Rh\bigr),
 \label{eq:V25-complex-witness}\\
 W^{\rm range}_{c,d}
 &:=
 \bigl(c,d,d-Kc\bigr),
 \label{eq:V25-range-witness}\\
 W^{\rm red}_{N}
 &:=
 \bigl(\widehat H,\widehat h_{\min},\sigma_H,e_H\bigr),
 \label{eq:V25-reduced-witness}\\
 W^{\rm KKT}_{N}
 &:=
 \bigl(\widehat{\mathbb K},
       \widehat v_{\min},
       s_{\min}(\widehat{\mathbb K}),
       \gamma_{25}^{\rm act}\bigr),
 \label{eq:V25-KKT-witness}\\
 W^{\rm prop}_{n}
 &:=
 \bigl(C(0),E_i,D_iE_u,C(t)\bigr),
 \label{eq:V25-propagation-witness}\\
 W^{\rm inc}_{n}
 &:=
 \bigl(D\mathcal S_{\rm KKT},
       \mathcal E_{\rm causal},
       \text{aligned difference}\bigr).
 \label{eq:V25-action-incidence-witness}
\end{align}
These distinguish a bad defining row, failure of the differential complex,
incompatible target data, a reduced soft mode, a KKT soft mode, forced
constraint production, and failure of action-to-causal incidence.


% -----------------------------------------------------------------------------
\subsection{V25 theorem ledger and next upstream row}
\label{subsec:V25-status}
% -----------------------------------------------------------------------------

\paragraph{New conclusions proved in V25.}
Theorem~\ref{thm:V25-root-equivalence} gives exact nonlinear equivalence
between a reduced equation and its defining-jet enlargement.
Theorem~\ref{thm:V25-triangular-factorization} proves that the causal
two-variable enlargement has no new kernel and gives its explicit inverse
cost.  Corollary~\ref{cor:V25-reduced-read-error} propagates finite errors
in the main, jet, and defining maps.  Theorem
\ref{thm:V25-compatible-range} identifies the exact target graph of
redundant curl rows.

Lemma~\ref{lem:V25-covariant-curl} and
Proposition~\ref{prop:V25-defining-propagation} give the exact finite
source of the jet constraints.  Theorem~\ref{thm:V25-Cartan-identities}
assembles the coframe, curvature, torsion, Yang--Mills, Higgs, fermion, and
ghost definition/Bianchi complex.  Theorem
\ref{thm:V25-augmented-contraction} supplies the cofinal enlarged causal
roots on the bounded mixed-level branch.

For the common action, Theorem~\ref{thm:V25-KKT-critical} proves critical
point equivalence and Theorem~\ref{thm:V25-KKT-congruence} reduces the
complete signed KKT Hessian by an exact congruence to the reduced Hessian
plus a unit saddle block.  This supplies the action inverse floor without
confusing it with the causal symmetrizer.

\paragraph{Imported conclusions not repeated.}
The action/response incidence machinery, Kato support curvature, curvature
and stress landing, branch-typed graph inverse, boundary flux, adjoint
range, time and field domain charts, and seven-sector nonlinear master
margin are imported from the attached sources and V1--V24.  V25 constructs
the first-order constraint and multiplier algebra needed to use those
results at mixed Sobolev levels.

\paragraph{Authoritative physical status.}
The attached files do not select and print one complete first-order
SM--Einstein variable list with mixed graph levels.  They do not give the
physical blocks \(A,B,C,J\), the differential-complex tails, the KKT
multipliers, the congruence margin, or a same-frame incidence between the
KKT action gradient and the causal first-order equations.  Nor do they
populate the incoming boundary conditions for every new defining
constraint.  Therefore neither \(\gamma_{25}^{\rm act}\) nor the final
KKT V23 margin can be evaluated, and the full continuum is not claimed.

\paragraph{Next upstream row.}
The next bottleneck is the action-to-causal and boundary incidence of the
enlarged state.  One must construct the equation multiplier taking the
KKT Euler rows to a symmetric-hyperbolic evolution plus zero-speed
definition rows, prove that it is boundedly invertible on the compatible
target, and identify the incoming boundary subspace without prescribing
independent curl data.  Its finite card must combine the V20 flux
projector with the V25 target graph and propagate its singular and
high-frequency tails into \(\gamma_{25}^{\rm act}\).  If no invertible
equation multiplier exists, the action Hessian and causal graph remain
distinct preconditioners and the V23 common-action radius must be proved
with the former while existence is proved with the latter.

\begin{firewall}[A KKT floor is not a causal symmetrizer]
\label{fire:V25-status}
The KKT Hessian is necessarily indefinite because of its unit saddle
block.  Its nonzero singular floor proves an action graph inverse, not a
positive time matrix or a dissipative boundary flux.  Conversely, a
positive prolonged symmetrizer does not prove common-action incidence.
The full programme requires both rows on the same enlarged variables.
\end{firewall}

% END APPENDED VERSION V25 MATERIAL



% =============================================================================
\section{Version V26: action-to-causal incidence on compatible boundary traces}
\label{sec:V26-action-causal}
% =============================================================================

\subsection{Purpose and nonduplication boundary}
\label{subsec:V26-purpose}

V25 separates the signed KKT action Hessian from the prolonged causal
operator and proves that the curl/Bianchi target is a compatibility graph.
V18--V22 supply abstract causal energy, boundary, range, and moving-domain
theorems once the occurring causal operator and boundary space are known.
What remains is an incidence theorem between those two operators on their
actual graph targets, followed by restriction of the boundary flux to the
V25 compatible trace space.

V26 proves both implications with quantitative finite-to-infinite errors.
It does not infer a causal symmetrizer from the KKT Hessian.  Instead it
requires an equation multiplier and measures its failure.  It also avoids
overprescribing the jet curls: the V20 flux theorem is applied after
pullback to the compatible trace coordinates, not to an artificial direct
sum of independent main and curl data.


% -----------------------------------------------------------------------------
\subsection{Exact and approximate equation multipliers}
\label{subsec:V26-linear-incidence}
% -----------------------------------------------------------------------------

Let \(\mathscr X\) be the enlarged KKT unknown graph space,
\(\mathscr Y_{\rm act}\) its compatible action target, and
\(\mathscr Y_{\rm cau}\) the compatible source--initial--boundary target of
the causal first-order system.  Let
\[
 K:\mathscr X\longrightarrow\mathscr Y_{\rm act},
 \qquad
 P:\mathscr X\longrightarrow\mathscr Y_{\rm cau},
 \qquad
 Q:\mathscr Y_{\rm act}\longrightarrow\mathscr Y_{\rm cau}.
\]
Here \(K\) is the V25 KKT Hessian, \(P\) is the prolonged causal graph
operator, and \(Q\) is the equation multiplier.  All three maps include
the declared quotient, trace, line, and V24 moving-domain frames.

Assume
\begin{equation}
 \|K^{-1}\|\le\gamma_K^{-1},\qquad
 \|Q^{-1}\|\le\gamma_Q^{-1},
 \qquad
 \gamma_K,\gamma_Q>0.
 \label{eq:V26-two-floors}
\end{equation}
Define the incidence defect
\begin{equation}
 R_{\rm inc}:=P-QK,
 \qquad
 \varepsilon_{\rm inc}:=\|R_{\rm inc}\|.
 \label{eq:V26-linear-incidence-defect}
\end{equation}

\begin{theorem}[Equation-multiplier incidence and causal inverse]
\label{thm:V26-equation-multiplier}
Put
\begin{equation}
 \eta_{\rm inc}
 :=
 \|K^{-1}Q^{-1}R_{\rm inc}\|
 \le
 \frac{\varepsilon_{\rm inc}}{\gamma_K\gamma_Q}.
 \label{eq:V26-incidence-radius}
\end{equation}
If \(\eta_{\rm inc}<1\), then \(P\) is a bounded isomorphism and
\begin{align}
 P
 &=QK\left(I+K^{-1}Q^{-1}R_{\rm inc}\right),
 \label{eq:V26-incidence-factorization}\\
 P^{-1}
 &=
 \left(I+K^{-1}Q^{-1}R_{\rm inc}\right)^{-1}
 K^{-1}Q^{-1},
 \label{eq:V26-causal-inverse}\\
 \|P^{-1}\|
 &\le
 \frac1{(1-\eta_{\rm inc})\gamma_K\gamma_Q}.
 \label{eq:V26-causal-inverse-bound}
\end{align}
Moreover,
\begin{equation}
 \left\|P^{-1}-K^{-1}Q^{-1}\right\|
 \le
 \frac{\eta_{\rm inc}}
 {(1-\eta_{\rm inc})\gamma_K\gamma_Q}.
 \label{eq:V26-inverse-incidence-error}
\end{equation}
When \(R_{\rm inc}=0\), the two graph inverses are exactly incident:
\(P^{-1}=K^{-1}Q^{-1}\).
\end{theorem}

\begin{proof}
The factorization is direct:
\[
 QK\left(I+K^{-1}Q^{-1}R_{\rm inc}\right)
 =QK+R_{\rm inc}=P.
\]
If \(\eta_{\rm inc}<1\), the last factor is invertible by its Neumann
series, which gives \eqref{eq:V26-causal-inverse}.  Multiplication of the
three inverse bounds proves
\eqref{eq:V26-causal-inverse-bound}.  If
\(T=K^{-1}Q^{-1}R_{\rm inc}\), then
\[
 (I+T)^{-1}-I=-(I+T)^{-1}T.
\]
Its norm is at most \(\eta_{\rm inc}/(1-\eta_{\rm inc})\); multiply by
\(\|K^{-1}Q^{-1}\|\) to obtain
\eqref{eq:V26-inverse-incidence-error}.
\end{proof}

\begin{firewall}[A large incidence upper bound is inconclusive]
\label{fire:V26-large-incidence}
Failure of \(\eta_{\rm inc}<1\) does not prove that \(P\) has a kernel.
It says that the declared KKT inverse and equation multiplier do not
certify \(P\) by this perturbative factorization.  The correct next
objects are the exact singular vectors of \(P\) or a sharper structured
factorization, together with the block of \(R_{\rm inc}\) attaining the
large response.
\end{firewall}

A finite equation-multiplier read requires its own tail.  Let
\(\widehat Q_N\) be finite rank apart from the identity on an independently
declared common complement.  Suppose
\begin{equation}
 \sigma_{Q,N}:=s_{\min}(\widehat Q_N)>0,
 \qquad
 \|Q-\widehat Q_N\|\le e_{Q,N}.
 \label{eq:V26-finite-Q-data}
\end{equation}

\begin{lemma}[Finite-to-infinite equation-multiplier floor]
\label{lem:V26-Q-floor}
If
\begin{equation}
 \gamma_{Q,N}:=\sigma_{Q,N}-e_{Q,N}>0,
 \label{eq:V26-Q-margin}
\end{equation}
then \(Q\) is invertible and
\begin{equation}
 \|Q^{-1}\|\le\gamma_{Q,N}^{-1}.
 \label{eq:V26-Q-inverse-bound}
\end{equation}
If the KKT floor is \(\gamma_K\), a sufficient measured incidence condition
is
\begin{equation}
 \boxed{
 \varepsilon_{\rm inc}^{\rm cert}
 <
 \gamma_K\gamma_{Q,N}.}
 \label{eq:V26-certified-incidence}
\end{equation}
\end{lemma}

\begin{proof}
Apply the Neumann perturbation lemma to
\(\widehat Q_N^{-1}(Q-\widehat Q_N)\).  The inverse bound is
\((\sigma_{Q,N}-e_{Q,N})^{-1}\).  Substitute it and
\(\varepsilon_{\rm inc}\le\varepsilon_{\rm inc}^{\rm cert}\) in
\eqref{eq:V26-incidence-radius}.
\end{proof}

The tail \(e_{Q,N}\) contains the finite read error, target leakage,
high-domain tail, adjacent-cutoff tail, and V24 frame derivative error,
combined by the V23 block majorant.  A finite singular value without these
terms is not \(\gamma_{Q,N}\).


% -----------------------------------------------------------------------------
\subsection{Nonlinear residual incidence and its connection term}
\label{subsec:V26-nonlinear-incidence}
% -----------------------------------------------------------------------------

Let \(\mathcal A_n(z)\) be the pulled KKT action gradient and
\(\mathcal P_n(z)\) the pulled prolonged causal residual.  An occurring
nonlinear equation multiplier is a \(C^1\) map
\[
 Q_n(z):
 \mathscr Y_{\rm act}\longrightarrow\mathscr Y_{\rm cau}.
\]
Define the exact nonlinear incidence residual
\begin{equation}
 r_n^{\rm inc}(z)
 :=
 \mathcal P_n(z)-Q_n(z)\mathcal A_n(z).
 \label{eq:V26-nonlinear-incidence}
\end{equation}

\begin{lemma}[Derivative of nonlinear action-to-causal incidence]
\label{lem:V26-nonlinear-derivative}
For every direction \(h\),
\begin{equation}
 \boxed{
 D\mathcal P_n(z)h
 =
 \bigl(DQ_n(z)[h]\bigr)\mathcal A_n(z)
 +Q_n(z)D\mathcal A_n(z)h
 +Dr_n^{\rm inc}(z)h.}
 \label{eq:V26-nonlinear-incidence-derivative}
\end{equation}
At an action root \(z_n\),
\begin{align}
 \mathcal P_n(z_n)&=r_n^{\rm inc}(z_n),
 \label{eq:V26-causal-value-at-action-root}\\
 D\mathcal P_n(z_n)
 &=
 Q_n(z_n)D\mathcal A_n(z_n)
 +Dr_n^{\rm inc}(z_n).
 \label{eq:V26-causal-derivative-at-root}
\end{align}
In particular, if \(r_n^{\rm inc}\) vanishes identically and \(Q_n(z)\)
is invertible, the action and causal residuals have exactly the same zero
set.
\end{lemma}

\begin{proof}
Differentiate \eqref{eq:V26-nonlinear-incidence} and apply the product
rule.  Equations
\eqref{eq:V26-causal-value-at-action-root}--%
\eqref{eq:V26-causal-derivative-at-root} follow after setting
\(\mathcal A_n(z_n)=0\).  If the incidence residual is zero and \(Q_n\)
is invertible, \(Q_n\mathcal A_n=0\) is equivalent to
\(\mathcal A_n=0\).
\end{proof}

Fix a base point \(z_*\), put \(Q_*=Q_n(z_*)\), and define on a radius
\(R\) ball
\begin{align}
 q_1(R)&:=\sup_z\|DQ_n(z)\|,
 &
 a_0(R)&:=\sup_z\|\mathcal A_n(z)\|,
 \label{eq:V26-Q-action-bounds}\\
 q_\Delta(R)&:=\sup_z\|Q_n(z)-Q_*\|,
 &
 k_0(R)&:=\sup_z\|D\mathcal A_n(z)\|,
 \label{eq:V26-Q-Hessian-bounds}\\
 r_1(R)&:=\sup_z\|Dr_n^{\rm inc}(z)\|.
 \label{eq:V26-incidence-derivative-bound}
\end{align}
Then
\begin{equation}
 \boxed{
 \sup_z
 \|D\mathcal P_n(z)-Q_*D\mathcal A_n(z)\|
 \le
 q_1(R)a_0(R)+q_\Delta(R)k_0(R)+r_1(R).}
 \label{eq:V26-nonlinear-master-incidence}
\end{equation}
The first term is the equation-multiplier analogue of the V23 frame
connection.  It vanishes at an exact action root but not throughout a
contraction ball.

\begin{theorem}[Cofinal landing of approximate nonlinear incidence]
\label{thm:V26-cofinal-incidence}
Let \(z_n\to z_\infty\) be the V23--V25 common-action roots.  Suppose
\(\mathcal P_n\) converges locally uniformly in the causal target norm to
\(\mathcal P_\infty\), and
\begin{equation}
 \|r_n^{\rm inc}(z_n)\|_{\mathscr Y_{\rm cau}}\longrightarrow0.
 \label{eq:V26-value-incidence-landing}
\end{equation}
Then
\begin{equation}
 \mathcal P_\infty(z_\infty)=0.
 \label{eq:V26-limit-causal-equation}
\end{equation}
If, in addition, \(D\mathcal A_n(z_n)\) and \(Q_n(z_n)\) have uniform
floors \(\gamma_K,\gamma_Q\) and
\begin{equation}
 \frac{\|Dr_n^{\rm inc}(z_n)\|}
 {\gamma_K\gamma_Q}
 \le\eta_*<1,
 \label{eq:V26-root-linear-incidence}
\end{equation}
then every linearized causal graph at \(z_n\) is invertible with
\begin{equation}
 \|\bigl(D\mathcal P_n(z_n)\bigr)^{-1}\|
 \le
 \frac1{(1-\eta_*)\gamma_K\gamma_Q}.
 \label{eq:V26-root-causal-inverse}
\end{equation}
\end{theorem}

\begin{proof}
Equation \eqref{eq:V26-causal-value-at-action-root} and
\eqref{eq:V26-value-incidence-landing} give
\(\mathcal P_n(z_n)\to0\).  Local uniform convergence and
\(z_n\to z_\infty\) give
\(\mathcal P_\infty(z_\infty)=0\).  At a root, use
\eqref{eq:V26-causal-derivative-at-root} in
Theorem~\ref{thm:V26-equation-multiplier}, with
\(K=D\mathcal A_n(z_n)\), \(Q=Q_n(z_n)\), and
\(R_{\rm inc}=Dr_n^{\rm inc}(z_n)\).
\end{proof}

\begin{firewall}[Value incidence and derivative incidence are independent]
\label{fire:V26-two-incidences}
A small \(r_n^{\rm inc}(z_n)\) proves that an action root nearly solves the
causal equation.  It does not bound \(Dr_n^{\rm inc}(z_n)\), which decides
the causal graph inverse.  Conversely, a small derivative residual can
coexist with a nonzero constant value residual.  Both rows are required
in Theorem~\ref{thm:V26-cofinal-incidence}.
\end{firewall}


% -----------------------------------------------------------------------------
\subsection{Pull the boundary flux back to the compatible trace space}
\label{subsec:V26-compatible-boundary}
% -----------------------------------------------------------------------------

Let \(\mathscr B\) be the total first-order trace fibre with self-adjoint
normal flux \(\mathfrak F_\nu\).  Let \(\mathscr C_\partial\) be the
independent compatible trace-coordinate space and
\begin{equation}
 T_\partial:\mathscr C_\partial\longrightarrow\mathscr B
 \label{eq:V26-boundary-compatibility-frame}
\end{equation}
a bounded injective frame with closed range.  Its columns impose the V25
defining, curl, Cartan, and multiplier trace identities.  Assume
\begin{equation}
 s_\partial\|c\|
 \le\|T_\partial c\|
 \le t_\partial\|c\|.
 \label{eq:V26-boundary-frame-bounds}
\end{equation}
Define the compatible trace fibre and pulled flux by
\begin{equation}
 \mathscr B_{\rm comp}:=\operatorname{Ran}T_\partial,
 \qquad
 \mathfrak F_{\rm comp}:=
 T_\partial^*\mathfrak F_\nu T_\partial.
 \label{eq:V26-pulled-compatible-flux}
\end{equation}

\begin{theorem}[Maximal boundary dissipation on a compatibility graph]
\label{thm:V26-compatible-flux}
The operator \(\mathfrak F_{\rm comp}\) is self-adjoint.  For every linear
subspace \(\mathscr M\subset\mathscr C_\partial\):
\begin{enumerate}[label=\textup{(\roman*)},leftmargin=2.8em]
 \item \(\mathscr M\) is nonnegative for
       \(\mathfrak F_{\rm comp}\) if and only if
       \(T_\partial\mathscr M\) is nonnegative for
       \(\mathfrak F_\nu\);
 \item \(\mathscr M\) is maximal nonnegative in
       \(\mathscr C_\partial\) if and only if
       \(T_\partial\mathscr M\) is maximal nonnegative within
       \(\mathscr B_{\rm comp}\);
 \item if
       \(\langle\mathfrak F_{\rm comp}c,c\rangle
       \ge\beta\|c\|^2\) on a coercive part of \(\mathscr M\), then
       \begin{equation}
        \langle\mathfrak F_\nu T_\partial c,T_\partial c\rangle
        \ge\frac{\beta}{t_\partial^2}
        \|T_\partial c\|^2
        \label{eq:V26-compatible-physical-margin}
       \end{equation}
       there.
\end{enumerate}
Thus the V20 incoming/outgoing/zero-speed construction must be performed
for \(\mathfrak F_{\rm comp}\).  Its image is the physical boundary
space; no independent curl datum is prescribed.
\end{theorem}

\begin{proof}
Self-adjointness is immediate.  The identity
\[
 \langle\mathfrak F_\nu T_\partial c,T_\partial c\rangle
 =
 \langle\mathfrak F_{\rm comp}c,c\rangle
\]
proves the first assertion.  Since \(T_\partial\) is a bijection from
\(\mathscr C_\partial\) onto \(\mathscr B_{\rm comp}\), strict inclusion
of a nonnegative subspace on one side is equivalent to strict inclusion
of its image or preimage on the other; this proves maximality.  Finally,
\(\|T_\partial c\|\le t_\partial\|c\|\) gives
\eqref{eq:V26-compatible-physical-margin}.
\end{proof}

A frequent case has a main trace \(v\) and an induced constraint trace
\(L_\partial v\):
\begin{equation}
 T_\partial v=\binom{v}{L_\partial v},
 \qquad
 \mathfrak F_\nu=
 \begin{pmatrix}
  \mathfrak F_m&0\\0&\mathfrak F_c
 \end{pmatrix}.
 \label{eq:V26-induced-constraint-trace}
\end{equation}
Then
\begin{equation}
 \mathfrak F_{\rm comp}
 =
 \mathfrak F_m+L_\partial^*\mathfrak F_cL_\partial.
 \label{eq:V26-induced-compatible-flux}
\end{equation}

\begin{corollary}[Explicit constraint-flux erosion of a main reserve]
\label{cor:V26-flux-erosion}
Suppose on a proposed main allowed space \(\mathscr M_m\),
\[
 \langle\mathfrak F_mv,v\rangle\ge\beta_m\|v\|^2,
 \qquad
 \mathfrak F_c\succeq-\kappa_cI.
\]
Then
\begin{equation}
 \boxed{
 \langle\mathfrak F_{\rm comp}v,v\rangle
 \ge
 \left(\beta_m-\kappa_c\|L_\partial\|^2\right)\|v\|^2.}
 \label{eq:V26-flux-erosion}
\end{equation}
A positive right side is a sufficient compatible boundary reserve.  If
the defining and curl variables are genuinely zero speed,
\(\mathfrak F_c=0\), they cause no flux erosion but their algebraic
compatibility is still required.
\end{corollary}

\begin{proof}
Use \eqref{eq:V26-induced-compatible-flux} and
\[
 \langle\mathfrak F_cL_\partial v,L_\partial v\rangle
 \ge-\kappa_c\|L_\partial v\|^2.
\]
\end{proof}


% -----------------------------------------------------------------------------
\subsection{Finite compatible-flux certificate}
\label{subsec:V26-finite-boundary}
% -----------------------------------------------------------------------------

Let \(\widehat T_\partial\) and
\(\widehat{\mathfrak F}_\nu\) be finite reads in the same frame.  Write
\begin{gather}
 \|T_\partial-\widehat T_\partial\|\le e_T,\qquad
 \|\mathfrak F_\nu-\widehat{\mathfrak F}_\nu\|\le e_F,
 \label{eq:V26-boundary-read-errors}\\
 t=\|T_\partial\|,\quad
 \widehat t=\|\widehat T_\partial\|,\quad
 f=\|\mathfrak F_\nu\|,\quad
 \widehat f=\|\widehat{\mathfrak F}_\nu\|.
 \label{eq:V26-boundary-read-norms}
\end{gather}
Define
\[
 \widehat{\mathfrak F}_{\rm comp}
 :=
 \widehat T_\partial^*
 \widehat{\mathfrak F}_\nu
 \widehat T_\partial.
\]

\begin{lemma}[Error in the pulled compatible flux]
\label{lem:V26-compatible-flux-error}
The exact bound is
\begin{equation}
 \boxed{
 \|\mathfrak F_{\rm comp}
   -\widehat{\mathfrak F}_{\rm comp}\|
 \le
 e_Tft+\widehat t\,e_Ft+\widehat t\,\widehat f\,e_T.}
 \label{eq:V26-compatible-flux-error}
\end{equation}
If a proposed coordinate boundary space \(\mathscr M\) satisfies
\[
 \langle\widehat{\mathfrak F}_{\rm comp}c,c\rangle
 \ge\widehat\beta\|c\|^2
 \quad(c\in\mathscr M),
\]
then the exact compatible reserve is at least
\begin{equation}
 \beta_{\rm comp}^{\rm cert}
 :=
 \widehat\beta
 -e_Tft-\widehat t\,e_Ft-\widehat t\,\widehat f\,e_T.
 \label{eq:V26-compatible-certified-reserve}
\end{equation}
A positive value certifies nonnegative, and on the strict component
coercive, physical flux after applying
Theorem~\ref{thm:V26-compatible-flux}.
\end{lemma}

\begin{proof}
Insert two intermediate terms:
\begin{align*}
 T_\partial^*\mathfrak F_\nu T_\partial
 -\widehat T_\partial^*
  \widehat{\mathfrak F}_\nu\widehat T_\partial
 &=
 (T_\partial-\widehat T_\partial)^*
 \mathfrak F_\nu T_\partial\\
 &\quad+
 \widehat T_\partial^*
 (\mathfrak F_\nu-\widehat{\mathfrak F}_\nu)T_\partial\\
 &\quad+
 \widehat T_\partial^*
 \widehat{\mathfrak F}_\nu
 (T_\partial-\widehat T_\partial).
\end{align*}
Taking norms proves \eqref{eq:V26-compatible-flux-error}.  The quadratic
form perturbation inequality then gives
\eqref{eq:V26-compatible-certified-reserve}.
\end{proof}

The error \(e_T\) includes the V24 field-domain chord and the V25
defining/curl compatibility tail.  The error \(e_F\) includes the V19
symbol extraction, V20 flux polarization, and V22 time transport tails.
If the coordinate allowed space is itself reconstructed by a spectral
projector, its V20/V24 projector error is an additional term; it cannot
be absorbed into \(e_T\) unless the corresponding subspace transport has
actually been applied.


% -----------------------------------------------------------------------------
\subsection{Principal-symbol obstruction to a local equation multiplier}
\label{subsec:V26-symbol-incidence}
% -----------------------------------------------------------------------------

The graph multiplier must also be local in the sense declared by the
physical reduction.  A zeroth-order multiplier has a sharp symbol test.

\begin{theorem}[Symbol incidence forced by a local equation multiplier]
\label{thm:V26-symbol-incidence}
Let \(K\) and \(P\) be first-order differential operators between equal
rank bundles, and let \(Q(x)\) be a zeroth-order bundle map.  If
\begin{equation}
 P=QK+R_0
 \label{eq:V26-local-equation-multiplier}
\end{equation}
with \(R_0\) of order zero, then for every real covector \(\xi\),
\begin{equation}
 \boxed{
 \sigma_1(P)(x,\xi)
 =
 Q(x)\sigma_1(K)(x,\xi).}
 \label{eq:V26-principal-symbol-incidence}
\end{equation}
If \(Q(x)\) is invertible, then
\begin{align}
 \ker\sigma_1(P)(x,\xi)
 &=\ker\sigma_1(K)(x,\xi),
 \label{eq:V26-symbol-kernel-incidence}\\
 \det\sigma_1(P)(x,\xi)
 &=\det Q(x)\det\sigma_1(K)(x,\xi).
 \label{eq:V26-characteristic-incidence}
\end{align}
Thus a rank, right-kernel, or characteristic-set mismatch at one
\((x,\xi)\) rules out an invertible local zeroth-order equation
multiplier.

More quantitatively, if
\[
 \|Q(x)^{-1}\|\le\gamma_Q^{-1},\qquad
 \|\sigma_1(P)-Q\sigma_1(K)\|\le\epsilon_\sigma,
\]
then
\begin{equation}
 s_{\min}\bigl(\sigma_1(P)(x,\xi)\bigr)
 \ge
 \gamma_Qs_{\min}\bigl(\sigma_1(K)(x,\xi)\bigr)
 -\epsilon_\sigma.
 \label{eq:V26-symbol-floor-transfer}
\end{equation}
\end{theorem}

\begin{proof}
The order-one part of \(QK\) is \(Q\sigma_1(K)\), while \(R_0\) has zero
principal symbol.  This proves
\eqref{eq:V26-principal-symbol-incidence}.  Left multiplication by an
invertible matrix preserves the right kernel and multiplies the
determinant, proving the next two identities.  Finally, for every unit
vector \(v\),
\begin{align*}
 \|\sigma_1(P)v\|
 &\ge
 \|Q\sigma_1(K)v\|-\epsilon_\sigma\\
 &\ge
 \gamma_Q\|\sigma_1(K)v\|-\epsilon_\sigma.
\end{align*}
Take the infimum over \(v\).
\end{proof}

\begin{assessment}[Order mismatch is a graph-level decision]
\label{ass:V26-order-mismatch}
If the KKT Euler operator and prolonged causal operator have different
orders in the selected variables, their equation multiplier cannot be an
invertible zeroth-order map satisfying
\eqref{eq:V26-local-equation-multiplier}.  A differential multiplier must
then be placed between shifted Sobolev target levels and its kernel,
range, boundary trace, and high-frequency tail tested anew.  An inverse
defined separately at each nonzero covector is generally
pseudodifferential and cannot be declared a local field equation without
that additional proof.
\end{assessment}


% -----------------------------------------------------------------------------
\subsection{Cofinal action, causal, and compatible-boundary handoff}
\label{subsec:V26-cofinal}
% -----------------------------------------------------------------------------

\begin{theorem}[Equation-multiplier completion of the V25 handoff]
\label{thm:V26-cofinal-handoff}
Assume on one cofinal family and one compact cylinder:
\begin{enumerate}[label=\textup{(E\arabic*)},leftmargin=3em]
 \item the V25 KKT action gradient has a uniform inverse floor
       \(\gamma_K>0\), satisfies the V23 nonlinear root conditions, and
       produces roots \(z_n\to z_\infty\);
 \item the equation multiplier has a finite-to-infinite floor
       \(\gamma_Q>0\), preserves every V25 compatible target row, and its
       principal symbols satisfy \eqref{eq:V26-principal-symbol-incidence}
       up to a vanishing measured tail;
 \item the value incidence obeys
       \eqref{eq:V26-value-incidence-landing}, while the derivative incidence
       obeys \eqref{eq:V26-root-linear-incidence};
 \item the prolonged operator has the V18 symmetrizer and coefficient
       bounds, and its V19--V22 interior, compatible-boundary, adjoint-range,
       moving-domain, corner, and constraint cards pass;
 \item the boundary frame has a uniform column floor, its pulled compatible
       flux passes the V20 incoming/zero-speed construction with
       \(\beta_{\rm comp}^{\rm cert}>0\) on every strict incoming component,
       and no independent curl data are prescribed;
 \item the V16--V17 Palatini, stress, product, line, and OS-to-Lorentz rows
       pass on these same roots.
\end{enumerate}
Then:
\begin{enumerate}[label=\textup{(\roman*)},leftmargin=2.8em]
 \item the linearized prolonged causal graph is a uniformly bounded
       isomorphism, with bound \eqref{eq:V26-root-causal-inverse};
 \item the limiting action root solves the limiting causal first-order
       system with the compatible boundary condition;
 \item the V25 defining, curl, Bianchi, gauge, Gauss, Hamiltonian, momentum,
       and matter-integrability constraints vanish in their declared norms;
 \item the primary limit satisfies every conclusion of
       Theorem~\ref{thm:V23-full-continuum-implication}, including the
       distributional Einstein equation
       \eqref{eq:V23-full-Einstein-equation}.
\end{enumerate}
\end{theorem}

\begin{proof}
Items \textup{(E1)--(E3)} and
Theorem~\ref{thm:V26-cofinal-incidence} give the causal linearized inverse
and the limiting causal equation.  The symbol incidence ensures that this
operator is the occurring local prolongation rather than a separately
fitted graph map.  Items \textup{(E4)--(E5)},
Theorem~\ref{thm:V26-compatible-flux}, and the V20--V22 results give the
causal range with the induced physical boundary and corner data.

The defining constraints are components of the enlarged roots.  The V25
compatibility identities then annihilate the curl and Bianchi rows, while
the imported subsidiary estimates annihilate the gauge, Gauss, and
Einstein constraints.  Item \textup{(E6)} supplies the remaining common
action, stress, and continuation passages.  The V23 full-continuum
implication now applies.
\end{proof}

\begin{firewall}[Exact symbol incidence is not full graph incidence]
\label{fire:V26-symbol-not-graph}
Equation \eqref{eq:V26-principal-symbol-incidence} controls only the
highest-order response.  Lower-order connection terms, KKT multipliers,
initial data, boundary traces, and target compatibility can still make
\(R_{\rm inc}\ne0\).  The full graph residual
\eqref{eq:V26-linear-incidence-defect} and its nonlinear value and
derivative versions remain required.
\end{firewall}


% -----------------------------------------------------------------------------
\subsection{Executable V26 incidence and compatible-boundary card}
\label{subsec:V26-card}
% -----------------------------------------------------------------------------

A V26 card is evaluated after the V25 KKT and causal reductions have been
constructed, but before either is used as the V23 preconditioner.

\begin{enumerate}[label=\textup{(\arabic*)},leftmargin=3em]
 \item Freeze the enlarged unknown, action target, causal
       source--initial--boundary target, compatible curl graphs, Sobolev
       levels, physical quotient, and V24 frames.
 \item Store the KKT Hessian \(K_n\), its V25 finite-to-infinite floor
       \(\gamma_{K,n}\), and the prolonged causal graph \(P_n\) on exactly
       the same unknown.
 \item Issue source-complete action-target commands for \(Q_n\), including
       held-out mixed sector and multiplier directions.  Store its finite
       floor, read error, target leakage, high-domain tail, and adjacent
       tail, then compute \(\gamma_{Q,N}\).
 \item At the V19 oscillatory scales compare
       \(\sigma_1(P_n)\) with \(Q_n\sigma_1(K_n)\) on a finite
       base--covector--polarization net plus its Lipschitz and high-frequency
       tails.  Return any rank, right-kernel, or characteristic mismatch
       before testing lower-order incidence.
 \item Compute the full graph residual
       \(R_{\rm inc}=P_n-Q_nK_n\), including initial and boundary blocks.
       On the nonlinear ball separately retain
       \(r_n^{\rm inc}(z)\), \(Dr_n^{\rm inc}(z)\), \(DQ_n(z)\), and all
       terms in \eqref{eq:V26-nonlinear-master-incidence}.
 \item Build the V25 compatible boundary frame \(T_{\partial,n}\) from
       independent trace coordinates.  Test its column floor and every
       defining, curl, Cartan, and multiplier trace incidence on held-out
       data.
 \item Pull the occurring normal flux back by
       \eqref{eq:V26-pulled-compatible-flux}.  Store \(e_T,e_F\), compute
       \(\beta_{\rm comp}^{\rm cert}\), and apply the V20
       incoming/outgoing/zero-speed construction to this pulled matrix.
 \item Run the V21 adjoint boundary map and V22 field/time domain transport
       on the compatible trace space.  Initial and corner data must lie in
       the same compatibility graph; no curl or Bianchi datum is chosen
       independently.
 \item Report the causal inverse bound
       \eqref{eq:V26-root-causal-inverse}, the value-incidence cofinal tail,
       the constraint boundary debit, and the V23--V25 root and stress
       tails on one sequence.
\end{enumerate}

For \(\sigma_{Q,N}>0\), define the dimensionless margins
\begin{align}
 \mu_Q
 &:=
 1-\frac{e_{Q,N}}{\sigma_{Q,N}},
 \label{eq:V26-Q-relative-margin}\\
 \mu_{\rm inc}
 &:=
 1-\frac{\varepsilon_{\rm inc}^{\rm cert}}
          {\gamma_K\gamma_{Q,N}},
 \label{eq:V26-incidence-relative-margin}\\
 \mu_\partial
 &:=
 \frac{\beta_{\rm comp}^{\rm cert}}
 {1+\|\widehat{\mathfrak F}_{\rm comp}\|}.
 \label{eq:V26-boundary-relative-margin}
\end{align}
The compact pass number is
\begin{equation}
 \boxed{
 \mathfrak m_{26}
 :=
 \min\{\mu_Q,\mu_{\rm inc},\mu_\partial\}.}
 \label{eq:V26-master-margin}
\end{equation}
A positive value certifies the three finite inverse, incidence, and strict
boundary reserves.  Vanishing nonlinear value incidence and the cofinal
action/stress sums remain separate convergence requirements.

Canonical failure witnesses include
\begin{align}
 W_N^Q
 &:=
 \bigl(\widehat Q_N,\widehat y_{\min},
       \sigma_{Q,N},e_{Q,N}\bigr),
 \label{eq:V26-Q-witness}\\
 W_n^{\rm sym}
 &:=
 \bigl(x,\xi,v,
       [\sigma_1(P_n)-Q_n\sigma_1(K_n)]v\bigr),
 \label{eq:V26-symbol-witness}\\
 W_n^{\rm inc}
 &:=
 \bigl(z,h,r_n^{\rm inc}(z),
       Dr_n^{\rm inc}(z)h\bigr),
 \label{eq:V26-incidence-witness}\\
 W_n^{\partial{\rm comp}}
 &:=
 \bigl(c,T_{\partial,n}c,
       \text{defining/curl trace residual}\bigr),
 \label{eq:V26-boundary-compatibility-witness}\\
 W_n^{\partial{\rm flux}}
 &:=
 \bigl(c,\widehat{\mathfrak F}_{\rm comp}c,
       \beta_{\rm comp}^{\rm cert}\bigr).
 \label{eq:V26-boundary-flux-witness}
\end{align}
They distinguish a multiplier soft mode, symbol mismatch, value or
derivative incidence failure, incompatible trace, and compatible incoming
flux failure.


% -----------------------------------------------------------------------------
\subsection{V26 theorem ledger and next upstream row}
\label{subsec:V26-status}
% -----------------------------------------------------------------------------

\paragraph{New conclusions proved in V26.}
Theorem~\ref{thm:V26-equation-multiplier} gives the exact graph
factorization and inverse bound relating the V25 KKT Hessian to the
prolonged causal operator, including a measured incidence defect.
Lemma~\ref{lem:V26-Q-floor} gives the finite-to-infinite multiplier floor.
Lemma~\ref{lem:V26-nonlinear-derivative} separates nonlinear value
incidence, derivative incidence, and the equation-frame connection term.
Theorem~\ref{thm:V26-cofinal-incidence} lands the causal equation on the
action roots and supplies its linearized inverse.

Theorem~\ref{thm:V26-compatible-flux} pulls the total normal flux back to
the V25 compatible trace fibre and proves maximality there.
Corollary~\ref{cor:V26-flux-erosion} computes the constraint-flux cost,
and Lemma~\ref{lem:V26-compatible-flux-error} propagates the finite frame
and flux errors.  Theorem~\ref{thm:V26-symbol-incidence} gives a sharp
rank, kernel, and characteristic obstruction to a local zeroth-order
equation multiplier.

\paragraph{Imported conclusions not repeated.}
The KKT action and jet-complex algebra, common-action root, principal
symbol extraction, hyperbolic energy, incoming graph theorem, adjoint
range, moving-domain chart, and stress/Einstein landing remain the results
of V16--V25 and the attached sources.  V26 supplies their missing
action-to-causal and compatible-boundary incidence.

\paragraph{Authoritative physical status.}
The supplied manuscripts do not print an equation multiplier between one
complete KKT first-order SM--Einstein action target and one prolonged
causal target.  They do not populate its singular floor, its principal
symbol incidence, its full value and derivative residuals, or the
compatible boundary frame containing all jet and multiplier traces.
Thus \(\mathfrak m_{26}\) cannot be evaluated from the attached tables.
The hypotheses of Theorem~\ref{thm:V26-cofinal-handoff} are not physically
verified, and the full continuum is not claimed.

\paragraph{Next upstream row.}
The equation multiplier should next be derived rather than fitted.
Starting from the selected first-order KKT action, isolate the
velocity--momentum Legendre block after gauge and constraint reduction.
A uniform hyperregularity floor should construct the local order-zero
map from Euler rows to evolution rows and its positive symmetrizer.
The proof must retain the fermionic first-order block and all gauge null
directions separately.  A remaining Legendre kernel after the physical
quotient is an explicit obstruction; a covector-dependent inverse without
a locality proof is not an admissible replacement.  The resulting
multiplier must then be inserted into the compatible flux card above.

\begin{firewall}[V26 is an incidence theorem, not a populated multiplier]
\label{fire:V26-status}
Positive KKT and causal inverse floors do not construct \(Q\).  Two
invertible operators can have different local symbols, boundary traces,
and physical equations.  The multiplier, its locality and target
compatibility, and both incidence residuals are independent physical
rows until the Legendre construction or an equivalent occurring
calculation supplies them.
\end{firewall}

% END APPENDED VERSION V26 MATERIAL

% ============================================================================
% APPENDED VERSION V27 MATERIAL
% ============================================================================

\section{V27: Partial Legendre Convexity, Hamilton--Pontryagin Incidence,
and the Differential Equation Multiplier}
\label{sec:V27-legendre-hp}

\subsection{Purpose, inherited input, and the exact gap}
\label{subsec:V27-purpose}

V26 isolated a principal-symbol obstruction: a pointwise order-zero matrix
cannot in general turn an Euler--Lagrange or KKT target into a prolonged
first-order causal target.  Its next row asked whether the multiplier could
instead be derived from the action.  This note carries out that derivation
for a gauge-reduced bosonic block whose velocity Hessian is uniformly
positive and whose spatial Hessian has the hyperbolic sign.  It also records
exactly why the same argument cannot be applied to an unreduced gauge block,
a first-order fermion block, an indefinite ghost block, or the unreduced
DeWitt conformal direction.

One result is imported and not reproved.  The typed matter Legendre theorem
in the attached grand-tensor manuscript treats a finite Lagrangian
\(\mathscr L(z,v)\), proves the reciprocal identities for
\(p=D_v\mathscr L\) and its Hamiltonian, and proposes the finite residual
\(\|A_hK_h-I\|^2+\|K_hA_h-I\|^2\).  The new work below is the spatial-jet
extension: it computes the full partial-Hamiltonian Hessian, proves its
positive floor, constructs the symmetric-hyperbolic principal matrices,
and derives the first-order differential map between action residuals and
causal residuals.  Thus none of the finite Legendre reciprocity already in
the attachment is counted as a V27 advance.

Fix a spacetime cylinder \(I\times\Sigma\), a local physical bundle
\(Y_{\rm phys}\to I\times\Sigma\), and a bosonic field value \(y\).  Write
\[
 v=\partial_t y,
 \qquad
 w=(w_1,\ldots,w_d),
 \qquad
 w_i=\partial_i y.
\]
All adjoints in this note are taken in the declared fibre metrics and
density.  Gauge fixing and constraint reduction are part of the declaration
of \(Y_{\rm phys}\); they are not conclusions of the Legendre calculation.
The calculation is therefore conditional on a physical slice having
already been constructed and compared with the V25 compatible jet complex.

% -----------------------------------------------------------------------------
\subsection{Uniform partial Legendre transform on the physical bosonic fibre}
\label{subsec:V27-partial-legendre}
% -----------------------------------------------------------------------------

Let
\[
 L=L(y,v,w)
\]
be a real \(C^3\) first-order Lagrangian density on an open state set
\(\mathcal O\).  Aggregate the spatial slots and define
\begin{equation}
 G:=L_{vv},
 \qquad
 C:=L_{vw},
 \qquad
 D:=L_{ww}.
 \label{eq:V27-GCD}
\end{equation}
Thus \(C\) maps a spatial-jet variation to a velocity covector and
\(C^*=L_{wv}\).  We impose the following quantitative physical-fibre
hypothesis.

\begin{assumption}[Reduced Legendre-hyperbolicity window]
\label{ass:V27-legendre-window}
For every \((y,v,w)\in\mathcal O\), the admissible velocity fibre is convex
and
\begin{equation}
 g_* I\preceq G\preceq g^* I,
 \qquad
 -D\succeq d_* I,
 \qquad
 \|C\|\le c_*,
 \label{eq:V27-legendre-window}
\end{equation}
with \(g_*,g^*,d_*>0\), uniformly in the state and the cofinal level.
\end{assumption}

The sign of \(D\) is the familiar wave sign: for
\(L=\frac12|v|^2-\frac12|w|^2-V(y)\), one has \(G=I\) and \(D=-I\).
The mixed block \(C\) is allowed to be nonzero.

\begin{lemma}[Strong monotonicity and inverse velocity bound]
\label{lem:V27-velocity-inverse}
Under Assumption~\ref{ass:V27-legendre-window}, for fixed \((y,w)\) the map
\begin{equation}
 \mathcal L_{y,w}:v\longmapsto p=L_v(y,v,w)
 \label{eq:V27-legendre-map}
\end{equation}
is injective on each convex admissible velocity fibre, is a local
\(C^2\)-diffeomorphism, and its inverse on its image satisfies
\begin{equation}
 \|v(p_1)-v(p_0)\|
 \le g_*^{-1}\|p_1-p_0\|.
 \label{eq:V27-inverse-lipschitz}
\end{equation}
\end{lemma}

\begin{proof}
For \(\delta v=v_1-v_0\), the fundamental theorem of calculus and
\eqref{eq:V27-legendre-window} give
\[
 \langle L_v(v_1)-L_v(v_0),\delta v\rangle
 =\int_0^1
   \langle G(v_0+s\delta v)\delta v,\delta v\rangle\,ds
 \ge g_*\|\delta v\|^2.
\]
If the two momenta agree, then \(\delta v=0\), proving injectivity.
Cauchy--Schwarz in the same inequality gives
\(g_*\|\delta v\|\le\|p_1-p_0\|\), which is
\eqref{eq:V27-inverse-lipschitz}.  Since \(G\) is invertible at every
point, the inverse-function theorem supplies a local \(C^2\) inverse.
Injectivity makes the local inverses agree on overlaps, hence defines the
inverse on the image.
\end{proof}

Define the partial Hamiltonian
\begin{equation}
 H(y,p,w)
 :=\langle p,v(y,p,w)\rangle-L(y,v(y,p,w),w).
 \label{eq:V27-partial-H}
\end{equation}
The standard first derivative identities are
\begin{equation}
 H_p=v,
 \qquad
 H_w=-L_w,
 \qquad
 H_y=-L_y.
 \label{eq:V27-H-first}
\end{equation}
They are included to fix signs.  Their finite-dimensional velocity-only
counterparts are part of the imported result described above.

\begin{theorem}[Spatial-jet Hessian and quantitative convexity]
\label{thm:V27-Hessian}
Let \(U=(p,w)\).  Under Assumption~\ref{ass:V27-legendre-window},
\begin{equation}
 E:=H_{UU}
 =
 \begin{pmatrix}
  G^{-1} & -G^{-1}C\\
  -C^*G^{-1} & -D+C^*G^{-1}C
 \end{pmatrix}.
 \label{eq:V27-Hessian}
\end{equation}
Moreover, for \(\zeta=(\delta p,\delta w)\),
\begin{equation}
 \langle E\zeta,\zeta\rangle
 =
 \langle G^{-1}(\delta p-C\delta w),
                 \delta p-C\delta w\rangle
 +\langle -D\delta w,\delta w\rangle.
 \label{eq:V27-square-completion}
\end{equation}
Consequently
\begin{equation}
 E\succeq e_* I,
 \qquad
 e_*:=
 \frac{\min\{(g^*)^{-1},d_*\}}{(1+c_*)^2}>0.
 \label{eq:V27-energy-floor}
\end{equation}
\end{theorem}

\begin{proof}
Differentiate \(p=L_v(y,v,w)\) at fixed \(y\):
\[
 dp=G\,dv+C\,dw,
 \qquad
 dv=G^{-1}dp-G^{-1}C\,dw.
\]
Using \eqref{eq:V27-H-first} gives
\[
 H_{pp}=G^{-1},
 \qquad
 H_{pw}=-G^{-1}C,
 \qquad
 H_{ww}=-D+C^*G^{-1}C,
\]
which is \eqref{eq:V27-Hessian}.  Direct expansion proves
\eqref{eq:V27-square-completion}.

Let
\(S(\delta p,\delta w)=(\delta p-C\delta w,\delta w)\).
Then \(\|S^{-1}\|\le1+\|C\|\le1+c_*\).  Hence
\[
 \|S\zeta\|\ge(1+c_*)^{-1}\|\zeta\|.
\]
Since \(G^{-1}\succeq(g^*)^{-1}I\) and \(-D\succeq d_*I\),
\eqref{eq:V27-square-completion} yields
\[
 \langle E\zeta,\zeta\rangle
 \ge\min\{(g^*)^{-1},d_*\}\|S\zeta\|^2
 \ge e_*\|\zeta\|^2.
\]
\end{proof}

\begin{remark}[Why the mixed block does not consume the positive floor]
\label{rem:V27-mixed-block}
A direct estimate of the two off-diagonal terms in
\eqref{eq:V27-Hessian} would impose an unnecessary smallness condition on
\(C\).  The square completion shows that no such condition is required.
The price is only the shear-conditioning factor \((1+c_*)^2\).
\end{remark}

% -----------------------------------------------------------------------------
\subsection{Convex extension and the exact Friedrichs symmetrizer}
\label{subsec:V27-friedrichs}
% -----------------------------------------------------------------------------

For \(U=(p,w_1,\ldots,w_d)\), introduce the constant self-adjoint exchange
maps
\begin{equation}
 J^i(a,b_1,\ldots,b_d)
 :=(b_i,0,\ldots,0,\overset{i}{a},0,\ldots,0).
 \label{eq:V27-Ji}
\end{equation}
Set
\begin{equation}
 F^i(U;y):=-J^iD_UH(U;y).
 \label{eq:V27-flux}
\end{equation}
In components the first-order system is
\begin{align}
 \partial_t p-\partial_iH_{w_i}+H_y&=0,
 \label{eq:V27-p-evolution}\\
 \partial_t w_i-\partial_iH_p&=0,
 \label{eq:V27-w-evolution}\\
 \partial_t y-H_p&=0.
 \label{eq:V27-y-definition}
\end{align}
The first two rows have conservation form
\begin{equation}
 \partial_tU+\partial_iF^i(U;y)=(-H_y,0,\ldots,0).
 \label{eq:V27-conservation-form}
\end{equation}

\begin{theorem}[Partial Hamiltonian as a convex extension]
\label{thm:V27-friedrichs}
Under Assumption~\ref{ass:V27-legendre-window}, the principal
linearization of \eqref{eq:V27-conservation-form} is symmetrized by
\(E=H_{UU}\).  More precisely,
\begin{equation}
 A^0=E,
 \qquad
 A^i=E\,D_UF^i=-EJ^iE,
 \label{eq:V27-symmetric-matrices}
\end{equation}
where \(A^0\succeq e_*I\) and every \(A^i\) is self-adjoint.  The augmented
\((y,U)\) system is symmetric hyperbolic after assigning any uniformly
positive fibre metric to the algebraic-in-space row
\eqref{eq:V27-y-definition}.
\end{theorem}

\begin{proof}
Since \(J^i\) is constant and self-adjoint,
\[
 D_UF^i=-J^iH_{UU}=-J^iE.
\]
Left multiplication of the linearized \(U\) equation by \(E\) therefore
gives \eqref{eq:V27-symmetric-matrices}.  The product \(-EJ^iE\) is
self-adjoint, while Theorem~\ref{thm:V27-Hessian} gives the positive time
floor.  Dependence on \(y\), derivatives of the coefficients, and \(H_y\)
contribute only lower-order terms to the frozen principal calculation.
The \(y\) row contains \(\partial_t y\) and no spatial derivative of \(y\),
so its chosen positive metric adds a positive diagonal time block and a
zero diagonal spatial block.
\end{proof}

For a conormal \(n=(n_i)\), the resulting bosonic normal flux is
\begin{equation}
 A(n):=n_iA^i=-EJ(n)E,
 \qquad
 J(n):=n_iJ^i.
 \label{eq:V27-normal-flux}
\end{equation}
This supplies the concrete symmetrizer and flux input required by the V26
compatible trace pullback.  It does not choose a maximally nonnegative
boundary graph; that remains a separate boundary certificate.

% -----------------------------------------------------------------------------
\subsection{Hamilton--Pontryagin action and the differential readout}
\label{subsec:V27-HP}
% -----------------------------------------------------------------------------

Treat \(y,v,w,p,\lambda\) as independent and consider the local
Hamilton--Pontryagin action
\begin{equation}
 \mathcal S_{\rm HP}
 =\int_{I\times\Sigma}
 \left[
 L(y,v,w)
 +\langle p,\partial_ty-v\rangle
 +\sum_i\langle\lambda^i,\partial_i y-w_i\rangle
 \right].
 \label{eq:V27-HP-action}
\end{equation}
Ignoring the already-declared boundary term, its KKT residual rows are
\begin{align}
 C_t&:=\partial_ty-v,
 & C_i&:=\partial_i y-w_i,
 \label{eq:V27-CtCi}\\
 A_v&:=L_v-p,
 & A_{w_i}&:=L_{w_i}-\lambda^i,
 \label{eq:V27-AvAw}\\
 E_y&:=L_y-\partial_tp-\partial_i\lambda^i.
 \label{eq:V27-Ey}
\end{align}
The algebraic rows \(A_v=0\), \(A_w=0\) eliminate
\(p=L_v\) and \(\lambda^i=L_{w_i}\).  Define the causal residuals
\begin{align}
 E_p&:=\partial_tp+\partial_iL_{w_i}-L_y,
 \label{eq:V27-Ep}\\
 E^{\rm dyn}_{w_i}&:=\partial_tw_i-\partial_iv.
 \label{eq:V27-Ewdyn}
\end{align}

\begin{theorem}[Exact differential equation incidence]
\label{thm:V27-differential-incidence}
For arbitrary smooth independent fields, without imposing any equation,
the residual identities
\begin{align}
 E_p&=-E_y+\partial_iA_{w_i},
 \label{eq:V27-incidence-Ep}\\
 E^{\rm dyn}_{w_i}&=\partial_iC_t-\partial_tC_i
 \label{eq:V27-incidence-Ew}
\end{align}
hold identically.  In particular, the action-to-causal equation readout is
a local differential operator of order one.  Its principal symbol is
allowed to depend linearly on the covector, exactly as required by the V26
zeroth-order symbol obstruction.
\end{theorem}

\begin{proof}
From \eqref{eq:V27-Ey} and \eqref{eq:V27-AvAw},
\[
 -E_y+\partial_iA_{w_i}
 =\partial_tp+\partial_i\lambda^i-L_y
   +\partial_i(L_{w_i}-\lambda^i)
 =E_p.
\]
Commutation of coordinate derivatives and \eqref{eq:V27-CtCi} give
\[
 \partial_iC_t-\partial_tC_i
 =\partial_i\partial_ty-\partial_iv
  -\partial_t\partial_i y+\partial_tw_i
 =E^{\rm dyn}_{w_i}.
\]
No field equation was used.
\end{proof}

The theorem gives more than agreement of zero sets: it gives the missing
operator at nonzero residual.  To state its inverse correctly, one must
retain the defining and algebraic rows rather than discard them.

Let \(q\ge0\).  On the chosen cylinder or coordinate patch, choose constants
\(d_x,d_t\) such that
\begin{equation}
 \|\partial_i f_i\|_{H^q}\le d_x\|f\|_{H^{q+1}},
 \qquad
 \|(\partial_i f)_i\|_{H^q}\le d_x\|f\|_{H^{q+1}},
 \qquad
 \|\partial_t f\|_{H^q}\le d_t\|f\|_{H^{q+1}}.
 \label{eq:V27-derivative-bounds}
\end{equation}
Use the action target
\begin{equation}
 \mathcal A_q
 :=H^q(E_y)\oplus
 H^{q+1}(C_t)\oplus H^{q+1}(C_i)
 \oplus H^{q+1}(A_{w_i}).
 \label{eq:V27-action-target}
\end{equation}
Define
\begin{equation}
 \mathcal Q_q(E_y,C_t,C_i,A_w)
 :=(E_p,C_t,C_i,A_w,E_w^{\rm dyn}),
 \label{eq:V27-Qq}
\end{equation}
where the first and last entries are given by
\eqref{eq:V27-incidence-Ep}--\eqref{eq:V27-incidence-Ew}.  Its codomain is
equipped with the displayed product norm: \(E_p,E_w^{\rm dyn}\) lie in
\(H^q\), while the retained rows keep their \(H^{q+1}\) norms.

\begin{theorem}[Compatible target graph and bounded inverse]
\label{thm:V27-target-graph}
The range \(\mathcal C_q^{\rm comp}:=\operatorname{ran}\mathcal Q_q\) is
the closed compatible graph characterized by
\begin{equation}
 E^{\rm dyn}_{w_i}=\partial_iC_t-\partial_tC_i.
 \label{eq:V27-causal-compatibility}
\end{equation}
The map
\(\mathcal Q_q:\mathcal A_q\to\mathcal C_q^{\rm comp}\) is a bounded
bijection.  Its inverse reconstructs the discarded Euler row by
\begin{equation}
 E_y=-E_p+\partial_iA_{w_i}.
 \label{eq:V27-Qq-inverse}
\end{equation}
For the product norms just declared,
\begin{align}
 \|\mathcal Q_q\|
 &\le
 \left[
 \max\{2,1+2d_x^2,1+2d_t^2\}
 \right]^{1/2},
 \label{eq:V27-Q-upper}\\
 \|\mathcal Q_q^{-1}\|
 &\le
 \left[
 \max\{2,1+2d_x^2\}
 \right]^{1/2}.
 \label{eq:V27-Q-inverse-upper}
\end{align}
\end{theorem}

\begin{proof}
The compatibility identity follows from
Theorem~\ref{thm:V27-differential-incidence}.  Conversely, given a tuple
satisfying \eqref{eq:V27-causal-compatibility}, define \(E_y\) by
\eqref{eq:V27-Qq-inverse}.  Then
\eqref{eq:V27-incidence-Ep} holds and the retained entries are unchanged,
so the tuple lies in the range.  This also proves injectivity and gives the
inverse.  The range is closed because it is the kernel of the bounded map
\[
 (E_p,C_t,C_i,A_w,E_w)\longmapsto
 E_w-\partial_iC_t+\partial_tC_i.
\]

For the upper bound, use \((a+b)^2\le2a^2+2b^2\) and
\eqref{eq:V27-derivative-bounds}:
\begin{align*}
 \|E_p\|_{H^q}^2
 &\le2\|E_y\|_{H^q}^2+2d_x^2\|A_w\|_{H^{q+1}}^2,\\
 \|E_w\|_{H^q}^2
 &\le2d_x^2\|C_t\|_{H^{q+1}}^2
      +2d_t^2\|C_i\|_{H^{q+1}}^2.
\end{align*}
Adding the retained-row norms yields \eqref{eq:V27-Q-upper}.  Applying the
same estimate to \eqref{eq:V27-Qq-inverse}, and again adding the retained
rows, yields \eqref{eq:V27-Q-inverse-upper}.
\end{proof}

\begin{corollary}[No derivative loss after choosing the correct target]
\label{cor:V27-no-loss}
The first-order multiplier does not define a bounded automorphism of one
unshifted \(H^q\) product target.  It does define a bounded isomorphism from
the shifted action target \(\mathcal A_q\) onto the compatible causal graph
\(\mathcal C_q^{\rm comp}\).  Therefore the derivative in the equation
multiplier is a target-design issue, not an analytic loss in the causal
solution estimate.
\end{corollary}

\begin{proof}
The positive assertion is
Theorem~\ref{thm:V27-target-graph}.  For the negative assertion, the symbol
of \(\partial_iA_{w_i}\) grows like \(|\xi|\), so it is unbounded from an
unshifted \(H^q\) row to \(H^q\) unless that row is restricted by additional
regularity or a graph norm.
\end{proof}

The algebraic row \(A_v=L_v-p\) was omitted from \(\mathcal A_q\) only
because it is retained unchanged by the full multiplier.  Adding an identity
factor on \(A_v\), on gauge-fixing rows, and on the V25 curl/Bianchi rows
preserves all conclusions and makes the complete target map square after
the compatible range condition is imposed.

% -----------------------------------------------------------------------------
\subsection{Equivalence with the partial Hamiltonian equations}
\label{subsec:V27-equivalence}
% -----------------------------------------------------------------------------

\begin{theorem}[Root equivalence on the compatible first-order jet]
\label{thm:V27-root-equivalence}
Assume \(A_v=A_w=C_t=C_i=0\).  Then the Hamilton--Pontryagin Euler row
\(E_y=0\) is equivalent to the partial Hamiltonian system
\eqref{eq:V27-p-evolution}--\eqref{eq:V27-y-definition}, and the spatial
jet propagation row \eqref{eq:V27-w-evolution} follows identically.
Conversely, a solution of the partial Hamiltonian system with
\(p=L_v\), \(w_i=\partial_i y\), and \(\lambda^i=L_{w_i}\) is a stationary
point of \(\mathcal S_{\rm HP}\) modulo the declared boundary term.
\end{theorem}

\begin{proof}
The defining row \(C_t=0\) and \eqref{eq:V27-H-first} give
\(\partial_ty=v=H_p\).  Differentiating \(w_i=\partial_i y\) gives
\(\partial_tw_i=\partial_i v=\partial_iH_p\).  By
Theorem~\ref{thm:V27-differential-incidence}, \(A_w=0\) gives
\(E_p=-E_y\), so the two remaining evolution and Euler rows are equivalent.
This proves the forward implication.  The reverse implication substitutes
the stated algebraic definitions into
\eqref{eq:V27-CtCi}--\eqref{eq:V27-Ey}; the Hamilton equations set every
bulk variation to zero.  Boundary stationarity is exactly the boundary
condition already separated in V25--V26.
\end{proof}

% -----------------------------------------------------------------------------
\subsection{Kernel firewalls: fermions, gauge directions, ghosts, and gravity}
\label{subsec:V27-firewalls}
% -----------------------------------------------------------------------------

Theorem~\ref{thm:V27-Hessian} is deliberately not asserted for every field
in the Standard Model--Einstein action.

\begin{proposition}[Affine velocity dependence is a Legendre obstruction]
\label{prop:V27-affine-obstruction}
If a nonzero velocity subspace \(N\) satisfies
\begin{equation}
 L(y,v+\eta,w)=L(y,v,w)+\ell(y,w)[\eta]
 \quad\text{for every }\eta\in N\text{ and all nearby }v,
 \label{eq:V27-affine-velocity}
\end{equation}
where the linear coefficient \(\ell\) is independent of \(v\), then
\(G\eta=0\) for every \(\eta\in N\).  Hence no positive \(g_*\) exists on a
space containing \(N\), the Legendre map is not locally invertible there,
and the convex-extension construction above cannot supply that block's
symmetrizer.
\end{proposition}

\begin{proof}
Differentiate \eqref{eq:V27-affine-velocity} with respect to the base
velocity \(v\).  Its right-hand increment is independent of \(v\), hence
\(L_v(y,v+\eta,w)=L_v(y,v,w)\) for every \(\eta\in N\).  Differentiating at
\(\eta=0\) gives \(G\eta=0\).  This contradicts
\(G\succeq g_*I\) for every \(g_*>0\), and the inverse-function theorem
cannot apply.
\end{proof}

First-order Dirac kinetic terms are affine in the fermion velocity.  Their
causal block must therefore retain the independent first-order Dirac
symmetrizer and the V20 fermionic boundary domain.  A Yang--Mills potential
before gauge fixing contains a gauge or multiplier direction with a
degenerate velocity Hessian; it must be fixed or quotiented before
Assumption~\ref{ass:V27-legendre-window} is tested.  A BRST ghost block has
a signed or graded kinetic pairing and requires its own declared hyperbolic
energy; positivity cannot be inferred from the bosonic Legendre block.

For gravity, the attached manuscript records the ADM/DeWitt kinetic
coefficient and its Legendre reciprocity.  The DeWitt form has a conformal
indefinite direction before the Hamiltonian constraint and gauge reduction.
Thus a positive physical \(G\) for gravity must be demonstrated on the
actual reduced slice used by the continuum theorem.  Neither the algebraic
reciprocity of the coefficient nor invertibility of an indefinite KKT
Hessian proves \eqref{eq:V27-legendre-window}.

\begin{firewall}[V27 convexity has a declared physical domain]
\label{fire:V27-physical-domain}
The positive Hessian \(E\) is proved only on a bosonic field slice for which
\(G\succ0\) and \(-D\succ0\) have been certified.  Fermionic first-order
rows, gauge-null rows, ghosts, and unreduced gravitational directions are
separate blocks.  Silently placing any of them inside \(G\) would turn a
conditional theorem into a false one.
\end{firewall}

% -----------------------------------------------------------------------------
\subsection{Finite certificates for the Legendre and incidence rows}
\label{subsec:V27-finite}
% -----------------------------------------------------------------------------

At cofinal level \(n\), let \(\widehat G_n,\widehat C_n,\widehat D_n\) be
the assembled physical-fibre matrices and let certified operator errors
obey
\begin{equation}
 \|G_n-\widehat G_n\|\le\varepsilon_{G,n},
 \quad
 \|C_n-\widehat C_n\|\le\varepsilon_{C,n},
 \quad
 \|D_n-\widehat D_n\|\le\varepsilon_{D,n}.
 \label{eq:V27-GCD-errors}
\end{equation}
Define
\begin{align}
 g_{*,n}^{\rm cert}
 &:=\lambda_{\min}(\widehat G_n)-\varepsilon_{G,n},
 \label{eq:V27-gmin-cert}\\
 g_n^{*,\rm cert}
 &:=\|\widehat G_n\|+\varepsilon_{G,n},
 \label{eq:V27-gmax-cert}\\
 d_{*,n}^{\rm cert}
 &:=\lambda_{\min}(-\widehat D_n)-\varepsilon_{D,n},
 \label{eq:V27-dmin-cert}\\
 c_n^{*,\rm cert}
 &:=\|\widehat C_n\|+\varepsilon_{C,n}.
 \label{eq:V27-cmax-cert}
\end{align}

\begin{lemma}[Finite-to-physical convexity floor]
\label{lem:V27-finite-floor}
If \(g_{*,n}^{\rm cert}>0\) and \(d_{*,n}^{\rm cert}>0\), then the exact
level-\(n\) partial Hamiltonian Hessian satisfies
\begin{equation}
 E_n\succeq e_{*,n}^{\rm cert}I,
 \qquad
 e_{*,n}^{\rm cert}
 :=
 \frac{\min\{(g_n^{*,\rm cert})^{-1},d_{*,n}^{\rm cert}\}}
      {(1+c_n^{*,\rm cert})^2}.
 \label{eq:V27-finite-energy-floor}
\end{equation}
\end{lemma}

\begin{proof}
Weyl's inequality and \eqref{eq:V27-GCD-errors} give
\(G_n\succeq g_{*,n}^{\rm cert}I\),
\(G_n\preceq g_n^{*,\rm cert}I\),
\(-D_n\succeq d_{*,n}^{\rm cert}I\), and
\(\|C_n\|\le c_n^{*,\rm cert}\).  Apply
Theorem~\ref{thm:V27-Hessian}.
\end{proof}

The imported reciprocal residual tests only the \(pp\) corner of
\eqref{eq:V27-Hessian}.  A V27 certificate must additionally output
\begin{align}
 R_{H,n}
 &:=\widehat E_n-
 \begin{pmatrix}
  \widehat G_n^{-1}&-\widehat G_n^{-1}\widehat C_n\\
  -\widehat C_n^*\widehat G_n^{-1}&
  -\widehat D_n+\widehat C_n^*\widehat G_n^{-1}\widehat C_n
 \end{pmatrix},
 \label{eq:V27-Hessian-residual}\\
 R_{i,n}^{\rm sym}
 &:=\widehat A_n^i-(\widehat A_n^i)^*,
 \label{eq:V27-sym-residual}\\
 R_{p,n}^{\rm inc}
 &:=\widehat E_{p,n}+\widehat E_{y,n}
    -\partial_i\widehat A_{w_i,n},
 \label{eq:V27-p-inc-residual}\\
 R_{w_i,n}^{\rm inc}
 &:=\widehat E^{\rm dyn}_{w_i,n}
   -\partial_i\widehat C_{t,n}+\partial_t\widehat C_{i,n}.
 \label{eq:V27-w-inc-residual}
\end{align}

\begin{proposition}[Measured differential incidence]
\label{prop:V27-measured-incidence}
Suppose the finite incidence residuals in
\eqref{eq:V27-p-inc-residual}--\eqref{eq:V27-w-inc-residual} satisfy
\begin{equation}
 \|R_{p,n}^{\rm inc}\|_{H^q}
 +\sum_i\|R_{w_i,n}^{\rm inc}\|_{H^q}
 \le\varepsilon_{Q,n}.
 \label{eq:V27-incidence-error}
\end{equation}
Then the measured action--causal residual pair lies within
\(\varepsilon_{Q,n}\) of the exact incidence graph generated by
\(\mathcal Q_q\).  If the retained rows are
unchanged, the reconstructed action Euler row obeys
\begin{equation}
 \|\widehat E_{y,n}-
   (-\widehat E_{p,n}+\partial_i\widehat A_{w_i,n})\|_{H^q}
 \le\varepsilon_{Q,n}.
 \label{eq:V27-reconstruction-error}
\end{equation}
\end{proposition}

\begin{proof}
Subtract the exact graph identities
\eqref{eq:V27-incidence-Ep}--\eqref{eq:V27-incidence-Ew} from the measured
rows.  Holding the measured action rows fixed and correcting the two causal
rows by these defects produces an exact incidence pair, so the product-target
distance to the incidence graph is bounded by their norm.  The first defect is precisely the left-hand side of
\eqref{eq:V27-reconstruction-error}.
\end{proof}

For numerical work one may equivalently use the component declarations
\begin{equation}
 \begin{split}
 R_{pp,n}&=\widehat E_{pp,n}-\widehat G_n^{-1},\\
 R_{pw,n}&=\widehat E_{pw,n}+\widehat G_n^{-1}\widehat C_n,\\
 R_{ww,n}&=\widehat E_{ww,n}+\widehat D_n
 -\widehat C_n^*\widehat G_n^{-1}\widehat C_n.
 \end{split}
 \label{eq:V27-Hessian-residual-components}
\end{equation}

% -----------------------------------------------------------------------------
\subsection{Cofinal handoff to the full sector assembly}
\label{subsec:V27-cofinal}
% -----------------------------------------------------------------------------

\begin{theorem}[Legendre--incidence handoff]
\label{thm:V27-cofinal-handoff}
Consider a cofinal family of gauge-reduced bosonic action blocks and retain
the non-Legendre sectors as independent first-order blocks.  Assume:
\begin{enumerate}
 \item the certified quantities in
 \eqref{eq:V27-gmin-cert}--\eqref{eq:V27-cmax-cert} have uniform limits
 \[
  \inf_n g_{*,n}^{\rm cert}>0,
  \qquad
  \inf_n d_{*,n}^{\rm cert}>0,
  \qquad
  \sup_n(g_n^{*,\rm cert}+c_n^{*,\rm cert})<\infty;
 \]
 \item the Hessian, symmetry, and differential-incidence residuals converge
 to zero in their declared operator and shifted Sobolev norms;
 \item the fermion and ghost blocks possess their separately certified
 positive causal symmetrizers and compatible domains;
 \item the V23 off-diagonal sector coupling test, the V25 jet-complex and
 action-root tests, and the V26 compatible-boundary flux test pass with a
 positive common margin.
\end{enumerate}
Then the bosonic cofinal limit has a uniformly positive partial-Hamiltonian
symmetrizer; its action residual target is boundedly isomorphic to the
compatible causal target graph; and the assembled causal operator is the
differential equation-frame image of the common-action KKT system.  All
V26 causal inverse and boundary conclusions therefore apply with the
equation multiplier \(\mathcal Q_q\) constructed above.
\end{theorem}

\begin{proof}
Lemma~\ref{lem:V27-finite-floor} and the uniform hypotheses give
\(\inf_ne_{*,n}^{\rm cert}>0\).  Cofinal convergence of the blocks and
Theorem~\ref{thm:V27-Hessian} produce the positive limit symmetrizer;
Theorem~\ref{thm:V27-friedrichs} gives the symmetric principal matrices.
Theorems~\ref{thm:V27-differential-incidence} and
\ref{thm:V27-target-graph}, together with the vanishing incidence residual,
give the action-to-causal target isomorphism.  The separately symmetrized
first-order sectors and the bosonic block are assembled by the V23
off-diagonal theorem.  The V25 compatible jet and common-action hypotheses
identify the same physical roots, while the V26 pullback places the normal
flux on the compatible boundary fibre.  The V26 handoff theorem then gives
the stated inverse and boundary conclusions.
\end{proof}

The theorem is a handoff, not a physical closure.  Its first and third
hypotheses have not been populated by the attached finite tables.

% -----------------------------------------------------------------------------
\subsection{Executable V27 certificate card}
\label{subsec:V27-card}
% -----------------------------------------------------------------------------

At every cofinal level and sampled background, the certificate should
record
\begin{equation}
 \begin{split}
 \mathfrak m_{27,n}:=\min\{&
 g_{*,n}^{\rm cert},
 d_{*,n}^{\rm cert},
 e_{*,n}^{\rm cert},
 \lambda_{\min}(\widehat A^0_{f,n})-\varepsilon_{f,n},
 \lambda_{\min}(\widehat A^0_{gh,n})-\varepsilon_{gh,n},\\
 &\mathfrak m_{23,n}^{\rm coupling},
 \mathfrak m_{25,n}^{\rm jet/action},
 \mathfrak m_{26,n}^{\rm boundary/incidence}
 \},
 \end{split}
 \label{eq:V27-master-margin}
\end{equation}
together with the non-margin residual vector
\begin{equation}
 \mathfrak r_{27,n}:=
 \bigl(
 \|R_{H,n}\|,
 \max_i\|R_{i,n}^{\rm sym}\|,
 \|R_{p,n}^{\rm inc}\|_{H^q},
 \max_i\|R_{w_i,n}^{\rm inc}\|_{H^q},
 \|\mathcal Q_{q,n}\|,
 \|\mathcal Q_{q,n}^{-1}\|
 \bigr).
 \label{eq:V27-residual-vector}
\end{equation}
A pass requires \(\mathfrak m_{27,n}>0\), controlled multiplier norms, and
residual convergence under cofinal refinement.  The row must be evaluated
on the same physical projector and collar frame used by V24--V26.

Useful failure witnesses are
\begin{align}
 W_n^{\rm Legendre}
 &:=(\eta_n,\|\eta_n\|=1,
     \langle\widehat G_n\eta_n,\eta_n\rangle,
     \varepsilon_{G,n}),
 \label{eq:V27-Legendre-witness}\\
 W_n^{\rm spatial}
 &:=(\omega_n,\|\omega_n\|=1,
     -\langle\widehat D_n\omega_n,\omega_n\rangle,
     \varepsilon_{D,n}),
 \label{eq:V27-spatial-witness}\\
 W_n^{\rm incidence}
 &:=(R_{p,n}^{\rm inc},R_{w,n}^{\rm inc},
     \text{responsible KKT and defining rows}),
 \label{eq:V27-incidence-witness}\\
 W_n^{\rm sector}
 &:=(\text{sector name},\text{unremoved null direction},
     \text{declared replacement symmetrizer}).
 \label{eq:V27-sector-witness}
\end{align}
These distinguish a velocity kernel, the wrong spatial sign, a failed
differential readout, and an incorrectly classified first-order sector.

% -----------------------------------------------------------------------------
\subsection{V27 theorem ledger and next upstream row}
\label{subsec:V27-status}
% -----------------------------------------------------------------------------

\paragraph{New conclusions proved in V27.}
Theorem~\ref{thm:V27-Hessian} extends the imported velocity-only Legendre
result to the complete spatial-jet Hessian and gives an explicit positive
floor without a small mixed-block assumption.  Theorem~\ref{thm:V27-friedrichs}
constructs the exact convex-extension symmetrizer and normal flux.
Theorem~\ref{thm:V27-differential-incidence} derives the missing
action-to-causal multiplier from the Hamilton--Pontryagin action, and
Theorem~\ref{thm:V27-target-graph} proves that it is a bounded isomorphism
on the correctly shifted compatible target.  Theorem~\ref{thm:V27-root-equivalence}
proves equality of action and Hamiltonian roots on the compatible jet.
Proposition~\ref{prop:V27-affine-obstruction} supplies an explicit kernel
witness for sectors to which the construction cannot apply.

\paragraph{Imported conclusions not repeated.}
Finite velocity Legendre reciprocity, the ADM/DeWitt coefficient, the KKT
action and jet complex, the seven-sector coupling theorem, causal inverse,
and compatible-boundary pullback remain results of the attached sources
and V16--V26.  V27 uses them only at the stated handoff points.

\paragraph{Authoritative physical status.}
The attached manuscripts do not provide the gauge-reduced matrices
\(G,C,D\) for every bosonic Standard Model--Einstein background, their
uniform cofinal floors, the separate positive fermion and ghost time
matrices in the same frame, or the V27 differential-incidence residuals.
They also do not establish that the gravitational constraint quotient used
by the continuum construction removes every negative DeWitt direction.
Therefore \(\mathfrak m_{27,n}\) is not populated, the hypotheses of
Theorem~\ref{thm:V27-cofinal-handoff} are not physically verified, and the
full Standard Model--Einstein continuum is not claimed.

\paragraph{Next upstream row.}
The next bottleneck is the time matrix of the fully coupled system.  The
bosonic convex-extension matrices, Dirac time matrix, ghost block, and
gauge/constraint reduction must be placed in one field-dependent frame.
One must prove a positive floor for the complete coupled \(A^0\), identify
which Standard Model and Einstein couplings are genuinely lower order, and
bound every mixed principal block that remains.  The required result is a
block Schur or comparison theorem with finite-to-continuum erosion, followed
by a check that its characteristic cone is compatible with the V22 moving
boundary.  If that coupled floor fails, the calculation must move upstream
to the gauge and constraint projector rather than alter the causal inverse.

\begin{firewall}[V27 derives the multiplier but does not populate the model]
\label{fire:V27-status}
Hamilton--Pontryagin incidence removes the logical gap left by fitting an
unknown order-zero equation multiplier.  It does not certify the physical
Legendre signs, the reduced gravity slice, the first-order sector energies,
or the complete boundary flux.  Those are measured model rows and remain
open until the common finite assembly supplies them with uniform margins.
\end{firewall}

% END APPENDED VERSION V27 MATERIAL

% ============================================================================
% APPENDED VERSION V28 MATERIAL
% ============================================================================

\section{V28: Coupled Time-Matrix Positivity, Characteristic-Cone Control,
and Moving-Boundary Separation}
\label{sec:V28-coupled-cone}

\subsection{Purpose and nonduplication boundary}
\label{subsec:V28-purpose}

V27 constructs a positive partial-Hamiltonian symmetrizer for each certified
bosonic block and keeps the fermion, ghost, gauge, and reduced-gravity blocks
separate.  A direct sum of positive time matrices is positive, but the
physical first-order assembly need not be a direct sum: a common frame,
gauge reduction, or derivative interaction can introduce mixed principal
blocks.  This note proves the missing coupled time-matrix and characteristic
cone estimates, then transfers their spectral gap to the V22 moving boundary.

The result does not repeat the V23 seven-sector theorem.  V23 controls the
inverse of the complete nonlinear residual derivative in a transported graph
norm by a block norm majorant.  V28 instead treats the Hermitian principal
coefficient \(A^0\), the generalized eigenvalues of
\((A(n),A^0)\), and the inertia of the moving normal flux
\(A(n)-v_\partial A^0\).  The supported Schur complements in the attached
manuscripts concern positive Gram or source-ancestry forms and do not supply
this PDE time-matrix calculation.

Let
\begin{equation}
 \mathcal X=\bigoplus_{a=1}^N\mathcal X_a
 \label{eq:V28-sector-space}
\end{equation}
be the physical first-order fibre after the V25 compatible-jet restriction.
All statements are pointwise in the background and spacetime unless a
uniformity quantifier is written explicitly.  Put
\begin{equation}
 A^\mu=D^\mu+K^\mu,
 \qquad
 D^\mu:=\operatorname{diag}(A_1^\mu,\ldots,A_N^\mu),
 \label{eq:V28-principal-split}
\end{equation}
where \(K^\mu\) contains the mixed sector blocks in the common physical
frame and, by definition, \(K^\mu_{aa}=0\).  Any diagonal principal
correction is absorbed into \(A_a^\mu\) before the split.  The V27 bosonic Hessian, separately supplied Dirac and ghost
symmetrizers, and reduced gauge/gravity time blocks enter as the diagonal
matrices \(A_a^0\).

% -----------------------------------------------------------------------------
\subsection{Relative mixed-block theorem for the complete time matrix}
\label{subsec:V28-time-matrix}
% -----------------------------------------------------------------------------

Assume
\begin{equation}
 A_a^0=(A_a^0)^*,
 \qquad
 e_a I\preceq A_a^0\preceq E_a I,
 \qquad e_a>0.
 \label{eq:V28-sector-time-floors}
\end{equation}
Let \(D^0\) be their direct sum and suppose \(K^0=(K^0)^*\).  Define the
dimensionless relative time coupling
\begin{equation}
 R^0:=(D^0)^{-1/2}K^0(D^0)^{-1/2},
 \qquad
 \rho_0:=\|R^0\|.
 \label{eq:V28-relative-time}
\end{equation}

\begin{theorem}[Coupled positive time matrix]
\label{thm:V28-coupled-time}
If \(\rho_0<1\), then
\begin{equation}
 (1-\rho_0)D^0\preceq A^0\preceq(1+\rho_0)D^0.
 \label{eq:V28-relative-order}
\end{equation}
In particular, with \(e_{\min}:=\min_ae_a\) and
\(E_{\max}:=\max_aE_a\),
\begin{align}
 A^0&\succeq a_*I,
 &a_*&:=(1-\rho_0)e_{\min}>0,
 \label{eq:V28-time-floor}\\
 \|A^0\|&\le(1+\rho_0)E_{\max},
 &\|(A^0)^{-1}\|&\le a_*^{-1}.
 \label{eq:V28-time-upper-inverse}
\end{align}
\end{theorem}

\begin{proof}
For \(u\in\mathcal X\), put \(z=(D^0)^{1/2}u\).  Then
\[
 \langle A^0u,u\rangle
 =\langle(I+R^0)z,z\rangle.
\]
Self-adjointness and \(\|R^0\|=\rho_0\) imply
\((1-\rho_0)\|z\|^2\le
\langle(I+R^0)z,z\rangle\le(1+\rho_0)\|z\|^2\), proving
\eqref{eq:V28-relative-order}.  The sector floors and ceilings then give
\eqref{eq:V28-time-floor}--\eqref{eq:V28-time-upper-inverse}.
\end{proof}

There is also a certificate that avoids forming a square root.  For
\(a\ne b\), suppose
\begin{equation}
 \|K^0_{ab}\|\le\kappa_{ab}=\kappa_{ba},
 \qquad K^0_{ba}=(K^0_{ab})^*,
 \label{eq:V28-block-bounds}
\end{equation}
and form the real comparison matrix
\begin{equation}
 M^0_{aa}=e_a,
 \qquad
 M^0_{ab}=-\kappa_{ab}\quad(a\ne b).
 \label{eq:V28-comparison-matrix}
\end{equation}

\begin{lemma}[Sector comparison floor]
\label{lem:V28-comparison-floor}
If \(m_*:=\lambda_{\min}(M^0)>0\), then \(A^0\succeq m_*I\).
\end{lemma}

\begin{proof}
Write \(x_a=\|u_a\|\).  Equations
\eqref{eq:V28-sector-time-floors} and \eqref{eq:V28-block-bounds} give
\[
 \langle A^0u,u\rangle
 \ge\sum_ae_ax_a^2-2\sum_{a<b}\kappa_{ab}x_ax_b
 =x^TM^0x
 \ge m_*\sum_ax_a^2.
\]
\end{proof}

\begin{corollary}[Sharp two-block scalar majorant]
\label{cor:V28-two-block}
For
\begin{equation}
 A^0=
 \begin{pmatrix}B&X\\X^*&F\end{pmatrix},
 \qquad B\succeq bI,
 \quad F\succeq fI,
 \quad\|X\|\le x,
 \label{eq:V28-two-block-matrix}
\end{equation}
one has
\begin{equation}
 A^0\succeq
 \frac{b+f-\sqrt{(b-f)^2+4x^2}}{2}\,I.
 \label{eq:V28-two-block-floor}
\end{equation}
The displayed floor is positive exactly when \(x^2<bf\).
\end{corollary}

\begin{proof}
Apply Lemma~\ref{lem:V28-comparison-floor} to
\(\begin{psmallmatrix}b&-x\\-x&f\end{psmallmatrix}\) and compute its
smaller eigenvalue.  Positivity is equivalent to positive trace and
determinant; the trace is positive and the determinant is \(bf-x^2\).
\end{proof}

The relative criterion and the comparison criterion are alternatives.
The first is usually sharper when a reliable block square root is available;
the second supplies a transparent sector witness when it is not.

% -----------------------------------------------------------------------------
\subsection{Mixed spatial blocks and a certified characteristic cone}
\label{subsec:V28-cone}
% -----------------------------------------------------------------------------

Assume every \(A^i=(A^i)^*\).  For a unit spatial covector \(n\), write
\begin{equation}
 D(n):=n_iD^i,
 \qquad
 K(n):=n_iK^i,
 \qquad
 A(n):=D(n)+K(n).
 \label{eq:V28-normal-symbol}
\end{equation}
Normalize in the diagonal time metric:
\begin{equation}
 B(n):=(D^0)^{-1/2}D(n)(D^0)^{-1/2},
 \qquad
 R(n):=(D^0)^{-1/2}K(n)(D^0)^{-1/2}.
 \label{eq:V28-normalized-space}
\end{equation}
Define
\begin{equation}
 c_{\rm diag}:=\sup_{|n|=1}\|B(n)\|,
 \qquad
 \rho_x:=\sup_{|n|=1}\|R(n)\|,
 \qquad
 c_*:=\frac{c_{\rm diag}+\rho_x}{1-\rho_0}.
 \label{eq:V28-speed-bound}
\end{equation}

\begin{theorem}[Full-system speed and covector-cone bound]
\label{thm:V28-cone-bound}
Under the hypotheses of Theorem~\ref{thm:V28-coupled-time}, every generalized
eigenvalue \(\lambda_j(n)\) of
\begin{equation}
 A(n)r=\lambda A^0r
 \label{eq:V28-generalized-eigenproblem}
\end{equation}
is real and satisfies
\begin{equation}
 |\lambda_j(n)|\le c_*.
 \label{eq:V28-speed-interval}
\end{equation}
Consequently the principal symbol
\begin{equation}
 \sigma(\tau,\xi)=\tau A^0+A(\xi)
 \label{eq:V28-full-symbol}
\end{equation}
is invertible whenever \(|\tau|>c_*|\xi|\).  Hence the characteristic
covectors lie inside the certified cone
\begin{equation}
 |\tau|\le c_*|\xi|.
 \label{eq:V28-characteristic-cone}
\end{equation}
\end{theorem}

\begin{proof}
The pair \((A(n),A^0)\) is Hermitian definite, so its generalized
eigenvalues are real.  For a nonzero \(r\), let
\(z=(D^0)^{1/2}r\).  The Rayleigh quotient satisfies
\[
 \frac{|\langle A(n)r,r\rangle|}{\langle A^0r,r\rangle}
 \le
 \frac{(\|B(n)\|+\|R(n)\|)\|z\|^2}
      {(1-\rho_0)\|z\|^2}
 \le c_*.
\]
The min--max characterization proves \eqref{eq:V28-speed-interval}.
If \(\xi\ne0\), write \(n=\xi/|\xi|\).  Singularity of
\eqref{eq:V28-full-symbol} would give a generalized eigenvalue
\(-\tau/|\xi|\), contradicting \eqref{eq:V28-speed-interval} when
\(|\tau|>c_*|\xi|\).  The case \(\xi=0\) follows from invertibility of
\(A^0\).
\end{proof}

\begin{remark}[Direct sums and cone broadening]
\label{rem:V28-cone-broadening}
When \(K^\mu=0\), \(c_*\) reduces to the largest sector speed in the
diagonal metric.  Mixed time blocks enlarge the bound by
\((1-\rho_0)^{-1}\); mixed spatial blocks add \(\rho_x\).  This is a
rigorous enclosure, not a claim that either debit equals the true change of
the fastest characteristic.
\end{remark}

% -----------------------------------------------------------------------------
\subsection{Which interactions can enter the mixed principal bank}
\label{subsec:V28-interaction-audit}
% -----------------------------------------------------------------------------

Let \(\mathcal E_a(u,\partial u)\) denote the selected first-order equation
row.  The mixed principal incidence is the derivative
\begin{equation}
 \mathscr P^\mu_{ab}
 :=D_{\partial_\mu u_b}\mathcal E_a,
 \qquad a\ne b.
 \label{eq:V28-principal-incidence}
\end{equation}
It must be computed after the V25 jet enlargement and after the V27
differential equation multiplier.  Inspecting an unprolonged Lagrangian is
not enough, because elimination of a connection or a constraint variable
can move a derivative between blocks.

\begin{proposition}[Lower-order interaction criterion]
\label{prop:V28-lower-order}
An interaction term contributes no mixed principal block to the declared
first-order system if, after expressing it in the retained jet variables,
its linearization contains no derivative of the varied field.  Equivalently,
its contribution to every off-diagonal derivative
\(\mathscr P^\mu_{ab}\) in \eqref{eq:V28-principal-incidence} vanishes.
Such a term may change the lower-order Lipschitz and contraction constants,
but it changes neither \(K^0\), \(K(n)\), nor the cone bound
\eqref{eq:V28-speed-bound}.
\end{proposition}

\begin{proof}
The principal coefficient is, by definition, the coefficient multiplying
\(\partial_\mu\delta u_b\) in the linearized row \(\delta\mathcal E_a\).
The stated absence makes this coefficient zero.  The remaining terms
multiply \(\delta u_b\) without differentiation and therefore enter the
zeroth-order operator used in the V23 nonlinear estimate.
\end{proof}

For the minimal Standard Model written in retained first-order jets, the
Higgs polynomial potential and Yukawa monomials contain no derivatives and
therefore pass Proposition~\ref{prop:V28-lower-order}.  Expanding a minimal
gauge derivative gives a derivative of the matter field plus an algebraic
gauge-potential action; variation of the gauge potential in that algebraic
term contains no derivative of the gauge variation.  The corresponding
currents are lower order relative to a second-order wave block and become
algebraic after their first derivatives are retained as jet variables.

This observation is conditional on the declared first-order formulation.
Eliminating the spin connection can expose derivatives of a tetrad in the
Dirac row.  Nonminimal curvature--Higgs terms, Pauli terms, higher-derivative
effective operators, or derivative gauge conditions can also create genuine
mixed principal blocks.  Their \(\mathscr P^\mu_{ab}\) must be placed in
\(K^\mu\); naming an interaction ``minimal'' is not a substitute for the
incidence calculation.

\begin{firewall}[A coupling constant is not a principal-block norm]
\label{fire:V28-coupling-order}
The magnitude of a Yukawa or gauge coupling does not by itself measure
\(\rho_0\) or \(\rho_x\).  Derivative order and the chosen jet variables
decide whether that coupling enters the principal bank.  Lower-order
couplings are charged to V23; only the nonzero matrices in
\eqref{eq:V28-principal-incidence} are charged to V28.
\end{firewall}

% -----------------------------------------------------------------------------
\subsection{Moving-boundary normal flux and spectral separation}
\label{subsec:V28-moving-boundary}
% -----------------------------------------------------------------------------

Let a boundary face have outward unit spatial conormal \(n\) and signed
normal speed \(v_\partial\) in the chosen time chart.  Pullback to a fixed
collar changes the normal principal matrix to
\begin{equation}
 \mathcal F_\partial
 :=A(n)-v_\partial A^0.
 \label{eq:V28-moving-flux}
\end{equation}
Introduce the energy-normalized velocity matrix
\begin{equation}
 V(n):=(A^0)^{-1/2}A(n)(A^0)^{-1/2}.
 \label{eq:V28-velocity-matrix}
\end{equation}
It is self-adjoint, its eigenvalues are the
\(\lambda_j(n)\) of \eqref{eq:V28-generalized-eigenproblem}, and
\begin{equation}
 (A^0)^{-1/2}\mathcal F_\partial(A^0)^{-1/2}
 =V(n)-v_\partial I.
 \label{eq:V28-normalized-moving-flux}
\end{equation}

\begin{theorem}[Moving-boundary gap and constant inertia]
\label{thm:V28-moving-gap}
Define
\begin{equation}
 \gamma_\partial
 :=\min_j|\lambda_j(n)-v_\partial|.
 \label{eq:V28-boundary-gap}
\end{equation}
If \(\gamma_\partial>0\), the boundary is noncharacteristic, and the
positive, negative, and zero inertias of \(\mathcal F_\partial\) are
\begin{equation}
 \begin{split}
 n_+(\mathcal F_\partial)&=\#\{j:\lambda_j(n)>v_\partial\},\\
 n_-(\mathcal F_\partial)&=\#\{j:\lambda_j(n)<v_\partial\},\\
 n_0(\mathcal F_\partial)&=0.
 \end{split}
 \label{eq:V28-boundary-inertia}
\end{equation}
Along any continuous background and boundary path on which
\(\gamma_\partial\ge\gamma_*>0\), these inertias are constant and the
incoming/outgoing spectral projectors are continuous.  If the coefficients
are \(C^1\), the projectors are \(C^1\), so they satisfy the constant-rank
input of the V22 Kato transport.
\end{theorem}

\begin{proof}
Equation \eqref{eq:V28-normalized-moving-flux} is a congruence, so Sylvester
inertia is preserved.  Its eigenvalues are
\(\lambda_j(n)-v_\partial\), proving
\eqref{eq:V28-boundary-inertia}.  A uniform gap prevents an eigenvalue from
crossing zero, hence the inertias are constant.  The positive and negative
projectors are Riesz integrals on contours separated from zero.  Continuous,
respectively \(C^1\), dependence of the matrix transfers through those
integrals to the projectors.
\end{proof}

\begin{theorem}[Protected characteristic centre]
\label{thm:V28-protected-centre}
Let
\(\mathcal B_\partial:=V(n)-v_\partial I\) be the normalized moving flux.
Suppose there is a \(C^1\), fixed-rank orthogonal projector \(P_0\) such that
\begin{equation}
 \mathcal B_\partial P_0=0,
 \qquad
 \|\mathcal B_\partial z\|\ge\gamma_Z\|z\|
 \quad(z\in\operatorname{Ran}(I-P_0))
 \label{eq:V28-protected-centre}
\end{equation}
for a uniform \(\gamma_Z>0\).  Then
\(\operatorname{Ker}\mathcal F_\partial=(A^0)^{-1/2}
\operatorname{Ran}P_0\), the zero inertia has the constant value
\(\operatorname{rank}P_0\), and the positive, negative, and centre
projectors are \(C^1\).  Thus the V22 transport applies even though the
boundary is characteristic, provided the centre trace is retained in the
compatible constraint bundle.
\end{theorem}

\begin{proof}
Self-adjointness and \(\mathcal B_\partial P_0=0\) imply
\(P_0\mathcal B_\partial=0\), so \(\operatorname{Ran}P_0\) and its
orthogonal complement reduce \(\mathcal B_\partial\).  The lower bound on
the complement makes the compressed operator invertible and places its
spectrum outside \((-\gamma_Z,\gamma_Z)\).  Hence the kernel is exactly the
declared centre.  Riesz integrals on three mutually separated contours give
the \(C^1\) projectors.  Congruence by \((A^0)^{1/2}\) gives the stated
kernel of the unnormalized flux.
\end{proof}
\begin{corollary}[Supersonic sufficient condition]
\label{cor:V28-supersonic-boundary}
If \(v_\partial>c_*\), then \(\mathcal F_\partial\prec0\).  If
\(v_\partial<-c_*\), then \(\mathcal F_\partial\succ0\).  In either case
\begin{equation}
 \gamma_\partial\ge |v_\partial|-c_*.
 \label{eq:V28-supersonic-gap}
\end{equation}
\end{corollary}

\begin{proof}
Theorem~\ref{thm:V28-cone-bound} places every \(\lambda_j(n)\) in
\([-c_*,c_*]\).  Apply Theorem~\ref{thm:V28-moving-gap}.
\end{proof}

The corollary is only a sufficient test.  A timelike physical boundary may
have incoming and outgoing modes and still possess a positive gap; in that
case the V20 incoming graph, pulled back to the compatible trace fibre in
V26, is the relevant domain.

% -----------------------------------------------------------------------------
\subsection{Finite angular, frame, and boundary-gap certificates}
\label{subsec:V28-finite}
% -----------------------------------------------------------------------------

At cofinal level \(q\), let \(\widehat D_q^\mu,\widehat K_q^\mu\) be the
assembled matrices in the common V24 frame.  To avoid confusing the cofinal
index with the Sobolev exponent in V27, this note uses \(q\) only as a
finite-level label in the present subsection.

Let
\begin{equation}
 \|D_q^0-\widehat D_q^0\|\le\varepsilon_{D0,q},
 \qquad
 \|K_q^0-\widehat K_q^0\|\le\varepsilon_{K0,q},
 \label{eq:V28-time-errors}
\end{equation}
and let the diagonal sector floors already include the V27 and separate
first-order sector errors.  A direct full-matrix certificate is
\begin{equation}
 a_{*,q}^{\rm cert}
 :=\lambda_{\min}(\widehat A_q^0)
   -\varepsilon_{D0,q}-\varepsilon_{K0,q}.
 \label{eq:V28-direct-time-cert}
\end{equation}
Alternatively, inflate every measured off-diagonal block bound by its route
and tail error, build \(\widehat M_q^0\) as in
\eqref{eq:V28-comparison-matrix}, and use
\begin{equation}
 m_{*,q}^{\rm cert}:=\lambda_{\min}(\widehat M_q^0).
 \label{eq:V28-comparison-cert}
\end{equation}
Either positive number is a rigorous time floor when its input bounds are
upper certificates.

For the angular supremum, take a \(\delta_q\)-net
\(\mathcal N_q\subset S^{d-1}\).  Because the normalized spatial symbol is
linear in \(n\), any certified Lipschitz constant
\begin{equation}
 L_{B,q}\ge
 \sup_{|h|=1}\left\|
 (\widehat D_q^0)^{-1/2}\widehat A_q(h)
 (\widehat D_q^0)^{-1/2}
 \right\|
 \label{eq:V28-angular-Lipschitz}
\end{equation}
gives
\begin{equation}
 \widehat c_q^{\rm net}
 :=\max_{n\in\mathcal N_q}
 \left\|
 (\widehat A_q^0)^{-1/2}\widehat A_q(n)
 (\widehat A_q^0)^{-1/2}
 \right\|
 +L_{V,q}\delta_q+\varepsilon_{V,q},
 \label{eq:V28-net-speed}
\end{equation}
where \(L_{V,q}\) includes the certified change from the diagonal to the
full time metric and \(\varepsilon_{V,q}\) includes matrix, frame, and
cofinal-tail errors.  The notation separates what must be bounded; it does
not assert \(L_{V,q}=L_{B,q}\).

At a sampled boundary point, set
\begin{equation}
 \widehat\gamma_{\partial,q}
 :=\min_j|\lambda_j(\widehat V_q(n))-\widehat v_{\partial,q}|,
 \qquad
 \gamma_{\partial,q}^{\rm cert}
 :=\widehat\gamma_{\partial,q}
   -\varepsilon_{V,q}-\varepsilon_{v,q}
   -L_{\partial,q}\delta_{\partial,q}.
 \label{eq:V28-boundary-gap-cert}
\end{equation}
Here \(\delta_{\partial,q}\) is the joint background, time, and boundary
net radius and \(L_{\partial,q}\) its certified modulus.

For the protected-centre route, let \(\widehat P_{0,q}\) be the orthogonal
projector onto the declared constraint centre, put
\(\widehat{\mathcal B}_{\partial,q}
 =\widehat V_q(n)-\widehat v_{\partial,q}I\), and define
\begin{align}
 \widehat\gamma_{Z,q}
 &:=s_{\min}\!\left(
  (I-\widehat P_{0,q})\widehat{\mathcal B}_{\partial,q}
  (I-\widehat P_{0,q})
  \big|_{\operatorname{Ran}(I-\widehat P_{0,q})}
 \right),
 \label{eq:V28-centre-gap-measured}\\
 \gamma_{Z,q}^{\rm cert}
 &:=\widehat\gamma_{Z,q}
   -\varepsilon_{V,q}-\varepsilon_{v,q}
   -L_{\partial,q}\delta_{\partial,q},
 \label{eq:V28-centre-gap-cert}\\
 r_{Z,q}
 &:=\|\widehat{\mathcal B}_{\partial,q}\widehat P_{0,q}\|
   +\|\widehat P_{0,q}^2-\widehat P_{0,q}\|
   +\|\widehat P_{0,q}-(\widehat P_{0,q})^*\|.
 \label{eq:V28-centre-residual}
\end{align}
An exact symbolic incidence
\(\mathcal B_{\partial,q}P_{0,q}=0\) together with
\(\gamma_{Z,q}^{\rm cert}>0\) invokes
Theorem~\ref{thm:V28-protected-centre} at that level.  A merely small
\(r_{Z,q}\) certifies only a near-zero cluster.  It gives an exact centre in
the cofinal limit only when \(r_{Z,q}\to0\), the projectors converge with
constant rank, and the operator tails vanish.
\begin{lemma}[Finite-to-continuum inertia certificate]
\label{lem:V28-inertia-cert}
If \(a_{*,q}^{\rm cert}>0\) and
\(\gamma_{\partial,q}^{\rm cert}>0\), then the exact time matrix is
positive and the exact moving normal flux has the same inertia as the
measured flux throughout the certified net cell.  If, more strongly, the
total normalized flux error \(\varepsilon\) is less than half the measured
gap \(\widehat\gamma_{\partial,q}\), then the positive and negative spectral
projectors obey the Davis--Kahan bound
\begin{equation}
 \|\Pi_{\partial,q}-\widehat\Pi_{\partial,q}\|
 \le\frac{2\varepsilon}{\widehat\gamma_{\partial,q}}.
 \label{eq:V28-projector-error}
\end{equation}
\end{lemma}

\begin{proof}
Weyl's inequality applied to \eqref{eq:V28-direct-time-cert} proves
positivity.  Apply it again to the normalized measured and exact fluxes.
The error debit in \eqref{eq:V28-boundary-gap-cert} leaves every eigenvalue
on the same side of zero throughout the cell, so inertia is unchanged.
For the projector estimate, use the Riesz-projector or sin-theta theorem
with the two spectral clusters separated by the interval about zero; the
assumption \(\varepsilon<\widehat\gamma/2\) leaves separation at least
\(\widehat\gamma/2\), giving the displayed conservative constant.
\end{proof}

The self-adjoint defects
\begin{equation}
 \Delta_q^\mu:=\widehat A_q^\mu-(\widehat A_q^\mu)^*
 \label{eq:V28-selfadjoint-defect}
\end{equation}
must also be reported.  Replacing a measured matrix by its Hermitian part is
useful for diagnosis but is not a physical proof: a nonvanishing continuum
skew principal block is a derivative-order error that cannot be absorbed by
the zeroth-order V23 estimate.  The required cofinal row is
\begin{equation}
 \max_\mu\|\Delta_q^\mu\|\longrightarrow0
 \label{eq:V28-selfadjoint-limit}
\end{equation}
after all route, frame, and tail errors are included.

% -----------------------------------------------------------------------------
\subsection{Cofinal coupled-cone handoff}
\label{subsec:V28-cofinal}
% -----------------------------------------------------------------------------

\begin{theorem}[Coupled principal and moving-domain handoff]
\label{thm:V28-cofinal-handoff}
Suppose a cofinal Standard Model--Einstein first-order assembly satisfies:
\begin{enumerate}
 \item the V27 sector time floors are uniform and every independent
 first-order sector has its declared positive symmetrizer;
 \item either \(\sup_q\rho_{0,q}<1\) in the relative test or
 \(\inf_qm_{*,q}^{\rm cert}>0\) in the comparison test;
 \item \(\sup_q(c_{{\rm diag},q}+\rho_{x,q})<\infty\), all principal
 self-adjoint defects vanish, the coefficient bank is \(C^1\) in the V24
 frame, and the frame and high-mode tails are summable;
 \item on the transported boundary, either
 \(\inf_q\gamma_{\partial,q}^{\rm cert}>0\), or the protected-centre route
 has constant projector rank, \(\inf_q\gamma_{Z,q}^{\rm cert}>0\),
 \(r_{Z,q}\to0\), and summably convergent centre projectors;
 \item the V23 nonlinear, V25 compatible-jet, V26 compatible-flux, and V27
 action-to-causal incidence margins remain positive after the V28 debits.
\end{enumerate}
Then the cofinal principal system has a uniformly positive time matrix,
real characteristics contained in a uniform cone, and a moving normal flux
with constant inertia.  Its incoming projector has the regularity and
constant rank required by the V22 domain transport.  The V22 causal
evolution family and the V23 nonlinear inverse therefore apply to the same
coupled physical system and the same compatible boundary domain.
\end{theorem}

\begin{proof}
Theorem~\ref{thm:V28-coupled-time} or
Lemma~\ref{lem:V28-comparison-floor} gives the uniform time floor.  Cofinal
convergence and \eqref{eq:V28-selfadjoint-limit} yield self-adjoint limiting
principal matrices.  Theorem~\ref{thm:V28-cone-bound} gives the uniform cone.
Lemma~\ref{lem:V28-inertia-cert} and summability stabilize the boundary
inertia and projector; Theorem~\ref{thm:V28-moving-gap} supplies its
constant rank and regularity.  These are precisely the principal and domain
inputs missing from the V22 handoff.  The final item identifies that causal
system with the nonlinear common-action assembly of V23 and V25--V27.
\end{proof}

% -----------------------------------------------------------------------------
\subsection{Executable V28 certificate and status}
\label{subsec:V28-card-status}
% -----------------------------------------------------------------------------

The finite card must store the diagonal sector floors and ceilings, every
mixed principal block \(K_{ab}^\mu\), its adjoint partner, and the frame in
which they were compared.  Its principal margin is
\begin{equation}
 \mathfrak m_{28,q}:=
 \min\left\{
  a_{*,q}^{\rm cert},
  m_{*,q}^{\rm cert},
  1-\rho_{0,q}^{\rm cert},
  \gamma_{{\rm route},q}^{\rm cert},
  \mathfrak m_{27,q}
 \right\},
 \label{eq:V28-master-margin}
\end{equation}
where \(\gamma_{{\rm route},q}^{\rm cert}\) is the noncharacteristic gap
\(\gamma_{\partial,q}^{\rm cert}\) or the protected-centre gap
\(\gamma_{Z,q}^{\rm cert}\).  Only the chosen time-floor route is included:
if the direct or
comparison route is used, the unused alternative is omitted rather than set
to zero.  The non-margin output is
\begin{equation}
 \mathfrak r_{28,q}:=
 \left(
  \max_\mu\|\Delta_q^\mu\|,
  \widehat c_q^{\rm net},
  L_{V,q}\delta_q,
  \varepsilon_{V,q},
  \|\Pi_{\partial,q}-\widehat\Pi_{\partial,q}\|,
  \text{frame and principal tails}
 \right).
 \label{eq:V28-residual-card}
\end{equation}

Failure witnesses are
\begin{align}
 W_q^{A^0}
 &:=(u_q,\|u_q\|=1,
     \langle\widehat A_q^0u_q,u_q\rangle,
     \text{responsible sector blocks}),
 \label{eq:V28-time-witness}\\
 W_q^{\rm cone}
 &:=(n_q,r_q,\widehat\lambda_q,
     \text{angular and tail debit}),
 \label{eq:V28-cone-witness}\\
 W_q^{\partial}
 &:=(t_q,y_q,n_q,\widehat v_{\partial,q},
     \widehat\lambda_{j,q},
     \gamma_{\partial,q}^{\rm cert}),
 \label{eq:V28-boundary-witness}\\
 W_q^{\rm order}
 &:=(a,b,\mu,\mathscr P^\mu_{ab},
     \text{retained or eliminated jet variable}).
 \label{eq:V28-order-witness}
\end{align}

\paragraph{New conclusions proved in V28.}
Theorem~\ref{thm:V28-coupled-time} proves positivity of the complete time
matrix from a dimensionless relative coupling, while
Lemma~\ref{lem:V28-comparison-floor} and
Corollary~\ref{cor:V28-two-block} provide blockwise alternatives.
Theorem~\ref{thm:V28-cone-bound} converts the mixed principal bank into a
uniform characteristic-cone enclosure.  Theorem~\ref{thm:V28-moving-gap}
identifies the exact moving-boundary flux, its inertia, and the gap required
by the V22 projector transport.  Lemma~\ref{lem:V28-inertia-cert} makes that
gap finite and stable under frame, net, and tail errors.

\paragraph{Imported conclusions not repeated.}
The V23 full nonlinear inverse, V22 Kato domain transport, V26 compatible
flux pullback, V27 partial-Hamiltonian symmetrizer, and the source-Gram Schur
theorems in the attachments are used only as inputs.  V28 supplies their
missing coupled principal and characteristic incidence.

\paragraph{Authoritative physical status.}
The attachments do not print one common-frame bank
\(\{A_a^\mu,K_{ab}^\mu\}\) for the gauge-reduced Einstein, Yang--Mills,
Higgs, fermion, ghost, source, and constraint sectors.  They do not classify
every interaction by \eqref{eq:V28-principal-incidence}, populate
\(\rho_0,\rho_x\), or give an angular and moving-boundary spectral-gap card.
Thus \(\mathfrak m_{28,q}\) is not physically evaluable from the supplied
tables, Theorem~\ref{thm:V28-cofinal-handoff} remains conditional, and the
full Standard Model--Einstein continuum is not claimed.

\paragraph{Next upstream row.}
The next note must choose and write an explicit gauge-fixed minimal
Einstein--Standard-Model principal system in one variable list.  It should
derive, rather than name, the harmonic-Einstein, Lorenz--Yang--Mills,
Higgs, Dirac, and ghost frozen symbols; classify the spin-connection and
stress-current terms after first-order prolongation; and output the actual
off-diagonal matrices in
\eqref{eq:V28-principal-incidence}.  Any surviving negative time direction
must be traced to a gauge or constraint projector.  Only after that
calculation can the V28 margins be populated and compared with the physical
boundary speed.

\begin{firewall}[V28 controls a supplied principal bank]
\label{fire:V28-status}
The relative and comparison theorems rigorously propagate sector data; they
do not manufacture that data.  A block omitted from the principal-incidence
audit is not bounded by setting its norm to zero.  The full-continuum claim
therefore waits for the explicit gauge-fixed common-frame symbol and its
cofinal certificate.
\end{firewall}

% END APPENDED VERSION V28 MATERIAL

% ============================================================================
% APPENDED VERSION V29 MATERIAL
% ============================================================================

\section{V29: Explicit Minimal Einstein--Standard-Model Frozen Symbols
and the Einstein--Dirac Principal Obstruction}
\label{sec:V29-explicit-symbols}

\subsection{Purpose, imported Einstein row, and conventions}
\label{subsec:V29-purpose}

V28 reduces the coupled causal problem to a common-frame principal bank
\(\{A_a^\mu,K_{ab}^\mu\}\).  V18 has already derived the de Donder-reduced
Einstein symbol, and V19 has already proved how an occurring common-action
response can be tested against a declared symbol.  Those results are
imported.  V29 writes the remaining minimal classical sector symbols in one
frozen Lorentz frame and audits the mixed derivative terms created by the
first-order reduction.

The outcome is partly positive and partly obstructive.  Harmonic Einstein,
Lorenz Yang--Mills, Higgs, and the commuting Faddeev--Popov representative
reduce to copies of one wave symbol.  The Dirac block has the same metric
null cone and a positive first-order symmetrizer.  Polynomial, Yukawa, and
minimal gauge interactions are lower order in retained jet variables.  On a
nonzero spinor background, however, the Dirac stress tensor contains a
derivative of the spinor variation.  In a naive first-order Einstein--Dirac
system this is a genuine upper triangular principal block.  V29 computes its
location and proves the simultaneous symmetrizer equation which must remove
it.  It is not set to zero by terminology.

Use signature \((-+++ )\).  Let \(T\) be a future unit timelike vector and
\(e_1,e_2,e_3\) an orthonormal spatial frame at a frozen background point.
Choose normal coordinates and compatible bundle frames at that point; their
connection coefficients and derivatives of coefficients are lower order in
the frozen symbol.  Write a coordinate time vector as
\begin{equation}
 \partial_t=N T+\beta^ie_i,
 \qquad N>0,
 \label{eq:V29-lapse-shift}
\end{equation}
with lapse \(N\) and shift \(\beta\).  Clifford multiplication obeys
\begin{equation}
 \gamma(\zeta)\gamma(\eta)+\gamma(\eta)\gamma(\zeta)
 =2g^{-1}(\zeta,\eta)I.
 \label{eq:V29-Clifford}
\end{equation}

% -----------------------------------------------------------------------------
\subsection{One wave first-order block for all normally hyperbolic sectors}
\label{subsec:V29-wave-block}
% -----------------------------------------------------------------------------

Let \(u\) take values in a finite-rank real or complex bundle and suppose
its frozen second-order equation has principal part
\begin{equation}
 -\nabla_T^2u+\sum_{i=1}^3\nabla_{e_i}^2u.
 \label{eq:V29-wave-principal}
\end{equation}
Introduce
\begin{equation}
 \pi:=\nabla_Tu,
 \qquad
 w_i:=\nabla_{e_i}u.
 \label{eq:V29-wave-jets}
\end{equation}
Modulo lower-order curvature and commutator rows, the dynamic variables
\(U_W=(\pi,w_1,w_2,w_3)\) satisfy
\begin{equation}
 \nabla_TU_W+\sum_iA_W^i\nabla_{e_i}U_W=\operatorname{l.o.t.},
 \qquad
 A_W(n)=
 \begin{pmatrix}
  0&-n^T\otimes I\\
  -n\otimes I&0
 \end{pmatrix}.
 \label{eq:V29-wave-first-order}
\end{equation}

\begin{lemma}[Wave symmetrizer and frozen speeds]
\label{lem:V29-wave-speeds}
The matrices \(A_W^i\) are self-adjoint in the direct-sum fibre metric and
the time matrix is \(A_W^0=I\).  For \(|n|=1\),
\begin{equation}
 \operatorname{spec}A_W(n)=\{-1,+1,0\},
 \label{eq:V29-wave-spectrum}
\end{equation}
where the zero eigenspace consists of the two spatial-jet components
orthogonal to \(n\).  The defining row \(\nabla_Tu-\pi=0\) adds an algebraic
in-space variable and does not change the nonzero speeds.
\end{lemma}

\begin{proof}
Self-adjointness is visible in \eqref{eq:V29-wave-first-order}.  Decompose
\(w=nw_n+w_\perp\).  The matrix acts on \((\pi,w_n)\) by
\(\begin{psmallmatrix}0&-1\\-1&0\end{psmallmatrix}\) and annihilates
\(w_\perp\), proving \eqref{eq:V29-wave-spectrum}.  The defining row has
only a time derivative at principal order.
\end{proof}

\begin{proposition}[Minimal normally hyperbolic sector bank]
\label{prop:V29-wave-sectors}
After the indicated gauge reductions, the following frozen principal parts
are copies of \eqref{eq:V29-wave-principal}:
\begin{enumerate}
 \item the de Donder-reduced Einstein perturbation on
 \(S^2T^*M\), up to the invertible scalar row factor already computed in
 V18;
 \item the background-Lorenz reduced Yang--Mills perturbation on
 \(T^*M\otimes\mathfrak g_{\rm SM}\);
 \item the gauge-covariant Higgs equation on its realified representation
 bundle;
 \item the commuting coefficient representative of the Faddeev--Popov
 ghost equation on \(\mathfrak g_{\rm SM}\).
\end{enumerate}
Thus their first-order reductions have the positive time matrix and speeds
of Lemma~\ref{lem:V29-wave-speeds} in the frozen orthonormal frame.
\end{proposition}

\begin{proof}
The Einstein item is Theorem~\ref{thm:V18-reduced-symbol}.  For Yang--Mills,
linearize \(D^\mu F_{\mu\nu}\) at a background connection.  Its two-derivative
part on a perturbation \(a\) is
\[
 \nabla^\mu\nabla_\mu a_\nu
 -\nabla_\nu\nabla^\mu a_\mu.
\]
Adding the background-Lorenz gauge row
\(\nabla_\nu\nabla^\mu a_\mu\) cancels the mixed term.  Curvature and
background gauge commutators have at most one derivative of \(a\).
The Higgs kinetic Euler row is \(D^\mu D_\mu\phi\); gauge potentials and
bundle connections affect only its lower-order frozen coefficients.  The
linearization of the Lorenz gauge condition along an infinitesimal gauge
parameter is the Faddeev--Popov operator \(D^\mu D_\mu\), with the same
principal part.  Applying the jet reduction
\eqref{eq:V29-wave-jets} proves the last statement.
\end{proof}

The ghost conclusion is analytic.  Grassmann parity, BRST pairing, and a
Krein or signed physical interpretation are not converted into a positive
Hilbert energy by replacing ghost coefficients with commuting test vectors.
That physical-domain row remains the separate V27 ghost certificate.

% -----------------------------------------------------------------------------
\subsection{Dirac block and its light cone}
\label{subsec:V29-Dirac}
% -----------------------------------------------------------------------------

In the adapted tetrad, multiply the Dirac equation by its standard positive
time-current matrix.  With a convention-dependent harmless overall sign,
its principal part becomes
\begin{equation}
 \nabla_T\psi+\alpha^i\nabla_{e_i}\psi,
 \qquad
 \alpha^i:=\gamma(T)\gamma(e_i),
 \label{eq:V29-Dirac-first-order}
\end{equation}
in the positive spinor inner product induced by \(T\).

\begin{lemma}[Dirac symmetrizer and characteristic speeds]
\label{lem:V29-Dirac-speeds}
The matrices \(\alpha^i\) are self-adjoint in the \(T\)-spinor product and,
for every unit spatial \(n\),
\begin{equation}
 \alpha(n)^2=I,
 \qquad
 \operatorname{spec}\alpha(n)\subset\{-1,+1\}.
 \label{eq:V29-alpha-square}
\end{equation}
Consequently \(A_D^0=I\), the Dirac speeds are \(\pm1\), and its
characteristic covectors satisfy \(g^{-1}(\xi,\xi)=0\).  Fermion masses,
the Higgs Yukawa mass matrix, and algebraic gauge-potential terms do not
change this principal symbol.
\end{lemma}

\begin{proof}
The defining adjoint property of the conserved Dirac current makes
\(\alpha^i\) self-adjoint.  Since \(n\perp T\), the Clifford relations give
anticommutation of \(\gamma(T)\) and \(\gamma(n)\), together with
\(\gamma(T)^2=-I\) and \(\gamma(n)^2=I\).  Hence
\[
 [\gamma(T)\gamma(n)]^2
 =-\gamma(T)^2\gamma(n)^2=I.
\]
Self-adjointness then gives spectrum in \(\{-1,+1\}\).  Equivalently,
the unmultiplied Dirac symbol obeys
\(\gamma(\xi)^2=g^{-1}(\xi,\xi)I\), so it is invertible off the metric null
cone.  The listed mass and potential terms contain no derivative of
\(\psi\).  The same conclusion holds on each chiral Weyl block after
the standard realification used for the causal energy.
\end{proof}

\begin{corollary}[Lapse--shift speed bank]
\label{cor:V29-lapse-shift-speeds}
In the coordinates \eqref{eq:V29-lapse-shift}, the wave and Dirac dynamic
blocks can be written with time matrix \(I\).  In a spatial unit direction
\(n\), their coordinate speeds belong to
\begin{equation}
 \{-\beta\cdot n-N,-\beta\cdot n,
   -\beta\cdot n+N\}
 \label{eq:V29-coordinate-speeds}
\end{equation}
for a wave jet and to
\(\{-\beta\cdot n-N,-\beta\cdot n+N\}\) for a Dirac spinor.  Thus the
block-diagonal bank has
\begin{equation}
 c_{\rm diag}\le |\beta|+N.
 \label{eq:V29-coordinate-cone-bound}
\end{equation}
\end{corollary}

\begin{proof}
Substitute \(T=N^{-1}(\partial_t-\beta^i\partial_i)\) in the orthonormal
systems and multiply by \(N\).  Their spatial matrices become
\(-\beta^iI+NA^i\).  Apply
Lemma~\ref{lem:V29-wave-speeds} and
Lemma~\ref{lem:V29-Dirac-speeds}.
\end{proof}

% -----------------------------------------------------------------------------
\subsection{Minimal interactions and the one derivative-stress exception}
\label{subsec:V29-interactions}
% -----------------------------------------------------------------------------

Retain the wave jets for the metric, gauge potential, and Higgs field.
Then the following terms are algebraic in the first-order state: the Higgs
potential, Yukawa monomials, gauge commutators in \(F\), covariant-potential
pieces in \(D\Phi\) and \(D\psi\), Yang--Mills and Higgs currents, the
Yang--Mills and Higgs stresses, and connection coefficients expressed
through the retained metric or tetrad jet.  Their linearizations contribute
to the V23 lower-order matrix, not to the V28 mixed principal bank.

The symmetric Dirac stress is different.  Its derivative part has the form
\begin{equation}
 T^{D}_{\mu\nu}\big|_{\nabla\psi}
 =\frac{i}{4}\left(
  \overline\psi\,\gamma_{(\mu}\nabla_{\nu)}\psi
  -(\nabla_{(\mu}\overline\psi)\gamma_{\nu)}\psi
 \right)
 \label{eq:V29-Dirac-stress}
\end{equation}
up to the fixed real-action convention.  At a background spinor
\(\psi_0\), the derivative part of its linearization is
\begin{equation}
 \sigma_1(\delta T^D)(\xi)\dot\psi
 =c_D\left(
  \overline\psi_0\gamma_{(\mu}\xi_{\nu)}\dot\psi
  -\overline{\dot\psi}\,\xi_{(\mu}\gamma_{\nu)}\psi_0
 \right),
 \qquad c_D\ne0,
 \label{eq:V29-Dirac-stress-symbol}
\end{equation}
where the precise grouping of conjugate realified components depends on the
chosen real spinor coordinates.  What is invariant is that the expression
is linear in \(\xi\), linear in \(\dot\psi\), and proportional to
\(\psi_0\); it is generically nonzero when \(\psi_0\ne0\).

\begin{proposition}[Einstein--Dirac triangular principal block]
\label{prop:V29-triangular-block}
Use the wave first-order reduction for the reduced Einstein row and the
unprolonged first-order Dirac row.  Absorb the variation of the spin
connection into the retained metric/tetrad jet.  After using the Dirac row
to eliminate a time derivative from \eqref{eq:V29-Dirac-stress-symbol}, the
frozen dynamic principal matrices have the form
\begin{equation}
 P^0=
 \begin{pmatrix}I&0\\0&I\end{pmatrix},
 \qquad
 P^i=
 \begin{pmatrix}
  G^i&B^i(\psi_0)\\
  0&D^i
 \end{pmatrix}.
 \label{eq:V29-triangular-principal}
\end{equation}
Here \(G^i\) is the Einstein wave-jet matrix, \(D^i\) is the Dirac matrix,
and \(B^i(\psi_0)\) is the spatial part of the stress symbol.  It vanishes
at \(\psi_0=0\) and is not forced to vanish at a general spinor background.
All other minimal bosonic interactions described above give zero
off-diagonal principal blocks in this retained-jet formulation.
\end{proposition}

\begin{proof}
The reduced Einstein equation is second order in the metric.  Its
first-order momentum equation contains the stress source.  Since
\eqref{eq:V29-Dirac-stress-symbol} differentiates \(\dot\psi\), that source
is first order in the Dirac variable and therefore has the same first-order
status as the Dirac equation after reduction.  Eliminating its \(T\)-part
with \eqref{eq:V29-Dirac-first-order} leaves matrices multiplying spatial
derivatives of \(\dot\psi\), denoted \(B^i(\psi_0)\).  The Dirac variation
of the spin connection contains the first metric/tetrad jet but no
derivative of that retained jet, so the lower-left principal block is zero.
The remaining interaction classification follows from the retained-jet
criterion of Proposition~\ref{prop:V28-lower-order}.
\end{proof}

Equation \eqref{eq:V29-Dirac-stress-symbol} is displayed only to identify
derivative order and background dependence.  Its component signs must be
regenerated from the exact real spinor and adjoint convention used in the
finite common action; V29 does not use those signs in a positivity claim.

% -----------------------------------------------------------------------------
\subsection{The exact simultaneous symmetrizer equation}
\label{subsec:V29-symmetrizer}
% -----------------------------------------------------------------------------

The block triangular matrices in
\eqref{eq:V29-triangular-principal} are not self-adjoint in the direct-sum
metric unless \(B^i=0\).  A field-dependent positive symmetrizer may still
exist.  Write
\begin{equation}
 H=
 \begin{pmatrix}H_G&X\\X^*&H_D\end{pmatrix},
 \qquad H=H^*>0.
 \label{eq:V29-general-symmetrizer}
\end{equation}

\begin{theorem}[Einstein--Dirac simultaneous symmetrizer test]
\label{thm:V29-sylvester}
Assume \(H_GG^i=(H_GG^i)^*\) and
\(H_DD^i=(H_DD^i)^*\) for the uncoupled blocks.  Then
\(HP^i=(HP^i)^*\) for every spatial \(i\) if and only if
\begin{align}
 (G^i)^*X-XD^i&=H_GB^i,
 \label{eq:V29-Sylvester}\
 X^*B^i+H_DD^i&=(X^*B^i+H_DD^i)^*
 \label{eq:V29-bottom-symmetry}
\end{align}
for every \(i\).  The same single matrix \(X\) must solve all directions.
If \(P^0=I\), the symmetrized time matrix is \(H\); hence positivity is
equivalent to
\begin{equation}
 H_G>0,
 \qquad
 H_D-X^*H_G^{-1}X>0.
 \label{eq:V29-H-Schur}
\end{equation}
In the normalized case \(H_G=H_D=I\), the sufficient and necessary Schur
condition is \(\|X\|<1\).
\end{theorem}

\begin{proof}
Direct multiplication gives
\[
 HP^i=
 \begin{pmatrix}
  H_GG^i&H_GB^i+XD^i\\
  X^*G^i&X^*B^i+H_DD^i
 \end{pmatrix}.
\]
The upper-left block is self-adjoint by hypothesis.  Equality of the
upper-right block with the adjoint of the lower-left block is exactly
\eqref{eq:V29-Sylvester}; self-adjointness of the lower-right block is
\eqref{eq:V29-bottom-symmetry}.  These conditions are also sufficient.
Since \(P^0=I\), the time coefficient after multiplication is \(H\).
The positive-block Schur-complement criterion gives
\eqref{eq:V29-H-Schur}.  With identity diagonal blocks it reduces to
\(I-X^*X>0\), equivalent to \(\|X\|<1\).
\end{proof}

\begin{corollary}[Sharp defect of the direct-sum symmetrizer]
\label{cor:V29-direct-sum-defect}
For \(H=\operatorname{diag}(H_G,H_D)\), the spatial self-adjointness defect
is
\begin{equation}
 HP^i-(HP^i)^*
 =
 \begin{pmatrix}
  0&H_GB^i\\
  -(B^i)^*H_G&0
 \end{pmatrix},
 \qquad
 \|HP^i-(HP^i)^*\|=\|H_GB^i\|.
 \label{eq:V29-triangular-defect}
\end{equation}
Thus the direct-sum V28 bank is symmetric at a general background only if
the stress block vanishes in the chosen formulation.
\end{corollary}

\begin{proof}
Set \(X=0\) in the preceding block product.  A matrix
\(\begin{psmallmatrix}0&C\\-C^*&0\end{psmallmatrix}\) has norm \(\|C\|\),
as follows by squaring its absolute value.
\end{proof}

\begin{firewall}[Squaring Dirac does not solve the mixed first-order row]
\label{fire:V29-Dirac-square}
The identity \(\gamma(\xi)^2=g^{-1}(\xi,\xi)I\) proves the Dirac
characteristic cone.  It does not remove the derivative Dirac stress from
the Einstein momentum equation, produce the common matrix \(X\) in
\eqref{eq:V29-Sylvester}, or preserve the first-order fermion boundary
domain.  A squared Dirac replacement would therefore bypass the V20 domain
and is not used here.
\end{firewall}

% -----------------------------------------------------------------------------
\subsection{Two legitimate upstream closures of the stress block}
\label{subsec:V29-upstream-closures}
% -----------------------------------------------------------------------------

\begin{proposition}[Zero-spinor-background closure]
\label{prop:V29-zero-spinor}
On a background with \(\psi_0=0\), the derivative stress block
\(B^i(\psi_0)\) vanishes.  For the minimal gauge-fixed classical system,
the frozen principal bank described in V29 is block diagonal, has positive
time matrix \(I\), and has coordinate cone bounded by
\eqref{eq:V29-coordinate-cone-bound}.  Its wave-jet zero-speed modes must
still be assigned to the V25 constraint complex and the V28 protected
boundary centre.
\end{proposition}

\begin{proof}
Every term in \eqref{eq:V29-Dirac-stress-symbol} contains \(\psi_0\) or its
adjoint, so it vanishes.  Proposition~\ref{prop:V29-triangular-block} removes
the remaining mixed principal blocks.  The sector time matrices and speeds
are those of Lemmas~\ref{lem:V29-wave-speeds} and
\ref{lem:V29-Dirac-speeds}, followed by
Corollary~\ref{cor:V29-lapse-shift-speeds}.  The zero wave-jet eigenvalues
are explicit in \eqref{eq:V29-wave-spectrum}.
\end{proof}

At a general spinor background, there are two admissible routes.
\begin{enumerate}
 \item Keep the unprolonged Dirac variable and solve the simultaneous system
 \eqref{eq:V29-Sylvester}--\eqref{eq:V29-bottom-symmetry} for one positive
 \(H\).  The solution must be incident with the V27 equation multiplier and
 the V26 compatible boundary flux.
 \item Introduce a spinor jet \(\chi_\mu=\nabla_\mu\psi\), so the stress is
 algebraic, and prolong the Dirac equation together with its curl and
 defining constraints.  This route must prove a positive symmetrizer for the
 enlarged system, propagation of the spinor-jet constraints, equivalence of
 roots, and preservation of the V20 first-order boundary trace.  Merely
 adding \(\chi\) as an independent value is insufficient.
\end{enumerate}
These routes are alternatives.  Combining the favorable statement from one
with the domain statement from the other would be a target mismatch.

% -----------------------------------------------------------------------------
\subsection{Finite common-action certificate for the explicit bank}
\label{subsec:V29-finite}
% -----------------------------------------------------------------------------

At cofinal level \(q\), run the V19 packet or impulse-symbol acquisition on
the field ordering
\begin{equation}
 (g,\partial g;A,F;\Phi,D\Phi;\psi;c,\partial c;
   \text{constraints and multipliers}).
 \label{eq:V29-field-order}
\end{equation}
After the V24 frame and V27 equation multiplier are applied, retain every
principal column and form the target residuals
\begin{align}
 R_{W,q}^\mu
 &:=\widehat A_{W,q}^\mu-A_{W,\rm tar}^\mu,
 \label{eq:V29-wave-residual}\\
 R_{D,q}^\mu
 &:=\widehat A_{D,q}^\mu-A_{D,\rm tar}^\mu,
 \label{eq:V29-Dirac-residual}\\
 R_{B,q}^\mu
 &:=\widehat B_q^\mu-B_{\rm stress}^\mu(\widehat\psi_{0,q}),
 \label{eq:V29-stress-residual}\\
 R_{0B,q}
 &:=\widehat B_q^0,
 \label{eq:V29-time-elimination-residual}\\
 R_{{\rm low},q}^\mu
 &:=\max_{(a,b)\in\mathcal L}
   \|\widehat{\mathscr P}_{ab,q}^\mu\|,
 \label{eq:V29-lower-order-residual}
\end{align}
where \(\mathcal L\) lists the interactions declared lower order above.
The time-elimination residual checks that use of the Dirac equation really
removed the stress time derivative in the same target frame.

For the direct-sum route, the exact obstruction is
\begin{equation}
 d_{{\rm tri},q}
 :=\max_i\|H_{G,q}\widehat B_q^i\|.
 \label{eq:V29-triangular-certificate}
\end{equation}
For the general symmetrizer route, fit or derive one \(X_q\) and retain
\begin{align}
 r_{{\rm Syl},q}
 &:=\max_i\|
  (\widehat G_q^i)^*X_q-X_q\widehat D_q^i
  -H_{G,q}\widehat B_q^i
 \|,
 \label{eq:V29-Sylvester-residual}\\
 r_{{\rm bot},q}
 &:=\max_i\|
 X_q^*\widehat B_q^i+H_{D,q}\widehat D_q^i
 -(X_q^*\widehat B_q^i+H_{D,q}\widehat D_q^i)^*
 \|,
 \label{eq:V29-bottom-residual}\\
 h_{*,q}^{\rm cert}
 &:=\lambda_{\min}\!left(
 H_{D,q}-X_q^*H_{G,q}^{-1}X_q
 \right)-\varepsilon_{H,q}.
 \label{eq:V29-H-floor-cert}
\end{align}
Here \(\varepsilon_{H,q}\) includes the inverse, frame, packet, and
high-mode errors.  A direction-dependent \(X_q(n)\) does not pass: the same
matrix must make every \(r_{{\rm Syl},q}^i\) and
\(r_{{\rm bot},q}^i\) small.

\begin{theorem}[Cofinal explicit-symbol handoff]
\label{thm:V29-cofinal-handoff}
Assume the V19 occurrence residuals for
\eqref{eq:V29-wave-residual}--\eqref{eq:V29-lower-order-residual} vanish
uniformly, the lapse has a positive floor, and the frame and tails are
summable.  Assume also that the uncoupled \(H_{G,q},H_{D,q}\) have
uniform positive floors and ceilings and converge in the common frame.
Suppose either:
\begin{enumerate}
 \item \(\psi_{0,q}\to0\) in a norm controlling
 \(B^i(\psi_{0,q})\), with \(d_{{\rm tri},q}\to0\); or
 \item \(X_q\to X\), \(\inf_qh_{*,q}^{\rm cert}>0\), and
 \(r_{{\rm Syl},q}+r_{{\rm bot},q}\to0\).
\end{enumerate}
Then the cofinal minimal gauge-fixed principal bank is symmetric hyperbolic,
has the metric characteristic cone, and supplies the populated analytic
sector input of V28.  The V28 moving-boundary handoff applies if its separate
normal-flux gap or protected-centre certificate also passes.
\end{theorem}

\begin{proof}
Proposition~\ref{prop:V29-wave-sectors} and
Lemma~\ref{lem:V29-Dirac-speeds} identify all diagonal limiting symbols.
The vanishing lower-order residual prevents a falsely omitted mixed
principal block.  On the first route,
Proposition~\ref{prop:V29-zero-spinor} gives the direct-sum symmetrizer.  On
the second, Theorem~\ref{thm:V29-sylvester}, convergence of the residuals,
and the positive Schur floor give a limiting positive simultaneous
symmetrizer.  The triangular determinant factors as the product of the Einstein-wave and
Dirac determinants, so the coupled characteristic set is their common metric
null cone.  V28 then supplies the quantitative common-frame cone and boundary
margins.  The final boundary assertion is exactly the remaining V28
hypothesis.
\end{proof}

% -----------------------------------------------------------------------------
\subsection{V29 ledger and next upstream row}
\label{subsec:V29-status}
% -----------------------------------------------------------------------------

\paragraph{New conclusions proved in V29.}
Proposition~\ref{prop:V29-wave-sectors} derives one explicit wave-jet bank
for the reduced Einstein, Yang--Mills, Higgs, and commuting ghost equations.
Lemma~\ref{lem:V29-Dirac-speeds} derives the positive Dirac time block and
its metric null cone.  Proposition~\ref{prop:V29-triangular-block} identifies
the derivative Dirac stress as the exceptional minimal mixed principal row.
Theorem~\ref{thm:V29-sylvester} gives necessary and sufficient equations for
one positive simultaneous symmetrizer, and
Corollary~\ref{cor:V29-direct-sum-defect} gives its sharp direct-sum failure
witness.

\paragraph{Imported conclusions not repeated.}
The reduced Einstein symbol, common-action packet extraction, first-order
Dirac boundary space, compatible jet complex, differential equation
multiplier, and coupled-cone theorem remain the results of V18--V28.  V29
uses them to place the standard frozen sector formulas in the occurring
target frame.

\paragraph{Authoritative physical status.}
The attached files do not contain the common-frame packet columns in
\eqref{eq:V29-wave-residual}--\eqref{eq:V29-lower-order-residual}, a
realified finite Dirac-stress block \(B_q^i\), or a common positive solution
\(X_q\) of all equations \eqref{eq:V29-Sylvester}.  They also do not select
the zero-spinor-background route as the physical continuum history.
Therefore the alternatives of Theorem~\ref{thm:V29-cofinal-handoff} are not
populated, and the full Standard Model--Einstein continuum is not claimed.

\paragraph{Next upstream row.}
The obstruction is now localized.  The next note should derive the
Einstein--Dirac block from a real tetrad/Palatini Hamilton--Pontryagin action
with independent spin connection.  It must decide whether the common-action
multiplier generates a matching lower-left block, produces a positive
solution of the simultaneous Sylvester equations, or requires a compatible
spinor-jet prolongation.  This calculation must include the V20 boundary
form; otherwise a bulk symmetrizer could destroy the admitted fermion trace.

\begin{firewall}[V29 identifies the difficult physical block]
\label{fire:V29-status}
Agreement of every diagonal sector with the metric light cone does not prove
that the coupled Einstein--Dirac first-order system is symmetric hyperbolic.
The derivative stress block is zero on a zero-spinor background and generally
nonzero otherwise.  The full-continuum route must certify its simultaneous
symmetrizer or replace the variable system with a proved equivalent one.
\end{firewall}

% END APPENDED VERSION V29 MATERIAL


% =============================================================================
% START APPENDED VERSION V30 MATERIAL
% =============================================================================

\section{V30: Covariant spinor-jet prolongation, algebraic Cartan torsion,
and the differentiated fermion domain}
\label{sec:V30}

\paragraph{Purpose and correction of scope.}
V29 isolated a nonzero upper-right principal block when the Einstein wave
equation and the unprolonged Dirac equation were forced into one uniformly
first-order triangular matrix.  That computation is correct for that
variable list, but it is not a no-go theorem for Einstein--Dirac evolution.
The Einstein equation is second order in the metric and only first order in
the spinor, whereas the Dirac equation is first order in the spinor.
Consequently its natural principal object is graded.  A compatible
first-order prolongation promotes the covariant derivative of the spinor to
a jet value.  The stress then becomes algebraic in the enlarged state.  The
price is an exact defining constraint, its propagation equation, and a
curvature-incidence row.  This note proves those statements and treats the
fermion boundary trace without replacing the first-order Dirac domain.

The connection-stationarity and torsion alternative proved in the attached
grand tensor manuscript are imported rather than repeated.  In particular,
on a nondegenerate coframe the Palatini--Holst connection equation has the
unique spin-sourced Einstein--Cartan torsion, with a quantitative inverse
bound.  V30 uses that result to shift the V16 Levi--Civita incidence target
to a spin-dependent Einstein--Cartan target.

\paragraph{External consistency check, not an imported proof.}
M\"uller and Nowaczyk construct a metric-natural spinor jet bundle and prove
that the first-order prolongation of the reduced
Einstein--Dirac--Maxwell operator is symmetric hyperbolic; see
\url{https://arxiv.org/abs/1504.01034}.  Gro{\ss}e and Murro prove
well-posedness for the Dirac operator on globally hyperbolic manifolds with
timelike boundary after imposing an admissible first-order boundary
condition; see \url{https://arxiv.org/abs/1806.06544}.  These references
confirm that the prolongation and domain questions have legitimate
solutions in the ordinary continuum theory.  Every implication used below
is nevertheless proved from the displayed identities and the earlier
V20--V28 certificates.

% -----------------------------------------------------------------------------
\subsection{The graded differential order and the V29 defect}
\label{subsec:V30-graded-order}
% -----------------------------------------------------------------------------

Write the gauge-fixed Einstein--Cartan--Dirac equations schematically as
\begin{align}
 \mathcal G(g,\partial g,\partial^2g)
   &=\mathcal T(g,\partial g,\psi,\nabla\psi),
 \label{eq:V30-Einstein-order}\\
 \mathcal E(g,\partial g,\psi,\nabla\psi)
   &:=\gamma^\mu\nabla^\omega_\mu\psi-M(\Phi)\psi=0.
 \label{eq:V30-Dirac-row}
\end{align}
Here \(\omega\) is either the independent Cartan connection or its
algebraically eliminated Einstein--Cartan value, and \(M(\Phi)\) includes
the mass and Yukawa endomorphisms.  The Yang--Mills representation is
included in \(\nabla^\omega\).

\begin{definition}[Graded principal order]
\label{def:V30-graded-principal}
Assign weight two to the metric equation and weight one to the Dirac
equation.  In the metric column retain order two in \(g\); in the spinor
column retain order one in \(\psi\).  A mixed entry is graded principal only
when it reaches the retained order of its row and column.  Equivalently, the
state is prolonged to
\begin{equation}
 U_{\rm gr}
 =
 \bigl(g,k_\mu=\partial_\mu g;\,
       \psi,\chi_a=\nabla_a\psi;\,
       A,F;\Phi,D\Phi;\omega,R,T\bigr),
 \label{eq:V30-graded-state}
\end{equation}
and a derivative already displayed in \(U_{\rm gr}\) is evaluated as a
value, subject to its defining and curl constraints.
\end{definition}

\begin{proposition}[Location of the V29 triangular block]
\label{prop:V30-defect-location}
In the variables of Definition~\ref{def:V30-graded-principal}, the
derivative stress block \(B^i\) of
\eqref{eq:V29-triangular-principal} is not a principal derivative block.  It is
the derivative of the algebraic map
\begin{equation}
 \mathcal T_{\rm D}
 =
 \mathcal T_{\rm D}(g,k,\psi,\chi)
 \label{eq:V30-stress-algebraic}
\end{equation}
with respect to the jet value \(\chi\).  The V29 block reappears only after
one eliminates \(\chi_i\) by substituting \(\chi_i=\nabla_i\psi\) before
forming a uniform first-order symbol.
\end{proposition}

\begin{proof}
The Dirac stress has the local form
\[
 \mathcal T_{{\rm D},\mu\nu}
 =
 \frac12\operatorname{Re}
 \left\langle\psi,
 \gamma_{(\mu}\nabla_{\nu)}^\omega\psi\right\rangle
 +\mathcal A_{\mu\nu}(g,\psi),
\]
up to convention-dependent algebraic multiples of the Dirac equation.
Replacing every covariant derivative by the corresponding component of
\(\chi\) gives \eqref{eq:V30-stress-algebraic}.  Hence
\(\partial\mathcal T_{\rm D}/\partial(\partial_i\chi)=0\).  If instead
\(\chi_i\) is eliminated, the chain rule gives
\[
 \frac{\partial\mathcal T_{\rm D}}{\partial(\partial_i\psi)}
 =
 \frac{\partial\mathcal T_{\rm D}}{\partial\chi_i},
\]
which is exactly the frozen block computed in
\eqref{eq:V29-Dirac-stress-symbol}.  Thus V29 detected an elimination
artifact of a genuine derivative coupling, not a false calculation.
\end{proof}

\begin{firewall}[Grading is not permission to discard equations]
\label{fire:V30-grading}
Calling \(\chi\) a value does not make it independent physical data.  A
valid prolongation must impose \(\chi-\nabla\psi=0\), propagate that
constraint, reproduce the curvature commutator, and preserve the original
Dirac boundary trace.  Without those rows,
\eqref{eq:V30-stress-algebraic} is only a change of notation.
\end{firewall}

% -----------------------------------------------------------------------------
\subsection{Algebraic Einstein--Cartan elimination}
\label{subsec:V30-Cartan}
% -----------------------------------------------------------------------------

Let \(e\) be a nondegenerate coframe.  Decompose a metric-compatible
connection as
\begin{equation}
 \omega=\omega^{\rm LC}(e)+K,
 \qquad
 T^I=K^I{}_J\wedge e^J.
 \label{eq:V30-contorsion}
\end{equation}
Denote by \(\mathscr J_e\) the linear map \(K\mapsto K\wedge e\), and by
\(\mathscr C_e\) the Cartan map used in the attached
connection-stationarity theorem.  The imported theorem says that, after the
invertible Palatini--Holst map has been applied, the connection Euler row is
an algebraic equation
\begin{equation}
 \mathscr C_e(T)+\Sigma(e,\psi)=r_\omega
 \label{eq:V30-Cartan-row}
\end{equation}
and that \(\mathscr C_e\) is injective with a certified inverse floor on the
retained nondegenerate-coframe class.  The sign of \(\Sigma\) is absorbed
into its definition.

\begin{proposition}[Spin-dependent connection target]
\label{prop:V30-EC-target}
Assume
\[
 \|\mathscr C_e^{-1}\|\le c_C^{-1},
 \qquad
 \|\mathscr J_e^{-1}\|\le c_J^{-1},
 \qquad c_C,c_J>0.
\]
Define
\begin{align}
 T_{\rm EC}(e,\psi)
 &:=-\mathscr C_e^{-1}\Sigma(e,\psi),
 \label{eq:V30-EC-torsion}\\
 K_{\rm EC}(e,\psi)
 &:=\mathscr J_e^{-1}T_{\rm EC}(e,\psi),
 \label{eq:V30-EC-contorsion}\\
 \omega_{\rm EC}(e,\psi)
 &:=\omega^{\rm LC}(e)+K_{\rm EC}(e,\psi).
 \label{eq:V30-EC-connection}
\end{align}
Then every approximate connection satisfying
\eqref{eq:V30-Cartan-row} obeys
\begin{align}
 \|T-T_{\rm EC}\|
 &\le c_C^{-1}\|r_\omega\|,
 \label{eq:V30-torsion-error}\\
 \|\omega-\omega_{\rm EC}\|
 &\le(c_Cc_J)^{-1}\|r_\omega\|.
 \label{eq:V30-connection-error}
\end{align}
On the spinless branch \(\Sigma=0\), this reduces to the V16
Levi--Civita target.
\end{proposition}

\begin{proof}
Subtract
\(\mathscr C_e(T_{\rm EC})+\Sigma=0\) from
\eqref{eq:V30-Cartan-row}.  Applying the inverse bound gives
\eqref{eq:V30-torsion-error}.  Equation
\eqref{eq:V30-contorsion} and the inverse bound for \(\mathscr J_e\) then
give \eqref{eq:V30-connection-error}.  If \(\Sigma=0\), injectivity gives
\(T_{\rm EC}=K_{\rm EC}=0\).
\end{proof}

\begin{corollary}[Cartan torsion is lower order in the prolonged Dirac row]
\label{cor:V30-torsion-lower-order}
After substituting \(\omega_{\rm EC}\),
\begin{equation}
 \nabla^{\omega_{\rm EC}}_\mu\psi
 =
 \nabla^{\rm LC}_\mu\psi
 +\frac14K_{{\rm EC},\mu IJ}\gamma^{IJ}\psi
 \label{eq:V30-spin-connection-shift}
\end{equation}
has the same principal spinor symbol as the Levi--Civita Dirac operator.
If the spin current is algebraic in \(\psi\), the added term is algebraic
and generates only the familiar local nonlinear fermion interaction after
elimination of torsion.
\end{corollary}

\begin{proof}
The difference of two metric-compatible spin connections acts on spinors
by the zeroth-order endomorphism displayed in
\eqref{eq:V30-spin-connection-shift}.  Proposition~\ref{prop:V30-EC-target}
makes \(K_{\rm EC}\) an algebraic function of \(e\) and \(\psi\).  It
therefore contributes no covector \(\xi\) to the principal symbol.
\end{proof}

\begin{firewall}[The independent connection is not declared hyperbolic]
\label{fire:V30-connection}
The Palatini connection Euler row used here is algebraic after the Cartan
map is inverted.  V30 does not invent a propagating connection degree of
freedom.  Lorentzian evolution remains the reduced metric-wave system of
V18 and V29; \(\omega\), \(T\), and \(R\) are retained incidence variables
whose defining and Bianchi rows belong to the V25 compatible complex.
\end{firewall}

% -----------------------------------------------------------------------------
\subsection{Exact covariant spinor-jet incidence}
\label{subsec:V30-spinor-jet}
% -----------------------------------------------------------------------------

Let \(\nabla\) denote the total spin, gauge, and tensor connection after
the Cartan choice has been made, and let
\begin{equation}
 \mathcal D:=\gamma^\mu\nabla_\mu-M.
 \label{eq:V30-D}
\end{equation}
Metric and Clifford compatibility give \(\nabla_a\gamma^\mu=0\).  For a
spinor-valued covariant Hessian, write
\begin{equation}
 [\nabla_a,\nabla_b]_{\rm H}\psi
 :=
 (\nabla_a\nabla\psi)_b-(\nabla_b\nabla\psi)_a.
 \label{eq:V30-Hessian-commutator}
\end{equation}
This convention includes the connection on the derivative covector index.
Introduce independent vector-spinors \(\chi_a\) and the residuals
\begin{align}
 E&:=\mathcal D\psi,
 \label{eq:V30-E}\\
 C_a&:=\chi_a-\nabla_a\psi,
 \label{eq:V30-C}\\
 E_a&:=
 \mathcal D\chi_a
 +\gamma^\mu[\nabla_a,\nabla_\mu]_{\rm H}\psi
 -(\nabla_aM)\psi.
 \label{eq:V30-Ea}
\end{align}
The covariant derivative in \(\mathcal D\chi_a\) acts on both the spinor
and covector indices.  The constraint \(C_a\) is the component version of
the V25 fermion-jet defect \(C_\Xi\).  The new row \(E_a\) and identity
below are the contracted, differentiated Dirac equation; they are distinct
from the exterior integrability identity
\eqref{eq:V25-fermion-integrability-identity} and do not repeat its proof.

\begin{theorem}[Exact prolongation identity]
\label{thm:V30-exact-incidence}
For every smooth \((\psi,\chi)\), without imposing a field equation,
\begin{equation}
 \boxed{E_a-\nabla_aE=\mathcal DC_a.}
 \label{eq:V30-exact-incidence}
\end{equation}
Consequently:
\begin{enumerate}
 \item if \(E=0\) and \(C_a=0\), then \(E_a=0\);
 \item if \(E=0\) and \(E_a=0\), then each \(C_a\) solves
       \(\mathcal DC_a=0\);
 \item if the homogeneous Dirac problem for \(C_a\) is unique and its
       initial and incoming boundary data vanish, then
       \(E=E_a=0\) implies \(C_a=0\).
\end{enumerate}
\end{theorem}

\begin{proof}
Using \(\nabla_a\gamma^\mu=0\),
\begin{align*}
 \nabla_aE
 &=
 \gamma^\mu\nabla_a\nabla_\mu\psi
 -(\nabla_aM)\psi-M\nabla_a\psi\\
 &=
 \gamma^\mu\nabla_\mu\nabla_a\psi
 +\gamma^\mu[\nabla_a,\nabla_\mu]_{\rm H}\psi
 -(\nabla_aM)\psi-M\nabla_a\psi\\
 &=
 \mathcal D(\nabla_a\psi)
 +\gamma^\mu[\nabla_a,\nabla_\mu]_{\rm H}\psi
 -(\nabla_aM)\psi.
\end{align*}
Subtracting this identity from \eqref{eq:V30-Ea} proves
\eqref{eq:V30-exact-incidence}.  The first two consequences are immediate.
For the third, the difference is a homogeneous Dirac solution with zero
admissible data, so uniqueness gives \(C_a=0\).
\end{proof}

For the covariant-Hessian convention
\eqref{eq:V30-Hessian-commutator}, the commutator can be written, with
representation conventions suppressed, as
\begin{equation}
 [\nabla_a,\nabla_b]_{\rm H}\psi
 =
 \frac14R_{abIJ}\gamma^{IJ}\psi
 +\rho(F_{ab})\psi
 -T^c{}_{ab}\nabla_c\psi.
 \label{eq:V30-spinor-commutator}
\end{equation}
In a noncoordinate frame the structure coefficient is included in the last
term.  On the compatible jet, replace \(\nabla_c\psi\) by \(\chi_c\).

\begin{corollary}[Algebraicity of the prolonged source]
\label{cor:V30-prolonged-source}
If \(R\), \(F\), \(T\), \(\chi\), and \(D\Phi\) are retained values and
\(M=M(\Phi)\), then every term in
\[
 \gamma^\mu[\nabla_a,\nabla_\mu]_{\rm H}\psi-(\nabla_aM)\psi
\]
is algebraic in the V30 state.  Hence the principal spinor symbol of each
row \(E_a\) is another copy of
\begin{equation}
 \sigma_{\mathcal D}(\xi)=\gamma^\mu\xi_\mu.
 \label{eq:V30-Dirac-symbol}
\end{equation}
\end{corollary}

\begin{proof}
Equation~\eqref{eq:V30-spinor-commutator} expresses the commutator through
the retained curvature, torsion, spinor, and spinor-jet values.
Furthermore
\((\nabla_aM)\psi=(D_\Phi M)(D_a\Phi)\psi\), plus an algebraic gauge
commutator when \(M\) is represented covariantly.  None differentiates a
retained value.  The only derivative of \(\chi_a\) is the leading
\(\gamma^\mu\nabla_\mu\chi_a\).
\end{proof}

\begin{proposition}[Quantitative constraint propagation]
\label{prop:V30-constraint-energy}
Suppose the realified Dirac row has a symmetrizer \(H_D\) for which
\[
 h_D I\le H_DA_D^0\le H_D^+I,
 \qquad
 H_DA_D^j=(H_DA_D^j)^*,
 \qquad h_D>0,
\]
and suppose the lower coefficients and first derivatives of the principal
coefficients are bounded by \(L_D\).  On a time slab, impose a homogeneous
maximally nonnegative incoming trace for every \(C_a\).  Then
\begin{equation}
 \|C(t)\|_{L^2}^2
 \le
 \exp(C_Dt)
 \left(
  \|C(0)\|_{L^2}^2+
  \int_0^t\|\mathcal DC(s)\|_{L^2}^2\,\mathrm ds
 \right),
 \label{eq:V30-C-energy}
\end{equation}
where \(C_D\) depends only on \(h_D^{-1}\), \(H_D^+\), and \(L_D\).
In particular, exact prolonged roots with compatible zero data have
\(C=0\).
\end{proposition}

\begin{proof}
Multiply the realified equation for \(C\) by \(C^*H_D\), take the real
part, and integrate by parts.  Symmetry converts the principal spatial
terms into the boundary flux and coefficient-derivative terms.  The
maximally nonnegative incoming condition gives the favorable sign for the
flux.  Cauchy--Schwarz and
\(2ab\le a^2+b^2\) bound the forcing by
\(\|\mathcal DC\|_{L^2}^2\) plus the energy.  The positive time floor makes
the energy equivalent to \(\|C\|_{L^2}^2\), and Gronwall gives
\eqref{eq:V30-C-energy}.
\end{proof}

% -----------------------------------------------------------------------------
\subsection{Bulk symmetric hyperbolicity on the compatible jet}
\label{subsec:V30-hyperbolic}
% -----------------------------------------------------------------------------

\begin{theorem}[Direct-sum Dirac prolongation]
\label{thm:V30-direct-sum}
Fix a finite set of jet directions \(\mathcal I\).  Assume:
\begin{enumerate}
 \item the V29 reduced bosonic wave-jet bank is written in the V27 equation
       frame with positive symmetrizer \(H_G\);
 \item the V25 retained \(R,F,T,D\Phi\) values satisfy their defining
       constraints and have a symmetric-hyperbolic propagation bank with
       symmetrizer \(H_J\);
 \item the V29 realified Dirac bank has symmetrizer \(H_D\);
 \item the source terms in the differentiated Dirac rows are represented
       by \eqref{eq:V30-spinor-commutator}, with no omitted derivative
       column.
\end{enumerate}
Then the principal matrix for
\[
 (U_G,U_J,\psi,(\chi_a)_{a\in\mathcal I})
\]
has no Einstein derivative-stress block.  If the curvature-incidence rows
in item \(2\) make the remaining differentiated-coefficient terms
lower-order in this state, its top-order bank is the direct sum
\begin{equation}
 \mathscr P_{30}(\xi)
 =
 \mathscr P_G(\xi)\oplus\mathscr P_J(\xi)
 \oplus
 \bigoplus_{a\in\{0\}\cup\mathcal I}
 \mathscr P_D(\xi),
 \label{eq:V30-direct-sum-symbol}
\end{equation}
with symmetrizer
\begin{equation}
 H_{30}
 =
 H_G\oplus H_J
 \oplus
 \bigoplus_{a\in\{0\}\cup\mathcal I}H_D.
 \label{eq:V30-direct-sum-H}
\end{equation}
Its time floor is the minimum of the sector floors and its characteristic
cone is the union of the sector cones.  If every sector has the metric cone,
so does the enlarged system.
\end{theorem}

\begin{proof}
Proposition~\ref{prop:V30-defect-location} makes the Dirac stress algebraic
in the state, so differentiating a spinor is absent from the Einstein
principal row.  Corollary~\ref{cor:V30-prolonged-source} gives one Dirac
symbol for \(\psi\) and each \(\chi_a\).  By item \(2\), the curvature and
other differentiated coefficients are values with compatible propagation;
item \(4\) rules out a hidden derivative column.  Therefore the top-order
matrix is \eqref{eq:V30-direct-sum-symbol}.  Blockwise multiplication proves
symmetry of \(H_{30}\mathscr P_{30}^\mu\), and the least block time floor is
positive.  The determinant of a direct sum is the product of the block
determinants, which proves the cone statement.
\end{proof}

\begin{corollary}[Removal of the V29 Schur obstruction]
\label{cor:V30-removes-Sylvester}
Under Theorem~\ref{thm:V30-direct-sum}, the Einstein--Dirac coupling has
\(\rho_0=\rho_x=0\) in the V28 comparison matrix.  No matrix \(X\) solving
\eqref{eq:V29-Sylvester} is required.  The result holds at a general
spinor background.
\end{corollary}

\begin{proof}
The off-diagonal principal blocks are zero in
\eqref{eq:V30-direct-sum-symbol}.  Thus the corresponding V28 coupling
norms vanish.  The size of \(\psi\) affects lower-order coefficients but
does not alter this principal conclusion.
\end{proof}

\begin{firewall}[The curvature hypothesis is a real occurrence row]
\label{fire:V30-curvature-occurrence}
The V25 Cartan and Bianchi identities prove compatibility of retained jets;
they do not by themselves prove that the finite common action produces the
specific symmetric-hyperbolic curvature propagation bank required in
Theorem~\ref{thm:V30-direct-sum}.  If \(R\) is eliminated before
\eqref{eq:V30-Ea} is formed, derivatives of the metric jet reappear.  The
curvature occurrence and its equation multiplier must therefore be read
from the same finite action, or the alternative mixed-order energy estimate
must be certified directly.
\end{firewall}

% -----------------------------------------------------------------------------
\subsection{Root equivalence and the action constraint}
\label{subsec:V30-roots}
% -----------------------------------------------------------------------------

\begin{theorem}[Root equivalence of the compatible prolongation]
\label{thm:V30-root-equivalence}
Let \(\mathcal U\) be a region on which the Dirac initial-boundary problem
is unique.  Consider all original reduced bosonic and matter rows together
with the enlarged rows
\begin{equation}
 E=0,\qquad E_a=0,\qquad
 \mathcal C_{\rm jet}=0,\qquad
 \mathcal C_{\rm Cartan}=0,
 \label{eq:V30-enlarged-roots}
\end{equation}
where \(\mathcal C_{\rm jet}\) contains the V25 defining and integrability
constraints and in particular \(C_a\), and
\(\mathcal C_{\rm Cartan}\) contains
\eqref{eq:V30-Cartan-row}.  Give \(C_a\) zero initial and incoming data.
Then projection
\[
 (g,\omega,\psi,\chi,R,F,T,D\Phi,\ldots)
 \longmapsto(g,\omega,\psi,A,\Phi,\ldots)
\]
is a bijection between compatible enlarged roots and roots of the original
reduced Einstein--Cartan--Standard-Model equations, with
\(\omega=\omega_{\rm EC}\).
\end{theorem}

\begin{proof}
Start with an original root.  Set
\(\chi_a=\nabla_a\psi\), take \(R,F,T,D\Phi\) from their geometric
definitions, and set \(\omega=\omega_{\rm EC}\).  All defining constraints
vanish, the Cartan row vanishes by the imported torsion theorem, and
Theorem~\ref{thm:V30-exact-incidence} gives \(E_a=0\).  This constructs a
compatible enlarged root.

Conversely, let an enlarged root be given.  Proposition
\ref{prop:V30-constraint-energy} and
\eqref{eq:V30-exact-incidence} give \(C_a=0\), so the independent spinor
jets are the derivatives of \(\psi\).  The V25 defining constraints
identify the curvature, gauge curvature, torsion, and Higgs jets with their
geometric values.  Proposition~\ref{prop:V30-EC-target} identifies
\(\omega\) with \(\omega_{\rm EC}\).  Substitution in the remaining rows
gives precisely the original reduced equations.  Both constructions are
forced, hence inverse.
\end{proof}

Let \(\widetilde S_{\rm SMEC}[g,\omega,\psi,\chi,A,\Phi,\ldots]\)
denote the first-order Standard-Model--Einstein--Cartan action obtained by
replacing every occurrence of \(\nabla\psi\) and
\(\nabla\overline\psi\) in the matter Lagrangian by the corresponding
independent jet value.  The action-level realization uses a Lagrange
multiplier \(\lambda^a\) for \(C_a\):
\begin{equation}
 S_{\rm jet}
 =
 \widetilde S_{\rm SMEC}[g,\omega,\psi,\chi,A,\Phi,\ldots]
 +
 \operatorname{Re}\int_{\mathcal U}
 \langle\lambda^a,\chi_a-\nabla_a\psi\rangle\,\mathrm dV_g.
 \label{eq:V30-jet-action}
\end{equation}
Variation in \(\lambda\) imposes \(C=0\); variation in \(\chi\) determines
the multiplier contribution; and elimination on the constraint surface
returns the original action because
\(\widetilde S_{\rm SMEC}|_{\chi=\nabla\psi}=S_{\rm SMEC}\).  The V25
KKT congruence is the applicable Hessian result.  It is not re-proved here.

\begin{proposition}[No spurious action roots under regular reduction]
\label{prop:V30-action-roots}
Assume the constraint derivative
\[
 D_{(\psi,\chi)}C=(-\nabla,I)
\]
has the V25 right-inverse floor on the chosen Sobolev chart and the
multiplier row of \eqref{eq:V30-jet-action} is included.  Then critical
points of \(S_{\rm jet}\) satisfying the regular KKT equations project to
critical points of \(S_{\rm SMEC}\), and every critical point of
\(S_{\rm SMEC}\) has the canonical lift
\(\chi=\nabla\psi\).
\end{proposition}

\begin{proof}
On \(C=0\), an admissible variation of the original fields has the unique
tangent lift
\(\delta\chi=\delta(\nabla\psi)\).  Pairing the KKT equation with every such
tangent variation annihilates the multiplier term and leaves the first
variation of \(S_{\rm SMEC}\).  Hence the projection is critical.
Conversely, lift an original critical point by \(\chi=\nabla\psi\).  The
right-inverse floor and the standard annihilator/range identity solve for
the multiplier that cancels the normal component of the enlarged first
variation.  Thus the full KKT row holds.  This is exactly the regular
constraint reduction encoded abstractly by Theorem
\ref{thm:V25-KKT-critical}.
\end{proof}

% -----------------------------------------------------------------------------
\subsection{The differentiated first-order boundary domain}
\label{subsec:V30-boundary}
% -----------------------------------------------------------------------------

Let \(\mathcal B\) be a timelike boundary with outward spacelike unit normal
\(n\).  Let \(\Pi(x)\) be the orthogonal projector onto the homogeneous
allowed Dirac trace space
\(\mathcal M_x=\operatorname{Ran}\Pi(x)\), as acquired in V20 and transported
in V22--V26.  The original condition is
\begin{equation}
 (I-\Pi)\psi|_{\mathcal B}=0.
 \label{eq:V30-original-boundary}
\end{equation}

\begin{theorem}[Tangentially differentiated trace]
\label{thm:V30-differentiated-trace}
Let \(\tau\) be tangent to \(\mathcal B\) and
\(\chi_\tau=\nabla_\tau\psi\).  Every \(C^1\) solution of
\eqref{eq:V30-original-boundary} satisfies the exact affine trace identity
\begin{equation}
 \boxed{
 (I-\Pi)\chi_\tau=(\nabla_\tau\Pi)\psi
 }\quad\hbox{on }\mathcal B.
 \label{eq:V30-tangential-boundary}
\end{equation}
Its homogeneous top-order trace space is another copy of
\(\mathcal M\).  If \(\mathcal M\) is maximal nonnegative for the Dirac
flux form
\[
 \mathfrak b(u,v)
 :=\langle u,H_DA_D(n)v\rangle,
\]
then \(\mathcal M^{\oplus N}\) is maximal nonnegative for the direct-sum
flux of \(N\) tangential jet copies.
\end{theorem}

\begin{proof}
Differentiate \((I-\Pi)\psi=0\) covariantly along \(\tau\):
\[
 0
 =
 \nabla_\tau((I-\Pi)\psi)
 =
 -(\nabla_\tau\Pi)\psi+(I-\Pi)\nabla_\tau\psi,
\]
which proves \eqref{eq:V30-tangential-boundary}.  The right side contains no
derivative of \(\chi_\tau\), so it is a lower-order boundary source in the
differentiated energy estimate; the homogeneous variation lies in
\(\mathcal M\).

The direct-sum form is
\(\mathfrak b_N((u_j),(u_j))=\sum_j\mathfrak b(u_j,u_j)\), hence it is
nonnegative on \(\mathcal M^{\oplus N}\).  In finite dimension a maximal
nonnegative subspace of a Hermitian form with inertia \((p,z,m)\) has
dimension \(p+z\) and contains the null space.  The direct-sum form has
inertia \((Np,Nz,Nm)\), while \(\mathcal M^{\oplus N}\) has dimension
\(N(p+z)\) and contains its null space.  It therefore cannot have a proper
nonnegative extension and is maximal.
\end{proof}

\begin{proposition}[Normal jet reconstruction]
\label{prop:V30-normal-reconstruction}
Split the Dirac row at \(\mathcal B\) as
\[
 E
 =
 \gamma(n)\nabla_n\psi
 +\gamma(\tau^\alpha)\nabla_{\tau^\alpha}\psi-M\psi.
\]
Because \(g(n,n)=1\), Clifford multiplication satisfies
\(\gamma(n)^2=I\) up to the fixed sign convention and is invertible.  On a
Dirac root the normal jet is therefore determined by
\begin{equation}
 \chi_n
 =
 \gamma(n)^{-1}
 \left(
 M\psi-\gamma(\tau^\alpha)\chi_{\tau^\alpha}
 \right).
 \label{eq:V30-normal-reconstruction}
\end{equation}
It is not an independently prescribable boundary datum.
\end{proposition}

\begin{proof}
Set \(E=0\), substitute
\(\chi_{\tau^\alpha}=\nabla_{\tau^\alpha}\psi\), and multiply by
\(\gamma(n)^{-1}\).
\end{proof}

\begin{corollary}[Boundary debit from a varying projector]
\label{cor:V30-boundary-debit}
Let \(L_\Pi:=\sup_{\mathcal B,\tau}\|\nabla_\tau\Pi\|\), and let
\(\mathfrak L\) be a bounded right inverse of
\((I-\Pi)\) on its range.  Subtract the lift
\[
 \widetilde\chi_\tau
 :=
 \chi_\tau-\mathfrak L((\nabla_\tau\Pi)\psi).
\]
Then \(\widetilde\chi_\tau\in\mathcal M\), and the differentiated boundary
energy has the same homogeneous flux sign as the original Dirac problem,
plus a lower-order debit bounded by
\begin{equation}
 C_{\partial}\|\mathfrak L\|L_\Pi
 \|\psi\|_{\partial}\|\chi_\tau\|_{\partial}
 +
 C_{\partial}\|\mathfrak L\|^2L_\Pi^2
 \|\psi\|_{\partial}^2.
 \label{eq:V30-projector-debit}
\end{equation}
\end{corollary}

\begin{proof}
Equation~\eqref{eq:V30-tangential-boundary} gives
\((I-\Pi)\widetilde\chi_\tau=0\).  Expand the quadratic flux of
\(\chi_\tau=\widetilde\chi_\tau+
\mathfrak L((\nabla_\tau\Pi)\psi)\).  The first term has the admitted sign.
The cross term is bounded by Cauchy--Schwarz and the operator norms, and the
square of the lift gives the second term in
\eqref{eq:V30-projector-debit}.
\end{proof}

\begin{firewall}[Normal differentiation is recovered, not prescribed]
\label{fire:V30-normal-domain}
One may differentiate the admitted boundary condition tangentially.
Differentiating it in the outward normal direction would impose extra
boundary data and generally overdetermine the first-order Dirac problem.
Normal spinor jets must be reconstructed from
\eqref{eq:V30-normal-reconstruction}, or controlled by the original Dirac
equation in the usual tangential energy hierarchy.  A finite boundary
certificate that assigns an independent projector to every normal jet does
not pass V30.
\end{firewall}

% -----------------------------------------------------------------------------
\subsection{Finite and cofinal certificate}
\label{subsec:V30-finite}
% -----------------------------------------------------------------------------

At cofinal level \(q\), enlarge the V29 field order by
\[
 (\psi_q,\chi_{a,q},R_q,F_q,T_q,D\Phi_q,\lambda_q^a)
\]
and acquire every row from the same common action and the same V24 frame.
In addition to the earlier V19--V29 residuals, store
\begin{align}
 r_{C,q}
 &:=
 \|\chi_{a,q}-\nabla_{a,q}\psi_q\|,
 \label{eq:V30-rC}\\
 r_{{\rm inc},q}
 &:=
 \|E_{a,q}-\nabla_{a,q}E_q-\mathcal D_qC_{a,q}\|,
 \label{eq:V30-rinc}\\
 r_{{\rm comm},q}
 &:=
 \left\|
 [\nabla_{a,q},\nabla_{b,q}]_{\rm H}\psi_q
 -\frac14R_{abIJ,q}\gamma_q^{IJ}\psi_q
 -\rho(F_{ab,q})\psi_q
 +T^c{}_{ab,q}\chi_{c,q}
 \right\|,
 \label{eq:V30-rcomm}\\
 r_{M,q}
 &:=
 \|\nabla_{a,q}M_q-(D_\Phi M)_q(D_a\Phi)_q\|,
 \label{eq:V30-rM}\\
 r_{{\rm Cartan},q}
 &:=
 \|\mathscr C_{e_q}(T_q)+\Sigma(e_q,\psi_q)\|,
 \label{eq:V30-rCartan}\\
 r_{\omega{\rm EC},q}
 &:=
 \|\omega_q-\omega^{\rm LC}(e_q)-K_{\rm EC}(e_q,\psi_q)\|,
 \label{eq:V30-romega}\\
 r_{{\rm alg},q}
 &:=
 \max_\mu
 \left\|
 \frac{\partial\mathcal T_{{\rm D},q}}
 {\partial(\partial_\mu\chi_q)}
 \right\|,
 \label{eq:V30-ralg}\\
 r_{{\rm sum},q}
 &:=
 \max_\mu
 \left\|
 \widehat{\mathscr P}_{30,q}^\mu
 -
 \left(
 \widehat{\mathscr P}_{G,q}^\mu\oplus
 \widehat{\mathscr P}_{J,q}^\mu\oplus
 \bigoplus\widehat{\mathscr P}_{D,q}^\mu
 \right)
 \right\|.
 \label{eq:V30-rsum}
\end{align}
The derivative in \eqref{eq:V30-ralg} is taken after \(\chi\) has been
declared an independent jet coordinate.  It must vanish; testing only the
formula obtained after eliminating \(\chi\) would reproduce V29 instead.

For a timelike boundary also store
\begin{align}
 r_{\partial\tau,q}
 &:=
 \max_\tau
 \|(I-\Pi_q)\chi_{\tau,q}
       -(\nabla_{\tau,q}\Pi_q)\psi_q\|,
 \label{eq:V30-rbtan}\\
 r_{\partial n,q}
 &:=
 \left\|
 \chi_{n,q}
 -
 \gamma_q(n)^{-1}
 \left(M_q\psi_q-\gamma_q(\tau^\alpha)
                   \chi_{\tau^\alpha,q}\right)
 \right\|,
 \label{eq:V30-rbnorm}\\
 r_{\partial{\rm flux},q}
 &:=
 \left\|
 \mathfrak b_{30,q}|_{\mathcal M_{30,q}}
 -
 \bigoplus\mathfrak b_{D,q}|_{\mathcal M_q}
 \right\|.
 \label{eq:V30-rbflux}
\end{align}
The product trace \(\mathcal M_{30,q}\) includes only the original spinor
and tangential differentiated copies; the normal component is reconstructed
by \eqref{eq:V30-rbnorm}.

Let
\begin{align}
 c_{C,q}^{\rm cert}
 &:=
 \sigma_{\min}(\mathscr C_{e_q})-\varepsilon_{C,q},
 \label{eq:V30-Cartan-floor}\\
 c_{J,q}^{\rm cert}
 &:=
 \sigma_{\min}(\mathscr J_{e_q})-\varepsilon_{J,q},
 \label{eq:V30-contorsion-floor}\\
 h_{30,q}^{\rm cert}
 &:=
 \min(h_{G,q}^{\rm cert},h_{J,q}^{\rm cert},
      h_{D,q}^{\rm cert})-\varepsilon_{30,q},
 \label{eq:V30-time-floor}\\
 \beta_{30,q}^{\rm cert}
 &:=
 \beta_{D,q}^{\rm cert}
 -C_\partial\bigl(
    \|\mathfrak L_q\|\,\|\nabla_\tau\Pi_q\|\,\eta_q
    +r_{\partial{\rm flux},q}\bigr)
 \label{eq:V30-boundary-floor}
\end{align}
be the certified Cartan, contorsion, time, and boundary reserves.  Here
\(\eta_q\) is the trace-radius debit needed to absorb the lower-order term
\eqref{eq:V30-projector-debit}.  It is zero for a covariantly constant
projector.

\begin{theorem}[Cofinal spinor-jet handoff]
\label{thm:V30-cofinal-handoff}
Assume:
\begin{enumerate}
 \item the V19 occurrence, V24 frame, V25 compatible-complex, V26
       equation-incidence, V27 action, and V28 cone hypotheses hold for the
       enlarged V30 field list;
 \item all residuals
       \eqref{eq:V30-rC}--\eqref{eq:V30-rsum} tend to zero uniformly and
       the cofinal tails are summable;
 \item the three floors
       \(c_{C,q}^{\rm cert}\), \(c_{J,q}^{\rm cert}\), and
       \(h_{30,q}^{\rm cert}\) have positive uniform lower bounds;
 \item at a timelike boundary,
       \eqref{eq:V30-rbtan}--\eqref{eq:V30-rbflux} tend to zero,
       \(\inf_q\beta_{30,q}^{\rm cert}>0\), and the V25 subsidiary
       constraints have compatible incoming data.
\end{enumerate}
Then:
\begin{enumerate}
 \item the cofinal connection is the spin-dependent
       Einstein--Cartan connection, with error controlled by
       \eqref{eq:V30-connection-error};
 \item the enlarged spinor jets are equivalent to derivatives of the
       original spinor at action and causal roots;
 \item the V29 derivative-stress block is absent from the enlarged
       principal bank, so the Einstein--Dirac top-order coupling needs no
       Sylvester correction;
 \item the limiting enlarged bank is symmetric hyperbolic with the metric
       characteristic cone and preserves the differentiated first-order
       fermion domain;
 \item the V28 moving-boundary handoff applies after its remaining
       boundary-speed or protected-centre margin is inserted.
\end{enumerate}
\end{theorem}

\begin{proof}
The positive Cartan and contorsion floors and
\eqref{eq:V30-rCartan} give the limiting connection by
Proposition~\ref{prop:V30-EC-target}.  Vanishing incidence residuals pass
\eqref{eq:V30-exact-incidence} to the limit.  The constraint energy
\eqref{eq:V30-C-energy}, together with the initial and incoming
compatibility, gives \(C=0\); Theorem~\ref{thm:V30-root-equivalence} then
proves root equivalence.

Vanishing algebraicity and direct-sum residuals give
\eqref{eq:V30-direct-sum-symbol}.  The positive time floor and
Theorem~\ref{thm:V30-direct-sum} give symmetric hyperbolicity and the metric
cone.  The tangential and normal boundary residuals pass
\eqref{eq:V30-tangential-boundary} and
\eqref{eq:V30-normal-reconstruction} to the limit.
Theorem~\ref{thm:V30-differentiated-trace} and the positive reserve
\eqref{eq:V30-boundary-floor} preserve the first-order allowed trace.
Finally, the V26 incidence theorem transfers the action root to the causal
row, while the separate V28 boundary margin supplies the claimed moving
boundary statement.
\end{proof}

\begin{firewall}[Conditional continuum status]
\label{fire:V30-status}
Theorem~\ref{thm:V30-cofinal-handoff} closes the logical
Einstein--Dirac derivative-stress obstruction, but the attached files do
not contain the populated spinor-jet packet columns, the common-action
curvature propagation row, the differentiated allowed-trace columns, or
their positive finite floors.  Those are occurrence data, not consequences
of the continuum identities.  The full Standard Model--Einstein continuum
is therefore not claimed in V30.
\end{firewall}

% -----------------------------------------------------------------------------
\subsection{V30 ledger and next upstream row}
\label{subsec:V30-status}
% -----------------------------------------------------------------------------

\paragraph{New conclusions proved in V30.}
Theorem~\ref{thm:V30-exact-incidence} gives the exact covariant
spinor-jet residual identity.  Proposition~\ref{prop:V30-constraint-energy}
turns it into quantitative constraint propagation.  Proposition
\ref{prop:V30-defect-location} proves that the V29 triangular stress block
is removed by a compatible jet value, while
Theorem~\ref{thm:V30-direct-sum} states and proves the precise
curvature-incidence hypotheses under which the enlarged bank is
direct-sum symmetric hyperbolic.  Proposition~\ref{prop:V30-EC-target}
shifts the connection target from Levi--Civita to Einstein--Cartan.
Theorem~\ref{thm:V30-differentiated-trace} and Proposition
\ref{prop:V30-normal-reconstruction} preserve the original first-order
fermion boundary domain.

\paragraph{Imported conclusions not repeated.}
Connection stationarity, injectivity of the Cartan map, uniqueness and the
finite bound for spin-sourced torsion, the weak Palatini insertion, V25 jet
and Bianchi compatibility, V26 action/causal incidence, V27 Legendre
positivity, and the V28 moving-boundary theorem remain earlier results.

\paragraph{Authoritative physical status.}
No attached manuscript prints the finite Hessian columns for
\((\psi,\chi,\lambda)\) in \eqref{eq:V30-jet-action}, a realified
curvature-jet evolution symbol incident with that Hessian, or the
tangentially differentiated V20 boundary closure.  V30 supplies a finite
test that such data can pass.  It does not manufacture the missing data or
infer them from continuum plausibility.

\paragraph{Next upstream row.}
The next iteration should construct the common-action occurrence card for
\eqref{eq:V30-jet-action}.  It must derive, from actual packet and boundary
impulses, the multiplier block, the curvature commutator column, the
direct-sum residual \eqref{eq:V30-rsum}, and the tangential trace
\eqref{eq:V30-rbtan}.  If the curvature row is not generated with a
positive hyperbolic floor, the correct upstream alternative is a
Douglis--Nirenberg mixed-order energy certificate for
\eqref{eq:V30-Einstein-order}--\eqref{eq:V30-Dirac-row}; returning to the
unprolonged uniform first-order triangular bank would recreate the V29
obstruction.

% END APPENDED VERSION V30 MATERIAL


% =============================================================================
% START APPENDED VERSION V31 MATERIAL
% =============================================================================

\section{V31: Fermion KKT incidence, skew action symbol, and the
action-compatible boundary fork}
\label{sec:V31}

\paragraph{Purpose and relation to earlier notes.}
V25 proves the abstract KKT congruence for a jet constraint, V27 proves a
Hamilton--Pontryagin differential readout for a general first-order
bosonic chart, and V30 proves the covariant differentiated Dirac identity.
V31 does not repeat those results.  It computes the missing fermion-specific
composition.  The calculation has three consequences.  First, the affine
fermion Lagrangian has zero jet-velocity Hessian, but the KKT saddle Hessian
can still be invertible.  Second, the action Euler row is converted to the
causal Dirac row by an explicit zero-order skew multiplier, after which the
V30 prolongation is a second factored differential incidence.  Third, a
conservative scalar action and a timelike-boundary Dirac domain are
compatible only on an isotropic, hence zero-flux, trace space unless
additional boundary degrees of freedom carry the missing flux.

\paragraph{External consistency check.}
The MIT and related boundary conditions treated by Gro{\ss}e and Murro have
zero normal Dirac flux and yield a symmetric positive hyperbolic
initial-boundary problem; see
\url{https://arxiv.org/abs/1806.06544}.  V31 uses no theorem from that paper
in its proof.  The citation records that the zero-flux branch isolated
below is a populated continuum boundary class rather than a formal empty
case.

% -----------------------------------------------------------------------------
\subsection{Realified affine fermion action and its KKT residual}
\label{subsec:V31-KKT}
% -----------------------------------------------------------------------------

Let \(S_{\mathbb R}\) be the realification of the spinor representation,
including all Standard Model internal indices.  Work in one V24 common
frame and let \(z\) denote the realified spinor.  Let
\(\chi_a\) be independent covariant jet values and let
\(\ell(z,\chi;e,\omega,A,\Phi)\) be the real fermion Lagrangian density.
The other fields are parameters in the following fibrewise calculation.
Put
\begin{equation}
 p^a:=D_{\chi_a}\ell.
 \label{eq:V31-momentum}
\end{equation}
The formal covariant adjoint is fixed by
\begin{equation}
 \int_{\mathcal U}\langle\lambda^a,\nabla_a u\rangle\,\mathrm dV
 =
 \int_{\mathcal U}\langle\nabla_a^\dagger\lambda^a,u\rangle\,\mathrm dV
 +
 \int_{\partial\mathcal U}
 \langle\lambda^n,u\rangle\,\mathrm dS,
 \qquad
 \lambda^n:=n_a\lambda^a.
 \label{eq:V31-adjoint-convention}
\end{equation}
Thus the leading part of \(\nabla^\dagger\) is minus the covariant
divergence.

Consider the constrained action
\begin{equation}
 S_{\rm K,D}[z,\chi,\lambda]
 :=
 \int_{\mathcal U}
 \left(
  \ell(z,\chi)
  +\langle\lambda^a,\chi_a-\nabla_a z\rangle
 \right)\mathrm dV.
 \label{eq:V31-KKT-action}
\end{equation}
Its interior residuals are
\begin{align}
 K_z
 &:=
 D_z\ell-\nabla_a^\dagger\lambda^a,
 \label{eq:V31-Kz}\\
 K_{\chi_a}
 &:=
 p^a+\lambda^a,
 \label{eq:V31-Kchi}\\
 C_a
 &:=
 \chi_a-\nabla_a z.
 \label{eq:V31-C}
\end{align}
The reduced Euler residual, evaluated before imposing \(C=0\), is
\begin{equation}
 E_{\rm A}(z,\chi)
 :=
 D_z\ell+\nabla_a^\dagger p^a.
 \label{eq:V31-action-Euler}
\end{equation}
The subscript A distinguishes the formally variational row from the causal
Dirac normalization.

\begin{theorem}[Exact fermion KKT-to-Euler incidence]
\label{thm:V31-KKT-incidence}
For arbitrary smooth independent \((z,\chi,\lambda)\),
\begin{equation}
 \boxed{
 E_{\rm A}
 =
 K_z+\nabla_a^\dagger K_{\chi_a}.}
 \label{eq:V31-KKT-incidence}
\end{equation}
No action equation or jet constraint is used.  The inverse residual
reconstruction is
\begin{equation}
 K_z
 =
 E_{\rm A}-\nabla_a^\dagger K_{\chi_a}.
 \label{eq:V31-Kz-inverse}
\end{equation}
\end{theorem}

\begin{proof}
Substitute \eqref{eq:V31-Kz} and \eqref{eq:V31-Kchi}:
\[
 K_z+\nabla_a^\dagger K_{\chi_a}
 =
 D_z\ell-\nabla_a^\dagger\lambda^a
 +\nabla_a^\dagger(p^a+\lambda^a)
 =
 D_z\ell+\nabla_a^\dagger p^a.
\]
This is \eqref{eq:V31-action-Euler}; rearrangement gives
\eqref{eq:V31-Kz-inverse}.
\end{proof}

\begin{corollary}[Affine fermions do not need a velocity Legendre inverse]
\label{cor:V31-no-Legendre}
Suppose \(\ell\) is affine in \(\chi\), so
\begin{equation}
 D_{\chi\chi}^2\ell=0.
 \label{eq:V31-affine-jet}
\end{equation}
The fermion velocity Legendre map is then noninvertible.  Nevertheless, the
V25 KKT Hessian remains invertible exactly when its reduced action Hessian
is invertible; its constraint-multiplier saddle block supplies the missing
unit pairing.  Therefore the V27 positive velocity-Hessian criterion is a
bosonic route and is not a required fermion hypothesis.
\end{corollary}

\begin{proof}
The velocity Hessian is the V25 block \(C\), which is zero under
\eqref{eq:V31-affine-jet}.  The exact V25 congruence
\eqref{eq:V25-KKT-congruence} then has the diagonal form
\[
 H_{\rm red}\oplus
 \begin{pmatrix}0&I\\I&0\end{pmatrix}.
\]
The second summand is an invertible self-adjoint unitary.  Hence the KKT
Hessian is invertible precisely when \(H_{\rm red}\) is, even though the
velocity Legendre map has zero derivative.
\end{proof}

\begin{firewall}[A fermion is not a failed bosonic Legendre sector]
\label{fire:V31-Legendre}
Requiring a positive \(D_{\chi\chi}^2\ell\) would change the first-order
Dirac action into a second-order spinor theory.  The correct finite
certificate is affine-jet incidence plus the KKT action floor and the
skew action-to-causal multiplier derived below.
\end{firewall}

% -----------------------------------------------------------------------------
\subsection{The skew action symbol and the causal Dirac multiplier}
\label{subsec:V31-skew-symbol}
% -----------------------------------------------------------------------------

At a frozen point write
\begin{equation}
 B^a:=D^2_{z\chi_a}\ell,
 \qquad
 C^{ab}:=D^2_{\chi_a\chi_b}\ell.
 \label{eq:V31-BC}
\end{equation}
Adjoints are taken in the same real Hilbert frame as the V25 Hessian.

\begin{proposition}[Principal Euler symbol of a first-order real action]
\label{prop:V31-action-symbol}
The principal part of the linearization of
\(E_{\rm A}(z,\nabla z)\) is
\begin{equation}
 (B^a-(B^a)^*)\nabla_a\delta z
 -
 C^{ab}\nabla_a\nabla_b\delta z.
 \label{eq:V31-action-principal}
\end{equation}
Consequently, under the affine condition
\eqref{eq:V31-affine-jet}, the action Euler coefficient
\begin{equation}
 J_{\rm A}^a:=B^a-(B^a)^*
 \label{eq:V31-skew-coefficient}
\end{equation}
is skew-adjoint.  This skewness is forced by symmetry of the action
Hessian and the skew formal adjoint of a first derivative.
\end{proposition}

\begin{proof}
Linearize \eqref{eq:V31-action-Euler}.  The derivative of \(D_z\ell\)
through the jet argument contributes
\(B^a\nabla_a\delta z\).  The derivative of \(p^a\) through \(z\) is
\((B^a)^*\delta z\), and the leading part of
\(\nabla_a^\dagger\) is \(-\nabla_a\); hence it contributes
\(-(B^a)^*\nabla_a\delta z\).  The derivative of \(p^a\) through
\(\chi_b=\nabla_bz\) contributes
\(-C^{ab}\nabla_a\nabla_b\delta z\).  Coefficient derivatives are lower
order.  This proves \eqref{eq:V31-action-principal}.  When \(C=0\),
taking the adjoint of \eqref{eq:V31-skew-coefficient} gives
\((J_{\rm A}^a)^*=-J_{\rm A}^a\).
\end{proof}

The real action Euler row is not yet the positive-time causal row.  Let
\(\mathfrak j_D\) be an invertible, zeroth-order real bundle map, normally
the real complex structure and the fixed temporal Clifford normalization,
and define
\begin{equation}
 E_D:=\mathfrak j_D E_{\rm A}.
 \label{eq:V31-causal-multiplier}
\end{equation}
The coefficient test below controls the principal row.  Full nonlinear
incidence additionally requires the zeroth-order and algebraic residual
\begin{equation}
 R_{D,0}:=E_D-\mathcal Dz
 \label{eq:V31-full-Dirac-residual}
\end{equation}
to vanish in the same field frame.

\begin{theorem}[Action-to-Dirac coefficient criterion]
\label{thm:V31-action-Dirac}
Assume \eqref{eq:V31-affine-jet}.  The causal row
\eqref{eq:V31-causal-multiplier} has target principal coefficients
\(A_D^a\) if and only if
\begin{equation}
 \boxed{
 A_D^a
 =
 \mathfrak j_D\bigl(B^a-(B^a)^*\bigr)
 \quad\hbox{for every }a.}
 \label{eq:V31-action-Dirac-coefficients}
\end{equation}
If the V29 Dirac symmetrizer \(H_D\) satisfies
\[
 H_DA_D^0>0,
 \qquad
 H_DA_D^i=(H_DA_D^i)^*,
\]
then \(\mathfrak j_D\) is the required common-action equation multiplier
for the fermion block.  Conversely, failure of
\eqref{eq:V31-action-Dirac-coefficients} at one coefficient rules out that
multiplier in the chosen field and equation frame.
\end{theorem}

\begin{proof}
Proposition~\ref{prop:V31-action-symbol} gives
\(\sigma_1(E_{\rm A})(\xi)=
iJ_{\rm A}^a\xi_a\), with the harmless Fourier \(i\) absorbed into the
realification convention.  Left multiplication by \(\mathfrak j_D\)
therefore gives the coefficients on the right of
\eqref{eq:V31-action-Dirac-coefficients}.  Equality coefficient by
coefficient is necessary and sufficient for equality for every covector.
The stated positivity and symmetry conditions are exactly the V29
symmetric-hyperbolic criterion in that equation frame.  A coefficient
mismatch cannot be repaired by lower-order terms.
\end{proof}

\begin{corollary}[Characteristic incidence]
\label{cor:V31-characteristic-incidence}
If \(\mathfrak j_D\) is invertible and
\eqref{eq:V31-action-Dirac-coefficients} holds, then
\[
 \ker\sigma_1(E_D)(\xi)
 =
 \ker\sigma_1(E_{\rm A})(\xi)
\]
and both rows have the V29 metric null cone.  Thus the multiplier changes
the energy normalization but not the characteristic covectors.
\end{corollary}

\begin{proof}
Left multiplication by an invertible map preserves the right kernel.
The V29 Clifford calculation identifies the characteristic set of
\(A_D^a\xi_a\) with the metric null cone.
\end{proof}

% -----------------------------------------------------------------------------
\subsection{Factored incidence from the common action to the V30 system}
\label{subsec:V31-factored}
% -----------------------------------------------------------------------------

Use the V30 jet defect \(C_a\) and define the prolonged causal residual by
\begin{equation}
 E_{D,a}
 :=
 \nabla_aE_D+\mathcal D C_a.
 \label{eq:V31-prolonged-causal}
\end{equation}
When \(R_{D,0}=0\), this is the V30 row
\eqref{eq:V30-Ea}, by Theorem~\ref{thm:V30-exact-incidence}.
Define the residual map
\begin{equation}
 \mathcal Q_{31}(K_z,K_\chi,C)
 :=
 (E_D,K_\chi,C,E_{D,a}),
 \label{eq:V31-Q}
\end{equation}
where
\begin{align}
 E_D
 &:=
 \mathfrak j_D
 \left(K_z+\nabla_a^\dagger K_{\chi_a}\right),
 \label{eq:V31-Q-E}\\
 E_{D,a}
 &:=
 \nabla_aE_D+\mathcal DC_a.
 \label{eq:V31-Q-Ea}
\end{align}

\begin{theorem}[Exact compatible graph and inverse]
\label{thm:V31-compatible-graph}
The map \(\mathcal Q_{31}\) factors as two differential maps of order one:
\[
 (K_z,K_\chi,C)
 \longmapsto(E_D,K_\chi,C)
 \longmapsto(E_D,K_\chi,C,E_{D,a}).
\]
Its total differential order is two in \(K_\chi\).  Its range is the
compatible graph
\begin{equation}
 \mathcal C_{31}^{\rm comp}
 =
 \left\{
 (E_D,K_\chi,C,E_{D,a}):
 E_{D,a}=\nabla_aE_D+\mathcal DC_a
 \right\}.
 \label{eq:V31-compatible-graph}
\end{equation}
On this graph it is a bijection, with inverse
\begin{equation}
 \boxed{
 K_z
 =
 \mathfrak j_D^{-1}E_D
 -\nabla_a^\dagger K_{\chi_a}.}
 \label{eq:V31-Q-inverse}
\end{equation}
\end{theorem}

\begin{proof}
The first factor is Theorem~\ref{thm:V31-KKT-incidence} followed by the
zeroth-order map \(\mathfrak j_D\).  The second factor is the exact
definition \eqref{eq:V31-Q-Ea}.  Hence every image lies in
\eqref{eq:V31-compatible-graph}.  Conversely, for a tuple in that graph,
define \(K_z\) by \eqref{eq:V31-Q-inverse}.  Then
\eqref{eq:V31-Q-E} and \eqref{eq:V31-Q-Ea} hold, so the tuple is in the
range.  The retained \(K_\chi\) and \(C\) rows and
\eqref{eq:V31-Q-inverse} also prove injectivity.
\end{proof}

Let \(s\ge1\).  On a fixed compact chart suppose
\begin{align}
 \|\nabla^\dagger u\|_{H^s}
 &\le d_s\|u\|_{H^{s+1}},
 \label{eq:V31-derivative-bound1}\\
 \|\nabla u\|_{H^{s-1}}+\|\mathcal Du\|_{H^{s-1}}
 &\le d_s\|u\|_{H^s},
 \label{eq:V31-derivative-bound2}\\
 \|\mathfrak j_D\|
 &\le j_+,
 \qquad
 \|\mathfrak j_D^{-1}\|\le j_-.
 \label{eq:V31-j-bounds}
\end{align}
Equip the action residual space with
\begin{equation}
 \mathcal A_{31}^s
 :=
 H^s(K_z)\oplus H^{s+1}(K_\chi)\oplus H^s(C)
 \label{eq:V31-action-space}
\end{equation}
and the compatible graph with the product norm in
\begin{equation}
 H^s(E_D)\oplus H^{s+1}(K_\chi)
 \oplus H^s(C)\oplus H^{s-1}(E_{D,a}).
 \label{eq:V31-causal-space}
\end{equation}

\begin{theorem}[No loss on the shifted fermion graph]
\label{thm:V31-shifted-bounds}
With the preceding norms,
\(\mathcal Q_{31}:\mathcal A_{31}^s\to
\mathcal C_{31}^{\rm comp}\) and its inverse are bounded.  One may take
\begin{align}
 \|\mathcal Q_{31}\|^2
 &\le
 \max\left\{
  2j_+^2(1+2d_s^2),\,
  1+2j_+^2d_s^2+4j_+^2d_s^4,\,
  1+2d_s^2
 \right\},
 \label{eq:V31-Q-bound}\\
 \|\mathcal Q_{31}^{-1}\|^2
 &\le
 \max\left\{
  2j_-^2,\,
  1+2d_s^2,\,
  1
 \right\}.
 \label{eq:V31-Q-inverse-bound}
\end{align}
\end{theorem}

\begin{proof}
Equations \eqref{eq:V31-Q-E} and
\eqref{eq:V31-derivative-bound1} give
\[
 \|E_D\|_{H^s}^2
 \le
 2j_+^2\|K_z\|_{H^s}^2
 +2j_+^2d_s^2\|K_\chi\|_{H^{s+1}}^2.
\]
Equations \eqref{eq:V31-Q-Ea} and
\eqref{eq:V31-derivative-bound2} give
\[
 \|E_{D,a}\|_{H^{s-1}}^2
 \le
 2d_s^2\|E_D\|_{H^s}^2
 +2d_s^2\|C\|_{H^s}^2.
\]
Add the retained \(K_\chi\) and \(C\) norms and collect the coefficient of
each input norm.  This gives \eqref{eq:V31-Q-bound}.  For the inverse,
\eqref{eq:V31-Q-inverse} gives
\[
 \|K_z\|_{H^s}^2
 \le
 2j_-^2\|E_D\|_{H^s}^2
 +2d_s^2\|K_\chi\|_{H^{s+1}}^2.
\]
Adding the retained rows proves \eqref{eq:V31-Q-inverse-bound}.
\end{proof}

\begin{corollary}[Root equivalence including the action algebraic row]
\label{cor:V31-root-equivalence}
On the compatible graph,
\[
 K_z=K_\chi=C=0
 \quad\Longleftrightarrow\quad
 E_D=K_\chi=C=E_{D,a}=0.
\]
If \(R_{D,0}=0\), the causal residual is the physical Dirac residual;
hence the V30 prolonged causal root and the fermion KKT action root agree
provided the algebraic row \(K_\chi=0\) is retained.
\end{corollary}

\begin{proof}
The forward implication follows from \eqref{eq:V31-Q-E} and
\eqref{eq:V31-Q-Ea}.  Conversely,
\eqref{eq:V31-Q-inverse} gives \(K_z=0\); the other action residuals are
retained explicitly.
\end{proof}

\begin{firewall}[The multiplier is factored and graph-valued]
\label{fire:V31-factor}
The composite map is not an invertible zeroth-order square matrix on an
ambient unshifted target.  It is a product of two first-order differential
maps and is invertible only onto the compatible graph while retaining
\(K_\chi\) and \(C\).  Fitting one unconstrained matrix to the action and
causal packet banks would test the wrong object.
\end{firewall}

% -----------------------------------------------------------------------------
\subsection{Boundary variation and the zero-flux theorem}
\label{subsec:V31-boundary}
% -----------------------------------------------------------------------------

Equation~\eqref{eq:V31-adjoint-convention} gives the boundary part of the
first variation of \eqref{eq:V31-KKT-action}:
\begin{equation}
 \delta S_{\rm K,D}\big|_{\partial\mathcal U}
 =
 -\int_{\partial\mathcal U}
 \langle\lambda^n,\delta z\rangle\,\mathrm dS.
 \label{eq:V31-KKT-boundary-variation}
\end{equation}
On the algebraic row \(K_\chi=0\), this becomes
\begin{equation}
 \delta S_{\rm K,D}\big|_{\partial\mathcal U}
 =
 \int_{\partial\mathcal U}
 \langle p^n,\delta z\rangle\,\mathrm dS,
 \label{eq:V31-reduced-boundary-variation}
\end{equation}
which is exactly the boundary term of the reduced first-order action.

Let \(\mathcal M\) be a linear allowed trace space and let
\(\mathcal M^\circ\) denote its annihilator in the real boundary duality.
If a boundary scalar density \(b(z)\) is added, action differentiability
on traces \(z,\delta z\in\mathcal M\) requires
\begin{equation}
 p^n+D b(z)\in\mathcal M^\circ.
 \label{eq:V31-boundary-conormal}
\end{equation}

\begin{proposition}[Exact KKT boundary incidence]
\label{prop:V31-boundary-incidence}
The KKT and reduced boundary covectors agree on \(K_\chi=0\):
\begin{equation}
 -\lambda^n=p^n.
 \label{eq:V31-boundary-covector-incidence}
\end{equation}
Consequently the KKT action and the reduced action impose the same
conormal condition \eqref{eq:V31-boundary-conormal}.  No independent
boundary value may be assigned to \(\lambda^n\) after the algebraic row is
enforced.
\end{proposition}

\begin{proof}
Contract \(K_{\chi_a}=p^a+\lambda^a=0\) with the outward normal.
Substitution in \eqref{eq:V31-KKT-boundary-variation} gives
\eqref{eq:V31-reduced-boundary-variation}.
\end{proof}

Let \(\mathfrak g_n\) be the real Green form of the linearized action Euler
operator:
\begin{equation}
 \int_{\mathcal U}
 \bigl(
  \langle u,Lv\rangle-\langle Lu,v\rangle
 \bigr)\mathrm dV
 =
 \int_{\partial\mathcal U}\mathfrak g_n(u,v)\,\mathrm dS.
 \label{eq:V31-Green-form}
\end{equation}
The V29 causal normal flux is
\begin{equation}
 \mathfrak b_n(u,v)
 :=
 \langle u,H_DA_D(n)v\rangle.
 \label{eq:V31-causal-flux}
\end{equation}
Boundary action/causal incidence means that, after the same realification
and equation normalization, there is a fixed nonzero scalar \(\zeta_D\)
(the conventional complex factor in the standard Dirac normalization)
such that
\begin{equation}
 \mathfrak g_n(u,v)=\zeta_D\mathfrak b_n(u,v).
 \label{eq:V31-Green-flux-incidence}
\end{equation}
This equality is a V26 boundary-incidence datum and is certified below;
it is not inferred from the equality of bulk characteristic cones.

\begin{theorem}[Conservative action forces an isotropic Dirac trace]
\label{thm:V31-zero-flux}
Assume the following on a timelike boundary component:
\begin{enumerate}
 \item the bulk-plus-boundary functional is a real \(C^2\) scalar action;
 \item its boundary density depends locally on boundary fields but contains
       no independent propagating boundary reservoir;
 \item the tangent trace space for both field variations is
       \(\mathcal M\), and it is closed under the realified complex
       structure when the spinor representation is complex;
 \item the action Green form and causal flux obey
       \eqref{eq:V31-Green-flux-incidence} in the common frame.
\end{enumerate}
Then
\begin{equation}
 \boxed{
 \mathfrak b_n(u,v)=0
 \quad\hbox{for every }u,v\in\mathcal M.}
 \label{eq:V31-isotropic-trace}
\end{equation}
In particular, if \(\mathcal M\) is also maximally nonnegative for the
causal flux, it is a maximal zero-flux subspace.  A strictly dissipative
trace space containing some \(u\) with
\(\mathfrak b_n(u,u)>0\) cannot be the boundary domain of that conservative
bulk action.
\end{theorem}

\begin{proof}
Take two variations \(u,v\) whose traces lie in \(\mathcal M\).
Symmetry of the second derivative of a real \(C^2\) scalar functional gives
\[
 D^2S[u,v]-D^2S[v,u]=0.
\]
The interior difference is the Green identity
\eqref{eq:V31-Green-form}.  The Hessian of any scalar boundary density is
symmetric, so it contributes zero to this antisymmetrized second
variation.  For a complex spinor, testing also with the realified complex
multiple of either variation annihilates both real polarizations.
Therefore \(\mathfrak g_n(u,v)=0\) for all admitted traces.  Since
\(\zeta_D\ne0\), \eqref{eq:V31-Green-flux-incidence} proves
\eqref{eq:V31-isotropic-trace}.  Its diagonal
then vanishes.  Hence a trace with strictly positive flux contradicts the
necessary isotropy.
\end{proof}

\begin{corollary}[Balanced neutral incoming/outgoing graph]
\label{cor:V31-neutral-graph}
Let the self-adjoint causal flux operator have spectral decomposition
\(S_+\oplus S_0\oplus S_-\) and inertia \((p,z,p)\).  If
\(\mathcal M\) is maximal totally isotropic for \(\mathfrak b_n\), then
\(S_0\subset\mathcal M\) and there is an isometry
\begin{equation}
 U:(S_+,\mathfrak b_n)\longrightarrow
 (S_-,-\mathfrak b_n)
 \label{eq:V31-neutral-isometry}
\end{equation}
such that
\begin{equation}
 \mathcal M
 =
 S_0\oplus\{u_++Uu_+:u_+\in S_+\}.
 \label{eq:V31-neutral-graph}
\end{equation}
Thus a zero-flux allowed trace is generally a neutral graph pairing
incoming and outgoing modes; it is not the kernel of the flux operator.
The V26 compatible-graph theorem transports this domain.  The V28
protected-centre theorem applies separately only to the genuine operator
kernel \(S_0\).
\end{corollary}

\begin{proof}
Maximal isotropy implies \(S_0\subset\mathcal M\), because adjoining a
missing null vector would preserve isotropy.  Work in the quotient by
\(S_0\).  The projection of \(\mathcal M/S_0\) to \(S_+\) is injective:
a vector with zero positive component and nonzero negative component has
strictly negative quadratic flux.  The same argument applies to the
projection to \(S_-\).  A maximal isotropic subspace in balanced inertia
has dimension \(p\), so both projections are isomorphisms.  Hence it is
the graph of a map \(U:S_+\to S_-\).  Isotropy of two graph vectors gives
\(
\mathfrak b_n(u_+,v_+)+\mathfrak b_n(Uu_+,Uv_+)=0
\), which is exactly \eqref{eq:V31-neutral-isometry}.
\end{proof}

\begin{proposition}[Boundary-reservoir alternative]
\label{prop:V31-reservoir}
Let \(y\) be an independent boundary field with Green form
\(\mathfrak b_{\partial}\).  A conservative enlarged action may admit a
trace relation
\(\mathcal L\subset S_{\mathbb R}|_{\partial}\oplus Y\) provided
\begin{equation}
 \mathfrak b_n(u,v)
 +\mathfrak b_{\partial}(y,w)
 =0
 \quad
 \text{for all }(u,y),(v,w)\in\mathcal L.
 \label{eq:V31-reservoir-balance}
\end{equation}
Its bulk projection can have nonzero flux, but that flux is exactly
balanced by the boundary sector.  Eliminating \(y\) with a retarded
condition may produce an effective dissipative bulk boundary rule, but the
resulting rule is no longer the Euler domain of the local conservative
bulk action alone.
\end{proposition}

\begin{proof}
Apply the proof of Theorem~\ref{thm:V31-zero-flux} to the total action.
The antisymmetrized second variation is the sum of the bulk and boundary
Green forms, which must vanish on the total tangent trace relation.  This
is \eqref{eq:V31-reservoir-balance}.  Setting \(v=u,w=y\) shows that any
bulk flux is the negative of boundary flux.  A retarded elimination chooses
a time direction and replaces the reversible enlarged system by its
effective open subsystem, so it is not a local scalar variation of the
bulk field alone.
\end{proof}

\begin{firewall}[Maximal dissipation and action stationarity are separate]
\label{fire:V31-boundary-fork}
Well-posedness of a maximally dissipative Dirac boundary condition does not
prove that it is the boundary domain of the common scalar action.  The
action route must certify zero-flux isotropy, or it must print the boundary
reservoir and the balance law \eqref{eq:V31-reservoir-balance}.  Omitting
the reservoir while retaining a strict flux debit breaks common-action
incidence.
\end{firewall}

% -----------------------------------------------------------------------------
\subsection{Finite occurrence card}
\label{subsec:V31-finite}
% -----------------------------------------------------------------------------

At cofinal level \(q\), acquire the realified Hessian blocks
\[
 \widehat B_q^a
 =
 D^2_{z\chi_a}\widehat\ell_q,
 \qquad
 \widehat C_q^{ab}
 =
 D^2_{\chi_a\chi_b}\widehat\ell_q
\]
from the same protected action root as the V25 KKT card.  Store
\begin{align}
 r_{{\rm aff},q}
 &:=
 \max_{a,b}\|\widehat C_q^{ab}\|,
 \label{eq:V31-raff}\\
 r_{{\rm skew},q}
 &:=
 \max_a
 \|
 \widehat J_{{\rm A},q}^a
 -(\widehat B_q^a-(\widehat B_q^a)^*)
 \|,
 \label{eq:V31-rskew}\\
 r_{{\rm Dirac},q}
 &:=
 \max_a
 \|
 \widehat A_{D,q}^a
 -\widehat{\mathfrak j}_{D,q}
   (\widehat B_q^a-(\widehat B_q^a)^*)
 \|,
 \label{eq:V31-rDirac}\\
 r_{D,0,q}
 &:=
 \|\widehat E_{D,q}-\widehat{\mathcal D}_q\widehat z_q\|,
 \label{eq:V31-rD0}\\
 j_{*,q}^{\rm cert}
 &:=
 s_{\min}(\widehat{\mathfrak j}_{D,q})
 -\varepsilon_{j,q},
 \label{eq:V31-jfloor}\\
 r_{{\rm KKT}\to E,q}
 &:=
 \|
 \widehat E_{{\rm A},q}
 -\widehat K_{z,q}
 -\nabla_{a,q}^\dagger\widehat K_{\chi_a,q}
 \|,
 \label{eq:V31-rKETE}\\
 r_{{\rm prol},q}
 &:=
 \max_a
 \|
 \widehat E_{D,a,q}
 -\nabla_{a,q}\widehat E_{D,q}
 -\widehat{\mathcal D}_q\widehat C_{a,q}
 \|.
 \label{eq:V31-rprol}
\end{align}
The action-to-causal symbol is acquired factor by factor.  At a frozen
covector \(\xi\), the two target relations are
\begin{align}
 \widehat e_D
 &=
 \widehat{\mathfrak j}_D
 \bigl(
  \widehat k_z+
  \sigma(\nabla^\dagger)(\xi)\widehat k_\chi
 \bigr),
 \label{eq:V31-symbol-factor1}\\
 \widehat e_{D,a}
 &=
 \sigma(\nabla_a)(\xi)\widehat e_D
 +\gamma(\xi)\widehat c_a.
 \label{eq:V31-symbol-factor2}
\end{align}
Their composition is quadratic in \(\xi\) on the \(K_\chi\) column.  A
finite test must retain that degree-two coefficient rather than fit a
zeroth-order matrix.

For the conservative boundary branch, let \(W_q\) be an orthonormal basis
of the allowed trace and acquire the action Green matrix \(G_{{\rm A},q}\)
and causal flux matrix \(F_{D,q}\).  Store
\begin{align}
 r_{{\rm Green},q}
 &:=
 \|G_{{\rm A},q}-\zeta_{D,q}F_{D,q}\|,
 \label{eq:V31-rGreen}\\
 r_{{\rm iso},q}
 &:=
 \|W_q^*G_{{\rm A},q}W_q\|,
 \label{eq:V31-riso}\\
 r_{{\rm conormal},q}
 &:=
 \|
 (I-P_{\mathcal M_q^\circ})
 (\widehat p_q^n+D\widehat b_q)
 \|,
 \label{eq:V31-rconormal}\\
 r_{{\rm neutral},q}
 &:=
 \|W_q^*F_{D,q}W_q\|.
 \label{eq:V31-rneutral}
\end{align}
Maximality is certified by the V26/V28 dimension, inertia, and compatible
range tests.  For the reservoir branch, replace
\eqref{eq:V31-riso} by
\begin{equation}
 r_{{\rm reservoir},q}
 :=
 \left\|
 W_{{\rm tot},q}^*
 (G_{{\rm A},q}\oplus G_{\partial,q})
 W_{{\rm tot},q}
 \right\|.
 \label{eq:V31-rreservoir}
\end{equation}

\begin{theorem}[Cofinal fermion action-incidence handoff]
\label{thm:V31-cofinal-handoff}
Assume:
\begin{enumerate}
 \item the V25 KKT action floors and the V30 Cartan, jet, curvature, and
       causal time floors have positive uniform lower bounds;
 \item the residuals
       \eqref{eq:V31-raff}--\eqref{eq:V31-rprol} tend uniformly to zero,
       \(\inf_qj_{*,q}^{\rm cert}>0\), and the shifted graph constants in
       \eqref{eq:V31-Q-bound}--\eqref{eq:V31-Q-inverse-bound} are uniform;
 \item the finite symbols satisfy both factored relations
       \eqref{eq:V31-symbol-factor1}--\eqref{eq:V31-symbol-factor2} on a
       cofinal directional net with summable extraction errors;
 \item either the conservative boundary residuals
       \eqref{eq:V31-rGreen}--\eqref{eq:V31-rneutral} vanish and the V26
       compatible neutral-graph certificate passes; any genuine operator
       kernel must also pass V28 protected-centre separately; or a populated
       boundary
       reservoir satisfies \eqref{eq:V31-rreservoir} and has its own
       positive causal time floor.
\end{enumerate}
Then the fermion KKT action roots, reduced action Euler roots, causal Dirac
roots, and V30 prolonged roots are incident on one compatible graph.  The
fermion characteristic cone is the metric null cone.  On the conservative
boundary branch the admitted fermion trace is a zero-flux neutral
incoming/outgoing graph; on the reservoir branch its nonzero bulk flux is
balanced by the occurring boundary sector.
\end{theorem}

\begin{proof}
The V25 KKT congruence and
Corollary~\ref{cor:V31-no-Legendre} give the action Hessian inverse without
a fermion velocity Legendre inverse.  Vanishing affine, skew, Dirac, and full-row residuals passes
\eqref{eq:V31-action-Dirac-coefficients} and
\eqref{eq:V31-full-Dirac-residual} to the limit.
The positive \(\mathfrak j_D\) inverse floor gives characteristic incidence
by Corollary~\ref{cor:V31-characteristic-incidence}.

The two vanishing differential residuals, their factored symbol
relations, and the uniform shifted bounds pass
Theorems~\ref{thm:V31-KKT-incidence} and
\ref{thm:V31-compatible-graph} to the cofinal limit.  Corollary
\ref{cor:V31-root-equivalence} then identifies all four root
descriptions.  V29 supplies the metric null cone.

On the first boundary branch, Green incidence, conormal incidence, and
vanishing isotropy residual apply
Theorem~\ref{thm:V31-zero-flux};
Corollary~\ref{cor:V31-neutral-graph} and the V26 compatible-flux theorem
supply the boundary-domain handoff, with V28 used separately for any true
kernel.  On the second branch,
Proposition~\ref{prop:V31-reservoir} and the vanishing total Green residual
give exact flux balance, while the additional time floor makes the enlarged
boundary sector causal.
\end{proof}

\begin{firewall}[V31 status]
\label{fire:V31-status}
V31 proves the exact action-to-Dirac-to-prolonged incidence and the
conservative boundary obstruction.  The attached files do not print
\(\widehat B_q^a\), \(\widehat{\mathfrak j}_{D,q}\), the factored impulse
columns, the action Green matrix on the V20 trace, or a boundary reservoir.
Therefore Theorem~\ref{thm:V31-cofinal-handoff} is not physically populated
and the full Standard Model--Einstein continuum is not claimed.
\end{firewall}

% -----------------------------------------------------------------------------
\subsection{V31 ledger and next upstream row}
\label{subsec:V31-status}
% -----------------------------------------------------------------------------

\paragraph{New conclusions proved in V31.}
Theorem~\ref{thm:V31-KKT-incidence} gives the exact fermion
KKT-to-Euler residual map.  Proposition~\ref{prop:V31-action-symbol}
derives the skew first-order action symbol, and
Theorem~\ref{thm:V31-action-Dirac} gives the coefficientwise causal
multiplier test.  Theorems~\ref{thm:V31-compatible-graph} and
\ref{thm:V31-shifted-bounds} compose this with the V30 prolongation on the
correct shifted compatible graph.  Theorem~\ref{thm:V31-zero-flux}
proves that a conservative common action selects a zero-flux fermion trace,
while Proposition~\ref{prop:V31-reservoir} identifies the required local
conservative route to a nonzero projected bulk flux.

\paragraph{Imported conclusions not repeated.}
The abstract KKT congruence, Hamilton--Pontryagin incidence, covariant
spinor-jet identity, Dirac cone, compatible boundary pullback, and protected
centre remain V25--V30 results.  V31 supplies their missing
fermion-specific composition and boundary compatibility test.

\paragraph{Authoritative physical status.}
The notes now distinguish three separate finite objects: the symmetric KKT
Hessian, the skew variational Dirac symbol, and the positive-time causal
Dirac symbol.  Their exact incidence is testable by
\eqref{eq:V31-rskew}--\eqref{eq:V31-rprol}.  No attached source contains
those measured blocks in one common frame, so their occurrence remains
open.

\paragraph{Next upstream row.}
The next iteration should decide the boundary fork from the actual
common-action data.  It should derive the finite second-variation Green
matrix of the Palatini--Standard-Model action, compare its null subspace
with the V20 fermion trace, and either certify the conservative zero-flux
neutral graph or construct an explicit boundary reservoir with a
positive Hamiltonian and exact flux balance.  Only after that decision is
it meaningful to combine the fermion sector with the full gauge,
gravitational, and subsidiary boundary bank.

% END APPENDED VERSION V31 MATERIAL


% =============================================================================
% START APPENDED VERSION V32 MATERIAL
% =============================================================================

\section{V32: Explicit Standard Model fermion Green form,
equivariant neutral graphs, and the Higgs boundary intertwiner}
\label{sec:V32}

\paragraph{Purpose and scope.}
V31 proves that a conservative fermion action requires a totally isotropic
allowed trace and that such a trace is a neutral graph between the positive
and negative normal-flux spaces.  It leaves the graph abstract.  V32
computes its projector, adds gauge equivariance, and applies the resulting
representation criterion to the chiral Standard Model.  The conclusion is
a precise fork.  A constant bag-type graph that pairs the unbroken
left- and right-handed Standard Model multiplets is forbidden by their
inequivalent gauge representations.  A nonzero Higgs field and full-rank
Yukawa maps can supply a covariant intertwiner by polar decomposition.
The minimal model without a right-handed neutrino retains a rank defect in
the lepton block and therefore needs a different boundary choice or an
additional boundary/bulk degree of freedom.

No finite occurrence is inferred from the continuum formulas.  The last
subsection turns every formula into a same-frame packet certificate.

% -----------------------------------------------------------------------------
\subsection{Normalized Green matrix and all neutral graph projectors}
\label{subsec:V32-neutral}
% -----------------------------------------------------------------------------

Let \(F=H_DA_D(n)\) be the self-adjoint normal-flux matrix on the
realified or complex fermion trace fibre \(E\).  First remove the genuine
kernel \(E_0=\ker F\), which remains subject to the separate V28
protected-centre test.  On \(E_0^\perp\), let
\[
 E_+:=\operatorname{Ran}\mathbf 1_{(0,\infty)}(F),
 \qquad
 E_-:=\operatorname{Ran}\mathbf 1_{(-\infty,0)}(F).
\]
Define positive operators
\[
 F_+:=F|_{E_+},
 \qquad
 F_-:=-F|_{E_-}.
\]
The V26 flux normalization
\begin{equation}
 \mathcal N
 :=
 F_+^{1/2}\oplus F_-^{1/2}
 \label{eq:V32-flux-normalization}
\end{equation}
converts the nonzero flux to
\begin{equation}
 F_{\rm nor}
 =
 \mathcal N^{-*}F\mathcal N^{-1}
 =
 \begin{pmatrix}I&0\\0&-I\end{pmatrix}.
 \label{eq:V32-normal-flux}
\end{equation}

\begin{theorem}[Classification of maximal neutral graphs]
\label{thm:V32-neutral-classification}
Assume
\(\dim E_+=\dim E_-=m\).  A subspace
\(\mathcal M\subset E_+\oplus E_-\) is maximal totally isotropic for
\(F_{\rm nor}\) if and only if there is a unitary map
\(U:E_+\to E_-\) such that
\begin{equation}
 \mathcal M=\mathcal M_U
 :=
 \{u\oplus Uu:u\in E_+\}.
 \label{eq:V32-neutral-graph}
\end{equation}
The ordinary Hilbert-orthogonal projector onto this graph is
\begin{equation}
 \boxed{
 \Pi_U
 =
 \frac12
 \begin{pmatrix}
  I&U^*\\
  U&I
 \end{pmatrix}.}
 \label{eq:V32-graph-projector}
\end{equation}
Moreover,
\begin{align}
 \Pi_U^*&=\Pi_U,
 &
 \Pi_U^2&=\Pi_U,
 \label{eq:V32-projector-identities}\\
 W_U^*F_{\rm nor}W_U&=0,
 &
 W_U&:=2^{-1/2}\binom{I}{U}.
 \label{eq:V32-neutral-basis}
\end{align}
\end{theorem}

\begin{proof}
If \(U\) is unitary, then for \(u,v\in E_+\),
\[
 \left\langle
  u\oplus Uu,F_{\rm nor}(v\oplus Uv)
 \right\rangle
 =
 \langle u,v\rangle-\langle Uu,Uv\rangle=0.
\]
The graph has dimension \(m\), the maximal possible isotropic dimension for
inertia \((m,0,m)\), so it is maximal.

Conversely, let \(\mathcal M\) be maximal isotropic.  Projection
\(\pi_+:\mathcal M\to E_+\) is injective, since
\((0,w)\in\mathcal M\) would have flux \(-\|w\|^2\).  Maximality gives
\(\dim\mathcal M=m\), so \(\pi_+\) is bijective.  Hence
\(\mathcal M\) is the graph of \(U=\pi_-\pi_+^{-1}\).  Isotropy gives
\(\langle Uu,Uv\rangle=\langle u,v\rangle\), and equality of dimensions
makes \(U\) unitary.

Direct multiplication using \(U^*U=UU^*=I\) proves
\eqref{eq:V32-projector-identities}.  The last identity follows from
\[
 W_U^*F_{\rm nor}W_U
 =
 \tfrac12(I-U^*U)=0.
\]
\end{proof}

\begin{corollary}[Projector derivative]
\label{cor:V32-projector-derivative}
For a \(C^1\) unitary graph \(U(x)\),
\begin{equation}
 \nabla_\tau\Pi_U
 =
 \frac12
 \begin{pmatrix}
  0&(\nabla_\tau U)^*\\
  \nabla_\tau U&0
 \end{pmatrix},
 \qquad
 \|\nabla_\tau\Pi_U\|
 =
 \frac12\|\nabla_\tau U\|.
 \label{eq:V32-projector-derivative}
\end{equation}
Thus the V30 differentiated-boundary debit is controlled directly by the
derivative of the neutral intertwiner.
\end{corollary}

\begin{proof}
Differentiate \eqref{eq:V32-graph-projector}.  A block matrix
\(\left(\begin{smallmatrix}0&A^*\\A&0\end{smallmatrix}\right)\)
has norm \(\|A\|\), as follows by squaring it.
\end{proof}

\begin{proposition}[Conormal action condition on a neutral graph]
\label{prop:V32-conormal}
Suppose the reduced action boundary momentum is incident with the Green
matrix in the sense
\begin{equation}
 \langle p^n(z),v\rangle
 =
 \zeta_D\langle z,F_{\rm nor}v\rangle
 \label{eq:V32-momentum-Green}
\end{equation}
for a fixed nonzero normalization \(\zeta_D\).  If
\(z,v\in\mathcal M_U\), then the boundary first variation vanishes:
\[
 \langle p^n(z),v\rangle=0.
\]
Hence the essential neutral-graph trace satisfies the V31 conormal
condition with no additional scalar boundary potential.
\end{proposition}

\begin{proof}
Theorem~\ref{thm:V32-neutral-classification} gives
\(\langle z,F_{\rm nor}v\rangle=0\) for all graph vectors.  Substitute this
in \eqref{eq:V32-momentum-Green}.
\end{proof}

% -----------------------------------------------------------------------------
\subsection{Gauge-equivariant neutral graphs}
\label{subsec:V32-equivariant}
% -----------------------------------------------------------------------------

Let \(G\) be a compact internal gauge group acting unitarily by
\(\rho_\pm\) on \(E_\pm\).  Assume \(F_{\rm nor}\) commutes with this action.

\begin{theorem}[Equivariant neutral-graph criterion]
\label{thm:V32-equivariant-graph}
The neutral graph \(\mathcal M_U\) is \(G\)-invariant if and only if
\begin{equation}
 U\rho_+(g)=\rho_-(g)U
 \qquad(g\in G).
 \label{eq:V32-intertwiner}
\end{equation}
An equivariant maximal neutral graph exists if and only if the unitary
representations \((E_+,\rho_+)\) and \((E_-,\rho_-)\) are isomorphic.
Equivalently, every irreducible \(G\)-representation occurs with the same
multiplicity in \(E_+\) and \(E_-\).
\end{theorem}

\begin{proof}
The image of \(u\oplus Uu\) under \(g\) is
\(\rho_+(g)u\oplus\rho_-(g)Uu\).  It belongs to the same graph for every
\(u\) precisely when its second component equals
\(U\rho_+(g)u\), which is \eqref{eq:V32-intertwiner}.  Thus an invariant
graph is the graph of a unitary intertwiner.  Conversely, every unitary
intertwiner gives an invariant graph.  Complete reducibility of a finite
dimensional unitary representation and Schur's lemma identify existence of
an invertible intertwiner with equality of all irreducible
multiplicities.  Its polar factor is a unitary intertwiner.
\end{proof}

\begin{corollary}[Infinitesimal certificate]
\label{cor:V32-infinitesimal}
If \(G\) is connected with Lie algebra representation
\(\rho_{\pm,*}\), condition \eqref{eq:V32-intertwiner} is equivalent to
\begin{equation}
 U\rho_{+,*}(X)-\rho_{-,*}(X)U=0
 \qquad(X\in\mathfrak g).
 \label{eq:V32-infinitesimal-intertwiner}
\end{equation}
It is enough to test a basis of \(\mathfrak g\).
\end{corollary}

\begin{proof}
Differentiate \eqref{eq:V32-intertwiner} along
\(\exp(tX)\).  Conversely, the infinitesimal equality exponentiates on the
identity component, which is all of a connected group.
\end{proof}

\paragraph{Boundary-covariant chiral bag ansatz.}
On a timelike boundary, restriction of a four-dimensional left or right
Weyl spinor gives the same complex irreducible boundary spin module
\(\mathscr S_\partial\), up to a fixed normal-Clifford identification.
Consider a local boundary graph whose spin part is this fixed
\(\operatorname{Spin}(1,2)\)-intertwiner and whose internal part is
\begin{equation}
 U_{\rm int}:R_R\longrightarrow R_L.
 \label{eq:V32-chiral-bag-map}
\end{equation}
This is the chiral bag-type ansatz considered below.  Other local
conditions using extra tangential structure are not ruled out by its
failure.

\begin{proposition}[No constant unbroken-Standard-Model chiral bag map]
\label{prop:V32-SM-no-constant}
For one generation, let
\begin{align}
 R_L
 &:=
 (\mathbf3,\mathbf2)_{1/6}
 \oplus(\mathbf1,\mathbf2)_{-1/2},
 \label{eq:V32-RL}\\
 R_R^{\rm min}
 &:=
 (\mathbf3,\mathbf1)_{2/3}
 \oplus(\mathbf3,\mathbf1)_{-1/3}
 \oplus(\mathbf1,\mathbf1)_{-1},
 \label{eq:V32-RR-min}\\
 R_R^\nu
 &:=
 R_R^{\rm min}\oplus(\mathbf1,\mathbf1)_0.
 \label{eq:V32-RR-nu}
\end{align}
Here the entries are representations of
\(SU(3)_c\times SU(2)_L\) with the subscript denoting hypercharge.
Neither \(R_R^{\rm min}\) nor \(R_R^\nu\) is isomorphic to \(R_L\).
Therefore no constant
\(SU(3)_c\times SU(2)_L\times U(1)_Y\)-equivariant invertible map
\eqref{eq:V32-chiral-bag-map} exists, for any number of generations.
\end{proposition}

\begin{proof}
The \(SU(2)_L\)-restriction of \(R_L\) is a direct sum of four copies of
the fundamental doublet when colour multiplicity is counted.  The
\(SU(2)_L\)-restriction of either right-handed space is a sum of trivial
representations.  Their characters, hence their irreducible
multiplicities, differ.  They cannot be isomorphic.  Adding generations
multiplies both multiplicity lists by the same positive integer and does
not remove the mismatch.  Theorem~\ref{thm:V32-equivariant-graph} then
rules out the constant equivariant chiral bag graph.
\end{proof}

\begin{firewall}[Scope of the representation obstruction]
\label{fire:V32-no-go-scope}
Proposition~\ref{prop:V32-SM-no-constant} concerns a local, boundary-Lorentz
covariant bag graph that pairs the complete unbroken left and right chiral
multiplets by a constant internal map.  It does not rule out a boundary
condition with a Higgs-dependent intertwiner, extra tangential structure,
charge-conjugate doubling when the representation permits it, or an
explicit boundary reservoir.  Those alternatives must be printed rather
than silently identified with the constant bag ansatz.
\end{firewall}

% -----------------------------------------------------------------------------
\subsection{The Higgs--Yukawa intertwiner}
\label{subsec:V32-Higgs}
% -----------------------------------------------------------------------------

Let \(N_g\) be the number of generations.  Let the Higgs doublet have
\(\Phi\in(\mathbf1,\mathbf2)_{1/2}\), and put
\(\widetilde\Phi:=i\sigma_2\overline\Phi\), which transforms as
\((\mathbf1,\mathbf2)_{-1/2}\).  Let
\(Y_u,Y_d,Y_e\) and, when present, \(Y_\nu\) be the generation-space
Yukawa matrices.  Define
\begin{align}
 \mathcal Y_{\Phi,q}
 &:R_u\oplus R_d\longrightarrow R_Q,
 &
 \mathcal Y_{\Phi,q}(u,d)
 &:=
 \widetilde\Phi\otimes Y_u u
 +\Phi\otimes Y_d d,
 \label{eq:V32-Yq}\\
 \mathcal Y_{\Phi,\ell}
 &:R_\nu\oplus R_e\longrightarrow R_L^\ell,
 &
 \mathcal Y_{\Phi,\ell}(\nu,e)
 &:=
 \widetilde\Phi\otimes Y_\nu\nu
 +\Phi\otimes Y_e e.
 \label{eq:V32-Yl}
\end{align}
The colour identity is implicit in the quark block.  Set
\begin{equation}
 \mathcal Y_\Phi
 :=
 \mathcal Y_{\Phi,q}\oplus\mathcal Y_{\Phi,\ell}.
 \label{eq:V32-Yfull}
\end{equation}

\begin{lemma}[Orthogonal Higgs frame]
\label{lem:V32-Higgs-frame}
For every \(\Phi\in\mathbb C^2\),
\begin{equation}
 \langle\widetilde\Phi,\Phi\rangle=0,
 \qquad
 \|\widetilde\Phi\|=\|\Phi\|.
 \label{eq:V32-Higgs-orthogonal}
\end{equation}
If \(\Phi\ne0\), the matrix
\begin{equation}
 Q_\Phi
 :=
 \|\Phi\|^{-1}
 \begin{pmatrix}\widetilde\Phi&\Phi\end{pmatrix}
 \label{eq:V32-Higgs-unitary}
\end{equation}
is unitary.
\end{lemma}

\begin{proof}
Write \(\Phi=(a,b)^T\).  Then
\(\widetilde\Phi=(\overline b,-\overline a)^T\), up to the fixed harmless
sign convention.  Direct multiplication gives zero inner product and
equal squared norm \(|a|^2+|b|^2\).  The two normalized columns are
orthonormal.
\end{proof}

\begin{theorem}[Singular values and the Higgs--Yukawa floor]
\label{thm:V32-Yukawa-floor}
Assume the right-handed neutrino block is present.  The singular values of
\(\mathcal Y_\Phi\) are
\begin{equation}
 \|\Phi\|\,
 \left(
  s(Y_u)\cup s(Y_d)\cup s(Y_\nu)\cup s(Y_e)
 \right),
 \label{eq:V32-Yukawa-singular-values}
\end{equation}
with the quark values repeated by colour.  Hence
\begin{equation}
 s_{\min}(\mathcal Y_\Phi)
 =
 \|\Phi\|\,y_*,
 \qquad
 y_*:=
 \min_{f\in\{u,d,\nu,e\}}s_{\min}(Y_f).
 \label{eq:V32-Yukawa-floor}
\end{equation}
In particular, \(\mathcal Y_\Phi\) is invertible exactly when
\(\Phi\ne0\) and all four Yukawa matrices are invertible.
\end{theorem}

\begin{proof}
By Lemma~\ref{lem:V32-Higgs-frame}, multiplication by
\((\widetilde\Phi,\Phi)\) is \(\|\Phi\|\) times a unitary map from the
up/down or neutrino/charged-lepton two-summand space to the weak doublet.
Thus \(\mathcal Y_{\Phi,q}\) is unitarily equivalent to
\(\|\Phi\|\operatorname{diag}(Y_u,Y_d)\), with a colour identity, and
\(\mathcal Y_{\Phi,\ell}\) is unitarily equivalent to
\(\|\Phi\|\operatorname{diag}(Y_\nu,Y_e)\).  Singular values are invariant
under unitary multiplication and combine under direct sum.  This proves
\eqref{eq:V32-Yukawa-singular-values} and
\eqref{eq:V32-Yukawa-floor}.
\end{proof}

\begin{corollary}[Minimal-model lepton rank defect]
\label{cor:V32-minimal-rank}
Without a right-handed neutrino, the lepton map has domain dimension
\(N_g\) and target dimension \(2N_g\), and
\begin{equation}
 \operatorname{rank}\mathcal Y_{\Phi,\ell}\le N_g<2N_g.
 \label{eq:V32-minimal-rank}
\end{equation}
Therefore it cannot supply an invertible neutral-graph intertwiner for the
complete lepton doublet, even when \(\Phi\ne0\) and \(Y_e\) is invertible.
\end{corollary}

\begin{proof}
Only the charged-lepton summand is present in the domain.  Rank is at most
the domain dimension.  The strict dimension inequality proves
noninvertibility.
\end{proof}

Assume now the floor
\begin{equation}
 \eta_H:=\|\Phi\|y_*>0.
 \label{eq:V32-etaH}
\end{equation}
Define the left polar factor
\begin{equation}
 U_\Phi
 :=
 \mathcal Y_\Phi
 (\mathcal Y_\Phi^*\mathcal Y_\Phi)^{-1/2}.
 \label{eq:V32-polar-unitary}
\end{equation}

\begin{theorem}[Gauge-covariant Higgs polar intertwiner]
\label{thm:V32-Higgs-polar}
Under \eqref{eq:V32-etaH}, \(U_\Phi:R_R^\nu\to R_L\) is unitary.  If
\(g\in SU(3)_c\times SU(2)_L\times U(1)_Y\), then
\begin{align}
 \mathcal Y_{g\Phi}
 &=
 \rho_L(g)\mathcal Y_\Phi\rho_R(g)^{-1},
 \label{eq:V32-Y-covariance}\\
 U_{g\Phi}
 &=
 \rho_L(g)U_\Phi\rho_R(g)^{-1}.
 \label{eq:V32-U-covariance}
\end{align}
Thus the Higgs-dependent graph transforms covariantly even though no
constant unbroken-gauge intertwiner exists.
\end{theorem}

\begin{proof}
Unitarity follows directly:
\[
 U_\Phi^*U_\Phi
 =
 (\mathcal Y_\Phi^*\mathcal Y_\Phi)^{-1/2}
 \mathcal Y_\Phi^*\mathcal Y_\Phi
 (\mathcal Y_\Phi^*\mathcal Y_\Phi)^{-1/2}
 =I.
\]
The domain and target have equal dimension, so \(U_\Phi U_\Phi^*=I\).
Each term in \eqref{eq:V32-Yq}--\eqref{eq:V32-Yl} has matching colour,
weak-isospin, and hypercharge; this gives
\eqref{eq:V32-Y-covariance}.  Since the representations are unitary,
\[
 \mathcal Y_{g\Phi}^*\mathcal Y_{g\Phi}
 =
 \rho_R(g)
 (\mathcal Y_\Phi^*\mathcal Y_\Phi)
 \rho_R(g)^{-1}.
\]
Uniqueness and covariance of the positive inverse square root then give
\eqref{eq:V32-U-covariance}.
\end{proof}

\begin{proposition}[Derivative bound for the Higgs graph]
\label{prop:V32-polar-derivative}
Let \(Y=\mathcal Y_\Phi\), assume
\(s_{\min}(Y)\ge\eta_H>0\), and let \(\dot Y\) be any covariant
tangential derivative.  Then
\begin{equation}
 \|\dot U_\Phi\|
 \le
 \left(
  \eta_H^{-1}
  +\|Y\|^2\eta_H^{-3}
 \right)\|\dot Y\|.
 \label{eq:V32-polar-derivative}
\end{equation}
Consequently the graph projector obeys
\begin{equation}
 \|\nabla_\tau\Pi_{U_\Phi}\|
 \le
 \frac12
 \left(
  \eta_H^{-1}
  +\|\mathcal Y_\Phi\|^2\eta_H^{-3}
 \right)
 \|\nabla_\tau\mathcal Y_\Phi\|.
 \label{eq:V32-Higgs-projector-bound}
\end{equation}
\end{proposition}

\begin{proof}
Put \(A=Y^*Y\), so \(A\ge\eta_H^2I\).  The functional-calculus formula
\[
 A^{-1/2}
 =
 \frac1\pi\int_0^\infty
 t^{-1/2}(A+tI)^{-1}\,\mathrm dt
\]
gives
\[
 \|D(A^{-1/2})[\dot A]\|
 \le
 \frac12\eta_H^{-3}\|\dot A\|.
\]
Indeed, differentiate the resolvent under the integral and evaluate
\(\pi^{-1}\int_0^\infty
t^{-1/2}(\eta_H^2+t)^{-2}\,\mathrm dt=(2\eta_H^3)^{-1}\).
Since
\(\|\dot A\|\le2\|Y\|\|\dot Y\|\), differentiating
\(U=YA^{-1/2}\) gives
\[
 \|\dot U\|
 \le
 \eta_H^{-1}\|\dot Y\|
 +\|Y\|\frac12\eta_H^{-3}
   2\|Y\|\|\dot Y\|,
\]
which is \eqref{eq:V32-polar-derivative}.  Apply
Corollary~\ref{cor:V32-projector-derivative}.
\end{proof}

\begin{firewall}[Higgs zeros and zero Yukawa singular values are domain changes]
\label{fire:V32-Higgs-floor}
The polar graph \eqref{eq:V32-polar-unitary} is undefined when the Higgs
field vanishes or a required Yukawa singular value is zero.  Near that
locus the derivative bound blows up and the allowed trace need not have
constant rank.  One cannot use a broken-phase Higgs graph across an
unbroken-phase boundary segment without a separate chart, a partial
isometry plus complementary boundary data, or additional fields.
\end{firewall}

% -----------------------------------------------------------------------------
\subsection{From the internal polar map to the full fermion trace}
\label{subsec:V32-full-trace}
% -----------------------------------------------------------------------------

Let \(C_{\rm sp}:E_+^{\rm sp}\to E_-^{\rm sp}\) be the normalized
spin-boundary isometry supplied by the chosen local covariant bag chart.
On the branch with \(\eta_H>0\) define
\begin{equation}
 U_{\rm SM}(\Phi)
 :=
 C_{\rm sp}\otimes U_\Phi.
 \label{eq:V32-full-U}
\end{equation}

\begin{theorem}[Higgs-populated conservative fermion trace]
\label{thm:V32-Higgs-trace}
Assume:
\begin{enumerate}
 \item the normalized chiral bag chart factors as
       \(E_+=E_+^{\rm sp}\otimes R_R^\nu\) and
       \(E_-=E_-^{\rm sp}\otimes R_L\);
 \item the right-handed representation contains the neutrino block and all
       four generation-space Yukawa matrices have positive singular floors;
 \item \(\|\Phi\|\) has a positive boundary floor;
 \item \(C_{\rm sp}\) is unitary in the V26 normalized flux metric and is
       covariant under changes of the boundary spin frame;
 \item the V31 Green/causal normalization
       \eqref{eq:V31-Green-flux-incidence} holds.
\end{enumerate}
Then \(U_{\rm SM}(\Phi)\) is a gauge- and spin-covariant unitary neutral
intertwiner.  The projector
\begin{equation}
 \Pi_{\rm SM}(\Phi)
 =
 \frac12
 \begin{pmatrix}
  I&U_{\rm SM}(\Phi)^*\\
  U_{\rm SM}(\Phi)&I
 \end{pmatrix}
 \label{eq:V32-SM-projector}
\end{equation}
defines a constant-rank, maximal zero-flux trace.  Its spinor action
boundary variation vanishes for allowed variations at fixed \(\Phi\).
For general coupled variations the remaining term is exactly the mixed
Higgs covector of Proposition~\ref{prop:V32-moving-domain}.  Its
differentiated spacetime trace is controlled by
\eqref{eq:V32-Higgs-projector-bound} together with
\(\nabla C_{\rm sp}\).
\end{theorem}

\begin{proof}
Theorem~\ref{thm:V32-Higgs-polar} gives a gauge-covariant unitary internal
map.  Tensoring it with the covariant unitary spin map gives a unitary
intertwiner between the full normalized positive and negative flux spaces.
Theorem~\ref{thm:V32-neutral-classification} proves maximal isotropy and
\eqref{eq:V32-SM-projector}.  Proposition~\ref{prop:V32-conormal} removes
the spinor action boundary variation at fixed Higgs.  Differentiating the
tensor product and
using
\[
 \|\nabla(C_{\rm sp}\otimes U_\Phi)\|
 \le
 \|\nabla C_{\rm sp}\|
 +\|\nabla U_\Phi\|
\]
for unitary factors, followed by
Corollary~\ref{cor:V32-projector-derivative} and
Proposition~\ref{prop:V32-polar-derivative}, gives the derivative claim.
\end{proof}

\begin{proposition}[Mixed Higgs covector of the moving trace domain]
\label{prop:V32-moving-domain}
Let \(\Pi(\Phi)=\Pi_{\rm SM}(\Phi)\), let \(\psi\) obey
\((I-\Pi(\Phi))\psi=0\), and let \(\mathfrak L_\Phi\) be a right inverse
of \(I-\Pi(\Phi)\) on its range.  A tangent variation satisfies
\begin{equation}
 (I-\Pi)\delta\psi
 =
 (D_\Phi\Pi[\delta\Phi])\psi.
 \label{eq:V32-moving-domain-variation}
\end{equation}
Hence it has the decomposition
\begin{equation}
 \delta\psi
 =
 \delta\psi_\parallel
 +\mathfrak L_\Phi
   (D_\Phi\Pi[\delta\Phi])\psi,
 \qquad
 \delta\psi_\parallel\in\operatorname{Ran}\Pi.
 \label{eq:V32-moving-domain-split}
\end{equation}
The neutral part contributes no Dirac action boundary variation.  The
remaining variation is the Higgs covector
\begin{equation}
 \mathcal J_{\Phi,\partial}[\delta\Phi]
 :=
 \left\langle
  p_\psi^n,
  \mathfrak L_\Phi(D_\Phi\Pi[\delta\Phi])\psi
 \right\rangle.
 \label{eq:V32-mixed-Higgs-covector}
\end{equation}
If the Higgs boundary momentum and scalar boundary density contribute
\(p_\Phi^n+D_\Phi b\), the full action is stationary on an allowed Higgs
tangent space \(T_\Phi\mathcal H_\partial\) if and only if
\begin{equation}
 \bigl(p_\Phi^n+D_\Phi b+\mathcal J_{\Phi,\partial}\bigr)
 \big|_{T_\Phi\mathcal H_\partial}=0.
 \label{eq:V32-mixed-conormal}
\end{equation}
\end{proposition}

\begin{proof}
Differentiate \((I-\Pi(\Phi))\psi=0\) to obtain
\eqref{eq:V32-moving-domain-variation}.  Subtracting the displayed right
inverse lift gives \eqref{eq:V32-moving-domain-split}.  The first summand
lies in the neutral graph, so Proposition~\ref{prop:V32-conormal}
annihilates it.  Substitution of the second summand in the Dirac boundary
variation gives \eqref{eq:V32-mixed-Higgs-covector}.  Adding the Higgs and
boundary-density variations proves the necessary and sufficient conormal
condition \eqref{eq:V32-mixed-conormal}.
\end{proof}

\begin{firewall}[A field-dependent domain has mixed variations]
\label{fire:V32-moving-domain}
Gauge covariance of \(\Pi_{\rm SM}(\Phi)\) does not by itself make the
coupled action differentiable.  The finite history must also populate
\eqref{eq:V32-mixed-conormal}.  Treating \(\Phi\) as fixed in the spinor
boundary calculation while varying it in the Higgs action omits a real
boundary row.
\end{firewall}

\begin{corollary}[The minimal Standard Model cannot use the full Higgs graph]
\label{cor:V32-minimal-fork}
For the minimal field content without a right-handed neutrino,
Theorem~\ref{thm:V32-Higgs-trace} fails its first hypothesis by the rank
defect \eqref{eq:V32-minimal-rank}.  A full conservative boundary
certificate must therefore choose at least one of the following and print
it explicitly:
\begin{enumerate}
 \item add a right-handed neutrino or another representation completing the
       lepton intertwiner;
 \item give the unpaired neutrino sector a separate gauge- and
       boundary-covariant neutral trace using additional boundary
       structure;
 \item couple the rank-defect sector to a boundary reservoir satisfying
       \eqref{eq:V31-reservoir-balance};
 \item restrict the physical history to a boundaryless or purely
       characteristic setting for that sector.
\end{enumerate}
\end{corollary}

\begin{proof}
Corollary~\ref{cor:V32-minimal-rank} rules out an invertible polar factor on
the complete lepton trace.  The listed alternatives exhaust the immediate
ways to restore equal trace rank, provide the missing neutral partner, move
the flux into an explicit enlarged sector, or remove the timelike boundary
domain from the problem.
\end{proof}

% -----------------------------------------------------------------------------
\subsection{Finite common-action and representation certificate}
\label{subsec:V32-finite}
% -----------------------------------------------------------------------------

At cofinal level \(q\), first normalize the occurring fermion flux by its
certified positive and negative square roots.  Acquire
\[
 \widehat F_{{\rm nor},q},\quad
 \widehat C_{{\rm sp},q},\quad
 \widehat U_{\Phi,q},\quad
 \widehat U_{{\rm SM},q},\quad
 \widehat\Pi_q,\quad
 \widehat\Phi_q,\quad
 \widehat Y_{f,q},\quad
 \widehat{\mathcal Y}_q
\]
in the same V24 field, spin, gauge, and equation frame.  Store
\begin{align}
 r_{F,q}
 &:=
 \left\|
 \widehat F_{{\rm nor},q}
 -
 \begin{pmatrix}I&0\\0&-I\end{pmatrix}
 \right\|,
 \label{eq:V32-rF}\\
 r_{U,q}
 &:=
 \|\widehat U_{{\rm SM},q}^*\widehat U_{{\rm SM},q}-I\|
 +\|\widehat U_{{\rm SM},q}\widehat U_{{\rm SM},q}^*-I\|,
 \label{eq:V32-rU}\\
 r_{\Pi,q}
 &:=
 \left\|
 \widehat\Pi_q
 -\frac12
 \begin{pmatrix}
  I&\widehat U_{{\rm SM},q}^*\\
  \widehat U_{{\rm SM},q}&I
 \end{pmatrix}
 \right\|,
 \label{eq:V32-rPi}\\
 r_{{\rm neutral},q}
 &:=
 \|\widehat W_q^*
       \widehat F_{{\rm nor},q}\widehat W_q\|,
 \qquad
 \widehat W_q:=2^{-1/2}\binom I{\widehat U_{{\rm SM},q}},
 \label{eq:V32-rneutral}\\
 r_{{\rm cov},q}
 &:=
 \max_{X\in\mathcal X_{\mathfrak g}}
 \left\|
 D_\Phi\widehat U_{{\rm SM},q}
  [\rho_{\Phi,q}(X)\widehat\Phi_q]
 -\rho_{-,q}(X)\widehat U_{{\rm SM},q}
 +\widehat U_{{\rm SM},q}\rho_{+,q}(X)
 \right\|,
 \label{eq:V32-rcov}\\
 r_{\mathcal Y,q}
 &:=
 \|\widehat{\mathcal Y}_q-
       \mathcal Y_{\widehat\Phi_q}
       (\widehat Y_{u,q},\widehat Y_{d,q},
        \widehat Y_{\nu,q},\widehat Y_{e,q})\|,
 \label{eq:V32-rY}\\
 r_{{\rm tensor},q}
 &:=
 \|\widehat U_{{\rm SM},q}
    -\widehat C_{{\rm sp},q}\otimes\widehat U_{\Phi,q}\|,
 \label{eq:V32-rtensor}\\
 r_{{\rm polar},q}
 &:=
 \left\|
 \widehat U_{\Phi,q}-
 \widehat{\mathcal Y}_q
 (\widehat{\mathcal Y}_q^*\widehat{\mathcal Y}_q)^{-1/2}
 \right\|,
 \label{eq:V32-rpolar}\\
 r_{\nabla\Pi,q}
 &:=
 \max_\tau
 \left\|
 \nabla_{\tau,q}\widehat\Pi_q
 -\frac12
 \begin{pmatrix}
  0&(\nabla_{\tau,q}\widehat U_{{\rm SM},q})^*\\
  \nabla_{\tau,q}\widehat U_{{\rm SM},q}&0
 \end{pmatrix}
 \right\|,
 \label{eq:V32-rnablaPi}\\
 r_{{\rm domvar},q}
 &:=
 \sup_{\substack{(\delta\Phi,\delta\psi)\in\widehat{\mathcal T}_{\partial,q}\\ \|(\delta\Phi,\delta\psi)\|=1}}
 \|(I-\widehat\Pi_q)\delta\psi
 -(D_\Phi\widehat\Pi_q[\delta\Phi])\widehat\psi_q\|,
 \label{eq:V32-rdomvar}\\
 r_{{\rm mix},q}
 &:=
 \left\|
 (\widehat p_{\Phi,q}^n+D_\Phi\widehat b_q
  +\widehat{\mathcal J}_{\Phi,\partial,q})
 \big|_{T_{\widehat\Phi_q}\mathcal H_{\partial,q}}
 \right\|.
 \label{eq:V32-rmix}
\end{align}
Here \(\mathcal X_{\mathfrak g}\) is a fixed orthonormal basis of the
Standard Model gauge Lie algebra.  The derivative term in
\eqref{eq:V32-rcov} is essential: it distinguishes covariance of the
field-dependent Higgs graph from invariance of a constant graph.  The
representation card separately stores the irreducible multiplicity
vectors; their inequality is the exact failure witness for the constant
equivariant branch.

Define the floors
\begin{align}
 h_{\Phi,q}^{\rm cert}
 &:=
 \|\widehat\Phi_q\|-\varepsilon_{\Phi,q},
 \label{eq:V32-Higgs-floor}\\
 y_{*,q}^{\rm cert}
 &:=
 \min_f s_{\min}(\widehat Y_{f,q})-\varepsilon_{Y,q},
 \label{eq:V32-Y-floor}\\
 \eta_{H,q}^{\rm cert}
 &:=
 h_{\Phi,q}^{\rm cert}y_{*,q}^{\rm cert}
 -\varepsilon_{{\rm prod},q},
 \label{eq:V32-eta-cert}\\
 \beta_{{\rm graph},q}^{\rm cert}
 &:=
 \beta_{30,q}^{\rm cert}
 -C_\partial
 \left[
  r_{\nabla\Pi,q}
  +\frac12\|\nabla_\tau\widehat C_{{\rm sp},q}\|
  +\frac12
   \left(
    (\eta_{H,q}^{\rm cert})^{-1}
    +\|\widehat{\mathcal Y}_q\|^2
      (\eta_{H,q}^{\rm cert})^{-3}
   \right)
   \|\nabla_\tau\widehat{\mathcal Y}_q\|
 \right].
 \label{eq:V32-boundary-reserve}
\end{align}
The last expression is used only when
\(\eta_{H,q}^{\rm cert}>0\).

\begin{theorem}[Cofinal Higgs-boundary handoff]
\label{thm:V32-cofinal-handoff}
Assume:
\begin{enumerate}
 \item the V31 action/Dirac/Green incidence residuals vanish and their
       KKT, equation-multiplier, and causal time floors remain positive;
 \item the flux has balanced positive and negative ranks after removal of a
       separately certified V28 kernel;
 \item the residuals
       \eqref{eq:V32-rF}--\eqref{eq:V32-rmix} tend uniformly to zero,
       and the cofinal extraction tails are summable; the representation
       multiplicity card records the constant-graph failure while
       \eqref{eq:V32-rcov} certifies the moving Higgs graph;
 \item the field content contains the right-handed neutrino block,
       \(\inf_q\eta_{H,q}^{\rm cert}>0\), and
       \(\inf_q\beta_{{\rm graph},q}^{\rm cert}>0\);
 \item the V20 complementing count and the V26 compatible trace pullback
       pass for \(\widehat\Pi_q\), and the Higgs boundary tangent space in
       \eqref{eq:V32-rmix} has its required complementing and energy floor.
\end{enumerate}
Then the cofinal Standard Model fermion trace is the
gauge- and spin-covariant projector
\(\Pi_{\rm SM}(\Phi)\) of \eqref{eq:V32-SM-projector}.  It is
constant-rank and maximal neutral.  The mixed conormal residual makes it
incident with the full conservative fermion--Higgs action, and the
derivative reserve makes it compatible with the differentiated V30
spinor-jet boundary domain.  Together with V30--V31, this closes the
fermion bulk/action/trace chain on that Higgs-floor branch.
\end{theorem}

\begin{proof}
The positive Higgs--Yukawa floor and vanishing Yukawa residual apply
Theorems~\ref{thm:V32-Yukawa-floor} and
\ref{thm:V32-Higgs-polar}.  The unitary, projector, tensor, polar, covariance, and neutral residuals
pass Theorems~\ref{thm:V32-neutral-classification} and
\ref{thm:V32-Higgs-polar} to the limit.  The unequal representation
multiplicities retain Proposition~\ref{prop:V32-SM-no-constant} as the
failure certificate for the constant branch.  The V31 Green incidence and
Proposition~\ref{prop:V32-conormal} give action compatibility.  The
moving-domain identity and vanishing mixed conormal residual give the
full fermion--Higgs action boundary stationarity.  The projector derivative
identity, polar derivative bound, and positive reserve
\eqref{eq:V32-boundary-reserve} control the V30 differentiated trace.  V20
supplies the first-order complementing count and V26 transports
the compatible graph into the physical trace frame.  The already-certified
V30--V31 root and symbol incidence completes the stated chain.
\end{proof}

\begin{firewall}[V32 physical status]
\label{fire:V32-status}
The attached sources do not choose a right-handed-neutrino completion,
print a positive boundary Higgs floor, show full-rank finite Yukawa
matrices in the common packet frame, identify the V20 allowed trace with
\eqref{eq:V32-SM-projector}, or populate the mixed Higgs conormal row
\eqref{eq:V32-mixed-conormal}.  Theorem~\ref{thm:V32-cofinal-handoff} is
therefore a proved conditional branch, not a populated full
Standard Model--Einstein continuum.
\end{firewall}

% -----------------------------------------------------------------------------
\subsection{V32 ledger and next upstream row}
\label{subsec:V32-status}
% -----------------------------------------------------------------------------

\paragraph{New conclusions proved in V32.}
Theorem~\ref{thm:V32-neutral-classification} computes every maximal neutral
graph and its projector.  Theorem~\ref{thm:V32-equivariant-graph} turns
gauge invariance into equality of representation multiplicities.
Proposition~\ref{prop:V32-SM-no-constant} proves the constant unbroken
chiral-bag obstruction.  Theorems~\ref{thm:V32-Yukawa-floor} and
\ref{thm:V32-Higgs-polar} construct the Higgs-dependent covariant
alternative with an exact singular floor.  Proposition
\ref{prop:V32-polar-derivative} supplies the projector regularity needed by
V30.  Proposition~\ref{prop:V32-moving-domain} supplies the mixed Higgs
conormal row, while Corollary~\ref{cor:V32-minimal-rank} isolates the
minimal-model neutrino rank defect.

\paragraph{Imported conclusions not repeated.}
The abstract neutral-graph requirement, KKT/Dirac incidence, spinor-jet
prolongation, compatible trace transport, and first-order complementing
condition remain V20 and V25--V31 results.  V32 supplies their explicit
Standard Model representation and Higgs--Yukawa content.

\paragraph{Authoritative physical status.}
The fermion boundary problem has split into decidable branches.  A constant
unbroken-gauge chiral bag graph fails.  A Higgs polar graph passes only with
a nonvanishing boundary Higgs, full-rank Yukawa matrices, and a
right-handed-neutrino completion.  The minimal field content needs a
separate neutrino boundary construction or a reservoir.  None of these
choices is populated by the attached finite histories.

\paragraph{Next upstream row.}
The next iteration should join the selected fermion neutral graph to the
Yang--Mills Gauss, BRST, harmonic Einstein, Cartan, and spinor-jet
subsidiary boundary system.  It must compute the full incoming constraint
matrix and prove that the physical allowed trace is invariant under the
subsidiary flow.  If the minimal neutrino branch is retained, that note must
also construct its missing boundary partner before claiming a common
constraint-preserving graph.

% END APPENDED VERSION V32 MATERIAL

% =============================================================================
% START APPENDED VERSION V33 MATERIAL
% =============================================================================

\section{V33: Full incoming subsidiary matrix, neutral boundary repair,
and cofinal constraint preservation}
\label{sec:V33}

\paragraph{Purpose and relation to the preceding notes.}
V20 defined the induced constraint-boundary defect
\(\mathfrak D_C\) and proved that its vanishing is sufficient for
subsidiary constraint propagation.  V25 identified the Cartan--Standard-Model
defining complex and proved that its Bianchi and integrability rows are
compatible target identities, rather than independent equations.  V30--V32
then constructed the differentiated fermion domain and, on the stated
Higgs-floor branch, an action-compatible maximal neutral fermion graph.
What is still missing is the incidence calculation between that physical
boundary graph and \emph{all} incoming subsidiary rows.  This note supplies
the exact finite matrix which must be calculated, proves the range and
singular-floor criteria for repairing a boundary graph without leaving its
neutral/action-compatible chart, and gives a cofinal nonlinear theorem.
It does not presume that the attached histories contain the entries of that
matrix.

\subsection{Independent subsidiary target and the induced operator defect}
\label{subsec:V33-defect}

Fix a cofinal level \(q\), a boundary point, and a tangential frequency in
the compact frequency card used by V20.  All spaces below are the finite
realifications of the corresponding complex trace spaces.  Let
\(\mathcal A_q\) be the free physical boundary-amplitude space and let
\(\mathcal J_{\partial,q}\) be the boundary-jet space after the declared
first-order reduction.  A boundary chart is a \(C^2\) family
\begin{equation}
 T_q(z):\mathcal A_q\longrightarrow\mathcal J_{\partial,q},
 \qquad z\in B_r(0)\subset\mathcal Z_q,
 \label{eq:V33-trace-chart}
\end{equation}
whose range obeys the main-system boundary equations.  Normal derivatives
which occur in a constraint are eliminated using the same main equations;
changing that elimination changes \(T_q\) and must be recorded.

Choose an independent subsidiary family
\begin{equation}
 \mathfrak Q_q
 =
 \{\mathrm{Ein},\mathrm{YM},\mathrm{BRST},
   \mathrm{def},\mathrm{spin}\}_q .
 \label{eq:V33-independent-family}
\end{equation}
Here the entries mean, in the selected reduction, the harmonic/Einstein
constraint variables, the Yang--Mills gauge and Gauss variables, the
BRST gauge-fixing variables, the independent metric or
Cartan--Standard-Model defining defects, and the prolonged spinor-jet
defect.  The list is a block list, not a claim that every named component
is independent.  In particular the V25 rows
\[
 B_{\mathsf R},\quad B_{\mathsf T},\quad B_F,\quad
 B_\Pi,\quad B_\Xi
\]
are omitted after the compatible-range reduction
of Theorem~\ref{thm:V25-Cartan-identities}.  Hamiltonian, momentum, and
Gauss readouts are included only to the rank left after their exact
incidence identities have been imposed.

For \(a\in\mathfrak Q_q\), let
\[
 \mathcal C_{\partial,q}^a:
 \mathcal J_{\partial,q}\longrightarrow\mathcal S_q^a
\]
be the actual frozen boundary trace of the subsidiary variable.  Its
normalized boundary splitting is
\[
 P_{a,q}^++P_{a,q}^0+P_{a,q}^-=I,
 \qquad
 K_{a,q}:\operatorname{Ran}P_{a,q}^+
             \longrightarrow\operatorname{Ran}P_{a,q}^- ,
\]
with \(\|K_{a,q}\|\le\kappa_{a,q}<1\).  Define
\begin{align}
 Q_{a,q}
 &:=
 P_{a,q}^- -K_{a,q}P_{a,q}^+,
 \label{eq:V33-incoming-row}\\
 \mathcal S_{q}^-
 &:=
 \bigoplus_{a\in\mathfrak Q_q}\operatorname{Ran}P_{a,q}^-,
 \qquad
 Q_q:=\bigoplus_{a\in\mathfrak Q_q}Q_{a,q},
 \label{eq:V33-subsidiary-sum}\\
 \mathfrak D_q(z)
 &:=
 Q_q\mathcal C_{\partial,q}(z)T_q(z)
 \in
 \operatorname{Hom}_{\mathbb R}(\mathcal A_q,\mathcal S_q^-).
 \label{eq:V33-full-defect}
\end{align}
Equation \eqref{eq:V33-full-defect} is the full version of
\eqref{eq:V20-constraint-boundary-defect}.  Its vanishing is an operator
identity: it must hold for every free physical amplitude, rather than only
for the particular trace of one computed history.

\begin{proposition}[Exact derivative of the incoming defect]
\label{prop:V33-defect-derivative}
For every \(\dot z\in\mathcal Z_q\),
\begin{align}
 D\mathfrak D_q(z)[\dot z]
 ={}&
 (DQ_q(z)[\dot z])\mathcal C_{\partial,q}(z)T_q(z)
 \nonumber\\
 &+
 Q_q(z)(D\mathcal C_{\partial,q}(z)[\dot z])T_q(z)
 \nonumber\\
 &+
 Q_q(z)\mathcal C_{\partial,q}(z)DT_q(z)[\dot z].
 \label{eq:V33-defect-derivative}
\end{align}
If only the physical graph is varied while the frozen subsidiary splitting
and coefficients are fixed, the first two terms vanish.  They may not be
dropped when \(z\) also changes the background field, the normal, the
normalization, or the subsidiary contraction.
\end{proposition}

\begin{proof}
Differentiate the three-factor composition in
\eqref{eq:V33-full-defect}.  All spaces are finite-dimensional at level
\(q\), so the ordinary product rule applies without a domain qualification.
\end{proof}

\begin{firewall}[Compatible identities are not extra incoming rows]
\label{fire:V33-no-duplicate-rows}
Appending the V25 Bianchi and matter-integrability defects to
\(\mathfrak Q_q\) after retaining their defining defects inserts linearly
dependent rows into the target.  The resulting zero singular value is a
manufactured cokernel.  The full incoming matrix below is full only after
the compatible target has been quotiented by those exact identities.
\end{firewall}

\subsection{The neutral fermion chart and the moving Higgs column}
\label{subsec:V33-neutral-chart}

On the V32 Higgs-floor branch put
\[
 U_0(\Phi):=U_{\rm SM}(\Phi),\qquad
 W_0(\Phi):=2^{-1/2}\binom I{U_0(\Phi)}.
\]
Let \(\mathfrak u_+\) denote the real vector space of skew-Hermitian
endomorphisms of the normalized positive-flux space.  An exact neutral
chart through \(U_0\) is
\begin{equation}
 U(\Phi,\Omega)
 :=
 U_0(\Phi)\exp\Omega,
 \qquad \Omega\in\mathfrak u_+,
 \qquad
 W(\Phi,\Omega):=2^{-1/2}\binom I{U(\Phi,\Omega)}.
 \label{eq:V33-unitary-chart}
\end{equation}

\begin{lemma}[Neutrality is exact throughout the chart]
\label{lem:V33-neutral-chart}
For every \(\Omega^*=-\Omega\),
\begin{equation}
 U(\Phi,\Omega)^*U(\Phi,\Omega)=I,
 \qquad
 W(\Phi,\Omega)^*F_{\rm nor}W(\Phi,\Omega)=0.
 \label{eq:V33-exact-neutrality}
\end{equation}
At \(\Omega=0\), a Higgs/tangent variation \((h,\dot\Omega)\) gives
\begin{align}
 \dot U
 &=
 D_\Phi U_0(\Phi)[h]+U_0(\Phi)\dot\Omega,
 \label{eq:V33-U-tangent}\\
 U_0^*\dot U+\dot U^*U_0&=0,
 \label{eq:V33-unitary-tangent}\\
 \dot W
 &=2^{-1/2}
 \binom 0{D_\Phi U_0(\Phi)[h]+U_0(\Phi)\dot\Omega}.
 \label{eq:V33-W-tangent}
\end{align}
\end{lemma}

\begin{proof}
The exponential of a skew-Hermitian matrix is unitary, and \(U_0\) is
unitary by Theorem~\ref{thm:V32-Higgs-trace}.  This proves the first
identity.  Substitution into
\(F_{\rm nor}=\operatorname{diag}(I,-I)\) proves the second.  Differentiate
\(U_0(\Phi)\exp\Omega\) at \(\Omega=0\) to obtain
\eqref{eq:V33-U-tangent}.  Differentiating \(U_0^*U_0=I\), and using
\(\dot\Omega^*=-\dot\Omega\), gives
\eqref{eq:V33-unitary-tangent}.  The last formula is the derivative of
\(W\).
\end{proof}

Gauge covariance restricts the allowable \(\Omega\)-columns.  If the
correction is field-dependent, require
\begin{equation}
 \Omega(g\Phi)
 =
 \rho_+(g)\Omega(\Phi)\rho_+(g)^{-1}.
 \label{eq:V33-Omega-covariance}
\end{equation}
Then
\begin{equation}
 U(g\Phi,\Omega(g\Phi))
 =
 \rho_-(g)U(\Phi,\Omega(\Phi))\rho_+(g)^{-1}.
 \label{eq:V33-corrected-covariance}
\end{equation}
For a field-independent correction,
\eqref{eq:V33-Omega-covariance} reduces to the commutant condition
\([\Omega,\rho_+(g)]=0\).  Thus the parameter space used in the rank test
must be the covariant subspace
\(\mathfrak u_{+,{\rm cov}}\subset\mathfrak u_+\), not the full unitary
tangent space.

The boundary-jet lift in \eqref{eq:V33-trace-chart} must also contain the
V32 moving-domain map.  In amplitude coordinates
\((h,\chi_\parallel)\), its fermion component is
\begin{equation}
 T_{{\rm mov},q}(z)(h,\chi_\parallel)
 =
 \chi_\parallel+
 \mathfrak L_{\Phi,z}
 \bigl(D_\Phi\Pi_z[h]\bigr)\psi_q,
 \qquad
 \chi_\parallel\in\operatorname{Ran}\Pi_z.
 \label{eq:V33-moving-lift}
\end{equation}
Consequently the Higgs column of
\eqref{eq:V33-defect-derivative} contains
\begin{equation}
 Q_q\mathcal C_{\partial,q}\,
 \iota_\psi\mathfrak L_{\Phi,0}
 \bigl(D_\Phi\Pi_0[h]\bigr)\psi_q,
 \label{eq:V33-moving-constraint-column}
\end{equation}
in addition to the ordinary Higgs variation of the coefficients.
Here \(\iota_\psi\) inserts the spinor trace in the complete boundary jet.
This is the constraint counterpart of the mixed action covector
\eqref{eq:V32-mixed-Higgs-covector}.

Let \(\mathfrak A_{\partial,q}(z)\) collect the independent action conormal
residuals not already solved identically by the chart.  Its
fermion--Higgs component is
\begin{equation}
 \mathfrak A_{\Phi\psi,q}(z)
 :=
 \left.
 \bigl(
 p_{\Phi,q}^n+D_\Phi b_q+
 \mathcal J_{\Phi,\partial,q}(\Pi_z)
 \bigr)
 \right|_{T_{\Phi_q}\mathcal H_{\partial,q}}.
 \label{eq:V33-action-row}
\end{equation}
Define the combined boundary incidence residual
\begin{equation}
 \mathfrak B_q(z)
 :=
 \bigl(\mathfrak D_q(z),\mathfrak A_{\partial,q}(z)\bigr)
 \in\mathcal Y_q,
 \label{eq:V33-combined-residual}
\end{equation}
where
\[
 \mathcal Y_q
 :=
 \operatorname{Hom}_{\mathbb R}(\mathcal A_q,\mathcal S_q^-)
 \oplus\mathcal A_{\partial,q}^*.
\]
Using \(\mathfrak B_q\), rather than \(\mathfrak D_q\) alone, prevents a
constraint repair from destroying the common-action boundary condition.

\begin{firewall}[A free unitary correction can break gauge covariance]
\label{fire:V33-covariant-columns}
The skew-Hermitian condition on \(\Omega\) preserves zero flux, but it does
not imply \eqref{eq:V33-Omega-covariance}.  A singular floor computed with
noncovariant unitary columns is therefore not a physical repair margin.
The finite matrix must be restricted to covariant columns before its rank
or smallest singular value is reported.
\end{firewall}

\subsection{The full real matrix and its exact obstruction}
\label{subsec:V33-matrix}

Choose orthonormal real bases \(\{Z_\alpha\}\) of the admissible parameter
space \(\mathcal Z_q^{\rm ad}\) and \(\{E_\mu\}\) of \(\mathcal Y_q\).
Admissibility means that the chart already enforces the exact flux
condition, gauge covariance, the declared complementing count, and any
boundary equation solved by parametrization.  Put
\begin{align}
 (\mathbf d_q)_\mu
 &:=
 \langle E_\mu,\mathfrak B_q(0)\rangle_{\mathbb R},
 \label{eq:V33-defect-vector}\\
 (\mathbf L_q)_{\mu\alpha}
 &:=
 \left\langle
 E_\mu,D\mathfrak B_q(0)[Z_\alpha]
 \right\rangle_{\mathbb R}.
 \label{eq:V33-full-matrix}
\end{align}
For complex matrices the real inner product is
\(\langle X,Y\rangle_{\mathbb R}:=\operatorname{Re}\operatorname{tr}(X^*Y)\).
Thus \(\mathbf L_q\) is an ordinary real matrix, even for fermion and ghost
columns.

With row blocks
\[
 \mathrm{Ein},\ \mathrm{YM},\ \mathrm{BRST},\
 \mathrm{def},\ \mathrm{spin},\ \mathrm{act}
\]
and column blocks
\[
 g/e,\ \omega,\ A,\ \Phi,\ \psi,\ c,\bar c,b,
\]
the matrix has the schematic form
\begin{equation}
 \mathbf L_q=
 \begin{pmatrix}
 L_{\mathrm{Ein},g}&L_{\mathrm{Ein},\omega}
   &L_{\mathrm{Ein},A}&L_{\mathrm{Ein},\Phi}
   &L_{\mathrm{Ein},\psi}&L_{\mathrm{Ein},\mathrm{gh}}\\
 L_{\mathrm{YM},g}&L_{\mathrm{YM},\omega}
   &L_{\mathrm{YM},A}&L_{\mathrm{YM},\Phi}
   &L_{\mathrm{YM},\psi}&L_{\mathrm{YM},\mathrm{gh}}\\
 L_{\mathrm{BRST},g}&L_{\mathrm{BRST},\omega}
   &L_{\mathrm{BRST},A}&L_{\mathrm{BRST},\Phi}
   &L_{\mathrm{BRST},\psi}&L_{\mathrm{BRST},\mathrm{gh}}\\
 L_{\mathrm{def},g}&L_{\mathrm{def},\omega}
   &L_{\mathrm{def},A}&L_{\mathrm{def},\Phi}
   &L_{\mathrm{def},\psi}&L_{\mathrm{def},\mathrm{gh}}\\
 L_{\mathrm{spin},g}&L_{\mathrm{spin},\omega}
   &L_{\mathrm{spin},A}&L_{\mathrm{spin},\Phi}
   &L_{\mathrm{spin},\psi}&L_{\mathrm{spin},\mathrm{gh}}\\
 L_{\mathrm{act},g}&L_{\mathrm{act},\omega}
   &L_{\mathrm{act},A}&L_{\mathrm{act},\Phi}
   &L_{\mathrm{act},\psi}&L_{\mathrm{act},\mathrm{gh}}
 \end{pmatrix}_q .
 \label{eq:V33-block-matrix}
\end{equation}
The ghost column abbreviates \(c,\bar c,b\).  No off-diagonal block in
\eqref{eq:V33-block-matrix} is declared zero merely because it joins
different named sectors.  For example,
\eqref{eq:V33-moving-constraint-column} contributes to
\(L_{\mathrm{spin},\Phi}\), and variation of the stress/current
coefficients contributes to the Einstein and Yang--Mills mixed blocks.
Structural zeros require an identity in the common frame.

\begin{theorem}[Tangent obstruction, least correction, and dimension count]
\label{thm:V33-linear-repair}
Let \(L:\mathcal X\to\mathcal Y\) be the real linear map represented by
\(\mathbf L_q\), and let \(d=\mathfrak B_q(0)\).  The first-order
action-compatible neutral repair equation
\begin{equation}
 Lz=-d
 \label{eq:V33-linear-repair}
\end{equation}
is solvable if and only if
\begin{equation}
 P_{\ker L^*}d=0.
 \label{eq:V33-left-obstruction}
\end{equation}
When it is solvable, its unique minimum-norm solution is
\begin{equation}
 z_{\rm lin}=-L^\dagger d,
 \label{eq:V33-minimum-repair}
\end{equation}
where \(L^\dagger\) is the Moore--Penrose inverse.  If
\[
 \rho_L:=
 s_{\min}\!\left(
 L\big|_{(\ker L)^\perp}:
 (\ker L)^\perp\longrightarrow\operatorname{Ran}L
 \right)>0,
\]
then
\begin{equation}
 \|z_{\rm lin}\|\le\rho_L^{-1}\|d\|.
 \label{eq:V33-linear-bound}
\end{equation}
In particular, surjectivity requires
\begin{equation}
 \dim_{\mathbb R}\mathcal X
 \ge
 \dim_{\mathbb R}\mathcal Y.
 \label{eq:V33-dimension-count}
\end{equation}
\end{theorem}

\begin{proof}
In finite-dimensional real Hilbert spaces,
\(\operatorname{Ran}L=(\ker L^*)^\perp\).  Hence \(-d\) lies in the range
of \(L\) exactly when \eqref{eq:V33-left-obstruction} holds.  The
Moore--Penrose inverse maps a range vector to the unique solution
orthogonal to \(\ker L\), proving \eqref{eq:V33-minimum-repair}.  The
inverse norm on that orthogonal complement is \(\rho_L^{-1}\), which gives
\eqref{eq:V33-linear-bound}.  A surjective map cannot have a smaller
domain dimension than target dimension.
\end{proof}

\begin{corollary}[Explicit left-null failure witness]
\label{cor:V33-left-witness}
If \(y\in\ker L^*\) and
\(\langle y,d\rangle_{\mathbb R}\ne0\), no infinitesimal variation inside
the declared neutral, gauge-covariant, action-compatible boundary chart
can remove the incoming constraint defect.  Moreover
\begin{equation}
 \inf_z\|d+Lz\|
 =
 \|P_{\ker L^*}d\|.
 \label{eq:V33-best-linear-residual}
\end{equation}
\end{corollary}

\begin{proof}
Orthogonally decompose \(d=P_{\operatorname{Ran}L}d+
P_{\ker L^*}d\).  The first component can be cancelled and the second is
orthogonal to every \(Lz\).
\end{proof}

This witness has a direct physical use.  If the minimal Standard Model
branch omits the right-handed neutrino and no reservoir or separate
neutral partner is introduced, the admissible fermion column space is
smaller than the V32 full-Higgs chart.  Equation
\eqref{eq:V33-dimension-count} and the left-null vectors of the resulting
matrix decide whether the remaining bosonic and boundary-reservoir
columns can repair the loss.  The missing column may not be inserted
formally.

\subsection{Coupled Schur construction without sectorwise splicing}
\label{subsec:V33-Schur}

The full matrix can sometimes be certified through a block factorization.
The following result states precisely when such a calculation is valid.

\begin{theorem}[Triangular and small-coupling boundary inverses]
\label{thm:V33-block-inverse}
Suppose selected complements give decompositions
\(\mathcal X=\bigoplus_{j=1}^m\mathcal X_j\) and
\(\mathcal Y=\bigoplus_{j=1}^m\mathcal Y_j\), and write
\[
 L=D+N,\qquad D=\operatorname{diag}(D_1,\ldots,D_m),
\]
where every
\(D_j:\mathcal X_j\to\mathcal Y_j\) is bijective.

\begin{enumerate}
\item If \(N\) is strictly block lower triangular, then \(ND^{-1}\) is
      nilpotent of order at most \(m\), and
      \begin{equation}
       L^{-1}
       =
       D^{-1}\sum_{\ell=0}^{m-1}(-ND^{-1})^\ell.
       \label{eq:V33-triangular-inverse}
      \end{equation}
\item For an arbitrary off-diagonal coupling \(E\), if
      \begin{equation}
       \eta:=\|D^{-1}E\|<1,
       \label{eq:V33-block-smallness}
      \end{equation}
      then \(D+E\) is invertible and
      \begin{equation}
       \|(D+E)^{-1}\|
       \le\frac{\|D^{-1}\|}{1-\eta},
       \qquad
       s_{\min}(D+E)
       \ge\frac{1-\eta}{\|D^{-1}\|}.
       \label{eq:V33-block-floor}
      \end{equation}
\end{enumerate}
\end{theorem}

\begin{proof}
Since \(N\) is strictly block lower triangular, \(ND^{-1}\) has the same
property and \((ND^{-1})^m=0\).  The factorization
\[
 L=(I+ND^{-1})D
\]
and the finite geometric sum give
\eqref{eq:V33-triangular-inverse}.  In the second case,
\[
 D+E=D(I+D^{-1}E).
\]
The Neumann series for \(I+D^{-1}E\) converges under
\eqref{eq:V33-block-smallness}; its norm is at most
\((1-\eta)^{-1}\).  This proves the inverse bound.  The singular-floor
bound is its reciprocal.
\end{proof}

\begin{firewall}[Diagonal floors do not certify the coupled boundary row]
\label{fire:V33-no-sector-splice}
Positive singular floors for the Einstein, Yang--Mills, spinor, and action
diagonal blocks do not imply a positive floor for
\eqref{eq:V33-block-matrix}.  One must prove an exact triangular order,
verify \eqref{eq:V33-block-smallness}, or compute the singular value of the
full realified matrix.  Independent sector certificates cannot be
multiplied or minimized after ignoring the mixed columns.
\end{firewall}

\subsection{Nonlinear exact repair in a neutral chart}
\label{subsec:V33-nonlinear}

The linear test is decisive for tangent compatibility, but the continuum
programme requires an exact boundary graph.  The next theorem converts a
surjective tangent matrix into such a graph while remaining inside the
exact chart \eqref{eq:V33-unitary-chart}.

\begin{theorem}[Quantitative nonlinear boundary repair]
\label{thm:V33-nonlinear-repair}
Let \(\mathcal X\) be the full admissible real parameter space and let
\(L_*^{\rm full}:\mathcal X\to\mathcal Y\) be a surjective reference
linearization.  Put
\[
 \mathcal X_0=(\ker L_*^{\rm full})^\perp,\qquad
 L_*:=L_*^{\rm full}|_{\mathcal X_0}:
 \mathcal X_0\longrightarrow\mathcal Y.
\]
Then \(L_*\) is bijective.  Let
\(\mathfrak B:B_r(0)\subset\mathcal X_0\to\mathcal Y\) be \(C^1\),
put \(R_*:=L_*^{-1}\), and put \(\kappa_*:=\|R_*\|\).  Assume
\begin{align}
 \sup_{\|z\|\le r}
 \|R_*(D\mathfrak B(z)-L_*)\|
 &\le\theta<1,
 \label{eq:V33-derivative-radius}\\
 \|R_*\mathfrak B(0)\|
 &\le(1-\theta)r.
 \label{eq:V33-center-condition}
\end{align}
Then there is a unique \(z_\sharp\in B_r(0)\cap\mathcal X_0\) such that
\begin{equation}
 \mathfrak B(z_\sharp)=0.
 \label{eq:V33-exact-boundary-root}
\end{equation}
It satisfies
\begin{equation}
 \|z_\sharp\|
 \le
 \frac{\|R_*\mathfrak B(0)\|}{1-\theta}.
 \label{eq:V33-root-size}
\end{equation}
If the chart is \eqref{eq:V33-unitary-chart} with covariant parameter
space \eqref{eq:V33-Omega-covariance}, the repaired fermion graph remains
exactly maximal neutral and gauge covariant.  Because
\(\mathfrak A_{\partial}\) is a component of \(\mathfrak B\), the repaired
graph also obeys the mixed Higgs conormal equation
\eqref{eq:V33-action-row}.
\end{theorem}

\begin{proof}
Define
\[
 \mathcal T(z):=z-R_*\mathfrak B(z).
\]
For \(z,w\in B_r(0)\), the fundamental theorem of calculus and
\eqref{eq:V33-derivative-radius} give
\[
 \|\mathcal T(z)-\mathcal T(w)\|
 \le\theta\|z-w\|.
\]
Also
\[
 \|\mathcal T(z)\|
 \le\|\mathcal T(0)\|+\theta\|z\|
 \le(1-\theta)r+\theta r=r,
\]
so \(\mathcal T\) is a strict contraction of the closed ball.  Its unique
fixed point \(z_\sharp\) satisfies
\(R_*\mathfrak B(z_\sharp)=0\).  The inverse \(R_*\) is injective, hence
\(\mathfrak B(z_\sharp)=0\).  The fixed-point estimate
\[
 \|z_\sharp\|
 \le\|\mathcal T(0)\|+\theta\|z_\sharp\|
\]
proves \eqref{eq:V33-root-size}.  Exact neutrality and covariance follow
from Lemma~\ref{lem:V33-neutral-chart} and
\eqref{eq:V33-corrected-covariance}; the action conclusion follows from
the second component of \eqref{eq:V33-combined-residual}.
\end{proof}

\begin{corollary}[Second-derivative sufficient condition]
\label{cor:V33-second-derivative}
If \(L_*=D\mathfrak B(0)|_{\mathcal X_0}\) and
\[
 \sup_{\|z\|\le r}\|D^2\mathfrak B(z)\|\le M,
\]
then \eqref{eq:V33-derivative-radius} holds with
\begin{equation}
 \theta=\kappa_*Mr.
 \label{eq:V33-quadratic-theta}
\end{equation}
Thus
\[
 \kappa_*Mr<1,\qquad
 \|L_*^{-1}\mathfrak B(0)\|
 \le(1-\kappa_*Mr)r
\]
are directly checkable sufficient conditions for an exact repair.
\end{corollary}

\begin{proof}
The mean-value formula gives
\(\|D\mathfrak B(z)-D\mathfrak B(0)\|\le M\|z\|\).
Multiply by \(R_*\) and use \(\|z\|\le r\).
\end{proof}

\subsection{Cofinal stability and subsidiary propagation}
\label{subsec:V33-cofinal}

Transport all cutoff spaces into the common V24 frame and fix one reference
inverse \(R_*\) on \(\mathcal X_0\).  This avoids comparing pseudoinverses
whose kernels may rotate with \(q\).

\begin{theorem}[Cofinal boundary graph and constraint preservation]
\label{thm:V33-cofinal}
Suppose for every \(q\ge q_0\):
\begin{enumerate}
\item the admissible chart \(T_q(z)\) has the same real parameter and target
      spaces after transport, and it includes
      \eqref{eq:V33-moving-lift} and \eqref{eq:V33-action-row};
\item the V20 main and subsidiary complementing conditions hold uniformly,
      with \(\sup_{a,q}\kappa_{a,q}<1\);
\item for one bijective \(L_*:\mathcal X_0\to\mathcal Y\),
      \begin{equation}
       \sup_{q\ge q_0}\sup_{\|z\|\le r}
       \|R_*(D\mathfrak B_q(z)-L_*)\|
       \le\theta<1,
       \label{eq:V33-uniform-radius}
      \end{equation}
      and
      \begin{equation}
       \sup_{q\ge q_0}\|R_*\mathfrak B_q(0)\|
       \le(1-\theta)r;
       \label{eq:V33-uniform-center}
      \end{equation}
\item the adjacent boundary action satisfies
      \begin{equation}
       \epsilon_q:=
       \sup_{\|z\|\le r}
       \|\mathfrak B_{q+1}(z)-\mathfrak B_q(z)\|,
       \qquad
       \sum_{q\ge q_0}\epsilon_q<\infty.
       \label{eq:V33-adjacent-boundary}
      \end{equation}
\end{enumerate}
Then there are unique repaired parameters
\(z_q\in B_r(0)\cap\mathcal X_0\) with
\(\mathfrak B_q(z_q)=0\), and
\begin{equation}
 \|z_\infty-z_q\|
 \le
 \frac{\|R_*\|}{1-\theta}
 \sum_{j\ge q}\epsilon_j.
 \label{eq:V33-boundary-tail}
\end{equation}
The limiting allowed trace is gauge covariant, action compatible, and
maximal neutral on the V32 Higgs-floor branch.

If, in addition, the independent subsidiary systems have compatible zero
initial data and interior residuals \(R_{C,q}\to0\) in the uniform V20
energy norm, then
\begin{equation}
 \|c_q\|_{\rm sub}
 \le C_T\bigl(
 \|c_q(0)\|+\|R_{C,q}\|
 \bigr)\longrightarrow0.
 \label{eq:V33-subsidiary-landing}
\end{equation}
For exact zero initial data and \(R_{C,q}=0\), every cutoff subsidiary
constraint vanishes identically.
\end{theorem}

\begin{proof}
Theorem~\ref{thm:V33-nonlinear-repair} applies uniformly to each
\(\mathfrak B_q\).  Let
\(\mathcal T_q(z)=z-R_*\mathfrak B_q(z)\).  Since
\(\mathcal T_q\) is a \(\theta\)-contraction,
\begin{align*}
 \|z_{q+1}-z_q\|
 &\le
 \|\mathcal T_{q+1}(z_{q+1})
       -\mathcal T_{q+1}(z_q)\|
 +\|\mathcal T_{q+1}(z_q)-\mathcal T_q(z_q)\|\\
 &\le
 \theta\|z_{q+1}-z_q\|+\|R_*\|\epsilon_q.
\end{align*}
Therefore
\[
 \|z_{q+1}-z_q\|
 \le\frac{\|R_*\|}{1-\theta}\epsilon_q.
\]
Summation proves convergence and \eqref{eq:V33-boundary-tail}.  Uniform
convergence of \(\mathfrak B_q\) on the ball passes
\(\mathfrak B_q(z_q)=0\) to the limit.  The chart identities give the
neutrality, covariance, and action claims.

The first component of \(\mathfrak B_q(z_q)=0\) is
\(\mathfrak D_q(z_q)=0\).  Hence the boundary datum in
\eqref{eq:V20-constraint-boundary-estimate} vanishes for every free main
trace.  Applying that estimate to the direct sum of the independent
subsidiary systems gives \eqref{eq:V33-subsidiary-landing}.  The exact
claim follows from uniqueness.  The V25 compatibility identities then
annihilate the omitted Bianchi and integrability defects once the defining
defects vanish.
\end{proof}

\begin{firewall}[Boundary repair does not remove an interior Ward debit]
\label{fire:V33-interior-debit}
The equation \(\mathfrak D_q(z_q)=0\) removes the incoming boundary source
in the subsidiary energy estimate.  It does not set the interior
gauge--Ward, differentiated Einstein, cutoff-incidence, or collar residuals
to zero.  Those remain the \(R_{C,q}\) debit in
\eqref{eq:V33-subsidiary-landing} and must be controlled in the V23 target
regularity.
\end{firewall}

\subsection{Finite certified card}
\label{subsec:V33-finite-card}

At level \(q\), acquire in one transported frame
\[
 \widehat{\mathbf d}_q,\quad
 \widehat{\mathbf L}_q,\quad
 \widehat T_q,\quad
 \widehat{\mathcal C}_{\partial,q},\quad
 \widehat Q_q,\quad
 \widehat\Pi_q,\quad
 \widehat{D_\Phi\Pi}_q,
\]
together with orthonormal real bases for the covariant parameter columns
and the independent target rows.  Store rigorous errors
\begin{equation}
 \|\mathbf d_q-\widehat{\mathbf d}_q\|\le\varepsilon_{d,q},
 \qquad
 \|\mathbf L_q-\widehat{\mathbf L}_q\|\le\varepsilon_{L,q}.
 \label{eq:V33-matrix-errors}
\end{equation}
Let
\begin{align}
 \widehat\rho_q
 &:=
 s_{\min}^{+}(\widehat{\mathbf L}_q),
 \label{eq:V33-observed-floor}\\
 \rho_q^{\rm cert}
 &:=
 \widehat\rho_q-\varepsilon_{L,q},
 \label{eq:V33-certified-floor}\\
 r_{{\rm left},q}
 &:=
 \left\|
 (I-\widehat{\mathbf L}_q\widehat{\mathbf L}_q^\dagger)
 \widehat{\mathbf d}_q
 \right\|,
 \label{eq:V33-left-residual}\\
 \widehat z_{{\rm lin},q}
 &:=
 -\widehat{\mathbf L}_q^\dagger\widehat{\mathbf d}_q,
 \label{eq:V33-recorded-correction}\\
 r_{{\rm lin},q}^{\rm cert}
 &:=
 r_{{\rm left},q}
 +\varepsilon_{d,q}
 +\varepsilon_{L,q}\|\widehat z_{{\rm lin},q}\|.
 \label{eq:V33-linear-certified-residual}
\end{align}
Here \(s_{\min}^{+}\) is the smallest nonzero singular value.  Full row
rank is reported separately; \(s_{\min}^{+}>0\) by itself does not imply
surjectivity.

\begin{proposition}[Rigorous perturbation bounds for the finite card]
\label{prop:V33-certified-card}
Let \(\widehat r_q=\operatorname{rank}\widehat{\mathbf L}_q\).
The ordered exact singular values satisfy
\begin{equation}
 \sigma_{\widehat r_q}(\mathbf L_q)
 \ge
 \sigma_{\widehat r_q}(\widehat{\mathbf L}_q)
 -\varepsilon_{L,q}
 =\widehat\rho_q-\varepsilon_{L,q}.
 \label{eq:V33-Weyl-floor}
\end{equation}
If \(\widehat{\mathbf L}_q\) has full row rank and
\(\widehat\rho_q>\varepsilon_{L,q}\), then
\(\mathbf L_q\) also has full row rank and
\(s_{\min}(\mathbf L_q)\ge\rho_q^{\rm cert}>0\).
Moreover,
\begin{equation}
 \|\mathbf d_q+\mathbf L_q\widehat z_{{\rm lin},q}\|
 \le r_{{\rm lin},q}^{\rm cert}.
 \label{eq:V33-certified-linear-bound}
\end{equation}
If
\(\sup_{\|z\|\le\|\widehat z_{{\rm lin},q}\|}
\|D^2\mathfrak B_q(z)\|\le M_q\), then
\begin{equation}
 \|\mathfrak B_q(\widehat z_{{\rm lin},q})\|
 \le
 r_{{\rm lin},q}^{\rm cert}
 +\frac12M_q\|\widehat z_{{\rm lin},q}\|^2.
 \label{eq:V33-certified-nonlinear-bound}
\end{equation}
\end{proposition}

\begin{proof}
The singular-value perturbation inequality gives
\eqref{eq:V33-Weyl-floor}.  By definition,
\[
 \widehat{\mathbf d}_q+
 \widehat{\mathbf L}_q\widehat z_{{\rm lin},q}
 =
 (I-\widehat{\mathbf L}_q\widehat{\mathbf L}_q^\dagger)
 \widehat{\mathbf d}_q.
\]
Add and subtract the hatted vector and matrix, then use
\eqref{eq:V33-matrix-errors} to prove
\eqref{eq:V33-certified-linear-bound}.  Taylor's formula with integral
remainder gives the additional quadratic term in
\eqref{eq:V33-certified-nonlinear-bound}.
\end{proof}

A complete V33 certificate records:
\begin{enumerate}
\item the independent-row incidence rank which removes the V25 compatible
      identities;
\item the dimensions in \eqref{eq:V33-dimension-count}, both before and
      after gauge-covariant restriction of the columns;
\item the full block matrix \eqref{eq:V33-block-matrix}, rather than only
      its diagonal blocks;
\item \(r_{{\rm left},q}\), \(\rho_q^{\rm cert}\), the full-row-rank flag,
      the nonlinear radius \(M_q\), and the center inequality of
      Theorem~\ref{thm:V33-nonlinear-repair};
\item the exact flux, gauge-covariance, moving-domain, and mixed-conormal
      residuals from V32 evaluated on the repaired graph;
\item the adjacent debit \(\epsilon_q\) and the interior subsidiary debit
      in \eqref{eq:V33-subsidiary-landing}.
\end{enumerate}

\begin{firewall}[A numerical pseudoinverse is not an exact root]
\label{fire:V33-pseudoinverse}
The vector \(-\widehat{\mathbf L}_q^\dagger\widehat{\mathbf d}_q\) is only
a tangent correction.  An exact constraint-preserving boundary graph
requires the full-row-rank and nonlinear radius hypotheses of
Theorem~\ref{thm:V33-nonlinear-repair}, or a separate exact construction.
The certified residual \eqref{eq:V33-certified-nonlinear-bound} must not be
reported as zero merely because it is small.
\end{firewall}

\subsection{V33 ledger and next upstream row}
\label{subsec:V33-ledger}

\paragraph{New conclusions proved in V33.}
Equation \eqref{eq:V33-full-defect} assembles every independent incoming
subsidiary row as one operator on all free physical amplitudes.
Proposition~\ref{prop:V33-defect-derivative} gives its exact moving-frame
derivative.  Lemma~\ref{lem:V33-neutral-chart} constructs an exact neutral
unitary chart, and \eqref{eq:V33-moving-constraint-column} inserts the V32
Higgs-domain motion into the constraint matrix.  Theorem
\ref{thm:V33-linear-repair} gives the exact left-null obstruction and the
minimum neutral correction.  Theorems~\ref{thm:V33-block-inverse} and
\ref{thm:V33-nonlinear-repair} turn the full coupled matrix into a
quantitative exact repair criterion.  Theorem~\ref{thm:V33-cofinal} proves
cofinal convergence of the repaired graphs and invokes V20 to propagate
the independent constraints.

\paragraph{Imported conclusions not repeated.}
The subsidiary energy estimate is V18/V20, the compatible Cartan--SM
identities are V25, the differentiated spinor domain is V30, the
action/Dirac Green incidence is V31, and the explicit Higgs--Yukawa neutral
graph and mixed conormal are V32.  V33 supplies their previously missing
boundary-incidence matrix and repair theorem.

\paragraph{Authoritative physical status.}
The attached finite histories do not print the matrices
\(\mathcal C_{\partial,q}^a\), the normalized subsidiary splittings
\(P_{a,q}^{\pm,0}\), the independent-row quotient, the covariant admissible
parameter basis, or the full real matrices \(\mathbf L_q\).  They also do
not choose one of the V32 neutrino branches.  Therefore no positive V33
singular floor, vanishing left obstruction, or nonlinear repair radius is
currently populated.  The conditional theorem closes the logical
boundary-preservation implication; it does not establish the full
Standard Model--Einstein continuum from the supplied histories.

\paragraph{Next upstream row.}
The next iteration should select one concrete enlarged first-order
Einstein--Yang--Mills--Higgs--Dirac--BRST reduction and derive, from its
actual equations, the principal subsidiary Green matrices and the maps
\(\mathcal C_{\partial,q}^a\).  Those formulas determine the structural
zeros and mixed blocks of \eqref{eq:V33-block-matrix}.  If no
right-handed-neutrino or boundary-reservoir choice is made, the calculation
must first isolate the minimal-neutrino left-null witness rather than
assuming the missing covariant column exists.

% END APPENDED VERSION V33 MATERIAL

% =============================================================================
% START APPENDED VERSION V34 MATERIAL
% =============================================================================

\section{V34: Harmonic--Yang--Mills subsidiary Green matrices,
constraint readout reduction, and the spinor boundary identity}
\label{sec:V34}

\paragraph{Purpose and relation to V33.}
V33 defined the full incoming subsidiary defect in an arbitrary
first-order reduction and proved the exact range, singular-floor, nonlinear
repair, and cofinal criteria.  Its matrix was deliberately left abstract.
This note selects the Lorentzian metric-potential wave branch already used
in V18, V25, V29, and V30.  It derives the harmonic-Einstein and
Lorenz--Yang--Mills subsidiary equations, computes their normal Green
matrices and incoming trace maps, proves the non-glancing right-inverse
criterion, and determines which advertised constraints are readouts rather
than independent incoming fields.  It then computes the V30 spinor-jet
boundary defect on the V32 moving neutral graph.

The harmonic-Einstein part is consistent with the constraint-preserving
Sommerfeld formulation described in
\url{https://arxiv.org/abs/gr-qc/0612051} and with the frozen-coefficient
generalized-harmonic boundary analysis in
\url{https://arxiv.org/abs/0707.2797}.  The use of Lorenz gauge to obtain a
Yang--Mills or Yang--Mills--Higgs wave system is consistent with
\url{https://arxiv.org/abs/1309.1977} and
\url{https://arxiv.org/abs/1310.3361}.  The proofs below are self-contained
at the principal and differential-identity levels used here.

\subsection{One selected enlarged reduction}
\label{subsec:V34-selected-reduction}

Let \((M,g)\) have signature \((-+++)\) and a timelike boundary
\(\mathcal T\).  At a frozen boundary point choose an orthonormal frame
\[
 T,\quad n,\quad e_1,\quad e_2,
\]
where \(T\) is future timelike and tangent to \(\mathcal T\), \(n\) is the
outward unit spacelike normal, and \(e_A\) are spacelike and tangent to
\(\mathcal T\).  Connection coefficients vanish at the frozen point.
Curvature, torsion, background gauge curvature, second fundamental forms,
and derivatives of the frame are lower order in the following principal
calculation.

The selected branch has:
\begin{enumerate}
\item the generalized harmonic metric equation with gauge covector \(C_\nu\);
\item a connection-potential Yang--Mills equation in background Lorenz gauge
      with adjoint-valued gauge scalar \(\Gamma\);
\item wave jets for the metric, gauge potential, Higgs, internal ghosts, and
      the diffeomorphism ghost;
\item the V30 covariant Dirac prolongation
      \((\psi,\chi_a)\), with \(\chi_a\) an independent spinor field;
\item enough second jets of the metric and gauge potential to realize the
      first-order states of \(C\) and \(\Gamma\);
\item the V25 defining rows, but not their duplicate Bianchi or
      integrability consequences.
\end{enumerate}
The choice in item \(5\) is forced by differential order:
\(C\) and \(\Gamma\) are first derivatives of the main fields, whereas their
incoming wave traces contain one further derivative.  A first-jet main
trace alone cannot carry the V33 map
\(\mathcal C_{\partial,q}\) in the energy trace class.

Write the physical Einstein and Yang--Mills residuals as
\begin{align}
 E_{\mu\nu}
 &:=
 G_{\mu\nu}+\Lambda g_{\mu\nu}-T_{\mu\nu},
 \label{eq:V34-physical-Einstein}\\
 Y_\nu
 &:=
 D^\mu F_{\mu\nu}-J_\nu .
 \label{eq:V34-physical-YM}
\end{align}
Use the reduced residuals
\begin{align}
 \mathscr E_{\mu\nu}
 &:=
 E_{\mu\nu}
 -\nabla_{(\mu}C_{\nu)}
 +\frac12g_{\mu\nu}\nabla_\alpha C^\alpha,
 \label{eq:V34-reduced-Einstein}\\
 \mathscr Y_\nu
 &:=
 Y_\nu+D_\nu\Gamma .
 \label{eq:V34-reduced-YM}
\end{align}
For the nonlinear local formula, \(\Gamma\) is the chosen covariant Lorenz
gauge scalar in a fixed gauge chart.  For the frozen linearized symbol at a
background \((\bar g,\bar A)\), the actual choices are
\begin{align}
 C_\nu(h)
 &=
 \bar\nabla^\mu h_{\mu\nu}
 -\frac12\bar\nabla_\nu\operatorname{tr}_{\bar g}h,
 \label{eq:V34-linear-C}\\
 \gamma(a)
 &:=
 \bar D^\mu a_\mu .
 \label{eq:V34-linear-gamma}
\end{align}
A background Lorenz gauge
\(\bar D^\mu(A-\bar A)_\mu\) has
\eqref{eq:V34-linear-gamma} as its derivative.  At the anchor
\(A=\bar A\), metric variation of this gauge function has no principal
column.  This fixes a structural zero which would be false for an
unrecorded field-dependent gauge choice.

\subsection{Forced interior identities and the independent constraint list}
\label{subsec:V34-forced-identities}

The V18 identity already gives
\begin{equation}
 \nabla^\alpha\nabla_\alpha C_\nu
 +\operatorname{Ric}_{\nu\mu}C^\mu
 =
 -2\left(
 \nabla^\mu T_{\mu\nu}
 +\nabla^\mu\mathscr E_{\mu\nu}
 \right).
 \label{eq:V34-Einstein-subsidiary}
\end{equation}
The Yang--Mills analogue is as follows.

\begin{lemma}[Covariant double divergence of curvature]
\label{lem:V34-double-divergence}
For every connection curvature \(F\),
\begin{equation}
 D^\nu D^\mu F_{\mu\nu}=0.
 \label{eq:V34-double-divergence}
\end{equation}
\end{lemma}

\begin{proof}
Antisymmetry of \(F_{\mu\nu}\) and interchange of the dummy indices give
\[
 2D^\nu D^\mu F_{\mu\nu}
 =
 [D^\nu,D^\mu]F_{\mu\nu}.
\]
The spacetime-curvature terms are Ricci contractions with the antisymmetric
two-form and cancel.  The internal curvature term is the complete
contraction of
\([F^{\nu\mu},F_{\mu\nu}]\).  In an orthonormal frame it is a sum of
commutators of each Lie-algebra component with a scalar multiple of itself,
and hence vanishes.  The identity is tensorial, so the pointwise
calculation proves \eqref{eq:V34-double-divergence}.
\end{proof}

\begin{proposition}[Forced Lorenz--Yang--Mills subsidiary equation]
\label{prop:V34-YM-subsidiary}
The reduced residual \eqref{eq:V34-reduced-YM} satisfies
\begin{equation}
 \boxed{
 D^\nu D_\nu\Gamma
 =
 D^\nu\mathscr Y_\nu+D^\nu J_\nu .}
 \label{eq:V34-YM-subsidiary}
\end{equation}
Thus \(\Gamma\) obeys an adjoint-bundle wave equation whose forcing is the
differentiated reduced Yang--Mills residual plus the actual current Ward
defect.  At finite cutoff neither term may be set to zero before its
corresponding V23 residual row is controlled.
\end{proposition}

\begin{proof}
Take \(D^\nu\) of \eqref{eq:V34-reduced-YM}, substitute
\eqref{eq:V34-physical-YM}, and use
Lemma~\ref{lem:V34-double-divergence}:
\[
 D^\nu\mathscr Y_\nu
 =
 -D^\nu J_\nu+D^\nu D_\nu\Gamma.
\]
Rearrangement gives \eqref{eq:V34-YM-subsidiary}.
\end{proof}

Gauge invariance of the matter action expresses \(D^\nu J_\nu\) as a
linear incidence readout of the Higgs, fermion, and ghost Euler residuals.
Diffeomorphism invariance similarly expresses
\(\nabla^\mu T_{\mu\nu}\) through all matter and gauge residuals.  These
Noether identities remove circularity only when evaluated in the same
action and equation frame.  Equations
\eqref{eq:V34-Einstein-subsidiary} and
\eqref{eq:V34-YM-subsidiary} retain the two Ward terms explicitly.

\begin{proposition}[Hamiltonian, momentum, and Gauss are readouts]
\label{prop:V34-constraint-readouts}
Let
\[
 \mathcal H:=E(T,T),\qquad
 \mathcal M_i:=E(T,e_i),\qquad
 \mathcal G:=Y(T).
\]
Then
\begin{align}
 E_{\mu\nu}
 &=
 \mathscr E_{\mu\nu}
 +\nabla_{(\mu}C_{\nu)}
 -\frac12g_{\mu\nu}\nabla_\alpha C^\alpha,
 \label{eq:V34-E-readout}\\
 Y_\nu
 &=
 \mathscr Y_\nu-D_\nu\Gamma,
 \label{eq:V34-Y-readout}\\
 \mathcal G
 &=
 \mathscr Y(T)-D_T\Gamma.
 \label{eq:V34-Gauss-readout}
\end{align}
Consequently, on a reduced root, \(\mathcal H\) and \(\mathcal M_i\) are
linear readouts of \(\nabla C\), and \(\mathcal G=-D_T\Gamma\).
If the homogeneous subsidiary problems give \(C=0\) and \(\Gamma=0\),
then all three physical constraint families vanish.  They are not
additional incoming subsidiary wave fields in this potential reduction.
\end{proposition}

\begin{proof}
Equations \eqref{eq:V34-E-readout} and
\eqref{eq:V34-Y-readout} are rearrangements of the definitions of the
reduced residuals.  Contract with \(T\otimes T\), \(T\otimes e_i\), and
\(T\), respectively.  At a reduced root the displayed residual terms
vanish.  A subsidiary solution which vanishes identically has vanishing
covariant derivative, proving the last implication.
\end{proof}

\begin{firewall}[No duplicate Gauss or ADM incoming row]
\label{fire:V34-no-duplicate-physical-constraints}
In the selected potential-wave branch, adding independent incoming rows
for \(\mathcal H,\mathcal M_i,\mathcal G\) after adding the full states of
\(C\) and \(\Gamma\) duplicates
\eqref{eq:V34-E-readout}--\eqref{eq:V34-Gauss-readout}.  The
selected reduction therefore forbids a second nonzero-speed Gauss row.
\end{firewall}

\subsection{Universal wave subsidiary Green matrix}
\label{subsec:V34-wave-Green}

Let \(u\) be a realified bundle-valued field with principal equation
\[
 -D_T^2u+D_n^2u+\sum_{A=1}^2D_A^2u=0.
\]
Set
\[
 v:=D_Tu,\qquad s:=D_nu,\qquad r_A:=D_Au,\qquad
 U_u:=(v,s,r_1,r_2).
\]
The principal first-order system is
\begin{equation}
 D_TU_u+A_W^nD_nU_u+\sum_AA_W^AD_AU_u=0,
 \qquad
 A_W^n=
 \begin{pmatrix}
  0&-I&0\\
  -I&0&0\\
  0&0&0_{2}
 \end{pmatrix}.
 \label{eq:V34-wave-normal-matrix}
\end{equation}
The identity is the symmetrizer.  Define characteristic variables
\begin{equation}
 z_+:=2^{-1/2}(v-s),\qquad
 z_-:=2^{-1/2}(v+s),\qquad z_0:=(r_1,r_2).
 \label{eq:V34-wave-characteristics}
\end{equation}

\begin{theorem}[Normal splitting, flux, and incoming row]
\label{thm:V34-wave-Green}
The spectral projectors of \eqref{eq:V34-wave-normal-matrix} are
\begin{align}
 P_W^+
 &=
 \frac12
 \begin{pmatrix}
  I&-I&0\\-I&I&0\\0&0&0
 \end{pmatrix},
 \label{eq:V34-Pplus}\\
 P_W^-
 &=
 \frac12
 \begin{pmatrix}
  I&I&0\\I&I&0\\0&0&0
 \end{pmatrix},
 \label{eq:V34-Pminus}\\
 P_W^0
 &=
 \begin{pmatrix}
  0&0&0\\0&0&0\\0&0&I_2
 \end{pmatrix}.
 \label{eq:V34-Pzero}
\end{align}
The outward Green flux is
\begin{equation}
 \langle U_u,A_W^nU_u\rangle_{\mathbb R}
 =
 \|z_+\|^2-\|z_-\|^2.
 \label{eq:V34-wave-flux}
\end{equation}
For a contraction \(K:\operatorname{Ran}P_W^+\to
\operatorname{Ran}P_W^-\), the incoming boundary condition is
\begin{equation}
 q_K(U_u)
 :=
 z_--Kz_+
 =
 2^{-1/2}\bigl((I-K)v+(I+K)s\bigr)=0.
 \label{eq:V34-wave-incoming-row}
\end{equation}
If \(\|K\|\le1\), this boundary graph is maximally nonnegative.  If
\(\|K\|\le\kappa<1\), it has flux reserve
\begin{equation}
 \langle U_u,A_W^nU_u\rangle_{\mathbb R}
 \ge(1-\kappa^2)\|z_+\|^2.
 \label{eq:V34-wave-flux-reserve}
\end{equation}
The absorbing choice \(K=0\) is
\begin{equation}
 (D_T+D_n)u=0.
 \label{eq:V34-wave-Sommerfeld}
\end{equation}
\end{theorem}

\begin{proof}
On \((v,s)\), \(A_W^n\) is
\(\begin{psmallmatrix}0&-I\\-I&0\end{psmallmatrix}\).
Its \(+1\), \(-1\), and zero eigenspaces give
\eqref{eq:V34-Pplus}--\eqref{eq:V34-Pzero}.
Substitution of \eqref{eq:V34-wave-characteristics} gives
\eqref{eq:V34-wave-flux}.  On the graph \(z_-=Kz_+\),
\[
 \|z_+\|^2-\|z_-\|^2
 \ge(1-\|K\|^2)\|z_+\|^2.
\]
The positive and negative ranks agree, so the graph is maximal.
Equation \eqref{eq:V34-wave-incoming-row} is direct algebra, and \(K=0\)
gives \eqref{eq:V34-wave-Sommerfeld}.
\end{proof}

\begin{corollary}[Nonexpansive subsidiary graphs still propagate zero]
\label{cor:V34-nonexpansive-energy}
The V20 subsidiary energy estimate remains valid for
\(\|K\|\le1\), with no strictly positive boundary reserve required.
The same statement holds for any normalized symmetric-hyperbolic
boundary flux \(\operatorname{diag}(I,-I,0)\) whose incoming graph
satisfies \(z_-=Kz_+\).  In particular a homogeneous maximal neutral
graph with \(\|K\|=1\) still propagates zero initial constraints when
the interior subsidiary source vanishes.
\end{corollary}

\begin{proof}
For the wave block, the boundary term in the energy identity has the
favorable sign by \eqref{eq:V34-wave-flux}.  In a general normalized
block the identical calculation gives
\(\|z_+\|^2-\|Kz_+\|^2\ge0\).  Strict positivity is not used
in the Gronwall uniqueness argument; it is needed only for a coercive
trace estimate.
\end{proof}

\subsection{Actual harmonic and Lorenz trace maps}
\label{subsec:V34-actual-trace}

Let
\begin{align}
 H^g_{\lambda\mid\mu\nu}
 &:=
 \bar\nabla_\lambda h_{\mu\nu},
 &
 K^g_{\sigma\lambda\mid\mu\nu}
 &:=
 \bar\nabla_\sigma H^g_{\lambda\mid\mu\nu},
 \label{eq:V34-metric-two-jet}\\
 H^A_{\lambda\mid\nu}
 &:=
 \bar D_\lambda a_\nu,
 &
 K^A_{\sigma\lambda\mid\nu}
 &:=
 \bar D_\sigma H^A_{\lambda\mid\nu}.
 \label{eq:V34-gauge-two-jet}
\end{align}
At principal order define the contractions
\begin{align}
 (\mathsf H_EK^g_\sigma)_\nu
 &:=
 \eta^{\lambda\mu}
 K^g_{\sigma\lambda\mid\mu\nu}
 -\frac12
 K^g_{\sigma\nu\mid\alpha}{}^\alpha,
 \label{eq:V34-HE-two-jet}\\
 \mathsf H_YK^A_\sigma
 &:=
 \eta^{\lambda\nu}
 K^A_{\sigma\lambda\mid\nu}.
 \label{eq:V34-HY-two-jet}
\end{align}
Then the concrete V33 subsidiary trace maps are
\begin{align}
 \mathcal C_\partial^E(K^g)
 &=
 \bigl(
 \mathsf H_EK^g_T,\,
 \mathsf H_EK^g_n,\,
 \mathsf H_EK^g_1,\,
 \mathsf H_EK^g_2
 \bigr),
 \label{eq:V34-Cpartial-E}\\
 \mathcal C_\partial^Y(K^A)
 &=
 \bigl(
 \mathsf H_YK^A_T,\,
 \mathsf H_YK^A_n,\,
 \mathsf H_YK^A_1,\,
 \mathsf H_YK^A_2
 \bigr).
 \label{eq:V34-Cpartial-Y}
\end{align}
Indeed these are
\((D_TC,D_nC,D_AC)\) and
\((D_T\gamma,D_n\gamma,D_A\gamma)\), respectively.

For a real frozen covector
\[
 \zeta=\zeta_TT^\flat+\zeta_nn^\flat+\zeta_Ae^A
 \ne0,
\]
define the first-symbol contractions
\begin{align}
 \bigl(\mathsf H_E(\zeta)h\bigr)_\nu
 &:=
 \zeta^\mu h_{\mu\nu}
 -\frac12\zeta_\nu\operatorname{tr}_\eta h,
 \label{eq:V34-HE-symbol}\\
 \mathsf H_Y(\zeta)a
 &:=
 \zeta^\mu a_\mu.
 \label{eq:V34-HY-symbol}
\end{align}
On a plane-wave two-jet,
\[
 D_\sigma C=\zeta_\sigma\mathsf H_E(\zeta)h,
 \qquad
 D_\sigma\gamma=\zeta_\sigma\mathsf H_Y(\zeta)a.
\]
Hence the full incoming matrices for contractions \(K_E,K_Y\) are
\begin{align}
 \mathsf M_E^{K_E}(\zeta)
 &:=
 \mathsf R_{K_E}(\zeta)\mathsf H_E(\zeta),
 \label{eq:V34-ME}\\
 \mathsf M_Y^{K_Y}(\zeta)
 &:=
 \mathsf R_{K_Y}(\zeta)\mathsf H_Y(\zeta),
 \label{eq:V34-MY}\\
 \mathsf R_K(\zeta)
 &:=
 2^{-1/2}
 \left[
  (\zeta_T+\zeta_n)I+
  (\zeta_n-\zeta_T)K
 \right].
 \label{eq:V34-RK}
\end{align}
For \(K=0\), these reduce to
\begin{equation}
 \mathsf M_E^0(\zeta)
 =
 2^{-1/2}(\zeta_T+\zeta_n)\mathsf H_E(\zeta),
 \qquad
 \mathsf M_Y^0(\zeta)
 =
 2^{-1/2}(\zeta_T+\zeta_n)\mathsf H_Y(\zeta).
 \label{eq:V34-Sommerfeld-matrices}
\end{equation}
These are explicit matrices after choosing the standard symmetric-tensor,
one-form, and Lie-algebra bases.  No finite-history coefficient is hidden
in their frozen principal part.

\begin{lemma}[Surjectivity of the harmonic and Lorenz contractions]
\label{lem:V34-H-surjective}
For every nonzero real covector \(\zeta\),
\[
 \mathsf H_E(\zeta):
 \operatorname{Sym}^2(\mathbb R^{1,3})^*
 \longrightarrow(\mathbb R^{1,3})^*,
 \qquad
 \mathsf H_Y(\zeta):
 (\mathbb R^{1,3})^*\otimes\mathfrak g_{\rm SM}
 \longrightarrow\mathfrak g_{\rm SM}
\]
are surjective.
\end{lemma}

\begin{proof}
The adjoint of \(\mathsf H_Y(\zeta)\) maps an adjoint scalar \(x\) to
\(\zeta^\sharp\otimes x\); it has zero kernel because \(\zeta\ne0\).

For the Einstein map, suppose a covector \(v\) annihilates its range.
Then, for every symmetric \(h\),
\[
 \left\langle
 \zeta_{(\mu}v_{\nu)}
 -\frac12(\zeta\mathbin{\cdot}v)\eta_{\mu\nu},
 h^{\mu\nu}
 \right\rangle=0.
\]
Thus the symmetric tensor in the first slot vanishes.  Taking its Lorentz
trace in four dimensions gives
\(\zeta\mathbin{\cdot}v=0\), and hence
\(\zeta_{(\mu}v_{\nu)}=0\).  Choose \(w\) with
\(\zeta(w)\ne0\).  Evaluation on \(w,w\) gives \(v(w)=0\); evaluation on
\(w,y\) for \(y\in\ker\zeta\) gives \(v(y)=0\).  Therefore \(v=0\).
Both adjoints are injective, so both finite-dimensional maps are
surjective.
\end{proof}

Equip coefficient spaces with the Euclidean norms of the frozen
orthonormal bases and put
\begin{equation}
 h_E(\zeta):=s_{\min}^+(\mathsf H_E(\zeta)),
 \qquad
 h_Y(\zeta):=s_{\min}^+(\mathsf H_Y(\zeta)).
 \label{eq:V34-H-floors}
\end{equation}
Lemma~\ref{lem:V34-H-surjective} makes both positive for
\(\zeta\ne0\).  On any compact normalized real frequency set avoiding
\(\zeta=0\), their infima are positive.

\begin{proposition}[Exact singular values of the two contractions]
\label{prop:V34-exact-H-floors}
In the orthonormal coefficient norms just declared,
\begin{align}
 \mathsf H_E(\zeta)\mathsf H_E(\zeta)^*
 &=
 \frac12|\zeta|_{\rm E}^2I_4
 +\frac12(\eta\zeta)\otimes(\eta\zeta),
 \label{eq:V34-HE-Gram}\\
 \mathsf H_Y(\zeta)\mathsf H_Y(\zeta)^*
 &=|\zeta|_{\rm E}^2I_{\mathfrak g_{\rm SM}}.
 \label{eq:V34-HY-Gram}
\end{align}
Consequently
\begin{equation}
 h_E(\zeta)=2^{-1/2}|\zeta|_{\rm E},
 \qquad
 h_Y(\zeta)=|\zeta|_{\rm E}.
 \label{eq:V34-exact-H-floors}
\end{equation}
\end{proposition}

\begin{proof}
Use the Frobenius-orthonormal basis of symmetric tensors consisting of
the four diagonal matrix units and
\(2^{-1/2}(E_{\mu\nu}+E_{\nu\mu})\) for \(\mu<\nu\).
Expanding \eqref{eq:V34-HE-symbol} in that basis gives
\eqref{eq:V34-HE-Gram}.  The Euclidean vector \(\eta\zeta\) has
length \(|\zeta|_{\rm E}\), so the Gram eigenvalues are
\(|\zeta|_{\rm E}^2\) once and
\(\tfrac12|\zeta|_{\rm E}^2\) three times.  For the Lorenz map,
each Lie-algebra component is contraction with \(\eta\zeta\), giving
\eqref{eq:V34-HY-Gram}.  Taking square roots proves
\eqref{eq:V34-exact-H-floors}.
\end{proof}

\begin{theorem}[Non-glancing right inverse and singular floor]
\label{thm:V34-incoming-floor}
Assume \(\|K\|\le\kappa\) and define
\begin{equation}
 \delta_K(\zeta)
 :=
 2^{-1/2}
 \left(
 |\zeta_T+\zeta_n|
 -\kappa|\zeta_n-\zeta_T|
 \right).
 \label{eq:V34-deltaK}
\end{equation}
If \(\delta_K(\zeta)>0\), then
\begin{equation}
 s_{\min}(\mathsf R_K(\zeta))
 \ge\delta_K(\zeta).
 \label{eq:V34-RK-floor}
\end{equation}
Consequently
\(\mathsf M_E^{K_E}(\zeta)\) and
\(\mathsf M_Y^{K_Y}(\zeta)\) are surjective whenever the corresponding
\(\delta_{K_a}(\zeta)\) is positive, and on their kernel complements
\begin{align}
 s_{\min}^+(\mathsf M_E^{K_E}(\zeta))
 &\ge
 \delta_{K_E}(\zeta)h_E(\zeta),
 \label{eq:V34-ME-floor}\\
 s_{\min}^+(\mathsf M_Y^{K_Y}(\zeta))
 &\ge
 \delta_{K_Y}(\zeta)h_Y(\zeta).
 \label{eq:V34-MY-floor}
\end{align}
Right inverses are
\begin{align}
 \mathsf S_E(\zeta)
 &:=
 \mathsf H_E(\zeta)^\dagger
 \mathsf R_{K_E}(\zeta)^{-1},
 \label{eq:V34-SE}\\
 \mathsf S_Y(\zeta)
 &:=
 \mathsf H_Y(\zeta)^\dagger
 \mathsf R_{K_Y}(\zeta)^{-1},
 \label{eq:V34-SY}
\end{align}
with norms at most the reciprocals of the right sides of
\eqref{eq:V34-ME-floor} and \eqref{eq:V34-MY-floor}.
\end{theorem}

\begin{proof}
For every fibre vector \(x\),
\begin{align*}
 \|\mathsf R_K(\zeta)x\|
 &\ge
 2^{-1/2}\left(
 |\zeta_T+\zeta_n|\|x\|
 -|\zeta_n-\zeta_T|\|Kx\|
 \right)\\
 &\ge\delta_K(\zeta)\|x\|.
\end{align*}
This proves \eqref{eq:V34-RK-floor}.  By
Lemma~\ref{lem:V34-H-surjective},
\(\mathsf H_a\mathsf H_a^\dagger=I\).  Therefore
\[
 \mathsf M_a\mathsf S_a
 =
 \mathsf R_{K_a}\mathsf H_a\mathsf H_a^\dagger
 \mathsf R_{K_a}^{-1}=I.
\]
The product lower bounds follow by first applying the singular floor of
\(\mathsf R_{K_a}\) and then that of
\(\mathsf H_a\) on its kernel complement.  The same multiplication gives
the right-inverse norm estimates.
\end{proof}

\begin{corollary}[Populated principal Einstein--Yang--Mills block]
\label{cor:V34-populated-EY-block}
Let \(\mathfrak G_{\rm SM}\) have real dimension
\(8+3+1=12\).  On a compact normalized real frequency set
\(\mathcal K_{\rm ng}\), suppose
\[
 \inf_{\zeta\in\mathcal K_{\rm ng}}\delta_{K_E}(\zeta)>0,
 \qquad
 \inf_{\zeta\in\mathcal K_{\rm ng}}\delta_{K_Y}(\zeta)>0.
\]
Then the independent wave-subsidiary incoming map has \(4+12=16\)
rows and is surjective.  Its analytic principal floor is
\begin{equation}
 \rho_{EY}
 :=
 \inf_{\zeta\in\mathcal K_{\rm ng}}
 \min\{
 \delta_{K_E}(\zeta)h_E(\zeta),
 \delta_{K_Y}(\zeta)h_Y(\zeta)
 \}>0.
 \label{eq:V34-EY-floor}
\end{equation}
\end{corollary}

\begin{proof}
Take the direct sum of Theorem~\ref{thm:V34-incoming-floor} on the
four-dimensional harmonic covector and on the twelve-dimensional
Standard Model adjoint bundle.  Compactness makes the positive pointwise
lower bounds uniform.
\end{proof}

\begin{firewall}[The real Green floor is not the complex Lopatinski test]
\label{fire:V34-real-versus-Lopatinski}
Equation \eqref{eq:V34-EY-floor} is the real-frequency principal Green
floor for the symmetric-hyperbolic energy problem.  A full
pseudodifferential complementing certificate must also evaluate the
stable normal roots for complex Laplace frequency and include glancing
charts.  Positivity of \(\rho_{EY}\) must not be reported as that stronger
uniform Lopatinski determinant.
\end{firewall}

\subsection{Zero-speed defining defects}
\label{subsec:V34-zero-speed}

\begin{lemma}[Wave-jet definition and curl constraints have zero speed]
\label{lem:V34-wave-jet-zero}
For the principal wave first-order system
\begin{align}
 D_Tu-v&=0,
 \label{eq:V34-wave-u}\\
 D_Tv-D^iw_i&=0,
 \label{eq:V34-wave-v}\\
 D_Tw_i-D_iv&=0,
 \label{eq:V34-wave-w}
\end{align}
define
\begin{align}
 c_i&:=w_i-D_iu,
 \label{eq:V34-wave-definition-defect}\\
 z_{ij}&:=D_iw_j-D_jw_i.
 \label{eq:V34-wave-curl-defect}
\end{align}
At principal order,
\begin{equation}
 D_Tc_i=0,\qquad D_Tz_{ij}=0.
 \label{eq:V34-wave-defect-propagation}
\end{equation}
Their normal matrix is zero, so \(P^-=P^+=0\) and \(P^0=I\).
\end{lemma}

\begin{proof}
Use \eqref{eq:V34-wave-u} and \eqref{eq:V34-wave-w}:
\[
 D_Tc_i=D_Tw_i-D_iD_Tu=D_iv-D_iv=0.
\]
Similarly,
\[
 D_Tz_{ij}=D_iD_Tw_j-D_jD_Tw_i
 =D_iD_jv-D_jD_iv=0
\]
at principal order.  Curvature commutators are algebraic lower-order
terms in the V25 retained-curvature state.
\end{proof}

\begin{corollary}[No incoming V25 definition block on this branch]
\label{cor:V34-no-incoming-defining}
The metric, gauge-potential, Higgs, and ghost wave-jet defining and curl
defects contribute no nonzero-speed incoming row to the V34 principal
subsidiary matrix.  Their compatible initial values still have to vanish,
and their lower-order forced transport debits remain part of the V25/V23
constraint norm.
\end{corollary}

\begin{proof}
Apply Lemma~\ref{lem:V34-wave-jet-zero} componentwise and then use the V25
compatible-range identities to remove duplicate curls.
\end{proof}

\subsection{The spinor-jet row is killed by the differentiated projector}
\label{subsec:V34-spinor-row}

Let
\[
 C_a^\psi:=\chi_a-D_a\psi.
\]
At the boundary let \(\Pi=\Pi_{\rm SM}(\Phi)\) be the V32 maximal neutral
projector.  The main spinor and differentiated spinor domains are
\begin{align}
 (I-\Pi)\psi&=0,
 \label{eq:V34-psi-domain}\\
 (I-\Pi)\chi_a&=(D_a\Pi)\psi.
 \label{eq:V34-chi-domain}
\end{align}

\begin{theorem}[Exact spinor subsidiary boundary identity]
\label{thm:V34-spinor-boundary}
On \eqref{eq:V34-psi-domain}--\eqref{eq:V34-chi-domain},
\begin{equation}
 (I-\Pi)C_a^\psi=0.
 \label{eq:V34-Cpsi-in-graph}
\end{equation}
Equivalently, if
\(\Pi\) is the characteristic graph
\(\{u_+\oplus U_{\rm SM}u_+\}\), then
\begin{equation}
 \boxed{
 \bigl(P_D^- -U_{\rm SM}P_D^+\bigr)C_a^\psi=0.}
 \label{eq:V34-spin-incoming-zero}
\end{equation}
Thus the spinor-jet block of the V33 incoming defect vanishes identically
on the correctly differentiated V32 domain.  Its graph is conservative:
\(\|U_{\rm SM}\|=1\), so Corollary~\ref{cor:V34-nonexpansive-energy},
rather than a strict boundary reserve, supplies constraint uniqueness.
\end{theorem}

\begin{proof}
Differentiate \((I-\Pi)\psi=0\) covariantly:
\begin{equation}
 (I-\Pi)D_a\psi=(D_a\Pi)\psi.
 \label{eq:V34-differentiated-psi-domain}
\end{equation}
Subtract \eqref{eq:V34-differentiated-psi-domain} from
\eqref{eq:V34-chi-domain} to obtain
\eqref{eq:V34-Cpsi-in-graph}.  Membership in the graph of
\(U_{\rm SM}\) is exactly
\eqref{eq:V34-spin-incoming-zero}.  V32 proves that \(U_{\rm SM}\) is
unitary, so the normal Dirac flux of this graph is zero.  The V30
homogeneous Dirac equation for \(C_a^\psi\), its energy estimate, and
Corollary~\ref{cor:V34-nonexpansive-energy} propagate zero.
\end{proof}

\begin{firewall}[The inhomogeneous differentiated trace is essential]
\label{fire:V34-spinor-inhomogeneous}
Imposing \((I-\Pi)\chi_a=0\) when \(D_a\Pi\ne0\) gives
\[
 (I-\Pi)C_a^\psi=-(D_a\Pi)\psi,
\]
which is a genuine incoming spinor-constraint source.  The homogeneous
graph for \(\chi_a\) is valid only when the projector is covariantly
constant along that direction.
\end{firewall}

\subsection{BRST rows are Euler incidences, not new subsidiary waves}
\label{subsec:V34-BRST}

Let \(c_{\rm YM}\) be the internal ghost and \(c_{\rm diff}\) the
diffeomorphism ghost of the selected gauge-fixed continuum action.  At
principal order,
\begin{align}
 sA_\mu&=D_\mu c_{\rm YM}+\mathcal L_{c_{\rm diff}}A_\mu,
 \label{eq:V34-BRST-A}\\
 sg_{\mu\nu}&=\mathcal L_{c_{\rm diff}}g_{\mu\nu},
 \label{eq:V34-BRST-g}\\
 s\Gamma&=\mathcal M_{\rm YM}c_{\rm YM}
          +\operatorname{l.o.t.},
 \label{eq:V34-BRST-Gamma}\\
 sC_\nu&=(\Box_g\delta_\nu{}^\mu+
          \operatorname{Ric}_\nu{}^\mu)c_{{\rm diff},\mu}
          +\operatorname{l.o.t.},
 \label{eq:V34-BRST-C}
\end{align}
where \(\mathcal M_{\rm YM}\) is the Faddeev--Popov operator with principal
symbol \(\Box_g\) on the adjoint bundle.  The two ghost equations are
main Euler rows with the wave Green matrix
\eqref{eq:V34-wave-normal-matrix}.

\begin{proposition}[On-shell BRST preservation of the constraint boundary]
\label{prop:V34-BRST-boundary}
Let
\[
 B_EC:=q_{K_E}(D_TC,D_nC,D_AC),\qquad
 B_Y\Gamma:=q_{K_Y}(D_T\Gamma,D_n\Gamma,D_A\Gamma).
\]
At principal order,
\begin{align}
 s(B_EC)
 &=
 B_E\bigl((\Box_g+\operatorname{Ric})c_{\rm diff}\bigr)
 +\operatorname{l.o.t.},
 \label{eq:V34-BRST-BE}\\
 s(B_Y\Gamma)
 &=
 B_Y(\mathcal M_{\rm YM}c_{\rm YM})
 +\operatorname{l.o.t.}.
 \label{eq:V34-BRST-BY}
\end{align}
Hence exact ghost Euler equations and the lower-order incidence identities
make the homogeneous constraint boundary conditions BRST invariant.
Off shell, the right sides are boundary readouts of the ghost Euler
residuals and belong in the V23/V33 incidence debit.  They are not
independent incoming constraint variables.
\end{proposition}

\begin{proof}
Apply the linear boundary operators \(B_E,B_Y\) to
\eqref{eq:V34-BRST-Gamma}--\eqref{eq:V34-BRST-C}.  At principal order the
frozen coefficients commute with the BRST derivative.  The remaining
commutators differentiate coefficients at most once and are lower order
in the selected prolonged state.
\end{proof}

Gauge covariance of the moving fermion graph also supplies its BRST
tangency.  Infinitesimal covariance of \(\Pi(\Phi)\) is
\begin{equation}
 D_\Phi\Pi[\rho_\Phi(X)\Phi]
 =
 [\rho_\psi(X),\Pi].
 \label{eq:V34-Pi-infinitesimal-covariance}
\end{equation}
For \(\Pi\psi=\psi\), this gives
\begin{equation}
 (I-\Pi)\rho_\psi(X)\psi
 =
 D_\Phi\Pi[\rho_\Phi(X)\Phi]\psi.
 \label{eq:V34-BRST-domain-tangent}
\end{equation}
Equation \eqref{eq:V34-BRST-domain-tangent} is precisely the V32
moving-domain tangent equation for the gauge/BRST variation
\((\delta\Phi,\delta\psi)=(\rho_\Phi(X)\Phi,\rho_\psi(X)\psi)\).

\begin{firewall}[Analytic ghost waves do not settle the ghost action domain]
\label{fire:V34-ghost-domain}
The commuting-coefficient ghost representative has the wave Green matrix
computed above.  This does not choose a Grassmann/Krein realization or
prove that ghost, antighost, and Nakanishi--Lautrup boundary conditions
make the gauge-fixed action BRST invariant.  That common-action boundary
row remains the V27/V33 ghost-domain obligation.
\end{firewall}

\subsection{The assembled V34 principal subsidiary bank}
\label{subsec:V34-assembled}

Let \(N_\chi\) be the number of retained spinor-jet directions.  Before
restricting to the allowed main trace, the normal principal matrix of the
independent subsidiary evolution is
\begin{equation}
 A_{\rm sub}^n
 =
 \bigl(A_W^n\otimes I_4\bigr)
 \oplus
 \bigl(A_W^n\otimes I_{12}\bigr)
 \oplus
 \bigl(\alpha(n)\otimes I_{N_\chi}\bigr)
 \oplus0_{\rm def}.
 \label{eq:V34-full-normal-bank}
\end{equation}
The blocks are respectively harmonic Einstein, Lorenz Yang--Mills,
spinor-jet definition, and zero-speed wave-jet definitions.  On a
plane-wave main two-jet, its incoming defect is
\begin{equation}
 \mathcal C_{{\rm in},34}(\zeta)
 =
 \begin{pmatrix}
  \mathsf M_E^{K_E}(\zeta)&0&0\\
  0&\mathsf M_Y^{K_Y}(\zeta)&0\\
  0&0&
  (P_D^--U_{\rm SM}P_D^+)
  (\chi_a-\zeta_a\psi)
 \end{pmatrix},
 \label{eq:V34-full-incoming-map}
\end{equation}
with no incoming zero-speed row.  On the differentiated V32 domain the
last block vanishes by Theorem~\ref{thm:V34-spinor-boundary}.  Therefore
the effective independent nonzero-speed incoming principal target has
the sixteen rows certified by
Corollary~\ref{cor:V34-populated-EY-block}.

\begin{theorem}[Principal population of the V33 constraint matrix]
\label{thm:V34-principal-population}
On the selected V34 branch, assume:
\begin{enumerate}
\item the real non-glancing frequency card has
      \(\rho_{EY}>0\) in \eqref{eq:V34-EY-floor};
\item the V25 two-jet defining rows make
      \eqref{eq:V34-Cpartial-E}--\eqref{eq:V34-Cpartial-Y} genuine traces
      in the declared graph norm;
\item the V32 Higgs-floor branch and the differentiated domain
      \eqref{eq:V34-chi-domain} hold;
\item the Hamiltonian, momentum, Gauss, compatible Bianchi, and BRST rows
      are quotiented by the exact readout identities proved above.
\end{enumerate}
Then the constraint component of the V33 frozen real matrix has a
surjective sixteen-row Einstein--Yang--Mills diagonal block with right
inverse norm at most \(\rho_{EY}^{-1}\).  The spinor incoming block is
structurally zero on allowed traces, and the remaining defining defects
have zero normal speed.  All principal off-diagonal Einstein--Yang--Mills
subsidiary blocks vanish in the anchor background gauge
\eqref{eq:V34-linear-gamma}; stress, current, curvature, and frame
couplings enter the lower-order and nonlinear V33 radius.
\end{theorem}

\begin{proof}
The wave diagonal blocks and their right inverses are
Corollary~\ref{cor:V34-populated-EY-block}.  Theorem
\ref{thm:V34-spinor-boundary} removes the spinor row after restriction to
the allowed trace.  Corollary~\ref{cor:V34-no-incoming-defining} removes
the zero-speed rows.  Proposition~\ref{prop:V34-constraint-readouts} and
Proposition~\ref{prop:V34-BRST-boundary} remove the duplicate readout
rows.  V29 and V30 show that, after spinor prolongation, stress and current
couplings are algebraic in the retained state; the anchor background
Lorenz gauge removes a principal metric column from \(\gamma\).
Thus the displayed principal block is diagonal.  Its direct-sum right
inverse has norm bounded by the reciprocal of the least diagonal floor.
\end{proof}

\begin{firewall}[Surjectivity before action restriction is not V33 repair]
\label{fire:V34-action-intersection}
Theorem~\ref{thm:V34-principal-population} proves surjectivity from the
full metric and gauge two-jet amplitudes.  V33 requires surjectivity after
those columns are restricted to the physical, flux-compatible,
gauge-covariant, and common-action conormal chart.  The restriction can
remove columns or create a left-null witness.  The analytic floor
\(\rho_{EY}\) is therefore an upstream input to the V33 matrix, not yet
the final boundary-repair floor.
\end{firewall}

\subsection{Finite and cofinal V34 card}
\label{subsec:V34-finite}

At cutoff \(q\), in the common V24 boundary frame, acquire
\[
 \widehat A_{W,q}^n,\quad
 \widehat P_{W,q}^{\pm,0},\quad
 \widehat{\mathsf H}_{E,q}(\zeta),\quad
 \widehat{\mathsf H}_{Y,q}(\zeta),\quad
 \widehat K_{E,q},\quad
 \widehat K_{Y,q},
\]
and the matrices of
\eqref{eq:V34-Cpartial-E}, \eqref{eq:V34-Cpartial-Y}, and
\eqref{eq:V34-spin-incoming-zero}.  Store
\begin{align}
 r_{A,q}^{34}
 &:=
 \left\|
 \widehat A_{W,q}^n-
 \begin{pmatrix}0&-I&0\\-I&0&0\\0&0&0_2\end{pmatrix}
 \right\|,
 \label{eq:V34-rA}\\
 r_{P,q}^{34}
 &:=
 \sum_{\epsilon\in\{+,-,0\}}
 \|\widehat P_{W,q}^\epsilon-P_W^\epsilon\|,
 \label{eq:V34-rP}\\
 r_{{\rm read},q}^{34}
 &:=
 \|\,\widehat E-\widehat{\mathscr E}
      -\nabla_{(\cdot}\widehat C_{\cdot)}
      +\tfrac12\widehat g\,\nabla\mathbin{\cdot}\widehat C\,\|
 \nonumber\\
 &\quad+
 \|\,\widehat Y-\widehat{\mathscr Y}
      +\widehat D\widehat\Gamma\,\|,
 \label{eq:V34-rread}\\
 r_{{\rm div},q}^{34}
 &:=
 \|\widehat D^\nu\widehat D^\mu\widehat F_{\mu\nu}\|,
 \label{eq:V34-rdiv}\\
 r_{{\rm spin},q}^{34}
 &:=
 \max_a
 \|(I-\widehat\Pi_q)
   (\widehat\chi_{a,q}-\widehat D_a\widehat\psi_q)\|,
 \label{eq:V34-rspin}\\
 r_{{\rm BRST},q}^{34}
 &:=
 \sup_{\|X\|=1}
 \|D_\Phi\widehat\Pi_q[\rho_\Phi(X)\widehat\Phi_q]
   -[\rho_\psi(X),\widehat\Pi_q]\|.
 \label{eq:V34-rBRST}
\end{align}
The exact algebraic identities make the ideal values zero.  At a numerical
cutoff they measure incidence and transport error.

Let \(\mathcal K_{{\rm ng},q}\) be the printed real frequency net and let
the net-to-continuum Lipschitz errors be
\(\varepsilon_{\zeta,E,q},\varepsilon_{\zeta,Y,q}\).  Define
\begin{align}
 \widehat h_{E,q}
 &:=
 \min_{\zeta\in\mathcal K_{{\rm ng},q}}
 s_{\min}^+(\widehat{\mathsf H}_{E,q}(\zeta)),
 \label{eq:V34-hEhat}\\
 \widehat h_{Y,q}
 &:=
 \min_{\zeta\in\mathcal K_{{\rm ng},q}}
 s_{\min}^+(\widehat{\mathsf H}_{Y,q}(\zeta)),
 \label{eq:V34-hYhat}\\
 \widehat\delta_{a,q}
 &:=
 \min_{\zeta\in\mathcal K_{{\rm ng},q}}
 2^{-1/2}\left(
 |\zeta_T+\zeta_n|
 -\|\widehat K_{a,q}\||\zeta_n-\zeta_T|
 \right),
 \label{eq:V34-deltahat}\\
 \rho_{34,q}^{\rm cert}
 &:=
 \min_{a\in\{E,Y\}}
 \left[
 (\widehat\delta_{a,q}-\varepsilon_{\zeta,a,q})
 (\widehat h_{a,q}-\varepsilon_{H,a,q})
 \right]
 -\varepsilon_{{\rm frame},q}.
 \label{eq:V34-certified-floor}
\end{align}
Every factor in brackets must be positive before the product is used.

\begin{theorem}[Cofinal handoff to the action-restricted V33 matrix]
\label{thm:V34-cofinal-handoff}
Assume on each compact cylinder:
\begin{enumerate}
\item the selected V34 variable list and background gauges are used at
      every cutoff in one transported frame;
\item
      \(\inf_q\rho_{34,q}^{\rm cert}>0\);
\item the residuals
      \eqref{eq:V34-rA}--\eqref{eq:V34-rBRST} and the second-jet defining
      residuals tend uniformly to zero with summable adjacent tails;
\item the V18/V23 interior Ward and differentiated main-residual debits
      tend to zero in the subsidiary source norms;
\item the V32 Higgs/Yukawa floor and differentiated projector branch
      remains valid.
\end{enumerate}
Then the cofinal real-frequency principal incoming subsidiary map has the
explicit form \eqref{eq:V34-full-incoming-map}.  Its independent
Einstein--Yang--Mills block is uniformly surjective, the spinor-jet
incoming defect vanishes on the allowed trace, and no duplicate
Hamiltonian, momentum, Gauss, compatible Bianchi, or BRST incoming row is
present.  This data may be inserted into the constraint component of the
V33 full action-restricted matrix.
\end{theorem}

\begin{proof}
The perturbation inequality for singular values and
\eqref{eq:V34-certified-floor} preserve the uniform diagonal floor from
Theorem~\ref{thm:V34-incoming-floor}.  Summable transport tails give a
common limiting frame and matrices.  The vanishing readout, double
divergence, spinor, BRST, and defining residuals pass the exact identities
to the limit.  The forced interior identities then have vanishing sources
in the V23 subsidiary norm.  The assembled statement follows from
Theorem~\ref{thm:V34-principal-population}.
\end{proof}

\begin{firewall}[V34 physical status]
\label{fire:V34-status}
The analytic principal matrices in
\eqref{eq:V34-ME}--\eqref{eq:V34-MY} are now explicit.  The attached
histories do not provide a complex-frequency/glancing complementing card,
an action-compatible set of the remaining gravitational-radiation and
Yang--Mills physical boundary columns, or a ghost/antighost action domain.
They also do not select the V32 neutrino branch.  Therefore the full V33
action-restricted singular floor and the complete
Standard Model--Einstein continuum remain unproved.
\end{firewall}

\subsection{V34 ledger and next upstream row}
\label{subsec:V34-ledger}

\paragraph{New conclusions proved in V34.}
Proposition~\ref{prop:V34-YM-subsidiary} derives the forced
Lorenz--Yang--Mills wave equation.  Proposition
\ref{prop:V34-constraint-readouts} proves that Hamiltonian, momentum, and
Gauss constraints are readouts on the selected potential branch.
Theorem~\ref{thm:V34-wave-Green} computes the subsidiary normal matrix,
projectors, flux, and incoming contraction.  Equations
\eqref{eq:V34-ME}--\eqref{eq:V34-MY} are the actual harmonic and Lorenz
incoming matrices.  Theorem~\ref{thm:V34-incoming-floor} constructs their
right inverses and non-glancing floors.  Lemma
\ref{lem:V34-wave-jet-zero} removes zero-speed definition defects from
the incoming target.  Theorem~\ref{thm:V34-spinor-boundary} proves that
the differentiated V32 projector kills the spinor-jet incoming defect.
Proposition~\ref{prop:V34-BRST-boundary} places BRST boundary variation in
the ghost Euler incidence row.

\paragraph{Imported conclusions not repeated.}
The reduced Einstein identity is V18, the wave and Dirac bulk symbols are
V29, the compatible defining complex is V25, the spinor prolongation is
V30, and the moving neutral fermion graph is V32.  V33 supplies the
nonlinear action-restricted repair theorem into which the V34 matrices
feed.

\paragraph{Next upstream row.}
The next iteration should complete the boundary column side of the
matrix.  It must choose action-compatible physical boundary data for the
six metric components left after the four harmonic rows, for the
Yang--Mills polarizations left after the twelve Lorenz rows, and for the
Higgs and ghost/antighost sectors.  It should then compute the
complex-frequency stable-root determinant, including glancing charts, and
evaluate the V33 left-null obstruction after the common-action conormal
restriction.  A failure should be returned as an explicit stable mode or
left-null covector, not hidden by the positive real Green floor.

% END APPENDED VERSION V34 MATERIAL

% =============================================================================
% START APPENDED VERSION V35 MATERIAL
% =============================================================================

\section{V35: Complex stable-root atlas, bosonic action boundary test,
and the glancing obstruction}
\label{sec:V35}

\paragraph{Purpose and relation to V34.}
V34 computed the real normal Green matrices and the sixteen independent
harmonic-Einstein and Lorenz--Yang--Mills incoming rows.  It did not prove
the complex-frequency complementing condition or test the remaining
boundary columns against the common action.  V35 performs both
calculations for the frozen half-space model selected in V34.  The result
separates two issues which must not be conflated: the constraint
Sommerfeld rows remain uniformly complementing at glancing, whereas
natural closed-action completions can lose their uniform stable-root
floor.  A finite microlocal pivot atlas repairs the algebraic rank, but
its action incidence is an additional symplectic condition.

The BRST boundary discussion is consistent with the systematic local
construction in
\url{https://arxiv.org/abs/gr-qc/9610023}; the known tension between
BRST-invariant gravitational boundary data and strong ellipticity is also
described in \url{https://arxiv.org/abs/1301.0717}.  No result from those
papers is used without proof in the frozen calculation below.

\subsection{Stable normal root and its sharp Sommerfeld factor}
\label{subsec:V35-stable-root}

Use an inward coordinate \(x\ge0\), so the outward normal derivative is
\(D_n=-\partial_x\).  For a scalar component of any V34 wave block, take
the Laplace--Fourier ansatz
\begin{equation}
 u(t,x,y)
 =
 e^{st+i k\cdot y-\lambda x}u_0,
 \qquad
 \operatorname{Re}s>0,\quad k\in\mathbb R^2.
 \label{eq:V35-stable-ansatz}
\end{equation}
The wave equation gives
\begin{equation}
 \lambda^2=s^2+|k|^2.
 \label{eq:V35-root-equation}
\end{equation}
Let \(\lambda(s,k)\) be the unique square root with
\(\operatorname{Re}\lambda>0\).  Its limiting value on
\(\operatorname{Re}s=0\) is always taken from the open right half-plane.
The normalized closed frequency hemisphere is
\begin{equation}
 \overline{\mathbb S}_+
 :=
 \left\{
 (s,k):
 \operatorname{Re}s\ge0,\quad
 |s|^2+|k|^2=1
 \right\}.
 \label{eq:V35-normalized-hemisphere}
\end{equation}
The glancing set is
\begin{equation}
 \mathcal G
 :=
 \{(s,k)\in\overline{\mathbb S}_+:
 s=\pm i|k|,\ \lambda=0\}.
 \label{eq:V35-glancing-set}
\end{equation}

\begin{lemma}[Stable branch sign]
\label{lem:V35-branch-sign}
For the branch above,
\begin{equation}
 \operatorname{Re}(s\overline\lambda)\ge0.
 \label{eq:V35-positive-cross}
\end{equation}
\end{lemma}

\begin{proof}
Write \(s=a+ib\) and \(\lambda=c+id\).  In the open half-plane
\(a,c>0\).  The imaginary part of
\(\lambda^2=s^2+|k|^2\) gives \(cd=ab\).  Hence
\[
 \operatorname{Re}(s\overline\lambda)
 =ac+bd
 =a\frac{c^2+b^2}{c}\ge0.
\]
Pass to the closed hemisphere by continuity.
\end{proof}

\begin{theorem}[Sharp uniform Sommerfeld lower bound]
\label{thm:V35-Sommerfeld-factor}
On \(\overline{\mathbb S}_+\),
\begin{equation}
 \boxed{
 |s+\lambda|
 \ge
 \max\{|s|,|k|\}
 \ge2^{-1/2}.}
 \label{eq:V35-Sommerfeld-lower}
\end{equation}
The constant \(2^{-1/2}\) is sharp and is attained at glancing.
In particular \(s+\lambda\) has no zero on the normalized stable
hemisphere.
\end{theorem}

\begin{proof}
Lemma~\ref{lem:V35-branch-sign} gives
\[
 |s+\lambda|^2
 =
 |s|^2+|\lambda|^2
 +2\operatorname{Re}(s\overline\lambda)
 \ge |s|^2.
\]
Also
\[
 |\lambda|^2
 =
 |s^2+|k|^2|
 \ge\bigl||k|^2-|s|^2\bigr|.
\]
If \(|k|\ge|s|\), this and the preceding expansion give
\[
 |s+\lambda|^2
 \ge |s|^2+|k|^2-|s|^2=|k|^2.
\]
If \(|s|\ge|k|\), the first bound already gives
\(|s+\lambda|\ge|s|\).  Normalization yields
\(\max\{|s|,|k|\}\ge2^{-1/2}\).  At
\(s=\pm i/\sqrt2\), \(|k|=1/\sqrt2\), one has
\(\lambda=0\), so equality holds.
\end{proof}

\begin{firewall}[Glancing does not annihilate the Sommerfeld factor]
\label{fire:V35-glancing-Sommerfeld}
At glancing, the stable normal exponent \(\lambda\) vanishes, but the
Sommerfeld factor tends to the nonzero number \(s\).  A determinant which
loses its floor there does so through a different boundary row, typically
a pure Neumann factor \(\lambda\), rather than through \(s+\lambda\).
\end{firewall}

\subsection{Complex harmonic and Lorenz contraction floors}
\label{subsec:V35-complex-contractions}

Put
\begin{equation}
 \zeta(s,k)
 :=
 (s,\lambda,ik_1,ik_2)
 \in\mathbb C^4,
 \qquad
 |\zeta|_{\rm E}^2
 :=
 |s|^2+|\lambda|^2+|k|^2.
 \label{eq:V35-complex-zeta}
\end{equation}
On \(\overline{\mathbb S}_+\),
\begin{equation}
 |\zeta|_{\rm E}\ge1.
 \label{eq:V35-zeta-floor}
\end{equation}
Let
\[
 \bar h_{\mu\nu}
 :=
 h_{\mu\nu}
 -\frac12\eta_{\mu\nu}\operatorname{tr}_\eta h
\]
be trace reversal.  In four dimensions this is an invertible involution,
and the harmonic symbol is
\begin{equation}
 \mathsf D_E(\zeta)\bar h
 :=
 \zeta^\mu\bar h_{\mu\nu}.
 \label{eq:V35-trace-reversed-harmonic}
\end{equation}
The Lorenz symbol is
\begin{equation}
 \mathsf D_Y(\zeta)a
 :=
 \zeta^\mu a_\mu.
 \label{eq:V35-complex-Lorenz}
\end{equation}

\begin{proposition}[Exact complex Gram matrices]
\label{prop:V35-complex-Gram}
With the Hermitian coefficient norms induced by the frozen orthonormal
component bases,
\begin{align}
 \mathsf D_E(\zeta)\mathsf D_E(\zeta)^*
 &=
 \frac12|\zeta|_{\rm E}^2I_4
 +\frac12\overline{(\eta\zeta)}\,
             \overline{(\eta\zeta)}^{\,*},
 \label{eq:V35-complex-E-Gram}\\
 \mathsf D_Y(\zeta)\mathsf D_Y(\zeta)^*
 &=
 |\zeta|_{\rm E}^2I_{\mathfrak g_{\rm SM}}.
 \label{eq:V35-complex-Y-Gram}
\end{align}
Thus
\begin{align}
 s_{\min}^+(\mathsf D_E(\zeta))
 &=2^{-1/2}|\zeta|_{\rm E},
 \label{eq:V35-complex-E-floor}\\
 s_{\min}^+(\mathsf D_Y(\zeta))
 &=|\zeta|_{\rm E}.
 \label{eq:V35-complex-Y-floor}
\end{align}
Both maps are surjective at every stable and glancing frequency.
\end{proposition}

\begin{proof}
Use the Frobenius-orthonormal basis of complex symmetric matrices:
the diagonal matrix units and
\(2^{-1/2}(E_{\mu\nu}+E_{\nu\mu})\) for \(\mu<\nu\).
Direct multiplication gives
\eqref{eq:V35-complex-E-Gram}.  Since
\(|\eta\zeta|_{\rm E}=|\zeta|_{\rm E}\), its eigenvalues are
\(|\zeta|_{\rm E}^2\) once and
\(\tfrac12|\zeta|_{\rm E}^2\) three times.
The Lorenz map is componentwise contraction with \(\eta\zeta\), which
gives \eqref{eq:V35-complex-Y-Gram}.  The singular values are the square
roots of the Gram eigenvalues.  Positivity proves surjectivity.
\end{proof}

For the absorbing subsidiary choice \(K=0\), the stable incoming rows are
\begin{align}
 \mathsf L_E^{\rm Som}(s,k)
 &:=
 2^{-1/2}(s+\lambda)\mathsf D_E(\zeta),
 \label{eq:V35-E-Sommerfeld-symbol}\\
 \mathsf L_Y^{\rm Som}(s,k)
 &:=
 2^{-1/2}(s+\lambda)\mathsf D_Y(\zeta).
 \label{eq:V35-Y-Sommerfeld-symbol}
\end{align}
For the conservative constraint-Dirichlet choice, they are simply
\begin{equation}
 \mathsf L_E^{\rm D}:=\mathsf D_E(\zeta),
 \qquad
 \mathsf L_Y^{\rm D}:=\mathsf D_Y(\zeta).
 \label{eq:V35-constraint-Dirichlet}
\end{equation}

\begin{theorem}[Uniform complex constraint complementing floors]
\label{thm:V35-constraint-Lopatinski}
On the complete normalized stable hemisphere, including
\(\mathcal G\),
\begin{align}
 s_{\min}^+(\mathsf L_E^{\rm Som})
 &\ge\frac1{2\sqrt2},
 \label{eq:V35-E-Som-floor}\\
 s_{\min}^+(\mathsf L_Y^{\rm Som})
 &\ge\frac12,
 \label{eq:V35-Y-Som-floor}\\
 s_{\min}^+(\mathsf L_E^{\rm D})
 &\ge\frac1{\sqrt2},
 \label{eq:V35-E-D-floor}\\
 s_{\min}^+(\mathsf L_Y^{\rm D})
 &\ge1.
 \label{eq:V35-Y-D-floor}
\end{align}
Therefore neither the harmonic nor the Lorenz constraint row has a
complex stable-root or glancing rank defect on this frozen branch.
\end{theorem}

\begin{proof}
Multiply the bounds in
Theorem~\ref{thm:V35-Sommerfeld-factor},
\eqref{eq:V35-zeta-floor}, and
Proposition~\ref{prop:V35-complex-Gram}.  The extra
\(2^{-1/2}\) is the characteristic normalization in
\eqref{eq:V35-E-Sommerfeld-symbol}--%
\eqref{eq:V35-Y-Sommerfeld-symbol}.
\end{proof}

\subsection{A four-chart stable-amplitude completion}
\label{subsec:V35-pivot-atlas}

The reduced metric stable space has ten amplitudes
\(\bar h_{\mu\nu}=\bar h_{\nu\mu}\).  The harmonic rows account for four;
six supplementary rows are required.  The Yang--Mills potential has four
amplitudes per generator; its Lorenz row accounts for one, leaving three.
The following construction gives an exact microlocal completion without
calling those supplementary components physical.

For \(r\in\{T,n,1,2\}\), let
\begin{align}
 S_{E,r}\bar h
 &:=
 (\bar h_{\mu\nu})_{\mu,\nu\ne r},
 \label{eq:V35-SEr}\\
 S_{Y,r}a
 &:=
 (a_\mu)_{\mu\ne r}.
 \label{eq:V35-SYr}
\end{align}
Thus \(S_{E,r}\) has six coordinate components and \(S_{Y,r}\) has three
per Lie-algebra component.  Define
\begin{align}
 \mathcal B_{E,r}(\zeta)\bar h
 &:=
 \bigl(\mathsf D_E(\zeta)\bar h,S_{E,r}\bar h\bigr),
 \label{eq:V35-BEr}\\
 \mathcal B_{Y,r}(\zeta)a
 &:=
 \bigl(\mathsf D_Y(\zeta)a,S_{Y,r}a\bigr).
 \label{eq:V35-BYr}
\end{align}

\begin{theorem}[Pivot determinants and reconstruction]
\label{thm:V35-pivot-atlas}
If \(\zeta^r\ne0\), both maps in
\eqref{eq:V35-BEr}--\eqref{eq:V35-BYr} are bijective.  In coordinate
bases ordered with the \(r\)-columns last,
\begin{align}
 |\det\mathcal B_{E,r}(\zeta)|
 &=|\zeta^r|^4,
 \label{eq:V35-det-E}\\
 |\det\mathcal B_{Y,r}(\zeta)|
 &=|\zeta^r|^{\dim\mathfrak g_{\rm SM}}
 =|\zeta^r|^{12}.
 \label{eq:V35-det-Y}
\end{align}
If \(q_r:=|\zeta|_{\rm E}/|\zeta^r|\), their inverses obey
\begin{align}
 \|\mathcal B_{E,r}^{-1}\|
 &\le
 C_E(1+|\zeta^r|^{-1})(1+q_r)^2,
 \label{eq:V35-inverse-E}\\
 \|\mathcal B_{Y,r}^{-1}\|
 &\le
 C_Y(1+|\zeta^r|^{-1}+q_r),
 \label{eq:V35-inverse-Y}
\end{align}
for fixed dimensional constants \(C_E,C_Y\).
\end{theorem}

\begin{proof}
Write \(d=S_{E,r}\bar h\), \(b=\mathsf D_E(\zeta)\bar h\), and
\(x_\nu=\bar h_{r\nu}\).  For \(\nu\ne r\),
\begin{equation}
 x_\nu
 =
 \frac{
 b_\nu-\sum_{\mu\ne r}\zeta^\mu d_{\mu\nu}
 }{\zeta^r}.
 \label{eq:V35-reconstruct-cross}
\end{equation}
Then
\begin{equation}
 x_r
 =
 \frac{
 b_r-\sum_{\mu\ne r}\zeta^\mu x_\mu
 }{\zeta^r}.
 \label{eq:V35-reconstruct-diagonal}
\end{equation}
This proves bijectivity.  In the ordered coordinate matrix, the four
constraint rows restricted to the four \(x\)-columns are triangular with
diagonal \(\zeta^r\); the selector contributes an identity, proving
\eqref{eq:V35-det-E}.  Equations
\eqref{eq:V35-reconstruct-cross}--%
\eqref{eq:V35-reconstruct-diagonal} and Cauchy--Schwarz give
\eqref{eq:V35-inverse-E}.

For Yang--Mills, if \(d=S_{Y,r}a\) and
\(b=\mathsf D_Y(\zeta)a\), then
\[
 a_r
 =
 \frac{b-\sum_{\mu\ne r}\zeta^\mu d_\mu}{\zeta^r}.
\]
This proves the second determinant and inverse bound componentwise on the
adjoint bundle.
\end{proof}

\begin{corollary}[A stable atlas covering glancing]
\label{cor:V35-stable-atlas}
The four sets
\begin{equation}
 \mathcal U_r
 :=
 \left\{
 (s,k)\in\overline{\mathbb S}_+:
 |\zeta^r|\ge\frac12|\zeta|_{\rm E}
 \right\}
 \label{eq:V35-pivot-charts}
\end{equation}
cover \(\overline{\mathbb S}_+\).  On \(\mathcal U_r\),
\(q_r\le2\), so the inverse bounds in
Theorem~\ref{thm:V35-pivot-atlas} are uniform.  At a glancing point
\(\zeta^n=\lambda=0\), at least one of \(T,1,2\) supplies a valid pivot.
\end{corollary}

\begin{proof}
For four complex components,
\[
 |\zeta|_{\rm E}
 \le2\max_r|\zeta^r|.
\]
This proves the cover and \(q_r\le2\).  At glancing,
\(|s|^2+|k|^2=1\), so a nonnormal component is nonzero.
\end{proof}

Choose smooth homogeneous cutoffs \(\chi_r\) subordinate to slight
enlargements of the sets \(\mathcal U_r\).  Quantizing each
\(\chi_rS_{a,r}\) gives a chartwise order-zero pseudodifferential
completion and parametrix for the constraint rows.  A single global
boundary operator follows only if the target identifications agree on
overlaps; that transition condition is retained explicitly in the finite
certificate below.  Thus the atlas proves algebraic complex-frequency
complementing in local charts.  It is neither a local boundary condition
nor automatically a variational subspace.

\begin{firewall}[A pivot atlas is not a physical polarization choice]
\label{fire:V35-pivot-not-physical}
The six \(S_{E,r}\) and three \(S_{Y,r}\) components depend on the
frequency chart.  They need not coincide with gravitational radiation,
Yang--Mills polarizations, or an action conormal.  The positive
determinants \eqref{eq:V35-det-E}--\eqref{eq:V35-det-Y} settle stable
algebraic rank only.
\end{firewall}

\subsection{Exact boundary action test}
\label{subsec:V35-action-test}

Let \(Q\) be the realified boundary configuration fibre of one quadratic
wave block, with conormal momentum \(p=D_nu\).  On
\(\mathcal P=Q\oplus Q\), use the Green symplectic form
\begin{equation}
 \omega((u,p),(v,q))
 :=
 \langle p,v\rangle-\langle u,q\rangle.
 \label{eq:V35-Green-symplectic}
\end{equation}
A homogeneous linear boundary relation is
\begin{equation}
 \mathcal L_{A,B}
 :=
 \{(u,p):Au+Bp=0\},
 \qquad
 M:=(A\ B).
 \label{eq:V35-linear-boundary-relation}
\end{equation}

\begin{theorem}[Finite Lagrangian boundary criterion]
\label{thm:V35-Lagrangian-criterion}
Suppose \(\dim Q=m\) and \(M:Q\oplus Q\to Q\) has rank \(m\).
Then \(\mathcal L_{A,B}\) is Lagrangian for
\eqref{eq:V35-Green-symplectic} if and only if
\begin{equation}
 \boxed{AB^*=BA^*.}
 \label{eq:V35-AB-condition}
\end{equation}
Equivalently, with
\[
 \mathbb J=
 \begin{pmatrix}0&-I\\I&0\end{pmatrix},
\]
the condition is
\begin{equation}
 M\mathbb JM^*=0.
 \label{eq:V35-MJM}
\end{equation}
\end{theorem}

\begin{proof}
The rank assumption makes
\(\dim\ker M=m=\tfrac12\dim\mathcal P\).  The symplectic orthogonal of
\(\ker M\) is
\(\mathbb J\operatorname{Ran}M^*\), up to the harmless sign convention
in \(\mathbb J\).  Thus \(\ker M\) is isotropic, and hence Lagrangian,
exactly when every vector \(\mathbb JM^*y\) lies in \(\ker M\).  This is
\(M\mathbb JM^*=0\).  Direct multiplication gives
\[
 M\mathbb JM^*=BA^*-AB^*,
\]
which proves \eqref{eq:V35-AB-condition}.
\end{proof}

\begin{corollary}[Robin and Sommerfeld distinction]
\label{cor:V35-Robin-Sommerfeld}
A graph \(p=Ru\) is Lagrangian if and only if \(R=R^*\).  A real
time-frequency Sommerfeld graph
\[
 p=-D_Tu
\]
has \(R=-i\tau I\), which is skew-Hermitian for
\(\tau\ne0\), and is therefore not a closed-action Lagrangian graph.
Its strict causal dissipation must be balanced by an explicitly enlarged
boundary sector if the total action is to remain conservative.
\end{corollary}

\begin{proof}
Take \(A=-R\), \(B=I\) in
\eqref{eq:V35-AB-condition}.  Formal Fourier transformation in boundary
time sends \(D_T\) to \(i\tau\).
\end{proof}

For later use, every contraction has a canonical conservative causal
dilation.  Define
\[
 D_K:=(I-K^*K)^{1/2},
 \qquad
 D_{K^*}:=(I-KK^*)^{1/2}.
\]

\begin{proposition}[Julia dilation of a dissipative constraint graph]
\label{prop:V35-Julia}
For a contraction \(K:E_+\to E_-\), the block operator
\begin{equation}
 \mathcal U_K
 :=
 \begin{pmatrix}
  K&D_{K^*}\\
  D_K&-K^*
 \end{pmatrix}
 :
 E_+\oplus E_-
 \longrightarrow
 E_-\oplus E_+
 \label{eq:V35-Julia-operator}
\end{equation}
is unitary.  Hence
\begin{equation}
 \binom{z_-^{\rm phys}}{z_-^{\rm res}}
 =
 \mathcal U_K
 \binom{z_+^{\rm phys}}{z_+^{\rm res}}
 \label{eq:V35-dilated-graph}
\end{equation}
is a maximal neutral graph for the sum of the physical and reservoir
causal fluxes.  With \(z_+^{\rm res}=0\), the physical relation is
\(z_-^{\rm phys}=Kz_+^{\rm phys}\), and the missing flux is carried by
\(z_-^{\rm res}=D_Kz_+^{\rm phys}\).
\end{proposition}

\begin{proof}
Functional calculus and the polar decomposition give the intertwining
identity
\[
 KD_K=D_{K^*}K.
\]
The diagonal blocks of
\(\mathcal U_K^*\mathcal U_K\) are
\[
 K^*K+D_K^2=I,
 \qquad
 D_{K^*}^2+KK^*=I,
\]
and the off-diagonal blocks vanish by the intertwining identity.
The identical calculation for
\(\mathcal U_K\mathcal U_K^*\) proves unitarity.  Equality of the total
incoming and outgoing norms proves neutrality and maximality.
\end{proof}

\begin{firewall}[A causal unitary dilation still needs action incidence]
\label{fire:V35-dilation-action}
Proposition~\ref{prop:V35-Julia} balances the normalized causal flux.
It becomes a common-action boundary condition only after an actual
reservoir action, its Green form, and the V27 equation multiplier identify
that unitary graph with an action Lagrangian.  A formal auxiliary copy of
the trace space is not a physical boundary reservoir.
\end{firewall}

\subsection{A BRST-relative Yang--Mills boundary branch}
\label{subsec:V35-relative-YM}

For one adjoint component write
\[
 a_\parallel=(a_T,a_1,a_2),\qquad
 p_\mu=D_na_\mu.
\]
Consider the local relative conditions
\begin{equation}
 a_\parallel=0,
 \qquad
 \Gamma=p_n+D^aa_a=0.
 \label{eq:V35-relative-YM}
\end{equation}
Because a vanishing boundary trace has vanishing tangential derivative,
the second condition reduces to \(p_n=0\) on the first.  For the internal
ghost quartet use
\begin{equation}
 c_{\rm YM}=0,\qquad
 \bar c_{\rm YM}=0,\qquad
 b_{\rm YM}=0
 \quad\hbox{on }\mathcal T.
 \label{eq:V35-relative-ghost}
\end{equation}

\begin{theorem}[Action and BRST incidence of the relative branch]
\label{thm:V35-relative-incidence}
For the frozen gauge-fixed Yang--Mills wave action:
\begin{enumerate}
\item the boundary subspace
      \eqref{eq:V35-relative-YM} is Lagrangian for the Yang--Mills wave
      Green form;
\item Dirichlet ghost and antighost data kill the ghost Green form, while
      the auxiliary \(b\)-row is algebraic;
\item the linearized BRST rules preserve
      \eqref{eq:V35-relative-YM}--%
      \eqref{eq:V35-relative-ghost} on the ghost Euler equation.
\end{enumerate}
\end{theorem}

\begin{proof}
Allowed variations obey
\(\delta a_\parallel=0\) and \(\delta p_n=0\).
Thus
\[
 \langle p,\delta a\rangle
 =
 \langle p_\parallel,\delta a_\parallel\rangle
 +\langle p_n,\delta a_n\rangle=0.
\]
The subspace has four independent conditions in the eight-dimensional
wave phase space, so it is Lagrangian.  This also follows from
Theorem~\ref{thm:V35-Lagrangian-criterion}.

The second-order ghost Green form pairs ghost and antighost Dirichlet
traces, so it vanishes when both variations have zero trace.  The
Nakanishi--Lautrup field has no normal kinetic momentum in the standard
gauge-fixing quartet.

Finally,
\[
 s a_\parallel=D_\parallel c_{\rm YM}=0
\]
because the ghost trace is zero,
\(s\bar c_{\rm YM}=b_{\rm YM}=0\), and
\(sc_{\rm YM}=-\tfrac12[c_{\rm YM},c_{\rm YM}]=0\).
The V34 identity \(s\Gamma=\mathcal M_{\rm YM}c_{\rm YM}\) vanishes on
the ghost Euler equation.  Thus every boundary row is preserved.
\end{proof}

The complex stable boundary matrix of
\eqref{eq:V35-relative-YM} is, in the ordered variables
\((a_\parallel,a_n)\),
\begin{equation}
 \mathcal B_{\rm rel}(s,k)
 =
 \begin{pmatrix}
  I_3&0\\
  ik^\parallel&\lambda
 \end{pmatrix},
 \qquad
 \det\mathcal B_{\rm rel}=\lambda.
 \label{eq:V35-relative-determinant}
\end{equation}
For the full Standard Model adjoint bundle the determinant is
\(\lambda^{12}\).

\begin{proposition}[Explicit Yang--Mills glancing quasimode]
\label{prop:V35-YM-glancing}
The relative branch has no uniform stable-root singular floor on
\(\overline{\mathbb S}_+\).  Let
\[
 \tau_0=|k_0|=2^{-1/2},\qquad
 s_\epsilon=\epsilon+i\tau_0,\qquad
 r_\epsilon=(|s_\epsilon|^2+|k_0|^2)^{1/2},
\]
and normalize
\[
 \widetilde s_\epsilon=s_\epsilon/r_\epsilon,
 \qquad
 \widetilde k_\epsilon=k_0/r_\epsilon,
 \qquad
 \widetilde\lambda_\epsilon
 =
 r_\epsilon^{-1}
 (s_\epsilon^2+|k_0|^2)^{1/2}.
\]
Then
\(\operatorname{Re}\widetilde s_\epsilon>0\),
\((\widetilde s_\epsilon,\widetilde k_\epsilon)
\in\overline{\mathbb S}_+\), and
\(\widetilde\lambda_\epsilon\to0\).  For the unit stable amplitude
\(a=n^\flat X\), \(\|X\|=1\),
\begin{equation}
 \|\mathcal B_{\rm rel}
   (\widetilde s_\epsilon,\widetilde k_\epsilon)a\|
 =
 |\widetilde\lambda_\epsilon|
 \longrightarrow0.
 \label{eq:V35-YM-quasimode}
\end{equation}
\end{proposition}

\begin{proof}
One has
\[
 s_\epsilon^2+|k_0|^2
 =
 \epsilon^2+2i\epsilon\tau_0,
\]
so its stable square root tends to zero like
\(\epsilon^{1/2}\).  Normalization does not alter that conclusion.
For \(a=n^\flat X\), the three Dirichlet rows vanish and the Lorenz row
is \(\widetilde\lambda_\epsilon X\).  This proves
\eqref{eq:V35-YM-quasimode}.
\end{proof}

\begin{firewall}[Energy well-posedness and a uniform determinant differ]
\label{fire:V35-energy-versus-uniform}
The relative Yang--Mills branch is a standard homogeneous
Dirichlet--Neumann energy boundary condition and has no incoming energy
source.  Proposition~\ref{prop:V35-YM-glancing} says that its
inhomogeneous stable-root inverse is not uniform in the unweighted
boundary norm at glancing.  It does not produce a nonzero growing mode
with \(\operatorname{Re}s>0\).
\end{firewall}

\subsection{The naive gravitational completion fails the action test}
\label{subsec:V35-gravity-action}

Take the normal pivot \(r=n\) in
Theorem~\ref{thm:V35-pivot-atlas}.  Its six supplementary rows fix
\begin{equation}
 \bar h_{ab}=0,
 \qquad a,b\in\{T,1,2\},
 \label{eq:V35-six-Dirichlet}
\end{equation}
and the other four rows impose \(C_\nu=0\).  In boundary phase
variables
\[
 q_\nu:=\bar h_{n\nu},
 \qquad
 p_\nu:=D_n\bar h_{n\nu},
\]
the constraint rows reduce to
\begin{align}
 p_a&=0,
 \qquad a\in\{T,1,2\},
 \label{eq:V35-gravity-pa}\\
 p_n+D^aq_a&=0.
 \label{eq:V35-gravity-pn}
\end{align}

\begin{proposition}[Explicit gravitational action obstruction]
\label{prop:V35-gravity-action-obstruction}
The boundary subspace
\eqref{eq:V35-six-Dirichlet}--%
\eqref{eq:V35-gravity-pn} is not isotropic for the minimal wave Green
form at a nonzero tangential frequency.  Hence the normal-pivot
complementing boundary matrix is not a closed-action boundary condition
for that quadratic representative.
\end{proposition}

\begin{proof}
Choose compactly supported boundary test fields \(f_a\) and \(g\) such
that
\(\int_{\mathcal T}(D^af_a)g\ne0\).  Consider two allowed tangent
vectors.  The first has
\[
 q_a=f_a,\quad q_n=0,\quad
 p_a=0,\quad p_n=-D^af_a.
\]
The second has
\[
 \widetilde q_a=0,\quad
 \widetilde q_n=g,\quad
 \widetilde p_\nu=0.
\]
All fixed \(\bar h_{ab}\) variations vanish.  Their Green symplectic
pairing is
\[
 \omega((q,p),(\widetilde q,\widetilde p))
 =
 -\int_{\mathcal T}(D^af_a)g,
\]
which is nonzero by construction.  Thus the subspace is not isotropic.
\end{proof}

The same branch has stable determinant
\begin{equation}
 \det\mathcal B_{E,n}(\zeta)=\lambda^4,
 \label{eq:V35-gravity-normal-determinant}
\end{equation}
so it also loses its unweighted uniform stable-root floor at glancing.
Switching to one of the nonnormal charts in
Corollary~\ref{cor:V35-stable-atlas} repairs the determinant.  It does not
by itself repair the action problem: the different chart relation must
receive its own symplectic test.

\begin{corollary}[No configuration-only boundary density repairs the
nonsymmetric conormal]
\label{cor:V35-no-config-repair}
On the free \(q_\nu\)-variables, the relation
\eqref{eq:V35-gravity-pa}--\eqref{eq:V35-gravity-pn} has the form
\(p=Rq\), where
\[
 (Rq)_a=0,\qquad
 (Rq)_n=-D^aq_a.
\]
The tangential differential operator \(R\) is not formally self-adjoint.
Therefore it cannot be the Euler derivative of a quadratic boundary
functional depending only on \(q\).  A variational repair must alter the
boundary relation or add boundary variables.
\end{corollary}

\begin{proof}
The formal adjoint has the missing components
\[
 (R^*q)_a=D_aq_n,\qquad (R^*q)_n=0,
\]
so \(R\ne R^*\).  The Hessian of a scalar quadratic boundary functional
is formally self-adjoint.  Equivalently,
Corollary~\ref{cor:V35-Robin-Sommerfeld} rules out this graph.
\end{proof}

\begin{firewall}[This is not a no-go theorem for Einstein boundaries]
\label{fire:V35-not-Einstein-no-go}
Proposition~\ref{prop:V35-gravity-action-obstruction} concerns the naive
minimal wave Green form and the normal-pivot six-component Dirichlet
completion.  The Gibbons--Hawking--York term, a different generalized
harmonic boundary gauge, a boundary multiplier, or additional boundary
fields changes the conormal relation.  Those alternatives must be
computed; they are not ruled out here.
\end{firewall}

\subsection{The action-restricted V33 decision}
\label{subsec:V35-V33-decision}

At a fixed complex frequency, let
\[
 \mathcal S_+(s,k)
\]
be the stable amplitude space of the complete prolonged main system.
Let \(B_{\rm con}\) contain the harmonic, Lorenz, and differentiated
spinor rows already certified in V34, and let \(B_{\rm act}\) be a
full-rank boundary relation which passes
\eqref{eq:V35-AB-condition} in the common action frame.  Define
\begin{equation}
 \mathbb B_{35}(s,k)
 :=
 \left.
 \binom{B_{\rm con}(s,k)}{B_{\rm act}(s,k)}
 \right|_{\mathcal S_+(s,k)}.
 \label{eq:V35-full-stable-boundary}
\end{equation}

Assume that the total number of independent boundary rows equals
\(\dim\mathcal S_+(s,k)\) and that the chosen stable chart identifies
the domain and target with fixed equal-dimensional fibres.

\begin{theorem}[Stable-root alternative after action restriction]
\label{thm:V35-stable-alternative}
On a compact homogeneous complex-frequency chart, exactly one of the
following holds:
\begin{enumerate}
\item
      \begin{equation}
       \rho_{35}
       :=
       \inf_{(s,k)}
       s_{\min}\mathbb B_{35}(s,k)>0,
       \label{eq:V35-full-floor}
      \end{equation}
      in which case the frozen boundary problem has a uniform
      action-compatible stable-root inverse on that chart;
\item there is a sequence of unit stable amplitudes \(v_j\) and frequencies
      \((s_j,k_j)\) such that
      \begin{equation}
       \|\mathbb B_{35}(s_j,k_j)v_j\|\longrightarrow0.
       \label{eq:V35-stable-witness}
      \end{equation}
\end{enumerate}
For a finite chart matrix the second alternative is equivalently returned
by right singular vectors at the vanishing singular floor.  A left
singular vector gives the V33 adjoint obstruction to solving inhomogeneous
boundary data.
\end{theorem}

\begin{proof}
The smallest singular value is continuous wherever the stable bundle has
been trivialized.  On a compact chart it has either a positive minimum or
a sequence tending to zero.  In the latter case choose corresponding
unit right singular vectors.  Compactness gives a convergent subsequence
of frequencies and the asserted near-kernel.  The singular-value
decomposition supplies the left witness and identifies it with the
orthogonal complement of the attained boundary range in the zero-floor
limit.
\end{proof}

Theorem~\ref{thm:V35-constraint-Lopatinski} proves that a witness in
\eqref{eq:V35-stable-witness} cannot arise from the constraint rows alone.
Proposition~\ref{prop:V35-YM-glancing} supplies such a witness for the
relative Yang--Mills action completion.  Proposition
\ref{prop:V35-gravity-action-obstruction} supplies a separate failure of
action incidence for the naive gravitational completion.

\subsection{Finite V35 certificate and cofinal handoff}
\label{subsec:V35-finite}

At cutoff \(q\), acquire on a homogeneous complex-frequency net
\[
 \widehat\lambda_q,\quad
 \widehat{\mathcal S}_{+,q},\quad
 \widehat{\mathsf D}_{E,q},\quad
 \widehat{\mathsf D}_{Y,q},\quad
 \widehat{\mathbb B}_{35,q},\quad
 \widehat A_q,\quad\widehat B_q
\]
in the common V24 equation and boundary-action frame.  Store
\begin{align}
 r_{\lambda,q}
 &:=
 \sup_{\overline{\mathbb S}_+}
 |\widehat\lambda_q^2-s^2-|k|^2|,
 \label{eq:V35-rlambda}\\
 r_{{\rm branch},q}
 &:=
 \sup_{\overline{\mathbb S}_+}
 \left[
 \bigl(-\operatorname{Re}\widehat\lambda_q\bigr)_+
 +\bigl(-\operatorname{Re}(s\overline{\widehat\lambda_q})\bigr)_+
 \right],
 \label{eq:V35-rbranch}\\
 r_{{\rm Som},q}
 &:=
 \left(
 2^{-1/2}
 -\inf_{\overline{\mathbb S}_+}
 |s+\widehat\lambda_q|
 \right)_+,
 \label{eq:V35-rSom}\\
 r_{{\rm Gram},q}
 &:=
 \sup_{\overline{\mathbb S}_+}
 \left\|
 \widehat{\mathsf D}_{E,q}
 \widehat{\mathsf D}_{E,q}^*
 -\frac12|\widehat\zeta_q|_{\rm E}^2I
 -\frac12\overline{(\eta\widehat\zeta_q)}\,
          \overline{(\eta\widehat\zeta_q)}^{\,*}
 \right\|
 \nonumber\\
 &\quad+
 \sup_{\overline{\mathbb S}_+}
 \left\|
 \widehat{\mathsf D}_{Y,q}
 \widehat{\mathsf D}_{Y,q}^*
 -|\widehat\zeta_q|_{\rm E}^2I
 \right\|,
 \label{eq:V35-rGram}\\
 \widehat\alpha_{{\rm act},q}
 &:=
 \inf_{\overline{\mathbb S}_+}
 s_{\min}\!\left(
   \begin{array}{cc}\widehat A_q&\widehat B_q\end{array}
 \right),
 \label{eq:V35-action-rank-floor}\\
 r_{{\rm Lag},q}
 &:=
 \sup_{\overline{\mathbb S}_+}
 \|\widehat A_q\widehat B_q^*
   -\widehat B_q\widehat A_q^*\|,
 \label{eq:V35-rLag}\\
 \widehat\rho_{35,q}
 &:=
 \min_{\rm net}
 s_{\min}(\widehat{\mathbb B}_{35,q}).
 \label{eq:V35-rhohat}
\end{align}
The net-to-continuum error must include the Hölder
\(\lambda\sim\operatorname{dist}(\cdot,\mathcal G)^{1/2}\) behavior near
glancing; an ordinary global Lipschitz estimate for \(\lambda\) is false.
Let that certified interpolation and frame error be
\(\varepsilon_{35,q}\), and define
\begin{equation}
 \rho_{35,q}^{\rm cert}
 :=
 \widehat\rho_{35,q}-\varepsilon_{35,q}.
 \label{eq:V35-rhocert}
\end{equation}

\begin{theorem}[Cofinal stable-boundary handoff]
\label{thm:V35-cofinal-handoff}
Suppose:
\begin{enumerate}
\item the V35 complex constraint residuals
      \eqref{eq:V35-rlambda}--\eqref{eq:V35-rGram} tend uniformly to zero;
\item \(r_{{\rm Lag},q}\to0\),
      \(\inf_q\widehat\alpha_{{\rm act},q}>0\), and the ghost,
      Higgs, and fermion action-domain residuals from V27/V32
      tend to zero in the same frame;
\item
      \(\inf_q\rho_{35,q}^{\rm cert}>0\);
\item the stable projectors, chart transitions, and boundary operators
      have summable adjacent tails, including the glancing
      square-root chart.
\end{enumerate}
Then the limiting frozen boundary family is simultaneously
constraint-preserving, action Lagrangian, and uniformly complementing on
the complete complex stable hemisphere.  Its inverse and chart
transitions furnish the boundary component required by the V33 nonlinear
repair theorem.

If instead a normalized sequence satisfies
\eqref{eq:V35-stable-witness}, the positive handoff fails and the sequence,
together with its left singular vectors, is the required obstruction
certificate.  The relative Yang--Mills conditions
\eqref{eq:V35-relative-YM} fall in this second alternative in the
unweighted norm by Proposition~\ref{prop:V35-YM-glancing}.
\end{theorem}

\begin{proof}
The vanishing root, branch, Sommerfeld, and Gram residuals pass
Theorem~\ref{thm:V35-constraint-Lopatinski} to the common limit.
The uniform action-rank floor and Lagrangian residual pass
Theorem~\ref{thm:V35-Lagrangian-criterion}.  The positive certified floor
and the singular-value perturbation inequality give a uniform inverse for
\(\mathbb B_{35}\).  Summable adjacent tails identify the stable bundles
and inverses across cutoffs.  This is precisely the fixed-frame boundary
input used by Theorem~\ref{thm:V33-cofinal}.

In the zero-floor case, Theorem~\ref{thm:V35-stable-alternative} supplies
the right and left singular witnesses.  Equation
\eqref{eq:V35-YM-quasimode} is the explicit witness for the relative
branch.
\end{proof}

\begin{firewall}[V35 physical status]
\label{fire:V35-status}
The complex harmonic and Lorenz constraint rows are now uniformly
certified through glancing, and the microlocal pivot atlas has an explicit
inverse.  The relative Yang--Mills/ghost branch passes the frozen action
and BRST tests but has an unweighted glancing singular sequence.  The
normal-pivot gravitational completion fails the action symplectic test.
The attached histories contain neither a replacement gravitational
boundary action nor a physical Julia reservoir, and still do not choose
the V32 neutrino branch.  The full action-restricted V33 floor and the
complete Standard Model--Einstein continuum are therefore not claimed.
\end{firewall}

\subsection{V35 ledger and next upstream row}
\label{subsec:V35-ledger}

\paragraph{New conclusions proved in V35.}
Theorem~\ref{thm:V35-Sommerfeld-factor} gives the sharp stable
Sommerfeld lower bound.  Proposition~\ref{prop:V35-complex-Gram} extends
the V34 contraction floors to complex stable and glancing frequencies.
Theorem~\ref{thm:V35-constraint-Lopatinski} proves uniform complementing
for the constraint rows.  Theorem~\ref{thm:V35-pivot-atlas} constructs
the four-chart completion and exact determinants.  Theorem
\ref{thm:V35-Lagrangian-criterion} gives the finite common-action test.
Theorem~\ref{thm:V35-relative-incidence} constructs the local
BRST-relative Yang--Mills/ghost branch, while Proposition
\ref{prop:V35-YM-glancing} returns its exact glancing witness.
Proposition~\ref{prop:V35-gravity-action-obstruction} returns the
symplectic obstruction for the naive gravitational completion.

\paragraph{Imported conclusions not repeated.}
V34 supplies the interior subsidiary equations, real Green matrices,
constraint readout reduction, and differentiated spinor boundary
identity.  V33 supplies the full action-restricted nonlinear repair and
left-null framework.  V27 supplies the Hamilton--Pontryagin
action-to-causal incidence, while V31--V32 supply the fermion neutral
graph and moving Higgs conormal.

\paragraph{Next upstream row.}
The next iteration should repair the gravitational boundary action rather
than change the already positive constraint symbol.  It should compute
the quadratic generalized-harmonic Einstein action including the
Gibbons--Hawking--York and gauge-fixing boundary terms, derive its exact
conormal matrix, and test that matrix with
\eqref{eq:V35-AB-condition}.  If no local Lagrangian completion has a
uniform glancing floor, it should construct a genuine boundary field
action realizing the Julia dilation and include its stress, BRST, and
cofinal tails.  The Higgs mixed conormal and the V32 neutrino branch must
then be joined to the same boundary action.

% END APPENDED VERSION V35 MATERIAL

% =============================================================================
% START APPENDED VERSION V36 MATERIAL
% =============================================================================

\section{V36: GHY--DeWitt conormal incidence, the normal gauge kernel,
and the action--Lopatinski tradeoff}
\label{sec:V36}

\paragraph{Purpose and relation to V35.}
V35 proved that the harmonic constraint rows themselves remain uniformly
surjective at complex stable and glancing frequencies.  It then tested the
algebraically convenient normal pivot
\(\bar h_{ab}=0\), \(C_\mu=0\) against a minimal wave Green form and found
an action defect.  The missing upstream distinction is that the
Gibbons--Hawking--York Dirichlet variable is the induced perturbation
\(h_{ab}\), not its trace reverse \(\bar h_{ab}\).  This note derives the
quadratic de Donder action metric and the GHY boundary one-form in the V34
frame, proves that
\[
 h_{ab}=0,\qquad C_\mu=0
\]
is an action-Lagrangian boundary relation, and computes its complete stable
matrix.  The action defect disappears, but a stronger analytic obstruction
appears: a normal boundary diffeomorphism gives an exact stable kernel at
every frequency in the open Laplace half-plane.  At glancing the kernel
dimension increases from one to two.

The standard GHY variation and canonical boundary momentum are reviewed in
\url{https://arxiv.org/abs/1605.01603}.  The de Donder quadratic action and
BRST boundary construction are consistent with
\url{https://arxiv.org/abs/gr-qc/9610023}.  The known Euclidean
non-ellipticity of fixed intrinsic metric plus DeWitt gauge, and its relation
to residual gauge freedom, are described in
\url{https://arxiv.org/abs/1301.0717}.  Those references provide context;
the Lorentzian stable-root calculation and every rank statement used below
are proved directly.

\subsection{Trace reversal and the principal gauge-fixed action}
\label{subsec:V36-DeWitt-action}

Retain the V34 frozen orthonormal frame
\((T,n,e_1,e_2)\), with
\(\eta=\operatorname{diag}(-1,1,1,1)\), and let Greek indices range over
all four directions.  Lower-case Latin indices
\(a,b\in\{T,1,2\}\) are tangent to the timelike boundary.  All fields and
test variations in this note are compactly supported in a boundary chart,
or vanish at the temporal and spatial edge of that chart, so tangential
integration by parts has no corner remainder.

For a symmetric covariant perturbation \(h\), define
\begin{equation}
 (\mathcal P h)_{\mu\nu}
 :=\bar h_{\mu\nu}
 :=h_{\mu\nu}
   -\frac12\eta_{\mu\nu}h,
 \qquad
 h:=\eta^{\rho\sigma}h_{\rho\sigma}.
 \label{eq:V36-trace-reversal}
\end{equation}
The DeWitt fibre form is
\begin{equation}
 \mathbb G(h,k)
 :=h_{\mu\nu}k^{\mu\nu}
   -\frac12hk
 =\bar h^{\mu\nu}k_{\mu\nu}.
 \label{eq:V36-DeWitt-form}
\end{equation}

\begin{lemma}[Four-dimensional trace reversal and DeWitt signature]
\label{lem:V36-DeWitt-involution}
On \(\operatorname{Sym}^2(\mathbb R^{1,3})^*\),
\begin{equation}
 \mathcal P^2=I,
 \qquad
 \mathcal P^*=\mathcal P,
 \qquad
 \|\mathcal P h\|_{\rm E}=\|h\|_{\rm E},
 \label{eq:V36-P-involution}
\end{equation}
where the adjoint and norm in this display use the Frobenius-orthonormal
component basis from V35.  The bilinear form \(\mathbb G\) is nondegenerate
and has signature \((6,4)\).
\end{lemma}

\begin{proof}
Taking an \(\eta\)-trace in \eqref{eq:V36-trace-reversal} gives
\(\operatorname{tr}_\eta\bar h=-h\), hence applying \(\mathcal P\) again
returns \(h\).  In the Euclidean coefficient space, put
\(e=\operatorname{diag}(-1,1,1,1)\).  Since \(\|e\|_{\rm E}^2=4\),
\[
 \mathcal P=I-\frac12 e\otimes e
\]
is the orthogonal reflection which is \(-I\) on \(\mathbb Re\) and
\(+I\) on \(e^\perp\).  This proves
\eqref{eq:V36-P-involution}.

On the four diagonal components, \(\mathbb G\) is
\(I-\tfrac12e\otimes e\), so it has three positive and one negative
direction.  Of the six off-diagonal components, the three with one time
index are negative and the three with two spatial indices are positive.
Thus the total signature is \((6,4)\), proving nondegeneracy.
\end{proof}

At the frozen point the linear harmonic covector is the V34 map
\begin{equation}
 C_\nu(h)
 :=\partial^\mu h_{\mu\nu}
   -\frac12\partial_\nu h
 =\partial^\mu\bar h_{\mu\nu}.
 \label{eq:V36-linear-C}
\end{equation}
Use the following first-derivative representative of the principal
Fierz--Pauli density:
\begin{align}
 \mathcal L_{\rm FP}^{\rm prin}
 &:=-\frac12\partial_\lambda h_{\mu\nu}
                    \partial^\lambda h^{\mu\nu}
   +(\partial_\mu h^{\mu\nu})(\partial^\lambda h_{\lambda\nu})
 \nonumber\\
 &\quad
   -(\partial_\mu h^{\mu\nu})\partial_\nu h
   +\frac12\partial_\lambda h\partial^\lambda h.
 \label{eq:V36-FP-density}
\end{align}
Its total divergences at the physical boundary are not discarded; the GHY
variation below fixes their boundary representative.

\begin{proposition}[Exact de Donder square cancellation]
\label{prop:V36-deDonder-square}
With unit de Donder gauge parameter,
\begin{equation}
 \boxed{
 \mathcal L_{\rm FP}^{\rm prin}-C_\nu C^\nu
 =-\frac12\partial_\lambda h_{\mu\nu}
                 \partial^\lambda\bar h^{\mu\nu}.}
 \label{eq:V36-minimal-density}
\end{equation}
The principal Euler operator is
\(\mathcal P\Box h\), equivalently \(\Box h\), and therefore agrees with
the V34 reduced wave block up to its fixed nonzero scalar normalization.
\end{proposition}

\begin{proof}
Put \(D_\nu=\partial^\mu h_{\mu\nu}\).  Then
\[
 C_\nu C^\nu
 =D_\nu D^\nu-D^\nu\partial_\nu h
   +\frac14\partial_\lambda h\partial^\lambda h.
\]
The middle three terms of \eqref{eq:V36-FP-density} minus this square leave
\(\tfrac14\partial h\cdot\partial h\).  Hence the left side of
\eqref{eq:V36-minimal-density} is
\[
 -\frac12\partial h_{\mu\nu}\cdot\partial h^{\mu\nu}
 +\frac14\partial h\cdot\partial h
 =-\frac12\partial h_{\mu\nu}\cdot
                  \partial\bar h^{\mu\nu}.
\]
Self-adjointness of \(\mathcal P\), followed by one integration by parts,
gives the Euler operator \(\mathcal P\Box h\).  Lemma
\ref{lem:V36-DeWitt-involution} makes this equivalent to \(\Box h=0\).
\end{proof}

\begin{firewall}[The DeWitt action is not the V27 positive energy]
\label{fire:V36-DeWitt-not-positive}
Lemma~\ref{lem:V36-DeWitt-involution} gives four negative fibre directions.
Thus action self-adjointness or a Lagrangian boundary relation does not by
itself give the positive Friedrichs estimate of V27.  The V24 equation
multiplier and the reduced physical quotient must still identify the action
Green form with the V34 causal flux.  V36 uses \(\mathbb G\) only for the
variational incidence which V35 left open.
\end{firewall}

\subsection{GHY plus gauge-fixing boundary one-form}
\label{subsec:V36-GHY-conormal}

Let \(\gamma\) be the metric induced on a non-null boundary and choose the
outward spacelike normal.  With the common normalization
\begin{equation}
 S_D[g]
 :=\frac1{2\kappa}\int_M\sqrt{-g}\,R
   +\frac1\kappa\int_{\partial M}\sqrt{|\gamma|}\,K,
 \label{eq:V36-Dirichlet-action}
\end{equation}
the GHY cancellation gives
\begin{equation}
 \delta S_D
 =-\frac1{2\kappa}\int_M\sqrt{-g}\,G^{\mu\nu}\delta g_{\mu\nu}
 +\frac1{2\kappa}\int_{\partial M}\sqrt{|\gamma|}\,
   (K^{ab}-K\gamma^{ab})\delta\gamma_{ab},
 \label{eq:V36-GHY-variation}
\end{equation}
with the bulk sign tied to variation of the covariant metric.  Reversing the
normal reverses both displayed boundary signs and changes none of the
Lagrangian or rank conclusions.

For completeness, add the background-harmonic gauge-fixing functional
\begin{equation}
 S_{\rm gf}[g]
 :=-\frac1{2\kappa}\int_M\sqrt{-g}\,C_\mu C^\mu.
 \label{eq:V36-gf-action}
\end{equation}
At the flat anchor \(C=0\), its quadratic boundary variation is determined
by
\begin{equation}
 \delta C_\nu
 =\partial^\mu\delta h_{\mu\nu}
  -\frac12\partial_\nu\delta h.
 \label{eq:V36-delta-C}
\end{equation}
Define
\begin{equation}
 \mathcal Q^{\mu\nu}(C)
 :=n^{(\mu}C^{\nu)}
   -\frac12\eta^{\mu\nu}C_n.
 \label{eq:V36-QC}
\end{equation}

\begin{theorem}[Frozen GHY--de Donder conormal]
\label{thm:V36-GHY-conormal}
At a planar flat anchor, the boundary one-form of the quadratic Hessian of
\(S_D+S_{\rm gf}\) is
\begin{equation}
 \boxed{
 \Theta_{36}(h;\dot h)
 =\frac1{2\kappa}\int_{\partial M}
 \left[
   \Pi_{(1)}^{ab}(h)\dot h_{ab}
   -2\mathcal Q^{\mu\nu}(C(h))\dot h_{\mu\nu}
 \right],}
 \label{eq:V36-boundary-one-form}
\end{equation}
where
\begin{align}
 K^{(1)}_{ab}(h)
 &:=\frac12
 \left(D_nh_{ab}-D_ah_{nb}-D_bh_{na}\right),
 \label{eq:V36-linear-K}\\
 \Pi_{(1)}^{ab}(h)
 &:=K_{(1)}^{ab}(h)-\eta^{ab}K_{(1)}(h).
 \label{eq:V36-linear-BY}
\end{align}
Thus the ten-component action conormal coefficient is
\begin{equation}
 \mathfrak p_{36}^{\mu\nu}(h)
 =e_a^\mu e_b^\nu\Pi_{(1)}^{ab}(h)
  -2\mathcal Q^{\mu\nu}(C(h)).
 \label{eq:V36-action-conormal}
\end{equation}
\end{theorem}

\begin{proof}
The Palatini identity writes the variation of
\(\sqrt{-g}R\) as the Einstein tensor term plus the divergence of a term
linear in \(\nabla\delta g\).  Varying
\(2\sqrt{|\gamma|}K\) cancels that normal derivative and leaves the
Brown--York coefficient in \eqref{eq:V36-GHY-variation}.  Linearizing
\(K_{ab}=\tfrac12\mathcal L_n\gamma_{ab}\) while allowing the lapse and
shift components of \(h\) gives \eqref{eq:V36-linear-K}, hence
\eqref{eq:V36-linear-BY}.

At \(C=0\), variation of \eqref{eq:V36-gf-action} has boundary part
\[
 -\frac1\kappa\int_{\partial M}
 n^\mu C^\nu
 \left(\dot h_{\mu\nu}
       -\frac12\eta_{\mu\nu}\dot h\right).
\]
Because \(\dot h\) is symmetric, its coefficient is exactly
\(-2\mathcal Q^{\mu\nu}(C)\) after the common factor
\((2\kappa)^{-1}\) is removed.  Adding the GHY term proves
\eqref{eq:V36-boundary-one-form} and
\eqref{eq:V36-action-conormal}.
\end{proof}

\begin{firewall}[Principal conormal scope]
\label{fire:V36-conormal-scope}
Equation~\eqref{eq:V36-action-conormal} is the exact quadratic conormal at
the planar flat anchor selected for the V34--V35 symbol calculation.
Background curvature, nonzero extrinsic curvature, generalized-harmonic
source derivatives, matter stress, and corner terms contribute lower-order
or separate corner rows.  They do not alter the principal rank computation
below, but they must be restored before a nonlinear common-action boundary
theorem is claimed.
\end{firewall}
\subsection{The fixed-induced-metric action branch}
\label{subsec:V36-action-branch}

Consider the homogeneous principal boundary relation
\begin{equation}
 \mathcal L_{\rm GHY}
 :=\left\{
 h_{ab}=0,\quad C_\mu(h)=0
 \right\}.
 \label{eq:V36-GHY-branch}
\end{equation}
It has six essential induced-metric rows and four gauge-fixing rows.  The
normal derivatives of \(h_{ab}\) remain free; the four equations
\(C_\mu=0\) solve for four independent combinations of the remaining normal
derivatives.

\begin{theorem}[Action-Lagrangian incidence of the GHY branch]
\label{thm:V36-GHY-Lagrangian}
On the real boundary field space with compact tangential support,
\(\mathcal L_{\rm GHY}\) is Lagrangian for the quadratic
GHY--de Donder Green form.  Equivalently, after passing to the canonical
DeWitt momentum, its ten boundary rows satisfy the V35 condition
\(AB^*=BA^*\) at every real tangential Fourier frequency.
\end{theorem}

\begin{proof}
For an allowed field and allowed variation,
\(C(h)=0\) and \(\dot h_{ab}=0\).  Both terms in
\eqref{eq:V36-boundary-one-form} therefore vanish.  Antisymmetrizing the
second variation shows that the Green symplectic form vanishes on
\(\mathcal L_{\rm GHY}\).  The six trace rows are independent.  In the
boundary phase variables \((h,D_nh)\), the normal part of
\(C_\nu\) is \(D_n\bar h_{n\nu}\); trace reversal is invertible, so those
four rows are independent modulo the first six.  Hence the relation has
codimension ten in a twenty-dimensional phase fibre.  It is isotropic and
half-dimensional, hence Lagrangian.
\end{proof}

The cancellation can also be seen without appealing to the variational
argument.  On \(h_{ab}=0\), put
\begin{equation}
 v_a:=h_{na},
 \qquad
 \phi:=h_{nn},
 \qquad
 \rho:=D_n(h^a{}_a).
 \label{eq:V36-v-phi}
\end{equation}
The boundary trace \(h^a{}_a\) is zero, but its normal derivative
\(\rho\) is free.  The harmonic rows are therefore
\begin{align}
 D_nv_a-\frac12D_a\phi&=0,
 \label{eq:V36-Ca-phase}\\
 \frac12(D_n\phi-\rho)+D^av_a&=0.
 \label{eq:V36-Cn-phase}
\end{align}
For a second allowed variation \((w_a,\psi)\), write
\(\sigma=D_n(k^a{}_a)\) for its free normal trace derivative.  The
restricted DeWitt Green form is
\begin{align}
 \omega_{\mathbb G}
 &=\int_{\partial M}
 \left[
  2(D_nv^a\,w_a-v^aD_nw_a)
  +\frac12\bigl((D_n\phi-\rho)\psi
                   -\phi(D_n\psi-\sigma)\bigr)
 \right].
 \label{eq:V36-restricted-Green}
\end{align}
Substitution of \eqref{eq:V36-Ca-phase}--%
\eqref{eq:V36-Cn-phase} gives
\[
 \int_{\partial M}
 \left[
  D^a\phi\,w_a-v^aD_a\psi
  -(D^av_a)\psi+\phi D^aw_a
 \right]=0
\]
after two tangential integrations by parts.  This explicitly verifies the
relative factors in the action conormal, including the normal derivative of
the fixed induced trace.

\begin{lemma}[Metric-domain BRST tangent at the flat anchor]
\label{lem:V36-BRST-tangent}
For the linearized diffeomorphism ghost \(c_\mu\),
\begin{align}
 s h_{\mu\nu}&=2\partial_{(\mu}c_{\nu)},
 \label{eq:V36-BRST-h}\\
 s C_\nu&=\Box c_\nu.
 \label{eq:V36-BRST-C}
\end{align}
At \(K_{ab}=0\), Dirichlet tangential ghost trace \(c_a=0\) preserves
\(h_{ab}=0\); the ghost Euler equation preserves \(C_\nu=0\).
\end{lemma}

\begin{proof}
Insert \eqref{eq:V36-BRST-h} in \eqref{eq:V36-linear-C}:
\[
 sC_\nu
 =\partial^\mu(\partial_\mu c_\nu+\partial_\nu c_\mu)
  -\partial_\nu\partial^\mu c_\mu
 =\Box c_\nu.
\]
On the boundary,
\(s h_{ab}=2D_{(a}c_{b)}+2K_{ab}c_n\).  At the planar anchor the second
term vanishes, while a zero boundary trace has zero tangential derivative.
The two claims follow.
\end{proof}

\begin{firewall}[The ghost domain does not remove a metric kernel]
\label{fire:V36-ghost-domain}
Lemma~\ref{lem:V36-BRST-tangent} proves only tangency of the metric rows.
A full BRST action domain must also choose compatible ghost, antighost, and
Nakanishi--Lautrup rows and test their Green form.  Whether the normal ghost
trace is allowed determines whether the stable mode below is quotiented as
a gauge transformation or retained as a boundary edge mode.  Either choice
leaves the unquotiented metric boundary matrix singular; changing the ghost
domain alone does not supply its missing metric row.
\end{firewall}

\subsection{Exact complex stable matrix and its rank defect}
\label{subsec:V36-stable-matrix}

Use the V35 stable ansatz and root
\begin{equation}
 u=e^{st+ik\cdot y-\lambda x}u_0,
 \qquad
 \lambda^2=s^2+|k|^2,
 \qquad
 \operatorname{Re}s>0,
 \quad\operatorname{Re}\lambda>0.
 \label{eq:V36-stable-root}
\end{equation}
Let
\begin{equation}
 d_a:=(s,ik_1,ik_2),
 \qquad
 d^ad_a=-s^2-|k|^2=-\lambda^2.
 \label{eq:V36-tangential-null}
\end{equation}
For a stable exponential, \(h_{ab}=0\) also gives \(D_nh_{ab}=0\), so
\(\rho=0\).  The remaining four amplitudes are
\((v_a,\phi)\).  Equations
\eqref{eq:V36-Ca-phase}--\eqref{eq:V36-Cn-phase} have stable matrix
\begin{equation}
 \mathcal H_D(s,k)
 :=
 \begin{pmatrix}
  \lambda I_3&-\frac12d\\
  d^\sharp&\frac12\lambda
 \end{pmatrix},
 \label{eq:V36-HD}
\end{equation}
where \(d\) is a three-component column and
\(d^\sharp=(-s,ik_1,ik_2)\) is a row.  Thus
\(d^\sharp d=-\lambda^2\).

\begin{theorem}[Exact open-half-plane rank, right kernel, and left witness]
\label{thm:V36-open-rank}
For every normalized stable frequency with \(\operatorname{Re}s>0\),
\(\lambda\ne0\) and
\begin{equation}
 \operatorname{rank}\mathcal H_D=3.
 \label{eq:V36-HD-rank}
\end{equation}
An exact right kernel and left kernel are
\begin{align}
 r_D(s,k)
 &:=
 \binom{d/(2\lambda)}{1},
 &\mathcal H_Dr_D&=0,
 \label{eq:V36-right-kernel}\\
 \ell_D(s,k)
 &:=
 \left(-d^\sharp/\lambda,\ 1\right),
 &\ell_D\mathcal H_D&=0.
 \label{eq:V36-left-kernel}
\end{align}
Consequently the complete ten-by-ten stable boundary matrix obtained by
adjoining the six induced-metric selectors has rank nine, determinant zero,
and smallest singular value zero at every open stable frequency.
\end{theorem}

\begin{proof}
If \(\lambda=0\), then \(s^2+|k|^2=0\), which forces \(s\) to be purely
imaginary and contradicts \(\operatorname{Re}s>0\).  The upper-left block
of \(\mathcal H_D\) therefore has rank three.  Direct multiplication and
\eqref{eq:V36-tangential-null} give
\begin{align*}
 \lambda\frac d{2\lambda}-\frac d2&=0,\\
 d^\sharp\frac d{2\lambda}+\frac\lambda2
 &=\frac{-\lambda^2}{2\lambda}+\frac\lambda2=0,
\end{align*}
proving the right kernel.  Similarly, the first three columns in
\(\ell_D\mathcal H_D\) vanish, and its last entry is
\[
 \frac{d^\sharp d}{2\lambda}+\frac\lambda2=0.
\]
Thus the rank is at most three and at least three, proving
\eqref{eq:V36-HD-rank}.

Order the full amplitude columns as the six \(h_{ab}\) components followed
by \((v,\phi)\).  The six selector rows form an identity block and the
Schur remainder is \(\mathcal H_D\).  Hence the full rank is
\(6+3=9\).  Its determinant and smallest singular value vanish.
\end{proof}

For inhomogeneous induced data \(g_{ab}\), eliminate the first six columns.
If the resulting harmonic right-hand side is \(f=(f_a,f_n)\), the exact
compatibility condition is
\begin{equation}
 f_n-\lambda^{-1}d^\sharp f_a=0.
 \label{eq:V36-inhomogeneous-compatibility}
\end{equation}
For the original ten rows the corresponding V33 left witness is obtained by
placing \(\ell_D\) in the last four row slots and adding the unique six
selector coefficients that cancel its action on the eliminated
\(h_{ab}\)-columns.

\begin{theorem}[Glancing rank drop]
\label{thm:V36-glancing-rank}
At a V35 glancing frequency, \(\lambda=0\) and \(d\ne0\).  Then
\begin{equation}
 \operatorname{rank}\mathcal H_D=2,
 \qquad
 \ker\mathcal H_D
 =\left\{
 (v,\phi):\ \phi=0,\ d^\sharp v=0
 \right\},
 \label{eq:V36-glancing-kernel}
\end{equation}
so the reduced kernel has dimension two.  The full ten-row boundary matrix
has rank eight.
\end{theorem}

\begin{proof}
Putting \(\lambda=0\) in \eqref{eq:V36-HD} gives
\[
 \mathcal H_D(0,d)
 =\begin{pmatrix}0&-d/2\\d^\sharp&0\end{pmatrix}.
\]
The last column has rank one because \(d\ne0\), and the last row has rank
one because \(d^\sharp\ne0\).  They act between disjoint domain and target
summands, so the total rank is two.  The displayed kernel follows.  Adding
the six independent induced selectors gives full rank eight.
\end{proof}
\subsection{Pure-gauge origin and the factor-through-constraint no-go}
\label{subsec:V36-gauge-origin}

\begin{proposition}[The stable kernel is a normal gauge mode]
\label{prop:V36-pure-gauge}
Fix an open stable frequency and let
\begin{equation}
 \xi_n
 :=\frac1{2\lambda}
 e^{st+ik\cdot y-\lambda x},
 \qquad
 \xi_a:=0.
 \label{eq:V36-normal-gauge-vector}
\end{equation}
Using the derivative convention in which \(D_n\) has stable symbol
\(\lambda\), set
\begin{equation}
 h^{(\xi)}_{\mu\nu}
 :=2\partial_{(\mu}\xi_{\nu)}.
 \label{eq:V36-gauge-mode}
\end{equation}
Then \(h^{(\xi)}\) solves the V34 reduced homogeneous wave equation,
obeys \(h^{(\xi)}_{ab}=0\), and satisfies
\(C_\nu(h^{(\xi)})=0\) identically in the half-space.  Its boundary
amplitude is precisely \(r_D\) in \eqref{eq:V36-right-kernel}.
\end{proposition}

\begin{proof}
The scalar factor in \eqref{eq:V36-normal-gauge-vector} obeys the wave
equation by \eqref{eq:V36-stable-root}; derivatives commute with the
frozen wave operator, so every component of \(h^{(\xi)}\) also obeys it.
Because \(\xi_a=0\), the tangential--tangential components vanish.  The
remaining amplitudes are
\[
 h^{(\xi)}_{nn}=2\lambda\frac1{2\lambda}=1,
 \qquad
 h^{(\xi)}_{na}=d_a\frac1{2\lambda},
\]
which is \(r_D\).  Finally, for an arbitrary vector \(\xi\),
\[
 C_\nu(2\partial_{(\cdot}\xi_{\cdot)})
 =\Box\xi_\nu.
\]
The chosen \(\xi\) is a wave, hence \(C_\nu=0\) throughout the
half-space.
\end{proof}

\begin{theorem}[No homogeneous subsidiary operator can repair fixed
intrinsic metric]
\label{thm:V36-factor-no-go}
Let \(\mathcal R\) be any homogeneous linear boundary differential or
pseudodifferential operator whose metric dependence factors through the
harmonic covector and its jets:
\begin{equation}
 \mathcal R(h)=\mathcal T(C(h),\nabla C(h),\ldots).
 \label{eq:V36-factor-through-C}
\end{equation}
Then the boundary system
\begin{equation}
 h_{ab}=0,
 \qquad
 \mathcal R(h)=0
 \label{eq:V36-general-subsidiary-branch}
\end{equation}
admits the nonzero decaying stable mode
\eqref{eq:V36-gauge-mode} at every frequency with
\(\operatorname{Re}s>0\).  Thus replacing \(C=0\) by a Dirichlet, Robin,
Sommerfeld, or higher-jet homogeneous boundary condition on \(C\) alone
cannot restore the metric Lopatinski inverse.
\end{theorem}

\begin{proof}
Proposition~\ref{prop:V36-pure-gauge} gives
\(h^{(\xi)}_{ab}=0\) and \(C(h^{(\xi)})\equiv0\) in the entire
half-space.  Every jet of \(C\) therefore vanishes, so
\(\mathcal R(h^{(\xi)})=0\).  The mode is nonzero and decays normally
because \(\operatorname{Re}\lambda>0\).  It is an exact stable kernel,
which rules out injectivity and hence any Lopatinski inverse.
\end{proof}

\begin{firewall}[Gauge in the bulk versus gauge at the boundary]
\label{fire:V36-bulk-boundary-gauge}
The word ``gauge'' in Proposition~\ref{prop:V36-pure-gauge} records the
bulk differential identity \(h=2\partial_{(\mu}\xi_{\nu)}\).  It does not
by itself license quotienting this mode.  If the boundary gauge group
requires \(\xi_n|_{\partial M}=0\), the mode is not an admissible gauge
transformation and remains a genuine nonunique solution of the gauge-fixed
boundary problem.  If nonzero normal boundary diffeomorphisms are allowed,
the quotient, ghost domain, boundary embedding variable, and V24 frame must
all remove the same direction.  Neither choice produces an unrecorded
inverse of the ten-row matrix.
\end{firewall}

\begin{corollary}[No constant-rank full-kernel quotient through glancing]
\label{cor:V36-no-kernel-bundle}
The exact stable kernel of the fixed-intrinsic-metric branch has dimension
one for \(\operatorname{Re}s>0\) and dimension two at glancing.  Therefore
the full kernels do not form a constant-rank vector bundle on the closed
stable hemisphere.  A quotient by the entire kernel cannot satisfy the V24
fixed-domain chart hypothesis through glancing without an additional
boundary gauge or edge-field resolution.
\end{corollary}

\begin{proof}
The dimensions are Theorems~\ref{thm:V36-open-rank} and
\ref{thm:V36-glancing-rank}.  Fibre dimension is locally constant for a
vector bundle.  Since it jumps from one to two, the family of full kernels
is not such a bundle.  Quotienting only the open-half-plane line leaves an
additional glancing null direction and therefore does not give a uniform
inverse on the closed hemisphere.
\end{proof}

Nonzero background \(K_{ab}\) can lift the exact planar mode at finite
frequency because
\(s h_{ab}=2D_{(a}c_{b)}+2K_{ab}c_n\).  It cannot repair the principal
boundary symbol: under high-frequency rescaling, the \(K\)-terms are lower
order and the homogeneous matrix remains \eqref{eq:V36-HD}.  Hence V36
returns a principal uniform obstruction, not a claim that every curved
finite-frequency matrix has an exact kernel.

\subsection{Why the V35 pivot and the GHY selector fail differently}
\label{subsec:V36-selector-tradeoff}

On the GHY essential subspace \(h_{ab}=0\), one has
\begin{equation}
 h=\phi,
 \qquad
 \bar h_{ab}=-\frac12\eta_{ab}\phi,
 \qquad
 \bar h_{na}=v_a,
 \qquad
 \bar h_{nn}=\frac12\phi.
 \label{eq:V36-trace-reversed-GHY}
\end{equation}
Thus the GHY selector \(h_{ab}=0\) and the V35 selector
\(\bar h_{ab}=0\) are different six-planes.  Trace reversal is invertible,
but it does not carry one boundary plane into itself.

\begin{proposition}[Exact selector tradeoff]
\label{prop:V36-selector-tradeoff}
Let \(Q_Dh=(h_{ab})\) and \(Q_{\bar D}h=(\bar h_{ab})\).  On the normal
gauge vector \(r_D\),
\begin{align}
 Q_Dr_D&=0,
 \label{eq:V36-QD-kernel}\\
 \|Q_{\bar D}r_D\|_{\rm E}&=\frac{\sqrt3}{2}.
 \label{eq:V36-Qbar-kernel}
\end{align}
However,
\begin{equation}
 \|r_D\|_{\rm E}^2
 =1+\frac{|d|_{\rm E}^2}{2|\lambda|^2},
 \qquad
 \frac{\|Q_{\bar D}r_D\|_{\rm E}}
      {\|r_D\|_{\rm E}}
 \longrightarrow0
 \quad\hbox{as }\lambda\to0.
 \label{eq:V36-Qbar-angle}
\end{equation}
Therefore the GHY selector is exactly action-compatible but does not see the
open stable gauge line.  The V35 trace-reversed selector sees that line
pointwise for \(\lambda\ne0\), but loses a uniform transverse angle at
glancing; it also fails the V35 action test.
\end{proposition}

\begin{proof}
The first two claims follow from \(h_{ab}=0\) and
\eqref{eq:V36-trace-reversed-GHY}: the three diagonal tangential components
of \(\bar h\) each have absolute value \(1/2\), while the off-diagonal
components vanish.  In the Frobenius tensor norm, the mixed components
\(h_{na}=d_a/(2\lambda)\) occur twice, giving
\[
 \|r_D\|_{\rm E}^2
 =1+2\frac{|d|_{\rm E}^2}{4|\lambda|^2}.
\]
On approach to glancing, \(|d|_{\rm E}\) stays nonzero while
\(|\lambda|\to0\), which proves \eqref{eq:V36-Qbar-angle}.  The action
statements are Theorem~\ref{thm:V36-GHY-Lagrangian} and Proposition
\ref{prop:V35-gravity-action-obstruction}.
\end{proof}

\begin{theorem}[Necessary gauge transversality for every completion]
\label{thm:V36-gauge-transversality}
Let \(S(s,k)\) be any six-row supplementary metric selector and define the
full stable matrix
\begin{equation}
 \mathcal B_S(s,k)
 :=\binom{S(s,k)}{\mathsf D_E(\zeta)\mathcal P}.
 \label{eq:V36-general-completion}
\end{equation}
At an open stable frequency, injectivity of \(\mathcal B_S\) requires
\begin{equation}
 S(s,k)r_D(s,k)\ne0.
 \label{eq:V36-open-transversality}
\end{equation}
At glancing, injectivity requires that \(S\) be injective on the
2-dimensional plane in \eqref{eq:V36-glancing-kernel}.  A uniform stable
floor requires a positive lower angle on the normalized versions of both
families.
\end{theorem}

\begin{proof}
The harmonic block annihilates \(r_D\) by
Theorem~\ref{thm:V36-open-rank}.  If the supplementary block also
annihilates it, then the full matrix has a nonzero kernel.  This proves
\eqref{eq:V36-open-transversality}.  At glancing the same argument applies
to every vector in \(\ker\mathcal H_D\).  For the quantitative assertion,
if \(u\) is a unit vector in either kernel, then
\[
 s_{\min}(\mathcal B_S)
 \le\|\mathcal B_Su\|=\|Su\|.
\]
Taking the infimum proves that a vanishing transverse angle forces a
vanishing stable floor.
\end{proof}

\begin{corollary}[The next search is finite and action-restricted]
\label{cor:V36-next-search}
A candidate replacing the GHY selector must simultaneously satisfy:
\begin{enumerate}
\item the action-Lagrangian equation of V35 in the DeWitt canonical frame;
\item nonzero open-frequency action on \(r_D\);
\item injectivity on the glancing plane
      \(\{(v,0):d^\sharp v=0\}\);
\item BRST tangency with a recorded ghost and boundary-edge domain;
\item a positive V34 causal floor after the V24 equation multiplier.
\end{enumerate}
Conditions 2 and 3 are explicit finite singular-value tests.  Theorems
\ref{thm:V36-factor-no-go} and \ref{thm:V36-gauge-transversality} show that
altering only the homogeneous subsidiary row cannot satisfy them.
\end{corollary}

\begin{firewall}[No contradiction with V35]
\label{fire:V36-no-V35-contradiction}
V35 proved uniform surjectivity of the four-row harmonic contraction and a
chartwise algebraic completion.  V36 proves singularity of a particular
ten-row intersection: the harmonic rows plus the GHY induced-metric
selector.  A surjective four-row map can have a singular ten-row completion.
The two results address different matrices and are fully compatible.
\end{firewall}
\subsection{Finite V36 card and the cofinal negative decision}
\label{subsec:V36-finite-card}

At cutoff \(q\), store in the common V24 trace frame
\begin{equation}
 \widehat{\mathcal P}_q,
 \quad
 \widehat{\mathbb G}_q,
 \quad
 \widehat\Theta_{36,q},
 \quad
 \widehat C_q,
 \quad
 \widehat Q_{D,q},
 \quad
 \widehat{\mathcal B}_{D,q},
 \quad
 \widehat r_{D,q},
 \quad
 \widehat\ell_{D,q}.
 \label{eq:V36-finite-packet}
\end{equation}
The last two vectors are normalized right and left candidates for the full
stable boundary matrix, including the six selector components in the left
candidate; the latter is stored as a column whose Hermitian adjoint is the
algebraic left row.  The formulas containing \(1/\lambda\) are used only
where \(\lambda\ne0\).  Chartwise one stores the equivalent polynomial
representatives obtained by multiplication by \(\lambda\), followed by
Euclidean normalization.  Thus the open right line extends to the
distinguished glancing subline spanned by \((d/2,0)\) inside
\eqref{eq:V36-glancing-kernel}; this statement does not identify that line
with the entire two-dimensional glancing kernel.  On real tangential Fourier
covectors let
\begin{equation}
 \widehat{\mathbb J}_{\mathbb G,q}
 :=
 \begin{pmatrix}
 0&-\widehat{\mathbb G}_q\\
 \widehat{\mathbb G}_q&0
 \end{pmatrix},
 \label{eq:V36-DeWitt-J}
\end{equation}
and write the phase-space boundary rows as
\(\widehat M_q=(\widehat A_q\ \widehat B_q)\).  Record
\begin{align}
 r_{\mathcal P,q}
 &:=\|\widehat{\mathcal P}_q^2-I\|
    +\|\widehat{\mathcal P}_q^*-\widehat{\mathcal P}_q\|,
 \label{eq:V36-rP}\\
 r_{\mathbb G,q}
 &:=\|\widehat{\mathbb G}_q^2-I\|
    +\|\widehat{\mathbb G}_q^*-\widehat{\mathbb G}_q\|,
 \label{eq:V36-rG}\\
 r_{{\rm Lag},q}
 &:=\sup_{\tau,k\in\mathbb R}
 \|\widehat M_q(i\tau,k)
   \widehat{\mathbb J}_{\mathbb G,q}
   \widehat M_q(i\tau,k)^*\|,
 \label{eq:V36-rLag}\\
 r_{{\rm root},q}
 &:=\sup_{\overline{\mathbb S}_+}
 \left|\widehat\lambda_q^2-s^2-|k|^2\right|,
 \label{eq:V36-rroot}\\
 r_{{\rm R},q}
 &:=\sup_{\overline{\mathbb S}_+}
 \|\widehat{\mathcal B}_{D,q}\widehat r_{D,q}\|,
 \label{eq:V36-rR}\\
 r_{{\rm L},q}
 &:=\sup_{\overline{\mathbb S}_+}
 \|\widehat\ell_{D,q}^*
       \widehat{\mathcal B}_{D,q}\|,
 \label{eq:V36-rL}\\
 \widehat\rho_{10,q}
 &:=\inf_{\overline{\mathbb S}_+}
 s_{\min}(\widehat{\mathcal B}_{D,q}).
 \label{eq:V36-rho10}
\end{align}
The real-frequency supremum in \eqref{eq:V36-rLag} is homogeneous and is
taken on the normalized real cotangent sphere.  Hermitian adjoints there
represent formal tangential adjoints.  They are not substituted for the
complex Laplace transpose in
\eqref{eq:V36-left-kernel}.

For rank localization also store, on compact charts separated from
\(\mathcal G\), the ninth singular value, and on glancing the eighth:
\begin{align}
 \widehat\rho_{9,q}(\delta)
 &:=\inf_{\substack{\overline{\mathbb S}_+\\
                    |\lambda|\ge\delta}}
 s_9(\widehat{\mathcal B}_{D,q}),
 \label{eq:V36-rho9}\\
 \widehat\rho_{8,q}^{\mathcal G}
 &:=\inf_{\mathcal G}
 s_8(\widehat{\mathcal B}_{D,q}).
 \label{eq:V36-rho8}
\end{align}
Here singular values are in descending order.  Positive values of these two
quantities certify that the only losses are the one open stable direction
and the second glancing direction already identified; they do not repair
the tenth singular value.

\begin{theorem}[Cofinal obstruction from a normalized gauge witness]
\label{thm:V36-cofinal-obstruction}
Suppose the finite GHY--de Donder matrices converge uniformly in the common
V24 frame and there are normalized witnesses with
\(r_{{\rm R},q}\to0\).  Then
\begin{equation}
 \limsup_{q\to\infty}\widehat\rho_{10,q}=0.
 \label{eq:V36-zero-cofinal-floor}
\end{equation}
If also \(r_{{\rm L},q}\to0\), the limiting right failure carries a V33
adjoint compatibility obstruction.  If
\(r_{\mathcal P,q}+r_{\mathbb G,q}+r_{{\rm Lag},q}\to0\) and the action-row
rank has a uniform floor, then the limiting boundary relation is action
Lagrangian despite this analytic failure.  Therefore the fixed-intrinsic
GHY branch cannot enter the positive-floor alternative of V35 or V33.
\end{theorem}

\begin{proof}
For every normalized right candidate,
\[
 s_{\min}(\widehat{\mathcal B}_{D,q})
 \le
 \|\widehat{\mathcal B}_{D,q}\widehat r_{D,q}\|.
\]
Taking the frequency infimum and then the limit proves
\eqref{eq:V36-zero-cofinal-floor}.  The left residual converges to a vector
orthogonal to the limiting range and therefore gives the V33 compatibility
functional.  The involution, action-rank, and Lagrangian residuals pass the
finite criterion of Theorem~\ref{thm:V35-Lagrangian-criterion}; uniform
matrix convergence supplies the required limit.  This establishes the
simultaneous action-positive and stable-negative decision.
\end{proof}

\begin{proposition}[Certified selector-angle obstruction]
\label{prop:V36-angle-certificate}
For a proposed supplementary selector \(S_q\), let
\(\widehat K_{D,1,q}(s,k)\) be the certified continuation of the
one-dimensional line spanned by \(r_D\) on the open half-plane, and let
\(\widehat K_{D,2,q}(s,k)\) be the certified continuation of the residual
two-plane in \eqref{eq:V36-glancing-kernel}.  Define
\begin{align}
 \alpha_{1,q}
 &:=\inf_{\operatorname{Re}s>0}
 \inf_{\substack{u\in\widehat K_{D,1,q}\\\|u\|=1}}
 \|S_qu\|,
 \label{eq:V36-alpha1}\\
 \alpha_{2,q}^{\mathcal G}
 &:=\inf_{\mathcal G}
 \inf_{\substack{u\in\widehat K_{D,2,q}\\\|u\|=1}}
 \|S_qu\|.
 \label{eq:V36-alpha2}
\end{align}
If either certified lower bound tends to zero, the full stable floor tends
to zero.  Positive values are necessary but not sufficient, because the
selector must also complement the rest of the six-dimensional harmonic
kernel.  A positive V33 handoff requires uniform positive lower bounds for
both displayed angles, the full complementing floor, and the action and BRST
residuals.
\end{proposition}

\begin{proof}
This is the singular-value inequality in the proof of
Theorem~\ref{thm:V36-gauge-transversality}, first on the open stable line
and then on the glancing kernel plane.  Approximate kernels add their
recorded residuals to the right-hand side; convergence sends those errors
to zero.
\end{proof}

\begin{firewall}[V36 physical status]
\label{fire:V36-status}
The quadratic GHY--de Donder conormal and its action-Lagrangian
fixed-intrinsic boundary relation are now populated.  That branch has an
exact normal-gauge stable kernel and a two-dimensional glancing kernel, so
its full stable floor is zero.  V36 neither deletes the mode by an
unrecorded quotient nor calls it a physical graviton.  It does not yet
construct an alternative action boundary density, an edge-mode action, or
a curved nonlinear BRST domain.  The full Standard Model--Einstein
continuum remains unproved.
\end{firewall}

\subsection{V36 ledger and next upstream row}
\label{subsec:V36-ledger}

\paragraph{New conclusions proved in V36.}
Lemma~\ref{lem:V36-DeWitt-involution} fixes the exact action fibre and its
signature.  Proposition~\ref{prop:V36-deDonder-square} derives the minimal
gauge-fixed wave Hessian.  Theorem~\ref{thm:V36-GHY-conormal} computes the
quadratic GHY--de Donder conormal, and Theorem
\ref{thm:V36-GHY-Lagrangian} proves action incidence of the induced-metric
branch.  Theorems~\ref{thm:V36-open-rank} and
\ref{thm:V36-glancing-rank} return its exact rank-nine and rank-eight
stable decisions.  Proposition~\ref{prop:V36-pure-gauge} identifies the
normal gauge generator.  Theorem~\ref{thm:V36-factor-no-go} proves that no
homogeneous condition factoring only through \(C\) repairs it.  Theorem
\ref{thm:V36-gauge-transversality} gives the finite necessary test for any
replacement selector.

\paragraph{Imported conclusions not repeated.}
V34 supplies the reduced metric wave equation, harmonic subsidiary system,
and independent incoming-row count.  V35 supplies the complex stable root,
constraint contraction floors, pivot atlas, and finite Lagrangian
criterion.  V24 and V27 supply the common trace frame and the distinction
between the indefinite action Hessian and the reduced positive causal
symmetrizer.  V33 supplies the right/left obstruction handoff.

\paragraph{Next upstream row.}
The next iteration should test the first genuinely different variational
branch rather than modify the subsidiary operator.  In the same Lorentzian
stable-root frame it should derive the boundary Hessian produced by a
boundary gauge functional such as
\[
 F_a(h)=D^bh_{ab}-\frac12D_ah^b{}_b,
\]
or by fixing the linearized Brown--York/extrinsic-curvature momentum.  It
must compute the resulting ten-row stable determinant, the two selector
angles \eqref{eq:V36-alpha1}--\eqref{eq:V36-alpha2}, the DeWitt action
Lagrangian residual, and the BRST ghost incidence.  The Euclidean
strong-ellipticity claim in \url{https://arxiv.org/abs/1301.0717} is not a
Lorentzian proof and must be recalculated.  If every local branch still
loses the glancing angle, introduce an explicit boundary embedding or
reservoir field with its own action and stress before returning to the
Higgs conormal and V32 neutrino choice.

% END APPENDED VERSION V36 MATERIAL


% BEGIN APPENDED VERSION V37 MATERIAL

\clearpage
\section{V37: conformal--mean-curvature polarization, open complementing
determinant, and glancing collapse}
\label{sec:V37-conformal-mean-curvature}

\subsection{Purpose, sign ledger, and relation to V36}
\label{subsec:V37-purpose}

V36 separated two questions which had previously been conflated.  The
fixed-induced-metric GHY branch is Lagrangian for the quadratic action Green
form, but its Lorentzian stable boundary matrix has an exact normal-gauge
kernel.  The present note tests the two alternatives left in the V36 ledger.
First, it proves that a boundary gauge functional depending only on the
intrinsic tensor \(h_{ab}\), including

\begin{equation}
 F_a(h):=D^bh_{ab}-\frac12D_a h^b{}_b,
 \label{eq:V37-intrinsic-F}
\end{equation}

cannot meet the V36 normal-gauge transversality test.  Second, it passes to
the genuinely different conformal--mean-curvature polarization: fix the
trace-free induced metric and the mean curvature, and retain the four
harmonic rows.  This branch is action-Lagrangian and has a nonzero
open-half-plane determinant.  Its determinant is nevertheless

\begin{equation}
 \det\mathcal M_{CM}=\frac13\lambda^5,
 \label{eq:V37-headline-determinant}
\end{equation}

in the reduced trace chart, and the full boundary rank drops from ten to
seven at glancing.  Thus the branch repairs the exact open kernel without
producing a uniform Kreiss--Lopatinski floor.

Conformal class plus mean curvature is standard geometric boundary data in
the Riemannian Einstein problem; see
\url{https://arxiv.org/abs/1103.1152} and
\url{https://arxiv.org/abs/1805.11559}.  Lorentzian cavity results require a
separate analysis.  In particular, the explicit ill-posedness construction
in \url{https://arxiv.org/abs/2505.20410} and the revised geometric-boundary
analysis in \url{https://arxiv.org/abs/2503.12599} show why a Euclidean
ellipticity statement cannot be imported as a Lorentzian uniform-boundary
estimate.  The calculation below is local, planar, and in the V35 harmonic
stable frame; it neither imports nor reproves either global result.

There is one sign convention which must be made explicit.  Put

\begin{equation}
 \Pi^{ab}:=K^{ab}-K\gamma^{ab},
 \qquad
 \Pi:=\gamma_{ab}\Pi^{ab}=-2K
 \label{eq:V37-Pi-trace}
\end{equation}

for a three-dimensional timelike boundary.  Introduce
\(\varepsilon_D\in\{+1,-1\}\) by writing the metric part of the Dirichlet
boundary one-form as

\begin{equation}
 \Theta_D^{\rm met}
 =\frac{\varepsilon_D}{2\kappa}
   \int_{\partial M}\sqrt{|\gamma|}\,
       \Pi^{ab}\delta\gamma_{ab}.
 \label{eq:V37-epsilon-D}
\end{equation}

The displayed convention in V36 corresponds to \(\varepsilon_D=+1\).
Equation (2.4) of \url{https://arxiv.org/abs/1605.01603}, with its stated
choice of curvature and normal signs and with a covariant induced-metric
variation, corresponds to \(\varepsilon_D=-1\).  Retaining
\(\varepsilon_D\) prevents an orientation or curvature convention from
being mistaken for a boundary-rank statement.  Every homogeneous boundary
relation and every stable rank below is independent of this sign.

\begin{firewall}[What V37 corrects and what it preserves]
\label{fire:V37-sign-firewall}
The source cited above carries a minus sign in front of
\((K^{ab}-K\gamma^{ab})\delta\gamma_{ab}\), whereas
\eqref{eq:V36-GHY-variation} used the opposite boundary-one-form convention.
V37 therefore does not use that sign without a parameter.  This convention
audit changes neither V36's branch equations \(h_{ab}=0,C_\mu=0\), nor its
rank-nine/rank-eight calculation, nor its Lagrangian conclusion: the
boundary one-form vanishes on that branch for either value of
\(\varepsilon_D\).  It does change the coefficient with which one writes the
York generating term, so that coefficient is derived below rather than
copied from V36.
\end{firewall}

\subsection{The intrinsic-jet no-go}
\label{subsec:V37-intrinsic-no-go}

Let \(J_{\partial}^{N}h^{\mathsf T}\) denote the boundary tangential jet of
the induced perturbation \(h^{\mathsf T}=(h_{ab})\) through order \(N\).
An intrinsic homogeneous boundary operator is a linear operator of the form

\begin{equation}
 \mathfrak F_N(h)
 =L_N(D_a)h_{bc},
 \label{eq:V37-intrinsic-jet-operator}
\end{equation}

with no \(h_{na}\), \(h_{nn}\), or normal derivative among its arguments.
Equation~\eqref{eq:V37-intrinsic-F} is the first-order example selected in
the V36 ledger.

\begin{theorem}[Intrinsic tangential jets miss the V36 normal gauge line]
\label{thm:V37-intrinsic-no-go}
For every \(N<\infty\) and every operator
\eqref{eq:V37-intrinsic-jet-operator},

\begin{equation}
 \mathfrak F_N\bigl(h^{(\xi)}\bigr)=0
 \label{eq:V37-intrinsic-annihilation}
\end{equation}

on the exact normal pure-gauge stable mode of
Proposition~\ref{prop:V36-pure-gauge}.  In particular,
\(F_a(h^{(\xi)})=0\).  Adding or substituting only such intrinsic rows cannot
give a positive selector angle on the open V36 gauge line.
\end{theorem}

\begin{proof}
The V36 generator has \(\xi_a=0\) at the planar anchor.  Hence

\[
 h^{(\xi)}_{ab}=2D_{(a}\xi_{b)}=0
\]

identically as a function on the boundary.  Every tangential derivative of
this zero tensor is zero, proving
\eqref{eq:V37-intrinsic-annihilation}.  The necessary angle condition is
Theorem~\ref{thm:V36-gauge-transversality}.
\end{proof}

\begin{corollary}[Role left for the intrinsic harmonic functional]
\label{cor:V37-F-role}
The row \(F_a=0\) may fix tangential reparametrization inside a larger mixed
or edge-field boundary system.  It cannot by itself replace the missing
normal-gauge row.  Any viable local selector must use a normal/extrinsic
quantity, or must enlarge the boundary field space so that the normal
diffeomorphism has a recorded boundary representative.
\end{corollary}

\subsection{York generating term and conformal--mean-curvature action}
\label{subsec:V37-York-action}

For an induced-metric variation define

\begin{equation}
 \vartheta:=\gamma^{ab}\delta\gamma_{ab},
 \qquad
 \delta\gamma^{\rm TF}_{ab}
 :=\delta\gamma_{ab}-\frac13\gamma_{ab}\vartheta.
 \label{eq:V37-trace-splitting}
\end{equation}

Because \(\Pi=-2K\), the exact algebraic splitting is

\begin{equation}
 \Pi^{ab}\delta\gamma_{ab}
 =\Pi_{\rm TF}^{ab}\delta\gamma^{\rm TF}_{ab}
  -\frac23K\vartheta.
 \label{eq:V37-Pi-splitting}
\end{equation}

Also

\begin{equation}
 \delta\bigl(\sqrt{|\gamma|}K\bigr)
 =\sqrt{|\gamma|}
   \left(\delta K+\frac12K\vartheta\right).
 \label{eq:V37-delta-rootgamma-K}
\end{equation}

Starting with the V36 Dirichlet plus de Donder action, define the
sign-covariant York action

\begin{equation}
 S_{CM}^{(\varepsilon_D)}
 :=S_D+S_{\rm gf}
   +\frac{2\varepsilon_D}{3\kappa}
      \int_{\partial M}\sqrt{|\gamma|}\,K.
 \label{eq:V37-CM-action}
\end{equation}

For \(\varepsilon_D=-1\), the total coefficient of the GHY term is one
third of the Dirichlet coefficient, the usual four-dimensional conformal
choice.  For the internal V36 boundary-one-form convention
\(\varepsilon_D=+1\), the same generating transformation is written with
total coefficient five thirds.  The two descriptions produce the same
fixed-\([\gamma]\), fixed-\(K\) polarization after the sign convention is
used consistently.

\begin{theorem}[Conformal--mean-curvature first variation and Green form]
\label{thm:V37-CM-one-form}
At principal order near the planar flat anchor, the boundary part of the
nonlinear first variation of \eqref{eq:V37-CM-action} is

\begin{equation}
 \boxed{
 \Theta_{CM}(g;\delta g)
 =\frac1{2\kappa}\int_{\partial M}\sqrt{|\gamma|}
 \left[
  \varepsilon_D\Pi_{\rm TF}^{ab}\delta\gamma^{\rm TF}_{ab}
  +\frac{4\varepsilon_D}{3}\,\delta K
  -2\mathcal Q^{\mu\nu}(C)\delta g_{\mu\nu}
 \right].}
 \label{eq:V37-CM-one-form}
\end{equation}

No undifferentiated trace-variation term remains.  Its frozen Hessian Green
two-form on perturbations \(h,k\) is

\begin{align}
 2\kappa\,\Omega_{CM}(h,k)
 &=\int_{\partial M}
 \bigg[
 \varepsilon_D\!
 \left(
  \Pi_{(1),\rm TF}^{ab}(h)k^{\rm TF}_{ab}
  -\Pi_{(1),\rm TF}^{ab}(k)h^{\rm TF}_{ab}
 \right)
 \nonumber\\
 &\qquad
 +\frac{2\varepsilon_D}{3}
 \left(
  (h^a{}_a)K_{(1)}(k)
  -(k^a{}_a)K_{(1)}(h)
 \right)
 \nonumber\\
 &\qquad
 -2\left(
  \mathcal Q^{\mu\nu}(C(h))k_{\mu\nu}
  -\mathcal Q^{\mu\nu}(C(k))h_{\mu\nu}
 \right)
 \bigg].
 \label{eq:V37-CM-Green-form}
\end{align}
\end{theorem}

\begin{proof}
Multiply \eqref{eq:V37-Pi-splitting} by
\(\varepsilon_D/(2\kappa)\).  The added term in
\eqref{eq:V37-CM-action}, after extracting the same common factor, is

\[
 \frac1{2\kappa}
 \left[
  \frac{4\varepsilon_D}{3}\delta K
  +\frac{2\varepsilon_D}{3}K\vartheta
 \right].
\]

Its last summand cancels the trace term in
\eqref{eq:V37-Pi-splitting}.  The de Donder contribution is the last term of
\eqref{eq:V36-boundary-one-form} and is unaffected by the boundary
generating function.  This proves \eqref{eq:V37-CM-one-form}.

To obtain the Hessian Green form, differentiate the one-form in two field
directions and antisymmetrize.  At the flat anchor,
\(\delta_h\sqrt{|\gamma|}=(h^a{}_a)/2\).  Hence the
\((4\varepsilon_D/3)\sqrt{|\gamma|}\,\delta K\) term contributes
\[
 \frac{2\varepsilon_D}{3}
 \left((h^a{}_a)K_{(1)}(k)-(k^a{}_a)K_{(1)}(h)\right).
\]
The antisymmetrized linearizations of the other two terms give the first
and third lines of \eqref{eq:V37-CM-Green-form}.
\end{proof}

Define the homogeneous conformal--mean-curvature branch

\begin{equation}
 \mathcal L_{CM}
 :=\left\{
 h^{\rm TF}_{ab}=0,
 \quad K_{(1)}(h)=0,
 \quad C_\mu(h)=0
 \right\}.
 \label{eq:V37-CM-branch}
\end{equation}

There are five trace-free induced-metric rows, one mean-curvature row, and
four harmonic rows.

\begin{theorem}[Action-Lagrangian incidence of the CM branch]
\label{thm:V37-CM-Lagrangian}
For compactly supported real tangential data,
\(\mathcal L_{CM}\) is Lagrangian for the quadratic DeWitt Green form of
\(S_{CM}^{(\varepsilon_D)}\).  The conclusion holds for both values of
\(\varepsilon_D\).
\end{theorem}

\begin{proof}
If \(h,k\) are tangent to \(\mathcal L_{CM}\), then both have zero
trace-free induced metric, zero linearized mean curvature, and zero harmonic
constraint.  Every term in \eqref{eq:V37-CM-Green-form} vanishes, proving
isotropy.

It remains to check dimension.  In boundary phase coordinates
\((h,D_nh)\), the five trace-free rows are independent configuration rows.
Modulo them, the normal-principal parts of the remaining five rows solve in
succession for

\[
 D_n(h^a{}_a),\qquad D_nh_{na}\ (a=0,1,2),\qquad D_nh_{nn}.
\]

Thus the ten rows are independent in the twenty-dimensional phase fibre.
The relation is isotropic and half-dimensional, hence Lagrangian.
\end{proof}

The Green cancellation can be checked directly.  Put

\begin{equation}
 h_{ab}=\frac{\tau}{3}\eta_{ab},
 \qquad v_a=h_{na},
 \qquad \phi=h_{nn}.
 \label{eq:V37-reduced-fields}
\end{equation}

The phase rows in \eqref{eq:V37-CM-branch} give

\begin{align}
 D_n\tau&=2D^av_a,
 \label{eq:V37-phase-tau}\\
 D_nv_a&=\frac16D_a\tau+\frac12D_a\phi,
 \label{eq:V37-phase-v}\\
 D_n\phi&=0.
 \label{eq:V37-phase-phi}
\end{align}

For a second allowed field write
\(k_{ab}=\sigma\eta_{ab}/3\), \(w_a=k_{na}\), and
\(\psi=k_{nn}\).  The DeWitt bilinear form restricted to configuration
traces is

\begin{equation}
 \mathbb G(h,k)
 =2v^aw_a-\frac16\tau\sigma
  -\frac12\tau\psi-\frac12\phi\sigma+\frac12\phi\psi.
 \label{eq:V37-restricted-DeWitt}
\end{equation}

The free trace-free part of \(D_nh_{ab}\) pairs with
\(k^{\rm TF}_{ab}=0\) and drops out.  Substituting
\eqref{eq:V37-phase-tau}--\eqref{eq:V37-phase-phi} into
\(\mathbb G(D_nh,k)-\mathbb G(h,D_nk)\) gives

\begin{align}
 &\frac13\left(D^a\tau\,w_a-v^aD_a\sigma
              -(D^av_a)\sigma+\tau D^aw_a\right)
 \nonumber\\
 &\qquad
 +D^a\phi\,w_a-v^aD_a\psi
 -(D^av_a)\psi+\phi D^aw_a.
 \label{eq:V37-Green-divergence}
\end{align}

The first line is one third of
\(D^a(\tau w_a-v_a\sigma)\); the second is
\(D^a(\phi w_a-v_a\psi)\).  Its integral vanishes, independently verifying
Theorem~\ref{thm:V37-CM-Lagrangian} and all relative factors in
\eqref{eq:V37-phase-tau}--\eqref{eq:V37-phase-phi}.

\subsection{BRST tangency and its off-shell limit}
\label{subsec:V37-BRST}

\begin{proposition}[Planar linear BRST tangent]
\label{prop:V37-BRST-tangent}
At a flat totally geodesic boundary, the linearized BRST transformations
obey

\begin{align}
 s h^{\rm TF}_{ab}
 &=2\bigl(D_{(a}c_{b)}\bigr)^{\rm TF},
 \label{eq:V37-BRST-TF}\\
 sK_{(1)}&=-D^aD_a c_n,
 \label{eq:V37-BRST-K}\\
 sC_\mu&=\Box c_\mu.
 \label{eq:V37-BRST-C}
\end{align}

Consequently the Dirichlet ghost trace \(c_\mu|_{\partial M}=0\) preserves
the first two metric rows, and the bulk ghost equation preserves the third.
With

\begin{equation}
 c_\mu=0,
 \qquad \bar c_\mu=0,
 \qquad B_\mu=0
 \quad\hbox{on }\partial M,
 \label{eq:V37-quartet-boundary}
\end{equation}

the linear transformations \(s\bar c_\mu=B_\mu\), \(sB_\mu=0\) are tangent
to the auxiliary quartet rows.
\end{proposition}

\begin{proof}
The first and third equations are the trace-free part of
\eqref{eq:V36-BRST-h} and \eqref{eq:V36-BRST-C}.  For the second, insert
\(h_{ab}=D_ac_b+D_bc_a\) and
\(h_{na}=D_nc_a+D_ac_n\) into

\[
 K_{(1)}=\frac12\left(D_nh^a{}_a-2D^ah_{na}\right).
\]

Commutation of flat derivatives cancels the terms containing \(c_a\) and
leaves \(-D^aD_ac_n\).  A field whose boundary trace is zero has zero
tangential derivatives there.  This proves tangency of the first two rows;
the other statements follow from the displayed transformations.
\end{proof}

\begin{firewall}[On-shell tangency is not an off-shell BRST action domain]
\label{fire:V37-BRST-offshell}
Dirichlet data for \(c_\mu\) do not imply
\((\Box c_\mu)|_{\partial M}=0\) for an arbitrary off-shell ghost.  The
preservation of \(C_\mu=0\) above uses the ghost Euler equation.  Imposing
both \(c_\mu=0\) and \(\Box c_\mu=0\) as independent boundary data would not
be the ordinary half-dimensional domain of a second-order ghost operator.
Thus Proposition~\ref{prop:V37-BRST-tangent} populates the linear on-shell
BRST tangent and the auxiliary rows, but it does not claim a nonlinear
off-shell BV--BFV construction.
\end{firewall}

\subsection{Exact Lorentzian stable matrix}
\label{subsec:V37-stable-matrix}

Use the V35--V36 stable ansatz

\begin{equation}
 e^{st+ik\cdot y-\lambda x},
 \qquad
 \lambda^2=s^2+|k|^2,
 \qquad
 \operatorname{Re}s>0,
 \quad \operatorname{Re}\lambda>0,
 \label{eq:V37-stable-root}
\end{equation}

and retain

\begin{equation}
 d=(s,ik_1,ik_2)^T,
 \qquad
 d^\sharp=(-s,ik_1,ik_2),
 \qquad
 d^\sharp d=-\lambda^2.
 \label{eq:V37-d-symbol}
\end{equation}

After the five rows \(h^{\rm TF}_{ab}=0\), the amplitudes are
\((\tau,v_0,v_1,v_2,\phi)\).  The stable symbols of the remaining rows are

\begin{align}
 K_{(1)}&=\frac\lambda2\tau-d^\sharp v,
 \label{eq:V37-K-symbol}\\
 C_a&=\lambda v_a-\frac16d_a\tau-\frac12d_a\phi,
 \label{eq:V37-Ca-symbol}\\
 C_n&=d^\sharp v+\frac\lambda2(\phi-\tau).
 \label{eq:V37-Cn-symbol}
\end{align}

Thus, with rows ordered as \((K,C_a,C_n)\), the reduced five-by-five matrix
is

\begin{equation}
 \mathcal M_{CM}(s,k)
 :=
 \begin{pmatrix}
  \lambda/2&-d^\sharp&0\\
  -d/6&\lambda I_3&-d/2\\
  -\lambda/2&d^\sharp&\lambda/2
 \end{pmatrix}.
 \label{eq:V37-MCM}
\end{equation}

\begin{theorem}[Open complementing determinant and explicit inverse]
\label{thm:V37-open-determinant}
For every stable frequency with \(\operatorname{Re}s>0\),

\begin{equation}
 \det\mathcal M_{CM}=\frac13\lambda^5\ne0.
 \label{eq:V37-det-MCM}
\end{equation}

Hence the complete ten-row matrix has rank ten at every open stable
frequency.  If the data in the row order of
\eqref{eq:V37-MCM} are \((g,f_a,f_n)\), its solution is

\begin{align}
 \phi&=\frac{2(f_n+g)}{\lambda},
 \label{eq:V37-inverse-phi}\\
 \tau&=\frac{3}{2\lambda}
 \left(\frac{d^\sharp f}{\lambda}-f_n\right),
 \label{eq:V37-inverse-tau}\\
 v&=\frac1\lambda
 \left(f+\frac d6\tau+\frac d2\phi\right).
 \label{eq:V37-inverse-v}
\end{align}
\end{theorem}

\begin{proof}
Add the first row of \eqref{eq:V37-MCM} to the last.  The new last row is
\((0,0,\lambda/2)\).  After eliminating its last column, the remaining
Schur complement in \((\tau,v)\) has determinant

\[
 \det(\lambda I_3)
 \left(\frac\lambda2-\frac{d^\sharp d}{6\lambda}\right)
 =\lambda^3\left(\frac\lambda2+\frac\lambda6\right)
 =\frac23\lambda^4.
\]

Multiplication by the last pivot \(\lambda/2\) proves
\eqref{eq:V37-det-MCM}.  The stable root cannot have \(\lambda=0\) in the
open Laplace half-plane.  The same row operation gives
\(\lambda\phi/2=f_n+g\).  Substitution into the \(C_a\) rows and then the
\(K\) row gives \eqref{eq:V37-inverse-phi}--%
\eqref{eq:V37-inverse-v}.  The five trace-free selector rows form an
identity block in the chosen chart, so the full rank is \(5+5=10\).
\end{proof}

\begin{firewall}[Pointwise complementing is not uniform complementing]
\label{fire:V37-pointwise-not-uniform}
The factors \(\lambda^{-1}\) and \(\lambda^{-2}\) in
\eqref{eq:V37-inverse-phi}--\eqref{eq:V37-inverse-v} are part of the exact
boundary inverse.  They are not harmless chart notation: glancing belongs
to the closure of the normalized stable hemisphere.  Equation
\eqref{eq:V37-det-MCM} therefore proves pointwise open injectivity, not a
uniform closed-hemisphere estimate.
\end{firewall}

\subsection{Glancing kernel, cokernel, and rank-seven decision}
\label{subsec:V37-glancing}

\begin{theorem}[Exact glancing kernel and cokernel]
\label{thm:V37-glancing-rank}
At a glancing frequency \(\lambda=0\), \(d\ne0\), the reduced matrix has

\begin{equation}
 \operatorname{rank}\mathcal M_{CM}=2.
 \label{eq:V37-glancing-reduced-rank}
\end{equation}

Its right kernel is

\begin{equation}
 \ker\mathcal M_{CM}
 =\left\{
  (\tau,v,\phi):
  d^\sharp v=0,
  \quad \tau+3\phi=0
 \right\},
 \qquad \dim=3.
 \label{eq:V37-glancing-right-kernel}
\end{equation}

With algebraic transpose pairing and left coefficients ordered as
\((\alpha,\beta^T,\gamma)\) for \((K,C_a,C_n)\), the left kernel is

\begin{equation}
 \ker\mathcal M_{CM}^{T}
 =\left\{
  (\alpha,\beta^T,\gamma):
  \gamma=\alpha,
  \quad \beta^Td=0
 \right\},
 \qquad \dim=3.
 \label{eq:V37-glancing-left-kernel}
\end{equation}

The complete ten-row boundary matrix has rank seven and a
three-dimensional stable kernel at glancing.
\end{theorem}

\begin{proof}
Putting \(\lambda=0\) in \eqref{eq:V37-MCM} leaves

\begin{equation}
 \mathcal M_{CM}(0,d)
 =\begin{pmatrix}
  0&-d^\sharp&0\\
  -d/6&0&-d/2\\
  0&d^\sharp&0
 \end{pmatrix}.
 \label{eq:V37-MCM-glancing}
\end{equation}

The first and last rows span one row, while the three middle rows are
multiples of the nonzero column \(d\) and span one further row.  They act on
disjoint column combinations, so the rank is exactly two.  The equations
are \(d^\sharp v=0\) and \(d(\tau+3\phi)=0\), proving
\eqref{eq:V37-glancing-right-kernel}.

A left combination vanishes precisely when its coefficient on the
\(v\)-columns is \((\gamma-\alpha)d^\sharp=0\) and its coefficient on the
\((\tau,\phi)\)-columns is proportional to \(\beta^Td=0\).  This proves
\eqref{eq:V37-glancing-left-kernel}.  Restoring the five independent
trace-free selector rows gives full rank \(5+2=7\).
\end{proof}

For the original ten rows, a reduced left vector in
\eqref{eq:V37-glancing-left-kernel} has a unique full lift: add five
trace-free-selector coefficients so that its action on the eliminated
\(h^{\rm TF}_{ab}\)-columns vanishes.  This is the left-kernel Schur lift
through the selector identity block.  The right vectors are embedded by
setting their five trace-free amplitudes to zero.

For inhomogeneous data \((g,f,f_n)\), the glancing compatibility conditions
are therefore

\begin{equation}
 g+f_n=0,
 \qquad
 \beta^Tf=0
 \quad\hbox{for every }\beta\hbox{ with }\beta^Td=0.
 \label{eq:V37-glancing-compatibility}
\end{equation}

Equivalently, \(f\) must lie in the algebraic line spanned by \(d\).  These
are three independent V33 left obstructions: one from \(K+C_n\), and two
from the quotient of the \(C_a\) data by \(\operatorname{span}\{d\}\).

\begin{corollary}[The V36 glancing plane survives]
\label{cor:V37-V36-plane-survives}
The two-dimensional V36 glancing plane

\[
 \left\{(\tau,v,\phi):\tau=\phi=0,\ d^\sharp v=0\right\}
\]

is contained in \eqref{eq:V37-glancing-right-kernel}.  The CM branch adds
the scalar glancing direction \((\tau,v,\phi)=(-3,0,1)\).
\end{corollary}

\begin{firewall}[No physical polarization count from the raw kernel]
\label{fire:V37-glancing-physical}
The three vectors in \eqref{eq:V37-glancing-right-kernel} are algebraic
stable boundary amplitudes.  V37 identifies the distinguished limit of the
V36 normal gauge line below, but it does not classify every remaining
glancing vector as gauge or physical.  Such a classification requires the
recorded ghost/edge domain and the V24 physical quotient.
\end{firewall}

\subsection{Exact selector angle and the generalized conformal family}
\label{subsec:V37-angle-family}

Let \(S_{CM}h=(h^{\rm TF}_{ab},K_{(1)}(h))\) denote the six geometric rows.
On the V36 normal gauge vector \(r_D\),

\begin{equation}
 h^{\rm TF}_{ab}(r_D)=0,
 \qquad
 K_{(1)}(r_D)=\frac\lambda2.
 \label{eq:V37-SCM-on-rD}
\end{equation}

Using the full symmetric-tensor Euclidean norm from V36 gives the exact
normalized angle

\begin{equation}
 \frac{\|S_{CM}r_D\|_{\rm E}}{\|r_D\|_{\rm E}}
 =\frac{|\lambda|^2}
        {\sqrt{4|\lambda|^2+2|d|_{\rm E}^2}}
 \longrightarrow0
 \quad\hbox{as }\lambda\to0.
 \label{eq:V37-CM-angle}
\end{equation}

\begin{theorem}[Open repair with quadratic angle loss]
\label{thm:V37-open-repair-angle-loss}
The conformal--mean-curvature selector detects the V36 normal gauge line at
every open stable frequency, but its normalized angle tends to zero
quadratically in \(|\lambda|\) relative to nonzero limiting \(d\).  At
glancing it annihilates the entire V36 residual plane.  Consequently the
full CM boundary matrix has

\begin{equation}
 \inf_{\overline{\mathbb S}_+}
 s_{\min}(\mathcal B_{CM})=0.
 \label{eq:V37-zero-uniform-floor}
\end{equation}
\end{theorem}

\begin{proof}
For \(r_D=(0,d/(2\lambda),1)\) in the reduced trace chart,

\[
 K_{(1)}(r_D)
 =-\frac{d^\sharp d}{2\lambda}=\frac\lambda2.
\]

The full tensor norm is
\(\|r_D\|_{\rm E}^2=1+|d|_{\rm E}^2/(2|\lambda|^2)\), which proves
\eqref{eq:V37-CM-angle}.  At glancing, every V36 residual vector has
\(h^{\rm TF}_{ab}=0\) and \(K_{(1)}=-d^\sharp v=0\).  The singular-value
witness inequality and Theorem~\ref{thm:V37-glancing-rank} prove
\eqref{eq:V37-zero-uniform-floor}.
\end{proof}

The generalized conformal boundary family often fixes a density

\begin{equation}
 \mathcal H_p(\gamma,K):=|\gamma|^pK
 \label{eq:V37-generalized-density}
\end{equation}

together with the conformal class.  This is the family considered, in
Euclidean and Lorentzian cavity settings, in
\url{https://arxiv.org/abs/2402.04308} and
\url{https://arxiv.org/abs/2505.20410}.

\begin{proposition}[Principal \(p\)-universality]
\label{prop:V37-p-universality}
At a background \((\gamma_0,K_0)\),

\begin{equation}
 \delta\mathcal H_p
 =|\gamma_0|^p
  \left(\delta K+pK_0\gamma_0^{ab}\delta\gamma_{ab}\right).
 \label{eq:V37-density-linearization}
\end{equation}

At the planar anchor \(K_0=0\), every \(p\) has exactly the stable matrix
\eqref{eq:V37-MCM}.  At a curved anchor the \(pK_0\)-term is order zero and
does not change the high-frequency principal boundary symbol.  No fixed
finite \(p\) restores the missing glancing principal rank.
\end{proposition}

\begin{proof}
Differentiate the determinant density and use
\(\delta\log|\gamma|=\gamma^{ab}\delta\gamma_{ab}\).  The first assertion is
exact.  Mean-curvature variation is first order in the metric perturbation,
whereas multiplication by the background \(K_0\) contains no derivative.
The latter therefore disappears from the homogeneous principal symbol.
\end{proof}

\begin{firewall}[Relation to the Lorentzian cavity literature]
\label{fire:V37-literature-scope}
Proposition~\ref{prop:V37-p-universality} is a local principal-symbol
statement at a planar frozen boundary.  It is consistent with the
high-angular-frequency instability established on spherical and cylindrical
cavities in \url{https://arxiv.org/abs/2505.20410}, but it is not a proof of
that global instability.  Conversely, the dense-range and Holmgren-type
results in the revised 2026 version of
\url{https://arxiv.org/abs/2503.12599}, which also records corner-angle data,
do not supply the uniform harmonic stable floor required by V24--V35.  The
claims concern different estimates and are not identified.
\end{firewall}

\subsection{Finite V37 card and cofinal decision}
\label{subsec:V37-finite-card}

At cutoff \(q\), store in the common V24 frame

\begin{equation}
 \widehat\varepsilon_{D,q},
 \quad
 \widehat P_{{\rm TF},q},
 \quad
 \widehat K_q,
 \quad
 \widehat C_q,
 \quad
 \widehat\Theta_{CM,q},
 \quad
 \widehat{\mathcal M}_{CM,q},
 \quad
 \widehat{\mathcal B}_{CM,q},
 \quad
 \widehat R_{\mathcal G,q},
 \quad
 \widehat L_{\mathcal G,q}.
 \label{eq:V37-finite-packet}
\end{equation}

The columns of \(R_{\mathcal G,q}\) are the full embeddings of a basis
of \eqref{eq:V37-glancing-right-kernel}.  The rows of
\(L_{\mathcal G,q}\) include the five Schur-lift coefficients and
approximate full left witnesses whose reduced components form a basis of
\eqref{eq:V37-glancing-left-kernel}.  Record

\begin{align}
 r_{{\rm det},q}
 &:=\sup_{\overline{\mathbb S}_+}
 \left|3\det\widehat{\mathcal M}_{CM,q}
       -\widehat\lambda_q^5\right|,
 \label{eq:V37-rdet}\\
 r_{{\rm R}\mathcal G,q}
 &:=\sup_{\mathcal G}
 \|\widehat{\mathcal B}_{CM,q}
       \widehat R_{\mathcal G,q}\|,
 \label{eq:V37-rRG}\\
 r_{{\rm L}\mathcal G,q}
 &:=\sup_{\mathcal G}
 \|\widehat L_{\mathcal G,q}
       \widehat{\mathcal B}_{CM,q}\|,
 \label{eq:V37-rLG}\\
 \widehat\rho_{10,q}^{CM}
 &:=\inf_{\overline{\mathbb S}_+}
 s_{\min}(\widehat{\mathcal B}_{CM,q}),
 \label{eq:V37-rho10}\\
 \widehat\rho_{7,q}^{\mathcal G}
 &:=\inf_{\mathcal G}
 s_7(\widehat{\mathcal B}_{CM,q}),
 \label{eq:V37-rho7}\\
 \widehat\alpha_{D,q}^{CM}
 &:=\inf_{\operatorname{Re}s>0}
 \frac{\|\widehat S_{CM,q}\widehat r_{D,q}\|}
      {\|\widehat r_{D,q}\|}.
 \label{eq:V37-alphaD}
\end{align}

Singular values are ordered decreasingly.  A positive certified
\(\widehat\rho_{7,q}^{\mathcal G}\), together with vanishing eighth through
tenth glancing values, localizes the exact rank-seven decision.  On compact
charts \(|\lambda|\ge\delta>0\), also store the inverse residual

\begin{equation}
 r_{{\rm inv},q}(\delta)
 :=\sup_{|\lambda|\ge\delta}
 \left\|
 \widehat{\mathcal M}_{CM,q}^{-1}
 \widehat{\mathcal M}_{CM,q}-I_5
 \right\|.
 \label{eq:V37-rinv}
\end{equation}

The inverse is never evaluated at glancing.

\begin{theorem}[Cofinal CM obstruction]
\label{thm:V37-cofinal-obstruction}
Assume uniform convergence in the common frame, normalized columns of
\(\widehat R_{\mathcal G,q}\), and

\[
 r_{{\rm R}\mathcal G,q}\longrightarrow0.
\]

Then

\begin{equation}
 \limsup_{q\to\infty}\widehat\rho_{10,q}^{CM}=0.
 \label{eq:V37-cofinal-zero-floor}
\end{equation}

If the V35 action-rank and Lagrangian residuals also tend to zero, the limit
is an action-Lagrangian boundary relation with no uniform stable inverse.
If \(r_{{\rm L}\mathcal G,q}\to0\), it carries the three independent
glancing compatibility functionals in
\eqref{eq:V37-glancing-compatibility}.
\end{theorem}

\begin{proof}
For every normalized column \(u\) of
\(\widehat R_{\mathcal G,q}\),

\[
 s_{\min}(\widehat{\mathcal B}_{CM,q})
 \le
 \|\widehat{\mathcal B}_{CM,q}u\|.
\]

Take the glancing infimum and the limit.  This proves
\eqref{eq:V37-cofinal-zero-floor}.  The action statement is the V35 finite
Lagrangian criterion applied to the one-form
\eqref{eq:V37-CM-one-form}.  Convergence of the normalized left rows gives
annihilators of the limiting range, and
Theorem~\ref{thm:V37-glancing-rank} identifies their dimension and
compatibility equations.
\end{proof}

\subsection{Boundary-embedding seed and the next upstream row}
\label{subsec:V37-edge-seed}

The preceding failures have complementary meanings.  Purely intrinsic
metric jets never see the normal boundary diffeomorphism.  Mean curvature
sees it for \(\lambda\ne0\), but with an angle that collapses at glancing.
The next branch should therefore record the boundary displacement rather
than silently forbid or quotient it.

Introduce a scalar normal embedding perturbation \(\zeta\) with linear BRST
law

\begin{equation}
 s\zeta=-c_n|_{\partial M}.
 \label{eq:V37-zeta-BRST}
\end{equation}

\begin{proposition}[Gauge-invariant normal extrinsic seed]
\label{prop:V37-edge-seed}
At the planar flat anchor,

\begin{equation}
 \mathcal K_\zeta
 :=K_{(1)}-D^aD_a\zeta
 \label{eq:V37-Kzeta}
\end{equation}

is invariant under the linear BRST differential:

\begin{equation}
 s\mathcal K_\zeta=0.
 \label{eq:V37-Kzeta-invariant}
\end{equation}
\end{proposition}

\begin{proof}
Equations~\eqref{eq:V37-BRST-K} and
\eqref{eq:V37-zeta-BRST} give

\[
 s\mathcal K_\zeta
 =-D^aD_ac_n-D^aD_a(-c_n)=0.
\]
\end{proof}

\begin{firewall}[The seed is not yet an edge-mode theorem]
\label{fire:V37-edge-seed-status}
Adding \(\zeta\) increases the boundary phase space.  A valid continuum row
must supply its boundary kinetic term, canonical momentum, stress tensor,
ghost partner, quotient projector, and causal symmetrizer.  Merely replacing
\(K\) by \eqref{eq:V37-Kzeta} makes a gauge-invariant expression; it does not
by itself prove a positive physical-quotient floor.  No such floor is
claimed in V37.
\end{firewall}

The next iteration should derive the minimal local boundary action for
\(\zeta\) (and, if required, tangential embedding components), antisymmetrize
its coupled bulk--boundary Green form, and compute the extended stable
matrix before quotienting.  The decisive finite tests are:

\begin{enumerate}
\item constant rank of the coupled gauge orbit through glancing;
\item a BRST-compatible projector onto the enlarged physical trace bundle;
\item a positive V34 causal singular floor on that quotient;
\item equality between the edge stress obtained by metric variation and the
      boundary term in the coupled Bianchi identity;
\item preservation of the V31 Higgs source conormal and the V32 neutrino
      branch when the gravitational edge action is restored.
\end{enumerate}

If no local boundary kinetic coefficient passes these tests, the failure is
then specific enough to justify a nonlocal transparent reservoir; it is not
yet justified by V36 or V37 alone.

\subsection{V37 ledger}
\label{subsec:V37-ledger}

\paragraph{New conclusions proved in V37.}
Theorem~\ref{thm:V37-intrinsic-no-go} eliminates all purely intrinsic
tangential-jet repairs of the V36 normal gauge line.
Theorem~\ref{thm:V37-CM-one-form} derives the sign-audited York boundary
one-form, and Theorem~\ref{thm:V37-CM-Lagrangian} proves the resulting
action polarization.  Theorem~\ref{thm:V37-open-determinant} proves the
exact open determinant and inverse.  Theorem~\ref{thm:V37-glancing-rank}
proves the rank-seven glancing collapse and its right and left kernels.
Theorem~\ref{thm:V37-open-repair-angle-loss} proves the exact selector-angle
loss, while Proposition~\ref{prop:V37-p-universality} extends the principal
negative decision to the generalized conformal density family.  Proposition
\ref{prop:V37-edge-seed} supplies the first gauge-invariant embedding
combination for the next branch.

\paragraph{Imported conclusions not repeated.}
V35 supplies the stable root, normalized hemisphere, pivot atlas, and
finite Lagrangian criterion.  V36 supplies the GHY normal gauge line, its
glancing plane, and the necessary selector-angle test.  V24, V27, V33, and
V34 supply respectively the common trace bundle, the distinction between
action and causal forms, the right/left obstruction handoff, and the
reduced wave row count.  No Euclidean or spherical-cavity result is used as
a substitute for the planar Lorentzian calculation above.

\paragraph{Five-version cross-program audit cadence.}
At V40 and every fifth SM--Einstein version thereafter, the newest retained
Renewal NS stability and Yang--Mills mass-gap notes in the shared
\texttt{latex\_notes} folder are to be read before the following version is
chosen.  The audit must record only dependencies that change this programme's
proof obligations or direction, must cite the precise imported result, and
must not alter either companion programme's files.

\paragraph{Honest continuum status.}
The conformal--mean-curvature branch is action-Lagrangian and pointwise
complementing in the open Laplace half-plane, but it has zero uniform stable
floor and rank seven at glancing.  Its linear BRST tangent is on-shell, not
a nonlinear off-shell BV--BFV domain.  The boundary embedding action and
physical quotient are not yet constructed.  Therefore the full
Standard Model--Einstein continuum is not proved.

% END APPENDED VERSION V37 MATERIAL



% BEGIN APPENDED VERSION V38 MATERIAL

\clearpage
\section{V38: normal-embedding exact sequence, constant-rank gauge orbit,
and the surviving physical glancing kernel}
\label{sec:V38-normal-embedding}

\subsection{Purpose and the precise edge-mode question}
\label{subsec:V38-purpose}

V37 produced the invariant scalar
\(\mathcal K_\zeta=K_{(1)}-D^aD_a\zeta\), but deliberately stopped
before calling it an edge action or a physical quotient.  This note performs
the finite linear tests that can be decided without such a claim.  There are
two distinct questions:

\begin{enumerate}
\item does the normal embedding scalar turn the collapsing normal gauge line
      into a constant-rank orbit through glancing; and
\item after quotienting that orbit, does the conformal--mean-curvature
      operator acquire a positive closed-hemisphere floor?
\end{enumerate}

The answers are respectively yes and no.  The first answer is an exact split
sequence with a uniformly transverse \(\zeta\)-coordinate.  The second is an
exact factorization: in gauge-invariant variables the quotient operator is
the V37 matrix \(\mathcal M_{CM}\) itself.  Thus its three-dimensional
glancing kernel is physical with respect to this one-dimensional normal
gauge quotient.

Embedding and edge fields are standard tools for constructing local phase
spaces with boundaries; see
\url{https://arxiv.org/abs/1601.04744} and
\url{https://arxiv.org/abs/1706.05061}.  A boundary action can also generate
edge dynamics, as explained for gauge-theory models in
\url{https://arxiv.org/abs/1912.06025}.  The review and reassessment
\url{https://arxiv.org/abs/2312.01918} is a useful warning that gauge
invariance of a presymplectic form alone is not the reduction theorem: its
degeneracy and pullback to the reduced space must also be checked.  These
papers motivate the construction below.  None of their conclusions is used
as a substitute for the stable-symbol or Green-form calculations proved
here.

This note extends rather than repeats the V31 reservoir and V35 dilation
rows.  No reservoir is introduced: first the smallest local Stueckelberg
extension is tested exactly.  It also does not revisit the V37 determinant
calculation; that result is imported only through the displayed matrix
\eqref{eq:V37-MCM} and its already proved rank.

\subsection{Stable normal orbit and gauge-invariant quotient coordinates}
\label{subsec:V38-exact-sequence}

Retain the V37 stable conventions
\eqref{eq:V37-stable-root}--\eqref{eq:V37-d-symbol}.  After the five
trace-free induced-metric selectors, write the extended trace amplitude as

\begin{equation}
 x=(\tau,v,\phi,\zeta)\in
 \mathbb C\oplus\mathbb C^3\oplus\mathbb C\oplus\mathbb C
 =:\mathsf X_{\rm ext}.
 \label{eq:V38-extended-trace}
\end{equation}

For a pure normal stable ghost of boundary amplitude \(q\), the linear BRST
symbol is

\begin{equation}
 s_qx
 =q\,g(s,k),
 \qquad
 g(s,k):=(0,d,2\lambda,-1).
 \label{eq:V38-gauge-column}
\end{equation}

Indeed, the V36 convention gives
\(s_qv=dq\), \(s_q\phi=2\lambda q\), and \(s_q\tau=0\), while
\eqref{eq:V37-zeta-BRST} gives \(s_q\zeta=-q\).  Define

\begin{align}
 V&:=v+d\zeta,
 &\Phi&:=\phi+2\lambda\zeta,
 \label{eq:V38-invariant-variables}\
 Q(s,k)x&:=(\tau,V,\Phi),
 &R(\tau,V,\Phi)&:=(\tau,V,\Phi,0).
 \label{eq:V38-QR}
\end{align}

Here and below \(R\) places \(V\) in the three \(v\)-slots and \(\Phi\)
in the \(\phi\)-slot.

\begin{theorem}[Split normal-embedding exact sequence]
\label{thm:V38-split-sequence}
At every point of the closed normalized stable hemisphere,

\begin{equation}
 0\longrightarrow\mathbb C
 \mathop{\longrightarrow}^{q\mapsto qg}
 \mathsf X_{\rm ext}
 \mathop{\longrightarrow}^{Q}
 \mathbb C^5
 \longrightarrow0
 \label{eq:V38-exact-sequence}
\end{equation}

is a split exact sequence, with right splitting \(R\).  In particular,

\begin{equation}
 Qg=0,
 \qquad QR=I_5,
 \qquad
 x=RQx+g(-\zeta).
 \label{eq:V38-splitting-identities}
\end{equation}

The normal gauge orbit therefore has constant rank one through glancing,
and \((\tau,V,\Phi)\) are global coordinates on its algebraic quotient.
\end{theorem}

\begin{proof}
The first and last components of \(Qg\) vanish immediately, and its middle
component is \(d+d(-1)=0\).  Thus \(\operatorname{ran}g\subset\ker Q\).
The identity \(QR=I_5\) proves that \(Q\) is onto.  If \(Qx=0\), then
\(\tau=0\), \(v=-d\zeta\), and \(\phi=-2\lambda\zeta\); hence
\(x=g(-\zeta)\).  This proves exactness.  Direct substitution proves the
last identity in \eqref{eq:V38-splitting-identities}.

Finally, the last component of \(g\) is \(-1\), so \(g\) never vanishes,
including at \(\lambda=0\).  Its span is consequently a rank-one subbundle,
and the displayed splitting identifies the quotient with the fixed
five-dimensional target of \(Q\).
\end{proof}

The splitting gives two useful projectors.  The oblique projector

\begin{equation}
 P_0:=RQ,
 \qquad
 P_0(\tau,v,\phi,\zeta)
 =(\tau,v+d\zeta,\phi+2\lambda\zeta,0)
 \label{eq:V38-oblique-projector}
\end{equation}

satisfies \(P_0^2=P_0\), \(\ker P_0=\operatorname{span}\{g\}\), and
\(\operatorname{ran}P_0=\{\zeta=0\}\).  For the coordinate Euclidean
Hermitian metric, put

\begin{equation}
 N:=g^\dagger g=1+\lVert d\rVert_{\rm E}^2+4|\lambda|^2,
 \qquad
 P_\perp:=I_6-\frac{gg^\dagger}{N}.
 \label{eq:V38-orthogonal-projector}
\end{equation}

Then \(P_\perp\) is the continuous orthogonal projector onto \(g^\perp\).
Both projectors kill the BRST orbit \eqref{eq:V38-gauge-column}.

\begin{corollary}[Uniform edge-coordinate transversality]
\label{cor:V38-edge-angle}
On the normalized hemisphere \eqref{eq:V35-normalized-hemisphere}, the
coordinate functional \(\ell_\zeta(x)=\zeta\) obeys

\begin{equation}
 \frac{|\ell_\zeta(gq)|}{\lVert gq\rVert_{\rm E}}
 =\frac1{\sqrt{N}}
 \ge\frac1{\sqrt6}
 \qquad(q\ne0).
 \label{eq:V38-edge-angle}
\end{equation}

Thus the embedding coordinate supplies a uniform selector for the normal
gauge line even though the V37 mean-curvature selector loses its angle at
glancing.
\end{corollary}

\begin{proof}
Normalization gives \(\lVert d\rVert_{\rm E}^2=|s|^2+|k|^2=1\).  The root
identity and the triangle inequality give
\(|\lambda|^2=|s^2+|k|^2|\le |s|^2+|k|^2=1\).  Hence \(N\le6\).  The
equality before the inequality follows from the last component of
\(gq\).
\end{proof}

\subsection{The quotient operator is exactly the V37 operator}
\label{subsec:V38-factorization}

The stable symbol of the invariant extrinsic scalar is

\begin{align}
 \mathcal K_\zeta
 &=\frac\lambda2\tau-d^\sharp v+\lambda^2\zeta
 =\frac\lambda2\tau-d^\sharp V,
 \label{eq:V38-Kzeta-symbol}
\end{align}

because \(d^\sharp d=-\lambda^2\).  Substitution of
\(v=V-d\zeta\) and \(\phi=\Phi-2\lambda\zeta\) into the harmonic rows
\eqref{eq:V37-Ca-symbol}--\eqref{eq:V37-Cn-symbol} gives

\begin{align}
 C_a&=\lambda V_a-\frac16d_a\tau-\frac12d_a\Phi,
 \label{eq:V38-Ca-invariant}\
 C_n&=d^\sharp V+\frac\lambda2(\Phi-\tau).
 \label{eq:V38-Cn-invariant}
\end{align}

Every \(\zeta\)-term cancels.  Let \(\mathcal B_{\rm ext}\) denote the
five-row reduced matrix with rows
\((\mathcal K_\zeta,C_a,C_n)\) on the six columns
\((\tau,v,\phi,\zeta)\).

\begin{theorem}[Exact quotient factorization and ranks]
\label{thm:V38-factorization-rank}
The extended reduced matrix factors as

\begin{equation}
 \boxed{\mathcal B_{\rm ext}=\mathcal M_{CM}Q.}
 \label{eq:V38-B-factorization}
\end{equation}

Consequently:

\begin{enumerate}
\item if \(\operatorname{Re}s>0\), then
\[
 \operatorname{rank}\mathcal B_{\rm ext}=5,
 \qquad
 \ker\mathcal B_{\rm ext}=\operatorname{span}\{g\};
 \]
\item at glancing \(\lambda=0,d\ne0\), then
\[
 \operatorname{rank}\mathcal B_{\rm ext}=2,
 \qquad
 \dim\ker\mathcal B_{\rm ext}=4;
 \]
\item after restoring the five trace-free induced-metric rows, the full
      ten-by-eleven matrix has rank ten in the open half-plane and rank
      seven at glancing;
\item the left kernel is zero in the open half-plane and at glancing is the
      same three-dimensional space
      \eqref{eq:V37-glancing-left-kernel} as in V37.
\end{enumerate}
\end{theorem}

\begin{proof}
Equations~\eqref{eq:V38-Kzeta-symbol}--\eqref{eq:V38-Cn-invariant}
are precisely the rows of \(\mathcal M_{CM}\) applied to
\(Qx=(\tau,V,\Phi)\), proving the factorization.  Since \(Q\) is onto,
\(\operatorname{ran}(\mathcal M_{CM}Q)=\operatorname{ran}\mathcal M_{CM}\),
so the two matrices have the same rank and the same left kernel.

In the open half-plane Theorem~\ref{thm:V37-open-determinant} makes
\(\mathcal M_{CM}\) invertible.  Hence
\(\ker\mathcal B_{\rm ext}=\ker Q=\operatorname{span}\{g\}\).
At glancing Theorem~\ref{thm:V37-glancing-rank} gives rank two.  Rank-nullity
on the six-dimensional extended trace fibre gives nullity four.  The five
trace-free selectors retain their independent identity block and contain no
\(\zeta\)-column, giving full ranks \(5+5=10\) and \(5+2=7\).
The left-kernel assertion follows from surjectivity of \(Q\) and the V37
left-kernel calculation.
\end{proof}

The splitting makes the glancing kernel more informative than a dimension
count.  From \eqref{eq:V37-glancing-right-kernel},

\begin{equation}
 \ker\mathcal B_{\rm ext}\big|_{\mathcal G}
 =R\left\{(\tau,V,\Phi):
       d^\sharp V=0,\ \tau+3\Phi=0\right\}
  \oplus\operatorname{span}\{g\}.
 \label{eq:V38-glancing-kernel-split}
\end{equation}

The first summand has dimension three and the second has dimension one.
The directness follows because \(R\)-images have \(\zeta=0\), whereas a
nonzero multiple of \(g\) has nonzero \(\zeta\)-component.

\begin{corollary}[A constant-rank gauge bundle does not give a physical
floor]
\label{cor:V38-no-physical-floor}
The operator induced by \(\mathcal B_{\rm ext}\) on
\(\mathsf X_{\rm ext}/\operatorname{span}\{g\}\) is, under the global
quotient chart \(Q\), exactly \(\mathcal M_{CM}\).  It is invertible at
every open stable frequency, but at glancing it has the three-dimensional
kernel

\begin{equation}
 \left\{(\tau,V,\Phi):d^\sharp V=0,\ \tau+3\Phi=0\right\}.
 \label{eq:V38-physical-glancing-kernel}
\end{equation}

Therefore no continuous choice of complement to the normal gauge line can
give this branch a positive closed-hemisphere least-singular-value floor.
\end{corollary}

\begin{proof}
The quotient identification follows from the split exact sequence and
\eqref{eq:V38-B-factorization}.  The open and glancing kernel statements
are the V37 determinant and kernel theorems in the quotient coordinates.
At a glancing point the induced operator is noninjective, so its least
singular value is zero for every pair of nonsingular domain and target
frames.  Continuity then rules out a positive uniform lower bound on the
closed hemisphere.  Changing the complement changes only the quotient
frame; it cannot remove the nonzero quotient classes in
\eqref{eq:V38-physical-glancing-kernel}.
\end{proof}

\begin{firewall}[What the scalar edge mode repairs]
\label{fire:V38-repair-scope}
The scalar \(\zeta\) repairs the bundle defect identified in
Corollary~\ref{cor:V36-no-kernel-bundle}: the \emph{normal gauge orbit} now
has constant rank and a uniform selector.  It does not repair the full
stable kernel, whose dimension still jumps from one to four before
quotienting and from zero to three after quotienting.  Calling all four
glancing vectors gauge would contradict the exact sequence
\eqref{eq:V38-exact-sequence}; three of them have nonzero quotient class.
\end{firewall}

\subsection{Lowest-order local scalar action: uniqueness and residual}
\label{subsec:V38-local-action}

It remains to ask whether \(\mathcal K_\zeta\) is the Euler row of a local
boundary action.  Work at the frozen flat anchor, take real fields with
compact tangential support, and restrict first to the lowest-order scalar
ansatz

\begin{equation}
 S_{\partial}^{(2)}[h,\zeta]
 :=\int_{\partial M}
 \left(
  \frac a2D_a\zeta D^a\zeta
  +b\zeta K_{(1)}(h)
  +\frac{m^2}{2}\zeta^2
 \right)\,\dd^3y .
 \label{eq:V38-action-ansatz}
\end{equation}

The constants are real.  This is a classification only within the displayed
ansatz: one tangential derivative per scalar factor, a zeroth-order mass,
and one linear mean-curvature coupling.

\begin{theorem}[Unique invariant Euler row in the minimal scalar ansatz]
\label{thm:V38-minimal-action}
The \(\zeta\)-Euler expression of \eqref{eq:V38-action-ansatz} is

\begin{equation}
 E_\zeta=-aD^aD_a\zeta+bK_{(1)}+m^2\zeta.
 \label{eq:V38-zeta-Euler}
\end{equation}

Under \(\delta_\xi\zeta=-\xi_n\) and
\(\delta_\xi K_{(1)}=-D^aD_a\xi_n\), it is invariant for every compactly
supported \(\xi_n\) if and only if

\begin{equation}
 a=b,
 \qquad m^2=0.
 \label{eq:V38-action-coefficients}
\end{equation}

For a nonzero action the resulting row is

\begin{equation}
 E_\zeta=a\mathcal K_\zeta.
 \label{eq:V38-E-is-Kzeta}
\end{equation}
\end{theorem}

\begin{proof}
Tangential integration by parts gives
\eqref{eq:V38-zeta-Euler}.  Its gauge variation is

\[
 \delta_\xi E_\zeta
 =(a-b)D^aD_a\xi_n-m^2\xi_n.
\]

For this differential operator to annihilate every compactly supported
test function, both its second-order and zeroth-order coefficients must
vanish.  This is \eqref{eq:V38-action-coefficients}.  Substitution then
gives \eqref{eq:V38-E-is-Kzeta}.
\end{proof}

The Euler row is gauge invariant, but the same statement is false for the
quadratic action itself.

\begin{proposition}[Exact residual gauge variation]
\label{prop:V38-action-residual}
For the nontrivial coefficient choice \(a=b\), \(m^2=0\), the frozen action
obeys

\begin{equation}
 \delta_\xi S_{\partial}^{(2)}
 =-a\int_{\partial M}\xi_n K_{(1)}(h)\,\dd^3y.
 \label{eq:V38-action-residual}
\end{equation}

Thus the minimal scalar action is not a standalone off-shell
diffeomorphism-invariant boundary action.
\end{proposition}

\begin{proof}
Varying the two terms and using the transformation laws gives

\[
 \delta_\xi S_{\partial}^{(2)}
 =a\int_{\partial M}
 \left(-D_a\zeta D^a\xi_n-\xi_nK_{(1)}
       -\zeta D^aD_a\xi_n\right)\dd^3y.
\]

Integrating the first term by parts gives
\(+a\int\zeta D^aD_a\xi_n\), which cancels the last term.  The remaining
term is exactly \eqref{eq:V38-action-residual}.
\end{proof}

For signature \((-++)\), the free kinetic part of
\eqref{eq:V38-action-ansatz} has canonical energy

\begin{equation}
 H_\zeta=-\frac a2\int
 \left((D_t\zeta)^2+|D_y\zeta|^2\right)\dd^2y.
 \label{eq:V38-scalar-energy}
\end{equation}

Hence its unreduced scalar energy is positive for \(a<0\).  This sign test
does not establish physical positivity: \(\zeta\) is a gauge coordinate in
\eqref{eq:V38-exact-sequence}, and the quotient retains the independent
kernel \eqref{eq:V38-physical-glancing-kernel}.

The first variation of the frozen action can be recorded without hiding
the metric coupling:

\begin{equation}
 \delta S_{\partial}^{(2)}
 =\int_{\partial M}
 \left(E_\zeta\,\delta\zeta
       +b\zeta K_{(1)}(\delta h)\right)\dd^3y
 +\int_{\partial M}D_a\left(aD^a\zeta\,\delta\zeta\right)\dd^3y.
 \label{eq:V38-action-first-variation}
\end{equation}

The last divergence is the scalar corner current.  The middle term shows
why the invariant Euler row alone is not the full variational boundary
condition: metric variation produces an additional normal-derivative
one-form.

\subsection{Green-form obstruction for the naive replacement}
\label{subsec:V38-Green-obstruction}

Define the naive Stueckelbergized CM relation by

\begin{equation}
 \mathcal L_{CM,\zeta}^{\rm naive}
 :=\left\{
 h^{\rm TF}_{ab}=0,
 \quad \mathcal K_\zeta=0,
 \quad C_\mu=0
 \right\}.
 \label{eq:V38-naive-relation}
\end{equation}

It is the relation whose stable matrix was analyzed above.  Its symbol and
its action incidence are different questions.

\begin{theorem}[The naive relation is not isotropic for the bulk CM form]
\label{thm:V38-naive-not-isotropic}
For compactly supported boundary jets, the restriction of the V37 metric
Green form to \(\mathcal L_{CM,\zeta}^{\rm naive}\) is

\begin{equation}
 2\kappa\,\Omega_{CM}(h,k)
 =\frac{2\varepsilon_D}{3}
 \int_{\partial M}
 \left(\tau_hD^aD_a\zeta_k
       -\tau_kD^aD_a\zeta_h\right)\dd^3y.
 \label{eq:V38-naive-Green-form}
\end{equation}

This bilinear form is not identically zero.  Therefore replacing
\(K_{(1)}=0\) by \(\mathcal K_\zeta=0\) in the V37 branch does not by itself
produce an action-Lagrangian relation.
\end{theorem}

\begin{proof}
On the relation, the trace-free and harmonic terms in
\eqref{eq:V37-CM-Green-form} vanish.  The remaining term is

\[
 \frac{2\varepsilon_D}{3}
 \int(\tau_hK_{(1)}(k)-\tau_kK_{(1)}(h)).
\]

The row \(\mathcal K_\zeta=0\) gives
\(K_{(1)}(h)=D^aD_a\zeta_h\) and the same identity for \(k\), proving
\eqref{eq:V38-naive-Green-form}.

To see nonvanishing, choose a real compactly supported smooth function
\(g\) with \(D^aD_ag\ne0\).  At the boundary-jet level take
\(\zeta_h=0\), \(\tau_h=D^aD_ag\), \(\zeta_k=g\), and \(\tau_k=0\).
The normal derivatives of the metric components can be chosen to satisfy
\(C_\mu=0\) and \(K_{(1)}=D^aD_a\zeta\), just as the V37 phase rows solve
for those derivatives.  The right-hand side is a nonzero multiple of
\(\int(D^aD_ag)^2\dd^3y\).  Hence the relation is not isotropic.
\end{proof}

The theorem does not rule out an extended action.  It proves that such an
action must contribute its own boundary/corner phase form and metric
one-form, and that these contributions must be included before a Lagrangian
statement is made.  Proposition~\ref{prop:V38-action-residual} supplies an
equally sharp Noether matching target: a gravitational embedding completion
must cancel \(+a\int\xi_nK_{(1)}\) against the residual
\eqref{eq:V38-action-residual}.  Until that cancellation is derived from a
single action, equality between an edge Hilbert stress and the boundary term
in the Bianchi identity is unproved.

\begin{firewall}[No stress or BRST overclaim]
\label{fire:V38-stress-firewall}
Equation~\eqref{eq:V38-E-is-Kzeta} is an invariant Euler expression, not an
off-shell BV--BFV construction.  Equation~\eqref{eq:V38-action-residual}
shows the missing off-shell term explicitly, and
Theorem~\ref{thm:V38-naive-not-isotropic} shows the missing Green-form term
explicitly.  A stress tensor computed from the intrinsic kinetic term alone
would omit the variation of \(\zeta K\) and the embedding map.  It therefore
cannot yet be identified with the gravitational Bianchi flux.
\end{firewall}

\subsection{Finite certificate and the next upstream row}
\label{subsec:V38-certificate}

A finite V38 certificate at a sample \((s,k)\) stores

\begin{align}
 r_{Qg}&:=\|Qg\|,
 &r_{QR}&:=\|QR-I_5\|,
 \label{eq:V38-residual-QR}\
 r_{P}&:=\|P_\perp^2-P_\perp\|,
 &r_{BP}&:=\|\mathcal B_{\rm ext}-\mathcal M_{CM}Q\|,
 \label{eq:V38-residual-factor}\
 r_{\rm dec}(x)&:=\|x-RQx-g(-\zeta)\|.
 \label{eq:V38-residual-decomposition}
\end{align}

The exact proofs make all five residuals zero.  Numerically, one should also
store rank five for a point with \(\operatorname{Re}s>0\), rank two for the
reduced glancing matrix, and the selector ratio in
\eqref{eq:V38-edge-angle}.  These checks diagnose implementation and frame
errors; they do not replace the exact factorization theorem.

The next upstream object is the full linearized embedding map rather than a
second scalar coefficient search.  Introduce tangential embedding
components \(X_a\) together with \(X_n=\zeta\), form the pulled-back induced
metric and extrinsic curvature from one covariant action, and derive all
metric, embedding, ghost, and corner rows by variation.  Its frozen complex
must decide whether the two transverse \(V\)-directions and the scalar
\(\tau+3\Phi\) direction in
\eqref{eq:V38-physical-glancing-kernel} are removed by genuine additional
gauge or are physical boundary channels.  If the full embedding quotient
again induces \(\mathcal M_{CM}\), the calculation has reached the point
where the V31 conservative boundary reservoir, with its own stress and
Green form, is required; the abstract V35 Julia dilation alone is still not
an action.

The V31 Higgs source conormal and the V32 neutrino branch are spectators in
this frozen gravitational calculation.  They must be rechecked only after a
single coupled embedding action passes the action, quotient, causal-floor,
and Bianchi tests.  Adding them before the gravitational glancing kernel is
resolved would hide rather than close the present obstruction.

\subsection{V38 ledger}
\label{subsec:V38-ledger}

\paragraph{New conclusions proved in V38.}
Theorem~\ref{thm:V38-split-sequence} constructs a global split quotient for
the normal embedding scalar, and Corollary~\ref{cor:V38-edge-angle} proves a
uniform selector angle.  Theorem~\ref{thm:V38-factorization-rank} proves
that the extended operator factors through the V37 CM matrix, with exact
open and glancing ranks.  Corollary~\ref{cor:V38-no-physical-floor} proves
that the three-dimensional quotient glancing kernel survives.  Within the
lowest-order scalar ansatz, Theorem~\ref{thm:V38-minimal-action} classifies
the unique coefficients that produce \(\mathcal K_\zeta\), while
Proposition~\ref{prop:V38-action-residual} computes the uncancelled gauge
variation.  Theorem~\ref{thm:V38-naive-not-isotropic} proves that the naive
row replacement is not Lagrangian for the original bulk CM Green form.

\paragraph{Imported conclusions not repeated.}
V36 supplies the pure normal gauge mode and the fixed-induced-metric bundle
defect.  V37 supplies the CM stable matrix, determinant, Green form, and
glancing kernels.  V31 supplies the conservative bulk--boundary Green
cancellation required of any later reservoir, and V35 supplies only the
abstract unitary dilation and finite certificate framework.  The edge-mode
literature cited in the purpose subsection supplies context, not any rank,
positivity, or action-incidence conclusion used in the proofs.

\paragraph{Five-version companion audit.}
V38 is not an audit version.  In accordance with the V37 cadence, the newest
retained NS stability and Yang--Mills mass-gap notes in the shared folder
will be inspected at V40, and only precise results that change this
programme's dependencies or next bottleneck will be imported.  Neither
companion file is modified here.

\paragraph{Honest continuum status.}
The normal embedding scalar now has a rigorous constant-rank gauge orbit and
quotient projector.  Its quotient nevertheless retains an exact
three-dimensional glancing kernel, and the minimal scalar action has an
uncancelled normal-diffeomorphism variation.  A full embedding action,
extended Green form, stress/Bianchi identity, and positive physical causal
floor remain unconstructed.  The full Standard Model--Einstein continuum is
therefore not proved.

% END APPENDED VERSION V38 MATERIAL



% BEGIN APPENDED VERSION V39 MATERIAL

\clearpage
\section{V39: full embedding exact sequence, pure-pullback no-go, and the
three-channel glancing rank budget}
\label{sec:V39-full-embedding}

\subsection{V38 typesetting erratum and purpose}
\label{subsec:V39-erratum-purpose}

Four intended display row terminators in V38 were emitted as single
control-space tokens.  This joined adjacent formulas typographically but
did not change any definition, label, proof, or rank calculation.  For an
unambiguous reading, the affected formulas are reproduced here without new
labels:

\begin{equation*}
 V=v+d\zeta,
 \qquad \Phi=\phi+2\lambda\zeta,
 \qquad Qx=(\tau,V,\Phi),
 \qquad R(\tau,V,\Phi)=(\tau,V,\Phi,0).
\end{equation*}

\begin{equation*}
 C_a=\lambda V_a-\frac16d_a\tau-\frac12d_a\Phi,
 \qquad
 C_n=d^\sharp V+\frac\lambda2(\Phi-\tau).
\end{equation*}

\begin{equation*}
 r_{Qg}=\lVert Qg\rVert,
 \qquad r_{QR}=\lVert QR-I_5\rVert.
\end{equation*}

\begin{equation*}
 r_P=\lVert P_\perp^2-P_\perp\rVert,
 \qquad
 r_{BP}=\lVert\mathcal B_{\rm ext}-\mathcal M_{CM}Q\rVert,
 \qquad
 r_{\rm dec}(x)=\lVert x-RQx-g(-\zeta)\rVert.
\end{equation*}

The corresponding V38 labels remain
\eqref{eq:V38-invariant-variables}, \eqref{eq:V38-QR},
\eqref{eq:V38-Ca-invariant}, \eqref{eq:V38-Cn-invariant}, and
\eqref{eq:V38-residual-QR}--\eqref{eq:V38-residual-decomposition}.

V38 showed that a normal embedding scalar makes the normal gauge orbit a
line bundle but leaves three nonzero quotient classes at glancing.  The
next question is whether the omitted tangential embedding components change
that conclusion.  This note constructs the complete four-component frozen
embedding complex.  The result is definitive for a pure pullback extension:
all four diffeomorphism directions form a uniformly complemented gauge
bundle, but the quotient operator is exactly the original V37
conformal--mean-curvature operator.  The three glancing classes therefore
remain physical with respect to the full diffeomorphism quotient.

The final subsection converts that negative result into a positive design
constraint.  Any regular, eliminable boundary reservoir which repairs the
rank-two glancing matrix must carry at least three independent stable
boundary amplitudes.  The lower bound is sharp at the algebraic level and
matches the two transverse-vector classes plus one scalar class found in
V38.

\subsection{The four-component stable embedding complex}
\label{subsec:V39-embedding-complex}

Let

\begin{equation}
 p:=(d,\lambda)\in\mathbb C^4
 \label{eq:V39-p-covector}
\end{equation}

be the stable derivative covector in tangential-normal order.  For a
boundary embedding amplitude

\begin{equation}
 X=(X_a,X_n)=(X_a,\zeta)\in\mathbb C^4,
 \label{eq:V39-X-field}
\end{equation}

define the frozen symmetric-gradient map
\(\mathcal S_p:\mathbb C^4\to\operatorname{Sym}^2\mathbb C^4\) by

\begin{equation}
 (\mathcal S_pX)_{ab}:=d_aX_b+d_bX_a,
 \label{eq:V39-Sp-ab}
\end{equation}

\begin{equation}
 (\mathcal S_pX)_{na}:=\lambda X_a+d_aX_n,
 \qquad
 (\mathcal S_pX)_{nn}:=2\lambda X_n.
 \label{eq:V39-Sp-normal}
\end{equation}

The extended metric-embedding fibre and its invariant representative are

\begin{equation}
 \mathsf E_p
 :=\operatorname{Sym}^2\mathbb C^4\oplus\mathbb C^4,
 \qquad
 \widehat h:=Q_4(h,X):=h+\mathcal S_pX.
 \label{eq:V39-Q4-definition}
\end{equation}

For a four-component stable diffeomorphism amplitude \(q\), define

\begin{equation}
 G_pq:=(\mathcal S_pq,-q)\in\mathsf E_p.
 \label{eq:V39-Gp-definition}
\end{equation}

This is the symbol of \(s_qh=\mathcal S_pq\), \(s_qX=-q\).  The sign
agrees with \eqref{eq:V37-zeta-BRST} in the normal component.

\begin{theorem}[Full embedding split exact sequence]
\label{thm:V39-full-exact-sequence}
For every frequency in the closed normalized stable hemisphere,

\begin{equation}
 0\longrightarrow\mathbb C^4
 \mathop{\longrightarrow}^{G_p}
 \mathsf E_p
 \mathop{\longrightarrow}^{Q_4}
 \operatorname{Sym}^2\mathbb C^4
 \longrightarrow0
 \label{eq:V39-full-exact-sequence}
\end{equation}

is split exact.  A splitting is

\begin{equation}
 R_4(H):=(H,0),
 \qquad Q_4R_4=I_{10},
 \label{eq:V39-R4}
\end{equation}

and every \((h,X)\in\mathsf E_p\) has the unique decomposition

\begin{equation}
 (h,X)=R_4Q_4(h,X)+G_p(-X).
 \label{eq:V39-full-decomposition}
\end{equation}

In particular, the full stable diffeomorphism orbit has constant rank four
through glancing.
\end{theorem}

\begin{proof}
The definitions give

\[
 Q_4G_pq=\mathcal S_pq+\mathcal S_p(-q)=0,
 \]

so \(\operatorname{ran}G_p\subset\ker Q_4\).  Surjectivity follows from
\(Q_4R_4=I_{10}\).  Conversely, if \(Q_4(h,X)=0\), then
\(h=-\mathcal S_pX\), whence
\((h,X)=G_p(-X)\).  Thus the sequence is exact.  The same identity gives
\eqref{eq:V39-full-decomposition}.  Finally, the second component of
\(G_pq\) is \(-q\), so \(G_pq=0\) implies \(q=0\) at every frequency.
Therefore \(G_p\) has rank four everywhere.
\end{proof}

Equip symmetric tensors with the Euclidean Frobenius norm in the stable
coordinate frame and \(\mathsf E_p\) with the orthogonal product norm.  The
edge-coordinate selector is

\begin{equation}
 \ell_X(h,X):=X.
 \label{eq:V39-edge-selector}
\end{equation}

\begin{proposition}[Uniform four-orbit transversality]
\label{prop:V39-four-orbit-angle}
On \(\overline{\mathbb S}_+\),

\begin{equation}
 \frac{\lVert\ell_X(G_pq)\rVert}{\lVert G_pq\rVert}
 \ge\frac13
 \qquad(q\ne0).
 \label{eq:V39-four-orbit-angle}
\end{equation}

Moreover,

\begin{equation}
 G_p^\dagger G_p=I_4+\mathcal S_p^\dagger\mathcal S_p\ge I_4,
 \label{eq:V39-Gram-lower}
\end{equation}

so the orthogonal gauge projector

\begin{equation}
 P_{G,p}:=G_p(G_p^\dagger G_p)^{-1}G_p^\dagger
 \label{eq:V39-gauge-projector}
\end{equation}

and physical representative projector \(I-P_{G,p}\) are continuous and
uniformly bounded on the closed hemisphere.
\end{proposition}

\begin{proof}
For rank-one tensors,
\(\lVert p\otimes q\rVert_{\rm F}=\lVert p\rVert\lVert q\rVert\).
Hence

\[
 \lVert\mathcal S_pq\rVert_{\rm F}
 \le2\lVert p\rVert\lVert q\rVert.
\]

On the normalized hemisphere,
\(\lVert d\rVert^2=1\) and \(|\lambda|^2\le1\), so
\(\lVert p\rVert^2\le2\).  It follows that

\[
 \lVert G_pq\rVert^2
 =\lVert\mathcal S_pq\rVert_{
 \rm F}^2+\lVert q\rVert^2
 \le9\lVert q\rVert^2.
\]

Since \(\lVert\ell_X(G_pq)\rVert=\lVert q\rVert\), this proves
\eqref{eq:V39-four-orbit-angle}.  The Gram identity follows directly from
the product norm.  Its smallest eigenvalue is at least one, so its inverse
exists with norm at most one.  Continuity of \(p\), \(G_p\), and matrix
inversion on this uniformly positive Gram family proves the projector
statement.
\end{proof}

\subsection{Pulled-back geometric rows and exact quotient preservation}
\label{subsec:V39-pullback-factorization}

Let \(\mathcal A_{CM}(p)\) denote the full ten-row V37 stable boundary
operator on a metric amplitude \(H\): five rows
\(H_{ab}^{\rm TF}=0\), the row \(K_{(1)}(H)=0\), and the four harmonic
rows \(C_\mu(H)=0\).  Define its pure embedding pullback by

\begin{equation}
 \mathcal A_{\rm emb}(p):=\mathcal A_{CM}(p)Q_4(p).
 \label{eq:V39-Aemb-definition}
\end{equation}

The next lemma verifies that this algebraic definition is exactly the
linearized geometric pullback at the planar anchor.

\begin{lemma}[Frozen pullback identities]
\label{lem:V39-pullback-identities}
For the stable extension \(X e^{st+ik\cdot y-\lambda x}\),

\begin{equation}
 K_{(1)}(\widehat h)
 =K_{(1)}(h)-D^aD_aX_n,
 \label{eq:V39-K-pullback}
\end{equation}

\begin{equation}
 C_\mu(\widehat h)
 =C_\mu(h)+(\lambda^2+d^\sharp d)X_\mu
 =C_\mu(h).
 \label{eq:V39-C-pullback}
\end{equation}

The induced row is

\begin{equation}
 \widehat h_{ab}^{\rm TF}
 =\left(h_{ab}+d_aX_b+d_bX_a\right)^{\rm TF}.
 \label{eq:V39-induced-pullback}
\end{equation}
\end{lemma}

\begin{proof}
The induced formula is the \(ab\)-component of
\eqref{eq:V39-Q4-definition}.  For the extrinsic row, use the V37 symbol
\(K_{(1)}(H)=\lambda H^a{}_a/2-d^\sharp H_{n\cdot}\).  The symmetric
gradient contributes

\[
 \lambda d^\sharp X
 -d^\sharp(\lambda X+dX_n)
 =-d^\sharp dX_n
 =-D^aD_aX_n.
\]

For a pure symmetric gradient, the harmonic covector is the wave operator
on its generating vector.  Its stable symbol is
\(\lambda^2+d^\sharp d\), which vanishes by
\eqref{eq:V37-d-symbol}.  This proves all three statements.
\end{proof}

\begin{theorem}[Pure-pullback rank theorem]
\label{thm:V39-pure-pullback-rank}
The full embedding boundary operator has

\begin{equation}
 \operatorname{rank}\mathcal A_{\rm emb}=10,
 \qquad
 \dim\ker\mathcal A_{\rm emb}=4
 \label{eq:V39-open-rank-nullity}
\end{equation}

at every open stable frequency.  At glancing it has

\begin{equation}
 \operatorname{rank}\mathcal A_{\rm emb}=7,
 \qquad
 \dim\ker\mathcal A_{\rm emb}=7.
 \label{eq:V39-glancing-rank-nullity}
\end{equation}

More precisely, at every frequency,

\begin{equation}
 \ker\mathcal A_{\rm emb}
 =R_4(\ker\mathcal A_{CM})
 \mathbin\oplus\operatorname{ran}G_p.
 \label{eq:V39-kernel-splitting}
\end{equation}

The left kernel of \(\mathcal A_{\rm emb}\) equals the left kernel of
\(\mathcal A_{CM}\): it is zero in the open half-plane and has dimension
three at glancing.
\end{theorem}

\begin{proof}
Surjectivity of \(Q_4\) gives

\[
 \operatorname{ran}(\mathcal A_{CM}Q_4)
 =\operatorname{ran}\mathcal A_{CM}.
\]

The V37 full ranks are ten in the open half-plane and seven at glancing,
so \eqref{eq:V39-open-rank-nullity} and
\eqref{eq:V39-glancing-rank-nullity} follow by rank-nullity on the
fourteen-dimensional extended fibre.

For the kernel, decompose an arbitrary extended vector by
\eqref{eq:V39-full-decomposition}.  Applying \(\mathcal A_{\rm emb}\)
kills the gauge summand and gives
\(\mathcal A_{CM}Q_4(h,X)\) on the split representative.  This proves
\eqref{eq:V39-kernel-splitting}; the sum is direct because \(R_4\)-images
have zero embedding component while every nonzero gauge image has a
nonzero embedding component.  Finally, a left row annihilates
\(\mathcal A_{CM}Q_4\) if and only if it annihilates
\(\mathcal A_{CM}\), again because \(Q_4\) is onto.
\end{proof}

\begin{corollary}[The full diffeomorphism quotient retains three physical
glancing classes]
\label{cor:V39-full-quotient-no-go}
The quotient

\begin{equation}
 \mathsf E_p/\operatorname{ran}G_p
 \mathop{\cong}^{Q_4}
 \operatorname{Sym}^2\mathbb C^4
 \label{eq:V39-full-quotient-chart}
\end{equation}

intertwines the induced embedding operator with
\(\mathcal A_{CM}\).  After the five trace-free rows are eliminated, its
glancing kernel is

\begin{equation}
 \mathsf K_{\rm phys}
 =\left\{(\tau,V,\Phi):
 d^\sharp V=0,\ \tau+3\Phi=0\right\},
 \qquad \dim\mathsf K_{\rm phys}=3.
 \label{eq:V39-Kphys}
\end{equation}

Thus adding all four embedding components, with no new gauge-invariant
dynamics, cannot create a positive physical closed-hemisphere floor.
\end{corollary}

\begin{proof}
The quotient isomorphism is Theorem~\ref{thm:V39-full-exact-sequence}.
The intertwining relation is the definition
\eqref{eq:V39-Aemb-definition}.  The reduced glancing kernel is
Theorem~\ref{thm:V37-glancing-rank}, with the invariant amplitudes denoted
by \((\tau,V,\Phi)\).  Since these are nonzero quotient classes, no choice
of representative or gauge projector can remove them.  The induced least
singular value is therefore zero at glancing.
\end{proof}

The physical kernel has a canonical sector decomposition at each glancing
frequency:

\begin{equation}
 \mathsf K_{\rm vec}
 :=\{(0,V,0):d^\sharp V=0\},
 \qquad \dim\mathsf K_{\rm vec}=2,
 \label{eq:V39-vector-kernel}
\end{equation}

\begin{equation}
 \mathsf K_{\rm sc}
 :=\operatorname{span}\{(-3,0,1)\},
 \qquad
 \mathsf K_{\rm phys}=\mathsf K_{\rm vec}\oplus\mathsf K_{\rm sc}.
 \label{eq:V39-scalar-kernel}
\end{equation}

The corresponding left obstruction splits into

\begin{equation}
 \mathsf L_{\rm vec}
 :=\{(0,\beta^T,0):\beta^Td=0\},
 \qquad \dim\mathsf L_{\rm vec}=2,
 \label{eq:V39-vector-cokernel}
\end{equation}

\begin{equation}
 \mathsf L_{\rm sc}
 :=\operatorname{span}\{(1,0,1)\},
 \qquad
 \ker\mathcal M_{CM}^T
 =\mathsf L_{\rm vec}\oplus\mathsf L_{\rm sc}.
 \label{eq:V39-scalar-cokernel}
\end{equation}

These decompositions identify which quotient sectors a genuine boundary
dynamics must connect.

\subsection{What a pure-pullback quadratic action does and does not do}
\label{subsec:V39-pullback-action}

The quotient preservation is not limited to boundary rows.  Let
\(\mathcal H\) be any finite frozen metric trace space, let
\(A:\mathcal H\to\mathcal H^*\) be symmetric for the chosen bilinear
pairing, and define

\begin{equation}
 F_A^{\rm emb}(h,X)
 :=\frac12\left\langle Q_4(h,X),A Q_4(h,X)\right\rangle.
 \label{eq:V39-pullback-quadratic-action}
\end{equation}

\begin{proposition}[Linear Noether identity for a pure pullback]
\label{prop:V39-pullback-Noether}
The quadratic functional \eqref{eq:V39-pullback-quadratic-action} is exactly
invariant under \((\delta h,\delta X)=G_pq\).  If its Euler covectors are
denoted by \((E_h,E_X)\), then

\begin{equation}
 E_X=\mathcal S_p^\dagger E_h,
 \qquad
 \mathcal S_p^\dagger E_h-E_X=0.
 \label{eq:V39-linear-Noether-identity}
\end{equation}

Its Hessian has \(\operatorname{ran}G_p\) in its radical, and the induced
Hessian on the quotient is exactly \(A\).
\end{proposition}

\begin{proof}
Theorem~\ref{thm:V39-full-exact-sequence} gives \(Q_4G_pq=0\), so the
gauge variation of \(F_A^{\rm emb}\) vanishes.  Differentiating the
quadratic functional gives

\[
 E_h=A\widehat h,
 \qquad
 E_X=\mathcal S_p^\dagger A\widehat h
      =\mathcal S_p^\dagger E_h,
\]

which is the displayed Noether identity.  The Hessian is
\(Q_4^\dagger A Q_4\), so it annihilates every vector in \(\ker Q_4\).
On the split quotient representative \(R_4H\), it evaluates to
\(\langle H,AH\rangle\), proving the last assertion.
\end{proof}

\begin{firewall}[Gauge completion is not new physical dynamics]
\label{fire:V39-pullback-firewall}
Proposition~\ref{prop:V39-pullback-Noether} gives a valid frozen linear
Noether identity, but the quotient Hessian is unchanged.  It therefore
cannot cancel the V38 residual by merely renaming the old metric data, and
it cannot lift \(\mathsf K_{\rm phys}\).  A nonlinear Einstein embedding
action must also account for the moving integration domain, corner terms,
ghosts, and metric stress.  No nonlinear BV--BFV or stress/Bianchi equality
is claimed here.
\end{firewall}

\subsection{Exact lower bound on regular boundary reservoir channels}
\label{subsec:V39-rank-budget}

The quotient no-go says that the next successful row must change the
physical quotient operator.  The following rank argument determines the
minimum regular finite boundary content needed to do so.

Fix a glancing frequency and abbreviate

\begin{equation}
 M_0:=\mathcal M_{CM}(0,d):\mathbb C^5\to\mathbb C^5,
 \qquad \operatorname{rank}M_0=2.
 \label{eq:V39-M0}
\end{equation}

Let \(y\in\mathbb C^r\) be physical boundary amplitudes with coupled square
symbol

\begin{equation}
 \mathfrak M(x,y)
 :=\bigl(M_0x+Uy,\;Vx+Dy\bigr),
 \label{eq:V39-coupled-symbol}
\end{equation}

where \(U:\mathbb C^r\to\mathbb C^5\),
\(V:\mathbb C^5\to\mathbb C^r\), and
\(D:\mathbb C^r\to\mathbb C^r\).  Assume that \(D\) is invertible at
this frequency, so the boundary variables are regularly eliminable.

\begin{theorem}[Three-channel lower bound]
\label{thm:V39-three-channel-lower-bound}
If \(\mathfrak M\) is invertible at glancing, then

\begin{equation}
 r\ge3.
 \label{eq:V39-r-at-least-three}
\end{equation}

Equivalently, no regularly eliminable scalar or two-amplitude boundary
reservoir can repair the V37 glancing rank defect.
\end{theorem}

\begin{proof}
Solving the second component of \(\mathfrak M(x,y)=0\) gives

\[
 y=-D^{-1}Vx.
\]

The Schur complement on the metric quotient is therefore

\begin{equation}
 S=M_0-UD^{-1}V.
 \label{eq:V39-Schur-complement}
\end{equation}

Invertibility of \(D\) implies that \(\mathfrak M\) is invertible if and
only if \(S\) is invertible.  But

\[
 \operatorname{rank}(UD^{-1}V)\le r
\]

and the rank inequality gives

\[
 \operatorname{rank}S
 \le\operatorname{rank}M_0+operatorname{rank}(UD^{-1}V)
 \le2+r.
\]

An invertible endomorphism of \(\mathbb C^5\) has rank five.  Hence
\(5\le2+r\), which is \eqref{eq:V39-r-at-least-three}.
\end{proof}

\begin{proposition}[Algebraic sharpness of the channel count]
\label{prop:V39-rank-budget-sharp}
At a fixed glancing frequency there exists a rank-three correction
\(\Delta:\mathbb C^5\to\mathbb C^5\) such that \(M_0+\Delta\) is
invertible.  Thus the lower bound in
Theorem~\ref{thm:V39-three-channel-lower-bound} is sharp as a pointwise
linear-algebra statement.
\end{proposition}

\begin{proof}
Choose a complement \(\mathsf X_1\) of \(\ker M_0\) in \(\mathbb C^5\)
and a complement \(\mathsf C\) of \(\operatorname{ran}M_0\) in the target.
Both \(\ker M_0\) and \(\mathsf C\) have dimension three, while

\[
 M_0|_{\mathsf X_1}:\mathsf X_1\longrightarrow
 \operatorname{ran}M_0
\]

is an isomorphism.  Choose any isomorphism
\(T:\ker M_0\to\mathsf C\), and define \(\Delta=T\) on
\(\ker M_0\) and \(\Delta=0\) on \(\mathsf X_1\).  Then
\(\operatorname{rank}\Delta=3\), and relative to the two direct-sum
decompositions the map \(M_0+\Delta\) is the direct sum of the two
isomorphisms \(M_0|_{\mathsf X_1}\) and \(T\).  It is therefore
invertible.
\end{proof}

One may realize any factorization \(\Delta=-UD^{-1}V\) with three
auxiliary amplitudes at a single frequency by choosing bases in the two
three-dimensional defect spaces.  This observation proves only algebraic
feasibility.  A continuum reservoir must provide \(U,V,D\) from one local
or explicitly controlled nonlocal action, remain regular on every stable
chart, carry a positive causal energy on the physical quotient, and produce
the same stress term in metric variation and in the boundary Bianchi
identity.

\begin{firewall}[Singular reservoirs are outside the rank proof]
\label{fire:V39-singular-reservoir}
Theorem~\ref{thm:V39-three-channel-lower-bound} assumes that \(D\) is
invertible at glancing.  If a proposed boundary field has its own glancing
root, Schur elimination is illegal there and the complete
\((5+r)\)-dimensional stable system must be analyzed.  Such a singular
channel is not a counterexample to the theorem and is not automatically a
repair; it may add further kernel or incoming modes.
\end{firewall}

\subsection{Finite certificate and V40 handoff}
\label{subsec:V39-certificate-handoff}

A finite full-embedding certificate stores

\begin{equation}
 r_{QG}:=\lVert Q_4G_p\rVert,
 \qquad
 r_{QR}^{(4)}:=\lVert Q_4R_4-I_{10}\rVert,
 \qquad
 r_{\rm dec}^{(4)}:=\lVert z-R_4Q_4z-G_p(-X)\rVert.
 \label{eq:V39-exact-sequence-residuals}
\end{equation}

\begin{equation}
 r_{\rm fact}^{(4)}
 :=\lVert\mathcal A_{\rm emb}-\mathcal A_{CM}Q_4\rVert,
 \qquad
 r_{P_G}:=\lVert P_{G,p}^2-P_{G,p}\rVert.
 \label{eq:V39-factor-projector-residuals}
\end{equation}

The exact arguments make these residuals zero.  Open and glancing samples
must return ranks \((10,10)\) for \(\mathcal A_{CM}\) and
\(\mathcal A_{\rm emb}\) in the open half-plane, and ranks \((7,7)\) at
glancing.  For any proposed regular reservoir, the certificate must also
record the rank of \(D\), the rank of \(UD^{-1}V\), and the five singular
values of the Schur complement \eqref{eq:V39-Schur-complement}.

V40 is the scheduled five-version companion audit.  Before choosing a
three-channel action, it must read the newest retained Renewal NS stability
and Yang--Mills mass-gap notes and identify any proved boundary reservoir,
Green-form, stress, or glancing-sector result that changes the present
requirements.  The audit should then construct or reject a concrete
three-amplitude quotient dynamics rather than append another pure embedding
field.

\subsection{V39 ledger}
\label{subsec:V39-ledger}

\paragraph{New conclusions proved in V39.}
Theorem~\ref{thm:V39-full-exact-sequence} constructs the full
four-component embedding exact sequence, and
Proposition~\ref{prop:V39-four-orbit-angle} gives a uniform gauge selector
and projector.  Lemma~\ref{lem:V39-pullback-identities} derives the frozen
geometric pullback.  Theorem~\ref{thm:V39-pure-pullback-rank} and
Corollary~\ref{cor:V39-full-quotient-no-go} prove that a pure pullback leaves
the V37 quotient operator and its three physical glancing classes unchanged.
Proposition~\ref{prop:V39-pullback-Noether} supplies the exact frozen linear
Noether identity for any quadratic pullback action.  Finally,
Theorem~\ref{thm:V39-three-channel-lower-bound} proves the minimum regular
reservoir channel count, and Proposition~\ref{prop:V39-rank-budget-sharp}
proves pointwise algebraic sharpness.

\paragraph{Imported conclusions not repeated.}
V37 supplies the CM open and glancing ranks, right kernel, and left kernel.
V38 supplies the scalar exact sequence and establishes that the three
glancing quotient classes are not part of the normal gauge line.  V31 and
V35 continue to impose, respectively, the conservative Green cancellation
and the distinction between an abstract dilation and a boundary action.

\paragraph{Honest continuum status.}
All four frozen diffeomorphism directions now form a constant-rank,
uniformly complemented gauge bundle.  This completes the linear algebra of
the pure embedding quotient, not the continuum programme: the quotient
still has two transverse-vector and one scalar glancing class.  At least
three regular physical boundary amplitudes are necessary to lift them, and
no covariant action, causal symmetrizer, stress/Bianchi identity, or coupled
SM source compatibility has yet been proved.  The full Standard
Model--Einstein continuum is therefore not proved.

% END APPENDED VERSION V39 MATERIAL



% BEGIN APPENDED VERSION V40 MATERIAL

\clearpage
\section{V40: five-version companion audit and a rank-three transparent
glancing port}
\label{sec:V40-audit-transparent-port}

\subsection{Scheduled cross-program audit}
\label{subsec:V40-cross-audit}

This is the five-version audit required by the V37/V39 ledger.  Before the
V40 direction was chosen, the following retained companion files in the
shared \texttt{latex\_notes} folder were read without modification:

\begin{enumerate}
\item \texttt{Renewal\_NS\_Stability\_Research\_Notes\_V12.tex},
      SHA--256
      \[
      \texttt{f12bfd29b52945beaffb1d295a37a149570c03d275664d6fb0dcb2bdfc4d3079};
      \]
\item \texttt{Renewal\_Yang\_Mills\_Mass\_Gap\_Research\_Notes\_V21.tex},
      SHA--256
      \[
      \texttt{7c80642af5d8e683bf073090bf11d10403cbe2f17c0f878bf0603723ec314876}.
      \]
\end{enumerate}

Only conclusions that alter the present proof obligations are recorded.
No companion theorem is transplanted across a type mismatch.

\paragraph{Navier--Stokes import.}
The NS V12 theorem labelled
\texttt{thm:v12-heat-edge-trichotomy} applies to the fully assembled
exterior source after its positive action has been acquired.  Its preceding
firewall proves that the V34 raw history residual is not that complete
source.  The corollary labelled \texttt{cor:v12-endpoint-summable} and the
proposition labelled \texttt{prop:v12-native-response-no-go} then prove
that the native heat-Duhamel endpoint response can be summable at doubled
records and cannot by itself close the inherited critical core.

For the Einstein boundary problem this changes the reservoir incidence
gate.  A boundary amplitude may be called a physical source channel only
after it is identified with the complete Green/Bianchi flux of the coupled
action.  A response estimate for a projected glancing residual is not that
identity.  Moreover a passive Duhamel response, even when sharply bounded,
does not supply the required nonsummable or coercive continuum debit.

\paragraph{Yang--Mills import.}
The YM V21 theorem labelled \texttt{thm:V21-actual-law-split} replaces the
local-Gibbs Witten split on the actual Perron law by a positive comparison
minus the enlarged defect
\(\Theta_\eta+S_{\vartheta,-}\).  The lemma
\texttt{lem:V21-negative-Hessian} proves that
\(S_{\vartheta,-}\ne0\) whenever the density defect is nonconstant, and
the assessment \texttt{ass:V21-corrected-carrier} requires that channel in
the complete source Gram.  The firewall
\texttt{fire:V21-repair-locality} states that the new channel need not
inherit finite propagation because it depends on the global Perron
eigenfunction.

This affects the later quantum continuum audit, but it does not provide one
of the three Lorentzian stable amplitudes required by V39.  A Hessian defect
on a reduced configuration law and a stable boundary channel are different
objects.  It would be circular to count the former in the V39 rank budget
or to assign it locality.  Thus V40 keeps the glancing repair classical,
retains all coefficients explicitly, and tests nonlocality rather than
hiding it.

\begin{protocol}[Result of the V40 audit]
\label{prot:V40-audit-result}
The next Einstein boundary construction must meet all of the following
requirements:
\begin{enumerate}
\item at least three independent physical stable amplitudes at glancing;
\item exact incidence with the complete bulk Green/Bianchi flux before any
      response or capacity estimate;
\item a local action or a declared transparent nonlocal operator, with no
      inference of locality from a reduced-state comparison;
\item at the quantum cofinal stage, enlargement by every actual-law density
      defect such as \(S_{\vartheta,-}^{1/2}\), without identifying that
      enlargement with classical boundary rank.
\end{enumerate}
The next companion audit is V45.
\end{protocol}

\subsection{A canonical rank-two vector projection at glancing}
\label{subsec:V40-glancing-projection}

Return to the reduced V37 quotient coordinates

\begin{equation}
 x=(\tau,V,\Phi)\in\mathbb C\oplus\mathbb C^3\oplus\mathbb C.
 \label{eq:V40-reduced-x}
\end{equation}

Let

\begin{equation}
 \eta:=\operatorname{diag}(-1,1,1),
 \qquad d^\sharp=d^T\eta.
 \label{eq:V40-eta}
\end{equation}

On the normalized glancing set,
\(s=\epsilon i|k|\), \(\epsilon\in\{-1,1\}\),
\(|s|^2+|k|^2=1\), and therefore

\begin{equation}
 |k|^2=\frac12,
 \qquad
 q_d:=d^Td=s^2-|k|^2=-1.
 \label{eq:V40-qd-glancing}
\end{equation}

Define

\begin{equation}
 P_d:=I_3-\frac{(\eta d)d^\sharp}{q_d}.
 \label{eq:V40-Pd}
\end{equation}

\begin{lemma}[Exact vector-defect projector]
\label{lem:V40-vector-projector}
At every normalized glancing frequency,

\begin{equation}
 P_d^2=P_d,
 \qquad
 \operatorname{ran}P_d=\ker d^\sharp,
 \qquad
 \ker P_d=\operatorname{span}\{\eta d\},
 \qquad
 \operatorname{rank}P_d=2.
 \label{eq:V40-Pd-properties}
\end{equation}

Moreover,

\begin{equation}
 \mathbb C^3
 =\operatorname{span}\{d\}
 \mathbin\oplus \eta(\ker d^\sharp).
 \label{eq:V40-output-splitting}
\end{equation}
\end{lemma}

\begin{proof}
Since \(d^\sharp\eta d=d^Td=q_d\), direct multiplication gives
\(P_d^2=P_d\), \(d^\sharp P_d=0\), and \(P_d\eta d=0\).  Because
\(q_d=-1\ne0\), the displayed kernel is one-dimensional, so the range has
dimension two and equals \(\ker d^\sharp\).

For the last statement, suppose \(W\in\ker d^\sharp\) and
\(\eta W=ad\).  Then \(W=a\eta d\), and applying \(d^\sharp\) gives
\(0=aq_d\), hence \(a=0\).  Thus
\(\eta(\ker d^\sharp)\) meets \(\operatorname{span}\{d\}\) trivially.
Their dimensions are two and one, proving the direct sum.
\end{proof}

The same formula defines a smooth order-zero rational symbol in a
neighborhood of glancing on which \(q_d\ne0\).  It is not a global stable
symbol: \(q_d=s^2-|k|^2\) vanishes at other frequencies.

\subsection{An explicit sharp rank-three correction}
\label{subsec:V40-rank-three-correction}

Order the target of \(\mathcal M_{CM}\) as \((K,C_a,C_n)\).  Define

\begin{equation}
 \Delta_d(\tau,V,\Phi)
 :=\bigl(\Phi,\,\eta P_dV,\,\Phi\bigr).
 \label{eq:V40-Delta}
\end{equation}

The correction acts on the V39 defect sectors by

\begin{equation}
 \Delta_d(-3,0,1)=(1,0,1),
 \label{eq:V40-scalar-pairing}
\end{equation}

\begin{equation}
 \Delta_d(0,V,0)=(0,\eta V,0)
 \quad\hbox{when }d^\sharp V=0.
 \label{eq:V40-vector-pairing}
\end{equation}

Thus it pairs the scalar right defect with the scalar left obstruction and
the two vector right defects with a complement to the vector range.

Let

\begin{equation}
 \sigma:=s+\lambda.
 \label{eq:V40-sigma}
\end{equation}

By Theorem~\ref{thm:V35-Sommerfeld-factor},
\(|\sigma|\ge2^{-1/2}\) on the normalized closed stable hemisphere.

\begin{theorem}[Sharp glancing repair]
\label{thm:V40-sharp-glancing-repair}
At every normalized glancing frequency,

\begin{equation}
 \operatorname{rank}\Delta_d=3,
 \label{eq:V40-Delta-rank}
\end{equation}

and for every nonzero \(\gamma\in\mathbb C\),

\begin{equation}
 \mathcal M_{\gamma,d}
 :=\mathcal M_{CM}(0,d)+\gamma\sigma\Delta_d
 \label{eq:V40-corrected-M}
\end{equation}

is invertible.  Consequently its least singular value has a positive
minimum on the compact normalized glancing set.
\end{theorem}

\begin{proof}
The vector part \(\eta P_d\) has rank two.  The scalar part has rank one and
lands in the \((K,C_n)\)-plane, while the vector part lands only in the
\(C_a\)-slots.  Hence \(\Delta_d\) has rank three.

Suppose \(\mathcal M_{\gamma,d}x=0\).  At glancing, the \(K\) and \(C_n\)
rows are

\begin{equation}
 -d^\sharp V+\gamma\sigma\Phi=0,
 \qquad
 d^\sharp V+\gamma\sigma\Phi=0.
 \label{eq:V40-K-Cn-equations}
\end{equation}

Adding and subtracting, and using \(\gamma\sigma\ne0\), gives
\(\Phi=0\) and \(d^\sharp V=0\).  The three \(C_a\) rows reduce to

\begin{equation}
 -\frac16d\tau+\gamma\sigma\eta V=0.
 \label{eq:V40-Ca-equation}
\end{equation}

The two summands lie in the direct summands of
\eqref{eq:V40-output-splitting}.  Hence both vanish.  Since
\(d\ne0\), \(\tau=0\), and since \(\eta\) is invertible, \(V=0\).
Thus the kernel is zero, proving invertibility.

The matrix depends continuously on the glancing frequency.  A continuous
positive least singular value on a compact set has a positive minimum.
\end{proof}

This proof realizes the algebraic sharpness statement of V39 with a
rotationally covariant formula; it does not choose arbitrary right and left
singular-vector frames.

\subsection{Three-amplitude transparent realization}
\label{subsec:V40-transparent-realization}

Define the rank-three glancing port fibre

\begin{equation}
 \mathsf Y_d:=\ker d^\sharp\oplus\mathbb C.
 \label{eq:V40-Yd}
\end{equation}

For \(x=(\tau,V,\Phi)\), set

\begin{equation}
 J_dx:=(P_dV,\Phi)\in\mathsf Y_d,
 \label{eq:V40-Jd}
\end{equation}

and, for \((W,a)\in\mathsf Y_d\), set

\begin{equation}
 U_d(W,a):=(a,\eta W,a).
 \label{eq:V40-Ud}
\end{equation}

Then \(\Delta_d=U_dJ_d\).  Introduce the independent port amplitude
\(y\in\mathsf Y_d\) and the eight-row system

\begin{equation}
 \mathfrak T_{\gamma,d}(x,y)
 :=\left(
 \mathcal M_{CM}x+\gamma\sigma U_dy,
 \ \sigma(y-J_dx)
 \right).
 \label{eq:V40-transparent-system}
\end{equation}

\begin{corollary}[Regular three-amplitude realization]
\label{cor:V40-three-amplitude-realization}
At every normalized glancing frequency,
\(\mathfrak T_{\gamma,d}\) is an invertible endomorphism of
\(\mathbb C^5\oplus\mathsf Y_d\) whenever \(\gamma\ne0\).  Eliminating the
port gives exactly

\begin{equation}
 \mathcal M_{CM}+\gamma\sigma\Delta_d.
 \label{eq:V40-Schur-correction}
\end{equation}

The port block has inverse norm at most \(\sqrt2\) on the normalized stable
hemisphere wherever the microlocal fibre \(\mathsf Y_d\) is defined.
\end{corollary}

\begin{proof}
The second component of \eqref{eq:V40-transparent-system} has diagonal
block \(\sigma I_{\mathsf Y_d}\).  It is invertible and its inverse norm is
at most \(\sqrt2\).  Its equation gives \(y=J_dx\).  Substitution in the
first component gives \eqref{eq:V40-Schur-correction}, which is invertible
by Theorem~\ref{thm:V40-sharp-glancing-repair}.  The Schur-complement
criterion proves invertibility of the full system.
\end{proof}

Because \(q_d=-1\) on glancing, \(P_d\), \(J_d\), \(U_d\), and the
corrected matrix extend continuously to a common neighborhood of the
glancing set.  Compactness and Theorem~\ref{thm:V40-sharp-glancing-repair}
give a smaller neighborhood with a positive singular floor.  On every
closed chart separated from glancing, V37 already gives a positive floor
for \(\mathcal M_{CM}\).  Thus V40 supplies a two-chart stable parametrix:
the corrected port near glancing and the original CM inverse away from it.
It does not yet supply a single global boundary domain on their overlap.

\subsection{Locality, action reciprocity, and source-incidence tests}
\label{subsec:V40-action-tests}

The word ``transparent'' in
Corollary~\ref{cor:V40-three-amplitude-realization} is essential.  The
projection \eqref{eq:V40-Pd} contains
\((s^2-|k|^2)^{-1}\), and \(\sigma=s+\sqrt{s^2+|k|^2}\) contains the stable
Dirichlet-to-Neumann root.  Hence the port is not the symbol of a finite
tangential-jet boundary density.  It is defined microlocally near glancing,
where its denominator is separated from zero.  No global time-domain
causality or finite-propagation statement follows from the calculation.

There is also an action-reciprocity defect in the displayed coordinates.
With the coordinate Euclidean pairing, the off-diagonal blocks in
\eqref{eq:V40-transparent-system} are

\begin{equation}
 \gamma\sigma U_d
 \quad\hbox{and}\quad
 -\sigma J_d.
 \label{eq:V40-offdiagonal-blocks}
\end{equation}

They are not adjoints: the scalar component of \(U_d\) lands in both the
\(K\) and \(C_n\) rows, whereas the coordinate adjoint of \(J_d\) lands
only in the \(\Phi\)-dual column.  Therefore
\eqref{eq:V40-transparent-system} is not the Hessian system of a quadratic
boundary functional under that naive pairing.

This failure does not prove that no action realization exists.  The correct
pairing is the V37 DeWitt Green form enlarged by a boundary phase form, not
the coordinate Euclidean pairing.  Proposition~\ref{prop:V31-reservoir}
requires the exact identity

\begin{equation}
 \mathfrak b_n(u,v)+\mathfrak b_\partial(y,w)=0
 \label{eq:V40-required-Green-balance}
\end{equation}

on the coupled relation.  V40 has not constructed
\(\mathfrak b_\partial\), so it has not passed this action test.

The companion audit adds a second independent gate.  Even after a Green
balance is found, the port map \(J_d\) must be derived from the actual
boundary term in the coupled Bianchi identity.  The NS V12 source-incidence
result forbids identifying \(J_dx\) with that complete source merely because
it detects every glancing kernel vector.  Detection, action incidence, and
physical source incidence are three separate statements.

\begin{firewall}[Status of the V40 port]
\label{fire:V40-port-status}
The V40 port proves that three amplitudes are sufficient for a uniform
\emph{microlocal glancing rank repair}.  It is rational in the tangential
symbol, uses the square-root stable root, is not globally defined, and has
not been derived from the DeWitt Green pairing or the Bianchi stress.  It is
therefore neither a local conservative boundary action nor a completed
Einstein boundary condition.  These limitations are proved or displayed,
not suppressed.
\end{firewall}

\subsection{Finite certificate and V41 upstream row}
\label{subsec:V40-certificate-handoff}

At a normalized glancing sample store

\begin{equation}
 r_{P^2}:=\lVert P_d^2-P_d\rVert,
 \qquad
 r_{dP}:=\lVert d^\sharp P_d\rVert,
 \qquad
 r_q:=|q_d+1|.
 \label{eq:V40-projector-residuals}
\end{equation}

\begin{equation}
 \rho_{\Delta}:=\operatorname{rank}\Delta_d,
 \qquad
 \rho_{\gamma}:=s_{\min}(\mathcal M_{CM}+\gamma\sigma\Delta_d),
 \qquad
 r_{UJ}:=\lVert\Delta_d-U_dJ_d\rVert.
 \label{eq:V40-correction-certificate}
\end{equation}

A 512-point double-precision sweep over both signs of \(s\) and 256 angular
directions of \(k\), with \(\gamma=1\), returned

\begin{equation}
 \max r_{P^2}=2.87\times10^{-16},
 \qquad
 \max r_{dP}=2.91\times10^{-16},
 \qquad
 \min\rho_\gamma=1.48735\times10^{-1}.
 \label{eq:V40-numerical-projector-certificate}
\end{equation}

Every sample had

\begin{equation}
 \operatorname{rank}P_d=2,
 \qquad
 \operatorname{rank}\Delta_d=3,
 \qquad
 \operatorname{rank}(\mathcal M_{CM}+\sigma\Delta_d)=5.
 \label{eq:V40-numerical-ranks}
\end{equation}

The exact proof is Theorem~\ref{thm:V40-sharp-glancing-repair}; the sweep is
only an implementation and frame check.

V41 should go upstream from coordinate reciprocity to the actual V37
DeWitt boundary phase form.  It must find a rank-three port pair
\((J_{\rm in},J_{\rm out})\) which both induces an isomorphism between the
right and left glancing defects and is reciprocal for a nondegenerate
boundary Green form.  It should then derive the port stress and compare its
normal diffeomorphism identity with the Bianchi boundary term.  If no
three-amplitude local phase form can meet both conditions, the obstruction
will justify a doubled Hamiltonian reservoir or a nonlocal transparent
dilation.

\subsection{V40 ledger}
\label{subsec:V40-ledger}

\paragraph{New conclusions proved in V40.}
The scheduled companion audit imports the NS complete-source incidence gate
and the YM actual-law defect/locality gate without crossing their hypotheses.
Lemma~\ref{lem:V40-vector-projector} constructs the canonical rank-two
vector-defect projection.  Theorem~\ref{thm:V40-sharp-glancing-repair}
gives an explicit rank-three correction and proves its uniform invertibility
on the complete normalized glancing set.  Corollary
\ref{cor:V40-three-amplitude-realization} factors it through exactly three
regular port amplitudes.  The action subsection proves that the displayed
system is nonlocal and fails naive coordinate reciprocity, while retaining
the precise V31 Green-balance test for the next iteration.

\paragraph{Imported conclusions not repeated.}
V39 supplies the three-channel lower bound and the vector/scalar defect
decomposition.  V37 supplies the CM matrix and the uniform Sommerfeld
factor.  NS V12 supplies only its complete-source and summable-native-response
statements.  YM V21 supplies only its actual-law density-defect and locality
statements.  No NS regularity estimate or YM mass-gap estimate is asserted
for the gravitational boundary system.

\paragraph{Honest continuum status.}
The algebraic lower bound of V39 is now attained by an explicit microlocal
three-amplitude stable port, so glancing rank is no longer the immediate
unknown.  The port is not yet a local action, a global causal boundary
operator, or a complete Bianchi source.  Its Green form, stress, ghost
complex, overlap domain, and cofinal SM/YM/NS compatibility remain open.
The full Standard Model--Einstein continuum is therefore not proved.

% END APPENDED VERSION V40 MATERIAL



% BEGIN V40 LATE COMPANION AUDIT ADDENDUM

\subsection{Late-arriving NS V13 audit addendum}
\label{subsec:V40-late-NS13-audit}

During the V40 construction the retained NS companion advanced from V12 to

\begin{center}
\texttt{Renewal\_NS\_Stability\_Research\_Notes\_V13.tex},
\end{center}

with SHA--256

\begin{center}
\texttt{11ddff5af57c02e01651ba0db202e32da5e8359b96dffe71ae548528d4f09ef3}.
\end{center}

The new file was read before V40 was closed.  Its theorem labelled
\texttt{thm:v13-record-exterior-incidence} proves that a first doubling of
the critical record forces a fully assembled exterior response of order
\(R\) and exterior critical-source action of order at least \(\nu R^2\) on
an intrinsic final parabolic subinterval.  Its corollary
\texttt{cor:v13-doubled-record-source-divergence} produces pairwise disjoint
complete-source cards with divergent total critical action along an
unbounded doubled-record chronology.  The same note proves that leverage
alone need not yield high-parent residence and that the covariant-variation
packet retains an unpaid signed quadratic--cubic cancellation.

This new result strengthens, rather than reverses, the V40 audit.  NS V13
shows that complete-source incidence can be forced by an intrinsic event
once the event-to-source identity is proved.  The gravitational analogue is
therefore more specific: V41 should seek an intrinsic boundary event,
constraint identity, or energy-flux threshold which forces nonzero
occupancy of the complete three-port source.  The algebraic fact that
\(J_dx\) detects every glancing defect does not supply such an event or
identity.  No Navier--Stokes estimate is imported into the gravitational
symbol, and the NS file was not modified.

\begin{firewall}[Adjusted V41 incidence target]
\label{fire:V40-adjusted-V41-target}
The V40 transparent port is a candidate carrier.  It becomes a physical
source only after a single coupled action or Bianchi identity proves that
the actual boundary flux lands in that carrier.  An intrinsic trigger may
then force port action, as the record doubling does in NS V13.  Neither the
trigger nor the landing identity has yet been constructed for Einstein--SM.
\end{firewall}

% END V40 LATE COMPANION AUDIT ADDENDUM



% BEGIN V40 LATE YANG--MILLS V22 AUDIT ADDENDUM

\subsection{Late-arriving Yang--Mills V22 audit addendum}
\label{subsec:V40-late-YM22-audit}

Before V40 was delivered, the Yang--Mills companion advanced to

\begin{center}
\texttt{Renewal\_Yang\_Mills\_Mass\_Gap\_Research\_Notes\_V22.tex},
\end{center}

with SHA--256

\begin{center}
\texttt{42ce37c2b26857f7591fe0195a492818200b56ef2930797e77b23d0f6004a367}.
\end{center}

The new appendix was read without modifying it.  Its theorem labelled
\texttt{thm:V22-strong-density-bound} proves, on the branch where the V9
and V21 complete transfers are literally the same operator, the
volume-uniform estimate

\begin{equation}
 \lVert S_{\vartheta,-}\rVert_{\rm op}
 \le C_{\rm den}|\beta|.
 \label{eq:V40-imported-YM22-density-bound}
\end{equation}

Its corollary \texttt{cor:V22-physical-density-scale} converts this to a
uniform physical stock when the still-explicit ratio
\(|\gamma_{\alpha,j}|/\alpha\) is bounded.  The threshold card
\texttt{prop:V22-threshold-card} supplies a literal source-compressed bad
event for the positive density excess.

These statements refine the future quantum ledger: conditional on the
declared complete-transfer incidence, the magnitude of the added density
channel need no longer be left wholly unknown.  They do not turn that
channel into a classical Lorentzian boundary amplitude.  V22 explicitly
leaves open the complete transfer/fibre incidence and the locality of
\((D_\alpha)_+^{1/2}\) across spectral zero.  Consequently the V40
transparent port cannot borrow locality or action incidence from
\eqref{eq:V40-imported-YM22-density-bound}, and the V41 DeWitt/Bianchi
direction is unchanged.

% END V40 LATE YANG--MILLS V22 AUDIT ADDENDUM



% BEGIN APPENDED VERSION V41 MATERIAL

\clearpage
\section{V41: DeWitt-reciprocal rank-three graph, fixed-scale lift, and
the first-order local-action obstruction}
\label{sec:V41-DeWitt-rank-three}

\subsection{Purpose and change of coordinates}
\label{subsec:V41-purpose}

V40 constructed a first-order rank-three correction of the stable quotient
matrix.  It repaired every normalized glancing frequency, but its projector
contained \((s^2-|k|^2)^{-1}\), its Sommerfeld factor contained the stable
square root, and its off-diagonal port blocks were not reciprocal in the
coordinate pairing.  The present note moves upstream to the actual reduced
DeWitt phase pairing.

There are two conclusions.  First, a completely local rank-three
configuration functional produces a reciprocal Lagrangian graph and lifts
the glancing kernel at every fixed frequency scale.  Second, that functional
is zeroth order, so its normalized principal correction decays like
\(\rho^{-1}\) at frequency \(\rho\).  A real local quadratic action in only
three configuration ports cannot replace it by a first-order principal
term: its first-order Euler symbol is a three-by-three skew matrix and has
rank at most two.  Thus V41 identifies the exact reason that V40 obtained
principal rank only by using a transparent nonlocal symbol.

All metric amplitudes below are the full-embedding invariant amplitudes of
V39; hats are suppressed.  Consequently a functional of these amplitudes
is invariant under the frozen four-component gauge orbit.  No new quotient
or gauge direction is introduced.

\subsection{The reduced DeWitt phase graph}
\label{subsec:V41-phase-graph}

Put

\begin{equation}
 x:=(\tau,v_0,v_1,v_2,\phi)^T,
 \qquad
 p:=D_nx.
 \label{eq:V41-x-p}
\end{equation}

Equation~\eqref{eq:V37-restricted-DeWitt} is the symmetric bilinear form

\begin{equation}
 \mathbb G(x,z)=x^THz,
 \label{eq:V41-G-H}
\end{equation}

where, in the block order \((\tau,v,\phi)\),

\begin{equation}
 H=
 \begin{pmatrix}
 -1/6&0&-1/2\cr
 0&2\eta&0\cr
 -1/2&0&1/2
 \end{pmatrix},
 \qquad
 \eta=\operatorname{diag}(-1,1,1).
 \label{eq:V41-DeWitt-matrix}
\end{equation}

The scalar block has determinant \(-1/3\), and \(2\eta\) is invertible, so
\(H\) is nondegenerate.  The V37 phase rows
\eqref{eq:V37-phase-tau}--\eqref{eq:V37-phase-phi} are the graph

\begin{equation}
 p=A(D)x,
 \label{eq:V41-A-graph}
\end{equation}

with

\begin{equation}
 A(D)(\tau,v,\phi)
 :=\left(
 2D^av_a,
 \ D_a\left(\frac\tau6+\frac\phi2\right),
 \ 0
 \right).
 \label{eq:V41-A-definition}
\end{equation}

The stable graph matrix is

\begin{equation}
 L_0(s,k):=\lambda I_5-A(d).
 \label{eq:V41-L0}
\end{equation}

It is row-equivalent to the V37 matrix.  More precisely, its first row is
\(2K\), its next three rows are \(C_a\), and its last row is
\(2(K+C_n)\).  Therefore it has the same open determinant decision and the
same three-dimensional glancing kernel as \(\mathcal M_{CM}\).

For two phase vectors \((x,p)\) and \((z,q)\), use the radial Green form

\begin{equation}
 \Omega_H((x,p),(z,q))
 :=\mathbb G(p,z)-\mathbb G(x,q).
 \label{eq:V41-phase-Green}
\end{equation}

V37 proved that the graph \eqref{eq:V41-A-graph} is Lagrangian after
tangential integration: the remaining integrand is the divergence
\eqref{eq:V37-Green-divergence}.

\subsection{A local reciprocal rank-three generating functional}
\label{subsec:V41-local-generator}

Choose the boundary time direction used by the hyperbolic foliation and let

\begin{equation}
 Jx:=(v_1,v_2,\phi)^T\in\mathbb R^3.
 \label{eq:V41-J}
\end{equation}

Equivalently, \(J\) selects the spatial projection of the mixed normal
trace and the normal-normal trace.  Define

\begin{equation}
 S:=J^TJ,
 \qquad
 B:=H^{-1}S.
 \label{eq:V41-S-B}
\end{equation}

Direct inversion of the scalar DeWitt block gives

\begin{equation}
 B(\tau,v_0,v_1,v_2,\phi)
 =\left(-\frac32\phi,\ 0,\ \frac12v_1,\ \frac12v_2,
 \frac12\phi\right).
 \label{eq:V41-B-explicit}
\end{equation}

In particular,

\begin{equation}
 \operatorname{rank}B=3,
 \qquad
 HB=S=S^T=B^TH.
 \label{eq:V41-B-reciprocity}
\end{equation}

For a real constant \(\mu\), consider the local quadratic boundary
generator

\begin{equation}
 F_\mu[x]
 :=\frac\mu2\int_{\partial M}
 \left(v_1^2+v_2^2+\phi^2\right)\dd^3y.
 \label{eq:V41-Fmu}
\end{equation}

Its variation is

\begin{equation}
 \delta F_\mu
 =\mu\int_{\partial M}\mathbb G(Bx,\delta x)\dd^3y.
 \label{eq:V41-Fmu-variation}
\end{equation}

Choosing the sign of the boundary generating term so that it shifts the
old graph forward gives

\begin{equation}
 \mathcal L_\mu
 :=\{(x,p):p=A(D)x+\mu Bx\}.
 \label{eq:V41-Lmu-graph}
\end{equation}

\begin{theorem}[DeWitt-reciprocal Lagrangian graph]
\label{thm:V41-Lagrangian-graph}
For compactly supported real tangential data, \(\mathcal L_\mu\) is
Lagrangian for \(\Omega_H\) for every real \(\mu\).  The added graph term
is local, rank three, and is the Hessian of the single functional
\eqref{eq:V41-Fmu}.
\end{theorem}

\begin{proof}
The V37 graph with \(\mu=0\) is Lagrangian.  The additional contribution to
the Green pairing of two graph vectors is

\[
 \mu\bigl(\mathbb G(Bx,z)-\mathbb G(x,Bz)\bigr)
 =\mu x^T(B^TH-HB)z=0
\]

by \eqref{eq:V41-B-reciprocity}.  Thus the shifted graph is isotropic.
It remains a graph over all five configuration variables and hence is
half-dimensional in the ten-dimensional phase fibre.  It is therefore
Lagrangian.  Equations~\eqref{eq:V41-B-explicit} and
\eqref{eq:V41-Fmu-variation} prove the other assertions.
\end{proof}

The functional is also frozen-gauge invariant.  Indeed, its arguments are
components of \(Q_4(h,X)=\widehat h\), so
\(Q_4G_p=0\) from Theorem~\ref{thm:V39-full-exact-sequence}.  Its metric and
embedding Euler covectors consequently obey the linear Noether identity of
Proposition~\ref{prop:V39-pullback-Noether}.  This is an actual Euler-source
incidence for the quadratic boundary functional, rather than an arbitrary
stable selector.  It is not yet the complete nonlinear Einstein Bianchi
source.

\subsection{Exact fixed-scale glancing lift}
\label{subsec:V41-fixed-scale-lift}

The stable matrix of \eqref{eq:V41-Lmu-graph} is

\begin{equation}
 L_\mu(s,k):=\lambda I_5-A(d)-\mu B.
 \label{eq:V41-Lmu}
\end{equation}

\begin{theorem}[Local rank-three lift at fixed scale]
\label{thm:V41-fixed-scale-lift}
For every nonzero real \(\mu\), \(L_\mu(0,d)\) is invertible at every
normalized glancing frequency.  Its least singular value has a positive
minimum on the compact normalized glancing set.
\end{theorem}

\begin{proof}
Let \(L_\mu(0,d)x=0\).  The last phase row and
\eqref{eq:V41-B-explicit} give

\begin{equation}
 -\frac\mu2\phi=0,
 \label{eq:V41-lift-phi}
\end{equation}

so \(\phi=0\).  The \(v_0\)-row is

\begin{equation}
 -d_0\left(\frac\tau6+\frac\phi2\right)=0.
 \label{eq:V41-lift-tau}
\end{equation}

At glancing \(d_0=s=\pm i/\sqrt2\ne0\), hence \(\tau=0\).  The two
spatial \(v\)-rows now give

\begin{equation}
 -\frac\mu2v_i=0,
 \qquad i=1,2,
 \label{eq:V41-lift-vsp}
\end{equation}

so \(v_1=v_2=0\).  Finally the \(\tau\)-row reduces to
\(-2d^\sharp v=2sv_0=0\), and therefore \(v_0=0\).  The kernel is zero.
Continuity and compactness give the uniform fixed-scale minimum.
\end{proof}

This lift is reciprocal and local, unlike the V40 transparent port.  The
next theorem explains why it is nevertheless not a principal continuum
repair.

\subsection{The conic high-frequency loss}
\label{subsec:V41-conic-loss}

Let \(\widehat p=(\widehat s,\widehat k)\) be a normalized glancing
frequency, let \(\widehat\lambda=0\), and scale the tangential frequency by
\(\rho>0\).  Homogeneity gives

\begin{equation}
 L_\mu(\rho\widehat p)
 =\rho L_0(\widehat p)-\mu B.
 \label{eq:V41-scaled-Lmu}
\end{equation}

\begin{theorem}[A local zeroth-order lift has zero principal floor]
\label{thm:V41-zero-principal-floor}
For every fixed \(\mu\) and every unit vector
\(r\in\ker L_0(\widehat p)\),

\begin{equation}
 \left\|\rho^{-1}L_\mu(\rho\widehat p)r\right\|
 =\frac{|\mu|}{\rho}\lVert Br\rVert.
 \label{eq:V41-kernel-decay}
\end{equation}

Consequently

\begin{equation}
 s_{\min}\!\left(\rho^{-1}L_\mu(\rho\widehat p)\right)
 \le\frac{|\mu|\lVert B\rVert}{\rho}
 \longrightarrow0.
 \label{eq:V41-zero-floor}
\end{equation}

The order-one principal boundary symbol is still \(L_0\) and retains its
rank-two glancing value.
\end{theorem}

\begin{proof}
Apply \eqref{eq:V41-scaled-Lmu} to
\(r\in\ker L_0(\widehat p)\), divide by \(\rho\), and take norms.  The
least singular value is bounded above by the norm on any unit test vector,
which proves \eqref{eq:V41-zero-floor}.  A zeroth-order term is absent from
the homogeneous order-one principal symbol, proving the last statement.
\end{proof}

Thus the fixed-scale inverse loses one tangential derivative on the old
glancing kernel.  Choosing \(\mu\) proportional to \(\rho\) would restore
the order, but such a choice is an order-one operator rather than a fixed
local coefficient.

\subsection{No three-configuration-port first-order local action}
\label{subsec:V41-first-order-no-go}

The preceding loss suggests replacing the mass generator by a first-order
boundary density.  The following theorem rules out that repair within the
smallest real configuration-only class.

Let \(z=Jx\in\mathbb R^3\), and consider a frozen real local quadratic
functional whose first-order principal part is

\begin{equation}
 F_1[z]
 :=\frac12\int_{\partial M}z^TC^aD_az\,\dd^3y,
 \label{eq:V41-first-order-functional}
\end{equation}

with constant real three-by-three matrices \(C^a\).  Compact support removes
tangential boundary-of-boundary terms.

\begin{theorem}[Odd-port first-order rank obstruction]
\label{thm:V41-odd-port-obstruction}
The Euler operator of \eqref{eq:V41-first-order-functional} has principal
part

\begin{equation}
 \mathcal E_1(D)
 =\frac12\sum_a\bigl(C^a-(C^a)^T\bigr)D_a.
 \label{eq:V41-E1}
\end{equation}

At every real tangential covector its complex Fourier symbol has rank at
most two.  Therefore the induced DeWitt-reciprocal correction

\begin{equation}
 H^{-1}J^T\mathcal E_1(D)J
 \label{eq:V41-induced-first-order-correction}
\end{equation}

has rank at most two and cannot lift the rank-two five-by-five glancing
matrix to rank five.  No real local first-order quadratic action in only
the three configuration ports \((v_1,v_2,\phi)\) supplies the V39 principal
rank budget.
\end{theorem}

\begin{proof}
Vary \eqref{eq:V41-first-order-functional}.  Integrating the derivative on
the second occurrence of \(\delta z\) gives

\[
 \delta F_1
 =\frac12\int\delta z^T
 \sum_a\bigl(C^a-(C^a)^T\bigr)D_az\,\dd^3y,
\]

which proves \eqref{eq:V41-E1}.  For a covector \(\xi\), put

\[
 C(\xi):=\frac12\sum_a
 \bigl(C^a-(C^a)^T\bigr)\xi_a.
\]

This is a real skew-symmetric three-by-three matrix.  Every odd-dimensional
skew-symmetric matrix has zero determinant and even rank, hence rank at most
two.  Multiplication by \(i\), \(H^{-1}J^T\), and \(J\) cannot increase
rank.  The V39 rank inequality then gives total glancing rank at most
\(2+2=4<5\), proving the obstruction.
\end{proof}

\begin{firewall}[Scope of the first-order no-go]
\label{fire:V41-no-go-scope}
Theorem~\ref{thm:V41-odd-port-obstruction} concerns real, local,
configuration-only, first-order quadratic actions on three ports.  It does
not rule out a second-order Wentzell action, four or more configuration
ports, a doubled Hamiltonian phase space, or a nonlocal self-adjoint symbol
such as \(|D_\partial|I_3\).  Each alternative changes the stable row count,
the characteristic set, or locality and must be analyzed as a new coupled
system.
\end{firewall}

\subsection{Source incidence and the next upstream system}
\label{subsec:V41-source-handoff}

The quadratic functional \eqref{eq:V41-Fmu} improves the V40 source status
in one limited but exact sense.  Its rank-three force is the derivative of a
single frozen gauge-invariant functional, and its metric and embedding
forces obey the corresponding linear Noether identity.  There is no
unlisted ``raw'' component inside this particular force.  However, it is a
lower-order force and depends on the chosen boundary observer splitting.
Metric variation of that observer, the unit normal, the moving embedding,
and the volume density has not been included.  Hence it is not yet the
complete boundary stress in the nonlinear Bianchi identity required by the
V40 companion audit.

The smallest next candidate is a doubled Hamiltonian reservoir with three
configuration ports and three conjugate momenta.  Its first-order kinetic
matrix can be a nondegenerate six-by-six skew form, avoiding the odd-rank
obstruction while retaining three physical incoming channel amplitudes.
V42 should write one local quadratic boundary action for that phase field,
derive its Green form and reciprocal coupling to the DeWitt graph, and
analyze the full stable matrix at both the gravitational glancing cone and
the reservoir's own characteristic cone.  Schur elimination is permitted
only where the reservoir block is invertible.

\subsection{Finite certificate}
\label{subsec:V41-certificate}

The finite fixed-scale card stores

\begin{equation}
 r_{HB}:=\lVert HB-J^TJ\rVert,
 \qquad
 r_{B^TH}:=\lVert B^TH-HB\rVert,
 \qquad
 \rho_B:=\operatorname{rank}B.
 \label{eq:V41-reciprocity-certificate}
\end{equation}

\begin{equation}
 \rho_{L,\mathcal G}:=operatorname{rank}L_\mu(0,d),
 \qquad
 \sigma_{L,\mathcal G}:=s_{\min}(L_\mu(0,d)).
 \label{eq:V41-lift-certificate}
\end{equation}

For \(\mu=1\), a 512-point double-precision sweep over both glancing signs
and 256 angular directions returned

\begin{equation}
 r_{HB}=0,
 \qquad
 r_{B^TH}=0,
 \qquad
 \rho_B=3,
 \qquad
 \rho_{L,\mathcal G}=5,
 \label{eq:V41-numerical-ranks}
\end{equation}

and

\begin{equation}
 \min_{\mathcal G}\sigma_{L,\mathcal G}
 =9.09317\times10^{-2}.
 \label{eq:V41-numerical-floor}
\end{equation}

Theorem~\ref{thm:V41-fixed-scale-lift} is the exact proof.  The frequency
scaling test must additionally show the \(O(\rho^{-1})\) normalized decay
in \eqref{eq:V41-zero-floor}; a fixed-scale numerical floor is not a
principal estimate.

\subsection{V41 ledger}
\label{subsec:V41-ledger}

\paragraph{New conclusions proved in V41.}
The V37 CM relation is rewritten as the DeWitt phase graph
\(p=A(D)x\).  Theorem~\ref{thm:V41-Lagrangian-graph} constructs a local
rank-three reciprocal graph shift from one quadratic boundary functional.
Theorem~\ref{thm:V41-fixed-scale-lift} proves exact fixed-scale glancing
invertibility.  Theorem~\ref{thm:V41-zero-principal-floor} proves that this
zeroth-order repair loses one derivative and leaves the principal glancing
rank unchanged.  Theorem~\ref{thm:V41-odd-port-obstruction} proves that a
three-configuration-port real local first-order action cannot supply the
missing principal rank.

\paragraph{Imported conclusions not repeated.}
V37 supplies the restricted DeWitt form, phase rows, and CM glancing kernel.
V39 supplies the full embedding quotient and rank-three lower bound.  V40
supplies the nonlocal first-order rank repair and the current NS/YM source
incidence gates.  No companion-program theorem is reinterpreted as a
DeWitt boundary action.

\paragraph{Honest continuum status.}
A local gauge-invariant reciprocal rank-three boundary generator now exists
and has an exact frozen Noether source, but it is lower order and has zero
principal causal floor.  The minimal three-port first-order
configuration-only action is ruled out.  A doubled Hamiltonian or
second-order boundary dynamics, its complete stress/Bianchi identity, and
its global stable estimate remain unconstructed.  The full Standard
Model--Einstein continuum is therefore not proved.

% END APPENDED VERSION V41 MATERIAL



% BEGIN APPENDED VERSION V42 MATERIAL

\clearpage
\section{V42: local Hamiltonian three-port, exact Green balance, and the
surface-wave sign dichotomy}
\label{sec:V42-Hamiltonian-port}

\subsection{Purpose and correction of the V41 handoff}
\label{subsec:V42-purpose}

V41 proved that a first-order real action on the three metric port
coordinates alone has rank at most two.  It proposed a doubled Hamiltonian
reservoir as the next upstream object.  The minimal doubling does not require
six new fields: the three metric defect coordinates can supply one half of
the canonical phase space.  V42 therefore adds only three conjugate boundary
fields.  This attains the V39 channel lower bound exactly.

The resulting boundary action is local, has a nonnegative subluminal port
energy, is reciprocal for the DeWitt graph, and satisfies the V31
bulk--boundary Green balance at the frozen linear level.  Its full
eight-by-eight symbol is invertible on the gravitational glancing cone and
on the reservoir's own characteristic cone.  A complete stable analysis
nevertheless reveals a new obstruction: the two possible Green-orientation
signs give either an open-half-plane growing transverse mode or a purely
imaginary decaying transverse surface wave.  Thus V42 constructs the local
action and proves its exact failure mode; it does not claim a global causal
floor.

Retain the V41 variables, DeWitt matrix, phase graph, and port map

\begin{equation}
 x=(\tau,v_0,v_1,v_2,\phi)^T,
 \qquad
 p=D_nx,
 \qquad
 z:=Jx=(v_1,v_2,\phi)^T.
 \label{eq:V42-x-z}
\end{equation}

All components of \(x\) are the full-embedding invariant components of
V39.  The new real boundary field

\begin{equation}
 y=(y_1,y_2,y_3)^T
 \label{eq:V42-y}
\end{equation}

is BRST inert at the frozen level.

\subsection{One local boundary action and all of its Euler rows}
\label{subsec:V42-boundary-action}

Fix a boundary propagation speed

\begin{equation}
 0<c<1
 \label{eq:V42-subluminal-c}
\end{equation}

and let \(D_i\), \(i=1,2\), denote the spatial derivatives in the chosen
boundary hyperbolic foliation.  Define

\begin{equation}
 S_{\partial,c}^{(2)}[x,y]
 :=\int_{\partial M}
 \left[
  y^TD_tz-\frac12y^Ty
  -\frac{c^2}{2}D_iz^TD_iz
 \right]\dd^3y.
 \label{eq:V42-boundary-action}
\end{equation}

The absence of a mass term makes the principal system homogeneous.  A mass
can be restored as a lower-order term only after the conic estimate is
settled.

\begin{proposition}[Exact first variation]
\label{prop:V42-first-variation}
For compact tangential support,

\begin{equation}
 \delta S_{\partial,c}^{(2)}
 =\int_{\partial M}
 \left[
  \delta y^T(D_tz-y)
  +\delta z^T(-D_ty+c^2\Delta_yz)
 \right]\dd^3y
 \label{eq:V42-first-variation}
\end{equation}

up to the time-corner one-form

\begin{equation}
 \Theta_{\partial}(\delta z)=\int y^T\delta z\,\dd^2y.
 \label{eq:V42-corner-one-form}
\end{equation}

Consequently the port Euler equations are

\begin{equation}
 D_tz-y=0,
 \qquad
 -D_ty+c^2\Delta_yz=0
 \label{eq:V42-free-port-Euler}
\end{equation}

when the port is uncoupled from the bulk metric force.
\end{proposition}

\begin{proof}
Vary the three terms in \eqref{eq:V42-boundary-action}.  Integrating
\(y^TD_t\delta z\) once in time gives
\(-D_ty^T\delta z\) and the corner term
\eqref{eq:V42-corner-one-form}.  Integrating the spatial kinetic variation
once gives \(c^2\Delta_yz^T\delta z\).  The variation in \(y\) is
\(\delta y^T(D_tz-y)\).  This is exactly
\eqref{eq:V42-first-variation}.
\end{proof}

The port Hamiltonian is

\begin{equation}
 E_{\partial,c}[z,y]
 :=\frac12\int
 \left(|y|^2+c^2|D_iz|^2\right)\dd^2y
 \ge0.
 \label{eq:V42-port-energy}
\end{equation}

Its characteristic speed is \(c<1\), so its cone lies strictly inside the
bulk light cone in the frozen boundary frame.  This statement concerns the
port sector only; it is not yet a symmetrizer for the complete gauge-reduced
Einstein system.

\subsection{Reciprocal coupling to the DeWitt graph}
\label{subsec:V42-DeWitt-coupling}

Put

\begin{equation}
 C:=H^{-1}J^T.
 \label{eq:V42-C}
\end{equation}

Using \eqref{eq:V41-DeWitt-matrix},

\begin{equation}
 C(y_1,y_2,y_3)
 =\left(-\frac32y_3,\ 0,\ \frac12y_1,\ \frac12y_2,\
 \frac12y_3\right)^T,
 \qquad
 CJ=B.
 \label{eq:V42-C-explicit}
\end{equation}

Let \(\varepsilon\in\{+1,-1\}\) record the relative Green orientation of
the boundary action and the bulk outward conormal.  Define the coupled
relation

\begin{equation}
 D_tz-y=0,
 \label{eq:V42-coupled-y-row}
\end{equation}

\begin{equation}
 H(p-A(D)x)
 =\varepsilon J^T\left(D_ty-c^2\Delta_yz\right).
 \label{eq:V42-coupled-metric-row}
\end{equation}

The sign is retained rather than guessed from a normal-orientation
convention.  Both choices are tested below.

\begin{theorem}[Exact frozen bulk--boundary Green balance]
\label{thm:V42-Green-balance}
Let \((x,y)\) and \((x',y')\) satisfy
\eqref{eq:V42-coupled-y-row}--\eqref{eq:V42-coupled-metric-row}.  Modulo the
V37 tangential divergence of the unperturbed CM graph, their bulk radial
Green form is

\begin{equation}
 \Omega_H
 =\varepsilon\left(D_t\mathcal J^0+D_i\mathcal J^i\right),
 \label{eq:V42-Green-divergence}
\end{equation}

where

\begin{equation}
 \mathcal J^0:=y^Tz'-z^Ty',
 \label{eq:V42-time-current}
\end{equation}

\begin{equation}
 \mathcal J^i
 :=-c^2\left((D_iz)^Tz'-z^TD_iz'\right).
 \label{eq:V42-space-current}
\end{equation}

Thus the enlarged relation obeys the V31 conservative reservoir identity
after the boundary current is included with the opposite induced
orientation.
\end{theorem}

\begin{proof}
Write

\[
 f:=D_ty-c^2\Delta_yz.
\]

The difference between the shifted and unshifted bulk Green forms is, by
\eqref{eq:V42-coupled-metric-row},

\[
 \varepsilon(f^Tz'-z^Tf').
\]

The mass-free spatial operator is formally self-adjoint.  Direct
differentiation, using \(D_tz=y\) and \(D_tz'=y'\), gives

\[
 D_t(y^Tz'-z^Ty')
 =D_ty^Tz'-z^TD_ty'
\]

and

\[
 D_i\mathcal J^i
 =-c^2\left((\Delta_yz)^Tz'-z^T\Delta_yz'\right).
\]

Their sum is \(f^Tz'-z^Tf'\), proving
\eqref{eq:V42-Green-divergence}.  Integration over a boundary slab gives
the V31 cancellation with the boundary current assigned the opposite
Stokes orientation.
\end{proof}

The same equations give the port power identity

\begin{equation}
 \frac{\dd}{\dd t}E_{\partial,c}
 =\int (D_tz)^T\left(D_ty-c^2\Delta_yz\right)\dd^2y.
 \label{eq:V42-port-power}
\end{equation}

Thus the force in the coupled metric row is the force whose work changes
the port energy.  This is the complete force of the frozen quadratic port
action, not a projected response fitted after the fact.

\subsection{The full stable matrix and its Schur form}
\label{subsec:V42-stable-matrix}

Use the stable ansatz of V37.  Define

\begin{equation}
 a_c(s,k):=s^2+c^2|k|^2.
 \label{eq:V42-ac}
\end{equation}

The full eight-row stable operator is

\begin{equation}
 \mathfrak L_{\varepsilon,c}(s,k)(x,y)
 :=\left(
 L_0x-\varepsilon C\left(sy+c^2|k|^2Jx\right),
 \ sJx-y
 \right).
 \label{eq:V42-full-stable-operator}
\end{equation}

The \(y\)-diagonal block is \(-I_3\), so exact elimination is legal and
gives

\begin{equation}
 y=sJx,
 \qquad
 L_{\rm eff}^{\varepsilon,c}x
 =\left(L_0-\varepsilon a_cB\right)x=0.
 \label{eq:V42-effective-operator}
\end{equation}

This Schur calculation is used only for exact rank.  The eliminated row is
second order in \(z\), so a global estimate must use Wentzell/Douglis--
Nirenberg weights rather than treating
\eqref{eq:V42-effective-operator} as an ordinary first-order boundary row.

\subsection{The two declared characteristic cones are repaired}
\label{subsec:V42-two-cones}

\begin{theorem}[Gravitational glancing is lifted locally]
\label{thm:V42-gravitational-glancing}
Let \(0<c<1\).  On the gravitational glancing cone

\begin{equation}
 s=\pm i|k|,
 \qquad \lambda=0,
 \qquad k\ne0,
 \label{eq:V42-gravitational-cone}
\end{equation}

the full operator \(\mathfrak L_{\varepsilon,c}\) is invertible for both
values of \(\varepsilon\).
\end{theorem}

\begin{proof}
On \eqref{eq:V42-gravitational-cone},

\begin{equation}
 a_c=-(1-c^2)|k|^2\ne0.
 \label{eq:V42-ac-gravitational}
\end{equation}

The Schur complement is the V41 matrix
\(L_\mu=L_0-\mu B\) with
\(\mu=\varepsilon a_c\ne0\).  The component proof of
Theorem~\ref{thm:V41-fixed-scale-lift} uses only \(\mu\ne0\), so it gives
invertibility.  Since the eliminated block is \(-I_3\), the full operator
is invertible as well.
\end{proof}

\begin{theorem}[The port characteristic cone is also lifted]
\label{thm:V42-port-cone}
On the nonzero port cone

\begin{equation}
 s=\pm ic|k|,
 \qquad k\ne0,
 \label{eq:V42-port-cone}
\end{equation}

one has \(a_c=0\) and

\begin{equation}
 \lambda^2=(1-c^2)|k|^2>0.
 \label{eq:V42-lambda-port-cone}
\end{equation}

Hence \(\mathfrak L_{\varepsilon,c}\) is invertible for both signs.
\end{theorem}

\begin{proof}
Equation~\eqref{eq:V42-port-cone} gives \(a_c=0\), so the Schur complement
\eqref{eq:V42-effective-operator} is exactly \(L_0\).  The displayed value
of \(\lambda^2\) is positive because \(0<c<1\), and the stable branch has
\(\lambda>0\).  The V37 determinant therefore makes \(L_0\) invertible.
Schur equivalence proves the full assertion.
\end{proof}

The two theorems show that cone coincidence is not the remaining problem.
The obstruction lies on a hybrid transverse branch between the declared
cones.

\subsection{Exact orientation dichotomy}
\label{subsec:V42-sign-dichotomy}

Rotate the spatial frequency so that \(k=(r,0)\), \(r\ge0\), and consider
the transverse amplitude

\begin{equation}
 x_{\rm T}:=(0,0,0,1,0)^T.
 \label{eq:V42-transverse-vector}
\end{equation}

For this vector the effective equation has the single factor

\begin{equation}
 f_{\varepsilon,c}(s,r)
 :=\lambda-\frac{\varepsilon}{2}\left(s^2+c^2r^2\right).
 \label{eq:V42-transverse-factor}
\end{equation}

The remaining port amplitude is fixed by \(y=sJx_{\rm T}\).

\begin{theorem}[Growing mode for the positive orientation]
\label{thm:V42-growing-mode}
For \(\varepsilon=+1\), every \(0<c<1\) admits the exact open stable mode

\begin{equation}
 k=0,
 \qquad s=2,
 \qquad \lambda=2,
 \qquad x=x_{\rm T},
 \qquad y=2Jx_{\rm T}.
 \label{eq:V42-growing-mode}
\end{equation}

Thus this orientation is not even pointwise complementing in the open
Laplace half-plane.
\end{theorem}

\begin{proof}
At \(k=0\), the stable root is \(\lambda=s\) for real \(s>0\).  Hence

\[
 f_{+,c}(2,0)=2-\frac{2^2}{2}=0.
\]

All nontransverse rows vanish on \(x_{\rm T}\), and the second block of
\eqref{eq:V42-full-stable-operator} vanishes for the displayed \(y\).
The mode is nonzero and decays normally as \(e^{-2x}\).
\end{proof}

The opposite sign removes this growing mode but cannot give a closed-
half-plane floor.

\begin{theorem}[Exact decaying surface-wave family for the negative
orientation]
\label{thm:V42-surface-wave}
Let \(\varepsilon=-1\), \(0<c<1\), and choose any \(r>0\).  Put

\begin{equation}
 A_r:=\sqrt{1+(1-c^2)r^2},
 \qquad
 \lambda_r:=A_r-1>0,
 \label{eq:V42-Ar-lambda}
\end{equation}

and

\begin{equation}
 q_r:=2-c^2r^2-2A_r<0,
 \qquad
 s_r:=\pm i\sqrt{-q_r}.
 \label{eq:V42-qr-sr}
\end{equation}

Then

\begin{equation}
 \lambda_r^2=s_r^2+r^2,
 \qquad
 s_r^2+c^2r^2=-2\lambda_r,
 \label{eq:V42-surface-identities}
\end{equation}

and

\begin{equation}
 f_{-,c}(s_r,r)=0.
 \label{eq:V42-surface-zero}
\end{equation}

Consequently \((x_{
m T},s_rJx_{
m T})\) is a nonzero stable boundary
mode with purely imaginary Laplace frequency and strictly positive normal
decay.  The negative orientation has no uniform closed-half-plane
Lopatinski floor.
\end{theorem}

\begin{proof}
First,

\[
 q_r+r^2
 =2+(1-c^2)r^2-2A_r
 =(A_r-1)^2=\lambda_r^2.
\]

Also

\[
 q_r+c^2r^2=2-2A_r=-2\lambda_r.
\]

These prove \eqref{eq:V42-surface-identities}.  To see \(q_r<0\), note
that \(A_r>1\); the identity
\(q_r+r^2=(A_r-1)^2<r^2\), which follows from
\(A_r-1<\sqrt{1-c^2}\,r<r\), gives \(q_r<0\).  Finally,

\[
 f_{-,c}(s_r,r)
 =\lambda_r+\frac12(s_r^2+c^2r^2)
 =\lambda_r-\lambda_r=0.
\]

The stable root is \(\lambda_r>0\), so the mode decays normally.  Its
Laplace frequency is purely imaginary because \(q_r<0\).
\end{proof}

\begin{corollary}[The local conservative port does not close the causal
floor]
\label{cor:V42-no-global-floor}
Neither relative Green orientation in
\eqref{eq:V42-coupled-metric-row} yields a uniform stable inverse for the
action \eqref{eq:V42-boundary-action}: one has an open growing mode and the
other has a closed-boundary surface-wave family.
\end{corollary}

\begin{proof}
Combine Theorems~\ref{thm:V42-growing-mode} and
\ref{thm:V42-surface-wave}.
\end{proof}

\subsection{Why V42 does not contradict the V41 odd-rank theorem}
\label{subsec:V42-no-contradiction}

The principal cross term in the action is

\begin{equation}
 \int y^TD_tz.
 \label{eq:V42-cross-principal}
\end{equation}

On the combined six-dimensional defect-port space \((z,y)\), its
first-order Euler coefficient is the canonical skew form

\begin{equation}
 \mathbb J_6=
 \begin{pmatrix}
 0&-I_3\cr
 I_3&0
 \end{pmatrix},
 \qquad
 \det\mathbb J_6=1.
 \label{eq:V42-J6}
\end{equation}

V41 considered a three-dimensional configuration space and hence an odd
skew matrix of rank at most two.  V42 adds three independent conjugate
fields and obtains an even, nondegenerate skew matrix.  Eliminating \(y\)
turns the system into the second-order Wentzell operator
\eqref{eq:V42-effective-operator}; it does not produce a forbidden
first-order three-by-three configuration symbol.

\begin{firewall}[Action success versus causal failure]
\label{fire:V42-status}
V42 passes three gates that V40 did not: locality of the frozen action,
DeWitt reciprocity, and exact V31 Green balance with a nonnegative port
energy.  It also passes both declared cone-rank tests.  It fails the global
stable test by explicit modes.  Therefore its action is a valid diagnostic
reservoir, not an admissible final Einstein boundary condition.
\end{firewall}

\subsection{Finite certificate and the V43 upstream row}
\label{subsec:V42-certificate-handoff}

For a proposed implementation store

\begin{equation}
 r_C:=\lVert HC-J^T\rVert,
 \qquad
 r_{CJ}:=\lVert CJ-B\rVert,
 \qquad
 r_{\rm Green}:=\lVert\Omega_H-\varepsilon D_a\mathcal J^a\rVert.
 \label{eq:V42-action-residuals}
\end{equation}

On each characteristic sample also store the full eight-by-eight rank and
least singular value.  The exact failure witnesses require the residuals

\begin{equation}
 r_+:=\lVert\mathfrak L_{+,c}(2,0)(x_{\rm T},2Jx_{\rm T})\rVert,
 \label{eq:V42-growing-residual}
\end{equation}

\begin{equation}
 r_-(r):=
 \left\lVert
 \mathfrak L_{-,c}(s_r,(r,0))
 (x_{\rm T},s_rJx_{\rm T})
 \right\rVert.
 \label{eq:V42-surface-residual}
\end{equation}

Both are zero by the preceding theorems.

As a roundoff certificate, take \(c=1/2\), sample 256 spatial
directions and both temporal signs on each normalized cone, and evaluate the
full matrix \eqref{eq:V42-full-stable-operator}.  Both orientations have
rank eight at every sample.  The least sampled singular value is
\(0.0482690\) on gravitational glancing; on the port cone it is
\(0.3107603\) for \(\varepsilon=+1\) and \(0.3029099\) for
\(\varepsilon=-1\).  At \(r=1\), the negative-orientation witness has
\[
 q_r=-0.895751311065,\qquad
 \lambda_r=0.322875655532,
\]
and its full residual is zero to working precision; the positive witness
also has zero residual.  The Schur residual is below
\(4.3\times10^{-15}\).  These samples audit the displayed formulas and do
not replace the exact proofs above.

V43 should retain the local Green-balanced phase space and alter the
impedance rather than returning to a raw selector.  The first tests are:
add the most general local gyroscopic term compatible with the positive
port Hamiltonian; compute its transverse dispersion polynomial; and decide
whether it can move the surface family into the stable left half-plane
without recreating an open growing root.  If conservative finite-dimensional
ports necessarily retain imaginary surface spectrum, the next admissible
route is a larger conservative dilation whose retarded bulk projection is
dissipative, exactly as distinguished in Proposition~\ref{prop:V31-reservoir}.
The nonlinear metric/observer variation and stress/Bianchi identity remain
behind that stable decision.

\subsection{V42 ledger}
\label{subsec:V42-ledger}

\paragraph{New conclusions proved in V42.}
Proposition~\ref{prop:V42-first-variation} derives the complete Euler rows
of a local three-field Hamiltonian port.  Theorem~\ref{thm:V42-Green-balance}
proves exact frozen V31 Green cancellation and
\eqref{eq:V42-port-energy} gives its nonnegative subluminal energy.
Theorems~\ref{thm:V42-gravitational-glancing} and
\ref{thm:V42-port-cone} prove invertibility on both declared characteristic
cones.  Theorems~\ref{thm:V42-growing-mode} and
\ref{thm:V42-surface-wave} prove the complete orientation dichotomy by
explicit stable modes.

\paragraph{Imported conclusions not repeated.}
V41 supplies \(H,J,B\), the reciprocal graph criterion, and the fixed-scale
rank-three lift.  V39 supplies the three-channel lower bound.  V40 supplies
the complete-source incidence requirement from the companion audit.  No NS
or YM estimate is used to infer a gravitational stable bound.

\paragraph{Honest continuum status.}
A minimal local gauge-invariant Hamiltonian boundary action with exactly
three new fields, a positive port energy, and an exact frozen Green balance
is now constructed.  Its global stable inverse is disproved for both
orientation signs.  A stable conservative dilation or justified retarded
dissipative projection, nonlinear covariant stress, ghost completion, and
the coupled Standard-Model continuum remain open.  The full Standard
Model--Einstein continuum is therefore not proved.

% END APPENDED VERSION V42 MATERIAL



% BEGIN APPENDED VERSION V43 MATERIAL

\clearpage
\section{V43: the complete constant-gyroscope test and a
high-frequency surface-spectrum no-go}
\label{sec:V43-gyroscope-no-go}

\subsection{Purpose and scope}
\label{subsec:V43-purpose}

V42 constructed the minimal three-port Hamiltonian action and proved that
its two Green orientations fail in different ways.  Its V43 handoff asked
whether the failure could be removed by a local gyroscopic term without
changing the positive port Hamiltonian.  This note answers that question
for the complete constant-coefficient quadratic class at the same
differential order.

The answer is negative.  Every real three-by-three gyroscope leaves the
positive orientation with an exact open-half-plane growing mode.  For the
negative orientation, an exact response-map determinant shows that
high-frequency, purely imaginary, normally decaying surface modes survive
for every gyroscope and every \(0<c<1\).  The exceptional value
\(c^2=3/4\), where the leading longitudinal response vanishes, is treated
separately rather than removed from the statement.

This is a no-go result for frequency-independent finite local gyroscopes.
It does not rule out a larger conservative dilation or the retarded
positive-real impedance obtained by projecting an infinite bath.

Retain the V42 port variables
\[
 z=Jx=(v_1,v_2,\phi)^T,\qquad y\in\mathbb R^3,
\]
the DeWitt matrices \(H,C=H^{-1}J^T\), and the speed \(0<c<1\).

\subsection{The most general constant gyroscopic term}
\label{subsec:V43-general-gyro}

Let
\begin{equation}
 \Gamma^T=-\Gamma,\qquad \Gamma\in{\rm Mat}_{3\times3}(\mathbb R),
 \label{eq:V43-Gamma-skew}
\end{equation}
and define
\begin{equation}
 S_{\partial,c,\Gamma}^{(2)}
 :=S_{\partial,c}^{(2)}
 +\frac12\int_{\partial M}z^T\Gamma D_tz\,\dd^3y.
 \label{eq:V43-gyro-action}
\end{equation}

\begin{proposition}[Completeness at first time order]
\label{prop:V43-gyro-complete}
Modulo a time-corner functional, every real constant quadratic term that
is bilinear in \(z,D_tz\) is of the form added in
\eqref{eq:V43-gyro-action} for one and only one skew matrix \(\Gamma\).
The term changes neither the Hamiltonian
\eqref{eq:V42-port-energy} nor its nonnegativity.
\end{proposition}

\begin{proof}
Write the general term as
\[
 \frac12\int z^TKD_tz\,\dd^3y
\]
and decompose \(K=K_{\rm sym}+K_{\rm skew}\).  Since
\[
 z^TK_{\rm sym}D_tz
 =\frac12D_t(z^TK_{\rm sym}z),
\]
the symmetric part is a time-corner term.  The remaining matrix
\(\Gamma:=K_{\rm skew}\) is skew and is unique.  A first-order symplectic
term changes the canonical one-form but contributes no term to the
Legendre energy.  Equivalently,
\((D_tz)^T\Gamma D_tz=0\), so the power identity below has exactly the V42
energy.
\end{proof}

The frozen full-embedding invariance of \(z\) and BRST inertness of \(y\)
make \eqref{eq:V43-gyro-action} frozen-gauge invariant.  No new selector is
inserted.

\begin{proposition}[Euler rows, power, and Green current]
\label{prop:V43-Euler-Green}
The first variation of \eqref{eq:V43-gyro-action} gives
\begin{equation}
 D_tz-y=0,
 \qquad
 -D_ty+\Gamma D_tz+c^2\Delta_yz=0
 \label{eq:V43-free-Euler}
\end{equation}
for the uncoupled port.  Put
\begin{equation}
 f_\Gamma:=D_ty-\Gamma D_tz-c^2\Delta_yz.
 \label{eq:V43-force}
\end{equation}
The coupled metric row is
\begin{equation}
 H(p-A(D)x)=\varepsilon J^Tf_\Gamma,
 \qquad \varepsilon\in\{+1,-1\}.
 \label{eq:V43-coupled-row}
\end{equation}
Its port power identity is
\begin{equation}
 \frac{\dd}{\dd t}E_{\partial,c}
 =\int(D_tz)^Tf_\Gamma\,\dd^2y.
 \label{eq:V43-power}
\end{equation}
For two solutions, the V42 Green identity remains exact after replacing
its time current by
\begin{equation}
 {\cal J}^0_\Gamma
 :=y^Tz'-z^Ty'+z^T\Gamma z'.
 \label{eq:V43-Green-current}
\end{equation}
Thus the gyroscopic action is conservative and DeWitt reciprocal at the
frozen level.
\end{proposition}

\begin{proof}
The variation of the added term is
\[
 \delta\left(\frac12z^T\Gamma D_tz\right)
 =\delta z^T\Gamma D_tz
\]
up to its time corner; skewness is used after integrating the second
variation term by parts.  This proves \eqref{eq:V43-free-Euler} and
\eqref{eq:V43-force}.  In the energy derivative the new force contributes
\[
 -(D_tz)^T\Gamma D_tz=0,
\]
which proves \eqref{eq:V43-power}.

The additional term in the skew Green pairing is
\[
 (-\Gamma D_tz)^Tz'-z^T(-\Gamma D_tz')
 =(D_tz)^T\Gamma z'+z^T\Gamma D_tz'
 =D_t(z^T\Gamma z').
\]
Adding this derivative to the V42 current proves the last assertion.
\end{proof}

\subsection{Stable symbol and the exact bulk response}
\label{subsec:V43-response}

For the stable ansatz define
\begin{equation}
 W_\Gamma(s,k)
 :=a_c(s,k)I_3-s\Gamma,
 \qquad
 a_c(s,k)=s^2+c^2|k|^2.
 \label{eq:V43-W}
\end{equation}
The full eight-row operator is
\begin{equation}
 {\mathfrak L}_{\varepsilon,c,\Gamma}(x,y)
 :=
 \left(
 L_0x-\varepsilon C\bigl(
 sy-s\Gamma Jx+c^2|k|^2Jx
 \bigr),
 \ sJx-y
 \right).
 \label{eq:V43-full-operator}
\end{equation}
Eliminating the invertible \(y\)-block gives
\begin{equation}
 L_{\rm eff}^{\varepsilon,c,\Gamma}
 =L_0-\varepsilon CW_\Gamma J.
 \label{eq:V43-effective}
\end{equation}

Fix a nonzero spatial frequency and rotate coordinates so that
\(k=(r,0)\), \(r>0\).  Use the stable bulk root
\begin{equation}
 \lambda^2=s^2+r^2.
 \label{eq:V43-bulk-root}
\end{equation}

\begin{lemma}[Exact three-port response map]
\label{lem:V43-response}
Whenever \(\lambda\ne0\), the solution of
\[
 L_0x=Ch
\]
has port trace
\begin{equation}
 Jx=R(s,r)h,
 \qquad
 R(s,r)=\operatorname{diag}(d,b,b),
 \label{eq:V43-R}
\end{equation}
where
\begin{equation}
 b:=\frac1{2\lambda},
 \qquad
 d:=\frac{4\lambda^2-r^2}{8\lambda^3}.
 \label{eq:V43-b-d}
\end{equation}
In particular \(L_0\) is invertible on the stable sheet when
\(\lambda>0\).
\end{lemma}

\begin{proof}
Write \(h=(h_1,h_2,h_3)^T\) and
\[
 \chi:=\frac{\tau}{6}+\frac{\phi}{2}.
\]
The last, transverse, temporal, and longitudinal rows of \(L_0x=Ch\)
give, respectively,
\begin{equation}
 \lambda\phi=\frac{h_3}{2},\qquad
 \lambda v_2=\frac{h_2}{2},\qquad
 \lambda v_0=s\chi,\qquad
 \lambda v_1-ir\chi=\frac{h_1}{2}.
 \label{eq:V43-response-rows}
\end{equation}
Since \(\tau=6\chi-3\phi\), substitution into the first row gives
\[
 8\lambda\chi-\frac{ir}{\lambda}h_1=0,
 \qquad
 \chi=\frac{ir}{8\lambda^2}h_1,
\]
where \eqref{eq:V43-bulk-root} was used.  Therefore
\[
 v_1
 =\frac{h_1}{2\lambda}-\frac{r^2h_1}{8\lambda^3}
 =\frac{4\lambda^2-r^2}{8\lambda^3}h_1,
 \quad
 v_2=\frac{h_2}{2\lambda},
 \quad
 \phi=\frac{h_3}{2\lambda}.
\]
This is \eqref{eq:V43-R}.  With \(h=0\), these equations give
\(Jx=0\), then \(\chi=v_0=\tau=0\), proving injectivity and hence
invertibility of the square matrix \(L_0\).
\end{proof}

Resolve \(\Gamma\) in the rotated port basis
\((v_L,v_T,\phi)\):
\begin{equation}
 \Gamma=
 \begin{pmatrix}
 0&\gamma_{LT}&\gamma_{L\phi}\cr
 -\gamma_{LT}&0&\gamma_{T\phi}\cr
 -\gamma_{L\phi}&-\gamma_{T\phi}&0
 \end{pmatrix}.
 \label{eq:V43-Gamma-components}
\end{equation}
Set
\begin{equation}
 \alpha:=\gamma_{T\phi},
 \qquad
 \beta^2:=\gamma_{LT}^2+\gamma_{L\phi}^2.
 \label{eq:V43-alpha-beta}
\end{equation}

\begin{lemma}[Exact negative-orientation determinant]
\label{lem:V43-determinant}
For \(\varepsilon=-1\), put
\begin{equation}
 P_1:=1+a_cd,
 \qquad
 P_2:=1+a_cb.
 \label{eq:V43-P1-P2}
\end{equation}
Then \(L_{\rm eff}^{-,c,\Gamma}\) is singular exactly when
\begin{equation}
 F_r(s):=
 P_1P_2^2+s^2\left(
 b^2P_1\alpha^2+dbP_2\beta^2
 \right)=0.
 \label{eq:V43-F-exact}
\end{equation}
This identity remains valid when \(d=0\).
\end{lemma}

\begin{proof}
For the negative orientation,
\[
 L_{\rm eff}^{-,c,\Gamma}=L_0+CW_\Gamma J.
\]
The matrix determinant lemma and Lemma~\ref{lem:V43-response} give
\[
 \det L_{\rm eff}^{-,c,\Gamma}
 =\det L_0\,\det(I_3+R W_\Gamma).
\]
Direct expansion of the last three-by-three determinant yields
\eqref{eq:V43-F-exact}.  One can also multiply
\[
 \det\!\left(
 \operatorname{diag}(d^{-1},b^{-1},b^{-1})
 +a_cI_3-s\Gamma
 \right)
\]
by \(db^2\); the apparent inverse of \(d\) cancels and produces
\eqref{eq:V43-F-exact}.  The direct determinant is polynomial in \(d\),
so the formula also holds at \(d=0\).
\end{proof}

\subsection{The positive orientation always grows}
\label{subsec:V43-positive}

\begin{theorem}[Exact growing mode for every gyroscope]
\label{thm:V43-positive-growing}
Let \(0<c<1\) and let \(\Gamma\) be any real skew three-by-three matrix.
The positive orientation has a nonzero stable mode at
\begin{equation}
 k=0,\qquad s=\lambda=2.
 \label{eq:V43-positive-frequency}
\end{equation}
Hence no constant gyroscope makes the positive orientation pointwise
complementing in the open Laplace half-plane.
\end{theorem}

\begin{proof}
Odd-dimensional skew matrices are singular.  Choose
\(0\ne\zeta=(\zeta_1,\zeta_2,\zeta_3)^T\in\ker\Gamma\), and set
\begin{equation}
 x_\zeta:=(-3\zeta_3,0,\zeta_1,\zeta_2,\zeta_3)^T,
 \qquad
 y_\zeta:=2\zeta.
 \label{eq:V43-positive-vector}
\end{equation}
Then \(Jx_\zeta=\zeta\).  At \eqref{eq:V43-positive-frequency}, direct
substitution in the five V41 graph rows gives
\begin{equation}
 L_0x_\zeta=C(4\zeta).
 \label{eq:V43-positive-L0}
\end{equation}
The port force is
\[
 sy_\zeta-s\Gamma Jx_\zeta
 =4\zeta.
\]
Thus both blocks of
\({\mathfrak L}_{+,c,\Gamma}(x_\zeta,y_\zeta)\) vanish.  The mode has
\({\rm Re}\,s=2>0\) and normal decay rate \(\lambda=2\).
\end{proof}

\subsection{The negative orientation always has surface spectrum}
\label{subsec:V43-negative}

Put
\begin{equation}
 L:=\sqrt{1-c^2}>0,
 \qquad
 D:=\frac{4L^2-1}{8L^3}
 =\frac{3-4c^2}{8(1-c^2)^{3/2}}.
 \label{eq:V43-L-D}
\end{equation}
For real \(w\) in a fixed bounded interval and sufficiently large \(r\),
define
\begin{equation}
 \omega_r(w):=\sqrt{c^2r^2-rw},
 \qquad
 \lambda_r(w):=\sqrt{L^2r^2+rw},
 \qquad
 s_r(w):=i\omega_r(w).
 \label{eq:V43-scaled-frequency}
\end{equation}
Then \(a_c(s_r(w),r)=rw\).  If
\[
 \ell_r(w):=\frac{\lambda_r(w)}r
 =\sqrt{L^2+\frac wr},
 \qquad
 D_r(w):=\frac{4\ell_r(w)^2-1}{8\ell_r(w)^3},
\]
the exact response coefficients are
\begin{equation}
 b=\frac1{2r\ell_r},
 \qquad
 d=\frac{D_r}{r},
 \qquad
 P_1=1+wD_r,
 \qquad
 P_2=1+\frac{w}{2\ell_r}.
 \label{eq:V43-scaled-response}
\end{equation}

\begin{lemma}[Uniform high-frequency determinant]
\label{lem:V43-uniform-limit}
On every compact real \(w\)-interval,
\begin{equation}
 F_r(s_r(w))\longrightarrow Q_{c,\Gamma}(w)
 \label{eq:V43-uniform-convergence}
\end{equation}
uniformly as \(r\to\infty\), where
\begin{align}
 Q_{c,\Gamma}(w)
 :=\;&(1+Dw)\left(1+\frac{w}{2L}\right)^2
 \nonumber\\
 &-\frac{c^2}{4L^2}(1+Dw)\alpha^2
 -\frac{c^2D}{2L}
 \left(1+\frac{w}{2L}\right)\beta^2.
 \label{eq:V43-Q}
\end{align}
\end{lemma}

\begin{proof}
Equations \eqref{eq:V43-scaled-frequency} and
\eqref{eq:V43-scaled-response} give
\[
 s_r(w)^2b^2
 =\frac{-c^2+w/r}{4\ell_r(w)^2},
 \qquad
 s_r(w)^2db
 =\frac{(-c^2+w/r)D_r(w)}{2\ell_r(w)}.
\]
On a compact \(w\)-interval,
\(\ell_r\to L\) and \(D_r\to D\) uniformly.  Substitute these limits into
\eqref{eq:V43-F-exact}.  All denominators remain bounded away from zero
for sufficiently large \(r\), so the convergence is uniform.
\end{proof}

\begin{theorem}[High-frequency surface mode for every gyroscope]
\label{thm:V43-negative-surface}
Let \(0<c<1\) and let \(\Gamma\) be any real constant skew
three-by-three matrix.  For the negative orientation there are arbitrarily
large \(r\), purely imaginary frequencies \(s_r\), and stable roots
\(\lambda_r>0\) such that
\begin{equation}
 \det L_{\rm eff}^{-,c,\Gamma}(s_r,(r,0))=0.
 \label{eq:V43-negative-zero}
\end{equation}
After a spatial rotation when required, these give nonzero normally
decaying surface modes.  Consequently the negative orientation has no
uniform closed-half-plane Lopatinski floor.
\end{theorem}

\begin{proof}
First suppose \(D\ne0\), equivalently \(c^2\ne3/4\).  The real polynomial
\(Q_{c,\Gamma}\) in \eqref{eq:V43-Q} is cubic, with leading coefficient
\[
 \frac{D}{4L^2}\ne0.
\]
Choose \(W>0\) so large that
\(Q_{c,\Gamma}(-W)\) and \(Q_{c,\Gamma}(W)\) have opposite signs.
Lemma~\ref{lem:V43-uniform-limit} gives the same strict sign difference
for \(F_r(s_r(-W))\) and \(F_r(s_r(W))\) when \(r\) is large.  The
intermediate value theorem supplies
\(w_r\in(-W,W)\) with \(F_r(s_r(w_r))=0\).
For such \(r\), both radicands in
\eqref{eq:V43-scaled-frequency} are positive.  Hence \(s_r(w_r)\) is
purely imaginary and \(\lambda_r(w_r)>0\).  Lemmas
\ref{lem:V43-response} and \ref{lem:V43-determinant} turn this determinant
zero into a nonzero stable mode.  Since \(r\) can be arbitrarily large,
no uniform floor exists.

It remains to treat \(D=0\), namely \(c^2=3/4\) and \(L=1/2\).  Regard the
spatial--scalar entries
\((\gamma_{L\phi},\gamma_{T\phi})\) as the components of one spatial
covector.  If that covector is nonzero, choose the spatial frequency
direction so that
\(\alpha=\gamma_{T\phi}\ne0\).  The limiting polynomial becomes
\begin{equation}
 Q_{c,\Gamma}(w)
 =\left(1+\frac{w}{2L}\right)^2
 -\frac{c^2\alpha^2}{4L^2}.
 \label{eq:V43-resonant-Q}
\end{equation}
It has two distinct real roots
\[
 w=-2L\mathbin{\pm}c|\alpha|.
\]
Each root has a small real interval on which \(Q_{c,\Gamma}\) changes
sign.  Uniform convergence transfers that sign change to \(F_r\) for
large \(r\), and the preceding argument applies.

Finally suppose the spatial--scalar covector vanishes.  Then
\(\gamma_{L\phi}=\gamma_{T\phi}=0\), so \(\alpha=0\), and the exact
determinant factors as
\begin{equation}
 F_r(s)
 =P_2\left(P_1P_2+s^2db\,\gamma_{LT}^2\right).
 \label{eq:V43-resonant-factor}
\end{equation}
The exact V42 surface frequency satisfies \(P_2=1+a_c/(2\lambda)=0\)
for every \(r>0\).  It therefore remains a determinant zero even when
\(\gamma_{LT}\ne0\).  This completes the resonant case and the theorem.
\end{proof}

\begin{corollary}[Constant local gyroscopes cannot close V42]
\label{cor:V43-gyro-no-go}
For every \(0<c<1\) and every frequency-independent real skew
\(\Gamma\), neither Green orientation of
\eqref{eq:V43-coupled-row} gives a uniform stable inverse.  The positive
orientation has the exact growing mode of
Theorem~\ref{thm:V43-positive-growing}; the negative orientation has the
surface spectrum of Theorem~\ref{thm:V43-negative-surface}.
\end{corollary}

\begin{firewall}[Exact scope of the no-go theorem]
\label{fire:V43-scope}
Corollary~\ref{cor:V43-gyro-no-go} closes the complete constant quadratic
gyroscopic class with one time derivative and the V42 positive
Hamiltonian.  It does not close frequency-dependent positive-real
impedances, auxiliary fields with spatially dispersive couplings, or
infinite conservative baths.  Those are distinct operator classes and
must be tested with their own Green and causal estimates.
\end{firewall}

\subsection{Finite certificate and the V44 upstream row}
\label{subsec:V43-certificate}

A direct numerical audit used the exact five-by-five \(L_0\), the
eight-by-eight operator \eqref{eq:V43-full-operator}, and random real skew
matrices.  The largest residuals were
\begin{equation}
 \|JL_0^{-1}C-R\|=2.1\times10^{-16},
 \quad
 |\det(I+RW_\Gamma)-F_r|=2.9\times10^{-15},
 \quad
 r_{\rm Schur}=6.8\times10^{-15}.
 \label{eq:V43-numerical-residuals}
\end{equation}
Fifty random kernel witnesses for
Theorem~\ref{thm:V43-positive-growing} had residual below
\(9.3\times10^{-16}\).  At \(c=1/2\), \(r=2000\), three unrelated
gyroscopes produced real \(w\)-root brackets
\[
 [-3.00,-2.95],\qquad[-3.35,-3.30],\qquad[-2.70,-2.65].
\]
The resonant simple-root and exact-factor cases were also checked.  These
computations audit the formulas; the proof of
Theorem~\ref{thm:V43-negative-surface} is the uniform-limit and
intermediate-value argument above.

V44 must now go upstream from finite constant gyroscopes.  The next
candidate is a causal retarded impedance
\begin{equation}
 f_{\rm ret}(s,k)=Z(s,k)z,
 \qquad
 {\rm Re}\,\langle z,Z(s,k)z\rangle\ge0
 \quad({\rm Re}\,s>0),
 \label{eq:V43-positive-real-target}
\end{equation}
derived from a conservative dilation rather than postulated as a fitted
selector.  The required V44 gates are: an explicit spectral measure and
bath action; a proof that retarded elimination gives an analytic
positive-real \(Z\); exact DeWitt/BRST source incidence; and a
three-channel stable determinant with a nonzero conic floor.  If a scalar
bath cannot meet the last gate, the spectral measure must be
matrix-valued before returning to nonlinear metric stress and the
Standard-Model sectors.

\subsection{V43 ledger}
\label{subsec:V43-ledger}

\paragraph{New conclusions proved in V43.}
Proposition~\ref{prop:V43-gyro-complete} identifies the entire constant
first-time-order energy-preserving gyroscopic class.
Proposition~\ref{prop:V43-Euler-Green} proves its Euler, power, and Green
identities.  Lemma~\ref{lem:V43-response} derives the exact three-port
bulk response, and Lemma~\ref{lem:V43-determinant} reduces the negative
orientation to one scalar determinant.
Theorems~\ref{thm:V43-positive-growing} and
\ref{thm:V43-negative-surface} prove failure for both orientations and
include the resonant speed \(c^2=3/4\).

\paragraph{Imported conclusions not repeated.}
V41 supplies \(H,J,C\) and the reciprocal graph.  V42 supplies the
minimal Hamiltonian port, the ungyroscoped surface branch, and the
eight-row Schur structure.  V43 rederives only the response needed to
treat arbitrary \(\Gamma\); it does not repeat the earlier gauge quotient
or channel lower bound.

\paragraph{Honest continuum status.}
The finite constant local gyroscopic repair is now ruled out rigorously.
A retarded positive-real impedance with a conservative ancestry has not
yet been constructed, and no global stable floor, nonlinear covariant
stress/Bianchi identity, ghost completion, or coupled Standard-Model
continuum is claimed.  The full Standard Model--Einstein continuum
therefore remains open.

% END APPENDED VERSION V43 MATERIAL



% BEGIN APPENDED VERSION V44 MATERIAL

\clearpage
\section{V44: conservative half-line bath, positive-real impedance, and
the complete scalar-bath sign obstruction}
\label{sec:V44-scalar-bath}

\subsection{Purpose and upstream change}
\label{subsec:V44-purpose}

V43 ruled out every frequency-independent constant gyroscope attached to
the three metric defect coordinates.  Its handoff required a retarded
impedance with a conservative ancestry.  V44 constructs the simplest such
object: a three-component wave field on an auxiliary half-line whose
boundary trace is \(z=Jx\).  Retarded elimination gives the homogeneous
square-root impedance
\[
 Z_{\eta,v}(s,k)
 =\eta\sqrt{s^2+v^2|k|^2}\,I_3.
\]
This map is analytic and positive-real in the open Laplace half-plane, has
a nonnegative spectral measure, and comes from a positive local bath
energy.

Those properties are necessary but not sufficient.  The exact V43 bulk
response shows that the two Green orientations again fail for different
reasons.  The positive orientation always has a transverse
closed-half-plane zero or an open growing root.  The negative orientation
always has a longitudinal open root or a normally decaying surface mode.
V44 proves this for every coupling \(\eta>0\) and every strictly
subluminal bath speed \(0<v<1\).

Thus the scalar bath is not adopted as the final boundary condition.  Its
construction and failure identify the next required change: the bath
spectral measure must be matrix-valued and adapted to the longitudinal
Einstein response.

\subsection{A local conservative dilation}
\label{subsec:V44-dilation}

Let \(a\in[0,\infty)\) be an auxiliary bath coordinate and let
\[
 q=q(a,t,y)\in\mathbb R^3.
\]
Restrict the configuration space by the trace constraint
\begin{equation}
 q(0,t,y)=z(t,y)=Jx(t,y).
 \label{eq:V44-trace}
\end{equation}
For constants \(\eta>0\) and \(0<v<1\), define
\begin{equation}
 S_{\rm bath}[q]
 :=\frac{\eta}{2}\int_{\partial M}\int_0^\infty
 \left(
 |D_tq|^2-|\partial_aq|^2-v^2|D_iq|^2
 \right)\dd a\,\dd^3y.
 \label{eq:V44-bath-action}
\end{equation}

\begin{proposition}[Bath equation, traction, and positive energy]
\label{prop:V44-bath-variation}
The Euler equation in the bath interior is
\begin{equation}
 -D_t^2q+\partial_a^2q+v^2\Delta_yq=0.
 \label{eq:V44-bath-wave}
\end{equation}
The boundary variation at \(a=0\) is
\begin{equation}
 \eta\,\partial_aq(0)^T\delta z.
 \label{eq:V44-boundary-variation}
\end{equation}
Consequently the force exerted on the metric port is
\begin{equation}
 f_{\rm bath}:=-\eta\,\partial_aq(0).
 \label{eq:V44-bath-force}
\end{equation}
The bath energy
\begin{equation}
 E_{\rm bath}
 :=\frac{\eta}{2}\int_{\mathbb R^2}\int_0^\infty
 \left(
 |D_tq|^2+|\partial_aq|^2+v^2|D_iq|^2
 \right)\dd a\,\dd^2y
 \label{eq:V44-bath-energy}
\end{equation}
is nonnegative, and its boundary power is
\begin{equation}
 \frac{\dd}{\dd t}E_{\rm bath}
 =\int_{\mathbb R^2}(D_tz)^Tf_{\rm bath}\,\dd^2y
 \label{eq:V44-bath-power}
\end{equation}
when the field decays at \(a=\infty\).
\end{proposition}

\begin{proof}
Vary \eqref{eq:V44-bath-action} and integrate once in \(t,y,a\).
The \(a\)-integration gives
\[
 \left[-\eta\,\partial_aq^T\delta q\right]_{a=0}^{a=\infty}
 =\eta\,\partial_aq(0)^T\delta z,
\]
which proves \eqref{eq:V44-bath-wave}--\eqref{eq:V44-bath-force}.
Differentiate \eqref{eq:V44-bath-energy}, use
\eqref{eq:V44-bath-wave}, and integrate the spatial terms.  The only
remaining term is
\[
 -\eta\int(D_tq(0))^T\partial_aq(0)\,\dd^2y
 =\int(D_tz)^Tf_{\rm bath}\,\dd^2y.
\]
This proves the power identity and positivity is immediate from
\eqref{eq:V44-bath-energy}.
\end{proof}

Couple the bath to the DeWitt graph by
\begin{equation}
 H(p-A(D)x)=\varepsilon J^Tf_{\rm bath},
 \qquad \varepsilon\in\{+1,-1\}.
 \label{eq:V44-coupled-graph}
\end{equation}
Because the trace constraint uses the full-embedding invariant variable
\(z=Jx\), variation of the single action produces exactly the incidence
\(J^Tf_{\rm bath}\).  The frozen BRST/Noether identity of V39 therefore
holds without a fitted selector.

\begin{proposition}[Conservative Green ancestry]
\label{prop:V44-Green-ancestry}
For two bath fields \(q,q'\), the boundary traction pairing is
\begin{equation}
 f_{\rm bath}^Tz'-z^Tf_{\rm bath}'
 =-\eta\left[
 (\partial_aq)^Tq'-q^T\partial_aq'
 \right]_{a=0}.
 \label{eq:V44-bath-Green}
\end{equation}
It is the \(a=0\) flux of the formally self-adjoint wave operator
\eqref{eq:V44-bath-wave}.  With the opposite induced Stokes orientation,
it cancels the shifted DeWitt radial Green form in
\eqref{eq:V44-coupled-graph}.
\end{proposition}

\begin{proof}
Substitute \eqref{eq:V44-trace} and
\eqref{eq:V44-bath-force}.  The right side is precisely the boundary term
obtained by subtracting the two integrated bath wave pairings.  The
metric contribution is
\(\varepsilon(f_{\rm bath}^Tz'-z^Tf_{\rm bath}')\), as in V42.
Assigning the bath boundary the opposite Stokes orientation gives exact
cancellation.
\end{proof}

\subsection{Retarded elimination and the spectral measure}
\label{subsec:V44-retarded}

For a Laplace--Fourier mode, put
\begin{equation}
 \kappa_v(s,k):=\sqrt{s^2+v^2|k|^2},
 \qquad {\rm Re}\,\kappa_v>0
 \quad\hbox{when }{\rm Re}\,s>0.
 \label{eq:V44-kappa}
\end{equation}
The unique bath solution that decays as \(a\to\infty\) is
\begin{equation}
 \widehat q(a;s,k)=e^{-a\kappa_v(s,k)}\widehat z(s,k).
 \label{eq:V44-retarded-solution}
\end{equation}
Therefore
\begin{equation}
 \widehat f_{\rm bath}
 =Z_{\eta,v}(s,k)\widehat z,
 \qquad
 Z_{\eta,v}(s,k):=\eta\kappa_v(s,k)I_3.
 \label{eq:V44-impedance}
\end{equation}

\begin{theorem}[Analytic positive-real impedance with positive measure]
\label{thm:V44-positive-real}
The impedance \eqref{eq:V44-impedance} is analytic for
\({\rm Re}\,s>0\) and satisfies
\begin{equation}
 {\rm Re}\,\langle u,Z_{\eta,v}(s,k)u\rangle
 =\eta\,{\rm Re}\,\kappa_v(s,k)|u|^2>0
 \label{eq:V44-positive-real}
\end{equation}
for \(u\ne0\).  If \(m=v|k|\), then
\begin{equation}
 \kappa_v(s,k)
 =s+\frac{2s}{\pi}\int_0^m
 \frac{\sqrt{m^2-\omega^2}}{s^2+\omega^2}\,\dd\omega,
 \qquad {\rm Re}\,s>0.
 \label{eq:V44-spectral-representation}
\end{equation}
Thus its oscillator spectral density
\begin{equation}
 \rho_m(\omega)
 :=\frac2\pi\sqrt{m^2-\omega^2}\,
 {\bf 1}_{[0,m]}(\omega)
 \label{eq:V44-spectral-density}
\end{equation}
is nonnegative.
\end{theorem}

\begin{proof}
If \(s=x+iy\) with \(x>0\), then \(s^2+m^2\) is never on the closed
negative real axis: its imaginary part is \(2xy\) when \(y\ne0\), and it
is positive when \(y=0\).  The principal square root is consequently
analytic and has strictly positive real part.  This proves
\eqref{eq:V44-positive-real}.

For real \(s>0\), the substitution \(\omega=m\sin\theta\), followed by
one elementary rational integral, gives
\[
 \int_0^m\frac{\sqrt{m^2-\omega^2}}{s^2+\omega^2}\,\dd\omega
 =\frac{\pi}{2}\left(\frac{\sqrt{s^2+m^2}}s-1\right).
\]
Rearrangement proves \eqref{eq:V44-spectral-representation} on the
positive real axis.  Both sides are analytic in the right half-plane, so
the identity theorem proves it there.  Formula
\eqref{eq:V44-spectral-density} is visibly nonnegative.
\end{proof}

The extended bath system is conservative.  Selecting
\eqref{eq:V44-retarded-solution} is a causal radiation condition, and its
projection onto \(z\) is passive.  This is the conservative-dilation
versus retarded-projection distinction required in V31 and V43.

\subsection{Exact three-channel determinant}
\label{subsec:V44-determinant}

Use the V43 response map for \(k=(r,0)\):
\[
 R=\operatorname{diag}(d,b,b),\qquad
 b=\frac1{2\lambda},\qquad
 d=\frac{4\lambda^2-r^2}{8\lambda^3},\qquad
 \lambda^2=s^2+r^2.
\]
After retarded bath elimination, the metric stable matrix is
\begin{equation}
 L_{\rm bath}^{\varepsilon}
 :=L_0-\varepsilon C Z_{\eta,v}J.
 \label{eq:V44-effective}
\end{equation}

\begin{lemma}[Longitudinal--transverse factorization]
\label{lem:V44-factorization}
Whenever \(\lambda\ne0\),
\begin{equation}
 \frac{\det L_{\rm bath}^{\varepsilon}}{\det L_0}
 ={\cal L}_\varepsilon\,{\cal T}_\varepsilon^2,
 \label{eq:V44-factorization}
\end{equation}
where
\begin{equation}
 {\cal L}_\varepsilon
 :=1-\varepsilon\eta\kappa_vd,
 \qquad
 {\cal T}_\varepsilon
 :=1-\varepsilon\frac{\eta\kappa_v}{2\lambda}.
 \label{eq:V44-L-T}
\end{equation}
\end{lemma}

\begin{proof}
The matrix determinant lemma gives
\[
 \det L_{\rm bath}^{\varepsilon}
 =\det L_0\det(I_3-\varepsilon RZ_{\eta,v}).
\]
Both \(R\) and \(Z_{\eta,v}\) are diagonal, so the last determinant is
exactly \eqref{eq:V44-factorization}.
\end{proof}

For \(r>0\), homogeneity permits \(r=1\).  Introduce
\begin{equation}
 Y:=\frac{\lambda^2}{r^2},
 \qquad
 \delta:=1-v^2\in(0,1).
 \label{eq:V44-Y-delta}
\end{equation}
On the stable branches,
\begin{equation}
 \frac{\kappa_v^2}{r^2}=Y-\delta.
 \label{eq:V44-kappa-Y}
\end{equation}

\subsection{The positive orientation: an explicit transverse zero}
\label{subsec:V44-positive}

\begin{theorem}[Positive-orientation scalar-bath obstruction]
\label{thm:V44-positive-obstruction}
For every \(\eta>0\) and \(0<v<1\), the positive orientation
\(\varepsilon=+1\) fails the stable determinant test.

More precisely:
\begin{enumerate}
 \item if \(0<\eta<2\), \({\cal T}_+=0\) at a purely imaginary
 closed-half-plane frequency with \(\lambda\ne0\);
 \item if \(\eta=2\), \({\cal T}_+=0\) at \(r=0\) for every
 \({\rm Re}\,s>0\);
 \item if \(\eta>2\) and \(\eta v<2\), \({\cal T}_+=0\) at an open
 real growing frequency;
 \item if \(\eta>2\) and \(\eta v\ge2\), \({\cal T}_+=0\) at a purely
 imaginary frequency, with positive normal decay when \(\eta v>2\).
\end{enumerate}
Thus this orientation never has a nonzero uniform conic floor.
\end{theorem}

\begin{proof}
The equation \({\cal T}_+=0\) is
\begin{equation}
 \frac{\kappa_v}{\lambda}=\frac2\eta.
 \label{eq:V44-positive-T-zero}
\end{equation}
If \(0<\eta<2\), set
\begin{equation}
 Y_-:=-\frac{\eta^2\delta}{4-\eta^2}<0.
 \label{eq:V44-Y-minus}
\end{equation}
Take
\[
 s=ir\sqrt{1-Y_-},\qquad
 \lambda=ir\sqrt{-Y_-},\qquad
 \kappa_v=ir\sqrt{\delta-Y_-},
\]
as boundary values of the stable branches.  Squaring
\eqref{eq:V44-positive-T-zero} gives
\(\eta^2(Y-\delta)=4Y\), which is exactly
\eqref{eq:V44-Y-minus}; the displayed branches have the same positive
imaginary orientation, so the unsquared equation also holds.

If \(\eta=2\), choose \(r=0\).  Then
\(\kappa_v=\lambda=s\) in the open half-plane and
\({\cal T}_+=1-\eta/2=0\).

Let \(\eta>2\) and put
\begin{equation}
 Y_+:=\frac{\eta^2\delta}{\eta^2-4}>\delta.
 \label{eq:V44-Y-plus}
\end{equation}
The positive real square roots of \(Y_+\) and \(Y_+-\delta\) satisfy
\eqref{eq:V44-positive-T-zero}.  Moreover,
\[
 Y_+>1,\quad Y_+=1,\quad Y_+<1
\]
according as
\[
 \eta v<2,\quad \eta v=2,\quad \eta v>2.
\]
In the first case take \(s=r\sqrt{Y_+-1}>0\), giving an open growing
stable mode.  In the second take \(s=0\).  In the third take
\(s=ir\sqrt{1-Y_+}\); here
\(\lambda=r\sqrt{Y_+}>0\), so the determinant zero is a normally decaying
surface mode.  Lemma~\ref{lem:V44-factorization} proves every assertion.
\end{proof}

\subsection{The negative orientation: longitudinal cubic dichotomy}
\label{subsec:V44-negative}

The negative longitudinal equation \({\cal L}_-=0\) is
\begin{equation}
 8Y^{3/2}
 +\eta(Y-\delta)^{1/2}(4Y-1)=0,
 \label{eq:V44-negative-branch}
\end{equation}
with square roots taken on the stable branches.  Squaring gives
\begin{equation}
 P_{\eta,\delta}(Y)
 :=64Y^3-\eta^2(Y-\delta)(4Y-1)^2=0.
 \label{eq:V44-cubic}
\end{equation}

\begin{lemma}[Exact discriminant and branch assignment]
\label{lem:V44-discriminant}
The discriminant of \eqref{eq:V44-cubic}, with its quadratic
interpretation when \(\eta=2\), is
\begin{equation}
 \Delta_{\eta,\delta}
 =256\eta^4\left[
 \eta^2(1-4\delta)^3-432\delta^2
 \right].
 \label{eq:V44-discriminant}
\end{equation}
Whenever \(\Delta_{\eta,\delta}<0\), the nonreal conjugate root pair of
\(P_{\eta,\delta}\) satisfies the unsquared negative-orientation equation
\eqref{eq:V44-negative-branch}.  It therefore produces open stable roots
\(s\) with \({\rm Re}\,s>0\).
\end{lemma}

\begin{proof}
Expanding \eqref{eq:V44-cubic} gives
\[
 16(4-\eta^2)Y^3
 +8\eta^2(1+2\delta)Y^2
 -\eta^2(1+8\delta)Y
 +\eta^2\delta.
\]
Substitution into the elementary cubic discriminant formula gives
\eqref{eq:V44-discriminant}.  At \(\eta=2\) the cubic coefficient
vanishes; the same expression is the discriminant of the remaining
finite quadratic pair.

At a root define
\[
 \sigma(Y)
 :=\frac{8Y^{3/2}}
 {\eta(Y-\delta)^{1/2}(4Y-1)}\in\{+1,-1\}.
\]
When \(\Delta<0\), the conjugate pair is simple.  Along that pair
\(\sigma\) is continuous and takes values in a discrete set, so it cannot
change sign without a collision.  For \(\eta\downarrow0\),
\[
 Y_\pm
 =\left(\frac{\eta^2\delta}{64}\right)^{1/3}
 e^{\pm i\pi/3}+o(\eta^{2/3}).
\]
On the stable square-root branches, direct substitution gives
\(\sigma(Y_\pm)=-1+o(1)\), hence exactly \(-1\) at the roots.
This is \eqref{eq:V44-negative-branch}.  The finite conjugate pair
continues through the degree drop at \(\eta=2\), so the assignment
persists throughout every \(\Delta<0\) component.

For nonreal \(Y\), choose
\[
 \lambda=r\sqrt Y,\qquad s=r\sqrt{Y-1}
\]
with positive real parts.  Equation \eqref{eq:V44-kappa-Y} then fixes the
stable \(\kappa_v\) branch.  Thus \({\rm Re}\,s>0\), and the determinant
factor \({\cal L}_-\) vanishes.
\end{proof}

\begin{theorem}[Negative-orientation scalar-bath obstruction]
\label{thm:V44-negative-obstruction}
For every \(\eta>0\) and \(0<v<1\), the negative orientation
\(\varepsilon=-1\) has either an open longitudinal growing root or a
purely imaginary normally decaying surface mode.  Hence it never has a
uniform stable inverse.
\end{theorem}

\begin{proof}
Recall \(\delta=1-v^2\).

If \(\delta\ge1/4\), then
\((1-4\delta)^3\le0\), so
\(\Delta_{\eta,\delta}<0\) for every \(\eta>0\).
Lemma~\ref{lem:V44-discriminant} gives an open stable root.

Suppose \(0<\delta<1/4\), and define
\begin{equation}
 \eta_*^2
 :=\frac{432\delta^2}{(1-4\delta)^3},
 \qquad
 Y_*:=\frac{3\delta}{2(1+2\delta)}.
 \label{eq:V44-threshold}
\end{equation}
For \(0<\eta<\eta_*\), the discriminant is negative and the same lemma
again gives an open root.

For \(\eta\ge\eta_*\), consider on
\(\delta<Y<1/4\) the positive function
\begin{equation}
 {\cal E}_\delta(Y)
 :=\frac{8Y^{3/2}}
 {\sqrt{Y-\delta}(1-4Y)}.
 \label{eq:V44-E-delta}
\end{equation}
It diverges at both endpoints.  Its logarithmic derivative vanishes
exactly when
\[
 \frac3Y-\frac1{Y-\delta}+\frac8{1-4Y}=0,
\]
whose unique solution is \(Y=Y_*\).  Direct evaluation gives
\[
 {\cal E}_\delta(Y_*)^2
 =\frac{432\delta^2}{(1-4\delta)^3}
 =\eta_*^2.
\]
Thus \({\cal E}_\delta(Y)=\eta\) has one double solution at
\(\eta=\eta_*\) and two solutions when \(\eta>\eta_*\).  Each solution
satisfies \eqref{eq:V44-negative-branch}, because \(4Y-1<0\).
Set
\[
 s=ir\sqrt{1-Y},\qquad
 \lambda=r\sqrt Y>0,\qquad
 \kappa_v=r\sqrt{Y-\delta}>0.
\]
This is a purely imaginary, normally decaying longitudinal surface mode.
All parameter ranges are exhausted.
\end{proof}

\begin{corollary}[Passivity alone does not close the Einstein port]
\label{cor:V44-scalar-no-go}
The scalar impedance \(Z_{\eta,v}=\eta\kappa_vI_3\) has a positive local
conservative dilation, a nonnegative spectral measure, and a strictly
positive-real retarded map.  Nevertheless neither Green orientation has
a uniform stable inverse, for any \(\eta>0\) and \(0<v<1\).
\end{corollary}

\begin{proof}
Combine Theorems~\ref{thm:V44-positive-real},
\ref{thm:V44-positive-obstruction}, and
\ref{thm:V44-negative-obstruction}.
\end{proof}

\begin{firewall}[What the scalar no-go proves]
\label{fire:V44-scope}
Corollary~\ref{cor:V44-scalar-no-go} does not say that passive retarded
closure is impossible.  It proves that an isotropic scalar spectral
measure cannot accommodate the unequal longitudinal and transverse
Einstein response eigenvalues \(d,b,b\).  A matrix-valued measure may
alter that eigenvalue matching while retaining positivity and
conservative ancestry; it requires a new determinant proof.
\end{firewall}

\subsection{Finite certificate and V45 companion-audit handoff}
\label{subsec:V44-certificate}

A direct audit compared the five-by-five determinant of
\eqref{eq:V44-effective} with
\eqref{eq:V44-factorization} at one hundred random complex frequencies,
both orientations, and random \((v,\eta)\).  The largest ratio residual
was \(2.4\times10^{-14}\).  Five-hundred-point Gauss quadrature checked
\eqref{eq:V44-spectral-representation} to \(5.1\times10^{-10}\).

For the negative orientation, representative open roots were
\[
 0.0182643+0.9864604i,\qquad
 0.0458812+0.9358102i.
\]
For \(\delta=0.1\), the exact threshold is
\(\eta_*=\sqrt{20}=4.472135955\) and \(Y_*=1/8\).
At \(\eta=1.5\eta_*\), the two surface roots were bracketed at
\[
 Y=0.1048665809,\qquad Y=0.1708940327,
\]
with unsquared residual below \(2.3\times10^{-16}\).
These calculations audit the formulas and do not replace the proofs.

V45 is a scheduled five-version synchronization point.  Before choosing
the matrix-valued bath, it must inspect the newest retained
Navier--Stokes and Yang--Mills notes and import only exact companion
results that change the source-incidence, positivity, or contraction
requirements.  The post-audit construction should then specify a
positive semidefinite matrix spectral density
\(\rho(\omega,k)\), derive its retarded \(Z(s,k)\) from a local
conservative bath, and test
\[
 \det(I_3-\varepsilon R(s,k)Z(s,k))
\]
on the entire normalized closed stable set.  Nonlinear metric stress,
Bianchi closure, ghosts, and the Standard-Model sectors remain behind
that causal boundary gate.

\subsection{V44 ledger}
\label{subsec:V44-ledger}

\paragraph{New conclusions proved in V44.}
Propositions~\ref{prop:V44-bath-variation} and
\ref{prop:V44-Green-ancestry} construct the local positive-energy
conservative dilation and its exact DeWitt Green incidence.
Theorem~\ref{thm:V44-positive-real} proves analyticity, strict
positive-realness, and the positive spectral representation.
Lemma~\ref{lem:V44-factorization} gives the exact three-channel stable
determinant.  Theorems~\ref{thm:V44-positive-obstruction} and
\ref{thm:V44-negative-obstruction} prove the complete scalar-bath
orientation obstruction.

\paragraph{Imported conclusions not repeated.}
V39 supplies frozen full-embedding invariance and source incidence.
V41 supplies \(H,J,C\).  V43 supplies the diagonal bulk response
\(R=\operatorname{diag}(d,b,b)\).  V44 uses those results without
repeating their quotient or response derivations.

\paragraph{Honest continuum status.}
A causal positive-real impedance with a local conservative ancestry has
now been constructed, but the scalar spectral measure fails the global
stable determinant.  A matrix-valued bath has not yet passed positivity,
Green, and conic-floor gates.  No nonlinear covariant stress/Bianchi
identity, ghost completion, or coupled Standard-Model continuum is
claimed.  The full Standard Model--Einstein continuum remains open.

% END APPENDED VERSION V44 MATERIAL



% BEGIN APPENDED VERSION V45 MATERIAL

\clearpage
\section{V45: scheduled NS/YM synchronization, complete matrix-bath
incidence, and the surviving noncommuting branch}
\label{sec:V45-matrix-bath}

\subsection{Frozen companion audit before direction choice}
\label{subsec:V45-audit}

This is the scheduled five-version synchronization required after V40.
The companion files were read before choosing the V45 matrix ansatz.  The
finalized states used in this audit are
\begin{align*}
 &\texttt{Renewal\_NS\_Stability\_Research\_Notes\_V17.tex},\\
 &\qquad\texttt{SHA256}=
 \texttt{c9717f98a2d7b8be59846def1097b8b92bde27230aac8f7a53d947125870d1ae},
 \\
 &\texttt{Renewal\_Yang\_Mills\_Mass\_Gap\_Research\_Notes\_V25.tex},\\
 &\qquad\texttt{SHA256}=
 \texttt{e85faf51d07b7b4364e46204fe9b5eb05a0ffcb6eb778abcde10ff07b944baa8}.
\end{align*}
The hashes freeze the exact companion states.  No later companion result
is silently imported.

The NS V15--V17 appendices prove three type-compatible design lessons.
First, a signed source must be paired with the contemporaneous carrier at
the same output frequency before taking positive parts.  Second, a
fixed-scale or negative-order response has no uniform principal floor.
Third, a full source-regularity ladder and a separately paid occupation
ledger can both be exact while failing to meet on the same record
history.  These are not Einstein estimates.  For V45 they require the
bath response, source incidence, and principal determinant to be computed
on one complete three-port vector before channelwise inequalities.

The YM V25 appendix supplies a second algebraic lesson.  Its exact
source-deficit formula retains the full cross term through a Schur
completion; its spectral and weighted-locality pencils must be populated
on the same complete carrier; and loss of either uniform margin returns a
specific witness rather than permission to weaken the norm.  Again, no YM
gap estimate is imported.  V45 reproves the corresponding matrix
identities for the Einstein port.

\begin{firewall}[Companion type boundary]
\label{fire:V45-audit-scope}
The companion notes change the assembly order and the failure ledger.
They do not prove positivity, causality, or a stable estimate for the
Einstein symbol.  Every SM statement below follows from the V41--V44
DeWitt response and a new direct matrix calculation.
\end{firewall}

\subsection{The complete common-speed matrix bath}
\label{subsec:V45-matrix-action}

Let
\begin{equation}
 M=M^T\succeq0,\qquad M\in{\rm Mat}_{3\times3}(\mathbb R),
 \label{eq:V45-M-positive}
\end{equation}
and replace the scalar internal metric in V44 by \(M\):
\begin{equation}
 S_{{\rm bath},M}[q]
 :=\frac12\int_{\partial M}\int_0^\infty
 \left[
 (D_tq)^TM(D_tq)-(\partial_aq)^TM(\partial_aq)
 -v^2(D_iq)^TM(D_iq)
 \right]\dd a\,\dd^3y,
 \label{eq:V45-matrix-action}
\end{equation}
with \(q(0)=z=Jx\) and \(0<v<1\).

\begin{proposition}[Positive matrix measure and exact incidence]
\label{prop:V45-matrix-bath}
The action \eqref{eq:V45-matrix-action} has nonnegative energy
\begin{equation}
 E_{{\rm bath},M}
 =\frac12\int\!\!\int
 \left[
 (D_tq)^TM(D_tq)+(\partial_aq)^TM(\partial_aq)
 +v^2(D_iq)^TM(D_iq)
 \right]\dd a\,\dd^2y.
 \label{eq:V45-matrix-energy}
\end{equation}
Retarded elimination gives
\begin{equation}
 f_M=Z_M(s,k)z,
 \qquad
 Z_M(s,k)=\kappa_v(s,k)M.
 \label{eq:V45-matrix-impedance}
\end{equation}
Its spectral density is
\begin{equation}
 \rho_{m,M}(\omega)
 =\frac2\pi\sqrt{m^2-\omega^2}\,
 {\bf1}_{[0,m]}(\omega)M\succeq0,
 \qquad m=v|k|.
 \label{eq:V45-matrix-density}
\end{equation}
For \({\rm Re}\,s>0\),
\begin{equation}
 {\rm Re}\,\langle u,Z_M(s,k)u\rangle
 =({\rm Re}\,\kappa_v)u^TMu\ge0.
 \label{eq:V45-matrix-positive-real}
\end{equation}
The metric variation is exactly
\(J^Tf_M\), and the conservative bath Green flux is the V44 flux with
the Euclidean pairing replaced by the \(M\)-pairing.
\end{proposition}

\begin{proof}
The V44 integration by parts applies componentwise with the constant
symmetric matrix \(M\).  Its boundary traction is
\(-M\partial_aq(0)\).  The retarded solution remains
\(q(a)=e^{-a\kappa_v}z\), proving
\eqref{eq:V45-matrix-impedance}.  Multiply the scalar spectral
representation of V44 by \(M\) to obtain
\eqref{eq:V45-matrix-density}.  Positivity of the quadratic energy and
\eqref{eq:V45-matrix-positive-real} follow from \(M\succeq0\).
Finally, the trace constraint \(q(0)=Jx\) makes the boundary variation
\(J^Tf_M\); no component is inserted after variation.
\end{proof}

The common bath history is kept as a three-vector throughout.  Diagonal
bounds on its components are not taken before the determinant.

\subsection{Full rank is forced at gravitational glancing}
\label{subsec:V45-rank}

For either Green orientation, the retarded effective matrix is
\begin{equation}
 L_M^\varepsilon
 :=L_0-\varepsilon C\kappa_vMJ.
 \label{eq:V45-effective}
\end{equation}

\begin{theorem}[Rank-deficient matrix measures cannot repair glancing]
\label{thm:V45-rank-three}
If \(\operatorname{rank}M<3\), then
\(L_M^\varepsilon\) is singular at every nonzero gravitational glancing
frequency, for both \(\varepsilon=\pm1\).  Hence a common-speed matrix
bath that repairs the complete V37 defect must have
\begin{equation}
 M\succ0.
 \label{eq:V45-M-definite}
\end{equation}
\end{theorem}

\begin{proof}
At gravitational glancing, V37 gives
\(\dim\ker L_0=3\).  The update
\(C\kappa_vMJ\) has rank at most \(\operatorname{rank}M<3\).
Its restriction to the three-dimensional space \(\ker L_0\) therefore
has a nonzero kernel vector \(x\).  For that vector,
\[
 L_0x=0,\qquad C\kappa_vMJx=0,
\]
so \(L_M^\varepsilon x=0\).  A positive semidefinite three-by-three
matrix of rank three is positive definite, proving the last assertion.
\end{proof}

This is the SM version of complete-source incidence: three scalar
margins do not compensate for a missing port direction.

\subsection{Exact full-Gram determinant and Schur completion}
\label{subsec:V45-Schur}

Fix \(k=(r,0)\).  In the port basis
\((v_L;v_T,\phi)\), partition
\begin{equation}
 M=
 \begin{pmatrix}
 A&u^T\cr
 u&N
 \end{pmatrix},
 \qquad
 A\in\mathbb R,\quad u\in\mathbb R^2,\quad
 N=N^T\in{\rm Mat}_{2\times2}(\mathbb R).
 \label{eq:V45-M-block}
\end{equation}
For \(M\succ0\),
\begin{equation}
 N\succ0,\qquad
 m_{\rm sc}:=A-u^TN^{-1}u
 =\frac{\det M}{\det N}>0.
 \label{eq:V45-static-Schur}
\end{equation}
Use the V43 response and put
\begin{equation}
 \xi:=\kappa_vd,\qquad
 \theta:=\kappa_vb.
 \label{eq:V45-xi-theta}
\end{equation}
Define the complete Gram invariants
\begin{equation}
 B:=\operatorname{tr}N,\qquad
 C_*:=A\operatorname{tr}N-|u|^2,\qquad
 D:=\det N,\qquad E:=\det M.
 \label{eq:V45-invariants}
\end{equation}
Positive definiteness implies \(A,B,C_*,D,E>0\).

\begin{theorem}[Complete matrix determinant]
\label{thm:V45-complete-determinant}
For both orientations,
\begin{align}
 F_\varepsilon(s,k)
 &:=\frac{\det L_M^\varepsilon}{\det L_0}
 =\det\!\left(I_3-\varepsilon
 \operatorname{diag}(\xi,\theta,\theta)M\right)
 \nonumber\\
 &=1-\varepsilon(A\xi+B\theta)
 +(C_*\xi\theta+D\theta^2)
 -\varepsilon E\xi\theta^2.
 \label{eq:V45-invariant-determinant}
\end{align}
Where \(I_2-\varepsilon\theta N\) is invertible, the same determinant is
\begin{equation}
 F_\varepsilon
 =\det(I_2-\varepsilon\theta N)
 \left[1-\varepsilon\xi\,m_\varepsilon(\theta)\right],
 \label{eq:V45-Schur-determinant}
\end{equation}
with the exact dynamic Schur coefficient
\begin{equation}
 m_\varepsilon(\theta)
 :=A+\varepsilon\theta
 u^T(I_2-\varepsilon\theta N)^{-1}u.
 \label{eq:V45-dynamic-Schur}
\end{equation}
The invariant polynomial \eqref{eq:V45-invariant-determinant} remains valid
at the apparent Schur poles.
\end{theorem}

\begin{proof}
The V43 matrix determinant lemma gives
\[
 \frac{\det L_M^\varepsilon}{\det L_0}
 =\det(I_3-\varepsilon RM\kappa_v).
\]
Insert \(R=\operatorname{diag}(d,b,b)\).  Expanding by principal minors
gives: the trace \(A\xi+B\theta\), the sum of two-by-two principal minors
\(C_*\xi\theta+D\theta^2\), and the determinant
\(E\xi\theta^2\).  This proves
\eqref{eq:V45-invariant-determinant}.

Alternatively, take the Schur complement of the transverse block
\(I_2-\varepsilon\theta N\).  Its scalar complement is
\[
 1-\varepsilon A\xi
 -\xi\theta u^T(I_2-\varepsilon\theta N)^{-1}u
 =1-\varepsilon\xi m_\varepsilon(\theta),
\]
proving \eqref{eq:V45-Schur-determinant}.  Multiplication by the
transverse determinant cancels every apparent pole and returns the
invariant polynomial.
\end{proof}

\begin{lemma}[Positive complete Schur limit]
\label{lem:V45-positive-Schur-limit}
For the negative orientation,
\begin{equation}
 m_-(\theta)
 =A-\theta u^T(I_2+\theta N)^{-1}u,
 \qquad
 \lim_{|\theta|\to\infty}m_-(\theta)=m_{\rm sc}>0
 \label{eq:V45-negative-Schur-limit}
\end{equation}
along sectors avoiding the finite transverse poles.
\end{lemma}

\begin{proof}
Factor
\[
 (I_2+\theta N)^{-1}
 =\theta^{-1}(N+\theta^{-1}I_2)^{-1}
\]
and pass to the limit in
\eqref{eq:V45-dynamic-Schur}.  Equation
\eqref{eq:V45-static-Schur} gives strict positivity.
\end{proof}

The quantity \(m_{\rm sc}\) is the exact cross-incidence deficit.  Dropping
\(u\) before the determinant replaces it by \(A\) and loses the complete
square, just as the audited companion arguments warned.  Conversely,
positive off-diagonal incidence can reduce \(A\) only to the strictly
positive Schur complement; it cannot cancel the leading longitudinal
coupling.

Indeed, near bulk glancing,
\[
 |\theta|\asymp|\lambda|^{-1},\qquad
 |\xi|\asymp|\lambda|^{-3},
\]
and
\begin{equation}
 F_-(s,k)
 \sim E\,\xi\theta^2.
 \label{eq:V45-glancing-leading}
\end{equation}
Thus full positive incidence preserves, rather than removes, the
longitudinal singular winding seen in V44.  Equation
\eqref{eq:V45-glancing-leading} is an asymptotic obstruction indicator,
not yet a proof of a global zero.

\subsection{Two exact exclusions}
\label{subsec:V45-exclusions}

\begin{corollary}[A commuting bath is already ruled out]
\label{cor:V45-block-no-go}
Suppose \(M\succ0\) but \(u=0\) in
\eqref{eq:V45-M-block}.  Then neither orientation has a uniform stable
inverse.
\end{corollary}

\begin{proof}
For \(\varepsilon=-1\),
\[
 F_-=(1+A\xi)\det(I_2+\theta N).
\]
The first factor is exactly the V44 negative longitudinal factor with
scalar coupling \(\eta=A>0\), so V44 supplies an open root or surface
mode.

For \(\varepsilon=+1\),
\[
 F_+=(1-A\xi)\det(I_2-\theta N).
\]
Diagonalize the positive matrix \(N\).  Each eigenvalue \(n>0\) supplies
the V44 positive transverse factor \(1-n\theta\), and V44 gives a closed
zero or growing root.  In either orientation the full determinant
vanishes.
\end{proof}

Thus any remaining negative-orientation candidate must have genuine
same-history longitudinal--transverse incidence \(u\ne0\).

The positive orientation can be excluded for the entire positive Gram,
without assuming \(u=0\), throughout a natural subluminal speed range.

\begin{theorem}[Positive-orientation full-Gram path obstruction]
\label{thm:V45-positive-path-no-go}
Let \(M\succ0\) and
\begin{equation}
 0<v\le\frac{\sqrt3}{2}.
 \label{eq:V45-speed-range}
\end{equation}
Then the positive-orientation determinant \(F_+\) has either an open
real-frequency zero or a zero on the closed imaginary-frequency boundary.
Consequently it has no uniform conic floor.
\end{theorem}

\begin{proof}
Set \(r=1\) by homogeneity and write
\[
 {\cal K}(s)
 :=M^{1/2}\kappa_v(s)
 \operatorname{diag}(d(s),b(s),b(s))M^{1/2}.
\]
Whenever the diagonal coefficients are real and positive,
\({\cal K}\) is positive self-adjoint and
\[
 F_+(s)=\det(I_3-{\cal K}(s)).
\]

Consider three continuous paths.

First let \(s=i\omega\), \(0\le\omega\le v\).  Here
\(\kappa_v,\lambda,b\) are nonnegative real.  Moreover
\[
 d=\frac{3-4\omega^2}{8(1-\omega^2)^{3/2}}\ge0
\]
by \eqref{eq:V45-speed-range}.  At \(\omega=v\),
\({\cal K}=0\); at \(\omega=0\), denote the matrix by
\({\cal K}_0\).  If \({\cal K}_0\) has an eigenvalue at least one,
continuity from the zero matrix gives \(F_+=0\) on this path.

Assume therefore that every eigenvalue of \({\cal K}_0\) is less than
one.  On the real path \(s=\sigma\ge0\), all response coefficients are
positive and
\[
 {\cal K}(0)={\cal K}_0,\qquad
 \lim_{\sigma\to\infty}{\cal K}(\sigma)=\frac12M.
\]
If \(M/2\) has an eigenvalue greater than one, the real path crosses one
and gives an open growing root.  If it has an eigenvalue equal to one,
the same zero occurs exactly at \(r=0\), where
\(\kappa_vd=\kappa_vb=1/2\).

It remains only the case \(M/2<I_3\).  Take the high imaginary path
\(s=i\omega\), \(\omega>1\), with boundary values from
\({\rm Re}\,s>0\).  Both \(\kappa_v/\lambda\) and
\(\kappa_vd\) are positive real.  As \(\omega\downarrow1\), all
eigenvalues of \({\cal K}\) tend to \(+\infty\); as
\(\omega\to\infty\), they tend to the eigenvalues of \(M/2\), all below
one.  Continuity forces an eigenvalue through one.  Hence \(F_+=0\) on
the high boundary path.  The three cases exhaust all possible endpoint
inertias.
\end{proof}

\begin{firewall}[The exact surviving branch]
\label{fire:V45-surviving}
V45 does not prove a stable matrix bath.  It proves only that:
\begin{enumerate}
 \item rank-deficient positive measures fail gravitational glancing;
 \item block-diagonal full measures fail both orientations;
 \item for \(0<v\le\sqrt3/2\), every positive definite full Gram fails
 the positive orientation; and
 \item the only common-speed constant-Gram branch not decided here is the
 negative orientation with \(M\succ0\) and \(u\ne0\).
\end{enumerate}
High bath speeds in the positive orientation and the complete
negative noncommuting determinant require further analysis.  No sampled
matrix is promoted to a theorem.
\end{firewall}

\subsection{Angular locality is a separate gate}
\label{subsec:V45-locality}

The response splitting used above depends on the frequency direction.
On the spatial vector part of \(z\),
\begin{equation}
 P_L(k)=\widehat k\,\widehat k^T,\qquad
 P_T(k)=I_2-P_L(k).
 \label{eq:V45-LT-projectors}
\end{equation}
A matrix spectral density tuned directly to \(d\) and \(b\) will generally
contain \(P_L(k)\) and \(P_T(k)\).  These are order-zero Fourier
multipliers and are nonlocal Riesz-transform combinations in physical
space; they are also undefined at \(k=0\) until a compatible extension is
specified.

Therefore a conic determinant margin for a \(k\)-dependent
\(\rho(\omega,k)\) does not by itself prove a local continuum boundary
action.  V45 keeps two independent rows:
\begin{equation}
 \inf_{\Sigma_{\rm st}}s_{\min}
 \bigl(I_3-\varepsilon RZ\bigr)>0,
 \qquad
 {\cal L}_{\rm loc}[\rho]<\infty,
 \label{eq:V45-two-margins}
\end{equation}
where the second row must be populated by an explicit kernel decay or
commutator estimate.  This is the type-compatible content of the YM V25
edge/locality distinction.  Failure of either row must return its
frequency and polarization witness.

\subsection{Finite certificate and V46 acquisition order}
\label{subsec:V45-certificate}

A direct audit compared
\eqref{eq:V45-invariant-determinant} and
\eqref{eq:V45-Schur-determinant} with the full three-by-three determinant
at one hundred random complex frequencies and positive definite Grams.
The largest residuals were
\[
 4.5\times10^{-14}
 \quad\hbox{and}\quad
 1.5\times10^{-14},
\]
respectively.  Thirty random rank-two Grams gave glancing rank at most
four for the five-row metric matrix, as
Theorem~\ref{thm:V45-rank-three} predicts.  One hundred random positive
Grams at \(v=1/2\) were classified by the three paths in
Theorem~\ref{thm:V45-positive-path-no-go}: 76 entered through the low
boundary path, 20 through the real path, and 4 through the high boundary
path.  These counts audit the exhaustive endpoint logic; they are not the
proof.

For the unresolved negative noncommuting branch, eight random positive
Grams at \(v=1/2\) all had open roots in a numerical scan, with examples
\[
 0.0259722+0.9260657i,\qquad
 0.0445885+0.9112130i.
\]
This is evidence for the next theorem, not a universal conclusion.

For V46 set
\begin{equation}
 Y:=\frac{\lambda^2}{r^2},\qquad
 \delta:=1-v^2,\qquad
 \theta^2=\frac{Y-\delta}{4Y},\qquad
 \xi=\frac{4Y-1}{4Y}\theta.
 \label{eq:V45-Y-parametrization}
\end{equation}
Then the remaining negative determinant is the exact algebraic function
\begin{equation}
 F_-(Y)
 =U_M(Y)+\theta(Y)V_M(Y),
 \label{eq:V45-even-odd}
\end{equation}
where
\begin{align}
 U_M(Y)&:=1+\theta^2(C_*h+D),\nonumber\\
 V_M(Y)&:=Ah+B+E h\theta^2,\qquad
 h:=\frac{4Y-1}{4Y}.
 \label{eq:V45-U-V}
\end{align}
The elimination polynomial
\begin{equation}
 {\cal P}_M(Y):=U_M(Y)^2-\theta(Y)^2V_M(Y)^2
 \label{eq:V45-elimination-polynomial}
\end{equation}
must be cleared of denominators, factored, and accompanied by an
unsquared branch assignment.  V46 should decide whether positivity of
\(M\) forces an open or boundary zero.  If it does not, it must give an
exact positive Gram and a verified closed-cone floor.  Only after that
decision should a direction-dependent matrix measure be introduced, and
then both margins in \eqref{eq:V45-two-margins} must be proved.

\subsection{V45 ledger}
\label{subsec:V45-ledger}

\paragraph{New conclusions proved in V45.}
The scheduled NS/YM audit is frozen by filename and hash with a strict
type firewall.  Proposition~\ref{prop:V45-matrix-bath} constructs the
complete positive matrix measure and preserves one common bath history.
Theorem~\ref{thm:V45-rank-three} proves full Gram rank is necessary.
Theorem~\ref{thm:V45-complete-determinant} gives the exact invariant and
Schur-completed determinants.  Corollary~\ref{cor:V45-block-no-go} rules
out commuting Grams, and
Theorem~\ref{thm:V45-positive-path-no-go} rules out every full positive
Gram in the positive orientation for
\(0<v\le\sqrt3/2\).

\paragraph{Imported conclusions not repeated.}
V43 supplies \(R=\operatorname{diag}(d,b,b)\).  V44 supplies the
half-line bath and scalar obstruction.  The NS and YM notes contribute
only assembly requirements; all matrix-bath identities and exclusions are
proved anew here.

\paragraph{Honest continuum status.}
The positive matrix bath has conservative ancestry, complete source
incidence, and an exact determinant, but no global stable floor has been
proved.  The negative noncommuting branch remains open, as do angular
locality, nonlinear covariant stress/Bianchi closure, ghosts, and the
coupled Standard-Model sectors.  The full Standard Model--Einstein
continuum is not proved.

% END APPENDED VERSION V45 MATERIAL



% BEGIN APPENDED VERSION V46 MATERIAL

\clearpage
\section{V46: ratio-plane uniformization and the complete
common-speed positive-matrix bath no-go}
\label{sec:V46-common-speed-no-go}

\subsection{Purpose and exact advance}
\label{subsec:V46-purpose}

V45 reduced the common-speed matrix bath to one unresolved branch: the
negative Green orientation with a positive definite Gram and genuine
longitudinal--transverse incidence.  It left the exact algebraic
elimination card \(\mathcal P_M(Y)\) and required an unsquared branch
assignment.

V46 closes that branch and, at the same time, removes the speed
restriction from the V45 positive-orientation theorem.  The useful
coordinate is not \(Y\) itself but the retarded response ratio
\[
 \theta=\frac{\kappa_v}{2\lambda}.
\]
In this coordinate the full three-channel determinant is a real quintic.
The positive orientation fails because this quintic always has a negative
real root.  The negative orientation fails because positivity of the bath
Gram forces a strict violation of the third Hurwitz condition.  A
ratio-plane mapping lemma turns those algebraic roots into open or
closed-boundary stable-frequency zeros without squaring the determinant.

The result covers every \(M\succeq0\), every \(0<v<1\), and both Green
orientations within the V45 common-speed half-line bath.  It does not
cover sums of baths with different speeds or a frequency-dependent matrix
spectral measure.

\subsection{The retarded ratio plane}
\label{subsec:V46-ratio-plane}

Fix \(r=|k|>0\) and define
\begin{equation}
 \lambda(s,r):=\sqrt{s^2+r^2},
 \qquad
 \kappa_v(s,r):=\sqrt{s^2+v^2r^2},
 \qquad
 \Theta_v(s,r):=\frac{\kappa_v(s,r)}{2\lambda(s,r)},
 \label{eq:V46-Theta}
\end{equation}
using the branches with positive real part for
\({\rm Re}\,s>0\).  Put
\begin{equation}
 \delta:=1-v^2\in(0,1).
 \label{eq:V46-delta}
\end{equation}

\begin{lemma}[Exact ratio-to-frequency map]
\label{lem:V46-ratio-map}
Let \(\theta\) satisfy \({\rm Re}\,\theta>0\) and
\(\theta\ne1/2\).  Set
\begin{equation}
 \frac{s^2}{r^2}
 =\frac{4\theta^2-v^2}{1-4\theta^2}.
 \label{eq:V46-inverse-ratio}
\end{equation}
There is a unique choice of \(s\) with \({\rm Re}\,s>0\) when
\({\rm Im}\,\theta\ne0\).  When \(\theta>0\) is real, the same relation is
realized either by an open real frequency or by the corresponding stable
boundary value described below.  In each case,
\begin{equation}
 \Theta_v(s,r)=\theta.
 \label{eq:V46-ratio-realization}
\end{equation}
More precisely, for real \(\theta>0\):
\begin{enumerate}
 \item \(v/2<\theta<1/2\) gives a real open frequency \(s>0\);
 \item \(0<\theta<v/2\) gives a low purely imaginary boundary frequency;
 \item \(\theta>1/2\) gives a high purely imaginary boundary frequency;
 \item \(\theta=v/2\) gives \(s=0\), while
 \(\theta=1/2\) is the \(r=0\) homogeneous direction.
\end{enumerate}
A nonzero purely imaginary \(\theta\) also gives a stable imaginary
boundary value.  Hence every nonzero response-ratio value with
\({\rm Re}\,\theta\ge0\), including \(\theta=1/2\) through the homogeneous
direction, is realized at an open or closed-boundary stable frequency.
\end{lemma}

\begin{proof}
Write \(w=2\theta\) and \(z=s^2/r^2\).  The squared ratio equation is
\[
 w^2=\frac{z+v^2}{z+1},
 \qquad
 z=\frac{w^2-v^2}{1-w^2}.
\]
If \(w\) is in the first open quadrant, then \(w^2\) is in the upper
half-plane.  The real Moebius map
\[
 W\longmapsto\frac{W-v^2}{1-W}
\]
has determinant \(1-v^2=\delta>0\), so it maps the upper half-plane to
itself.  Its principal square root \(s/r=\sqrt z\) lies in the first
quadrant.  The functions
\[
 \frac{\lambda^2}{r^2}=z+1=\frac{\delta}{1-w^2},
 \qquad
 \frac{\kappa_v^2}{r^2}=z+v^2=w^2(z+1)
\]
then have their stable square-root branches, and analytic continuation
from \(v<w<1\) fixes \(\kappa_v/\lambda=w\), not \(-w\).  Complex
conjugation proves the lower-quadrant statement.

For positive real \(w\), the sign of
\((w^2-v^2)/(1-w^2)\) gives the four cases listed in the theorem.
For imaginary \(w\), the same quotient is negative real and supplies the
boundary value.  The endpoint \(w=v\) gives \(s=0\).  For \(w=1\), take
\(r=0\) and any \(s>0\); then \(\lambda=\kappa_v=s\).
\end{proof}

\subsection{The exact quintic without squared branches}
\label{subsec:V46-quintic}

Retain the V45 block notation
\[
 M=\begin{pmatrix}A&u^T\cr u&N\end{pmatrix}\succ0,
\]
and its invariants
\begin{equation}
 B=\operatorname{tr}N,\qquad
 C_*=A\operatorname{tr}N-|u|^2,\qquad
 D=\det N,\qquad E=\det M.
 \label{eq:V46-invariants}
\end{equation}
All five invariants \(A,B,C_*,D,E\) are strictly positive.  Let
\begin{equation}
 \alpha:=\delta-\frac14.
 \label{eq:V46-alpha}
\end{equation}

The V45 variables obey
\begin{equation}
 \theta=\frac{\kappa_v}{2\lambda},
 \qquad
 Y=\frac{\lambda^2}{r^2}
 =\frac{\delta}{1-4\theta^2},
 \qquad
 h:=\frac{4Y-1}{4Y}
 =\frac{\theta^2+\alpha}{\delta},
 \qquad
 \xi=h\theta.
 \label{eq:V46-uniformization}
\end{equation}

\begin{theorem}[Exact common-speed determinant quintic]
\label{thm:V46-quintic}
Define
\begin{align}
 p_M(\theta)
 :=\;&E\theta^5+C_*\theta^4
 +(A+E\alpha)\theta^3
 +(C_*\alpha+D\delta)\theta^2
 \nonumber\\
 &+(A\alpha+B\delta)\theta+\delta.
 \label{eq:V46-p}
\end{align}
Then the full retarded determinants of V45 satisfy
\begin{equation}
 \delta F_-(s,k)=p_M(\Theta_v(s,k)),
 \qquad
 \delta F_+(s,k)=p_M(-\Theta_v(s,k)).
 \label{eq:V46-both-signs}
\end{equation}
No squaring or extraneous branch is present.
\end{theorem}

\begin{proof}
The V45 negative determinant is
\[
 F_-
 =1+A\xi+B\theta+C_*\xi\theta+D\theta^2
   +E\xi\theta^2.
\]
Insert \(\xi=h\theta\) and
\(h=(\theta^2+\alpha)/\delta\), then multiply by \(\delta\):
\begin{align*}
 \delta F_-
 ={}&\delta
 +A\theta(\theta^2+\alpha)+B\delta\theta
 +C_*\theta^2(\theta^2+\alpha)+D\delta\theta^2\\
 &+E\theta^3(\theta^2+\alpha).
\end{align*}
Collecting powers gives \eqref{eq:V46-p}.  Replacing
\(\theta\) by \(-\theta\) reverses precisely the odd principal-minor
terms, which is the V45 positive determinant.  This proves
\eqref{eq:V46-both-signs}.
\end{proof}

The V45 eliminated \(Y\)-polynomial is the product
\[
 \delta^2F_-(\theta)F_-(-\theta)
 =p_M(\theta)p_M(-\theta).
\]
The quintic formulation retains which factor belongs to which Green
orientation and is therefore stronger than the squared card.

\subsection{The positive orientation fails for every speed}
\label{subsec:V46-positive}

\begin{theorem}[Universal positive-orientation root]
\label{thm:V46-positive-root}
Let \(M\succ0\) and \(0<v<1\).  The positive-orientation determinant
\(F_+\) has an open or closed-boundary stable-frequency zero.
Consequently the V45 speed restriction
\(v\le\sqrt3/2\) is unnecessary for the common-speed class.
\end{theorem}

\begin{proof}
The quintic has
\[
 p_M(0)=\delta>0,
 \qquad
 \lim_{\theta\to-\infty}p_M(\theta)=-\infty
\]
because its leading coefficient is \(E>0\).  The intermediate value
theorem gives a negative real root \(\theta_-<0\).  Put
\(\theta_+:=-\theta_->0\).  By
\eqref{eq:V46-both-signs},
\[
 \delta F_+(\theta_+)=p_M(-\theta_+)=p_M(\theta_-)=0.
\]
Lemma~\ref{lem:V46-ratio-map} realizes \(\theta_+\) as an open real
frequency or a closed imaginary boundary frequency, including the
homogeneous endpoint if \(\theta_+=1/2\).  Thus the stable determinant
fails.
\end{proof}

This proof preserves the complete matrix cross term.  It does not
diagonalize \(M\), and it covers the high-speed branch left open in V45.

\subsection{The negative orientation violates Hurwitz stability}
\label{subsec:V46-negative}

Write the quintic coefficients as
\begin{equation}
 \begin{split}
 a_5&=E,\qquad a_4=C_*,
 \qquad a_3=A+E\alpha,\\
 a_2&=C_*\alpha+D\delta,
 \qquad a_1=A\alpha+B\delta,
 \qquad a_0=\delta.
 \end{split}
 \label{eq:V46-coefficients}
\end{equation}
Its third Hurwitz determinant is
\begin{equation}
 \Delta_3
 :=a_4a_3a_2-a_4^2a_1-a_5a_2^2+a_5a_4a_0.
 \label{eq:V46-Hurwitz-three}
\end{equation}

\begin{lemma}[Strict negative Gram identity]
\label{lem:V46-S-negative}
Put
\begin{equation}
 S_M:=AD-BC_*+E.
 \label{eq:V46-S}
\end{equation}
Then \(S_M<0\).  Moreover,
\begin{equation}
 \boxed{\quad
 \Delta_3
 =\delta\left(C_*S_M-EDa_2\right).
 \quad}
 \label{eq:V46-Delta-three}
\end{equation}
\end{lemma}

\begin{proof}
An orthogonal rotation in the two-dimensional transverse block leaves the
response unchanged and diagonalizes
\[
 N=\operatorname{diag}(n_1,n_2),\qquad
 u=(p,q)^T.
\]
Then
\[
 C_*=A(n_1+n_2)-p^2-q^2,\quad
 D=n_1n_2,\quad
 E=An_1n_2-p^2n_2-q^2n_1.
\]
Direct substitution gives
\begin{equation}
 S_M
 =-n_1(An_1-p^2)-n_2(An_2-q^2)<0,
 \label{eq:V46-S-diagonal}
\end{equation}
because the two displayed principal minors of \(M\) are positive.

Substitute \eqref{eq:V46-coefficients} into
\eqref{eq:V46-Hurwitz-three}.  The terms containing
\(A C_*^2\alpha\) and \(E C_*^2\alpha^2\) cancel, leaving
\[
 \Delta_3
 =\delta\left[
 C_*(AD-BC_*+E)-ED(C_*\alpha+D\delta)
 \right],
\]
which is \eqref{eq:V46-Delta-three}.
\end{proof}

\begin{theorem}[Universal negative-orientation root]
\label{thm:V46-negative-root}
Let \(M\succ0\) and \(0<v<1\).  The negative-orientation determinant
\(F_-\) has an open or closed-boundary stable-frequency zero.  Therefore
the noncommuting branch left by V45 has no uniform stable inverse.
\end{theorem}

\begin{proof}
Assume, for contradiction, that no such zero exists.  By
Lemma~\ref{lem:V46-ratio-map} and
Theorem~\ref{thm:V46-quintic}, every root of \(p_M\) must then lie in the
strict left half \(\theta\)-plane.  Thus \(p_M\) is a Hurwitz polynomial.

A real Hurwitz polynomial with positive leading coefficient has all
coefficients positive: factor it into negative real linear factors and
quadratic factors
\(\theta^2+2a\theta+(a^2+b^2)\) with \(a>0\).
Hence the coefficient
\[
 a_2=C_*\alpha+D\delta
\]
must be positive.  The Routh--Hurwitz criterion also requires
\(\Delta_3>0\).  But Lemma~\ref{lem:V46-S-negative} gives
\[
 \Delta_3
 =\delta(C_*S_M-EDa_2)<0,
\]
because every factor outside \(S_M\) is positive and \(S_M<0\).
This contradiction proves that \(p_M\) has a root with
\({\rm Re}\,\theta\ge0\).  It cannot be \(\theta=0\), since
\(p_M(0)=\delta\).  Lemma~\ref{lem:V46-ratio-map} converts it into the
claimed stable-frequency obstruction.
\end{proof}

The proof also covers the case \(a_2\le0\): then coefficient positivity,
already necessary for Hurwitz stability, fails before the third minor is
needed.

\subsection{Complete common-speed no-go and its boundary}
\label{subsec:V46-complete}

\begin{corollary}[Complete common-speed positive-matrix bath no-go]
\label{cor:V46-complete-no-go}
Let \(0<v<1\), let \(M=M^T\succeq0\), and let
\(Z_M=\kappa_vM\) be the V45 common-speed retarded impedance.
For both \(\varepsilon=+1\) and \(\varepsilon=-1\), the coupled Einstein
stable matrix lacks a uniform closed-half-plane inverse.

If \(\operatorname{rank}M<3\), gravitational glancing remains singular.
If \(M\succ0\), Theorems~\ref{thm:V46-positive-root} and
\ref{thm:V46-negative-root} supply the respective determinant zeros.
\end{corollary}

\begin{proof}
The rank-deficient case is Theorem~\ref{thm:V45-rank-three}.  A
positive semidefinite three-by-three matrix of full rank is positive
definite, so the two V46 theorems exhaust the remaining case.
\end{proof}

\begin{firewall}[Exact scope]
\label{fire:V46-scope}
Corollary~\ref{cor:V46-complete-no-go} closes one large but specific
class: a constant positive internal Gram multiplying one common
square-root bath impedance.  It does not close
\[
 Z(s,k)=\sum_{\ell=1}^L\kappa_{v_\ell}(s,k)M_\ell,
 \qquad M_\ell\succeq0,
\]
with distinct speeds, nor a continuous matrix-valued measure whose
eigenspaces vary with frequency.  Those classes no longer reduce to one
quintic in the ratio \(\theta\).  They must retain the complete
source-incidence and angular-locality rows identified in V45.
\end{firewall}

\subsection{Finite certificate and V47 upstream row}
\label{subsec:V46-certificate}

A direct audit used one thousand random positive definite Grams and
random \(0<v<1\).  It verified
\eqref{eq:V46-both-signs} against the full matrix determinant to
\(1.1\times10^{-12}\), the ratio inverse
\eqref{eq:V46-inverse-ratio} to \(1.6\times10^{-10}\), and the normalized
physical determinant at computed roots to \(4.0\times10^{-15}\).
The Hurwitz identity \eqref{eq:V46-Delta-three} agreed to
\(9.1\times10^{-13}\).  Among those samples, \(a_2\le0\) occurred 80
times; in the other 920 cases the third Hurwitz determinant was strictly
negative.  The largest sampled value of \(S_M\) was
\(-0.1627975\).

Representative negative-orientation roots were
\[
 \theta=0.1454046+0.8826956i,
 \qquad
 s=0.0317463+0.8835773i,
\]
and
\[
 \theta=0.2626377+0.5110638i,
 \qquad
 s=0.0388978+0.9387303i.
\]
Positive-orientation witnesses included both an open real root
\(s=1.1811726\) and high imaginary boundary roots.  These samples audit
the exact theorem; they are not used in its proof.

V47 must break the one-ratio mechanism rather than retune \(M\).  The
minimal next ansatz is
\begin{equation}
 Z_{(2)}(s,k)
 :=\kappa_{v_1}(s,k)M_1+\kappa_{v_2}(s,k)M_2,
 \qquad
 M_1,M_2\succeq0,\quad v_1\ne v_2.
 \label{eq:V46-two-speed-target}
\end{equation}
It has an explicit direct-sum conservative bath and a positive matrix
spectral measure.  V47 should derive its exact determinant in the V43
response basis, test rank-three incidence at both bath cones and
gravitational glancing, and decide whether the two independent ratios can
satisfy a uniform stable floor.  If the construction uses
\(P_L(k),P_T(k)\), it must also prove the locality row
\eqref{eq:V45-two-margins}; a frequencywise margin alone is insufficient.

\subsection{V46 ledger}
\label{subsec:V46-ledger}

\paragraph{New conclusions proved in V46.}
Lemma~\ref{lem:V46-ratio-map} gives the exact physical realization of the
retarded ratio half-plane.  Theorem~\ref{thm:V46-quintic} replaces the
squared V45 elimination card by an unsquared determinant quintic.
Theorem~\ref{thm:V46-positive-root} removes the prior bath-speed
restriction.  Lemma~\ref{lem:V46-S-negative} and
Theorem~\ref{thm:V46-negative-root} prove the universal third-Hurwitz-minor
obstruction.  Corollary~\ref{cor:V46-complete-no-go} closes the entire
common-speed constant positive-matrix bath class.

\paragraph{Imported conclusions not repeated.}
V43 supplies the exact response \(R\).  V44 supplies the conservative
square-root bath.  V45 supplies the complete Gram invariants and the
rank-deficient glancing theorem.  V46 starts from those cards and proves
the new ratio map and Hurwitz obstruction without repeating their
derivations.

\paragraph{Honest continuum status.}
Every constant positive matrix multiplying one common bath impedance is
now ruled out.  A multi-speed or genuinely frequency-dependent positive
matrix measure has not passed the stable and locality gates.  Nonlinear
covariant stress/Bianchi closure, ghosts, and the coupled Standard-Model
sectors remain open.  The full Standard Model--Einstein continuum is not
proved.

% END APPENDED VERSION V46 MATERIAL



% BEGIN APPENDED VERSION V47 MATERIAL

\clearpage
\section{V47: the two-speed direct-sum bath, exact mixed determinant,
and the covariance bottleneck}
\label{sec:V47-two-speed}

\subsection{Purpose and acquisition order}
\label{subsec:V47-purpose}

V46 proves that no constant positive matrix multiplying one common
square-root impedance can stabilize the Einstein three-port response.
The next local passive enlargement is therefore a direct sum of two
half-line baths with distinct speeds.  V47 constructs that enlargement
before testing it, keeps all mixed principal minors in the determinant,
and separates three questions which cannot be identified:
\begin{enumerate}
 \item collective incidence at gravitational glancing;
 \item stable invertibility of the resulting two-ratio symbol; and
 \item locality together with tangential rotational covariance.
\end{enumerate}

The exact advance is fourfold.  First, the two-speed impedance is derived
from a local positive action.  Second, its determinant is written without
discarding either matrix cross term.  Third, a self-adjoint path argument
gives exact necessary inequalities for any positive-orientation
candidate.  Fourth, tangential covariance is imposed on the local action.
It forces constant bath Grams to commute with rotations and reduces the
negative-orientation problem to one unsquared two-radical longitudinal
equation.  A degree-six elimination card is derived for that equation,
with its branch filter left explicit.

V47 does not assert that the degree-six card has, or lacks, a physical
root for every pair of positive couplings.  That is the next exact
bottleneck.

\subsection{Local conservative realization}
\label{subsec:V47-action}

Let \(0<v_1<v_2<1\), let \(M_j=M_j^T\succeq0\), and introduce independent
three-component half-line fields \(q_j(a,y)\), \(j=1,2\), with the common
boundary trace
\begin{equation}
 q_1(0,y)=q_2(0,y)=z(y)=Jx(y).
 \label{eq:V47-common-trace}
\end{equation}
Define
\begin{align}
 S_{{\rm bath},(2)}[q_1,q_2]
 :=\frac12\sum_{j=1}^2\int_{\partial M}\int_0^\infty
 \bigl[&(D_tq_j)^TM_j(D_tq_j)
 -(\partial_aq_j)^TM_j(\partial_aq_j)
 \nonumber\\
 &-v_j^2(D_iq_j)^TM_j(D_iq_j)\bigr]
 \dd a\,\dd^3y.
 \label{eq:V47-action}
\end{align}

\begin{proposition}[Direct-sum positive bath]
\label{prop:V47-direct-sum}
The action \eqref{eq:V47-action} has nonnegative energy
\begin{equation}
 E_{{\rm bath},(2)}
 =\frac12\sum_{j=1}^2\int\!\!\int
 \left[(D_tq_j)^TM_j(D_tq_j)
 +(\partial_aq_j)^TM_j(\partial_aq_j)
 +v_j^2(D_iq_j)^TM_j(D_iq_j)\right]
 \dd a\,\dd^2y.
 \label{eq:V47-energy}
\end{equation}
Retarded elimination gives the total boundary force
\begin{equation}
 f_{(2)}=Z_{(2)}(s,k)z,
 \qquad
 Z_{(2)}(s,k)=\kappa_1M_1+\kappa_2M_2,
 \qquad
 \kappa_j:=\sqrt{s^2+v_j^2|k|^2}.
 \label{eq:V47-impedance}
\end{equation}
For \({\rm Re}\,s>0\),
\begin{equation}
 {\rm Re}\,\langle u,Z_{(2)}u\rangle
 =\sum_{j=1}^2({\rm Re}\,\kappa_j)u^TM_ju\ge0.
 \label{eq:V47-positive-real}
\end{equation}
The matrix spectral density is the positive sum
\begin{equation}
 \rho_{(2)}(\omega,k)
 =\frac2\pi\sum_{j=1}^2
 \sqrt{v_j^2|k|^2-\omega^2}\,
 {\bf1}_{[0,v_j|k|]}(\omega)M_j\succeq0.
 \label{eq:V47-density}
\end{equation}
If the matrices are constant in the physical port frame, the realization
is local and its Green current is the sum of the two V45 bath currents.
\end{proposition}

\begin{proof}
Vary the two summands independently.  Their decaying stable solutions are

\[
 q_j(a)=e^{-a\kappa_j}z,
\]

and their boundary tractions are
\(-M_j\partial_aq_j(0)=\kappa_jM_jz\).  The common trace variation sums
the two tractions, proving \eqref{eq:V47-impedance}.  The energy and Green
current follow by the same integration by parts as in V45, now summed over
the direct-sum index.  The scalar spectral representation applied to each
speed gives \eqref{eq:V47-density}.  Positivity of \(M_j\) proves the two
positivity statements.  Only differential operators and constant
matrices occur in \eqref{eq:V47-action}, proving locality in the physical
port frame.
\end{proof}

\subsection{Collective glancing incidence and the two bath cones}
\label{subsec:V47-cones}

Put \(r=|k|>0\),
\begin{equation}
 \lambda:=\sqrt{s^2+r^2},\qquad
 Y:=\frac{\lambda^2}{r^2},\qquad
 \delta_j:=1-v_j^2.
 \label{eq:V47-Y-delta}
\end{equation}

\begin{theorem}[Collective rank is necessary and sufficient for incidence]
\label{thm:V47-collective-rank}
At gravitational glancing \(s=ir+0\), define
\begin{equation}
 G_{\rm g}:=\sqrt{\delta_1}M_1+\sqrt{\delta_2}M_2.
 \label{eq:V47-Gg}
\end{equation}
Then
\begin{equation}
 \ker G_{\rm g}=\ker(M_1+M_2)
 =\ker M_1\cap\ker M_2.
 \label{eq:V47-kernel-identity}
\end{equation}
If \(M_1+M_2\) is not positive definite, the two-speed effective Einstein
matrix is singular at gravitational glancing for both Green orientations.
Thus collective full incidence requires
\begin{equation}
 M_1+M_2\succ0.
 \label{eq:V47-collective-positive}
\end{equation}
Conversely, \eqref{eq:V47-collective-positive} makes the bath update full
rank at glancing.  This converse is an incidence statement only; it is not
a stable-inverse theorem.
\end{theorem}

\begin{proof}
For positive semidefinite matrices and positive weights,

\[
 x^TG_{\rm g}x=0
 \quad\Longleftrightarrow\quad
 x^TM_1x=x^TM_2x=0
 \quad\Longleftrightarrow\quad
 M_1x=M_2x=0.
\]

The same equivalence holds for \(M_1+M_2\), proving
\eqref{eq:V47-kernel-identity}.  At gravitational glancing,
\(\kappa_j=ir\sqrt{\delta_j}\), so
\(Z_{(2)}=irG_{\rm g}\).  If its rank is less than three, the rank argument
of Theorem~\ref{thm:V45-rank-three} leaves a nonzero vector in the
three-dimensional gravitational kernel untouched.  This proves
singularity.  If the sum is positive definite, the kernel identity makes
\(G_{\rm g}\) positive definite and hence the update has full port rank.
\end{proof}

The two bath cones do not impose the same rank condition.  Define
\begin{equation}
 \theta_j:=\frac{\kappa_j}{2\lambda},
 \qquad
 h:=\frac{4\lambda^2-r^2}{4\lambda^2}.
 \label{eq:V47-theta-h}
\end{equation}
On the stable boundary their exact values are
\begin{align}
 s=iv_1r+0:\quad&
 \theta_1=0,
 \qquad
 \theta_2=\frac{\sqrt{v_2^2-v_1^2}}{2\sqrt{1-v_1^2}},
 \qquad h=1-\frac1{4\delta_1},
 \label{eq:V47-lower-cone}\\
 s=iv_2r+0:\quad&
 \theta_2=0,
 \qquad
 \theta_1=\frac{i\sqrt{v_2^2-v_1^2}}{2\sqrt{1-v_2^2}},
 \qquad h=1-\frac1{4\delta_2}.
 \label{eq:V47-upper-cone}
\end{align}
At the lower cone only \(M_2\) remains in the impedance, while at the
upper cone only \(M_1\) remains.  Since \(\lambda=r\sqrt{1-v_j^2}\ne0\)
at either cone, \(L_0\) is invertible there.  Rank deficiency of the
surviving matrix therefore does not by itself prove a determinant zero;
the exact determinant below must be evaluated.  This distinguishes bath
thresholds from the three-dimensional gravitational glancing kernel.

\subsection{The exact mixed determinant}
\label{subsec:V47-determinant}

For a fixed direction rotate to the response basis
\((v_L;v_T,\phi)\) and write
\begin{equation}
 M_j=\begin{pmatrix}A_j&u_j^T\cr u_j&N_j\end{pmatrix},
 \qquad B_j:=\operatorname{tr}N_j,
 \qquad C_j:=A_jB_j-|u_j|^2,
 \qquad D_j:=\det N_j,
 \qquad E_j:=\det M_j.
 \label{eq:V47-blocks}
\end{equation}
The mixed principal-minor coefficients are
\begin{align}
 C_{12}&:=A_1B_2+A_2B_1-2u_1^Tu_2,
 \nonumber\\
 D_{12}&:=\operatorname{tr}\bigl((\operatorname{adj}N_1)N_2\bigr),
 \nonumber\\
 E_{112}&:=\operatorname{tr}\bigl((\operatorname{adj}M_1)M_2\bigr),
 \qquad
 E_{122}:=\operatorname{tr}\bigl((\operatorname{adj}M_2)M_1\bigr).
 \label{eq:V47-mixed-minors}
\end{align}
For positive semidefinite \(M_1,M_2\), all these mixed coefficients are
nonnegative.  This follows either from the positive mixed discriminants
or by expanding the relevant two-by-two and three-by-three principal
minors of \(t_1M_1+t_2M_2\succeq0\) for \(t_1,t_2\ge0\).

\begin{theorem}[Two-speed full-Gram determinant]
\label{thm:V47-mixed-determinant}
Set
\begin{equation}
 T(\theta_1,\theta_2):=\theta_1M_1+\theta_2M_2,
 \qquad H(h):=\operatorname{diag}(h,1,1).
 \label{eq:V47-TH}
\end{equation}
For either orientation,
\begin{align}
 F_\varepsilon
 &:=\frac{\det L_{(2)}^\varepsilon}{\det L_0}
 =\det\bigl(I_3-\varepsilon H(h)T(\theta_1,\theta_2)\bigr)
 \nonumber\\
 &=1-\varepsilon\bigl[h(A_1\theta_1+A_2\theta_2)
 +(B_1\theta_1+B_2\theta_2)\bigr]
 \nonumber\\
 &\quad+h\bigl(C_1\theta_1^2+C_{12}\theta_1\theta_2
 +C_2\theta_2^2\bigr)
 +\bigl(D_1\theta_1^2+D_{12}\theta_1\theta_2
 +D_2\theta_2^2\bigr)
 \nonumber\\
 &\quad-\varepsilon h\bigl(E_1\theta_1^3
 +E_{112}\theta_1^2\theta_2
 +E_{122}\theta_1\theta_2^2+E_2\theta_2^3\bigr).
 \label{eq:V47-mixed-determinant}
\end{align}
The two ratios lie on the exact stable-sheet curve
\begin{equation}
 \theta_j^2=\frac14\left(1-\frac{\delta_j}{Y}\right),
 \qquad
 \frac{1-4\theta_1^2}{\delta_1}
 =\frac{1-4\theta_2^2}{\delta_2}
 =\frac1Y,
 \qquad
 h=1-\frac1{4Y}.
 \label{eq:V47-ratio-curve}
\end{equation}
No determinant squaring has occurred.
\end{theorem}

\begin{proof}
The matrix determinant lemma and the V43 response give

\[
 F_\varepsilon
 =\det\left(I_3-\varepsilon
 \operatorname{diag}(d,b,b)(\kappa_1M_1+\kappa_2M_2)\right).
\]

Since \(d\kappa_j=h\theta_j\) and \(b\kappa_j=\theta_j\), this is the
first line of \eqref{eq:V47-mixed-determinant}.  Expand by principal
minors of \(T\).  Its longitudinal trace coefficient is
\(A_1\theta_1+A_2\theta_2\), its transverse trace coefficient is
\(B_1\theta_1+B_2\theta_2\), and its two quadratic principal-minor sums
have the \(C\)- and \(D\)-coefficients in
\eqref{eq:V47-blocks,eq:V47-mixed-minors}.  Finally,

\[
 \det T=E_1\theta_1^3+E_{112}\theta_1^2\theta_2
 +E_{122}\theta_1\theta_2^2+E_2\theta_2^3.
\]

This proves the determinant.  The identities
\(\kappa_j^2=\lambda^2-\delta_jr^2\) prove
\eqref{eq:V47-ratio-curve} with the inherited retarded branches.
\end{proof}

The curve \eqref{eq:V47-ratio-curve}, rather than a free pair of complex
variables, is the exact two-speed response surface.  Eliminating one
ratio introduces a new square root.  This is precisely why the V46
one-variable quintic and its third Hurwitz minor do not transfer to V47.

\subsection{Positive orientation: a necessary gain chamber}
\label{subsec:V47-positive}

The positive orientation is already sharply constrained without solving
the full two-ratio equation.  All matrix inequalities in this subsection
are in the fixed response basis for the chosen nonzero direction \(k\).
Put
\begin{equation}
 M_\Sigma:=M_1+M_2,
 \qquad M_v:=v_1M_1+v_2M_2,
 \qquad H_0:=\operatorname{diag}(3/4,1,1).
 \label{eq:V47-endpoint-matrices}
\end{equation}

\begin{theorem}[Necessary chamber for a positive-orientation candidate]
\label{thm:V47-positive-chamber}
Assume \(M_j\succeq0\) and that the positive-orientation determinant has
no zero on the open stable cone or its closed frequency boundary.  Then
all of the following strict conditions are necessary:
\begin{align}
 &M_\Sigma\succ2I_3,
 \label{eq:V47-high-condition}\\
 &M_v\succ\operatorname{diag}(8/3,2,2),
 \label{eq:V47-zero-condition}\\
 &v_1<\frac{\sqrt3}{2}.
 \label{eq:V47-slow-condition}
\end{align}
Moreover, define
\begin{equation}
 h_1:=1-\frac1{4(1-v_1^2)},
 \qquad
 \eta_{21}:=\frac{\sqrt{v_2^2-v_1^2}}{2\sqrt{1-v_1^2}},
 \qquad
 H_1:=\operatorname{diag}(h_1,1,1).
 \label{eq:V47-lower-data}
\end{equation}
Then \(h_1>0\) and the faster bath alone must satisfy
\begin{equation}
 \eta_{21}H_1^{1/2}M_2H_1^{1/2}\succ I_3,
 \quad\hbox{equivalently}\quad
 M_2\succ\eta_{21}^{-1}H_1^{-1}.
 \label{eq:V47-cone-condition}
\end{equation}
In particular \(M_2\succ0\).  These inequalities are necessary, not
sufficient, for stable invertibility.
\end{theorem}

\begin{proof}
Theorem~\ref{thm:V47-collective-rank} first gives
\(M_\Sigma\succ0\).  On the high imaginary boundary write
\(s=i\omega r+0\), \(\omega>1\).  Then
\begin{equation}
 \theta_j(\omega)
 =\frac{\sqrt{\omega^2-v_j^2}}{2\sqrt{\omega^2-1}}>0,
 \qquad
 h(\omega)=1+\frac1{4(\omega^2-1)}>0.
 \label{eq:V47-high-path}
\end{equation}
Thus \(H^{1/2}TH^{1/2}\) is positive self-adjoint and
\begin{equation}
 F_+=\det\bigl(I_3-H^{1/2}TH^{1/2}\bigr).
 \label{eq:V47-selfadjoint-F}
\end{equation}
As \(\omega\downarrow1\), its least eigenvalue tends to \(+\infty\),
because \(T\sim(2\sqrt{\omega^2-1})^{-1}G_{\rm g}\) and
\(G_{\rm g}\succ0\).  As \(\omega\to\infty\), the matrix tends to
\(M_\Sigma/2\).  If the latter had an eigenvalue below one, continuity of
the ordered eigenvalues would give a boundary zero.  Equality gives the
exact homogeneous zero at \(r=0\), where \(H=I_3\) and
\(\theta_1=\theta_2=1/2\).  Hence \eqref{eq:V47-high-condition} is
necessary.

On the real path \(s=\sigma r\), \(\sigma\ge0\), the same self-adjoint
form holds.  It tends to \(M_\Sigma/2\succ I_3\), while at \(\sigma=0\)
it is
\begin{equation}
 K_0=H_0^{1/2}\frac{M_v}{2}H_0^{1/2}.
 \label{eq:V47-K0}
\end{equation}
Absence of an eigenvalue crossing through one forces \(K_0\succ I_3\).
Congruence by \(H_0^{-1/2}\) gives
\eqref{eq:V47-zero-condition}.

Now follow the low imaginary boundary from \(\omega=0\).  As long as
\(\omega\le v_1\), both ratios are nonnegative real.  Moreover
\begin{equation}
 h(\omega)=\frac{3-4\omega^2}{4(1-\omega^2)}.
 \label{eq:V47-low-h}
\end{equation}
If \(v_1\ge\sqrt3/2\), this path reaches \(h=0\) before or at the lower bath cone.
At that point \(H^{1/2}TH^{1/2}\) has a zero eigenvalue.  Since all three
eigenvalues start above one by \eqref{eq:V47-K0}, continuity gives a
positive-orientation zero.  Therefore \eqref{eq:V47-slow-condition} is
necessary.

When \(v_1<\sqrt3/2\), the path instead reaches the lower bath cone with
\(H=H_1\succ0\), \(\theta_1=0\), and
\(\theta_2=\eta_{21}\).  Its endpoint matrix is

\[
 K_1=\eta_{21}H_1^{1/2}M_2H_1^{1/2}.
\]

If its least eigenvalue were at most one, the path from
\(K_0\succ I_3\) would again cross one, including equality at the cone.
Thus \(K_1\succ I_3\), which is exactly
\eqref{eq:V47-cone-condition}.
\end{proof}

The chamber is severe: the total high-frequency gain, the zero-frequency
weighted gain, and the faster bath acting alone at the lower threshold
must all dominate the identity.  Passing these endpoint tests supplies no
control of the complex segment between the two bath cones.

\subsection{Local rotational covariance forces a scalar bottleneck}
\label{subsec:V47-covariance}

The determinant calculation permits arbitrary matrices in a fixed port
frame.  The continuum programme also requires covariance under rotations
of the two tangential coordinates.  Let
\begin{equation}
 \rho(Q):=\operatorname{diag}(Q,1),
 \qquad Q\in SO(2),
 \label{eq:V47-rho}
\end{equation}
acting on the physical port \(z=(\mathbf v,\phi)\).

\begin{theorem}[Commutant of local tangential rotations]
\label{thm:V47-covariant-commutant}
Suppose the distinct-speed local action \eqref{eq:V47-action} is invariant
under every constant tangential rotation and rotations do not change the
bath speeds.  Then each constant symmetric Gram has the form
\begin{equation}
 M_j=\operatorname{diag}(a_jI_2,c_j),
 \qquad a_j,c_j\ge0.
 \label{eq:V47-covariant-M}
\end{equation}
Conversely, matrices of this form make the direct-sum bath tangentially
rotationally invariant and local.
\end{theorem}

\begin{proof}
Distinct speeds are distinct eigenvalues of the tangential principal
symbol, so an invariance cannot mix the two bath labels.  Hence each Gram
obeys

\[
 \rho(Q)^TM_j\rho(Q)=M_j
 \qquad\hbox{for every }Q\in SO(2).
\]

Write \(M_j=\left(\begin{smallmatrix}S&w\\w^T&c\end{smallmatrix}\right)\).
The displayed identity says \(Q^TSQ=S\) and \(Q^Tw=w\) for every rotation.
The only fixed vector is zero, and the only symmetric two-tensor fixed by
all planar rotations is a scalar multiple of \(I_2\).  Positivity gives
the nonnegative coefficients in \eqref{eq:V47-covariant-M}.  The converse
is immediate in \eqref{eq:V47-action}.
\end{proof}

After rotating \(k\) to the longitudinal axis, define
\begin{equation}
 z_v:=a_1\theta_1+a_2\theta_2,
 \qquad
 z_\phi:=c_1\theta_1+c_2\theta_2.
 \label{eq:V47-zv-zphi}
\end{equation}
The exact determinant factorizes as
\begin{equation}
 F_\varepsilon
 =(1-\varepsilon h z_v)
  (1-\varepsilon z_v)
  (1-\varepsilon z_\phi).
 \label{eq:V47-covariant-factorization}
\end{equation}

\begin{proposition}[Exact negative-orientation reduction]
\label{prop:V47-negative-reduction}
In the local rotationally covariant class, the open and closed-boundary
zeros of \(F_-\) can arise only from
\begin{equation}
 g_-(s,k):=1+h(s,k)z_v(s,k).
 \label{eq:V47-gminus}
\end{equation}
The two factors \(1+z_v\) and \(1+z_\phi\) never vanish on the stable
frequency closure.  If either \(a_1=0\) or \(a_2=0\), the remaining
longitudinal factor is the already excluded one-speed scalar bath.
Therefore a genuinely new covariant negative-orientation candidate must
have
\begin{equation}
 a_1a_2>0,
 \qquad v_1\ne v_2.
 \label{eq:V47-genuine-two-speed}
\end{equation}
\end{proposition}

\begin{proof}
For \({\rm Re}\,s>0\), the branch argument behind
Lemma~\ref{lem:V46-ratio-map} gives
\({\rm Re}\,\theta_j>0\).  One direct proof is to apply the positive
determinant Moebius map

\[
 s^2/r^2\longmapsto
 \frac{s^2/r^2+v_j^2}{s^2/r^2+1}=4\theta_j^2
\]

to the upper and lower half-planes and then select the square root with
positive real part.  On the imaginary boundary, each \(\theta_j\) is
positive real below its cone, positive imaginary between its cone and
gravitational glancing, and positive real above gravitational glancing.
Thus nonnegative sums \(z_v,z_\phi\) cannot equal \(-1\) in the open cone
or on its boundary.  Factorization
\eqref{eq:V47-covariant-factorization} proves the first assertion.  If one
vector weight vanishes, \(g_-\) contains only one square-root speed and is
the V44 negative longitudinal factor; V46 also includes it in the
common-speed no-go.  Hence \eqref{eq:V47-genuine-two-speed} is necessary
for a new branch.
\end{proof}

This theorem makes the locality issue concrete.  Frequencywise matrices
with longitudinal--transverse cross terms can evade the commutant only by
using \(P_L(k)\), \(P_T(k)\), or an equivalent angular selector.  Such a
selector is a nonlocal Riesz multiplier and must satisfy the second margin
in \eqref{eq:V45-two-margins}; it is not supplied by the local direct-sum
action.

\subsection{The branch-safe two-radical card}
\label{subsec:V47-radical-card}

In the genuine covariant branch let \(a_1,a_2>0\).  On the stable sheet,
\begin{equation}
 \theta_j(Y)=\frac{\sqrt{Y-\delta_j}}{2\sqrt Y},
 \qquad
 h(Y)=\frac{4Y-1}{4Y},
 \label{eq:V47-Y-branches}
\end{equation}
where every square root is inherited from
\(\lambda=\sqrt{s^2+r^2}\) and \(\kappa_j\).

\begin{theorem}[Unsquared equation and degree-six elimination]
\label{thm:V47-elimination}
The physical negative longitudinal factor vanishes exactly when
\begin{equation}
 {\cal U}(Y):=
 8Y^{3/2}+(4Y-1)
 \left[a_1\sqrt{Y-\delta_1}+a_2\sqrt{Y-\delta_2}\right]=0
 \label{eq:V47-unsquared-U}
\end{equation}
on the retarded stable sheet.  Define
\begin{align}
 P(Y)&:=64Y^3-(4Y-1)^2
 \left[a_1^2(Y-\delta_1)+a_2^2(Y-\delta_2)\right],
 \label{eq:V47-P}\\
 {\cal Q}(Y)&:=P(Y)^2
 -4a_1^2a_2^2(4Y-1)^4
 (Y-\delta_1)(Y-\delta_2).
 \label{eq:V47-Q}
\end{align}
Then every physical zero satisfies the real polynomial equation
\begin{equation}
 {\cal Q}(Y)=0,
 \qquad \deg {\cal Q}\le6.
 \label{eq:V47-Qzero}
\end{equation}
The converse is false without the original branch test
\eqref{eq:V47-unsquared-U}.
\end{theorem}

\begin{proof}
Substitute \eqref{eq:V47-Y-branches} into \(1+hz_v=0\) and multiply by
\(8Y^{3/2}\).  This gives \eqref{eq:V47-unsquared-U} without squaring.
If it holds, squaring once and separating the mixed radical gives

\[
 P(Y)=2a_1a_2(4Y-1)^2
 \sqrt{(Y-\delta_1)(Y-\delta_2)}.
\]

Squaring again gives \eqref{eq:V47-Q}.  Both terms have degree at most
six.  Each squaring loses a sign and sheet assignment, so a root of
\({\cal Q}\) is only a candidate until the retarded roots are reconstructed
and \({\cal U}=0\) is checked.
\end{proof}

\begin{firewall}[What the polynomial does and does not prove]
\label{fire:V47-polynomial-scope}
Equation \eqref{eq:V47-Qzero} is a finite acquisition card, not a no-go
theorem.  A root \(Y\) is physically relevant only if
\(Y=1+s^2/r^2\) for the stable choice \({\rm Re}\,s\ge0\), all three
retarded square-root signs agree, and
\eqref{eq:V47-unsquared-U} vanishes.  Counting all six polynomial roots
without these filters would repeat the squared-branch error avoided in
V46.
\end{firewall}

\subsection{Finite certificate and V48 handoff}
\label{subsec:V47-certificate}

A direct numerical audit used one thousand random positive definite pairs,
random distinct speeds, complex frequencies with \({\rm Re}\,s>0\), and
both Green orientations.  The largest discrepancy between the full
three-by-three determinant and \eqref{eq:V47-mixed-determinant} was
\(3.8\times10^{-11}\); the ratio-curve residual in
\eqref{eq:V47-ratio-curve} was \(1.1\times10^{-14}\).

Three thousand random positive pairs were then passed through the
ordered necessary tests in Theorem~\ref{thm:V47-positive-chamber}.  The
first failing rows were: 457 at the high boundary, 238 on the real path,
36 at the \(h=0\) low boundary, and 813 at the lower bath cone.  The
remaining 1456 samples entered the necessary chamber; none is thereby
declared stable.

As a branch-filter audit, the covariant data
\[
 (a_1,a_2;v_1,v_2)=(0.8,1.3;0.35,0.78)
\]
gave the conjugate physical roots
\begin{equation}
 Y=0.1636885\pm0.0820775i,
 \qquad
 s/r=0.0448218\pm0.9155984i,
 \label{eq:V47-sample-root}
\end{equation}
with normalized unsquared residual \(1.6\times10^{-16}\).  Two other
positive parameter sets also produced filtered open roots, with largest
normalized residual \(7.2\times10^{-13}\).  These examples show that the
degree-six candidates can be genuine; they do not establish a universal
root theorem.

V48 should attack \eqref{eq:V47-unsquared-U} upstream on the slit
\(Y\)-plane.  The required result is an argument-principle or exact
root-location theorem which either:
\begin{enumerate}
 \item proves that every \(a_1,a_2>0\) and
 \(0<v_1<v_2<1\) has an open or boundary stable root; or
 \item exhibits a rigorously certified root-free parameter chamber and a
 uniform lower bound for \(|g_-|\).
\end{enumerate}
Only after that dichotomy is decided should the programme enlarge to
three speeds or to a genuinely matrix-valued continuous spectral measure.
If a noncovariant angular Gram is reconsidered, its physical-space kernel
and commutator bounds must be proved at the same time.

\subsection{V47 ledger}
\label{subsec:V47-ledger}

\paragraph{New conclusions proved in V47.}
Proposition~\ref{prop:V47-direct-sum} constructs the local passive
two-speed bath.  Theorem~\ref{thm:V47-collective-rank} identifies its exact
collective glancing carrier.  Theorem~\ref{thm:V47-mixed-determinant}
retains every mixed principal minor on the two-ratio response curve.
Theorem~\ref{thm:V47-positive-chamber} gives four strict necessary
positive-orientation inequalities.  Theorem~\ref{thm:V47-covariant-commutant}
classifies local rotationally covariant constant Grams.
Proposition~\ref{prop:V47-negative-reduction} reduces their negative
orientation to one longitudinal factor, and
Theorem~\ref{thm:V47-elimination} gives its branch-safe finite card.

\paragraph{Imported conclusions not repeated.}
V43 supplies the three-port response, V45 supplies the matrix determinant
method and the independent locality margin, and V46 supplies the complete
one-speed no-go and ratio branch discipline.  V47 uses those cards without
rederiving their earlier sector calculations.

\paragraph{Honest continuum status.}
The minimal two-speed bath is now local, passive, and algebraically
explicit.  Stable invertibility has not been proved.  The covariant
negative branch is reduced to \eqref{eq:V47-unsquared-U}; the positive
branch has only necessary endpoint inequalities.  Nonlinear covariant
stress/Bianchi closure, ghosts, and the coupled Standard-Model sectors
remain open.  The full Standard Model--Einstein continuum is not proved.

% END APPENDED VERSION V47 MATERIAL



% BEGIN APPENDED VERSION V48 MATERIAL

\clearpage
\section{V48: a Nyquist-index no-go for every local covariant
negative square-root bath mixture}
\label{sec:V48-nyquist-no-go}

\subsection{Purpose and exact advance}
\label{subsec:V48-purpose}

V47 reduced the local rotationally covariant negative orientation to the
longitudinal factor
\[
 1+d(s,r)\sum_{\ell}a_\ell\kappa_{v_\ell}(s,r).
\]
For two speeds it also produced a degree-six elimination polynomial, but
correctly left the physical root count open because two squarings obscure
the retarded sheet.

V48 resolves that bottleneck without further elimination.  The relevant
object is the image of the imaginary-frequency boundary under the
unsquared longitudinal factor.  A boundary zero can occur only on a
short real segment above the zero of the longitudinal response.  If no
such zero occurs, a right-hand indentation around gravitational glancing
rotates the factor by \(3\pi/2\).  The complete Nyquist contour then has
winding number two and forces exactly two open right-half-plane zeros.

The proof applies to every finite positive mixture of subluminal
square-root baths and, under a compact-support hypothesis, to a positive
measure of bath speeds.  It therefore closes the negative orientation for
the entire local rotationally covariant direct-sum class, rather than only
the two-speed example.

\subsection{The positive speed measure and its boundary values}
\label{subsec:V48-measure}

Let \(\mu\) be a finite nonzero positive Borel measure with
\begin{equation}
 \operatorname{supp}\mu\subset[v_-,v_+]\Subset(0,1).
 \label{eq:V48-support}
\end{equation}
Define
\begin{align}
 \lambda(s,r)&:=\sqrt{s^2+r^2},
 \qquad
 \kappa_v(s,r):=\sqrt{s^2+v^2r^2},
 \nonumber\\
 d(s,r)&:=\frac{4\lambda^2-r^2}{8\lambda^3},
 \qquad
 {\mathfrak z}_\mu(s,r):=
 \int\kappa_v(s,r)\,\dd\mu(v),
 \label{eq:V48-zmu}\\
 g_\mu(s,r)&:=1+d(s,r){\mathfrak z}_\mu(s,r).
 \label{eq:V48-gmu}
\end{align}
All square roots are retarded: their real parts are positive for
\({\rm Re}\,s>0\).  The factor is homogeneous of degree zero, so set
\(r=1\) until the final corollary.

For \(0\le\omega<1\), put
\begin{align}
 A_\mu(\omega)
 &:=\int_{(\omega,1)}
       \sqrt{v^2-\omega^2}\,\dd\mu(v),
 \nonumber\\
 B_\mu(\omega)
 &:=\int_{(0,\omega)}
       \sqrt{\omega^2-v^2}\,\dd\mu(v),
 \label{eq:V48-AB}\\
 D_-(\omega)
 &:=\frac{3-4\omega^2}
 {8(1-\omega^2)^{3/2}}.
 \label{eq:V48-Dminus}
\end{align}
For \(\omega>1\), put
\begin{align}
 C_\mu(\omega)
 &:=\int\sqrt{\omega^2-v^2}\,\dd\mu(v),
 \nonumber\\
 D_+(\omega)
 &:=\frac{4\omega^2-3}
 {8(\omega^2-1)^{3/2}}.
 \label{eq:V48-CDplus}
\end{align}

\begin{lemma}[Exact stable-boundary trace]
\label{lem:V48-boundary-trace}
The upper stable boundary values are
\begin{align}
 g_\mu(i\omega+0,1)
 &=1+D_-(\omega)A_\mu(\omega)
   +iD_-(\omega)B_\mu(\omega),
 &&0\le\omega<1,
 \label{eq:V48-low-trace}\\
 g_\mu(i\omega+0,1)
 &=1+D_+(\omega)C_\mu(\omega)>1,
 &&\omega>1.
 \label{eq:V48-high-trace}
\end{align}
In particular, with
\begin{equation}
 \omega_0:=\frac{\sqrt3}{2},
 \qquad
 G_\mu:=\int\sqrt{1-v^2}\,\dd\mu(v)>0,
 \label{eq:V48-omega0-G}
\end{equation}
one has
\begin{align}
 g_\mu(0,1)&=1+\frac38\int v\,\dd\mu(v)>1,
 \nonumber\\
 g_\mu(i\omega_0+0,1)&=1,
 \nonumber\\
 g_\mu(i\omega+0,1)
 &=1-\frac{iG_\mu}{8(1-\omega^2)^{3/2}}
   +o\!\left((1-\omega^2)^{-3/2}\right)
 \quad(\omega\uparrow1).
 \label{eq:V48-endpoint-trace}
\end{align}
The lower boundary values are the complex conjugates.
\end{lemma}

\begin{proof}
For \(0\le\omega<1\), the stable bulk root is
\(\lambda=\sqrt{1-\omega^2}>0\).  A bath root is positive real when
\(v>\omega\), positive imaginary when \(v<\omega\), and zero when
\(v=\omega\).  Thus
\[
 {\mathfrak z}_\mu(i\omega+0,1)
 =A_\mu(\omega)+iB_\mu(\omega).
\]
Direct substitution in the V43 response gives
\(d(i\omega+0,1)=D_-(\omega)\), proving
\eqref{eq:V48-low-trace}.

For \(\omega>1\),
\[
 \lambda=i\sqrt{\omega^2-1},
 \qquad
 {\mathfrak z}_\mu=iC_\mu(\omega),
 \qquad
 d=-iD_+(\omega).
\]
This proves \eqref{eq:V48-high-trace}.  The first two endpoint identities
follow by setting \(\omega=0\) and \(\omega=\omega_0\).  As
\(\omega\uparrow1\), \(A_\mu\to0\), \(B_\mu\to G_\mu\), and
\[
 D_-(\omega)
 =-\frac1{8(1-\omega^2)^{3/2}}
  +O\!\left((1-\omega^2)^{-1/2}\right),
\]
which proves the last line.  Real coefficients and the retarded branch
give conjugation symmetry.
\end{proof}

\subsection{The only possible imaginary-axis zero}
\label{subsec:V48-boundary-zero}

Let
\begin{equation}
 v_*:=\inf\operatorname{supp}\mu.
 \label{eq:V48-vstar}
\end{equation}

\begin{lemma}[Boundary-zero classification]
\label{lem:V48-boundary-zero}
The factor \(g_\mu\) has no imaginary-axis zero above gravitational
glancing, below \(\omega_0\), or at a frequency where
\(B_\mu(\omega)>0\).  A positive imaginary-axis zero exists exactly when
there is
\begin{equation}
 \omega\in(\omega_0,v_*]
 \quad\hbox{such that}\quad
 1+D_-(\omega)A_\mu(\omega)=0.
 \label{eq:V48-boundary-equation}
\end{equation}
The negative-frequency zero is its conjugate.  In particular, if
\(v_*\le\omega_0\), no boundary zero is possible.
\end{lemma}

\begin{proof}
Equation \eqref{eq:V48-high-trace} excludes \(\omega>1\).  For
\(0\le\omega<1\), a zero in \eqref{eq:V48-low-trace} requires
\[
 D_-(\omega)B_\mu(\omega)=0.
\]
At \(D_-=0\), the factor equals one.  If \(D_->0\), its real part is at
least one.  It remains \(D_-<0\), equivalently
\(\omega>\omega_0\), together with \(B_\mu(\omega)=0\).  Positivity of
the measure gives
\[
 B_\mu(\omega)=0
 \quad\Longleftrightarrow\quad
 \omega\le v_*.
\]
On that interval the factor is real and its vanishing is precisely
\eqref{eq:V48-boundary-equation}.  Conjugation gives the negative
frequency statement.
\end{proof}

\subsection{The exact Nyquist index}
\label{subsec:V48-index}

\begin{theorem}[Boundary zero or exactly two open roots]
\label{thm:V48-index-two}
Let \(\mu\) satisfy \eqref{eq:V48-support}.  Exactly one of the
following cases occurs:
\begin{enumerate}
 \item \(g_\mu\) has at least one conjugate pair of imaginary-axis zeros
 described by Lemma~\ref{lem:V48-boundary-zero}; or
 \item \(g_\mu\) has no imaginary-axis zero and has exactly two zeros,
 counted with multiplicity, in \({\rm Re}\,s>0\).
\end{enumerate}
In the second case the two zeros are nonreal and form a conjugate pair.
Consequently \(g_\mu\) never has a uniform inverse on the closed stable
frequency half-plane.
\end{theorem}

\begin{proof}
If the boundary equation holds, the first alternative is proved.  Assume
from now on that there is no boundary zero.

The function \({\mathfrak z}_\mu\), and hence \(g_\mu\), is analytic in
\({\rm Re}\,s>0\).  The only singularities relevant to its closed
boundary trace are the square-root thresholds; \(s=\pm i\) are the
gravitational glancing points where \(d\) diverges.  The compact support
in \((0,1)\) gives uniform dominated convergence on compact subsets and
at every nonsingular boundary point.

First traverse the positive imaginary boundary upward.  On
\(0\le\omega\le\omega_0\), equation \eqref{eq:V48-low-trace} lies in the
closed right half-plane and returns to \(1\) at \(\omega_0\), so its net
argument change is zero.  On
\(\omega_0<\omega<1\), the factor remains positive real while
\(B_\mu=0\), because a real sign change would be the excluded boundary
zero.  Once \(B_\mu>0\), its imaginary part is strictly negative and it
remains in the open lower half-plane.  By
\eqref{eq:V48-endpoint-trace}, its argument tends to \(-\pi/2\) as
\(\omega\uparrow1\).

Resolve gravitational glancing by the right-hand semicircle
\begin{equation}
 s=i+\epsilon e^{i\varphi},
 \qquad -\frac\pi2\le\varphi\le\frac\pi2.
 \label{eq:V48-glancing-arc}
\end{equation}
Uniformly on this arc,
\begin{align}
 \lambda
 &=(2i\epsilon e^{i\varphi})^{1/2}(1+o(1))
 =\sqrt{2\epsilon}\,
   e^{i(\pi/4+\varphi/2)}(1+o(1)),
 \nonumber\\
 {\mathfrak z}_\mu&=iG_\mu+o(1),
 \nonumber\\
 g_\mu
 &=-\frac{iG_\mu}{8\lambda^3}(1+o(1)).
 \label{eq:V48-glancing-asymptotic}
\end{align}
The continuous argument of the leading term changes from \(-\pi/2\) to
\(-2\pi\).  The glancing arc therefore contributes \(-3\pi/2\).
Above glancing, \eqref{eq:V48-high-trace} is positive real.  Its
continuous argument stays at \(-2\pi\).  Hence the complete indented
positive half-axis path \(\gamma_+\), oriented from \(0\) to \(+i\infty\),
satisfies
\begin{equation}
 \Delta_{\gamma_+}\arg g_\mu=-2\pi.
 \label{eq:V48-upper-index}
\end{equation}

For \(R>1\) and \(\eta>0\), use the positively oriented boundary of
the truncated half-disk with vertical diameter \({\rm Re}\,s=\eta\),
and then take \(R\to\infty\) and \(\eta\downarrow0\).  The large arc
contributes zero because, uniformly there,
\begin{equation}
 g_\mu(s,1)\longrightarrow
 1+\frac{\mu((0,1))}{2}>0.
 \label{eq:V48-infinity}
\end{equation}
Away from \(s=\pm i\), dominated convergence and the assumed absence
of a boundary zero make the shifted vertical trace converge to the
boundary curves above.  Near \(s=\pm i\), it converges to the right-hand
glancing arcs used in \eqref{eq:V48-glancing-arc}.  Thus the limiting
vertical boundary is oriented downward: the reverse of \(\gamma_+\)
contributes \(+2\pi\), and conjugation gives a second \(+2\pi\) on the
path from \(0\) to \(-i\infty\).  The complete threshold interval adds
no separate indentation or winding; its nonzero trace is already confined
to the half-planes used above.  The total contour change is therefore
\begin{equation}
 \Delta_{\partial{\cal D}}\arg g_\mu=4\pi.
 \label{eq:V48-total-index}
\end{equation}

There are no poles in the open right half-plane.  The argument principle
applied first at finite radius and positive shift, then in the limits,
therefore gives exactly two open zeros counted with multiplicity.
For real \(s>0\), all roots, \(d\), and the measure are positive, so
\[
 g_\mu(s,1)>1.
\]
Neither open zero is real.  Conjugation symmetry makes them a nonreal
pair.
\end{proof}

\begin{remark}[Why the fractional glancing turn is decisive]
\label{rem:V48-fractional-turn}
The index does not come from a polynomial degree.  It comes from the
\(\lambda^{-3}\) longitudinal response at gravitational glancing while
the positive bath impedance approaches \(iG_\mu\ne0\).  The local
behavior \(-iG_\mu/(8\lambda^3)\), together with the half-turn of
\(\lambda\), creates the \(3\pi/2\) rotation.  Adding positive bath
speeds changes \(G_\mu\) but cannot change that exponent or its sign.
\end{remark}

\subsection{Finite mixtures and the complete covariant negative no-go}
\label{subsec:V48-covariant-corollary}

\begin{corollary}[Every finite positive speed mixture fails]
\label{cor:V48-finite-no-go}
Let
\begin{equation}
 {\mathfrak z}(s,r)
 =\sum_{\ell=1}^L a_\ell
 \sqrt{s^2+v_\ell^2r^2},
 \qquad
 a_\ell>0,\quad 0<v_\ell<1.
 \label{eq:V48-finite-mixture}
\end{equation}
Then \(1+d{\mathfrak z}\) has either an imaginary-axis zero or exactly two
open right-half-plane zeros.  This includes every two-speed longitudinal
factor left by V47.
\end{corollary}

\begin{proof}
Apply Theorem~\ref{thm:V48-index-two} to the atomic measure
\(\mu=\sum_\ell a_\ell\delta_{v_\ell}\).
\end{proof}

\begin{corollary}[Complete local covariant negative-orientation no-go]
\label{cor:V48-covariant-negative-no-go}
Consider any finite direct sum of the local positive half-line baths of
V47, or a compactly supported positive speed measure, and impose
tangential rotational covariance.  If its collective bath incidence is
full at gravitational glancing, the negative-orientation Einstein
effective matrix has an open or closed-boundary determinant zero.
Therefore no member of this class has a uniform stable inverse.
\end{corollary}

\begin{proof}
Theorem~\ref{thm:V47-covariant-commutant} makes the physical-port
impedance
\begin{equation}
 Z(s,r)
 =\operatorname{diag}
 \bigl({\mathfrak z}_\mu(s,r)I_2,
       {\mathfrak z}_\nu(s,r)\bigr)
 \label{eq:V48-covariant-Z}
\end{equation}
for nonnegative vector and scalar speed measures \(\mu,\nu\).  In the
direct-integral case, covariance commutes with multiplication by \(v^2\),
so the same planar commutant classification holds almost everywhere in
the speed variable.  Full glancing incidence forces \(\mu\ne0\) and
\(\nu\ne0\).  In the response basis the negative determinant factors as
\begin{equation}
 F_-=
 \bigl(1+d{\mathfrak z}_\mu\bigr)
 \bigl(1+b{\mathfrak z}_\mu\bigr)
 \bigl(1+b{\mathfrak z}_\nu\bigr).
 \label{eq:V48-negative-factorization}
\end{equation}
The first factor is covered by
Theorem~\ref{thm:V48-index-two}.  If collective incidence were not full,
Theorem~\ref{thm:V47-collective-rank} would already give the glancing
zero.  These cases exhaust the class.
\end{proof}

\begin{firewall}[Scope of the V48 closure]
\label{fire:V48-scope}
Corollary~\ref{cor:V48-covariant-negative-no-go} closes positive sums and
positive speed measures of the V47 square-root half-line bath in the
negative Green orientation.  It does not cover derivative boundary
couplings, gyroscopic auxiliary dynamics, indefinite internal energies,
or angularly varying matrix eigenspaces generated by additional local
tensor fields.  It also does not decide the strong-gain positive
orientation left by Theorem~\ref{thm:V47-positive-chamber}.
\end{firewall}

\subsection{Finite certificate and V49 upstream row}
\label{subsec:V48-certificate}

A deterministic audit examined six hundred random two-speed positive
mixtures.  In every case without a low boundary zero, the branch-filtered
degree-six card of V47 returned exactly two open roots; the largest
normalized unsquared residual was \(1.8\times10^{-10}\).  One hundred
additional high-speed, strong-coupling samples targeted
\eqref{eq:V48-boundary-equation}; 95 produced the predicted boundary
crossing.

Twenty random positive mixtures with between two and seven speeds were
evaluated on a right-shifted vertical contour with logarithmic resolution
at gravitational glancing.  The largest error in the upper-path phase
change \(-2\pi\) was \(4.4\times10^{-11}\), and the smallest sampled
contour modulus was \(0.7529\).  These computations audit the branch and
orientation bookkeeping; the proof of Theorem~\ref{thm:V48-index-two} is
the exact argument-principle calculation.

V49 must now test the only local covariant square-root branch not closed:
the positive orientation inside the strict chamber of
Theorem~\ref{thm:V47-positive-chamber}.  In the measure notation its
factors are
\begin{equation}
 F_+
 =\bigl(1-d{\mathfrak z}_\mu\bigr)
  \bigl(1-b{\mathfrak z}_\mu\bigr)
  \bigl(1-b{\mathfrak z}_\nu\bigr).
 \label{eq:V48-positive-target}
\end{equation}
The acquisition order is:
\begin{enumerate}
 \item compute the separate Nyquist index of each scalar factor under the
 V47 strict endpoint inequalities;
 \item decide whether one positive speed measure makes all three factors
 root-free with a common conic floor; and
 \item if the index still forces a root, move upstream to a local
 derivative or gyroscopic bath coupling rather than adding further
 square-root speeds.
\end{enumerate}

\subsection{V48 ledger}
\label{subsec:V48-ledger}

\paragraph{New conclusions proved in V48.}
Lemma~\ref{lem:V48-boundary-trace} gives the exact unsquared boundary
curve for a positive speed measure.
Lemma~\ref{lem:V48-boundary-zero} classifies every possible imaginary-axis
zero.  Theorem~\ref{thm:V48-index-two} proves the universal dichotomy:
a boundary zero, or otherwise exactly two open roots.
Corollaries~\ref{cor:V48-finite-no-go,cor:V48-covariant-negative-no-go}
close every finite-speed and compact-measure local covariant
negative-orientation square-root bath.

\paragraph{Imported conclusions not repeated.}
V43 supplies the longitudinal response \(d\).  V47 supplies the local
direct-sum action, the covariance commutant, and the exact determinant
factorization.  V48 begins at the remaining scalar factor and proves its
root count by an unsquared contour argument.

\paragraph{Honest continuum status.}
The negative orientation is now excluded for the complete local
rotationally covariant positive square-root mixture class.  The strong
positive orientation, more general local auxiliary dynamics, nonlinear
covariant stress/Bianchi closure, ghosts, and the coupled Standard-Model
sectors remain open.  The full Standard Model--Einstein continuum is not
proved.

% END APPENDED VERSION V48 MATERIAL



% BEGIN APPENDED VERSION V49 MATERIAL

\clearpage
\section{V49: exact strong-gain stability for the positive bath
and an explicit local covariant port construction}
\label{sec:V49-positive-construction}

\subsection{Purpose and exact advance}
\label{subsec:V49-purpose}

V48 closes every local rotationally covariant positive square-root mixture
in the negative Green orientation.  The remaining square-root branch is
the positive orientation inside the strict endpoint chamber isolated in
V47.  V49 decides that branch at the port-Schur level.

There is a genuine difference between weak and strong positive feedback.
Below the endpoint thresholds, a real, high-boundary, or low-cone zero is
forced.  Above all thresholds, the scalar Nyquist index is zero.  Exact
moment and first-cone inequalities are both necessary and sufficient for
the longitudinal and transverse factors to be zero-free on the closed
stable frequency half-plane.

The inequalities are nonempty.  Two local baths at speeds \(1/2\) and
\(9/10\), with constant rotationally covariant positive Grams, give an
explicit zero-free three-port Schur matrix with a uniform conic floor.
This is a port-level construction.  The direct full Einstein operator at
gravitational glancing and its normalized inverse estimate remain to be
proved before it can be promoted to a continuum closure.

\subsection{Profiles and monotonicity}
\label{subsec:V49-profiles}

Retain the positive speed measure \(\mu\) and the functions
\({\mathfrak z}_\mu,d\) from V48.  Also put
\begin{equation}
 b(s,r):=\frac1{2\lambda(s,r)}.
 \label{eq:V49-b}
\end{equation}
The positive longitudinal and transverse scalar factors are
\begin{equation}
 \ell_\mu(s,r):=1-d(s,r){\mathfrak z}_\mu(s,r),
 \qquad
 t_\mu(s,r):=1-b(s,r){\mathfrak z}_\mu(s,r).
 \label{eq:V49-factors}
\end{equation}

Define the mass, first speed moment, and first occupied speed by
\begin{equation}
 m_0(\mu):=\mu((0,1)),
 \qquad
 m_1(\mu):=\int v\,\dd\mu(v),
 \qquad
 v_\mu:=\inf\operatorname{supp}\mu.
 \label{eq:V49-moments}
\end{equation}
At the first cone set
\begin{align}
 A_\mu^*
 &:=\int_{(v_\mu,1)}
       \sqrt{v^2-v_\mu^2}\,\dd\mu(v),
 \label{eq:V49-Astar}\\
 T_\mu^*
 &:=\frac{A_\mu^*}{2\sqrt{1-v_\mu^2}},
 \label{eq:V49-Tstar}\\
 L_\mu^*
 &:=\frac{3-4v_\mu^2}
 {8(1-v_\mu^2)^{3/2}}\,A_\mu^*.
 \label{eq:V49-Lstar}
\end{align}

\begin{lemma}[The five monotone response profiles]
\label{lem:V49-monotonicity}
Let \(c=v^2\in(0,1)\).  With \(x\ge0\), the real-frequency profiles
\begin{align}
 q_T^{\rm R}(x;c)
 &:=\frac{\sqrt{x+c}}{2\sqrt{x+1}},
 \nonumber\\
 q_L^{\rm R}(x;c)
 &:=\frac{(4x+3)\sqrt{x+c}}{8(x+1)^{3/2}}
 \label{eq:V49-real-profiles}
\end{align}
are strictly increasing from \(v/2,3v/8\), respectively, to \(1/2\).
For \(x>1\), the high-boundary profiles
\begin{align}
 q_T^{\rm H}(x;c)
 &:=\frac{\sqrt{x-c}}{2\sqrt{x-1}},
 \nonumber\\
 q_L^{\rm H}(x;c)
 &:=\frac{(4x-3)\sqrt{x-c}}{8(x-1)^{3/2}}
 \label{eq:V49-high-profiles}
\end{align}
are strictly decreasing from \(+\infty\) to \(1/2\).
Finally,
\begin{equation}
 q_L^{\rm low}(x;c)
 :=\frac{(3-4x)\sqrt{c-x}}{8(1-x)^{3/2}}
 \label{eq:V49-low-profile}
\end{equation}
is strictly decreasing for
\(0\le x<\min(c,3/4)\).
\end{lemma}

\begin{proof}
The squares of \(q_T^{\rm R}\) and \(q_T^{\rm H}\), apart from the
positive constant \(1/4\), are
\[
 \frac{x+c}{x+1},
 \qquad
 \frac{x-c}{x-1}.
\]
Their derivative numerators are \(1-c>0\) and \(c-1<0\).

For the remaining three profiles, logarithmic differentiation and
multiplication by the positive denominators leave the respective affine
numerators
\begin{align}
 N_{\rm R}(x;c)&=(6-4c)x+3-c>0,
 \nonumber\\
 N_{\rm H}(x;c)&=(4c-6)x+3-c<0
 \quad(x>1),
 \nonumber\\
 N_{\rm low}(x;c)&=(6-4c)x+c-3.
 \label{eq:V49-derivative-numerators}
\end{align}
The first two signs are immediate.  For the last, if \(c\le3/4\), its
maximum on \(0\le x\le c\) is
\[
 N_{\rm low}(c;c)=-(4c-3)(c-1)\le0.
\]
If \(c\ge3/4\), its maximum on \(0\le x\le3/4\) is
\[
 N_{\rm low}(3/4;c)=\frac32-2c\le0,
\]
with strict inequality in the open domain.  The endpoint values and
limits follow directly.
\end{proof}

\subsection{The transverse positive factor}
\label{subsec:V49-transverse}

\begin{theorem}[Exact transverse strong-gain criterion]
\label{thm:V49-transverse}
Let \(\mu\) satisfy the V48 compact positive support hypothesis.  The
factor \(t_\mu\) has no zero on
\({\rm Re}\,s\ge0\), \((s,r)\ne(0,0)\), if and only if
\begin{equation}
 \boxed{\quad
 m_0(\mu)>2,\qquad
 m_1(\mu)>2,\qquad
 T_\mu^*>1.
 \quad}
 \label{eq:V49-transverse-criterion}
\end{equation}
When these conditions hold, there is a uniform homogeneous floor
\begin{equation}
 \inf_{\substack{{\rm Re}\,s\ge0,\ r\ge0\\|s|^2+r^2=1}}
 |t_\mu(s,r)|>0.
 \label{eq:V49-transverse-floor}
\end{equation}
\end{theorem}

\begin{proof}
Set \(r=1\).  On the high imaginary boundary,
\[
 t_\mu(i\omega+0,1)
 =1-\int q_T^{\rm H}(\omega^2;v^2)\,\dd\mu(v),
 \qquad \omega>1.
\]
Lemma~\ref{lem:V49-monotonicity} shows that this real function increases
from \(-\infty\) to \(1-m_0/2\).  Thus \(m_0<2\) gives a high-boundary
zero, while \(m_0=2\) gives the homogeneous \(r=0\) zero.  Avoiding both
requires \(m_0>2\).

On the positive real axis,
\[
 t_\mu(s,1)
 =1-\int q_T^{\rm R}(s^2;v^2)\,\dd\mu(v).
\]
It decreases from \(1-m_1/2\) to \(1-m_0/2<0\).  Hence \(m_1<2\) gives an
open real zero and \(m_1=2\) gives the \(s=0\) zero.  Avoidance requires
\(m_1>2\).

On \(0\le\omega\le v_\mu\), every bath root is real and
\[
 t_\mu(i\omega+0,1)
 =1-\frac{A_\mu(\omega)}{2\sqrt{1-\omega^2}}.
\]
Each integrand is strictly decreasing in \(\omega^2\).  The subtracted
quantity therefore decreases from \(m_1/2>1\) to \(T_\mu^*\).  A low
boundary zero is absent exactly when \(T_\mu^*>1\).

These arguments prove necessity.  Assume all three strict conditions.
They also prove that the factor is negative on the real axis, the high
boundary, and the low boundary up to the first cone.  After the first
occupied cone and below gravitational glancing, its imaginary part is
strictly negative:
\[
 t_\mu(i\omega+0,1)
 =1-\frac{A_\mu(\omega)}{2\sqrt{1-\omega^2}}
 -i\frac{B_\mu(\omega)}{2\sqrt{1-\omega^2}}.
\]
Thus there is no boundary zero.

For the open root count, start the continuous upper-boundary argument at
\(-\pi\).  After the first cone the curve remains in the lower half-plane
and tends to argument \(-\pi/2\) as \(\omega\uparrow1\).  On the
right-hand glancing arc,
\begin{equation}
 t_\mu=-\frac{iG_\mu}{2\lambda}(1+o(1)),
 \label{eq:V49-T-glancing}
\end{equation}
so its argument returns from \(-\pi/2\) to \(-\pi\).  The high boundary
and the large semicircle end at \(1-m_0/2<0\).  Hence the upper path has
zero net argument change.  Conjugation gives zero total Nyquist winding.
There are no open poles, so the argument principle gives no open zeros.

Finally normalize by \(|s|^2+r^2=1\).  If the asserted infimum were zero,
a convergent minimizing sequence would limit either to a finite zero or
to gravitational glancing.  The first is excluded, while
\eqref{eq:V49-T-glancing} diverges at glancing.  This proves the uniform
floor.
\end{proof}

\subsection{The longitudinal positive factor}
\label{subsec:V49-longitudinal}

\begin{theorem}[Exact longitudinal strong-gain criterion]
\label{thm:V49-longitudinal}
Let \(\mu\) satisfy the same hypothesis and put
\(\omega_0=\sqrt3/2\).  The factor \(\ell_\mu\) has no zero on the closed
stable frequency half-plane if and only if
\begin{equation}
 \boxed{\quad
 m_0(\mu)>2,\qquad
 m_1(\mu)>\frac83,\qquad
 v_\mu<\omega_0,\qquad
 L_\mu^*>1.
 \quad}
 \label{eq:V49-longitudinal-criterion}
\end{equation}
Under these conditions,
\begin{equation}
 \inf_{\substack{{\rm Re}\,s\ge0,\ r\ge0\\|s|^2+r^2=1}}
 |\ell_\mu(s,r)|>0.
 \label{eq:V49-longitudinal-floor}
\end{equation}
\end{theorem}

\begin{proof}
On the high boundary,
\[
 \ell_\mu(i\omega+0,1)
 =1-\int q_L^{\rm H}(\omega^2;v^2)\,\dd\mu(v).
\]
This increases from \(-\infty\) to \(1-m_0/2\), so absence of a high or
homogeneous zero requires \(m_0>2\).  On the real axis,
\[
 \ell_\mu(s,1)
 =1-\int q_L^{\rm R}(s^2;v^2)\,\dd\mu(v),
\]
which decreases from \(1-3m_1/8\) to \(1-m_0/2<0\).  Absence of a real
or zero-frequency root therefore requires \(m_1>8/3\).

Suppose \(v_\mu\ge\omega_0\).  All bath roots remain real from
\(\omega=0\) through \(\omega_0\), while
\[
 \ell_\mu(0,1)=1-\frac38m_1<0,
 \qquad
 \ell_\mu(i\omega_0+0,1)=1.
\]
The intermediate value theorem gives a boundary zero.  Hence
\(v_\mu<\omega_0\) is necessary.

For \(0\le\omega\le v_\mu<\omega_0\),
\begin{equation}
 \ell_\mu(i\omega+0,1)
 =1-D_-(\omega)A_\mu(\omega).
 \label{eq:V49-low-L}
\end{equation}
Lemma~\ref{lem:V49-monotonicity} makes
\(D_-A_\mu\) strictly decreasing from \(3m_1/8>1\) to \(L_\mu^*\).
Thus the first-cone path is zero-free exactly when \(L_\mu^*>1\).
This completes necessity.

Assume all four inequalities.  The preceding monotonicity shows that
\(\ell_\mu\) is negative on the real axis, high boundary, and first-cone
segment.  Between the first cone and \(\omega_0\),
\[
 \ell_\mu
 =1-D_-A_\mu-iD_-B_\mu
\]
lies in the lower half-plane.  It reaches \(1\) at \(\omega_0\).  For
\(\omega_0<\omega<1\), \(D_-<0\), so
\[
 \ell_\mu
 =1+|D_-|A_\mu+i|D_-|B_\mu
\]
lies in the open first quadrant and tends to argument \(\pi/2\).
At glancing,
\begin{equation}
 \ell_\mu=\frac{iG_\mu}{8\lambda^3}(1+o(1)).
 \label{eq:V49-L-glancing}
\end{equation}
The right-hand arc changes the argument from \(\pi/2\) to \(-\pi\), a
change of \(-3\pi/2\).  Before the arc, the path from the initial negative
real value to \(1\) through the lower half-plane and then to \(+i\infty\)
contributes \(+3\pi/2\).  The upper path therefore has net change zero.
The high boundary and the large arc remain negative because \(m_0>2\).
Conjugation gives total Nyquist index zero.  The argument principle proves
that there is no open zero.

The compactness argument used in
Theorem~\ref{thm:V49-transverse}, now with the divergence
\eqref{eq:V49-L-glancing}, proves
\eqref{eq:V49-longitudinal-floor}.
\end{proof}

\begin{remark}[The longitudinal cone condition dominates the vector
transverse condition]
\label{rem:V49-dominance}
When \(v_\mu<\omega_0\),
\begin{equation}
 L_\mu^*
 =h(v_\mu)T_\mu^*,
 \qquad
 h(v_\mu)=\frac{3-4v_\mu^2}{4(1-v_\mu^2)}\in(0,1).
 \label{eq:V49-LT-relation}
\end{equation}
Thus \(L_\mu^*>1\) implies \(T_\mu^*>1\).  Also
\(m_1>8/3\) implies \(m_1>2\).  A vector speed measure satisfying the
longitudinal criterion automatically stabilizes its transverse vector
factor.
\end{remark}

\subsection{An explicit local positive-orientation port}
\label{subsec:V49-explicit}

Take two speeds
\begin{equation}
 v_1=\frac12,\qquad v_2=\frac9{10},
 \label{eq:V49-speeds}
\end{equation}
and the vector and scalar measures
\begin{equation}
 \mu=\delta_{1/2}+5\delta_{9/10},
 \qquad
 \nu=\delta_{1/2}+3\delta_{9/10}.
 \label{eq:V49-measures}
\end{equation}
Equivalently, the two local V47 bath Grams in the physical port frame are
\begin{equation}
 M_1=I_3,\qquad
 M_2=\operatorname{diag}(5,5,3).
 \label{eq:V49-Grams}
\end{equation}

\begin{theorem}[Explicit zero-free positive port-Schur matrix]
\label{thm:V49-explicit-port}
The action \eqref{eq:V47-action} with
\eqref{eq:V49-speeds,eq:V49-Grams} is local, rotationally covariant, and
has positive energy.  For the positive Green orientation its response
Schur matrix is
\begin{equation}
 {\cal S}_+
 =\operatorname{diag}
 \bigl(\ell_\mu,t_\mu,t_\nu\bigr)
 \label{eq:V49-Schur}
\end{equation}
in the \((v_L,v_T,\phi)\) basis, and there is \(c_{\rm port}>0\) such that
\begin{equation}
 s_{\min}({\cal S}_+(s,k))\ge c_{\rm port}
 \label{eq:V49-port-floor}
\end{equation}
on the normalized closed stable cone away from the zero covector.
\end{theorem}

\begin{proof}
For \(\mu\),
\begin{equation}
 m_0=6,\qquad m_1=5,\qquad v_\mu=\frac12,
 \qquad A_\mu^*=\sqrt{14}.
 \label{eq:V49-mu-data}
\end{equation}
Consequently
\begin{equation}
 L_\mu^*=\frac{2\sqrt{42}}9>1,
 \qquad
 T_\mu^*=\frac{\sqrt{42}}3>1.
 \label{eq:V49-mu-cones}
\end{equation}
Theorem~\ref{thm:V49-longitudinal} and
Remark~\ref{rem:V49-dominance} therefore give uniform floors for
\(\ell_\mu\) and \(t_\mu\).

For \(\nu\),
\begin{equation}
 m_0=4,\qquad m_1=\frac{16}{5},\qquad
 T_\nu^*=\frac{\sqrt{42}}5>1.
 \label{eq:V49-nu-data}
\end{equation}
Theorem~\ref{thm:V49-transverse} gives a uniform floor for \(t_\nu\).
The matrices \eqref{eq:V49-Grams} have the covariant commutant form of
Theorem~\ref{thm:V47-covariant-commutant}; V47 then supplies locality,
positive energy, and the diagonal factorization.  Taking the minimum of
the three scalar floors proves \eqref{eq:V49-port-floor}.
\end{proof}

\begin{firewall}[Port stability is not yet the full Einstein estimate]
\label{fire:V49-port-scope}
Theorem~\ref{thm:V49-explicit-port} proves a uniform inverse for the exact
three-port Schur matrix wherever the V43 response reduction is defined.
It does not yet prove the normalized inverse estimate for the unreduced
Einstein-plus-bath operator at \(\lambda=0\).  The determinant lemma uses
\(L_0^{-1}\), which is singular there, and the divergence of the Schur
factors alone does not quantify the cancellation in the reconstructed
five-row metric variables.  That direct glancing calculation is required
before this port construction can be promoted to a linear continuum
theorem.
\end{firewall}

\subsection{Finite certificate and V50 handoff}
\label{subsec:V49-certificate}

For the explicit measures \eqref{eq:V49-measures}, the exact criterion
values were
\[
 (m_0,m_1,v_\mu,L_\mu^*,T_\mu^*)
 =(6,5,0.5,1.4401646,2.1602469)
\]
and
\[
 (m_0,m_1,v_\nu,T_\nu^*)
 =(4,3.2,0.5,1.2961481).
\]
A right-shifted Nyquist audit gave upper-path phase errors below
\(3.4\times10^{-11}\).  A normalized complex-frequency grid found a
minimum of the three factor moduli \(0.17512\); this sampled number audits
the exact zero-index proof and is not substituted for the compactness
argument establishing \(c_{\rm port}>0\).  Ten thousand random points
also verified the three derivative signs in
\eqref{eq:V49-derivative-numerators}.

V50 must begin with the scheduled audit of the newest NS and YM notes.
It should then return to the original, unreduced V43/V47 matrices and:
\begin{enumerate}
 \item insert the explicit Grams \eqref{eq:V49-Grams} before using
 \(L_0^{-1}\);
 \item compute the exact full operator and kernel map at
 \(\lambda=0\);
 \item prove or refute a uniform normalized smallest-singular-value bound
 through gravitational glancing; and
 \item retain the local bath variables so the estimate controls the
 conservative continuum realization, rather than only its retarded
 impedance.
\end{enumerate}
If the direct estimate fails, V50 should identify the missing source row
and move upstream to the boundary coupling, rather than weakening the
port norm.

\subsection{V49 ledger}
\label{subsec:V49-ledger}

\paragraph{New conclusions proved in V49.}
Lemma~\ref{lem:V49-monotonicity} proves all response-profile signs.
Theorems~\ref{thm:V49-transverse,thm:V49-longitudinal} give exact
necessary-and-sufficient strong-gain criteria and uniform scalar floors.
Theorem~\ref{thm:V49-explicit-port} supplies an explicit local, positive,
rotationally covariant two-speed bath with a uniformly invertible
positive-orientation port-Schur matrix.

\paragraph{Imported conclusions not repeated.}
V47 supplies the direct-sum action, covariance commutant, and exact
three-port factorization.  V48 supplies the speed-measure boundary traces
and contour discipline.  V49 changes the Green sign, proves the new
zero-index criteria, and constructs a parameter point satisfying them.

\paragraph{Honest continuum status.}
A stable passive local port has now been constructed for the positive
orientation.  The full unreduced Einstein glancing estimate, dynamical
bath norm, nonlinear covariant stress/Bianchi closure, ghosts, and the
coupled Standard-Model sectors remain open.  The full Standard
Model--Einstein continuum is not proved.

% END APPENDED VERSION V49 MATERIAL


% BEGIN APPENDED VERSION V50 MATERIAL

\clearpage
\section{V50: direct Einstein--bath glancing inverse, weighted outgoing
continuum, and the total-energy threshold obstruction}
\label{sec:V50-direct-bath}

\subsection{Purpose, acquisition order, and the scheduled companion audit}
\label{subsec:V50-purpose-audit}

V49 constructed an explicit local, positive-energy, rotationally covariant
two-speed bath whose positive-orientation three-port Schur matrix is
uniformly invertible.  Its proof deliberately stopped before gravitational
glancing because the determinant lemma used \(L_0^{-1}\), which has a
fifth-order pole at \(\lambda=0\).  V50 removes that restriction.  It inserts
the two bath Grams into the original five metric rows, computes the limiting
matrix without division by \(\lambda\), and proves a uniform least-singular-
value floor for the complete metric trace.

The local half-line fields are then restored.  They reveal a distinction
which the finite port determinant cannot see.  The outgoing limiting-
absorption graph is uniformly controlled in every polynomially weighted
local-energy norm with weight exponent greater than \(1/2\).  Its ordinary
unweighted total bath energy, however, diverges as the Laplace frequency
approaches the bath continuous spectrum.  Thus V50 proves the direct metric
floor and a full weighted outgoing trace estimate, and also proves that a
closed-cone estimate in the unweighted half-line energy norm is false.

As required at every fifth SM version, the newest completed companion
versions were read before this direction was fixed.  The frozen read-only
mathematical snapshot was
\begin{center}
\begin{tabular}{p{0.54\textwidth}p{0.38\textwidth}}
\texttt{Renewal\_NS\_Stability\_Research\_Notes\_V24.tex}
&\texttt{7ddd e763 9706 3317 76a9 4f50 3aaf b0cc}\newline
 \texttt{d44a 9c8b 9868 4548 8deb 6140 5cb7 c4df}\\[0.4em]
\texttt{Renewal\_Yang\_Mills\_Mass\_Gap\_Research\_Notes\_V29.tex}
&\texttt{3d9d 1d31 67a2 74a3 95e6 1fba 8a85 ee06}\newline
 \texttt{cac3 febd 56aa ac3c 0888 4ae6 6098 4cbb}
\end{tabular}
\end{center}
The NS file advanced from V23 to V24 during the audit; the displayed V24
file is the authoritative completed snapshot.  During final validation a file
named \texttt{Renewal\_NS\_Stability\_Research\_Notes\_V25.tex} appeared,
but it was byte-identical to the displayed V24 hash and contained no V25
section.  It therefore supplied no additional mathematical row.  The V24
Theorem~\texttt{thm:v24-overlap-separation} is an exact alternative for the
two \emph{assembled} compensated helicity profiles, and its later balance
shows that separate absolute estimates return to endpoint enstrophy and
palinstrophy.  The type-correct consequence here is to assemble the two bath
tractions at their common trace before estimating them and to keep the
unweighted energy debit visible.  No Navier--Stokes flux estimate is used in
the Einstein symbol.

The YM V29 note proves fixed-order Gaussian excess and every finite-window
restriction of its uniform fusion-cumulant inequality, while leaving the
single constant uniform in cumulant order, representation height, and
support size open.  The type-correct consequence here is that a finite port
matrix or a frequency grid cannot establish the continuum estimate.  The
proof below instead uses exact invertibility on the whole closed normalized
cone and compactness, and states separately the infinite-dimensional bath
norm for which uniformity fails.  No Wilson multiplier or Yang--Mills
cumulant bound is imported.

\begin{firewall}[Companion results remain in their own types]
\label{fire:V50-companion-types}
The NS assembled-profile identities and the YM fusion estimates do not
become estimates for a Lorentzian Einstein boundary symbol.  They determine
the audit discipline only: sum the complete common-trace source, preserve
signed versus absolute ledgers, and distinguish every finite restriction
from the uniform continuum norm.  All analytic assertions below are proved
directly for the V41/V47 operator.
\end{firewall}

\subsection{The explicit local bath before Schur reduction}
\label{subsec:V50-full-system}

Retain the V41 metric variables and port map
\begin{equation}
 x=(\tau,v_0,v_1,v_2,\phi)^T,
 \qquad Jx=(v_1,v_2,\phi)^T,
 \qquad C=H^{-1}J^T.
 \label{eq:V50-x-J-C}
\end{equation}
Thus
\begin{equation}
 C(f_1,f_2,f_3)^T
 =\left(-\frac32f_3,0,\frac12f_1,\frac12f_2,
 \frac12f_3\right)^T.
 \label{eq:V50-C-explicit}
\end{equation}
Use the V49 data
\begin{equation}
 v_1=\frac12,
 \qquad v_2=\frac9{10},
 \qquad
 M_1=I_3,
 \qquad M_2=\operatorname{diag}(5,5,3),
 \label{eq:V50-explicit-data}
\end{equation}
and the retarded roots
\begin{equation}
 \kappa_j(s,k):=\sqrt{s^2+v_j^2|k|^2},
 \qquad \operatorname{Re}\kappa_j>0
 \quad(\operatorname{Re}s>0).
 \label{eq:V50-kappa}
\end{equation}
Their boundary values are taken by limiting absorption from
\(\operatorname{Re}s>0\).  Put
\begin{equation}
 Z_*(s,k):=\kappa_1M_1+\kappa_2M_2,
 \qquad
 L_*(s,k):=L_0(s,k)-CZ_*(s,k)J.
 \label{eq:V50-Z-Lstar}
\end{equation}

The unreduced outgoing stable system with metric datum \(F\in\mathbb C^5\)
is
\begin{align}
 (-\partial_a^2+\kappa_j^2)q_j&=0,
 &q_j(0)&=Jx,
 \qquad j=1,2,
 \label{eq:V50-bath-ODE}\
 L_0x-C\sum_{j=1}^2\bigl(-M_j\partial_aq_j(0)\bigr)&=F.
 &&
 \label{eq:V50-metric-traction}
\end{align}
For \(\operatorname{Re}s>0\), the finite-energy/outgoing condition gives
\begin{equation}
 q_j(a)=e^{-\kappa_ja}Jx,
 \qquad
 -M_j\partial_aq_j(0)=\kappa_jM_jJx.
 \label{eq:V50-outgoing-solution}
\end{equation}
Hence \eqref{eq:V50-metric-traction} is exactly \(L_*x=F\).  This is an
elimination by solving the local half-line Euler equation; the fields
\(q_j\) remain part of the state below.

\begin{lemma}[Exact determinant normalization of the metric graph]
\label{lem:V50-det-L0}
For the V41 graph matrix,
\begin{equation}
 \det L_0(s,k)=\frac43\lambda^5,
 \qquad \lambda^2=s^2+|k|^2.
 \label{eq:V50-det-L0}
\end{equation}
\end{lemma}

\begin{proof}
The V41 rows of \(L_0\) are, in terms of the V37 rows,
\[
 2K,\qquad C_a,\qquad 2(K+C_n).
\]
On the two rows \(K,C_n\), the row transformation is

\[
 \begin{pmatrix}2&0\\2&2\end{pmatrix},
\]
whose determinant is four.  The three \(C_a\) rows are unchanged.  Multiply
the V37 identity \(\det\mathcal M_{CM}=\lambda^5/3\) by four.
\end{proof}

\subsection{Direct gravitational-glancing matrix and inverse}
\label{subsec:V50-glancing}

Let \(r=|k|>0\), rotate \(k\) to \((r,0)\), and write the metric variables
as \((\tau,v_0,v_L,v_T,\phi)\).  At either gravitational boundary point
put
\begin{equation}
 s=\sigma ir,
 \qquad \lambda=0,
 \qquad \sigma\in\{+1,-1\}.
 \label{eq:V50-glancing-frequency}
\end{equation}
The two retarded bath roots and their assembled impedance are
\begin{align}
 \kappa_1&=\sigma ir\frac{\sqrt3}{2},
 &\kappa_2&=\sigma ir\frac{\sqrt{19}}{10},
 \label{eq:V50-glancing-kappas}\
 Z_*&=\operatorname{diag}(\xi_v^\sigma,\xi_v^\sigma,
 \xi_\phi^\sigma),
 &\xi_v^\sigma&:=\sigma ir g_v,
 \qquad \xi_\phi^\sigma:=\sigma ir g_\phi,
 \label{eq:V50-glancing-Z}
\end{align}
where
\begin{equation}
 g_v:=\frac{\sqrt3+\sqrt{19}}2,
 \qquad
 g_\phi:=\frac{5\sqrt3+3\sqrt{19}}{10}.
 \label{eq:V50-gains}
\end{equation}

\begin{theorem}[Exact direct glancing inverse]
\label{thm:V50-direct-glancing}
At every frequency \eqref{eq:V50-glancing-frequency}, the unreduced
five-row metric matrix is
\begin{equation}
 L_*^\sigma=
 \begin{pmatrix}
 0&2s&-2ir&0&\frac32\xi_\phi^\sigma\\
 -s/6&0&0&0&-s/2\\
 -ir/6&0&-\xi_v^\sigma/2&0&-ir/2\\
 0&0&0&-\xi_v^\sigma/2&0\\
 0&0&0&0&-\xi_\phi^\sigma/2
 \end{pmatrix}.
 \label{eq:V50-direct-glancing-matrix}
\end{equation}
It is invertible, with
\begin{equation}
 \boxed{\quad
 \det L_*^\sigma
 =\frac{r^2(\xi_v^\sigma)^2\xi_\phi^\sigma}{24}
 =-\frac{\sigma i}{24}r^5g_v^2g_\phi\ne0.
 \quad}
 \label{eq:V50-direct-glancing-det}
\end{equation}
More explicitly, if the row datum is
\(f=(f_\tau,f_0,f_L,f_T,f_\phi)^T\), the solution of
\(L_*^\sigma x=f\) is
\begin{align}
 \phi&=-\frac{2f_\phi}{\xi_\phi^\sigma},
 &v_T&=-\frac{2f_T}{\xi_v^\sigma},
 \label{eq:V50-glancing-inverse-one}\
 \tau&=-\frac{6f_0}{s}+\frac{6f_\phi}{\xi_\phi^\sigma},
 &v_L&=-\frac2{\xi_v^\sigma}
 \left(f_L-\frac{ir}{s}f_0\right),
 \label{eq:V50-glancing-inverse-two}\
 v_0&=\frac{f_\tau+2irv_L+3f_\phi}{2s}.
 \label{eq:V50-glancing-inverse-three}
\end{align}
\end{theorem}

\begin{proof}
At \(\lambda=0\), the five V41 rows are
\[
 -A(d)x=
 \left(-2d^\sharp v,
 -d_a\left(\frac\tau6+\frac\phi2\right),0\right).
\]
For \(d=(s,ir,0)^T\), insert \eqref{eq:V50-C-explicit} and
\eqref{eq:V50-glancing-Z} into \(L_0-CZ_*J\).  This gives
\eqref{eq:V50-direct-glancing-matrix} without using \(L_0^{-1}\).

Expand its determinant first along the last row and then along the
transverse row.  The remaining three-by-three determinant is
\[
 \det\begin{pmatrix}
 0&2s&-2ir\\
 -s/6&0&0\\
 -ir/6&0&-\xi_v^\sigma/2
 \end{pmatrix}
 =-\frac{s^2\xi_v^\sigma}{6}
 =\frac{r^2\xi_v^\sigma}{6}.
\]
Multiplication by the two pivots
\(-\xi_v^\sigma/2\) and \(-\xi_\phi^\sigma/2\) proves the first equality in
\eqref{eq:V50-direct-glancing-det}; substitution of
\eqref{eq:V50-glancing-Z} proves the second.

For the inverse, the last and transverse rows give
\eqref{eq:V50-glancing-inverse-one}.  Put
\(\chi=\tau/6+\phi/2\).  The temporal row gives
\(\chi=-f_0/s\), and the longitudinal row then gives the formula for
\(v_L\).  Recovering \(\tau=6\chi-3\phi\) gives the first formula in
\eqref{eq:V50-glancing-inverse-two}.  The first row finally gives
\eqref{eq:V50-glancing-inverse-three}.  Every denominator is nonzero
because \(r>0\), \(s=\sigma ir\), and \(g_v,g_\phi>0\).
\end{proof}

\begin{remark}[The fifth-order pole cancels exactly]
\label{rem:V50-pole-cancellation}
For \(\lambda\ne0\), the V43 response and V49 factorization give
\begin{equation}
 \frac{\det L_*}{\det L_0}
 =(1-d\mathfrak z_\mu)(1-b\mathfrak z_\mu)
  (1-b\mathfrak z_\nu),
 \label{eq:V50-factorization}
\end{equation}
where
\[
 b=\frac1{2\lambda},
 \quad d=\frac{4\lambda^2-r^2}{8\lambda^3},
 \quad \mathfrak z_\mu=\kappa_1+5\kappa_2,
 \quad \mathfrak z_\nu=\kappa_1+3\kappa_2.
\]
Using \eqref{eq:V50-det-L0}, the leading \(\lambda^{-5}\) term on the
right has the finite product
\[
 \lim_{\lambda\to0}\det L_0
 \frac{r^2\mathfrak z_\mu^2\mathfrak z_\nu}{32\lambda^5}
 =\frac{r^2(\xi_v^\sigma)^2\xi_\phi^\sigma}{24},
\]
which is exactly \eqref{eq:V50-direct-glancing-det}.  Thus the divergent
port factors encode a finite nonzero direct determinant; they are not a
norm estimate by themselves.
\end{remark}

\subsection{Uniform closed-cone floor for all five metric variables}
\label{subsec:V50-uniform-metric-floor}

Let
\begin{equation}
 \overline{\mathbb S}_+
 :=\{(s,k):\operatorname{Re}s\ge0,
 |s|^2+|k|^2=1\},
 \label{eq:V50-normalized-cone}
\end{equation}
with every square root defined by the retarded boundary value.  Rotational
covariance identifies all directions of \(k\) with the calculation above.

\begin{theorem}[Uniform direct Einstein metric estimate]
\label{thm:V50-uniform-metric-floor}
There is \(c_E>0\) such that
\begin{equation}
 \boxed{\qquad
 \|L_*(s,k)x\|\ge c_E\|x\|
 \qquad}
 \label{eq:V50-uniform-metric-floor}
\end{equation}
for every \((s,k)\in\overline{\mathbb S}_+\) and every
\(x\in\mathbb C^5\).  In particular, the V49 explicit positive bath repairs
the complete five-variable gravitational glancing kernel and has a uniform
normalized metric inverse through both bath thresholds and gravitational
glancing.
\end{theorem}

\begin{proof}
The retarded roots \(\kappa_j\) are continuous on the closed stable sheet,
including their zero values at the two bath thresholds.  Hence \(L_*\) is a
continuous matrix family on the compact set
\eqref{eq:V50-normalized-cone}.

If \(\lambda\ne0\), Lemma~\ref{lem:V50-det-L0} makes \(L_0\) invertible and
the determinant lemma gives \eqref{eq:V50-factorization}.  The vector
measure \(\mu=\delta_{1/2}+5\delta_{9/10}\) satisfies the V49 exact
longitudinal and transverse strong-gain criteria, and the scalar measure
\(\nu=\delta_{1/2}+3\delta_{9/10}\) satisfies its transverse criterion.
Therefore none of the three factors vanishes anywhere on this part of the
closed cone.  If \(\lambda=0\), then \(r>0\) and
Theorem~\ref{thm:V50-direct-glancing} gives a nonzero determinant.  Thus
\(L_*\) is invertible at every point of the compact cone.

The least singular value is a continuous function of the matrix entries.
It is positive pointwise and therefore has a positive minimum \(c_E\) on
the compact cone.  This proves \eqref{eq:V50-uniform-metric-floor}.
\end{proof}

\subsection{The unweighted conservative bath norm has no closed-cone floor}
\label{subsec:V50-energy-obstruction}

For a stable outgoing bath amplitude define its ordinary half-line energy
quantity
\begin{equation}
 \mathcal E_j[q_j;s,k]
 :=\int_0^\infty\left(
 |\partial_aq_j|_{M_j}^2+
 (|s|^2+v_j^2r^2)|q_j|_{M_j}^2
 \right)\dd a,
 \qquad |u|_M^2:=u^*Mu.
 \label{eq:V50-unweighted-energy}
\end{equation}
It is the positive Fourier-amplitude counterpart of the V47 conserved bath
energy.  For \eqref{eq:V50-outgoing-solution}, exact integration gives
\begin{equation}
 \mathcal E_j
 =\frac{|\kappa_j|^2+|s|^2+v_j^2r^2}
 {2\operatorname{Re}\kappa_j}
 |Jx|_{M_j}^2,
 \qquad \operatorname{Re}s>0.
 \label{eq:V50-energy-exact}
\end{equation}

\begin{theorem}[Threshold divergence of total bath energy]
\label{thm:V50-energy-divergence}
No constant \(C\) can make
\begin{equation}
 \|x\|^2+\sum_{j=1}^2\mathcal E_j[q_j;s,k]
 \le C\|F\|^2
 \label{eq:V50-false-energy-estimate}
\end{equation}
hold for all normalized open stable frequencies and all solutions of
\eqref{eq:V50-bath-ODE}--\eqref{eq:V50-metric-traction}.  More precisely,
approaching gravitational glancing with
\(s=\delta+\sigma ir\), \(\delta\downarrow0\), gives
\begin{equation}
 \kappa_j
 =\sigma ir\sqrt{1-v_j^2}
  +\frac{\delta}{\sqrt{1-v_j^2}}+O(\delta^2)
 \label{eq:V50-kappa-expansion}
\end{equation}
at fixed scale, and hence
\begin{equation}
 \mathcal E_j
 =\left(
 \frac{r^2\sqrt{1-v_j^2}}{\delta}+O(1)
 \right)|Jx|_{M_j}^2.
 \label{eq:V50-energy-divergence}
\end{equation}
The same leading law holds after normalization to
\eqref{eq:V50-normalized-cone}.
\end{theorem}

\begin{proof}
Equation~\eqref{eq:V50-energy-exact} follows from
\[
 \int_0^\infty e^{-2a\operatorname{Re}\kappa_j}\dd a
 =\frac1{2\operatorname{Re}\kappa_j}.
\]
Differentiate the identity
\(\kappa_j^2=s^2+v_j^2r^2\) at
\(s=\sigma ir\).  Since
\(\kappa_j=\sigma ir\sqrt{1-v_j^2}\) there,
\[
 \frac{\partial\kappa_j}{\partial s}
 =\frac{s}{\kappa_j}
 =\frac1{\sqrt{1-v_j^2}},
\]
which proves \eqref{eq:V50-kappa-expansion}.  The numerator in
\eqref{eq:V50-energy-exact} tends to
\[
 (1-v_j^2)r^2+r^2+v_j^2r^2=2r^2.
\]
This proves \eqref{eq:V50-energy-divergence}.

To disprove \eqref{eq:V50-false-energy-estimate} with bounded data, choose
a normalized glancing point and an \(x_0\) with \(Jx_0\ne0\).  Put
\(F_0=L_*x_0\) at that point.  Along normalized open frequencies tending
to it, set \(x_\delta=L_*^{-1}F_0\).  Theorem~\ref{thm:V50-uniform-metric-floor}
and continuity give \(x_\delta\to x_0\), so
\(Jx_\delta\to Jx_0\).  Both explicit Grams are positive definite, and at
least one term in \eqref{eq:V50-energy-divergence} therefore tends to
infinity while \(F_0\) stays fixed.  This contradicts every finite \(C\).
Normalization changes all frequencies by a factor tending to one and does
not change the leading divergence.
\end{proof}

\begin{firewall}[What failed is the total-energy resolvent norm]
\label{fire:V50-energy-scope}
Theorem~\ref{thm:V50-energy-divergence} does not restore a metric zero: the
five-row matrix retains the positive floor
\eqref{eq:V50-uniform-metric-floor}.  At \(\operatorname{Re}s=0\), the
half-line solutions are oscillatory generalized eigenfunctions of the
conservative bath and have infinite total spatial energy.  A uniform
closed-cone resolvent into the unweighted energy Hilbert space is therefore
the wrong statement at continuous spectrum.  The outgoing limiting-
absorption topology is constructed next; it is not relabelled as conserved
total energy.
\end{firewall}

\subsection{A uniform weighted outgoing realization with the bath fields
retained}
\label{subsec:V50-weighted-outgoing}

For \(\alpha>1/2\), put
\begin{equation}
 \langle a\rangle=(1+a^2)^{1/2},
 \qquad I_\alpha:=\int_0^\infty\langle a\rangle^{-2\alpha}\dd a<\infty,
 \label{eq:V50-weight-integral}
\end{equation}
and define
\begin{equation}
 \|q_j\|_{LE_\alpha(s,k;M_j)}^2
 :=\int_0^\infty\langle a\rangle^{-2\alpha}
 \left(
 |\partial_aq_j|_{M_j}^2+
 (|s|^2+v_j^2r^2)|q_j|_{M_j}^2
 \right)\dd a.
 \label{eq:V50-weighted-LE}
\end{equation}
The outgoing condition means the limiting-absorption branch
\(q_j=e^{-\kappa_ja}Jx\), including at \(\operatorname{Re}s=0\); membership
in the weighted space alone is not used to choose a radiation sign.

\begin{theorem}[Uniform full outgoing trace-graph estimate]
\label{thm:V50-weighted-graph}
For every \(\alpha>1/2\), all normalized closed stable frequencies, and all
outgoing states satisfying \eqref{eq:V50-bath-ODE}--%
\eqref{eq:V50-metric-traction},
\begin{equation}
 \boxed{\quad
 \|x\|^2+\sum_{j=1}^2
 \|q_j\|_{LE_\alpha(s,k;M_j)}^2
 \le \frac{1+12I_\alpha}{c_E^2}\,\|F\|^2.
 \quad}
 \label{eq:V50-weighted-graph-estimate}
\end{equation}
This estimate includes the two local half-line fields and remains valid at
both bath thresholds and at gravitational glancing as an outgoing
limiting-absorption statement.
\end{theorem}

\begin{proof}
On the normalized cone,
\[
 |\kappa_j|^2=|s^2+v_j^2r^2|
 \le |s|^2+v_j^2r^2\le1.
\]
Also \(\operatorname{Re}\kappa_j\ge0\), so
\(|e^{-\kappa_ja}|\le1\).  Therefore
\begin{align*}
 \sum_{j=1}^2\|q_j\|_{LE_\alpha}^2
 &\le2I_\alpha\sum_{j=1}^2|Jx|_{M_j}^2\\
 &=2I_\alpha(Jx)^*(M_1+M_2)Jx\\
 &\le12I_\alpha\|x\|^2,
\end{align*}
because \(M_1+M_2=\operatorname{diag}(6,6,4)\) and
\(\|J\|=1\).  Theorem~\ref{thm:V50-uniform-metric-floor} gives
\(\|x\|\le c_E^{-1}\|F\|\).  Combining the two inequalities proves
\eqref{eq:V50-weighted-graph-estimate}.  The proof used no division by
\(\operatorname{Re}\kappa_j\), so it includes both thresholds and the
oscillatory glancing boundary values.
\end{proof}

\begin{assessment}[Exact scope of the linear continuum result]
\label{ass:V50-linear-scope}
Theorem~\ref{thm:V50-uniform-metric-floor} closes the V49 direct
five-variable glancing gate.  Theorem~\ref{thm:V50-weighted-graph} retains
the local bath fields and controls the complete homogeneous outgoing trace
graph in a standard limiting-absorption local-energy topology.  These are
linear frozen principal-symbol statements.

They do not yet prove the inhomogeneous half-line limiting-absorption
resolvent, a time-domain maximal evolution generator, or a nonlinear
covariant Einstein boundary action with complete stress, Bianchi, ghost,
and corner rows.  The weight in \eqref{eq:V50-weighted-LE} is an analytic
radiation topology, not the conserved V47 total energy.  The Standard-Model
matter, Higgs, neutrino, and Yang--Mills sectors have not yet been coupled to
this completed linear trace graph.
\end{assessment}

\subsection{Finite certificate and V51 upstream target}
\label{subsec:V50-certificate-next}

At normalized gravitational glancing, \(r=2^{-1/2}\).  Direct
double-precision evaluation of \eqref{eq:V50-direct-glancing-matrix} gave,
for both signs,
\begin{equation}
 |\det L_*^\sigma|=0.148498625868288,
 \qquad s_{\min}(L_*^\sigma)=0.106445626447148.
 \label{eq:V50-numerical-glancing}
\end{equation}
The determinant agrees with \eqref{eq:V50-direct-glancing-det} to
\(1.2\times10^{-16}\).  A \(121\times241\) normalized
\((r,\arg s)\) grid gave the sampled value
\begin{equation}
 \min_{\rm grid}s_{\min}(L_*)=0.0219006.
 \label{eq:V50-sampled-floor}
\end{equation}
This is only an implementation audit; the positive constant in
Theorem~\ref{thm:V50-uniform-metric-floor} comes from exact invertibility
and compactness.  Along the upper glancing approach the numerical products
\(\delta\mathcal E_j/|Jx|_{M_j}^2\) converged to
\begin{equation}
 \frac{\sqrt3}{4}=0.4330127
 \quad\hbox{and}\quad
 \frac{\sqrt{19}}{20}=0.2179449,
 \label{eq:V50-energy-asymptotic-check}
\end{equation}
as predicted by \eqref{eq:V50-energy-divergence}.

V51 should remain upstream of nonlinear SM coupling long enough to settle
the radiation space itself:
\begin{enumerate}
\item solve the inhomogeneous matrix half-line equations with outgoing
boundary values and prove a two-weight limiting-absorption estimate,
including the adjoint incoming graph;
\item derive the corresponding Green-flux identity and identify the exact
trace spaces at the two bath thresholds, without treating weighted local
energy as conserved total energy;
\item use that estimate to construct the time-domain generator or
scattering evolution for the linear Einstein--bath system and verify the
full embedding quotient and ghost domain; and
\item only after those rows close, return to the nonlinear covariant
stress/Bianchi identity and then couple the Higgs, fermion, and Yang--Mills
sources through their complete action incidence.
\end{enumerate}
If the inhomogeneous threshold estimate loses a trace derivative, V51 must
record the loss in the Einstein source norm and move to the bulk--bath
coupling; it must not hide the loss by strengthening the bath datum after
the fact.

\subsection{V50 ledger}
\label{subsec:V50-ledger}

\paragraph{New conclusions proved in V50.}
The explicit V49 Grams were inserted into the original V41 metric rows.
Theorem~\ref{thm:V50-direct-glancing} computes their exact glancing matrix,
determinant, and inverse without \(L_0^{-1}\).  Theorem~%
\ref{thm:V50-uniform-metric-floor} proves a uniform normalized inverse for
all five metric variables on the entire closed stable cone.
Theorem~\ref{thm:V50-energy-divergence} proves that the corresponding
unweighted total bath-energy resolvent estimate is false at continuous
spectrum.  Theorem~\ref{thm:V50-weighted-graph} retains both local
half-line fields and proves their uniform outgoing trace-graph estimate in
weighted local energy.

\paragraph{Imported conclusions not repeated.}
V41 supplies \(L_0,C,J\) and the DeWitt graph.  V47 supplies the local
positive direct-sum bath and its common trace.  V49 supplies the exact
zero-free factor criteria for \(\lambda\ne0\) and the explicit Grams.  The
NS V24 and YM V29 companion results affect the completeness and uniformity
discipline only, as recorded in
Firewall~\ref{fire:V50-companion-types}.

\paragraph{Honest continuum status.}
The frozen linear Einstein metric trace now has a rigorous uniform
closed-cone floor for an explicit local conservative bath, and the complete
homogeneous outgoing bath graph is controlled in weighted local energy.
The unweighted half-line energy resolvent is rigorously excluded at the
continuous spectrum.  Inhomogeneous limiting absorption, time-domain
well-posedness, nonlinear covariant stress/Bianchi closure, ghosts, and the
coupled Standard-Model sectors remain open.  The full Standard
Model--Einstein continuum is therefore not proved.

\section*{Version V50 append-only record}
\addcontentsline{toc}{section}{Version V50 append-only record}
The complete V49 mathematical source was retained before this continuation.
Its SHA--256 digest was
\begin{center}\scriptsize\ttfamily
f93274386c93c686501eb29fb513d98b21eee9eca322949ba986b43f3a038bf0.
\end{center}
No mathematical content of V49 was altered.  On the next iteration, retain
this full file, append before its unique active document-closing command,
increment the filename to \texttt{\_V51.tex}, and remove only the
superseded V50 file from this exact SM note series.  The next scheduled
companion audit is V55.

% END APPENDED VERSION V50 MATERIAL


% BEGIN APPENDED VERSION V51 MATERIAL

\clearpage
\section{V51: sharp two-weight limiting absorption, source-complete
traction, and the adjoint radiation graph}
\label{sec:V51-limiting-absorption}

\subsection{Purpose and nonduplication}
\label{subsec:V51-purpose}

V50 proved a uniform closed-cone inverse for the five metric amplitudes and
for the homogeneous outgoing bath trace graph.  It also proved that the
unweighted half-line energy of a boundary-forced generalized eigenfunction
diverges at continuous spectrum.  The missing object was therefore not
another finite port determinant.  It was the inhomogeneous half-line
resolvent, including the force which an interior bath source returns to the
Einstein trace.

V51 computes that object exactly.  The Dirichlet half-line Green kernel is
uniform down to \(\kappa=0\), but its threshold kernel
\(\min\{a,b\}\) imposes a sharp two-weight polygon.  If the source has
weight \(\alpha\) and the outgoing field is measured with reciprocal weight
\(-\beta\), the uniform \(H^2\) limiting-absorption estimate holds exactly
when
\[
 \alpha>\frac12,\qquad
 \beta>\frac12,\qquad
 \alpha+\beta\ge2.
\]
The boundary load itself requires only \(\alpha>1/2\).  Thus there is no
hidden tangential derivative loss, but there is an unavoidable half-line
moment trade at the zero bath thresholds.

The two matrix baths are then reassembled at their common metric trace.  An
interior source \(g_j\) changes the \(j\)-th traction by the exact load
\(-M_j\int e^{-\kappa_j a}g_j(a)\,\dd a\).  Combining this load with the V50
metric floor gives a uniform inhomogeneous estimate for the five Einstein
variables, both bath fields, and both tractions.  Finally the incoming
adjoint kernel and its Green identity are proved.  These steps retain the
local continuum variables throughout and do not reinterpret weighted local
energy as the conserved V47 bath energy.

This is not a scheduled companion-audit version.  The next NS/YM audit
remains V55.

\subsection{The scalar Dirichlet radiation atlas}
\label{subsec:V51-scalar-atlas}

For \(\gamma\in\mathbb R\), define
\begin{equation}
 L^2_\gamma(\mathbb R_+)
 :=\left\{f:
 \|f\|_\gamma^2
 :=\int_0^\infty\langle a\rangle^{2\gamma}|f(a)|^2\,\dd a
 <\infty\right\},
 \qquad
 \langle a\rangle=(1+a^2)^{1/2}.
 \label{eq:V51-weighted-L2}
\end{equation}
Vector-valued spaces use the Euclidean component norm.  Let
\begin{equation}
 \mathcal K
 :=\{\kappa\in\mathbb C:
 \operatorname{Re}\kappa\ge0,\ |\kappa|\le1\}.
 \label{eq:V51-kappa-set}
\end{equation}
For \(\kappa\ne0\), put
\begin{equation}
 G_\kappa(a,b)
 :=\frac{e^{-\kappa|a-b|}-e^{-\kappa(a+b)}}{2\kappa}
 =e^{-\kappa\max\{a,b\}}
   \frac{\sinh(\kappa\min\{a,b\})}{\kappa},
 \label{eq:V51-Gk}
\end{equation}
and define its threshold value by
\begin{equation}
 G_0(a,b):=\min\{a,b\}.
 \label{eq:V51-G0}
\end{equation}
The corresponding Dirichlet resolvent and boundary-load functional are
\begin{align}
 (R_D(\kappa)g)(a)
 &:=\int_0^\infty G_\kappa(a,b)g(b)\,\dd b,
 \label{eq:V51-RD}\\
 \ell_\kappa(g)
 &:=\int_0^\infty e^{-\kappa b}g(b)\,\dd b.
 \label{eq:V51-load}
\end{align}
All boundary values with \(\operatorname{Re}\kappa=0\) are inherited from
\(\operatorname{Re}\kappa>0\).

\begin{proposition}[Exact inhomogeneous half-line solution and traction]
\label{prop:V51-exact-halfline}
Let \(\operatorname{Re}\kappa>0\), \(z\in\mathbb C^m\), and let \(g\) be
smooth with compact support.  The unique decaying solution of
\begin{equation}
 (-\partial_a^2+\kappa^2)q=g,
 \qquad q(0)=z,
 \label{eq:V51-halfline-equation}
\end{equation}
is
\begin{equation}
 q(a)=e^{-\kappa a}z+R_D(\kappa)g(a).
 \label{eq:V51-halfline-solution}
\end{equation}
Its normal derivative and outward traction for a constant symmetric matrix
\(M\) are
\begin{align}
 q'(0)&=-\kappa z+\ell_\kappa(g),
 \label{eq:V51-normal-derivative}\\
 t_M(q):=-Mq'(0)
 &=\kappa Mz-M\ell_\kappa(g).
 \label{eq:V51-exact-traction}
\end{align}
The same formulas define the limiting-absorption solution at every
\(\kappa\in\mathcal K\), including
\begin{equation}
 q(a)=z+\int_0^\infty\min\{a,b\}g(b)\,\dd b,
 \qquad
 t_M(q)=-M\int_0^\infty g(b)\,\dd b
 \label{eq:V51-threshold-solution}
\end{equation}
at \(\kappa=0\).
\end{proposition}

\begin{proof}
For \(a\ne b\), both exponentials in \eqref{eq:V51-Gk} solve the
homogeneous equation.  The kernel vanishes at \(a=0\), is continuous at
\(a=b\), and its first derivative has jump
\[
 \partial_aG_\kappa(b+0,b)-\partial_aG_\kappa(b-0,b)=-1.
\]
Therefore
\((-\partial_a^2+\kappa^2)G_\kappa(\,\cdot\,,b)=\delta_b\).
The two exponential terms decay when
\(\operatorname{Re}\kappa>0\), which proves
\eqref{eq:V51-halfline-solution}.  The difference of two decaying solutions
with zero trace is a multiple of \(\sinh(\kappa a)\) and hence must vanish.

For \(0<a<b\),
\[
 \partial_aG_\kappa(a,b)=e^{-\kappa b}\cosh(\kappa a),
\]
so \(\partial_aG_\kappa(0,b)=e^{-\kappa b}\).  Differentiating
\eqref{eq:V51-halfline-solution} at zero proves
\eqref{eq:V51-normal-derivative} and
\eqref{eq:V51-exact-traction}.  Finally
\(\sinh(\kappa m)/\kappa\to m\), so
\(G_\kappa(a,b)\to\min\{a,b\}\).  The threshold formulas follow by taking
the limiting-absorption value.
\end{proof}

\subsection{Uniform kernel bounds and the weighted Hardy edge}
\label{subsec:V51-Hardy}

\begin{lemma}[Uniform pointwise Green bounds]
\label{lem:V51-Green-bounds}
For every \(\kappa\in\mathcal K\) and \(a,b\ge0\),
\begin{equation}
 |G_\kappa(a,b)|\le\min\{a,b\},
 \qquad
 |\partial_aG_\kappa(a,b)|\le1
 \quad(a\ne b).
 \label{eq:V51-Green-bounds}
\end{equation}
Moreover, for every \(\alpha>1/2\),
\begin{equation}
 |\ell_\kappa(g)|
 \le I_\alpha^{1/2}\|g\|_\alpha,
 \qquad
 I_\alpha:=\int_0^\infty
 \langle b\rangle^{-2\alpha}\,\dd b,
 \label{eq:V51-load-bound}
\end{equation}
uniformly in \(\kappa\in\mathcal K\).
\end{lemma}

\begin{proof}
Write \(m=\min\{a,b\}\), \(A=\max\{a,b\}\), and
\(x=\operatorname{Re}\kappa\ge0\).  Since
\[
 \frac{\sinh(\kappa m)}{\kappa}
 =\int_0^m\cosh(\kappa t)\,\dd t,
 \qquad
 |\cosh(\kappa t)|\le e^{xt},
\]
the second expression in \eqref{eq:V51-Gk} gives
\[
 |G_\kappa(a,b)|
 \le\int_0^m e^{-x(A-t)}\,\dd t\le m.
\]
If \(a<b\), then
\(\partial_aG_\kappa=e^{-\kappa b}\cosh(\kappa a)\), whose modulus is at
most one.  If \(a>b\), then
\[
 \partial_aG_\kappa=-e^{-\kappa a}\sinh(\kappa b).
\]
The identity
\(|\sinh(x+iy)|^2=\sinh^2x+\sin^2y\le\cosh^2x\)
again gives a bound by one.  The threshold bounds follow by continuity.
Finally \(|e^{-\kappa b}|\le1\), and weighted Cauchy--Schwarz gives
\eqref{eq:V51-load-bound}.
\end{proof}

The zero kernel has the exact decomposition
\begin{equation}
 (R_D(0)g)(a)
 =\int_0^a b\,g(b)\,\dd b
  +a\int_a^\infty g(b)\,\dd b.
 \label{eq:V51-zero-decomposition}
\end{equation}
The two terms are the forward and backward weighted Hardy operators.

\begin{lemma}[The two Hardy criteria used at threshold]
\label{lem:V51-two-Hardy}
Let \(u,v>0\).  The estimates
\begin{align}
 \left\|\int_0^a f(b)\,\dd b\right\|_{L^2(u\,\dd a)}
 &\le2A^{1/2}\|f\|_{L^2(v\,\dd b)},
 \label{eq:V51-forward-Hardy}\\
 \left\|\int_a^\infty f(b)\,\dd b\right\|_{L^2(u\,\dd a)}
 &\le2B^{1/2}\|f\|_{L^2(v\,\dd b)}
 \label{eq:V51-backward-Hardy}
\end{align}
hold whenever, respectively,
\begin{align}
 A&:=\sup_{R>0}
 \left(\int_R^\infty u(a)\,\dd a\right)
 \left(\int_0^R v(b)^{-1}\,\dd b\right)<\infty,
 \label{eq:V51-Hardy-A}\\
 B&:=\sup_{R>0}
 \left(\int_0^R u(a)\,\dd a\right)
 \left(\int_R^\infty v(b)^{-1}\,\dd b\right)<\infty.
 \label{eq:V51-Hardy-B}
\end{align}
\end{lemma}

\begin{proof}
Both integral operators are positive, so it is enough to use
\(f\ge0\).  For the forward operator set
\[
 U(R)=\int_R^\infty u(a)\,\dd a,
 \qquad
 V(R)=\int_0^R v(b)^{-1}\,\dd b,
 \qquad
 F(R)=\int_0^R f(b)\,\dd b.
\]
Cauchy--Schwarz with the factors \(V(b)^{1/4}\) and
\(V(b)^{-1/4}\) gives
\[
 F(R)^2
 \le
 2V(R)^{1/2}
 \int_0^R v(b)f(b)^2V(b)^{1/2}\,\dd b,
\]
because
\[
 \int_0^R v(b)^{-1}V(b)^{-1/2}\,\dd b
 =2V(R)^{1/2}.
\]
Fubini therefore yields
\begin{align*}
 \int_0^\infty u(R)F(R)^2\,\dd R
 &\le2\int_0^\infty v(b)f(b)^2V(b)^{1/2}
 \left(\int_b^\infty u(R)V(R)^{1/2}\,\dd R\right)\dd b.
\end{align*}
Since \(U(R)V(R)\le A\), integration by parts gives
\begin{align*}
 \int_b^\infty u(R)V(R)^{1/2}\,\dd R
 &=U(b)V(b)^{1/2}
 +\frac12\int_b^\infty
 U(R)V(R)^{-1/2}\,\dd V(R)\\
 &\le A V(b)^{-1/2}
 +\frac A2\int_{V(b)}^\infty t^{-3/2}\,\dd t\\
 &=2A V(b)^{-1/2}.
\end{align*}
If \(V(\infty)<\infty\), stopping the last integral there only improves
the bound.  Substitution proves
\[
 \int_0^\infty u|F|^2
 \le4A\int_0^\infty v|f|^2,
\]
which is \eqref{eq:V51-forward-Hardy}.

For the backward operator use
\[
 \widetilde U(R)=\int_0^R u(a)\,\dd a,
 \qquad
 \widetilde V(R)=\int_R^\infty v(b)^{-1}\,\dd b,
 \qquad
 \widetilde F(R)=\int_R^\infty f(b)\,\dd b.
\]
The same calculation with the integration directions reversed gives
\[
 \int_0^\infty u|\widetilde F|^2
 \le4B\int_0^\infty v|f|^2.
\]
All integrations may first be performed on
\((\varepsilon,\varepsilon^{-1})\); monotone convergence handles zero or
infinite endpoint values.  Density removes the compact-support
restriction.
\end{proof}

\subsection{The sharp two-weight limiting-absorption theorem}
\label{subsec:V51-sharp-LAP}

Write \(H^2_{-\beta}\) for the weighted norm
\begin{equation}
 \|u\|_{H^2_{-\beta}}^2
 :=\sum_{m=0}^2
 \int_0^\infty\langle a\rangle^{-2\beta}
 |\partial_a^m u(a)|^2\,\dd a.
 \label{eq:V51-weighted-H2}
\end{equation}

\begin{theorem}[Sharp uniform half-line weight polygon]
\label{thm:V51-sharp-weight-polygon}
The uniform estimate
\begin{equation}
 \sup_{\kappa\in\mathcal K}
 \|R_D(\kappa)g\|_{H^2_{-\beta}}
 \le C_{\alpha,\beta}\|g\|_\alpha
 \label{eq:V51-uniform-LAP}
\end{equation}
holds for all compactly supported smooth \(g\), and hence by completion, if
and only if
\begin{equation}
 \boxed{\qquad
 \alpha>\frac12,\qquad
 \beta>\frac12,\qquad
 \alpha+\beta\ge2.
 \qquad}
 \label{eq:V51-sharp-polygon}
\end{equation}
On this same region, \(R_D(\kappa)\) extends continuously in the weak
operator topology to the imaginary \(\kappa\)-axis, including
\(\kappa=0\).  Separately,
\begin{equation}
 \sup_{\kappa\in\mathcal K}
 \|\ell_\kappa\|_{L^2_\alpha\to\mathbb C}<\infty
 \quad\Longleftrightarrow\quad
 \alpha>\frac12.
 \label{eq:V51-sharp-load}
\end{equation}
\end{theorem}

\begin{proof}
Assume \eqref{eq:V51-sharp-polygon}.  By
Lemma~\ref{lem:V51-Green-bounds}, the zeroth-order part of
\(R_D(\kappa)\) is pointwise dominated by the positive threshold operator
\(R_D(0)\).  Apply Lemma~\ref{lem:V51-two-Hardy} to the first term of
\eqref{eq:V51-zero-decomposition} with
\[
 u(a)=\langle a\rangle^{-2\beta},
 \qquad
 v(b)=\frac{\langle b\rangle^{2\alpha}}{b^2}.
\]
The product in \eqref{eq:V51-Hardy-A} is finite.  At infinity its two
factors have product bounded by a constant multiple of
\[
 R^{1-2\beta}
 \begin{cases}
 R^{3-2\alpha},&\alpha<3/2,\\
 1+\log R,&\alpha=3/2,\\
 1,&\alpha>3/2,
 \end{cases}
 \label{eq:V51-forward-asymptotic}
\]
and the conditions in \eqref{eq:V51-sharp-polygon} make every case
bounded.  The origin causes no singularity.

For the second term of \eqref{eq:V51-zero-decomposition}, use the backward
Hardy inequality with
\[
 u(a)=a^2\langle a\rangle^{-2\beta},
 \qquad
 v(b)=\langle b\rangle^{2\alpha}.
\]
The product in \eqref{eq:V51-Hardy-B} is finite; at infinity it is bounded
by a constant multiple of
\[
 \begin{cases}
 R^{3-2\beta},&\beta<3/2,\\
 1+\log R,&\beta=3/2,\\
 1,&\beta>3/2
 \end{cases}
 R^{1-2\alpha}.
 \label{eq:V51-backward-asymptotic}
\]
Again \eqref{eq:V51-sharp-polygon} is exactly sufficient.  This proves the
uniform \(L^2_\alpha\to L^2_{-\beta}\) bound.

The derivative kernel is bounded by one.  Its weighted kernel
\[
 \langle a\rangle^{-\beta}
 \partial_aG_\kappa(a,b)
 \langle b\rangle^{-\alpha}
\]
is uniformly Hilbert--Schmidt because \(\alpha,\beta>1/2\).  Hence the
first derivative obeys the same type of estimate.  Finally
\[
 \partial_a^2R_D(\kappa)g
 =\kappa^2R_D(\kappa)g-g.
\]
Here \(|\kappa|\le1\), and \(L^2_\alpha\) embeds in
\(L^2_{-\beta}\).  This proves \eqref{eq:V51-uniform-LAP}.
The pointwise kernel convergence, the uniform bound, and density give the
stated weak operator continuity.

For necessity, the supremum includes \(\kappa=0\).  Choose a compactly
supported \(g\) with \(\int b\,g(b)\,\dd b\ne0\).  Then
\(R_D(0)g\) is a nonzero constant beyond the support of \(g\), forcing
\(\beta>1/2\).  The kernel is symmetric, so boundedness
\(L^2_\alpha\to L^2_{-\beta}\) implies by duality boundedness
\(L^2_\beta\to L^2_{-\alpha}\).  The same constant-tail test forces
\(\alpha>1/2\).

For \(N\ge1\), set
\[
 g_N(b)=N^{-\alpha-1/2}\mathbf1_{[N,2N]}(b).
\]
Its \(L^2_\alpha\) norm is bounded above and below independently of \(N\).
For \(a\ge2N\),
\[
 R_D(0)g_N(a)=\int_N^{2N}b\,g_N(b)\,\dd b
 =\frac32N^{3/2-\alpha}.
\]
Measuring this constant on \(2N\le a\le4N\) gives
\[
 \|R_D(0)g_N\|_{-\beta}
 \ge c_{\alpha,\beta}N^{2-\alpha-\beta}.
\]
Uniform boundedness therefore requires \(\alpha+\beta\ge2\).
This proves necessity of the polygon.

The load assertion follows from \eqref{eq:V51-load-bound}.  At
\(\kappa=0\), \(\ell_0(g)=\int g\); by the Riesz representation theorem
this is bounded on \(L^2_\alpha\) exactly when
\(\langle a\rangle^{-\alpha}\in L^2\), namely
\(\alpha>1/2\).
\end{proof}

\begin{remark}[The diagonal choice]
\label{rem:V51-diagonal-weight}
Taking \(\alpha=\beta\) gives the simple sharp diagonal condition
\(\alpha=\beta\ge1\).  The endpoint \((1,1)\) is permitted by the Hardy
estimate even though the corresponding weighted kernel is not
Hilbert--Schmidt.  Replacing the Hardy argument by a Hilbert--Schmidt
estimate would incorrectly discard this endpoint.
\end{remark}

\subsection{The source-complete Einstein--bath traction system}
\label{subsec:V51-coupled-system}

Return to the explicit V50 speeds and Grams.  On the normalized stable cone
let
\begin{equation}
 \kappa_j=\sqrt{s^2+v_j^2r^2},
 \qquad
 v_1=\frac12,\quad v_2=\frac9{10},
 \qquad
 M_1=I_3,\quad M_2=\operatorname{diag}(5,5,3).
 \label{eq:V51-two-baths}
\end{equation}
Give the two local bath equations independent interior sources
\(g_j\in L^2_\alpha(\mathbb R_+;\mathbb C^3)\):
\begin{align}
 (-\partial_a^2+\kappa_j^2)q_j&=g_j,
 &q_j(0)&=Jx,
 \label{eq:V51-sourced-baths}\\
 L_0x-C(t_1+t_2)&=F,
 &t_j&:=-M_jq_j'(0).
 \label{eq:V51-sourced-metric}
\end{align}
Define the complete returned source
\begin{equation}
 \Lambda(g_1,g_2)
 :=\sum_{j=1}^2M_j\ell_{\kappa_j}(g_j).
 \label{eq:V51-complete-load}
\end{equation}

\begin{theorem}[Exact reduction with every source retained]
\label{thm:V51-source-complete-reduction}
For every open stable frequency, and by limiting absorption on the closed
cone, the system \eqref{eq:V51-sourced-baths}--%
\eqref{eq:V51-sourced-metric} is equivalent to
\begin{align}
 q_j&=e^{-\kappa_ja}Jx+R_D(\kappa_j)g_j,
 \label{eq:V51-sourced-q}\\
 t_j&=\kappa_jM_jJx-M_j\ell_{\kappa_j}(g_j),
 \label{eq:V51-sourced-t}\\
 L_*x&=F-C\Lambda(g_1,g_2),
 \label{eq:V51-sourced-Lstar}
\end{align}
where \(L_*=L_0-C(\kappa_1M_1+\kappa_2M_2)J\) is the V50 matrix.  No
component of either interior source disappears under reduction.
\end{theorem}

\begin{proof}
Apply Proposition~\ref{prop:V51-exact-halfline} componentwise to each
bath.  This gives \eqref{eq:V51-sourced-q} and
\eqref{eq:V51-sourced-t}.  Substitution into the metric row gives
\[
 L_0x-C\sum_j\kappa_jM_jJx
 +C\sum_jM_j\ell_{\kappa_j}(g_j)=F.
\]
The first two terms are \(L_*x\); rearrangement proves
\eqref{eq:V51-sourced-Lstar}.  The source dependence occurs through the
literal sum \eqref{eq:V51-complete-load}, so there is no trace-only or
single-bath replacement.
\end{proof}

\begin{theorem}[Uniform inhomogeneous Einstein--bath estimate]
\label{thm:V51-full-inhomogeneous-estimate}
Assume the sharp weight conditions
\eqref{eq:V51-sharp-polygon}.  There is a constant
\(C_{\alpha,\beta}^{\rm EB}<\infty\), independent of the normalized closed
stable frequency, such that every outgoing solution satisfies
\begin{align}
 \|x\|
 +\sum_{j=1}^2\left(
  \|q_j\|_{H^2_{-\beta}}+|t_j|
 \right)
 \le C_{\alpha,\beta}^{\rm EB}
 \left(
  \|F\|+\sum_{j=1}^2\|g_j\|_\alpha
 \right).
 \label{eq:V51-full-inhomogeneous-estimate}
\end{align}
The estimate includes both zero bath thresholds, gravitational glancing,
and the upper and lower continuous-spectrum boundary.  The only
half-line weight loss is the sharp polygon
\eqref{eq:V51-sharp-polygon}.
\end{theorem}

\begin{proof}
The load bound \eqref{eq:V51-load-bound}, the fixed matrices \(C,M_j\), and
\eqref{eq:V51-sourced-Lstar} give
\[
 \|L_*x\|
 \le\|F\|+C_\alpha\sum_j\|g_j\|_\alpha.
\]
The V50 uniform metric floor therefore yields
\begin{equation}
 \|x\|
 \le c_E^{-1}
 \left(\|F\|+C_\alpha\sum_j\|g_j\|_\alpha\right).
 \label{eq:V51-x-bound}
\end{equation}

For \(\beta>1/2\), the homogeneous extension obeys
\[
 \sup_{\kappa\in\mathcal K}
 \|e^{-\kappa a}z\|_{H^2_{-\beta}}
 \le C_\beta|z|,
\]
because \(|\kappa|\le1\), \(|e^{-\kappa a}|\le1\), and
\(\langle a\rangle^{-\beta}\in L^2\).  Combine this with
Theorem~\ref{thm:V51-sharp-weight-polygon} in
\eqref{eq:V51-sourced-q}.  Since \(\|J\|=1\),
\[
 \|q_j\|_{H^2_{-\beta}}
 \le C_{\alpha,\beta}
 \left(\|x\|+\|g_j\|_\alpha\right).
\]
Finally \eqref{eq:V51-sourced-t}, \(|\kappa_j|\le1\), and
\eqref{eq:V51-load-bound} give the same bound for \(|t_j|\).
Insert \eqref{eq:V51-x-bound} and sum the two bath labels.
\end{proof}

\begin{corollary}[Scale-covariant form and absence of tangential derivative
loss]
\label{cor:V51-scaled-estimate}
Let
\[
 \rho=(|s|^2+r^2)^{1/2}>0
\]
without normalizing the frequency.  Define
\begin{align}
 \|g\|_{\mathcal Y_{\alpha,\rho}}^2
 &:=\rho^{-3}\int_0^\infty
 \langle\rho a\rangle^{2\alpha}|g(a)|^2\,\dd a,
 \label{eq:V51-scaled-source}\\
 \|q\|_{\mathcal X_{-\beta,\rho}}^2
 &:=\rho\int_0^\infty\langle\rho a\rangle^{-2\beta}|q|^2\,\dd a
 +\rho^{-1}\int_0^\infty\langle\rho a\rangle^{-2\beta}|q'|^2\,\dd a
 \nonumber\\
 &\quad+\rho^{-3}\int_0^\infty
 \langle\rho a\rangle^{-2\beta}|q''|^2\,\dd a.
 \label{eq:V51-scaled-receiver}
\end{align}
Then
\begin{equation}
 \|x\|+\sum_{j=1}^2
 \left(\|q_j\|_{\mathcal X_{-\beta,\rho}}
       +\rho^{-1}|t_j|\right)
 \le C_{\alpha,\beta}^{\rm EB}
 \left(
  \rho^{-1}\|F\|
  +\sum_{j=1}^2\|g_j\|_{\mathcal Y_{\alpha,\rho}}
 \right).
 \label{eq:V51-scaled-estimate}
\end{equation}
Thus the sourced bath creates no additional tangential derivative loss
beyond the natural order-one metric row and order-two bath equation.
\end{corollary}

\begin{proof}
Put \(A=\rho a\), \(Q_j(A)=q_j(A/\rho)\), and
\[
 G_j(A)=\rho^{-2}g_j(A/\rho).
\]
After division by \(\rho^2\), the bath equation has normalized root
\(\widehat\kappa_j=\kappa_j/\rho\).  The traction is
\(t_j=\rho\widehat t_j\), while homogeneity of the principal Einstein
matrix gives \(L_*=\rho\widehat L_*\) and metric datum
\(\widehat F=F/\rho\).  Direct change of variables gives
\[
 \|G_j\|_\alpha=\|g_j\|_{\mathcal Y_{\alpha,\rho}},
 \qquad
 \|Q_j\|_{H^2_{-\beta}}
 =\|q_j\|_{\mathcal X_{-\beta,\rho}}.
\]
Apply Theorem~\ref{thm:V51-full-inhomogeneous-estimate} to the normalized
system.
\end{proof}

\begin{firewall}[The exact weight loss cannot be removed by the metric
floor]
\label{fire:V51-weight-loss}
The V50 constant \(c_E\) controls the five finite metric amplitudes after
the complete bath load is formed.  It cannot improve the half-line
threshold kernel \(\min\{a,b\}\).  If
\(\alpha\le1/2\), the returned source \(\ell_0(g)\) is not even a bounded
trace functional.  If \(\beta\le1/2\), a compact source produces a constant
threshold tail outside the receiver space.  If
\(\alpha+\beta<2\), the translated block sequence in the proof of
Theorem~\ref{thm:V51-sharp-weight-polygon} makes the field norm diverge.
Strengthening the Einstein norm does not pay any of these three failures.
\end{firewall}

\subsection{Incoming adjoint and exact Green incidence}
\label{subsec:V51-adjoint}

For a boundary value of the retarded root, call \(\kappa\) outgoing and
\(\overline\kappa\) incoming.  Define
\begin{equation}
 E_\kappa z:=e^{-\kappa a}z,
 \qquad
 \mathcal Q_\kappa(z,g):=E_\kappa z+R_D(\kappa)g.
 \label{eq:V51-Qk}
\end{equation}

\begin{theorem}[Adjoint radiation graph and Green identity]
\label{thm:V51-adjoint-Green}
On the sharp weight polygon,
\begin{equation}
 R_D(\kappa)^*=R_D(\overline\kappa):
 L^2_\beta\longrightarrow L^2_{-\alpha}
 \label{eq:V51-resolvent-adjoint}
\end{equation}
under the unweighted \(L^2\) pairing.  Let
\[
 q=\mathcal Q_\kappa(z,g),
 \qquad
 p=\mathcal Q_{\overline\kappa}(w,h),
\]
where \(g\in L^2_\alpha\), \(h\in L^2_\beta\), and initially take compact
sources.  For every real symmetric matrix \(M\),
\begin{align}
 \int_0^\infty
 \left(p^*Mg-h^*Mq\right)\dd a
 &=
 w^*Mq'(0)-p'(0)^*Mz
 \nonumber\\
 &=t_M(p)^*z-w^*t_M(q).
 \label{eq:V51-Green-identity}
\end{align}
The identity extends by density and limiting absorption to the imaginary
\(\kappa\)-axis, including \(\kappa=0\).  Consequently the source load in
\eqref{eq:V51-exact-traction} is the complete variational boundary
incidence of the inhomogeneous bath equation.
\end{theorem}

\begin{proof}
The kernel is symmetric in \(a,b\) and
\[
 \overline{G_\kappa(a,b)}=G_{\overline\kappa}(a,b).
\]
Fubini first for compact sources proves
\eqref{eq:V51-resolvent-adjoint}; the weighted bounds extend it to the
stated dual spaces.

For \(\operatorname{Re}\kappa>0\), both radiation solutions decay.  Since
\[
 p^*M(-q''+\kappa^2q)
 -(-p''+\overline\kappa^2p)^*Mq
 =\partial_a\left(p'^*Mq-p^*Mq'\right),
\]
integration over the half-line leaves only the endpoint \(a=0\) and gives
the first line of \eqref{eq:V51-Green-identity}.  The second line is the
definition \(t_M(u)=-Mu'(0)\).  The pairings are continuous between
\(L^2_\alpha\) and \(L^2_{-\alpha}\), and between
\(L^2_\beta\) and \(L^2_{-\beta}\).  Theorem~%
\ref{thm:V51-sharp-weight-polygon} and
Lemma~\ref{lem:V51-Green-bounds} therefore permit
\(\operatorname{Re}\kappa\downarrow0\).  This proves the boundary identity
for the outgoing/incoming limiting pair, including the threshold kernel.
\end{proof}

\begin{corollary}[Complete two-bath adjoint incidence]
\label{cor:V51-two-bath-adjoint}
Summing \eqref{eq:V51-Green-identity} with \(M=M_j\) over \(j=1,2\)
produces exactly the common-trace boundary covector
\[
 t_1+t_2
 =Z_*Jx-\Lambda(g_1,g_2)
\]
appearing in \eqref{eq:V51-sourced-metric}.  Because
\(HC=J^T\), its contraction with the DeWitt metric row is the contraction
with the common bath trace.  No separate source Gram or trace-only
replacement is needed for the frozen linear Green identity.
\end{corollary}

\begin{proof}
The first assertion is the sum of
\eqref{eq:V51-sourced-t}.  The V41 identity
\(C=H^{-1}J^T\) gives \(HC=J^T\).  Hence, for a metric row force \(f\),
\[
 \mathbb G(Cf,x)=f^*Jx
\]
with the corresponding Hermitian polarization.  Apply this to the summed
traction and combine it with
Theorem~\ref{thm:V51-adjoint-Green}.
\end{proof}

\subsection{What is still required for a time-domain generator}
\label{subsec:V51-generator-gate}

\begin{assessment}[Frequency-domain closure and generator boundary]
\label{ass:V51-generator-boundary}
The scale-covariant estimate
\eqref{eq:V51-scaled-estimate} supplies the exact inhomogeneous boundary
resolvent which V50 lacked.  It controls the complete returned source,
includes the adjoint incoming graph, and is uniform through both bath
thresholds.  It is therefore suitable input for a Kreiss or semigroup
construction of the coupled linear problem.

It is not by itself that construction.  The matrix \(L_*(s,k)\) is the
stable boundary trace of the bulk Einstein wave equation after selecting
the decaying normal root.  A time-domain generator also requires the bulk
inhomogeneous Einstein resolvent, a common energy domain for the
full-embedding quotient and the harmonic ghost quartet, and a proof that
the boundary Green relation is maximal on that domain.  Those data are not
contained in a five-by-five stable matrix.  Declaring maximal
dissipativity from \eqref{eq:V51-scaled-estimate} alone would repeat the
finite-to-continuum error excluded by the V50 audit.
\end{assessment}

The next upstream step is now precise.  V52 should attach the V51 boundary
resolvent to the frozen harmonic Einstein bulk system on a half-space,
retain the trace-free and full-embedding rows, and prove an inhomogeneous
Kreiss estimate in a declared Sobolev/weighted bath graph norm.  It should
then identify the adjoint domain from
\eqref{eq:V51-Green-identity}.  Only if that domain is maximal may it be
promoted to a time-domain generator.  Failure at zero tangential frequency
must be reported separately from the already closed gravitational
glancing cone.

\subsection{Finite certificate}
\label{subsec:V51-certificate}

For the exponential source \(g(a)=e^{-\varrho a}c\),
\(\operatorname{Re}\varrho>0\), the exact scalar card is
\begin{align}
 \ell_\kappa(g)
 &=\frac{c}{\kappa+\varrho},
 \label{eq:V51-exp-load}\\
 R_D(\kappa)g(a)
 &=\frac{e^{-\varrho a}-e^{-\kappa a}}
 {\kappa^2-\varrho^2}\,c,
 \qquad \kappa\ne\varrho,
 \label{eq:V51-exp-resolvent}\\
 R_D(0)g(a)
 &=\frac{1-e^{-\varrho a}}{\varrho^2}\,c.
 \label{eq:V51-exp-threshold}
\end{align}
The removable value at \(\kappa=\varrho\) is
\(a e^{-\varrho a}c/(2\varrho)\).  Direct substitution makes the ODE
residual, boundary residual, and traction residual identically zero.  These
formulas test the sign in \eqref{eq:V51-exact-traction} and the continuous
threshold value without numerical quadrature.

A separate dyadic card uses \(g_N\) from the necessity proof.  Its source
norm remains between two positive constants while its receiver norm grows
as \(N^{2-\alpha-\beta}\).  This is an exact regression test for the third
face of \eqref{eq:V51-sharp-polygon}; a finite-grid Green-kernel test cannot
replace it.

\subsection{V51 ledger and V52 acquisition order}
\label{subsec:V51-ledger}

\paragraph{New conclusions proved in V51.}
Proposition~\ref{prop:V51-exact-halfline} gives the full sourced half-line
solution and traction.  Theorem~\ref{thm:V51-sharp-weight-polygon} proves
the necessary-and-sufficient two-weight region for a uniform
\(H^2\) limiting-absorption resolvent, including the exact trace threshold.
Theorems~\ref{thm:V51-source-complete-reduction} and~\ref{thm:V51-full-inhomogeneous-estimate}
assemble both bath sources into the V50 metric inverse and control all
metric variables, bath fields, and tractions.  Corollary~%
\ref{cor:V51-scaled-estimate} proves the scale-covariant estimate with no
extra tangential derivative loss.  Theorem~\ref{thm:V51-adjoint-Green}
constructs the incoming adjoint and proves its exact boundary Green
incidence.

\paragraph{Imported conclusions not repeated.}
V41 supplies \(HC=J^T\), V47 supplies the two local positive bath actions
and common trace, and V50 supplies the uniform five-variable metric floor.
V51 begins with the inhomogeneous half-line Euler equations and proves the
new resolvent, source, weight, and adjoint statements directly.

\paragraph{Open rows after V51.}
\begin{enumerate}
\item The bulk inhomogeneous harmonic Einstein resolvent has not yet been
combined with \eqref{eq:V51-scaled-estimate}.
\item Maximality of the full Einstein--bath Green domain and the resulting
time-domain generator remain unproved.
\item The full-embedding ghost domain, nonlinear covariant bath stress,
Bianchi identity, and corner current remain open.
\item Higgs, fermion, neutrino, and Yang--Mills sources have not yet been
attached to the complete action incidence.
\end{enumerate}

\paragraph{Honest continuum status.}
The explicit local bath now has a rigorous, sharp, source-complete
limiting-absorption theory at the frozen boundary-symbol level, including
its incoming adjoint.  This closes the inhomogeneous half-line row left by
V50.  It does not prove a maximal time evolution or the nonlinear coupled
Standard Model--Einstein construction.  The full Standard Model--Einstein
continuum is therefore not proved.

\section*{Version V51 append-only record}
\addcontentsline{toc}{section}{Version V51 append-only record}
The complete V50 mathematical source was retained before this continuation.
Its SHA--256 digest was
\begin{center}\scriptsize\ttfamily
67482e0544826d3da57c181581ffc3cd91daf39dc29f247d7a8ea201833c24cb.
\end{center}
No mathematical content of V50 was altered.  On the next iteration, retain
this full file, append before its unique active document-closing command,
increment the filename to \texttt{\_V52.tex}, and remove only the
superseded V51 file from this exact SM note series.  The next scheduled
companion audit remains V55.

% END APPENDED VERSION V51 MATERIAL

% BEGIN APPENDED VERSION V52 MATERIAL

\clearpage
\section{V52: full harmonic-Einstein half-space resolvent, maximal normal
Green relation, and the positive-generator gate}
\label{sec:V52-halfspace}

\subsection{Purpose and exact scope}
\label{subsec:V52-purpose}

V51 solved the sourced auxiliary half-line equations and proved the sharp
two-weight limiting-absorption polygon.  It stopped before attaching that
calculus to the Einstein bulk.  V52 performs this attachment for the
frozen planar generalized-harmonic system already fixed in V36--V41.

The principal de Donder Einstein equation is componentwise the wave
equation after the invertible trace-reversal map.  After Laplace transform
in time and Fourier transform in the two boundary-space directions, its
normal equation is therefore the same Dirichlet half-line equation as in
V51, now with root
\(\lambda=\sqrt{s^2+|k|^2}\).  This observation is used with all ten metric
components present.  The five trace-free induced components have
homogeneous Dirichlet trace.  The remaining five components carry the V41
DeWitt graph and the two local V47 baths.  An interior Einstein source
changes the metric conormal by the exact load
\(\ell_\lambda(f)\); it is inserted before the V50 inverse is used.

The resulting theorem is a source-complete, scale-covariant frozen
half-space resolvent through the bulk continuous spectrum, both bath
thresholds, and gravitational glancing.  V52 also proves that the
homogeneous metric--two-bath boundary relation is a maximal Lagrangian
subspace of the complete normal Green phase space.  This identifies the
formal incoming adjoint domain without a missing boundary covector.

Neither statement supplies a positive Hilbert-space time generator.  The
Einstein action form has the V36 DeWitt signature \((6,4)\), and the
limiting-absorption estimate uses reciprocal spatial weights.  V52
therefore separates maximality of the normal Green relation from maximal
dissipativity of a positive physical evolution.

This is not a fifth-version companion audit.  The next scheduled NS/YM
audit remains V55.

\subsection{The complete frozen normal equations}
\label{subsec:V52-normal-equations}

Let \(\xi\ge0\) be the inward half-space coordinate, so the outward normal
derivative at \(\xi=0\) is
\begin{equation}
 D_n=-\partial_\xi.
 \label{eq:V52-outward-normal}
\end{equation}
Use the stable tangential transform
\begin{equation}
 e^{st+ik\cdot y},
 \qquad
 \operatorname{Re}s\ge0,
 \qquad
 r:=|k|,
 \qquad
 \lambda:=\sqrt{s^2+r^2},
 \label{eq:V52-stable-data}
\end{equation}
with the retarded root inherited from \(\operatorname{Re}s>0\).
The V36 principal Euler equation
\(\mathcal P\Box h=\mathcal P F_E\) is equivalent, after applying
\(\mathcal P^2=I\), to ten scalar normal equations.  Resolve the invariant
metric amplitude into
\begin{equation}
 \widehat h(\xi)
 =\bigl(u^{\rm TF}(\xi),x(\xi)\bigr),
 \qquad
 x=(\tau,v_0,v_1,v_2,\phi)^T,
 \label{eq:V52-metric-splitting}
\end{equation}
where \(u^{\rm TF}\in\mathbb C^5\) contains the trace-free induced
components and \(x\in\mathbb C^5\) is the V41 reduced trace block.  Write
the transformed, trace-reversed equation source in the same splitting as
\((f^{\rm TF},f)\).  The bulk equations are
\begin{align}
 (-\partial_\xi^2+\lambda^2)u^{\rm TF}
 &=f^{\rm TF},
 &u^{\rm TF}(0)&=0,
 \label{eq:V52-TF-equation}\\
 (-\partial_\xi^2+\lambda^2)x
 &=f,
 &x(0)&=x_0.
 \label{eq:V52-trace-equation}
\end{align}
The homogeneous trace-free boundary value is the actual CM action row,
not an elimination of its normal momentum.

Retain the two V51 sourced baths
\begin{align}
 (-\partial_a^2+\kappa_j^2)q_j&=g_j,
 &q_j(0)&=Jx_0,
 \qquad j=1,2,
 \label{eq:V52-bath-equations}\\
 \kappa_j&=\sqrt{s^2+v_j^2r^2},
 &(v_1,v_2)&=\left(\frac12,\frac9{10}\right),
 \label{eq:V52-bath-roots}
\end{align}
with \(M_1=I_3\), \(M_2=\operatorname{diag}(5,5,3)\), and
\begin{equation}
 t_j:=-M_j\partial_aq_j(0).
 \label{eq:V52-bath-traction}
\end{equation}
The inhomogeneous reduced metric boundary row is
\begin{equation}
 D_nx(0)-A(d)x_0-C(t_1+t_2)=F_\partial,
 \label{eq:V52-metric-boundary-row}
\end{equation}
where \(A(d)\), \(C=H^{-1}J^T\), and the row orientation are exactly those
of V41 and the positive V50 branch.

\begin{firewall}[Source normalization]
\label{fire:V52-source-normalization}
The symbol \(f\) in \eqref{eq:V52-trace-equation} is the equation-normalized
source after applying the invertible trace-reversal map.  The physical
linearized stress source is obtained by the inverse fixed fibre map.
No positivity of \(f\) is assumed, and no Bianchi or matter conservation
identity is inferred from the normal resolvent.
\end{firewall}

\subsection{Exact full reduction before estimates}
\label{subsec:V52-exact-reduction}

Use the V51 operators \(R_D(\zeta)\) and
\(\ell_\zeta\) for either normal coordinate.  For the bulk source put
\begin{equation}
 \mathcal L_\lambda f
 :=\ell_\lambda(f)
 =\int_0^\infty e^{-\lambda\xi}f(\xi)\,\dd\xi,
 \label{eq:V52-bulk-load}
\end{equation}
componentwise, and retain the complete bath load
\begin{equation}
 \Lambda_g
 :=\sum_{j=1}^2M_j\ell_{\kappa_j}(g_j).
 \label{eq:V52-bath-load}
\end{equation}

\begin{theorem}[Exact source-complete half-space reduction]
\label{thm:V52-exact-reduction}
For every open stable frequency, and by limiting absorption on the closed
stable sheet, equations
\eqref{eq:V52-TF-equation}--\eqref{eq:V52-metric-boundary-row} are
equivalent to
\begin{align}
 u^{\rm TF}(\xi)
 &=R_D(\lambda)f^{\rm TF}(\xi),
 \label{eq:V52-TF-solution}\\
 x(\xi)
 &=e^{-\lambda\xi}x_0+R_D(\lambda)f(\xi),
 \label{eq:V52-x-solution}\\
 q_j(a)
 &=e^{-\kappa_ja}Jx_0+R_D(\kappa_j)g_j(a),
 \label{eq:V52-q-solution}\\
 t_j
 &=\kappa_jM_jJx_0-M_j\ell_{\kappa_j}(g_j),
 \label{eq:V52-t-solution}\\
 L_*(s,k)x_0
 &=F_\partial+\mathcal L_\lambda f-C\Lambda_g.
 \label{eq:V52-boundary-reduction}
\end{align}
Here \(L_*=L_0-CZ_*J\) is the direct five-row matrix of V50.  In
particular, the bulk Einstein source, both interior bath sources, and the
external boundary datum occur in the final metric equation with their
exact signs.
\end{theorem}

\begin{proof}
Apply Proposition~\ref{prop:V51-exact-halfline} to
\eqref{eq:V52-TF-equation} and
\eqref{eq:V52-trace-equation}.  Since \(D_n=-\partial_\xi\),
\begin{align}
 D_nu^{\rm TF}(0)
 &=-\ell_\lambda(f^{\rm TF}),
 \label{eq:V52-TF-conormal}\\
 D_nx(0)
 &=\lambda x_0-\ell_\lambda(f).
 \label{eq:V52-x-conormal}
\end{align}
The same proposition applied in the \(a\)-coordinate gives
\eqref{eq:V52-q-solution} and \eqref{eq:V52-t-solution}.  Insert these
identities into \eqref{eq:V52-metric-boundary-row}:
\[
 \bigl(\lambda I_5-A(d)-CZ_*J\bigr)x_0
 -\ell_\lambda(f)+C\Lambda_g=F_\partial.
\]
The matrix in parentheses is \(L_*\).  Rearrangement proves
\eqref{eq:V52-boundary-reduction}.  Every elimination used an explicit
Green kernel; none used a formal inverse at \(\lambda=0\) or
\(\kappa_j=0\).
\end{proof}

\begin{remark}[The bulk source sign is an outward-normal effect]
\label{rem:V52-bulk-source-sign}
The V51 bath traction is
\(-M_jq_j'(0)=\kappa_jM_jJx_0-M_j\ell_{\kappa_j}(g_j)\).
For the Einstein half-space the same derivative formula is followed by
\(D_n=-\partial_\xi\), producing
\(D_nx(0)=\lambda x_0-\ell_\lambda(f)\).  Moving this last load to the
right of the boundary row gives the plus sign in
\eqref{eq:V52-boundary-reduction}.  This sign is fixed before the V50
matrix inverse is invoked.
\end{remark}

\subsection{The full normalized half-space estimate}
\label{subsec:V52-normalized-estimate}

On a normal half-line use the V51 norms
\[
 \|g\|_{L^2_\alpha}^2
 =\int_0^\infty\langle z\rangle^{2\alpha}|g(z)|^2\,\dd z,
 \qquad
 \|u\|_{H^2_{-\beta}}^2
 =\sum_{m=0}^2\int_0^\infty
 \langle z\rangle^{-2\beta}|\partial_z^m u(z)|^2\,\dd z.
\]
For the metric norm, sum these componentwise over all ten components.

Here and below \emph{outgoing} is part of the solution class: it means
the boundary value of \(E_\zeta c+R_D(\zeta)g\) obtained from
\(\operatorname{Re}\zeta>0\).  At \(\zeta=0\) this prescription excludes
the independent affine homogeneous solution \(az\).  When
\(\beta>3/2\), that affine function can itself lie in
\(H^2_{-\beta}\), so membership in the weighted space alone is not a
radiation condition.

\begin{theorem}[Frozen harmonic-Einstein--bath resolvent]
\label{thm:V52-full-resolvent}
Assume
\begin{equation}
 \alpha>\frac12,\qquad
 \beta>\frac12,\qquad
 \alpha+\beta\ge2.
 \label{eq:V52-weight-polygon}
\end{equation}
There is \(C_{\alpha,\beta}^{\rm EEB}<\infty\) such that, for every
normalized closed stable frequency
\begin{equation}
 \operatorname{Re}s\ge0,
 \qquad |s|^2+r^2=1,
 \label{eq:V52-normalization}
\end{equation}
the system \eqref{eq:V52-TF-equation}--%
\eqref{eq:V52-metric-boundary-row} has a unique outgoing
limiting-absorption solution and
\begin{align}
 &|x_0|
 +\|\widehat h\|_{H^2_{-\beta}(\dd\xi)}
 +|D_n\widehat h(0)|
 \nonumber\\
 &\qquad
 +\sum_{j=1}^2
 \left(
  \|q_j\|_{H^2_{-\beta}(\dd a)}+|t_j|
 \right)
 \nonumber\\
 &\quad\le
 C_{\alpha,\beta}^{\rm EEB}
 \left(
  |F_\partial|
  +\|f^{\rm TF}\|_{L^2_\alpha(\dd\xi)}
  +\|f\|_{L^2_\alpha(\dd\xi)}
  +\sum_{j=1}^2\|g_j\|_{L^2_\alpha(\dd a)}
 \right).
 \label{eq:V52-full-resolvent}
\end{align}
The estimate is uniform at \(\lambda=0\), at \(\kappa_j=0\), and on every
nonzero imaginary normal root.  The weight region
\eqref{eq:V52-weight-polygon} is sharp already for any one Einstein or
bath component.
\end{theorem}

\begin{proof}
On \eqref{eq:V52-normalization},
\[
 |\lambda|^2=|s^2+r^2|\le1,
 \qquad
 |\kappa_j|^2=|s^2+v_j^2r^2|\le1.
\]
The V51 load bound therefore applies uniformly to
\(\mathcal L_\lambda f\) and to every term in \(\Lambda_g\).
Equation~\eqref{eq:V52-boundary-reduction} and the V50 metric floor give
\begin{equation}
 |x_0|
 \le C_\alpha
 \left(
  |F_\partial|+\|f\|_{L^2_\alpha}
  +\sum_j\|g_j\|_{L^2_\alpha}
 \right).
 \label{eq:V52-x0-bound}
\end{equation}

The sharp V51 half-line theorem bounds
\(R_D(\lambda)f^{\rm TF}\),
\(R_D(\lambda)f\), and
\(R_D(\kappa_j)g_j\) in \(H^2_{-\beta}\).
Because \(\beta>1/2\), the homogeneous extensions
\(e^{-\lambda\xi}x_0\) and
\(e^{-\kappa_ja}Jx_0\) have uniformly bounded
\(H^2_{-\beta}\) norms.  Equations
\eqref{eq:V52-TF-solution}--\eqref{eq:V52-q-solution} and
\eqref{eq:V52-x0-bound} therefore control all three half-line fields.

The conormals \eqref{eq:V52-TF-conormal}--%
\eqref{eq:V52-x-conormal} are bounded by
\(|x_0|\) and the V51 load estimate.  The bath tractions are bounded by
\eqref{eq:V52-t-solution}.  This proves
\eqref{eq:V52-full-resolvent}.  Uniqueness follows in reverse: zero data
make \eqref{eq:V52-boundary-reduction} equal to \(L_*x_0=0\), so V50 gives
\(x_0=0\); the unique outgoing scalar half-line solutions then vanish.

Sharpness follows by setting all components but one to zero and using the
three necessity tests in
Theorem~\ref{thm:V51-sharp-weight-polygon}.  The boundary coupling cannot
improve a compactly supported constant tail, the unbounded threshold load,
or the translated-block sequence of a decoupled trace-free Dirichlet
component.
\end{proof}

\begin{firewall}[Weighted Kreiss scope]
\label{fire:V52-weighted-Kreiss}
Theorem~\ref{thm:V52-full-resolvent} is a complete frozen
Laplace--Fourier half-space estimate in explicitly declared reciprocal
normal weights.  It is stronger than a homogeneous Lopatinski determinant
and includes interior forcing.  It is not the unweighted
\((\operatorname{Re}s)^{-1}\) Hilbert resolvent estimate required by the
Hille--Yosida theorem.  V50 already proves why such an unweighted estimate
cannot hold for generalized bath eigenfunctions on the closed imaginary
axis.
\end{firewall}

\subsection{Scale-covariant full estimate}
\label{subsec:V52-scaled}

For a nonzero tangential frequency put
\begin{equation}
 \rho=(|s|^2+r^2)^{1/2}.
 \label{eq:V52-rho}
\end{equation}
Use the V51 scale-covariant source and receiver norms
\(\mathcal Y_{\alpha,\rho}\) and
\(\mathcal X_{-\beta,\rho}\) on both the \(\xi\)- and \(a\)-half-lines.

\begin{corollary}[No extra derivative loss in the full half-space system]
\label{cor:V52-scaled-resolvent}
For every \(\rho>0\),
\begin{align}
 &|x_0|
 +\|\widehat h\|_{\mathcal X_{-\beta,\rho}(\dd\xi)}
 +\rho^{-1}|D_n\widehat h(0)|
 \nonumber\\
 &\qquad
 +\sum_{j=1}^2
 \left(
  \|q_j\|_{\mathcal X_{-\beta,\rho}(\dd a)}
  +\rho^{-1}|t_j|
 \right)
 \nonumber\\
 &\quad\le
 C_{\alpha,\beta}^{\rm EEB}
 \left(
  \rho^{-1}|F_\partial|
  +\|f^{\rm TF}\|_{\mathcal Y_{\alpha,\rho}(\dd\xi)}
  +\|f\|_{\mathcal Y_{\alpha,\rho}(\dd\xi)}
  +\sum_{j=1}^2
    \|g_j\|_{\mathcal Y_{\alpha,\rho}(\dd a)}
 \right).
 \label{eq:V52-scaled-resolvent}
\end{align}
Thus the complete frozen Einstein--bath boundary problem has the natural
order-one boundary and order-two interior scaling.  The only loss is the
sharp normal-weight polygon.
\end{corollary}

\begin{proof}
Rescale both normal coordinates by
\(\Xi=\rho\xi\), \(A=\rho a\).  Rescale every second-order equation source
by \(\rho^{-2}\), every normal traction by \(\rho^{-1}\), and the
order-one boundary datum by \(\rho^{-1}\).  Then
\[
 \widehat\lambda=\lambda/\rho,
 \qquad
 \widehat\kappa_j=\kappa_j/\rho,
 \qquad
 |\widehat s|^2+\widehat r^2=1.
\]
The change-of-variable identities are exactly those proved in
Corollary~\ref{cor:V51-scaled-estimate}.  Apply
Theorem~\ref{thm:V52-full-resolvent} to the normalized system.
\end{proof}

\subsection{The full-embedding quotient estimate}
\label{subsec:V52-embedding-quotient}

At a stable frequency let the V39 extended boundary amplitude be
\((h,X)\), with
\begin{equation}
 \widehat h=Q_4(h,X)=h+\mathcal S_pX,
 \qquad
 G_pq=(\mathcal S_pq,-q).
 \label{eq:V52-Q4}
\end{equation}
The CM rows and the bath trace are functions of \(\widehat h\).

\begin{corollary}[Exact quotient resolvent and canonical representative]
\label{cor:V52-quotient-resolvent}
Define the frozen quotient norm by
\begin{equation}
 \|[(h,X)]\|_Q:=\|Q_4(h,X)\|.
 \label{eq:V52-quotient-norm}
\end{equation}
The estimates
\eqref{eq:V52-full-resolvent} and
\eqref{eq:V52-scaled-resolvent} hold verbatim for the extended
metric--embedding class in this quotient norm.  The canonical splitting
\begin{equation}
 R_4(\widehat h)=(\widehat h,0)
 \label{eq:V52-canonical-splitting}
\end{equation}
realizes a representative with the same controlled metric norm.  No
estimate holds for an arbitrary raw representative without a gauge norm,
because the exact gauge orbit \(G_pq\) is invisible to all quotient rows.
\end{corollary}

\begin{proof}
V39 proves
\[
 Q_4G_p=0,\qquad Q_4R_4=I,
\qquad
 (h,X)=R_4Q_4(h,X)+G_p(-X).
\]
The V39 pulled-back geometric rows and the V47 common bath trace depend
only on \(Q_4(h,X)\).  Apply the V52 estimate to this invariant amplitude.
The displayed decomposition proves both the canonical representative
claim and the impossibility of controlling a free gauge summand from
quotient data.
\end{proof}

\begin{firewall}[On-shell ghost tangency is not an off-shell quotient norm]
\label{fire:V52-ghost-quotient}
The canonical representative \(X=0\) corresponds at the boundary to
removing the V39 embedding gauge coordinate.  V37 proves tangency of the
CM rows for Dirichlet ghosts only after the ghost Euler equation is used.
Corollary~\ref{cor:V52-quotient-resolvent} therefore proves the analytic
metric quotient estimate; it does not construct the off-shell
ghost--antighost--Nakanishi--Lautrup action domain or identify a positive
BRST cohomology norm.
\end{firewall}

\subsection{Maximality of the complete normal Green relation}
\label{subsec:V52-maximal-Green}

At one real tangential frequency, or after tangential integration for
compact support, write the normal phase coordinates as follows:
\begin{itemize}
\item \((u,\pi)\in\mathbb C^5\oplus\mathbb C^5\) for the trace-free
      induced metric and its conormal;
\item \((x,p)\in\mathbb C^5\oplus\mathbb C^5\) for the reduced metric
      trace and normal derivative;
\item \((z_j,t_j)\in\mathbb C^3\oplus\mathbb C^3\) for the two bath
      traces and tractions.
\end{itemize}
Use the total normal Green form
\begin{align}
 \Omega_{\rm tot}
 &:=
 \pi^*u'-u^*\pi'
 +p^*Hx'-x^*Hp'
 \nonumber\\
 &\quad
 +\sum_{j=1}^2
 \left(z_j^*t_j'-t_j^*z_j'\right).
 \label{eq:V52-total-Green}
\end{align}
The bath sign is the orientation induced from the boundary
\(a=0\); it is opposite to the metric outward conormal orientation.

Define the homogeneous normal relation
\begin{equation}
 \mathcal L_{\rm EEB}
 :=\left\{
 \begin{array}{l}
 u=0,\\
 z_1=Jx,\quad z_2=Jx,\\
 p=A(D)x+C(t_1+t_2)
 \end{array}
 \right\}.
 \label{eq:V52-Green-relation}
\end{equation}
No condition is imposed separately on \(\pi,t_1,t_2\).

\begin{theorem}[Maximal normal Green relation]
\label{thm:V52-maximal-Green}
For the explicit positive Grams of V50,
\(\mathcal L_{\rm EEB}\) is Lagrangian for
\(\Omega_{\rm tot}\).  Equivalently,
\begin{equation}
 \mathcal L_{\rm EEB}^{\perp_{\Omega_{\rm tot}}}
 =\mathcal L_{\rm EEB}.
 \label{eq:V52-self-adjoint-relation}
\end{equation}
Thus the frozen homogeneous incoming adjoint has exactly the same complete
normal boundary relation; no additional bath-difference traction or metric
source row is missing.
\end{theorem}

\begin{proof}
The matrices \(H,M_1,M_2\) are nondegenerate, so
\(\Omega_{\rm tot}\) is nondegenerate on the
\[
 10+10+6+6=32
\]
dimensional normal phase space.  Let two phase vectors lie in
\(\mathcal L_{\rm EEB}\).  The trace-free contribution vanishes because
\(u=u'=0\).  V37 proves, after tangential integration, that
\[
 (A(D)x)^*Hx'-x^*HA(D)x'=0.
\]
The identity \(HC=J^T\) gives the remaining metric contribution
\[
 (t_1+t_2)^*Jx'
 -(Jx)^*(t_1'+t_2').
\]
The two bath contributions, using \(z_j=Jx\), sum to its negative:
\[
 (Jx)^*(t_1'+t_2')
 -(t_1+t_2)^*Jx'.
\]
Hence \(\Omega_{\rm tot}\) vanishes on the relation.

The equations \(u=0\) give five independent conditions.  The two common
trace equations give six more, and the metric momentum graph gives five
independent conditions after those configuration equations are imposed.
Thus \(\mathcal L_{\rm EEB}\) has codimension sixteen and dimension
sixteen.  An isotropic half-dimensional subspace of a nondegenerate
symplectic space is Lagrangian, proving
\eqref{eq:V52-self-adjoint-relation}.
\end{proof}

\begin{corollary}[Outgoing/incoming adjoint closure with interior sources]
\label{cor:V52-adjoint-closure}
Pair an outgoing solution of the V52 transformed equations with the
incoming solution obtained by conjugating all three retarded normal roots.
The V51 Green identity, applied once to each of the ten Einstein
components and once to each bath, leaves only the external metric boundary
data pairing.  For homogeneous \(F_\partial=0\), every internal boundary
term cancels by Theorem~\ref{thm:V52-maximal-Green}.  Hence the incoming
adjoint domain has no boundary row beyond
\eqref{eq:V52-Green-relation}.
\end{corollary}

\begin{proof}
Apply Theorem~\ref{thm:V51-adjoint-Green} with \(M=I_5\) to the trace-free
block, with the reduced DeWitt matrix \(H\) to the trace block, and with
\(M=M_j\) to the baths.  The scalar Green identity is valid for any
constant symmetric nondegenerate coefficient matrix because the normal
operator is scalar in the fibre.  Sum the identities.  The trace-free
Dirichlet terms vanish, and the remaining internal boundary expression is
exactly \(\Omega_{\rm tot}\) on
\(\mathcal L_{\rm EEB}\).  Theorem~\ref{thm:V52-maximal-Green} annihilates
it and proves maximality of the homogeneous adjoint domain.
\end{proof}

\begin{assessment}[What maximality now means]
\label{ass:V52-maximality-scope}
Theorem~\ref{thm:V52-maximal-Green} is the missing normal-domain
maximality statement for the local quadratic principal action.  It is
stronger than merely showing cancellation on selected solutions: the
dimension count proves that the relation equals its full Green
orthogonal.  Together with the V52 resolvent, it closes the frozen
frequency-domain bulk--boundary adjoint row.

The theorem concerns the normal Green phase space after tangential
integration.  Because \(A(D)\) contains a tangential time derivative, the
same formula must be rewritten in a first-order time Hamiltonian phase
space before it can be called the domain of a time-evolution generator.
\end{assessment}

\subsection{The positive-generator obstruction and next upstream row}
\label{subsec:V52-generator}

\begin{theorem}[Why the present proof does not imply maximal
dissipativity]
\label{thm:V52-no-Hilbert-generator-yet}
The V52 results do not establish that the coupled time-evolution operator
is maximal dissipative on a positive Hilbert energy space.  Two independent
requirements remain:
\begin{enumerate}
\item the V36 metric action fibre form has signature \((6,4)\), so the
      Green self-adjointness in
      \eqref{eq:V52-self-adjoint-relation} is not positivity of the
      harmonic Einstein energy;
\item the estimate \eqref{eq:V52-full-resolvent} maps positive normal
      source weights to negative normal receiver weights and therefore is
      not an unweighted Hille--Yosida resolvent bound.
\end{enumerate}
Consequently neither the Lagrangian dimension count nor the weighted
closed-cone resolvent can be promoted, without an additional physical
symmetrizer and constraint quotient, to a unitary or contraction
semigroup theorem.
\end{theorem}

\begin{proof}
The first statement is the exact signature calculation of
Lemma~\ref{lem:V36-DeWitt-involution}; adding the nonnegative bath energies
does not remove a negative metric subspace on which all bath traces vanish.
The second statement follows directly from the sharp V51 polygon.  At the
continuous spectrum, a boundary-forced outgoing wave is not in the
unweighted half-line energy space, as proved in V50.  The Hille--Yosida
criterion instead requires a resolvent on one fixed positive Hilbert space
with the appropriate \((\operatorname{Re}s)^{-1}\) bound.  Neither
condition is supplied by an indefinite Green form or by a map between two
different weighted spaces.  Hence the claimed implication is unavailable.
\end{proof}

V53 must now move upstream to the positive physical phase space rather
than repeat the normal determinant.  It should:
\begin{enumerate}
\item write the first-order-in-time harmonic Einstein system together with
      the two local baths and the boundary embedding variables;
\item separate the propagated harmonic constraints and pure gauge
      directions by an explicit bounded projector compatible with the
      V52 domain;
\item compute the induced energy on the quotient and decide whether it is
      positive and whether the boundary relation is maximal
      dissipative or skew-adjoint there;
\item if positivity fails, return the exact negative quotient vector and
      move to the boundary action or gauge-fixing functional rather than
      infer a semigroup from the weighted resolvent; and
\item only after a positive linear generator is proved, restore the
      nonlinear covariant bath stress, Bianchi identity, complete ghost
      domain, and Standard-Model sources.
\end{enumerate}

\subsection{Finite certificates}
\label{subsec:V52-certificates}

For an exponential Einstein source
\[
 f(\xi)=e^{-\varrho\xi}c,
 \qquad \operatorname{Re}\varrho>0,
\]
the exact bulk load and zero-trace solution are
\begin{equation}
 \mathcal L_\lambda f=\frac{c}{\lambda+\varrho},
 \qquad
 R_D(\lambda)f(\xi)
 =\frac{e^{-\varrho\xi}-e^{-\lambda\xi}}
 {\lambda^2-\varrho^2}c.
 \label{eq:V52-exp-bulk-card}
\end{equation}
At \(\lambda=\varrho\), the displayed quotient denotes its continuous
limit \(\xi e^{-\lambda\xi}c/(2\lambda)\).
The outward conormal is
\(-\mathcal L_\lambda f\).  Inserting this card and the V51 exponential
bath cards into \eqref{eq:V52-boundary-reduction} makes the ten bulk ODE
residuals, six common-trace residuals, five metric graph residuals, and all
three traction signs exact algebraic zeros.

The dimension certificate for
Theorem~\ref{thm:V52-maximal-Green} is
\begin{equation}
 \dim\mathcal P_{\rm tot}=32,
 \qquad
 \operatorname{rank}B_{\rm EEB}=5+6+5=16,
 \qquad
 \dim\ker B_{\rm EEB}=16.
 \label{eq:V52-dimension-card}
\end{equation}
A finite implementation must also store the residual
\[
 \sup_{\substack{\zeta,\zeta'\in\ker B_{\rm EEB}\\
                  \|\zeta\|=\|\zeta'\|=1}}
 |\Omega_{\rm tot}(\zeta,\zeta')|,
\]
which is exactly zero by the cancellation proof.  Rank sixteen without
this Green residual would not certify a Lagrangian domain.

\subsection{V52 ledger}
\label{subsec:V52-ledger}

\paragraph{New conclusions proved in V52.}
Theorem~\ref{thm:V52-exact-reduction} inserts the complete interior
Einstein source into the direct V50 metric equation with the exact
outward-normal sign.  Theorem~\ref{thm:V52-full-resolvent} proves a sharp,
source-complete frozen half-space estimate for all ten harmonic metric
components, both local baths, and every normal traction.  Corollary~%
\ref{cor:V52-scaled-resolvent} proves the natural all-frequency scaling.
Corollary~\ref{cor:V52-quotient-resolvent} lifts the estimate to the exact
V39 full-embedding quotient.  Theorem~\ref{thm:V52-maximal-Green} proves
that the complete homogeneous metric--bath normal relation is Lagrangian
and maximal, and Corollary~\ref{cor:V52-adjoint-closure} identifies its
incoming adjoint boundary domain.

\paragraph{Imported conclusions not repeated.}
V36 supplies the componentwise principal wave equation and DeWitt
signature.  V37/V41 supply the CM graph, \(H,A,C,J\), and its tangential
Green cancellation.  V39 supplies the split full-embedding quotient.
V47 supplies the two local bath actions.  V50 supplies the uniform direct
five-row metric floor, and V51 supplies the sharp scalar
limiting-absorption and adjoint kernels.

\paragraph{Honest continuum status.}
The frozen linear harmonic Einstein--bath half-space problem now has a
rigorous source-complete weighted resolvent and a maximal normal Green
relation, including the full embedding quotient.  A positive physical
Hilbert symmetrizer, an unweighted time-domain generator, the off-shell
ghost domain, nonlinear covariant stress/Bianchi closure, and the coupled
Higgs, fermion, neutrino, and Yang--Mills sectors remain open.  The full
Standard Model--Einstein continuum is therefore not proved.

\section*{Version V52 append-only record}
\addcontentsline{toc}{section}{Version V52 append-only record}
The complete V51 mathematical source was retained before this continuation.
Its SHA--256 digest was
\begin{center}\scriptsize\ttfamily
0b450614ca0a5032a3632f082d753576d05e4a526d79383aee93d59baa460b04.
\end{center}
No mathematical content of V51 was altered.  On the next iteration, retain
this full file, append before its unique active document-closing command,
increment the filename to \texttt{\_V53.tex}, and remove only the
superseded V52 file from this exact SM note series.  The next scheduled
companion audit remains V55.

% END APPENDED VERSION V52 MATERIAL

% BEGIN APPENDED VERSION V53 MATERIAL

\clearpage
\section{V53: null-shell radiative cohomology, TT positive energy, and
the positive-bath orientation obstruction}
\label{sec:V53-radiative-energy}

\subsection{Purpose and nonduplication boundary}
\label{subsec:V53-purpose}

V52 closed the frozen source-complete normal resolvent and its maximal
DeWitt Green relation.  It did not produce a positive time generator.
The first item in its handoff was a first-order harmonic-Einstein
symmetrizer, but that calculation already exists in V29: every
de Donder-reduced metric component is a copy of the wave block with
positive identity time matrix.  Repeating that block would not address
the missing physical quotient.

V53 instead performs two operations which were not present in V29 or
V39.  First, it computes the exact null-shell cohomology of the harmonic
metric amplitude by residual diffeomorphisms and identifies it with the
two transverse-traceless polarizations.  This supplies a positive
radiative symmetrizer on the frozen interior quotient.  Second, it
restricts the V50 port incidence to that radiative fibre and derives the
time-domain power sign of the two V44 Green orientations.  The stable
positive orientation selected in V49--V52 adds, rather than cancels, the
positive metric and bath exchange powers.  The opposite orientation
admits a positive exchange Gram, but the complete V48 theorem already
excludes its stable square-root bath class.

All statements below are linear, frozen, and at a flat anchor.  The
radiative projector is an interior Fourier multiplier.  Its preservation
of the complete half-space boundary domain is tested rather than assumed.

\subsection{The imported first-order wave state and its constraint shell}
\label{subsec:V53-first-order}

Let \(\bar h=\mathcal Ph\) be trace reversal and, on a spacelike slice,
write
\begin{equation}
 \pi_{\mu\nu}:=\partial_t\bar h_{\mu\nu},
 \qquad
 w_{i\mu\nu}:=\partial_i\bar h_{\mu\nu}.
 \label{eq:V53-wave-jets}
\end{equation}
The specialization of the V29 wave bank is
\begin{align}
 \partial_t\bar h&=\pi,
 &
 \partial_t\pi&=\partial^iw_i,
 &
 \partial_tw_i&=\partial_i\pi.
 \label{eq:V53-first-order-wave}
\end{align}
Its positive gauge-fixed energy is
\begin{equation}
 E_{\rm gf}[\bar h]
 :=\frac12\int_{\mathbb R^3}
 \left(|\pi|_{\rm E}^2+\sum_{i=1}^3|w_i|_{\rm E}^2\right)\dd^3x.
 \label{eq:V53-gf-energy}
\end{equation}
This is precisely the identity symmetrizer of
Lemma~\ref{lem:V29-wave-speeds}, applied to the Einstein item of
Proposition~\ref{prop:V29-wave-sectors}.  It is recorded here only to fix
the variables used in the quotient.

The harmonic covector and the first-order jet defects are
\begin{align}
 C_\nu&=-\pi_{0\nu}+\partial^i\bar h_{i\nu}
       =-\pi_{0\nu}+w^i{}_{i\nu},
 \label{eq:V53-harmonic-constraint}\\
 {\cal J}_{i\mu\nu}&:=w_{i\mu\nu}-\partial_i\bar h_{\mu\nu},
 &
 {\cal K}_{ij\mu\nu}&:=\partial_iw_{j\mu\nu}
                      -\partial_jw_{i\mu\nu}.
 \label{eq:V53-jet-constraints}
\end{align}
The evolution \eqref{eq:V53-first-order-wave} preserves
\({\cal J}={\cal K}=0\), and \(C\) obeys the subsidiary wave equation
already derived in V34.  Thus \(C|_{t=0}=\partial_tC|_{t=0}=0\) selects
the harmonic solution shell.  The positive form
\eqref{eq:V53-gf-energy} is not yet physical: a nonzero residual gauge
wave has positive Euclidean jet norm.

\subsection{Exact null-shell radiative cohomology}
\label{subsec:V53-null-cohomology}

Fix a nonzero spatial covector \(\xi\in\mathbb R^3\), put
\(\omega=|\xi|\), and choose one of the two null branches
\begin{equation}
 k_\mu=(-\omega,\xi_i),
 \qquad
 k^\mu=(\omega,\xi^i),
 \qquad
 k^\mu k_\mu=0.
 \label{eq:V53-null-covector}
\end{equation}
Let
\begin{equation}
 Z_k
 :=\left\{\bar a\in\operatorname{Sym}^2\mathbb C^4:
             k^\mu\bar a_{\mu\nu}=0\right\}.
 \label{eq:V53-harmonic-amplitudes}
\end{equation}
For a residual wave-gauge amplitude \(\epsilon\in\mathbb C^4\), omit the
irrelevant common Fourier factor \(i\) and define
\begin{equation}
 (\Gamma_k\epsilon)_{\mu\nu}
 :=k_\mu\epsilon_\nu+k_\nu\epsilon_\mu
   -\eta_{\mu\nu}k^\rho\epsilon_\rho.
 \label{eq:V53-residual-gauge-map}
\end{equation}
This is the trace-reversed form of
\(\delta h_{\mu\nu}=2\partial_{(\mu}\epsilon_{\nu)}\).

Let
\begin{equation}
 P_i{}^j(\xi):=\delta_i{}^j-\widehat\xi_i\widehat\xi^j,
 \qquad \widehat\xi:=\xi/|\xi|,
 \label{eq:V53-transverse-projector}
\end{equation}
and define the spatial TT projector
\begin{equation}
 (\Pi_{\rm TT}b)_{ij}
 :=P_i{}^kP_j{}^\ell b_{k\ell}
   -\frac12P_{ij}P^{k\ell}b_{k\ell}.
 \label{eq:V53-TT-projector}
\end{equation}
Its range \({\cal T}_\xi\) is the two-dimensional space of symmetric
spatial tensors satisfying
\(\xi^ib_{ij}=0\) and \(\delta^{ij}b_{ij}=0\).
Set
\begin{equation}
 T_k:Z_k\longrightarrow{\cal T}_\xi,
 \qquad
 T_k\bar a:=\Pi_{\rm TT}(\bar a_{ij}).
 \label{eq:V53-T-map}
\end{equation}

\begin{theorem}[Split radiative exact sequence]
\label{thm:V53-radiative-exact-sequence}
For every nonzero null covector \eqref{eq:V53-null-covector},
\begin{equation}
 0\longrightarrow\mathbb C^4
 \mathop{\longrightarrow}^{\Gamma_k}
 Z_k
 \mathop{\longrightarrow}^{T_k}
 {\cal T}_\xi
 \longrightarrow0
 \label{eq:V53-radiative-exact-sequence}
\end{equation}
is split exact.  In particular,
\begin{equation}
 \dim Z_k=6,\qquad
 \operatorname{rank}\Gamma_k=4,\qquad
 Z_k/\operatorname{ran}\Gamma_k
 \cong{\cal T}_\xi,\qquad
 \dim{\cal T}_\xi=2.
 \label{eq:V53-radiative-dimensions}
\end{equation}
A norm-one splitting sends a TT tensor \(b\) to the four-tensor with
spatial block \(b\) and all time components zero.
\end{theorem}

\begin{proof}
Contraction of \eqref{eq:V53-residual-gauge-map} gives
\[
 k^\mu(\Gamma_k\epsilon)_{\mu\nu}
 =k^2\epsilon_\nu
  +k_\nu(k^\rho\epsilon_\rho)
  -k_\nu(k^\rho\epsilon_\rho)=0,
\]
so \(\operatorname{ran}\Gamma_k\subset Z_k\).
The harmonic contraction has rank four by
Proposition~\ref{prop:V35-complex-Gram}; hence \(\dim Z_k=10-4=6\).

Rotate the spatial frame so that
\(\xi=(0,0,\omega)\).  Then the harmonic equations are
\[
 \bar a_{3\nu}=-\bar a_{0\nu}.
\]
On the six free coordinates
\(\bar a_{00},\bar a_{01},\bar a_{02},
  \bar a_{11},\bar a_{12},\bar a_{22}\),
the gauge shifts include
\begin{align*}
 \delta\bar a_{00}&=\omega(-\epsilon_0+\epsilon_3),&
 \delta\bar a_{01}&=-\omega\epsilon_1,&
 \delta\bar a_{02}&=-\omega\epsilon_2,\\
 \delta(\bar a_{11}+\bar a_{22})
 &=-2\omega(\epsilon_0+\epsilon_3).
\end{align*}
They set uniquely
\(\bar a_{00},\bar a_{01},\bar a_{02}\), and
\(\bar a_{11}+\bar a_{22}\) to zero.  The harmonic equations then set
every component with an index \(0\) or \(3\) to zero.  The two survivors
are
\(\bar a_{11}=-\bar a_{22}\) and \(\bar a_{12}\), exactly the TT
amplitudes.  This also proves that \(\Gamma_k\) is injective.

Directly, the spatial part of a gauge amplitude is
\[
 (\Gamma_k\epsilon)_{ij}
 =\xi_i\epsilon_j+\xi_j\epsilon_i
  -\delta_{ij}k^\rho\epsilon_\rho.
\]
The two longitudinal terms are killed by \(P\), and the last term is
killed by the transverse trace subtraction.  Thus
\(T_k\Gamma_k=0\).  Conversely, a TT spatial tensor with zero time
components lies in \(Z_k\), so \(T_k\) is onto and the asserted inclusion
is a right inverse.  Exactness and all dimensions follow.
\end{proof}

\begin{corollary}[Uniform TT representative and positive radiative energy]
\label{cor:V53-positive-radiative-energy}
On the unit direction sphere, \(\Pi_{\rm TT}\) is a smooth family of
orthogonal rank-two projectors with operator norm one.  Defining it
arbitrarily at \(\xi=0\) gives a bounded Fourier multiplier on every
\(L^2\)-Sobolev space.  It commutes with the scalar wave evolution.
The quotient energy
\begin{equation}
 E_{\rm rad}([\bar h])
 :=\frac12\int_{\mathbb R^3}
 \left(
  |\partial_t\Pi_{\rm TT}\bar h^{\rm sp}|^2
  +|\nabla\Pi_{\rm TT}\bar h^{\rm sp}|^2
 \right)\dd^3x
 \label{eq:V53-radiative-energy}
\end{equation}
is positive definite on the nonzero-frequency radiative quotient and is
conserved for the frozen interior wave equation.
\end{corollary}

\begin{proof}
Equation~\eqref{eq:V53-TT-projector} is the orthogonal projection onto
the trace-free symmetric square of the two-plane
\(\widehat\xi^\perp\), so its rank is two, its norm is one, and it depends
smoothly on \(\widehat\xi\).  A bounded direction multiplier commutes
with the Sobolev weight and with the scalar Fourier wave multiplier.
Theorem~\ref{thm:V53-radiative-exact-sequence} says that its kernel on
\(Z_k\) is exactly the residual gauge image.  Therefore
\eqref{eq:V53-radiative-energy} vanishes on a radiative class only when
that class is zero.  Conservation is the V29 wave-energy identity
restricted to the invariant TT range.
\end{proof}

\begin{firewall}[Zero frequency and the half-space]
\label{fire:V53-TT-scope}
There is no direction-independent rank-two continuation of
\(\Pi_{\rm TT}\) through \(\xi=0\).  Assigning a value on that
measure-zero set defines an \(L^2\) multiplier but does not classify
spatially homogeneous moduli.  Moreover, \(\Pi_{\rm TT}\) contains the
Riesz multiplier \(\xi_i\xi_j/|\xi|^2\).  On a half-space it requires a
choice of normal realization and is not automatically a local
boundary-preserving projector.  Corollary
\ref{cor:V53-positive-radiative-energy} therefore proves the frozen
interior physical symmetrizer, not invariance of the V52 boundary domain.
\end{firewall}

\subsection{The V39 quotient is not the radiative quotient}
\label{subsec:V53-two-quotients}

The V39 split exact sequence has quotient
\[
 \mathsf E_p/\operatorname{ran}G_p
 \mathop{\cong}^{Q_4}\operatorname{Sym}^2\mathbb C^4,
\]
which has dimension ten.  It removes the redundancy between a metric and
its embedding representative.  By contrast,
\eqref{eq:V53-radiative-exact-sequence} is imposed only after the
harmonic and null-shell equations and then divides by residual wave
diffeomorphisms; its quotient has dimension two.

\begin{proposition}[Exact non-substitution of the two quotients]
\label{prop:V53-quotient-distinction}
The V39 projector \(I-P_{G,p}\) cannot by itself serve as the radiative
projector of V53.  Its quotient map is onto a ten-dimensional metric
fibre, whereas \(T_k\) has the four-dimensional kernel
\(\operatorname{ran}\Gamma_k\) inside the six-dimensional harmonic
solution fibre.  A boundary embedding theory may decide to retain some
of those transformations as edge symmetries, but that is additional
physical data and is not supplied by the algebraic identity \(Q_4G_p=0\).
\end{proposition}

\begin{proof}
The first target dimension and kernel are
Theorem~\ref{thm:V39-full-exact-sequence}.  The second are
Theorem~\ref{thm:V53-radiative-exact-sequence}.  Since their targets and
kernels have different dimensions, the maps cannot coincide or replace
one another.  The last statement records that an edge symmetry is a
choice of boundary phase space, not a consequence of a finite-dimensional
split exact sequence.
\end{proof}

\subsection{The bath port has a nonzero radiative projection}
\label{subsec:V53-radiative-port}

Choose the physical spatial normal \(n=e_3\).  By tangential rotation a
unit propagation direction can be written
\begin{equation}
 \widehat\xi=(\sin\theta,0,\cos\theta).
 \label{eq:V53-oblique-direction}
\end{equation}
Put
\[
 e=(0,1,0),\qquad
 m=(\cos\theta,0,-\sin\theta),
\]
and use the Frobenius-orthonormal TT basis
\begin{equation}
 E_+:=\frac{e\otimes e-m\otimes m}{\sqrt2},
 \qquad
 E_\times:=\frac{e\otimes m+m\otimes e}{\sqrt2}.
 \label{eq:V53-TT-basis}
\end{equation}
For \(a=(a_+,a_\times)^T\), the V41 port trace
\(Jx=(h_{31},h_{32},h_{33})^T\) restricts to
\begin{equation}
 J(a_+E_++a_\times E_\times)=R_\theta a,
 \qquad
 R_\theta
 =\frac{\sin\theta}{\sqrt2}
 \begin{pmatrix}
  \cos\theta&0\\
  0&-1\\
  -\sin\theta&0
 \end{pmatrix}.
 \label{eq:V53-radiative-trace}
\end{equation}

For a port traction \(y=(y_1,y_2,y_3)^T\), the spatial block of the V50
metric incidence \(Cy\) is
\begin{equation}
 (Cy)^{\rm sp}
 =\frac12
 \begin{pmatrix}
 -y_3&0&y_1\\
 0&-y_3&y_2\\
 y_1&y_2&y_3
 \end{pmatrix}.
 \label{eq:V53-C-spatial}
\end{equation}

\begin{theorem}[Exact radiative port incidence]
\label{thm:V53-radiative-port}
In the basis \eqref{eq:V53-TT-basis},
\begin{equation}
 \Pi_{\rm TT}(Cy)^{\rm sp}
 =R_\theta^*y,
 \qquad
 R_\theta^*R_\theta
 =\frac{\sin^2\theta}{2}I_2.
 \label{eq:V53-radiative-adjoint}
\end{equation}
Thus the bath port is invisible to radiative waves at exact normal
incidence, but has rank two and acts on both physical polarizations at
every oblique direction \(0<\theta\le\pi/2\).
\end{theorem}

\begin{proof}
Equation~\eqref{eq:V53-radiative-trace} follows by taking the
\((31),(32),(33)\) entries of \(E_+\) and \(E_\times\).  Contracting
\eqref{eq:V53-C-spatial} with the two TT tensors gives
\begin{align*}
 \langle E_+,(Cy)^{\rm sp}\rangle_{\rm F}
 &=\frac{\sin\theta\cos\theta\,y_1
         -\sin^2\theta\,y_3}{\sqrt2},\\
 \langle E_\times,(Cy)^{\rm sp}\rangle_{\rm F}
 &=-\frac{\sin\theta\,y_2}{\sqrt2}.
\end{align*}
These are the two components of \(R_\theta^*y\).  Multiplication of the
displayed matrix gives
\(R_\theta^*R_\theta=(\sin^2\theta/2)I_2\), proving the rank statement.
\end{proof}

This theorem is an incidence calculation on the radiative fibre.  A
single travelling TT plane wave need not satisfy the five V52
trace-free induced Dirichlet rows; a boundary solution combines incoming
and outgoing amplitudes.  The conclusion used below is narrower and
exact: the V50 force \(C y\) has a nonzero physical TT projection, so its
orientation defect cannot be dismissed as entirely pure gauge.

\subsection{Positive product energy and the orientation sign}
\label{subsec:V53-power-sign}

Let \(p=D_nx\), \(t_\Sigma=t_1+t_2\), and consider any constant real
symmetric positive definite matrix \(K\) on the reduced five-component
metric fibre.  Its bulk wave energy is
\begin{equation}
 E_{{\rm met},K}
 :=\frac12\int_{\xi\ge0}
 \left(
  x_t^TKx_t+\partial_\xi x^TK\partial_\xi x
  +\partial_Ax^TK\partial_Ax
 \right)\dd\xi\,\dd^2y.
 \label{eq:V53-K-metric-energy}
\end{equation}
For the homogeneous wave equation,
\begin{equation}
 \frac{\dd}{\dd t}E_{{\rm met},K}
 =\int_{\partial M}x_t^TKp\,\dd^2y.
 \label{eq:V53-metric-power}
\end{equation}
The sum of the two positive V47 bath energies obeys
\begin{equation}
 \frac{\dd}{\dd t}E_{{\rm bath},(2)}
 =\int_{\partial M}(Jx_t)^Tt_\Sigma\,\dd^2y.
 \label{eq:V53-bath-power}
\end{equation}

For the positive orientation retained in V50--V52,
\begin{equation}
 p=A(D)x+Ct_\Sigma.
 \label{eq:V53-positive-graph}
\end{equation}
Consequently
\begin{equation}
 \frac{\dd}{\dd t}
 \left(E_{{\rm met},K}+E_{{\rm bath},(2)}\right)
 =
 \int_{\partial M}
 \left[
 x_t^TKA(D)x
 +x_t^T(KC+J^T)t_\Sigma
 \right]\dd^2y.
 \label{eq:V53-positive-power}
\end{equation}

\begin{theorem}[Positive-orientation product-energy obstruction]
\label{thm:V53-positive-energy-obstruction}
There is no \(K=K^T\succ0\) for which the metric--bath exchange term in
\eqref{eq:V53-positive-power} cancels.  More strongly, for every such
\(K\), that exchange term takes both signs on smooth compatible
instantaneous boundary jets.  Hence the V49--V52 positive orientation is
neither conservative nor dissipative in any bare product of a constant
positive metric wave energy \eqref{eq:V53-K-metric-energy} and the
positive V47 bath energy.
\end{theorem}

\begin{proof}
From the explicit V50 matrices,
\begin{equation}
 JC=\frac12I_3.
 \label{eq:V53-JC-half}
\end{equation}
Cancellation for arbitrary \(x_t,t_\Sigma\) would require
\(KC=-J^T\).  If that identity held, then for every nonzero
\(y\in\mathbb R^3\),
\[
 (Cy)^TK(Cy)
 =-y^TC^TJ^Ty
 =-y^T(JC)^Ty
 =-\frac12|y|^2,
\]
contradicting \(K\succ0\).

Thus \(B:=KC+J^T\ne0\).  Choose \(v\) with \(B^Tv\ne0\) and take
\(t_\Sigma=\pm R B^Tv\).  For a fixed boundary jet with \(x_t=v\), the
exchange term is
\(\pm R|B^Tv|^2\), whereas \(x_t^TKA(D)x\) is independent of \(R\).
Both signs occur for large \(R\).  Such traction jets are compatible:
near \(a=0\), a smooth bath profile can prescribe its trace and normal
derivative independently, while the metric graph then prescribes \(p\).
This proves the assertion.
\end{proof}

\begin{corollary}[The sign survives on the radiative quotient]
\label{cor:V53-radiative-sign}
At every oblique direction, the positive-orientation radiative incidence
contributes
\[
 (R_\theta\dot a)^Tt_\Sigma
\]
to the TT metric power.  The positive bath power contributes the same
quantity.  Their exchange sum is therefore
\begin{equation}
 2(R_\theta\dot a)^Tt_\Sigma,
 \label{eq:V53-doubled-radiative-power}
\end{equation}
rather than zero.  At normal incidence \(R_\theta=0\), both terms vanish.
\end{corollary}

\begin{proof}
Theorem~\ref{thm:V53-radiative-port} gives
\(C_{\rm rad}=R_\theta^*\) in the Euclidean TT energy.  Pairing the
positive metric incidence \(+R_\theta^*t_\Sigma\) with \(\dot a\) gives
\((R_\theta\dot a)^Tt_\Sigma\).  The trace constraint is
\(z=R_\theta a\), so
\eqref{eq:V53-bath-power} gives the identical second term.
\end{proof}

\subsection{The opposite orientation and its exact tradeoff}
\label{subsec:V53-negative-orientation}

For the opposite Green orientation the graph would be
\begin{equation}
 p=A(D)x-Ct_\Sigma.
 \label{eq:V53-negative-graph}
\end{equation}
Its product-energy exchange cancels when \(KC=J^T\).  That algebraic
condition does have positive solutions.

\begin{proposition}[A positive exchange-Gram family]
\label{prop:V53-positive-K-family}
For arbitrary \(\gamma,\delta>0\), define, in the coordinate order
\((\tau,v_0,v_1,v_2,\phi)\),
\begin{equation}
 K_{\gamma,\delta}
 :=
 \begin{pmatrix}
 \gamma&0&0&0&3\gamma\\
 0&\delta&0&0&0\\
 0&0&2&0&0\\
 0&0&0&2&0\\
 3\gamma&0&0&0&2+9\gamma
 \end{pmatrix}.
 \label{eq:V53-K-family}
\end{equation}
Then
\begin{equation}
 K_{\gamma,\delta}\succ0,
 \qquad
 K_{\gamma,\delta}C=J^T.
 \label{eq:V53-K-properties}
\end{equation}
Thus the traction exchange in
\eqref{eq:V53-negative-graph} exactly cancels the positive V47 bath
power.
\end{proposition}

\begin{proof}
For \(x=(\tau,v_0,v_1,v_2,\phi)\),
\begin{equation}
 x^TK_{\gamma,\delta}x
 =\gamma(\tau+3\phi)^2+\delta v_0^2
  +2v_1^2+2v_2^2+2\phi^2,
 \label{eq:V53-K-square}
\end{equation}
which is positive for \(x\ne0\).  The first two columns of \(C\) are
\(\tfrac12e_{v_1}\) and \(\tfrac12e_{v_2}\), and the third is
\(-\tfrac32e_\tau+\tfrac12e_\phi\).  Multiplication by
\eqref{eq:V53-K-family} sends them to
\(e_{v_1},e_{v_2},e_\phi\), the three columns of \(J^T\).
For \eqref{eq:V53-negative-graph}, the metric exchange is therefore
\(-x_t^TJ^Tt_\Sigma\), the negative of
\eqref{eq:V53-bath-power}.
\end{proof}

The positive exchange identity is still not a complete generator
symmetrizer.  Put
\begin{equation}
 W:=\tau+3\phi,
 \qquad
 v_\perp:=(v_1,v_2).
 \label{eq:V53-W-vperp}
\end{equation}

\begin{lemma}[Intrinsic graph power residual]
\label{lem:V53-intrinsic-power}
For every \(\gamma,\delta>0\),
\begin{align}
 x_t^TK_{\gamma,\delta}A(D)x
 &=
 \left(\frac{\delta}{6}-2\gamma\right)W_t(v_0)_t
 +2\gamma W_t\operatorname{div}_yv_\perp
 \nonumber\\
 &\quad
 +\frac13(v_\perp)_t\cdot\nabla_yW.
 \label{eq:V53-A-power-residual}
\end{align}
The integral of this expression over the tangential plane takes both
signs on smooth compactly supported instantaneous jets.  Consequently
\(K_{\gamma,\delta}\) solves the traction-exchange sign but does not make
the bare metric boundary graph dissipative.
\end{lemma}

\begin{proof}
The V41 operator is
\[
 A(D)x
 =
 \left(
 -2(v_0)_t+2\operatorname{div}_yv_\perp,\,
 \frac16W_t,\,
 \frac16\partial_1W,\,
 \frac16\partial_2W,\,
 0
 \right)^T.
\]
Multiplication by \eqref{eq:V53-K-family} and contraction with \(x_t\)
gives \eqref{eq:V53-A-power-residual}.  If
\(\delta\ne12\gamma\), take a nonzero compactly supported \(W_t\), set
\((v_0)_t=\pm W_t\), and set \(v_\perp=(v_\perp)_t=0\).  The two
integrals are nonzero and have opposite signs.  If
\(\delta=12\gamma\), take any compactly supported \(v_\perp\) with
nonzero divergence, set \((v_\perp)_t=0\), and set
\(W_t=\pm\operatorname{div}_yv_\perp\).  The integral is then
\[
 \pm2\gamma\int_{\mathbb R^2}
 |\operatorname{div}_yv_\perp|^2\,\dd^2y.
\]
Both signs occur, proving the claim.
\end{proof}

\begin{theorem}[Exact orientation tradeoff for the retained bath class]
\label{thm:V53-orientation-tradeoff}
For the finite positive half-line baths used in V47--V52:
\begin{enumerate}
\item the positive orientation has the uniform stable metric floor of
      Theorem~\ref{thm:V50-uniform-metric-floor}, but fails the positive
      product-energy exchange test of
      Theorem~\ref{thm:V53-positive-energy-obstruction};
\item the negative orientation admits the positive exchange Grams
      \eqref{eq:V53-K-family}, but lies in the complete local covariant
      negative-orientation obstruction of
      Corollary~\ref{cor:V48-covariant-negative-no-go};
\item neither orientation presently supplies both a stable half-space
      inverse and a positive dissipative or conservative time-generator
      energy.
\end{enumerate}
\end{theorem}

\begin{proof}
The first item combines the cited V50 theorem with
Theorem~\ref{thm:V53-positive-energy-obstruction}.  The second combines
Proposition~\ref{prop:V53-positive-K-family} with the V48 result, whose
hypotheses include every finite direct sum of the local positive
square-root half-line baths used here.  The third is their conjunction;
equation~\eqref{eq:V53-A-power-residual} also shows why exchange
cancellation alone would not finish the negative branch.
\end{proof}

\begin{firewall}[Scope of the energy obstruction]
\label{fire:V53-energy-scope}
Theorem~\ref{thm:V53-positive-energy-obstruction} rules out the bare
block product of a constant positive metric fibre Gram and the V47 bath
energy.  It does not rule out a positive energy with metric--bath cross
terms, an additional boundary storage field, a derivative port
coupling, or a nonlocal physical TT realization.  Any such repair must
print its time matrix, prove its positive floor, and preserve the V50
closed-cone inverse.  The TT theorem shows that a repair cannot simply
declare the offending force to be pure gauge at oblique incidence.
\end{firewall}

\subsection{Finite certificates and the V54 upstream target}
\label{subsec:V53-certificates}

A null-shell certificate in the frame
\(k_\mu=(-1,0,0,1)\) stores the matrices
\[
 D_k:\operatorname{Sym}^2\mathbb C^4\to\mathbb C^4,
 \qquad
 \Gamma_k:\mathbb C^4\to\operatorname{Sym}^2\mathbb C^4,
 \qquad
 T_k:Z_k\to\mathbb C^2,
\]
and verifies
\begin{equation}
 \operatorname{rank}D_k=4,\quad
 \operatorname{rank}\Gamma_k=4,\quad
 D_k\Gamma_k=0,\quad
 \operatorname{rank}T_k=2,\quad
 T_k\Gamma_k=0.
 \label{eq:V53-cohomology-card}
\end{equation}
The port card stores
\begin{equation}
 JC-\frac12I_3=0,\qquad
 R_\theta^*R_\theta-\frac{\sin^2\theta}{2}I_2=0,\qquad
 K_{\gamma,\delta}C-J^T=0,
 \label{eq:V53-port-card}
\end{equation}
together with the square decomposition
\eqref{eq:V53-K-square}.  These are exact algebraic residuals, not
frequency sampling claims.

V54 must now project the complete \emph{boundary domain}, rather than
only the port incidence, into physical characteristic variables.  It
should:
\begin{enumerate}
\item pair incoming and outgoing TT amplitudes so that the five
      trace-free induced rows and the reduced \(x\)-graph are imposed
      simultaneously;
\item determine the half-space realization of the Riesz/TT projector and
      test preservation of the harmonic subsidiary domain;
\item introduce the smallest explicit boundary storage or cross-energy
      block capable of reversing the doubled power
      \eqref{eq:V53-doubled-radiative-power};
\item recompute the complete stable determinant after that modification,
      since reversing only the V47 force sign is excluded by V48; and
\item return a concrete negative vector or unstable root if the proposed
      augmented time matrix or stable floor fails.
\end{enumerate}

\subsection{V53 ledger}
\label{subsec:V53-ledger}

\paragraph{New conclusions proved in V53.}
Theorem~\ref{thm:V53-radiative-exact-sequence} computes the exact
two-polarization cohomology of nonzero null harmonic amplitudes by
residual wave diffeomorphisms.  Corollary
\ref{cor:V53-positive-radiative-energy} supplies its bounded TT
representative and positive frozen interior energy.  Proposition
\ref{prop:V53-quotient-distinction} proves that the V39 embedding
quotient is not this radiative quotient.  Theorem
\ref{thm:V53-radiative-port} computes the exact physical projection of
the bath trace and force.  Theorem
\ref{thm:V53-positive-energy-obstruction} and Corollary
\ref{cor:V53-radiative-sign} prove the positive-orientation power defect.
Proposition~\ref{prop:V53-positive-K-family} constructs the opposite
orientation's positive exchange Gram, and Theorem
\ref{thm:V53-orientation-tradeoff} proves the resulting
stability--energy tradeoff.

\paragraph{Imported conclusions not repeated.}
V29 supplies the positive first-order wave symmetrizer.  V34/V35 supply
the propagated harmonic constraint and its rank-four symbol.  V39
supplies the embedding exact sequence.  V41 supplies \(A,C,J\), V47
supplies the two positive bath energies, V48 excludes the complete
negative square-root orientation, and V50 supplies the positive
orientation's uniform direct metric floor.

\paragraph{Honest continuum status.}
The frozen interior radiative quotient now has an explicit positive
symmetrizer, and the selected bath's failure to extend it to a positive
product-energy boundary generator has an exact sign witness.  A
boundary-preserving TT projector, a stable positive augmented boundary
energy, the off-shell ghost domain, nonlinear covariant stress/Bianchi
closure, and the coupled Higgs, fermion, neutrino, and Yang--Mills
sectors remain open.  The full Standard Model--Einstein continuum is
therefore not proved.

\section*{Version V53 append-only record}
\addcontentsline{toc}{section}{Version V53 append-only record}
The complete V52 mathematical source was retained before this
continuation.  Its SHA--256 digest was
\begin{center}\scriptsize\ttfamily
852b7054eadf391296ce2ea001a65494064af28c3ea3401c09612938cdbfcd01.
\end{center}
No mathematical content of V52 was altered.  On the next iteration,
retain this full file, append before its unique active
document-closing command, increment the filename to
\texttt{\_V54.tex}, and remove only the superseded V53 file from this
exact SM note series.  The next scheduled companion audit remains V55.

% END APPENDED VERSION V53 MATERIAL

% BEGIN APPENDED VERSION V54 MATERIAL

\clearpage
\section{V54: boundary TT noninvariance, passive sign dichotomy, and
the minimal mediator's principal-order failure}
\label{sec:V54-boundary-energy}

\subsection{Purpose and the four proposed repairs}
\label{subsec:V54-purpose}

V53 constructed the positive two-polarization TT quotient in the frozen
interior and proved that the V49--V52 positive bath orientation doubles
the radiative exchange power at oblique incidence.  Four apparently
small repairs remained possible:
\begin{enumerate}
\item restrict the whole-space TT projector to the half-space and project
      the existing V52 domain;
\item reverse a trace or force sign while retaining the effective V50
      stable matrix;
\item add a metric--traction cross term to the energy; or
\item place the bath behind a minimal positive boundary storage element.
\end{enumerate}
V54 tests all four.  The first three fail by exact finite or functional
identities.  The fourth produces a legitimate positive-energy
interconnection and lies outside the direct-trace hypothesis of V48, but
its effective impedance is order zero at gravitational glancing.  It
therefore cannot retain the normalized principal floor.

This is an upstream result rather than another determinant fit.  It
identifies what V55 must seek in the scheduled companion audit: a
boundary-adapted physical projector or a passive mediator with a genuine
order-one principal response.

\subsection{The whole-space TT multiplier does not preserve the
essential boundary plane}
\label{subsec:V54-TT-boundary}

Let \(n=e_3\) be the spatial boundary normal and let
\(Q_{\rm TF}\) select the trace-free part of the induced symmetric tensor
on the timelike boundary with coordinates \((t,e_1,e_2)\).  The essential
V52 metric row is
\begin{equation}
 Q_{\rm TF}h=0.
 \label{eq:V54-essential-plane}
\end{equation}
For a real spatial direction
\begin{equation}
 \widehat\xi=(\sin\theta,0,\cos\theta),
 \qquad 0<\theta\le\frac{\pi}{2},
 \label{eq:V54-oblique-direction}
\end{equation}
let \(\Pi_{\rm TT}(\widehat\xi)\) be the whole-space projector
\eqref{eq:V53-TT-projector}.

\begin{theorem}[Exact TT boundary-plane counterexample]
\label{thm:V54-TT-noninvariance}
The kernel of \(Q_{\rm TF}\) is not invariant under
\(\Pi_{\rm TT}(\widehat\xi)\) at any oblique direction.  More precisely,
for the spatial tensor
\begin{equation}
 b:=n\otimes n
 \label{eq:V54-normal-tensor}
\end{equation}
with every time component zero,
\begin{equation}
 Q_{\rm TF}b=0,
 \qquad
 Q_{\rm TF}\Pi_{\rm TT}(\widehat\xi)b\ne0.
 \label{eq:V54-TT-counterexample}
\end{equation}
Thus the V53 whole-space multiplier cannot simply be restricted to the
V52 domain and called a boundary-preserving physical projector.
\end{theorem}

\begin{proof}
Use the orthonormal transverse frame from V53,
\[
 e=(0,1,0),
 \qquad
 m=(\cos\theta,0,-\sin\theta),
 \qquad
 E_+=\frac{e\otimes e-m\otimes m}{\sqrt2}.
\]
The tensor \(b\) has no induced component, proving
\(Q_{\rm TF}b=0\).  Its only TT coefficient is
\(\langle E_+,b\rangle=-\sin^2\theta/\sqrt2\), so
\begin{equation}
 \Pi_{\rm TT}b=-\frac{\sin^2\theta}{\sqrt2}E_+.
 \label{eq:V54-projected-normal}
\end{equation}
In particular,
\[
 (\Pi_{\rm TT}b)_{22}=-\frac{\sin^2\theta}{2}\ne0,
 \qquad
 (\Pi_{\rm TT}b)_{00}=0.
\]
If its induced trace-free part vanished, its induced block would be a
multiple of \(\operatorname{diag}(-1,1,1)\).  The zero \(00\) entry would
force that multiple, and hence the \(22\) entry, to vanish.  This is a
contradiction, proving \eqref{eq:V54-TT-counterexample}.
\end{proof}

\begin{firewall}[What the counterexample decides]
\label{fire:V54-TT-scope}
The tensor \eqref{eq:V54-normal-tensor} is a configuration-space test and
need not itself be a harmonic null solution.  Theorem
\ref{thm:V54-TT-noninvariance} decides the operator-domain question:
the standard full-space projector does not preserve even the essential
configuration row.  It does not exclude a different York/Hodge projector
defined by an elliptic half-space problem with boundary conditions.  Such
a projector must print its normal realization, prove a uniform Sobolev
bound, and commute with the harmonic subsidiary evolution on its domain.
\end{firewall}

\subsection{No static sign transformer separates passivity from the
stable orientation}
\label{subsec:V54-sign-transformer}

Use the positive metric Gram of
Proposition~\ref{prop:V53-positive-K-family}, so
\(KC=J^T\), and put \(z=Jx\).  Generalize the direct trace and metric
incidence by two nonzero real gains:
\begin{equation}
 q_j(0)=\sigma z,
 \qquad
 p=A(D)x+\varepsilon Ct_\Sigma,
 \qquad
 t_\Sigma=t_1+t_2.
 \label{eq:V54-static-transformer}
\end{equation}
The retarded bath law is then
\begin{equation}
 t_\Sigma=\sigma Z_*z.
 \label{eq:V54-transformed-traction}
\end{equation}

\begin{theorem}[Static transformer sign dichotomy]
\label{thm:V54-static-sign-dichotomy}
For \eqref{eq:V54-static-transformer}:
\begin{align}
 {\cal P}_{\rm exchange}
 &=(\varepsilon+\sigma)z_t^Tt_\Sigma,
 \label{eq:V54-transformer-power}\\
 L_{\varepsilon,\sigma}
 &=L_0-\varepsilon\sigma\,CZ_*J.
 \label{eq:V54-transformer-stable}
\end{align}
Consequently exact positive-energy exchange cancellation requires
\(\sigma=-\varepsilon\), and then
\begin{equation}
 \varepsilon\sigma=-\varepsilon^2<0.
 \label{eq:V54-passive-negative-sign}
\end{equation}
The V50 principal repair has the opposite sign
\(\varepsilon\sigma>0\).  No real static trace/force transformer can
retain that repair while making the metric--bath exchange passive.
\end{theorem}

\begin{proof}
The metric exchange power is
\[
 x_t^TK(\varepsilon Ct_\Sigma)
 =\varepsilon z_t^Tt_\Sigma.
\]
The bath boundary velocity is
\(\partial_tq_j(0)=\sigma z_t\), so the summed V47 power is
\(\sigma z_t^Tt_\Sigma\).  This proves
\eqref{eq:V54-transformer-power}.  For a stable metric amplitude,
\(p=\lambda x\).  Substitute
\eqref{eq:V54-transformed-traction} into the second row of
\eqref{eq:V54-static-transformer} and use
\(L_0=\lambda I-A(d)\) to obtain
\eqref{eq:V54-transformer-stable}.  Power cancellation and
\eqref{eq:V54-passive-negative-sign} are immediate.
\end{proof}

For the sign choices of the retained notes,
\((\varepsilon,\sigma)=(1,1)\) gives the stable V50 matrix and doubled
power, while either \((-1,1)\) or \((1,-1)\) gives passive exchange and
the V48 negative-orientation symbol.  The dichotomy is therefore
independent of whether one describes the sign change as a force reversal
or an ideal trace transformer.

\subsection{Bath traction is not an energy-space trace}
\label{subsec:V54-traction-trace}

The most direct cross-energy response to
\eqref{eq:V53-doubled-radiative-power} would contain the traction
\(t=-M\partial_aq(0)\).  The V47 configuration energy controls one normal
derivative in \(L^2\), but it does not control the boundary value of that
derivative.

\begin{theorem}[Unbounded traction trace on the bath energy space]
\label{thm:V54-traction-unbounded}
Let \(M\succ0\) and let the bath energy space have configuration norm
\begin{equation}
 \|q\|_{{\cal E}_v}^2
 :=\int_{\mathbb R^2}\int_0^\infty
 \left(
 q^TMq+\partial_aq^TM\partial_aq
 +v^2\partial_Aq^TM\partial_Aq
 \right)\dd a\,\dd^2y.
 \label{eq:V54-bath-energy-norm}
\end{equation}
The map
\begin{equation}
 q\longmapsto-M\partial_aq(0)
 \label{eq:V54-traction-map}
\end{equation}
has no bounded extension from this energy space to
\(L^2(\mathbb R^2;\mathbb R^3)\), even on fields with zero boundary
trace.  Therefore a nonzero quadratic cross term
\(\int z^Tt\,\dd^2y\) is not a bounded energy correction on the natural
V47 Hilbert state.
\end{theorem}

\begin{proof}
Choose \(e\in\mathbb R^3\setminus\{0\}\),
\(\varphi\in C_c^\infty(\mathbb R^2)\setminus\{0\}\), and
\(\psi\in C_c^\infty([0,\infty))\) with
\(\psi(0)=0\), \(\psi'(0)=1\).  Put
\begin{equation}
 q_N(a,y)
 :=N^{-1/2}\psi(Na)\varphi(y)e.
 \label{eq:V54-boundary-layer}
\end{equation}
Then \(q_N(0)=0\), while
\begin{equation}
 -M\partial_aq_N(0,y)
 =-N^{1/2}\varphi(y)Me.
 \label{eq:V54-layer-traction}
\end{equation}
The normal-gradient part of
\eqref{eq:V54-bath-energy-norm} is exactly
\[
 (e^TMe)\|\varphi\|_{L^2}^2
 \int_0^\infty|\psi'(u)|^2\,\dd u,
\]
independent of \(N\).  The undifferentiated and tangential-gradient parts
are \(O(N^{-2})\).  Hence \(\|q_N\|_{{\cal E}_v}\) is bounded while the
traction norm grows like \(N^{1/2}\).  This proves unboundedness.

To test a cross term at a prescribed nonzero trace \(z\), add
\eqref{eq:V54-boundary-layer} to any fixed smooth extension of \(z\).
The energy remains bounded and the added traction is unbounded in a
chosen direction.  Thus the bilinear cross term is not continuous.
\end{proof}

\begin{firewall}[Domain strengthening is a new theorem obligation]
\label{fire:V54-traction-scope}
On the operator domain, an equation may define traction weakly in a
negative trace space; on the V51 weighted \(H^2\) resolvent domain it is
an ordinary boundary value.  Neither fact makes it a bounded coordinate
of the conserved \(H^1\oplus L^2\) energy state.  A cross-energy repair
using \(t\) must enlarge the phase space or prove a stronger graph norm
whose generator is closed and whose positive energy is equivalent to
that norm.
\end{firewall}

\subsection{A minimal passive three-channel mediator}
\label{subsec:V54-mediator}

The V39 channel theorem shows that a regularly eliminable glancing repair
needs at least three boundary amplitudes.  Introduce exactly three,
\(y(t,\mathbf y)\in\mathbb R^3\), and replace the common bath trace by
\begin{equation}
 q_1(0)=q_2(0)=y.
 \label{eq:V54-mediated-trace}
\end{equation}
Let \(m>0\), \(\chi>0\), set \(z=Jx\), and define
\begin{equation}
 r:=\chi(z-y).
 \label{eq:V54-spring-force}
\end{equation}
Use the negative-feedback metric row and mediator equation
\begin{align}
 p&=A(D)x-Cr,
 \label{eq:V54-mediated-metric}\\
 m y_{tt}+\chi(y-z)+t_\Sigma&=0.
 \label{eq:V54-mediated-oscillator}
\end{align}
The mediator storage is
\begin{equation}
 E_{\rm med}
 :=\frac12\int_{\mathbb R^2}
 \left(m|y_t|^2+\chi|y-z|^2\right)\dd^2y.
 \label{eq:V54-mediator-energy}
\end{equation}

\begin{proposition}[Exact positive exchange balance]
\label{prop:V54-mediated-energy}
For any \(K_{\gamma,\delta}\) in
\eqref{eq:V53-K-family}, the traction exchange among the metric, mediator,
and two baths cancels exactly:
\begin{equation}
 \left.\frac{\dd}{\dd t}
 \left(
 E_{{\rm met},K_{\gamma,\delta}}
 +E_{\rm med}+E_{{\rm bath},(2)}
 \right)\right|_{\rm exchange}
 =0.
 \label{eq:V54-mediated-power-balance}
\end{equation}
The only remaining metric boundary power is the intrinsic graph term
\(x_t^TK_{\gamma,\delta}A(D)x\) computed in
Lemma~\ref{lem:V53-intrinsic-power}.
\end{proposition}

\begin{proof}
Since \(K_{\gamma,\delta}C=J^T\), the force term in
\eqref{eq:V54-mediated-metric} contributes
\(-z_t^Tr\) to the metric power.  Differentiate
\eqref{eq:V54-mediator-energy} and use
\eqref{eq:V54-mediated-oscillator}:
\begin{align*}
 \frac{\dd}{\dd t}E_{\rm med}
 &=
 \int\left[
  my_t^Ty_{tt}+\chi(y-z)^T(y_t-z_t)
 \right]\dd^2y\\
 &=
 \int\left[-y_t^Tt_\Sigma+z_t^Tr\right]\dd^2y.
\end{align*}
The bath power is \(y_t^Tt_\Sigma\) by
\eqref{eq:V53-bath-power} with trace \(y\).  The three exchange terms
sum to zero.  The \(A(D)\) term is unaffected.
\end{proof}

This mediator is positive and uses the minimum three configuration
channels allowed by Theorem~\ref{thm:V39-three-channel-lower-bound} when
the mediator is regularly eliminable.  It is not yet claimed to arise
from the original V36 DeWitt boundary action; it is a candidate positive
time-energy interconnection.

\subsection{Retarded mediator impedance and the principal-order no-go}
\label{subsec:V54-mediator-symbol}

For a Laplace--Fourier mode, outgoing elimination of the two baths gives
\(t_\Sigma=Z_*y\).  Equation
\eqref{eq:V54-mediated-oscillator} becomes
\begin{equation}
 D_{\rm med}(s,k)y=\chi z,
 \qquad
 D_{\rm med}:=ms^2I_3+\chi I_3+Z_*(s,k).
 \label{eq:V54-mediator-denominator}
\end{equation}
Where this matrix is invertible,
\begin{align}
 y&=\chi D_{\rm med}^{-1}z,
 \label{eq:V54-y-response}\\
 r&=Z_{\rm med}(s,k)z,
 \qquad
 Z_{\rm med}
 :=\chi(ms^2I_3+Z_*)D_{\rm med}^{-1}.
 \label{eq:V54-mediated-impedance}
\end{align}

\begin{lemma}[Open-half-plane mediator invertibility]
\label{lem:V54-mediator-invertible}
For \(\operatorname{Re}s>0\), the homogeneous mediator--bath subsystem
with \(z=0\) has only the zero finite-energy outgoing solution.  Hence
\(D_{\rm med}(s,k)\) is invertible there.
\end{lemma}

\begin{proof}
With \(z=0\), the sum of
\eqref{eq:V54-mediator-energy} and the two positive V47 bath energies is
conserved: Proposition~\ref{prop:V54-mediated-energy} applies with the
metric port clamped.  A nonzero mode proportional to \(e^{st}\) would
make this positive energy proportional to
\(e^{2\operatorname{Re}s\,t}\), contradicting conservation when
\(\operatorname{Re}s>0\).  Zero energy forces \(y=y_t=0\) and both bath
fields to vanish.  The outgoing bath elimination is exact in the open
half-plane, so a kernel vector of \(D_{\rm med}\) would give precisely
such a mode.  No kernel exists.
\end{proof}

The stable reduced metric matrix for
\eqref{eq:V54-mediated-metric} is
\begin{equation}
 L_{\rm med}(s,k)
 :=L_0(s,k)+CZ_{\rm med}(s,k)J.
 \label{eq:V54-mediated-metric-symbol}
\end{equation}

\begin{theorem}[Minimal mediator loses the glancing principal floor]
\label{thm:V54-mediator-principal-no-go}
Along either gravitational glancing ray
\begin{equation}
 s=\pm i\rho r,\qquad k=\rho r k_0,\qquad |k_0|=1,
 \qquad \lambda=0,\qquad \rho\longrightarrow\infty,
 \label{eq:V54-glancing-ray}
\end{equation}
with fixed \(r>0\),
\begin{equation}
 Z_{\rm med}(s,k)\longrightarrow\chi I_3,
 \qquad
 \rho^{-1}Z_{\rm med}(s,k)\longrightarrow0.
 \label{eq:V54-mediator-order-zero}
\end{equation}
Consequently
\begin{equation}
 \rho^{-1}L_{\rm med}(s,k)
 \longrightarrow L_0(\pm ir,rk_0),
 \label{eq:V54-mediated-glancing-limit}
\end{equation}
The limiting matrix has rank two.  Therefore
\begin{equation}
 \inf_{\rho\ge1}
 s_{\min}\!\left(\rho^{-1}L_{\rm med}
   (\pm i\rho r,\rho r k_0)\right)=0,
 \label{eq:V54-no-mediated-floor}
\end{equation}
and this positive minimal mediator does not give a uniform principal
Kreiss--Lopatinski floor.
\end{theorem}

\begin{proof}
At gravitational glancing the two bath roots have boundary values
\begin{equation}
 \kappa_j
 =\pm i\rho r\sqrt{1-v_j^2},
 \label{eq:V54-bath-glancing-roots}
\end{equation}
with the sign fixed by limiting absorption.  Hence
\[
 Z_*(s,k)=O(\rho),
 \qquad
 ms^2I_3=-m\rho^2r^2I_3.
\]
Factor the common \(-m\rho^2r^2\) from numerator and denominator in
\eqref{eq:V54-mediated-impedance}.  Their quotient tends to \(I_3\),
which proves \eqref{eq:V54-mediator-order-zero}.

The graph matrix \(L_0\) is homogeneous of degree one.  Divide
\eqref{eq:V54-mediated-metric-symbol} by \(\rho\) and use
\eqref{eq:V54-mediator-order-zero} to obtain
\eqref{eq:V54-mediated-glancing-limit}.  The V37 glancing computation,
equivalently the reduced matrix \(M_0\) in
\eqref{eq:V39-M0}, gives rank two.  Singular values are continuous, so
the three singular values which vanish at the limit force
\eqref{eq:V54-no-mediated-floor}.
\end{proof}

The same order-zero conclusion holds formally at \(m=0\):
\(Z_{\rm med}=\chi Z_*(\chi I+Z_*)^{-1}\to\chi I\) on any glancing
component where the leading bath matrix is invertible.  Adding positive
finite inertia therefore does not cause the loss; the constant spring is
subprincipal relative to the order-one defect.

\begin{firewall}[What the mediator proves and what it does not]
\label{fire:V54-mediator-scope}
Proposition~\ref{prop:V54-mediated-energy} proves a genuine positive
exchange identity and Lemma~\ref{lem:V54-mediator-invertible} excludes an
open growing mode in the clamped mediator--bath subsystem.  Theorem
\ref{thm:V54-mediator-principal-no-go} proves that this subsystem cannot
repair the gravitational glancing defect uniformly at high frequency.
It does not exclude an order-one passive mediator, a boundary continuum,
or a boundary-adapted physical quotient on which the rank-two limit is
the correct number of radiative rows.  A time-generator theorem for the
mediator would also have to define the Wentzell-type domain on which
\(t_\Sigma\) is meaningful and prove that the resulting operator is
closed and maximal; the power identity alone does not do so.
\end{firewall}

\subsection{Finite certificates and the V55 companion audit}
\label{subsec:V54-certificates}

The boundary-projector card takes \(b=n\otimes n\) and verifies
\begin{equation}
 Q_{\rm TF}b=0,
 \qquad
 (\Pi_{\rm TT}b)_{22}=-\frac12\sin^2\theta.
 \label{eq:V54-TT-card}
\end{equation}
The sign card stores
\begin{equation}
 {\cal P}_{\rm exchange}
 =(\varepsilon+\sigma)z_t^Tt_\Sigma,
 \qquad
 {\rm sign}_{\rm eff}=\varepsilon\sigma.
 \label{eq:V54-sign-card}
\end{equation}
The traction card stores the bounded normal-gradient energy and the
\(N^{1/2}\) traction growth of
\eqref{eq:V54-boundary-layer}.

For the mediator, a finite certificate verifies
\begin{equation}
 r_{\rm power}
 :=
 \left|
 -z_t^Tr+\frac{\dd}{\dd t}E_{\rm med}
 +y_t^Tt_\Sigma
 \right|=0,
 \label{eq:V54-power-card}
\end{equation}
the rank-two glancing limit of \(L_0\), and the decay of
\(s_{\min}(\rho^{-1}L_{\rm med})\).  For
\(m=0.7\), \(\chi=1.3\), \(r=1\), \(k_0=(1,0)\), and the V49 bath data, the values at
\(\rho=10,100,1000,10000\) are
\begin{equation}
 2.066371\times10^{-2},\quad
 2.055622\times10^{-3},\quad
 2.055482\times10^{-4},\quad
 2.055480\times10^{-5}.
 \label{eq:V54-numerical-smin}
\end{equation}
These numbers illustrate the proved limit; they are not used to prove it.

V55 is the next scheduled five-version audit.  Before selecting another
mediator, it must read the newest completed NS and Yang--Mills notes in
the shared \texttt{latex\_notes} folder and test whether they contain:
\begin{enumerate}
\item a boundary Hodge projector with a proved domain and commutator
      estimate;
\item a passive order-one traction/velocity impedance with local or
      conservative-continuum ancestry;
\item a boundary storage norm that controls the conormal trace without
      replacing the energy topology by the V51 resolvent graph norm; or
\item a physical characteristic reduction in which the rank-two
      glancing limit is already the complete radiative boundary count.
\end{enumerate}
If none is available, V55 should derive the principal symbol required of
an order-one passive mediator before proposing further fields.

\subsection{V54 ledger}
\label{subsec:V54-ledger}

\paragraph{New conclusions proved in V54.}
Theorem~\ref{thm:V54-TT-noninvariance} gives an exact counterexample to
restricting the whole-space TT projector to the V52 essential boundary
plane.  Theorem~\ref{thm:V54-static-sign-dichotomy} proves that no static
trace/force sign transformer separates passive exchange from the
negative effective bath sign.  Theorem
\ref{thm:V54-traction-unbounded} proves that bath traction is not a
bounded coordinate of the natural positive energy state.
Proposition~\ref{prop:V54-mediated-energy} constructs a minimal
three-channel positive exchange mediator, while Theorem
\ref{thm:V54-mediator-principal-no-go} proves that its response is
subprincipal and loses the glancing floor.

\paragraph{Imported conclusions not repeated.}
V39 supplies the three-channel lower bound and the rank-two glancing
matrix.  V47 supplies the positive bath energy and traction power.  V48
supplies the complete direct-trace negative-orientation obstruction.  V50
supplies the stable positive-orientation floor.  V53 supplies the
positive reduced metric Grams, the TT radiative incidence, and the
doubled-power identity.

\paragraph{Honest continuum status.}
The elementary projector, transformer, cross-energy, and finite-storage
repairs have now been decided.  A boundary-adapted physical projector or
an order-one passive mediator remains necessary before a positive stable
Einstein--bath generator can be claimed.  The off-shell ghost domain,
nonlinear covariant stress/Bianchi closure, and the coupled Higgs,
fermion, neutrino, and Yang--Mills sectors also remain open.  The full
Standard Model--Einstein continuum is therefore not proved.

\section*{Version V54 append-only record}
\addcontentsline{toc}{section}{Version V54 append-only record}
The complete V53 mathematical source was retained before this
continuation.  Its SHA--256 digest was
\begin{center}\scriptsize\ttfamily
7f44056b971cbc6c13734526128146ed9ed5ca460b7dbfdee680a0aa037571c8.
\end{center}
No mathematical content of V53 was altered.  On the next iteration,
retain this full file, append before its unique active
document-closing command, increment the filename to
\texttt{\_V55.tex}, and remove only the superseded V54 file from this
exact SM note series.  The next iteration must perform the scheduled
companion audit before selecting its new construction.

% END APPENDED VERSION V54 MATERIAL

% BEGIN APPENDED VERSION V55 MATERIAL

\clearpage
\section{V55: companion audit, exact glancing port determinant, and
the complete Ohmic principal-mediator obstruction}
\label{sec:V55-audit-principal-port}

\subsection{Scheduled NS/Yang--Mills companion audit}
\label{subsec:V55-companion-audit}

V55 begins with the five-version audit required by the user and by the
V54 ledger.  The shared \texttt{latex\_notes} folder was inspected
without modifying any companion file.

The newest completed Navier--Stokes note was
\(\texttt{Renewal\_NS\_Stability\_Research\_Notes\_V29.tex}\),
with SHA--256 digest
\begin{center}\scriptsize\ttfamily
77e7087398ae9ca94469a91464cee3e0c1780d2f7513cbc4620d9881dd795a05.
\end{center}
Its V29 appendix identifies the Huber selector as the gradient of an
exact soft spectral-tail meet, proves a cutoff-radius-weighted occupation
bound paid by ordinary Leray dissipation, and derives parabolic lifetime
and persistence-or-recent-influx alternatives.  Applying these results
to the V28 low-total residual pays its long-lived terminal self part, but
the radius-dependent recent-entry intervals, initial tail, and frozen
viscous difference still require a simultaneous signed assembly.

The newest completed Yang--Mills note was
\(\texttt{Renewal\_Yang\_Mills\_Mass\_Gap\_Research\_Notes\_V32.tex}\),
with SHA--256 digest
\begin{center}\scriptsize\ttfamily
f8e4d8a509c2a247cd6deef4f48009b0c4b3f3ba67f2e5eabf6429b31a409a42.
\end{center}
Its V32 appendix proves a quadratic Wilson tail debit, an exact ternary
replicated-coordinate inequality, and ternary marked-tree control on
every fixed thermal window.  The constant is not uniform as the thermal
window expands, so the high-to-low Racah branch remains open.

\begin{assessment}[Transfer decision from the scheduled audit]
\label{ass:V55-companion-transfer}
Neither audited appendix contains a boundary Hodge/TT projector, a
traction trace estimate on an \(H^1\) energy space, a passive order-one
boundary impedance, or a physical incoming/outgoing gravitational
reduction.  No theorem from either file can populate a missing V54
operator row.

There is a relevant proof discipline.  NS V29 pays long persistence but
does not thereby assemble radius-dependent recent-entry intervals; YM
V32 controls every fixed thermal window but not thermal escape.  In the
present programme, the analogous forbidden shortcut would be to promote
the order-zero V54 mediator, or a pointwise glancing repair, to a uniform
principal inverse.  The normalized order-zero correction vanishes, and
glancing invertibility alone does not exclude a stable-cone zero.  The
V55 direction therefore stays with the fallback prescribed in V54:
derive the required principal port symbol and test the simplest passive
order-one realization on the full stable cone.
\end{assessment}

\begin{firewall}[Audit scope]
\label{fire:V55-audit-scope}
The NS and Yang--Mills conclusions above are used only to decide whether
a proved companion mechanism transfers.  Their tail, Wilson-link, and
cumulant estimates are not reinterpreted as gravitational boundary
estimates.  Both companion programmes retain their stated open
residuals, and no claim about Navier--Stokes regularity or the
Yang--Mills mass gap is made here.
\end{firewall}

\subsection{Exact three-port determinant at gravitational glancing}
\label{subsec:V55-glancing-determinant}

Rotate the tangential spatial covector to \(k=(r,0)\), \(r>0\), and
consider either gravitational glancing sheet
\begin{equation}
 s=\pm ir,\qquad \lambda=0.
 \label{eq:V55-glancing-point}
\end{equation}
In the reduced coordinate order
\((\tau,v_0,v_1,v_2,\phi)\), the plus-sheet V41 matrix is
\begin{equation}
 L_{\rm g}:=L_0(ir,(r,0))
 =
 \begin{pmatrix}
 0&2ir&-2ir&0&0\\
 -ir/6&0&0&0&-ir/2\\
 -ir/6&0&0&0&-ir/2\\
 0&0&0&0&0\\
 0&0&0&0&0
 \end{pmatrix}.
 \label{eq:V55-Lg}
\end{equation}
It has rank two.  A right-kernel frame and left-kernel frame are
\begin{align}
 R_{\rm g}
 &:=
 \begin{pmatrix}
 0&0&-3\\
 1&0&0\\
 1&0&0\\
 0&1&0\\
 0&0&1
 \end{pmatrix},
 \label{eq:V55-right-frame}\\
 \Lambda_{\rm g}
 &:=
 \begin{pmatrix}
 0&1&-1&0&0\\
 0&0&0&1&0\\
 1&0&0&0&1
 \end{pmatrix}.
 \label{eq:V55-left-frame}
\end{align}
The V41 port maps obey
\begin{equation}
 JR_{\rm g}=I_3,
 \qquad
 \Lambda_{\rm g}C
 =\operatorname{diag}\left(-\frac12,\frac12,-1\right).
 \label{eq:V55-defect-incidence}
\end{equation}

\begin{theorem}[Exact glancing determinant for an arbitrary port matrix]
\label{thm:V55-glancing-determinant}
For every complex \(3\times3\) matrix \(Z\),
\begin{align}
 \det(L_{\rm g}+CZJ)
 &=-\frac{r^2}{24}\det Z,
 \label{eq:V55-glancing-plus-det}\\
 \det(L_{\rm g}-CZJ)
 &= \frac{r^2}{24}\det Z.
 \label{eq:V55-glancing-minus-det}
\end{align}
The same formula holds on the opposite temporal sheet and after
tangential rotation.  Thus an order-one three-port correction repairs
the five-row metric matrix at gravitational glancing if and only if its
principal port matrix is invertible there.
\end{theorem}

\begin{proof}
Multiplication of \eqref{eq:V55-Lg} by
\eqref{eq:V55-right-frame} gives zero, and multiplication on the left by
\eqref{eq:V55-left-frame} also gives zero.  Direct multiplication by the
explicit V50 matrices \(C,J\) gives
\eqref{eq:V55-defect-incidence}.  Hence the correction maps the complete
three-dimensional right defect to the complete left defect precisely
through the nonsingular diagonal matrix in
\eqref{eq:V55-defect-incidence} followed by \(Z\).

For the normalization constant, insert
\[
 C=
 \begin{pmatrix}
 0&0&-3/2\\0&0&0\\1/2&0&0\\0&1/2&0\\0&0&1/2
 \end{pmatrix},
 \qquad
 J=
 \begin{pmatrix}
 0&0&1&0&0\\0&0&0&1&0\\0&0&0&0&1
 \end{pmatrix}
\]
into \(L_{\rm g}\pm CZJ\) and expand the \(5\times5\) determinant.
Only terms containing two entries of \(L_{\rm g}\) and three entries of
\(Z\) survive.  Their signed sum is respectively
\(-r^2\det Z/24\) and \(r^2\det Z/24\).
Repeating the same expansion after replacing \(ir\) by \(-ir\) gives the
other temporal sheet with the same factor; this step does not impose a
reality condition on the arbitrary complex matrix \(Z\).  Orthogonal
tangential rotation changes neither determinant.  This proves all
claims.
\end{proof}

\begin{corollary}[Necessary principal mediator card]
\label{cor:V55-principal-card}
Let a causal mediator have a continuous degree-one leading impedance
\begin{equation}
 \rho^{-1}Z_{\rm med}(\rho s,\rho k)
 \longrightarrow Z^{(1)}(s,k)
 \label{eq:V55-principal-limit}
\end{equation}
uniformly on the normalized closed stable cone.  For the passive
negative-feedback metric graph
\begin{equation}
 L^{(1)}_{\rm tot}=L_0+CZ^{(1)}J
 \label{eq:V55-passive-principal-matrix}
\end{equation}
to have a uniform inverse, it is necessary that
\begin{equation}
 \inf_{(s,k)\in{\cal G}}
 s_{\min}\bigl(Z^{(1)}(s,k)\bigr)>0.
 \label{eq:V55-glancing-port-floor}
\end{equation}
It is also necessary that
\(\det(L_0+CZ^{(1)}J)\) have no zero anywhere else on the closed stable
hemisphere.  Glancing invertibility alone is not sufficient.
\end{corollary}

\begin{proof}
Homogeneity of \(L_0\) and
\eqref{eq:V55-principal-limit} give the normalized limit
\eqref{eq:V55-passive-principal-matrix}.  Theorem
\ref{thm:V55-glancing-determinant} makes this matrix singular whenever
\(Z^{(1)}\) is singular at glancing.  Continuity and compactness upgrade
pointwise invertibility there to the uniform lower bound
\eqref{eq:V55-glancing-port-floor}.  A zero away from glancing directly
violates the required complete stable-cone inverse, proving the final
statement.
\end{proof}

For a passive displacement-to-traction law
\(t=Z_{\rm med}(s,k)z\), the time-domain supply is \(z_t^Tt\).  Its
frequency card must therefore include
\begin{equation}
 \operatorname{Re}\!
 \left(\overline s\,y^*Z_{\rm med}(s,k)y\right)\ge0
 \qquad
 (\operatorname{Re}s>0)
 \label{eq:V55-passivity-card}
\end{equation}
for every port vector \(y\), in addition to causality, the reality
condition, order-one growth, and
\eqref{eq:V55-glancing-port-floor}.

\subsection{The zero-speed conservative bath is the Ohmic endpoint}
\label{subsec:V55-zero-speed}

The most direct way to meet
\eqref{eq:V55-glancing-port-floor} is to set the tangential bath speed to
zero.  Let \(M=M^T\succeq0\) and use
\begin{equation}
 S_{0,M}[q]
 :=\frac12\int_{\partial M}\int_0^\infty
 \left(q_t^TMq_t-\partial_aq^TM\partial_aq\right)
 \dd a\,\dd^3y,
 \qquad
 q(0)=z.
 \label{eq:V55-zero-speed-action}
\end{equation}
The outgoing root in the open half-plane is
\(\sqrt{s^2}=s\), and therefore
\begin{equation}
 t=-M\partial_aq(0)=sMz.
 \label{eq:V55-Ohmic-impedance}
\end{equation}

\begin{proposition}[Positive conservative ancestry and principal rank]
\label{prop:V55-Ohmic-ancestry}
The action \eqref{eq:V55-zero-speed-action} has nonnegative energy
\begin{equation}
 E_{0,M}
 =\frac12\int\!\!\int
 \left(q_t^TMq_t+\partial_aq^TM\partial_aq\right)
 \dd a\,\dd^2y
 \label{eq:V55-zero-speed-energy}
\end{equation}
and boundary power \(z_t^Tt\).  Its retarded impedance
\(Z_0=sM\) is homogeneous of order one and obeys
\begin{equation}
 \operatorname{Re}\left(\overline s\,y^*Z_0y\right)
 =|s|^2y^*My\ge0.
 \label{eq:V55-Ohmic-passivity}
\end{equation}
At gravitational glancing it repairs the metric rank exactly when
\(M\) is positive definite.
\end{proposition}

\begin{proof}
The energy and boundary-power calculation is the V44 integration by
parts with its tangential-gradient term omitted.  The open-half-plane
outgoing solution is \(q(a)=e^{-sa}z\), which proves
\eqref{eq:V55-Ohmic-impedance} and
\eqref{eq:V55-Ohmic-passivity}.  At glancing \(s=\pm ir\ne0\), so
\(sM\) is invertible exactly when \(M\) is.  Apply
Theorem~\ref{thm:V55-glancing-determinant}.
\end{proof}

Passing glancing is not the full determinant test.  The remainder of V55
proves that every positive \(M\) fails elsewhere.

\subsection{The endpoint ratio map}
\label{subsec:V55-ratio-map}

Put
\begin{equation}
 \theta:=\frac{s}{2\lambda},
 \qquad
 \lambda^2=s^2+r^2.
 \label{eq:V55-theta}
\end{equation}

\begin{lemma}[Stable realization of the zero-speed ratio]
\label{lem:V55-endpoint-ratio}
Every nonzero \(\theta\) with
\(\operatorname{Re}\theta\ge0\) is realized by an open or
closed-boundary stable frequency with \(r\ge0\).  If
\(\operatorname{Re}\theta>0\) and
\(\operatorname{Im}\theta\ne0\), the realizing frequency has
\(\operatorname{Re}s>0\).
\end{lemma}

\begin{proof}
Write \(w=2\theta\).  For \(r>0\), inversion of
\eqref{eq:V55-theta} gives
\begin{equation}
 \frac{s^2}{r^2}=\frac{w^2}{1-w^2},
 \qquad
 \frac{\lambda^2}{r^2}=\frac1{1-w^2}.
 \label{eq:V55-ratio-inversion}
\end{equation}
If \(w\) lies in the first open quadrant, then \(w^2\) lies in the upper
half-plane.  The real Moebius map \(W\mapsto W/(1-W)\) has positive
determinant and maps the upper half-plane to itself.  Its principal
square root gives \(\operatorname{Re}s>0\), and the compatible square
root for \(\lambda\) has positive real part.  Analytic continuation from
\(0<w<1\) fixes \(s/\lambda=w\).  Complex conjugation treats the fourth
quadrant.

For real \(w\), the intervals \(0<w<1\) and \(w>1\) give respectively
an open positive real frequency and a purely imaginary closed-boundary
frequency.  The value \(w=1\) is realized at \(r=0\), where
\(s=\lambda>0\).  For nonzero purely imaginary \(w\),
\eqref{eq:V55-ratio-inversion} gives purely imaginary \(s\) and positive
real \(\lambda\).  These cases exhaust
\(\operatorname{Re}\theta\ge0\).
\end{proof}

\subsection{Exact Ohmic quintic}
\label{subsec:V55-Ohmic-quintic}

In the port basis \((v_L,v_T,\phi)\), write
\begin{equation}
 M=
 \begin{pmatrix}
 A&u^T\\
 u&N
 \end{pmatrix}\succ0
 \label{eq:V55-M-block}
\end{equation}
and retain the V46 invariants
\begin{equation}
 B:=\operatorname{tr}N,\qquad
 C_*:=A\operatorname{tr}N-|u|^2,\qquad
 D:=\det N,\qquad
 E:=\det M.
 \label{eq:V55-M-invariants}
\end{equation}
They are strictly positive.  Define
\begin{align}
 p_M^{(0)}(\theta)
 :=\;&E\theta^5+C_*\theta^4
 +\left(A+\frac34E\right)\theta^3
 +\left(\frac34C_*+D\right)\theta^2
 \nonumber\\
 &+\left(\frac34A+B\right)\theta+1.
 \label{eq:V55-Ohmic-quintic}
\end{align}

\begin{theorem}[Exact zero-speed determinant polynomial]
\label{thm:V55-Ohmic-quintic}
Where \(\lambda\ne0\), the passive negative-feedback determinant
satisfies
\begin{equation}
 \frac{\det(L_0+CsMJ)}{\det L_0}
 =p_M^{(0)}\!\left(\frac{s}{2\lambda}\right).
 \label{eq:V55-Ohmic-determinant}
\end{equation}
No squared branch or extraneous root is introduced.
\end{theorem}

\begin{proof}
The V45 response matrix is
\[
 R=\operatorname{diag}\left(
 \frac{4\lambda^2-r^2}{8\lambda^3},
 \frac1{2\lambda},
 \frac1{2\lambda}
 \right).
\]
The matrix determinant lemma gives the negative-feedback ratio
\[
 F_-=\det(I_3+RsM).
\]
This is exactly the V46 common-matrix algebra with
\(\kappa_v=s\), hence with
\[
 \delta=1,\qquad
 \alpha=\delta-\frac14=\frac34,\qquad
 \theta=\frac{s}{2\lambda}.
\]
Substitution into the unsquared V46 expansion
\eqref{eq:V46-p} gives
\eqref{eq:V55-Ohmic-quintic} and
\eqref{eq:V55-Ohmic-determinant} directly.  The calculation uses the
analytic open-half-plane branch \(\sqrt{s^2}=s\), so no sign was chosen
after squaring.
\end{proof}

\subsection{Strict Hurwitz obstruction}
\label{subsec:V55-Hurwitz}

Let \(a_5,\ldots,a_0\) be the coefficients of
\eqref{eq:V55-Ohmic-quintic}.  Thus
\begin{equation}
 \begin{split}
 a_5&=E,\quad a_4=C_*,\quad
 a_3=A+\frac34E,\\
 a_2&=\frac34C_*+D,\quad
 a_1=\frac34A+B,\quad a_0=1.
 \end{split}
 \label{eq:V55-Ohmic-coefficients}
\end{equation}
Define
\begin{equation}
 S_M:=AD-BC_*+E.
 \label{eq:V55-SM}
\end{equation}

\begin{lemma}[Endpoint third Hurwitz determinant]
\label{lem:V55-Hurwitz-negative}
For every \(M\succ0\),
\begin{equation}
 S_M<0,
 \qquad
 \Delta_3
 :=a_4a_3a_2-a_4^2a_1-a_5a_2^2+a_5a_4a_0
 =C_*S_M-EDa_2<0.
 \label{eq:V55-Delta3}
\end{equation}
\end{lemma}

\begin{proof}
Orthogonally diagonalize the positive transverse block
\(N=\operatorname{diag}(n_1,n_2)\) and write \(u=(p,q)^T\).
The V46 calculation gives
\[
 S_M
 =-n_1(An_1-p^2)-n_2(An_2-q^2)<0,
\]
because both displayed principal minors of \(M\) are positive.
Direct substitution of
\eqref{eq:V55-Ohmic-coefficients} into \(\Delta_3\) cancels the
quadratic \(\tfrac34\)-terms and gives the middle identity in
\eqref{eq:V55-Delta3}.  Since
\(C_*,E,D,a_2>0\), both remaining summands are strictly negative.
\end{proof}

\begin{theorem}[Complete positive Ohmic mediator no-go]
\label{thm:V55-Ohmic-no-go}
Let \(M=M^T\succeq0\) and couple the zero-speed conservative bath
\eqref{eq:V55-zero-speed-action} through the passive negative-feedback
metric graph \(L_0+CsMJ\).  Then the complete normalized stable matrix
has no uniform inverse.
\begin{enumerate}
\item If \(\operatorname{rank}M<3\), it remains singular at
      gravitational glancing.
\item If \(M\succ0\), its determinant has an open or closed-boundary
      stable-frequency zero away from the glancing rank defect.
\end{enumerate}
Thus no constant positive Ohmic matrix supplies both positive energy and
the complete stable Einstein boundary floor.
\end{theorem}

\begin{proof}
The first item is Theorem~\ref{thm:V55-glancing-determinant} with
\(Z=sM\).

For the second, suppose the determinant had no stable zero.  By
Lemma~\ref{lem:V55-endpoint-ratio} and
Theorem~\ref{thm:V55-Ohmic-quintic}, every root of
\(p_M^{(0)}\) would have to lie in the strict left half-plane.  The
polynomial would be Hurwitz, so the Routh--Hurwitz criterion would require
\(\Delta_3>0\).  Lemma~\ref{lem:V55-Hurwitz-negative} gives the
contradiction \(\Delta_3<0\).  Hence a nonzero root satisfies
\(\operatorname{Re}\theta\ge0\), and
Lemma~\ref{lem:V55-endpoint-ratio} realizes it as an open or
closed-boundary stable frequency.  Equation
\eqref{eq:V55-Ohmic-determinant} then makes the metric determinant zero.
\end{proof}

\begin{corollary}[Scalar cubic witness]
\label{cor:V55-scalar-cubic}
For \(M=\gamma I_3\), \(\gamma>0\), the longitudinal factor is
\begin{equation}
 1+\frac{\gamma}{8}w(w^2+3),
 \qquad w=\frac{s}{\lambda}.
 \label{eq:V55-scalar-cubic}
\end{equation}
The equation
\begin{equation}
 w^3+3w+\frac8\gamma=0
 \label{eq:V55-cubic-equation}
\end{equation}
has one negative real root and a complex-conjugate pair with strictly
positive real part.  The latter give two open right-half-plane metric
zeros.
\end{corollary}

\begin{proof}
The cubic derivative \(3w^2+3\) is positive on the real line, so there
is exactly one real root, and it is negative.  Its discriminant is
\[
 -108-\frac{1728}{\gamma^2}<0,
\]
so the other two roots are nonreal.  Their sum is the negative of the
real root, hence each has positive real part.  Apply the construction in
Lemma~\ref{lem:V55-endpoint-ratio}.
\end{proof}

\begin{firewall}[What the Ohmic no-go excludes]
\label{fire:V55-Ohmic-scope}
Theorem~\ref{thm:V55-Ohmic-no-go} includes every constant positive
matrix, including longitudinal--transverse mixing, and not merely the
rotationally diagonal scalar case.  It does not exclude a
frequency-dependent matrix-valued principal impedance, derivative
couplings with an independent boundary state, or a boundary-adapted
radiative quotient.  Any next candidate must still satisfy the exact
passivity card \eqref{eq:V55-passivity-card}; changing the polynomial by
an active fitted symbol would not solve the positive-generator problem.
\end{firewall}

\subsection{Finite certificates and the V56 construction target}
\label{subsec:V55-certificates}

At \(r=1\), exact rational polynomial expansion gives
\begin{equation}
 \frac{\det(L_{\rm g}+CZJ)}{\det Z}=-\frac1{24},
 \qquad
 \frac{\det(L_{\rm g}-CZJ)}{\det Z}=\frac1{24}.
 \label{eq:V55-exact-det-card}
\end{equation}
The kernel card stores
\begin{equation}
 L_{\rm g}R_{\rm g}=0,\quad
 \Lambda_{\rm g}L_{\rm g}=0,\quad
 JR_{\rm g}=I_3,\quad
 \Lambda_{\rm g}C=
 \operatorname{diag}(-1/2,1/2,-1).
 \label{eq:V55-kernel-card}
\end{equation}
For a positive \(M\), a finite endpoint card verifies
\begin{equation}
 \Delta_3-(C_*S_M-EDa_2)=0,
 \qquad
 S_M<0,\qquad \Delta_3<0,
 \label{eq:V55-Hurwitz-card}
\end{equation}
then maps a root with
\(\operatorname{Re}\theta\ge0\) through
\eqref{eq:V55-ratio-inversion} and checks the full \(5\times5\)
determinant.

The successful principal class is now sharply constrained.  V56 should
construct the smallest explicit \(Z^{(1)}(s,k)\) which:
\begin{enumerate}
\item is analytic for \(\operatorname{Re}s>0\), obeys the reality
      symmetry, is homogeneous of degree one, and satisfies
      \eqref{eq:V55-passivity-card};
\item has the glancing floor \eqref{eq:V55-glancing-port-floor};
\item is not a positive square-root speed mixture excluded by V48 and
      is not the constant Ohmic factor \(sM\) excluded by V55;
\item arises from a printed local symmetric-hyperbolic boundary system
      or a positive conservative continuum, rather than a fitted
      frequency response; and
\item makes \(\det(L_0+CZ^{(1)}J)\) zero-free on the complete normalized
      stable hemisphere.
\end{enumerate}
The first candidate should include a noncommuting tangential derivative
or gyroscopic boundary state.  If its energy or determinant fails, V56
must return the exact negative vector or stable root.

\subsection{V55 ledger}
\label{subsec:V55-ledger}

\paragraph{New conclusions proved in V55.}
The scheduled companion audit found no transferable boundary mechanism
and recorded the completed NS V29 and YM V32 states without modifying
either companion file.
Theorem~\ref{thm:V55-glancing-determinant} computes the exact determinant
of an arbitrary three-port correction at gravitational glancing, and
Corollary~\ref{cor:V55-principal-card} derives the necessary order-one
port floor.  Proposition~\ref{prop:V55-Ohmic-ancestry} constructs the
simplest positive conservative order-one bath.  Lemma
\ref{lem:V55-endpoint-ratio}, Theorem
\ref{thm:V55-Ohmic-quintic}, and Lemma
\ref{lem:V55-Hurwitz-negative} extend the exact common-matrix determinant
calculus to the zero-speed endpoint.  Theorem
\ref{thm:V55-Ohmic-no-go} proves that every constant positive Ohmic
matrix still has a complete stable-cone obstruction.

\paragraph{Imported conclusions not repeated.}
V39 supplies the three-dimensional glancing defect, V41 supplies
\(L_0,C,J\), V44/V47 supply conservative half-line bath ancestry, V46
supplies the unsquared common-matrix determinant algebra, V48 excludes
positive speed mixtures in the passive orientation, and V53/V54 supply
the positive metric exchange Gram and the order-zero mediator failure.
The NS V29 and YM V32 results are audited but not imported as
gravitational estimates.

\paragraph{Honest continuum status.}
The necessary principal passive-port card is now explicit, and the
constant Ohmic realization has been completely excluded.  No
frequency-dependent noncommuting passive mediator, boundary-preserving
physical projector, positive stable Einstein time generator, off-shell
ghost domain, nonlinear covariant stress/Bianchi closure, or coupled
Standard-Model matter boundary system has yet been proved.  The full
Standard Model--Einstein continuum is therefore not proved.

\section*{Version V55 append-only record}
\addcontentsline{toc}{section}{Version V55 append-only record}
The complete V54 mathematical source was retained before this
continuation.  Its SHA--256 digest was
\begin{center}\scriptsize\ttfamily
a964602e50cb1660001cd58d80755701c60f978f9e1cf3f147dd6fe88088cbc9.
\end{center}
No mathematical content of V54 was altered.  On the next iteration,
retain this full file, append before its unique active
document-closing command, increment the filename to
\texttt{\_V56.tex}, and remove only the superseded V55 file from this
exact SM note series.  The next scheduled companion audit is V60.

% END APPENDED VERSION V55 MATERIAL
% BEGIN APPENDED VERSION V56 MATERIAL

\clearpage
\section{V56: a positive time--tangential continuum, exact matrix
impedance, and its two-root glancing-index obstruction}
\label{sec:V56-derivative-bath}

\subsection{Purpose and upstream choice}
\label{subsec:V56-purpose}

V55 proved that the passive three-port graph requires an invertible
order-one impedance at gravitational glancing, but that the simplest
Ohmic law \(Z=sM\), \(M\succeq0\), always has a stable determinant zero.
The prescribed next class was a frequency-dependent noncommuting
impedance with a printed local positive ancestry.  This appendix tests
the smallest rotationally covariant derivative coupling with those
properties.

The new bath is not a constant boundary gyroscope of V43.  Its
time--tangential cross derivative is distributed throughout the
auxiliary half-line, and its elimination produces two different
frequency-shifted square roots whose longitudinal--scalar eigenspaces
are mixed.  It is also outside the V48 covariant commutant: in the
\((v_L,v_T,\phi)\) frame its impedance has a nonzero \(L\)-to-\(\phi\)
entry.

The construction meets the local-action, positive-energy, passivity, and
glancing-rank gates.  Nevertheless an unsquared Nyquist calculation
proves that its passive mixed determinant has exactly two open
right-half-plane zeros for every positive coupling strength in a
nonempty subluminal parameter chamber.  Thus derivative mixing changes
the matrix response but does not change the fatal glancing index.

\subsection{The local derivative-coupled half-line action}
\label{subsec:V56-local-action}

Let
\[
 q=(u_1,u_2,\chi)^T
\]
be a real three-component field on the auxiliary half-line \(a\ge0\).
The pair \(u=(u_1,u_2)\) transforms as a tangential spatial vector and
\(\chi\) as a scalar.  For \(A=1,2\), define the symmetric matrices
\begin{equation}
 {\cal B}^1
 =\beta
 \begin{pmatrix}0&0&1\\0&0&0\\1&0&0\end{pmatrix},
 \qquad
 {\cal B}^2
 =\beta
 \begin{pmatrix}0&0&0\\0&0&1\\0&1&0\end{pmatrix},
 \qquad \beta>0.
 \label{eq:V56-B-matrices}
\end{equation}
Equivalently,
\begin{equation}
 q_t^T{\cal B}^A\partial_Aq
 =\beta\left(
 u_{A,t}\partial_A\chi+\chi_t\partial_Au_A
 \right),
 \label{eq:V56-cross-scalar}
\end{equation}
with summation over \(A\).  This is a local \(SO(2)\)-scalar.  For
\(\eta,c>0\), set
\begin{equation}
 \begin{split}
 S_{\eta,c,\beta}[q]
 :=\frac{\eta}{2}\int_{\partial M}\int_0^\infty
 \bigl(&|q_t|^2-|q_a|^2-c^2|\nabla_yq|^2\\
 &+2q_t^T{\cal B}^A\partial_Aq\bigr)
 \,\dd a\,\dd^3y,
 \qquad q(0)=z:=Jx .
 \end{split}
 \label{eq:V56-action}
\end{equation}

\begin{proposition}[Euler equation, positive energy, and boundary power]
\label{prop:V56-energy}
The Euler equation of \eqref{eq:V56-action} is
\begin{equation}
 q_{tt}-q_{aa}-c^2\Delta_yq
 +2{\cal B}^A\partial_t\partial_Aq=0.
 \label{eq:V56-Euler}
\end{equation}
Its canonical energy is
\begin{equation}
 E_{\eta,c,\beta}
 =\frac{\eta}{2}\int\!\!\int
 \left(|q_t|^2+|q_a|^2+c^2|\nabla_yq|^2\right)
 \dd a\,\dd^2y\ge0.
 \label{eq:V56-energy}
\end{equation}
If
\begin{equation}
 t:=-\eta q_a(0),
 \label{eq:V56-traction}
\end{equation}
then, for decaying data or vanishing tangential flux at infinity,
\begin{equation}
 \frac{\dd}{\dd t}E_{\eta,c,\beta}
 =\int_{\partial M}z_t^Tt\,\dd^2y.
 \label{eq:V56-power}
\end{equation}
The cross derivative changes the equation but cancels from the energy.

For a tangential covector of length \(r\), the three zero-normal-
frequency characteristic speeds have magnitudes
\begin{equation}
 c,\qquad
 v_-:=\sqrt{c^2+\beta^2}-\beta,\qquad
 v_+:=\sqrt{c^2+\beta^2}+\beta.
 \label{eq:V56-speeds}
\end{equation}
In particular they are real, \(v_-v_+=c^2\), and the bath cones are
strictly inside gravitational glancing when \(v_+<1\).
\end{proposition}

\begin{proof}
Since the matrices \({\cal B}^A\) are symmetric, variation of the two
derivative factors in \eqref{eq:V56-cross-scalar} gives the same mixed
term.  This proves \eqref{eq:V56-Euler}.  The momentum is
\[
 \frac{\partial{\cal L}}{\partial q_t}
 =\eta(q_t+{\cal B}^A\partial_Aq).
\]
The Legendre expression \(q_t^T\partial{\cal L}/\partial q_t-{\cal L}\)
cancels the cross term and gives \eqref{eq:V56-energy}.

Multiplying \eqref{eq:V56-Euler} by \(\eta q_t^T\) gives the usual time
derivative of \eqref{eq:V56-energy}.  The mixed contribution is the
tangential divergence
\[
 2\eta q_t^T{\cal B}^A\partial_Aq_t
 =\eta\partial_A(q_t^T{\cal B}^Aq_t).
\]
The normal flux then gives
\eqref{eq:V56-power} with \eqref{eq:V56-traction}.

For a rotated tangential frequency \(k=(r,0)\), the matrix
\({\cal B}(k):=k_A{\cal B}^A\) has eigenvalues
\(\beta r,0,-\beta r\).  At zero normal frequency the corresponding
dispersion equations are
\[
 -\omega^2+c^2r^2-2\beta r\omega=0,\quad
 -\omega^2+c^2r^2=0,\quad
 -\omega^2+c^2r^2+2\beta r\omega=0.
\]
Their positive and negative roots have precisely the magnitudes in
\eqref{eq:V56-speeds}.
\end{proof}

The passive Einstein connection uses the V53 sign
\begin{equation}
 p=A(D)x-Ct,\qquad z=Jx.
 \label{eq:V56-passive-graph}
\end{equation}
For every positive exchange Gram \(K_{\gamma,\delta}\) of V53,
\(K_{\gamma,\delta}C=J^T\), so the metric traction power and
\eqref{eq:V56-power} cancel exactly.  The intrinsic graph power residual
of Lemma~\ref{lem:V53-intrinsic-power} is still present; the statement
here is only the exact exchange cancellation.

\subsection{Retarded matrix square root and strict passivity}
\label{subsec:V56-retarded-impedance}

Use a stable Laplace frequency \(\operatorname{Re}s>0\) and a real
tangential covector \(k\).  The decaying equation in \(a\) is
\begin{equation}
 -q_{aa}
 +\left(s^2I_3+c^2|k|^2I_3+2is{\cal B}(k)\right)q=0.
 \label{eq:V56-normal-equation}
\end{equation}
Let \(Q(s,k)\) be the matrix square root of the parenthesis whose three
eigenvalues have positive real part.  Then
\begin{equation}
 q(a)=e^{-aQ(s,k)}z,\qquad
 Z_{\eta,c,\beta}(s,k)=\eta Q(s,k),\qquad
 t=Z_{\eta,c,\beta}z.
 \label{eq:V56-impedance}
\end{equation}

For \(k=(r,0)\), use the port order \((v_L,v_T,\phi)\) and put
\begin{align}
 \kappa_0&:=\sqrt{s^2+c^2r^2},
 \nonumber\\
 \kappa_\pm&:=
 \sqrt{s^2+c^2r^2\mathbin{\pm}2i\beta sr},
 \label{eq:V56-kappas}\\
 S_\kappa&:=\kappa_++\kappa_-,
 \qquad
 D_\kappa:=\kappa_+-\kappa_- .
 \nonumber
\end{align}
All roots are the retarded branches.  Direct diagonalization gives
\begin{equation}
 Q(s,(r,0))
 =
 \begin{pmatrix}
 S_\kappa/2&0&D_\kappa/2\\
 0&\kappa_0&0\\
 D_\kappa/2&0&S_\kappa/2
 \end{pmatrix}.
 \label{eq:V56-Q-explicit}
\end{equation}
The off-diagonal entry is nonzero when \(\beta sr\ne0\).

\begin{theorem}[Analytic causal impedance with positive conservative
ancestry]
\label{thm:V56-passivity}
Assume \(v_+<1\).  The matrix \(Q\) is homogeneous of degree one,
analytic for \(\operatorname{Re}s>0\), obeys the real-field symmetry
\begin{equation}
 Q(\overline s,-k)=\overline{Q(s,k)},
 \label{eq:V56-reality}
\end{equation}
and extends at \(k=0\) as \(Q(s,0)=sI_3\).  Moreover, for every
\(z\in\mathbb C^3\),
\begin{equation}
 \operatorname{Re}\left(
 \overline s\,z^*Z_{\eta,c,\beta}(s,k)z
 \right)
 =
 \eta(\operatorname{Re}s)
 \int_0^\infty
 \left(
 |q_a|^2+(|s|^2+c^2|k|^2)|q|^2
 \right)\dd a
 \ge0.
 \label{eq:V56-strict-passivity}
\end{equation}
Thus \eqref{eq:V56-impedance} satisfies the V55 displacement-to-traction
passivity card and is derived from the positive local action
\eqref{eq:V56-action}.
\end{theorem}

\begin{proof}
The factors
\begin{equation}
 \begin{split}
 \kappa_+^2
 &=(s-iv_-r)(s+iv_+r),\\
 \kappa_-^2
 &=(s-iv_+r)(s+iv_-r)
 \end{split}
 \label{eq:V56-factorized-kappas}
\end{equation}
have their only zeros on the imaginary boundary.  If the imaginary part
of either radicand vanishes in the open half-plane, its real part is
strictly positive.  Hence neither radicand meets the nonpositive real
axis there, and the square roots with positive real part are analytic.
Equations \eqref{eq:V56-kappas}--\eqref{eq:V56-Q-explicit} prove
homogeneity, reality, and the extension at \(k=0\).

Multiply \eqref{eq:V56-normal-equation} by
\(\overline s\,q^*\), integrate in \(a\), and use
\(q_a(0)=-Qz\).  Integration by parts gives
\[
 \overline s\,z^*Qz
 =
 \int_0^\infty
 \overline s\left(
 |q_a|^2+(s^2+c^2|k|^2)|q|^2
 +2is\,q^*{\cal B}(k)q
 \right)\dd a.
\]
The last integrand is purely imaginary because
\({\cal B}(k)\) is Hermitian and
\(\overline s(2is)=2i|s|^2\).  Taking real parts yields
\eqref{eq:V56-strict-passivity}.
\end{proof}

\subsection{The derivative mixing repairs gravitational glancing}
\label{subsec:V56-glancing}

At upper gravitational glancing \(s=ir\), \(r>0\), define
\begin{equation}
 g_+:=\sqrt{1-c^2+2\beta},\qquad
 g_-:=\sqrt{1-c^2-2\beta},\qquad
 g_0:=\sqrt{1-c^2}.
 \label{eq:V56-glancing-g}
\end{equation}
The condition \(v_+<1\) implies \(g_+,g_-,g_0>0\), and
\begin{equation}
 \kappa_+=irg_+,\qquad
 \kappa_-=irg_-,\qquad
 \kappa_0=irg_0,\qquad
 \det Q=-ir^3g_+g_-g_0\ne0.
 \label{eq:V56-glancing-Q}
\end{equation}

\begin{corollary}[Exact full-rank glancing repair]
\label{cor:V56-glancing-repair}
For the passive graph \eqref{eq:V56-passive-graph},
\begin{equation}
 \det\left(L_{\rm g}+C Z_{\eta,c,\beta}(ir,(r,0))J\right)
 =
 \frac{i\eta^3r^5}{24}g_+g_-g_0\ne0.
 \label{eq:V56-glancing-determinant}
\end{equation}
The opposite temporal sheet and every tangential direction are also
invertible.  Hence the obstruction proved below is not a failure of the
V55 glancing port floor.
\end{corollary}

\begin{proof}
Insert
\(\det Z_{\eta,c,\beta}=\eta^3\det Q\) from
\eqref{eq:V56-glancing-Q} into the exact V55 identity
\(\det(L_{\rm g}+CZJ)=-r^2\det Z/24\).
Reality and rotational covariance give the other sheets and directions.
\end{proof}

\subsection{Exact passive determinant factorization}
\label{subsec:V56-determinant}

Away from \(\lambda=0\), retain the V43 response coefficients
\begin{equation}
 \lambda=\sqrt{s^2+r^2},\qquad
 b=\frac1{2\lambda},\qquad
 d=\frac{4\lambda^2-r^2}{8\lambda^3}.
 \label{eq:V56-response}
\end{equation}
The response matrix is \(R=\operatorname{diag}(d,b,b)\).

\begin{theorem}[Transverse and mixed determinant factors]
\label{thm:V56-factorization}
For the passive metric matrix
\begin{equation}
 L_{\rm der}:=L_0+C Z_{\eta,c,\beta}J
 \label{eq:V56-effective-matrix}
\end{equation}
and \(\operatorname{Re}s>0\),
\begin{equation}
 \frac{\det L_{\rm der}}{\det L_0}
 =T_{\eta,c}(s,r)H_{\eta,c,\beta}(s,r),
 \label{eq:V56-full-factorization}
\end{equation}
where
\begin{align}
 T_{\eta,c}
 &:=1+\eta b\kappa_0,
 \label{eq:V56-T-factor}\\
 H_{\eta,c,\beta}
 &:=1+\frac{\eta}{2}(d+b)(\kappa_++\kappa_-)
       +\eta^2db\,\kappa_+\kappa_- .
 \label{eq:V56-H-factor}
\end{align}
No square root has been squared or eliminated in this identity.
\end{theorem}

\begin{proof}
The matrix determinant lemma and the V43 identity
\(JL_0^{-1}C=R\) give
\[
 \frac{\det L_{\rm der}}{\det L_0}
 =\det(I_3+\eta RQ).
\]
The transverse row of \eqref{eq:V56-Q-explicit} gives
\eqref{eq:V56-T-factor}.  In the \((L,\phi)\) block, direct expansion
gives
\begin{align*}
 &\left(1+\frac{\eta dS_\kappa}{2}\right)
  \left(1+\frac{\eta bS_\kappa}{2}\right)
 -\frac{\eta^2dbD_\kappa^2}{4}\\
 &\qquad
 =1+\frac{\eta}{2}(d+b)S_\kappa
   +\eta^2db\,\kappa_+\kappa_-,
\end{align*}
because \(S_\kappa^2-D_\kappa^2=4\kappa_+\kappa_-\).
This proves the result.
\end{proof}

The transverse factor has no stable zero.  Indeed,
\[
 T_{\eta,c}=1+\eta\Theta_c,\qquad
 \Theta_c=\frac{\kappa_0}{2\lambda}.
\]
In the open half-plane both retarded roots have positive real part, so
their quotient cannot be a negative real number: an identity
\(\kappa_0=-a\lambda\), \(a>0\), would give opposite signs for their real
parts.  On the upper imaginary boundary the quotient is positive real
for \(0\le\omega<c\), positive imaginary for \(c<\omega<1\), and positive
real for \(\omega>1\); the lower trace is its conjugate.  It therefore
never equals \(-1/\eta\), including at the finite thresholds.

\subsection{The complete imaginary-boundary trace}
\label{subsec:V56-boundary}

Put
\begin{equation}
 \omega_{\rm g}:=\frac{\sqrt3}{2}
 \label{eq:V56-response-zero}
\end{equation}
and impose the nonempty chamber
\begin{equation}
 0<v_-<v_+<\omega_{\rm g}<1.
 \label{eq:V56-low-cone-chamber}
\end{equation}
Set \(r=1\) by homogeneity and take upper stable boundary values
\(s=i\omega+0\).

\begin{lemma}[No mixed-factor boundary zero]
\label{lem:V56-boundary-nonzero}
Under \eqref{eq:V56-low-cone-chamber},
\[
 H_{\eta,c,\beta}(i\omega+0,1)\ne0
 \qquad(\omega\ge0,\ \omega\ne1)
\]
for every \(\eta>0\).  More precisely:
\begin{enumerate}
\item for \(0\le\omega<v_-\), the factor is real and greater than one;
\item for \(v_-<\omega<v_+\), its real part is greater than one and its
      imaginary part is positive;
\item for \(v_+<\omega<\omega_{\rm g}\), its imaginary part is positive;
\item for \(\omega_{\rm g}<\omega<1\), its real part is greater than one;
\item for \(\omega>1\), it is real and greater than one.
\end{enumerate}
The same statements hold at the finite interval endpoints by
continuity.
\end{lemma}

\begin{proof}
Below gravitational glancing, \(b>0\), while \(d>0\) for
\(\omega<\omega_{\rm g}\) and \(d<0\) for
\(\omega_{\rm g}<\omega<1\).

For \(0\le\omega<v_-\), both \(\kappa_\pm\) are positive real, so every
nonconstant term in \eqref{eq:V56-H-factor} is positive.  For
\(v_-<\omega<v_+\), write
\[
 \kappa_+=i\alpha,\qquad
 \kappa_-=\rho,\qquad \alpha,\rho>0.
\]
Since this whole interval lies below \(\omega_{\rm g}\), put
\(K=\eta(d+b)/2>0\) and \(L=\eta^2db>0\).  Then
\[
 H=1+K\rho+i(K\alpha+L\alpha\rho),
\]
which proves the second assertion.

For \(v_+<\omega<1\), write
\(\kappa_+=i\alpha\), \(\kappa_-=i\rho\).
Below \(\omega_{\rm g}\),
\[
 \operatorname{Im}H
 =\frac{\eta}{2}(d+b)(\alpha+\rho)>0.
\]
Above \(\omega_{\rm g}\), \(db<0\) and
\(\kappa_+\kappa_-=-\alpha\rho\), hence
\[
 \operatorname{Re}H
 =1+\eta^2|d|b\alpha\rho>1.
\]

Finally, for \(\omega>1\), write
\[
 b=-iB,\quad d=-iD,\quad
 \kappa_+=i\alpha,\quad\kappa_-=i\rho,
 \qquad B,D,\alpha,\rho>0.
\]
Both nonconstant terms in \eqref{eq:V56-H-factor} are then positive
real.  Continuity covers the thresholds and \(d=0\).
\end{proof}

\subsection{The fourth-order glancing turn and two open roots}
\label{subsec:V56-index}

The mixed factor has the reality symmetry
\begin{equation}
 H_{\eta,c,\beta}(\overline s,r)
 =\overline{H_{\eta,c,\beta}(s,r)}.
 \label{eq:V56-H-reality}
\end{equation}
Indeed, conjugation interchanges \(\kappa_+\) and \(\kappa_-\), while
\eqref{eq:V56-H-factor} uses only their sum and product.

\begin{theorem}[Exact two-root obstruction]
\label{thm:V56-two-root-no-go}
Assume \eqref{eq:V56-low-cone-chamber}.  For every \(\eta>0\), the
factor \(H_{\eta,c,\beta}(\,\cdot\,,r)\) has exactly two zeros, counted
with multiplicity, in \(\operatorname{Re}s>0\) for each \(r>0\).
They are nonreal and form a conjugate pair.  The transverse factor has
no stable zero, so the full passive matrix
\(L_0+CZ_{\eta,c,\beta}J\) has exactly the same two open determinant
zeros.

Consequently the positive local continuum \eqref{eq:V56-action}
repairs gravitational glancing and satisfies the passive exchange law,
but it cannot supply a stable Einstein boundary generator.
\end{theorem}

\begin{proof}
Set \(r=1\).  The factor \(H\) is analytic in the open right half-plane:
all three retarded roots are analytic there and \(\lambda\ne0\).
Lemma~\ref{lem:V56-boundary-nonzero} excludes every finite
imaginary-boundary zero away from gravitational glancing.

Follow the upper imaginary boundary from \(0\) to \(i\infty\).  Before
\(\omega_{\rm g}\), the proof of
Lemma~\ref{lem:V56-boundary-nonzero} keeps \(H\) on the positive real
axis, in the first quadrant, or in the open upper half-plane.  At
\(\omega_{\rm g}\) it is in the first quadrant.  Between
\(\omega_{\rm g}\) and \(1\), it lies in the open right half-plane.
As \(\omega\uparrow1\), its positive real fourth-order term dominates,
so its continuous argument returns to zero.  Thus the boundary part
below glancing has net argument change zero.  Above glancing, \(H>1\)
is positive real.

It remains to compute the right-hand glancing indentation
\[
 s=i+\epsilon e^{i\varphi},
 \qquad -\frac{\pi}{2}\le\varphi\le\frac{\pi}{2}.
\]
Uniformly there,
\begin{align}
 \lambda^2&=2i\epsilon e^{i\varphi}(1+o(1)),
 \nonumber\\
 \kappa_+\kappa_-&=-g_+g_-+o(1),
 \nonumber\\
 H&=\frac{\eta^2g_+g_-}{16\lambda^4}(1+o(1)).
 \label{eq:V56-glancing-asymptotic}
\end{align}
The continuous argument of \(\lambda\) increases from \(0\) to
\(\pi/2\).  Therefore \(\lambda^{-4}\) contributes exactly
\(-2\pi\).  The complete upper path has
\begin{equation}
 \Delta_{\gamma_+}\arg H=-2\pi.
 \label{eq:V56-upper-index}
\end{equation}

On a large right-half-plane semicircle,
\[
 H(s,1)\longrightarrow
 1+\eta+\frac{\eta^2}{4}
 =\left(1+\frac{\eta}{2}\right)^2>0
\]
uniformly.  Apply the argument principle first on a shifted truncated
half-disk and then let the radius tend to infinity and the shift tend to
zero.  The positively oriented vertical boundary traverses the upper
path in reverse, contributing \(+2\pi\); the conjugate lower path
contributes another \(+2\pi\).  The large arc contributes zero.
Finite bath thresholds contribute no winding because their boundary
values are nonzero and continuous.  The total change is \(4\pi\), so
\(H\) has exactly two open zeros and no poles.

For real \(s>0\), \(\kappa_-\) is the conjugate of \(\kappa_+\), while
\(b,d>0\).  Hence \(H(s,1)>1\), so neither zero is real.
Equation \eqref{eq:V56-H-reality} makes the two zeros a conjugate pair.
The zero-free transverse factor and
Theorem~\ref{thm:V56-factorization} transfer the count to the complete
metric determinant.  Homogeneity gives the statement for every \(r>0\).
\end{proof}

\begin{remark}[Why noncommutation does not change this index]
\label{rem:V56-index-mechanism}
The longitudinal and scalar channels are genuinely mixed, but their
two-by-two determinant contains the product \(db\det Q_{L\phi}\).
Near gravitational glancing this term is a positive multiple of
\(\lambda^{-4}\).  Its full \(-2\pi\) turn is already an integer
Nyquist obstruction.  The matrix mixing changes the finite boundary
curve and the locations of the two roots, but not this leading turn.
\end{remark}

\subsection{An exact parameter point and finite certificate}
\label{subsec:V56-certificate}

The chamber \eqref{eq:V56-low-cone-chamber} is nonempty.  Take
\begin{equation}
 c^2=\frac38,\qquad
 \beta=\frac18,\qquad
 \eta=1.
 \label{eq:V56-rational-point}
\end{equation}
Then
\begin{equation}
 v_-=\frac12,\qquad
 v_+=\frac34,\qquad
 c=\sqrt{\frac38},
 \label{eq:V56-rational-speeds}
\end{equation}
so \(v_+<\sqrt3/2\).  At \(r=1\), one of the two open zeros is
\begin{equation}
 s_*
 =0.04309076283651219425
  +0.94101937146740150118\,i,
 \label{eq:V56-root-certificate}
\end{equation}
with
\begin{equation}
 \lambda_*
 =0.35927389979459602678
  +0.11286442623204293350\,i.
 \label{eq:V56-lambda-certificate}
\end{equation}
Direct seventy-digit evaluation gives
\[
 |H(s_*,1)|<1.1\times10^{-70},
 \qquad
 |\det(L_0+CQJ)(s_*,(1,0))|<3.6\times10^{-72}.
\]
At upper glancing the exact determinant is
\[
 \det(L_{\rm g}+CQJ)
 =0.01886898013826749814\,i,
\]
in agreement with \eqref{eq:V56-glancing-determinant}.  An adaptive
right-half-plane contour gives winding number two.  These values audit
the exact factorization and branch assignments; the proof of
Theorem~\ref{thm:V56-two-root-no-go} is the boundary sign and argument-
principle calculation.

\begin{firewall}[Exact scope of the V56 no-go]
\label{fire:V56-scope}
Theorem~\ref{thm:V56-two-root-no-go} closes the direct Dirichlet
half-line action \eqref{eq:V56-action} in the chamber
\(v_+<\sqrt3/2\), for every positive gain.  It does not close:
\begin{enumerate}
\item unequal positive kinetic or stiffness Grams that make the
      longitudinal--scalar eigenvectors depend on \(s\);
\item derivative coupling in the boundary trace or traction itself;
\item a passive dynamic canonical transformer between \(z\) and the
      bath trace;
\item a matrix-valued continuum whose glancing impedance has a different
      determinant phase; or
\item the half-space York/Hodge projector left open in V54.
\end{enumerate}
No conclusion beyond the printed action and parameter chamber is
inferred from the two-root index.
\end{firewall}

\subsection{V56 acquisition ledger and V57 target}
\label{subsec:V56-ledger}

\paragraph{New conclusions proved in V56.}
Proposition~\ref{prop:V56-energy} constructs the first local positive
time--tangential derivative bath in this programme and proves its exact
energy and flux.  Theorem~\ref{thm:V56-passivity} derives its analytic
matrix square-root impedance and the V55 passivity inequality.
Corollary~\ref{cor:V56-glancing-repair} proves that the noncommuting
order-one port repairs the full three-dimensional glancing defect.
Theorem~\ref{thm:V56-factorization} retains the mixed source incidence
in an unsquared determinant.  Lemma~\ref{lem:V56-boundary-nonzero} and
Theorem~\ref{thm:V56-two-root-no-go} prove that the passive branch
nevertheless has exactly two open roots.

\paragraph{Imported conclusions not repeated.}
V41 and V43 supply \(L_0,C,J\) and the exact bulk response
\(R=\operatorname{diag}(d,b,b)\).  V46 supplies only the scalar
transverse ratio-plane fact.  V53 fixes the positive-energy exchange
sign, and V55 supplies the arbitrary-port glancing determinant.
The V56 matrix impedance, boundary classification, and fourth-order
index are new.

\begin{programme}[V57 upstream dynamic-transformer target]
\label{prog:V57-target}
Do not add another direct Dirichlet square-root bath with the same
glancing phase.  Instead:
\begin{enumerate}
\item derive the most general minimal local canonical boundary
      transformer with a positive time matrix and a three-port trace;
\item eliminate it without dropping its matrix numerator, obtaining the
      homogeneous order-one matrix Moebius transform of a positive
      continuum impedance;
\item prove its passivity from the combined storage identity rather than
      from a fitted frequency response;
\item compute the determinant phase of its glancing principal matrix and
      require compensation of the V56 fourth-order turn; and
\item test the complete closed stable cone.  Return an explicit negative
      energy vector, boundary zero, or open root if any gate fails.
\end{enumerate}
A successful transformer must also leave a local covariant source
incidence that can later be varied with respect to the metric.
\end{programme}

\paragraph{Honest continuum status.}
A frequency-dependent noncommuting passive order-one mediator with local
positive ancestry has now been constructed and completely tested, but it
has two open stable-frequency determinant zeros.  Therefore a positive
stable Einstein--bath generator is still missing.  The boundary
York/Hodge domain, off-shell ghost complex, nonlinear covariant
stress/Bianchi closure, and the coupled Higgs, fermion, neutrino, and
Yang--Mills sectors remain open.  The full Standard Model--Einstein
continuum is not proved.

\section*{Version V56 append-only record}
\addcontentsline{toc}{section}{Version V56 append-only record}
The complete V55 mathematical source was retained before this
continuation.  Its SHA--256 digest was
\begin{center}\scriptsize\ttfamily
bac7765373dedc4cc99a808d30a53b029776d567c7c3130abd92c1fdce328e1e.
\end{center}
No mathematical content of V55 was altered.  On the next iteration,
retain this full file, append before its unique active
document-closing command, increment the filename to
\texttt{\_V57.tex}, and remove only the superseded V56 file from this
exact SM note series.  The next scheduled companion audit remains V60.

% END APPENDED VERSION V56 MATERIAL

% BEGIN APPENDED VERSION V57 MATERIAL

\clearpage
\section{V57: the complete principal canonical transformer and its
unavoidable low-cone boundary zero}
\label{sec:V57-canonical-transformer}

\subsection{Purpose and the upstream reduction}
\label{subsec:V57-purpose}

V56 constructed a local positive time--tangential continuum whose
three-port impedance is order one, passive, noncommuting, and invertible
at gravitational glancing.  Its mixed determinant nevertheless has two
open roots because the direct Dirichlet-to-Neumann graph retains a
fourth-order glancing turn.  The V57 handoff therefore asked for a
dynamic canonical transformer between the metric port and a positive
continuum.

This appendix derives the complete constant principal transformer with
three external coordinates \(z\) and three bath coordinates \(y\).
Modulo time-corner terms, every real quadratic local action with one
time derivative is governed by a skew six-by-six matrix.  Its general
port relation contains two self-gyroscopes and one cross transformer.
Eliminating a positive square-root bath gives the homogeneous matrix
linear-fractional response
\[
 Z_{\rm tr}
 =s\Gamma_z+s^2G(\kappa_vM-s\Gamma_y)^{-1}G^T.
\]
The formula keeps the complete numerator and is passive by an exact
power identity.

The result is again negative, but at an earlier and simpler frequency
than V56.  For every regular transformer \(G\in GL(3,\mathbb R)\), all
self-gyroscopes, every \(M\succ0\), and every bath speed below the zero
of the longitudinal Einstein response, a determinant zero occurs on
the low imaginary-frequency boundary before the first internal
transformer pole.  Thus neither a glancing calculation nor a fitted
gain can rescue this canonical class.  The repeated three-port
obstruction directs V58 upstream to the boundary physical quotient.

\subsection{The complete first-order canonical time form}
\label{subsec:V57-time-form}

Let
\begin{equation}
 w:=\binom{z}{y}\in\mathbb R^6,
 \qquad z,y\in\mathbb R^3,
 \label{eq:V57-w}
\end{equation}
and consider a real constant local quadratic density containing exactly
one time derivative.  Modulo a time-corner density it has the form
\begin{equation}
 {\cal L}_{\rm tr}^{(1)}
 =\frac12w^T\Omega w_t,
 \qquad
 \Omega^T=-\Omega.
 \label{eq:V57-canonical-density}
\end{equation}
Write its blocks as
\begin{equation}
 \Omega=
 \begin{pmatrix}
 \Gamma_z&G\\
 -G^T&-\Gamma_y
 \end{pmatrix},
 \qquad
 \Gamma_z^T=-\Gamma_z,\quad
 \Gamma_y^T=-\Gamma_y.
 \label{eq:V57-Omega-blocks}
\end{equation}
Define the external effort \(r\) and the effort \(t\) delivered to the
bath by assigning the two Euler covectors as \((r,-t)\).  Then
\begin{equation}
 r=\Gamma_z z_t+Gy_t,
 \qquad
 t=G^Tz_t+\Gamma_y y_t.
 \label{eq:V57-transformer-relations}
\end{equation}

\begin{proposition}[Completeness and exact lossless power]
\label{prop:V57-transformer-complete}
Modulo time corners, equations
\eqref{eq:V57-Omega-blocks}--\eqref{eq:V57-transformer-relations}
describe every real constant quadratic six-coordinate action with one
time derivative.  For all real port velocities,
\begin{equation}
 z_t^Tr=y_t^Tt.
 \label{eq:V57-lossless-power}
\end{equation}
Thus the transformer stores no Hamiltonian energy and transfers power
without loss.  It is regular on all three port directions when
\begin{equation}
 G\in GL(3,\mathbb R).
 \label{eq:V57-G-regular}
\end{equation}
\end{proposition}

\begin{proof}
For a general coefficient matrix \(K\), decompose
\(K=K_{\rm sym}+K_{\rm skew}\).  The symmetric part satisfies
\[
 \frac12w^TK_{\rm sym}w_t
 =\frac14\partial_t(w^TK_{\rm sym}w)
\]
and is a time corner.  The unique surviving coefficient is the skew
matrix \(\Omega:=K_{\rm skew}\).  Every real skew six-by-six matrix has
the block form \eqref{eq:V57-Omega-blocks}.  Variation of
\eqref{eq:V57-canonical-density} gives \(\Omega w_t\), which is exactly
\eqref{eq:V57-transformer-relations} with the declared effort signs.

Skewness gives
\[
 z_t^T\Gamma_zz_t=0,\qquad
 y_t^T\Gamma_yy_t=0,
\]
and the two cross products
\(z_t^TGy_t\) and \(y_t^TG^Tz_t\) agree.  This proves
\eqref{eq:V57-lossless-power}.  Finally, \(G\) must be invertible for the
cross relation to carry an arbitrary three-component port velocity in
both directions.
\end{proof}

\begin{corollary}[No positive canonical time matrix in the minimal
coordinate class]
\label{cor:V57-no-positive-time-form}
The principal Euler time matrix of
\eqref{eq:V57-canonical-density} is skew and therefore cannot be
symmetric positive definite.  In fact,
\begin{equation}
 w^T\Omega w=0
 \qquad\hbox{for every real }w.
 \label{eq:V57-skew-zero}
\end{equation}
Adding a positive constant Hamiltonian
\(\frac12w^TKw\), \(K\succ0\), adds the zeroth-order Euler term \(Kw\)
but does not change the homogeneous degree-one transformer
\eqref{eq:V57-transformer-relations}.  A genuinely positive evolution
time matrix therefore requires an enlarged momentum/state system; it
cannot be obtained by relabelling the minimal canonical time form.
\end{corollary}

\begin{proof}
Equation \eqref{eq:V57-skew-zero} is the defining quadratic identity for
a real skew matrix.  A symmetric positive definite real matrix has a
strictly positive quadratic form on nonzero vectors, so the two
properties are incompatible.  A Hamiltonian without derivatives
contributes only its ordinary gradient to the Euler equation and hence
is subprincipal relative to \(\Omega\partial_t\).
\end{proof}

This corollary is an algebraic classification of the stated
six-coordinate first-order action.  It does not exclude a doubled
symmetric-hyperbolic state with its own positive storage.

\subsection{Connection to a positive square-root continuum}
\label{subsec:V57-bath-connection}

Attach the V45 positive common-speed half-line bath at the \(y\)-port:
\begin{equation}
 t=Z_b(s,k)y,\qquad
 Z_b(s,k):=\kappa_v(s,k)M,
 \qquad
 \kappa_v:=\sqrt{s^2+v^2|k|^2},
 \label{eq:V57-load}
\end{equation}
where \(0<v<1\), \(M=M^T\succ0\), and the square root is retarded.
The bath energy is positive and
\begin{equation}
 \operatorname{Re}
 \left(\overline s\,y^*Z_b(s,k)y\right)>0
 \quad
 (\operatorname{Re}s>0,\ y\ne0).
 \label{eq:V57-load-strict}
\end{equation}

For a Laplace--Fourier mode,
\eqref{eq:V57-transformer-relations} and \eqref{eq:V57-load} give
\begin{equation}
 \left(\kappa_vM-s\Gamma_y\right)y=sG^Tz.
 \label{eq:V57-internal-equation}
\end{equation}

\begin{lemma}[Open-half-plane internal invertibility]
\label{lem:V57-internal-invertible}
The matrix
\begin{equation}
 D_y(s,k):=\kappa_v(s,k)M-s\Gamma_y
 \label{eq:V57-Dy}
\end{equation}
is invertible for every \(\operatorname{Re}s>0\).
\end{lemma}

\begin{proof}
If \(D_yy=0\), then \(Z_by=s\Gamma_yy\).  Multiplication by
\(\overline s\,y^*\) gives
\[
 \overline s\,y^*Z_by
 =|s|^2y^*\Gamma_yy.
\]
The right side is purely imaginary because \(\Gamma_y\) is real skew,
whereas the real part of the left side is strictly positive by
\eqref{eq:V57-load-strict} unless \(y=0\).  Hence the kernel is zero.
\end{proof}

Elimination of \(y\) is now exact:
\begin{equation}
 y=sD_y^{-1}G^Tz,
 \qquad
 r=Z_{\rm tr}(s,k)z,
 \label{eq:V57-y-elimination}
\end{equation}
with
\begin{equation}
 \boxed{\quad
 Z_{\rm tr}(s,k)
 =s\Gamma_z+s^2G
   \left(\kappa_vM-s\Gamma_y\right)^{-1}G^T.
 \quad}
 \label{eq:V57-transformer-impedance}
\end{equation}

\begin{theorem}[Homogeneous matrix transform and inherited passivity]
\label{thm:V57-transformer-passive}
The impedance \eqref{eq:V57-transformer-impedance} is analytic in the
open Laplace half-plane, homogeneous of degree one, and obeys the
real-field symmetry.  It also satisfies
\begin{equation}
 \operatorname{Re}
 \left(\overline s\,z^*Z_{\rm tr}(s,k)z\right)
 =
 \operatorname{Re}
 \left(\overline s\,y^*Z_b(s,k)y\right)
 \ge0.
 \label{eq:V57-transformer-passivity}
\end{equation}
The identity follows from the combined local transformer--bath power,
not from a fitted transfer function.
\end{theorem}

\begin{proof}
Lemma~\ref{lem:V57-internal-invertible} proves analyticity.
Every term in \(D_y\) has degree one, so
\(s^2GD_y^{-1}G^T\), and hence \(Z_{\rm tr}\), has degree one.
Reality follows from the real coefficient matrices and the retarded
root.

For the complex mode, subtract the two port supplies obtained from
\eqref{eq:V57-transformer-relations}:
\begin{align*}
 \overline s\,z^*r-\overline s\,y^*t
 =|s|^2\bigl(
 z^*\Gamma_zz+z^*Gy-y^*G^Tz-y^*\Gamma_yy
 \bigr).
\end{align*}
Each self term is purely imaginary, and the two cross terms are complex
conjugates with opposite signs.  The real part therefore vanishes.
Using \(r=Z_{\rm tr}z\) and \(t=Z_by\) proves
\eqref{eq:V57-transformer-passivity}.
\end{proof}

The passive Einstein graph remains
\begin{equation}
 p=A(D)x-Cr,\qquad z=Jx.
 \label{eq:V57-passive-metric-graph}
\end{equation}
Because \eqref{eq:V57-lossless-power} transfers the bath power to the
external port, the V53 positive exchange Gram cancels the metric
traction exchange exactly.  The V53 intrinsic \(A(D)\)-power residual
remains separate.

\subsection{The augmented stable matrix}
\label{subsec:V57-augmented-matrix}

For \(k=(r,0)\), retain
\begin{equation}
 R(s,r)=\operatorname{diag}(d,b,b),\qquad
 b=\frac1{2\lambda},\qquad
 d=\frac{4\lambda^2-r^2}{8\lambda^3},\qquad
 \lambda=\sqrt{s^2+r^2}.
 \label{eq:V57-R}
\end{equation}
The direct metric--transformer--bath system is the eight-by-eight block
matrix
\begin{equation}
 {\mathfrak L}_{\rm tr}(s,k)
 :=
 \begin{pmatrix}
 L_0+Cs\Gamma_zJ&CsG\\
 -sG^TJ&\kappa_vM-s\Gamma_y
 \end{pmatrix}.
 \label{eq:V57-augmented}
\end{equation}
When \(D_y\) is invertible, its Schur complement is
\begin{equation}
 L_{\rm tr}
 =L_0+CZ_{\rm tr}J.
 \label{eq:V57-effective-metric}
\end{equation}
For \(\lambda\ne0\),
\begin{equation}
 \frac{\det L_{\rm tr}}{\det L_0}
 =\det\left(I_3+RZ_{\rm tr}\right).
 \label{eq:V57-response-determinant}
\end{equation}
These identities retain all three transformer directions; no diagonal
or scalar reduction has been taken.

\begin{proof}
Taking the Schur complement of the lower-right block in
\eqref{eq:V57-augmented} gives
\[
 L_0+Cs\Gamma_zJ
 +Cs^2GD_y^{-1}G^TJ
 =L_0+CZ_{\rm tr}J.
\]
The matrix determinant lemma and \(JL_0^{-1}C=R\) then give
\eqref{eq:V57-response-determinant}.
\end{proof}

\subsection{The first-positive-interval lemma}
\label{subsec:V57-first-positive}

Let
\begin{equation}
 \omega_{\rm g}:=\frac{\sqrt3}{2},
 \qquad
 0<v<\omega_{\rm g},
 \label{eq:V57-low-speed}
\end{equation}
and set \(r=1\) by homogeneity.  On the low upper imaginary boundary
\(s=i\omega+0\), \(0\le\omega<v\), put
\begin{equation}
 \kappa(\omega):=\sqrt{v^2-\omega^2},
 \qquad
 {\mathsf H}_y(\omega)
 :=\kappa(\omega)M-i\omega\Gamma_y.
 \label{eq:V57-Hy}
\end{equation}
This is Hermitian.  At \(\omega=0\),
\({\mathsf H}_y(0)=vM\succ0\).

\begin{lemma}[The positive interval ends in a zero eigenvalue]
\label{lem:V57-positive-interval}
There is a number \(\omega_*\in(0,v]\) such that
\begin{equation}
 {\mathsf H}_y(\omega)\succ0
 \quad(0\le\omega<\omega_*),
 \qquad
 \lambda_{\min}{\mathsf H}_y(\omega)\longrightarrow0
 \quad(\omega\uparrow\omega_*).
 \label{eq:V57-positive-interval}
\end{equation}
If \(\Gamma_y\ne0\), then \(\omega_*<v\).  If
\(\Gamma_y=0\), then \(\omega_*=v\).
\end{lemma}

\begin{proof}
Positive definiteness at zero and continuity give a nonempty positive
interval.  If \(\Gamma_y=0\), then
\({\mathsf H}_y=\kappa M\), which is positive for \(\omega<v\) and tends
to zero at \(v\).

Suppose \(\Gamma_y\ne0\).  A nonzero real three-by-three skew matrix has
eigenvalues \(0,\pm i\gamma\) for some \(\gamma>0\).  Consequently the
Hermitian matrix \(-i\Gamma_y\) has eigenvalues \(0,\gamma,-\gamma\).
As \(\omega\uparrow v\),
\[
 {\mathsf H}_y(\omega)\longrightarrow-iv\Gamma_y,
\]
which has a negative eigenvalue.  Hence positive definiteness is lost at
a first \(\omega_*<v\).  Continuity of the least eigenvalue gives
\eqref{eq:V57-positive-interval}.
\end{proof}

\subsection{The universal low-cone boundary zero}
\label{subsec:V57-boundary-zero}

For \(0\le\omega<\omega_*\), define
\begin{equation}
 Z_{\rm tr}^{\rm low}(\omega)
 :=i\omega\Gamma_z
 -\omega^2G{\mathsf H}_y(\omega)^{-1}G^T.
 \label{eq:V57-low-impedance}
\end{equation}
Both terms are Hermitian.  Since \(\omega<v<\omega_{\rm g}\), the metric
response
\begin{equation}
 R_\omega
 :=R(i\omega+0,1)
 =\operatorname{diag}(d_\omega,b_\omega,b_\omega)
 \label{eq:V57-low-R}
\end{equation}
is real positive definite.  Put
\begin{equation}
 {\mathsf K}(\omega)
 :=I_3+R_\omega^{1/2}
 Z_{\rm tr}^{\rm low}(\omega)
 R_\omega^{1/2}.
 \label{eq:V57-K}
\end{equation}
Then \({\mathsf K}(\omega)\) is Hermitian and
\begin{equation}
 \det{\mathsf K}(\omega)
 =\det\left(I_3+R_\omega
 Z_{\rm tr}^{\rm low}(\omega)\right).
 \label{eq:V57-K-determinant}
\end{equation}

\begin{theorem}[Boundary zero before the first transformer pole]
\label{thm:V57-boundary-zero}
Let \(G\in GL(3,\mathbb R)\), \(M\succ0\),
\(\Gamma_z^T=-\Gamma_z\), and
\(\Gamma_y^T=-\Gamma_y\).  If \(0<v<\sqrt3/2\), then there exists
\begin{equation}
 0<\omega_0<\omega_*\le v
 \label{eq:V57-zero-location}
\end{equation}
such that
\begin{equation}
 \det{\mathfrak L}_{\rm tr}(i\omega_0+0,(1,0))=0.
 \label{eq:V57-augmented-zero}
\end{equation}
At this frequency,
\[
 \lambda=\sqrt{1-\omega_0^2}>0,\qquad
 \kappa_v=\sqrt{v^2-\omega_0^2}>0,\qquad
 {\mathsf H}_y(\omega_0)\succ0.
\]
Thus the zero is a normally decaying stable-boundary mode and occurs
strictly before any internal transformer pole or bath threshold.
No regular member of the complete canonical class
\eqref{eq:V57-Omega-blocks} supplies a uniform passive
Kreiss--Lopatinski floor.
\end{theorem}

\begin{proof}
At \(\omega=0\), equation \eqref{eq:V57-low-impedance} gives
\({\mathsf K}(0)=I_3\).  We show that its least eigenvalue becomes
negative before \(\omega_*\).

Let
\[
 P(\omega):=R_\omega^{1/2}G.
\]
The matrices \(R_\omega\) remain uniformly positive and bounded on
\([0,\omega_*]\), because \(\omega_*\le v<\omega_{\rm g}\).  Hence
\(P(\omega)\) and \(P(\omega)^{-1}\) remain bounded.  From
Lemma~\ref{lem:V57-positive-interval},
\[
 \lambda_{\max}\left(
 P{\mathsf H}_y^{-1}P^T
 \right)\longrightarrow+\infty
 \qquad(\omega\uparrow\omega_*).
\]
For example, if \(e(\omega)\) is a unit least-eigenvalue vector of
\({\mathsf H}_y\), testing on
\(P(\omega)^{-T}e(\omega)\) gives this divergence after normalization.

The gyroscopic term
\[
 i\omega R_\omega^{1/2}\Gamma_zR_\omega^{1/2}
\]
is Hermitian and uniformly bounded.  Weyl's inequality applied to
\eqref{eq:V57-K} therefore gives
\begin{equation}
 \lambda_{\min}{\mathsf K}(\omega)
 \longrightarrow-\infty
 \qquad(\omega\uparrow\omega_*).
 \label{eq:V57-K-negative}
\end{equation}
Since \(\lambda_{\min}{\mathsf K}(0)=1\), continuity supplies
\(\omega_0\in(0,\omega_*)\) with
\(\lambda_{\min}{\mathsf K}(\omega_0)=0\).

At \(\omega_0\), the internal block \(D_y={\mathsf H}_y\) is positive
and invertible.  Equations
\eqref{eq:V57-response-determinant} and
\eqref{eq:V57-K-determinant} make the effective metric determinant zero,
and the Schur determinant identity for
\eqref{eq:V57-augmented} gives
\eqref{eq:V57-augmented-zero}.  The inequalities in
\eqref{eq:V57-zero-location} give the two strictly positive normal roots.
\end{proof}

\begin{remark}[Why the low-cone crossing is upstream of glancing]
\label{rem:V57-low-before-glancing}
V56 failed through a fourth-order turn at \(\omega=1\).  The canonical
transformer fails earlier.  On the interval \(0<\omega<v\), the positive
bath impedance is real, but the two time derivatives in the transformed
term contribute \(s^2=-\omega^2\).  The resulting negative Hermitian
direction diverges as the internal positive block approaches its first
zero.  Since the Einstein response is positive on
\(0<\omega<\sqrt3/2\), an eigenvalue must cross \(-1\) before that
endpoint.
\end{remark}


\subsection{The reciprocal square-root endpoint}

The gyro-free endpoint already separates upper glancing success from low-cone failure.  Set
\[
 \Gamma_z=\Gamma_y=0,
 \qquad
 N:=GM^{-1}G^\top .
\]
Since \(G\in GL(3,\mathbb R)\) and \(M\succ0\), one has \(N\succ0\), and the transformer impedance becomes
\[
 Z_{\rm rec}(s,k)
   =\frac{s^2}{\kappa_v(s,k)}\,N .
\]
Thus this endpoint is reciprocal and is generated by the same positive scalar bath in every internal component, while allowing an arbitrary positive definite port metric \(N\).

\begin{corollary}[Exact upper-glancing transversality]
\label{cor:v57-reciprocal-glancing}
Let \(|k|=r>0\), \(0<v<1\), and approach the upper gravitational glancing point \(s=ir+0\) on the physical branch.  Then
\[
 Z_{\rm rec}(ir,k)
   =\frac{ir}{\sqrt{1-v^2}}\,N
\]
and
\[
 \det\!\bigl(L_g(ir,k)+C Z_{\rm rec}(ir,k)J\bigr)
  =\frac{i r^5\det N}{24(1-v^2)^{3/2}}\ne0 .
\]
\end{corollary}

\begin{proof}
At \(s=ir+0\),
\[
 \kappa_v=\sqrt{-r^2+v^2r^2+0}
          =ir\sqrt{1-v^2}
\]
on the branch continued from \(\Re s>0\).  Hence
\[
 \frac{s^2}{\kappa_v}
 =\frac{-r^2}{ir\sqrt{1-v^2}}
 =\frac{ir}{\sqrt{1-v^2}}.
\]
The exact upper-glancing determinant identity proved in the V55 note, applied with
\(Z=ir(1-v^2)^{-1/2}N\), gives
\[
 \det(L_g+CZJ)
  =\frac{i r^5}{24}
     \det\!\left((1-v^2)^{-1/2}N\right),
\]
which is the displayed expression.  Positivity of \(N\) makes its determinant
strictly positive.
\end{proof}

This corollary is useful because it prevents a misleading diagnosis: the local
reciprocal transformer can remove the upper-glancing rank defect exactly.  Its
failure is instead forced below the first bath threshold by the preceding
low-cone theorem.

\subsection{An exact rational certificate}

Take the completely isotropic choice
\[
 |k|=1,\qquad v=\frac12,\qquad
 G=M=I_3,\qquad \Gamma_z=\Gamma_y=0 .
\]
For \(s=i\omega+0\), \(0<\omega<1/2\), put
\[
 x:=\omega^2,\qquad
 q(x):=\frac{x}{\sqrt{\frac14-x}}.
\]
Then \(Z_{\rm tr}^{\rm low}=-qI_3\).  The response eigenvalues from
\eqref{eq:V57-R} give
\[
 \frac{\det(L_0+CZ_{\rm tr}J)}{\det L_0}
   =(1-q\,d)(1-q\,b)^2,
 \qquad
 d=\frac{3-4x}{8(1-x)^{3/2}},\quad
 b=\frac{1}{2\sqrt{1-x}}.
\]
The transverse equation \(qb=1\) is
\[
 x=2\sqrt{\left(\frac14-x\right)(1-x)}.
\]
Both sides are positive for \(0<x<1/4\).  Squaring and simplifying gives
\[
 3x^2-5x+1=0,
\]
whose only root in that interval is
\[
 x_T=\frac{5-\sqrt{13}}6,\qquad
 \omega_T=\sqrt{\frac{5-\sqrt{13}}6}
          \simeq0.48208725429739563 .
\]
Consequently the determinant has an exact double transverse zero at
\(s=i\omega_T+0\).  It is not simultaneously longitudinal, since
\[
 1-q(x_T)d(x_T)
   =1-\frac{d(x_T)}{b(x_T)}
   =\frac{1}{4(1-x_T)}>0.
\]
The internal block is still positive definite there because
\(\sqrt{1/4-x_T}>0\).  Both gravitational normal roots are also strictly
positive because
\[
 0<\omega_T<\frac12<\frac{\sqrt3}{2}<1.
\]
At the unit upper-glancing point, Corollary~\ref{cor:v57-reciprocal-glancing} gives
the nonzero value
\[
 \det(L_g+C Z_{\rm tr}J)
   =\frac{i}{9\sqrt3}.
\]
Thus the same explicit local completion passes the exact upper-glancing test
and fails the low-cone test.

\subsection{Scope firewall and the upstream turn}

The result in this note closes the following class at principal order:
a real constant quadratic canonical boundary transformer on the six
coordinates \((z,y)\), containing one time derivative, coupled through an
invertible three-port matrix \(G\) to one positive common-speed scalar bath
with \(M\succ0\) and \(0<v<\sqrt3/2\).  Arbitrary constant gyroscopic
self-couplings \(\Gamma_z,\Gamma_y\) are included.  Every member of this class
has a boundary zero below its first internal threshold.

Several genuinely different possibilities remain outside the theorem:

\begin{enumerate}
\item a rank-deficient cross block \(G\), combined with metric-side
gyroscopic terms so that the total boundary correction can nevertheless
have the required rank;
\item several bath speeds or a matrix-valued noncommuting normal operator,
rather than one common square root;
\item an enlarged Hamiltonian state with independent momenta and a positive
evolution time matrix;
\item normal-derivative trace or traction couplings that are not equivalent
to the canonical coordinate time form analysed here;
\item restriction to the actual half-space physical quotient before the
three-channel determinant is imposed.
\end{enumerate}

The fifth item is now logically prior.  The V52--V57 sequence has repeatedly
shown that demanding a passive completion on the entire three-dimensional
metric port produces an obstruction, even when the completion repairs the
glancing rank.  Continuing to vary another full three-port mediator would
therefore risk circular exploration.  The next note will go upstream and
construct the boundary-adapted Hodge--York projector from its elliptic
half-space problem.  The concrete tasks are:

\begin{enumerate}
\item define the projector and its Sobolev domain with the boundary
conditions made explicit;
\item prove boundedness, idempotence, and compatibility with the essential
plane and harmonic subsidiary constraints;
\item compute the rank of the physical incoming boundary space;
\item only after that rank is fixed, pull the Einstein response and the
matter traction to the physical quotient and test the resulting reduced
passive determinant.
\end{enumerate}

If the physical count is two, the universal three-port zero proved here is
not the correct final obstruction; if the count remains three, the theorem
identifies a robust reason to enlarge the state or the bath class.

\subsection{V57 result ledger}

The new deductions are:

\begin{enumerate}
\item modulo time corners, every real constant quadratic one-time-derivative
action on \((z,y)\) is governed by one skew canonical matrix;
\item its regular three-port form is parametrised by
\((\Gamma_z,G,\Gamma_y)\), and it obeys the exact lossless power identity;
\item connection to the positive V45 bath yields the analytic homogeneous
linear-fractional impedance
\[
 s\Gamma_z+s^2G(\kappa_vM-s\Gamma_y)^{-1}G^\top;
\]
\item the full augmented determinant and its Schur complement agree
exactly while the internal block is invertible in \(\Re s>0\);
\item for every regular member with \(v<\sqrt3/2\), continuity and an
eigenvalue blow-up force a low-cone boundary zero before the first internal
pole;
\item the reciprocal endpoint has nonzero exact upper-glancing determinant,
and the isotropic \(v=1/2\) example has the closed-form double zero
\(\omega_T^2=(5-\sqrt{13})/6\).
\end{enumerate}

These statements do not construct the full Standard Model--Einstein
continuum.  They remove one precisely stated principal local completion
class and determine the upstream boundary quotient calculation that must
precede the next mediator design.

\subsection{Append-only record}

This V57 material was appended to the exact V56 parent whose SHA--256 digest
is
\[
\texttt{53e412d3e6ac08b4b440d1c972d8ddca4e3ac86e9fceff4eea50961e4b0d66dc}.
\]
The next research note is V58.  The next cross-program inspection of the
newest completed Yang--Mills and Navier--Stokes notes is due at V60.

% END APPENDED VERSION V57 MATERIAL

% BEGIN APPENDED VERSION V58 MATERIAL

\clearpage
\section{V58: a boundary Dirichlet--York projector, the rank-two
radiative glancing quotient, and the harmonic-incidence obstruction}
\label{sec:V58-boundary-York}

\subsection{Purpose and correction of the boundary count}
\label{subsec:V58-purpose}

V57 classified the complete constant canonical transformer between the
three V41 metric ports and one positive common-speed bath.  It proved
that every regular member develops a low-cone boundary zero before its
first internal pole.  That theorem is exact for the five-row graph to
which the transformer was attached.  It did not decide whether all
three graph directions survive the harmonic and residual-diffeomorphism
quotient.

This note goes upstream, as required by the V57 handoff.  It first
constructs the Dirichlet realization of the spatial York projector on a
half-space, including its weak domain, coercivity, Sobolev bounds, and
boundary trace.  A trace-reversed gauge lift preserves the full
conformal essential plane.  The lift commutes with the scalar wave
evolution only on an explicit second-order compatibility domain; outside
that domain its harmonic-subsidiary defect is printed rather than
discarded.

The stable glancing kernel can then be classified exactly.  Its
three dimensions split into one residual-gauge line and two independent
TT classes.  In the V41 port coordinates the gauge line is the
tangential-longitudinal port, while the transverse-tangential and
normal-normal ports represent the two radiative classes.  Thus the
radiative glancing defect has rank two, rather than the rank three used
before the V53 radiative quotient.

There is a second, independent issue.  The V41 incidence \(C\) maps
every nonzero three-port traction into a nonzero boundary value of the
harmonic constraint.  In fact the induced constraint-source map has
rank three.  Consequently the V57 low-cone zeros do not lie in the
homogeneous harmonic subsidiary domain.  They remain valid zeros of the
declared coupled graph, but they cannot be promoted to physical
Einstein instabilities until a boundary Bianchi/constraint compensator
is constructed.  Merely deleting the gauge port would leave this
defect unchanged.

\subsection{The Dirichlet York problem on the spatial half-space}
\label{subsec:V58-Dirichlet-York}

Let
\[
 \Omega:=\mathbb R^3_+
 =\{x=(x_1,x_2,x_3):x_3>0\},
 \qquad n=e_3,
\]
and let \(S^2_0\) denote the real trace-free symmetric \(3\times3\)
tensors.  For a vector field \(X\), define the conformal Killing
operator
\begin{equation}
 (\mathbb L X)_{ij}
 :=\partial_iX_j+\partial_jX_i
   -\frac23\delta_{ij}\partial^kX_k .
 \label{eq:V58-conformal-Killing}
\end{equation}
Its values lie in \(S^2_0\).  Use
\(\dot H^1_0(\Omega;\mathbb R^3)\), the completion of
\(C_c^\infty(\Omega;\mathbb R^3)\) in the gradient norm.  This choice
fixes the normal realization: the vector trace is zero on
\(\partial\Omega\), and no free boundary conformal Killing field is
included.

For smooth trace-free \(h\), integration by parts against a zero-trace
vector gives
\begin{equation}
 \mathbb L^*h=-2\partial^ih_{ij}.
 \label{eq:V58-L-adjoint}
\end{equation}
The Dirichlet vector Laplacian associated with the York complex is
\begin{equation}
 {\cal A}_D:=\mathbb L^*\mathbb L
 =-2\left(\Delta I_3+\frac13\nabla\operatorname{div}\right),
 \qquad X|_{\partial\Omega}=0.
 \label{eq:V58-vector-Laplacian}
\end{equation}

\begin{lemma}[Exact Dirichlet York coercivity]
\label{lem:V58-York-coercivity}
For every \(X\in\dot H^1_0(\Omega;\mathbb R^3)\),
\begin{equation}
 \|\mathbb L X\|_{L^2}^2
 =2\|\nabla X\|_{L^2}^2
  +\frac23\|\operatorname{div}X\|_{L^2}^2.
 \label{eq:V58-York-coercivity}
\end{equation}
In particular,
\begin{equation}
 \|\nabla X\|_{L^2}
 \le\frac1{\sqrt2}\|\mathbb L X\|_{L^2},
 \label{eq:V58-York-Korn}
\end{equation}
and \({\cal A}_D\) has no nonzero finite-gradient Dirichlet kernel.
\end{lemma}

\begin{proof}
For a compactly supported smooth vector field, expand
\(\mathbb L X:\mathbb L X\).  The symmetric-gradient part gives
\[
 \int_\Omega
 |\partial_iX_j+\partial_jX_i|^2
 =
 2\|\nabla X\|_2^2
 +2\int_\Omega\partial_iX_j\,\partial_jX_i.
\]
Integration by parts, with the zero trace eliminating the boundary
term, gives
\[
 \int_\Omega\partial_iX_j\,\partial_jX_i
 =\|\operatorname{div}X\|_2^2.
\]
Subtracting the trace in
\eqref{eq:V58-conformal-Killing} changes the coefficient of the last
term from \(2\) to \(2/3\).  This proves
\eqref{eq:V58-York-coercivity} on the dense smooth class and hence on
\(\dot H^1_0\).  Inequality \eqref{eq:V58-York-Korn} and the kernel
claim follow immediately.
\end{proof}

For \(h\in L^2(\Omega;S^2_0)\), let \(X_h\in\dot H^1_0\) be the weak
solution of
\begin{equation}
 \int_\Omega\mathbb L X_h:\mathbb L Y
 =
 \int_\Omega h:\mathbb L Y
 \qquad
 \text{for every }Y\in\dot H^1_0.
 \label{eq:V58-York-weak}
\end{equation}
Define
\begin{equation}
 \Pi_Y^D h:=h-\mathbb L X_h .
 \label{eq:V58-York-projector}
\end{equation}

\begin{theorem}[Half-space Dirichlet--York projection]
\label{thm:V58-York-projector}
Equation \eqref{eq:V58-York-weak} has a unique solution and
\(\Pi_Y^D\) is the \(L^2\)-orthogonal projector
\begin{equation}
 L^2(\Omega;S^2_0)
 =
 \ker\mathbb L^*
 \mathbin{\widehat\oplus}
 \mathbb L\dot H^1_0(\Omega;\mathbb R^3).
 \label{eq:V58-York-decomposition}
\end{equation}
More explicitly,
\begin{align}
 (\Pi_Y^D)^2&=\Pi_Y^D,
 &
 (\Pi_Y^D)^*&=\Pi_Y^D,
 \label{eq:V58-York-idempotent}\\
 \mathbb L^*\Pi_Y^Dh&=0,
 &
 \|\Pi_Y^Dh\|_2&\le\|h\|_2,
 \label{eq:V58-York-range-bound}\\
 \|\mathbb L X_h\|_2&\le\|h\|_2,
 &
 \|\nabla X_h\|_2&\le2^{-1/2}\|h\|_2.
 \label{eq:V58-York-solution-bound}
\end{align}
The range of \(\mathbb L\) is closed, so no closure is hidden in
\eqref{eq:V58-York-decomposition}.
\end{theorem}

\begin{proof}
The form on the left of \eqref{eq:V58-York-weak} is continuous and,
by Lemma~\ref{lem:V58-York-coercivity}, coercive in the
\(\dot H^1_0\) norm.  The right side is bounded by
\(\|h\|_2\|\mathbb L Y\|_2\).  Lax--Milgram gives a unique \(X_h\).

Equation \eqref{eq:V58-York-weak} says precisely that
\[
 \langle\Pi_Y^Dh,\mathbb L Y\rangle_{L^2}=0
\]
for every \(Y\).  Hence \(\Pi_Y^Dh\in\ker\mathbb L^*\).  Taking
\(Y=X_h\) gives
\[
 \|\mathbb L X_h\|_2^2
 =\langle h,\mathbb L X_h\rangle,
\]
and therefore the first estimate in
\eqref{eq:V58-York-solution-bound}; the second follows from
\eqref{eq:V58-York-Korn}.  The two summands in
\eqref{eq:V58-York-projector} are orthogonal, so Pythagoras gives
\(\|\Pi_Y^Dh\|_2\le\|h\|_2\).

If \(h\in\ker\mathbb L^*\), the weak right side is zero and uniqueness
gives \(X_h=0\).  If \(h=\mathbb L X\) with
\(X\in\dot H^1_0\), uniqueness gives \(X_h=X\).  These two observations
prove idempotence, the range and kernel identities, and
self-adjointness.

Finally, suppose \(\mathbb L X_m\) is Cauchy in \(L^2\).  The coercive
estimate makes \(X_m\) Cauchy in \(\dot H^1_0\), with limit \(X\), and
continuity gives \(\mathbb L X_m\to\mathbb L X\).  Thus the range is
closed.
\end{proof}

\begin{proposition}[Homogeneous Sobolev boundedness]
\label{prop:V58-York-Sobolev}
For every integer \(m\ge0\), the projector extends to the homogeneous
Sobolev realization and obeys
\begin{equation}
 \|\Pi_Y^Dh\|_{\dot H^m}
 +\|\mathbb L X_h\|_{\dot H^m}
 +\|\nabla X_h\|_{\dot H^m}
 \le C_m\|h\|_{\dot H^m}
 \label{eq:V58-York-Sobolev}
\end{equation}
for smooth data in the Dirichlet compatibility class, and hence by
completion on that class.  At \(m=0\), the preceding theorem gives
\(C_0\le1+1+2^{-1/2}\).
\end{proposition}

\begin{proof}
Tangential difference quotients commute with
\({\cal A}_D\) and preserve the zero trace.  Applying
\eqref{eq:V58-York-coercivity} to each tangential derivative gives the
estimate for all derivatives containing at most one unresolved normal
derivative.  The coefficient of \(\partial_3^2X\) in
\[
 \Delta X+\frac13\nabla\operatorname{div}X
\]
is \(\operatorname{diag}(1,1,4/3)\), which is invertible.  Equation
\eqref{eq:V58-vector-Laplacian} therefore solves successively for every
remaining pair of normal derivatives in terms of tangential derivatives
and derivatives of \(\mathbb L^*h\).  Induction gives
\[
 \|\nabla^{m+1}X_h\|_2\le C_m\|\nabla^mh\|_2.
\]
Substitution in \eqref{eq:V58-York-projector} proves
\eqref{eq:V58-York-Sobolev}.  Density of smooth compatible data gives
the extension.  This proof uses the homogeneous scale, so it does not
silently invoke a false Poincare inequality on the unbounded half-space.
\end{proof}

\begin{corollary}[Preservation of the spatial conformal boundary plane]
\label{cor:V58-spatial-boundary-preservation}
Let \(Q_{\partial,2}^{\rm TF}\) take the trace-free part of the
tangential \(2\times2\) block \(h_{AB}\), \(A,B=1,2\).  Whenever the
traces are defined,
\begin{equation}
 Q_{\partial,2}^{\rm TF}\Pi_Y^Dh
 =
 Q_{\partial,2}^{\rm TF}h.
 \label{eq:V58-spatial-boundary-preservation}
\end{equation}
\end{corollary}

\begin{proof}
The Dirichlet trace of \(X_h\) is zero, so all of its tangential
derivatives have zero trace.  Hence
\[
 (\mathbb L X_h)_{AB}\big|_{\partial\Omega}
 =-\frac23\delta_{AB}\partial_3(X_h)_3.
\]
Its tangential trace-free part is zero.  Apply
\eqref{eq:V58-York-projector}.
\end{proof}

This is the half-space normal realization missing from the whole-space
multiplier in V54.  It is a genuine bounded projector, but its spatial
statement alone is not yet the full Lorentzian conformal row.

\subsection{The trace-reversed lift and its exact commutator}
\label{subsec:V58-gauge-lift}

Let \(\bar h_{\mu\nu}\) be the trace-reversed metric perturbation and put
\[
 h_0^{\rm sp}
 :=\bar h_{ij}-\frac13\delta_{ij}\bar h^k{}_k .
\]
Solve \eqref{eq:V58-York-weak} for \(X=X_{h_0^{\rm sp}}\), extend it
to a spacetime covector by
\[
 \widetilde X_0=0,\qquad \widetilde X_i=X_i,
\]
and use the trace-reversed gauge operator already fixed in V53,
\begin{equation}
 (\Gamma\widetilde X)_{\mu\nu}
 :=
 \partial_\mu\widetilde X_\nu
 +\partial_\nu\widetilde X_\mu
 -\eta_{\mu\nu}\partial^\rho\widetilde X_\rho.
 \label{eq:V58-Gamma}
\end{equation}
Define the lifted York representative
\begin{equation}
 {\boldsymbol\Pi}_Y^D\bar h
 :=\bar h-\Gamma\widetilde X.
 \label{eq:V58-lifted-projector}
\end{equation}

The spatial trace-free part of \(\Gamma\widetilde X\) is exactly
\(\mathbb L X\).  Consequently
\begin{equation}
 \left({\boldsymbol\Pi}_Y^D\bar h\right)_{\rm sp}^{\rm TF}
 =\Pi_Y^D h_0^{\rm sp}.
 \label{eq:V58-lift-spatial-part}
\end{equation}

\begin{theorem}[Boundary preservation and idempotence of the lifted
representative]
\label{thm:V58-lift-boundary}
On smooth fields for which the displayed traces exist,
\begin{equation}
 ({\boldsymbol\Pi}_Y^D)^2={\boldsymbol\Pi}_Y^D.
 \label{eq:V58-lift-idempotent}
\end{equation}
Moreover, if \(Q_{\rm TF}\) is the full Lorentzian induced
trace-free row of V52--V54, then
\begin{equation}
 Q_{\rm TF}{\boldsymbol\Pi}_Y^D\bar h
 =Q_{\rm TF}\bar h.
 \label{eq:V58-full-boundary-preservation}
\end{equation}
Thus the lifted projector preserves the essential boundary plane
\(Q_{\rm TF}\bar h=0\).
\end{theorem}

\begin{proof}
Equation \eqref{eq:V58-lift-spatial-part} and
Theorem~\ref{thm:V58-York-projector} show that the spatial trace-free
part after one application lies in \(\ker\mathbb L^*\).  Its York
solution is therefore zero.  A second application makes no change,
which proves \eqref{eq:V58-lift-idempotent}.

Let \(a,b\in\{0,1,2\}\) be tangential spacetime indices.  The Dirichlet
condition holds for \(X\) at every time, so the boundary traces of
\(\partial_a\widetilde X_b\) vanish.  At the planar boundary,
\begin{equation}
 (\Gamma\widetilde X)_{ab}\big|_{\partial\Omega}
 =-\gamma_{ab}\partial_3X_3 .
 \label{eq:V58-Gamma-boundary-trace}
\end{equation}
This is a pure induced-trace tensor.  Its \(Q_{\rm TF}\) image is zero,
proving \eqref{eq:V58-full-boundary-preservation}.  Trace reversal adds
only a multiple of the induced metric and therefore does not change
this conclusion.
\end{proof}

The spatial projector is bounded on every homogeneous Sobolev level by
Proposition~\ref{prop:V58-York-Sobolev}.  The full lift contains
\(\partial_tX\).  For \(m\ge1\), applying that proposition to
\(\partial_t h_0^{\rm sp}\) gives the graph estimate
\begin{equation}
 \|\Gamma\widetilde X\|_{\dot H_x^m}
 \le C_m\left(
 \|h_0^{\rm sp}\|_{\dot H_x^m}
 +\|\partial_t h_0^{\rm sp}\|_{\dot H_x^{m-1}}
 \right).
 \label{eq:V58-lift-Sobolev}
\end{equation}
The homogeneous formulation is essential at zero spatial frequency:
there is no claimed \(L^2_x\) bound for the undifferentiated
\(\partial_tX\) from an \(L^2_x\) configuration alone.

Let
\begin{equation}
 {\cal C}_\nu(\bar h):=\partial^\mu\bar h_{\mu\nu}
 \label{eq:V58-harmonic-covector}
\end{equation}
be the harmonic covector.  The gauge-complex identity is
\begin{equation}
 {\cal C}(\Gamma\widetilde X)=\Box\widetilde X.
 \label{eq:V58-complex-identity}
\end{equation}
It follows without any domain assumption that
\begin{equation}
 {\cal C}({\boldsymbol\Pi}_Y^D\bar h)
 =
 {\cal C}(\bar h)-\Box\widetilde X.
 \label{eq:V58-harmonic-defect}
\end{equation}

\begin{definition}[Wave-compatible York domain]
\label{def:V58-compatible-domain}
For \(m\ge2\), let \({\mathscr D}_Y^m\) consist of wave fields
\(\bar h\) with the regularity needed in
\eqref{eq:V58-lift-Sobolev}, whose York vector \(X\) obeys
\begin{equation}
 X|_{\partial\Omega}=0,\qquad
 (\Box X)|_{\partial\Omega}=0,\qquad
 \Box X\in\dot H^1_0(\Omega;\mathbb R^3).
 \label{eq:V58-second-compatibility}
\end{equation}
The second condition is the ordinary second Dirichlet compatibility
condition for the vector wave.  It is a domain condition, rather than a
consequence of the first trace alone.
\end{definition}

\begin{theorem}[Exact wave and subsidiary preservation]
\label{thm:V58-wave-commutation}
If \(\bar h\in{\mathscr D}_Y^m\) satisfies the componentwise wave
equation
\begin{equation}
 \Box\bar h=0,
 \label{eq:V58-metric-wave}
\end{equation}
then
\begin{equation}
 \Box X=0,\qquad
 \Box{\boldsymbol\Pi}_Y^D\bar h=0,\qquad
 {\cal C}({\boldsymbol\Pi}_Y^D\bar h)={\cal C}(\bar h).
 \label{eq:V58-wave-preservation}
\end{equation}
In particular, the lifted York representative preserves the harmonic
subsidiary shell on \({\mathscr D}_Y^m\).

Without \eqref{eq:V58-second-compatibility}, the exact defect is the
nonzero vector \(W:=\Box X\), which satisfies
\begin{equation}
 {\cal A}_D W=0
 \quad\hbox{in }\Omega,
 \qquad
 W|_{\partial\Omega}=(\Box X)|_{\partial\Omega}.
 \label{eq:V58-Poisson-defect}
\end{equation}
Thus failure of commutation is an explicit homogeneous vector-Poisson
boundary lift.
\end{theorem}

\begin{proof}
All operators have constant coefficients in the interior.  From
\({\cal A}_DX=\mathbb L^*h_0^{\rm sp}\) and
\eqref{eq:V58-metric-wave},
\[
 {\cal A}_D(\Box X)
 =\mathbb L^*(\Box h_0^{\rm sp})=0.
\]
On \({\mathscr D}_Y^m\), \(W=\Box X\) has zero Dirichlet trace and lies
in \(\dot H^1_0\).  Testing \({\cal A}_DW=0\) against \(W\) and using
Lemma~\ref{lem:V58-York-coercivity} gives
\(\mathbb L W=0\), hence \(\nabla W=0\).  Its zero trace and
finite-gradient realization give \(W=0\).

The gauge operator commutes with \(\Box\), so
\[
 \Box{\boldsymbol\Pi}_Y^D\bar h
 =\Box\bar h-\Gamma(\Box\widetilde X)=0.
\]
Equation \eqref{eq:V58-complex-identity} then proves subsidiary
preservation.  If the second boundary trace is not zero, the same
interior calculation still gives
\({\cal A}_DW=0\), with the displayed boundary value.  This is
\eqref{eq:V58-Poisson-defect}.
\end{proof}

\begin{firewall}[What the half-space projector does]
\label{fire:V58-projector-scope}
The operator \(\Pi_Y^D\) is a genuine orthogonal half-space projector,
and its lifted representative preserves the V52 essential plane.
The wave-commutation theorem holds on the printed compatibility domain;
it does not assert that the old V41 port graph automatically preserves
that domain.  The lift fixes the spatial Dirichlet gauge directions
with zero boundary trace.  It does not by itself quotient nonzero
boundary edge transformations, the temporal gauge coordinate, or the
off-shell ghost--antighost complex.
\end{firewall}

\subsection{The radiative quotient of the glancing kernel}
\label{subsec:V58-glancing-quotient}

Rotate the tangential covector to \(k=(r,0)\), take the plus glancing
sheet \(s=ir\), \(\lambda=0\), and use the V55 right-kernel frame.
Write its three columns as
\begin{align}
 g_\parallel&:=(0,1,1,0,0)^T,
 \label{eq:V58-gparallel}\\
 r_\times&:=(0,0,0,1,0)^T,
 \label{eq:V58-rcross}\\
 r_+&:=(-3,0,0,0,1)^T.
 \label{eq:V58-rplus}
\end{align}
They form a basis of \(\ker L_{\rm g}\).

At this glancing point the spatial propagation direction is \(e_1\).
Use the TT basis on its transverse plane
\begin{equation}
 E_+:=\frac{e_2\otimes e_2-e_3\otimes e_3}{\sqrt2},
 \qquad
 E_\times:=\frac{e_2\otimes e_3+e_3\otimes e_2}{\sqrt2}.
 \label{eq:V58-glancing-TT-basis}
\end{equation}
Let \(T_{\rm g}\) be the V53 spatial TT quotient map, restricted to
\(\ker L_{\rm g}\).  Adding a trace-reversal multiple of
\(\delta_{ij}\) does not affect \(T_{\rm g}\).

\begin{theorem}[Exact rank-two radiative glancing quotient]
\label{thm:V58-rank-two-glancing}
The restricted TT map satisfies
\begin{align}
 T_{\rm g}g_\parallel&=0,
 \label{eq:V58-T-gauge}\\
 T_{\rm g}r_\times&=\sqrt2\,E_\times,
 \label{eq:V58-T-cross}\\
 T_{\rm g}r_+&=-\sqrt2\,E_+.
 \label{eq:V58-T-plus}
\end{align}
Consequently
\begin{equation}
 0\longrightarrow\operatorname{span}\{g_\parallel\}
 \longrightarrow\ker L_{\rm g}
 \mathop{\longrightarrow}^{T_{\rm g}}
 \operatorname{span}\{E_+,E_\times\}
 \longrightarrow0
 \label{eq:V58-glancing-exact-sequence}
\end{equation}
is split exact.  In particular, the residual-gauge quotient of the
three-dimensional glancing defect has dimension two.

The same conclusion holds on the opposite temporal sheet and after
tangential rotation.
\end{theorem}

\begin{proof}
The reduced variables mean
\[
 h_{ab}=\frac{\tau}{3}\eta_{ab},\qquad
 h_{3a}=v_a,\qquad h_{33}=\phi.
\]
The spatial block of \(g_\parallel\) has only the longitudinal entry
\(h_{13}=h_{31}=1\), which the TT projector for direction \(e_1\)
annihilates.  The spatial block of \(r_\times\) has
\(h_{23}=h_{32}=1\), giving the Frobenius coefficient
\(\sqrt2\) on \(E_\times\).  The spatial block of \(r_+\) is
\[
 \operatorname{diag}(-1,-1,1).
\]
Its transverse \(e_2,e_3\) block has coefficient \(-\sqrt2\) on
\(E_+\).  These calculations prove
\eqref{eq:V58-T-gauge}--\eqref{eq:V58-T-plus}.

The map is therefore onto and its kernel on the three-dimensional
domain is exactly \(\operatorname{span}\{g_\parallel\}\).  A splitting
sends \(E_\times\) to \(r_\times/\sqrt2\) and \(E_+\) to
\(-r_+/\sqrt2\).  This proves exactness.  The V53 radiative exact
sequence identifies a harmonic null amplitude with zero TT image as a
residual wave diffeomorphism, so the kernel line has the stated gauge
meaning.  Complex conjugation and tangential orthogonal covariance give
the other cases.
\end{proof}

\begin{corollary}[The physical glancing port is two-dimensional]
\label{cor:V58-physical-port}
The V41 trace \(Jx=(v_1,v_2,\phi)^T\) obeys
\begin{equation}
 Jg_\parallel=e_1,\qquad
 Jr_\times=e_2,\qquad
 Jr_+=e_3.
 \label{eq:V58-J-kernel-basis}
\end{equation}
After the radiative quotient, the physical glancing port is therefore
\begin{equation}
 {\cal Z}_{\rm phys,g}
 =\operatorname{span}\{e_2,e_3\},
 \qquad
 P_{\rm phys,g}=\operatorname{diag}(0,1,1)
 \label{eq:V58-physical-port-projector}
\end{equation}
in the aligned frame.  For a general nonzero tangential spatial
covector \(k\in\mathbb R^2\), this becomes
\begin{equation}
 P_{\rm phys,g}(k)
 =
 \begin{pmatrix}
 I_2-\widehat k\widehat k^T&0\\
 0&1
 \end{pmatrix},
 \qquad
 \operatorname{rank}P_{\rm phys,g}=2.
 \label{eq:V58-rotated-port-projector}
\end{equation}
\end{corollary}

\begin{proof}
The three identities in \eqref{eq:V58-J-kernel-basis} are the exact
identity \(JR_{\rm g}=I_3\) from V55, with the columns named in
\eqref{eq:V58-gparallel}--\eqref{eq:V58-rplus}.  Theorem
\ref{thm:V58-rank-two-glancing} removes the first column and retains the
other two.  Tangential rotation changes the first two port coordinates
into longitudinal and transverse components and leaves the
normal-normal coordinate fixed, giving
\eqref{eq:V58-rotated-port-projector}.
\end{proof}

\begin{corollary}[Revised regular channel lower bound]
\label{cor:V58-two-channel-lower-bound}
A regularly eliminable auxiliary system whose correction is required
to be injective on the radiative glancing defect needs at least two
independent amplitudes.  A rank-two auxiliary can act injectively on
that kernel at a fixed glancing frequency.  The earlier three-channel lower bound of V39
continues to hold on the five-dimensional embedding quotient before
the residual-radiative quotient is taken.
\end{corollary}

\begin{proof}
Theorem~\ref{thm:V58-rank-two-glancing} gives a two-dimensional
radiative kernel.  A correction factored through an \(r\)-dimensional
regular auxiliary space has rank at most \(r\) on that kernel.  It
cannot be injective when \(r<2\).  Conversely, choose an isomorphism
from the two-dimensional kernel to \(\mathbb C^2\) and an injection
of \(\mathbb C^2\) into the boundary target.  Their composition acts
injectively on the kernel.  Complete invertibility still requires the
physical target quotient.  V39 used the three-dimensional kernel
before the additional exact sequence
\eqref{eq:V58-glancing-exact-sequence}; its statement is therefore not
contradicted.
\end{proof}

The corollary is a channel count, not a two-channel continuum theorem.
It neither chooses an analytic complement away from glancing nor proves
a positive boundary energy.

\subsection{The harmonic source carried by the old port incidence}
\label{subsec:V58-constraint-incidence}

Let \(p=D_nx\), retain the V41 graph \(p=A(D)x\), and write a general
graph residual in the coordinate order
\((\tau,v_0,v_1,v_2,\phi)\) as
\begin{equation}
 f:=p-A(D)x=(f_\tau,f_{v_0},f_{v_1},f_{v_2},f_\phi)^T.
 \label{eq:V58-graph-residual}
\end{equation}
V41 proved that these five rows are, without a singular row operation,
\begin{equation}
 f=\bigl(2K,\ C_0,\ C_1,\ C_2,\ 2(K+C_n)\bigr)^T.
 \label{eq:V58-row-identification}
\end{equation}
It follows that the harmonic boundary source contained in \(f\) is
\begin{equation}
 {\mathsf S}_{\cal C}f
 :=
 \begin{pmatrix}
 f_{v_0}\\ f_{v_1}\\ f_{v_2}\\
 (f_\phi-f_\tau)/2
 \end{pmatrix}
 =
 \begin{pmatrix}C_0\\C_1\\C_2\\C_n\end{pmatrix}.
 \label{eq:V58-constraint-source-map}
\end{equation}

\begin{lemma}[Constraint-preserving graph forces form one line]
\label{lem:V58-constraint-force-line}
The map \({\mathsf S}_{\cal C}:\mathbb C^5\to\mathbb C^4\) has rank
four and
\begin{equation}
 \ker{\mathsf S}_{\cal C}
 =
 \operatorname{span}\{(1,0,0,0,1)^T\}.
 \label{eq:V58-constraint-force-kernel}
\end{equation}
Thus a graph force preserves the homogeneous harmonic rows precisely
when it changes \(K\) and \(K+C_n\) by the same amount and changes none
of the \(C_a\) rows.
\end{lemma}

\begin{proof}
The four displayed coordinates in
\eqref{eq:V58-constraint-source-map} are independent.  Their common
kernel is
\[
 f_{v_0}=f_{v_1}=f_{v_2}=0,\qquad f_\phi=f_\tau,
\]
which is exactly \eqref{eq:V58-constraint-force-kernel}.
\end{proof}

The explicit V41 reciprocal incidence is
\begin{equation}
 C=
 \begin{pmatrix}
 0&0&-3/2\\
 0&0&0\\
 1/2&0&0\\
 0&1/2&0\\
 0&0&1/2
 \end{pmatrix}.
 \label{eq:V58-C-matrix}
\end{equation}
Consequently
\begin{equation}
 {\mathsf S}_{\cal C}C
 =
 \begin{pmatrix}
 0&0&0\\
 1/2&0&0\\
 0&1/2&0\\
 0&0&1
 \end{pmatrix}.
 \label{eq:V58-SC-C}
\end{equation}

\begin{theorem}[Injective harmonic-incidence obstruction]
\label{thm:V58-harmonic-incidence}
The map \({\mathsf S}_{\cal C}C:\mathbb C^3\to\mathbb C^4\) has rank
three.  Hence, for either sign,
\begin{equation}
 p=A(D)x\pm Cr
 \label{eq:V58-forced-graph}
\end{equation}
preserves the homogeneous harmonic boundary rows only when \(r=0\).

Let \(Z(s,k)\) be any three-port impedance and suppose
\(\lambda\ne0\).  Every nonzero vector in
\begin{equation}
 \ker\bigl(L_0(s,k)+CZ(s,k)J\bigr)
 \label{eq:V58-old-kernel}
\end{equation}
has a nonzero induced harmonic boundary source.  In particular, all
low-cone zero modes supplied by Theorem
\ref{thm:V57-boundary-zero} lie outside the domain
\({\cal C}_\mu=0\).
\end{theorem}

\begin{proof}
Matrix \eqref{eq:V58-SC-C} is injective, proving the first assertion.

Let \(x\ne0\) lie in \eqref{eq:V58-old-kernel} and put
\(r=ZJx\).  If \(r=0\), then \(L_0x=0\).  For \(\lambda\ne0\), the V37 open-determinant theorem and the fixed
nonsingular row equivalence between its matrix and \(L_0\) make
\(L_0\) invertible.  This would force \(x=0\), a
contradiction.  Hence \(r\ne0\), and injectivity of
\({\mathsf S}_{\cal C}C\) gives a nonzero harmonic source.  The V57
low-cone theorem has \(\lambda=\sqrt{1-\omega_0^2}>0\), so it is
included.
\end{proof}

\begin{corollary}[The exact V57 witness is off the subsidiary shell]
\label{cor:V58-V57-witness}
For the V57 isotropic example
\[
 v=\frac12,\qquad G=M=I_3,\qquad
 \Gamma_z=\Gamma_y=0,
\]
let
\[
 x_T=\frac{5-\sqrt{13}}6,\qquad
 \omega_T=\sqrt{x_T},\qquad
 \lambda_T=\sqrt{1-x_T},\qquad
 q_T=\frac{x_T}{\sqrt{\frac14-x_T}}.
\]
Then \(q_T=2\lambda_T\), and for
\(e_{v_2}=(0,0,0,1,0)^T\),
\begin{equation}
 \left(L_0(i\omega_T,(1,0))-q_TCJ\right)e_{v_2}=0.
 \label{eq:V58-explicit-old-zero}
\end{equation}
Its harmonic boundary covector has
\begin{equation}
 C_2=\lambda_T\ne0.
 \label{eq:V58-explicit-C2}
\end{equation}
Thus even the cross-polarized configuration vector in the exact double
zero is not a homogeneous harmonic solution.
\end{corollary}

\begin{proof}
The polynomial \(3x_T^2-5x_T+1=0\) is equivalent to
\[
 x_T^2=4\left(\frac14-x_T\right)(1-x_T),
\]
and all factors are positive.  Taking the positive square root gives
\(q_T=2\lambda_T\).  For \(k=(1,0)\), the V41 operator has
\[
 L_0e_{v_2}=\lambda_Te_{v_2},\qquad
 CJe_{v_2}=\frac12e_{v_2},
\]
which proves \eqref{eq:V58-explicit-old-zero}.  Its graph residual is
\(q_TCJe_{v_2}=\lambda_Te_{v_2}\).
Equations \eqref{eq:V58-row-identification} and
\eqref{eq:V58-constraint-source-map} identify the \(v_2\)-row with
\(C_2\), proving \eqref{eq:V58-explicit-C2}.
\end{proof}

\begin{firewall}[Status of the V55--V57 determinant obstructions]
\label{fire:V58-old-determinant-scope}
The V55--V57 determinant and passivity theorems remain valid for their
printed five-row metric graphs.  Theorem
\ref{thm:V58-harmonic-incidence} changes their physical
interpretation: the old incidence replaces nonzero harmonic boundary
rows by bath traction.  No preceding note derived a bath stress or
constraint field whose Bianchi identity supplies exactly
\eqref{eq:V58-SC-C}.  Until such a compensator is constructed, neither
an old zero nor an old stable floor is a theorem on the homogeneous
physical harmonic quotient.
\end{firewall}

\subsection{What the quotient changes and what it does not}
\label{subsec:V58-consequences}

The two V58 conclusions must be applied in order.

First, Corollary~\ref{cor:V58-physical-port} removes the
tangential-longitudinal glancing port from the radiative quotient.  A
future physical correction should therefore be tested first on two
radiative characteristic channels, not required a priori to invert an
unphysical three-by-three port matrix.

Second, Theorem~\ref{thm:V58-harmonic-incidence} prevents the immediate
replacement
\[
 Z\longmapsto P_{\rm phys,g}ZP_{\rm phys,g}.
\]
The old incidence \(C\) sends both retained columns \(e_2,e_3\) to
nonzero harmonic boundary sources.  Projecting the impedance cannot
change the injective matrix \({\mathsf S}_{\cal C}C\).  A legitimate
two-channel construction must instead alter the boundary incidence or
supply a constraint field whose boundary source cancels it exactly.

There are two structurally distinct routes:

\begin{enumerate}
\item retain \({\cal C}_\mu=0\) as four independent subsidiary rows and
couple a two-component positive continuum through physical
characteristic or Brown--York shear variables whose incidence lies in
the quotient Green form; or
\item enlarge the boundary state by a constraint/ghost compensator
\(\chi_\mu\) and derive, from one action and one Noether identity, the
opposite source
\[
 {\cal C}(\chi)=-{\mathsf S}_{\cal C}Cr.
\]
The enlarged subsidiary system must then have its own closed maximal
domain and positive physical quotient.
\end{enumerate}

The first route is smaller and will be taken in V59.  The object to
construct is the first-order boundary characteristic quotient of the
V53 wave state:
\[
 U^\pm_{\mu\nu}:=\pi_{\mu\nu}\pm w_{3\mu\nu}.
\]
V59 must apply the boundary Dirichlet--York representative to the
constraint shell, identify the two TT incoming combinations and their
outgoing partners, and derive their normal energy flux.  It must then
attach any positive continuum only to this two-component flux pair
while keeping the four \({\cal C}_\mu\) rows homogeneous.  The decisive
test is the full augmented characteristic matrix; a two-by-two
impedance determinant alone will not establish domain preservation.

\subsection{Finite exact certificates}
\label{subsec:V58-certificates}

Let
\[
 R_{\rm g}:=(g_\parallel,r_\times,r_+).
\]
In the ordered TT coordinates \((E_+,E_\times)\), the quotient card is
\begin{equation}
 [T_{\rm g}]R_{\rm g}
 =
 \begin{pmatrix}
 0&0&-\sqrt2\\
 0&\sqrt2&0
 \end{pmatrix},
 \qquad
 \operatorname{rank}([T_{\rm g}]R_{\rm g})=2,
 \qquad
 JR_{\rm g}=I_3.
 \label{eq:V58-TT-card}
\end{equation}
The port-projector card is
\begin{equation}
 P_{\rm phys,g}^2=P_{\rm phys,g},\qquad
 P_{\rm phys,g}g_{\rm port}=0,\qquad
 \operatorname{rank}P_{\rm phys,g}=2,
 \quad g_{\rm port}=Jg_\parallel.
 \label{eq:V58-port-card}
\end{equation}
The constraint card is
\begin{equation}
 \operatorname{rank}{\mathsf S}_{\cal C}=4,\qquad
 \operatorname{rank}({\mathsf S}_{\cal C}C)=3,\qquad
 \ker({\mathsf S}_{\cal C}C)=\{0\}.
 \label{eq:V58-constraint-card}
\end{equation}
Finally, the V57 witness card stores
\begin{equation}
 q_T-2\lambda_T=0,\qquad
 (L_0-q_TCJ)e_{v_2}=0,\qquad
 {\mathsf S}_{\cal C}(q_TCJe_{v_2})
 =(0,0,\lambda_T,0)^T\ne0.
 \label{eq:V58-witness-card}
\end{equation}

For the York block, a finite Fourier card at a nonzero real covector
\(\xi\) stores the symbol \(\mathbb L(\xi)\) and verifies
\begin{equation}
 \mathbb L(\xi)^*\mathbb L(\xi)
 =2|\xi|^2I_3+\frac23\xi\xi^T,
 \qquad
 \lambda_{\min}
 \bigl(\mathbb L(\xi)^*\mathbb L(\xi)\bigr)
 =2|\xi|^2.
 \label{eq:V58-York-symbol-card}
\end{equation}
This is the pointwise counterpart of
Lemma~\ref{lem:V58-York-coercivity}; the half-space projector itself is
defined by the weak Dirichlet problem, not by importing the whole-space
multiplier.

\subsection{V58 result ledger}
\label{subsec:V58-ledger}

The new deductions are:

\begin{enumerate}
\item the spatial half-space York complex with zero vector trace has an
exact coercive identity, closed range, a norm-one \(L^2\) projector,
and homogeneous Sobolev bounds;
\item its trace-reversed gauge lift is idempotent and preserves the full
Lorentzian conformal essential plane;
\item the lift commutes with the componentwise metric wave equation and
preserves the harmonic subsidiary covector on the explicit
second-Dirichlet-compatibility domain, while its off-domain commutator
is the vector-Poisson lift \eqref{eq:V58-Poisson-defect};
\item the three-dimensional V37/V55 glancing kernel has one residual
gauge line and an exact two-dimensional TT quotient;
\item the physical glancing port is the transverse-tangential plus
normal-normal plane, so the regular physical channel lower bound is
two after the radiative quotient;
\item the old reciprocal incidence \(C\) injects every nonzero port
traction into the harmonic boundary rows; and
\item the V57 exact low-cone witness, and more generally every nonzero
low-cone kernel of the old graph, lies outside the homogeneous harmonic
subsidiary domain.
\end{enumerate}

These results do not construct the full Standard Model--Einstein
continuum.  They define the missing half-space York realization and
settle the glancing physical count, but they also expose a prior
constraint-incidence mismatch.  A positive two-channel characteristic
boundary generator, its complete domain, its stress/Bianchi identity,
the off-shell ghost complex, nonlinear Einstein closure, and the
coupled Standard-Model sectors remain open.

\subsection{Append-only record}

This V58 material was appended to the exact V57 parent whose SHA--256
digest is
\[
\texttt{27c2641254957d353d870394a9169845bf5895c87c46bc26945821596f6c1f3f}.
\]
The next research note is V59.  The next cross-program inspection of
the newest completed Yang--Mills and Navier--Stokes notes remains due
at V60.

% END APPENDED VERSION V58 MATERIAL

% BEGIN APPENDED VERSION V59 MATERIAL

\clearpage
\section{V59: a positive two-polarization radiative--continuum
generator, its zero-free interface symbol, and the remaining metric lift}
\label{sec:V59-radiative-generator}

\subsection{Purpose and nonduplication boundary}
\label{subsec:V59-purpose}

V58 constructed a boundary Dirichlet--York representative and proved two
facts which change the continuation.  The glancing defect has two
radiative TT classes after the residual wave-gauge line is removed, but
the old V41 incidence injects every nonzero port traction into the
harmonic boundary covector.  A two-by-two truncation of the old
three-port determinant would therefore repeat the constraint mismatch.

This note builds the physical quotient system first.  The two TT
amplitudes form a wave field on the Einstein half-space.  A
two-component positive continuum on a second half-line is attached by
one geometric trace equality and one conormal balance.  Both conditions
come from a single closed quadratic form.  The coupled form defines a
nonnegative self-adjoint spatial operator, and its abstract wave
generator is skew-adjoint on the homogeneous finite-energy space.  Thus
the construction has a positive conserved energy, a complete operator
domain, and a unitary time evolution on the radiative quotient.

The Laplace--Fourier interface matrix is
\[
 {\cal B}_{\rm rad}(s,k)
 =\bigl(\lambda+m\kappa_v\bigr)I_2.
\]
It is zero-free on the complete normalized closed stable cone and has a
strict uniform floor, including gravitational glancing and the
continuum threshold.  Its characteristic scattering matrix is
contractive in the open Laplace half-plane.

The result is a completed positive linear
Einstein-radiation--continuum generator \emph{on the two-polarization
quotient}.  It is not yet a generator for the full ten-component metric.
The inverse lift must impose four homogeneous subsidiary rows and a
gauge/embedding complement while reproducing the quotient conormal
balance from one covariant boundary action.  That lift, together with
the scheduled NS/Yang--Mills audit, is the V60 task.

\subsection{The two-polarization characteristic quotient}
\label{subsec:V59-characteristic-quotient}

Let \({\mathscr T}\to\mathbb R^2_k\setminus\{0\}\) denote the
rank-two Hermitian radiative bundle identified by the V53 null-shell
exact sequence and by the V58 glancing exact sequence.  On any
measurable orthonormal polarization chart, write the Fourier
representative of a section as
\[
 \widehat a=(a_+,a_\times)^T,
\]
and suppress the hat below.  Every \(y\)-integral and Sobolev space
below is Plancherel shorthand for the corresponding measurable direct
integral of the fibres \({\mathscr T}_k\); the real Einstein fields form
the invariant reality subspace.  In particular, tangential derivatives
act by the scalar multiplier \(ik\) and do not differentiate a moving
polarization frame.  All operators below are scalar multiples of
\(I_{\mathscr T}\), so the construction is invariant under measurable
unitary changes of polarization frame.  No global continuous choice of
\(E_+,E_\times\) is assumed.

On the wave-compatible York domain
\({\mathscr D}_Y^m\) of V58, the radiative quotient obeys two copies of
the scalar wave equation.  In the metric half-space
\[
 \Omega_x:=\{(y,x):y\in\mathbb R^2,\ x>0\},
\]
write
\begin{equation}
 a_{tt}=\partial_x^2a+\Delta_ya.
 \label{eq:V59-radiative-wave}
\end{equation}
The outward conormal is
\begin{equation}
 p:=D_na=-\partial_xa\big|_{x=0}.
 \label{eq:V59-radiative-conormal}
\end{equation}
For a smooth finite-energy solution,
\begin{equation}
 \frac{\dd}{\dd t}E_{\rm rad}
 =\int_{\mathbb R^2}
 \operatorname{Re}\langle a_t(0,y),p(y)\rangle_{\mathscr T}\,\dd^2y,
 \label{eq:V59-radiative-power}
\end{equation}
where
\begin{equation}
 E_{\rm rad}
 :=\frac12\int_{\Omega_x}
 \left(|a_t|^2+|\partial_xa|^2+|\nabla_ya|^2\right)
 \dd x\,\dd^2y.
 \label{eq:V59-radiative-energy}
\end{equation}
This is the V53 positive quotient energy, now written in the V58
boundary realization.

Define the normal characteristic variables
\begin{equation}
 u_{\rm in}:=a_t+p=a_t-\partial_xa,
 \qquad
 u_{\rm out}:=a_t-p=a_t+\partial_xa.
 \label{eq:V59-characteristics}
\end{equation}
Then
\begin{equation}
 \operatorname{Re}\langle a_t,p\rangle_{\mathscr T}
 =\frac14\left(|u_{\rm in}|^2-|u_{\rm out}|^2\right).
 \label{eq:V59-characteristic-flux}
\end{equation}
Thus the physical incoming and outgoing boundary bundles both have
rank two.  At gravitational glancing their fibres are the two classes
\((r_+,r_\times)\) identified in V58; the normal energy flux tends to
zero there, but the bundle rank does not jump.

\begin{proposition}[Exact quotient wave and flux]
\label{prop:V59-quotient-flux}
On the V58 wave-compatible York domain, the equations
\eqref{eq:V59-radiative-wave}--\eqref{eq:V59-characteristic-flux}
are independent of the choice of measurable polarization frame.
Their energy is positive definite on the nonzero-frequency radiative
quotient, and the boundary form has exactly two positive and two
negative characteristic squares.
\end{proposition}

\begin{proof}
V53 identifies the harmonic null cohomology with the rank-two TT fibre,
and its wave operator and Euclidean quotient energy are scalar in that
fibre.  V58 proves that the boundary York representative preserves the
essential plane and commutes with the wave equation on
\({\mathscr D}_Y^m\).  A unitary frame change preserves the fibre norm,
the scalar wave operator, and every displayed pairing.

Differentiate \eqref{eq:V59-radiative-energy}, integrate the spatial
derivatives by parts, and use the outward normal
\(D_n=-\partial_x\).  The interior terms cancel by
\eqref{eq:V59-radiative-wave}, leaving
\eqref{eq:V59-radiative-power}.  Expanding the two squares in
\eqref{eq:V59-characteristic-flux} proves the characteristic
decomposition and the rank statement.
\end{proof}

\subsection{A single positive quotient action and its interface}
\label{subsec:V59-action}

Introduce a second normal coordinate \(\rho>0\) and a continuum field
\[
 b=b(t,y,\rho)\in{\mathscr T}.
\]
Fix
\begin{equation}
 m>0,\qquad 0<v<1.
 \label{eq:V59-positive-parameters}
\end{equation}
The quotient action is
\begin{align}
 S_{\rm q}[a,b]
 &:=\frac12\int_{\mathbb R_t}\int_{\Omega_x}
 \left(|a_t|^2-|\partial_xa|^2-|\nabla_ya|^2\right)
 \dd x\,\dd^2y\,\dd t
 \nonumber\\
 &\quad
 +\frac m2\int_{\mathbb R_t}\int_{\mathbb R^2}
 \int_0^\infty
 \left(|b_t|^2-|\partial_\rho b|^2
       -v^2|\nabla_yb|^2\right)
 \dd\rho\,\dd^2y\,\dd t,
 \label{eq:V59-quotient-action}
\end{align}
with configuration domain constrained by
\begin{equation}
 b(t,y,0)=a(t,y,0).
 \label{eq:V59-trace-continuity}
\end{equation}
Set
\begin{equation}
 t_b:=-m\partial_\rho b\big|_{\rho=0}.
 \label{eq:V59-bath-traction}
\end{equation}

\begin{theorem}[Euler equations and natural conormal balance]
\label{thm:V59-action-interface}
Stationarity of \eqref{eq:V59-quotient-action} on the common-trace
domain gives
\begin{align}
 a_{tt}&=\partial_x^2a+\Delta_ya,
 \label{eq:V59-Euler-a}\\
 b_{tt}&=\partial_\rho^2b+v^2\Delta_yb,
 \label{eq:V59-Euler-b}\\
 b(0)&=a(0),
 \label{eq:V59-Euler-trace}\\
 p+t_b&=0.
 \label{eq:V59-traction-balance}
\end{align}
No fitted transfer function or additional boundary force is used.
\end{theorem}

\begin{proof}
Interior integration by parts gives
\eqref{eq:V59-Euler-a} and \eqref{eq:V59-Euler-b}.  Variations in the
configuration domain have one common boundary trace
\[
 \delta b(0)=\delta a(0)=:\delta z.
\]
The metric spatial term contributes
\(-\operatorname{Re}\langle p,\delta z\rangle_{\mathscr T}\), while
the continuum normal term contributes
\(-\operatorname{Re}\langle t_b,\delta z\rangle_{\mathscr T}\).
Their sum vanishes for every \(\delta z\) exactly when
\eqref{eq:V59-traction-balance} holds.
\end{proof}

The conserved positive energy is
\begin{align}
 E_{\rm q}
 &:=E_{\rm rad}
 +\frac m2\int_{\mathbb R^2}\int_0^\infty
 \left(
 |b_t|^2+|\partial_\rho b|^2
 +v^2|\nabla_yb|^2
 \right)\dd\rho\,\dd^2y.
 \label{eq:V59-total-energy}
\end{align}

\begin{corollary}[Exact positive exchange cancellation]
\label{cor:V59-energy-balance}
Every smooth solution of
\eqref{eq:V59-Euler-a}--\eqref{eq:V59-traction-balance} satisfies
\begin{equation}
 \frac{\dd}{\dd t}E_{\rm q}=0.
 \label{eq:V59-energy-conservation}
\end{equation}
\end{corollary}

\begin{proof}
Equation \eqref{eq:V59-radiative-power} gives the metric boundary
power
\(\operatorname{Re}\langle a_t,p\rangle_{\mathscr T}\).  Integration
by parts in the continuum gives
\(\operatorname{Re}\langle b_t(0),t_b\rangle_{\mathscr T}\).
Trace continuity gives \(b_t(0)=a_t(0)\), and the conormal balance
makes their sum zero.
\end{proof}

\subsection{Closed form, complete domain, and the unitary group}
\label{subsec:V59-closed-form}

Let
\begin{equation}
 {\mathcal H}_q
 :=
 L^2(\Omega_x;{\mathscr T})
 \oplus
 L^2(\mathbb R^2_y\times\mathbb R^+_\rho;
       {\mathscr T},m\,\dd y\,\dd\rho)
 \label{eq:V59-configuration-Hilbert}
\end{equation}
with its natural positive inner product.  Define
\begin{equation}
 {\mathcal V}_q
 :=
 \left\{
 (a,b)\in H^1(\Omega_x;{\mathscr T})
 \oplus H^1(\mathbb R^2\times\mathbb R^+;{\mathscr T}):
 \operatorname{Tr}_xa=\operatorname{Tr}_\rho b
 \right\}.
 \label{eq:V59-form-domain}
\end{equation}
The continuum \(H^1\) norm is understood with \(v\nabla_yb\);
because \(v>0\), it is equivalent to the ordinary \(H^1\) norm.  On
this domain put
\begin{align}
 {\mathfrak q}_v[(a,b)]
 &:=
 \int_{\Omega_x}\left(
 |\partial_xa|^2+|\nabla_ya|^2\right)\dd x\,\dd^2y
 \nonumber\\
 &\quad
 +m\int_{\mathbb R^2}\int_0^\infty
 \left(
 |\partial_\rho b|^2+v^2|\nabla_yb|^2
 \right)\dd\rho\,\dd^2y.
 \label{eq:V59-closed-form}
\end{align}

\begin{theorem}[Positive self-adjoint quotient operator]
\label{thm:V59-self-adjoint-operator}
The form \(({\mathfrak q}_v,{\mathcal V}_q)\) is densely defined,
closed, and nonnegative on \({\mathcal H}_q\).  Its associated operator
\({\mathcal A}_q\) is nonnegative and self-adjoint.  Among pairs with
\(H^2\) regularity, membership in its operator domain is equivalent to
\begin{equation}
 b(0)=a(0),\qquad
 D_na-m\partial_\rho b(0)=0.
 \label{eq:V59-operator-interface}
\end{equation}
In the interiors,
\begin{equation}
 {\mathcal A}_q(a,b)
 =
 \left(
 -\partial_x^2a-\Delta_ya,\,
 -\partial_\rho^2b-v^2\Delta_yb
 \right).
 \label{eq:V59-spatial-operator}
\end{equation}
Moreover,
\begin{equation}
 \ker{\mathcal A}_q=\{0\}.
 \label{eq:V59-operator-kernel}
\end{equation}
\end{theorem}

\begin{proof}
Smooth compactly supported pairs whose two traces vanish belong to
\({\mathcal V}_q\) and are dense in \({\mathcal H}_q\), so the form
domain is dense.  The trace maps
\(H^1(\mathbb R^3_+)\to H^{1/2}(\mathbb R^2)\) are continuous.
Therefore the common-trace condition is a closed linear condition in
the product \(H^1\) space.  The norm
\[
 \|(a,b)\|_{{\mathcal H}_q}^2
 +{\mathfrak q}_v[(a,b)]
\]
is equivalent on \({\mathcal V}_q\) to the product \(H^1\) norm, using
\(m>0\) and \(v>0\).  Hence the form is closed and nonnegative.

The representation theorem for closed nonnegative forms gives a unique
nonnegative self-adjoint operator \({\mathcal A}_q\).  Integration by
parts against a common-trace test pair gives the two interior operators
in \eqref{eq:V59-spatial-operator}.  Its boundary coefficient is
\[
 \left\langle D_na-m\partial_\rho b(0),\delta z
 \right\rangle_{\mathscr T}
 =\left\langle p+t_b,\delta z\right\rangle_{\mathscr T}.
\]
It vanishes for all common trace variations precisely when
\eqref{eq:V59-operator-interface} holds.  This proves the asserted
\(H^2\) characterization; the representation-form identity
characterizes the full operator domain.

If \({\mathcal A}_q(a,b)=0\), then
\({\mathfrak q}_v[(a,b)]=0\).  Both fields have zero spatial gradient
and are therefore constant on their connected half-spaces.  The only
constant \(L^2\) fields on these infinite domains are zero.  This proves
\eqref{eq:V59-operator-kernel}.
\end{proof}

Let \(\dot{\mathcal V}_q\) be the completion of
\({\mathcal V}_q\) in the norm
\({\mathfrak q}_v^{1/2}\), and set
\begin{equation}
 {\mathcal E}_q:=\dot{\mathcal V}_q\oplus{\mathcal H}_q,
 \qquad
 \|(u,\dot u)\|_{{\mathcal E}_q}^2
 ={\mathfrak q}_v[u]+\|\dot u\|_{{\mathcal H}_q}^2.
 \label{eq:V59-energy-space}
\end{equation}
The zero-kernel result makes this a norm rather than a seminorm.
For \(u\in\dot{\mathcal V}_q\), the condition
\({\mathcal A}_qu\in{\mathcal H}_q\) means that the distribution
induced by the form is represented by an element of
\({\mathcal H}_q\); equivalently, it is the spectral-calculus
extension of \({\mathcal A}_q\) to the homogeneous completion.

\begin{theorem}[Skew-adjoint finite-energy generator]
\label{thm:V59-skew-adjoint-generator}
The wave operator
\begin{equation}
 {\mathcal G}_q
 \binom{u}{w}
 :=
 \binom{w}{-{\mathcal A}_qu},
 \qquad
 D({\mathcal G}_q)
 =\left\{(u,w)\in\dot{\mathcal V}_q\oplus{\mathcal H}_q:
 w\in{\mathcal V}_q,\ {\mathcal A}_qu\in{\mathcal H}_q\right\},
 \label{eq:V59-wave-generator}
\end{equation}
has a unique skew-adjoint realization on
\({\mathcal E}_q\).  It generates a strongly continuous unitary group
and its squared norm is \(2E_{\rm q}\).  In particular, the quotient
interface has a complete maximal time-generator domain; the power
identity is not being used as a substitute for maximality.
\end{theorem}

\begin{proof}
By Theorem~\ref{thm:V59-self-adjoint-operator},
\({\mathcal A}_q\) is nonnegative, self-adjoint, and injective.  Its
spectral resolution defines
\(\cos(t{\mathcal A}_q^{1/2})\) and
\({\mathcal A}_q^{-1/2}\sin(t{\mathcal A}_q^{1/2})\) on the homogeneous
energy completion.  The solution is
\begin{align*}
 u(t)
 &=
 \cos(t{\mathcal A}_q^{1/2})u_0
 +{\mathcal A}_q^{-1/2}
  \sin(t{\mathcal A}_q^{1/2})w_0,\\
 w(t)
 &=
 -{\mathcal A}_q^{1/2}
  \sin(t{\mathcal A}_q^{1/2})u_0
 +\cos(t{\mathcal A}_q^{1/2})w_0.
\end{align*}
On each positive spectral fibre this is a rotation of
\(({\mathcal A}_q^{1/2}u,w)\), so it is unitary in
\eqref{eq:V59-energy-space}.  The group law and strong continuity
follow from the spectral calculus.  The spectral characterization
of its infinitesimal generator gives exactly the domain in
\eqref{eq:V59-wave-generator}; Stone's theorem makes that generator
skew-adjoint.  The form identity identifies
\(\|{\mathcal A}_q^{1/2}u\|^2={\mathfrak q}_v[u]\), which proves the
energy statement.
\end{proof}

\begin{firewall}[The level at which maximality is proved]
\label{fire:V59-generator-scope}
Theorems~\ref{thm:V59-self-adjoint-operator} and
\ref{thm:V59-skew-adjoint-generator} prove a complete positive
generator for the two-polarization quotient and its continuum.  The
configuration variable \(a\) is a section of the radiative bundle, so
the action is generally pseudodifferential when pulled back to an
unreduced metric.  No local ten-component metric action or off-shell
BRST domain is inferred from the quotient form.
\end{firewall}

\subsection{The exact stable interface matrix}
\label{subsec:V59-stable-symbol}

For a Laplace--Fourier mode
\[
 e^{st+ik\cdot y},
 \qquad \operatorname{Re}s\ge0,\qquad r:=|k|,
\]
use the retarded roots
\begin{equation}
 \lambda:=\sqrt{s^2+r^2},
 \qquad
 \kappa_v:=\sqrt{s^2+v^2r^2}.
 \label{eq:V59-roots}
\end{equation}
The outgoing homogeneous fields are
\begin{equation}
 a(x)=e^{-\lambda x}a_0,\qquad
 b(\rho)=e^{-\kappa_v\rho}a_0.
 \label{eq:V59-outgoing-fields}
\end{equation}
Their conormals are
\[
 p=\lambda a_0,\qquad
 t_b=m\kappa_va_0.
\]
Thus \eqref{eq:V59-traction-balance} is
\begin{equation}
 {\cal B}_{\rm rad}(s,k)a_0=0,
 \qquad
 {\cal B}_{\rm rad}
 :=(\lambda+m\kappa_v)I_2.
 \label{eq:V59-interface-matrix}
\end{equation}

\begin{theorem}[Complete closed-cone radiative floor]
\label{thm:V59-closed-cone-floor}
For every \(m>0\) and \(0<v<1\),
\begin{equation}
 \det{\cal B}_{\rm rad}(s,k)\ne0
 \label{eq:V59-no-interface-zero}
\end{equation}
at every nonzero open or closed stable frequency.  On the normalized
closed stable cone
\begin{equation}
 \operatorname{Re}s\ge0,\qquad |s|^2+r^2=1,
 \label{eq:V59-normalized-cone}
\end{equation}
there is a constant \(c_{m,v}>0\) such that
\begin{equation}
 s_{\min}({\cal B}_{\rm rad}(s,k))
 =|\lambda+m\kappa_v|
 \ge c_{m,v}.
 \label{eq:V59-uniform-floor}
\end{equation}
In particular, the two radiative glancing classes and the continuum
threshold are both controlled.
\end{theorem}

\begin{proof}
If \(\operatorname{Re}s>0\), the retarded square roots have strictly
positive real parts.  Therefore
\[
 \operatorname{Re}(\lambda+m\kappa_v)>0,
\]
so their sum cannot vanish.

On the closed boundary write \(s=i\omega+0\).  By symmetry it suffices
to take \(\omega\ge0\).  If \(0\le\omega<vr\), both roots are
nonnegative real and are not simultaneously zero.  If
\(vr<\omega<r\), then \(\lambda>0\) is real and
\(\kappa_v=i\sqrt{\omega^2-v^2r^2}\), so their sum has positive real
part.  If \(\omega>r\), both roots are positive imaginary and their
positive linear combination cannot vanish.  At the two endpoints,
\begin{align}
 \omega=vr:\quad&
 \lambda=r\sqrt{1-v^2},\quad \kappa_v=0,
 \label{eq:V59-bath-threshold}\\
 \omega=r:\quad&
 \lambda=0,\quad
 \kappa_v=ir\sqrt{1-v^2}.
 \label{eq:V59-grav-glancing}
\end{align}
The displayed nonzero values cover the thresholds when \(r>0\).
If \(r=0\), analytic continuation gives
\(\lambda=\kappa_v=s\), and the sum is \((1+m)s\ne0\).
If \(s=0\) and \(r>0\), it is \(r(1+mv)>0\).
This exhausts the closed boundary and proves
\eqref{eq:V59-no-interface-zero}.

The normalized cone is compact, both retarded roots are continuous on
it, and the continuous function
\(|\lambda+m\kappa_v|\) has no zero there.  It therefore has the
strictly positive minimum in \eqref{eq:V59-uniform-floor}.
\end{proof}

\subsection{Characteristic scattering and strict passivity}
\label{subsec:V59-scattering}

Eliminating the outgoing continuum alone gives
\[
 t_b=m\kappa_va(0),\qquad
 p=-m\kappa_va(0).
\]
At the metric boundary,
\begin{align}
 u_{\rm in}
 &=(s-m\kappa_v)a(0),
 \label{eq:V59-uin-symbol}\\
 u_{\rm out}
 &=(s+m\kappa_v)a(0).
 \label{eq:V59-uout-symbol}
\end{align}
The denominator in the scattering relation is nonzero in the open
half-plane, and
\begin{equation}
 u_{\rm in}={\mathcal S}_{m,v}(s,k)u_{\rm out},
 \qquad
 {\mathcal S}_{m,v}
 :=\frac{s-m\kappa_v}{s+m\kappa_v}I_2.
 \label{eq:V59-scattering-matrix}
\end{equation}

\begin{lemma}[Strict retarded power inequality]
\label{lem:V59-retarded-power}
For \(\operatorname{Re}s>0\),
\begin{equation}
 \operatorname{Re}(\overline s\,\kappa_v)>0.
 \label{eq:V59-retarded-positive}
\end{equation}
Consequently
\begin{equation}
 |s+m\kappa_v|^2-|s-m\kappa_v|^2
 =4m\operatorname{Re}(\overline s\,\kappa_v)>0.
 \label{eq:V59-Cayley-gap}
\end{equation}
\end{lemma}

\begin{proof}
Write \(s=\sigma+i\omega\), with \(\sigma>0\), and
\(\kappa_v=u+iq\), with \(u>0\).  The imaginary part of
\(\kappa_v^2=s^2+v^2r^2\) gives
\[
 uq=\sigma\omega.
\]
Hence
\[
 \operatorname{Re}(\overline s\,\kappa_v)
 =\sigma u+\omega q
 =\frac{\sigma}{u}(u^2+\omega^2)>0.
\]
Expanding the two moduli gives
\eqref{eq:V59-Cayley-gap}.
\end{proof}

\begin{theorem}[Strictly contractive physical scattering]
\label{thm:V59-scattering-contraction}
For every open stable frequency,
\begin{equation}
 \|{\mathcal S}_{m,v}(s,k)\|<1.
 \label{eq:V59-strict-contraction}
\end{equation}
Its closed-boundary limits have norm at most one.  Equality occurs
exactly where the continuum has no propagating normal energy flux; a
strict loss of metric boundary flux is energy transported into the
explicit conservative continuum, not an unrecorded dissipative term.
\end{theorem}

\begin{proof}
The denominator \(s+m\kappa_v\) has positive real part and is nonzero.
Lemma~\ref{lem:V59-retarded-power} gives
\[
 \left|\frac{s-m\kappa_v}{s+m\kappa_v}\right|<1.
\]
For a nonzero boundary frequency \(s=i\omega+0\), if
\(|\omega|\le vr\), then \(\kappa_v\) is real and nonnegative, and
\[
 |i\omega-m\kappa_v|=|i\omega+m\kappa_v|.
\]
This is the evanescent or threshold regime and has no propagating
normal continuum flux.  If \(|\omega|>vr\), write
\(\kappa_v=i\,\operatorname{sgn}(\omega)
\sqrt{\omega^2-v^2r^2}\).  Both terms in the denominator then have the
same imaginary sign, and
\[
 \left|\frac{s-m\kappa_v}{s+m\kappa_v}\right|
 =
 \frac{\left||\omega|-m\sqrt{\omega^2-v^2r^2}\right|}
      {|\omega|+m\sqrt{\omega^2-v^2r^2}}
 <1.
\]
This proves both the closed-boundary bound and its equality statement.
The full metric--continuum energy is conserved by
Corollary~\ref{cor:V59-energy-balance}; the apparent contraction is
therefore precisely the outgoing radiation into the retained
\(\rho\)-continuum.
\end{proof}

\subsection{A homogeneous subsidiary block}
\label{subsec:V59-subsidiary}

The V34 harmonic covector satisfies the componentwise wave equation
\[
 \Box{\cal C}_\mu=0.
\]
To record the constraint domain without feeding it the physical
traction, introduce the four-component subsidiary block
\begin{align}
 c_{tt}&=\partial_x^2c+\Delta_yc,
 \label{eq:V59-constraint-wave}\\
 c|_{x=0}&=0.
 \label{eq:V59-constraint-Dirichlet}
\end{align}
Its positive Dirichlet form is
\begin{equation}
 {\mathfrak q}_{\cal C}[c]
 :=\int_{\Omega_x}
 \left(|\partial_xc|^2+|\nabla_yc|^2\right)
 \dd x\,\dd^2y,
 \qquad
 D({\mathfrak q}_{\cal C})=H^1_0(\Omega_x;\mathbb C^4).
 \label{eq:V59-constraint-form}
\end{equation}

\begin{theorem}[Invariant zero-constraint subspace]
\label{thm:V59-constraint-invariance}
The direct sum of the quotient form
\({\mathfrak q}_v\) and the Dirichlet form
\({\mathfrak q}_{\cal C}\) defines a nonnegative self-adjoint spatial
operator and a skew-adjoint finite-energy wave generator.  The closed
subspace
\begin{equation}
 c=0,\qquad c_t=0
 \label{eq:V59-zero-constraint-subspace}
\end{equation}
is invariant.  Thus the two-polarization continuum can be evolved
together with a homogeneous harmonic subsidiary system without the
rank-three source \({\mathsf S}_{\cal C}Cr\) found in V58.
\end{theorem}

\begin{proof}
The Dirichlet Laplacian on four components is nonnegative and
self-adjoint by the same closed-form argument as
Theorem~\ref{thm:V59-self-adjoint-operator}.  Direct sums preserve
self-adjointness, nonnegativity, and the spectral wave construction.
The zero initial state is the unique solution of
\eqref{eq:V59-constraint-wave}--\eqref{eq:V59-constraint-Dirichlet},
so \eqref{eq:V59-zero-constraint-subspace} is invariant.  The quotient
interface contains only \((a,b)\); no term occurs in the \(c\)-equation.
\end{proof}

\begin{firewall}[A subsidiary block is not yet the metric identity]
\label{fire:V59-subsidiary-scope}
Theorem~\ref{thm:V59-constraint-invariance} proves a compatible
positive abstract block.  To identify \(c\) with
\({\cal C}(\bar h)\), the inverse York/BRST lift must prove that the
ten-component metric boundary rows imply
\eqref{eq:V59-constraint-Dirichlet} and the quotient interface
simultaneously.  V58 proves the required projector identity on its
compatibility domain, but it does not provide this complete inverse
boundary lift.  No such identification is claimed in V59.
\end{firewall}

\subsection{Exact threshold and witness comparisons}
\label{subsec:V59-comparisons}

At gravitational glancing \(s=ir\), \(r>0\),
\begin{equation}
 {\cal B}_{\rm rad}(ir,k)
 =imr\sqrt{1-v^2}\,I_2,
 \qquad
 \left|\det{\cal B}_{\rm rad}\right|
 =m^2r^2(1-v^2)>0.
 \label{eq:V59-glancing-value}
\end{equation}
At the continuum threshold \(s=ivr\),
\begin{equation}
 {\cal B}_{\rm rad}(ivr,k)
 =r\sqrt{1-v^2}\,I_2.
 \label{eq:V59-threshold-value}
\end{equation}
At zero temporal frequency,
\begin{equation}
 {\cal B}_{\rm rad}(0,k)
 =r(1+mv)I_2.
 \label{eq:V59-static-value}
\end{equation}
These three values show directly that the interface does not trade one
cone defect for another.

For the exact V57 numbers \(v=1/2\) and
\[
 \omega_T^2=\frac{5-\sqrt{13}}6,
 \qquad r=1,
\]
the V59 physical interface factor is
\begin{equation}
 \lambda_T+m\kappa_T
 =
 \sqrt{1-\omega_T^2}
 +m\sqrt{\frac14-\omega_T^2}>0.
 \label{eq:V59-old-witness-removed}
\end{equation}
Thus the old low-cone zero is absent from the new quotient interface
for every \(m>0\).  This comparison does not remove the V57
determinant; V58 proved that its kernel did not satisfy the homogeneous
subsidiary row.

\subsection{Scope firewall and the V60 lift}
\label{subsec:V59-lift-scope}

V59 closes the following linear problem: the positive two-polarization
Einstein radiative quotient, coupled to a positive local
half-line continuum, has a single quadratic action, a closed
self-adjoint spatial realization, a skew-adjoint finite-energy
generator, a conserved energy, a contractive retarded scattering
matrix, and a uniform zero-free stable interface.

The construction does not yet close the ten-component Einstein system.
The remaining lift must supply, from one common boundary action:

\begin{enumerate}
\item a bounded map from the V58 York/TT quotient trace and conormal back
to the metric--embedding first-order state;
\item four metric rows which imply the homogeneous subsidiary boundary
condition \({\cal C}|_{\partial M}=0\);
\item the complementary gauge/embedding rows with constant rank through
glancing and the second Dirichlet compatibility required in V58;
\item equality of the quotient conormal balance with the TT projection
of the Brown--York/embedding variation;
\item the corresponding off-shell ghost, antighost, and multiplier
domain; and
\item a stress/Noether identity showing that the continuum boundary
momentum enters the nonlinear Bianchi equation with the same sign.
\end{enumerate}

V60 is a scheduled five-version audit.  It must first inspect the newest
completed NS and Yang--Mills notes in the shared folder.  It should
import a result only if it supplies one of these six missing operations.
After that audit, V60 should attempt the first three lift rows at the
frozen linear level and return an exact rank or commutator obstruction
if they cannot coexist.

\subsection{Finite certificates}
\label{subsec:V59-certificates}

For \(m=1\), \(v=1/2\), the exact cone values are
\begin{equation}
 {\cal B}_{\rm rad}(ir,k)=\frac{i\sqrt3}{2}rI_2,
 \qquad
 {\cal B}_{\rm rad}\left(\frac{ir}{2},k\right)
 =\frac{\sqrt3}{2}rI_2,
 \qquad
 {\cal B}_{\rm rad}(0,k)=\frac32rI_2.
 \label{eq:V59-exact-values-card}
\end{equation}
A normalized closed-cone certificate stores
\begin{equation}
 \rho_{\rm rad}(m,v)
 :=\inf_{\substack{\operatorname{Re}s\ge0\\|s|^2+r^2=1}}
 |\lambda+m\kappa_v|>0.
 \label{eq:V59-floor-card}
\end{equation}
An open-half-plane scattering card stores
\begin{equation}
 r_{\rm Cayley}
 :=
 |s+m\kappa_v|^2-|s-m\kappa_v|^2
 -4m\operatorname{Re}(\overline s\,\kappa_v)=0,
 \qquad
 \|{\cal S}_{m,v}\|<1.
 \label{eq:V59-scattering-card}
\end{equation}
The form-domain card stores the common trace and conormal residuals
\begin{equation}
 r_{\rm tr}:=\|\operatorname{Tr}_xa-\operatorname{Tr}_\rho b\|,
 \qquad
 r_{\rm con}:=\|D_na-m\partial_\rho b(0)\|,
 \label{eq:V59-domain-card}
\end{equation}
which vanish on \(D({\cal A}_q)\).

\subsection{V59 result ledger}
\label{subsec:V59-ledger}

The new deductions are:

\begin{enumerate}
\item the boundary-normal characteristic form of the radiative quotient
has exactly two incoming and two outgoing squares;
\item trace continuity and conormal balance follow from one positive
two-polarization quotient action;
\item the common-trace quadratic form is dense, closed, nonnegative, and
has a self-adjoint spatial operator with the complete interface domain;
\item the associated finite-energy wave generator is skew-adjoint and
defines a unitary group;
\item the stable interface
\((\lambda+m\kappa_v)I_2\) has a uniform floor on the complete
normalized closed stable cone;
\item its retarded characteristic scattering is strictly contractive in
the open half-plane and records all lost metric flux in the explicit
continuum; and
\item a direct-sum homogeneous harmonic wave block preserves the zero
constraint subspace without the old rank-three constraint source.
\end{enumerate}

These statements construct the positive stable radiative-continuum
generator which was missing after V53, but only on the physical
two-polarization quotient.  The full metric/embedding/ghost lift,
nonlinear stress/Bianchi identity, and coupled Higgs, fermion, neutrino,
and Yang--Mills sectors remain open.  The full Standard Model--Einstein
continuum is not proved.

\subsection{Append-only record}

This V59 material was appended to the exact V58 parent whose SHA--256
digest is
\[
\texttt{c792a8785cca6635103cf579bfe4b3c2c101464662e73ad0fa60ca03a1115822}.
\]
The next research note is V60.  V60 begins with the scheduled
cross-program inspection of the newest completed Yang--Mills and
Navier--Stokes notes before selecting the full metric lift.

% END APPENDED VERSION V59 MATERIAL
% BEGIN APPENDED VERSION V60 MATERIAL

\clearpage
\section{V60: companion audit, exact frozen TT trace lift, and the
rank-two Bianchi-compensator obstruction}
\label{sec:V60-frozen-lift}

\subsection{Purpose, scheduled audit, and nonduplication boundary}
\label{subsec:V60-purpose}

V59 constructed a positive self-adjoint and skew-adjoint
radiative--continuum system on the two-polarization quotient.  It did
not identify that quotient system with a boundary domain for the full
ten-component metric.  In particular, V59 left six precise operations:
a trace/conormal lift, four homogeneous harmonic rows, a
gauge/embedding complement with the V58 second compatibility, the
Brown--York projection identity, an off-shell BRST domain, and a
nonlinear stress/Bianchi identity.

This is the scheduled five-version companion audit.  The newest
Navier--Stokes file inspected was
\texttt{Renewal\_NS\_Stability\_Research\_Notes\_V31.tex}.  Its V31
appendix proves that flat dyadic marking recovers critical coarea only
after the exact first-moment weight is restored, and that the available
nonlinear high-parent formation estimate is paid by
\(\int\|\Lambda u\|_2^4\), which is continuation-class stock.  Its
relevant lesson here is a firewall: a correct positive estimate must
not be used to supply a missing domain map.

During the audit the Yang--Mills programme replaced its V38 filename by
\texttt{Renewal\_Yang\_Mills\_Mass\_Gap\_Research\_Notes\_V39.tex}
and completed the V39 appendix.  V39 evaluates the remaining
thin-interior Berenstein--Zelevinsky interval, proves a universal
squared one-link cold tail, closes the disjoint-overlap and global
ternary marked-tree rows, and derives an exact all-order connected
replica/tree decomposition.  Its all-order analytic fibre estimate
remains open.  The relevant lesson is to identify and compute the exact
finite-dimensional residual before attempting a global compiler.

Neither companion theorem supplies a York inverse, a metric boundary
Noether identity, or a harmonic constraint row.  No NS or Yang--Mills
estimate is imported as an Einstein lemma.  The audit instead changes
the acquisition order: V60 first isolates the exact finite-dimensional
residual of the V59 metric lift.  The outcome is a positive trace
splitting and a rank-two conormal obstruction.  This prevents the V41
five-row incidence from being reused under a different name.

\subsection{Both-sheet glancing coordinates}
\label{subsec:V60-glancing-coordinates}

Use the V41 reduced coordinate order
\[
 x=(\tau,v_0,v_\parallel,v_\perp,\phi)^T
\]
after rotating a nonzero tangential covector to
\(k=r e_1\), \(r>0\).  Let
\[
 s=\epsilon ir,\qquad \epsilon\in\{+1,-1\},\qquad \lambda=0,
\]
and write
\[
 {\mathscr K}_\epsilon
 :=\ker L_0(\epsilon ir,r e_1).
\]
The V41 formula \eqref{eq:V41-A-definition} gives
\begin{equation}
 {\mathscr K}_\epsilon
 =
 \left\{
 (-3\gamma,\alpha,\epsilon\alpha,\beta,\gamma)^T:
 \alpha,\beta,\gamma\in\mathbb C
 \right\}.
 \label{eq:V60-glancing-kernel}
\end{equation}
Thus define
\begin{align}
 g_\epsilon&:=(0,1,\epsilon,0,0)^T,
 \label{eq:V60-gauge-vector}\\
 r_\times&:=(0,0,0,1,0)^T,
 \label{eq:V60-cross-vector}\\
 r_+&:=(-3,0,0,0,1)^T.
 \label{eq:V60-plus-vector}
\end{align}
For \(\epsilon=+1\), these are the V58 vectors
\eqref{eq:V58-gparallel}--\eqref{eq:V58-rplus}.  Equation
\eqref{eq:V60-glancing-kernel} proves the opposite-sheet signs rather
than inferring them by symmetry.

Define the residual-gauge coordinate and the normalized radiative
trace by
\begin{align}
 G_\epsilon x
 &:=
 \frac12(v_0+\epsilon v_\parallel),
 \label{eq:V60-gauge-coordinate}\\
 Q_{\rm rad}x
 &:=
 \binom{-\sqrt2\,\phi}{\sqrt2\,v_\perp}
 =:\binom{a_+}{a_\times}.
 \label{eq:V60-radiative-trace}
\end{align}
The normalization agrees with the V58 TT images:
\[
 Q_{\rm rad}g_\epsilon=0,\qquad
 Q_{\rm rad}r_\times=(0,\sqrt2)^T,\qquad
 Q_{\rm rad}r_+=(-\sqrt2,0)^T.
\]
Define a right inverse
\begin{equation}
 R_{\rm rad}
 \binom{a_+}{a_\times}
 :=
 \begin{pmatrix}
 3a_+/\sqrt2\\
 0\\
 0\\
 a_\times/\sqrt2\\
 -a_+/\sqrt2
 \end{pmatrix}.
 \label{eq:V60-radiative-right-inverse}
\end{equation}

\begin{theorem}[Exact both-sheet frozen trace lift]
\label{thm:V60-trace-lift}
For each \(\epsilon=\pm1\), the map
\begin{equation}
 (G_\epsilon,Q_{\rm rad}):
 {\mathscr K}_\epsilon\longrightarrow
 \mathbb C\oplus\mathbb C^2
 \label{eq:V60-kernel-coordinate-map}
\end{equation}
is an isomorphism with inverse
\begin{equation}
 (\alpha,a)\longmapsto
 \alpha g_\epsilon+R_{\rm rad}a.
 \label{eq:V60-kernel-coordinate-inverse}
\end{equation}
More explicitly,
\begin{align}
 Q_{\rm rad}R_{\rm rad}&=I_2,
 &
 G_\epsilon R_{\rm rad}&=0,
 \label{eq:V60-splitting-identities}\\
 x&=g_\epsilon G_\epsilon x+
 R_{\rm rad}Q_{\rm rad}x
 \qquad(x\in{\mathscr K}_\epsilon).
 \label{eq:V60-kernel-decomposition}
\end{align}
The operator
\begin{equation}
 P_\epsilon:=I_5-g_\epsilon G_\epsilon
 \label{eq:V60-gauge-complement-projector}
\end{equation}
is the Euclidean orthogonal projector of rank four onto
\(\ker G_\epsilon\), has norm one, and obeys
\begin{equation}
 P_\epsilon{\mathscr K}_\epsilon
 =\operatorname{span}\{r_\times,r_+\},
 \qquad
 P_\epsilon|_{{\mathscr K}_\epsilon}
 =R_{\rm rad}Q_{\rm rad}.
 \label{eq:V60-projector-on-kernel}
\end{equation}
In the displayed Euclidean norms,
\begin{equation}
 \|Q_{\rm rad}\|=\sqrt2,\qquad
 \|R_{\rm rad}\|=\sqrt5.
 \label{eq:V60-lift-norms}
\end{equation}
After tangential rotation, the same formulas with
\(v_\parallel=\widehat k\cdot v_{\rm sp}\) and
\(v_\perp=\widehat k^\perp\cdot v_{\rm sp}\) form a smooth
constant-rank family on each glancing sheet.
\end{theorem}

\begin{proof}
At glancing, \(L_0=-A(d)\).  The three \(v\)-rows of
\eqref{eq:V41-A-definition} give
\[
 d_a(\tau/6+\phi/2)=0.
\]
Since \(d\ne0\), this is \(\tau=-3\phi\).  The first row gives
\[
 -s v_0+irv_\parallel=0,
\]
and hence \(v_\parallel=\epsilon v_0\).  The variables
\(v_0,v_\perp,\phi\) are free, proving
\eqref{eq:V60-glancing-kernel}.

Direct substitution gives
\eqref{eq:V60-splitting-identities}.  If
\(x=\alpha g_\epsilon+\beta r_\times+\gamma r_+\), then
\(G_\epsilon x=\alpha\) and
\(R_{\rm rad}Q_{\rm rad}x=\beta r_\times+\gamma r_+\).
This proves the isomorphism and
\eqref{eq:V60-kernel-decomposition}.

The covector defining \(G_\epsilon\) is
\(g_\epsilon^*/\|g_\epsilon\|^2\).  Therefore
\(g_\epsilon G_\epsilon\) is the orthogonal rank-one projector onto
\(\operatorname{span}\{g_\epsilon\}\), which proves all assertions
about \(P_\epsilon\).  The two columns of \(R_{\rm rad}\) are
orthogonal and have squared norms \(5\) and \(1/2\); the two rows of
\(Q_{\rm rad}\) are orthogonal and have squared norm \(2\).  This proves
\eqref{eq:V60-lift-norms}.  Tangential rotation is unitary and preserves
all ranks and norms.
\end{proof}

\begin{corollary}[The V59 configuration trace has a bounded frozen lift]
\label{cor:V60-V59-trace-lift}
At either gravitational glancing sheet, the V59 common-trace condition
can be lifted as
\begin{equation}
 b(0)=Q_{\rm rad}x,\qquad
 x_{\rm rad}=R_{\rm rad}b(0),
 \label{eq:V60-lifted-common-trace}
\end{equation}
after the gauge condition \(G_\epsilon x=0\) is imposed on
\({\mathscr K}_\epsilon\).  This is a bounded right inverse of the
radiative trace and does not use the old three-port bath.
\end{corollary}

\begin{firewall}[Trace lifting is not conormal lifting]
\label{fire:V60-trace-not-conormal}
Theorem~\ref{thm:V60-trace-lift} solves the finite-dimensional
configuration trace problem at glancing.  It does not say that
\(R_{\rm rad}\) may be inserted as a force into the V41 graph.  Metric
conormals are constrained by the four V41 harmonic residual rows.  The
next theorem computes this independent obstruction.
\end{firewall}

\subsection{The natural radiative conormal and its constraint image}
\label{subsec:V60-conormal-obstruction}

Let
\begin{equation}
 z_{\rm phys}:=(v_\perp,\phi)^T,
 \qquad
 a=Uz_{\rm phys},
 \qquad
 U:=
 \begin{pmatrix}
 0&-\sqrt2\\
 \sqrt2&0
 \end{pmatrix}.
 \label{eq:V60-U-map}
\end{equation}
Thus \(U^*U=2I_2\).  If
\(\theta=(\theta_+,\theta_\times)^T\) is a V59 TT conormal, equality of
virtual work gives the physical V41 port force
\begin{equation}
 r_{\rm phys}=U^*\theta.
 \label{eq:V60-virtual-work-force}
\end{equation}
Embed it in the old three-port order
\((v_\parallel,v_\perp,\phi)\) by
\begin{equation}
 \iota_{\rm phys}
 :=
 \begin{pmatrix}
 0&0\\
 1&0\\
 0&1
 \end{pmatrix}.
 \label{eq:V60-physical-port-inclusion}
\end{equation}
Using the reciprocal incidence \(C\) from
\eqref{eq:V58-C-matrix}, define
\begin{equation}
 C_{\rm rad}:=C\iota_{\rm phys}U^*.
 \label{eq:V60-radiative-incidence-definition}
\end{equation}
In the residual order
\((f_\tau,f_{v_0},f_{v_\parallel},f_{v_\perp},f_\phi)\),
direct multiplication gives
\begin{equation}
 C_{\rm rad}
 =
 \begin{pmatrix}
 3/\sqrt2&0\\
 0&0\\
 0&0\\
 0&1/\sqrt2\\
 -1/\sqrt2&0
 \end{pmatrix}
 =R_{\rm rad}.
 \label{eq:V60-radiative-incidence}
\end{equation}
The equality with \(R_{\rm rad}\) is a consequence of the chosen
DeWitt virtual-work normalization; it does not identify configuration
and momentum domains.

Recall from \eqref{eq:V58-constraint-source-map} that
\[
 {\mathsf S}_{\cal C}f
 =(f_{v_0},f_{v_\parallel},f_{v_\perp},
   (f_\phi-f_\tau)/2)^T.
\]
Therefore
\begin{equation}
 D_{\rm rad}
 :={\mathsf S}_{\cal C}C_{\rm rad}
 =
 \begin{pmatrix}
 0&0\\
 0&0\\
 0&1/\sqrt2\\
 -\sqrt2&0
 \end{pmatrix}.
 \label{eq:V60-radiative-constraint-image}
\end{equation}

\begin{theorem}[No two-polarization conormal lift in five homogeneous
metric rows]
\label{thm:V60-five-row-no-go}
The natural reciprocal radiative incidence satisfies
\begin{equation}
 \operatorname{rank}C_{\rm rad}
 =\operatorname{rank}D_{\rm rad}=2.
 \label{eq:V60-natural-ranks}
\end{equation}
In particular, every nonzero TT conormal produces a nonzero harmonic
boundary source.  The plus traction sources \(C_n\), while the cross
traction sources \(C_\perp\), which is \(C_2\) in the aligned frame.

More generally, at a fixed frequency let
\(L:\mathbb C^2\to\mathbb C^5\) be any linear additive lift of the two
TT conormals into the V41 residual \(f=p-A(D)x\).  The map \(L\) may
depend on that frequency.  If all four harmonic rows remain
homogeneous,
\begin{equation}
 {\mathsf S}_{\cal C}L=0,
 \label{eq:V60-homogeneous-lift-condition}
\end{equation}
then
\begin{equation}
 \operatorname{rank}L\le1.
 \label{eq:V60-lift-rank-bound}
\end{equation}
Consequently no injective two-polarization conormal lift exists in the
old five residual rows.  The conclusion is unchanged by the choice of
graph sign or by an invertible renormalization of the TT variables.
\end{theorem}

\begin{proof}
Equation \eqref{eq:V60-radiative-constraint-image} has two nonzero
entries in different rows and different columns, proving
\eqref{eq:V60-natural-ranks} and the polarization statement.

For the general claim, V58 proved
\[
 \ker{\mathsf S}_{\cal C}
 =\operatorname{span}\{(1,0,0,0,1)^T\}.
\]
Condition \eqref{eq:V60-homogeneous-lift-condition} puts
\(\operatorname{ran}L\) inside this one-dimensional kernel, so
\(\operatorname{rank}L\le1\).  An injective map from \(\mathbb C^2\)
has rank two, a contradiction.  Multiplying \(L\) by a sign or
precomposing it with an isomorphism does not change its rank or the
kernel inclusion.
\end{proof}

\begin{corollary}[The obstruction survives removal of the gauge port]
\label{cor:V60-gauge-removal-no-help}
Deleting the longitudinal gauge port \(v_\parallel\) does not repair
the conormal lift.  The physical columns
\((v_\perp,\phi)\) already have the injective constraint image
\eqref{eq:V60-radiative-constraint-image}.  Thus the obstruction is
carried by the two TT polarizations themselves, not by the residual
gauge line \(g_\epsilon\).
\end{corollary}

\subsection{A minimal algebraic Bianchi target}
\label{subsec:V60-minimal-compensator}

Let
\begin{equation}
 E_{2n}:\mathbb C^2\longrightarrow\mathbb C^4,
 \qquad
 E_{2n}\binom{q_\perp}{q_n}
 :=
 \begin{pmatrix}0\\0\\q_\perp\\q_n\end{pmatrix},
 \label{eq:V60-compensator-inclusion}
\end{equation}
and set
\begin{equation}
 K_{\rm B}
 :=
 \begin{pmatrix}
 0&1/\sqrt2\\
 -\sqrt2&0
 \end{pmatrix}.
 \label{eq:V60-Bianchi-matrix}
\end{equation}
Then
\begin{equation}
 D_{\rm rad}=E_{2n}K_{\rm B},
 \qquad
 \det K_{\rm B}=1.
 \label{eq:V60-Bianchi-factorization}
\end{equation}
For a metric residual \(f\in\mathbb C^5\) and a two-component
compensator Euler covector \(\zeta\in\mathbb C^2\), define the extended
constraint source
\begin{equation}
 \widetilde{\mathsf S}_{\cal C}(f,\zeta)
 :=
 {\mathsf S}_{\cal C}f+E_{2n}\zeta.
 \label{eq:V60-extended-constraint-map}
\end{equation}

\begin{theorem}[Two algebraic compensator components are necessary and sufficient]
\label{thm:V60-minimal-compensator}
For either graph sign \(\sigma\in\{+1,-1\}\) and every TT conormal
\(\theta\),
\begin{equation}
 (f,\zeta)
 =
 \left(\sigma C_{\rm rad}\theta,
       -\sigma K_{\rm B}\theta\right)
 \label{eq:V60-compensated-incidence}
\end{equation}
obeys
\begin{equation}
 \widetilde{\mathsf S}_{\cal C}(f,\zeta)=0.
 \label{eq:V60-compensated-constraint}
\end{equation}
This compensation is rank-minimal.  Precisely, if \(Z\) is a finite
dimensional compensator space and maps
\[
 B:Z\to\mathbb C^4,\qquad
 M:\mathbb C^2\to Z
\]
satisfy
\begin{equation}
 D_{\rm rad}+BM=0,
 \label{eq:V60-general-compensation}
\end{equation}
then
\begin{equation}
 \dim Z\ge2,\qquad \operatorname{rank}B\ge2.
 \label{eq:V60-compensator-lower-bound}
\end{equation}
Thus two components, landing specifically in the
\((C_\perp,C_n)\) constraint plane, are both sufficient algebraically
and necessary by rank.
\end{theorem}

\begin{proof}
Substitute \eqref{eq:V60-compensated-incidence} into
\eqref{eq:V60-extended-constraint-map} and use
\eqref{eq:V60-Bianchi-factorization}:
\[
 \widetilde{\mathsf S}_{\cal C}(f,\zeta)
 =\sigma D_{\rm rad}\theta
  -\sigma E_{2n}K_{\rm B}\theta=0.
\]
Conversely, \eqref{eq:V60-general-compensation} implies
\[
 \operatorname{rank}(BM)
 =\operatorname{rank}D_{\rm rad}=2.
\]
Since \(\operatorname{rank}(BM)\le
\min\{\operatorname{rank}B,\dim Z\}\), both inequalities in
\eqref{eq:V60-compensator-lower-bound} follow.  The displayed
construction uses \(Z=\mathbb C^2\), \(B=E_{2n}\), and
\(M=-K_{\rm B}\), so the lower bound is attained.
\end{proof}

\begin{firewall}[Algebraic Bianchi target versus covariant action]
\label{fire:V60-compensator-scope}
Theorem~\ref{thm:V60-minimal-compensator} identifies the exact two
Euler covectors that a boundary stress or embedding sector must supply.
It does not construct that sector.  In particular,
\(\zeta=-K_{\rm B}\theta\) has not yet been derived from a common
covariant boundary action, and no nonlinear Noether identity has been
proved.  Declaring \(\zeta\) by hand would repeat the arbitrary-source
problem corrected in V58.  V61 must derive
\eqref{eq:V60-compensated-incidence} as one Euler/Noether incidence or
return an action-level obstruction.
\end{firewall}

\subsection{The gauge complement does not replace York compatibility}
\label{subsec:V60-compatibility}

The projector \(P_\epsilon\) supplies a constant-rank algebraic
complement on the frozen glancing fibre.  It does not imply the V58
domain conditions
\[
 X|_{\partial\Omega}=0,\qquad
 (\Box X)|_{\partial\Omega}=0,\qquad
 \Box X\in\dot H^1_0.
\]
The distinction already appears on a compactly supported pure-gauge
test.

\begin{proposition}[Explicit zero-trace failure of second compatibility]
\label{prop:V60-second-compatibility-test}
Let \(\xi>0\) be the half-space normal coordinate.  Choose
\(\chi\in C_c^\infty([0,\infty))\) with \(\chi=1\) near zero,
a nonzero
\(\psi\in C_c^\infty(\mathbb R_t\times\mathbb R^2_y)\), and a constant
spatial vector \(e\).  Put
\begin{equation}
 X(t,y,\xi):=\xi^2\chi(\xi)\psi(t,y)e,
 \qquad
 \widetilde X_0=0,
 \qquad
 \bar h:=\Gamma\widetilde X.
 \label{eq:V60-pure-gauge-test}
\end{equation}
Then
\begin{align}
 X|_{\xi=0}&=0,
 &
 D_nX|_{\xi=0}&=0,
 &
 \bar h|_{\xi=0}&=0,
 \label{eq:V60-zero-first-traces}\\
 (\Box X)|_{\xi=0}&=2\psi e\ne0,
 &
 {\cal C}(\bar h)|_{\xi=0}&=2\psi e\ne0.
 \label{eq:V60-nonzero-second-trace}
\end{align}
The spatial trace-free part of \(\bar h\) is \(\mathbb L X\), so the
Dirichlet York vector of \(\bar h\) is exactly \(X\).  Consequently the
essential configuration trace and every zero-order glancing complement
row may vanish while the second York compatibility fails.
\end{proposition}

\begin{proof}
The factor \(\xi^2\) makes \(X\) and all of its first derivatives vanish
at \(\xi=0\), proving the first three identities in
\eqref{eq:V60-zero-first-traces}.  Tangential and temporal terms in
\(\Box X\) retain the factor \(\xi^2\), while
\[
 \partial_\xi^2(\xi^2\chi(\xi))|_{\xi=0}=2.
\]
This proves the first identity in
\eqref{eq:V60-nonzero-second-trace}.  The gauge-complex identity
\eqref{eq:V58-complex-identity} proves the second.

The spatial trace-free part of
\(\Gamma\widetilde X\) is \(\mathbb L X\).  Substitution in the weak
York equation \eqref{eq:V58-York-weak} shows that \(X\) is its solution;
uniqueness in Theorem~\ref{thm:V58-York-projector} makes it the York
vector.  Since the full metric trace is zero at the boundary, the
essential trace and every algebraic row depending only on that trace
vanish there.  The nonzero value of \(\Box X\) proves the last claim.
\end{proof}

\begin{corollary}[Typed frozen candidate domain]
\label{cor:V60-typed-domain}
A valid metric lift must keep the intersections
\begin{equation}
 {\mathscr D}_{\rm cand}^m
 \subset
 {\mathscr D}_Y^m
 \cap\ker\!\left({\cal C}|_{\partial\Omega}\right)
 \cap\ker\!\left(G_\epsilon\operatorname{Tr}_{\partial}x\right)
 \cap\ker Q_{\rm TF}
 \label{eq:V60-candidate-domain}
\end{equation}
as typed domain conditions.  The last two algebraic rows do not imply
the first two.  At glancing,
\(Q_{\rm rad}\) maps the remaining trace kernel isomorphically to the
two V59 polarizations by Theorem~\ref{thm:V60-trace-lift}.  Its conormal
can be attached only after the rank-two compensation of
Theorem~\ref{thm:V60-minimal-compensator} is derived from an action.
\end{corollary}

\subsection{Consequence of the audit and the V61 acquisition order}
\label{subsec:V60-next}

The companion audit and the frozen calculation select one upstream
problem.  Adding the V59 traction directly to the old five metric rows
is impossible by Theorem~\ref{thm:V60-five-row-no-go}.  Enlarging the
old three-port bath or deleting its gauge column cannot change the
one-dimensional kernel of \({\mathsf S}_{\cal C}\).  The next admissible
step is therefore a boundary action with two explicit compensator
Euler coordinates.

V61 should seek fields \(\zeta_\perp,\zeta_n\) in the
metric--embedding boundary sector and a single quadratic frozen action
whose variation has the incidence
\eqref{eq:V60-compensated-incidence}.  It must then prove:

\begin{enumerate}
\item the linear boundary Noether identity
\({\mathsf S}_{\cal C}f+E_{2n}\zeta=0\) off shell;
\item equality of its TT metric virtual work with the V59 conormal
pairing;
\item a closed positive energy/domain after the new coordinates are
coupled to the V59 continuum;
\item preservation of the V58 second compatibility; and
\item constant rank not only at glancing but on the complete closed
stable cone.
\end{enumerate}

If no local quadratic metric--embedding action can realize the
invertible matrix \(K_{\rm B}\), V61 should print the corresponding
Euler/Noether rank obstruction and move further upstream to a boundary
Stueckelberg or edge-mode pair.  A positive power identity by itself is
not a substitute for these maps, exactly as the NS V31 audit warns.

\subsection{Finite certificates}
\label{subsec:V60-certificates}

In the ordered glancing basis
\((g_\epsilon,r_\times,r_+)\),
\begin{equation}
 \left[(G_\epsilon,Q_{\rm rad})\right]
 =
 \begin{pmatrix}
 1&0&0\\
 0&0&-\sqrt2\\
 0&\sqrt2&0
 \end{pmatrix},
 \qquad
 \det\left[(G_\epsilon,Q_{\rm rad})\right]=2.
 \label{eq:V60-trace-certificate}
\end{equation}
The conormal certificate is
\begin{equation}
 {\mathsf S}_{\cal C}C_{\rm rad}
 =
 \begin{pmatrix}
 0&0\\
 0&0\\
 0&1/\sqrt2\\
 -\sqrt2&0
 \end{pmatrix},
 \qquad
 \left({\mathsf S}_{\cal C}C_{\rm rad}\right)^*
 \left({\mathsf S}_{\cal C}C_{\rm rad}\right)
 =
 \begin{pmatrix}2&0\\0&1/2\end{pmatrix}.
 \label{eq:V60-conormal-certificate}
\end{equation}
Thus its singular values are \(\sqrt2\) and \(1/\sqrt2\).
The minimal compensation card is
\begin{equation}
 K_{\rm B}
 =
 \begin{pmatrix}0&1/\sqrt2\\-\sqrt2&0\end{pmatrix},
 \qquad
 \det K_{\rm B}=1,
 \qquad
 D_{\rm rad}-E_{2n}K_{\rm B}=0.
 \label{eq:V60-compensation-card}
\end{equation}
The last sign is written for the pair
\((f,\zeta)=(C_{\rm rad}\theta,-K_{\rm B}\theta)\) through
\(\widetilde{\mathsf S}_{\cal C}(f,\zeta)
=D_{\rm rad}\theta-E_{2n}K_{\rm B}\theta\).

\subsection{V60 result ledger and claim boundary}
\label{subsec:V60-ledger}

Relative to V1--V59, this note proves:

\begin{enumerate}
\item the scheduled NS/Yang--Mills audit supplies no missing Einstein
domain map and redirects the work to the exact frozen residual;
\item on both glancing sheets, the residual-gauge coordinate and the
two normalized TT traces give an explicit isomorphism of the
three-dimensional glancing kernel;
\item the TT configuration trace has a bounded right inverse of norm
\(\sqrt5\), and the gauge complement is an orthogonal constant-rank
projector on each glancing sheet;
\item the natural reciprocal TT conormal has an injective rank-two
harmonic source, supported exactly in the
\((C_\perp,C_n)\) plane;
\item no injective two-component conormal lift of any kind can take
values in the old five residual rows while all four harmonic rows stay
homogeneous;
\item two compensator Euler coordinates are necessary and sufficient
for algebraic cancellation of that source, with the exact invertible
matrix \(K_{\rm B}\); and
\item an explicit compactly supported pure-gauge field proves that the
algebraic glancing complement cannot replace the V58 second
compatibility condition.
\end{enumerate}

The V59 positive quotient generator remains valid and is now equipped
with an exact frozen configuration trace lift.  A full
metric--embedding conormal lift is not proved.  The compensator has not
been derived from a covariant action, the off-shell BRST domain is
still open, and neither the nonlinear Bianchi identity nor the coupled
Higgs, fermion, neutrino, and Yang--Mills sectors has been closed.  The
full Standard Model--Einstein continuum is not proved.

\subsection{Append-only record}

This V60 material was appended to the exact V59 parent whose SHA--256
digest is
\[
\texttt{14bc3ae24191429ab9e2f744c4b5c67cdd725f9b4ddcdf8a49b81ac13978bcd9}.
\]
The next research note is V61.  It should construct the rank-two
boundary Noether compensator from one metric--embedding action or prove
that this requires an additional boundary edge-mode pair.

% END APPENDED VERSION V60 MATERIAL
% BEGIN APPENDED VERSION V61 MATERIAL

\clearpage
\section{V61: pure-embedding no-go, a minimal edge trace action, and the
DeWitt duality--glancing tangency defect}
\label{sec:V61-edge-action}

\subsection{Purpose and nonduplication boundary}
\label{subsec:V61-purpose}

V60 solved the configuration trace problem on both gravitational
glancing sheets and proved a separate conormal obstruction.  The
radiative trace has the bounded right inverse \(R_{\rm rad}\), but every
injective two-component force in the old five metric residual rows
sources the harmonic constraint.  For the natural DeWitt-reciprocal
incidence, the source lies exactly in the
\((C_\perp,C_n)\) plane and factors as
\[
 {\mathsf S}_{\cal C}C_{\rm rad}=E_{2n}K_{\rm B},
 \qquad \det K_{\rm B}=1.
\]
V60 therefore identified two compensator Euler covectors as the next
action-level acquisition.

This note first tests the four embedding variables already constructed
in V39.  They cannot supply the compensator: a TT functional annihilates
the pure symmetric-gradient embedding directions, so its embedding
Euler covector is zero.  The note then introduces a genuinely new
two-component boundary edge variable and derives the V60 compensator
from one quadratic trace action.  The compensator is no longer declared
as an external source.

There remains a precise obstruction.  The unique DeWitt-dual shift
which turns the edge Ward identity into the full
\((C_\perp,C_n)\) residual identity is not tangent to the old glancing
kernel in the plus sector.  A second shift is tangent and gives the
same identity on the radiative Euler range, but not on an arbitrary
five-row residual.  This one-dimensional defect is printed explicitly.
Thus V61 constructs the minimal edge incidence and its positive trace
penalty, but does not claim a full metric--edge generator.

\subsection{Why the existing V39 embedding Euler row vanishes}
\label{subsec:V61-pure-embedding}

Let \({\cal S}_p\) be the V39 frozen symmetric-gradient map and
\(Q_4(h,X)=h+{\cal S}_pX\) its metric--embedding quotient.  On a
nonzero null covector, let \({\cal T}_p\) be the V53 spatial TT map,
with the V60 normalization \(Q_{\rm rad}\) on the reduced glancing
kernel.

\begin{lemma}[TT annihilates the pure embedding direction]
\label{lem:V61-TT-pure-embedding}
At every real nonzero null frequency, and in particular on either
gravitational glancing sheet,
\begin{equation}
 {\cal T}_p{\cal S}_p=0.
 \label{eq:V61-TT-Sp-zero}
\end{equation}
Consequently, if a differentiable frozen functional has the form
\begin{equation}
 F(h,X,b)=\Phi\!\left({\cal T}_pQ_4(h,X),b\right),
 \label{eq:V61-pure-pullback-functional}
\end{equation}
then its Euler covector in the existing four-component embedding
variable vanishes identically:
\begin{equation}
 E_X={\cal S}_p^*{\cal T}_p^*E_{\rm TT}=0.
 \label{eq:V61-existing-embedding-Euler-zero}
\end{equation}
\end{lemma}

\begin{proof}
Let \(\xi\) be the spatial propagation covector and let
\[
 P_\xi=I-\widehat\xi\widehat\xi^*
\]
be the transverse projector.  The spatial block of a pure symmetric
gradient is
\[
 ({\cal S}_pq)_{ij}=\xi_iq_j+\xi_jq_i
\]
up to the common Fourier convention.  Both indices in the TT
projection are multiplied by \(P_\xi\), and
\(P_\xi\xi=0\).  Its transverse trace is zero for the same reason.
Hence the transverse trace-free projection vanishes, proving
\eqref{eq:V61-TT-Sp-zero}.  The chain rule applied to
\eqref{eq:V61-pure-pullback-functional} then gives
\eqref{eq:V61-existing-embedding-Euler-zero}.
\end{proof}

\begin{theorem}[Pure-pullback embeddings cannot realize the V60
compensator]
\label{thm:V61-pure-embedding-no-go}
Let a nonzero TT conormal \(\theta\in\mathbb C^2\) act through a
functional of the form
\eqref{eq:V61-pure-pullback-functional}.  In the V41 equation-normalized
five-row residual, its metric force is, up to the common orientation
sign,
\begin{equation}
 f=C_{\rm rad}\theta,
 \label{eq:V61-pure-embedding-metric-force}
\end{equation}
and therefore
\begin{equation}
 {\mathsf S}_{\cal C}f
 =E_{2n}K_{\rm B}\theta\ne0.
 \label{eq:V61-pure-embedding-source}
\end{equation}
The existing V39 embedding Euler covector is zero by
\eqref{eq:V61-existing-embedding-Euler-zero}.  Hence no projection of
that Euler covector can cancel
\eqref{eq:V61-pure-embedding-source}.  A pure pullback through \(Q_4\)
does not supply either V60 compensator component.
\end{theorem}

\begin{proof}
The exact reduced DeWitt identity proved below in
\eqref{eq:V61-HC-Qstar} identifies the TT variational covector with
\(C_{\rm rad}\theta\).  V60 computed
\({\mathsf S}_{\cal C}C_{\rm rad}=E_{2n}K_{\rm B}\), and
\(K_{\rm B}\) is invertible, so the displayed source is nonzero for
every nonzero \(\theta\).  Lemma~\ref{lem:V61-TT-pure-embedding}
gives the zero embedding Euler covector.  A linear image of zero is
zero and cannot cancel an injective rank-two source.
\end{proof}

\begin{firewall}[Diffeomorphism Noether identity versus the harmonic
graph source]
\label{fire:V61-Noether-distinction}
V39 correctly proves
\(E_X={\cal S}_p^*E_h\) for a pure metric--embedding pullback.
For the TT coupling, both sides vanish in the embedding direction
because \({\cal T}_p{\cal S}_p=0\).  This diffeomorphism Noether
identity is not the V60 harmonic residual identity
\({\mathsf S}_{\cal C}f+E_{2n}\zeta=0\).  Conflating the two identities
would hide the nonzero matrix
\({\mathsf S}_{\cal C}C_{\rm rad}\).
\end{firewall}

\subsection{Two canonical lifts of the edge shift}
\label{subsec:V61-two-lifts}

Retain the reduced DeWitt matrix \(H\) from
\eqref{eq:V41-DeWitt-matrix}, in the aligned coordinate order
\((\tau,v_0,v_\parallel,v_\perp,\phi)\):
\begin{equation}
 H=
 \begin{pmatrix}
 -1/6&0&0&0&-1/2\\
 0&-2&0&0&0\\
 0&0&2&0&0\\
 0&0&0&2&0\\
 -1/2&0&0&0&1/2
 \end{pmatrix}.
 \label{eq:V61-H-matrix}
\end{equation}
Let \(E_{2n}\), \(K_{\rm B}\), \(Q_{\rm rad}\),
\(R_{\rm rad}\), and \({\mathsf S}_{\cal C}\) be the V60 maps.  Direct
multiplication gives the additional reciprocity identity
\begin{equation}
 HC_{\rm rad}=Q_{\rm rad}^*,
 \qquad
 C_{\rm rad}=H^{-1}Q_{\rm rad}^*.
 \label{eq:V61-HC-Qstar}
\end{equation}

The constraint-dual edge lift is
\begin{equation}
 B_{\cal C}
 :=H^{-1}{\mathsf S}_{\cal C}^*E_{2n}
 =
 \begin{pmatrix}
 0&0\\
 0&0\\
 0&0\\
 1/2&0\\
 0&1
 \end{pmatrix}.
 \label{eq:V61-constraint-dual-lift}
\end{equation}
It obeys
\begin{align}
 B_{\cal C}^*H&=E_{2n}^*{\mathsf S}_{\cal C},
 \label{eq:V61-constraint-duality}\\
 Q_{\rm rad}B_{\cal C}&=K_{\rm B}^*.
 \label{eq:V61-constraint-lift-trace}
\end{align}

There is also a lift tangent to the V60 radiative glancing plane:
\begin{equation}
 B_{\rm T}
 :=R_{\rm rad}K_{\rm B}^*
 =
 \begin{pmatrix}
 0&-3\\
 0&0\\
 0&0\\
 1/2&0\\
 0&1
 \end{pmatrix}.
 \label{eq:V61-tangent-lift}
\end{equation}
It satisfies
\begin{equation}
 Q_{\rm rad}B_{\rm T}=K_{\rm B}^*,
 \qquad
 \operatorname{ran}B_{\rm T}
 \subset\operatorname{span}\{r_\times,r_+\}
 \subset{\mathscr K}_\epsilon
 \quad(\epsilon=\pm1).
 \label{eq:V61-tangent-properties}
\end{equation}
Their difference is the \(Q_{\rm rad}\)-null map
\begin{equation}
 N:=B_{\rm T}-B_{\cal C}
 =
 \begin{pmatrix}
 0&-3\\
 0&0\\
 0&0\\
 0&0\\
 0&0
 \end{pmatrix},
 \qquad Q_{\rm rad}N=0.
 \label{eq:V61-lift-difference}
\end{equation}

\begin{theorem}[DeWitt duality and glancing tangency cannot hold
simultaneously]
\label{thm:V61-duality-tangency-no-go}
There is no map \(B:\mathbb C^2\to\mathbb C^5\) satisfying both
\begin{align}
 B^*H&=E_{2n}^*{\mathsf S}_{\cal C},
 \label{eq:V61-required-duality}\\
 \operatorname{ran}B&\subset{\mathscr K}_\epsilon
 \label{eq:V61-required-tangency}
\end{align}
on either glancing sheet.  The obstruction has rank one and occurs only
in the scalar/plus column.  More precisely,
\begin{equation}
 \Delta
 :=B_{\rm T}^*H-E_{2n}^*{\mathsf S}_{\cal C}
 =N^*H
 =
 \begin{pmatrix}
 0&0&0&0&0\\
 1/2&0&0&0&3/2
 \end{pmatrix}.
 \label{eq:V61-duality-defect}
\end{equation}
Although \(\Delta\ne0\), it annihilates the radiative Euler range:
\begin{equation}
 \Delta C_{\rm rad}=0.
 \label{eq:V61-defect-annihilates-radiative-range}
\end{equation}
\end{theorem}

\begin{proof}
Because \(H\) is nondegenerate,
\eqref{eq:V61-required-duality} is equivalent, after taking adjoints,
to
\[
 HB={\mathsf S}_{\cal C}^*E_{2n}.
\]
It has the unique solution \(B=B_{\cal C}\).  The first column of
\(B_{\cal C}\) is \(r_\times/2\) and lies in
\({\mathscr K}_\epsilon\).  Its second column is
\((0,0,0,0,1)^T\), which violates the glancing relation
\(\tau=-3\phi\).  Thus the unique dual lift is not tangent.

The matrix \(B_{\rm T}\) replaces that second column by
\(r_+=(-3,0,0,0,1)^T\), so it is tangent by V60.  Subtracting the two
matrices gives \(N\); multiplication by \(H\) gives
\eqref{eq:V61-duality-defect}.  Finally, the plus column of
\(C_{\rm rad}\) obeys \(f_\tau+3f_\phi=0\), while its cross column has
both entries zero.  Hence \(\Delta C_{\rm rad}=0\).
\end{proof}

\begin{corollary}[The tangent lift has the correct Ward identity on the
TT coupling range]
\label{cor:V61-tangent-range-Ward}
For every \(\theta\in\mathbb C^2\),
\begin{equation}
 B_{\rm T}^*H C_{\rm rad}\theta
 =
 E_{2n}^*{\mathsf S}_{\cal C}C_{\rm rad}\theta
 =K_{\rm B}\theta.
 \label{eq:V61-tangent-range-identity}
\end{equation}
Thus the tangent and constraint-dual lifts give the same compensator on
the two-dimensional TT Euler range, although only \(B_{\cal C}\)
represents the constraint covector on an arbitrary five-row residual.
\end{corollary}

\subsection{A single edge trace action}
\label{subsec:V61-trace-action}

Introduce a new edge variable
\[
 e=(e_\perp,e_n)^T\in\mathbb C^2
\]
and write \(z\in\mathbb C^2\) for the continuum boundary trace.  Define
the invariant mismatch
\begin{equation}
 w(x,e,z)
 :=z-Q_{\rm rad}x+K_{\rm B}^*e.
 \label{eq:V61-invariant-mismatch}
\end{equation}
For a two-component multiplier \(\vartheta\), let
\begin{equation}
 {\cal F}_{\rm tr}[x,e,z,\vartheta]
 :=
 \int_{\partial M}
 \operatorname{Re}\langle\vartheta,
 w(x,e,z)\rangle_{\mathbb C^2}\,\dd^3y.
 \label{eq:V61-trace-action}
\end{equation}
All statements below may be read fibrewise after tangential
Laplace--Fourier transformation.

\begin{theorem}[One action produces the metric force and both
compensator covectors]
\label{thm:V61-trace-action-incidence}
The Euler covectors of \eqref{eq:V61-trace-action}, on the realification
of the complex fibre, are
\begin{align}
 E_x&=-Q_{\rm rad}^*\vartheta,
 &
 E_e&=K_{\rm B}\vartheta,
 \label{eq:V61-trace-Euler-xe}\\
 E_z&=\vartheta,
 &
 E_\vartheta&=w.
 \label{eq:V61-trace-Euler-ztheta}
\end{align}
If the metric residual is normalized by
\begin{equation}
 f:=H^{-1}E_x,
 \qquad
 \zeta:=E_e,
 \label{eq:V61-equation-normalization}
\end{equation}
then
\begin{align}
 f&=-C_{\rm rad}\vartheta,
 &
 \zeta&=K_{\rm B}\vartheta,
 \label{eq:V61-action-compensated-pair}\\
 {\mathsf S}_{\cal C}f+E_{2n}\zeta&=0.
 \label{eq:V61-action-constraint-identity}
\end{align}
Thus \({\cal F}_{\rm tr}\) realizes exactly the
\(\sigma=-1\) incidence of
\eqref{eq:V60-compensated-incidence}.  The common-trace equation is
\begin{equation}
 z=Q_{\rm rad}x-K_{\rm B}^*e.
 \label{eq:V61-edge-common-trace}
\end{equation}
In the edge gauge \(e=0\), it is the V59 trace
\(z=Q_{\rm rad}x\).
\end{theorem}

\begin{proof}
Differentiate the real bilinear pairing in
\eqref{eq:V61-trace-action}.  The four derivatives are exactly
\eqref{eq:V61-trace-Euler-xe}--%
\eqref{eq:V61-trace-Euler-ztheta}.  Equation
\eqref{eq:V61-HC-Qstar} gives
\(f=-C_{\rm rad}\vartheta\).  Therefore V60's factorization gives
\[
 {\mathsf S}_{\cal C}f+E_{2n}\zeta
 =-E_{2n}K_{\rm B}\vartheta
  +E_{2n}K_{\rm B}\vartheta=0.
\]
Variation of \(\vartheta\) gives
\eqref{eq:V61-edge-common-trace}; setting \(e=0\) gives the last claim.
\end{proof}

\begin{theorem}[Exact edge Ward identities]
\label{thm:V61-edge-Ward}
For either \(B=B_{\cal C}\) or \(B=B_{\rm T}\), the transformation
\begin{equation}
 \delta_\eta x=B\eta,\qquad
 \delta_\eta e=\eta,\qquad
 \delta_\eta z=\delta_\eta\vartheta=0
 \label{eq:V61-edge-shift}
\end{equation}
leaves \(w\) and \({\cal F}_{\rm tr}\) invariant.  Its Euler identity is
\begin{equation}
 B^*E_x+E_e=0.
 \label{eq:V61-edge-Ward-general}
\end{equation}
For \(B=B_{\cal C}\), this is identically
\begin{equation}
 E_{2n}^*{\mathsf S}_{\cal C}f+\zeta=0
 \label{eq:V61-edge-Ward-constraint}
\end{equation}
for the equation normalization \(E_x=Hf\).  For \(B=B_{\rm T}\), the
same formula holds on the Euler range of
\({\cal F}_{\rm tr}\) by
\eqref{eq:V61-defect-annihilates-radiative-range}.
\end{theorem}

\begin{proof}
Both lifts obey \(Q_{\rm rad}B=K_{\rm B}^*\).  Hence
\[
 \delta_\eta w
 =-Q_{\rm rad}B\eta+K_{\rm B}^*\eta=0.
\]
Differentiating the invariant functional in \(\eta\) gives
\eqref{eq:V61-edge-Ward-general}.  For the constraint-dual lift, use
\(B_{\cal C}^*H=E_{2n}^*{\mathsf S}_{\cal C}\).  For the tangent lift,
subtract the two identities.  Their difference is
\(\Delta f\), which vanishes for
\(f=-C_{\rm rad}\vartheta\) by
\eqref{eq:V61-defect-annihilates-radiative-range}.
\end{proof}

The same mismatch has the nonnegative penalty
\begin{equation}
 {\cal F}_{\mu}^{+}[x,e,z]
 :=
 \frac{\mu}{2}\int_{\partial M}|w(x,e,z)|^2\,\dd^3y,
 \qquad \mu>0.
 \label{eq:V61-positive-penalty}
\end{equation}

\begin{proposition}[Positive trace storage and closed hard-trace domain]
\label{prop:V61-positive-trace-storage}
The Hessian of \({\cal F}_{\mu}^{+}\) is nonnegative and has fibre rank
two.  Its Euler formulas are
\eqref{eq:V61-trace-Euler-xe}--%
\eqref{eq:V61-trace-Euler-ztheta} with
\(\vartheta=\mu w\), except that there is no independent
\(E_\vartheta\).  It obeys both edge Ward identities and the exact
constraint identity \eqref{eq:V61-action-constraint-identity}.

At the Sobolev trace level, if
\(x,z,e\in H^{1/2}(\partial M;\mathbb C^2)\) in their corresponding
finite fibres, then
\begin{equation}
 {\cal D}_{\rm tr}
 :=\{(x,e,z):w(x,e,z)=0\}
 \label{eq:V61-hard-trace-domain}
\end{equation}
is a closed linear subspace.  Its edge gauge \(e=0\) is isomorphic to
the V59 common-trace domain.
\end{proposition}

\begin{proof}
The mismatch map is onto because its coefficient on \(z\) is \(I_2\).
Thus its Gram Hessian has rank two and is nonnegative.  Substitution
\(\vartheta=\mu w\) in the preceding variational calculation proves the
Euler and Ward identities.  The trace maps are continuous, and
\(Q_{\rm rad}\) and \(K_{\rm B}\) are bounded finite matrices.
Therefore \(w\) is continuous on the product trace space and its kernel
is closed.  If \(e=0\), the equation \(w=0\) is
\(z=Q_{\rm rad}x\), the V59 common trace.
\end{proof}

\begin{firewall}[The trace action is not yet a total edge dynamics]
\label{fire:V61-trace-action-scope}
The multiplier action gives the exact hard trace and derives
\(\zeta\) from the same functional as the metric force.  The positive
penalty gives a nonnegative boundary storage term.  Because
\(Q_{\rm rad}\) is the direction-dependent TT/York multiplier, these are
frozen Fourier-fibre functionals and generally pseudodifferential in
the tangential variables; no local covariant density is claimed.
Neither construction yet supplies a kinetic or BRST quartet action for
\(e\).  If \(e\) were
varied as an ordinary unconstrained physical field with no additional
term, its total Euler equation would force its compensator covector to
vanish.  If it is treated as a gauge coordinate, the shift
\eqref{eq:V61-edge-shift} must be extended to the entire metric--edge
action and its ghost domain.  That total invariance and maximal
generator domain are not proved in V61.
\end{firewall}

\subsection{The minimal glancing edge exact sequence}
\label{subsec:V61-edge-exact-sequence}

The tangent lift gives a precise frozen quotient which retains exactly
the two V59 polarizations.  On either glancing sheet define
\begin{equation}
 {\cal Q}_{E,\epsilon}:
 {\mathscr K}_\epsilon\oplus\mathbb C^2
 \longrightarrow\mathbb C^2,
 \qquad
 {\cal Q}_{E,\epsilon}(x,e)
 :=Q_{\rm rad}x-K_{\rm B}^*e,
 \label{eq:V61-edge-quotient-map}
\end{equation}
and
\begin{equation}
 {\cal G}_{E,\epsilon}(\alpha,\eta)
 :=
 \left(\alpha g_\epsilon+B_{\rm T}\eta,\eta\right).
 \label{eq:V61-edge-gauge-map}
\end{equation}

\begin{theorem}[Split exact minimal edge completion at glancing]
\label{thm:V61-edge-exact-sequence}
For each \(\epsilon=\pm1\),
\begin{equation}
 0\longrightarrow\mathbb C\oplus\mathbb C^2
 \mathop{\longrightarrow}^{{\cal G}_{E,\epsilon}}
 {\mathscr K}_\epsilon\oplus\mathbb C^2
 \mathop{\longrightarrow}^{{\cal Q}_{E,\epsilon}}
 \mathbb C^2
 \longrightarrow0
 \label{eq:V61-edge-exact-sequence}
\end{equation}
is split exact.  A splitting is
\begin{equation}
 a\longmapsto(R_{\rm rad}a,0).
 \label{eq:V61-edge-splitting}
\end{equation}
Thus the new edge pair introduces two shift directions but no additional
physical glancing polarization.  The quotient is canonically the V59
radiative fibre.
\end{theorem}

\begin{proof}
The identities
\(Q_{\rm rad}g_\epsilon=0\) and
\(Q_{\rm rad}B_{\rm T}=K_{\rm B}^*\) give
\({\cal Q}_{E,\epsilon}{\cal G}_{E,\epsilon}=0\).
The edge component of
\({\cal G}_{E,\epsilon}(\alpha,\eta)\) is \(\eta\), so a zero image first
gives \(\eta=0\), and then
\(\alpha g_\epsilon=0\) gives \(\alpha=0\).  Hence the first map is
injective.

The V60 identity \(Q_{\rm rad}R_{\rm rad}=I_2\) proves surjectivity and
\eqref{eq:V61-edge-splitting}.  Conversely, suppose
\({\cal Q}_{E,\epsilon}(x,e)=0\).  V60 decomposes
\[
 x=g_\epsilon G_\epsilon x+
   R_{\rm rad}Q_{\rm rad}x
  =g_\epsilon G_\epsilon x+
   R_{\rm rad}K_{\rm B}^*e
  =g_\epsilon G_\epsilon x+B_{\rm T}e.
\]
Therefore
\((x,e)={\cal G}_{E,\epsilon}(G_\epsilon x,e)\), proving exactness.
The quotient dimension is \(5-3=2\), and the splitting identifies it
with the radiative fibre.
\end{proof}

\begin{corollary}[Positive V59 energy on the edge quotient]
\label{cor:V61-edge-positive-quotient}
Pull back the V59 radiative fibre norm by
\({\cal Q}_{E,\epsilon}\).  Its radical is exactly
\(\operatorname{ran}{\cal G}_{E,\epsilon}\), and it descends to the
original positive two-polarization norm on the quotient.  Consequently
the minimal edge completion creates no negative or zero-norm physical
glancing mode.
\end{corollary}

\begin{proof}
The radical of a pulled-back positive norm is the kernel of the
surjective pullback map.  Theorem~\ref{thm:V61-edge-exact-sequence}
identifies that kernel with
\(\operatorname{ran}{\cal G}_{E,\epsilon}\).  The splitting is a right
inverse, so the induced quotient norm is exactly the V59 norm.
\end{proof}

\begin{firewall}[Frozen quotient versus a spacetime edge bundle]
\label{fire:V61-edge-quotient-scope}
Theorem~\ref{thm:V61-edge-exact-sequence} is a complete
finite-dimensional theorem on the two glancing sheets.  It does not
construct a local edge bundle over the whole closed stable cone, nor a
spacetime differential operator whose symbol is \(B_{\rm T}\).
The V59 self-adjoint generator acts on the radiative quotient; pulling
back its fibre norm does not by itself make a closed generator on raw
metric and edge representatives.
\end{firewall}

\subsection{V62 acquisition order}
\label{subsec:V61-next}

V61 changes the next question from existence of a compensator incidence
to extension of its symmetry and domain.  The required work is now:

\begin{enumerate}
\item extend the tangent lift \(B_{\rm T}\) away from glancing and test
constant rank on the complete closed stable cone;
\item decide whether the scalar duality defect
\[
 \Delta f=\binom{0}{(f_\tau+3f_\phi)/2}
\]
is supplied by the existing mean-curvature/normal-embedding Euler row
or requires one additional scalar edge coordinate;
\item construct the edge ghost, antighost, and multiplier rows so that
the shift Ward identity is off-shell tangent to the total domain;
\item combine the hard trace action with the V59 bulk form without
forcing \(\vartheta=0\), and prove closedness and maximality; and
\item compare the resulting full stable symbol with the V59 zero-free
factor \(\lambda+m\kappa_v\).
\end{enumerate}

If the scalar defect cannot be absorbed by the geometric
mean-curvature polarization, V62 should prove the corresponding
three-edge lower bound before adding a field.  The nonlinear
Brown--York/Bianchi identity remains downstream of this frozen
action-domain problem.

\subsection{Finite certificates}
\label{subsec:V61-certificates}

The action reciprocity card is
\begin{equation}
 HC_{\rm rad}-Q_{\rm rad}^*=0,
 \qquad
 B_{\cal C}^*H-E_{2n}^*{\mathsf S}_{\cal C}=0,
 \qquad
 Q_{\rm rad}B_{\cal C}-K_{\rm B}^*=0.
 \label{eq:V61-reciprocity-card}
\end{equation}
The tangency card is
\begin{equation}
 B_{\rm T}=R_{\rm rad}K_{\rm B}^*,
 \qquad
 Q_{\rm rad}B_{\rm T}=K_{\rm B}^*,
 \qquad
 \operatorname{rank}B_{\rm T}=2.
 \label{eq:V61-tangency-card}
\end{equation}
The scalar defect card is
\begin{equation}
 \Delta
 =
 \begin{pmatrix}
 0&0&0&0&0\\
 1/2&0&0&0&3/2
 \end{pmatrix},
 \qquad
 \operatorname{rank}\Delta=1,
 \qquad
 \Delta C_{\rm rad}=0.
 \label{eq:V61-defect-card}
\end{equation}
For the trace action, the residual card is
\begin{equation}
 r_{\rm tr}:=z-Q_{\rm rad}x+K_{\rm B}^*e,
 \qquad
 r_{\cal C}:={\mathsf S}_{\cal C}f+E_{2n}\zeta,
 \qquad
 r_{\cal C}=0.
 \label{eq:V61-action-card}
\end{equation}

\subsection{V61 result ledger and claim boundary}
\label{subsec:V61-ledger}

Relative to V1--V60, this note proves:

\begin{enumerate}
\item TT projection annihilates the existing V39 pure embedding
directions, so their Euler covector cannot supply the rank-two V60
compensator;
\item the exact DeWitt identity
\(HC_{\rm rad}=Q_{\rm rad}^*\) converts TT virtual work into the V60
equation-normalized metric incidence;
\item the constraint-dual shift \(B_{\cal C}\) and glancing-tangent shift
\(B_{\rm T}\) both induce \(K_{\rm B}^*\) on the radiative trace;
\item no two-component shift is simultaneously the full DeWitt dual of
\((C_\perp,C_n)\) and tangent to the old glancing kernel; the difference
is the explicit rank-one scalar row \(\Delta\);
\item the scalar defect vanishes identically on the TT Euler range;
\item one quadratic multiplier action derives the metric force, both
edge compensator covectors, the hard common trace, and the exact V60
constraint cancellation;
\item a nonnegative penalty realizes the same Ward identity and its
hard-trace zero set is a closed Sobolev trace domain; and
\item the new two-component edge pair fits into a split exact glancing
sequence whose physical quotient is exactly the positive V59
two-polarization fibre.
\end{enumerate}

A genuine action-level source for the two V60 compensators has now been
constructed, but only in a frozen boundary edge completion.  The shift
is not yet a BRST symmetry of the full metric--edge bulk action, the
rank-one scalar defect has not been geometrically absorbed, and no
closed raw-representative generator or complete-cone stable inverse is
proved.  The nonlinear stress/Bianchi identity and the coupled Higgs,
fermion, neutrino, and Yang--Mills sectors remain open.  The full
Standard Model--Einstein continuum is not proved.

\subsection{Append-only record}

This V61 material was appended to the exact V60 parent whose SHA--256
digest is
\[
\texttt{3722934dd4516002b40b4c265e7a619fb234b0d444eb0a015874276205094a90}.
\]
The next research note is V62.  It should resolve the rank-one
duality--tangency defect and extend the edge symmetry beyond the frozen
glancing fibre before attempting the total generator.

% END APPENDED VERSION V61 MATERIAL
% BEGIN APPENDED VERSION V62 MATERIAL

\clearpage
\section{V62: canonical angular TT trace, rotating two-edge Ward identity,
and the normal-incidence essential-row obstruction}
\label{sec:V62-angular-TT}

\subsection{Purpose and correction of the next acquisition}
\label{subsec:V62-purpose}

V61 constructed a two-component edge trace action and found a rank-one
difference between two extensions of its edge shift.  The
constraint-dual lift \(B_{\cal C}\) represents the
\((C_\perp,C_n)\) covector on every five-row residual but is not tangent
to the old plus glancing direction.  The tangent lift \(B_{\rm T}\)
preserves the glancing kernel and gives the same Ward identity on the
TT Euler range.  Their difference is
\[
 \Delta f=\binom{0}{(f_\tau+3f_\phi)/2},
 \qquad
 \Delta C_{\rm rad}=0.
\]

It would be premature to add a third edge field for this scalar.
The quotient trace \(Q_{\rm rad}\) in V60--V61 was defined only by its
values on the glancing kernel.  Its extension to all five reduced
configuration coordinates was a convenient choice, not the canonical
spatial TT projection.  V62 computes the canonical extension and shows
that the scalar difference is precisely an off-kernel extension
ambiguity.  It vanishes on the physical coupling range and supplies no
third independent traction.

The calculation then continues along the complete real propagating
angular cone.  The two-edge Ward identity extends exactly at every
oblique angle, but its metric trace and constraint incidence both
vanish at normal incidence.  This rank loss is not a defect of the edge
factorization: the V37 conformal essential row has removed the induced
TT boundary components.  Consequently the full V59 radiative continuum
cannot be lifted uniformly while that essential row is kept
Dirichlet.  V63 must move upstream to a mixed induced-TT boundary
polarization.

\subsection{The canonical TT extension of the glancing trace}
\label{subsec:V62-canonical-glancing-extension}

At plus glancing, the spatial block of a reduced V37 configuration is
\begin{equation}
 h^{\rm sp}(x)=
 \begin{pmatrix}
 \tau/3&0&v_\parallel\\
 0&\tau/3&v_\perp\\
 v_\parallel&v_\perp&\phi
 \end{pmatrix}.
 \label{eq:V62-reduced-spatial-block}
\end{equation}
The propagation direction is \(e_\parallel\), and use the V58
transverse basis \((e_\perp,n)\).  Direct contraction with
\(E_+,E_\times\) gives the canonical full-fibre TT map
\begin{equation}
 Q_{\rm g}^{\rm TT}x
 :=
 \binom{(\tau/3-\phi)/\sqrt2}{\sqrt2\,v_\perp}.
 \label{eq:V62-canonical-glancing-map}
\end{equation}
On the glancing kernel \(\tau=-3\phi\), this is exactly
\(Q_{\rm rad}x=(-\sqrt2\phi,\sqrt2v_\perp)^T\).

The two extensions differ on the full five-dimensional fibre by
\begin{equation}
 Q_{\rm g}^{\rm TT}-Q_{\rm rad}
 =
 \frac1{3\sqrt2}
 \begin{pmatrix}
 1&0&0&0&3\\
 0&0&0&0&0
 \end{pmatrix}.
 \label{eq:V62-extension-difference}
\end{equation}

\begin{proposition}[The V61 scalar is an off-kernel extension row]
\label{prop:V62-extension-row}
The difference
\eqref{eq:V62-extension-difference} vanishes on
\({\mathscr K}_\epsilon\) for both glancing sheets.  Hence
\(Q_{\rm g}^{\rm TT}\) and \(Q_{\rm rad}\) define the same radiative
quotient trace, the same bounded splitting \(R_{\rm rad}\), and the
same virtual work on variations tangent to
\({\mathscr K}_\epsilon\).

The nonzero row is proportional to the annihilator
\[
 \chi(x):=\tau+3\phi
\]
of the glancing kernel.  Therefore changing between the two extensions
can change an equation-normalized off-kernel force and its DeWitt dual
without changing any physical glancing trace.
\end{proposition}

\begin{proof}
Every vector in \({\mathscr K}_\epsilon\) obeys
\(\tau+3\phi=0\) by
\eqref{eq:V60-glancing-kernel}.  Substitution in
\eqref{eq:V62-extension-difference} proves equality of the two maps on
that kernel.  The V60 splitting takes values in the kernel, so both
maps compose with \(R_{\rm rad}\) to \(I_2\).  Covectors which differ
by a multiple of \(\chi\) have equal pairing with every tangent
variation.  The last statement follows because the DeWitt
normalization applies \(H^{-1}\) to the chosen full-fibre covector.
\end{proof}

\begin{corollary}[No third edge follows from the rank-one defect]
\label{cor:V62-no-third-edge}
The V61 identity \(\Delta C_{\rm rad}=0\) shows that its rank-one
duality--tangency difference annihilates every Euler covector generated
by the two-polarization trace action.  Together with
Proposition~\ref{prop:V62-extension-row}, this proves that V61 supplied
no lower bound for a third physical edge component.  Two edge
components remain necessary by V60 and sufficient for the
coupling-specific Ward identity by V61.

This does not make \(\Delta\) vanish on an arbitrary five-row residual.
It says only that adding a third field to cancel a covector which is
zero on the actual TT Euler range would not be justified.
\end{corollary}

\subsection{TT projection on the real propagating angular cone}
\label{subsec:V62-angular-trace}

Let
\begin{equation}
 \widehat\xi
 =(\mathfrak s,0,\mathfrak c),
 \qquad
 \mathfrak s:=\sin\theta,\quad
 \mathfrak c:=\cos\theta,\quad
 0\le\theta\le\frac{\pi}{2},
 \label{eq:V62-angular-direction}
\end{equation}
where \(\theta=\pi/2\) is gravitational glancing and \(\theta=0\)
is normal incidence.  Choose the oriented transverse frame
\begin{equation}
 e=(0,1,0),
 \qquad
 t=(-\mathfrak c,0,\mathfrak s),
 \label{eq:V62-transverse-frame}
\end{equation}
and
\begin{equation}
 E_+(\theta):=\frac{e\otimes e-t\otimes t}{\sqrt2},
 \qquad
 E_\times(\theta):=
 \frac{e\otimes t+t\otimes e}{\sqrt2}.
 \label{eq:V62-angular-TT-basis}
\end{equation}
The sign of \(E_\times\) is chosen to agree with V58 at glancing; it is
the opposite of the optional cross-frame sign used in the V53 port
display.

Contracting \eqref{eq:V62-reduced-spatial-block} with this basis defines
\begin{equation}
 Q_\theta^{\rm TT}x
 :=
 \begin{pmatrix}
 \dfrac1{\sqrt2}
 \left[
 \mathfrak s^2\left(\dfrac{\tau}{3}-\phi\right)
 +2\mathfrak s\mathfrak c\,v_\parallel
 \right]\\[7pt]
 \sqrt2\,\mathfrak s\,v_\perp
 \end{pmatrix}.
 \label{eq:V62-angular-TT-map}
\end{equation}

\begin{theorem}[Exact angular TT trace and singular values]
\label{thm:V62-angular-trace}
The map \(Q_\theta^{\rm TT}:\mathbb C^5\to\mathbb C^2\) is the
spatial TT projection of every reduced V37 configuration
\eqref{eq:V62-reduced-spatial-block}.  Its two rows are orthogonal and
\begin{equation}
 Q_\theta^{\rm TT}(Q_\theta^{\rm TT})^*
 =
 \begin{pmatrix}
 \mathfrak s^2
 \left(2-\dfrac{13}{9}\mathfrak s^2\right)&0\\
 0&2\mathfrak s^2
 \end{pmatrix}.
 \label{eq:V62-angular-trace-Gram}
\end{equation}
Consequently,
\begin{align}
 \operatorname{rank}Q_\theta^{\rm TT}
 &=2
 &&(0<\theta\le\pi/2),
 \label{eq:V62-oblique-trace-rank}\\
 Q_0^{\rm TT}&=0,
 \label{eq:V62-normal-trace-zero}\\
 s_{\min}(Q_\theta^{\rm TT})
 &=
 \mathfrak s
 \sqrt{2-\frac{13}{9}\mathfrak s^2}.
 \label{eq:V62-angular-smallest-singular}
\end{align}
Every right inverse \(R_\theta\) for \(0<\theta\le\pi/2\) satisfies
\begin{equation}
 \|R_\theta\|
 \ge
 \frac1{\mathfrak s
 \sqrt{2-\frac{13}{9}\mathfrak s^2}}
 \longrightarrow\infty
 \qquad(\theta\downarrow0).
 \label{eq:V62-right-inverse-blowup}
\end{equation}
\end{theorem}

\begin{proof}
The \(t\)-\(t\) component of
\eqref{eq:V62-reduced-spatial-block} is
\[
 h_{tt}
 =\mathfrak c^2\frac{\tau}{3}
  +\mathfrak s^2\phi
  -2\mathfrak s\mathfrak c\,v_\parallel.
\]
Hence
\[
 \langle E_+,h^{\rm sp}\rangle
 =\frac{h_{ee}-h_{tt}}{\sqrt2}
 =\frac{
 \mathfrak s^2(\tau/3-\phi)
 +2\mathfrak s\mathfrak c\,v_\parallel}{\sqrt2}.
\]
Also
\[
 \langle E_\times,h^{\rm sp}\rangle
 =\sqrt2\,h_{et}
 =\sqrt2\,\mathfrak s\,v_\perp.
\]
This proves \eqref{eq:V62-angular-TT-map}.

The rows have disjoint coordinate support.  The squared norm of the
plus row is
\[
 \frac{\mathfrak s^4}{18}
 +2\mathfrak s^2\mathfrak c^2
 +\frac{\mathfrak s^4}{2}
 =
 \mathfrak s^2
 \left(2-\frac{13}{9}\mathfrak s^2\right),
\]
and the cross row has squared norm \(2\mathfrak s^2\).  Since
\(0\le\mathfrak s\le1\), the first is the smaller nonzero eigenvalue.
This proves the Gram, rank, and singular-value formulas.  For a
surjective matrix, every right inverse has norm at least the reciprocal
of its least singular value, giving
\eqref{eq:V62-right-inverse-blowup}.
\end{proof}

\subsection{DeWitt incidence and the rotating constraint plane}
\label{subsec:V62-angular-incidence}

Define the equation-normalized angular metric incidence by
\begin{equation}
 C_\theta^{\rm TT}
 :=H^{-1}(Q_\theta^{\rm TT})^*.
 \label{eq:V62-angular-incidence-definition}
\end{equation}
Direct inversion of the V41 scalar DeWitt block gives
\begin{equation}
 C_\theta^{\rm TT}
 =
 \frac1{\sqrt2}
 \begin{pmatrix}
 \mathfrak s^2&0\\
 0&0\\
 \mathfrak s\mathfrak c&0\\
 0&\mathfrak s\\
 -\mathfrak s^2&0
 \end{pmatrix},
 \qquad
 HC_\theta^{\rm TT}=(Q_\theta^{\rm TT})^*.
 \label{eq:V62-angular-incidence}
\end{equation}
Its harmonic source is
\begin{equation}
 D_\theta^{\rm TT}
 :={\mathsf S}_{\cal C}C_\theta^{\rm TT}
 =
 \frac1{\sqrt2}
 \begin{pmatrix}
 0&0\\
 \mathfrak s\mathfrak c&0\\
 0&\mathfrak s\\
 -\mathfrak s^2&0
 \end{pmatrix},
 \label{eq:V62-angular-constraint-source}
\end{equation}
in the row order
\((C_0,C_\parallel,C_\perp,C_n)\).

Define the isometric moving inclusion
\begin{equation}
 E_\theta
 :=
 \begin{pmatrix}
 0&0\\
 0&\mathfrak c\\
 1&0\\
 0&-\mathfrak s
 \end{pmatrix},
 \qquad
 E_\theta^*E_\theta=I_2,
 \label{eq:V62-moving-constraint-plane}
\end{equation}
and
\begin{equation}
 K_\theta
 :=
 \frac{\mathfrak s}{\sqrt2}
 \begin{pmatrix}0&1\\1&0\end{pmatrix}.
 \label{eq:V62-angular-compensator}
\end{equation}

\begin{theorem}[Exact two-channel angular factorization]
\label{thm:V62-angular-factorization}
For every \(0\le\theta\le\pi/2\),
\begin{equation}
 D_\theta^{\rm TT}=E_\theta K_\theta.
 \label{eq:V62-angular-factorization}
\end{equation}
For \(0<\theta\le\pi/2\), both
\(C_\theta^{\rm TT}\) and \(D_\theta^{\rm TT}\) have rank two, and
\begin{equation}
 \det K_\theta=-\frac{\mathfrak s^2}{2}.
 \label{eq:V62-angular-K-determinant}
\end{equation}
The cross polarization sources \(C_\perp\).  The plus polarization
sources the unit constraint direction
\begin{equation}
 \mathfrak c\,C_\parallel-\mathfrak s\,C_n.
 \label{eq:V62-plus-constraint-direction}
\end{equation}
At glancing this is \(-C_n\); at normal incidence both source
amplitudes vanish.
\end{theorem}

\begin{proof}
Multiplication of
\eqref{eq:V62-moving-constraint-plane} and
\eqref{eq:V62-angular-compensator} gives
\eqref{eq:V62-angular-constraint-source}.  The columns of
\(E_\theta\) are orthonormal.  Thus
\(\operatorname{rank}D_\theta^{\rm TT}
=\operatorname{rank}K_\theta\), which is two for
\(\mathfrak s>0\) and zero for \(\mathfrak s=0\).  The two nonzero
columns of \(C_\theta^{\rm TT}\) have disjoint support, so it has the
same oblique rank.  The determinant and polarization statements are
read directly from the two matrices.
\end{proof}

The constraint-dual angular edge lift is
\begin{equation}
 B_{{\cal C},\theta}
 :=H^{-1}{\mathsf S}_{\cal C}^*E_\theta
 =
 \begin{pmatrix}
 0&0\\
 0&0\\
 0&\mathfrak c/2\\
 1/2&0\\
 0&-\mathfrak s
 \end{pmatrix}.
 \label{eq:V62-angular-edge-lift}
\end{equation}

\begin{corollary}[Angular DeWitt Ward identities]
\label{cor:V62-angular-Ward-matrices}
For all \(\theta\),
\begin{align}
 B_{{\cal C},\theta}^*H
 &=E_\theta^*{\mathsf S}_{\cal C},
 \label{eq:V62-angular-constraint-duality}\\
 Q_\theta^{\rm TT}B_{{\cal C},\theta}
 &=K_\theta^*.
 \label{eq:V62-angular-shift-trace}
\end{align}
No inverse of \(K_\theta\) is used in either identity.
\end{corollary}

\begin{proof}
The first identity follows from the definition and self-adjointness of
\(H\).  For the second, use
\[
 Q_\theta^{\rm TT}B_{{\cal C},\theta}
 =(C_\theta^{\rm TT})^*
 {\mathsf S}_{\cal C}^*E_\theta
 =(D_\theta^{\rm TT})^*E_\theta
 =K_\theta^*E_\theta^*E_\theta
 =K_\theta^*.
\]
\end{proof}

\subsection{The angular two-edge trace action}
\label{subsec:V62-angular-action}

For \(e,z,\vartheta\in\mathbb C^2\), put
\begin{equation}
 w_\theta
 :=z-Q_\theta^{\rm TT}x+K_\theta^*e
 \label{eq:V62-angular-mismatch}
\end{equation}
and
\begin{equation}
 {\cal F}_{{\rm tr},\theta}
 :=
 \int_{\partial M}
 \operatorname{Re}\langle\vartheta,w_\theta\rangle
 \,\dd^3y.
 \label{eq:V62-angular-trace-action}
\end{equation}

\begin{theorem}[Two edges give the exact Ward identity at every
propagating angle]
\label{thm:V62-angular-action}
The action \eqref{eq:V62-angular-trace-action} is invariant under
\begin{equation}
 \delta_\eta x=B_{{\cal C},\theta}\eta,
 \qquad
 \delta_\eta e=\eta,
 \qquad
 \delta_\eta z=\delta_\eta\vartheta=0.
 \label{eq:V62-angular-edge-shift}
\end{equation}
Its Euler covectors are
\begin{align}
 E_x&=-(Q_\theta^{\rm TT})^*\vartheta,
 &
 \zeta:=E_e&=K_\theta\vartheta,
 \label{eq:V62-angular-Euler-xe}\\
 E_z&=\vartheta,
 &
 E_\vartheta&=w_\theta.
 \label{eq:V62-angular-Euler-ztheta}
\end{align}
With \(f:=H^{-1}E_x\),
\begin{align}
 f&=-C_\theta^{\rm TT}\vartheta,
 \label{eq:V62-angular-metric-force}\\
 {\mathsf S}_{\cal C}f+E_\theta\zeta&=0,
 \label{eq:V62-angular-full-Ward}\\
 B_{{\cal C},\theta}^*E_x+\zeta&=0.
 \label{eq:V62-angular-edge-Ward}
\end{align}
Thus the same two edge components compensate the rotating constraint
plane at every oblique angle.  At normal incidence all metric and edge
source terms vanish.
\end{theorem}

\begin{proof}
Corollary~\ref{cor:V62-angular-Ward-matrices} gives
\[
 \delta_\eta w_\theta
 =-Q_\theta^{\rm TT}B_{{\cal C},\theta}\eta
   +K_\theta^*\eta=0.
\]
This proves invariance.  Direct variation gives the four Euler
covectors.  The DeWitt identity in
\eqref{eq:V62-angular-incidence} gives
\(f=-C_\theta^{\rm TT}\vartheta\), and the factorization
\eqref{eq:V62-angular-factorization} gives
\[
 {\mathsf S}_{\cal C}f+E_\theta\zeta
 =-E_\theta K_\theta\vartheta
   +E_\theta K_\theta\vartheta=0.
\]
The edge Ward identity follows either by direct substitution or by
differentiating the invariant action.  When \(\mathfrak s=0\),
both \(Q_\theta^{\rm TT}\) and \(K_\theta\) vanish.
\end{proof}

\begin{corollary}[Positive angular trace penalty]
\label{cor:V62-angular-positive-penalty}
For every \(\mu>0\),
\begin{equation}
 {\cal F}_{\mu,\theta}^{+}
 :=
 \frac{\mu}{2}\int_{\partial M}|w_\theta|^2\,\dd^3y
 \label{eq:V62-angular-positive-penalty}
\end{equation}
is nonnegative and satisfies
\eqref{eq:V62-angular-full-Ward} with
\(\vartheta=\mu w_\theta\).  For fixed \(\theta\), its hard-trace
zero set \(w_\theta=0\) is a closed Sobolev trace subspace.
\end{corollary}

\subsection{Why the normal-incidence zero is an essential-row
obstruction}
\label{subsec:V62-normal-obstruction}

At normal incidence, the physical transverse plane is the tangent
spatial plane \(\operatorname{span}\{e_\parallel,e_\perp\}\).  A
normal-incidence TT tensor has a nonzero trace-free \(2\times2\)
tangential block.  The V37 conformal essential condition sets the
entire induced trace-free \(3\times3\) boundary tensor to zero.  On its
reduced configuration plane, the remaining tangential spatial block is
\((\tau/3)I_2\), whose two-dimensional TT projection vanishes.  This is
exactly \eqref{eq:V62-normal-trace-zero}.

\begin{theorem}[No uniform V59 trace lift on the fixed CM essential
plane]
\label{thm:V62-no-uniform-CM-lift}
There is no uniformly bounded family of right inverses for
\(Q_\theta^{\rm TT}\) on
\(0<\theta\le\pi/2\).  No continuous right inverse exists at
\(\theta=0\), because \(Q_0^{\rm TT}=0\).

Consequently a boundary construction which keeps
\(Q_{\rm TF}h=0\) as a homogeneous essential row cannot reproduce the
V59 two-component common trace uniformly on the full propagating cone.
At exact normal incidence it forces the continuum trace in
\eqref{eq:V62-angular-mismatch} to satisfy
\begin{equation}
 z=0
 \label{eq:V62-normal-continuum-Dirichlet}
\end{equation}
rather than the V59 relation \(z=a\) for an arbitrary radiative trace.
\end{theorem}

\begin{proof}
The right-inverse lower bound
\eqref{eq:V62-right-inverse-blowup} diverges as
\(\theta\downarrow0\), proving the first statement.  A zero map has no
right inverse onto \(\mathbb C^2\), proving the second.

The reduced plane
\eqref{eq:V62-reduced-spatial-block} is obtained only after imposing
the V37 induced conformal row.  The angular TT trace is therefore the
complete physical trace available inside that fixed plane.  Its normal
value is zero.  At \(\theta=0\), also \(K_\theta=0\), so the hard trace
\(w_\theta=0\) becomes \(z=0\).  V59 instead takes the two-polarization
field itself as the common trace.  The two domains cannot agree for an
arbitrary nonzero normal-incidence radiative boundary value.
\end{proof}

\begin{firewall}[The normal zero is not a bulk polarization loss]
\label{fire:V62-normal-scope}
Theorem~\ref{thm:V62-no-uniform-CM-lift} does not say that
normal-incidence gravitational radiation has no TT polarizations.
It says that the V37 boundary polarization has fixed those two induced
TT components to zero before the reduced five coordinates are formed.
The rank loss belongs to the chosen essential boundary row, consistent
with the V53 theorem that the old mixed-normal port is radiatively
invisible at normal incidence.
\end{firewall}

\subsection{V63 acquisition order: free the induced TT pair}
\label{subsec:V62-next}

The two-edge construction is sufficient on every oblique propagating
direction and requires no third edge on its Euler range.  The next
upstream step is to change the metric boundary polarization itself.
At the frozen quadratic level V63 should:

\begin{enumerate}
\item decompose the five-dimensional induced trace-free boundary fibre
into the two spatial TT components seen at normal incidence and a
three-dimensional complementary conformal/gauge sector;
\item retain the complementary three components as homogeneous
essential rows, but replace Dirichlet data on the TT pair by the V59
common trace and conormal balance;
\item use the V37 Green form
\[
 \Pi_{\rm TF}^{ab}\,\delta\gamma_{ab}^{\rm TF}
\]
to derive the TT conormal incidence from the same CM action, rather
than inserting \(C_\theta^{\rm TT}\) by hand;
\item recompute the harmonic source and edge Ward maps on this enlarged
seven-coordinate trace space; and
\item test constant rank at normal incidence, glancing, the continuum
threshold, and the static face.
\end{enumerate}

This mixed induced-TT polarization is the first candidate capable of
matching the V59 trace at both normal incidence and glancing.  If its
three-component complement cannot be chosen continuously over the
direction sphere, V63 should state the bundle obstruction and use
projector-valued rows rather than a global frame.

\subsection{Finite certificates}
\label{subsec:V62-certificates}

The canonical glancing-extension card is
\begin{equation}
 Q_{\rm g}^{\rm TT}-Q_{\rm rad}
 =
 \frac1{3\sqrt2}
 \begin{pmatrix}1&0&0&0&3\\0&0&0&0&0\end{pmatrix},
 \qquad
 (Q_{\rm g}^{\rm TT}-Q_{\rm rad})
 {\mathscr K}_\epsilon=0.
 \label{eq:V62-extension-card}
\end{equation}
The angular Gram card is
\begin{equation}
 Q_\theta^{\rm TT}(Q_\theta^{\rm TT})^*
 =
 \operatorname{diag}\!\left(
 \mathfrak s^2(2-13\mathfrak s^2/9),
 2\mathfrak s^2\right).
 \label{eq:V62-Gram-card}
\end{equation}
The action cards are
\begin{align}
 HC_\theta^{\rm TT}-(Q_\theta^{\rm TT})^*&=0,
 \label{eq:V62-HC-card}\\
 {\mathsf S}_{\cal C}C_\theta^{\rm TT}
 -E_\theta K_\theta&=0,
 \label{eq:V62-factor-card}\\
 B_{{\cal C},\theta}^*H
 -E_\theta^*{\mathsf S}_{\cal C}&=0,
 \label{eq:V62-dual-card}\\
 Q_\theta^{\rm TT}B_{{\cal C},\theta}
 -K_\theta^*&=0.
 \label{eq:V62-shift-card}
\end{align}
At normal incidence,
\begin{equation}
 \operatorname{rank}Q_0^{\rm TT}
 =\operatorname{rank}C_0^{\rm TT}
 =\operatorname{rank}K_0=0.
 \label{eq:V62-normal-rank-card}
\end{equation}

\subsection{V62 result ledger and claim boundary}
\label{subsec:V62-ledger}

Relative to V1--V61, this note proves:

\begin{enumerate}
\item the convenient V60 glancing trace and the canonical full-fibre TT
trace differ only by the scalar annihilator \(\tau+3\phi\), so they
define the same quotient trace;
\item the V61 rank-one scalar defect vanishes on the actual TT Euler
range and does not supply a lower bound for a third edge field;
\item the exact TT projection of the V37 reduced configuration plane is
the angular map \eqref{eq:V62-angular-TT-map};
\item this map has rank two at every oblique propagating angle, rank
zero at normal incidence, and every family of right inverses blows like
\(1/\sin\theta\);
\item its DeWitt metric force sources the rotating two-dimensional
constraint plane
\(\operatorname{span}\{C_\perp,
\mathfrak cC_\parallel-\mathfrak sC_n\}\);
\item the same two edge components give an exact action-level Ward
identity throughout the real propagating angular cone;
\item the two-channel incidence vanishes, rather than becoming
singular, at normal incidence; and
\item the fixed V37 conformal essential row is exactly what prevents a
uniform lift of the V59 two-polarization common trace.
\end{enumerate}

No full stable-cone metric lift is proved.  The angular theorem covers
the real propagating cone, not open Laplace frequencies, evanescent
frequencies, or the static face.  The induced TT pair has not yet been
freed in the CM action domain, and the edge shift has no completed BRST
quartet or raw-representative generator.  The nonlinear
Brown--York/Bianchi identity and the coupled Higgs, fermion, neutrino,
and Yang--Mills sectors remain open.  The full Standard Model--Einstein
continuum is not proved.

\subsection{Append-only record}

This V62 material was appended to the exact V61 parent whose SHA--256
digest is
\[
\texttt{a51f380bdbeff2b1356d2c4e0e1895b2bf9d84efb5a019fb1bc78ea877837c76}.
\]
The next research note is V63.  It should replace Dirichlet data on the
induced TT pair by the V59 trace/conormal relation while retaining and
rechecking the remaining CM, harmonic, York, and edge rows.

% END APPENDED VERSION V62 MATERIAL

% BEGIN APPENDED VERSION V63 MATERIAL

\clearpage
\section{V63: a uniformly split induced radiative graph, the
mean-curvature blow-up, and an exact real-cone metric--continuum lift}
\label{sec:V63-induced-radiative-graph}

\subsection{Purpose and the upstream correction}
\label{subsec:V63-purpose}

V62 proved that the canonical TT trace of the five-coordinate V37
conformal--mean-curvature plane has rank two at every oblique real
propagating angle and rank zero at normal incidence.  The loss occurs
because all five components of the induced trace-free boundary metric
were set to zero before that reduced plane was formed.  Its provisional
handoff was therefore to free the two induced spatial TT components
while retaining the three-component complement.

That operation is necessary at normal incidence, but it is not by
itself sufficient at gravitational glancing.  The induced restriction
of the V53 cross tensor vanishes when the ray is tangent to the
boundary.  V63 computes this restriction and proves an exact
\(1/\cos\theta\) inverse loss.  It then constructs a bounded residual
wave-gauge completion which rotates one induced time--space component
into the cross trace.  The resulting direction-dependent two-plane has
a uniform configuration splitting on the complete real angular
interval.

A second obstruction appears if the V37 row \(K_{(1)}=0\) is kept.
For the plus polarization, every null-gauge representative satisfying
that row has a normal gauge amplitude proportional to
\(\sin^2\theta/\cos\theta\).  Thus the tentative V62 prescription would
only move the glancing singularity from the TT trace into the gauge
lift.  The correct upstream operation is to reverse the
mean-curvature Legendre choice on the radiative graph: retain the
physical induced trace as part of the common configuration trace and
allow \(K_{(1)}\) to enter its Brown--York conormal.

The main result below is an exact frozen theorem on every nonzero real
null covector.  A two-component continuum trace determines all six
induced metric components through a uniformly split graph; the four
harmonic rows remain homogeneous; and variation of the same
GHY--de Donder action gives the V59 radiative conormal.  With the
continuum traction, the complete twenty-four-dimensional normal Green
relation is Lagrangian and half-dimensional.  Analytic continuation to
the open Laplace cone, a Sobolev realization of the angular projector,
and the off-shell ghost/edge domain remain for V64.

\subsection{A five-dimensional induced trace-free chart}
\label{subsec:V63-induced-chart}

Fix the planar boundary frame
\[
 (u,e_1,e_2,n),
\]
where \(u\) is future timelike, \(e_1,e_2\) are spatial and tangent,
and \(n\) is the outward spacelike normal.  The induced metric is
\(\eta_\partial=\operatorname{diag}(-1,1,1)\) in
\((u,e_1,e_2)\).  Let
\[
 {cal V}_{\partial}^{\rm TF}
 :=\left\{v_{ab}=v_{ba}:\eta_\partial^{ab}v_{ab}=0\right\}.
\]
This real fibre has dimension five.

Use the following Euclidean-Frobenius orthonormal coefficient basis:
\begin{align}
 F_+&:=\frac{e_1\otimes e_1-e_2\otimes e_2}{\sqrt2},
 &
 F_\times&:=\frac{e_1\otimes e_2+e_2\otimes e_1}{\sqrt2},
 \label{eq:V63-fixed-spatial-TT}\
 F_{01}&:=\frac{u^\flat\otimes e_1+e_1\otimes u^\flat}{\sqrt2},
 &
 F_{02}&:=\frac{u^\flat\otimes e_2+e_2\otimes u^\flat}{\sqrt2},
 \label{eq:V63-time-space-basis}\
 F_s&:=\frac{2u^\flat\otimes u^\flat
                 +e_1\otimes e_1+e_2\otimes e_2}{\sqrt6}.
 \label{eq:V63-scalar-TF-basis}
\end{align}
Every displayed tensor has zero Lorentzian induced trace.  The first
two span the spatial two-dimensional trace-free plane selected at
normal incidence.  The last three give an explicit complementary
plane.

\begin{proposition}[Exact normal-incidence split]
\label{prop:V63-normal-split}
Let
\begin{equation}
 {mathbb P}_{N}
 :=F_+\otimes_EF_+ +F_\times\otimes_EF_\times
 \label{eq:V63-normal-projector}
\end{equation}
with respect to the Euclidean coefficient inner product.  Then
\[
 {mathbb P}_{N}^2={\mathbb P}_{N}={\mathbb P}_{N}^*,
 \qquad
 \operatorname{rank}{\mathbb P}_{N}=2,
\]
and
\begin{equation}
 {cal V}_{\partial}^{\rm TF}
 =\operatorname{span}\{F_+,F_\times\}
 \mathbin{\mathop{\oplus}^{\perp_E}}
 \operatorname{span}\{F_{01},F_{02},F_s\}.
 \label{eq:V63-normal-direct-sum}
\end{equation}
Thus the proposed two-plus-three split exists globally in the fixed
planar foliation.  The issue is its incidence on oblique radiative
classes, not a dimension defect.
\end{proposition}

\begin{proof}
The five matrices in
\eqref{eq:V63-fixed-spatial-TT}--%
\eqref{eq:V63-scalar-TF-basis} are visibly linearly independent.
Their squared Frobenius norms are one and their pairwise products are
zero.  Since \(\dim{\cal V}_{\partial}^{\rm TF}=5\), they are an
orthonormal basis.  The projector and direct-sum identities follow.
\end{proof}

\subsection{The induced restriction loses one glancing polarization}
\label{subsec:V63-induced-restriction}

Use the V62 real null direction
\begin{equation}
 \widehat\xi=(\mathfrak s,0,\mathfrak c),
 \qquad
 \mathfrak s=\sin\theta,
 \quad
 \mathfrak c=\cos\theta,
 \quad
 0\leq\theta\leq\frac\pi2,
 \label{eq:V63-real-direction}
\end{equation}
and a nonzero spatial wave number \(\varrho>0\).  The Fourier factor is
chosen as
\begin{equation}
 \exp\!\left(i\varrho(-t+\mathfrak s y^1+
                         \mathfrak c x_n)\right).
 \label{eq:V63-null-Fourier-factor}
\end{equation}
All Green identities below are real bilinear identities and are then
complexified algebraically; superscript \(T\) denotes bilinear transpose.
Euclidean stars are used only for the finite Gram and projector cards.
Thus
\[
 D_n=i\varrho\mathfrak c,
 \qquad
 D_0=-i\varrho,
 \qquad
 D_1=i\varrho\mathfrak s,
 \qquad
 D_2=0.
\]
Let \(E_+(\theta),E_\times(\theta)\) be the V62 spatial TT basis and
put
\begin{equation}
 h^{\rm TT}(a)
 :=a_+E_+(\theta)+a_\times E_\times(\theta),
 \qquad a=(a_+,a_\times)^T.
 \label{eq:V63-TT-amplitude}
\end{equation}
It is the norm-one V53 splitting of the radiative cohomology class.

Write
\[
 \tau_\partial:=\eta_\partial^{ab}h_{ab},
 \qquad
 u_{ab}:=h_{ab}-\frac13\tau_\partial\eta_{\partial,ab}.
\]
Direct restriction of \eqref{eq:V63-TT-amplitude} gives
\begin{equation}
 I_\theta a
 :=u(h^{\rm TT}(a))
 =\left[-\frac{1+\mathfrak c^2}{2}F_+
       +\frac{\mathfrak s^2}{\sqrt{12}}F_s\right]a_+
  -\mathfrak c F_\times a_\times.
 \label{eq:V63-induced-TF-restriction}
\end{equation}
The induced trace is
\begin{equation}
 \tau_\partial(h^{\rm TT}(a))
 =\frac{\mathfrak s^2}{\sqrt2}a_+.
 \label{eq:V63-induced-trace}
\end{equation}

\begin{theorem}[Exact induced-only glancing loss]
\label{thm:V63-induced-loss}
The columns of \(I_\theta:\mathbb C^2\to
{\cal V}_{\partial}^{\rm TF}\) are orthogonal and
\begin{equation}
 I_\theta^*I_\theta
 =\operatorname{diag}\!\left(
  \frac{1+\mathfrak c^2+\mathfrak c^4}{3},
  \mathfrak c^2\right).
 \label{eq:V63-induced-Gram}
\end{equation}
Consequently
\begin{align}
 \operatorname{rank}I_\theta&=2
 &&(0\leq\theta<\pi/2),
 \label{eq:V63-induced-oblique-rank}\
 \operatorname{rank}I_{\pi/2}&=1,
 \label{eq:V63-induced-glancing-rank}\
 s_{\min}(I_\theta)&=\mathfrak c
 &&(0\leq\theta<\pi/2).
 \label{eq:V63-induced-smallest-singular}
\end{align}
Every induced-only right inverse therefore has norm at least
\(1/\mathfrak c\) and diverges at gravitational glancing.  The same
rank-one glancing loss occurs after applying the fixed projector
\({\mathbb P}_N\).
\end{theorem}

\begin{proof}
The plus column in \eqref{eq:V63-induced-TF-restriction} is supported
on \(F_+,F_s\), while the cross column is supported on \(F_\times\).
Its squared plus norm is
\[
 \frac{(1+\mathfrak c^2)^2}{4}
 +\frac{\mathfrak s^4}{12}
 =\frac{1+\mathfrak c^2+\mathfrak c^4}{3}.
\]
The cross norm is \(\mathfrak c\).  Moreover
\[
 \frac{1+\mathfrak c^2+\mathfrak c^4}{3}
 -\mathfrak c^2
 =\frac{(1-\mathfrak c^2)^2}{3}\geq0,
\]
so the cross value is the least singular value before glancing.  At
\(\mathfrak c=0\), the V62 cross tensor has one normal and one
\(e_2\) index, hence its entire induced restriction is zero.  Applying
\({\mathbb P}_N\) leaves coefficients
\(-(1+\mathfrak c^2)/2,-\mathfrak c\), so it cannot restore the lost
column.
\end{proof}

\begin{firewall}[Representative scope of the induced no-go]
\label{fire:V63-induced-scope}
Theorem~\ref{thm:V63-induced-loss} concerns the canonical V53 TT
representative.  It does not prove that every representative of the
radiative class has zero induced cross trace at glancing.  A residual
null gauge can change the induced tensor while leaving the V53
cohomology class fixed.  The next construction uses exactly that
freedom and records it rather than treating the induced restriction as
gauge invariant.
\end{firewall}

\subsection{A bounded cross-gauge completion and moving two-plane}
\label{subsec:V63-cross-completion}

Let the dimensionless null covector be
\[
 \ell_\mu=(-1,\mathfrak s,0,\mathfrak c).
\]
For the cross amplitude define
\begin{equation}
 g^{\times}_\mu
 :=-\frac{\mathfrak s}{\sqrt2(1+\mathfrak c)}
     a_\times\,\delta_{\mu2},
 \qquad
 \xi^{\times}_\mu:=\frac{g^{\times}_\mu}{i\varrho},
 \label{eq:V63-cross-gauge-vector}
\end{equation}
and add the underlying metric gauge perturbation
\begin{equation}
 h^g_{\mu\nu}
 :=\ell_\mu g^{\times}_\nu+ell_\nu g^{\times}_\mu.
 \label{eq:V63-cross-gauge-metric}
\end{equation}
The half-angle coefficient
\begin{equation}
 \mathfrak t:=\frac{\mathfrak s}{1+\mathfrak c}
 =\tan\frac\theta2
 \label{eq:V63-half-angle}
\end{equation}
is continuous and lies in \([0,1]\) on the complete angular interval.

Because \(\ell\) is null, V36 gives
\(\delta C_\mu=\Box\xi^{\times}_\mu=0\).  The gauge vector has no
normal component, so the V37 identity
\(\delta K_{(1)}=-D^aD_a\xi_n\) gives
\(\delta K_{(1)}=0\).  The ungauged cross tensor also has
\(K_{(1)}=0\).

The nonzero induced coefficients of the completed cross tensor are
\begin{align}
 \sqrt2\,(h^{\rm TT}+h^g)_{12}&=-a_\times,
 &
 \sqrt2\,(h^{\rm TT}+h^g)_{02}&=\mathfrak t a_\times.
 \label{eq:V63-completed-cross-components}
\end{align}
Hence define the moving columns
\begin{align}
 G_+(\theta)
 &:=-\frac{1+\mathfrak c^2}{2}F_+
   +\frac{\mathfrak s^2}{\sqrt{12}}F_s,
 \label{eq:V63-G-plus}\
 G_\times(\theta)
 &:=-F_\times+\mathfrak t F_{02},
 \label{eq:V63-G-cross}
\end{align}
and the injection
\begin{equation}
 G_\theta a:=G_+(\theta)a_+ +G_\times(\theta)a_\times.
 \label{eq:V63-moving-injection}
\end{equation}

\begin{theorem}[Uniformly split gauge-completed induced trace]
\label{thm:V63-uniform-G}
For every \(0\leq\theta\leq\pi/2\),
\begin{equation}
 G_\theta^*G_\theta
 =\operatorname{diag}\!\left(
 \frac{1+\mathfrak c^2+\mathfrak c^4}{3},
 1+\mathfrak t^2\right).
 \label{eq:V63-G-Gram}
\end{equation}
In particular,
\begin{equation}
 s_{\min}(G_\theta)\geq\frac1{\sqrt3},
 \qquad
 Q_\theta^G
 :=(G_\theta^*G_\theta)^{-1}G_\theta^*,
 \qquad
 Q_\theta^GG_\theta=I_2,
 \qquad
 \|Q_\theta^G\|\leq\sqrt3.
 \label{eq:V63-G-right-inverse}
\end{equation}
The rank-two projector
\begin{equation}
 {\mathbb P}_\theta^G:=G_\theta Q_\theta^G
 \label{eq:V63-moving-projector}
\end{equation}
is continuous at normal incidence and glancing.

A continuous complementary frame is given, before normalization, by
\begin{align}
 H_+(\theta)
 &:=\frac{\mathfrak s^2}{\sqrt{12}}F_+
    +\frac{1+\mathfrak c^2}{2}F_s,
 \label{eq:V63-H-plus}\
 H_\times(\theta)
 &:=\mathfrak t F_\times+F_{02},
 \label{eq:V63-H-cross}\
 H_0&:=F_{01}.
 \label{eq:V63-H-zero}
\end{align}
Thus there is no bundle obstruction on this planar half-angle interval.
\end{theorem}

\begin{proof}
The two columns of \(G_\theta\) have disjoint support in the fixed
orthonormal basis.  The first squared norm is the plus entry already
computed in \eqref{eq:V63-induced-Gram}; the second is
\(1+\mathfrak t^2\).  The first lies in \([1/3,1]\), and the second in
\([1,2]\), proving the singular-value floor and the right-inverse
bound.  The coefficients are continuous because
\(1+\mathfrak c\geq1\).

Direct products give
\[
 (G_+,H_+)_E=(G_\times,H_\times)_E=0.
\]
The three \(H\)-vectors are mutually orthogonal, are nonzero at both
endpoints, and are orthogonal to both \(G\)-columns.  They therefore
span the rank-three complement continuously.  Formula
\eqref{eq:V63-moving-projector} is the orthogonal projector onto the
moving two-plane.
\end{proof}

\begin{remark}[What the moving plane costs]
\label{rem:V63-moving-cost}
The fixed normal-incidence plane of
Proposition~\ref{prop:V63-normal-split} is local in the chosen boundary
foliation but loses a cross column at glancing.  The moving plane keeps
a uniform floor by using the direction multiplier
\(\mathfrak t=\mathfrak s/(1+\mathfrak c)\) and a residual-gauge
representative.  When pulled back from Fourier space it is therefore a
projector-valued order-zero boundary operator.  V63 proves its finite
real-cone algebra, not its full Sobolev or causal realization.
\end{remark}

\subsection{Why retaining zero mean curvature still fails}
\label{subsec:V63-K-no-go}

For the plus TT representative, direct substitution in the V36
extrinsic-curvature formula gives
\begin{align}
 K^{\rm TT}_{11}
 &=-\frac{i\varrho\mathfrak c(2-\mathfrak c^2)}{2\sqrt2}a_+,
 &
 K^{\rm TT}_{22}
 &=\frac{i\varrho\mathfrak c}{2\sqrt2}a_+,
 \label{eq:V63-plus-K-components}\
 K_{(1)}(h^{\rm TT}_+)
 &=-\frac{i\varrho\mathfrak c\mathfrak s^2}{2\sqrt2}a_+.
 \label{eq:V63-plus-mean-curvature}
\end{align}
The cross value is zero.

Let a residual null gauge have dimensionless normal amplitude
\(g_n=i\varrho\xi_n\).  The tangential wave operator has symbol
\(\varrho^2\mathfrak c^2\), so the V37 transformation law is
\begin{equation}
 \delta K_{(1)}
 =-D^aD_a\xi_n
 =i\varrho\mathfrak c^2 g_n.
 \label{eq:V63-gauge-K-incidence}
\end{equation}

\begin{theorem}[The fixed mean-curvature row forces a glancing gauge
blow-up]
\label{thm:V63-K-blowup}
Fix a nonzero plus amplitude and \(0<\mathfrak c\leq1\).  Every
harmonic null representative of the same V53 radiative class which
satisfies
\[
 K_{(1)}=0
\]
has normal gauge component
\begin{equation}
 g_n=\frac{\mathfrak s^2}{2\sqrt2\,\mathfrak c}a_+.
 \label{eq:V63-required-normal-gauge}
\end{equation}
Consequently
\begin{equation}
 \|g\|\geq
 \frac{\mathfrak s^2}{2\sqrt2\,\mathfrak c}|a_+|
 \longrightarrow\infty
 \qquad(\mathfrak c\downarrow0).
 \label{eq:V63-normal-gauge-blowup}
\end{equation}
There is no uniformly bounded physical-class lift which both retains
the V37 homogeneous mean-curvature row and approaches glancing through
the real null cone.
\end{theorem}

\begin{proof}
The V53 split exact sequence says that two harmonic null amplitudes in
the same TT class differ by a unique residual null-gauge amplitude.
Tangential gauge components do not occur in
\eqref{eq:V63-gauge-K-incidence}.  Equating the sum of
\eqref{eq:V63-plus-mean-curvature} and
\eqref{eq:V63-gauge-K-incidence} to zero gives
\eqref{eq:V63-required-normal-gauge}.  Its absolute value gives the
lower bound.  At exact glancing both terms in the mean-curvature
equation vanish, but that endpoint equation cannot provide a bounded
continuation of the unique oblique solution.
\end{proof}

\begin{firewall}[Relation to the V37 and V60 endpoint kernels]
\label{fire:V63-K-scope}
Theorem~\ref{thm:V63-K-blowup} does not contradict the existence of the
bounded V60 glancing representatives \(r_+,r_\times\).  It proves that
the representative selected by the homogeneous mean-curvature row for
\(\mathfrak c>0\) has no bounded limit to that endpoint.  The obstruction
is the same principal angle loss seen in V37, now evaluated on the
physical plus class rather than on an unclassified stable vector.
\end{firewall}

\subsection{The physical induced-trace graph}
\label{subsec:V63-physical-graph}

The blow-up disappears if the induced trace is retained instead of
forcing its conjugate mean curvature to zero.  Define
\begin{equation}
 L_\theta
 :=\begin{pmatrix}\mathfrak s^2/\sqrt2&0\end{pmatrix}
 :\mathbb C^2\longrightarrow\mathbb C
 \label{eq:V63-L-trace}
\end{equation}
and the six-component induced configuration graph
\begin{equation}
 {\mathscr G}_\theta
 :=\left\{(u,\tau_\partial):
 u=G_\theta z,\quad \tau_\partial=L_\theta z,
 \quad z\in\mathbb C^2\right\}.
 \label{eq:V63-induced-configuration-graph}
\end{equation}
The plus representative is left in the norm-one V53 TT gauge.  The
cross representative is completed by
\eqref{eq:V63-cross-gauge-vector}.  Both then satisfy
\begin{equation}
 C_\mu=0,
 \qquad
 (u,\tau_\partial)=(G_\theta a,L_\theta a),
 \label{eq:V63-representative-graph}
\end{equation}
without any singular normal gauge.

To derive the conormal, use the V37 algebraic splitting of the
Brown--York pairing,
\begin{equation}
 \Pi^{ab}\delta\gamma_{ab}
 =\Pi_{\rm TF}^{ab}\delta u_{ab}
  -\frac23K_{(1)}\delta\tau_\partial.
 \label{eq:V63-BY-splitting}
\end{equation}
For a trace-free tensor \(W\), let
\[
 W^\vee\Pi_{\rm TF}:=\Pi_{\rm TF}^{ab}W_{ab};
\]
this uses the Lorentzian raised-index Brown--York pairing, not the
Euclidean coefficient product used in the Gram matrices.  Define the
two-component action conormal on the graph by
\begin{equation}
 r_\theta(h)
 :=\varepsilon_D\left[
 G_\theta^\vee\Pi_{(1),\rm TF}(h)
 -\frac23L_\theta^T K_{(1)}(h)
 \right].
 \label{eq:V63-radiative-BY-conormal}
\end{equation}
Here \(G_\theta^\vee\Pi\) means contraction with each of the two
columns of \(G_\theta\).

\begin{proposition}[The conormal is forced by the GHY action]
\label{prop:V63-conormal-from-action}
On the harmonic shell \(C_\mu=0\), restriction of the V36
GHY--de Donder boundary one-form to variations tangent to
\({\mathscr G}_\theta\) is
\begin{equation}
 \Theta_{D,\theta}^{\rm met}(h;\delta z)
 =\frac1{2\kappa}\int_{\partial M}
 r_\theta(h)^T\delta z.
 \label{eq:V63-restricted-one-form}
\end{equation}
For two graph-tangent amplitudes, the metric Green form is
\begin{equation}
 2\kappa\,\Omega_{D,\theta}(h,k)
 =r_\theta(h)^Tz(k)-z(h)^Tr_\theta(k).
 \label{eq:V63-restricted-Green}
\end{equation}
Thus no TT conormal incidence is inserted by hand: both terms in
\eqref{eq:V63-radiative-BY-conormal} are forced by the same GHY
variation.
\end{proposition}

\begin{proof}
On \(C_\mu=0\), the gauge-fixing coefficient in the V36 boundary
one-form vanishes.  A graph-tangent variation obeys
\(\delta u=G_\theta\delta z\) and
\(\delta\tau_\partial=L_\theta\delta z\).  Substitution in
\eqref{eq:V63-BY-splitting} gives
\eqref{eq:V63-restricted-one-form}.  Antisymmetrizing its linearization
gives \eqref{eq:V63-restricted-Green}.  Equivalently, substitute the
same graph into the V37 Green form; the term containing
\(\tau_\partial K_{(1)}\) produces exactly the second summand of
\eqref{eq:V63-radiative-BY-conormal}.
\end{proof}

\subsection{Exact recovery of the V59 trace and conormal}
\label{subsec:V63-V59-recovery}

The following calculation is the decisive finite certificate.  For the
ungauged plus representative,
\begin{align}
 G_+^\vee\Pi_{(1),\rm TF}
 &=\frac{i\varrho\mathfrak c}{6}
   \left(2+2\mathfrak c^2-\mathfrak c^4\right)a_+,
 \label{eq:V63-plus-BY-part}\
 -\frac23L_+K_{(1)}
 &=\frac{i\varrho\mathfrak c}{6}
   \mathfrak s^4a_+.
 \label{eq:V63-plus-trace-K-part}
\end{align}
Since
\[
 2+2\mathfrak c^2-\mathfrak c^4+\mathfrak s^4=3,
\]
the two pieces sum to \(i\varrho\mathfrak c a_+/2\).

For the gauge-completed cross representative, the only conormal
components needed are
\begin{align}
 K_{12}&=-\frac{i\varrho}{2\sqrt2}a_\times,
 &
 K_{02}&=\frac{i\varrho\mathfrak s}{2\sqrt2}a_\times,
 &
 K_{(1)}&=0.
 \label{eq:V63-cross-K-components}
\end{align}
Raising the time index reverses the sign of the \(02\) contraction.
Using \(G_\times=-F_\times+\mathfrak tF_{02}\) gives
\begin{equation}
 G_\times^\vee\Pi_{(1),\rm TF}
 =\frac{i\varrho}{2}
  \left(1-\mathfrak t\mathfrak s\right)a_\times
 =\frac{i\varrho\mathfrak c}{2}a_\times.
 \label{eq:V63-cross-BY-part}
\end{equation}

\begin{theorem}[Exact real-cone trace/conormal lift]
\label{thm:V63-exact-real-cone-lift}
For every \(\varrho>0\) and
\(0\leq\theta\leq\pi/2\), the representatives in
\eqref{eq:V63-representative-graph} have graph trace
\begin{equation}
 z=a
 \label{eq:V63-trace-is-amplitude}
\end{equation}
and action conormal
\begin{equation}
 r_\theta
 =\frac{\varepsilon_D}{2}
   i\varrho\mathfrak c\,a.
 \label{eq:V63-action-conormal-value}
\end{equation}
With the equation normalization
\begin{equation}
 p_{\rm rad}:=2\varepsilon_D r_\theta,
 \label{eq:V63-equation-normalization}
\end{equation}
one obtains
\begin{equation}
 p_{\rm rad}=i\varrho\mathfrak c\,a=D_na,
 \label{eq:V63-V59-conormal}
\end{equation}
which is exactly the V59 outward radiative conormal in this Fourier
orientation.  The equality holds at normal incidence and extends
continuously to the zero conormal at gravitational glancing.
\end{theorem}

\begin{proof}
The configuration equality follows from
\eqref{eq:V63-representative-graph} and
\(Q_\theta^GG_\theta=I_2\).  Equations
\eqref{eq:V63-plus-BY-part}--%
\eqref{eq:V63-plus-trace-K-part} give the plus column of
\eqref{eq:V63-action-conormal-value};
\eqref{eq:V63-cross-BY-part} gives the cross column.  Multiplication by
\(2\varepsilon_D\), with \(\varepsilon_D^2=1\), gives
\eqref{eq:V63-V59-conormal}.  The endpoint values are obtained by
putting \(\mathfrak c=1\) and \(\mathfrak c=0\).
\end{proof}

\begin{corollary}[The trace term is essential in the plus channel]
\label{cor:V63-trace-term-essential}
Deleting the
\(-\tfrac23L_\theta^TK_{(1)}\) term from
\eqref{eq:V63-radiative-BY-conormal} would leave the angle-dependent
plus coefficient
\[
 \frac{i\varrho\mathfrak c}{6}
 (2+2\mathfrak c^2-\mathfrak c^4).
\]
It would not equal the scalar radiative conormal.  The induced trace is
therefore not a dispensable complement variable: it is the part of the
GHY virtual work which repairs the plus normalization without the
singular gauge vector
\eqref{eq:V63-required-normal-gauge}.
\end{corollary}

\subsection{The maximal frozen metric--continuum Green relation}
\label{subsec:V63-maximal-relation}

At a fixed nonzero real null covector, adjoin a two-component continuum
boundary phase pair \((b,t)\).  Normalize \(t\) in the same action units
as \(r_\theta\), with total boundary Green form
\begin{equation}
 \Omega_{\rm tot}
 :=\Omega_D
 +\frac1{2\kappa}
 \left(t^Tb'-b^Tt'\right).
 \label{eq:V63-total-Green}
\end{equation}
The corresponding V59 equation-normalized traction is
\begin{equation}
 t_b:=2\varepsilon_D t.
 \label{eq:V63-bath-normalization}
\end{equation}

Define
\begin{equation}
 {\cal L}_{63}(\theta)
 :=\left\{
 \begin{array}{l}
 u=G_\theta b,\\
 \tau_\partial=L_\theta b,\\
 C_\mu=0,\\
 t+r_\theta=0.
 \end{array}
 \right.
 \label{eq:V63-full-relation}
\end{equation}
There are five equations in the first row, one in the second, four in
the third, and two in the fourth.

\begin{theorem}[Maximal action Green relation on the real null cone]
\label{thm:V63-maximal-Green}
For every \(\varrho>0\) and
\(0\leq\theta\leq\pi/2\),
\({\cal L}_{63}(\theta)\) is Lagrangian in the
\(24\)-dimensional metric--continuum normal phase fibre.  More
precisely,
\begin{equation}
 \operatorname{rank}B_{63}=5+1+4+2=12,
 \qquad
 \dim\ker B_{63}=12,
 \qquad
 {\cal L}_{63}(\theta)^{\perp_{\Omega_{\rm tot}}}
 ={\cal L}_{63}(\theta).
 \label{eq:V63-maximal-card}
\end{equation}
On the radiative representatives of
Theorem~\ref{thm:V63-exact-real-cone-lift}, its last row is exactly
\begin{equation}
 p_{\rm rad}+t_b=0,
 \label{eq:V63-V59-balance}
\end{equation}
the V59 conormal balance, while the first two rows give the common trace
\(b=a\).
\end{theorem}

\begin{proof}
Take two elements of \({\cal L}_{63}(\theta)\).  The homogeneous
harmonic rows annihilate the de Donder term in the V36 Green form.
Proposition~\ref{prop:V63-conormal-from-action} and the common
configuration trace give the metric contribution
\[
 \frac1{2\kappa}
 \left(r_\theta^Tb'-b^Tr_\theta'\right).
\]
The continuum contribution, with \(t=-r_\theta\), is its negative.
Thus \({\cal L}_{63}(\theta)\) is isotropic.

The metric normal phase fibre has dimension twenty and the continuum
pair has dimension four.  The six induced-configuration equations are
independent because their metric configuration block is the identity
on the six induced components.  Modulo those rows, the four harmonic
conditions have the same independent normal pivots used in the V36 and
V37 dimension proofs.  Finally the two traction equations are
independent because their continuum-traction block is \(I_2\).  Hence
the row rank is twelve.  An isotropic half-dimensional subspace of a
nondegenerate symplectic fibre is Lagrangian, proving maximality.

On the physical representatives,
\(Q_\theta^Gu=b=a\) and
Theorem~\ref{thm:V63-exact-real-cone-lift} gives
\(2\varepsilon_Dr_\theta=p_{\rm rad}\).  Multiplying the last row by
\(2\varepsilon_D\) gives
\eqref{eq:V63-V59-balance}.
\end{proof}

\begin{assessment}[What has now been lifted]
\label{ass:V63-lift-status}
V59's two-polarization common trace and conormal balance now have an
exact representative in the full ten-component metric normal phase
fibre at every nonzero real propagating angle.  The four harmonic rows
are part of the same maximal relation, and the conormal is the
Brown--York variation of the GHY--de Donder action on the induced
configuration graph.  V60's rank-two compensator is not needed for
this on-shell real-cone relation because the physical conormal is no
longer inserted into the old five-row CM graph.

This is stronger than the V62 angular edge identity but narrower in
frequency: it is a frozen null-shell theorem.  It does not yet prove
that \({\cal L}_{63}(\theta)\) is the boundary domain of a closed
positive time generator, nor that the moving cross representative is
an off-shell BRST orbit with a recorded edge variable.
\end{assessment}

\subsection{Endpoint and finite certificates}
\label{subsec:V63-certificates}

At normal incidence,
\begin{equation}
 \mathfrak c=1,quad \mathfrak s=\mathfrak t=0,qquad
 G_0=(-F_+,-F_\times),qquad
 L_0=0,qquad
 r_0=\frac{\varepsilon_D i\varrho}{2}I_2.
 \label{eq:V63-normal-card}
\end{equation}
At gravitational glancing,
\begin{equation}
 \mathfrak c=0,quad \mathfrak s=\mathfrak t=1,qquad
 G_{\rm g}=left(
 -\frac12F_++\frac1{\sqrt{12}}F_s,
 -F_\times+F_{02}
 \right),qquad
 L_{\rm g}=\begin{pmatrix}1/\sqrt2&0\end{pmatrix},qquad
 r_{\rm g}=0.
 \label{eq:V63-glancing-card}
\end{equation}
Both endpoint injections have rank two.

The configuration certificates are
\begin{align}
 I_\theta^*I_\theta
 &-\operatorname{diag}\!\left(
  (1+\mathfrak c^2+\mathfrak c^4)/3,
  \mathfrak c^2\right)=0,
 \label{eq:V63-I-card}\
 G_\theta^*G_\theta
 &-\operatorname{diag}\!\left(
  (1+\mathfrak c^2+\mathfrak c^4)/3,
  1+\mathfrak t^2\right)=0,
 \label{eq:V63-G-card}\
 Q_\theta^GG_\theta-I_2&=0,
 \qquad
 s_{\min}(G_\theta)\geq1/\sqrt3.
 \label{eq:V63-splitting-card}
\end{align}
The gauge and conormal certificates are
\begin{align}
 K_{(1)}(h_+^{\rm TT})
 +i\varrho\mathfrak c^2
 \frac{\mathfrak s^2}{2\sqrt2\mathfrak c}a_+&=0,
 \label{eq:V63-K-cancellation-card}\
 G_+^\vee\Pi_{\rm TF}
 -\frac23L_+K_{(1)}
 -\frac{i\varrho\mathfrak c}{2}a_+&=0,
 \label{eq:V63-plus-conormal-card}\
 G_\times^\vee\Pi_{\rm TF}
 -\frac{i\varrho\mathfrak c}{2}a_\times&=0,
 \label{eq:V63-cross-conormal-card}\
 p_{\rm rad}-2\varepsilon_Dr_\theta&=0.
 \label{eq:V63-normalization-card}
\end{align}
Finally the maximality certificate is
\begin{equation}
 \dim{\cal P}_{\rm met+bath}=24,qquad
 \operatorname{rank}B_{63}=12,qquad
 \Omega_{\rm tot}|_{\ker B_{63}}=0.
 \label{eq:V63-Lagrangian-card}
\end{equation}

\subsection{V64 acquisition order}
\label{subsec:V63-next}

V63 removes the real-cone trace and conormal obstruction.  The next
iteration must decide whether this finite graph is an actual stable
boundary domain.  In order:

\begin{enumerate}
\item analytically continue \(G_\theta,L_\theta\), and the cross-gauge
completion from real null directions to the retarded
Laplace--Fourier roots \((s,k,\lambda)\), using projector-valued symbols
rather than a chosen global polarization frame;
\item prove bounded multiplier estimates for the continued graph and
its complement, including the denominator corresponding to
\(1+\mathfrak c\);
\item form the complete augmented stable matrix for the six induced
graph rows, four harmonic rows, and two continuum traction rows;
\item test its exact ranks and singular floors at normal incidence,
gravity glancing, the continuum threshold, and the static face;
\item realize the cross residual gauge as a recorded boundary
edge/ghost transformation and recheck the second Dirichlet
compatibility of the V58 York projector; and
\item only after those tests, combine the V59 positive quotient form
with the full metric constraint quotient and attempt a closed positive
time-generator domain.
\end{enumerate}

If the analytic continuation of the half-angle graph develops a pole
or loses complementarity, V64 should return that exact symbol and move
upstream to a first-order characteristic graph.  It should not restore
the homogeneous mean-curvature row, whose physical plus-class blow-up
is already proved in Theorem~\ref{thm:V63-K-blowup}.  The next scheduled
inspection of the newest Yang--Mills and Navier--Stokes notes remains
V65, not V64.

\subsection{V63 result ledger and claim boundary}
\label{subsec:V63-ledger}

Relative to V1--V62, this note proves:

\begin{enumerate}
\item the induced trace-free restriction of the canonical V53 TT
representative has a one-dimensional cross loss at gravitational
glancing and every induced-only inverse blows like
\(1/\cos\theta\);
\item a bounded explicit residual null gauge supplies the missing cross
trace and produces a continuous, uniformly split two-plane with
right-inverse norm at most \(\sqrt3\);
\item the associated three-dimensional complement has an explicit
continuous frame on the full planar angular interval;
\item retaining \(K_{(1)}=0\) forces the physical plus representative's
normal gauge amplitude to blow like \(1/\cos\theta\);
\item replacing that row by the physical induced-trace graph folds the
trace--mean-curvature pair into the Brown--York conormal with no
singular gauge;
\item the same GHY--de Donder variation gives exactly the V59 common
trace and outward radiative conormal at every real propagating angle;
and
\item after adjoining the continuum traction, the complete frozen
metric--continuum normal relation has rank twelve in dimension
twenty-four, is isotropic, and is therefore Lagrangian and maximal.
\end{enumerate}

No open-half-plane complementing determinant or uniform stable-cone
floor is proved.  No Sobolev domain, positive full-metric time
generator, off-shell BRST/BV--BFV complex, nonlinear
Brown--York/Bianchi identity, or coupled Higgs, fermion, neutrino, and
Yang--Mills closure is claimed.  The moving graph is presently a
nonzero real-frequency symbol.  The full Standard Model--Einstein
continuum is not proved.

\subsection{Append-only record}

This V63 material was appended to the exact V62 parent whose SHA--256
digest is
\[
\texttt{cca082ab176e2c4503b54d541cc988847f96071619df4c07b2ed91de668ac02a}.
\]
The next research note is V64.  It should analytically continue the
induced graph, assemble the full stable matrix, and decide the
closed-cone floor before the scheduled V65 companion audit.

% END APPENDED VERSION V63 MATERIAL

% BEGIN APPENDED VERSION V64 MATERIAL

\clearpage
\section{V64: a holomorphic half-angle null frame, the exact induced-graph
stable kernel, and a uniformly complementing axial metric lift}
\label{sec:V64-holomorphic-axial-lift}

\subsection{Purpose and correction of the stable acquisition}
\label{subsec:V64-purpose}

V63 constructed a uniformly split induced configuration graph on every
nonzero real null direction and proved that its GHY--de Donder conormal
is exactly the V59 radiative conormal.  It also proved maximality of the
corresponding normal Green relation.  Neither statement is a stable
complementing theorem.  The first task of V64 is therefore to continue
the graph to the retarded Laplace cone and compute the complete stable
kernel rather than infer a determinant from the Lagrangian dimension
count.

Two corrections are required.  Directly replacing
\(\mathfrak c\) by \(\lambda/s\) in the ungauged V63 plus column creates
an artificial second-order pole at the static face \(s=0\).  Applying
the same residual null-gauge completion to both TT polarizations replaces
that quotient by the half-angle Cayley symbol
\[
 \chi=\frac{i|k|}{s+\lambda}.
\]
It is holomorphic in the open half-plane and bounded by one on the
closed stable cone.  This gives a uniformly split analytic TT frame at
normal incidence, gravitational glancing, the continuum threshold, and
the static face.

The second correction is structural.  Six induced-configuration rows
and four harmonic rows retain the V36 normal-gauge stable line at every
open frequency.  At gravitational glancing they retain the full
two-dimensional V36 Dirichlet--harmonic kernel.  The V59 traction rows
cannot remove either zero-trace line.  Thus the V63 induced graph remains
a valid action Green relation but its unquotiented stable matrix has rank
eleven in the open cone and rank ten at gravitational glancing.

V64 supplies a uniformly transverse replacement.  A fixed null reference
\(q\) defines four axial rows \(q^\mu h_{\mu\nu}=0\).  Their incidence on
a residual gauge amplitude has determinant \(2(s+\lambda)^4\).  Together
with the four harmonic rows, two holomorphic TT common-trace rows, and two
V59 traction rows, they form an invertible twelve-by-twelve
metric--continuum stable matrix with a positive singular-value floor on
the complete normalized closed cone.  Pullback through the V39 full
embedding exact sequence gives both a uniformly complementing
sixteen-row split representative and a four-orbit quotient formulation.
The axial selector's off-shell action/BRST realization is not supplied;
that is the scheduled V65 audit and construction task.

\subsection{Retarded roots and the half-angle Cayley symbol}
\label{subsec:V64-half-angle}

Rotate the tangential spatial covector to \(k=(\rho,0)\), \(\rho\geq0\),
and use the V37--V59 retarded roots
\begin{equation}
 \lambda^2=s^2+\rho^2,
 \qquad
 \kappa_v^2=s^2+v^2\rho^2,
 \qquad
 \operatorname{Re}\lambda>0,
 \quad
 \operatorname{Re}\kappa_v>0
 \quad(\operatorname{Re}s>0),
 \label{eq:V64-retarded-roots}
\end{equation}
with their limiting-absorption boundary values.  Here
\(m>0\) and \(0<v<1\) are the V59 continuum parameters.  Exclude only
the zero frequency \((s,\rho)=(0,0)\).

Define
\begin{equation}
 \alpha_+(s,\rho):=s+\lambda,
 \qquad
 \chi(s,\rho):=\frac{i\rho}{s+\lambda}.
 \label{eq:V64-chi-definition}
\end{equation}
The root identity gives
\begin{equation}
 \chi^2=\frac{s-\lambda}{s+\lambda}.
 \label{eq:V64-chi-Cayley}
\end{equation}

\begin{lemma}[Strict retarded root pairing]
\label{lem:V64-root-pairing}
For \(\operatorname{Re}s>0\),
\begin{equation}
 \operatorname{Re}(\overline s\,\lambda)>0,
 \qquad
 |s+\lambda|^2-|s-\lambda|^2
 =4\operatorname{Re}(\overline s\,\lambda)>0.
 \label{eq:V64-root-pairing}
\end{equation}
\end{lemma}

\begin{proof}
Write \(s=\sigma+i\omega\) and \(\lambda=u+iq\), with
\(\sigma,u>0\).  The imaginary part of
\(\lambda^2=s^2+\rho^2\) gives \(uq=\sigma\omega\).  Hence
\[
 \operatorname{Re}(\overline s\lambda)
 =\sigma u+\omega q
 =\frac{\sigma}{u}(u^2+\omega^2)>0.
\]
Expansion of the two moduli gives the second equality.
\end{proof}

\begin{theorem}[Holomorphic half-angle bound]
\label{thm:V64-chi-bound}
The denominator \(s+\lambda\) has no zero at any nonzero open or closed
stable frequency.  The symbol \(\chi\) is holomorphic for
\(\operatorname{Re}s>0\), continuous in the limiting-absorption sense on
the normalized closed cone, and
\begin{align}
 |\chi|&<1 &&(\operatorname{Re}s>0),
 \label{eq:V64-chi-open}\\
 |\chi|&\leq1 &&(\operatorname{Re}s\geq0).
 \label{eq:V64-chi-closed}
\end{align}
Moreover there is a constant \(c_+>0\) such that
\begin{equation}
 |s+\lambda|\geq c_+
 \quad\hbox{when}\quad
 \operatorname{Re}s\geq0,
 \quad |s|^2+\rho^2=1.
 \label{eq:V64-spluslambda-floor}
\end{equation}
\end{theorem}

\begin{proof}
In the open half-plane both \(s\) and \(\lambda\) have positive real
part.  On the closed boundary, \(s+\lambda=0\) would imply
\(\lambda=-s\), hence \(\rho=0\).  The retarded continuation at
\(\rho=0\) has \(\lambda=s\), so a nonzero frequency cannot satisfy
that equation.

By \eqref{eq:V64-chi-Cayley} and
Lemma~\ref{lem:V64-root-pairing},
\[
 |\chi|^2
 =\left|\frac{s-\lambda}{s+\lambda}\right|<1
\]
in the open half-plane.  Limiting absorption gives the closed inequality.
The normalized closed cone, with its retarded-root sheet, is compact.
The continuous nonzero function \(|s+\lambda|\) therefore has the positive
minimum in \eqref{eq:V64-spluslambda-floor}.
\end{proof}

The endpoint values are
\begin{align}
 \rho=0:&\quad \lambda=s,\quad \chi=0,
 \label{eq:V64-chi-normal}\\
 s=\pm i\rho,\ \lambda=0:&\quad \chi=\pm1,
 \label{eq:V64-chi-glancing}\\
 s=0,\ \lambda=\rho:&\quad \chi=i,
 \label{eq:V64-chi-static}\\
 s=\pm iv\rho,\ \kappa_v=0:&\quad
 |\chi|=1.
 \label{eq:V64-chi-threshold}
\end{align}
Thus the half-angle chart does not exchange the normal or glancing pole
for a static or continuum-threshold pole.

\subsection{A bounded complex null frame}
\label{subsec:V64-null-frame}

Use the stable derivative covector, in the outward normal frame,
\begin{equation}
 p_\mu:=(s,i\rho,0,\lambda),
 \qquad
 p^\mu p_\mu=-s^2-\rho^2+\lambda^2=0.
 \label{eq:V64-null-covector}
\end{equation}
Choose the fixed null reference
\begin{equation}
 q^\mu:=(1,0,0,1),
 \qquad
 q_\mu=(-1,0,0,1),
 \qquad
 q^\mu q_\mu=0,
 \qquad
 p_\mu q^\mu=s+\lambda.
 \label{eq:V64-null-reference}
\end{equation}
Let \(e^\mu=(0,0,1,0)\), and define
\begin{equation}
 T^\mu:=(\chi,-1,0,\chi),
 \qquad
 T_\mu=(-\chi,-1,0,\chi).
 \label{eq:V64-bounded-transverse-vector}
\end{equation}
Then
\begin{equation}
 p_\mu T^\mu=0,
 \qquad
 q_\mu T^\mu=0,
 \qquad
 T^\mu T_\mu=1,
 \qquad
 e^\mu T_\mu=0.
 \label{eq:V64-null-frame-identities}
\end{equation}
The first equality is exactly
\((s+\lambda)\chi-i\rho=0\).

Define the covariant polarization tensors
\begin{align}
 {\mathcal E}_+
 &:=\frac{e\otimes e-T\otimes T}{\sqrt2},
 \label{eq:V64-E-plus}\\
 {\mathcal E}_\times
 &:=\frac{e\otimes T+T\otimes e}{\sqrt2}.
 \label{eq:V64-E-cross}
\end{align}
Here the tensors on the right use the covectors \(e_\mu,T_\mu\).

\begin{proposition}[Holomorphic TT splitting]
\label{prop:V64-TT-splitting}
For every nonzero stable frequency,
\begin{align}
 p^\mu{\mathcal E}_{A,\mu\nu}&=0,
 &
 \eta^{\mu\nu}{\mathcal E}_{A,\mu\nu}&=0,
 \label{eq:V64-E-harmonic}\\
 {\mathcal E}_A^{\mu\nu}{\mathcal E}_{B,\mu\nu}
 &=\delta_{AB},
 &&A,B\in\{+,\times\}.
 \label{eq:V64-E-orthonormal}
\end{align}
The entries are holomorphic polynomials of degree at most two in \(\chi\)
and are uniformly bounded on the normalized closed cone.  They give a
split representative of the V53 null cohomology on the retarded complex
null shell.
\end{proposition}

\begin{proof}
Equations \eqref{eq:V64-null-frame-identities} say that \(e,T\) are
orthonormal spacelike vectors transverse to both \(p\) and the reference
\(q\).  The trace, transversality, and bilinear products of the two
standard TT combinations follow by direct contraction.  Their entries
contain only \(1,\chi,\chi^2\), which are holomorphic and bounded by
Theorem~\ref{thm:V64-chi-bound}.  The V53 exact-sequence proof is
algebraic for every nonzero complex null covector, and these two tensors
supply its right splitting.
\end{proof}

On the upper and lower propagating boundary sheets, this frame differs
from the V63 frame only by the corresponding residual null-gauge and
reality orientation.  It therefore represents the same two TT cohomology
classes.  No equality of raw representatives across the two Fourier
orientations is asserted.
\subsection{Uniform induced graph and its GHY conormal}
\label{subsec:V64-induced-graph}

Use the V63 induced trace-free basis
\((F_+,F_\times,F_{01},F_{02},F_s)\).  The induced trace-free parts of
\({\mathcal E}_+,{\mathcal E}_\times\) are the columns
\begin{align}
 {\mathcal G}_+
 &:=-F_++\chi F_{01}-\frac{\chi^2}{\sqrt3}F_s,
 \label{eq:V64-G-plus}\\
 {\mathcal G}_\times
 &:=-F_\times+\chi F_{02},
 \label{eq:V64-G-cross}
\end{align}
and their induced traces are
\begin{equation}
 {\mathcal L}_\chi
 :=\begin{pmatrix}\chi^2/\sqrt2&0\end{pmatrix}.
 \label{eq:V64-L-trace}
\end{equation}
Let \({\mathcal G}_\chi:\mathbb C^2\to
{\cal V}_\partial^{\rm TF}\) denote the two-column injection.

\begin{theorem}[Uniform stable induced splitting]
\label{thm:V64-uniform-induced-splitting}
With the Euclidean coefficient norm,
\begin{equation}
 {\mathcal G}_\chi^*{\mathcal G}_\chi
 =\operatorname{diag}\!\left(
 1+|\chi|^2+\frac{|\chi|^4}{3},
 1+|\chi|^2\right).
 \label{eq:V64-G-Gram}
\end{equation}
Hence, on the complete normalized closed cone,
\begin{equation}
 1\leq s_{\min}({\mathcal G}_\chi)
 \leq s_{\max}({\mathcal G}_\chi)
 \leq\sqrt{\frac73},
 \qquad
 \|({\mathcal G}_\chi^*{\mathcal G}_\chi)^{-1}
       {\mathcal G}_\chi^*\|\leq1.
 \label{eq:V64-G-uniform-bounds}
\end{equation}
The six-component graph
\begin{equation}
 u={\mathcal G}_\chi z,
 \qquad
 \tau_\partial={\mathcal L}_\chi z
 \label{eq:V64-six-induced-graph}
\end{equation}
is holomorphic and bounded at all four distinguished faces.
\end{theorem}

\begin{proof}
The plus column is supported on \(F_+,F_{01},F_s\), while the cross
column is supported on \(F_\times,F_{02}\).  The five basis vectors are
Euclidean orthonormal, which gives the diagonal Gram.  The bounds follow
from \(|\chi|\leq1\).
\end{proof}

The next identity derives the traction without a component expansion.
For harmonic, spacetime trace-free tensors \(E_{\mu\nu}\) and
\(F_{\mu\nu}\) satisfying
\[
 p^\mu E_{\mu\nu}=p^\mu F_{\mu\nu}=0,
\]
let \(E_{ab},F_{ab}\) be their induced blocks and use the V36 linearized
extrinsic curvature with \(D_n=\lambda\).

\begin{lemma}[Transverse Brown--York contraction]
\label{lem:V64-BY-contraction}
For the two tensors in
\eqref{eq:V64-E-plus}--\eqref{eq:V64-E-cross},
\begin{equation}
 \left(K_{(1)}({\mathcal E}_A)^{ab}
 -K_{(1)}({\mathcal E}_A)\eta_\partial^{ab}\right)
 {\mathcal E}_{B,ab}
 =\frac\lambda2\delta_{AB}.
 \label{eq:V64-BY-contraction}
\end{equation}
The action conormal is \(\varepsilon_D\) times the left side of
\eqref{eq:V64-BY-contraction}.  Consequently the V63 equation-normalized
conormal on the graph is
\begin{equation}
 p_{\rm rad}=\lambda z.
 \label{eq:V64-stable-radiative-conormal}
\end{equation}
\end{lemma}

\begin{proof}
Write \(d_a=(s,i\rho,0)\).  Transversality of \(E\) and \(F\) gives
\begin{equation}
 d^aE_{ab}=-\lambda E_{nb},
 \qquad
 d^aF_{ab}=-\lambda F_{nb},
 \label{eq:V64-transverse-components}
\end{equation}
and trace-freeness gives
\(\tau_E+E_{nn}=\tau_F+F_{nn}=0\).  From
\[
 K_{ab}(E)=\frac12(\lambda E_{ab}-d_aE_{nb}-d_bE_{na})
\]
one obtains the polarized identity
\begin{align*}
 K(E)^{ab}F_{ab}
 &=\frac\lambda2
   \left(E^{ab}F_{ab}+2E^{na}F_{na}\right),\\
 K(E)&=\frac\lambda2(\tau_E+2E_{nn}).
\end{align*}
Therefore
\begin{align*}
 (K(E)^{ab}-K(E)\eta^{ab})F_{ab}
 &=\frac\lambda2
 \left(E^{ab}F_{ab}+2E^{na}F_{na}
       -\tau_E\tau_F-2E_{nn}\tau_F\right)\\
 &=\frac\lambda2
 \left(E^{ab}F_{ab}+2E^{na}F_{na}+E_{nn}F_{nn}\right)\\
 &=\frac\lambda2 E^{\mu\nu}F_{\mu\nu}.
\end{align*}
The middle equality uses
\(E_{nn}=-\tau_E\) and \(F_{nn}=-\tau_F\).  Taking
\(E={\mathcal E}_A\), \(F={\mathcal E}_B\), and using
\eqref{eq:V64-E-orthonormal} proves the displayed matrix identity.
V63's equation normalization multiplies the action conormal by
\(2\varepsilon_D\), giving
\eqref{eq:V64-stable-radiative-conormal}.
\end{proof}

For reference, the cross component calculation is
\begin{align}
 K_{02}({\mathcal E}_\times)
 &=-\frac{(s+\lambda)\chi}{2\sqrt2}
 =-\frac{i\rho}{2\sqrt2},
 \label{eq:V64-cross-K02}\\
 K_{12}({\mathcal E}_\times)
 &=-\frac{\lambda+i\rho\chi}{2\sqrt2},
 \label{eq:V64-cross-K12}
\end{align}
and the Lorentzian sign in the \(02\) contraction makes
\({\mathcal G}_\times^\vee\Pi_{\rm TF}=\lambda/2\), in agreement with
the invariant proof.

\subsection{The induced graph is still not complementing}
\label{subsec:V64-induced-kernel}

Let \({\mathcal B}_{\rm ind}\) be the twelve-row stable operator on a
ten-component metric amplitude \(h\) and a two-component outgoing
continuum amplitude \(b\) whose rows are
\begin{equation}
 (h_{ab}^{\rm TF},\tau_\partial)
 =({\mathcal G}_\chi b,{\mathcal L}_\chi b),
 \qquad
 C_\mu(h)=0,
 \qquad
 \lambda b+m\kappa_v b=0.
 \label{eq:V64-induced-stable-operator}
\end{equation}
The last row is the GHY conormal after the first two groups are imposed.

For \(\lambda\ne0\), the V36 normal gauge line has the polynomial
representative
\begin{equation}
 g_D(s,\rho)
 :=\left(
 h_{ab}=0,
 \quad h_{na}=\frac12d_a,
 \quad h_{nn}=\lambda,
 \quad b=0
 \right).
 \label{eq:V64-polynomial-gauge-vector}
\end{equation}
It is \(\lambda r_D\) in the V36 normalization and extends continuously
to gravitational glancing.

\begin{proposition}[Zero Brown--York radiative work of the Dirichlet
kernel]
\label{prop:V64-gauge-zero-work}
The vector \(g_D\) satisfies all rows of
\eqref{eq:V64-induced-stable-operator}.  Its Brown--York conormal paired
with either \({\mathcal E}_A\) is zero.
\end{proposition}

\begin{proof}
The induced metric and continuum trace vanish.  The V36 identities give
\[
 C_a=\lambda(d_a/2)-d_a\lambda/2=0,
 \qquad
 C_n=d^\sharp d/2+\lambda^2/2=0.
\]
For the unscaled V36 line,
\(K_{ab}=-d_ad_b/(2\lambda)\) and \(K=\lambda/2\).
Transversality gives
\(d^ad^b{\mathcal E}_{A,ab}=\lambda^2{\mathcal E}_{A,nn}\), while
trace-freeness gives
\(\tau_\partial=-{\mathcal E}_{A,nn}\).  Hence
\[
 (K^{ab}-K\eta^{ab}){\mathcal E}_{A,ab}
 =-\frac\lambda2
   ({\mathcal E}_{A,nn}+\tau_\partial)=0.
\]
Scaling by \(\lambda\) preserves the zero.  Thus the traction row also
vanishes.
\end{proof}

\begin{theorem}[Exact rank of the induced stable relation]
\label{thm:V64-induced-rank}
At every nonzero stable frequency with \(\lambda\ne0\),
\begin{equation}
 \operatorname{rank}{\mathcal B}_{\rm ind}=11,
 \qquad
 \ker{\mathcal B}_{\rm ind}=\operatorname{span}\{g_D\}.
 \label{eq:V64-induced-open-rank}
\end{equation}
At gravitational glancing \(\lambda=0\), \(d\ne0\),
\begin{equation}
 \operatorname{rank}{\mathcal B}_{\rm ind}=10,
 \qquad
 \ker{\mathcal B}_{\rm ind}
 =\left\{(h_{ab}=0,\ h_{na}=v_a,\ h_{nn}=0,\ b=0):
 d^\sharp v=0\right\}.
 \label{eq:V64-induced-glancing-rank}
\end{equation}
Thus the V63 action Green relation has zero unquotiented stable floor even
after the holomorphic frame correction.
\end{theorem}

\begin{proof}
First omit the two traction rows.  Given a solution of the six graph rows,
subtract \({\mathcal E}_+b_++{\mathcal E}_\times b_\times\).  The result
has zero induced block and remains harmonic.  Writing its remaining
components as \(v_a=h_{na}\), \(\phi=h_{nn}\), the four harmonic rows are
exactly the V36 matrix
\[
 {\mathcal H}_D\binom v\phi
 =\begin{pmatrix}
   \lambda I_3&-d/2\\ d^\sharp&\lambda/2
  \end{pmatrix}\binom v\phi,
 \qquad
 d^\sharp d=-\lambda^2.
\]
For \(\lambda\ne0\), its upper block solves
\(v=d\phi/(2\lambda)\), and the last row then vanishes identically.
Thus it has rank three and kernel \(\operatorname{span}\{r_D\}\) at every
open or closed frequency with \(\lambda\ne0\).  Restoring the two free
TT graph amplitudes shows that the first ten rows on the twelve amplitudes
have rank nine and a three-dimensional kernel consisting of those two
lifted graph directions plus the Dirichlet gauge line.

On the two lifted graph directions the remaining rows are
\[
 (\lambda+m\kappa_v)b.
\]
V59 proves this scalar factor is nonzero at every nonzero open or closed
stable frequency.  The two traction rows therefore remove exactly the
two graph directions.  Proposition~\ref{prop:V64-gauge-zero-work} shows
that they leave \(g_D\), proving \eqref{eq:V64-induced-open-rank}.

At \(\lambda=0\), \(d\ne0\), the two nonzero parts of
\({\mathcal H}_D\) are the independent rank-one maps
\(\phi\mapsto-d\phi/2\) and \(v\mapsto d^\sharp v\).  Hence the fixed
induced plus harmonic metric rows have rank eight and the displayed
two-dimensional kernel.  The two lifted TT graph directions again make
the first ten-row kernel dimension four.  For every vector in that
glancing kernel,
\(K=-d^\sharp v=0\), and
\(d^a{\mathcal E}_{A,ab}=0\); its radiative Brown--York contraction is
therefore zero.  The V59 factor at gravitational glancing is
\(m\kappa_v\ne0\), so the traction rows remove the two nonzero-\(b\)
directions and leave exactly the two zero-trace vectors in
\eqref{eq:V64-induced-glancing-rank}.
\end{proof}

\begin{firewall}[Lagrangian does not imply complementing]
\label{fire:V64-Lagrangian-scope}
Theorem~\ref{thm:V64-induced-rank} does not invalidate the V63 maximal
Green theorem.  A half-dimensional isotropic boundary relation can
intersect the stable subspace.  V63 proved variational maximality in the
full normal phase fibre; V64 proves that its intersection has dimension
one in the open stable cone and two at gravitational glancing.  Both
statements are simultaneously true.
\end{firewall}

\subsection{A uniformly transverse null-reference selector}
\label{subsec:V64-axial-selector}

For an underlying metric amplitude define the four axial rows
\begin{equation}
 ({\mathcal A}_qh)_\nu:=q^\mu h_{\mu\nu},
 \label{eq:V64-axial-row}
\end{equation}
with \(q\) from \eqref{eq:V64-null-reference}.  A residual metric gauge
amplitude \(\epsilon\in\mathbb C^4\) acts by
\begin{equation}
 (\Sigma_p\epsilon)_{\mu\nu}
 :=p_\mu\epsilon_\nu+p_\nu\epsilon_\mu.
 \label{eq:V64-underlying-gauge-map}
\end{equation}
Because \(p^2=0\), this map takes values in the harmonic amplitude space.
Its axial incidence is
\begin{equation}
 {\mathcal A}_q\Sigma_p
 =(s+\lambda)I_4+p\otimes q
 =:{\mathcal M}_q.
 \label{eq:V64-axial-gauge-incidence}
\end{equation}
Here \((q\epsilon)=q^\mu\epsilon_\mu\).

\begin{theorem}[Exact axial gauge transversality]
\label{thm:V64-axial-transversality}
For every nonzero stable frequency,
\begin{align}
 \det{\mathcal M}_q&=2(s+\lambda)^4,
 \label{eq:V64-axial-determinant}\\
 {\mathcal M}_q^{-1}
 &=\frac1{s+\lambda}
 \left(I_4-\frac{p\otimes q}{2(s+\lambda)}\right).
 \label{eq:V64-axial-inverse}
\end{align}
On the normalized closed cone, the inverse norm is uniformly bounded.
Thus the axial row is transverse to all four residual null-gauge
directions at normal incidence, glancing, the continuum threshold, and
the static face.
\end{theorem}

\begin{proof}
The matrix determinant lemma, with
\(q p=s+\lambda\), gives
\[
 \det\bigl((s+\lambda)I_4+p\otimes q\bigr)
 =(s+\lambda)^4
 \left(1+\frac{qp}{s+\lambda}\right)
 =2(s+\lambda)^4.
\]
The Sherman--Morrison formula gives
\eqref{eq:V64-axial-inverse}.  Theorem~\ref{thm:V64-chi-bound} gives a
uniform denominator floor, while \(p\) is bounded on the normalized cone.
\end{proof}

Define the holomorphic TT quotient row
\begin{equation}
 {\mathcal T}_\chi h
 :=\binom{{\mathcal E}_+^{\mu\nu}h_{\mu\nu}}
              {{\mathcal E}_\times^{\mu\nu}h_{\mu\nu}}.
 \label{eq:V64-TT-row}
\end{equation}
Transversality gives
\begin{equation}
 {\mathcal T}_\chi\Sigma_p=0,
 \qquad
 {\mathcal T}_\chi({\mathcal E}_+,{\mathcal E}_\times)=I_2.
 \label{eq:V64-TT-gauge-invariance}
\end{equation}

\begin{lemma}[Axial TT decomposition of the harmonic shell]
\label{lem:V64-axial-TT-decomposition}
Let
\[
 Z_p:=\{h\in\operatorname{Sym}^2\mathbb C^4:C_p(h)=0\}.
\]
Every \(h\in Z_p\) has a unique decomposition
\begin{equation}
 h={\mathcal E}_+a_+ +{\mathcal E}_\times a_\times
   +\Sigma_p\epsilon.
 \label{eq:V64-harmonic-decomposition}
\end{equation}
If also \({\mathcal A}_qh=0\), then \(\epsilon=0\) and
\(a={\mathcal T}_\chi h\).
\end{lemma}

\begin{proof}
This is the V53 split exact sequence written for the underlying metric
rather than its trace reversal.  Trace reversal is an invertible fibre
map, and the two tensors in
Proposition~\ref{prop:V64-TT-splitting} give the two-dimensional right
splitting.  Hence \(Z_p\) is the direct sum of their span and
\(\operatorname{ran}\Sigma_p\).  Both polarization tensors are
orthogonal to \(q\), so applying \({\mathcal A}_q\) to
\eqref{eq:V64-harmonic-decomposition} gives
\({\mathcal M}_q\epsilon=0\).  Theorem~\ref{thm:V64-axial-transversality}
forces \(\epsilon=0\), and
\eqref{eq:V64-TT-gauge-invariance} gives the coefficient formula.
\end{proof}
\subsection{The complete axial metric--continuum stable matrix}
\label{subsec:V64-axial-matrix}

On a ten-component stable metric amplitude \(h\) and a two-component
outgoing continuum amplitude \(b\), define
\begin{equation}
 {\mathcal B}_{64}(s,\rho)
 \binom{h}{b}
 :=\begin{pmatrix}
 C_p(h)\\
 {\mathcal A}_qh\\
 {\mathcal T}_\chi h-b\\
 \lambda{\mathcal T}_\chi h+m\kappa_v b
 \end{pmatrix}.
 \label{eq:V64-complete-boundary-matrix}
\end{equation}
The row orders are respectively four, four, two, and two.  The last row
is the V64 GHY conormal of
Lemma~\ref{lem:V64-BY-contraction} plus the V59 outgoing continuum
traction.

\begin{theorem}[Complete closed-cone complementing theorem]
\label{thm:V64-complete-floor}
For every \(m>0\), \(0<v<1\), and every nonzero open or closed stable
frequency,
\begin{equation}
 \ker{\mathcal B}_{64}(s,\rho)=\{0\},
 \qquad
 \operatorname{rank}{\mathcal B}_{64}=12.
 \label{eq:V64-complete-rank}
\end{equation}
On the normalized closed cone there is a constant
\(\rho_{64}(m,v)>0\) such that
\begin{equation}
 s_{\min}({\mathcal B}_{64}(s,\rho))
 \geq\rho_{64}(m,v).
 \label{eq:V64-complete-floor}
\end{equation}
This includes normal incidence, gravitational glancing, the continuum
threshold, and the static face.
\end{theorem}

\begin{proof}
Suppose \({\mathcal B}_{64}(h,b)=0\).  The first row puts \(h\) in
\(Z_p\), and the second row together with
Lemma~\ref{lem:V64-axial-TT-decomposition} gives
\[
 h={\mathcal E}_+a_+ +{\mathcal E}_\times a_\times.
\]
The third row gives \(a=b\).  The fourth then becomes
\[
 (\lambda+m\kappa_v)b=0.
\]
The V59 closed-cone theorem says this scalar factor never vanishes at a
nonzero stable frequency.  Hence \(b=a=0\) and \(h=0\).  The square
matrix is injective and therefore has rank twelve.

All entries of the normalized matrix are continuous on the compact
closed cone: Theorem~\ref{thm:V64-chi-bound} controls \(\chi\), and the
retarded roots have their fixed limiting-absorption values.  A continuous
family of invertible finite matrices on a compact set has a positive
least-singular-value minimum, which is
\eqref{eq:V64-complete-floor}.
\end{proof}

In a basis adapted to the harmonic contraction, the four gauge columns,
and the two TT columns, the nonzero determinant factors are
\begin{equation}
 \det(C_p|_{Z_p^\perp}),
 \qquad
 2(s+\lambda)^4,
 \qquad
 1,
 \qquad
 (\lambda+m\kappa_v)^2.
 \label{eq:V64-adapted-determinant-factors}
\end{equation}
The first factor denotes any certified V35 complement of the harmonic
kernel.  Equation \eqref{eq:V64-adapted-determinant-factors} is an
adapted-frame certificate, not a claim that its value is invariant under
arbitrary row normalization.

\begin{corollary}[Exact recovery of the V59 quotient symbol]
\label{cor:V64-V59-symbol}
Eliminating the four harmonic and four axial gauge components, the
remaining common trace is \(a=b\), and the Schur complement of
\({\mathcal B}_{64}\) is exactly
\begin{equation}
 {\mathcal B}_{\rm rad}
 =(\lambda+m\kappa_v)I_2.
 \label{eq:V64-V59-Schur}
\end{equation}
Thus the full stable floor is a genuine lift of the V59 quotient
interface, rather than a different two-channel determinant.
\end{corollary}

\subsection{Full embedding pullback and a split representative}
\label{subsec:V64-embedding}

Reuse the V39 full embedding exact sequence without changing its
variables:
\begin{equation}
 \widehat h=Q_4(h,X):=h+{\mathcal S}_pX,
 \qquad
 G_pq=({\mathcal S}_pq,-q).
 \label{eq:V64-Q4-recall}
\end{equation}
Define the quotient boundary operator
\begin{equation}
 {\mathcal B}_{64}^{Q}(h,X,b)
 :={\mathcal B}_{64}(Q_4(h,X),b).
 \label{eq:V64-quotient-pullback}
\end{equation}

\begin{proposition}[Embedding quotient kernel and floor]
\label{prop:V64-embedding-quotient}
At every nonzero stable frequency,
\begin{equation}
 \ker{\mathcal B}_{64}^{Q}
 =\left\{G_pq:q\in\mathbb C^4\right\}
 \quad\hbox{with }b=0.
 \label{eq:V64-embedding-kernel}
\end{equation}
The induced operator on the quotient by \(\operatorname{ran}G_p\) is
exactly \({\mathcal B}_{64}\) and has the uniform floor of
Theorem~\ref{thm:V64-complete-floor}.  The orbit bundle is uniformly
complemented by the V39 projector.
\end{proposition}

\begin{proof}
V39 gives \(Q_4G_p=0\), so the displayed orbit lies in the kernel.
Conversely, if the pullback vanishes, Theorem~\ref{thm:V64-complete-floor}
forces \(Q_4(h,X)=0\) and \(b=0\).  Exactness of the V39 sequence then
puts \((h,X)\) in \(\operatorname{ran}G_p\).  The quotient statement and
uniform complement are precisely the V39 splitting and orbit-angle
bound combined with the V64 floor.
\end{proof}

For a square split representative, append the four edge-coordinate rows
\(X=0\):
\begin{equation}
 \widetilde{\mathcal B}_{64}(h,X,b)
 :=\binom{{\mathcal B}_{64}(Q_4(h,X),b)}{X}.
 \label{eq:V64-split-embedding-matrix}
\end{equation}

\begin{theorem}[Uniform sixteen-row embedding representative]
\label{thm:V64-embedding-floor}
The matrix \(\widetilde{\mathcal B}_{64}\) has rank sixteen at every
nonzero open or closed stable frequency.  On the normalized cone,
\begin{equation}
 \inf s_{\min}(\widetilde{\mathcal B}_{64})>0.
 \label{eq:V64-embedding-floor}
\end{equation}
\end{theorem}

\begin{proof}
The triangular coordinate map
\[
 (h,X,b)\longmapsto(Q_4(h,X),X,b)
\]
is invertible with inverse
\((\widehat h,X,b)\mapsto(\widehat h-{\mathcal S}_pX,X,b)\).
Both maps are uniformly bounded on the normalized cone because
\({\mathcal S}_p\) is linear in the bounded stable covector.  In these
coordinates, \(\widetilde{\mathcal B}_{64}\) is the direct sum of
\({\mathcal B}_{64}\) and \(I_4\).  Theorem
\ref{thm:V64-complete-floor} proves the rank and uniform floor.
\end{proof}

\begin{firewall}[The split edge row is not yet a BRST action domain]
\label{fire:V64-embedding-scope}
The row \(X=0\) is an algebraic choice of the V39 split representative.
It is not invariant under \(sX=-c\) and has not been derived from a
boundary gauge fermion together with ghost, antighost, and
Nakanishi--Lautrup rows.  Likewise, the axial row is uniformly transverse
and gives the correct quotient, but V64 has not supplied its multiplier
action or proved off-shell tangency.  Theorem
\ref{thm:V64-embedding-floor} is therefore a stable-symbol theorem, not a
complete BV--BFV or positive time-generator theorem.
\end{firewall}

\subsection{Finite certificates and endpoint matrix tests}
\label{subsec:V64-certificates}

The analytic cards are
\begin{align}
 (s+\lambda)\chi-i\rho&=0,
 &
 \chi^2-\frac{s-\lambda}{s+\lambda}&=0,
 \label{eq:V64-chi-card}\\
 p^2=q^2=p\cdot T=q\cdot T&=0,
 &
 T^2=e^2&=1,
 \label{eq:V64-frame-card}\\
 {\mathcal E}_A:{\mathcal E}_B&=\delta_{AB},
 &
 p\mathbin\lrcorner{\mathcal E}_A&=0.
 \label{eq:V64-polarization-card}
\end{align}
The induced and conormal cards are
\begin{align}
 {\mathcal G}_\chi^*{\mathcal G}_\chi
 &-\operatorname{diag}\!\left(
 1+|\chi|^2+|\chi|^4/3,
 1+|\chi|^2\right)=0,
 \label{eq:V64-Gram-card}\\
 2\left(K^{ab}-K\eta^{ab}\right){\mathcal E}_{A,ab}
 -\lambda&=0.
 \label{eq:V64-BY-card}
\end{align}
The axial cards are
\begin{align}
 \det{\mathcal M}_q-2(s+\lambda)^4&=0,
 \label{eq:V64-axial-det-card}\\
 {\mathcal M}_q^{-1}{\mathcal M}_q-I_4&=0.
 \label{eq:V64-axial-inverse-card}
\end{align}
The two rank cards are
\begin{align}
 \operatorname{rank}{\mathcal B}_{\rm ind}
 &=\begin{cases}11,&\lambda\ne0,\\10,&\lambda=0,\ d\ne0,
 \end{cases}
 \label{eq:V64-induced-rank-card}\\
 \operatorname{rank}{\mathcal B}_{64}&=12,
 \qquad
 \operatorname{rank}\widetilde{\mathcal B}_{64}=16.
 \label{eq:V64-complete-rank-card}
\end{align}

For \(m=1\), \(v=1/2\), a finite validator should test at least:
\begin{enumerate}
\item an open rational-root point
\((s,\rho,\lambda)=(3/4,1,5/4)\), where \(\chi=i/2\);
\item normal incidence \((s,\rho,\lambda)=(1,0,1)\);
\item the static face \((s,\rho,\lambda)=(0,1,1)\), where \(\chi=i\);
\item gravitational glancing \((s,\rho,\lambda)=(i,1,0)\), where
\(\chi=1\); and
\item the continuum threshold
\((s,\rho,\kappa_v)=(i/2,1,0)\).
\end{enumerate}
At every sample the twelve singular values of \({\mathcal B}_{64}\)
must be nonzero.  The same samples applied to
\({\mathcal B}_{\rm ind}\) must return the ranks in
\eqref{eq:V64-induced-rank-card}; otherwise the calculation has hidden
the exact gauge/cohomology kernel rather than repaired it.

\subsection{V65 acquisition order and scheduled companion audit}
\label{subsec:V64-next}

V65 is a five-version audit.  Before changing the axial or embedding
rows, it must inspect the newest completed Yang--Mills and Navier--Stokes
notes in the shared folder and record only results which alter one of the
following remaining gates.

After that audit, V65 should:
\begin{enumerate}
\item derive a boundary multiplier or gauge-fermion action whose Euler
rows produce the axial selector on the invariant V39 representative
\(\widehat h=Q_4(h,X)\);
\item replace the algebraic row \(X=0\) by a complete ghost, antighost,
multiplier, and edge domain with the same sixteen-row quotient floor;
\item prove the corresponding Noether/Ward identity and normal Green
cancellation off shell;
\item determine whether the axial boundary symbol is a causal Sobolev
multiplier in a first-order time realization, rather than relying only on
Laplace holomorphy;
\item combine that domain with the V59 positive quotient form and test a
closed positive full-metric constraint quotient generator; and
\item only then begin the nonlinear Brown--York stress/Bianchi lift and
restore the coupled Standard-Model source sectors.
\end{enumerate}

If no finite BRST gauge fermion produces the axial rows without doubling
normal ghost data, V65 should retain the exact stable quotient theorem and
move upstream to a characteristic edge action.  It should not return to
the induced graph as an unquotiented stable boundary matrix, because
Theorem~\ref{thm:V64-induced-rank} has fixed its exact kernel dimensions.

\subsection{V64 result ledger and claim boundary}
\label{subsec:V64-ledger}

Relative to V1--V63, this note proves:
\begin{enumerate}
\item the half-angle Cayley symbol \(\chi=i\rho/(s+\lambda)\) is
holomorphic in the open half-plane, bounded by one on the complete closed
stable cone, and has no normal, glancing, threshold, or static pole;
\item it gives an explicit bounded complex null frame, two holomorphic TT
polarizations, and a uniformly split six-component induced graph;
\item the GHY Brown--York conormal of either normalized polarization is
exactly \(\lambda/2\), so the equation-normalized traction is the V59
conormal \(\lambda a\);
\item the V63 induced action relation nevertheless has stable rank eleven
away from gravitational glancing and rank ten at glancing, with its exact
one- and two-dimensional zero-trace kernels;
\item the null-reference axial gauge incidence has determinant
\(2(s+\lambda)^4\) and a uniformly bounded inverse;
\item the twelve-row harmonic, axial, common-TT-trace, and continuum
traction matrix is invertible with a positive closed-cone singular-value
floor and has the exact V59 interface as its Schur complement; and
\item the V39 embedding pullback has precisely the four gauge-orbit kernel,
while its split sixteen-row representative has a uniform closed-cone
floor.
\end{enumerate}

The axial and split-embedding rows have not yet been obtained from a
complete boundary gauge fermion or off-shell BRST action.  No positive
full-metric time generator, nonlinear stress/Bianchi identity, curved
boundary atlas, or coupled Higgs, fermion, neutrino, and Yang--Mills
closure is proved.  The full Standard Model--Einstein continuum is not
proved.

\subsection{Append-only record}

This V64 material was appended to the exact V63 parent whose SHA--256
digest is
\[
\texttt{05b79e94d18dd2a9868a891988f4ae1164c364872db5599bb5082fa5f72437b6}.
\]
The next research note is V65.  It begins with the required inspection of
the newest retained Yang--Mills and Navier--Stokes notes, then attempts the
action and BRST realization of the uniformly complementing axial lift.

% END APPENDED VERSION V64 MATERIAL

% BEGIN APPENDED VERSION V65 MATERIAL

\clearpage
\section{V65: companion audit, the axial action defect, a full TT
configuration graph, and characteristic BRST localization}
\label{sec:V65-full-graph}

\subsection{Purpose and the five-version companion audit}
\label{subsec:V65-purpose-audit}

V64 solved the closed stable-symbol problem but deliberately did not call its
axial matrix an action domain.  Its twelve rows were four harmonic rows,
four null-reference axial rows, two common-TT rows, and two continuum
traction rows.  Their stable matrix is uniformly invertible, yet that fact
does not decide whether the corresponding twelve-dimensional normal-phase
relation is isotropic for the V36 GHY--de Donder Green form.

The compulsory five-version audit was performed before choosing the V65
branch and repeated at the final V65 snapshot.  The newest completed
companion files inspected were
\begin{center}
\begin{tabular}{p{0.56\textwidth}p{0.36\textwidth}}
\texttt{Renewal\_Yang\_Mills\_Mass\_Gap\_Research\_Notes\_V46.tex}
 &\scriptsize\ttfamily 1349a5d285324394b159dfdcf15e72c\newline
 bfb3df73d4844273c59ff87118fa49e90\\[0.6em]
\texttt{Renewal\_NS\_Stability\_Research\_Notes\_V31.tex}
 &\scriptsize\ttfamily
 f1121b62238ad5491bb7cf03635ea993\newline
 4942edf573ed8faaa9ee79e1d38a6696
\end{tabular}
\end{center}
The completed Yang--Mills V46 resolves the normalization defect found at V45.
The exact Peter--Weyl product coefficient is \(d_{\rm in}/d_T\); the output
factor \(d_T\) cancels only upon taking the Fourier-algebra norm, leaving raw
multiplicity-resolved pinching mass.  That identity restores the one-link and
disjoint-overlap estimates and closes the pair-only ternary atlas without an
artificial representation-formation law.  The mixed common-junction spectator
sector and global ternary MTA remain open.  The transferable instruction is
therefore precise: preserve the exact coefficient, defect operators,
multiplicity contraction, and final resolvent until their common norm is
taken; do not manufacture a scalar weight before that assembly.

The Navier--Stokes file proves an exact dyadic formation ladder and then
proves its boundary: a marked bounded family does not recover the unmarked
critical first moment without restoring the missing scale, and the resulting
universal estimate is the old continuation stock.  The remaining possible
gain must retain a correlation which the separate norm estimates discard.

Neither companion supplies an Einstein boundary form, a de Donder ghost
complex, or a Standard-Model stress identity.  Their type-correct effect on
V65 is methodological and decisive: the normalized V64 stable matrix must be
compared with the actual V36 conormal before it is promoted to an action, and
the retarded half-angle multiplier must retain its characteristic
realization rather than be treated as a harmless algebraic mark.

\begin{theorem}[V65 companion transfer decision]
\label{thm:V65-companion-decision}
The companion audit changes the acquisition in exactly three ways.
\begin{enumerate}
\item The V36 Green form, not the rank of the normalized stable matrix,
      decides action compatibility.
\item A frozen conormal graph may use the V64 retarded splitting, but a
      spacetime claim must retain the forward characteristic transport and
      its formal-adjoint sector.
\item When the axial and harmonic rows fail the first test, V65 moves
      upstream to the full TT configuration graph instead of appending a
      formal gauge-fermion label to those rows.
\end{enumerate}
No Yang--Mills or Navier--Stokes inequality is imported.
\end{theorem}

\begin{firewall}[Companion claims remain in their own types]
\label{fire:V65-companion-types}
The V45 representation-dimension law is not a metric conormal, and the V31
formation mark is not a causal boundary projector.  V65 uses only their
normalization discipline.  In particular, the provisional V45
formation-first claims are not imported: the later V45 coefficient
counterexample and withdrawal control the audit.
\end{firewall}

\subsection{Canonical gauge coordinate in the V36 boundary one-form}
\label{subsec:V65-canonical-gauge-coordinate}

Write the boundary indices as \(a\in\{0,1,2\}\), let \(n=3\), and retain
\[
 \tau_\partial:=\eta_\partial^{ab}h_{ab},
 \qquad
 {\mathcal Q}^{\mu\nu}(C)
 =n^{(\mu}C^{\nu)}-\frac12\eta^{\mu\nu}C_n.
\]
Define four configuration coordinates
\begin{equation}
 y_a(h):=-2h_{na},
 \qquad
 y_n(h):=\tau_\partial-h_{nn}.
 \label{eq:V65-canonical-y}
\end{equation}

\begin{lemma}[Exact canonical splitting of the de Donder one-form]
\label{lem:V65-canonical-splitting}
For every variation \(\dot h\),
\begin{equation}
 -2{\mathcal Q}^{\mu\nu}(C)\dot h_{\mu\nu}
 =C^a\dot y_a+C_n\dot y_n.
 \label{eq:V65-QC-canonical}
\end{equation}
Consequently the V36 quadratic boundary one-form is
\begin{equation}
 2\kappa\Theta_{36}(h;\dot h)
 =\int_{\partial M}
 \left(\Pi_{(1)}^{ab}(h)\dot h_{ab}
       +C^a(h)\dot y_a+C_n(h)\dot y_n\right).
 \label{eq:V65-canonical-one-form}
\end{equation}
The normal-derivative Legendre map
\begin{equation}
 D_nh\longmapsto(\Pi_{(1)}^{ab},C^a,C_n)
 \label{eq:V65-Legendre-map}
\end{equation}
is an isomorphism at every fixed tangential jet.
\end{lemma}

\begin{proof}
Using symmetry of \(\dot h\) and \(n^\mu=\delta_n^\mu\),
\begin{align*}
 -2{\mathcal Q}^{\mu\nu}(C)\dot h_{\mu\nu}
 &=-2C^a\dot h_{na}-2C_n\dot h_{nn}
   +C_n(\dot\tau_\partial+\dot h_{nn})\\
 &=C^a\dot y_a+C_n\dot y_n.
\end{align*}
Substitution in the V36 one-form proves
\eqref{eq:V65-canonical-one-form}.  The normal principal part of the
minimal density is the nondegenerate DeWitt form
\(\mathbb G(D_nh,D_n\dot h)\).  V36 proved that its fibre matrix has
signature \((6,4)\), hence is invertible.  The displayed variables are just
the coefficients of that normal Legendre transform, proving the last claim.
\end{proof}

Thus \(C=0\) is a Neumann condition on the four \(y\)-coordinates.  The six
induced components are the complementary configuration coordinates.  This
identifies the precise error in treating four additional axial
configuration rows as though they could be adjoined without changing the
Green polarization.

\subsection{The V64 axial--harmonic relation is not isotropic}
\label{subsec:V65-V64-nonisotropic}

At normal incidence \(\chi=0\), let \({\mathcal T}_0\) be contraction with
the two V64 tensors.  Consider the putative normal-phase relation whose rows
are
\begin{equation}
 {\mathcal A}_qh=0,
 \qquad C(h)=0,
 \qquad {\mathcal T}_0h=b,
 \qquad {\mathcal T}_0^\vee\Pi_{(1)}+t=0.
 \label{eq:V65-putative-V64-action-relation}
\end{equation}
The last row denotes the action-normalized version of the two V64 traction
rows; its common nonzero normalization will not affect isotropy.

\begin{theorem}[Exact action defect of the V64 row set]
\label{thm:V65-V64-action-defect}
The relation \eqref{eq:V65-putative-V64-action-relation} is not isotropic for
the V36 metric Green form plus the continuum canonical form.  Hence no
quadratic boundary action has exactly the twelve V64 bosonic rows as a
Lagrangian conormal domain.
\end{theorem}

\begin{proof}
The equations \(q^\mu h_{\mu\nu}=0\), with
\(q^\mu=(1,0,0,1)\), solve the normal configuration components uniquely in
terms of the induced block:
\begin{equation}
 h_{na}=-h_{0a},
 \qquad h_{nn}=h_{00}.
 \label{eq:V65-axial-lift}
\end{equation}
Let \(J_q\gamma\) denote this lift.  In the induced order \((0,1,2)\), take
\begin{equation}
 \gamma_U=\operatorname{diag}(0,1,1),
 \qquad
 W^{ab}=\operatorname{diag}(0,1,1).
 \label{eq:V65-action-defect-witness}
\end{equation}
Both are orthogonal to the normal-incidence plus and cross TT tensors.

By the Legendre isomorphism in
Lemma~\ref{lem:V65-canonical-splitting}, there are boundary jets \(U,V\)
with
\begin{align*}
 &h_U=J_q\gamma_U,\quad C(U)=0,\quad\Pi(U)=0,\quad b_U=t_U=0,\\
 &h_V=0,\quad C(V)=0,\quad\Pi(V)=W,\quad b_V=t_V=0.
\end{align*}
They satisfy every row in
\eqref{eq:V65-putative-V64-action-relation}.  Equation
\eqref{eq:V65-canonical-one-form} gives
\begin{equation}
 2\kappa\Omega_{36}\(U,V\)
 =-W^{ab}\gamma_{U,ab}=-2\ne0.
 \label{eq:V65-action-defect-value}
\end{equation}
The continuum contribution is zero.  This proves non-isotropy.  A Hessian
conormal domain is Lagrangian, so the last assertion follows.
\end{proof}

\begin{assessment}[What the defect changes]
\label{ass:V65-defect-meaning}
V64's rank-twelve theorem remains correct as a stable-symbol theorem.  The
defect is variational: the four axial configuration rows and four harmonic
Neumann rows eliminate both halves of the \(y,C\) gauge sector, while the
four induced directions orthogonal to the TT trace retain both their
configuration and momentum.  The dimensions happen to balance, but the
Green form does not vanish.  The repair must redistribute the rows, not
renormalize them.
\end{assessment}
\subsection{The full TT configuration graph}
\label{subsec:V65-full-TT-graph}

Let
\begin{equation}
 {\mathcal E}_\chi:
 \mathbb C^2\longrightarrow\operatorname{Sym}^2\mathbb C^4,
 \qquad
 {\mathcal E}_\chi b
 :={\mathcal E}_+b_++{\mathcal E}_\times b_\times
 \label{eq:V65-full-E-map}
\end{equation}
be the bounded V64 holomorphic polarization map.  Its columns obey
\begin{equation}
 p\mathbin\lrcorner{\mathcal E}_A=0,
 \qquad
 \operatorname{tr}_\eta{\mathcal E}_A=0,
 \qquad
 q\mathbin\lrcorner{\mathcal E}_A=0,
 \qquad
 {\mathcal E}_A:{\mathcal E}_B=\delta_{AB}.
 \label{eq:V65-E-identities}
\end{equation}
Replace the V64 row set by the ten configuration equations
\begin{equation}
 {\mathcal D}_{65}(h,b):=h-{\mathcal E}_\chi b=0.
 \label{eq:V65-full-configuration-graph}
\end{equation}

\begin{proposition}[Induced graph plus axial completion equals the full graph]
\label{prop:V65-full-graph-equivalence}
At every nonzero stable frequency,
\eqref{eq:V65-full-configuration-graph} is equivalent to
\begin{equation}
 h_{ab}=({\mathcal E}_\chi b)_{ab},
 \qquad
 {\mathcal A}_qh=0.
 \label{eq:V65-induced-axial-equivalence}
\end{equation}
It implies the V64 induced trace-free and trace graph.  Conversely, those
six induced equations together with the four axial equations determine all
ten components of \(h\) and give
\eqref{eq:V65-full-configuration-graph}.
\end{proposition}

\begin{proof}
The forward implication follows from
\(q\mathbin\lrcorner{\mathcal E}_A=0\).  Conversely the induced equations
fix \(h_{ab}\).  The four axial equations then give
\eqref{eq:V65-axial-lift}.  The tensor
\({\mathcal E}_\chi b\) has the same induced block and satisfies the same
axial equations, so uniqueness of that lift makes the two full tensors
equal.
\end{proof}

The full graph has a uniform configuration angle.  Indeed, if the metric
components carry the Euclidean Frobenius norm, then 
\begin{equation}
 \|{\mathcal E}_+\|_{\rm E}\le2\sqrt2,
 \qquad
 \|{\mathcal E}_\times\|_{\rm E}\le\sqrt6,
 \qquad
 \|{\mathcal E}_\chi\|\le\sqrt{14},
 \label{eq:V65-E-Euclidean-bound}
\end{equation}
because \(\|T\|_{\rm E}^2=1+2|\chi|^2\le3\).  The row matrix
\([I_{10},-{\mathcal E}_\chi]\) has all ten nonzero singular values at
least one.

\subsection{The frozen conormal action and maximal Green relation}
\label{subsec:V65-conormal-action}

Let \((\mathfrak p_h,\mathfrak p_b)\) be the metric and continuum action
conormals in the V36/V63 common normalization.  At one frozen oriented
frequency, write \({\mathcal E}_\chi^\vee\) for the bilinear transpose.
Introduce a symmetric-tensor multiplier \(\Lambda\) and the frozen
quadratic boundary functional
\begin{equation}
 F_{65}^{\rm fr}(h,b,\Lambda)
 :=\Lambda^\vee(h-{\mathcal E}_\chi b).
 \label{eq:V65-frozen-multiplier-action}
\end{equation}
Its Euler equations, after adding the bulk and continuum boundary
one-forms, are
\begin{equation}
 h-{\mathcal E}_\chi b=0,
 \qquad
 \mathfrak p_h+\Lambda=0,
 \qquad
 \mathfrak p_b-{\mathcal E}_\chi^\vee\Lambda=0.
 \label{eq:V65-multiplier-Euler-rows}
\end{equation}
Eliminating \(\Lambda\) gives
\begin{equation}
 {\mathcal L}_{65}^{\rm fr}
 :=\left\{,h={\mathcal E}_\chi b,\quad
 \mathfrak p_b+{\mathcal E}_\chi^\vee\mathfrak p_h=0,\right\}.
 \label{eq:V65-frozen-conormal-relation}
\end{equation}

\begin{theorem}[Maximal frozen action relation]
\label{thm:V65-frozen-Lagrangian}
For each fixed nonzero stable frequency,
\({\mathcal L}_{65}^{\rm fr}\) is Lagrangian in the
\(24\)-dimensional metric--continuum normal phase fibre.  It is the
conormal bundle of the two-dimensional configuration graph
\eqref{eq:V65-full-configuration-graph}.
\end{theorem}

\begin{proof}
For two tangent variations, \(\dot h={\mathcal E}_\chi\dot b\).  The
canonical one-form restricts to
\[
 \mathfrak p_h^\vee\dot h+\mathfrak p_b^T\dot b
 =({\mathcal E}_\chi^\vee\mathfrak p_h+\mathfrak p_b)^T\dot b=0.
\]
Its antisymmetrization therefore vanishes.  The ten configuration rows are
independent because their metric block is \(I_{10}\); modulo them, the two
conormal rows are independent because their continuum-momentum block is
\(I_2\).  The relation has codimension twelve and dimension twelve, so its
isotropy is maximal.  Equations
\eqref{eq:V65-multiplier-Euler-rows} also prove the conormal-bundle claim
directly.
\end{proof}

This theorem corrects the V64 variational defect without changing the
physical frame.  It replaces the separate harmonic and axial rows by the
single full configuration graph.  On stable solutions the graph itself
forces the metric amplitude to be TT, so the harmonic contraction vanishes
there; no harmonic row is being silently added to the action relation.

\subsection{Complete stable matrix and exact V59 Schur complement}
\label{subsec:V65-stable-matrix}

Let \(\mathfrak p_{36}^{\rm eq}(h)\) be the full ten-component V36 metric
conormal after the common V63 equation normalization.  Thus it includes the
Brown--York induced block and the de Donder term
\(-2{\mathcal Q}(C)\), with the same fixed invertible scalar multiplier on
both.  Define the twelve-row stable matrix
\begin{equation}
 {\mathcal B}_{65}(s,\rho)
 \binom h b
 :=\binom{
 h-{\mathcal E}_\chi b
 }{
 {\mathcal E}_\chi^\vee\mathfrak p_{36}^{\rm eq}(h)
       +m\kappa_v b
 }.
 \label{eq:V65-stable-matrix}
\end{equation}
The row orders are ten and two.

\begin{lemma}[Full-graph conormal card]
\label{lem:V65-full-conormal-card}
For every \(a\in\mathbb C^2\),
\begin{equation}
 C_p({\mathcal E}_\chi a)=0,
 \qquad
 {\mathcal E}_\chi^\vee
 \mathfrak p_{36}^{\rm eq}({\mathcal E}_\chi a)
 =\lambda a.
 \label{eq:V65-full-conormal-card}
\end{equation}
\end{lemma}

\begin{proof}
The first identity is transversality and trace-freeness in
\eqref{eq:V65-E-identities}.  Hence the de Donder part of the full conormal
vanishes.  The polarized V64 Brown--York identity gives
\(\lambda\delta_{AB}/2\) in action units.  Multiplication by the common
V63 equation factor \(2\varepsilon_D\), including the
\(\varepsilon_D\) already present in the action conormal, gives
\(\lambda\delta_{AB}\).  Linearity proves the result.
\end{proof}

\begin{theorem}[Action-compatible complete closed-cone floor]
\label{thm:V65-complete-floor}
For every \(m>0\), \(0<v<1\), and every nonzero open or closed stable
frequency,
\begin{equation}
 \ker{\mathcal B}_{65}=\{0\},
 \qquad
 \operatorname{rank}{\mathcal B}_{65}=12.
 \label{eq:V65-complete-rank}
\end{equation}
On the normalized closed cone there is
\(\delta_{65}(m,v)>0\) such that
\begin{equation}
 s_{\min}({\mathcal B}_{65}(s,\rho))\ge\delta_{65}(m,v).
 \label{eq:V65-complete-floor}
\end{equation}
In coordinates adapted to the configuration graph, the only physical
Schur factor is
\begin{equation}
 (\lambda+m\kappa_v)I_2.
 \label{eq:V65-V59-Schur}
\end{equation}
\end{theorem}

\begin{proof}
Put \(r=h-{\mathcal E}_\chi b\), so
\(h=r+{\mathcal E}_\chi b\).  Lemma
\ref{lem:V65-full-conormal-card} changes
\eqref{eq:V65-stable-matrix} to the block-triangular system
\begin{equation}
 r=0,
 \qquad
 {\mathcal E}_\chi^\vee\mathfrak p_{36}^{\rm eq}(r)
 +(\lambda+m\kappa_v)b=0.
 \label{eq:V65-triangular-system}
\end{equation}
The V59 theorem says the scalar factor is nonzero at every nonzero closed
stable frequency.  Thus \(r=b=0\), proving injectivity and rank.

On the normalized cone, the changes
\((h,b)\leftrightarrow(r,b)\) are uniformly bounded by
\eqref{eq:V65-E-Euclidean-bound}.  The conormal coefficient in the lower
left block is continuous and bounded because \(p,\chi\) are bounded.  V59
gives a positive minimum of
\(|\lambda+m\kappa_v|\).  The inverse of the triangular matrix is therefore
uniformly bounded, proving \eqref{eq:V65-complete-floor} and the Schur
statement.
\end{proof}

\begin{corollary}[The V64 residual kernels are removed in action form]
\label{cor:V65-kernels-removed}
The V36 normal gauge line and both members of its glancing plane have
nonzero full configuration defect unless they vanish.  In particular, the
physical zero-induced cross vector at glancing is acquired by the normal
components of \eqref{eq:V65-full-configuration-graph}; the repair does not
classify it as gauge or discard its TT cohomology class.
\end{corollary}

\begin{proof}
A vector in the kernel of the full configuration defect has
\(h={\mathcal E}_\chi b\).  The traction row and V59 factor then force
\(b=0\), and hence \(h=0\).  This applies at glancing as well as away from
it.  The full graph includes the V64 cross normal components, so their
physical class is represented rather than quotiented out.
\end{proof}
\subsection{Local characteristic realization of the half-angle frame}
\label{subsec:V65-characteristic-localization}

The full graph is holomorphic and causal as a stable multiplier, but its
formula contains the square-root root \(\lambda\).  The half-angle identity
shows exactly how to retain the missing characteristic scale.  Put
\begin{equation}
 \alpha_+:=s+\lambda,
 \qquad
 y=\binom{y_+}{y_\times},
 \qquad z\in\mathbb C,
 \label{eq:V65-edge-variables}
\end{equation}
and impose the three transport rows
\begin{equation}
 \alpha_+y=i\rho b,
 \qquad
 \alpha_+z=i\rho y_+.
 \label{eq:V65-transport-symbol}
\end{equation}
Define the local algebraic metric chart, in the order \((0,1,2,n)\), by
\begin{equation}
 {\mathfrak E}(b,y,z)
 :=\frac1{\sqrt2}
 \begin{pmatrix}
 -z&-y_+&-y_\times&z\\
 -y_+&-b_+&-b_\times&y_+\\
 -y_\times&-b_\times&b_+&y_\times\\
 z&y_+&y_\times&-z
 \end{pmatrix}.
 \label{eq:V65-local-metric-chart}
\end{equation}

\begin{theorem}[Exact characteristic localization]
\label{thm:V65-characteristic-localization}
At every nonzero open or closed stable frequency,
\eqref{eq:V65-transport-symbol} has the unique solution
\begin{equation}
 y=\chi b,
 \qquad z=\chi^2b_+,
 \label{eq:V65-edge-solution}
\end{equation}
and
\begin{equation}
 {\mathfrak E}(b,y,z)={\mathcal E}_\chi b.
 \label{eq:V65-local-equals-E}
\end{equation}
Moreover
\begin{equation}
 q^\mu{\mathfrak E}_{\mu\nu}=0
 \label{eq:V65-local-axial-identity}
\end{equation}
off the transport shell, for arbitrary \((b,y,z)\).
\end{theorem}

\begin{proof}
V64 proved that \(\alpha_+\ne0\) and has a positive floor on the normalized
closed cone.  Division in the first transport row gives \(y=\chi b\), and
the second gives \(z=\chi y_+=\chi^2b_+\).  Substitution in
\eqref{eq:V65-local-metric-chart} gives, component by component,
\[
 \frac1{\sqrt2}
 \begin{pmatrix}
 -\chi^2b_+&-\chi b_+&-\chi b_\times&\chi^2b_+\\
 -\chi b_+&-b_+&-b_\times&\chi b_+\\
 -\chi b_\times&-b_\times&b_+&\chi b_\times\\
 \chi^2b_+&\chi b_+&\chi b_\times&-\chi^2b_+
 \end{pmatrix},
\]
which is the expansion of the two V64 tensors.  Adding the zeroth and normal
rows of \eqref{eq:V65-local-metric-chart} gives zero, proving the last
identity without using the transport equations.
\end{proof}

In physical derivatives, the transport rows are the local characteristic
cascade
\begin{equation}
 (D_0+D_n)y=D_1b,
 \qquad
 (D_0+D_n)z=D_1y_+,
 \label{eq:V65-characteristic-cascade}
\end{equation}
because \(q\cdot p=s+\lambda\) and \(D_1\) has symbol \(i\rho\).
Thus the two powers of \(\chi\) require two successive outgoing null
transports, not a hidden inverse-scale estimate.

Define the fifteen-row localized stable system on
\((h,b,y,z)\in\mathbb C^{10}\oplus\mathbb C^2\oplus
\mathbb C^2\oplus\mathbb C\) by
\begin{equation}
 {\mathcal B}_{65}^{\rm loc}
 \begin{pmatrix}h\\b\\y\\z\end{pmatrix}
 :=\begin{pmatrix}
 h-{\mathfrak E}(b,y,z)\\
 \alpha_+y-i\rho b\\
 \alpha_+z-i\rho y_+\\
 (\lambda+m\kappa_v)b
 \end{pmatrix}.
 \label{eq:V65-local-stable-matrix}
\end{equation}

\begin{theorem}[Uniform local characteristic stable floor]
\label{thm:V65-local-floor}
The matrix \({\mathcal B}_{65}^{\rm loc}\) has rank fifteen at every
nonzero open or closed stable frequency and a positive least-singular-value
floor on the normalized closed cone.  In triangular coordinates its
diagonal factors are
\begin{equation}
 I_{10},
 \qquad
 \alpha_+ I_2,
 \qquad
 \alpha_+,
 \qquad
 (\lambda+m\kappa_v)I_2.
 \label{eq:V65-local-diagonal-factors}
\end{equation}
\end{theorem}

\begin{proof}
The two transport equations first give
\eqref{eq:V65-edge-solution}.  The metric equation then gives
\(h={\mathcal E}_\chi b\), and the final row forces \(b=0\) by V59.
Back substitution gives \(y=z=h=0\), so the square matrix is injective.
For the quantitative statement use, successively, the transport residuals,
the metric defect, and \(b\) as coordinates.  This is a bounded triangular
change on the normalized cone because \(\rho,\chi,{\mathfrak E}\) are
bounded and \(|\alpha_+|\) has the V64 positive floor.  The inverse diagonal
factors in \eqref{eq:V65-local-diagonal-factors} are uniformly bounded by
that floor and the V59 floor.
\end{proof}

\subsection{V39 pullback and an algebraic BRST quartet}
\label{subsec:V65-BRST-pullback}

Use the V39 invariant representative
\begin{equation}
 \widehat h=Q_4(h,X):=h+{\mathcal S}_pX,
 \qquad
 G_p\epsilon=({\mathcal S}_p\epsilon,-\epsilon).
 \label{eq:V65-Q4}
\end{equation}
Introduce a ghost \(c\), antighost \(\bar c\), and
Nakanishi--Lautrup field \(B\), all with four components, and define
\begin{equation}
 sh={\mathcal S}_pc,
 \qquad sX=-c,
 \qquad sc=0,
 \qquad s\bar c=B,
 \qquad sB=0.
 \label{eq:V65-BRST-differential}
\end{equation}
Let \(b,y,z\) and the physical multiplier fields be BRST invariant.  Then
\begin{equation}
 s^2=0,
 \qquad s\widehat h=0,
 \qquad
 s(\widehat h-{\mathfrak E}(b,y,z))=0.
 \label{eq:V65-BRST-invariance}
\end{equation}

The edge-coordinate gauge fermion and its exact term are
\begin{equation}
 \Psi_X:=\bar c^TX,
 \qquad
 s\Psi_X=B^TX+\bar c^Tc.
 \label{eq:V65-X-gauge-fermion}
\end{equation}
The signs use the odd Leibniz rule and \(sX=-c\).

\begin{proposition}[Contractible frozen gauge quartet]
\label{prop:V65-contractible-quartet}
The pairs \((X,-c)\) and \((\bar c,B)\) are BRST doublets.  The polynomial
BRST cohomology of the quartet is the constants, and the Euler rows of
\eqref{eq:V65-X-gauge-fermion} contain the algebraic conditions
\begin{equation}
 X=0,
 \qquad c=0,
 \qquad \bar c=0,
 \label{eq:V65-quartet-Euler-rows}
\end{equation}
with \(B\) supplying the \(X\)-Euler multiplier.  No normal ghost derivative
or second ghost boundary trace is introduced.
\end{proposition}

\begin{proof}
Define a contracting homotopy on the generators by
\(\iota c=-X\), \(\iota B=\bar c\), and
\(\iota X=\iota\bar c=0\), extended as a graded derivation.  Then
\(s\iota+\iota s\) is the quartet polynomial-degree operator.  Every closed
polynomial of positive quartet degree is therefore exact.  Variation of
\(B,\bar c,c\) in \eqref{eq:V65-X-gauge-fermion} gives respectively the
three displayed rows; variation of \(X\) determines \(B\) together with the
invariant physical action.
\end{proof}

A local frozen multiplier density which generates the forward graph and
transport rows is
\begin{align}
 F_{65}^{\rm loc}
 :=&\ \Lambda^\vee
 \bigl(Q_4(h,X)-{\mathfrak E}(b,y,z)\bigr)
 \nonumber\\
 &+\alpha^T\bigl((D_0+D_n)y-D_1b\bigr)
 +\beta\bigl((D_0+D_n)z-D_1y_+\bigr)
 +s\Psi_X,
 \label{eq:V65-local-multiplier-density}
\end{align}
with multipliers \(\Lambda,\alpha,\beta\) inert under \(s\).  Equations
\eqref{eq:V65-BRST-invariance} and \(s^2=0\) prove
\begin{equation}
 sF_{65}^{\rm loc}=0
 \label{eq:V65-local-BRST-invariance}
\end{equation}
off shell.  Variation of the multipliers gives the full metric graph and
the two characteristic transports; variation of \(B\) gives \(X=0\).
Variation of \(h,b,y,z\) supplies the formal-adjoint multiplier rows.  After
eliminating those multipliers, the physical conormal row is
\eqref{eq:V65-frozen-conormal-relation}.

The corresponding Schur-reduced nineteen-row stable matrix is
\begin{equation}
 \widetilde{\mathcal B}_{65}^{\rm loc,Q}
 \begin{pmatrix}h\\X\\b\\y\\z\end{pmatrix}
 :=\begin{pmatrix}
 Q_4(h,X)-{\mathfrak E}(b,y,z)\\
 X\\
 \alpha_+y-i\rho b\\
 \alpha_+z-i\rho y_+\\
 (\lambda+m\kappa_v)b
 \end{pmatrix}.
 \label{eq:V65-local-Q-matrix}
\end{equation}

\begin{theorem}[BRST-split full-embedding floor]
\label{thm:V65-local-Q-floor}
Before the four \(X\)-rows are imposed, the stable kernel of the other
fifteen rows is exactly
\begin{equation}
 \{(G_p\epsilon,0,0,0):\epsilon\in\mathbb C^4\}.
 \label{eq:V65-Q-orbit-kernel}
\end{equation}
After the \(X\)-rows are included,
\begin{equation}
 \operatorname{rank}\widetilde{\mathcal B}_{65}^{\rm loc,Q}=19
 \label{eq:V65-Q-rank}
\end{equation}
at every nonzero closed stable frequency, with a positive normalized-cone
singular-value floor.  Its physical Schur complement is still exactly
\((\lambda+m\kappa_v)I_2\).
\end{theorem}

\begin{proof}
The rows without \(X=0\) depend on \((h,X)\) only through \(Q_4\).  V39
exactness and Theorem~\ref{thm:V65-local-floor} therefore identify their
kernel with \(\ker Q_4=\operatorname{ran}G_p\), with all physical and
transport amplitudes zero.  The row \(X=0\) is injective on this orbit.
The triangular coordinate map
\[
 (h,X,b,y,z)\longmapsto(Q_4(h,X),X,b,y,z)
\]
and its inverse are uniformly bounded on the normalized cone by the V39
estimate.  In these coordinates the matrix is the direct sum of
\({\mathcal B}_{65}^{\rm loc}\) and \(I_4\), proving the rank, floor, and
Schur assertions.
\end{proof}
\subsection{Why the adjoint characteristic sector cannot be discarded}
\label{subsec:V65-causal-adjoint}

The local multiplier density is an off-shell BRST-invariant quadratic
system, but it contains both forward constraints and the formal-adjoint
multiplier equations.  Eliminating all characteristic and multiplier fields
would return the retarded symbol \(\chi\) in one block and its advanced
formal adjoint in the transposed block.  The two blocks cannot be identified
as one causal self-adjoint operator.

\begin{theorem}[One-sided causal self-adjoint elimination no-go]
\label{thm:V65-causal-selfadjoint-no-go}
Fix \(\rho>0\).  Let \(K\) be a time-translation-invariant convolution
operator with distribution kernel \(k(t)\).  If
\begin{equation}
 \operatorname{supp}k\subset[0,\infty)
 \label{eq:V65-causal-support}
\end{equation}
and \(K\) is formally self-adjoint on a single real field, then
\(\operatorname{supp}k\subset\{0\}\), and its Laplace symbol is a polynomial
in \(s\).  Consequently \(K\) cannot have the V64 symbol
\begin{equation}
 \chi(s,\rho)=\frac{i\rho}{s+\sqrt{s^2+\rho^2}}
 \label{eq:V65-chi-nonlocal}
\end{equation}
for nonzero \(\rho\).
\end{theorem}

\begin{proof}
The formal-adjoint convolution kernel is
\(k^\dagger(t)=k(-t)^*\).  Self-adjointness and
\eqref{eq:V65-causal-support} therefore put the support of \(k\) in both
closed half-lines, hence at \(t=0\).  A distribution supported at one point
is a finite sum of derivatives of the Dirac mass.  Its Laplace transform is
a polynomial in \(s\).  For fixed \(\rho>0\), the function in
\eqref{eq:V65-chi-nonlocal} has square-root branch points at
\(s=\pm i\rho\) and is not a polynomial.  This is a contradiction.
\end{proof}

\begin{corollary}[Necessary doubled or local characteristic realization]
\label{cor:V65-doubled-characteristic}
A real variational realization of the retarded V65 graph must retain either
\begin{enumerate}
\item the forward characteristic fields together with their formal-adjoint
      multiplier fields and compatible corner data; or
\item an equivalent doubled retarded/advanced formulation.
\end{enumerate}
It cannot replace both sectors by one causal self-adjoint half-angle
operator.  The local cascade
\eqref{eq:V65-characteristic-cascade} realizes the first alternative at the
principal frozen level.
\end{corollary}

\begin{firewall}[What the local action does not yet supply]
\label{fire:V65-local-action-scope}
Equation \eqref{eq:V65-local-multiplier-density} is a local quadratic
multiplier density with an exact off-shell BRST differential.  V65 has not
yet selected the incoming/outgoing corner traces for
\((\alpha,\beta)\), proved maximal dissipativity of the combined
forward--adjoint characteristic generator, or matched its real-time Green
form to the positive V59 quotient energy.  Nor has it proved that the full
time-domain boundary relation propagates the nonlinear harmonic constraints.
The stable floors in Theorems~\ref{thm:V65-complete-floor} and
\ref{thm:V65-local-Q-floor} do not imply those statements.
\end{firewall}

\subsection{Finite algebraic certificates}
\label{subsec:V65-certificates}

The action-defect card at normal incidence is
\begin{equation}
 {\mathcal T}_0J_q\gamma_U=0,
 \qquad
 {\mathcal T}_0^\vee W=0,
 \qquad
 2\kappa\Omega_{36}\(U,V\)=-2.
 \label{eq:V65-action-defect-card}
\end{equation}
The full-graph cards are
\begin{align}
 {\mathcal A}_q{\mathcal E}_\chi&=0,
 &
 {\mathcal E}_\chi^\vee{\mathcal E}_\chi&=I_2,
 \label{eq:V65-full-graph-cards}\\
 {\mathcal E}_\chi^\vee
 \mathfrak p_{36}^{\rm eq}{\mathcal E}_\chi&=\lambda I_2,
 &
 \operatorname{rank}{\mathcal B}_{65}&=12.
 \label{eq:V65-conormal-rank-card}
\end{align}
The local cards are
\begin{align}
 (s+\lambda)y-i\rho b&=0,
 &
 (s+\lambda)z-i\rho y_+&=0,
 \label{eq:V65-transport-card}\\
 {\mathfrak E}(b,y,z)-{\mathcal E}_\chi b&=0,
 &
 \operatorname{rank}{\mathcal B}_{65}^{\rm loc}&=15,
 \label{eq:V65-local-rank-card}\\
 \operatorname{rank}\widetilde{\mathcal B}_{65}^{\rm loc,Q}&=19,
 &
 sF_{65}^{\rm loc}&=0.
 \label{eq:V65-BRST-rank-card}
\end{align}
The last identity follows from \(sF_{65}^{\rm loc}=0\), not from an
on-shell ghost equation.

A finite validator should test the open rational-root point, normal
incidence, static face, both gravitational glancing orientations, continuum
threshold, and an evanescent boundary point.  At each sample it must verify
the component identity \eqref{eq:V65-local-equals-E}, the conormal card,
and ranks \(12,15,19\).  It must also verify the exact normal-incidence
Green witness separately; an invertible stable matrix cannot detect that
action defect.

\subsection{V66 acquisition order}
\label{subsec:V65-next}

V66 should retain the full configuration graph and proceed in this order:
\begin{enumerate}
\item construct the normal Green form for the forward fields
      \((y,z)\) and adjoint fields \((\alpha,\beta)\) on a finite collar;
\item choose characteristic incoming, outgoing, and corner traces which make
      that first-order operator closed and prove its maximal dissipative or
      skew-adjoint realization;
\item combine the resulting edge Green cancellation with the V36
      metric and V59 continuum forms, and prove the positive physical
      quotient estimate after removing the contractible quartet;
\item derive the boundary evolution of the harmonic constraint from the
      linearized Bianchi identity and decide whether the full graph gives a
      homogeneous maximally dissipative constraint domain;
\item if the constraint row fails, go upstream to a generalized-harmonic
      source or an additional characteristic constraint edge, without
      returning to the non-isotropic V64 row set;
\item only after these linear generator gates close, covariantize the graph,
      derive the nonlinear Brown--York stress identity, and restore Higgs,
      fermion, neutrino, and Yang--Mills sources.
\end{enumerate}
The next compulsory companion audit is V70.

\subsection{V65 result ledger and claim boundary}
\label{subsec:V65-ledger}

Relative to V1--V64, this note proves:
\begin{enumerate}
\item the required V45 Yang--Mills and V31 Navier--Stokes companion audit,
      with no cross-type theorem import;
\item the exact canonical decomposition of the V36 de Donder boundary term
      into the four gauge pairs \((y,C)\);
\item a normal-incidence witness proving that the twelve V64 axial plus
      harmonic rows are not an action-isotropic relation;
\item the ten-row full TT configuration graph and its twelve-dimensional
      frozen conormal Lagrangian;
\item a complete action-compatible twelve-by-twelve stable matrix with a
      positive closed-cone floor and the exact V59 Schur factor;
\item an explicit two-stage outgoing characteristic realization of every
      \(\chi\) and \(\chi^2\) entry, with a uniform fifteen-row floor;
\item a V39-invariant local multiplier density, an off-shell nilpotent
      algebraic BRST quartet, and a uniformly split nineteen-row stable
      representative; and
\item the support-theoretic obstruction to eliminating the retarded and
      formal-adjoint characteristic sectors into one causal self-adjoint
      half-angle operator.
\end{enumerate}

The frozen conormal action and the local multiplier density are rigorous
linear principal constructions.  A closed positive forward--adjoint collar
generator, its corner domain, harmonic-constraint propagation, a nonlinear
BRST/BV--BFV master action, nonlinear Brown--York/Bianchi stress closure,
curved boundary atlas, and coupled Standard-Model source sectors are not
proved.  The full Standard Model--Einstein continuum is not proved.

\subsection{Append-only record}

This V65 material was appended to the exact V64 parent whose SHA--256 digest
is
\[
 \texttt{2ed0dcd226c75cfe7ad1008ca692097a4b8265dac557247036c4e2434c59326e}.
\]
The next research note is V66 and begins with the forward--adjoint
characteristic Green form and corner-domain calculation.

% END APPENDED VERSION V65 MATERIAL

% BEGIN APPENDED VERSION V66 MATERIAL

\clearpage
\section{V66: harmonic-residual identification, spin-completed causal lift,
and forward--adjoint corner domains}
\label{sec:V66-spin-completion}

\subsection{Purpose and correction of the V65 acquisition}
\label{subsec:V66-purpose}

V65 produced a local algebraic chart for the full TT representative and
realized its factors \(\chi\) and \(\chi^2\) by the outgoing cascade
\[
 (D_0+D_n)y=D_1b,
 \qquad
 (D_0+D_n)z=D_1y_+.
\]
That construction was proved only on the stable shell.  Its multiplier
Euler equations contain a forward transport and a formal adjoint transport,
but V65 deliberately left their corner domains and their status as a time
generator open.  The first calculation below shows that the naive forward
corner for both sectors cannot give a positive forward semigroup.

The repair is upstream.  The three outgoing rows are first identified with
the harmonic constraint of the V65 metric chart.  The missing reverse rows
are then supplied by the bulk wave equation.  In the plus channel the result
is the spin-one symmetric first-order system; in the cross channel it is the
spin-one-half system.  These systems have a positive passive realization on
a half-collar, their stable Weyl function is exactly
\(\operatorname{diag}(\chi^2,\chi)\), and their adjoints have a precise
terminal-value interpretation.  The construction also identifies the V59
Dirichlet subsidiary field with the actual linearized harmonic constraint
on the full metric graph.

No V59 quotient-generator theorem is reproved here.  V66 uses its positive
form and its scalar Schur factor as established inputs.  The next compulsory
Yang--Mills/Navier--Stokes audit remains V70.

\subsection{The outgoing cascade is exactly the harmonic residual}
\label{subsec:V66-harmonic-card}

Retain the rotated tangential chart \(k=(\rho,0)\), the outward derivative
\(D_n=-\partial_x\), and the V65 variables
\[
 w=(b_+,y_+,z,b_\times,y_\times).
\]
It is convenient to rewrite the V65 algebraic metric chart as
\begin{equation}
 \mathfrak E(w)=\frac1{\sqrt2}
 \begin{pmatrix}
 -z&-y_+&-y_\times&z\\
 -y_+&-b_+&-b_\times&y_+\\
 -y_\times&-b_\times&b_+&y_\times\\
 z&y_+&y_\times&-z
 \end{pmatrix},
 \label{eq:V66-metric-chart}
\end{equation}
in the order \((0,1,2,n)\).  Define
\begin{align}
 R_+&:=(D_0+D_n)y_+-D_1b_+,
 &R_\times&:=(D_0+D_n)y_\times-D_1b_\times,
 \label{eq:V66-R-y}\\
 R_z&:=(D_0+D_n)z-D_1y_+.
 \label{eq:V66-R-z}
\end{align}

\begin{theorem}[Exact harmonic-residual card]
\label{thm:V66-harmonic-card}
For arbitrary off-shell fields \(w\), without using a root or a field
equation,
\begin{align}
 \operatorname{tr}_\eta\mathfrak E(w)&=0,
 &q^\mu\mathfrak E(w)_{\mu\nu}&=0,
 \label{eq:V66-algebraic-trace-axial}\\
 \mathcal C_\nu(\mathfrak E(w))
 &=\frac1{\sqrt2}(R_z,R_+,R_\times,-R_z)_\nu.
 \label{eq:V66-exact-C-card}
\end{align}
Consequently the three V65 transport rows are equivalent to
\(\mathcal C(\mathfrak E(w))=0\).  They are not an additional collection
of four independent boundary conditions.
\end{theorem}

\begin{proof}
The diagonal entries in \eqref{eq:V66-metric-chart} give
\(-h_{00}+h_{11}+h_{22}+h_{nn}=0\).  Its zeroth and normal rows add to zero,
which is the displayed null-reference axial identity.  Since the trace
vanishes, \(\mathcal C_\nu=D^\mu h_{\mu\nu}\).  With
\(D^\mu=(-D_0,D_1,0,D_n)\), direct contraction gives
\begin{align*}
 \sqrt2\,\mathcal C_0
 &=D_0z-D_1y_++D_nz=R_z,\\
 \sqrt2\,\mathcal C_1
 &=D_0y_+-D_1b_++D_ny_+=R_+,\\
 \sqrt2\,\mathcal C_2
 &=D_0y_\times-D_1b_\times+D_ny_\times=R_\times,\\
 \sqrt2\,\mathcal C_n&=-R_z.
\end{align*}
The map \((R_z,R_+,R_\times)\mapsto\mathcal C\) is injective, proving the
last statement.
\end{proof}

The Euclidean Frobenius norm also has the exact algebraic form
\begin{equation}
 \|\mathfrak E(w)\|_{
m E}^2
 =|b_+|^2+|b_\times|^2+2|y_+|^2+2|y_\times|^2+2|z|^2.
 \label{eq:V66-chart-norm}
\end{equation}
Thus the chart has no hidden indefinite algebraic direction; positivity of
a time generator is a separate differential question.

\subsection{Why the triangular multiplier does not close forward in time}
\label{subsec:V66-triangular-no-go}

The obstruction is already present in one scalar transport.  On
\(I=(0,L)\), let
\begin{equation}
 A_Lf=f',\qquad D(A_L)=\{f\in H^1(I):f(L)=0\}.
 \label{eq:V66-left-shift-generator}
\end{equation}
This is the spatial generator of
\((\partial_t-\partial_x)f=0\) with its incoming value fixed at \(x=L\).

\begin{theorem}[Forward multiplier corner obstruction]
\label{thm:V66-multiplier-corner-no-go}
The operator \(A_L\) is maximal dissipative and generates the left-shift
contraction semigroup.  Its Hilbert adjoint is
\begin{equation}
 A_L^*g=-g',\qquad D(A_L^*)=\{g\in H^1(I):g(0)=0\}.
 \label{eq:V66-shift-adjoint}
\end{equation}
The Euler equation of a multiplier for the forward transport evolves the
multiplier with
\begin{equation}
 B:=-A_L^*=\partial_x,qquad D(B)=\{g\in H^1(I):g(0)=0\}.
 \label{eq:V66-bad-forward-adjoint}
\end{equation}
The operator \(B\) is not the generator of a strongly continuous forward
semigroup on \(L^2(I)\).  In particular, the forward field and this formal
adjoint cannot be made into the desired positive joint forward collar
generator merely by imposing the two variational corner conditions.
\end{theorem}

\begin{proof}
The semigroup generated by \(A_L\) is
\[
 (T(t)f)(x)=
 \begin{cases}f(x+t),&x+t<L,\\0,&x+t\ge L.\end{cases}
\]
It is a contraction, and its generator is exactly
\eqref{eq:V66-left-shift-generator}.  Integration by parts gives
\eqref{eq:V66-shift-adjoint}.

For \(\mu>0\), the solution of \((\mu-B)u=f\), \(u(0)=0\), is
\begin{equation}
 u(x)=-e^{\mu x}\int_0^xe^{-\mu r}f(r)\,\dd r.
 \label{eq:V66-bad-resolvent}
\end{equation}
Taking \(f=1\) shows that the norm of this resolvent is bounded below by
\(c_L e^{\mu L}/\mu^{3/2}\) for all sufficiently large \(\mu\).  This violates the Hille--Yosida
resolvent bound of every strongly continuous semigroup generator.  The
same failure is visible from
\(\operatorname{Re}\langle Bg,g\rangle=|g(L)|^2/2\), whose ratio to
\(\|g\|^2\) is unbounded on boundary layers at \(L\).
\end{proof}

This theorem sharpens the V65 causality firewall.  The formal adjoint is a
terminal-value field, or the forward subsystem must be completed before a
positive state-space realization is claimed.

\subsection{The minimal spin completion}
\label{subsec:V66-spin-system}

Put
\begin{equation}
 J_3=\begin{pmatrix}1&0&0\\0&0&0\\0&0&-1\end{pmatrix},
 \qquad
 J_1=\frac1{\sqrt2}
 \begin{pmatrix}0&1&0\\1&0&1\\0&1&0\end{pmatrix},
 \label{eq:V66-spin-one-matrices}
\end{equation}
and
\begin{equation}
 \sigma_3=\begin{pmatrix}1&0\\0&-1\end{pmatrix},
 \qquad
 \sigma_1=\begin{pmatrix}0&1\\1&0\end{pmatrix}.
 \label{eq:V66-Pauli-matrices}
\end{equation}
Both pairs are real symmetric.  With
\begin{equation}
 U_+:=(b_+,\sqrt2y_+,z)^T,
 \qquad U_\times:=(b_\times,y_\times)^T,
 \label{eq:V66-spin-fields}
\end{equation}
define the spin-completed equations
\begin{align}
 (D_0I_3-D_nJ_3-D_1J_1)U_+&=0,
 \label{eq:V66-spin-one-system}\\
 (D_0I_2-D_n\sigma_3-D_1\sigma_1)U_\times&=0.
 \label{eq:V66-spin-half-system}
\end{align}
Their component form is
\begin{align}
 (D_0-D_n)b_+-D_1y_+&=0,
 \label{eq:V66-plus-reverse}\\
 2D_0y_+-D_1(b_++z)&=0,
 \label{eq:V66-plus-middle}\\
 (D_0+D_n)z-D_1y_+&=0,
 \label{eq:V66-plus-forward-z}\\
 (D_0-D_n)b_\times-D_1y_\times&=0,
 \label{eq:V66-cross-reverse}\\
 (D_0+D_n)y_\times-D_1b_\times&=0.
 \label{eq:V66-cross-forward}
\end{align}
\begin{theorem}[Constraint propagation inside the spin completion]
\label{thm:V66-spin-constraint}
Every smooth solution of \eqref{eq:V66-spin-one-system} satisfies
\begin{equation}
 D_0R_+=0.
 \label{eq:V66-Rplus-conserved}
\end{equation}
If \(R_+=0\) on one time slice, it remains zero and the spin equations
imply all three V65 outgoing rows.  On that constraint subspace,
\begin{equation}
 (D_0^2-D_n^2-D_1^2)f=0,
 \qquad f\in\{b_+,y_+,z,b_\times,y_\times\}.
 \label{eq:V66-component-wave}
\end{equation}
Thus the extra rows are the first-order wave completion of the harmonic
cascade, rather than new physical boundary data.
\end{theorem}

\begin{proof}
Equations \eqref{eq:V66-plus-reverse}--%
\eqref{eq:V66-plus-forward-z} give
\[
 D_0b_+=D_nb_++D_1y_+,
 \quad
 D_0y_+=\tfrac12D_1(b_++z),
 \quad
 D_0z=-D_nz+D_1y_+.
\]
Using commutation of the frozen derivatives,
\begin{align*}
 D_0R_+
 &=D_0^2y_++D_nD_0y_+-D_1D_0b_+\\
 &=\tfrac12D_1D_n(b_+-z)+D_1^2y_+
   +\tfrac12D_1D_n(b_++z)
   -D_1D_nb_+-D_1^2y_+=0.
\end{align*}
When \(R_+=0\), subtracting it from
\eqref{eq:V66-plus-middle} gives
\((D_0-D_n)y_+=D_1z\).  Hence \((b_+,y_+)\) and
\((y_+,z)\) each obey the two opposite first-order factors; composing them
gives \eqref{eq:V66-component-wave}.  The same factorization follows
directly from the two cross equations.
\end{proof}

\begin{corollary}[Invariant harmonic spin subspace]
\label{cor:V66-invariant-harmonic-spin}
The closed initial constraint \(R_+=0\), together with
\eqref{eq:V66-spin-one-system}--\eqref{eq:V66-spin-half-system}, makes
\(h=\mathfrak E(w)\) trace-free, null-reference axial, harmonic, and
componentwise wave.  No independent fourth harmonic datum is present because
\(\mathcal C_n=-\mathcal C_0\) on the algebraic chart.
\end{corollary}

\subsection{Exact retarded kernel and preservation of the V65 floor}
\label{subsec:V66-stable-kernel}

On a decaying Laplace--Fourier mode, \(D_0=s\), \(D_1=i\rho\), and
\(D_n=\lambda\).  In the unscaled coordinates
\((b_+,y_+,z;b_\times,y_\times)\), the two spin matrices are
\begin{align}
 M_+(s,\rho)&=
 \begin{pmatrix}
 s-\lambda&-i\rho&0\\
 -i\rho&2s&-i\rho\\
 0&-i\rho&s+\lambda
 \end{pmatrix},
 \label{eq:V66-M-plus}\\
 M_\times(s,\rho)&=
 \begin{pmatrix}s-\lambda&-i\rho\\-i\rho&s+\lambda\end{pmatrix}.
 \label{eq:V66-M-cross}
\end{align}

\begin{theorem}[Exact spin Weyl kernel]
\label{thm:V66-spin-Weyl-kernel}
At every nonzero open or limiting closed stable frequency,
\begin{align}
 \ker M_+&=\operatorname{span}\bigl((1,\chi,\chi^2)^T\bigr),
 &\operatorname{rank}M_+&=2,
 \label{eq:V66-plus-kernel}\\
 \ker M_\times&=\operatorname{span}\bigl((1,\chi)^T\bigr),
 &\operatorname{rank}M_\times&=1.
 \label{eq:V66-cross-kernel}
\end{align}
Consequently the five spin rows have rank three and exactly the same kernel
as the three triangular V65 transport rows.  In particular the replacement
is valid at normal incidence, both gravitational glancing sheets, the static
face, and the continuum threshold.
\end{theorem}

\begin{proof}
The identities
\begin{equation}
 (s+\lambda)\chi=i\rho,
 \qquad
 i\rho\chi=s-\lambda,
 \qquad
 2s\chi=i\rho(1+\chi^2)
 \label{eq:V66-half-angle-identities}
\end{equation}
follow from \(\lambda^2=s^2+\rho^2\).  They show directly that the displayed
vectors lie in the two kernels.  The determinant of \(M_\times\) vanishes,
while its entry \(s+\lambda\) never vanishes, so its rank is one.  For the
plus block, the lower-right two-by-two minor equals
\[
 2s(s+\lambda)+\rho^2=(s+\lambda)^2\ne0.
\]
Its determinant vanishes on the root shell, hence its rank is exactly two.
The V65 triangular matrix has the same two kernel vectors and rank three
because each outgoing row contains the nonzero factor \(s+\lambda\).
Equality of the kernels follows.
\end{proof}

The normalized variables expose the symmetric-square origin of the plus
block.  If
\[
 v=(1,\chi)^T,
\]
then the plus stable vector is the symmetric square
\((1,\sqrt2\chi,\chi^2)^T\).  The matrices \(J_3,J_1\) are the corresponding
spin-one representatives of \(\sigma_3,\sigma_1\).  This explains why the
previous triangular chain acquired a zero-speed constraint component when
it was completed.

Define the rectangular, Schur-reduced stable certificate on
\((h,w)\in\mathbb C^{10}\oplus\mathbb C^5\) by
\begin{equation}
 \widehat{\mathcal B}_{66}^{\rm spin}
 \binom hw
 :=\begin{pmatrix}
 h-\mathfrak E(w)\\
 M_+(b_+,y_+,z)^T\\
 M_\times(b_\times,y_\times)^T\\
 (\lambda+m\kappa_v)(b_+,b_\times)^T
 \end{pmatrix}.
 \label{eq:V66-spin-stable-certificate}
\end{equation}
It has seventeen rows and fifteen columns.

\begin{theorem}[Spin completion preserves the complete stable floor]
\label{thm:V66-spin-stable-floor}
For every \(m>0\), \(0<v<1\), and every nonzero open or limiting closed
stable frequency,
\begin{equation}
 \ker\widehat{\mathcal B}_{66}^{\rm spin}=0,
 \qquad
 \operatorname{rank}\widehat{\mathcal B}_{66}^{\rm spin}=15.
 \label{eq:V66-spin-full-rank}
\end{equation}
On the normalized closed cone its least singular value has a positive
minimum.  The physical Schur factor remains exactly
\((\lambda+m\kappa_v)I_2\).
\end{theorem}

\begin{proof}
A vector in the kernel first obeys the spin equations.  By
Theorem~\ref{thm:V66-spin-Weyl-kernel}, it has
\[
 y=\chi b,
 \qquad z=\chi^2b_+.
\]
The graph row then gives \(h=\mathcal E_\chi b\), and the last two rows give
\((\lambda+m\kappa_v)b=0\).  V59 proves that this factor is nonzero on the
complete nonzero closed stable cone, so \(b=0\) and then \(w=h=0\).  This
proves full column rank.

All matrix entries are continuous on the normalized retarded-root cone.
That cone is compact, and the least singular value never vanishes there;
therefore it has a positive minimum.  The elimination just performed shows
that no factor other than the V59 scalar occurs in the physical Schur
complement.
\end{proof}
\subsection{Passive half-collar generator and the adjoint terminal domain}
\label{subsec:V66-passive-generator}

The stable calculation becomes causal only after a state-space domain is
specified.  Work first in one rotated tangential Plancherel chart; the
second tangential coordinate is a passive direct-integral parameter.  Set
\begin{equation}
 \mathbf U=(b_+,\sqrt2y_+,z;b_\times,y_\times)^T,
 \qquad
 \mathscr J_n=J_3\oplus\sigma_3,
 \qquad
 \mathscr J_1=J_1\oplus\sigma_1,
 \label{eq:V66-combined-spin-field}
\end{equation}
and on \(x>0\) define
\begin{equation}
 \mathscr A_e:=\mathscr J_nD_n+\mathscr J_1D_1
 =-\mathscr J_n\partial_x+\mathscr J_1D_1.
 \label{eq:V66-edge-generator-expression}
\end{equation}
The normal incoming and outgoing traces are
\begin{equation}
 \Gamma_{\rm in}\mathbf U=(b_+,b_\times)^T,
 \qquad
 \Gamma_{\rm out}\mathbf U=(z,y_\times)^T.
 \label{eq:V66-edge-traces}
\end{equation}
The middle spin-one component \(y_+\) has zero normal characteristic speed
and requires no independent normal trace.

Let \(\mathcal H_e=L^2(\mathbb R^2_y\times\mathbb R_+^x;\mathbb C^5)\).
The graph domains below mean the closures of smooth compactly supported
fields for which the indicated characteristic trace vanishes.  Equivalently,
they are the standard Friedrichs graph domains after tangential Fourier
transformation:
\begin{align}
 D(\mathscr A_{\rm in})
 &:=\{\mathbf U:\mathbf U,\mathscr A_e\mathbf U\in\mathcal H_e,
                 \ \Gamma_{\rm in}\mathbf U=0\},
 \label{eq:V66-Ain-domain}\\
 D(\mathscr A_{\rm out})
 &:=\{\mathbf U:\mathbf U,\mathscr A_e\mathbf U\in\mathcal H_e,
                 \ \Gamma_{\rm out}\mathbf U=0\}.
 \label{eq:V66-Aout-domain}
\end{align}
Only the nonzero-speed traces occur in these definitions.

\begin{lemma}[Exact edge Green form]
\label{lem:V66-edge-Green}
For smooth graph-domain fields,
\begin{align}
 &\langle\mathscr A_e\mathbf U,\mathbf V\rangle_{
 \mathcal H_e}
 +\langle\mathbf U,\mathscr A_e\mathbf V\rangle_{
 \mathcal H_e}
 \nonumber\\
 &\qquad=
 \langle\Gamma_{\rm in}\mathbf U,
        \Gamma_{\rm in}\mathbf V\rangle_{\partial}
 -\langle\Gamma_{\rm out}\mathbf U,
        \Gamma_{\rm out}\mathbf V\rangle_{\partial}.
 \label{eq:V66-edge-Green-form}
\end{align}
In particular,
\begin{equation}
 \operatorname{Re}\langle\mathscr A_e\mathbf U,\mathbf U\rangle
 =\frac12\left(
 \|\Gamma_{\rm in}\mathbf U\|_\partial^2
 -\|\Gamma_{\rm out}\mathbf U\|_\partial^2\right).
 \label{eq:V66-edge-flux}
\end{equation}
\end{lemma}

\begin{proof}
The tangential part cancels because \(D_1\) is skew-adjoint and
\(\mathscr J_1\) is Hermitian.  Integrating
\(-\mathscr J_n\partial_x\) on the half-line leaves
\(\mathbf U(0)^*\mathscr J_n\mathbf V(0)\).  Since
\(\mathscr J_n=\operatorname{diag}(1,0,-1,1,-1)\), this is exactly the
right side of \eqref{eq:V66-edge-Green-form}.  Density extends the identity
to the graph traces.
\end{proof}

\begin{theorem}[Maximal dissipative incoming realization]
\label{thm:V66-maximal-dissipative}
The operator \(\mathscr A_{\rm in}\) is maximal dissipative on
\(\mathcal H_e\),
\begin{equation}
 \mathscr A_{\rm in}^*=-\mathscr A_{\rm out},
 \label{eq:V66-edge-adjoint}
\end{equation}
and both \(\mathscr A_{\rm in}\) and its adjoint generate contraction
semigroups.  Thus the homogeneous forward edge problem and the dual
terminal problem have complete Hilbert-space domains.
\end{theorem}

\begin{proof}
Equation \eqref{eq:V66-edge-flux} is nonpositive on
\(D(\mathscr A_{\rm in})\), so \(\mathscr A_{\rm in}\) is dissipative.
The Green form and the independent incoming/outgoing trace ranges give
\eqref{eq:V66-edge-adjoint}.  This is the usual annihilator calculation:
the incoming trace of a field in \(D(\mathscr A_{\rm in})\) vanishes, while
its outgoing trace is arbitrary, forcing the outgoing trace in the adjoint
domain to vanish.

For \(\zeta>0\),
\[
 \|(\zeta-\mathscr A_{\rm in})\mathbf U\|\,\|\mathbf U\|
 \ge
 \operatorname{Re}\langle
 (\zeta-\mathscr A_{\rm in})\mathbf U,\mathbf U\rangle
 \ge\zeta\|\mathbf U\|^2.
\]
Hence the range is closed.  Its orthogonal complement is
\(\ker(\zeta+\mathscr A_{\rm out})\).  On the outgoing domain,
\[
 \operatorname{Re}\langle
 (\zeta+\mathscr A_{\rm out})\mathbf V,\mathbf V\rangle
 =\zeta\|\mathbf V\|^2
  +\tfrac12\|\Gamma_{\rm in}\mathbf V\|_\partial^2,
\]
so this kernel is zero.  The range is therefore dense and closed, hence all
of \(\mathcal H_e\).  Lumer--Phillips proves maximal dissipativity and the
semigroup assertions.  The same argument applies to the adjoint.
\end{proof}

For a smooth boundary-controlled solution
\begin{equation}
 \partial_t\mathbf U=\mathscr A_e\mathbf U,
 \qquad
 \Gamma_{\rm in}\mathbf U=g,
 \label{eq:V66-controlled-edge}
\end{equation}
Lemma~\ref{lem:V66-edge-Green} gives the exact storage identity
\begin{equation}
 \frac{\dd}{\dd t}\frac12\|\mathbf U(t)\|_{\mathcal H_e}^2
 =\frac12\left(\|g(t)\|_\partial^2
 -\|\Gamma_{\rm out}\mathbf U(t)\|_\partial^2\right).
 \label{eq:V66-passive-storage}
\end{equation}
For zero initial state this extends by density to
\begin{equation}
 \|\mathbf U(T)\|_{\mathcal H_e}^2
 +\int_0^T\|\Gamma_{\rm out}\mathbf U\|_\partial^2\,\dd t
 =\int_0^T\|g\|_\partial^2\,\dd t.
 \label{eq:V66-passive-estimate}
\end{equation}
The edge lift is therefore a passive causal system, including its zero-speed
constraint component.

The correct variational meaning of the formal adjoint can now be stated
without pretending that it is a second forward copy.  Put
\begin{equation}
 P_e:=\partial_t-\mathscr A_e,
 \qquad
 P_e^\dagger:=-\partial_t+\mathscr A_e.
 \label{eq:V66-P-and-adjoint}
\end{equation}

\begin{proposition}[Forward--adjoint slab and corner domain]
\label{prop:V66-forward-adjoint-corners}
On \([0,T]\times\mathbb R_+\), smooth fields obey
\begin{align}
 &\int_0^T\langle P_e\mathbf U,\mathbf V\rangle\,\dd t
 -\int_0^T\langle\mathbf U,P_e^\dagger\mathbf V\rangle\,\dd t
 \nonumber\\
 &\quad=
 [\langle\mathbf U,\mathbf V\rangle]_{0}^{T}
 -\int_0^T\left(
 \langle\Gamma_{\rm in}\mathbf U,\Gamma_{\rm in}\mathbf V\rangle
 -\langle\Gamma_{\rm out}\mathbf U,\Gamma_{\rm out}\mathbf V\rangle
 \right)\dd t.
 \label{eq:V66-spacetime-Green}
\end{align}
A forward variation with
\begin{equation}
 \delta\mathbf U(0)=0,
 \qquad \Gamma_{\rm in}\delta\mathbf U=0
 \label{eq:V66-forward-corner}
\end{equation}
and an adjoint field with
\begin{equation}
 \mathbf V(T)=0,
 \qquad \Gamma_{\rm out}\mathbf V=0
 \label{eq:V66-adjoint-corner}
\end{equation}
annihilate every term on the right.  The forward equation is governed by
\(\mathscr A_{\rm in}\); after the reversal \(\tau=T-t\), the adjoint is
governed by
\(\mathscr A_{\rm in}^*=-\mathscr A_{\rm out}\).  Hence both halves are
well posed in their correct time directions.
\end{proposition}

\begin{proof}
Integrate the time derivative and use
Lemma~\ref{lem:V66-edge-Green}.  Conditions
\eqref{eq:V66-forward-corner}--\eqref{eq:V66-adjoint-corner} remove the
temporal and characteristic boundary pairings separately.  The semigroup
statement is Theorem~\ref{thm:V66-maximal-dissipative} after time reversal.
For nonhomogeneous smooth data the necessary zeroth-order corner
compatibility is
\(\Gamma_{\rm in}\mathbf U(0)=g(0)\); higher regularity requires its
successive equation derivatives.  No adjoint initial datum is prescribed.
\end{proof}
\subsection{The retarded Weyl function is the V65 full TT frame}
\label{subsec:V66-Weyl-function}

For \(\operatorname{Re}s>0\), an \(L^2\) decaying half-collar mode is
proportional to \(e^{-\lambda x}\), so \(D_n=\lambda\).  The incoming trace
is the two-polarization amplitude \(b=(b_+,b_\times)\).

\begin{theorem}[Passive causal realization of the half-angle frame]
\label{thm:V66-passive-Weyl}
The unique decaying spin solution with incoming value \(b\) satisfies
\begin{equation}
 y_+=\chi b_+,
 \qquad z=\chi^2b_+,
 \qquad y_\times=\chi b_\times.
 \label{eq:V66-Weyl-state}
\end{equation}
Its outgoing Weyl function is
\begin{equation}
 \Gamma_{\rm out}\mathbf U
 =\mathscr S_e(s,\rho)\Gamma_{\rm in}\mathbf U,
 \qquad
 \mathscr S_e(s,\rho)=
 \begin{pmatrix}\chi^2&0\\0&\chi\end{pmatrix}.
 \label{eq:V66-edge-Weyl-function}
\end{equation}
It is holomorphic and strictly contractive for
\(\operatorname{Re}s>0\), and its limiting closed-cone norm is at most one.
Moreover
\begin{equation}
 \mathfrak E(w)=\mathcal E_\chi b,
 \label{eq:V66-Weyl-equals-TT}
\end{equation}
so the causal state-space output is exactly the V64--V65 full TT
representative, not an approximation to it.
\end{theorem}

\begin{proof}
Theorem~\ref{thm:V66-spin-Weyl-kernel} gives
\eqref{eq:V66-Weyl-state} and uniqueness of each stable spin vector.
Substitution gives \eqref{eq:V66-edge-Weyl-function}.  V64 proves
\(|\chi|<1\) in the open half-plane, \(|\chi|\le1\) on the limiting cone,
and holomorphy there, so the same statements hold for
\(\mathscr S_e\).  Finally, substituting
\eqref{eq:V66-Weyl-state} in \eqref{eq:V66-metric-chart} reproduces the
component expansion of \(\mathcal E_+b_++\mathcal E_\times b_\times\)
proved in V65.
\end{proof}

The state and metric norms obey the exact stable identities
\begin{align}
 \|\mathbf U\|^2
 &= (1+2|\chi|^2+|\chi|^4)|b_+|^2
   +(1+|\chi|^2)|b_\times|^2,
 \label{eq:V66-state-norm}\\
 \|\mathcal E_\chi b\|_{
m E}^2
 &= (1+2|\chi|^2+2|\chi|^4)|b_+|^2
   +(1+2|\chi|^2)|b_\times|^2.
 \label{eq:V66-metric-norm}
\end{align}
Consequently
\begin{equation}
 \|b\|^2\le\|\mathcal E_\chi b\|_{
m E}^2\le5\|b\|^2,
 \qquad
 \|\mathbf U\|^2\le4\|b\|^2.
 \label{eq:V66-uniform-lift-bounds}
\end{equation}
The positive storage estimate \eqref{eq:V66-passive-estimate} is the
time-domain counterpart of these stable multiplier bounds.

\subsection{Identification and propagation of the actual harmonic constraint}
\label{subsec:V66-constraint-propagation}

Let \(\widehat h=Q_4(h,X)\) be the V39 invariant metric.  Impose the local
configuration graph
\begin{equation}
 \widehat h=\mathfrak E(w)
 \label{eq:V66-Q4-spin-graph}
\end{equation}
and the spin system, with \(R_+=0\) on the initial corner.  By
Theorems~\ref{thm:V66-harmonic-card} and
\ref{thm:V66-spin-constraint}, the boundary trace of the actual harmonic
covector is then homogeneous.

\begin{theorem}[Linear harmonic constraint closes on the full graph]
\label{thm:V66-linear-constraint-closure}
Assume the planar linearized generalized-harmonic Einstein equation for
\(\widehat h\), with a source whose linearized stress tensor is conserved.
Assume the spin-completed boundary graph
\eqref{eq:V66-Q4-spin-graph}, the initial spin constraint \(R_+=0\), and
Einstein initial data satisfying
\begin{equation}
 \mathcal C(\widehat h)|_{t=0}=0,
 \qquad
 \partial_t\mathcal C(\widehat h)|_{t=0}=0.
 \label{eq:V66-constraint-corner-data}
\end{equation}
Then
\begin{equation}
 \mathcal C(\widehat h)=0
 \label{eq:V66-constraint-zero}
\end{equation}
throughout the linear evolution.  The subsidiary Dirichlet block of V59 is
therefore the actual metric harmonic constraint on this full-graph domain.
\end{theorem}

\begin{proof}
The spin equations propagate \(R_+=0\); the cross and \(z\) forward rows
are already among those equations.  The exact card
\eqref{eq:V66-exact-C-card} therefore gives
\begin{equation}
 \mathcal C(\widehat h)|_{x=0}=0.
 \label{eq:V66-constraint-boundary}
\end{equation}
Taking the divergence of the reduced linear Einstein equation and using the
linearized Bianchi identity together with source conservation gives
\begin{equation}
 \Box\mathcal C_\nu(\widehat h)=0
 \label{eq:V66-constraint-wave}
\end{equation}
componentwise.  Multiply by
\(\partial_t\overline{\mathcal C}_\nu\), integrate over the half-space, and
use \eqref{eq:V66-constraint-boundary}.  The standard Dirichlet wave energy
is conserved.  Its initial value is zero by
\eqref{eq:V66-constraint-corner-data}; hence both derivatives of
\(\mathcal C\) vanish.  The homogeneous energy-space convention removes the
constant zero mode, so \(\mathcal C=0\).  This is precisely the V59
Dirichlet subsidiary evolution, now with the subsidiary symbol
\(c_{\mathcal C}=\mathcal C(\widehat h)\).
\end{proof}

\begin{corollary}[Positive physical quotient after the algebraic quartet]
\label{cor:V66-positive-physical-quotient}
On the zero-constraint spin domain, the pullback of the V59 form
\(\mathfrak q_v\oplus\mathfrak q_{\mathcal C}\) has
\(\mathfrak q_{\mathcal C}=0\).  The V65 ghost quartet
\((X,c_{\rm gh},\bar c_{\rm gh},B)\) is contractible; here \(c_{\rm gh}\)
is distinct from the subsidiary field \(c_{\mathcal C}\).  Its algebraic
gauge slice has \(X=0\).  Therefore the BRST cohomology energy of the linear
planar domain is exactly the positive V59 radiative--continuum energy.  No
negative metric or edge polarization survives in the physical quotient.
\end{corollary}

\begin{proof}
The theorem identifies the V59 subsidiary variable with the actual
constraint and makes it zero.  V59 proves positivity and injectivity of
\(\mathfrak q_v\) on the two-polarization quotient.  V65 proves that the
quartet has only constant cohomology and that the full graph represents the
same two TT amplitudes.  Removing these contractible variables leaves
exactly \(\mathfrak q_v\).
\end{proof}

\begin{firewall}[What V66 does and does not close]
\label{fire:V66-scope}
Theorem~\ref{thm:V66-maximal-dissipative} closes the homogeneous edge
operator and Proposition~\ref{prop:V66-forward-adjoint-corners} closes the
forward/terminal adjoint pairing for prescribed incoming edge data.  It does
not yet construct a single autonomous domain in which the spin incoming
trace, the V36 canonical conormal, and the V59 continuum traction exchange
exactly the same positive power.  That port identification is required
before the passive storage may be added to the V59 energy or discarded as a
pure realization state.

The operator \(D_1\) in this note is defined raywise after rotating
\(k\) to \((\rho,0)\).  The measurable direct integral is rigorous, but a
finite-rank covariant differential spin bundle over all tangential
directions has not yet been constructed.  No nonlinear generalized-harmonic
source, curved corner, Brown--York stress theorem, BV--BFV master equation,
or coupled Higgs, fermion, neutrino, and Yang--Mills generator is claimed.
\end{firewall}

\subsection{Finite certificates}
\label{subsec:V66-certificates}

A finite validator used the unscaled order
\((b_+,y_+,z,b_\times,y_\times)\).  It checked one hundred random exact
harmonic cards, symmetry of the spin matrices, conservation
\(D_0R_+=0\), the two stable kernel vectors, rank three of the five spin
rows, equality with the V65 triangular kernel, and rank fifteen of
\eqref{eq:V66-spin-stable-certificate}.  For the diagnostic choice
\((m,v)=(0.8,0.5)\), the least singular values of the seventeen-by-fifteen
certificate were
\begin{center}
\begin{tabular}{c|ccccccc}
point&open&normal&static&glancing (+)&glancing (-)&threshold&evanescent\\
\hline
\(s_{\min}\)&
0.77808&0.82962&0.41795&0.23241&0.23241&0.29323&0.39678
\end{tabular}
\end{center}
The smallest value in two hundred deterministic open-cone samples was
\(0.23729\).  These samples detect transcription errors; compactness and
the analytic kernel proof establish the uniform theorem.

\subsection{V67 acquisition order}
\label{subsec:V66-next}

V67 should preserve the spin completion and proceed in this order:
\begin{enumerate}
\item construct a covariant angular spin bundle whose local matrices reduce
      to \(J_1,J_3,\sigma_1,\sigma_3\) in every rotated ray chart;
\item express the V36 canonical Green form and the V59 wave characteristics
      in the edge port variables \((\Gamma_{\rm in},\Gamma_{\rm out})\);
\item determine the unique boundary counterterm or port interconnection, if
      it exists, which cancels the metric, continuum, and realization fluxes
      without changing the Schur factor \(\lambda+m\kappa_v\);
\item prove maximality of the resulting closed autonomous generator by a
      closed form or a resolvent range theorem; if the port signatures do
      not match, go upstream to a covariant Friedrichs first-order lift
      rather than modifying the already correct stable determinant;
\item extend the algebraic BRST quartet to the resulting operator domain and
      compute the linear BFV boundary charge before any nonlinear claim;
\item only after these gates close, restore curvature, Brown--York/Bianchi
      stress balance, and the Standard-Model source sectors.
\end{enumerate}
The next compulsory companion audit remains V70.

\subsection{V66 result ledger and claim boundary}
\label{subsec:V66-ledger}

Relative to V1--V65, V66 proves:
\begin{enumerate}
\item the exact identity between the three V65 outgoing rows and the four
      component harmonic covector on the trace-free axial chart;
\item a Hille--Yosida obstruction to treating the naive formal-adjoint
      multiplier as a second positive forward transport;
\item the minimal symmetric spin-one/spin-one-half completion and conservation
      of its single plus-channel subsidiary constraint;
\item exact equality of the completed stable kernel with the V65 half-angle
      graph at every open and closed stable endpoint;
\item a seventeen-by-fifteen full-column-rank certificate with the unchanged
      V59 Schur factor and a positive normalized-cone floor;
\item a maximal dissipative half-collar realization, its adjoint domain,
      passive storage identity, and compatible forward/terminal corner data;
\item the exact Weyl transfer
      \(\operatorname{diag}(\chi^2,\chi)\) and recovery of the full TT metric
      frame; and
\item propagation of the actual linear harmonic constraint and its
      identification with the V59 Dirichlet subsidiary block.
\end{enumerate}

The closed autonomous metric--edge--continuum port, global differential
angular bundle, nonlinear BRST/BV--BFV action, curved-corner stress identity,
and coupled Standard-Model continuum remain open.  The full Standard
Model--Einstein continuum is not proved.

\subsection{Append-only record}

This V66 material was appended to the exact V65 parent whose SHA--256 digest
is
\[
 \texttt{b4507ef744099ad8b0be4bf3a6d4ffc47e3716860af8e2686aac10c980a1dfb9}.
\]
The next research note is V67 and begins with the covariant angular spin
bundle and the V36/V59 port-signature calculation.

% END APPENDED VERSION V66 MATERIAL

% BEGIN APPENDED VERSION V67 MATERIAL

\clearpage
\section{V67: polarization-covariant spin dilation, exact V36--V59 port
signature, and a differentiated unitary generator}
\label{sec:V67-port-dilation}

\subsection{Purpose and the cross-port correction}
\label{subsec:V67-purpose}

V66 proved that the raywise plus spin-one and cross spin-one-half systems
have a maximal dissipative incoming realization and that their stable state
reconstructs the V65 full TT tensor.  That result is correct as a passive
realization of the half-angle entries.  It does not yet identify its edge
flux with the physical V36--V59 wave flux.

The physical comparison detects one missing state.  The V59 metric wave has
characteristic scattering multiplier \(\chi^2 I_2\), whereas the V66 edge
output is \(\operatorname{diag}(\chi^2,\chi)\).  On the open cone the two
maps have different singular values, so no flux-unitary port change can
identify them.  V67 repairs this at the acquisition point: add a single
cross variable \(z_\times\), replace the cross spin-one-half block by a
second copy of the spin-one block, and retain the V65 metric as an
observation which does not use \(z_\times\).

The resulting six-state system is scalar on the rank-two V59 radiative
bundle.  It therefore patches across arbitrary measurable polarization
frames.  After one characteristic derivative its Green flux is exactly the
V36 metric wave flux.  The V59 common-trace and conormal conditions then
become the graph of an explicit unitary two-medium scattering matrix.  This
produces a skew-adjoint differentiated radiative--continuum generator.  A
local ten-component metric potential domain is still a later gate.

No companion audit is due at V67; the next Yang--Mills/Navier--Stokes audit
remains V70.

\subsection{The angular spin bundle}
\label{subsec:V67-angular-bundle}

Let \(\mathscr T\to\mathbb R^2_k\setminus\{0\}\) be the measurable
rank-two Hermitian radiative bundle of V59.  Write
\begin{equation}
 \mathscr S:=\mathscr T\otimes\mathbb C^3,
 \qquad
 \mathbf U=(b,\sqrt2y,z)^T,
 \label{eq:V67-spin-bundle}
\end{equation}
where \(b,y,z\) are sections of \(\mathscr T\).  The added component is the
cross entry of \(z=(z_+,z_\times)^T\).  Let
\begin{equation}
 \mathcal D_\parallel:=i|D_y|
 \label{eq:V67-Dparallel}
\end{equation}
be the skew-adjoint scalar multiplier with symbol \(i\rho\) in the rotated
ray chart.  With the V66 matrices \(J_3,J_1\), define
\begin{equation}
 \mathcal P_{67}\mathbf U
 :=\left(D_0I_3-D_nJ_3-\mathcal D_\parallel J_1\right)\mathbf U=0,
 \label{eq:V67-bundle-spin-system}
\end{equation}
where the displayed matrices act on the three-stage factor and the identity
acts on \(\mathscr T\).

In components, for each polarization index \(A\in\{+,\times\}\),
\begin{align}
 (D_0-D_n)b_A-\mathcal D_\parallel y_A&=0,
 \label{eq:V67-spin-top}\\
 2D_0y_A-\mathcal D_\parallel(b_A+z_A)&=0,
 \label{eq:V67-spin-middle}\\
 (D_0+D_n)z_A-\mathcal D_\parallel y_A&=0.
 \label{eq:V67-spin-bottom}
\end{align}
Define the two-component subsidiary variable
\begin{equation}
 R_A:=(D_0+D_n)y_A-\mathcal D_\parallel b_A.
 \label{eq:V67-spin-constraint}
\end{equation}

\begin{theorem}[Polarization-covariant spin-one dilation]
\label{thm:V67-angular-spin-dilation}
Equation \eqref{eq:V67-bundle-spin-system} is a well-defined operator on
\(\mathscr S\), independent of the measurable orthonormal polarization
frame used to write its two components.  It satisfies
\begin{equation}
 D_0R=0.
 \label{eq:V67-two-constraint-propagation}
\end{equation}
On the closed invariant subspace \(R=0\), all six components obey the
scalar wave equation.  At every nonzero open or limiting closed stable
frequency,
\begin{equation}
 \ker\mathcal P_{67}(s,\rho,\lambda)
 =\left\{a\otimes(1,\sqrt2\chi,\chi^2)^T:
                 a\in\mathscr T_k\right\},
 \label{eq:V67-stable-spin-bundle}
\end{equation}
and the stable kernel has dimension two.
\end{theorem}

\begin{proof}
A change of measurable polarization frame acts by the same unitary map on
\(b,y,z\).  Every coefficient in
\eqref{eq:V67-bundle-spin-system} is scalar on \(\mathscr T\), so it
commutes with that change.  The V66 computation of \(D_0R_+=0\) uses only
commutation of the derivatives and applies componentwise, proving
\eqref{eq:V67-two-constraint-propagation}.  With \(R=0\), the top and
constraint equations are opposite first-order factors for \((b,y)\), and
the constraint and bottom equations are the opposite factors for
\((y,z)\); their products give the wave equation.

At a stable mode the stage matrix is
\[
 M_+(s,\rho)=
 \begin{pmatrix}
 s-\lambda&-i\rho&0\\
 -i\rho&2s&-i\rho\\
 0&-i\rho&s+\lambda
 \end{pmatrix}.
\]
V66 proves that it has rank two and kernel
\(\operatorname{span}\bigl((1,\chi,\chi^2)^T\bigr)\).  Tensoring with the
two-dimensional identity gives \eqref{eq:V67-stable-spin-bundle}.
\end{proof}

For \(k\ne0\), choose the canonical oriented plane-of-incidence frame
\(\widehat k=k/|k|\), \(e=\epsilon\widehat k\).  In that frame let
\begin{equation}
 \mathfrak E^{(6)}_{\widehat k}(b,y,z)
 :=\mathfrak E_{\widehat k}(b,y,z_+),
 \label{eq:V67-six-state-observation}
\end{equation}
where the right side is the V65 algebraic metric chart.  The new
\(z_\times\) is an unobserved realization state.  On the stable kernel,
\begin{equation}
 \mathfrak E^{(6)}_{\widehat k}(b,\chi b,\chi^2b)
 =\mathcal E_\chi b.
 \label{eq:V67-observation-recovers-TT}
\end{equation}
Under a spatial rotation, the canonical plane-of-incidence frame and the
metric tensor rotate together, so this observation is equivariant on the
punctured tangential cone.  At normal incidence the stable value \(z=0\)
and the physical rank-two plane have the continuous V64 limit.

\begin{firewall}[Angular covariance versus differential locality]
\label{fire:V67-angular-locality}
The bundle operator is a rigorous Plancherel direct-integral operator, but
\(\mathcal D_\parallel=i|D_y|\) is pseudodifferential in physical boundary
coordinates.  A first-order scalar tangential symbol on \(\mathscr T\)
which is invariant under all planar rotations would have the form
\(k_AB^A\) with an invariant covector \((B^1,B^2)\); rotational invariance
forces both coefficients to vanish.  Hence the nonzero symbol
\(i|k|J_1\) cannot be replaced by a local first-order operator while the
stage factor is kept rotationally trivial.  A local covariant realization
must give the stage factor a nontrivial angular representation or add
fields.  V67 does not hide this remaining locality gate.
\end{firewall}

\subsection{The exact defect of the V66 cross port}
\label{subsec:V67-cross-port-defect}

For the V36/V59 radiative amplitude \(b\), define the metric wave
characteristics
\begin{equation}
 u_{\rm in}:=(D_0+D_n)b,
 \qquad
 u_{\rm out}:=(D_0-D_n)b.
 \label{eq:V67-metric-characteristics}
\end{equation}
On a stable mode,
\begin{equation}
 u_{\rm in}=(s+\lambda)b,
 \qquad
 u_{\rm out}=(s-\lambda)b
             =(s+\lambda)\chi^2b.
 \label{eq:V67-stable-metric-scattering}
\end{equation}
Thus the physical metric scattering map is \(\chi^2I_2\).

\begin{theorem}[No flux-unitary identification of the minimal V66 port]
\label{thm:V67-V66-port-no-go}
At every open stable frequency with \(0<|\chi|<1\), there are no unitary
input and output changes of basis which transform
\begin{equation}
 \operatorname{diag}(\chi^2,\chi)
 \label{eq:V67-old-edge-transfer}
\end{equation}
into \(\chi^2I_2\).  Consequently the minimal five-state V66 realization
cannot carry the V59 physical characteristic flux in both polarizations.
Among algebraic observations linear in the old cross states
\((b_\times,y_\times)\), one additional state is necessary to obtain the
multiplier \(\chi^2b_\times\) for every open frequency.
\end{theorem}

\begin{proof}
Unitary left and right multiplication preserve singular values.  The two
singular values of \eqref{eq:V67-old-edge-transfer} are
\((|\chi|,|\chi|^2)\), while those of \(\chi^2I_2\) are
\((|\chi|^2,|\chi|^2)\).  They differ when \(0<|\chi|<1\).

On the old stable cross kernel, every frequency-independent algebraic
linear observation has the form
\((a+c\chi)b_\times\).  If this were \(\chi^2b_\times\) on a nonempty open
set, the polynomial identity \(a+c\chi=\chi^2\) would hold for infinitely
many \(\chi\), which is impossible.  The new bottom state has exactly the
required stable value \(z_\times=\chi^2b_\times\).
\end{proof}
\subsection{Six-state stable certificate}
\label{subsec:V67-six-state-certificate}

In a canonical plus/cross frame use the order
\begin{equation}
 w_6=(b_+,y_+,z_+,b_\times,y_\times,z_\times).
 \label{eq:V67-six-state-order}
\end{equation}
Let \(M_{67}=M_+\oplus M_+\), and let
\(\mathfrak E^{(6)}\) be the ten-by-six observation matrix in
\eqref{eq:V67-six-state-observation}.  Define the eighteen-row,
sixteen-column stable certificate
\begin{equation}
 \widehat{\mathcal B}_{67}^{\rm spin}
 \binom h{w_6}
 :=\begin{pmatrix}
 h-\mathfrak E^{(6)}w_6\\
 M_{67}w_6\\
 (\lambda+m\kappa_v)(b_+,b_\times)^T
 \end{pmatrix}.
 \label{eq:V67-stable-certificate}
\end{equation}

\begin{theorem}[Complete stable floor after the port dilation]
\label{thm:V67-stable-floor}
For every \(m>0\), \(0<v<1\), and every nonzero open or limiting closed
stable frequency,
\begin{equation}
 \ker\widehat{\mathcal B}_{67}^{\rm spin}=0,
 \qquad
 \operatorname{rank}\widehat{\mathcal B}_{67}^{\rm spin}=16.
 \label{eq:V67-stable-rank}
\end{equation}
Its least singular value has a positive minimum on the normalized closed
retarded-root cone.  Elimination of the six realization variables leaves
exactly the V59 Schur factor
\begin{equation}
 (\lambda+m\kappa_v)I_2.
 \label{eq:V67-unchanged-Schur}
\end{equation}
\end{theorem}

\begin{proof}
The two spin-one blocks have total rank four and their common kernel is
\[
 y=\chi b,
 \qquad z=\chi^2b.
\]
On that kernel, \eqref{eq:V67-observation-recovers-TT} turns the first ten
rows into \(h=\mathcal E_\chi b\).  The last two rows and the V59 zero-free
theorem force \(b=0\), after which the spin kernel and graph rows force
\(y=z=h=0\).  This proves full column rank and the Schur statement.
Continuity of the matrix entries, compactness of the normalized retarded
cone, and pointwise full column rank give the positive least-singular-value
minimum.
\end{proof}

The extra state does not alter the metric representative, the action
conormal, or the BRST quartet.  It changes only the realization port which
was too small in the cross channel.

\subsection{The V36 metric Green form in spin ports}
\label{subsec:V67-metric-port}

V63 proved, in the equation normalization retained by V65, that restriction
of the V36 GHY--de Donder one-form to the physical radiative graph has
configuration trace \(b\) and conormal
\begin{equation}
 p_{\rm rad}=D_nb.
 \label{eq:V67-V36-conormal}
\end{equation}
Therefore its real power is
\begin{equation}
 \operatorname{Re}\langle D_0b,D_nb\rangle_{\mathscr T}
 =\frac14\left(\|u_{\rm in}\|^2-\|u_{\rm out}\|^2\right).
 \label{eq:V67-V36-characteristic-power}
\end{equation}

Let
\begin{equation}
 \mathcal P_+:=D_0+D_n,
 \qquad
 \mathbf W:=\mathcal P_+\mathbf U.
 \label{eq:V67-differentiated-spin-state}
\end{equation}
The derivatives commute at the planar anchor, so \(\mathbf W\) obeys the
same spin-one system.  Its characteristic traces are
\begin{equation}
 \Gamma_{\rm in}\mathbf W=\mathcal P_+b,
 \qquad
 \Gamma_{\rm out}\mathbf W=\mathcal P_+z.
 \label{eq:V67-W-traces}
\end{equation}

\begin{theorem}[Exact V36--spin port identity]
\label{thm:V67-V36-spin-port}
On every solution of the six-state spin system,
\begin{equation}
 \Gamma_{\rm in}\mathbf W=u_{\rm in},
 \qquad
 \Gamma_{\rm out}\mathbf W=u_{\rm out}.
 \label{eq:V67-exact-port-identification}
\end{equation}
Consequently the differentiated spin storage
\begin{equation}
 E_{e,1}:=\frac14\|\mathbf W\|_{\mathcal H_e}^2
 \label{eq:V67-differentiated-storage}
\end{equation}
has boundary power exactly equal to
\eqref{eq:V67-V36-characteristic-power}:
\begin{equation}
 \frac{\dd}{\dd t}E_{e,1}
 =\frac14\left(\|u_{\rm in}\|^2-\|u_{\rm out}\|^2\right).
 \label{eq:V67-spin-is-metric-power}
\end{equation}
This is an equality of Green signatures, including both polarizations and
all normal, glancing, static, threshold, and evanescent limits.
\end{theorem}

\begin{proof}
The incoming equality is the definition of \(u_{\rm in}\).  The top and
bottom spin equations give
\[
 (D_0-D_n)b=\mathcal D_\parallel y=(D_0+D_n)z,
\]
which is the outgoing equality.  Apply the V66 edge Green identity to
\(\mathbf W\).  Its unscaled storage is
\(\|\mathbf W\|^2/2\); multiplying that identity by one half gives
\eqref{eq:V67-spin-is-metric-power}.  Equation
\eqref{eq:V67-V36-characteristic-power} proves exact agreement with the
metric Green form.  The formulas are polynomial differential identities,
so their limiting-root values require no division.
\end{proof}

On a stable mode, the theorem reduces to
\begin{equation}
 \Gamma_{\rm in}\mathbf W=(s+\lambda)b,
 \qquad
 \Gamma_{\rm out}\mathbf W=(s+\lambda)z
 =(s-\lambda)b,
 \label{eq:V67-stable-derivative-ports}
\end{equation}
which explains both the missing factor in the V66 cross port and the exact
physical scattering \(\chi^2I_2\).
\subsection{The continuum port and its unitary interface matrix}
\label{subsec:V67-continuum-port}

Write the V59 continuum coordinate as \(r>0\), and let
\(\varphi(t,y,r)\) be its \(\mathscr T\)-valued field.  Introduce the
standard positive first-order variables
\begin{equation}
 V_0:=\sqrt m\,D_0\varphi,
 \qquad
 V_r:=\sqrt m\,\partial_r\varphi,
 \qquad
 V_A:=\sqrt m\,vD_A\varphi.
 \label{eq:V67-continuum-first-order-variables}
\end{equation}
They obey
\begin{align}
 D_0V_0&=\partial_rV_r+vD^AV_A,
 \label{eq:V67-continuum-V0}\\
 D_0V_r&=\partial_rV_0,
 \label{eq:V67-continuum-Vr}\\
 D_0V_A&=vD_AV_0.
 \label{eq:V67-continuum-VA}
\end{align}
The positive continuum energy is
\begin{equation}
 E_c=\frac12\int_{r>0}
 \left(|V_0|^2+|V_r|^2+|V_A|^2\right)\,\dd r\,\dd^2y.
 \label{eq:V67-continuum-energy}
\end{equation}
Define its boundary characteristics
\begin{equation}
 c_{\rm in}:=V_0-V_r,
 \qquad
 c_{\rm out}:=V_0+V_r.
 \label{eq:V67-continuum-characteristics}
\end{equation}
If \(t_b=-m\partial_r\varphi\) is the V59 traction, then
\begin{equation}
 c_{\rm in}=\sqrt m\,D_0\varphi+\frac{t_b}{\sqrt m},
 \qquad
 c_{\rm out}=\sqrt m\,D_0\varphi-\frac{t_b}{\sqrt m},
 \label{eq:V67-continuum-port-physical}
\end{equation}
and integration by parts gives
\begin{equation}
 \frac{\dd E_c}{\dd t}
 =\operatorname{Re}\langle D_0\varphi,t_b\rangle
 =\frac14\left(\|c_{\rm in}\|^2-\|c_{\rm out}\|^2\right).
 \label{eq:V67-continuum-port-power}
\end{equation}

The V59 interface is
\begin{equation}
 D_0b=D_0\varphi,
 \qquad
 p_{\rm rad}+t_b=0.
 \label{eq:V67-velocity-traction-interface}
\end{equation}
The first equation is the time derivative of common configuration trace;
that trace itself is fixed by one initial corner condition.

\begin{theorem}[Exact unitary V59 port graph]
\label{thm:V67-unitary-port-graph}
For every \(m>0\), define
\begin{equation}
 S_m:=\frac1{1+m}
 \begin{pmatrix}
 1-m&2\sqrt m\\
 2\sqrt m&m-1
 \end{pmatrix}.
 \label{eq:V67-interface-scattering}
\end{equation}
Then
\begin{equation}
 S_m^*=S_m,
 \qquad S_m^2=I_2.
 \label{eq:V67-Sm-unitary}
\end{equation}
The physical interface \eqref{eq:V67-velocity-traction-interface} is
equivalent to
\begin{equation}
 \binom{u_{\rm out}}{c_{\rm out}}
 =(S_m\otimes I_{\mathscr T})
 \binom{u_{\rm in}}{c_{\rm in}}.
 \label{eq:V67-unitary-port-relation}
\end{equation}
Consequently this boundary graph is maximal neutral for the total
incoming-minus-outgoing Green form and cancels the metric and continuum
powers exactly.
\end{theorem}

\begin{proof}
Put \(r_m=\sqrt m\), \(w=D_0b=D_0\varphi\), and
\(t=t_b=-p_{\rm rad}\).  Then
\begin{align*}
 u_{\rm in}&=w-t,&u_{\rm out}&=w+t,\\
 c_{\rm in}&=r_mw+t/r_m,&c_{\rm out}&=r_mw-t/r_m.
\end{align*}
Solving the first column for \((w,t)\) and substituting in the second gives
\eqref{eq:V67-unitary-port-relation}.  Direct multiplication gives
\eqref{eq:V67-Sm-unitary}.  Conversely the same invertible calculation
recovers common velocity and opposite traction.

Unitarity gives
\[
 \|u_{\rm in}\|^2+\|c_{\rm in}\|^2
 =\|u_{\rm out}\|^2+\|c_{\rm out}\|^2.
\]
Equations \eqref{eq:V67-spin-is-metric-power} and
\eqref{eq:V67-continuum-port-power} therefore cancel.  The graph of a
unitary between equal-dimensional incoming and outgoing spaces is
half-dimensional and equals its Green orthogonal, proving maximal
neutrality.
\end{proof}

\subsection{A skew-adjoint differentiated radiative--continuum generator}
\label{subsec:V67-differentiated-generator}

Let \(\mathscr A_s\) be the maximal two-copy spin-one differential
expression acting on \(\mathbf W\), and let \(\mathscr A_c\) be the maximal
continuum expression in
\eqref{eq:V67-continuum-V0}--\eqref{eq:V67-continuum-VA}.  Their graph
traces are respectively \((u_{\rm in},u_{\rm out})\) and
\((c_{\rm in},c_{\rm out})\).  Equip
\begin{equation}
 \mathcal H_{67}:=\mathcal H_s\oplus\mathcal H_c
 \label{eq:V67-generator-Hilbert-space}
\end{equation}
with the scaled inner product
\begin{equation}
 \langle(\mathbf W,\mathbf V),(\mathbf W',\mathbf V')\rangle_{67}
 :=\frac12\langle\mathbf W,\mathbf W'\rangle_{\mathcal H_s}
   +\langle\mathbf V,\mathbf V'\rangle_{\mathcal H_c}.
 \label{eq:V67-generator-inner-product}
\end{equation}
Define
\begin{equation}
 \mathscr G_{67}(\mathbf W,\mathbf V)
 :=(\mathscr A_s\mathbf W,\mathscr A_c\mathbf V)
 \label{eq:V67-generator-expression}
\end{equation}
on the graph domain cut out by
\eqref{eq:V67-unitary-port-relation}.

\begin{theorem}[Skew-adjoint differentiated port generator]
\label{thm:V67-skew-adjoint-generator}
The operator \(\mathscr G_{67}\) is skew-adjoint on
\(\mathcal H_{67}\).  It generates a strongly continuous unitary group and
conserves the positive energy
\begin{equation}
 E_{67}=\frac14\|\mathbf W\|_{\mathcal H_s}^2
       +\frac12\|\mathbf V\|_{\mathcal H_c}^2.
 \label{eq:V67-total-positive-energy}
\end{equation}
The boundary domain is maximal: no additional characteristic condition may
be imposed without reducing the adjoint domain.
\end{theorem}

\begin{proof}
The Green identity of the spin expression, with the factor one half in
\eqref{eq:V67-generator-inner-product}, and the continuum integration by
parts give
\begin{align*}
 &\langle\mathscr G_{67}X,Y\rangle_{67}
 +\langle X,\mathscr G_{67}Y\rangle_{67}\\
 &\quad=\frac12\left(
 \langle U_{\rm in},Y_{\rm in}\rangle
 -\langle U_{\rm out},Y_{\rm out}\rangle
ight),
\end{align*}
where \(U_{\rm in}=(u_{\rm in},c_{\rm in})\) and
\(U_{\rm out}=(u_{\rm out},c_{\rm out})\).  The unitary graph makes this
zero, so the operator is skew-symmetric.

The nonzero characteristic traces of both constant-coefficient Friedrichs
systems are onto their boundary trace spaces: smooth compact boundary data
can be extended independently into a thin collar, with all zero-speed
components set to zero.  Therefore the adjoint domain is obtained by taking
the Green orthogonal of the boundary graph.  Theorem
\ref{thm:V67-unitary-port-graph} says that this orthogonal is the same graph.
Thus \(\mathscr G_{67}^*=-\mathscr G_{67}\).  Stone's theorem gives the
unitary group, and its squared Hilbert norm is twice
\eqref{eq:V67-total-positive-energy}.
\end{proof}

\begin{proposition}[Smooth potential compatibility and recovery of the V59 Schur factor]
\label{prop:V67-potential-compatibility}
Let a smooth solution of the differentiated generator lie in the range of
smooth spin potentials \(\mathbf U\) with \(R=0\), and let its continuum
first-order state be the gradient of a smooth potential \(\varphi\).  If
\begin{equation}
 b|_{r=0}=\varphi|_{r=0}
 \label{eq:V67-initial-common-trace}
\end{equation}
holds at one time, the differential constraints and common trace remain
valid for that smooth potential solution.  For a zero-initial-data stable
mode, elimination of the two half-space solutions gives
\begin{equation}
 (\lambda+m\kappa_v)b=0.
 \label{eq:V67-port-Schur-recovery}
\end{equation}
\end{proposition}

\begin{proof}
The spin subsidiary variable is time-constant by
\eqref{eq:V67-two-constraint-propagation}.  The curl and mixed-derivative
constraints of \eqref{eq:V67-continuum-first-order-variables} propagate
because the constant derivatives commute.  The port graph is equivalent to
common boundary velocity and opposite traction; common velocity preserves
\eqref{eq:V67-initial-common-trace}.

For a stable solution, \(p_{\rm rad}=\lambda b\), while the decaying
continuum field has \(t_b=m\kappa_v\varphi(0)\).  Common trace and opposite
traction give \eqref{eq:V67-port-Schur-recovery}, exactly the V59 factor.
No density, closed-range, or inverse estimate for the potential map is used.
\end{proof}
\begin{firewall}[Differentiated generator versus the metric potential]
\label{fire:V67-differentiated-scope}
Theorem~\ref{thm:V67-skew-adjoint-generator} is a complete positive generator
for the differentiated spin state \(\mathbf W\) and the continuum strain
state.  It does not yet prove that the range of
\(\mathbf U\mapsto(D_0+D_n)\mathbf U\) is a closed invariant subspace with
a bounded inverse modulo the homogeneous zero mode.  Accordingly, it is not
yet a generator theorem for the undifferentiated ten-component metric
potential \(\widehat h=\mathfrak E^{(6)}\mathbf U\).

The V36 conormal identity used in
Theorem~\ref{thm:V67-V36-spin-port} is the established radiative-graph
identity.  A differentiable local action on the enlarged six-state field,
including the unobserved \(z_\times\), its adjoint multiplier, the V65 BRST
quartet, and all corner terms, has not been constructed.  The angular
operator is still pseudodifferential, and the real Fourier involution must
be recorded in any local angular dilation.  These are the V68 gates.
\end{firewall}

\subsection{Finite algebraic certificates}
\label{subsec:V67-certificates}

A finite validator checked:
\begin{enumerate}
\item rank two and the vector \((1,\chi,\chi^2)\) for each spin-one block;
\item propagation of both subsidiary variables \(D_0R_A=0\);
\item rank sixteen of the eighteen-by-sixteen certificate
      \eqref{eq:V67-stable-certificate};
\item the stable port identity
      \((s+\lambda)\chi^2=s-\lambda\);
\item orthogonality and involutivity of \(S_m\) for one hundred positive
      values of \(m\);
\item equivalence of the scattering graph with common velocity and opposite
      traction, and exact cancellation of the two fluxes; and
\item the strict V66 cross mismatch at the rational open point, where
      \(|\chi|=1/2\) but \(|\chi|^2=1/4\).
\end{enumerate}
For \((m,v)=(0.8,0.5)\), the least singular values of
\eqref{eq:V67-stable-certificate} were
\begin{center}
\begin{tabular}{c|ccccccc}
point&open&normal&static&glancing (+)&glancing (-)&threshold&evanescent\\
\hline
\(s_{\min}\)&
0.77808&0.82962&0.41795&0.23241&0.23241&0.29323&0.39678
\end{tabular}
\end{center}
The smallest value in two hundred deterministic open-cone samples was
\(0.19981\).  The samples are transcription checks; the kernel and
compactness arguments prove the rank and floor.

\subsection{V68 acquisition order}
\label{subsec:V67-next}

V68 should preserve the six-state port dilation and proceed in this order:
\begin{enumerate}
\item characterize the range and kernel of the characteristic derivative
      \(\mathbf U\mapsto(D_0+D_n)\mathbf U\) on the spin constraint space,
      and determine the correct homogeneous potential completion;
\item prove or disprove equivalence of the differentiated energy
      \eqref{eq:V67-total-positive-energy} with the V59 wave energy on that
      completed potential quotient;
\item construct an angularly local differential dilation with a recorded
      real structure, or prove the minimal extra representation required to
      replace \(i|D_y|\);
\item pull the unitary port graph back through
      \(\widehat h=Q_4(h,X)=\mathfrak E^{(6)}\mathbf U\), compute every V36
      canonical boundary term, and supply the adjoint and corner rows from
      one quadratic action;
\item compute the resulting linear BFV charge and prove BRST invariance of
      the closed operator domain;
\item only after those linear gates close, restore curvature, the nonlinear
      Brown--York/Bianchi identity, and the Higgs, fermion, neutrino, and
      Yang--Mills source sectors.
\end{enumerate}
If the potential derivative has an uncontrolled kernel, V68 must go upstream
to a standard first-order wave phase space rather than add a coercive term
which would change the physical Schur factor.  The next compulsory companion
audit remains V70.

\subsection{V67 result ledger and claim boundary}
\label{subsec:V67-ledger}

Relative to V1--V66, V67 proves:
\begin{enumerate}
\item a polarization-frame-covariant two-copy spin-one bundle on the
      punctured tangential Fourier cone;
\item propagation of its two subsidiary constraints and exact recovery of
      the V65 TT observation on the stable kernel;
\item a singular-value obstruction to identifying the five-state V66 cross
      port with the physical V59 port;
\item necessity, within algebraic linear observations of the old cross
      state, of one additional \(\chi^2\) realization variable;
\item an eighteen-by-sixteen complete stable certificate with a positive
      closed-cone floor and unchanged V59 Schur factor;
\item exact equality of the differentiated spin Green signature and the
      V36 radiative characteristic power;
\item the explicit self-adjoint unitary matrix \(S_m\) expressing the V59
      common-trace and traction interface; and
\item a skew-adjoint positive generator for the differentiated spin ports
      and continuum strains on that maximal unitary boundary graph.
\end{enumerate}

The undifferentiated metric potential generator, local real angular
dilation, complete V36/BRST/BFV action domain, nonlinear stress identity,
and coupled Standard-Model continuum remain open.  The full Standard
Model--Einstein continuum is not proved.

\subsection{Append-only record}

This V67 material was appended to the exact V66 parent whose SHA--256 digest
is
\[
 \texttt{decd71390935d4f40844d8be4d924dbb4668957fd93f5df6ab267a1e4e0569e2}.
\]
The next research note is V68 and begins with the characteristic-derivative
range and homogeneous potential-energy calculation.

% END APPENDED VERSION V67 MATERIAL

% BEGIN APPENDED VERSION V68 MATERIAL

\clearpage
\section{V68: grazing obstruction to the one-sided derivative, a local real
wave-phase dilation, and exact V59 energy equivalence}
\label{sec:V68-local-phase}

\subsection{Purpose and upstream correction}
\label{subsec:V68-purpose}

V67 constructed a correct unitary metric--continuum port graph and a
skew-adjoint generator for a differentiated spin state.  It explicitly left
open whether the one-sided characteristic derivative
\(\mathbf U\mapsto(D_0+D_n)\mathbf U\) has closed range and a bounded inverse
on the physical potential space.  V68 begins with that test.

The test fails at normal grazing.  On one wave sheet the multiplier
\(D_0+D_n\) tends quadratically to zero as the tangential frequency tends to
zero.  Its kernel has measure zero on the planar direct integral, but its
range is not closed and its norm is not equivalent to the V59 wave energy.
Thus the V67 differentiated generator cannot by itself be promoted to a
metric-potential generator.

The repair is the standard wave phase, written in characteristic form.  For
each radiative polarization retain the velocity, normal gradient, and both
tangential gradient components.  This adds the transverse gradient omitted
by the raywise spin system.  The resulting eight-state system is local,
real, symmetric hyperbolic, and exactly isometric to the V59 radiative
energy.  Its boundary ports are the V59 wave characteristics, so the unitary
matrix \(S_m\) of V67 supplies a skew-adjoint autonomous generator on the
closed potential-compatible subspace.

No companion audit is due at V68.  The next scheduled Yang--Mills and
Navier--Stokes audit remains V70.

\subsection{The one-sided characteristic derivative has nonclosed range}
\label{subsec:V68-grazing-obstruction}

Use full spatial Fourier variables \((\xi,k)\), where \(D_n\) has symbol
\(-i\xi\), \(\rho=|k|\), and
\begin{equation}
 \omega(\xi,\rho):=\sqrt{\xi^2+\rho^2}.
 \label{eq:V68-wave-frequency}
\end{equation}
On the positive wave sheet \(D_0=i\omega\), the V67 one-sided derivative
has symbol
\begin{equation}
 \mathfrak d_+(\xi,
ho)=i(\omega-\xi).
 \label{eq:V68-dplus-symbol}
\end{equation}
For \(1\le\xi\le2\),
\begin{equation}
 0<\omega-\xi
 =\frac{\rho^2}{\omega+\xi}
 \le\frac{\rho^2}{2}.
 \label{eq:V68-grazing-quadratic}
\end{equation}

\begin{theorem}[Normal-grazing nonclosed-range obstruction]
\label{thm:V68-nonclosed-range}
On the planar finite-energy direct integral, the restriction of
\(D_0+D_n\) to the positive constrained spin wave sheet is injective, but
its range is not closed and its inverse on the range is unbounded.  There is
no constant \(c>0\) such that
\begin{equation}
 \|(D_0+D_n)\mathbf U\|\ge c\|\mathbf U\|
 \label{eq:V68-no-coercive-dplus}
\end{equation}
for all band-limited physical spin states.  On a tangential torus the exact
\(k=0\) outgoing mode is an actual kernel, so the obstruction is stronger.
\end{theorem}

\begin{proof}
Restrict to the closed positive-sheet spectral subspace with
\(1\le\xi\le2\) and \(0<\rho<1\).  The multiplier
\(\omega-\xi\) is positive almost everywhere, so the map is injective on
this planar \(L^2\) subspace.  Choose normalized states \(U_j\) supported in
pairwise disjoint annuli
\[
 \varepsilon_j<\rho<2\varepsilon_j,
 \qquad \varepsilon_j\downarrow0,
\]
and in the constrained spin eigenline.  Equation
\eqref{eq:V68-grazing-quadratic} gives
\begin{equation}
 \|(D_0+D_n)U_j\|
 \le2\varepsilon_j^2\|U_j\|longrightarrow0.
 \label{eq:V68-grazing-sequence}
\end{equation}
Thus the map is not bounded below.  An injective bounded operator on this
band-limited Hilbert subspace has closed range only if its inverse on the
range is bounded, by the closed-range theorem.  Hence the range is not
closed.  When \(k=0\) is a genuine Fourier mode,
\(\omega=\xi>0\) makes \(\mathfrak d_+=0\) exactly.
\end{proof}

\begin{corollary}[The V67 differentiated norm is not the V59 potential norm]
\label{cor:V68-V67-not-equivalent}
The V67 energy \(\|(D_0+D_n)\mathbf U\|^2\) is positive on its own state
space but is not equivalent to the undifferentiated finite-energy wave norm.
No completion using only that norm can recover the missing normal-grazing
potential continuously.
\end{corollary}

\begin{proof}
The normalized sequence in
\eqref{eq:V68-grazing-sequence} has unit wave-phase norm and differentiated
norm tending to zero.  Equivalent Hilbert norms must bound each other by a
positive constant.
\end{proof}

\subsection{The exact local wave-phase map}
\label{subsec:V68-phase-map}

Let \(a(t,y,x)\) be one component of the V59 radiative potential and write
\begin{equation}
 v_a:=D_0a,
 \qquad g_n:=D_na,
 \qquad g_A:=D_Aa,
 \quad A\in\{1,2\}.
 \label{eq:V68-wave-jet}
\end{equation}
Define
\begin{equation}
 B:=\frac{v_a+g_n}{2},
 \qquad
 Z:=\frac{v_a-g_n}{2},
 \qquad
 Y_A:=\frac{g_A}{\sqrt2},
 \label{eq:V68-local-phase-map}
\end{equation}
and put \(\Phi=(B,Y_1,Y_2,Z)^T\).  Apply this construction componentwise
to the rank-two bundle \(\mathscr T\).

\begin{theorem}[Exact phase isometry and characteristic ports]
\label{thm:V68-phase-isometry}
The map \(\mathcal J:(a,D_0a)\mapsto\Phi\) satisfies
\begin{equation}
 \|\Phi\|^2
 =\frac12\left(
 |D_0a|^2+|D_na|^2+|D_1a|^2+|D_2a|^2\right).
 \label{eq:V68-phase-isometry}
\end{equation}
Thus, with the direct integral and both polarizations, its squared norm is
exactly the V59 radiative energy.  Its nonzero normal traces are
\begin{equation}
 u_{\rm in}=2B=D_0a+D_na,
 \qquad
 u_{\rm out}=2Z=D_0a-D_na.
 \label{eq:V68-exact-characteristic-ports}
\end{equation}
There is no grazing loss because both characteristic halves and the full
tangential gradient are retained.
\end{theorem}

\begin{proof}
The parallelogram identity gives
\[
 |B|^2+|Z|^2
 =\frac12\left(|v_a|^2+|g_n|^2\right),
\]
while \(|Y_1|^2+|Y_2|^2=(|g_1|^2+|g_2|^2)/2\).  Their sum is
\eqref{eq:V68-phase-isometry}.  Adding and subtracting the definitions of
\(B,Z\) gives \eqref{eq:V68-exact-characteristic-ports}.
\end{proof}
The potential range has an intrinsic local description.  Define the mixed
and tangential curl constraints
\begin{align}
 K_A(\Phi)&:=\sqrt2D_nY_A-D_A(B-Z),
 \qquad A=1,2,
 \label{eq:V68-mixed-constraints}\\
 K_{12}(\Phi)&:=D_1Y_2-D_2Y_1.
 \label{eq:V68-tangential-curl}
\end{align}
Let \(\mathcal R_\Phi\) be their common distributional kernel in the local
phase Hilbert space.

\begin{theorem}[Closed potential range]
\label{thm:V68-closed-potential-range}
Equip the phase space with
\begin{equation}
 \langle\Phi,\Psi\rangle_{\rm ph}:=2\langle\Phi,\Psi\rangle_{L^2}.
 \label{eq:V68-phase-inner-product}
\end{equation}
Then \(\mathcal J\) extends to a unitary map from the homogeneous V59
radiative energy space onto \(\mathcal R_\Phi\).  In particular,
\(\mathcal R_\Phi\) is closed, and \(a\) is recovered uniquely modulo the
homogeneous constant mode.
\end{theorem}

\begin{proof}
Equations \eqref{eq:V68-local-phase-map} give
\begin{equation}
 B+Z=D_0a,
 \qquad B-Z=D_na,
 \qquad \sqrt2Y_A=D_Aa.
 \label{eq:V68-inverse-jet}
\end{equation}
Commutation of the derivatives proves
\(K_A=K_{12}=0\), so the range lies in \(\mathcal R_\Phi\).
Theorem~\ref{thm:V68-phase-isometry} says exactly that the V59 Hilbert norm
is the norm induced by \eqref{eq:V68-phase-inner-product}.

Conversely, if \(\Phi\in\mathcal R_\Phi\), the spatial covector field
\[
 g=(B-Z,\sqrt2Y_1,\sqrt2Y_2)
\]
is curl-free in the distributional sense.  The half-space is simply
connected.  The homogeneous Sobolev Poincare lemma, equivalently division
by the nonzero spatial Fourier covector, gives
\(a\in\dot H^1\), unique modulo constants, with \(Da=g\).  Set
\(D_0a=B+Z\).  This reconstructs \eqref{eq:V68-local-phase-map} and proves
surjectivity.  An isometric image of a complete Hilbert space is closed.
\end{proof}

\subsection{A local real symmetric-hyperbolic dilation}
\label{subsec:V68-local-system}

The phase variables obey the first-order system
\begin{align}
 D_0B&=D_nB+\frac1{\sqrt2}D^AY_A,
 \label{eq:V68-local-B}\\
 D_0Y_A&=\frac1{\sqrt2}D_A(B+Z),
 \label{eq:V68-local-Y}\\
 D_0Z&=-D_nZ+\frac1{\sqrt2}D^AY_A.
 \label{eq:V68-local-Z}
\end{align}
Its normal matrix and tangential symbol are
\begin{equation}
 A_n=\operatorname{diag}(1,0,0,-1),
 \label{eq:V68-normal-matrix}
\end{equation}
and
\begin{equation}
 A(\zeta)=\frac1{\sqrt2}
 \begin{pmatrix}
 0&\zeta_1&\zeta_2&0\\
 \zeta_1&0&0&\zeta_1\\
 \zeta_2&0&0&\zeta_2\\
 0&\zeta_1&\zeta_2&0
 \end{pmatrix}.
 \label{eq:V68-tangential-matrix}
\end{equation}
Both are real symmetric, and
\(A(R\zeta)\) is obtained by the ordinary vector rotation of
\((Y_1,Y_2)\).  Thus the system is local and covariant under planar
rotations; no multiplier \(|D_y|\) occurs.

\begin{theorem}[Local phase equivalence to the radiative wave]
\label{thm:V68-local-wave-equivalence}
For every finite-energy wave solution
\begin{equation}
 D_0^2a=D_n^2a+D^AD_Aa,
 \label{eq:V68-radiative-wave}
\end{equation}
the phase map \eqref{eq:V68-local-phase-map} solves
\eqref{eq:V68-local-B}--\eqref{eq:V68-local-Z}.  Conversely, every solution
of the local phase system whose initial state lies in
\(\mathcal R_\Phi\) remains in \(\mathcal R_\Phi\) and is the phase map of
a unique homogeneous-energy solution of \eqref{eq:V68-radiative-wave}.
More precisely,
\begin{equation}
 D_0K_A=0,
 \qquad D_0K_{12}=0.
 \label{eq:V68-constraint-propagation}
\end{equation}
\end{theorem}

\begin{proof}
Insert \eqref{eq:V68-local-phase-map} and the wave equation.  For example,
\begin{align*}
 D_0B
 &=\frac12(D_0^2a+D_nD_0a)\\
 &=D_nB+\frac12D^AD_Aa
 =D_nB+\frac1{\sqrt2}D^AY_A,
\end{align*}
and the other two equations follow by the same differentiation.

For a phase solution, subtracting the \(Z\)-equation from the
\(B\)-equation and differentiating the \(Y\)-equation give
\begin{equation}
 D_0(B-Z)=D_n(B+Z),
 \qquad
 \sqrt2D_0Y_A=D_A(B+Z).
 \label{eq:V68-gradient-evolution}
\end{equation}
These identities imply \eqref{eq:V68-constraint-propagation} by commuting
derivatives.  Theorem~\ref{thm:V68-closed-potential-range} reconstructs
\(a\).  Equation \eqref{eq:V68-gradient-evolution} says that the spatial
gradient of \(D_0a-(B+Z)\) vanishes; the homogeneous convention removes
that constant.  Adding the \(B,Z\) equations now gives
\(D_0^2a=D_n^2a+D^AD_Aa\).
\end{proof}

\begin{corollary}[Exact local radiative energy and Green flux]
\label{cor:V68-local-energy}
For the two-polarization system, define
\begin{equation}
 E_{\rm ph}:=\int_{x>0}
 \left(|B|_{\mathscr T}^2+|Y|_{\mathscr T\otimes\mathbb R^2}^2
       +|Z|_{\mathscr T}^2\right)\,\dd x\,\dd^2y.
 \label{eq:V68-phase-energy}
\end{equation}
On \(\mathcal R_\Phi\), this is exactly the V59 radiative energy, and
\begin{equation}
 \frac{\dd E_{\rm ph}}{\dd t}
 =\int_{x=0}\left(|B|^2-|Z|^2\right)\,\dd^2y
 =\frac14\int_{x=0}
 \left(|u_{\rm in}|^2-|u_{\rm out}|^2\right)\,\dd^2y.
 \label{eq:V68-phase-flux}
\end{equation}
\end{corollary}

\begin{proof}
The energy equality is Theorem~\ref{thm:V68-phase-isometry}.  Multiply the
three local equations by \(2\overline B,2\overline Y_A,2\overline Z\), take
real parts, and integrate.  The tangential terms cancel because the matrices
\eqref{eq:V68-tangential-matrix} are symmetric.  Normal integration by
parts leaves \(|B|^2-|Z|^2\) at \(x=0\).  The last equality is
\eqref{eq:V68-exact-characteristic-ports}.
\end{proof}
\subsection{Stable reduction and a local phase certificate}
\label{subsec:V68-stable-certificate}

Rotate only for this finite symbol calculation so that \(k=(\rho,0)\).
For one polarization the local phase matrix and its curl row are
\begin{align}
 L_{68}(s,\rho,\lambda)
 &:={}
 \begin{pmatrix}
 s-\lambda&-i\rho/\sqrt2&0&0\\
 -i\rho/\sqrt2&s&0&-i\rho/\sqrt2\\
 0&0&s&0\\
 0&-i\rho/\sqrt2&0&s+\lambda
 \end{pmatrix},
 \label{eq:V68-local-stable-block}\\
 c_{68}(\rho)&:=\begin{pmatrix}0&0&i\rho&0\end{pmatrix}.
 \label{eq:V68-curl-row}
\end{align}
The column order is \((B,Y_1,Y_2,Z)\).

\begin{lemma}[Exact local stable phase line]
\label{lem:V68-local-stable-line}
At every nonzero open or limiting closed stable frequency,
\begin{equation}
 \ker\binom{L_{68}}{c_{68}}
 =\operatorname{span}\biggl(
 \begin{pmatrix}1\\\sqrt2\chi\\0\\\chi^2\end{pmatrix}
 \biggr),
 \qquad
 \operatorname{rank}\binom{L_{68}}{c_{68}}=3.
 \label{eq:V68-local-stable-kernel}
\end{equation}
Equivalently, for a radiative potential amplitude \(a\),
\begin{equation}
 \Phi(a)=
 \begin{pmatrix}
 (s+\lambda)a/2\\ i\rho a/\sqrt2\\0\\ (s-\lambda)a/2
 \end{pmatrix}.
 \label{eq:V68-stable-phase-map}
\end{equation}
\end{lemma}

\begin{proof}
Substitution uses
\((s+\lambda)\chi=i\rho\),
\((s+\lambda)\chi^2=s-\lambda\), and
\(2s\chi=i\rho(1+\chi^2)\).  If \(s\ne0\), the transverse row gives
\(Y_2=0\), and the remaining three-by-three block is the V66 spin-one block
under the invertible middle-coordinate scaling; it has rank two.  If
\(s=0\), then \(\rho>0\), the curl row gives \(Y_2=0\), and the
longitudinal block again has rank two.  This covers every nonzero closed
endpoint and proves the rank and kernel.  Multiplying the normalized kernel
by \((s+\lambda)a/2\) gives \eqref{eq:V68-stable-phase-map}.
\end{proof}

For two polarizations define the twelve-row, eight-column phase certificate
\begin{equation}
 \mathcal C_{68}:=
 \begin{pmatrix}
 L_{68}&0\\0&L_{68}\\
 c_{68}&0\\0&c_{68}\\
 (\lambda+m\kappa_v)e_B^T&0\\
 0&(\lambda+m\kappa_v)e_B^T
 \end{pmatrix},
 \qquad e_B=(1,0,0,0)^T.
 \label{eq:V68-phase-certificate}
\end{equation}

\begin{theorem}[Uniform local phase floor]
\label{thm:V68-local-phase-floor}
The matrix \(\mathcal C_{68}\) has rank eight at every nonzero open or
limiting closed stable frequency and a positive least-singular-value floor
on the normalized closed cone.  Its only physical zero test is again
\(\lambda+m\kappa_v\).  In potential normalization,
\(B=(s+\lambda)a/2\); the factor \((s+\lambda)/2\) is invertible with a
uniform normalized-cone floor, so this certificate has exactly the V59
Schur zero set.
\end{theorem}

\begin{proof}
By Lemma~\ref{lem:V68-local-stable-line}, the first ten rows leave one
amplitude in each polarization.  The last two rows kill those amplitudes by
the V59 zero-free theorem.  Hence the kernel is zero and the rank is eight.
The matrix is continuous on the compact normalized retarded-root cone, so
its nonzero least singular value has a positive minimum.  V64 supplies the
stated floor for \(s+\lambda\).
\end{proof}

The local phase also records the characteristic derivative of the metric.
In the aligned plus/cross frame,
\begin{equation}
 \mathfrak E\left(B,\frac1{\sqrt2}(Y_{+,1},Y_{\times,1})^T,Z_+\right)
 =\frac{s+\lambda}{2}\mathcal E_\chi a
 \label{eq:V68-metric-characteristic-observation}
\end{equation}
on the stable phase line.  The transverse gradient and \(Z_\times\) are
realization variables needed for local covariance and the physical port;
they do not change the V65 metric representative.

\subsection{The correct local unitary generator}
\label{subsec:V68-local-generator}

Let \(\mathscr A_{\rm ph}\) be the maximal local differential expression
\eqref{eq:V68-local-B}--\eqref{eq:V68-local-Z} for both polarizations, and
retain the V67 continuum first-order expression \(\mathscr A_c\).  Use the
boundary ports
\begin{equation}
 u_{\rm in}=2B,
 \qquad u_{\rm out}=2Z,
 \qquad
 c_{\rm in}=V_0-V_r,
 \qquad c_{\rm out}=V_0+V_r.
 \label{eq:V68-all-ports}
\end{equation}
On
\(\mathcal H_{68}=\mathcal H_{\rm ph}\oplus\mathcal H_c\), use the inner
product
\begin{equation}
 \langle(\Phi,V),(\Psi,W)\rangle_{68}
 :=2\langle\Phi,\Psi\rangle_{L^2}
   +\langle V,W\rangle_{\mathcal H_c}.
 \label{eq:V68-total-inner-product}
\end{equation}
Define
\begin{equation}
 \mathscr G_{68}(\Phi,V)
 :=(\mathscr A_{\rm ph}\Phi,\mathscr A_cV)
 \label{eq:V68-generator-expression}
\end{equation}
on the maximal graph domain satisfying the V67 unitary relation
\begin{equation}
 \binom{u_{\rm out}}{c_{\rm out}}
 =(S_m\otimes I_{\mathscr T})
 \binom{u_{\rm in}}{c_{\rm in}}.
 \label{eq:V68-unitary-boundary-domain}
\end{equation}

\begin{theorem}[Local skew-adjoint phase--continuum generator]
\label{thm:V68-skew-adjoint-generator}
The operator \(\mathscr G_{68}\) is skew-adjoint on
\(\mathcal H_{68}\) and generates a strongly continuous unitary group.
Its positive conserved energy is
\begin{equation}
 E_{68}=E_{\rm ph}+E_c
 =\frac12\|(\Phi,V)\|_{68}^2.
 \label{eq:V68-total-energy}
\end{equation}
The closed potential-compatible subspace consisting of
\begin{enumerate}
\item \(\Phi\in\mathcal R_\Phi\),
\item continuum states arising from a gradient potential \(\varphi\), and
\item one common initial boundary trace \(a|_{x=0}=\varphi|_{r=0}\) at \(t=0\),
\end{enumerate}
is reducing for the unitary group.  On that subspace,
\(\mathscr G_{68}\) is unitarily equivalent to the V59
radiative--continuum wave generator.
\end{theorem}

\begin{proof}
The Green form of the phase expression, multiplied by the factor two in
\eqref{eq:V68-total-inner-product}, is
\[
 2\langle B,B'\rangle-2\langle Z,Z'\rangle
 =\frac12\left(
 \langle u_{\rm in},u'_{\rm in}\rangle
 -\langle u_{\rm out},u'_{\rm out}\rangle\right).
\]
The continuum Green form is
\[
 \frac12\left(
 \langle c_{\rm in},c'_{\rm in}\rangle
 -\langle c_{\rm out},c'_{\rm out}\rangle\right).
\]
The graph of \(S_m\otimes I_{\mathscr T}\) is maximal neutral because
\(S_m\) is unitary.  As in V67, the nonzero characteristic trace map of the
maximal constant-coefficient systems is onto; taking the Green orthogonal
therefore gives the same boundary graph.  Hence
\(\mathscr G_{68}^*=-\mathscr G_{68}\).  Stone's theorem gives the unitary
group and conservation of \eqref{eq:V68-total-energy}.

The local phase constraints propagate by
\eqref{eq:V68-constraint-propagation}; the continuum curl and
mixed-derivative constraints propagate by the commuting first-order
system.  The unitary port graph is equivalent to common boundary velocity
and opposite traction.  It therefore preserves an initially common
potential trace.  These statements hold for positive and negative time, so
the closed compatible subspace is reducing.

Finally, Theorem~\ref{thm:V68-closed-potential-range} intertwines the
radiative wave equation with the phase system, preserves the exact V59
energy, and identifies its characteristics by
\eqref{eq:V68-exact-characteristic-ports}.  The continuum variables and
interface are the standard V59 first-order reduction.  This proves unitary
equivalence on the compatible subspace.
\end{proof}

\begin{corollary}[Resolution of the V67 potential-energy gate]
\label{cor:V68-V67-gate-resolved}
The differentiated V67 generator remains a valid auxiliary unitary system,
but it is not needed to define the physical potential energy.  The local
phase generator \(\mathscr G_{68}\) has a closed potential range and its
energy is exactly, rather than merely comparable to, the V59 positive
radiative--continuum energy.
\end{corollary}
\begin{firewall}[Local quotient phase versus the full metric action]
\label{fire:V68-scope}
V68 removes the one-sided derivative obstruction, supplies a real local
angular system, and proves exact equivalence with the V59 positive quotient
generator.  The local fields remain sections of the radiative bundle
\(\mathscr T\).  A local bundle map from their full Cauchy phase to the
undifferentiated ten-component metric \(\widehat h\), valid off the stable
shell and at all corners, has not yet been constructed.

Equation \eqref{eq:V68-metric-characteristic-observation} identifies a
stable characteristic derivative of the V65 metric representative.  It is
not a substitute for the V36 canonical one-form on a metric potential
domain.  In particular, V68 has not shown that the transverse gradient and
\(Z_\times\) form BRST doublets, has not supplied their adjoint multiplier
rows from one local action, and has not computed a BFV boundary charge.  No
nonlinear curvature, Brown--York/Bianchi stress, or coupled Standard-Model
source theorem is inferred.
\end{firewall}

\subsection{Finite certificates}
\label{subsec:V68-certificates}

The V68 validator checked two hundred random phase jets for the exact
isometry and port identities, and two hundred random local states for
propagation of all three spatial integrability constraints.  It checked the
local stable kernel, curl row, and rank eight of
\eqref{eq:V68-phase-certificate} at the open rational point, normal
incidence, static face, both glancing sheets, continuum threshold,
evanescent face, and a non-axis-aligned tangential covector.  For
\((m,v)=(0.8,0.5)\), the corresponding least singular values were
\begin{center}
\begin{tabular}{c|cccccccc}
point&open&normal&static&gl. (+)&gl. (-)&threshold&evanescent&angled\\
\hline
\(s_{\min}\)&
1.02241&1.00000&0.51502&0.30438&0.30438&0.37318&0.49430&0.90444
\end{tabular}
\end{center}
The smallest value in three hundred deterministic random open-cone samples
was \(0.07998\).  For \(\xi=1\) and
\(\rho=10^{-1},3\cdot10^{-2},10^{-2},3\cdot10^{-3}\), the normalized
one-sided factors were respectively
\[
 4.96\cdot10^{-3},\quad4.50\cdot10^{-4},\quad
 5.00\cdot10^{-5},\quad4.50\cdot10^{-6},
\]
confirming the quadratic grazing loss.  The analytic proofs, rather than
these samples, establish the theorems.

\subsection{V69 acquisition order}
\label{subsec:V68-next}

V69 should retain the local phase generator and proceed in this order:
\begin{enumerate}
\item construct the full metric Cauchy-phase map from
      \((a,D_0a)\) and \(\Phi\) without dividing by the one-sided
      characteristic derivative;
\item pull the V36 canonical one-form through
      \(\widehat h=Q_4(h,X)\) and the local phase constraints, keeping the
      transverse gradient and \(Z_\times\) rows explicit;
\item decide whether those extra local states are contractible realization
      pairs or require a genuine positive boundary storage term;
\item derive all adjoint and corner rows from one quadratic local action and
      prove that its operator domain is invariant under the V65 BRST
      differential;
\item compute the linear BFV charge, its Green pairing, and its nilpotency on
      the closed domain;
\item only after these gates close, restore curvature and the nonlinear
      Brown--York/Bianchi stress identity before coupling the Higgs,
      fermion, neutrino, and Yang--Mills sectors.
\end{enumerate}
If the full metric phase cannot be reconstructed locally from the radiative
bundle, V69 must go upstream to the ten-component generalized-harmonic wave
phase rather than reintroduce a half-angle inverse.  The next compulsory
companion audit remains V70.

\subsection{V68 result ledger and claim boundary}
\label{subsec:V68-ledger}

Relative to V1--V67, V68 proves:
\begin{enumerate}
\item injectivity but nonclosed range of the V67 one-sided characteristic
derivative on the planar positive wave sheet, with an exact grazing
sequence;
\item failure of the differentiated V67 norm to control the physical wave
potential;
\item an exact local phase map retaining velocity, normal gradient, and both
tangential gradients for each polarization;
\item unitary identification of its closed integrability subspace with the
V59 radiative energy space;
\item a real, rotationally covariant, symmetric-hyperbolic differential
system with exact propagation of the curl and mixed-gradient constraints;
\item exact V59 characteristic ports and positive radiative Green flux;
\item a local twelve-by-eight stable certificate with a positive closed-cone
floor and the unchanged physical Schur zero set; and
\item a skew-adjoint local phase--continuum generator whose compatible
potential subspace is unitarily equivalent to the V59 generator.
\end{enumerate}

The full ten-component metric potential action, BRST/BFV operator domain,
nonlinear Einstein stress closure, and coupled Standard-Model continuum
remain open.  The full Standard Model--Einstein continuum is not proved.

\subsection{Append-only record}

This V68 material was appended to the exact V67 parent whose SHA--256 digest
is
\[
 \texttt{5bf86a3b64d91a79c5780f48c0d6b61bcd39debb9892cf7fa0a81f9c2977a7a5}.
\]
The next research note is V69 and begins with the full metric Cauchy-phase
pullback and V36 canonical-form calculation.

% END APPENDED VERSION V68 MATERIAL

% BEGIN APPENDED VERSION V69 MATERIAL

\clearpage
\section{V69: ten-component characteristic phase, the differential
$Q_4$ jet repair, and the kinematic BFV charge}
\label{sec:V69-metric-jet-BFV}

\subsection{Purpose, upstream decision, and nonduplication boundary}
\label{subsec:V69-purpose}

V68 constructed a local real first-order phase for the two positive
radiative polarizations and proved exact equivalence with the V59
radiative--continuum energy.  Its remaining local-action problem cannot be
solved by applying the stable-fibre map $Q_4$ directly to that energy state.
There are two independent reasons.  First, the V68 phase contains the full
first derivative of the radiative potential, whereas the V36 canonical
one-form pairs a ten-component metric potential with its conormal.  Recovering
an undifferentiated potential pointwise from its derivative would reintroduce
the low-frequency inverse excluded in V68.  Second, the differential
Stueckelberg map
\[
 \widehat h=h+2D_{(\mu}X_{\nu)}
\]
maps the first jet of $X$ into the metric potential, but maps its second jet
into the metric phase.  The finite-dimensional V39 stable splitting does not
record this Sobolev-order shift.

V29 and V53 already derive the componentwise generalized-harmonic wave jet,
its positive Euclidean symmetrizer, and the harmonic subsidiary shell.  They
are imported here and are not reproved.  V36 already derives the exact
GHY--de Donder conormal, V39 the configuration-level split exact sequence,
V53 the two-dimensional radiative cohomology, and V65 the algebraic BRST
quartet.  The new work in V69 is:
\begin{enumerate}
\item the characteristic recombination of the full ten-component V53 phase,
      including an undifferentiated homogeneous-energy metric and its closed
      potential range;
\item an exact local expression for every V36 canonical coefficient in that
      phase, without division by a one-sided characteristic;
\item a sharp same-energy derivative-loss theorem for the differential
      $Q_4$ pullback;
\item a split exact second-jet replacement of $Q_4$, with a frequency-free
      triangular inverse;
\item the prolonged contractible BRST complex and its explicit cotangent
      BFV moment map; and
\item a classification of the extra V68 local variables by that complex.
\end{enumerate}
The BFV calculation below is kinematic: it is the exact charge on the
jet-extended boundary cotangent bundle.  A single spacetime action with all
ghost, multiplier, adjoint, and corner domains remains a later gate.

No companion audit is due in V69.  The mandatory inspection of the newest
Yang--Mills and Navier--Stokes notes occurs before V70 is written.

\subsection{The full characteristic metric phase}
\label{subsec:V69-full-phase}

Work at the planar flat anchor on
$\Omega=\mathbb R^2_y\times\mathbb R_+^x$, with
$D_n=-\partial_x$ and $A\in\{1,2\}$.  Let
$H=(H_{\mu\nu})$ be a real symmetric covariant tensor.  Trace reversal is
written
\begin{equation}
 \overline U_{\mu\nu}
 :=U_{\mu\nu}-\frac12\eta_{\mu\nu}
                  \eta^{\rho\sigma}U_{\rho\sigma}
 \label{eq:V69-trace-reversal}
\end{equation}
for every symmetric tensor $U$.

Define the tensor-valued characteristic phase
\begin{equation}
 B:=\frac{D_0H+D_nH}{2},\qquad
 Z:=\frac{D_0H-D_nH}{2},\qquad
 Y_A:=\frac{D_AH}{\sqrt2},
 \qquad \Phi_H:=(B,Y_1,Y_2,Z).
 \label{eq:V69-full-characteristic-map}
\end{equation}
Thus $B,Y_A,Z$ take values in
$\operatorname{Sym}^2(\mathbb R^{1,3})^*$; there are forty derivative
components.  Introduce the derivative selectors
\begin{equation}
 \mathscr D_0\Phi_H:=B+Z,\qquad
 \mathscr D_n\Phi_H:=B-Z,\qquad
 \mathscr D_A\Phi_H:=\sqrt2Y_A.
 \label{eq:V69-derivative-selectors}
\end{equation}
The compatible potential range is
\begin{equation}
 \mathcal R_{10}:=\left\{(H,\Phi_H):
  D_nH=\mathscr D_n\Phi_H,\quad
  D_AH=\mathscr D_A\Phi_H\right\},
 \label{eq:V69-full-potential-range}
\end{equation}
where $H$ is taken modulo its homogeneous constant mode.

\begin{theorem}[Ten-component characteristic phase isometry]
\label{thm:V69-full-phase-isometry}
Let $\mathcal E_{10}$ be the homogeneous Cauchy energy space of symmetric
tensors,
\begin{equation}
 \|(H,V)\|_{\mathcal E_{10}}^2
 :=\int_\Omega
   \left(|V|_{\rm E}^2+|D_nH|_{\rm E}^2
                  +|D_1H|_{\rm E}^2+|D_2H|_{\rm E}^2\right).
 \label{eq:V69-full-wave-norm}
\end{equation}
The map $(H,V)\mapsto(H,\Phi_H)$, with $V=D_0H$, is unitary from
$\mathcal E_{10}$ onto $\mathcal R_{10}$ when the latter is given the norm
\begin{equation}
 \|(H,\Phi_H)\|_{69}^2
 :=2\int_\Omega
 \left(|B|_{\rm E}^2+|Y_1|_{\rm E}^2
       +|Y_2|_{\rm E}^2+|Z|_{\rm E}^2\right).
 \label{eq:V69-full-phase-norm}
\end{equation}
In particular $\mathcal R_{10}$ is closed, the inverse uses
\begin{equation}
 V=B+Z,\qquad D_nH=B-Z,\qquad D_AH=\sqrt2Y_A,
 \label{eq:V69-full-phase-inverse}
\end{equation}
and no one-sided characteristic is divided out.
\end{theorem}

\begin{proof}
The parallelogram identity in the Euclidean symmetric-tensor fibre gives
\[
 2(|B|_{\rm E}^2+|Z|_{\rm E}^2)
 =|V|_{\rm E}^2+|D_nH|_{\rm E}^2,
 \]
and $2|Y_A|_{\rm E}^2=|D_AH|_{\rm E}^2$.  This proves equality of
\eqref{eq:V69-full-wave-norm} and \eqref{eq:V69-full-phase-norm}.
Equations \eqref{eq:V69-full-phase-inverse} are obtained by adding and
subtracting the definitions.  Conversely, the spatial part of
\eqref{eq:V69-full-phase-inverse} is already imposed in
\eqref{eq:V69-full-potential-range}; it reconstructs the homogeneous
potential, while the first equation reconstructs its velocity.  An
isometric image of the complete energy space is closed.
\end{proof}

The tensor-valued V68 system is
\begin{align}
 D_0H&=B+Z,
 \label{eq:V69-H-evolution}\\
 D_0B&=D_nB+\frac1{\sqrt2}D^AY_A,
 \label{eq:V69-B-evolution}\\
 D_0Y_A&=\frac1{\sqrt2}D_A(B+Z),
 \label{eq:V69-Y-evolution}\\
 D_0Z&=-D_nZ+\frac1{\sqrt2}D^AY_A.
 \label{eq:V69-Z-evolution}
\end{align}
It is the characteristic coordinate form of the V29/V53 wave bank, not a
new field equation.

\begin{proposition}[Exact propagation of the full potential graph]
\label{prop:V69-full-potential-propagation}
Set
\begin{equation}
 F_n:=B-Z-D_nH,\qquad
 F_A:=\sqrt2Y_A-D_AH.
 \label{eq:V69-potential-defects}
\end{equation}
Every solution of \eqref{eq:V69-H-evolution}--%
\eqref{eq:V69-Z-evolution} satisfies
\begin{equation}
 D_0F_n=0,\qquad D_0F_A=0.
 \label{eq:V69-potential-defect-propagation}
\end{equation}
Consequently $\mathcal R_{10}$ is reducing for the two-sided free wave
group, and its restriction is unitarily equivalent to the V53
ten-component wave evolution.
\end{proposition}

\begin{proof}
Subtract the $Z$ equation from the $B$ equation to obtain
$D_0(B-Z)=D_n(B+Z)$.  Equation \eqref{eq:V69-H-evolution} then gives
$D_0F_n=0$.  The $Y_A$ equation similarly gives
$D_0(\sqrt2Y_A)=D_A(B+Z)=D_AD_0H$, proving the other identities.
They hold for positive and negative time.  Theorem
\ref{thm:V69-full-phase-isometry} supplies the unitary identification.
\end{proof}

\subsection{The harmonic shell is an algebraic phase row}
\label{subsec:V69-harmonic-phase}

Trace reversal commutes with every derivative selector.  Define
\begin{equation}
 \mathscr C_\nu(\Phi_H)
 :=-(\overline B+\overline Z)_{0\nu}
   +(\overline B-\overline Z)_{n\nu}
   +\sqrt2\,(\overline Y_A)_{A\nu}.
 \label{eq:V69-phase-harmonic-row}
\end{equation}

\begin{lemma}[Exact harmonic readout]
\label{lem:V69-harmonic-readout}
On $\mathcal R_{10}$,
\begin{equation}
 \mathscr C_\nu(\Phi_H)
 =D^\mu\overline H_{\mu\nu}=C_\nu(H).
 \label{eq:V69-harmonic-readout}
\end{equation}
Hence the harmonic constraint is a local zero-order row of the derivative
phase.  If $H$ obeys the free reduced metric wave equation, then
\begin{equation}
 \Box\mathscr C_\nu=0.
 \label{eq:V69-harmonic-wave}
\end{equation}
The closed initial subspace
$\mathscr C|_{t=0}=D_0\mathscr C|_{t=0}=0$ is reducing.
\end{lemma}

\begin{proof}
In the chosen orthonormal frame,
$D^\mu=(-D_0,D_1,D_2,D_n)$.  Substitute
\eqref{eq:V69-derivative-selectors}; the result is exactly
\eqref{eq:V69-phase-harmonic-row}.  Commuting the constant-coefficient wave
operator with $D^\mu$ proves \eqref{eq:V69-harmonic-wave}.  Uniqueness for
the homogeneous wave equation proves invariance in both time directions.
Closedness is understood in the natural graph norm for the displayed
constraint and its time derivative.
\end{proof}

This does not replace the V53 radiative quotient.  The harmonic phase has
six solution amplitudes at a nonzero null covector; its residual wave-gauge
image has dimension four, and only their quotient has the two positive TT
classes.  V69 retains that distinction.

\subsection{Exact V36 canonical coefficients in characteristic variables}
\label{subsec:V69-V36-pullback}

For a boundary index $a\in\{0,1,2\}$ put
\begin{equation}
 \mathscr D_a\Phi_H=
 \begin{cases}
  B+Z,&a=0,\\
  \sqrt2Y_a,&a=1,2.
 \end{cases}
 \label{eq:V69-boundary-derivative-selector}
\end{equation}
Define
\begin{align}
 K^{69}_{ab}(\Phi_H)
 &:=\frac12\left[
   (B-Z)_{ab}-(\mathscr D_a\Phi_H)_{nb}
              -(\mathscr D_b\Phi_H)_{na}\right],
 \label{eq:V69-phase-K}\\
 \Pi_{69}^{ab}(\Phi_H)
 &:=(K^{69})^{ab}-\eta_\partial^{ab}
                         \eta_\partial^{cd}K^{69}_{cd},
 \label{eq:V69-phase-Pi}\\
 \mathfrak p_{69}^{\mu\nu}(\Phi_H)
 &:=e_a^\mu e_b^\nu\Pi_{69}^{ab}
 -2\left(n^{(\mu}\mathscr C^{\nu)}
          -\frac12\eta^{\mu\nu}\mathscr C_n\right).
 \label{eq:V69-phase-conormal}
\end{align}

\begin{theorem}[Local pullback of the GHY--de Donder one-form]
\label{thm:V69-V36-phase-pullback}
For every smooth $(H,\Phi_H)\in\mathcal R_{10}$,
\begin{equation}
 K^{69}_{ab}=K^{(1)}_{ab}(H),\qquad
 \Pi_{69}^{ab}=\Pi_{(1)}^{ab}(H),\qquad
 \mathfrak p_{69}=\mathfrak p_{36}(H).
 \label{eq:V69-V36-coefficient-equality}
\end{equation}
Consequently the complete V36 one-form is
\begin{equation}
 \Theta_{36}(H;\dot H)
 =\frac1{2\kappa}\int_{\partial\Omega}
       \mathfrak p_{69}^{\mu\nu}(\Phi_H)\dot H_{\mu\nu}.
 \label{eq:V69-V36-pulled-one-form}
\end{equation}
Equivalently, with the V65 coordinates
\begin{equation}
 y_a(H)=-2H_{na},\qquad
 y_n(H)=\eta_\partial^{ab}H_{ab}-H_{nn},
 \label{eq:V69-canonical-y}
\end{equation}
one has
\begin{equation}
 2\kappa\Theta_{36}
 =\int_{\partial\Omega}
 \left(\Pi_{69}^{ab}\dot H_{ab}
       +\mathscr C^a\dot y_a+\mathscr C_n\dot y_n\right).
 \label{eq:V69-canonical-splitting}
\end{equation}
At fixed $B+Z$, $Y_A$, and $H$, the map
\begin{equation}
 B-Z\longmapsto(\Pi_{69}^{ab},\mathscr C^a,\mathscr C_n)
 \label{eq:V69-characteristic-Legendre-map}
\end{equation}
is a ten-dimensional linear isomorphism.
\end{theorem}

\begin{proof}
On the potential range,
$B-Z=D_nH$ and
$\mathscr D_a\Phi_H=D_aH$.  Substitution in
\eqref{eq:V69-phase-K} gives the V36 formula
\[
 K^{(1)}_{ab}
 =\tfrac12(D_nH_{ab}-D_aH_{nb}-D_bH_{na}).
\]
Lemma \ref{lem:V69-harmonic-readout} identifies the de Donder term.
Equations \eqref{eq:V69-phase-Pi}--%
\eqref{eq:V69-phase-conormal} then give all three identities in
\eqref{eq:V69-V36-coefficient-equality}, and the V36 variational formula
gives \eqref{eq:V69-V36-pulled-one-form}.  The canonical splitting is the
algebraic identity of V65 with $h=H$ and $C=\mathscr C$.

At fixed tangential and time jet, varying $B-Z$ is exactly varying $D_nH$.
The final assertion is therefore the V65 normal Legendre isomorphism in the
characteristic coordinates.  No factor $D_0+D_n$ is inverted.
\end{proof}

For two phase fields, the antisymmetrization of
\eqref{eq:V69-V36-pulled-one-form} gives the full Green pairing
\begin{equation}
 2\kappa\Omega_{36}\bigl((H,\Phi_H),(K,\Phi_K)\bigr)
 =\int_{\partial\Omega}
 \left(\mathfrak p_{69}(\Phi_H)^{\mu\nu}K_{\mu\nu}
       -H_{\mu\nu}\mathfrak p_{69}(\Phi_K)^{\mu\nu}\right).
 \label{eq:V69-full-Green-pairing}
\end{equation}
Thus the transverse gradients are present in the conormal through the
tangential derivatives in $K^{69}$ and $\mathscr C$; they are not invisible
auxiliary coefficients.

\subsection{Why the ten-component positive wave port is not the action port}
\label{subsec:V69-inertia}

The principal de Donder power is governed by the V36 DeWitt form
\begin{equation}
 \mathbb G(U,V)=\overline U^{\mu\nu}V_{\mu\nu}.
 \label{eq:V69-DeWitt-form}
\end{equation}
The characteristic identity is
\begin{equation}
 \mathbb G(D_0H,D_nH)
 =\mathbb G(B,B)-\mathbb G(Z,Z).
 \label{eq:V69-DeWitt-characteristic-power}
\end{equation}

\begin{proposition}[Block-preserving positive-port obstruction]
\label{prop:V69-positive-port-obstruction}
There are no invertible real maps $T_{\rm in},T_{\rm out}$ on the
ten-component symmetric-tensor fibre such that, for all $B,Z$,
\begin{equation}
 \mathbb G(B,B)-\mathbb G(Z,Z)
 =|T_{\rm in}B|_{\rm E}^2-|T_{\rm out}Z|_{\rm E}^2.
 \label{eq:V69-impossible-positive-port}
\end{equation}
Hence the positive Euclidean ten-wave phase cannot be identified with the
V36 action port by a change of variables which preserves the incoming and
outgoing characteristic halves separately.  On the real TT radiative
quotient, however, $\mathbb G$ restricts to the positive Frobenius form, so
the V67--V68 two-polarization port identity is unaffected.
\end{proposition}

\begin{proof}
Setting $Z=0$ in \eqref{eq:V69-impossible-positive-port} would make
$\mathbb G(B,B)=|T_{\rm in}B|_{\rm E}^2$ positive definite.  V36 proves that
$\mathbb G$ has inertia $(6,4)$, a contradiction to Sylvester's law.
Equivalently, setting $B=0$ yields the same contradiction for the outgoing
block.  A TT representative is trace free and purely spatial in the V53
splitting; on it $\mathbb G(U,U)=|U|_{\rm E}^2$.  This proves the last
statement.
\end{proof}

The proposition locates the correct role of the full ten-wave symmetrizer.
It controls the gauge-fixed hyperbolic equation.  It is not a positive
realization of all V36 canonical directions.  The four negative DeWitt
directions must be removed or paired by the gauge/ghost complex before a
positive physical action-domain statement is possible.

\subsection{The differential $Q_4$ loses one energy derivative}
\label{subsec:V69-Q4-loss}

At the flat anchor the physical-space form of V39's invariant representative
is
\begin{equation}
 Q_4(h,X)_{\mu\nu}
 :=h_{\mu\nu}+D_\mu X_\nu+D_\nu X_\mu.
 \label{eq:V69-differential-Q4}
\end{equation}

\begin{theorem}[Sharp same-energy obstruction for the $Q_4$ phase]
\label{thm:V69-Q4-derivative-loss}
For every finite Sobolev index $s$, the map
\begin{equation}
 (h,X)\longmapsto
 \left(Q_4(h,X),D Q_4(h,X)\right)
 \label{eq:V69-Q4-phase-map}
\end{equation}
is not bounded from
$H^{s+1}\times H^{s+1}$ to the metric first-jet space
$H^{s+1}\times H^s$, when both input fields are charged at the same
first-derivative order.  More concretely, there is a sequence with
\begin{equation}
 \|DX_N\|_{H^s}=1,\qquad
 \|D(Q_4(0,X_N))\|_{H^s}\longrightarrow\infty.
 \label{eq:V69-Q4-loss-sequence}
\end{equation}
Thus the uniformly split V39 stable-fibre map does not, by itself, define a
bounded map of the V53/V69 energy spaces.
\end{theorem}

\begin{proof}
On a flat torus choose a nonzero constant covector $e_\nu$ and
\[
 X_N(x)=a_Ne^{iNx_1}e_\nu,
 \qquad
 a_N=N^{-1}(1+N^2)^{-s/2}.
\]
Up to the fixed Fourier normalization,
$\|DX_N\|_{H^s}=1$.  Choose $e_\nu$ so that
$D_1X_{N,\nu}+D_\nu X_{N,1}$ is nonzero.  Then
$D Q_4(0,X_N)$ contains a nonzero second derivative of $X_N$, and its
$H^s$ norm is bounded below by a fixed positive multiple of $N$.
This proves \eqref{eq:V69-Q4-loss-sequence}.  Taking $h=0$ proves the
unboundedness assertion.  The argument is the exact $Q_4$ specialization of
the V24 Fourier separator.
\end{proof}

There are two honest repairs.  One may charge $X$ at one higher Sobolev
order, or retain the required derivatives as independent jet variables and
impose their holonomy constraints.  V69 takes the latter route because it
keeps every boundary coefficient local and makes the BRST prolongation
explicit.

\subsection{A split exact second-jet replacement of $Q_4$}
\label{subsec:V69-jet-Q4}

Let
\begin{equation}
 P_{\beta\mid\mu\nu}\sim D_\beta h_{\mu\nu},\qquad
 J_{\alpha\nu}\sim D_\alpha X_\nu,
 \qquad
 K_{\beta\alpha\nu}\sim D_\beta J_{\alpha\nu}
 \label{eq:V69-jet-fields}
\end{equation}
be independent jet coordinates.  The indices $\alpha,\beta$ range over all
four spacetime directions.  Define the symmetrization maps
\begin{align}
 (\Sigma J)_{\mu\nu}&:=J_{\mu\nu}+J_{\nu\mu},
 \label{eq:V69-Sigma-J}\\
 (\Sigma K)_{\beta\mid\mu\nu}
 &:=K_{\beta\mu\nu}+K_{\beta\nu\mu}.
 \label{eq:V69-Sigma-K}
\end{align}
The jet invariant is
\begin{equation}
 Q_4^{(2)}(h,P,X,J,K)
 :=(H,W):=(h+\Sigma J,\ P+\Sigma K).
 \label{eq:V69-Q4-jet-map}
\end{equation}
Here $H$ has ten components and $W$ has forty; the source has
$10+40+4+16+64=134$ components.

For a jet gauge parameter
$q=(q^{(0)},q^{(1)},q^{(2)})$ of dimensions $4,16,64$, put
\begin{equation}
 G_4^{(2)}q
 :=(\Sigma q^{(1)},\Sigma q^{(2)},
                   -q^{(0)},-q^{(1)},-q^{(2)}).
 \label{eq:V69-jet-gauge-map}
\end{equation}

\begin{theorem}[Second-jet split exact sequence]
\label{thm:V69-jet-exact-sequence}
The frequency-free sequence
\begin{equation}
 0\longrightarrow\mathbb R^{84}
 \mathop{\longrightarrow}^{G_4^{(2)}}
 \mathbb R^{134}
 \mathop{\longrightarrow}^{Q_4^{(2)}}
 \mathbb R^{50}
 \longrightarrow0
 \label{eq:V69-jet-exact-sequence}
\end{equation}
is split exact.  A splitting is
\begin{equation}
 R_4^{(2)}(H,W):=(H,W,0,0,0),
 \qquad Q_4^{(2)}R_4^{(2)}=I_{50}.
 \label{eq:V69-jet-splitting}
\end{equation}
The square triangular coordinate map
\begin{equation}
 \mathcal T_{69}(h,P,X,J,K)
 :=(h+\Sigma J,P+\Sigma K,X,J,K)
 \label{eq:V69-triangular-coordinate-map}
\end{equation}
has inverse
\begin{equation}
 \mathcal T_{69}^{-1}(H,W,X,J,K)
 =(H-\Sigma J,W-\Sigma K,X,J,K).
 \label{eq:V69-triangular-inverse}
\end{equation}
Consequently both maps have universal finite operator bounds and the jet
splitting has no stable-root, glancing, or low-frequency denominator.
\end{theorem}

\begin{proof}
Equations \eqref{eq:V69-Q4-jet-map} and
\eqref{eq:V69-jet-gauge-map} give
$Q_4^{(2)}G_4^{(2)}=0$.  The last three components of
$G_4^{(2)}q$ show that it is injective.  The displayed right inverse proves
surjectivity of $Q_4^{(2)}$.

If $Q_4^{(2)}(h,P,X,J,K)=0$, then
$h=-\Sigma J$ and $P=-\Sigma K$.  Taking
\[
 q^{(0)}=-X,\qquad q^{(1)}=-J,
 \qquad q^{(2)}=-K
\]
in \eqref{eq:V69-jet-gauge-map} reproduces
$(h,P,X,J,K)$.  Thus the kernel is the gauge image.  Equations
\eqref{eq:V69-triangular-coordinate-map}--%
\eqref{eq:V69-triangular-inverse} are direct substitutions.  They contain
only the fixed symmetrization maps, so their norms are bounded by numerical
constants independent of every Fourier parameter.
\end{proof}

The holonomic subspace which returns the differential map is cut out by
\begin{align}
 F^{(1)}_{\alpha\nu}&:=J_{\alpha\nu}-D_\alpha X_\nu,
 \label{eq:V69-first-holonomy}\\
 F^{(2)}_{\beta\alpha\nu}&:=K_{\beta\alpha\nu}
                         -D_\beta J_{\alpha\nu},
 \label{eq:V69-second-holonomy}\\
 F^h_{\beta\mid\mu\nu}&:=P_{\beta\mid\mu\nu}
                         -D_\beta h_{\mu\nu}.
 \label{eq:V69-metric-holonomy}
\end{align}

\begin{corollary}[Exact recovery of the differential phase]
\label{cor:V69-holonomic-Q4}
On the common zero set of
\eqref{eq:V69-first-holonomy}--\eqref{eq:V69-metric-holonomy},
\begin{equation}
 H=Q_4(h,X),\qquad W=DH.
 \label{eq:V69-holonomic-invariant}
\end{equation}
Conversely every sufficiently regular differential $Q_4$ configuration has
this jet lift.  Thus the jet extension removes the same-space unbounded map
by changing coordinates, not by asserting a false Sobolev estimate.
\end{corollary}

\begin{proof}
Substitute $J=DX$, $K=DJ$, and $P=Dh$ in
\eqref{eq:V69-Q4-jet-map}.  Then
$H=h+\Sigma(DX)=Q_4(h,X)$ and
$W=Dh+\Sigma(DJ)=D(h+\Sigma DX)=DH$.  Taking the displayed derivatives of
any regular $(h,X)$ gives the converse lift.
\end{proof}

The forty components of $W$ are converted to
$(B,Y_1,Y_2,Z)$ by the fixed invertible characteristic transformation
\eqref{eq:V69-full-characteristic-map}.  Hence
$Q_4^{(2)}$ lands directly in the full phase of
Theorem~\ref{thm:V69-full-phase-isometry}.

\subsection{Prolonged BRST complex and holonomy tangency}
\label{subsec:V69-prolonged-BRST}

Introduce ghost jets
\begin{equation}
 c_\nu,\qquad c^{(1)}_{\alpha\nu},\qquad
 c^{(2)}_{\beta\alpha\nu}
 \label{eq:V69-ghost-jets}
\end{equation}
with the same $4+16+64$ component count as the jet gauge parameter.  Define
the odd differential
\begin{align}
 sh&=\Sigma c^{(1)},&
 sP&=\Sigma c^{(2)},
 \label{eq:V69-BRST-hP}\\
 sX&=-c,&sJ&=-c^{(1)},&sK&=-c^{(2)},
 \label{eq:V69-BRST-XJK}\\
 sc&=0,&sc^{(1)}&=0,&sc^{(2)}&=0.
 \label{eq:V69-BRST-ghosts}
\end{align}

\begin{theorem}[Jet-level invariant quotient and cohomology]
\label{thm:V69-jet-BRST-cohomology}
The differential satisfies $s^2=0$ and
\begin{equation}
 sH=0,\qquad sW=0.
 \label{eq:V69-BRST-invariant-HW}
\end{equation}
On the polynomial jet algebra, its cohomology is precisely the polynomial
algebra in the fifty invariant coordinates $(H,W)$:
\begin{equation}
 H^\bullet(s)\cong\mathbb R[H,W].
 \label{eq:V69-jet-cohomology}
\end{equation}
In particular the prolonged gauge variables introduce no new physical
cohomology.
\end{theorem}

\begin{proof}
Nilpotency follows because all ghost jets are closed.  The cancellation
$s(h+\Sigma J)=\Sigma c^{(1)}-\Sigma c^{(1)}=0$ and its $P,K$ analogue prove
\eqref{eq:V69-BRST-invariant-HW}.  Use the invertible coordinates
\eqref{eq:V69-triangular-coordinate-map}.  In them $s$ is zero on $(H,W)$
and consists of the doublets
\[
 (X,-c),\qquad(J,-c^{(1)}),\qquad(K,-c^{(2)}).
\]
The homotopy defined by
$\iota c=-X$, $\iota c^{(1)}=-J$, and
$\iota c^{(2)}=-K$, and zero on the first member of each pair, satisfies
$s\iota+\iota s=N_{\rm jet}$, the total degree in the doublet variables.
Every closed polynomial of positive doublet degree is exact.  The
degree-zero polynomials are exactly $\mathbb R[H,W]$.
\end{proof}

Define the ghost holonomy defects
\begin{align}
 R^{(1)}_{\alpha\nu}
 &:=c^{(1)}_{\alpha\nu}-D_\alpha c_\nu,
 \label{eq:V69-first-ghost-holonomy}\\
 R^{(2)}_{\beta\alpha\nu}
 &:=c^{(2)}_{\beta\alpha\nu}
                  -D_\beta c^{(1)}_{\alpha\nu}.
 \label{eq:V69-second-ghost-holonomy}
\end{align}

\begin{proposition}[BRST tangency of the holonomic phase]
\label{prop:V69-holonomy-tangency}
The holonomy defects obey
\begin{align}
 sF^{(1)}&=-R^{(1)},
 \label{eq:V69-sF1}\\
 sF^{(2)}&=-R^{(2)},
 \label{eq:V69-sF2}\\
 sF^h_{\beta\mid\mu\nu}
 &=R^{(2)}_{\beta\mu\nu}+R^{(2)}_{\beta\nu\mu},
 \label{eq:V69-sFh}\\
 sR^{(1)}&=sR^{(2)}=0.
 \label{eq:V69-sR}
\end{align}
Therefore the common holonomic subspace $F=R=0$ is BRST invariant.  Any
boundary relation written only in $(H,W)$, including the pulled V36
coefficient relation, is automatically tangent to this BRST vector field.
\end{proposition}

\begin{proof}
Apply $s$ to \eqref{eq:V69-first-holonomy}--%
\eqref{eq:V69-metric-holonomy}, commute $s$ with the fixed derivatives, and
use \eqref{eq:V69-BRST-hP}--\eqref{eq:V69-BRST-ghosts}.  This gives the four
displayed identities.  Equations \eqref{eq:V69-BRST-invariant-HW} then prove
the final assertion.
\end{proof}

For completeness, add antighosts
$\bar c,\bar c^{(1)},\bar c^{(2)}$ and multipliers
$L,L^{(1)},L^{(2)}$, with
\begin{equation}
 s\bar c=L,\qquad s\bar c^{(1)}=L^{(1)},\qquad
 s\bar c^{(2)}=L^{(2)},\qquad sL=sL^{(1)}=sL^{(2)}=0.
 \label{eq:V69-antighost-doublets}
\end{equation}
The algebraic jet gauge fermion
\begin{equation}
 \Psi_{69}:=\langle\bar c,X\rangle
 +\langle\bar c^{(1)},J\rangle
 +\langle\bar c^{(2)},K\rangle
 \label{eq:V69-jet-gauge-fermion}
\end{equation}
has the exact term
\begin{align}
 s\Psi_{69}
 ={}&\langle L,X\rangle+\langle\bar c,c\rangle
 +\langle L^{(1)},J\rangle+\langle\bar c^{(1)},c^{(1)}\rangle
 \nonumber\\
 &+\langle L^{(2)},K\rangle+\langle\bar c^{(2)},c^{(2)}\rangle,
 \label{eq:V69-exact-gauge-fixing}
\end{align}
with the same odd-Leibniz sign convention as V65.  Its Euler slice sets the
jet gauge variables and ghosts to zero, leaving exactly $(H,W)$.

\subsection{The kinematic jet-BFV charge}
\label{subsec:V69-BFV-charge}

Let $\mathcal Z_{69}$ denote the graded jet coordinate space consisting of
the fields in
\eqref{eq:V69-BRST-hP}--\eqref{eq:V69-antighost-doublets}.  On its graded
cotangent bundle use canonical coordinates $(q^I,p_I)$ and the canonical
graded symplectic form $\omega_{\rm can}$.  Let $Q_{69}$ be the odd vector
field defined by the displayed BRST differential.

\begin{definition}[Kinematic BFV moment map]
\label{def:V69-kinematic-BFV}
Define
\begin{equation}
 \boxed{\ \Omega_{69}^{\rm kin}:=
  \sum_I(-1)^{|q^I|}\langle p_I,Q_{69}q^I\rangle\ }
 \label{eq:V69-BFV-charge}
\end{equation}
with the canonical sign chosen so that its Hamiltonian vector field on base
coordinates is $Q_{69}$.  In components its nonzero base part is
\begin{align}
 \Omega_{69}^{\rm kin}
 ={}&\langle p_h,\Sigma c^{(1)}\rangle
     +\langle p_P,\Sigma c^{(2)}\rangle
     -\langle p_X,c\rangle
     -\langle p_J,c^{(1)}\rangle
     -\langle p_K,c^{(2)}\rangle
 \nonumber\\
 &+\langle p_{\bar c},L\rangle
  +\langle p_{\bar c^{(1)}},L^{(1)}\rangle
  +\langle p_{\bar c^{(2)}},L^{(2)}\rangle,
 \label{eq:V69-BFV-charge-components}
\end{align}
up to the single global canonical sign fixed in
\eqref{eq:V69-BFV-charge}.
\end{definition}

\begin{theorem}[Hamiltonian generation, Green incidence, and nilpotency]
\label{thm:V69-BFV-nilpotency}
The charge has ghost number one and satisfies
\begin{equation}
 \iota_{Q_{69}^{T^*}}\omega_{\rm can}
 =\delta\Omega_{69}^{\rm kin},
 \qquad
 \{\Omega_{69}^{\rm kin},\Omega_{69}^{\rm kin}\}=0.
 \label{eq:V69-BFV-master}
\end{equation}
Here $Q_{69}^{T^*}$ is the cotangent lift.  Moreover the pullback of the V36
canonical one-form through $Q_4^{(2)}$ is horizontal and invariant:
\begin{equation}
 \iota_{Q_{69}}(Q_4^{(2)})^*\Theta_{36}=0,
 \qquad
 \mathcal L_{Q_{69}}(Q_4^{(2)})^*\Theta_{36}=0.
 \label{eq:V69-BFV-Green-incidence}
\end{equation}
Thus the linear jet charge is compatible with the exact V36 Green pairing
before any boundary gauge slice is imposed.
\end{theorem}

\begin{proof}
The cotangent-lift moment-map identity is the defining property of
\eqref{eq:V69-BFV-charge}; a direct differentiation gives the first equation
in \eqref{eq:V69-BFV-master}.  The Hamiltonian vector field of the self
bracket is twice the graded commutator of the lifted vector field with
itself.  This vanishes because $Q_{69}^2=0$.  The bracket is therefore a
constant of ghost number two.  There is no nonzero constant of ghost number
two, so it is zero.

By \eqref{eq:V69-BRST-invariant-HW}, both the potential and derivative
arguments of the pulled V36 one-form are annihilated by $Q_{69}$.  Its
contraction with the gauge vector is therefore zero.  Cartan's formula and
the same invariance give the Lie-derivative identity.  Antisymmetrizing the
one-form gives \eqref{eq:V69-full-Green-pairing}, proving the asserted Green
incidence.
\end{proof}

\begin{firewall}[Kinematic BFV charge versus a BFV boundary theory]
\label{fire:V69-BFV-scope}
Theorem~\ref{thm:V69-BFV-nilpotency} is an exact finite-jet cotangent
calculation.  It does not yet identify all $p_I$ with conormals derived from
one spacetime action, choose the ghost and multiplier incoming/outgoing
domains, or prove that the resulting unbounded operator is closed and
skew-adjoint.  Normal integration by parts in the holonomy constraints
produces additional multiplier traces, and temporal integration produces
initial/final corner terms.  Until those terms and domains are written and
their maximality is proved, $\Omega_{69}^{\rm kin}$ must not be called the
full dynamical BFV charge of the coupled system.
\end{firewall}

\subsection{Classification of the extra V68 local variables}
\label{subsec:V69-extra-state-classification}

The prolonged cohomology gives a definite answer to one V68 acquisition
question.  In an aligned stable chart V68's transverse gradient satisfies
$Y_2=0$ on one plane wave, but in spacetime it is the local field
$D_2a/\sqrt2$.  The cross outgoing field is
$Z_\times=(D_0-D_n)a_\times/2$ and is half of the physical outgoing port.

\begin{proposition}[The V68 completion fields are not jet-gauge doublets]
\label{prop:V69-V68-fields-not-doublets}
The transverse gradients $Y_{A,\perp}$ and the outgoing cross component
$Z_\times$ are BRST-invariant radiative phase coordinates and are not
members of the contractible $(X,J,K)$ jet sector.  Any eventual
BRST-invariant graph realizing their V53 class in the full metric phase must
map them into the invariant coordinates $(H,W)$; V69 does not assume that
still-missing local graph.  On the compatible radiative range:
\begin{enumerate}
\item $Y_{A,\perp}$ is fixed by the potential integrability constraints and
      has no independent normal characteristic trace;
\item $Z_\times$ is the physical outgoing characteristic trace; and
\item both contributions are already included in the V68 bulk energy, so
      neither requires an additional positive boundary storage term.
\end{enumerate}
Declaring either field BRST exact would enlarge the BRST differential and
change the V53/V59 physical cohomology; it is not a consequence of the V65
or V69 gauge complex.
\end{proposition}

\begin{proof}
V65 declares the radiative variables BRST invariant, and V69 prolongs the
differential only on the metric-embedding jet sector and its ghosts.  Thus
adjoining the V68 phase to the complex in
Theorem~\ref{thm:V69-jet-BRST-cohomology} does not contract any of its
components.  This conclusion precedes, and does not assume, a local graph
from the radiative phase to $(H,W)$.  The V68 potential-range equations
identify the transverse gradients with $D_Aa/\sqrt2$, while its normal
matrix assigns zero normal speed to the $Y_A$ block.  Its port theorem identifies
$2Z_\times=u_{{\rm out},\times}$.  Finally, the exact V68 phase isometry
contains $|Y|^2+|Z|^2$ in the already conserved bulk energy.  These facts
prove all three assertions.
\end{proof}

\subsection{Finite algebraic certificates}
\label{subsec:V69-certificates}

The new constructions reduce to the following exact cards:
\begin{align}
 2(|B|_{\rm E}^2+|Y|_{\rm E}^2+|Z|_{\rm E}^2)
 &=|D_0H|_{\rm E}^2+|D_nH|_{\rm E}^2+|D_AH|_{\rm E}^2,
 \label{eq:V69-phase-card}\\
 \mathfrak p_{69}(\Phi_H)&=\mathfrak p_{36}(H),
 \label{eq:V69-conormal-card}\\
 Q_4^{(2)}G_4^{(2)}&=0,
 \qquad
 \operatorname{rank}G_4^{(2)}=84,
 \qquad
 \operatorname{rank}Q_4^{(2)}=50,
 \label{eq:V69-jet-rank-card}\\
 \operatorname{rank}\mathcal T_{69}&=134,
 \qquad
 (\mathcal T_{69})^{-1}\mathcal T_{69}=I_{134},
 \label{eq:V69-triangular-card}\\
 s^2&=0,\qquad s(H,W)=0,
 \qquad
 \{\Omega_{69}^{\rm kin},\Omega_{69}^{\rm kin}\}=0.
 \label{eq:V69-BFV-card}
\end{align}
A finite verifier accompanying this iteration constructs the
$50\times134$ and $134\times84$ matrices in an orthonormal symmetric-tensor
basis, checks their ranks and exact product, checks the square triangular
inverse, samples the phase and harmonic identities, and recovers the V36
DeWitt inertia $(6,4)$.  The analytic arguments above establish the results;
the finite calculation guards the component conventions.

\subsection{V70 acquisition order and mandatory companion audit}
\label{subsec:V69-next}

Before V70 chooses a direction, it must inspect the newest completed
Yang--Mills and Navier--Stokes notes in the shared folder, record their exact
versions and digests, and import only type-correct information.  It should
then proceed in this order:
\begin{enumerate}
\item add multiplier fields for
      \eqref{eq:V69-first-holonomy}--\eqref{eq:V69-metric-holonomy} and derive
      their normal and temporal boundary terms from one quadratic local
      spacetime action;
\item prolong the V34 diffeomorphism-ghost wave bank to the same second-jet
      order and select complementary ghost, antighost, and multiplier
      characteristic domains;
\item prove maximality of the combined forward--adjoint graph and identify
      the canonical momenta $p_I$ in
      \eqref{eq:V69-BFV-charge-components};
\item restrict that domain to the harmonic/radiative cohomology and verify
      that its positive Green power is exactly the V68--V59 unitary port,
      with no contribution from the contractible jet quartet;
\item only then promote $\Omega_{69}^{\rm kin}$ to a dynamical linear BFV
      charge and test the corner master equation;
\item after the linear action gates close, restore curvature and the
      Brown--York/Bianchi stress identity before coupling the Higgs,
      fermion, neutrino, and Yang--Mills sources.
\end{enumerate}
If the multiplier traces cannot be made complementary without mixing the
positive and negative DeWitt characteristic sectors, V70 must keep the
Krein/BRST sector explicitly and reduce only after cohomology.  It must not
replace the action form by the Euclidean ten-wave norm.

\subsection{V69 result ledger and claim boundary}
\label{subsec:V69-ledger}

Relative to V1--V68, this note proves:
\begin{enumerate}
\item a closed ten-component characteristic potential range with exact V53
      energy and no one-sided inverse;
\item a local algebraic harmonic readout and its reducing zero-constraint
      shell;
\item the exact V36 GHY--de Donder conormal and canonical splitting in full
      characteristic variables;
\item the DeWitt-inertia obstruction to a block-preserving positive port on
      all ten metric components;
\item a sharp one-derivative obstruction for the same-energy differential
      $Q_4$ phase map;
\item a frequency-free $84\to134\to50$ split exact jet replacement, with a
      bounded triangular inverse;
\item the prolonged contractible BRST complex and BRST tangency of its
      holonomic subspace;
\item an explicit nilpotent kinematic BFV moment map compatible with the V36
      Green pairing; and
\item proof that the transverse-gradient and cross outgoing variables added
      in V68 are invariant physical phase components, not disposable jet
      doublets or extra boundary storage.
\end{enumerate}

The full forward--adjoint multiplier action, unbounded ghost/antighost
operator domain, dynamical BFV corner charge, nonlinear curvature and
Brown--York/Bianchi stress identity, and coupled Standard-Model continuum
remain open.  The full Standard Model--Einstein continuum is not proved.

\subsection{Append-only record}
\label{subsec:V69-append-record}

V69 was appended to the complete V68 source.  The SHA-256 digest of the
exact parent file was
\[
 \texttt{5974f22fa90a18901e604dec2c330348ac0c41b4029e5a127304077fadee4f0d}.
\]
The next research note is V70 and begins with the mandatory Yang--Mills and
Navier--Stokes companion audit, followed by the holonomy-multiplier action
and dynamical-domain calculation.

% END APPENDED VERSION V69 MATERIAL

% BEGIN APPENDED VERSION V70 MATERIAL

\clearpage
\section{V70: companion-audited holonomy action, a closed KKT graph,
and the action-derived boundary BFV polarization}
\label{sec:V70-holonomy-action}

\subsection{Mandatory V70 companion audit}
\label{subsec:V70-companion-audit}

The shared-folder audit was performed before choosing the V70 construction.
The exact read-only snapshots were
\begin{center}
\begin{tabular}{@{}lll@{}}
\toprule
programme&file status&SHA--256\\
\midrule
Yang--Mills&
\(\texttt{Renewal\_Yang\_Mills\_Mass\_Gap\_Research\_Notes\_V53.tex}\)&
\scriptsize\ttfamily
6c6b0fa56b227a9d21ef12b1d43aecbcd33ad3b1da2d8f2f3850f6d4cb458332\\
Navier--Stokes&
\(\texttt{Renewal\_NS\_Stability\_Research\_Notes\_V31.tex}\)&
\scriptsize\ttfamily
f1121b62238ad5491bb7cf03635ea9934942edf573ed8faaa9ee79e1d38a6696\\
\bottomrule
\end{tabular}
\end{center}
The Yang--Mills filename advanced from V52 to V53 while the audit was in
progress and was refreshed again before the V70 construction was fixed.
The final audit snapshot has \(41761\) lines.  Its V53 material proves a
universal one-sided transform gap, closes the four quartic star contributions,
then uses degree-adapted routing to close all twelve path contributions and the
complete exclusive-pair quartic marked-tree inequality.  Its final firewall
keeps triple-active, four-active, and embedded four-block re-entry outside that
theorem, so no full quartic or mass-gap conclusion is inferred.  The snapshot
closes with an explicit V53 claim firewall and append-only record.
The Navier--Stokes V31 snapshot has \(23052\) lines.

The retained Yang--Mills V52--V53 material first forms one aggregate invariant
output projector and only then takes a trace norm.  Its exact cost is the
total carrier dimension, rather than a sum of one-copy costs.  It also keeps
the auxiliary defect and its output first moment assembled until the final
estimate.  V53 additionally fixes the global dimension orientation before
using the local transforms.  Its exclusive-pair quartic row is closed, while
the higher-incidence re-entry and mass-gap gates remain open in the audited
snapshot.

Navier--Stokes V31 proves that flat annular marking is removed by an exact
first moment and that a weaker source norm merely restores the missing scale
weight.  Its useful remaining route keeps recent-entry times, signed triad
work, and active parent tails correlated before taking norms.  The
global-regularity theorem remains explicitly open.

\begin{theorem}[Type-correct V70 transfer decision]
\label{thm:V70-companion-transfer}
No Yang--Mills representation estimate and no Navier--Stokes formation
estimate is a metric boundary or BRST theorem.  Their audit changes the V70
construction only through the following incidence rules:
\begin{enumerate}
\item form the complete invariant metric-jet defect before assigning
      separate multiplier rows;
\item keep the two-stage holonomy chain and its formal-adjoint multipliers in
      one quadratic action until the boundary Green form is computed;
\item fix one global face orientation and derive all component adjoint
      conditions from it, rather than choosing signs component by component;
      and
\item reduce by the contractible BRST sector only after that assembled
      boundary form is shown to be Lagrangian.
\end{enumerate}
In particular a multiplier must not be introduced for a metric-jet defect
which is already a bounded linear combination of the invariant wave-jet
defect and the second embedding holonomy.
\end{theorem}

\begin{proof}
The four statements are decisions about order of assembly.  The first is
the direct analogue of forming the complete Yang--Mills invariant projector
before compression.  The second and third retain the exact incidence
correlation and global orientation, as the peer audits require before a weak
norm or local split is taken.  Neither step imports a cross-programme
inequality.  The final
assertion is proved algebraically in
Proposition~\ref{prop:V70-independent-defects}.
\end{proof}

\begin{firewall}[Strict companion type boundary]
\label{fire:V70-companion-types}
Carrier dimensions, Schur taxes, annular first moments, and Duhamel source
weights do not occur in the V70 metric action.  The peer files provide
counterexample-aware assembly discipline only.  Their open mass-gap and
regularity gates remain in their own programmes.
\end{firewall}

\subsection{The independent defect bank}
\label{subsec:V70-independent-bank}

Retain the V69 jet coordinates
\[
 H=h+\Sigma J,\qquad W=P+\Sigma K.
\]
Besides the two gauge holonomies define the invariant and representative
metric-jet defects
\begin{align}
 F^{(1)}&:=J-DX,
 \label{eq:V70-F1}\\
 F^{(2)}&:=K-DJ,
 \label{eq:V70-F2}\\
 F^H&:=W-DH,
 \label{eq:V70-FH}\\
 F^h&:=P-Dh.
 \label{eq:V70-Fh}
\end{align}
Here \(D\) is the complete spacetime first derivative with all tensor
indices retained.  The symmetrization \(\Sigma\) is the frequency-free map
of V69.

\begin{proposition}[Exact independent-defect identity]
\label{prop:V70-independent-defects}
Off shell, without a field equation or a holonomy assumption,
\begin{equation}
 \boxed{\quad F^H=F^h+\Sigma F^{(2)}.\quad}
 \label{eq:V70-defect-identity}
\end{equation}
Consequently the independent defining rows are
\begin{equation}
 F^H=0,\qquad F^{(1)}=0,\qquad F^{(2)}=0.
 \label{eq:V70-independent-rows}
\end{equation}
They imply \(F^h=0\).  A separate multiplier for \(F^h\) would duplicate a
row and produce a reducible antighost incidence.
\end{proposition}

\begin{proof}
Use \(H=h+\Sigma J\) and \(W=P+\Sigma K\):
\[
 W-DH
 =P-Dh+\Sigma(K-DJ)
 =F^h+\Sigma F^{(2)}.
\]
This is \eqref{eq:V70-defect-identity}; rearrangement proves the remaining
claims.
\end{proof}

The invariant row \(F^H=0\) is the ten-component wave-jet definition already
handled by V34, V53, and V69.  V70 therefore constructs one action only for
the genuinely independent gauge chain \(F^{(1)},F^{(2)}\).

Introduce ghost jets \(c,c^{(1)},c^{(2)}\) and put
\begin{align}
 R^{(1)}&:=c^{(1)}-Dc,
 \label{eq:V70-R1}\\
 R^{(2)}&:=c^{(2)}-Dc^{(1)}.
 \label{eq:V70-R2}
\end{align}
The prolonged V69 differential gives
\begin{equation}
 sF^{(1)}=-R^{(1)},\qquad
 sF^{(2)}=-R^{(2)},\qquad
 sF^H=0,\qquad sR^{(1)}=sR^{(2)}=0.
 \label{eq:V70-defect-BRST}
\end{equation}

\subsection{One BRST-exact action for both holonomy layers}
\label{subsec:V70-holonomy-action}

Let \(M_T=[0,T]\times\Omega\) be a compactly truncated slab; fields may
first be taken smooth, with the Hilbert closures specified below.  Introduce
odd antighost tensors \(\eta^{(1)},\eta^{(2)}\) of the same types as
\(F^{(1)},F^{(2)}\), and even multiplier tensors
\(L^{(1)},L^{(2)}\).  Set
\begin{equation}
 s\eta^{(r)}=L^{(r)},\qquad sL^{(r)}=0,
 \qquad r=1,2.
 \label{eq:V70-nonminimal-doublets}
\end{equation}
With the real fibre contractions understood, define
\begin{equation}
 \Psi_{\rm hol}
 :=\int_{M_T}
 \left(\langle\eta^{(1)},F^{(1)}\rangle
       +\langle\eta^{(2)},F^{(2)}\rangle\right)
 \label{eq:V70-holonomy-gauge-fermion}
\end{equation}
and
\begin{align}
 S_{\rm hol}
 :=s\Psi_{\rm hol}
 =\int_{M_T}\bigl[
 &\langle L^{(1)},J-DX\rangle
 +\langle L^{(2)},K-DJ\rangle
 \nonumber\\
 &+\langle\eta^{(1)},c^{(1)}-Dc\rangle
 +\langle\eta^{(2)},c^{(2)}-Dc^{(1)}\rangle
 \bigr].
 \label{eq:V70-holonomy-action}
\end{align}

\begin{theorem}[Off-shell invariance and exact Euler content]
\label{thm:V70-holonomy-Euler}
The action \eqref{eq:V70-holonomy-action} is local, quadratic, and
BRST invariant off shell:
\begin{equation}
 sS_{\rm hol}=0.
 \label{eq:V70-action-invariance}
\end{equation}
Its multiplier and antighost Euler equations are precisely
\begin{equation}
 F^{(1)}=F^{(2)}=0,\qquad
 R^{(1)}=R^{(2)}=0.
 \label{eq:V70-holonomy-Euler-constraints}
\end{equation}
The remaining algebraic top-jet equations imply, successively,
\begin{equation}
 L^{(2)}=L^{(1)}=0,\qquad
 \eta^{(2)}=\eta^{(1)}=0.
 \label{eq:V70-adjoint-zero}
\end{equation}
Thus the holonomy action carries no additional on-shell bulk mode and no
independent positive storage.
\end{theorem}

\begin{proof}
The odd Leibniz rule and
\eqref{eq:V70-defect-BRST} give
\[
 s\langle L^{(r)},F^{(r)}\rangle
 =-\langle L^{(r)},R^{(r)}\rangle,
 \qquad
 s\langle\eta^{(r)},R^{(r)}\rangle
 =\langle L^{(r)},R^{(r)}\rangle.
\]
The terms cancel for each \(r\), proving
\eqref{eq:V70-action-invariance}.  Variation of \(L^{(r)}\) and
\(\eta^{(r)}\) gives
\eqref{eq:V70-holonomy-Euler-constraints}.

Variation of \(K\) gives \(L^{(2)}=0\).  The \(J\)-equation is
\(L^{(1)}-D^\dagger L^{(2)}=0\), hence \(L^{(1)}=0\).
The remaining \(X\)-equation is then automatic.  The odd equations have the
same triangular form: variation of \(c^{(2)}\) gives
\(\eta^{(2)}=0\), and variation of \(c^{(1)}\) then gives
\(\eta^{(1)}-D^\dagger\eta^{(2)}=0\).  This proves
\eqref{eq:V70-adjoint-zero}.  No norm or equation of motion was used to
discard a multiplier.
\end{proof}

\subsection{Every spatial and temporal boundary term}
\label{subsec:V70-boundary-form}

Let \(N_\alpha\) be the outward conormal to an arbitrary face of the slab.
Use left variations in the odd sector.  Integration by parts in
\eqref{eq:V70-holonomy-action} gives the boundary one-form
\begin{align}
 \Theta_{\rm hol}
 =-\int_{\partial M_T}\bigl[
 &(N\mathbin{\lrcorner}L^{(1)})^\nu\,\delta X_\nu
 +(N\mathbin{\lrcorner}L^{(2)})^{\alpha\nu}
                                  \,\delta J_{\alpha\nu}
 \nonumber\\
 &+(N\mathbin{\lrcorner}\eta^{(1)})^\nu\,\delta c_\nu
 +(N\mathbin{\lrcorner}\eta^{(2)})^{\alpha\nu}
                                  \,\delta c^{(1)}_{\alpha\nu}
 \bigr].
 \label{eq:V70-holonomy-boundary-one-form}
\end{align}
There is no independent \(K\) or \(c^{(2)}\) boundary coordinate because
the action contains no derivative of either top jet.

Define the face momenta
\begin{align}
 \pi_X&:=-N\mathbin{\lrcorner}L^{(1)},&
 \pi_J&:=-N\mathbin{\lrcorner}L^{(2)},
 \label{eq:V70-even-face-momenta}\\
 \pi_c&:=-N\mathbin{\lrcorner}\eta^{(1)},&
 \pi_{c^{(1)}}&:=-N\mathbin{\lrcorner}\eta^{(2)}.
 \label{eq:V70-odd-face-momenta}
\end{align}
Then
\begin{equation}
 \Theta_{\rm hol}
 =\int_{\partial M_T}
 \left(\langle\pi_X,\delta X\rangle
       +\langle\pi_J,\delta J\rangle
       +\langle\pi_c,\delta c\rangle
       +\langle\pi_{c^{(1)}},\delta c^{(1)}\rangle\right),
 \label{eq:V70-canonical-boundary-one-form}
\end{equation}
with the graded order fixed as displayed.

\begin{proposition}[Complete variational polarization]
\label{prop:V70-boundary-polarization}
On the spatial face choose
\begin{equation}
 \pi_X=\pi_J=\pi_c=\pi_{c^{(1)}}=0.
 \label{eq:V70-spatial-multiplier-domain}
\end{equation}
At the initial and terminal time faces choose respectively
\begin{align}
 X=J=c=c^{(1)}&=0 &&\text{at }t=0,
 \label{eq:V70-initial-coordinate-domain}\\
 \pi_X=\pi_J=\pi_c=\pi_{c^{(1)}}&=0 &&\text{at }t=T.
 \label{eq:V70-terminal-momentum-domain}
\end{align}
Each is a maximal Lagrangian polarization of the corresponding graded
boundary phase.  All terms in
\eqref{eq:V70-holonomy-boundary-one-form} vanish for allowed variations,
and the three faces meet without an uncancelled corner pairing.
\end{proposition}

\begin{proof}
The exterior derivative of
\eqref{eq:V70-canonical-boundary-one-form} is the direct sum of canonical
coordinate--momentum symplectic pairs.  Setting all momenta to zero, or all
coordinates to zero, gives an isotropic half-dimensional subspace and hence
a Lagrangian subspace.  Equation
\eqref{eq:V70-spatial-multiplier-domain} kills the spatial term.
At \(t=0\) the coordinate variations vanish, and at \(t=T\) the momenta
vanish.  Their restrictions to a common corner kill both possible
polarizations there.
\end{proof}

\subsection{Closedness and maximality of the holonomy KKT graph}
\label{subsec:V70-closed-KKT}

The variational calculation becomes an operator statement as follows.
Let \(\mathscr X\) be the \(L^2\) space of triples
\(Q=(X,J,K)\), and let \(\mathscr Y\) be the \(L^2\) space of pairs with
the types of \((F^{(1)},F^{(2)})\).  Define
\begin{equation}
 \mathcal D_2(X,J,K):=(J-DX,\ K-DJ).
 \label{eq:V70-chain-operator}
\end{equation}
Its domain consists of \(X,J\in H^1(M_T)\), \(K\in L^2(M_T)\), with the
zero initial traces in
\eqref{eq:V70-initial-coordinate-domain}; no coordinate trace is fixed on
the terminal or spatial faces.

\begin{lemma}[The two-stage defect operator is closed]
\label{lem:V70-D2-closed}
The operator
\(\mathcal D_2:D(\mathcal D_2)\subset\mathscr X\to\mathscr Y\)
is densely defined and closed.  Its Hilbert adjoint has the formal interior
expression
\begin{equation}
 \mathcal D_2^*(L^{(1)},L^{(2)})
 =\left(-D^\dagger L^{(1)},
        L^{(1)}-D^\dagger L^{(2)},L^{(2)}\right),
 \label{eq:V70-D2-adjoint}
\end{equation}
and its domain is exactly
\begin{multline}
 D(\mathcal D_2^*)
 =\bigl\{(L^{(1)},L^{(2)})\in\mathscr Y:
 D^\dagger L^{(1)},D^\dagger L^{(2)}\in L^2,\\
 N\mathbin{\lrcorner}L^{(1)}
 =N\mathbin{\lrcorner}L^{(2)}=0
 \text{ on the terminal and spatial faces}\bigr\}.
 \label{eq:V70-adjoint-face-domain}
\end{multline}
The normal traces in this formula are the weak traces of the corresponding
\(H(\operatorname{div})\) tensors.  There is no multiplier restriction at
the initial face.
\end{lemma}

\begin{proof}
Smooth triples satisfying the initial trace condition are dense.  Suppose
\(Q_m\to Q\) in \(\mathscr X\) and
\(\mathcal D_2Q_m\to(F_1,F_2)\) in \(\mathscr Y\).  Then
\[
 DX_m=J_m-(\mathcal D_2Q_m)_1\longrightarrow J-F_1,
 \qquad
 DJ_m=K_m-(\mathcal D_2Q_m)_2\longrightarrow K-F_2.
\]
Closedness of the weak derivative gives \(X,J\in H^1\), the same initial
traces, and \(\mathcal D_2Q=(F_1,F_2)\).  Hence \(\mathcal D_2\) is closed.

For smooth fields, integration by parts gives
\begin{align*}
 \langle\mathcal D_2Q,L\rangle
 ={}&\langle X,-D^\dagger L^{(1)}\rangle
 +\langle J,L^{(1)}-D^\dagger L^{(2)}\rangle
 +\langle K,L^{(2)}\rangle\\
 &-\int_{\partial M_T}
 \left(\langle N\mathbin{\lrcorner}L^{(1)},X\rangle
       +\langle N\mathbin{\lrcorner}L^{(2)},J\rangle\right).
\end{align*}
The initial coordinate traces vanish, while terminal and spatial traces are
arbitrary.  Their annihilator is exactly
\eqref{eq:V70-adjoint-face-domain}.  This proves the adjoint formula and
domain.
\end{proof}

Define the KKT Hessian
\begin{equation}
 \mathcal K_{70}
 :=\begin{pmatrix}0&\mathcal D_2^*\\
                   \mathcal D_2&0\end{pmatrix},
 \qquad
 D(\mathcal K_{70})
 :=D(\mathcal D_2)\oplus D(\mathcal D_2^*).
 \label{eq:V70-KKT-operator}
\end{equation}

\begin{theorem}[Self-adjoint maximal holonomy Hessian]
\label{thm:V70-KKT-selfadjoint}
The operator \(\mathcal K_{70}\) is self-adjoint and closed on
\(\mathscr X\oplus\mathscr Y\).  It is the Hessian operator of the even part
of \(S_{\rm hol}\) with precisely the initial, terminal, and spatial
polarizations
\eqref{eq:V70-spatial-multiplier-domain}--%
\eqref{eq:V70-terminal-momentum-domain}.  The commuting-coefficient
representative of the odd block has the same closed graph and complementary
adjoint domain.
\end{theorem}

\begin{proof}
Let \(D=\mathcal D_2\).  For every closed densely defined operator,
\[
 \begin{pmatrix}0&D^*\\D&0\end{pmatrix}^*
 =
 \begin{pmatrix}0&D^*\\D&0\end{pmatrix}
\]
on \(D(D)\oplus D(D^*)\).  Indeed, testing first against
\((x,0)\) forces the second adjoint component into \(D(D^*)\), and testing
against \((0,y)\) forces the first into \(D(D)\); the two resulting actions
are \(D^*\) and \(D\).  Thus no proper symmetric extension exists.
Lemma~\ref{lem:V70-D2-closed} supplies the required hypotheses and identifies
the boundary domains.  The quadratic polarization of
\(\langle L,\mathcal D_2Q\rangle\) has exactly this off-diagonal Hessian.
The odd statement concerns the analytic graph only; it does not turn
Grassmann parity into a positive Hilbert norm.
\end{proof}

\begin{corollary}[No hidden multiplier trace]
\label{cor:V70-no-hidden-trace}
The spatial multiplier-zero relation
\eqref{eq:V70-spatial-multiplier-domain} is not an ad hoc deletion.  It is
the exact adjoint annihilator of the free gauge-coordinate trace in
\(D(\mathcal D_2)\).  Conversely, imposing a gauge-coordinate Dirichlet
trace would select the complementary maximal adjoint domain.
\end{corollary}

\subsection{The action-derived boundary BFV charge}
\label{subsec:V70-boundary-BFV}

The BRST differential induces on every boundary face
\begin{align}
 sX&=-c,&sJ&=-c^{(1)},&sc&=sc^{(1)}=0,
 \label{eq:V70-boundary-BRST-coordinates}\\
 s\pi_X&=s\pi_J=0,&
 s\pi_c&=\pi_X,&s\pi_{c^{(1)}}&=\pi_J.
 \label{eq:V70-boundary-BRST-momenta}
\end{align}
The momentum identities follow directly from
\(s\eta^{(r)}=L^{(r)}\) and
\eqref{eq:V70-even-face-momenta}--%
\eqref{eq:V70-odd-face-momenta}.

Let \(\omega_{\partial,70}=\delta\Theta_{\rm hol}\) be the canonical graded
boundary symplectic form.

\begin{theorem}[Dynamical holonomy BFV generator]
\label{thm:V70-BFV-generator}
The action-derived boundary charge is
\begin{equation}
 \boxed{\quad
 \Omega_{\partial,70}^{\rm hol}
 :=-\int_{\partial M_T}
 \left(\langle\pi_X,c\rangle
       +\langle\pi_J,c^{(1)}\rangle\right).
 \quad}
 \label{eq:V70-boundary-BFV-charge}
\end{equation}
With the canonical graded bracket,
\begin{equation}
 \iota_{Q_{\partial,70}}\omega_{\partial,70}
 =\delta\Omega_{\partial,70}^{\rm hol},
 \qquad
 \{\Omega_{\partial,70}^{\rm hol},
   \Omega_{\partial,70}^{\rm hol}\}=0.
 \label{eq:V70-boundary-master-equation}
\end{equation}
It generates all transformations in
\eqref{eq:V70-boundary-BRST-coordinates}--%
\eqref{eq:V70-boundary-BRST-momenta}.
The spatial momentum polarization, the initial coordinate polarization, and
the terminal momentum polarization are BRST invariant.
\end{theorem}

\begin{proof}
The derivatives of \eqref{eq:V70-boundary-BFV-charge} with respect to
\(\pi_X,\pi_J\) give \(-c,-c^{(1)}\), while its graded derivatives with
respect to \(c,c^{(1)}\) give the momentum transformations in
\eqref{eq:V70-boundary-BRST-momenta}.  It is independent of the remaining
coordinates and momenta.  This proves the Hamiltonian identity.
The charge is linear in mutually commuting momenta and contains no conjugate
coordinate, so its self-bracket vanishes directly.

If all momenta vanish, their BRST images vanish because
\(s\pi_c=\pi_X\) and \(s\pi_{c^{(1)}}=\pi_J\).  If all initial coordinates
vanish, \(sX=-c\) and \(sJ=-c^{(1)}\) also vanish there.  This proves
invariance of all three polarizations.
\end{proof}

\begin{proposition}[Identification with the V69 kinematic charge]
\label{prop:V70-V69-charge-identification}
On the holonomic boundary phase, the V69 cotangent moment map reduces to
\eqref{eq:V70-boundary-BFV-charge}.  The top jets \(K,c^{(2)}\) are normal
derivatives or algebraic variables rather than independent boundary
coordinates, so their primary boundary momenta vanish.  The transformations
of the ghost momenta encode
\(s\eta^{(r)}=L^{(r)}\); no separate antighost momentum term is missing.
\end{proposition}

\begin{proof}
The boundary one-form
\eqref{eq:V70-canonical-boundary-one-form} identifies the only nonzero
canonical momenta.  Substitute these momenta in the component expression for
the V69 cotangent moment map.  The \(X,J\) terms give
\eqref{eq:V70-boundary-BFV-charge}; the \(K,c^{(2)}\) terms vanish on their
primary momentum constraints.  Finally
\[
 s\pi_c=-N\mathbin{\lrcorner}s\eta^{(1)}
       =-N\mathbin{\lrcorner}L^{(1)}=\pi_X
\]
and similarly at the second level.  These are exactly the momentum
Hamilton equations of the displayed charge.
\end{proof}

\subsection{Decoupling from the invariant metric Green form}
\label{subsec:V70-decoupling}

The triangular V69 coordinates split the metric-embedding jet into the
invariant variables \((H,W)\) and the gauge variables \((X,J,K)\).
The V36 boundary one-form depends only on \((H,W)\), while
\(\Theta_{\rm hol}\) is the canonical form of the gauge chain.

\begin{theorem}[Maximal product relation and zero cohomological power]
\label{thm:V70-product-relation}
Let \(\mathcal L_{\rm inv}\) be any Lagrangian boundary relation expressed
only in the invariant V36 metric variables and any invariant continuum
variables.  Then
\begin{equation}
 \mathcal L_{\rm inv}\times
 \{\pi_X=\pi_J=\pi_c=\pi_{c^{(1)}}=0\}
 \label{eq:V70-product-boundary-relation}
\end{equation}
is Lagrangian for the direct-sum boundary Green form.  It is BRST invariant.
On the Euler shell of \(S_{\rm hol}\), its holonomy Green form and BFV charge
both vanish.  Therefore adjoining the complete \(Q_4\) holonomy sector does
not change the positive V68/V59 energy on the radiative cohomology.
\end{theorem}

\begin{proof}
The exterior derivative of the total one-form is the direct sum of the
invariant Green form and \(\omega_{\partial,70}\).  The first factor is
Lagrangian by hypothesis; the second is Lagrangian by
Proposition~\ref{prop:V70-boundary-polarization}.  Their product is therefore
isotropic and half-dimensional.  BRST acts trivially on the invariant
variables and preserves the momentum-zero factor by
Theorem~\ref{thm:V70-BFV-generator}.

Theorem~\ref{thm:V70-holonomy-Euler} gives
\(L^{(r)}=\eta^{(r)}=0\) on the holonomy Euler shell, hence all momenta in
\eqref{eq:V70-even-face-momenta}--%
\eqref{eq:V70-odd-face-momenta} vanish.  Equations
\eqref{eq:V70-canonical-boundary-one-form} and
\eqref{eq:V70-boundary-BFV-charge} then vanish.  The remaining positive
energy is precisely the invariant radiative--continuum energy already proved
in V68.
\end{proof}

\begin{firewall}[What the product theorem assumes]
\label{fire:V70-product-scope}
Theorem~\ref{thm:V70-product-relation} removes the holonomy sector as an
obstruction once an invariant physical Lagrangian relation is supplied.  It
does not construct the missing local spacetime graph from the V68
two-polarization phase to the full \((H,W)\) metric phase.  Applying it to
the frozen V65 conormal relation is valid fibrewise; promoting that relation
to a local curved action remains a separate theorem.
\end{firewall}

\subsection{Why the KKT theorem is not a time-generator theorem}
\label{subsec:V70-KKT-firewall}

\begin{firewall}[Spacetime self-adjointness versus positive evolution]
\label{fire:V70-not-generator}
Theorem~\ref{thm:V70-KKT-selfadjoint} proves maximality of the spacetime
variational Hessian with complementary face domains.  The KKT form is
indefinite and the top jets are algebraic.  This does not make
\(\mathcal K_{70}\) a positive Hamiltonian or a skew-adjoint time generator.
No such generator is needed for a contractible constraint sector: its
physical Euler shell has zero multiplier state.  The positive evolution
still comes from the invariant V68 wave--continuum generator.  A full action
claim requires that generator's local metric graph and V36 counterterm to be
placed in the same spacetime functional.
\end{firewall}

\subsection{Finite algebraic certificates}
\label{subsec:V70-certificates}

For a finite matrix \(D\), the two-stage defect representative is
\begin{equation}
 \mathsf D_2=
 \begin{pmatrix}
 -D&I&0\\
 0&-D&I
 \end{pmatrix}.
 \label{eq:V70-finite-chain}
\end{equation}
Its KKT matrix and boundary charge matrix are
\begin{equation}
 \mathsf K_{70}=
 \begin{pmatrix}0&\mathsf D_2^*\\
                 \mathsf D_2&0\end{pmatrix},
 \qquad
 \mathsf Q_{\partial,70}^2=0.
 \label{eq:V70-finite-KKT}
\end{equation}
The exact cards are
\begin{align}
 F^H-F^h-\Sigma F^{(2)}&=0,
 \label{eq:V70-defect-card}\\
 sS_{\rm hol}&=0,
 \label{eq:V70-action-card}\\
 \mathsf K_{70}^*&=\mathsf K_{70},
 \label{eq:V70-KKT-card}\\
 \omega_{\partial,70}|_{\pi=0}&=0,
 \qquad
 \dim\{\pi=0\}=\frac12\dim\mathcal P_{\partial,70},
 \label{eq:V70-Lagrangian-card}\\
 \{\Omega_{\partial,70}^{\rm hol},
   \Omega_{\partial,70}^{\rm hol}\}&=0.
 \label{eq:V70-BFV-card}
\end{align}
A finite verifier accompanying this note checks
\eqref{eq:V70-defect-card} in random tensor jets, constructs
\(\mathsf D_2\) for random derivative matrices, verifies the KKT symmetry and
rank, verifies the two boundary Lagrangian polarizations, and checks the
nilpotent boundary BRST matrix and its Hamiltonian realization.  The
operator proofs, rather than the finite samples, establish closedness and
maximality.

\subsection{V71 acquisition order}
\label{subsec:V70-next}

No companion audit is due in V71; the next mandatory audit is V75.  V71
should preserve the independent holonomy action and proceed in this order:
\begin{enumerate}
\item write the invariant first-order parent action for \((H,W)\) whose
      elimination is the V36 de Donder wave action;
\item compute the unique quadratic boundary counterterm taking its minimal
      DeWitt wave one-form to the exact GHY--de Donder one-form
      \eqref{eq:V69-V36-pulled-one-form};
\item intersect its boundary phase with the harmonic reducing shell and the
      V68 positive radiative phase, keeping the V53 residual-gauge quotient
      explicit;
\item construct a local differential radiative-to-metric graph, or prove the
      first exact obstruction to such a graph; a pseudodifferential TT
      projector must not be relabelled local;
\item combine the resulting invariant relation with
      \eqref{eq:V70-product-boundary-relation} and the V59 continuum
      interface, then prove the full linear action's corner master equation;
\item only after this linear metric gate closes, restore curvature and the
      nonlinear Brown--York/Bianchi stress identity before adding the Higgs,
      fermion, neutrino, and Yang--Mills source blocks.
\end{enumerate}
If the local radiative graph fails, V71 must retain the full harmonic metric
phase and perform the residual-gauge reduction at the operator level rather
than returning to the half-angle multiplier.

\subsection{V70 result ledger and claim boundary}
\label{subsec:V70-ledger}

Relative to V1--V69, this note proves:
\begin{enumerate}
\item the mandatory, digest-recorded audit of the current Yang--Mills and
      Navier--Stokes peer files, including the mid-audit YM version change;
\item the exact dependency
      \(F^H=F^h+\Sigma F^{(2)}\), eliminating a redundant multiplier and
      antighost row;
\item one local quadratic BRST-exact action for both independent
      differential \(Q_4\) holonomy layers;
\item all spatial, initial, terminal, and corner terms of that action, with
      complementary Lagrangian polarizations;
\item closedness of the two-stage defect operator and self-adjoint maximality
      of its infinite-dimensional KKT Hessian;
\item the action-derived boundary BFV symplectic form and nilpotent charge;
\item exact identification of that charge with the appropriate reduction of
      the V69 kinematic cotangent charge; and
\item vanishing of the holonomy Green power and BFV charge on cohomology, so
      the contractible sector does not alter the V68/V59 positive energy.
\end{enumerate}

The invariant first-order metric parent action, its exact V36 boundary
counterterm, a local radiative-to-full-metric graph, the complete linear
metric--continuum corner master equation, nonlinear curvature and
Brown--York/Bianchi stress closure, and the coupled Standard-Model continuum
remain open.  The full Standard Model--Einstein continuum is not proved.

\subsection{Append-only record}
\label{subsec:V70-append-record}

V70 was appended to the complete V69 source.  The SHA--256 digest of the
exact parent file was
\[
 \texttt{72c785b8579932aaed0467c309ffe4ea5d6247fab9fac71ff645ace4be63880a}.
\]
The next research note is V71 and begins with the invariant first-order
metric parent action and the exact V36 boundary-counterterm calculation.

% END APPENDED VERSION V70 MATERIAL

% BEGIN APPENDED VERSION V71 MATERIAL

\clearpage
\section{V71: the invariant minimal parent, the GHY ancestry obstruction,
and the no-local-splitting theorem}
\label{sec:V71-parent-ancestry}

\subsection{Purpose, nonduplication boundary, and the upstream turn}
\label{subsec:V71-purpose}

V70 closed the independent differential \(Q_4\) holonomy sector by one
BRST-exact action and a maximal KKT graph.  It deliberately left the invariant
metric pair \((H,W)\) untouched.  The next proposed step was to start from a
minimal first-order DeWitt parent and add a quadratic boundary counterterm
whose variation is the exact V36 GHY--de Donder conormal.

This note tests that proposal before using it.  There are three conclusions.
First, the invariant minimal parent exists and eliminates exactly to the V36
minimal de Donder density.  Second, after that minimalization there is no local
quadratic boundary functional of the metric configuration and its tangential
jets which converts the parent boundary one-form into the V36 one-form.  The
normal Legendre matrices already disagree.  This is not a failure of the GHY
term: it says that the second-derivative Einstein divergence and the GHY term
must be retained together until their normal-jet cancellation has occurred.
Third, the missing map from two fixed radiative scalar amplitudes to a full
harmonic metric cannot be a finite-order local differential right inverse at
the V59/V68 energy order.  The correct endpoint is therefore the full local
harmonic complex followed by operator-level residual-gauge reduction.

The note also gives an exact local repair of the first obstruction on an
enlarged boundary jet cotangent bundle.  That repair identifies the boundary
variables which a covariant Einstein--GHY parent must carry in V72.  It is a
variational incidence theorem, not a positive-generator theorem.

No companion audit is due in V71.  The next mandatory read-only Yang--Mills
and Navier--Stokes audit remains V75.

\subsection{The invariant first-order minimal parent}
\label{subsec:V71-minimal-parent}

Work at the same planar flat anchor and with the same outward spacelike normal
as V36.  Write
\begin{equation}
 \Gamma_T:=[0,T]\times\partial\Omega
 \label{eq:V71-spatial-face}
\end{equation}
for the timelike spatial face; the initial and terminal faces are kept separate.
For symmetric tensors \(A,B\), retain
\begin{equation}
 \mathbb G(A,B)
 :=A_{\mu\nu}B^{\mu\nu}
   -\frac12(\operatorname{tr}_\eta A)
                (\operatorname{tr}_\eta B)
 =\overline A^{\mu\nu}B_{\mu\nu}.
 \label{eq:V71-DeWitt-form}
\end{equation}
For one-form-valued symmetric tensors put
\begin{equation}
 \mathbb G_4(U,V)
 :=\eta^{\alpha\beta}\mathbb G(U_\alpha,V_\beta).
 \label{eq:V71-spacetime-DeWitt-form}
\end{equation}
The V69 variables \(H\) and \(W\) are BRST invariant.  In the parent action
\(W\) is initially independent of \(DH\):
\begin{equation}
 S_{71}^{\rm min}[H,W]
 :=\frac1{2\kappa}\int_{M_T}
 \left\{
  \frac12\mathbb G_4(W,W)-\mathbb G_4(W,DH)
 \right\}.
 \label{eq:V71-minimal-parent-action}
\end{equation}

\begin{theorem}[Exact parent elimination and invariant Euler rows]
\label{thm:V71-parent-elimination}
The independent Euler equations of \eqref{eq:V71-minimal-parent-action} are
\begin{equation}
 W=DH,
 \qquad
 D_\alpha\overline W^{\alpha\mu\nu}=0.
 \label{eq:V71-parent-Euler}
\end{equation}
Eliminating \(W\) gives
\begin{equation}
 S_{71}^{\rm min}[H,DH]
 =-\frac1{4\kappa}\int_{M_T}
   \mathbb G_4(DH,DH),
 \label{eq:V71-eliminated-action}
\end{equation}
which is exactly the V36 density
\(-\tfrac12D_\alpha H_{\mu\nu}D^\alpha\overline H^{\mu\nu}\)
with the common action factor \((2\kappa)^{-1}\).  The reduced Euler equation
is \(\mathcal P\Box H=0\), equivalently \(\Box H=0\).  Off shell,
\begin{equation}
 sS_{71}^{\rm min}=0.
 \label{eq:V71-parent-BRST}
\end{equation}
\end{theorem}

\begin{proof}
Variation of \(W\) gives
\[
 \mathbb G_4(\delta W,W-DH)=0.
\]
The DeWitt form and the spacetime metric are nondegenerate, hence
\(W=DH\).  Substitution leaves one half minus one copy of the same quadratic
form, proving \eqref{eq:V71-eliminated-action}.  Variation of \(H\), followed
by one integration by parts, gives the second row of
\eqref{eq:V71-parent-Euler}.  With \(W=DH\), it is
\(D_\alpha\overline{D^\alpha H}=\mathcal P\Box H=0\).  V36 proves
\(\mathcal P^2=I\), so this is equivalent to the ten wave equations.  Finally
\(sH=sW=0\) in the triangular V69 variables, proving
\eqref{eq:V71-parent-BRST} without use of an Euler equation.
\end{proof}

\begin{proposition}[Complete minimal-parent boundary potential]
\label{prop:V71-minimal-boundary-potential}
For arbitrary smooth variations, the spatial-face term of the parent action is
\begin{equation}
 \Theta_{71}^{\rm min}
 =-\frac1{2\kappa}\int_{\Gamma_T}
       \mathbb G(W_n,\delta H)
 =\frac1{2\kappa}\int_{\Gamma_T}
       (\mathfrak p_{\rm min}(W))^{\mu\nu}\delta H_{\mu\nu},
 \label{eq:V71-minimal-boundary-potential}
\end{equation}
where
\begin{equation}
 \mathfrak p_{\rm min}(W)^{\mu\nu}
 :=-\overline{W_n}^{\mu\nu}.
 \label{eq:V71-minimal-momentum}
\end{equation}
On the algebraic shell \(W=DH\), this is the ordinary minimal DeWitt wave
potential.  No \(\delta W\) boundary term occurs.
\end{proposition}

\begin{proof}
Only the term \(-\mathbb G_4(W,DH)\) contains a derivative of a varied
configuration.  Its integration by parts gives
\(-\mathbb G(W_n,\delta H)\) on the outward face.  The trace-reversal identity
in \eqref{eq:V71-DeWitt-form} gives
\eqref{eq:V71-minimal-momentum}.  Since \(W\) is algebraic in the parent bulk
Lagrangian, variation of \(W\) creates no boundary derivative.
\end{proof}

\subsection{The exact V36 first-jet conormal map}
\label{subsec:V71-V36-jet-map}

Let
\begin{equation}
 r=(r_0,r_1,r_2,r_n)
 \in T^*M|_{\partial M}\otimes\operatorname{Sym}^2T^*M
 \label{eq:V71-boundary-first-jet}
\end{equation}
be an independent first jet.  Replace every \(D_\alpha H\) in the V36
formulas by \(r_\alpha\) and define
\begin{align}
 C_\nu(r)
 &:=\eta^{\alpha\beta}(r_\beta)_{\alpha\nu}
       -\frac12(r_\nu)^\rho{}_{\rho},
 \label{eq:V71-jet-C}\\
 K_{ab}(r)
 &:=\frac12\left[(r_n)_{ab}-(r_a)_{nb}-(r_b)_{na}\right],
 \label{eq:V71-jet-K}\\
 \Pi^{ab}(r)
 &:=K^{ab}(r)-\eta_\partial^{ab}\eta_\partial^{cd}K_{cd}(r),
 \label{eq:V71-jet-Pi}\\
 (\mathcal A_{36}r)^{\mu\nu}
 &:=e_a^\mu e_b^\nu\Pi^{ab}(r)
   -2\left(n^{(\mu}C^{\nu)}(r)
          -\frac12\eta^{\mu\nu}C_n(r)\right).
 \label{eq:V71-A36}
\end{align}
Thus \(\mathcal A_{36}(DH)=\mathfrak p_{36}(H)\) by V36.  The minimal
first-jet conormal map is
\begin{equation}
 \mathcal A_{\rm min}r:=-\overline{r_n},
 \qquad
 \Delta_{71}:=\mathcal A_{36}-\mathcal A_{\rm min}.
 \label{eq:V71-Delta}
\end{equation}

The difference is already visible in the coefficient of the normal jet.  In
the independent symmetric-component order
\begin{equation}
 (00,01,02,0n,11,12,1n,22,2n,nn),
 \label{eq:V71-component-order}
\end{equation}
let \(M_{36}\) be the matrix satisfying
\begin{equation}
 u^TM_{36}v
 =\left\langle\mathcal A_{36}(0,0,0,v),u\right\rangle.
 \label{eq:V71-M36-definition}
\end{equation}
Here off-diagonal coordinates represent both symmetric entries, and the
right side is the full tensor contraction.  Direct substitution gives
\begin{equation}
\scriptsize
 M_{36}=
 \begin{pmatrix}
 -\frac12&0&0&0&1&0&0&1&0&-\frac12\\
 0&-1&0&0&0&0&0&0&0&0\\
 0&0&-1&0&0&0&0&0&0&0\\
 0&0&0&2&0&0&0&0&0&0\\
 1&0&0&0&-\frac12&0&0&-1&0&\frac12\\
 0&0&0&0&0&1&0&0&0&0\\
 0&0&0&0&0&0&-2&0&0&0\\
 1&0&0&0&-1&0&0&-\frac12&0&\frac12\\
 0&0&0&0&0&0&0&0&-2&0\\
 -\frac12&0&0&0&\frac12&0&0&\frac12&0&-\frac12
 \end{pmatrix}.
 \label{eq:V71-M36-matrix}
\end{equation}

\begin{lemma}[Normal Legendre comparison card]
\label{lem:V71-normal-comparison}
The matrix \(M_{36}\) is symmetric,
\begin{equation}
 \det M_{36}=1,
 \qquad
 \operatorname{inertia}(M_{36})=(4,6,0)
 \label{eq:V71-M36-inertia}
\end{equation}
for the displayed outward-normal convention.  Reversing the conormal changes
the sign and interchanges the positive and negative counts.  It is not the
normal matrix of \(\mathcal A_{\rm min}\), nor is it any scalar multiple of
that matrix.
\end{lemma}

\begin{proof}
Symmetry is visible in \eqref{eq:V71-M36-matrix}.  Reorder it as the
direct sum of the \((00,11,22,nn)\) block and the six one-dimensional
blocks carried by \((01,02,0n,12,1n,2n)\).  The leading principal minors of
the four-dimensional block, including the zeroth minor, are
\begin{equation}
 1,\quad-\frac12,\quad-\frac34,\quad-\frac58,\quad\frac18.
 \label{eq:V71-diagonal-principal-minors}
\end{equation}
Hence its exact \(LDL^T\) pivots have signs \((- ,+,+,-)\), so that block
has inertia \((2,2,0)\).  The six scalar blocks are
\((-1,-1,2,1,-2,-2)\), contributing \((2,4,0)\).  Their product is
\(8\), while the four-block determinant is \(1/8\); therefore the full
determinant is one and the inertia is \((4,6,0)\).  For a component proof
of the last
assertion, let \(E_{\mu\nu}\) denote the symmetric tensor with both displayed
off-diagonal entries equal to one.  With all tangential derivatives zero,
\begin{center}
\begin{tabular}{@{}ccl@{}}
\toprule
normal jet and test&minimal contraction&V36 contraction\\
\midrule
\(r_n=u=E_{12}\)&\(-2\)&\(1\)\\
\(r_n=u=E_{0n}\)&\(2\)&\(2\)\\
\bottomrule
\end{tabular}
\end{center}
The first minimal value is \(-\mathbb G(E_{12},E_{12})=-2\); the second is
\(-\mathbb G(E_{0n},E_{0n})=2\).  The V36 values follow respectively from
\(K_{12}=1/2\) and from \(C_0=1\), \(C^0=-1\).  The two ratios are
incompatible, proving both inequality and the absence of a scalar
calibration.
\end{proof}

\subsection{Why the post-minimal GHY counterterm does not exist}
\label{subsec:V71-counterterm-obstruction}

\begin{theorem}[Post-minimal boundary-counterterm obstruction]
\label{thm:V71-no-configuration-counterterm}
There is no local quadratic boundary functional
\begin{equation}
 B[H]=\int_{\Gamma_T}
       b(H,D_aH,\ldots,D_{a_1}\cdots D_{a_m}H)
 \label{eq:V71-boundary-functional-class}
\end{equation}
with finitely many tangential derivatives such that, for arbitrary boundary
first jets,
\begin{equation}
 \Theta_{71}^{\rm min}+\delta B=\Theta_{36}.
 \label{eq:V71-forbidden-counterterm-identity}
\end{equation}
The conclusion remains true if a local dependence on \(D_nH\) is allowed but
the final one-form is required to contain no \(\delta(D_nH)\) term.
\end{theorem}

\begin{proof}
After tangential integrations by parts, \(\delta B\) for a functional in
\eqref{eq:V71-boundary-functional-class} has the form
\[
 \int_{\Gamma_T}\langle E_B(H),\delta H\rangle,
\]
where \(E_B\) contains tangential derivatives only.  It cannot change the
coefficient of the freely variable normal jet \(D_nH\).  Equality
\eqref{eq:V71-forbidden-counterterm-identity} would therefore force equality
of the two normal Legendre matrices.  Lemma
\ref{lem:V71-normal-comparison} disproves that equality.

Now allow \(b\) to depend on \(D_nH\).  Its first variation contains
\[
 \left\langle\frac{\partial b}{\partial(D_nH)},
                   \delta(D_nH)\right\rangle.
\]
Neither side of the desired difference
\(\Theta_{36}-\Theta_{71}^{\rm min}\) contains a normal-jet variation.
Because \(H\) and \(D_nH\) are independent boundary-jet coordinates, the
coefficient must vanish identically.  Thus the functional is independent of
\(D_nH\), up to a term constant in that coordinate, and the first argument
applies.
\end{proof}

\begin{corollary}[The order of the Einstein--GHY cancellation is forced]
\label{cor:V71-order-forced}
The V36 conormal cannot be recovered by first replacing the quadratic
Einstein--gauge-fixed action with the minimal first-derivative density and
then adjoining a metric-only boundary scalar.  A valid local parent must do
one of the following:
\begin{enumerate}
\item retain the second-derivative Einstein divergence and the GHY variation
      in the same parent until their \(\delta(D_nH)\) terms cancel; or
\item enlarge the boundary phase by independent first-jet coordinates and
      their conjugate momenta.
\end{enumerate}
\end{corollary}

\begin{proof}
Theorem~\ref{thm:V71-no-configuration-counterterm} rules out the proposed
post-minimal scalar.  The original Einstein--GHY derivation in V36 performs
the cancellation before passing to its final first-jet conormal, which is the
first alternative.  The second alternative is constructed next.
\end{proof}

\begin{firewall}[No contradiction with the GHY theorem]
\label{fire:V71-GHY-consistency}
The GHY term cancels the normal derivative of a metric variation arising from
the second-derivative Einstein--Hilbert bulk action.  Equation
\eqref{eq:V71-minimal-parent-action} has already discarded that
second-derivative representative and has no \(\delta W\) boundary term.
Theorem~\ref{thm:V71-no-configuration-counterterm} concerns this order of
operations.  It neither withdraws the V36 conormal nor asserts that the
Einstein--GHY action is nonexistent.
\end{firewall}

\subsection{Exact jet-cotangent lift of the V36 potential}
\label{subsec:V71-jet-cotangent-lift}

The obstruction is removed locally by retaining the boundary first jet as an
independent coordinate.  Let \(\varpi^\alpha\) be its cotangent momentum and
define the adjoint blocks \(\Delta_{71,\alpha}^{\vee}\) by
\begin{equation}
 \langle H,\Delta_{71}r\rangle
 =\sum_{\alpha=0,1,2,n}
   \langle\Delta_{71,\alpha}^{\vee}H,r_\alpha\rangle.
 \label{eq:V71-Delta-adjoint}
\end{equation}
The local bilinear generating functional is
\begin{equation}
 F_{71}^{\rm jet}[H,r]
 :=\frac1{2\kappa}\int_{\Gamma_T}
      \langle H,\Delta_{71}r\rangle.
 \label{eq:V71-jet-generator}
\end{equation}

\begin{theorem}[Exact enlarged-phase recovery of the V36 one-form]
\label{thm:V71-jet-lift}
Adjoin to the minimal parent boundary potential the canonical jet term and
the variation of \eqref{eq:V71-jet-generator}:
\begin{equation}
 \Theta_{71}^{\rm ext}
 :=\Theta_{71}^{\rm min}
   +\frac1{2\kappa}\int_{\Gamma_T}
        \langle\varpi^\alpha,\delta r_\alpha\rangle
   +\delta F_{71}^{\rm jet}.
 \label{eq:V71-extended-potential}
\end{equation}
On the algebraic momentum row
\begin{equation}
 \varpi^\alpha=-\Delta_{71,\alpha}^{\vee}H
 \label{eq:V71-jet-momentum-row}
\end{equation}
and the local compatibility rows
\begin{equation}
 r_a=D_aH\quad(a=0,1,2),
 \qquad r_n=W_n,
 \label{eq:V71-jet-compatibility}
\end{equation}
one has the exact identity
\begin{equation}
 \boxed{\quad
 \Theta_{71}^{\rm ext}\big|_{\eqref{eq:V71-jet-momentum-row}
                              ,\eqref{eq:V71-jet-compatibility}}
 =\Theta_{36}(H;\delta H).
 \quad}
 \label{eq:V71-exact-V36-lift}
\end{equation}
The transformation is local and BRST invariant.  In the enlarged cotangent
phase it is an exact symplectic transformation, and
\eqref{eq:V71-jet-momentum-row} is the transformed auxiliary zero-section.
\end{theorem}

\begin{proof}
The generator varies as
\begin{equation}
 2\kappa\,\delta F_{71}^{\rm jet}
 =\int_{\Gamma_T}
   \left(
    \langle\Delta_{71}r,\delta H\rangle
    +\sum_\alpha
      \langle\Delta_{71,\alpha}^{\vee}H,\delta r_\alpha\rangle
   \right).
 \label{eq:V71-generator-variation}
\end{equation}
The second term cancels the jet cotangent term on
\eqref{eq:V71-jet-momentum-row}.  The coefficient of \(\delta H\) becomes
\[
 \mathcal A_{\rm min}r+\Delta_{71}r=\mathcal A_{36}r.
\]
The compatibility rows then turn this into
\(\mathfrak p_{36}(H)\), which proves
\eqref{eq:V71-exact-V36-lift}.

Adding \(\delta F_{71}^{\rm jet}\) changes the canonical momenta by
\begin{equation}
 p_H\longmapsto p_H+\Delta_{71}r,
 \qquad
 \varpi^\alpha\longmapsto
 \varpi^\alpha+\Delta_{71,\alpha}^{\vee}H.
 \label{eq:V71-canonical-shift}
\end{equation}
Because the added one-form is exact, its field-space exterior derivative is
zero; hence the shift is symplectic.  The momentum row is precisely zero of
the shifted \(\varpi\).  Finally \(H,W\) are BRST invariant, and the
auxiliary jet and cotangent variables may be assigned zero BRST variation.
Every displayed row and the generator are then invariant off shell.
\end{proof}

\begin{corollary}[The repair adds no radiative class]
\label{cor:V71-jet-no-cohomology}
The variables \((r,\varpi)\) in Theorem~\ref{thm:V71-jet-lift} are fixed by
local compatibility and algebraic momentum rows.  They add no freely
propagating bulk solution and no new radiative quotient class.  They do add a
boundary/corner cotangent sector whose full temporal-edge polarization must
be derived from the covariant parent action.
\end{corollary}

\begin{proof}
Given the bulk pair \((H,W)\), equation
\eqref{eq:V71-jet-compatibility} fixes \(r\) uniquely, and
\eqref{eq:V71-jet-momentum-row} then fixes \(\varpi\) uniquely.  There is no
homogeneous auxiliary amplitude.  The last sentence records that the
canonical term has boundary-of-boundary data; V71 has not supplied a temporal
kinetic law for it.
\end{proof}

\subsection{The same-energy local splitting obstruction}
\label{subsec:V71-local-splitting-obstruction}

The V53 radiative exact sequence is split by the direction-dependent spatial
TT projector.  That is a bounded Fourier multiplier in the interior, but it
contains \(\xi_i\xi_j/|\xi|^2\).  The following theorem shows that this
nonlocality cannot be removed by a different finite-order differential
representative at the same energy order.

Let \(E(D)\) be a translation-invariant finite-order real differential
operator from a fixed two-component amplitude \(a\) to a symmetric metric
perturbation.  Its polynomial symbol is \(E(ik)\).  Two symbols which differ
by a right multiple of \(k^\mu k_\mu\) are identified on wave solutions.

\begin{theorem}[No finite-order local same-energy radiative splitting]
\label{thm:V71-no-local-radiative-splitting}
There is no \(E(D)\) with all of the following properties on every nonzero
real null covector \(k\):
\begin{enumerate}
\item \(E(ik)z\) is harmonic,
\begin{equation}
 k^\mu\overline{E(ik)z}_{\mu\nu}=0;
 \label{eq:V71-local-split-harmonic}
\end{equation}
\item its class is a right inverse for the V53 radiative quotient, so in
      particular it is nonzero for every \(z\ne0\); and
\item it acts at the same homogeneous energy order, equivalently there are
      constants \(0<c\le C<\infty\) such that
\begin{equation}
 c|z|\le\|E(irk)z\|\le C|z|
 \quad\text{for every }r>0
 \label{eq:V71-same-energy-bound}
\end{equation}
      after \(k\) is normalized on either null sheet.
\end{enumerate}
The conclusion also rules out a local differential map from the V68
fixed two-component potential phase to the full \((H,W)\) metric phase which
is bounded and bounded below in the V59/V68 energy.
\end{theorem}

\begin{proof}
Restrict the polynomial symbol to one null ray \(rk\).  Condition
\eqref{eq:V71-same-energy-bound} makes every matrix entry bounded for
\(0<r<\infty\).  Its positive-degree homogeneous coefficients must therefore
vanish on that ray.  As the ray was arbitrary, all positive-degree terms
vanish on the real null cone.  That cone is Zariski dense in the irreducible
quadratic hypersurface \(k^2=0\), so its polynomial vanishing ideal is
generated by \(k^2\).  Those terms act as zero on wave solutions.  The
effective on-shell symbol is therefore a constant map
\(E_0:\mathbb R^2\to\operatorname{Sym}^2\!\bigl((\mathbb R^{1,3})^*\bigr)\).

Fix \(z\).  The harmonic row now says
\begin{equation}
 k^\mu\overline{E_0z}_{\mu\nu}=0
 \label{eq:V71-constant-harmonic}
\end{equation}
for every real null \(k\).  Four null covectors can be chosen to span
\((\mathbb R^{1,3})^*\); for example
\begin{equation}
 (1,1,0,0),\quad(1,-1,0,0),\quad
 (1,0,1,0),\quad(1,0,0,1).
 \label{eq:V71-spanning-null-covectors}
\end{equation}
Equation \eqref{eq:V71-constant-harmonic} for this spanning set forces every
column of \(\overline{E_0z}\) to vanish.  Trace reversal is an involution, so
\(E_0z=0\).  This holds for both input columns, contradicting the lower bound
and the right-inverse property.

For the phase formulation, a bounded same-energy local reconstruction has an
on-shell polynomial symbol with exactly the bound
\eqref{eq:V71-same-energy-bound}: both source and target energies carry the
same factor \(|k|^2\).  The preceding argument applies componentwise.
\end{proof}

\begin{corollary}[Sharp alternatives for a radiative representative]
\label{cor:V71-radiative-alternatives}
A representative of the two radiative classes in the full harmonic metric
must use at least one of the following:
\begin{enumerate}
\item a direction-dependent order-zero Fourier multiplier, such as the V53
      TT projector;
\item a positive-order differential potential with the corresponding
      derivative loss and a different energy domain; or
\item the unreduced local harmonic metric complex, followed by quotienting
      the residual wave-gauge image at the operator level.
\end{enumerate}
No item may be relabelled as a bounded local two-scalar graph at V59 energy.
\end{corollary}

\begin{proof}
Theorem~\ref{thm:V71-no-local-radiative-splitting} excludes a finite-order
local representative with no derivative loss.  V53 exhibits the first
alternative and explicitly records its Riesz multiplier.  Positive-order
polynomial symbols grow under radial scaling unless the source norm pays the
same order.  Retaining the exact harmonic complex avoids selecting a
representative and gives the third alternative.
\end{proof}

\subsection{The correct V71 harmonic/radiative endpoint}
\label{subsec:V71-operator-quotient}

Let \(\mathscr Z_{71}\) be the closed reducing subspace of the V69 potential
range on which
\begin{equation}
 \mathscr C=0,
 \qquad D_0\mathscr C=0,
 \label{eq:V71-harmonic-reducing-space}
\end{equation}
and let \(\mathscr G_{71}\) denote the closure, in the V69 homogeneous
energy norm, of the residual wave-gauge image.  Fibrewise at every nonzero
null covector,
V53 gives
\begin{equation}
 0\longrightarrow\mathbb C^4
 \mathop{\longrightarrow}^{\Gamma_k} Z_k
 \longrightarrow Z_k/\operatorname{ran}\Gamma_k
 \longrightarrow 0,
 \qquad
 \dim(Z_k/\operatorname{ran}\Gamma_k)=2.
 \label{eq:V71-retained-radiative-complex}
\end{equation}

\begin{proposition}[Local complex, nonlocal representative]
\label{prop:V71-local-complex-endpoint}
The V69 ten-component phase equations, the harmonic rows
\eqref{eq:V71-harmonic-reducing-space}, and the residual gauge differential
are local.  Their physical interior state is the quotient
\begin{equation}
 \mathscr H_{71}^{\rm rad}:=\mathscr Z_{71}/\mathscr G_{71}.
 \label{eq:V71-radiative-Hilbert-quotient}
\end{equation}
The V53 TT multiplier gives a bounded interior Hilbert representative and
identifies its positive energy with the V68/V59 two-polarization energy.
There is no bounded finite-order local splitting of
\eqref{eq:V71-radiative-Hilbert-quotient} into the full metric phase.
Consequently V71 retains the quotient itself and does not impose a local
radiative-to-metric graph.
\end{proposition}

\begin{proof}
Locality and reduction of the harmonic shell are the V69 phase and harmonic
propagation identities.  The fibre dimensions and exact residual image are
V53.  Its norm-one TT projector gives the bounded interior direct-integral
representative and positive quotient norm.  The definition by closure makes
\(\mathscr G_{71}\) a closed subspace, hence the Hilbert quotient is well
defined.  Theorem
\ref{thm:V71-no-local-radiative-splitting} excludes the asserted local
splitting.  Retaining the quotient is therefore forced within the declared
same-energy and locality class.
\end{proof}

\begin{firewall}[Boundary gauge charges remain to be reduced]
\label{fire:V71-boundary-gauge-charges}
Proposition~\ref{prop:V71-local-complex-endpoint} is an interior
operator-level quotient.  A residual diffeomorphism with nonzero boundary
trace need not be a null vector of the GHY boundary symplectic form; it can
carry a boundary charge.  V71 therefore does not identify the inverse image
of the V59 quotient relation with a Lagrangian subspace of the unreduced V36
boundary phase.  V72 must derive the boundary residual-gauge moment map and
its BFV reduction before making that claim.
\end{firewall}

\subsection{What can already be combined with V70}
\label{subsec:V71-combination}

\begin{theorem}[Invariant direct-sum action before physical reduction]
\label{thm:V71-direct-sum-action}
The quadratic functional
\begin{equation}
 S_{71}^{\rm pre}
 :=S_{71}^{\rm min}+S_{\rm hol}
 \label{eq:V71-preaction}
\end{equation}
is local and BRST invariant off shell.  Its Euler system is the invariant
metric wave equation together with the contractible V70 holonomy rows.  Its
boundary potential is the direct sum of
\(\Theta_{71}^{\rm min}\) and \(\Theta_{\rm hol}\).  Replacing the first
summand by the exact jet-cotangent lift
\eqref{eq:V71-extended-potential} recovers the V36 conormal without changing
the V70 Euler or BFV charge.
\end{theorem}

\begin{proof}
Theorem~\ref{thm:V71-parent-elimination} and V70 give locality and separate
off-shell BRST invariance.  The variables \((H,W)\) and
\((X,J,K,L,\eta,c)\) are independent triangular coordinates, so the Euler
rows and integrations by parts split.  Theorem~\ref{thm:V71-jet-lift} changes
only the invariant boundary cotangent sector.  V70's holonomy boundary form
and charge contain none of \((H,W,r,\varpi)\), hence are unchanged.
\end{proof}

\begin{firewall}[Why this is not yet the linear SM--Einstein action]
\label{fire:V71-not-full-action}
The minimal parent has the correct reduced bulk wave equation, while the
jet-cotangent lift has the correct V36 spatial conormal.  V71 has not derived
both from one covariant Einstein--Hilbert--GHY parent on a slab with temporal
edges.  It has also not reduced nonvanishing boundary gauge charges or proved
the full corner master equation.  Theorem
\ref{thm:V71-direct-sum-action} is therefore a pre-reduction linear action
statement.  It does not yet authorize coupling the V59 continuum traction to
a local full-metric boundary domain.
\end{firewall}

\subsection{Finite algebraic certificates}
\label{subsec:V71-certificates}

The exact finite cards for this note are
\begin{align}
 \operatorname{EL}_W(S_{71}^{\rm min})&=W-DH,
 \label{eq:V71-parent-card}\\
 \det M_{36}&=1,
 \qquad \operatorname{inertia}(M_{36})=(4,6,0),
 \label{eq:V71-normal-card}\\
 M_{36}&\ne M_{\rm min},
 \label{eq:V71-mismatch-card}\\
 \Theta_{71}^{\rm ext}\big|_{\varpi=-\Delta^\vee H}
 &=\frac1{2\kappa}\int
   \langle\mathcal A_{36}r,\delta H\rangle,
 \label{eq:V71-lift-card}\\
 \operatorname{span}\{k_1,k_2,k_3,k_4\}
 &=(\mathbb R^{1,3})^*
 \label{eq:V71-null-span-card}
\end{align}
for the four null covectors in
\eqref{eq:V71-spanning-null-covectors}.  A finite verifier reconstructs
\(M_{36}\) directly from \eqref{eq:V71-jet-C}--%
\eqref{eq:V71-A36}, checks its exact rational determinant, numerical inertia,
and both mismatch witnesses.  It samples the parent elimination and the
jet-cotangent cancellation, and checks that the four null covectors have rank
four.  The polynomial scaling and variational arguments in the proofs give
the analytic no-go statements.

\subsection{V72 acquisition order}
\label{subsec:V71-next}

V72 should go upstream rather than trying another metric-only counterterm:
\begin{enumerate}
\item start from the quadratic second-derivative Einstein--Hilbert action,
      its GHY term, and the de Donder gauge-fixing term before discarding a
      divergence;
\item introduce an independent connection or complete first jet and derive
      its bulk, spatial-face, initial-face, terminal-face, and corner
      variations in one calculation;
\item identify the jet-cotangent rows of
      Theorem~\ref{thm:V71-jet-lift} inside that covariant parent and eliminate
      only algebraic pairs whose edge variation has vanished;
\item derive the residual boundary-gauge moment map, separate proper gauge
      transformations from charged boundary symmetries, and perform the BFV
      reduction on the full harmonic metric phase;
\item pull the V59 positive radiative--continuum relation back only after that
      reduction, then take its direct product with the V70 holonomy relation;
\item prove the spatial and temporal corner master equation for this complete
      linear parent before restoring curvature or Standard-Model sources.
\end{enumerate}
A local two-scalar metric representative is no longer an acquisition target;
Theorem~\ref{thm:V71-no-local-radiative-splitting} replaces it by the local
complex and operator-quotient route.

\subsection{V71 result ledger and claim boundary}
\label{subsec:V71-ledger}

Relative to V1--V70, this note proves:
\begin{enumerate}
\item a local BRST-invariant first-order parent for the invariant metric pair
      \((H,W)\), with exact elimination to the V36 minimal de Donder density;
\item its complete minimal boundary potential and algebraic normal momentum;
\item the explicit V36 first-jet conormal map and its invertible
      ten-dimensional normal matrix;
\item an exact normal-symbol obstruction to recovering the V36 conormal from
      the already-minimal parent by any metric-only quadratic boundary
      counterterm;
\item the order-of-operations alternative which keeps the Einstein divergence
      with GHY, together with an exact local jet-cotangent lift realizing the
      second alternative;
\item a same-energy scaling obstruction to every finite-order local
      differential splitting of two fixed radiative amplitudes into a full
      harmonic metric;
\item the resulting local-harmonic-complex and operator-quotient endpoint;
      and
\item the BRST-invariant direct sum of the invariant minimal parent and the
      V70 contractible holonomy action.
\end{enumerate}

The covariant Einstein--Hilbert--GHY first-order ancestry action, its temporal
edge and corner terms, the residual boundary-gauge moment map, the reduced
metric--continuum Lagrangian relation, the complete linear corner master
equation, nonlinear curvature and Brown--York/Bianchi stress closure, and the
coupled Standard-Model continuum remain open.  The full Standard
Model--Einstein continuum is not proved.

\subsection{Append-only record}
\label{subsec:V71-append-record}

V71 was appended to the complete V70 source.  The SHA--256 digest of the exact
parent file was
\[
 \texttt{3b05e612185a8f74230122baea9104aead7a589bf1b9527bf1b37b9a7cf1ea9f}.
\]
The next research note is V72 and begins with the unminimalized quadratic
Einstein--Hilbert--GHY ancestry parent and its full slab/corner variation.

% END APPENDED VERSION V71 MATERIAL

% BEGIN APPENDED VERSION V72 MATERIAL

\clearpage
\section{V72: residual metric-gauge moment map, the exact corner cocycle,
and the second boundary BFV factor}
\label{sec:V72-residual-moment-map}

\subsection{Purpose and the distinction forced by V71}
\label{subsec:V72-purpose}

V71 proved that the already-minimal invariant parent cannot be converted to
the V36 GHY--de Donder potential by a metric-only boundary scalar.  It also
proved that the two radiative classes have no bounded finite-order local
splitting into a full harmonic metric at V59 energy.  The prescribed upstream
route is therefore to retain the full harmonic metric complex and reduce its
residual gauge action on the actual V36 boundary phase.

One separation is essential.  The V69/V70 embedding differential, denoted
\(s_{\rm emb}\) here, obeys
\begin{equation}
 s_{\rm emb}H=s_{\rm emb}W=0.
 \label{eq:V72-embedding-invariance}
\end{equation}
The V53 residual wave-gauge transformation changes \(H\).  It is not another
component of \(s_{\rm emb}\).  V72 constructs its Hamiltonian action and
shows that the corresponding boundary charge has no central term for
tangentially edge-vanishing parameters.  For unrestricted edge data the
failure is an explicit codimension-two cocycle.  This is the missing input
which must be carried into the unminimalized Einstein--Hilbert--GHY ancestry
parent; it is computed before a corner master equation is claimed.

No companion audit is due in V72.  The next mandatory Yang--Mills and
Navier--Stokes audit remains V75.

\subsection{All faces of the minimal invariant parent}
\label{subsec:V72-all-faces}

Let
\begin{equation}
 M_T=[0,T]\times\Omega,
 \qquad
 \partial M_T=\Sigma_T\cup(-\Sigma_0)\cup\Gamma_T,
 \qquad
 \Gamma_T=[0,T]\times\partial\Omega,
 \label{eq:V72-slab-decomposition}
\end{equation}
with the displayed induced orientations.  Write \(N_F\) for the outward
conormal on a face \(F\).  V71's parent action gives, before imposing a face
domain,
\begin{equation}
 \Theta_{71,F}^{\rm min}
 =-\frac1{2\kappa}\int_F
      \mathbb G(N_F\mathbin{\lrcorner}W,\delta H).
 \label{eq:V72-parent-face-potential}
\end{equation}

\begin{proposition}[Oriented full-slab Green identity]
\label{prop:V72-full-slab-Green}
For two smooth parent fields \((H,W)\) and \((K,V)\), the antisymmetrized
first variation is
\begin{align}
 2\kappa\,\Omega_{71}^{\rm min}
 =-\sum_{F\in\{\Sigma_T,-\Sigma_0,\Gamma_T\}}
 \int_F\bigl[
 &\mathbb G(N_F\mathbin{\lrcorner}W,K)
 \nonumber\\[-2mm]
 &-\mathbb G(H,N_F\mathbin{\lrcorner}V)
 \bigr].
 \label{eq:V72-full-slab-Green}
\end{align}
There is no intrinsic codimension-two term in this first-order bulk Green
identity.  Corner terms arise only after a face expression is integrated by
parts tangentially, or after an explicit face action is added.
\end{proposition}

\begin{proof}
Apply Stokes' theorem once to the cross term
\(-\mathbb G_4(W,DH)\) in the V71 parent and use the oriented boundary in
\eqref{eq:V72-slab-decomposition}.  Antisymmetrization gives
\eqref{eq:V72-full-slab-Green}.  A first-order bulk integration has support on
codimension-one faces.  No tangential face integration has yet been made, so
there is no boundary-of-boundary contribution.
\end{proof}

On \(\Gamma_T\), V71's jet-cotangent lift replaces
\(\Theta_{71,\Gamma}^{\rm min}\) by the exact V36 potential.  It does not
supply cap or edge actions.  V72 consequently derives the residual symmetry
first on the spatial face and records its edge defect explicitly.

\subsection{The second, residual differential}
\label{subsec:V72-residual-differential}

For a covector field \(\epsilon\), define the linearized metric-gauge map
\begin{equation}
 (\mathcal G\epsilon)_{\mu\nu}
 :=D_\mu\epsilon_\nu+D_\nu\epsilon_\mu.
 \label{eq:V72-metric-gauge-map}
\end{equation}
Let \(\gamma_\mu\) be an odd residual ghost and set
\begin{align}
 q_{\rm res}H&=\mathcal G\gamma,
 &q_{\rm res}W&=D(\mathcal G\gamma),
 \label{eq:V72-residual-on-HW}\\
 q_{\rm res}\gamma&=0,
 &q_{\rm res}(\text{V70 variables})&=0.
 \label{eq:V72-residual-on-ghost}
\end{align}
The embedding differential acts trivially on \(\gamma\).

\begin{proposition}[Commuting linear bicomplex and residual shell]
\label{prop:V72-bicomplex}
The two odd differentials satisfy
\begin{equation}
 s_{\rm emb}^2=q_{\rm res}^2=0,
 \qquad
 s_{\rm emb}q_{\rm res}+q_{\rm res}s_{\rm emb}=0.
 \label{eq:V72-bicomplex}
\end{equation}
Moreover
\begin{equation}
 C_\nu(\mathcal G\gamma)=\Box\gamma_\nu.
 \label{eq:V72-gauge-C-identity}
\end{equation}
Thus \(q_{\rm res}\) preserves the harmonic shell and the reduced metric wave
solution space when
\begin{equation}
 \Box\gamma=0.
 \label{eq:V72-residual-ghost-shell}
\end{equation}
This is an on-shell residual differential.  It is not yet an off-shell bulk
Faddeev--Popov complex.
\end{proposition}

\begin{proof}
Nilpotency of \(q_{\rm res}\) follows from
\(q_{\rm res}\gamma=0\), and the anticommutator vanishes because the two
differentials act on disjoint triangular factors while both annihilate the
other ghost.  For the harmonic covector,
\begin{align*}
 C_\nu(\mathcal G\gamma)
 &=D^\mu(D_\mu\gamma_\nu+D_\nu\gamma_\mu)
   -\frac12D_\nu(2D^\rho\gamma_\rho)\\
 &=\Box\gamma_\nu
\end{align*}
by commutation at the flat anchor.  The metric wave operator also commutes
with \(\mathcal G\), proving preservation on
\eqref{eq:V72-residual-ghost-shell}.  No unconstrained ghost action has been
introduced, so the final scope statement follows.
\end{proof}

\subsection{Exact V36 conormal of a residual gauge wave}
\label{subsec:V72-pure-gauge-conormal}

Let \(a,b\in\{0,1,2\}\) be tangent to \(\Gamma_T\), and put
\begin{equation}
 \Box_\partial:=D^aD_a.
 \label{eq:V72-boundary-wave-operator}
\end{equation}
For a scalar \(f\) define the tangential trace-adjusted Hessian
\begin{equation}
 \mathscr P^{ab}(f)
 :=-D^aD^bf+\eta_\partial^{ab}\Box_\partial f.
 \label{eq:V72-P-operator}
\end{equation}
It obeys the exact differential identity
\begin{equation}
 D_a\mathscr P^{ab}(f)=0.
 \label{eq:V72-P-divergence}
\end{equation}

\begin{theorem}[Pure residual-gauge GHY conormal]
\label{thm:V72-pure-gauge-conormal}
If \(\Box\epsilon=0\), then the V36 conormal of
\(\mathcal G\epsilon\) is
\begin{equation}
 \boxed{\quad
 \mathfrak p_{36}^{\mu\nu}(\mathcal G\epsilon)
 =e_a^\mu e_b^\nu\mathscr P^{ab}(\epsilon_n).
 \quad}
 \label{eq:V72-pure-gauge-p36}
\end{equation}
In particular, its gauge-fixing part vanishes, all mixed and normal conormal
components vanish, and its induced block is tangentially divergence free.
\end{theorem}

\begin{proof}
Proposition~\ref{prop:V72-bicomplex} gives
\(C(\mathcal G\epsilon)=0\).  Commuting derivatives in the linearized
extrinsic curvature gives
\begin{align}
 K_{ab}^{(1)}(\mathcal G\epsilon)
 &=\frac12\bigl[
 D_n(D_a\epsilon_b+D_b\epsilon_a)
 -D_a(D_n\epsilon_b+D_b\epsilon_n)
 \nonumber\\
 &\hspace{31mm}
 -D_b(D_n\epsilon_a+D_a\epsilon_n)
 \bigr]
 =-D_aD_b\epsilon_n.
 \label{eq:V72-pure-gauge-K}
\end{align}
Its tangential trace is \(-\Box_\partial\epsilon_n\), so the Brown--York
combination is exactly \(\mathscr P^{ab}(\epsilon_n)\).  Substitution in the
V36 conormal proves \eqref{eq:V72-pure-gauge-p36}.  Finally
\[
 D_a\mathscr P^{ab}(f)
 =-\Box_\partial D^bf+D^b\Box_\partial f=0,
\]
which proves the divergence statement.
\end{proof}

\subsection{Hamiltonian residual action on the V36 boundary phase}
\label{subsec:V72-moment-map}

Write a point of the unconstrained linear V36 boundary phase as
\begin{equation}
 z=(H,p),
 \qquad
 \omega_{36}(z_1,z_2)
 :=\frac1{2\kappa}\int_{\Gamma_T}
     \left(p_1^{\mu\nu}H_{2,\mu\nu}
          -H_{1,\mu\nu}p_2^{\mu\nu}\right).
 \label{eq:V72-canonical-V36-phase}
\end{equation}
A residual parameter has the lifted phase vector
\begin{equation}
 \mathcal V_\epsilon
 :=\left(\mathcal G\epsilon,
          \mathfrak p_{36}(\mathcal G\epsilon)\right).
 \label{eq:V72-lifted-gauge-vector}
\end{equation}
Define
\begin{equation}
 \mu_\epsilon(z)
 :=\omega_{36}(\mathcal V_\epsilon,z)
 =\frac1{2\kappa}\int_{\Gamma_T}
 \left[
  \mathfrak p_{36}(\mathcal G\epsilon)^{\mu\nu}H_{\mu\nu}
  -(\mathcal G\epsilon)_{\mu\nu}p^{\mu\nu}
 \right].
 \label{eq:V72-residual-moment-map}
\end{equation}

\begin{theorem}[Exact residual boundary moment map]
\label{thm:V72-residual-moment-map}
For every residual wave parameter \(\Box\epsilon=0\),
\begin{equation}
 \iota_{\mathcal V_\epsilon}\omega_{36}
 =\delta\mu_\epsilon.
 \label{eq:V72-Hamiltonian-identity}
\end{equation}
Thus the V53 residual metric-gauge action is Hamiltonian on the full V36
boundary phase.  Its moment map is not the V70 holonomy charge: it acts on the
invariant coordinates \((H,p)\), whereas the V70 charge acts only on the
contractible embedding-jet cotangent sector.
\end{theorem}

\begin{proof}
The phase form is constant and the lifted gauge vector is independent of
\(z\).  Therefore differentiation of
\eqref{eq:V72-residual-moment-map} in a direction \(\delta z\) gives
\(\omega_{36}(\mathcal V_\epsilon,\delta z)\), which is precisely
\eqref{eq:V72-Hamiltonian-identity}.  The final distinction follows from
\eqref{eq:V72-embedding-invariance} and the explicit support of the V70 BFV
charge.
\end{proof}

On the V36 conormal graph \(p=\mathfrak p_{36}(H)\), the moment map has the
fully field-based form
\begin{equation}
 2\kappa\,\mu_\epsilon(H)
 =\int_{\Gamma_T}
 \left[
  \mathscr P^{ab}(\epsilon_n)H_{ab}
  -(D_\mu\epsilon_\nu+D_\nu\epsilon_\mu)
      \mathfrak p_{36}^{\mu\nu}(H)
 \right].
 \label{eq:V72-moment-map-on-graph}
\end{equation}
This formula keeps normal derivatives of the parameter visible; they cannot
be removed by tangential integration by parts.

\subsection{The exact edge cocycle}
\label{subsec:V72-edge-cocycle}

Let
\begin{equation}
 \partial\Gamma_T=\mathcal E_T\cup(-\mathcal E_0),
 \qquad
 \mathcal E_t:=\{t\}\times\partial\Omega,
 \label{eq:V72-edge-decomposition}
\end{equation}
and write \(\nu_a\) for the outward conormal to an edge within
\(\Gamma_T\).  Assume \(\partial\Omega\) itself has no boundary; otherwise
its additional edges are included in the same notation.

\begin{theorem}[Residual-gauge central cocycle is a corner term]
\label{thm:V72-corner-cocycle}
For two residual wave parameters \(\epsilon,\zeta\),
\begin{align}
 \mathcal K_{72}(\epsilon,\zeta)
 &:=\{\mu_\epsilon,\mu_\zeta\}
 =\omega_{36}(\mathcal V_\epsilon,\mathcal V_\zeta)
 \nonumber\\
 &=\frac1\kappa\int_{\partial\Gamma_T}
  \nu_a\left[
   \mathscr P^{ab}(\epsilon_n)\zeta_b
   -\mathscr P^{ab}(\zeta_n)\epsilon_b
  \right].
 \label{eq:V72-corner-cocycle}
\end{align}
It is bilinear and antisymmetric.  It is a central Lie-algebra two-cocycle
because the residual linear gauge algebra is abelian.  It is not identically
zero in the local residual-wave jet algebra.  In particular it vanishes on
the edge domain
\begin{equation}
 \mathscr E_0
 :=\left\{\epsilon:\Box\epsilon=0,
     \ \epsilon_b|_{\partial\Gamma_T}=0
     \text{ for }b\in T\Gamma_T\right\}.
 \label{eq:V72-edge-zero-domain}
\end{equation}
\end{theorem}

\begin{proof}
Theorem~\ref{thm:V72-pure-gauge-conormal} gives
\begin{align*}
 2\kappa\,\omega_{36}
   (\mathcal V_\epsilon,\mathcal V_\zeta)
 =\int_{\Gamma_T}\bigl[
 &\mathscr P^{ab}(\epsilon_n)
        (D_a\zeta_b+D_b\zeta_a)\\
 &-(D_a\epsilon_b+D_b\epsilon_a)
        \mathscr P^{ab}(\zeta_n)
 \bigr].
\end{align*}
Symmetry of \(\mathscr P\) makes this twice
\[
 \int_{\Gamma_T}
 \left[\mathscr P^{ab}(\epsilon_n)D_a\zeta_b
       -\mathscr P^{ab}(\zeta_n)D_a\epsilon_b\right].
\]
Integrate tangentially and use
\eqref{eq:V72-P-divergence}; only the oriented boundary of the face remains.
This is \eqref{eq:V72-corner-cocycle}.  Antisymmetry is visible.  Every
Chevalley--Eilenberg differential of a constant antisymmetric form on an
abelian algebra is zero, so it is a two-cocycle.

For a nonzero local witness, use flat coordinates \((t,y^1,y^2,x_n)\) and
take
\begin{equation}
 \epsilon_n=t y^1,\qquad
 \zeta_1=t^2+x_n^2,
 \qquad\text{all other components zero}.
 \label{eq:V72-cocycle-witness}
\end{equation}
Both parameters solve the four-dimensional wave equation.  On \(x_n=0\),
\(\mathscr P^{01}(\epsilon_n)=1\).  The final-edge density in
\eqref{eq:V72-corner-cocycle} is therefore \(\nu_0T^2/\kappa\), while
the initial-edge density is zero.  Thus the local corner form is nonzero.
The edge condition \eqref{eq:V72-edge-zero-domain} kills both terms
pointwise.
\end{proof}

\begin{corollary}[Proper edge algebra versus charged corner algebra]
\label{cor:V72-proper-versus-corner}
The moment map is equivariant on \(\mathscr E_0\), and the corresponding
Hamiltonians Poisson commute.  If tangential residual parameters are allowed
to have arbitrary edge traces, their Hamiltonians realize the central
extension determined by \(\mathcal K_{72}\).  Such transformations cannot be
quotiented as ordinary proper gauge transformations without either imposing
an edge domain or adjoining corner degrees of freedom whose symplectic
cocycle cancels \eqref{eq:V72-corner-cocycle}.
\end{corollary}

\begin{proof}
For an abelian Hamiltonian action, equivariance is exactly the vanishing of
its constant moment-map cocycle.  Apply
Theorem~\ref{thm:V72-corner-cocycle}.  Symplectic reduction by a centrally
extended action requires the central value to be included or cancelled,
which gives the stated alternatives.
\end{proof}

\subsection{The residual boundary BFV factor}
\label{subsec:V72-residual-BFV}

Promote the residual parameter to the odd ghost \(\gamma\), retain its
cotangent momentum, and define the charge
\begin{equation}
 \Omega_{72}^{\rm res}(z,\gamma)
 :=\mu_\gamma(z).
 \label{eq:V72-residual-BFV-charge}
\end{equation}
The charge is understood on the residual wave-ghost shell
\eqref{eq:V72-residual-ghost-shell}; its cotangent Hamilton equation sends the
ghost momentum to the moment-map constraint.

\begin{theorem}[Conditional residual master equation]
\label{thm:V72-residual-master}
On the edge domain \(\mathscr E_0\),
\begin{equation}
 \{\Omega_{72}^{\rm res},\Omega_{72}^{\rm res}\}=0.
 \label{eq:V72-residual-master}
\end{equation}
It generates \(q_{\rm res}H=\mathcal G\gamma\) and the corresponding conormal
translation.  For unrestricted edge ghosts, polarization of its self-bracket
is exactly \(\mathcal K_{72}\); hence the master-equation defect is supported
on \(\partial\Gamma_T\), not in the bulk of the face.
\end{theorem}

\begin{proof}
The residual algebra is abelian, so the BFV self-bracket is the polarization
of the moment-map cocycle.  Theorem~\ref{thm:V72-corner-cocycle} makes that
cocycle zero on \(\mathscr E_0\), proving
\eqref{eq:V72-residual-master}.  The Hamiltonian identity in
Theorem~\ref{thm:V72-residual-moment-map} gives the metric and conormal
transformations.  The same cocycle formula identifies the unrestricted
defect and its support.
\end{proof}

Let \(\Omega_{70}^{\rm hol}\) be the action-derived V70 holonomy charge.

\begin{corollary}[Two commuting boundary BFV factors]
\label{cor:V72-total-BFV}
On the product of the V70 holonomy phase with the V36 residual phase, and on
\(\mathscr E_0\),
\begin{equation}
 \Omega_{72}^{\rm tot}
 :=\Omega_{70}^{\rm hol}+\Omega_{72}^{\rm res},
 \qquad
 \{\Omega_{72}^{\rm tot},\Omega_{72}^{\rm tot}\}=0.
 \label{eq:V72-total-BFV-charge}
\end{equation}
The first factor removes the differential embedding-jet quartet; the second
implements the residual harmonic-metric quotient.  They must not be merged
before their distinct moment maps and domains are imposed.
\end{corollary}

\begin{proof}
V70 proves the holonomy self-bracket is zero.  Theorem
\ref{thm:V72-residual-master} proves the residual self-bracket is zero on the
edge domain.  The holonomy charge contains only
\((X,J,c,c^{(1)})\) and their momenta, whereas the residual charge contains
\((H,p,\gamma)\).  The product symplectic form has zero cross brackets, so the
mixed bracket vanishes.
\end{proof}

\subsection{Compatibility with the V71 jet-cotangent lift}
\label{subsec:V72-jet-equivariance}

Extend the residual differential to the V71 boundary jet variables by
\begin{align}
 q_{\rm res}r_\alpha
 &=D_\alpha(\mathcal G\gamma),
 \label{eq:V72-q-r}\\
 q_{\rm res}\varpi^\alpha
 &=-\Delta_{71,\alpha}^{\vee}(\mathcal G\gamma).
 \label{eq:V72-q-varpi}
\end{align}

\begin{proposition}[Equivariance of the exact conormal lift]
\label{prop:V72-jet-equivariance}
The compatibility rows and momentum row of V71 obey
\begin{align}
 q_{\rm res}(r_a-D_aH)&=0,
 &q_{\rm res}(r_n-W_n)&=0,
 \label{eq:V72-q-compatibility}\\
 q_{\rm res}
 \left(\varpi^\alpha+\Delta_{71,\alpha}^{\vee}H\right)&=0.
 \label{eq:V72-q-momentum-row}
\end{align}
Consequently the pullback of the residual moment map through the V71
jet-cotangent relation is the same function
\eqref{eq:V72-residual-moment-map}.  No auxiliary jet charge is added.
\end{proposition}

\begin{proof}
Insert \eqref{eq:V72-residual-on-HW},
\eqref{eq:V72-q-r}, and \eqref{eq:V72-q-varpi}; all three identities cancel
term by term.  V71 proves that the constrained extended potential pulls back
to \(\Theta_{36}\).  Its exterior derivative therefore pulls back to
\(\omega_{36}\), and the Hamiltonian contraction is the same moment map.
The auxiliary coordinates are uniquely fixed by their rows, so there is no
independent translation generator.
\end{proof}

\subsection{Reduction order forced by the corner calculation}
\label{subsec:V72-reduction-order}

\begin{theorem}[Necessary linear reduction order]
\label{thm:V72-reduction-order}
Within the local same-energy programme, a boundary realization of the V59
radiative--continuum relation must proceed in the following order:
\begin{enumerate}
\item retain the full V36 harmonic metric boundary phase and the V71
      first-jet cotangent lift;
\item choose an edge domain which kills \(\mathcal K_{72}\), or add a corner
      phase with the opposite cocycle;
\item impose the residual moment-map constraint
      \(\mu_\epsilon=0\) for proper parameters and quotient their Hamiltonian
      action;
\item identify the resulting two radiative classes with the V53/V59 quotient;
      and only then
\item take the product with the V70 holonomy reduction and the positive V59
      continuum interface.
\end{enumerate}
Skipping the second step leaves a nonzero BFV master-equation defect.  Skipping
the third confuses the six-dimensional harmonic amplitude space with its
two-dimensional radiative quotient.
\end{theorem}

\begin{proof}
The corner alternative is Corollary
\ref{cor:V72-proper-versus-corner}.  Hamiltonian reduction requires the zero
level of the moment map before quotienting the proper action.  V53 proves that
the residual quotient, rather than the harmonic kernel itself, has two
radiative dimensions.  V70 and Corollary~\ref{cor:V72-total-BFV} show that the
holonomy factor is independent and can be taken afterward.  Each omitted
step therefore violates one of the proved incidence, dimension, or master
identities.
\end{proof}

\begin{firewall}[What V72 has not identified]
\label{fire:V72-reduction-scope}
Theorem~\ref{thm:V72-reduction-order} gives necessary steps.  V72 has not yet
proved that the zero level of \(\mu\), divided by
\(\mathscr E_0\), is globally isomorphic to the V59 boundary quotient on a
half-space.  It has not supplied a corner phase cancelling
\(\mathcal K_{72}\) for unrestricted ghosts.  It also has not derived the
V71 jet-cotangent sector from a single unminimalized
Einstein--Hilbert--GHY action.  Those are construction gates, not consequences
of the moment-map calculation.
\end{firewall}

\subsection{Finite algebraic certificates}
\label{subsec:V72-certificates}

For a Fourier covector \(k\), the exact cards are
\begin{align}
 C(\mathcal G_k\epsilon)&=k^2\epsilon,
 \label{eq:V72-C-card}\\
 K_{ab}(\mathcal G_k\epsilon)&=-k_ak_b\epsilon_n,
 \label{eq:V72-K-card}\\
 k_a\mathscr P^{ab}_k(\epsilon_n)&=0,
 \label{eq:V72-divergence-card}\\
 \omega_{36}(\mathcal V_\epsilon,\mathcal V_\zeta)&=0
 \quad\text{on an edge-free Fourier face},
 \label{eq:V72-isotropy-card}\\
 q_{\rm res}
 \left(\varpi+\Delta_{71}^{\vee}H\right)&=0.
 \label{eq:V72-jet-card}
\end{align}
A finite verifier samples nonzero real null covectors and arbitrary residual
parameters, reconstructs \(\mathcal G_k\), \(C\), \(K\), and
\(\mathscr P_k\), and checks all first three identities.  It checks isotropy
of the lifted gauge vectors at each tangential Fourier mode, the Hamiltonian
moment-map derivative, preservation of the V71 jet rows, and zero cross
bracket with a disjoint V70 canonical factor.  The tangential Stokes proof,
not the Fourier samples, establishes the exact corner formula.

\subsection{V73 acquisition order}
\label{subsec:V72-next}

V73 should now return to the ancestry action with the residual corner data
fixed:
\begin{enumerate}
\item write the quadratic Einstein--Hilbert Jacobi density before integration
      by parts, add the quadratic GHY and de Donder terms, and introduce an
      independent complete first jet or connection;
\item derive its algebraic jet equations and every face potential, verifying
      that elimination gives both the V36 bulk equation and spatial conormal;
\item compute its cap potentials and the induced edge one-form rather than
      borrowing the spatial formula on a timelike cap;
\item decide between the edge-zero residual domain and a minimal corner
      cotangent field whose cocycle is
      \(-\mathcal K_{72}\);
\item construct the off-shell residual ghost--antighost action and show that
      its boundary BFV charge reduces to
      \(\Omega_{72}^{\rm res}\) on the residual ghost shell;
\item only after the full linear master equation holds, identify the reduced
      two-polarization V59 relation and join the continuum traction.
\end{enumerate}
If the connection parent produces a different edge cocycle, V73 must compare
it by an explicit coboundary calculation; it may not discard either corner
term.

\subsection{V72 result ledger and claim boundary}
\label{subsec:V72-ledger}

Relative to V1--V71, this note proves:
\begin{enumerate}
\item the oriented all-face Green identity of the invariant minimal parent;
\item the separate residual differential on \((H,W)\), its commutation with
      the V69 embedding differential, and its exact harmonic-shell condition;
\item the full V36 conormal of a residual gauge wave and its divergence-free
      induced block;
\item the Hamiltonian residual moment map on the unreduced V36 boundary
      phase;
\item the exact codimension-two central cocycle for unrestricted edge data;
\item the nilpotent residual BFV factor on the tangentially edge-zero domain
      and its zero cross bracket with the V70 holonomy charge;
\item equivariance of the V71 jet-cotangent lift under the residual action;
      and
\item the necessary reduction order before the V59 radiative--continuum
      relation can be pulled back to a full metric domain.
\end{enumerate}

The unminimalized Einstein--Hilbert--GHY first-order ancestry action, cap and
corner action terms, an unrestricted corner-cocycle cancellation, the
off-shell residual ghost action, the globally identified reduced
metric--continuum Lagrangian relation, the complete linear corner master
equation, nonlinear curvature and Brown--York/Bianchi stress closure, and the
coupled Standard-Model continuum remain open.  The full Standard
Model--Einstein continuum is not proved.

\subsection{Append-only record}
\label{subsec:V72-append-record}

V72 was appended to the complete V71 source.  The SHA--256 digest of the exact
parent file was
\[
 \texttt{dee586c7741bb852896fc9b108c190965d248ad370e7379013a365ea49055767}.
\]
The next research note is V73 and begins with the unintegrated quadratic
Einstein--Hilbert--GHY ancestry parent and its face/corner variation.

% END APPENDED VERSION V72 MATERIAL


% =============================================================================
% BEGIN APPENDED VERSION V73 MATERIAL
% =============================================================================
\clearpage
\section*{V73 research note: the torsion-free Palatini ancestor, the GHY
boundary-polarization obstruction, and an exact compatible-jet repair}
\addcontentsline{toc}{section}{V73 Palatini ancestor and GHY boundary polarization}
\label{sec:V73-Palatini-GHY-ancestor}

\subsection*{Purpose, source comparison, and direction of attack}

V71 proved that the minimal wave parent cannot be converted into the V36
GHY conormal by a metric-only quadratic boundary counterterm.  V72 then
identified the residual harmonic moment map and its genuine corner cocycle.
The next missing link is therefore upstream of either calculation: one must
retain the connection, or an equivalent complete first jet, before performing
the integrations by parts which produce the minimal wave density.

The attached Grand Tensor manuscript and the newest available attached
Einstein--Standard Model ancestry manuscript were searched for a local
quadratic connection parent with its complete planar face polarization.  They
contain the global Palatini--Holst, GHY, Brown--York, and spectral-tail
interfaces, but do not print the calculation below.  This note does not repeat
their continuum handoff theorems.  It supplies the missing flat-anchor
calculation and distinguishes two boundary operations which cannot be
identified:
\begin{enumerate}[label=\textup{(\roman*)},leftmargin=2.8em]
\item cancelling arbitrary independent connection variations at a face;
\item using the GHY representative and varying the boundary connection as the
      Levi--Civita jet of the boundary metric variation.
\end{enumerate}
The distinction is forced by an explicit one-component witness.  An enlarged
compatible-jet parent then restores free variation without changing the
reduced Einstein--GHY action.

All statements are at the same planar Minkowski anchor and with the same
curvature and outward-normal convention as V36.  The calculation is local on
one face.  Summing it with the Stokes orientation gives every smooth side and
cap.  The joints of a piecewise slab remain a separate corner row and are not
silently set to zero.

\subsection{The unintegrated torsion-free connection parent}
\label{subsec:V73-bulk-parent}

Let (H_{\mu\nu}=H_{\nu\mu}), let
\(overline H_{\mu\nu}=H_{\mu\nu}-\tfrac12\eta_{\mu\nu}H\), and let
\(A^\alpha{}_{\mu\nu}=A^\alpha{}_{\nu\mu}\) be an independent torsion-free
connection perturbation.  Put
\begin{align}
 \mathcal R^{(1)}_{\mu\nu}(A)
 &:=D_\alpha A^\alpha{}_{\mu\nu}
    -D_\nu A^\alpha{}_{\mu\alpha},
 \label{eq:V73-linear-Ricci}\\
 \mathcal Q(A)
 &:=\eta^{\mu\nu}
 \left(
 A^\alpha{}_{\alpha\beta}A^\beta{}_{\mu\nu}
 -A^\alpha{}_{\nu\beta}A^\beta{}_{\mu\alpha}
 \right),
 \label{eq:V73-connection-quadratic}\\
 \mathcal L_{\rm Pal}^{(2)}(H,A)
 &:=-\overline H^{\mu\nu}\mathcal R^{(1)}_{\mu\nu}(A)
    +\mathcal Q(A).
 \label{eq:V73-Palatini-parent-density}
\end{align}
This is the coefficient of order two in
\(\sqrt{-g}\,g^{\mu\nu}R_{\mu\nu}(A)\), after the order-one background
divergence is separated.  Indeed,
\begin{equation}
 \sqrt{-g}\,g^{\mu\nu}
 =\eta^{\mu\nu}-\overline H^{\mu\nu}+O(H^2),
 \label{eq:V73-densitized-inverse-expansion}
\end{equation}
and the two quadratic connection contractions in the Ricci tensor give
\eqref{eq:V73-connection-quadratic}.

For an oriented face \(F\) with outward conormal \(n\), define the
free-connection transgression
\begin{equation}
 b_{\rm can}(H,A;n)
 :=n_\alpha\overline H^{\mu\nu}A^\alpha{}_{\mu\nu}
   -n_\nu\overline H^{\mu\nu}A^\alpha{}_{\mu\alpha}.
 \label{eq:V73-canonical-transgression}
\end{equation}

\begin{theorem}[Exact algebraic Palatini elimination]
\label{thm:V73-Palatini-elimination}
On a flat slab, set
\begin{equation}
 S_{\rm can}^{(2)}[H,A]
 :=\frac1{2\kappa}
 \left{
 \int_\Omega\mathcal L_{\rm Pal}^{(2)}(H,A)
 +\sum_F\int_F b_{\rm can}(H,A;n_F)
 \right}.
 \label{eq:V73-canonical-parent-action}
\end{equation}
Then arbitrary torsion-free variations of \(A\) have no face term.  The
connection Euler equation is algebraic and has the unique solution
\begin{equation}
 \boxed{
 A^\alpha{}_{\mu\nu}=\Gamma^\alpha{}_{\mu\nu}(H)
 :=\frac12\eta^{\alpha\rho}
 \left(D_\mu H_{\rho\nu}+D_\nu H_{\rho\mu}
       -D_\rho H_{\mu\nu}\right).}
 \label{eq:V73-Levi-Civita-elimination}
\end{equation}
The reduced bulk Euler operator is the linearized Einstein tensor with the
covariant-metric sign of V36:
\begin{equation}
 \delta_H S_{\rm can}^{(2)}\big|_{A=\Gamma(H),\,\mathrm{bulk}}
 =-\frac1{2\kappa}\int_\Omega
 G_{(1)}^{\mu\nu}(H)\,\delta H_{\mu\nu}.
 \label{eq:V73-reduced-bulk-Einstein}
\end{equation}
\end{theorem}

\begin{proof}
Integrating the two derivative terms in
\eqref{eq:V73-Palatini-parent-density} once gives
\begin{align}
 &\int_\Omega\mathcal L_{\rm Pal}^{(2)}(H,A)
 +\sum_F\int_F b_{\rm can}(H,A;n_F)
 \nonumber\\
 &\quad=\int_\Omega
 \left[
 (D_\alpha\overline H^{\mu\nu})A^\alpha{}_{\mu\nu}
 -(D_\nu\overline H^{\mu\nu})A^\alpha{}_{\mu\alpha}
 +\mathcal Q(A)
 \right].
 \label{eq:V73-algebraic-parent}
\end{align}
Thus no derivative of \(A\) and no boundary variation of \(A\) remains.
The symmetric-connection Euler equation is the linearization of densitized
metric compatibility.  If
\(T_\lambda:=A^\rho{}_{\lambda\rho}\), it reads
\begin{equation}
 D_\lambda\overline H^{\mu\nu}
 =A^\mu{}_{\lambda\rho}\eta^{\rho\nu}
  +A^\nu{}_{\lambda\rho}\eta^{\mu\rho}
  -T_\lambda\eta^{\mu\nu}.
 \label{eq:V73-densitized-compatibility}
\end{equation}
Contracting with \(\eta_{\mu\nu}\) in four dimensions gives
\(T_\lambda=\tfrac12D_\lambda H\).  Substitution into
\eqref{eq:V73-densitized-compatibility} gives, after lowering indices,
\begin{equation}
 A_{\mu\lambda\nu}+A_{\nu\lambda\mu}=D_\lambda H_{\mu\nu}.
 \label{eq:V73-linear-metric-compatibility}
\end{equation}
Cyclically permuting \(\lambda,\mu,\nu\), using torsion freedom in the
last two indices, and taking the usual sum-minus combination proves
\eqref{eq:V73-Levi-Civita-elimination}.  The same calculation with zero
right-hand side proves uniqueness; equivalently, the forty-dimensional
quadratic form in the independent components of \(A\) is nondegenerate.

Only the term
\(-\overline H^{\mu\nu}\mathcal R^{(1)}_{\mu\nu}(A)\) depends explicitly
on \(H\) before elimination.  Trace reversal is self-adjoint, so at
\(A=\Gamma(H)\) its \(H\)-variation is
\(-G_{(1)}^{\mu\nu}(H)\delta H_{\mu\nu}\).  The envelope identity removes
the \(A\)-variation because the algebraic connection equation holds.  This
proves \eqref{eq:V73-reduced-bulk-Einstein}.
\end{proof}

\begin{corollary}[Exact completed-square identity]
\label{cor:V73-connection-completed-square}
Let \(C:=A-\Gamma(H)\).  The algebraic density in
\eqref{eq:V73-algebraic-parent} satisfies
\begin{equation}
 \mathcal L_{\rm alg}(H,A)
 =\mathcal L_{\rm alg}(H,\Gamma(H))+\mathcal Q(C).
 \label{eq:V73-connection-square}
\end{equation}
There is no linear term in \(C\).
\end{corollary}

\begin{proof}
The density is affine in \(A\) plus the homogeneous quadratic form
\(\mathcal Q(A)\).  The linear term in its translation by
\(\Gamma(H)\) is precisely the connection Euler equation, which vanishes by
Theorem~\ref{thm:V73-Palatini-elimination}.  Homogeneity leaves
\(\mathcal Q(C)\).
\end{proof}

\subsection{Why the free-connection transgression is not GHY}
\label{subsec:V73-boundary-obstruction}

Work first on the V36 spacelike-normal side face.  Decompose
\begin{equation}
 q_{ab}:=H_{ab},\qquad
 u_a:=H_{an},\qquad
 \phi:=H_{nn},\qquad
 \tau:=\eta^{ab}q_{ab},
 \label{eq:V73-face-decomposition}
\end{equation}
and retain
\begin{equation}
 K^{(1)}_{ab}
 =\frac12(D_nq_{ab}-D_au_b-D_bu_a),
 \qquad
 \Pi_{(1)}^{ab}=K_{(1)}^{ab}-\eta^{ab}K_{(1)}.
 \label{eq:V73-K-Pi}
\end{equation}

\begin{lemma}[Complete face one-form of the canonical transgression]
\label{lem:V73-canonical-face-form}
After \(A=\Gamma(H)\), the \(H\)-variation of
\eqref{eq:V73-canonical-transgression} is
\begin{align}
 \vartheta_{\rm can}(H;\delta H)
 &=-K_{(1)}^{ab}\delta q_{ab}
   +\frac12K_{(1)}(\delta\tau+\delta\phi)
 \nonumber\\
 &\quad
   +\frac12\delta u^aD_a(\phi-\tau)
   +\frac14(\delta\tau-\delta\phi)D_n\tau.
 \label{eq:V73-canonical-face-one-form}
\end{align}
It obeys the exact tangential identity
\begin{align}
 \vartheta_{\rm can}
 +\Pi_{(1)}^{ab}\delta q_{ab}
 &=\frac12(\delta\tau-\delta\phi)D^au_a
   +\frac12\delta u^aD_a(\phi-\tau).
 \label{eq:V73-canonical-versus-BY}\\
 \int_F\left(
 \vartheta_{\rm can}+\Pi_{(1)}^{ab}\delta q_{ab}\right)
 &=\delta\int_F\frac12u^aD_a(\phi-\tau)
 \nonumber\\
 &\quad
 -\int_{\partial F}\frac12\nu_au^a\delta(\phi-\tau).
 \label{eq:V73-canonical-tangential-exact}
\end{align}
Thus the canonical free-\(A\) polarization is aligned with the opposite
Brown--York sign, up to a tangentially exact term and an explicit edge
one-form.  It is not the V36 GHY polarization.
\end{lemma}

\begin{proof}
The trace of the linearized connection is
\begin{equation}
 \Gamma^\alpha{}_{\mu\alpha}(H)=\frac12D_\mu H,
 \label{eq:V73-Gamma-trace}
\end{equation}
while
\begin{equation}
 \Gamma^n{}_{ab}=-K^{(1)}_{ab},\qquad
 \Gamma^n{}_{an}=\frac12D_a\phi,\qquad
 \Gamma^n{}_{nn}=\frac12D_n\phi.
 \label{eq:V73-Gamma-face-components}
\end{equation}
Insert these identities and
\begin{align*}
 \delta\overline H^{ab}
 &=\delta q^{ab}-\frac12\eta^{ab}(\delta\tau+\delta\phi),\\
 \delta\overline H^{an}&=\delta u^a,\\
 \delta\overline H^{nn}&=\frac12(\delta\phi-\delta\tau)
\end{align*}
into
\begin{equation}
 \delta_Hb_{\rm can}
 =\delta\overline H^{\mu\nu}\Gamma^n{}_{\mu\nu}
  -\delta\overline H^{\mu n}
       \Gamma^\alpha{}_{\mu\alpha}.
 \label{eq:V73-bcan-variation}
\end{equation}
This gives \eqref{eq:V73-canonical-face-one-form}.  Using
\(K_{(1)}=\tfrac12D_n\tau-D^au_a\) gives
\eqref{eq:V73-canonical-versus-BY}.  Tangential integration by parts gives
\eqref{eq:V73-canonical-tangential-exact}, including its sign at the edge.
\end{proof}

\begin{proposition}[A transverse-traceless sign witness]
\label{prop:V73-TT-boundary-witness}
At a face point take
\begin{equation}
 H=0,\qquad D_aH=0,\qquad u=\phi=\tau=0,
 \qquad K_{11}^{(1)}=k,\quad K_{22}^{(1)}=-k,
 \label{eq:V73-TT-jet}
\end{equation}
with every other component zero, and choose
\(\delta q_{11}=k\), \(\delta q_{22}=-k\).  Then
\begin{equation}
 \vartheta_{\rm can}=-2k^2,
 \qquad
 \Pi_{(1)}^{ab}\delta q_{ab}=2k^2.
 \label{eq:V73-TT-sign-witness}
\end{equation}
In particular no boundary density independent of \(A\), local in \(H\) and
its tangential jets, can change the canonical free-\(A\) face form into the
V36 Brown--York face form for all first jets.
\end{proposition}

\begin{proof}
The jet is tangentially constant and trace free, so
\(K_{(1)}=0\).  Equation~\eqref{eq:V73-canonical-face-one-form} leaves
\(-K^{ab}\delta q_{ab}=-2k^2\), whereas the Brown--York contraction has
the opposite sign.  At the chosen point an \(A\)-independent quadratic
intrinsic boundary density has zero first variation: every undifferentiated
factor \(H\) vanishes and every tangential jet vanishes.  Tangential
divergences also vanish locally.  It therefore cannot repair
\eqref{eq:V73-TT-sign-witness}.  A normal-jet term would change the boundary
order and is outside the asserted intrinsic class.
\end{proof}

\subsection{The genuine quadratic GHY representative}
\label{subsec:V73-GHY-representative}

Let \(F\) be any planar non-null face, let \(n^2=\sigma\in\{-1,1\}\), and
write \(\gamma\) for its background induced metric.  Define
\begin{equation}
 q_{ab}=e_a^\mu e_b^\nu H_{\mu\nu},\qquad
 \tau=\gamma^{ab}q_{ab},\qquad
 \phi=H_{\mu\nu}n^\mu n^\nu,
 \qquad
 A^\perp{}_{ab}=n_\alpha A^\alpha{}_{\mu\nu}e_a^\mu e_b^\nu.
 \label{eq:V73-general-face-data}
\end{equation}

\begin{lemma}[Quadratic GHY coefficient on every planar face]
\label{lem:V73-quadratic-GHY}
After the order-one background divergence and its GHY cancellation are
paired, the part of the order-two density \(2\sqrt{|\gamma|}K\) which is
bilinear in \(H\) and the order-one connection is
\begin{equation}
 \boxed{
 b_{\rm GHY}^{(2)}(H,A;n)
 =2q^{ab}A^\perp{}_{ab}
  -(\tau+\sigma\phi)\gamma^{ab}A^\perp{}_{ab}.}
 \label{eq:V73-GHY-quadratic-density}
\end{equation}
The formula applies to a timelike side \(\sigma=1\) and to either
oriented spacelike cap \(\sigma=-1\); the chosen outward conormal is already
contained in \(A^\perp\).
\end{lemma}

\begin{proof}
For a fixed planar hypersurface the normalized conormal varies by
\begin{equation}
 \delta n_\mu=\frac\sigma2\phi\,n_\mu,
 \label{eq:V73-normal-variation}
\end{equation}
the tangential inverse metric varies by \(-q^{ab}\), and the induced volume
varies by \(\tfrac12\tau\).  Since the anchor has zero extrinsic curvature,
the connection part of the first variation is
\(K^{(1)}=-\gamma^{ab}A^\perp{}_{ab}\).  Multiplying its three first-order
variations gives
\begin{align*}
 [2\sqrt{|\gamma|}K]_{H A}
 &=2q^{ab}A^\perp{}_{ab}
   -\tau\gamma^{ab}A^\perp{}_{ab}
   -\sigma\phi\gamma^{ab}A^\perp{}_{ab},
\end{align*}
which is \eqref{eq:V73-GHY-quadratic-density}.  Terms containing the
order-two connection cancel the order-two background Einstein divergence,
just as the order-one pair cancels; they do not enter the bilinear parent.
\end{proof}

\begin{proposition}[GHY does not permit arbitrary independent face connection
variation]
\label{prop:V73-GHY-free-A-obstruction}
For the bulk parent \eqref{eq:V73-Palatini-parent-density}, the coefficient
of an arbitrary face variation \(\delta A\) after adding
\eqref{eq:V73-GHY-quadratic-density} is
\begin{align}
 \mathfrak m_F(H;\delta A)
 &:=-n_\alpha\overline H^{\mu\nu}
       \delta A^\alpha{}_{\mu\nu}
   +n_\nu\overline H^{\mu\nu}
       \delta A^\alpha{}_{\mu\alpha}
   +\delta_A b_{\rm GHY}^{(2)}(H,A;n).
 \label{eq:V73-GHY-A-mismatch}
\end{align}
This covector is not identically zero.  On the V36 side face, take
\(q=u=0\), \(\phi\ne0\), and vary only
\(\delta A^n{}_{00}\).  Then
\begin{equation}
 \boxed{
 \mathfrak m_F(H;\delta A)
 =\frac12\phi\,\delta A^n{}_{00}\ne0.}
 \label{eq:V73-GHY-A-witness}
\end{equation}
Moreover, among face densities affine in \(A\), cancellation of arbitrary
\(\delta A\) uniquely selects \(b_{\rm can}\) up to an \(A\)-independent
density.  Consequently the GHY and free-connection polarizations cannot be
identified.
\end{proposition}

\begin{proof}
The first two terms of \eqref{eq:V73-GHY-A-mismatch} are the Stokes term
from varying the two derivatives of \(A\) in the bulk parent.  For the
displayed pure-normal \(H\),
\(\overline H^{00}=\phi/2\), so the bulk contribution to
\(\delta A^n{}_{00}\) is \(-\phi/2\).  Equation
\eqref{eq:V73-GHY-quadratic-density} contributes
\(-\phi\eta^{00}\delta A^n{}_{00}=\phi\delta A^n{}_{00}\).
Their sum is \eqref{eq:V73-GHY-A-witness}.

If a density \(b(H,A)\) cancels arbitrary \(\delta A\), its derivative with
respect to every independent component of \(A\) must be the negative of the
bulk Stokes coefficient.  Integration in the affine \(A\)-space gives
\(b=b_{\rm can}+f(H,D_{\!\parallel}H,\ldots)\).  Proposition
\ref{prop:V73-TT-boundary-witness} excludes the required conversion by such
an intrinsic \(f\).  This proves both assertions.
\end{proof}

\begin{firewall}[The obstruction is a boundary-polarization statement]
\label{fire:V73-boundary-polarization-scope}
Propositions~\ref{prop:V73-TT-boundary-witness} and
\ref{prop:V73-GHY-free-A-obstruction} do not obstruct Einstein gravity or the
GHY action.  They show that one cannot simultaneously demand the V36 GHY
representative and arbitrary independent connection variations on the face
without adding a compatibility relation or additional face cotangent data.
The familiar metric derivation uses the compatibility relation
\(A=\Gamma(H)\) at the boundary and therefore never makes the forbidden
identification.
\end{firewall}

\subsection{Exact compatible-jet repair and reduction to V36}
\label{subsec:V73-compatible-repair}

Let \(r_{\lambda\mu\nu}=r_{\lambda\nu\mu}\) be the complete independent
metric first jet retained in V71 and define the algebraic Christoffel map
\begin{equation}
 \Gamma^\alpha{}_{\mu\nu}(r)
 :=\frac12\eta^{\alpha\rho}
 (r_{\mu\rho\nu}+r_{\nu\rho\mu}-r_{\rho\mu\nu}).
 \label{eq:V73-Christoffel-jet-map}
\end{equation}
On a face introduce a multiplier
\(\Lambda_\alpha{}^{\mu\nu}=\Lambda_\alpha{}^{\nu\mu}\) and the pairing
\begin{equation}
 b_{\rm comp}(A,r,\Lambda)
 :=\Lambda_\alpha{}^{\mu\nu}
 \left(A^\alpha{}_{\mu\nu}
       -\Gamma^\alpha{}_{\mu\nu}(r)\right).
 \label{eq:V73-compatible-boundary-pairing}
\end{equation}
The inherited V71 jet-cotangent relation imposes \(r=DH\) without discarding
its boundary momentum.

\begin{theorem}[Free enlarged parent with exact Einstein--GHY reduction]
\label{thm:V73-enlarged-GHY-parent}
Consider
\begin{equation}
 S_{73}^{(2)}[H,A,r,\Lambda]
 :=\frac1{2\kappa}
 \left{
 \int_\Omega\mathcal L_{\rm Pal}^{(2)}(H,A)
 +\sum_F\int_F
 \left(b_{\rm GHY}^{(2)}(H,A;n_F)
       +b_{\rm comp}(A,r,\Lambda)\right)
 \right},
 \label{eq:V73-enlarged-parent-action}
\end{equation}
together with the V71 complete-jet relation \(r=DH\).  Then:
\begin{enumerate}[label=\textup{(\alph*)},leftmargin=2.8em]
\item \(\delta\Lambda\) gives \(A=\Gamma(r)\) on every face;
\item \(\delta A\) is arbitrary on every face and gives
\begin{equation}
 \Lambda=-\mathfrak m_F(H),
 \label{eq:V73-Lambda-elimination}
\end{equation}
where \(\mathfrak m_F\) is the explicit covector in
\eqref{eq:V73-GHY-A-mismatch};
\item the bulk connection equation gives \(A=\Gamma(H)\), and the jet
      equation makes the face constraint identical to its trace;
\item eliminating \(A,r,\Lambda\) gives exactly the flat-anchor quadratic
      coefficient of the Einstein--Hilbert plus GHY action on every smooth
      face;
\item after adjoining the V36 de Donder term
      \(-\frac1{2\kappa}\int C_\mu(H)C^\mu(H)\), the smooth-face one-form is
\begin{equation}
 \boxed{
 \Theta_{73,F}
 =\frac1{2\kappa}\int_F
 \left[
 \Pi_{(1)}^{ab}(H)\,\delta q_{ab}
 -2\mathcal Q^{\mu\nu}(C(H))\,\delta H_{\mu\nu}
 \right],}
 \label{eq:V73-exact-V36-reduction}
\end{equation}
which is the V36 conormal with the actual orientation of \(F\).
\end{enumerate}
\end{theorem}

\begin{proof}
Items (a) and (b) follow by varying the two factors in
\eqref{eq:V73-compatible-boundary-pairing}; the remaining coefficient of
\(\delta A\) is exactly \eqref{eq:V73-GHY-A-mismatch}.  Theorem
\ref{thm:V73-Palatini-elimination} and the V71 jet relation give (c).

For (d), expand the covariant Einstein--Hilbert integrand through order two.
Equation~\eqref{eq:V73-densitized-inverse-expansion}, the Ricci decomposition,
and cancellation of the order-one background divergence give precisely
\eqref{eq:V73-Palatini-parent-density}.  Lemma
\ref{lem:V73-quadratic-GHY} is precisely the corresponding order-two GHY
coefficient.  The multiplier imposes, but does not alter, the graph
\(A=\Gamma(DH)\).  Elementary constrained elimination therefore returns the
quadratic coefficient of the original covariant action.  No integration by
parts to the minimal wave representative is used.

The GHY variation of that reduced action is the Brown--York term proved in
Theorem~\ref{thm:V36-GHY-conormal}.  The de Donder functional is independent
of the connection variables and contributes the already computed
\(-2\mathcal Q(C)\) coefficient.  Their sum proves
\eqref{eq:V73-exact-V36-reduction}.
\end{proof}

\begin{corollary}[Residual gauge compatibility]
\label{cor:V73-residual-compatible-parent}
For the V72 residual transformation
\begin{equation}
 (\mathcal G\varepsilon)_{\mu\nu}
 =D_\mu\varepsilon_\nu+D_\nu\varepsilon_\mu,
 \label{eq:V73-residual-metric-shift}
\end{equation}
the induced connection shift is
\begin{equation}
 \Gamma^\alpha{}_{\mu\nu}(\mathcal G\varepsilon)
 =D_\mu D_\nu\varepsilon^\alpha.
 \label{eq:V73-residual-connection-shift}
\end{equation}
Hence the compatible-jet face relation is preserved by the residual complex.
On the ghost shell \(\Box\varepsilon=0\), eliminating the parent leaves the
V72 residual moment map and its corner cocycle unchanged.
\end{corollary}

\begin{proof}
Insert \eqref{eq:V73-residual-metric-shift} into
\eqref{eq:V73-Levi-Civita-elimination}.  Commutation of flat derivatives
cancels the two crossed terms and leaves
\eqref{eq:V73-residual-connection-shift}.  The multiplier relation is
therefore equivariant.  Constrained elimination changes neither the reduced
symplectic form nor the reduced residual fundamental vector field, so the
moment map and the cocycle computed from them in V72 are preserved.
\end{proof}

\begin{firewall}[Smooth faces do not solve the joint master equation]
\label{fire:V73-corner-still-open}
Equation~\eqref{eq:V73-exact-V36-reduction} is exact on every smooth face.
For a piecewise smooth slab, variations of the two face normals meet at a
joint.  The quadratic Hayward or boost-angle coefficient, its ghost
completion, and its comparison with the nonzero V72 cocycle have not been
computed here.  The new multiplier
\(\Lambda=-\mathfrak m_F(H)\) identifies the missing face cotangent datum, but
does not by itself cancel the V72 corner cocycle.  Edge-zero residual
parameters remain a valid nilpotent domain; unrestricted residual parameters
still require a corner extension.
\end{firewall}

\subsection{What V73 closes and the next acquisition}
\label{subsec:V73-ledger}

\paragraph{New exact conclusions.}
\begin{enumerate}[label=\textup{(\arabic*)},leftmargin=2.8em]
\item The torsion-free Palatini quadratic parent
      \eqref{eq:V73-Palatini-parent-density} has a unique algebraic connection
      equation and eliminates exactly to the Levi--Civita jet.
\item The boundary density which permits arbitrary independent connection
      variations is \eqref{eq:V73-canonical-transgression}; its complete face
      one-form is \eqref{eq:V73-canonical-face-one-form}.
\item The transverse-traceless witness
      \eqref{eq:V73-TT-sign-witness} proves that this canonical polarization is
      not the V36 GHY polarization.
\item The genuine quadratic GHY coefficient is
      \eqref{eq:V73-GHY-quadratic-density} on both sides and caps.  The scalar
      witness \eqref{eq:V73-GHY-A-witness} proves that it does not permit
      arbitrary independent face-connection variation.
\item The compatible-jet multiplier
      \eqref{eq:V73-compatible-boundary-pairing} resolves the conflict.  Its
      elimination returns the exact Einstein--GHY quadratic action, and adding
      de Donder gauge fixing returns the full V36 conormal.
\item The construction is equivariant under the V72 residual differential and
      therefore transports, rather than erases, the residual corner cocycle.
\end{enumerate}

\paragraph{Next forced step.}
The first open local row is now sharply delimited.  On a two-face joint one
must expand the oriented normal-pair angle and induced joint volume to second
order, express the result on the compatible jet graph, and compute its
field-space exterior derivative.  That joint two-form must then be compared
term by term with \(K_{72}\).  If the geometric joint variables do not cancel
that cocycle, one must adjoin the smallest ghost corner cotangent pair and
solve the linear corner master equation.  Only after this joint calculation
should the parent be inserted into the global reduced metric-continuum
Lagrangian relation.

\begin{assessment}[V73 continuum status]
The full Standard--Model--Einstein continuum remains unproved.  V73 closes
the previously missing local torsion-free connection ancestry and gives an
exact enlarged face parent whose elimination reproduces the V36
Einstein--GHY--de Donder conormal.  It also proves why the naive free-connection
and GHY boundary prescriptions could not be merged.  The piecewise-slab joint
action, unrestricted residual corner master equation, nonlinear curved-anchor
extension, common renormalized Standard Model stress domain, and global
cofinal continuum comparison remain open.
\end{assessment}

\section*{V73 exact checks and append-only record}
\addcontentsline{toc}{section}{V73 exact checks and append-only record}

The forty independent components of a torsion-free four-dimensional
connection were enumerated in the basis
\(A^\alpha{}_{\mu\nu}\), \(\mu\le\nu\).  A direct component
assembly of the quadratic density in
\eqref{eq:V73-algebraic-parent} has full rank forty.  On a fixed seeded test
jet, solving its Euler system gave
\begin{equation}
 \max_{\alpha,\mu,\nu}
 \left|A^\alpha{}_{\mu\nu}
       -\Gamma^\alpha{}_{\mu\nu}(H)\right|
 =2.23\times10^{-16}.
 \label{eq:V73-numerical-connection-check}
\end{equation}
This numerical check is not used as a proof; it audits the signs and index
placements of the analytic derivation in
Theorem~\ref{thm:V73-Palatini-elimination}.  The two independent boundary
witnesses \eqref{eq:V73-TT-sign-witness} and
\eqref{eq:V73-GHY-A-witness} were also evaluated directly from the component
formulas.

For the installed manuscript, the V72 byte prefix was retained exactly up to
its sole live final document terminator.  The V73 material was appended after
that prefix, a single live \verb|\end{document}| was restored, and the
candidate was checked for duplicate V73 labels, balanced environments,
balanced braces outside comments, and exact predecessor-prefix equality
before the predecessor was removed.

% =============================================================================
% END APPENDED VERSION V73 MATERIAL
% =============================================================================

% =============================================================================
% BEGIN APPENDED VERSION V74 MATERIAL
% =============================================================================
\clearpage
\section*{V74 research note: the quadratic side--cap joint, the failure of a
geometric scalar to cancel the residual cocycle, and the minimal corner
cotangent extension}
\addcontentsline{toc}{section}{V74 quadratic joint and corner cotangent extension}
\label{sec:V74-joint-corner-extension}

\subsection*{Purpose and exact relation to V72--V73}

V73 derived the local torsion-free Palatini ancestor and the compatible-jet
face polarization which eliminates to the V36 Einstein--GHY--de Donder
conormal.  Its first unresolved local row was the codimension-two joint.
V72 had already proved that unrestricted residual wave gauges carry the
nonzero central form \(\mathcal K_{72}\).  The two questions are logically
distinct:
\begin{enumerate}[label=\textup{(\roman*)},leftmargin=2.8em]
\item what is the actual quadratic Hayward or boost-angle coefficient at an
      orthogonal side--cap joint?
\item can its field-space variation cancel \(\mathcal K_{72}\), and if not,
      what is the smallest local corner phase which can?
\end{enumerate}
This note answers both questions.  The joint coefficient is computed directly
from normalized conormals.  Its Hessian is symmetric, hence its field-space
curvature vanishes.  It cannot cancel the antisymmetric V72 cocycle.  A
three-component local cotangent presentation is then constructed, and its
canonical quotient is proved minimal at every finite jet or Fourier fibre.

Before appending this note, a strengthened source audit found malformed inline
math openings in the newly installed V73 block and four older doubled openings.
They were repaired without changing any formula or prose.  The corrected V73
source, whose digest is recorded below, is the exact byte prefix of V74.  The
audit now counts opening and closing inline delimiters in addition to the
environment, brace, label, reference, and document-terminator checks.

No companion-program audit is due in V74.  The next mandatory comparison with
the newest Yang--Mills and Navier--Stokes notes remains V75.

\subsection{Normalized normals and the exact quadratic joint coefficient}
\label{subsec:V74-joint-expansion}

Fix one spacelike codimension-two joint \(J\) at the intersection of a
timelike side and a spacelike cap.  Let \(N\) be the outward spacelike unit
normal to the side and \(U\) the chosen outward timelike unit normal to the
cap:
\begin{equation}
 N^2=1,\qquad U^2=-1,\qquad N\mathbin{\cdot}U=0.
 \label{eq:V74-background-normals}
\end{equation}
Let \(E_A\), \(A=1,2\), be an orthonormal frame tangent to \(J\).  Along the
straight metric path
\begin{equation}
 g_\varepsilon=\eta+\varepsilon H
 \label{eq:V74-metric-path}
\end{equation}
retain the two hypersurfaces as fixed level sets and normalize their conormals
with \(g_\varepsilon\).  Write those conormals as
\(N^\flat_\varepsilon\) and \(U^\flat_\varepsilon\), and put
\begin{equation}
 z_\varepsilon
 :=g_\varepsilon^{-1}
   (N^\flat_\varepsilon,U^\flat_\varepsilon),
 \qquad
 \chi_\varepsilon:=\operatorname{arsinh}z_\varepsilon.
 \label{eq:V74-boost-definition}
\end{equation}
Changing the joint orientation changes the overall sign of \(\chi\); it
changes none of the exactness or cocycle conclusions below.

Use the component abbreviations
\begin{equation}
 a:=H_{NN},\quad b:=H_{UU},\quad c:=H_{NU},\quad
 v_A:=H_{NA},\quad w_A:=H_{UA},\quad
 \rho:=H_A{}^A.
 \label{eq:V74-joint-components}
\end{equation}

\begin{lemma}[Second-order normalized-normal expansion]
\label{lem:V74-normal-expansion}
At the orthogonal flat anchor,
\begin{align}
 z_\varepsilon
 &=-\varepsilon c
   +\varepsilon^2
    \left[
      v_Aw^A+\frac12(a-b)c
    \right]
   +O(\varepsilon^3),
 \label{eq:V74-z-expansion}\\
 \chi_\varepsilon
 &=-\varepsilon c
   +\varepsilon^2
    \left[
      v_Aw^A+\frac12(a-b)c
    \right]
   +O(\varepsilon^3),
 \label{eq:V74-chi-expansion}\\
 \sqrt{\det\sigma_\varepsilon}
 &=1+\frac{\varepsilon}{2}\rho+O(\varepsilon^2),
 \label{eq:V74-joint-volume-expansion}
\end{align}
where \(\sigma_\varepsilon\) is the metric induced on \(J\).
\end{lemma}

\begin{proof}
The inverse metric has the Neumann expansion
\begin{equation}
 g_\varepsilon^{-1}
 =\eta^{-1}-\varepsilon H^\sharp
  +\varepsilon^2H^\sharp\mathbin{\circ}H^\sharp
  +O(\varepsilon^3).
 \label{eq:V74-inverse-expansion}
\end{equation}
Normalization of the fixed side and cap conormals gives, to the order needed,
\begin{align}
 N^\flat_\varepsilon
 &=\left(1+\frac{\varepsilon}{2}a\right)N^\flat
   +O(\varepsilon^2),
 \label{eq:V74-side-normalization}\\
 U^\flat_\varepsilon
 &=\left(1-\frac{\varepsilon}{2}b\right)U^\flat
   +O(\varepsilon^2).
 \label{eq:V74-cap-normalization}
\end{align}
The first inverse-metric correction contributes \(-c\).  The second
inverse-metric correction is
\begin{equation}
 H_{N\lambda}H^\lambda{}_U
 =(a-b)c+v_Aw^A.
 \label{eq:V74-H-square-cross}
\end{equation}
Multiplication by the two normalization factors subtracts
\(\tfrac12(a-b)c\), proving \eqref{eq:V74-z-expansion}.  Since
\(\operatorname{arsinh}z=z+O(z^3)\), there is no quadratic correction from
the inverse hyperbolic sine, which proves \eqref{eq:V74-chi-expansion}.
The standard determinant derivative on the two-dimensional joint gives
\eqref{eq:V74-joint-volume-expansion}.
\end{proof}

\begin{theorem}[Exact quadratic Hayward density at an orthogonal joint]
\label{thm:V74-quadratic-Hayward}
Let the signed non-null joint action be
\begin{equation}
 S_{\rm joint}[g]
 :=\frac1\kappa\sum_J\varsigma_J
   \int_J\sqrt{\det\sigma}\,\chi,
 \qquad \varsigma_J\in\{-1,1\}.
 \label{eq:V74-Hayward-action}
\end{equation}
Its coefficient of order \(\varepsilon^2\) along
\eqref{eq:V74-metric-path} is
\begin{equation}
 \boxed{
 S_{\rm joint}^{(2)}[H]
 =\frac1\kappa\sum_J\varsigma_J\int_J
 \left[
 H_{NA}H_U{}^A
 +\frac12
  \left(H_{NN}-H_{UU}-H_A{}^A\right)H_{NU}
 \right].}
 \label{eq:V74-quadratic-joint-action}
\end{equation}
This is the missing side--cap coefficient in the V73 piecewise-slab action.
\end{theorem}

\begin{proof}
Multiply \eqref{eq:V74-chi-expansion} by
\eqref{eq:V74-joint-volume-expansion}.  The order-two coefficient is
\begin{equation}
 v_Aw^A+\frac12(a-b)c-\frac12\rho c,
 \end{equation}
which is the integrand in \eqref{eq:V74-quadratic-joint-action}.  The signs
\(\varsigma_J\) implement the Stokes orientation at the initial and final
joints.
\end{proof}

For two metric perturbations \(H,K\), define the polarization
\begin{align}
 \mathfrak j_{74}(H,K)
 &:=\frac12
 \left(H_{NA}K_U{}^A+K_{NA}H_U{}^A\right)
 \nonumber\\
 &\quad
 +\frac14
 \left[
  (H_{NN}-H_{UU}-H_A{}^A)K_{NU}
 +(K_{NN}-K_{UU}-K_A{}^A)H_{NU}
 \right].
 \label{eq:V74-joint-polarization}
\end{align}
It is symmetric and satisfies
\(\mathfrak j_{74}(H,H)\) equal to the integrand of
\eqref{eq:V74-quadratic-joint-action}.

\subsection{Why the geometric joint scalar cannot cancel
\(\mathcal K_{72}\)}
\label{subsec:V74-Hayward-no-cocycle}

\begin{theorem}[The Hayward corner shift has zero field-space curvature]
\label{thm:V74-Hayward-exact}
Let
\begin{equation}
 \Theta_{\rm joint}^{(2)}:=\delta S_{\rm joint}^{(2)}.
 \label{eq:V74-Hayward-potential}
\end{equation}
Then
\begin{equation}
 \boxed{
 \omega_{\rm joint}^{(2)}
 :=\delta\Theta_{\rm joint}^{(2)}=0.}
 \label{eq:V74-Hayward-two-form-zero}
\end{equation}
For residual gauge vectors \(\mathcal V_\epsilon,\mathcal V_\zeta\),
\begin{equation}
 \omega_{\rm joint}^{(2)}
   (\mathcal V_\epsilon,\mathcal V_\zeta)=0.
 \label{eq:V74-Hayward-residual-zero}
\end{equation}
Consequently the geometric joint scalar cannot cancel the nonzero V72 form
\(\mathcal K_{72}\).
\end{theorem}

\begin{proof}
The Hessian of \eqref{eq:V74-quadratic-joint-action} is twice the symmetric
polarization \eqref{eq:V74-joint-polarization}.  Hence, for commuting
variations \(K,L\),
\begin{equation}
 \delta_K\delta_L S_{\rm joint}^{(2)}
 -\delta_L\delta_K S_{\rm joint}^{(2)}=0.
 \label{eq:V74-commuting-Hessian}
\end{equation}
This is precisely \eqref{eq:V74-Hayward-two-form-zero}.  Residual metric
gauge transformations form an abelian linear algebra, so their fundamental
vectors commute and \eqref{eq:V74-Hayward-residual-zero} follows.
The V72 witness \eqref{eq:V72-cocycle-witness} has
\(\mathcal K_{72}\ne0\).  A zero two-form cannot equal
\(-\mathcal K_{72}\).
\end{proof}

\begin{corollary}[No scalar joint counterterm solves the unrestricted master
equation]
\label{cor:V74-no-scalar-corner-repair}
Adding any local scalar joint functional \(B_J(H,DH,\ldots)\) to the action
changes the corner potential by \(\delta B_J\) and leaves its field-space
two-form unchanged.  It therefore cannot alter the central term
\eqref{eq:V72-corner-cocycle}.  In particular, no choice of boost-angle
reference subtraction can solve the unrestricted residual master equation.
\end{corollary}

\begin{proof}
Nilpotence of the field-space differential gives
\(\delta^2B_J=0\).  Equivalently, for an abelian residual algebra with
constant central coefficients, the Chevalley--Eilenberg differential of a
scalar-action shift has zero antisymmetric polarization.  Apply
Theorem~\ref{thm:V74-Hayward-exact}.
\end{proof}

\begin{firewall}[Variational differentiability and cocycle cancellation are
different rows]
\label{fire:V74-joint-scope}
The joint action \eqref{eq:V74-quadratic-joint-action} is required for the
piecewise Einstein--GHY variational problem.  Its failure to cancel
\(\mathcal K_{72}\) does not make it optional.  Conversely, a new corner
symplectic phase may cancel the residual cocycle but does not replace the
geometric boost-angle action.  Both rows must be retained.
\end{firewall}

\subsection{A local corner cotangent phase with the opposite cocycle}
\label{subsec:V74-corner-cotangent}

On each oriented edge \(J\subset\partial\Gamma_T\), define the two trace maps
on residual wave parameters:
\begin{align}
 \mathsf T_J(\epsilon)_b
 &:=\epsilon_b|_J,
 \label{eq:V74-tangential-trace}\\
 \mathsf R_J(\epsilon)^b
 &:=\left.\nu_a\mathscr P^{ab}(\epsilon_n)\right|_J.
 \label{eq:V74-normal-Hessian-trace}
\end{align}
Then V72 can be written
\begin{equation}
 \mathcal K_{72}(\epsilon,\zeta)
 =\frac1\kappa\sum_J\int_J
 \left[
  \mathsf R_J(\epsilon)^b\mathsf T_J(\zeta)_b
 -\mathsf R_J(\zeta)^b\mathsf T_J(\epsilon)_b
 \right].
 \label{eq:V74-K72-trace-form}
\end{equation}

Introduce even corner fields \(X_b,\Pi^b\) with canonical form
\begin{equation}
 \omega_{\rm cor}
 \bigl((X_1,\Pi_1),(X_2,\Pi_2)\bigr)
 :=\frac1\kappa\sum_J\int_J
 \left(\Pi_1^bX_{2,b}-X_{1,b}\Pi_2^b\right).
 \label{eq:V74-corner-symplectic-form}
\end{equation}
Define the residual corner translation by
\begin{equation}
 \boxed{
 \mathcal V_\epsilon^{\rm cor}X_b
 =\mathsf T_J(\epsilon)_b,\qquad
 \mathcal V_\epsilon^{\rm cor}\Pi^b
 =-\mathsf R_J(\epsilon)^b.}
 \label{eq:V74-corner-residual-action}
\end{equation}

\begin{theorem}[Exact opposite corner cocycle and corner moment map]
\label{thm:V74-opposite-cocycle}
The corner translations are Hamiltonian with
\begin{equation}
 \mu_\epsilon^{\rm cor}(X,\Pi)
 :=-\frac1\kappa\sum_J\int_J
 \left[
  \mathsf R_J(\epsilon)^bX_b
 +\mathsf T_J(\epsilon)_b\Pi^b
 \right].
 \label{eq:V74-corner-moment-map}
\end{equation}
They satisfy
\begin{align}
 \iota_{\mathcal V_\epsilon^{\rm cor}}\omega_{\rm cor}
 &=\delta\mu_\epsilon^{\rm cor},
 \label{eq:V74-corner-Hamiltonian-identity}\\
 \omega_{\rm cor}
  (\mathcal V_\epsilon^{\rm cor},
   \mathcal V_\zeta^{\rm cor})
 &=-\mathcal K_{72}(\epsilon,\zeta).
 \label{eq:V74-opposite-corner-cocycle}
\end{align}
The construction is local and uses at most the two tangential derivatives
already present in \(\mathscr P\).
\end{theorem}

\begin{proof}
For an arbitrary variation \((\delta X,\delta\Pi)\),
\begin{align}
 \omega_{\rm cor}
  (\mathcal V_\epsilon^{\rm cor},(\delta X,\delta\Pi))
 &=
 \frac1\kappa\sum_J\int_J
 \left[
  -\mathsf R_J(\epsilon)^b\delta X_b
  -\mathsf T_J(\epsilon)_b\delta\Pi^b
 \right],
\end{align}
which is the differential of \eqref{eq:V74-corner-moment-map}.  Evaluating
on \(\mathcal V_\zeta^{\rm cor}\) gives
\begin{align}
 \frac1\kappa\sum_J\int_J
 \left[
  -\mathsf R_J(\epsilon)^b\mathsf T_J(\zeta)_b
  +\mathsf T_J(\epsilon)_b\mathsf R_J(\zeta)^b
 \right],
\end{align}
the negative of \eqref{eq:V74-K72-trace-form}.  Locality is immediate from
\eqref{eq:V74-normal-Hessian-trace}.
\end{proof}

\begin{theorem}[Fibrewise minimality after removal of null directions]
\label{thm:V74-corner-minimality}
Fix a finite tangential jet fibre, or one tangential Fourier fibre, at \(J\).
Let
\begin{equation}
 \mathscr W_J
 :=\operatorname{span}
 \left\{
  \bigl(\mathsf T_J(\epsilon),-\mathsf R_J(\epsilon)\bigr)
  :\epsilon\ \text{is a residual jet}
 \right\}
 \label{eq:V74-corner-translation-span}
\end{equation}
and quotient \(\mathscr W_J\) by the radical of the restriction of
\(\omega_{\rm cor}\).  The resulting symplectic carrier has dimension
\begin{equation}
 \dim\mathscr W_J^{\rm red}
 =\operatorname{rank}\mathcal K_{72,J}.
 \label{eq:V74-minimal-rank}
\end{equation}
Every symplectic corner realization whose residual fundamental vectors have
pairing \(-\mathcal K_{72,J}\) has dimension at least this rank.  Hence the
reduced carrier obtained from \((X_b,\Pi^b)\) is minimal.  The unreduced
three-pair presentation is a convenient local differential realization of
that minimal quotient.
\end{theorem}

\begin{proof}
By Theorem~\ref{thm:V74-opposite-cocycle}, the matrix of
\(\omega_{\rm cor}|_{\mathscr W_J}\) in any spanning family of residual
translations is \(-\mathcal K_{72,J}\).  Quotienting by its radical produces
a nondegenerate skew form whose dimension is the rank of that matrix.  In any
other symplectic realization, the span of the residual fundamental vectors
has the same restricted skew matrix.  Its ambient symplectic space must
contain a subspace of dimension at least the rank.  The quotient constructed
above attains the bound.
\end{proof}

\begin{firewall}[Meaning of minimality]
\label{fire:V74-minimality-scope}
Theorem~\ref{thm:V74-corner-minimality} is a fibrewise linear symplectic
minimality statement.  It does not prove a preferred nonlinear corner
quantization, a global polarization over a curved joint, or a positive corner
Hamiltonian.  It proves that no smaller linear symplectic carrier can realize
the required opposite bilinear form after null directions are removed.
\end{firewall}

\subsection{The unrestricted residual BFV charge}
\label{subsec:V74-unrestricted-BFV}

Let \(\gamma\) be the odd residual ghost of V72, still subject in this note to
the residual wave equation.  Extend its differential to the corner variables
by
\begin{equation}
 q_{\rm res}X_b=\gamma_b|_J,\qquad
 q_{\rm res}\Pi^b
 =-\nu_a\mathscr P^{ab}(\gamma_n)|_J,\qquad
 q_{\rm res}\gamma=0.
 \label{eq:V74-corner-BRST-differential}
\end{equation}
Flatness and \(q_{\rm res}\gamma=0\) give \(q_{\rm res}^2=0\) on the corner
coordinates.

On the product phase
\begin{equation}
 (\mathcal P_{36},\omega_{36})
 \times(\mathcal P_{\rm cor},\omega_{\rm cor}),
 \label{eq:V74-product-residual-phase}
\end{equation}
put
\begin{equation}
 \widetilde\mu_\epsilon
 :=\mu_\epsilon+\mu_\epsilon^{\rm cor}.
 \label{eq:V74-equivariant-total-moment-map}
\end{equation}

\begin{theorem}[Unrestricted equivariant moment map and residual master
equation]
\label{thm:V74-unrestricted-master}
For arbitrary smooth residual wave parameters, with no edge-zero condition,
\begin{align}
 \{\widetilde\mu_\epsilon,\widetilde\mu_\zeta\}
 &=0,
 \label{eq:V74-total-moment-commutation}\\
 \Omega_{74}^{\rm res}
 &:=\widetilde\mu_\gamma,
 \label{eq:V74-unrestricted-BFV-charge}\\
 \boxed{
 \{\Omega_{74}^{\rm res},\Omega_{74}^{\rm res}\}=0.}
 \label{eq:V74-unrestricted-master-equation}
\end{align}
The charge generates the V72 metric and conormal residual transformations
together with \eqref{eq:V74-corner-BRST-differential}.
\end{theorem}

\begin{proof}
The product bracket is the sum of the two factor brackets.  The V36 factor is
\(\mathcal K_{72}(\epsilon,\zeta)\) by
Theorem~\ref{thm:V72-corner-cocycle}; the corner factor is its negative by
Theorem~\ref{thm:V74-opposite-cocycle}.  This proves
\eqref{eq:V74-total-moment-commutation}.  The residual algebra is abelian, so
polarization of the BFV self-bracket is precisely that moment-map bracket.
It vanishes for every pair of ghost jets, proving
\eqref{eq:V74-unrestricted-master-equation}.  The two Hamiltonian identities
give the generated transformations.
\end{proof}

Let \(\Omega_{70}^{\rm hol}\) denote the independent holonomy charge retained
in V72.

\begin{corollary}[Unrestricted product with the holonomy factor]
\label{cor:V74-total-linear-charge}
On the product of the V70 holonomy phase, the V36 metric phase, and the V74
corner phase,
\begin{equation}
 \Omega_{74}^{\rm tot}
 :=\Omega_{70}^{\rm hol}+\Omega_{74}^{\rm res},
 \qquad
 \boxed{
 \{\Omega_{74}^{\rm tot},\Omega_{74}^{\rm tot}\}=0}
 \label{eq:V74-total-linear-master}
\end{equation}
for unrestricted residual wave ghosts at the edge.
\end{corollary}

\begin{proof}
The holonomy charge has zero self-bracket by V70.  It is supported on the
contractible embedding-jet variables, whereas the new corner charge is
supported on \((X,\Pi,\gamma)\) and the V36 variables.  The product form has
zero cross brackets.  Apply
Theorem~\ref{thm:V74-unrestricted-master}.
\end{proof}

\begin{corollary}[No new bulk radiative class]
\label{cor:V74-no-bulk-mode}
The variables \((X,\Pi)\) are supported only on the codimension-two joint and
carry an affine residual action.  They add no solution of the bulk wave
operator and no class to the V59 two-polarization bulk quotient.  Their role
is to make the unrestricted edge action Hamiltonian and equivariant.
\end{corollary}

\begin{proof}
Equations \eqref{eq:V74-corner-residual-action} and
\eqref{eq:V74-corner-BRST-differential} contain no normal evolution equation
and no bulk kinetic operator.  The corner phase is absent from the bulk
solution complex.  Its only coupling in
\eqref{eq:V74-equivariant-total-moment-map} is through boundary traces of the
residual parameter.
\end{proof}

\subsection{Reduction order after the joint calculation}
\label{subsec:V74-reduction-order}

\begin{theorem}[Updated local linear reduction order]
\label{thm:V74-updated-reduction-order}
The unrestricted local linear Einstein boundary reduction may now proceed in
the following order:
\begin{enumerate}[label=\textup{(\arabic*)},leftmargin=2.8em]
\item use the V73 compatible-jet Palatini parent with the V74 quadratic joint
      scalar on every side--cap intersection;
\item take the product with the V74 corner cotangent phase and impose the
      equivariant constraints \(\widetilde\mu_\epsilon=0\);
\item quotient the unrestricted residual Hamiltonian action generated by
      \(\widetilde\mu\);
\item identify the reduced harmonic metric classes with the V53/V59
      two-polarization quotient at the operator level;
\item only then pull back the positive V59 continuum relation and take its
      product with the V70 holonomy reduction.
\end{enumerate}
The corner-cocycle obstruction no longer forces the tangential edge-zero
domain.  The remaining local gate before item~\textup{(3)} is an off-shell
ghost--antighost action and a closed analytic trace domain for the maps
\(\mathsf T_J,\mathsf R_J\).
\end{theorem}

\begin{proof}
V73 proves the first-order ancestry and smooth-face reduction.  Theorem
\ref{thm:V74-quadratic-Hayward} supplies its missing joint scalar.
Theorem~\ref{thm:V74-unrestricted-master} makes the unrestricted residual
action equivariant, so symplectic reduction is no longer obstructed by
\(\mathcal K_{72}\).  V71 proves that the metric quotient must be taken at
the operator level rather than by a local two-scalar splitting.  V72 proves
that this reduction precedes the V59 pullback and the holonomy product.  These
facts give the stated order.
\end{proof}

\begin{firewall}[On-shell scope of the new master equation]
\label{fire:V74-master-scope}
Equation~\eqref{eq:V74-unrestricted-master-equation} is the complete
unrestricted corner cancellation for the linear residual boundary algebra,
but the ghost still obeys \(\Box\gamma=0\).  V74 has not constructed the
off-shell bulk antighost, Nakanishi--Lautrup field, their face conditions, or
the joint terms required by their variational problem.  It therefore does not
claim the full off-shell BV--BFV master equation.
\end{firewall}

\subsection{V74 ledger and next acquisition}
\label{subsec:V74-ledger}

\paragraph{New exact conclusions.}
\begin{enumerate}[label=\textup{(\arabic*)},leftmargin=2.8em]
\item The normalized side and cap conormals give the exact boost expansion
      \eqref{eq:V74-chi-expansion}.
\item The quadratic Hayward coefficient is
      \eqref{eq:V74-quadratic-joint-action}.
\item Its field-space two-form vanishes, so neither it nor any scalar corner
      counterterm can cancel the V72 central cocycle.
\item The local corner cotangent action
      \eqref{eq:V74-corner-residual-action} has symplectic cocycle exactly
      \(-\mathcal K_{72}\).
\item After quotienting null directions, its dimension equals the rank of the
      V72 cocycle and is minimal among linear symplectic realizations.
\item The combined residual moment map is equivariant and its BFV charge
      satisfies the unrestricted edge master equation on the residual
      wave-ghost shell.
\item The independent V70 holonomy charge continues to commute with this
      enlarged residual factor.
\end{enumerate}

\paragraph{Next forced step.}
V75 must first perform the scheduled comparison with the newest Yang--Mills
and Navier--Stokes notes.  Subject to that audit, the next construction is the
off-shell de Donder ghost quartet on the V73 parent: introduce antighost and
Nakanishi--Lautrup fields, derive their side, cap, and joint variations, prove
that the induced boundary charge restricts to
\(\Omega_{74}^{\rm res}\), and populate a Sobolev trace domain on which
\(\mathsf T_J\) and the second-order map \(\mathsf R_J\) are closed.  Only
after that off-shell gate should the zero-level quotient be compared globally
with the V59 radiative phase.

\begin{assessment}[V74 continuum status]
The full Standard--Model--Einstein continuum remains unproved.  V74 closes
the flat orthogonal side--cap joint coefficient and the unrestricted residual
corner cocycle at the linear on-shell level.  The off-shell BV--BFV ghost
action, closed trace domains, global reduced metric--continuum Lagrangian
relation, nonlinear curved joints, renormalized common Standard Model stress
domain, and cofinal interacting continuum remain open.
\end{assessment}

\section*{V74 exact checks and append-only record}
\addcontentsline{toc}{section}{V74 exact checks and append-only record}

The joint coefficient was checked independently by forming
\(g_\varepsilon^{-1}\), normalizing the fixed conormals numerically, and
fitting the odd and even parts of
\(\sqrt{\det\sigma_\varepsilon}\operatorname{arsinh}z_\varepsilon\).
At step sizes \(10^{-3},5\cdot10^{-4},2.5\cdot10^{-4}\), the errors in the
predicted quadratic coefficient were respectively
\begin{equation}
 8.05\cdot10^{-6},\qquad
 2.01\cdot10^{-6},\qquad
 5.03\cdot10^{-7},
 \label{eq:V74-joint-numerical-check}
\end{equation}
showing the expected factor-four second-order convergence.  A seeded
three-component test of \eqref{eq:V74-opposite-corner-cocycle} returned
\begin{equation}
 \mathcal K_{72}=-1.8337528179634308,\qquad
 \mathcal K_{\rm cor}=1.8337528179634308,\qquad
 \mathcal K_{72}+\mathcal K_{\rm cor}=0.
 \label{eq:V74-cocycle-numerical-check}
\end{equation}
These checks audit the expansions and signs; the analytic proofs do not
depend on floating-point evidence.

The corrected authoritative V73 prefix used for this append had SHA--256
\begin{center}
\texttt{e3dc8eff45318ae338b3a05bf54bccbc39e009eb0f9ace51091e130025503d11}.
\end{center}
Before installation, the candidate was required to preserve that prefix
byte-for-byte, contain one active document terminator, balance all
environments and braces, have no duplicate label or unresolved internal
reference, and have equal counts of valid \verb|\(| and \verb|\)| inline
math delimiters.  The exact V73 predecessor is removed only after the
installed V74 hash and those structural conditions are verified.

% =============================================================================
% END APPENDED VERSION V74 MATERIAL
% =============================================================================

% =============================================================================
% BEGIN APPENDED VERSION V75 MATERIAL
% =============================================================================
\clearpage
\section*{V75 research note: scheduled NS/YM audit, the jointly assembled
de Donder quartet, and an off-shell residual boundary master equation}
\addcontentsline{toc}{section}{V75 companion audit and off-shell de Donder quartet}
\label{sec:V75-off-shell-quartet}

\subsection{Mandatory five-version companion audit}
\label{subsec:V75-companion-audit}

The V75 direction was selected only after reading the complete newest
appendices in the two companion programmes.  The snapshots used were
\begin{center}
\small
\begin{tabular}{@{}p{0.48\textwidth}rp{0.34\textwidth}@{}}
\toprule
\textbf{source}&\textbf{lines}&\textbf{SHA--256}\\
\midrule
\texttt{Renewal\_NS\_Stability\_Research\_Notes\_V31.tex}
&23052&\texttt{f1121b62238ad5491bb7cf03635ea9934942edf573ed8faaa9ee79e1d38a6696}\\
\texttt{Renewal\_Yang\_Mills\_Mass\_Gap\_Research\_Notes\_V63.tex}
&46597&\texttt{8a41e5bf4c8b19933dc2fb2ab8bea2f9ae29f0903c9b86e309c5b58a491840d5}\\
\bottomrule
\end{tabular}
\end{center}

The YM V63 filename is byte-identical to the audited V62 content and contains
no V63 appendix marker; it is recorded as the current shared-folder snapshot,
not as an additional mathematical result.

NS V31 proves that flat dyadic marks reconstruct critical coarea only after
the first radial moment is retained.  Separate norm estimates lose the
correlation among recent-entry times, signed transfer, and the active
high-parent tail, and return the continuation-class stock.  It supplies no
gravity or ghost estimate.

YM V60--V62 prove a parallel assembly restriction in a different problem.
V60 requires every resolution sharing an input row to enter one joint raw
functional before trace norm.  V61 replaces separate taxes by a partial trace
of the aggregate projector, and V62 proves a rectangular layer-cake identity
while showing that independently optimized local orientations can
underestimate the coherent joint coefficient by an unbounded factor.  Their
JREC capacity remains open and supplies no Einstein, BRST, or Sobolev theorem.

\begin{assessment}[V75 audit transfer]
\label{ass:V75-audit-transfer}
No analytic inequality is imported from either companion programme.  Their
common proved lesson changes the construction order: the \(B\)-field,
metric gauge functional, ghost Green term, every face transgression, and the
V74 corner phase are retained in one quadratic complex before the algebraic
\(B\)-equation or any boundary reduction is applied.  Eliminating \(B\) in
the bulk and reconstructing its face term afterward would repeat the
forbidden loss of shared data.
\end{assessment}

\begin{firewall}[Companion separation]
\label{fire:V75-companion-separation}
The NS restricted formation occupation and the YM JREC functional are not
field-theoretic gauge-fixing estimates.  V75 uses only their independently
proved order-of-assembly warning.  No NS regularity or Yang--Mills mass-gap
claim is transferred to the Einstein programme.
\end{firewall}

\subsection{The off-shell de Donder quartet as one BRST-exact functional}
\label{subsec:V75-quartet-action}

Retain the V36 linear harmonic covector
\begin{equation}
 C_\nu(H)
 :=D^\mu H_{\mu\nu}-\frac12D_\nu H.
 \label{eq:V75-harmonic-functional}
\end{equation}
Introduce an odd ghost \(\gamma_\mu\), an odd antighost \(\bar\gamma_\mu\), and an
even Nakanishi--Lautrup field \(B_\mu\).  Define
\begin{equation}
 sH_{\mu\nu}=D_\mu \gamma_\nu+D_\nu \gamma_\mu,\qquad
 s\gamma_\mu=0,\qquad
 s\bar\gamma_\mu=B_\mu,\qquad
 sB_\mu=0.
 \label{eq:V75-BRST-differential}
\end{equation}
Then \(s^2=0\) off shell and
\begin{equation}
 sC_\nu(H)=\Box \gamma_\nu.
 \label{eq:V75-sC-box}
\end{equation}

Choose the gauge fermion
\begin{equation}
 \Psi_{75}
 :=\frac1{2\kappa}\int_\Omega
 \bar\gamma^\nu\left(C_\nu(H)+\frac14B_\nu\right).
 \label{eq:V75-gauge-fermion}
\end{equation}

\begin{theorem}[Exact off-shell quartet action]
\label{thm:V75-quartet-action}
With left BRST differentiation and the odd Leibniz rule,
\begin{equation}
 \boxed{
 S_{\rm qt}^{(2)}
 :=s\Psi_{75}
 =\frac1{2\kappa}\int_\Omega
 \left[
  B^\nu C_\nu(H)+\frac14B_\nu B^\nu
  -\bar\gamma^\nu\Box \gamma_\nu
 \right].}
 \label{eq:V75-quartet-action}
\end{equation}
It is BRST invariant without using a field equation:
\begin{equation}
 sS_{\rm qt}^{(2)}=0.
 \label{eq:V75-off-shell-invariance}
\end{equation}
The algebraic \(B\)-equation and ghost equations are
\begin{equation}
 B_\nu=-2C_\nu(H),\qquad
 \Box \gamma_\nu=0,\qquad
 \Box\bar\gamma_\nu=0.
 \label{eq:V75-quartet-equations}
\end{equation}
Eliminating \(B\) gives exactly the V36 normalization
\begin{equation}
 S_{\rm qt}^{(2)}\big|_{B=-2C}
 =\frac1{2\kappa}\int_\Omega
 \left[-C_\nu C^\nu-\bar\gamma^\nu\Box \gamma_\nu\right].
 \label{eq:V75-B-elimination}
\end{equation}
\end{theorem}

\begin{proof}
Equation~\eqref{eq:V75-sC-box} follows by commuting flat derivatives:
\begin{align*}
 C_\nu(\mathcal G\gamma)
 &=D^\mu(D_\mu \gamma_\nu+D_\nu \gamma_\mu)
   -D_\nu D^\mu \gamma_\mu
 =\Box \gamma_\nu.
\end{align*}
Applying \(s\) to \eqref{eq:V75-gauge-fermion} gives
\eqref{eq:V75-quartet-action}.  The two nonzero variations of its integrand
are
\begin{equation}
 s(B^\nu C_\nu)=B^\nu\Box \gamma_\nu,\qquad
 s(-\bar\gamma^\nu\Box \gamma_\nu)=-B^\nu\Box \gamma_\nu,
\end{equation}
so they cancel pointwise, proving
\eqref{eq:V75-off-shell-invariance}.  Varying \(B\) gives
\(C+\tfrac12B=0\).  The antighost and ghost variations give the two wave
equations.  Finally
\begin{equation}
 (-2C)\mathbin{\cdot}C+\frac14(-2C)^2=-C^2,
\end{equation}
which proves \eqref{eq:V75-B-elimination}.
\end{proof}

\begin{corollary}[The quartet does not change the physical bulk equation]
\label{cor:V75-quartet-contractible}
The pair \((\bar\gamma,B)\) is a BRST doublet.  Its local BRST cohomology is
contractible, and eliminating \(B\) adds the de Donder gauge-fixing and
Faddeev--Popov wave operator without adding a metric solution class to the
V59 radiative quotient.
\end{corollary}

\begin{proof}
The contracting homotopy on polynomials of positive
\((\bar\gamma,B)\)-degree sends \(B\) to \(\bar\gamma\) and \(\bar\gamma\) to zero.
Its graded anticommutator with \(s\) is the doublet number operator.  Hence
positive doublet degree has zero cohomology.  The remaining equations in
\eqref{eq:V75-quartet-equations} are gauge and ghost equations, not an
additional metric polarization.
\end{proof}

\subsection{All-face Stokes form and the unreduced conormal}
\label{subsec:V75-all-face-conormal}

The peer-audit lesson is implemented by one exact Stokes identity, not by
separate estimates.  On the complete oriented slab,
\begin{align}
 S_{\rm qt}^{(2)}
 &=\frac1{2\kappa}
 \left\{
 \int_\Omega
 \left[
  B^\nu C_\nu(H)+\frac14B^2
  +D_\mu\bar\gamma^\nu D^\mu \gamma_\nu
 \right]
 -\sum_F\int_F\bar\gamma^\nu D_{n_F}\gamma_\nu
 \right\}.
 \label{eq:V75-joint-Stokes-representative}
\end{align}
The face sum includes every timelike side and both spacelike caps with their
actual outward conormals.

For a face \(F\), retain the V36 tensor
\begin{equation}
 \mathcal Q_F^{\mu\nu}(B)
 :=n_F^{(\mu}B^{\nu)}
   -\frac12\eta^{\mu\nu}(B\mathbin{\cdot}n_F).
 \label{eq:V75-QB}
\end{equation}

\begin{theorem}[Complete quartet face one-form]
\label{thm:V75-quartet-face-form}
The boundary one-form of \eqref{eq:V75-quartet-action}, with the bulk ghost
equation written on \(\gamma\), is
\begin{equation}
 \boxed{
 \Theta_{{\rm qt},F}
 =\frac1{2\kappa}\int_F
 \left[
  \mathcal Q_F^{\mu\nu}(B)\,\delta H_{\mu\nu}
  +(D_{n_F}\bar\gamma^\nu)\delta \gamma_\nu
  -\bar\gamma^\nu D_{n_F}\delta \gamma_\nu
 \right].}
 \label{eq:V75-quartet-face-one-form}
\end{equation}
No tangential integration by parts is used, so the quartet creates no
additional codimension-two scalar in this polarization.  On the algebraic
shell \(B=-2C(H)\), its metric coefficient is
\begin{equation}
 \mathcal Q_F(B)=-2\mathcal Q_F(C(H)),
 \label{eq:V75-shell-Q}
\end{equation}
the gauge-fixing contribution in the V36 conormal.
\end{theorem}

\begin{proof}
Variation of \(B^\nu C_\nu(H)\), followed by one bulk integration by parts
in the \(H\)-variation, gives the face coefficient
\(\mathcal Q_F^{\mu\nu}(B)\delta H_{\mu\nu}\).  Green's identity gives
\begin{align}
 -\int_\Omega\bar\gamma^\nu\Box\delta \gamma_\nu
 &=-\int_\Omega(\Box\bar\gamma^\nu)\delta \gamma_\nu
 \nonumber\\
 &\quad+\sum_F\int_F
 \left[
  (D_{n_F}\bar\gamma^\nu)\delta \gamma_\nu
  -\bar\gamma^\nu D_{n_F}\delta \gamma_\nu
 \right].
 \label{eq:V75-ghost-Green-identity}
\end{align}
The \(\delta\bar\gamma\) variation contains no differentiated variation in the
unintegrated representative.  Adding the terms proves
\eqref{eq:V75-quartet-face-one-form}.  Equation
\eqref{eq:V75-shell-Q} follows immediately from the \(B\)-equation.
\end{proof}

\begin{corollary}[Off-shell Einstein--quartet face conormal]
\label{cor:V75-off-shell-conormal}
On the V73 compatible-connection graph, the ungauge-fixed Einstein--GHY
coefficient and the unreduced quartet coefficient give
\begin{equation}
 \boxed{
 \mathfrak p_{75}^{\mu\nu}(H,B)
 =e_a^\mu e_b^\nu\Pi_{(1)}^{ab}(H)
  +\mathcal Q_F^{\mu\nu}(B).}
 \label{eq:V75-off-shell-conormal}
\end{equation}
Algebraic elimination of \(B\) gives
\begin{equation}
 \mathfrak p_{75}(H,-2C(H))=\mathfrak p_{36}(H).
 \label{eq:V75-conormal-reduction}
\end{equation}
\end{corollary}

\begin{proof}
The V73 gravitational face term is
\(\Pi_{(1)}^{ab}\delta H_{ab}\).  Add the metric part of
\eqref{eq:V75-quartet-face-one-form}.  Substitution of
\eqref{eq:V75-shell-Q} gives exactly
\eqref{eq:V36-action-conormal}.
\end{proof}

\subsection{The off-shell conormal graph is BRST invariant}
\label{subsec:V75-off-shell-graph}

\begin{lemma}[Pure gauge Brown--York block without a wave equation]
\label{lem:V75-off-shell-pure-gauge-BY}
For every smooth one-form \(\gamma\), without assuming \(\Box \gamma=0\),
\begin{equation}
 \Pi_{(1)}^{ab}(\mathcal G\gamma)
 =\mathscr P^{ab}(\gamma_n)
 =-D^aD^b\gamma_n+\eta_\partial^{ab}\Box_\partial \gamma_n.
 \label{eq:V75-off-shell-pure-gauge-BY}
\end{equation}
\end{lemma}

\begin{proof}
The cancellation in the V72 calculation of
\(K_{ab}^{(1)}(\mathcal G\gamma)\) uses only commutation of flat derivatives:
\begin{equation}
 K_{ab}^{(1)}(\mathcal G\gamma)=-D_aD_b\gamma_n.
\end{equation}
Taking its tangential trace gives
\eqref{eq:V75-off-shell-pure-gauge-BY}.  No four-dimensional wave equation
enters.
\end{proof}

\begin{theorem}[Invariant enlarged conormal graph]
\label{thm:V75-invariant-conormal-graph}
Extend \(s\) to the metric boundary momentum by
\begin{equation}
 sp^{\mu\nu}
 :=e_a^\mu e_b^\nu\mathscr P^{ab}(\gamma_n).
 \label{eq:V75-off-shell-p-shift}
\end{equation}
Then
\begin{equation}
 s\left[p-\mathfrak p_{75}(H,B)\right]=0
 \label{eq:V75-conormal-graph-invariance}
\end{equation}
for arbitrary \(\gamma\).  After imposing \(B=-2C(H)\), this shell is preserved
precisely when \(\Box \gamma=0\), and its restriction is the V72 lifted residual
vector.
\end{theorem}

\begin{proof}
Lemma~\ref{lem:V75-off-shell-pure-gauge-BY} gives
\(s\Pi_{(1)}=\mathscr P(\gamma_n)\), while \(sB=0\) gives
\(s\mathcal Q_F(B)=0\).  Thus the BRST variation of
\eqref{eq:V75-off-shell-conormal} equals
\eqref{eq:V75-off-shell-p-shift}, proving
\eqref{eq:V75-conormal-graph-invariance}.

The algebraic shell is \(B+2C(H)=0\).  Its variation is
\begin{equation}
 s(B+2C(H))=2\Box \gamma.
 \label{eq:V75-shell-preservation}
\end{equation}
It is therefore invariant exactly on the residual wave-ghost shell.  There
\(\mathcal Q_F(\Box \gamma)=0\), so
\eqref{eq:V75-off-shell-p-shift} agrees with the V72 pure-gauge V36 conormal.
\end{proof}

\begin{firewall}[Why \(B\) must not be eliminated before the face graph]
\label{fire:V75-no-early-B-elimination}
If \(B=-2C(H)\) is inserted before the boundary graph is formed, its BRST
variation becomes \(-2\Box \gamma\), and the momentum shift appears to require the
residual wave equation.  The independent \(B\)-field makes
\eqref{eq:V75-conormal-graph-invariance} an off-shell identity.  The wave
equation is needed only when the algebraically reduced V36 graph is selected.
This is the exact shared-data loss warned against by the V75 companion audit.
\end{firewall}

\subsection{A closed Sobolev trace domain for the V74 corner maps}
\label{subsec:V75-trace-domain}

Let \(\Omega\) be a compact flat slab with planar faces and joints.  Use
standard local Sobolev spaces on the slab; the following is a sufficient
domain, not a claim of energy-minimal regularity.

\begin{theorem}[Closed high-regularity ghost and corner trace domain]
\label{thm:V75-closed-trace-domain}
Fix \(s>3\) and set
\begin{equation}
 \mathscr D_{75}^s:=H^s(\Omega;T^*\Omega).
 \label{eq:V75-ghost-domain}
\end{equation}
For every codimension-two joint \(J\subset\partial\Gamma_T\), the maps
\begin{align}
 \mathsf T_J:\gamma&\longmapsto \gamma_b|_J,
 &
 \mathsf T_J&:\mathscr D_{75}^s\longrightarrow
 H^{s-1}(J)\hookrightarrow L^2(J),
 \label{eq:V75-T-trace-bound}\\
 \mathsf R_J:\gamma&\longmapsto
 \nu_a\mathscr P^{ab}(\gamma_n)|_J,
 &
 \mathsf R_J&:\mathscr D_{75}^s\longrightarrow
 H^{s-3}(J)\hookrightarrow L^2(J)
 \label{eq:V75-R-trace-bound}
\end{align}
are bounded and hence closed.  The residual shell
\begin{equation}
 \mathscr Z_{75}^s
 :=\ker\!\left(
 \Box:H^s(\Omega)\longrightarrow H^{s-2}(\Omega)
 \right)
 \label{eq:V75-closed-wave-shell}
\end{equation}
is closed.  With \(X,\Pi\in L^2(J)\), every pairing in the V74 corner
symplectic form and moment map is continuous on
\(\mathscr D_{75}^s\).
\end{theorem}

\begin{proof}
The trace from a four-dimensional slab to a codimension-two joint loses one
Sobolev derivative.  Thus \(\gamma|_J\in H^{s-1}(J)\).  The operator
\(\mathscr P\) has tangential order two, so its joint trace lies in
\(H^{s-3}(J)\).  Since \(s>3\), both spaces embed continuously in \(L^2(J)\).
Bounded linear maps between Hilbert spaces are closed.  The wave operator is
bounded from \(H^s\) to \(H^{s-2}\), so its kernel is closed.  Cauchy--Schwarz
then proves continuity of the \(L^2\) corner pairings.
\end{proof}

\begin{assessment}[Trace-domain scope]
\label{ass:V75-trace-scope}
Theorem~\ref{thm:V75-closed-trace-domain} closes the existence of one common
analytic domain for the local corner algebra.  It does not prove that
\(s>3\) is sharp, nor that the same traces extend to the eventual finite-energy
Lorentzian domain.  Lowering regularity requires a hyperbolic graph-trace
theorem which retains the second tangential normal-ghost Hessian.
\end{assessment}

\subsection{Off-shell boundary BFV charge}
\label{subsec:V75-off-shell-BFV}

On the unreduced metric phase define, for an arbitrary ghost \(\gamma\),
\begin{equation}
 \mu_{75}^{\rm met}(\gamma;H,p)
 :=\frac1{2\kappa}\int_{\Gamma_T}
 \left[
  \mathscr P^{ab}(\gamma_n)H_{ab}
  -(\mathcal G\gamma)_{\mu\nu}p^{\mu\nu}
 \right].
 \label{eq:V75-off-shell-metric-moment}
\end{equation}
This is the moment map of
\((\mathcal G\gamma,e_ae_b\mathscr P^{ab}(\gamma_n))\) on
\((\mathcal P_{36},\omega_{36})\), now interpreted on the enlarged
\((H,B)\) conormal graph.

Retain the V74 corner phase and moment map
\(\mu_\gamma^{\rm cor}\).  Complete the nonminimal boundary fields to a canonical
graded cotangent phase with momenta
\(\pi_{\bar\gamma},\pi_B,\pi_\gamma\) and form
\begin{equation}
 \omega_{\rm nm}
 :=\frac1{2\kappa}\int_{\Gamma_T}
 \left[
  \delta\pi_{\bar\gamma}^\nu\wedge\delta\bar\gamma_\nu
  +\delta\pi_B^\nu\wedge\delta B_\nu
  +\delta\pi_\gamma^\nu\wedge\delta \gamma_\nu
 \right]_{\rm graded}.
 \label{eq:V75-nonminimal-phase}
\end{equation}
The nonminimal cotangent-lift charge is
\begin{equation}
 \Omega_{75}^{\rm nm}
 :=\frac1{2\kappa}\int_{\Gamma_T}\pi_{\bar\gamma}^\nu B_\nu.
 \label{eq:V75-nonminimal-charge}
\end{equation}
It generates \(s\bar\gamma=B\), \(sB=s\gamma=0\) and the corresponding cotangent
transformations.

\begin{theorem}[Off-shell equivariant boundary charge]
\label{thm:V75-off-shell-boundary-master}
Define
\begin{equation}
 \boxed{
 \Omega_{75}^{\rm off}
 :=\mu_{75}^{\rm met}(\gamma)
   +\mu_\gamma^{\rm cor}
   +\Omega_{75}^{\rm nm}.}
 \label{eq:V75-off-shell-boundary-charge}
\end{equation}
On the domain \(\mathscr D_{75}^s\), \(s>3\), this charge is well defined for
arbitrary ghosts and satisfies
\begin{equation}
 \boxed{
 \{\Omega_{75}^{\rm off},\Omega_{75}^{\rm off}\}=0.}
 \label{eq:V75-off-shell-boundary-master}
\end{equation}
It preserves the enlarged conormal graph
\eqref{eq:V75-conormal-graph-invariance}.
\end{theorem}

\begin{proof}
The calculation of Theorem~\ref{thm:V72-corner-cocycle} used the wave
equation only to remove the gauge-fixing part of the already reduced V36
conormal.  For the enlarged conormal, Lemma
\ref{lem:V75-off-shell-pure-gauge-BY} gives directly
\begin{equation}
 \{\mu_{75}^{\rm met}(\gamma),\mu_{75}^{\rm met}(\zeta)\}
 =\mathcal K_{72}(\gamma,\zeta)
 \label{eq:V75-off-shell-metric-cocycle}
\end{equation}
for arbitrary ghost test jets \(\gamma,\zeta\).  The V74 corner phase contributes
\(-\mathcal K_{72}(\gamma,\zeta)\).  Thus their sum is equivariant without
\(\Box \gamma=0\).

The cotangent lift of the doublet differential has zero self-bracket because
\(s^2\bar\gamma=sB=0\).  It has zero cross bracket with the first two terms:
those terms depend on \(\gamma\) but not on \(\pi_\gamma\), and are independent of
\(\bar\gamma,B,\pi_{\bar\gamma},\pi_B\).  Polarization proves
\eqref{eq:V75-off-shell-boundary-master}.  Theorem
\ref{thm:V75-invariant-conormal-graph} proves graph preservation, and
Theorem~\ref{thm:V75-closed-trace-domain} proves continuity of all corner
terms.
\end{proof}

\begin{theorem}[Exact reduction to the V74 residual charge]
\label{thm:V75-reduction-to-V74}
Restrict \(\Omega_{75}^{\rm off}\) to
\begin{equation}
 B=-2C(H),\qquad \Box \gamma=0,\qquad \pi_{\bar\gamma}=0,
 \label{eq:V75-reduction-shell}
\end{equation}
and to the V73 compatible connection and jet graph.  Then
\begin{equation}
 \boxed{
 \Omega_{75}^{\rm off}\big|_{\eqref{eq:V75-reduction-shell}}
 =\Omega_{74}^{\rm res}.}
 \label{eq:V75-charge-reduction}
\end{equation}
Taking the product with the V70 holonomy phase gives
\begin{equation}
 \Omega_{75}^{\rm tot}
 :=\Omega_{70}^{\rm hol}+\Omega_{75}^{\rm off},
 \qquad
 \{\Omega_{75}^{\rm tot},\Omega_{75}^{\rm tot}\}=0
 \label{eq:V75-total-master}
\end{equation}
off shell on the enlarged product phase.
\end{theorem}

\begin{proof}
Equation~\eqref{eq:V75-conormal-reduction} identifies the metric phase with
the V36 graph.  On \(\Box \gamma=0\),
\eqref{eq:V75-off-shell-pure-gauge-BY} is exactly the V72 conormal shift, so
\(\mu_{75}^{\rm met}(\gamma)=\mu_\gamma\).  The corner term is already
\(\mu_\gamma^{\rm cor}\), and the nonminimal term vanishes on
\(\pi_{\bar\gamma}=0\).  This proves \eqref{eq:V75-charge-reduction}.  The V70
charge has zero self-bracket and disjoint canonical support, so its cross
bracket with \(\Omega_{75}^{\rm off}\) vanishes.  Apply
Theorem~\ref{thm:V75-off-shell-boundary-master}.
\end{proof}

\begin{firewall}[Algebraic BFV closure is not the full BV--BFV gluing theorem]
\label{fire:V75-BVBFV-scope}
Equation~\eqref{eq:V75-off-shell-boundary-master} closes the local linear
boundary charge on an explicit high-regularity domain.  V75 has not built the
bulk antifield symplectic form, proved the BV--BFV pullback identity on a
curved piecewise boundary, lowered the domain to finite energy, or established
global gluing across caps.  Those are separate requirements for a full
off-shell BV--BFV theorem.
\end{firewall}

\subsection{V75 result ledger and next acquisition}
\label{subsec:V75-ledger}

\paragraph{New exact conclusions.}
\begin{enumerate}[label=\textup{(\arabic*)},leftmargin=2.8em]
\item The mandatory NS V31/YM V63-snapshot audit was performed on digest-recorded
      snapshots; it transferred only a joint-assembly rule.
\item The gauge fermion \eqref{eq:V75-gauge-fermion} produces an off-shell
      nilpotent de Donder quartet with the exact V36 normalization after
      \(B\)-elimination.
\item The complete oriented face potential is
      \eqref{eq:V75-quartet-face-one-form}.
\item Keeping \(B\) independent gives the enlarged conormal
      \eqref{eq:V75-off-shell-conormal}, whose graph is BRST invariant for an
      arbitrary ghost.
\item The algebraically reduced V36 graph is invariant precisely on the
      residual wave-ghost shell, explaining the V72 condition without making
      it an off-shell assumption.
\item The sufficient Sobolev domain \(H^s\), \(s>3\), makes both V74 corner
      trace maps bounded and the wave-ghost shell closed.
\item The off-shell boundary charge
      \eqref{eq:V75-off-shell-boundary-charge} has zero self-bracket and
      restricts exactly to the unrestricted V74 residual charge.
\item The V70 holonomy factor remains a commuting factor of the off-shell
      enlarged charge.
\end{enumerate}

\paragraph{Next forced step.}
The local algebra now supports the reduction which V74 ordered.  V76 should
construct the half-space harmonic metric quotient on the zero level of
\(\widetilde\mu\), prove that the quotient Green form is symplectically
isomorphic to the two radiative V53/V59 phase, and identify the image of the
V59 positive continuum relation.  If a global local differential section is
needed, V71's no-splitting theorem requires replacing it by an operator
quotient or a controlled pseudodifferential chart.  In parallel, the
finite-energy graph-trace problem and the full curved BV--BFV gluing identity
remain open and must not be inferred from the \(s>3\) model.

\begin{assessment}[V75 continuum status]
The full Standard--Model--Einstein continuum remains unproved.  V75 closes
the flat linear off-shell de Donder quartet, its complete face conormal, a
high-regularity corner trace domain, and the enlarged residual boundary
master equation.  The global radiative quotient identification, pullback of
the V59 positive continuum relation, finite-energy trace extension, curved
BV--BFV gluing, nonlinear Brown--York/Bianchi stress identity, renormalized
common Standard Model stress domain, and cofinal interacting continuum remain
open.
\end{assessment}

\section*{V75 exact checks and append-only record}
\addcontentsline{toc}{section}{V75 exact checks and append-only record}

The identities used above were checked in a component Fourier-symbol script
with exact integer arithmetic.  For seeded integer covectors and ghosts it
verified
\begin{equation}
 C(\mathcal G_k\gamma)=k^2\gamma,\qquad
 K_{ab}^{(1)}(\mathcal G_k\gamma)=-k_ak_b\gamma_n,\qquad
 s\mathfrak p_{75}=e_ae_b\mathscr P^{ab}(\gamma_n)
 \label{eq:V75-component-cards}
\end{equation}
component by component.  It also checked the polynomial identity
\begin{equation}
 B\mathbin{\cdot}C+\frac14B^2=-C^2
 \quad\text{at }B=-2C
 \label{eq:V75-B-card}
\end{equation}
and the cancellation of the metric and corner skew forms for arbitrary
integer trace vectors.  These calculations audit signs and factors; the
proofs above are exact.

The V74 predecessor used for this append had SHA--256
\begin{center}
\texttt{5c5745933c3853e111170b2e5396280410a9ba07b4ab92f479df18faa2f03a6f}.
\end{center}
The V75 candidate retains its bytes up to the sole live final document
terminator, appends this material, and restores exactly one active
\verb|\end{document}|.  Validation includes balanced environments and
braces, unique labels, resolved internal references, valid inline and display
math delimiters, exact predecessor-prefix equality, and a final inventory
which removes only the superseded SM V74 file.

% =============================================================================
% END APPENDED VERSION V75 MATERIAL
% =============================================================================
% =============================================================================
% BEGIN APPENDED VERSION V76 MATERIAL
% =============================================================================
\clearpage
\section*{V76 research note: zero-level residual reduction, minimal corner
dressing, and the exact V59 pullback on the joint-free energy domain}
\addcontentsline{toc}{section}{V76 zero-level quotient and V59 pullback}
\label{sec:V76-zero-level-radiative-reduction}

\subsection*{Purpose and nonduplication boundary}

V75 ended with a precise instruction: form the zero-level quotient of the
harmonic metric boundary phase by the residual action and compare the result
with the two-polarization V53/V59 phase.  Several ingredients of that
comparison already exist and will not be rebuilt here.  V53 proved the
nonzero-null-frequency exact sequence and its norm-one TT splitting.  V59
constructed the positive radiative--continuum action and its skew-adjoint
finite-energy generator.  V63 proved the exact action normalization of the
radiative trace and conormal, V68 gave a local potential-compatible phase
which is unitary to V59, and V71 proved that no bounded finite-order local
metric splitting can realize the same quotient at the same energy order.
V72--V75 supplied the missing boundary moment map, corner cocycle, minimal
opposite-cocycle carrier, and off-shell quartet.

The issue left between those results is not another choice of polarization
frame.  It is the following exact reduction question.  Given a radiative
class, when can one choose its metric representative and its V74 corner data
so that the total residual moment map is zero?  If such a choice exists, is it
unique modulo the residual orbit, and what symplectic form does it carry?
This note answers those questions by a general linear reduction lemma and
then applies it to the renewal Einstein phase.

There are two geometries which must be kept distinct.  The V59 time generator
lives on
\(
 \mathbb R_t\times\mathbb R_y^2\times\mathbb R_x^+
\)
and has no finite temporal joints.  On its smooth joint-free core the corner
cocycle vanishes and the reduction is exactly the V53/V59 radiative phase.
A compact slab has caps and side--cap joints.  There the V74 carrier is
essential.  The reduction remains exact, but its image is characterized by
an explicit proper-charge obstruction.  Proving that this obstruction
vanishes on the eventual curved finite-energy interface is a later gluing
problem and is not silently assumed below.

\subsection{A linear corner-dressing theorem}
\label{subsec:V76-abstract-reduction}

We first isolate the algebra which was distributed through V72--V75.  Let
\((M,\omega_M)\) and \((C,\omega_C)\) be finite-dimensional symplectic
vector spaces, let \(\mathfrak g\) be a vector space of residual parameters,
and let
\begin{equation}
 A:\mathfrak g\longrightarrow M,
 \qquad
 B:\mathfrak g\longrightarrow C
 \label{eq:V76-AB-maps}
\end{equation}
be linear maps.  Put
\begin{equation}
 K(\epsilon,\zeta):=\omega_M(A\epsilon,A\zeta)
 \label{eq:V76-metric-cocycle}
\end{equation}
and assume
\begin{equation}
 \omega_C(B\epsilon,B\zeta)=-K(\epsilon,\zeta).
 \label{eq:V76-opposite-cocycle-assumption}
\end{equation}
The combined fundamental vector and its span are
\begin{equation}
 V\epsilon:=(A\epsilon,B\epsilon),
 \qquad
 W:=\operatorname{ran}V\subset M\oplus C.
 \label{eq:V76-combined-orbit}
\end{equation}
The product form is denoted by
\(\omega:=\omega_M\oplus\omega_C\).

Define the radical of the metric cocycle and the corner transpose by
\begin{align}
 \operatorname{Rad}K
 &:=%
 \{\epsilon\in\mathfrak g:
       K(\epsilon,\zeta)=0\ \hbox{for every }\zeta\in\mathfrak g\},
 \label{eq:V76-radK}\\
 B^\flat:C&\longrightarrow\mathfrak g^*,
 &
 (B^\flat c)(\epsilon)&:=\omega_C(B\epsilon,c).
 \label{eq:V76-B-flat}
\end{align}
We say that \(C\) is the \emph{minimal corner carrier} when it is the
symplectic quotient of \(\operatorname{span}B(\mathfrak g)\) by the radical
of the restricted form.  Thus, after this quotient, \(B\) is onto and
\begin{equation}
 \ker B=\operatorname{Rad}K.
 \label{eq:V76-kernel-B-radK}
\end{equation}

\begin{lemma}[Isotropic orbit and zero moment level]
\label{lem:V76-isotropic-reduction}
Under \eqref{eq:V76-opposite-cocycle-assumption}, \(W\) is isotropic.  The
linear moment map
\begin{equation}
 \widetilde\mu_\epsilon(z):=\omega(V\epsilon,z)
 \label{eq:V76-abstract-moment}
\end{equation}
has simultaneous zero level
\begin{equation}
 \widetilde\mu^{-1}(0)=W^{\perp_\omega}.
 \label{eq:V76-zero-level-orthogonal}
\end{equation}
Consequently
\begin{equation}
 P_{\rm red}:=W^{\perp_\omega}/W
 \label{eq:V76-reduced-phase}
\end{equation}
carries the nondegenerate reduced symplectic form.
\end{lemma}

\begin{proof}
For two residual parameters,
\begin{align*}
 \omega(V\epsilon,V\zeta)
 &=\omega_M(A\epsilon,A\zeta)
   +\omega_C(B\epsilon,B\zeta)\\
 &=K(\epsilon,\zeta)-K(\epsilon,\zeta)=0.
\end{align*}
Thus \(W\subset W^{\perp_\omega}\).  Equation
\eqref{eq:V76-zero-level-orthogonal} is the definition of simultaneous
vanishing of \eqref{eq:V76-abstract-moment}.  The form descends because
changing either representative by an element of \(W\) does not change its
pairing with a vector in \(W^{\perp_\omega}\).  If a class pairs to zero
with every reduced class, its representative lies in
\((W^{\perp_\omega})^{\perp_\omega}=W\), so the class is zero.
\end{proof}

Now let \((R,\omega_R)\) be a candidate radiative phase and suppose that
\begin{equation}
 q:M\longrightarrow R,
 \qquad
 qA=0,
 \qquad
 \ker q=\operatorname{ran}A,
 \qquad
 q\ \hbox{is onto}.
 \label{eq:V76-radiative-quotient-hypotheses}
\end{equation}
For \(r\in R\), choose any \(m\in M\) with \(qm=r\) and define
\begin{equation}
 \ell_m(\epsilon):=\omega_M(A\epsilon,m).
 \label{eq:V76-metric-charge-functional}
\end{equation}

\begin{lemma}[Representative-independent obstruction]
\label{lem:V76-obstruction-class}
The class
\begin{equation}
 \chi(r):=[\ell_m]
 \in\mathfrak g^*/\operatorname{ran}B^\flat
 \label{eq:V76-obstruction-map}
\end{equation}
does not depend on the choice of \(m\).  For a minimal finite-dimensional
corner carrier,
\begin{equation}
 \operatorname{ran}B^\flat
 =\operatorname{Ann}(\operatorname{Rad}K),
 \label{eq:V76-Bflat-range}
\end{equation}
and hence
\begin{equation}
 \chi(r)=0
 \quad\Longleftrightarrow\quad
 \ell_m(\epsilon)=0
 \quad\hbox{for every }\epsilon\in\operatorname{Rad}K.
 \label{eq:V76-proper-charge-criterion}
\end{equation}
\end{lemma}

\begin{proof}
If \(m'=m+A\alpha\), then
\begin{equation}
 \ell_{m'}(\epsilon)-\ell_m(\epsilon)
 =K(\epsilon,\alpha)
 =-(B^\flat B\alpha)(\epsilon).
 \label{eq:V76-obstruction-choice-change}
\end{equation}
Thus the quotient class is independent of the lift.

If \(\epsilon\in\ker B\), every element of
\(\operatorname{ran}B^\flat\) vanishes on \(\epsilon\).  Therefore
\(\operatorname{ran}B^\flat\subset\operatorname{Ann}(\ker B)\).  Minimality
makes \(B\) onto and \(\omega_C\) nondegenerate, so \(B^\flat\) is
injective.  Hence
\begin{align*}
 \dim\operatorname{ran}B^\flat
 &=\dim C
  =\dim\mathfrak g-\dim\ker B\\
 &=\dim\operatorname{Ann}(\ker B).
\end{align*}
The inclusion is equality.  Use
\eqref{eq:V76-kernel-B-radK} to obtain
\eqref{eq:V76-Bflat-range} and
\eqref{eq:V76-proper-charge-criterion}.
\end{proof}

Put
\begin{equation}
 R_{\rm adm}:=\ker\chi.
 \label{eq:V76-admissible-radiative-phase}
\end{equation}

\begin{theorem}[Minimal corner dressing and exact reduced image]
\label{thm:V76-corner-dressing}
Assume \eqref{eq:V76-opposite-cocycle-assumption}, minimality of \(C\),
and \eqref{eq:V76-radiative-quotient-hypotheses}.  Then the map
\begin{equation}
 \overline q:P_{\rm red}\longrightarrow R,
 \qquad
 \overline q[(m,c)]:=qm
 \label{eq:V76-descended-radiative-map}
\end{equation}
is well defined and is a linear bijection
\begin{equation}
 \boxed{\quad P_{\rm red}\cong R_{\rm adm}.\quad}
 \label{eq:V76-reduction-bijection}
\end{equation}
More explicitly, a lift \(m\) of \(r\) lies over the zero moment level if
and only if there is a corner variable \(c_m\) satisfying
\begin{equation}
 B^\flat c_m=-\ell_m.
 \label{eq:V76-dressing-equation}
\end{equation}
For \(r\in R_{\rm adm}\), this \(c_m\) exists and is unique.  Replacing
\(m\) by \(m+A\alpha\) replaces it by
\begin{equation}
 c_{m+A\alpha}=c_m+B\alpha,
 \label{eq:V76-dressing-gauge-change}
\end{equation}
so the dressed pair changes by exactly one combined residual orbit vector.
The reduced form transported to \(R_{\rm adm}\) is
\begin{equation}
 \boxed{
 \omega_{\rm adm}(r_1,r_2)
 =\omega_M(m_1,m_2)+\omega_C(c_{m_1},c_{m_2}).}
 \label{eq:V76-transported-form}
\end{equation}
\end{theorem}

\begin{proof}
The map is constant on combined orbits because \(qA=0\).  A pair
\((m,c)\) belongs to the zero level exactly when, for every \(\epsilon\),
\begin{equation}
 0=\omega(V\epsilon,(m,c))
   =\ell_m(\epsilon)+(B^\flat c)(\epsilon),
\end{equation}
which is \eqref{eq:V76-dressing-equation}.  Lemma
\ref{lem:V76-obstruction-class} proves existence precisely for
\(r\in R_{\rm adm}\).  If \(B^\flat c=0\), then \(c\) is symplectically
orthogonal to \(B(\mathfrak g)=C\), so \(c=0\).  Thus the dressing is
unique.

Equation \eqref{eq:V76-obstruction-choice-change} and
\(B^\flat B\alpha=-K(\,cdot\,,\alpha)\) give
\eqref{eq:V76-dressing-gauge-change}.  This proves surjectivity of
\(\overline q\) onto \(R_{\rm adm}\) and independence of the dressed
class from the metric lift.

For injectivity, suppose \(qm=0\).  Then \(m=A\alpha\).  The zero-level
condition gives
\begin{align*}
 0
 &=K(\epsilon,\alpha)+\omega_C(B\epsilon,c)\\
 &=\omega_C(B\epsilon,c-B\alpha)
 \qquad\hbox{for every }\epsilon.
\end{align*}
Surjectivity of \(B\) and nondegeneracy of \(\omega_C\) imply
\(c=B\alpha\).  Hence \((m,c)=V\alpha\) and its reduced class is zero.
Finally, \eqref{eq:V76-transported-form} is the product form evaluated on
the unique dressed representatives.  Equation
\eqref{eq:V76-dressing-gauge-change} and the zero-level condition show
directly that it is independent of both lifts.
\end{proof}

\begin{corollary}[Closed Hilbert version]
\label{cor:V76-Hilbert-reduction}
The conclusions of Lemma~\ref{lem:V76-isotropic-reduction} remain valid for
strong symplectic Hilbert spaces when \(A,B\) are bounded and \(W\) is
replaced by its closed range.  The conclusions of
Theorem~\ref{thm:V76-corner-dressing} remain valid if, in addition,
\(B\) has closed range and the finite-dimensional annihilator identity
\eqref{eq:V76-Bflat-range} is replaced by the corresponding closed-range
identity.  On a dense smooth core the weak condition
\begin{equation}
 \omega(V\epsilon,z)=0
 \quad\hbox{for every residual }\epsilon
 \label{eq:V76-weak-moment-zero}
\end{equation}
is exactly the integral moment-map equation; its Hilbert closure defines the
zero level without requiring pointwise joint traces.
\end{corollary}

\begin{proof}
Continuity preserves isotropy under closure.  For a closed subspace of a
strong symplectic Hilbert space,
\((W^{\perp_\omega})^{\perp_\omega}=W\), so the reduced form is
nondegenerate.  The remaining proof is unchanged after using the closed
range theorem in place of the dimension count.  On the smooth core,
\eqref{eq:V76-weak-moment-zero} is
\eqref{eq:V76-abstract-moment}; continuity then supplies its unique closed
extension.
\end{proof}

\begin{firewall}[What the abstract theorem does and does not assert]
\label{fire:V76-abstract-scope}
Theorem~\ref{thm:V76-corner-dressing} proves an exact quotient theorem; it
does not assert that \(R_{\rm adm}=R\).  That equality is a separate
proper-charge statement.  Nor does
\eqref{eq:V76-transported-form} automatically equal a previously chosen
radiative form: the possible difference is an explicit corner term and must
be checked in the geometry under study.
\end{firewall}

\subsection{Application to the V75 Einstein and V74 corner phases}
\label{subsec:V76-Einstein-application}

Return to a compact flat slab and first work on a fixed finite
tangential jet fibre or tangential Fourier fibre at each joint.  The traces
are taken from the smooth \(H^s\)-domain of V75, \(s>3\).  On the
algebraically reduced harmonic shell put
\begin{align}
 A\epsilon
 &:=%
 \left(
  \mathcal G\epsilon,
  e_a^{\mu}e_b^{\nu}\mathscr P^{ab}(\epsilon_n)
 \right),
 \label{eq:V76-Einstein-A}\\
 B\epsilon
 &:=%
 \bigoplus_J
 \left(
  \mathsf T_J(\epsilon),-\mathsf R_J(\epsilon)
 \right).
 \label{eq:V76-Einstein-B}
\end{align}
Here the first map is the V75 unrestricted conormal lift and the second is
the V74 corner translation.  Let \(C_{74}^{\min}\) be the direct sum over
joints of the carriers
\(\mathscr W_J^{\rm red}\) from
Theorem~\ref{thm:V74-corner-minimality}.  By construction,
\begin{equation}
 B:\mathfrak g_{76}^s\longrightarrow C_{74}^{\min}
 \quad\hbox{is onto},
 \label{eq:V76-B-onto}
\end{equation}
after quotienting residual parameters which induce the same translation.

\begin{proposition}[V74--V75 are exactly the hypotheses of the dressing
theorem]
\label{prop:V76-hypotheses-realized}
With
\begin{equation}
 M=\mathcal P_{36},
 \qquad C=C_{74}^{\min},
 \qquad
 \omega_M=\omega_{36},
 \qquad \omega_C=\omega_{\rm cor},
 \label{eq:V76-phase-identification}
\end{equation}
one has
\begin{align}
 \omega_{36}(A\epsilon,A\zeta)
 &=\mathcal K_{72}(\epsilon,\zeta),
 \label{eq:V76-metric-K}\\
 \omega_{\rm cor}(B\epsilon,B\zeta)
 &=-\mathcal K_{72}(\epsilon,\zeta).
 \label{eq:V76-corner-minus-K}
\end{align}
The simultaneous zero level of the resulting linear moment map is exactly
\begin{equation}
 \widetilde\mu_\epsilon
 =\mu_\epsilon+\mu_\epsilon^{\rm cor}=0
 \quad\hbox{for every }\epsilon,
 \label{eq:V76-V74-zero-level}
\end{equation}
and its quotient is a symplectic reduction.  This statement remains true
off shell on the V75 enlarged \((H,B)\) conormal graph after replacing the
residual parameter by an arbitrary ghost jet; the physical radiative
identification below is taken only after
\(B=-2C(H)\), \(C(H)=0\), and \(\Box\epsilon=0\).
\end{proposition}

\begin{proof}
Equation \eqref{eq:V76-metric-K} is
Theorem~\ref{thm:V72-corner-cocycle}, and
\eqref{eq:V76-corner-minus-K} is
Theorem~\ref{thm:V74-opposite-cocycle}.  Their normalizations already use
the same \(\kappa\).  The V74 total moment map is the sum in
\eqref{eq:V76-V74-zero-level}, so Lemma
\ref{lem:V76-isotropic-reduction} applies.  V75 proves the same two cocycle
identities for arbitrary ghost jets on the enlarged conormal graph.  Its
algebraic and harmonic reduction is precisely the last displayed shell.
\end{proof}

The radical which appears in the abstract theorem is now concrete:
\begin{equation}
 \mathfrak g_{76}^{\rm prop}
 :=\operatorname{Rad}\mathcal K_{72}
 =\{\epsilon:B\epsilon=0\ \hbox{in }C_{74}^{\min}\}.
 \label{eq:V76-proper-residual-algebra}
\end{equation}
It includes every residual parameter whose pair of V74 joint traces is null
in the reduced corner carrier.  For a harmonic metric trace--conormal datum
\(m=(H,p)\), its charge functional is the exact V72 expression
\begin{equation}
 \ell_m(\epsilon)
 =\frac1{2\kappa}\int_{\Gamma_T}
 \left[
  \mathscr P^{ab}(\epsilon_n)H_{ab}
  -(\mathcal G\epsilon)_{\mu\nu}p^{\mu\nu}
 \right].
 \label{eq:V76-explicit-charge-functional}
\end{equation}

\begin{theorem}[Exact finite-slab fibre lifting criterion]
\label{thm:V76-finite-slab-criterion}
Let \(q\) be any radiative quotient of the corresponding finite harmonic
metric conormal fibre for which \(\ker q=\operatorname{ran}A\).  A radiative
class \(r\)
has a representative on the V74 zero level if and only if
\begin{equation}
 \boxed{
 \ell_m(\epsilon)=0
 \quad\hbox{for every }
 \epsilon\in\mathfrak g_{76}^{\rm prop}.}
 \label{eq:V76-finite-slab-proper-condition}
\end{equation}
When this holds, the reduced corner variable is uniquely determined by
\begin{equation}
 \frac1\kappa\sum_J\int_J
 \left[
  -\mathsf R_J(\epsilon)^bX_b
  -\mathsf T_J(\epsilon)_b\Pi^b
 \right]
 =-\ell_m(\epsilon)
 \label{eq:V76-explicit-corner-dressing}
\end{equation}
for every residual \(\epsilon\), modulo the combined residual orbit.
Thus the minimal corner carrier cancels every charged component detected by
\(\mathcal K_{72}\); it cannot cancel a nonzero charge on
\(\operatorname{Rad}\mathcal K_{72}\), because all of its Hamiltonians
vanish there.
\end{theorem}

\begin{proof}
Equation \eqref{eq:V76-finite-slab-proper-condition} is
\eqref{eq:V76-proper-charge-criterion} with
\(K=\mathcal K_{72}\).  The left side of
\eqref{eq:V76-explicit-corner-dressing} is
\((B^\flat(X,\Pi))(\epsilon)\) in the signs and normalization of
\eqref{eq:V74-corner-moment-map}.  Hence the displayed equation is exactly
\eqref{eq:V76-dressing-equation}.  Existence, uniqueness, and the orbit
statement follow from Theorem~\ref{thm:V76-corner-dressing}.  If
\(\epsilon\) lies in the radical, then \(B\epsilon=0\) in the minimal
carrier, so the left side vanishes for every corner datum.  This proves the
last assertion.
\end{proof}

\begin{firewall}[The minimal cocycle carrier is not a universal charge
reservoir]
\label{fire:V76-no-universal-reservoir}
V74 proved fibrewise minimality for realizing the opposite two-cocycle.  It did not
prove that the resulting corner phase can absorb an arbitrary element of
\(\mathfrak g^*\).  Theorem~\ref{thm:V76-finite-slab-criterion} gives the
exact distinction: its image is
\(\operatorname{Ann}(\operatorname{Rad}\mathcal K_{72})\).  Any finite-slab
claim that the whole V59 phase survives reduction must therefore prove
\eqref{eq:V76-finite-slab-proper-condition}; cocycle cancellation alone is
not that proof.
\end{firewall}

\subsection{The nonzero-frequency harmonic conormal quotient}
\label{subsec:V76-frequency-quotient}

Fix a nonzero real null covector \(k=(-|\xi|,\xi)\), or a nonzero point of
the V63 stable boundary cone after analytic continuation.  Let
\(\mathscr M_k^{\rm har}\) be the harmonic metric solution--conormal fibre:
its configuration amplitude belongs to \(Z_k\), and its momentum is the
V36 conormal of that solution.  Define
\begin{equation}
 q_k:\mathscr M_k^{\rm har}\longrightarrow\mathscr R_k,
 \qquad
 q_k(H,p):=(a,p_{\rm rad}),
 \label{eq:V76-fibre-radiative-map}
\end{equation}
where
\begin{equation}
 a=T_k\overline H,
 \qquad
 p_{\rm rad}=2\varepsilon_D r_\theta(H)
 \label{eq:V76-fibre-trace-conormal}
\end{equation}
uses the V53 quotient and the V63 action conormal.  At real frequency,
Theorem~\ref{thm:V63-exact-real-cone-lift} gives
\begin{equation}
 p_{\rm rad}=D_na.
 \label{eq:V76-exact-radiative-conormal}
\end{equation}

\begin{theorem}[Exact nonzero-frequency metric/conormal sequence]
\label{thm:V76-frequency-exact-sequence}
For every such nonzero fibre,
\begin{equation}
 0\longrightarrow\mathfrak g_k
 \mathop{\longrightarrow}^{A_k}
 \mathscr M_k^{\rm har}
 \mathop{\longrightarrow}^{q_k}
 \mathscr R_k
 \longrightarrow0
 \label{eq:V76-conormal-exact-sequence}
\end{equation}
is exact.  In particular, the conormal does not create an additional
kernel beyond the four residual wave-gauge amplitudes, and the quotient has
exactly the two V53 polarizations with the V59 conormal normalization.
\end{theorem}

\begin{proof}
On configurations, exactness is
Theorem~\ref{thm:V53-radiative-exact-sequence}.  Thus \(q_k(H,p)=0\)
implies
\(\overline H=\Gamma_k\epsilon\), equivalently
\(H=\mathcal G_k\epsilon\).  Because the datum lies on the linear conormal
graph, its momentum is then the conormal of this pure gauge field.  The V75
identity
\begin{equation*}
 \mathfrak p_{36}(\mathcal G\epsilon)
 =e_ae_b\mathscr P^{ab}(\epsilon_n)
\end{equation*}
on the residual shell makes the full pair exactly \(A_k\epsilon\).  This
proves \(\ker q_k=\operatorname{ran}A_k\).

For surjectivity, use the V53 norm-one TT splitting for the configuration.
V63 computes the restriction of the actual GHY one-form to that graph and
proves \eqref{eq:V76-exact-radiative-conormal}, including normal incidence
and the continuous glancing endpoint.  Hence every radiative trace--conormal
pair has a metric conormal lift.  No independent momentum kernel remains
because momentum is fixed by the conormal graph.
\end{proof}

\begin{corollary}[The quotient uses no local metric section]
\label{cor:V76-no-local-section-used}
The map \(q_k\) and the exact sequence
\eqref{eq:V76-conormal-exact-sequence} are intrinsic quotient statements.
They may be assembled by measurable order-zero polarization charts, but no
finite-order local right inverse is required.  Any claim of such a right
inverse at V59 energy would contradict
Theorem~\ref{thm:V71-no-local-radiative-splitting}.
\end{corollary}

\begin{proof}
The quotient map \(T_k\) is defined without selecting a representative.
V63 supplies one chart for calculating the induced conormal, while changes
of polarization frame act unitarily on both \(a\) and \(p_{\rm rad}\).  The
no-splitting theorem then gives the last statement.
\end{proof}

\subsection{Joint-free Hilbert quotient and the V68 phase}
\label{subsec:V76-joint-free-Hilbert}

Consider now the geometry used by the V59 generator:
\begin{equation}
 \Omega_x=\mathbb R_y^2\times\mathbb R_x^+,
 \qquad
 \Gamma_T=\mathbb R_t\times\mathbb R_y^2.
 \label{eq:V76-V59-geometry}
\end{equation}
Start with smooth solutions and residual waves whose tangential support is
compact, and then take the homogeneous energy closure.  There is no
boundary of \(\Gamma_T\), so
\begin{equation}
 \partial\Gamma_T=\varnothing,
 \qquad
 \mathcal K_{72}=0,
 \qquad
 C_{74}^{\min}=\{0\}.
 \label{eq:V76-no-joint-cocycle}
\end{equation}

Let \(\mathscr Z_{76}^{\rm en}\) be the V71 closed harmonic reducing space
on this geometry and let \(\mathscr G_{76}^{\rm en}\) be the closure of the
residual wave-gauge image.  Define the weak zero level
\begin{equation}
 \mathscr N_{76}^{\rm en}
 :=\left\{z\in\mathscr Z_{76}^{\rm en}:
   \omega_{36}(A\epsilon,z)=0
   \text{ for every }A\epsilon\in\mathscr G_{76}^{\rm en}
 \right\}.
 \label{eq:V76-energy-zero-level}
\end{equation}
On the smooth core this is the V72 moment equation.  The definition by
symplectic orthogonality gives its closed finite-energy extension without
using the second-order joint trace \(\mathsf R_J\), because there is no
joint.

\begin{lemma}[The on-shell gauge Green form is a joint term]
\label{lem:V76-gauge-Green-joint}
Let \(H\) solve the ungauge-fixed linearized Einstein equation and the
de Donder condition, and let \(\epsilon\) obey \(\Box\epsilon=0\).  On the
V36 solution--conormal graph there is a local antisymmetric tensor
\(\mathcal U_\epsilon^{\mu\nu}(H)=-\mathcal U_\epsilon^{\nu\mu}(H)\) such
that
\begin{equation}
 \omega_{36}(A\epsilon,z_H)
 =\frac1{2\kappa}\int_{\partial\Gamma_T}
   n_\mu\nu_\nu\mathcal U_\epsilon^{\mu\nu}(H).
 \label{eq:V76-gauge-Green-joint}
\end{equation}
Here \(n\) is the conormal to \(\Gamma_T\) and \(\nu\) is its outward
conormal at the joint.  In particular, this moment vanishes when
\(\partial\Gamma_T=\varnothing\) and the fields have sufficient decay.
\end{lemma}

\begin{proof}
Let \(\mathcal E^{\mu\nu}(H)\) be the formally self-adjoint linearized
Einstein operator before gauge fixing, and let
\(j_{\mathcal E}^{\mu}(H_1,H_2)\) be its Lagrange current, with the order of
its arguments chosen to agree with \(\omega_{36}\).  Formal
self-adjointness gives
\begin{equation}
 D_\mu j_{\mathcal E}^{\mu}(\mathcal G\epsilon,H)
 =(\mathcal G\epsilon)_{\mu\nu}\mathcal E^{\mu\nu}(H)
  -\mathcal E^{\mu\nu}(\mathcal G\epsilon)H_{\mu\nu}.
 \label{eq:V76-Einstein-Lagrange-identity}
\end{equation}
Linearized diffeomorphism invariance and the linearized Bianchi identity are
the off-shell identities
\begin{equation}
 \mathcal E(\mathcal G\epsilon)=0,
 \qquad
 D_\mu\mathcal E^{\mu\nu}(H)=0.
 \label{eq:V76-linear-gauge-Bianchi}
\end{equation}
Thus the right side of
\eqref{eq:V76-Einstein-Lagrange-identity} is
\begin{equation}
 2D_\mu\!\left(\epsilon_\nu\mathcal E^{\mu\nu}(H)\right).
 \label{eq:V76-Bianchi-integration}
\end{equation}
The difference
\(
 j_{\mathcal E}^{\mu}(\mathcal G\epsilon,H)
 -2\epsilon_\nu\mathcal E^{\mu\nu}(H)
\)
is therefore an identically conserved local linear differential current.
Applying the elementary flat algebraic Poincare homotopy to that current,
or equivalently carrying out the two integrations by parts in
\eqref{eq:V76-Einstein-Lagrange-identity}, gives a local antisymmetric
\(\mathcal U_\epsilon^{\mu\nu}(H)\) with
\begin{equation}
 j_{\mathcal E}^{\mu}(\mathcal G\epsilon,H)
 =2\epsilon_\nu\mathcal E^{\mu\nu}(H)
  +D_\nu\mathcal U_\epsilon^{\mu\nu}(H).
 \label{eq:V76-superpotential-identity}
\end{equation}
This is an identity of jets and does not use a boundary condition.

On the Einstein shell the first term on the right vanishes.  The de Donder
gauge-fixing Green current also vanishes because
\(
 C(H)=0
\)
and
\(
 C(\mathcal G\epsilon)=\Box\epsilon=0
\).
Adding GHY changes the face potential by a field-space exact variation and
therefore does not change its antisymmetrized Green form; changes between
face representatives are tangential divergences and are absorbed into the
same superpotential.  Hence the normal integral of
\eqref{eq:V76-superpotential-identity} is precisely the V36 pairing in
\eqref{eq:V76-gauge-Green-joint}.  Stokes' theorem on \(\Gamma_T\) proves
that equation and its joint-free consequence.
\end{proof}

\begin{lemma}[TT representatives lie on the joint-free zero level]
\label{lem:V76-TT-zero-moment}
Every nonzero-frequency V53 TT solution--conormal representative belongs to
\(\mathscr N_{76}^{\rm en}\).  Moreover, every element of
\(\mathscr Z_{76}^{\rm en}\) is the sum of such a TT representative and an
element of \(\mathscr G_{76}^{\rm en}\).  Consequently
\begin{equation}
 \mathscr N_{76}^{\rm en}=\mathscr Z_{76}^{\rm en}.
 \label{eq:V76-zero-level-whole-harmonic-space}
\end{equation}
\end{lemma}

\begin{proof}
Lemma~\ref{lem:V76-gauge-Green-joint} makes the residual moment of every
smooth harmonic solution vanish on the compact-support core of
\eqref{eq:V76-V59-geometry}; this includes every TT representative.  At one
nonzero null covector the same fact has the direct algebraic certificate
\begin{equation}
 \langle\Gamma_k\epsilon,E\rangle_D
 =2\epsilon_\nu k_\mu E^{\mu\nu}
  -(k\mathbin{\cdot}\epsilon)\operatorname{tr}E=0
 \label{eq:V76-gauge-TT-orthogonality}
\end{equation}
for every trace-free transverse polarization \(E\).  The first statement,
which follows from the full superpotential identity rather than a selected
normal Fourier branch, is the one used for the half-space closure.

The V53 exact sequence decomposes each nonzero null fibre into its TT class
and a residual amplitude.  Its norm-one TT multiplier is bounded in the
homogeneous energy norm.  Measurable direct integration, followed by
closure, gives the asserted decomposition of
\(\mathscr Z_{76}^{\rm en}\).  Gauge--gauge pairings vanish by
\eqref{eq:V76-no-joint-cocycle}, and gauge--TT pairings vanish by the first
part.  Therefore every harmonic energy datum satisfies
\eqref{eq:V76-energy-zero-level}.
\end{proof}

Let \(\mathscr E_{\rm rad}\) denote the metric part of the V59 homogeneous
energy phase.  Its elements are \((a,a_t)\), with norm squared
\begin{equation}
 \|(a,a_t)\|_{\mathscr E_{\rm rad}}^2
 =\int_{\Omega_x}
  \left(|a_t|^2+|\partial_xa|^2+|\nabla_ya|^2\right).
 \label{eq:V76-radiative-energy-norm}
\end{equation}

\begin{theorem}[Unitary zero-level quotient]
\label{thm:V76-unitary-radiative-quotient}
The fibre maps \(q_k\) assemble to a unitary symplectic map
\begin{equation}
 U_{76}:
 \mathscr N_{76}^{\rm en}/\mathscr G_{76}^{\rm en}
 \mathop{\longrightarrow}^{\cong}
 \mathscr E_{\rm rad}.
 \label{eq:V76-unitary-quotient-map}
\end{equation}
The reduced Green form is exactly the two-scalar V59 Green form with
\(p_{\rm rad}=D_na\), and the reduced norm is exactly
\eqref{eq:V76-radiative-energy-norm}.  The spatially homogeneous
zero-frequency sector is excluded from this assertion.
\end{theorem}

\begin{proof}
Lemma~\ref{lem:V76-TT-zero-moment} reduces the quotient to the V71 Hilbert
quotient.  Fibrewise injectivity and surjectivity are
Theorem~\ref{thm:V76-frequency-exact-sequence}.  The V53 TT projector has
operator norm one and identifies the quotient energy with the TT energy, so
the direct-integral map is isometric and onto; it is therefore unitary.

On the TT representative, Proposition
\ref{prop:V63-conormal-from-action} gives the restricted metric Green form,
and Theorem~\ref{thm:V63-exact-real-cone-lift} converts its conormal to
\(p_{\rm rad}=D_na\).  This is exactly the V59 scalar-wave boundary form.
Gauge directions are the zero-level quotient directions, so the equality is
independent of the TT representative.  V53 proved that no continuous
rank-two extension exists at \(\xi=0\); assigning that measure-zero fibre an
arbitrary value does not classify homogeneous moduli.  Hence the final
restriction is necessary.
\end{proof}

Let
\begin{equation}
 \mathcal J_{68}(a,D_0a)
 =\left(
  \frac{D_0a+D_na}{2},
  \frac{D_1a}{\sqrt2},
  \frac{D_2a}{\sqrt2},
  \frac{D_0a-D_na}{2}
 \right)
 =:(B,Y_1,Y_2,Z)
 \label{eq:V76-V68-map}
\end{equation}
be the V68 local phase map, componentwise on the rank-two radiative bundle.

\begin{corollary}[Local phase after, and only after, quotienting]
\label{cor:V76-local-phase-after-quotient}
The composite
\begin{equation}
 \mathcal J_{68}\circ U_{76}:
 \mathscr N_{76}^{\rm en}/\mathscr G_{76}^{\rm en}
 \longrightarrow\mathscr R_\Phi
 \label{eq:V76-quotient-to-local-phase}
\end{equation}
is unitary onto the closed V68 potential-compatible range
\(\mathscr R_\Phi\).  On that range the quotient dynamics is the local
symmetric-hyperbolic V68 system.  This does not provide a local map back to
the ten-component metric; that inverse remains the nonlocal or quotient
operation required by V71.
\end{corollary}

\begin{proof}
V68 proves that \(\mathcal J_{68}\) is a unitary map from the V59
radiative energy phase onto \(\mathscr R_\Phi\), including its propagated
curl and mixed constraints.  Compose it with
Theorem~\ref{thm:V76-unitary-radiative-quotient}.  A local inverse into the
full metric would contradict
Theorem~\ref{thm:V71-no-local-radiative-splitting}.
\end{proof}

The result can be summarized by the commuting array
\begin{equation}
 \begin{array}{ccccc}
 \mathscr N_{76}^{\rm en}
 &\mathop{\longrightarrow}^{q}&
 \mathscr E_{\rm rad}
 &\mathop{\longrightarrow}^{\mathcal J_{68}}&
 \mathscr R_\Phi\\
 \big\downarrow\rlap{\scriptstyle\pi}
 &&\big\Vert&&\big\Vert\\
 \mathscr N_{76}^{\rm en}/\mathscr G_{76}^{\rm en}
 &\mathop{\longrightarrow}^{U_{76}}&
 \mathscr E_{\rm rad}
 &\mathop{\longrightarrow}^{\mathcal J_{68}}&
 \mathscr R_\Phi.
 \end{array}
 \label{eq:V76-commuting-quotient-diagram}
\end{equation}
Here the upper \(q\) is constant on residual orbits.  Equivalently,
\begin{equation}
 q=U_{76}\pi,
 \qquad
 \mathcal J_{68}q=\mathcal J_{68}U_{76}\pi.
 \label{eq:V76-commuting-equalities}
\end{equation}

\subsection{Exact pullback of the positive continuum relation}
\label{subsec:V76-continuum-pullback}

Let \(P_b\) be the two-polarization continuum boundary phase with canonical
pair \((b,t_b)\).  On
\(\mathscr E_{\rm rad}\times P_b\), let \(L_{59}\) be the V59 relation
\begin{equation}
 \operatorname{Tr}_\rho b=\operatorname{Tr}_x a,
 \qquad
 p_{\rm rad}+t_b=0.
 \label{eq:V76-V59-relation}
\end{equation}
The V59 action and V63 Green calculation prove that this relation is
Lagrangian in the radiative--continuum phase.

Define its reduced metric pullback by
\begin{equation}
 L_{76}^{\rm red}
 :=\left\{
  ([z],b,t_b):
  (U_{76}[z],b,t_b)\in L_{59}
 \right\}
 \label{eq:V76-reduced-pullback}
\end{equation}
and its zero-level inverse image by
\begin{equation}
 \widetilde L_{76}
 :=\left\{
  (z,b,t_b):
  z\in\mathscr N_{76}^{\rm en},
  (U_{76}\pi z,b,t_b)\in L_{59}
 \right\}.
 \label{eq:V76-unreduced-inverse-image}
\end{equation}

\begin{theorem}[The pulled-back interface is Lagrangian]
\label{thm:V76-pulled-interface-Lagrangian}
On the joint-free V59 energy geometry:
\begin{enumerate}[label=\textup{(\alph*)},leftmargin=2.8em]
\item \(L_{76}^{\rm red}\) is Lagrangian in the reduced
      metric--continuum phase;
\item \(\widetilde L_{76}\) is Lagrangian in the unreduced product phase and
      contains the residual orbit \(\mathscr G_{76}^{\rm en}\) as a null
      subspace of its zero-level presentation;
\item its quotient by that orbit is exactly \(L_{76}^{\rm red}\), and its
      physical rows are precisely
\begin{equation}
 \operatorname{Tr}_\rho b
 =\operatorname{Tr}_x U_{76}[z],
 \qquad
 D_nU_{76}[z]+t_b=0.
 \label{eq:V76-physical-interface-rows}
\end{equation}
\end{enumerate}
No local metric representative is used in any of these statements.
\end{theorem}

\begin{proof}
The map \(U_{76}\) is symplectic by
Theorem~\ref{thm:V76-unitary-radiative-quotient}.  Therefore the inverse
image of the V59 Lagrangian relation is Lagrangian, proving (a).

For (b), use the general inverse-image argument for symplectic reduction.
Let \(W=\mathscr G_{76}^{\rm en}\) and let \(\pi:W^{\perp}/W\) be the
quotient map.  Two vectors in \(\widetilde L_{76}\) have zero pairing
because their reduced classes lie in the Lagrangian
\(L_{76}^{\rm red}\); hence the inverse image is isotropic.  Conversely, if
a vector is symplectically orthogonal to \(\widetilde L_{76}\), its pairing
with \(W\subset\widetilde L_{76}\) vanishes, so it belongs to
\(W^\perp\).  Its reduced class is orthogonal to
\(L_{76}^{\rm red}\) and therefore belongs to that Lagrangian.  The original
vector lies in \(\widetilde L_{76}\).  Thus the inverse image equals its
symplectic orthogonal and is Lagrangian.  The quotient assertion in (c) is
the definition of the inverse image, and
\eqref{eq:V76-physical-interface-rows} is
\eqref{eq:V76-V59-relation} with the exact conormal
\eqref{eq:V76-exact-radiative-conormal}.
\end{proof}

\begin{corollary}[Skew-adjoint generator on the full harmonic metric
quotient]
\label{cor:V76-reduced-skew-generator}
Let \(\mathcal E_q\) and \(\mathcal G_q\) be the V59 total energy space and
skew-adjoint generator.  Extend \(U_{76}\) by the identity on the continuum
factor and write this unitary map as \(\mathbb U_{76}\).  Then
\begin{equation}
 \mathcal G_{76}^{\rm red}
 :=\mathbb U_{76}^{-1}\mathcal G_q\mathbb U_{76},
 \qquad
 D(\mathcal G_{76}^{\rm red})
 :=\mathbb U_{76}^{-1}D(\mathcal G_q)
 \label{eq:V76-conjugated-generator}
\end{equation}
is skew-adjoint and generates a strongly continuous unitary group on the
zero-level harmonic metric quotient times the continuum energy phase.  Its
positive norm is exactly twice the V59 energy
\(E_q\).
\end{corollary}

\begin{proof}
Unitary conjugation preserves adjoints, domains, and strong continuity of
the generated group.  The V59 theorem gives
\(\mathcal G_q^*=-\mathcal G_q\), so
\begin{align*}
 (\mathcal G_{76}^{\rm red})^*
 &=\mathbb U_{76}^{-1}\mathcal G_q^*\mathbb U_{76}
 =-\mathcal G_{76}^{\rm red}.
\end{align*}
The norm statement follows from unitarity and the V59 energy identity.
\end{proof}

\begin{assessment}[What is now meant by the full metric generator]
\label{ass:V76-full-metric-generator-scope}
Corollary~\ref{cor:V76-reduced-skew-generator} advances the V59 result from
an abstract two-polarization phase to the operator-level quotient of the
full local harmonic metric complex.  It is not a skew-adjoint generator on
all unreduced ten metric components, because residual gauge waves have been
divided out.  This quotient is the physical object selected by the V53 exact
sequence and by the V71 no-splitting theorem.
\end{assessment}

\subsection{The remaining finite-slab comparison is an explicit pair of
tests}
\label{subsec:V76-finite-slab-defect}

For a compact slab, let \(R_{76}^{\rm adm}\) be the admissible radiative
space defined by \eqref{eq:V76-finite-slab-proper-condition}.  For
\(r_i\in R_{76}^{\rm adm}\), choose metric lifts \(m_i\) and their unique
corner dressings \(c_i\).  Define the comparison form
\begin{equation}
 \Delta_{76}(r_1,r_2)
 :=\omega_{36}(m_1,m_2)
   +\omega_{\rm cor}(c_1,c_2)
   -\omega_{59}^{\rm rad}(r_1,r_2).
 \label{eq:V76-finite-slab-defect-form}
\end{equation}

\begin{proposition}[Well-defined finite-slab obstruction pair]
\label{prop:V76-finite-slab-obstruction-pair}
The obstruction map \(\chi\) of
\eqref{eq:V76-obstruction-map} and the skew form \(\Delta_{76}\) of
\eqref{eq:V76-finite-slab-defect-form} are independent of every metric lift
and corner representative.  The whole V59 relation pulls back
symplectically through the compact-slab V74 reduction if and only if
\begin{align}
 \chi(r)&=0
 &&\text{for every metric datum }r\text{ occurring in }L_{59},
 \label{eq:V76-image-test}\\
 \Delta_{76}(r_1,r_2)&=0
 &&\text{for every two such metric data}.
 \label{eq:V76-form-test}
\end{align}
On the joint-free domain both conditions hold by
Theorems~\ref{thm:V76-unitary-radiative-quotient} and
\ref{thm:V76-pulled-interface-Lagrangian}.
\end{proposition}

\begin{proof}
Lift independence of \(\chi\) is
Lemma~\ref{lem:V76-obstruction-class}.  The first two terms in
\eqref{eq:V76-finite-slab-defect-form} are the reduced form
\eqref{eq:V76-transported-form}, so their sum is also lift independent.
The first condition is exactly the assertion that every required radiative
datum lies in the image of the reduction.  Once that holds, the second is
exactly the assertion that the reduced form agrees with the V59 form.  A
symplectic identification pulls a Lagrangian relation back to a Lagrangian
relation, while failure of either image or form equality prevents that
identification.  The last statement follows from the already proved
joint-free equalities.
\end{proof}

\begin{firewall}[No finite-energy joint theorem has been smuggled in]
\label{fire:V76-no-joint-trace-claim}
The weak Hilbert zero level
\eqref{eq:V76-energy-zero-level} solves the V59 geometry because that
geometry has no codimension-two joint.  It does not extend the second-order
map \(\mathsf R_J\) to the minimal V59 energy space on a compact slab.
V75's \(H^s\), \(s>3\), trace theorem remains the available joint domain.
The finite-energy graph-trace extension and the evaluation of
\eqref{eq:V76-image-test}--\eqref{eq:V76-form-test} on curved glued slabs
remain open.
\end{firewall}

\subsection{V76 result ledger and next acquisition}
\label{subsec:V76-ledger}

\paragraph{New exact conclusions.}
\begin{enumerate}[label=\textup{(\arabic*)},leftmargin=2.8em]
\item The opposite V72/V74 cocycles make the combined residual orbit
      isotropic, and the simultaneous moment zero level is its symplectic
      orthogonal.
\item The obstruction to lifting a radiative class is the
      representative-independent charge class
      \(\chi\in\mathfrak g^*/\operatorname{ran}B^\flat\).
\item Fibrewise, for the V74 minimal carrier, lifting is equivalent to
      vanishing of the metric moment on
      \(\operatorname{Rad}\mathcal K_{72}\).  When it exists, the corner
      dressing is unique and changes equivariantly with the metric lift.
\item The nonzero-frequency harmonic metric--conormal sequence has kernel
      exactly the four residual wave gauges and quotient exactly the two V53
      polarizations with the V63/V59 conormal normalization.
\item On the joint-free V59 generator geometry, TT representatives have
      zero residual moment, and the zero-level Hilbert quotient is unitarily
      symplectically isomorphic to the V59 radiative energy phase.
\item Composing with V68 gives a local symmetric-hyperbolic realization only
      after quotienting; no prohibited local metric section is introduced.
\item The inverse image of the V59 trace/traction relation is Lagrangian on
      the zero level, and its quotient is exactly the positive V59 relation.
\item Unitary conjugation supplies a skew-adjoint positive
      metric-quotient--continuum generator with the complete V59 domain.
\item For compact slabs, the remaining global comparison is reduced
      fibrewise to the explicit proper-charge test \eqref{eq:V76-image-test} and form test
      \eqref{eq:V76-form-test}.
\end{enumerate}

\paragraph{Next forced step.}
V77 should evaluate the proper-charge functional
\eqref{eq:V76-explicit-charge-functional} on the complete V63/V65 physical
graph at a finite side--cap joint.  The calculation must retain the same
metric, residual, corner, and continuum traces until the joint integration
is complete.  If the radical charge vanishes, one should compute
\(\Delta_{76}\) and prove the compact-slab pullback.  If it does not vanish,
the correct upstream response is to identify the missing proper-edge
constraint or enlarge the edge cotangent phase; it is not to choose a local
metric splitting.  The V80 iteration remains the next mandatory NS/YM
companion audit.

\begin{assessment}[V76 continuum status]
The full Standard--Model--Einstein continuum remains unproved.  V76 closes
the zero-level radiative quotient and the positive continuum pullback on the
joint-free flat finite-energy generator geometry.  It also gives an exact
finite-slab fibrewise corner-dressing theorem and isolates the only two tests needed
for the compact-joint comparison.  The finite-energy joint trace, compact
slab proper-charge and form tests, curved BV--BFV gluing, nonlinear
Brown--York/Bianchi stress identity, common renormalized Standard Model
stress domain, and cofinal interacting continuum remain open.
\end{assessment}

\section*{V76 exact checks and append-only record}
\addcontentsline{toc}{section}{V76 exact checks and append-only record}

The signs in Theorem~\ref{thm:V76-corner-dressing} were audited by finite
rational symplectic matrices satisfying
\begin{equation}
 A^T\Omega_MA+B^T\Omega_CB=0.
 \label{eq:V76-matrix-isotropy-card}
\end{equation}
For every seeded admissible lift, the computed dressing obeyed
\begin{equation}
 B^T\Omega_Cc=-A^T\Omega_Mm,
 \label{eq:V76-matrix-dressing-card}
\end{equation}
and replacing \(m\) by \(m+A\alpha\) replaced \(c\) by
\(c+B\alpha\).  The calculation is an audit of the displayed signs; the
proofs above are exact.

The V75 predecessor used for this append had SHA--256
\begin{center}
\texttt{f2f98dd968b7d2e76967b93fc2a1758379d8617f3ab1e4c1e05dfe55295bc31e}.
\end{center}
The V76 candidate retains its bytes up to the sole live final document
terminator, appends this material, and restores exactly one active
\verb|\end{document}|.  Validation includes balanced environments and
braces, unique labels, resolved internal references, valid inline and display
math delimiters, exact predecessor-prefix equality, and a final inventory
which removes only the superseded SM V75 file.

% =============================================================================
% END APPENDED VERSION V76 MATERIAL
% =============================================================================





% =============================================================================
% BEGIN APPENDED VERSION V77 MATERIAL
% =============================================================================
\clearpage
\section*{V77 research note: the exact Einstein joint charge, a
charge-carrying radical missed by cocycle minimality, and why the full V74
corner cotangent phase must be retained}
\addcontentsline{toc}{section}{V77 joint charge and charge-complete V74 phase}
\label{sec:V77-joint-charge-completion}

\subsection*{Purpose and correction of the next compact-slab gate}

V76 proved that the joint-free V59 generator is the zero-level quotient of
the full harmonic metric complex, and it reduced the compact-slab comparison
to two tests.  The first was the image test: the V72 metric moment must vanish
on the radical of the V74 corner cocycle, because the minimal opposite-cocycle
carrier has no Hamiltonian in that radical.  V76 did not assume that this
test passed.

This note evaluates it.  The calculation must distinguish the rank of a
two-cocycle from the rank of the complete charge map.  At a nonzero joint
Fourier mode, the V72 trace cocycle has a transverse tangential radical.  The
Einstein joint charge is generally nonzero in precisely that direction.  An
explicit transverse-traceless \(3\)-\(4\)-\(5\) wave gives a nonzero witness.
Thus V74's carrier remains minimal for cancelling the cocycle, but it is too
small for quotienting every charged residual transformation.

The necessary canonical pair is already present in V74's unreduced
three-pair presentation.  It is the transverse projection of
\((X_b,\Pi^b)\).  Its residual translation is isotropic, so it does not
contribute to the cocycle rank, but its conjugate momentum absorbs the
transverse Brown--York charge.  The error would be to replace the full V74
phase by its cocycle-minimal quotient before imposing the moment equations.
The construction is proved on every nonzero flat joint Fourier fibre.  The
zero joint mode, finite-energy joint traces, and curved gluing remain
separate problems.

\subsection{An improved Brown--York joint charge}
\label{subsec:V77-improved-joint-charge}

Let \(\Gamma\) be a planar timelike face with spacelike conormal \(n\),
induced flat metric \(\eta_\partial\), and boundary
\(\partial\Gamma=\bigcup_JJ\).  Indices \(a,b\) are tangential to
\(\Gamma\), and \(\nu_a\) is the outward conormal to a joint within
\(\Gamma\).  Write
\begin{equation}
 q_{ab}:=H_{ab},
 \qquad
 \tau:=\eta_\partial^{ab}q_{ab},
 \qquad
 f:=\epsilon_n.
 \label{eq:V77-face-data}
\end{equation}
On the harmonic Einstein shell the V36 conormal is the Brown--York momentum
\(\Pi^{ab}_{(1)}\).  The V72 moment therefore reads
\begin{equation}
 2\kappa\,\mu_\epsilon(H)
 =\int_\Gamma
 \left[
  \mathscr P^{ab}(f)q_{ab}
  -2D_a\epsilon_b\,\Pi^{ab}_{(1)}
 \right],
 \label{eq:V77-face-moment-start}
\end{equation}
where
\begin{equation}
 \mathscr P^{ab}(f)
 =-D^aD^bf+\eta_\partial^{ab}\Box_\partial f.
 \label{eq:V77-P-operator}
\end{equation}

\begin{lemma}[Linear face constraints]
\label{lem:V77-face-constraints}
If \(H\) solves the vacuum linearized Einstein equation on the flat
background, then its induced face data obey
\begin{align}
 D_a\Pi^{ab}_{(1)}&=0,
 \label{eq:V77-momentum-constraint}\\
 D_aD_bq^{ab}-\Box_\partial\tau&=0.
 \label{eq:V77-scalar-constraint}
\end{align}
These identities do not require a choice of residual gauge parameter.
\end{lemma}

\begin{proof}
The linearized Gauss--Codazzi identities on a planar anchor give
\begin{align*}
 G^{(1)}_{nb}(H)&=D_aK_{(1)}{}^a{}_b-D_bK_{(1)}
                 =D_a\Pi_{(1)}{}^a{}_b,\\
 2G^{(1)}_{nn}(H)&=D_aD_bq^{ab}-\Box_\partial\tau.
\end{align*}
The anchor extrinsic curvature and intrinsic curvature vanish, so there are
no lower-order terms.  The vacuum equations set both projections of
\(G^{(1)}\) to zero and prove
\eqref{eq:V77-momentum-constraint}--%
\eqref{eq:V77-scalar-constraint}.
\end{proof}

\begin{theorem}[Exact improved joint charge]
\label{thm:V77-improved-joint-charge}
For every vacuum harmonic metric field \(H\) and residual wave parameter
\(\epsilon\), the V72 face moment is the sum of the joint charges
\begin{equation}
 \boxed{
 \mu_\epsilon(H)
 =\frac1{2\kappa}\sum_J\int_J\nu_a
 \left[
  f\left(D_bq^{ab}-D^a\tau\right)
  +(D^af)\tau
  -(D_bf)q^{ab}
  -2\epsilon_b\Pi^{ab}_{(1)}
 \right].}
 \label{eq:V77-improved-joint-charge}
\end{equation}
In particular, this gives an explicit representative of the
superpotential asserted abstractly in
Lemma~\ref{lem:V76-gauge-Green-joint}.
\end{theorem}

\begin{proof}
Integrate the first term of
\eqref{eq:V77-face-moment-start} twice.  Flat tangential derivatives
commute, and one obtains
\begin{align}
 \int_\Gamma\mathscr P^{ab}(f)q_{ab}
 &=
 \int_\Gamma f\left(
  -D_aD_bq^{ab}+\Box_\partial\tau
 \right)
 \nonumber\\
 &\quad
 +\sum_J\int_J\nu_a
 \left[
  fD_bq^{ab}
  -fD^a\tau
  +(D^af)\tau
  -(D_bf)q^{ab}
 \right].
 \label{eq:V77-P-integration}
\end{align}
The bulk term vanishes by
\eqref{eq:V77-scalar-constraint}.  Symmetry of
\(\Pi_{(1)}\) gives
\begin{align}
 -2\int_\Gamma D_a\epsilon_b\,\Pi^{ab}_{(1)}
 &=
 2\int_\Gamma\epsilon_bD_a\Pi^{ab}_{(1)}
 -2\sum_J\int_J\nu_a\epsilon_b\Pi^{ab}_{(1)}.
 \label{eq:V77-Pi-integration}
\end{align}
The remaining bulk term vanishes by
\eqref{eq:V77-momentum-constraint}.  Adding the two boundary expressions and
dividing by \(2\kappa\) proves
\eqref{eq:V77-improved-joint-charge}.
\end{proof}

\begin{corollary}[Exact relation to the V76 proper-charge functional]
\label{cor:V77-V76-charge-functional}
On the physical conormal graph, the functional
\(\ell_m(\epsilon)\) in
\eqref{eq:V76-explicit-charge-functional} equals the right side of
\eqref{eq:V77-improved-joint-charge}.  Thus the V76 image obstruction is a
joint functional of the residual zero and first jets and the physical
metric and Brown--York first jets.
\end{corollary}

\begin{proof}
Equation~\eqref{eq:V77-face-moment-start} is the V72 moment on the harmonic
V36 conormal graph.  Theorem~\ref{thm:V77-improved-joint-charge} is an exact
integration of that expression, so it is the same functional.
\end{proof}

\subsection{The nonzero joint-frequency cocycle radical}
\label{subsec:V77-cocycle-radical}

Take a temporal joint \(J=\{t=t_J\}\times\Sigma\), with local flat
coordinates \(y=(y^1,y^2)\) on \(\Sigma\).  Put
\begin{equation}
 \nu_a=\sigma_J\delta_a^0,
 \qquad \sigma_J\in\{-1,1\},
 \label{eq:V77-joint-orientation}
\end{equation}
and fix a nonzero joint Fourier covector
\(k=(k_1,k_2)\).  A dot denotes \(D_0\) at the joint.  For
\(f=\epsilon_n\), the V74 normal-Hessian trace has symbol
\begin{align}
 \mathsf R_J(\epsilon)^0
 &=\sigma_J|k|^2 f,
 \label{eq:V77-R0-symbol}\\
 \mathsf R_J(\epsilon)^A
 &=\sigma_J\,i k^A\dot f.
 \label{eq:V77-RA-symbol}
\end{align}
The irrelevant common Fourier density and the reality pairing are
suppressed.

Let
\begin{equation}
 e_\parallel^A:=\frac{k^A}{|k|},
 \qquad
 e_\perp^A:=\varepsilon^{A}{}_B e_\parallel^B,
 \label{eq:V77-parallel-perp-frame}
\end{equation}
and decompose the tangential trace as
\begin{equation}
 \mathsf T_J(\epsilon)
 =(T_0,T_\parallel,T_\perp).
 \label{eq:V77-T-decomposition}
\end{equation}

\begin{lemma}[Exact trace image and rank-four cocycle]
\label{lem:V77-trace-rank}
At a fixed nonzero joint Fourier covector, the residual wave-jet trace image
is
\begin{equation}
 \mathscr W_{J,k}
 \cong
 \mathbb C^3_T\oplus
 \operatorname{span}_{\mathbb C}\{R_0,R_\parallel\}.
 \label{eq:V77-trace-image}
\end{equation}
The values
\((T_0,T_\parallel,T_\perp,f,\dot f)\) are independent residual Cauchy
jets.  In the ordered coordinates
\begin{equation}
 (T_0,T_\parallel,T_\perp,R_0,R_\parallel),
 \label{eq:V77-trace-coordinate-order}
\end{equation}
the V72 cocycle, apart from its common \(1/\kappa\) factor and orientation,
has matrix
\begin{equation}
 K_{J,k}=
 \begin{pmatrix}
 0&0&0&-1&0\\
 0&0&0&0&-1\\
 0&0&0&0&0\\
 1&0&0&0&0\\
 0&1&0&0&0
 \end{pmatrix}.
 \label{eq:V77-K-matrix}
\end{equation}
Consequently
\begin{equation}
 \operatorname{rank}K_{J,k}=4,
 \qquad
 \operatorname{Rad}K_{J,k}
 =\operatorname{span}\{T_\perp\}.
 \label{eq:V77-rank-and-radical}
\end{equation}
\end{lemma}

\begin{proof}
Equations
\eqref{eq:V77-R0-symbol}--\eqref{eq:V77-RA-symbol} show that the range of
\(\mathsf R_J\) is the two-plane spanned by its time component and the
spatial component parallel to \(k\).  At a noncharacteristic time slice,
the wave equation constrains the second time jet, not the values of
\(\epsilon_\mu\) and \(D_0\epsilon_\mu\).  Hence the three tangential values,
\(f\), and \(\dot f\) can be prescribed independently in the local residual
wave-jet algebra.  Since \(k\ne0\), the map
\((f,\dot f)\mapsto(R_0,R_\parallel)\) is invertible.

In these coordinates the V74 trace form is
\begin{equation*}
 K_{J,k}(v,w)
 =R_0(v)T_0(w)+R_\parallel(v)T_\parallel(w)
  -R_0(w)T_0(v)-R_\parallel(w)T_\parallel(v).
\end{equation*}
This is exactly \eqref{eq:V77-K-matrix}.  Its two canonical \(2\)-by-\(2\)
blocks have rank four, and its only trace radical is the displayed
transverse coordinate.
\end{proof}

\begin{corollary}[What V74 minimality removes]
\label{cor:V77-V74-rank-four-carrier}
At \(k\ne0\), the V74 minimal opposite-cocycle carrier has dimension four.
It contains the two canonical pairs dual to
\((T_0,R_0)\) and \((T_\parallel,R_\parallel)\).  The transverse trace
\(T_\perp\) is zero in that quotient, even though it is a nonzero residual
boundary transformation.
\end{corollary}

\begin{proof}
Apply Theorem~\ref{thm:V74-corner-minimality} and
\eqref{eq:V77-rank-and-radical}.  Quotienting the five-dimensional trace
image by its one-dimensional radical leaves the rank-four symplectic
carrier.
\end{proof}

\subsection{The radical carries a physical Einstein charge}
\label{subsec:V77-radical-charge}

Restrict the improved charge to a residual trace in
\(\operatorname{Rad}K_{J,k}\).  Equations
\eqref{eq:V77-R0-symbol}--\eqref{eq:V77-RA-symbol} and \(k\ne0\) imply
\begin{equation}
 f=0,\qquad \dot f=0
 \label{eq:V77-radical-normal-zero}
\end{equation}
at the joint.  Every tangential derivative of \(f\) also vanishes there.
The only surviving trace is
\begin{equation}
 \epsilon_0=0,\qquad
 \epsilon_\parallel=0,\qquad
 \epsilon_\perp=T_\perp.
 \label{eq:V77-radical-trace}
\end{equation}
Theorem~\ref{thm:V77-improved-joint-charge} therefore gives
\begin{equation}
 \boxed{
 \mu_\epsilon(H)\big|_{\operatorname{Rad}K_{J,k}}
 =-\frac1\kappa\int_J
 T_\perp\,
 e_{\perp b}\nu_a\Pi^{ab}_{(1)}(H).}
 \label{eq:V77-radical-charge}
\end{equation}

\begin{theorem}[A nonzero \(3\)-\(4\)-\(5\) TT witness]
\label{thm:V77-TT-radical-witness}
The functional \eqref{eq:V77-radical-charge} is not identically zero on the
nonzero-frequency V59 radiative graph.

More precisely, use coordinates
\((t,y,z,x)\), take the side \(x=0\) with \(D_n=\partial_x\), let
\(y\) have period \(2\pi/3\), and put
\begin{equation}
 \phi(t,y,x):=\cos(5t)\cos(3y)\cos(4x).
 \label{eq:V77-345-scalar}
\end{equation}
Define a metric perturbation by
\begin{equation}
 H_{yz}=H_{zy}:=\partial_x\phi,
 \qquad
 H_{xz}=H_{zx}:=-\partial_y\phi,
 \qquad
 H_{\mu\nu}=0\ \hbox{otherwise}.
 \label{eq:V77-345-TT-field}
\end{equation}
Then \(H\) is trace free, divergence free, and solves the wave equation.
At the joint
\begin{equation}
 t_J=\frac{\pi}{10}
 \label{eq:V77-witness-joint-time}
\end{equation}
its Brown--York momentum obeys
\begin{equation}
 \Pi^{0z}_{(1)}\big|_{x=0,t=t_J}
 =-\frac{15}{2}\sin(3y).
 \label{eq:V77-witness-momentum}
\end{equation}
The residual wave
\begin{equation}
 \epsilon_z
 :=\cos\!\left(5(t-t_J)\right)\sin(3y)\cos(4x),
 \qquad
 \epsilon_\mu=0\quad(\mu\ne z)
 \label{eq:V77-witness-residual}
\end{equation}
has joint trace \(T_\perp=\sin(3y)\), \(R=0\), and hence lies in
\(\operatorname{Rad}K_{J,k}\) for \(k=(3,0)\).  If the \(z\)-period has
length \(L_z\) and \(\nu_0=\sigma_J\), then
\begin{equation}
 \boxed{
 \mu_\epsilon(H)
 =\sigma_J\frac{5\pi L_z}{2\kappa}\ne0.}
 \label{eq:V77-nonzero-radical-charge}
\end{equation}
\end{theorem}

\begin{proof}
The dispersion identity \(5^2=3^2+4^2\) proves
\(\Box\phi=0\).  Derivatives commute with the wave operator, so every
component in \eqref{eq:V77-345-TT-field} is a wave.  The field is
off-diagonal and hence trace free.  Its only possibly nonzero spatial
divergence is the \(z\)-component,
\begin{equation}
 \partial_yH_{yz}+\partial_xH_{xz}
 =\partial_y\partial_x\phi-\partial_x\partial_y\phi=0.
 \label{eq:V77-TT-divergence-check}
\end{equation}
All time components vanish, so the four-dimensional harmonic condition
also holds.

At \(x=0\),
\begin{equation}
 H_{xz}=3\cos(5t)\sin(3y).
\end{equation}
The only required extrinsic-curvature component is
\begin{equation}
 K^{(1)}_{0z}
 =-\frac12\partial_tH_{xz}
 =\frac{15}{2}\sin(5t)\sin(3y).
\end{equation}
Its trace is zero, and raising the time index reverses the sign.  At
\(t_J=\pi/10\) this proves
\eqref{eq:V77-witness-momentum}.

The same dispersion identity proves \(\Box\epsilon_z=0\).  At the joint its
only tangential value is the spatial direction \(z\), which is perpendicular
to \(k=(3,0)\).  Its normal component is zero, so
\(\mathsf R_J(\epsilon)=0\).  Lemma~\ref{lem:V77-trace-rank} places it in
the cocycle radical.  Finally substitute
\eqref{eq:V77-witness-momentum} in
\eqref{eq:V77-radical-charge} and use
\begin{equation}
 \int_0^{2\pi/3}\sin^2(3y)\,\dd y=\frac{\pi}{3}.
 \label{eq:V77-sine-square-integral}
\end{equation}
The result is \eqref{eq:V77-nonzero-radical-charge}, with the joint
orientation carried by \(\sigma_J\).
\end{proof}

\begin{corollary}[The V76 image test fails after premature cocycle minimalization]
\label{cor:V77-image-test-fails}
At a nonzero joint frequency, the cocycle-minimal quotient of the V74
opposite-cocycle carrier does not make every V59 radiative datum admissible.  In the notation of
Proposition~\ref{prop:V76-finite-slab-obstruction-pair},
\begin{equation}
 \chi\ne0
 \label{eq:V77-chi-nonzero}
\end{equation}
on the physical graph.  Therefore the compact-slab V59 pullback cannot be claimed after discarding
the radical merely because it contributes zero to the cocycle.
\end{corollary}

\begin{proof}
The residual parameter in
Theorem~\ref{thm:V77-TT-radical-witness} lies in
\(\operatorname{Rad}K_{J,k}\), while its metric moment is nonzero.
Lemma~\ref{lem:V76-obstruction-class} identifies this value with the
restriction of \(\chi\) to the radical.  No Hamiltonian on the rank-four cocycle-minimal quotient can cancel it.
The full unreduced V74 phase is examined next.
\end{proof}

\begin{firewall}[This does not alter the joint-free V76 theorem]
\label{fire:V77-joint-free-consistency}
The witness is supported at a finite temporal joint.  On the V59 generator
geometry \(\partial\Gamma_T=\varnothing\), the two oriented joint
evaluations are absent and
Theorem~\ref{thm:V77-improved-joint-charge} gives zero.  Thus
Theorem~\ref{thm:V76-unitary-radiative-quotient} and the joint-free
skew-adjoint generator remain unchanged.
\end{firewall}

\subsection{The retained V74 transverse pair and charge completeness}
\label{subsec:V77-retained-transverse-pair}

At \(k\ne0\), let
\begin{equation}
 P_{\perp,k}{}^A{}_B
 :=\delta^A{}_B-\frac{k^Ak_B}{|k|^2}
 \label{eq:V77-transverse-projector}
\end{equation}
be the rank-one orthogonal projector on the joint.  Apply it to the spatial
components of the \emph{existing} V74 corner fields:
\begin{equation}
 X_\perp^A:=P_{\perp,k}{}^A{}_BX^B,
 \qquad
 \Pi_{{\rm cor},\perp}^A
 :=P_{\perp,k}{}^A{}_B\Pi^B.
 \label{eq:V77-existing-transverse-pair}
\end{equation}
The corresponding missing residual trace is
\begin{equation}
 \mathsf S_{J,k}(\epsilon)^A
 :=P_{\perp,k}{}^A{}_B\,
   \mathsf T_J(\epsilon)^B.
 \label{eq:V77-S-trace}
\end{equation}

\begin{proposition}[Symplectic decomposition of the full V74 presentation]
\label{prop:V77-V74-full-decomposition}
At a nonzero joint Fourier mode, the full three-pair V74 phase decomposes
symplectically as
\begin{equation}
 \mathcal P_{{\rm cor},J,k}^{\rm full}
 \cong
 \mathcal P_{{\rm cor},J,k}^{\rm coc,min}
 \oplus
 \operatorname{span}\{X_\perp,\Pi_{{\rm cor},\perp}\},
 \label{eq:V77-V74-symplectic-decomposition}
\end{equation}
where the first summand has dimension four and realizes the nondegenerate
part of \(-K_{J,k}\), while the second is one canonical pair with form
\begin{equation}
 \omega_\perp
 \bigl((X_1,\Pi_1),(X_2,\Pi_2)\bigr)
 =\frac1\kappa
 \left(
  \Pi_1\mathbin{\cdot}X_2-X_1\mathbin{\cdot}\Pi_2
 \right).
 \label{eq:V77-transverse-symplectic-form}
\end{equation}
The V74 residual action and moment map on the second summand are
\begin{align}
 \mathcal V_\epsilon X_\perp
 &=\mathsf S_{J,k}(\epsilon),
 &
 \mathcal V_\epsilon\Pi_{{\rm cor},\perp}&=0,
 \label{eq:V77-existing-transverse-action}\\
 \mu_\epsilon^\perp
 &=-\frac1\kappa
   \mathsf S_{J,k}(\epsilon)
   \mathbin{\cdot}\Pi_{{\rm cor},\perp}.
 \label{eq:V77-existing-transverse-moment}
\end{align}
Thus the transverse pair was present in V74 and retained in V75; it is lost
only if one performs the radical quotient from the V74 minimality theorem
before solving the moment equation.
\end{proposition}

\begin{proof}
Use the orthogonal joint decomposition
\begin{equation*}
 \mathbb C^3
 =\operatorname{span}\{e_0,e_\parallel\}
  \oplus\operatorname{span}\{e_\perp\}
\end{equation*}
for both \(X_b\) and \(\Pi^b\).  The canonical V74 form is the orthogonal
direct sum of the corresponding cotangent pairs.  By
\eqref{eq:V77-R0-symbol}--\eqref{eq:V77-RA-symbol}, the residual momentum
shift \(-\mathsf R_J(\epsilon)\) has no transverse component.  Projecting
\eqref{eq:V74-corner-residual-action} therefore gives
\eqref{eq:V77-existing-transverse-action}.  Projecting
\eqref{eq:V74-corner-moment-map} gives
\eqref{eq:V77-existing-transverse-moment}.  The remaining two pairs realize
the two canonical blocks of \eqref{eq:V77-K-matrix}; the last pair realizes
its radical translation and conjugate momentum.
\end{proof}

\begin{lemma}[The retained transverse action has zero cocycle]
\label{lem:V77-transverse-zero-cocycle}
The transverse action is Hamiltonian and satisfies
\begin{align}
 \iota_{\mathcal V_\epsilon^\perp}\omega_\perp
 &=\delta\mu_\epsilon^\perp,
 \label{eq:V77-transverse-Hamiltonian-identity}\\
 \omega_\perp(
  \mathcal V_\epsilon^\perp,
  \mathcal V_\zeta^\perp)&=0.
 \label{eq:V77-transverse-zero-cocycle}
\end{align}
Consequently retaining this pair does not alter the V74 cancellation of the
V72 cocycle.
\end{lemma}

\begin{proof}
For a variation \((\delta X_\perp,\delta\Pi_\perp)\),
\begin{equation*}
 \omega_\perp(
  (\mathsf S(\epsilon),0),
  (\delta X_\perp,\delta\Pi_\perp))
 =-\frac1\kappa
   \mathsf S(\epsilon)\mathbin{\cdot}\delta\Pi_\perp,
\end{equation*}
which is the variation of
\eqref{eq:V77-existing-transverse-moment}.  Two residual vectors both have
zero momentum component and hence have zero canonical pairing.
\end{proof}

\begin{theorem}[The full V74 zero level cancels the radical charge]
\label{thm:V77-full-V74-cancels-radical}
Keep the full V74 three-pair presentation in the V75 product phase.  On the
physical Einstein conormal graph, the zero-level equation in the transverse
radical is exactly
\begin{equation}
 \boxed{
 \Pi_{{\rm cor},\perp}^A
 =-P_{\perp,k}{}^A{}_B\,
   \nu_a\Pi^{aB}_{(1)}(H).}
 \label{eq:V77-corner-Pi-shell}
\end{equation}
It cancels the charge \eqref{eq:V77-radical-charge}, including the witness
\eqref{eq:V77-nonzero-radical-charge}, without changing the total residual
cocycle or the V75 off-shell master equation.
\end{theorem}

\begin{proof}
For a radical trace, the metric term
\eqref{eq:V77-radical-charge} and the projected V74 moment
\eqref{eq:V77-existing-transverse-moment} sum to
\begin{equation}
 -\frac1\kappa
 \mathsf S(\epsilon)\mathbin{\cdot}
 \left(
  P_{\perp,k}\nu_a\Pi^{a\cdot}_{(1)}
  +\Pi_{{\rm cor},\perp}
 \right).
\end{equation}
The transverse trace is arbitrary, so its coefficient vanishes exactly on
\eqref{eq:V77-corner-Pi-shell}.  Lemma
\ref{lem:V77-transverse-zero-cocycle} shows that the retained pair adds no
cocycle.  It is already a canonical factor in the V74 moment map and hence
in \(\Omega_{75}^{\rm off}\), whose master equation was proved in V75.
\end{proof}

\begin{theorem}[The nonzero-frequency image obstruction is removed by the
full V74 phase]
\label{thm:V77-image-obstruction-removed}
At every nonzero joint Fourier covector, every harmonic radiative
trace--conormal class has a representative on the zero level of the full
V74/V75 residual moment map.  More precisely:
\begin{enumerate}[label=\textup{(\alph*)},leftmargin=2.8em]
\item the rank-four V74 sector solves the charge equation on the two
      nondegenerate cocycle blocks;
\item equation \eqref{eq:V77-corner-Pi-shell} solves it on the transverse
      trace radical;
\item any residual jet with
      \(\mathsf T_J=\mathsf R_J=0\) has zero physical joint charge;
\item after imposing the zero level and dividing by the residual action,
      \((X_\perp,\Pi_{{\rm cor},\perp})\) contributes no independent
      radiative fibre.
\end{enumerate}
Thus the V76 compact-slab image test
\eqref{eq:V76-image-test} passes at \(k\ne0\) precisely when the full V74
presentation retained by V75 is used.
\end{theorem}

\begin{proof}
Part (a) is the invertibility of the two canonical blocks in
\eqref{eq:V77-K-matrix}.  Part (b) is
Theorem~\ref{thm:V77-full-V74-cancels-radical}.

For (c), vanishing of \(\mathsf R_J\) at \(k\ne0\) implies
\(f=\dot f=0\) by
\eqref{eq:V77-R0-symbol}--\eqref{eq:V77-RA-symbol}.  Every term containing
\(f\) or its tangential first derivatives in
\eqref{eq:V77-improved-joint-charge} then vanishes.  Vanishing of
\(\mathsf T_J\) kills the Brown--York term, so the whole joint charge is
zero.

For (d), the zero level fixes \(\Pi_{{\rm cor},\perp}\) uniquely as a
function of the metric momentum.  The residual translation is onto the
one-dimensional \(X_\perp\) line, so quotienting removes \(X_\perp\).
There is no remaining homogeneous coordinate or momentum in this pair.
Together the four statements give a zero-level lift for every radiative
class, unique modulo the full residual orbit.
\end{proof}

\begin{theorem}[Cocycle-minimal versus charge-complete minimal dimension]
\label{thm:V77-charge-complete-minimality}
At one nonzero joint frequency, a symplectic corner realization which both
cancels \(-K_{J,k}\) and represents the transverse radical faithfully with
an arbitrary Hamiltonian charge has dimension at least six.  The full V74
three-pair phase has dimension six and attains this lower bound.
\end{theorem}

\begin{proof}
The nondegenerate part of \(K_{J,k}\) has rank four, so the images of four
quotient generators span a four-dimensional symplectic subspace in every
minimal cocycle realization.  A fundamental vector representing the
radical must pair to zero with all four of those images.  In a
four-dimensional carrier their symplectic orthogonal is zero, so a faithful
radical vector cannot lie there.  It requires a new nonzero orthogonal
vector.  Nondegeneracy of the ambient skew form then requires a conjugate
vector as well, giving at least two further dimensions.  The lower bound is
therefore six.  The decomposition
\eqref{eq:V77-V74-symplectic-decomposition} shows that the full V74 phase
consists of exactly those three canonical pairs and attains the bound.
\end{proof}

\begin{corollary}[The retained pair adds no form defect on its
zero-coordinate slice]
\label{cor:V77-no-new-form-defect}
Every full V74 residual orbit at \(k\ne0\) has a unique representative with
\begin{equation}
 X_\perp=0.
 \label{eq:V77-Xperp-zero-slice}
\end{equation}
On this slice the pullback of \(\omega_\perp\) to physical variations
vanishes, although \(\Pi_{{\rm cor},\perp}\) varies with the metric through
\eqref{eq:V77-corner-Pi-shell}.  Hence retaining the radical pair removes
the image obstruction without adding a term to the remaining V76 comparison
form \(\Delta_{76}\).
\end{corollary}

\begin{proof}
Surjectivity of the residual shift on the transverse line gives the unique
gauge \(X_\perp=0\).  Tangent vectors to this slice have the form
\((0,\delta\Pi_{{\rm cor},\perp})\).  Two such vectors have zero pairing in
\eqref{eq:V77-transverse-symplectic-form}.  Thus this canonical pair has
zero pullback two-form on the slice.
\end{proof}

\begin{proposition}[The V75 off-shell charge already contains the repair]
\label{prop:V77-V75-already-complete}
On the V75 \(H^s\)-domain, \(s>3\), the projection of its BRST differential
to the transverse pair is
\begin{equation}
 sX_\perp=P_{\perp,k}\mathsf T_J(\gamma),
 \qquad
 s\Pi_{{\rm cor},\perp}=0.
 \label{eq:V77-existing-transverse-BRST}
\end{equation}
The corresponding term of \(\Omega_{75}^{\rm off}\) is
\begin{equation}
 -\frac1\kappa
 \left\langle
  P_{\perp,k}\mathsf T_J(\gamma),
  \Pi_{{\rm cor},\perp}
 \right\rangle_J.
 \label{eq:V77-existing-transverse-charge}
\end{equation}
No new V77 field or BFV term is required.  The full V75 master equation and
its commuting product with the V70 holonomy charge remain valid.
\end{proposition}

\begin{proof}
Project the V74/V75 identities
\(sX_b=\gamma_b|_J\) and
\(s\Pi^b=-\mathsf R_J(\gamma)^b\).  The transverse projection of
\(\mathsf R_J\) vanishes by
\eqref{eq:V77-RA-symbol}.  Projection of the V74 corner moment gives
\eqref{eq:V77-existing-transverse-charge}.  V75 proved the master equation
for the full, unreduced V74 phase; no alteration is being made here.
\end{proof}

\begin{firewall}[Why the V74 radical quotient must be delayed]
\label{fire:V77-delay-radical-quotient}
Theorem~\ref{thm:V74-corner-minimality} is correct as a statement about the
smallest carrier of the two-cocycle.  Theorem
\ref{thm:V77-TT-radical-witness} shows that its radical can nevertheless
carry a nonzero Hamiltonian on the physical metric graph.  Therefore one
must impose the full V74 moment equations and quotient their residual action
before discarding any cocycle-null corner pair.  This is the same
shared-data rule used in V75 for the Nakanishi--Lautrup field: algebraic
elimination before the complete graph is formed loses a necessary equation.
\end{firewall}

\begin{firewall}[Frequency and locality scope]
\label{fire:V77-projector-scope}
The projector \eqref{eq:V77-transverse-projector} is used only to prove the
nonzero-frequency decomposition and minimality.  The underlying full V74
fields \((X_b,\Pi^b)\), their action, and their moment map are local and need
no projector.  V77 has not classified the zero joint mode, for which
\eqref{eq:V77-R0-symbol}--\eqref{eq:V77-RA-symbol} lose rank, nor proved a
finite-energy joint trace or the nonlinear curved counterpart of the full
moment equation.
\end{firewall}
\subsection{What remains of the compact-slab pullback}
\label{subsec:V77-remaining-form-test}

The V76 compact comparison had two logically independent requirements.
Corollary~\ref{cor:V77-image-test-fails} shows that the first fails for the
cocycle-minimal quotient, while
Theorem~\ref{thm:V77-image-obstruction-removed} proves that the full V74
carrier retained in V75 passes it at every \(k\ne0\).  The second requirement remains the
form identity
\begin{equation}
 \Delta_{76}(r_1,r_2)=0.
 \label{eq:V77-remaining-Delta-test}
\end{equation}
Corollary~\ref{cor:V77-no-new-form-defect} proves that the retained transverse
pair contributes zero to this test in its canonical zero-coordinate slice.
Thus V78 must compute only the original V74 dressing on the rank-four trace
quotient and compare its canonical two-form with the joint term left by the
V63/V59 Green identity.

\begin{firewall}[Image completion is not yet symplectic identification]
\label{fire:V77-image-versus-form}
Solving the moment equation for every radiative datum proves surjectivity of
the reduced relation.  It does not by itself prove that the transported
two-form equals the V59 two-form.  That equality is exactly
\eqref{eq:V77-remaining-Delta-test}; it remains open and will not be inferred
from equivariance or dimension counts.
\end{firewall}

\subsection{V77 result ledger and next acquisition}
\label{subsec:V77-ledger}

\paragraph{New exact conclusions.}
\begin{enumerate}[label=\textup{(\arabic*)},leftmargin=2.8em]
\item The V72 metric moment on a vacuum harmonic conormal graph is the
      explicit Brown--York joint charge
      \eqref{eq:V77-improved-joint-charge}.
\item At every nonzero temporal-joint Fourier mode, the residual trace image
      has dimension five, the cocycle has rank four, and its trace radical is
      the one-dimensional spatially transverse tangential trace.
\item The exact \(3\)-\(4\)-\(5\) TT wave
      \eqref{eq:V77-345-TT-field} has a nonzero charge under that radical,
      proving that the V74 minimal cocycle carrier fails the V76 compact
      image test.
\item The third canonical pair already present in the full V74 phase is
      necessary and sufficient to absorb the missed charge.  Its residual
      action has zero cocycle, so the V74 cancellation is preserved.
\item The full V74 zero-level equation fixes its transverse corner momentum
      to the transverse Brown--York joint momentum, and quotienting removes
      its coordinate.
      It adds no independent radiative fibre.
\item The full charge-complete V74 carrier makes every nonzero-frequency
      radiative datum admissible without changing the nilpotent V75 boundary
      charge.
\item In the zero-coordinate slice the retained pair contributes no reduced
      two-form, leaving only the original V74 dressing in the compact
      symplectic comparison.
\end{enumerate}

\paragraph{Next forced step.}
V78 should solve the V74 rank-four dressing equation explicitly in the
\((T_0,T_\parallel;R_0,R_\parallel)\) basis, substitute that solution into
\(\omega_{\rm cor}\), and compute the full defect
\(\Delta_{76}\) against the V63/V59 Green form.  If it vanishes, the
nonzero-frequency compact-slab Lagrangian pullback and generator gluing can
be closed.  If it does not, the missing term must be sought upstream in the
geometric joint polarization rather than cancelled by an arbitrary scalar.
The zero mode and a local curved realization remain separate.  The next
mandatory NS/YM companion audit remains V80.

\begin{assessment}[V77 continuum status]
The full Standard--Model--Einstein continuum remains unproved.  V77
evaluates the first compact-slab reduction test, proves that the V74
cocycle-minimal phase misses a charged transverse radical, and proves that
the unreduced three-pair phase already retained by V75 is the smallest
charge-complete realization.  The compact image problem is
therefore closed at nonzero joint frequency.  The compact form identity,
zero joint mode, finite-energy joint trace, nonlinear curved corner completion,
nonlinear Brown--York/Bianchi stress identity, common renormalized Standard
Model stress domain, and cofinal interacting continuum remain open.
\end{assessment}

\section*{V77 exact checks and append-only record}
\addcontentsline{toc}{section}{V77 exact checks and append-only record}

The trace-space matrix \eqref{eq:V77-K-matrix} was row-reduced over the
integers: its rank is four and its kernel is generated by the third basis
vector.  The \(3\)-\(4\)-\(5\) witness was checked componentwise:
\begin{equation}
 5^2-3^2-4^2=0,\qquad
 \partial_yH_{yz}+\partial_xH_{xz}=0,\qquad
 \Pi^{0z}_{(1)}=-\frac{15}{2}\sin(5t)\sin(3y).
 \label{eq:V77-witness-check-card}
\end{equation}
At \(t=\pi/10\), exact integration over one \(y\)-period gives
\eqref{eq:V77-nonzero-radical-charge}.  These checks audit the explicit
coefficients; the proofs above are analytic.

The V76 predecessor used for this append had SHA--256
\begin{center}
\texttt{86790a2c1cf8d3f463675c77b104f13b9be6956f3b13000d9a6f3d643acee7a1}.
\end{center}
The V77 candidate retains its bytes up to the sole live final document
terminator, appends this material, and restores exactly one active
\verb|\end{document}|.  Validation includes balanced environments and
braces, unique labels, resolved internal references, valid inline and display
math delimiters, exact predecessor-prefix equality, and a final inventory
which removes only the superseded SM V76 file.

% =============================================================================
% END APPENDED VERSION V77 MATERIAL
% =============================================================================




% =============================================================================
% BEGIN APPENDED VERSION V78 MATERIAL
% =============================================================================
\clearpage
\section*{V78 research note: explicit V74 dressing, the nonzero compact
corner two-form, and the forced cap--interface gluing equation}
\addcontentsline{toc}{section}{V78 explicit dressing and compact form defect}
\label{sec:V78-dressing-form-defect}

\subsection*{Purpose and relation to V76--V77}

V76 separated the compact-slab comparison into an image test and a
symplectic-form test.  V77 evaluated the image test.  It proved that the
rank-four cocycle-minimal quotient discards a charge-carrying transverse
direction, while the full local three-pair V74 phase retained by V75 is the
smallest charge-complete carrier.  Every nonzero-frequency radiative datum
therefore has a zero-level lift.

It remains to compute the form transported by that lift.  This note solves
all six V74 corner variables on the physical V63 graph, fixes the residual
transverse coordinate, and substitutes the result into the V74 canonical
form.  The transverse pair found in V77 contributes zero, but the two
nondegenerate cocycle pairs leave a scalar plus-polarization endpoint
two-form.  That form is generically nonzero on the full two-branch radiative
phase.  Hence the completed compact side reduction is not symplectically
identical to the bare V59 side phase by itself.

This does not affect the joint-free V59 generator.  It identifies the next
upstream equation: the complete oriented cap and continuum endpoint
potential must contribute the opposite two-form.  A scalar Hayward
counterterm cannot do so because its field-space curvature is zero.  The
calculation below gives the exact form which the cap--interface gluing
construction must reproduce.

\subsection{Physical temporal-joint reduction of the V77 charge}
\label{subsec:V78-physical-joint-charge}

Work on a temporal joint
\(J=\{t=t_J\}\times\Sigma\) of a flat timelike face and put
\begin{equation}
 \nu_a=\sigma_J\delta_a^0,
 \qquad
 \sigma_J\in\{-1,1\}.
 \label{eq:V78-joint-orientation}
\end{equation}
Let
\begin{equation}
 \mathcal Q:=(-\Delta_J)^{1/2}
 \label{eq:V78-Q-operator}
\end{equation}
on the mean-zero joint sector.  Thus
\(\mathcal Q^{-1}\) and \(\Delta_J^{-1}\) are bounded on each closed
nonzero-frequency cone and are defined fibrewise throughout this note.
Set
\begin{equation}
 \mathcal L_A:=D_A\Delta_J^{-1}.
 \label{eq:V78-longitudinal-lift}
\end{equation}
On the flat joint,
\begin{align}
 D^A\mathcal L_A&=I,
 \label{eq:V78-L-right-inverse}\\
 \sum_A\mathcal L_A^*\mathcal L_A&=\mathcal Q^{-2}.
 \label{eq:V78-L-Gram}
\end{align}

For the induced metric write
\begin{equation}
 \varsigma:=q_A{}^A.
 \label{eq:V78-spatial-joint-trace}
\end{equation}
The V63 physical graph has
\begin{equation}
 q_{00}=0,
 \qquad
 D_Aq^{0A}=0,
 \qquad
 \tau=\varsigma.
 \label{eq:V78-physical-induced-identities}
\end{equation}
Indeed the plus representative has \(q_{0A}=0\), while the completed cross
representative has only the component transverse to the joint Fourier
covector, so its tangential divergence vanishes.  The cross representative
is trace free.

\begin{lemma}[Physical form of the improved joint charge]
\label{lem:V78-physical-joint-charge}
On the V63 graph, the V77 charge becomes
\begin{equation}
 \boxed{
 \mu_\epsilon(H)
 =\frac{\sigma_J}{2\kappa}\int_J
 \left[
  fD_0\varsigma-(D_0f)\varsigma
  -2\epsilon_b\Pi^{0b}_{(1)}
 \right].}
 \label{eq:V78-physical-joint-charge}
\end{equation}
For a nonzero joint Fourier mode the V74 normal trace is
\begin{align}
 \mathsf R^0&=\sigma_J\mathcal Q^2f,
 \label{eq:V78-R0-operator}\\
 \mathsf R^A&=\sigma_JD_0D^Af.
 \label{eq:V78-RA-operator}
\end{align}
\end{lemma}

\begin{proof}
Insert \eqref{eq:V78-physical-induced-identities} in
\eqref{eq:V77-improved-joint-charge}.  The two terms containing
\(D_Aq^{0A}\) vanish, \(q^{00}=0\), and
\(\tau=\varsigma\).  The remaining normal terms are the Wronskian in
\eqref{eq:V78-physical-joint-charge}.  Equations
\eqref{eq:V78-R0-operator}--\eqref{eq:V78-RA-operator} are
\eqref{eq:V77-R0-symbol}--\eqref{eq:V77-RA-symbol} before choosing a
longitudinal Fourier frame.
\end{proof}

\subsection{Exact solution of the full V74 dressing equation}
\label{subsec:V78-dressing-solution}

Denote the full V74 corner momentum by
\(\Pi_{\rm cor}^b\), reserving \(\Pi_{(1)}^{ab}\) for the Brown--York
momentum.  Its moment map is
\begin{equation}
 \mu_\epsilon^{\rm cor}
 =-\frac1\kappa\int_J
 \left[
  \mathsf R^b(\epsilon)X_b
  +\mathsf T_b(\epsilon)\Pi_{\rm cor}^b
 \right].
 \label{eq:V78-full-corner-moment}
\end{equation}

\begin{theorem}[All nonzero-frequency dressing rows]
\label{thm:V78-complete-dressing}
On the V63 physical graph and at every nonzero joint frequency, the zero
level
\begin{equation}
 \mu_\epsilon+\mu_\epsilon^{\rm cor}=0
 \quad\hbox{for every residual }\epsilon
 \label{eq:V78-zero-level}
\end{equation}
is equivalent to
\begin{align}
 X_0
 &=\frac12\mathcal Q^{-2}D_0\varsigma,
 \label{eq:V78-X0-dressing}\\
 D^AX_A&=\frac12\varsigma,
 \label{eq:V78-divX-dressing}\\
 \Pi_{\rm cor}^b
 &=-\sigma_J\Pi_{(1)}^{0b}.
 \label{eq:V78-Picor-dressing}
\end{align}
The transverse residual action fixes the unique quotient representative
\begin{equation}
 X_A=\frac12\mathcal L_A\varsigma,
 \qquad
 P_\perp{}^A{}_BX^B=0.
 \label{eq:V78-X-longitudinal-gauge}
\end{equation}
Equations \eqref{eq:V78-X0-dressing}--%
\eqref{eq:V78-X-longitudinal-gauge} solve the entire six-component V74
corner phase, including the transverse momentum row proved separately in
V77.
\end{theorem}

\begin{proof}
First vary the arbitrary tangential trace
\(\mathsf T_b(\epsilon)=\epsilon_b|_J\).  Its coefficient in
\eqref{eq:V78-physical-joint-charge} and
\eqref{eq:V78-full-corner-moment} is
\begin{equation*}
 -\frac1\kappa
 \epsilon_b
 \left(\sigma_J\Pi_{(1)}^{0b}+\Pi_{\rm cor}^b\right).
\end{equation*}
It vanishes for every \(\epsilon_b\) exactly when
\eqref{eq:V78-Picor-dressing} holds.

For the normal residual component, insert
\eqref{eq:V78-R0-operator}--\eqref{eq:V78-RA-operator}.  After dividing by
the common factor \(\sigma_J/\kappa\), the coefficient is
\begin{equation}
 \frac12\int_J
 \left[fD_0\varsigma-(D_0f)\varsigma\right]
 -\int_J
 \left[\mathcal Q^2f\,X_0
       +(D_0D^Af)X_A\right].
 \label{eq:V78-normal-dressing-functional}
\end{equation}
Integrate the last spatial derivative on the closed joint.  The independent
coefficients of \(f\) and \(D_0f\) are
\begin{equation}
 \frac12D_0\varsigma-\mathcal Q^2X_0,
 \qquad
 -\frac12\varsigma+D^AX_A.
 \label{eq:V78-normal-dressing-coefficients}
\end{equation}
Their vanishing gives
\eqref{eq:V78-X0-dressing} and
\eqref{eq:V78-divX-dressing}.  The longitudinal solution of the second row
is \(\frac12\mathcal L_A\varsigma\) by
\eqref{eq:V78-L-right-inverse}.  Its transverse part is exactly the V77
radical coordinate and is removed uniquely by the transverse residual
translation.  This proves
\eqref{eq:V78-X-longitudinal-gauge}.
\end{proof}

\begin{corollary}[No hidden homogeneous corner variable at nonzero frequency]
\label{cor:V78-no-hidden-corner-variable}
For fixed nonzero joint frequency, the zero level fixes
\(X_0\), the longitudinal part of \(X_A\), and all components of
\(\Pi_{\rm cor}^b\).  The sole undetermined coordinate is the transverse
part of \(X_A\), and the residual quotient removes it.  Hence the full V74
corner phase adds no independent nonzero-frequency physical class.
\end{corollary}

\begin{proof}
This is the component count in
Theorem~\ref{thm:V78-complete-dressing}: five variables are fixed by the
moment rows and the sixth is one residual orbit coordinate.
\end{proof}

\subsection{Pullback of the V74 canonical form}
\label{subsec:V78-corner-form-pullback}

Use the V74 canonical potential
\begin{equation}
 \vartheta_{\rm cor}
 :=\frac1\kappa\int_J
  \Pi_{\rm cor}^b\,\delta X_b,
 \qquad
 \delta\vartheta_{\rm cor}=\omega_{\rm cor}.
 \label{eq:V78-corner-potential}
\end{equation}
The transverse term vanishes on
\eqref{eq:V78-X-longitudinal-gauge}, in agreement with
Corollary~\ref{cor:V77-no-new-form-defect}.

\begin{lemma}[Brown--York rows on the physical plus trace]
\label{lem:V78-BY-plus-rows}
On the V53 TT plus representative, and hence on the plus part of the V63
physical graph,
\begin{align}
 \Pi_{(1)}^{00}
 &=-\frac12D_n\varsigma,
 \label{eq:V78-Pi00-plus}\\
 P_\parallel{}^A{}_B\Pi_{(1)}^{0B}
 &=\frac12D_0\mathcal L^A D_n\varsigma.
 \label{eq:V78-Pi0A-plus}
\end{align}
The completed cross representative has
\begin{equation}
 \varsigma=0,\qquad
 \Pi_{(1)}^{00}=0,\qquad
 P_\parallel{}^A{}_B\Pi_{(1)}^{0B}=0.
 \label{eq:V78-cross-no-longitudinal-row}
\end{equation}
\end{lemma}

\begin{proof}
For the plus representative \(H_{0\mu}=0\), its spatial trace is zero, and
spatial transversality gives
\begin{equation}
 H_{nn}=-\varsigma,
 \qquad
 D^AH_{An}=D_n\varsigma.
 \label{eq:V78-plus-TT-relations}
\end{equation}
The longitudinal solution of the second identity is
\(H_{An}=\mathcal L_AD_n\varsigma\).  The linear mean curvature is
\begin{equation*}
 K_{(1)}
 =\frac12D_n\varsigma-D^AH_{An}
 =-\frac12D_n\varsigma.
\end{equation*}
Since \(K_{00}^{(1)}=0\),
\(\Pi^{00}_{(1)}=K_{(1)}\), proving
\eqref{eq:V78-Pi00-plus}.  Also
\begin{equation*}
 K_{0A}^{(1)}=-\frac12D_0H_{An},
 \qquad
 \Pi^{0A}_{(1)}=-K_{0A}^{(1)}
 =\frac12D_0\mathcal L^AD_n\varsigma,
\end{equation*}
which proves \eqref{eq:V78-Pi0A-plus}.

For the V63 cross representative, the induced trace is zero.  Its only
time--space Brown--York component is transverse to the joint covector, as
recorded in \eqref{eq:V63-cross-K-components}.  The mean curvature and
\(00\) component vanish.  This gives
\eqref{eq:V78-cross-no-longitudinal-row}.
\end{proof}

\begin{theorem}[Exact pulled-back corner potential and two-form]
\label{thm:V78-pulled-corner-form}
On the complete V63 physical graph, in the quotient gauge
\eqref{eq:V78-X-longitudinal-gauge}, the V74 corner potential is
\begin{equation}
 \boxed{
 \vartheta_{\rm cor}^{\rm phys}
 =\frac{\sigma_J}{4\kappa}
 \left[
  \left\langle
   D_n\varsigma,\mathcal Q^{-2}\delta D_0\varsigma
  \right\rangle_J
  -
  \left\langle
   D_0D_n\varsigma,\mathcal Q^{-2}\delta\varsigma
  \right\rangle_J
 \right].}
 \label{eq:V78-pulled-corner-potential}
\end{equation}
Its field-space exterior derivative is
\begin{align}
 \omega_{\rm cor}^{\rm phys}(s_1,s_2)
 &=\frac{\sigma_J}{4\kappa}
 \bigl[
  \langle D_ns_1,\mathcal Q^{-2}D_0s_2\rangle_J
 -\langle D_ns_2,\mathcal Q^{-2}D_0s_1\rangle_J
 \nonumber\\
 &\hspace{5.2em}
 -\langle D_0D_ns_1,\mathcal Q^{-2}s_2\rangle_J
 +\langle D_0D_ns_2,\mathcal Q^{-2}s_1\rangle_J
 \bigr],
 \label{eq:V78-pulled-corner-two-form}
\end{align}
where \(s_i=\delta_i\varsigma\).  The cross polarization and every
plus--cross pairing vanish in this form; only the plus trace occurs.
\end{theorem}

\begin{proof}
Substitute \eqref{eq:V78-X0-dressing},
\eqref{eq:V78-X-longitudinal-gauge}, and
\eqref{eq:V78-Picor-dressing} in
\eqref{eq:V78-corner-potential}.  Equations
\eqref{eq:V78-Pi00-plus}--\eqref{eq:V78-Pi0A-plus} give
\begin{align*}
 \Pi_{\rm cor}^0\delta X_0
 &=\frac{\sigma_J}{4}
   (D_n\varsigma)\,
   \mathcal Q^{-2}\delta D_0\varsigma,\\
 \left\langle
  \Pi_{{\rm cor},\parallel}^A,\delta X_A
 \right\rangle_J
 &=-\frac{\sigma_J}{4}
 \left\langle
  D_0\mathcal L^AD_n\varsigma,
  \mathcal L_A\delta\varsigma
 \right\rangle_J\\
 &=-\frac{\sigma_J}{4}
 \left\langle
  D_0D_n\varsigma,
  \mathcal Q^{-2}\delta\varsigma
 \right\rangle_J,
\end{align*}
where the last equality is \eqref{eq:V78-L-Gram}.  The transverse potential
is zero because \(X_\perp=0\).  The cross polarization has no
\(\varsigma\), \(00\), or longitudinal momentum row by
\eqref{eq:V78-cross-no-longitudinal-row}.  This proves
\eqref{eq:V78-pulled-corner-potential}.  Antisymmetrizing its linearization
in two variations gives
\eqref{eq:V78-pulled-corner-two-form}.
\end{proof}

\begin{firewall}[A one-branch calculation cannot test this two-form]
\label{fire:V78-one-branch-insufficient}
On a single frozen normal characteristic branch and one fixed polarization,
all trace and momentum components are scalar multiples of one amplitude.
Every alternating two-form on that one-dimensional amplitude line vanishes.
Equation~\eqref{eq:V78-pulled-corner-two-form} tests the full physical phase,
where independent incoming and outgoing normal data, or equivalently
independent Cauchy position and velocity data, are present.  A zero obtained
after selecting one branch is therefore not evidence that the compact form
defect vanishes.
\end{firewall}

\subsection{An exact nonzero two-branch witness}
\label{subsec:V78-two-branch-witness}

Take a flat joint torus with coordinates \(y,z\), let \(y\) have period
\(2\pi/q\), let the \(z\)-period have length \(L_z\), and choose
\begin{equation}
 q>0,\qquad \rho>0,\qquad
 \omega:=\sqrt{q^2+\rho^2}.
 \label{eq:V78-witness-frequencies}
\end{equation}
Consider the two scalar plus traces
\begin{align}
 s_1(t,y,x)
 &:=\cos(qy)\sin(\rho x)\cos(\omega t),
 \label{eq:V78-witness-s1}\\
 s_2(t,y,x)
 &:=\cos(qy)\cos(\rho x)\sin(\omega t).
 \label{eq:V78-witness-s2}
\end{align}
Both solve the scalar wave equation.  At \(t=0,x=0\),
\begin{align}
 s_1&=0,&D_0s_1&=0,&D_ns_1&=\rho\cos(qy),
 &D_0D_ns_1&=0,
 \label{eq:V78-witness-jet1}\\
 s_2&=0,&D_0s_2&=\omega\cos(qy),&D_ns_2&=0,
 &D_0D_ns_2&=0.
 \label{eq:V78-witness-jet2}
\end{align}

\begin{theorem}[The compact corner two-form is nonzero]
\label{thm:V78-nonzero-corner-form}
The traces \eqref{eq:V78-witness-s1}--%
\eqref{eq:V78-witness-s2} are realized by superpositions of the V53 TT plus
polarization with normal frequencies \(\rho\) and \(-\rho\).  Their V74
corner pairing is
\begin{equation}
 \boxed{
 \omega_{\rm cor}^{\rm phys}(s_1,s_2)
 =\sigma_J
  \frac{\pi L_z\,\rho\omega}{4\kappa q^3}\ne0.}
 \label{eq:V78-nonzero-corner-value}
\end{equation}
For the \(3\)-\(4\)-\(5\) choice
\((q,\rho,\omega)=(3,4,5)\), the value is
\begin{equation}
 \omega_{\rm cor}^{\rm phys}(s_1,s_2)
 =\sigma_J\frac{5\pi L_z}{27\kappa}.
 \label{eq:V78-345-corner-value}
\end{equation}
\end{theorem}

\begin{proof}
At fixed nonzero \(q,\rho\), the V53 TT plus splitting has joint spatial
trace
\(\varsigma=(q^2/\omega^2)a_+/\sqrt2\).  This coefficient is nonzero, so
each scalar wave \(s_i\) has a unique plus amplitude.  The sine and cosine
normal factors are the real sums and differences of the two normal
frequencies.  Thus both are genuine two-branch physical variations.

Insert \eqref{eq:V78-witness-jet1}--%
\eqref{eq:V78-witness-jet2} into
\eqref{eq:V78-pulled-corner-two-form}.  Only its first term survives.
Moreover,
\begin{equation}
 \mathcal Q^{-2}\cos(qy)=q^{-2}\cos(qy),
 \qquad
 \int_0^{2\pi/q}\cos^2(qy)\,\dd y=\frac{\pi}{q}.
 \label{eq:V78-witness-integrals}
\end{equation}
Multiplication by the \(z\)-length gives
\eqref{eq:V78-nonzero-corner-value}.  Substitution of \(q=3\),
\(\rho=4\), and \(\omega=5\) gives
\eqref{eq:V78-345-corner-value}.
\end{proof}

\subsection{Evaluation of the V76 form test}
\label{subsec:V78-Delta-evaluation}

Let \(\omega_{59}^{\rm rad}\) be the V59 radiative Green form and let
\(\omega_{36}^{\rm phys}\) be the V36 metric Green form restricted to the
V63 physical graph.

\begin{lemma}[The side metric contribution already has V59 normalization]
\label{lem:V78-side-form-equality}
On every nonzero real propagating fibre,
\begin{equation}
 \omega_{36}^{\rm phys}=\omega_{59}^{\rm rad}.
 \label{eq:V78-side-form-equality}
\end{equation}
The identity extends by linearity to two-branch data and by the V65
holomorphic frame to the connected nonzero stable sheet.
\end{lemma}

\begin{proof}
Proposition~\ref{prop:V63-conormal-from-action} gives the metric Green form
on the physical graph.  Theorem~\ref{thm:V63-exact-real-cone-lift} gives
\(p_{\rm rad}=D_na\) in exact V59 normalization, proving
\eqref{eq:V78-side-form-equality} on each real branch.  Both sides are
bilinear, so the identity holds for their sums.  V65 expresses the same
graph and conormal through a holomorphic nonzero-frequency frame; the
identity theorem gives the stable continuation away from its declared
branch set.
\end{proof}

\begin{theorem}[The compact side-plus-corner form test fails]
\label{thm:V78-Delta-nonzero}
For the full V74 dressing retained in V75, the V76 comparison form on the
compact physical side graph is
\begin{equation}
 \boxed{
 \Delta_{76}
 =\omega_{\rm cor}^{\rm phys}.}
 \label{eq:V78-Delta-equals-corner}
\end{equation}
It is not identically zero.  In particular,
\begin{equation}
 \Delta_{76}(s_1,s_2)
 =\sigma_J
  \frac{\pi L_z\,\rho\omega}{4\kappa q^3}
 \label{eq:V78-Delta-witness}
\end{equation}
on the two-branch witness of
Theorem~\ref{thm:V78-nonzero-corner-form}.  Therefore the inverse image of
the bare V59 side relation is not Lagrangian in the completed compact
side-plus-corner reduced form by itself.
\end{theorem}

\begin{proof}
The definition \eqref{eq:V76-finite-slab-defect-form} is
\begin{equation*}
 \Delta_{76}
 =\omega_{36}^{\rm phys}
  +\omega_{\rm cor}^{\rm phys}
  -\omega_{59}^{\rm rad}.
\end{equation*}
The transverse V74 pair has zero pullback by
Corollary~\ref{cor:V77-no-new-form-defect}.  Lemma
\ref{lem:V78-side-form-equality} cancels the first and third terms, leaving
\eqref{eq:V78-Delta-equals-corner}.  The nonzero value is
\eqref{eq:V78-nonzero-corner-value}.  If the inverse image of the bare V59
relation were Lagrangian, its form would vanish on every two tangent
vectors; the witness contradicts isotropy.
\end{proof}

\begin{corollary}[No scalar joint counterterm can repair the defect]
\label{cor:V78-no-scalar-repair}
No scalar functional of the existing joint fields, including the V74
quadratic Hayward term, can cancel
\(\omega_{\rm cor}^{\rm phys}\) by being added to the action.
\end{corollary}

\begin{proof}
A scalar joint functional \(F\) changes the field-space potential by
\(\delta F\), whose exterior derivative is
\(\delta^2F=0\).  Theorem~\ref{thm:V78-nonzero-corner-form} gives a nonzero
two-form.  V74 independently computed zero curvature for the quadratic
Hayward scalar, which is the geometric instance of the same argument.
\end{proof}

\begin{firewall}[What has failed]
\label{fire:V78-failure-scope}
The failure is not the V59 positive action, the V63 conormal normalization,
the V74 master equation, or the V77 image completion.  It is the attempted
identification of a compact timelike side together with its independent
corner cotangent phase with the bare joint-free V59 side phase.  A compact
variational problem also has oriented cap potentials and continuum endpoint
data.  Those terms were absent from \(\Delta_{76}\) and must be assembled
before any global gluing conclusion.
\end{firewall}

\subsection{The forced upstream cap--interface equation}
\label{subsec:V78-cap-interface-equation}

Let
\(\omega_{{\rm cap+cont},J}^{\rm edge}\) denote the codimension-two form
obtained by restricting the complete oriented Einstein cap potential, the
V71/V73 compatible-jet cotangent data, and the V59 continuum cap potential
to the common interface joint.  This notation is a target for the next
derivation, not an assumed form.

\begin{theorem}[Necessary and sufficient nonzero-frequency gluing equation]
\label{thm:V78-forced-gluing-equation}
On the nonzero-frequency physical graph, the full compact
metric--continuum relation can have the bare V59 reduced symplectic form if
and only if the complete cap--interface calculation gives
\begin{equation}
 \boxed{
 \omega_{{\rm cap+cont},J}^{\rm edge}
 =-\omega_{\rm cor}^{\rm phys}.}
 \label{eq:V78-forced-cap-cancellation}
\end{equation}
Equivalently, its plus-trace block must be the negative of
\eqref{eq:V78-pulled-corner-two-form}; its cross and plus--cross blocks must
vanish.
\end{theorem}

\begin{proof}
The side metric and V59 forms already agree by
\eqref{eq:V78-side-form-equality}.  The full V74 zero-level dressing adds
exactly \(\omega_{\rm cor}^{\rm phys}\) by
\eqref{eq:V78-Delta-equals-corner}.  The only omitted pieces of the compact
action are the cap and continuum endpoint potentials and their compatible
jet restrictions.  Additivity of the oriented Green form makes the total
defect
\begin{equation*}
 \omega_{\rm cor}^{\rm phys}
 +\omega_{{\rm cap+cont},J}^{\rm edge}.
\end{equation*}
It vanishes exactly under
\eqref{eq:V78-forced-cap-cancellation}.  The polarization statement follows
because Theorem~\ref{thm:V78-pulled-corner-form} has only a plus block.
\end{proof}

\begin{firewall}[The cancellation equation is not yet a proof of
cancellation]
\label{fire:V78-gluing-not-assumed}
Equation~\eqref{eq:V78-forced-cap-cancellation} is necessary and sufficient,
but V78 has not derived its left side.  Defining a symbol for the missing cap
form does not establish the identity.  V79 must vary one combined
side--cap--continuum action, retain all shared joint data until the final
antisymmetrization, and either prove the equality or exhibit a second
nonzero witness.
\end{firewall}

\subsection{V78 result ledger and next acquisition}
\label{subsec:V78-ledger}

\paragraph{New exact conclusions.}
\begin{enumerate}[label=\textup{(\arabic*)},leftmargin=2.8em]
\item On the V63 graph, the normal part of the improved Einstein joint charge
      is the scalar Wronskian
      \(fD_0\varsigma-(D_0f)\varsigma\).
\item The complete V74 zero-level dressing is
      \eqref{eq:V78-X0-dressing}--%
      \eqref{eq:V78-Picor-dressing}; its transverse coordinate is the only
      orbit variable.
\item The plus TT Brown--York rows are
      \eqref{eq:V78-Pi00-plus}--%
      \eqref{eq:V78-Pi0A-plus}; the cross polarization has no longitudinal
      dressing block.
\item Pullback of the full V74 canonical form gives the explicit endpoint
      two-form \eqref{eq:V78-pulled-corner-two-form}.
\item The two-branch scalar waves
      \eqref{eq:V78-witness-s1}--\eqref{eq:V78-witness-s2} prove that this
      form is nonzero.  The exact \(3\)-\(4\)-\(5\) value is
      \eqref{eq:V78-345-corner-value}.
\item Because the V63 side Green form already equals the V59 form, the V76
      compact defect is exactly the nonzero V74 corner form.
\item A scalar joint counterterm cannot cancel the defect.
\item The remaining global question is the concrete cap--interface identity
      \eqref{eq:V78-forced-cap-cancellation}.
\end{enumerate}

\paragraph{Next forced step.}
V79 must derive
\(\omega_{{\rm cap+cont},J}^{\rm edge}\) from the complete oriented action.
It should start from the V73 compatible-jet Einstein--GHY parent on the
spacelike cap, retain the V74 geometric Hayward scalar and full cotangent
phase, and add the V59 continuum cap potential before identifying the common
trace.  The calculation must then be restricted to the V63/V65 physical
graph and compared term by term with
\eqref{eq:V78-pulled-corner-two-form}.  If cancellation fails, the
nonzero witness \eqref{eq:V78-Delta-witness} identifies the missing plus
endpoint channel.  The next mandatory NS/YM audit remains V80.

\begin{assessment}[V78 continuum status]
The full Standard--Model--Einstein continuum remains unproved.  V78 closes
the explicit nonzero-frequency V74 dressing and evaluates the second V76
compact-slab test: the side-plus-corner form differs from the bare V59 form
by the nonzero plus endpoint two-form
\eqref{eq:V78-pulled-corner-two-form}.  The exact cap--continuum
cancellation, zero joint mode, finite-energy joint trace, nonlinear curved
gluing, Brown--York/Bianchi stress identity, common renormalized Standard
Model stress domain, and cofinal interacting continuum remain open.
\end{assessment}

\section*{V78 exact checks and append-only record}
\addcontentsline{toc}{section}{V78 exact checks and append-only record}

The operator identities used in the dressing were checked at a Fourier mode
with
\(\Delta_J=-q^2\):
\begin{equation}
 D^AD_A\Delta_J^{-1}=1,\qquad
 \sum_A(D_A\Delta_J^{-1})^*
        (D_A\Delta_J^{-1})=q^{-2}.
 \label{eq:V78-operator-check-card}
\end{equation}
The witness uses
\begin{equation}
 \omega^2=q^2+\rho^2,\qquad
 \int_0^{2\pi/q}\cos^2(qy)\,\dd y=\frac{\pi}{q},
 \label{eq:V78-witness-check-card}
\end{equation}
and gives \eqref{eq:V78-nonzero-corner-value} directly.  These are exact
symbol checks; the preceding integration-by-parts arguments prove the
operator statements.

The V77 predecessor used for this append had SHA--256
\begin{center}
\texttt{632e82214a0ff970236cc58f7de384454f800263b9a3d90309289b1abb07648c}.
\end{center}
The V78 candidate retains its bytes up to the sole live final document
terminator, appends this material, and restores exactly one active
\verb|\end{document}|.  Validation includes balanced environments and
braces, unique labels, resolved internal references, valid inline and display
math delimiters, exact predecessor-prefix equality, and a final inventory
which removes only the superseded SM V77 file.

% =============================================================================
% END APPENDED VERSION V78 MATERIAL
% =============================================================================
\clearpage
% =============================================================================
% BEGIN APPENDED VERSION V79 MATERIAL
% =============================================================================
\section{V79: the cap transgression cancels the nonzero-frequency corner form}
\label{sec:V79-cap-transgression}

V78 isolated a nonzero joint two-form and left one precise alternative: the
oriented cap and continuum endpoint calculation must produce its negative,
or the proposed compact metric--continuum relation is not isotropic.  This
note performs that calculation on the smooth nonzero-joint-frequency core.
The result is a cancellation, but it is not caused by a scalar counterterm
on the corner.  It comes from the Green boundary term which separates the
Einstein transverse--traceless cap norm from the scalar V59 cap norm when the
plus polarization is written in the local joint-trace coordinate
\(\varsigma\).  At the level of potentials, that Green term and the V74
corner potential add to an exact field-space variation.  Their two-forms
therefore cancel.

The calculation is deliberately limited to the flat linear physical graph
already constructed in V63--V76.  It does not settle the zero joint mode,
finite-energy trace closure, curved corners, the nonlinear constraints, or
the interacting Standard Model stress domain.

\subsection{Oriented cap potentials from the two actions}
\label{subsec:V79-oriented-cap-potentials}

Fix a temporal cap \(\Sigma_t\) and write its outward unit conormal as
\(u=\sigma_J\partial_t\), where \(\sigma_J=+1\) on a terminal cap and
\(\sigma_J=-1\) on an initial cap.  Near its intersection \(J\) with the
timelike side, choose the V59 spatial collar \(x\geq0\), with
\begin{equation}
 D_n=-\partial_x,
 \qquad J=\{x=0\}.
 \label{eq:V79-collar-orientation}
\end{equation}
Thus \(D_n\) is the outward side derivative.  All joint fields below are
smooth and have compact support in the remaining joint variables, or are
periodic as in the V78 witness.  We work on the subspace on which
\(Q=(-\Delta_J)^{1/2}\) is invertible.

\begin{lemma}[Einstein--GHY cap potential on the harmonic physical shell]
\label{lem:V79-Einstein-cap-potential}
Let \(H_{ij}\) be the spatial transverse--traceless part of a physical
linearized metric on \(\Sigma_t\).  In the V63 Einstein action units, the
oriented cap potential and its two-form are
\begin{align}
 \Theta_{E,\Sigma_t}(H)
 &=\frac{\sigma_J}{4\kappa}
   \int_{\Sigma_t}D_0H^{ij}\,\delta H_{ij},
 \label{eq:V79-Einstein-cap-theta}\\
 \Omega_{E,\Sigma_t}(H_1,H_2)
 &=\frac{\sigma_J}{4\kappa}
   \int_{\Sigma_t}
   \bigl(D_0H_1^{ij}H_{2,ij}-D_0H_2^{ij}H_{1,ij}\bigr).
 \label{eq:V79-Einstein-cap-omega}
\end{align}
The gauge-fixing part makes no contribution on the harmonic shell.
\end{lemma}

\begin{proof}
The all-face variation of the Einstein--Hilbert, GHY, and de Donder parent
derived in V36 and retained in V73 is
\begin{equation*}
 \Theta_{36}=\frac1{2\kappa}\int
 \bigl(\Pi_{(1)}^{ab}\delta h_{ab}
       -2Q(C)^{\mu\nu}\delta h_{\mu\nu}\bigr).
\end{equation*}
On the harmonic shell \(C=0\).  On a temporal cap, the induced variables are
the spatial components and
\begin{equation*}
 K_{ij}^{(1)}=\frac{\sigma_J}{2}D_0H_{ij},
 \qquad K^{(1)}=0,
 \qquad
 \Pi_{(1)}^{ij}=\frac{\sigma_J}{2}D_0H^{ij},
\end{equation*}
because \(H\) is spatially trace free.  Substitution proves
\eqref{eq:V79-Einstein-cap-theta}.  Antisymmetrizing its field-space
variation proves \eqref{eq:V79-Einstein-cap-omega}.
\end{proof}

Multiply the V59 quotient action by the common positive factor
\((2\kappa)^{-1}\), as in the V63 normalization.  Its metric radiative
amplitude is denoted by \(a\), and its continuum amplitude by \(b(\rho)\).

\begin{lemma}[The V59 cap has no intrinsic codimension-two curvature]
\label{lem:V79-V59-cap-no-edge}
The oriented V59 cap potential is
\begin{equation}
 \Theta_{59,\Sigma_t}
 =\frac{\sigma_J}{4\kappa}
   \left(
    \int_{\Omega_x}D_0a\,\delta a\,\dd x\,\dd^2y
    +m\int_{\mathbb R^2_y}\int_0^\infty
       D_0b(\rho)\,\delta b(\rho)\,\dd\rho\,\dd^2y
   \right).
 \label{eq:V79-V59-cap-theta}
\end{equation}
Restricting the common interface trace \(b(0)=a|_J\) creates no additional
joint two-form.  In particular, the continuum bath cannot by itself cancel
the V78 corner form.
\end{lemma}

\begin{proof}
Variation of the time kinetic terms in the V59 action gives
\eqref{eq:V79-V59-cap-theta}.  Neither term contains a normal spatial
derivative on the cap.  Pullback to the linear trace subspace
\(b(0)=a|_J\) commutes with \(\delta\), and no integration by parts in
\(x\) or \(\rho\) is required.  Consequently its exterior derivative is
the pullback of the two bulk canonical forms and has no distribution or
separate integral supported on \(J\).
\end{proof}

\subsection{The plus TT Green identity at the joint}
\label{subsec:V79-plus-green-identity}

At a nonzero joint frequency set
\begin{equation}
 B:=Q^{-2},
 \qquad
 A_+:=1-BD_n^2.
 \label{eq:V79-BA-operators}
\end{equation}
The operators \(B,D_0,D_n\) commute on the flat collar, and \(B\) is
self-adjoint on the nonzero-frequency joint core.  Let
\(L_A=D_A\Delta_J^{-1}\), as in V63 and V78.

\begin{lemma}[Local reconstruction of the plus TT tensor]
\label{lem:V79-plus-reconstruction}
Let \(s=\varsigma\) be the tangential trace of the plus polarization on the
side.  In a joint frame with one longitudinal tangential direction and one
transverse direction, its spatial TT tensor is
\begin{align}
 H_{nn}(s)&=-s,
 &H_{nA}(s)&=L_A D_ns,
 \label{eq:V79-plus-H-normal}\\
 H_{\parallel\parallel}(s)&=BD_n^2s,
 &H_{\perp\perp}(s)&=s-BD_n^2s,
 \label{eq:V79-plus-H-tangential}
\end{align}
with all other plus components zero.  Its norm-one V53 radiative amplitude
is
\begin{equation}
 a_+(s)=\sqrt2\,A_+s.
 \label{eq:V79-plus-amplitude}
\end{equation}
For two smooth plus traces \(s_1,s_2\), one has the exact cap Green identity
\begin{equation}
 \boxed{
 \begin{split}
 \langle H(s_1),H(s_2)\rangle_{\Sigma_t}
 &=\langle a_+(s_1),a_+(s_2)\rangle_{\Sigma_t}\\
 &\quad+
 \langle s_1,BD_ns_2\rangle_J
 +\langle D_ns_1,Bs_2\rangle_J .
 \end{split}}
 \label{eq:V79-plus-cap-Green-identity}
\end{equation}
\end{lemma}

\begin{proof}
The first row is the V63 boundary-stable TT graph.  Tangential
transversality gives the longitudinal component in the second row, and
spatial tracelessness gives the transverse component.  At a normal Fourier
frequency \(\rho\), \(A_+\) has symbol
\(1+\rho^2/q^2\).  The V63 trace formula
\(s=q^2a_+/\bigl(\sqrt2(q^2+\rho^2)\bigr)\) therefore proves
\eqref{eq:V79-plus-amplitude}.

Write primes for \(D_n\).  Direct contraction of the symmetric tensor gives
\begin{align*}
 \langle H(s_1),H(s_2)\rangle
 =\int_{x\geq0}\bigl(&2s_1s_2-s_1Bs_2''-s_1''Bs_2
       +2s_1'Bs_2'\\
 &\hspace{6em}+2s_1''B^2s_2''\bigr).
\end{align*}
On the other hand,
\begin{equation*}
 \langle a_+(s_1),a_+(s_2)\rangle
 =2\int_{x\geq0}(s_1-Bs_1'')(s_2-Bs_2'').
\end{equation*}
The first integrand minus the second is exactly
\begin{equation}
 D_n\bigl(s_1Bs_2'+s_1'Bs_2\bigr).
 \label{eq:V79-total-normal-derivative}
\end{equation}
Integration over the collar, using compact support at \(x=\infty\) and
\eqref{eq:V79-collar-orientation}, proves
\eqref{eq:V79-plus-cap-Green-identity}.  No normal-frequency branch choice
or formal inversion of a trace map is used.
\end{proof}

\begin{remark}[Why this term was invisible in the joint-free quotient]
\label{rem:V79-quotient-versus-local-trace}
The V53/V76 norm-one quotient compares \(H\) and \(a_+\) after the normal
Fourier transform and hence sees the bulk equality.  Formula
\eqref{eq:V79-plus-cap-Green-identity} is the transgression left when the
same equality is written on a cap with boundary in the local side variable
\(s\).  Thus it does not contradict
Theorem~\ref{thm:V76-unitary-radiative-quotient}.  It supplies the joint term
which that joint-free theorem was not intended to contain.
\end{remark}

\subsection{Exact cancellation with the V74 corner phase}
\label{subsec:V79-exact-cancellation}

Apply \eqref{eq:V79-plus-cap-Green-identity} with
\(s_1=D_0s\) and \(s_2=\delta s\).  The difference between the Einstein
plus cap potential and the V59 plus cap potential is the joint one-form
\begin{equation}
 \vartheta_{\rm cap}^{\rm edge}
 =\frac{\sigma_J}{4\kappa}
 \left(
  \langle D_0s,BD_n\delta s\rangle_J
  +\langle D_0D_ns,B\delta s\rangle_J
 \right).
 \label{eq:V79-cap-edge-potential}
\end{equation}

\begin{theorem}[The forced V78 cap equation holds]
\label{thm:V79-forced-equation-proved}
On the smooth nonzero-joint-frequency physical graph,
\begin{equation}
 \boxed{
 \omega_{\rm cap}^{\rm edge}
 =-\omega_{\rm cor}^{\rm phys}.}
 \label{eq:V79-cap-corner-two-form-cancellation}
\end{equation}
Since Lemma~\ref{lem:V79-V59-cap-no-edge} gives no additional continuum
joint curvature,
\begin{equation}
 \boxed{
 \omega_{{\rm cap+cont},J}^{\rm edge}
 =-\omega_{\rm cor}^{\rm phys},}
 \label{eq:V79-V78-forced-equation-solved}
\end{equation}
which is precisely the necessary and sufficient equation
\eqref{eq:V78-forced-cap-cancellation}.
\end{theorem}

\begin{proof}
Antisymmetrizing \eqref{eq:V79-cap-edge-potential} on two solutions gives
\begin{align}
 \omega_{\rm cap}^{\rm edge}(s_1,s_2)
 =\frac{\sigma_J}{4\kappa}\bigl[&
 \langle D_0s_1,BD_ns_2\rangle_J
 -\langle D_0s_2,BD_ns_1\rangle_J
 \notag\\
 &+\langle D_0D_ns_1,Bs_2\rangle_J
 -\langle D_0D_ns_2,Bs_1\rangle_J\bigr].
 \label{eq:V79-cap-edge-two-form}
\end{align}
Self-adjointness of \(B\) rewrites its first line as
\begin{equation*}
 \langle D_ns_2,BD_0s_1\rangle_J
 -\langle D_ns_1,BD_0s_2\rangle_J.
\end{equation*}
Comparison with the V78 formula
\eqref{eq:V78-pulled-corner-two-form} shows that both its first pair and its
second pair have the opposite sign.  This proves
\eqref{eq:V79-cap-corner-two-form-cancellation}.  The continuum assertion is
Lemma~\ref{lem:V79-V59-cap-no-edge}.
\end{proof}

The cancellation is even sharper at the level of potentials.

\begin{corollary}[Action-level generating functional]
\label{cor:V79-action-generator}
Define
\begin{equation}
 F_{79,J}(s)
 :=\frac{\sigma_J}{4\kappa}
   \langle D_0s,BD_ns\rangle_J.
 \label{eq:V79-generating-functional}
\end{equation}
Then
\begin{equation}
 \boxed{
 \vartheta_{\rm cap}^{\rm edge}
 +\vartheta_{\rm cor}^{\rm phys}
 =\delta F_{79,J}.}
 \label{eq:V79-potential-exactness}
\end{equation}
Thus the cap transgression and the full V74 corner dressing change the V59
potential only by a canonical generating function.
\end{corollary}

\begin{proof}
Insert the V78 corner potential
\begin{equation*}
 \vartheta_{\rm cor}^{\rm phys}
 =\frac{\sigma_J}{4\kappa}
 \left(
  \langle D_ns,B\delta D_0s\rangle_J
  -\langle D_0D_ns,B\delta s\rangle_J
 \right)
\end{equation*}
into \eqref{eq:V79-cap-edge-potential}.  The two terms containing
\(D_0D_ns\) cancel.  The remaining terms are
\begin{equation*}
 \frac{\sigma_J}{4\kappa}
 \left(
  \langle D_0s,BD_n\delta s\rangle_J
  +\langle D_ns,B D_0\delta s\rangle_J
 \right),
\end{equation*}
which is \(\delta F_{79,J}\) by self-adjointness of \(B\).
\end{proof}

\begin{remark}[Consistency with the V78 no-counterterm result]
\label{rem:V79-no-contradiction-counterterm}
Corollary~\ref{cor:V78-no-scalar-repair} remains valid: a scalar supported
at the joint cannot change the nonzero corner two-form.  Here the opposite
two-form first comes from the Einstein cap Green transgression
\eqref{eq:V79-cap-edge-two-form}.  Only after the two curvatures cancel does
their sum of potentials become the exact variation
\eqref{eq:V79-potential-exactness}.
\end{remark}

\subsection{Compact isotropy and composition}
\label{subsec:V79-compact-isotropy}

\begin{theorem}[Nonzero-frequency compact physical relation is isotropic]
\label{thm:V79-compact-isotropy}
Consider the flat linear compact slab on the smooth nonzero-joint-frequency
core.  Use the V76 reduced radiative quotient for the two physical
polarizations, the V59 positive continuum relation, the V63 physical side
graph, and the full V74 zero-level corner dressing.  The pullback of the
sum of the side, oriented cap, continuum endpoint, and corner symplectic
forms is zero.  With the image/completeness statement already proved in
Theorem~\ref{thm:V76-pulled-interface-Lagrangian}, the resulting reduced
compact relation is Lagrangian on this core.
\end{theorem}

\begin{proof}
The metric side form equals the V59 radiative side form by
\eqref{eq:V78-side-form-equality}.  The cross polarization has
\(s=0\) and hence contributes neither
\eqref{eq:V78-pulled-corner-two-form} nor
\eqref{eq:V79-cap-edge-two-form}; its bulk cap form is the V53/V76
norm-one quotient form.  In the plus block, Theorem
\ref{thm:V79-forced-equation-proved} cancels the only remaining joint
curvature.  Lemma~\ref{lem:V79-V59-cap-no-edge} identifies the continuum
endpoint form without another joint term.  Therefore the total pullback
vanishes.  The V76 frequency exact sequence and image theorem supply the
maximality statement on the same nonzero-frequency quotient core.
\end{proof}

\begin{corollary}[Oppositely oriented caps glue without a joint anomaly]
\label{cor:V79-oriented-gluing}
Let two flat linear slabs be glued across a common temporal cap, with
matching physical and continuum data.  Their V79 generating functionals
and their cap/corner transgression forms cancel pairwise on the common
joint.
\end{corollary}

\begin{proof}
The terminal cap of the first slab has \(\sigma_J=+1\), while the initial
cap of the second has \(\sigma_J=-1\).  Every expression
\eqref{eq:V79-cap-edge-potential},
\eqref{eq:V79-cap-edge-two-form}, and
\eqref{eq:V79-generating-functional} is linear in \(\sigma_J\).  Matching
the shared data therefore makes their two contributions equal and
opposite.  The ordinary bulk cap potentials cancel by the same oriented
Stokes rule.
\end{proof}

\begin{firewall}[Exact scope of the V79 compact theorem]
\label{fire:V79-scope}
Theorem~\ref{thm:V79-compact-isotropy} is a flat, linear,
nonzero-joint-frequency statement on the already constructed reduced
physical quotient.  It proves the cap equation left open in V78 and repairs
the compact form test there.  It does not construct a local half-space
metric section: the cross and bulk pieces continue to use the operator
quotient of V76, while the plus joint transgression is computed directly in
the local trace chart.  It also does not cover \(Q=0\), nonsmooth joint
traces, a nonlinear BV--BFV master action, or interacting Standard Model
stress.
\end{firewall}

\subsection{The V78 witness as a sign and normalization check}
\label{subsec:V79-witness-check}

\begin{corollary}[Exact cancellation on the \(3\)-\(4\)-\(5\) witness]
\label{cor:V79-345-cancellation}
For the two V78 plus waves with
\((q,\rho,\omega)=(3,4,5)\),
\begin{equation}
 \omega_{\rm cap}^{\rm edge}(s_1,s_2)
 =-\sigma_J\frac{5\pi L_z}{27\kappa},
 \qquad
 \omega_{\rm cor}^{\rm phys}(s_1,s_2)
 = \sigma_J\frac{5\pi L_z}{27\kappa}.
 \label{eq:V79-345-cancellation}
\end{equation}
\end{corollary}

\begin{proof}
The second value is \eqref{eq:V78-345-corner-value}.  The first follows
from \eqref{eq:V79-cap-corner-two-form-cancellation}.  This checks the cap
orientation, the factor \(1/(4\kappa)\), and the sign of the normal Green
term simultaneously.
\end{proof}

\subsection{V79 result ledger and next acquisition}
\label{subsec:V79-ledger}

\paragraph{New exact conclusions.}
\begin{enumerate}[label=\textup{(\arabic*)},leftmargin=2.8em]
\item The oriented Einstein--GHY cap potential on the harmonic TT shell is
      \eqref{eq:V79-Einstein-cap-theta}.
\item The V59 radiative and continuum cap potential is
      \eqref{eq:V79-V59-cap-theta}; the trace condition creates no separate
      continuum joint curvature.
\item The plus TT reconstruction from \(s=\varsigma\) is
      \eqref{eq:V79-plus-H-normal}--%
      \eqref{eq:V79-plus-H-tangential}, and its radiative amplitude is
      \(a_+=\sqrt2(1-Q^{-2}D_n^2)s\).
\item Their cap norms differ by the exact joint Green term
      \eqref{eq:V79-plus-cap-Green-identity}.
\item The curvature of that cap transgression is the negative of the V78
      corner curvature, proving the forced equation
      \eqref{eq:V79-V78-forced-equation-solved}.
\item At the potential level, cap plus corner equals the exact variation of
      \(F_{79,J}\) in \eqref{eq:V79-generating-functional}.
\item The compact flat nonzero-frequency reduced physical relation is
      Lagrangian, and its joint contributions compose under opposite cap
      orientations.
\end{enumerate}

\paragraph{Next forced step.}
V80 is an audit version.  Before extending the proof, it must inspect the
newest Yang--Mills and Navier--Stokes notes in the shared folder and record
any useful common domain, boundary-symbol, or cofinal-closure mechanisms.
It must then analyze the zero joint mode \(Q=0\), where the inverse
\(Q^{-2}\) used in V74, V78, and V79 is unavailable.  The first task is to
separate spatially homogeneous gauge charges from the normally propagating
TT polarizations and decide whether the zero-mode cap/corner phase is
trivial, finite dimensional, or requires a new global boundary variable.

\begin{assessment}[V79 continuum status]
The full Standard--Model--Einstein continuum remains unproved.  V79 closes
the flat linear nonzero-frequency compact symplectic gluing problem left by
V78: the oriented Einstein cap transgression cancels the V74 corner
two-form exactly, and the V59 continuum endpoint creates no additional
joint curvature.  The zero joint mode, finite-energy joint trace, nonlinear
curved BV--BFV gluing, Brown--York/Bianchi stress identity, common
renormalized Standard Model stress domain, and cofinal interacting continuum
remain open.
\end{assessment}

\section*{V79 exact checks and append-only record}
\addcontentsline{toc}{section}{V79 exact checks and append-only record}

The central algebraic check is the polynomial identity
\begin{align}
 &2s_1s_2-s_1Bs_2''-s_1''Bs_2+2s_1'Bs_2'
   +2s_1''B^2s_2''
 \notag\\
 &\quad
 -2(s_1-Bs_1'')(s_2-Bs_2'')
 =D_n(s_1Bs_2'+s_1'Bs_2),
 \label{eq:V79-algebra-check-card}
\end{align}
which proves the cap Green term before any use of the wave equation.  The
potential check is
\begin{equation}
 \vartheta_{\rm cap}^{\rm edge}
 +\vartheta_{\rm cor}^{\rm phys}
 =\frac{\sigma_J}{4\kappa}\delta
  \langle D_0s,Q^{-2}D_ns\rangle_J.
 \label{eq:V79-potential-check-card}
\end{equation}
Taking \(\delta\) gives
\eqref{eq:V79-cap-corner-two-form-cancellation} because \(\delta^2=0\).

The V78 predecessor used for this append has SHA--256
\begin{center}
\texttt{0580b36579e4cd3a7221f57dd357f64de07af97b352f91b8d415ddf1edda8dd5}.
\end{center}
The V79 candidate retains its bytes up to the sole live final document
terminator, appends this material, and restores exactly one active
\verb|\end{document}|.  Validation includes balanced environments and
braces, unique labels, resolved internal references, valid inline and display
math delimiters, exact predecessor-prefix equality, and a final inventory
which removes only the superseded SM V78 file.

% =============================================================================
% END APPENDED VERSION V79 MATERIAL
% =============================================================================
\clearpage
% =============================================================================
% BEGIN APPENDED VERSION V80 MATERIAL
% =============================================================================
\section{V80: companion audit and the exact joint-zero-mode split}
\label{sec:V80-zero-mode}

V79 solved the nonzero-joint-frequency cap equation.  The inverse
\(Q^{-2}\) used there cannot be continued to the constant joint mode, and
the vanishing of the trace variable \(\varsigma\) at normal incidence does
not imply that the two radiative polarizations vanish.  This audit version
therefore separates the spectral projection onto the constant joint mode
before applying any inverse.  It then computes the metric charge, corner
cocycle, side form, and cap form directly on that sector.

The result is simple but structural.  The constant joint mode has two
normal-incidence TT amplitudes.  Its V72 joint cocycle and its physical
improved charge are both zero, so its minimal anomaly-cancelling corner
carrier is the zero space.  The TT polarization matrices are constant and
orthonormal, making the Einstein side and cap forms equal to the V59 forms
without a Green transgression.  Thus the exact joint zero mode closes by a
direct chart, while the infrared neighborhood of zero remains a separate
finite-energy problem.

\subsection{Scheduled NS/Yang--Mills audit}
\label{subsec:V80-audit}

The following two snapshots were read before the V80 calculation.
\begin{center}
\begin{tabular}{p{0.42\textwidth}p{0.13\textwidth}p{0.34\textwidth}}
\toprule
companion file & bytes & SHA--256\\
\midrule
\texttt{Renewal\_NS\_Stability\_Research\_Notes\_V31.tex}
& \(887537\)
& \texttt{f1121b62238ad5491bb7cf03635ea9934942edf573ed8faaa9ee79e1d38a6696}\\
\texttt{Renewal\_Yang\_Mills\_Mass\_Gap\_Research\_Notes\_V71.tex}
& \(1867690\)
& \texttt{7ccb5a4af32affb3967cacb73f78bef038ec125bf2a025b7010c4c34bbd37e50}\\
\bottomrule
\end{tabular}
\end{center}

The NS V31 snapshot proves that flat dyadic probability marks recover the
unweighted critical coarea only after multiplication by the exact radius
first moment.  It also gives a balanced-helicity plane wave on which the
signed hinge diagnostics vanish while the positive total coarea remains
nonzero.  The relevant discipline for the present note is that a vanishing
reduced diagnostic does not erase a physical sector, and that a missing
spectral weight cannot be hidden in a normalization.

The Yang--Mills V71 snapshot computes an exact Racah output marginal before
estimating its conditional symbol.  The isolated singlet and fixed output
windows receive an inverse-dimension debit, whereas a growing thermal
mesoscopic window retains order-one cumulative mass.  The relevant
discipline is to isolate an exceptional spectral atom before applying a
coarse inverse or supremum; normalization of the complement does not prove
control of the atom.

\begin{firewall}[Companion results are not imported across equation types]
\label{fire:V80-audit-types}
No Navier--Stokes coarea estimate and no Yang--Mills representation bound is
used as an Einstein estimate.  The audit changes only the acquisition
order: first apply the exact spectral projection, then compute the
exceptional sector in its own variables, and only afterwards use an inverse
on the complementary sector.
\end{firewall}

\subsection{The constant and mean-zero joint sectors}
\label{subsec:V80-spectral-split}

For a compact flat joint \(J\), let
\begin{equation}
 P_0f:=\frac1{|J|}\int_Jf,
 \qquad
 P_\perp:=I-P_0.
 \label{eq:V80-joint-projections}
\end{equation}
Then \(Q=(-\Delta_J)^{1/2}\) is zero on \(P_0L^2(J)\) and is invertible on
the smooth \(P_\perp\) core.  On a noncompact joint, the single Fourier
point \(q=0\) is absent from the \(L^2\) direct integral; the discussion
below classifies the corresponding generalized fibre and the compact-joint
mode.

\begin{lemma}[The V72 cocycle vanishes on the constant joint mode]
\label{lem:V80-zero-cocycle}
For every residual wave jet with constant joint dependence,
\begin{equation}
 P_0R^0=0,
 \qquad
 P_0R^A=0,
 \qquad
 \mathcal K_{72}|_{P_0\mathfrak g}=0.
 \label{eq:V80-zero-R-and-cocycle}
\end{equation}
Consequently the minimal symplectic carrier required to realize the
opposite cocycle on this sector is the zero vector space.
\end{lemma}

\begin{proof}
The V77 normal-trace formulas are
\begin{equation*}
 R^0=\sigma_JQ^2f,
 \qquad
 R^A=\sigma_JD_0D^Af.
\end{equation*}
Both vanish after \(P_0\), because \(Q^2P_0=0\) and \(D^AP_0=0\).
The V72 cocycle is the alternating pairing of residual trace data with
these normal traces, so it vanishes.  A zero cocycle is already represented
Hamiltonianly without adding a symplectic vector space; its minimal opposite
carrier therefore has dimension zero.
\end{proof}

\begin{definition}[The two zero-mode corner presentations]
\label{def:V80-zero-corner-presentations}
The cocycle-minimal presentation on \(P_0\) has no corner symplectic
carrier, because Lemma~\ref{lem:V80-zero-cocycle} gives rank zero.  The full
V74 charge-complete presentation may instead be retained.  On its constant
mode it specializes to
\begin{align}
 \mathcal V_\epsilon^{\rm cor}X_b&=P_0\epsilon_b,
 &
 \mathcal V_\epsilon^{\rm cor}\Pi_{\rm cor}^{\,b}&=0,
 \label{eq:V80-zero-corner-action}\\
 \mu_{\epsilon,0}^{\rm cor}
 &=-\frac1\kappa\int_J
   (P_0\epsilon_b)\Pi_{\rm cor}^{\,b}.
 \label{eq:V80-zero-corner-moment}
\end{align}
This is a cotangent charge resolution with zero cocycle.  V80 retains this
full presentation in the off-shell product phase and uses the
zero-dimensional statement only after physical reduction.
\end{definition}

\begin{proposition}[The full zero-mode corner pair reduces without a
physical spectator]
\label{prop:V80-zero-corner-reduces}
The constant residual trace map
\(\epsilon\mapsto P_0\epsilon_b\) is onto the three joint components.
On the normal-incidence physical TT graph, the combined zero-level equation
sets
\begin{equation}
 P_0\Pi_{\rm cor}^{\,b}=0,
 \label{eq:V80-zero-corner-momentum}
\end{equation}
and quotient by \eqref{eq:V80-zero-corner-action} removes
\(P_0X_b\).  Hence no independent constant corner class or corner two-form
survives physical reduction.  The result agrees with the zero-dimensional
cocycle-minimal carrier on this graph.
\end{proposition}

\begin{proof}
At a fixed temporal joint jet, a residual wave has freely prescribable
boundary value \(\epsilon_b\); its wave equation determines second normal
or temporal derivatives and does not constrain that value.  Thus the trace
map is onto.  Lemma~\ref{lem:V80-zero-charge} proves that the metric moment
vanishes on the physical TT graph.  The combined moment
\(\mu_\epsilon+\mu_{\epsilon,0}^{\rm cor}=0\) for every constant trace then
gives \eqref{eq:V80-zero-corner-momentum}.  The residual translations are
transitive in \(P_0X_b\), so the quotient admits the gauge
\(P_0X_b=0\).  The corner potential
\(\kappa^{-1}\int\Pi_{\rm cor}^{\,b}\delta X_b\) and its two-form therefore
pull back to zero.
\end{proof}

\begin{firewall}[The mean-zero inverse has no infrared estimate]
\label{fire:V80-no-Q-limit}
The spectral split makes every occurrence of \(Q^{-2}\) in V78--V79 an
operator on \(P_\perp\); the full zero-mode resolution
\eqref{eq:V80-zero-corner-action} uses no inverse.  This does not give a
bound uniform as nonzero \(|q|\downarrow0\).  On a fixed compact joint the
first positive eigenvalue supplies a finite constant; a family of growing
joints or the noncompact finite-energy closure requires a separate weighted
infrared domain.
\end{firewall}

\subsection{Direct normal-incidence TT chart}
\label{subsec:V80-normal-chart}

Choose an oriented orthonormal basis \(e_1,e_2\) tangent to \(J\) and put
\begin{align}
 E_+^0&:=\frac{e_1\otimes e_1-e_2\otimes e_2}{\sqrt2},
 &
 E_\times^0&:=\frac{e_1\otimes e_2+e_2\otimes e_1}{\sqrt2}.
 \label{eq:V80-normal-TT-basis}
\end{align}
They obey
\begin{equation}
 E_\alpha^0:E_\beta^0=\delta_{\alpha\beta},
 \qquad
 \operatorname{tr}_J E_\alpha^0=0.
 \label{eq:V80-normal-TT-Gram}
\end{equation}

\begin{lemma}[Complete constant-joint physical metric]
\label{lem:V80-normal-TT-field}
Every constant-joint, nonzero-normal-frequency physical class has a TT
representative
\begin{equation}
 H_{AB}=a_+E_{+,AB}^0+a_\times E_{\times,AB}^0,
 \qquad
 H_{n\mu}=H_{0\mu}=0,
 \label{eq:V80-normal-TT-field}
\end{equation}
where \(a_+,a_\times\) are scalar solutions of the one-normal-dimensional
wave equation.  In particular,
\begin{equation}
 \varsigma:=H_A{}^A=0
 \label{eq:V80-zero-trace-nonzero-radiation}
\end{equation}
for both independent polarizations.
\end{lemma}

\begin{proof}
At normal incidence the spatial propagation covector is parallel to the
side normal.  Spatial transversality removes every component with a normal
index, and tracelessness leaves the two-dimensional trace-free symmetric
tangential block.  The two tensors
\eqref{eq:V80-normal-TT-basis} are an orthonormal basis of that block.
They are exactly the \(\theta=0\) endpoint of the V62/V63 angular TT basis,
and the V53 quotient has two classes there.  Their traces vanish by
\eqref{eq:V80-normal-TT-Gram}, proving
\eqref{eq:V80-zero-trace-nonzero-radiation}.
\end{proof}

\begin{remark}[The NS audit nullspace lesson is realized exactly]
\label{rem:V80-null-diagnostic}
Equation~\eqref{eq:V80-zero-trace-nonzero-radiation} is the gravitational
analogue of the audit warning, not an identification of the two theories.
The scalar diagnostic \(\varsigma\) has a two-dimensional physical
nullspace at normal incidence.  Therefore the V79 plus-trace chart must be
replaced there by \((a_+,a_\times)\); taking the limit of its inverse formula
would incorrectly erase both modes.
\end{remark}

\subsection{Charge and Green forms on the zero joint mode}
\label{subsec:V80-zero-charge-forms}

\begin{lemma}[Brown--York rows and the improved charge vanish where needed]
\label{lem:V80-zero-charge}
For the TT field \eqref{eq:V80-normal-TT-field}, the side extrinsic
curvature and Brown--York momentum obey
\begin{align}
 K_{AB}^{(1)}&=\frac12D_nH_{AB},
 &K^{(1)}&=0,
 \label{eq:V80-normal-K}\\
 \Pi_{(1)}^{00}&=0,
 &\Pi_{(1)}^{0A}&=0.
 \label{eq:V80-normal-Pi-rows}
\end{align}
For every constant-joint residual parameter, the V77/V78 improved joint
charge restricts to
\begin{equation}
 \mu_\epsilon(H)=0.
 \label{eq:V80-zero-physical-charge}
\end{equation}
Thus every normal-incidence physical TT datum lies on the undressed metric
zero level.
\end{lemma}

\begin{proof}
In the flat side formula
\begin{equation*}
 K_{ab}^{(1)}
 =\frac12(D_nh_{ab}-D_ah_{nb}-D_bh_{na}),
\end{equation*}
all mixed normal components vanish and all joint derivatives vanish.  This
gives \eqref{eq:V80-normal-K}; its trace is zero because \(H_A{}^A=0\).
The time rows in \eqref{eq:V80-normal-Pi-rows} follow immediately.

The improved charge on the V63 graph is
\begin{equation*}
 \mu_\epsilon(H)
 =\frac{\sigma_J}{2\kappa}\int_J
 \left[
  fD_0\varsigma-(D_0f)\varsigma
  -2\epsilon_b\Pi_{(1)}^{0b}
 \right].
\end{equation*}
Equations \eqref{eq:V80-zero-trace-nonzero-radiation} and
\eqref{eq:V80-normal-Pi-rows} make every term zero.
\end{proof}

\begin{theorem}[Exact side and cap form matching at normal incidence]
\label{thm:V80-zero-form-matching}
On the field \eqref{eq:V80-normal-TT-field}, the Einstein side potential
restricts to
\begin{equation}
 \Theta_{E,\Gamma}^{q=0}
 =\frac1{4\kappa}\int_\Gamma
 \sum_{\alpha\in\{+,\times\}}
 D_na_\alpha\,\delta a_\alpha,
 \label{eq:V80-zero-side-potential}
\end{equation}
and the oriented cap potential restricts to
\begin{equation}
 \Theta_{E,\Sigma_t}^{q=0}
 =\frac{\sigma_J}{4\kappa}\int_{\Omega_x}
 \sum_{\alpha\in\{+,\times\}}
 D_0a_\alpha\,\delta a_\alpha\,\dd x\,\dd^2y.
 \label{eq:V80-zero-cap-potential}
\end{equation}
These are exactly the action-normalized V59 radiative side and cap
potentials.  There is no codimension-two transgression:
\begin{equation}
 \vartheta_{\rm cap}^{\rm edge}|_{q=0}=0,
 \qquad
 \omega_{\rm cor}^{\rm phys}|_{q=0}=0.
 \label{eq:V80-zero-no-edge-forms}
\end{equation}
\end{theorem}

\begin{proof}
Insert \eqref{eq:V80-normal-K} in the V36/GHY one-form.  Since
\(\Pi^{AB}_{(1)}=D_nH^{AB}/2\), orthonormality
\eqref{eq:V80-normal-TT-Gram} gives
\eqref{eq:V80-zero-side-potential}.  Lemma
\ref{lem:V79-Einstein-cap-potential} and the same Gram identity give
\eqref{eq:V80-zero-cap-potential}.  No normal integration by parts is
needed: the polarization matrices are constant and the scalar amplitudes
are already the norm-one V59 variables.  This proves the first equality in
\eqref{eq:V80-zero-no-edge-forms}.  For the second, use
Proposition~\ref{prop:V80-zero-corner-reduces}: on the physical zero level
one may take \(P_0X=P_0\Pi_{\rm cor}=0\), so the corner potential and
two-form vanish.  No undefined \(Q^{-2}\) formula is evaluated.
\end{proof}

\begin{theorem}[The compact \(q=0\), nonzero-normal-spectrum relation is
Lagrangian]
\label{thm:V80-zero-mode-Lagrangian}
On the smooth constant-joint sector with nonzero normal spectral parameter,
the inverse image of the V59 radiative--continuum relation under the V63
normal-incidence metric quotient is Lagrangian in the reduced compact
Einstein--continuum phase.  No independent corner canonical class survives
on this sector.
\end{theorem}

\begin{proof}
The V63 exact conormal lift includes normal incidence and identifies both
TT amplitudes and their side conormals with V59.  Theorem
\ref{thm:V80-zero-form-matching} identifies the oriented cap forms, while
Lemma~\ref{lem:V79-V59-cap-no-edge} supplies the continuum cap form and the
common trace.  Lemmas \ref{lem:V80-zero-cocycle} and
\ref{lem:V80-zero-charge}, together with
Proposition~\ref{prop:V80-zero-corner-reduces}, show that the full corner
resolution leaves no physical class.  Hence the complete
pullback form vanishes.  Maximality follows from the two-dimensional V53
TT quotient and the V59 self-adjoint interface relation: the metric lift
adds only residual gauge directions, which are removed by reduction.
\end{proof}

\subsection{The fully homogeneous spectral point}
\label{subsec:V80-total-zero}

The case \(q=0\) with nonzero normal spectrum is physical normal-incidence
radiation.  It must be distinguished from the total spatial frequency
\((q,\rho)=(0,0)\).

\begin{proposition}[No dynamical total-zero class in the V59 half-space
energy phase]
\label{prop:V80-total-zero-excluded}
On the V59 metric half-space and auxiliary continuum half-line, a spatially
constant radiative configuration is not a nonzero element of the stated
\(L^2\)/homogeneous energy phase.  In the operator kernel, trace continuity
and finite continuum energy force the total-zero pair to vanish.
\end{proposition}

\begin{proof}
A nonzero constant radiative field on the infinite metric half-space is not
in \(L^2\), while a spatially constant nonzero velocity has infinite kinetic
energy.  In the homogeneous configuration completion, zero-gradient
constants represent the zero seminorm class rather than an additional
radiative degree of freedom.  For an operator-kernel pair, the V59 closed
form has zero value only when \(a\) and \(b\) have zero spatial gradients.
The continuum half-line admits no nonzero constant \(L^2\) function, so
\(b=0\).  Trace continuity gives \(a|_J=0\), and constancy gives \(a=0\).
\end{proof}

\begin{firewall}[Compact global moduli are a separate sector]
\label{fire:V80-global-moduli}
If every spatial direction is compactified and the auxiliary continuum is
removed or its domain is changed, constant flat-metric moduli and global
Killing parameters may survive.  Proposition
\ref{prop:V80-total-zero-excluded} does not classify that different global
problem.  The V59 half-space continuum used in this programme has the
domain stated above, on which the total-zero dynamical kernel is absent.
\end{firewall}

\subsection{V80 result ledger and next acquisition}
\label{subsec:V80-ledger}

\paragraph{New exact conclusions.}
\begin{enumerate}[label=\textup{(\arabic*)},leftmargin=2.8em]
\item The scheduled snapshots of NS V31 and Yang--Mills V71 were inspected,
      with hashes recorded, without importing estimates across equation
      types.
\item The joint spectral splitting \(P_0\oplus P_\perp\) is made before
      any use of \(Q^{-2}\).
\item On \(P_0\), the V72 cocycle vanishes and the cocycle-minimal opposite
      symplectic carrier has dimension zero.
\item The full V74 constant mode remains a charge-complete cotangent
      resolution: its momentum is zero on the physical graph and its
      coordinate is removed by residual translation, so it leaves no
      independent reduced class.
\item The constant-joint physical sector still has two independent TT
      amplitudes even though \(\varsigma=0\).
\item Its Brown--York time rows and improved residual charge vanish, so the
      two TT modes lie on the undressed metric zero level.
\item Its side and cap Green potentials equal the V59 potentials directly,
      with neither the V78 corner form nor the V79 cap transgression.
\item The compact \(q=0\), nonzero-normal-spectrum physical relation is
      Lagrangian.  The total spatial zero point is absent from the V59
      half-space energy/operator kernel.
\end{enumerate}

\paragraph{Next forced step.}
V81 must address the neighborhood of \(q=0\), not recompute the exact zero
fibre.  It should identify the sharp Sobolev or weighted Fourier domain on
which the V78/V79 expressions containing \(Q^{-2}\) are finite, test that
domain against the V59 homogeneous energy space, and decide whether the
joint cap generator is closable.  If the weight is stronger than V59
energy, the attack must move upstream to an unintegrated local corner
potential or add an explicit infrared boundary variable; it must not hide
the loss behind the compact-joint spectral gap.

\begin{assessment}[V80 continuum status]
The full Standard--Model--Einstein continuum remains unproved.  V80 closes
the exact joint-zero-mode classification for the flat linear compact
physical relation and shows why it cannot be obtained as a
\(Q^{-2}\)-limit.  The near-zero finite-energy joint domain, nonlinear
curved BV--BFV gluing, Brown--York/Bianchi stress identity, common
renormalized Standard Model stress domain, and cofinal interacting continuum
remain open.
\end{assessment}

\section*{V80 exact checks and append-only record}
\addcontentsline{toc}{section}{V80 exact checks and append-only record}

The zero-mode check card is
\begin{equation}
 Q^2P_0=0,\qquad D_AP_0=0,\qquad
 \operatorname{tr}_JE_+^0=\operatorname{tr}_JE_\times^0=0,\qquad
 E_\alpha^0:E_\beta^0=\delta_{\alpha\beta}.
 \label{eq:V80-zero-check-card}
\end{equation}
It proves separately that the cocycle trace, scalar dressing coordinate,
and cap Green transgression vanish, while the physical rank remains two.

The V79 predecessor used for this append has SHA--256
\begin{center}
\texttt{03951eaebd4da33c451c89de0f6a06eebf043c36f3d6cddcefd36683bdd8b26e}.
\end{center}
The V80 candidate retains its bytes up to the sole live final document
terminator, appends this material, and restores exactly one active
\verb|\end{document}|.  Validation includes balanced environments and
braces, unique labels, resolved new internal references, valid inline and
display math delimiters, exact predecessor-prefix equality, and a final
inventory which removes only the superseded SM V79 file.

% =============================================================================
% END APPENDED VERSION V80 MATERIAL
% =============================================================================
\clearpage
% =============================================================================
% BEGIN APPENDED VERSION V81 MATERIAL
% =============================================================================
\section{V81: sharp infrared domain of the cap generator and the upstream
combined-form closure}
\label{sec:V81-infrared-domain}

V80 treated the exact joint zero fibre without an inverse.  The remaining
question is whether the V78/V79 formulas are continuous as nonzero joint
frequencies approach zero in the V59 homogeneous energy norm.  This note
answers that question sharply on the physical plus range.

The factor \(Q^{-2}\) is partly cancelled by the TT trace multiplier:
\(\varsigma\) itself has a factor \(|q|^2\).  After this cancellation, the
V79 generating functional still has fibre norm proportional to
\(|q|^{-1}\).  Hence it has no bounded extension to the full unweighted V59
energy phase.  A normalized smooth family localized at
\(|q|\asymp\varepsilon\) makes it grow like \(\varepsilon^{-1}\).
The sharp symmetric repair is one half infrared derivative on each Cauchy
factor.  This weighted domain is dense but proper.

The obstruction concerns separated potentials.  The cap and corner
two-forms cancel exactly on the smooth punctured core, so their \emph{sum}
has the unique continuous extension zero on the reduced physical energy
graph.  The next construction must therefore return upstream to the local
unseparated Einstein--GHY--corner current and define its weak reduction
before solving the dressing equation.  Introducing a new infrared boundary
field is not yet justified.

\subsection{Physical-range multipliers}
\label{subsec:V81-physical-multipliers}

Use unitary tangential and normal Fourier transforms, with covector
\(q\in\mathbb R^2\), \(r=|q|>0\), and normal variable
\(\rho\in\mathbb R\) on the doubled collar.  Thus restriction to the joint
is
\(f|_J=(2\pi)^{-1/2}\int_{\mathbb R}\widehat f(\rho)\,\dd\rho\).
The exact constants below refer to this doubled-collar normalization.  Set
\begin{equation}
 \lambda_r(\rho)^2:=r^2+\rho^2.
 \label{eq:V81-lambda}
\end{equation}

\begin{lemma}[Exact physical cancellation of the formal \(Q^{-2}\) factor]
\label{lem:V81-physical-multipliers}
For the plus amplitude \(a=a_+\), the V63/V79 relation gives
\begin{align}
 \widehat s(q,\rho)
 &=m_r(\rho)\widehat a(q,\rho),
 &
 m_r(\rho)&:=\frac{r^2}{\sqrt2\lambda_r(\rho)^2},
 \label{eq:V81-s-multiplier}\\
 r^{-2}\widehat{D_ns}(q,\rho)
 &=n_r(\rho)\widehat a(q,\rho),
 &
 n_r(\rho)&:=\frac{i\rho}{\sqrt2\lambda_r(\rho)^2}.
 \label{eq:V81-BDns-multiplier}
\end{align}
The same first multiplier sends the cap velocity
\(v=D_0a\) to \(D_0s\).  Thus no \(r^{-2}\) factor remains on the physical
range at fixed nonzero total spatial frequency.  The map nevertheless
cannot be reconstructed from \(s\) at \(r=0\), because V80 gives \(s=0\)
for both independent normal-incidence amplitudes.
\end{lemma}

\begin{proof}
Equation~\eqref{eq:V79-plus-amplitude} is
\begin{equation*}
 a=\sqrt2(1-Q^{-2}D_n^2)s.
\end{equation*}
At the Fourier point \((q,\rho)\), its multiplier is
\(\sqrt2(r^2+\rho^2)/r^2\).  Inversion for \(r>0\) gives
\eqref{eq:V81-s-multiplier}.  Multiplication by \(i\rho r^{-2}\) gives
\eqref{eq:V81-BDns-multiplier}.  The operators commute with \(D_0\).
The final statement is Theorem~\ref{thm:V80-zero-form-matching} and
Equation~\eqref{eq:V80-zero-trace-nonzero-radiation}.
\end{proof}

At one temporal cap, \((a,v)\) are independent Cauchy variables.  Inserting
Lemma~\ref{lem:V81-physical-multipliers} into the V79 generator defines the
bilinear form
\begin{equation}
 \mathfrak F_{81}(v,a)
 :=\frac{\sigma_J}{4\kappa}
 \int_{\mathbb R^2_q}
 \overline{T_v(q)}T_a(q)\,\dd^2q,
 \label{eq:V81-F-bilinear}
\end{equation}
where
\begin{align}
 T_v(q)&:=(2\pi)^{-1/2}\int_{\mathbb R}m_r(\rho)\widehat v(q,\rho)\,\dd\rho,
 \label{eq:V81-Tv}\\
 T_a(q)&:=(2\pi)^{-1/2}\int_{\mathbb R}n_r(\rho)\widehat a(q,\rho)\,\dd\rho.
 \label{eq:V81-Ta}
\end{align}
Taking the real part gives the real-field functional
\(F_{79,J}\); the complexified form is convenient for norm calculations.

\subsection{Exact fibre norm and the energy counterexample}
\label{subsec:V81-fibre-norm}

For fixed \(q\), define the Cauchy energy norms
\begin{equation}
 E_v(q):=\int_{\mathbb R}|\widehat v(q,\rho)|^2\,\dd\rho,
 \qquad
 E_a(q):=\int_{\mathbb R}
 \lambda_r(\rho)^2|\widehat a(q,\rho)|^2\,\dd\rho.
 \label{eq:V81-fibre-energies}
\end{equation}

\begin{theorem}[Exact \(r^{-1}\) fibre norm]
\label{thm:V81-exact-fibre-norm}
The two trace functionals have exact squared norms
\begin{align}
 \|T_v(q)\|^2_{E_v(q)^*}
 &=(2\pi)^{-1}\int_{\mathbb R}|m_r(\rho)|^2\,\dd\rho
 =\frac{r}{8},
 \label{eq:V81-Tv-norm}\\
 \|T_a(q)\|^2_{E_a(q)^*}
 &=(2\pi)^{-1}\int_{\mathbb R}
   \frac{|n_r(\rho)|^2}{\lambda_r(\rho)^2}\,\dd\rho
 =\frac1{32r^3}.
 \label{eq:V81-Ta-norm}
\end{align}
Consequently the exact norm of the fibre bilinear
\((v,a)\mapsto\overline{T_v(q)}T_a(q)\) is
\begin{equation}
 \boxed{
 \|T_v(q)\|\,\|T_a(q)\|
 =\frac1{16r}.}
 \label{eq:V81-fibre-product-norm}
\end{equation}
\end{theorem}

\begin{proof}
The elementary integrals
\begin{align*}
 \int_{\mathbb R}\frac{\dd\rho}{(r^2+\rho^2)^2}
 &=\frac{\pi}{2r^3},\\
 \int_{\mathbb R}\frac{\rho^2\,\dd\rho}
 {(r^2+\rho^2)^3}
 &=\frac{\pi}{8r^3}
\end{align*}
and the common factor \((2\pi)^{-1}\) from joint restriction give
\eqref{eq:V81-Tv-norm}--\eqref{eq:V81-Ta-norm} after inserting
\eqref{eq:V81-s-multiplier}--\eqref{eq:V81-BDns-multiplier}.
Cauchy--Schwarz is sharp: its extremizers are proportional to
\(m_r\) for the velocity and to
\(\overline{n_r}\lambda_r^{-2}\) for the position.  Multiplication of the
two exact norms proves \eqref{eq:V81-fibre-product-norm}.
\end{proof}

\begin{theorem}[The V79 generator is unbounded on V59 energy]
\label{thm:V81-energy-unbounded}
There is no constant \(C\) such that
\begin{equation}
 |\mathfrak F_{81}(v,a)|
 \le C
 \left(\int E_v(q)\,\dd^2q\right)^{1/2}
 \left(\int E_a(q)\,\dd^2q\right)^{1/2}
 \label{eq:V81-false-energy-bound}
\end{equation}
on the smooth punctured physical core.  More precisely, there are smooth
real Cauchy data \((a_\varepsilon,v_\varepsilon)\), supported in
\(\varepsilon<|q|<2\varepsilon\), with both displayed energy norms equal to
one and
\begin{equation}
 |\mathfrak F_{81}(v_\varepsilon,a_\varepsilon)|
 \ge\frac{c}{\kappa\varepsilon}
 \label{eq:V81-unbounded-packet}
\end{equation}
for an absolute \(c>0\).
\end{theorem}

\begin{proof}
Choose \(\chi_\varepsilon\in C_c^\infty(\mathbb R^2)\) supported in the
stated annulus with \(\|\chi_\varepsilon\|_2=1\).  At each
\(q\) in its support, take the normalized equality cases from the proof of
Theorem~\ref{thm:V81-exact-fibre-norm}:
\begin{align*}
 \widehat v_\varepsilon(q,\rho)
 &=\chi_\varepsilon(q)
   \frac{m_r(\rho)}{\|m_r\|_{L^2_\rho}},\\
 \widehat a_\varepsilon(q,\rho)
 &=\chi_\varepsilon(q)
   \frac{\overline{n_r(\rho)}}
        {\lambda_r(\rho)^2
         \left(\displaystyle\int_{\mathbb R}
          |n_r(\eta)|^2\lambda_r(\eta)^{-2}\dd\eta\right)^{1/2}}.
\end{align*}
Their total velocity and position energies are one.  Choose the harmless
common phase so that the integrand in
\eqref{eq:V81-F-bilinear} is nonnegative.  Then
\begin{equation*}
 |\mathfrak F_{81}|
 =\frac1{4\kappa}
  \int|\chi_\varepsilon(q)|^2\frac1{16|q|}\,\dd^2q
 \ge\frac1{128\kappa\varepsilon}.
\end{equation*}
Symmetrizing the data under
\((q,\rho)\mapsto(-q,-\rho)\) gives real fields and changes only an absolute
constant.  Cutting off the normal profiles and passing to the limit gives
Schwartz normal data with the same lower scaling.

The position extremizer is odd in the normal frequency, so its metric trace
at \(x=0\) is zero.  To place the packet in the full smooth V59
trace--conormal core, use an auxiliary continuum collar coordinate
\(\zeta>0\).  Choose a smooth \(b_\varepsilon\), supported in
\(0<\zeta<\delta_\varepsilon\), whose value and first \(\zeta\)-derivative
at zero realize respectively the common trace and the conormal balance;
choose its cap velocity trace to equal that of \(v_\varepsilon\).  A scaled
polynomial cutoff realizes these finitely many boundary jets.  Its V59
position and velocity energy is bounded by a constant times
\(\delta_\varepsilon\) times the squared prescribed jets, while higher graph
norms may grow.  Taking \(\delta_\varepsilon\) sufficiently small keeps the
total continuum energy bounded without changing
\(\mathfrak F_{81}\).  Renormalization by a fixed factor therefore gives the
same \(\varepsilon^{-1}\) lower bound in the full V59 energy phase.  This
contradicts \eqref{eq:V81-false-energy-bound}.
\end{proof}

\begin{firewall}[The physical multiplier cancellation is insufficient]
\label{fire:V81-cancellation-insufficient}
It is correct that the \(r^{-2}\) symbol cancels on the physical trace
range.  Theorem~\ref{thm:V81-energy-unbounded} proves that replacing this
fact by a uniform energy estimate would still lose the residual
\(r^{-1}\) fibre norm.  The exact \(q=0\) chart of V80 and the punctured
limit are therefore different domain statements.
\end{firewall}

\subsection{Sharp weighted infrared domains}
\label{subsec:V81-weighted-domain}

For \(0\le\theta\le1\), define
\begin{equation}
 \|(v,a)\|_{\mathscr E_{\rm IR}^{\theta}}^2
 :=\int_{q,\rho}
 \left(
  r^{-2\theta}|\widehat v|^2
  +r^{-2(1-\theta)}\lambda_r^2|\widehat a|^2
 \right)\dd\rho\,\dd^2q.
 \label{eq:V81-IR-domain}
\end{equation}

\begin{theorem}[One sharp infrared power, with arbitrary split]
\label{thm:V81-weighted-bound}
For every \(0\le\theta\le1\),
\begin{equation}
 \boxed{
 |\mathfrak F_{81}(v,a)|
 \le\frac1{64\kappa}
 \left(\int r^{-2\theta}E_v(q)\,\dd^2q\right)^{1/2}
 \left(\int r^{-2(1-\theta)}E_a(q)\,\dd^2q\right)^{1/2}.}
 \label{eq:V81-weighted-F-bound}
\end{equation}
The symmetric choice \(\theta=1/2\) places \(Q^{-1/2}\) on the velocity and
on the spatial-energy variable \(\lambda a\).  The spaces
\(\mathscr E_{\rm IR}^{\theta}\) are dense in the V59 energy phase and are
proper.  The total power one is sharp for this scale family.
\end{theorem}

\begin{proof}
Theorem~\ref{thm:V81-exact-fibre-norm} gives
\begin{equation*}
 |\mathfrak F_{81}|
 \le\frac1{64\kappa}
 \int r^{-1}\sqrt{E_v(q)E_a(q)}\,\dd^2q.
\end{equation*}
Factor
\(r^{-1}=r^{-\theta}r^{-(1-\theta)}\) and apply Cauchy--Schwarz in
\(q\), proving \eqref{eq:V81-weighted-F-bound}.  Smooth data supported away
from \(q=0\) are dense in the unweighted direct integral, so every displayed
weighted domain is dense.  The packets of
Theorem~\ref{thm:V81-energy-unbounded} have bounded unweighted energy but
their required weighted product grows as \(\varepsilon^{-1}\).  Hence the
unweighted space is strictly larger and the total power cannot be lowered
uniformly on this family.
\end{proof}

\begin{remark}[The freshest Yang--Mills audit lesson]
\label{rem:V81-YM-product-range}
The Yang--Mills companion advanced beyond the V80 snapshot while V80 was
being installed.  Its new product-state reduction reinforces the same
acquisition rule: estimate an operator on the actual admissible range before
taking its ambient operator norm.  V81 does exactly that in
Lemma~\ref{lem:V81-physical-multipliers}.  The remaining
\(r^{-1}\) loss is therefore a genuine physical-range loss, not the earlier
formal \(Q^{-2}\) overestimate.  No Yang--Mills inequality is used here.
\end{remark}

\subsection{What extends to the energy phase}
\label{subsec:V81-combined-extension}

\begin{theorem}[Continuous extension of the combined physical defect]
\label{thm:V81-combined-defect-extension}
Let \(\mathscr D_{\rm sm}\) be the smooth physical Cauchy core with compact
tangential Fourier support away from \(q=0\), augmented by the direct
\(q=0\) chart of V80.  On \(\mathscr D_{\rm sm}\), let
\begin{equation}
 \Delta_{81}^{\rm tot}
 :=\omega_{36}^{\rm side+cap}
   +\omega_{74}^{\rm cor}
   +\omega_{59}^{\rm cont}
   -\omega_{59}^{\rm total}.
 \label{eq:V81-total-defect}
\end{equation}
Then
\begin{equation}
 \Delta_{81}^{\rm tot}=0
 \quad\hbox{on }\mathscr D_{\rm sm},
 \label{eq:V81-core-defect-zero}
\end{equation}
and it has the unique continuous extension equal to zero on the V76/V59
reduced physical energy graph.
\end{theorem}

\begin{proof}
For \(q\ne0\), the side forms agree by
\eqref{eq:V78-side-form-equality} and the cap and corner forms cancel by
\eqref{eq:V79-cap-corner-two-form-cancellation}.  The continuum endpoint is
the V59 form by Lemma~\ref{lem:V79-V59-cap-no-edge}.  At \(q=0\), Theorem
\ref{thm:V80-zero-form-matching} gives the same conclusion directly.  This
proves \eqref{eq:V81-core-defect-zero}.

The V76 radiative quotient is unitary in the V59 homogeneous energy norm
away from the excluded total spatial zero class, and V80 treats the compact
normal-incidence fibre.  Smooth punctured data are dense in that direct
integral.  The zero bilinear form has a unique continuous extension to the
completion, namely zero.
\end{proof}

\begin{firewall}[A zero extended defect is not a closed separated BFV
potential]
\label{fire:V81-no-separated-energy-potential}
Theorem~\ref{thm:V81-combined-defect-extension} proves an energy-level
statement about the \emph{combined reduced physical two-form}.  It does not
extend the separate V74 corner potential, the cap edge potential, or the
generator \(F_{79,J}\) to all V59 energy data.  Theorem
\ref{thm:V81-energy-unbounded} forbids the last extension.  Nor does a zero
physical defect define the off-shell ghost/antifield domain required by a
closed nonlinear BV--BFV action.
\end{firewall}

\subsection{Upstream consequence and V81 ledger}
\label{subsec:V81-ledger}

\begin{proposition}[The next construction should remain unsplit]
\label{prop:V81-upstream-unsplit}
On the full V59 energy phase, an action-domain construction cannot use
\(F_{79,J}\) as a continuous canonical transformation.  Any successful
finite-energy formulation must instead do at least one of the following:
\begin{enumerate}[label=\textup{(\roman*)},leftmargin=3.5em]
\item retain the local Einstein cap current and full V74 cotangent variables
      until their sum is formed, then close the combined weak form before
      gauge fixing; or
\item restrict the action domain to a weighted space such as
      \(\mathscr E_{\rm IR}^{1/2}\) and prove that the evolution preserves
      that domain.
\end{enumerate}
Adding an infrared edge field is warranted only if both routes fail for a
proved reason.
\end{proposition}

\begin{proof}
Continuity of a canonical transformation generated by \(F_{79,J}\) would
imply a bound of the form \eqref{eq:V81-false-energy-bound}, contradicted by
Theorem~\ref{thm:V81-energy-unbounded}.  The weighted alternative is
continuous by Theorem~\ref{thm:V81-weighted-bound}.  The combined weak-form
alternative is not excluded because Theorem
\ref{thm:V81-combined-defect-extension} proves that its physical
antisymmetrization already closes as zero.
\end{proof}

\paragraph{New exact conclusions.}
\begin{enumerate}[label=\textup{(\arabic*)},leftmargin=2.8em]
\item On the physical plus range, the formal \(Q^{-2}\) multiplier reduces
      to the exact symbols \eqref{eq:V81-s-multiplier}--%
      \eqref{eq:V81-BDns-multiplier}.
\item The two cap trace functionals have exact squared norms
      \(r/8\) and \(1/(32r^3)\); their bilinear product norm is
      \(1/(16r)\).
\item A normalized smooth infrared packet proves that the V79 generator is
      unbounded on the V59 energy phase.
\item The generator is continuous on the sharp family of weighted domains
      \eqref{eq:V81-IR-domain}; the symmetric domain uses one half infrared
      derivative on each Cauchy factor.
\item The complete reduced physical two-form defect nevertheless extends
      uniquely as zero to the energy completion.
\item The unresolved finite-energy problem is the off-shell, separated
      action domain, so the attack moves upstream to the unsplit local
      current.
\end{enumerate}

\paragraph{Next forced step.}
V82 must write the Einstein cap current, V74 corner cotangent current, and
V59 endpoint current as one local weak Green identity before solving the
moment equation.  It should identify a closed graph norm in which the
combined current is continuous and the residual gauge action is closable.
The physical restriction must reproduce
Theorem~\ref{thm:V81-combined-defect-extension}; any proposed separated
domain must be tested against the packet
\eqref{eq:V81-unbounded-packet}.  Only after this weak domain exists should
the programme return to the Brown--York/Bianchi stress identity.

\begin{assessment}[V81 continuum status]
The full Standard--Model--Einstein continuum remains unproved.  V81 gives
the sharp infrared diagnosis of the V79 generator and closes the combined
reduced physical two-form on the V59 energy completion.  A closed off-shell
compact BV--BFV action domain, nonlinear curved gluing,
Brown--York/Bianchi stress identity, common renormalized Standard Model
stress domain, and cofinal interacting continuum remain open.
\end{assessment}

\section*{V81 exact checks and append-only record}
\addcontentsline{toc}{section}{V81 exact checks and append-only record}

The exact one-dimensional integrals used in the fibre norm are
\begin{equation}
 \int_{\mathbb R}(r^2+\rho^2)^{-2}\dd\rho
 =\frac{\pi}{2r^3},
 \qquad
 \int_{\mathbb R}\rho^2(r^2+\rho^2)^{-3}\dd\rho
 =\frac{\pi}{8r^3}.
 \label{eq:V81-integral-check-card}
\end{equation}
With the two joint-restriction factors \((2\pi)^{-1/2}\), they give
\(\sqrt{(r/8)(1/(32r^3))}=1/(16r)\), the exact residual
infrared loss.

The V80 predecessor used for this append has SHA--256
\begin{center}
\texttt{16035310193d0c9068f44145ffc82fd6a503425afd89059d413fd4f9b1e7e6c0}.
\end{center}
The V81 candidate retains its bytes up to the sole live final document
terminator, appends this material, and restores exactly one active
\verb|\end{document}|.  Validation includes balanced environments and
braces, unique labels, resolved new internal references, valid inline and
display math delimiters, exact predecessor-prefix equality, and a final
inventory which removes only the superseded SM V80 file.

% =============================================================================
% END APPENDED VERSION V81 MATERIAL
% =============================================================================
\clearpage
% =============================================================================
% BEGIN APPENDED VERSION V82 MATERIAL
% =============================================================================
\section{V82: an unsplit Sobolev--Green corner phase and closed residual
reduction}
\label{sec:V82-unsplit-domain}

V81 proved that the separated cap generator is not continuous in the V59
energy norm, even after restriction to the physical TT range.  That result
does not obstruct the original local V72--V75 system: the latter contains
the metric Green current and the full V74 corner cotangent variables before
the moment equation is solved.  This note gives that unsplit system a closed
Hilbert realization on a compact flat slab.

The key regularity gain comes from the exact divergence identity
\(D_a\mathscr P^{ab}(f)=0\).  If the residual normal component has
\(H^2\) regularity on a face, then \(\mathscr P(f)\in L^2\) and belongs
columnwise to \(H(\operatorname{div})\).  Its edge-normal trace is therefore
in \(H^{-1/2}\), where it pairs canonically with the
\(H^{1/2}\) residual edge value.  The V72 cocycle and the V74 opposite
cocycle consequently extend as continuous duality pairings.  No inverse
Laplacian, polarization frame, or infrared spectral gap is used.

This constructs a closed off-shell residual/corner reduction domain and
extends the algebraic master identity to it.  It does not yet prove that the
full nonlinear Einstein--Standard Model evolution preserves this graph
domain, nor that the separate cap and corner potentials are continuous in
the unweighted V59 energy norm.

\subsection{Face and joint Hilbert phases}
\label{subsec:V82-Hilbert-phases}

Let \(\Gamma\) be one compact smooth timelike face of a flat slab and let
\(J=\partial\Gamma\) be the oriented union of its spacelike joints.  Write
\(E_\Gamma\) for the finite-rank bundle of symmetric spacetime two-tensors
restricted to \(\Gamma\), and \(E_J:=T^*\Gamma|_J\) for the three residual
tangential components.  All Sobolev and (H(operatorname{div})) norms use
a fixed positive auxiliary product metric.  Lorentzian index signs act only
as bounded constant bundle maps and do not change the trace theorem.  At the
critical half order, (H^{-1/2}(Gamma)) means by definition the continuous
dual of the chosen restriction-space realization of (H^{1/2}(Gamma));
this avoids any (H^{1/2}_{00}) convention ambiguity on a face with edge.

Define the metric boundary phase
\begin{equation}
 \mathcal M_{82}(\Gamma)
 :=H^{1/2}(\Gamma;E_\Gamma)
   \oplus H^{-1/2}(\Gamma;E_\Gamma),
 \label{eq:V82-metric-phase}
\end{equation}
with points \(z=(H,p)\) and canonical form
\begin{equation}
 \omega_{M,82}(z_1,z_2)
 :=\frac1{2\kappa}
 \left(
  \langle p_1,H_2\rangle_{-1/2,1/2}
  -\langle p_2,H_1\rangle_{-1/2,1/2}
 \right).
 \label{eq:V82-metric-form}
\end{equation}
Define the joint cotangent phase
\begin{equation}
 \mathcal C_{82}(J)
 :=H^{1/2}(J;E_J)
   \oplus H^{-1/2}(J;E_J^*),
 \label{eq:V82-corner-phase}
\end{equation}
with points \(c=(X,\Pi)\) and form
\begin{equation}
 \omega_{C,82}(c_1,c_2)
 :=\frac1\kappa
 \left(
  \langle\Pi_1,X_2\rangle_{-1/2,1/2}
  -\langle\Pi_2,X_1\rangle_{-1/2,1/2}
 \right).
 \label{eq:V82-corner-form}
\end{equation}
For several faces and joints, take the finite oriented direct sum.

\begin{lemma}[The canonical Sobolev forms are strong and nondegenerate]
\label{lem:V82-canonical-forms}
The forms \eqref{eq:V82-metric-form} and
\eqref{eq:V82-corner-form} are continuous, strongly nondegenerate Hilbert
symplectic forms.  Their product form
\begin{equation}
 \omega_{82}:=\omega_{M,82}\oplus\omega_{C,82}
 \label{eq:V82-product-form}
\end{equation}
is strongly nondegenerate.
\end{lemma}

\begin{proof}
For every real Hilbert Sobolev space \(V=H^{1/2}\), its continuous dual is
\(V^*=H^{-1/2}\) under the displayed duality.  The canonical form on
\(V\oplus V^*\) sends \((x,\xi)\) to the dual functional
\((y,\eta)\mapsto\langle\xi,y\rangle-\langle\eta,x\rangle\).  The Riesz
isomorphisms of \(V\) and \(V^*\) show that this map to
\((V\oplus V^*)^*\) is bounded and onto.  Multiplication by the nonzero
constants \((2\kappa)^{-1}\) and \(\kappa^{-1}\), and finite direct sums,
preserve this property.
\end{proof}

\subsection{Residual Cauchy jets and the weak normal trace}
\label{subsec:V82-residual-jets}

On the face, retain the boundary Cauchy jet of an arbitrary off-shell ghost
or residual parameter:
\begin{equation}
 e_\mu:=\epsilon_\mu|_\Gamma,
 \qquad
 d_\mu:=D_n\epsilon_\mu|_\Gamma.
 \label{eq:V82-residual-Cauchy-jet}
\end{equation}
Set
\begin{equation}
 \mathfrak G_{82}(\Gamma)
 :=H^2(\Gamma;T^*M|_\Gamma)
   \oplus H^1(\Gamma;T^*M|_\Gamma),
 \qquad \epsilon_\Gamma:=(e,d).
 \label{eq:V82-ghost-jet-space}
\end{equation}
This is a space of face jets; no wave equation is imposed.

The boundary value of the metric gauge variation is the bounded first-order
map
\begin{equation}
 G_\Gamma(e,d):=(\mathcal G\epsilon)|_\Gamma
 \in H^1(\Gamma;E_\Gamma).
 \label{eq:V82-G-map}
\end{equation}
Indeed, its tangential components use \(D_ae_b\), and its components with a
normal index use \(d_b+D_be_n\) or \(2d_n\).  Define
\begin{equation}
 P_\Gamma(e,d)^{\mu\nu}
 :=e_a^\mu e_b^\nu\mathscr P^{ab}(e_n)
 \in L^2(\Gamma;E_\Gamma).
 \label{eq:V82-P-map}
\end{equation}

\begin{lemma}[Bounded joint trace and normal-Hessian trace]
\label{lem:V82-TR-bounded}
The maps
\begin{align}
 T_J(e,d)_b&:=e_b|_J,
 &T_J&:\mathfrak G_{82}(\Gamma)longrightarrow H^{1/2}(J;E_J),
 \label{eq:V82-T-map}\\
 R_J(e,d)^b&:=\nu_a\mathscr P^{ab}(e_n)|_J,
 &R_J&:\mathfrak G_{82}(\Gamma)longrightarrow H^{-1/2}(J;E_J^*)
 \label{eq:V82-R-map}
\end{align}
are bounded.  The second line is the weak normal trace of an
\(H(\operatorname{div})\) tensor and agrees with the classical V74 trace on
smooth jets.
\end{lemma}

\begin{proof}
The ordinary trace theorem sends \(H^2(\Gamma)\) continuously to
\(H^{3/2}(J)\), which embeds continuously in \(H^{1/2}(J)\).  This proves
\eqref{eq:V82-T-map}.

Put \(S^{ab}:=\mathscr P^{ab}(e_n)\).  Since \(e_n\in H^2(\Gamma)\), one
has \(S\in L^2(\Gamma)\).  The distributional identity
\begin{equation}
 D_aS^{ab}=0
 \label{eq:V82-P-divergence-weak}
\end{equation}
holds by commutation of the flat tangential derivatives.  For each fixed
\(b\), the column \(S^{\cdot b}\) therefore belongs to
\(H(\operatorname{div};\Gamma)\), with norm bounded by
\(C\|e_n\|_{H^2}\).  The normal-trace theorem for
\(H(\operatorname{div})\) gives
\begin{equation*}
 \|\nu_aS^{ab}\|_{H^{-1/2}(J)}
 \le C\left(
  \|S^{\cdot b}\|_{L^2(\Gamma)}
  +\|D_aS^{ab}\|_{L^2(\Gamma)}
 \right)
 \le C\|e_n\|_{H^2(\Gamma)}.
\end{equation*}
This proves \eqref{eq:V82-R-map}.  Density of smooth tensors in the graph
space and uniqueness of the weak normal trace prove agreement with the
classical formula.
\end{proof}

\begin{corollary}[Weak Green identity for the trace-adjusted Hessian]
\label{cor:V82-weak-Green}
For \(\epsilon_\Gamma,\zeta_\Gamma\in\mathfrak G_{82}(\Gamma)\),
\begin{equation}
 \int_\Gamma
 \mathscr P^{ab}(e_n)(D_a\zeta_b+D_b\zeta_a)
 =2\langle R_J(\epsilon),T_J(\zeta)\rangle_{-1/2,1/2},
 \label{eq:V82-weak-Green-identity}
\end{equation}
with the oriented sum over components of \(J\).
\end{corollary}

\begin{proof}
For smooth jets, symmetry of \(\mathscr P\), integration by parts, and
\eqref{eq:V72-P-divergence} give the identity.  The left side is continuous
under \(H^2\times H^2\) convergence, and the right side is continuous by
Lemma~\ref{lem:V82-TR-bounded}.  Smooth jets are dense in
\(\mathfrak G_{82}\), so the identity extends.
\end{proof}

\subsection{Bounded unsplit residual action and cocycle cancellation}
\label{subsec:V82-unsplit-action}

Define the metric and corner fundamental-vector maps
\begin{align}
 A_{82}\epsilon_\Gamma
 &:=(G_\Gamma(e,d),P_\Gamma(e,d))
 \in\mathcal M_{82}(\Gamma),
 \label{eq:V82-A-map}\\
 B_{82}\epsilon_\Gamma
 &:=(T_J(e,d),-R_J(e,d))
 \in\mathcal C_{82}(J).
 \label{eq:V82-B-map}
\end{align}
The Sobolev embeddings \(H^1\hookrightarrow H^{1/2}\) and
\(L^2\hookrightarrow H^{-1/2}\), together with
Lemma~\ref{lem:V82-TR-bounded}, show that both maps are bounded.

\begin{theorem}[Continuous metric cocycle and exact opposite corner
cocycle]
\label{thm:V82-cocycle-cancellation}
For all \(\epsilon,\zeta\in\mathfrak G_{82}(\Gamma)\),
\begin{align}
 \omega_{M,82}(A_{82}\epsilon,A_{82}\zeta)
 &=\frac1\kappa\left[
  \langle R_J(\epsilon),T_J(\zeta)\rangle
  -\langle R_J(\zeta),T_J(\epsilon)\rangle
 \right],
 \label{eq:V82-metric-cocycle}\\
 \omega_{C,82}(B_{82}\epsilon,B_{82}\zeta)
 &=-\omega_{M,82}(A_{82}\epsilon,A_{82}\zeta).
 \label{eq:V82-opposite-cocycle}
\end{align}
Consequently the combined map
\begin{equation}
 V_{82}:=(A_{82},B_{82}):
 \mathfrak G_{82}\longrightarrow
 \mathcal P_{82}:=\mathcal M_{82}\oplus\mathcal C_{82}
 \label{eq:V82-combined-map}
\end{equation}
has isotropic range.  All displayed maps and bilinear forms are continuous;
none contains \(Q^{-1}\) or \(Q^{-2}\).
\end{theorem}

\begin{proof}
Insert \eqref{eq:V82-A-map} in the metric form.  The two volume terms are
the weak versions of the two terms in the V72 pure-gauge pairing.  Apply
Corollary~\ref{cor:V82-weak-Green} first with
\((\epsilon,\zeta)\), then with \((\zeta,\epsilon)\), including the factor
\((2\kappa)^{-1}\).  This gives
\eqref{eq:V82-metric-cocycle}.  Substitution of
\eqref{eq:V82-B-map} in the canonical corner form gives its negative
directly.  Their sum is zero, proving isotropy.  Boundedness was established
above, and the formulas are differential and duality-local.
\end{proof}

For \(Z=(H,p,X,\Pi)\in\mathcal P_{82}\), define
\begin{align}
 \widetilde\mu_{82,\epsilon}(Z)
 &:=\omega_{82}(V_{82}\epsilon,Z)
 \notag\\
 &=\frac1{2\kappa}\left[
   \langle P_\Gamma(e,d),H\rangle
  -\langle p,G_\Gamma(e,d)\rangle
 \right]
 \notag\\
 &\quad-\frac1\kappa\left[
   \langle R_J(e,d),X\rangle
  +\langle\Pi,T_J(e,d)\rangle
 \right].
 \label{eq:V82-combined-moment}
\end{align}

\begin{corollary}[Continuous equivariant moment map]
\label{cor:V82-continuous-moment}
The map
\begin{equation}
 (\epsilon,Z)\longmapsto\widetilde\mu_{82,\epsilon}(Z)
 \label{eq:V82-moment-bilinear}
\end{equation}
is continuous on \(\mathfrak G_{82}\times\mathcal P_{82}\), and
\begin{equation}
 \{\widetilde\mu_{82,\epsilon},
    \widetilde\mu_{82,\zeta}\}=0.
 \label{eq:V82-moment-commutation}
\end{equation}
It is the Sobolev extension of the V74/V75 residual moment map.
\end{corollary}

\begin{proof}
Every term in \eqref{eq:V82-combined-moment} is one of the continuous
\(H^{-1/2},H^{1/2}\) dualities already established.  Its Hamiltonian vector
is the constant vector \(V_{82}\epsilon\).  The bracket of two linear
Hamiltonians is therefore
\(\omega_{82}(V_{82}\epsilon,V_{82}\zeta)=0\) by
Theorem~\ref{thm:V82-cocycle-cancellation}.  Agreement on smooth data with
\eqref{eq:V74-equivariant-total-moment-map} proves the last assertion.
\end{proof}

\subsection{Closed-orbit Hilbert reduction}
\label{subsec:V82-closed-reduction}

The range of a bounded operator need not be closed.  To avoid a hidden
Hausdorff assumption, put
\begin{equation}
 W_{82}:=\overline{\operatorname{ran}V_{82}}^{\mathcal P_{82}}.
 \label{eq:V82-closed-orbit}
\end{equation}

\begin{theorem}[Closed weak zero level and nondegenerate reduced form]
\label{thm:V82-closed-reduction}
The closed orbit space \(W_{82}\) is isotropic.  The simultaneous weak zero
level of the continuous moments is
\begin{equation}
 \mathcal N_{82}
 :=\left\{Z:\widetilde\mu_{82,\epsilon}(Z)=0
        \text{ for all }\epsilon\in\mathfrak G_{82}\right\}
 =W_{82}^{\perp_{\omega_{82}}}.
 \label{eq:V82-weak-zero-level}
\end{equation}
It is closed, contains \(W_{82}\), and the quotient
\begin{equation}
 \mathcal P_{82}^{\rm red}:=\mathcal N_{82}/W_{82}
 \label{eq:V82-reduced-phase}
\end{equation}
is a Hilbert space carrying the continuous nondegenerate reduced form
induced by \(\omega_{82}\).
\end{theorem}

\begin{proof}
Theorem~\ref{thm:V82-cocycle-cancellation} makes
\(\operatorname{ran}V_{82}\) isotropic.  Continuity of \(\omega_{82}\)
makes its closure isotropic.  The common kernel of the continuous linear
functionals \(Z\mapsto\omega_{82}(V_{82}\epsilon,Z)\) is closed and equals
the symplectic orthogonal of the closure of their span, which proves
\eqref{eq:V82-weak-zero-level}.  Isotropy gives
\(W_{82}\subset W_{82}^{\perp_{\omega_{82}}}\), so the form descends.

Because \(\omega_{82}\) is strongly nondegenerate and \(W_{82}\) is a
closed subspace,
\begin{equation}
 (W_{82}^{\perp_{\omega_{82}}})^{\perp_{\omega_{82}}}=W_{82}.
 \label{eq:V82-double-orthogonal}
\end{equation}
If a reduced class pairs to zero with every reduced class, its
representative lies in the left side of
\eqref{eq:V82-double-orthogonal}, hence in \(W_{82}\).  Thus the reduced
form is nondegenerate.  Both numerator and denominator are closed Hilbert
subspaces, so their quotient is Hilbert.
\end{proof}

\begin{remark}[Closure of the orbit is part of the result]
\label{rem:V82-orbit-closure}
V82 does not assert that \(\operatorname{ran}V_{82}\) itself is closed.
Quotienting by its closure is the analytic version of dividing by residual
directions which are limits in the boundary graph norm.  A later Fredholm or
closed-range theorem may identify the closure with the range, but the weak
reduction above does not depend on that additional claim.
\end{remark}

\subsection{Extension of the off-shell residual master identity}
\label{subsec:V82-master-extension}

Let \(\gamma\in\mathfrak G_{82}\) be odd.  Define the residual part of the
boundary charge by
\begin{equation}
 \Omega_{82}^{\rm res}(Z,\gamma)
 :=\widetilde\mu_{82,\gamma}(Z).
 \label{eq:V82-residual-charge}
\end{equation}
The nonminimal V75 doublet sector may be completed in any matched Sobolev
cotangent pair \(Y\oplus Y^*\) for which
\(\langle\pi_{\bar\gamma},B\rangle\) is continuous; denote its unchanged
cotangent-lift charge by \(\Omega_{75}^{\rm nm}\).

\begin{theorem}[Sobolev-continuous off-shell master equation]
\label{thm:V82-master-equation}
On \(\mathcal P_{82}\times\mathfrak G_{82}\) and a matched nonminimal
cotangent domain,
\begin{equation}
 \Omega_{82}^{\rm off}
 :=\Omega_{82}^{\rm res}+\Omega_{75}^{\rm nm}
 \label{eq:V82-offshell-charge}
\end{equation}
is continuous in the boundary phase variables and satisfies
\begin{equation}
 \boxed{
 \{\Omega_{82}^{\rm off},\Omega_{82}^{\rm off}\}=0.}
 \label{eq:V82-master-equation}
\end{equation}
No residual wave equation and no inverse tangential Laplacian is used.
\end{theorem}

\begin{proof}
Continuity of the residual charge is
Corollary~\ref{cor:V82-continuous-moment}.  Its self-bracket is zero by
\eqref{eq:V82-moment-commutation}.  The nonminimal charge is the cotangent
lift of the square-zero doublet differential
\(s\bar\gamma=B\), \(sB=0\), so its self-bracket is zero.  As in V75, the
two charges use disjoint conjugate momenta and their cross bracket vanishes.
The smooth master identity therefore holds term by term and extends by
continuity to the completed domain.
\end{proof}

\begin{firewall}[What the completed master equation does not yet say]
\label{fire:V82-master-scope}
Equation~\eqref{eq:V82-master-equation} is a boundary residual/corner master
identity on the specified flat Sobolev phase.  It does not construct the
bulk nonlinear BV Laplacian, prove quantum anomaly cancellation, or prove
invariance of this domain under a nonlinear Einstein--Standard Model
evolution.  The matched nonminimal factor is routine cotangent analysis;
the unresolved analytic content is the interacting bulk-to-boundary graph.
\end{firewall}

\subsection{Regression against the V81 infrared packet}
\label{subsec:V82-infrared-regression}

\begin{proposition}[The unsplit formulation contains no separated
\(r^{-1}\) multiplier]
\label{prop:V82-packet-regression}
Every smooth V81 packet belongs to the smooth core of
\(\mathcal P_{82}\) together with smooth corner data.  On that core, the
combined residual cocycle is identically zero by
\eqref{eq:V82-opposite-cocycle}, independently of its tangential Fourier
support.  In particular, forming the V82 combined moment and bracket does
not produce the separated scalar functional
\(\mathfrak F_{81}\) or its \(r^{-1}\) fibre norm.
\end{proposition}

\begin{proof}
For each \(\varepsilon>0\), the cut-off packet and its continuum boundary
jet are smooth and have compact tangential frequency support away from zero,
so all classical traces belong to the stated Sobolev spaces.  The V82
bracket is computed by the local weak Green identity
\eqref{eq:V82-weak-Green-identity}; its metric edge term is canceled by the
canonical corner term before any zero-level equation is solved.  The
operator \(Q^{-2}\) appears only after choosing the V78 dressing gauge and
eliminating \((X,\Pi)\).  Since neither operation occurs in
\eqref{eq:V82-combined-moment}, the fibre norm
\eqref{eq:V81-fibre-product-norm} is absent.
\end{proof}

\begin{assessment}[Meaning of the regression]
\label{ass:V82-regression-meaning}
Proposition~\ref{prop:V82-packet-regression} does not prove a uniform bound
of every V82 trace by the V59 energy norm; V81 rules out such a statement
for the eliminated generator.  It proves that the unbounded factor is
created by separation and gauge solution, not by the local cocycle
cancellation.  The V82 graph norm is therefore a valid off-shell closure
domain, while equality of its reduced physical completion with the entire
V59 energy phase remains to be proved.
\end{assessment}

\subsection{V82 result ledger and next acquisition}
\label{subsec:V82-ledger}

\paragraph{New exact conclusions.}
\begin{enumerate}[label=\textup{(\arabic*)},leftmargin=2.8em]
\item The metric boundary phase closes in the canonical pair
      \(H^{1/2}(\Gamma)\oplus H^{-1/2}(\Gamma)\), and the joint phase closes
      in the analogous \(H^{1/2}(J)\oplus H^{-1/2}(J)\) pair.
\item Off-shell residual face Cauchy jets in \(H^2\oplus H^1\) map boundedly
      to both phases.
\item The divergence-free trace-adjusted Hessian has a bounded weak
      \(H^{-1/2}\) edge-normal trace by the
      \(H(\operatorname{div})\) theorem.
\item The weak Green identity \eqref{eq:V82-weak-Green-identity} extends the
      V72 corner calculation without requiring \(s>3\) bulk regularity.
\item The metric cocycle and full V74 opposite cocycle cancel continuously
      before any dressing equation or inverse Laplacian is used.
\item Closing the combined residual orbit gives the closed zero level and
      the nondegenerate Hilbert reduced phase
      \eqref{eq:V82-reduced-phase}.
\item The V75 off-shell residual master identity extends continuously to
      this Sobolev boundary domain.
\item The construction passes the V81 packet regression in the precise
      sense that no separated \(r^{-1}\) generator is formed.
\end{enumerate}

\paragraph{Next forced step.}
V83 must compare the physical part of
\(\mathcal P_{82}^{\rm red}\) with the V59 energy phase.  The required map
is the closure of the V63/V76 radiative quotient on the compact face--joint
graph.  It must prove either that every V59 energy datum has a representative
in the V82 weak zero level with graph-norm approximation, or identify the
exact proper closed subspace.  The proof must keep the corner variables
uneliminated and must test surjectivity on the V81 infrared packets.  Only a
closed-range/density theorem, not the smooth fibrewise bijection, can settle
this question.

\begin{assessment}[V82 continuum status]
The full Standard--Model--Einstein continuum remains unproved.  V82
constructs the first closed unsplit Sobolev realization of the flat
off-shell residual/corner master system and its Hilbert reduction, avoiding
the V81 infrared-divergent generator.  Identification of that reduced graph
with the full V59 energy phase, nonlinear curved gluing, the
Brown--York/Bianchi stress identity, the common renormalized Standard Model
stress domain, and the cofinal interacting continuum remain open.
\end{assessment}

\section*{V82 exact checks and append-only record}
\addcontentsline{toc}{section}{V82 exact checks and append-only record}

The regularity and duality check card is
\begin{equation}
 H^2(\Gamma)\xrightarrow{\ \mathscr P\ }L^2(\Gamma)
 \cap H(\operatorname{div};\Gamma)
 \xrightarrow{\ \nu\cdot\ }H^{-1/2}(J),
 \qquad
 H^2(\Gamma)\xrightarrow{\ \operatorname{Tr}\ }H^{3/2}(J)
 \hookrightarrow H^{1/2}(J).
 \label{eq:V82-regularity-check-card}
\end{equation}
The sign check is
\begin{equation}
 \omega_{M,82}(A_{82}\epsilon,A_{82}\zeta)
 +\omega_{C,82}(B_{82}\epsilon,B_{82}\zeta)=0,
 \label{eq:V82-sign-check-card}
\end{equation}
which is the weak Sobolev form of the V72/V74 cocycle cancellation.

The V81 predecessor used for this append has SHA--256
\begin{center}
\texttt{11c89e8a86f71281cfc173a546e7a6ffcc7f8621335e1e19254d624174a40606}.
\end{center}
The V82 candidate retains its bytes up to the sole live final document
terminator, appends this material, and restores exactly one active
\verb|\end{document}|.  Validation includes balanced environments and
braces, unique labels, resolved new internal references, valid inline and
display math delimiters, exact predecessor-prefix equality, and a final
inventory which removes only the superseded SM V81 file.

% =============================================================================
% END APPENDED VERSION V82 MATERIAL
% =============================================================================
\clearpage
% =============================================================================
% BEGIN APPENDED VERSION V83 MATERIAL
% =============================================================================
\section{V83: the critical trace balance, a supercritical closed phase,
and the exact energy gate}
\label{sec:V83-critical-scale}

V82 supplies a closed unsplit residual/corner phase at the sufficient face
regularity \(H^2\oplus H^1\).  This note asks for the lowest Sobolev orders
at which the same local Green cancellation could close.  Formal order
counting selects the face jet \(H^{3/2}\oplus H^{1/2}\), but its
normal-Hessian trace lies at the endpoint of the standard fractional normal
trace theorem.  The endpoint must not be inserted by continuity without a
separate theorem.

We therefore prove the construction at every strictly supercritical order
\(H^{3/2+\delta}\oplus H^{1/2+\delta}\), \(0<\delta<1/2\), using exact
dual Sobolev phases obtained by continuous downcasting.  We then isolate the
remaining comparison with the V59 physical energy completion as one
necessary-and-sufficient quotient-norm estimate.  Finally, an exact Fourier
calculation shows why the explicit V78 dressing section cannot prove or
disprove that estimate.

\subsection{Append-only typographic clarification for V82}
\label{subsec:V83-V82-errata}

The append-only V82 text contains five lost TeX escape marks in prose and
two displayed arrows.  The strings
\[
 H(\texttt{operatorname\{div\}}),\qquad
 (H^{-1/2}(\texttt{Gamma})),\qquad
 (H^{1/2}(\texttt{Gamma}))
\]
there mean, respectively,
\[
 H(\operatorname{div}),\qquad H^{-1/2}(\Gamma),\qquad
 H^{1/2}(\Gamma),
\]
and the two occurrences of
\(\mathfrak G_{82}(\Gamma)\texttt{longrightarrow}\) in
the V82 \(T_J,R_J\) display mean
\(\mathfrak G_{82}(\Gamma)\longrightarrow\).
These are typographic corrections only.  The definitions and proofs used
below are the mathematical formulas just displayed; no earlier byte is
altered.

\subsection{Formal balance and the endpoint obstruction}
\label{subsec:V83-formal-balance}

Let the residual value \(e=\epsilon|_\Gamma\) have face order \(H^\alpha\)
and its normal derivative \(d=D_n\epsilon|_\Gamma\) have order
\(H^{\alpha-1}\).  The four local maps used by V72--V75 then have the formal
orders
\begin{equation}
\begin{array}{c|c}
\text{map}&\text{Sobolev order}\\ \hline
G_\Gamma(e,d)=(\mathcal G\epsilon)|_\Gamma&\alpha-1\\
P_\Gamma(e,d)=\mathscr P(e_n)&\alpha-2\\
T_J(e,d)=e|_J&\alpha-\tfrac12\\
R_J(e,d)=\nu\mathbin{\lrcorner}\mathscr P(e_n)|_J&
          \alpha-\tfrac52 .
\end{array}
\label{eq:V83-order-table}
\end{equation}
Here the last line is the expected normal trace order for a divergence-free
tensor of order \(\alpha-2\).

\begin{proposition}[Unique formal exact-duality exponent]
\label{prop:V83-unique-balance}
If both pairs in \eqref{eq:V83-order-table} are required to occur at their
undiminished orders and to be exact Sobolev duals, then
\begin{equation}
 \alpha=\frac32.
 \label{eq:V83-critical-alpha}
\end{equation}
The resulting formal phase orders are
\begin{equation}
 \bigl(H^{1/2}\oplus H^{-1/2}\bigr)_\Gamma
 \oplus
 \bigl(H^1\oplus H^{-1}\bigr)_J .
 \label{eq:V83-formal-critical-phase}
\end{equation}
\end{proposition}

\begin{proof}
Exact duality of the metric orders gives
\((\alpha-1)+(\alpha-2)=0\).  Exact duality of the joint orders gives
\((\alpha-\tfrac12)+(\alpha-\tfrac52)=0\).  Both equations are
\(2\alpha-3=0\), proving \eqref{eq:V83-critical-alpha}; substitution gives
\eqref{eq:V83-formal-critical-phase}.
\end{proof}

\begin{firewall}[The formal endpoint is not yet a trace theorem]
\label{fire:V83-endpoint}
For \(\alpha=3/2\), the tensor
\(\mathscr P(e_n)\) has order \(H^{-1/2}\).  The standard fractional
normal-trace construction applies for tensor order \(s>-1/2\), whereas
\(s=-1/2\) is its lower endpoint.  The V82 \(H(\operatorname{div})\)
argument also does not apply, because the tensor need not be in \(L^2\).
Thus \eqref{eq:V83-formal-critical-phase} is an order balance, not a proved
closed phase.  An endpoint Besov or Lions--Magenes theorem, or a
counterexample, is still required to decide that exact scale.
\end{firewall}

\subsection{The fractional normal trace above the endpoint}
\label{subsec:V83-fractional-trace}

Fix for the rest of the note
\begin{equation}
 0<\delta<\frac12.
 \label{eq:V83-delta-range}
\end{equation}
Sobolev spaces on a face with edge are restriction spaces, and negative
orders denote the continuous duals of the indicated positive-order spaces.

\begin{lemma}[Fractional divergence-normal trace]
\label{lem:V83-fractional-normal}
Let \(S\in H^{-1/2+\delta}(\Gamma;T\Gamma\otimes E)\) and suppose
\(D_aS^{ab}=0\) in distributions.  There is a unique
\[
 \gamma_\nu S\in H^{-1+\delta}(J;E)
\]
which equals \(\nu_aS^{ab}|_J\) for smooth \(S\), and
\begin{equation}
 \|\gamma_\nu S\|_{H^{-1+\delta}(J)}
 \le C_\delta\|S\|_{H^{-1/2+\delta}(\Gamma)}.
 \label{eq:V83-fractional-normal-bound}
\end{equation}
\end{lemma}

\begin{proof}
The trace map
\[
 \gamma_0:H^{3/2-\delta}(\Gamma;E^*)
 \longrightarrow H^{1-\delta}(J;E^*)
\]
has a bounded right inverse \(\mathcal E_\delta\).  For
\(\phi\in H^{1-\delta}(J;E^*)\), set
\begin{equation}
 \langle\gamma_\nu S,\phi\rangle_{-1+\delta,\,1-\delta}
 :=
 \langle S,D(\mathcal E_\delta\phi)\rangle_
 {-1/2+\delta,\,1/2-\delta}.
 \label{eq:V83-normal-definition}
\end{equation}
If two liftings have the same trace, their difference belongs to the kernel
of \(\gamma_0\).  Since
\(1/2<3/2-\delta<3/2\), smooth fields compactly supported away from \(J\)
are dense in that kernel.  Distributional integration by parts and
\(D_aS^{ab}=0\) make the right side of
\eqref{eq:V83-normal-definition} vanish on the difference.  The definition
is therefore independent of the lifting.

Moreover,
\[
 |\langle\gamma_\nu S,\phi\rangle|
 \le
 \|S\|_{H^{-1/2+\delta}}
 \|\mathcal E_\delta\phi\|_{H^{3/2-\delta}}
 \le C_\delta\|S\|_{H^{-1/2+\delta}}
                \|\phi\|_{H^{1-\delta}},
\]
which proves \eqref{eq:V83-fractional-normal-bound} and uniqueness.
For smooth \(S\), ordinary integration by parts identifies
\eqref{eq:V83-normal-definition} with the classical normal trace.
\end{proof}

\subsection{A proved supercritical canonical phase}
\label{subsec:V83-supercritical-phase}

Define the residual face-jet space
\begin{equation}
 \mathfrak G_{83}^{\delta}(\Gamma)
 :=
 H^{3/2+\delta}(\Gamma;T^*M|_\Gamma)
 \oplus
 H^{1/2+\delta}(\Gamma;T^*M|_\Gamma).
 \label{eq:V83-residual-space}
\end{equation}
For \((e,d)\in\mathfrak G_{83}^{\delta}\), the first-order metric gauge
trace and the second-order trace-adjusted Hessian satisfy
\begin{align}
 G_\Gamma(e,d)&\in H^{1/2+\delta}(\Gamma;E_\Gamma),
 \label{eq:V83-G-order}\\
 P_\Gamma(e,d)=\mathscr P(e_n)&\in
 H^{-1/2+\delta}(\Gamma;E_\Gamma).
 \label{eq:V83-P-order}
\end{align}
The flat identity \(D_a\mathscr P^{ab}(e_n)=0\) holds distributionally.
The ordinary trace theorem and
Lemma~\ref{lem:V83-fractional-normal} give bounded maps
\begin{align}
 T_J(e,d)&:=e|_J\in H^{1+\delta}(J;E_J),
 \label{eq:V83-T-order}\\
 R_J(e,d)&:=\gamma_\nu\mathscr P(e_n)
             \in H^{-1+\delta}(J;E_J^*).
 \label{eq:V83-R-order}
\end{align}

Use the continuous embeddings
\[
 H^{1/2+\delta}(\Gamma)\hookrightarrow H^{1/2-\delta}(\Gamma),
 \qquad
 H^{1+\delta}(J)\hookrightarrow H^{1-\delta}(J)
\]
to define exact cotangent phases
\begin{align}
 \mathcal M_{83}^{\delta}(\Gamma)
 &:=
 H^{1/2-\delta}(\Gamma;E_\Gamma)
 \oplus H^{-1/2+\delta}(\Gamma;E_\Gamma^*),
 \label{eq:V83-metric-phase}\\
 \mathcal C_{83}^{\delta}(J)
 &:=
 H^{1-\delta}(J;E_J)
 \oplus H^{-1+\delta}(J;E_J^*).
 \label{eq:V83-corner-phase}
\end{align}
For points \(z=(H,p)\) and \(c=(X,\Pi)\), set
\begin{align}
 \omega_{M,83}^{\delta}(z_1,z_2)
 &:=
 \frac1{2\kappa}
 \left(
 \langle p_1,H_2\rangle_{-1/2+\delta,\,1/2-\delta}
 -
 \langle p_2,H_1\rangle_{-1/2+\delta,\,1/2-\delta}
 \right),
 \label{eq:V83-metric-form}\\
 \omega_{C,83}^{\delta}(c_1,c_2)
 &:=
 \frac1\kappa
 \left(
 \langle\Pi_1,X_2\rangle_{-1+\delta,\,1-\delta}
 -
 \langle\Pi_2,X_1\rangle_{-1+\delta,\,1-\delta}
 \right).
 \label{eq:V83-corner-form}
\end{align}

\begin{lemma}[Weak Green identity at supercritical order]
\label{lem:V83-weak-Green}
For \(\epsilon,\zeta\in\mathfrak G_{83}^{\delta}\),
\begin{equation}
 \left\langle
 \mathscr P^{ab}(e_n),D_a\zeta_b+D_b\zeta_a
 \right\rangle_{\Gamma}
 =
 2\langle R_J(\epsilon),T_J(\zeta)\rangle_
 {-1+\delta,\,1-\delta},
 \label{eq:V83-weak-Green}
\end{equation}
with the oriented sum over joint components.
\end{lemma}

\begin{proof}
The left side is a continuous
\(H^{-1/2+\delta},H^{1/2-\delta}\) duality because
\(D\zeta\in H^{1/2+\delta}\hookrightarrow H^{1/2-\delta}\).
The right side is continuous by
\eqref{eq:V83-T-order}--\eqref{eq:V83-R-order} and the downcasting above.
For smooth jets, symmetry of \(\mathscr P\), its divergence identity, and
integration by parts give the formula.  Smooth jets are dense in
\(\mathfrak G_{83}^{\delta}\), so the identity extends.
\end{proof}

Put
\begin{equation}
 \mathcal P_{83}^{\delta}
 :=
 \mathcal M_{83}^{\delta}\oplus\mathcal C_{83}^{\delta},
 \qquad
 \omega_{83}^{\delta}
 :=
 \omega_{M,83}^{\delta}\oplus\omega_{C,83}^{\delta},
 \label{eq:V83-product-phase}
\end{equation}
and define
\begin{equation}
 V_{83}^{\delta}\epsilon
 :=
 \bigl((G_\Gamma\epsilon,P_\Gamma\epsilon),
       (T_J\epsilon,-R_J\epsilon)\bigr).
 \label{eq:V83-combined-action}
\end{equation}

\begin{theorem}[Supercritical closed residual/corner reduction]
\label{thm:V83-supercritical-reduction}
For every \(\delta\) satisfying \eqref{eq:V83-delta-range}:
\begin{enumerate}[label=\textup{(\roman*)},leftmargin=2.8em]
\item \(\omega_{83}^{\delta}\) is a continuous strongly nondegenerate
      Hilbert symplectic form;
\item \(V_{83}^{\delta}:\mathfrak G_{83}^{\delta}\to
      \mathcal P_{83}^{\delta}\) is bounded and has isotropic range;
\item with
      \[
       W_{83}^{\delta}
       :=\overline{\operatorname{ran}V_{83}^{\delta}}^{
                    \mathcal P_{83}^{\delta}},
       \qquad
       \mathcal N_{83}^{\delta}
       :=(W_{83}^{\delta})^{\perp_{\omega_{83}^{\delta}}},
      \]
      the quotient
      \begin{equation}
       (\mathcal P_{83}^{\delta})^{\rm red}
       :=\mathcal N_{83}^{\delta}/W_{83}^{\delta}
       \label{eq:V83-reduced-phase}
      \end{equation}
      is a Hilbert space with a continuous nondegenerate reduced form;
\item the V75 residual plus matched nonminimal charge extends continuously
      to this phase and obeys the off-shell master identity.
\end{enumerate}
\end{theorem}

\begin{proof}
Each phase is \(H^s\oplus(H^s)^*\) with its canonical form, so the Riesz
isomorphism proves (i).  The mapping orders
\eqref{eq:V83-G-order}--\eqref{eq:V83-R-order} and the two continuous
embeddings prove boundedness.

Apply Lemma~\ref{lem:V83-weak-Green} to \((\epsilon,\zeta)\) and
\((\zeta,\epsilon)\).  With the factor \((2\kappa)^{-1}\), it gives
\begin{equation}
 \omega_{M,83}^{\delta}
 \bigl((G\epsilon,P\epsilon),(G\zeta,P\zeta)\bigr)
 =
 \frac1\kappa
 \left(
  \langle R\epsilon,T\zeta\rangle
  -
  \langle R\zeta,T\epsilon\rangle
 \right).
 \label{eq:V83-metric-cocycle}
\end{equation}
Substitution of \((T,-R)\) into
\eqref{eq:V83-corner-form} gives the negative of
\eqref{eq:V83-metric-cocycle}.  Thus the range is isotropic.

The closed-orbit argument of
Theorem~\ref{thm:V82-closed-reduction} now applies verbatim: continuity
makes \(W_{83}^{\delta}\) closed and isotropic, the common moment zero level
is its symplectic orthogonal, and strong nondegeneracy gives the double
orthogonal identity.  This proves (iii).  The residual moment is the
continuous linear Hamiltonian
\[
 \widetilde\mu_{83,\epsilon}^{\delta}(Z)
 :=\omega_{83}^{\delta}(V_{83}^{\delta}\epsilon,Z).
\]
Its bracket with the moment for \(\zeta\) is the vanished combined cocycle.
The V75 nonminimal charge is the same square-zero cotangent lift and has
zero cross bracket.  This proves (iv).
\end{proof}

\begin{assessment}[What the supercritical theorem settles]
\label{ass:V83-supercritical-meaning}
Theorem~\ref{thm:V83-supercritical-reduction} lowers the sufficient V82 face
value order from \(H^2\) to every \(H^{3/2+\delta}\), while retaining a
closed strong phase and the local cocycle cancellation.  It does not assert
the endpoint \(\delta=0\), uniformity of constants as
\(\delta\downarrow0\), or equality of the reduced phase with the V59 energy
completion.
\end{assessment}

\subsection{Exact failure of the explicit V78 dressing section}
\label{subsec:V83-dressing-section}

On the V81 TT fibre, the V78 scalar dressing coordinate has the form
\begin{equation}
 \widehat X_0(q)
 =
 \frac1{\sqrt{2\pi}}
 \int_{\mathbb R}
 \frac{\widehat v(q,\rho)}
      {2\sqrt2\,(r^2+\rho^2)}\,d\rho,
 \qquad r=|q|>0.
 \label{eq:V83-X0-functional}
\end{equation}

\begin{proposition}[Exact fibre norm of the dressing section]
\label{prop:V83-X0-unbounded}
As a functional on \(L^2(\mathbb R_\rho)\), the squared norm of
\eqref{eq:V83-X0-functional} is
\begin{equation}
 \|\widehat X_0(q)\|_{\mathrm{op}}^2
 =\frac1{32r^3}.
 \label{eq:V83-X0-fibre-norm}
\end{equation}
Consequently there are smooth V59 energy-unit packets supported where
\(r\asymp\varepsilon\) for which
\begin{equation}
 \|X_0\|_{L^2(J)}
 \ge c\varepsilon^{-3/2}.
 \label{eq:V83-X0-packet-growth}
\end{equation}
The explicit V78 section is therefore unbounded from the V59 energy phase
to \(H^{1-\delta}(J)\) for every
\(0<\delta<1/2\).
\end{proposition}

\begin{proof}
Plancherel and
\[
 \int_{\mathbb R}\frac{d\rho}{(r^2+\rho^2)^2}
 =\frac{\pi}{2r^3}
\]
give
\[
 \frac1{2\pi}\frac18
 \int_{\mathbb R}\frac{d\rho}{(r^2+\rho^2)^2}
 =\frac1{32r^3},
\]
proving \eqref{eq:V83-X0-fibre-norm}.  Choose the normal-frequency profile
proportional to the Riesz representer in
\eqref{eq:V83-X0-functional}, and choose a tangential profile of unit
\(L^2_q\) norm supported in
\(\{\varepsilon\le r\le2\varepsilon\}\).  After the V59 energy
normalization used in V81, the output multiplier is bounded below by
\(c\varepsilon^{-3/2}\) throughout that annulus.  This proves
\eqref{eq:V83-X0-packet-growth}.  The inhomogeneous
\(H^{1-\delta}\) norm dominates the \(L^2\) norm, proving the final claim.
\end{proof}

\begin{firewall}[A bad section does not determine the quotient norm]
\label{fire:V83-section-versus-quotient}
Proposition~\ref{prop:V83-X0-unbounded} concerns one chosen representative
of a residual equivalence class.  The norm in the reduced Hilbert phase is
the infimum over all representatives.  An unbounded section can coexist
with a bounded quotient map if a different residual representative cancels
its large low-frequency component.  Thus the proposition rules out the V78
section as a comparison tool; it does not rule out the V59-to-reduction
map.
\end{firewall}

\subsection{The exact energy-to-reduction gate}
\label{subsec:V83-energy-gate}

Let \(\mathscr D_{59}^{\rm sm}\) denote the smooth compact-core physical
data used in V59, modulo the already identified smooth gauge kernel, and let
\(\|\cdot\|_{\mathcal E_{59}}\) be its physical energy norm.  The V63/V76
smooth reduction assigns to \(u\in\mathscr D_{59}^{\rm sm}\) a smooth
representative
\[
 Z_u^\delta\in\mathcal N_{83}^{\delta}.
\]
Changing that representative changes it by
\(\operatorname{ran}V_{83}^{\delta}\), so define
\begin{equation}
 S_{83}^{\delta,\rm sm}u
 :=[Z_u^\delta]\in(\mathcal P_{83}^{\delta})^{\rm red}.
 \label{eq:V83-smooth-comparison-map}
\end{equation}

\begin{theorem}[Necessary and sufficient quotient estimate]
\label{thm:V83-energy-gate}
For fixed \(0<\delta<1/2\), the map
\eqref{eq:V83-smooth-comparison-map} extends uniquely to a bounded map
\[
 S_{83}^{\delta}:\mathcal E_{59}
 \longrightarrow(\mathcal P_{83}^{\delta})^{\rm red}
\]
if and only if there is a constant \(C_\delta<\infty\) such that
\begin{equation}
 \boxed{
 \inf_{w\in W_{83}^{\delta}}
 \|Z_u^\delta+w\|_{\mathcal P_{83}^{\delta}}
 \le C_\delta\|u\|_{\mathcal E_{59}}
 \quad
 (u\in\mathscr D_{59}^{\rm sm}).}
 \label{eq:V83-quotient-gate}
\end{equation}
If, in addition, there is \(c_\delta>0\) such that
\begin{equation}
 c_\delta\|u\|_{\mathcal E_{59}}
 \le
 \inf_{w\in W_{83}^{\delta}}
 \|Z_u^\delta+w\|_{\mathcal P_{83}^{\delta}}
 \label{eq:V83-reverse-gate}
\end{equation}
on the smooth core and the range of
\(S_{83}^{\delta,\rm sm}\) is dense, then
\(S_{83}^{\delta}\) is a topological isomorphism onto the full reduced
phase.
\end{theorem}

\begin{proof}
Because \(W_{83}^{\delta}\) is closed, the quotient norm is exactly
\[
 \|[Z_u^\delta]\|_{(\mathcal P_{83}^{\delta})^{\rm red}}
 =
 \inf_{w\in W_{83}^{\delta}}
 \|Z_u^\delta+w\|_{\mathcal P_{83}^{\delta}}.
\]
Thus \eqref{eq:V83-quotient-gate} is precisely boundedness of
\(S_{83}^{\delta,\rm sm}\) on the dense subspace
\(\mathscr D_{59}^{\rm sm}\).  The standard completion theorem gives the
unique bounded extension, and restriction of any bounded extension gives
the estimate.  This proves the equivalence.

The reverse estimate passes to the completion and makes the extended map
bounded below.  Its range is then closed.  Density of the smooth range makes
that closed range the whole target, and the inverse is bounded by
\(c_\delta^{-1}\).
\end{proof}

\begin{corollary}[Symplectic identification once the two-sided gate holds]
\label{cor:V83-symplectic-gate}
Suppose the hypotheses giving a topological isomorphism in
Theorem~\ref{thm:V83-energy-gate} hold.  If the V63/V76 smooth comparison
obeys
\[
 (S_{83}^{\delta,\rm sm})^*
 \omega_{83}^{\delta,\rm red}
 =\omega_{59}
\]
on \(\mathscr D_{59}^{\rm sm}\), then
\(S_{83}^{\delta}\) is a symplectic Hilbert isomorphism.
\end{corollary}

\begin{proof}
Both forms and the extended map are continuous.  Equality on the dense
smooth core therefore passes to the completion.
\end{proof}

\begin{assessment}[Why \eqref{eq:V83-quotient-gate} is the bottleneck]
\label{ass:V83-gate-bottleneck}
The V82 construction proves that the quotient on the left of
\eqref{eq:V83-quotient-gate} exists.  V81 and
Proposition~\ref{prop:V83-X0-unbounded} prove that two natural separated
representatives have divergent infrared norms.  Neither result computes
the infimum over the closed residual orbit.  The next calculation must
therefore solve the minimum-norm residual normal equation itself, rather
than choose another explicit dressing formula.
\end{assessment}

\subsection{V83 result ledger and next acquisition}
\label{subsec:V83-ledger}

\paragraph{New exact conclusions.}
\begin{enumerate}[label=\textup{(\arabic*)},leftmargin=2.8em]
\item The unique formal undiminished duality scale is
      \(H^{3/2}\oplus H^{1/2}\) for residual face jets,
      \((H^{1/2}\oplus H^{-1/2})_\Gamma\) for metric data, and
      \((H^1\oplus H^{-1})_J\) for corner data.
\item The required normal-Hessian map at that scale is an endpoint trace;
      V82 does not prove it.
\item For every \(0<\delta<1/2\), the fractional normal trace is bounded
      from divergence-free \(H^{-1/2+\delta}\) tensors to
      \(H^{-1+\delta}\) joint data.
\item The downcast phases
      \[
       (H^{1/2-\delta}\oplus H^{-1/2+\delta})_\Gamma
       \oplus
       (H^{1-\delta}\oplus H^{-1+\delta})_J
      \]
      are strong Hilbert symplectic phases carrying the continuous
      unsplit V72--V75 cancellation and closed residual reduction.
\item The V78 \(X_0\) dressing functional has exact squared fibre norm
      \(1/(32r^3)\) and is unbounded on V59 infrared energy packets.
\item Bounded extension from the V59 energy completion is equivalent
      exactly to the quotient estimate \eqref{eq:V83-quotient-gate};
      the reverse estimate and dense range give a full topological
      identification.
\end{enumerate}

\paragraph{Next forced step.}
V84 must calculate the infimum in \eqref{eq:V83-quotient-gate} on the flat
Fourier fibres.  For each \(r>0\), it must write the normal equation for the
minimum-norm residual correction in
\(\mathcal P_{83}^{\delta}\), compute its least singular value on the V81
infrared packet, and distinguish a true quotient divergence from the
unbounded V78 section.  The calculation should also track
\(C_\delta\) as \(\delta\downarrow0\); a uniform estimate would support an
endpoint completion, while divergence would locate the missing infrared
weight.

\begin{assessment}[V83 continuum status]
\label{ass:V83-continuum-status}
The full Standard--Model--Einstein continuum remains unproved.  V83 proves
a strict supercritical family of closed unsplit residual/corner phases and
reduces their comparison with the V59 physical phase to the exact
quotient-norm gate \eqref{eq:V83-quotient-gate}.  The critical endpoint,
the quotient estimate, nonlinear curved gluing, the Brown--York/Bianchi
stress identity, the common renormalized Standard Model stress domain, and
the cofinal interacting continuum remain open.
\end{assessment}

\section*{V83 exact checks and append-only record}
\addcontentsline{toc}{section}{V83 exact checks and append-only record}

\begin{equation}
\boxed{
\begin{gathered}
 0<\delta<\tfrac12,\qquad
 \mathfrak G_{83}^{\delta}
 =H^{3/2+\delta}\oplus H^{1/2+\delta},\\
 \mathcal P_{83}^{\delta}
 =
 (H^{1/2-\delta}\oplus H^{-1/2+\delta})_\Gamma
 \oplus
 (H^{1-\delta}\oplus H^{-1+\delta})_J,\\
 \omega_{83}^{\delta}
 (V_{83}^{\delta}\epsilon,V_{83}^{\delta}\zeta)=0,\\
 \|\widehat X_0(q)\|_{\mathrm{op}}^2=\frac1{32|q|^3},\\
 \|S_{83}^{\delta,\rm sm}u\|_{\rm quotient}
 =
 \inf_{w\in W_{83}^{\delta}}
 \|Z_u^\delta+w\|_{\mathcal P_{83}^{\delta}}.
\end{gathered}}
\label{eq:V83-check-card}
\end{equation}

The V82 predecessor used for this append has SHA--256
\begin{center}
\texttt{1267f9b643de50c2b15ef33e7912a78f43a50a216e726ac2e4116bba0d6ccdf9}.
\end{center}
The V83 candidate retains its bytes up to the sole live final document
terminator, appends this material, and restores exactly one active
\verb|\end{document}|.  Validation includes balanced environments and
braces, unique labels, resolved new internal references, valid inline and
display math delimiters, exact predecessor-prefix equality, and a final
inventory which removes only the superseded SM V82 file.

% =============================================================================
% END APPENDED VERSION V83 MATERIAL
% =============================================================================

\clearpage
% =============================================================================
% BEGIN APPENDED VERSION V84 MATERIAL
% =============================================================================
\section{V84: the orbit-invariant cap difference, failure of the
inhomogeneous quotient gate, and the forced homogeneous repair}
\label{sec:V84-cap-invariant}

V83 reduced comparison with the V59 energy phase to the quotient estimate
\eqref{eq:V83-quotient-gate}.  An exact least-squares representative cannot
be found merely by solving the V78 dressing equations: the residual range
need not be closed, and an unbounded section need not imply an unbounded
quotient.  This note first states the correct Hilbert projection problem.

The decisive step is then dual rather than primal.  On a flat side between
two temporal joints there is a bounded linear cap-difference observable
which annihilates the entire closed residual orbit.  Its value on the
physical plus lift is controlled by the common V59 interface trace.  The
trace has an exact \(r^{-1/2}\) energy norm in the coupled
metric--continuum form.  Scaling therefore produces energy-unit data whose
reduced inhomogeneous norm diverges like \(r^{-1/2}\).

This proves that every V83 inhomogeneous supercritical comparison map fails
to extend to the full homogeneous V59 energy phase.  It also identifies the
repair: the joint configuration topology must vanish at least as
\(r^{1/2}\) at tangential zero frequency.  The homogeneous Sobolev orders
already suggested by V83 have this property; their zero mode must be
adjoined through the separate V80 chart.

\subsection{The closed-orbit least-squares problem}
\label{subsec:V84-projection}

Let \(\mathcal H\) be a Hilbert space, let
\(V:\mathcal G\to\mathcal H\) be bounded, and put
\(W=\overline{\operatorname{ran}V}\).  Write \(P_W\) for the orthogonal
projection defined by the positive Hilbert inner product, independently of
any symplectic form on \(\mathcal H\).

\begin{lemma}[Exact quotient minimizer and regularized normal equation]
\label{lem:V84-projection}
For every \(Z\in\mathcal H\),
\begin{align}
 \inf_{w\in W}\|Z+w\|_{\mathcal H}
 &=\|(I-P_W)Z\|_{\mathcal H},
 \label{eq:V84-projection-distance}\\
 w_*&=-P_WZ
 \label{eq:V84-projection-minimizer}
\end{align}
is the unique minimizing orbit vector.  A minimizing residual parameter
\(\epsilon_*\) with \(w_*=V\epsilon_*\) exists if and only if
\(P_WZ\in\operatorname{ran}V\).  When it exists, it obeys
\begin{equation}
 V^*(Z+V\epsilon_*)=0.
 \label{eq:V84-normal-equation}
\end{equation}
Without a closed-range hypothesis, the regularized equations
\begin{equation}
 (V^*V+\eta I)\epsilon_\eta=-V^*Z,
 \qquad \eta>0,
 \label{eq:V84-regularized-normal}
\end{equation}
satisfy
\begin{equation}
 V\epsilon_\eta\longrightarrow-P_WZ
 \quad\hbox{as }\eta\downarrow0.
 \label{eq:V84-regularized-limit}
\end{equation}
\end{lemma}

\begin{proof}
The orthogonal decomposition
\[
 Z+w=(I-P_W)Z+(P_WZ+w)
\]
has its two terms in \(W^\perp\) and \(W\).  Pythagoras proves
\eqref{eq:V84-projection-distance} and
\eqref{eq:V84-projection-minimizer}.  The second assertion is then
immediate, and differentiation of
\(\|Z+V\epsilon\|^2\) gives
\eqref{eq:V84-normal-equation}.

For the last assertion, the polar decomposition reduces the claim to the
spectral theorem for the positive operator \(V^*V\).  Equivalently,
\[
 V\epsilon_\eta
 =-V(V^*V+\eta I)^{-1}V^*Z
 \longrightarrow-P_{\overline{\operatorname{ran}V}}Z
\]
strongly, because the scalar multiplier
\(\lambda/(\lambda+\eta)\) tends to \(1\) on the positive spectrum and is
bounded by \(1\).
\end{proof}

\begin{assessment}[Why a fibre pseudoinverse was not assumed]
\label{ass:V84-no-pseudoinverse}
For \(V=V_{83}^{\delta}\), V82--V83 prove boundedness but not closed range.
Thus the quotient minimizer always exists in the closed orbit
\(W_{83}^{\delta}\), while an exact residual parameter realizing it may
fail to exist.  A formula using
\((V^*V)^{-1}V^*\) would silently assume the missing closed-range estimate.
The annihilator argument below bounds the quotient itself and avoids that
assumption.
\end{assessment}

\subsection{A bounded observable annihilating the residual orbit}
\label{subsec:V84-invariant}

Take the flat timelike side
\[
 \Gamma_T=[0,T]\times\mathbb R_y^2,\qquad T>0,
\]
with temporal joints \(J_0\) and \(J_T\).  Write
\(X_0^{(0)},X_0^{(T)}\) for the time component of the V74 corner coordinate
at the two joints and \(H_{00}\) for the time--time component of the
induced metric variable on the side.  Define
\begin{equation}
 \mathcal I_T(H,p,X,\Pi)
 :=
 X_0^{(T)}-X_0^{(0)}
 -\frac12\int_0^T H_{00}(t,\cdot)\,\dd t .
 \label{eq:V84-cap-invariant}
\end{equation}

\begin{lemma}[Bounded residual-invariant cap difference]
\label{lem:V84-cap-invariant}
For each \(0<\delta<1/2\), the map
\begin{equation}
 \mathcal I_T:\mathcal P_{83}^{\delta}
 \longrightarrow H^{1/2-\delta}(\mathbb R^2_y)
 \label{eq:V84-invariant-map}
\end{equation}
is bounded and
\begin{equation}
 \mathcal I_TV_{83}^{\delta}\epsilon=0
 \quad
 (\epsilon\in\mathfrak G_{83}^{\delta}).
 \label{eq:V84-invariant-annihilates}
\end{equation}
Consequently it vanishes on \(W_{83}^{\delta}\), and every
\(Z\in\mathcal N_{83}^{\delta}\) obeys
\begin{equation}
 \|[Z]\|_{(\mathcal P_{83}^{\delta})^{\rm red}}
 \ge
 C_{\delta,T}^{-1}
 \|\mathcal I_TZ\|_{H^{1/2-\delta}}.
 \label{eq:V84-invariant-lower-bound}
\end{equation}
\end{lemma}

\begin{proof}
The two corner evaluations in \eqref{eq:V84-cap-invariant} lie in
\(H^{1-\delta}(\mathbb R^2)\), which embeds continuously in
\(H^{1/2-\delta}(\mathbb R^2)\).  Time integration is bounded from
\(H^{1/2-\delta}(\Gamma_T)\) to
\(H^{1/2-\delta}(\mathbb R^2)\).  To see the latter statement, extend in
time, Fourier transform, pair the temporal transform with the
\(L^2\) transform of the indicator of \([0,T]\), and use
\(\langle q\rangle^{1/2-\delta}
\le\langle(\tau,q)\rangle^{1/2-\delta}\).
This proves boundedness.

Under a residual translation,
\[
 \delta_\epsilon X_0^{(j)}=\epsilon_0|_{J_j},
 \qquad
 \delta_\epsilon H_{00}=2D_0\epsilon_0.
\]
The fundamental theorem of calculus gives
\[
 \delta_\epsilon\mathcal I_T
 =
 \epsilon_0(T)-\epsilon_0(0)
 -\int_0^TD_0\epsilon_0\,\dd t=0.
\]
Density proves the identity at the stated Sobolev orders.  Continuity then
extends it from \(\operatorname{ran}V_{83}^{\delta}\) to its closure.
For any \(w\in W_{83}^{\delta}\),
\[
 \|\mathcal I_TZ\|
 =\|\mathcal I_T(Z+w)\|
 \le C_{\delta,T}\|Z+w\|_{\mathcal P_{83}^{\delta}}.
\]
Taking the infimum proves \eqref{eq:V84-invariant-lower-bound}.
\end{proof}

\subsection{Exact V59 norm of the common interface trace}
\label{subsec:V84-interface-trace}

At tangential frequency \(q\), put \(r=|q|>0\).  The V59 position form for
one radiative polarization is
\begin{align}
 \mathfrak q_r[a,b]
 &:=
 \int_0^\infty
 \left(|\partial_xa|^2+r^2|a|^2\right)\dd x
 \notag\\
 &\quad+
 m\int_0^\infty
 \left(|\partial_\rho b|^2+v^2r^2|b|^2\right)\dd\rho,
 \qquad a(0)=b(0).
 \label{eq:V84-V59-fibre-form}
\end{align}

\begin{proposition}[Sharp common-trace inequality]
\label{prop:V84-sharp-trace}
For every pair in the V59 fibre form domain,
\begin{equation}
 |a(0)|^2=|b(0)|^2
 \le\frac1{r(1+mv)}\mathfrak q_r[a,b].
 \label{eq:V84-sharp-trace-bound}
\end{equation}
The constant is sharp, and
\begin{equation}
 \inf_{\substack{a(0)=b(0)=z}}
 \mathfrak q_r[a,b]
 =
 r(1+mv)|z|^2.
 \label{eq:V84-trace-capacity}
\end{equation}
The unique minimizer is
\begin{equation}
 a(x)=ze^{-rx},
 \qquad
 b(\rho)=ze^{-vr\rho}.
 \label{eq:V84-trace-minimizer}
\end{equation}
\end{proposition}

\begin{proof}
For \(a(\infty)=0\) in the form sense,
\[
 |a(0)|^2
 =-2\operatorname{Re}\int_0^\infty
   \overline a\,\partial_xa\,\dd x
 \le
 r^{-1}\int_0^\infty|\partial_xa|^2\,\dd x
 +r\int_0^\infty|a|^2\,\dd x.
\]
Multiplication by \(r\) gives a lower bound \(r|a(0)|^2\) for the first
integral in \eqref{eq:V84-V59-fibre-form}.  Applying the same argument to
\(b\) with \(vr\), and then multiplying by \(m\), gives
\(mvr|b(0)|^2\).  The common-trace condition proves
\eqref{eq:V84-sharp-trace-bound}.

Both scalar inequalities are equalities exactly for the decaying
exponentials in \eqref{eq:V84-trace-minimizer}.  Direct substitution gives
\(r|z|^2+mvr|z|^2\), proving
\eqref{eq:V84-trace-capacity} and uniqueness.
\end{proof}

\subsection{The physical cap difference on the plus graph}
\label{subsec:V84-physical-difference}

Let \(s=\varsigma\) be the plus trace of V78 and let \(a\) be the V59
radiative plus amplitude.  Equations
\eqref{eq:V78-X0-dressing} and
\eqref{eq:V79-plus-amplitude} imply a local evolution identity which does
not use the separated inverse as an estimate.

\begin{lemma}[Derivative of the dressed time coordinate]
\label{lem:V84-X0-derivative}
On every nonzero tangential Fourier fibre of the physical plus graph,
\begin{equation}
 D_0X_0=-\frac1{2\sqrt2}a|_{x=0}.
 \label{eq:V84-X0-derivative}
\end{equation}
If the initial radiative velocity vanishes, then \(X_0^{(0)}=0\), and
\begin{equation}
 X_0^{(T)}-X_0^{(0)}
 =
 -\frac1{2\sqrt2}\int_0^Ta(t,0)\,\dd t.
 \label{eq:V84-X0-time-integral}
\end{equation}
\end{lemma}

\begin{proof}
At frequency \(r\), V78 gives
\[
 X_0=\frac12r^{-2}D_0s.
\]
The V79 physical relation is
\[
 r^2a=\sqrt2\,(r^2-D_n^2)s.
\]
Because \(s\) obeys the radiative wave equation,
\[
 D_0^2s=D_n^2s-r^2s=-\frac{r^2}{\sqrt2}a.
\]
Differentiating \(X_0\) proves
\eqref{eq:V84-X0-derivative}.  The multiplier relation commutes with
\(D_0\), so zero initial velocity gives \(D_0s=0\) and hence
\(X_0^{(0)}=0\).  Time integration gives
\eqref{eq:V84-X0-time-integral}.
\end{proof}

The V59 fibre operators have an exact dilation.  Let \(A_r\) be the
nonnegative self-adjoint operator represented by
\(\mathfrak q_r\), including the common-trace interface.  If
\[
 (U_rf)(X):=r^{-1/2}f(X/r)
\]
is applied to both half-line components, then
\begin{equation}
 U_rA_rU_r^{-1}=r^2A_1.
 \label{eq:V84-operator-dilation}
\end{equation}
This follows directly by changing variables in
\eqref{eq:V84-V59-fibre-form}; the common-trace domain is preserved.

\begin{theorem}[A true \(r^{-1/2}\) quotient obstruction]
\label{thm:V84-quotient-obstruction}
Fix \(T>0\) and \(0<\delta<1/2\).  There are smooth physical plus packets
\(u_\varepsilon\), supported in
\(\{\varepsilon<|q|<2\varepsilon\}\), with
\begin{equation}
 \|u_\varepsilon\|_{\mathcal E_{59}}=1
 \label{eq:V84-energy-unit}
\end{equation}
and zero initial velocity, for which every V83 zero-level representative
\(Z_\varepsilon^\delta\) satisfies
\begin{equation}
 \|[Z_\varepsilon^\delta]\|_
 {(\mathcal P_{83}^{\delta})^{\rm red}}
 \ge c_{\delta,T,m,v}\varepsilon^{-1/2}.
 \label{eq:V84-quotient-growth}
\end{equation}
Consequently the upper estimate
\eqref{eq:V83-quotient-gate} fails for every fixed
\(0<\delta<1/2\).
\end{theorem}

\begin{proof}
At unit tangential frequency take the sharp trace minimizer
\[
 u_1:=
 \frac1{\sqrt{1+mv}}
 \left(e^{-x},e^{-v\rho}\right).
\]
Proposition~\ref{prop:V84-sharp-trace} gives
\(\mathfrak q_1[u_1]=1\) and common trace
\((1+mv)^{-1/2}\).  Define the position datum at frequency \(r\) by
\begin{equation}
 u_r(x,\rho)
 :=
 r^{-1/2}u_1(rx,r\rho).
 \label{eq:V84-scaled-position}
\end{equation}
Then \(\mathfrak q_r[u_r]=1\), while its common trace is
\[
 r^{-1/2}(1+mv)^{-1/2}.
\]

Let the V59 wave start from \(u_r\) with zero velocity.  By
\eqref{eq:V84-operator-dilation}, its metric boundary trace has the form
\begin{equation}
 a_r(t,0)=r^{-1/2}F(rt),
 \qquad
 F(s):=
 \left[\cos(s\sqrt{A_1})u_1\right]_a(0).
 \label{eq:V84-scaled-boundary-evolution}
\end{equation}
The wave group is strongly continuous in the form norm, and boundary
evaluation is continuous in that norm by
Proposition~\ref{prop:V84-sharp-trace}.  Hence \(F\) is continuous and
\[
 F(0)=\frac1{\sqrt{1+mv}}.
\]
There is \(r_0>0\) such that for \(0<r<r_0\) and \(0\le t\le T\),
\[
 \operatorname{Re}F(rt)
 \ge\frac1{2\sqrt{1+mv}}.
\]
Lemma~\ref{lem:V84-X0-derivative} therefore gives
\begin{equation}
 |X_0^{(T)}-X_0^{(0)}|
 \ge
 \frac{T}{4\sqrt{2(1+mv)}}r^{-1/2}.
 \label{eq:V84-cap-growth-fibre}
\end{equation}

Choose a smooth tangential function \(\chi_\varepsilon\) supported in the
stated annulus with \(L^2_q\) norm one and use
\(\chi_\varepsilon(q)u_{|q|}\) as the fibre datum.  Its total V59 position
energy is one.  Smooth form-core approximation, followed by the unitary
V59 evolution, preserves the lower bound up to a fixed factor.  The
V63 plus lift has \(H_{00}=0\), so
\[
 \mathcal I_TZ_\varepsilon^\delta
 =X_0^{(T)}-X_0^{(0)}.
\]
On the shrinking annulus the inhomogeneous
\(H^{1/2-\delta}\) norm is uniformly equivalent to \(L^2_q\).
Equations \eqref{eq:V84-invariant-lower-bound} and
\eqref{eq:V84-cap-growth-fibre} prove
\eqref{eq:V84-quotient-growth}.  The right side of
\eqref{eq:V83-quotient-gate} is one, so that estimate cannot hold.
\end{proof}

\begin{firewall}[The obstruction uses the coupled V59 interface]
\label{fire:V84-coupled-interface}
Theorem~\ref{thm:V84-quotient-obstruction} is not a free doubled-collar
surrogate.  Its initial datum belongs to the common-trace form domain of the
coupled V59 metric and continuum half-lines; the factor \(1+mv\) is the
exact interface capacity.  The subsequent evolution is the V59
self-adjoint wave group and uses its exact dilation
\eqref{eq:V84-operator-dilation}.
\end{firewall}

\subsection{The infrared weight forced by the obstruction}
\label{subsec:V84-forced-weight}

The preceding proof depends on the low-frequency behavior of the norm in
which the invariant \(\mathcal I_T\) is measured.  Let a proposed spatial
joint configuration norm have radial Fourier multiplier
\(\vartheta(r)\) at \(r\ll1\).

\begin{corollary}[Necessary low-frequency weight]
\label{cor:V84-necessary-weight}
If the cap-difference map is bounded from a proposed reduced phase to the
norm with multiplier \(\vartheta(r)\), and the full V59 energy phase maps
boundedly into that reduced phase, then necessarily
\begin{equation}
 \sup_{0<r<r_0}\frac{\vartheta(r)}{\sqrt r}<\infty.
 \label{eq:V84-weight-criterion}
\end{equation}
For a homogeneous Sobolev norm
\(\dot H^\beta(\mathbb R^2)\), this requires
\begin{equation}
 \beta\ge\frac12.
 \label{eq:V84-homogeneous-threshold}
\end{equation}
\end{corollary}

\begin{proof}
Apply the two assumed bounded maps to the scaled energy-unit data in
Theorem~\ref{thm:V84-quotient-obstruction}.  Their cap difference has
unweighted size bounded below by \(cr^{-1/2}\), so its proposed norm is
bounded below by \(c\vartheta(r)r^{-1/2}\).  Uniform boundedness gives
\eqref{eq:V84-weight-criterion}.  For
\(\vartheta(r)=r^\beta\), the quotient is bounded near zero exactly when
\(\beta\ge1/2\).
\end{proof}

\begin{assessment}[The forced upstream change]
\label{ass:V84-upstream-change}
V83 used inhomogeneous joint configuration spaces, whose multiplier tends
to one as \(r\downarrow0\); Theorem
\ref{thm:V84-quotient-obstruction} therefore rules them out for the full
V59 homogeneous energy phase.  The order \(H^{1-\delta}\) was not itself
too weak.  Its homogeneous version has multiplier \(r^{1-\delta}\), which
satisfies \eqref{eq:V84-weight-criterion} for every
\(0<\delta<1/2\).  The next construction must replace the punctured
inhomogeneous boundary topology by a homogeneous one and attach the exact
\(q=0\) sector from V80 as a separate finite-dimensional chart.  It must
then recheck the fractional normal trace and strong cotangent duality in
that hybrid topology.
\end{assessment}

\subsection{V84 result ledger and next acquisition}
\label{subsec:V84-ledger}

\paragraph{New exact conclusions.}
\begin{enumerate}[label=\textup{(\arabic*)},leftmargin=2.8em]
\item The quotient minimizer is the Hilbert projection onto the closed
      residual orbit.  A minimizing residual parameter need not exist
      unless the residual range is closed; Tikhonov normal equations
      converge to the closed-orbit minimizer.
\item The cap difference \(\mathcal I_T\) in
      \eqref{eq:V84-cap-invariant} is bounded and annihilates the full closed
      residual orbit.
\item The coupled V59 common trace has exact squared energy norm
      \(1/[r(1+mv)]\), with the exponential minimizer
      \eqref{eq:V84-trace-minimizer}.
\item On the plus graph,
      \(D_0X_0=-a(0)/(2\sqrt2)\).
\item Exact V59 dilation turns the sharp trace minimizer into energy-unit
      packets whose orbit-invariant cap difference grows as
      \(r^{-1/2}\).
\item Hence the inhomogeneous quotient gate of V83 fails for every
      \(0<\delta<1/2\); this is a quotient obstruction, not merely an
      unbounded dressing section.
\item Any replacement topology must weight the cap-difference coordinate
      by at least \(r^{1/2}\) near zero.  Homogeneous
      \(\dot H^\beta\) meets this necessary condition precisely for
      \(\beta\ge1/2\).
\end{enumerate}

\paragraph{Next forced step.}
V85 is an audit version.  It must inspect the newest Yang--Mills and
Navier--Stokes notes before fixing the repair.  It should then construct the
hybrid phase consisting of the punctured homogeneous cotangent pairs
\[
 (\dot H^{1/2-\delta}\oplus\dot H^{-1/2+\delta})_\Gamma
 \oplus
 (\dot H^{1-\delta}\oplus\dot H^{-1+\delta})_J
\]
together with the exact V80 \(q=0\) chart.  The proof must specify quotient
conventions for homogeneous constants, establish the corresponding
fractional Green trace, and test both the V81 and V84 packets.  Only after
that closed hybrid phase exists should the minimum-norm comparison be
reopened.

\begin{assessment}[V84 continuum status]
\label{ass:V84-continuum-status}
The full Standard--Model--Einstein continuum remains unproved.  V84 proves
that the inhomogeneous V83 reduced phase is strictly too strong in the
infrared for the homogeneous V59 energy completion and determines the
necessary low-frequency weight.  A closed homogeneous-plus-zero-mode
residual phase, its full V59 quotient identification, nonlinear curved
gluing, the Brown--York/Bianchi stress identity, the common renormalized
Standard Model stress domain, and the cofinal interacting continuum remain
open.
\end{assessment}

\section*{V84 exact checks and append-only record}
\addcontentsline{toc}{section}{V84 exact checks and append-only record}

\begin{equation}
\boxed{
\begin{gathered}
 \mathcal I_T
 =X_0^{(T)}-X_0^{(0)}
  -\tfrac12\int_0^TH_{00}\,\dd t,
 \qquad
 \mathcal I_TW_{83}^{\delta}=0,\\
 \inf_{a(0)=b(0)=z}\mathfrak q_r[a,b]
 =r(1+mv)|z|^2,\\
 D_0X_0=-\frac{a(0)}{2\sqrt2},
 \qquad
 \|[Z_\varepsilon^\delta]\|_{\rm red}
 \ge c\varepsilon^{-1/2},\\
 \vartheta(r)=O(r^{1/2})
 \quad\text{is necessary at }r\downarrow0.
\end{gathered}}
\label{eq:V84-check-card}
\end{equation}

The V83 predecessor used for this append has SHA--256
\begin{center}
\texttt{5871582e4f49c665c001fc8fc9c69ad6ebeea5abf8f641ac6cff55eb5e3e01a2}.
\end{center}
The V84 candidate retains its bytes up to the sole live final document
terminator, appends this material, and restores exactly one active
\verb|\end{document}|.  Validation includes balanced environments and
braces, unique labels, resolved new internal references, valid inline and
display math delimiters, exact predecessor-prefix equality, and a final
inventory which removes only the superseded SM V83 file.

% =============================================================================
% END APPENDED VERSION V84 MATERIAL
% =============================================================================

\clearpage
% =============================================================================
% BEGIN APPENDED VERSION V85 MATERIAL
% =============================================================================
\section{V85: companion audit, the special Hessian trace, and a punctured
homogeneous residual phase}
\label{sec:V85-homogeneous-phase}

V84 proved that the inhomogeneous V83 quotient cannot receive the full V59
homogeneous energy phase continuously.  Its orbit-invariant cap difference
grows like \(r^{-1/2}\) on energy-unit packets.  The necessary repair is an
infrared-vanishing boundary topology, but simply replacing every \(H^s\) by
\(\dot H^s\) is not automatically valid: homogeneous spaces at different
orders do not have the global embeddings used in V83, and a generic
fractional normal trace on a finite two-cap side does not split into two
independently controlled homogeneous traces.

This mandatory audit version resolves both issues on the flat linear
Einstein system.  The residual jet is placed in the intersection of the
orders immediately above and below the critical scale.  More importantly,
the normal trace is not treated as that of an arbitrary divergence-free
tensor.  The trace-adjusted Hessian has the exact V77/V78 rows
\[
 R^0=\sigma Q^2f,\qquad R^A=\sigma D_0D^Af,
\]
so ordinary value and first-normal-derivative trace theorems give the
required homogeneous corner momentum directly.  This produces a strong
punctured off-shell Hilbert phase and its closed residual reduction.  The
exact \(q=0\) physical chart remains the separate V80 chart.

\subsection{Mandatory V85 companion audit}
\label{subsec:V85-audit}

Two reads after the concurrently running companion updates gave the stable
snapshots
\begin{center}
\begin{tabular}{p{0.43\textwidth}r r p{0.28\textwidth}}
\toprule
companion file & lines & bytes & SHA--256\\
\midrule
\texttt{Renewal\_NS\_Stability\_Research\_Notes\_V31.tex}
& \(23053\) & \(887537\) &
\texttt{f1121b62238ad5491bb7cf03635ea9934942edf573ed8faaa9ee79e1d38a6696}\\
\texttt{Renewal\_Yang\_Mills\_Mass\_Gap\_Research\_Notes\_V82.tex}
& \(57538\) & \(2040892\) &
\texttt{c86ed98fa5ac62fc6dd898300ee0a623d696a727c7c608cd6984f4444e53f312}\\
\bottomrule
\end{tabular}
\end{center}
Each snapshot had exactly one active final document terminator.  Neither
companion file was modified.

The NS snapshot is unchanged from the V80 audit.  Its terminal V31 result
still says that removing a flat probability mark costs exactly one dyadic
radius moment; heat formation supplies that moment only at the critical
source order.  This supports the V84 conclusion that a missing spectral
weight must appear explicitly in the topology or the estimate.

The current Yang--Mills V82 chapter first uses the V81 exchange-parity
projection and then proves that every affine function of the pair Casimir has
zero even leakage.  Only the best affine remainder, equivalently a
second-divided-difference curvature, remains in that channel.  The relevant
acquisition rule is to remove the exact structural image before estimating
the remainder.  For the present problem, this means retaining the complete
trace-adjusted-Hessian symbol before applying an ambient normal-trace bound.

\begin{firewall}[Audit separation]
\label{fire:V85-audit-separation}
No Navier--Stokes first-moment estimate and no Yang--Mills Racah or parity
identity is used as an Einstein inequality.  The audit changes only the
proof order: expose the exact correlated or symmetry-reduced object before
estimating its components.
\end{firewall}

\subsection{Punctured homogeneous conventions}
\label{subsec:V85-homogeneous-conventions}

On \(\mathbb R^d\), let \(\mathscr S_\bullet(\mathbb R^d;E)\) be the
Schwartz sections whose Fourier transform vanishes in a neighborhood of the
origin.  For real \(s\), define
\begin{equation}
 \|u\|_{\dot H^s(\mathbb R^d;E)}^2
 :=
 \int_{\mathbb R^d}|\xi|^{2s}|\widehat u(\xi)|_E^2\,\dd\xi
 \label{eq:V85-homogeneous-norm}
\end{equation}
and take the completion of \(\mathscr S_\bullet\), with the usual quotient
by zero-seminorm polynomials when required.  Negative orders are the
continuous duals of the displayed positive-order spaces.  On a flat face
or slab, use the restriction norm obtained by the infimum over whole-space
extensions.  On a compact flat joint, the same notation means the spectral
completion of the mean-zero subspace.  The constant joint mode is never
inserted into an inverse power of \(Q=(-\Delta_J)^{1/2}\).

These definitions have three immediate properties used below:
\begin{align}
 D&:\dot H^s\longrightarrow\dot H^{s-1}
 \quad\hbox{is bounded},
 \label{eq:V85-derivative-map}\\
 \gamma_0&:\dot H^s(\Gamma)\longrightarrow
 \dot H^{s-1/2}(J)
 \quad(s>1/2),
 \label{eq:V85-value-trace}\\
 \gamma_1&:\dot H^s(\Gamma)\longrightarrow
 \dot H^{s-3/2}(J)
 \quad(s>3/2)
 \label{eq:V85-normal-derivative-trace}
\end{align}
are bounded on the flat restriction spaces.  The last map is the trace of
the derivative normal to \(J\) within \(\Gamma\).  The strict inequalities
in \eqref{eq:V85-value-trace}--%
\eqref{eq:V85-normal-derivative-trace} are essential.

\subsection{Why the generic two-cap lifting argument fails}
\label{subsec:V85-generic-trace-failure}

The failure below is not an obstruction to the special Einstein map.  It
shows why the generic divergence-normal proof from V83 cannot simply be
copied into independent homogeneous joint spaces.

Fix \(T>0\), \(r>0\), and \(1/2<s<3/2\).  On one tangential Fourier fibre,
equip scalar functions of \(t\in\mathbb R\) with
\begin{equation}
 \|u\|_{\dot H_r^s}^2
 :=
 \frac1{2\pi}\int_{\mathbb R}
 (\tau^2+r^2)^s|\widehat u(\tau)|^2\,\dd\tau.
 \label{eq:V85-parameter-Hs}
\end{equation}
The endpoint difference functional is
\[
 L_-u:=u(T)-u(0).
\]

\begin{lemma}[Exact norm and finite zero-frequency limit of the difference
trace]
\label{lem:V85-difference-trace}
The squared dual norm of \(L_-\) is
\begin{equation}
 J_{s,T}^-(r)
 :=
 \frac1{2\pi}\int_{\mathbb R}
 \frac{|e^{i\tau T}-1|^2}{(\tau^2+r^2)^s}\,\dd\tau.
 \label{eq:V85-difference-functional-norm}
\end{equation}
Moreover,
\begin{equation}
 0<
 J_{s,T}^-(0)
 :=
 \frac1{2\pi}\int_{\mathbb R}
 \frac{|e^{i\tau T}-1|^2}{|\tau|^{2s}}\,\dd\tau
 <\infty,
 \qquad
 J_{s,T}^-(r)\longrightarrow J_{s,T}^-(0).
 \label{eq:V85-difference-limit}
\end{equation}
\end{lemma}

\begin{proof}
Fourier inversion writes
\[
 L_-u
 =
 \frac1{2\pi}\int_{\mathbb R}
 (e^{i\tau T}-1)\widehat u(\tau)\,\dd\tau.
\]
Cauchy--Schwarz in \eqref{eq:V85-parameter-Hs} gives
\eqref{eq:V85-difference-functional-norm}, and equality is attained by the
Riesz representer.  Near \(\tau=0\), the integrand in
\eqref{eq:V85-difference-limit} is comparable to
\(T^2|\tau|^{2-2s}\), which is integrable because \(s<3/2\).  At infinity
it is bounded by \(4|\tau|^{-2s}\), which is integrable because \(s>1/2\).
Dominated convergence proves the limit; positivity is immediate.
\end{proof}

\begin{proposition}[No product-homogeneous right inverse for two independent
caps]
\label{prop:V85-no-product-lifting}
Put \(\beta=s-\tfrac12>0\).  The two-cap trace
\[
 u\longmapsto(u(0),u(T))
\]
has no right inverse bounded from
\(\dot H^\beta(\mathbb R_y^2)\oplus
 \dot H^\beta(\mathbb R_y^2)\) to the homogeneous face space
\(\dot H^s(\mathbb R_t\times\mathbb R_y^2)\).
\end{proposition}

\begin{proof}
Take a tangential function \(f_r\) of \(L^2\) norm one supported where
\(r\le|q|\le2r\), and prescribe the antisymmetric traces
\((f_r,-f_r)\).  Their product boundary norm is comparable to
\(r^\beta\).  Every extension \(u\) has
\(|L_-u|=2\) fibrewise on the normalized profile.  By
Lemma~\ref{lem:V85-difference-trace}, its bulk norm is bounded below by a
positive constant independent of small \(r\).  The ratio of every lifting
norm to the boundary product norm therefore grows at least as
\(r^{-\beta}\).
\end{proof}

\begin{firewall}[The failed lifting is not the Einstein normal trace]
\label{fire:V85-special-trace-needed}
Proposition~\ref{prop:V85-no-product-lifting} forbids deriving two
independent homogeneous normal traces for an arbitrary negative-order
divergence-free tensor by a product-space lifting.  The Einstein tensor
\(\mathscr P(e_n)\) is a trace-adjusted scalar Hessian.  Its joint rows
factor through ordinary value and derivative traces, as proved next.
\end{firewall}

\subsection{The special trace-adjusted-Hessian factorization}
\label{subsec:V85-special-Hessian}

Let \(f=e_n\) be the component of the residual value normal to the timelike
face.  At a temporal joint \(J_j\), write
\(\nu=\sigma_jD_0\) for the outward normal within the face.  The exact flat
formulas already established in V77/V78 are
\begin{align}
 R_j^0(f)&=\sigma_jQ^2(f|_{J_j}),
 \label{eq:V85-R0-factor}\\
 R_j^A(f)&=\sigma_jD^A\bigl((D_0f)|_{J_j}\bigr).
 \label{eq:V85-RA-factor}
\end{align}

\begin{lemma}[Homogeneous boundedness of the special Hessian trace]
\label{lem:V85-special-R-bounded}
For \(0<\delta<1/2\), the map
\begin{equation}
 f\longmapsto R_j(f)
 \quad\hbox{is bounded from}\quad
 \dot H^{3/2+\delta}(\Gamma_T)
 \quad\hbox{to}\quad
 \dot H^{-1+\delta}(J_j;E_J^*).
 \label{eq:V85-special-R-map}
\end{equation}
The estimate holds jointly for the two caps.
\end{lemma}

\begin{proof}
The value trace theorem gives
\[
 f|_{J_j}\in\dot H^{1+\delta}(J_j),
\]
and \(Q^2\) maps this space to
\(\dot H^{-1+\delta}(J_j)\).  Because
\(3/2+\delta>3/2\), the strict normal-derivative trace theorem gives
\[
 (D_0f)|_{J_j}\in\dot H^\delta(J_j).
\]
One joint derivative maps the latter to
\(\dot H^{-1+\delta}(J_j)\).  Equations
\eqref{eq:V85-R0-factor}--\eqref{eq:V85-RA-factor} prove the estimate.
There are only two caps, so their direct-sum estimate has the same order.
\end{proof}

\begin{remark}[The gain comes from the exact image]
\label{rem:V85-special-image}
An arbitrary tensor of order \(-1/2+\delta\) need not have the product
normal trace in Lemma~\ref{lem:V85-special-R-bounded}.  The result uses the
fact that the time row of \(\mathscr P(f)\) contains either two joint
derivatives of the value trace or one joint derivative of the first normal
derivative trace.  This is precisely the structural projection which the
companion audit requires us to take before estimating.
\end{remark}

\subsection{The punctured homogeneous phase}
\label{subsec:V85-punctured-phase}

Fix \(0<\delta<1/2\).  Define the two-sided residual jet space
\begin{align}
 \dot{\mathfrak G}_{85}^{\delta}(\Gamma)
 &:=
 \left(
  \dot H^{3/2-\delta}(\Gamma;T^*M|_\Gamma)
  \cap
  \dot H^{3/2+\delta}(\Gamma;T^*M|_\Gamma)
 \right)
 \notag\\
 &\quad\oplus
 \left(
  \dot H^{1/2-\delta}(\Gamma;T^*M|_\Gamma)
  \cap
  \dot H^{1/2+\delta}(\Gamma;T^*M|_\Gamma)
 \right),
 \label{eq:V85-residual-intersection}
\end{align}
with the sum of the four Hilbert norms.  The first factor contains
\(e=\epsilon|_\Gamma\), the second \(d=D_n\epsilon|_\Gamma\).

The lower orders in \eqref{eq:V85-residual-intersection} place the
configuration traces in
\begin{align}
 G_\Gamma(e,d)&\in
 \dot H^{1/2-\delta}(\Gamma;E_\Gamma),
 \label{eq:V85-G-map}\\
 T_J(e,d)=e|_J&\in
 \dot H^{1-\delta}(J;E_J),
 \label{eq:V85-T-map}
\end{align}
while the upper order of \(e_n\) gives
\begin{align}
 P_\Gamma(e,d)=\mathscr P(e_n)&\in
 \dot H^{-1/2+\delta}(\Gamma;E_\Gamma^*),
 \label{eq:V85-P-map}\\
 R_J(e,d)&\in
 \dot H^{-1+\delta}(J;E_J^*)
 \label{eq:V85-R-map}
\end{align}
by Lemma~\ref{lem:V85-special-R-bounded}.

Define the punctured metric and corner phases
\begin{align}
 \dot{\mathcal M}_{85}^{\delta}
 &:=
 \dot H^{1/2-\delta}(\Gamma;E_\Gamma)
 \oplus
 \dot H^{-1/2+\delta}(\Gamma;E_\Gamma^*),
 \label{eq:V85-metric-phase}\\
 \dot{\mathcal C}_{85}^{\delta}
 &:=
 \bigoplus_{J_j}
 \left[
  \dot H^{1-\delta}(J_j;E_J)
  \oplus
  \dot H^{-1+\delta}(J_j;E_J^*)
 \right].
 \label{eq:V85-corner-phase}
\end{align}
Use the exact duality pairings to define
\begin{align}
 \dot\omega_{M,85}^{\delta}
 ((H_1,p_1),(H_2,p_2))
 &:=
 \frac1{2\kappa}
 \left(
  \langle p_1,H_2\rangle-\langle p_2,H_1\rangle
 \right),
 \label{eq:V85-metric-form}\\
 \dot\omega_{C,85}^{\delta}
 ((X_1,\Pi_1),(X_2,\Pi_2))
 &:=
 \frac1\kappa
 \left(
  \langle\Pi_1,X_2\rangle-\langle\Pi_2,X_1\rangle
 \right).
 \label{eq:V85-corner-form}
\end{align}

\begin{lemma}[Homogeneous weak Green identity on the special image]
\label{lem:V85-homogeneous-Green}
For \(\epsilon,\zeta\in\dot{\mathfrak G}_{85}^{\delta}\),
\begin{equation}
 \left\langle
 \mathscr P^{ab}(e_n),D_a\zeta_b+D_b\zeta_a
 \right\rangle_\Gamma
 =
 2\sum_j
 \langle R_j(\epsilon),T_j(\zeta)\rangle_
 {-1+\delta,\,1-\delta}.
 \label{eq:V85-homogeneous-Green}
\end{equation}
\end{lemma}

\begin{proof}
The left side is a continuous
\(\dot H^{-1/2+\delta},\dot H^{1/2-\delta}\) duality by
\eqref{eq:V85-P-map} and the lower-order part of
\eqref{eq:V85-residual-intersection}.  The right side is continuous by
\eqref{eq:V85-T-map} and \eqref{eq:V85-R-map}.  On the annular Schwartz
core the identity is ordinary integration by parts, symmetry of
\(\mathscr P\), and \(D_a\mathscr P^{ab}=0\).  That core is dense in the
intersection graph norm, so the identity extends.
\end{proof}

Put
\begin{align}
 \dot{\mathcal P}_{85}^{\delta}
 &:=
 \dot{\mathcal M}_{85}^{\delta}
 \oplus\dot{\mathcal C}_{85}^{\delta},
 \label{eq:V85-product-phase}\\
 \dot V_{85}^{\delta}\epsilon
 &:=
 \bigl((G_\Gamma\epsilon,P_\Gamma\epsilon),
       (T_J\epsilon,-R_J\epsilon)\bigr).
 \label{eq:V85-residual-action}
\end{align}

\begin{theorem}[Closed punctured homogeneous reduction]
\label{thm:V85-homogeneous-reduction}
For every \(0<\delta<1/2\):
\begin{enumerate}[label=\textup{(\roman*)},leftmargin=2.8em]
\item the product of
      \eqref{eq:V85-metric-form} and
      \eqref{eq:V85-corner-form} is a continuous strongly nondegenerate
      Hilbert symplectic form on
      \(\dot{\mathcal P}_{85}^{\delta}\);
\item the map \(\dot V_{85}^{\delta}\) is bounded and its range is
      isotropic;
\item with
      \[
       \dot W_{85}^{\delta}
       :=
       \overline{\operatorname{ran}\dot V_{85}^{\delta}},
       \qquad
       \dot{\mathcal N}_{85}^{\delta}
       :=
       (\dot W_{85}^{\delta})^{\perp_{\dot\omega_{85}^{\delta}}},
      \]
      the quotient
      \begin{equation}
       (\dot{\mathcal P}_{85}^{\delta})^{\rm red}
       :=
       \dot{\mathcal N}_{85}^{\delta}/\dot W_{85}^{\delta}
       \label{eq:V85-reduced-phase}
      \end{equation}
      is Hilbert and carries the continuous nondegenerate reduced form;
\item the V75 residual and matched nonminimal charge extends continuously
      and satisfies the off-shell master identity on this punctured phase.
\end{enumerate}
\end{theorem}

\begin{proof}
Each phase in \eqref{eq:V85-metric-phase}--%
\eqref{eq:V85-corner-phase} is \(V\oplus V^*\) with its canonical form,
proving (i).  Equations \eqref{eq:V85-G-map}--%
\eqref{eq:V85-R-map} prove boundedness.  Applying
Lemma~\ref{lem:V85-homogeneous-Green} in both orders gives
\begin{align*}
 &\dot\omega_{M,85}^{\delta}
 \bigl((G\epsilon,P\epsilon),(G\zeta,P\zeta)\bigr)\\
 &\qquad=
 \frac1\kappa\sum_j
 \left(
  \langle R_j\epsilon,T_j\zeta\rangle
  -
  \langle R_j\zeta,T_j\epsilon\rangle
 \right).
\end{align*}
The corner form evaluated on \((T,-R)\) is its negative, proving (ii).
The closed-orbit and double-orthogonal proof of
Theorem~\ref{thm:V83-supercritical-reduction} gives (iii).  The residual
moments are continuous linear Hamiltonians with brackets equal to the
vanished combined cocycle.  Adding the unchanged square-zero V75
nonminimal cotangent lift proves (iv).
\end{proof}

\subsection{The exact zero mode and the two infrared regressions}
\label{subsec:V85-zero-and-regressions}

On the planar direct integral, \(q=0\) is absent as an \(L^2\) fibre and
the total constant configuration is removed by the homogeneous quotient.
On a compact joint, split \(P_0\oplus P_\perp\) before using
\eqref{eq:V85-homogeneous-norm}.  Let
\(\mathcal E_{80}^0\) denote the V80 \(q=0\), nonzero-normal-spectrum V59
energy phase of the two amplitudes \(a_+,a_\times\).

\begin{proposition}[Hybrid physical chart]
\label{prop:V85-hybrid-chart}
The correct physical target after the punctured reduction is
\begin{equation}
 \mathcal R_{85}^{\delta}
 :=
 (\dot{\mathcal P}_{85}^{\delta})^{\rm red}
 \oplus\mathcal E_{80}^0
 \label{eq:V85-hybrid-target}
\end{equation}
on a compact joint.  The second summand carries exactly the V59 side and
cap forms.  Its full V74 constant corner pair has zero momentum on the
physical graph and its coordinate is removed by constant residual
translation, so no independent corner class is added.  The total spatial
zero class remains absent.
\end{proposition}

\begin{proof}
This is the direct-sum form of
Theorem~\ref{thm:V80-zero-form-matching},
Theorem~\ref{thm:V80-zero-mode-Lagrangian}, and
Proposition~\ref{prop:V80-total-zero-excluded}.  The spectral sectors are
orthogonal, and no inverse \(Q^{-1}\) is evaluated on \(P_0\).
\end{proof}

\begin{proposition}[Regression against the V81 and V84 packets]
\label{prop:V85-packet-regression}
On the punctured smooth core:
\begin{enumerate}[label=\textup{(\roman*)},leftmargin=2.8em]
\item the combined metric and corner residual cocycles vanish identically,
      so the V81 separated \(r^{-1}\) generator is not formed;
\item the explicit V78 coordinate section remains unbounded, since its
      \(L^2_q\) size \(r^{-3/2}\) has
      \(\dot H^{1-\delta}\) size \(r^{-1/2-\delta}\);
\item the V84 orbit-invariant physical cap difference has
      \(\dot H^{1-\delta}\) size
      \begin{equation}
       r^{1-\delta}r^{-1/2}
       =r^{1/2-\delta},
       \label{eq:V85-V84-regression}
      \end{equation}
      which is bounded and tends to zero as \(r\downarrow0\).
\end{enumerate}
Thus the particular quotient lower bound which ruled out the
inhomogeneous phase no longer diverges, but these statements do not yet
prove the full V59 quotient estimate.
\end{proposition}

\begin{proof}
Item (i) is Theorem~\ref{thm:V85-homogeneous-reduction} on annular smooth
data.  Item (ii) multiplies the exact growth from
\eqref{eq:V83-X0-packet-growth} by the homogeneous joint weight
\(r^{1-\delta}\).  Item (iii) applies the same weight to
\eqref{eq:V84-cap-growth-fibre}.  Because \(0<\delta<1/2\), its exponent is
positive.
\end{proof}

\begin{firewall}[Packet regression is not the quotient gate]
\label{fire:V85-regression-scope}
Equation~\eqref{eq:V85-V84-regression} removes one proved obstruction.  It
does not bound the full homogeneous metric--corner representative modulo
\(\dot W_{85}^{\delta}\), establish closed range of the residual action, or
identify the quotient with all V59 data.  The explicit V78 section remains
unusable, exactly as required by V83.
\end{firewall}

\subsection{V85 result ledger and next acquisition}
\label{subsec:V85-ledger}

\paragraph{New exact conclusions.}
\begin{enumerate}[label=\textup{(\arabic*)},leftmargin=2.8em]
\item The mandatory stable NS V31 and Yang--Mills V82 snapshots were read
      and hashed without modifying either companion.
\item The two-cap trace into independent product homogeneous spaces has no
      bounded right inverse; its antisymmetric channel has a finite
      nonzero zero-frequency extension cost.
\item The Einstein normal trace avoids that generic obstruction because
      its exact rows factor as \(Q^2f\) and \(D^A(D_0f)\).
\item A residual intersection of the orders above and below the critical
      exponent maps boundedly to the exact homogeneous cotangent pairs.
\item The punctured homogeneous metric and corner phases are strong, their
      residual orbit is isotropic, and closing that orbit gives a
      nondegenerate Hilbert reduction with the continuous V75 master
      identity.
\item The compact \(q=0\) physical sector is adjoined by the direct V80
      two-polarization chart; no constant corner class survives and the
      total spatial zero remains absent.
\item The V84 invariant obstruction decays like \(r^{1/2-\delta}\) in the
      new joint topology, while the explicit V78 section remains
      unbounded.  The full quotient estimate is still open.
\end{enumerate}

\paragraph{Next forced step.}
V86 must formulate the homogeneous analogue of
\eqref{eq:V83-quotient-gate} and compute it on the complete V59 scaled
interface family, not on the V78 section.  The most efficient route is the
dual one: classify all bounded orbit annihilators of
\(\dot V_{85}^{\delta}\) on a fixed \(r=1\) fibre, use the exact V59
dilation to transport them to \(r>0\), and prove either a uniform
observability inequality or a new invariant counterexample.  The
two-sided residual intersection and the direct V80 zero chart must remain
fixed during that calculation.

\begin{assessment}[V85 continuum status]
\label{ass:V85-continuum-status}
The full Standard--Model--Einstein continuum remains unproved.  V85
constructs a rigorous punctured homogeneous off-shell residual/corner phase
and its closed reduction, using the special trace-adjusted-Hessian image
rather than a false generic product trace.  Identification with the full
V59 energy phase, uniform control as \(\delta\downarrow0\), nonlinear
curved gluing, the Brown--York/Bianchi stress identity, the common
renormalized Standard Model stress domain, and the cofinal interacting
continuum remain open.
\end{assessment}

\section*{V85 exact checks and append-only record}
\addcontentsline{toc}{section}{V85 exact checks and append-only record}

\begin{equation}
\boxed{
\begin{gathered}
 J_{s,T}^-(r)
 =\frac1{2\pi}\int_{\mathbb R}
 \frac{|e^{i\tau T}-1|^2}{(\tau^2+r^2)^s}\,\dd\tau
 \longrightarrow J_{s,T}^-(0)\in(0,\infty),\\
 R^0=\sigma Q^2(f|_J),\qquad
 R^A=\sigma D^A((D_0f)|_J),\\
 \dot{\mathfrak G}_{85}^{\delta}
 =
 (\dot H^{3/2-\delta}\cap\dot H^{3/2+\delta})
 \oplus
 (\dot H^{1/2-\delta}\cap\dot H^{1/2+\delta}),\\
 \dot\omega_{85}^{\delta}
 (\dot V_{85}^{\delta}\epsilon,\dot V_{85}^{\delta}\zeta)=0,
 \qquad
 r^{1-\delta}r^{-1/2}=r^{1/2-\delta}.
\end{gathered}}
\label{eq:V85-check-card}
\end{equation}

The V84 predecessor used for this append has SHA--256
\begin{center}
\texttt{524822d271fe9385b6525ba6916970cc23b285e1e38381781515d6e477db4d68}.
\end{center}
The V85 candidate retains its bytes up to the sole live final document
terminator, appends this material, and restores exactly one active
\verb|\end{document}|.  Validation includes balanced environments and
braces, unique labels, resolved new internal references, valid inline and
display math delimiters, exact predecessor-prefix equality, and a final
inventory which removes only the superseded SM V84 file.

% =============================================================================
% END APPENDED VERSION V85 MATERIAL
% =============================================================================

\clearpage
% =============================================================================
% BEGIN APPENDED VERSION V86 MATERIAL
% =============================================================================
\section{V86: a critical anisotropic boundary phase and the exact
corner-erasing residual correction}
\label{sec:V86-anisotropic-phase}

V85 constructed a rigorous punctured homogeneous phase by using the special
trace-adjusted-Hessian image.  Its isotropic face topology is a valid
off-shell Hilbert topology, but V59 supplies a more precise topology for
finite-energy time histories.  At every time, the common metric--continuum
trace costs exactly \(r(1+mv)\) in the V59 position form.  After integration
in time, the natural critical face configuration space is therefore
\(L^2_t\dot H^{1/2}_y\), with dual
\(L^2_t\dot H^{-1/2}_y\).

This note builds the corresponding critical anisotropic residual/corner
phase.  No \(\delta>0\) loss is needed because time regularity is stated
separately and the corner momentum is obtained from the exact Hessian rows.
The note then constructs an explicit residual orbit vector which removes all
three V78 corner coordinates on the physical plus graph.  Although its
parameter contains \(Q^{-2}\), its metric gauge image has only three
components, two of which are rewritten directly in physical variables.
Thus the remaining V59 quotient problem becomes a definite metric trace and
conormal estimate rather than a corner-coordinate estimate.

\subsection{The V59 energy-selected history norm}
\label{subsec:V86-energy-history}

For a joint function \(z(y)\), write
\[
 \|z\|_{\dot H_y^{1/2}}^2
 :=\int_{\mathbb R^2}|q|\,|\widehat z(q)|^2\,\dd^2q.
\]
For a side history \(z(t,y)\), define
\begin{equation}
 \|z\|_{\mathcal X_\Gamma}^2
 :=
 \int_0^T\|z(t,\cdot)\|_{\dot H_y^{1/2}}^2\,\dd t,
 \qquad
 \mathcal X_\Gamma:=
 L^2(0,T;\dot H_y^{1/2}).
 \label{eq:V86-face-config-space}
\end{equation}
Its continuous dual is
\begin{equation}
 \mathcal X_\Gamma^*
 =
 L^2(0,T;\dot H_y^{-1/2}).
 \label{eq:V86-face-momentum-space}
\end{equation}

\begin{theorem}[Exact V59 interface-history estimate]
\label{thm:V86-V59-history}
Let \((a,b)\) be one polarization of a V59 finite-energy solution and let
\(z(t,y)=a(t,0,y)=b(t,0,y)\) be its common interface trace.  Then
\begin{equation}
 (1+mv)\|z\|_{\mathcal X_\Gamma}^2
 \le
 2T\,E_{\rm q}(0).
 \label{eq:V86-V59-history-bound}
\end{equation}
The spatial exponent \(1/2\) is sharp under V59 dilation.
\end{theorem}

\begin{proof}
At each \(q\ne0\) and almost every \(t\),
Proposition~\ref{prop:V84-sharp-trace} gives
\[
 r(1+mv)|\widehat z(t,q)|^2
 \le\mathfrak q_r[a(t),b(t)].
\]
Integrate in \(q\) and \(t\).  The position form is at most twice the
conserved energy \(E_{\rm q}(t)=E_{\rm q}(0)\), which proves
\eqref{eq:V86-V59-history-bound}.  The scaled minimizers
\eqref{eq:V84-scaled-position} have unit position energy and trace of size
\(r^{-1/2}\).  Equation~\eqref{eq:V84-scaled-boundary-evolution} and
continuity of its nonzero unit-frequency trace show that this size persists
on the fixed interval \([0,T]\) when \(r\) is small.  Replacing
\(\dot H^{1/2}\) by \(\dot H^\beta\) multiplies that history by
\(r^\beta\), which is uniformly bounded near zero only when
\(\beta\ge1/2\).  Hence the exponent is sharp.
\end{proof}

\begin{remark}[Time and tangential regularity are not merged]
\label{rem:V86-anisotropic-choice}
The estimate \eqref{eq:V86-V59-history-bound} is an
\(L^2_t\dot H_y^{1/2}\) estimate on a fixed time interval.  It does not
assert an isotropic \(\dot H^{1/2}_{t,y}\) trace theorem.  Keeping these
orders separate is necessary for an exact comparison with the conserved
V59 energy.
\end{remark}

\subsection{A critical anisotropic residual domain}
\label{subsec:V86-residual-domain}

All spaces below are punctured at \(q=0\), using the homogeneous conventions
of Subsection~\ref{subsec:V85-homogeneous-conventions}.  Let
\(\epsilon_a\) denote the components tangent to the timelike face,
\(f:=\epsilon_n\) its face-normal component, and
\(d_\mu=D_n\epsilon_\mu|_\Gamma\) the independent normal jet.  Define
\begin{align}
 \mathfrak E_{\rm tan}
 &:=
 H^1(0,T;\dot H_y^{1/2})
 \cap
 L^2(0,T;\dot H_y^{3/2}),
 \label{eq:V86-tangential-residual}\\
 \mathfrak F_{\rm nor}
 &:=
 H^1(0,T;\dot H_y^{3/2})
 \cap
 H^2(0,T;\dot H_y^{1/2})
 \cap
 H^2(0,T;\dot H_y^{-1/2}),
 \label{eq:V86-normal-residual}\\
 \mathfrak D_{\rm jet}
 &:=
 L^2(0,T;\dot H_y^{1/2}).
 \label{eq:V86-normal-jet}
\end{align}
Every intersection has its graph norm.  Put
\begin{equation}
 \mathfrak G_{86}
 :=
 \mathfrak E_{\rm tan}^{\,\oplus3}
 \oplus\mathfrak F_{\rm nor}
 \oplus\mathfrak D_{\rm jet}^{\,\oplus4}.
 \label{eq:V86-residual-space}
\end{equation}
The finite multiplicities only record components and may be replaced by the
corresponding bundles.

\begin{lemma}[Bounded anisotropic face and cap maps]
\label{lem:V86-maps-bounded}
On \(\mathfrak G_{86}\), the four residual maps obey
\begin{align}
 G_\Gamma(e,d)&\in\mathcal X_\Gamma(E_\Gamma),
 \label{eq:V86-G-map}\\
 P_\Gamma(e,d)=\mathscr P(f)&\in
 \mathcal X_\Gamma^*(E_\Gamma^*),
 \label{eq:V86-P-map}\\
 T_j(e,d)=e|_{J_j}&\in
 \dot H_y^{1/2}(J_j;E_J),
 \label{eq:V86-T-map}\\
 R_j(e,d)&\in
 \dot H_y^{-1/2}(J_j;E_J^*).
 \label{eq:V86-R-map}
\end{align}
All maps are bounded.
\end{lemma}

\begin{proof}
For the tangential residual components, one time derivative is controlled
by the first factor in \eqref{eq:V86-tangential-residual}, and one joint
derivative is controlled by its second factor.  The mixed face components
also contain \(d+D f\): the \(d\) term is
\eqref{eq:V86-normal-jet}, while \(D_0f\) and \(D_Af\) lie in
\(\mathcal X_\Gamma\) by \eqref{eq:V86-normal-residual}.  This proves
\eqref{eq:V86-G-map}.

In flat coordinates the trace-adjusted Hessian consists of terms of the
three types
\[
 D_A D_Bf,\qquad D_0D_Af,\qquad D_0^2f,
\]
together with their trace combination.  The first is in
\(L^2_t\dot H_y^{-1/2}\) by the
\(\dot H_y^{3/2}\) factor, the second by the
\(H^2_t\dot H_y^{1/2}\) factor, and the third by the
\(H^2_t\dot H_y^{-1/2}\) factor.  This proves
\eqref{eq:V86-P-map}.

The Hilbert-valued time trace theorem sends
\(H^1_t\dot H_y^{1/2}\) to endpoint
\(\dot H_y^{1/2}\).  It applies to every tangential component and, using
the second factor in \eqref{eq:V86-normal-residual}, to \(f\).  Hence
\eqref{eq:V86-T-map}.  Finally, the exact rows
\eqref{eq:V85-R0-factor}--\eqref{eq:V85-RA-factor} give
\[
 R^0=\sigma Q^2f|_{J_j},\qquad
 R^A=\sigma D^A(D_0f)|_{J_j}.
\]
The \(H^1_t\dot H_y^{3/2}\) factor controls the first endpoint in
\(\dot H_y^{-1/2}\), and
\(H^2_t\dot H_y^{1/2}\) controls the second.  This proves
\eqref{eq:V86-R-map}.
\end{proof}

\subsection{The critical anisotropic canonical phase}
\label{subsec:V86-canonical-phase}

Define
\begin{align}
 \mathcal M_{86}
 &:=
 \mathcal X_\Gamma(E_\Gamma)
 \oplus\mathcal X_\Gamma^*(E_\Gamma^*),
 \label{eq:V86-metric-phase}\\
 \mathcal C_{86}
 &:=
 \bigoplus_{j\in\{0,T\}}
 \left[
  \dot H_y^{1/2}(J_j;E_J)
  \oplus
  \dot H_y^{-1/2}(J_j;E_J^*)
 \right],
 \label{eq:V86-corner-phase}\\
 \mathcal P_{86}
 &:=
 \mathcal M_{86}\oplus\mathcal C_{86}.
 \label{eq:V86-product-phase}
\end{align}
Give the two factors their canonical duality forms,
\begin{align}
 \omega_{M,86}((H_1,p_1),(H_2,p_2))
 &:=
 \frac1{2\kappa}
 \left(\langle p_1,H_2\rangle-\langle p_2,H_1\rangle\right),
 \label{eq:V86-metric-form}\\
 \omega_{C,86}((X_1,\Pi_1),(X_2,\Pi_2))
 &:=
 \frac1\kappa
 \left(\langle\Pi_1,X_2\rangle-\langle\Pi_2,X_1\rangle\right).
 \label{eq:V86-corner-form}
\end{align}
Let \(\omega_{86}\) be their product and set
\begin{equation}
 V_{86}\epsilon
 :=
 \bigl((G_\Gamma\epsilon,P_\Gamma\epsilon),
       (T_J\epsilon,-R_J\epsilon)\bigr).
 \label{eq:V86-residual-action}
\end{equation}

\begin{theorem}[Critical anisotropic closed reduction]
\label{thm:V86-anisotropic-reduction}
The phase \((\mathcal P_{86},\omega_{86})\) is strongly symplectic,
\(V_{86}:\mathfrak G_{86}\to\mathcal P_{86}\) is bounded, and
\begin{equation}
 \omega_{86}(V_{86}\epsilon,V_{86}\zeta)=0.
 \label{eq:V86-cocycle-cancellation}
\end{equation}
With
\begin{equation}
 W_{86}:=\overline{\operatorname{ran}V_{86}},
 \qquad
 \mathcal N_{86}:=W_{86}^{\perp_{\omega_{86}}},
 \qquad
 \mathcal P_{86}^{\rm red}:=\mathcal N_{86}/W_{86},
 \label{eq:V86-reduction}
\end{equation}
the reduced space is Hilbert with a continuous nondegenerate form.  The V75
residual plus matched nonminimal charge extends continuously and satisfies
the off-shell master identity on this punctured critical phase.
\end{theorem}

\begin{proof}
The phase is a finite direct sum of \(V\oplus V^*\) canonical Hilbert
phases, so its form is strong.  Boundedness is
Lemma~\ref{lem:V86-maps-bounded}.  On the smooth punctured core, the flat
Green identity gives
\begin{equation}
 \langle\mathscr P(f),\mathcal G\zeta\rangle_\Gamma
 =
 2\sum_j\langle R_j(\epsilon),T_j(\zeta)\rangle.
 \label{eq:V86-Green-identity}
\end{equation}
Both sides are continuous in the spaces of
Lemma~\ref{lem:V86-maps-bounded}, so density extends the identity.  Apply it
in both orders.  The metric cocycle is the alternating joint pairing and
the corner form on \((T,-R)\) is its negative, proving
\eqref{eq:V86-cocycle-cancellation}.

Closure preserves isotropy.  The common continuous moment zero level is
\(W_{86}^{\perp_{\omega_{86}}}\), and the strong double-orthogonal identity
proves nondegeneracy of the quotient form.  The master-equation statement
then follows exactly as in
Theorem~\ref{thm:V85-homogeneous-reduction}.
\end{proof}

\begin{remark}[The direct zero chart is unchanged]
\label{rem:V86-zero-chart}
On a compact joint, the physical target is
\(\mathcal P_{86}^{\rm red}\oplus\mathcal E_{80}^0\).  On the planar direct
integral, the point \(q=0\) is absent.  No inverse tangential operator acts
on the V80 chart.
\end{remark}

\subsection{An exact orbit vector erasing the V78 corner coordinate}
\label{subsec:V86-corner-erasing}

On the smooth nonzero-frequency physical plus graph, write
\[
 s:=\varsigma,\qquad
 \phi:=\frac12Q^{-2}s.
\]
Define a residual parameter tangent to the timelike face by
\begin{equation}
 e_0:=-D_0\phi,
 \qquad
 e_A:=D_A\phi,
 \qquad
 e_n:=0,
 \qquad
 d:=0.
 \label{eq:V86-erasing-residual}
\end{equation}

\begin{proposition}[Exact corner erasure and gauge image]
\label{prop:V86-corner-erasure}
The residual parameter \eqref{eq:V86-erasing-residual} satisfies, on every
temporal joint,
\begin{equation}
 T_Je=-X^{\rm V78},
 \label{eq:V86-erases-X}
\end{equation}
where \(X^{\rm V78}\) is the complete longitudinal V78 physical dressing.
It has \(P_\Gamma(e)=R_J(e)=0\), and its nonzero metric gauge components are
\begin{align}
 (\mathcal Ge)_{00}&=\frac1{\sqrt2}a,
 \label{eq:V86-erasing-G00}\\
 (\mathcal Ge)_{0A}&=0,
 \label{eq:V86-erasing-G0A}\\
 (\mathcal Ge)_{AB}&=D_AD_BQ^{-2}s.
 \label{eq:V86-erasing-GAB}
\end{align}
Consequently the V78 zero-level representative is equivalent in the
residual quotient to one with
\begin{equation}
 X=0,
 \qquad
 \Pi_{\rm cor}\ \hbox{unchanged},
 \qquad
 H\longmapsto H+\mathcal Ge.
 \label{eq:V86-erased-representative}
\end{equation}
\end{proposition}

\begin{proof}
Since \(\Delta_J=-Q^2\), the V78 coordinate formulas are
\[
 X_0=\frac12Q^{-2}D_0s=D_0\phi,
 \qquad
 X_A=\frac12D_A\Delta_J^{-1}s=-D_A\phi.
\]
Thus \eqref{eq:V86-erasing-residual} gives \(e_b=-X_b\), proving
\eqref{eq:V86-erases-X}.  The face-normal component \(e_n\) is zero, so the
special Brown--York gauge momentum \(P_\Gamma=\mathscr P(e_n)\), and hence
its joint normal trace \(R_J\), vanish.  Therefore the corner momentum does
not move.

Direct differentiation gives
\[
 (\mathcal Ge)_{00}=-2D_0^2\phi,\qquad
 (\mathcal Ge)_{0A}=D_0D_A\phi-D_AD_0\phi=0,\qquad
 (\mathcal Ge)_{AB}=2D_AD_B\phi.
\]
Lemma~\ref{lem:V84-X0-derivative} is equivalent to
\(-2D_0^2\phi=a/\sqrt2\), which proves
\eqref{eq:V86-erasing-G00}.  The last expression is
\eqref{eq:V86-erasing-GAB}.  Adding the orbit vector
\(V_{86}e\) proves \eqref{eq:V86-erased-representative}.
\end{proof}

\begin{assessment}[What remains after corner erasure]
\label{ass:V86-after-erasure}
The large \(X_0\) and longitudinal \(X_A\) coordinates are removed
simultaneously, including their cap correlation.  The remaining metric
correction has no mixed time--space row.  Its time--time row is the V59
amplitude itself and is therefore measured by the energy-selected history
norm of Theorem~\ref{thm:V86-V59-history}.  The only nonlocal expression
left in the configuration correction is the order-zero joint Riesz block
\(D_AD_BQ^{-2}s\).  The next quotient test must combine this block with the
original V63 \(H_{AB}\) before estimating; estimating it separately would
repeat the forbidden section argument.
\end{assessment}

\subsection{The reduced comparison gate after erasure}
\label{subsec:V86-reduced-gate}

For a smooth V59 physical datum \(u\), let
\[
 Z_u^{\rm er}
 :=
 \bigl(H_u+\mathcal Ge_u,\ p_u,\ 0,\ \Pi_u\bigr)
\in\mathcal N_{86}
\]
denote the representative \eqref{eq:V86-erased-representative}.  It is
defined on the punctured smooth core; changing any preliminary lift changes
its class only by \(W_{86}\).

\begin{proposition}[Exact remaining quotient criterion]
\label{prop:V86-remaining-gate}
The smooth comparison extends boundedly from the full V59 energy completion
to \(\mathcal P_{86}^{\rm red}\) if and only if
\begin{equation}
 \inf_{w\in W_{86}}
 \|Z_u^{\rm er}+w\|_{\mathcal P_{86}}
 \le C\|u\|_{\mathcal E_{59}}
 \label{eq:V86-remaining-gate}
\end{equation}
on the smooth punctured core.  In testing this estimate, the configuration
part of \(Z_u^{\rm er}\) is exactly
\begin{equation}
 \left(
  H_{00}+\frac a{\sqrt2},\
  H_{0A},\
  H_{AB}+D_AD_BQ^{-2}s
 \right),
 \label{eq:V86-erased-config-block}
\end{equation}
with \(H_{0A}=0\) on the plus representative.  No corner configuration
term remains.
\end{proposition}

\begin{proof}
The equivalence is the quotient completion argument of
Theorem~\ref{thm:V83-energy-gate}, now in
\(\mathcal P_{86}\).  Proposition~\ref{prop:V86-corner-erasure} gives
\eqref{eq:V86-erased-config-block}.
\end{proof}

\begin{firewall}[The conormal and full metric block are still open]
\label{fire:V86-gate-open}
Theorem~\ref{thm:V86-V59-history} controls the common radiative interface
trace, not yet every component of
\eqref{eq:V86-erased-config-block}.  It also does not prove the hidden
boundary regularity
\(p_u\in L^2_t\dot H_y^{-1/2}\), nor a bound for the unchanged corner
momentum \(\Pi_u\).  These components and their residual correlations must
be included before \eqref{eq:V86-remaining-gate} can be claimed.
\end{firewall}

\subsection{V86 result ledger and next acquisition}
\label{subsec:V86-ledger}

\paragraph{New exact conclusions.}
\begin{enumerate}[label=\textup{(\arabic*)},leftmargin=2.8em]
\item The V59 conserved energy controls its common interface history in
      \(L^2_t\dot H_y^{1/2}\), with the exact coefficient \(1+mv\);
      the spatial half order is sharp.
\item A critical anisotropic residual domain maps boundedly to the exact
      canonical pairs
      \[
       (L^2_t\dot H_y^{1/2}\oplus
        L^2_t\dot H_y^{-1/2})_\Gamma
       \oplus
       (\dot H_y^{1/2}\oplus\dot H_y^{-1/2})_{J_0\sqcup J_T}.
      \]
\item The special Hessian factorization supplies all cap momenta, the weak
      Green identity is continuous, and the metric and corner cocycles
      cancel on this critical phase.
\item Closing the residual orbit gives a nondegenerate Hilbert reduction
      and a continuous off-shell V75 master identity.
\item The residual parameter \eqref{eq:V86-erasing-residual} removes the
      complete V78 corner coordinate without moving its momentum.
\item Its mixed metric row vanishes, its \(00\) row is \(a/\sqrt2\), and
      its spatial row is the order-zero Riesz block
      \(D_AD_BQ^{-2}s\).
\item The full comparison is reduced to
      \eqref{eq:V86-remaining-gate} on the erased representative; metric
      conormal and corner-momentum estimates remain open.
\end{enumerate}

\paragraph{Next forced step.}
V87 must insert the explicit V63 plus polarization matrix into
\eqref{eq:V86-erased-config-block}.  It should determine whether
\(H_{AB}+D_AD_BQ^{-2}s\) is an energy-bounded trace of the common V59
amplitude or retains an infrared singular component.  In the same Fourier
frame it must combine the Brown--York conormal \(p_u\) with the unchanged
\(\Pi_u\) before taking the
\(L^2_t\dot H_y^{-1/2}\oplus\dot H_y^{-1/2}\) norm.  If either combined
block is unbounded, its orbit-annihilating functional must be exhibited;
another section estimate is insufficient.

\begin{assessment}[V86 continuum status]
\label{ass:V86-continuum-status}
The full Standard--Model--Einstein continuum remains unproved.  V86
constructs an energy-matched critical anisotropic off-shell phase and
removes the divergent corner configuration by an exact residual orbit
vector.  The remaining linear flat gate is the combined metric
trace--conormal estimate \eqref{eq:V86-remaining-gate}.  Nonlinear curved
gluing, the Brown--York/Bianchi stress identity, the common renormalized
Standard Model stress domain, and the cofinal interacting continuum also
remain open.
\end{assessment}

\section*{V86 exact checks and append-only record}
\addcontentsline{toc}{section}{V86 exact checks and append-only record}

\begin{equation}
\boxed{
\begin{gathered}
 (1+mv)\|a|_\Gamma\|_{L^2_t\dot H_y^{1/2}}^2
 \le2T E_{\rm q}(0),\\
 \mathcal P_{86}
 =
 (L^2_t\dot H_y^{1/2}\oplus L^2_t\dot H_y^{-1/2})_\Gamma
 \oplus
 (\dot H_y^{1/2}\oplus\dot H_y^{-1/2})_{J_0\sqcup J_T},\\
 e_0=-\tfrac12Q^{-2}D_0s,\qquad
 e_A=\tfrac12D_AQ^{-2}s,\qquad T_Je=-X^{\rm V78},\\
 (\mathcal Ge)_{00}=a/\sqrt2,\qquad
 (\mathcal Ge)_{0A}=0,\qquad
 (\mathcal Ge)_{AB}=D_AD_BQ^{-2}s.
\end{gathered}}
\label{eq:V86-check-card}
\end{equation}

The V85 predecessor used for this append has SHA--256
\begin{center}
\texttt{c18f6ac51ffac9b10bb0451a152bcc47258d7a7c1e09acb74258590c894486c6}.
\end{center}
The V86 candidate retains its bytes up to the sole live final document
terminator, appends this material, and restores exactly one active
\verb|\end{document}|.  Validation includes balanced environments and
braces, unique labels, resolved new internal references, valid inline and
display math delimiters, exact predecessor-prefix equality, and a final
inventory which removes only the superseded SM V85 file.

% =============================================================================
% END APPENDED VERSION V86 MATERIAL
% =============================================================================

\clearpage
% =============================================================================
% BEGIN APPENDED VERSION V87 MATERIAL
% =============================================================================
\section{V87: exact plus-sector orbit normal form, a conormal quotient
obstruction, and the forced port completion}
\label{sec:V87-orbit-normal-form}

V86 removed the V78 corner configuration and reduced its gauge image to a
metric block.  This note completes the same operation on the plus-sector
corner momentum.  A face-normal residual scalar removes the momentum without
changing the metric configuration.  Direct substitution of the V62 TT
matrix then cancels every incidence-angle coefficient.  The resulting
orbit representative contains only the V59 trace \(a\) and conormal
\(D_na\).

This normal form makes the remaining topology question decisive.  The V59
radiative conormal is a residual-quotient invariant.  An incoming packet
with fixed nonzero normal frequency and tangential frequency tending to zero
has order-one \(L^2\) conormal trace, while the canonical
\(\dot H_y^{-1/2}\) momentum norm diverges.  Thus the V86 critical canonical
phase, like V83 before it, is a valid off-shell test phase but cannot be the
physical completion of the full V59 energy space.

The correct physical boundary completion is the characteristic port space
of V67/V68.  Its topology is \(L^2\), its interface relation is the exact
unitary matrix \(S_m\), and it keeps trace and conormal in their energy-flux
combination.  The off-shell canonical phase and the physical energy
completion must therefore be related as a rigged test/solution pair rather
than identified by equivalent Hilbert norms.

\subsection{A residual correction erasing the corner momentum}
\label{subsec:V87-momentum-erasure}

Continue on the smooth nonzero-frequency plus graph.  Its physical V78
corner momentum is
\begin{equation}
 \Pi_{\rm cor}^{\,b}
 =-\sigma_J\Pi_{(1)}^{0b}.
 \label{eq:V87-physical-corner-momentum}
\end{equation}
Set
\begin{equation}
 f:=\frac12Q^{-2}D_ns,
 \qquad
 e_n:=f,\qquad e_a:=0,\qquad
 d_a:=-D_af,\qquad d_n:=0.
 \label{eq:V87-momentum-residual}
\end{equation}

\begin{lemma}[The correction has zero metric configuration image]
\label{lem:V87-zero-G}
The residual jet \eqref{eq:V87-momentum-residual} obeys
\begin{equation}
 G_\Gamma(e,d)=0.
 \label{eq:V87-zero-G}
\end{equation}
\end{lemma}

\begin{proof}
On a flat face, the tangential--tangential gauge trace uses only
\(D_ae_b+D_be_a\), which is zero.  A mixed component is
\(d_a+D_ae_n=-D_af+D_af=0\), and the normal--normal component is
\(2d_n=0\).
\end{proof}

\begin{proposition}[Exact momentum erasure]
\label{prop:V87-momentum-erasure}
For the residual jet \eqref{eq:V87-momentum-residual},
\begin{equation}
 R_J(e,d)=\Pi_{\rm cor}.
 \label{eq:V87-R-equals-Pi}
\end{equation}
Consequently addition of its orbit vector sends
\begin{equation}
 \Pi_{\rm cor}\longmapsto0,\qquad
 p\longmapsto p+\mathscr P(f),
 \label{eq:V87-momentum-shift}
\end{equation}
while leaving the metric and the already erased corner coordinate fixed.
\end{proposition}

\begin{proof}
The V78 plus rows are
\[
 \Pi_{(1)}^{00}=-\frac12D_ns,
 \qquad
 \Pi_{(1)}^{0A}=\frac12D_0\mathcal L^AD_ns,
 \qquad
 \mathcal L^A=-D^AQ^{-2}.
\]
Therefore
\begin{align*}
 \Pi_{\rm cor}^{\,0}
 &=\frac{\sigma_J}{2}D_ns
 =\sigma_JQ^2f,\\
 \Pi_{\rm cor}^{\,A}
 &=\frac{\sigma_J}{2}D_0D^AQ^{-2}D_ns
 =\sigma_JD_0D^Af.
\end{align*}
These are exactly the special normal-trace rows
\eqref{eq:V85-R0-factor}--\eqref{eq:V85-RA-factor}, proving
\eqref{eq:V87-R-equals-Pi}.  The corner part of the residual action is
\((T,-R)\); here \(T_a=e_a=0\), so
\(\Pi_{\rm cor}\mapsto\Pi_{\rm cor}-R=0\).  The metric momentum shifts by
\(P_\Gamma=\mathscr P(f)\), and
Lemma~\ref{lem:V87-zero-G} proves the remaining assertions.
\end{proof}

\subsection{Exact angular cancellation}
\label{subsec:V87-angular-cancellation}

Choose the V62 frame with tangential covector along \(y^1\):
\begin{equation}
 q=(\lambda\mathfrak s,0),\qquad
 D_n=i\lambda\mathfrak c,\qquad
 \mathfrak s^2+\mathfrak c^2=1.
 \label{eq:V87-angular-symbols}
\end{equation}
The norm-one plus tensor
\(E_+=(e\otimes e-t\otimes t)/\sqrt2\) has induced components
\begin{equation}
 H_{11}=-\frac{\mathfrak c^2}{\sqrt2}a,\qquad
 H_{22}=\frac1{\sqrt2}a,\qquad
 s=H_A{}^A=\frac{\mathfrak s^2}{\sqrt2}a.
 \label{eq:V87-plus-induced-components}
\end{equation}

\begin{lemma}[Configuration normal form]
\label{lem:V87-config-normal-form}
After the corner-erasing correction of
Proposition~\ref{prop:V86-corner-erasure}, the induced metric components
are
\begin{equation}
 H'_{ab}
 =
 \frac a{\sqrt2}
 \begin{pmatrix}
  1&0&0\\
  0&-1&0\\
  0&0&1
 \end{pmatrix}_{ab},
 \qquad a,b\in\{0,1,2\}.
 \label{eq:V87-config-normal-form}
\end{equation}
This matrix is independent of \(\mathfrak s,\mathfrak c\).
\end{lemma}

\begin{proof}
The correction has \(00\) component \(a/\sqrt2\), zero mixed component,
and spatial component \(D_AD_BQ^{-2}s\) by
Proposition~\ref{prop:V86-corner-erasure}.  At
\eqref{eq:V87-angular-symbols}, its only nonzero spatial entry is
\[
 D_1^2Q^{-2}s=-s.
\]
Thus
\[
 H'_{11}
 =-\frac{\mathfrak c^2}{\sqrt2}a
  -\frac{\mathfrak s^2}{\sqrt2}a
 =-\frac a{\sqrt2},
\qquad
 H'_{22}=\frac a{\sqrt2}.
\]
Together with the \(00\) and mixed rows, this is
\eqref{eq:V87-config-normal-form}.
\end{proof}

For the momentum calculation, the nonzero Brown--York components of the
ungauged plus tensor are
\begin{align}
 \Pi_{(1)}^{00}
 &=-\frac{i\lambda\mathfrak c\mathfrak s^2}{2\sqrt2}a,
 &
 \Pi_{(1)}^{01}
 &=-\frac{i\lambda\mathfrak c\mathfrak s}{2\sqrt2}a,
 \notag\\
 \Pi_{(1)}^{11}
 &=-\frac{i\lambda\mathfrak c}{2\sqrt2}a,
 &
 \Pi_{(1)}^{22}
 &=\frac{i\lambda\mathfrak c(2-\mathfrak c^2)}
          {2\sqrt2}a.
 \label{eq:V87-plus-BY-matrix}
\end{align}
All omitted components vanish.  These formulas follow directly from
\[
 K_{ab}^{(1)}
 =\frac12(D_nH_{ab}-D_aH_{nb}-D_bH_{na}),
 \qquad
 \Pi_{(1)}^{ab}=K_{(1)}^{ab}-K_{(1)}\eta_\partial^{ab},
\]
and agree with the time rows already checked in V78.

At the same Fourier point,
\begin{equation}
 f=\frac{i\mathfrak c}{2\sqrt2\lambda}a.
 \label{eq:V87-f-symbol}
\end{equation}
Using the V72 sign convention
\(\mathscr P^{ab}(f)=-D^aD^bf+
\eta_\partial^{ab}\Box_\partial f\), one obtains
\begin{align}
 \mathscr P^{00}(f)
 &=\frac{i\lambda\mathfrak c\mathfrak s^2}{2\sqrt2}a,
 &
 \mathscr P^{01}(f)
 &=\frac{i\lambda\mathfrak c\mathfrak s}{2\sqrt2}a,
 \notag\\
 \mathscr P^{11}(f)
 &=\frac{i\lambda\mathfrak c}{2\sqrt2}a,
 &
 \mathscr P^{22}(f)
 &=\frac{i\lambda\mathfrak c^3}{2\sqrt2}a.
 \label{eq:V87-P-matrix}
\end{align}

\begin{theorem}[Complete plus-sector orbit normal form]
\label{thm:V87-plus-normal-form}
After applying successively the residual corrections
\eqref{eq:V86-erasing-residual} and
\eqref{eq:V87-momentum-residual}, the full plus-sector zero-level
representative is
\begin{align}
 H'_{ab}
 &=\frac a{\sqrt2}\operatorname{diag}(1,-1,1)_{ab},
 \label{eq:V87-final-H}\\
 p'^{ab}
 &=\frac{D_na}{\sqrt2}\,
   \delta_2^a\delta_2^b,
 \label{eq:V87-final-p}\\
 X'&=0,\qquad \Pi'_{\rm cor}=0.
 \label{eq:V87-final-corner}
\end{align}
All formulas are independent of incidence angle and extend by rotation of
the joint frame.
\end{theorem}

\begin{proof}
Lemma~\ref{lem:V87-config-normal-form} gives
\eqref{eq:V87-final-H}, and
Propositions~\ref{prop:V86-corner-erasure} and
\ref{prop:V87-momentum-erasure} give
\eqref{eq:V87-final-corner}.  Adding
\eqref{eq:V87-plus-BY-matrix} and \eqref{eq:V87-P-matrix} cancels the
\(00\), \(01\), and \(11\) entries.  The remaining entry is
\[
 p'^{22}
 =
 \frac{i\lambda\mathfrak c}{2\sqrt2}
 \left(2-\mathfrak c^2+\mathfrak c^2\right)a
 =
 \frac{i\lambda\mathfrak c}{\sqrt2}a
 =
 \frac{D_na}{\sqrt2}.
\]
Rotational covariance replaces \(\delta_2\) by the unit joint vector
orthogonal to \(q\).
\end{proof}

\subsection{A radiative-conormal quotient invariant}
\label{subsec:V87-conormal-invariant}

Let \(\mathcal Q_{\rm rad}\) be the V53/V63 radiative quotient map on metric
trace--conormal data.  By the exact sequence,
\begin{equation}
 \mathcal Q_{\rm rad}A\epsilon=0.
 \label{eq:V87-Qrad-gauge-zero}
\end{equation}
Let \(\mathcal Q_{\rm con}\) denote its conormal component.  On the
plus-sector normal form,
\begin{equation}
 \mathcal Q_{\rm con}(H',p')=D_na.
 \label{eq:V87-Qcon-normal-form}
\end{equation}

\begin{lemma}[Bounded conormal annihilator on a normal-frequency band]
\label{lem:V87-conormal-annihilator}
After restricting the normal spectral variable to
\(1\le|\xi|\le2\), the map
\begin{equation}
 \Lambda_{\rm con}Z
 :=
 \mathcal Q_{\rm con}(H,p)
 \label{eq:V87-conormal-functional}
\end{equation}
is a bounded order-zero map from the V86 metric phase to
\(L^2_t\dot H_y^{-1/2}\).  It annihilates
\(\operatorname{ran}V_{86}\) and hence \(W_{86}\).  Therefore
\begin{equation}
 \|[Z]\|_{\mathcal P_{86}^{\rm red}}
 \ge C^{-1}
 \|\Lambda_{\rm con}Z\|_{L^2_t\dot H_y^{-1/2}}
 \label{eq:V87-conormal-lower}
\end{equation}
on that band.
\end{lemma}

\begin{proof}
The V63 quotient and splitting matrices are uniformly bounded on the
complete angular interval.  On the fixed normal band, all remaining
normal multipliers are bounded above and below, so the conormal row is an
order-zero matrix on the V86 momentum topology.  Equation
\eqref{eq:V87-Qrad-gauge-zero} proves annihilation of the residual range;
continuity gives annihilation of its closure.  Apply the same quotient
annihilator argument as in
Lemma~\ref{lem:V84-cap-invariant} to obtain
\eqref{eq:V87-conormal-lower}.
\end{proof}

\subsection{The normal-incidence packet}
\label{subsec:V87-normal-packet}

Use the V67 characteristic variables
\[
 u_{\rm in}=D_0a+D_na,\qquad
 u_{\rm out}=D_0a-D_na
\]
and the continuum variables \(c_{\rm in},c_{\rm out}\).  The exact interface
graph is
\begin{equation}
 \binom{u_{\rm out}}{c_{\rm out}}
 =
 \frac1{1+m}
 \begin{pmatrix}
  1-m&2\sqrt m\\
  2\sqrt m&m-1
 \end{pmatrix}
 \binom{u_{\rm in}}{c_{\rm in}}.
 \label{eq:V87-port-graph}
\end{equation}
For \(c_{\rm in}=0\),
\begin{equation}
 D_na
 =\frac{u_{\rm in}-u_{\rm out}}2
 =\frac{m}{1+m}u_{\rm in}.
 \label{eq:V87-conormal-from-incoming}
\end{equation}

\begin{theorem}[The V86 quotient gate fails on incoming packets]
\label{thm:V87-V86-gate-fails}
There are smooth V59 energy-unit packets \(u_\varepsilon\), with
tangential support
\(\{\varepsilon<|q|<2\varepsilon\}\), normal spectral support
\(\{1<|\xi|<2\}\), and no incoming continuum wave, such that
\begin{equation}
 \|[Z_{u_\varepsilon}]\|_{\mathcal P_{86}^{\rm red}}
 \ge c_{T,m}\varepsilon^{-1/2}.
 \label{eq:V87-V86-quotient-growth}
\end{equation}
Hence the V86 comparison estimate
\eqref{eq:V86-remaining-gate} is false.
\end{theorem}

\begin{proof}
Choose a smooth incoming metric characteristic pulse with normal Fourier
support in \((1,2)\), translate it so that a fixed positive portion crosses
the interface during \((0,T)\), and normalize its incoming energy to one.
Choose independently
\(\chi_\varepsilon\in C_c^\infty(\mathbb R^2_q)\), supported in the stated
annulus with \(\|\chi_\varepsilon\|_2=1\).  The wave dispersion
\(\omega=(\xi^2+|q|^2)^{1/2}\) has
\(\partial_\xi\omega\) uniformly bounded above and below on these supports.
Plancherel in the characteristic coordinate therefore gives
\begin{equation}
 \|u_{\rm in,\varepsilon}\|_{L^2((0,T)\times\mathbb R^2)}
 \ge c_T>0
 \label{eq:V87-incoming-trace-lower}
\end{equation}
uniformly for small \(\varepsilon\).  The initial bulk energy is one after
a fixed normalization.

Set \(c_{\rm in}=0\).  Equation
\eqref{eq:V87-conormal-from-incoming} gives an \(L^2\) lower bound
\[
 \|D_na_\varepsilon\|_{L^2_{t,y}}
 \ge\frac{m}{1+m}c_T.
\]
On the tangential annulus,
\[
 \|D_na_\varepsilon\|_{L^2_t\dot H_y^{-1/2}}
 \ge c_{T,m}\varepsilon^{-1/2}.
\]
The conormal is the quotient invariant
\eqref{eq:V87-Qcon-normal-form}.  Apply
\eqref{eq:V87-conormal-lower} to prove
\eqref{eq:V87-V86-quotient-growth}.  Smooth compact packets are obtained
by the chosen frequency cutoffs and the standard V59 form core.
\end{proof}

\begin{firewall}[The obstruction is physical, not a dressing norm]
\label{fire:V87-physical-obstruction}
The lower bound uses the radiative quotient conormal, which annihilates
every residual orbit vector.  It therefore cannot be removed by a different
corner or metric representative.  The failure concerns the attempt to put
the physical conormal in the canonical dual
\(\dot H_y^{-1/2}\) topology.  It does not challenge the V59 energy
generator or the exact cancellation of the total Green flux.
\end{firewall}

\subsection{The forced characteristic-port completion}
\label{subsec:V87-port-completion}

Define the physical boundary port Hilbert space
\begin{equation}
 \mathcal U_{87}
 :=
 L^2((0,T)\times\mathbb R^2;\mathscr T)^{\oplus4},
 \qquad
 U=(u_{\rm in},c_{\rm in},u_{\rm out},c_{\rm out}),
 \label{eq:V87-port-space}
\end{equation}
with Green form
\begin{align}
 \mathfrak B_{87}(U,U')
 &:=
 \frac14\left[
  \langle u_{\rm in},u'_{\rm in}\rangle
 +\langle c_{\rm in},c'_{\rm in}\rangle
 \right.
 \notag\\
 &\hspace{7em}\left.
 -\langle u_{\rm out},u'_{\rm out}\rangle
 -\langle c_{\rm out},c'_{\rm out}\rangle
 \right].
 \label{eq:V87-port-Green-form}
\end{align}

\begin{proposition}[Closed maximal-neutral physical interface]
\label{prop:V87-port-interface}
The graph \(\mathcal L_{87}\subset\mathcal U_{87}\) defined by
\eqref{eq:V87-port-graph} is closed and maximal neutral for
\(\mathfrak B_{87}\).  The incoming packet of
Theorem~\ref{thm:V87-V86-gate-fails} has uniformly bounded port norm.
\end{proposition}

\begin{proof}
The matrix in \eqref{eq:V87-port-graph} is self-adjoint and squares to the
identity by \eqref{eq:V67-Sm-unitary}.  Its graph is therefore closed, the
incoming and outgoing norms agree, and the graph is neutral.  A unitary
graph between equal Hilbert spaces is maximal neutral.  The packet has
unit incoming norm; unitarity gives the same outgoing norm.
\end{proof}

\begin{assessment}[Two completions with different jobs]
\label{ass:V87-two-completions}
The V86 canonical phase remains a valid dense off-shell domain for the
residual moment map and master equation.  Theorem
\ref{thm:V87-V86-gate-fails} proves that its reduced Hilbert norm is not
equivalent to the full V59 physical energy norm.  The latter reaches the
boundary continuously through the characteristic port graph
\(\mathcal L_{87}\), where trace and conormal occur in the energy-flux
combination.  The next step is therefore a rigged construction: use the
canonical phase as a test domain, the V59 energy space as the central
physical Hilbert space, and the port topology for boundary propagation.
\end{assessment}

\subsection{V87 result ledger and next acquisition}
\label{subsec:V87-ledger}

\paragraph{New exact conclusions.}
\begin{enumerate}[label=\textup{(\arabic*)},leftmargin=2.8em]
\item A face-normal residual scalar removes the complete plus corner
      momentum without changing the metric configuration.
\item The two residual corrections put the plus zero-level representative
      in the angle-independent normal form
      \eqref{eq:V87-final-H}--\eqref{eq:V87-final-corner}.
\item The remaining metric variables are exactly the V59 trace \(a\) and
      conormal \(D_na\).
\item The radiative conormal is a bounded residual-orbit annihilator on a
      fixed normal band.
\item Smooth incoming energy-unit packets make its
      \(L^2_t\dot H_y^{-1/2}\) norm grow like
      \(\varepsilon^{-1/2}\), proving failure of the V86 quotient gate.
\item The same packets have bounded \(L^2\) characteristic-port norm.
\item The V67 interface matrix defines a closed maximal-neutral port graph,
      which is the correct physical boundary completion.
\end{enumerate}

\paragraph{Next forced step.}
V88 must construct the rigged bridge explicitly.  It should choose a dense
physical core
\[
 \mathscr D_{\rm BFV}
 \subset\mathcal E_{59}
 \subset\mathscr D_{\rm BFV}^*
\]
on which the V86 canonical moment map is defined, prove invariance of the
core under the V59 unitary group or replace it by an integrated graph
domain, and show that the combined residual/corner Green form descends
continuously to the V87 port graph even though the individual canonical
trace norms do not.  The cross polarization and its transverse corner
momentum must be included in that bridge.

\begin{assessment}[V87 continuum status]
\label{ass:V87-continuum-status}
The full Standard--Model--Einstein continuum remains unproved.  V87 solves
the flat plus-sector residual normal form and proves that no strong
canonical trace/conormal quotient used so far is norm-equivalent to the full
V59 energy phase.  It identifies the closed unitary characteristic-port
graph as the physical completion.  The rigged off-shell/physical bridge,
the cross transverse row, nonlinear curved gluing, the Brown--York/Bianchi
stress identity, the common renormalized Standard Model stress domain, and
the cofinal interacting continuum remain open.
\end{assessment}

\section*{V87 exact checks and append-only record}
\addcontentsline{toc}{section}{V87 exact checks and append-only record}

\begin{equation}
\boxed{
\begin{gathered}
 f=\tfrac12Q^{-2}D_ns,\qquad
 G_\Gamma(f,-Df)=0,\qquad R_J(f)=\Pi_{\rm cor},\\
 H'_{ab}=\frac a{\sqrt2}\operatorname{diag}(1,-1,1)_{ab},
 \qquad
 p'^{ab}=\frac{D_na}{\sqrt2}\delta_2^a\delta_2^b,
 \qquad X'=\Pi'=0,\\
 D_na=\frac{m}{1+m}u_{\rm in}\quad(c_{\rm in}=0),\\
 \|[Z_{u_\varepsilon}]\|_{\mathcal P_{86}^{\rm red}}
 \ge c\varepsilon^{-1/2},
 \qquad
 \|(u_{\rm in},c_{\rm in},u_{\rm out},c_{\rm out})\|_{L^2}=O(1).
\end{gathered}}
\label{eq:V87-check-card}
\end{equation}

The V86 predecessor used for this append has SHA--256
\begin{center}
\texttt{ea93b753f2764d169940d2b0af1ad396c8966f9a53a79fd589d0871d1739d35c}.
\end{center}
The V87 candidate retains its bytes up to the sole live final document
terminator, appends this material, and restores exactly one active
\verb|\end{document}|.  Validation includes balanced environments and
braces, unique labels, resolved new internal references, valid inline and
display math delimiters, exact predecessor-prefix equality, and a final
inventory which removes only the superseded SM V86 file.

% =============================================================================
% END APPENDED VERSION V87 MATERIAL
% =============================================================================


\clearpage
% =============================================================================
% BEGIN APPENDED VERSION V88 MATERIAL
% =============================================================================
\section{V88: the rigged BFV--energy bridge, weak characteristic traces,
and closure of the cross corner row}
\label{sec:V88-rigged-bridge}

V87 proved that the strong V86 canonical quotient cannot be the Hilbert
completion of the V59 energy phase.  The obstruction was a genuine
radiative conormal, not a choice of corner dressing.  The appropriate
replacement is therefore a rigged construction.  Smooth compatible states
retain all canonical face and cap traces and support the V75 master identity;
the V59 energy phase is the pivot Hilbert space; arbitrary energy solutions
have characteristic boundary values only by transposition.  On the dense
subspace where those values are square integrable, the transposed trace is
the V87 port trace.

There is a second distinction which must be kept explicit.  Skew-adjointness
of the V68 generator does not imply an unweighted (L^2) boundary-observation
estimate on its whole energy space.  High-frequency grazing packets rule out
such an estimate.  The unitary port graph is nevertheless meaningful on the
operator core, on its closed trace domain, and distributionally on every
finite-energy solution.  This is sufficient for the weak Green identity.

Finally, this note includes the cross polarization omitted in V87.  Its V78
corner configuration is zero, while its transverse corner momentum is a
rotated tangential derivative of the cross amplitude.  The sharp V84 trace
capacity puts that row exactly in the energy-controlled
\(\dot H_y^{-1/2}\) cap space.  Thus it creates no analogue of the plus
conormal obstruction.

\subsection{The physical generator and a smooth compatible core}
\label{subsec:V88-physical-core}

Let \(\mathcal H_{\rm phys}\) be the closed potential-compatible reducing
subspace of \(\mathcal H_{68}\) in
Theorem~\ref{thm:V68-skew-adjoint-generator}.  Through the V68 phase
isometry it is the V59 radiative--continuum energy space.  Write
\begin{equation}
 \mathscr G:=\mathscr G_{68}|_{\mathcal H_{\rm phys}},
 \qquad
 U(t):=e^{t\mathscr G},
 \label{eq:V88-physical-group}
\end{equation}
so \(\mathscr G^*=-\mathscr G\) and \(U(t)\) is unitary.

Let \(\mathscr C_{\rm comp}\) consist of smooth potential-compatible phase
and continuum states which
\begin{enumerate}[label=\textup{(\roman*)},leftmargin=2.8em]
\item have compact tangential Fourier support in an annulus
      \(n^{-1}\le |q|\le n\) for some \(n\);
\item are smooth up to the interface and compactly supported in the two
      closed normal half-lines, with support allowed to meet the interface;
\item satisfy the V67 port graph and all compatibility conditions obtained
      by repeatedly applying \(\mathscr G\).
\end{enumerate}
Finite sums over annuli are understood.  The last condition is the usual
infinite compatibility condition at the interface; it is imposed on the
jets, not as an extra dynamical equation.

Denote by \(\Gamma U(t)X\) the four strong characteristic traces of a state
in this core, ordered as
\begin{equation}
 \Gamma U(t)X
 =\bigl(u_{\rm in},c_{\rm in},u_{\rm out},c_{\rm out}\bigr)(t).
 \label{eq:V88-strong-port-trace}
\end{equation}
Let \(\langle Q\rangle=(I-\Delta_y)^{1/2}\).  For \(N\ge1\), set
\begin{align}
 p_N(X)^2
 &:={}
 \sum_{j+\ell\le N}
 \|\langle Q\rangle^j\mathscr G^\ell X\|_{\mathcal H_{\rm phys}}^2
 \notag\\
 &\quad+
 \sum_{j+\ell\le N}
 \int_{-N}^{N}
 \|\langle Q\rangle^j\partial_t^\ell
       \Gamma U(t)X\|_{L_y^2}^{,2}\,\dd t.
 \label{eq:V88-BFV-seminorms}
\end{align}

\begin{definition}[Rigged BFV test space]
\label{def:V88-BFV-space}
The space \(\mathscr D_{\rm BFV}\) is the completion of
\(\mathscr C_{\rm comp}\) for the increasing family
\(\{p_N\}_{N\ge1}\), after identifying a Cauchy sequence with its limit in
\(\mathcal H_{\rm phys}\).
\end{definition}

\begin{theorem}[Dense invariant rigging]
\label{thm:V88-dense-rigging}
The identification in Definition~\ref{def:V88-BFV-space} is injective and
gives continuous dense embeddings
\begin{equation}
 \mathscr D_{\rm BFV}
 \lhook\joinrel\longrightarrow
 \mathcal H_{\rm phys}
 \lhook\joinrel\longrightarrow
 \mathscr D_{\rm BFV}',
 \label{eq:V88-rigged-triple}
\end{equation}
where the second map uses the pivot inner product.  The group \(U(s)\)
preserves \(\mathscr D_{\rm BFV}\) and acts continuously on its Frechet
topology.  On every finite slab and after taking sufficiently many
seminorms, its metric, residual, and cap traces lie in the V86 spaces.
Consequently the V86 moment map and the V75 master identity are defined on
\(\mathscr D_{\rm BFV}\).
\end{theorem}

\begin{proof}
The bulk term with \(j=\ell=0\) in every \(p_N\) controls the physical
Hilbert norm.  It remains to exclude a sequence converging to zero in the
pivot while its port traces converge to a nonzero limit.  Let \(X_k\) be
such a sequence and let \(g\) be a smooth compactly supported four-port
test function.  The nonzero characteristic trace map of the maximal V68
and continuum expressions has a smooth collar right inverse \(E\).  The
spacetime Green formula gives
\begin{equation}
 \int_{\mathbb R}\mathfrak B_{87}
      \bigl(\Gamma U(t)X_k,g(t)\bigr)\,\dd t
 =\int_{\mathbb R}
 \left\langle U(t)X_k,
  -\partial_tEg(t)-\mathscr G_{\max}^*Eg(t)
 \right\rangle\,\dd t.
 \label{eq:V88-closability-Green}
\end{equation}
The right side tends to zero because \(U(t)\) is unitary and the lifted
test field is smooth with compact time support.  The port Green form is
nondegenerate, so every distributional pairing with the proposed trace
limit is zero.  The limit is zero.  The same argument after applying
\(\langle Q\rangle^j\partial_t^\ell\) proves simultaneous closability of
all trace operators in \eqref{eq:V88-BFV-seminorms}.  Hence the completed
space injects into the pivot, is complete, and has the asserted first
embedding.  The pivot pairing gives the second continuous embedding.

Smooth vectors for the skew-adjoint generator are dense by spectral
cutoff.  Tangential translations strongly commute with \(\mathscr G\),
because the coefficients and the matrix \(S_m\) are independent of \(y\).
Applying compact cutoffs in \(i\mathscr G\) and in \(\langle Q\rangle\),
then the standard normal collar core, approximates every compatible energy
state by \(\mathscr C_{\rm comp}\).  Removing \(q=0\) does not change an
\(L^2\) direct integral.  This proves density.

For \(X\in\mathscr C_{\rm comp}\),
\(\mathscr G^\ell U(s)X=U(s)\mathscr G^\ell X\), tangential multipliers
commute with \(U(s)\), and the trace history is translated by \(s\).
Thus, for every \(N\),
\begin{equation}
 p_N(U(s)X)\le C_{N,s}\,p_{N+\lceil|s|\rceil+1}(X).
 \label{eq:V88-core-invariance-bound}
\end{equation}
Completion proves invariance and continuity.

Finally, the seminorms control arbitrarily many time and tangential
derivatives of the phase and port variables.  The local first-order
equations recover the noncharacteristic normal jets; potential
compatibility recovers the metric jet.  Hilbert-valued time trace and the
fractional normal trace theorem used in V83--V86 then put the face and cap
data in \eqref{eq:V86-G-map}--\eqref{eq:V86-R-map}.  The V86 theorem applies
on every finite slab.
\end{proof}

\begin{remark}[Why the trace seminorms are part of the rigging]
\label{rem:V88-trace-seminorms-needed}
Graph powers of \(\mathscr G\) and tangential regularity alone do not assert
an unweighted observation estimate at grazing.  The second line of
\eqref{eq:V88-BFV-seminorms} records the boundary regularity actually used
by the canonical construction.  It is not claimed to be bounded by the
V59 energy norm.
\end{remark}

\subsection{There is no full-energy \(L^2\) port observation}
\label{subsec:V88-no-full-observation}

Let \(\mathcal O_T\) be the strong port observation on the smooth core,
\begin{equation}
 \mathcal O_TX:=\Gamma U(\cdot)X|_{(0,T)}.
 \label{eq:V88-port-observation}
\end{equation}

\begin{theorem}[High-frequency grazing obstruction]
\label{thm:V88-grazing-observation-obstruction}
For every \(T>0\), \(\mathcal O_T\) is not bounded from
\(\mathcal H_{\rm phys}\) to
\(L^2((0,T)\times\mathbb R_y^2)^{\oplus4}\).  More precisely, there are
smooth compatible unit-energy packets \(X_N\) for which
\begin{equation}
 \|\mathcal O_TX_N\|_{L^2_{t,y}}
 \ge c_{T,m}N^{1/2}.
 \label{eq:V88-grazing-port-growth}
\end{equation}
\end{theorem}

\begin{proof}
Use one radiative polarization and an incoming metric half-wave with no
incoming continuum half-wave.  Choose a tangential carrier
\(q=N^2e_1+\eta\), with \(\eta\) in a fixed compact box, and put
\(\xi=N\zeta\), with \(1<\zeta<2\).  On this box the wave frequency and
inward normal group speed satisfy
\begin{equation}
 \omega=(|q|^2+\xi^2)^{1/2}\asymp N^2,
 \qquad
 v_n=\partial_\xi\omega\asymp N^{-1}.
 \label{eq:V88-grazing-scaling}
\end{equation}
Let \(A(\eta,\zeta)\) be a nonzero smooth product cutoff in those boxes and
choose the phase-state Fourier amplitude \(N^{-1/2}A(\eta,\zeta)\).
The Jacobian \(\dd\xi=N\dd\zeta\) makes its bulk \(L^2\) phase norm, and
hence its V59 energy, independent of \(N\).  Translate the rescaled normal
profile so that it meets the interface during \((0,T)\).  At \(x=0\),
after removing the unit-modulus carrier, the incoming trace is
\(N^{1/2}F_N(t,y)\), where the exact phase expansion is
\begin{equation}
 \omega(N^2e_1+\eta,N\zeta)-N^2-\eta_1
 \longrightarrow \frac{\zeta^2}{2}
 \label{eq:V88-paraxial-phase-limit}
\end{equation}
uniformly with all derivatives on the fixed support.  Plancherel in
\(y\) and dominated convergence in \((\eta,\zeta,t)\) give
\begin{equation}
 \int_0^T\|u_{{\rm in},N}(t,0)\|_{L_y^2}^2\,\dd t
 =N\left(C_{A,T}+o(1)\right),
 \qquad C_{A,T}>0,
 \label{eq:V88-transport-trace-growth}
\end{equation}
after choosing the translation so the limiting paraxial trace is nonzero.
This is an exact rescaled Fourier limit, so no geometric-optics remainder
is being assumed.  The exact V67 relation with
\(c_{\rm in}=0\) gives all outgoing ports by the fixed unitary matrix
\(S_m\), so it neither removes nor enlarges the order of the lower bound.
Taking square roots proves \eqref{eq:V88-grazing-port-growth}.
\end{proof}

\begin{firewall}[Generator maximality is not observation admissibility]
\label{fire:V88-maximality-versus-observation}
Theorem~\ref{thm:V68-skew-adjoint-generator} says that the interface graph
gives the maximal skew-adjoint operator domain.  It does not say that the
operator-domain trace extends boundedly to the entire pivot Hilbert space.
Theorem~\ref{thm:V88-grazing-observation-obstruction} separates these two
statements.  V87 used a non-grazing packet and remains valid.
\end{firewall}

The closability proof in \eqref{eq:V88-closability-Green} gives a useful
intermediate Hilbert space.  Define
\begin{equation}
 \mathcal E_{{\rm port},T}
 :=D(\overline{\mathcal O_T}),
 \qquad
 \|X\|_{{\rm port},T}^2
 :=\|X\|_{\mathcal H_{\rm phys}}^2
   +\|\overline{\mathcal O_T}X\|_{L^2_{t,y}}^2.
 \label{eq:V88-closed-port-domain}
\end{equation}

\begin{corollary}[Closed dense trace domain]
\label{cor:V88-closed-trace-domain}
The space \(\mathcal E_{{\rm port},T}\) is a dense Hilbert subspace of
\(\mathcal H_{\rm phys}\), the port trace is bounded for its graph norm,
and its trace lies in \(\mathcal L_{87}\).  The inclusion is strict.
\end{corollary}

\begin{proof}
The observation is densely defined on \(\mathscr C_{\rm comp}\) and
closable by \eqref{eq:V88-closability-Green}.  The domain of a closed
operator with its graph norm is Hilbert and dense because it contains the
dense core.  The matrix relation is preserved under \(L^2\) limits.  If the
domain were all of \(\mathcal H_{\rm phys}\), the closed graph theorem
would make \(\overline{\mathcal O_T}\) bounded, contradicting
Theorem~\ref{thm:V88-grazing-observation-obstruction}.
\end{proof}

\subsection{Weak ports on every energy solution}
\label{subsec:V88-weak-ports}

Let \(\mathscr T_\partial\) be the Frechet space of smooth compactly
supported four-port test functions, with the direct-limit convention in
time and the Schwartz topology in \(y\).  Extend the form
\(\mathfrak B_{87}\) to distribution--test pairs by duality.

\begin{theorem}[Distributional characteristic trace]
\label{thm:V88-weak-characteristic-trace}
There is a unique continuous history-trace map
\begin{equation}
 \boldsymbol\Gamma_{\rm w}:
 \mathcal H_{\rm phys}\longrightarrow\mathscr T_\partial'
 \label{eq:V88-weak-trace-map}
\end{equation}
with the following properties.
\begin{enumerate}[label=\textup{(\roman*)},leftmargin=2.8em]
\item On \(\mathscr D_{\rm BFV}\), it agrees with the strong
      characteristic trace.
\item For every \(X\in\mathcal H_{\rm phys}\) and every port test \(g\),
      \begin{equation}
       \int_{\mathbb R}\mathfrak B_{87}
       \bigl((\boldsymbol\Gamma_{\rm w}X)(t),g(t)\bigr)\,\dd t
       =\int_{\mathbb R}
       \left\langle U(t)X,
       -\partial_tEg-\mathscr G_{\max}^*Eg
       \right\rangle\,\dd t.
       \label{eq:V88-transposition-definition}
      \end{equation}
\item Its values satisfy the physical interface relation in distributions:
      \begin{equation}
       \binom{u_{\rm out}}{c_{\rm out}}
       =(S_m\otimes I_{\mathscr T})
       \binom{u_{\rm in}}{c_{\rm in}}.
       \label{eq:V88-distributional-port-graph}
      \end{equation}
\item On \(\mathcal E_{{\rm port},T}\), its restriction to \((0,T)\)
      is the \(L^2\) trace \(\overline{\mathcal O_T}\).
\end{enumerate}
\end{theorem}

\begin{proof}
For a smooth state, spacetime integration by parts gives
\eqref{eq:V88-transposition-definition}.  The right side obeys
\begin{equation}
 \left|\text{right side of \eqref{eq:V88-transposition-definition}}\right|
 \le \|X\|_{\mathcal H_{\rm phys}}
 \int_{\mathbb R}
 \|\partial_tEg+\mathscr G_{\max}^*Eg\|\,\dd t,
 \label{eq:V88-weak-trace-bound}
\end{equation}
because \(U(t)\) is unitary.  This is a continuous seminorm of \(g\).
The nondegenerate port form and the collar right inverse therefore define
a unique element of \(\mathscr T_\partial'\), independent of the chosen
lift: two lifts differ by a zero-trace test field, whose Green boundary term
vanishes.  Density extends the construction to every energy state and
proves continuity and property (i).

Approximate \(X\) in the pivot by \(X_k\in\mathscr C_{\rm comp}\).  Every
strong trace obeys the V67 matrix relation.  Bound
\eqref{eq:V88-weak-trace-bound} gives distributional convergence, and a
constant matrix preserves that convergence, proving (iii).  If
\(X\in\mathcal E_{{\rm port},T}\), take an approximating sequence in the
graph norm of \(\overline{\mathcal O_T}\).  Its strong traces converge both
in \(L^2\) and in distributions, so the limits agree.  This proves (iv).
\end{proof}

\begin{theorem}[Rigged Green cancellation]
\label{thm:V88-rigged-Green-cancellation}
For \(X\in\mathcal H_{\rm phys}\) and
\(Y\in\mathscr D_{\rm BFV}\), the distribution--test pairing
\begin{equation}
 \int_{\mathbb R}
 \mathfrak B_{87}\bigl(
   (\boldsymbol\Gamma_{\rm w}X)(t),\Gamma U(t)Y
 \bigr)\,\dd t
 \label{eq:V88-rigged-Green-pairing}
\end{equation}
is well defined after a compact time cutoff and is continuous in both
rigged arguments.  It vanishes identically.  Thus the physical Green form
descends to the V87 port graph with one energy leg and one BFV test leg;
no product of two distributional conormals is used.
\end{theorem}

\begin{proof}
Theorem~\ref{thm:V88-weak-characteristic-trace} makes the first boundary
value a distribution, while the seminorms
\eqref{eq:V88-BFV-seminorms} make the second a smooth test trace.  Hence the
pairing exists and the bounds
\eqref{eq:V88-weak-trace-bound} and \eqref{eq:V88-BFV-seminorms} prove
continuity.

Write the incoming distribution of \(X\) as \(\alpha\), the incoming test
trace of \(Y\) as \(\beta\), and put \(S=S_m\otimes I_{\mathscr T}\).
Equation~\eqref{eq:V88-distributional-port-graph} and the strong graph
condition give outgoing traces \(S\alpha\) and \(S\beta\).  Therefore the
integrand is
\begin{equation}
 \frac14\left(
  \langle\alpha,\beta\rangle
  -\langle S\alpha,S\beta\rangle
 \right)=0,
 \label{eq:V88-distributional-neutrality}
\end{equation}
because \(S^*S=I\).  The identity remains valid after any compact time
cutoff.
\end{proof}

\begin{corollary}[Precise relation to the canonical construction]
\label{cor:V88-canonical-relation}
The V86 moment map and master identity remain strong identities on
\(\mathscr D_{\rm BFV}\).  A general V59 energy state is represented in
\(\mathscr D_{\rm BFV}'\), and its dynamical boundary identity is
Theorem~\ref{thm:V88-rigged-Green-cancellation}.  No continuous comparison map
from \(\mathcal H_{\rm phys}\) into the strong quotient
\(\mathcal P_{86}^{\rm red}\) is required or asserted.
\end{corollary}

\begin{proof}
The first assertion is Theorem~\ref{thm:V88-dense-rigging}.  The pivot
embedding in \eqref{eq:V88-rigged-triple} gives the dual representation.
Theorem~\ref{thm:V88-rigged-Green-cancellation} supplies the descended
boundary identity.  The final statement is forced by
Theorem~\ref{thm:V87-V86-gate-fails}.
\end{proof}

\subsection{The cross transverse corner row}
\label{subsec:V88-cross-corner}

Fix the oriented quarter-turn \(J:T^*J\to T^*J\) by declaring that, in the
V63 aligned frame, \(J(D_1)=D_2\).  For the cross amplitude
\(a_\times\), equations \eqref{eq:V63-cross-K-components} and
\eqref{eq:V78-Picor-dressing} give
\begin{equation}
 \varsigma_\times=0,
 \qquad X_\times=0,
 \qquad \Pi_{{\rm cor},\times}^{0}=0,
 \qquad
 \Pi_{{\rm cor},\times}^{A}
 =\frac{\sigma_J}{2\sqrt2}
   J^A{}_B D^B a_\times.
 \label{eq:V88-cross-corner-row}
\end{equation}
Here the sign of \(J\) is fixed by the displayed aligned-frame convention,
so no orientation sign is hidden.

\begin{proposition}[Energy bound and zero pulled-back cross form]
\label{prop:V88-cross-corner-bound}
At either temporal cap,
\begin{align}
 \|\Pi_{{\rm cor},\times}\|_{\dot H_y^{-1/2}}^2
 &=\frac18
   \|a_\times|_\Gamma\|_{\dot H_y^{1/2}}^2,
 \label{eq:V88-cross-corner-norm}\\
 \|\Pi_{{\rm cor},\times}\|_{\dot H_y^{-1/2}}^2
 &\le\frac{1}{4(1+mv)}E_{\rm q}.
 \label{eq:V88-cross-corner-energy-bound}
\end{align}
The cross contribution to the pulled-back V74 corner potential and
two-form is zero.
\end{proposition}

\begin{proof}
The rotation \(J\) is orthogonal and the multiplier \(D_y\) has length
\(|q|\).  Plancherel applied to \eqref{eq:V88-cross-corner-row} gives
\[
 \int |q|^{-1}|\widehat\Pi_{{\rm cor},\times}|^2\,\dd^2q
 =\frac18\int |q||\widehat a_\times|^2\,\dd^2q,
\]
which is \eqref{eq:V88-cross-corner-norm}.  The pointwise V84 common-trace
capacity gives
\((1+mv)\|a_\times|_\Gamma\|_{\dot H^{1/2}}^2\le2E_{\rm q}\);
this proves \eqref{eq:V88-cross-corner-energy-bound}.

The cross corner configuration is identically zero in the quotient gauge,
including along cross variations.  Hence
\(\Pi_{{\rm cor},\times}^A\delta X_A=0\).  There are no time or
longitudinal terms by \eqref{eq:V88-cross-corner-row}.  Its field-space
exterior derivative is therefore zero, agreeing with
Theorem~\ref{thm:V78-pulled-corner-form}.
\end{proof}

\begin{assessment}[The flat two-polarization bridge is now closed]
\label{ass:V88-flat-bridge-status}
On the planar nonzero-frequency linear theory, both polarizations now fit
one coherent structure.  The BFV moment map is defined on the dense
invariant test rigging; the V59 energy is the pivot; every energy solution
has a distributional port satisfying the exact unitary interface graph;
and the cross transverse cap momentum is energy controlled.  What fails is
only the stronger claim that all energy states possess square-integrable
ports or lie in the V86 strong canonical quotient.
\end{assessment}

\subsection{V88 result ledger and next acquisition}
\label{subsec:V88-ledger}

\paragraph{New exact conclusions.}
\begin{enumerate}[label=\textup{(\arabic*)},leftmargin=2.8em]
\item The compatible V68/V59 energy generator has a dense invariant BFV
      test rigging carrying all V86 face and cap maps.
\item The canonical moment map and V75 master identity remain strong on
      that test space, while the energy Hilbert space embeds in its dual.
\item High-frequency grazing packets disprove an unweighted
      energy-to-\(L^2\)-port observation bound.
\item The strong port observation is closable and has a dense, strict
      Hilbert graph domain.
\item Every energy solution has a unique distributional characteristic
      trace defined by transposition.
\item The exact V67 matrix relation persists in distributions, and unitary
      neutrality gives the rigged Green cancellation with one test leg.
\item The cross transverse corner momentum is exactly
      \(\sigma_J J D a_\times/(2\sqrt2)\), is bounded in
      \(\dot H_y^{-1/2}\) by the V59 energy, and contributes zero to the
      pulled-back corner two-form.
\end{enumerate}

\paragraph{Next forced step.}
The remaining obstruction is no longer the flat radiative boundary
topology.  V89 should move upstream to the curved Einstein identity which
must make the interface stress compatible with the subsidiary system.  It
must derive, from the GHY--de Donder variational formula and the contracted
Bianchi identity, the exact boundary divergence law for the Brown--York
momentum in the presence of a continuum traction.  Only after that identity
is fixed can a Standard Model stress tensor be inserted without creating an
independent boundary constraint.  The derivation should separate geometric
identities from analytic trace assumptions and should state the Sobolev
orders needed for the rigged pairing.

\begin{assessment}[V88 continuum status]
\label{ass:V88-continuum-status}
The full Standard--Model--Einstein continuum remains unproved.  V88 closes
the flat linear two-polarization BFV/energy bridge in the rigged sense and
identifies the exact limit of strong port regularity.  Curved nonlinear
gluing, the Brown--York/Bianchi stress identity, a common renormalized
Standard Model stress domain, and the cofinal interacting continuum remain
open.
\end{assessment}

\section*{V88 exact checks and append-only record}
\addcontentsline{toc}{section}{V88 exact checks and append-only record}

\begin{equation}
\boxed{
\begin{gathered}
 \mathscr D_{\rm BFV}\subset\mathcal H_{\rm phys}
 \subset\mathscr D_{\rm BFV}',\\
 \boldsymbol\Gamma_{\rm w}\mathcal H_{\rm phys}\subset\mathscr T_\partial',
 \qquad
 \boldsymbol\Gamma_{{\rm w,out}}=S_m\boldsymbol\Gamma_{{\rm w,in}},\\
 \mathfrak B_{87}(\boldsymbol\Gamma_{\rm w}X,\Gamma Y)=0,
 \qquad X\in\mathcal H_{\rm phys},\quad Y\in\mathscr D_{\rm BFV},\\
 \Pi_{{\rm cor},\times}^{A}
 =\frac{\sigma_J}{2\sqrt2}J^A{}_BD^Ba_\times,
 \qquad
 \|\Pi_{{\rm cor},\times}\|_{\dot H^{-1/2}}^2
 \le\frac{E_{\rm q}}{4(1+mv)}.
\end{gathered}}
\label{eq:V88-check-card}
\end{equation}

The V87 predecessor used for this append has SHA--256
\begin{center}
\texttt{cb26d8122bf3a19853fc1fababec3d1c3c58068b98469731af1a895a0c3b77a0}.
\end{center}
The V88 candidate retains its bytes up to the sole live final document
terminator, appends this material, and restores exactly one active
\verb|\end{document}|.  Validation includes balanced environments and
braces, unique labels, resolved new internal references, valid inline and
display math delimiters, exact predecessor-prefix equality, and a final
inventory which removes only the superseded SM V87 file.

% =============================================================================
% END APPENDED VERSION V88 MATERIAL
% =============================================================================


\clearpage
% =============================================================================
% BEGIN APPENDED VERSION V89 MATERIAL
% =============================================================================
\section{V89: curved Brown--York--Codazzi closure, the compulsory clock,
and the exact traction-equivariance Ward identity}
\label{sec:V89-curved-Ward}

V88 closed the planar radiative boundary problem as a rigged
BFV--energy--port system.  The next question is geometric rather than
topological.  A continuum traction is a covector in the two-polarization
field space; the contracted Bianchi identity is a spacetime covector
identity.  Equating those objects without a metric-variation incidence is
ill typed and reproduces the V58/V60 subsidiary source.

This note derives the missing identity in three steps.  First, Codazzi gives
the exact divergence of the Brown--York momentum.  Second, a covariant
auxiliary-cylinder action gives both the V59 traction and a boundary stress.
The speed \(0<v<1\) forces a timelike clock field into that action.  Third,
the common-trace map \(z=\mathcal A(\gamma)\) converts the field-space
traction into a metric stress through the adjoint of
\(D\mathcal A\).  Diffeomorphism equivariance of \(\mathcal A\) is exactly
the identity which cancels the traction term in the boundary Ward law.

The result does not assume a nonlinear curved TT extractor exists.  It
proves a necessary and sufficient Ward gate for any proposed extractor and
exhibits its failure as an explicit distributional source.  This separates
the geometric identity from the analytic construction still required.

\subsection{Conventions and the exact Codazzi row}
\label{subsec:V89-Codazzi}

Let \((M,g)\) be a four-dimensional Lorentzian manifold with smooth
timelike boundary \(\Gamma\).  Let \(n\) be the outward spacelike unit
normal, \(n^\mu n_\mu=1\), and put
\begin{equation}
 \gamma_{\mu\nu}:=g_{\mu\nu}-n_\mu n_\nu,
 \qquad
 K_{ab}:=\gamma_a{}^\mu\gamma_b{}^\nu\nabla_\mu n_\nu,
 \qquad
 \Pi^{ab}:=K^{ab}-K\gamma^{ab}.
 \label{eq:V89-BY-conventions}
\end{equation}
These are the nonlinear conventions whose flat linearization is
\eqref{eq:V36-linear-K}--\eqref{eq:V36-linear-BY}.

\begin{theorem}[Brown--York--Codazzi divergence identity]
\label{thm:V89-BY-Codazzi}
For every smooth metric with the conventions
\eqref{eq:V89-BY-conventions},
\begin{equation}
 \boxed{
 D_a\Pi^a{}_b
 =n^\mu\gamma_b{}^\nu G_{\mu\nu}.}
 \label{eq:V89-Codazzi-row}
\end{equation}
If the bulk Einstein equation is normalized as
\begin{equation}
 G_{\mu\nu}+\Lambda g_{\mu\nu}=\kappa T_{\mu\nu}^{\rm bulk},
 \label{eq:V89-Einstein-normalization}
\end{equation}
then
\begin{equation}
 D_a\Pi^a{}_b
 =\kappa n^\mu\gamma_b{}^\nu T_{\mu\nu}^{\rm bulk}.
 \label{eq:V89-BY-matter-flux}
\end{equation}
If \(\kappa\) is absorbed into the manuscript's stress convention, the
same statement holds with the displayed \(\kappa\) removed.
\end{theorem}

\begin{proof}
The contracted Codazzi equation for
\(K_{ab}=\gamma_a{}^\mu\gamma_b{}^\nu\nabla_\mu n_\nu\) is
\begin{equation}
 D_aK^a{}_b-D_bK
 =n^\mu\gamma_b{}^\nu R_{\mu\nu}.
 \label{eq:V89-contracted-Codazzi}
\end{equation}
Metric compatibility gives
\(D_a(K^a{}_b-K\delta^a_b)=D_a\Pi^a{}_b\).  Moreover
\(n^\mu\gamma_b{}^\nu g_{\mu\nu}=0\), so replacing the Ricci tensor by the
Einstein tensor does not change the mixed projection.  This proves
\eqref{eq:V89-Codazzi-row}.  The cosmological term in
\eqref{eq:V89-Einstein-normalization} has the same vanishing mixed
projection, proving \eqref{eq:V89-BY-matter-flux}.
\end{proof}

\begin{remark}[Orientation check]
\label{rem:V89-orientation-check}
Reversing \(n\) reverses \(K\), \(\Pi\), and the mixed bulk flux, so both
sides of \eqref{eq:V89-Codazzi-row} change sign.  The identity is therefore
compatible with the orientation convention used in the V36 and V78 cap
formulas.
\end{remark}

\subsection{Why a subluminal continuum needs a clock}
\label{subsec:V89-clock-necessity}

Let \(\Theta\) be a boundary clock with timelike gradient and define
\begin{equation}
 u_a:=-\frac{D_a\Theta}{\sqrt{-D_c\Theta D^c\Theta}},
 \qquad u_au^a=-1,
 \qquad s^{ab}:=\gamma^{ab}+u^au^b.
 \label{eq:V89-clock-frame}
\end{equation}
For \(0<v<1\), put
\begin{equation}
 h_v^{ab}:=-u^au^b+v^2s^{ab}
 =v^2\gamma^{ab}+(v^2-1)u^au^b.
 \label{eq:V89-effective-cone}
\end{equation}
At the flat clock \(\Theta=t\), this tensor is
\(\operatorname{diag}(-1,v^2,v^2)\).

\begin{proposition}[The clock is compulsory when \(v\ne1\)]
\label{prop:V89-clock-compulsory}
No diffeomorphism-covariant symmetric contravariant two-tensor constructed
pointwise from \(\gamma\) alone has, at a Lorentzian flat point, principal
cone \(-\omega^2+v^2|q|^2=0\) with \(0<v<1\).  An additional timelike datum,
such as \(u\) in \eqref{eq:V89-clock-frame}, is necessary.
\end{proposition}

\begin{proof}
At a flat point, naturality under the stabilizer of \(\gamma\) makes any
such rank-two tensor Lorentz invariant.  The only Lorentz-invariant element
of \(\operatorname{Sym}^2(T\Gamma)\) is a scalar multiple of
\(\gamma^{ab}\).  Its characteristic cone is the light cone of \(\gamma\),
independent of the scalar.  The cone with speed \(v\ne1\) has a preferred
timelike eigendirection and is not Lorentz invariant.  Formula
\eqref{eq:V89-effective-cone} supplies precisely that additional direction.
\end{proof}

\begin{firewall}[The V59 parameter \(v\) carries geometric data]
\label{fire:V89-v-clock-scope}
The fixed-background V59 theory may treat \(v\) as a constant coefficient.
A curved diffeomorphism-covariant completion must also specify the clock
which distinguishes time from the two tangential spatial directions.  If
the clock is held external, its Euler term remains in the Ward identity.  A
closed Einstein system must make the clock dynamical or prove that the
corresponding source vanishes.
\end{firewall}

\subsection{A covariant cylinder action and its traction}
\label{subsec:V89-cylinder-action}

Let \(\mathscr T\) be the real rank-two radiative bundle with its fibre
metric, and let
\(\varphi:\Gamma\times\mathbb R_r^+\to\mathscr T\).  Parallel transport in
\(\mathscr T\) is understood in every derivative.  Consider
\begin{equation}
 S_c[\gamma,\Theta,\varphi]
 :=-\frac m2\int_\Gamma\sqrt{|\gamma|}\int_0^\infty
 \left(
  |\partial_r\varphi|^2
  +h_v^{ab}\langle D_a\varphi,D_b\varphi\rangle
 \right)\dd r.
 \label{eq:V89-cylinder-action}
\end{equation}
Assume decay at \(r=\infty\), set \(z:=\varphi|_{r=0}\), and define
\begin{equation}
 t_b:=-m\partial_r\varphi|_{r=0}.
 \label{eq:V89-cylinder-traction}
\end{equation}

\begin{proposition}[Euler equation and exact V59 flat limit]
\label{prop:V89-cylinder-Euler}
Variation of \(\varphi\) gives
\begin{equation}
 \mathcal E_\varphi
 :=m\left[
  \partial_r^2\varphi+D_a(h_v^{ab}D_b\varphi)
 \right]=0
 \label{eq:V89-cylinder-Euler}
\end{equation}
and the boundary one-form
\begin{equation}
 -\int_\Gamma\sqrt{|\gamma|}\,
 \langle t_b,\delta z\rangle.
 \label{eq:V89-cylinder-boundary-one-form}
\end{equation}
At \(\gamma=\operatorname{diag}(-1,1,1)\), \(\Theta=t\), and constant
bundle frame, \eqref{eq:V89-cylinder-Euler} is
\begin{equation}
 \partial_t^2\varphi
 =\partial_r^2\varphi+v^2\Delta_y\varphi,
 \label{eq:V89-V59-flat-cylinder}
\end{equation}
and \eqref{eq:V89-cylinder-traction} is the V59 traction.
\end{proposition}

\begin{proof}
Integrate the \(r\)- and tangential derivative terms in
\(\delta_\varphi S_c\) once.  The outward normal of the auxiliary half-line
at \(r=0\) is \(-\partial_r\), so its coefficient is
\(m\partial_r\varphi=-t_b\), proving
\eqref{eq:V89-cylinder-boundary-one-form}.  The interior coefficient is
\eqref{eq:V89-cylinder-Euler}.  Substitution of
\eqref{eq:V89-effective-cone} at the flat clock gives
\(\partial_r^2\varphi-\partial_t^2\varphi+v^2\Delta_y\varphi=0\), which is
\eqref{eq:V89-V59-flat-cylinder}.
\end{proof}

Define the fixed-trace cylinder stress and the integrated clock Euler
coefficient by the on-shell-in-\(\varphi\) variation
\begin{equation}
 \delta S_c
 =\frac12\int_\Gamma\sqrt{|\gamma|}\,
    \tau_c^{ab}\delta\gamma_{ab}
  +\int_\Gamma\sqrt{|\gamma|}\,
    \mathcal E_\Theta\delta\Theta
  -\int_\Gamma\sqrt{|\gamma|}\,
    \langle t_b,\delta z\rangle.
 \label{eq:V89-cylinder-variation}
\end{equation}
The adjective ``fixed-trace'' means that \(z\) is varied independently in
this formula.

\begin{theorem}[Cylinder Ward identity]
\label{thm:V89-cylinder-Ward}
If \(\mathcal E_\varphi=0\), diffeomorphism invariance of
\eqref{eq:V89-cylinder-action} gives
\begin{equation}
 \boxed{
 D_a\tau_c^a{}_b
 =\mathcal E_\Theta D_b\Theta
  -\langle t_b,D_bz\rangle.}
 \label{eq:V89-cylinder-Ward}
\end{equation}
This is an off-clock-shell identity; no clock equation has been used.
\end{theorem}

\begin{proof}
Let \(\zeta\) be a compactly supported tangential vector field.  Under its
infinitesimal diffeomorphism,
\begin{equation*}
 \delta_\zeta\gamma_{ab}=2D_{(a}\zeta_{b)},
 \qquad
 \delta_\zeta\Theta=\zeta^aD_a\Theta,
 \qquad
 \delta_\zeta z=\zeta^aD_az.
\end{equation*}
Insert these variations in \eqref{eq:V89-cylinder-variation}.  Invariance
and one tangential integration by parts give
\[
 0=\int_\Gamma\sqrt{|\gamma|}\,\zeta^b
 \left[-D_a\tau_c^a{}_b
       +\mathcal E_\Theta D_b\Theta
       -\langle t_b,D_bz\rangle\right].
\]
Arbitrariness of \(\zeta\) proves \eqref{eq:V89-cylinder-Ward}.
\end{proof}

\subsection{The common-trace incidence and its equivariance defect}
\label{subsec:V89-equivariant-incidence}

Let \(\mathcal A\) be a proposed curved radiative trace map from induced
metrics to sections of \(\mathscr T\), and impose
\begin{equation}
 z=\mathcal A(\gamma).
 \label{eq:V89-common-trace-map}
\end{equation}
For a field-space covector \(t\), define the metric adjoint
\((D\mathcal A_\gamma)^\dagger t\) by
\begin{equation}
 \int_\Gamma\sqrt{|\gamma|}\,
 \langle t,D\mathcal A_\gamma[h]\rangle
 =\frac12\int_\Gamma\sqrt{|\gamma|}\,
 \bigl((D\mathcal A_\gamma)^\dagger t\bigr)^{ab}h_{ab}.
 \label{eq:V89-incidence-adjoint}
\end{equation}
The factor \(1/2\) matches the stress convention.

Define the infinitesimal equivariance defect by
\begin{equation}
 \mathfrak d_{\mathcal A}(\zeta)
 :=D\mathcal A_\gamma[\mathcal L_\zeta\gamma]
   -\mathcal L_\zeta z,
 \qquad z=\mathcal A(\gamma),
 \label{eq:V89-equivariance-defect}
\end{equation}
and its distributional adjoint by
\begin{equation}
 \langle\mathfrak d_{\mathcal A}^*t,\zeta\rangle
 :=\int_\Gamma\sqrt{|\gamma|}\,
 \langle t,\mathfrak d_{\mathcal A}(\zeta)\rangle.
 \label{eq:V89-equivariance-defect-adjoint}
\end{equation}

\begin{lemma}[Exact adjoint-incidence divergence]
\label{lem:V89-incidence-divergence}
For every smooth \(t\),
\begin{equation}
 D_a\bigl((D\mathcal A_\gamma)^\dagger t\bigr)^a{}_b
 =-\langle t,D_bz\rangle
  -(\mathfrak d_{\mathcal A}^*t)_b.
 \label{eq:V89-incidence-divergence}
\end{equation}
In particular, the traction incidence has the required divergence exactly
when \(\mathcal A\) is infinitesimally diffeomorphism-equivariant.
\end{lemma}

\begin{proof}
Put \(h=\mathcal L_\zeta\gamma\) in
\eqref{eq:V89-incidence-adjoint}.  Symmetry of the adjoint tensor and one
integration by parts give
\begin{align*}
 \int\langle t,D\mathcal A_\gamma[\mathcal L_\zeta\gamma]\rangle
 &=\int (D\mathcal A_\gamma)^\dagger t^{ab}D_a\zeta_b\\
 &=-\int D_a\bigl((D\mathcal A_\gamma)^\dagger t\bigr)^a{}_b\zeta^b.
\end{align*}
On the other hand, \eqref{eq:V89-equivariance-defect} makes the left side
equal to
\[
 \int\langle t,D_bz\rangle\zeta^b
 +\langle\mathfrak d_{\mathcal A}^*t,\zeta\rangle.
\]
Equality for arbitrary \(\zeta\) proves
\eqref{eq:V89-incidence-divergence}.
\end{proof}

After imposing \eqref{eq:V89-common-trace-map}, the chain rule in
\eqref{eq:V89-cylinder-variation} gives the effective continuum stress
\begin{equation}
 \tau_{\rm eff}^{ab}
 :=\tau_c^{ab}
   -\bigl((D\mathcal A_\gamma)^\dagger t_b\bigr)^{ab}.
 \label{eq:V89-effective-stress}
\end{equation}

\begin{theorem}[Traction-equivariance Ward closure]
\label{thm:V89-traction-Ward-closure}
On the cylinder Euler equation,
\begin{equation}
 \boxed{
 D_a\tau_{\rm eff}^a{}_b
 =\mathcal E_\Theta D_b\Theta
  +(\mathfrak d_{\mathcal A}^*t_b)_b.}
 \label{eq:V89-effective-Ward}
\end{equation}
Therefore the continuum contribution is divergence free on the clock shell
if and only if its trace map has zero equivariance defect against the
attained traction covectors.  The field-space term
\(-\langle t_b,D_bz\rangle\) is not an additional boundary constraint; it
is cancelled by the divergence of the adjoint metric incidence.
\end{theorem}

\begin{proof}
Take the divergence of \eqref{eq:V89-effective-stress}.  Substitute
Theorem~\ref{thm:V89-cylinder-Ward} and
Lemma~\ref{lem:V89-incidence-divergence}.  The two copies of
\(\langle t_b,D_bz\rangle\) cancel and leave
\eqref{eq:V89-effective-Ward}.
\end{proof}

\begin{corollary}[Exact nonlinear Bianchi compatibility gate]
\label{cor:V89-Bianchi-gate}
Suppose a free induced-metric variation gives the interface equation
\begin{equation}
 \Pi^{ab}+\kappa\bigl(
   \tau_{\rm eff}^{ab}+\tau_{\rm SM}^{ab}
 \bigr)=0,
 \label{eq:V89-interface-metric-equation}
\end{equation}
where \(\tau_{\rm SM}\) is any additional boundary Standard Model stress.
Then the divergence of this equation is exactly
\begin{align}
 0={}&n^\mu\gamma_b{}^\nu T_{\mu\nu}^{\rm bulk}
 +\mathcal E_\Theta D_b\Theta
 +(\mathfrak d_{\mathcal A}^*t_b)_b
 +D_a\tau_{\rm SM}^a{}_b.
 \label{eq:V89-full-boundary-Ward-gate}
\end{align}
Thus a homogeneous Einstein subsidiary boundary row follows from the bulk,
clock, continuum, and Standard Model Euler equations precisely when their
normal/tangential fluxes obey \eqref{eq:V89-full-boundary-Ward-gate}; no
independent traction source may be inserted.
\end{corollary}

\begin{proof}
Apply \(D_a\) to \eqref{eq:V89-interface-metric-equation}, use
\eqref{eq:V89-BY-matter-flux} and
\eqref{eq:V89-effective-Ward}, and divide by \(\kappa\).
\end{proof}

\begin{firewall}[What the flat V59 action does not yet provide]
\label{fire:V89-missing-curved-extractor}
The V59 action proves the field equations, common trace, traction balance,
positive energy, and the exact port graph on a fixed flat background.  It
does not define the nonlinear map \(\mathcal A(\gamma)\), its metric
adjoint, or its equivariance defect, and it does not supply the clock Euler
equation required by \(v<1\).  Consequently
\eqref{eq:V89-effective-Ward} is a proved gate, not a claim that the curved
continuum has already passed it.
\end{firewall}

\subsection{Relation to the V58/V60 source and the V63 conormal}
\label{subsec:V89-prior-source-relation}

The old V41 incidence fed a two-polarization traction directly into five
metric residual rows.  V58 and V60 computed a nonzero harmonic source
\(D_{\rm rad}\).  In the present notation, such a source is the frozen
symbol of a nonzero Ward defect: a raw field-space traction was used without
the adjoint of an equivariant metric trace map.

By contrast, V63 derives its radiative conormal by restricting the single
GHY--de Donder one-form to the physical induced-metric graph.  At the flat
anchor, this is precisely the chain-rule mechanism in
\eqref{eq:V89-incidence-adjoint}.  V63 proves that the resulting two
polarizations both have equation-normalized conormal \(D_na\).  It does not
construct a nonlinear curved \(\mathcal A\), so no earlier theorem implies
\(\mathfrak d_{\mathcal A}=0\) away from the anchor.

\begin{proposition}[Linear anchor consistency]
\label{prop:V89-linear-anchor-consistency}
At the planar harmonic anchor, restricting the GHY variation through the
V63 physical trace graph gives the adjoint traction incidence of
\eqref{eq:V89-incidence-adjoint}.  Hence its virtual work is
\begin{equation}
 \frac1{2\kappa}\int_\Gamma
 \langle D_na,\delta a\rangle_{\mathscr T},
 \label{eq:V89-flat-virtual-work}
\end{equation}
and the continuum boundary variation is its opposite when
\(p_{\rm rad}+t_b=0\).  The V59 port power cancellation is the time
derivative of this common-action pairing on smooth solutions.
\end{proposition}

\begin{proof}
Proposition~\ref{prop:V63-conormal-from-action} identifies the restricted
GHY one-form, and Theorem~\ref{thm:V63-exact-real-cone-lift} gives
\(p_{\rm rad}=D_na\).  Proposition~\ref{prop:V89-cylinder-Euler} gives the
continuum boundary coefficient \(-t_b\delta z\).  Common trace makes
\(\delta z=\delta a\), while traction balance makes the two coefficients
opposite.  Differentiating along a solution gives the V67 characteristic
power identity.
\end{proof}

\subsection{Strong and rigged analytic meanings}
\label{subsec:V89-analytic-orders}

\begin{proposition}[A sufficient strong Sobolev regime]
\label{prop:V89-Sobolev-regime}
Let \(g\in H^s(M)\) with \(s>7/2\), uniformly nondegenerate up to the
boundary, and suppose the bulk Einstein equation holds with
\(T^{\rm bulk}\in H^{s-2}(M)\).  Then
\begin{equation}
 \gamma\in H^{s-1/2}(\Gamma),
 \quad
 \Pi\in H^{s-3/2}(\Gamma),
 \quad
 D\Pi, nT^{\rm bulk}\in H^{s-5/2}(\Gamma),
 \label{eq:V89-strong-orders}
\end{equation}
and \eqref{eq:V89-Codazzi-row} holds in
\(H^{s-5/2}(\Gamma)\).  If the first variations of
\(S_c\) and \(\mathcal A\) define continuous covectors at these orders,
then \eqref{eq:V89-effective-Ward} and
\eqref{eq:V89-full-boundary-Ward-gate} hold in the same distribution space.
\end{proposition}

\begin{proof}
The trace theorem gives the first three inclusions in
\eqref{eq:V89-strong-orders}: one normal derivative is lost in \(K\), and
one tangential derivative is lost in its divergence.  Since \(s>7/2\), the
boundary coefficient spaces used in the connection and product terms are
Sobolev algebras, so the nonlinear contractions are continuous.  The
curvature trace has order \(s-5/2\).  Smooth approximation proves Codazzi at
that order.  The remaining identities are variational distribution
identities; continuity of the stated first variations permits the same
density passage.
\end{proof}

\begin{corollary}[Rigged formulation below the strong regime]
\label{cor:V89-rigged-Ward}
On the V88 rigging, equations \eqref{eq:V89-cylinder-Ward} and
\eqref{eq:V89-effective-Ward} remain meaningful with one physical energy
leg and one \(\mathscr D_{\rm BFV}\) test variation.  The traction and
adjoint-incidence terms are then defined by transposition.  No pointwise
Brown--York conormal trace of an arbitrary energy solution is required.
\end{corollary}

\begin{proof}
The definition \eqref{eq:V89-incidence-adjoint} is already a test-dual
definition.  The weak characteristic trace and its Green pairing are
continuous by Theorem~\ref{thm:V88-weak-characteristic-trace}.  Insert a
BFV test diffeomorphism in the proofs of
Theorem~\ref{thm:V89-cylinder-Ward} and
Lemma~\ref{lem:V89-incidence-divergence}; every integration by parts is then
a distribution--test duality.  Density of the smooth compatible core gives
the asserted extension.
\end{proof}

\subsection{V89 result ledger and next acquisition}
\label{subsec:V89-ledger}

\paragraph{New exact conclusions.}
\begin{enumerate}[label=\textup{(\arabic*)},leftmargin=2.8em]
\item The nonlinear Brown--York divergence is exactly the mixed Einstein
      tensor, with orientation and stress normalization fixed.
\item A subluminal cylinder cone with \(0<v<1\) requires a timelike clock;
      the metric alone cannot define it covariantly.
\item The covariant cylinder action reproduces the V59 wave equation and
      traction at the flat clock.
\item Its fixed-trace Ward law contains exactly
      \(-\langle t_b,D_bz\rangle\) and the clock Euler source.
\item The adjoint of the common-trace metric variation has divergence
      \(-\langle t_b,D_bz\rangle\) precisely for an equivariant trace map.
\item The two traction terms cancel in the effective continuum stress; the
      only remaining sources are the clock equation and the explicit
      equivariance defect.
\item Combining this Ward law with Codazzi gives the exact nonlinear
      boundary Bianchi gate \eqref{eq:V89-full-boundary-Ward-gate}.
\item The identities have both a sufficient strong Sobolev realization and
      a V88 distribution--test realization.
\end{enumerate}

\paragraph{Next forced step and scheduled audit.}
V90 is the next five-version audit.  It must first inspect the newest NS and
Yang--Mills notes in the shared folder.  It should then build the Standard
Model side of \eqref{eq:V89-full-boundary-Ward-gate}: one common
gauge--Higgs--fermion stress convention, its off-shell diffeomorphism Ward
identity including gauge Euler terms, and the Sobolev/renormalized domain on
which its normal flux pairs with the V88 rigging.  Independently, the curved
radiative extractor must eventually supply a quantitative bound for
\(\mathfrak d_{\mathcal A}\), with exact vanishing as the target.

\begin{assessment}[V89 continuum status]
\label{ass:V89-continuum-status}
The full Standard--Model--Einstein continuum remains unproved.  V89 derives
the curved Brown--York/Codazzi and continuum Ward identities and isolates
the exact nonlinear equivariance gate.  A dynamical clock, a curved
equivariant radiative trace map, the common renormalized Standard Model
stress domain, nonlinear existence, and the cofinal interacting limit
remain open.
\end{assessment}

\section*{V89 exact checks and append-only record}
\addcontentsline{toc}{section}{V89 exact checks and append-only record}

\begin{equation}
\boxed{
\begin{gathered}
 D_a\Pi^a{}_b=n^\mu\gamma_b{}^\nu G_{\mu\nu},\\
 D_a\tau_c^a{}_b
 =\mathcal E_\Theta D_b\Theta-\langle t_b,D_bz\rangle,\\
 D_a\bigl((D\mathcal A)^\dagger t_b\bigr)^a{}_b
 =-\langle t_b,D_bz\rangle
  -(\mathfrak d_{\mathcal A}^*t_b)_b,\\
 D_a\tau_{\rm eff}^a{}_b
 =\mathcal E_\Theta D_b\Theta
  +(\mathfrak d_{\mathcal A}^*t_b)_b.
\end{gathered}}
\label{eq:V89-check-card}
\end{equation}

The V88 predecessor used for this append has SHA--256
\begin{center}
\texttt{e56bfa0b20945052c9a7430dd823997fe06944866406487c8629543a7897c339}.
\end{center}
The V89 candidate retains its bytes up to the sole live final document
terminator, appends this material, and restores exactly one active
\verb|\end{document}|.  Validation includes balanced environments and
braces, unique labels, resolved new internal references, valid inline and
display math delimiters, exact predecessor-prefix equality, and a final
inventory which removes only the superseded SM V88 file.

% =============================================================================
% END APPENDED VERSION V89 MATERIAL
% =============================================================================


\clearpage
% =============================================================================
% BEGIN APPENDED VERSION V90 MATERIAL
% =============================================================================
\section{V90: compulsory NS/YM audit, the Brown--York sign correction,
and a common Standard Model Ward domain}
\label{sec:V90-SM-Ward-domain}

V90 begins with the scheduled inspection of the two companion programmes.
It then corrects one orientation sign in V89 by deriving the Codazzi row
directly from the already fixed V36 action variation.  With that correction,
the Standard Model boundary Ward identity has exactly the sign required to
cancel the mixed bulk stress in the Einstein boundary equation.

The main construction is a representation-independent Standard Model Ward
bank.  Gauge, Higgs, and fermion Euler covectors are kept together until
after the gauge and diffeomorphism identities are formed.  This produces a
single bulk stress and a single boundary flux identity.  A final theorem
states sufficient common-domain and renormalized-form hypotheses under which
the same identity survives removal of a quantum regulator and pairs with the
V88 BFV rigging.  Those hypotheses are explicit gates; the present note does
not infer them from sectorwise finite values.

\subsection{The V90 companion audit}
\label{subsec:V90-companion-audit}

The shared-folder snapshots were read twice, before and after the audit.
They were stable and had the following exact inventory:
\begin{center}
\begin{tabular}{@{}p{0.49\linewidth}r p{0.31\linewidth}@{}}
\toprule
file & bytes & SHA--256\\
\midrule
\texttt{Renewal\_NS\_Stability\_Research\_Notes\_V31.tex}
&887537&\texttt{f1121b62238ad5491bb7cf03635ea9934942edf573ed8faaa9ee79e1d38a6696}\\
\texttt{Renewal\_Yang\_Mills\_Mass\_Gap\_Research\_Notes\_V88.tex}
&2150040&\texttt{4b76be6cb95e3c194528c333eb3467adfd35ccf00378213e9d05eb8ba1ab0744}\\
\bottomrule
\end{tabular}
\end{center}

The NS file is byte-for-byte the V31 snapshot reviewed at V85.  It therefore
adds no new result at this checkpoint.  No NS theorem is reimported.

The YM file advanced from V82 to V88.  Its new exact result relevant to proof
order is the balanced ternary multiplicity
\begin{equation}
 m_{R_m,R_m,Q_{2m}}(T_k)=k+1,
 \qquad 0\le k\le m.
 \label{eq:V90-YM-multiplicity-audit}
\end{equation}
After retaining that full multiplicity space, its normalized thermal trace
is bounded below by a positive constant and above by \(C\sqrt m\).  The
remaining factor is an ancestry average, not an output-count error.  Thus YM
V88 supplies no continuum or mass-gap theorem for import into the present
programme.

\begin{proposition}[Transferable audit rule]
\label{prop:V90-audit-rule}
In the Standard Model stress construction, regulator traces and counterterm
limits must be taken only after all gauge representations, generations,
chiralities, and common Euler rows have been assembled in one Ward
functional.  A sectorwise trace estimate cannot certify the common Ward
identity unless its counterterms are restrictions of one invariant action.
\end{proposition}

\begin{proof}
Equation~\eqref{eq:V90-YM-multiplicity-audit} is an exact example in which
choosing one ancestry vector before tracing misses a factor growing with the
representation label.  The Standard Model Ward identity below is linear in
the complete first variation, whereas stress renormalization is performed on
quadratic composite fields.  Summing restrictions of one invariant
counterterm action preserves the first-variation identity.  Independently
renormalized sector stresses need not have counterterms whose gauge and
diffeomorphism variations cancel.  Hence the asserted proof order is
necessary for the deduction used here.
\end{proof}

\subsection{Append-only correction of the V89 Codazzi sign}
\label{subsec:V90-Codazzi-correction}

V89 printed
\(D_a\Pi^a{}_b=+n^\mu\gamma_b{}^\nu G_{\mu\nu}\).  That sign is
incompatible with the V36 variational convention, which is the convention
governing this manuscript.  The correction is recorded here without
altering the V89 bytes.

\begin{theorem}[Action-fixed Brown--York divergence sign]
\label{thm:V90-corrected-Codazzi}
With the V36 first variation
\begin{equation}
 \delta S_D
 =-\frac1{2\kappa}\int_M\sqrt{-g}\,
     G^{\mu\nu}\delta g_{\mu\nu}
  +\frac1{2\kappa}\int_\Gamma\sqrt{|\gamma|}\,
     \Pi^{ab}\delta\gamma_{ab},
 \label{eq:V90-V36-variation-recalled}
\end{equation}
the exact tangential identity is
\begin{equation}
 \boxed{
 D_a\Pi^a{}_b
 =-n^\mu\gamma_b{}^\nu G_{\mu\nu}.}
 \label{eq:V90-corrected-Codazzi}
\end{equation}
Consequently, under
\(G_{\mu\nu}+\Lambda g_{\mu\nu}=\kappa T_{\mu\nu}^{\rm bulk}\),
\begin{equation}
 D_a\Pi^a{}_b
 =-\kappa n^\mu\gamma_b{}^\nu T_{\mu\nu}^{\rm bulk}.
 \label{eq:V90-corrected-matter-flux}
\end{equation}
\end{theorem}

\begin{proof}
Let \(\zeta\) be tangent to \(\Gamma\), compactly supported away from its
temporal corners, and put \(\delta g=\mathcal L_\zeta g\).  Diffeomorphism
invariance makes \eqref{eq:V90-V36-variation-recalled} zero.  The contracted
Bianchi identity and one integration by parts give
\begin{align*}
 -\int_M G^{\mu\nu}\nabla_\mu\zeta_\nu
 &=-\int_\Gamma n^\mu G_{\mu\nu}\zeta^\nu,\\
 \int_\Gamma\Pi^{ab}D_a\zeta_b
 &=-\int_\Gamma(D_a\Pi^a{}_b)\zeta^b.
\end{align*}
The common factor is nonzero, so arbitrariness of \(\zeta\) gives
\eqref{eq:V90-corrected-Codazzi}.  The cosmological term has zero mixed
normal--tangential projection, proving
\eqref{eq:V90-corrected-matter-flux}.
\end{proof}

\begin{corollary}[Corrected V89 boundary gate]
\label{cor:V90-corrected-V89-gate}
The V89 cylinder Ward identity, incidence-divergence identity, and effective
stress identity are unchanged.  Only equations
\eqref{eq:V89-Codazzi-row}, \eqref{eq:V89-BY-matter-flux},
\eqref{eq:V89-full-boundary-Ward-gate}, and the first line of
\eqref{eq:V89-check-card} require the orientation correction.  The corrected
full gate is
\begin{equation}
 \boxed{
 0=-n^\mu\gamma_b{}^\nu T_{\mu\nu}^{\rm bulk}
 +\mathcal E_\Theta D_b\Theta
 +(\mathfrak d_{\mathcal A}^*t_b)_b
 +D_a\tau_{\rm SM}^a{}_b.}
 \label{eq:V90-corrected-full-gate}
\end{equation}
All later uses in this manuscript refer to
\eqref{eq:V90-corrected-Codazzi}--\eqref{eq:V90-corrected-full-gate}.
\end{corollary}

\begin{proof}
The cylinder and incidence proofs use only tangential diffeomorphism
invariance of the auxiliary action and do not invoke Codazzi.  Taking the
divergence of the interface equation
\eqref{eq:V89-interface-metric-equation} with
\eqref{eq:V90-corrected-matter-flux} changes only the sign of the mixed bulk
term and gives \eqref{eq:V90-corrected-full-gate}.
\end{proof}

\subsection{One classical Standard Model Ward bank}
\label{subsec:V90-classical-SM-Ward}

Let \(P\to M\) be the Standard Model gauge bundle with compact gauge group
and invariant inner product.  Let \(A\) be a connection, \(F\) its
curvature, \(\Phi\) the Higgs multiplet in its unitary representation, and
\(\Psi\) the direct sum of all chiral fermion and generation bundles.  The
spin connection and gauge connection are both included in
\(\nabla^A\Psi\).  Yukawa matrices and the Higgs potential are arbitrary
smooth gauge-invariant coefficients.  This notation keeps the full field
content while avoiding a choice of gauge basis.

Let \(S_{\rm SM}[g,A,\Phi,\Psi,\overline\Psi]\) be the standard local
gauge- and diffeomorphism-invariant classical action.  Define its Euler
covectors and Hilbert stress by
\begin{align}
 \delta S_{\rm SM}
 ={}&\frac12\int_M\sqrt{-g}\,
 T_{\rm SM}^{\mu\nu}\delta g_{\mu\nu}
 +\langle\mathcal E_A,\delta A\rangle
 +\langle\mathcal E_\Phi,\delta\Phi\rangle
 \notag\\
 &+\langle\mathcal E_\Psi,\delta\Psi\rangle
 +\langle\mathcal E_{\overline\Psi},
          \delta\overline\Psi\rangle
 +\Theta_{\partial,\rm SM}(\delta).
 \label{eq:V90-SM-first-variation}
\end{align}
Pairings include spacetime integration and the Hermitian real part.  The
boundary term is retained rather than discarded sector by sector.

For a gauge parameter \(\alpha\), write
\begin{align}
 \delta_\alpha A&=D_A\alpha,
 &\delta_\alpha\Phi&=-\rho_\Phi(\alpha)\Phi,
 &\delta_\alpha\Psi&=-\rho_\Psi(\alpha)\Psi.
 \label{eq:V90-gauge-variations}
\end{align}
The conjugate spinor has the dual variation.

\begin{lemma}[Complete gauge Noether identity]
\label{lem:V90-gauge-Noether}
For every compactly supported \(\alpha\),
\begin{align}
 0={}&
 \langle\mathcal E_A,D_A\alpha\rangle
 -\langle\mathcal E_\Phi,
ho_\Phi(\alpha)\Phi\rangle
 -\langle\mathcal E_\Psi,
ho_\Psi(\alpha)\Psi\rangle
 \notag\\
 &-\langle\mathcal E_{\overline\Psi},
       \overline\Psi\rho_\Psi(\alpha)^*\rangle.
 \label{eq:V90-gauge-Noether}
\end{align}
This single identity contains all gauge factors and matter
representations.
\end{lemma}

\begin{proof}
Insert \eqref{eq:V90-gauge-variations} in
\eqref{eq:V90-SM-first-variation}.  Compact support removes the boundary
term, the metric variation is zero, and gauge invariance makes the first
variation vanish.  The displayed identity is what remains.
\end{proof}

For a vector field \(\xi\), use the gauge-covariant diffeomorphism
variations
\begin{align}
 \mathbb L_\xi^A A&:=\iota_\xi F,
 &\mathbb L_\xi^A\Phi&:=\nabla_\xi^A\Phi,
 &\mathbb L_\xi^A\Psi&:=\mathbb L_\xi^{K,A}\Psi,
 \label{eq:V90-covariant-Lie}
\end{align}
where \(\mathbb L^{K,A}\) is the gauge-covariant Kosmann derivative.  The
ordinary Lie derivative of \(A\) differs from \(\iota_\xi F\) by
\(D_A(\iota_\xi A)\); Lemma~\ref{lem:V90-gauge-Noether} removes that pure
gauge difference, including its action on \(\Phi\) and \(\Psi\).

\begin{theorem}[Off-shell classical Standard Model Ward identity]
\label{thm:V90-classical-SM-Ward}
For every compactly supported vector field \(\xi\),
\begin{align}
 0={}&\int_M\sqrt{-g}\,
 T_{\rm SM}^{\mu\nu}\nabla_\mu\xi_\nu
 +\langle\mathcal E_A,\iota_\xi F\rangle
 +\langle\mathcal E_\Phi,\nabla_\xi^A\Phi\rangle
 \notag\\
 &+\langle\mathcal E_\Psi,
           \mathbb L_\xi^{K,A}\Psi\rangle
 +\langle\mathcal E_{\overline\Psi},
           \mathbb L_\xi^{K,A}\overline\Psi\rangle.
 \label{eq:V90-offshell-SM-Ward}
\end{align}
In particular, all Standard Model Euler equations imply
\begin{equation}
 \nabla_\mu T_{\rm SM}^{\mu}{}_{\nu}=0
 \label{eq:V90-onshell-SM-conservation}
\end{equation}
as a distribution.
\end{theorem}

\begin{proof}
Apply an infinitesimal diffeomorphism to
\eqref{eq:V90-SM-first-variation}.  The metric variation is
\(2\nabla_{(\mu}\xi_{\nu)}\).  Replace the ordinary transformations of the
gauge and charged fields by \eqref{eq:V90-covariant-Lie}; their difference
is the gauge transformation with parameter \(\iota_\xi A\), whose complete
Euler pairing vanishes by Lemma~\ref{lem:V90-gauge-Noether}.  Diffeomorphism
invariance gives \eqref{eq:V90-offshell-SM-Ward}.  On the Euler shell, one
integration by parts against arbitrary compactly supported \(\xi\) gives
\eqref{eq:V90-onshell-SM-conservation}.
\end{proof}

\begin{firewall}[Why the fermion row is left in Kosmann form]
\label{fire:V90-Kosmann-form}
Expanding the Kosmann derivative off shell moves derivatives of \(\xi\)
between the fermion Euler covectors and the Belinfante spin current.  The
compact identity \eqref{eq:V90-offshell-SM-Ward} is independent of that
integration-by-parts choice and uses the symmetric Hilbert stress already
selected by metric variation.  Dropping the Kosmann term before imposing the
Dirac equations would give an incomplete off-shell Ward row.
\end{firewall}

\subsection{Boundary Standard Model flux with the corrected sign}
\label{subsec:V90-boundary-SM-Ward}

Let \(S_{\partial,\rm SM}\) be a gauge- and diffeomorphism-invariant
boundary action, including boundary conditions and allowed local
counterterms.  Define its stress by
\begin{equation}
 \delta S_{\partial,\rm SM}
 =\frac12\int_\Gamma\sqrt{|\gamma|}\,
   \tau_{\rm SM}^{ab}\delta\gamma_{ab}
 +\langle e_{\partial},\delta\phi_{\partial}\rangle.
 \label{eq:V90-boundary-SM-variation}
\end{equation}
Let \(\pi_{\partial}\) be the boundary Euler covector from
\(\Theta_{\partial,\rm SM}\) in
\eqref{eq:V90-SM-first-variation}, and set
\begin{equation}
 \mathcal E_{\partial}:=\pi_{\partial}+e_{\partial}.
 \label{eq:V90-total-boundary-Euler}
\end{equation}
The symbol \(\mathbb L_b^{A,K}\phi_{\partial}\) denotes the coefficient of
a tangential vector field \(\zeta^b\) in the covariant Lie derivative of all
boundary fields.

\begin{theorem}[Exact Standard Model boundary Ward identity]
\label{thm:V90-boundary-SM-Ward}
On the bulk Standard Model Euler equations,
\begin{equation}
 \boxed{
 D_a\tau_{\rm SM}^a{}_b
 =n^\mu\gamma_b{}^\nu T_{\mu\nu}^{\rm SM}
  +\mathcal W_{\partial,b}^{\rm SM},}
 \label{eq:V90-boundary-SM-Ward}
\end{equation}
where the off-shell boundary Ward source is defined weakly by
\begin{equation}
 \int_\Gamma\sqrt{|\gamma|}\,
 \mathcal W_{\partial,b}^{\rm SM}\zeta^b
 :=\left\langle\mathcal E_{\partial},
       \mathbb L_\zeta^{A,K}\phi_{\partial}\right\rangle.
 \label{eq:V90-boundary-Ward-source}
\end{equation}
Thus the total boundary Euler equation implies
\begin{equation}
 D_a\tau_{\rm SM}^a{}_b
 =n^\mu\gamma_b{}^\nu T_{\mu\nu}^{\rm SM}.
 \label{eq:V90-onshell-boundary-flux}
\end{equation}
\end{theorem}

\begin{proof}
For a tangential \(\zeta\), the on-bulk-shell diffeomorphism variation of
the bulk action gives
\begin{equation}
 \left\langle\pi_{\partial},
       \mathbb L_\zeta^{A,K}\phi_{\partial}\right\rangle
 =-\int_\Gamma\sqrt{|\gamma|}\,
 n^\mu T_{\mu\nu}^{\rm SM}\zeta^\nu.
 \label{eq:V90-bulk-boundary-Noether}
\end{equation}
The boundary action Ward identity gives
\begin{equation}
 \int_\Gamma\sqrt{|\gamma|}\,
 (D_a\tau_{\rm SM}^a{}_b)\zeta^b
 =\left\langle e_{\partial},
       \mathbb L_\zeta^{A,K}\phi_{\partial}\right\rangle.
 \label{eq:V90-surface-Noether}
\end{equation}
Substitute
\(e_{\partial}=\mathcal E_{\partial}-\pi_{\partial}\) and use
\eqref{eq:V90-bulk-boundary-Noether}.  This gives
\eqref{eq:V90-boundary-SM-Ward}--\eqref{eq:V90-boundary-Ward-source}.
The total boundary Euler equation removes the last term.
\end{proof}

\begin{corollary}[Cancellation in the corrected Einstein gate]
\label{cor:V90-SM-cancels-gate}
Suppose \(T^{\rm bulk}=T^{\rm SM}+T^{\rm other}\).  Inserting
\eqref{eq:V90-boundary-SM-Ward} in
\eqref{eq:V90-corrected-full-gate} cancels the Standard Model mixed bulk
flux exactly and leaves
\begin{equation}
 0=-n^\mu\gamma_b{}^\nu T_{\mu\nu}^{\rm other}
 +\mathcal E_\Theta D_b\Theta
 +(\mathfrak d_{\mathcal A}^*t_b)_b
 +\mathcal W_{\partial,b}^{\rm SM}.
 \label{eq:V90-reduced-boundary-gate}
\end{equation}
On the clock, equivariance, and Standard Model boundary Euler shells, the
remaining homogeneous condition is just the normal flux balance for any
other bulk sector.
\end{corollary}

\begin{proof}
The Standard Model terms in the corrected gate are
\(-nT^{\rm SM}+D\tau_{\rm SM}\).  Equation
\eqref{eq:V90-boundary-SM-Ward} makes their sum
\(\mathcal W_{\partial}^{\rm SM}\), proving
\eqref{eq:V90-reduced-boundary-gate}.
\end{proof}

\subsection{A common renormalized stress-domain theorem}
\label{subsec:V90-renormalized-domain}

The classical identity does not license a pointwise quantum stress tensor.
Let \(\mathscr H_{\rm SM}\) be the physical quantum Hilbert space and let
\(H_{\rm SM}\ge0\) be its Hamiltonian.  Fix the common form domain
\begin{equation}
 \mathscr D_{\rm SM}:=D((I+H_{\rm SM})^{1/2}),
 \qquad
 \mathscr D_{\rm SM}\subset\mathscr H_{\rm SM}
 \subset\mathscr D_{\rm SM}'.
 \label{eq:V90-SM-form-rigging}
\end{equation}
For a regulator \(N\), let
\(T_N^{\rm ren}(f)\) be the complete smeared, renormalized Standard Model
stress quadratic form, after all fields and a single invariant counterterm
action have been combined.  Let \(\mathcal E_{I,N}(r)\) denote the similarly
smeared Euler insertions.

\begin{theorem}[Sufficient common-domain renormalized Ward gate]
\label{thm:V90-renormalized-Ward-gate}
Assume the following for every compactly supported smooth test tensor \(f\)
and Euler test bank \(r\).
\begin{enumerate}[label=\textup{(\roman*)},leftmargin=2.8em]
\item Every \(T_N^{\rm ren}(f)\) and \(\mathcal E_{I,N}(r)\) is a continuous
      sesquilinear form on the same \(\mathscr D_{\rm SM}\).
\item There are regulator-independent constants such that
      \begin{align}
       |T_N^{\rm ren}(f)[\psi,\chi]|
       &\le C_f\|\psi\|_{\mathscr D_{\rm SM}}
                    \|\chi\|_{\mathscr D_{\rm SM}},
       \label{eq:V90-uniform-stress-form-bound}\\
       |\mathcal E_{I,N}(r)[\psi,\chi]|
       &\le C_r\|\psi\|_{\mathscr D_{\rm SM}}
                    \|\chi\|_{\mathscr D_{\rm SM}}.
       \label{eq:V90-uniform-Euler-form-bound}
      \end{align}
\item The stress and Euler forms are Cauchy in
      \(\mathcal B(\mathscr D_{\rm SM},\mathscr D_{\rm SM}')\), locally
      uniformly in bounded sets of test functions.
\item At each regulator, the stress and Euler forms are first variations of
      one gauge- and diffeomorphism-invariant regulated action; equivalently,
      they obey the regulated analogue of
      \eqref{eq:V90-offshell-SM-Ward} with defect \(R_N(\xi)\).
\item The Ward defects satisfy
      \begin{equation}
       \|R_N(\xi)\|_{\mathcal B(\mathscr D_{\rm SM},
                                  \mathscr D_{\rm SM}')}
       \longrightarrow0
       \label{eq:V90-Ward-defect-limit}
      \end{equation}
      for every test vector field, locally uniformly in its test topology.
\end{enumerate}
Then there are unique operator-valued distributions
\begin{equation}
 T_{\rm SM}^{\rm ren},\ \mathcal E_I^{\rm ren}
 \in\mathcal D'(M;
 \mathcal B(\mathscr D_{\rm SM},\mathscr D_{\rm SM}'))
 \label{eq:V90-renormalized-distributions}
\end{equation}
which are the regulator limits and obey the exact weak identity
\begin{align}
 0={}&T_{\rm SM}^{\rm ren}(\nabla_{(\mu}\xi_{\nu)})
 +\langle\mathcal E_A^{\rm ren},\iota_\xi F\rangle
 +\langle\mathcal E_\Phi^{\rm ren},\nabla_\xi^A\Phi\rangle
 \notag\\
 &+\langle\mathcal E_\Psi^{\rm ren},
       \mathbb L_\xi^{K,A}\Psi\rangle
 +\langle\mathcal E_{\overline\Psi}^{\rm ren},
       \mathbb L_\xi^{K,A}\overline\Psi\rangle.
 \label{eq:V90-renormalized-Ward}
\end{align}
On the common Euler-form shell, the renormalized stress is conserved as an
operator-valued distribution.
\end{theorem}

\begin{proof}
Completeness of
\(\mathcal B(\mathscr D_{\rm SM},\mathscr D_{\rm SM}')\) and the Cauchy
hypothesis give a unique limit for every test function.  Local uniformity in
the test topology makes these limits distributions, while
\eqref{eq:V90-uniform-stress-form-bound}--%
\eqref{eq:V90-uniform-Euler-form-bound} retain the common form domain.  Pass
to the limit in the regulated Ward identity.  Every stress and Euler term
converges in the form-operator norm, and
\eqref{eq:V90-Ward-defect-limit} removes the remainder.  This proves
\eqref{eq:V90-renormalized-Ward}.  Vanishing of the Euler forms leaves the
distributional divergence equation.
\end{proof}

\begin{firewall}[Finite sector values do not prove the common gate]
\label{fire:V90-sectorwise-renormalization}
Separate existence of gauge, Higgs, and fermion stress expectations does not
imply hypotheses \textup{(i)}--\textup{(v)}.  Their domains may differ, and
finite counterterms may carry uncancelled gauge or diffeomorphism variation.
The common invariant action and common form convergence are the content of
the gate, not bookkeeping after the limit.
\end{firewall}

\subsection{Normal flux and the V88 test dual}
\label{subsec:V90-normal-flux-gate}

An operator-valued distribution on \(M\) does not automatically restrict to
\(\Gamma\).  Let \(N^*\Gamma\) be the conormal bundle of the embedding.

\begin{proposition}[Two rigorous routes to the boundary flux]
\label{prop:V90-boundary-flux-routes}
The renormalized normal flux
\begin{equation}
 j_b^{\rm SM}:=
 n^\mu\gamma_b{}^\nu T_{\mu\nu}^{\rm ren}|_\Gamma
 \label{eq:V90-renormalized-normal-flux}
\end{equation}
is well defined on \(\mathscr D_{\rm SM}\) by either of the following
independent sufficient constructions.
\begin{enumerate}[label=\textup{(\alph*)},leftmargin=2.8em]
\item Every matrix-element distribution of \(T_{\rm SM}^{\rm ren}\) has
      wavefront set disjoint from \(N^*\Gamma\), locally uniformly for
      bounded subsets of \(\mathscr D_{\rm SM}\).  Then the distributional
      pullback to \(\Gamma\) exists.
\item The common regulated action has a boundary first variation whose flux
      forms converge directly in
      \(\mathcal B(\mathscr D_{\rm SM},\mathscr D_{\rm SM}')\), so
      \eqref{eq:V90-renormalized-normal-flux} is defined variationally rather
      than by restricting a bulk product.
\end{enumerate}
If, in addition, for some BFV seminorm \(p_N\),
\begin{equation}
 |\langle j^{\rm SM},\zeta\rangle_Gamma[\psi,\chi]|
 \le C\,p_N(\zeta)\,
 \|\psi\|_{\mathscr D_{\rm SM}}
 \|\chi\|_{\mathscr D_{\rm SM}},
 \label{eq:V90-flux-BFV-bound}
\end{equation}
then
\begin{equation}
 j^{\rm SM}
 \in\mathcal B(\mathscr D_{\rm SM},\mathscr D_{\rm SM}')
   \widehat\otimes\mathscr D_{\rm BFV}',
 \label{eq:V90-flux-BFV-dual}
\end{equation}
and the boundary Ward identity pairs continuously with the V88 physical
energy/BFV rigging.
\end{proposition}

\begin{proof}
Part \textup{(a)} is the distributional pullback theorem: transversality of
the wavefront set to the conormal is exactly its hypothesis.  Part
\textup{(b)} defines a continuous boundary distribution from the limiting
first variation and therefore needs no bulk restriction.  Estimate
\eqref{eq:V90-flux-BFV-bound} is precisely continuity in the BFV test
topology and the common quantum form topology; the universal property of the
completed tensor product gives \eqref{eq:V90-flux-BFV-dual}.
\end{proof}

\begin{assessment}[Classical closure and quantum acquisition boundary]
\label{ass:V90-Ward-status}
The complete classical Standard Model Euler bank now cancels its mixed bulk
stress in the corrected Einstein boundary gate.  At the quantum level, V90
reduces the same claim to a common invariant counterterm action, common form
convergence, vanishing Ward defect, and a valid normal-flux construction.
None of those analytic hypotheses may be replaced by separate finite sector
expectations.
\end{assessment}

\subsection{V90 result ledger and next acquisition}
\label{subsec:V90-ledger}

\paragraph{New exact conclusions.}
\begin{enumerate}[label=\textup{(\arabic*)},leftmargin=2.8em]
\item The required V90 audit found NS V31 unchanged and recorded the stable
      YM V88 multiplicity/ancestry result without importing a mass claim.
\item The V36 action convention fixes the Brown--York divergence with a
      minus mixed-Einstein sign; V89's opposite printed sign is corrected.
\item The complete gauge Euler bank permits a gauge-covariant
      diffeomorphism identity without dropping any charged sector.
\item The classical Standard Model Hilbert stress is conserved on the full
      Euler shell.
\item Its boundary stress obeys
      \(D\tau_{\rm SM}=nT^{\rm SM}+\mathcal W_\partial^{\rm SM}\), so its
      mixed bulk flux cancels exactly in the corrected Einstein gate.
\item A common-form convergence theorem carries the full off-shell Ward
      identity through a regulator limit.
\item The normal quantum flux requires either the wavefront pullback
      condition or a convergent boundary first variation, followed by the
      explicit V88 dual bound \eqref{eq:V90-flux-BFV-bound}.
\end{enumerate}

\paragraph{Next forced step.}
V91 must attack one of the two genuinely analytic quantum gates rather than
repeat the classical Ward algebra.  The most upstream choice is the normal
flux: select a concrete Hadamard/point-split regulator for the gauge, Higgs,
and fermion fields on a smooth timelike boundary, compute the conormal
wavefront intersection, and decide whether route \textup{(a)} of
Proposition~\ref{prop:V90-boundary-flux-routes} is available.  If the
intersection is nonempty, return to route \textup{(b)} and construct the
renormalized boundary first variation directly.  The curved radiative
equivariance defect and the dynamical clock equation remain parallel gates.

\begin{assessment}[V90 continuum status]
\label{ass:V90-continuum-status}
The full Standard--Model--Einstein continuum remains unproved.  V90 fixes
the Brown--York orientation, closes the complete classical Standard Model
Ward cancellation, and states a rigorous common-domain quantum passage.
The required uniform form estimates, boundary flux construction, dynamical
clock, curved equivariant extractor, nonlinear existence, and cofinal
interacting limit remain open.
\end{assessment}

\section*{V90 exact checks and append-only record}
\addcontentsline{toc}{section}{V90 exact checks and append-only record}

\begin{equation}
\boxed{
\begin{gathered}
 D_a\Pi^a{}_b=-n^\mu\gamma_b{}^\nu G_{\mu\nu},\\
 D_a\tau_{\rm SM}^a{}_b
 =n^\mu\gamma_b{}^\nu T_{\mu\nu}^{\rm SM}
  +\mathcal W_{\partial,b}^{\rm SM},\\
 -nT^{\rm SM}+D\tau_{\rm SM}=\mathcal W_\partial^{\rm SM},\\
 T_N^{\rm ren},\mathcal E_{I,N}^{\rm ren}
 \longrightarrow
 T_{\rm SM}^{\rm ren},\mathcal E_I^{\rm ren}
 \quad\hbox{in }\mathcal B(\mathscr D_{\rm SM},\mathscr D_{\rm SM}').
\end{gathered}}
\label{eq:V90-check-card}
\end{equation}

The V89 predecessor used for this append has SHA--256
\begin{center}
\texttt{a4a669e8746e4b09e45552ac0574d03ecdfb1f827aaefbe05f3fbbde97db74e5}.
\end{center}
The V90 candidate retains its bytes up to the sole live final document
terminator, appends this material, and restores exactly one active
\verb|\end{document}|.  Validation includes balanced environments and
braces, unique labels, resolved new internal references, valid inline and
display math delimiters, exact predecessor-prefix equality, and a final
inventory which removes only the superseded SM V89 file.  The audit was
read-only and changed neither companion file.

% =============================================================================
% END APPENDED VERSION V90 MATERIAL
% =============================================================================


\clearpage
% =============================================================================
% BEGIN APPENDED VERSION V91 MATERIAL
% =============================================================================
\section{V91: timelike Hadamard restriction, the boundary-diagonal
obstruction, and the forced variational flux route}
\label{sec:V91-boundary-Hadamard}

V90 gave two sufficient routes to a renormalized Standard Model normal
flux.  The first was a wavefront-transverse pullback of an already
renormalized bulk stress; the second was direct convergence of the boundary
first variation.  This note tests the first route at the most upstream
point-split level.

Two pullbacks must not be confused.  A Hadamard two-point distribution may
be restricted from \(M\times M\) to \(\Gamma\times\Gamma\), because a
timelike hypersurface has spacelike conormal whereas the characteristic
covectors are null.  A stress expectation then requires a second pullback:
coincidence along the diagonal of \(\Gamma\times\Gamma\).  The latter
conormal contains tangential null pairs.  Whenever the selected composite
has a nonzero restricted Hadamard symbol there, the two sets intersect, and
the coincidence limit exists only after an adequate boundary subtraction.

A massless scalar image calculation makes the analytic failure quantitative.
After subtracting the ordinary Minkowski Hadamard singularity, the reflected
term has a second-normal-derivative composite proportional to \(x^{-4}\).
It has no canonical restriction to \(x=0\); its extensions differ by
boundary-supported distributions.  Hence the bulk Hadamard condition alone
does not establish route \textup{(a)} of V90.  The common invariant boundary
counterterm action in route \textup{(b)} is forced before the full quantum
flux can be claimed.

\subsection{Wavefront geometry of a timelike boundary}
\label{subsec:V91-wavefront-geometry}

Let \(i:\Gamma\hookrightarrow M\) be a smooth timelike hypersurface.  Its
conormal bundle is
\begin{equation}
 N^*\Gamma
 =\{(x,\lambda n_x^\flat):x\in\Gamma,\ \lambda\ne0\}.
 \label{eq:V91-boundary-conormal}
\end{equation}
Every nonzero element of \(N^*\Gamma\) is spacelike.  For the product
embedding \(I=i\times i\),
\begin{equation}
 N^*(\Gamma\times\Gamma)
 =N^*\Gamma\times N^*\Gamma
 \quad\hbox{with either factor allowed to vanish}.
 \label{eq:V91-product-conormal}
\end{equation}

For a free normally hyperbolic scalar, Dirac, or gauge-fixed vector operator,
write the Hadamard wavefront relation as
\begin{equation}
 \mathcal C_+
 :=\left\{
 (x,k;x',-k'):
 (x,k)\sim(x',k'),\quad k\in\overline V_+^*,\quad k^2=0
 \right\}.
 \label{eq:V91-Hadamard-relation}
\end{equation}
The bundle indices and polarizations do not alter the cotangent relation.
Reflected boundary singularities propagate along broken null
bicharacteristics and still have null covectors in each slot.

\begin{theorem}[First timelike pullback is well defined]
\label{thm:V91-first-pullback}
Let \(W\in\mathcal D'(M\times M;E\boxtimes E^*)\) satisfy
\begin{equation}
 \operatorname{WF}(W)\subset\mathcal C_+
 \label{eq:V91-Hadamard-WF-bound}
\end{equation}
or the corresponding finite union of direct and reflected null relations.
Then the pullback
\begin{equation}
 W_\Gamma:=I^*W\in
 \mathcal D'(\Gamma\times\Gamma;i^*E\boxtimes i^*E^*)
 \label{eq:V91-boundary-two-point}
\end{equation}
exists uniquely.  The same statement holds after applying any finite-order
bidifferential operator before restriction.
\end{theorem}

\begin{proof}
The distributional pullback theorem requires
\begin{equation}
 \operatorname{WF}(W)\cap N^*(\Gamma\times\Gamma)=\varnothing.
 \label{eq:V91-first-transversality}
\end{equation}
Every nonzero covector in a factor of
\eqref{eq:V91-product-conormal} is spacelike.  Every nonzero covector in the
corresponding factor of \(\mathcal C_+\) is null.  They cannot be equal, so
\eqref{eq:V91-first-transversality} holds.  Broken propagation does not
change nullity.  Differential operators do not enlarge wavefront sets,
which proves the last assertion.
\end{proof}

\begin{remark}[Glancing does not spoil the first pullback]
\label{rem:V91-glancing-first-pullback}
A null covector may be tangent to a timelike boundary, but it cannot be
normal to it.  Glancing therefore belongs to the restricted wavefront set on
\(T^*(\Gamma\times\Gamma)\); it does not violate
\eqref{eq:V91-first-transversality}.
\end{remark}

\subsection{The second pullback is the real obstruction}
\label{subsec:V91-diagonal-obstruction}

Let \(\Delta_\Gamma:\Gamma\to\Gamma\times\Gamma\) be the diagonal map.  Its
conormal bundle is
\begin{equation}
 N^*\Delta_\Gamma
 =\{(x,\eta;x,-\eta):\eta\in T_x^*\Gamma,\ \eta\ne0\}.
 \label{eq:V91-diagonal-conormal}
\end{equation}
Unlike \(N^*\Gamma\), this bundle contains null pairs because the induced
metric on \(\Gamma\) is Lorentzian.

\begin{theorem}[Criterion for the boundary coincidence obstruction]
\label{thm:V91-diagonal-no-pullback}
Let \(P\) be the finite-order bidifferential operator selecting a proposed
boundary composite from a Hadamard two-point function.  If the restricted
principal Hadamard symbol of \(PW\) is nonzero at some tangential null
covector, then
\begin{equation}
 \operatorname{WF}(I^*PW)
 \cap N^*\Delta_\Gamma\ne\varnothing.
 \label{eq:V91-diagonal-intersection}
\end{equation}
Consequently \(\Delta_\Gamma^*I^*PW\) is not defined by the distributional
pullback theorem.  This hypothesis holds for the direct scalar Hadamard
term with \(P=I\), and it must be checked after any boundary-condition
cancellation and after applying the stress bidifferential operator.  A
subtraction which removes every surviving boundary-diagonal symbol is
necessary.
\end{theorem}

\begin{proof}
Choose a tangential null covector at which the restricted principal symbol
of \(PW\) is nonzero.  Its lift to \(T_x^*M\) is null because it annihilates
the spacelike normal and the induced metric is the restriction of \(g\).
The local Hadamard relation at coincidence contains
\((x,\eta;x,-\eta)\), and the nonzero principal symbol prevents \(P\) from
removing that wavefront element.  Pullback to \(\Gamma\times\Gamma\) retains
the pair.  It is also an element of
\eqref{eq:V91-diagonal-conormal}, proving
\eqref{eq:V91-diagonal-intersection}.  The pullback criterion therefore
fails.  Point splitting supplies a defined coincidence limit only after the
surviving singular part has been subtracted.
\end{proof}

\begin{firewall}[Restriction and renormalized coincidence are different]
\label{fire:V91-two-pullbacks}
Theorem~\ref{thm:V91-first-pullback} permits boundary two-point
distributions and boundary Green forms.  It does not define a local composite
field on the boundary.  Theorem~\ref{thm:V91-diagonal-no-pullback} concerns
the additional coincidence map needed for \(T_{\mu\nu}\), \(\Phi^*\Phi\),
currents, and other Wick polynomials.
\end{firewall}

\subsection{A flat image witness for the missing boundary subtraction}
\label{subsec:V91-image-witness}

Consider a massless scalar on the Minkowski half-space \(x>0\) with a flat
Dirichlet wall at \(x=0\).  Let \(\bar x'\) be reflection of \(x'\) across
the wall.  The image two-point distribution is
\begin{equation}
 W_D(X,X')=W_0(X,X')-W_0(X,\bar X'),
 \label{eq:V91-Dirichlet-image}
\end{equation}
where
\begin{equation}
 W_0(X,X')
 =\frac1{4\pi^2}
 \frac1{-(t-t'-i0)^2+|y-y'|^2+(x-x')^2}.
 \label{eq:V91-Minkowski-two-point}
\end{equation}
Subtracting the ordinary bulk Hadamard parametrix removes the first term but
leaves
\begin{equation}
 W_{\rm im}(X,X'):=-W_0(X,\bar X').
 \label{eq:V91-image-remainder}
\end{equation}

\begin{proposition}[Exact wall divergence of a point-split derivative]
\label{prop:V91-image-divergence}
At equal tangential coordinates and positive \(x=x'\),
\begin{align}
 W_{\rm im}(X,X)&=-\frac1{16\pi^2x^2},
 \label{eq:V91-image-square}\\
 \left.
 \partial_x\partial_{x'}W_{\rm im}(X,X')
 \right|_{X'=X}
 &=-\frac3{32\pi^2x^4}.
 \label{eq:V91-image-derivative}
\end{align}
The second expression is one of the point-split derivative composites from
which the scalar stress is assembled.  Thus ordinary bulk Hadamard
subtraction does not produce a distribution with a canonical boundary
restriction.
\end{proposition}

\begin{proof}
At equal tangential coordinates the reflected separation has squared length
\((x+x')^2\), so
\begin{equation*}
 W_{\rm im}(X,X')
 =-\frac1{4\pi^2(x+x')^2}.
\end{equation*}
Set \(x'=x\) to obtain \eqref{eq:V91-image-square}.  Two derivatives give
\begin{equation*}
 \partial_x\partial_{x'}
 \left[-\frac1{4\pi^2(x+x')^2}\right]
 =-\frac3{2\pi^2(x+x')^4};
\end{equation*}
setting \(x'=x\) proves \eqref{eq:V91-image-derivative}.
The function \(x^{-4}\) is not locally integrable at zero and has no
ordinary pullback to the wall.
\end{proof}

\begin{remark}[The witness is upstream of stress-tensor improvements]
\label{rem:V91-improvement-scope}
Improvement terms may cancel particular components in special flat states.
They do not make the composite derivative in
\eqref{eq:V91-image-derivative} canonically restrictable, nor do they supply
the gauge, Higgs, fermion, and clock counterterms required by the common V90
Ward identity.  The witness proves failure of the bulk-subtraction premise,
not a universal nonzero value for every improved normal flux.
\end{remark}

\subsection{Extension ambiguity is a boundary counterterm}
\label{subsec:V91-extension-ambiguity}

Let \(x\) be a boundary defining function.  On \(x>0\), the distribution
\(x^{-4}\) has scaling degree four in the single transverse variable.

\begin{theorem}[Finite local extension ambiguity]
\label{thm:V91-extension-ambiguity}
Extensions of \(x^{-4}\) from \((0,\infty)\) across \(x=0\) which
preserve scaling degree four exist.  Any two such extensions differ by
\begin{equation}
 c_0\delta(x)+c_1\delta'(x)+c_2\delta''(x)+c_3\delta'''(x).
 \label{eq:V91-extension-difference}
\end{equation}
With tangential tensor coefficients, covariance and gauge invariance restrict
the allowed combinations but do not remove the need to choose them.  These
terms are precisely variations of local boundary counterterms and their
normal-jet extensions.
\end{theorem}

\begin{proof}
The standard finite-part construction, obtained by subtracting the Taylor
polynomial of a test function through order three before integrating against
\(x^{-4}\), defines an extension.  The difference of two extensions is
supported at \(x=0\).  Preservation of scaling degree four bounds the difference to order at
most three.  A distribution supported at one point with that order is a
linear combination of
\(\delta,\delta',\delta'',\delta'''\).  The order bound is the degree of
divergence \(4-1=3\).  Tensor covariance and gauge invariance act linearly on
the coefficient space, leaving the invariant subspace corresponding to
admissible local counterterms.
\end{proof}

\begin{corollary}[Bulk Hadamard data do not select the boundary stress]
\label{cor:V91-bulk-Hadamard-insufficient}
The bulk Hadamard singularity and the bulk Ward identity do not determine a
unique renormalized stress at a timelike boundary.  Boundary normalization
conditions, imposed through one invariant counterterm action, are additional
data.
\end{corollary}

\begin{proof}
Proposition~\ref{prop:V91-image-divergence} leaves an \(x^{-4}\) boundary
singularity after bulk subtraction.  Theorem
\ref{thm:V91-extension-ambiguity} gives its nonunique local extensions.
Only a boundary renormalization prescription selects the coefficients.
\end{proof}

\subsection{The variational route and preservation of the Ward identity}
\label{subsec:V91-variational-route}

Let \(S_{N}^{\rm bulk}\) be a gauge- and diffeomorphism-invariant regulated
Standard Model action and let
\begin{equation}
 S_N^{\rm tot}
 :=S_N^{\rm bulk}+S_N^{\partial},
 \label{eq:V91-total-regulated-action}
\end{equation}
where \(S_N^\partial\) is a local invariant boundary action depending on
the induced fields and the normal jets permitted by the chosen boundary
problem.  Its coefficients include the extensions in
\eqref{eq:V91-extension-difference} in covariant form.

\begin{theorem}[Boundary first variation is the correct flux primitive]
\label{thm:V91-boundary-first-variation}
Assume that, on the common form domain \(\mathscr D_{\rm SM}\), the complete
boundary first variations of \(S_N^{\rm tot}\) converge in
\(\mathcal B(\mathscr D_{\rm SM},\mathscr D_{\rm SM}')\) against every
smooth compact boundary test variation.  Then their limit defines
\begin{equation}
 (\tau_{\rm SM}^{\rm ren},
  \mathcal E_{\partial}^{\rm ren},
  j_{\rm SM}^{\rm ren})
 \label{eq:V91-renormalized-boundary-bank}
\end{equation}
as boundary operator-valued distributions.  If each
\(S_N^{\rm tot}\) is gauge and diffeomorphism invariant and the regulated
Ward defects converge to zero, the limit obeys
\begin{equation}
 D_a(\tau_{\rm SM}^{\rm ren})^a{}_b
 =n^\mu\gamma_b{}^\nu
   (T_{\mu\nu}^{\rm SM})^{\rm ren}
 +\mathcal W_{\partial,b}^{\rm ren}
 \label{eq:V91-renormalized-boundary-Ward}
\end{equation}
without taking the boundary pullback of an independently renormalized bulk
composite.
\end{theorem}

\begin{proof}
The convergence hypothesis defines each coefficient of the common boundary
first variation as an element of
\(\mathcal D'(\Gamma;\mathcal B(\mathscr D_{\rm SM},
\mathscr D_{\rm SM}'))\).  At finite \(N\), apply the proof of
Theorem~\ref{thm:V90-boundary-SM-Ward} to the single action
\eqref{eq:V91-total-regulated-action}.  All terms belong to the same form
domain.  Passing to the form-operator limit and removing the Ward defect
gives \eqref{eq:V91-renormalized-boundary-Ward}.  Since the normal flux is a
coefficient of the convergent first variation, no separate pullback is used.
\end{proof}

\begin{firewall}[Counterterms must be varied before restriction]
\label{fire:V91-vary-before-restrict}
It is not sufficient to subtract a numerical boundary divergence from a
stress expectation after computing it.  The subtraction must be the first
variation of the same invariant boundary action in all metric, gauge, Higgs,
fermion, and clock directions.  Only then do the Euler and stress
counterterms satisfy the common Ward identity at each regulator.
\end{firewall}

\subsection{The exact V88 continuity target}
\label{subsec:V91-BFV-target}

Let \(\mathscr C_{\partial,\rm BFV}\) be the boundary traces of the V88
smooth compatible core.  For a boundary first-variation bank \(B_N\), put
\begin{equation}
 \|B_N\|_{r}
 :=\sup_{\substack{p_r(\zeta)\le1\\
                    \|\psi\|_{\mathscr D_{\rm SM}}\le1\\
                    \|\chi\|_{\mathscr D_{\rm SM}}\le1}}
 |B_N(\zeta)[\psi,\chi]|.
 \label{eq:V91-boundary-bank-norm}
\end{equation}

\begin{proposition}[One estimate closes the flux-to-BFV passage]
\label{prop:V91-one-BFV-estimate}
If there is an \(r\) such that
\begin{equation}
 \sup_N\|B_N\|_r<\infty,
 \qquad
 \|B_N-B_M\|_r\longrightarrow0,
 \label{eq:V91-BFV-Cauchy-gate}
\end{equation}
then the limiting bank in
\eqref{eq:V91-renormalized-boundary-bank} lies in
\begin{equation}
 \mathcal B(\mathscr D_{\rm SM},\mathscr D_{\rm SM}')
 \widehat\otimes\mathscr D_{\rm BFV}',
 \label{eq:V91-boundary-bank-BFV-dual}
\end{equation}
and its Ward identity pairs continuously with every V88 energy/BFV Green
identity.  Conversely, convergence only on each fixed smooth boundary field
does not imply \eqref{eq:V91-boundary-bank-BFV-dual}.
\end{proposition}

\begin{proof}
The first part is completion in the operator norm
\eqref{eq:V91-boundary-bank-norm}; it is exactly the estimate
\eqref{eq:V90-flux-BFV-bound} for the whole boundary Euler/stress bank.
The completed tensor-product conclusion follows as in
Proposition~\ref{prop:V90-boundary-flux-routes}.  Pointwise convergence on a
test space need not be equicontinuous in any fixed seminorm, so the converse
statement follows from the uniform boundedness gap between weak and strong
operator convergence.
\end{proof}

\begin{assessment}[Decision at the V90 flux fork]
\label{ass:V91-flux-decision}
For free Hadamard two-point functions, the first restriction to a timelike
boundary is valid.  The boundary coincidence required by the stress is not.
The image witness proves that ordinary bulk subtraction leaves a divergent
normal jet.  Therefore bulk Hadamard regularity alone does not verify route
\textup{(a)} of Proposition~\ref{prop:V90-boundary-flux-routes}.  The active
route is \textup{(b)}: construct one invariant boundary counterterm action,
take its complete first variation, and prove the BFV Cauchy estimate
\eqref{eq:V91-BFV-Cauchy-gate}.
\end{assessment}

\subsection{V91 result ledger and next acquisition}
\label{subsec:V91-ledger}

\paragraph{New exact conclusions.}
\begin{enumerate}[label=\textup{(\arabic*)},leftmargin=2.8em]
\item Hadamard two-point functions and all finite derivatives restrict to a
      timelike boundary because null wavefront covectors cannot be boundary
      conormals.
\item The boundary diagonal has null conormals; every surviving nonzero
      tangential Hadamard symbol intersects them, so coincidence is a
      separate obstruction.
\item The Dirichlet image remainder has exact singularities
      \(-1/(16\pi^2x^2)\) and
      \(-3/(32\pi^2x^4)\) in a stress derivative composite.
\item Extensions of the \(x^{-4}\) term differ by four boundary-supported
      normal-jet distributions before covariance restrictions.
\item Bulk Hadamard subtraction does not select the boundary stress or its
      normal flux.
\item A convergent complete boundary first variation defines the flux
      directly and preserves the V90 Ward identity when its action is
      invariant.
\item The single norm gate \eqref{eq:V91-BFV-Cauchy-gate} is sufficient to
      land the entire renormalized boundary bank in the V88 test dual.
\end{enumerate}

\paragraph{Next forced step.}
V92 must classify the local gauge- and diffeomorphism-invariant boundary
counterterms at the engineering orders reached by the gauge, Higgs, fermion,
clock, and gravitational normal jets.  It should quotient that finite list
by tangential divergences, compute the complete first-variation matrix, and
identify which coefficients are fixed by cancellation of the image
singularities and which remain physical boundary renormalization conditions.
Only after this common variation is assembled should one attempt
\eqref{eq:V91-BFV-Cauchy-gate}.

\begin{assessment}[V91 continuum status]
\label{ass:V91-continuum-status}
The full Standard--Model--Einstein continuum remains unproved.  V91 decides
the upstream normal-flux fork: ordinary bulk Hadamard subtraction is
insufficient at a timelike boundary, and the variational boundary
renormalization route is required.  The invariant counterterm bank, its
uniform form estimates, the dynamical clock, curved radiative equivariance,
nonlinear existence, and the cofinal interacting limit remain open.
\end{assessment}

\section*{V91 exact checks and append-only record}
\addcontentsline{toc}{section}{V91 exact checks and append-only record}

\begin{equation}
\boxed{
\begin{gathered}
 \operatorname{WF}(W)\cap N^*(\Gamma\times\Gamma)=\varnothing,
 \qquad I^*W\ \hbox{exists},\\
 \operatorname{WF}(I^*PW)\cap N^*\Delta_\Gamma\ne\varnothing
 \quad\hbox{when the tangential symbol survives},\\
 W_{\rm im}(X,X)=-\frac1{16\pi^2x^2},
 \qquad
 \partial_x\partial_{x'}W_{\rm im}|_{X'=X}
 =-\frac3{32\pi^2x^4},\\
 B_N\ \hbox{Cauchy in the BFV form norm}
 \Longrightarrow
 B_{\rm ren}\in
 \mathcal B(\mathscr D_{\rm SM},\mathscr D_{\rm SM}')
 \widehat\otimes\mathscr D_{\rm BFV}'.
\end{gathered}}
\label{eq:V91-check-card}
\end{equation}

The V90 predecessor used for this append has SHA--256
\begin{center}
\texttt{e396b0f4a3679439afc8691878a407c3ed00fba7c0325d13941eb1b8eabfa835}.
\end{center}
The V91 candidate retains its bytes up to the sole live final document
terminator, appends this material, and restores exactly one active
\verb|\end{document}|.  Validation includes balanced environments and
braces, unique labels, resolved new internal references, valid inline and
display math delimiters, exact predecessor-prefix equality, and a final
inventory which removes only the superseded SM V90 file.

% =============================================================================
% END APPENDED VERSION V91 MATERIAL
% =============================================================================


\clearpage
% =============================================================================
% BEGIN APPENDED VERSION V92 MATERIAL
% =============================================================================
\section{V92: the clock-symmetry obstruction, normal-jet variational
order, and the finite counterterm-policy gate}
\label{sec:V92-counterterm-policy}

V91 showed that the boundary extension of a four-dimensional stress
composite has four transverse distributional ambiguity directions before
symmetry and variational restrictions are imposed.  It then requested a
finite invariant counterterm list.  That request is not yet well posed for
the field content declared in V89.  The clock \(\Theta\) was required to
make the speed \(v<1\) covariant, but no internal clock symmetry, engineering
dimension, or allowed normal-jet order was fixed.

This note resolves that upstream ambiguity.  Without clock shift symmetry,
there are infinitely many independent zero-derivative invariant boundary
actions \(\int F(\Theta)\).  Independently, an action using at most \(r\)
normal jets can produce Euler distributions containing at most
\(\delta_\Gamma^{(r)}\).  Thus the four extensions in V91 are not all
available in a first-normal-jet Einstein boundary theory.  The intrinsic
first-jet distributional normal form is computed explicitly.

The result is a finite-policy theorem.  A finite counterterm matrix exists
after one imposes clock shift symmetry, polynomial engineering order, a
finite jet order, and the exact gauge/diffeomorphism symmetries.  Divergent
parts are fixed by cancellation only within the image of the resulting
variation matrix; finite parts in that image remain physical boundary
renormalization conditions.  A regulator demanding a coefficient outside
that image falsifies the policy and forces a higher-jet action or a different
boundary realization.

\subsection{No finite bank without a clock symmetry}
\label{subsec:V92-clock-no-go}

Let the clock be a real scalar with timelike gradient, as in
\eqref{eq:V89-clock-frame}.  Diffeomorphisms act by pullback and therefore
permit, for every smooth real function \(F\), the local boundary action
\begin{equation}
 C_F[\gamma,\Theta]
 :=\int_\Gamma\sqrt{|\gamma|}\,F(\Theta).
 \label{eq:V92-clock-potential-counterterm}
\end{equation}

\begin{theorem}[Infinite clock-potential obstruction]
\label{thm:V92-infinite-clock-bank}
Modulo tangential divergences and additive constants independent of all
fields, the family \(\{C_F:F\in C^\infty(\mathbb R)\}\) is
infinite dimensional.  Consequently diffeomorphism and Standard Model gauge
invariance alone do not yield a finite boundary counterterm bank at any
fixed derivative order.
\end{theorem}

\begin{proof}
A tangential divergence has identically zero Euler--Lagrange
derivative.  The Euler derivative of
\(\sqrt{|\gamma|}F(\Theta)\) with respect to \(\Theta\) is
\(\sqrt{|\gamma|}F'(\Theta)\).  Hence a finite relation among the
counterterm germs, modulo tangential divergences and field-independent
constants, would imply
\begin{equation}
 \sum_{j=1}^N a_jF_j'(c)=0
 \label{eq:V92-function-relation}
\end{equation}
for every value \(c\) in the clock range.  This jet test may be made while
keeping a fixed timelike first derivative of \(\Theta\); it does not set the
clock gradient to zero.  Taking \(F_j(c)=c^j\), \(j\ge1\), makes the derivatives
\(j c^{j-1}\) linearly independent.  Thus arbitrarily large finite
subfamilies of the actions themselves are independent even after the stated
quotient, and the quotient space is infinite dimensional.  Gauge invariance
is automatic because \(\Theta\) is a gauge singlet.
\end{proof}

\begin{corollary}[A clock choice is part of renormalization data]
\label{cor:V92-clock-choice}
A finite V91 counterterm matrix requires at least one of the following extra
inputs: a shift symmetry \(\Theta\mapsto\Theta+c\), a fixed finite-dimensional
clock potential family, or a positive power-counting dimension for
\(\Theta\) together with polynomial locality and a dimension bound.  None was contained in V59 or V89.
\end{corollary}

\begin{proof}
Each listed input removes or finitely truncates the independent family in
Theorem~\ref{thm:V92-infinite-clock-bank}.  Without one of them that theorem
applies.
\end{proof}

\begin{firewall}[No silent clock model]
\label{fire:V92-no-silent-clock-model}
The clock was forced by covariance of the subluminal cone, not supplied as a
completed Standard Model field.  Imposing shift symmetry may be a useful
branch, but it is additional physical structure.  The continuation below
states it explicitly whenever it is used.
\end{firewall}

\subsection{Normal-jet order controls boundary distribution order}
\label{subsec:V92-normal-jet-order}

Let \(x\ge0\) be a Gaussian boundary coordinate with \(\Gamma=\{x=0\}\)
and \(D_n=\partial_x\) at the anchor.  Let \(q\) denote any finite bundle of
bulk fields, including metric components.  Write
\begin{equation}
 q_j:=D_n^jq|_\Gamma,
 \qquad 0\le j\le r,
 \label{eq:V92-normal-jets}
\end{equation}
and consider a local boundary action
\begin{equation}
 C[q]
 =\int_\Gamma L\bigl(
  y,q_0,\ldots,q_r,D_aq_0,\ldots
 \bigr)\,\dd\mu_\gamma,
 \label{eq:V92-r-jet-action}
\end{equation}
with finite tangential jet order.  Tangential integrations by parts are
included in the Euler derivatives below.

\begin{theorem}[Variational support-order theorem]
\label{thm:V92-support-order}
The bulk distribution representing the first variation of
\eqref{eq:V92-r-jet-action} has the form
\begin{equation}
 \boxed{
 \mathcal E_C
 =\sum_{j=0}^r(-D_n)^j
  \left(
   \delta_\Gamma\,\frac{\delta L}{\delta q_j}
  \right),}
 \label{eq:V92-distributional-Euler}
\end{equation}
up to the algebraic metric variation of the measure and of the unit normal,
which has the same maximal normal order.  In particular,
\(\mathcal E_C\) contains no derivative of \(\delta_\Gamma\) of order
greater than \(r\).

Conversely, if a required regulator subtraction has a nonzero
\(\delta_\Gamma^{(j)}\) coefficient with \(j>r\), it cannot be the first
variation of an action whose largest normal jet is \(q_r\).
\end{theorem}

\begin{proof}
After tangential integrations by parts, the first variation is
\begin{equation}
 \delta C
 =\sum_{j=0}^r\int_\Gamma
 \left\langle\frac{\delta L}{\delta q_j},
 D_n^j\delta q\right\rangle\dd\mu_\gamma.
 \label{eq:V92-boundary-first-variation-jets}
\end{equation}
Represent the boundary integral by \(\delta_\Gamma\) in the collar and move
each normal derivative from \(\delta q\) onto the coefficient distribution.
This gives \eqref{eq:V92-distributional-Euler}.  Differentiating
\(\delta_\Gamma\) at most \(j\) times proves the upper order.  Variation of
the normal and measure contains at most the normal order already present in
the Lagrangian and cannot raise this maximum.

For the converse, distributions supported on \(\Gamma\) have a unique
finite expansion in transverse derivatives once a defining function is
fixed.  Formula \eqref{eq:V92-distributional-Euler} has zero coefficient in
every order above \(r\); it therefore cannot equal a distribution with a
nonzero higher coefficient.
\end{proof}

\begin{corollary}[First-jet restriction of the V91 ambiguity]
\label{cor:V92-first-jet-ambiguity}
If the common Einstein--Standard Model boundary action is restricted to the
induced fields and their first normal jets, the V91 ambiguity
\eqref{eq:V91-extension-difference} is variationally admissible only in the
subspace
\begin{equation}
 c_0\delta_\Gamma+c_1D_n\delta_\Gamma.
 \label{eq:V92-first-jet-extension-space}
\end{equation}
The coefficients of \(D_n^2\delta_\Gamma\) and
\(D_n^3\delta_\Gamma\) must be set to zero by the renormalization policy, or
the action must be enlarged to second and third normal jets respectively.
\end{corollary}

\begin{proof}
Apply Theorem~\ref{thm:V92-support-order} with \(r=1\) to the four independent
orders in \eqref{eq:V91-extension-difference}.
\end{proof}

\begin{remark}[Refinement of the V91 wording]
\label{rem:V92-refines-V91}
V91 correctly identified all scaling-degree-preserving distributional
extensions.  Calling all four directions boundary-counterterm variations
implicitly allowed boundary actions with normal jets through order three.
Corollary~\ref{cor:V92-first-jet-ambiguity} records the smaller subspace
available to a first-jet variational theory.
\end{remark}

\subsection{The intrinsic first-normal-jet variation matrix}
\label{subsec:V92-first-jet-matrix}

For \(r=1\), set
\begin{equation}
 A_0:=\frac{\delta L}{\delta q_0},
 \qquad
 A_1:=\frac{\delta L}{\delta q_1},
 \label{eq:V92-first-jet-Euler-coefficients}
\end{equation}
where \(A_0\) includes every tangential Euler derivative.  The symbol
\(D_n\) below acts on distribution densities by duality:
\begin{equation}
 \left\langle D_n(\delta_\Gamma A),f\right\rangle
 :=-\int_\Gamma\langle A,D_nf|_\Gamma\rangle\,\dd\mu_\gamma.
 \label{eq:V92-normal-distribution-definition}
\end{equation}
This definition uses only the boundary coefficient \(A\); it does not choose
an extension of \(A\) off \(\Gamma\).  Equation
\eqref{eq:V92-distributional-Euler} becomes
\begin{equation}
 \boxed{\mathcal E_C
 =\delta_\Gamma A_0-D_n(\delta_\Gamma A_1).}
 \label{eq:V92-first-jet-distribution}
\end{equation}

\begin{proposition}[Exact intrinsic jet-coefficient map]
\label{prop:V92-first-jet-matrix}
Relative to the fixed Gaussian collar, write a first-order supported
distribution density in the intrinsic normal form
\begin{equation}
 \mathcal E_C
 =\delta_\Gamma B_0+D_n(\delta_\Gamma B_1).
 \label{eq:V92-intrinsic-normal-form}
\end{equation}
For the abstract jet variables \(q_0=q|_\Gamma\) and
\(q_1=D_nq|_\Gamma\), the coefficient map is exactly
\begin{equation}
 \boxed{
 \binom{B_0}{B_1}
 =\begin{pmatrix}I&0\\0&-I\end{pmatrix}
  \binom{A_0}{A_1}.}
 \label{eq:V92-variation-matrix}
\end{equation}
It is an isomorphism at the level of unrestricted jet coefficients, with
\(A_0=B_0\) and \(A_1=-B_1\).  It does not assert that an arbitrary pair is
the first variation of one local invariant scalar.
\end{proposition}

\begin{proof}
Comparison of \eqref{eq:V92-first-jet-distribution} with
\eqref{eq:V92-intrinsic-normal-form} proves the matrix formula.  Uniqueness
follows by testing first with functions having prescribed boundary value and
zero normal derivative, and then with functions having zero boundary value
and prescribed normal derivative.  No quantity \(D_nA_1\) appears: such a
quantity would depend on an arbitrary off-boundary extension of \(A_1\).
\end{proof}

Geometric parameterizations require a separate chain rule.  For example, in
Gaussian coordinates
\begin{equation}
 K_{ab}=\tfrac12D_n\gamma_{ab},
 \qquad
 \delta K_{ab}=\tfrac12D_n\delta\gamma_{ab}
                 +\mathcal L_{ab}(\delta\gamma,\delta g_{n\bullet}),
 \label{eq:V92-K-chain-rule}
\end{equation}
where \(\mathcal L\) contains no normal derivative of the variation.  Thus
an Euler coefficient conjugate to \(K\) contributes
\(-\tfrac12D_n(\delta_\Gamma A_K)\) to the principal transverse part and a
collar-dependent but intrinsic zero-order term after the full metric and
normal variations are combined.  This changes the concrete columns of
\(\mathbb V_1\), but cannot create \(D_n^2\delta_\Gamma\).

The coefficient isomorphism in \eqref{eq:V92-variation-matrix} still does not
make arbitrary \((B_0,B_1)\) admissible.  The pair must arise as the complete
Euler derivative of one local invariant scalar.  This imposes Helmholtz
symmetry, the geometric chain rules, and the Ward identities.

\subsection{Ward image of an invariant counterterm action}
\label{subsec:V92-Ward-image}

Collect the boundary fields into
\begin{equation}
 \mathbf q=(\gamma,A_\parallel,\Phi,Psi,\overline\Psi,\Theta),
 \qquad
 \mathbf q_1=D_n\mathbf q.
 \label{eq:V92-complete-boundary-jet}
\end{equation}
Let \(C_\alpha[\mathbf q,\mathbf q_1]\),
\(1\le\alpha\le N\), be invariant local first-jet actions and let
\(V_\alpha\) denote their complete first-variation columns, including
metric stress and every Euler component.  Define
\begin{equation}
 \mathbb V_1:
 \mathbb R^N\longrightarrow\mathscr B_0\oplus\mathscr B_1,
 \qquad
 c\longmapsto\sum_{\alpha=1}^Nc_\alpha V_\alpha,
 \label{eq:V92-counterterm-variation-map}
\end{equation}
where \(\mathscr B_j\) is the coefficient space of
\(D_n^j\delta_\Gamma\) in the complete field bank.

\begin{theorem}[Counterterm columns lie in the common Ward kernel]
\label{thm:V92-counterterm-Ward-kernel}
Let \(\mathbb W\) be the combined gauge and tangential-diffeomorphism Ward
operator acting on the complete supported first-variation bank, tested with
parameters compactly supported away from corners (or with the declared
corner completion included).  Then
\begin{equation}
 \boxed{\mathbb W\mathbb V_1=0.}
 \label{eq:V92-Ward-matrix-zero}
\end{equation}
Thus cancellation of a regulated boundary divergence \(R_N\) by invariant
first-jet counterterms is possible only if
\begin{equation}
 R_N\in\operatorname{Ran}\mathbb V_1
 \subseteq\ker\mathbb W.
 \label{eq:V92-counterterm-range-gate}
\end{equation}
The divergent part of the coefficient vector is fixed up to
\(\ker\mathbb V_1\); finite vectors in the same range are boundary
renormalization conditions and are not selected by bulk Hadamard data.
\end{theorem}

\begin{proof}
Gauge and diffeomorphism invariance of each \(C_\alpha\) makes its first
variation vanish on the corresponding infinitesimal orbit.  After the
integrations by parts used in V90, this is exactly
\(\mathbb WV_\alpha=0\).  Linearity gives
\eqref{eq:V92-Ward-matrix-zero}.  A subtraction produced by these actions is
by definition in their variation range, proving the necessary condition
\eqref{eq:V92-counterterm-range-gate}.  If two coefficient vectors cancel
the same divergence, their difference lies in \(\ker\mathbb V_1\).  Adding a
finite vector changes the finite boundary action while preserving the Ward
identity; that choice is not fixed by the singular bulk parametrix.
\end{proof}

\begin{corollary}[A computable failure certificate]
\label{cor:V92-counterterm-failure-certificate}
If a continuous functional \(\Lambda\) on the supported variation bank
satisfies
\begin{equation}
 \Lambda\mathbb V_1=0,
 \qquad \Lambda R_N\ne0,
 \label{eq:V92-counterterm-annihilator}
\end{equation}
then no allowed invariant first-jet counterterm cancels \(R_N\).  This is the
counterterm analogue of the residual-orbit annihilators in V84 and V87.
\end{corollary}

\begin{proof}
Membership of \(R_N\) in \(\operatorname{Ran}\mathbb V_1\) would imply
\(\Lambda R_N=0\), contradicting
\eqref{eq:V92-counterterm-annihilator}.
\end{proof}

\subsection{A sufficient policy for finite generation}
\label{subsec:V92-finite-policy}

Consider the following explicit policy:
\begin{enumerate}[label=\textup{(P\arabic*)},leftmargin=3.2em]
\item \(\Theta\mapsto\Theta+c\) is a symmetry, and the clock enters only
      through the normalized \(u\) in \eqref{eq:V89-clock-frame} and its
      derivatives.
\item Counterterms are local polynomials in a finite field and jet bank,
      have boundary engineering dimension at most three, and use at most one
      normal derivative.
\item Standard Model gauge invariance and tangential diffeomorphism
      invariance are exact.  Any gauge Chern--Simons term is treated as a
      separate quantized, gauge-invariant-modulo-integers sector rather than
      as an ordinary real coefficient.
\item Tangential divergences, algebraic tensor identities, and terms which
      vanish identically under the declared boundary conditions are
      quotiented before forming \(\mathbb V_1\).
\end{enumerate}

\begin{theorem}[Conditional finite-generation theorem]
\label{thm:V92-finite-generation}
Under policies \textup{(P1)}--\textup{(P4)}, the vector space of admissible
counterterm germs is finite dimensional.  It therefore admits a finite basis
\(\{C_\alpha\}_{\alpha=1}^N\) and a finite complete variation matrix
\(\mathbb V_1\).  The conclusion fails without \textup{(P1)} by
Theorem~\ref{thm:V92-infinite-clock-bank}.
\end{theorem}

\begin{proof}
Every derivative has positive engineering dimension, so the dimension-three
bound permits only finitely many differentiated factors and finitely many
jet multiindices.  All Standard Model fields other than the normalized clock
direction have fixed positive engineering dimension, bounding their
polynomial degree.  Undifferentiated copies of \(u\) carry no dimension, but
\(u_au^a=-1\); after this identity, copies of \(u\) can only occupy the
finitely many free tensor slots supplied by the other factors.  The field
representations and tensor ranks are finite, so only finitely many invariant
contractions remain.  Within this bounded polynomial jet bank, the spans of algebraic identities
and tangential divergences are subspaces of a finite-dimensional space;
quotienting by them preserves finite dimensionality.  The last statement is Theorem
\ref{thm:V92-infinite-clock-bank}.
\end{proof}

\begin{firewall}[The policy is not yet the physical counterterm basis]
\label{fire:V92-policy-scope}
Policies \textup{(P1)}--\textup{(P4)} prove that a finite computation is
possible.  They do not select CP, baryon/lepton-number, neutrino, bag, or
clock-boundary conditions, and they do not enumerate the resulting invariant
basis.  Those choices change \(N\), \(\operatorname{Ran}\mathbb V_1\), and
the physical finite renormalization parameters.
\end{firewall}

\subsection{Accessible and inaccessible parts of the image divergence}
\label{subsec:V92-image-matching}

Let the regulated form of the V91 scalar image singularity have supported
ambiguity vector
\begin{equation}
 R_N^{\rm im}
 =\sum_{j=0}^3r_{j,N}D_n^j\delta_\Gamma.
 \label{eq:V92-image-ambiguity-vector}
\end{equation}

\begin{proposition}[Exact first-jet matching decision]
\label{prop:V92-first-jet-matching}
Under the first-jet policy, cancellation by the common invariant action
requires
\begin{equation}
 r_{2,N}=r_{3,N}=0
 \label{eq:V92-high-jet-zero-test}
\end{equation}
after the bulk finite-part prescription, and
\begin{equation}
 (r_{0,N},r_{1,N})
 \in\operatorname{Ran}\mathbb V_1.
 \label{eq:V92-low-jet-range-test}
\end{equation}
These two conditions are jointly necessary.  If they hold and the selected
preimage coefficients satisfy the V90 common form estimates, their divergent
parts are fixed by the cancellation equation, while the finite kernel and
finite homogeneous solutions remain boundary renormalization data.
\end{proposition}

\begin{proof}
Corollary~\ref{cor:V92-first-jet-ambiguity} proves
\eqref{eq:V92-high-jet-zero-test}; Theorem
\ref{thm:V92-counterterm-Ward-kernel} proves
\eqref{eq:V92-low-jet-range-test}.  Solving the finite linear cancellation
equation fixes a coset modulo \(\ker\mathbb V_1\).  Adding a finite solution
of the homogeneous Ward-compatible equation changes only the finite
renormalization condition.  The V90 form estimates are independently needed
to pass from coefficient cancellation to the operator-valued boundary
limit.
\end{proof}

\begin{assessment}[What V92 changes]
\label{ass:V92-policy-decision}
The next computation cannot be an unconditional finite counterterm table.
The clock sector makes that table infinite until a clock symmetry or finite
potential family is declared.  After an explicit finite policy is chosen,
the renormalization problem becomes a finite range test for the complete
first-variation matrix, preceded by the high-normal-jet zero test
\eqref{eq:V92-high-jet-zero-test}.  This is strictly upstream of estimating
the V91 BFV norm.
\end{assessment}

\subsection{V92 result ledger and next acquisition}
\label{subsec:V92-ledger}

\paragraph{New exact conclusions.}
\begin{enumerate}[label=\textup{(\arabic*)},leftmargin=2.8em]
\item Diffeomorphism and gauge invariance permit infinitely many clock
      potentials unless an additional clock restriction is imposed.
\item A boundary action with normal jets through order \(r\) produces
      supported Euler distributions only through
      \(D_n^r\delta_\Gamma\).
\item A first-jet action accesses only the \(\delta_\Gamma\) and
      \(D_n\delta_\Gamma\) directions of the four V91 extension
      ambiguities.
\item The intrinsic first-jet coefficient map is the exact diagonal map
      \eqref{eq:V92-variation-matrix}; geometric variables then apply their
      explicit chain rules.
\item Complete invariant counterterm columns lie in the joint gauge and
      diffeomorphism Ward kernel.
\item A left annihilator of \(\mathbb V_1\) gives a rigorous certificate
      that a regulator divergence cannot be cancelled within the policy.
\item Policies \textup{(P1)}--\textup{(P4)} are sufficient for finite
      generation, but remain an explicit conditional branch.
\item The exact matching gate is
      \eqref{eq:V92-high-jet-zero-test} plus
      \eqref{eq:V92-low-jet-range-test}, followed by the V90 common form
      estimate.
\end{enumerate}

\paragraph{Next forced step.}
V93 should take the conservative shift-symmetric, first-normal-jet branch as
a declared test policy and construct its invariant basis in blocks.  It
should begin with the scalar/Higgs and pure geometric terms, compute their
metric, clock, and Higgs first variations together, and quotient tangential
divergences before adding fermion bilinears or the quantized gauge
Chern--Simons sector.  Any dependence on CP, neutrino, or bag boundary data
must remain visible as a branch rather than being silently fixed.

\begin{assessment}[V92 continuum status]
\label{ass:V92-continuum-status}
The full Standard--Model--Einstein continuum remains unproved.  V92 proves
that the counterterm classification is infinite without a clock symmetry,
derives the exact normal-jet variation order, and reduces every declared
finite policy to a Ward-compatible range problem.  The physical policy,
explicit invariant basis, uniform quantum form bounds, curved radiative
equivariance, nonlinear existence, and cofinal interacting limit remain
open.
\end{assessment}

\section*{V92 exact checks and append-only record}
\addcontentsline{toc}{section}{V92 exact checks and append-only record}

\begin{equation}
\boxed{
\begin{gathered}
 \{\int_\Gamma\sqrt{|\gamma|}F(\Theta):F\in C^\infty\}
 \quad\hbox{is infinite dimensional},\\
 \mathcal E_C
 =\sum_{j=0}^r(-D_n)^j
  \left(\delta_\Gamma\frac{\delta L}{\delta q_j}\right),\\
 \binom{B_0}{B_1}
 =\begin{pmatrix}I&0\\0&-I\end{pmatrix}
  \binom{A_0}{A_1},
 \qquad \mathbb W\mathbb V_1=0,\\
 r_{2,N}=r_{3,N}=0,
 \qquad (r_{0,N},r_{1,N})\in\operatorname{Ran}\mathbb V_1.
\end{gathered}}
\label{eq:V92-check-card}
\end{equation}

The V91 predecessor used for this append has SHA--256
\begin{center}
\texttt{60b38ee2083d1974cc88891154b95b8e65e311bfaaf676ddd81dd163591e6bd7}.
\end{center}
The V92 candidate retains its bytes up to the sole live final document
terminator, appends this material, and restores exactly one active
\verb|\end{document}|.  Validation includes balanced environments and
braces, unique labels, resolved new internal references, valid inline and
display math delimiters, exact predecessor-prefix equality, and a final
inventory which removes only the superseded SM V91 file.

% =============================================================================
% END APPENDED VERSION V92 MATERIAL
% =============================================================================

\clearpage
% =============================================================================
% BEGIN APPENDED VERSION V93 MATERIAL
% =============================================================================
\section{V93: the algebraic extrinsic--Higgs counterterm block, its complete
first variation, and two anchor range obstructions}
\label{sec:V93-algebraic-counterterms}

V92 proved that a finite counterterm calculation requires an explicit
policy and proposed the shift-symmetric, dimension-three, first-normal-jet
branch as a conservative test.  This note constructs the first nontrivial
block of that branch.  It treats terms which are algebraic in the extrinsic
curvature and the normalized clock direction, contain an undifferentiated
Standard Model Higgs doublet and at most one normal Higgs derivative, and
contain no tangential derivative or gauge curvature.  Intrinsic curvature,
derivatives of the clock direction, tangential Higgs kinetic terms, gauge
curvature, fermions, and ghosts are therefore outside this block and remain
for later versions.

The classification is pointwise and exact.  The clock reduces the stabilizer
of the Lorentzian boundary metric to the spatial group in two dimensions.
Invariant theory for one scalar pair, one spatial vector, and one spatial
traceless symmetric tensor then gives a finite basis through cubic order.
The complete first variation shows why its metric, clock, Higgs, and normal
gauge columns cannot be varied independently.  At the flat zero-Higgs anchor
the resulting normal-symbol image has explicit missing directions, giving
two regulator failure certificates before any BFV norm estimate is attempted.

\subsection{Declared block and engineering dimensions}
\label{subsec:V93-declared-block}

Retain the V89 timelike unit clock direction \(u\), and define the positive
spatial projector
\begin{equation}
 q_{ab}:=\gamma_{ab}+u_au_b,
 \qquad q_a{}^bu_b=0,
 \qquad q_a{}^a=2.
 \label{eq:V93-spatial-projector}
\end{equation}
Decompose the extrinsic curvature into spatial irreducibles:
\begin{align}
 \kappa&:=K_{ab}u^au^b,
 &\vartheta&:=q^{ab}K_{ab},
 \notag\\
 b_a&:=q_a{}^cK_{cd}u^d,
 &\sigma_{ab}&:=q_a{}^cq_b{}^dK_{cd}
                    -\frac12\vartheta q_{ab}.
 \label{eq:V93-K-decomposition-data}
\end{align}
Thus \(b_au^a=0\), \(\sigma_{ab}u^b=0\), and
\(q^{ab}\sigma_{ab}=0\).  Put
\begin{equation}
 p:=b_ab^a,
 \qquad s:=\sigma_{ab}\sigma^{ab},
 \qquad t:=b^a\sigma_{ab}b^b.
 \label{eq:V93-even-K-invariants}
\end{equation}
If a spatial orientation is selected, let
\(\varepsilon^{ab}:=\varepsilon^{cab}u_c\) and define the parity-odd cubic
\begin{equation}
 \widetilde t:=\varepsilon^{ab}b_a\sigma_{bc}b^c.
 \label{eq:V93-odd-K-invariant}
\end{equation}

Let \(H_{\mathbb R}\) be the realification of the Higgs representation,
with invariant real inner product \(g_H\) and complex structure \(J\), so
\(g_H(Jv,w)=-g_H(v,Jw)\).  With
\(\Xi:=D_n^A\Phi\), set
\begin{equation}
 X:=g_H(\Phi,\Phi),
 \qquad Y_R:=g_H(\Phi,\Xi),
 \qquad Y_I:=g_H(J\Phi,\Xi).
 \label{eq:V93-Higgs-invariants}
\end{equation}
The term \(Y_I\) is odd under complex conjugation in this realified
description.  Its admission is a CP/boundary-orientation choice, so it will
remain displayed with an independent switch.

Use the four-dimensional engineering dimensions on boundary restrictions:
\begin{equation}
 [K]=[D]=1,
 \qquad [u]=0,
 \qquad [\Phi]=1,
 \qquad [\Xi]=2.
 \label{eq:V93-engineering-dimensions}
\end{equation}
Let \(\mathfrak C_{K\Phi}^{\le3}\) be the real vector space of local scalar
polynomial germs of dimension at most three which use only
\(\gamma,u,K,\Phi,\Xi\), are invariant under the Standard Model gauge group
and tangential diffeomorphisms, and are quotiented by algebraic identities.
No tangential integration-by-parts relation occurs because this block has no
tangential derivative.

\subsection{Exact invariant basis through dimension three}
\label{subsec:V93-invariant-basis}

\begin{theorem}[Algebraic extrinsic--Higgs basis]
\label{thm:V93-algebraic-basis}
On the spatial-parity-even branch and with \(Y_I\) excluded, a basis of
\(\mathfrak C_{K\Phi}^{\le3}\) is
\begin{align}
 \mathcal B_{K,0}&=\{1\},
 \notag\\
 \mathcal B_{K,1}&=\{\kappa,\vartheta\},
 \notag\\
 \mathcal B_{K,2}&=
 \{\kappa^2,\kappa\vartheta,\vartheta^2,p,s\},
 \notag\\
 \mathcal B_{K,3}&=
 \{\kappa^3,\kappa^2\vartheta,\kappa\vartheta^2,
   \vartheta^3,\kappa p,\vartheta p,\kappa s,\vartheta s,t\},
 \notag\\
 \mathcal B_{\Phi}&=\{X,\kappa X,\vartheta X,Y_R\}.
 \label{eq:V93-even-basis}
\end{align}
It has dimension \(21\).  On an oriented branch without spatial parity,
\(\widetilde t\) supplies one additional independent cubic generator.  On
the CP-unrestricted Higgs branch, \(Y_I\) supplies one additional independent
dimension-three generator.  Allowing both gives dimension \(23\).
\end{theorem}

\begin{proof}
At a point, fix \((\gamma,u)\).  Its spatial orthogonal complement is a
two-dimensional Euclidean representation.  The decomposition
\eqref{eq:V93-K-decomposition-data} identifies \(K\) with two trivial
representations \((\kappa,\vartheta)\), one vector \(b\), and one traceless
symmetric tensor \(\sigma\).  The parity-even orthogonal invariants through
degree three are: arbitrary scalar monomials in
\((\kappa,\vartheta)\); the two quadratic invariants \(p\) and \(s\), each
possibly multiplied by one scalar at cubic order; and the cubic contraction
\(t\).  A traceless symmetric two-by-two matrix obeys
\(\sigma^2=\tfrac12s q\), so it has no independent cubic trace.  Every other
complete contraction reduces to the displayed list.  With only rotations
imposed, the single new pseudoscalar at this order is \(\widetilde t\).

The Standard Model gauge group acts transitively on spheres in the single
Higgs-doublet fibre; its derivative-free real polynomial invariants are
polynomials in \(X\).  The dimension bound permits only \(X\), and its only
products with the extrinsic variables are \(\kappa X\) and
\(\vartheta X\).  A scalar linear in \(\Xi\) must pair \(\Xi\) with
\(\Phi\); the two real pairings are precisely \(Y_R\) and \(Y_I\).

For independence, first vary \(\kappa\) and \(\vartheta\).  Then set
\(\sigma=0\) and vary \(b\), set \(b=0\) and vary \(\sigma\), and finally
hold their norms fixed while changing the relative spatial angle; this
isolates \(t\) and, on the oriented branch, \(\widetilde t\).  Varying
\(\Phi\) and independently prescribing the components of \(\Xi\) along
\(\Phi\) and \(J\Phi\) isolates the Higgs generators.  Hence no nonzero
linear relation exists.
\end{proof}

\begin{corollary}[Coefficient normal form]
\label{cor:V93-coefficient-normal-form}
Every germ in the full oriented, CP-unrestricted block has a unique form
\begin{align}
 f={}&P_{\le3}(\kappa,\vartheta)
 +(a_0+a_\kappa\kappa+a_\vartheta\vartheta)p
 +(c_0+c_\kappa\kappa+c_\vartheta\vartheta)s
 \notag\\
 &+d\,t+\widetilde d\,\widetilde t
 +(e_0+e_\kappa\kappa+e_\vartheta\vartheta)X
 +rY_R+zY_I,
 \label{eq:V93-coefficient-form}
\end{align}
where \(P_{\le3}\) is a polynomial of total degree at most three.  Spatial
parity sets \(\widetilde d=0\), and the declared CP-even branch sets
\(z=0\).
\end{corollary}

\begin{proof}
This is the coordinate expansion in the basis of
Theorem~\ref{thm:V93-algebraic-basis}.
\end{proof}

\subsection{The coupled first-variation column}
\label{subsec:V93-complete-variation}

Consider the action
\begin{equation}
 C_f:=\int_\Gamma\sqrt{|\gamma|}\,f.
 \label{eq:V93-block-action}
\end{equation}
First treat \(\gamma,K,u,\Phi,\Xi\) as independent pointwise variables and
write
\begin{equation}
 \delta(\sqrt{|\gamma|}f)
 =\sqrt{|\gamma|}\left(
  \frac12\widehat\tau^{ab}\delta\gamma_{ab}
  +P^{ab}\delta K_{ab}
  +M_a\delta u^a
  +g_H(H_0,\delta\Phi)
  +g_H(Q,\delta\Xi)
 \right).
 \label{eq:V93-independent-first-variation}
\end{equation}
The hat records that the normalization of \(u\), the definition of \(K\),
and \(\Xi=D_n^A\Phi\) have not yet been varied.

\begin{proposition}[Explicit extrinsic and Higgs momenta]
\label{prop:V93-explicit-momenta}
For the coefficient form \eqref{eq:V93-coefficient-form}, the parity-even
part of the extrinsic momentum is
\begin{align}
 P^{ab}={}&f_\kappa u^au^b+f_\vartheta q^{ab}
 +2f_pu^{(a}b^{b)}+2f_s\sigma^{ab}
 \notag\\
 &+f_t\left(
  2u^{(a}\sigma^{b)}{}_cb^c+b^ab^b-\frac12p q^{ab}
 \right)
 +\widetilde d\,
   \frac{\partial\widetilde t}{\partial K_{ab}}.
 \label{eq:V93-extrinsic-momentum}
\end{align}
Here every partial derivative holds \((\gamma,u,\Phi,\Xi)\) fixed.  The
Higgs momenta are
\begin{align}
 H_0&=2f_X\Phi+r\Xi-zJ\Xi,
 &Q&=r\Phi+zJ\Phi.
 \label{eq:V93-Higgs-momenta}
\end{align}
These formulas include all mixed \(X\kappa\) and \(X\vartheta\) columns.
\end{proposition}

\begin{proof}
At fixed \((\gamma,u)\), differentiation of
\eqref{eq:V93-K-decomposition-data} gives
\begin{align*}
 \delta_K\kappa&=u^au^b\delta K_{ab},
 &\delta_K\vartheta&=q^{ab}\delta K_{ab},\\
 \delta_Kp&=2u^{(a}b^{b)}\delta K_{ab},
 &\delta_Ks&=2\sigma^{ab}\delta K_{ab},\\
 \delta_Kt&=left(
  2u^{(a}\sigma^{b)}{}_cb^c+b^ab^b-\frac12p q^{ab}
 \right)\delta K_{ab}.
\end{align*}
The chain rule proves \eqref{eq:V93-extrinsic-momentum}.  Next,
\(\delta X=2g_H(\Phi,\delta\Phi)\),
\(\delta Y_R=g_H(\Xi,\delta\Phi)+g_H(\Phi,\delta\Xi)\), and
\(\delta Y_I=-g_H(J\Xi,\delta\Phi)+g_H(J\Phi,\delta\Xi)\).
Substitution proves \eqref{eq:V93-Higgs-momenta}.
\end{proof}

Put
\begin{equation}
 N:=\sqrt{-D_a\Theta D^a\Theta}.
 \label{eq:V93-clock-lapse}
\end{equation}
At fixed \(\gamma\), the clock variation of the normalized direction is
\begin{equation}
 \delta_\Theta u^a=-N^{-1}q^{ab}D_b\delta\Theta.
 \label{eq:V93-clock-direction-variation}
\end{equation}
For a Lie-algebra-valued normal gauge variation \(\alpha_n=\delta A_n\),
define the moment-map covector
\begin{equation}
 \langle j_n,\alpha_n\rangle
 :=g_H\bigl(Q,\rho_\Phi(\alpha_n)\Phi\bigr).
 \label{eq:V93-normal-gauge-current}
\end{equation}

\begin{theorem}[Complete supported first variation of the block]
\label{thm:V93-complete-first-variation}
After enforcing \(u=u(\gamma,\Theta)\),
\(K=K(g)\), and \(\Xi=D_n^A\Phi\), the supported Euler columns of
\eqref{eq:V93-block-action} have the intrinsic normal forms
\begin{align}
 \mathcal T_f^{ab}
 &=\delta_\Gamma\tau_{f,0}^{ab}
   -D_n(\delta_\Gamma P^{ab}),
 \label{eq:V93-metric-supported-column}\\
 \mathcal E_{\Phi,f}
 &=\delta_\Gamma H_0-D_n^A(\delta_\Gamma Q),
 \label{eq:V93-Higgs-supported-column}\\
 \mathcal E_{\Theta,f}
 &=\delta_\Gamma D_b\bigl(N^{-1}q^{ab}M_a\bigr),
 \label{eq:V93-clock-supported-column}\\
 \mathcal E_{A_n,f}&=\delta_\Gamma j_n.
 \label{eq:V93-gauge-supported-column}
\end{align}
The metric convention is
\(\delta C_f=\tfrac12\langle\mathcal T_f,\delta\gamma\rangle+\cdots\).
The zero-order tensor \(\tau_{f,0}\) is the unique sum of
\(\widehat\tau\) and the algebraic chain-rule terms from the normalization of
\(u\), the lower-order part of \(\delta K\), the metric dependence of
\(D_n^A\), and the boundary measure.  In particular, the normal derivative
coefficient is exactly \(-P^{ab}\) in the stress convention and no
\(D_n^2\delta_\Gamma\) term occurs.
\end{theorem}

\begin{proof}
Equation \eqref{eq:V92-K-chain-rule} makes the principal metric-jet term
\(\tfrac12P^{ab}D_n\delta\gamma_{ab}\).  Distributional integration by
parts and the factor \(1/2\) in the stress convention give
\eqref{eq:V93-metric-supported-column}; every omitted metric chain-rule term
is algebraic in the variation and therefore belongs to
\(\tau_{f,0}\).  Covariant normal integration by parts of
\(g_H(Q,D_n^A\delta\Phi)\) gives
\eqref{eq:V93-Higgs-supported-column}.  The connection part of
\(\delta\Xi\) is exactly \eqref{eq:V93-normal-gauge-current}.

Finally insert \eqref{eq:V93-clock-direction-variation} into the \(M\)-term
of \eqref{eq:V93-independent-first-variation} and integrate tangentially:
\[
 -\int_\Gamma\sqrt{|\gamma|}\,
 N^{-1}M_aq^{ab}D_b\delta\Theta
 =\int_\Gamma\sqrt{|\gamma|}\,
 D_b(N^{-1}q^{ab}M_a)\delta\Theta.
\]
This proves the clock column.  The four formulas arise from the same action,
so their gauge and diffeomorphism Ward identities hold jointly by
Theorem~\ref{thm:V92-counterterm-Ward-kernel}.
\end{proof}

\begin{firewall}[The normal Higgs term has a gauge column]
\label{fire:V93-no-dropped-gauge-column}
The \(Y_I\) generator is gauge invariant as a complete action, but its
variation with respect to \(A_n\) need not vanish.  Keeping the Higgs column
while discarding \eqref{eq:V93-gauge-supported-column} would break the
complete Noether identity.  The \(Y_R\) contribution to \(j_n\) vanishes by
skew-adjointness of the unitary gauge representation, whereas the \(Y_I\)
contribution is the displayed moment map.
\end{firewall}

\subsection{Injectivity and anchor-level range failures}
\label{subsec:V93-range-tests}

\begin{theorem}[No variational kernel in the algebraic block]
\label{thm:V93-block-injectivity}
The complete first-variation map on
\(\mathfrak C_{K\Phi}^{\le3}\) is injective.  Thus, if a regulated
divergence lies in this block's complete variation range, its coefficient
vector is unique.
\end{theorem}

\begin{proof}
Suppose every first-variation column vanishes as a polynomial germ.  The
Higgs normal coefficient in \eqref{eq:V93-Higgs-supported-column} is
\(Q=r\Phi+zJ\Phi\); varying nonzero \(\Phi\) gives \(r=z=0\).  The metric
normal coefficient then gives \(P=\partial f/\partial K=0\) for all
\((K,\Phi)\).  Hence \(f\) is independent of \(K\), so every coefficient of
the nonconstant extrinsic generators and of \(\kappa X,\vartheta X\)
vanishes.  The remaining polynomial is \(c+e_0X\).  Its Higgs variation
forces \(e_0=0\), and its metric measure variation forces \(c=0\).  Therefore
the kernel is zero.
\end{proof}

At the flat extrinsic anchor \(K=0\), formula
\eqref{eq:V93-extrinsic-momentum} reduces to
\begin{equation}
 P^{ab}|_{K=0}
 =(A_0+A_XX)u^au^b+(B_0+B_XX)q^{ab}
 \label{eq:V93-flat-anchor-metric-image}
\end{equation}
for four constants determined by the coefficients of
\(\kappa,\vartheta,\kappa X,\vartheta X\).  Likewise
\begin{equation}
 Q\in\operatorname{span}_{\mathbb R}\{\Phi,J\Phi\}.
 \label{eq:V93-Higgs-normal-image}
\end{equation}

\begin{corollary}[Two explicit range-obstruction certificates]
\label{cor:V93-anchor-obstructions}
For the V93 block, the following regulated first-normal coefficients cannot
be cancelled.
\begin{enumerate}[label=\textup{(\alph*)},leftmargin=2.8em]
\item At \(K=0\), any metric \(D_n\delta_\Gamma\) coefficient with a nonzero
      spatial traceless part or a nonzero mixed \(u\)-spatial part.
\item At any \(\Phi\), any Higgs \(D_n^A\delta_\Gamma\) coefficient with a
      component real-orthogonal to
      \(\operatorname{span}\{\Phi,J\Phi\}\).  In particular, at
      \(\Phi=0\) the V93 Higgs normal image is zero.
\end{enumerate}
Each forbidden component defines a left annihilator of this block's
variation matrix in the sense of
\eqref{eq:V92-counterterm-annihilator}.
\end{corollary}

\begin{proof}
Part \textup{(a)} follows from
\eqref{eq:V93-flat-anchor-metric-image}, whose two tensors are purely scalar
under the spatial stabilizer.  Part \textup{(b)} follows from
\eqref{eq:V93-Higgs-normal-image}.  Pairing a proposed regulator coefficient
with any missing irreducible component gives a functional which annihilates
every V93 column but not that coefficient.
\end{proof}

\begin{assessment}[What the first finite block decides]
\label{ass:V93-block-decision}
The V92 finite-policy gate is now concrete for the algebraic
extrinsic--Higgs sector: its invariant basis, complete coupled variation,
kernel, and two anchor range obstructions are known.  This block cannot by
itself absorb a general boundary stress or Higgs image divergence.  A
regulator producing either forbidden component must be tested against the
tangential-derivative and intrinsic-curvature blocks before the first-jet
policy is rejected.
\end{assessment}

\subsection{V93 result ledger and next acquisition}
\label{subsec:V93-ledger}

\paragraph{New exact conclusions.}
\begin{enumerate}[label=\textup{(\arabic*)},leftmargin=2.8em]
\item The parity-even algebraic \((K,u,\Phi,D_n^A\Phi)\) block has the
      explicit \(21\)-element basis \eqref{eq:V93-even-basis}.
\item Dropping spatial parity and the CP-even restriction adds the two
      separately identified generators \(\widetilde t\) and \(Y_I\).
\item The extrinsic and Higgs momenta are
      \eqref{eq:V93-extrinsic-momentum}--\eqref{eq:V93-Higgs-momenta}.
\item Metric, clock, Higgs, and normal gauge columns are the coupled
      first variation \eqref{eq:V93-metric-supported-column}--
      \eqref{eq:V93-gauge-supported-column} of one invariant action.
\item The complete variation matrix has trivial kernel on this block.
\item Equations \eqref{eq:V93-flat-anchor-metric-image} and
      \eqref{eq:V93-Higgs-normal-image} give two computable range failures.
\end{enumerate}

\paragraph{Next forced step.}
V94 should classify the tangential two-derivative scalar block generated by
intrinsic curvature, the hypersurface-orthogonal clock congruence, and
covariant Higgs derivatives.  It must quotient the Raychaudhuri and
tangential integration-by-parts identities before counting generators, and
then determine which of the V93 anchor annihilators survive.  Gauge curvature
and fermion bilinears should remain separate until that quotient is fixed.

\begin{assessment}[V93 continuum status]
\label{ass:V93-continuum-status}
The full Standard--Model--Einstein continuum remains unproved.  V93 closes
one finite counterterm block but not the tangential derivative, curvature,
gauge, fermion, ghost, quantum form-domain, curved radiative equivariance,
nonlinear existence, or cofinal interacting-limit gates.
\end{assessment}

\section*{V93 exact checks and append-only record}
\addcontentsline{toc}{section}{V93 exact checks and append-only record}

\begin{equation}
\boxed{
\begin{gathered}
 \dim\mathfrak C_{K\Phi}^{\le3,+}=21,
 \qquad
 \dim\mathfrak C_{K\Phi}^{\le3,\mathrm{or},\mathrm{CP\,free}}=23,\\
 \mathcal T_f^{ab}=\delta_\Gamma\tau_{f,0}^{ab}
        -D_n(\delta_\Gamma P^{ab}),
 \qquad
 \mathcal E_{\Phi,f}=\delta_\Gamma H_0-D_n^A(\delta_\Gamma Q),\\
 \ker\mathbb V_{K\Phi}=0,
 \qquad
 P^{ab}|_{K=0}\in\operatorname{span}\{u^au^b,q^{ab}\},
 \qquad
 Q\in\operatorname{span}_{\mathbb R}\{\Phi,J\Phi\}.
\end{gathered}}
\label{eq:V93-check-card}
\end{equation}

The V92 predecessor used for this append has SHA--256
\begin{center}
\texttt{41a98730c0480ed52760c076a1fd81ad3fd8329f648781d1cdfea49e6968e95f}.
\end{center}
The V93 candidate retains its bytes up to the sole live final document
terminator, appends this material, and restores exactly one active
\verb|\end{document}|.  Validation includes balanced environments and
braces, unique labels, resolved new internal references, valid inline and
display math delimiters, exact predecessor-prefix equality, and a final
inventory which removes only the superseded SM V92 file.

% =============================================================================
% END APPENDED VERSION V93 MATERIAL
% =============================================================================

\clearpage
% =============================================================================
% BEGIN APPENDED VERSION V94 MATERIAL
% =============================================================================
\section{V94: the hypersurface-orthogonal tangential seed quotient and the
persistence of the first-normal range obstruction}
\label{sec:V94-tangential-seed}

V93 classified the algebraic extrinsic--Higgs counterterm block and found
missing first-normal metric and Higgs directions at the flat anchor.  Before
enlarging the action by every dimension-three mixed monomial, the V92 policy
requires the tangential integration-by-parts quotient to be fixed.  This
note performs that quotient for the intrinsic two-derivative clock--geometry
seed and for the one-tangential-derivative Higgs drift.

The result is small.  Hypersurface orthogonality leaves acceleration,
expansion, and shear.  Raychaudhuri makes the contraction
\(R_{ab}u^au^b\) dependent modulo a divergence.  The parity-even pure seed
therefore has four generators.  The Higgs drift has two real generators,
one of which is equivalent to \((D_au^a)X\).  None contains a normal jet, so
none changes the \(D_n\delta_\Gamma\) image.  The two V93 anchor
annihilators consequently survive this enlargement exactly.

\subsection{Clock-congruence decomposition and Raychaudhuri quotient}
\label{subsec:V94-clock-decomposition}

Continue to use the V89 normalized clock direction and the V93 spatial
projector.  Define
\begin{equation}
 a_b:=u^aD_au_b,
 \qquad
 \theta:=D_au^a,
 \qquad
 \chi_{ab}:=q_a{}^cq_b{}^dD_{(c}u_{d)}
             -\frac12\theta q_{ab}.
 \label{eq:V94-aether-data}
\end{equation}
Because \(u_a\) is proportional to the exact one-form \(D_a\Theta\), its
projected vorticity vanishes.  Hence
\begin{equation}
 D_au_b=-u_aa_b+\frac12\theta q_{ab}+\chi_{ab}.
 \label{eq:V94-aether-decomposition}
\end{equation}

\begin{lemma}[Three quadratic first-derivative invariants]
\label{lem:V94-three-aether-invariants}
Every parity-even scalar quadratic in \(D u\) is a linear combination of
\begin{equation}
 a^2:=a_aa^a,
 \qquad \theta^2,
 \qquad \chi^2:=\chi_{ab}\chi^{ab}.
 \label{eq:V94-three-aether-scalars}
\end{equation}
These three germs are independent modulo tangential divergences.
\end{lemma}

\begin{proof}
Insert \eqref{eq:V94-aether-decomposition} into a complete contraction of
two copies of \(D_au_b\).  Orthogonality and tracelessness remove all cross
terms, leaving only \eqref{eq:V94-three-aether-scalars}.  Independence may
be tested in foliation-adapted local jets: the spatial gradient of the lapse
sets \(a\), while the trace and traceless parts of the spatial second
fundamental form set \(\theta\) and \(\chi\) independently.  Compactly
supported high-frequency realizations of these jets make the three
quadratic principal forms independent.  A tangential divergence has zero
Euler principal form, so no nonzero linear combination is a divergence.
\end{proof}

The timelike Raychaudhuri identity in boundary dimension three is
\begin{equation}
 u^aD_a\theta
 =-\frac12\theta^2-\chi^2-R_{ab}[\gamma]u^au^b+D_aa^a.
 \label{eq:V94-Raychaudhuri}
\end{equation}

\begin{proposition}[Curvature contraction quotient]
\label{prop:V94-curvature-quotient}
Modulo tangential divergences,
\begin{equation}
 R_{ab}[\gamma]u^au^b
 \equiv \frac12\theta^2-\chi^2.
 \label{eq:V94-Ruu-quotient}
\end{equation}
Thus \(R_{ab}u^au^b\) is not an additional counterterm generator in the
quotient used by V92.
\end{proposition}

\begin{proof}
Solve \eqref{eq:V94-Raychaudhuri} for \(R_{ab}u^au^b\).  The term
\(D_aa^a\) is a divergence, and
\begin{equation*}
 u^aD_a\theta=D_a(\theta u^a)-\theta^2
 \equiv-\theta^2.
\end{equation*}
Substitution gives \eqref{eq:V94-Ruu-quotient}.
\end{proof}

\begin{theorem}[Pure tangential seed basis]
\label{thm:V94-pure-seed-basis}
Consider parity-even scalar germs of engineering dimension two made from
\((\gamma,u)\), with either one curvature tensor, two first derivatives of
\(u\), or one second derivative of \(u\).  Modulo tangential divergences and
the unit/hypersurface-orthogonality identities, a basis is
\begin{equation}
 \boxed{\mathcal B_{\mathrm{tan},2}
 =\{R[\gamma],a^2,\theta^2,\chi^2\}.}
 \label{eq:V94-pure-seed-basis}
\end{equation}
\end{theorem}

\begin{proof}
A curvature contraction is a combination of \(R\) and \(R_{ab}u^au^b\).
Proposition~\ref{prop:V94-curvature-quotient} removes the second.  Lemma
\ref{lem:V94-three-aether-invariants} handles terms quadratic in first
derivatives.  Integrating a term linear in a second derivative once reduces
it to the same quadratic bank plus a divergence.  The four displayed
classes are independent: \(R\) has the nonzero three-dimensional
Einstein--Hilbert metric Euler derivative, while the other three have the
independent clock-foliation principal forms isolated in the proof of
Lemma~\ref{lem:V94-three-aether-invariants}.
\end{proof}

\begin{remark}[Why the dimension matters]
\label{rem:V94-dimension-three-boundary}
The intrinsic boundary has dimension three.  The Einstein--Hilbert density
is therefore not topological, so its Euler derivative distinguishes
\(R[\gamma]\) from a divergence.  The corresponding assertion would fail
for a two-dimensional boundary.
\end{remark}

\subsection{The tangential Higgs drift quotient}
\label{subsec:V94-Higgs-drift}

At engineering dimension three, the only gauge-invariant real scalars
linear in one tangential Higgs derivative and using no other positive-
dimension factor are
\begin{equation}
 Z_R:=u^ag_H(\Phi,D_a^A\Phi),
 \qquad
 Z_I:=u^ag_H(J\Phi,D_a^A\Phi).
 \label{eq:V94-Higgs-drifts}
\end{equation}

\begin{lemma}[Real drift integration by parts]
\label{lem:V94-real-drift-IBP}
With \(X=g_H(\Phi,\Phi)\),
\begin{equation}
 Z_R=\frac12u^aD_aX
 \equiv-\frac12\theta X.
 \label{eq:V94-real-drift-quotient}
\end{equation}
The complex-structure drift \(Z_I\) has no analogous reduction because it
is the contraction of \(u\) with the Higgs phase current.
\end{lemma}

\begin{proof}
Metric and gauge compatibility give \(D_aX=2g_H(\Phi,D_a^A\Phi)\).
Then
\(D_a(Xu^a)=u^aD_aX+\theta X\), proving the quotient formula.  The
one-form \(g_H(J\Phi,D_a^A\Phi)\) is not the derivative of a gauge-invariant
real scalar; at a nonzero Higgs value its phase-direction jet can be varied
independently of \(D_aX\).  Hence the same integration-by-parts reduction
does not apply to \(Z_I\).
\end{proof}

Consequently the new Higgs-drift quotient is represented by
\begin{equation}
 \mathcal B_{\Phi,\mathrm{tan},3}=\{\theta X,Z_I\},
 \label{eq:V94-Higgs-drift-basis}
\end{equation}
with \(Z_I\) again retained as an explicit CP-sensitive branch.  The term
\(\theta X\) is present even though \(\theta\) alone is a divergence.

For later variation tests, let
\begin{equation}
 C_{\lambda,z}^{\mathrm{drift}}
 :=\int_\Gamma\sqrt{|\gamma|}\,
   (\lambda\theta X+zZ_I).
 \label{eq:V94-drift-action}
\end{equation}

\begin{proposition}[Higgs and clock Euler columns of the drift]
\label{prop:V94-drift-variation}
Ignoring only the algebraically determined metric stress, the Euler
coefficients of \eqref{eq:V94-drift-action} are
\begin{align}
 E_\Phi={}&2\lambda\theta\Phi
 -z\bigl(2J D_u^A\Phi+\theta J\Phi\bigr),
 \label{eq:V94-drift-Higgs-Euler}\\
 E_\Theta^{(\lambda)}={}&
 -\lambda D_b\bigl(N^{-1}q^{ab}D_aX\bigr),
 \label{eq:V94-drift-clock-Euler}
\end{align}
while the \(zZ_I\) clock coefficient is obtained from the same normalized-
direction formula with
\(M_a=z g_H(J\Phi,D_a^A\Phi)\).  The tangential gauge current is defined by
\begin{equation}
 \langle j^a,\alpha\rangle
 :=z u^a g_H(J\Phi,\rho_\Phi(\alpha)\Phi).
 \label{eq:V94-drift-gauge-current}
\end{equation}
All these contributions are proportional to \(\delta_\Gamma\); none has a
normal derivative of that distribution.
\end{proposition}

\begin{proof}
For the \(\lambda\)-term, vary \(X\) and integrate the variation of
\(D_au^a\) tangentially.  Using
\eqref{eq:V93-clock-direction-variation} gives
\eqref{eq:V94-drift-clock-Euler}.  For the \(z\)-term,
skew-adjointness of \(J\), covariant integration by parts, and
\(D_au^a=\theta\) give
\begin{equation*}
 \delta_\Phi\int\sqrt{|\gamma|}Z_I
 =-\int\sqrt{|\gamma|}\,
 g_H(2J D_u^A\Phi+\theta J\Phi,\delta\Phi).
\end{equation*}
The gauge and clock variations give the stated momenta.  Every derivative
integrated by parts is tangential, proving the final support-order claim.
\end{proof}

\subsection{Persistence of the V93 normal-symbol obstruction}
\label{subsec:V94-persistent-obstruction}

Let \(\mathbb V_{K\Phi}\) be the V93 variation map and let
\(\mathbb V_{\mathrm{seed}}\) contain the four actions from
\eqref{eq:V94-pure-seed-basis} and the two branches in
\eqref{eq:V94-Higgs-drift-basis}.

\begin{theorem}[Tangential seeds do not enlarge the first-normal image]
\label{thm:V94-normal-image-persistence}
The projection of \(\operatorname{Ran}\mathbb V_{\mathrm{seed}}\) onto the
coefficient space of \(D_n\delta_\Gamma\) is zero.  Therefore
\begin{equation}
 \operatorname{pr}_1\operatorname{Ran}
 (\mathbb V_{K\Phi}\oplus\mathbb V_{\mathrm{seed}})
 =\operatorname{pr}_1\operatorname{Ran}\mathbb V_{K\Phi}.
 \label{eq:V94-normal-image-unchanged}
\end{equation}
In particular, both range obstructions in
Corollary~\ref{cor:V93-anchor-obstructions} remain valid after the V94
enlargement.
\end{theorem}

\begin{proof}
Every V94 generator depends on boundary fields and tangential jets only.
Its first variation requires only tangential integrations by parts and is a
coefficient times \(\delta_\Gamma\).  The V92 support-order formula thus has
\(A_1=0\) for every field.  Projection onto the first-normal coefficient is
zero, which proves \eqref{eq:V94-normal-image-unchanged} and the last claim.
\end{proof}

\begin{corollary}[The next block is forced]
\label{cor:V94-next-block-forced}
Repair of a forbidden V93 metric normal coefficient requires a
dimension-three generator which contains \(K\), for example \(K\) times one
of the dimension-two tangential seeds, or a term containing a tangential
derivative of \(K\).  Repair of a forbidden Higgs normal coefficient
requires a generator containing \(D_n^A\Phi\) with additional admissible
fibre structure.  Purely intrinsic counterterms cannot repair either defect.
\end{corollary}

\begin{proof}
This is the contrapositive of
Theorem~\ref{thm:V94-normal-image-persistence}, combined with the first-jet
support-order theorem \ref{thm:V92-support-order}.
\end{proof}

\begin{firewall}[No premature rejection of the first-jet policy]
\label{fire:V94-no-premature-rejection}
The V93 annihilators now rule out the algebraic and pure tangential seed
blocks, but not the whole dimension-three first-jet policy.  Mixed
\(K R\), \(K(Du)^2\), derivative-of-\(K\), gauge, and fermion structures
have not yet been classified.  A regulator component in a V93 missing
direction rejects the policy only if it also lies outside those remaining
blocks.
\end{firewall}

\subsection{V94 result ledger and next acquisition}
\label{subsec:V94-ledger}

\paragraph{New exact conclusions.}
\begin{enumerate}[label=\textup{(\arabic*)},leftmargin=2.8em]
\item Hypersurface orthogonality gives the decomposition
      \eqref{eq:V94-aether-decomposition} and three independent quadratic
      clock invariants.
\item Raychaudhuri gives the quotient identity
      \eqref{eq:V94-Ruu-quotient}.
\item The parity-even pure two-derivative seed basis is the four-element
      list \eqref{eq:V94-pure-seed-basis}.
\item The tangential Higgs drift quotient is represented by
      \(\theta X\) and the CP-sensitive phase drift \(Z_I\).
\item Their Higgs, clock, and gauge Euler columns are explicitly coupled by
      \eqref{eq:V94-drift-Higgs-Euler}--
      \eqref{eq:V94-drift-gauge-current}.
\item The V94 block has zero first-normal projection, so both V93 anchor
      annihilators persist.
\end{enumerate}

\paragraph{Next forced step.}
V95 is an audit version.  It must first inspect the newest Yang--Mills and
Navier--Stokes notes in the shared folder.  It should then classify the
dimension-three mixed metric-normal block generated by \(K\) times the V94
seeds and by tangential derivatives of \(K\), quotienting Codazzi and
integration-by-parts identities before testing whether the spatial
traceless and mixed normal directions enter the range.

\begin{assessment}[V94 continuum status]
\label{ass:V94-continuum-status}
The full Standard--Model--Einstein continuum remains unproved.  V94 closes
the intrinsic two-derivative and Higgs-drift seed quotient and proves that
it cannot repair the current first-normal range defects.  The mixed
dimension-three, gauge, fermion, ghost, quantum form-domain, curved
radiative-equivariance, nonlinear-existence, and cofinal interacting-limit
gates remain open.
\end{assessment}

\section*{V94 exact checks and append-only record}
\addcontentsline{toc}{section}{V94 exact checks and append-only record}

\begin{equation}
\boxed{
\begin{gathered}
 R_{ab}u^au^b\equiv\frac12\theta^2-\chi^2,
 \qquad
 \mathcal B_{\mathrm{tan},2}=\{R,a^2,\theta^2,\chi^2\},\\
 Z_R\equiv-\frac12\theta X,
 \qquad
 \mathcal B_{\Phi,\mathrm{tan},3}=\{\theta X,Z_I\},\\
 \operatorname{pr}_1\operatorname{Ran}\mathbb V_{\mathrm{seed}}=0,
 \qquad
 \operatorname{pr}_1\operatorname{Ran}
 (\mathbb V_{K\Phi}\oplus\mathbb V_{\mathrm{seed}})
 =\operatorname{pr}_1\operatorname{Ran}\mathbb V_{K\Phi}.
\end{gathered}}
\label{eq:V94-check-card}
\end{equation}

The V93 predecessor used for this append has SHA--256
\begin{center}
\texttt{3795303205fe58586169dc8dd473157a6eedeaa2bff73e071daaf84fdf1f11d1}.
\end{center}
The V94 candidate retains its bytes up to the sole live final document
terminator, appends this material, and restores exactly one active
\verb|\end{document}|.  Validation includes balanced environments and
braces, unique labels, resolved new internal references, valid inline and
display math delimiters, exact predecessor-prefix equality, and a final
inventory which removes only the superseded SM V93 file.

% =============================================================================
% END APPENDED VERSION V94 MATERIAL
% =============================================================================

\clearpage
% =============================================================================
% BEGIN APPENDED VERSION V95 MATERIAL
% =============================================================================
\section{V95: companion audit, the common-coefficient range theorem, and
the forced Higgs-normal/gauge-current compatibility}
\label{sec:V95-common-coefficient}

V94 showed that purely tangential counterterms do not enlarge the
first-normal image and proposed the mixed dimension-three block next.  This
is also the fifth version after the V90 companion audit, so the shared
Yang--Mills and Navier--Stokes notes were inspected before that larger
classification was begun.

The audit exposes an upstream linear-algebra issue.  A common invariant
counterterm action has one coefficient vector.  Solving its metric, clock,
Higgs, and gauge cancellation equations separately can produce different
coefficient vectors and therefore does not prove cancellation by one action.
This note gives the exact common-preimage criterion.  Applied to the V93
normal Higgs generators, it yields a pointwise compatibility equation
between the Higgs \(D_n^A\delta_\Gamma\) coefficient and the normal gauge
current.  This equation is a new left-annihilator test for every regulator.

\subsection{Required V95 companion audit}
\label{subsec:V95-audit}

The read-only snapshot used for this audit was
\begin{center}
\begin{tabular}{@{}lll@{}}
programme & file and byte count & SHA--256\\
\hline
Yang--Mills &
\texttt{Renewal\_Yang\_Mills\_Mass\_Gap\_Research\_Notes\_V94.tex},
\(2260562\) &
\texttt{b4e48dc6ef69e39aeb804b49832d2d8e7c23d0511100b685dd3ff772784068f9}\\
Navier--Stokes &
\texttt{Renewal\_NS\_Stability\_Research\_Notes\_V31.tex},
\(887537\) &
\texttt{f1121b62238ad5491bb7cf03635ea9934942edf573ed8faaa9ee79e1d38a6696}
\end{tabular}
\end{center}
The NS file is unchanged from the V90 audit.  The YM file advanced from V88
to V94.

The YM V94 result which transfers structurally is its common-effect rule:
two noncommuting orientation effects are added on the same original product
state before the final supremum.  Separate suprema lose their compatibility.
Its spectator quotient also removes identity spectators before estimating
the active capacity.  The NS V31 result which transfers structurally is the
balanced-helicity nullspace: a signed polarization statistic can vanish
while a positive total critical coarea remains.  Thus a projected or signed
sector test need not see the complete obstruction.

\begin{firewall}[Audit results are analogies, not imported estimates]
\label{fire:V95-audit-scope}
No Yang--Mills capacity bound or Navier--Stokes coarea estimate is used as an
Einstein--Standard Model theorem.  The only imported design rule is to keep
the common underlying datum before projection.  In the present programme
that datum is the single coefficient vector of one invariant counterterm
action.  All statements below are proved directly in its finite-dimensional
variation space.
\end{firewall}

\subsection{One coefficient vector for every field sector}
\label{subsec:V95-stacked-map}

Let \(C\) be a finite-dimensional declared counterterm coefficient space.
For each complete field sector \(i\in I\), let
\begin{equation}
 V_i:C\longrightarrow B_i
 \label{eq:V95-sector-variation-map}
\end{equation}
be the corresponding projection of the first-variation map, and put
\begin{equation}
 V:=\bigoplus_{i\in I}V_i:
 C\longrightarrow\bigoplus_{i\in I}B_i.
 \label{eq:V95-stacked-variation-map}
\end{equation}
A regulator supplies a target
\(R=(R_i)_{i\in I}\).

\begin{theorem}[Common-coefficient cancellation criterion]
\label{thm:V95-common-preimage}
Cancellation by one counterterm action is equivalent to
\begin{equation}
 \boxed{R\in\operatorname{Ran}V.}
 \label{eq:V95-common-range}
\end{equation}
The sectorwise conditions
\begin{equation}
 R_i\in\operatorname{Ran}V_i
 \qquad(i\in I)
 \label{eq:V95-sectorwise-ranges}
\end{equation}
are necessary but not sufficient.

For two sectors, choose any sectorwise preimages
\(c_1,c_2\in C\), so \(V_1c_1=R_1\) and \(V_2c_2=R_2\).  A common preimage
exists if and only if
\begin{equation}
 c_1-c_2\in\ker V_1+\ker V_2.
 \label{eq:V95-two-sector-coset-test}
\end{equation}
This condition is independent of the chosen sectorwise preimages.
\end{theorem}

\begin{proof}
Equation \eqref{eq:V95-common-range} is the definition of one vector
\(c\in C\) satisfying every equation \(V_ic=R_i\).  It immediately implies
\eqref{eq:V95-sectorwise-ranges}.  Conversely, separate preimages may differ,
so the sectorwise statement alone does not construct such a \(c\).

For two sectors, a common vector has the form
\(c=c_1+k_1=c_2+k_2\), with \(k_i\in\ker V_i\).  Hence
\(c_1-c_2=k_2-k_1\in\ker V_1+\ker V_2\).  If the latter membership holds,
write \(c_1-c_2=k_2-k_1\); then
\(c_1+k_1=c_2+k_2\) is a common preimage.  Replacing \(c_i\) by another
preimage adds an element of \(\ker V_i\), leaving the membership condition
unchanged.
\end{proof}

\begin{corollary}[Stacked annihilator certificate]
\label{cor:V95-stacked-annihilator}
If \(\Lambda\in(\bigoplus_iB_i)^*\) satisfies
\begin{equation}
 \Lambda V=0,
 \qquad \Lambda R\ne0,
 \label{eq:V95-stacked-annihilator}
\end{equation}
then no common counterterm action cancels \(R\), even if every sectorwise
condition \eqref{eq:V95-sectorwise-ranges} holds.
\end{corollary}

\begin{proof}
If \(R=Vc\), then \(\Lambda R=\Lambda Vc=0\), a contradiction.
\end{proof}

\begin{corollary}[Exact spectator quotient]
\label{cor:V95-spectator-quotient}
The coefficient directions invisible to every selected sector form
\begin{equation}
 K_{\rm sp}:=\bigcap_{i\in I}\ker V_i=\ker V.
 \label{eq:V95-common-kernel}
\end{equation}
The stacked map factors injectively through \(C/K_{\rm sp}\), and range
testing is unchanged by this quotient.  For the complete V93 algebraic
block, \(K_{\rm sp}=0\) by
Theorem~\ref{thm:V93-block-injectivity}.
\end{corollary}

\begin{proof}
The kernel of a direct-sum map is the intersection of the component
kernels.  The first isomorphism theorem gives the injective factor map and
preserves its range.  The final statement is the V93 injectivity theorem.
\end{proof}

\subsection{A forced Higgs-normal/gauge-current relation}
\label{subsec:V95-Higgs-gauge-compatibility}

Return to the two V93 normal Higgs generators
\(rY_R+zY_I\).  Their Higgs normal momentum and normal gauge current are
\begin{align}
 Q&=r\Phi+zJ\Phi,
 \label{eq:V95-Q-rz}\\
 \langle j_n,\alpha\rangle
 &=z\,g_H(J\Phi,\rho_\Phi(\alpha)\Phi),
 \label{eq:V95-j-rz}
\end{align}
because \(g_H(\Phi,\rho_\Phi(\alpha)\Phi)=0\) for the skew-adjoint unitary
representation.

\begin{theorem}[Normal Higgs--gauge compatibility]
\label{thm:V95-Higgs-gauge-compatibility}
Let a regulator target in this V93 subblock be written intrinsically as
\begin{equation}
 \mathcal E_{\Phi,N}
 =\delta_\Gamma H_N-D_n^A(\delta_\Gamma Q_N),
 \qquad j_{n,N}\in\mathfrak g^*.
 \label{eq:V95-regulated-Higgs-gauge-pair}
\end{equation}
If \(\Phi\ne0\), cancellation by one real combination of \(Y_R,Y_I\)
requires unique real numbers \(q_R,q_I\) such that
\begin{equation}
 Q_N=q_R\Phi+q_IJ\Phi
 \label{eq:V95-Q-decomposition}
\end{equation}
and the simultaneous gauge equation
\begin{equation}
 \boxed{
 \langle j_{n,N},\alpha\rangle
 =q_Ig_H(J\Phi,\rho_\Phi(\alpha)\Phi)
 \quad\hbox{for every }\alpha\in\mathfrak g.}
 \label{eq:V95-forced-gauge-current}
\end{equation}
At \(\Phi=0\), both accessible quantities vanish.  Thus separate Higgs and
gauge subtractions do not establish a common invariant subtraction.
\end{theorem}

\begin{proof}
Equation \eqref{eq:V93-Higgs-momenta} gives
\eqref{eq:V95-Q-rz}.  For nonzero \(\Phi\), the real vectors
\(\Phi,J\Phi\) are orthogonal and have equal positive norm, so the
coefficients in \eqref{eq:V95-Q-decomposition} are unique and force
\(r=q_R\), \(z=q_I\).  Substitution in
\eqref{eq:V93-normal-gauge-current} gives
\eqref{eq:V95-forced-gauge-current}.  At \(\Phi=0\), both formulas are zero.
\end{proof}

\begin{corollary}[Pointwise left annihilator]
\label{cor:V95-Higgs-gauge-annihilator}
For every gauge parameter \(\alpha\), the functional
\begin{equation}
 \Lambda_{\alpha,\Phi}(Q,j)
 :=\langle j,\alpha\rangle
   -g_H(Q,\rho_\Phi(\alpha)\Phi)
 \label{eq:V95-explicit-Higgs-gauge-annihilator}
\end{equation}
annihilates the complete \((Y_R,Y_I)\) variation image.  Therefore a target
with \(\Lambda_{\alpha,\Phi}(Q_N,j_{n,N})\ne0\) for one \(\alpha\) fails the
common V93 range test.
\end{corollary}

\begin{proof}
Insert \(Q=r\Phi+zJ\Phi\).  Skew-adjointness removes the \(r\)-term, and
the remaining term equals \eqref{eq:V95-j-rz}.
\end{proof}

\begin{assessment}[Direction after the audit]
\label{ass:V95-audit-direction}
The mixed dimension-three block must be classified as complete stacked
columns, not as independent metric, clock, Higgs, and gauge lists.  The V94
YM audit justifies retaining one common datum before projection; the theorem
above supplies the direct SM proof.  Any regulator calculation should now
print its full target vector and test the stacked annihilators before a
sector norm is estimated.
\end{assessment}

\subsection{V95 result ledger and next acquisition}
\label{subsec:V95-ledger}

\paragraph{New exact conclusions.}
\begin{enumerate}[label=\textup{(\arabic*)},leftmargin=2.8em]
\item The required YM/NS companion audit is recorded with exact file hashes.
\item Common cancellation is the stacked range condition
      \eqref{eq:V95-common-range}; sectorwise range membership is weaker.
\item The two-sector compatibility test is the kernel-sum condition
      \eqref{eq:V95-two-sector-coset-test}.
\item The common invisible coefficient space is exactly
      \(\bigcap_i\ker V_i\), and it is zero for the complete V93 block.
\item The V93 \(Y_R,Y_I\) columns force the regulator compatibility
      \eqref{eq:V95-forced-gauge-current}.
\item Equation \eqref{eq:V95-explicit-Higgs-gauge-annihilator} is a
      pointwise computable failure certificate.
\end{enumerate}

\paragraph{Next forced step.}
V96 should classify the parity-even mixed dimension-three generators formed
from one extrinsic-curvature factor and the V94 dimension-two seeds.  It must
compute their entire stacked first variations and test the spatial
traceless and mixed metric normal directions at \(K=0\).  Derivative-of-
\(K\), gauge, and fermion blocks should be added only after this algebraic
mixed image is known.

\begin{assessment}[V95 continuum status]
\label{ass:V95-continuum-status}
The full Standard--Model--Einstein continuum remains unproved.  V95 closes
the common-coefficient compatibility gate and one coupled Higgs--gauge
annihilator, but the remaining counterterm blocks, regulator target,
uniform quantum form bounds, curved radiative equivariance, nonlinear
existence, and cofinal interacting limit remain open.
\end{assessment}

\section*{V95 exact checks and append-only record}
\addcontentsline{toc}{section}{V95 exact checks and append-only record}

\begin{equation}
\boxed{
\begin{gathered}
 R\in\operatorname{Ran}\bigoplus_iV_i,
 \qquad
 c_1-c_2\in\ker V_1+\ker V_2,\\
 \ker\bigoplus_iV_i=\bigcap_i\ker V_i,
 \qquad
 \Lambda V=0,\ \Lambda R\ne0\Longrightarrow R\notin\operatorname{Ran}V,\\
 Q_N=q_R\Phi+q_IJ\Phi,
 \qquad
 \langle j_{n,N},\alpha\rangle
 =q_Ig_H(J\Phi,\rho_\Phi(\alpha)\Phi).
\end{gathered}}
\label{eq:V95-check-card}
\end{equation}

The V94 predecessor used for this append has SHA--256
\begin{center}
\texttt{b8194cf93631df92c6442ab9bbb4a775bf1b95a4b31f9af4dd15f8761fdf291e}.
\end{center}
The V95 candidate retains its bytes up to the sole live final document
terminator, appends this material, and restores exactly one active
\verb|\end{document}|.  Validation includes balanced environments and
braces, unique labels, resolved new internal references, valid inline and
display math delimiters, exact predecessor-prefix equality, and a final
inventory which removes only the superseded SM V94 file.

% =============================================================================
% END APPENDED VERSION V95 MATERIAL
% =============================================================================

\clearpage
% =============================================================================
% BEGIN APPENDED VERSION V96 MATERIAL
% =============================================================================
\section{V96: the mixed metric-normal tensor seed, its exact pointwise
image, and the derivative-of-\(K\) closure obstruction}
\label{sec:V96-mixed-normal-seed}

V95 required the next counterterm calculation to keep one common
coefficient vector while enlarging the metric first-normal image.  The
obvious enlargement is linear in the extrinsic curvature and quadratic in
tangential clock jets or linear in intrinsic curvature.  A scalar-only
enumeration is insufficient: an action
\(\int K_{ab}T^{ab}\) has first-normal metric coefficient \(T^{ab}\), so
the correct object to classify is the complete symmetric tensor seed.

This note gives that pointwise classification.  There are sixteen
parity-even algebraic tensor generators at tangential engineering dimension
two.  Their scalar, mixed, and spatial-traceless parts give an exact rank
test for the V93 missing normal directions.  The classification is not yet
the integration-by-parts quotient.  Raychaudhuri shows explicitly that the
subspace with no tangential derivative of \(K\) is not closed under that
quotient: eliminating \(R_{ab}u^au^b\) creates \(D_aK_{bc}\).  The
derivative-of-\(K\) block is therefore forced upstream before a final basis
count can be claimed.

\subsection{Curvature and clock tensor irreducibles}
\label{subsec:V96-tensor-irreducibles}

Decompose the intrinsic Ricci tensor relative to the V93 clock splitting:
\begin{align}
 \rho&:=R_{ab}[\gamma]u^au^b,
 &\varrho&:=q^{ab}R_{ab}[\gamma],
 \notag\\
 \jmath_a&:=q_a{}^cR_{cd}[\gamma]u^d,
 &\mathcal S_{ab}&:=q_a{}^cq_b{}^dR_{cd}[\gamma]
                    -\frac12\varrho q_{ab}.
 \label{eq:V96-Ricci-irreducibles}
\end{align}
Thus \(\jmath\) is spatial and \(\mathcal S\) is spatial, symmetric, and
traceless.  For a spatial symmetric tensor, use angular brackets for the
trace-free projection,
\begin{equation}
 v^{\langle a}w^{b\rangle}
 :=q^a{}_cq^b{}_dv^{(c}w^{d)}
   -\frac12q^{ab}q_{cd}v^cw^d.
 \label{eq:V96-spatial-TF}
\end{equation}

Let \(\mathscr T_2^+\) be the real vector space of parity-even symmetric
contravariant two-tensor polynomial germs of engineering dimension two
which are linear in intrinsic curvature or quadratic in \(D u\), contain no
derivative of curvature, and are reduced only by pointwise algebraic
identities.  Tangential integration by parts is deliberately postponed.

\begin{theorem}[Sixteen algebraic tensor seeds]
\label{thm:V96-sixteen-tensor-seeds}
A basis of \(\mathscr T_2^+\) is the following list:
\begin{align}
 \mathcal B_R={}&\{
 \rho u^au^b,\ \varrho u^au^b,
 \rho q^{ab},\ \varrho q^{ab},
 u^{(a}\jmath^{b)},\ \mathcal S^{ab}\},
 \label{eq:V96-curvature-tensor-basis}\\
 \mathcal B_{Du,\mathrm{sc}}={}&\{
 a^2u^au^b,\ \theta^2u^au^b,\ \chi^2u^au^b,
 a^2q^{ab},\ \theta^2q^{ab},\ \chi^2q^{ab}\},
 \label{eq:V96-aether-scalar-tensor-basis}\\
 \mathcal B_{Du,\mathrm{mix}}={}&\{
 u^{(a}\theta a^{b)},\
 u^{(a}\chi^{b)}{}_ca^c\},
 \label{eq:V96-aether-mixed-tensor-basis}\\
 \mathcal B_{Du,\mathrm{TF}}={}&\{
 a^{\langle a}a^{b\rangle},\ \theta\chi^{ab}\}.
 \label{eq:V96-aether-TF-tensor-basis}
\end{align}
Consequently \(\dim\mathscr T_2^+=16\).
\end{theorem}

\begin{proof}
Under the spatial orthogonal group, a symmetric boundary tensor splits into
two scalar slots \(u^au^b,q^{ab}\), one spatial-vector slot
\(u^{(a}v^{b)}\), and one spatial-traceless slot.  The Ricci decomposition
\eqref{eq:V96-Ricci-irreducibles} supplies respectively two curvature
scalars, one vector, and one traceless tensor, giving the six elements of
\(\mathcal B_R\).

For two factors from \eqref{eq:V94-aether-decomposition}, the scalar
invariants are \(a^2,\theta^2,\chi^2\).  The spatial vectors are
\(\theta a\) and \(\chi a\).  The spatial traceless tensors are
\(a^{\langle a}a^{b\rangle}\) and \(\theta\chi^{ab}\).  In two spatial
dimensions,
\begin{equation*}
 \chi^a{}_c\chi^{cb}=\frac12\chi^2q^{ab},
\end{equation*}
so it adds no generator.  Orthogonality, symmetry, and parity exclude all
remaining contractions at this degree.

Independence follows by the irreducible splitting.  Within each slot, the
curvature jets \((\rho,\varrho,\jmath,\mathcal S)\), lapse-gradient jet
\(a\), and trace/traceless foliation jets \((\theta,\chi)\) may be prescribed
independently at a point subject to the displayed algebraic symmetries.
Polynomial evaluation then isolates every listed coefficient.
\end{proof}

\begin{firewall}[This is an algebraic associated-graded basis]
\label{fire:V96-associated-graded-only}
Theorem~\ref{thm:V96-sixteen-tensor-seeds} does not assert that all sixteen
classes survive the quotient by tangential divergences.  That quotient
mixes this bank with derivatives of \(K\), as proved below.  The theorem is
the exact algebraic input and normal-symbol image before that mixing.
\end{firewall}

\subsection{Mixed actions and their first-normal image}
\label{subsec:V96-mixed-actions}

For \(T\in\mathscr T_2^+\), define
\begin{equation}
 C_T:=\int_\Gamma\sqrt{|\gamma|}\,K_{ab}T^{ab}.
 \label{eq:V96-mixed-action}
\end{equation}
It is parity even, has engineering dimension three, and uses one normal
metric jet.  Let \(\mathbb V_{K\mathrm{tan}}\) be the complete stacked
variation map of the sixteen actions, including their metric, clock, and
all induced matter columns.

\begin{theorem}[Exact metric first-normal symbol]
\label{thm:V96-normal-symbol}
At \(K=0\), the metric stress variation of \eqref{eq:V96-mixed-action} has
first-normal part
\begin{equation}
 \boxed{
 \operatorname{pr}_1\mathcal T_T^{ab}
 =-D_n(\delta_\Gamma T^{ab}).}
 \label{eq:V96-normal-symbol}
\end{equation}
Therefore the coefficient image of the full sixteen-generator bank is
exactly the evaluated span of \(\mathscr T_2^+\).  Algebraic variation of
\(T\) contributes only to the zero-normal coefficient at \(K=0\).
\end{theorem}

\begin{proof}
At fixed tangential jets,
\(\delta(K_{ab}T^{ab})=T^{ab}\delta K_{ab}\) when \(K=0\); the term
\(K_{ab}\delta T^{ab}\) vanishes.  Use
\(\delta K_{ab}=\tfrac12D_n\delta\gamma_{ab}+\mathcal L_{ab}\) from
\eqref{eq:V92-K-chain-rule}.  Distributional integration by parts and the
factor \(1/2\) in the stress convention give
\eqref{eq:V96-normal-symbol}.  Linearity proves the image statement.
\end{proof}

Decompose a proposed symmetric first-normal coefficient as
\begin{equation}
 Z^{ab}=A u^au^b+Bq^{ab}+2u^{(a}v^{b)}+z^{ab},
 \qquad v^au_a=0,
 \qquad z^a{}_a=z^{ab}u_b=0.
 \label{eq:V96-target-decomposition}
\end{equation}

\begin{corollary}[Pointwise rank test]
\label{cor:V96-pointwise-rank-test}
At a fixed anchor, the V96 normal-symbol equations for
\eqref{eq:V96-target-decomposition} are solvable precisely when
\begin{align}
 (A,B)&\in
 \operatorname{span}_{\mathbb R}\{(
ho,0),(\varrho,0),(a^2,0),
 (\theta^2,0),(\chi^2,0),
 (0,\rho),(0,\varrho),(0,a^2),(0,\theta^2),(0,\chi^2)\},
 \label{eq:V96-scalar-rank-test}\\
 v&\in\operatorname{span}_{\mathbb R}
 \{\jmath,\theta a,\chi a\},
 \label{eq:V96-mixed-rank-test}\\
 z&\in\operatorname{span}_{\mathbb R}
 \{\mathcal S,a^{\langle\cdot}a^{\cdot\rangle},\theta\chi\}.
 \label{eq:V96-TF-rank-test}
\end{align}
If all dimension-two seed jets vanish, the entire V96 normal-symbol image is
zero.  Hence the V93 mixed and traceless annihilators persist at the exactly
flat constant-clock anchor, while they can be lifted at a nonflat anchor
only when the spans \eqref{eq:V96-mixed-rank-test}--
\eqref{eq:V96-TF-rank-test} have the required ranks.
\end{corollary}

\begin{proof}
Project the basis
\eqref{eq:V96-curvature-tensor-basis}--
\eqref{eq:V96-aether-TF-tensor-basis} onto the four irreducible slots in
\eqref{eq:V96-target-decomposition}.  The displayed spans are exactly the
images in those slots.  If every seed jet is zero, every basis tensor
vanishes.
\end{proof}

\begin{remark}[Scalar slot simplification]
\label{rem:V96-scalar-slot}
At one point, if at least one of
\(\rho,\varrho,a^2,\theta^2,\chi^2\) is nonzero, the independent
\(u^au^b\) and \(q^{ab}\) coefficients make the scalar slot all of
\(\mathbb R^2\).  The vector and traceless slots retain the genuine rank
conditions \eqref{eq:V96-mixed-rank-test}--
\eqref{eq:V96-TF-rank-test}.
\end{remark}

\subsection{Why the no-derivative-of-\(K\) bank is not quotient closed}
\label{subsec:V96-quotient-closure}

Let \(L\) stand for either scalar extrinsic component
\(\kappa\) or \(\vartheta\).  Multiplying the V94 Raychaudhuri identity by
\(L\) and integrating by parts changes the earlier scalar quotient.

\begin{proposition}[Weighted Raychaudhuri identity]
\label{prop:V96-weighted-Raychaudhuri}
Modulo a tangential divergence,
\begin{equation}
 \boxed{
 L\rho
 \equiv
 \theta D_uL-a^aD_aL
 +\frac12L\theta^2-L\chi^2.}
 \label{eq:V96-weighted-Raychaudhuri}
\end{equation}
The first two terms contain tangential derivatives of \(K\).  Consequently
the span of actions \(K_{ab}T^{ab}\) with algebraic \(T\in\mathscr T_2^+\)
is not closed under the V92 tangential-divergence quotient.
\end{proposition}

\begin{proof}
Solving \eqref{eq:V94-Raychaudhuri} for \(\rho\) gives
\begin{equation*}
 \rho=-D_u\theta-\frac12\theta^2-\chi^2+D_aa^a.
\end{equation*}
For the last term,
\(L D_aa^a\equiv-a^aD_aL\).  For the first,
\begin{equation*}
 -L D_u\theta
 =-D_a(L\theta u^a)+\theta D_uL+L\theta^2.
\end{equation*}
Combining the terms proves \eqref{eq:V96-weighted-Raychaudhuri}.  Since
\(D_aL\) contains \(D_aK_{bc}\), the quotient leaves the algebraic bank.
\end{proof}

\begin{corollary}[Derivative-of-\(K\) block is compulsory]
\label{cor:V96-DK-compulsory}
A final dimension-three quotient basis containing the curvature generators
\(\kappa\rho\) or \(\vartheta\rho\) must also contain the derivative terms
\begin{equation}
 \theta D_u\kappa,\quad a^aD_a\kappa,
 \quad\theta D_u\vartheta,\quad a^aD_a\vartheta,
 \label{eq:V96-forced-DK-terms}
\end{equation}
or equivalent covariant combinations.  Eliminating either side of
\eqref{eq:V96-weighted-Raychaudhuri} before adjoining the other silently
changes the quotient space.
\end{corollary}

\begin{proof}
Apply Proposition~\ref{prop:V96-weighted-Raychaudhuri} to
\(L=\kappa\) and \(L=\vartheta\).
\end{proof}

\begin{firewall}[Normal-symbol solvability is not stacked solvability]
\label{fire:V96-symbol-not-stacked}
The pointwise tests \eqref{eq:V96-scalar-rank-test}--
\eqref{eq:V96-TF-rank-test} concern only the metric first-normal projection.
Each action \eqref{eq:V96-mixed-action} also has zero-order metric and clock
columns, and may mix with other sectors after quotienting.  Cancellation by
one action still requires the stacked V95 condition
\eqref{eq:V95-common-range}.
\end{firewall}

\subsection{V96 result ledger and next acquisition}
\label{subsec:V96-ledger}

\paragraph{New exact conclusions.}
\begin{enumerate}[label=\textup{(\arabic*)},leftmargin=2.8em]
\item The parity-even algebraic symmetric tensor seed has the sixteen-element
      basis \eqref{eq:V96-curvature-tensor-basis}--
      \eqref{eq:V96-aether-TF-tensor-basis}.
\item The mixed action \(\int K_{ab}T^{ab}\) has exact first-normal metric
      coefficient \(-T^{ab}\) at \(K=0\).
\item The scalar, mixed, and traceless pointwise images are the spans
      \eqref{eq:V96-scalar-rank-test}--
      \eqref{eq:V96-TF-rank-test}.
\item The V93 annihilators persist at the exactly flat constant-clock anchor
      but can lift at curved anchors satisfying the displayed rank tests.
\item Weighted Raychaudhuri gives
      \eqref{eq:V96-weighted-Raychaudhuri} and proves that the algebraic
      mixed bank is not closed under tangential integration by parts.
\item The four derivative-of-\(K\) structures
      \eqref{eq:V96-forced-DK-terms} are compulsory in the next quotient.
\end{enumerate}

\paragraph{Next forced step.}
V97 should construct the full parity-even derivative-of-\(K\) scalar bank at
dimension three, combine it with the sixteen V96 mixed actions, and compute
the quotient generated by weighted Raychaudhuri, contracted Bianchi, and
tangential integration by parts.  Only after that quotient is fixed should
the stacked metric/clock variation matrix and its kernel be computed.

\begin{assessment}[V96 continuum status]
\label{ass:V96-continuum-status}
The full Standard--Model--Einstein continuum remains unproved.  V96 computes
the mixed metric normal symbol and exposes the derivative-of-\(K\) closure
obstruction.  The final dimension-three quotient, gauge/fermion/ghost
blocks, regulator target, uniform quantum form bounds, curved radiative
equivariance, nonlinear existence, and cofinal interacting limit remain
open.
\end{assessment}

\section*{V96 exact checks and append-only record}
\addcontentsline{toc}{section}{V96 exact checks and append-only record}

\begin{equation}
\boxed{
\begin{gathered}
 \dim\mathscr T_2^+=16,
 \qquad
 \operatorname{pr}_1\mathcal T_T^{ab}
 =-D_n(\delta_\Gamma T^{ab}),\\
 v\in\operatorname{span}\{\jmath,\theta a,\chi a\},
 \qquad
 z\in\operatorname{span}\{\mathcal S,
 a^{\langle\cdot}a^{\cdot\rangle},\theta\chi\},\\
 L\rho\equiv\theta D_uL-a^aD_aL
 +\frac12L\theta^2-L\chi^2.
\end{gathered}}
\label{eq:V96-check-card}
\end{equation}

The V95 predecessor used for this append has SHA--256
\begin{center}
\texttt{246638b210380cad8d275ee91cc427dffe2eb7a3a548a65ca115f7206d17c9bf}.
\end{center}
The V96 candidate retains its bytes up to the sole live final document
terminator, appends this material, and restores exactly one active
\verb|\end{document}|.  Validation includes balanced environments and
braces, unique labels, resolved new internal references, valid inline and
display math delimiters, exact predecessor-prefix equality, and a final
inventory which removes only the superseded SM V95 file.

% =============================================================================
% END APPENDED VERSION V96 MATERIAL
% =============================================================================

\clearpage
% =============================================================================
% BEGIN APPENDED VERSION V97 MATERIAL
% =============================================================================
\section{V97: the divergence-module obstruction, the weighted
Raychaudhuri quotient orbit, and its scalar-only normal image}
\label{sec:V97-divergence-module}

V96 proved that multiplying the V94 Raychaudhuri quotient by an extrinsic
scalar creates tangential derivatives of \(K\).  This note identifies the
algebraic reason and computes the smallest quotient block in which the
problem closes.  Tangential divergences are a linear subspace of local
densities, but they are not an ideal under multiplication by a nonconstant
field.  Therefore an unweighted counterterm identity cannot be multiplied
term by term without retaining its Leibniz correction.

For each of the two extrinsic scalars \(L=\kappa,\vartheta\), seven natural
dimension-three germs form one weighted-Raychaudhuri orbit.  One pointwise
Raychaudhuri relation and two independent integration-by-parts relations
reduce this formal orbit to four classes.  All four have scalar metric
first-normal symbols.  Thus this entire orbit can modify the two scalar
normal coefficients, but it cannot supply the spatial-vector or
spatial-traceless directions isolated in V96.

\subsection{Divergences are not a multiplicative ideal}
\label{subsec:V97-not-an-ideal}

Write \(f\equiv g\) when their density difference is a tangential
divergence and hence has zero action against compactly supported variations
away from corners.

\begin{theorem}[Leibniz obstruction to multiplying quotient identities]
\label{thm:V97-Leibniz-obstruction}
If
\begin{equation}
 f-g=D_aV^a,
 \label{eq:V97-unweighted-divergence}
\end{equation}
then for every scalar germ \(L\),
\begin{equation}
 \boxed{
 Lf-Lg=D_a(LV^a)-(D_aL)V^a,
 \qquad
 Lf\equiv Lg-(D_aL)V^a.}
 \label{eq:V97-Leibniz-correction}
\end{equation}
Consequently multiplication descends to the divergence quotient only for
multipliers satisfying \((D_aL)V^a\equiv0\) for every representative
current involved.  It does not descend for a general dynamical
\(L=\kappa\) or \(\vartheta\).
\end{theorem}

\begin{proof}
Apply the covariant Leibniz rule to \(D_a(LV^a)\).  The first term on the
right of \eqref{eq:V97-Leibniz-correction} is a divergence; the second is
the unavoidable correction.  A multiplier defines a well-defined map on
equivalence classes exactly when it maps every zero class to a zero class,
which is the stated condition.
\end{proof}

\begin{corollary}[Why the V94 quotient cannot be multiplied]
\label{cor:V97-no-naive-Raychaudhuri-product}
The V94 relation
\(\rho\equiv\tfrac12\theta^2-\chi^2\) cannot be multiplied by
\(L=\kappa\) or \(\vartheta\) without the correction
\begin{equation}
 \theta D_uL-a^aD_aL.
 \label{eq:V97-Raychaudhuri-Leibniz-correction}
\end{equation}
The corrected identity is exactly
\eqref{eq:V96-weighted-Raychaudhuri}.
\end{corollary}

\begin{proof}
In the proof of Proposition~\ref{prop:V94-curvature-quotient}, the discarded
current is \(a^a-\theta u^a\).  Formula
\eqref{eq:V97-Leibniz-correction} contributes
\(-D_aL(a^a-\theta u^a)=\theta D_uL-a^aD_aL\), giving the result.
\end{proof}

\subsection{Derivative transfer and the seven-generator orbit}
\label{subsec:V97-Raychaudhuri-orbit}

For either \(L\in\{\kappa,\vartheta\}\), introduce
\begin{align}
 A_L&:=L\rho,
 &B_L&:=L\theta^2,
 &C_L&:=L\chi^2,
 \notag\\
 U_L&:=L D_u\theta,
 &V_L&:=L D_aa^a,
 &P_L&:=\theta D_uL,
 &Q_L&:=a^aD_aL.
 \label{eq:V97-seven-generators}
\end{align}
All have engineering dimension three and are parity even.

\begin{lemma}[Universal derivative transfer]
\label{lem:V97-derivative-transfer}
For any tensor coefficient \(A^{abc}\),
\begin{equation}
 \int_\Gamma\sqrt{|\gamma|}\,A^{abc}D_aK_{bc}
 \equiv
 -\int_\Gamma\sqrt{|\gamma|}\,K_{bc}D_aA^{abc}.
 \label{eq:V97-DK-transfer}
\end{equation}
Thus a representative with a tangential derivative on \(K\) can be changed
to one without it, provided the coefficient bank is enlarged by the
derivatives of \(A\).  This is a choice of representative, not permission to
discard the resulting terms.
\end{lemma}

\begin{proof}
Integrate \(D_a(A^{abc}K_{bc})\) tangentially.  Compact support away from
corners, or inclusion of the declared corner action, removes the boundary of
the boundary.
\end{proof}

\begin{proposition}[Three exact orbit relations]
\label{prop:V97-three-orbit-relations}
The seven germs \eqref{eq:V97-seven-generators} satisfy
\begin{align}
 A_L+U_L+\frac12B_L+C_L-V_L&=0,
 \label{eq:V97-pointwise-Raychaudhuri-relation}\\
 P_L+U_L+B_L&\equiv0,
 \label{eq:V97-u-IBP-relation}\\
 Q_L+V_L&\equiv0.
 \label{eq:V97-a-IBP-relation}
\end{align}
The first relation is pointwise; the last two are tangential
integration-by-parts identities.
\end{proposition}

\begin{proof}
Multiplying \eqref{eq:V94-Raychaudhuri} by \(L\) and moving every term to
the left gives \eqref{eq:V97-pointwise-Raychaudhuri-relation}.  Next,
\begin{align*}
 D_a(L\theta u^a)&=P_L+U_L+B_L,\\
 D_a(La^a)&=Q_L+V_L.
\end{align*}
Both right sides are divergences, proving the quotient relations.
\end{proof}

\begin{theorem}[Formal weighted-Raychaudhuri quotient]
\label{thm:V97-four-class-orbit}
In the seven-dimensional formal span generated by
\eqref{eq:V97-seven-generators}, quotienting by the three relations
\eqref{eq:V97-pointwise-Raychaudhuri-relation}--
\eqref{eq:V97-a-IBP-relation} leaves dimension four.  Two convenient bases
of the quotient are
\begin{equation}
 \{[A_L],[B_L],[C_L],[U_L]\}
 \quad\hbox{and}\quad
 \{[A_L],[B_L],[C_L],[V_L]\}.
 \label{eq:V97-two-orbit-bases}
\end{equation}
The transition between them is
\begin{equation}
 [V_L]=[A_L]+[U_L]+\frac12[B_L]+[C_L].
 \label{eq:V97-orbit-transition}
\end{equation}
This is the complete quotient only inside the displayed formal orbit;
additional differential identities from the remaining tensor sectors have
not yet been imposed.
\end{theorem}

\begin{proof}
In the ordered coordinates
\((A_L,B_L,C_L,U_L,V_L,P_L,Q_L)\), the three relation rows have pivots in
\(A_L,P_L,Q_L\), respectively, and are linearly independent.  Hence their
quotient has dimension \(7-3=4\).  Eliminating \(V_L,P_L,Q_L\) gives the
first basis; eliminating \(U_L,P_L,Q_L\) gives the second.
Equation \eqref{eq:V97-orbit-transition} is the first relation solved for
\(V_L\).
\end{proof}

\begin{firewall}[Formal rank is not yet global local-cohomology rank]
\label{fire:V97-formal-rank-scope}
Theorem~\ref{thm:V97-four-class-orbit} proves the rank of the explicitly
printed Raychaudhuri/Leibniz subsystem.  It does not assert that no further
contracted-Bianchi, Codazzi, or tensor integration-by-parts relation meets
this span after the full dimension-three bank is added.
\end{firewall}

\subsection{Normal-symbol image of the orbit}
\label{subsec:V97-orbit-normal-image}

Every germ in \eqref{eq:V97-seven-generators} is one extrinsic scalar
\(L\) times a tangential scalar, except \(P_L,Q_L\), whose derivative of
\(L\) can be transferred by Lemma~\ref{lem:V97-derivative-transfer}.

\begin{theorem}[The weighted orbit is scalar in the metric normal symbol]
\label{thm:V97-scalar-normal-symbol}
At \(K=0\), the first-normal metric coefficient of every action in the
weighted-Raychaudhuri orbit lies in
\begin{equation}
 \operatorname{span}_{\mathbb R}\{u^au^b,q^{ab}\}.
 \label{eq:V97-scalar-normal-image}
\end{equation}
Consequently adjoining the complete orbit can alter the scalar test
\eqref{eq:V96-scalar-rank-test}, but it cannot alter the mixed and
spatial-traceless rank tests
\eqref{eq:V96-mixed-rank-test}--\eqref{eq:V96-TF-rank-test}.
\end{theorem}

\begin{proof}
For \(L=\kappa\), variation with respect to \(K_{ab}\) gives
\(u^au^b\); for \(L=\vartheta\), it gives \(q^{ab}\).  A tangential scalar
multiplier preserves that tensor type.  For \(P_L,Q_L\), first use
\eqref{eq:V97-u-IBP-relation}--\eqref{eq:V97-a-IBP-relation} to choose a
representative linear in \(L\) without \(D L\).  The normal symbol is
invariant under adding a tangential divergence to the action, so the same
conclusion follows.
\end{proof}

\begin{corollary}[Mixed and traceless generators remain indispensable]
\label{cor:V97-mixed-TF-indispensable}
Any regulator target whose metric first-normal coefficient has a component
outside \(\operatorname{span}\{u u,q\}\) still requires the V96 vector or
traceless tensor generators.  Weighted scalar Raychaudhuri counterterms
cannot cancel it, regardless of representative.
\end{corollary}

\begin{proof}
Project Theorem~\ref{thm:V97-scalar-normal-symbol} onto the orthogonal
spatial-vector and spatial-traceless irreducible slots.
\end{proof}

\begin{assessment}[What the quotient calculation changes]
\label{ass:V97-quotient-decision}
The derivative-of-\(K\) terms forced by V96 are required while performing
the quotient, but Lemma~\ref{lem:V97-derivative-transfer} permits a final
representative without \(D K\) only after the differentiated coefficients
have been retained.  The V94 scalar quotient cannot be multiplied naively.
The correct weighted orbit has four formal classes for each of
\(\kappa,\vartheta\), and it does not address the genuinely mixed or
traceless normal symbols.
\end{assessment}

\subsection{V97 result ledger and next acquisition}
\label{subsec:V97-ledger}

\paragraph{New exact conclusions.}
\begin{enumerate}[label=\textup{(\arabic*)},leftmargin=2.8em]
\item The Leibniz formula \eqref{eq:V97-Leibniz-correction} proves that the
      tangential-divergence quotient is not a module over arbitrary local
      scalar germs.
\item Derivatives of \(K\) transfer according to
      \eqref{eq:V97-DK-transfer}, but their differentiated coefficients must
      remain in the bank.
\item Each weighted-Raychaudhuri orbit has the three exact relations
      \eqref{eq:V97-pointwise-Raychaudhuri-relation}--
      \eqref{eq:V97-a-IBP-relation}.
\item Their formal quotient has dimension four and the two bases
      \eqref{eq:V97-two-orbit-bases}.
\item The entire orbit has scalar-only metric first-normal symbol and cannot
      repair the V96 mixed or traceless rank failures.
\end{enumerate}

\paragraph{Next forced step.}
V98 should construct the corresponding vector and spatial-traceless
differential orbits.  It must decompose the contracted Bianchi identity for
\((\rho,\varrho,\jmath,\mathcal S)\), combine it with derivatives of
\((a,\theta,\chi)\), and compute which relations survive after multiplication
by \(b_a\) and \(\sigma_{ab}\).  The resulting quotient must then be inserted
into the stacked V95 variation map.

\begin{assessment}[V97 continuum status]
\label{ass:V97-continuum-status}
The full Standard--Model--Einstein continuum remains unproved.  V97 closes
the scalar weighted-Raychaudhuri quotient subsystem but not the vector and
traceless differential orbits, the remaining field counterterms, regulator
target, uniform quantum form bounds, curved radiative equivariance,
nonlinear existence, or cofinal interacting limit.
\end{assessment}

\section*{V97 exact checks and append-only record}
\addcontentsline{toc}{section}{V97 exact checks and append-only record}

\begin{equation}
\boxed{
\begin{gathered}
 Lf\equiv Lg-(D_aL)V^a
 \quad\text{when }f-g=D_aV^a,\\
 A_L+U_L+\frac12B_L+C_L-V_L=0,
 \quad P_L+U_L+B_L\equiv0,
 \quad Q_L+V_L\equiv0,\\
 \dim\mathscr R_L^{\rm formal}=4,
 \qquad
 \operatorname{pr}_1\mathscr R_L
 \subset\operatorname{span}\{u^au^b,q^{ab}\}.
\end{gathered}}
\label{eq:V97-check-card}
\end{equation}

The V96 predecessor used for this append has SHA--256
\begin{center}
\texttt{720d6d29f4129acde4d3f357a1036c0d641042350dd602cc66a7f403e7d31f00}.
\end{center}
The V97 candidate retains its bytes up to the sole live final document
terminator, appends this material, and restores exactly one active
\verb|\end{document}|.  Validation includes balanced environments and
braces, unique labels, resolved new internal references, valid inline and
display math delimiters, exact predecessor-prefix equality, and a final
inventory which removes only the superseded SM V96 file.

% =============================================================================
% END APPENDED VERSION V97 MATERIAL
% =============================================================================

\clearpage
% =============================================================================
% BEGIN APPENDED VERSION V98 MATERIAL
% =============================================================================
\section{V98: the missing lower-dimensional \(KDu\) block, the corrected
Codazzi interface, and the intrinsic Bianchi relation}
\label{sec:V98-KDu-Codazzi}

V97 organized the weighted-Raychaudhuri orbit at engineering dimension
three.  That calculation exposes a lower-dimensional block which must be
inserted first.  Since both \(K\) and \(D u\) have dimension one, their
bilinear scalar contractions already have boundary dimension two.  They
were outside the algebraic V93 block and the purely tangential V94 block.

This note classifies the parity-even \(KDu\) scalars, computes their exact
metric first-normal symbols, and applies the corrected V90 Codazzi sign.
One linear combination is the normal--clock contraction of the bulk Ricci
tensor modulo a tangential divergence.  Replacing the boundary expression
by that bulk curvature jet would violate the declared first-normal-jet
policy, so the boundary representative must be retained.  Separately, the
intrinsic contracted Bianchi identity supplies one exact dimension-three
relation among the V96 curvature and clock irreducibles.

\subsection{Four bilinear invariants and their normal symbols}
\label{subsec:V98-KDu-basis}

Use the decompositions \eqref{eq:V93-K-decomposition-data} and
\eqref{eq:V94-aether-decomposition}.  Let \(\mathfrak C_{KDu}^{2,+}\) be
the pointwise parity-even scalar germs bilinear in \(K\) and \(D u\).

\begin{theorem}[The \(KDu\) basis]
\label{thm:V98-KDu-basis}
A basis of \(\mathfrak C_{KDu}^{2,+}\) is
\begin{equation}
 \boxed{
 \mathcal B_{KDu}
 =\{\kappa\theta,\ \vartheta\theta,\ a^ab_a,
       \ \sigma^{ab}\chi_{ab}\}.}
 \label{eq:V98-KDu-basis}
\end{equation}
The four classes occupy respectively the two scalar, spatial-vector, and
spatial-traceless pairings of the clock stabilizer.
\end{theorem}

\begin{proof}
The extrinsic tensor decomposes into
\((\kappa,\vartheta,b,\sigma)\), and \(D u\) decomposes into
\((\theta,a,\chi)\).  A parity-even scalar pairs equal irreducible spatial
types.  The two extrinsic scalars may each pair with \(\theta\), while the
vector and traceless tensors have the unique pairings \(a\cdot b\) and
\(\sigma:\chi\).  Orthogonality of inequivalent irreducibles proves spanning
and independence.
\end{proof}

For constants \((c_\kappa,c_\vartheta,c_b,c_\sigma)\), put
\begin{equation}
 f_{KDu}:=c_\kappa\kappa\theta+c_\vartheta\vartheta\theta
          +c_ba^ab_a+c_\sigma\sigma^{ab}\chi_{ab}.
 \label{eq:V98-KDu-action-density}
\end{equation}

\begin{proposition}[Exact first-normal coefficient]
\label{prop:V98-KDu-normal-symbol}
At \(K=0\), the extrinsic momentum of
\(\int_\Gamma\sqrt{|\gamma|}f_{KDu}\) is
\begin{equation}
 \boxed{
 P^{ab}_{KDu}
 =c_\kappa\theta u^au^b+c_\vartheta\theta q^{ab}
  +c_bu^{(a}a^{b)}+c_\sigma\chi^{ab}.}
 \label{eq:V98-KDu-normal-symbol}
\end{equation}
Accordingly its metric stress contains
\(-D_n(\delta_\Gamma P^{ab}_{KDu})\).  If
\(\theta\ne0\), both scalar normal slots are accessible.  The mixed slot is
the line spanned by \(a\), and the spatial-traceless slot is the line spanned
by \(\chi\).
\end{proposition}

\begin{proof}
Differentiate \eqref{eq:V98-KDu-action-density} with respect to \(K_{ab}\)
at fixed \((\gamma,u,D u)\).  The four derivatives are respectively
\(\theta u^au^b\), \(\theta q^{ab}\),
\(u^{(a}a^{b)}\), and \(\chi^{ab}\).  The stress statement follows from
Theorem~\ref{thm:V96-normal-symbol}.
\end{proof}

\begin{corollary}[Refined flat-anchor statement]
\label{cor:V98-refined-flat-anchor}
At an exactly constant-clock anchor, \(D_au_b=0\), the entire V98 normal
image vanishes, so the V93 mixed and traceless annihilators persist.  At a
nonconstant clock, the first missing directions can already enter at
engineering dimension two; the dimension-three V96 tensors are not their
first possible source.
\end{corollary}

\begin{proof}
Constant \(u\) gives \(a=\theta=\chi=0\) in
\eqref{eq:V98-KDu-normal-symbol}.  The second statement follows whenever
\(a\) or \(\chi\) is nonzero.
\end{proof}

\subsection{Corrected Codazzi and the bulk-jet interface}
\label{subsec:V98-Codazzi-interface}

With the V89 definitions and the V90 orientation correction,
\begin{equation}
 D_aK^a{}_b-D_bK
 =-n^\mu\gamma_b{}^\nu R_{\mu\nu}[g].
 \label{eq:V98-corrected-Codazzi}
\end{equation}
Here \(K=\vartheta-\kappa\).

\begin{theorem}[Clock-contracted Codazzi identity]
\label{thm:V98-clock-Codazzi}
Modulo a tangential divergence,
\begin{equation}
 \boxed{
 a^ab_a+\left(\frac12\vartheta-\kappa\right)\theta
 -\sigma^{ab}\chi_{ab}
 \equiv
 -n^\mu u^\nu R_{\mu\nu}[g].}
 \label{eq:V98-clock-Codazzi-quotient}
\end{equation}
Thus the four V98 generators have one geometric interface relation only
after the normal--tangential bulk Ricci jet is admitted.
\end{theorem}

\begin{proof}
Contract \eqref{eq:V98-corrected-Codazzi} with \(u^b\).  Tangential
integration by parts gives
\begin{equation*}
 u^b(D_aK^a{}_b-D_bK)
 \equiv-K^a{}_bD_au^b+K\theta.
\end{equation*}
Direct substitution of the two irreducible decompositions yields
\begin{equation*}
 K^{ab}D_au_b=-a^ab_a+\frac12\vartheta\theta
               +\sigma^{ab}\chi_{ab}.
\end{equation*}
Using \(K=\vartheta-\kappa\) proves
\eqref{eq:V98-clock-Codazzi-quotient} with the corrected V90 sign.
\end{proof}

\begin{firewall}[Codazzi does not remove a first-jet boundary generator]
\label{fire:V98-Codazzi-policy}
The right side of \eqref{eq:V98-clock-Codazzi-quotient} is a restriction of
the four-dimensional Ricci tensor and contains second metric jets.  Policy
\textup{(P2)} of V92 permits at most one normal derivative in the boundary
action.  Therefore one may not quotient the left side to zero or replace it
by the right side while claiming to remain inside that policy.  The V98
boundary combination must be retained unless the counterterm bank is
explicitly enlarged to bulk-curvature boundary restrictions or the bulk
field equation supplies a justified substitution.
\end{firewall}

\begin{corollary}[Vacuum-anchor conditional relation]
\label{cor:V98-vacuum-Codazzi}
On a branch where the mixed bulk Ricci component
\(n^\mu u^\nu R_{\mu\nu}\) is proved to vanish at the anchor,
\begin{equation}
 a\cdot b+\left(\frac12\vartheta-\kappa\right)\theta
 -\sigma:\chi\equiv0.
 \label{eq:V98-vacuum-KDu-relation}
\end{equation}
This is conditional on that bulk equation; it is not an off-shell
counterterm identity.
\end{corollary}

\begin{proof}
Set the right side of \eqref{eq:V98-clock-Codazzi-quotient} to zero only
under the stated hypothesis.
\end{proof}

\subsection{The intrinsic dimension-three Bianchi relation}
\label{subsec:V98-intrinsic-Bianchi}

Let \(G^{ab}[\gamma]=R^{ab}-\tfrac12R\gamma^{ab}\) be the intrinsic
three-dimensional Einstein tensor.

\begin{theorem}[Clock-contracted intrinsic Bianchi quotient]
\label{thm:V98-clock-Bianchi}
The contracted Bianchi identity implies the parity-even dimension-three
relation
\begin{equation}
 \boxed{
 -a^a\jmath_a+\frac12\theta\rho
 +\chi^{ab}\mathcal S_{ab}\equiv0.}
 \label{eq:V98-clock-Bianchi-quotient}
\end{equation}
This relation belongs to the purely tangential curvature--clock bank and
must be imposed before counting its cubic generators.
\end{theorem}

\begin{proof}
Since \(D_aG^{ab}=0\),
\begin{equation*}
 D_a(u_bG^{ab})=(D_au_b)G^{ab}
\end{equation*}
is a tangential divergence.  Insert
\eqref{eq:V94-aether-decomposition}.  The acceleration term is
\(-a\cdot\jmath\).  Because
\(q_{ab}G^{ab}=\rho\), the expansion term is
\(\tfrac12\theta\rho\).  Tracelessness leaves
\(\chi:\mathcal S\), proving the formula.
\end{proof}

\begin{assessment}[Reordered counterterm acquisition]
\label{ass:V98-reordered-acquisition}
The first-normal range analysis must now proceed in engineering-dimension
order.  The V98 \(KDu\) block precedes the V96 dimension-three tensor seeds.
Its complete stacked metric and clock variations must be combined with the
Codazzi interface without using the bulk equation off shell.  The intrinsic
Bianchi relation independently reduces the pure tangential cubic bank.
\end{assessment}

\subsection{V98 result ledger and next acquisition}
\label{subsec:V98-ledger}

\paragraph{New exact conclusions.}
\begin{enumerate}[label=\textup{(\arabic*)},leftmargin=2.8em]
\item The missing parity-even \(KDu\) block has the four-element basis
      \eqref{eq:V98-KDu-basis}.
\item Its exact first-normal metric coefficient is
      \eqref{eq:V98-KDu-normal-symbol}, supplying scalar, mixed, and
      traceless directions at nonconstant clocks.
\item Corrected Codazzi gives the interface relation
      \eqref{eq:V98-clock-Codazzi-quotient} with the V90 sign.
\item That relation cannot eliminate a boundary generator inside the V92
      first-normal-jet policy because its right side is a bulk second jet.
\item Contracted intrinsic Bianchi gives the independent tangential cubic
      quotient \eqref{eq:V98-clock-Bianchi-quotient}.
\end{enumerate}

\paragraph{Next forced step.}
V99 should compute the complete clock and zero-order metric variation of the
four V98 actions, stack it with their normal symbols, and determine its
kernel both off shell and on the conditional mixed-Ricci-zero branch.  It
should then revisit the V96 rank test in engineering-dimension order before
adding gauge or fermion counterterms.

\begin{assessment}[V98 continuum status]
\label{ass:V98-continuum-status}
The full Standard--Model--Einstein continuum remains unproved.  V98 inserts
the missing lower-dimensional mixed block and fixes two geometric quotient
identities.  Its stacked variation, the remaining counterterm sectors,
regulator target, uniform quantum form bounds, curved radiative
equivariance, nonlinear existence, and cofinal interacting limit remain
open.
\end{assessment}

\section*{V98 exact checks and append-only record}
\addcontentsline{toc}{section}{V98 exact checks and append-only record}

\begin{equation}
\boxed{
\begin{gathered}
 \mathcal B_{KDu}=\{\kappa\theta,\vartheta\theta,a\cdot b,
 \sigma:\chi\},\\
 P^{ab}_{KDu}=c_\kappa\theta u^au^b+c_\vartheta\theta q^{ab}
 +c_bu^{(a}a^{b)}+c_\sigma\chi^{ab},\\
 a\cdot b+\left(\frac12\vartheta-\kappa\right)\theta
 -\sigma:\chi\equiv-n^\mu u^\nu R_{\mu\nu}[g],\\
 -a\cdot\jmath+\frac12\theta\rho+\chi:\mathcal S\equiv0.
\end{gathered}}
\label{eq:V98-check-card}
\end{equation}

The V97 predecessor used for this append has SHA--256
\begin{center}
\texttt{f4349badbc2b412388a247aa3937ba94e80627311e522d92e6ff9eca42aa6d9d}.
\end{center}
The V98 candidate retains its bytes up to the sole live final document
terminator, appends this material, and restores exactly one active
\verb|\end{document}|.  Validation includes balanced environments and
braces, unique labels, resolved new internal references, valid inline and
display math delimiters, exact predecessor-prefix equality, and a final
inventory which removes only the superseded SM V97 file.

% =============================================================================
% END APPENDED VERSION V98 MATERIAL
% =============================================================================

\clearpage
% =============================================================================
% BEGIN APPENDED VERSION V99 MATERIAL
% =============================================================================
\section{V99: the complete clock operator of the \(KDu\) block, generic
injectivity, pointwise rank strata, and the conditional Codazzi kernel}
\label{sec:V99-KDu-stacked-kernel}

V98 inserted the missing four-dimensional coefficient space
\(\mathfrak C_{KDu}^{2,+}\), computed its metric first-normal symbol, and
identified a Codazzi relation only after a bulk mixed-Ricci condition is
used.  This note separates three notions which must not be merged: the
off-shell polynomial-germ variation map, its evaluation at one background,
and the quotient restricted to a bulk-equation branch.

The off-shell map is injective.  Its pointwise metric-normal evaluation has
rank equal to the number of nonzero clock irreducible slots and vanishes at
a constant clock.  At the joint flat anchor \(K=0\), \(D u=0\), the entire
first variation of every bilinear action vanishes.  This is a background
rank drop, not an off-shell counterterm identity.  If the mixed bulk Ricci
component is imposed to vanish and variations are tangent to that branch,
corrected Codazzi produces one conditional null direction and a
three-dimensional quotient.

\subsection{Independent-jet first variation and clock Euler operator}
\label{subsec:V99-clock-operator}

Let
\begin{equation}
 f=c_\kappa\kappa\theta+c_\vartheta\vartheta\theta
   +c_ba^ab_a+c_\sigma\sigma^{ab}\chi_{ab}
 \label{eq:V99-density}
\end{equation}
and first regard \((\gamma,K,u,D u)\) as independent boundary jet
variables.  Define
\begin{equation}
 M_b:=\left(\frac{\partial f}{\partial u^b}\right)_{\gamma,K,D u},
 \qquad
 N^{ab}:=\frac{\partial f}{\partial(D_au_b)}.
 \label{eq:V99-MN-definition}
\end{equation}

\begin{proposition}[Explicit derivative momentum]
\label{prop:V99-derivative-momentum}
The momentum conjugate to \(D_au_b\) is
\begin{equation}
 \boxed{
 N^{ab}=(c_\kappa\kappa+c_\vartheta\vartheta)\gamma^{ab}
       +c_bu^ab^b+c_\sigma\sigma^{ab}.}
 \label{eq:V99-N-momentum}
\end{equation}
After tangential integration by parts, the unconstrained Euler covector of
the unit direction is
\begin{equation}
 \mathcal U_b:=M_b-D_aN^{ab}.
 \label{eq:V99-u-Euler}
\end{equation}
\end{proposition}

\begin{proof}
At fixed \((\gamma,K,u)\),
\begin{align*}
 \delta\theta&=\gamma^{ab}D_a\delta u_b,\\
 \delta a^b&=u^aD_a\delta u^b,\\
 \sigma^{ab}\delta\chi_{ab}&=\sigma^{ab}D_a\delta u_b.
\end{align*}
The projections in the last identity disappear because \(\sigma\) is
spatial and traceless.  Reading off the derivative of the variation gives
\eqref{eq:V99-N-momentum}; one integration by parts gives
\eqref{eq:V99-u-Euler}.
\end{proof}

The physical unit direction is not independent but is generated by the
clock.  Recall
\(\delta_\Theta u^b=-N_\Theta^{-1}q^{bc}D_c\delta\Theta\), where
\(N_\Theta=\sqrt{-D\Theta\cdot D\Theta}\).

\begin{theorem}[Complete clock Euler coefficient]
\label{thm:V99-clock-Euler}
The clock Euler coefficient of the complete action
\(C_f=\int_\Gamma\sqrt{|\gamma|}f\) is
\begin{equation}
 \boxed{
 E_{\Theta,f}
 =D_c\left(N_\Theta^{-1}q^{cb}mathcal U_b\right),
 \qquad
 \mathcal U_b=M_b-D_aN^{ab},}
 \label{eq:V99-clock-Euler}
\end{equation}
and its supported bulk form is \(\delta_\Gamma E_{\Theta,f}\).  It has no
normal derivative of \(\delta_\Gamma\).  Together with
\eqref{eq:V98-KDu-normal-symbol}, this is the clock/metric-normal part of one
stacked V95 column, not two independently adjustable subtractions.
\end{theorem}

\begin{proof}
After the first tangential integration by parts, the direction variation is
\(\int\sqrt{|\gamma|}\,\mathcal U_b\delta u^b\).  Substitute the normalized
clock variation and integrate once more tangentially.  This gives
\eqref{eq:V99-clock-Euler}.  No normal integration by parts occurs, proving
the support statement.
\end{proof}

\begin{remark}[Zero-order metric column]
\label{rem:V99-zero-order-metric}
The same variation has the intrinsic form
\begin{equation}
 \mathcal T_f^{ab}
 =\delta_\Gamma\tau_{f,0}^{ab}
  -D_n(\delta_\Gamma P_{KDu}^{ab}),
 \label{eq:V99-full-metric-form}
\end{equation}
where \(P_{KDu}\) is explicit in
\eqref{eq:V98-KDu-normal-symbol}.  The coefficient \(\tau_{f,0}\) is fixed,
not free: it is the Hilbert variation of the measure, all index projectors,
the Levi--Civita connection in \(D u\), and the lower-order part of
\(\delta K\).  Its expanded coordinate formula depends on the chosen
lapse--shift parameterization, while the distribution
\eqref{eq:V99-full-metric-form} and its Ward identity are intrinsic.  V99
therefore retains it as the uniquely defined Hilbert derivative rather than
printing a gauge-dependent component expansion.
\end{remark}

\subsection{Off-shell injectivity and pointwise rank strata}
\label{subsec:V99-rank-strata}

Let \(V_{KDu}\) be the complete off-shell first-variation map from the four
coefficients in \eqref{eq:V99-density}.

\begin{theorem}[Off-shell polynomial-germ injectivity]
\label{thm:V99-offshell-injectivity}
As a map of local polynomial germs,
\begin{equation}
 \boxed{\ker V_{KDu}=0.}
 \label{eq:V99-offshell-kernel}
\end{equation}
It already suffices to use the metric first-normal polynomial
\eqref{eq:V98-KDu-normal-symbol}.
\end{theorem}

\begin{proof}
Suppose the first-normal coefficient vanishes for every admissible clock
jet.  The four summands in \eqref{eq:V98-KDu-normal-symbol} occupy the two
scalar, spatial-vector, and spatial-traceless irreducible slots.  Prescribe
\(\theta\), \(a\), and \(\chi\) independently at a point.  The scalar slots
give \(c_\kappa=c_\vartheta=0\), the vector slot gives \(c_b=0\), and the
traceless slot gives \(c_\sigma=0\).  Hence the complete variation kernel is
zero.
\end{proof}

At a fixed background \(x\), denote evaluation of the normal-symbol map by
\(P_x:\mathbb R^4\to\operatorname{Sym}^2(T_x\Gamma)\).

\begin{theorem}[Exact pointwise kernel]
\label{thm:V99-pointwise-kernel}
The pointwise kernel is the direct coordinate span
\begin{equation}
 \boxed{
 \ker P_x
 =\operatorname{span}\{e_\kappa,e_\vartheta:\theta(x)=0\}
  \oplus\operatorname{span}\{e_b:a(x)=0\}
  \oplus\operatorname{span}\{e_\sigma:\chi(x)=0\}.}
 \label{eq:V99-pointwise-kernel}
\end{equation}
Here the notation means that the indicated coordinate vectors are included
exactly when their condition holds.  In particular, \(P_x\) is injective
when \(\theta(x)\ne0\), \(a(x)\ne0\), and \(\chi(x)\ne0\).
\end{theorem}

\begin{proof}
The four images in \eqref{eq:V98-KDu-normal-symbol} lie in mutually
inequivalent irreducible slots, except for the two scalar images, which lie
in the independent tensors \(u u\) and \(q\).  Each coordinate is therefore
lost precisely when its multiplying clock jet vanishes.
\end{proof}

\begin{corollary}[Joint flat-anchor first-order blindness]
\label{cor:V99-flat-anchor-blindness}
At \(K=0\) and \(D u=0\), the complete first variation of every V98 action
vanishes, including \(P_{KDu}\), \(E_{\Theta,f}\), and
\(\tau_{f,0}\).  The four actions begin at bilinear order in perturbations
about this anchor.
\end{corollary}

\begin{proof}
Each density is bilinear in \((K,D u)\).  Its differential is
\((D u)\cdot\delta K+K\cdot\delta(D u)\), which vanishes when both factors
are zero.  Equations \eqref{eq:V98-KDu-normal-symbol} and
\eqref{eq:V99-clock-Euler} give the same conclusion componentwise.
\end{proof}

\begin{firewall}[Background blindness is not a counterterm identity]
\label{fire:V99-background-not-kernel}
Corollary~\ref{cor:V99-flat-anchor-blindness} concerns the derivative of the
evaluation map at one background.  It does not contradict the off-shell
injectivity \eqref{eq:V99-offshell-kernel}.  A flat-background regulator
calculation cannot determine these coefficients at first order; curved
backgrounds or second variations are required.
\end{firewall}

\subsection{The conditional Codazzi quotient}
\label{subsec:V99-conditional-Codazzi-kernel}

Define the coefficient vector
\begin{equation}
 k_C:=(-1,\tfrac12,1,-1)
 \label{eq:V99-Codazzi-vector}
\end{equation}
in the ordered basis
\((\kappa\theta,\vartheta\theta,a\cdot b,\sigma:\chi)\).

\begin{theorem}[One conditional on-shell null direction]
\label{thm:V99-conditional-Codazzi-kernel}
On a branch where
\(n^\mu u^\nu R_{\mu\nu}=0\) and all allowed variations are tangent to that
constraint, the action represented by \(k_C\) is a tangential divergence.
Let \(V_{KDu}^{\rm Cod}\) denote the variation map after quotienting
only by this printed Codazzi relation.  Its coefficient space has dimension
three and
\begin{equation}
 \boxed{\ker V_{KDu}^{\rm Cod}=\operatorname{span}\{k_C\}.}
 \label{eq:V99-conditional-kernel}
\end{equation}
Off shell, this vector is not a kernel direction.
\end{theorem}

\begin{proof}
Equation \eqref{eq:V98-vacuum-KDu-relation} is precisely the linear
combination with coefficient vector \(k_C\).  On the stated branch it is a
divergence, and tangent variations preserve its vanishing action.  The four
pointwise generators are independent by
Theorem~\ref{thm:V98-KDu-basis}, so quotienting by this one nonzero relation
leaves dimension three.  Off shell,
\eqref{eq:V98-clock-Codazzi-quotient} identifies the same combination with
the generally nonzero bulk Ricci restriction, so it is not a kernel.
\end{proof}

\begin{assessment}[Stacked decision after V99]
\label{ass:V99-stacked-decision}
The V98 block has no off-shell coefficient redundancy, but its evaluated
normal symbol is stratified by the clock background.  The conditional
Codazzi quotient is a different map with one null direction.  Regulator
matching must declare which map is being used, retain the clock Euler column
\eqref{eq:V99-clock-Euler}, and never use the conditional kernel in an
off-shell Ward identity.
\end{assessment}

\subsection{V99 result ledger and next acquisition}
\label{subsec:V99-ledger}

\paragraph{New exact conclusions.}
\begin{enumerate}[label=\textup{(\arabic*)},leftmargin=2.8em]
\item The derivative momentum and complete clock Euler operator are
      \eqref{eq:V99-N-momentum}--\eqref{eq:V99-clock-Euler}.
\item The off-shell polynomial-germ variation map is injective.
\item Its exact background-dependent normal-symbol kernel is
      \eqref{eq:V99-pointwise-kernel}.
\item At \(K=0\), \(D u=0\), the whole V98 first variation vanishes and the
      actions first appear at bilinear perturbative order.
\item On the mixed-Ricci-zero tangent branch only, corrected Codazzi produces
      the one-dimensional kernel \eqref{eq:V99-conditional-kernel}.
\end{enumerate}

\paragraph{Next forced step.}
V100 is a compulsory companion-audit version.  After inspecting the newest
Yang--Mills and Navier--Stokes notes, it should derive the second variation
of the V98 block at the joint flat anchor.  That Hessian is the first object
which can determine these four coefficients in the flat perturbative theory
and must be stacked with the V93/V96 first-order columns without double
charging conditional Codazzi.

\begin{assessment}[V99 continuum status]
\label{ass:V99-continuum-status}
The full Standard--Model--Einstein continuum remains unproved.  V99 closes
the V98 clock operator and kernel distinctions, but the flat-anchor Hessian,
remaining counterterm sectors, regulator target, uniform quantum form
bounds, curved radiative equivariance, nonlinear existence, and cofinal
interacting limit remain open.
\end{assessment}

\section*{V99 exact checks and append-only record}
\addcontentsline{toc}{section}{V99 exact checks and append-only record}

\begin{equation}
\boxed{
\begin{gathered}
 N^{ab}=(c_\kappa\kappa+c_\vartheta\vartheta)\gamma^{ab}
       +c_bu^ab^b+c_\sigma\sigma^{ab},\\
 E_{\Theta,f}=D_c\left[N_\Theta^{-1}q^{cb}
 (M_b-D_aN^{ab})\right],\\
 \ker V_{KDu}=0,
 \qquad
 \ker V_{KDu}^{\rm Cod}
 =\operatorname{span}\{(-1,\tfrac12,1,-1)\}.
\end{gathered}}
\label{eq:V99-check-card}
\end{equation}

The V98 predecessor used for this append has SHA--256
\begin{center}
\texttt{a4921cdee3b206dd144c4daa1cb75f286d6be99426354abd0fea5afeff472f55}.
\end{center}
The V99 candidate retains its bytes up to the sole live final document
terminator, appends this material, and restores exactly one active
\verb|\end{document}|.  Validation includes balanced environments and
braces, unique labels, resolved new internal references, valid inline and
display math delimiters, exact predecessor-prefix equality, and a final
inventory which removes only the superseded SM V98 file.

% =============================================================================
% END APPENDED VERSION V99 MATERIAL
% =============================================================================

\clearpage
% =============================================================================
% BEGIN APPENDED VERSION V100 MATERIAL
% =============================================================================
\section{V100: companion audit, the flat-anchor \(KDu\) Hessian, and the
exact signed-kernel test before positive majorization}
\label{sec:V100-flat-Hessian}

V99 proved that every \(KDu\) first variation vanishes at the joint anchor
\(K=0\), \(D u=0\).  The first coefficient-sensitive object there is the
second variation.  Because this is the fifth version after V95, the newest
Yang--Mills and Navier--Stokes notes were inspected before computing it.

The audit reinforces a rule which can be proved directly here: retain the
exact signed operator before replacing it by positive separate-leg squares.
The V98 Hessian is an off-diagonal operator between the extrinsic and clock
jet spaces.  Its scalar block maps two extrinsic scalars to one clock
expansion scalar and therefore has a compulsory kernel line for every fixed
nonzero coefficient vector.  A positive majorant which charges that line
cannot factor through the Hessian.  This is an exact finite-dimensional
failure certificate, not an analogy imported from either companion
programme.

\subsection{Required V100 companion audit}
\label{subsec:V100-audit}

The read-only snapshot was
\begin{center}
\begin{tabular}{@{}lll@{}}
programme & file and byte count & SHA--256\\
\hline
Yang--Mills &
\texttt{Renewal\_Yang\_Mills\_Mass\_Gap\_Research\_Notes\_V97.tex},
\(2320055\) &
\texttt{84499369b26baeb36672ef2c0f6f299ae9c238d10bb241520857d71a8357bb8d}\\
Navier--Stokes &
\texttt{Renewal\_NS\_Stability\_Research\_Notes\_V31.tex},
\(887537\) &
\texttt{f1121b62238ad5491bb7cf03635ea9934942edf573ed8faaa9ee79e1d38a6696}
\end{tabular}
\end{center}
The NS file is unchanged from V95.  The YM file advanced from V94 to V97.

YM V97 exhibits an exact channel on which its complete signed cumulant
vanishes while a separated three-square covariance majorant is positive.
It proves the kernel-inclusion criterion for factorization and deletes
disconnected channels before positivity.  NS V31 retains the complete
signed transfer on the total branch because its polarization hinge has a
balanced-helicity nullspace.  Both observations warn against replacing a
connected object by separately positive components before checking its
kernel.

\begin{firewall}[The audit supplies a test order, not an SM estimate]
\label{fire:V100-audit-scope}
No companion estimate is used below.  The SM Hessian, kernel, and
factorization failure are derived from the V98 action itself.  The audit
only determines the order: compute the signed Hessian, delete its exact
kernel, and only then introduce a positive norm or regulator majorant.
\end{firewall}

\subsection{Exact second variation at the joint flat anchor}
\label{subsec:V100-second-variation}

Fix a flat boundary metric, a constant clock direction \(u\), and
\(K=0\).  Let an extrinsic perturbation be decomposed as
\begin{equation}
 k\longleftrightarrow
 (\dot\kappa,\dot\vartheta,\dot b,\dot\sigma),
 \label{eq:V100-k-perturbation}
\end{equation}
and a hypersurface-orthogonal clock-derivative perturbation as
\begin{equation}
 w\longleftrightarrow(\dot\theta,\dot a,\dot\chi).
 \label{eq:V100-w-perturbation}
\end{equation}
The vector and traceless components are spatial.

For a coefficient vector
\(c=(c_\kappa,c_\vartheta,c_b,c_\sigma)\), define the bilinear pairing
\begin{equation}
 B_c(k,w):=
 \dot\theta(c_\kappa\dot\kappa+c_\vartheta\dot\vartheta)
 +c_b\dot a\cdot\dot b+c_\sigma\dot\chi:\dot\sigma.
 \label{eq:V100-bilinear-pairing}
\end{equation}

\begin{theorem}[Flat-anchor Hessian]
\label{thm:V100-flat-Hessian}
For two perturbations \((k_1,w_1)\), \((k_2,w_2)\), the Hessian of the V98
action density is
\begin{equation}
 \boxed{
 D^2f_c[(k_1,w_1),(k_2,w_2)]
 =B_c(k_1,w_2)+B_c(k_2,w_1).}
 \label{eq:V100-Hessian-bilinear}
\end{equation}
In particular,
\begin{equation}
 D^2f_c[(k,w),(k,w)]=2B_c(k,w).
 \label{eq:V100-Hessian-quadratic}
\end{equation}
No measure, projector, or connection correction survives at this order.
\end{theorem}

\begin{proof}
Every V98 density is bilinear in \((K,D u)\).  At the joint anchor its value
and first differential vanish.  Differentiating once in each leg gives
exactly the polarization of
\eqref{eq:V100-bilinear-pairing}.  A second variation of the measure or a
projector is multiplied by the zero background density; a first variation
of either is multiplied by the zero first differential.  Connection and
normalization corrections likewise contain a background factor \(K\) or
\(D u\).  Hence only the displayed bilinear term remains.
\end{proof}

Let
\begin{align}
 X&:=\mathbb R^2\oplus V\oplus S_0^2,
 &Y&:=\mathbb R\oplus V\oplus S_0^2,
 \label{eq:V100-jet-spaces}
\end{align}
be the extrinsic and clock-derivative irreducible spaces.  Define
\begin{equation}
 A_c:X\longrightarrow Y,
 \qquad
 A_c(\dot\kappa,\dot\vartheta,\dot b,\dot\sigma)
 :=(c_\kappa\dot\kappa+c_\vartheta\dot\vartheta,
    c_b\dot b,c_\sigma\dot\sigma).
 \label{eq:V100-A-map}
\end{equation}
Then \(B_c(k,w)=\langle A_ck,w\rangle_Y\).

\begin{corollary}[Off-diagonal Hessian operator]
\label{cor:V100-off-diagonal-Hessian}
With the direct-sum inner product, the Hessian operator is
\begin{equation}
 \boxed{
 H_c=\begin{pmatrix}0&A_c^*\\A_c&0\end{pmatrix},
 \qquad
 H_c^2=\begin{pmatrix}A_c^*A_c&0\\0&A_cA_c^*\end{pmatrix}.}
 \label{eq:V100-Hessian-operator}
\end{equation}
Thus the canonical positive square of the exact Hessian preserves its
kernel; an arbitrary separate-leg norm need not.
\end{corollary}

\begin{proof}
The polarization formula \eqref{eq:V100-Hessian-bilinear} is the quadratic
form of the displayed block operator.  Matrix multiplication gives its
square.
\end{proof}

\subsection{Exact Hessian kernel and coefficient identifiability}
\label{subsec:V100-Hessian-kernel}

\begin{theorem}[Kernel of one combined Hessian]
\label{thm:V100-Hessian-kernel}
For every coefficient vector \(c\),
\begin{equation}
 \ker H_c=\ker A_c\oplus\ker A_c^*.
 \label{eq:V100-Hessian-kernel}
\end{equation}
If \((c_\kappa,c_\vartheta)\ne(0,0)\) and
\(c_b c_\sigma\ne0\), then
\begin{equation}
 \boxed{
 \ker H_c
 =\operatorname{span}\{(-c_\vartheta,c_\kappa,0,0;0)\}.}
 \label{eq:V100-generic-Hessian-kernel}
\end{equation}
If a displayed coefficient block vanishes, the corresponding full
irreducible domain and codomain blocks join the kernel.
\end{theorem}

\begin{proof}
For an off-diagonal operator,
\(H_c(k,w)=0\) exactly when \(A_ck=0\) and \(A_c^*w=0\).  In the generic
case the vector and traceless blocks are scalar multiples of identities.
The scalar map \(\mathbb R^2\to\mathbb R\) has kernel perpendicular to
\((c_\kappa,c_\vartheta)\) and its adjoint is injective.  This proves the
formula and the degeneration statement.
\end{proof}

\begin{theorem}[The four-action Hessian bank is coefficient-injective]
\label{thm:V100-Hessian-bank-injective}
The map from \(c\in\mathbb R^4\) to the bilinear form
\(D^2f_c\) is injective, even though every single combined Hessian has the
scalar perturbation kernel \eqref{eq:V100-generic-Hessian-kernel}.
\end{theorem}

\begin{proof}
If the bilinear form vanishes for all perturbations, choose independently
\((\dot\kappa,\dot\theta)\),
\((\dot\vartheta,\dot\theta)\),
\((\dot b,\dot a)\), and
\((\dot\sigma,\dot\chi)\).  Equation
\eqref{eq:V100-bilinear-pairing} then forces the four coefficients to zero.
The kernel of a fixed operator acts in perturbation space and must not be
confused with the kernel of the coefficient-to-Hessian map.
\end{proof}

\subsection{Factorization test for positive majorants}
\label{subsec:V100-majorant-factorization}

\begin{lemma}[Finite-dimensional kernel criterion]
\label{lem:V100-kernel-factorization}
For linear maps \(H:E\to F\) and \(\mathcal A:E\to Z\), there is a linear
\(C:F\to Z\) with \(\mathcal A=CH\) if and only if
\begin{equation}
 \ker H\subseteq\ker\mathcal A.
 \label{eq:V100-kernel-inclusion}
\end{equation}
\end{lemma}

\begin{proof}
Necessity is immediate.  Under the kernel inclusion, define
\(C(Hx)=\mathcal A x\) on \(\operatorname{Ran}H\); it is well defined.
Extend it arbitrarily to a complement of that range.
\end{proof}

\begin{proposition}[A full separate-leg square cannot factor]
\label{prop:V100-separate-square-failure}
Assume the generic coefficient conditions of
Theorem~\ref{thm:V100-Hessian-kernel}.  Let
\(\mathcal A_\oplus:X\oplus Y\to X\oplus Y\) be the identity analysis map,
whose Gram form charges the full separate-leg norm
\begin{equation}
 \|k\|_X^2+\|w\|_Y^2.
 \label{eq:V100-full-leg-majorant}
\end{equation}
Then no \(C\) satisfies \(\mathcal A_\oplus=CH_c\).  In particular, the
positive form \eqref{eq:V100-full-leg-majorant} is not the square of an
analysis map factoring through the exact Hessian.
\end{proposition}

\begin{proof}
The nonzero scalar perturbation in
\eqref{eq:V100-generic-Hessian-kernel} is killed by \(H_c\) and retained by
the identity.  Thus
\(\ker H_c\not\subseteq\ker\mathcal A_\oplus\), and
Lemma~\ref{lem:V100-kernel-factorization} applies.
\end{proof}

\begin{corollary}[Exact connected positive replacement]
\label{cor:V100-connected-positive-replacement}
The canonical positive operator \(H_c^2\) in
\eqref{eq:V100-Hessian-operator} deletes exactly the signed Hessian kernel.
Any regulator majorant intended to represent the V98 flat Hessian must be
tested against \eqref{eq:V100-kernel-inclusion}; charging the scalar kernel
line is a disconnected addition and cannot determine a V98 coefficient.
\end{corollary}

\begin{proof}
For self-adjoint \(H_c\), \(\ker H_c^2=\ker H_c\).  The failure statement
is Proposition~\ref{prop:V100-separate-square-failure}.
\end{proof}

\begin{firewall}[Young's inequality is a bound, not a factorization]
\label{fire:V100-Young-majorant}
The estimate
\(2|B_c(k,w)|\le\epsilon\|A_ck\|^2+epsilon^{-1}\|w\|^2\)
may be useful after projecting to the active image.  It does not identify a
full separate-leg square with the signed Hessian, and it must not be used to
assign counterterm coefficients on \(\ker H_c\).
\end{firewall}

\begin{assessment}[Audit-adjusted direction]
\label{ass:V100-audit-direction}
The flat-anchor regulator must be compared with the four-action signed
Hessian bank before positive point-split or norm majorants are separated.
The coefficient bank is identifiable, but each fixed combined Hessian has a
perturbation kernel which its canonical square preserves.  This is the
flat-background analogue of retaining one common stacked first variation in
V95.
\end{assessment}

\subsection{V100 result ledger and next acquisition}
\label{subsec:V100-ledger}

\paragraph{New exact conclusions.}
\begin{enumerate}[label=\textup{(\arabic*)},leftmargin=2.8em]
\item The required YM/NS audit is recorded with exact hashes.
\item The flat-anchor Hessian is the signed polarization
      \eqref{eq:V100-Hessian-bilinear}.
\item Its operator and canonical positive square are
      \eqref{eq:V100-Hessian-operator}.
\item The generic perturbation kernel is the compulsory scalar line
      \eqref{eq:V100-generic-Hessian-kernel}.
\item The coefficient-to-Hessian map for the four actions is injective.
\item The full separate-leg square fails the kernel factorization test, while
      \(H_c^2\) preserves the exact connected kernel.
\end{enumerate}

\paragraph{Next forced step.}
V101 should write the regulator's flat boundary quadratic divergence as a
block operator on \(X\oplus Y\), project out its disconnected kernel
components, and test membership in the four-dimensional signed Hessian
bank.  It must keep the V99 Codazzi-generated conditional quotient separate
from this off-shell flat calculation.

\begin{assessment}[V100 continuum status]
\label{ass:V100-continuum-status}
The full Standard--Model--Einstein continuum remains unproved.  V100 closes
the flat-anchor Hessian and its exact kernel test, but the regulator
quadratic target, remaining counterterm sectors, uniform quantum form
bounds, curved radiative equivariance, nonlinear existence, and cofinal
interacting limit remain open.
\end{assessment}

\section*{V100 exact checks and append-only record}
\addcontentsline{toc}{section}{V100 exact checks and append-only record}

\begin{equation}
\boxed{
\begin{gathered}
 D^2f_c[(k_1,w_1),(k_2,w_2)]
 =B_c(k_1,w_2)+B_c(k_2,w_1),\\
 H_c=\begin{pmatrix}0&A_c^*\\A_c&0\end{pmatrix},
 \qquad
 H_c^2=\operatorname{diag}(A_c^*A_c,A_cA_c^*),\\
 \ker H_c=\ker A_c\oplus\ker A_c^*,
 \qquad
 \ker H_c\not\subseteq\ker\mathcal A_\oplus.
\end{gathered}}
\label{eq:V100-check-card}
\end{equation}

The V99 predecessor used for this append has SHA--256
\begin{center}
\texttt{4adfea9c923a58e375b9dd7b2b1870771d6d67cb11f98c127486f017eddf4f9e}.
\end{center}
The V100 candidate retains its bytes up to the sole live final document
terminator, appends this material, and restores exactly one active
\verb|\end{document}|.  Validation includes balanced environments and
braces, unique labels, resolved new internal references, valid inline and
display math delimiters, exact predecessor-prefix equality, and a final
inventory which removes only the superseded SM V99 file.

% =============================================================================
% END APPENDED VERSION V100 MATERIAL
% =============================================================================

\clearpage
% =============================================================================
% BEGIN APPENDED VERSION V101 MATERIAL
% =============================================================================
\section{V101: exact flat quadratic regulator matching for the \(KDu\)
Hessian bank and its orthogonal failure certificate}
\label{sec:V101-regulator-matching}

V100 computed the first nonzero flat-anchor variation of the four V98
actions and warned that a separate-leg positive square need not factor
through their signed Hessian.  This note solves the corresponding finite
range problem completely.  A general self-adjoint quadratic regulator
target is decomposed into diagonal and off-diagonal blocks on the extrinsic
and clock-derivative jet spaces.  The \(KDu\) bank has zero diagonal blocks,
and its off-diagonal block is exactly the full parity-even
\(O(2)\)-equivariant homomorphism space between the two jet representations.

Consequently every parity-even equivariant off-diagonal target has unique
\(KDu\) coefficients.  Diagonal pieces, maps between inequivalent spatial
irreducibles, and nonscalar endomorphisms of the vector or traceless sectors
are explicit residuals.  They are not discarded: they are left-annihilator
certificates showing which divergence must be assigned to another
counterterm block or which regulator symmetry has failed.

\subsection{The target operator and the four-dimensional Hessian space}
\label{subsec:V101-target-space}

Use the V100 real Hilbert spaces
\begin{equation}
 X=\mathbb R^2\oplus V\oplus S_0^2,
 \qquad
 Y=\mathbb R\oplus V\oplus S_0^2,
 \qquad E=X\oplus Y.
 \label{eq:V101-jet-spaces}
\end{equation}
Here \(V\) is the two-dimensional spatial vector representation and
\(S_0^2\) the two-dimensional spatial traceless-symmetric representation.
Let a regulated flat boundary quadratic divergence be represented by a
self-adjoint operator
\begin{equation}
 R=\begin{pmatrix}
    R_{XX}&R_{YX}^*\\
    R_{YX}&R_{YY}
   \end{pmatrix}
 \quad\text{on }E.
 \label{eq:V101-general-target}
\end{equation}

Define \(\mathscr A\subset\operatorname{Hom}(X,Y)\) by
\begin{equation}
 \mathscr A:=\left\{
 A_c=\operatorname{diag}
 \bigl((c_\kappa,c_\vartheta),c_bI_V,c_\sigma I_{S_0^2}\bigr):
 c\in\mathbb R^4\right\},
 \label{eq:V101-A-space}
\end{equation}
where \((c_\kappa,c_\vartheta):\mathbb R^2\to\mathbb R\) is the scalar row
map.  The Hessian space is
\begin{equation}
 \mathscr H_{KDu}:=left\{
 H_c=\begin{pmatrix}0&A_c^*\\A_c&0\end{pmatrix}:A_c\in\mathscr A
 \right\}.
 \label{eq:V101-Hessian-space}
\end{equation}

\begin{theorem}[Exact membership criterion]
\label{thm:V101-membership}
The regulator target \(R\) is the Hessian of one V98 counterterm
combination if and only if
\begin{equation}
 \boxed{
 R_{XX}=0,
 \qquad R_{YY}=0,
 \qquad R_{YX}\in\mathscr A.}
 \label{eq:V101-membership}
\end{equation}
When these conditions hold, the coefficient vector \(c\) is unique.
\end{theorem}

\begin{proof}
Every V98 Hessian has the block form
\eqref{eq:V100-Hessian-operator}, so the conditions are necessary.  If they
hold, \(R_{YX}=A_c\) for a coefficient vector \(c\), and
\eqref{eq:V101-general-target} equals \(H_c\).  Uniqueness follows from the
coefficient injectivity in
Theorem~\ref{thm:V100-Hessian-bank-injective}.
\end{proof}

\subsection{Equivariance makes the cross-block criterion complete}
\label{subsec:V101-equivariance}

Let spatial \(O(2)\) act trivially on the scalar summands and in its standard
spin-one and spin-two representations on \(V\) and \(S_0^2\), respectively.

\begin{theorem}[Equivariant cross-block completeness]
\label{thm:V101-equivariant-completeness}
The space of parity-even \(O(2)\)-equivariant real maps from \(X\) to \(Y\)
is exactly \(\mathscr A\):
\begin{equation}
 \boxed{
 \operatorname{Hom}_{O(2)}(X,Y)=\mathscr A.}
 \label{eq:V101-equivariant-Hom}
\end{equation}
Therefore an \(O(2)\)-equivariant target with zero diagonal blocks always
has a unique V98 counterterm preimage.
\end{theorem}

\begin{proof}
The scalar domain contains two copies of the trivial representation and the
scalar codomain one, so its equivariant homomorphisms form an arbitrary
real row \((c_\kappa,c_\vartheta)\).  The real spatial vector and traceless
symmetric representations are irreducible and inequivalent under
\(O(2)\).  Schur's lemma makes each equivariant endomorphism a real scalar
multiple of the identity and kills maps between inequivalent irreducibles.
These are exactly the four parameters in \eqref{eq:V101-A-space}.
\end{proof}

\begin{remark}[Why spatial parity is included]
\label{rem:V101-parity}
For the connected rotation group \(SO(2)\), a real two-dimensional rotation
representation also admits the complex-structure intertwiner.  Spatial
parity removes that additional pseudoscalar map.  If parity is not imposed,
the vector and spin-two blocks each acquire an independent orientation-odd
intertwiner and the four-action bank must be enlarged before claiming
completeness.
\end{remark}

\subsection{Orthogonal projection and coefficient extraction}
\label{subsec:V101-projection}

Equip all operator spaces with their Hilbert--Schmidt inner products.  Let
\(P_s^X,P_V^X,P_S^X\) and \(P_s^Y,P_V^Y,P_S^Y\) be the irreducible
projections.  Define
\begin{align}
 C_s&:=P_s^YR_{YX}P_s^X:\mathbb R^2\to\mathbb R,
 \notag\\
 C_V&:=P_V^YR_{YX}P_V^X:V\to V,
 &C_S&:=P_S^YR_{YX}P_S^X:S_0^2\to S_0^2.
 \label{eq:V101-diagonal-cross-blocks}
\end{align}

\begin{proposition}[Best \(KDu\) coefficient projection]
\label{prop:V101-best-projection}
The Hilbert--Schmidt orthogonal projection of \(R_{YX}\) onto
\(\mathscr A\) is
\begin{equation}
 \boxed{
 \Pi_{\mathscr A}R_{YX}
 =\operatorname{diag}\left(
 C_s,
 \frac{\operatorname{tr}C_V}{2}I_V,
 \frac{\operatorname{tr}C_S}{2}I_{S_0^2}
 \right).}
 \label{eq:V101-A-projection}
\end{equation}
Thus the extracted coefficients are the two entries of \(C_s\) and
\begin{equation}
 c_b=\frac12\operatorname{tr}C_V,
 \qquad
 c_\sigma=\frac12\operatorname{tr}C_S.
 \label{eq:V101-coefficient-extraction}
\end{equation}
\end{proposition}

\begin{proof}
Different irreducible source--target blocks are Hilbert--Schmidt orthogonal.
The scalar-to-scalar block already belongs to its full allowed space.  The
orthogonal projection of an endomorphism of a two-dimensional space onto
the line spanned by its identity is
\((\operatorname{tr}C/2)I\).  Applying this independently to \(V\) and
\(S_0^2\) proves the formula.
\end{proof}

Define the fitted Hessian
\begin{equation}
 H_{\rm fit}:=
 \begin{pmatrix}
 0&(\Pi_{\mathscr A}R_{YX})^*\\
 \Pi_{\mathscr A}R_{YX}&0
 \end{pmatrix}
 \label{eq:V101-fitted-Hessian}
\end{equation}
and the residual \(R_\perp:=R-H_{\rm fit}\).

\begin{theorem}[Orthogonal failure certificate]
\label{thm:V101-orthogonal-certificate}
The decomposition
\begin{equation}
 \boxed{R=H_{\rm fit}+R_\perp}
 \label{eq:V101-target-decomposition}
\end{equation}
is Hilbert--Schmidt orthogonal, and
\begin{equation}
 R\in\mathscr H_{KDu}
 \quad\Longleftrightarrow\quad R_\perp=0.
 \label{eq:V101-residual-zero-test}
\end{equation}
If \(R_\perp\ne0\), the functional
\begin{equation}
 \Lambda_{R_\perp}(T):=langle R_\perp,T\rangle_{\rm HS}
 \label{eq:V101-left-annihilator}
\end{equation}
annihilates every V98 Hessian but satisfies
\(\Lambda_{R_\perp}(R)=\|R_\perp\|_{\rm HS}^2>0\).
\end{theorem}

\begin{proof}
The diagonal blocks of every Hessian in \(\mathscr H_{KDu}\) are zero, and
\(\Pi_{\mathscr A}\) is the orthogonal projection of the cross block.
Hence \(R_\perp\) is orthogonal to the whole Hessian space.  This proves the
decomposition and zero test.  Orthogonality gives
\(\Lambda_{R_\perp}(H_c)=0\) and
\(\langle R_\perp,R\rangle=\|R_\perp\|^2\).
\end{proof}

\begin{corollary}[Disposition of disconnected positive pieces]
\label{cor:V101-diagonal-disposition}
Every nonzero diagonal block \(R_{XX}\) or \(R_{YY}\) belongs to
\(R_\perp\).  It cannot be cancelled by a V98 \(KDu\) counterterm and must
be matched against another invariant action block.  It may be deleted only
if the exact regulated observable, before majorization, proves that block is
zero.
\end{corollary}

\begin{proof}
This is the diagonal part of
Theorems~\ref{thm:V101-membership} and
\ref{thm:V101-orthogonal-certificate}.  The final sentence distinguishes an
exact zero from a positive upper bound, as required by the V100 kernel test.
\end{proof}

\begin{firewall}[Flat Hessian matching is not curved stacked matching]
\label{fire:V101-flat-not-curved}
The coefficients extracted by \eqref{eq:V101-coefficient-extraction} also
fix the curved first-normal and clock columns
\eqref{eq:V98-KDu-normal-symbol} and \eqref{eq:V99-clock-Euler}.  Matching a
flat quadratic target does not prove that those curved regulator targets
match the same coefficients.  The latter remains the common stacked V95
range test.
\end{firewall}

\begin{assessment}[Closed flat \(KDu\) range problem]
\label{ass:V101-flat-range-closed}
Within the parity-even flat \(KDu\) Hessian sector, the counterterm range is
now completely characterized.  Equivariance makes every off-diagonal target
accessible, coefficients are extracted by
\eqref{eq:V101-A-projection}, and every inaccessible component has the
explicit annihilator \eqref{eq:V101-left-annihilator}.  This closes the
finite-dimensional range calculation for this block without claiming the
regulator actually satisfies the criterion.
\end{assessment}

\subsection{V101 result ledger and next acquisition}
\label{subsec:V101-ledger}

\paragraph{New exact conclusions.}
\begin{enumerate}[label=\textup{(\arabic*)},leftmargin=2.8em]
\item The V98 Hessian membership criterion is
      \eqref{eq:V101-membership}.
\item The four cross-block parameters are exactly all parity-even
      \(O(2)\)-equivariant homomorphisms from \(X\) to \(Y\).
\item Formula \eqref{eq:V101-A-projection} gives the unique best-fit
      coefficients for an arbitrary target.
\item The residual zero test and left annihilator are
      \eqref{eq:V101-residual-zero-test}--
      \eqref{eq:V101-left-annihilator}.
\item Diagonal positive pieces are assigned to other action blocks unless
      an exact pre-majorant identity deletes them.
\end{enumerate}

\paragraph{Next forced step.}
V102 should derive a concrete point-split flat regulator target for the
clock/extrinsic boundary composite, before any separate-square majorant, and
apply \eqref{eq:V101-residual-zero-test}.  If no such composite is selected
by the current action, the programme must go upstream and derive it from the
second variation of the V89 common cylinder plus Einstein boundary action.

\begin{assessment}[V101 continuum status]
\label{ass:V101-continuum-status}
The full Standard--Model--Einstein continuum remains unproved.  V101 closes
the abstract flat \(KDu\) Hessian range problem, but the physical regulator
target, remaining counterterm sectors, uniform quantum form bounds, curved
radiative equivariance, nonlinear existence, and cofinal interacting limit
remain open.
\end{assessment}

\section*{V101 exact checks and append-only record}
\addcontentsline{toc}{section}{V101 exact checks and append-only record}

\begin{equation}
\boxed{
\begin{gathered}
 R\in\mathscr H_{KDu}
 \Longleftrightarrow R_{XX}=R_{YY}=0\ \text{and}\ R_{YX}\in\mathscr A,\\
 \operatorname{Hom}_{O(2)}(X,Y)=\mathscr A,\\
 R=H_{\rm fit}+R_\perp,
 \qquad
 R_\perp=0\Longleftrightarrow R\in\mathscr H_{KDu},\\
 \Lambda_{R_\perp}(H_c)=0,
 \qquad
 \Lambda_{R_\perp}(R)=\|R_\perp\|_{\rm HS}^2.
\end{gathered}}
\label{eq:V101-check-card}
\end{equation}

The V100 predecessor used for this append has SHA--256
\begin{center}
\texttt{6660bd3054f4161add67717a6bbe68c2d6bef82fe792cb5931e415bde62dad44}.
\end{center}
The V101 candidate retains its bytes up to the sole live final document
terminator, appends this material, and restores exactly one active
\verb|\end{document}|.  Validation includes balanced environments and
braces, unique labels, resolved new internal references, valid inline and
display math delimiters, exact predecessor-prefix equality, and a final
inventory which removes only the superseded SM V100 file.

% =============================================================================
% END APPENDED VERSION V101 MATERIAL
% =============================================================================

\clearpage
% =============================================================================
% BEGIN APPENDED VERSION V102 MATERIAL
% =============================================================================
\section{V102: absence of a tree-level \(KDu\) source in the V73/V89
subparent and the boundary-realization obstruction to a quantum target}
\label{sec:V102-physical-target}

V101 solved the abstract flat quadratic range problem for the four
\(KDu\) counterterms and requested a physical regulator target.  Returning
to the action shows that no such target is present in the minimal classical
subparent currently printed.  The V73 Einstein--GHY sector contains the
metric normal jet but no clock congruence.  The V89 cylinder sector contains
the clock congruence but no extrinsic curvature.  Additivity therefore makes
their mixed \(K\)--\(D u\) Hessian vanish identically.

A nonzero target could still be generated by additional boundary couplings
or by quantum loops.  The latter is not determined by the bulk differential
expression alone: it depends on the boundary realization.  An exact
half-line heat-kernel calculation gives opposite leading boundary
coefficients for Dirichlet and Neumann conditions.  Thus the programme must
select the complete gauge-fixed operator/boundary-condition pair before a
point-split or heat-kernel target can be inserted into the V101 range test.

\subsection{The printed classical subparent has zero mixed Hessian}
\label{subsec:V102-tree-zero}

Let \(S_{E,G}\) denote the V73 Einstein--Hilbert--GHY compatible-jet parent
after its auxiliary variables are eliminated, and let \(S_c\) be the V89
cylinder action.  For the present calculation retain the common induced
metric but display their other arguments:
\begin{equation}
 S_{\rm sub}[\gamma,K;u,\varphi]
 :=S_{E,G}[\gamma,K]+S_c[\gamma,u,\varphi].
 \label{eq:V102-subparent}
\end{equation}
The common-trace substitution \(z=\mathcal A(\gamma)\) does not introduce
either \(K\) into \(S_c\) or \(u\) into \(S_{E,G}\).

\begin{theorem}[Tree-level cross-Hessian vanishing]
\label{thm:V102-tree-cross-zero}
For arbitrary compactly supported variations of the independent boundary
jets,
\begin{equation}
 \boxed{
 D_KD_{D u}S_{\rm sub}=0.}
 \label{eq:V102-tree-cross-zero}
\end{equation}
In particular, at the flat constant-clock anchor the V101 off-diagonal
target of this classical subparent is zero and the unique fitted V98
coefficient vector is
\begin{equation}
 c_\kappa=c_\vartheta=c_b=c_\sigma=0.
 \label{eq:V102-tree-coefficients-zero}
\end{equation}
\end{theorem}

\begin{proof}
The reduced Einstein--GHY action is independent of \(u\) and its tangential
jets.  The cylinder density \eqref{eq:V89-cylinder-action} depends on
\(\gamma\), \(u\), and \(\varphi\), but contains no \(K\).  Hence each
summand in \eqref{eq:V102-subparent} is annihilated by one of the two mixed
derivatives.  Their sum satisfies \eqref{eq:V102-tree-cross-zero}.  The V101
coefficient-to-Hessian map is injective, so a zero target has the unique
preimage \eqref{eq:V102-tree-coefficients-zero}.
\end{proof}

\begin{corollary}[What can generate a nonzero target]
\label{cor:V102-nonzero-source}
A nonzero \(KDu\) quadratic target requires at least one of:
\begin{enumerate}[label=\textup{(\alph*)},leftmargin=2.8em]
\item an additional classical boundary density coupling \(K\) to the clock
      congruence;
\item elimination of an auxiliary field whose solution depends jointly on
      \(K\) and \(D u\); or
\item a quantum effective action for a fully specified boundary realization.
\end{enumerate}
None of these sources is supplied by the additive subparent
\eqref{eq:V102-subparent} itself.
\end{corollary}

\begin{proof}
This is the contrapositive of Theorem~\ref{thm:V102-tree-cross-zero}, with
the three listed mechanisms recording how the separated dependence could be
changed before taking the Hessian.
\end{proof}

\begin{firewall}[This is not a no-renormalization theorem]
\label{fire:V102-not-no-renormalization}
Equation \eqref{eq:V102-tree-cross-zero} concerns one classical subparent.
It does not say that Standard Model, ghost, graviton, clock, or cylinder
loops fail to generate \(KDu\) counterterms.  It says that such coefficients
cannot be inferred from the classical mixed Hessian currently printed.
\end{firewall}

\subsection{Exact boundary-condition dependence of a heat coefficient}
\label{subsec:V102-half-line-witness}

Consider the nonnegative scalar operator \(P=-\partial_x^2\) on the
half-line.  Its Dirichlet and Neumann heat kernels are obtained by images:
\begin{align}
 K_D(t;x,y)&=(4\pi t)^{-1/2}
 \left(e^{-(x-y)^2/(4t)}-e^{-(x+y)^2/(4t)}\right),
 \label{eq:V102-Dirichlet-kernel}\\
 K_N(t;x,y)&=(4\pi t)^{-1/2}
 \left(e^{-(x-y)^2/(4t)}+e^{-(x+y)^2/(4t)}\right).
 \label{eq:V102-Neumann-kernel}
\end{align}
Let \(K_{\mathbb R}(t;x,x)=(4\pi t)^{-1/2}\) be the full-line diagonal
density.

\begin{theorem}[Opposite boundary heat coefficients]
\label{thm:V102-boundary-heat-sign}
The integrated boundary corrections per endpoint are exactly
\begin{align}
 \int_0^\infty
  \left(K_D(t;x,x)-K_{\mathbb R}(t;x,x)\right)\,\dd x
 &=-\frac14,
 \label{eq:V102-Dirichlet-quarter}\\
 \int_0^\infty
  \left(K_N(t;x,x)-K_{\mathbb R}(t;x,x)\right)\,\dd x
 &=+\frac14.
 \label{eq:V102-Neumann-quarter}
\end{align}
Thus the same bulk differential expression has different, indeed opposite,
leading boundary heat coefficients under two admissible realizations.
\end{theorem}

\begin{proof}
On the diagonal, the image term is
\(\mp(4\pi t)^{-1/2}e^{-x^2/t}\).  Since
\begin{equation*}
 \int_0^\infty e^{-x^2/t}\,\dd x=\frac12\sqrt{\pi t},
\end{equation*}
its integral is
\(\mp(4\pi t)^{-1/2}\tfrac12\sqrt{\pi t}=\mp\tfrac14\).
\end{proof}

\begin{corollary}[The quantum boundary target needs the realization]
\label{cor:V102-realization-required}
No boundary heat-kernel or determinant divergence is uniquely determined
by a bulk differential expression alone.  At minimum the operator domain and
boundary conditions must be specified.  For a gauge theory the physical
input is the complete gauge-fixed field/ghost operator together with a
compatible boundary domain; sectorwise boundary choices do not define one
BRST or Ward-compatible effective action.
\end{corollary}

\begin{proof}
Theorem~\ref{thm:V102-boundary-heat-sign} gives two self-adjoint domains for
the same expression with distinct boundary coefficients.  Hence the
expression alone cannot determine the target.  The gauge statement follows
because the determinant and its Ward identity use the combined complex and
its common domain, as already required by the V90 common-action gate.
\end{proof}

\subsection{Exact data contract for the next regulator computation}
\label{subsec:V102-data-contract}

\begin{proposition}[Necessary regulator input]
\label{prop:V102-regulator-input}
Before applying the V101 range test to a quantum target, the following data
must be fixed on one declared branch:
\begin{enumerate}[label=\textup{(R\arabic*)},leftmargin=3.2em]
\item the gauge-fixed graviton, Standard Model, clock/cylinder, and ghost
      differential operators, including all mixing terms;
\item one common boundary domain, with the Higgs/Yukawa and neutrino branch,
      gauge/ghost compatibility, and corner completion stated;
\item a regulator and subtraction prescription preserving the combined
      gauge and tangential-diffeomorphism Ward identities up to a printed
      defect;
\item the exact pre-majorant quadratic boundary kernel on \(X\oplus Y\);
\item its diagonal/off-diagonal and \(O(2)\)-irreducible decomposition as in
      \eqref{eq:V101-general-target}.
\end{enumerate}
Without \textup{(R1)}--\textup{(R5)}, a numerical or symbolic
\(KDu\) coefficient is not defined by the current programme.
\end{proposition}

\begin{proof}
Items \textup{(R1)}--\textup{(R2)} define the operator realization; their
necessity is witnessed by
Theorem~\ref{thm:V102-boundary-heat-sign}.  Item \textup{(R3)} is the V90
common Ward requirement.  Items \textup{(R4)}--\textup{(R5)} are exactly the
input required by the V100 kernel test and V101 membership theorem.
\end{proof}

\begin{assessment}[Upstream bottleneck after V102]
\label{ass:V102-upstream-bottleneck}
The abstract counterterm range calculation is ahead of the physical target.
The minimal printed classical subparent has zero \(KDu\) cross Hessian, while
the quantum coefficient depends on boundary-domain choices which remain
branched in V32 and V90.  The programme should therefore select and test one
complete linearized boundary realization before attempting a heat-kernel
coefficient table.
\end{assessment}

\subsection{V102 result ledger and next acquisition}
\label{subsec:V102-ledger}

\paragraph{New exact conclusions.}
\begin{enumerate}[label=\textup{(\arabic*)},leftmargin=2.8em]
\item The V73 Einstein--GHY plus V89 cylinder subparent has the exact
      tree-level cross-Hessian zero \eqref{eq:V102-tree-cross-zero}.
\item Its unique fitted V98 coefficient vector is zero.
\item The half-line kernels
      \eqref{eq:V102-Dirichlet-kernel}--
      \eqref{eq:V102-Neumann-kernel} give opposite boundary coefficients
      \(-1/4\) and \(+1/4\).
\item Therefore the quantum boundary target requires a complete operator
      realization, not merely a bulk differential expression.
\item Conditions \textup{(R1)}--\textup{(R5)} are the exact data contract
      before V101 can be applied physically.
\end{enumerate}

\paragraph{Next forced step.}
V103 should choose the conservative linearized branch already constructed
in V31--V36: the V32 Higgs-floor fermion graph where defined, its explicitly
declared neutrino alternative, the V36 GHY--de Donder gravitational domain,
and compatible gauge/ghost boundary data.  It should test strong
ellipticity after Euclideanization sector by sector while retaining the
common BRST domain, and stop with a failure certificate if no single branch
satisfies all domain conditions.

\begin{assessment}[V102 continuum status]
\label{ass:V102-continuum-status}
The full Standard--Model--Einstein continuum remains unproved.  V102 proves
the absence of a classical \(KDu\) source in the printed subparent and the
necessity of boundary-realization data for a quantum target.  Selection of
that realization, heat coefficients, common quantum form bounds, curved
radiative equivariance, nonlinear existence, and the cofinal interacting
limit remain open.
\end{assessment}

\section*{V102 exact checks and append-only record}
\addcontentsline{toc}{section}{V102 exact checks and append-only record}

\begin{equation}
\boxed{
\begin{gathered}
 D_KD_{D u}(S_{E,G}+S_c)=0,
 \qquad c_\kappa=c_\vartheta=c_b=c_\sigma=0,\\
 \int_0^\infty(K_D-K_{\mathbb R})(t;x,x)\,\dd x=-\frac14,
 \qquad
 \int_0^\infty(K_N-K_{\mathbb R})(t;x,x)\,\dd x=+\frac14,\\
 \text{quantum boundary target}
 \Longrightarrow\text{complete operator/domain/regulator data}.
\end{gathered}}
\label{eq:V102-check-card}
\end{equation}

The V101 predecessor used for this append has SHA--256
\begin{center}
\texttt{7be03ccbda2c834e30ccb6f5be66c8fa22cc2c83345e05ffa9491414acddfc3a}.
\end{center}
The V102 candidate retains its bytes up to the sole live final document
terminator, appends this material, and restores exactly one active
\verb|\end{document}|.  Validation includes balanced environments and
braces, unique labels, resolved new internal references, valid inline and
display math delimiters, exact predecessor-prefix equality, and a final
inventory which removes only the superseded SM V101 file.

% =============================================================================
% END APPENDED VERSION V102 MATERIAL
% =============================================================================

\clearpage
% =============================================================================
% BEGIN APPENDED VERSION V103 MATERIAL
% =============================================================================
\section{V103: assembly of the conservative boundary branch and the exact
common-floor failure inherited from the Yang--Mills/ghost glancing mode}
\label{sec:V103-common-domain-failure}

V102 proved that quantum boundary coefficients require one complete
operator/domain realization.  The manuscript already contains the relevant
sector decisions, but they were proved at different stages.  V64 supplies a
uniformly complementing metric--continuum axial block.  V32 supplies a
conditional Higgs-dependent maximal-neutral fermion graph only after a
right-handed-neutrino completion and positive Higgs/Yukawa floors.  V35
proves that the relative Yang--Mills/ghost branch has an unweighted glancing
singular sequence.

This note assembles those results without recomputing their sector symbols.
At principal order the frozen Standard Model blocks are a direct sum; the
gauge, Higgs, and Yukawa couplings which mix sectors are lower order.  The
least singular value of a direct sum is the minimum of the sector values.
Consequently the V35 Yang--Mills/ghost witness embeds into the complete
candidate domain and forces its unweighted common complementing floor to
zero, even if every conditional fermion hypothesis is granted.  The
operator realization requested by V102 is therefore not yet available on
this conservative branch.

\subsection{Declared conservative branch}
\label{subsec:V103-declared-branch}

Consider the following attempted principal boundary realization:
\begin{enumerate}[label=\textup{(B\arabic*)},leftmargin=3.2em]
\item the V64 metric--continuum axial stable matrix
      \(\mathcal B_{64}\);
\item the V35 relative Yang--Mills/ghost boundary matrix
      \(\mathcal B_{\rm YMgh}\) in its unweighted stable norm;
\item the V32 Higgs-polar fermion graph \(\mathcal B_F(\Phi)\), conditional
      on a right-handed-neutrino completion,
      \(\eta_H^{\rm cert}>0\), and
      \(\beta_{\rm graph}^{\rm cert}>0\);
\item the scalar/Higgs boundary block \(\mathcal B_H\) satisfying its V20
      complementing and V32 mixed-conormal conditions.
\end{enumerate}
At frozen principal order write
\begin{equation}
 \mathcal B_{\rm cons}
 :=\mathcal B_{64}\oplus\mathcal B_{\rm YMgh}
   \oplus\mathcal B_F(\Phi)\oplus\mathcal B_H.
 \label{eq:V103-conservative-direct-sum}
\end{equation}
This direct sum is only a candidate.  Its lower-order coupled realization
would still have to obey the V90 common Ward and V95 common-coefficient
conditions.

\begin{lemma}[Direct-sum singular floor]
\label{lem:V103-direct-sum-floor}
For finite-dimensional Hilbert-space maps \(B_i:E_i\to F_i\),
\begin{equation}
 \boxed{
 s_{\min}\!\left(\bigoplus_iB_i\right)
 =\min_i s_{\min}(B_i).}
 \label{eq:V103-direct-sum-floor}
\end{equation}
The same statement holds pointwise for a continuous family of frozen
boundary symbols.
\end{lemma}

\begin{proof}
For \(x=(x_i)\),
\begin{equation*}
 \left\|\bigoplus_iB_ix_i\right\|^2
 =\sum_i\|B_ix_i\|^2
 \ge\left(\min_i s_{\min}(B_i)^2\right)\sum_i\|x_i\|^2.
\end{equation*}
This gives the lower bound.  Taking a unit vector supported in a sector
whose least singular value realizes the minimum gives equality.  Apply the
identity at each frozen frequency for a family.
\end{proof}

\subsection{Two independent branch obstructions}
\label{subsec:V103-two-obstructions}

\begin{proposition}[Minimal-neutrino branch is unavailable]
\label{prop:V103-minimal-neutrino-obstruction}
With minimal Standard Model field content and no right-handed neutrino, the
full V32 Higgs-polar fermion graph cannot be used in
\eqref{eq:V103-conservative-direct-sum}.  Its lepton intertwiner has unequal
trace ranks, so no invertible polar factor defines the required maximal
neutral graph.
\end{proposition}

\begin{proof}
This is exactly the rank defect in
Corollary~\ref{cor:V32-minimal-fork}.  V103 uses it only to reject the
minimal-content choice in the assembled domain; it does not repeat the V32
representation calculation.
\end{proof}

Thus the conservative test must either grant the conditional
right-handed-neutrino/Higgs-floor branch or construct one of the other V32
neutrino alternatives.  Grant the former temporarily to isolate the next
obstruction.

\begin{theorem}[Common unweighted floor fails on the relative gauge branch]
\label{thm:V103-common-floor-failure}
Assume all V32 fermion/Higgs hypotheses and the positive V64
metric--continuum floor.  On the V35 relative Yang--Mills/ghost branch,
there is a normalized stable-frequency sequence \(\zeta_j\) and unit gauge
amplitudes \(y_j\) such that
\begin{equation}
 \|\mathcal B_{\rm YMgh}(\zeta_j)y_j\|\longrightarrow0.
 \label{eq:V103-YM-singular-sequence}
\end{equation}
Embedding \(y_j\) with zero amplitudes in all other sectors gives
\begin{equation}
 \boxed{
 \inf_{\zeta\in\overline{\mathbb S}_+}
 s_{\min}(\mathcal B_{\rm cons}(\zeta))=0.}
 \label{eq:V103-common-floor-zero}
\end{equation}
Hence the conservative branch is not uniformly complementing in the
unweighted norm and does not supply the realization required by the current
V102 heat-kernel branch.  No assertion is made here about every possible
Euclidean continuation or weighted boundary calculus.
\end{theorem}

\begin{proof}
The sequence \eqref{eq:V103-YM-singular-sequence} is the V35 relative
Yang--Mills glancing witness recorded in
Proposition~\ref{prop:V35-YM-glancing} and the V35 physical-status firewall.
Set every non-gauge component of the full amplitude to zero.  The direct-sum
image norm is then exactly the Yang--Mills/ghost image norm and tends to
zero.  Equivalently, apply Lemma~\ref{lem:V103-direct-sum-floor}.  A uniform
complementing floor is necessary for this standard strongly elliptic
boundary realization, so its failure prevents the stated heat-kernel step.
\end{proof}

\begin{corollary}[Positive gravity cannot repair a decoupled gauge kernel]
\label{cor:V103-gravity-cannot-repair}
The positive V64 floor \(\rho_{64}(m,v)>0\) does not lift the V35 gauge
singular sequence in the principal direct sum.  Nor do positive conditional
fermion or Higgs floors lift it.
\end{corollary}

\begin{proof}
The embedded witness has zero amplitude in those sectors.  Their boundary
operators therefore contribute zero to its image, regardless of their
individual lower bounds.
\end{proof}

\begin{firewall}[Lower-order mixing is not a principal repair]
\label{fire:V103-lower-order-no-repair}
Gauge--Higgs masses, Yukawa matrices, background connections, and finite
extrinsic curvature enter lower-order frozen terms.  They may move isolated
finite-frequency eigenvalues but do not change the homogeneous glancing
principal symbol or create a uniform high-frequency floor.  A repair must
change the principal Yang--Mills/ghost boundary realization, its norm, or
the field content at the boundary.
\end{firewall}

\subsection{Exact alternatives after the failure}
\label{subsec:V103-alternatives}

\begin{proposition}[Four immediate programme responses]
\label{prop:V103-immediate-responses}
Within the current programme, the failure
\eqref{eq:V103-common-floor-zero} has at least the following four immediate, logically distinct responses:
\begin{enumerate}[label=\textup{(\alph*)},leftmargin=2.8em]
\item construct a different local action/BRST-compatible principal
      Yang--Mills/ghost boundary graph with a positive unweighted glancing
      floor;
\item prove complementing in a weighted anisotropic boundary norm and
      rebuild the heat calculus in that norm;
\item add a genuine boundary reservoir whose principal coupling pairs the
      glancing gauge mode and include its stress and ghosts;
\item abandon the timelike-boundary heat-kernel branch for that sector and
      use a boundaryless or characteristic realization.
\end{enumerate}
Merely adjoining the V64, V32, or scalar/Higgs blocks cannot constitute a
fifth response.
\end{proposition}

\begin{proof}
Theorem~\ref{thm:V103-common-floor-failure} locates the zero sequence in the
principal Yang--Mills/ghost block.  Each listed response changes one of the data responsible for the failed
principal estimate: the boundary map, the norm and calculus, the principal
field content, or the presence of the boundary problem itself.  Corollary
\ref{cor:V103-gravity-cannot-repair} excludes unchanged decoupled sectors.
\end{proof}

\begin{assessment}[Boundary-realization decision]
\label{ass:V103-domain-decision}
The data contract of V102 cannot yet be populated on the conservative local
branch.  Minimal fermion content fails before complementing is tested.  The
conditional right-handed-neutrino branch removes that rank defect but leaves
the independent V35 Yang--Mills/ghost glancing obstruction.  A heat-kernel
counterterm table produced without repairing or explicitly changing this
gauge domain would not belong to one common Standard Model--Einstein
realization.
\end{assessment}

\subsection{V103 result ledger and next acquisition}
\label{subsec:V103-ledger}

\paragraph{New exact conclusions.}
\begin{enumerate}[label=\textup{(\arabic*)},leftmargin=2.8em]
\item The conservative principal domain is assembled in
      \eqref{eq:V103-conservative-direct-sum} without repeating sector
      calculations.
\item Direct-sum least singular values obey
      \eqref{eq:V103-direct-sum-floor}.
\item The minimal-neutrino branch is excluded by the V32 rank defect.
\item Even granting the completed Higgs-floor branch, the embedded V35
      gauge witness proves the common-floor failure
      \eqref{eq:V103-common-floor-zero}.
\item Positive metric--continuum, fermion, and Higgs floors cannot lift a
      witness supported wholly in the gauge sector.
\item Proposition~\ref{prop:V103-immediate-responses} records the four
      available upstream repairs.
\end{enumerate}

\paragraph{Next forced step.}
V104 should return to the V35 Yang--Mills/ghost glancing symbol and test the
smallest principal reservoir completion.  It should derive the extended
Green form, require an action-Lagrangian and BRST-invariant graph, compute
the exact stable matrix, and either prove a uniform floor or exhibit the
surviving singular sequence.  The fermion branch choice remains explicit
and must not be silently fixed during that repair.

\begin{assessment}[V103 continuum status]
\label{ass:V103-continuum-status}
The full Standard--Model--Einstein continuum remains unproved.  V103 proves
that the current conservative common boundary realization fails its
unweighted Yang--Mills/ghost glancing floor.  A principal gauge-domain
repair, physical branch selection, quantum coefficients, common form
bounds, curved radiative equivariance, nonlinear existence, and the cofinal
interacting limit remain open.
\end{assessment}

\section*{V103 exact checks and append-only record}
\addcontentsline{toc}{section}{V103 exact checks and append-only record}

\begin{equation}
\boxed{
\begin{gathered}
 s_{\min}(\bigoplus_iB_i)=\min_i s_{\min}(B_i),\\
 \|\mathcal B_{\rm YMgh}(\zeta_j)y_j\|\to0,
 \qquad
 \inf_{\overline{\mathbb S}_+}s_{\min}(\mathcal B_{\rm cons})=0,\\
 \rho_{64}>0\quad\text{does not lift a gauge-supported singular sequence}.
\end{gathered}}
\label{eq:V103-check-card}
\end{equation}

The V102 predecessor used for this append has SHA--256
\begin{center}
\texttt{075d93ed8d71a05e1f611cb8b899fa9c77e6f413e54aeef1bdaadbe95bf02065}.
\end{center}
The V103 candidate retains its bytes up to the sole live final document
terminator, appends this material, and restores exactly one active
\verb|\end{document}|.  Validation includes balanced environments and
braces, unique labels, resolved new internal references, valid inline and
display math delimiters, exact predecessor-prefix equality, and a final
inventory which removes only the superseded SM V102 file.

% =============================================================================
% END APPENDED VERSION V103 MATERIAL
% =============================================================================


\clearpage
% =============================================================================
% BEGIN APPENDED VERSION V104 MATERIAL
% =============================================================================
\section{V104: the reciprocal boundary-reservoir test and a BRST-protected
Yang--Mills glancing vector}
\label{sec:V104-reservoir-no-go}

V103 isolated the relative Yang--Mills/ghost block as the first failed
unweighted complementing factor in the conservative common domain.  The next
candidate repair is to attach a boundary field to the unconstrained normal
gauge trace.  This note carries out that test for the smallest variational
class which is closed under taking finitely many reservoir components.

The result is a precise obstruction rather than a repair.  A reciprocal
quadratic boundary coupling has enough algebraic freedom to replace the bare
Neumann factor at glancing.  However, if the relative Dirichlet ghost domain
is kept and the coupled boundary Euler rows are required to be invariant under
the linearized BRST differential, the same identities which express BRST
incidence produce an explicit stable near-kernel.  Thus adding a passive
finite reservoir to the V35 relative branch cannot by itself populate the
V102 operator/domain card.  The proof is at frozen principal order and does
not assert a no-go theorem for boundary BRST multiplets with their own ghost
equations or for a changed ghost boundary graph.

\subsection{The reciprocal variational class}
\label{subsec:V104-reciprocal-class}

Work first with one realified adjoint component.  Retain the V35 tangential
relative condition
\begin{equation}
 a_\parallel=0
 \label{eq:V104-relative-tangential}
\end{equation}
and write (a:=a_n), (p:=D_na_n).  Let (E) be a finite-dimensional real
inner-product space and let (eta\in E) be a boundary reservoir.  Consider
the local quadratic boundary action
\begin{equation}
 S_\partial[a,\eta]
 :=
 \frac12\int_{\mathcal T}\langle\eta,L\eta\rangle_E
 +\int_{\mathcal T}\langle\eta,Ta\rangle_E,
 \qquad L=L^\dagger,
 \label{eq:V104-boundary-action}
\end{equation}
where (L) is a formally self-adjoint tangential differential operator on
(E) and (T) is a tangential differential operator from the scalar trace
to (E).  Formal adjoints include tangential integration by parts.  Combining
the bulk conormal term with \eqref{eq:V104-boundary-action} gives the two
boundary Euler rows
\begin{align}
 F_1&:=\Gamma+T^\dagger\eta=0,
 \label{eq:V104-first-row}\\
 F_2&:=L\eta+Ta=0.
 \label{eq:V104-second-row}
\end{align}
On \eqref{eq:V104-relative-tangential}, the principal part of the Lorenz row
is (Gamma=p), as in \eqref{eq:V35-relative-YM}.

\begin{proposition}[Extended Green identity]
\label{prop:V104-extended-Green}
The coupled rows \eqref{eq:V104-first-row}--%
\eqref{eq:V104-second-row} are variationally reciprocal.  If
((a,p,\eta)) and ((b,q,\xi)) satisfy their homogeneous linearizations,
then
\begin{equation}
 \int_{\mathcal T}
 \bigl(\langle p,b\rangle-\langle a,q\rangle\bigr)
 =
 \int_{\mathcal T}
 \bigl(\langle\eta,L\xi\rangle_E
       -\langle L\eta,\xi\rangle_E\bigr).
 \label{eq:V104-extended-Green}
\end{equation}
The right-hand side is the integrated tangential Green concomitant of (L).
It vanishes on a closed boundary cylinder, or equals the reservoir endpoint
symplectic flux when the time faces are retained.  Hence the bulk Green flux
plus that reservoir endpoint flux is zero.
\end{proposition}

\begin{proof}
The first row gives (p=-T^\dagger\eta) and
(q=-T^\dagger\xi).  Formal adjunction therefore gives
\[
 \int_{\mathcal T}(\langle p,b\rangle-\langle a,q\rangle)
 =
 \int_{\mathcal T}
 \bigl(-\langle\eta,Tb\rangle_E
       +\langle Ta,\xi\rangle_E\bigr).
\]
The second row gives (Tb=-L\xi) and (Ta=-L\eta).  Substitution proves
\eqref{eq:V104-extended-Green}.  Formal self-adjointness makes its integrand
a tangential divergence, which proves the last assertion by Stokes' theorem.
\end{proof}

Let (ell(\zeta)), (t(\zeta)), and (t^\dagger(\zeta)) denote the
homogeneous principal symbols of (L,T,T^\dagger), in the same complex
stable-frequency convention as V35.  After the three identity rows for
(a_\parallel), the new stable matrix on the amplitudes ((a,\eta)) is
\begin{equation}
 \mathcal B_{\rm res}(\zeta)
 =
 \begin{pmatrix}
   \lambda(\zeta)&t^\dagger(\zeta)\\
   t(\zeta)&\ell(\zeta)
 \end{pmatrix}.
 \label{eq:V104-reservoir-matrix}
\end{equation}
When (ell) is invertible, block elimination gives
\begin{equation}
 \det\mathcal B_{\rm res}
 =\det\ell\,
   \bigl(\lambda-t^\dagger\ell^{-1}t\bigr).
 \label{eq:V104-Schur-determinant}
\end{equation}
Thus reciprocity alone does not force the old factor (lambda): a nonzero
Schur correction could algebraically lift it.  BRST incidence is the decisive
additional condition.

\subsection{The free normal ghost jet}
\label{subsec:V104-free-ghost-jet}

At linearized principal order, write
\begin{equation}
 \mathsf s a_\mu=D_\mu c,
 \qquad \mathsf s c=0,
 \qquad c|_{\mathcal T}=0.
 \label{eq:V104-linear-BRST}
\end{equation}
The V34 identity is
(mathsf s\Gamma=\mathcal M_{\rm YM}c).  The relative ghost boundary
condition kills (c) and all its tangential derivatives at the boundary, but
it does not kill its first normal jet.

\begin{lemma}[Dirichlet data leave the first normal ghost jet free]
\label{lem:V104-free-normal-ghost}
For the frozen principal second-order ghost equation, the boundary jet
conditions
\begin{equation}
 c|_{\mathcal T}=0,
 \qquad
 (\mathcal M_{\rm YM}c)|_{\mathcal T}=0
 \label{eq:V104-ghost-jet-conditions}
\end{equation}
allow (\rho:=(D_nc)|_{\mathcal T}) to be prescribed arbitrarily at a
boundary point.
\end{lemma}

\begin{proof}
In boundary normal coordinates, the frozen principal ghost operator has a
nonzero (D_n^2) coefficient and contains no term which determines (D_nc)
from the Dirichlet trace.  Prescribe the first jet by
(c=0), (D_ac=0), and (D_nc=\rho).  Equation
\eqref{eq:V104-ghost-jet-conditions} determines the second normal derivative
(D_n^2c); all remaining second tangential and mixed jets may be set to zero.
This produces a formal solution jet for every (\rho).  Equivalently, the
polynomial (x_n\rho), followed by the required quadratic normal correction,
has the stated principal boundary jet.
\end{proof}

At a fixed nonzero real tangential covector, the most general linear
principal response of the bosonic reservoir to the ghost jet which survives
Dirichlet restriction can therefore be written
\begin{equation}
 \mathsf s\eta=r\rho.
 \label{eq:V104-reservoir-BRST}
\end{equation}
Terms proportional to (c) or its tangential derivatives vanish on the
relative ghost domain.  The symbol (r) may be matrix-valued and
frequency-dependent; locality and homogeneity make it bounded on each compact
normalized chart.

\begin{proposition}[Exact BRST incidence identities]
\label{prop:V104-BRST-identities}
If both boundary Euler rows
\eqref{eq:V104-first-row}--\eqref{eq:V104-second-row} are preserved by
\eqref{eq:V104-linear-BRST}--\eqref{eq:V104-reservoir-BRST} on the ghost
Euler equation, then their principal symbols obey
\begin{align}
 t^\dagger r&=0,
 \label{eq:V104-BRST-first}\\
 \ell r+t&=0.
 \label{eq:V104-BRST-second}
\end{align}
\end{proposition}

\begin{proof}
On the ghost Euler equation,
(mathsf s\Gamma=\mathcal M_{\rm YM}c=0).  Applying (mathsf s) to
(F_1) therefore leaves (t^\dagger r\rho).  Lemma
\ref{lem:V104-free-normal-ghost} makes (\rho) arbitrary, proving
\eqref{eq:V104-BRST-first}.  Applying (mathsf s) to (F_2) gives
((\ell r+t)\rho), because (mathsf s a=D_nc=\rho).  The same freedom
proves \eqref{eq:V104-BRST-second}.
\end{proof}

\subsection{The protected glancing vector}
\label{subsec:V104-protected-vector}

\begin{theorem}[Reciprocal-reservoir BRST obstruction]
\label{thm:V104-reservoir-obstruction}
Consider any finite-dimensional reservoir of the form
\eqref{eq:V104-boundary-action}.  Keep the V35 relative ghost domain and
suppose that both coupled Euler rows are BRST-invariant at principal order.
Then
\begin{equation}
 \mathcal B_{\rm res}(\zeta)
 \binom{1}{r(\zeta)}
 =
 \binom{\lambda(\zeta)}{0}.
 \label{eq:V104-protected-vector}
\end{equation}
Consequently the normalized V35 glancing sequence satisfies
\begin{equation}
 \inf_{\zeta\in\overline{\mathbb S}_+}
 s_{\min}\mathcal B_{\rm res}(\zeta)=0.
 \label{eq:V104-reservoir-floor-zero}
\end{equation}
The conclusion does not require (ell) to be invertible or positive.
\end{theorem}

\begin{proof}
Multiply \eqref{eq:V104-reservoir-matrix} by the displayed vector.  Its
upper component is
(lambda+t^\dagger r=\lambda) by
\eqref{eq:V104-BRST-first}; its lower component is
(t+\ell r=0) by \eqref{eq:V104-BRST-second}.  This proves
\eqref{eq:V104-protected-vector}.

Evaluate it on the normalized frequencies
((\widetilde s_\epsilon,\widetilde k_\epsilon)) of Proposition
\ref{prop:V35-YM-glancing}.  On a compact normalized chart, (r) is bounded,
so
\[
 v_\epsilon
 :=
 \frac{(1,r(\widetilde\zeta_\epsilon))}
      {\sqrt{1+\|r(\widetilde\zeta_\epsilon)\|^2}}
\]
is a unit stable amplitude and
\[
 \|\mathcal B_{\rm res}(\widetilde\zeta_\epsilon)v_\epsilon\|
 =
 \frac{|\widetilde\lambda_\epsilon|}
      {\sqrt{1+\|r(\widetilde\zeta_\epsilon)\|^2}}
 \le |\widetilde\lambda_\epsilon|
 \longrightarrow0.
\]
This proves \eqref{eq:V104-reservoir-floor-zero}.  No inverse or sign of
(ell) was used.
\end{proof}

\begin{corollary}[A passive reservoir decouples]
\label{cor:V104-passive-decouples}
If the reservoir is BRST-inert, (r=0), then BRST invariance of its Euler row
forces (t=0).  The matrix \eqref{eq:V104-reservoir-matrix} is block diagonal
and the original normal gauge vector remains the V35 quasimode.
\end{corollary}

\begin{proof}
Equation \eqref{eq:V104-BRST-second} becomes (t=0).  The assertion then
follows directly from \eqref{eq:V104-reservoir-matrix}.
\end{proof}

\begin{corollary}[Schur cancellation is forced when the reservoir is
invertible]
\label{cor:V104-Schur-cancellation}
If (ell) is invertible, the BRST identities imply
\begin{equation}
 t^\dagger\ell^{-1}t=0,
 \qquad
 \det\mathcal B_{\rm res}=\lambda\det\ell.
 \label{eq:V104-Schur-cancellation}
\end{equation}
Thus the algebraic correction in \eqref{eq:V104-Schur-determinant} cancels
exactly on every BRST-compatible member of this class.
\end{corollary}

\begin{proof}
Equation \eqref{eq:V104-BRST-second} gives
(r=-\ell^{-1}t).  Substitute this into
\eqref{eq:V104-BRST-first}, then use
\eqref{eq:V104-Schur-determinant}.
\end{proof}

Tensoring with the twelve-dimensional Standard Model adjoint fibre gives the
same near-kernel independently in each gauge direction.  Embedding one such
direction and setting all metric, continuum, Higgs, and fermion amplitudes to
zero reproduces the V103 common-domain obstruction.

\begin{firewall}[Scope of the reservoir obstruction]
\label{fire:V104-scope}
Theorem~\ref{thm:V104-reservoir-obstruction} applies to a finite tangential
reservoir whose quadratic action has the reciprocal form
\eqref{eq:V104-boundary-action}, whose BRST variation is exhausted by the
free normal jet of the unchanged Dirichlet ghost, and whose two displayed
Euler rows must each be preserved.  It does not cover a boundary BRST complex
with additional ghost, antighost, multiplier, or conormal equations; a change
of the bulk ghost boundary condition; a nonlocal Calder\'on graph; or a
weighted anisotropic calculus.  Any claimed escape must print the new Green
form and every new BRST boundary row, rather than merely changing the lower
right block (ell).
\end{firewall}

\subsection{V104 ledger and the forced V105 audit}
\label{subsec:V104-ledger}

\paragraph{New exact conclusions.}
\begin{enumerate}[label=\textup{(\arabic*)},leftmargin=2.8em]
\item Proposition~\ref{prop:V104-extended-Green} derives the extended Green
      identity of the reciprocal reservoir directly from its boundary action.
\item Equation \eqref{eq:V104-reservoir-matrix} is the exact enlarged stable
      matrix; \eqref{eq:V104-Schur-determinant} identifies the apparent
      algebraic glancing repair.
\item Lemma~\ref{lem:V104-free-normal-ghost} proves that the Dirichlet ghost
      leaves (D_nc) free at the boundary principal jet.
\item Proposition~\ref{prop:V104-BRST-identities} derives the two independent
      BRST incidence identities.
\item Theorem~\ref{thm:V104-reservoir-obstruction} combines those identities
      into the explicit protected vector \eqref{eq:V104-protected-vector} and
      proves the zero unweighted floor.
\item Corollary~\ref{cor:V104-Schur-cancellation} shows that the reciprocal
      Schur correction vanishes identically when the reservoir block is
      invertible.
\end{enumerate}

\paragraph{Next forced step.}
V105 is a five-version checkpoint.  It must first inspect the newest shared
Yang--Mills and Navier--Stokes notes and record their exact versions and
hashes.  It should then test the smallest genuine boundary BRST complex which
changes one hypothesis of Theorem~\ref{thm:V104-reservoir-obstruction}.  The
candidate must include its ghost and antighost Green form, multiplier rows,
nilpotency/incidence identities, stress contribution, and complete stable
matrix.  If those data are unavailable or circular, the programme should go
upstream to the choice of ghost boundary domain rather than attach another
passive boson.

\begin{assessment}[V104 continuum status]
\label{ass:V104-continuum-status}
The full Standard--Model--Einstein continuum remains unproved.  V104 rules
out the entire finite reciprocal-reservoir class above as a repair of the
V35 relative-domain glancing floor.  A changed boundary BRST complex or ghost
domain, the fermion branch selection, a common quantum realization and its
counterterms, curved radiative equivariance, nonlinear existence, and the
cofinal interacting limit remain open.
\end{assessment}

\section*{V104 exact checks and append-only record}
\addcontentsline{toc}{section}{V104 exact checks and append-only record}

The decisive symbolic identities are
\begin{equation}
 \boxed{
 t^\dagger r=0,
 \qquad
 \ell r+t=0,
 \qquad
 \begin{pmatrix}\lambda&t^\dagger\\t&\ell\end{pmatrix}
 \binom{1}{r}
 =\binom{\lambda}{0}.}
 \label{eq:V104-final-check}
\end{equation}
They are exact consequences of variation and BRST incidence, not numerical
singular-value samples.  The append operation preserves the complete V103
prefix byte for byte before its former live terminator.  The assembled V104
file is checked for one live final terminator, balanced environments and
braces, duplicate labels, newly introduced unresolved references, display
math delimiters, and control characters before the exact V103 predecessor is
removed.

% =============================================================================
% END APPENDED VERSION V104 MATERIAL
% =============================================================================


\clearpage
% =============================================================================
% BEGIN APPENDED VERSION V105 MATERIAL
% =============================================================================
\section{V105: companion audit, a BRST--Sommerfeld glancing repair, and its
exact action obstruction}
\label{sec:V105-Sommerfeld-audit}

V104 proved that reciprocal finite reservoirs cannot repair the V35 relative
branch while its Dirichlet ghost domain is left unchanged.  This fifth
iteration performs the required read-only companion audit and changes the
upstream datum which generated that obstruction: the ghost boundary graph.
At frozen principal order, a common Sommerfeld row for the gauge and ghost
fields is BRST-invariant and has the sharp V35 uniform stable-root floor.  The
same row is not isotropic for the quadratic wave Green form.  Thus the
complementing obstruction is removed, but the common-action gate is not.

This distinction implements both lessons of the companion audit.  The exact
signed boundary operator is kept before taking a singular-value norm, and the
positive floor is tested independently of the signed Green pairing.  No
companion theorem is imported as a Standard Model proof; all claims below are
proved in the frozen gauge/ghost system.

\subsection{Mandatory V105 companion audit}
\label{subsec:V105-audit}

The read-only snapshots used for this iteration are
\begin{center}
\begin{tabular}{|p{0.48\textwidth}|r|p{0.31\textwidth}|}
\hline
Companion file & Bytes & SHA--256\\
\hline
\texttt{Renewal\_NS\_Stability\_Research\_Notes\_V31.tex}
& (887537) &
\texttt{F1121B62238AD5491BB7CF03635EA9934942EDF573ED8FAAA9EE79E1D38A6696}\\
\hline
\texttt{Renewal\_Yang\_Mills\_Mass\_Gap\_Research\_Notes\_V99.tex}
& (2360927) &
\texttt{C40128565676B32DEFCCB8BC13C7186CA99C91D49FEE6F36B7D107D9A5A92D71}\\
\hline
\end{tabular}
\end{center}

The Navier--Stokes note is byte-identical to the V100 snapshot.  Its terminal
V31 balanced-helicity example has zero signed hinge quantities while its
positive total critical coarea remains nonzero.  The Yang--Mills programme has
advanced from V97 to V99.  V98 reconstructs the signed physical collision
operator before its absolute carrier and proves that a conditional Racah POVM
cannot be inserted as a physical joint pinching.  V99 proves fixed-label
Cartan visibility results but also proves that their positive floor is not
uniform as representation labels grow.

\begin{assessment}[Transfer used in V105]
\label{ass:V105-audit-transfer}
The shared structural lesson is an order of operations.  Preserve the exact
signed operator and its incidence identities; then compute the positive norm
or singular floor on that same operator.  A signed cancellation is not a
positive estimate, and a positive fixed-card estimate is not a uniform
continuum estimate.  V105 applies this order directly to the frozen
Yang--Mills boundary symbol.  It does not use the companion collision or fluid
identities as estimates for the Standard Model system.
\end{assessment}

\subsection{A changed gauge and ghost boundary graph}
\label{subsec:V105-Sommerfeld-graph}

Work at a flat timelike boundary and use the V35 normalized stable variables
\((s,k,\lambda)\), with
\begin{equation}
 \lambda^2=s^2+|k|^2,
 \qquad \operatorname{Re}\lambda>0
 \quad\hbox{for }\operatorname{Re}s>0.
 \label{eq:V105-stable-root}
\end{equation}
Let
\begin{equation}
 B_+:=D_T+D_n.
 \label{eq:V105-Sommerfeld-operator}
\end{equation}
For every Standard Model adjoint component impose the principal rows
\begin{equation}
 B_+a_\mu=0,
 \qquad
 B_+c=0,
 \qquad
 B_+\bar c=0,
 \qquad
 B_+b=0.
 \label{eq:V105-BRST-Sommerfeld-domain}
\end{equation}
The last row is a compatibility row when the Nakanishi--Lautrup field is kept
algebraic; it is not counted as an additional wave stable amplitude.

\begin{proposition}[Principal BRST incidence]
\label{prop:V105-BRST-incidence}
For the frozen linearized rules
\begin{equation}
 \mathsf s a_\mu=D_\mu c,
 \qquad
 \mathsf s c=0,
 \qquad
 \mathsf s\bar c=b,
 \qquad
 \mathsf s b=0,
 \label{eq:V105-linearized-BRST}
\end{equation}
the domain \eqref{eq:V105-BRST-Sommerfeld-domain} is a BRST subcomplex at
principal order.
\end{proposition}

\begin{proof}
Frozen principal derivatives commute.  Therefore
\[
 \mathsf s(B_+a_\mu)=B_+D_\mu c=D_\mu(B_+c)=0.
\]
The ghost row is invariant because (mathsf sc=0).  Moreover,
\[
 \mathsf s(B_+\bar c)=B_+b=0,
 \qquad
 \mathsf s(B_+b)=0.
\]
This proves closure of every displayed row.
\end{proof}

\begin{theorem}[Exact uniform stable-root floor]
\label{thm:V105-Sommerfeld-floor}
On a decaying stable amplitude, every differential row in
\eqref{eq:V105-BRST-Sommerfeld-domain} has scalar symbol (s+\lambda).
Consequently, on the normalized closed stable hemisphere,
\begin{equation}
 \boxed{
 s_{\min}\mathcal B_{\rm Som}(s,k)
 =|s+\lambda|
 \ge \max\{|s|,|k|\}
 \ge2^{-1/2}.}
 \label{eq:V105-Sommerfeld-floor}
\end{equation}
The bound holds through both glancing sheets and is unchanged after tensoring
with the twelve-dimensional Standard Model adjoint fibre.
\end{theorem}

\begin{proof}
Substitution of the stable normal exponential turns (B_+) into multiplication
by (s+\lambda).  The gauge, ghost, and antighost blocks are diagonal copies
of that scalar.  Hence every singular value is (|s+\lambda|).  The sharp
inequality is Theorem~\ref{thm:V35-Sommerfeld-factor}, whose proof used only
\eqref{eq:V105-stable-root}; tensoring with an identity changes multiplicity
but not singular values.
\end{proof}

\begin{corollary}[The V35 relative quasimode is removed]
\label{cor:V105-relative-witness-removed}
For the normal gauge amplitude (a=n^\flat X) used in
Proposition~\ref{prop:V35-YM-glancing}, the new boundary image has norm
\begin{equation}
 \|\mathcal B_{\rm Som}a\|=|s+\lambda|\|X\|
 \ge2^{-1/2}\|X\|.
 \label{eq:V105-old-witness-lifted}
\end{equation}
Thus that explicit glancing sequence is not a near-kernel of the changed
domain.
\end{corollary}

\begin{proof}
The normal component now carries the same Sommerfeld row as every other gauge
component.  Apply \eqref{eq:V105-Sommerfeld-floor}.
\end{proof}

\subsection{The signed Green test}
\label{subsec:V105-action-test}

The positive calculation above does not decide whether the boundary graph is
the stationary domain of the common quadratic action.  For one real wave
component let (q=a|_{\mathcal T}) and (p=D_na|_{\mathcal T}).  The
Sommerfeld row is the graph
\begin{equation}
 \mathcal L_{\rm Som}
 :=\{(q,p):p=-D_Tq\}.
 \label{eq:V105-Sommerfeld-phase-graph}
\end{equation}
Use the V35 wave Green form integrated over a boundary cylinder,
\begin{equation}
 \Omega((q,p),(v,r))
 :=\int_{\mathcal T}
 \bigl(\langle p,v\rangle-\langle q,r\rangle\bigr).
 \label{eq:V105-integrated-Green-form}
\end{equation}

\begin{theorem}[Sommerfeld action obstruction]
\label{thm:V105-Sommerfeld-action-obstruction}
The graph \eqref{eq:V105-Sommerfeld-phase-graph} is not isotropic for
\eqref{eq:V105-integrated-Green-form}.  In particular, it is not a Lagrangian
boundary domain for the closed quadratic wave action, and no real quadratic
boundary density depending only on (q) can produce its principal conormal
relation.
\end{theorem}

\begin{proof}
Take real compactly supported time profiles (f,g) and set
\((q,p)=(f,-D_Tf)\), \((v,r)=(g,-D_Tg)\).  Integration by parts in time gives
\begin{align}
 \Omega((q,p),(v,r))
 &=\int_{\mathcal T}
   \bigl(-\langle D_Tf,g\rangle+\langle f,D_Tg\rangle\bigr)\notag\\
 &=2\int_{\mathcal T}\langle f,D_Tg\rangle.
 \label{eq:V105-Sommerfeld-nonisotropy}
\end{align}
Choose (f,g) so that the last integral is nonzero.  Hence the graph is not
isotropic.

For the final assertion, the Hessian of a real scalar boundary functional is
formally self-adjoint.  The operator (-D_T) in
\eqref{eq:V105-Sommerfeld-phase-graph} is formally skew-adjoint on compactly
supported profiles.  Its nonzero principal part therefore cannot be that
Hessian.  Equivalently, the would-be density
(\tfrac12\langle q,-D_Tq\rangle) is a time divergence and has zero interior
Euler derivative.
\end{proof}

\begin{corollary}[Direct sums do not cure the action defect]
\label{cor:V105-direct-sum-action-defect}
The complete gauge and ghost direct sum contains the one-component witness
\eqref{eq:V105-Sommerfeld-nonisotropy}.  Therefore adjoint multiplicity,
positive metric--continuum floors, and conditional fermion/Higgs floors do not
turn \eqref{eq:V105-BRST-Sommerfeld-domain} into a common-action Lagrangian.
\end{corollary}

\begin{proof}
Support both tangent vectors in one real gauge component and set every other
component to zero.  The full Green form restricts to
\eqref{eq:V105-integrated-Green-form} on that subspace.
\end{proof}

\begin{firewall}[Complementing success is not common-action success]
\label{fire:V105-two-gates}
Theorem~\ref{thm:V105-Sommerfeld-floor} is a genuine uniform unweighted
stable-root estimate for the changed frozen domain.  Theorem
\ref{thm:V105-Sommerfeld-action-obstruction} is a genuine failure of action
incidence for the same signed boundary operator.  Neither statement cancels
the other.  Calling the branch acceptable because it is maximally dissipative
would change the programme's common closed-action requirement and must be
declared explicitly.
\end{firewall}

\subsection{V105 decision and next acquisition}
\label{subsec:V105-decision}

\begin{proposition}[Exact two-gate outcome]
\label{prop:V105-two-gate-outcome}
For the frozen BRST--Sommerfeld branch,
\begin{equation}
 \boxed{
 \rho_{\rm stable}\ge2^{-1/2},
 \qquad
 \mathcal L_{\rm Som}\text{ is not action-isotropic}.}
 \label{eq:V105-two-gate-outcome}
\end{equation}
Hence changing the ghost domain removes the V35 unweighted glancing witness
but does not yet provide the V102 common operator/domain realization.
\end{proposition}

\begin{proof}
Combine Theorems~\ref{thm:V105-Sommerfeld-floor} and
\ref{thm:V105-Sommerfeld-action-obstruction}.
\end{proof}

\paragraph{Next forced step.}
V106 should construct the smallest conservative dilation of the Sommerfeld
port, with actual boundary variables or a second propagating side carrying
the missing Green flux.  It must print the enlarged action, BRST rules, ghost
and antighost rows, and full stable matrix.  The key test is simultaneous:
the dilation must make the total Green graph Lagrangian without reintroducing
a (lambda)-near-kernel.  If a finite tangential dilation collapses to the
V104 protected vector, the programme should go further upstream to a
two-sided or boundaryless gauge realization.

\begin{assessment}[V105 continuum status]
\label{ass:V105-continuum-status}
The full Standard--Model--Einstein continuum remains unproved.  V105 supplies
an exact BRST-compatible glancing repair for the positive frozen gauge/ghost
stable matrix, and proves that this repaired domain is not the Lagrangian
domain of the closed quadratic wave action.  A conservative action dilation,
the fermion branch selection, a common quantum realization and counterterms,
curved radiative equivariance, nonlinear existence, and the cofinal
interacting limit remain open.
\end{assessment}

\section*{V105 exact checks and append-only record}
\addcontentsline{toc}{section}{V105 exact checks and append-only record}

The exact checkpoint identities are
\begin{equation}
 \boxed{
 \mathsf s(B_+a_\mu)=D_\mu(B_+c),
 \qquad
 s_{\min}\mathcal B_{\rm Som}=|s+\lambda|\ge2^{-1/2},
 \qquad
 \Omega|_{\mathcal L_{\rm Som}}
 =2\int_{\mathcal T}\langle f,D_Tg\rangle\not\equiv0.}
 \label{eq:V105-final-check}
\end{equation}
The append operation preserves the complete V104 prefix byte for byte before
its former live terminator.  The assembled V105 file is checked for one live
final terminator, balanced environments and braces, duplicate labels, newly
introduced unresolved references, display math delimiters, and control
characters before the exact V104 predecessor is removed.  The two audited
companion files were read only.

% =============================================================================
% END APPENDED VERSION V105 MATERIAL
% =============================================================================

\clearpage
% =============================================================================
% BEGIN APPENDED VERSION V106 MATERIAL
% =============================================================================
\section{V106: the minimal conservative Sommerfeld dilation and its shifted
boundary-characteristic kernel}
\label{sec:V106-conservative-dilation}

V105 changed the ghost boundary graph and obtained a sharp positive
Sommerfeld stable-root floor, but the corresponding one-field graph was not
action-isotropic.  The smallest real local action capable of carrying the
missing skew time flux uses a second wave collar.  This note constructs that
doubled action, proves its Lagrangian and BRST incidence, and diagonalizes its
complete stable matrix.

The conservative dilation passes the two algebraic gates which V105 could not
pass simultaneously: it is an actual stationary action domain and a BRST
subcomplex.  Nevertheless its stable matrix has an exact kernel on a shifted
boundary-characteristic sheet.  The construction therefore moves the V35
glancing loss rather than removing the unweighted complementing obstruction.

\subsection{Doubled real wave action}
\label{subsec:V106-doubled-action}

For one real adjoint component let (a) be the original wave field and let
(z) be a second propagating collar field.  Put
\begin{equation}
 q:=\binom{a}{z},
 \qquad
 p:=\binom{D_na}{D_nz},
 \qquad
 J:=\begin{pmatrix}0&-1\\1&0\end{pmatrix},
 \qquad J^T=-J,quad J^2=-I_2.
 \label{eq:V106-doubled-variables}
\end{equation}
The normal coordinates on the two collars are chosen inward from their common
boundary, so their quadratic Green forms add in the displayed (p)-frame.
For a real parameter (alpha), add the boundary density
\begin{equation}
 S_{\partial,\alpha}[q]
 :=
 \frac{\alpha}{2}
 \int_{\mathcal T}q^TJ D_Tq.
 \label{eq:V106-symplectic-boundary-action}
\end{equation}
This density is real and is not a total time derivative when both components
are independent.

\begin{lemma}[Exact boundary Euler derivative]
\label{lem:V106-boundary-variation}
For compactly supported variations on the time faces,
\begin{equation}
 \delta S_{\partial,\alpha}
 =\alpha\int_{\mathcal T}(\delta q)^TJ D_Tq.
 \label{eq:V106-boundary-variation}
\end{equation}
Consequently the combined bulk--boundary stationary row is
\begin{equation}
 C_\alpha(q,p):=p+\alpha JD_Tq=0.
 \label{eq:V106-conservative-row}
\end{equation}
\end{lemma}

\begin{proof}
Vary \eqref{eq:V106-symplectic-boundary-action} and integrate the term
(q^TJ D_T\delta q) by parts in time.  Since (J^T=-J), the integrated term
equals the first varied term rather than cancelling it.  Their sum is
\eqref{eq:V106-boundary-variation}.  Adding the bulk conormal variation
(\int p^T\delta q) gives \eqref{eq:V106-conservative-row}, up to the fixed
common conormal sign convention already used in V35--V105.
\end{proof}

\begin{theorem}[The doubled graph is Lagrangian]
\label{thm:V106-Lagrangian}
Let
\begin{equation}
 \mathcal L_\alpha
 :=\{(q,p):p=-\alpha JD_Tq\}.
 \label{eq:V106-Lagrangian-graph}
\end{equation}
Then (mathcal L_\alpha) is Lagrangian for the integrated doubled Green form
\begin{equation}
 \Omega_2((q,p),(v,r))
 :=\int_{\mathcal T}(p^Tv-q^Tr).
 \label{eq:V106-doubled-Green}
\end{equation}
\end{theorem}

\begin{proof}
The operator (JD_T) is formally self-adjoint because both (J) and
(D_T) are formally skew-adjoint and commute.  A graph (p=Rq) is
Lagrangian for \eqref{eq:V106-doubled-Green} when (R=R^\dagger).
For completeness, substitute (p=-\alpha JD_Tq) and
(r=-\alpha JD_Tv):
\[
 \Omega_2
 =-\alpha\int (JD_Tq)^Tv
   +\alpha\int q^TJD_Tv=0,
\]
where one time integration by parts and (J^T=-J) cancel the two terms.
The graph has half the dimension of the doubled boundary phase space, so
isotropy is maximal.
\end{proof}

\subsection{Doubled BRST complex}
\label{subsec:V106-BRST}

Let (gamma=(c,d)^T) be the ghosts of the original and reservoir gauge
fields, with corresponding doubled antighost and multiplier columns
(\bar\gamma) and (b_\gamma).  At frozen linearized order set
\begin{equation}
 \mathsf s q_\mu=D_\mu\gamma,
 \qquad
 \mathsf s\gamma=0,
 \qquad
 \mathsf s\bar\gamma=b_\gamma,
 \qquad
 \mathsf s b_\gamma=0.
 \label{eq:V106-doubled-BRST}
\end{equation}
Impose the same conservative principal row on every member of the quartet:
\begin{equation}
 C_\alpha(q_\mu,D_nq_\mu)=0,
 \quad
 C_\alpha(\gamma,D_n\gamma)=0,
 \quad
 C_\alpha(\bar\gamma,D_n\bar\gamma)=0,
 \quad
 C_\alpha(b_\gamma,D_nb_\gamma)=0.
 \label{eq:V106-doubled-BRST-domain}
\end{equation}
As before, the multiplier row is a prolongation compatibility row when the
multiplier is eliminated algebraically.

\begin{proposition}[BRST incidence of the conservative graph]
\label{prop:V106-BRST-incidence}
The domain \eqref{eq:V106-doubled-BRST-domain} is preserved by
\eqref{eq:V106-doubled-BRST} at frozen principal order.
\end{proposition}

\begin{proof}
The constant matrix (J) commutes with spacetime derivatives.  Hence
\[
 \mathsf s C_\alpha(q_\mu,D_nq_\mu)
 =D_\mu C_\alpha(\gamma,D_n\gamma)=0.
\]
The remaining three rows form the chain
(C_\alpha(\gamma)\mapsto0),
(C_\alpha(\bar\gamma)\mapsto C_\alpha(b_\gamma)), and
(C_\alpha(b_\gamma)\mapsto0).
\end{proof}

\subsection{Exact stable spectrum}
\label{subsec:V106-stable-spectrum}

On the V35 stable root, \eqref{eq:V106-conservative-row} has the two-by-two
matrix
\begin{equation}
 \mathcal B_\alpha(s,k)
 :=\lambda I_2+\alpha sJ.
 \label{eq:V106-stable-matrix}
\end{equation}

\begin{theorem}[Exact singular values and shifted kernel]
\label{thm:V106-shifted-kernel}
The matrix \eqref{eq:V106-stable-matrix} is normal and has singular values
\begin{equation}
 \sigma_\pm(s,k)=|\lambda\pm i\alpha s|.
 \label{eq:V106-exact-singular-values}
\end{equation}
For every (alpha>0), define
\begin{equation}
 \tau_\alpha:=(2+\alpha^2)^{-1/2},
 \qquad
 s_\alpha=i\tau_\alpha,
 \qquad
 |k_\alpha|=\sqrt{1+\alpha^2}\,\tau_\alpha.
 \label{eq:V106-kernel-frequency}
\end{equation}
Then (|s_\alpha|^2+|k_\alpha|^2=1), the stable branch has
(lambda_\alpha=\alpha\tau_\alpha>0), and
\begin{equation}
 \sigma_+(s_\alpha,k_\alpha)
 =|\lambda_\alpha+i\alpha s_\alpha|=0.
 \label{eq:V106-exact-kernel}
\end{equation}
For (alpha=0), the matrix is (lambda I_2) and has the original glancing
zero.  Therefore, for every real (alpha),
\begin{equation}
 \boxed{
 \inf_{\overline{\mathbb S}_+}
 s_{\min}\mathcal B_\alpha=0.}
 \label{eq:V106-zero-floor}
\end{equation}
\end{theorem}

\begin{proof}
The real skew matrix (J) is unitarily diagonalizable over (mathbb C),
with eigenvalues (+i) and (-i).  Since
(mathcal B_\alpha) is a polynomial in (J), it is normal and its
eigenvalues are (lambda\pm i\alpha s).  This proves
\eqref{eq:V106-exact-singular-values}.

At \eqref{eq:V106-kernel-frequency}, normalization is immediate and
\[
 \lambda_\alpha^2
 =s_\alpha^2+|k_\alpha|^2
 =-\tau_\alpha^2+(1+\alpha^2)\tau_\alpha^2
 =\alpha^2\tau_\alpha^2.
\]
The stable branch therefore has
(lambda_\alpha=\alpha\tau_\alpha).  Since
(i\alpha s_\alpha=-\alpha\tau_\alpha),
\eqref{eq:V106-exact-kernel} follows.  The (alpha=0) assertion is direct.
Taking the minimum over the normalized closed hemisphere proves
\eqref{eq:V106-zero-floor}.
\end{proof}

\begin{corollary}[Explicit stable vector]
\label{cor:V106-explicit-vector}
Let (e_+\in\mathbb C^2) be a unit vector with (Je_+=ie_+).  Then
\begin{equation}
 \mathcal B_\alpha(s_\alpha,k_\alpha)e_+=0.
 \label{eq:V106-explicit-vector}
\end{equation}
Approaching \((s_\alpha,k_\alpha)\) from
(\operatorname{Re}s>0) and normalizing the frequency gives a stable
quasimode with image norm tending to zero.
\end{corollary}

\begin{proof}
On the (+i) eigenspace of (J), the matrix acts by
(lambda+i\alpha s), which vanishes by
\eqref{eq:V106-exact-kernel}.  Continuity of the chosen stable square root
from the open right half-plane gives the quasimode statement.
\end{proof}

\begin{proposition}[No exponentially growing kernel]
\label{prop:V106-no-growing-kernel}
For real (k) and (alpha\ge0), a zero of
(lambda\pm i\alpha s) cannot have (operatorname{Re}s>0).
Thus Theorem~\ref{thm:V106-shifted-kernel} proves loss of a uniform
complementing floor, not an exponentially growing mode.
\end{proposition}

\begin{proof}
A zero implies, after squaring and using
(lambda^2=s^2+|k|^2),
\begin{equation}
 |k|^2=-(1+\alpha^2)s^2.
 \label{eq:V106-zero-locus}
\end{equation}
The left side is nonnegative real.  Hence (s^2) is nonpositive real.
This forces (s) to be purely imaginary or zero, and excludes
(operatorname{Re}s>0).
\end{proof}

Tensoring the doubled construction with all four gauge-potential components
and the Standard Model adjoint fibre only repeats the two singular values.
It cannot lift the kernel \eqref{eq:V106-explicit-vector}.

\begin{firewall}[What the doubled collar proves]
\label{fire:V106-scope}
The doubled collar is an actual local quadratic action model, not the formal
auxiliary copy excluded by the V35 dilation firewall.  It proves that the
minimal real symplectic conservative completion has a different
boundary-characteristic zero even though it is Lagrangian and BRST-invariant.
It does not rule out a higher-dimensional dispersive boundary action, a
nonlocal Calder\'on coupling, a weighted calculus, or a boundaryless global
gauge realization.  Those alternatives require their own complete stable
matrices.
\end{firewall}

\subsection{V106 three-branch ledger and next acquisition}
\label{subsec:V106-ledger}

The tested local branches now have the following exact status:
\begin{center}
\begin{tabular}{|l|c|c|c|}
\hline
Branch & BRST incidence & Action Lagrangian & Uniform stable floor\\
\hline
V35 relative & yes & yes & no\\
V105 Sommerfeld & yes & no & yes\\
V106 doubled symplectic & yes & yes & no\\
\hline
\end{tabular}
\end{center}
The table records proved properties of these three declared models; it is not
a no-go theorem for all boundary realizations.

\paragraph{Next forced step.}
V107 should decide between two upstream routes.  One route is to classify
the principal symbol of a general finite real self-adjoint boundary system
and test whether its determinant must intersect the stable characteristic
variety.  The other is to write the exact nonlocal Dirichlet-to-Neumann or
Calder\'on graph, including BRST and stress incidence, and determine whether
the V102 quantum calculus can be formulated for that graph.  Another ad hoc
two-by-two port rotation would repeat the family already exhausted by
Theorem~\ref{thm:V106-shifted-kernel}.

\begin{assessment}[V106 continuum status]
\label{ass:V106-continuum-status}
The full Standard--Model--Einstein continuum remains unproved.  V106 gives a
local conservative action dilation and proves its BRST incidence, but an
exact shifted boundary-characteristic kernel keeps its unweighted stable
floor at zero.  A viable common gauge realization, the fermion branch
selection, quantum coefficients and common form bounds, curved radiative
equivariance, nonlinear existence, and the cofinal interacting limit remain
open.
\end{assessment}

\section*{V106 exact checks and append-only record}
\addcontentsline{toc}{section}{V106 exact checks and append-only record}

The decisive exact identities are
\begin{equation}
 \boxed{
 (JD_T)^\dagger=JD_T,
 \qquad
 \sigma_\pm=|\lambda\pm i\alpha s|,
 \qquad
 \lambda_\alpha+i\alpha s_\alpha=0.}
 \label{eq:V106-final-check}
\end{equation}
The append operation preserves the complete V105 prefix byte for byte before
its former live terminator.  The assembled V106 file is checked for one live
final terminator, balanced environments and braces, duplicate labels, newly
introduced unresolved references, display math delimiters, and control
characters before the exact V105 predecessor is removed.

% =============================================================================
% END APPENDED VERSION V106 MATERIAL
% =============================================================================

\clearpage
% =============================================================================
% BEGIN APPENDED VERSION V107 MATERIAL
% =============================================================================
\section{V107: classification of all finite real time-local conservative
collar dilations}
\label{sec:V107-time-local-classification}

V106 treated the smallest two-collar symplectic action and found that its
stable zero moves from the bulk glancing sheet to a shifted
boundary-characteristic sheet.  Adding further copies could help only if a
larger real time-local action has principal blocks not already represented by
that calculation.  V107 proves that it does not: the skew normal form of the
most general such action is an orthogonal sum of V106 two-by-two blocks and
uncoupled Neumann blocks.

This is a classification within a declared class.  The boundary principal
action is local, real, first order in the distinguished boundary time, and
contains finitely many collar fields; it contains no principal tangential
spatial derivative.  The result does not cover spatially dispersive boundary
systems or nonlocal Calder\'on operators.

\subsection{The complete time-local principal action}
\label{subsec:V107-general-action}

For one real adjoint and spacetime component, let
(q=(q_1,\ldots,q_N)^T) collect (N\ge1) propagating collar traces and let
(p=(D_nq_1,\ldots,D_nq_N)^T) be their common conormal-frame momenta.  Start
from an arbitrary real constant matrix (C\in M_N(\mathbb R)) and the
principal boundary density
\begin{equation}
 S_{\partial,C}^{(1)}[q]
 :=\frac12\int_{\mathcal T}q^TC D_Tq.
 \label{eq:V107-general-first-order-action}
\end{equation}
Set
\begin{equation}
 A:=\frac12(C-C^T),
 \qquad A^T=-A.
 \label{eq:V107-skew-part}
\end{equation}

\begin{lemma}[Only the skew coefficient enters the Euler row]
\label{lem:V107-skew-completeness}
For variations compactly supported away from the time faces,
\begin{equation}
 \delta S_{\partial,C}^{(1)}
 =\int_{\mathcal T}(\delta q)^TA D_Tq.
 \label{eq:V107-general-variation}
\end{equation}
The symmetric part of (C) is a pure time divergence.  Hence every real
quadratic, constant-coefficient, first-time-derivative boundary action has the
same principal Euler row as one and only one real skew matrix (A).
\end{lemma}

\begin{proof}
Vary \eqref{eq:V107-general-first-order-action} and integrate the term
(q^TC D_T\delta q) by parts:
\[
 \delta S_{\partial,C}^{(1)}
 =\frac12\int_{\mathcal T}(\delta q)^T(C-C^T)D_Tq,
\]
which is \eqref{eq:V107-general-variation}.  Directly,
(q^T(C+C^T)D_Tq=D_T\bigl(q^T(C+C^T)q/2\bigr)), proving the divergence
statement.  Uniqueness follows because the Euler coefficient is exactly
((C-C^T)/2).
\end{proof}

Adding the collar bulk conormal variation gives the boundary row
\begin{equation}
 C_A(q,p):=p+A D_Tq=0.
 \label{eq:V107-general-boundary-row}
\end{equation}
A symmetric zero-order boundary Hessian may also be added, but it has lower
homogeneous order and cannot alter the principal stable floor considered
below.

\begin{proposition}[Action and power identities]
\label{prop:V107-action-power}
The graph
\begin{equation}
 \mathcal L_A:=\{(q,p):p=-A D_Tq\}
 \label{eq:V107-action-graph}
\end{equation}
is Lagrangian for the integrated direct-sum wave Green form.  Its instantaneous
normal wave power is zero:
\begin{equation}
 \langle p,D_Tq\rangle
 =-\langle A D_Tq,D_Tq\rangle=0.
 \label{eq:V107-zero-power}
\end{equation}
\end{proposition}

\begin{proof}
Since (A^T=-A) and (D_T^\dagger=-D_T), the product (AD_T) is formally
self-adjoint.  The graph of a formally self-adjoint conormal operator is
Lagrangian by the V35 finite action criterion; equivalently, substitution in
the Green form makes its two terms cancel after one time integration by
parts.  Equation \eqref{eq:V107-zero-power} follows from skew-symmetry of
(A).
\end{proof}

\subsection{Orthogonal normal form and the stable matrix}
\label{subsec:V107-normal-form}

The real skew normal-form theorem supplies an orthogonal matrix
\(O\in O(N)\), positive numbers \(\alpha_1,\ldots,\alpha_m\), and
\(r=N-2m\) such that
\begin{equation}
 O^TAO
 =\operatorname{diag}
 \bigl(\alpha_1J,\ldots,\alpha_mJ,0_r\bigr),
 \qquad
 J=\begin{pmatrix}0&-1\\1&0\end{pmatrix}.
 \label{eq:V107-skew-normal-form}
\end{equation}
On the V35 stable root, \eqref{eq:V107-general-boundary-row} becomes
\begin{equation}
 \mathcal B_A(s,k):=\lambda I_N+sA.
 \label{eq:V107-general-stable-matrix}
\end{equation}

\begin{theorem}[Exact finite time-local spectrum]
\label{thm:V107-exact-spectrum}
The stable matrix \eqref{eq:V107-general-stable-matrix} is normal.  Its
singular values, counted with multiplicity, are
\begin{equation}
 \left\{
 |\lambda+i\alpha_js|,
 |\lambda-i\alpha_js|:1\le j\le m
 \right\}
 \cup
 \left\{|\lambda|\text{ repeated }r\text{ times}\right\}.
 \label{eq:V107-singular-multiset}
\end{equation}
Thus every coupled two-plane is exactly a V106 block and every residual real
line is an uncoupled Neumann block.
\end{theorem}

\begin{proof}
Conjugation by (O) in \eqref{eq:V107-skew-normal-form} is unitary after
complexification and preserves singular values.  It reduces
(mathcal B_A) to the orthogonal sum
\[
 \bigoplus_{j=1}^m(\lambda I_2+\alpha_jsJ)
 \oplus\lambda I_r.
\]
Each two-plane matrix is normal because it and its adjoint are polynomials in
the normal matrix (J).  Its unitary eigenvalues are
(lambda\pm i\alpha_js), while the residual block is scalar.  This proves
normality of the direct sum and \eqref{eq:V107-singular-multiset}.
\end{proof}

\begin{theorem}[No finite time-local conservative dilation has a uniform
stable floor]
\label{thm:V107-no-uniform-floor}
For every (N\ge1) and every real matrix (C) in
\eqref{eq:V107-general-first-order-action},
\begin{equation}
 \boxed{
 \inf_{(s,k)\in\overline{\mathbb S}_+}
 s_{\min}\mathcal B_A(s,k)=0.}
 \label{eq:V107-zero-floor}
\end{equation}
More precisely:
\begin{enumerate}[label=\textup{(\roman*)},leftmargin=3.3em]
\item if (r>0), a residual line has the exact V35 glancing factor
      (lambda);
\item for every (alpha_j>0), its two-plane has an exact kernel at
      \begin{equation}
       s_j=i(2+\alpha_j^2)^{-1/2},
       \qquad
       |k_j|=\sqrt{1+\alpha_j^2}\,(2+\alpha_j^2)^{-1/2},
       \qquad
       \lambda_j=\alpha_j(2+\alpha_j^2)^{-1/2}.
       \label{eq:V107-block-kernel-frequency}
      \end{equation}
\end{enumerate}
If (N) is odd, case \textup{(i)} is unavoidable.  If (A) is invertible,
case \textup{(ii)} occurs in every normal two-plane.
\end{theorem}

\begin{proof}
If (r>0), choose a unit vector in the zero block of
\eqref{eq:V107-skew-normal-form}.  At either normalized V35 glancing point,
(lambda=0), so its boundary image vanishes.

For a nonzero block, the values in
\eqref{eq:V107-block-kernel-frequency} obey
(|s_j|^2+|k_j|^2=1) and
(lambda_j^2=s_j^2+|k_j|^2).  The stable branch has
(lambda_j>0), and
\[
 \lambda_j+i\alpha_js_j=0.
\]
Theorem~\ref{thm:V107-exact-spectrum} therefore gives a zero singular value.
Approach each boundary frequency from (operatorname{Re}s>0) to obtain a
normalized stable quasimode if the complementing test is stated only on the
open hemisphere.  A real skew matrix in odd dimension has nontrivial kernel,
which proves the parity assertion.  Invertibility is equivalent to (r=0).
\end{proof}

\begin{corollary}[Adding finitely many identical collar copies cannot repair
V106]
\label{cor:V107-no-copy-repair}
No orthogonal mixing of finitely many real wave collars through a local
first-time-derivative quadratic boundary action can simultaneously retain the
closed action and produce a positive unweighted stable-root floor.
\end{corollary}

\begin{proof}
Lemma~\ref{lem:V107-skew-completeness} exhausts the declared principal action
class, and Theorem~\ref{thm:V107-no-uniform-floor} proves the zero floor for
every member of that class.
\end{proof}

\subsection{BRST incidence does not remove the classified kernel}
\label{subsec:V107-BRST}

For (N) gauge-collar copies let (gamma\in\mathbb R^N) be the corresponding
ghost column and use
\begin{equation}
 \mathsf s q_\mu=D_\mu\gamma,
 \qquad \mathsf s\gamma=0.
 \label{eq:V107-multicollar-BRST}
\end{equation}
Impose (C_A=0) on gauge, ghost, antighost, and multiplier columns as in
V106.

\begin{proposition}[Complete time-local BRST incidence]
\label{prop:V107-BRST-incidence}
Every graph \eqref{eq:V107-general-boundary-row} is a frozen principal BRST
subcomplex under the matched ghost rows, but its gauge and ghost stable
matrices both contain the singular values
\eqref{eq:V107-singular-multiset}.  BRST closure therefore preserves rather
than lifts the kernel in Theorem~\ref{thm:V107-no-uniform-floor}.
\end{proposition}

\begin{proof}
The constant matrix (A) commutes with spacetime derivatives, so
\[
 \mathsf s C_A(q_\mu,D_nq_\mu)
 =D_\mu C_A(\gamma,D_n\gamma).
\]
The matched ghost row makes the right side zero; the antighost and multiplier
rows close exactly as in Proposition~\ref{prop:V106-BRST-incidence}.  The
same operator (C_A) acts on every column, so its stable singular multiset is
unchanged.
\end{proof}

\begin{firewall}[Exact scope of the classification]
\label{fire:V107-scope}
Theorem~\ref{thm:V107-no-uniform-floor} is complete for finite real collar
traces whose principal boundary action is first order only in the selected
boundary time.  It does not classify terms
(A^jD_j) with tangential spatial derivatives, auxiliary fields with a
different bulk dispersion, pseudodifferential square roots, or nonlocal
Calder\'on projectors.  Lower-order masses cannot alter the homogeneous zero
floor, but a different principal dispersion must be computed rather than
declared ineffective.
\end{firewall}

\subsection{V107 ledger and next upstream row}
\label{subsec:V107-ledger}

\paragraph{New exact conclusions.}
\begin{enumerate}[label=\textup{(\arabic*)},leftmargin=2.8em]
\item Lemma~\ref{lem:V107-skew-completeness} proves that the skew part of one
      real matrix is the entire first-time-derivative Euler datum.
\item Proposition~\ref{prop:V107-action-power} proves both Lagrangian
      incidence and zero boundary power for every member of the class.
\item Theorem~\ref{thm:V107-exact-spectrum} reduces the full stable matrix to
      two-plane V106 blocks and residual Neumann lines.
\item Theorem~\ref{thm:V107-no-uniform-floor} returns an exact normalized
      kernel frequency for every nonzero two-plane and glancing kernels for
      every residual line.
\item Proposition~\ref{prop:V107-BRST-incidence} proves that matched BRST
      rows close but do not remove the classified stable kernel.
\end{enumerate}

\paragraph{Next forced step.}
V108 should add the most general real tangential spatial principal pencil
(A^0D_T+A^1D_1+A^2D_2), derive its action and BRST restrictions, and test
its determinant on the full imaginary-frequency stable hemisphere.  If an
elliptic Clifford-type pencil passes that finite-dimensional test, its
boundary degrees of freedom, stress, and ghost signs must be identified.  If
all local pencils retain a kernel, the next branch is the exact nonlocal
Calder\'on graph.  Merely increasing the number of time-coupled copies has
been exhausted by V107.

\begin{assessment}[V107 continuum status]
\label{ass:V107-continuum-status}
The full Standard--Model--Einstein continuum remains unproved.  V107 rules
out every finite real time-local conservative collar dilation in the declared
class, including all orthogonal multi-copy extensions of V106.  Spatially
dispersive or nonlocal gauge domains, the fermion branch selection, common
quantum coefficients and form bounds, curved radiative equivariance,
nonlinear existence, and the cofinal interacting limit remain open.
\end{assessment}

\section*{V107 exact checks and append-only record}
\addcontentsline{toc}{section}{V107 exact checks and append-only record}

The exact classification checksum is
\begin{equation}
 \boxed{
 A=\frac12(C-C^T),
 \qquad
 O^TAO=\operatorname{diag}(\alpha_1J,\ldots,\alpha_mJ,0_r),
 \qquad
 \operatorname{Sing}(\lambda I+sA)
 =\{|\lambda\pm i\alpha_js|,|\lambda|^{[r]}\}.}
 \label{eq:V107-final-check}
\end{equation}
The append operation preserves the complete V106 prefix byte for byte before
its former live terminator.  The assembled V107 file is checked for one live
final terminator, balanced environments and braces, duplicate labels, newly
introduced unresolved references, display math delimiters, and control
characters before the exact V106 predecessor is removed.

% =============================================================================
% END APPENDED VERSION V107 MATERIAL
% =============================================================================

\clearpage
% =============================================================================
% BEGIN APPENDED VERSION V108 MATERIAL
% =============================================================================
\section{V108: the minimal real Clifford boundary action and its exact
stable-sheet intersection}
\label{sec:V108-Clifford-pencil}

V107 exhausted finite real conservative collars whose principal boundary
action differentiates only in the selected time direction.  The first
unexamined local escape is a spatially dispersive self-adjoint boundary
pencil.  V108 tests the maximally symmetric minimal example: the real
four-dimensional quaternionic representation of three anticommuting skew
generators.

The tangential Clifford symbol is invertible at every nonzero real boundary
covector and therefore removes the bare V35 glancing factor.  Its exact
stable determinant nevertheless vanishes for every coupling strength.  Below
the characteristic speed the zero lies on the imaginary-frequency boundary;
at the critical speed it lies on the static face; above that speed it becomes
an exponentially growing stable mode.  Thus this natural local elliptic
pencil also fails the common boundary realization.

\subsection{Quaternionic principal action}
\label{subsec:V108-quaternionic-action}

On (\mathbb R^4), define the real matrices
\begin{align}
 E_0&:=
 \begin{pmatrix}
 0&-1&0&0\\
 1&0&0&0\\
 0&0&0&-1\\
 0&0&1&0
 \end{pmatrix},
 &
 E_1&:=
 \begin{pmatrix}
 0&0&-1&0\\
 0&0&0&1\\
 1&0&0&0\\
 0&-1&0&0
 \end{pmatrix},
 \label{eq:V108-E01}\\
 E_2&:=
 \begin{pmatrix}
 0&0&0&-1\\
 0&0&-1&0\\
 0&1&0&0\\
 1&0&0&0
 \end{pmatrix}.
 \label{eq:V108-E2}
\end{align}
These are left multiplication by the three quaternionic units in a real
basis.  Direct multiplication gives
\begin{equation}
 E_a^T=-E_a,
 \qquad
 E_aE_b+E_bE_a=-2\delta_{ab}I_4,
 \qquad 0\le a,b\le2.
 \label{eq:V108-real-Clifford}
\end{equation}

For a four-component real collar trace (q), its conormal (p=D_nq), and a
coupling (c\ge0), take
\begin{equation}
 S_{\partial,c}[q]
 :=\frac c2\int_{\mathcal T}
 q^T(E_0D_T+E_1D_1+E_2D_2)q.
 \label{eq:V108-Clifford-action}
\end{equation}

\begin{proposition}[Action incidence]
\label{prop:V108-action-incidence}
The Euler row of the combined bulk and boundary action is
\begin{equation}
 C_c(q,p)
 :=p+c(E_0D_T+E_1D_1+E_2D_2)q=0.
 \label{eq:V108-Clifford-row}
\end{equation}
Its graph is Lagrangian for the four-copy wave Green form.
\end{proposition}

\begin{proof}
Each (E_a) is real skew and each constant tangential derivative is formally
skew on compactly supported boundary fields.  Therefore (E_aD_a) is
formally self-adjoint.  The variation calculation of Lemma
\ref{lem:V107-skew-completeness}, applied to each derivative, gives
\eqref{eq:V108-Clifford-row}.  Its conormal operator is formally
self-adjoint, so the graph is Lagrangian by the same direct Green-form
calculation as Proposition~\ref{prop:V107-action-power}.
\end{proof}

For a real oscillatory boundary covector ((\tau,k_1,k_2)\ne0), the
Fourier symbol of the tangential Euler operator is
\begin{equation}
 i c(\tau E_0+k_1E_1+k_2E_2).
 \label{eq:V108-real-tangential-symbol}
\end{equation}
Equation \eqref{eq:V108-real-Clifford} shows that its square is
(c^2(\tau^2+|k|^2)I_4).  Thus for (c>0) the tangential boundary operator
is elliptic.  In particular its real-frequency symbol does not vanish at the
bulk glancing covector.

\subsection{Exact complex stable determinant}
\label{subsec:V108-stable-determinant}

In the V35 Laplace--Fourier convention, put
\begin{equation}
 Q(s,k):=sE_0+ik_1E_1+ik_2E_2.
 \label{eq:V108-Q-symbol}
\end{equation}
Anticommutation gives
\begin{equation}
 Q(s,k)^2=(|k|^2-s^2)I_4,
 \qquad
 \operatorname{tr}Q=0.
 \label{eq:V108-Q-square}
\end{equation}
The stable boundary matrix is
\begin{equation}
 \mathcal B_c(s,k):=\lambda I_4+cQ(s,k),
 \qquad
 \lambda^2=s^2+|k|^2.
 \label{eq:V108-stable-matrix}
\end{equation}

\begin{theorem}[Exact determinant and zero locus]
\label{thm:V108-determinant}
For every complex stable frequency,
\begin{equation}
 \boxed{
 \det\mathcal B_c(s,k)
 =\left[(1+c^2)s^2+(1-c^2)|k|^2\right]^2.}
 \label{eq:V108-exact-determinant}
\end{equation}
For (c>0), every zero is produced by a two-dimensional complex eigenspace
of (Q).  For (c=0), the determinant is (lambda^4).
\end{theorem}

\begin{proof}
When (|k|^2-s^2\ne0), Equation \eqref{eq:V108-Q-square} and the zero trace
show that (Q) has eigenvalues
(+\sqrt{|k|^2-s^2}) and
(-\sqrt{|k|^2-s^2}), each twice.  Hence
\[
 \det(\lambda I_4+cQ)
 =\left(\lambda^2-c^2(|k|^2-s^2)\right)^2.
\]
Substitute (lambda^2=s^2+|k|^2) to obtain
\eqref{eq:V108-exact-determinant}.  Both sides are polynomials in (s,k), so
the identity extends across (|k|^2=s^2).  The eigenspace multiplicity follows
from the same factorization; the case (c=0) is immediate.
\end{proof}

\subsection{Subcritical, critical, and supercritical failures}
\label{subsec:V108-three-regimes}

\begin{theorem}[No Clifford coupling passes the stable test]
\label{thm:V108-no-coupling}
For every (c\ge0), the normalized stable matrix
\eqref{eq:V108-stable-matrix} has zero infimum singular value.
More precisely:
\begin{enumerate}[label=\textup{(\roman*)},leftmargin=3.3em]
\item If (0<c<1), set
      \begin{equation}
       s_c=i\sqrt{\frac{1-c^2}{2}},
       \qquad
       |k_c|=\sqrt{\frac{1+c^2}{2}},
       \qquad
       \lambda_c=c.
       \label{eq:V108-subcritical-frequency}
      \end{equation}
      Then ((s_c,k_c)\in\overline{\mathbb S}_+) and
      (ker\mathcal B_c(s_c,k_c)\ne0).
\item If (c=1), take (s_c=0), (|k_c|=1), and
      (lambda_c=1).  Then
      (ker\mathcal B_1(s_c,k_c)\ne0).
\item If (c>1), set
      \begin{equation}
       s_c=\sqrt{\frac{c^2-1}{2c^2}},
       \qquad
       |k_c|=\sqrt{\frac{c^2+1}{2c^2}},
       \qquad
       \lambda_c=1.
       \label{eq:V108-supercritical-frequency}
      \end{equation}
      Then (operatorname{Re}s_c>0), the frequency is normalized, and
      (ker\mathcal B_c(s_c,k_c)\ne0).  This is an exponentially growing
      stable mode.
\item If (c=0), the original V35 glancing vector is a kernel.
\end{enumerate}
Consequently
\begin{equation}
 \boxed{
 \inf_{\overline{\mathbb S}_+}s_{\min}\mathcal B_c=0
 \quad\text{for every }c\ge0.}
 \label{eq:V108-zero-floor}
\end{equation}
\end{theorem}

\begin{proof}
For (0<c<1), the displayed frequency satisfies
(|s_c|^2+|k_c|^2=1) and
\[
 \lambda_c^2=s_c^2+|k_c|^2
 =-\frac{1-c^2}{2}+\frac{1+c^2}{2}=c^2.
\]
It also makes the scalar factor in
\eqref{eq:V108-exact-determinant} vanish.  The stable branch has
(lambda_c=c>0), proving \textup{(i)}.  At (c=1), the determinant is
(4s^4), so the stated static frequency proves \textup{(ii)}.

For (c>1), normalization and (lambda_c=1) follow directly from
\eqref{eq:V108-supercritical-frequency}.  Moreover,
\[
 (1+c^2)s_c^2+(1-c^2)|k_c|^2=0,
\]
so Theorem~\ref{thm:V108-determinant} gives a kernel with
(operatorname{Re}s_c>0).  This proves \textup{(iii)}.  Part \textup{(iv)}
is the V35 Neumann factor.  The final infimum follows in all cases.
\end{proof}

\begin{corollary}[The elliptic tangential symbol does not imply
complementing]
\label{cor:V108-elliptic-not-complementing}
For every (0<c<1), the tangential Clifford operator is elliptic and its
stable matrix is invertible at bulk glancing, yet it has the distinct
imaginary-frequency kernel \eqref{eq:V108-subcritical-frequency}.  For
(c>1), increasing the boundary speed converts that boundary kernel into an
open-half-plane instability rather than eliminating it.
\end{corollary}

\begin{proof}
Tangential ellipticity was proved after
\eqref{eq:V108-real-tangential-symbol}.  The kernel locations and their
frequency types are Theorem~\ref{thm:V108-no-coupling}.
\end{proof}

\subsection{BRST extension and physical scope}
\label{subsec:V108-BRST}

Let the four collar copies carry a ghost column (gamma\in\mathbb R^4), an
antighost column, and a multiplier column.  Apply (C_c) from
\eqref{eq:V108-Clifford-row} to each and use
(mathsf s q_\mu=D_\mu\gamma), (mathsf s\gamma=0), with the usual
antighost--multiplier doublet.

\begin{proposition}[BRST closure and inherited determinant]
\label{prop:V108-BRST}
The matched Clifford rows form a frozen principal BRST subcomplex.  The gauge,
ghost, and antighost wave blocks all inherit the determinant
\eqref{eq:V108-exact-determinant}; hence BRST closure does not remove any of
the kernels in Theorem~\ref{thm:V108-no-coupling}.
\end{proposition}

\begin{proof}
All (E_a) are constant, so (C_c) commutes with every frozen spacetime
derivative.  Therefore
\[
 \mathsf s C_c(q_\mu,D_nq_\mu)
 =D_\mu C_c(\gamma,D_n\gamma)=0.
\]
The antighost row maps to the identical multiplier row.  Each differential
member has the same stable matrix \eqref{eq:V108-stable-matrix}, proving the
determinant statement.
\end{proof}

\begin{firewall}[Scope of the Clifford failure]
\label{fire:V108-scope}
V108 rules out the isotropic quaternionic pencil
(c(E_0D_T+E_1D_1+E_2D_2)) for every real coupling (c).  It does not prove
that every anisotropic skew matrix pencil has a stable zero.  Nor does the
four-copy auxiliary fibre yet have an identified Standard Model boundary
particle interpretation.  A different local pencil must pass the exact
stable determinant before its stress or quantum contribution is considered.
\end{firewall}

\subsection{V108 ledger and next acquisition}
\label{subsec:V108-ledger}

\paragraph{New exact conclusions.}
\begin{enumerate}[label=\textup{(\arabic*)},leftmargin=2.8em]
\item Equations \eqref{eq:V108-E01}--\eqref{eq:V108-real-Clifford} give an
      explicit minimal real tangential Clifford representation.
\item Proposition~\ref{prop:V108-action-incidence} proves that its boundary
      graph comes from a real local action and is Lagrangian.
\item Theorem~\ref{thm:V108-determinant} computes the complete complex stable
      determinant without a numerical sample.
\item Theorem~\ref{thm:V108-no-coupling} classifies the exact failure as an
      imaginary boundary mode for (0<c<1), a static mode for (c=1), and a
      growing mode for (c>1).
\item Proposition~\ref{prop:V108-BRST} proves closure of the matched BRST
      complex and persistence of the same determinant.
\end{enumerate}

\paragraph{Next forced step.}
V109 should derive a criterion for a general real skew tangential pencil
(A^0D_T+A^1D_1+A^2D_2), retaining the exact signed determinant before any
norm majorant.  It should either prove a stable-sheet intersection under
checkable spectral hypotheses or return an explicit anisotropic pencil with
a uniform floor.  If a passing pencil exists, its field representation and
stress must be identified before importing it.  If the criterion covers all
finite local pencils, the gauge branch should move to the nonlocal
Calder\'on realization.

\begin{assessment}[V108 continuum status]
\label{ass:V108-continuum-status}
The full Standard--Model--Einstein continuum remains unproved.  V108 proves
that the minimal tangentially elliptic real Clifford action cannot furnish
the missing gauge/ghost boundary domain at any coupling.  A general local or
nonlocal gauge realization, the fermion branch selection, common quantum
coefficients and form bounds, curved radiative equivariance, nonlinear
existence, and the cofinal interacting limit remain open.
\end{assessment}

\section*{V108 exact checks and append-only record}
\addcontentsline{toc}{section}{V108 exact checks and append-only record}

The exact determinant checksum is
\begin{equation}
 \boxed{
 Q^2=(|k|^2-s^2)I_4,
 \qquad
 \det(\lambda I_4+cQ)
 =\bigl((1+c^2)s^2+(1-c^2)|k|^2\bigr)^2.}
 \label{eq:V108-final-check}
\end{equation}
The append operation preserves the complete V107 prefix byte for byte before
its former live terminator.  The assembled V108 file is checked for one live
final terminator, balanced environments and braces, duplicate labels, newly
introduced unresolved references, display math delimiters, and control
characters before the exact V107 predecessor is removed.

% =============================================================================
% END APPENDED VERSION V108 MATERIAL
% =============================================================================

\clearpage
% =============================================================================
% BEGIN APPENDED VERSION V109 MATERIAL
% =============================================================================
\section{V109: spectral intersection criteria for general real skew
tangential boundary pencils}
\label{sec:V109-skew-pencil}

V108 computed one isotropic Clifford determinant.  A general real local
quadratic boundary action need not satisfy Clifford relations, so its stable
matrix cannot be reduced by the same scalar square identity.  V109 replaces
that special calculation with a variational spectral criterion valid for
arbitrary finite real skew coefficient matrices.

Two exact conclusions follow.  Every strictly subcharacteristic spatial
pencil intersects the stable sheet provided its full principal symbol is not
zero.  Independently, every purely spatial skew pencil fails for every
coefficient size: a small coefficient gives an imaginary-frequency boundary
kernel, a unit coefficient gives a static kernel, and a large coefficient
gives an exponentially growing mode.  The only class not decided here is a
mixed time--space pencil whose static spatial norm reaches or exceeds one and
whose time coefficient changes the open-half-plane spectrum.

\subsection{General action and imaginary-frequency reduction}
\label{subsec:V109-general-action}

Let (A^0,A^1,A^2\in M_N(\mathbb R)) satisfy
\begin{equation}
 (A^a)^T=-A^a,
 \qquad a=0,1,2,
 \label{eq:V109-skew-coefficients}
\end{equation}
and consider the principal boundary action
\begin{equation}
 S_\partial[q]
 :=\frac12\int_{\mathcal T}
 q^T(A^0D_T+A^1D_1+A^2D_2)q.
 \label{eq:V109-general-skew-action}
\end{equation}
The variation argument of V107 gives the action-Lagrangian boundary row
\begin{equation}
 p+(A^0D_T+A^1D_1+A^2D_2)q=0.
 \label{eq:V109-general-skew-row}
\end{equation}
Its complex stable matrix is
\begin{equation}
 \mathcal B_A(s,k)
 :=\lambda I_N+sA^0+i k_1A^1+i k_2A^2,
 \qquad
 \lambda^2=s^2+|k|^2.
 \label{eq:V109-general-stable-matrix}
\end{equation}

For (e=(e_1,e_2)\in\mathbb S^1), write
\begin{equation}
 A(e):=e_1A^1+e_2A^2.
 \label{eq:V109-spatial-pencil}
\end{equation}

\begin{lemma}[Hermitian boundary-axis form]
\label{lem:V109-Hermitian-axis}
At (s=i\tau), with (|k|\ge|\tau|), set
\begin{equation}
 K(\tau,k):=\tau A^0+k_1A^1+k_2A^2.
 \label{eq:V109-K-pencil}
\end{equation}
Then
\begin{equation}
 \mathcal B_A(i\tau,k)
 =\sqrt{|k|^2-\tau^2}\,I_N+iK(\tau,k)
 \label{eq:V109-Hermitian-stable-matrix}
\end{equation}
is Hermitian.  The eigenvalues of (iK) occur in pairs
(\pm\mu_j(\tau,k)), together with zero eigenvalues, and
\begin{equation}
 \max_j\mu_j(\tau,k)=\|K(\tau,k)\|_{\rm op}.
 \label{eq:V109-top-skew-value}
\end{equation}
Consequently the stable matrix is singular whenever
\begin{equation}
 \sqrt{|k|^2-\tau^2}
 =\|K(\tau,k)\|_{\rm op}>0.
 \label{eq:V109-intersection-equation}
\end{equation}
\end{lemma}

\begin{proof}
The matrix (K) is real skew, so (iK) is Hermitian.  Its nonzero complex
eigenvalues are real pairs (\pm\mu_j), because the eigenvalues of a real
skew matrix are pairs (\pm i\mu_j).  Its operator norm is the largest
(mu_j).  If \eqref{eq:V109-intersection-equation} holds, the eigenvalue
(-\|K\|_{\rm op}) of (iK) cancels the scalar stable root in
\eqref{eq:V109-Hermitian-stable-matrix}.
\end{proof}

\subsection{A general subcharacteristic intersection theorem}
\label{subsec:V109-subcharacteristic}

Define the static spatial speed
\begin{equation}
 a_{\rm sp}
 :=\sup_{e\in\mathbb S^1}\|A(e)\|_{\rm op}.
 \label{eq:V109-spatial-speed}
\end{equation}

\begin{theorem}[Strictly subcharacteristic pencils have a stable boundary
kernel]
\label{thm:V109-subcharacteristic-kernel}
Assume
\begin{equation}
 a_{\rm sp}<1.
 \label{eq:V109-subcharacteristic-hypothesis}
\end{equation}
If at least one of (A^0,A^1,A^2) is nonzero, then there exist
\(e\in\mathbb S^1\), \(\varepsilon\in\{-1,1\}\), and
\(r_*\) strictly between \(0\) and \(\varepsilon\) such that
\begin{equation}
 \sqrt{1-r_*^2}
 =\|r_*A^0+A(e)\|_{\rm op}>0.
 \label{eq:V109-radial-crossing}
\end{equation}
After the homogeneous normalization
\begin{equation}
 |k_*|=(1+r_*^2)^{-1/2},
 \qquad
 \tau_*=r_*|k_*|,
 \qquad
 k_*=|k_*|e,
 \label{eq:V109-normalized-crossing}
\end{equation}
the matrix (mathcal B_A(i\tau_*,k_*)) has a nonzero kernel.  Hence
\begin{equation}
 \inf_{\overline{\mathbb S}_+}s_{\min}\mathcal B_A=0.
 \label{eq:V109-subcharacteristic-zero-floor}
\end{equation}
If all three matrices vanish, the same conclusion holds at bulk glancing.
\end{theorem}

\begin{proof}
Because the pencil is not identically zero, there are (e) and
\(\varepsilon\in\{-1,1\}\) for which
\begin{equation}
 \varepsilon A^0+A(e)\ne0.
 \label{eq:V109-nonzero-null-pencil}
\end{equation}
Indeed, if both signs vanished for every (e), adding and subtracting the
two identities would give (A^0=0) and (A(e)=0) for every (e), hence
(A^1=A^2=0).

On the interval joining \(0\) to \(\varepsilon\), define the continuous
function
\begin{equation}
 f(r):=\sqrt{1-r^2}-\|rA^0+A(e)\|_{\rm op}.
 \label{eq:V109-crossing-function}
\end{equation}
At (r=0), hypothesis \eqref{eq:V109-subcharacteristic-hypothesis} gives
(f(0)=1-\|A(e)\|_{\rm op}>0).  At (r=\varepsilon),
\eqref{eq:V109-nonzero-null-pencil} gives
(f(\varepsilon)<0).  The intermediate value theorem supplies an interior
zero (r_*).  Its right side is positive because its left side is positive.

Under \eqref{eq:V109-normalized-crossing}, one has
(|s|^2+|k|^2=\tau_*^2+|k_*|^2=1), and homogeneity turns
\eqref{eq:V109-radial-crossing} into
\eqref{eq:V109-intersection-equation}.  Lemma
\ref{lem:V109-Hermitian-axis} supplies a kernel.  If the pencil is zero,
(mathcal B_A=\lambda I_N), which vanishes at glancing.
\end{proof}

\begin{corollary}[A small perturbation of the Neumann action cannot open a
uniform gap]
\label{cor:V109-small-action-no-gap}
No finite real action pencil with (a_{\rm sp}<1) can repair the V35
unweighted floor, regardless of the size of its time coefficient (A^0).
The kernel may move away from glancing, but it remains on the normalized
imaginary-frequency stable boundary.
\end{corollary}

\begin{proof}
Apply Theorem~\ref{thm:V109-subcharacteristic-kernel}.  The proof makes no
bound on (A^0).
\end{proof}

\subsection{Complete classification of purely spatial skew pencils}
\label{subsec:V109-pure-spatial}

Now set (A^0=0), without any smallness assumption.  If the spatial pencil
is nonzero, choose (e\in\mathbb S^1) and put
\begin{equation}
 \mu:=\|A(e)\|_{\rm op}>0.
 \label{eq:V109-mu-direction}
\end{equation}

\begin{theorem}[Every purely spatial real skew pencil fails]
\label{thm:V109-pure-spatial-no-go}
For (A^0=0), the stable matrix has an exact kernel in each of the following
regimes:
\begin{enumerate}[label=\textup{(\roman*)},leftmargin=3.3em]
\item If (0<\mu<1), let
      \begin{equation}
       |k|=(2-\mu^2)^{-1/2},
       \qquad
       s=i|k|\sqrt{1-\mu^2},
       \qquad
       \lambda=\mu|k|.
       \label{eq:V109-pure-subcritical}
      \end{equation}
      This is a normalized imaginary-frequency kernel.
\item If (mu=1), take (s=0), (k=e), and (lambda=1).  This is a
      normalized static kernel.
\item If (mu>1), let
      \begin{equation}
       |k|=\mu^{-1},
       \qquad
       s=\mu^{-1}\sqrt{\mu^2-1}>0,
       \qquad
       \lambda=1.
       \label{eq:V109-pure-supercritical}
      \end{equation}
      This is a normalized exponentially growing stable kernel.
\end{enumerate}
If (A^1=A^2=0), the stable matrix is (lambda I_N) and fails at
glancing.  Therefore every purely spatial finite real skew pencil has zero
uniform stable floor.
\end{theorem}

\begin{proof}
The Hermitian matrix (iA(e)) has an eigenvector (v) of eigenvalue
(-\mu).  At (k=|k|e), the stable matrix acts on that vector by
\begin{equation}
 \mathcal B_A(s,k)v=(\lambda-\mu|k|)v.
 \label{eq:V109-pure-eigenvalue}
\end{equation}

In case \textup{(i)},
(|s|^2+|k|^2=(2-\mu^2)|k|^2=1) and
(lambda^2=s^2+|k|^2=\mu^2|k|^2), so the stable branch has
(lambda=\mu|k|).  Case \textup{(ii)} is immediate.  In case
\textup{(iii)}, (s^2+|k|^2=1), hence (lambda=1), while
(mu|k|=1).  Equation \eqref{eq:V109-pure-eigenvalue} vanishes in all
three regimes.  The zero-pencil assertion is the bare Neumann calculation.
\end{proof}

\begin{corollary}[Anisotropy does not rescue a purely spatial action]
\label{cor:V109-anisotropy-no-rescue}
Theorem~\ref{thm:V109-pure-spatial-no-go} requires no commutation, Clifford,
or rotational symmetry among (A^1,A^2).  One nonzero spatial direction is
enough to produce the exact stable-sheet intersection.
\end{corollary}

\begin{proof}
Only the real skew spectral theorem for the single matrix (A(e)) was used.
\end{proof}

\subsection{BRST and action status}
\label{subsec:V109-BRST-status}

Apply the row \eqref{eq:V109-general-skew-row} to matched gauge, ghost,
antighost, and multiplier columns, with the same linearized BRST differential
as V106--V108.

\begin{proposition}[The spectral obstruction survives BRST closure]
\label{prop:V109-BRST}
Every general skew row above is action-Lagrangian and forms a frozen principal
BRST subcomplex under the matched ghost rows.  The gauge and ghost blocks have
the identical stable matrix \eqref{eq:V109-general-stable-matrix}, so the
kernels in Theorems~\ref{thm:V109-subcharacteristic-kernel} and
\ref{thm:V109-pure-spatial-no-go} survive in the full principal direct sum.
\end{proposition}

\begin{proof}
Action incidence follows because every (A^aD_a) is formally self-adjoint.
Constant coefficients commute with frozen spacetime derivatives, giving
\[
 \mathsf s C_A(q_\mu,D_nq_\mu)
 =D_\mu C_A(\gamma,D_n\gamma).
\]
The matched ghost row kills the right side, and the antighost maps to the
matched multiplier row.  Direct summation cannot lift a kernel supported in
one gauge or ghost component.
\end{proof}

\begin{firewall}[The undecided mixed supercharacteristic region]
\label{fire:V109-undecided-region}
V109 does not prove a no-go theorem for every mixed pencil.  The remaining
finite local region has (A^0\ne0) and
(a_{\rm sp}\ge1).  There the open-half-plane matrix need not be Hermitian,
and the one-direction argument for a purely spatial pencil no longer
diagonalizes the time coefficient.  A proof must analyze the exact complex
matrix pencil; replacing it by separate norms would discard the signed
spectral information needed to decide the determinant.
\end{firewall}

\subsection{V109 ledger and mandatory V110 audit}
\label{subsec:V109-ledger}

\paragraph{New exact conclusions.}
\begin{enumerate}[label=\textup{(\arabic*)},leftmargin=2.8em]
\item Lemma~\ref{lem:V109-Hermitian-axis} reduces the imaginary-frequency
      complementing test to the exact skew singular value equation
      \eqref{eq:V109-intersection-equation}.
\item Theorem~\ref{thm:V109-subcharacteristic-kernel} proves a stable-sheet
      intersection for every nonzero strictly subcharacteristic pencil.
\item Theorem~\ref{thm:V109-pure-spatial-no-go} classifies every purely
      spatial pencil by a boundary, static, or growing kernel according to
      one exact directional singular value.
\item Proposition~\ref{prop:V109-BRST} proves action and BRST incidence and
      embeds the same kernel into the matched gauge/ghost direct sum.
\end{enumerate}

\paragraph{Next forced step.}
V110 is a five-version checkpoint and must first audit the newest shared
Yang--Mills and Navier--Stokes notes, recording exact versions and hashes.
It should then address the remaining mixed supercharacteristic region or use
the audit to justify moving directly to the nonlocal Calder\'on graph.  Any
candidate local pencil must be tested on the exact complex stable matrix, and
any passing algebraic matrix must still identify its physical boundary fields,
stress, and BRST grading.

\begin{assessment}[V109 continuum status]
\label{ass:V109-continuum-status}
The full Standard--Model--Einstein continuum remains unproved.  V109 rules
out all strictly subcharacteristic real skew boundary pencils and every
purely spatial skew pencil.  The mixed supercharacteristic local region or a
nonlocal gauge realization, the fermion branch selection, common quantum
coefficients and form bounds, curved radiative equivariance, nonlinear
existence, and the cofinal interacting limit remain open.
\end{assessment}

\section*{V109 exact checks and append-only record}
\addcontentsline{toc}{section}{V109 exact checks and append-only record}

The exact spectral checksum is
\begin{equation}
 \boxed{
 \mathcal B_A(i\tau,k)
 =\sqrt{|k|^2-\tau^2}\,I+iK(\tau,k),
 \qquad
 \sqrt{|k|^2-\tau^2}=\|K(\tau,k)\|_{\rm op}
 \Longrightarrow\ker\mathcal B_A\ne0.}
 \label{eq:V109-final-check}
\end{equation}
The append operation preserves the complete V108 prefix byte for byte before
its former live terminator.  The assembled V109 file is checked for one live
final terminator, balanced environments and braces, duplicate labels, newly
introduced unresolved references, display math delimiters, and control
characters before the exact V108 predecessor is removed.

% =============================================================================
% END APPENDED VERSION V109 MATERIAL
% =============================================================================

\clearpage
% =============================================================================
% BEGIN APPENDED VERSION V110 MATERIAL
% =============================================================================
\section{V110: companion audit and complete classification of the two-collar
mixed skew pencil}
\label{sec:V110-two-collar-classification}

V109 left open mixed time--space pencils whose static spatial norm is at least
one.  This fifth iteration first performs the mandatory companion audit.  It
then closes the smallest remaining local class completely.  In two real
collar dimensions every skew coefficient is a multiple of one matrix (J),
so the time and spatial contributions can be retained in one exact signed
determinant.

The result is a full two-collar no-go theorem.  If the spatial coefficient is
below the exact mixed threshold, the determinant has a normalized
imaginary-frequency zero; above the threshold it has a normalized growing
zero.  Equality also has an imaginary-frequency zero.  Thus arbitrary
time--space mixing does not repair the minimal conservative dilation.

\subsection{Mandatory V110 companion audit}
\label{subsec:V110-audit}

The read-only companion snapshots are
\begin{center}
\begin{tabular}{|p{0.48\textwidth}|r|p{0.31\textwidth}|}
\hline
Companion file & Bytes & SHA--256\\
\hline
\texttt{Renewal\_NS\_Stability\_Research\_Notes\_V31.tex}
& (887537) &
\texttt{F1121B62238AD5491BB7CF03635EA9934942EDF573ED8FAAA9EE79E1D38A6696}\\
\hline
\texttt{Renewal\_Yang\_Mills\_Mass\_Gap\_Research\_Notes\_V102.tex}
& (2432402) &
\texttt{B4A616578B653A3B8A97C3DA94E2A1BACB94B583B2637364F36B2871D4B8303B}\\
\hline
\end{tabular}
\end{center}

The Navier--Stokes note is unchanged from the V105 audit.  Its V31 terminal
example still shows that a signed balanced-helicity readout may vanish while
the positive total critical coarea is nonzero.  The Yang--Mills note has
advanced from V99 to V102.  V100 retains the exact matrix-valued singular law
and shows that one Casimir moment does not control a near-unit sequence.  V101
separates aligned heat decay from crossing leakage.  V102 forces a heavy
internal-pair or off-cut carrier and retains the un-factorized joint pair/core
level profile; bounded auxiliary dimensions give no automatic decay.

\begin{assessment}[Transfer used in V110]
\label{ass:V110-audit-transfer}
The companion calculations again prescribe an order of proof, not an imported
estimate.  The exact signed object and all correlated coordinates must be
assembled before a positive norm, moment, or separated majorant is taken.
Accordingly V110 keeps the time coefficient, the spatial coefficient, the
stable square root, and their complex cross term in one determinant.  The
classification below is proved directly for the SM gauge boundary pencil.
\end{assessment}

\subsection{The complete two-dimensional skew action}
\label{subsec:V110-complete-action}

Let
\begin{equation}
 J:=\begin{pmatrix}0&-1\\1&0\end{pmatrix}.
 \label{eq:V110-J}
\end{equation}
Every real two-by-two skew matrix is a unique real multiple of (J).  Hence
the most general principal action in the V109 class on two collar traces is
\begin{equation}
 S_{\partial}[q]
 :=\frac12\int_{\mathcal T}
 q^TJ(\alpha D_T+\beta_1D_1+\beta_2D_2)q,
 \qquad
 \alpha\in\mathbb R,
 \quad \boldsymbol\beta=(\beta_1,\beta_2)\in\mathbb R^2.
 \label{eq:V110-general-two-action}
\end{equation}
Its stationary conormal row is
\begin{equation}
 p+J(\alpha D_T+\boldsymbol\beta\cdot D_\parallel)q=0.
 \label{eq:V110-general-two-row}
\end{equation}
As in V107--V109, the graph is Lagrangian because the real skew matrix \(J\)
multiplies formally skew tangential derivatives.

Let \(\beta:=|\boldsymbol\beta|\).  If \(\beta>0\), choose
\(e=\boldsymbol\beta/\beta\); if \(\beta=0\), choose any
\(e\in\mathbb S^1\).  Restricting to \(k=\kappa e\), \(\kappa>0\), loses no
generality for producing a counterexample.  The stable matrix is
\begin{equation}
 \mathcal B_{\alpha,\beta}(s,\kappa)
 :=\lambda I_2+(\alpha s+i\beta\kappa)J,
 \qquad
 \lambda^2=s^2+\kappa^2.
 \label{eq:V110-two-stable-matrix}
\end{equation}

\begin{theorem}[Exact correlated determinant]
\label{thm:V110-exact-determinant}
For every complex stable frequency,
\begin{align}
 \det\mathcal B_{\alpha,\beta}
 &=\lambda^2+(\alpha s+i\beta\kappa)^2
 \label{eq:V110-determinant-first}\\
 &=(1+\alpha^2)s^2
   +2i\alpha\beta\kappa s
   +(1-\beta^2)\kappa^2.
 \label{eq:V110-determinant-quadratic}
\end{align}
With \(w=s/\kappa\), its two roots are
\begin{equation}
 w_\pm
 =\frac{-i\alpha\beta
 \mathbin\pm\sqrt{\beta^2-1-\alpha^2}}
 {1+\alpha^2},
 \label{eq:V110-root-formula}
\end{equation}
where the square root is taken in either order; the unordered pair is
unambiguous.
\end{theorem}

\begin{proof}
For every scalar (z),
\(\det(\lambda I_2+zJ)=\lambda^2+z^2\), because \(J^2=-I_2\).  Set
\(z=\alpha s+i\beta\kappa\) and substitute
\(\lambda^2=s^2+\kappa^2\) to obtain
\eqref{eq:V110-determinant-first}--%
\eqref{eq:V110-determinant-quadratic}.  Division by (kappa^2) gives a
quadratic in (w).  Its discriminant is
\begin{equation}
 -4\alpha^2\beta^2
 -4(1+\alpha^2)(1-\beta^2)
 =4(\beta^2-1-\alpha^2),
 \label{eq:V110-discriminant}
\end{equation}
which proves \eqref{eq:V110-root-formula}.
\end{proof}

\subsection{The exact threshold classification}
\label{subsec:V110-threshold}

\begin{lemma}[The imaginary roots lie on the closed stable axis]
\label{lem:V110-imaginary-root-bound}
Assume \(\beta^2\le1+\alpha^2\), and put
\begin{equation}
 R:=\sqrt{1+\alpha^2-\beta^2},
 \qquad
 r_\pm:=\frac{-\alpha\beta\pm R}{1+\alpha^2}.
 \label{eq:V110-r-roots}
\end{equation}
Then \(w_\pm=ir_\pm\) and
\begin{equation}
 |r_\pm|\le1.
 \label{eq:V110-r-bound}
\end{equation}
\end{lemma}

\begin{proof}
The root identity follows from \eqref{eq:V110-root-formula}.  Cauchy--Schwarz
gives
\begin{align}
 (|\alpha\beta|+R)^2
 &\le(1+\alpha^2)(\beta^2+R^2)\notag\\
 &=(1+\alpha^2)^2.
 \label{eq:V110-Cauchy-bound}
\end{align}
Hence
\(|-\alpha\beta\pm R|\le1+\alpha^2\), proving
\eqref{eq:V110-r-bound}.
\end{proof}

\begin{theorem}[Complete two-collar mixed-pencil no-go]
\label{thm:V110-two-collar-no-go}
For every \((\alpha,\boldsymbol\beta)\in\mathbb R\times\mathbb R^2\),
the action-Lagrangian stable matrix generated by
\eqref{eq:V110-general-two-action} has zero uniform floor:
\begin{equation}
 \boxed{
 \inf_{\overline{\mathbb S}_+}
 s_{\min}\mathcal B_{\alpha,\beta}=0.}
 \label{eq:V110-zero-floor}
\end{equation}
The zero has the following exact location.
\begin{enumerate}[label=\textup{(\roman*)},leftmargin=3.3em]
\item If \(\beta^2\le1+\alpha^2\), choose either root \(r=r_\pm\) from
      \eqref{eq:V110-r-roots} and set
      \begin{equation}
       \kappa=(1+r^2)^{-1/2},
       \qquad s=ir\kappa,
       \qquad k=\kappa e.
       \label{eq:V110-normalized-imaginary-zero}
      \end{equation}
      Then ((s,k)\in\overline{\mathbb S}_+),
      (|s|^2+|k|^2=1), and
      \(\det\mathcal B_{\alpha,\beta}=0\).
\item If \(\beta^2>1+\alpha^2\), put
      \begin{equation}
       w_+:=
       \frac{\sqrt{\beta^2-1-\alpha^2}-i\alpha\beta}
       {1+\alpha^2},
       \qquad
       \kappa=(1+|w_+|^2)^{-1/2},
       \qquad
       s=\kappa w_+,
       \qquad k=\kappa e.
       \label{eq:V110-normalized-growing-zero}
      \end{equation}
      Then (operatorname{Re}s>0), the frequency is normalized, and
      \(\det\mathcal B_{\alpha,\beta}=0\).  This is an exponentially growing
      stable mode.
\end{enumerate}
\end{theorem}

\begin{proof}
In case \textup{(i)}, Lemma~\ref{lem:V110-imaginary-root-bound} gives
\(|r|\le1\).  Therefore the stable root
\[
 \lambda=\kappa\sqrt{1-r^2}
\]
is the nonnegative continuation from (operatorname{Re}s>0).  The
normalization is immediate, and \(s/\kappa=ir\) is a root of the exact
quadratic \eqref{eq:V110-determinant-quadratic}.  Hence the determinant
vanishes.  When (|r|=1), this zero lies at bulk glancing; otherwise it lies
on the non-glancing imaginary-frequency boundary.

In case \textup{(ii)}, the chosen (w_+) is the root in
\eqref{eq:V110-root-formula} with positive real part.  Homogeneity of
\eqref{eq:V110-determinant-quadratic} makes the normalized frequency a zero,
and \(\operatorname{Re}s=\kappa\operatorname{Re}w_+>0\).  At any determinant
zero, one of the two eigenvalues
\(\lambda\pm i(\alpha s+i\beta\kappa)\) vanishes, independently of the sign
chosen for the stable square root.  Thus both cases give a nonzero stable
kernel and prove \eqref{eq:V110-zero-floor}.
\end{proof}

\begin{corollary}[Time--space cancellation cannot rescue two collars]
\label{cor:V110-no-cancellation-rescue}
No choice of the time coefficient (alpha), spatial vector
\(\boldsymbol\beta\), or their relative size produces a uniformly
complementing two-collar conservative graph.  Increasing the spatial norm
past the mixed threshold \(\beta^2=1+\alpha^2\) creates a growing mode rather
than a positive margin.
\end{corollary}

\begin{proof}
This is the dichotomy in Theorem~\ref{thm:V110-two-collar-no-go}.
\end{proof}

\subsection{BRST direct sum and exact-carrier discipline}
\label{subsec:V110-BRST}

Use the row \eqref{eq:V110-general-two-row} on the two gauge collars and on
their matched ghost, antighost, and multiplier columns.

\begin{proposition}[BRST closure preserves the correlated zero]
\label{prop:V110-BRST}
The two-collar domain is action-Lagrangian and a frozen principal BRST
subcomplex.  The exact determinant
\eqref{eq:V110-determinant-quadratic} occurs in every differential member of
the quartet.  Hence the kernel in Theorem~\ref{thm:V110-two-collar-no-go}
embeds into the full Standard Model adjoint direct sum.
\end{proposition}

\begin{proof}
Action incidence follows from real skew coefficients multiplying formally
skew tangential derivatives.  Constant coefficients commute with the
linearized BRST differential, so the gauge row maps to the derivative of the
matched ghost row, the antighost row maps to the multiplier row, and the
remaining rows map to zero.  Support the stable kernel in one real adjoint
component and set every other Standard Model field to zero.
\end{proof}

\begin{firewall}[The determinant must not be separated]
\label{fire:V110-no-separation}
The cross term \(2i\alpha\beta\kappa s\) in
\eqref{eq:V110-determinant-quadratic} is part of the physical signed boundary
operator.  Bounding it by separate norms of the time and spatial pieces would
lose the root geometry and could neither prove nor refute complementing.
This is the boundary-symbol analogue of the V110 companion-audit warning: a
joint carrier cannot be replaced by independently optimized marginals.
\end{firewall}

\begin{firewall}[Scope of the two-collar classification]
\label{fire:V110-scope}
Theorem~\ref{thm:V110-two-collar-no-go} is complete for two real collar traces
and arbitrary constant first-order tangential coefficients.  In dimensions
\(N\ge3\), distinct skew coefficient matrices need not commute or share
two-planes, so the determinant is not generally a product of the quadratic
above.  V110 does not rule out those noncommuting pencils or nonlocal
Calder\'on graphs.
\end{firewall}

\subsection{V110 ledger and next acquisition}
\label{subsec:V110-ledger}

\paragraph{New exact conclusions.}
\begin{enumerate}[label=\textup{(\arabic*)},leftmargin=2.8em]
\item The mandatory audit records the exact NS V31 and YM V102 snapshots and
      imports only their order-of-operations lesson.
\item Equation \eqref{eq:V110-general-two-action} exhausts all real
      two-collar first-order tangential action coefficients.
\item Theorem~\ref{thm:V110-exact-determinant} retains and evaluates the full
      complex time--space cross term.
\item Lemma~\ref{lem:V110-imaginary-root-bound} places both subthreshold roots
      on the closed stable imaginary axis.
\item Theorem~\ref{thm:V110-two-collar-no-go} proves a normalized boundary or
      growing kernel for every coefficient pair.
\item Proposition~\ref{prop:V110-BRST} embeds the same exact zero into the
      action and BRST direct sum.
\end{enumerate}

\paragraph{Next forced step.}
V111 should test the smallest genuinely noncommuting skew pencil, which
requires at least four real collar components.  It must compute the exact
stable determinant or a certified Schur/Pfaffian factorization before taking
norms.  If every physically interpretable noncommuting candidate intersects
the stable sheet, the programme should stop enlarging local auxiliary fibres
and construct the nonlocal Calder\'on graph with its action, BRST, stress, and
quantum-domain data.

\begin{assessment}[V110 continuum status]
\label{ass:V110-continuum-status}
The full Standard--Model--Einstein continuum remains unproved.  V110 closes
the complete two-collar mixed local action class and finds no uniformly
complementing member.  Higher-dimensional noncommuting or nonlocal gauge
domains, the fermion branch selection, common quantum coefficients and form
bounds, curved radiative equivariance, nonlinear existence, and the cofinal
interacting limit remain open.
\end{assessment}

\section*{V110 exact checks and append-only record}
\addcontentsline{toc}{section}{V110 exact checks and append-only record}

The exact mixed-pencil checksum is
\begin{equation}
 \boxed{
 \det\mathcal B_{\alpha,\beta}
 =(1+\alpha^2)s^2+2i\alpha\beta\kappa s
  +(1-\beta^2)\kappa^2,
 \qquad
 w_\pm=\frac{-i\alpha\beta\pm
 \sqrt{\beta^2-1-\alpha^2}}{1+\alpha^2}.}
 \label{eq:V110-final-check}
\end{equation}
The append operation preserves the complete V109 prefix byte for byte before
its former live terminator.  The assembled V110 file is checked for one live
final terminator, balanced environments and braces, duplicate labels, newly
introduced unresolved references, display math delimiters, and control
characters before the exact V109 predecessor is removed.  The audited
companion files were read only.

% =============================================================================
% END APPENDED VERSION V110 MATERIAL
% =============================================================================

\clearpage
% =============================================================================
% BEGIN APPENDED VERSION V111 MATERIAL
% =============================================================================
\section{V111: correction of the noncommuting threshold, the odd-collar
glancing theorem, and the even Pfaffian gate}
\label{sec:V111-parity-Pfaffian}

The last acquisition sentence of V110 said that a genuinely noncommuting real
skew pencil requires at least four collar components.  That sentence is
incorrect: \(\mathfrak{so}(3)\) is already noncommutative.  The V110
two-dimensional classification and all of its proofs are unaffected.  V111
corrects the forward pointer and proves that the entire odd-dimensional class,
including dimension three, fails for a simpler parity reason.

At bulk glancing the scalar conormal factor \(\lambda\) vanishes.  What remains
of every real action-compatible first-order tangential pencil is a complex
skew matrix.  Such a matrix is singular in odd dimension.  In even dimension
its determinant is the square of its Pfaffian, giving an exact necessary gate
for any remaining local candidate.

\subsection{The general finite real action pencil}
\label{subsec:V111-general-pencil}

Let \(N\ge1\), and let \(A^0,A^1,A^2\in M_N(\mathbb R)\) be real skew
matrices.  Retain the V109 action row and stable matrix
\begin{align}
 p+(A^0D_T+A^1D_1+A^2D_2)q&=0,
 \label{eq:V111-action-row}\\
 \mathcal B_A(s,k)
 &:=\lambda I_N+sA^0+ik_1A^1+ik_2A^2,
 \qquad
 \lambda^2=s^2+|k|^2.
 \label{eq:V111-stable-matrix}
\end{align}
Fix \(\varepsilon\in\{-1,1\}\) and
\(e=(e_1,e_2)\in\mathbb S^1\).  The normalized glancing points are
\begin{equation}
 s_{\varepsilon}:=\frac{i\varepsilon}{\sqrt2},
 \qquad
 k_e:=\frac e{\sqrt2},
 \qquad
 \lambda=0.
 \label{eq:V111-normalized-glancing}
\end{equation}
Define the real skew null pencil
\begin{equation}
 G_{\varepsilon}(e)
 :=\varepsilon A^0+e_1A^1+e_2A^2.
 \label{eq:V111-null-pencil}
\end{equation}
Then
\begin{equation}
 \mathcal B_A(s_\varepsilon,k_e)
 =\frac{i}{\sqrt2}G_{\varepsilon}(e).
 \label{eq:V111-glancing-matrix}
\end{equation}

\begin{theorem}[Every odd-dimensional real action pencil fails at glancing]
\label{thm:V111-odd-no-go}
If \(N\) is odd, then for every choice of
\(A^0,A^1,A^2\),
\begin{equation}
 \det\mathcal B_A(s_\varepsilon,k_e)=0
 \qquad
 \text{for every }\varepsilon\in\{-1,1\},
 \ e\in\mathbb S^1.
 \label{eq:V111-odd-determinant-zero}
\end{equation}
Consequently
\begin{equation}
 \boxed{
 \inf_{\overline{\mathbb S}_+}
 s_{\min}\mathcal B_A=0.}
 \label{eq:V111-odd-zero-floor}
\end{equation}
This includes arbitrary noncommuting three-by-three skew coefficients.
\end{theorem}

\begin{proof}
For every \(e,\varepsilon\), the matrix
\(G_\varepsilon(e)\) is real skew.  In odd dimension,
\[
 \det G_\varepsilon(e)
 =\det G_\varepsilon(e)^T
 =\det(-G_\varepsilon(e))
 =(-1)^N\det G_\varepsilon(e)
 =-\det G_\varepsilon(e).
\]
Thus its determinant is zero.  Equation
\eqref{eq:V111-glancing-matrix} proves
\eqref{eq:V111-odd-determinant-zero}.  The normalized glancing points belong
to the closed stable hemisphere, so the least singular-value infimum is zero.
Approaching a fixed glancing point from \(\operatorname{Re}s>0\) and using
continuity of singular values gives a normalized stable quasimode if the
complementing condition is formulated only in the open half-plane.
\end{proof}

\begin{corollary}[Dimension three is excluded before noncommutation is used]
\label{cor:V111-three-excluded}
The smallest noncommuting case \(N=3\) cannot repair the gauge/ghost floor.
Its failure follows from odd skew parity and does not depend on commutators
\([A^a,A^b]\).
\end{corollary}

\begin{proof}
Apply Theorem~\ref{thm:V111-odd-no-go}.  Noncommuting elements exist in
\(\mathfrak{so}(3)\), but the determinant identity holds for every triple.
\end{proof}

\subsection{The explicit three-dimensional cross-product form}
\label{subsec:V111-cross-product}

Every \(A\in\mathfrak{so}(3)\) has a unique vector \(v\in\mathbb R^3\) such
that
\begin{equation}
 Ax=v\times x.
 \label{eq:V111-cross-product-map}
\end{equation}
Let \(v^a\) correspond to \(A^a\) and set
\begin{equation}
 w(s,k):=s v^0+ik_1v^1+ik_2v^2.
 \label{eq:V111-complex-cross-vector}
\end{equation}
Then the skew part of \(\mathcal B_A\) is the complex cross-product matrix
\([w]_\times\).

\begin{proposition}[Exact three-collar determinant]
\label{prop:V111-three-determinant}
For \(N=3\),
\begin{equation}
 \boxed{
 \det\mathcal B_A(s,k)
 =\lambda\bigl(\lambda^2+w(s,k)\mathbin{\cdot}w(s,k)\bigr),}
 \label{eq:V111-three-determinant}
\end{equation}
where the dot is the complex bilinear extension of the Euclidean inner
product.  The factor \(\lambda\) is present for every coefficient triple.
\end{proposition}

\begin{proof}
The complex cross-product identity is
\begin{equation}
 [w]_\times^2=w\,w^T-(w\mathbin{\cdot}w)I_3.
 \label{eq:V111-cross-square}
\end{equation}
It implies \([w]_\times w=0\), and on the bilinear orthogonal complement of
\(w\) its characteristic polynomial is
\(t^2+w\mathbin{\cdot}w\).  Therefore
\[
 \det(\lambda I_3+[w]_\times)
 =\lambda(\lambda^2+w\mathbin{\cdot}w).
\]
Polynomial continuation covers the isotropic cases
\(w\mathbin{\cdot}w=0\), proving the formula everywhere.
\end{proof}

\begin{remark}[The parity factor is the complete glancing diagnosis]
\label{rem:V111-parity-factor}
The second factor in \eqref{eq:V111-three-determinant} may create additional
boundary or growing roots, but it is irrelevant to the uniform-floor
decision: the first factor already vanishes on the whole bulk glancing set.
\end{remark}

\subsection{The exact even-dimensional Pfaffian gate}
\label{subsec:V111-Pfaffian}

Now let \(N=2m\).  For a real skew \(2m\times2m\) matrix \(G\), write
\(\operatorname{Pf}G\) for its Pfaffian, normalized by
\begin{equation}
 \det G=(\operatorname{Pf}G)^2.
 \label{eq:V111-Pfaffian-square}
\end{equation}

\begin{theorem}[Necessary and sufficient glancing gate in even dimension]
\label{thm:V111-Pfaffian-gate}
For \(N=2m\), the stable boundary matrix is invertible at every normalized
bulk glancing point if and only if
\begin{equation}
 \boxed{
 \inf_{\substack{\varepsilon\in\{-1,1\}\\ e\in\mathbb S^1}}
 \left|\operatorname{Pf}
 G_\varepsilon(e)\right|>0.}
 \label{eq:V111-uniform-Pfaffian-gate}
\end{equation}
Pointwise invertibility alone is equivalent to
\(\operatorname{Pf}G_\varepsilon(e)\ne0\) for every
\((\varepsilon,e)\).  Compactness of the two null circles makes the pointwise
and positive-infimum statements equivalent.
\end{theorem}

\begin{proof}
Equation \eqref{eq:V111-glancing-matrix} gives
\begin{equation}
 \det\mathcal B_A(s_\varepsilon,k_e)
 =\left(\frac{i}{\sqrt2}\right)^{2m}
 \bigl(\operatorname{Pf}G_\varepsilon(e)\bigr)^2.
 \label{eq:V111-glancing-Pfaffian}
\end{equation}
Thus the stable matrix is invertible at a glancing point exactly when the
corresponding Pfaffian is nonzero.  The Pfaffian is a polynomial in the matrix
entries, hence continuous in \(e\).  The domain
\(\{-1,1\}\times\mathbb S^1\) is compact, so a continuous nonvanishing
absolute value has a positive minimum.  Conversely, a positive minimum
implies pointwise invertibility.
\end{proof}

\begin{corollary}[The V108 Clifford pencil passes only the glancing gate]
\label{cor:V111-Clifford-glancing}
For the quaternionic matrices of V108,
\begin{equation}
 G_\varepsilon(e)^2=-2I_4,
 \qquad
 \left|\operatorname{Pf}G_\varepsilon(e)\right|=2
 \label{eq:V111-Clifford-Pfaffian}
\end{equation}
for every null direction.  Hence that pencil passes
\eqref{eq:V111-uniform-Pfaffian-gate}, even though V108 proved that its full
stable determinant vanishes away from bulk glancing.
\end{corollary}

\begin{proof}
The Clifford anticommutation relations give
\[
 G_\varepsilon(e)^2
 =-(\varepsilon^2+|e|^2)I_4=-2I_4.
\]
Therefore \(\det G_\varepsilon(e)=4\).  The Pfaffian square identity yields
absolute Pfaffian \(2\).  V108's determinant zero lies on its distinct
boundary-characteristic sheet, so there is no contradiction.
\end{proof}

\begin{firewall}[Passing the null-circle Pfaffian is not complementing]
\label{fire:V111-Pfaffian-not-full}
Equation \eqref{eq:V111-uniform-Pfaffian-gate} tests only bulk glancing.  A
candidate which passes it must still be tested on the entire complex stable
hemisphere.  The V108 Clifford pencil is the explicit regression: its
Pfaffian is uniformly positive on both null circles, yet its determinant has
imaginary, static, or growing zeros elsewhere depending on the coupling.
\end{firewall}

\subsection{BRST incidence and corrected acquisition order}
\label{subsec:V111-BRST-ledger}

As in V109--V110, impose the same action row on gauge, ghost, antighost, and
multiplier columns.

\begin{proposition}[Parity and Pfaffian gates persist in the BRST direct sum]
\label{prop:V111-BRST}
The matched rows form a frozen principal BRST subcomplex.  In odd collar
dimension the glancing kernel of Theorem~\ref{thm:V111-odd-no-go} occurs in
every differential quartet member.  In even dimension every candidate must
pass \eqref{eq:V111-uniform-Pfaffian-gate} before the remaining stable
frequencies or Standard Model sectors are assembled.
\end{proposition}

\begin{proof}
Constant coefficient action rows commute with the frozen BRST differential,
as proved in Proposition~\ref{prop:V109-BRST}.  The gauge and ghost blocks
therefore have identical principal stable matrices.  Direct sums preserve a
kernel supported in one member, while even-dimensional glancing
invertibility is exactly Theorem~\ref{thm:V111-Pfaffian-gate}.
\end{proof}

\begin{assessment}[Correction to the V110 forward pointer]
\label{ass:V111-correction}
The smallest genuinely noncommuting real skew pencil has \(N=3\), not
\(N=4\).  Theorem~\ref{thm:V111-odd-no-go} excludes \(N=3\) and every other
odd \(N\) before any commutator analysis.  Therefore \(N=4\) is the smallest
noncommuting dimension not already eliminated by parity.  This corrected
statement is the acquisition premise for V112.
\end{assessment}

\paragraph{Next forced step.}
V112 should classify or sharply test \(N=4\) pencils satisfying the uniform
null-circle Pfaffian gate.  It must retain the complete determinant on the
complex stable hemisphere.  The quaternionic V108 pencil is a required
regression because it passes the Pfaffian gate but fails globally.  If the
remaining four-dimensional class cannot be closed without an unstructured
coefficient search, the programme should go upstream to the exact nonlocal
Calder\'on graph rather than add further anonymous collar fields.

\begin{assessment}[V111 continuum status]
\label{ass:V111-continuum-status}
The full Standard--Model--Einstein continuum remains unproved.  V111 corrects
the noncommuting dimension threshold, rules out every odd-dimensional finite
real action pencil, and gives the exact even-dimensional Pfaffian glancing
gate.  A full even-dimensional or nonlocal gauge realization, the fermion
branch selection, common quantum coefficients and form bounds, curved
radiative equivariance, nonlinear existence, and the cofinal interacting
limit remain open.
\end{assessment}

\section*{V111 exact checks and append-only record}
\addcontentsline{toc}{section}{V111 exact checks and append-only record}

The exact parity and Pfaffian checks are
\begin{equation}
 \boxed{
 N\text{ odd}\Longrightarrow
 \det G_\varepsilon(e)=0,
 \qquad
 N=2m\Longrightarrow
 \det G_\varepsilon(e)
 =\bigl(\operatorname{Pf}G_\varepsilon(e)\bigr)^2.}
 \label{eq:V111-final-check}
\end{equation}
The append operation preserves the complete V110 prefix byte for byte before
its former live terminator.  The assembled V111 file is checked for one live
final terminator, balanced environments and braces, duplicate labels, newly
introduced unresolved references, display math delimiters, and control
characters before the exact V110 predecessor is removed.

% =============================================================================
% END APPENDED VERSION V111 MATERIAL
% =============================================================================

\clearpage
% =============================================================================
% BEGIN APPENDED VERSION V112 MATERIAL
% =============================================================================
\section{V112: complete stable classification of chiral four-collar
quaternionic pencils}
\label{sec:V112-chiral-four}

V111 proved that four is the smallest collar dimension which is both capable
of noncommuting skew coefficients and not already excluded by odd parity.  A
distinguished four-dimensional class is one chiral
\(\mathfrak{so}(4)\) summand.  V108 tested its isotropic orthonormal
coefficient triple.  V112 allows arbitrary real coefficient vectors, including
degenerate and anisotropic triples, and classifies the complete stable
determinant.

Every chiral pencil fails.  For each chosen spatial direction, one scalar
discriminant decides the type of failure.  A positive discriminant gives a
normalized growing mode.  A nonpositive discriminant gives a normalized
imaginary-frequency kernel, and a sharp Cauchy argument places its roots
inside the closed stable axis.  Thus neither anisotropy nor degeneracy within
one quaternionic chirality repairs the gauge/ghost floor.

\subsection{The chiral quaternionic action}
\label{subsec:V112-chiral-action}

Let \(E_1,E_2,E_3\in M_4(\mathbb R)\) be real skew matrices satisfying
\begin{equation}
 E_iE_j=-\delta_{ij}I_4+\epsilon_{ijk}E_k.
 \label{eq:V112-quaternion-relations}
\end{equation}
The matrices printed in V108, after relabelling
\((E_0,E_1,E_2)\) as \((E_1,E_2,E_3)\), give one such realization.  Their
span is one chiral \(\mathfrak{su}(2)\) summand of
\(\mathfrak{so}(4)\).

Choose arbitrary vectors
\begin{equation}
 u^0,u^1,u^2\in\mathbb R^3
 \label{eq:V112-coefficient-vectors}
\end{equation}
and define
\begin{equation}
 A^a:=u^a\mathbin{\cdot}E
 :=\sum_{\ell=1}^3u^a_\ell E_\ell,
 \qquad a=0,1,2.
 \label{eq:V112-chiral-coefficients}
\end{equation}
The real boundary action and its stationary row are
\begin{align}
 S_\partial[q]
 &:=
 \frac12\int_{\mathcal T}
 q^T(A^0D_T+A^1D_1+A^2D_2)q,
 \label{eq:V112-chiral-action}\\
 p+(A^0D_T+A^1D_1+A^2D_2)q&=0.
 \label{eq:V112-chiral-row}
\end{align}
Every coefficient is real skew, so the graph is action-Lagrangian by V109.

\subsection{Exact determinant reduction}
\label{subsec:V112-determinant}

For \(e=(e_1,e_2)\in\mathbb S^1\), put
\begin{equation}
 u(e):=e_1u^1+e_2u^2.
 \label{eq:V112-spatial-vector}
\end{equation}
Restrict to \(k=\kappa e\), \(\kappa>0\), and set
\begin{equation}
 z(s,\kappa,e):=s u^0+i\kappa u(e)\in\mathbb C^3.
 \label{eq:V112-complex-vector}
\end{equation}
The skew part of the stable matrix is
\begin{equation}
 S(s,\kappa,e)=z(s,\kappa,e)\mathbin{\cdot}E.
 \label{eq:V112-skew-symbol}
\end{equation}

\begin{lemma}[Quaternionic square]
\label{lem:V112-quaternion-square}
For every \(z\in\mathbb C^3\),
\begin{equation}
 (z\mathbin{\cdot}E)^2=-(z\mathbin{\cdot}z)I_4,
 \label{eq:V112-quaternion-square}
\end{equation}
where the dot on the right is complex bilinear.
\end{lemma}

\begin{proof}
Expand the square and use
\eqref{eq:V112-quaternion-relations}.  The symmetric coefficient
\(z_iz_j\) kills the antisymmetric \(\epsilon_{ijk}\) term, while the
\(-\delta_{ij}\) term gives \(-\sum_i z_i^2I_4\).
\end{proof}

\begin{theorem}[Exact chiral stable determinant]
\label{thm:V112-exact-determinant}
Let
\begin{equation}
 \mathcal B_u(s,\kappa,e)
 :=\lambda I_4+S(s,\kappa,e),
 \qquad
 \lambda^2=s^2+\kappa^2.
 \label{eq:V112-stable-matrix}
\end{equation}
Then
\begin{equation}
 \boxed{
 \det\mathcal B_u
 =\left(\lambda^2+z\mathbin{\cdot}z\right)^2.}
 \label{eq:V112-determinant-square}
\end{equation}
Writing
\begin{equation}
 a:=|u^0|^2,
 \qquad
 b(e):=u^0\mathbin{\cdot}u(e),
 \qquad
 c(e):=|u(e)|^2,
 \label{eq:V112-abc}
\end{equation}
and \(w=s/\kappa\), the scalar factor is
\begin{equation}
 P_e(w)
 :=(1+a)w^2+2ib(e)w+1-c(e).
 \label{eq:V112-directional-quadratic}
\end{equation}
\end{theorem}

\begin{proof}
Lemma~\ref{lem:V112-quaternion-square} implies that the two eigenvalues of
\(S\) are the opposite square roots of
\(-z\mathbin{\cdot}z\), each with multiplicity two.  Therefore
\[
 \det(\lambda I_4+S)
 =(\lambda^2+z\mathbin{\cdot}z)^2.
\]
Next,
\[
 z\mathbin{\cdot}z
 =a s^2+2ib(e)s\kappa-c(e)\kappa^2.
\]
Add \(\lambda^2=s^2+\kappa^2\), divide by \(\kappa^2\), and obtain
\eqref{eq:V112-directional-quadratic}.
\end{proof}

\subsection{Directional discriminant and stable roots}
\label{subsec:V112-discriminant}

Define
\begin{equation}
 \Delta(e):=(1+a)(c(e)-1)-b(e)^2.
 \label{eq:V112-discriminant}
\end{equation}
The two roots of \(P_e\) are
\begin{equation}
 w_\pm(e)
 =\frac{-ib(e)\pm\sqrt{\Delta(e)}}{1+a}.
 \label{eq:V112-root-formula}
\end{equation}

\begin{lemma}[Nonpositive-discriminant roots stay in the stable interval]
\label{lem:V112-root-bound}
If \(\Delta(e)\le0\), set
\begin{equation}
 R(e):=\sqrt{-\Delta(e)},
 \qquad
 r_\pm(e):=\frac{-b(e)\pm R(e)}{1+a}.
 \label{eq:V112-imaginary-roots}
\end{equation}
Then \(w_\pm=ir_\pm\) and
\begin{equation}
 |r_\pm(e)|\le1.
 \label{eq:V112-stable-r-bound}
\end{equation}
\end{lemma}

\begin{proof}
The root identity is immediate.  Since
\(\Delta(e)\le0\),
\begin{equation}
 (1+a)(c-1)\le b^2\le ac,
 \label{eq:V112-two-Cauchy-bounds}
\end{equation}
where the second inequality is Cauchy--Schwarz.  Subtracting the outer terms
gives \(c\le1+a\).  Hence
\(|b|\le\sqrt{ac}\le1+a\).  Moreover,
\begin{align}
 (1+a-|b|)^2-R^2
 &=(1+a)(a+c-2|b|)\notag\\
 &\ge0,
 \label{eq:V112-root-bound-square}
\end{align}
because \(2|b|\le a+c\).  Thus
\(|b|+R\le1+a\), which implies
\(|-b\pm R|\le1+a\) and proves
\eqref{eq:V112-stable-r-bound}.
\end{proof}

\begin{theorem}[Complete chiral four-collar no-go]
\label{thm:V112-chiral-no-go}
For every coefficient triple
\((u^0,u^1,u^2)\in(\mathbb R^3)^3\),
\begin{equation}
 \boxed{
 \inf_{\overline{\mathbb S}_+}
 s_{\min}\mathcal B_u=0.}
 \label{eq:V112-zero-floor}
\end{equation}
Indeed, fix any \(e\in\mathbb S^1\).
\begin{enumerate}[label=\textup{(\roman*)},leftmargin=3.3em]
\item If \(\Delta(e)\le0\), choose either \(r=r_\pm(e)\) and set
      \begin{equation}
       \kappa=(1+r^2)^{-1/2},
       \qquad
       s=ir\kappa,
       \qquad
       k=\kappa e.
       \label{eq:V112-normalized-boundary-zero}
      \end{equation}
      Then the frequency is normalized, belongs to the closed stable
      imaginary axis, and \(\det\mathcal B_u=0\).
\item If \(\Delta(e)>0\), choose
      \begin{equation}
       w_+(e):=\frac{\sqrt{\Delta(e)}-ib(e)}{1+a},
       \qquad
       \kappa=(1+|w_+(e)|^2)^{-1/2},
       \qquad
       s=\kappa w_+(e),
       \qquad
       k=\kappa e.
       \label{eq:V112-normalized-growing-zero}
      \end{equation}
      Then the frequency is normalized,
      \(\operatorname{Re}s>0\), and
      \(\det\mathcal B_u=0\).  This is an exponentially growing stable mode.
\end{enumerate}
\end{theorem}

\begin{proof}
If \(\Delta\le0\), Lemma~\ref{lem:V112-root-bound} gives
\(|r|\le1\).  The stable root at \(s=ir\kappa\) is
\(\lambda=\kappa\sqrt{1-r^2}\), with the limiting value zero allowed at
glancing.  Equation \eqref{eq:V112-normalized-boundary-zero} gives
\(|s|^2+|k|^2=1\), and \(s/\kappa=ir\) is a root of
\eqref{eq:V112-directional-quadratic}.  The determinant square in
\eqref{eq:V112-determinant-square} therefore vanishes.

If \(\Delta>0\), \(w_+\) in
\eqref{eq:V112-normalized-growing-zero} is a root of the same quadratic and
has positive real part.  Homogeneity gives a normalized determinant zero with
\(\operatorname{Re}s>0\).  These cases exhaust every real value of
\(\Delta(e)\), proving \eqref{eq:V112-zero-floor}.
\end{proof}

\begin{corollary}[Anisotropy inside one chirality cannot help]
\label{cor:V112-anisotropy}
Theorem~\ref{thm:V112-chiral-no-go} assumes no orthogonality, equal norm,
linear independence, or rotational symmetry among \(u^0,u^1,u^2\).  It
therefore contains the V108 isotropic Clifford pencil, every rank-deficient
chiral pencil, and all their anisotropic deformations.
\end{corollary}

\begin{proof}
All coefficient dependence enters only through the exact Gram scalars
\(a,b(e),c(e)\), with no restrictions beyond their definitions.
\end{proof}

\subsection{Relation to the Pfaffian glancing gate}
\label{subsec:V112-Pfaffian}

At the normalized glancing point of V111,
\begin{equation}
 z=\frac{i}{\sqrt2}\bigl(\varepsilon u^0+u(e)\bigr).
 \label{eq:V112-glancing-z}
\end{equation}

\begin{proposition}[Exact chiral Pfaffian]
\label{prop:V112-Pfaffian}
Up to the fixed orientation sign of the chosen quaternionic basis,
\begin{equation}
 \left|
 \operatorname{Pf}G_\varepsilon(e)
 \right|
 =
 \left|\varepsilon u^0+u(e)\right|^2.
 \label{eq:V112-Pfaffian-value}
\end{equation}
Thus the V111 uniform glancing gate holds precisely when
\begin{equation}
 \inf_{\varepsilon,e}
 |\varepsilon u^0+u(e)|>0.
 \label{eq:V112-chiral-glancing-gate}
\end{equation}
Even on that branch, Theorem~\ref{thm:V112-chiral-no-go} supplies a stable
zero away from glancing whenever the glancing zero itself is absent.
\end{proposition}

\begin{proof}
For \(x\in\mathbb R^3\), the quaternionic matrix
\(x\mathbin{\cdot}E\) squares to \(-|x|^2I_4\), so its determinant is
\(|x|^4\).  The Pfaffian square identity gives absolute Pfaffian
\(|x|^2\).  Set \(x=\varepsilon u^0+u(e)\).  Compactness gives the uniform
criterion.  The final assertion is Theorem~\ref{thm:V112-chiral-no-go}.
\end{proof}

\subsection{BRST incidence and the remaining four-dimensional class}
\label{subsec:V112-BRST}

\begin{proposition}[Chiral BRST direct sum]
\label{prop:V112-BRST}
Apply the same chiral action row to matched gauge, ghost, antighost, and
multiplier columns.  The resulting domain is action-Lagrangian and a frozen
principal BRST subcomplex, and every differential block has determinant
\eqref{eq:V112-determinant-square}.  Hence the kernel in
Theorem~\ref{thm:V112-chiral-no-go} embeds in the full Standard Model adjoint
direct sum.
\end{proposition}

\begin{proof}
Every \(A^a\) is real skew, giving action incidence.  Constant coefficients
commute with the frozen BRST derivatives, giving the matched quartet chain as
in Proposition~\ref{prop:V111-BRST}.  Direct-sum embedding preserves a kernel
supported in one real adjoint component.
\end{proof}

\begin{firewall}[The nonchiral \(N=4\) class remains]
\label{fire:V112-nonchiral-open}
Every real four-by-four skew matrix decomposes into commuting self-dual and
anti-self-dual parts.  V112 treats pencils whose three coefficient matrices
all lie in one fixed chiral summand.  A general pencil mixes both summands;
its determinant contains both chiral vector laws and is not the single square
\eqref{eq:V112-determinant-square}.  No conclusion about that mixed-chirality
class is asserted here.
\end{firewall}

\subsection{V112 ledger and next acquisition}
\label{subsec:V112-ledger}

\paragraph{New exact conclusions.}
\begin{enumerate}[label=\textup{(\arabic*)},leftmargin=2.8em]
\item Lemma~\ref{lem:V112-quaternion-square} reduces every chiral
      four-dimensional skew symbol to one complex bilinear scalar.
\item Theorem~\ref{thm:V112-exact-determinant} gives the complete stable
      determinant for arbitrary anisotropic coefficient vectors.
\item Lemma~\ref{lem:V112-root-bound} proves that every
      nonpositive-discriminant root lies on the closed stable imaginary axis.
\item Theorem~\ref{thm:V112-chiral-no-go} produces either that boundary
      kernel or an open-half-plane growing kernel for every coefficient
      triple.
\item Proposition~\ref{prop:V112-Pfaffian} evaluates the V111 glancing gate
      exactly and shows why passing it is insufficient.
\item Proposition~\ref{prop:V112-BRST} embeds the same failure in the
      action/BRST direct sum.
\end{enumerate}

\paragraph{Next forced step.}
V113 should use the commuting
\(\mathfrak{su}(2)_+\oplus\mathfrak{su}(2)_-\) decomposition of
\(\mathfrak{so}(4)\) to derive the exact mixed-chirality stable determinant.
It should retain both chiral vector laws in the same polynomial and decide
whether their interaction can remove all stable roots.  If not, the complete
four-collar local class is closed.  A surviving algebraic pencil must still
receive a physical boundary-field, stress, and ghost interpretation.

\begin{assessment}[V112 continuum status]
\label{ass:V112-continuum-status}
The full Standard--Model--Einstein continuum remains unproved.  V112 rules
out every chiral four-collar real action pencil, including all anisotropic
deformations of the V108 Clifford model.  The mixed-chirality four-collar or
nonlocal gauge realization, the fermion branch selection, common quantum
coefficients and form bounds, curved radiative equivariance, nonlinear
existence, and the cofinal interacting limit remain open.
\end{assessment}

\section*{V112 exact checks and append-only record}
\addcontentsline{toc}{section}{V112 exact checks and append-only record}

The exact chiral checksum is
\begin{equation}
 \boxed{
 \det\mathcal B_u
 =\kappa^4P_e(w)^2,
 \qquad
 P_e(w)=(1+a)w^2+2ib(e)w+1-c(e),
 \qquad
 \Delta(e)=(1+a)(c(e)-1)-b(e)^2.}
 \label{eq:V112-final-check}
\end{equation}
The append operation preserves the complete V111 prefix byte for byte before
its former live terminator.  The assembled V112 file is checked for one live
final terminator, balanced environments and braces, duplicate labels, newly
introduced unresolved references, display math delimiters, and control
characters before the exact V111 predecessor is removed.

% =============================================================================
% END APPENDED VERSION V112 MATERIAL
% =============================================================================

\clearpage
% =============================================================================
% BEGIN APPENDED VERSION V113 MATERIAL
% =============================================================================
\section{V113: polynomial root symmetry and the finite conservative
conormal-graph no-go theorem}
\label{sec:V113-polynomial-no-go}

V111 excluded every odd collar dimension by glancing parity.  V112 excluded
one full chiral half of the four-dimensional noncommuting class by an explicit
quadratic discriminant.  The remaining mixed-chirality calculation can be
resolved without expanding its quartic coefficients.  In every even
dimension, skew symmetry makes the characteristic polynomial even in the
normal root.  After substituting the bulk dispersion relation, the complete
stable determinant is a nonconstant polynomial in the Laplace frequency.  A
reality symmetry reflects its roots across the imaginary axis.  Therefore at
least one root lies in the closed right half-plane.

Together with V111, this proves a no-go theorem for every finite real
first-order conservative conormal graph in the declared class.  It does not
cover Lagrangian relations with Dirichlet configuration subspaces, boundary
fields with a different bulk dispersion, or nonlocal principal symbols.

\subsection{The declared finite conormal-graph class}
\label{subsec:V113-declared-class}

Let \(N\ge1\), let \(A^0,A^1,A^2\in M_N(\mathbb R)\) be real skew, and let
\begin{equation}
 A(e):=e_1A^1+e_2A^2,
 \qquad e\in\mathbb S^1.
 \label{eq:V113-spatial-pencil}
\end{equation}
The local quadratic boundary action fixes the full configuration trace
\(q\in\mathbb R^N\) freely and imposes the conormal graph
\begin{equation}
 p+(A^0D_T+A^1D_1+A^2D_2)q=0.
 \label{eq:V113-conormal-graph}
\end{equation}
For the fixed nonzero spatial frequency \(k=e\), define
\begin{align}
 S_e(s)&:=sA^0+iA(e),
 \label{eq:V113-Se}\\
 \mathcal B_e(s)&:=\lambda(s)I_N+S_e(s),
 \qquad
 \lambda(s)^2=s^2+1,
 \label{eq:V113-Be}
\end{align}
where \(\lambda\) is the stable square root continued from
\(\operatorname{Re}s>0\).

\begin{lemma}[Even characteristic polynomial of a skew matrix]
\label{lem:V113-even-characteristic}
If \(N\) is even and \(S^T=-S\), then
\begin{equation}
 \det(\lambda I_N+S)=\det(-\lambda I_N+S)
 =\det(\lambda I_N-S).
 \label{eq:V113-even-characteristic}
\end{equation}
Consequently \(\det(\lambda I_N+S)\) is a polynomial in \(\lambda^2\).
\end{lemma}

\begin{proof}
Since \(N\) is even,
\[
 \det(-\lambda I_N+S)
 =\det(\lambda I_N-S).
\]
Transposition and \(S^T=-S\) give
\[
 \det(\lambda I_N-S)
 =\det((\lambda I_N-S)^T)
 =\det(\lambda I_N+S).
\]
Thus the characteristic polynomial is even in \(\lambda\).
\end{proof}

\begin{proposition}[The stable determinant is a genuine polynomial]
\label{prop:V113-stable-polynomial}
Assume \(N\) is even.  There is a polynomial \(P_e\in\mathbb C[s]\) such that
\begin{equation}
 P_e(s)=\det\mathcal B_e(s)
 \label{eq:V113-polynomial-definition}
\end{equation}
for either square-root branch.  It has degree exactly \(N\), with positive
leading coefficient
\begin{equation}
 [s^N]P_e(s)=\det(I_N+A^0)>0.
 \label{eq:V113-leading-coefficient}
\end{equation}
\end{proposition}

\begin{proof}
Lemma~\ref{lem:V113-even-characteristic} expresses the determinant as a
polynomial in \(\lambda^2\) whose coefficients are polynomial functions of
the entries of \(S_e(s)\).  Substitution of
\(\lambda^2=s^2+1\) therefore gives a polynomial \(P_e(s)\), independent of
the sign of \(\lambda\).

Along positive real \(s\to+\infty\), the stable root satisfies
\(\lambda(s)/s\to1\), while
\[
 \frac{\mathcal B_e(s)}s
 \longrightarrow I_N+A^0.
\]
Hence
\[
 \frac{P_e(s)}{s^N}\longrightarrow\det(I_N+A^0).
\]
A real skew matrix has eigenvalues \(0\) and pairs
\(\pm i\alpha_j\), so
\[
 \det(I_N+A^0)=\prod_j(1+\alpha_j^2)>0.
\]
This proves the degree and leading-coefficient claim.
\end{proof}

\subsection{Root reflection and the right-half-plane alternative}
\label{subsec:V113-root-reflection}

\begin{lemma}[Imaginary-axis reflection symmetry]
\label{lem:V113-reflection}
The polynomial in Proposition~\ref{prop:V113-stable-polynomial} satisfies
\begin{equation}
 P_e(-\overline s)=\overline{P_e(s)}.
 \label{eq:V113-reflection}
\end{equation}
Thus every root \(s_0\) is accompanied by the root
\(-\overline{s_0}\).
\end{lemma}

\begin{proof}
One has
\begin{equation}
 S_e(-\overline s)=-\overline{S_e(s)}.
 \label{eq:V113-S-reflection}
\end{equation}
Choose at \(-\overline s\) the square root
\(\lambda'=\overline{\lambda(s)}\); the determinant is independent of the
root sign by Lemma~\ref{lem:V113-even-characteristic}.  Then
\begin{align}
 P_e(-\overline s)
 &=\det\bigl(\overline\lambda I_N-\overline{S_e(s)}\bigr)\notag\\
 &=\det\bigl(\overline\lambda I_N+\overline{S_e(s)}\bigr)\notag\\
 &=\overline{\det\bigl(\lambda I_N+S_e(s)\bigr)}.
 \label{eq:V113-reflection-proof}
\end{align}
The second equality again uses evenness under \(S\mapsto-S\), which follows
from transposition.  This proves \eqref{eq:V113-reflection}.
\end{proof}

\begin{theorem}[Every even-dimensional finite conormal graph has a stable
zero]
\label{thm:V113-even-no-go}
Let \(N\) be even.  For every real skew coefficient triple
\((A^0,A^1,A^2)\), there is a finite \(s_*\in\mathbb C\) with
\begin{equation}
 \operatorname{Re}s_*\ge0,
 \qquad
 P_e(s_*)=0.
 \label{eq:V113-right-root}
\end{equation}
After setting
\begin{equation}
 t_*:=(1+|s_*|^2)^{-1/2},
 \qquad
 \widehat s_*=t_*s_*,
 \qquad
 \widehat k_*=t_*e,
 \label{eq:V113-normalization}
\end{equation}
the normalized stable boundary matrix has a nonzero kernel.  If
\(\operatorname{Re}s_*>0\), this is an exponentially growing mode.  If
\(\operatorname{Re}s_*=0\), it is a boundary-axis mode and is the limit of
open-half-plane stable quasimodes.  In either case,
\begin{equation}
 \boxed{
 \inf_{\overline{\mathbb S}_+}
 s_{\min}\mathcal B_A=0.}
 \label{eq:V113-even-zero-floor}
\end{equation}
\end{theorem}

\begin{proof}
Proposition~\ref{prop:V113-stable-polynomial} shows that \(P_e\) is a
nonconstant degree-\(N\) polynomial.  The fundamental theorem of algebra
gives a root \(s_0\).  If \(\operatorname{Re}s_0\ge0\), take \(s_*=s_0\).
Otherwise Lemma~\ref{lem:V113-reflection} makes
\(-\overline{s_0}\) a root with positive real part; take that root.

The row \eqref{eq:V113-conormal-graph} is homogeneous of degree one, so
scaling \((s,k,\lambda)\) by \(t_*\) scales the stable matrix by \(t_*\) and
preserves its kernel.  Equation \eqref{eq:V113-normalization} gives
\(|\widehat s_*|^2+|\widehat k_*|^2=1\).
For a boundary-axis root, choose \(s_j\to s_*\) from
\(\operatorname{Re}s_j>0\), use the stable square-root continuation, and take
unit right singular vectors.  Continuity makes their image norms tend to
zero.  This proves \eqref{eq:V113-even-zero-floor}.
\end{proof}

\begin{theorem}[Finite conservative conormal-graph no-go]
\label{thm:V113-all-finite-no-go}
For every finite \(N\ge1\), no real local first-order boundary action whose
stationary domain is the full conormal graph
\eqref{eq:V113-conormal-graph} has a positive unweighted stable-root floor.
\end{theorem}

\begin{proof}
If \(N\) is odd, apply Theorem~\ref{thm:V111-odd-no-go}; its kernel lies at
bulk glancing.  If \(N\) is even, apply
Theorem~\ref{thm:V113-even-no-go}.  These alternatives exhaust all finite
dimensions.
\end{proof}

\begin{corollary}[The mixed-chirality four-collar class is closed]
\label{cor:V113-four-closed}
Every real four-collar pencil, including arbitrary mixtures of the
self-dual and anti-self-dual \(\mathfrak{so}(4)\) summands, has a stable
boundary zero.  The chiral discriminant in V112 gives its location for one
large subclass; Theorem~\ref{thm:V113-even-no-go} supplies existence for the
complete class.
\end{corollary}

\begin{proof}
Set \(N=4\) in Theorem~\ref{thm:V113-even-no-go}.  No chirality assumption
occurs there.
\end{proof}

\subsection{BRST, action, and the exact scope}
\label{subsec:V113-BRST-scope}

\begin{proposition}[The finite no-go survives the matched BRST construction]
\label{prop:V113-BRST}
Apply \eqref{eq:V113-conormal-graph} to finite gauge-collar, ghost,
antighost, and multiplier columns with the matched frozen BRST differential.
The domain is action-Lagrangian and BRST-closed, but every differential block
contains the stable zero of Theorem~\ref{thm:V113-all-finite-no-go}.  Hence
the full Standard Model adjoint direct sum has zero unweighted floor.
\end{proposition}

\begin{proof}
Real skew coefficients multiplying formally skew tangential derivatives give
the Lagrangian action graph.  Constant coefficients commute with frozen BRST
derivatives, so the quartet closes as in V109--V112.  The gauge and ghost
members have the same stable matrix, and a direct sum preserves a kernel
supported in one member.
\end{proof}

\begin{firewall}[The theorem is about full conormal graphs]
\label{fire:V113-scope}
Theorem~\ref{thm:V113-all-finite-no-go} assumes that every collar
configuration trace is free and that the action determines all conormals by a
finite local first-order tangential graph.  A general Lagrangian boundary
relation may instead contain Dirichlet configuration directions.  The V35
relative branch is of that mixed type; its remaining normal graph is
one-dimensional and is already excluded by glancing.  A changed mixed
Dirichlet graph, a field with different normal dispersion, an
order-changing Wentzell system, or a nonlocal Calder\'on relation is not
covered and must receive its own BRST and stable tests.
\end{firewall}

\begin{assessment}[Upstream decision after the finite graph theorem]
\label{ass:V113-upstream-decision}
Increasing the number of anonymous wave-collar reservoirs within the declared
conormal-graph class is now exhausted.  The next gauge-domain candidate must
change a structural datum: introduce and justify a BRST-compatible Dirichlet
subspace, change the reservoir dispersion/order, use a nonlocal
Dirichlet-to-Neumann graph, or abandon the timelike boundary realization.
Repeating another finite skew coefficient search would fall under
Theorem~\ref{thm:V113-all-finite-no-go}.
\end{assessment}

\subsection{V113 ledger and next acquisition}
\label{subsec:V113-ledger}

\paragraph{New exact conclusions.}
\begin{enumerate}[label=\textup{(\arabic*)},leftmargin=2.8em]
\item Lemma~\ref{lem:V113-even-characteristic} proves exact evenness in the
      normal root for every even-dimensional complex skew symbol.
\item Proposition~\ref{prop:V113-stable-polynomial} turns the complete stable
      determinant into a degree-\(N\) polynomial with positive leading
      coefficient.
\item Lemma~\ref{lem:V113-reflection} proves reflection of every root across
      the imaginary axis.
\item Theorem~\ref{thm:V113-even-no-go} converts those facts into a normalized
      right-half-plane or boundary-axis stable kernel.
\item Theorem~\ref{thm:V113-all-finite-no-go} combines the result with V111
      and closes every finite full conormal graph in the declared action
      class.
\item Proposition~\ref{prop:V113-BRST} embeds the obstruction into the matched
      gauge/ghost direct sum.
\end{enumerate}

\paragraph{Next forced step.}
V114 should go upstream to a mixed Dirichlet--conormal Lagrangian relation and
derive the BRST incidence constraint on its Dirichlet subspace.  It should
decide whether the normal ghost jet forces at least one free conormal graph
direction, which would reinsert the V113 obstruction, or identify a genuine
BRST boundary multiplet which closes that direction.  If locality again
forces a kernel, V115 must audit the companions and start the exact nonlocal
Calder\'on branch.

\begin{assessment}[V113 continuum status]
\label{ass:V113-continuum-status}
The full Standard--Model--Einstein continuum remains unproved.  V113 closes
every finite real first-order full conormal-graph dilation and therefore ends
the anonymous multi-collar search in that class.  Mixed Dirichlet, changed
dispersion, or nonlocal gauge domains, the fermion branch selection, common
quantum coefficients and form bounds, curved radiative equivariance,
nonlinear existence, and the cofinal interacting limit remain open.
\end{assessment}

\section*{V113 exact checks and append-only record}
\addcontentsline{toc}{section}{V113 exact checks and append-only record}

The decisive polynomial checksum is
\begin{equation}
 \boxed{
 P_e(s)=\det(\lambda I_N+sA^0+iA(e)),
 \quad
 \deg P_e=N,
 \quad
 [s^N]P_e=\det(I_N+A^0)>0,
 \quad
 P_e(-\overline s)=\overline{P_e(s)}.}
 \label{eq:V113-final-check}
\end{equation}
The append operation preserves the complete V112 prefix byte for byte before
its former live terminator.  The assembled V113 file is checked for one live
final terminator, balanced environments and braces, duplicate labels, newly
introduced unresolved references, display math delimiters, and control
characters before the exact V112 predecessor is removed.

% =============================================================================
% END APPENDED VERSION V113 MATERIAL
% =============================================================================

\clearpage
% =============================================================================
% BEGIN APPENDED VERSION V114 MATERIAL
% =============================================================================
\section{V114: mixed Dirichlet--conormal action relations and the free normal
BRST direction}
\label{sec:V114-mixed-Dirichlet}

V113 closed every finite local first-order action domain whose entire
configuration trace is free and whose conormal is fixed by a tangential
graph.  A general separated self-adjoint boundary relation may instead impose
Dirichlet conditions on part of the configuration fibre and a conormal graph
only on the complementary part.  V114 tests that escape while retaining the
V35 relative Dirichlet ghost domain.

The normal ghost jet is the decisive datum.  Dirichlet ghost trace kills all
tangential ghost derivatives but leaves \(D_nc\) arbitrary.  Therefore a
BRST-preserved gauge Dirichlet subspace cannot contain the normal gauge
configuration direction.  That direction must remain in the free conormal
sector.  The sector is nonempty, and the V113 polynomial theorem forces its
uniform stable floor to zero.  Thus regular mixed Dirichlet conditions do not
repair the local action branch unless the ghost domain or boundary BRST field
content is changed.

\subsection{Regular mixed Lagrangian relations}
\label{subsec:V114-regular-relation}

Let \(Q\) be a finite-dimensional real configuration fibre and use the
boundary phase space \(Q\oplus Q\) with Green form
\begin{equation}
 \omega((q,p),(v,r)):=\langle p,v\rangle-\langle q,r\rangle.
 \label{eq:V114-Green-form}
\end{equation}
Choose an orthogonal splitting
\begin{equation}
 Q=D\oplus F,
 \qquad
 P_D+P_F=I_Q,
 \label{eq:V114-DF-splitting}
\end{equation}
where \(D\) is the Dirichlet configuration subspace and \(F\) is the free
configuration subspace.  Let \(R\) be a formally self-adjoint tangential
operator on \(F\).  Define
\begin{equation}
 \mathcal L_{D,R}
 :=
 \left\{
 (q,p):
 P_Dq=0,\quad
 P_Fp+R P_Fq=0
 \right\}.
 \label{eq:V114-mixed-relation}
\end{equation}
The component \(P_Dp\) is unrestricted.

\begin{proposition}[Exact Lagrangian criterion for the regular split]
\label{prop:V114-Lagrangian}
The relation \(\mathcal L_{D,R}\) is Lagrangian for
\eqref{eq:V114-Green-form} if and only if \(R=R^\dagger\) on \(F\).
\end{proposition}

\begin{proof}
For two elements of \(\mathcal L_{D,R}\), the configuration traces lie in
\(F\), so the unrestricted \(D\)-momenta pair with zero.  The Green form
reduces to
\[
 \omega
 =\langle-Rq,v\rangle-\langle q,-Rv\rangle
 =\langle q,(R^\dagger-R)v\rangle.
\]
It vanishes for all \(q,v\in F\) exactly when \(R=R^\dagger\).
The relation has
\(\dim D\) free momentum coordinates and \(\dim F\) free configuration
coordinates, hence dimension \(\dim Q\).  Isotropy is therefore maximal.
\end{proof}

At real first tangential order, every local quadratic action on \(F\) has
\begin{equation}
 R=A^0D_T+A^1D_1+A^2D_2,
 \qquad
 (A^a)^T=-A^a,
 \label{eq:V114-R-pencil}
\end{equation}
up to lower-order symmetric terms, by the variation calculation in V107.

\subsection{Stable block decomposition}
\label{subsec:V114-stable-block}

For identical wave principal symbols on the configuration components, a
stable amplitude satisfies \(p=\lambda q\).  In the orthogonal coordinates
\(Q=D\oplus F\), the boundary matrix of
\eqref{eq:V114-mixed-relation} is
\begin{equation}
 \mathcal B_{D,R}(s,k)
 =
 \begin{pmatrix}
 I_D&0\\
 0&\lambda I_F+
       sA^0+ik_1A^1+ik_2A^2
 \end{pmatrix}.
 \label{eq:V114-block-stable-matrix}
\end{equation}

\begin{lemma}[Exact floor of the mixed relation]
\label{lem:V114-mixed-floor}
If \(F\ne\{0\}\), then
\begin{equation}
 s_{\min}\mathcal B_{D,R}
 =
 \min\left\{
 1,\,
 s_{\min}\bigl(
 \lambda I_F+sA^0+ik_1A^1+ik_2A^2
 \bigr)
 \right\}.
 \label{eq:V114-mixed-floor}
\end{equation}
Consequently every finite first-order relation
\(\mathcal L_{D,R}\) with \(F\ne\{0\}\) has zero uniform stable floor.
\end{lemma}

\begin{proof}
The first equality is the singular-value formula for an orthogonal direct
sum.  The free block is exactly the finite full conormal graph classified in
Theorem~\ref{thm:V113-all-finite-no-go}.  Its infimum singular value is zero,
so the same is true of the minimum.
\end{proof}

\begin{remark}[Why pure Dirichlet would pass only the stable matrix]
\label{rem:V114-pure-Dirichlet}
If \(F=\{0\}\), then the stable boundary matrix is \(I_D\) and has unit floor.
Whether this pure Dirichlet relation is a gauge/BRST domain is a separate
incidence question answered below.
\end{remark}

\subsection{BRST preservation forces a free normal direction}
\label{subsec:V114-BRST-normal}

Return to one real Yang--Mills adjoint component.  Let \(V_{\rm pot}\) be the
four-dimensional gauge-potential configuration fibre and let
\(n^\flat\in V_{\rm pot}\) denote the boundary-normal covector.  Extra collar
copies, if present, are included in \(Q\); the original normal gauge vector is
denoted
\begin{equation}
 \nu_X:=n^\flat X\in Q
 \label{eq:V114-normal-gauge-vector}
\end{equation}
for an adjoint amplitude \(X\).

Keep the relative ghost trace
\begin{equation}
 c|_{\mathcal T}=0.
 \label{eq:V114-Dirichlet-ghost}
\end{equation}
At frozen principal order,
\begin{equation}
 \mathsf s a_\mu=D_\mu c.
 \label{eq:V114-gauge-BRST}
\end{equation}
By Lemma~\ref{lem:V104-free-normal-ghost}, the boundary jet
\(\rho=(D_nc)|_{\mathcal T}\) is arbitrary on the ghost Euler equation, while
all tangential derivatives of the Dirichlet trace vanish.

\begin{theorem}[BRST incidence forces the normal gauge vector into \(F\)]
\label{thm:V114-normal-free}
Suppose the gauge Dirichlet row
\begin{equation}
 P_Da|_{\mathcal T}=0
 \label{eq:V114-gauge-Dirichlet-row}
\end{equation}
is preserved by the linearized BRST differential for every ghost Euler jet
satisfying \eqref{eq:V114-Dirichlet-ghost}.  Then
\begin{equation}
 P_D\nu_X=0
 \quad\text{for every adjoint amplitude }X,
 \qquad\text{hence}\qquad
 \nu_X\in F.
 \label{eq:V114-normal-in-F}
\end{equation}
In particular \(F\ne\{0\}\).
\end{theorem}

\begin{proof}
Apply \(\mathsf s\) to \eqref{eq:V114-gauge-Dirichlet-row}.  The tangential
components of \(Dc\) vanish because the boundary trace of \(c\) is zero.  The
only surviving principal jet is its normal component, so
\begin{equation}
 \mathsf s(P_Da)
 =P_D(n^\flat\rho).
 \label{eq:V114-BRST-Dirichlet-variation}
\end{equation}
Lemma~\ref{lem:V104-free-normal-ghost} permits arbitrary \(\rho=X\).
Preservation of the row therefore gives \(P_D\nu_X=0\) for every \(X\).
Orthogonality of \eqref{eq:V114-DF-splitting} then puts \(\nu_X\) in \(F\).
For nonzero \(X\), this vector is nonzero.
\end{proof}

\begin{corollary}[Pure gauge Dirichlet is incompatible with the unchanged
ghost domain]
\label{cor:V114-pure-Dirichlet-fails-BRST}
The pure Dirichlet choice \(D=Q\), \(F=\{0\}\) has a unit stable matrix but is
not preserved by BRST under the ghost condition
\eqref{eq:V114-Dirichlet-ghost}.
\end{corollary}

\begin{proof}
For \(D=Q\), one has \(P_D\nu_X=\nu_X\ne0\), contradicting the necessary
condition \eqref{eq:V114-normal-in-F}.
\end{proof}

\subsection{The combined local no-go theorem}
\label{subsec:V114-combined-no-go}

\begin{theorem}[No regular mixed local relation passes all three gates]
\label{thm:V114-mixed-no-go}
Consider a finite regular boundary relation of the form
\(\mathcal L_{D,R}\) for the frozen Yang--Mills wave block, with:
\begin{enumerate}[label=\textup{(\roman*)},leftmargin=3.3em]
\item \(R\) generated by a real local first-order quadratic tangential action;
\item the unchanged Dirichlet ghost trace
      \eqref{eq:V114-Dirichlet-ghost};
\item BRST preservation of the gauge Dirichlet rows.
\end{enumerate}
Then the relation is action-Lagrangian only when \(R=R^\dagger\), and on that
branch
\begin{equation}
 \boxed{
 \inf_{\overline{\mathbb S}_+}
 s_{\min}\mathcal B_{D,R}=0.}
 \label{eq:V114-three-gate-failure}
\end{equation}
Thus no member of this class is simultaneously action-Lagrangian,
BRST-preserved, and uniformly complementing.
\end{theorem}

\begin{proof}
Action incidence is Proposition~\ref{prop:V114-Lagrangian}.  BRST preservation
and Theorem~\ref{thm:V114-normal-free} force \(F\ne\{0\}\).  Lemma
\ref{lem:V114-mixed-floor}, using V113 on that free block, proves
\eqref{eq:V114-three-gate-failure}.
\end{proof}

\begin{corollary}[The normal direction carries the obstruction in every
adjoint component]
\label{cor:V114-adjoint-embedding}
Tensoring with the twelve-dimensional Standard Model adjoint fibre produces
at least twelve free normal gauge directions.  A stable near-kernel supported
in one such direction and its finite action closure embeds into the full
gauge/ghost principal direct sum.
\end{corollary}

\begin{proof}
Theorem~\ref{thm:V114-normal-free} applies independently to each adjoint
amplitude \(X\).  The free block contains their span, and direct-sum embedding
preserves the V113 singular sequence.
\end{proof}

\begin{firewall}[Scope of the regular mixed-relation theorem]
\label{fire:V114-scope}
Theorem~\ref{thm:V114-mixed-no-go} keeps the Dirichlet ghost trace and assumes
a regular orthogonal Dirichlet--conormal split with identical wave normal
dispersion on its free finite sector.  It does not cover a boundary
Stueckelberg field whose BRST variation absorbs \(D_nc\), a changed
Sommerfeld or nonlocal ghost graph, a singular frequency-dependent
Lagrangian relation, or a field with different principal normal order.
Those changes alter the hypotheses and must print their complete Green and
BRST rows.
\end{firewall}

\subsection{V114 ledger and mandatory V115 acquisition}
\label{subsec:V114-ledger}

\paragraph{New exact conclusions.}
\begin{enumerate}[label=\textup{(\arabic*)},leftmargin=2.8em]
\item Proposition~\ref{prop:V114-Lagrangian} proves the exact self-adjointness
      criterion for a regular mixed Dirichlet--conormal relation.
\item Equation \eqref{eq:V114-block-stable-matrix} and Lemma
      \ref{lem:V114-mixed-floor} isolate its Dirichlet unit block and free
      conormal block.
\item Theorem~\ref{thm:V114-normal-free} proves from the arbitrary normal
      ghost jet that every normal gauge direction lies in the free sector.
\item Corollary~\ref{cor:V114-pure-Dirichlet-fails-BRST} shows why the only
      empty-free-sector choice fails BRST.
\item Theorem~\ref{thm:V114-mixed-no-go} combines these facts with V113 and
      closes the regular finite local mixed-relation class under the unchanged
      ghost domain.
\end{enumerate}

\paragraph{Next forced step.}
V115 is a five-version checkpoint.  It must audit the newest shared
Yang--Mills and Navier--Stokes notes, then change one hypothesis of
Theorem~\ref{thm:V114-mixed-no-go}.  The least duplicative route is the exact
nonlocal Calder\'on/Dirichlet-to-Neumann graph with a matched ghost projector.
The note must derive its Green adjoint, BRST incidence, stable projection,
stress variation, and suitability for the V102 heat-kernel domain.  If a
local boundary Stueckelberg alternative is chosen instead, its new ghost and
multiplier fields must be explicit.

\begin{assessment}[V114 continuum status]
\label{ass:V114-continuum-status}
The full Standard--Model--Einstein continuum remains unproved.  V114 closes
the regular finite mixed Dirichlet--conormal action class while the relative
Dirichlet ghost domain is retained.  A changed ghost/BRST multiplet or
nonlocal gauge realization, the fermion branch selection, common quantum
coefficients and form bounds, curved radiative equivariance, nonlinear
existence, and the cofinal interacting limit remain open.
\end{assessment}

\section*{V114 exact checks and append-only record}
\addcontentsline{toc}{section}{V114 exact checks and append-only record}

The decisive mixed-relation checksum is
\begin{equation}
 \boxed{
 P_Da=0\text{ BRST-preserved}
 \Longrightarrow P_D(n^\flat X)=0
 \Longrightarrow n^\flat X\in F\ne\{0\}
 \Longrightarrow
 \inf s_{\min}\mathcal B_{D,R}=0.}
 \label{eq:V114-final-check}
\end{equation}
The append operation preserves the complete V113 prefix byte for byte before
its former live terminator.  The assembled V114 file is checked for one live
final terminator, balanced environments and braces, duplicate labels, newly
introduced unresolved references, display math delimiters, and control
characters before the exact V113 predecessor is removed.

% =============================================================================
% END APPENDED VERSION V114 MATERIAL
% =============================================================================

\clearpage
% =============================================================================
% BEGIN APPENDED VERSION V115 MATERIAL
% =============================================================================
\section{V115: companion audit, correction of the componentwise BRST claim,
and the Euclidean Dirichlet-to-Neumann realization}
\label{sec:V115-DtN}

V114 ended the regular local mixed-relation branch with the relative
Dirichlet ghost domain.  Before changing to a nonlocal graph, one incidence
error in V105--V113 must be corrected.  A boundary identity \(Bc=0\) does not
in general imply \(D_n(Bc)=0\).  Thus the earlier componentwise
Sommerfeld/Robin BRST arguments were valid for tangential gauge components but
did not justify preservation of the normal component.  Their action,
stable-matrix, determinant, and no-floor theorems remain valid as algebraic
boundary results; only the stated BRST closure of those componentwise graphs
is withdrawn.

The exact Dirichlet-to-Neumann operator supplies the missing identity on a
Riemannian product collar.  Its boundary row factorizes the ghost Euler
operator, so the ghost equation and boundary row together force the normal
derivative of the boundary defect to vanish.  The same nonlocal row is
action-Lagrangian, has a uniform parameter-stable floor, and defines a closed
nonnegative quadratic-form realization.  This is the first tested
gauge/ghost branch to pass those three frozen/product-collar gates
simultaneously.  Curved nonproduct incidence, harmonic zero modes, nonlocal
stress renormalization, and the full quantum coefficient table remain open.

\subsection{Mandatory V115 companion audit}
\label{subsec:V115-audit}

The read-only companion snapshots are
\begin{center}
\begin{tabular}{|p{0.48\textwidth}|r|p{0.31\textwidth}|}
\hline
Companion file & Bytes & SHA--256\\
\hline
\texttt{Renewal\_NS\_Stability\_Research\_Notes\_V31.tex}
& \(887537\) &
\texttt{F1121B62238AD5491BB7CF03635EA9934942EDF573ED8FAAA9EE79E1D38A6696}\\
\hline
\texttt{Renewal\_Yang\_Mills\_Mass\_Gap\_Research\_Notes\_V104.tex}
& \(2477653\) &
\texttt{B30D5EAB41246F519AF4112DA1D440941E36BE80BDE00041BD8EF1ECD8FEE281}\\
\hline
\end{tabular}
\end{center}

The Navier--Stokes note remains byte-identical to the V110 snapshot and
retains its balanced-helicity nullspace warning.  Yang--Mills has advanced
from V102 to V104.  V103 replaces a full pair average by the exact top
fractional Ky--Fan mean and identifies the
\(s_\alpha N_{\rm eff}/N^2\) threshold.  V104 retains two forced pair banks in
one product operator; the cold--cold exceptional rank is a product, and its
conditional symmetric threshold improves to order \(N\) only when the
physical tensor separation and row orientation are actually present.

\begin{assessment}[Transfer used in V115]
\label{ass:V115-audit-transfer}
The applicable lesson is to preserve the complete correlated operator before
forming a scalar majorant.  V115 therefore uses
\(\Lambda_\partial=\sqrt{\Delta_\partial}\) itself in the boundary equation,
factorization, form, and metric variation.  Replacing it by the scalar
inequality \(\|\Lambda_\partial q\|\le C\|q\|_{H^1}\) would erase the
factorization which repairs the normal BRST incidence.  No Yang--Mills
collision estimate or Navier--Stokes bound is imported.
\end{assessment}

\subsection{Correction of V105--V113 BRST incidence}
\label{subsec:V115-correction}

Let \(B=D_n+T_\parallel\) be any componentwise boundary operator whose
tangential part commutes with frozen derivatives.  For a tangential gauge
component \(a_j\),
\begin{equation}
 \mathsf s(Ba_j)=B(D_jc)=D_j(Bc),
 \label{eq:V115-tangential-BRST}
\end{equation}
and \(Bc=0\) does imply that the tangential derivative on the right vanishes.
For the normal component,
\begin{equation}
 \mathsf s(Ba_n)=B(D_nc)=D_n(Bc),
 \label{eq:V115-normal-BRST-defect}
\end{equation}
but the normal derivative of a vanishing boundary trace is unrestricted
without an additional equation or prolonged boundary row.

\begin{proposition}[Precise correction]
\label{prop:V115-BRST-correction}
The componentwise BRST-closure assertions in
Propositions~\ref{prop:V105-BRST-incidence},
\ref{prop:V106-BRST-incidence},
\ref{prop:V107-BRST-incidence},
\ref{prop:V108-BRST},
\ref{prop:V109-BRST},
\ref{prop:V110-BRST},
\ref{prop:V111-BRST},
\ref{prop:V112-BRST}, and
\ref{prop:V113-BRST}
require the additional normal condition
\begin{equation}
 D_n(Bc)|_{\mathcal T}=0
 \label{eq:V115-missing-normal-condition}
\end{equation}
or an equivalent ghost-Euler factorization.  Their proofs establish
tangential incidence but do not establish
\eqref{eq:V115-missing-normal-condition}.  All action and stable-symbol
statements in V105--V113 are independent of that missing implication and
remain unchanged.
\end{proposition}

\begin{proof}
Equations \eqref{eq:V115-tangential-BRST} and
\eqref{eq:V115-normal-BRST-defect} give the distinction.  A function may
vanish on a hypersurface while its normal derivative is arbitrary, as already
used in Lemma~\ref{lem:V104-free-normal-ghost}.  Therefore \(Bc=0\) alone
does not prove the normal row.  The action variations and stable determinants
in the cited versions did not use this implication.
\end{proof}

\begin{assessment}[Effect of the correction]
\label{ass:V115-correction-effect}
No previously claimed full continuum result depends on the erroneous step:
each affected local branch already failed either its action gate or its
uniform stable floor.  V114's main conclusion also remains valid.  If its free
conormal row is not BRST-preserved, the branch fails BRST immediately; if it
is preserved by extra data, its V113 stable floor is still zero.  The
correction changes the acquisition requirement for a successful branch: it
must print the normal ghost-Euler factorization explicitly.
\end{assessment}

\subsection{Riemannian product collar and exact square root}
\label{subsec:V115-product-collar}

Let
\begin{equation}
 M=[0,\infty)_r\times\Sigma
 \label{eq:V115-product-collar}
\end{equation}
near the boundary, with product Riemannian metric, and let
\(\Delta_{\Sigma,p}\ge0\) be the Hodge Laplacian on boundary \(p\)-forms.
Define
\begin{equation}
 \Lambda_p:=\Delta_{\Sigma,p}^{1/2}
 \label{eq:V115-Lambda-p}
\end{equation}
by spectral functional calculus.  The scalar ghost operator and its exact
factorization are
\begin{align}
 P_0&:=-\partial_r^2+\Delta_{\Sigma,0},
 \label{eq:V115-ghost-operator}\\
 B_0&:=-\partial_r+\Lambda_0,
 \label{eq:V115-ghost-boundary}\\
 P_0&=(\partial_r+\Lambda_0)B_0.
 \label{eq:V115-factorization}
\end{align}
For tangential one-forms set
\begin{equation}
 B_1:=-\partial_r+\Lambda_1.
 \label{eq:V115-oneform-boundary}
\end{equation}
Hodge functional calculus gives the exact intertwining identity
\begin{equation}
 d_\Sigma\Lambda_0=\Lambda_1d_\Sigma.
 \label{eq:V115-Hodge-intertwining}
\end{equation}

\begin{lemma}[The ghost equation supplies the missing normal row]
\label{lem:V115-normal-factorization}
If
\begin{equation}
 P_0c=0,
 \qquad
 B_0c|_{\partial M}=0,
 \label{eq:V115-ghost-domain}
\end{equation}
then
\begin{equation}
 \partial_r(B_0c)|_{\partial M}=0.
 \label{eq:V115-normal-defect-zero}
\end{equation}
\end{lemma}

\begin{proof}
Evaluate \eqref{eq:V115-factorization} on \(c\) at the boundary:
\[
 0=P_0c=(\partial_r+\Lambda_0)B_0c.
\]
The boundary trace \(B_0c\) is zero, so its tangential image
\(\Lambda_0B_0c\) is zero.  The remaining term is
\(\partial_r(B_0c)\), proving
\eqref{eq:V115-normal-defect-zero}.
\end{proof}

\subsection{Exact product-collar BRST subcomplex}
\label{subsec:V115-BRST}

Decompose the gauge potential trace into its tangential one-form
\(a_\parallel\) and normal scalar \(a_n\).  Impose
\begin{equation}
 B_1a_\parallel=0,
 \qquad
 B_0a_n=0,
 \qquad
 B_0c=0,
 \qquad
 B_0\bar c=0,
 \qquad
 B_0b=0,
 \label{eq:V115-full-boundary-domain}
\end{equation}
where the multiplier row is a compatibility row if \(b\) is eliminated
algebraically.

\begin{theorem}[On-shell BRST incidence of the Dirichlet-to-Neumann graph]
\label{thm:V115-BRST-incidence}
On the product collar, the domain
\eqref{eq:V115-full-boundary-domain} is preserved by the linearized BRST
differential on the ghost Euler equation:
\begin{equation}
 \mathsf s a_\parallel=d_\Sigma c,
 \qquad
 \mathsf s a_n=\partial_rc,
 \qquad
 \mathsf sc=0,
 \qquad
 \mathsf s\bar c=b,
 \qquad
 \mathsf sb=0.
 \label{eq:V115-product-BRST}
\end{equation}
\end{theorem}

\begin{proof}
For the tangential potential, product structure and
\eqref{eq:V115-Hodge-intertwining} give
\[
 \mathsf s(B_1a_\parallel)
 =B_1d_\Sigma c=d_\Sigma B_0c=0.
\]
For the normal potential, \(\partial_r\) commutes with
\(\Lambda_0\), so
\[
 \mathsf s(B_0a_n)
 =B_0\partial_rc=\partial_r(B_0c)=0
\]
by Lemma~\ref{lem:V115-normal-factorization}.  The ghost row maps to zero,
the antighost row maps to the displayed multiplier row, and the multiplier
row maps to zero.
\end{proof}

\subsection{Action, form domain, and the stable floor}
\label{subsec:V115-action-floor}

Use the outward conormal \(p=-\partial_rq\).  The boundary action on a
\(p\)-form trace is
\begin{equation}
 S_{\partial,p}[q]
 :=\frac12\int_\Sigma
 \langle q,\Lambda_pq\rangle.
 \label{eq:V115-boundary-action}
\end{equation}
Its Euler row is \(p+\Lambda_pq=0\), which is \(B_pq=0\).

\begin{proposition}[Lagrangian and closed-form realization]
\label{prop:V115-form-realization}
On a compact Riemannian manifold with product collar, the quadratic form
\begin{equation}
 \mathfrak q_p[u]
 :=
 \int_M|\nabla u|^2
 +\|\Lambda_p^{1/2}\gamma u\|_{L^2(\Sigma)}^2,
 \qquad
 \operatorname{Dom}\mathfrak q_p=H^1(M;\Lambda^p),
 \label{eq:V115-closed-form}
\end{equation}
is densely defined, nonnegative, and closed after adjoining the \(L^2\)
norm.  Its associated self-adjoint realization has boundary row
\begin{equation}
 -\partial_ru+\Lambda_pu=0.
 \label{eq:V115-associated-boundary-row}
\end{equation}
The corresponding boundary phase graph is Lagrangian.
\end{proposition}

\begin{proof}
The trace map sends \(H^1(M)\) continuously to
\(H^{1/2}(\Sigma)\), which is the form domain of the order-one positive
operator \(\Lambda_p\).  Hence the boundary term is finite and continuous in
the \(H^1\) norm.  Nonnegativity and completeness of \(H^1\) show that
\(\mathfrak q_p+\|\cdot\|_{L^2}^2\) is closed.  The representation theorem for
closed forms gives a self-adjoint nonnegative operator.  Green's identity and
variation of the boundary term give
\eqref{eq:V115-associated-boundary-row}.  Since \(\Lambda_p\) is
self-adjoint, the graph \(p=-\Lambda_pq\) is Lagrangian for the Green form.
\end{proof}

For the parameter problem \(P_p+z\), let
\(\mu=|\xi|\) be the boundary frequency and choose
\begin{equation}
 \lambda(\xi,z):=\sqrt{\mu^2+z},
 \qquad
 \operatorname{Re}\lambda\ge0,
 \qquad
 \operatorname{Re}z\ge0.
 \label{eq:V115-parameter-root}
\end{equation}
The stable boundary symbol is
\begin{equation}
 b_{\rm DtN}(\xi,z)=\lambda(\xi,z)+\mu.
 \label{eq:V115-stable-symbol}
\end{equation}

\begin{theorem}[Uniform parameter-stable floor]
\label{thm:V115-parameter-floor}
On the normalized parameter hemisphere
\begin{equation}
 \mu^2+|z|=1,
 \qquad
 \mu\ge0,
 \qquad
 \operatorname{Re}z\ge0,
 \label{eq:V115-parameter-normalization}
\end{equation}
one has
\begin{equation}
 \boxed{
 |b_{\rm DtN}(\xi,z)|
 \ge|\lambda(\xi,z)|
 \ge\max\{\mu,\sqrt{|z|}\}
 \ge2^{-1/2}.}
 \label{eq:V115-parameter-floor}
\end{equation}
At \(z=0\) and \(\mu=1\), the exact boundary value is \(2\).
Tensoring with gauge, ghost, form, and Standard Model adjoint identities does
not change the floor.
\end{theorem}

\begin{proof}
Since \(\operatorname{Re}\lambda\ge0\) and \(\mu\ge0\),
\[
 |\lambda+\mu|^2
 =|\lambda|^2+\mu^2+2\mu\operatorname{Re}\lambda
 \ge|\lambda|^2.
\]
Moreover,
\[
 |\lambda|^4
 =|\mu^2+z|^2
 =\mu^4+|z|^2+2\mu^2\operatorname{Re}z
 \ge\max\{\mu^4,|z|^2\}.
\]
Taking roots gives the middle inequality.  Under
\eqref{eq:V115-parameter-normalization}, at least one of
\(\mu,\sqrt{|z|}\) is at least \(2^{-1/2}\).  The \(z=0\) value is direct.
\end{proof}

\subsection{Metric variation and the quantum-domain boundary}
\label{subsec:V115-stress-quantum}

Let \(A_\gamma=\Delta_{\Sigma,p}(\gamma)\) on the orthogonal complement of
its harmonic kernel.  For a metric variation \(\dot\gamma\), write
\(\dot A=D_\gamma A[\dot\gamma]\).

\begin{proposition}[Exact Fr\'echet derivative of the boundary square root]
\label{prop:V115-square-root-variation}
On a fixed positive spectral subspace,
\begin{equation}
 D_\gamma\Lambda_p[\dot\gamma]
 =
 \frac1\pi\int_0^\infty
 t^{1/2}(A_\gamma+t)^{-1}
 \dot A
 (A_\gamma+t)^{-1}\,\dd t.
 \label{eq:V115-square-root-variation}
\end{equation}
Consequently the boundary metric variation contains the exact nonlocal term
\begin{equation}
 \frac12\left\langle
 q,D_\gamma\Lambda_p[\dot\gamma]q
 \right\rangle
 \label{eq:V115-nonlocal-stress-term}
\end{equation}
in addition to the boundary measure and bundle-metric variations.
\end{proposition}

\begin{proof}
Use the functional-calculus representation
\[
 A^{1/2}
 =\frac1\pi\int_0^\infty
 t^{-1/2}A(A+t)^{-1}\,\dd t.
\]
Since \(A(A+t)^{-1}=I-t(A+t)^{-1}\), its derivative in direction
\(\dot A\) is
\[
 t(A+t)^{-1}\dot A(A+t)^{-1}.
\]
Substitution gives \eqref{eq:V115-square-root-variation}.  Differentiating
\eqref{eq:V115-boundary-action} gives
\eqref{eq:V115-nonlocal-stress-term} and the stated elementary metric
variations.
\end{proof}

\begin{firewall}[This changes the quantum boundary calculus]
\label{fire:V115-nonlocal-calculus}
The realization in Proposition~\ref{prop:V115-form-realization} is a genuine
closed self-adjoint operator/domain pair, and
Theorem~\ref{thm:V115-parameter-floor} proves its frozen parameter
complementing estimate.  Its boundary operator is pseudodifferential and its
stress variation is nonlocal.  Therefore the local mixed-boundary heat
coefficients contemplated in V102 cannot be copied onto this branch.  The
appropriate pseudodifferential boundary heat calculus, harmonic-mode
projection, ghost grading, and curvature corrections must be acquired before
any Standard Model counterterm is claimed.
\end{firewall}

\begin{firewall}[Product collar is an actual hypothesis]
\label{fire:V115-product-scope}
The exact factorization \eqref{eq:V115-factorization} and commutation used in
Theorem~\ref{thm:V115-BRST-incidence} hold on the product principal model.
Extrinsic curvature, a varying normal connection, and nonproduct gauge
coefficients add commutators.  They are lower order for principal
complementing but matter for exact curved BRST incidence and the quantum
stress.  V115 does not declare them zero.
\end{firewall}

\subsection{V115 ledger and next acquisition}
\label{subsec:V115-ledger}

\paragraph{New exact conclusions.}
\begin{enumerate}[label=\textup{(\arabic*)},leftmargin=2.8em]
\item The mandatory audit records exact NS V31 and YM V104 snapshots.
\item Proposition~\ref{prop:V115-BRST-correction} corrects the unsupported
      normal-component implication in V105--V113 while preserving their
      independent action and stable-symbol results.
\item Lemma~\ref{lem:V115-normal-factorization} supplies the missing normal
      ghost row from the exact Dirichlet-to-Neumann factorization.
\item Theorem~\ref{thm:V115-BRST-incidence} proves the product-collar
      gauge/ghost BRST subcomplex.
\item Proposition~\ref{prop:V115-form-realization} constructs a closed
      nonnegative self-adjoint action realization.
\item Theorem~\ref{thm:V115-parameter-floor} proves a uniform nonlocal
      parameter-stable floor.
\item Proposition~\ref{prop:V115-square-root-variation} prints the exact
      nonlocal metric derivative required for the boundary stress.
\end{enumerate}

\paragraph{Next forced step.}
V116 should leave the flat/product principal level and compute the first
nonproduct commutator in
\[
 P_0-(\partial_r+\Lambda_0)(-\partial_r+\Lambda_0).
\]
It must identify the extrinsic-curvature and normal-connection terms, then
decide whether a lower-order correction to \(B_0\) restores the ghost
factorization and Hodge intertwining.  In parallel it should isolate harmonic
boundary modes as a finite-dimensional BRST complex.  Only after those two
steps should the pseudodifferential heat coefficient and nonlocal stress card
be populated.

\begin{assessment}[V115 continuum status]
\label{ass:V115-continuum-status}
The full Standard--Model--Einstein continuum remains unproved.  V115 corrects
the earlier componentwise BRST overclaim and constructs a product-collar
nonlocal gauge/ghost realization which passes action, on-shell BRST, closed
form, and uniform parameter-stable gates.  Curved factorization, harmonic
modes, the nonlocal heat/stress calculus, the fermion branch selection,
common quantum coefficients, nonlinear existence, and the cofinal
interacting limit remain open.
\end{assessment}

\section*{V115 exact checks and append-only record}
\addcontentsline{toc}{section}{V115 exact checks and append-only record}

The decisive factorization and floor checks are
\begin{equation}
 \boxed{
 P_0=(\partial_r+\Lambda_0)B_0,
 \qquad
 P_0c=B_0c=0
 \Longrightarrow\partial_r(B_0c)=0,
 \qquad
 |\lambda+\mu|\ge2^{-1/2}.}
 \label{eq:V115-final-check}
\end{equation}
The append operation preserves the complete V114 prefix byte for byte before
its former live terminator.  The assembled V115 file is checked for one live
final terminator, balanced environments and braces, duplicate labels, newly
introduced unresolved references, display math delimiters, and control
characters before the exact V114 predecessor is removed.  The audited
companion files were read only.

% =============================================================================
% END APPENDED VERSION V115 MATERIAL
% =============================================================================

\clearpage
% =============================================================================
% BEGIN APPENDED VERSION V116 MATERIAL
% =============================================================================
\section{V116: nonproduct collar Riccati defect, the first self-adjoint
correction, and harmonic-mode separation}
\label{sec:V116-Riccati-harmonic}

V115 constructed the exact product-collar boundary square root and identified
nonproduct geometry as the next obstruction.  V116 computes that obstruction
without replacing the nonlocal operator by a norm bound.  A varying boundary
measure is first removed by a unitary conjugation.  The exact
Dirichlet-to-Neumann factor then obeys an operator Riccati equation.  The
instantaneous square root fails by its normal derivative, and a unique
self-adjoint order-zero Sylvester correction cancels that first defect.

The same calculation explains why harmonic boundary modes must be separated.
The Sylvester inverse is elliptic on the positive spectral complement.  On
the harmonic bundle its denominator vanishes, but the harmonic--harmonic
numerator vanishes identically; the remaining finite-dimensional block is a
normal transport problem rather than a high-frequency square-root problem.

\subsection{Scalar Laplacian in boundary normal coordinates}
\label{subsec:V116-normal-coordinates}

Let
\begin{equation}
 g=dr^2+\gamma_r
 \label{eq:V116-collar-metric}
\end{equation}
in a collar of a compact boundary \(\Sigma=\{r=0\}\), with \(r\) increasing
inward.  Relative to the nonnegative tangential Laplacian
\(\Delta_{\gamma_r}\), the nonnegative scalar Laplacian has the normal form
\begin{equation}
 P_r
 =
 -\partial_r^2-h_r\partial_r+\Delta_{\gamma_r},
 \qquad
 h_r:=\partial_r\log J_r,
 \qquad
 d{\rm vol}_{\gamma_r}=J_r\,d{\rm vol}_{\gamma_0}.
 \label{eq:V116-raw-normal-form}
\end{equation}
The sign of \(h_r\) relative to the named mean curvature depends on the
chosen inward normal; Equation \eqref{eq:V116-raw-normal-form} fixes the
convention used below.

Define the unitary identification
\begin{equation}
 U_r:L^2(\Sigma,d{\rm vol}_{\gamma_r})
 \longrightarrow L^2(\Sigma,d{\rm vol}_{\gamma_0}),
 \qquad
 U_rf:=J_r^{1/2}f.
 \label{eq:V116-unitary-density-map}
\end{equation}

\begin{proposition}[Exact fixed-measure normal form]
\label{prop:V116-fixed-measure}
Conjugation by \(U_r\) gives
\begin{equation}
 \widetilde P
 :=U_rP_rU_r^{-1}
 =
 -\partial_r^2+\mathcal A_r,
 \label{eq:V116-fixed-normal-form}
\end{equation}
where
\begin{equation}
 \mathcal A_r
 :=
 U_r\Delta_{\gamma_r}U_r^{-1}
 +\frac12\partial_r h_r+\frac14h_r^2.
 \label{eq:V116-Ar}
\end{equation}
Every \(\mathcal A_r\) is formally self-adjoint on the fixed boundary Hilbert
space.  Its principal symbol is
\begin{equation}
 \sigma_2(\mathcal A_r)(y,\xi)=|\xi|_{\gamma_r}^2.
 \label{eq:V116-A-principal-symbol}
\end{equation}
\end{proposition}

\begin{proof}
Write \(u=U_r^{-1}v=J_r^{-1/2}v\).  Direct differentiation gives
\[
 U_r(-\partial_r^2-h_r\partial_r)U_r^{-1}v
 =
 -\partial_r^2v+
 \left(\frac12\partial_rh_r+\frac14h_r^2\right)v.
\]
The tangential term becomes
\(U_r\Delta_{\gamma_r}U_r^{-1}\).  Unitarity gives formal
self-adjointness, and conjugation by a scalar multiplication does not change
the order-two principal symbol.
\end{proof}

\subsection{The exact Riccati equation}
\label{subsec:V116-Riccati}

Let \(L_r\) be a tangential operator of order one and set
\begin{equation}
 B_L:=-\partial_r+L_r.
 \label{eq:V116-BL}
\end{equation}

\begin{lemma}[Factorization defect]
\label{lem:V116-factorization-defect}
One has the exact operator identity
\begin{equation}
 (\partial_r+L_r)B_L-\widetilde P
 =
 \partial_rL_r+L_r^2-\mathcal A_r.
 \label{eq:V116-factorization-defect}
\end{equation}
Thus exact factorization is equivalent to the operator Riccati equation
\begin{equation}
 \partial_rL_r+L_r^2=\mathcal A_r.
 \label{eq:V116-Riccati}
\end{equation}
\end{lemma}

\begin{proof}
Apply both factors to a test function.  The terms
\(L_r\partial_r\) cancel, while differentiating the varying operator in
\(\partial_r(L_ru)\) leaves \((\partial_rL_r)u\).  Comparison with
\eqref{eq:V116-fixed-normal-form} proves the identity.
\end{proof}

Assume temporarily that \(\mathcal A_r\) is positive on the spectral sector
under consideration and define the instantaneous square root
\begin{equation}
 \Lambda_r:=\mathcal A_r^{1/2}.
 \label{eq:V116-instantaneous-root}
\end{equation}

\begin{corollary}[The first nonproduct defect]
\label{cor:V116-first-defect}
The naive replacement \(L_r=\Lambda_r\) has exact factorization defect
\begin{equation}
 (\partial_r+\Lambda_r)(-\partial_r+\Lambda_r)-\widetilde P
 =\partial_r\Lambda_r.
 \label{eq:V116-first-defect}
\end{equation}
It is an order-one operator.  It vanishes on the V115 product collar and is
generically nonzero when \(\gamma_r\) varies.
\end{corollary}

\begin{proof}
Insert \(L_r=\Lambda_r\) into
\eqref{eq:V116-factorization-defect} and use
\(\Lambda_r^2=\mathcal A_r\).  Normal differentiation preserves
pseudodifferential order.
\end{proof}

\subsection{The order-zero Sylvester correction}
\label{subsec:V116-Sylvester}

Seek
\begin{equation}
 L_r^{(0)}:=\Lambda_r+C_r,
 \label{eq:V116-corrected-root}
\end{equation}
with \(C_r\) of order zero.  Cancellation of the order-one Riccati defect
requires
\begin{equation}
 \Lambda_rC_r+C_r\Lambda_r=-\partial_r\Lambda_r.
 \label{eq:V116-Sylvester-equation}
\end{equation}

\begin{theorem}[Exact self-adjoint first correction]
\label{thm:V116-first-correction}
On a positive spectral subspace of \(\Lambda_r\), Equation
\eqref{eq:V116-Sylvester-equation} has the unique solution
\begin{equation}
 \boxed{
 C_r
 =
 -\int_0^\infty
 e^{-t\Lambda_r}
 (\partial_r\Lambda_r)
 e^{-t\Lambda_r}\,\dd t.}
 \label{eq:V116-Sylvester-solution}
\end{equation}
It is formally self-adjoint and has pseudodifferential order zero.  The
remaining exact Riccati defect is
\begin{equation}
 (\partial_r+L_r^{(0)})(-\partial_r+L_r^{(0)})-\widetilde P
 =
 \partial_rC_r+C_r^2,
 \label{eq:V116-second-defect}
\end{equation}
which has order zero.
\end{theorem}

\begin{proof}
Let \(X=\partial_r\Lambda_r\).  Positivity makes the heat integral converge
on the declared spectral subspace.  Since \(\Lambda_r\) commutes with its
semigroup,
\begin{align}
 \Lambda_rC_r+C_r\Lambda_r
 &=
 -\int_0^\infty
 \left(
 \Lambda_re^{-t\Lambda_r}Xe^{-t\Lambda_r}
 +e^{-t\Lambda_r}Xe^{-t\Lambda_r}\Lambda_r
 \right)\dd t\notag\\
 &=
 \int_0^\infty
 \partial_t\left(
 e^{-t\Lambda_r}Xe^{-t\Lambda_r}
 \right)\dd t
 =-X.
 \label{eq:V116-Sylvester-check}
\end{align}
The homogeneous Sylvester equation has only zero solution because the sums of
positive spectral values are nonzero.  Since \(X=X^\dagger\), the integral is
self-adjoint.  In the symbolic calculus, the inverse Sylvester map divides an
order-one symbol by the order-one positive sum
\(|\xi|+|\xi|\), so \(C_r\) has order zero.

Finally insert \(L_r^{(0)}=\Lambda_r+C_r\) into
\eqref{eq:V116-factorization-defect}.  The terms
\(\partial_r\Lambda_r+\Lambda_rC_r+C_r\Lambda_r\) cancel by
\eqref{eq:V116-Sylvester-equation}; the remainder is exactly
\(\partial_rC_r+C_r^2\).
\end{proof}

\begin{corollary}[Formal recursive construction]
\label{cor:V116-recursion}
If \(D_r\) of order \(-1\) solves
\begin{equation}
 \Lambda_rD_r+D_r\Lambda_r
 =-(\partial_rC_r+C_r^2),
 \label{eq:V116-second-Sylvester}
\end{equation}
then \(L_r=\Lambda_r+C_r+D_r\) cancels the Riccati defect through order zero.
Iterating the same Sylvester inversion constructs the full classical
Dirichlet-to-Neumann symbol modulo a smoothing remainder.
\end{corollary}

\begin{proof}
The cross terms with \(\Lambda_r\) are the only order-zero contribution of
\(D_r\).  Equation \eqref{eq:V116-second-Sylvester} cancels
\eqref{eq:V116-second-defect}; all remaining terms have order at most
\(-1\).  Induction lowers the order by one at each step.
\end{proof}

\begin{proposition}[The corrected boundary graph remains variational]
\label{prop:V116-corrected-action}
At every finite symbolic order, choose the Sylvester correction
self-adjoint.  Then the boundary graph
\begin{equation}
 -\partial_rv+L_rv=0
 \label{eq:V116-corrected-boundary}
\end{equation}
is Lagrangian for the fixed-measure Green form.  In the original variable
\(u=U_r^{-1}v\), its tangential conormal operator at the boundary is
\begin{equation}
 L_r^{\rm raw}
 =
 U_r^{-1}L_rU_r-\frac12h_r.
 \label{eq:V116-raw-corrected-root}
\end{equation}
\end{proposition}

\begin{proof}
The graph of a self-adjoint conormal operator is Lagrangian.  For the second
identity, differentiate \(v=U_ru\):
\[
 -\partial_rv+L_rv
 =
 U_r\left[
 -\partial_ru+
 \left(U_r^{-1}L_rU_r-\frac12h_r\right)u
 \right].
\]
\end{proof}

\subsection{Harmonic modes and the singular Sylvester denominator}
\label{subsec:V116-harmonic}

Return to a smooth family of boundary Hodge Laplacians
\(\Delta_{\Sigma,p}(r)\) on a compact \(\Sigma\), and let
\begin{equation}
 \Pi_r:=\mathbf1_{\{0\}}(\Delta_{\Sigma,p}(r)),
 \qquad
 \mathcal H_r^p:=\operatorname{Ran}\Pi_r.
 \label{eq:V116-harmonic-projection}
\end{equation}
The dimension of \(\mathcal H_r^p\) is the Betti number \(b_p(\Sigma)\), hence
is constant.  Put \(\Lambda_r=\Delta_{\Sigma,p}(r)^{1/2}\) for this harmonic
analysis.

\begin{lemma}[The harmonic--harmonic numerator vanishes]
\label{lem:V116-harmonic-numerator}
For a differentiable family,
\begin{equation}
 \Pi_r(\partial_r\Lambda_r)\Pi_r=0.
 \label{eq:V116-harmonic-numerator}
\end{equation}
\end{lemma}

\begin{proof}
The spectral identity \(\Lambda_r\Pi_r=0\) gives, after differentiation,
\[
 (\partial_r\Lambda_r)\Pi_r+\Lambda_r(\partial_r\Pi_r)=0.
\]
Multiply on the left by \(\Pi_r\) and use
\(\Pi_r\Lambda_r=0\).  This yields
\eqref{eq:V116-harmonic-numerator}.
\end{proof}

\begin{theorem}[Spectral Sylvester formula with harmonic separation]
\label{thm:V116-harmonic-Sylvester}
Let \(\{\phi_j\}\) be an orthonormal eigenbasis of \(\Lambda_r\), with
\(\Lambda_r\phi_j=\mu_j\phi_j\), \(\mu_j\ge0\), and set
\(X=\partial_r\Lambda_r\).  The self-adjoint gauge choice
\(\Pi_rC_r\Pi_r=0\) gives
\begin{equation}
 \boxed{
 \langle\phi_i,C_r\phi_j\rangle
 =
 \begin{cases}
 -\dfrac{\langle\phi_i,X\phi_j\rangle}{\mu_i+\mu_j},
 &\mu_i+\mu_j>0,\\[8pt]
 0,&\mu_i=\mu_j=0.
 \end{cases}}
 \label{eq:V116-spectral-C}
\end{equation}
This solves \eqref{eq:V116-Sylvester-equation}.  The zero denominator creates
no compatibility obstruction because of
\eqref{eq:V116-harmonic-numerator}, but it leaves a finite-dimensional
harmonic transport freedom not determined by the elliptic square root.
\end{theorem}

\begin{proof}
Taking matrix elements of the Sylvester equation gives
\[
 (\mu_i+\mu_j)C_{ij}=-X_{ij}.
\]
This determines every entry with positive denominator.  If both eigenvalues
vanish, Lemma~\ref{lem:V116-harmonic-numerator} gives \(X_{ij}=0\), so the
equation is \(0=0\).  The declared gauge fixes the undetermined block to zero.
Self-adjointness follows from \(X=X^\dagger\) and the symmetric real
denominator.
\end{proof}

The natural connection on the harmonic bundle is the Kato connection
\begin{equation}
 \nabla_r^{\mathcal H}:=\Pi_r\partial_r
 \quad\text{on sections of }\mathcal H_r^p.
 \label{eq:V116-Kato-connection}
\end{equation}
It records the finite normal transport which the positive square-root
parametrix does not see.

\begin{proposition}[Boundary harmonic BRST content]
\label{prop:V116-harmonic-BRST}
For a connected boundary,
\begin{equation}
 \mathcal H^0(\Sigma;\mathfrak g_{\rm SM})
 \cong\mathfrak g_{\rm SM}
 \label{eq:V116-harmonic-zero}
\end{equation}
consists of constant adjoint ghosts and satisfies
\(d_\Sigma\mathcal H^0=0\).  These are finite global-gauge zero modes and
must be quotiented or retained as a finite ghost cohomology block.
The space
\begin{equation}
 \mathcal H^1(\Sigma;\mathfrak g_{\rm SM})
 \cong H^1(\Sigma;\mathbb R)\otimes\mathfrak g_{\rm SM}
 \label{eq:V116-harmonic-one}
\end{equation}
contains boundary flat/toron directions which are not images of harmonic
scalar ghosts at linearized principal order.  They require a separate
finite-dimensional physical/topological decision.
\end{proposition}

\begin{proof}
Hodge theory identifies harmonic forms with de Rham cohomology.  On a
connected compact boundary, harmonic functions are constant, so their
exterior derivative is zero.  Tensoring with the adjoint fibre gives
\eqref{eq:V116-harmonic-zero}.  The degree-one identification is the same
Hodge theorem, and exactness of a harmonic scalar derivative gives no
nonzero degree-one vector.
\end{proof}

\subsection{Principal complementing and the remaining exact incidence}
\label{subsec:V116-principal-status}

\begin{proposition}[Order-zero correction does not change the V115 principal
floor]
\label{prop:V116-principal-floor}
The corrected operator \(L_r=\Lambda_r+C_r+\Psi^{-1}+\cdots\) has the same
order-one principal symbol \(|\xi|_{\gamma_r}\) as \(\Lambda_r\).  Therefore
the V115 parameter-stable lower bound remains the principal complementing
bound for the nonproduct corrected graph.  This statement does not provide a
uniform low-frequency inverse on the harmonic block.
\end{proposition}

\begin{proof}
Theorem~\ref{thm:V116-first-correction} places \(C_r\) in order zero, and the
recursion lowers every later term.  Principal boundary symbols ignore those
orders, so Theorem~\ref{thm:V115-parameter-floor} applies unchanged.  Harmonic
modes have \(\xi=0\) and are outside that homogeneous principal argument.
\end{proof}

\begin{firewall}[Scalar Riccati factorization is not yet the full de Rham
Calder\'on projector]
\label{fire:V116-deRham-open}
V116 derives the scalar fixed-measure Riccati correction and the abstract
Hodge harmonic splitting.  For the gauge potential, a nonproduct collar
couples tangential and normal form components through second fundamental
form and connection terms.  Exact BRST incidence requires the block
Dirichlet-to-Neumann operator for the full de Rham gauge complex, not separate
scalar copies of \eqref{eq:V116-corrected-root}.  That block intertwining is
not claimed here.
\end{firewall}

\subsection{V116 ledger and next acquisition}
\label{subsec:V116-ledger}

\paragraph{New exact conclusions.}
\begin{enumerate}[label=\textup{(\arabic*)},leftmargin=2.8em]
\item Proposition~\ref{prop:V116-fixed-measure} removes the varying boundary
      density and prints the exact scalar normal operator.
\item Lemma~\ref{lem:V116-factorization-defect} identifies the
      Dirichlet-to-Neumann Riccati equation.
\item Corollary~\ref{cor:V116-first-defect} proves that the instantaneous
      square root fails exactly by \(\partial_r\Lambda_r\).
\item Theorem~\ref{thm:V116-first-correction} constructs the unique
      self-adjoint order-zero correction and computes its remaining defect.
\item Corollary~\ref{cor:V116-recursion} gives the complete formal symbolic
      recursion.
\item Theorem~\ref{thm:V116-harmonic-Sylvester} separates the harmonic
      zero-denominator block and proves its compatibility.
\item Proposition~\ref{prop:V116-harmonic-BRST} identifies global-gauge and
      boundary toron modes as distinct finite sectors.
\end{enumerate}

\paragraph{Next forced step.}
V117 should write the nonproduct Hodge--de Rham Laplacian in tangential/normal
block form, including second fundamental form and connection terms, and
derive the corresponding block Riccati equation.  It must test the exact
intertwining of the scalar ghost Calder\'on operator with the one-form
Calder\'on operator.  The Kato harmonic connection and its finite BRST
quotient must be included before any determinant or heat coefficient is
formed.

\begin{assessment}[V116 continuum status]
\label{ass:V116-continuum-status}
The full Standard--Model--Einstein continuum remains unproved.  V116
constructs the first nonproduct scalar Dirichlet-to-Neumann correction and
isolates the harmonic boundary sectors, while preserving the principal
parameter floor.  The full de Rham block factorization, curved BRST
intertwining, nonlocal heat/stress coefficients, the fermion branch
selection, nonlinear existence, and the cofinal interacting limit remain
open.
\end{assessment}

\section*{V116 exact checks and append-only record}
\addcontentsline{toc}{section}{V116 exact checks and append-only record}

The decisive Riccati and harmonic checks are
\begin{equation}
 \boxed{
 \partial_rL+L^2=\mathcal A,
 \qquad
 \Lambda C+C\Lambda=-\partial_r\Lambda,
 \qquad
 \Pi(\partial_r\Lambda)\Pi=0.}
 \label{eq:V116-final-check}
\end{equation}
The append operation preserves the complete V115 prefix byte for byte before
its former live terminator.  The assembled V116 file is checked for one live
final terminator, balanced environments and braces, duplicate labels, newly
introduced unresolved references, display math delimiters, and control
characters before the exact V115 predecessor is removed.

% =============================================================================
% END APPENDED VERSION V116 MATERIAL
% =============================================================================

\clearpage
% =============================================================================
% BEGIN APPENDED VERSION V117 MATERIAL
% =============================================================================
\section{V117: the exact scalar curvature subprincipal symbol and the first
curved ghost factor}
\label{sec:V117-subprincipal}

V116 derived the nonproduct Riccati equation and its abstract order-zero
Sylvester correction.  V117 computes that correction at the scalar principal
symbol level.  In boundary dimension three, the result is
\[
 \frac12K(\widehat\xi,\widehat\xi)-\frac12H
\]
in the original boundary density, where \(K\) is the second fundamental form,
\(H=\operatorname{tr}_\gamma K\), and
\(\widehat\xi=\xi/|\xi|_\gamma\).  This is the first local curvature datum in
the otherwise nonlocal Dirichlet-to-Neumann operator.

The calculation cancels the complete order-one nonproduct factorization
defect of the scalar ghost block.  It does not by itself prove the full
one-form BRST intertwining; the Hodge block also contains shape endomorphisms
and normal/tangential mixing.

\subsection{Geometric conventions and the root derivative}
\label{subsec:V117-geometric-conventions}

Retain the inward Gaussian collar
\(g=dr^2+\gamma_r\) and define
\begin{equation}
 K_{ij}:=\frac12\partial_r\gamma_{ij},
 \qquad
 H:=\gamma^{ij}K_{ij}.
 \label{eq:V117-KH}
\end{equation}
With the density convention of V116,
\begin{equation}
 h_r=\partial_r\log J_r=H_r.
 \label{eq:V117-h-equals-H}
\end{equation}
Let
\begin{equation}
 \lambda_1(r,y,\xi):=|\xi|_{\gamma_r}
 =\bigl(\gamma_r^{ij}\xi_i\xi_j\bigr)^{1/2}
 \label{eq:V117-principal-root}
\end{equation}
be the positive order-one symbol of the instantaneous square root
\(\Lambda_r=\mathcal A_r^{1/2}\).

\begin{lemma}[Normal derivative of the principal root]
\label{lem:V117-root-derivative}
For \(\xi\ne0\),
\begin{equation}
 \boxed{
 \partial_r\lambda_1
 =
 -\frac{K^{ij}\xi_i\xi_j}{\lambda_1}.}
 \label{eq:V117-root-derivative}
\end{equation}
\end{lemma}

\begin{proof}
Differentiate
\(\lambda_1^2=\gamma^{ij}\xi_i\xi_j\).  From
\(\partial_r\gamma_{ij}=2K_{ij}\), differentiation of
\(\gamma^{ik}\gamma_{kj}=\delta^i_j\) gives
\[
 \partial_r\gamma^{ij}=-2K^{ij}.
\]
Therefore
\[
 2\lambda_1\partial_r\lambda_1
 =-2K^{ij}\xi_i\xi_j,
\]
which proves \eqref{eq:V117-root-derivative}.
\end{proof}

\subsection{Fixed-measure and raw-density subprincipal terms}
\label{subsec:V117-two-subprincipals}

Let \(c_0^{\rm fix}\) be the scalar order-zero principal symbol of the V116
correction \(C_r\) on the fixed boundary Hilbert space.

\begin{theorem}[Exact first Sylvester symbol]
\label{thm:V117-fixed-correction}
The order-one part of the Sylvester equation
\[
 \Lambda C+C\Lambda=-\partial_r\Lambda
\]
is the scalar identity
\begin{equation}
 2\lambda_1c_0^{\rm fix}=-\partial_r\lambda_1.
 \label{eq:V117-symbol-Sylvester}
\end{equation}
Hence
\begin{equation}
 \boxed{
 c_0^{\rm fix}(r,y,\xi)
 =
 \frac{K^{ij}\xi_i\xi_j}{2|\xi|_{\gamma_r}^2}
 =
 \frac12K(\widehat\xi,\widehat\xi).}
 \label{eq:V117-fixed-correction}
\end{equation}
\end{theorem}

\begin{proof}
At order one, composition of the scalar order-one symbol \(\lambda_1\) with
the scalar order-zero symbol \(c_0^{\rm fix}\) is ordinary multiplication;
Poisson-bracket corrections have order zero and enter the next Riccati row.
Thus the two Sylvester products contribute
\(2\lambda_1c_0^{\rm fix}\).  Insert
Lemma~\ref{lem:V117-root-derivative} and divide by \(2\lambda_1\).
\end{proof}

Undoing the density conjugation uses
\eqref{eq:V116-raw-corrected-root}.  Scalar conjugation does not change the
order-zero principal symbol of \(C_r\), while it subtracts \(h_r/2=H/2\).

\begin{corollary}[Exact raw scalar subprincipal symbol]
\label{cor:V117-raw-subprincipal}
In the original \(L^2(\Sigma,d{\rm vol}_{\gamma_r})\) frame, the corrected
positive Dirichlet-to-Neumann symbol begins
\begin{equation}
 \boxed{
 \ell_{\rm raw}(r,y,\xi)
 =
 |\xi|_{\gamma_r}
 +\frac12K(\widehat\xi,\widehat\xi)
 -\frac12H
 +O(|\xi|^{-1}).}
 \label{eq:V117-raw-subprincipal}
\end{equation}
The sign is tied to the inward-normal convention
\eqref{eq:V117-KH}; reversing the normal changes both \(K\) and \(H\).
\end{corollary}

\begin{proof}
Combine \eqref{eq:V117-fixed-correction} with
\[
 L^{\rm raw}=U^{-1}LU-\frac12h
\]
from Proposition~\ref{prop:V116-corrected-action} and use
\eqref{eq:V117-h-equals-H}.  The next Riccati correction has order \(-1\).
\end{proof}

\subsection{Trace and anisotropic curvature channels}
\label{subsec:V117-curvature-split}

In boundary dimension three, decompose
\begin{equation}
 K=\frac H3\gamma+\mathring K,
 \qquad
 \operatorname{tr}_\gamma\mathring K=0.
 \label{eq:V117-K-split}
\end{equation}

\begin{proposition}[Exact mean/shear split]
\label{prop:V117-mean-shear}
The raw order-zero symbol in
\eqref{eq:V117-raw-subprincipal} is
\begin{equation}
 c_0^{\rm raw}(\widehat\xi)
 =
 -\frac H3+\frac12\mathring K(\widehat\xi,\widehat\xi).
 \label{eq:V117-mean-shear-symbol}
\end{equation}
Its angular average on the unit cotangent sphere is
\begin{equation}
 \fint_{\mathbb S_y^2}
 c_0^{\rm raw}(\widehat\xi)\,d\widehat\xi
 =-\frac H3.
 \label{eq:V117-angular-average}
\end{equation}
Moreover,
\begin{equation}
 -\frac H3+\frac12\lambda_{\min}(\mathring K)
 \le c_0^{\rm raw}(\widehat\xi)
 \le
 -\frac H3+\frac12\lambda_{\max}(\mathring K).
 \label{eq:V117-directional-bounds}
\end{equation}
\end{proposition}

\begin{proof}
Insert \eqref{eq:V117-K-split} into
\(\tfrac12K(\widehat\xi,\widehat\xi)-\tfrac12H\).
Since
\[
 \fint_{\mathbb S_y^2}
 \widehat\xi_i\widehat\xi_j\,d\widehat\xi
 =\frac13\gamma_{ij},
\]
the trace-free term has zero average.  The last bounds are the Rayleigh
quotient bounds for the symmetric endomorphism \(\mathring K\).
\end{proof}

\begin{corollary}[Umbilic and totally geodesic checks]
\label{cor:V117-geometric-checks}
If \(K=(H/3)\gamma\), then
\begin{equation}
 c_0^{\rm raw}=-\frac H3.
 \label{eq:V117-umbilic}
\end{equation}
If the boundary is totally geodesic, \(K=0\), then
\begin{equation}
 c_0^{\rm raw}=0,
 \label{eq:V117-totally-geodesic}
\end{equation}
recovering the V115 product symbol through order zero.
\end{corollary}

\begin{proof}
These are the two specializations of
\eqref{eq:V117-mean-shear-symbol}.
\end{proof}

\subsection{The corrected scalar ghost factor}
\label{subsec:V117-ghost-factor}

Let \(L_{\rm raw}^{[0]}\) be a self-adjoint pseudodifferential operator whose
order-one and order-zero symbols are those in
\eqref{eq:V117-raw-subprincipal}, and define
\begin{equation}
 B_{\rm gh}^{[0]}
 :=
 -\partial_r+L_{\rm raw}^{[0]}.
 \label{eq:V117-ghost-boundary-factor}
\end{equation}

\begin{theorem}[Cancellation of the complete order-one curved defect]
\label{thm:V117-order-one-cancellation}
After fixed-measure conjugation, the factorization using
\(L^{[0]}=\Lambda+C\) satisfies
\begin{equation}
 (\partial_r+L^{[0]})(-\partial_r+L^{[0]})
 -\widetilde P
 =
 \partial_rC+C^2\in\Psi^0(\Sigma).
 \label{eq:V117-order-zero-remainder}
\end{equation}
Thus the curvature symbol
\eqref{eq:V117-raw-subprincipal} cancels the entire order-one nonproduct
defect.  If a ghost satisfies
\begin{equation}
 \widetilde Pc=0,
 \qquad
 (-\partial_r+L^{[0]})c|_\Sigma=0,
 \label{eq:V117-truncated-ghost-domain}
\end{equation}
then its missing normal boundary defect is an order-zero residual:
\begin{equation}
 \partial_r\bigl((-\partial_r+L^{[0]})c\bigr)\big|_\Sigma
 =
 -L^{[0]}(-\partial_r+L^{[0]})c\big|_\Sigma
 -(\partial_rC+C^2)c\big|_\Sigma,
 \label{eq:V117-normal-residual}
\end{equation}
and hence equals
\(-(\partial_rC+C^2)c|_\Sigma\).
\end{theorem}

\begin{proof}
The factorization identity is
\eqref{eq:V116-second-defect}.  The boundary row kills the first term on the
right of \eqref{eq:V117-normal-residual}.  Apply the factorization to the
ghost Euler equation and solve for the normal derivative of the boundary
defect.  The remaining operator has order zero by
Theorem~\ref{thm:V116-first-correction}.
\end{proof}

\begin{firewall}[Order-zero residual is not exact BRST closure]
\label{fire:V117-residual-not-zero}
Equation \eqref{eq:V117-normal-residual} quantifies the remaining curved
normal BRST defect; it does not set it to zero.  The order-\(-1\) correction
from \eqref{eq:V116-second-Sylvester} cancels its order-zero symbol, and the
full Riccati/Calder\'on operator is required for exact closure.  Stopping at
\eqref{eq:V117-raw-subprincipal} would prove principal and subprincipal
incidence only.
\end{firewall}

\subsection{Parameter floor and action status}
\label{subsec:V117-floor-action}

\begin{proposition}[Lower-order curvature cannot change the principal
parameter floor]
\label{prop:V117-principal-floor}
The order-one principal boundary symbol remains
\(\lambda(\xi,z)+|\xi|_\gamma\).  Hence the normalized lower bound
\begin{equation}
 |\lambda+|\xi|_\gamma|\ge2^{-1/2}
 \label{eq:V117-principal-floor}
\end{equation}
from V115 is unchanged.  The finite-frequency inverse may depend on
\(H,\mathring K\), the harmonic block, and the lower Riccati terms.
\end{proposition}

\begin{proof}
Equation \eqref{eq:V117-raw-subprincipal} changes only order zero and lower
symbols.  Parameter complementing is determined by the homogeneous order-one
boundary symbol, so Theorem~\ref{thm:V115-parameter-floor} applies.
\end{proof}

\begin{proposition}[The subprincipal graph is action-compatible]
\label{prop:V117-action}
The real scalar symbol
\[
 \frac12K(\widehat\xi,\widehat\xi)-\frac12H
\]
is even in \(\xi\) and is the principal symbol of a self-adjoint order-zero
operator.  Choosing the full Sylvester correction by
\eqref{eq:V116-Sylvester-solution} makes the truncated boundary graph
Lagrangian.
\end{proposition}

\begin{proof}
The displayed expression is real and quadratic in
\(\widehat\xi\), hence even.  The operator construction in
Theorem~\ref{thm:V116-first-correction} is self-adjoint, and unitary
density conjugation preserves self-adjointness after the real
\(-H/2\) shift.  A self-adjoint conormal graph is Lagrangian.
\end{proof}

\begin{firewall}[The scalar symbol is not the one-form symbol]
\label{fire:V117-oneform-open}
For tangential and normal components of a one-form, the Hodge Laplacian has a
matrix normal form.  Its order-zero Dirichlet-to-Neumann symbol contains the
second fundamental form acting on form indices, not only the scalar
directional contraction in \eqref{eq:V117-raw-subprincipal}.  V117 therefore
does not reuse the scalar correction componentwise; doing so would repeat the
normal BRST error corrected in V115.
\end{firewall}

\subsection{V117 ledger and next acquisition}
\label{subsec:V117-ledger}

\paragraph{New exact conclusions.}
\begin{enumerate}[label=\textup{(\arabic*)},leftmargin=2.8em]
\item Lemma~\ref{lem:V117-root-derivative} computes the normal variation of
      the positive principal root from \(K\).
\item Theorem~\ref{thm:V117-fixed-correction} solves the scalar symbol
      Sylvester equation exactly.
\item Corollary~\ref{cor:V117-raw-subprincipal} adds the density correction
      and gives the complete raw order-zero symbol.
\item Proposition~\ref{prop:V117-mean-shear} separates its mean-curvature and
      trace-free directional channels.
\item Theorem~\ref{thm:V117-order-one-cancellation} proves cancellation of
      the full order-one factorization defect and prints the remaining
      order-zero normal BRST residual.
\item Propositions~\ref{prop:V117-principal-floor} and
      \ref{prop:V117-action} preserve the principal complementing and action
      gates.
\end{enumerate}

\paragraph{Next forced step.}
V118 should derive the tangential/normal block normal form of the Hodge
one-form Laplacian in the convention \eqref{eq:V117-KH}.  It must compute the
matrix order-zero Sylvester symbol, including the shape operator on one-form
indices, and test the boundary de Rham intertwining against the scalar symbol
\eqref{eq:V117-raw-subprincipal}.  Any unmatched term is the exact curved
BRST correction to be carried by the block Calder\'on graph.

\begin{assessment}[V117 continuum status]
\label{ass:V117-continuum-status}
The full Standard--Model--Einstein continuum remains unproved.  V117 computes
the exact first curvature correction to the scalar ghost
Dirichlet-to-Neumann operator and lowers its factorization defect by one full
pseudodifferential order.  The one-form block, exact curved BRST
intertwining, harmonic finite complex, nonlocal heat/stress coefficients,
fermion branch selection, nonlinear existence, and the cofinal interacting
limit remain open.
\end{assessment}

\section*{V117 exact checks and append-only record}
\addcontentsline{toc}{section}{V117 exact checks and append-only record}

The exact scalar curvature checksum is
\begin{equation}
 \boxed{
 \partial_r|\xi|_\gamma
 =-\frac{K^{ij}\xi_i\xi_j}{|\xi|_\gamma},
 \qquad
 c_0^{\rm raw}
 =\frac12K(\widehat\xi,\widehat\xi)-\frac12H.}
 \label{eq:V117-final-check}
\end{equation}
The append operation preserves the complete V116 prefix byte for byte before
its former live terminator.  The assembled V117 file is checked for one live
final terminator, balanced environments and braces, duplicate labels, newly
introduced unresolved references, display math delimiters, and control
characters before the exact V116 predecessor is removed.

% =============================================================================
% END APPENDED VERSION V117 MATERIAL
% =============================================================================

\clearpage
% =============================================================================
% BEGIN APPENDED VERSION V118 MATERIAL
% =============================================================================
\section{V118: the curved Hodge one-form collar block and subprincipal
de Rham incidence}
\label{sec:V118-hodge-block}

V117 computed the scalar curvature correction but left open the
normal--tangential coupling of a gauge one-form.  This note derives that
coupling directly from \(d\) and \(\delta\).  The resulting order-zero
Dirichlet-to-Neumann matrix contains both the shape operator and an
off-diagonal skew-Hermitian incidence pair.  These terms are forced: they
cancel the apparent curvature defect in the normal component of the
boundary gradient of a scalar ghost.

All formulas use the inward Gaussian collar and the convention
\(K=\tfrac12\partial_r\gamma\) fixed in
\eqref{eq:V117-KH}.  Write
\begin{equation}
 S_i{}^j:=K_i{}^j,
 \qquad
 (S\alpha)_i:=S_i{}^j\alpha_j
 \label{eq:V118-shape}
\end{equation}
for the shape endomorphism on tangential one-forms, and let \(D\),
\(d_\Sigma\), and \(\delta_\Sigma\) be the Levi--Civita derivative and the
de Rham operators of \((\Sigma,\gamma_r)\).

\subsection{Exact \(d\)-\(\delta\) collar algebra}
\label{subsec:V118-d-delta}

Decompose a one-form and a two-form as
\begin{equation}
 a=\alpha+\phi\,dr,
 \qquad
 \omega=\beta+dr\wedge\eta,
 \label{eq:V118-splittings}
\end{equation}
where \(\alpha,\eta\) are tangential one-forms and \(\beta\) is a
tangential two-form.

\begin{lemma}[Exact collar formulas for \(d\) and \(\delta\)]
\label{lem:V118-d-delta}
The decompositions \eqref{eq:V118-splittings} obey
\begin{align}
 da
 &=
 d_\Sigma\alpha
 +dr\wedge(\partial_r\alpha-d_\Sigma\phi),
 \label{eq:V118-d-oneform}\\
 \delta a
 &=
 \delta_\Sigma\alpha-\partial_r\phi-H\phi,
 \label{eq:V118-delta-oneform}\\
 (\delta\omega)_{\rm tan}
 &=
 \delta_\Sigma\beta-\partial_r\eta-H\eta+2S\eta,
 \label{eq:V118-delta-twoform-tan}\\
 (\delta\omega)_{\rm nor}
 &=
 -\delta_\Sigma\eta.
 \label{eq:V118-delta-twoform-nor}
\end{align}
\end{lemma}

\begin{proof}
In Gaussian coordinates the only mixed Christoffel symbols needed here are
\begin{equation}
 \Gamma^r_{ij}=-K_{ij},
 \qquad
 \Gamma^k_{ri}=S_i{}^k.
 \label{eq:V118-Christoffel}
\end{equation}
The exterior derivative is metric-independent, so
\eqref{eq:V118-d-oneform} follows by differentiating the two summands of
\(a\).  Next,
\[
 -\nabla^\mu a_\mu
 =
 -\gamma^{ij}(D_i\alpha_j+K_{ij}\phi)-\partial_r\phi,
\]
which is \eqref{eq:V118-delta-oneform}.

For the tangential component of \(\delta\omega\), use
\(\omega_{ri}=\eta_i\).  Direct substitution of
\eqref{eq:V118-Christoffel} gives
\[
 -\nabla^r\omega_{rj}
 =-\partial_r\eta_j+S_j{}^k\eta_k
\]
and
\[
 -\gamma^{ik}\nabla_i\omega_{kj}
 =
 (\delta_\Sigma\beta)_j-H\eta_j+S_j{}^k\eta_k.
\]
Their sum proves \eqref{eq:V118-delta-twoform-tan}.  Finally,
\[
 -\gamma^{ik}\nabla_i\omega_{kr}
 =
 D^i\eta_i+S^{ik}\beta_{ik}
 =
 -\delta_\Sigma\eta,
\]
because \(S^{ik}\) is symmetric and \(\beta_{ik}\) is antisymmetric.
\end{proof}

\subsection{The raw tangential--normal Hodge block}
\label{subsec:V118-raw-block}

Define the first-order tangential operator
\begin{equation}
 \mathcal C_K\alpha
 :=
 2K^{ij}D_i\alpha_j
 +(2D^iK_i{}^j-D^jH)\alpha_j.
 \label{eq:V118-CK}
\end{equation}

\begin{theorem}[Exact raw one-form normal form]
\label{thm:V118-raw-Hodge-block}
For the positive Hodge Laplacian
\(P_1=d\delta+\delta d\), one has
\begin{equation}
 \boxed{
 P_1
 \binom{\alpha}{\phi}
 =
 \begin{pmatrix}
 -\partial_r^2+(2S-HI)\partial_r+\Delta_{\Sigma,1}
 &
 -2S\,d_\Sigma-(d_\Sigma H)
 \\[2mm]
 \mathcal C_K
 &
 -\partial_r^2-H\partial_r+\Delta_{\Sigma,0}-\partial_rH
 \end{pmatrix}
 \binom{\alpha}{\phi}.}
 \label{eq:V118-raw-Hodge-block}
\end{equation}
Here \((d_\Sigma H)\) in the upper-right entry denotes multiplication by the
one-form \(d_\Sigma H\).
\end{theorem}

\begin{proof}
Put
\[
 \psi:=\delta_\Sigma\alpha-\partial_r\phi-H\phi,
 \qquad
 \eta:=\partial_r\alpha-d_\Sigma\phi.
\]
The tangential component of \(d\delta a+\delta da\), using
Lemma~\ref{lem:V118-d-delta}, is
\[
 d_\Sigma\psi+\delta_\Sigma d_\Sigma\alpha
 -\partial_r\eta-H\eta+2S\eta.
\]
The two occurrences of \(d_\Sigma\partial_r\phi\) cancel.  The two
mean-curvature gradient terms combine as
\[
 -d_\Sigma(H\phi)+H\,d_\Sigma\phi
 =-\phi\,d_\Sigma H.
\]
This proves the upper row of \eqref{eq:V118-raw-Hodge-block}.

The normal component is
\[
 \partial_r\psi-\delta_\Sigma\eta
 =
 -\partial_r^2\phi-H\partial_r\phi-(\partial_rH)\phi
 +\Delta_{\Sigma,0}\phi
 +[\partial_r,\delta_\Sigma]\alpha.
\]
Differentiating
\(\delta_\Sigma\alpha=-\gamma^{ij}D_i\alpha_j\), and using
\(\partial_r\gamma^{ij}=-2K^{ij}\), gives
\[
 [\partial_r,\delta_\Sigma]\alpha
 =
 2K^{ij}D_i\alpha_j
 +\gamma^{ij}(\partial_r\Gamma^k_{ij})\alpha_k.
\]
The standard metric-variation identity, which also follows by
differentiating the Christoffel formula, is
\[
 \gamma^{ij}\partial_r\Gamma^k_{ij}
 =
 2D^iK_i{}^k-D^kH.
\]
Thus the commutator is exactly \(\mathcal C_K\), proving the lower row.
\end{proof}

\begin{corollary}[The off-diagonal entries are formal adjoints]
\label{cor:V118-offdiagonal-adjoints}
At fixed \(r\),
\begin{equation}
 \bigl(-2S\,d_\Sigma-(d_\Sigma H)\bigr)^\dagger
 =
 \mathcal C_K.
 \label{eq:V118-offdiagonal-adjoints}
\end{equation}
\end{corollary}

\begin{proof}
For a tangential one-form \(\alpha\),
\[
 \bigl(-2S\,d_\Sigma\bigr)^\dagger\alpha
 =-2\delta_\Sigma(S\alpha)
 =
 2K^{ij}D_i\alpha_j+2(D^iK_i{}^j)\alpha_j.
\]
The adjoint of multiplication by \(-d_\Sigma H\) is contraction by
\(-d_\Sigma H\).  Adding the two expressions gives
\eqref{eq:V118-CK}.
\end{proof}

\subsection{The unitary normal connection and Schrödinger form}
\label{subsec:V118-unitary-connection}

The scalar density conjugation of V116 is only one part of the one-form
normal trivialization: the tangential covector metric also varies.  Define
\begin{equation}
 \mathscr D_r
 \binom{\alpha}{\phi}
 :=
 \binom{
 \partial_r\alpha+(\tfrac12H I-S)\alpha
 }{
 \partial_r\phi+\tfrac12H\phi
 }.
 \label{eq:V118-unitary-normal-connection}
\end{equation}

\begin{lemma}[Metric compatibility of \(\mathscr D_r\)]
\label{lem:V118-unitary}
For compactly supported boundary fields \(a,b\),
\begin{equation}
 \partial_r\langle a,b\rangle_{L^2(\Sigma,\gamma_r)}
 =
 \langle\mathscr D_ra,b\rangle_{L^2(\Sigma,\gamma_r)}
 +
 \langle a,\mathscr D_rb\rangle_{L^2(\Sigma,\gamma_r)}.
 \label{eq:V118-unitary-identity}
\end{equation}
\end{lemma}

\begin{proof}
For tangential covectors,
\[
 \partial_r\langle\alpha,\widetilde\alpha\rangle_\gamma
 =
 \langle\partial_r\alpha,\widetilde\alpha\rangle_\gamma
 +\langle\alpha,\partial_r\widetilde\alpha\rangle_\gamma
 -2\langle S\alpha,\widetilde\alpha\rangle_\gamma.
\]
The boundary density contributes
\(\partial_r d{\rm vol}_{\gamma_r}=H\,d{\rm vol}_{\gamma_r}\).
Splitting that last contribution equally between the two arguments gives
the tangential row of \eqref{eq:V118-unitary-normal-connection}.  The normal
coefficient has no inverse-metric variation, so only the two \(H/2\)
density terms remain.
\end{proof}

Let
\begin{equation}
 A_T:=\frac12H I-S.
 \label{eq:V118-AT}
\end{equation}

\begin{proposition}[Exact matrix Schrödinger reduction]
\label{prop:V118-Schrodinger}
In a parallel trivialization for the unitary connection
\(\mathscr D_r\), the one-form operator is
\begin{equation}
 P_1=-\mathscr D_r^2+\mathscr A_{1,r},
 \label{eq:V118-Schrodinger}
\end{equation}
where, in the raw splitting,
\begin{equation}
 \boxed{
 \mathscr A_{1,r}
 =
 \begin{pmatrix}
 \Delta_{\Sigma,1}+\partial_rA_T+A_T^2
 &
 -2S\,d_\Sigma-(d_\Sigma H)
 \\[1mm]
 \mathcal C_K
 &
 \Delta_{\Sigma,0}-\tfrac12\partial_rH+\tfrac14H^2
 \end{pmatrix}.}
 \label{eq:V118-A1}
\end{equation}
The off-diagonal order-one symbol is Hermitian, and
\(\mathscr A_{1,r}\) is formally self-adjoint in the fixed Hilbert
trivialization.
\end{proposition}

\begin{proof}
Expanding
\(-(\partial_r+A_T)^2\) gives the upper-left normal derivative terms in
\eqref{eq:V118-raw-Hodge-block} together with the zeroth-order term
\(-\partial_rA_T-A_T^2\).  Moving the latter term into
\(\mathscr A_{1,r}\) gives the first diagonal entry of
\eqref{eq:V118-A1}.  Similarly,
\[
 -(\partial_r+H/2)^2
 =
 -\partial_r^2-H\partial_r-\frac12\partial_rH-\frac14H^2,
\]
which gives the second diagonal entry.  The off-diagonal entries do not
contain \(\partial_r\) and are unchanged.  Formal self-adjointness follows
from Lemma~\ref{lem:V118-unitary},
Corollary~\ref{cor:V118-offdiagonal-adjoints}, and formal self-adjointness of
the Hodge Laplacian.
\end{proof}

\subsection{The complete extrinsic order-zero matrix}
\label{subsec:V118-order-zero-matrix}

Fix \(r\) and a point \(y\).  Choose geodesic tangential coordinates and a
parallel orthonormal form frame at \(y\).  For \(\xi\ne0\), set
\begin{equation}
 \rho:=|\xi|_\gamma,
 \qquad
 n:=\frac{\xi}{\rho},
 \qquad
 k:=K(n,n),
 \qquad
 v:=Sn.
 \label{eq:V118-rho-n-k-v}
\end{equation}
We call a symbol term \emph{extrinsic} if it vanishes after setting
\(K,H\), and their normal evolution to zero while keeping the instantaneous
boundary metric fixed.  This removes the ordinary intrinsic subprincipal
symbol, which is already present in a product collar.

\begin{theorem}[Extrinsic order-zero one-form Calderón symbol]
\label{thm:V118-oneform-subprincipal}
Let \(L_{1}^{\rm raw}\) denote the positive microlocal
Dirichlet-to-Neumann branch for \(P_1\).  Its complete extrinsic
order-zero symbol is
\begin{equation}
 \boxed{
 \ell_{1}^{\rm raw}(y,\xi)
 =
 \rho I_4+
 \begin{pmatrix}
 s_0 I_{T^*\Sigma}+S & -i\,v\\[1mm]
 i\,v^\dagger & s_0
 \end{pmatrix}
 +O(\rho^{-1}),}
 \qquad
 s_0:=\frac12(k-H).
 \label{eq:V118-oneform-subprincipal}
\end{equation}
The upper-right entry maps a scalar to a tangential one-form, and the
lower-left entry is its Hermitian adjoint.
\end{theorem}

\begin{proof}
The order-two symbol of \(\mathscr A_{1,r}\) is \(\rho^2I_4\).  By
\eqref{eq:V118-A1}, its extrinsic order-one symbol is
\begin{equation}
 a_1^K
 =
 \begin{pmatrix}
 0&-2i\rho\,v\\
 2i\rho\,v^\dagger&0
 \end{pmatrix}.
 \label{eq:V118-a1K}
\end{equation}
If the instantaneous positive square root has extrinsic expansion
\(\rho I_4+q_0^K+O(\rho^{-1})\), the order-one part of its square is
\(2\rho q_0^K\).  Therefore
\begin{equation}
 q_0^K
 =
 \begin{pmatrix}
 0&-i\,v\\
 i\,v^\dagger&0
 \end{pmatrix}.
 \label{eq:V118-q0K}
\end{equation}

The Riccati correction is independent of the form index at this order.
Indeed, the order-one Sylvester row is
\[
 2\rho c_0^K=-\partial_r\rho.
\]
Lemma~\ref{lem:V117-root-derivative} gives
\begin{equation}
 c_0^K=\frac12k\,I_4.
 \label{eq:V118-c0K}
\end{equation}
Finally, returning from the unitary normal connection
\eqref{eq:V118-unitary-normal-connection} to raw coefficient derivatives
subtracts
\[
 \begin{pmatrix}
 \tfrac12H I-S&0\\0&\tfrac12H
 \end{pmatrix}
\]
from the boundary root.  Adding this shift to
\eqref{eq:V118-q0K} and \eqref{eq:V118-c0K} proves
\eqref{eq:V118-oneform-subprincipal}.
\end{proof}

\begin{corollary}[Action compatibility]
\label{cor:V118-action}
The matrix in \eqref{eq:V118-oneform-subprincipal} is Hermitian.  Hence the
one-form Calderón graph is Lagrangian through extrinsic order zero; the exact
self-adjoint Calderón graph is action-compatible whenever the harmonic
finite sector is split off.
\end{corollary}

\begin{proof}
The shape operator is self-adjoint, \(s_0\) is real, and
\((-iv)^\dagger=iv^\dagger\).  The graph of a self-adjoint conormal operator
is Lagrangian.
\end{proof}

\subsection{Subprincipal de Rham incidence}
\label{subsec:V118-incidence}

For the scalar raw root of V117, write
\begin{equation}
 \ell_0^{\rm raw}=\rho+s_0+O(\rho^{-1}),
 \qquad
 s_0=\frac12(k-H).
 \label{eq:V118-scalar-root}
\end{equation}
The equality of the two symbols called \(s_0\) is deliberate.

\begin{theorem}[Curved gradient incidence through order one]
\label{thm:V118-gradient-incidence}
Let the boundary trace of a scalar gradient be represented symbolically by
\begin{equation}
 \mathcal G_0(\xi)u
 :=
 \binom{i\rho n}{\rho+s_0}u.
 \label{eq:V118-gradient-column}
\end{equation}
Then \eqref{eq:V118-oneform-subprincipal} gives
\begin{align}
 \bigl(\ell_1^{\rm raw}\mathcal G_0\bigr)_{\rm tan}
 &=
 i\rho^2n+i\rho s_0n+O(1),
 \label{eq:V118-incidence-tan}\\
 \bigl(\ell_1^{\rm raw}\mathcal G_0\bigr)_{\rm nor}
 &=
 \rho^2-H\rho+O(1).
 \label{eq:V118-incidence-nor}
\end{align}
These are precisely the degree-two and degree-one symbols of
\begin{equation}
 \binom{d_\Sigma L_0^{\rm raw}}
 {\Delta_{\Sigma,0}-H L_0^{\rm raw}}u.
 \label{eq:V118-gradient-evolution}
\end{equation}
\end{theorem}

\begin{proof}
For the tangential row, put \(v=Sn\) and calculate
\begin{align*}
 &\bigl(\rho I+s_0I+S\bigr)(i\rho n)
 -iv(\rho+s_0)
 \\
 &\hspace{2cm}
 =
 i\rho^2n+i\rho s_0n
 +i\rho Sn-i\rho v+O(1)
 =
 i\rho^2n+i\rho s_0n+O(1).
\end{align*}
For the normal row,
\begin{align*}
 iv^\dagger(i\rho n)+(\rho+s_0)^2
 &=
 -\rho\,K(n,n)+\rho^2+2\rho s_0+O(1)
 \\
 &=
 \rho^2+\rho(2s_0-k)+O(1)
 =
 \rho^2-H\rho+O(1).
\end{align*}
The scalar Euler equation
\[
 -\partial_r^2f-H\partial_rf+\Delta_{\Sigma,0}f=0
\]
gives
\(\partial_r^2f=(\Delta_{\Sigma,0}-H\partial_r)f\).
Together with
\(\partial_r(d_\Sigma f)=d_\Sigma(\partial_rf)\), this identifies
\eqref{eq:V118-gradient-evolution}.
\end{proof}

\begin{corollary}[Why componentwise scalar roots fail]
\label{cor:V118-componentwise-failure}
Replacing the matrix correction in
\eqref{eq:V118-oneform-subprincipal} by the componentwise scalar correction
\(s_0I_4\) leaves the normal incidence defect
\begin{equation}
 \rho\,K(n,n)
 \label{eq:V118-componentwise-defect}
\end{equation}
at order one.  Thus the componentwise prescription is compatible only in
directions with \(K(n,n)=0\), and not on a generic curved boundary.
\end{corollary}

\begin{proof}
Without the lower-left entry \(iv^\dagger\), the normal row is
\(\rho^2+2\rho s_0+O(1)=\rho^2+(k-H)\rho+O(1)\).  Subtract
\eqref{eq:V118-incidence-nor}.
\end{proof}

\subsection{Exact Calderón-range identity and BRST consequence}
\label{subsec:V118-exact-range}

The previous theorem is a local symbol computation.  The corresponding
exact identity holds on the range of compatible Calderón maps.

\begin{proposition}[Exact de Rham identity on the scalar Calderón range]
\label{prop:V118-exact-incidence}
Suppose the positive scalar and one-form Calderón maps
\(L_0^{\rm raw}\) and \(L_1^{\rm raw}\) exist on a common collar branch, with
their finite-dimensional zero modes projected out.  Then
\begin{equation}
 \boxed{
 L_1^{\rm raw}
 \binom{d_\Sigma}{L_0^{\rm raw}}
 =
 \binom{
 d_\Sigma L_0^{\rm raw}
 }{
 \Delta_{\Sigma,0}-H L_0^{\rm raw}
 }.}
 \label{eq:V118-exact-incidence}
\end{equation}
\end{proposition}

\begin{proof}
Let \(u\) be scalar boundary data and let \(f\) be its selected scalar
Calderón extension.  Thus
\[
 P_0f=0,\qquad f|_\Sigma=u,\qquad
 \partial_rf|_\Sigma=L_0^{\rm raw}u.
\]
Since the Hodge Laplacian commutes with \(d\),
\(P_1(df)=d(P_0f)=0\).  Its raw boundary value is
\[
 (df)|_\Sigma
 =
 \binom{d_\Sigma u}{L_0^{\rm raw}u}.
\]
By definition, \(L_1^{\rm raw}\) maps this boundary value to the raw normal
coefficient derivative of \(df\).  The tangential coefficient derivative is
\(d_\Sigma L_0^{\rm raw}u\).  The normal coefficient derivative is
\(\partial_r^2f|_\Sigma\), which the scalar Euler equation identifies as
\((\Delta_{\Sigma,0}-H L_0^{\rm raw})u\).  This proves
\eqref{eq:V118-exact-incidence}.
\end{proof}

\begin{corollary}[Exact on-shell BRST closure of the paired graphs]
\label{cor:V118-BRST-closure}
Under the hypotheses of Proposition~\ref{prop:V118-exact-incidence}, define
\begin{equation}
 B_p:=-\partial_r+L_p^{\rm raw},
 \qquad p=0,1.
 \label{eq:V118-Bp}
\end{equation}
If a ghost satisfies \(P_0c=0\) and \(B_0c|_\Sigma=0\), then
\begin{equation}
 B_1(dc)|_\Sigma=0.
 \label{eq:V118-BRST-closure}
\end{equation}
\end{corollary}

\begin{proof}
The scalar boundary equation identifies the boundary trace of \(dc\) with
the column on the left of \eqref{eq:V118-exact-incidence}.  The right side
of that identity is exactly the raw normal derivative of the same one-form.
Therefore \((-\partial_r+L_1^{\rm raw})dc\) vanishes at the boundary.
\end{proof}

\begin{firewall}[Scope of the exact identity]
\label{fire:V118-exact-scope}
Proposition~\ref{prop:V118-exact-incidence} is an exact range identity once
the compatible Calderón branches exist.  V118 does not yet construct those
branches globally across spectral crossings, nor does it quotient the
harmonic scalar and one-form sectors.  Those are finite-dimensional
connection problems, not errors in the local matrix symbol
\eqref{eq:V118-oneform-subprincipal}.
\end{firewall}

\subsection{Parameter floor and next finite complex}
\label{subsec:V118-floor-ledger}

\begin{proposition}[The matrix branch preserves the principal floor]
\label{prop:V118-principal-floor}
The principal boundary symbol of the one-form graph is
\((\lambda+\rho)I_4\).  Hence the V115 normalized estimate remains
\begin{equation}
 \|(\lambda+\rho)^{-1}I_4\|
 \le \sqrt2.
 \label{eq:V118-principal-floor}
\end{equation}
All curvature entries in \eqref{eq:V118-oneform-subprincipal} are order zero
and do not alter parameter complementing.
\end{proposition}

\begin{proof}
Theorem~\ref{thm:V118-oneform-subprincipal} has principal part \(\rho I_4\).
Apply the scalar inequality \eqref{eq:V117-principal-floor} to each fibre
component.
\end{proof}

\paragraph{New exact conclusions.}
\begin{enumerate}[label=\textup{(\arabic*)},leftmargin=2.8em]
\item Lemma~\ref{lem:V118-d-delta} and
      Theorem~\ref{thm:V118-raw-Hodge-block} derive the complete raw
      tangential--normal one-form collar operator.
\item Proposition~\ref{prop:V118-Schrodinger} identifies the unitary normal
      connection and the self-adjoint matrix Schrödinger family.
\item Theorem~\ref{thm:V118-oneform-subprincipal} computes the complete
      extrinsic order-zero Calderón matrix.
\item Theorem~\ref{thm:V118-gradient-incidence} proves that its shape and
      off-diagonal entries give the required subprincipal de Rham
      cancellations.
\item Corollary~\ref{cor:V118-componentwise-failure} prints the exact
      order-one error made by a componentwise scalar prescription.
\item Proposition~\ref{prop:V118-exact-incidence} and
      Corollary~\ref{cor:V118-BRST-closure} isolate the exact Calderón-range
      identity that yields on-shell curved BRST closure.
\end{enumerate}

\paragraph{Next forced step.}
V119 should construct Kato parallel transport for the moving harmonic
projectors \(\Pi_{0,r}\) and \(\Pi_{1,r}\), prove that the induced finite
de Rham complex is preserved by the collar connection, and identify the
constant-gauge quotient separately from harmonic one-form moduli.  This is
the missing finite-dimensional hypothesis in
Proposition~\ref{prop:V118-exact-incidence}.  Only after that quotient is
fixed should the nonlocal heat and stress tensors be formed.

\begin{assessment}[V118 continuum status]
\label{ass:V118-continuum-status}
The full Standard--Model--Einstein continuum remains unproved.  V118 closes
the local curved one-form symbol and proves exact BRST incidence
conditionally on compatible Calderón branches.  Global harmonic transport,
the finite BRST quotient, nonlocal heat/stress coefficients, fermion branch
selection, nonlinear existence, and the cofinal interacting limit remain
open.
\end{assessment}

\section*{V118 exact checks and append-only record}
\addcontentsline{toc}{section}{V118 exact checks and append-only record}

The curvature matrix and its incidence checksum are
\begin{equation}
 \boxed{
 M_K(n)=
 \begin{pmatrix}
 \tfrac12(k-H)I+S&-iSn\\
 i(Sn)^\dagger&\tfrac12(k-H)
 \end{pmatrix},
 \qquad
 2s_0-k=-H.}
 \label{eq:V118-final-check}
\end{equation}
The append operation preserves the complete V117 prefix byte for byte before
its former live terminator.  The assembled V118 file is checked for one live
final terminator, balanced environments and braces, duplicate labels, newly
introduced unresolved references, display math delimiters, and control
characters before the exact V117 predecessor is removed.

% =============================================================================
% END APPENDED VERSION V118 MATERIAL
% =============================================================================

\clearpage
% =============================================================================
% BEGIN APPENDED VERSION V119 MATERIAL
% =============================================================================
\section{V119: harmonic Kato transport, cohomology transport, and the
curved finite-sector Schur map}
\label{sec:V119-harmonic-transport}

V116 separated the zero denominator of the instantaneous Hodge square root,
and V118 derived the full curved one-form collar operator.  These two steps
do not yet imply that the harmonic forms constitute an invariant block of
the curved operator.  V119 resolves that distinction.

The harmonic bundle itself is smooth and admits unitary Kato transport.
Its topological labels carry the flat Gauss--Manin transport.  The two
transports agree only up to the changing Hodge metric and a unitary rotation.
Moreover, the full one-form Schrödinger family
\(\mathscr A_{1,r}\) of \eqref{eq:V118-A1} generally couples harmonic and
nonharmonic boundary forms.  The exact finite object is consequently a
Feshbach--Schur operator, rather than the naive compression
\(\boldsymbol\Pi\mathscr A_1\boldsymbol\Pi\).

\subsection{Smooth harmonic projectors and their gap}
\label{subsec:V119-smooth-projectors}

Let \(I\) be a compact collar interval and let \(\gamma_r\), \(r\in I\), be
a \(C^m\), \(m\ge2\), family of metrics on a fixed closed boundary
\(\Sigma\).  Use a smooth unitary trivialization of the varying
\(L^2\)-spaces, such as the normal connection construction of
\eqref{eq:V118-unitary-normal-connection}, and denote the conjugated Hodge
Laplacian on a fixed Hilbert space by \(\widehat\Delta_{p,r}\).
Set
\begin{equation}
 \Pi_{p,r}:=\mathbf1_{\{0\}}(\widehat\Delta_{p,r}),
 \qquad
 Q_{p,r}:=I-\Pi_{p,r},
 \qquad
 R_{p,r}:=Q_{p,r}
 \bigl(\widehat\Delta_{p,r}|_{\operatorname{Ran}Q_{p,r}}\bigr)^{-1}
 Q_{p,r}.
 \label{eq:V119-PQR}
\end{equation}

\begin{theorem}[Uniform harmonic gap and smooth projector]
\label{thm:V119-smooth-projector}
For each degree \(p\), the rank of \(\Pi_{p,r}\) is
\(b_p(\Sigma)\), independent of \(r\).  There is a number
\(g_{p,I}>0\) such that
\begin{equation}
 \operatorname{spec}(\widehat\Delta_{p,r})
 \cap(0,g_{p,I})=\varnothing
 \qquad (r\in I).
 \label{eq:V119-uniform-gap}
\end{equation}
The maps \(r\mapsto\Pi_{p,r}\) and \(r\mapsto R_{p,r}\) are \(C^{m-1}\)
on the corresponding Sobolev scales.  If \(\Gamma\) is a fixed positively
oriented circle enclosing zero and no positive spectrum, then
\begin{equation}
 \Pi_{p,r}
 =
 \frac{1}{2\pi i}\int_\Gamma
 (z-\widehat\Delta_{p,r})^{-1}\,dz.
 \label{eq:V119-Riesz-projector}
\end{equation}
\end{theorem}

\begin{proof}
Hodge theory identifies
\(\ker\widehat\Delta_{p,r}\) with \(H^p_{\rm dR}(\Sigma)\), so its dimension
is the metric-independent Betti number.  Ellipticity and compactness give
discrete spectrum with finite multiplicities.  The first positive
eigenvalue varies continuously with \(r\) by the min--max principle after
unitary trivialization.  It is positive at every \(r\), and its minimum on
the compact interval \(I\) is therefore the positive number \(g_{p,I}\).

The coefficients form a \(C^{m-1}\) family of elliptic operators with common
Sobolev domain after trivialization.  The resolvent identity makes
\((z-\widehat\Delta_{p,r})^{-1}\) \(C^{m-1}\) uniformly for
\(z\in\Gamma\).  Differentiation under the contour proves the projector
claim.  On \(Q_{p,r}\), the spectral bound
\(\widehat\Delta_{p,r}\ge g_{p,I}\) gives the reduced inverse and its smooth
dependence.
\end{proof}

\begin{lemma}[Exact derivative and reduced-resolvent formulas]
\label{lem:V119-projector-derivative}
With a dot denoting \(\partial_r\),
\begin{align}
 \dot\Pi_p
 &=
 -R_p\dot{\widehat\Delta}_p\Pi_p
 -\Pi_p\dot{\widehat\Delta}_pR_p,
 \label{eq:V119-Pdot}\\
 [\dot\Pi_p,\Pi_p]
 &=
 \Pi_p\dot{\widehat\Delta}_pR_p
 -R_p\dot{\widehat\Delta}_p\Pi_p.
 \label{eq:V119-Kato-generator-resolvent}
\end{align}
In particular \(\Pi_p\dot\Pi_p\Pi_p=0=Q_p\dot\Pi_pQ_p\).
\end{lemma}

\begin{proof}
Differentiating \(\Pi_p^2=\Pi_p\) gives the two zero diagonal blocks.
Differentiate
\(\widehat\Delta_p\Pi_p=0\), multiply on the left by \(R_p\), and use
\(R_p\widehat\Delta_p=Q_p\) to obtain
\[
 Q_p\dot\Pi_p\Pi_p=-R_p\dot{\widehat\Delta}_p\Pi_p.
\]
Taking adjoints gives
\(\Pi_p\dot\Pi_pQ_p=-\Pi_p\dot{\widehat\Delta}_pR_p\).
The sum of the two off-diagonal blocks proves
\eqref{eq:V119-Pdot}; their difference proves
\eqref{eq:V119-Kato-generator-resolvent}.
\end{proof}

\subsection{Unitary Kato transport}
\label{subsec:V119-Kato-transport}

Define the skew-adjoint generator
\begin{equation}
 G_{p,r}:=[\dot\Pi_{p,r},\Pi_{p,r}].
 \label{eq:V119-G}
\end{equation}

\begin{theorem}[Exact unitary transport of harmonic spaces]
\label{thm:V119-Kato-transport}
The initial-value problem
\begin{equation}
 \partial_rU_p(r,r_0)=G_{p,r}U_p(r,r_0),
 \qquad
 U_p(r_0,r_0)=I
 \label{eq:V119-Kato-ODE}
\end{equation}
has a unitary solution and satisfies
\begin{equation}
 \boxed{
 \Pi_{p,r}U_p(r,r_0)
 =
 U_p(r,r_0)\Pi_{p,r_0}.}
 \label{eq:V119-Kato-intertwining}
\end{equation}
If \(s(r)=U_p(r,r_0)s_0\) and
\(s_0\in\operatorname{Ran}\Pi_{p,r_0}\), then
\begin{equation}
 \Pi_{p,r}\partial_rs(r)=0
 \label{eq:V119-Kato-parallel}
\end{equation}
in the fixed unitary trivialization.
\end{theorem}

\begin{proof}
Since \(\Pi_p\) and \(\dot\Pi_p\) are self-adjoint,
\(G_p^\dagger=-G_p\), so the evolution is unitary.  Put
\(F(r)=\Pi_{p,r}U_p(r,r_0)-U_p(r,r_0)\Pi_{p,r_0}\).
Using \(\Pi\dot\Pi\Pi=0\), a direct differentiation gives
\(\dot F=G_pF\) and \(F(r_0)=0\).  Hence \(F=0\), proving
\eqref{eq:V119-Kato-intertwining}.  For \(s=\Pi s\),
\[
 \partial_rs=G_ps=(I-\Pi)\dot\Pi s,
\]
whose harmonic projection is zero.
\end{proof}

\begin{corollary}[No singular harmonic denominator]
\label{cor:V119-no-singular-denominator}
The undetermined harmonic--harmonic entry in
\eqref{eq:V116-spectral-C} is a choice of finite connection, not an
elliptic inverse.  The Kato choice is fixed by
\eqref{eq:V119-Kato-ODE}; the positive-spectrum Sylvester inverse remains
bounded by the gap \eqref{eq:V119-uniform-gap}.
\end{corollary}

\begin{proof}
Lemma~\ref{lem:V116-harmonic-numerator} proves compatibility of the zero
Sylvester row.  Theorem~\ref{thm:V119-Kato-transport} supplies the missing
finite transport.  If at least one square-root eigenvalue is positive, its
denominator is bounded below by \(\sqrt{g_{p,I}}\).
\end{proof}

\subsection{Kato transport is not Gauss--Manin transport}
\label{subsec:V119-two-connections}

Let
\begin{equation}
 j_{p,r}:H^p_{\rm dR}(\Sigma;\mathbb R)
 \longrightarrow\mathcal H^p_r
 \label{eq:V119-Hodge-isomorphism}
\end{equation}
send a cohomology class to its \(\gamma_r\)-harmonic representative.  The
Hodge metric on the fixed cohomology vector space is
\begin{equation}
 h_{p,r}(x,y)
 :=
 \langle j_{p,r}x,j_{p,r}y\rangle_{L^2(\gamma_r)}.
 \label{eq:V119-Hodge-metric}
\end{equation}

\begin{proposition}[Flat topological and metric transports]
\label{prop:V119-GM-versus-Kato}
There are two distinct structures:
\begin{enumerate}[label=\textup{(\alph*)},leftmargin=2.8em]
\item Gauss--Manin transport keeps \(x\in H^p_{\rm dR}(\Sigma)\) constant.
      It is flat and preserves the integral lattice, but it need not preserve
      the Hodge metric \(h_{p,r}\).
\item Kato transport, pulled back by \(j_{p,r}\), preserves the Hodge metric.
      In a Gauss--Manin-flat basis its connection matrix \(A_{p,r}^{K}\)
      satisfies
      \begin{equation}
       \dot h_p=(A_p^K)^\dagger h_p+h_pA_p^K
       \label{eq:V119-metric-compatibility}
      \end{equation}
      and therefore has the exact decomposition
      \begin{equation}
       \boxed{
       A_p^K=\frac12h_p^{-1}\dot h_p+\Omega_p,
       \qquad
       \Omega_p^\dagger h_p+h_p\Omega_p=0.}
       \label{eq:V119-Kato-GM-split}
      \end{equation}
\end{enumerate}
On a collar interval, the \(h_p\)-skew term can be removed by a unitary
choice of harmonic frame.  Along a multiparameter family it can carry Kato
curvature and cannot be discarded without a flatness test.
\end{proposition}

\begin{proof}
The underlying manifold and its de Rham cohomology are fixed, so constant
cohomology classes define the flat Gauss--Manin connection and preserve
\(H^p(\Sigma;\mathbb Z)\).  Theorem~\ref{thm:V119-Kato-transport} is unitary,
so its pullback is an isometry from \(h_{p,r_0}\) to \(h_{p,r}\).
Differentiating that isometry gives
\eqref{eq:V119-metric-compatibility}.  The operator
\(\tfrac12h_p^{-1}\dot h_p\) has the required \(h_p\)-symmetric part; the
remainder is \(h_p\)-skew.  A one-dimensional skew connection is removed by
solving its unitary ordinary differential equation.  In more than one
parameter, the commutator curvature is the obstruction to doing so around
loops.
\end{proof}

\begin{corollary}[The connected scalar harmonic line]
\label{cor:V119-scalar-harmonic-line}
If \(\Sigma\) is connected, then
\(\mathcal H_r^0\) is the constant line.  With
\begin{equation}
 {\rm Vol}_r:=\int_\Sigma d{\rm vol}_{\gamma_r},
 \qquad
 \overline H_r
 :=
 {\rm Vol}_r^{-1}\int_\Sigma H\,d{\rm vol}_{\gamma_r},
 \label{eq:V119-volume-H}
\end{equation}
one has
\begin{equation}
 \partial_r\log{\rm Vol}_r=\overline H_r,
 \qquad
 A_{0,r}^{K}=\frac12\overline H_r
 \label{eq:V119-scalar-Kato-connection}
\end{equation}
in the nonrotating frame.  Thus
\({\rm Vol}_r^{-1/2}\mathbf1\) is the normalized transported generator.
\end{corollary}

\begin{proof}
The first identity follows from
\(\partial_rd{\rm vol}_{\gamma_r}=H\,d{\rm vol}_{\gamma_r}\).
The Hodge metric on the constant cohomology generator is
\(h_{0,r}={\rm Vol}_r\).  Apply
\eqref{eq:V119-Kato-GM-split} with \(\Omega_0=0\).
\end{proof}

\begin{firewall}[The integral lattice uses Gauss--Manin transport]
\label{fire:V119-lattice}
Large-gauge identifications and Abelian toron periods are defined by the
integral cohomology lattice.  That lattice is constant under Gauss--Manin
transport, not under a generic metric Kato rotation.  Orthonormal determinant
frames may use Kato transport, but the finite moduli quotient must be
returned to the Gauss--Manin lattice before its volume is declared.
\end{firewall}

\subsection{The reduced de Rham complex}
\label{subsec:V119-reduced-complex}

Return temporarily to raw differential forms, where \(d_\Sigma\) is
metric-independent.

\begin{theorem}[Exact harmonic and reduced incidence identities]
\label{thm:V119-reduced-incidence}
For every \(p\),
\begin{align}
 d_\Sigma\Pi_p&=0,
 &
 \Pi_{p+1}d_\Sigma&=0,
 \label{eq:V119-harmonic-d-zero}\\
 d_\Sigma
 &=
 Q_{p+1}d_\Sigma Q_p,
 \label{eq:V119-QdQ}\\
 d_\Sigma R_p
 &=
 R_{p+1}d_\Sigma
 \label{eq:V119-Green-intertwining}
\end{align}
on smooth forms, with the last identity interpreted on the natural reduced
domains.  The analogous identity
\(\delta_\Sigma R_{p+1}=R_p\delta_\Sigma\) also holds.
\end{theorem}

\begin{proof}
A harmonic form is closed, proving \(d_\Sigma\Pi_p=0\).  If \(\eta\) is
harmonic of degree \(p+1\), then
\[
 \langle d_\Sigma\alpha,\eta\rangle
 =
 \langle\alpha,\delta_\Sigma\eta\rangle=0,
\]
which proves \(\Pi_{p+1}d_\Sigma=0\) and hence
\eqref{eq:V119-QdQ}.  Since
\(\Delta_{p+1}d_\Sigma=d_\Sigma\Delta_p\), both
\(d_\Sigma R_p\alpha\) and \(R_{p+1}d_\Sigma\alpha\) lie in the orthogonal
complement of harmonic forms and solve
\[
 \Delta_{p+1}u=d_\Sigma Q_p\alpha=d_\Sigma\alpha.
\]
Uniqueness on the reduced space proves
\eqref{eq:V119-Green-intertwining}.  Taking adjoints proves the codifferential
identity.
\end{proof}

\begin{corollary}[Finite harmonic de Rham differential]
\label{cor:V119-finite-differential}
The differential induced between harmonic representatives is zero:
\begin{equation}
 \mathcal H_r^p
 \xrightarrow{\ d_\Sigma\ }
 \mathcal H_r^{p+1},
 \qquad d_\Sigma|_{\mathcal H_r^p}=0.
 \label{eq:V119-finite-zero-complex}
\end{equation}
The cohomology labels themselves are preserved by Gauss--Manin transport,
while the nonzero spectral complex is preserved by
\eqref{eq:V119-QdQ}--\eqref{eq:V119-Green-intertwining}.
\end{corollary}

\subsection{Why the full curved operator does not preserve the harmonic
sum}
\label{subsec:V119-harmonic-leakage}

For the tangential--normal splitting of a one-form, define
\begin{equation}
 \boldsymbol\Pi_r
 :=
 \begin{pmatrix}\Pi_{1,r}&0\\0&\Pi_{0,r}\end{pmatrix},
 \qquad
 \boldsymbol Q_r:=I-\boldsymbol\Pi_r.
 \label{eq:V119-bold-PQ}
\end{equation}

\begin{proposition}[Exact mean-curvature leakage]
\label{prop:V119-mean-curvature-leakage}
Let \(\phi\in\mathcal H_r^0\).  The upper-right entry of
\(\mathscr A_{1,r}\) satisfies
\begin{align}
 \Pi_{1,r}
 \bigl(-2S\,d_\Sigma-d_\Sigma H\bigr)\phi
 &=0,
 \label{eq:V119-harmonic-cross-zero}\\
 Q_{1,r}
 \bigl(-2S\,d_\Sigma-d_\Sigma H\bigr)\phi
 &=-\phi\,d_\Sigma H.
 \label{eq:V119-harmonic-leakage}
\end{align}
Consequently, if \(d_\Sigma H\ne0\), then
\begin{equation}
 \boldsymbol Q_r\mathscr A_{1,r}\boldsymbol\Pi_r\ne0.
 \label{eq:V119-noninvariance}
\end{equation}
Thus the direct sum
\(\mathcal H_r^1\oplus dr\,\mathcal H_r^0\) is not an invariant subspace of
the generic curved one-form collar operator.
\end{proposition}

\begin{proof}
On a connected component, \(\phi\) is constant, so
\(d_\Sigma\phi=0\).  The remaining one-form is
\(-\phi\,d_\Sigma H\).  It is exact and therefore orthogonal to every
harmonic one-form:
\[
 \langle d_\Sigma H,\eta\rangle
 =
 \langle H,\delta_\Sigma\eta\rangle=0.
\]
This proves both projection formulas.  If \(d_\Sigma H\ne0\), the second is
nonzero for a nonzero constant \(\phi\), proving
\eqref{eq:V119-noninvariance}.
\end{proof}

\begin{firewall}[Correction to a naive harmonic-block reading]
\label{fire:V119-naive-harmonic-block}
The harmonic separation in
Theorem~\ref{thm:V116-harmonic-Sylvester} is exact for the instantaneous
operator \(\Lambda=\Delta_\Sigma^{1/2}\).  It does not assert that the full
curved matrix \(\mathscr A_1\) commutes with
\(\boldsymbol\Pi\).  Equation \eqref{eq:V119-harmonic-leakage} is an explicit
counterexample whenever the mean curvature is not tangentially constant.
\end{firewall}

\subsection{The exact finite Feshbach operator}
\label{subsec:V119-Feshbach}

Fix \(r\), put \(\mathbb A=\mathscr A_{1,r}\), and suppress \(r\) from
\(\boldsymbol\Pi,\boldsymbol Q\).  For a spectral parameter \(z\), assume
\begin{equation}
 \mathbb A_{QQ}(z)
 :=
 \boldsymbol Q(\mathbb A-z)\boldsymbol Q
 \label{eq:V119-AQQ}
\end{equation}
is invertible on \(\operatorname{Ran}\boldsymbol Q\).

\begin{theorem}[Finite-sector Feshbach reduction]
\label{thm:V119-Feshbach}
Define the finite-dimensional operator
\begin{equation}
 \boxed{
 \mathfrak F_1(z)
 :=
 \boldsymbol\Pi(\mathbb A-z)\boldsymbol\Pi
 -
 \boldsymbol\Pi\mathbb A\boldsymbol Q\,
 \mathbb A_{QQ}(z)^{-1}\,
 \boldsymbol Q\mathbb A\boldsymbol\Pi.}
 \label{eq:V119-Feshbach}
\end{equation}
Then
\begin{equation}
 \ker(\mathbb A-z)
 \cong
 \ker\mathfrak F_1(z).
 \label{eq:V119-kernel-equivalence}
\end{equation}
More precisely, if \(p\in\ker\mathfrak F_1(z)\), the corresponding full
vector is
\begin{equation}
 p+q,
 \qquad
 q=-\mathbb A_{QQ}(z)^{-1}
 \boldsymbol Q\mathbb A\boldsymbol\Pi p.
 \label{eq:V119-Feshbach-reconstruction}
\end{equation}
For real \(z\) in a self-adjoint resolvent interval,
\(\mathfrak F_1(z)\) is self-adjoint.  If
\(\mathbb A_{QQ}(z)>0\), the Schur correction is nonpositive:
\begin{equation}
 \langle p,
 [\mathfrak F_1(z)-\boldsymbol\Pi(\mathbb A-z)\boldsymbol\Pi]p\rangle
 =
 -\left\|
 \mathbb A_{QQ}(z)^{-1/2}
 \boldsymbol Q\mathbb A\boldsymbol\Pi p
 \right\|^2.
 \label{eq:V119-Schur-negative}
\end{equation}
\end{theorem}

\begin{proof}
Write \((\mathbb A-z)(p+q)=0\) in
\(\boldsymbol\Pi\oplus\boldsymbol Q\) blocks.  The \(Q\)-row uniquely gives
\eqref{eq:V119-Feshbach-reconstruction}.  Substitution into the
\(\Pi\)-row gives \(\mathfrak F_1(z)p=0\).  The converse is the same
calculation in reverse.  Adjointness follows from
\(\mathbb A^\dagger=\mathbb A\) and the self-adjoint inverse of the
\(QQ\)-block.  The last formula is obtained by moving one factor of the
positive inverse to each side.
\end{proof}

\begin{corollary}[Product-collar recovery]
\label{cor:V119-product-recovery}
If \(K=H=0\) and the collar coefficients are \(r\)-independent, then the
leakage \(\boldsymbol Q\mathbb A\boldsymbol\Pi\) vanishes and
\begin{equation}
 \mathfrak F_1(z)
 =
 \boldsymbol\Pi(\mathbb A-z)\boldsymbol\Pi.
 \label{eq:V119-product-Feshbach}
\end{equation}
Thus the finite sector decouples exactly in the V115 product model.
\end{corollary}

\begin{proof}
Under the hypotheses,
\(\mathscr A_1=\operatorname{diag}(\Delta_{\Sigma,1},
\Delta_{\Sigma,0})\), which commutes with
\(\boldsymbol\Pi\).
\end{proof}

\subsection{Finite BRST quotient ledger}
\label{subsec:V119-finite-BRST}

Let
\begin{equation}
 \mathfrak g_{\rm SM}
 =
 \mathfrak{su}(3)\oplus\mathfrak{su}(2)\oplus\mathfrak u(1)
 \label{eq:V119-gSM}
\end{equation}
at the Lie-algebra level.

\begin{proposition}[Constant ghosts and harmonic one-form moduli are
unpaired]
\label{prop:V119-finite-BRST}
On a connected boundary, the intrinsic harmonic part of the linearized
boundary BRST complex is
\begin{equation}
 H^0(\Sigma;\mathbb R)\otimes\mathfrak g_{\rm SM}
 \xrightarrow{\ 0\ }
 H^1(\Sigma;\mathbb R)\otimes\mathfrak g_{\rm SM}.
 \label{eq:V119-finite-BRST}
\end{equation}
Thus constant adjoint ghosts do not cancel harmonic one-form moduli.
If constant boundary gauge transformations belong to the declared gauge
group, their Lie algebra must be quotiented and the scalar ghost determinant
is a primed determinant.  If they are retained as boundary symmetries, they
form a finite edge-symmetry sector instead.  The degree-one harmonic modes
require their own finite moduli integral and global gauge quotient.
\end{proposition}

\begin{proof}
The first statement is Corollary~\ref{cor:V119-finite-differential} tensored
with \(\mathfrak g_{\rm SM}\).  A zero differential has no nonzero image, so
there is no ghost pairing with \(H^1\).  The two alternatives for constants
are the two possible definitions of the boundary gauge group.  In the
quotient alternative the zero eigenvalues are omitted from the Gaussian
determinant and replaced by the finite group-volume quotient.  In the
symmetry alternative they are not gauge redundancies and must be retained.
\end{proof}

\begin{firewall}[Linear torons are not automatically a non-Abelian torus]
\label{fire:V119-nonabelian-torons}
For the Abelian factor, Gauss--Manin transport preserves the integral period
lattice and the harmonic quotient is a genuine torus.  For the non-Abelian
factors, \(H^1(\Sigma)\otimes\mathfrak g\) is only the tangent space at a
flat background.  The nonlinear flat-connection moduli space can be
singular and includes the residual adjoint quotient.  V119 makes no global
smooth-torus claim for that sector.
\end{firewall}

\begin{assessment}[V119 finite-sector decision]
\label{ass:V119-finite-sector}
The topological harmonic complex and its transport are now exact, but the
full curved one-form operator generically leaks between harmonic and
nonharmonic boundary modes.  Therefore a determinant or heat trace may not
be split by merely deleting the tangential Hodge kernels.  One must first
use the exact Feshbach operator \(\mathfrak F_1(z)\), then apply the declared
global-gauge quotient and the Gauss--Manin lattice.
\end{assessment}

\paragraph{New exact conclusions.}
\begin{enumerate}[label=\textup{(\arabic*)},leftmargin=2.8em]
\item Theorem~\ref{thm:V119-smooth-projector} proves a uniform harmonic gap
      and smooth spectral projectors on a compact collar.
\item Lemma~\ref{lem:V119-projector-derivative} and
      Theorem~\ref{thm:V119-Kato-transport} construct the exact unitary Kato
      transport.
\item Proposition~\ref{prop:V119-GM-versus-Kato} separates metric Kato
      transport from lattice-preserving Gauss--Manin transport.
\item Theorem~\ref{thm:V119-reduced-incidence} proves exact de Rham
      intertwining on the reduced Green spaces.
\item Proposition~\ref{prop:V119-mean-curvature-leakage} proves that
      \(d_\Sigma H\) obstructs naive invariance of the harmonic one-form sum.
\item Theorem~\ref{thm:V119-Feshbach} replaces that false block separation
      by the exact finite-dimensional Schur operator.
\item Proposition~\ref{prop:V119-finite-BRST} fixes the linear finite BRST
      quotient alternatives without identifying harmonic one-forms with
      ghost images.
\end{enumerate}

\paragraph{Next forced step.}
V120 must begin with the scheduled NS/YM companion audit.  It should then
derive a BRST-compatible Feshbach reduction of the paired scalar and
one-form Calderón graphs: the reduced complement, finite Schur operator,
constant-gauge quotient, and harmonic one-form measure must be carried in
one determinant identity.  The first heat/stress variation is admissible
only after that identity is proved.

\begin{assessment}[V119 continuum status]
\label{ass:V119-continuum-status}
The full Standard--Model--Einstein continuum remains unproved.  V119
constructs harmonic transport and discovers the precise curved leakage that
prevents a naive zero-mode deletion.  The BRST-compatible Feshbach
determinant, nonlocal heat/stress coefficients, fermion branch selection,
nonlinear existence, and the cofinal interacting limit remain open.
\end{assessment}

\section*{V119 exact checks and append-only record}
\addcontentsline{toc}{section}{V119 exact checks and append-only record}

The harmonic transport and leakage checksum is
\begin{equation}
 \boxed{
 \dot U_p=[\dot\Pi_p,\Pi_p]U_p,
 \qquad
 \Pi_{1}
 \bigl[-(d_\Sigma H)\Pi_0\bigr]=0,
 \qquad
 Q_1\bigl[-(d_\Sigma H)\Pi_0\bigr]
 =-(d_\Sigma H)\Pi_0.}
 \label{eq:V119-final-check}
\end{equation}
The append operation preserves the complete V118 prefix byte for byte before
its former live terminator.  The assembled V119 file is checked for one live
final terminator, balanced environments and braces, duplicate labels, newly
introduced unresolved references, display math delimiters, and control
characters before the exact V118 predecessor is removed.

% =============================================================================
% END APPENDED VERSION V119 MATERIAL
% =============================================================================

\clearpage
% =============================================================================
% BEGIN APPENDED VERSION V120 MATERIAL
% =============================================================================
\section{V120: companion audit and the BRST-compatible Feshbach determinant}
\label{sec:V120-Feshbach-determinant}

V119 proved that the curved scalar and one-form collar families must be
reduced by a genuine Feshbach map before harmonic modes are quotiented.
V120 performs the scheduled companion audit and then gives the exact
finite-regulator algebra needed for that reduction.  The result separates
the complement determinant, the finite Schur determinant, the Hodge-volume
density, and any regulator defect in the de Rham chain.

\subsection{Scheduled NS/YM companion audit}
\label{subsec:V120-companion-audit}

The read-only audit was made on 5 September 2026.  The shared folder
contained the following companion snapshots:
\begin{align}
 &\texttt{Renewal\_NS\_Stability\_Research\_Notes\_V31.tex},
 \notag\\[-1mm]
 &\hspace{8mm}
 887537\ {\rm bytes},\qquad
 \operatorname{SHA256}
 =
 \texttt{F1121B62238AD5491BB7CF03635EA9934942EDF573ED8FAAA9EE79E1D38A6696},
 \label{eq:V120-NS-audit}\\
 &\texttt{Renewal\_Yang\_Mills\_Mass\_Gap\_Research\_Notes\_V106.tex},
 \notag\\[-1mm]
 &\hspace{8mm}
 2552798\ {\rm bytes},\qquad
 \operatorname{SHA256}
 =
 \texttt{0F510459BDF601424A5C7DE0C1C85B6658A79CDA22D172ECD97D97300AB1FD44}.
 \label{eq:V120-YM-audit}
\end{align}
The YM filesystem name ended in V106, while its sole live document tail
already contained a complete internally labelled V107 appendix.  V120
records both facts and does not rename or modify the companion file.

The NS note remains at its V31 stopping point.  Its relevant lesson is that
a restricted formation quantity must retain the correlation among active
times, signed work, and parent tails; replacing it too early by a global
positive bound returns the unresolved critical stock.  The newest YM
appendix proves an analogous finite-dimensional fact for recoupled
isotypic channels: the physical sector weight must remain attached to its
marginal, because conditional normalization followed by omission of the
weight creates a fictitious dimension loss.

\begin{programme}[V120 adjustment after the companion audit]
\label{prog:V120-audit-adjustment}
The SM finite sector will therefore be handled in the following order.
\begin{enumerate}[label=\textup{(\roman*)},leftmargin=3.5em]
\item retain the complete harmonic--complement coupling until the Schur map
      is formed;
\item use one regulator which preserves the scalar-to-one-form de Rham
      incidence, or print its exact chain defect;
\item take determinants only after the complement and finite factors have
      been separated by unit-determinant elimination;
\item carry the Hodge metric density with every finite factor, without
      independently normalizing harmonic channels and dropping the
      Jacobian.
\end{enumerate}
\end{programme}

\subsection{Exact two-sided Feshbach elimination}
\label{subsec:V120-two-sided-elimination}

The following algebra is stated for finite-dimensional regulated spaces.
It therefore has no zeta-determinant anomaly and no domain ambiguity.  Let
\(A_p\), \(p=0,1\), be regulated scalar and one-form operators, let
\(P_p\) be their declared finite-sector projections, and put \(Q_p=I-P_p\).
For the one-form operator one may take
\(P_1=\boldsymbol\Pi\) from \eqref{eq:V119-bold-PQ}; for the scalar operator
one takes \(P_0=\Pi_0\).  Fix \(z\) and abbreviate
\begin{align}
 A_{p,PP}(z)&:=P_p(A_p-z)P_p,
 &
 A_{p,PQ}&:=P_pA_pQ_p,
 \notag\\
 A_{p,QP}&:=Q_pA_pP_p,
 &
 B_p(z)&:=Q_p(A_p-z)Q_p.
 \label{eq:V120-blocks}
\end{align}
Assume \(B_p(z)\) is invertible.

\begin{definition}[Right lift, left residual map, and Schur operator]
\label{def:V120-eliminators}
In the decomposition \(P_p\oplus Q_p\), define
\begin{align}
 \mathcal R_p(z)
 &:=
 \begin{pmatrix}
 I&0\\
 -B_p(z)^{-1}A_{p,QP}&I
 \end{pmatrix},
 \label{eq:V120-R}\\
 \mathcal L_p(z)
 &:=
 \begin{pmatrix}
 I&-A_{p,PQ}B_p(z)^{-1}\\
 0&I
 \end{pmatrix},
 \label{eq:V120-L}\\
 F_p(z)
 &:=
 A_{p,PP}(z)
 -A_{p,PQ}B_p(z)^{-1}A_{p,QP}.
 \label{eq:V120-F}
\end{align}
\end{definition}

\begin{theorem}[Exact diagonal equivalence and determinant factorization]
\label{thm:V120-diagonal-equivalence}
The triangular eliminators have determinant one and satisfy
\begin{equation}
 \boxed{
 \mathcal L_p(z)(A_p-z)\mathcal R_p(z)
 =
 \begin{pmatrix}
 F_p(z)&0\\0&B_p(z)
 \end{pmatrix}.}
 \label{eq:V120-diagonal-equivalence}
\end{equation}
Consequently,
\begin{equation}
 \boxed{
 \det(A_p-z)=\det B_p(z)\det F_p(z).}
 \label{eq:V120-determinant-factorization}
\end{equation}
\end{theorem}

\begin{proof}
Right multiplication first gives
\[
 (A_p-z)\mathcal R_p
 =
 \begin{pmatrix}
 F_p&A_{p,PQ}\\
 0&B_p
 \end{pmatrix}.
\]
Left multiplication by \(\mathcal L_p\) kills the remaining upper-right
entry and proves \eqref{eq:V120-diagonal-equivalence}.  Both eliminators are
triangular with identity diagonal, so their determinants are one.  Taking
determinants proves \eqref{eq:V120-determinant-factorization}.
\end{proof}

\begin{corollary}[Kernel reconstruction]
\label{cor:V120-kernel-reconstruction}
The map
\begin{equation}
 J_p(z):P_p\mathcal H_p\longrightarrow\mathcal H_p,
 \qquad
 J_p(z)x
 :=
 \mathcal R_p(z)\binom{x}{0}
 =
 x-B_p(z)^{-1}A_{p,QP}x
 \label{eq:V120-J}
\end{equation}
is an isomorphism
\begin{equation}
 \ker F_p(z)\xrightarrow{\ \cong\ }\ker(A_p-z).
 \label{eq:V120-kernel-isomorphism}
\end{equation}
Thus true gauge or physical zero modes are kernels of the full Schur
operator, not kernels of the tangential Hodge Laplacian alone.
\end{corollary}

\begin{proof}
Apply \eqref{eq:V120-diagonal-equivalence} to
\(\binom{x}{0}\).  Since \(\mathcal L_p\) and \(\mathcal R_p\) are
invertible, the resulting vector is in the full kernel exactly when
\(F_p(z)x=0\).
\end{proof}

\begin{firewall}[No naive zero-mode deletion]
\label{fire:V120-no-naive-deletion}
A vector in \(\ker\Delta_{\Sigma,p}\) may acquire a nonharmonic tail through
\eqref{eq:V119-harmonic-leakage}.  Deleting it before forming \(F_p\) loses
both the reconstructed zero mode \eqref{eq:V120-J} and its Schur
self-energy.  Priming a determinant is legitimate only on
\(\ker F_p(0)\), after reconstruction has identified which vectors are
actual zero modes of the full regulated problem.
\end{firewall}

\subsection{Exact descent of a de Rham chain map}
\label{subsec:V120-chain-descent}

Let \(T:\mathcal H_0\to\mathcal H_1\) be the regulated de Rham incidence
operator.  In a bulk Hodge regulator this is \(d\); for boundary Cauchy data
it is the trace-gradient column represented in
\eqref{eq:V118-exact-incidence}.  Allow an explicit regulator defect
\begin{equation}
 \mathcal D
 :=
 (A_1-z)T-T(A_0-z).
 \label{eq:V120-chain-defect}
\end{equation}
Define two transformed incidence maps,
\begin{equation}
 T_{\rm sol}(z)
 :=
 \mathcal R_1(z)^{-1}T\mathcal R_0(z),
 \qquad
 T_{\rm res}(z)
 :=
 \mathcal L_1(z)T\mathcal L_0(z)^{-1}.
 \label{eq:V120-two-transformed-T}
\end{equation}
The first acts on Schur-lifted solution coordinates; the second acts on
their residual equations.

\begin{theorem}[BRST-compatible Feshbach identity]
\label{thm:V120-BRST-Feshbach}
Let
\[
 T_{\rm sol}^{PP}:=P_1T_{\rm sol}P_0,
 \qquad
 T_{\rm res}^{PP}:=P_1T_{\rm res}P_0.
\]
Then
\begin{equation}
 \boxed{
 F_1(z)T_{\rm sol}^{PP}(z)
 -
 T_{\rm res}^{PP}(z)F_0(z)
 =
 P_1\mathcal L_1(z)\mathcal D\,
 \mathcal R_0(z)P_0.}
 \label{eq:V120-BRST-Feshbach}
\end{equation}
In particular, if the regulator preserves the exact chain
\(\mathcal D=0\), then the finite Schur equations form an exact chain:
\begin{equation}
 F_1T_{\rm sol}^{PP}=T_{\rm res}^{PP}F_0.
 \label{eq:V120-exact-Schur-chain}
\end{equation}
Therefore \(T_{\rm sol}^{PP}\) maps
\(\ker F_0(z)\) into \(\ker F_1(z)\).
\end{theorem}

\begin{proof}
By Theorem~\ref{thm:V120-diagonal-equivalence},
\[
 \begin{pmatrix}F_1&0\\0&B_1\end{pmatrix}T_{\rm sol}
 -
 T_{\rm res}
 \begin{pmatrix}F_0&0\\0&B_0\end{pmatrix}
 =
 \mathcal L_1\mathcal D\mathcal R_0.
\]
Take the \(P_1P_0\) block.  This is exactly
\eqref{eq:V120-BRST-Feshbach}.  If \(\mathcal D=0\) and \(F_0x=0\), the
same equation gives \(F_1T_{\rm sol}^{PP}x=0\).
\end{proof}

\begin{corollary}[Compatibility with the exact curved Calderón identity]
\label{cor:V120-Calderon-descent}
For exact scalar and one-form Calderón branches satisfying
\eqref{eq:V118-exact-incidence}, the right side of
\eqref{eq:V120-BRST-Feshbach} vanishes before regulation.  A finite
regulator preserves exact on-shell BRST closure if
\begin{equation}
 (A_{1,N}-z)T_N=T_N(A_{0,N}-z).
 \label{eq:V120-regulator-chain-gate}
\end{equation}
If this gate fails, the right side of
\eqref{eq:V120-BRST-Feshbach} is the exact finite BRST anomaly; it may not
be hidden inside a norm estimate.
\end{corollary}

\begin{proof}
The unregulated identity is Proposition~\ref{prop:V118-exact-incidence}.
Apply Theorem~\ref{thm:V120-BRST-Feshbach} to the compressed operators.
\end{proof}

\begin{proposition}[A common Hodge spectral cutoff preserves the bulk
linear complex]
\label{prop:V120-common-Hodge-cutoff}
Let \(E_{p,N}\) be the sum of complete eigenspaces of
\(\Delta_{\Sigma,p}\) with eigenvalue at most \(\Lambda_N\), using the same
threshold in adjacent degrees.  Then
\begin{equation}
 E_{p+1,N}d_\Sigma=d_\Sigma E_{p,N}.
 \label{eq:V120-Hodge-cutoff-chain}
\end{equation}
Thus the intrinsic bulk Hodge complex has zero defect.  A cutoff applied
separately after adding curved lower-order collar terms need not have this
property and must be tested by
\eqref{eq:V120-regulator-chain-gate}.
\end{proposition}

\begin{proof}
If \(\Delta_p\alpha=\lambda\alpha\), then
\(\Delta_{p+1}d_\Sigma\alpha=d_\Sigma\Delta_p\alpha
=\lambda d_\Sigma\alpha\).  Hence \(d_\Sigma\) maps a complete eigenspace
at \(\lambda\) to the adjacent eigenspace at the same \(\lambda\), or to
zero.  Summing all eigenspaces below the common threshold proves the
identity.
\end{proof}

\subsection{The exact regulated determinant card}
\label{subsec:V120-determinant-card}

For positive regulated operators, define the scalar/one-form Gaussian
determinant factor
\begin{equation}
 \mathcal Z_N(z)
 :=
 \frac{\det(A_{0,N}-z)}
 {\det(A_{1,N}-z)^{1/2}}.
 \label{eq:V120-ZN}
\end{equation}
This is the gauge/ghost determinant card only; matter, fermion, Einstein,
interaction, and counterterm factors are not included.

\begin{theorem}[Complement times finite-sector factor]
\label{thm:V120-determinant-card}
Whenever the two \(QQ\) blocks are invertible,
\begin{equation}
 \boxed{
 \mathcal Z_N(z)
 =
 \frac{\det B_{0,N}(z)}
 {\det B_{1,N}(z)^{1/2}}
 \,
 \frac{\det F_{0,N}(z)}
 {\det F_{1,N}(z)^{1/2}}.}
 \label{eq:V120-determinant-card}
\end{equation}
No harmonic-channel normalization appears between the two factors.
\end{theorem}

\begin{proof}
Apply \eqref{eq:V120-determinant-factorization} for \(p=0,1\) and substitute
the two identities into \eqref{eq:V120-ZN}.
\end{proof}

\begin{firewall}[Finite identity before zeta regularization]
\label{fire:V120-zeta-order}
Equation \eqref{eq:V120-determinant-card} is an ordinary
finite-dimensional determinant identity.  Passing separately to zeta
determinants of \(B_p\) and \(F_p\) can introduce regularization choices
and local multiplicative anomalies.  The continuum programme must first
renormalize the common finite identity and prove that any splitting anomaly
is a declared local counterterm.  V120 makes no such limit claim.
\end{firewall}

\subsection{Hodge density and the normalization firewall}
\label{subsec:V120-Hodge-density}

Let \(V\) be one of the finite Schur spaces, let \(h\) be its Hodge metric,
and let \(\mathcal F\) be a positive quadratic form.  In a
Gauss--Manin-flat coordinate vector \(c\), write the matrices as \(h\) and
\(F\).

\begin{lemma}[Basis-invariant finite Gaussian factor]
\label{lem:V120-Hodge-Gaussian}
With the Hodge density
\begin{equation}
 d\mu_h(c):=(\det h)^{1/2}\,dc,
 \label{eq:V120-Hodge-density}
\end{equation}
one has
\begin{equation}
 \boxed{
 \int_V
 e^{-\frac12c^\dagger Fc}\,d\mu_h(c)
 =
 (2\pi)^{\dim V/2}
 \det(h^{-1}F)^{-1/2}.}
 \label{eq:V120-Hodge-Gaussian}
\end{equation}
The right side is invariant under every change of finite basis.
\end{lemma}

\begin{proof}
The ordinary Gaussian integral is
\((2\pi)^{n/2}(\det F)^{-1/2}\).  Multiplication by
\((\det h)^{1/2}\) gives
\[
 (2\pi)^{n/2}
 \left(\frac{\det F}{\det h}\right)^{-1/2}
 =
 (2\pi)^{n/2}\det(h^{-1}F)^{-1/2}.
\]
Under a basis change, \(h\) and \(F\) acquire the same two congruence
factors, which cancel in \(\det(h^{-1}F)\).
\end{proof}

\begin{corollary}[Kato frames and Gauss--Manin lattices give the same
finite density]
\label{cor:V120-two-frame-density}
One may evaluate \eqref{eq:V120-Hodge-Gaussian} in a Kato-orthonormal frame,
where \(h=I\), or in a Gauss--Manin frame, where the integral lattice is
fixed and the factor \((\det h)^{1/2}\) is explicit.  The values agree.
Rescaling each harmonic channel to unit norm and omitting the corresponding
Jacobian changes the answer.
\end{corollary}

\begin{proof}
This is the basis invariance in
Lemma~\ref{lem:V120-Hodge-Gaussian}.  Omitting the Jacobian deletes the
factor required for that invariance.
\end{proof}

\begin{firewall}[Disposition of constant ghosts]
\label{fire:V120-constant-ghosts}
For a connected boundary, the constant adjoint line is quotiented only if
its reconstructed vector \eqref{eq:V120-J} lies in the exact gauge kernel
\(\ker F_0(0)\) and the declared boundary gauge group contains that
constant transformation.  In that case use \(\det'F_0(0)\) and the finite
group-volume quotient.  If the constant transformation is an edge
symmetry, retain its finite measure.  Tangential harmonic one-forms are
not removed by the scalar ghost prime; their Schur form and global moduli
quotient are separate.
\end{firewall}

\subsection{Exact first variation in a Kato-fixed block frame}
\label{subsec:V120-first-variation}

Transport \(P_p(r)\) by Theorem~\ref{thm:V119-Kato-transport}, so the
projector is constant in the transported frame.  If \(G_p\) is the Kato
generator, the derivative of an operator in that frame is
\begin{equation}
 \nabla_r^KA_p
 :=
 \dot A_p+[A_p,G_p].
 \label{eq:V120-Kato-operator-derivative}
\end{equation}
Dots below mean this transported derivative, and \(z\) is held fixed.

\begin{proposition}[Derivative of the full Schur operator]
\label{prop:V120-Fdot}
The finite Schur derivative is
\begin{align}
 \dot F_p
 &=
 \dot A_{p,PP}
 -\dot A_{p,PQ}B_p^{-1}A_{p,QP}
 -A_{p,PQ}B_p^{-1}\dot A_{p,QP}
 \notag\\
 &\qquad
 +A_{p,PQ}B_p^{-1}\dot B_pB_p^{-1}A_{p,QP}.
 \label{eq:V120-Fdot}
\end{align}
Thus the first variation retains the derivative of the leakage, its
adjoint, and the complement resolvent; none may be replaced by a derivative
of the naive compression \(A_{PP}\).
\end{proposition}

\begin{proof}
Differentiate
\(F_p=A_{p,PP}-A_{p,PQ}B_p^{-1}A_{p,QP}\) and use
\(\partial_rB_p^{-1}=-B_p^{-1}\dot B_pB_p^{-1}\).
\end{proof}

\begin{theorem}[Exact regulated log-determinant variation]
\label{thm:V120-logdet-variation}
Away from finite zero modes,
\begin{align}
 \partial_r\log\mathcal Z_N
 &=
 \operatorname{Tr}\!\left(B_{0,N}^{-1}\dot B_{0,N}\right)
 -\frac12
 \operatorname{Tr}\!\left(B_{1,N}^{-1}\dot B_{1,N}\right)
 \notag\\
 &\quad
 +\operatorname{tr}\!\left(F_{0,N}^{-1}\dot F_{0,N}\right)
 -\frac12
 \operatorname{tr}\!\left(F_{1,N}^{-1}\dot F_{1,N}\right).
 \label{eq:V120-logdet-variation}
\end{align}
Here \(\operatorname{Tr}\) denotes the regulated complement trace and
\(\operatorname{tr}\) the finite Schur trace.  If exact gauge zero modes
are present, the same formula holds on the primed complement together with
the separately differentiated Hodge/group-volume factor.
\end{theorem}

\begin{proof}
Differentiate the logarithm of
\eqref{eq:V120-determinant-card} and use
\(\partial_r\log\det X=\operatorname{tr}(X^{-1}\dot X)\).
The Kato frame keeps the block projections fixed, so no unprinted projector
derivative occurs.
\end{proof}

\begin{assessment}[First admissible heat/stress input]
\label{ass:V120-first-variation}
Equation \eqref{eq:V120-logdet-variation}, with
\eqref{eq:V120-Fdot}, is the first finite expression from which a boundary
stress variation may be renormalized without losing the curved harmonic
leakage.  The complement traces contain the local high-frequency
divergences; the finite traces contain the global harmonic response.  A
continuum stress tensor requires a common subtraction scheme and the
metric variation of the Gauss--Manin/Hodge density.
\end{assessment}

\subsection{V120 ledger and next acquisition}
\label{subsec:V120-ledger}

\paragraph{New exact conclusions.}
\begin{enumerate}[label=\textup{(\arabic*)},leftmargin=2.8em]
\item The scheduled audit records the unchanged NS V31 snapshot and the YM
      file whose live content has advanced internally through V107.
\item Theorem~\ref{thm:V120-diagonal-equivalence} proves the exact
      unit-determinant two-sided Feshbach elimination.
\item Theorem~\ref{thm:V120-BRST-Feshbach} descends the de Rham chain and
      prints the exact finite regulator anomaly.
\item Proposition~\ref{prop:V120-common-Hodge-cutoff} gives one intrinsic
      chain-preserving cutoff and isolates the extra curved cutoff gate.
\item Theorem~\ref{thm:V120-determinant-card} separates the regulated
      gauge/ghost determinant into complement and finite Schur factors.
\item Lemma~\ref{lem:V120-Hodge-Gaussian} restores the Hodge density and
      forbids channel normalization without its Jacobian.
\item Proposition~\ref{prop:V120-Fdot} and
      Theorem~\ref{thm:V120-logdet-variation} give the exact finite metric
      variation with all leakage terms retained.
\end{enumerate}

\paragraph{Next forced step.}
V121 should build the first common heat-kernel subtraction for the
complement traces in \eqref{eq:V120-logdet-variation}.  It must compute
which boundary Seeley coefficients cancel in the BRST scalar/one-form
combination, keep the finite Schur/Hodge-density derivative separate, and
print every surviving local counterterm.  If a chain-preserving curved
regulator cannot be constructed, the work must move upstream to bound the
right side of \eqref{eq:V120-BRST-Feshbach} in the same subtraction norm.

\begin{assessment}[V120 continuum status]
\label{ass:V120-continuum-status}
The full Standard--Model--Einstein continuum remains unproved.  V120
constructs the exact finite determinant and its first variation while
preserving BRST incidence, harmonic leakage, and finite measure weights.
Common heat renormalization, the fermion branch, interacting nonlinear
existence, Einstein stress closure, and the cofinal limit remain open.
\end{assessment}

\section*{V120 exact checks and append-only record}
\addcontentsline{toc}{section}{V120 exact checks and append-only record}

The determinant and chain checksums are
\begin{equation}
 \boxed{
 \det(A_p-z)=\det B_p(z)\det F_p(z),
 \qquad
 F_1T_{\rm sol}^{PP}-T_{\rm res}^{PP}F_0
 =
 P_1\mathcal L_1\mathcal D\mathcal R_0P_0.}
 \label{eq:V120-final-check}
\end{equation}
The append operation preserves the complete V119 prefix byte for byte before
its former live terminator.  The assembled V120 file is checked for one live
final terminator, balanced environments and braces, duplicate labels, newly
introduced unresolved references, display math delimiters, and control
characters before the exact V119 predecessor is removed.

% =============================================================================
% END APPENDED VERSION V120 MATERIAL
% =============================================================================

\clearpage
% =============================================================================
% BEGIN APPENDED VERSION V121 MATERIAL
% =============================================================================
\section{V121: the fixed-graph heat-domain obstruction and the first
nonlocal heat coefficients}
\label{sec:V121-heat-obstruction}

V120 produced a correct abstract Feshbach determinant identity, conditional
on a scalar-to-one-form chain for the regulated operators.  Before assigning
heat coefficients to that identity, one must check the domains of the
unbounded realizations.  V121 performs that check and finds an obstruction
already in the product collar.

The boundary row
\(B=-\partial_r+\sqrt{\Delta_\Sigma}\) is a genuine parameter-elliptic,
self-adjoint fixed boundary condition.  It is not the interior decaying
Calderón graph.  Nevertheless, the V115 BRST argument proves incidence only
for the zero-energy ghost equation.  At a nonzero spectral value the normal
component acquires an exact boundary defect.  Therefore the fixed heat
semigroups do not form the de Rham complex required for a BRST heat
cancellation.

This note also computes the leading boundary heat coefficient of the fixed
nonlocal graph.  That coefficient is useful as a diagnostic and as a
counterterm for the fixed realization, but it is not labelled a
BRST-renormalized Standard Model coefficient.

\subsection{The fixed graph is elliptic and non-tautological}
\label{subsec:V121-fixed-graph}

Retain the product operator and boundary row
\begin{equation}
 P_p=-\partial_r^2+\Delta_{\Sigma,p},
 \qquad
 B_p=-\partial_r+\Lambda_p,
 \qquad
 \Lambda_p=\Delta_{\Sigma,p}^{1/2}.
 \label{eq:V121-fixed-graph}
\end{equation}

\begin{proposition}[Complementarity to the stable Calderón trace]
\label{prop:V121-complementarity}
For the parameter equation \((P_p+z)u=0\), a frozen stable mode has the form
\begin{equation}
 u(r,y)=e^{-\lambda r}e^{iy\cdot\xi}q,
 \qquad
 \lambda=\sqrt{|\xi|^2+z},
 \qquad
 \operatorname{Re}\lambda\ge0.
 \label{eq:V121-stable-mode}
\end{equation}
On this mode,
\begin{equation}
 B_pu|_{r=0}=(\lambda+|\xi|)q.
 \label{eq:V121-stable-B}
\end{equation}
Thus the fixed graph is complementary to the inward-decaying Calderón
range and has the V115 parameter floor.  In particular, it is not the
tautological boundary condition satisfied by every interior harmonic
extension.
\end{proposition}

\begin{proof}
The inward derivative of \eqref{eq:V121-stable-mode} is
\(-\lambda u\), while the principal symbol of \(\Lambda_p\) is
\(|\xi|\).  Substitution gives \eqref{eq:V121-stable-B}.  Its lower bound is
Theorem~\ref{thm:V115-parameter-floor}.
\end{proof}

\subsection{Exact spectral BRST defect}
\label{subsec:V121-spectral-defect}

Let a scalar ghost be separated as
\(c(r,y)=f(r)\psi(y)\), where
\begin{equation}
 \Delta_{\Sigma,0}\psi=\mu^2\psi,
 \qquad
 \mu\ge0.
 \label{eq:V121-boundary-mode}
\end{equation}
Write its gauge variation as
\begin{equation}
 dc=d_\Sigma c+(\partial_rc)\,dr.
 \label{eq:V121-dc-split}
\end{equation}

\begin{theorem}[The fixed graph is not an off-shell de Rham domain]
\label{thm:V121-spectral-domain-defect}
Suppose
\begin{equation}
 P_0c=\zeta c,
 \qquad
 B_0c|_\Sigma=0.
 \label{eq:V121-scalar-eigenmode}
\end{equation}
Then the tangential component of \(dc\) obeys the one-form boundary row,
but the normal component has the exact defect
\begin{align}
 B_1(d_\Sigma c)|_\Sigma&=0,
 \label{eq:V121-tangential-domain}\\
 B_0(\partial_rc)|_\Sigma&=\zeta\,\gamma c.
 \label{eq:V121-normal-domain-defect}
\end{align}
Consequently,
\begin{equation}
 d\,\operatorname{Dom}(P_{0,B_0})
 \not\subset
 \operatorname{Dom}(P_{1,B_1\oplus B_0})
 \label{eq:V121-domain-failure}
\end{equation}
whenever a nonzero spectral mode has nonzero boundary trace.
\end{theorem}

\begin{proof}
The tangential identity follows from
\(d_\Sigma\Lambda_0=\Lambda_1d_\Sigma\):
\[
 B_1d_\Sigma c=d_\Sigma B_0c=0.
\]
For the normal component, \(\partial_r\) commutes with \(\Lambda_0\), so
\[
 B_0\partial_rc=\partial_r(B_0c).
\]
The exact factorization
\[
 P_0=(\partial_r+\Lambda_0)B_0
\]
and the boundary equation \(B_0c=0\) give
\[
 \gamma(P_0c)=\partial_r(B_0c)|_\Sigma.
\]
Using \(P_0c=\zeta c\) proves
\eqref{eq:V121-normal-domain-defect}.  If
\(\zeta\gamma c\ne0\), the normal component violates its fixed boundary
row, proving \eqref{eq:V121-domain-failure}.
\end{proof}

\begin{corollary}[Correction to the heat interpretation of V120]
\label{cor:V121-V120-correction}
The abstract identities of V120 remain valid.  Their application to a
BRST heat determinant for the fixed graph
\eqref{eq:V121-fixed-graph} is not established: the chain defect
\(\mathcal D\) in \eqref{eq:V120-chain-defect} is nonzero on nonzero
spectral modes and its normal boundary trace is
\(\zeta\,\gamma c\).
\end{corollary}

\begin{proof}
Theorem~\ref{thm:V121-spectral-domain-defect} shows that the differential
expressions commute but the operator domains do not.  Intertwining of
unbounded self-adjoint operators requires both the differential identity and
domain incidence.  The latter fails, so the hypothesis
\eqref{eq:V120-regulator-chain-gate} cannot be inferred for the fixed heat
realizations.
\end{proof}

\begin{firewall}[On-shell closure is retained]
\label{fire:V121-onshell-retained}
At \(\zeta=0\), equation
\eqref{eq:V121-normal-domain-defect} vanishes and exactly recovers the V115
on-shell BRST theorem.  V121 withdraws no zero-energy incidence statement.
It only forbids its use as an off-shell domain identity for a determinant
or heat semigroup.
\end{firewall}

\subsection{Why the parameter-dependent root does not define the same heat
operator}
\label{subsec:V121-parameter-pencil}

For \(z\) in the resolvent sector define
\begin{equation}
 \Lambda_p(z):=(\Delta_{\Sigma,p}+z)^{1/2},
 \qquad
 B_p(z):=-\partial_r+\Lambda_p(z).
 \label{eq:V121-parameter-root}
\end{equation}

\begin{proposition}[Exact incidence for the boundary pencil]
\label{prop:V121-pencil-incidence}
The product parameter equation factorizes as
\begin{equation}
 P_p+z=(\partial_r+\Lambda_p(z))B_p(z).
 \label{eq:V121-parameter-factorization}
\end{equation}
Thus the V115--V118 on-shell incidence argument holds separately for every
fixed \(z\) when all boundary rows use the same parameter-dependent branch.
\end{proposition}

\begin{proof}
The operator \(\Lambda_p(z)\) is independent of \(r\) and
\(\Lambda_p(z)^2=\Delta_{\Sigma,p}+z\).  Expanding the two factors proves
\eqref{eq:V121-parameter-factorization}.  Hodge functional calculus gives
\(d_\Sigma\Lambda_0(z)=\Lambda_1(z)d_\Sigma\), so the tangential incidence
proof is unchanged; the factorization supplies the normal row.
\end{proof}

\begin{proposition}[The pencil is not the resolvent of the V115 fixed
realization]
\label{prop:V121-pencil-not-fixed-resolvent}
The domains
\begin{equation}
 \mathcal D_p(z):=
 \{u\in H^2:B_p(z)\gamma u=0\}
 \label{eq:V121-moving-domain}
\end{equation}
depend nontrivially on \(z\).  Hence the family obtained from
\((P_p+z,\mathcal D_p(z))\) is not the resolvent family of a single
operator on the original bulk Hilbert space.
\end{proposition}

\begin{proof}
On a boundary eigenvector with \(\Delta_{\Sigma,p}\psi=\mu^2\psi\),
\[
 [\Lambda_p(z_1)-\Lambda_p(z_2)]\psi
 =
 \bigl(\sqrt{\mu^2+z_1}-\sqrt{\mu^2+z_2}\bigr)\psi.
\]
This is nonzero for \(z_1\ne z_2\).  Therefore the normal trace required by
\(B_p(z)\gamma u=0\) changes with \(z\).  A fixed closed operator has one
parameter-independent domain, and its resolvent satisfies the resolvent
identity on that domain.  The moving-domain pencil cannot be identified
with that fixed resolvent without enlarging the state space.
\end{proof}

\begin{firewall}[No heat contour from a moving domain]
\label{fire:V121-no-moving-domain-heat}
The parameter-dependent pencil is valuable for complementing estimates and
on-shell BRST incidence.  It may not be inserted into the usual contour
formula for \(e^{-tA}\) as though it were \((A+z)^{-1}\) for a fixed
self-adjoint \(A\).  A fixed enlarged realization or a proof of the
appropriate operator-pencil calculus is required first.
\end{firewall}

\subsection{The natural off-shell repair loses the parameter floor}
\label{subsec:V121-natural-repair}

One may try to repair the normal gauge row by using the gauge-fixing
functional.  In the product collar,
\begin{equation}
 \delta a=\delta_\Sigma\alpha-\partial_r\phi
 \qquad
 (a=\alpha+\phi\,dr).
 \label{eq:V121-delta-a}
\end{equation}
Let \(A(\Lambda_0)\) and \(C(\Lambda_0)\) be order-zero scalar
pseudodifferential multipliers and consider
\begin{equation}
 N_{A,C}(a)
 :=
 A(\Lambda_0)B_0\phi+C(\Lambda_0)\delta a.
 \label{eq:V121-NAC}
\end{equation}

\begin{theorem}[Off-shell incidence versus parameter ellipticity]
\label{thm:V121-repair-nogo}
Assume the tangential row is \(B_1\alpha=0\).  If
\begin{equation}
 N_{A,C}(dc)=0
 \label{eq:V121-offshell-invariance}
\end{equation}
for every scalar boundary jet satisfying \(B_0c=0\), then
\begin{equation}
 C(\Lambda_0)=-A(\Lambda_0)
 \label{eq:V121-C-minus-A}
\end{equation}
and the normal row reduces to
\begin{equation}
 N_{A,-A}(a)
 =
 A(\Lambda_0)
 \bigl(\Lambda_0\phi-\delta_\Sigma\alpha\bigr).
 \label{eq:V121-invariant-row}
\end{equation}
Its stable principal symbol, after the tangential row has killed
\(\alpha\), is
\begin{equation}
 A(\mu)\mu\,\phi.
 \label{eq:V121-repair-stable-symbol}
\end{equation}
For a classical order-zero \(A\) regular at \(\mu=0\), this symbol has no
uniform parameter floor on
\(\mu^2+|z|=1\).
\end{theorem}

\begin{proof}
For a gauge variation \(a=dc\),
\[
 B_0\phi=B_0\partial_rc=\gamma P_0c,
 \qquad
 \delta(dc)=P_0c.
\]
The second normal jet of \(c\) is arbitrary after the first-order condition
\(B_0c=0\) is imposed, so \(\gamma P_0c\) is arbitrary.  Therefore
\eqref{eq:V121-offshell-invariance} forces \(A+C=0\), proving
\eqref{eq:V121-C-minus-A}.  Substitution of
\eqref{eq:V121-delta-a} into \(B_0\phi-\delta a\) cancels the two normal
derivatives and gives \eqref{eq:V121-invariant-row}.

On a stable parameter mode, let \(|\xi|=\mu\) and
\(\lambda=\sqrt{\mu^2+z}\).  The tangential boundary symbol is
\((\lambda+\mu)\alpha\), so \(\alpha=0\).  The remaining symbol of
\eqref{eq:V121-invariant-row} is \(A(\mu)\mu\phi\).
At \(\mu=0\), \(z=1\), it vanishes.  Continuity gives zero infimum on the
normalized parameter hemisphere.
\end{proof}

\begin{assessment}[Scope of the no-go]
\label{ass:V121-nogo-scope}
Theorem~\ref{thm:V121-repair-nogo} closes the most direct fixed first-order
repair built from \(B_0\phi\) and \(\delta a\).  It does not exclude an
enlarged boundary phase space, a Wentzell/dynamic boundary operator, or a
different Hodge boundary complex.  Any successful alternative must supply
both an off-shell domain chain and a uniform parameter floor.
\end{assessment}

\subsection{Frozen heat kernel of the fixed nonlocal graph}
\label{subsec:V121-frozen-heat}

Although the fixed graph is not an off-shell BRST complex, its own leading
heat coefficients can be computed exactly.  They provide a diagnostic for
any enlarged construction.

Let \(n=\dim\Sigma\), freeze the metric at a flat boundary point, and Fourier
transform tangentially.  At tangential length \(\rho=|\xi|\), the normal
operator is
\begin{equation}
 H_\rho=-\partial_x^2+\rho^2,
 \qquad x\ge0,
 \qquad
 \partial_xu(0)=\rho u(0).
 \label{eq:V121-halfline-operator}
\end{equation}

\begin{lemma}[Exact half-line boundary trace]
\label{lem:V121-halfline-trace}
After subtracting the infinite full-line bulk contribution, the heat trace
of \eqref{eq:V121-halfline-operator} is
\begin{equation}
 e^{-t\rho^2}b(\rho\sqrt t),
 \qquad
 \boxed{
 b(a)=-\frac14+\frac12e^{a^2}\operatorname{erfc}(a).}
 \label{eq:V121-b-function}
\end{equation}
\end{lemma}

\begin{proof}
For the half-line Laplacian with
\(\partial_xu(0)=\rho u(0)\), the heat kernel is
\begin{align*}
 k_\rho(x,y;t)
 &=
 e^{-t\rho^2}\Bigg[
 \frac{e^{-(x-y)^2/(4t)}+e^{-(x+y)^2/(4t)}}{\sqrt{4\pi t}}
 \\
 &\hspace{25mm}
 -\rho e^{\rho(x+y)+\rho^2t}
 \operatorname{erfc}\left(
 \frac{x+y}{2\sqrt t}+\rho\sqrt t
 \right)
 \Bigg].
\end{align*}
Direct differentiation at \(x=0\) verifies the boundary condition.
The diagonal image term integrates to \(1/4\).  With \(a=\rho\sqrt t\),
the last term integrates to
\[
 a e^{a^2}\int_0^\infty e^{2au}\operatorname{erfc}(u+a)\,du
 =
 \frac12\left[1-e^{a^2}\operatorname{erfc}(a)\right],
\]
where integration by parts gives
\[
 \int_0^\infty e^{2au}\operatorname{erfc}(u+a)\,du
 =
 \frac{e^{-a^2}-\operatorname{erfc}(a)}{2a}
\]
and the value at \(a=0\) follows by continuity.  Subtracting this term from
\(1/4\) proves \eqref{eq:V121-b-function}.
\end{proof}

\begin{theorem}[Universal leading boundary coefficient]
\label{thm:V121-leading-boundary-coefficient}
For a rank-\(r\) Laplace-type operator in dimension \(d=n+1\) with principal
boundary row
\(-\partial_r+|D_\Sigma|\), the leading two heat terms are
\begin{equation}
 \operatorname{Tr}(e^{-tA})
 =
 \frac{r\,{\rm Vol}(M)}{(4\pi t)^{(n+1)/2}}
 +
 \frac{r\,\beta_n\,{\rm Vol}(\Sigma)}{(4\pi t)^{n/2}}
 +o(t^{-n/2}),
 \label{eq:V121-leading-heat}
\end{equation}
where
\begin{equation}
 \boxed{
 \beta_n
 =
 -\frac14+
 \frac{\Gamma((n+1)/2)}
 {n\sqrt\pi\,\Gamma(n/2)}.}
 \label{eq:V121-beta-n}
\end{equation}
In boundary dimension three,
\begin{equation}
 \boxed{
 \beta_3=-\frac14+\frac{2}{3\pi}
 =\frac{8-3\pi}{12\pi}.}
 \label{eq:V121-beta-3}
\end{equation}
\end{theorem}

\begin{proof}
The bulk term is the usual frozen Euclidean coefficient.  For the boundary
term, integrate Lemma~\ref{lem:V121-halfline-trace} over tangential Fourier
variables and scale \(\eta=\sqrt t\,\xi\):
\begin{align*}
 C_n
 &:=
 \frac1{(2\pi)^n}\int_{\mathbb R^n}
 e^{-|\eta|^2}b(|\eta|)\,d\eta\\
 &=
 \frac{\omega_{n-1}}{(2\pi)^n}
 \left[
 -\frac18\Gamma(n/2)
 +\frac{\Gamma((n+1)/2)}{2n\sqrt\pi}
 \right].
\end{align*}
Here
\[
 \int_0^\infty r^{n-1}\operatorname{erfc}(r)\,dr
 =
 \frac{\Gamma((n+1)/2)}{n\sqrt\pi}.
\]
Since
\(C_n=(4\pi)^{-n/2}\beta_n\), simplifying with
\(\omega_{n-1}=2\pi^{n/2}/\Gamma(n/2)\) proves
\eqref{eq:V121-beta-n}.  Setting \(n=3\) gives
\eqref{eq:V121-beta-3}.  Localization and parabolic scaling show that
lower-order symbols and curvature contribute only after these two
homogeneous terms.
\end{proof}

\subsection{Diagnostic gauge/ghost counterterms}
\label{subsec:V121-diagnostic-counterterms}

At the Lie-algebra level,
\begin{equation}
 \dim\mathfrak g_{\rm SM}=8+3+1=12.
 \label{eq:V121-gauge-dimension}
\end{equation}
Let \(A_{0,B}\) be the fixed scalar realization and \(A_{1,B}\) the fixed
rank-four one-form realization, both with the common principal graph.
Define the candidate heat combination
\begin{equation}
 K_{\rm fix}(t)
 :=
 12\left[
 \operatorname{Tr}(e^{-tA_{0,B}})
 -\frac12\operatorname{Tr}(e^{-tA_{1,B}})
 \right].
 \label{eq:V121-Kfix}
\end{equation}

\begin{corollary}[The first two fixed-graph coefficients survive]
\label{cor:V121-fixed-coefficients}
In four bulk dimensions,
\begin{equation}
 \boxed{
 K_{\rm fix}(t)
 =
 -\frac{12\,{\rm Vol}(M)}{(4\pi t)^2}
 -
 \frac{12\beta_3\,{\rm Vol}(\Sigma)}{(4\pi t)^{3/2}}
 +o(t^{-3/2}).}
 \label{eq:V121-Kfix-leading}
\end{equation}
Thus the volume and boundary-area coefficients do not cancel in the
rank-weighted fixed-graph combination.
\end{corollary}

\begin{proof}
The scalar rank is one and the one-form rank is four, so
\[
 1-\frac12(4)=-1.
\]
Apply Theorem~\ref{thm:V121-leading-boundary-coefficient} and multiply by
the gauge dimension \eqref{eq:V121-gauge-dimension}.
\end{proof}

\begin{proposition}[First common subtraction for the fixed realization]
\label{prop:V121-first-subtraction}
The expression
\begin{equation}
 K_{\rm fix}^{[1/2]}(t)
 :=
 K_{\rm fix}(t)
 +\frac{12\,{\rm Vol}(M)}{(4\pi t)^2}
 +\frac{12\beta_3\,{\rm Vol}(\Sigma)}{(4\pi t)^{3/2}}
 \label{eq:V121-first-subtraction}
\end{equation}
satisfies
\begin{equation}
 K_{\rm fix}^{[1/2]}(t)=o(t^{-3/2}).
 \label{eq:V121-first-subtraction-remainder}
\end{equation}
The two printed subtraction densities are respectively the bulk
cosmological-volume and boundary-tension counterterms for this fixed
realization.
\end{proposition}

\begin{proof}
Subtract the two coefficients in
\eqref{eq:V121-Kfix-leading}.  Both are local integrals of the unit density
on the bulk or boundary.
\end{proof}

\begin{firewall}[These are diagnostic, not BRST coefficients]
\label{fire:V121-diagnostic-only}
Equations \eqref{eq:V121-Kfix-leading} and
\eqref{eq:V121-first-subtraction} are rigorous heat statements for the
fixed self-adjoint graph.  Because of
Theorem~\ref{thm:V121-spectral-domain-defect}, they do not establish a BRST
heat cancellation or a gauge-independent Standard Model counterterm.
Moreover, four-dimensional proper-time renormalization also requires the
subsequent \(t^{-1}\), \(t^{-1/2}\), and \(t^0\) coefficients.  V121 does
not claim that the first subtraction is sufficient.
\end{firewall}

\subsection{V121 ledger and next acquisition}
\label{subsec:V121-ledger}

\paragraph{New exact conclusions.}
\begin{enumerate}[label=\textup{(\arabic*)},leftmargin=2.8em]
\item Proposition~\ref{prop:V121-complementarity} confirms that the fixed
      nonlocal graph is a genuine parameter-elliptic realization rather
      than a tautological interior Calderón condition.
\item Theorem~\ref{thm:V121-spectral-domain-defect} computes the exact
      nonzero-energy normal BRST defect
      \(B_0\partial_rc=\zeta\gamma c\).
\item Corollary~\ref{cor:V121-V120-correction} blocks the unsupported use of
      the V120 determinant card as a BRST heat complex.
\item Propositions~\ref{prop:V121-pencil-incidence} and
      \ref{prop:V121-pencil-not-fixed-resolvent} separate the exact
      parameter pencil from a fixed heat realization.
\item Theorem~\ref{thm:V121-repair-nogo} proves that the natural off-shell
      gauge-fixing repair cancels the normal parameter root and loses its
      floor.
\item Lemma~\ref{lem:V121-halfline-trace} and
      Theorem~\ref{thm:V121-leading-boundary-coefficient} compute the exact
      leading nonlocal boundary heat constant
      \(\beta_3=-1/4+2/(3\pi)\).
\item Corollary~\ref{cor:V121-fixed-coefficients} prints the surviving
      volume and boundary-area terms of the fixed rank-weighted diagnostic.
\end{enumerate}

\paragraph{Next forced step.}
V122 should seek a fixed enlarged realization whose Schur complement
reproduces the parameter-dependent square-root pencil while its total
domain is independent of \(z\).  The first candidate is an auxiliary
half-cylinder realization of
\(\sqrt{\Delta_\Sigma+z}\).  Its scalar and one-form extensions must be
placed in one de Rham complex, and its action, self-adjointness, parameter
floor, and off-shell BRST domain must all be checked before returning to
the remaining heat coefficients.  If that extension fails, the programme
must compare absolute/relative Hodge domains with the Lorentzian stable
floor rather than continue the invalid fixed-graph heat cancellation.

\begin{assessment}[V121 continuum status]
\label{ass:V121-continuum-status}
The full Standard--Model--Einstein continuum remains unproved.  V121
identifies an off-shell quantum-domain obstruction hidden by the successful
zero-energy incidence and computes the first fixed-graph heat diagnostic.
A fixed enlarged BRST realization, complete heat subtraction, fermion
branch selection, nonlinear existence, Einstein stress closure, and the
cofinal interacting limit remain open.
\end{assessment}

\section*{V121 exact checks and append-only record}
\addcontentsline{toc}{section}{V121 exact checks and append-only record}

The spectral defect and heat checksum is
\begin{equation}
 \boxed{
 B_0(\partial_rc)|_\Sigma=\zeta\,\gamma c,
 \qquad
 \beta_3=-\frac14+\frac{2}{3\pi}.}
 \label{eq:V121-final-check}
\end{equation}
The append operation preserves the complete V120 prefix byte for byte before
its former live terminator.  The assembled V121 file is checked for one live
final terminator, balanced environments and braces, duplicate labels, newly
introduced unresolved references, display math delimiters, and control
characters before the exact V120 predecessor is removed.

% =============================================================================
% END APPENDED VERSION V121 MATERIAL
% =============================================================================

\clearpage
% =============================================================================
% BEGIN APPENDED VERSION V122 MATERIAL
% =============================================================================
\section{V122: a finite auxiliary collar gives a fixed off-shell de Rham
realization}
\label{sec:V122-auxiliary-collar}

V121 proved that freezing the square-root graph at \(z=0\) gives a valid
self-adjoint boundary realization but not an off-shell BRST heat complex.
The parameter-dependent square root repairs incidence, yet its moving
domain is not the resolvent of a fixed operator on the physical bulk alone.

V122 resolves this dichotomy on a product collar by retaining the geometric
system whose Schur complement is the boundary square root.  Attach a finite
auxiliary cylinder to the physical boundary and impose absolute Hodge
conditions only at the far end.  The resulting enlarged manifold carries a
fixed, nonnegative, self-adjoint Hodge Laplacian and an exact de Rham
complex.  Eliminating the auxiliary collar produces a
parameter-dependent boundary pencil on the physical side.  The tangential
and normal one-form channels acquire different hyperbolic functions, and
that difference cancels the V121 spectral defect.

\subsection{The finite collar and its scalar Dirichlet-to-Neumann maps}
\label{subsec:V122-scalar-cylinder}

Let
\begin{equation}
 X_R:=[0,R]_s\times\Sigma,
 \qquad
 g_X=ds^2+\gamma,
 \qquad R>0.
 \label{eq:V122-XR}
\end{equation}
For a nonnegative self-adjoint boundary Laplacian \(\Delta_\Sigma\) and
\(z\) with \(\operatorname{Re}z\ge0\), write
\begin{equation}
 \Lambda(z):=(\Delta_\Sigma+z)^{1/2},
 \qquad
 \operatorname{Re}\Lambda(z)\ge0.
 \label{eq:V122-Lambda-z}
\end{equation}
Functional calculus defines
\begin{align}
 S_R^N(z)
 &:=
 \Lambda(z)\tanh\!\bigl(R\Lambda(z)\bigr),
 \label{eq:V122-SN}\\
 S_R^D(z)
 &:=
 \Lambda(z)\coth\!\bigl(R\Lambda(z)\bigr),
 \label{eq:V122-SD}
\end{align}
where the continuous value of the second function at \(\Lambda=0\) is
\(R^{-1}\).

\begin{lemma}[Exact finite-cylinder solutions]
\label{lem:V122-cylinder-solutions}
For boundary data \(q\), the solutions of
\begin{equation}
 (-\partial_s^2+\Delta_\Sigma+z)v=0,
 \qquad
 v(0)=q
 \label{eq:V122-cylinder-equation}
\end{equation}
with respectively Neumann and Dirichlet conditions at \(s=R\) are
\begin{align}
 v_N(s)
 &=
 \frac{\cosh((R-s)\Lambda(z))}
 {\cosh(R\Lambda(z))}\,q,
 \label{eq:V122-vN}\\
 v_D(s)
 &=
 \frac{\sinh((R-s)\Lambda(z))}
 {\sinh(R\Lambda(z))}\,q.
 \label{eq:V122-vD}
\end{align}
Their outward conormal at \(s=0\) is
\begin{equation}
 -\partial_sv_N(0)=S_R^N(z)q,
 \qquad
 -\partial_sv_D(0)=S_R^D(z)q.
 \label{eq:V122-cylinder-DtN}
\end{equation}
\end{lemma}

\begin{proof}
The assertions hold on each spectral subspace
\(\Delta_\Sigma\psi=\mu^2\psi\), where they reduce to the elementary
solutions of
\(-v''+\lambda^2v=0\) with
\(\lambda=\sqrt{\mu^2+z}\).  Differentiation at \(s=0\) gives
\(\lambda\tanh(R\lambda)\) and
\(\lambda\coth(R\lambda)\).  Spectral functional calculus then proves the
operator identities.  At \(\lambda=0\), the Neumann solution is constant
and the Dirichlet solution is \((1-s/R)q\), giving the declared continuous
values.
\end{proof}

\begin{proposition}[Exact Schur energies]
\label{prop:V122-Schur-energies}
For real \(z\ge0\), the solutions in
Lemma~\ref{lem:V122-cylinder-solutions} uniquely minimize
\begin{equation}
 \mathcal E_z[v]
 :=
 \int_0^R
 \left(
 \|\partial_sv\|^2+
 \|\Delta_\Sigma^{1/2}v\|^2+
 z\|v\|^2
 \right)\,ds
 \label{eq:V122-cylinder-energy}
\end{equation}
with their stated endpoint data, and
\begin{align}
 \mathcal E_z[v_N]
 &=
 \langle q,S_R^N(z)q\rangle,
 \label{eq:V122-energy-N}\\
 \mathcal E_z[v_D]
 &=
 \langle q,S_R^D(z)q\rangle.
 \label{eq:V122-energy-D}
\end{align}
\end{proposition}

\begin{proof}
Strict convexity on the constrained affine space, modulo the constant
Neumann zero mode at \(z=0\), gives the minimizing Euler equation
\eqref{eq:V122-cylinder-equation}.  For a solution, integration by parts
gives
\[
 \mathcal E_z[v]
 =
 \langle v(R),\partial_sv(R)\rangle
 -\langle v(0),\partial_sv(0)\rangle.
\]
The first term vanishes for either the Neumann or Dirichlet endpoint
condition.  Insert \eqref{eq:V122-cylinder-DtN}.
\end{proof}

\subsection{A fixed enlarged scalar realization}
\label{subsec:V122-fixed-scalar}

Assume first that the physical metric is product near
\(\Sigma=\partial M\).  Glue \(X_R\) to \(M\) along
\(\{0\}\times\Sigma\) and call the resulting compact manifold
\begin{equation}
 \widehat M_R:=M\cup_\Sigma X_R.
 \label{eq:V122-Mhat}
\end{equation}
The interface is not a boundary.  The only new boundary is
\(\Sigma_R=\{R\}\times\Sigma\).

Equivalently, before gluing, write a scalar as a pair \((u,v)\) with
\begin{equation}
 \gamma_Mu=v(0),
 \qquad
 -\partial_ru-\partial_sv(0)=0.
 \label{eq:V122-scalar-transmission}
\end{equation}
The first row is continuity and the second is the sum of the two outward
conormals.

\begin{theorem}[Closed fixed scalar action]
\label{thm:V122-fixed-scalar}
The quadratic form
\begin{equation}
 \widehat{\mathfrak q}_{0,R}[u,v]
 :=
 \int_M|du|^2+\int_{X_R}|dv|^2
 \label{eq:V122-scalar-form}
\end{equation}
on \(H^1\)-pairs with the trace continuity in
\eqref{eq:V122-scalar-transmission}, together with Neumann condition at
\(\Sigma_R\), is densely defined, nonnegative, and closed after adding the
\(L^2\) norm.  Its associated operator is the scalar Laplacian on
\(\widehat M_R\) with a parameter-independent domain.
\end{theorem}

\begin{proof}
The trace-equality subspace is closed in
\(H^1(M)\oplus H^1(X_R)\).  On it,
\(\widehat{\mathfrak q}_{0,R}+\|\cdot\|_{L^2}^2\) is the squared \(H^1\)
norm up to equivalence, hence complete.  The representation theorem gives
a nonnegative self-adjoint operator.  Green's formula gives continuity and
the conormal balance at the interface; the natural far endpoint condition
is Neumann.  These are exactly the weak transmission conditions for the
Laplacian on the glued manifold.
\end{proof}

\begin{theorem}[The moving physical pencil is a fixed enlarged Schur
complement]
\label{thm:V122-fixed-Schur}
For a source supported in \(M\), solve
\begin{equation}
 (\widehat\Delta_{0,R}+z)(u,v)=(f,0).
 \label{eq:V122-enlarged-resolvent}
\end{equation}
Eliminating \(v\) gives the physical problem
\begin{align}
 (P_0+z)u&=f,
 \label{eq:V122-physical-equation}\\
 \bigl(-\partial_r+S_R^N(z)\bigr)\gamma_Mu&=0.
 \label{eq:V122-physical-SN}
\end{align}
Thus the moving-domain physical family is the Schur compression of one
fixed self-adjoint enlarged resolvent.
\end{theorem}

\begin{proof}
The auxiliary component has no source, so
Lemma~\ref{lem:V122-cylinder-solutions} gives
\(-\partial_sv(0)=S_R^N(z)\gamma_Mu\).  Substitute this into the conormal
balance \eqref{eq:V122-scalar-transmission}.  The physical bulk row is the
first component of \eqref{eq:V122-enlarged-resolvent}.
\end{proof}

\begin{corollary}[Recovery of the V115 boundary action]
\label{cor:V122-V115-action}
On the positive boundary spectrum,
\begin{equation}
 S_R^N(0)\longrightarrow\Lambda
 \quad\hbox{strongly as }R\to\infty,
 \label{eq:V122-SN-limit}
\end{equation}
and the minimized cylinder energy converges to
\(\langle q,\Lambda q\rangle\).  Therefore the V115 nonlocal boundary
action is the infinite-collar Schur energy, while each finite \(R\) retains
a compact fixed realization.
\end{corollary}

\begin{proof}
For every \(\mu>0\),
\(\mu\tanh(R\mu)\to\mu\).  At \(\mu=0\), both sides vanish.
Dominated spectral convergence on the form domain proves the strong form
limit.  Apply \eqref{eq:V122-energy-N}.
\end{proof}

\subsection{The full Hodge complex on the enlarged manifold}
\label{subsec:V122-Hodge-complex}

At the far boundary impose the absolute Hodge conditions
\begin{equation}
 \iota_{\nu_R}\omega=0,
 \qquad
 \iota_{\nu_R}d\omega=0.
 \label{eq:V122-absolute}
\end{equation}
For a scalar, the second row is Neumann.  For a one-form
\(\omega=\beta+\phi\,ds\) on the product cylinder, the two rows become
\begin{equation}
 \phi(R)=0,
 \qquad
 \partial_s\beta(R)=0.
 \label{eq:V122-oneform-far}
\end{equation}

\begin{theorem}[Fixed self-adjoint de Rham realization]
\label{thm:V122-fixed-deRham}
For every \(p\), the Hodge energy
\begin{equation}
 \widehat{\mathfrak q}_{p,R}[\omega]
 :=
 \|d\omega\|_{L^2(\widehat M_R)}^2
 +\|\delta\omega\|_{L^2(\widehat M_R)}^2
 \label{eq:V122-Hodge-form}
\end{equation}
with the absolute far-boundary conditions defines a nonnegative
self-adjoint Hodge Laplacian \(\widehat\Delta_{p,R}^{\rm abs}\).  On the
smooth core and, by closure, on the Hilbert-complex domains,
\begin{align}
 \widehat\Delta_{p+1,R}^{\rm abs}d
 &=
 d\widehat\Delta_{p,R}^{\rm abs},
 \label{eq:V122-offshell-chain}\\
 e^{-t\widehat\Delta_{p+1,R}^{\rm abs}}d
 &=
 d\,e^{-t\widehat\Delta_{p,R}^{\rm abs}}.
 \label{eq:V122-heat-chain}
\end{align}
Thus the enlarged scalar and one-form heat operators form an exact
off-shell de Rham complex.
\end{theorem}

\begin{proof}
The absolute Hodge form on a compact manifold with boundary is closed and
nonnegative.  The boundary conditions are preserved by \(d\): the first
condition for \(d\omega\) is the second condition for \(\omega\), and the
second is \(\iota_{\nu_R}d^2\omega=0\).  The differential identity
\(\Delta d=d\Delta\) holds on the smooth invariant core.  Closure gives the
Hilbert-complex operator identity.  The heat identity follows from spectral
functional calculus, first for bounded spectral truncations and then by
strong convergence on \(\operatorname{Dom}d\).
\end{proof}

\begin{corollary}[The harmonic sector is topological and discrete]
\label{cor:V122-harmonic-sector}
The enlarged harmonic spaces satisfy
\begin{equation}
 \ker\widehat\Delta_{p,R}^{\rm abs}
 \cong H^p_{\rm dR}(\widehat M_R)
 \cong H^p_{\rm dR}(M).
 \label{eq:V122-harmonic-cohomology}
\end{equation}
In particular, when \(M\) is connected, the scalar kernel is the single
constant line before tensoring with \(\mathfrak g_{\rm SM}\).
\end{corollary}

\begin{proof}
Absolute Hodge theory identifies the kernel with ordinary de Rham
cohomology.  Attaching a collar does not change homotopy type, so
\(\widehat M_R\) deformation retracts onto \(M\).
\end{proof}

\subsection{Tangential and normal Schur channels}
\label{subsec:V122-form-channels}

On the product cylinder, the Hodge Laplacian separates as
\begin{equation}
 \Delta_{X_R,p}
 \cong
 \bigl(-\partial_s^2+\Delta_{\Sigma,p}\bigr)
 \oplus
 ds\wedge
 \bigl(-\partial_s^2+\Delta_{\Sigma,p-1}\bigr).
 \label{eq:V122-Hodge-separation}
\end{equation}
The far absolute conditions are Neumann on the tangential summand and
Dirichlet on the normal summand.

\begin{proposition}[Exact one-form Schur pencil]
\label{prop:V122-oneform-Schur}
After eliminating the auxiliary cylinder, the physical one-form boundary
pencil in the product splitting \(a=\alpha+\phi\,dr\) is
\begin{equation}
 \boxed{
 \begin{aligned}
 \bigl(-\partial_r+S_{R,1}^N(z)\bigr)\alpha&=0,\\
 \bigl(-\partial_r+S_{R,0}^D(z)\bigr)\phi&=0,
 \end{aligned}}
 \label{eq:V122-oneform-Schur}
\end{equation}
where
\begin{align}
 S_{R,1}^N(z)
 &=
 \Lambda_1(z)\tanh(R\Lambda_1(z)),
 \notag\\
 S_{R,0}^D(z)
 &=
 \Lambda_0(z)\coth(R\Lambda_0(z)).
 \label{eq:V122-oneform-SND}
\end{align}
\end{proposition}

\begin{proof}
Equation \eqref{eq:V122-Hodge-separation} reduces the two components to
Lemma~\ref{lem:V122-cylinder-solutions}.  Smooth gluing identifies the
tangential traces directly.  If \(t\) is the signed normal coordinate,
then \(t=r\) on \(M\) and \(t=-s\) on \(X_R\); hence a normal coefficient
relative to \(ds\) has the opposite sign from its coefficient relative to
\(dr\).  The two sign changes in trace and conormal cancel, leaving the
second row of \eqref{eq:V122-oneform-Schur}.
\end{proof}

\begin{theorem}[Exact off-shell scalar-to-one-form incidence after
elimination]
\label{thm:V122-effective-incidence}
Let a product scalar mode satisfy
\begin{equation}
 (P_0+z)c=0,
 \qquad
 \bigl(-\partial_r+S_{R,0}^N(z)\bigr)c=0.
 \label{eq:V122-effective-scalar}
\end{equation}
Then \(dc=d_\Sigma c+(\partial_rc)\,dr\) satisfies both rows of
\eqref{eq:V122-oneform-Schur}.
\end{theorem}

\begin{proof}
The tangential row follows from functional calculus:
\[
 d_\Sigma S_{R,0}^N(z)=S_{R,1}^N(z)d_\Sigma.
\]
For the normal row, work on
\(\Delta_{\Sigma,0}\psi=\mu^2\psi\) and put
\(\lambda=\sqrt{\mu^2+z}\).  The scalar boundary row gives
\[
 \partial_rc
 =
 \lambda\tanh(R\lambda)c.
\]
The bulk equation gives
\(\partial_r^2c=\lambda^2c\).  Therefore
\begin{align*}
 \bigl(-\partial_r+\lambda\coth(R\lambda)\bigr)\partial_rc
 &=
 -\lambda^2c
 +\lambda\coth(R\lambda)
  \lambda\tanh(R\lambda)c\\
 &=0.
\end{align*}
The zero-eigenvalue limits follow continuously.  This proves the normal
row.
\end{proof}

\begin{corollary}[Cancellation of the V121 spectral defect]
\label{cor:V122-defect-cancelled}
The normal V121 defect arose from using the same frozen Neumann-type
square-root row on the scalar and normal one-form components.  In the
fixed enlarged realization, the scalar channel is Neumann at the far end
and the normal one-form channel is Dirichlet.  The exact identity
\begin{equation}
 \bigl[\lambda\coth(R\lambda)\bigr]
 \bigl[\lambda\tanh(R\lambda)\bigr]
 =
 \lambda^2
 \label{eq:V122-tanh-coth}
\end{equation}
cancels the spectral defect for every \(\lambda\).
\end{corollary}

\subsection{Uniform parameter floor}
\label{subsec:V122-floor}

\begin{theorem}[Finite-collar parameter floor]
\label{thm:V122-parameter-floor}
On a stable physical mode with
\(\lambda=\sqrt{\mu^2+z}\), \(\operatorname{Re}\lambda>0\), the scalar and
tangential boundary symbol and the normal boundary symbol are respectively
\begin{align}
 b_N(\lambda)
 &=
 \lambda\bigl[1+\tanh(R\lambda)\bigr]
 =
 \frac{2\lambda}{1+e^{-2R\lambda}},
 \label{eq:V122-bN}\\
 b_D(\lambda)
 &=
 \lambda\bigl[1+\coth(R\lambda)\bigr]
 =
 \frac{2\lambda}{1-e^{-2R\lambda}}.
 \label{eq:V122-bD}
\end{align}
For \(\mu^2+|z|=1\) and \(\operatorname{Re}z\ge0\),
\begin{equation}
 \boxed{
 |b_N(\lambda)|\ge|\lambda|\ge2^{-1/2},
 \qquad
 |b_D(\lambda)|\ge|\lambda|\ge2^{-1/2}.}
 \label{eq:V122-floor}
\end{equation}
\end{theorem}

\begin{proof}
The two hyperbolic identities give
\eqref{eq:V122-bN}--\eqref{eq:V122-bD}.  Since
\(|e^{-2R\lambda}|\le1\),
\[
 |1\pm e^{-2R\lambda}|\le2,
\]
so both displayed quotients have modulus at least \(|\lambda|\).
The root estimate is the one used in
Theorem~\ref{thm:V115-parameter-floor}.  On the normalized hemisphere the
origin is excluded, hence \(\operatorname{Re}\lambda>0\) and the
denominators are nonzero.
\end{proof}

\subsection{Comparison with the square-root symbols}
\label{subsec:V122-symbol-comparison}

\begin{proposition}[Finite-collar corrections are smoothing]
\label{prop:V122-smoothing}
At \(z=0\), on the positive spectral complement,
\begin{align}
 S_R^N(0)-\Lambda
 &=
 -2\Lambda(e^{2R\Lambda}+I)^{-1},
 \label{eq:V122-SN-smoothing}\\
 S_R^D(0)-\Lambda
 &=
 2\Lambda(e^{2R\Lambda}-I)^{-1}.
 \label{eq:V122-SD-smoothing}
\end{align}
Both differences are smoothing.  On the harmonic sector,
\begin{equation}
 S_R^N(0)\Pi=0,
 \qquad
 S_R^D(0)\Pi=R^{-1}\Pi.
 \label{eq:V122-harmonic-SND}
\end{equation}
Consequently the scalar and matrix principal/subprincipal symbols computed
in V117--V118 are unchanged by finite \(R\); only a smoothing plus
finite-rank sector is altered.
\end{proposition}

\begin{proof}
The scalar identities
\[
 \tanh x-1=-\frac{2}{e^{2x}+1},
 \qquad
 \coth x-1=\frac{2}{e^{2x}-1}
\]
give the formulas.  On the positive spectrum their multipliers decay
exponentially in \(\sqrt{\Delta_\Sigma}\), hence belong to
\(\Psi^{-\infty}\).  The values at zero were computed in
Lemma~\ref{lem:V122-cylinder-solutions}.  Smoothing and finite-rank
perturbations have zero classical homogeneous symbol at every order.
\end{proof}

\begin{firewall}[Nonproduct geometry remains to be extended]
\label{fire:V122-nonproduct-open}
V122 is exact for a product physical collar matched to a product auxiliary
cylinder.  A nonproduct metric requires a smooth continuation of the full
normal Hodge operator across the interface.  Reflecting only the scalar
instantaneous metric generally creates a jump in the second fundamental
form and would invalidate the V117--V118 curvature symbols.  The curved
extension must match the metric and connection jets or retain an explicit
transmission defect.
\end{firewall}

\subsection{Heat and determinant status}
\label{subsec:V122-heat-status}

\begin{assessment}[What the auxiliary collar fixes]
\label{ass:V122-what-fixed}
For every finite \(R\), the enlarged operator is compact-resolvent,
self-adjoint, action-derived, parameter-elliptic, and an exact off-shell
de Rham complex.  Its ordinary heat semigroup therefore supports genuine
BRST spectral pairing.  The physical square-root pencil appears only after
Schur elimination and is not mistaken for a fixed physical-domain
resolvent.
\end{assessment}

\begin{firewall}[The auxiliary collar has its own heat content]
\label{fire:V122-auxiliary-heat}
The heat trace of \(\widehat M_R\) includes the volume and far-boundary
coefficients of \(X_R\).  They are regulator data, not additional Standard
Model particles.  Before sending \(R\to\infty\), the programme must form a
relative heat trace or determinant which subtracts a declared reference
cylinder and retains the interface Schur determinant.  V122 does not count
the auxiliary collar as physical bulk.
\end{firewall}

\subsection{V122 ledger and next acquisition}
\label{subsec:V122-ledger}

\paragraph{New exact conclusions.}
\begin{enumerate}[label=\textup{(\arabic*)},leftmargin=2.8em]
\item Lemma~\ref{lem:V122-cylinder-solutions} and
      Proposition~\ref{prop:V122-Schur-energies} derive the finite-cylinder
      Neumann and Dirichlet Dirichlet-to-Neumann maps.
\item Theorems~\ref{thm:V122-fixed-scalar} and
      \ref{thm:V122-fixed-Schur} realize the moving physical pencil as the
      Schur complement of one fixed compact self-adjoint operator.
\item Theorem~\ref{thm:V122-fixed-deRham} constructs the exact off-shell
      Hodge heat complex.
\item Proposition~\ref{prop:V122-oneform-Schur} proves that tangential and
      normal one-form channels carry the Neumann and Dirichlet hyperbolic
      roots respectively.
\item Theorem~\ref{thm:V122-effective-incidence} and
      Corollary~\ref{cor:V122-defect-cancelled} cancel the V121 spectral
      defect by the identity \(\coth\,\tanh=1\).
\item Theorem~\ref{thm:V122-parameter-floor} preserves the normalized
      parameter floor for both channels.
\item Proposition~\ref{prop:V122-smoothing} proves that finite collar length
      changes no classical principal or subprincipal symbol.
\end{enumerate}

\paragraph{Next forced step.}
V123 should form the relative heat trace between the glued
\(\widehat M_R\) complex and a decoupled auxiliary reference, prove the
finite gluing determinant identity, and isolate the interface
Dirichlet-to-Neumann determinant without double-counting cylinder volume.
The scalar/one-form exact spectral pairing must be performed before the
relative trace is split.  The harmonic constant line and the
\(R^{-1}\) normal channel in \eqref{eq:V122-harmonic-SND} must remain
explicit.

\begin{assessment}[V122 continuum status]
\label{ass:V122-continuum-status}
The full Standard--Model--Einstein continuum remains unproved.  V122
constructs the first fixed enlarged realization that simultaneously passes
the action, parameter-floor, and off-shell de Rham-domain gates in the
product geometry.  Curved jet matching, relative heat gluing, fermion
boundary extension, nonlinear interacting existence, Einstein stress
closure, and the cofinal continuum remain open.
\end{assessment}

\section*{V122 exact checks and append-only record}
\addcontentsline{toc}{section}{V122 exact checks and append-only record}

The fixed-extension and incidence checksum is
\begin{equation}
 \boxed{
 S_R^N(z)=\Lambda(z)\tanh(R\Lambda(z)),
 \quad
 S_R^D(z)=\Lambda(z)\coth(R\Lambda(z)),
 \quad
 S_R^D(z)S_R^N(z)=\Delta_\Sigma+z.}
 \label{eq:V122-final-check}
\end{equation}
The append operation preserves the complete V121 prefix byte for byte before
its former live terminator.  The assembled V122 file is checked for one live
final terminator, balanced environments and braces, duplicate labels, newly
introduced unresolved references, display math delimiters, and control
characters before the exact V121 predecessor is removed.

% =============================================================================
% END APPENDED VERSION V122 MATERIAL
% =============================================================================

\clearpage
% =============================================================================
% BEGIN APPENDED VERSION V123 MATERIAL
% =============================================================================
\section{V123: complex-compatible relative heat gluing and the normalized
interface determinant}
\label{sec:V123-relative-gluing}

V122 replaced the moving physical square-root domain by a fixed Hodge
complex on the collar-extended manifold \(\widehat M_R\).  Its absolute heat
trace contains the auxiliary collar volume and its far-boundary
coefficients.  Those regulator contributions must be removed without
breaking the de Rham complex.

V123 chooses a complex-compatible decoupled reference.  On the physical
piece \(M\), impose relative Hodge conditions at the cut.  On the auxiliary
cylinder \(X_R\), impose absolute Hodge conditions at both ends.  The glued
and decoupled systems are then separately fixed de Rham complexes.  Their
relative heat trace has an exact longitudinal pairing, and heat locality
confines every surviving coefficient to the former interface.  A finite
Schur calculation gives the corresponding normalized interface
determinant.  The scalar matching operator is computed explicitly; the
one-form matching operator is kept as a matrix acquisition rather than
replaced by a componentwise scalar formula.

\subsection{The glued and decoupled Hodge complexes}
\label{subsec:V123-two-complexes}

Let
\(\widehat\Delta_{p,R}^{\rm gl}\) be the V122 Hodge Laplacian on
\(\widehat M_R\), with absolute conditions at the far boundary
\(\Sigma_R\).  At the artificial cut \(\Sigma\), define the relative
conditions on the physical piece by
\begin{equation}
 \nu^\flat\wedge\omega=0,
 \qquad
 \nu^\flat\wedge\delta\omega=0,
 \label{eq:V123-relative-BC}
\end{equation}
and the absolute conditions on the auxiliary side by
\begin{equation}
 \iota_\nu\omega=0,
 \qquad
 \iota_\nu d\omega=0.
 \label{eq:V123-absolute-BC}
\end{equation}
Write the corresponding operators as
\(\Delta_{M,p}^{\rm rel}\) and
\(\Delta_{X_R,p}^{\rm aa}\), where \(\mathrm{aa}\) means absolute at both
\(s=0\) and \(s=R\).  The decoupled operator is
\begin{equation}
 \widehat\Delta_{p,R}^{\rm dec}
 :=
 \Delta_{M,p}^{\rm rel}
 \oplus
 \Delta_{X_R,p}^{\rm aa}.
 \label{eq:V123-decoupled}
\end{equation}

\begin{proposition}[Both realizations are de Rham complexes]
\label{prop:V123-two-deRham-complexes}
The exterior derivative maps the scalar form domain into the one-form form
domain for both \(\widehat\Delta^{\rm gl}\) and
\(\widehat\Delta^{\rm dec}\).  On their smooth invariant cores,
\begin{equation}
 \widehat\Delta_{p+1,R}^{Y}d
 =
 d\widehat\Delta_{p,R}^{Y},
 \qquad
 Y\in\{{\rm gl},{\rm dec}\},
 \label{eq:V123-two-chain-identities}
\end{equation}
and the identities extend to their heat semigroups.
\end{proposition}

\begin{proof}
The glued assertion is Theorem~\ref{thm:V122-fixed-deRham}.  Absolute
conditions form a \(d\)-stable complex by the same proof.  Relative
conditions form the dual \(d\)-stable Hodge complex: for a scalar, the
first condition is Dirichlet; if its trace vanishes, the tangential trace of
its exterior derivative vanishes, which is the first relative one-form
row.  The second row is preserved on the operator core because
\(\delta d=\Delta_0\) and the relative scalar operator maps its domain back
to functions with zero boundary trace.  Equivalently, this is the standard
minimal de Rham extension and its Hodge Laplacian.  Direct sums preserve the
chain.  Spectral functional calculus gives the heat-semigroup extension.
\end{proof}

\subsection{The relative heat trace}
\label{subsec:V123-relative-heat}

Remove exact zero modes and define
\begin{equation}
 K_{p,R}^{\rm rel}(t)
 :=
 \operatorname{Tr}'e^{-t\widehat\Delta_{p,R}^{\rm gl}}
 -
 \operatorname{Tr}'e^{-t\widehat\Delta_{p,R}^{\rm dec}}.
 \label{eq:V123-relative-heat}
\end{equation}

\begin{theorem}[All relative heat coefficients are interface-supported]
\label{thm:V123-interface-locality}
At fixed finite \(R\), the small-time expansion of
\eqref{eq:V123-relative-heat} has no bulk coefficient from either \(M\) or
\(X_R\), and no coefficient supported at the far boundary \(\Sigma_R\).
Every surviving local coefficient is supported at the cut \(\Sigma\), where
the transmission condition is compared with the relative/absolute
decoupling.
\end{theorem}

\begin{proof}
The glued and decoupled differential operators have identical complete
symbols in the interiors of \(M\) and \(X_R\).  They also have identical
absolute boundary jets in a neighborhood of \(\Sigma_R\).  Heat
coefficients are local in the operator and boundary jets, so those
contributions cancel term by term in the difference.  The two realizations
differ only at \(\Sigma\): the glued problem has transmission across the
cut, whereas the reference has relative and absolute boundary conditions
on its two sides.  A partition of unity supported away from and near the
cut proves the claim.
\end{proof}

\begin{corollary}[The auxiliary volume is not counted as physical]
\label{cor:V123-no-auxiliary-volume}
The leading volume term proportional to
\(R\,{\rm Vol}(\Sigma)\) and every far-boundary heat coefficient cancel
exactly in \(K_{p,R}^{\rm rel}\).
\end{corollary}

\subsection{Exact longitudinal pairing before the split}
\label{subsec:V123-longitudinal-pairing}

For either \(Y={\rm gl}\) or \(Y={\rm dec}\), let
\(\mathcal H_{0,Y}^\perp\) be the orthogonal complement of scalar harmonic
forms and let
\(\mathcal E_{1,Y}=\overline{\operatorname{Ran}d}\) be the exact one-form
subspace.  Define
\begin{equation}
 U_Y
 :=
 d\bigl(\widehat\Delta_{0,R}^{Y}\bigr)^{-1/2}
 :
 \mathcal H_{0,Y}^\perp\longrightarrow\mathcal E_{1,Y}.
 \label{eq:V123-polar-d}
\end{equation}

\begin{theorem}[Unitary scalar--exact-one-form equivalence]
\label{thm:V123-longitudinal-equivalence}
The map \(U_Y\) is unitary onto \(\mathcal E_{1,Y}\) and
\begin{equation}
 \widehat\Delta_{1,R}^{Y}|_{\mathcal E_{1,Y}}
 =
 U_Y\widehat\Delta_{0,R}^{Y}|_{\mathcal H_{0,Y}^\perp}U_Y^\dagger.
 \label{eq:V123-unitary-equivalence}
\end{equation}
Consequently,
\begin{equation}
 \operatorname{Tr}_{\mathcal E_{1,Y}}
 e^{-t\widehat\Delta_{1,R}^{Y}}
 =
 \operatorname{Tr}'e^{-t\widehat\Delta_{0,R}^{Y}}
 \label{eq:V123-exact-heat-equality}
\end{equation}
for every \(t>0\).
\end{theorem}

\begin{proof}
On the reduced scalar space,
\[
 U_Y^\dagger U_Y
 =
 \Delta_{0,Y}^{-1/2}\delta d\Delta_{0,Y}^{-1/2}
 =I.
\]
Its range is dense in and closed onto
\(\overline{\operatorname{Ran}d}\), hence it is unitary onto
\(\mathcal E_{1,Y}\).  The chain identity
\(\Delta_{1,Y}d=d\Delta_{0,Y}\) gives
\eqref{eq:V123-unitary-equivalence}.  Apply heat functional calculus and
take the trace.
\end{proof}

\begin{corollary}[Exact relative longitudinal heat cancellation]
\label{cor:V123-relative-longitudinal}
Let
\[
 K_{1,R}^{\rm rel}
 =
 K_{1,{\rm ex},R}^{\rm rel}
 +
 K_{1,{\rm coex},R}^{\rm rel}
\]
be the exact/coexact orthogonal split after harmonic modes are removed.
Then
\begin{equation}
 \boxed{
 K_{1,{\rm ex},R}^{\rm rel}(t)
 =
 K_{0,R}^{\rm rel}(t).}
 \label{eq:V123-relative-exact-equality}
\end{equation}
For the Standard Model adjoint dimension \(12\), the relative
gauge/ghost heat combination is therefore
\begin{equation}
 \boxed{
 K_{{\rm BRST},R}^{\rm rel}(t)
 :=
 12\left(K_{0,R}^{\rm rel}
 -\frac12K_{1,R}^{\rm rel}\right)
 =
 6\left(
 K_{0,R}^{\rm rel}
 -K_{1,{\rm coex},R}^{\rm rel}
 \right).}
 \label{eq:V123-relative-BRST-heat}
\end{equation}
\end{corollary}

\begin{proof}
Subtract \eqref{eq:V123-exact-heat-equality} for the glued and decoupled
systems.  Insert the exact/coexact split and simplify the coefficients.
\end{proof}

\begin{firewall}[Order of operations]
\label{fire:V123-pair-before-split}
Equation \eqref{eq:V123-relative-exact-equality} uses the full glued complex
and the full decoupled complex.  It does not assert a separate BRST pairing
for an arbitrarily chosen physical Dirichlet determinant or for one scalar
factor of a raw one-form matching matrix.  Longitudinal pairing must be
performed before the relative determinant is factorized into interface
channels.
\end{firewall}

\subsection{Finite gluing algebra}
\label{subsec:V123-finite-gluing}

The determinant calculation is first made at a common finite regulator.
Let \(x\), \(y\), and \(q\) denote physical interior, auxiliary interior,
and interface variables.  Suppose the regulated positive quadratic form
has matrix
\begin{equation}
 \mathbb A(z)
 =
 \begin{pmatrix}
 A_M(z)&0&C_M\\
 0&A_X(z)&C_X\\
 C_M^\dagger&C_X^\dagger&D_M+D_X
 \end{pmatrix},
 \label{eq:V123-three-block}
\end{equation}
where \(A_M(z)\) and \(A_X(z)\) are the operators with homogeneous trace at
the cut.

\begin{theorem}[Exact finite gluing determinant]
\label{thm:V123-finite-gluing}
If \(A_M(z)\) and \(A_X(z)\) are invertible, define their boundary Schur
forms
\begin{align}
 N_M(z)&:=D_M-C_M^\dagger A_M(z)^{-1}C_M,
 \notag\\
 N_X(z)&:=D_X-C_X^\dagger A_X(z)^{-1}C_X.
 \label{eq:V123-two-DtN}
\end{align}
Then
\begin{equation}
 \boxed{
 \det\mathbb A(z)
 =
 \det A_M(z)\det A_X(z)
 \det\bigl(N_M(z)+N_X(z)\bigr).}
 \label{eq:V123-gluing-determinant}
\end{equation}
If a decoupled reference has the same two interior blocks and boundary
Schur form \(N_{\rm ref}(z)\), then
\begin{equation}
 \boxed{
 \frac{\det\mathbb A_{\rm gl}(z)}
 {\det\mathbb A_{\rm dec}(z)}
 =
 \det\mathcal W(z),
 \qquad
 \mathcal W(z)
 :=
 N_{\rm ref}(z)^{-1/2}
 [N_M(z)+N_X(z)]
 N_{\rm ref}(z)^{-1/2}.}
 \label{eq:V123-normalized-W}
\end{equation}
\end{theorem}

\begin{proof}
Eliminate \(x\) and \(y\) by two applications of
Theorem~\ref{thm:V120-diagonal-equivalence}.  Their Schur contributions add
on the common interface variable, giving
\eqref{eq:V123-gluing-determinant}.  Apply the same identity to the
reference.  The identical interior determinants cancel, leaving the ratio
of the two boundary determinants.  Congruence by
\(N_{\rm ref}^{-1/2}\) gives \eqref{eq:V123-normalized-W}.
\end{proof}

\begin{corollary}[No interface normalization without its determinant]
\label{cor:V123-no-interface-normalization-loss}
Changing to a boundary frame in which \(N_{\rm ref}=I\) is legitimate only
with the two factors \(N_{\rm ref}^{-1/2}\) in
\eqref{eq:V123-normalized-W}.  Replacing
\(N_M+N_X\) by a conditionally normalized channel and omitting
\(\det N_{\rm ref}\) changes the relative determinant.
\end{corollary}

\begin{proof}
Take determinants in \eqref{eq:V123-normalized-W}:
\[
 \det\mathcal W
 =
 \frac{\det(N_M+N_X)}{\det N_{\rm ref}}.
\]
The denominator is the Jacobian of the normalization.
\end{proof}

\subsection{The exact scalar interface determinant}
\label{subsec:V123-scalar-interface}

Let \(N_{M,0}(z)\) be the physical scalar Dirichlet-to-Neumann map with
outward conormal \(-\partial_r\).  The auxiliary scalar map with Neumann
condition at \(s=R\) is \(S_{R,0}^N(z)\).

\begin{theorem}[Normalized scalar gluing operator]
\label{thm:V123-scalar-W}
For \(z>0\), comparison of the glued scalar operator with the decoupled
physical-relative plus auxiliary-absolute scalar complex gives
\begin{equation}
 \boxed{
 \mathcal W_{0,R}(z)
 =
 S_{R,0}^N(z)^{-1/2}
 \bigl[N_{M,0}(z)+S_{R,0}^N(z)\bigr]
 S_{R,0}^N(z)^{-1/2}.}
 \label{eq:V123-scalar-W}
\end{equation}
At every common finite boundary regulator,
\begin{equation}
 \frac{
 \det(\widehat\Delta_{0,R}^{\rm gl}+z)}
 {
 \det(\Delta_{M,0}^{\rm rel}+z)
 \det(\Delta_{X_R,0}^{\rm aa}+z)}
 =
 \det\mathcal W_{0,R}(z).
 \label{eq:V123-scalar-relative-det}
\end{equation}
\end{theorem}

\begin{proof}
For scalars, relative boundary conditions on \(M\) are Dirichlet and
absolute conditions on \(X_R\) are Neumann.  If the interface value is
temporarily fixed to zero on the cylinder, the auxiliary interior
determinant is the Dirichlet-at-\(0\), Neumann-at-\(R\) determinant.
Releasing the interface value with no physical piece attached multiplies
that determinant by the auxiliary Schur form \(S_{R,0}^N(z)\).  Attaching
the physical piece replaces this Schur form by the conormal sum
\(N_{M,0}(z)+S_{R,0}^N(z)\).  The common physical Dirichlet and auxiliary
interior determinants cancel, so
Theorem~\ref{thm:V123-finite-gluing} gives
\eqref{eq:V123-scalar-W}--\eqref{eq:V123-scalar-relative-det}.
\end{proof}

\begin{proposition}[Scalar harmonic threshold]
\label{prop:V123-scalar-threshold}
On \(\ker\Delta_{\Sigma,0}\),
\begin{equation}
 S_{R,0}^N(z)
 =
 \sqrt z\,\tanh(R\sqrt z)
 =
 Rz+O(z^2)
 \qquad(z\downarrow0).
 \label{eq:V123-scalar-threshold}
\end{equation}
Thus \(\mathcal W_{0,R}(z)\) must be primed or combined with the global
gauge-volume factor before \(z=0\) is taken.
\end{proposition}

\begin{proof}
Use \(\tanh w=w-w^3/3+O(w^5)\).
\end{proof}

\subsection{The one-form matching matrix is not scalar}
\label{subsec:V123-oneform-open}

Let \(N_{M,1}(z)\) be the full physical one-form conormal response in the
tangential--normal trace splitting.  The raw auxiliary response from V122
is
\begin{equation}
 N_{X_R,1}^{\rm raw}(z)
 =
 \begin{pmatrix}
 S_{R,1}^N(z)&0\\
 0&S_{R,0}^D(z)
 \end{pmatrix}.
 \label{eq:V123-oneform-auxiliary-response}
\end{equation}

\begin{proposition}[Required one-form interface operator]
\label{prop:V123-oneform-W-requirement}
At a common finite trace regulator, the one-form relative determinant has
the form
\begin{equation}
 \frac{
 \det(\widehat\Delta_{1,R}^{\rm gl}+z)}
 {
 \det(\Delta_{M,1}^{\rm rel}+z)
 \det(\Delta_{X_R,1}^{\rm aa}+z)}
 =
 \det\mathcal W_{1,R}(z),
 \label{eq:V123-oneform-relative-det}
\end{equation}
where \(\mathcal W_{1,R}\) is the normalized Schur matching matrix obtained
from the full physical response, the auxiliary response
\eqref{eq:V123-oneform-auxiliary-response}, and the
relative/absolute reference polarizations.  It cannot be replaced by four
copies of \(\mathcal W_{0,R}\): its normal channel contains
\(S_{R,0}^D\), and its curved physical response contains the V118 shape
off-diagonal entries.
\end{proposition}

\begin{proof}
The finite determinant form follows from
Theorem~\ref{thm:V123-finite-gluing} after choosing the Hodge
relative/absolute boundary polarization.  Equations
\eqref{eq:V122-oneform-SND} and
\eqref{eq:V118-oneform-subprincipal} show respectively that the auxiliary
normal response and the curved physical response differ from a
componentwise scalar response.  Hence a scalar replacement changes the
Schur matrix.
\end{proof}

\begin{firewall}[The raw block is not yet the normalized Hodge matching
matrix]
\label{fire:V123-oneform-matrix-open}
Equation \eqref{eq:V123-oneform-auxiliary-response} is the raw
tangential--normal Dirichlet-to-Neumann response.  To obtain
\(\mathcal W_{1,R}\), V124 must write the finite symplectic change from raw
trace/conormal variables to the physical-relative and
auxiliary-absolute boundary polarizations.  Taking a determinant of the raw
diagonal block before that change would omit the reference Jacobian and
repeat the normalization error excluded in V120 and V123.
\end{firewall}

\subsection{Resolvent trace and relative determinant}
\label{subsec:V123-resolvent-trace}

\begin{theorem}[Finite Krein trace identity]
\label{thm:V123-Krein}
At every finite regulator and away from zero modes,
\begin{equation}
 \boxed{
 \operatorname{Tr}\!\left[
 (\mathbb A_{\rm gl}+z)^{-1}
 -(\mathbb A_{\rm dec}+z)^{-1}
 \right]
 =
 \operatorname{tr}\!\left[
 \mathcal W(z)^{-1}\partial_z\mathcal W(z)
 \right].}
 \label{eq:V123-Krein}
\end{equation}
\end{theorem}

\begin{proof}
Differentiate the logarithm of
\eqref{eq:V123-normalized-W}.  For a finite positive matrix,
\[
 \partial_z\log\det(\mathbb A+z)
 =
 \operatorname{Tr}(\mathbb A+z)^{-1}.
\]
The derivative of the right side is
\(\operatorname{tr}(\mathcal W^{-1}\partial_z\mathcal W)\).
\end{proof}

\begin{corollary}[Interface spectral data determine the relative heat]
\label{cor:V123-relative-heat-from-W}
The inverse Laplace transform of the right side of
\eqref{eq:V123-Krein} determines the regulated relative heat trace.
Therefore all interface heat coefficients can be extracted from the joint
large-\(z\), large-boundary-frequency symbol of the normalized matching
operator \(\mathcal W(z)\).
\end{corollary}

\begin{proof}
For positive finite matrices,
\[
 (\mathbb A+z)^{-1}
 =
 \int_0^\infty e^{-tz}e^{-t\mathbb A}\,dt.
\]
Subtract the two identities and use uniqueness of the Laplace transform.
\end{proof}

\subsection{V123 ledger and next acquisition}
\label{subsec:V123-ledger}

\paragraph{New exact conclusions.}
\begin{enumerate}[label=\textup{(\arabic*)},leftmargin=2.8em]
\item Proposition~\ref{prop:V123-two-deRham-complexes} constructs a
      de Rham-compatible decoupled reference.
\item Theorem~\ref{thm:V123-interface-locality} proves cancellation of all
      auxiliary bulk and far-boundary local heat coefficients.
\item Theorem~\ref{thm:V123-longitudinal-equivalence} and
      Corollary~\ref{cor:V123-relative-longitudinal} prove the exact
      scalar--longitudinal heat pairing before channel factorization.
\item Theorem~\ref{thm:V123-finite-gluing} proves the normalized finite
      interface determinant with its reference Jacobian.
\item Theorem~\ref{thm:V123-scalar-W} computes the scalar interface operator
      explicitly.
\item Proposition~\ref{prop:V123-scalar-threshold} identifies the global
      gauge threshold \(S_R^N(z)\sim Rz\).
\item Proposition~\ref{prop:V123-oneform-W-requirement} isolates the exact
      one-form matrix still required.
\item Theorem~\ref{thm:V123-Krein} expresses the relative resolvent and heat
      trace through the normalized interface response.
\end{enumerate}

\paragraph{Next forced step.}
V124 should construct the one-form boundary polarization explicitly in the
product collar.  It must convert raw tangential/normal values and conormals
to the relative/absolute reference variables, calculate
\(\mathcal W_{1,R}(z)\), and verify directly that its longitudinal
determinant channel matches \(\mathcal W_{0,R}(z)\).  Only the remaining
coexact interface determinant may enter the physical heat coefficient.
Curved shape mixing should then be reintroduced as the V118 order-zero
matrix perturbation.

\begin{assessment}[V123 continuum status]
\label{ass:V123-continuum-status}
The full Standard--Model--Einstein continuum remains unproved.  V123
constructs a genuine BRST-compatible relative heat problem and reduces its
finite spectral content to normalized interface determinants.  The
one-form Hodge matching matrix, curved relative coefficients, fermion
extension, nonlinear interacting existence, Einstein stress closure, and
the cofinal continuum remain open.
\end{assessment}

\section*{V123 exact checks and append-only record}
\addcontentsline{toc}{section}{V123 exact checks and append-only record}

The relative heat and scalar gluing checksum is
\begin{equation}
 \boxed{
 K_{1,{\rm ex},R}^{\rm rel}=K_{0,R}^{\rm rel},
 \qquad
 \mathcal W_{0,R}
 =
 (S_R^N)^{-1/2}(N_{M,0}+S_R^N)(S_R^N)^{-1/2}.}
 \label{eq:V123-final-check}
\end{equation}
The append operation preserves the complete V122 prefix byte for byte before
its former live terminator.  The assembled V123 file is checked for one live
final terminator, balanced environments and braces, duplicate labels, newly
introduced unresolved references, display math delimiters, and control
characters before the exact V122 predecessor is removed.

% =============================================================================
% END APPENDED VERSION V123 MATERIAL
% =============================================================================

\clearpage
% =============================================================================
% BEGIN APPENDED VERSION V124 MATERIAL
% =============================================================================
\section{V124: the product one-form boundary polarization and its
trace-class interface remainder}
\label{sec:V124-oneform-polarization}

V123 reduced the relative determinant to a normalized one-form interface
matrix but left its reference polarization implicit.  V124 writes that
polarization in raw tangential/normal trace variables on a product collar.
The physical relative reference fixes tangential values and releases the
normal value with Neumann conormal; the auxiliary absolute reference does
the converse.  This determines the reference determinant exactly.

In the frozen half-space model, the normalized one-form matching matrix is
a scalar multiple of the four-dimensional form fibre.  Its scalar is
identical to the V123 scalar gluing factor.  Hence the exact longitudinal
channel is visible directly, without replacing the one-form determinant by
a scalar ansatz.  The high-frequency limit of the matching matrix is
\(2I\).  Dividing by this universal local gluing normalization leaves an
identity plus a smoothing operator and therefore a genuine Fredholm
determinant after the harmonic threshold is removed.

\subsection{Raw product Green variables}
\label{subsec:V124-Green-variables}

Write a one-form on the physical product collar as
\begin{equation}
 a=\alpha+\phi\,dr,
 \qquad
 q_T:=\gamma\alpha,
 \qquad
 q_N:=\gamma\phi.
 \label{eq:V124-raw-q}
\end{equation}
With outward normal \(-\partial_r\), define component conormals
\begin{equation}
 p_T:=-\gamma(\partial_r\alpha),
 \qquad
 p_N:=-\gamma(\partial_r\phi).
 \label{eq:V124-raw-p}
\end{equation}
In a product collar the Hodge Laplacian is diagonal in this splitting, and
its Green form is
\begin{equation}
 \mathfrak G((q,p),(\widetilde q,\widetilde p))
 =
 \langle p_T,\widetilde q_T\rangle
 +\langle p_N,\widetilde q_N\rangle
 -\langle q_T,\widetilde p_T\rangle
 -\langle q_N,\widetilde p_N\rangle.
 \label{eq:V124-Green-form}
\end{equation}

\begin{lemma}[Relative conditions in raw variables]
\label{lem:V124-relative-raw}
At the physical cut, the relative Hodge conditions
\eqref{eq:V123-relative-BC} are equivalent to
\begin{equation}
 \boxed{
 q_T=0,
 \qquad
 p_N=0.}
 \label{eq:V124-relative-raw}
\end{equation}
\end{lemma}

\begin{proof}
The row \(\nu^\flat\wedge a=0\) kills the tangential trace \(q_T\).
In the product convention,
\[
 \delta a=\delta_\Sigma\alpha-\partial_r\phi.
\]
The tangential derivative of the zero trace \(q_T\) is zero, so
\(\gamma(\delta_\Sigma\alpha)=0\).  The second relative row
\(\gamma(\delta a)=0\) therefore becomes
\(-\gamma(\partial_r\phi)=p_N=0\).
\end{proof}

On the auxiliary collar write
\begin{equation}
 \widetilde a=\beta+\psi\,ds,
 \qquad
 \widetilde q_T:=\gamma\beta,
 \qquad
 \widetilde q_N:=\gamma\psi,
 \label{eq:V124-aux-q}
\end{equation}
with conormals at \(s=0\) defined using its outward normal
\(-\partial_s\).

\begin{lemma}[Absolute conditions in raw variables]
\label{lem:V124-absolute-raw}
At the auxiliary cut, the absolute Hodge conditions
\eqref{eq:V123-absolute-BC} are equivalent to
\begin{equation}
 \boxed{
 \widetilde p_T=0,
 \qquad
 \widetilde q_N=0.}
 \label{eq:V124-absolute-raw}
\end{equation}
\end{lemma}

\begin{proof}
The condition \(\iota_\nu\widetilde a=0\) kills the normal coefficient
\(\widetilde q_N\).  Since
\[
 d\widetilde a
 =
 d_\Sigma\beta
 +ds\wedge(\partial_s\beta-d_\Sigma\psi),
\]
the row \(\iota_\nu d\widetilde a=0\) is
\(\partial_s\beta-d_\Sigma\psi=0\) at the cut.  The tangential derivative
of the zero trace \(\widetilde q_N\) vanishes, so this becomes
\(\partial_s\beta=0\), or \(\widetilde p_T=0\).
\end{proof}

\begin{corollary}[Complementary Lagrangian polarizations]
\label{cor:V124-complementary-polarizations}
The physical relative reference releases \(q_N\) while fixing \(q_T\);
the auxiliary absolute reference releases \(\widetilde q_T\) while fixing
\(\widetilde q_N\).  Each row pair in
\eqref{eq:V124-relative-raw} and
\eqref{eq:V124-absolute-raw} is Lagrangian for the product Green form.
\end{corollary}

\begin{proof}
On either subspace every term of
\eqref{eq:V124-Green-form} vanishes.  Each subspace has one condition per
configuration component, so it is maximal isotropic.
\end{proof}

\subsection{The exact reference determinant}
\label{subsec:V124-reference-determinant}

At a common finite boundary regulator and \(z>0\), let the physical raw
Dirichlet-to-Neumann matrix be
\begin{equation}
 N_{M,1}(z)
 =
 \begin{pmatrix}
 A(z)&B(z)\\
 B(z)^\dagger&D(z)
 \end{pmatrix}
 :
 \binom{q_T}{q_N}
 \longmapsto
 \binom{p_T}{p_N}.
 \label{eq:V124-physical-DtN}
\end{equation}
Positivity of the physical Dirichlet problem gives \(N_{M,1}(z)>0\), hence
\(D(z)>0\).  The auxiliary raw response is
\begin{equation}
 S_{X,1}(z)
 =
 \begin{pmatrix}
 S_{R,1}^{N}(z)&0\\
 0&S_{R,0}^{D}(z)
 \end{pmatrix}.
 \label{eq:V124-aux-DtN}
\end{equation}

\begin{theorem}[Explicit normalized one-form matching matrix]
\label{thm:V124-oneform-W}
Define
\begin{align}
 F_{1,R}(z)
 &:=
 N_{M,1}(z)+S_{X,1}(z),
 \label{eq:V124-F1}\\
 R_{1,R}(z)
 &:=
 \begin{pmatrix}
 S_{R,1}^{N}(z)&0\\
 0&D(z)
 \end{pmatrix}.
 \label{eq:V124-R1}
\end{align}
Then the one-form relative determinant in
\eqref{eq:V123-oneform-relative-det} is
\begin{equation}
 \boxed{
 \det\mathcal W_{1,R}(z)
 =
 \frac{\det F_{1,R}(z)}
 {\det S_{R,1}^{N}(z)\det D(z)},}
 \label{eq:V124-det-W1}
\end{equation}
and one may take
\begin{equation}
 \boxed{
 \mathcal W_{1,R}(z)
 =
 R_{1,R}(z)^{-1/2}
 F_{1,R}(z)
 R_{1,R}(z)^{-1/2}.}
 \label{eq:V124-W1}
\end{equation}
\end{theorem}

\begin{proof}
Begin with the all-coefficient-Dirichlet interior determinants on both
pieces.  For the physical relative reference,
Lemma~\ref{lem:V124-relative-raw} keeps \(q_T=0\) and releases \(q_N\)
with \(p_N=0\).  The boundary Hessian for that released variable is the
normal block \(D(z)\), so its Gaussian Schur factor is \(\det D(z)\).

For the auxiliary absolute reference,
Lemma~\ref{lem:V124-absolute-raw} releases
\(\widetilde q_T\) with \(\widetilde p_T=0\) and keeps
\(\widetilde q_N=0\).  Its boundary Hessian is
\(S_{R,1}^{N}(z)\), so its factor is
\(\det S_{R,1}^{N}(z)\).

In the glued problem all common interface values are released.  Smooth
gluing identifies tangential values directly and identifies normal values
after the orientation sign described in
Proposition~\ref{prop:V122-oneform-Schur}.  The corresponding conormal sum
is the full Hessian \(F_{1,R}=N_{M,1}+S_{X,1}\).  The common interior
determinants cancel in the ratio, proving
\eqref{eq:V124-det-W1}.  Congruence by \(R_{1,R}^{-1/2}\) has precisely the
same determinant ratio and proves \eqref{eq:V124-W1}.
\end{proof}

\begin{corollary}[Off-diagonal shape data are retained]
\label{cor:V124-offdiagonal-retained}
The entries \(B(z)\) and \(B(z)^\dagger\) remain inside
\(\det F_{1,R}(z)\).  They do not appear in the reference denominator and
therefore cannot be discarded before the determinant is formed.
\end{corollary}

\begin{proof}
This is immediate from
\eqref{eq:V124-F1}--\eqref{eq:V124-det-W1}.  In a curved collar their
leading extrinsic values are the off-diagonal entries computed in
\eqref{eq:V118-oneform-subprincipal}.
\end{proof}

\subsection{Frozen half-space evaluation}
\label{subsec:V124-frozen-evaluation}

Freeze a flat product interface and a tangential frequency \(\xi\).  Put
\begin{equation}
 \lambda:=\sqrt{|\xi|^2+z},
 \qquad
 a_R:=\lambda\tanh(R\lambda),
 \qquad
 d_R:=\lambda\coth(R\lambda).
 \label{eq:V124-lambda-a-d}
\end{equation}
The physical stable Dirichlet-to-Neumann matrix and auxiliary response are
\begin{equation}
 N_{M,1}^{\rm fr}=\lambda I_4,
 \qquad
 S_{X,1}^{\rm fr}
 =
 \operatorname{diag}(a_RI_3,d_R).
 \label{eq:V124-frozen-responses}
\end{equation}

\begin{theorem}[Exact frozen one-form matching matrix]
\label{thm:V124-frozen-W1}
The reference and full Schur matrices are
\begin{align}
 R_{1,R}^{\rm fr}
 &=
 \operatorname{diag}(a_RI_3,\lambda),
 \label{eq:V124-frozen-R}\\
 F_{1,R}^{\rm fr}
 &=
 \operatorname{diag}((\lambda+a_R)I_3,\lambda+d_R).
 \label{eq:V124-frozen-F}
\end{align}
Consequently,
\begin{equation}
 \boxed{
 \mathcal W_{1,R}^{\rm fr}(z,\xi)
 =
 w_R(\lambda)I_4,
 \qquad
 w_R(\lambda)
 :=
 1+\coth(R\lambda)
 =
 \frac{2}{1-e^{-2R\lambda}}.}
 \label{eq:V124-frozen-W1}
\end{equation}
\end{theorem}

\begin{proof}
The tangential normalized entry is
\[
 \frac{\lambda+a_R}{a_R}
 =
 1+\coth(R\lambda).
\]
The normal normalized entry is
\[
 \frac{\lambda+d_R}{\lambda}
 =
 1+\coth(R\lambda).
\]
Thus all four entries coincide.  The exponential form is the elementary
identity for \(1+\coth\).
\end{proof}

\begin{corollary}[Direct scalar channel match]
\label{cor:V124-scalar-match}
The frozen scalar normalized matching operator of
\eqref{eq:V123-scalar-W} is
\begin{equation}
 \mathcal W_{0,R}^{\rm fr}(z,\xi)
 =
 \frac{\lambda+a_R}{a_R}
 =
 w_R(\lambda).
 \label{eq:V124-frozen-W0}
\end{equation}
Hence every exact longitudinal scalar mode occurs as an identical
one-form interface eigenchannel.  In particular,
\begin{equation}
 \det\mathcal W_{1,R}^{\rm fr}
 =
 \bigl(\det\mathcal W_{0,R}^{\rm fr}\bigr)^4
 \label{eq:V124-det-rank-law}
\end{equation}
at the frozen rank level.
\end{corollary}

\begin{proof}
The scalar physical response is \(\lambda\), its auxiliary absolute
reference response is \(a_R\), and its glued response is
\(\lambda+a_R\).  This proves \eqref{eq:V124-frozen-W0}.  Apply
Theorem~\ref{thm:V124-frozen-W1}.
\end{proof}

\begin{corollary}[Frozen BRST interface factor]
\label{cor:V124-frozen-BRST-factor}
Per Standard Model Lie-algebra generator, the frozen determinant
combination is
\begin{equation}
 \frac{\det\mathcal W_{0,R}^{\rm fr}}
 {(\det\mathcal W_{1,R}^{\rm fr})^{1/2}}
 =
 w_R(\lambda)^{-1}.
 \label{eq:V124-frozen-BRST-factor}
\end{equation}
One scalar channel cancels one half of the exact one-form channel; the
remaining power represents the coexact rank contribution of the
gauge/ghost determinant card.
\end{corollary}

\begin{proof}
Use \eqref{eq:V124-det-rank-law}:
\(w_R/(w_R^4)^{1/2}=w_R^{-1}\).
\end{proof}

\subsection{The universal factor two and the trace-class remainder}
\label{subsec:V124-factor-two}

\begin{proposition}[Uniform normalized interface floor]
\label{prop:V124-W-floor}
For \(\operatorname{Re}\lambda>0\),
\begin{equation}
 |w_R(\lambda)|\ge1.
 \label{eq:V124-W-floor}
\end{equation}
Moreover,
\begin{equation}
 w_R(\lambda)\longrightarrow2
 \qquad
 (|\lambda|\to\infty\text{ in a closed right half-sector}).
 \label{eq:V124-W-high-limit}
\end{equation}
\end{proposition}

\begin{proof}
By \eqref{eq:V124-frozen-W1},
\[
 |w_R(\lambda)|
 =
 \frac{2}{|1-e^{-2R\lambda}|}
 \ge1,
\]
because \(|1-e^{-2R\lambda}|\le2\).  Exponential decay of
\(e^{-2R\lambda}\) in a closed right half-sector gives the limit.
\end{proof}

Define the locally normalized interface factor
\begin{equation}
 \widetilde w_R(\lambda)
 :=
 \frac12w_R(\lambda)
 =
 (1-e^{-2R\lambda})^{-1}.
 \label{eq:V124-wtilde}
\end{equation}

\begin{theorem}[The finite-length interface remainder is trace class]
\label{thm:V124-trace-class}
On the positive spectral complement of a compact boundary,
\begin{equation}
 \widetilde{\mathcal W}_{0,R}^{\rm fr}
 :=
 \frac12\mathcal W_{0,R}^{\rm fr}
 =
 (I-e^{-2R\Lambda(z)})^{-1}
 \label{eq:V124-Wtilde-operator}
\end{equation}
satisfies
\begin{equation}
 \widetilde{\mathcal W}_{0,R}^{\rm fr}-I
 =
 e^{-2R\Lambda(z)}
 (I-e^{-2R\Lambda(z)})^{-1}
 \in\mathfrak S_1.
 \label{eq:V124-trace-class-difference}
\end{equation}
The same holds for the frozen one-form operator.  Its Fredholm determinant
is
\begin{equation}
 \boxed{
 \log\det_{\!F}\widetilde{\mathcal W}_{0,R}^{\rm fr}
 =
 -\sum_{\mu_j>0}
 \log\left(1-e^{-2R\sqrt{\mu_j^2+z}}\right).}
 \label{eq:V124-Fredholm-det}
\end{equation}
\end{theorem}

\begin{proof}
The eigenvalues of \(\Lambda(z)\) grow like
\(c\,j^{1/\dim\Sigma}\).  Hence
\(e^{-2R\Lambda(z)}\) is trace class, indeed smoothing.  On the positive
complement its norm is strictly below one for real \(z\ge0\), so the
geometric resolvent in
\eqref{eq:V124-trace-class-difference} is bounded.  A bounded operator
times a trace-class operator is trace class.  The Fredholm determinant is
the product of the eigenvalues
\((1-e^{-2R\sqrt{\mu_j^2+z}})^{-1}\); taking logarithms gives the sum.
\end{proof}

\begin{corollary}[The factor \(2\) is a local gluing normalization]
\label{cor:V124-factor-two-local}
At a boundary regulator of rank \(N_\partial\),
\begin{equation}
 \det\mathcal W_{0,R}^{\rm fr}
 =
 2^{N_\partial}
 \det\widetilde{\mathcal W}_{0,R}^{\rm fr}.
 \label{eq:V124-factor-two-det}
\end{equation}
The first factor diverges with boundary mode count and is independent of
\(R\), curvature, and spectral scale at principal order.  It belongs to
the local interface measure normalization.  The second factor is the
finite-length, trace-class collar response.
\end{corollary}

\begin{proof}
Equation \eqref{eq:V124-wtilde} gives one factor \(2\) per regulated
boundary mode.  The remaining determinant exists by
Theorem~\ref{thm:V124-trace-class}.
\end{proof}

\begin{firewall}[The local factor may not be silently deleted]
\label{fire:V124-factor-two-firewall}
Dividing by \(2I\) is a declared local gluing normalization.  In the
continuum determinant its regulated power must be absorbed into a boundary
measure or local counterterm using the same scalar/one-form regulator.
Deleting \(2^{N_\partial}\) independently in different form degrees would
break the relative BRST rank accounting.
\end{firewall}

\subsection{Harmonic threshold}
\label{subsec:V124-harmonic-threshold}

\begin{proposition}[Matched scalar and normal thresholds]
\label{prop:V124-harmonic-threshold}
For \(\mu=0\) and \(z\downarrow0\),
\begin{align}
 a_R&=Rz+O(z^2),
 &
 d_R&=R^{-1}+\frac R3z+O(z^2),
 \label{eq:V124-ad-threshold}\\
 w_R(\sqrt z)
 &=
 \frac1{R\sqrt z}+1+O(\sqrt z).
 \label{eq:V124-w-threshold}
\end{align}
The scalar, tangential-longitudinal, and normal normalized matching factors
have the same threshold singularity.  It is removed only after the exact
global-gauge prime and finite Hodge-volume quotient are applied.
\end{proposition}

\begin{proof}
Use the Taylor series
\(\tanh x=x-x^3/3+O(x^5)\) and
\(\coth x=x^{-1}+x/3+O(x^3)\), then apply
\eqref{eq:V124-frozen-W1}.
\end{proof}

\subsection{Curved matrix status}
\label{subsec:V124-curved-status}

\begin{assessment}[What is exact and what remains]
\label{ass:V124-product-status}
Equations \eqref{eq:V124-W1} and \eqref{eq:V124-det-W1} are exact for a
product collar with a general positive physical one-form
Dirichlet-to-Neumann matrix.  The scalar identity
\(\mathcal W_{1,R}^{\rm fr}=w_RI_4\) is the complete frozen principal model.
In a nonproduct collar, the V118 shape matrix perturbs the numerator and
the physical normal reference block \(D(z)\).  Its determinant contribution
must be computed after a jet-matched auxiliary continuation is chosen.
\end{assessment}

\begin{firewall}[Off-diagonal terms enter at second determinant order]
\label{fire:V124-offdiagonal-order}
At first variation around the diagonal frozen model, an off-diagonal shape
block has zero matrix trace; at second variation its product with the
adjoint contributes to \(\log\det\mathcal W_1\).  Therefore the absence of
an off-diagonal term in the first trace does not permit its removal from
the determinant or from the heat coefficient at the next order.
\end{firewall}

\subsection{V124 ledger and next acquisition}
\label{subsec:V124-ledger}

\paragraph{New exact conclusions.}
\begin{enumerate}[label=\textup{(\arabic*)},leftmargin=2.8em]
\item Lemmas~\ref{lem:V124-relative-raw} and
      \ref{lem:V124-absolute-raw} print the complementary physical and
      auxiliary Hodge polarizations.
\item Theorem~\ref{thm:V124-oneform-W} gives the missing normalized
      one-form matching matrix and its exact reference determinant.
\item Theorem~\ref{thm:V124-frozen-W1} proves that the flat matching matrix
      is \(w_RI_4\).
\item Corollary~\ref{cor:V124-scalar-match} verifies the scalar/exact-form
      interface channel directly.
\item Proposition~\ref{prop:V124-W-floor} proves the normalized matching
      floor and isolates the universal high-frequency factor \(2\).
\item Theorem~\ref{thm:V124-trace-class} proves that division by the common
      local factor leaves an identity plus a trace-class collar response.
\item Proposition~\ref{prop:V124-harmonic-threshold} matches the scalar and
      normal global-gauge thresholds.
\end{enumerate}

\paragraph{Next forced step.}
V125 must begin with the scheduled NS/YM companion audit.  It should then
choose a nonproduct auxiliary metric and connection whose interface jets
match the physical collar, insert the V118 order-zero shape matrix into
\eqref{eq:V124-W1}, and compute the first curved variation of the normalized
interface determinant.  The factor \(2\), the harmonic prime, and the
Gauss--Manin/Hodge density must remain common throughout.

\begin{assessment}[V124 continuum status]
\label{ass:V124-continuum-status}
The full Standard--Model--Einstein continuum remains unproved.  V124 closes
the product one-form interface polarization and separates its local gluing
normalization from a trace-class finite-collar response.  Curved jet
matching and heat coefficients, fermion extension, nonlinear interacting
existence, Einstein stress closure, and the cofinal continuum remain open.
\end{assessment}

\section*{V124 exact checks and append-only record}
\addcontentsline{toc}{section}{V124 exact checks and append-only record}

The one-form polarization and trace-class checksum is
\begin{equation}
 \boxed{
 \mathcal W_{1,R}^{\rm fr}
 =
 \left(1+\coth(R\lambda)\right)I_4,
 \qquad
 \frac12\mathcal W_{1,R}^{\rm fr}-I
 \in\mathfrak S_1.}
 \label{eq:V124-final-check}
\end{equation}
The append operation preserves the complete V123 prefix byte for byte before
its former live terminator.  The assembled V124 file is checked for one live
final terminator, balanced environments and braces, duplicate labels, newly
introduced unresolved references, display math delimiters, and control
characters before the exact V123 predecessor is removed.

% =============================================================================
% END APPENDED VERSION V124 MATERIAL
% =============================================================================

\clearpage
% =============================================================================
% BEGIN APPENDED VERSION V125 MATERIAL
% =============================================================================
\section{V125: companion audit, jet-matched curved gluing, and the first
BRST interface curvature symbol}
\label{sec:V125-curved-interface}

V124 computed the complete product boundary polarization and showed that
the normalized flat matching operator has principal value \(2I\).
V125 performs the scheduled companion audit, constructs a smooth
nonproduct auxiliary continuation, and computes the first extrinsic
curvature term of the normalized scalar/one-form determinant.

Smooth jet matching makes the physical and auxiliary order-zero
Dirichlet-to-Neumann curvature matrices exact negatives.  They therefore
cancel in the glued conormal sum, as they must at an artificial cut inside
a smooth manifold.  The relative/absolute reference is asymmetric,
however, and its normalization retains one local interface term.  In the
BRST scalar-minus-one-half-one-form combination, all directional
second-fundamental-form dependence cancels at this order and only the mean
curvature remains.

\subsection{Scheduled NS/YM companion audit}
\label{subsec:V125-audit}

The read-only audit on 5 September 2026 found
\begin{align}
 &\texttt{Renewal\_NS\_Stability\_Research\_Notes\_V31.tex},
 \notag\\[-1mm]
 &\hspace{8mm}
 887537\ {\rm bytes},\qquad
 \operatorname{SHA256}
 =
 \texttt{F1121B62238AD5491BB7CF03635EA9934942EDF573ED8FAAA9EE79E1D38A6696},
 \label{eq:V125-NS-audit}\\
 &\texttt{Renewal\_Yang\_Mills\_Mass\_Gap\_Research\_Notes\_V110.tex},
 \notag\\[-1mm]
 &\hspace{8mm}
 2634683\ {\rm bytes},\qquad
 \operatorname{SHA256}
 =
 \texttt{C4EB1444549FCEEA9B1800914087C97808B1552969F1C991047C0DC4EA7F94A2}.
 \label{eq:V125-YM-audit}
\end{align}
The YM filesystem name ended in V110, while its sole live tail contained a
complete internally labelled V111 appendix.  No companion file was
modified.

The NS conclusion is unchanged: a restricted correlated object must be
formed before it is replaced by a positive global bound.  YM V110 keeps two
signed covariance rows together until their common positive effect is
formed.  Its internally appended V111 then proves the exact upstream
Douglas criterion
\[
 K^\dagger K\le\lambda T
 \quad\Longleftrightarrow\quad
 K=KT_{\rm supp},\qquad
 \lambda_{\min}=\|KT^{\dagger/2}\|^2,
\]
and exhibits divergence when an apparently small response occupies a
small eigenvalue or leaks outside the support of \(T\).

\begin{assessment}[V125 audit transfer]
\label{ass:V125-audit-transfer}
For the SM interface determinant, the reference weight is the matrix
\(R_{1,R}\) in \eqref{eq:V124-R1}.  V125 therefore keeps the normalized
operator
\[
 R_{1,R}^{-1/2}F_{1,R}R_{1,R}^{-1/2}
\]
intact.  A trace estimate for \(F_{1,R}-R_{1,R}\) does not control this
operator near a harmonic threshold.  The exact support condition and the
inverse reference weight are tested before any determinant expansion.
No NS or YM estimate is imported.
\end{assessment}

\subsection{Smooth continuation across the artificial cut}
\label{subsec:V125-jet-matching}

Let the physical collar be
\begin{equation}
 g_M=dr^2+\gamma_r,
 \qquad r\ge0,
 \label{eq:V125-physical-collar}
\end{equation}
and let \(t\) be a signed normal coordinate, positive on \(M\).  On the
auxiliary side put \(s=-t\ge0\).

\begin{theorem}[Existence of a jet-matched auxiliary collar]
\label{thm:V125-jet-matched-extension}
For a smooth physical metric and compatible bundle connection near
\(\Sigma\), there is \(R_0>0\) and a smooth auxiliary metric and connection
on \([0,R_0]_s\times\Sigma\) such that the glued data are smooth in \(t\).
In a common boundary trivialization their jets obey
\begin{align}
 \partial_s^m\gamma_X(0)
 &=
 (-1)^m\partial_r^m\gamma_M(0),
 \label{eq:V125-metric-jets}\\
 \partial_s^m\nabla^X(0)
 &=
 (-1)^m\partial_r^m\nabla^M(0)
 \label{eq:V125-connection-jets}
\end{align}
for every \(m\ge0\).  The collar may be continued to arbitrary finite
length and made product near its far end without changing these interface
jets.
\end{theorem}

\begin{proof}
Extend the coefficient functions of \(\gamma_r\) and of the connection
smoothly from \(r\ge0\) to a two-sided interval in the signed coordinate
\(t\).  Smooth extension of half-line jets is obtained coefficientwise by
the standard Borel extension construction.  Positive definiteness is open,
so after shrinking the interval the extended symmetric tensor remains a
metric.  Define the auxiliary data by restriction to \(t=-s\).  The chain
rule gives \eqref{eq:V125-metric-jets}--%
\eqref{eq:V125-connection-jets}.  A cutoff supported away from \(s=0\)
interpolates to constant product data near the far end; compactness and a
convex interpolation inside the positive metric cone preserve positive
definiteness.
\end{proof}

\begin{corollary}[Opposite shape data on the two inward collars]
\label{cor:V125-opposite-shape}
With both \(r\) and \(s\) increasing into their respective pieces from the
cut,
\begin{equation}
 K_X=-K_M,
 \qquad
 H_X=-H_M,
 \qquad
 S_X=-S_M
 \label{eq:V125-opposite-K}
\end{equation}
at \(\Sigma\).
\end{corollary}

\begin{proof}
Take \(m=1\) in \eqref{eq:V125-metric-jets} and use
\(K=\tfrac12\partial_{\rm inward}\gamma\).  Trace with the common
interface metric and raise one index.
\end{proof}

\begin{proposition}[The enlarged curved complex is fixed and off-shell]
\label{prop:V125-curved-complex}
After the jet-matched gluing, the scalar and Hodge one-form Laplacians on
the enlarged compact manifold, with absolute conditions at the far
boundary, are fixed nonnegative self-adjoint operators and satisfy the
off-shell de Rham and heat identities
\eqref{eq:V122-offshell-chain}--\eqref{eq:V122-heat-chain}.
\end{proposition}

\begin{proof}
The interface is now a smooth interior hypersurface of the enlarged
manifold.  The proof of Theorem~\ref{thm:V122-fixed-deRham} applies without
a transmission singularity.  The far-boundary absolute conditions remain
unchanged.
\end{proof}

\subsection{Cancellation in the glued conormal sum}
\label{subsec:V125-sum-cancellation}

Fix \(y\in\Sigma\), \(\xi\ne0\), and use the notation
\begin{equation}
 \rho=|\xi|_\gamma,
 \qquad
 n=\xi/\rho,
 \qquad
 k=K_M(n,n),
 \qquad
 s_0=\frac12(k-H).
 \label{eq:V125-rho-k-s0}
\end{equation}
Let \(N_{M,p}\) and \(N_{X,p}\) be the outgoing
Dirichlet-to-Neumann maps of the two pieces.  Finite collar length changes
their classical symbols only by smoothing terms, as in
Proposition~\ref{prop:V122-smoothing}.

\begin{theorem}[Extrinsic subprincipal cancellation at a smooth cut]
\label{thm:V125-subprincipal-cancellation}
In a geometric symbol frame at the cut,
\begin{align}
 N_{M,0}
 &=
 \rho+s_0+O(\rho^{-1}),
 &
 N_{X,0}
 &=
 \rho-s_0+O(\rho^{-1}),
 \label{eq:V125-scalar-two-sides}\\
 N_{M,1}
 &=
 \rho I_4+M_K(n)+O(\rho^{-1}),
 &
 N_{X,1}
 &=
 \rho I_4-M_K(n)+O(\rho^{-1}),
 \label{eq:V125-oneform-two-sides}
\end{align}
where
\begin{equation}
 M_K(n)
 =
 \begin{pmatrix}
 s_0I_3+S&-iSn\\
 i(Sn)^\dagger&s_0
 \end{pmatrix}
 \label{eq:V125-MK}
\end{equation}
is the V118 extrinsic matrix.  Hence
\begin{align}
 N_{M,0}+N_{X,0}
 &=
 2\rho+O(\rho^{-1}),
 \label{eq:V125-scalar-sum}\\
 N_{M,1}+N_{X,1}
 &=
 2\rho I_4+O(\rho^{-1})
 \label{eq:V125-oneform-sum}
\end{align}
with no extrinsic order-zero term.
\end{theorem}

\begin{proof}
The physical scalar row is
\eqref{eq:V117-raw-subprincipal}.  The physical one-form row is
\eqref{eq:V118-oneform-subprincipal}.  Every displayed extrinsic
order-zero entry is linear in \(K,H,S\).  Apply
\eqref{eq:V125-opposite-K} on the auxiliary side.  The far-end correction
is smoothing and therefore does not enter the classical homogeneous
symbols.  Adding the two sides proves the last two identities.
\end{proof}

\begin{assessment}[Geometric meaning of the cancellation]
\label{ass:V125-smooth-cut-meaning}
The artificial interface is not a geometric boundary of the glued
manifold.  Its two outward conormal responses see opposite extrinsic
curvature.  The cancellation in
\eqref{eq:V125-scalar-sum}--\eqref{eq:V125-oneform-sum} is therefore the
symbolic statement that the glued numerator contains no spurious
one-sided mean-curvature counterterm.
\end{assessment}

\subsection{First curvature symbol of the normalized determinants}
\label{subsec:V125-log-symbol}

The relative reference is asymmetric.  For scalars its boundary
normalization is the auxiliary Neumann response \(N_{X,0}\).  For
one-forms, Theorem~\ref{thm:V124-oneform-W} normalizes by the auxiliary
tangential response and the physical normal response.  At classical
symbol level,
\begin{equation}
 R_{1}^{\rm ref}
 =
 \begin{pmatrix}
 \rho I_3-(s_0I_3+S)&0\\
 0&\rho+s_0
 \end{pmatrix}
 +O(\rho^{-1}).
 \label{eq:V125-reference-symbol}
\end{equation}

\begin{lemma}[First fibre-logarithm row]
\label{lem:V125-log-row}
If \(C\) is a fixed fibre endomorphism, then
\begin{align}
 \operatorname{tr}\log(\rho I+C)
 &=
 r\log\rho+\rho^{-1}\operatorname{tr}C+O(\rho^{-2}),
 \label{eq:V125-log-rho}\\
 \operatorname{tr}\log(2\rho I+C)
 &=
 r\log(2\rho)+(2\rho)^{-1}\operatorname{tr}C+O(\rho^{-2}).
 \label{eq:V125-log-2rho}
\end{align}
\end{lemma}

\begin{proof}
Factor out \(\rho I\) or \(2\rho I\) and use
\(\log(I+X)=X+O(X^2)\).
\end{proof}

\begin{theorem}[Scalar and one-form relative log symbols]
\label{thm:V125-relative-log-symbols}
After the common principal factor \(2\) is separated as in V124, the first
extrinsic homogeneous symbols of the relative interface logarithms are
\begin{align}
 \sigma_{-1}^{\rm ext}
 \bigl(\log\mathcal W_{0,R}\bigr)
 &=
 \frac{s_0}{\rho}
 =
 \frac{k-H}{2\rho},
 \label{eq:V125-scalar-log-symbol}\\
 \sigma_{-1}^{\rm ext}
 \bigl(\operatorname{tr}\log\mathcal W_{1,R}\bigr)
 &=
 \frac{k}{\rho}.
 \label{eq:V125-oneform-log-symbol}
\end{align}
\end{theorem}

\begin{proof}
For the scalar ratio,
\[
 \log\det\mathcal W_0
 =
 \log\det(N_{M,0}+N_{X,0})
 -\log\det N_{X,0}.
\]
The numerator has no extrinsic order-zero symbol by
\eqref{eq:V125-scalar-sum}.  The auxiliary denominator has order-zero
symbol \(-s_0\).  Lemma~\ref{lem:V125-log-row} therefore gives
\(+s_0/\rho\).

For one-forms, \eqref{eq:V124-det-W1} gives
\[
 \log\det\mathcal W_1
 =
 \log\det(N_{M,1}+N_{X,1})
 -\log\det(N_{X,1})_{TT}
 -\log\det(N_{M,1})_{NN}.
\]
The numerator again has no extrinsic order-zero symbol.  The tangential
auxiliary block has correction \(-(s_0I_3+S)\), while the physical normal
block has correction \(s_0\).  Thus the order-\(-1\) row is
\[
 \frac{\operatorname{tr}(s_0I_3+S)-s_0}{\rho}
 =
 \frac{(3s_0+H)-s_0}{\rho}
 =
 \frac{2s_0+H}{\rho}
 =
 \frac{k}{\rho}.
\]
\end{proof}

\begin{corollary}[First BRST interface curvature symbol]
\label{cor:V125-BRST-log-symbol}
Per Lie-algebra generator,
\begin{equation}
 \boxed{
 \sigma_{-1}^{\rm ext}
 \left[
 \log\mathcal W_{0,R}
 -\frac12\operatorname{tr}\log\mathcal W_{1,R}
 \right]
 =
 -\frac{H}{2\rho}.}
 \label{eq:V125-BRST-log-symbol}
\end{equation}
After multiplication by
\(\dim\mathfrak g_{\rm SM}=12\), the symbol is
\begin{equation}
 -\frac{6H}{\rho}.
 \label{eq:V125-SM-log-symbol}
\end{equation}
The directional curvature \(k=K(n,n)\) cancels completely at this order.
\end{corollary}

\begin{proof}
Insert
\eqref{eq:V125-scalar-log-symbol} and
\eqref{eq:V125-oneform-log-symbol}:
\[
 \frac{k-H}{2\rho}-\frac{k}{2\rho}
 =
 -\frac{H}{2\rho}.
\]
Multiply by twelve.
\end{proof}

\begin{corollary}[Sharp radial-cutoff divergence]
\label{cor:V125-H-divergence}
With a common sharp boundary covector cutoff \(|\xi|\le\Lambda\), the
order-\(-1\) Standard Model adjoint interface symbol contributes
\begin{equation}
 \boxed{
 -\frac{3\Lambda^2}{2\pi^2}
 \int_\Sigma H\,d{\rm vol}_\gamma}
 \label{eq:V125-H-divergence}
\end{equation}
to the regulated logarithmic interface determinant, up to the overall sign
convention used when converting \(\log Z\) to an effective action.
\end{corollary}

\begin{proof}
Use \eqref{eq:V125-SM-log-symbol} and
\[
 \frac1{(2\pi)^3}
 \int_{|\xi|\le\Lambda}\frac{d^3\xi}{|\xi|}
 =
 \frac{\Lambda^2}{4\pi^2}.
\]
Multiplication by \(-6H\) gives the result.
\end{proof}

\begin{firewall}[Interface term versus physical Einstein boundary term]
\label{fire:V125-interface-versus-physical}
The density in \eqref{eq:V125-H-divergence} belongs to the chosen
glued-versus-relative/absolute interface determinant.  The smooth glued
numerator itself has no order-zero extrinsic term.  A physical Einstein
boundary coefficient is obtained only after this interface contribution is
combined with the physical relative baseline and the declared
Gibbons--Hawking--York normalization.  V125 does not identify the displayed
coefficient with the final renormalized Newton or boundary coupling.
\end{firewall}

\subsection{The harmonic support gate}
\label{subsec:V125-harmonic-support}

Let \(R(z)\ge0\) be a finite reference matrix with
\begin{equation}
 R(z)|_{\operatorname{Ran}P}
 =
 zC+O(z^2),
 \qquad
 C>0,
 \qquad
 R(0)|_{\operatorname{Ran}Q}>0,
 \label{eq:V125-reference-threshold}
\end{equation}
where \(P\) is its threshold kernel and \(Q=I-P\).

\begin{proposition}[Exact support condition for normalized boundedness]
\label{prop:V125-support-gate}
Let \(F(z)\) be continuous at \(z=0\).  A necessary condition for
\begin{equation}
 R(z)^{-1/2}F(z)R(z)^{-1/2}
 \label{eq:V125-normalized-threshold}
\end{equation}
to remain bounded as \(z\downarrow0\) is
\begin{equation}
 \boxed{
 PF(0)P=0,
 \qquad
 QF(0)P=0,
 \qquad
 PF(0)Q=0.}
 \label{eq:V125-support-condition}
\end{equation}
If \eqref{eq:V125-support-condition} holds and
\(F(z)P=O(z)\), \(PF(z)=O(z)\), then the normalized family is bounded.
\end{proposition}

\begin{proof}
Relative to \(P\oplus Q\),
\[
 R(z)^{-1/2}
 =
 \begin{pmatrix}
 z^{-1/2}C^{-1/2}+O(z^{1/2})&0\\
 0&R_Q(0)^{-1/2}+O(z)
 \end{pmatrix}.
\]
A nonzero \(PFP\) block produces order \(z^{-1}\); a nonzero \(QFP\) or
\(PFQ\) block produces order \(z^{-1/2}\).  This proves necessity.  Under
the stated vanishing rates, the \(PP\) block is \(O(1)\), the off-diagonal
blocks are \(O(z^{1/2})\), and the \(QQ\) block is bounded.
\end{proof}

\begin{corollary}[Feshbach before the harmonic prime]
\label{cor:V125-Feshbach-before-prime}
At \(z=0\), the determinant expansion of
\(\mathcal W_{1,R}=R_{1,R}^{-1/2}F_{1,R}R_{1,R}^{-1/2}\) is legitimate on
the primed sector only after the V119 Feshbach reconstruction has tested
\eqref{eq:V125-support-condition}.  A curved leakage block such as
\eqref{eq:V119-harmonic-leakage} can otherwise create a
\(z^{-1/2}\) normalized divergence even when its unnormalized norm is
small.
\end{corollary}

\begin{proof}
Apply Proposition~\ref{prop:V125-support-gate} with the threshold projection
onto the harmonic reference kernel.  Equation
\eqref{eq:V119-harmonic-leakage} supplies an explicit possible
off-diagonal support block.  The Feshbach map determines whether it is part
of a reconstructed zero mode or a true coupling to the positive sector.
\end{proof}

\begin{firewall}[The curvature trace is a positive-frequency statement]
\label{fire:V125-positive-frequency}
The symbol calculation
\eqref{eq:V125-BRST-log-symbol} is homogeneous at \(\rho>0\).  It neither
proves nor assumes the harmonic support condition.  The high-frequency
curvature counterterm and the finite harmonic quotient are independent
rows which must both pass before a full determinant is defined.
\end{firewall}

\subsection{V125 ledger and next acquisition}
\label{subsec:V125-ledger}

\paragraph{New exact conclusions.}
\begin{enumerate}[label=\textup{(\arabic*)},leftmargin=2.8em]
\item The scheduled audit records the unchanged NS V31 file and the YM file
      whose live content has advanced internally through V111.
\item Theorem~\ref{thm:V125-jet-matched-extension} constructs a smooth
      auxiliary metric and connection with all interface jets matched.
\item Theorem~\ref{thm:V125-subprincipal-cancellation} proves cancellation
      of the two one-sided extrinsic order-zero matrices in the glued
      conormal sum.
\item Theorem~\ref{thm:V125-relative-log-symbols} computes the scalar and
      one-form relative logarithmic symbols.
\item Corollary~\ref{cor:V125-BRST-log-symbol} proves cancellation of
      directional curvature and leaves the exact mean-curvature symbol
      \(-H/(2|\xi|)\) per gauge generator.
\item Corollary~\ref{cor:V125-H-divergence} prints its common sharp-cutoff
      interface divergence.
\item Proposition~\ref{prop:V125-support-gate} imports the applicable
      support-weighted audit lesson as an exact harmonic normalization
      criterion.
\end{enumerate}

\paragraph{Next forced step.}
V126 should compute the order-\(-2\) symbol of the curved interface
logarithm.  It must retain the off-diagonal V118 shape pair through the
quadratic logarithm, include the order-\(-1\) Riccati symbols on both sides,
and separate the invariants \(H^2\), \(|K|^2\), and intrinsic boundary
curvature.  The calculation must be made on the positive-frequency
Feshbach complement; harmonic support remains a separate finite gate.

\begin{assessment}[V125 continuum status]
\label{ass:V125-continuum-status}
The full Standard--Model--Einstein continuum remains unproved.  V125
constructs curved jet matching and obtains the first BRST interface
curvature symbol while retaining the exact harmonic normalization gate.
The next heat coefficient, fermion extension, interacting nonlinear
existence, Einstein stress closure, and the cofinal continuum remain open.
\end{assessment}

\section*{V125 exact checks and append-only record}
\addcontentsline{toc}{section}{V125 exact checks and append-only record}

The curved gluing and BRST checksum is
\begin{equation}
 \boxed{
 M_{K_X}=-M_{K_M},
 \qquad
 \sigma_{-1}^{\rm ext}
 \left(\log\mathcal W_0-\tfrac12\operatorname{tr}\log\mathcal W_1\right)
 =
 -\frac{H}{2|\xi|}.}
 \label{eq:V125-final-check}
\end{equation}
The append operation preserves the complete V124 prefix byte for byte before
its former live terminator.  The assembled V125 file is checked for one live
final terminator, balanced environments and braces, duplicate labels, newly
introduced unresolved references, display math delimiters, and control
characters before the exact V124 predecessor is removed.

% =============================================================================
% END APPENDED VERSION V125 MATERIAL
% =============================================================================

\clearpage
% =============================================================================
% BEGIN APPENDED VERSION V126 MATERIAL
% =============================================================================
\section{V126: reflection parity and the universal quadratic interface
curvature row}
\label{sec:V126-quadratic-row}

V125 computed the order-\(-1\) logarithmic interface symbol from the
order-zero Dirichlet-to-Neumann matrices.  The next row normally depends on
the order-\(-1\) Riccati coefficient, tangential symbol composition, and
curvature squares.  V126 first isolates a sector in which the unknown
Riccati coefficient cancels.

Freeze the tangential jets at one interface point, retain arbitrary normal
jets, and use the jet-matched reflection of V125.  In this fibrewise sector
the positive conormal root has an expansion
\[
 \rho I\pm M+\rho^{-1}C+\cdots.
\]
Normal reflection changes the sign of \(M\) but not \(C\).  The even
coefficient \(C\) cancels from the normalized relative determinant.
The remaining quadratic expression is therefore determined entirely by
the V117--V118 order-zero matrices.

This calculation is not the full order-\(-2\) pseudodifferential symbol.
Tangential derivatives, Poisson brackets, intrinsic boundary curvature,
and bundle curvature occupy separate rows and remain open.

\subsection{Frozen Riccati parity}
\label{subsec:V126-Riccati-parity}

Let a tangential spectral channel of a fixed-measure collar operator have
\begin{equation}
 P=-\partial_r^2+\mathcal A(r),
 \qquad
 \mathcal A(r)
 =
 \rho(r)^2I+\mathcal A_1(r)+\mathcal A_0(r)+\cdots,
 \label{eq:V126-frozen-operator}
\end{equation}
where tangential composition is ordinary matrix multiplication in the
frozen sector.  Let its positive root solve
\begin{equation}
 \partial_rL+L^2=\mathcal A.
 \label{eq:V126-Riccati}
\end{equation}
Write
\begin{equation}
 L_M
 =
 \rho I+M+\rho^{-1}C+O(\rho^{-2}).
 \label{eq:V126-LM}
\end{equation}
On the reflected side use \(s=-r\) and
\(\mathcal A_X(s)=\mathcal A_M(-s)\) after the natural reflection of form
indices.

\begin{lemma}[First two Riccati rows]
\label{lem:V126-first-two-rows}
In the frozen matrix sector, the order-one and order-zero equations in
\eqref{eq:V126-Riccati} have the form
\begin{align}
 2\rho M
 &=
 \mathcal A_1-(\partial_r\rho)I,
 \label{eq:V126-M-row}\\
 2C
 &=
 \mathcal A_0-\partial_rM-M^2,
 \label{eq:V126-C-row}
\end{align}
after the factor \(\rho^{-1}\) in \eqref{eq:V126-LM} is accounted for.
\end{lemma}

\begin{proof}
Expand
\[
 \partial_rL+L^2
 =
 \rho^2I+
 [(\partial_r\rho)I+2\rho M]
 +
 [\partial_rM+M^2+2C]
 +O(\rho^{-1}).
\]
Match the order-one and order-zero terms with
\eqref{eq:V126-frozen-operator}.
\end{proof}

\begin{theorem}[Odd/even reflection parity]
\label{thm:V126-reflection-parity}
For the smooth jet-matched auxiliary side,
\begin{equation}
 \boxed{
 L_X
 =
 \rho I-M+\rho^{-1}C+O(\rho^{-2}).}
 \label{eq:V126-LX}
\end{equation}
Thus the order-zero root coefficient is odd and the order-\(-1\)
coefficient is even under normal reflection.
\end{theorem}

\begin{proof}
At the interface,
\(\partial_s\rho_X=-\partial_r\rho_M\).  The order-one coefficient
\(\mathcal A_1\) also changes by the reflected form-index convention so
that the full order-zero conormal matrix is the negative matrix already
proved in Theorem~\ref{thm:V125-subprincipal-cancellation}.  Hence
\(M_X=-M_M\).

As functions of the inward coordinates,
\(M_X(s)=-M_M(-s)\), so
\[
 \partial_sM_X(0)=\partial_rM_M(0).
\]
The reflected zeroth-order Schrödinger coefficient
\(\mathcal A_0\) is even at the smooth cut: terms with two normal
derivatives and terms quadratic in the first normal jet are unchanged.
Finally \(M_X^2=M_M^2\).  Equation
\eqref{eq:V126-C-row} therefore gives \(C_X=C_M\).
\end{proof}

\begin{firewall}[Scope of the parity theorem]
\label{fire:V126-parity-scope}
Theorem~\ref{thm:V126-reflection-parity} is an exact statement for the
tangentially frozen matrix Riccati recursion.  In the full
pseudodifferential product, order-\(-1\) also receives tangential
derivative and Poisson-bracket terms.  Their reflection parity must be
computed in the geometric symbol calculus; V126 does not assign them the
matrix \(C\).
\end{firewall}

\subsection{A universal logarithm lemma}
\label{subsec:V126-logarithm-lemma}

Let \(M\) and \(C\) be \(r\times r\) matrices independent of \(\rho\), and
put
\begin{equation}
 N_+(\rho)=\rho I+M+\rho^{-1}C,
 \qquad
 N_-(\rho)=\rho I-M+\rho^{-1}C.
 \label{eq:V126-Npm}
\end{equation}

\begin{lemma}[Sum and one-sided logarithms through order two]
\label{lem:V126-log-expansions}
As \(\rho\to\infty\),
\begin{align}
 \operatorname{tr}\log(N_++N_-)
 &=
 r\log(2\rho)+\rho^{-2}\operatorname{tr}C
 +O(\rho^{-3}),
 \label{eq:V126-log-sum}\\
 \operatorname{tr}\log N_-
 &=
 r\log\rho-\rho^{-1}\operatorname{tr}M
 \notag\\
 &\qquad
 +\rho^{-2}
 \left(
 \operatorname{tr}C-\frac12\operatorname{tr}M^2
 \right)
 +O(\rho^{-3}),
 \label{eq:V126-log-minus}\\
 \operatorname{tr}\log N_+
 &=
 r\log\rho+\rho^{-1}\operatorname{tr}M
 \notag\\
 &\qquad
 +\rho^{-2}
 \left(
 \operatorname{tr}C-\frac12\operatorname{tr}M^2
 \right)
 +O(\rho^{-3}).
 \label{eq:V126-log-plus}
\end{align}
\end{lemma}

\begin{proof}
The sum is
\[
 N_++N_-=2\rho(I+\rho^{-2}C).
\]
For each one-sided factor, extract \(\rho I\) and use
\[
 \operatorname{tr}\log(I+X)
 =
 \operatorname{tr}X-\frac12\operatorname{tr}X^2+O(\|X\|^3)
\]
with \(X=\pm\rho^{-1}M+\rho^{-2}C\).
\end{proof}

\subsection{Scalar quadratic row}
\label{subsec:V126-scalar-row}

For the scalar root, put
\begin{equation}
 m:=s_0=\frac12(k-H)
 \label{eq:V126-m}
\end{equation}
and write the reflected pair as
\begin{equation}
 N_{M,0}=\rho+m+\rho^{-1}c+O(\rho^{-2}),
 \qquad
 N_{X,0}=\rho-m+\rho^{-1}c+O(\rho^{-2}).
 \label{eq:V126-scalar-pair}
\end{equation}

\begin{proposition}[Scalar logarithm through its algebraic quadratic row]
\label{prop:V126-scalar-quadratic}
After the common principal factor \(2\) is removed,
\begin{equation}
 \boxed{
 \log\mathcal W_{0,R}
 =
 \frac{m}{\rho}
 +\frac{m^2}{2\rho^2}
 +O(\rho^{-3})}
 \label{eq:V126-scalar-quadratic}
\end{equation}
in the frozen jet-matched sector.  The even Riccati coefficient \(c\)
cancels.
\end{proposition}

\begin{proof}
The normalized scalar logarithm is
\[
 \log(N_{M,0}+N_{X,0})-\log N_{X,0}.
\]
Apply \eqref{eq:V126-log-sum} and
\eqref{eq:V126-log-minus} with rank one, then remove the constant
\(\log2\).  The two occurrences of \(c/\rho^2\) cancel.
\end{proof}

\subsection{One-form quadratic row}
\label{subsec:V126-oneform-row}

Use the V125 block notation
\begin{equation}
 M=
 \begin{pmatrix}
 M_T&U\\
 U^\dagger&m_N
 \end{pmatrix},
 \qquad
 M_T=s_0I_3+S,
 \qquad
 m_N=s_0.
 \label{eq:V126-M-blocks}
\end{equation}
The physical and auxiliary full responses are
\begin{equation}
 N_{M,1}=\rho I_4+M+\rho^{-1}C+O(\rho^{-2}),
 \qquad
 N_{X,1}=\rho I_4-M+\rho^{-1}C+O(\rho^{-2}).
 \label{eq:V126-oneform-pair}
\end{equation}

\begin{theorem}[One-form algebraic quadratic row]
\label{thm:V126-oneform-quadratic}
For the relative/absolute polarization of V124, the normalized one-form
logarithm obeys
\begin{equation}
 \boxed{
 \operatorname{tr}\log\mathcal W_{1,R}
 =
 \frac{\operatorname{tr}M_T-m_N}{\rho}
 +
 \frac{\operatorname{tr}M_T^2+m_N^2}{2\rho^2}
 +O(\rho^{-3})}
 \label{eq:V126-oneform-quadratic}
\end{equation}
after the rank-four principal factor \(2I_4\) is removed.  Every block of
the even coefficient \(C\) cancels from this row.
\end{theorem}

\begin{proof}
By \eqref{eq:V124-det-W1},
\begin{align*}
 \operatorname{tr}\log\mathcal W_{1,R}
 &=
 \operatorname{tr}\log(N_{M,1}+N_{X,1})\\
 &\quad
 -\operatorname{tr}\log(N_{X,1})_{TT}
 -\log(N_{M,1})_{NN}.
\end{align*}
Apply Lemma~\ref{lem:V126-log-expansions}.  The numerator contributes
\(\rho^{-2}\operatorname{tr}C\).  The tangential denominator contributes
\[
 -3\log\rho
 +\rho^{-1}\operatorname{tr}M_T
 -\rho^{-2}\operatorname{tr}C_T
 +\frac1{2\rho^2}\operatorname{tr}M_T^2.
\]
The normal denominator contributes
\[
 -\log\rho-\rho^{-1}m_N
 -\rho^{-2}C_N+\frac1{2\rho^2}m_N^2.
\]
Since \(\operatorname{tr}C=\operatorname{tr}C_T+C_N\), all \(C\)-terms
cancel.  Removing \(4\log2\) gives
\eqref{eq:V126-oneform-quadratic}.
\end{proof}

\begin{remark}[Disposition of the off-diagonal shape pair]
\label{rem:V126-offdiagonal-disposition}
The off-diagonal block \(U=-iSn\) cancels between the two sides in the
glued numerator before the logarithm is expanded.  It is absent from the
relative/absolute reference denominator.  Hence it contributes no
quadratic term in the smooth jet-matched frozen sector.  The V124 warning
about a second-order \(UU^\dagger\) term applies when the auxiliary
response is not jet matched and \(U\) survives in the numerator.
\end{remark}

\subsection{The universal BRST curvature square}
\label{subsec:V126-BRST-square}

\begin{theorem}[Directional quadratic curvature symbol]
\label{thm:V126-directional-quadratic}
Per gauge generator, the known algebraic order-\(-2\) row of the BRST
interface logarithm is
\begin{align}
 Q_K(n)
 &:=
 \frac12s_0^2
 -\frac14\left(
 \operatorname{tr}(s_0I_3+S)^2+s_0^2
 \right)
 \notag\\
 &=
 \boxed{
 \frac{H^2-K(n,n)^2}{8}
 -\frac14|K|^2.}
 \label{eq:V126-QK}
\end{align}
Thus
\begin{equation}
 \left[
 \log\mathcal W_{0,R}
 -\frac12\operatorname{tr}\log\mathcal W_{1,R}
 \right]_{\rm alg,-2}
 =
 \frac{Q_K(n)}{\rho^2}.
 \label{eq:V126-BRST-order-minus-two}
\end{equation}
\end{theorem}

\begin{proof}
The scalar quadratic coefficient is \(s_0^2/2\) by
Proposition~\ref{prop:V126-scalar-quadratic}.  The one-form coefficient is
one half of
\[
 \operatorname{tr}(s_0I_3+S)^2+s_0^2
 =
 4s_0^2+2s_0H+|S|^2.
\]
The determinant weight contributes minus one half, producing
\[
 -\frac12s_0^2-\frac12s_0H-\frac14|S|^2.
\]
Since \(s_0=(k-H)/2\), the first two terms simplify to
\((H^2-k^2)/8\).  Finally \(|S|^2=|K|^2\).
\end{proof}

\begin{lemma}[Angular fourth moment]
\label{lem:V126-angular-fourth}
In three boundary dimensions,
\begin{equation}
 \fint_{\mathbb S_y^2}K(n,n)^2\,dn
 =
 \frac{H^2+2|K|^2}{15}.
 \label{eq:V126-k-square-average}
\end{equation}
\end{lemma}

\begin{proof}
Rotational invariance gives
\[
 \fint_{\mathbb S^2}
 n_in_jn_kn_l\,dn
 =
 \frac1{15}
 (\gamma_{ij}\gamma_{kl}
 +\gamma_{ik}\gamma_{jl}
 +\gamma_{il}\gamma_{jk}).
\]
Contract with \(K^{ij}K^{kl}\).
\end{proof}

\begin{corollary}[Angularly averaged algebraic row]
\label{cor:V126-angular-Q}
The directional average is
\begin{equation}
 \boxed{
 \fint_{\mathbb S_y^2}Q_K(n)\,dn
 =
 \frac7{60}H^2-\frac4{15}|K|^2.}
 \label{eq:V126-Q-average}
\end{equation}
For the twelve Standard Model gauge generators, the averaged coefficient is
\begin{equation}
 \frac75H^2-\frac{16}{5}|K|^2.
 \label{eq:V126-SM-Q-average}
\end{equation}
\end{corollary}

\begin{proof}
Insert \eqref{eq:V126-k-square-average} into
\eqref{eq:V126-QK} and simplify.  Multiply by twelve for the second row.
\end{proof}

\begin{corollary}[Linear sharp-cutoff divergence of the algebraic row]
\label{cor:V126-linear-divergence}
With a common radial boundary cutoff \(|\xi|\le\Lambda\), the known
order-\(-2\) algebraic part contributes
\begin{equation}
 \boxed{
 \frac{\Lambda}{\pi^2}
 \int_\Sigma
 \left(
 \frac7{10}H^2-\frac85|K|^2
 \right)d{\rm vol}_\gamma}
 \label{eq:V126-linear-divergence}
\end{equation}
to the twelve-generator interface logarithm, before the derivative and
intrinsic order-\(-2\) rows are added.
\end{corollary}

\begin{proof}
In three dimensions,
\[
 \frac1{(2\pi)^3}
 \int_{|\xi|\le\Lambda}\frac{d^3\xi}{|\xi|^2}
 =
 \frac{\Lambda}{2\pi^2}.
\]
Multiply by \eqref{eq:V126-SM-Q-average}.
\end{proof}

\subsection{Missing rows at the same order}
\label{subsec:V126-missing-rows}

\begin{firewall}[The quadratic row is not the full heat coefficient]
\label{fire:V126-not-full-coefficient}
At homogeneous order \(-2\), the complete geometric symbol also contains:
\begin{enumerate}[label=\textup{(\roman*)},leftmargin=3.3em]
\item tangential derivatives \(D K\) and \(D H\);
\item Poisson-bracket terms from pseudodifferential composition;
\item intrinsic boundary Ricci and bundle-curvature terms;
\item the nonfrozen part of the order-\(-1\) Riccati root; and
\item possible local changes induced by the far-boundary reference before
      the smoothing limit is taken.
\end{enumerate}
Equation \eqref{eq:V126-linear-divergence} is only the universal algebraic
curvature-square contribution.  It is not the complete linear divergence
and is not assigned as a final Standard Model--Einstein counterterm.
\end{firewall}

\begin{assessment}[Parity gain]
\label{ass:V126-parity-gain}
The jet-matched reflection removes the unknown even Riccati coefficient
from the frozen relative determinant.  This is a genuine reduction of the
next heat problem: the remaining work at order \(-2\) is confined to the
geometric symbol-composition rows listed above, rather than an arbitrary
matrix square-root coefficient.
\end{assessment}

\subsection{V126 ledger and next acquisition}
\label{subsec:V126-ledger}

\paragraph{New exact conclusions.}
\begin{enumerate}[label=\textup{(\arabic*)},leftmargin=2.8em]
\item Theorem~\ref{thm:V126-reflection-parity} proves odd/even parity of
      the first two frozen Riccati corrections.
\item Lemma~\ref{lem:V126-log-expansions} gives the matrix logarithm rows
      needed for the normalized determinant.
\item Proposition~\ref{prop:V126-scalar-quadratic} and
      Theorem~\ref{thm:V126-oneform-quadratic} prove cancellation of the
      unknown even Riccati coefficient.
\item Theorem~\ref{thm:V126-directional-quadratic} computes the universal
      directional curvature-square symbol.
\item Corollary~\ref{cor:V126-angular-Q} reduces it to the invariants
      \(H^2\) and \(|K|^2\).
\item Corollary~\ref{cor:V126-linear-divergence} prints its common
      sharp-cutoff contribution while retaining the missing-row firewall.
\end{enumerate}

\paragraph{Next forced step.}
V127 should compute the tangential-derivative and Poisson-bracket part of
the order-\(-2\) symbol in geometric Weyl quantization.  It must test
whether its angular trace is a total boundary divergence on closed
\(\Sigma\).  The intrinsic Ricci and bundle-curvature rows should remain
separate until that test is complete.

\begin{assessment}[V126 continuum status]
\label{ass:V126-continuum-status}
The full Standard--Model--Einstein continuum remains unproved.  V126
computes the universal algebraic curvature-square row of the relative
BRST interface determinant.  Derivative and intrinsic heat rows, fermion
extension, interacting nonlinear existence, Einstein stress closure, and
the cofinal continuum remain open.
\end{assessment}

\section*{V126 exact checks and append-only record}
\addcontentsline{toc}{section}{V126 exact checks and append-only record}

The reflection-parity and quadratic checksum is
\begin{equation}
 \boxed{
 L_M=\rho I+M+\rho^{-1}C+\cdots,
 \quad
 L_X=\rho I-M+\rho^{-1}C+\cdots,
 \quad
 \fint Q_K=\frac7{60}H^2-\frac4{15}|K|^2.}
 \label{eq:V126-final-check}
\end{equation}
The append operation preserves the complete V125 prefix byte for byte before
its former live terminator.  The assembled V126 file is checked for one live
final terminator, balanced environments and braces, duplicate labels, newly
introduced unresolved references, display math delimiters, and control
characters before the exact V125 predecessor is removed.

% =============================================================================
% END APPENDED VERSION V126 MATERIAL
% =============================================================================

\clearpage
% =============================================================================
% BEGIN APPENDED VERSION V127 MATERIAL
% =============================================================================
\section{V127: Weyl-product derivative rows are phase-space divergences}
\label{sec:V127-Weyl-divergence}

V126 computed the algebraic order-\(-2\) curvature row in the
tangentially frozen Riccati calculus and left the first tangential
derivatives and Poisson brackets open.  V127 places the interface operators
on boundary half-densities and uses geometric Weyl quantization.  The first
star-product correction is a phase-space Poisson bracket.  With a common
covariant radial cutoff on a closed boundary, its integrated trace vanishes
exactly because the relevant Hamiltonian flow is tangent to the cutoff
surface.

This closes the integrated row linear in \(DK\) and \(DH\).  It does not
compute the pointwise representative of the local heat density, which
depends on symbol gauge, nor the second-derivative and intrinsic-curvature
terms in the Weyl calculus.

\subsection{Geometric Weyl conventions}
\label{subsec:V127-Weyl-conventions}

Work on half-densities over the closed three-manifold \(\Sigma\), with a
metric connection on the finite form/gauge bundle.  Let
\(\nabla_a^H\) denote the horizontal covariant derivative on
\(T^*\Sigma\), and let \(\partial_{\xi_a}\) be the vertical derivative.
For matrix symbols \(a,b\), define
\begin{equation}
 \{a,b\}
 :=
 \partial_{\xi_a}a\,\nabla_a^Hb
 -\nabla_a^Ha\,\partial_{\xi_a}b.
 \label{eq:V127-Poisson}
\end{equation}
The geometric Weyl product begins
\begin{equation}
 a\# b
 =
 ab+\frac{i}{2}\{a,b\}
 +O(\operatorname{ord}a+\operatorname{ord}b-2).
 \label{eq:V127-Weyl-product}
\end{equation}
Write
\begin{equation}
 \rho(y,\xi):=|\xi|_\gamma,
 \qquad
 \Omega_\Lambda
 :=
 \{(y,\xi):\rho(y,\xi)\le\Lambda\}.
 \label{eq:V127-radial-cutoff}
\end{equation}

\begin{lemma}[Trace of a scalar-matrix Poisson bracket]
\label{lem:V127-trace-bracket}
For a scalar symbol \(f\) and a matrix symbol \(M\),
\begin{equation}
 \operatorname{tr}\{fI,M\}
 =
 \{f,\operatorname{tr}M\}.
 \label{eq:V127-trace-bracket}
\end{equation}
\end{lemma}

\begin{proof}
The connection contribution to
\(\nabla^H M\) is a commutator with the bundle connection.  Its matrix trace
vanishes.  Scalar factors commute with all fibre matrices, so taking the
trace of \eqref{eq:V127-Poisson} gives the scalar bracket of the fibre trace.
\end{proof}

\subsection{Radial Hamiltonian flows have no cutoff flux}
\label{subsec:V127-Hamiltonian-flux}

\begin{theorem}[Exact radial Poisson integral]
\label{thm:V127-radial-Poisson-integral}
Let \(F\) be a smooth scalar function and let \(g\) be a classical symbol
for which the cutoff integral is defined.  Then
\begin{equation}
 \boxed{
 \int_{\Omega_\Lambda}
 \{F(\rho),g\}\,dyd\xi
 =0}
 \label{eq:V127-radial-Poisson-zero}
\end{equation}
for every \(\Lambda>0\), provided \(\Sigma\) is closed.
\end{theorem}

\begin{proof}
Let \(X_{F(\rho)}\) be the Hamiltonian vector field of \(F(\rho)\).
Hamiltonian fields preserve the symplectic volume, so
\[
 \{F(\rho),g\}
 =
 X_{F(\rho)}g
 =
 \operatorname{div}_{dyd\xi}(gX_{F(\rho)}).
\]
There is no base-boundary flux because \(\partial\Sigma=\varnothing\).
On the fibre boundary \(\rho=\Lambda\),
\[
 X_{F(\rho)}\rho
 =
 \{\rho,F(\rho)\}=0.
\]
Thus the Hamiltonian vector field is tangent to the radial cutoff surface
and its normal flux vanishes.  The divergence theorem proves
\eqref{eq:V127-radial-Poisson-zero}.
\end{proof}

\begin{corollary}[Matrix version]
\label{cor:V127-matrix-Poisson-zero}
For every matrix symbol \(M\),
\begin{equation}
 \int_{\Omega_\Lambda}
 \operatorname{tr}\{F(\rho)I,M\}\,dyd\xi
 =0.
 \label{eq:V127-matrix-Poisson-zero}
\end{equation}
\end{corollary}

\begin{proof}
Combine Lemma~\ref{lem:V127-trace-bracket} with
Theorem~\ref{thm:V127-radial-Poisson-integral}.
\end{proof}

\begin{remark}[Why the covariant cutoff matters]
\label{rem:V127-covariant-cutoff}
If the cutoff is imposed in coordinate components rather than by
\(\rho=|\xi|_\gamma\), the Hamiltonian field of \(F(\rho)\) need not be
tangent to its artificial faces.  Coordinate-face terms can then appear.
They are regulator artefacts and are absent from the common covariant
radial cutoff used here.
\end{remark}

\subsection{The first star correction to a logarithmic trace}
\label{subsec:V127-log-star}

Let
\begin{equation}
 A_\varepsilon
 :=
 \operatorname{Op}^W(\rho I+\varepsilon M)
 \label{eq:V127-Aepsilon}
\end{equation}
on a common finite spectral regulator, with \(\rho>0\).  Denote its
regulated trace by \(\operatorname{Tr}_\Lambda\).

\begin{lemma}[Variation of a regulated operator logarithm]
\label{lem:V127-log-variation}
For invertible \(A_\varepsilon\),
\begin{equation}
 \partial_\varepsilon
 \operatorname{Tr}_\Lambda\log A_\varepsilon
 =
 \operatorname{Tr}_\Lambda
 (A_\varepsilon^{-1}\partial_\varepsilon A_\varepsilon).
 \label{eq:V127-log-variation}
\end{equation}
\end{lemma}

\begin{proof}
In finite dimensions this is Jacobi's identity
\(\partial_\varepsilon\log\det A
=\operatorname{Tr}(A^{-1}\dot A)\).  The same formula applies to the
common finite spectral regulator before any asymptotic limit.
\end{proof}

\begin{proposition}[The order-\(-2\) first Weyl term integrates to zero]
\label{prop:V127-first-Weyl-zero}
At \(\varepsilon=0\), the order-\(-2\) first-product term in the symbol of
\(A_0^{-1}\partial_\varepsilon A_\varepsilon\) is
\begin{equation}
 \frac{i}{2}\{\rho^{-1}I,M\}.
 \label{eq:V127-first-Weyl-term}
\end{equation}
Its regulated phase-space trace is zero:
\begin{equation}
 \int_{\Omega_\Lambda}
 \operatorname{tr}
 \left[\frac{i}{2}\{\rho^{-1}I,M\}\right]dyd\xi
 =0.
 \label{eq:V127-first-Weyl-integral-zero}
\end{equation}
\end{proposition}

\begin{proof}
The order-\(-1\) inverse symbol of \(\operatorname{Op}^W(\rho I)\) is
\(\rho^{-1}I\).  Its first inverse correction vanishes because
\(\{\rho,\rho^{-1}\}=0\).  Apply the Weyl product
\eqref{eq:V127-Weyl-product} to \(\rho^{-1}I\) and \(M\).
The ordinary product has order \(-1\); the displayed bracket has order
\(-2\).  Corollary~\ref{cor:V127-matrix-Poisson-zero} with
\(F(\rho)=\rho^{-1}\) proves the integral identity.
\end{proof}

\subsection{Pointwise parity of the curvature bracket}
\label{subsec:V127-pointwise-parity}

At a point \(y\), choose geodesic coordinates and a parallel form frame.
Horizontal differentiation then obeys
\begin{equation}
 \nabla_a^H\rho=0,
 \qquad
 \partial_{\xi_a}\rho^{-1}
 =-\rho^{-2}n^a.
 \label{eq:V127-rho-derivatives}
\end{equation}

\begin{lemma}[Odd angular form]
\label{lem:V127-odd-angular}
If \(M(y,n)\) is even under \(n\mapsto-n\), then
\begin{equation}
 \operatorname{tr}\{\rho^{-1}I,M\}
 =
 -\rho^{-2}n^aD_a(\operatorname{tr}M)
 \label{eq:V127-bracket-local}
\end{equation}
is odd in \(n\).  Hence
\begin{equation}
 \int_{\mathbb S_y^2}
 \operatorname{tr}\{\rho^{-1}I,M\}\,dn=0
 \label{eq:V127-angular-bracket-zero}
\end{equation}
pointwise in \(y\).
\end{lemma}

\begin{proof}
Use \eqref{eq:V127-rho-derivatives} in
\eqref{eq:V127-Poisson}.  Parallel transport of \(\xi\) in the horizontal
derivative keeps \(n\) fixed at the chosen point.  The trace of \(M\) is
even, its base derivative remains even, and multiplication by \(n^a\)
makes the result odd.
\end{proof}

\begin{corollary}[Application to the V125 curvature matrices]
\label{cor:V127-curvature-angular-zero}
The scalar symbol
\(s_0=\tfrac12[K(n,n)-H]\), the tangential trace
\(\operatorname{tr}(s_0I_3+S)\), and the normal symbol \(s_0\) are all even
in \(n\).  Every first Weyl bracket generated from these order-zero
curvature symbols has zero angular trace.
\end{corollary}

\begin{proof}
\(K(n,n)\) is quadratic in \(n\), while \(H\) and
\(\operatorname{tr}S\) are independent of \(n\).  Apply
Lemma~\ref{lem:V127-odd-angular}.
\end{proof}

\subsection{No hidden first bracket in the Riccati square}
\label{subsec:V127-Riccati-star}

\begin{proposition}[Symmetric root products cancel their first bracket]
\label{prop:V127-Riccati-bracket-cancel}
For scalar \(\rho\) and matrix \(M\),
\begin{equation}
 \rho I\# M+M\#\rho I
 =
 2\rho M+O(\rho^{-1})
 \label{eq:V127-symmetric-product}
\end{equation}
with no first Poisson-bracket term.  Moreover,
\begin{equation}
 \operatorname{tr}\{M,M\}=0
 \label{eq:V127-MM-trace-zero}
\end{equation}
pointwise.
\end{proposition}

\begin{proof}
The first Weyl terms are
\[
 \frac{i}{2}\{\rho I,M\}
 +\frac{i}{2}\{M,\rho I\}=0
\]
by antisymmetry of the Poisson bracket and scalar commutation.
For the second identity,
\[
 \{M,M\}
 =
 [\partial_{\xi_a}M,\nabla_a^HM],
\]
whose fibre trace is zero.
\end{proof}

\begin{corollary}[The V126 even Riccati row has no linear \(DK\) trace]
\label{cor:V127-C-no-first-derivative-trace}
In geometric Weyl quantization, the order-zero equation determining the
order-\(-1\) root coefficient \(C\) has no first Poisson term linear in
\(D M\).  The first possible tangential-symbol corrections at that row
come from second Weyl derivatives of \(\rho\), intrinsic connection
curvature, and the nonfrozen zeroth-order operator coefficient.
\end{corollary}

\begin{proof}
The order-zero Riccati row receives the symmetric cross product of
\(\rho I\) and \(M\).  Proposition~\ref{prop:V127-Riccati-bracket-cancel}
removes its first bracket.  The bracket in \(M\# M\) is one symbolic order
lower and has zero trace in any case.
\end{proof}

\subsection{Integrated disposition of the derivative row}
\label{subsec:V127-disposition}

\begin{theorem}[No integrated order-\(-2\) term linear in \(DK\)]
\label{thm:V127-no-linear-DK}
For the jet-matched scalar/one-form relative determinant on a closed
\(\Sigma\), the integrated homogeneous order-\(-2\) contribution generated
by first Weyl products of the V117--V118 order-zero curvature matrices is
zero:
\begin{equation}
 \boxed{
 \left[
 \operatorname{Tr}_\Lambda
 \left(
 \log\mathcal W_{0,R}
 -\frac12\operatorname{tr}\log\mathcal W_{1,R}
 \right)
 \right]_{\textup{first Weyl},\, -2}
 =0.}
 \label{eq:V127-no-linear-DK}
\end{equation}
Equivalently, there is no integrated interface counterterm linear in one
tangential derivative of \(K\) or \(H\) at this order.
\end{theorem}

\begin{proof}
The glued numerator has no order-zero extrinsic curvature matrix by
Theorem~\ref{thm:V125-subprincipal-cancellation}.  The possible first Weyl
terms therefore arise from the scalar auxiliary denominator and the
tangential/normal one-form reference denominators.  Each has the form
\eqref{eq:V127-first-Weyl-term} with an even curvature matrix.  Its
phase-space trace vanishes by
Proposition~\ref{prop:V127-first-Weyl-zero}; equivalently its angular
average vanishes by
Corollary~\ref{cor:V127-curvature-angular-zero}.  The Riccati root itself
adds no first bracket at the required row by
Corollary~\ref{cor:V127-C-no-first-derivative-trace}.  Linear combination
with weights \(1\) and \(-1/2\) preserves zero.
\end{proof}

\begin{firewall}[Local density versus integrated coefficient]
\label{fire:V127-local-versus-integrated}
The pointwise Weyl symbol can contain the divergence
\eqref{eq:V127-bracket-local}.  V127 proves that its covariantly regulated
integral is zero on a closed boundary.  If the boundary has corners, an
edge, or a noncovariant cutoff, the flux term must be retained.  V127 does
not set the local representative to zero before integration.
\end{firewall}

\begin{assessment}[Order-\(-2\) status]
\label{ass:V127-order-two-status}
The algebraic \(H^2,|K|^2\) row in V126 is not modified by first Weyl
products after integration over a closed interface.  The remaining
integrated order-\(-2\) uncertainty is confined to intrinsic boundary
curvature, bundle curvature, and second-symbol-derivative terms.  These
cannot be inferred from the phase-space divergence theorem.
\end{assessment}

\subsection{V127 ledger and next acquisition}
\label{subsec:V127-ledger}

\paragraph{New exact conclusions.}
\begin{enumerate}[label=\textup{(\arabic*)},leftmargin=2.8em]
\item Theorem~\ref{thm:V127-radial-Poisson-integral} proves exact
      vanishing of radial Hamiltonian flux for the common covariant cutoff.
\item Proposition~\ref{prop:V127-first-Weyl-zero} identifies and removes
      the first star-product term in the integrated logarithmic trace.
\item Lemma~\ref{lem:V127-odd-angular} proves the equivalent pointwise
      angular-parity cancellation.
\item Proposition~\ref{prop:V127-Riccati-bracket-cancel} shows that the
      symmetric Riccati square has no first bracket at the even root row.
\item Theorem~\ref{thm:V127-no-linear-DK} proves absence of an integrated
      order-\(-2\) counterterm linear in \(DK\) or \(DH\).
\end{enumerate}

\paragraph{Next forced step.}
V128 should compute the intrinsic second-symbol-derivative row in
geometric Weyl calculus.  It must separate the scalar boundary curvature
\(R_\Sigma\), the trace of the one-form Weitzenbock endomorphism, and the
Standard Model adjoint bundle curvature.  Any connection-curvature term
whose fibre trace vanishes must be proved to vanish before the remaining
coefficient is combined with V126.

\begin{assessment}[V127 continuum status]
\label{ass:V127-continuum-status}
The full Standard--Model--Einstein continuum remains unproved.  V127
closes the integrated first-derivative/Poisson row of the order-\(-2\)
relative interface symbol.  Intrinsic heat rows, fermion extension,
interacting nonlinear existence, Einstein stress closure, and the cofinal
continuum remain open.
\end{assessment}

\section*{V127 exact checks and append-only record}
\addcontentsline{toc}{section}{V127 exact checks and append-only record}

The Weyl-divergence checksum is
\begin{equation}
 \boxed{
 \int_{\rho\le\Lambda}
 \operatorname{tr}\{\rho^{-1}I,M\}\,dyd\xi=0,
 \qquad
 [\operatorname{Tr}_\Lambda\log\mathcal W_{\rm BRST}]
 _{\textup{first Weyl},-2}=0.}
 \label{eq:V127-final-check}
\end{equation}
The append operation preserves the complete V126 prefix byte for byte before
its former live terminator.  The assembled V127 file is checked for one live
final terminator, balanced environments and braces, duplicate labels, newly
introduced unresolved references, display math delimiters, and control
characters before the exact V126 predecessor is removed.

% =============================================================================
% END APPENDED VERSION V127 MATERIAL
% =============================================================================

% =============================================================================
% BEGIN APPENDED VERSION V128 MATERIAL
% =============================================================================

\section{V128: intrinsic-root cancellation and closure of the integrated
order-minus-two interface row}
\label{sec:V128}

\subsection{Purpose and scope}
\label{subsec:V128-purpose}

V126 computed the part of the homogeneous order-\(-2\) interface logarithm
which is algebraic in the second fundamental form.  V127 then proved that
the first Weyl-product row linear in \(DK\) integrates to zero for the
covariant radial cutoff on a closed interface.  The unresolved terms at
that order were the scalar curvature of \(\Sigma\), the one-form
Weitzenbock endomorphism, the gauge-bundle curvature, and the second
symbol derivatives hidden in the Riccati root.  This note removes that
ambiguity by a product calibration.  The conclusion concerns the
\emph{relative interface determinant}; it does not remove the ordinary
local heat coefficients of the decoupled physical relative Hodge
problem.

We retain the jet-matched doubled collar of V125 and geometric Weyl
quantization from V127.  Put
\[
 \rho_j(y,\xi;z)=\bigl(|\xi|_{\gamma(y)}^2+z\bigr)^{1/2}
\]
on \(j\)-forms over \(\Sigma\), with the identity on the corresponding
bundle suppressed.  All fibre traces below include the form indices but
not the Standard Model adjoint multiplicity unless it is displayed.

\subsection{The order-minus-one Riccati root}
\label{subsec:V128-Riccati-root}

Let \(L_\pm\) denote the outgoing roots on the two reflected sides of the
interface.  In a parallel frame at a fixed boundary point their geometric
Weyl symbols have the form
\begin{equation}
 \ell_\pm
 =\rho I\mathbin{\pm}M+\rho^{-1}C_\pm+O(\rho^{-2}),
 \label{eq:V128-root-expansion}
\end{equation}
where \(M\) is the curvature matrix computed in V117--V118.  Write the
normal operator, after the common half-density conjugation, as
\begin{equation}
 -\partial_t^2+\mathcal A_1(t)\partial_t+\mathcal A_2(t,z).
 \label{eq:V128-normal-operator}
\end{equation}
The order-zero Riccati row is intrinsic to the complete boundary jet of
this operator.

\begin{lemma}[Geometric Weyl Riccati row]
\label{lem:V128-Riccati-row}
At the boundary, the homogeneous order-\(-1\) coefficient satisfies
\begin{equation}
 2\rho C_\pm
 =
 \mathcal A_{\pm,0}
 -\partial_\pm M_\pm
 -M_\pm^2
 -\mathcal B_2(\rho,\rho).
 \label{eq:V128-Riccati-row}
\end{equation}
Here \(\mathcal A_{\pm,0}\) is the homogeneous order-zero part of the
tangential operator, and \(\mathcal B_2\) is the universal second Weyl
bidifferential.  No first Weyl bracket occurs in
\eqref{eq:V128-Riccati-row}.
\end{lemma}

\begin{proof}
Substitute \eqref{eq:V128-root-expansion} into the factorization Riccati
equation for \eqref{eq:V128-normal-operator} and compare homogeneous
degree zero.  The ordinary product of \(\rho I\) with
\(\rho^{-1}C_\pm\), in both orders, gives \(2C_\pm\); multiplication by
\(\rho\) after the convention used in \eqref{eq:V128-root-expansion}
gives the left side of \eqref{eq:V128-Riccati-row}.  The remaining
degree-zero terms are the tangential coefficient, the normal derivative
of the order-zero root, the square of that root, and the second
bidifferential in \(\rho I\#\rho I\).  The two first brackets in
\(\rho I\# M+M\#\rho I\) cancel by
Proposition~\ref{prop:V127-Riccati-bracket-cancel}.  This proves the
displayed identity, independently of the coordinate representative of
\(\mathcal B_2\).
\end{proof}

\begin{lemma}[Reflection parity of the unknown root row]
\label{lem:V128-reflection-parity}
Under the jet reflection of V125, the scalar trace, the tangential
one-form trace, and the normal one-form entry in
\[
 C_M-C_X
\]
have zero integrated even intrinsic part.  Terms odd in \(DK\) or \(DH\)
are either off diagonal in the tangential/normal splitting or have zero
covariant radial trace by V127.
\end{lemma}

\begin{proof}
The reflected metric and connection jets satisfy
\(\partial_s^m=(-1)^m\partial_r^m\) at the cut.  In
\eqref{eq:V128-Riccati-row}, \(\mathcal A_0\) and
\(\mathcal B_2(\rho,\rho)\) depend only on the common intrinsic operator
and are even.  The normal derivative changes sign at the same time as
the order-zero root \(M\), so its diagonal trace is even as well.
Consequently the intrinsic diagonal trace agrees on the two sides and
drops out of \(C_M-C_X\).  The remaining first tangential derivative
terms are precisely the first-symbol terms classified in
Theorem~\ref{thm:V127-no-linear-DK}; their integrated trace is zero.
Tangential/normal off-diagonal terms do not contribute to a fibre trace.
\end{proof}

\begin{corollary}[Cancellation in both Hodge polarizations]
\label{cor:V128-root-trace-cancel}
With the relative/absolute polarization of V124,
\begin{align}
 \int_{\Omega_\Lambda}
 \operatorname{tr}(C_{M,0}-C_{X,0})\,dyd\xi&=0,
 \label{eq:V128-scalar-C-cancel}\\
 \int_{\Omega_\Lambda}
 \left[
 \operatorname{tr}_T(C_{M,1}-C_{X,1})
 -(C_{M,1}-C_{X,1})_{NN}
 \right]dyd\xi&=0
 \label{eq:V128-oneform-C-cancel}
\end{align}
for the even intrinsic row.
\end{corollary}

\begin{proof}
Apply Lemma~\ref{lem:V128-reflection-parity} to the scalar block and to
the tangential and normal diagonal blocks separately.  The relative
one-form reference uses their difference, so equality before taking that
linear combination proves \eqref{eq:V128-oneform-C-cancel}.
\end{proof}

\subsection{Exact product calibration}
\label{subsec:V128-product-calibration}

The parity calculation shows how the unknown Riccati coefficient drops
out.  The following exact model fixes the possible intrinsic invariant
which could otherwise be hidden by a change of symbol representative.

\begin{proposition}[Exact normalized product operators]
\label{prop:V128-product-operators}
Suppose the doubled metric and connection are products in a neighborhood
of \(\Sigma\), with arbitrary closed intrinsic metric and arbitrary fixed
bundle connection.  Let
\[
 \Lambda_j(z)=(\Delta_{\Sigma,j}+z)^{1/2}.
\]
For the finite auxiliary collar with the far-boundary conditions of
V122, the normalized interface operators are
\begin{align}
 \widetilde{\mathcal W}_{0,R}(z)
 &=(I-e^{-2R\Lambda_0(z)})^{-1},
 \label{eq:V128-product-W0}\\
 \widetilde{\mathcal W}_{1,R}(z)
 &=
 \operatorname{diag}\left(
 (I-e^{-2R\Lambda_1(z)})^{-1},
 (I-e^{-2R\Lambda_0(z)})^{-1}
 \right).
 \label{eq:V128-product-W1}
\end{align}
On the positive spectral subspace,
\begin{equation}
 \widetilde{\mathcal W}_{j,R}(z)-I\in\Psi^{-\infty}(\Sigma).
 \label{eq:V128-product-smoothing}
\end{equation}
\end{proposition}

\begin{proof}
The product Hodge Laplacian separates into tangential spectral modes.
For an eigenvalue \(\mu^2\) of \(\Delta_{\Sigma,j}\), write
\(\lambda=(\mu^2+z)^{1/2}\) with the branch from V119.  The physical
outgoing response is \(\lambda\).  The auxiliary Neumann response is
\(\lambda\tanh(R\lambda)\), while the auxiliary Dirichlet response is
\(\lambda\coth(R\lambda)\).  Dividing the conormal sums by the reference
response, as in V124, gives in either Hodge polarization
\[
 \frac{\lambda+\lambda\coth(R\lambda)}{2\lambda}
 =\frac{1}{1-e^{-2R\lambda}}.
\]
The tangential and normal one-form blocks use \(\Lambda_1\) and
\(\Lambda_0\), respectively, which proves
\eqref{eq:V128-product-W0}--\eqref{eq:V128-product-W1} by spectral
calculus.  For \(\mu\to\infty\),
\[
 (1-e^{-2R\lambda})^{-1}-1
 =\sum_{m\ge1}e^{-2mR\lambda}
\]
decays faster than every inverse power of \(\mu\), together with all
symbol derivatives.  This is the defining estimate for a smoothing
operator.
\end{proof}

\begin{corollary}[No local intrinsic interface coefficient in the
product case]
\label{cor:V128-no-product-local}
For every homogeneous finite-order symbol row,
\begin{equation}
 \sigma_{\mathrm{hom}}
 \log\widetilde{\mathcal W}_{0,R}
 =
 \sigma_{\mathrm{hom}}
 \log\widetilde{\mathcal W}_{1,R}
 =0
 \label{eq:V128-product-symbol-zero}
\end{equation}
in the product geometry.  In particular, the relative interface
logarithm has no local term proportional to \(R_\Sigma\), nor a term
linear in the intrinsic bundle curvature, when \(K=0\).
\end{corollary}

\begin{proof}
Equation~\eqref{eq:V128-product-smoothing} and holomorphic functional
calculus give
\(\log\widetilde{\mathcal W}_{j,R}\in\Psi^{-\infty}\) away from the
finite-dimensional threshold sector.  A smoothing operator has no
homogeneous symbol components.  Since the intrinsic geometry and
connection in Proposition~\ref{prop:V128-product-operators} were
arbitrary, the asserted local coefficients vanish.
\end{proof}

\subsection{Invariant classification and the curvature traces}
\label{subsec:V128-invariants}

\begin{lemma}[Weight-two scalar list]
\label{lem:V128-weight-two-list}
After fibre trace and angular integration on a closed three-dimensional
interface, a natural parity-even scalar of geometric weight two which can
occur in the present interface row is a linear combination of
\begin{equation}
 R_\Sigma,\qquad H^2,\qquad |K|^2,\qquad
 \operatorname{tr}_{E}\mathcal F_{ab}\,T^{ab}.
 \label{eq:V128-invariant-list}
\end{equation}
Here \(T^{ab}\) is formed by the angular average.  Total tangential
divergences integrate to zero.
\end{lemma}

\begin{proof}
Naturality and orthogonal invariance leave one contraction of intrinsic
curvature, the two quadratic contractions of the second fundamental
form, and a possible linear contraction of bundle curvature.  A single
tangential derivative of \(K\) has an unpaired index after even angular
integration and, more directly, was disposed of in V127.  Second
derivatives of a weight-zero scalar are divergences at this row.  These
are the standard invariant-theory possibilities at weight two.
\end{proof}

\begin{lemma}[Form and adjoint curvature traces]
\label{lem:V128-curvature-traces}
For the Hodge Laplacian on one-forms over \(\Sigma\),
\begin{equation}
 \operatorname{tr}_{T^*\Sigma}\operatorname{Ric}=R_\Sigma.
 \label{eq:V128-Ricci-trace}
\end{equation}
For every \(X\) in a compact Standard Model gauge Lie algebra,
\begin{equation}
 \operatorname{tr}_{\mathrm{ad}}(\operatorname{ad}_X)=0.
 \label{eq:V128-ad-trace-zero}
\end{equation}
Consequently no term linear in the gauge curvature survives the adjoint
fibre trace.
\end{lemma}

\begin{proof}
The first identity is the definition of scalar curvature as the trace of
the Ricci endomorphism.  For the second, choose an invariant positive
inner product on the compact Lie algebra.  Then
\(\operatorname{ad}_X\) is skew-adjoint and hence traceless.  Apply this
pointwise to each curvature component.
\end{proof}

\begin{theorem}[Closure of the integrated extrinsic order-minus-two row]
\label{thm:V128-order-two-closure}
For the jet-matched scalar/one-form relative interface determinant, with
the covariant radial regulator and on a closed \(\Sigma\), the complete
local homogeneous order-\(-2\) contribution per gauge generator is
\begin{equation}
 \boxed{
 \fint_{\mathbb S_y^2}
 \sigma_{-2}^{\mathrm{ext}}
 \left(
 \log\mathcal W_{0,R}
 -\frac12\operatorname{tr}\log\mathcal W_{1,R}
 \right)dn
 =
 \rho^{-2}
 \left(
 \frac{7}{60}H^2-\frac{4}{15}|K|^2
 \right).}
 \label{eq:V128-closed-order-two}
\end{equation}
There is no additional \(R_\Sigma\), first-derivative, or fibre-traced
linear gauge-curvature term in this \emph{extrinsic interface} row.
\end{theorem}

\begin{proof}
The algebraic \(H^2\) and \(|K|^2\) coefficients are the angular average
computed in Corollary~\ref{cor:V126-angular-Q}.  The first
symbol-derivative contribution is zero by
Theorem~\ref{thm:V127-no-linear-DK}.  The unknown even Riccati root
cancels between reflected sides by
Corollary~\ref{cor:V128-root-trace-cancel}.  It remains to exclude an
intrinsic invariant which might depend on the symbol representative.
Lemma~\ref{lem:V128-weight-two-list} gives the exhaustive list.
Corollary~\ref{cor:V128-no-product-local}, applied with arbitrary
intrinsic metric, sets the coefficient of \(R_\Sigma\) to zero.
Lemma~\ref{lem:V128-curvature-traces} removes a linear adjoint-curvature
term.  The two extrinsic coefficients remain those of V126, proving
\eqref{eq:V128-closed-order-two}.
\end{proof}

\begin{corollary}[Twelve-generator radial coefficient]
\label{cor:V128-twelve-generator}
For the Standard Model adjoint multiplicity
\(\dim\mathfrak g_{\rm SM}=12\), the sharp radial integral of the closed
extrinsic row is
\begin{equation}
 \boxed{
 \frac{\Lambda}{\pi^2}
 \int_\Sigma
 \left(
 \frac{7}{10}H^2-\frac{8}{5}|K|^2
 \right)d\operatorname{vol}_\Sigma.}
 \label{eq:V128-SM-radial-row}
\end{equation}
\end{corollary}

\begin{proof}
Multiply the per-generator angular coefficient
\eqref{eq:V128-closed-order-two} by \(12\) and use the radial phase-space
normalization already evaluated in
Corollary~\ref{cor:V126-linear-divergence}.
\end{proof}

\begin{firewall}[Interface subtraction versus physical heat]
\label{fire:V128-interface-versus-physical}
The qualification \emph{extrinsic interface} is essential.
Proposition~\ref{prop:V128-product-operators} proves that the normalized
gluing correction is smoothing in product geometry.  The decoupled
physical relative Hodge Laplacian still has its standard bulk and
boundary heat coefficients, including intrinsic and extrinsic curvature
terms.  Those baseline coefficients must be added separately before any
comparison with an Einstein--Hilbert or Gibbons--Hawking--York coupling.
\end{firewall}

\subsection{V128 ledger and next acquisition}
\label{subsec:V128-ledger}

\paragraph{New exact conclusions.}
\begin{enumerate}[label=\textup{(\arabic*)},leftmargin=2.8em]
\item Lemma~\ref{lem:V128-Riccati-row} isolates the complete geometric
      Weyl order-\(-1\) root equation.
\item Corollary~\ref{cor:V128-root-trace-cancel} proves cancellation of
      its even intrinsic trace in both Hodge polarizations.
\item Proposition~\ref{prop:V128-product-operators} computes the exact
      normalized scalar and one-form product operators and proves that
      their deviation from identity is smoothing.
\item Corollary~\ref{cor:V128-no-product-local} rules out an independent
      \(R_\Sigma\) coefficient in the relative interface determinant.
\item Theorem~\ref{thm:V128-order-two-closure} closes the integrated local
      order-\(-2\) extrinsic interface row.
\end{enumerate}

\paragraph{Next forced step.}
V129 must add the physical relative Hodge baseline kept separate in
Firewall~\ref{fire:V128-interface-versus-physical}.  The calculation
should be organized through exact/coexact cancellation, so that only the
coexact one-form boundary coefficient is acquired rather than
recomputing the longitudinal scalar sector.  The resulting boundary
invariants must then be compared with the Einstein and
Gibbons--Hawking--York tensor structures.  V130 is the next compulsory
cross-program audit.

\begin{assessment}[V128 continuum status]
\label{ass:V128-continuum-status}
The full Standard--Model--Einstein continuum remains unproved.  V128
closes one local relative-interface row; it does not establish the
physical relative heat baseline, fermion extension, interacting
nonlinear existence, Einstein stress closure, or the cofinal continuum.
\end{assessment}

\section*{V128 exact checks and append-only record}
\addcontentsline{toc}{section}{V128 exact checks and append-only record}

The product-calibration checksum is
\begin{equation}
 \boxed{
 \widetilde{\mathcal W}_{j,R}-I\in\Psi^{-\infty}
 \quad(K=0),\qquad
 \fint\sigma_{-2}^{\mathrm{ext}}\log\mathcal W_{\rm BRST}
 =
 \rho^{-2}
 \left(\frac{7}{60}H^2-\frac{4}{15}|K|^2\right).}
 \label{eq:V128-final-check}
\end{equation}
The append operation preserves the complete V127 prefix byte for byte before
its former live terminator.  The assembled V128 file is checked for one live
final terminator, balanced environments and braces, duplicate labels, newly
introduced unresolved references, display math delimiters, and control
characters before the exact V127 predecessor is removed.

% =============================================================================
% END APPENDED VERSION V128 MATERIAL
% =============================================================================

% =============================================================================
% BEGIN APPENDED VERSION V129 MATERIAL
% =============================================================================

\section{V129: the physical relative Hodge baseline and the
Einstein--GHY ratio}
\label{sec:V129}

\subsection{Why a separate baseline is necessary}
\label{subsec:V129-purpose}

V123--V128 computed a \emph{relative gluing correction}: the heat trace of
the collar-extended complex minus the heat traces of the decoupled
physical and auxiliary complexes.  Locality confines that difference to
the interface.  It does not include the ordinary heat coefficients of
the physical relative Hodge Laplacians on \(M\).  V129 now computes the
first three terms of that physical baseline.

The calculation has two purposes.  First, it verifies explicitly that
the gauge and ghost area divergences cancel.  Second, it determines the
relative coefficient of the bulk scalar-curvature term and its boundary
mean-curvature completion.  The result has the Einstein--Hilbert to
Gibbons--Hawking--York ratio \(1:2\) before the V128 interface correction
and before fermion, Higgs, and interacting contributions are added.

\subsection{Mixed heat convention}
\label{subsec:V129-heat-convention}

Let \(n=\partial_r\) be the inward unit normal of the V116 collar and put
\begin{equation}
 L_{ab}:=\langle\nabla_{e_a}e_b,n\rangle=-K_{ab},
 \qquad
 L:=\operatorname{tr}_\Sigma L=-H.
 \label{eq:V129-L-convention}
\end{equation}
Thus \(L\) is the mean curvature in the standard inward-normal heat
coefficient convention, while \(H\) is the V117 trace of
\(K_{ab}=\tfrac12\partial_r\gamma_{ab}\).

Write a Laplace-type operator as
\begin{equation}
 D=-\bigl(g^{\mu\nu}\nabla_\mu\nabla_\nu+E\bigr).
 \label{eq:V129-Laplace-type}
\end{equation}
For complementary orthogonal projectors
\(\Pi_-\) and \(\Pi_+=I-\Pi_-\), impose
\begin{equation}
 \Pi_-u|_\Sigma=0,
 \qquad
 \Pi_+(\nabla_n+S)u|_\Sigma=0,
 \qquad
 \chi:=\Pi_+-\Pi_-.
 \label{eq:V129-mixed-BC}
\end{equation}

\begin{theorem}[Heat coefficients through geometric weight two]
\label{thm:V129-mixed-heat}
For a smooth compact four-manifold with smooth boundary and the mixed
condition \eqref{eq:V129-mixed-BC},
\begin{align}
 \operatorname{Tr}e^{-tD}
 \sim (4\pi t)^{-2}\Bigg\{&
 \int_M\operatorname{tr}I
 \frac{\sqrt{4\pi t}}{4}
 \int_\Sigma\operatorname{tr}\chi
 \notag\\
 &+t\left[
 \int_M\operatorname{tr}\left(E+\frac16R\,I\right)
 \int_\Sigma\operatorname{tr}
 \left(\frac13L\,I+2S\Pi_+\right)
 \right]
 +O(t^{3/2})\Bigg\}.
 \label{eq:V129-mixed-heat}
\end{align}
\end{theorem}

\begin{proof}
The interior coefficients are obtained by the normal-coordinate
parametrix for \eqref{eq:V129-Laplace-type}.  At the boundary, freeze the
tangential variables and use the reflected half-line kernels.  Their
constant image term gives
\(\sqrt{4\pi t}\operatorname{tr}\chi/4\).  The first normal variation of
the volume density gives \(L\operatorname{tr}I/3\), while one insertion of
the Robin endomorphism gives \(2\operatorname{tr}(S\Pi_+)\).
Terms involving derivatives of \(\Pi_\pm\) or \(S\), bundle curvature, or
quadratic geometry have weight at least three and enter
\(O(t^{3/2})\).  These local calculations patch by a partition of unity,
which proves \eqref{eq:V129-mixed-heat}.
\end{proof}

\begin{remark}[Normalization check]
\label{rem:V129-halfline-check}
For one scalar Dirichlet component, \(\chi=-1\).  The second term of
\eqref{eq:V129-mixed-heat} is then
\[
 -\frac14(4\pi t)^{-3/2}\operatorname{Vol}(\Sigma),
\]
the standard half-line image contribution.  This fixes the normalization
without an appeal to a symbol convention.
\end{remark}

\subsection{Relative scalar and one-form data}
\label{subsec:V129-relative-data}

\begin{lemma}[Scalar relative data]
\label{lem:V129-scalar-data}
For the relative scalar Laplacian,
\begin{equation}
 \Pi_-=I,\qquad \Pi_+=0,\qquad
 \operatorname{tr}\chi=-1,\qquad E_0=0.
 \label{eq:V129-scalar-data}
\end{equation}
Consequently
\begin{align}
 K_0^{\rm rel}(t)
 \sim(4\pi t)^{-2}\Bigg\{&
 \int_M1
 \frac{\sqrt{4\pi t}}4\int_\Sigma(-1)
 \notag\\
 &+t\left[
 \frac16\int_MR+\frac13\int_\Sigma L
 \right]+O(t^{3/2})\Bigg\}.
 \label{eq:V129-scalar-heat}
\end{align}
\end{lemma}

\begin{proof}
Relative degree-zero boundary conditions are Dirichlet.  Insert the
displayed projectors and \(E_0=0\) into
Theorem~\ref{thm:V129-mixed-heat}.
\end{proof}

\begin{lemma}[One-form relative Robin endomorphism]
\label{lem:V129-oneform-Robin}
Decompose a one-form at the boundary into tangential and inward-normal
parts,
\begin{equation}
 a=\alpha+\phi\,n^\flat .
 \label{eq:V129-form-split}
\end{equation}
The relative Hodge conditions are the mixed conditions
\begin{equation}
 \alpha=0,\qquad
 (\nabla_n+S_1)\phi=0,
 \qquad
 S_1=H=-L
 \label{eq:V129-relative-oneform-BC}
\end{equation}
on the one-dimensional normal block.  Hence
\begin{equation}
 \operatorname{rank}\Pi_-=3,\qquad
 \operatorname{rank}\Pi_+=1,\qquad
 \operatorname{tr}\chi=-2.
 \label{eq:V129-oneform-ranks}
\end{equation}
\end{lemma}

\begin{proof}
The first relative condition \(n^\flat\wedge a=0\) gives
\(\alpha=0\).  On that boundary row,
\[
 \delta a=-\nabla_n\phi-H\phi .
\]
The second relative condition \(\delta a=0\) is therefore
\((\nabla_n+H)\phi=0\).  Equation~\eqref{eq:V129-L-convention} gives
\(S_1=H=-L\).  The tangential and normal fibre dimensions in four
dimensions are three and one, proving
\eqref{eq:V129-oneform-ranks}.
\end{proof}

\begin{lemma}[One-form bulk endomorphism and traced coefficients]
\label{lem:V129-oneform-traces}
For the Hodge Laplacian on one-forms in the convention
\eqref{eq:V129-Laplace-type},
\begin{equation}
 E_1=-\operatorname{Ric},
 \qquad
 \operatorname{tr}_{T^*M}
 \left(E_1+\frac16R\,I_4\right)
 =-\frac13R.
 \label{eq:V129-oneform-bulk-trace}
\end{equation}
Its relative boundary weight-two trace is
\begin{equation}
 \operatorname{tr}_{T^*M}
 \left(\frac13L\,I_4+2S_1\Pi_+\right)
 =
 \frac43L-2L=-\frac23L.
 \label{eq:V129-oneform-boundary-trace}
\end{equation}
\end{lemma}

\begin{proof}
The Weitzenbock identity is
\(\Delta_1=\nabla^*\nabla+\operatorname{Ric}\).  Rewriting it as
\eqref{eq:V129-Laplace-type} gives \(E_1=-\operatorname{Ric}\).
Taking the four-dimensional fibre trace yields
\(-R+4R/6=-R/3\).  For the boundary row, insert
\(\operatorname{rank}T^*M=4\), the rank-one normal projector, and
\(S_1=-L\) from Lemma~\ref{lem:V129-oneform-Robin}.
\end{proof}

\begin{proposition}[Relative one-form baseline]
\label{prop:V129-oneform-heat}
The one-form heat trace is
\begin{align}
 K_1^{\rm rel}(t)
 \sim(4\pi t)^{-2}\Bigg\{&
 4\int_M1
 \frac{\sqrt{4\pi t}}4\int_\Sigma(-2)
 \notag\\
 &+t\left[
 -\frac13\int_MR-\frac23\int_\Sigma L
 \right]+O(t^{3/2})\Bigg\}.
 \label{eq:V129-oneform-heat}
\end{align}
\end{proposition}

\begin{proof}
Combine Lemmas~\ref{lem:V129-oneform-Robin} and
\ref{lem:V129-oneform-traces} with
Theorem~\ref{thm:V129-mixed-heat}.
\end{proof}

\subsection{BRST combination and exact/coexact interpretation}
\label{subsec:V129-BRST}

Per gauge generator define the physical relative heat combination
\begin{equation}
 K_{\rm gh+g}^{\rm phys}(t)
 :=
 K_0^{\rm rel}(t)-\frac12K_1^{\rm rel}(t).
 \label{eq:V129-BRST-combination}
\end{equation}

\begin{theorem}[Area cancellation and Einstein--GHY completion]
\label{thm:V129-Einstein-GHY}
For the twelve-dimensional Standard Model adjoint fibre,
\begin{align}
 12K_{\rm gh+g}^{\rm phys}(t)
 \sim(4\pi t)^{-2}\Bigg\{&
 -12\int_M1
 \notag\\
 &+t\left[
 4\int_MR+8\int_\Sigma L
 \right]
 +O(t^{3/2})\Bigg\}.
 \label{eq:V129-SM-heat}
\end{align}
In particular,
\begin{equation}
 [t^{1/2}]\,12K_{\rm gh+g}^{\rm phys}(t)=0
 \label{eq:V129-area-cancel}
\end{equation}
and the weight-two boundary coefficient is exactly twice the bulk
coefficient:
\begin{equation}
 4\left(\int_MR+2\int_\Sigma L\right)
 =
 4\left(\int_MR-2\int_\Sigma H\right).
 \label{eq:V129-EH-GHY-ratio}
\end{equation}
\end{theorem}

\begin{proof}
Subtract one half of \eqref{eq:V129-oneform-heat} from
\eqref{eq:V129-scalar-heat}.  The ranks give
\[
 1-\frac12(4)=-1.
\]
The boundary image ranks give
\[
 -1-\frac12(-2)=0,
\]
which proves \eqref{eq:V129-area-cancel}.  At weight two the bulk and
boundary coefficients per generator are
\[
 \frac16-\frac12\left(-\frac13\right)=\frac13,
 \qquad
 \frac13-\frac12\left(-\frac23\right)=\frac23.
\]
Multiplication by \(12\) proves \eqref{eq:V129-SM-heat}; substituting
\(L=-H\) proves \eqref{eq:V129-EH-GHY-ratio}.
\end{proof}

\begin{proposition}[Exact/coexact reduction]
\label{prop:V129-coexact-reduction}
Away from zero modes, \(d\) gives a unitary spectral equivalence between
the relative scalar Laplacian and the exact relative one-form Laplacian.
Thus
\begin{equation}
 12\left(K_0^{\rm rel}-\frac12K_1^{\rm rel}\right)
 =
 6\left(K_0^{\rm rel}-K_{1,\rm coex}^{\rm rel}\right)
 -6b_1^{\rm rel},
 \label{eq:V129-coexact-identity}
\end{equation}
where \(b_1^{\rm rel}=\dim\ker\Delta_{1,\rm rel}\).  The last term is
constant in \(t\) and does not alter any coefficient displayed in
\eqref{eq:V129-SM-heat}.
\end{proposition}

\begin{proof}
If \(\Delta_0f=\lambda f\) with \(\lambda>0\), then
\(\Delta_1(df)=d\Delta_0f=\lambda\,df\), and
\(\lambda^{-1/2}d\) is isometric because
\(\|df\|^2=\langle f,\Delta_0f\rangle\).  Conversely an exact
\(\lambda\)-eigenform has this form.  Orthogonal Hodge decomposition
splits the one-form heat trace into exact, coexact, and harmonic parts.
Substitution into the left side gives
\eqref{eq:V129-coexact-identity}.
\end{proof}

\subsection{Proper-time coefficient and its limitation}
\label{subsec:V129-proper-time}

\begin{corollary}[Gauge-sector proper-time divergence]
\label{cor:V129-proper-time}
Define the Euclidean gauge-plus-ghost determinant contribution by
\begin{equation}
 \Gamma_{\rm gh+g}^{\rm phys}(\Lambda)
 :=
 \int_{\Lambda^{-2}}^\infty
 \frac{dt}{t}\,
 12K_{\rm gh+g}^{\rm phys}(t),
 \label{eq:V129-proper-time-def}
\end{equation}
with any fixed infrared regularization understood.  Its local ultraviolet
part through quadratic order is
\begin{equation}
 \boxed{
 \Gamma_{\rm gh+g}^{\rm phys}(\Lambda)
 =
 \frac1{(4\pi)^2}
 \left[
 -6\Lambda^4\int_M1
 +4\Lambda^2
 \left(\int_MR+2\int_\Sigma L\right)
 \right]
 +O(\log\Lambda).}
 \label{eq:V129-proper-time-divergence}
\end{equation}
\end{corollary}

\begin{proof}
Insert \eqref{eq:V129-SM-heat} into
\eqref{eq:V129-proper-time-def}.  The required elementary integrals are
\[
 \int_{\Lambda^{-2}}\!t^{-3}\,dt=\frac12\Lambda^4+O(1),
 \qquad
 \int_{\Lambda^{-2}}\!t^{-2}\,dt=\Lambda^2+O(1).
\]
This proves \eqref{eq:V129-proper-time-divergence}.
\end{proof}

\begin{firewall}[This is a sector coefficient, not Newton's constant]
\label{fire:V129-not-Newton}
Equation~\eqref{eq:V129-proper-time-divergence} is the physical
relative-boundary gauge-plus-ghost baseline in one proper-time
regularization.  The V128 interface term, fermions, Higgs fields,
renormalization conditions, and the microscopic renewal normalization
have not yet been combined with it.  Therefore its coefficient is not
identified here with the observed Newton constant.  The exact conclusion
is the local Einstein--GHY ratio in
\eqref{eq:V129-EH-GHY-ratio}.
\end{firewall}

\subsection{V129 ledger and next acquisition}
\label{subsec:V129-ledger}

\paragraph{New exact conclusions.}
\begin{enumerate}[label=\textup{(\arabic*)},leftmargin=2.8em]
\item Lemma~\ref{lem:V129-oneform-Robin} derives the shape-dependent
      relative Robin endomorphism \(S_1=-L\).
\item Proposition~\ref{prop:V129-oneform-heat} gives the physical
      relative one-form heat baseline through weight two.
\item Theorem~\ref{thm:V129-Einstein-GHY} proves cancellation of the
      boundary area term in the gauge-plus-ghost combination.
\item The same theorem proves the exact \(1:2\) Einstein--Hilbert to
      Gibbons--Hawking--York ratio at weight two.
\item Proposition~\ref{prop:V129-coexact-reduction} reconciles the local
      calculation with the exact/coexact cancellation used in V123.
\end{enumerate}

\paragraph{Next forced step.}
V130 is the compulsory cross-program audit.  It must inspect the newest
Yang--Mills and Navier--Stokes notes in the shared folder, record their
identities and extract only transferable proof mechanisms.  It should
then combine the physical coefficient
\eqref{eq:V129-EH-GHY-ratio} with the V128 relative-interface row while
keeping the different cutoff degrees separate.  A direct numerical
addition of \(\Lambda^2\) and \(\Lambda\) coefficients is forbidden.

\begin{assessment}[V129 continuum status]
\label{ass:V129-continuum-status}
The full Standard--Model--Einstein continuum remains unproved.  V129
supplies the physical bosonic gauge/ghost heat baseline and its
Einstein--GHY ratio.  The interface normalization, fermion and Higgs
sectors, regulator matching, interacting nonlinear construction,
Einstein stress closure, and cofinal continuum still require separate
proofs.
\end{assessment}

\section*{V129 exact checks and append-only record}
\addcontentsline{toc}{section}{V129 exact checks and append-only record}

The physical-baseline checksum is
\begin{equation}
 \boxed{
 [t^{1/2}]\,12K_{\rm gh+g}^{\rm phys}=0,
 \qquad
 [t]\,12K_{\rm gh+g}^{\rm phys}
 =
 4\left(\int_MR+2\int_\Sigma L\right).}
 \label{eq:V129-final-check}
\end{equation}
The append operation preserves the complete V128 prefix byte for byte before
its former live terminator.  The assembled V129 file is checked for one live
final terminator, balanced environments and braces, duplicate labels, newly
introduced unresolved references, display math delimiters, and control
characters before the exact V128 predecessor is removed.

% =============================================================================
% END APPENDED VERSION V129 MATERIAL
% =============================================================================

% =============================================================================
% BEGIN APPENDED VERSION V130 MATERIAL
% =============================================================================

\section{V130: compulsory companion audit, Krein-to-heat transfer, and
the boundary-variation mismatch}
\label{sec:V130}

\subsection{Purpose}
\label{subsec:V130-purpose}

V129 produced the physical relative gauge/ghost heat coefficient
\[
 4\left(\int_MR+2\int_\Sigma L\right),
\]
while V125 produced a sharp tangential-cutoff trace of the logarithmic
interface symbol.  Those two numbers use different regulators and cannot
be added.  The compulsory V130 companion audit reinforces the relevant
rule: keep the enlarged object through the operation which creates the
cancellation or moment, and expose the exact unmarking factor before
estimating it.

Accordingly, V130 returns to the finite Krein identity of V123,
differentiates the complete logarithmic gluing ratio in the spectral
parameter, and only then takes the inverse Laplace transform.  This gives
the interface heat coefficient in the same convention as V129.  The
combined coefficient fails the Einstein--GHY ratio.  The failure is a
precise local boundary term, and requiring the standard Dirichlet
gravitational variation fixes its unique weight-one subtraction.

\subsection{Mandatory V130 companion audit}
\label{subsec:V130-audit}

The read-only shared-folder snapshot on 5 September 2026 was
\begin{center}
\begin{tabular}{|p{0.49\textwidth}|r|p{0.31\textwidth}|}
\hline
Companion file & Bytes & SHA--256\\
\hline
\texttt{Renewal\_NS\_Stability\_Research\_Notes\_V31.tex}
& \(887537\) &
\texttt{F1121B62238AD5491BB7CF03635EA9934942EDF573ED8FAAA9EE79E1D38A6696}\\
\hline
\texttt{Renewal\_Yang\_Mills\_Mass\_Gap\_Research\_Notes\_V115.tex}
& \(2710319\) &
\texttt{E9A32C083DF3C8D5A4BB5728360EC2FAB19EE31AD5D3ECC30508FE69D59E01E8}\\
\hline
\end{tabular}
\end{center}
The NS file is byte-identical to the V125 audit.  Its V31 heat-formation
ladder proves that flat dyadic marking is removed by an exact first
radius moment; heat lifting is scale-free only at the critical source
rung, and moving the missing factor into a weighted weak norm merely
hides the same continuation stock.  Its sharp one-mode test rules out an
improvement from heat smoothing alone.

The Yang--Mills file has advanced to V115.  Its new off-cut analysis
factors the full probability space into an independent \(bc\)-only bank
and its complement before taking norms.  The complete ternary cumulant
then peels exactly into a mean times a reduced ternary cumulant plus an
off-cut covariance times a reduced binary covariance.  Estimating the
five original moments separately would erase that cancellation and lose
the endpoint debit.

\begin{assessment}[Transferable audit rule]
\label{ass:V130-audit-rule}
For the present SM boundary problem, the analogue of the NS unmarking
factor is the spectral derivative in the Krein formula, and the analogue
of the YM independent-bank peeling is the finite determinant ratio before
logarithmic or heat regularization.  V130 therefore retains
\(\mathcal W(z)\), differentiates \(\log\det\mathcal W(z)\), and performs
the inverse Laplace transform before extracting a cutoff coefficient.
No NS or YM estimate is imported.
\end{assessment}

\subsection{The common spectral object}
\label{subsec:V130-common-object}

At a common finite regulator define the total Standard Model adjoint
interface logarithm
\begin{equation}
 \mathscr B_R(z)
 :=
 12\left[
 \log\det\mathcal W_{0,R}(z)
 -\frac12\log\det\mathcal W_{1,R}(z)
 \right].
 \label{eq:V130-interface-log}
\end{equation}
Let
\begin{equation}
 K_{\rm BRST}^{\rm rel}(t)
 :=
 12\left[
 K_{0,R}^{\rm rel}(t)
 -\frac12K_{1,R}^{\rm rel}(t)
 \right].
 \label{eq:V130-relative-heat}
\end{equation}

\begin{theorem}[Signed Krein-to-heat identity]
\label{thm:V130-Krein-heat}
Away from the finite harmonic sector,
\begin{equation}
 \boxed{
 \partial_z\mathscr B_R(z)
 =
 \int_0^\infty e^{-tz}
 K_{\rm BRST}^{\rm rel}(t)\,dt.}
 \label{eq:V130-Krein-Laplace}
\end{equation}
Consequently the relative effective determinant
\begin{equation}
 \Gamma_{\rm BRST}^{\rm rel}
 :=
 \int_0^\infty\frac{dt}{t}
 K_{\rm BRST}^{\rm rel}(t)
 \label{eq:V130-relative-effective}
\end{equation}
equals \(-\mathscr B_R(0)\), up to the common ultraviolet subtraction,
infrared normalization, and harmonic quotient.
\end{theorem}

\begin{proof}
For each form degree, Theorem~\ref{thm:V123-Krein} identifies
\(\partial_z\log\det\mathcal W_j(z)\) with the glued resolvent trace minus
the physical and auxiliary reference resolvent traces.  The resolvent is
the Laplace transform of the heat semigroup.  Take the linear combination
with coefficients \(12\) and \(-6\) to obtain
\eqref{eq:V130-Krein-Laplace}.  Integrating in \(z\), or equivalently using
\(\log\det D=-\int_0^\infty t^{-1}\operatorname{Tr}e^{-tD}\,dt\), gives
the final sign.  The finite regulator makes both operations elementary;
the listed subtractions are needed only when the regulator is removed.
\end{proof}

\begin{firewall}[Why the V125 sharp trace is not yet a heat coefficient]
\label{fire:V130-sharp-not-heat}
Corollary~\ref{cor:V125-H-divergence} traces
\(\mathscr B_R(0)\) with a hard cutoff in the three tangential covariables.
Theorem~\ref{thm:V130-Krein-heat} first differentiates in \(z\), thereby
changing \(\rho^{-1}\) to \(\rho^{-3}\), and then performs a
four-dimensional heat regularization.  The two procedures need not have
the same power-divergence coefficient.  Their discrepancy is a regulator
conversion, not a contradiction.
\end{firewall}

\subsection{Exact inverse Laplace transform of the mean-curvature row}
\label{subsec:V130-H-transfer}

V125 gives the homogeneous symbol
\begin{equation}
 \sigma_{-1}\mathscr B_R(z)
 =
 -6H(y)\bigl(|\xi|^2+z\bigr)^{-1/2}.
 \label{eq:V130-B-H-symbol}
\end{equation}
The exponentially small finite-\(R\) terms are smoothing and do not alter
this homogeneous row.

\begin{lemma}[Spectral derivative of the curvature symbol]
\label{lem:V130-H-derivative}
The corresponding symbol of the relative resolvent trace is
\begin{equation}
 \partial_z\sigma_{-1}\mathscr B_R(z)
 =
 3H(y)\bigl(|\xi|^2+z\bigr)^{-3/2}.
 \label{eq:V130-H-resolvent}
\end{equation}
\end{lemma}

\begin{proof}
Differentiate \eqref{eq:V130-B-H-symbol} and use
\(\partial_z(q+z)^{-1/2}=-\tfrac12(q+z)^{-3/2}\).
\end{proof}

\begin{lemma}[Inverse Laplace and tangential Gaussian integrals]
\label{lem:V130-Laplace-Gaussian}
For \(q\ge0\),
\begin{align}
 \mathcal L^{-1}_{z\to t}(q+z)^{-3/2}
 &=
 \frac{2\sqrt t}{\sqrt\pi}e^{-tq},
 \label{eq:V130-inverse-Laplace}\\
 \frac1{(2\pi)^3}\int_{\mathbb R^3}e^{-t|\xi|^2}\,d\xi
 &=(4\pi t)^{-3/2}.
 \label{eq:V130-Gaussian}
\end{align}
\end{lemma}

\begin{proof}
The Gamma integral gives
\[
 (q+z)^{-3/2}
 =
 \frac1{\Gamma(3/2)}
 \int_0^\infty t^{1/2}e^{-t(q+z)}\,dt,
\]
and \(\Gamma(3/2)=\sqrt\pi/2\).  The second identity is the product of
three one-dimensional Gaussian integrals with the Fourier normalization
\((2\pi)^{-3}\).
\end{proof}

\begin{theorem}[Relative interface heat coefficient at weight two]
\label{thm:V130-relative-H-heat}
The mean-curvature part of the relative Standard Model adjoint heat
combination is
\begin{equation}
 \boxed{
 K_{\rm BRST}^{\rm rel}(t)\big|_H
 =
 \frac{3}{4\pi^2}\,t^{-1}
 \int_\Sigma H\,d\operatorname{vol}_\Sigma
 =
 (4\pi t)^{-2}t\,
 12\int_\Sigma H\,d\operatorname{vol}_\Sigma.}
 \label{eq:V130-relative-H-heat}
\end{equation}
Equivalently, in the V129 inward heat convention, its weight-two
boundary coefficient is
\begin{equation}
 A_{2,\partial}^{\rm rel}
 =
 -12\int_\Sigma L\,d\operatorname{vol}_\Sigma.
 \label{eq:V130-relative-L-coefficient}
\end{equation}
\end{theorem}

\begin{proof}
Apply Lemma~\ref{lem:V130-Laplace-Gaussian} to
\eqref{eq:V130-H-resolvent}, trace the tangential symbol, and integrate
over \(\Sigma\):
\begin{align*}
 &3H\,\frac{2\sqrt t}{\sqrt\pi}
 (4\pi t)^{-3/2}
 =
 \frac{3H}{4\pi^2}t^{-1}.
\end{align*}
Since
\((4\pi t)^{-2}t=(16\pi^2)^{-1}t^{-1}\), the bracket coefficient is
\(12H\).  Finally use \(L=-H\).
\end{proof}

\begin{corollary}[Proper-time conversion of the V125 sharp coefficient]
\label{cor:V130-factor-two-conversion}
With the lower proper-time cutoff \(t\ge\Lambda^{-2}\), the relative
effective action has the quadratic divergence
\begin{equation}
 \Gamma_{\rm BRST}^{\rm rel}(\Lambda)\big|_H
 =
 \frac{3\Lambda^2}{4\pi^2}
 \int_\Sigma H\,d\operatorname{vol}_\Sigma
 +O(\log\Lambda).
 \label{eq:V130-relative-effective-H}
\end{equation}
This has the effective-action sign opposite to
\(\mathscr B_R\) and one half the magnitude of the hard tangential trace
in Corollary~\ref{cor:V125-H-divergence}.
\end{corollary}

\begin{proof}
Insert \eqref{eq:V130-relative-H-heat} into
\eqref{eq:V130-relative-effective} and use
\(\int_{\Lambda^{-2}}t^{-2}\,dt=\Lambda^2+O(1)\).
The comparison with V125 is immediate.
\end{proof}

\subsection{The combined boundary coefficient}
\label{subsec:V130-combined}

Eliminating the auxiliary collar means taking the glued determinant
relative to the auxiliary determinant.  At the heat level this is the
sum of the physical relative trace and the relative gluing trace:
\begin{equation}
 K_{\rm eff}=K_{\rm phys}^{\rm rel}+K_{\rm BRST}^{\rm rel}.
 \label{eq:V130-effective-sum}
\end{equation}

\begin{theorem}[Exact weight-two coefficient after auxiliary elimination]
\label{thm:V130-combined-A2}
The Standard Model adjoint gauge/ghost sector constructed in V122--V129
has, after auxiliary elimination,
\begin{equation}
 \boxed{
 A_{2}^{\rm eff}
 =
 4\int_MR
 -4\int_\Sigma L
 =
 4\left(\int_MR-\int_\Sigma L\right).}
 \label{eq:V130-combined-A2}
\end{equation}
It therefore does not have the Dirichlet gravitational completion
\begin{equation}
 A_2^{\rm EH+GHY}
 =
 4\left(\int_MR+2\int_\Sigma L\right).
 \label{eq:V130-target-A2}
\end{equation}
The defect is
\begin{equation}
 \boxed{
 A_2^{\rm EH+GHY}-A_2^{\rm eff}
 =
 12\int_\Sigma L
 =
 -12\int_\Sigma H.}
 \label{eq:V130-A2-defect}
\end{equation}
\end{theorem}

\begin{proof}
Theorem~\ref{thm:V129-Einstein-GHY} gives the physical coefficient
\[
 A_2^{\rm phys}=4\int_MR+8\int_\Sigma L.
\]
Theorem~\ref{thm:V130-relative-H-heat} gives
\(A_{2,\partial}^{\rm rel}=-12\int_\Sigma L\), with no bulk term because
the relative heat trace is interface-supported.  Addition gives
\eqref{eq:V130-combined-A2}.  Subtract it from
\eqref{eq:V130-target-A2} to obtain \eqref{eq:V130-A2-defect}.
\end{proof}

\begin{corollary}[The defect cannot be tuned by collar length]
\label{cor:V130-no-R-tuning}
No choice of finite auxiliary length \(R>0\) changes
\eqref{eq:V130-A2-defect}.
\end{corollary}

\begin{proof}
By Proposition~\ref{prop:V128-product-operators}, all finite-\(R\)
corrections to the normalized high-frequency product operator are
exponentially small in \(R\rho\) and hence smoothing.  Local homogeneous
heat coefficients depend only on the finite symbol jet.  Therefore the
order-\(-1\) symbol \eqref{eq:V130-B-H-symbol}, and consequently
\eqref{eq:V130-A2-defect}, are \(R\)-independent.
\end{proof}

\begin{theorem}[Unique local weight-one variational subtraction]
\label{thm:V130-unique-counterterm}
Fix the bulk coefficient \(4\int_MR\), require locality, boundary
diffeomorphism invariance, and the Dirichlet Einstein variational ratio.
Then the unique smooth-boundary weight-one correction is
\begin{equation}
 \boxed{
 A_{2,\partial}^{\rm ct}
 =
 12\int_\Sigma L
 =
 -12\int_\Sigma H.}
 \label{eq:V130-counterterm-A2}
\end{equation}
For the lower proper-time cutoff its divergent action coefficient is
\begin{equation}
 \boxed{
 \Gamma_{\partial}^{\rm ct}(\Lambda)
 =
 \frac{3\Lambda^2}{4\pi^2}
 \int_\Sigma L\,d\operatorname{vol}_\Sigma
 =
 -\frac{3\Lambda^2}{4\pi^2}
 \int_\Sigma H\,d\operatorname{vol}_\Sigma.}
 \label{eq:V130-counterterm-action}
\end{equation}
\end{theorem}

\begin{proof}
On a smooth boundary the only local scalar of geometric weight one formed
from the metric jet is a constant multiple of the mean curvature \(L\);
tangential divergences have higher weight or integrate to zero.  The
required multiple is forced by the exact defect
\eqref{eq:V130-A2-defect}.  Multiplication by
\(\Lambda^2/(4\pi)^2\), as in V129, gives
\eqref{eq:V130-counterterm-action}.
\end{proof}

\begin{firewall}[Renormalization condition, not a derived microscopic
coupling]
\label{fire:V130-counterterm-scope}
Theorem~\ref{thm:V130-unique-counterterm} proves uniqueness only after the
Dirichlet Einstein variational ratio is imposed as a renormalization
condition.  It does not prove that the renewal microscopic measure
generates that subtraction automatically.  Without the subtraction, the
enlarged-collar gauge/ghost determinant has the ratio
\eqref{eq:V130-combined-A2}; with it, the weight-two metric variation is
Einstein--GHY compatible.  Determining which alternative is selected by
the microscopic continuum construction is now an explicit upstream
gate.
\end{firewall}

\subsection{Disposition of the V128 curvature-square row}
\label{subsec:V130-order-three}

\begin{assessment}[Do not mix cutoff degrees]
\label{ass:V130-do-not-mix}
The V128 interface term
\[
 \frac{\Lambda}{\pi^2}
 \int_\Sigma
 \left(\frac7{10}H^2-\frac85|K|^2\right)
\]
was also printed with a hard tangential cutoff and belongs to geometric
weight three in the four-dimensional heat ledger.  It must undergo its
own Krein derivative and inverse Laplace transform before it is combined
with the \(t^{3/2}\) physical relative Hodge coefficient.  It cannot
modify the weight-two defect \eqref{eq:V130-A2-defect}, and its coefficient
cannot be added to a \(\Lambda^2\) term.
\end{assessment}

\subsection{V130 ledger and next acquisition}
\label{subsec:V130-ledger}

\paragraph{New exact conclusions.}
\begin{enumerate}[label=\textup{(\arabic*)},leftmargin=2.8em]
\item The compulsory audit records the unchanged NS V31 note and the
      current YM V115 note.
\item Theorem~\ref{thm:V130-Krein-heat} fixes the determinant sign and the
      required order of spectral derivative and inverse Laplace transform.
\item Theorem~\ref{thm:V130-relative-H-heat} converts the V125
      order-\(-1\) logarithmic symbol into the common proper-time heat
      coefficient \(12\int_\Sigma H\).
\item Theorem~\ref{thm:V130-combined-A2} proves that auxiliary elimination
      changes the physical Einstein--GHY ratio to
      \(4(\int R-\int L)\).
\item Theorem~\ref{thm:V130-unique-counterterm} identifies the unique local
      weight-one subtraction which restores the Dirichlet gravitational
      ratio.
\end{enumerate}

\paragraph{Next forced step.}
V131 must repeat the Krein-to-heat transfer for the full directional
order-\(-2\) symbol of V128.  Unlike the mean-curvature row,
\(K(n,n)^2\) contains powers of
\(\xi/\sqrt{|\xi|^2+z}\); differentiating in \(z\) also differentiates
those angular factors.  The inverse Laplace transform must therefore be
performed before the angular average.  The result should be compared
with the physical mixed-boundary \(A_3\) coefficient, with the V130
weight-two counterterm held fixed.

\begin{assessment}[V130 continuum status]
\label{ass:V130-continuum-status}
The full Standard--Model--Einstein continuum remains unproved.  V130
finds and quantifies a boundary-variation obstruction in the
collar-enlarged gauge/ghost determinant and gives its unique local
Einstein--GHY subtraction.  Microscopic selection of that subtraction,
the weight-three boundary coefficient, fermion and Higgs sectors,
interacting nonlinear existence, Einstein stress closure, and the
cofinal continuum remain open.
\end{assessment}

\section*{V130 exact checks and append-only record}
\addcontentsline{toc}{section}{V130 exact checks and append-only record}

The common-regulator checksum is
\begin{equation}
 \boxed{
 A_{2,\partial}^{\rm rel}=-12\int_\Sigma L,
 \qquad
 A_2^{\rm eff}=4\left(\int_MR-\int_\Sigma L\right),
 \qquad
 A_2^{\rm ct}=12\int_\Sigma L.}
 \label{eq:V130-final-check}
\end{equation}
The append operation preserves the complete V129 prefix byte for byte before
its former live terminator.  The assembled V130 file is checked for one live
final terminator, balanced environments and braces, duplicate labels, newly
introduced unresolved references, display math delimiters, and control
characters before the exact V129 predecessor is removed.

% =============================================================================
% END APPENDED VERSION V130 MATERIAL
% =============================================================================

% =============================================================================
% BEGIN APPENDED VERSION V131 MATERIAL
% =============================================================================

\section{V131: joint parameter lift and the quadratic interface heat
coefficient}
\label{sec:V131}

\subsection{Purpose}
\label{subsec:V131-purpose}

V130 showed that a symbol traced at \(z=0\) with a hard tangential cutoff
is not itself the proper-time heat coefficient.  For the order-\(-1\)
mean-curvature row this conversion was scalar.  The V128 order-\(-2\)
row is subtler because its factor \(K(n,n)^2\) also depends on the joint
parameter direction.  V131 reconstructs that dependence from the
Riccati equation, differentiates before averaging, and computes the exact
Gaussian fourth moment.

The result is the weight-three \((H^2,|K|^2)\) coefficient of the
\emph{relative interface heat trace}.  The physical mixed-boundary
coefficient at the same weight remains separate and is the next
acquisition.

\subsection{Joint parameter form of the first curvature root}
\label{subsec:V131-joint-root}

Put
\begin{equation}
 \lambda=\lambda(\xi,z):=(|\xi|^2+z)^{1/2},
 \qquad
 v:=\frac{\xi}{\lambda},
 \qquad
 k_z:=K(v,v)
 =\frac{K^{ij}\xi_i\xi_j}{|\xi|^2+z}.
 \label{eq:V131-joint-variables}
\end{equation}
For \(z>0\), \(v\) lies in the open unit co-ball; at \(z=0\) it is the
unit direction \(n=\xi/|\xi|\).

\begin{lemma}[Joint scalar curvature root]
\label{lem:V131-scalar-root}
The scalar raw-density order-zero root for
\(\Delta_0+z\) is
\begin{equation}
 s_0(z)=\frac12(k_z-H).
 \label{eq:V131-s0}
\end{equation}
\end{lemma}

\begin{proof}
The principal root obeys
\(\lambda^2=\gamma^{ij}\xi_i\xi_j+z\).  Since \(z\) is independent of the
normal coordinate,
\[
 \partial_r\lambda
 =-\frac{K^{ij}\xi_i\xi_j}{\lambda}.
\]
The first scalar Sylvester correction in fixed boundary measure is
\(-(\partial_r\lambda)/(2\lambda)=k_z/2\), exactly as in V117.  Returning
to the raw boundary density subtracts \(H/2\).  This proves
\eqref{eq:V131-s0}.
\end{proof}

\begin{lemma}[Joint one-form diagonal blocks]
\label{lem:V131-oneform-root}
In the tangential/normal splitting, the diagonal order-zero blocks for
\(\Delta_1+z\) are
\begin{equation}
 M_T(z)=s_0(z)I_3+S,
 \qquad
 m_N(z)=s_0(z),
 \label{eq:V131-oneform-diagonal}
\end{equation}
where \(S=K^\sharp\) is the tangential shape endomorphism.  The
off-diagonal pair is linear in \(Sv\) and cancels from the jet-matched
relative logarithm before the quadratic expansion.
\end{lemma}

\begin{proof}
Adding \(zI\) changes only the scalar principal root from \(|\xi|\) to
\(\lambda\).  The shape endomorphism comes from the covariant
tangential/normal decomposition and is independent of \(z\).  The
normal derivative part is therefore the scalar
\eqref{eq:V131-s0}, while the tangential connection adds \(S\), giving
\eqref{eq:V131-oneform-diagonal}.  Reflection reverses the full
order-zero root and the two one-sided off-diagonal blocks cancel in the
glued numerator, as in
Remark~\ref{rem:V126-offdiagonal-disposition}.  They are absent from the
relative/absolute reference diagonal blocks.
\end{proof}

\begin{theorem}[Joint parameter lift of the closed quadratic row]
\label{thm:V131-joint-B2}
The Standard Model adjoint order-\(-2\) algebraic curvature-square symbol
of \(\mathscr B_R(z)\) is
\begin{align}
 \sigma_{-2}^{\rm quad}\mathscr B_R(z)
 &=
 \frac{12}{\lambda^2}
 \left[
 \frac{H^2-k_z^2}{8}
 -\frac14|K|^2
 \right]
 \notag\\
 &=
 \left(\frac32H^2-3|K|^2\right)
 (|\xi|^2+z)^{-1}
 \notag\\
 &\qquad
 -\frac32(K^{ij}\xi_i\xi_j)^2
 (|\xi|^2+z)^{-3}.
 \label{eq:V131-joint-B2}
\end{align}
At \(z=0\), angular averaging recovers
Corollary~\ref{cor:V126-angular-Q} and
Theorem~\ref{thm:V128-order-two-closure}.
\end{theorem}

\begin{proof}
The logarithm calculation in
Theorem~\ref{thm:V126-directional-quadratic} is algebraic in
\(s_0\), \(S\), and the principal root.  Substitute
\eqref{eq:V131-s0}--\eqref{eq:V131-oneform-diagonal}, replace
\(\rho\) by \(\lambda\), and multiply by the adjoint dimension \(12\).
This gives the first line.  Since
\[
 \frac{k_z^2}{\lambda^2}
 =
 (K^{ij}\xi_i\xi_j)^2(|\xi|^2+z)^{-3},
\]
expansion gives the final two lines.  At \(z=0\), \(v=n\); the stated
specialization then follows from the fourth angular moment in V126.
\end{proof}

\begin{firewall}[The joint lift is not obtained by freezing \(n\)]
\label{fire:V131-no-frozen-direction}
Replacing \(k_z\) by \(K(\xi/|\xi|,\xi/|\xi|)\) before differentiating in
\(z\) omits the derivative of \(\lambda^{-4}\) inside \(k_z^2\).  That
missing term is exactly the Gaussian fourth-moment contribution below.
The \(z=0\) angular average alone does not determine the heat coefficient.
\end{firewall}

\subsection{Spectral derivative before angular integration}
\label{subsec:V131-derivative}

\begin{proposition}[Quadratic relative resolvent symbol]
\label{prop:V131-resolvent-symbol}
Differentiating the joint logarithm gives
\begin{align}
 \partial_z\sigma_{-2}^{\rm quad}\mathscr B_R(z)
 &=
 -\left(\frac32H^2-3|K|^2\right)
 (|\xi|^2+z)^{-2}
 \notag\\
 &\qquad
 +\frac92(K^{ij}\xi_i\xi_j)^2
 (|\xi|^2+z)^{-4}.
 \label{eq:V131-resolvent-symbol}
\end{align}
\end{proposition}

\begin{proof}
Differentiate the last two lines of
\eqref{eq:V131-joint-B2}.  The powers contribute \(-1\) and \(-3\);
the latter multiplies the existing coefficient \(-3/2\).
\end{proof}

\begin{lemma}[Fourth Gaussian contraction]
\label{lem:V131-Gaussian-fourth}
In an orthonormal tangential frame,
\begin{equation}
 \frac1{(2\pi)^3}
 \int_{\mathbb R^3}
 (K^{ij}\xi_i\xi_j)^2e^{-t|\xi|^2}\,d\xi
 =
 (4\pi t)^{-3/2}
 \frac{H^2+2|K|^2}{4t^2}.
 \label{eq:V131-Gaussian-fourth}
\end{equation}
\end{lemma}

\begin{proof}
The normalized centered Gaussian with density proportional to
\(e^{-t|\xi|^2}\) has covariance
\(\mathbb E(\xi_i\xi_j)=\delta_{ij}/(2t)\).  Wick contraction gives
\[
 \mathbb E(\xi_i\xi_j\xi_k\xi_l)
 =
 \frac1{4t^2}
 (\delta_{ij}\delta_{kl}
  \delta_{ik}\delta_{jl}
  \delta_{il}\delta_{jk}).
\]
Contract with \(K^{ij}K^{kl}\) and multiply by the Gaussian mass
\((4\pi t)^{-3/2}\).
\end{proof}

\begin{lemma}[Two inverse Laplace rows]
\label{lem:V131-two-Laplace}
For \(q\ge0\),
\begin{align}
 \mathcal L^{-1}_{z\to t}(q+z)^{-2}
 &=t e^{-tq},
 \label{eq:V131-Laplace-two}\\
 \mathcal L^{-1}_{z\to t}(q+z)^{-4}
 &=\frac{t^3}{6}e^{-tq}.
 \label{eq:V131-Laplace-four}
\end{align}
\end{lemma}

\begin{proof}
Use
\[
 (q+z)^{-m}
 =
 \frac1{\Gamma(m)}
 \int_0^\infty t^{m-1}e^{-t(q+z)}\,dt
\]
for \(m=2,4\).
\end{proof}

\begin{theorem}[Relative curvature-square heat coefficient]
\label{thm:V131-relative-A3}
The complete integrated algebraic curvature-square contribution of the
relative Standard Model adjoint interface heat trace is
\begin{align}
 K_{\rm BRST}^{\rm rel}(t)\big|_{K^2}
 &=
 t(4\pi t)^{-3/2}
 \int_\Sigma
 \left(
 -\frac{21}{16}H^2
 +\frac{27}{8}|K|^2
 \right)d\operatorname{vol}_\Sigma
 \notag\\
 &=
 (4\pi t)^{-2}t^{3/2}
 \int_\Sigma
 \sqrt\pi
 \left(
 -\frac{21}{8}H^2
 +\frac{27}{4}|K|^2
 \right)d\operatorname{vol}_\Sigma.
 \label{eq:V131-relative-A3}
\end{align}
\end{theorem}

\begin{proof}
Apply Lemma~\ref{lem:V131-two-Laplace} to
\eqref{eq:V131-resolvent-symbol}.  The first row becomes
\[
 -t\left(\frac32H^2-3|K|^2\right)e^{-t|\xi|^2}.
\]
The second row becomes
\[
 \frac34t^3(K^{ij}\xi_i\xi_j)^2e^{-t|\xi|^2}.
\]
Use \eqref{eq:V130-Gaussian} for the first Gaussian integral and
Lemma~\ref{lem:V131-Gaussian-fourth} for the second.  After factoring
\(t(4\pi t)^{-3/2}\), the coefficients are
\begin{align*}
 H^2:\quad&-\frac32+\frac3{16}=-\frac{21}{16},\\
 |K|^2:\quad&3+\frac38=\frac{27}{8}.
 \end{align*}
Finally
\[
 t(4\pi t)^{-3/2}
 =
 2\sqrt\pi\,(4\pi t)^{-2}t^{3/2},
\]
which proves the second line.
\end{proof}

\begin{corollary}[Proper-time linear divergence]
\label{cor:V131-proper-time-linear}
With the common lower proper-time cutoff \(t\ge\Lambda^{-2}\), the
relative effective action has
\begin{equation}
 \boxed{
 \Gamma_{\rm BRST}^{\rm rel}(\Lambda)\big|_{K^2}
 =
 \frac{\Lambda}{64\pi^{3/2}}
 \int_\Sigma
 \left(
 -21H^2+54|K|^2
 \right)d\operatorname{vol}_\Sigma
 +O(\log\Lambda).}
 \label{eq:V131-proper-time-linear}
\end{equation}
\end{corollary}

\begin{proof}
The first line of \eqref{eq:V131-relative-A3} is proportional to
\(t^{-1/2}\).  Divide by \(t\) in the effective-action integral and use
\[
 \int_{\Lambda^{-2}}^1t^{-3/2}\,dt=2\Lambda+O(1).
\]
Substitution and simplification give
\eqref{eq:V131-proper-time-linear}.
\end{proof}

\begin{corollary}[Explicit regulator conversion]
\label{cor:V131-regulator-conversion}
The hard tangential trace in
Corollary~\ref{cor:V126-linear-divergence} and the proper-time coefficient
\eqref{eq:V131-proper-time-linear} are different linear combinations of
\(H^2\) and \(|K|^2\).  In particular, the conversion is not a common
scalar rescaling.
\end{corollary}

\begin{proof}
The hard trace is proportional to
\[
 \frac7{10}H^2-\frac85|K|^2,
\]
whereas the proper-time trace is proportional to
\(-21H^2+54|K|^2\).  Their coefficient ratios are unequal.  The reason is
the differentiated \(k_z^2\) factor isolated in
Firewall~\ref{fire:V131-no-frozen-direction}.
\end{proof}

\begin{firewall}[Scope of the word complete]
\label{fire:V131-complete-scope}
Theorem~\ref{thm:V131-relative-A3} is complete for the
\((H^2,|K|^2)\) sector of the jet-matched relative interface symbol
closed in V128.  It does not supply the physical mixed-boundary
\(A_3\) coefficient, corner terms, or a microscopic boundary
renormalization.  The V130 weight-two counterterm is held fixed and
cannot be used to alter this weight-three row.
\end{firewall}

\subsection{V131 ledger and next acquisition}
\label{subsec:V131-ledger}

\paragraph{New exact conclusions.}
\begin{enumerate}[label=\textup{(\arabic*)},leftmargin=2.8em]
\item Lemmas~\ref{lem:V131-scalar-root} and
      \ref{lem:V131-oneform-root} reconstruct the joint
      \((\xi,z^{1/2})\) curvature root.
\item Theorem~\ref{thm:V131-joint-B2} gives the joint parameter lift of
      the V128 quadratic logarithmic symbol.
\item Proposition~\ref{prop:V131-resolvent-symbol} differentiates the
      directional factor before angular integration.
\item Lemma~\ref{lem:V131-Gaussian-fourth} performs the required fourth
      Gaussian contraction.
\item Theorem~\ref{thm:V131-relative-A3} computes the common proper-time
      relative interface \(A_3\) curvature-square coefficient.
\item Corollary~\ref{cor:V131-regulator-conversion} proves that the hard
      tangential and proper-time conversions are not related by one
      scalar moment.
\end{enumerate}

\paragraph{Next forced step.}
V132 must compute the physical relative mixed-boundary \(A_3\)
coefficient for the scalar and one-form Hodge Laplacians in the exact
V129 inward-normal convention.  It must retain the Robin endomorphism
\(S_1=-L\), the tangential/normal projectors, and the ambient curvature
terms.  Only after forming
\(12(A_{3,0}^{\rm rel}-\tfrac12A_{3,1}^{\rm rel})\) may it add
\eqref{eq:V131-relative-A3}.  The comparison must distinguish intrinsic
\(R_\Sigma\), ambient normal curvature, \(H^2\), and \(|K|^2\).

\begin{assessment}[V131 continuum status]
\label{ass:V131-continuum-status}
The full Standard--Model--Einstein continuum remains unproved.  V131
closes the proper-time transfer of the relative quadratic interface row.
The physical weight-three baseline, microscopic boundary
renormalization, fermion and Higgs sectors, interacting nonlinear
existence, Einstein stress closure, and the cofinal continuum remain
open.
\end{assessment}

\section*{V131 exact checks and append-only record}
\addcontentsline{toc}{section}{V131 exact checks and append-only record}

The joint-parameter checksum is
\begin{equation}
 \boxed{
 A_{3,\partial}^{\rm rel}\big|_{K^2}
 =
 \sqrt\pi\int_\Sigma
 \left(-\frac{21}{8}H^2+\frac{27}{4}|K|^2\right),
 \qquad
 \Gamma_{\rm rel}\big|_{K^2}
 =
 \frac{\Lambda}{64\pi^{3/2}}
 \int_\Sigma(-21H^2+54|K|^2).}
 \label{eq:V131-final-check}
\end{equation}
The append operation preserves the complete V130 prefix byte for byte before
its former live terminator.  The assembled V131 file is checked for one live
final terminator, balanced environments and braces, duplicate labels, newly
introduced unresolved references, display math delimiters, and control
characters before the exact V130 predecessor is removed.

% =============================================================================
% END APPENDED VERSION V131 MATERIAL
% =============================================================================

% =============================================================================
% BEGIN APPENDED VERSION V132 MATERIAL
% =============================================================================

\section{V132: the physical mixed-boundary \(A_3\) coefficient and the
completed weight-three gauge/ghost row}
\label{sec:V132}

\subsection{Purpose and coefficient convention}
\label{subsec:V132-purpose}

V131 computed the relative-interface curvature-square contribution at
heat weight three.  V132 now computes the physical relative scalar and
one-form coefficient at the same weight.  The calculation uses the
standard universal mixed-boundary formula, including the tangential
derivative of the normal/tangential projector.  That derivative is
nonzero even though the projector is extended parallel to the normal
geodesics.

We retain the V129 convention
\[
 D=-\bigl(g^{\mu\nu}\nabla_\mu\nabla_\nu+E\bigr),
 \qquad
 \Pi_-u=0,\qquad
 \Pi_+(\nabla_n+S)u=0,
\]
where \(n\) is inward and
\[
 L_{ab}=\langle\nabla_{e_a}e_b,n\rangle=-K_{ab},
 \qquad L=-H.
\]
Write
\begin{equation}
 R_{nn}:=\sum_{a=1}^3R(e_a,n,e_a,n)=\operatorname{Ric}(n,n).
 \label{eq:V132-Rnn}
\end{equation}
Our heat-bracket normalization is
\begin{equation}
 \operatorname{Tr}e^{-tD}
 \sim
 (4\pi t)^{-2}
 \left[
 \cdots+t^{3/2}A_3(D,\mathcal B)+O(t^2)
 \right].
 \label{eq:V132-A3-normalization}
\end{equation}

\paragraph{Formula provenance.}
The universal coefficient below is the mixed-boundary \(a_3\) formula of
Branson--Gilkey, with the corrected convention collected in
D.~V.~Vassilevich, \emph{Heat kernel expansion: user's manual},
Physics Reports 388 (2003), Section 5.3.  It is restated here in the
normalization and notation used by this programme.

\subsection{Universal mixed coefficient}
\label{subsec:V132-universal}

\begin{theorem}[Universal unsmeared mixed-boundary \(A_3\)]
\label{thm:V132-universal-A3}
Let \(\chi=\Pi_+-\Pi_-\), extend the projectors parallel to the inward
normal, and let a colon denote the induced tangential covariant
derivative.  Then
\begin{align}
 A_3(D,\mathcal B)
 =\frac{\sqrt\pi}{192}\int_\Sigma
 \operatorname{tr}\Big[&
 96\chi E+16\chi R+8\chi R_{nn}
 \notag\\
 &+(13\Pi_+-7\Pi_-)L^2
 +(2\Pi_++10\Pi_-)|L|^2
 \notag\\
 &+96SL+192S^2
 -12\chi_{:a}\chi_{:a}
 \Big]d\operatorname{vol}_\Sigma .
 \label{eq:V132-universal-A3}
\end{align}
\end{theorem}

\begin{proof}
The local boundary parametrix at geometric weight two has the invariant
basis
\[
 \chi E,\ \chi R,\ \chi R_{nn},\
 \Pi_\pm L^2,\ \Pi_\pm|L|^2,\
 SL,\ S^2,\ \chi_{:a}\chi_{:a}.
\]
The half-line recursion and functorial product identities fix their
coefficients respectively as
\[
 96,\ 16,\ 8,\ (13,-7),\ (2,10),\
 96,\ 192,\ -12
\]
in the standard coefficient
\((4\pi)^{-3/2}/384\).
Converting from the standard expansion
\(\sum_k t^{(k-4)/2}a_k\) to
\eqref{eq:V132-A3-normalization} multiplies the boundary integrand by
\[
 (4\pi)^2\frac{(4\pi)^{-3/2}}{384}
 =\frac{\sqrt\pi}{192}.
\]
This gives \eqref{eq:V132-universal-A3}.
\end{proof}

\begin{firewall}[Projector derivatives cannot be discarded]
\label{fire:V132-projector-derivative}
Normal parallel extension gives \(\nabla_n\chi=0\), but it does not give
\(\chi_{:a}=0\).  The normal line rotates under tangential transport by
the second fundamental form.  Omitting the last term of
\eqref{eq:V132-universal-A3} changes the one-form
\(|L|^2\) coefficient and destroys the BRST checksum below.
\end{firewall}

\subsection{Relative scalar coefficient}
\label{subsec:V132-scalar}

\begin{proposition}[Scalar relative \(A_3\)]
\label{prop:V132-scalar-A3}
For the relative scalar Laplacian,
\begin{equation}
 \boxed{
 A_{3,0}^{\rm rel}
 =
 \frac{\sqrt\pi}{192}\int_\Sigma
 \left(
 -16R-8R_{nn}-7L^2+10|L|^2
 \right)d\operatorname{vol}_\Sigma.}
 \label{eq:V132-scalar-A3}
\end{equation}
\end{proposition}

\begin{proof}
For a relative scalar,
\[
 \Pi_-=I,\qquad \Pi_+=0,\qquad
 \chi=-I,\qquad S=0,\qquad E_0=0,
\qquad \chi_{:a}=0.
\]
Insert these values into
\eqref{eq:V132-universal-A3}.
\end{proof}

\subsection{Relative one-form traces}
\label{subsec:V132-oneform}

Let \(P_N=n^\flat\otimes n\) and \(P_T=I-P_N\) on
\(T^*M|_\Sigma\).  Relative one-form conditions have
\begin{equation}
 \Pi_+=P_N,\qquad
 \Pi_-=P_T,\qquad
 \chi=P_N-P_T=2P_N-I,\qquad
 S=-L\,P_N.
 \label{eq:V132-oneform-data}
\end{equation}

\begin{lemma}[Projector derivative trace]
\label{lem:V132-projector-trace}
For the one-form projectors in \eqref{eq:V132-oneform-data},
\begin{equation}
 \boxed{
 \operatorname{tr}_{T^*M}
 (\chi_{:a}\chi_{:a})
 =8|L|^2.}
 \label{eq:V132-chi-derivative}
\end{equation}
\end{lemma}

\begin{proof}
Metric compatibility and the definition of \(L\) give
\[
 \nabla_{e_a}n=-L_{ab}e_b.
\]
Put \(u_a:=\nabla_{e_a}n\).  Since
\(P_N=n^\flat\otimes n\),
\[
 (P_N)_{:a}=u_a^\flat\otimes n+n^\flat\otimes u_a,
 \qquad
 \chi_{:a}=2(P_N)_{:a}.
\]
For each \(a\), the square of the displayed off-diagonal rank-two
endomorphism has trace \(2|u_a|^2\).  Therefore
\[
 \operatorname{tr}(\chi_{:a}\chi_{:a})
 =4\sum_a2|u_a|^2
 =8\sum_{a,b}L_{ab}^2
 =8|L|^2.
\]
\end{proof}

\begin{lemma}[One-form curvature traces]
\label{lem:V132-oneform-curvature}
With \(E_1=-\operatorname{Ric}\),
\begin{align}
 \operatorname{tr}\chi&=-2,
 \label{eq:V132-tr-chi}\\
 \operatorname{tr}(\chi E_1)&=R-2R_{nn},
 \label{eq:V132-tr-chi-E}\\
 \operatorname{tr}(\chi R)&=-2R,
 \qquad
 \operatorname{tr}(\chi R_{nn})=-2R_{nn}.
 \label{eq:V132-tr-chi-curvature}
\end{align}
\end{lemma}

\begin{proof}
The normal line has \(\chi=+1\), while the three tangential covectors
have \(\chi=-1\), giving \eqref{eq:V132-tr-chi}.  Since
\(E_1=-\operatorname{Ric}\),
\[
 \operatorname{tr}(\chi E_1)
 =
 -R_{nn}
 +\sum_{a=1}^3\operatorname{Ric}(e_a,e_a)
 =R-2R_{nn}.
\]
The remaining two identities follow because \(R\) and \(R_{nn}\) act as
scalar endomorphisms.
\end{proof}

\begin{lemma}[One-form shape traces]
\label{lem:V132-oneform-shape}
The shape and Robin rows in
\eqref{eq:V132-universal-A3} have traces
\begin{align}
 \operatorname{tr}(13\Pi_+-7\Pi_-)&=-8,
 \label{eq:V132-shape-L2-projector}\\
 \operatorname{tr}(2\Pi_++10\Pi_-)&=32,
 \label{eq:V132-shape-norm-projector}\\
 96L\operatorname{tr}S+192\operatorname{tr}S^2
 &=96L^2.
 \label{eq:V132-Robin-shape}
\end{align}
\end{lemma}

\begin{proof}
The projector ranks are one and three.  Hence the first two traces are
\(13-21=-8\) and \(2+30=32\).  Since \(S=-LP_N\),
\(\operatorname{tr}S=-L\) and
\(\operatorname{tr}S^2=L^2\), giving
\(-96L^2+192L^2=96L^2\).
\end{proof}

\begin{theorem}[One-form relative \(A_3\)]
\label{thm:V132-oneform-A3}
The physical relative one-form coefficient is
\begin{equation}
 \boxed{
 A_{3,1}^{\rm rel}
 =
 \frac{\sqrt\pi}{192}\int_\Sigma
 \left(
 64R-208R_{nn}+88L^2-64|L|^2
 \right)d\operatorname{vol}_\Sigma.}
 \label{eq:V132-oneform-A3}
\end{equation}
\end{theorem}

\begin{proof}
Insert Lemmas~\ref{lem:V132-projector-trace},
\ref{lem:V132-oneform-curvature}, and
\ref{lem:V132-oneform-shape} into
\eqref{eq:V132-universal-A3}.  The curvature rows are
\[
 96(R-2R_{nn})-32R-16R_{nn}
 =64R-208R_{nn}.
\]
The \(L^2\) rows are \(-8L^2+96L^2=88L^2\).
The \(|L|^2\) rows are
\[
 32|L|^2-12(8|L|^2)=-64|L|^2.
\]
This proves the formula.
\end{proof}

\subsection{Physical BRST combination}
\label{subsec:V132-physical-BRST}

\begin{theorem}[Physical gauge/ghost \(A_3\)]
\label{thm:V132-physical-BRST-A3}
After the scalar-minus-one-half-one-form combination and multiplication
by the twelve Standard Model gauge generators,
\begin{align}
 A_{3,\rm BRST}^{\rm phys}
 &=
 12\left(
 A_{3,0}^{\rm rel}-\frac12A_{3,1}^{\rm rel}
 \right)
 \notag\\
 &=
 \sqrt\pi\int_\Sigma
 \left(
 -3R+6R_{nn}
 -\frac{51}{16}L^2
 +\frac{21}{8}|L|^2
 \right)d\operatorname{vol}_\Sigma
 \label{eq:V132-physical-ambient}\\
 &=
 \boxed{
 \sqrt\pi\int_\Sigma
 \left(
 -3R_\Sigma
 -\frac3{16}H^2
 -\frac38|K|^2
 \right)d\operatorname{vol}_\Sigma.}
 \label{eq:V132-physical-intrinsic}
\end{align}
\end{theorem}

\begin{proof}
Subtract one half of \eqref{eq:V132-oneform-A3} from
\eqref{eq:V132-scalar-A3}.  Before multiplying by \(12\), the bracket is
\[
 -48R+96R_{nn}-51L^2+42|L|^2
\]
times \(\sqrt\pi/192\).  Multiplication by \(12\) gives
\eqref{eq:V132-physical-ambient}.  The Gauss identity in the present
convention is
\begin{equation}
 R_\Sigma=R-2R_{nn}+L^2-|L|^2.
 \label{eq:V132-Gauss}
\end{equation}
Thus
\(-3R+6R_{nn}=-3R_\Sigma+3L^2-3|L|^2\).
Combine coefficients and use \(L^2=H^2\),
\(|L|^2=|K|^2\) to obtain
\eqref{eq:V132-physical-intrinsic}.
\end{proof}

\begin{corollary}[Product and totally geodesic checks]
\label{cor:V132-checks}
If \(K=0\), then
\begin{equation}
 A_{3,\rm BRST}^{\rm phys}
 =-3\sqrt\pi\int_\Sigma R_\Sigma.
 \label{eq:V132-product-check}
\end{equation}
If the ambient metric is flat, then
\begin{equation}
 R_\Sigma=H^2-|K|^2
 \label{eq:V132-flat-Gauss}
\end{equation}
and the two forms
\eqref{eq:V132-physical-ambient} and
\eqref{eq:V132-physical-intrinsic} agree identically.
\end{corollary}

\begin{proof}
Set \(K=0\) in \eqref{eq:V132-physical-intrinsic}.  For the second claim,
set \(R=R_{nn}=0\) in \eqref{eq:V132-Gauss} and substitute.
\end{proof}

\subsection{Addition of the relative interface}
\label{subsec:V132-completed}

\begin{theorem}[Completed auxiliary-eliminated weight-three coefficient]
\label{thm:V132-completed-A3}
Add the physical coefficient only after the BRST trace is formed, as
required by V123.  The auxiliary-eliminated gauge/ghost coefficient is
\begin{align}
 A_{3,\rm BRST}^{\rm eff}
 &:=
 A_{3,\rm BRST}^{\rm phys}
 +A_{3,\rm BRST}^{\rm rel}
 \notag\\
 &=
 \boxed{
 \sqrt\pi\int_\Sigma
 \left(
 -3R_\Sigma
 -\frac{45}{16}H^2
 +\frac{51}{8}|K|^2
 \right)d\operatorname{vol}_\Sigma.}
 \label{eq:V132-completed-A3}
\end{align}
\end{theorem}

\begin{proof}
Use \eqref{eq:V132-physical-intrinsic} and the relative coefficient
\eqref{eq:V131-relative-A3}.  The quadratic coefficients add as
\[
 -\frac3{16}-\frac{21}{8}
 =-\frac{45}{16},
 \qquad
 -\frac38+\frac{27}{4}
 =\frac{51}{8}.
\]
The relative interface coefficient has no \(R_\Sigma\) term by the V128
product calibration.
\end{proof}

\begin{corollary}[Completed proper-time linear divergence]
\label{cor:V132-completed-linear}
With \(t\ge\Lambda^{-2}\),
\begin{equation}
 \boxed{
 \Gamma_{\rm BRST}^{\rm eff}(\Lambda)\big|_{\rm wt\,3}
 =
 \frac{\Lambda}{8\pi^{3/2}}
 \int_\Sigma
 \left(
 -3R_\Sigma
 -\frac{45}{16}H^2
 +\frac{51}{8}|K|^2
 \right)d\operatorname{vol}_\Sigma
 +O(\log\Lambda).}
 \label{eq:V132-completed-linear}
\end{equation}
\end{corollary}

\begin{proof}
In \eqref{eq:V132-A3-normalization}, the \(A_3\) heat term is
\((4\pi)^{-2}t^{-1/2}A_3\).  Division by \(t\) and integration from
\(\Lambda^{-2}\) gives \(2\Lambda(4\pi)^{-2}A_3\).
Insert \eqref{eq:V132-completed-A3}.
\end{proof}

\begin{firewall}[Independent weight-three boundary couplings]
\label{fire:V132-weight-three-scope}
The invariants \(R_\Sigma\), \(H^2\), and \(|K|^2\) are independent
weight-three boundary counterterms.  The V130 requirement fixing the
weight-two Einstein--GHY ratio does not determine their finite
renormalized coefficients.  Nor does
\eqref{eq:V132-completed-A3} prove that a microscopic renewal regulator
selects the proper-time subtraction.  V132 supplies the exact
gauge/ghost divergence which that regulator must match.
\end{firewall}

\subsection{V132 ledger and next acquisition}
\label{subsec:V132-ledger}

\paragraph{New exact conclusions.}
\begin{enumerate}[label=\textup{(\arabic*)},leftmargin=2.8em]
\item Theorem~\ref{thm:V132-universal-A3} fixes the mixed-boundary
      coefficient in the programme's heat normalization.
\item Lemma~\ref{lem:V132-projector-trace} proves the essential projector
      derivative identity
      \(\operatorname{tr}\chi_{:a}\chi_{:a}=8|L|^2\).
\item Proposition~\ref{prop:V132-scalar-A3} and
      Theorem~\ref{thm:V132-oneform-A3} compute the physical relative
      scalar and one-form rows.
\item Theorem~\ref{thm:V132-physical-BRST-A3} reduces their Standard
      Model BRST combination to
      \(-3R_\Sigma-\tfrac3{16}H^2-\tfrac38|K|^2\).
\item Theorem~\ref{thm:V132-completed-A3} adds the V131 relative
      interface and closes the auxiliary-eliminated gauge/ghost
      weight-three coefficient.
\end{enumerate}

\paragraph{Next forced step.}
The next bottleneck is no longer a missing gauge/ghost heat coefficient.
V133 should move upstream to the fermion domain.  It must select a local,
strongly elliptic boundary realization for the squared Standard Model
Dirac operator, prove compatibility with the gauge representation and
the doubled collar, and compute its \(A_2\) boundary coefficient before
any attempt to cancel or combine it with V130.  If no chirality-compatible
local projector exists for the required field content, that obstruction
must be stated before proceeding to fermionic determinants.

\begin{assessment}[V132 continuum status]
\label{ass:V132-continuum-status}
The full Standard--Model--Einstein continuum remains unproved.  V132
closes the physical-plus-interface gauge/ghost heat ledger through
weight three.  Fermion and Higgs boundary realizations, microscopic
renormalization, interacting nonlinear existence, Einstein stress
closure, and the cofinal continuum remain open.
\end{assessment}

\section*{V132 exact checks and append-only record}
\addcontentsline{toc}{section}{V132 exact checks and append-only record}

The mixed-boundary checksum is
\begin{equation}
 \boxed{
 A_{3,\rm BRST}^{\rm phys}
 =\sqrt\pi\int_\Sigma
 \left(-3R_\Sigma-\frac3{16}H^2-\frac38|K|^2\right),
 \quad
 A_{3,\rm BRST}^{\rm eff}
 =\sqrt\pi\int_\Sigma
 \left(-3R_\Sigma-\frac{45}{16}H^2+\frac{51}{8}|K|^2\right).}
 \label{eq:V132-final-check}
\end{equation}
The append operation preserves the complete V131 prefix byte for byte before
its former live terminator.  The assembled V132 file is checked for one live
final terminator, balanced environments and braces, duplicate labels, newly
introduced unresolved references, display math delimiters, and control
characters before the exact V131 predecessor is removed.

% =============================================================================
% END APPENDED VERSION V132 MATERIAL
% =============================================================================

% =============================================================================
% BEGIN APPENDED VERSION V133 MATERIAL
% =============================================================================

\section{V133: Euclidean fermion transversality and the obstruction to a
mixed squared domain}
\label{sec:V133}

\subsection{What is inherited and what is new}
\label{subsec:V133-purpose}

V31--V32 already solve the Lorentzian action/flux problem at the finite
trace level.  In particular:
\begin{enumerate}[label=\textup{(\roman*)},leftmargin=3.3em]
\item every conservative maximal zero-flux fermion trace is a unitary
      graph;
\item no constant unbroken-Standard-Model chiral bag intertwiner exists;
\item on the branch with a right-handed neutrino, nonzero Higgs boundary
      floor, and invertible Yukawa matrices, the polar map
      \(U_{\rm SM}(\Phi)\) gives a gauge-covariant maximal neutral graph;
      and
\item the graph changes rank or loses derivative control at a Higgs or
      Yukawa zero.
\end{enumerate}
V133 does not repeat those results.  It asks whether the selected trace
is an elliptic Euclidean boundary condition and whether its squared
Dirac domain is of the ordinary mixed type used in V129--V132.

\subsection{Frozen Euclidean Dirac half-space}
\label{subsec:V133-halfspace}

Let \(D_F\) be a Euclidean Dirac-type operator on a Hermitian Clifford
module \(\mathcal E\).  Use inward normal coordinate \(r\ge0\) and the
Clifford convention
\begin{equation}
 c(\zeta)^*=-c(\zeta),
 \qquad
 c(\zeta)^2=-|\zeta|^2I.
 \label{eq:V133-Clifford}
\end{equation}
At a fixed boundary point and nonzero tangential covector \(\xi\), the
frozen principal equation is
\begin{equation}
 \left[c(n)\partial_r+i\,c(\xi)\right]\psi=0.
 \label{eq:V133-frozen-Dirac}
\end{equation}
Define
\begin{equation}
 J:=i\,c(n),
 \qquad
 A(\xi):=i\,c(n)c(\xi).
 \label{eq:V133-JA}
\end{equation}

\begin{lemma}[Flux involution and normal evolution]
\label{lem:V133-JA}
The operators in \eqref{eq:V133-JA} satisfy
\begin{equation}
 J^*=J,\quad J^2=I,\qquad
 A(\xi)^*=A(\xi),\quad A(\xi)^2=|\xi|^2I,\qquad
 JA(\xi)+A(\xi)J=0.
 \label{eq:V133-JA-identities}
\end{equation}
Moreover, \eqref{eq:V133-frozen-Dirac} is equivalent to
\begin{equation}
 \partial_r\psi=A(\xi)\psi.
 \label{eq:V133-normal-evolution}
\end{equation}
\end{lemma}

\begin{proof}
The adjoint and square identities follow directly from
\eqref{eq:V133-Clifford} and
\(c(n)c(\xi)+c(\xi)c(n)=0\).  Multiplying
\eqref{eq:V133-frozen-Dirac} on the left by \(-c(n)\) gives
\[
 \partial_r\psi-i\,c(n)c(\xi)\psi=0,
\]
which is \eqref{eq:V133-normal-evolution}.
\end{proof}

Let \(\mathcal E_\pm=\ker(J\mp I)\).  Anticommutation in
\eqref{eq:V133-JA-identities} makes \(A(\xi)\) off diagonal relative to
\(\mathcal E_+\oplus\mathcal E_-\).  Write
\begin{equation}
 A(\xi)
 =
 |\xi|
 \begin{pmatrix}
 0&C(\widehat\xi)^*\\
 C(\widehat\xi)&0
 \end{pmatrix},
 \qquad
 \widehat\xi=\xi/|\xi|.
 \label{eq:V133-Calderon-block}
\end{equation}

\begin{lemma}[Calderon graph]
\label{lem:V133-Calderon-graph}
For every \(\widehat\xi\), the map
\(C(\widehat\xi):\mathcal E_+\to\mathcal E_-\) is unitary.  The boundary
values of solutions of \eqref{eq:V133-frozen-Dirac} which decay as
\(r\to+\infty\) form the graph
\begin{equation}
 \mathcal C_-(\widehat\xi)
 =
 \left\{
 u\oplus[-C(\widehat\xi)u]:
 u\in\mathcal E_+
 \right\}.
 \label{eq:V133-decaying-graph}
\end{equation}
\end{lemma}

\begin{proof}
The square identity \(A(\xi)^2=|\xi|^2I\) applied to
\eqref{eq:V133-Calderon-block} gives
\(C^*C=I=CC^*\).  A decaying solution of
\eqref{eq:V133-normal-evolution} has initial value in the
\(-|\xi|\)-eigenspace of \(A(\xi)\).  The eigenvalue equation for
\(u\oplus v\) is
\[
 C^*v=-u,\qquad Cu=-v,
\]
so \(v=-Cu\), proving \eqref{eq:V133-decaying-graph}.
\end{proof}

\subsection{Exact Lopatinski criterion for a neutral graph}
\label{subsec:V133-Lopatinski}

Let \(U:\mathcal E_+\to\mathcal E_-\) be unitary and impose the local
boundary condition
\begin{equation}
 \psi|_\Sigma\in\mathcal M_U,
 \qquad
 \mathcal M_U:=\{u\oplus Uu:u\in\mathcal E_+\}.
 \label{eq:V133-U-domain}
\end{equation}
This is the Euclidean counterpart of the normalized neutral graph of
V32.

\begin{theorem}[Graph Lopatinski criterion]
\label{thm:V133-Lopatinski}
The frozen first-order boundary problem
\eqref{eq:V133-frozen-Dirac}--\eqref{eq:V133-U-domain} satisfies the
Lopatinski condition at \(\widehat\xi\) if and only if
\begin{equation}
 \boxed{
 U+C(\widehat\xi):
 \mathcal E_+\longrightarrow\mathcal E_-
 \quad\hbox{is invertible}.}
 \label{eq:V133-Lopatinski}
\end{equation}
On a compact boundary, it is uniformly strongly elliptic at principal
level if and only if
\begin{equation}
 \boxed{
 \delta_{\rm Lop}(U)
 :=
 \inf_{\substack{y\in\Sigma\\|\xi|=1}}
 s_{\min}\bigl(U(y)+C_y(\widehat\xi)\bigr)>0.}
 \label{eq:V133-Lopatinski-gap}
\end{equation}
\end{theorem}

\begin{proof}
The Lopatinski condition says that the prescribed boundary subspace has
zero intersection with the decaying Cauchy-data subspace and that their
dimensions sum to the full trace dimension.  Both graphs have dimension
\(\dim\mathcal E_+\).  A vector belongs to their intersection precisely
when
\[
 u\oplus Uu
 =
 u\oplus[-C(\widehat\xi)u],
\]
or \((U+C(\widehat\xi))u=0\).  Thus zero intersection is equivalent to
invertibility, proving \eqref{eq:V133-Lopatinski}.  The sphere bundle
\(\{(y,\xi):|\xi|=1\}\) is compact and the least singular value is
continuous.  Pointwise invertibility is uniform exactly when its infimum
is positive.
\end{proof}

\begin{corollary}[Neutrality does not imply ellipticity]
\label{cor:V133-neutral-not-elliptic}
There are maximal zero-flux unitary graphs which fail the Lopatinski
condition.
\end{corollary}

\begin{proof}
Fix one direction \(\widehat\xi_0\) and take
\(U=-C(\widehat\xi_0)\).  It is unitary, hence its graph is maximal
neutral by V32, but
\(U+C(\widehat\xi_0)=0\).
\end{proof}

\begin{firewall}[The missing Euclidean hypothesis in V32]
\label{fire:V133-V32-gap}
Theorem~\ref{thm:V32-Higgs-trace} proves unitarity, gauge covariance, and
zero flux of \(U_{\rm SM}(\Phi)\) on its Higgs-floor branch.  None of
those properties implies
\eqref{eq:V133-Lopatinski-gap}.  A fermion heat determinant may therefore
be formed from that graph only after the additional certificate
\[
 \delta_{\rm Lop}(U_{\rm SM}(\Phi))>0
\]
has been proved uniformly on the boundary and along the cofinal family.
\end{firewall}

\subsection{Gauge covariance of the ellipticity gate}
\label{subsec:V133-gauge}

\begin{proposition}[Gauge invariance of the Lopatinski singular values]
\label{prop:V133-gauge-invariance}
Suppose the gauge action is unitary, commutes with Clifford
multiplication, and acts by \(\rho_\pm(g)\) on
\(\mathcal E_\pm\).  If
\begin{align}
 U_{g\Phi}
 &=\rho_-(g)U_\Phi\rho_+(g)^{-1},
 \label{eq:V133-U-gauge}\\
 C_{g\cdot y}(\widehat\xi)
 &=\rho_-(g)C_y(\widehat\xi)\rho_+(g)^{-1},
 \label{eq:V133-C-gauge}
\end{align}
then every singular value of \(U_\Phi+C_y(\widehat\xi)\), and hence
\(\delta_{\rm Lop}\), is gauge invariant.
\end{proposition}

\begin{proof}
Equations \eqref{eq:V133-U-gauge}--\eqref{eq:V133-C-gauge} give
\[
 U_{g\Phi}+C_{g\cdot y}
 =
 \rho_-(g)(U_\Phi+C_y)\rho_+(g)^{-1}.
\]
Left and right multiplication by unitary matrices preserves singular
values.
\end{proof}

\begin{corollary}[Exact Standard Model branch structure]
\label{cor:V133-SM-branches}
The fermion boundary programme has the following rigorous branches.
\begin{enumerate}[label=\textup{(\alph*)},leftmargin=3em]
\item In the minimal model without a right-handed neutrino, the V32 polar
      graph is rank deficient before ellipticity is tested.
\item If \(\Phi=0\) or a required Yukawa singular value vanishes, the
      V32 graph loses its fixed-rank chart before ellipticity is tested.
\item On the right-handed-neutrino Higgs-floor branch, the remaining
      first-order Euclidean condition is exactly
      \eqref{eq:V133-Lopatinski-gap}; it is gauge invariant by
      Proposition~\ref{prop:V133-gauge-invariance}.
\end{enumerate}
\end{corollary}

\begin{proof}
The first two statements are
Proposition~\ref{prop:V32-SM-no-constant} together with the rank and floor
conclusions of V32.  The third is
Theorem~\ref{thm:V133-Lopatinski} applied to the graph already constructed
there, followed by Proposition~\ref{prop:V133-gauge-invariance}.
\end{proof}

\subsection{What squaring the graph domain actually produces}
\label{subsec:V133-square}

Let \(P_U\) be the orthogonal projector onto \(\mathcal M_U\) and put
\(Q_U=I-P_U\).  The first-order domain is
\begin{equation}
 Q_U\psi|_\Sigma=0.
 \label{eq:V133-first-domain}
\end{equation}
The natural squared domain is
\begin{equation}
 Q_U\psi|_\Sigma=0,
 \qquad
 Q_UD_F\psi|_\Sigma=0.
 \label{eq:V133-squared-domain}
\end{equation}

\begin{lemma}[Clifford exchange for a neutral graph]
\label{lem:V133-Clifford-exchange}
For a maximal neutral graph for the boundary Green form,
\begin{equation}
 Q_Uc(n)=c(n)P_U,
 \qquad
 P_Uc(n)=c(n)Q_U.
 \label{eq:V133-Clifford-exchange}
\end{equation}
\end{lemma}

\begin{proof}
Maximal neutrality says \(c(n)\mathcal M_U\) is orthogonal to
\(\mathcal M_U\).  Both spaces have half the trace dimension, so
\(c(n)\mathcal M_U=\mathcal M_U^\perp\).  Thus \(c(n)\) exchanges the
ranges of \(P_U\) and \(Q_U\), which is exactly
\eqref{eq:V133-Clifford-exchange}.
\end{proof}

Write the Dirac operator near the boundary as
\begin{equation}
 D_F=c(n)\bigl(\nabla_n+B_0+B_1\bigr),
 \label{eq:V133-product-Dirac}
\end{equation}
where \(B_1\) is the tangential first-order part and \(B_0\) is
zeroth order.

\begin{theorem}[Squared graph boundary operator]
\label{thm:V133-squared-graph}
On fields satisfying \(Q_U\psi=0\), the second row of
\eqref{eq:V133-squared-domain} is equivalent to
\begin{equation}
 P_U\nabla_n(P_U\psi)
 +P_UB_1P_U\psi
 +S_U P_U\psi=0,
 \label{eq:V133-oblique-square}
\end{equation}
where
\begin{equation}
 S_U
 :=
 P_UB_0P_U+P_U(\nabla_nP_U)P_U
 \label{eq:V133-SU}
\end{equation}
up to the fixed sign convention used to place \(c(n)\) on the left.
Thus the squared problem is an ordinary mixed
Dirichlet--Robin problem only if
\begin{equation}
 \boxed{
 P_UB_1P_U=0.}
 \label{eq:V133-no-oblique-condition}
\end{equation}
\end{theorem}

\begin{proof}
Use \(Q_U\psi=0\) to write \(\psi=P_U\psi\).  Then
\[
 Q_UD_F\psi
 =
 Q_Uc(n)(\nabla_n+B_0+B_1)P_U\psi.
\]
Apply \eqref{eq:V133-Clifford-exchange} and the invertibility of
\(c(n)\) to replace the leading factor by \(P_U\).  Expanding
\(\nabla_n(P_U\psi)\) gives the normal derivative and
\(P_U(\nabla_nP_U)P_U\).  The tangential principal part is exactly
\(P_UB_1P_U\).  It vanishes precisely under
\eqref{eq:V133-no-oblique-condition}; otherwise
\eqref{eq:V133-oblique-square} is an oblique boundary operator.
\end{proof}

\begin{corollary}[Gate for using the V132 mixed formula]
\label{cor:V133-mixed-formula-gate}
The universal mixed heat formula
\eqref{eq:V132-universal-A3} applies to the squared fermion graph only
after both
\begin{equation}
 \delta_{\rm Lop}(U)>0,
 \qquad
 P_UB_1P_U=0
 \label{eq:V133-two-fermion-gates}
\end{equation}
have been established.  If the second condition fails, a strongly
elliptic oblique-boundary heat calculus is required and its tangential
symbol contributes new local invariants.
\end{corollary}

\begin{proof}
The first condition supplies principal ellipticity.  The second removes
the tangential derivative from \eqref{eq:V133-oblique-square}, leaving a
mixed Dirichlet condition on \(Q_U\) and Robin condition on \(P_U\).
These are precisely the hypotheses of
Theorem~\ref{thm:V132-universal-A3}.  Without the second condition, its
boundary operator class is different.
\end{proof}

\begin{firewall}[Squaring does not select the fermion domain]
\label{fire:V133-square-not-selection}
The rule \(Q_U\psi=Q_UD_F\psi=0\) propagates a selected first-order
domain to \(D_F^2\); it does not choose \(U\).  It also does not repair a
failure of \(\delta_{\rm Lop}\), a Higgs zero, a missing right-handed
neutrino, or an unbroken-gauge representation mismatch.  Using a standard
bag \(A_2\) coefficient before these graph gates are proved would replace
the Standard Model boundary problem by a different theory.
\end{firewall}

\subsection{Finite and cofinal certificates}
\label{subsec:V133-cofinal}

\begin{proposition}[Stable finite-level ellipticity handoff]
\label{prop:V133-cofinal-handoff}
Let \(U_q\) and \(C_q\) be finite-level graph and Calderon maps, transported
to common boundary frames.  Suppose
\begin{equation}
 \inf_q\delta_{\rm Lop}(U_q)\ge\delta_*>0,
 \qquad
 \|U_q-U\|+\|C_q-C\|\longrightarrow0
 \label{eq:V133-cofinal-gap}
\end{equation}
uniformly on the unit boundary cotangent bundle.  Then
\begin{equation}
 \delta_{\rm Lop}(U)\ge\delta_*>0.
 \label{eq:V133-limit-gap}
\end{equation}
If in addition
\begin{equation}
 \|P_{U_q}B_{1,q}P_{U_q}\|\longrightarrow0,
 \label{eq:V133-cofinal-oblique}
\end{equation}
then the limiting squared domain satisfies
\(P_UB_1P_U=0\).
\end{proposition}

\begin{proof}
The least singular value is one-Lipschitz in operator norm:
\[
 |s_{\min}(A)-s_{\min}(B)|\le\|A-B\|.
\]
Apply this uniformly to
\(A=U_q+C_q\) and \(B=U+C\), then take the infimum and the limit.
For the second assertion, norm convergence of the projectors follows
from the explicit graph projector
\eqref{eq:V32-graph-projector}.  Pass to the limit in
\eqref{eq:V133-cofinal-oblique}.
\end{proof}

\begin{assessment}[Actual fermion bottleneck]
\label{ass:V133-fermion-bottleneck}
The fermion obstruction is now a pair of explicit, same-frame matrix
tests rather than an unspecified boundary choice:
\[
 \delta_{\rm Lop}(U_{\rm SM}(\Phi))>0,
 \qquad
 P_{U_{\rm SM}}B_1P_{U_{\rm SM}}=0.
\]
The first is required for a local Euclidean determinant.  The second
decides whether the square belongs to the ordinary mixed heat class or
to the oblique class.  Neither test is contained in the Lorentzian
zero-flux proof of V32.
\end{assessment}

\subsection{V133 ledger and next acquisition}
\label{subsec:V133-ledger}

\paragraph{New exact conclusions.}
\begin{enumerate}[label=\textup{(\arabic*)},leftmargin=2.8em]
\item Lemma~\ref{lem:V133-Calderon-graph} computes the decaying Euclidean
      Cauchy-data graph.
\item Theorem~\ref{thm:V133-Lopatinski} proves the exact uniform
      transversality criterion for a selected neutral graph.
\item Corollary~\ref{cor:V133-neutral-not-elliptic} proves that the V32
      zero-flux property alone is insufficient.
\item Proposition~\ref{prop:V133-gauge-invariance} proves gauge invariance
      of the new singular-value gate.
\item Theorem~\ref{thm:V133-squared-graph} derives the squared boundary
      operator and isolates the oblique term \(P_UB_1P_U\).
\item Proposition~\ref{prop:V133-cofinal-handoff} gives the uniform
      finite-to-continuum ellipticity handoff.
\end{enumerate}

\paragraph{Next forced step.}
V134 should evaluate the two V133 matrices for the actual Higgs-polar
factorization \(U_{\rm SM}=C_{\rm sp}\otimes U_\Phi\).  It must first
write the Euclidean spin Calderon map in the same chiral and flux frames
as \(C_{\rm sp}\).  If the spectral gap is positive and
\eqref{eq:V133-no-oblique-condition} holds, V134 may compute the fermion
\(A_2\) coefficient.  If the oblique block survives, the programme must
switch to the oblique heat formula instead of forcing the V132 mixed
coefficient.

\begin{assessment}[V133 continuum status]
\label{ass:V133-continuum-status}
The full Standard--Model--Einstein continuum remains unproved.  V133
connects the existing Lorentzian Higgs-polar trace to the precise
Euclidean ellipticity and squared-domain gates.  Those gates, the
fermion heat coefficient, Higgs sector, microscopic renormalization,
interacting nonlinear existence, Einstein stress closure, and the
cofinal continuum remain open.
\end{assessment}

\section*{V133 exact checks and append-only record}
\addcontentsline{toc}{section}{V133 exact checks and append-only record}

The fermion-domain checksum is
\begin{equation}
 \boxed{
 \delta_{\rm Lop}(U)
 =
 \inf_{y,|\xi|=1}s_{\min}(U+C_y(\widehat\xi))>0,
 \qquad
 P_UB_1P_U=0
 \ \Longleftrightarrow\
 \hbox{ordinary mixed square}.}
 \label{eq:V133-final-check}
\end{equation}
The append operation preserves the complete V132 prefix byte for byte before
its former live terminator.  The assembled V133 file is checked for one live
final terminator, balanced environments and braces, duplicate labels, newly
introduced unresolved references, display math delimiters, and control
characters before the exact V132 predecessor is removed.

% =============================================================================
% END APPENDED VERSION V133 MATERIAL
% =============================================================================

% =============================================================================
% BEGIN APPENDED VERSION V134 MATERIAL
% =============================================================================

\section{V134: the chiral Clifford-module gate and the determinant-line
fork}
\label{sec:V134}

\subsection{Why V133 must be applied one level upstream}
\label{subsec:V134-purpose}

V133 derived the correct ellipticity test for a local unitary graph on a
\emph{Hermitian Euclidean Clifford module}.  The V32 Higgs-polar graph
was constructed instead from the Lorentzian Standard Model action and
was conditionally factorized as
\[
 E_+^{\rm sp}\otimes R_R^\nu
 \longrightarrow
 E_-^{\rm sp}\otimes R_L.
\]
Before its Lopatinski matrix can be evaluated, one must prove that these
are the positive and negative flux spaces of one gauge-covariant
Euclidean Dirac-type principal symbol.  V134 proves that they cannot be
so identified in the unbroken chiral Standard Model: the Calderon map
itself would already be a unitary gauge intertwiner between \(R_R^\nu\)
and \(R_L\).

This does not invalidate the Lorentzian neutral-graph theorem.  It fixes
the Euclidean continuation order.  One must first choose either a
vectorlike Hermitian doubling, keeping track of its determinant
multiplicity and phase, or a genuinely chiral determinant-line boundary
problem.  Only then is there a defined Euclidean fermion operator to
which V133 and the heat calculus can be applied.

\subsection{Gauge covariance of the principal Calderon map}
\label{subsec:V134-Calderon-gauge}

Let \((\mathcal E,c,\rho)\) be a Hermitian Clifford module for a compact
gauge group \(G\), so that
\begin{equation}
 \rho(g)c(\zeta)=c(\zeta)\rho(g)
 \qquad(g\in G).
 \label{eq:V134-Clifford-gauge}
\end{equation}
Let \(J=i\,c(n)\), let
\(\mathcal E_\pm=\ker(J\mp I)\), and let
\(C(\widehat\xi):\mathcal E_+\to\mathcal E_-\) be the unitary Calderon
block from \eqref{eq:V133-Calderon-block}.

\begin{lemma}[Gauge invariance of the flux splitting]
\label{lem:V134-flux-gauge}
The spaces \(\mathcal E_\pm\) are \(G\)-invariant.
\end{lemma}

\begin{proof}
Equation~\eqref{eq:V134-Clifford-gauge} with \(\zeta=n\) gives
\(\rho(g)J=J\rho(g)\).  Therefore every spectral subspace of \(J\) is
invariant.
\end{proof}

\begin{theorem}[Principal Calderon intertwiner]
\label{thm:V134-Calderon-intertwiner}
For every nonzero tangential direction,
\begin{equation}
 \boxed{
 C(\widehat\xi)\rho_+(g)
 =
 \rho_-(g)C(\widehat\xi)
 \qquad(g\in G),}
 \label{eq:V134-Calderon-intertwiner}
\end{equation}
where \(\rho_\pm\) are the restrictions to \(\mathcal E_\pm\).
Consequently \((\mathcal E_+,\rho_+)\) and
\((\mathcal E_-,\rho_-)\) are unitarily isomorphic gauge
representations before any mass, Higgs, or Yukawa term is introduced.
\end{theorem}

\begin{proof}
The frozen normal evolution operator
\[
 A(\xi)=i\,c(n)c(\xi)
\]
commutes with \(\rho(g)\) by \eqref{eq:V134-Clifford-gauge}.  Relative to
the invariant splitting of Lemma~\ref{lem:V134-flux-gauge}, compare the
lower-left blocks in
\[
 A(\xi)\rho(g)=\rho(g)A(\xi).
\]
Using \eqref{eq:V133-Calderon-block} gives
\eqref{eq:V134-Calderon-intertwiner}.  The map \(C(\widehat\xi)\) is
unitary by Lemma~\ref{lem:V133-Calderon-graph}, hence is a unitary
representation isomorphism.
\end{proof}

\begin{corollary}[No \(R_R^\nu\)-to-\(R_L\) principal flux
factorization]
\label{cor:V134-no-chiral-flux-factor}
Suppose the spin factors in the V32 ansatz are fixed isomorphic boundary
spin modules.  If
\begin{equation}
 \mathcal E_+\cong E_+^{\rm sp}\otimes R_R^\nu,
 \qquad
 \mathcal E_-\cong E_-^{\rm sp}\otimes R_L
 \label{eq:V134-V32-factorization}
\end{equation}
as representations of the unbroken Standard Model gauge group, then
these spaces cannot be the flux splitting of a gauge-covariant Hermitian
Euclidean Clifford module.
\end{corollary}

\begin{proof}
The spin factors carry no Standard Model internal representation, so
Theorem~\ref{thm:V134-Calderon-intertwiner} would imply
\(R_R^\nu\cong R_L\).  Proposition~\ref{prop:V32-SM-no-constant} proves
that these representations are not isomorphic.
\end{proof}

\begin{firewall}[Family covariance does not alter the principal symbol]
\label{fire:V134-Higgs-lower-order}
The V32 identity
\[
 U_{g\Phi}=\rho_L(g)U_\Phi\rho_R(g)^{-1}
\]
is a valid covariance law for a Higgs-dependent family of boundary
graphs.  It does not make the Higgs-independent principal Clifford block
\(C(\widehat\xi)\) into a map from \(R_R^\nu\) to \(R_L\).
Yukawa and Higgs terms are zeroth order in the Dirac operator; they cannot
change the representation content of the principal flux spaces or the
Lopatinski subspace.
\end{firewall}

\subsection{The actual chiral kinetic bundles}
\label{subsec:V134-chiral-bundles}

For a unitary gauge representation \(V\), the Euclidean chiral Dirac
operator has the form
\begin{equation}
 D_V^+:
 \Gamma(S^+\otimes V)
 \longrightarrow
 \Gamma(S^-\otimes V),
 \qquad
 D_V^-=(D_V^+)^*.
 \label{eq:V134-chiral-Dirac}
\end{equation}
The internal representation is the \emph{same} on the two spin
chiralities in \eqref{eq:V134-chiral-Dirac}.  The physical Standard Model
is a direct sum of several such chiral arrows with different \(V\)'s and
orientations; it is not one self-adjoint Dirac operator which uses
\(R_L\) as one Clifford half and \(R_R\) as the other.

\begin{proposition}[Hermitian double of one chiral arrow]
\label{prop:V134-Hermitian-double}
For every chiral arrow \eqref{eq:V134-chiral-Dirac}, the operator
\begin{equation}
 \mathbb D_V
 :=
 \begin{pmatrix}
 0&D_V^-\\
 D_V^+&0
 \end{pmatrix}
 \quad\hbox{on}\quad
 (S^+\otimes V)\oplus(S^-\otimes V)
 \label{eq:V134-doubled-Dirac}
\end{equation}
is formally self-adjoint and has gauge-covariant Clifford principal
symbol.  Its square is
\begin{equation}
 \mathbb D_V^2
 =
 \begin{pmatrix}
 D_V^-D_V^+&0\\
 0&D_V^+D_V^-
 \end{pmatrix}.
 \label{eq:V134-doubled-square}
\end{equation}
\end{proposition}

\begin{proof}
Formal adjunction exchanges the two off-diagonal entries, proving
self-adjointness.  Clifford multiplication acts on the spin factors and
commutes with the common internal representation \(V\), giving
\eqref{eq:V134-Clifford-gauge}.  Direct block multiplication gives
\eqref{eq:V134-doubled-square}.
\end{proof}

\begin{lemma}[Doubling changes the determinant multiplicity]
\label{lem:V134-determinant-multiplicity}
At a finite regulator, assume \(D_V^+\) is an invertible square matrix and
\(D_V^-=(D_V^+)^*\).  Then
\begin{align}
 \det(D_V^-D_V^+)&=|\det D_V^+|^2,
 \label{eq:V134-one-chiral-modulus}\\
 \det\mathbb D_V^2&=|\det D_V^+|^4,
 \label{eq:V134-double-determinant}\\
 \frac12\log\det\mathbb D_V^2
 &=2\log|\det D_V^+|.
 \label{eq:V134-double-log}
\end{align}
\end{lemma}

\begin{proof}
The first identity is multiplicativity of the determinant and
\(\det(D_V^+)^*=\overline{\det D_V^+}\).
The two diagonal blocks in \eqref{eq:V134-doubled-square} have the same
nonzero determinant, proving the second and third identities.
\end{proof}

\begin{firewall}[The Hermitian double is not one physical Weyl
determinant]
\label{fire:V134-double-not-Weyl}
The self-adjoint double gives two copies of the logarithmic modulus in
\eqref{eq:V134-double-log} and contains no phase of
\(\det D_V^+\).  A single chiral fermion requires a determinant or
Pfaffian line, its orientation, and its gauge/anomaly transport.
Dividing the doubled heat coefficient by two recovers a modulus only
after the boundary domains of \(D_V^-D_V^+\) and \(D_V^+D_V^-\) have
been paired.  It does not reconstruct the missing phase.
\end{firewall}

\subsection{Two admissible Euclidean branches}
\label{subsec:V134-fork}

\begin{theorem}[Euclidean fermion fork]
\label{thm:V134-Euclidean-fork}
A Euclidean continuation of the chiral Standard Model fermion sector
compatible with the conclusions above must declare one of the following
two structures.
\begin{enumerate}[label=\textup{(\mathrm D)},leftmargin=4.0em]
\item A Hermitian doubled Clifford module, assembled representation by
      representation as in
      Proposition~\ref{prop:V134-Hermitian-double}, together with local
      graph domains satisfying the V133 Lopatinski test.  Its heat trace
      computes a modulus with the multiplicity
      \eqref{eq:V134-double-log}; the chiral phase is separate.
\end{enumerate}
\begin{enumerate}[label=\textup{(\mathrm C)},leftmargin=4.0em]
\item A chiral Fredholm boundary problem for the arrows \(D_V^+\),
      together with a determinant/Pfaffian line, a boundary condition
      such as a spectral projector when required, and a specified phase
      transport.  Its heat coefficients cannot be replaced automatically
      by one half of an arbitrary doubled bag problem.
\end{enumerate}
The V32 Higgs-polar graph alone specifies neither branch.
\end{theorem}

\begin{proof}
Corollary~\ref{cor:V134-no-chiral-flux-factor} excludes the proposed
identification of the unbroken left and right Standard Model
representations as the two flux halves of one Hermitian Clifford module.
Proposition~\ref{prop:V134-Hermitian-double} constructs a Hermitian
module for each chiral arrow, giving branch \textup{(D)} and its
determinant limitation.  Retaining the undoubled arrow
\eqref{eq:V134-chiral-Dirac} gives a Fredholm map between distinct
bundles rather than a self-adjoint operator, so its determinant is a
line-valued object with phase data; this is branch \textup{(C)}.
These exhaust the two ways used here to pass from a chiral arrow to a
Euclidean determinant object without changing its principal symbol.
\end{proof}

\subsection{Corrected use of the Higgs-polar graph}
\label{subsec:V134-Higgs-disposition}

\begin{proposition}[Where \(U_\Phi\) can enter after doubling]
\label{prop:V134-Higgs-after-double}
On branch \textup{(D)}, the principal flux spaces for every doubled
representation \(V\) are already isomorphic through the spin Calderon
map.  A Higgs--Yukawa field may enter:
\begin{enumerate}[label=\textup{(\alph*)},leftmargin=3em]
\item as the zeroth-order endomorphism in the doubled Dirac operator; or
\item as a background-dependent lower-order boundary endomorphism on a
      fixed-rank elliptic graph.
\end{enumerate}
It may not be used to create the missing principal representation
isomorphism.
\end{proposition}

\begin{proof}
Theorem~\ref{thm:V134-Calderon-intertwiner} supplies the principal
isomorphism independently of lower-order fields.  The Yukawa operator is
zeroth order, so adding it leaves \(A(\xi)\), \(C(\widehat\xi)\), and the
principal Lopatinski space unchanged.  A background-dependent projector
may alter lower-order boundary data provided its rank and principal
transversality remain fixed.
\end{proof}

\begin{corollary}[Revised V133 test object]
\label{cor:V134-revised-test}
The matrix in \eqref{eq:V133-Lopatinski-gap} must be formed from the graph
projector selected on the declared doubled representation-by-
representation Clifford module, or from the principal boundary symbol of
the declared chiral Fredholm problem.  It cannot be evaluated by
substituting the abstract factor
\(C_{\rm sp}\otimes U_\Phi:R_R^\nu\to R_L\) into a Calderon map whose
domain and target have not been identified with that same module.
\end{corollary}

\begin{proof}
The Lopatinski sum \(U+C(\widehat\xi)\) is defined only when both maps have
the same domain and codomain.  Corollary~\ref{cor:V134-no-chiral-flux-factor}
shows that the proposed unbroken chiral factorization does not provide
such a common principal bundle.  Either branch in
Theorem~\ref{thm:V134-Euclidean-fork} does.
\end{proof}

\begin{assessment}[Upstream obstruction selected]
\label{ass:V134-upstream-obstruction}
The missing datum is not a numerical bag phase.  It is the Euclidean
fermion determinant object itself.  The programme must choose between a
doubled local modulus calculation and a chiral determinant-line
calculation, state how the physical chiral phase is retained, and only
then compute a fermion boundary heat coefficient.  This is upstream of
the V133 transversality and oblique-symbol tests.
\end{assessment}

\subsection{Finite/cofinal data card}
\label{subsec:V134-card}

\begin{definition}[Euclidean fermion determinant card]
\label{def:V134-card}
At each finite level \(q\), record
\begin{equation}
 \mathfrak D_{F,q}
 =
 \left(
 \begin{array}{c}
 \text{chiral arrows and internal representations};\\
 \text{branch \(\mathrm D\) or \(\mathrm C\)};\\
 \text{boundary projector and its principal symbol};\\
 \delta_{{\rm Lop},q};\
 \|P_qB_{1,q}P_q\|;\\
 \text{determinant/Pfaffian multiplicity, phase, and zero-mode quotient}
 \end{array}
 \right).
 \label{eq:V134-card}
\end{equation}
\end{definition}

\begin{theorem}[No fermion heat handoff without the card]
\label{thm:V134-no-heat-without-card}
A cofinal fermion heat coefficient can be combined with V130--V132 only
if the cards \(\mathfrak D_{F,q}\) have a fixed branch, stable
representation multiplicities, a uniform positive Lopatinski floor, a
declared mixed or oblique squared-domain class, and compatible
determinant-line transport.  These conditions are necessary even if all
bulk Dirac symbols and Lorentzian neutral-graph residuals converge.
\end{theorem}

\begin{proof}
Changing the branch changes the determinant multiplicity by
Lemma~\ref{lem:V134-determinant-multiplicity} and may add phase data.
Changing representation multiplicity changes the bundle rank and all
heat coefficients.  Loss of the Lopatinski floor destroys uniform
ellipticity by Theorem~\ref{thm:V133-Lopatinski}.  A nonzero oblique block
changes the applicable heat formula by
Corollary~\ref{cor:V133-mixed-formula-gate}.  Finally, local heat
convergence controls the modulus but cannot identify a determinant-line
phase.  Each listed entry is therefore necessary.
\end{proof}

\subsection{V134 ledger and next acquisition}
\label{subsec:V134-ledger}

\paragraph{New exact conclusions.}
\begin{enumerate}[label=\textup{(\arabic*)},leftmargin=2.8em]
\item Theorem~\ref{thm:V134-Calderon-intertwiner} proves that the
      principal Calderon map is already a unitary gauge intertwiner.
\item Corollary~\ref{cor:V134-no-chiral-flux-factor} rules out the
      \(R_R^\nu\)-to-\(R_L\) factorization as the flux splitting of an
      unbroken gauge-covariant Hermitian Clifford module.
\item Proposition~\ref{prop:V134-Hermitian-double} constructs the correct
      representation-by-representation Hermitian double.
\item Lemma~\ref{lem:V134-determinant-multiplicity} proves its exact
      modulus multiplicity and exposes the missing chiral phase.
\item Theorem~\ref{thm:V134-Euclidean-fork} states the doubled-local and
      chiral-determinant-line branches without identifying them.
\item Theorem~\ref{thm:V134-no-heat-without-card} lists the necessary
      finite/cofinal handoff data.
\end{enumerate}

\paragraph{Next forced step.}
V135 is the next compulsory companion audit.  After the audit it should
choose one Euclidean fermion branch for calculation.  The conservative
choice is branch \textup{(D)} as a modulus computation: construct a
standard local bag projector separately on every doubled internal
representation, prove its V133 gap and mixed-square condition, and
divide the doubled modulus coefficient by the exact multiplicity in
\eqref{eq:V134-double-log}.  The determinant phase and anomaly transport
must remain an explicit unsolved card rather than being set to one.

\begin{assessment}[V134 continuum status]
\label{ass:V134-continuum-status}
The full Standard--Model--Einstein continuum remains unproved.  V134
identifies the Euclidean chiral-bundle obstruction that precedes a
fermion heat determinant.  Branch selection, phase/anomaly transport,
the fermion and Higgs coefficients, microscopic renormalization,
interacting nonlinear existence, Einstein stress closure, and the
cofinal continuum remain open.
\end{assessment}

\section*{V134 exact checks and append-only record}
\addcontentsline{toc}{section}{V134 exact checks and append-only record}

The chiral-bundle checksum is
\begin{equation}
 \boxed{
 C(\widehat\xi)\rho_+(g)=\rho_-(g)C(\widehat\xi),
 \qquad
 \frac12\log\det\mathbb D_V^2
 =2\log|\det D_V^+|.}
 \label{eq:V134-final-check}
\end{equation}
The append operation preserves the complete V133 prefix byte for byte before
its former live terminator.  The assembled V134 file is checked for one live
final terminator, balanced environments and braces, duplicate labels, newly
introduced unresolved references, display math delimiters, and control
characters before the exact V133 predecessor is removed.

% =============================================================================
% END APPENDED VERSION V134 MATERIAL
% =============================================================================

% =============================================================================
% BEGIN APPENDED VERSION V135 MATERIAL
% =============================================================================

\section{V135: companion audit, the doubled bag modulus, and the
three-generation weight-two cancellation}
\label{sec:V135}

\subsection{Purpose}
\label{subsec:V135-purpose}

V134 exposed the Euclidean fermion fork.  V135 performs the compulsory
companion audit and then selects branch \textup{(D)} for a controlled
calculation of the fermion \emph{modulus}.  Each physical Weyl arrow is
first paired with its adjoint to form one Hermitian Clifford module.
A standard spin-only bag projector is imposed on that doubled module.
The projector has a uniform principal boundary gap, its square is an
ordinary mixed problem, and its \(A_2\) coefficient has the
Einstein--GHY ratio.

The determinant multiplicity from V134 is retained through the
calculation.  After dividing by that exact multiplicity, the modulus of
the three-generation Standard Model with right-handed neutrinos cancels
the physical gauge/ghost weight-two coefficient of V129.  It does not
cancel the relative collar interface.  This shows that the V130
gauge-sector counterterm cannot be promoted to a full Standard Model
renormalization condition before the Higgs and remaining fermion phase
data are added.

\subsection{Mandatory V135 companion audit}
\label{subsec:V135-audit}

The read-only shared-folder snapshot on 5 September 2026 was
\begin{center}
\begin{tabular}{|p{0.49\textwidth}|r|p{0.31\textwidth}|}
\hline
Companion file & Bytes & SHA--256\\
\hline
\texttt{Renewal\_NS\_Stability\_Research\_Notes\_V31.tex}
& \(887537\) &
\texttt{F1121B62238AD5491BB7CF03635EA9934942EDF573ED8FAAA9EE79E1D38A6696}\\
\hline
\texttt{Renewal\_Yang\_Mills\_Mass\_Gap\_Research\_Notes\_V118.tex}
& \(2770257\) &
\texttt{6B8EB793BAEBB8DE046B12025B8BDF71300DD7BBCD4ED6DD8D84C6D1C8BFE8B1}\\
\hline
\end{tabular}
\end{center}
NS remains byte-identical to the V125 and V130 audits.  Its warning is
unchanged: moving a missing scale factor into a weighted norm does not
pay that factor.

YM V116--V118 now keeps two correlated pair labels inside one physical
positive completion.  The common outer blocks need not form a
rectangular label set.  Absolute values are taken only on genuine
orthogonal blocks, and the two hybrid amplitudes are recombined by a
coefficient-level triangle inequality rather than the false operator
order \(|A+B|\le|A|+|B|\).

\begin{assessment}[V135 audit adjustment]
\label{ass:V135-audit-adjustment}
The fermionic analogue is to retain \(D_V^+\) and \(D_V^-\) in the one
Hermitian block \(\mathbb D_V\) through boundary-domain selection,
squaring, and heat tracing.  Only after the doubled determinant identity
is applied may its coefficient be divided to obtain one Weyl modulus.
Computing two independently optimized chiral heat problems and then
declaring them conjugate would not establish a common determinant line
or boundary domain.
\end{assessment}

\subsection{A gauge-covariant bag on each doubled representation}
\label{subsec:V135-bag}

Let \(\Gamma\) be the Euclidean chirality involution on the four-component
spin module:
\begin{equation}
 \Gamma^*=\Gamma,\qquad
 \Gamma^2=I,\qquad
 \Gamma c(\zeta)+c(\zeta)\Gamma=0.
 \label{eq:V135-chirality}
\end{equation}
For either sign \(\varepsilon\in\{+1,-1\}\), define
\begin{equation}
 \chi_B:=\varepsilon c(n)\Gamma,
 \qquad
 \Pi_\pm^B:=\frac12(I\pm\chi_B).
 \label{eq:V135-bag-projectors}
\end{equation}
The boundary condition is
\begin{equation}
 \Pi_-^B\psi|_\Sigma=0.
 \label{eq:V135-bag-condition}
\end{equation}
All three operators act on spin indices and as the identity on the
internal representation \(V\).

\begin{lemma}[Bag involution and Green neutrality]
\label{lem:V135-bag-involution}
One has
\begin{equation}
 \chi_B^*=\chi_B,\qquad
 \chi_B^2=I,\qquad
 \chi_Bc(n)+c(n)\chi_B=0.
 \label{eq:V135-bag-identities}
\end{equation}
Consequently \(c(n)\) exchanges the two bag subspaces and the allowed
space \(\operatorname{Ran}\Pi_+^B\) is maximal neutral for the Dirac
Green form.
\end{lemma}

\begin{proof}
Using \(c(n)^*=-c(n)\) and
\(\Gamma c(n)=-c(n)\Gamma\),
\[
 (c(n)\Gamma)^*
 =\Gamma[-c(n)]
 =c(n)\Gamma.
\]
Also
\[
 (c(n)\Gamma)^2
 =-c(n)^2\Gamma^2=I.
\]
The anticommutator with \(c(n)\) vanishes by the same identities.
Thus \(c(n)\) maps the \(+1\) eigenspace of \(\chi_B\) unitarily onto its
\(-1\) eigenspace.  The Green pairing of two allowed vectors vanishes,
and the allowed rank is half the spinor rank, proving maximality.
\end{proof}

\begin{corollary}[Gauge covariance]
\label{cor:V135-bag-gauge}
On the Hermitian double
\((S^+\otimes V)\oplus(S^-\otimes V)\), the projectors
\(\Pi_\pm^B\otimes I_V\) commute with the gauge representation.
\end{corollary}

\begin{proof}
The bag projectors act only on the spin factor.  The doubled internal
representation is the same \(V\) on both spin chiralities by
Proposition~\ref{prop:V134-Hermitian-double}.
\end{proof}

\subsection{Uniform strong ellipticity}
\label{subsec:V135-ellipticity}

\begin{lemma}[Bag and normal evolution anticommute]
\label{lem:V135-bag-A}
For every tangential \(\xi\),
\begin{equation}
 \chi_BA(\xi)+A(\xi)\chi_B=0,
 \qquad
 A(\xi)=i\,c(n)c(\xi).
 \label{eq:V135-bag-A}
\end{equation}
\end{lemma}

\begin{proof}
Use anticommutation of \(\Gamma\) with both \(c(n)\) and \(c(\xi)\), and
anticommutation of \(c(n)\) with \(c(\xi)\).  Directly,
\[
 (c(n)\Gamma)(i c(n)c(\xi))
 =i\Gamma c(\xi),
\]
while
\[
 (i c(n)c(\xi))(c(n)\Gamma)
 =-i\Gamma c(\xi).
\]
\end{proof}

\begin{theorem}[Exact bag Lopatinski gap]
\label{thm:V135-bag-gap}
Let \(\mathcal C_-(\widehat\xi)\) be the decaying Cauchy-data space of
V133.  For every \(w\in\mathcal C_-(\widehat\xi)\),
\begin{equation}
 \boxed{
 \|\Pi_-^Bw\|^2=\frac12\|w\|^2.}
 \label{eq:V135-bag-gap}
\end{equation}
Hence the boundary map
\(\Pi_-^B:\mathcal C_-(\widehat\xi)\to\operatorname{Ran}\Pi_-^B\)
is an isomorphism with least singular value \(2^{-1/2}\), uniformly in
the boundary point, tangential direction, internal representation, and
bag sign.
\end{theorem}

\begin{proof}
On the decaying space,
\(A(\widehat\xi)w=-w\).  By
Lemma~\ref{lem:V135-bag-A},
\(\chi_Bw\) belongs to the \(+1\) eigenspace of
\(A(\widehat\xi)\), which is orthogonal to the \(-1\) eigenspace.
Thus \(\langle w,\chi_Bw\rangle=0\).  Since
\(\Pi_-^B=(I-\chi_B)/2\),
\[
 \|\Pi_-^Bw\|^2
 =
 \langle w,\Pi_-^Bw\rangle
 =
 \frac12\|w\|^2.
\]
The domain and target have equal half rank, so injectivity with the
displayed norm is bijectivity and gives the stated singular value.
\end{proof}

\begin{corollary}[The V133 transversality gate is passed]
\label{cor:V135-V133-gap-passed}
Every representation-wise doubled bag problem passes the principal
strong-ellipticity gate of Theorem~\ref{thm:V133-Lopatinski} with a
uniform, cutoff-independent floor.
\end{corollary}

\begin{proof}
The boundary-projection form of the Lopatinski condition is equivalent
to the graph form in V133.  Theorem~\ref{thm:V135-bag-gap} supplies its
uniform inverse bound.
\end{proof}

\subsection{The squared bag domain}
\label{subsec:V135-square}

For the massless doubled kinetic operator, impose
\begin{equation}
 \Pi_-^B\psi=0,
 \qquad
 \Pi_-^B\mathbb D_V\psi=0.
 \label{eq:V135-squared-bag-domain}
\end{equation}

\begin{theorem}[Ordinary mixed square and Robin endomorphism]
\label{thm:V135-bag-square}
The domain \eqref{eq:V135-squared-bag-domain} is equivalent to
\begin{equation}
 \boxed{
 \Pi_-^B\psi=0,
 \qquad
 \Pi_+^B
 \left(\nabla_n-\frac12L\right)
 \Pi_+^B\psi=0.}
 \label{eq:V135-bag-mixed}
\end{equation}
Thus
\begin{equation}
 S_B=-\frac12L\,\Pi_+^B,
 \qquad
 P_UB_1P_U=0.
 \label{eq:V135-bag-S}
\end{equation}
The squared problem is in the ordinary mixed class of V129--V132.
\end{theorem}

\begin{proof}
The tangential principal Dirac operator anticommutes with \(\chi_B\) by
Lemma~\ref{lem:V135-bag-A}, so it maps the allowed bag space into the
Dirichlet complement.  After the outer normal Clifford factor is moved
through the second boundary row, its compression to the allowed space
therefore vanishes.  This proves \(P_UB_1P_U=0\).

In an inward orthonormal collar frame, the remaining zeroth-order
spin-connection term in the Dirac normal decomposition is
\(-L/2\).  Commuting the bag projector through
\(\mathbb D_V\) then turns
\(\Pi_-^B\mathbb D_V\psi=0\) into the second row of
\eqref{eq:V135-bag-mixed}.  Equivalently, the shape contribution from
\(\nabla_a n=-L_{ab}e_b\) contracts to
\(-\tfrac12\operatorname{tr}L\) on the rank-half allowed spin space.
\end{proof}

\begin{firewall}[Yukawa terms are not included in the baseline square]
\label{fire:V135-Yukawa-later}
A Higgs--Yukawa endomorphism does not change
Theorem~\ref{thm:V135-bag-gap}, because it is lower order.  It does
change the zeroth-order bulk endomorphism and may add a projected
boundary term to \(S_B\).  V135 computes the massless kinetic curvature
baseline.  Higgs and Yukawa contributions must be inserted after their
own boundary-domain compatibility has been checked.
\end{firewall}

\subsection{Doubled heat coefficient}
\label{subsec:V135-heat}

\begin{lemma}[Lichnerowicz trace]
\label{lem:V135-Lichnerowicz}
For a four-component doubled spinor with internal representation
dimension \(d=\dim_\mathbb C V\), the curvature part of the squared
operator has
\begin{equation}
 E_F=-\frac14R\,I+\textup{(traceless Clifford--gauge curvature)},
 \label{eq:V135-EF}
\end{equation}
and therefore
\begin{equation}
 \operatorname{tr}
 \left(E_F+\frac16R\,I\right)
 =-\frac d3R.
 \label{eq:V135-bulk-trace}
\end{equation}
\end{lemma}

\begin{proof}
The Lichnerowicz formula is
\(\mathbb D_V^2=\nabla^*\nabla+R/4+\tfrac12c(F_V)\).
In the Laplace convention \eqref{eq:V129-Laplace-type}, this gives
\eqref{eq:V135-EF}.  The spin trace of a Clifford two-form is zero.
The remaining trace is
\[
 4d\left(-\frac14+\frac16\right)R
 =-\frac d3R.
\]
\end{proof}

\begin{lemma}[Bag boundary trace]
\label{lem:V135-boundary-trace}
The bag projector ranks are
\(\operatorname{rank}\Pi_+^B=\operatorname{rank}\Pi_-^B=2d\), and
\begin{equation}
 \operatorname{tr}
 \left(
 \frac13L\,I+2S_B
 \right)
 =
 -\frac{2d}{3}L.
 \label{eq:V135-boundary-trace}
\end{equation}
\end{lemma}

\begin{proof}
The two eigenspaces of the bag involution have equal spin rank two.
Using \eqref{eq:V135-bag-S},
\[
 \frac13L(4d)
 +2\left(-\frac12L\right)(2d)
 =
 \frac{4d}{3}L-2dL
 =
 -\frac{2d}{3}L.
\]
\end{proof}

\begin{theorem}[Doubled bag \(A_2\)]
\label{thm:V135-doubled-A2}
The massless doubled bag heat trace has curvature coefficient
\begin{equation}
 \boxed{
 A_{2,F}^{\rm dbl}(V)
 =
 -\frac d3
 \left(
 \int_MR+2\int_\Sigma L
 \right).}
 \label{eq:V135-doubled-A2}
\end{equation}
\end{theorem}

\begin{proof}
Insert Lemmas~\ref{lem:V135-Lichnerowicz} and
\ref{lem:V135-boundary-trace} into the universal weight-two formula
\eqref{eq:V129-mixed-heat}.  The boundary coefficient is twice the bulk
coefficient, giving the display.
\end{proof}

\subsection{One Weyl modulus and the Standard Model count}
\label{subsec:V135-Weyl-count}

\begin{proposition}[One Weyl modulus weight]
\label{prop:V135-one-Weyl}
For one physical complex Weyl arrow with internal dimension \(d\), the
effective modulus is
\begin{equation}
 -\log|\det D_V^+|
 =
 -\frac14\log\det\mathbb D_V^2
 =
 \frac14\int_0^\infty\frac{dt}{t}
 \operatorname{Tr}e^{-t\mathbb D_V^2},
 \label{eq:V135-Weyl-modulus}
\end{equation}
up to the common ultraviolet subtraction and zero-mode quotient.
Its weighted curvature coefficient is therefore
\begin{equation}
 \boxed{
 A_{2,F}^{\rm Weyl\,mod}(V)
 =
 -\frac d{12}
 \left(
 \int_MR+2\int_\Sigma L
 \right).}
 \label{eq:V135-Weyl-A2}
\end{equation}
\end{proposition}

\begin{proof}
Lemma~\ref{lem:V134-determinant-multiplicity} gives
\(\tfrac12\log\det\mathbb D_V^2=2\log|\det D_V^+|\).
The heat representation
\(\log\det P=-\int t^{-1}\operatorname{Tr}e^{-tP}\,dt\)
then proves \eqref{eq:V135-Weyl-modulus}.  Multiply
\eqref{eq:V135-doubled-A2} by \(1/4\).
\end{proof}

For one generation the internal dimensions of the Weyl arrows are
\begin{center}
\begin{tabular}{|c|c|}
\hline
Multiplet & Complex internal dimension\\
\hline
\(Q_L\) & \(3\cdot2=6\)\\
\(u_R\) & \(3\)\\
\(d_R\) & \(3\)\\
\(L_L\) & \(2\)\\
\(e_R\) & \(1\)\\
\(\nu_R\), when present & \(1\)\\
\hline
\end{tabular}
\end{center}

\begin{lemma}[Chiral dimension count]
\label{lem:V135-SM-count}
For \(N_g\) generations,
\begin{equation}
 d_\chi
 =
 \begin{cases}
 16N_g,&\text{with one right-handed neutrino per generation},\\
 15N_g,&\text{in the minimal model}.
 \end{cases}
 \label{eq:V135-SM-count}
\end{equation}
\end{lemma}

\begin{proof}
Sum the six entries in the table, or omit the final entry, and multiply
by \(N_g\).
\end{proof}

\begin{theorem}[Three-generation physical weight-two cancellation]
\label{thm:V135-three-generation-cancel}
For \(N_g=3\) with right-handed neutrinos, the total Weyl-modulus
coefficient is
\begin{equation}
 A_{2,F}^{\rm SM\,mod}
 =
 -4\left(
 \int_MR+2\int_\Sigma L
 \right).
 \label{eq:V135-SM-fermion-A2}
\end{equation}
It cancels the physical relative gauge/ghost coefficient of V129:
\begin{equation}
 \boxed{
 A_{2,\rm BRST}^{\rm phys}
 +A_{2,F}^{\rm SM\,mod}=0.}
 \label{eq:V135-physical-cancellation}
\end{equation}
\end{theorem}

\begin{proof}
Lemma~\ref{lem:V135-SM-count} gives \(d_\chi=48\).
Sum \eqref{eq:V135-Weyl-A2} over all Weyl arrows to obtain
\eqref{eq:V135-SM-fermion-A2}.  Theorem~\ref{thm:V129-Einstein-GHY}
gives the opposite physical gauge/ghost coefficient.
\end{proof}

\begin{corollary}[What remains after auxiliary elimination]
\label{cor:V135-after-interface}
Before adding Higgs/Yukawa curvature terms or the chiral phase, the
gauge/ghost plus Weyl-modulus coefficient after auxiliary elimination is
\begin{equation}
 \boxed{
 A_{2,\rm BRST+F}^{\rm eff}
 =
 -12\int_\Sigma L
 =
 12\int_\Sigma H.}
 \label{eq:V135-after-interface}
\end{equation}
It has no bulk scalar-curvature term.
\end{corollary}

\begin{proof}
Use the physical cancellation
\eqref{eq:V135-physical-cancellation}.  The remaining relative interface
coefficient is \(-12\int_\Sigma L\) by
\eqref{eq:V130-relative-L-coefficient}.
\end{proof}

\begin{corollary}[Minimal-model comparison]
\label{cor:V135-minimal-model}
For three generations without right-handed neutrinos,
\begin{equation}
 A_{2,F}^{\rm min\,mod}
 =
 -\frac{15}{4}
 \left(\int_MR+2\int_\Sigma L\right),
 \label{eq:V135-minimal-fermion}
\end{equation}
and after adding the auxiliary-eliminated gauge/ghost sector,
\begin{equation}
 A_{2,\rm BRST+F}^{\rm min,eff}
 =
 \frac14\int_MR-\frac{23}{2}\int_\Sigma L.
 \label{eq:V135-minimal-total}
\end{equation}
\end{corollary}

\begin{proof}
Now \(d_\chi=45\).  Apply
\eqref{eq:V135-Weyl-A2}, then add the gauge/ghost coefficient
\eqref{eq:V130-combined-A2}.
\end{proof}

\begin{firewall}[The V130 subtraction must wait for the full sector sum]
\label{fire:V135-counterterm-waits}
The V130 term \(+12\int_\Sigma L\) was uniquely forced only after fixing
the gauge/ghost bulk coefficient and imposing its Einstein--GHY ratio.
For the three-generation right-handed-neutrino modulus, the physical
bulk \(R\) coefficient cancels before the interface is added.
Adding the V130 subtraction at this stage would cancel
\eqref{eq:V135-after-interface} and leave zero weight-two curvature
coefficient, not a nonzero induced Newton coupling.  The Higgs sector,
nonminimal scalar coupling, fermion phase/anomaly data, and microscopic
bare gravitational coefficient must therefore be assembled before a
full Standard Model boundary counterterm is selected.
\end{firewall}

\begin{firewall}[Modulus is not the chiral determinant]
\label{fire:V135-modulus-only}
V135 computes only the local ultraviolet coefficient of the chiral
determinant modulus on branch \textup{(D)}.  The determinant/Pfaffian
line phase, eta contribution, global anomaly, and boundary transport are
not supplied by the bag heat trace.  The cancellation
\eqref{eq:V135-physical-cancellation} is a modulus-level local heat
identity and does not trivialize those phase data.
\end{firewall}

\subsection{V135 ledger and next acquisition}
\label{subsec:V135-ledger}

\paragraph{New exact conclusions.}
\begin{enumerate}[label=\textup{(\arabic*)},leftmargin=2.8em]
\item The compulsory audit records the unchanged NS V31 note and current
      YM V118 note.
\item Theorem~\ref{thm:V135-bag-gap} proves the exact uniform
      \(2^{-1/2}\) boundary-projection gap for the doubled bag.
\item Theorem~\ref{thm:V135-bag-square} proves that its square is ordinary
      mixed with \(S_B=-L/2\).
\item Theorem~\ref{thm:V135-doubled-A2} computes the doubled kinetic
      curvature coefficient with the Einstein--GHY ratio.
\item Proposition~\ref{prop:V135-one-Weyl} applies the exact determinant
      multiplicity to one Weyl modulus.
\item Theorem~\ref{thm:V135-three-generation-cancel} proves the local
      weight-two cancellation between the physical gauge/ghost and
      three-generation right-handed-neutrino Weyl modulus.
\item Corollary~\ref{cor:V135-after-interface} isolates the remaining
      collar-interface boundary coefficient.
\end{enumerate}

\paragraph{Next forced step.}
V136 should add the Higgs modulus with a declared local boundary
condition and arbitrary nonminimal curvature coupling \(\xi_H\).  It
must compute both bulk and boundary weight-two coefficients and solve
the algebraic Einstein--GHY matching equation only after adding
\eqref{eq:V135-after-interface}.  Dirichlet and Robin Higgs domains must
be kept as separate branches.  The chiral phase remains an independent
determinant-line card.

\begin{assessment}[V135 continuum status]
\label{ass:V135-continuum-status}
The full Standard--Model--Einstein continuum remains unproved.  V135
constructs a strongly elliptic doubled bag modulus and closes its local
weight-two coefficient, revealing an exact three-generation physical
gauge--fermion cancellation.  The Higgs domain and coupling, chiral
phase/anomaly transport, microscopic gravitational normalization,
interacting nonlinear existence, Einstein stress closure, and the
cofinal continuum remain open.
\end{assessment}

\section*{V135 exact checks and append-only record}
\addcontentsline{toc}{section}{V135 exact checks and append-only record}

The doubled-modulus checksum is
\begin{equation}
 \boxed{
 s_{\min}(\Pi_-^B|_{\mathcal C_-})=2^{-1/2},
 \quad
 A_{2,F}^{\rm Weyl\,mod}(V)
 =-\frac{\dim V}{12}\left(\int_MR+2\int_\Sigma L\right),
 \quad
 A_{2,\rm BRST+F}^{\rm eff}=-12\int_\Sigma L.}
 \label{eq:V135-final-check}
\end{equation}
The append operation preserves the complete V134 prefix byte for byte before
its former live terminator.  The assembled V135 file is checked for one live
final terminator, balanced environments and braces, duplicate labels, newly
introduced unresolved references, display math delimiters, and control
characters before the exact V134 predecessor is removed.

% =============================================================================
% END APPENDED VERSION V135 MATERIAL
% =============================================================================

% =============================================================================
% BEGIN APPENDED VERSION V136 MATERIAL
% =============================================================================

\section{V136: the Higgs Gaussian, its two local boundary branches, and
the exact weight-two matching defect}
\label{sec:V136}

\subsection{Purpose and scope}
\label{subsec:V136-purpose}

V135 left the auxiliary-eliminated gauge, ghost, and three-generation
right-handed-neutrino Weyl-modulus coefficient
\[
 -12\int_\Sigma L.
\]
The next missing local determinant is the Higgs Gaussian.  This version
computes it without identifying Dirichlet and Robin domains.  The Higgs
doublet is realified before the trace, so its fibre rank is four.  The
factor from the real bosonic Gaussian is also retained explicitly.

The result separates three logically different data: the bulk nonminimal
coupling \(\xi_H\), the geometric Robin coefficient \(\sigma_H\), and a
possible microscopic gravitational boundary coefficient.  Their matching
equation is exact.  It is not a dynamical selection rule until a microscopic
renormalization condition fixes one of those data.

\subsection{Realified Higgs Hessian and determinant weight}
\label{subsec:V136-Hessian}

Let \(\mathcal H_{\mathbb R}\) be the realification of the complex Standard
Model Higgs doublet.  Thus
\begin{equation}
 \operatorname{rank}_{\mathbb R}\mathcal H_{\mathbb R}=4.
 \label{eq:V136-Higgs-rank}
\end{equation}
On a fixed Euclidean gauge and metric background, write the block-diagonal
Higgs fluctuation operator as
\begin{equation}
 P_H
 =-D^\mu D_\mu+\xi_H R\,I_4+U_H
 =-\bigl(D^\mu D_\mu+E_H\bigr),
 \qquad
 E_H=-\xi_HR\,I_4-U_H,
 \label{eq:V136-Higgs-operator}
\end{equation}
where \(U_H=U_H^*\) is the Hessian of the noncurvature potential.  In
particular, no multiple of \(R I_4\) is hidden in \(U_H\).

For a positive infrared realization of \(P_H\), define its real bosonic
one-loop heat combination by
\begin{equation}
 K_H^{\rm eff}(t):=-\frac12\operatorname{Tr}e^{-tP_H}.
 \label{eq:V136-Higgs-effective-heat}
\end{equation}
This is the heat integrand of
\(\tfrac12\log\det P_H\).  A positive mass shift may be used to define the
large-time integral; it does not alter the local derivation below, although
its expansion is part of \(U_H\).

\begin{lemma}[Bulk Higgs coefficient]
\label{lem:V136-Higgs-bulk}
The weight-two bulk term in \(K_H^{\rm eff}\) is
\begin{equation}
 \boxed{
 A_{2,H}^{\rm bulk}
 =
 \int_M\left[
 \left(2\xi_H-\frac13\right)R
 +\frac12\operatorname{tr}_{\mathcal H_{\mathbb R}}U_H
 \right].}
 \label{eq:V136-Higgs-bulk}
\end{equation}
\end{lemma}

\begin{proof}
By Theorem~\ref{thm:V129-mixed-heat} and
\eqref{eq:V136-Higgs-operator}, the unweighted scalar heat trace has bulk
coefficient
\[
 \operatorname{tr}\left(E_H+\frac16R I_4\right)
 =4\left(\frac16-\xi_H\right)R-\operatorname{tr}U_H.
\]
Multiplication by the exact Gaussian weight \(-1/2\) in
\eqref{eq:V136-Higgs-effective-heat} gives
\eqref{eq:V136-Higgs-bulk}.
\end{proof}

\begin{firewall}[The block-diagonal hypothesis]
\label{fire:V136-Higgs-block}
Equation~\eqref{eq:V136-Higgs-operator} is exact at the symmetric background
\(\Phi_0=0\), and also in a background gauge whose quadratic gauge--Higgs
mixing has been cancelled.  At a general nonstationary Higgs background the
gauge and scalar fluctuations form one coupled Laplace-type operator.  Its
completed connection and endomorphism must be traced jointly.  V136 does not
discard such mixing; it computes the scalar block that can be inserted only
on a branch where the displayed Hessian is the actual diagonal block.
\end{firewall}

\subsection{Dirichlet branch}
\label{subsec:V136-Dirichlet}

On branch \(\mathrm{(H_D)}\), impose
\begin{equation}
 \varphi|_\Sigma=0,
 \qquad
 \Pi_-=I_4,
 \quad \Pi_+=0,
 \quad \chi=-I_4.
 \label{eq:V136-Higgs-Dirichlet}
\end{equation}
This domain is gauge covariant because the zero section is invariant under
the boundary gauge representation.

\begin{proposition}[Dirichlet area and curvature rows]
\label{prop:V136-Higgs-Dirichlet}
The boundary terms of the effective Higgs heat combination on
\(\mathrm{(H_D)}\) are
\begin{align}
 [t^{1/2}]K_{H,D}^{\rm eff}
 &=
 +\frac{\sqrt{4\pi t}}2\int_\Sigma1,
 \label{eq:V136-Dirichlet-area}\\
 A_{2,H,D}^{\partial}
 &=
 -\frac23\int_\Sigma L.
 \label{eq:V136-Dirichlet-boundary}
\end{align}
Consequently its geometric weight-two coefficient is
\begin{equation}
 \boxed{
 A_{2,H,D}^{\rm geom}
 =
 \left(2\xi_H-\frac13\right)\int_MR
 -\frac23\int_\Sigma L.}
 \label{eq:V136-Dirichlet-geometric}
\end{equation}
\end{proposition}

\begin{proof}
The unweighted boundary image trace in
\eqref{eq:V129-mixed-heat} is
\(\operatorname{tr}\chi=-4\), while its weight-two boundary trace is
\(4L/3\).  Multiplication by \(-1/2\) proves
\eqref{eq:V136-Dirichlet-area} and
\eqref{eq:V136-Dirichlet-boundary}.  Add the curvature part of
Lemma~\ref{lem:V136-Higgs-bulk}.
\end{proof}

\subsection{Robin branch}
\label{subsec:V136-Robin}

On branch \(\mathrm{(H_R)}\), impose
\begin{equation}
 (D_n+S_H)\varphi|_\Sigma=0,
 \qquad
 S_H=\sigma_HL\,I_4+T_H,
 \label{eq:V136-Higgs-Robin}
\end{equation}
where \(S_H=S_H^*\) commutes with the boundary gauge action and
\(T_H\) contains no component proportional to \(L I_4\).  Thus
\(\sigma_H\) records the entire geometric Robin coefficient.

\begin{proposition}[Robin area and curvature rows]
\label{prop:V136-Higgs-Robin}
On \(\mathrm{(H_R)}\),
\begin{align}
 [t^{1/2}]K_{H,R}^{\rm eff}
 &=
 -\frac{\sqrt{4\pi t}}2\int_\Sigma1,
 \label{eq:V136-Robin-area}\\
 A_{2,H,R}^{\partial}
 &=
 -\left(\frac23+4\sigma_H\right)\int_\Sigma L
 -\int_\Sigma\operatorname{tr}T_H.
 \label{eq:V136-Robin-boundary}
\end{align}
Its geometric weight-two coefficient is therefore
\begin{equation}
 \boxed{
 A_{2,H,R}^{\rm geom}
 =
 \left(2\xi_H-\frac13\right)\int_MR
 -\left(\frac23+4\sigma_H\right)\int_\Sigma L.}
 \label{eq:V136-Robin-geometric}
\end{equation}
\end{proposition}

\begin{proof}
For pure Robin data, \(\chi=I_4\), so the unweighted image trace is
\(+4\).  The unweighted weight-two boundary trace is
\[
 \frac43L+2\operatorname{tr}S_H
 =\left(\frac43+8\sigma_H\right)L+2\operatorname{tr}T_H.
\]
Multiply both expressions by \(-1/2\), then add the curvature part of
Lemma~\ref{lem:V136-Higgs-bulk}.
\end{proof}

\begin{corollary}[The Higgs boundary area divergence is independent]
\label{cor:V136-Higgs-area}
Neither local scalar branch preserves the V129 cancellation of the
boundary area coefficient.  Dirichlet and Robin give equal magnitudes and
opposite signs.  Hence a boundary cosmological, or surface-tension,
renormalization is a separate weight-one card and cannot be absorbed into
an Einstein--GHY coefficient.
\end{corollary}

\begin{proof}
Equations~\eqref{eq:V136-Dirichlet-area} and
\eqref{eq:V136-Robin-area} are nonzero and multiply a different power of
the proper-time cutoff from the weight-two curvature terms.
\end{proof}

\subsection{Which Robin coefficient is selected by a quadratic action?}
\label{subsec:V136-action-domain}

Fix the inward normal convention of \eqref{eq:V129-L-convention}.  For a
self-adjoint boundary endomorphism \(B_H\), consider
\begin{align}
 I_H^{(2)}[\varphi]
 :=&\frac12\int_M
 \left(
 |D\varphi|^2
 +\left\langle\varphi,
 (\xi_HR I_4+U_H)\varphi\right\rangle
 \right)
 \notag\\
 &+\frac12\int_\Sigma
 \langle\varphi,B_H\varphi\rangle.
 \label{eq:V136-quadratic-action}
\end{align}

\begin{lemma}[Action-domain relation]
\label{lem:V136-action-domain}
If the boundary value of \(\varphi\) is free, stationarity of
\eqref{eq:V136-quadratic-action} imposes
\begin{equation}
 (D_n+S_H)\varphi=0,
 \qquad S_H=-B_H.
 \label{eq:V136-action-Robin}
\end{equation}
In particular, completing the nonminimal bulk term
\(\tfrac12\xi_H\int_MR|\varphi|^2\) by the scalar GHY partner
\(\xi_H\int_\Sigma L|\varphi|^2\) gives
\begin{equation}
 B_H=2\xi_HL I_4,
 \qquad
 \boxed{\sigma_H=-2\xi_H.}
 \label{eq:V136-natural-Robin}
\end{equation}
\end{lemma}

\begin{proof}
Since the outward unit normal is \(-n\), integration by parts gives the
boundary variation
\[
 -\int_\Sigma
 \langle\delta\varphi,D_n\varphi\rangle
 +\int_\Sigma
 \langle\delta\varphi,B_H\varphi\rangle
 =
 -\int_\Sigma
 \langle\delta\varphi,(D_n-B_H)\varphi\rangle.
\]
Arbitrary boundary variations yield
\((D_n-B_H)\varphi=0\), which is
\eqref{eq:V136-action-Robin}.  The displayed scalar GHY partner equals
\(\tfrac12\int_\Sigma\langle\varphi,
 2\xi_HLI_4\varphi\rangle\), proving
\eqref{eq:V136-natural-Robin}.
\end{proof}

\begin{firewall}[Three different matching statements]
\label{fire:V136-three-matchings}
Gauge covariance permits every scalar \(\sigma_H\).  Action stationarity
relates \(\sigma_H\) to a chosen scalar boundary action.  Heat-coefficient
Einstein--GHY matching gives a third algebraic relation below.  None of
these implications may be reversed without the corresponding hypothesis.
\end{firewall}

\subsection{Combination with V135 and the exact defect}
\label{subsec:V136-combination}

For the three-generation model with one right-handed neutrino per
generation, add the V135 coefficient
\(-12\int_\Sigma L\).  Write
\begin{equation}
 A_{2,\rm pregrav}^{\rm geom}
 =C_R\int_MR+C_L\int_\Sigma L.
 \label{eq:V136-CRCL-def}
\end{equation}

\begin{theorem}[Full pre-gravitational weight-two coefficients]
\label{thm:V136-pregrav-coefficients}
The two local Higgs branches give
\begin{align}
 C_R&=2\xi_H-\frac13,
 \label{eq:V136-CR}\\
 C_L^{D}&=-\frac{38}{3},
 \label{eq:V136-CL-D}\\
 C_L^{R}&=-\frac{38}{3}-4\sigma_H.
 \label{eq:V136-CL-R}
\end{align}
Their Einstein--GHY defects
\(\mathfrak d:=C_L-2C_R\) are
\begin{equation}
 \boxed{
 \mathfrak d_D=-12-4\xi_H,
 \qquad
 \mathfrak d_R=-12-4(\xi_H+\sigma_H).}
 \label{eq:V136-defects}
\end{equation}
Thus, without a microscopic gravitational term or a new local boundary
subtraction, the coefficient has the Einstein--GHY ratio precisely when
\begin{equation}
 \boxed{
 \xi_H=-3\quad\text{on }\mathrm{(H_D)},
 \qquad
 \xi_H+\sigma_H=-3\quad\text{on }\mathrm{(H_R)}.}
 \label{eq:V136-no-counterterm-match}
\end{equation}
\end{theorem}

\begin{proof}
Add \eqref{eq:V135-after-interface} to
\eqref{eq:V136-Dirichlet-geometric} or
\eqref{eq:V136-Robin-geometric}.  This proves
\eqref{eq:V136-CR}--\eqref{eq:V136-CL-R}.  Direct subtraction gives
\eqref{eq:V136-defects}.  The Einstein--GHY ratio is exactly
\(\mathfrak d=0\), which gives
\eqref{eq:V136-no-counterterm-match}.
\end{proof}

\begin{corollary}[Required local boundary subtraction]
\label{cor:V136-required-subtraction}
If \(c_\Sigma\int_\Sigma L\) is added while the induced bulk coefficient
is held fixed, the unique matching values are
\begin{equation}
 \boxed{
 c_\Sigma^D=12+4\xi_H,
 \qquad
 c_\Sigma^R=12+4(\xi_H+\sigma_H).}
 \label{eq:V136-required-subtraction}
\end{equation}
In particular, the V130 gauge-sector value \(c_\Sigma=12\) is compatible
with the completed Standard Model coefficient precisely for
\begin{equation}
 \xi_H=0\quad\text{on }\mathrm{(H_D)},
 \qquad
 \xi_H+\sigma_H=0\quad\text{on }\mathrm{(H_R)}.
 \label{eq:V136-V130-compatibility}
\end{equation}
On the action-compatible Robin branch
\eqref{eq:V136-natural-Robin}, the latter condition forces
\(\xi_H=\sigma_H=0\).
\end{corollary}

\begin{proof}
Adding \(c_\Sigma\) changes the defect to
\(\mathfrak d+c_\Sigma\).  Set it to zero and use
\eqref{eq:V136-defects}.  Substituting \(c_\Sigma=12\) proves
\eqref{eq:V136-V130-compatibility}; substituting
\(\sigma_H=-2\xi_H\) proves the final statement.
\end{proof}

\begin{corollary}[Microscopic gravitational matching equation]
\label{cor:V136-microscopic-match}
Let the microscopic weight-two gravitational row be
\begin{equation}
 g_R\int_MR+g_L\int_\Sigma L.
 \label{eq:V136-microscopic-row}
\end{equation}
Then the total coefficient has the Einstein--GHY ratio if and only if
\begin{equation}
 \boxed{
 g_L-2g_R=
 \begin{cases}
 12+4\xi_H,&\mathrm{(H_D)},\\
 12+4(\xi_H+\sigma_H),&\mathrm{(H_R)}.
 \end{cases}}
 \label{eq:V136-microscopic-equation}
\end{equation}
This is one affine constraint on the microscopic gravitational data; it
does not determine \(g_R\), \(g_L\), and \(\xi_H\) separately.
\end{corollary}

\begin{proof}
The total defect is
\((g_L-2g_R)+(C_L-2C_R)\).  Set it to zero and apply
\eqref{eq:V136-defects}.
\end{proof}

\begin{corollary}[Action-compatible branch without subtraction]
\label{cor:V136-natural-no-subtraction}
On the action-compatible Robin branch \(\sigma_H=-2\xi_H\), and with no
microscopic gravitational defect, the induced coefficient has the
Einstein--GHY ratio if and only if
\begin{equation}
 \xi_H=3,
 \qquad \sigma_H=-6.
 \label{eq:V136-natural-no-subtraction}
\end{equation}
This algebraic value is not asserted to be physically selected.
\end{corollary}

\begin{proof}
Substitute \(\sigma_H=-2\xi_H\) into the Robin equation in
\eqref{eq:V136-no-counterterm-match}.
\end{proof}

\subsection{Proper-time form and separated lower-order rows}
\label{subsec:V136-proper-time}

\begin{proposition}[Higgs ultraviolet rows through weight two]
\label{prop:V136-proper-time}
With the proper-time convention of \eqref{eq:V129-proper-time-def}, the
Higgs determinant contributes
\begin{align}
 \Gamma_{H,D}(\Lambda)
 =\frac1{(4\pi)^2}\Bigg[&
 -\Lambda^4\int_M1
 +\frac{\sqrt{4\pi}}3\Lambda^3\int_\Sigma1
 \notag\\
 &+\Lambda^2\left{
 \left(2\xi_H-\frac13\right)\int_MR
 -\frac23\int_\Sigma L
 +\frac12\int_M\operatorname{tr}U_H
 \right}\Bigg]
 +O(\Lambda),
 \label{eq:V136-proper-time-D}
\end{align}
and
\begin{align}
 \Gamma_{H,R}(\Lambda)
 =\frac1{(4\pi)^2}\Bigg[&
 -\Lambda^4\int_M1
 -\frac{\sqrt{4\pi}}3\Lambda^3\int_\Sigma1
 \notag\\
 &+\Lambda^2\left{
 \left(2\xi_H-\frac13\right)\int_MR
 -\left(\frac23+4\sigma_H\right)\int_\Sigma L
 +\frac12\int_M\operatorname{tr}U_H
 -\int_\Sigma\operatorname{tr}T_H
 \right}\Bigg]
 +O(\Lambda).
 \label{eq:V136-proper-time-R}
\end{align}
The potential and nongeometric Robin rows are displayed separately and do
not enter \eqref{eq:V136-defects}.
\end{proposition}

\begin{proof}
The rank-four leading heat coefficient, multiplied by \(-1/2\), is
\(-2\int_M1\).  Its \(t^{-3}\) proper-time integral contributes
\(-\Lambda^4\int_M1\).  Equations
\eqref{eq:V136-Dirichlet-area} and
\eqref{eq:V136-Robin-area}, together with
\(\int_{\Lambda^{-2}}t^{-5/2}dt=(2/3)\Lambda^3+O(1)\), give the area
rows.  The weight-two integral is
\(\int_{\Lambda^{-2}}t^{-2}dt=\Lambda^2+O(1)\).  Insert
Lemma~\ref{lem:V136-Higgs-bulk} and Propositions
\ref{prop:V136-Higgs-Dirichlet}--\ref{prop:V136-Higgs-Robin}.
\end{proof}

\begin{firewall}[A coefficient identity is not induced gravity]
\label{fire:V136-not-induced-gravity}
The affine equations \eqref{eq:V136-required-subtraction} and
\eqref{eq:V136-microscopic-equation} classify weight-two local matching.
They do not prove a positive Newton constant, select a subtraction scheme,
or determine the finite renormalized coupling.  The chiral phase, the
coupled gauge--Higgs Hessian away from the diagonal branch, and the renewal
microscopic normalization remain outside this calculation.
\end{firewall}

\subsection{V136 ledger and next acquisition}
\label{subsec:V136-ledger}

\paragraph{New exact conclusions.}
\begin{enumerate}[label=\textup{(\arabic*)},leftmargin=2.8em]
\item The real Higgs fibre has rank four and its Gaussian contributes the
      exact determinant weight \(-1/2\).
\item Lemma~\ref{lem:V136-Higgs-bulk} computes the nonminimal bulk
      coefficient without hiding the potential Hessian.
\item Propositions~\ref{prop:V136-Higgs-Dirichlet} and
      \ref{prop:V136-Higgs-Robin} keep the two strongly elliptic local
      scalar domains separate and expose their opposite area rows.
\item Lemma~\ref{lem:V136-action-domain} derives the Robin coefficient from
      a declared quadratic boundary action.
\item Theorem~\ref{thm:V136-pregrav-coefficients} computes the complete
      gauge--ghost--fermion-modulus--Higgs weight-two defect after the V130
      interface conversion.
\item Corollaries~\ref{cor:V136-required-subtraction} and
      \ref{cor:V136-microscopic-match} solve the local and microscopic
      Einstein--GHY matching equations exactly.
\end{enumerate}

\paragraph{Next forced step.}
The Higgs potential in V136 is diagonal only on the stated background-gauge
branch.  V137 should form the full gauge--Higgs quadratic Hessian about a
general stationary background, identify the background gauge which removes
the first-order mixing, and prove that its ghost, Goldstone, and physical
scalar boundary domains share a common strongly elliptic realization.  If
that domain cannot be made compatible with the V129 relative gauge domain,
the programme must go upstream to the coupled matrix boundary symbol rather
than add diagonal heat coefficients.

\begin{assessment}[V136 continuum status]
\label{ass:V136-continuum-status}
The full Standard--Model--Einstein continuum remains unproved.  V136 closes
the diagonal Higgs Gaussian and its exact local weight-two matching defect.
The general coupled gauge--Higgs Hessian, chiral phase and anomaly transport,
microscopic gravitational normalization, interacting nonlinear existence,
Einstein stress closure, and cofinal continuum remain open.
\end{assessment}

\section*{V136 exact checks and append-only record}
\addcontentsline{toc}{section}{V136 exact checks and append-only record}

The Higgs matching checksum is
\begin{equation}
 \boxed{
 \mathfrak d_D=-12-4\xi_H,
 \qquad
 \mathfrak d_R=-12-4(\xi_H+\sigma_H),
 \qquad
 g_L-2g_R=-\mathfrak d.}
 \label{eq:V136-final-check}
\end{equation}
The append operation preserves the complete V135 prefix byte for byte before
its former live terminator.  The assembled V136 file is checked for one live
final terminator, balanced environments and braces, duplicate labels, newly
introduced unresolved references, display math delimiters, and control
characters before the exact V135 predecessor is removed.  The companion
files are not modified.

% =============================================================================
% END APPENDED VERSION V136 MATERIAL
% =============================================================================

% =============================================================================
% BEGIN APPENDED VERSION V137 MATERIAL
% =============================================================================

\section{V137: the coupled gauge--Higgs Hessian, the broken-boundary
Robin obstruction, and its matrix repair}
\label{sec:V137}

\subsection{Purpose}
\label{subsec:V137-purpose}

V136 computed the Higgs Gaussian on a diagonal Hessian branch and left the
general gauge--Higgs mixing as the next gate.  V137 now expands around a
general smooth background, uses one background gauge to cancel the entire
first-order off-diagonal symbol, and keeps the remaining zeroth-order
coupling inside one matrix Laplace operator.

There is a boundary distinction.  Higgs Dirichlet data are compatible with
the relative gauge and Dirichlet ghost traces.  A separate Higgs Robin row
is not BRST stable at a boundary point where the Higgs orbit map is nonzero.
The missing normal ghost jet cancels only after the Higgs Robin row is
coupled to the normal gauge trace.  The resulting self-adjoint matrix Robin
block is strongly elliptic and has the same weight-two trace as the
separately computed rows, because its off-diagonal entries have zero matrix
trace.

\subsection{Orbit map and background gauge}
\label{subsec:V137-background-gauge}

Let \(G\) be compact, let \(\mathfrak g\) carry an invariant inner product,
and let \(\rho:\mathfrak g\to\mathfrak{so}(\mathcal H_{\mathbb R})\) be
the realified Higgs representation.  At a background Higgs field \(\Phi\),
choose the orbit-map sign so that
\begin{equation}
 \tau_\Phi:\mathfrak g\longrightarrow\mathcal H_{\mathbb R},
 \qquad
 s\varphi=-\tau_\Phi c,
 \qquad
 sa=Dc.
 \label{eq:V137-orbit-map}
\end{equation}
Here \(a\), \(\varphi\), and \(c\) are respectively the gauge, Higgs, and
ghost fluctuations.  The adjoint of the orbit map is denoted
\(\tau_\Phi^*\).

Define the background gauge functional
\begin{equation}
 \mathcal F_\Phi(a,\varphi)
 :=D^*a-\tau_\Phi^*\varphi
 \label{eq:V137-gauge-functional}
\end{equation}
and add \(\tfrac12\|\mathcal F_\Phi\|^2\) to the Euclidean quadratic
action.

\begin{lemma}[Cancellation of the derivative mixing]
\label{lem:V137-derivative-cancellation}
The combined Yang--Mills--Higgs Hessian in the gauge
\eqref{eq:V137-gauge-functional} has scalar principal symbol
\begin{equation}
 \sigma_2(P_{AH})(x,\zeta)
 =|\zeta|^2
 I_{(T^*M\otimes\mathfrak g)\oplus\mathcal H_{\mathbb R}}.
 \label{eq:V137-scalar-principal-symbol}
\end{equation}
All uncancelled gauge--Higgs off-diagonal terms are bundle
endomorphisms, linear in \(D\Phi\) or in zeroth-order background data.
\end{lemma}

\begin{proof}
The linearized covariant Higgs derivative has the form
\[
 D\varphi+\tau_\Phi(a).
\]
The cross term in its squared norm is
\(\langle D\varphi,\tau_\Phi(a)\rangle\).  Keeping only the derivative
falling on the fluctuation and integrating by parts gives
\[
 \langle D^*a,\tau_\Phi^*\varphi\rangle.
\]
The cross term in
\(\tfrac12\|D^*a-\tau_\Phi^*\varphi\|^2\) is its negative.  Hence the
off-diagonal first-order fluctuation terms cancel exactly.  If the
derivative instead falls on \(\tau_\Phi\), it produces a coefficient
linear in \(D\Phi\) multiplying \(a\) and \(\varphi\), with no derivative
on either fluctuation.  The usual Feynman background gauge also cancels the
nonminimal \(D_\mu D_\nu\) vector symbol.  Both diagonal second-order
symbols are therefore \(|\zeta|^2I\), proving the claim.
\end{proof}

\begin{proposition}[Ghost operator]
\label{prop:V137-ghost-operator}
The linearization of \(\mathcal F_\Phi\) along
\eqref{eq:V137-orbit-map} is
\begin{equation}
 P_c c
 =D^*Dc+\tau_\Phi^*\tau_\Phi c
 +Z_c(D\Phi)c,
 \label{eq:V137-ghost-operator}
\end{equation}
where \(Z_c(D\Phi)\) is zeroth order and vanishes when the background
orbit map is covariantly constant.  In particular, \(P_c\) is of Laplace
type and has the same scalar principal symbol as the V129 relative ghost.
\end{proposition}

\begin{proof}
Apply \(s\) to \eqref{eq:V137-gauge-functional}.  The terms in which a
derivative hits the ghost are
\(D^*Dc+\tau_\Phi^*\tau_\Phi c\).  Derivatives of the background orbit map
are multiplication operators on \(c\), giving \(Z_c(D\Phi)\).
\end{proof}

\subsection{One matrix Laplace operator}
\label{subsec:V137-matrix-Laplace}

By Lemma~\ref{lem:V137-derivative-cancellation}, the bosonic Hessian can be
written in the convention of \eqref{eq:V129-Laplace-type} as
\begin{equation}
 P_{AH}
 =-\left(\nabla_{AH}^2+E_{AH}\right),
 \qquad
 E_{AH}=
 \begin{pmatrix}
 E_A&Z_\Phi^*\\
 Z_\Phi&E_H^{\rm bg}
 \end{pmatrix}.
 \label{eq:V137-matrix-Laplace}
\end{equation}
The diagonal scalar block contains
\begin{equation}
 E_H^{\rm bg}=-\xi_HR I_4-U_H-U_{\rm gf},
 \label{eq:V137-Higgs-diagonal-E}
\end{equation}
where \(U_{\rm gf}\) is the Goldstone potential generated by the gauge
functional.  The vector and ghost diagonal blocks likewise contain orbit
mass endomorphisms.  None of these orbit or potential terms is declared to
be a multiple of \(R\).

\begin{lemma}[Off-diagonal endomorphisms do not enter \(A_2\)]
\label{lem:V137-offdiag-trace}
For the matrix operator \eqref{eq:V137-matrix-Laplace},
\begin{equation}
 \operatorname{tr}E_{AH}
 =\operatorname{tr}E_A+
  \operatorname{tr}E_H^{\rm bg}.
 \label{eq:V137-offdiag-trace}
\end{equation}
Consequently the off-diagonal background-gradient term \(Z_\Phi\) makes
no contribution to the weight-two interior heat coefficient.
\end{lemma}

\begin{proof}
The fibre trace of a block matrix is the sum of its diagonal-block traces.
The universal interior coefficient in
Theorem~\ref{thm:V129-mixed-heat} is
\(\operatorname{tr}(E+RI/6)\), so only those diagonal traces occur.
\end{proof}

\begin{firewall}[Potential rows remain present]
\label{fire:V137-potentials-present}
Lemma~\ref{lem:V137-offdiag-trace} removes only the off-diagonal block from
the trace.  The orbit masses, Goldstone gauge-fixing potential, Higgs
potential Hessian, and background gauge curvature remain in the diagonal
operators.  Their field-dependent weight-two rows are not scalar-curvature
or mean-curvature coefficients and are not part of the Einstein--GHY defect
computed here.
\end{firewall}

\subsection{Dirichlet Higgs branch on a broken background}
\label{subsec:V137-Dirichlet}

Retain the V129 relative gauge and ghost traces and impose Higgs Dirichlet
data:
\begin{equation}
 a_\top=0,
 \qquad
 (D_n+H)a_n=0,
 \qquad
 c=0,
 \qquad
 \varphi=0
 \quad\text{on }\Sigma.
 \label{eq:V137-Dirichlet-domain}
\end{equation}

\begin{theorem}[Strong ellipticity of the Dirichlet Higgs branch]
\label{thm:V137-Dirichlet-ellipticity}
The matrix operator \eqref{eq:V137-matrix-Laplace} with
\eqref{eq:V137-Dirichlet-domain} is a self-adjoint strongly elliptic mixed
problem whenever its diagonal connections and endomorphisms are smooth and
self-adjoint.  Its pure \(R,L\) weight-two coefficient is exactly the
Dirichlet coefficient in Theorem~\ref{thm:V136-pregrav-coefficients}.
\end{theorem}

\begin{proof}
The boundary Green form depends only on the scalar principal part.  The
tangential gauge and Higgs blocks are Dirichlet, while the normal gauge
block is Robin with the self-adjoint scalar endomorphism \(H=-L\).
Zeroth-order off-diagonal bulk entries create no boundary Green term.
Thus the domain is self-adjoint.  Freezing at a boundary point, the decaying
normal solution of the parameter Laplacian is multiplied by
\(e^{-\kappa r}\).  Dirichlet evaluation is injective on the Dirichlet
blocks, and the Robin principal row is \(-\kappa\) on the complementary
block.  This proves the complementing condition uniformly on every closed
parameter cone away from the origin.

At weight two the off-diagonal bulk trace vanishes by
Lemma~\ref{lem:V137-offdiag-trace}.  The boundary projector and diagonal
Robin trace are exactly those used in V129 and in the Dirichlet branch of
V136.  Hence their pure geometric coefficients are unchanged.
\end{proof}

\begin{corollary}[First boundary BRST row]
\label{cor:V137-Dirichlet-BRST}
The Higgs Dirichlet trace is stable under the boundary BRST transformation
with ghost Dirichlet data:
\begin{equation}
 c|_\Sigma=0
 \quad\Longrightarrow\quad
 (s\varphi)|_\Sigma=-\tau_\Phi c|_\Sigma=0.
 \label{eq:V137-Dirichlet-BRST}
\end{equation}
\end{corollary}

\subsection{No-go for an uncoupled Higgs Robin row}
\label{subsec:V137-Robin-no-go}

Suppose instead that one appends to the V129 relative gauge/Dirichlet ghost
domain the separate Higgs Robin condition
\begin{equation}
 B_H^{\rm sep}\varphi
 :=(D_n+S_H)\varphi=0.
 \label{eq:V137-separate-Robin}
\end{equation}

\begin{theorem}[Broken-boundary Robin obstruction]
\label{thm:V137-Robin-obstruction}
At a boundary point \(y\) for which
\(\tau_{\Phi(y)}\ne0\), the separate Robin row
\eqref{eq:V137-separate-Robin} is not preserved by off-shell BRST
transformations whose ghost obeys only Dirichlet data.  More precisely,
\begin{equation}
 \left.s(B_H^{\rm sep}\varphi)\right|_y
 =-\tau_{\Phi(y)}(D_nc)|_y.
 \label{eq:V137-Robin-obstruction}
\end{equation}
The right side can be arbitrary in
\(\operatorname{Ran}\tau_{\Phi(y)}\).
\end{theorem}

\begin{proof}
Differentiate \(s\varphi=-\tau_\Phi c\):
\[
 (D_n+S_H)s\varphi
 =-(D_n\tau_\Phi)c-\tau_\Phi D_nc-S_H\tau_\Phi c.
\]
Both terms containing the undifferentiated ghost vanish at the boundary,
leaving \eqref{eq:V137-Robin-obstruction}.  The normal derivative of a
Dirichlet function is an unrestricted boundary jet: in a collar, for any
smooth compactly supported \(u(y)\), the ghost
\(c(r,y)=r\chi(r)u(y)\) has zero trace and normal derivative \(u\).
Choosing \(u\) outside the kernel of \(\tau_{\Phi(y)}\) proves the no-go.
\end{proof}

\begin{corollary}[The only uncoupled escape]
\label{cor:V137-uncoupled-escape}
The separate Robin row is stable for all Dirichlet ghost jets only on the
boundary-unbroken locus
\begin{equation}
 \tau_\Phi|_\Sigma=0.
 \label{eq:V137-boundary-unbroken}
\end{equation}
Imposing both \(c|_\Sigma=0\) and \(D_nc|_\Sigma=0\) is not an admissible
replacement, because it supplies two boundary rows for one second-order
ghost equation.
\end{corollary}

\begin{proof}
The first assertion follows from the arbitrariness of the Dirichlet normal
jet in Theorem~\ref{thm:V137-Robin-obstruction}.  The second follows from
the rank count for a scalar second-order boundary problem and can also be
seen on the half-line: a decaying solution has only one boundary amplitude,
so simultaneous value and normal-value constraints overdetermine it.
\end{proof}

\subsection{Coupled normal-gauge/Higgs Robin repair}
\label{subsec:V137-matrix-Robin}

Keep \(a_\top=0\) and \(c=0\), but put the normal gauge and Higgs traces in
one Robin block.  Define
\begin{equation}
 \mathcal B_+
 \binom{a_n}{\varphi}
 :=
 \left[
 D_n+
 \begin{pmatrix}
 H I_{\mathfrak g}&\tau_\Phi^*\\
 \tau_\Phi&S_H
 \end{pmatrix}
 \right]
 \binom{a_n}{\varphi}=0.
 \label{eq:V137-matrix-Robin}
\end{equation}
The lower row is
\begin{equation}
 (D_n+S_H)\varphi+\tau_\Phi a_n=0.
 \label{eq:V137-repaired-Higgs-row}
\end{equation}

\begin{lemma}[Cancellation of the missing ghost jet]
\label{lem:V137-ghost-jet-cancel}
Under \eqref{eq:V137-orbit-map}, the BRST variation of the repaired Higgs
row vanishes at the boundary for every Dirichlet ghost jet:
\begin{equation}
 \left.s\left((D_n+S_H)\varphi+\tau_\Phi a_n\right)\right|_\Sigma=0.
 \label{eq:V137-ghost-jet-cancel}
\end{equation}
\end{lemma}

\begin{proof}
The first term varies as in the proof of
Theorem~\ref{thm:V137-Robin-obstruction}, giving
\(-\tau_\Phi D_nc\) after the undifferentiated ghost is set to zero.
The variation of \(\tau_\Phi a_n\) contributes
\(+\tau_\Phi D_nc\); variations of the background orbit map are
proportional to \(c|_\Sigma\) and vanish.  The two normal-jet terms cancel
before any estimate is taken.
\end{proof}

\begin{theorem}[Self-adjoint strong ellipticity of the repaired block]
\label{thm:V137-matrix-Robin-ellipticity}
Assume \(S_H=S_H^*\).  The bosonic matrix operator
\eqref{eq:V137-matrix-Laplace}, with tangential gauge Dirichlet data and
the matrix Robin row \eqref{eq:V137-matrix-Robin}, is self-adjoint and
strongly elliptic at every fixed smooth cutoff.  The off-diagonal orbit
map does not alter the principal complementing determinant.
\end{theorem}

\begin{proof}
The Robin endomorphism in \eqref{eq:V137-matrix-Robin} is self-adjoint
because its two off-diagonal entries are adjoints.  The Green boundary
form therefore vanishes on the domain.  In the frozen parameter problem,
the matrix multiplying the decaying amplitude is
\[
 -\kappa I+
 \begin{pmatrix}
 H I&\tau_\Phi^*\\
 \tau_\Phi&S_H
 \end{pmatrix}.
\]
The second matrix is order zero in the joint covariable, whereas
\(-\kappa I\) is the homogeneous principal boundary symbol.  Hence the
principal complementing map is \(-\kappa I\), invertible for every
nonzero point of the elliptic parameter cone.  The tangential gauge block
is Dirichlet and is independently complementing.  Standard compactness of
the normalized parameter sphere gives the uniform principal inverse bound.
\end{proof}

\begin{proposition}[Cofinal semiboundedness criterion]
\label{prop:V137-cofinal-form}
Let \(q\) index the renewal cutoffs.  If the negative parts of the matrix
Robin endomorphisms satisfy a uniform bound
\begin{equation}
 \sup_q\left\|
 \begin{pmatrix}
 H_q I&\tau_{\Phi_q}^*\\
 \tau_{\Phi_q}&S_{H,q}
 \end{pmatrix}_{-}
 \right\|_{L^\infty(\Sigma_q)}<\infty
 \label{eq:V137-uniform-Robin-bound}
\end{equation}
on collars with a uniform trace inequality, then the repaired quadratic
forms are uniformly lower semibounded after one common infrared shift.
\end{proposition}

\begin{proof}
For every \(\epsilon>0\), the uniform collar trace inequality gives
\[
 \|u\|_{L^2(\Sigma_q)}^2
 \le \epsilon\|\nabla u\|_{L^2(M_q)}^2
 +C_\epsilon\|u\|_{L^2(M_q)}^2
\]
with \(C_\epsilon\) independent of \(q\).  Pairing the boundary value with
the negative part in \eqref{eq:V137-uniform-Robin-bound} and choosing
\(\epsilon\) small absorbs its gradient contribution.  The remaining
uniform multiple of the bulk \(L^2\) norm is removed by one common shift.
\end{proof}

\begin{firewall}[Principal ellipticity is not a cofinal form bound]
\label{fire:V137-ellipticity-versus-form}
Theorem~\ref{thm:V137-matrix-Robin-ellipticity} is insensitive to the size
of the order-zero orbit map at any fixed cutoff.  A sequence with
\(\|\tau_{\Phi_q}\|_\infty\to\infty\) can nevertheless lose every common
lower form bound.  Proposition~\ref{prop:V137-cofinal-form} is therefore a
separate cofinal hypothesis, not a consequence of strong ellipticity.
\end{firewall}

\subsection{Persistence of the V136 geometric trace}
\label{subsec:V137-trace-persistence}

\begin{theorem}[Coupled weight-two trace invariance]
\label{thm:V137-coupled-A2}
For the background-gauge Hessian and either admissible boundary branch:
\begin{enumerate}[label=\textup{(\roman*)},leftmargin=2.8em]
\item the Dirichlet Higgs domain \eqref{eq:V137-Dirichlet-domain}; or
\item the repaired matrix Robin domain \eqref{eq:V137-matrix-Robin}, with
      \(S_H=\sigma_HLI_4+T_H\),
\end{enumerate}
the pure \(R,L\) part of the one-loop gauge--ghost--Higgs heat combination
is the sum of the V129 gauge/ghost row and the corresponding V136 Higgs row.
After the V135 Weyl modulus and relative interface are included, its
Einstein--GHY defect remains
\begin{equation}
 \boxed{
 \mathfrak d_D=-12-4\xi_H,
 \qquad
 \mathfrak d_R=-12-4(\xi_H+\sigma_H).}
 \label{eq:V137-defect-persistence}
\end{equation}
\end{theorem}

\begin{proof}
In the interior, Lemma~\ref{lem:V137-offdiag-trace} removes the
off-diagonal endomorphism from \(A_2\).  The remaining orbit and potential
endomorphisms contain no declared multiple of \(R\), so the pure scalar
curvature trace is the diagonal V129 plus V136 trace.

At the boundary, the matrix Robin coefficient in
Theorem~\ref{thm:V129-mixed-heat} depends on
\(\operatorname{tr}S\).  The two orbit-map blocks in
\eqref{eq:V137-matrix-Robin} are off diagonal and have zero matrix trace.
The diagonal entries are exactly the normal-gauge value \(H=-L\) of V129
and the Higgs value \(\sigma_HL I_4+T_H\) of V136.  Dirichlet ranks are
also unchanged.  Therefore the separately derived pure \(R,L\) rows add
without approximation.  The V135 fermion and interface rows then give
\eqref{eq:V137-defect-persistence} by
Theorem~\ref{thm:V136-pregrav-coefficients}.
\end{proof}

\begin{firewall}[Extent of BRST closure]
\label{fire:V137-BRST-extent}
Lemma~\ref{lem:V137-ghost-jet-cancel} proves the new scalar Robin-row
incidence which failed in Theorem~\ref{thm:V137-Robin-obstruction}.  The
variation of the normal gauge-fixing row still involves the second normal
ghost jet, as it already does for the relative gauge domain.  Full off-shell
BRST closure requires the antighost and multiplier boundary rows and their
quadratic action.  V137 does not replace that extended complex by an
on-shell ghost equation.
\end{firewall}

\subsection{V137 ledger and next acquisition}
\label{subsec:V137-ledger}

\paragraph{New exact conclusions.}
\begin{enumerate}[label=\textup{(\arabic*)},leftmargin=2.8em]
\item Lemma~\ref{lem:V137-derivative-cancellation} proves that one
      background gauge removes the full first-order gauge--Higgs mixing.
\item Lemma~\ref{lem:V137-offdiag-trace} shows why the residual matrix
      coupling does not alter \(A_2\).
\item Theorem~\ref{thm:V137-Dirichlet-ellipticity} closes the general
      background Dirichlet Higgs branch.
\item Theorem~\ref{thm:V137-Robin-obstruction} gives an exact normal-ghost
      jet no-go for a separate Robin Higgs row on a broken boundary.
\item Lemma~\ref{lem:V137-ghost-jet-cancel} repairs that row before norms by
      coupling it to the normal gauge trace.
\item Theorem~\ref{thm:V137-matrix-Robin-ellipticity} and
      Proposition~\ref{prop:V137-cofinal-form} separate fixed-cutoff
      ellipticity from the uniform form bound needed cofinally.
\item Theorem~\ref{thm:V137-coupled-A2} proves persistence of the V136
      Einstein--GHY defect on both admissible coupled branches.
\end{enumerate}

\paragraph{Next forced step.}
The local modulus coefficients are now closed through weight two on two
explicit gauge--Higgs domains.  V138 should return to the chiral
determinant-line branch left open in V134--V135.  It must compute the local
four-dimensional Standard Model anomaly traces representation by
representation, including a right-handed neutrino, and distinguish
perturbative gauge and mixed gravitational cancellation from the global
\(SU(2)\) mod-two obstruction and from a boundary eta trivialization.

\begin{assessment}[V137 continuum status]
\label{ass:V137-continuum-status}
The full Standard--Model--Einstein continuum remains unproved.  V137 closes
the general-background gauge--Higgs weight-two trace and supplies a repaired
strongly elliptic Robin block.  The extended off-shell BRST boundary complex,
chiral determinant phase and global anomaly transport, microscopic
gravitational normalization, nonlinear existence, Einstein stress closure,
and cofinal interacting limit remain open.
\end{assessment}

\section*{V137 exact checks and append-only record}
\addcontentsline{toc}{section}{V137 exact checks and append-only record}

The coupled-domain checksum is
\begin{equation}
 \boxed{
 \sigma_2(P_{AH})=|\zeta|^2I,
 \quad
 s\bigl((D_n+S_H)\varphi+\tau_\Phi a_n\bigr)|_\Sigma=0,
 \quad
 \operatorname{tr}
 \begin{pmatrix}H I&\tau_\Phi^*\\
 \tau_\Phi&S_H\end{pmatrix}
 =H\dim\mathfrak g+\operatorname{tr}S_H.}
 \label{eq:V137-final-check}
\end{equation}
The append operation preserves the complete V136 prefix byte for byte before
its former live terminator.  The assembled V137 file is checked for one live
final terminator, balanced environments and braces, duplicate labels, newly
introduced unresolved references, display math delimiters, and control
characters before the exact V136 predecessor is removed.  No companion file
is modified.

% =============================================================================
% END APPENDED VERSION V137 MATERIAL
% =============================================================================

% =============================================================================
% BEGIN APPENDED VERSION V138 MATERIAL
% =============================================================================

\section{V138: the Standard Model chiral anomaly polynomial, the
traditional \(SU(2)\) parity test, and the determinant-line remainder}
\label{sec:V138}

\subsection{Purpose}
\label{subsec:V138-purpose}

V134 separated the local doubled modulus from the phase of the physical
chiral determinant.  V135--V137 used only the modulus branch.  V138 now
computes the representation-theoretic obstruction to a gauge-equivariant
phase.  All fermions are first written as left-handed Weyl fields.  This
fixes every conjugation sign in the cubic and mixed anomaly traces.

The finite sums prove cancellation of the perturbative gauge and mixed
gauge--gravitational anomaly polynomial for each generation.  The number of
\(SU(2)\) doublets is also even, so the traditional mod-two anomaly is
absent.  These conclusions flatten the local anomaly connection and remove
one specified global holonomy.  They do not canonically trivialize the
determinant line on a manifold with boundary.

\subsection{All-left-handed representation ledger}
\label{subsec:V138-representations}

Use the hypercharge convention \(Q_{\rm em}=T_3+Y\).  A right-handed
physical fermion is represented below by its left-handed charge conjugate.
For one generation, including a sterile right-handed neutrino, the chiral
representation is
\begin{equation}
 \begin{aligned}
 Q&=(\mathbf3,\mathbf2)_{1/6},
 &u^c&=(\overline{\mathbf3},\mathbf1)_{-2/3},
 &d^c&=(\overline{\mathbf3},\mathbf1)_{1/3},\\
 L&=(\mathbf1,\mathbf2)_{-1/2},
 &e^c&=(\mathbf1,\mathbf1)_{1},
 &\nu^c&=(\mathbf1,\mathbf1)_{0}.
 \end{aligned}
 \label{eq:V138-left-handed-ledger}
\end{equation}

Normalize the quadratic indices by
\begin{equation}
 T_{\mathbf3}^{SU(3)}=T_{\mathbf2}^{SU(2)}=\frac12
 \label{eq:V138-quadratic-index}
\end{equation}
and the cubic \(SU(3)\) index by
\begin{equation}
 A(\mathbf3)=1,
 \qquad A(\overline{\mathbf3})=-1.
 \label{eq:V138-cubic-index}
\end{equation}

\begin{lemma}[Degree-six anomaly density]
\label{lem:V138-anomaly-polynomial}
For a left-handed four-dimensional Weyl family in a complex
representation \(\mathcal R\), the local anomaly polynomial is the
degree-six component
\begin{equation}
 I_6(\mathcal R)
 =\left[\widehat A(TM)\,
 \operatorname{ch}_{\mathcal R}(F)\right]_6
 =\frac16\operatorname{tr}_{\mathcal R}
 \left(\frac{iF}{2\pi}\right)^3
 -\frac{p_1(TM)}{24}\operatorname{tr}_{\mathcal R}
 \left(\frac{iF}{2\pi}\right).
 \label{eq:V138-anomaly-polynomial}
\end{equation}
There is no pure gravitational six-form term for a spin-one-half Weyl
field in four dimensions.
\end{lemma}

\begin{proof}
Expand
\(\widehat A=1-p_1/24+O(8)\) and
\(\operatorname{ch}(F)=
 \operatorname{tr}\exp(iF/2\pi)\).  Total degree six can arise only from
the cubic Chern-character term or from the product of the degree-four
Pontryagin form with the degree-two Chern-character term.  This gives the
display and leaves no degree-six polynomial made solely from tangent-bundle
Pontryagin forms.
\end{proof}

\begin{remark}[Source check]
\label{rem:V138-source-check}
The determinant-line curvature interpretation used here is the local family
index theorem of Bismut--Freed.  The boundary phase warning below is the
eta-invariant and determinant-line mechanism of Dai--Freed.  The mod-two
test is the original \(SU(2)\) anomaly of Witten.  V138 uses these inputs
only after reducing the Standard Model part to the explicit finite traces
below.
\end{remark}

\subsection{Perturbative nonabelian and mixed anomalies}
\label{subsec:V138-nonabelian}

\begin{theorem}[Pure \(SU(3)^3\) cancellation]
\label{thm:V138-SU3-cubic}
For one Standard Model generation,
\begin{equation}
 \mathcal A_{333}
 =2A(\mathbf3)+A(\overline{\mathbf3})
  +A(\overline{\mathbf3})
 =2-1-1=0.
 \label{eq:V138-SU3-cubic}
\end{equation}
\end{theorem}

\begin{proof}
The quark doublet contains two left-handed color triplets.  The conjugated
up and down singlets are one color antitriplet each.  The leptons are color
singlets.  Insert the normalization \eqref{eq:V138-cubic-index}.
\end{proof}

\begin{proposition}[No perturbative \(SU(2)^3\) anomaly]
\label{prop:V138-SU2-cubic}
The local cubic \(SU(2)\) anomaly vanishes representation by
representation.
\end{proposition}

\begin{proof}
The symmetric invariant
\(d^{abc}=\operatorname{tr}(T^{(a}T^bT^{c)})\) vanishes for
\(\mathfrak{su}(2)\).  Equivalently, every finite-dimensional
\(SU(2)\) representation is real or pseudoreal and has zero cubic index.
This is a perturbative statement; the mod-two obstruction is treated
separately below.
\end{proof}

\begin{theorem}[Mixed nonabelian--hypercharge cancellation]
\label{thm:V138-mixed-nonabelian}
The two possibly nonzero coefficients are
\begin{align}
 \mathcal A_{33Y}
 &=2\left(\frac16\right)\left(\frac12\right)
 +\left(-\frac23\right)\left(\frac12\right)
 +\left(\frac13\right)\left(\frac12\right)
 =0,
 \label{eq:V138-SU3SU3Y}\\
 \mathcal A_{22Y}
 &=3\left(\frac16\right)\left(\frac12\right)
 +\left(-\frac12\right)\left(\frac12\right)
 =0.
 \label{eq:V138-SU2SU2Y}
\end{align}
Every coefficient with one nonabelian generator and otherwise abelian or
gravitational insertions vanishes separately because the nonabelian
generator is traceless.
\end{theorem}

\begin{proof}
For \(SU(3)^2Y\), multiply each hypercharge and color quadratic index by
the dimension of its weak representation.  This gives
\(1/6-1/3+1/6=0\).  For \(SU(2)^2Y\), multiply by color multiplicity,
giving \(1/4-1/4=0\).  A trace containing exactly one generator of a
semisimple factor factorizes through \(\operatorname{tr}T^a=0\), proving
the final assertion.
\end{proof}

\subsection{Cubic hypercharge and mixed gravitational traces}
\label{subsec:V138-abelian}

\begin{theorem}[Cubic hypercharge cancellation]
\label{thm:V138-hypercharge-cubic}
For one generation,
\begin{align}
 \mathcal A_{YYY}
 &=6\left(\frac16\right)^3
 +3\left(-\frac23\right)^3
 +3\left(\frac13\right)^3
 +2\left(-\frac12\right)^3
 +1^3+0^3
 \notag\\
 &=\frac1{36}-\frac89+\frac19-\frac14+1
 =0.
 \label{eq:V138-hypercharge-cubic}
\end{align}
\end{theorem}

\begin{proof}
The multiplicities in order are the dimensions of the color--weak fibres:
\(6,3,3,2,1,1\).  On the common denominator \(36\), the numerator is
\(1-32+4-9+36=0\).
\end{proof}

\begin{theorem}[Mixed gravitational--hypercharge cancellation]
\label{thm:V138-gravity-hypercharge}
The linear hypercharge trace is
\begin{equation}
 \mathcal A_{\mathrm{grav}\,Y}
 =6\left(\frac16\right)
 +3\left(-\frac23\right)
 +3\left(\frac13\right)
 +2\left(-\frac12\right)+1+0
 =0.
 \label{eq:V138-gravity-hypercharge}
\end{equation}
Hence the \(p_1(TM)F_Y\) term in
\eqref{eq:V138-anomaly-polynomial} cancels.
\end{theorem}

\begin{proof}
The displayed terms are \(1-2+1-1+1+0\).  Their sum is zero.
\end{proof}

\begin{corollary}[Generation-wise local anomaly cancellation]
\label{cor:V138-local-cancellation}
For every positive number \(N_g\) of identical generations, with or
without sterile right-handed neutrinos,
\begin{equation}
 \boxed{I_6^{\rm SM}=0.}
 \label{eq:V138-I6-zero}
\end{equation}
The sterile neutrino changes the modulus count in V135 but changes none of
the anomaly sums in V138.
\end{corollary}

\begin{proof}
Theorems~\ref{thm:V138-SU3-cubic},
\ref{thm:V138-mixed-nonabelian},
\ref{thm:V138-hypercharge-cubic}, and
\ref{thm:V138-gravity-hypercharge}, together with
Proposition~\ref{prop:V138-SU2-cubic}, exhaust the degree-six terms in the
product-group decomposition of \eqref{eq:V138-anomaly-polynomial}.
Each vanishes for one generation and hence for \(N_g\) copies.  The field
\(\nu^c\) is a gauge singlet of hypercharge zero.
\end{proof}

\begin{corollary}[Yukawa independence of the local anomaly]
\label{cor:V138-Yukawa-independence}
Higgs vacuum expectation values, Yukawa matrices, and fermion masses do not
alter \eqref{eq:V138-I6-zero} as long as the chiral representation and
gauge symmetry are unchanged.
\end{corollary}

\begin{proof}
The polynomial \eqref{eq:V138-anomaly-polynomial} depends on the principal
chiral symbol and the twisting representation.  Yukawa and mass operators
are zeroth-order deformations and do not change those representation traces.
\end{proof}

\subsection{The traditional \(SU(2)\) global anomaly}
\label{subsec:V138-Witten}

\begin{theorem}[Even-doublet test]
\label{thm:V138-Witten-parity}
For the simply connected \(SU(2)\) gauge factor on a spin four-manifold,
the traditional Witten mod-two anomaly vanishes for every number of
Standard Model generations.  The number of left-handed fundamental
doublets is
\begin{equation}
 N_{\mathbf2}=N_g(3+1)=4N_g\equiv0\pmod2.
 \label{eq:V138-doublet-parity}
\end{equation}
For three generations it is twelve.
\end{theorem}

\begin{proof}
The quark doublet occurs once for each of three colors and the lepton
doublet occurs once, giving four weak doublets per generation.  All
right-handed conjugate fields in \eqref{eq:V138-left-handed-ledger} are
weak singlets.  The mod-two index associated with the nontrivial element of
\(\pi_4(SU(2))\cong\mathbb Z_2\) is additive over doublets, so an even number
has zero parity.
\end{proof}

\begin{firewall}[Scope of the parity theorem]
\label{fire:V138-Witten-scope}
Theorem~\ref{thm:V138-Witten-parity} removes the traditional anomaly for
the fundamental doublets of the simply connected \(SU(2)\) factor on spin
backgrounds.  It is not a computation of
\(\Omega_5^{\mathrm{Spin}}(BG_{\rm SM})\) for every quotient
\[
 G_{\rm SM}=(SU(3)\times SU(2)\times U(1))/\Gamma,
\]
nor of possible spin--gauge structures on non-spin manifolds.  Those global
forms require their own bordism and eta-holonomy calculation.
\end{firewall}

\subsection{What has and has not happened to the determinant line}
\label{subsec:V138-determinant-line}

Let \(\mathcal L_{\rm ch}\) denote the tensor product of the physical Weyl
determinant lines over a smooth parameter chart of metrics and gauge
connections with a fixed admissible boundary setup.

\begin{theorem}[Local flatness]
\label{thm:V138-local-flatness}
On every chart on which the family of chiral Fredholm operators and its
boundary data are defined, the gauge and mixed gauge--gravitational
curvature of the Bismut--Freed connection on \(\mathcal L_{\rm ch}\)
vanishes for the Standard Model representation:
\begin{equation}
 \boxed{\Omega^{\mathcal L_{\rm ch}}_{\rm local}=0.}
 \label{eq:V138-local-flatness}
\end{equation}
Thus the perturbative anomaly supplies no local obstruction to a parallel
phase choice.
\end{theorem}

\begin{proof}
The family index theorem identifies the local determinant-line curvature,
up to the universal descent normalization, with the fibre integral of the
degree-six density \eqref{eq:V138-anomaly-polynomial}.  Corollary
\ref{cor:V138-local-cancellation} makes that density zero after summing the
physical chiral representations.
\end{proof}

\begin{proposition}[Flat does not mean canonically trivial]
\label{prop:V138-flat-not-trivial}
Equation~\eqref{eq:V138-local-flatness} does not by itself give a global
nonvanishing parallel section of \(\mathcal L_{\rm ch}\).  A flat line can
retain nontrivial holonomy around a loop in parameter space.  On a manifold
with boundary, a phase construction additionally requires compatible
boundary determinant-line or eta data and a gluing law.
\end{proposition}

\begin{proof}
A flat unitary line bundle is classified by its holonomy representation of
the fundamental group; zero curvature kills infinitesimal holonomy but not
the representation.  For Dirac families with boundary, the exponentiated
eta invariant naturally takes values in the inverse determinant line of the
boundary family.  A scalar phase results only after choosing and transporting
the required trivialization.  Neither the local bag modulus of V135 nor the
finite sums above provide that choice.
\end{proof}

\begin{corollary}[Exact status of the fermion phase card]
\label{cor:V138-phase-status}
The V134 fermion fork is narrowed as follows:
\begin{enumerate}[label=\textup{(\roman*)},leftmargin=2.8em]
\item the local perturbative curvature obstruction vanishes;
\item the traditional \(SU(2)\) mod-two holonomy vanishes;
\item quotient-group bordism holonomies and the boundary eta
      trivialization remain to be computed.
\end{enumerate}
Therefore the phase may not yet be set equal to one.
\end{corollary}

\subsection{V138 ledger and next acquisition}
\label{subsec:V138-ledger}

\paragraph{New exact conclusions.}
\begin{enumerate}[label=\textup{(\arabic*)},leftmargin=2.8em]
\item Equation~\eqref{eq:V138-left-handed-ledger} fixes the conjugation
      convention for every Standard Model Weyl multiplet.
\item Theorems~\ref{thm:V138-SU3-cubic}--
      \ref{thm:V138-gravity-hypercharge} evaluate every perturbative anomaly
      trace as an explicit rational sum.
\item Corollary~\ref{cor:V138-local-cancellation} proves generation-wise
      vanishing of the complete degree-six anomaly polynomial.
\item Theorem~\ref{thm:V138-Witten-parity} proves that the twelve-doublet
      three-generation model passes the traditional mod-two test.
\item Theorem~\ref{thm:V138-local-flatness} converts the polynomial
      cancellation into local flatness of the chiral determinant line.
\item Proposition~\ref{prop:V138-flat-not-trivial} identifies the precise
      remaining holonomy and boundary-trivialization gap.
\end{enumerate}

\paragraph{Next forced step.}
V139 should construct the boundary phase card rather than infer it from
local flatness.  It should define a self-adjoint odd-dimensional boundary
family, state the spectral-section or invertibility hypothesis, build the
exponentiated eta element in the inverse determinant line, and prove its
finite-cutoff gluing and gauge-transport laws.  The actual Standard Model
quotient \(\Gamma\) must remain a declared input until its five-dimensional
spin-bordism holonomies are evaluated.

\begin{assessment}[V138 continuum status]
\label{ass:V138-continuum-status}
The full Standard--Model--Einstein continuum remains unproved.  V138 closes
the perturbative anomaly polynomial and the traditional \(SU(2)\) mod-two
test.  Quotient-group global holonomy, boundary eta transport, extended
BRST closure, microscopic gravitational normalization, nonlinear existence,
Einstein stress closure, and the cofinal interacting limit remain open.
\end{assessment}

\section*{V138 exact checks and append-only record}
\addcontentsline{toc}{section}{V138 exact checks and append-only record}

The anomaly checksum is
\begin{equation}
 \boxed{
 \mathcal A_{333}=\mathcal A_{33Y}=\mathcal A_{22Y}
 =\mathcal A_{YYY}=\mathcal A_{\mathrm{grav}\,Y}=0,
 \qquad
 N_{\mathbf2}=4N_g\equiv0\pmod2.}
 \label{eq:V138-final-check}
\end{equation}
The append operation preserves the complete V137 prefix byte for byte before
its former live terminator.  The assembled V138 file is checked for one live
final terminator, balanced environments and braces, duplicate labels, newly
introduced unresolved references, display math delimiters, and control
characters before the exact V137 predecessor is removed.  The companion
files are read only.

% =============================================================================
% END APPENDED VERSION V138 MATERIAL
% =============================================================================

% =============================================================================
% BEGIN APPENDED VERSION V139 MATERIAL
% =============================================================================

\section{V139: the finite-cutoff chiral phase line, eta transport, and
the cofinal trivialization criterion}
\label{sec:V139}

\subsection{Purpose}
\label{subsec:V139-purpose}

V138 proved that the Standard Model anomaly polynomial vanishes and that
the traditional \(SU(2)\) mod-two loop has trivial parity.  Those results
make the chiral determinant line locally flat; they do not choose a scalar
phase.  V139 constructs the missing finite-cutoff object.

The construction has three levels.  First, an elliptic spectral boundary
condition makes the chiral Dirac family Fredholm and produces its
determinant line.  Second, Bismut--Freed transport expresses its holonomy by
a five-dimensional reduced eta invariant.  Third, because the four-manifold
has boundary, the Dai--Freed eta object is itself line-valued on a boundary
mapping torus.  A scalar phase exists only after that boundary line is
trivialized compatibly with gluing.  The final proposition states the exact
inverse-system condition required for a cofinal phase.

\subsection{Boundary family and spectral sections}
\label{subsec:V139-spectral-section}

Fix one finite renewal cutoff and one connected smooth parameter chart
\(\mathcal B\) of metrics, gauge connections, and Higgs--Yukawa data.  Let
\begin{equation}
 D_b^+:H^1(M;S^+\otimes E_b)longrightarrow
 L^2(M;S^-\otimes E_b),
 \qquad b\in\mathcal B,
 \label{eq:V139-chiral-family}
\end{equation}
be the physical all-left-handed Standard Model Dirac family.  In a product
collar, its doubled self-adjoint operator has normal form
\begin{equation}
 \mathbb D_b=c(n)(D_n+A_b+Z_b),
 \label{eq:V139-collar-Dirac}
\end{equation}
where \(A_b=A_b^*\) is the tangential principal Dirac family on the closed
three-manifold \(\Sigma\) and \(Z_b\) is zeroth order.

\begin{definition}[Spectral-section datum]
\label{def:V139-spectral-section}
A spectral-section datum on \(\mathcal B\) is a smooth family of orthogonal
pseudodifferential projections \(P_b\) with the same principal symbol as
the nonnegative spectral projection of \(A_b\), such that there is a
constant \(R>0\) for which
\begin{equation}
 A_bu>R u\Longrightarrow P_bu=u,
 \qquad
 A_bu<-R u\Longrightarrow P_bu=0
 \label{eq:V139-spectral-section}
\end{equation}
in the spectral resolution, uniformly on compact subsets of
\(\mathcal B\).  The associated APS-type domain is denoted
\(\operatorname{Dom}D_{b,P}^+\).
\end{definition}

\begin{proposition}[Uniform gap gives a canonical local datum]
\label{prop:V139-gap-APS}
If
\begin{equation}
 \inf_{b\in\mathcal B}\operatorname{dist}
 (0,\operatorname{spec}A_b)\ge\delta>0,
 \label{eq:V139-boundary-gap}
\end{equation}
then
\begin{equation}
 P_b^{\rm APS}=1_{[0,\infty)}(A_b)
 \label{eq:V139-APS-projector}
\end{equation}
is a smooth family of order-zero pseudodifferential projections and defines
a Fredholm family \(D_{b,P^{\rm APS}}^+\).
\end{proposition}

\begin{proof}
Choose a contour \(\Gamma\) separating the positive and negative spectra
at distance less than \(\delta\) from the imaginary axis.  Functional
calculus gives
\[
 P_b^{\rm APS}=\frac1{2\pi i}\int_\Gamma(z-A_b)^{-1}\,dz.
\]
The uniform gap makes the resolvent and all parameter derivatives bounded
on the contour, so the projection is smooth.  Its principal symbol is the
positive spectral projector of the tangential Dirac symbol.  The standard
Calderon comparison then shows that the APS boundary map is Fredholm; hence
the chiral family with this domain is Fredholm.
\end{proof}

\begin{firewall}[Gap closure is carried, not ignored]
\label{fire:V139-gap-closure}
If an eigenvalue of \(A_b\) crosses zero, the sharp projector
\eqref{eq:V139-APS-projector} is not a smooth constant-rank family.  One
must choose a spectral section and retain the finite-dimensional crossing
space.  Deleting that space would delete its spectral flow and can change
the determinant-line holonomy.
\end{firewall}

\subsection{The determinant line}
\label{subsec:V139-determinant-line}

For a Fredholm operator \(T\), use the convention
\begin{equation}
 \operatorname{Det}(T)
 :=\Lambda^{\rm top}\ker T\otimes
 \left(\Lambda^{\rm top}\operatorname{coker}T\right)^*.
 \label{eq:V139-Det-convention}
\end{equation}

\begin{theorem}[Finite-cutoff chiral determinant line]
\label{thm:V139-determinant-line}
Given a smooth spectral section \(P\), the fibres
\begin{equation}
 \mathcal L_{P,b}:=\operatorname{Det}(D_{b,P}^+)
 \label{eq:V139-determinant-line}
\end{equation}
form a smooth Hermitian line bundle \(\mathcal L_P\to\mathcal B\) with its
Quillen metric and Bismut--Freed connection.  The regularized chiral
determinant is naturally a section of this line rather than a complex-valued
function across kernel crossings.
\end{theorem}

\begin{proof}
On a neighbourhood where a finite-dimensional subspace contains every
small eigenmode, split the family into that finite block and an invertible
complement.  The complement has a smooth zeta determinant, while the finite
block contributes exactly the exterior-power line in
\eqref{eq:V139-Det-convention}.  On overlaps, ordinary determinant
multiplicativity gives the transition maps and their cocycle.  The heat
regularization of the invertible complement defines the Quillen norm and
connection.  These local constructions agree on overlaps, producing the
claimed line bundle.
\end{proof}

\begin{proposition}[Change of spectral section is finite dimensional]
\label{prop:V139-section-change}
Let \(P_0\) and \(P_1\) be two spectral sections.  Their difference is a
smooth family of finite-rank operators modulo a common high-energy
projection.  There is a finite-dimensional relative line
\(\mathcal R(P_1,P_0)\) and a canonical isomorphism
\begin{equation}
 \mathcal L_{P_1}
 \cong\mathcal L_{P_0}\otimes\mathcal R(P_1,P_0),
 \label{eq:V139-relative-line}
\end{equation}
with the cocycle law
\begin{equation}
 \mathcal R(P_2,P_1)\otimes\mathcal R(P_1,P_0)
 \cong\mathcal R(P_2,P_0).
 \label{eq:V139-relative-cocycle}
\end{equation}
\end{proposition}

\begin{proof}
Choose a common spectral cutoff beyond the transition bands of all three
sections.  Above and below it the projections agree.  The change of domain
is therefore supported on the finite sum of eigenspaces in the intervening
band.  Define \(\mathcal R(P_1,P_0)\) to be the determinant line of the
resulting finite-dimensional two-term relative complex.  Exact-sequence
multiplicativity of determinant lines gives
\eqref{eq:V139-relative-line}.  Splicing the relative complexes for
\((P_2,P_1)\) and \((P_1,P_0)\) cancels the common middle block and proves
\eqref{eq:V139-relative-cocycle}.
\end{proof}

\begin{corollary}[The V135 bag does not choose this line]
\label{cor:V139-bag-not-phase}
The local doubled bag domain of V135 determines a positive modulus for
\(\mathbb D_b^2\), but it does not specify a spectral section for the
physical chiral operator \(D_b^+\) or a trivialization of
\(\mathcal L_P\).
\end{corollary}

\begin{proof}
The bag projector acts on the Hermitian double and mixes the two spin
chiralities.  Its determinant identity in
Lemma~\ref{lem:V134-determinant-multiplicity} retains only
\(|\det D_b^+|\).  A spectral section instead polarizes the tangential
chiral boundary family and carries the finite crossing spaces in
Proposition~\ref{prop:V139-section-change}.  Neither datum determines the
other.
\end{proof}

\subsection{Closed-loop eta transport}
\label{subsec:V139-eta-transport}

For a self-adjoint Dirac operator \(B\) on a closed odd-dimensional
manifold, define
\begin{equation}
 \xi(B):=\frac{\eta(B)+\dim\ker B}{2},
 \qquad
 \tau(B):=\exp\bigl(-2\pi i\xi(B)\bigr).
 \label{eq:V139-reduced-eta}
\end{equation}
This fixes the chirality convention used below; reversing chirality
complex-conjugates every phase.

Let \(\gamma:S^1\to\mathcal B\) be a smooth loop.  Its mapping torus is a
five-manifold \(Y_\gamma\) carrying the Dirac operator assembled from the
four-dimensional family.  Let \(D_{Y_{\gamma,\epsilon}}\) denote the
operator for the metric in which the circle direction is rescaled by
\(\epsilon^{-2}\).

\begin{theorem}[Bismut--Freed holonomy formula]
\label{thm:V139-BF-holonomy}
If the four-manifold is closed, the holonomy of the determinant-line
connection around \(\gamma\) is
\begin{equation}
 \operatorname{Hol}_\gamma(\mathcal L)
 =\lim_{\epsilon\downarrow0}
 \tau(D_{Y_{\gamma,\epsilon}}).
 \label{eq:V139-BF-holonomy}
\end{equation}
The right side is invariant under subdivision and composition of loops.
\end{theorem}

\begin{proof}
This is the Bismut--Freed holonomy theorem for a family of even-dimensional
Dirac operators.  Its hypotheses are met by a smooth Fredholm family with
the Quillen connection.  The adiabatic limit identifies parallel transport
of the small-eigenvalue determinant with the reduced eta invariant of the
mapping torus, while the heat-regularized complement gives the local
connection integral.  Under gluing of mapping cylinders, the intermediate
determinant factors pair with their duals; exponentiation removes the
integer index ambiguity.  Hence the resulting scalar is multiplicative.
\end{proof}

\begin{corollary}[Effect of V138]
\label{cor:V139-flat-holonomy-character}
For the Standard Model representation, V138 makes the determinant-line
connection flat.  Its remaining closed-manifold anomaly is therefore a
unitary character
\begin{equation}
 \chi_{\rm anom}:\pi_1(\mathcal B)\longrightarrow U(1),
 \qquad
 \chi_{\rm anom}([\gamma])
 =\operatorname{Hol}_\gamma(\mathcal L).
 \label{eq:V139-anomaly-character}
\end{equation}
Theorem~\ref{thm:V138-Witten-parity} proves that this character is one on
the traditional \(SU(2)\) generator represented by fundamental doublets.
\end{corollary}

\begin{proof}
Flat parallel transport depends only on the homotopy class of a loop and is
multiplicative under concatenation, so it is a character.  The specified
\(SU(2)\) value is the mod-two eta holonomy computed in V138.
\end{proof}

\subsection{Why the renewal boundary makes the eta object line-valued}
\label{subsec:V139-boundary-eta}

For the programme, \(M\) has boundary \(\Sigma\).  Therefore
\begin{equation}
 \partial Y_\gamma=W_\gamma,
 \qquad
 W_\gamma:=\Sigma\mathbin{\rtimes_\gamma}S^1,
 \label{eq:V139-boundary-mapping-torus}
\end{equation}
and \(W_\gamma\) is four-dimensional.  Let
\(D_{W_\gamma}^+\) be its induced chiral Dirac operator.

\begin{theorem}[Dai--Freed line-valued eta element]
\label{thm:V139-DF-element}
For a five-dimensional Dirac manifold with boundary as in
\eqref{eq:V139-boundary-mapping-torus}, the exponentiated reduced eta
invariant has the canonical type
\begin{equation}
 \tau(Y_\gamma)
 \in\operatorname{Det}(D_{W_\gamma}^+)^{-1}.
 \label{eq:V139-DF-element}
\end{equation}
If two such bordisms are glued along an oppositely oriented common boundary,
their eta elements pair through the determinant line and its dual, yielding
the eta element of the glued bordism.
\end{theorem}

\begin{proof}
This is the Dai--Freed eta-invariant theorem.  With boundary, changing the
spectral boundary trivialization changes the numerical reduced eta invariant
by precisely the determinant of the finite boundary transition.  Treating
its exponential as an element of the inverse determinant line absorbs that
change and makes the object canonical.  In a gluing, the outgoing boundary
line of the first piece is dual to the incoming boundary line of the second;
their canonical pairing cancels the two transition factors.  The analytic
eta gluing formula then gives the stated multiplicativity.
\end{proof}

\begin{corollary}[A scalar renewal phase needs boundary data]
\label{cor:V139-boundary-trivialization}
A scalar holonomy around \(\gamma\) is obtained from
\eqref{eq:V139-DF-element} only after choosing a nonzero boundary vector
\begin{equation}
 \ell_\gamma\in\operatorname{Det}(D_{W_\gamma}^+)
 \label{eq:V139-boundary-vector}
\end{equation}
and pairing:
\begin{equation}
 \operatorname{Hol}^{\ell}_\gamma
 :=\langle\tau(Y_\gamma),\ell_\gamma\rangle\in\mathbb C^\times.
 \label{eq:V139-paired-holonomy}
\end{equation}
The vectors \(\ell_\gamma\) must respect disjoint union, orientation
reversal, and gluing.  An arbitrary choice loop by loop does not define a
determinant-line trivialization.
\end{corollary}

\begin{proof}
The first statement is the canonical pairing of a one-dimensional vector
space with its dual.  Parallel transport must multiply under loop
concatenation.  Theorem~\ref{thm:V139-DF-element} supplies this law for the
eta elements, so the paired scalars obey it exactly when the chosen boundary
vectors obey the dual gluing law.
\end{proof}

\begin{firewall}[Local anomaly cancellation is insufficient at a boundary]
\label{fire:V139-local-insufficient}
Equation~\eqref{eq:V138-I6-zero} removes the local curvature of the line.  It
does not create \(\ell_\gamma\), prove that
\(\operatorname{Det}(D_{W_\gamma}^+)\) is canonically trivial, or evaluate
its holonomy for the chosen quotient form of \(G_{\rm SM}\).  Replacing the
line-valued eta element by the number one would assume the remaining theorem.
\end{firewall}

\subsection{Gauge invariance as a holonomy criterion}
\label{subsec:V139-gauge-invariance}

\begin{theorem}[Flat-line trivialization criterion]
\label{thm:V139-flat-trivialization}
Let \(\mathcal B_0\) be a connected component and suppose the boundary eta
lines have been trivialized compatibly with gluing.  The flat chiral
determinant line admits a nowhere-zero parallel section on \(\mathcal B_0\)
if and only if
\begin{equation}
 \operatorname{Hol}^{\ell}_\gamma=1
 \quad\text{for every }[\gamma]\in\pi_1(\mathcal B_0).
 \label{eq:V139-trivial-holonomy}
\end{equation}
When the loops include large gauge transformations, such a section is the
required gauge-equivariant scalar phase.
\end{theorem}

\begin{proof}
If a nonzero parallel section exists, transport around a loop returns the
same vector, so every holonomy is one.  Conversely, fix a nonzero vector in
one fibre and transport it to any other point.  Two paths differ by a loop;
\eqref{eq:V139-trivial-holonomy} makes the result path independent.  The
transported vectors form a nonzero parallel section.  Applying the same
argument to paths ending at gauge-equivalent points gives gauge
equivariance.
\end{proof}

\begin{corollary}[Exact finite-cutoff phase gate]
\label{cor:V139-finite-phase-gate}
At one renewal cutoff, the physical chiral phase is a scalar precisely after
all of the following are supplied:
\begin{enumerate}[label=\textup{(\roman*)},leftmargin=2.8em]
\item a smooth spectral section, or a uniform boundary spectral gap;
\item compatible trivializations of the boundary mapping-torus determinant
      lines;
\item unit eta holonomy for every generator of the actual metric--gauge
      parameter-space fundamental group.
\end{enumerate}
V138 proves the infinitesimal condition and one traditional global
generator, but not the complete list in \textup{(iii)}.
\end{corollary}

\subsection{The cofinal phase system}
\label{subsec:V139-cofinal}

Let \(q\preceq q'\) be two renewal cutoffs.  Denote their chiral determinant
lines by \(\mathcal L_q\) and suppose the renewal comparison supplies unitary
line isomorphisms
\begin{equation}
 J_{q'q}:\mathcal L_q\longrightarrow\mathcal L_{q'}.
 \label{eq:V139-line-comparison}
\end{equation}

\begin{theorem}[Cofinal scalar-phase criterion]
\label{thm:V139-cofinal-phase}
A cutoff-independent nonzero parallel phase section
\(s=(s_q)_q\) with
\begin{equation}
 J_{q'q}s_q=s_{q'}
 \label{eq:V139-compatible-sections}
\end{equation}
exists if the following conditions hold:
\begin{enumerate}[label=\textup{(\roman*)},leftmargin=2.8em]
\item every \(\mathcal L_q\) has trivial paired eta holonomy;
\item the comparison maps are parallel and satisfy
\begin{equation}
 J_{q''q'}J_{q'q}=J_{q''q}
 \quad(q\preceq q'\preceq q'');
 \label{eq:V139-comparison-cocycle}
\end{equation}
\item the spectral-section relative lines and boundary eta trivializations
      used to construct \(J_{q'q}\) satisfy their determinant-line gluing
      cocycles.
\end{enumerate}
Conversely, any compatible nonzero parallel family
\eqref{eq:V139-compatible-sections} forces \textup{(i)} and the cocycle on
its one-dimensional span.
\end{theorem}

\begin{proof}
Choose a base cutoff \(q_0\) and a nonzero vector in one fibre.  Condition
\textup{(i)} and Theorem~\ref{thm:V139-flat-trivialization} extend it to a
parallel section \(s_{q_0}\).  For any later cutoff, set
\(s_q=J_{qq_0}s_{q_0}\).  Conditions \textup{(ii)}--\textup{(iii)} make
this definition independent of every intermediate cutoff and every spectral
section or gluing decomposition.  This proves existence and
\eqref{eq:V139-compatible-sections}.

Conversely, parallel transport of a compatible nonzero section around a
loop returns the same section, so the holonomy is one.  Applying two
successive comparison maps and the direct map to \(s_q\) gives the same
nonzero vector; hence the maps agree on its line, which is the claimed
cocycle necessity.
\end{proof}

\begin{firewall}[No cofinal phase from pointwise choices]
\label{fire:V139-no-pointwise-phase}
Choosing a unit vector in each one-dimensional \(\mathcal L_q\) separately
does not prove Theorem~\ref{thm:V139-cofinal-phase}.  The uncancelled ratios
are phases on comparison triangles and mapping-torus loops.  They are
exactly the cocycle and holonomy data which the continuum construction must
control.
\end{firewall}

\subsection{V139 ledger and next acquisition}
\label{subsec:V139-ledger}

\paragraph{New exact conclusions.}
\begin{enumerate}[label=\textup{(\arabic*)},leftmargin=2.8em]
\item Proposition~\ref{prop:V139-gap-APS} gives the canonical APS family on
      a uniformly gapped boundary branch.
\item Theorem~\ref{thm:V139-determinant-line} constructs the finite-cutoff
      chiral determinant line across kernel changes.
\item Proposition~\ref{prop:V139-section-change} retains spectral flow in a
      finite relative determinant line with an exact cocycle.
\item Theorem~\ref{thm:V139-BF-holonomy} identifies closed-loop phase
      transport with the five-dimensional eta invariant.
\item Theorem~\ref{thm:V139-DF-element} proves that the renewal boundary
      changes that eta invariant into an inverse-determinant-line element.
\item Theorem~\ref{thm:V139-flat-trivialization} states the exact
      finite-cutoff scalar phase criterion.
\item Theorem~\ref{thm:V139-cofinal-phase} lifts the criterion to the
      renewal inverse system and exposes every required cocycle.
\end{enumerate}

\paragraph{Next forced step.}
V140 is the next compulsory companion audit.  It should inspect the current
NS and Yang--Mills notes and then test the V139 phase gate against the
actual renewal comparison maps.  The immediate mathematical acquisition is
to declare the global form of \(G_{\rm SM}\), identify the corresponding
five-dimensional spin-bordism generators, and evaluate their Standard Model
eta characters.  If the source programme has not selected \(\Gamma\), V140
must branch by quotient rather than silently use the direct product.

\begin{assessment}[V139 continuum status]
\label{ass:V139-continuum-status}
The full Standard--Model--Einstein continuum remains unproved.  V139
constructs the conditional finite-cutoff and cofinal phase lines and gives
necessary and sufficient trivial-holonomy criteria.  The quotient-group eta
characters, compatible boundary trivializations, extended BRST complex,
microscopic gravitational normalization, nonlinear existence, Einstein
stress closure, and cofinal interacting limit remain open.
\end{assessment}

\section*{V139 exact checks and append-only record}
\addcontentsline{toc}{section}{V139 exact checks and append-only record}

The phase-line checksum is
\begin{equation}
 \boxed{
 \tau(Y_\gamma)\in\operatorname{Det}(D_{W_\gamma}^+)^{-1},
 \quad
 \operatorname{Hol}^{\ell}_\gamma
 =\langle\tau(Y_\gamma),\ell_\gamma\rangle,
 \quad
 J_{q''q'}J_{q'q}=J_{q''q}.}
 \label{eq:V139-final-check}
\end{equation}
The append operation preserves the complete V138 prefix byte for byte before
its former live terminator.  The assembled V139 file is checked for one live
final terminator, balanced environments and braces, duplicate labels, newly
introduced unresolved references, display math delimiters, and control
characters before the exact V138 predecessor is removed.  No companion file
is modified.

% =============================================================================
% END APPENDED VERSION V139 MATERIAL
% =============================================================================

% =============================================================================
% BEGIN APPENDED VERSION V140 MATERIAL
% =============================================================================

\section{V140: compulsory companion audit, the exact \(\mathbb Z_6\)
matter kernel, and direct cofinal phase defects}
\label{sec:V140}

\subsection{Purpose}
\label{subsec:V140-purpose}

V139 reduced the chiral phase to determinant-line holonomy and compatible
cutoff comparison maps.  V140 first performs the compulsory fifth-version
audit.  It then resolves one input which V139 deliberately left open: the
common central kernel of the complete Standard Model matter and Higgs
representation.  The maximal faithful global form is the \(\mathbb Z_6\)
quotient, while quotients by its subgroups are distinct non-effective global
branches.

The audit also sharpens the cofinal phase gate.  The relevant objects are
direct transition ratios and their gauge-invariant triangle defects, not
separate convergence estimates for arbitrary local phase representatives.
Along a cofinal chain, absolute summability of the exact adjacent phase
increments is sufficient.  Square summability is not; an explicit unitary
counterexample proves the missing first-moment requirement.

\subsection{Mandatory V140 companion audit}
\label{subsec:V140-audit}

The stable read-only shared-folder snapshot at
\(09{:}31{:}10\) Pacific time on 5 September 2026 was
\begin{center}
\begin{tabular}{|p{0.49\textwidth}|r|p{0.31\textwidth}|}
\hline
Companion file & Bytes & SHA--256\\
\hline
\texttt{Renewal\_NS\_Stability\_Research\_Notes\_V31.tex}
& \(887537\) &
\texttt{F1121B62238AD5491BB7CF03635EA9934942EDF573ED8FAAA9EE79E1D38A6696}\\
\hline
\texttt{Renewal\_Yang\_Mills\_Mass\_Gap\_Research\_Notes\_V122.tex}
& \(2848799\) &
\texttt{38199115BFE3F34F5D81DAC8BA4F5B466D34F261C6C0E7E18F73DA030490E88D}\\
\hline
\end{tabular}
\end{center}

The NS file remains byte-identical to the V135 audit.  Its V31 formation
theorem removes dyadic multiplicity exactly, but its generally available
payment is \(\int\|\Lambda u\|_2^4\), the continuation stock which the
argument is meant to obtain.  A formally one-use bound can therefore remain
analytically circular.  The sharper unresolved quantity is a restricted,
stopping-time-correlated occupation, not the unrestricted continuation norm.

YM V122 corrects a factor of two in its Wilson clock by returning to an exact
scalar normalization identity.  It then evaluates a transition-scale
finite difference directly through a coherent-state formula.  The direct
difference has a nonzero square-root limit even though the stronger separate
\(o(n^{-1})\) approximation contemplated in V121 is neither proved nor
needed.  V122 rejects only the stated scalar route and retains the complete
signed physical coefficient as a separate card.

\begin{assessment}[V140 audit adjustment]
\label{ass:V140-audit-adjustment}
For the phase system, the transferable rules are:
\begin{enumerate}[label=\textup{(\roman*)},leftmargin=2.8em]
\item derive the hypercharge and quotient normalization from the exact
      representation kernel before computing a global eta character;
\item compute the comparison triangle and adjacent eta ratio directly,
      rather than demand convergence of separately chosen phase sections;
\item require an independently controlled first moment of phase increments;
      a norm equivalent to the desired phase variation would be circular;
\item a failed scalar comparison does not by itself reject the complete
      line-valued gluing coefficient.
\end{enumerate}
\end{assessment}

\subsection{Integer hypercharge and the common central kernel}
\label{subsec:V140-central-kernel}

Put \(y:=6Y\) and let the covering group be
\begin{equation}
 \widetilde G
 :=SU(3)_c\times SU(2)_L\times U(1)_{\widehat Y},
 \label{eq:V140-covering-group}
\end{equation}
where \(e^{i\theta}\in U(1)_{\widehat Y}\) acts on integer charge \(y\)
as \(e^{iy\theta}\).  The all-left-handed fermion ledger and the Higgs are
\begin{equation}
 \begin{array}{c|rrrrrrr}
 &Q&u^c&d^c&L&e^c&\nu^c&H\\ \hline
 y&1&-4&2&-3&6&0&3.
 \end{array}
 \label{eq:V140-integer-hypercharges}
\end{equation}
Let
\begin{equation}
 \omega_3:=e^{2\pi i/3},
 \qquad
 k:=\left(\omega_3I_3,-I_2,e^{i\pi/3}\right)
 \in Z(\widetilde G).
 \label{eq:V140-k-generator}
\end{equation}

\begin{lemma}[The displayed generator acts trivially]
\label{lem:V140-k-trivial}
The element \(k\) acts as the identity on all seven multiplets in
\eqref{eq:V140-integer-hypercharges}.
\end{lemma}

\begin{proof}
On the multiplets in the table, the central phases are respectively
\begin{align*}
 Q:&\quad \omega_3(-1)e^{i\pi/3}=1,
 &u^c:&\quad \omega_3^{-1}e^{-4i\pi/3}=1,\\
 d^c:&\quad \omega_3^{-1}e^{2i\pi/3}=1,
 &L:&\quad (-1)e^{-i\pi}=1,\\
 e^c:&\quad e^{2\pi i}=1,
 &\nu^c:&\quad1,
 &H:&\quad(-1)e^{i\pi}=1.
\end{align*}
Thus \(k\) lies in the common kernel.
\end{proof}

\begin{theorem}[Exact Standard Model representation kernel]
\label{thm:V140-exact-kernel}
The common kernel of the complete fermion and Higgs representation of
\(\widetilde G\) is
\begin{equation}
 \boxed{
 K_{\rm SM}=\langle k\rangle\cong\mathbb Z_6.}
 \label{eq:V140-exact-kernel}
\end{equation}
\end{theorem}

\begin{proof}
Any element in the kernel must be central because the fundamental color and
weak representations occur.  Write it as
\begin{equation}
 z=\left(\omega_3^aI_3,(-1)^bI_2,e^{i\theta}\right),
 \qquad a\in\mathbb Z/3,
 \quad b\in\mathbb Z/2.
 \label{eq:V140-general-center}
\end{equation}
Trivial action on \(e^c\) gives \(e^{6i\theta}=1\), so
\(\theta=m\pi/3\) for some \(m\in\mathbb Z/6\).  Trivial action on
\(L\) gives
\[
 (-1)^b e^{-3i\theta}=(-1)^{b+m}=1,
\]
hence \(b\equiv m\pmod2\).  Trivial action on \(d^c\) gives
\[
 \omega_3^{-a}e^{2i\theta}
 =\omega_3^{m-a}=1,
\]
hence \(a\equiv m\pmod3\).  Therefore \(z=k^m\).  Conversely every
power of \(k\) is trivial by Lemma~\ref{lem:V140-k-trivial}.  This proves
the equality and the order-six claim.
\end{proof}

\begin{corollary}[All quotient branches and the faithful branch]
\label{cor:V140-quotient-branches}
The connected global gauge groups obtained from the same local Lie algebra
and on which all Standard Model fields are representations are
\begin{equation}
 G_\Gamma:=\widetilde G/\Gamma,
 \qquad
 \Gamma\in\{1,\mathbb Z_2,\mathbb Z_3,\mathbb Z_6\},
 \qquad \Gamma\le K_{\rm SM}.
 \label{eq:V140-quotient-branches}
\end{equation}
The action on the complete Standard Model field content is faithful exactly
for
\begin{equation}
 \boxed{G_{\rm SM}^{\rm faithful}=\widetilde G/\mathbb Z_6.}
 \label{eq:V140-faithful-group}
\end{equation}
For a proper subgroup \(\Gamma<K_{\rm SM}\), the residual finite center
\(K_{\rm SM}/\Gamma\) acts trivially on all listed local fields.
\end{corollary}

\begin{proof}
A representation of \(\widetilde G\) descends through a quotient precisely
when the quotient subgroup lies in its kernel.  The subgroups of the cyclic
group of order six are exactly those displayed.  The kernel after quotient
by \(\Gamma\) is \(K_{\rm SM}/\Gamma\), which is trivial precisely for
\(\Gamma=K_{\rm SM}\).
\end{proof}

\begin{firewall}[The local Lagrangian does not erase the global branch]
\label{fire:V140-global-branch}
All four groups in \eqref{eq:V140-quotient-branches} have the same Lie
algebra and the same perturbative anomaly polynomial.  Their principal
bundles, allowed flux sectors, one-form symmetries, spin-bordism groups, and
eta characters can differ.  Faithfulness selects the maximal quotient if
one requires an effective action on the printed fields; it does not prove
that the microscopic renewal construction has made that choice.
\end{firewall}

\begin{corollary}[Correct domain of the global anomaly character]
\label{cor:V140-bordism-domain}
On quotient branch \(G_\Gamma\), the residual phase character of V139 must
factor through the five-dimensional anomaly bordism data for that same
group,
\begin{equation}
 \chi_{\rm anom}^{\Gamma}:
 \Omega_5^{\rm Spin}(BG_\Gamma)\longrightarrow U(1),
 \label{eq:V140-bordism-character}
\end{equation}
with the appropriate refinement when the spacetime carries a combined
spin--gauge structure.  A computation for \(\Gamma=1\) cannot be silently
used for \(\Gamma=\mathbb Z_6\).
\end{corollary}

\begin{proof}
The mapping-torus eta invariant in
Theorem~\ref{thm:V139-BF-holonomy} is invariant under spin--gauge bordism
when the local polynomial vanishes.  Its bundle on the mapping torus is a
principal \(G_\Gamma\)-bundle, so its bordism class and the domain of the
character depend on \(BG_\Gamma\).
\end{proof}

\subsection{Direct transition ratios}
\label{subsec:V140-transition-ratios}

Let \(q\preceq q'\) be renewal cutoffs and let
\(J_{q'q}:\mathcal L_q\to\mathcal L_{q'}\) be the line comparison of V139.
Choose arbitrary local unit vectors \(s_q\in\mathcal L_q\) and define the
exact phase ratio
\begin{equation}
 J_{q'q}s_q=r_{q'q}s_{q'},
 \qquad r_{q'q}\in U(1).
 \label{eq:V140-transition-ratio}
\end{equation}
For \(q\preceq q'\preceq q''\), define its triangle defect
\begin{equation}
 \delta_{q''q'q}
 :=r_{q''q'}r_{q'q}r_{q''q}^{-1}\in U(1).
 \label{eq:V140-triangle-defect}
\end{equation}

\begin{lemma}[Gauge invariance of the triangle defect]
\label{lem:V140-defect-invariant}
Under arbitrary rephasings \(s_q\mapsto u_qs_q\), \(u_q\in U(1)\),
\begin{align}
 r_{q'q}&\longmapsto u_qr_{q'q}u_{q'}^{-1},
 \label{eq:V140-ratio-gauge}\\
 \delta_{q''q'q}&\longmapsto\delta_{q''q'q}.
 \label{eq:V140-defect-gauge}
\end{align}
Thus \(\delta\) can be computed directly without proving convergence of any
chosen representatives \(s_q\).
\end{lemma}

\begin{proof}
Insert \(u_qs_q\) into \eqref{eq:V140-transition-ratio} to obtain the first
formula.  Substitute the three transformed ratios into
\eqref{eq:V140-triangle-defect}.  Since the scalars commute, every \(u_q\),
\(u_{q'}\), and \(u_{q''}\) cancels with its inverse.
\end{proof}

\begin{theorem}[Exact comparison-cocycle test]
\label{thm:V140-cocycle-test}
The comparison maps satisfy
\begin{equation}
 J_{q''q'}J_{q'q}=J_{q''q}
 \label{eq:V140-map-cocycle}
\end{equation}
if and only if
\begin{equation}
 \boxed{\delta_{q''q'q}=1}
 \label{eq:V140-unit-defect}
\end{equation}
for every comparison triangle.  On a directed chain with a base cutoff,
unit triangle defects allow a rephasing for which every comparison ratio is
one.
\end{theorem}

\begin{proof}
Apply both sides of \eqref{eq:V140-map-cocycle} to the nonzero vector \(s_q\).
The left side gives
\(r_{q'q}r_{q''q'}s_{q''}\), while the right side gives
\(r_{q''q}s_{q''}\).  Equality is exactly
\eqref{eq:V140-unit-defect}.

Conversely, fix a base \(q_0\) and put
\(u_q=r_{qq_0}\), with \(u_{q_0}=1\).  The transformation law
\eqref{eq:V140-ratio-gauge} and the unit triangle relation give transformed
ratio one for every \(q\preceq q'\).  Hence all comparison maps identify the
rephased vectors.
\end{proof}

\begin{firewall}[A scalar section estimate can hide the obstruction]
\label{fire:V140-separate-sections}
The vectors \(s_q\) are gauge choices in one-dimensional fibres.  Bounds on
their individual arguments depend on those choices.  The triangle defect is
gauge invariant and is the direct analogue of the YM V122 finite difference:
it should be evaluated before any separate phase approximation.
\end{firewall}

\subsection{The exact first-moment requirement on a cofinal chain}
\label{subsec:V140-first-moment}

Let \(q_0\prec q_1\prec q_2\prec\cdots\) be a cofinal chain and abbreviate
\begin{equation}
 r_n:=r_{q_{n+1}q_n}\in U(1).
 \label{eq:V140-adjacent-ratio}
\end{equation}

\begin{theorem}[Absolute phase variation gives a cofinal limit]
\label{thm:V140-l1-phase}
If all triangle defects are one and
\begin{equation}
 \sum_{n=0}^\infty|1-r_n|<\infty,
 \label{eq:V140-l1-phase}
\end{equation}
then the ordered products
\begin{equation}
 R_N:=\prod_{n=0}^{N-1}r_n
 \label{eq:V140-phase-product}
\end{equation}
converge to a nonzero element of \(U(1)\).  After one limiting rephasing,
the transported determinant-line sections form the compatible cofinal phase
of Theorem~\ref{thm:V139-cofinal-phase} along the chain.
\end{theorem}

\begin{proof}
For \(M>N\), telescoping the product and using \(|r_n|=1\) gives
\begin{equation}
 |R_M-R_N|
 =\left|\prod_{n=N}^{M-1}r_n-1\right|
 \le\sum_{n=N}^{M-1}|r_n-1|.
 \label{eq:V140-product-Cauchy}
\end{equation}
The tail tends to zero by \eqref{eq:V140-l1-phase}, so \((R_N)\) is Cauchy.
Its limit has modulus one.  Unit triangle defects make every non-adjacent
comparison equal to the corresponding adjacent product by
Theorem~\ref{thm:V140-cocycle-test}; hence the limiting rephasing is
independent of subdivisions and supplies compatible transport.
\end{proof}

\begin{proposition}[Square summability is insufficient]
\label{prop:V140-l2-insufficient}
There is a sequence \(r_n\in U(1)\) for which
\begin{equation}
 \sum_{n=1}^\infty|1-r_n|^2<\infty
 \label{eq:V140-l2-finite}
\end{equation}
but the products \(R_N\) do not converge.  One may take
\begin{equation}
 r_n=e^{i/n}.
 \label{eq:V140-harmonic-phases}
\end{equation}
\end{proposition}

\begin{proof}
Since \(|1-e^{ix}|\le|x|\), the series in
\eqref{eq:V140-l2-finite} is bounded by \(\sum n^{-2}\).  On the other hand,
\begin{equation}
 R_N=\exp\left(i\sum_{n=1}^{N}\frac1n\right).
 \label{eq:V140-harmonic-product}
\end{equation}
If \(R_N\) converged, it would be Cauchy.  Choose \(M_N\) minimal with
\(\sum_{n=N+1}^{M_N}n^{-1}\ge\pi\).  The overshoot is at most
\(1/(N+1)\), so
\[
 |R_{M_N}-R_N|
 =\left|e^{i\sum_{n=N+1}^{M_N}1/n}-1\right|
 \longrightarrow2,
\]
contradicting the Cauchy property.
\end{proof}

\begin{corollary}[Phase-unmarking target]
\label{cor:V140-phase-unmarking}
A quadratic estimate such as
\(\sum|1-r_n|^2<\infty\) is not enough to construct the cofinal phase.
The next analytic target is either the absolute first moment
\eqref{eq:V140-l1-phase} or an exact signed/telescoping eta identity which
proves convergence of \eqref{eq:V140-phase-product} without taking absolute
values.
\end{corollary}

\begin{proof}
Sufficiency of the first alternative is
Theorem~\ref{thm:V140-l1-phase}; failure of the quadratic substitute is
Proposition~\ref{prop:V140-l2-insufficient}.  Conditional convergence can
occur through signed cancellation, so absolute summability is sufficient
but not asserted necessary.
\end{proof}

\begin{firewall}[The first moment must be independently paid]
\label{fire:V140-first-moment-circular}
Defining a control norm whose value is
\(\sum|1-r_n|\) and then assuming it finite would restate the desired
cofinal phase variation.  As in NS V31, the proof must bound that first
moment by renewal data already controlled independently, or retain an exact
signed eta telescope before absolute values.
\end{firewall}

\subsection{V140 ledger and next acquisition}
\label{subsec:V140-ledger}

\paragraph{New exact conclusions.}
\begin{enumerate}[label=\textup{(\arabic*)},leftmargin=2.8em]
\item The mandatory audit records exact NS V31 and YM V122 snapshots and
      changes the phase strategy at the comparison level.
\item Theorem~\ref{thm:V140-exact-kernel} proves that the common Standard
      Model fermion--Higgs kernel is exactly \(\mathbb Z_6\).
\item Corollary~\ref{cor:V140-quotient-branches} enumerates every subgroup
      quotient and identifies the faithful branch.
\item Corollary~\ref{cor:V140-bordism-domain} fixes the correct group-dependent
      domain for the remaining eta character.
\item Lemma~\ref{lem:V140-defect-invariant} and
      Theorem~\ref{thm:V140-cocycle-test} replace arbitrary individual phase
      convergence by exact gauge-invariant triangle defects.
\item Theorem~\ref{thm:V140-l1-phase} proves a usable sufficient cofinal
      criterion.
\item Proposition~\ref{prop:V140-l2-insufficient} proves sharply that
      quadratic phase control does not pay the required first moment.
\end{enumerate}

\paragraph{Next forced step.}
V141 should take the faithful branch
\(G_{\rm SM}^{\rm faithful}=\widetilde G/\mathbb Z_6\) and compute, or reduce
by an Atiyah--Hirzebruch spectral sequence, the torsion generators relevant
to \(\Omega_5^{\rm Spin}(BG_{\rm SM}^{\rm faithful})\).  Each surviving
generator must be evaluated by the full Standard Model eta character.  The
three proper-subgroup branches should remain listed but need not be computed
unless the microscopic renewal input selects a non-effective global form.

\begin{assessment}[V140 continuum status]
\label{ass:V140-continuum-status}
The full Standard--Model--Einstein continuum remains unproved.  V140 fixes
the exact faithful global gauge group and the direct cofinal phase defects.
The faithful-quotient bordism/eta character, a boundary eta trivialization,
extended BRST closure, microscopic gravitational normalization, nonlinear
existence, Einstein stress closure, and the cofinal interacting limit remain
open.
\end{assessment}

\section*{V140 exact checks and append-only record}
\addcontentsline{toc}{section}{V140 exact checks and append-only record}

The quotient-and-phase checksum is
\begin{equation}
 \boxed{
 K_{\rm SM}=\left\langle
 (\omega_3I_3,-I_2,e^{i\pi/3})\right\rangle\cong\mathbb Z_6,
 \quad
 \delta_{q''q'q}=r_{q''q'}r_{q'q}r_{q''q}^{-1},
 \quad
 \sum_n|1-r_n|<\infty\Longrightarrow\prod_nr_n\text{ converges}.}
 \label{eq:V140-final-check}
\end{equation}
The append operation preserves the complete V139 prefix byte for byte before
its former live terminator.  The assembled V140 file is checked for one live
final terminator, balanced environments and braces, duplicate labels, newly
introduced unresolved references, display math delimiters, and control
characters before the exact V139 predecessor is removed.  The audited
companion files are read only.

% =============================================================================
% END APPENDED VERSION V140 MATERIAL
% =============================================================================

% =============================================================================
% BEGIN APPENDED VERSION V141 MATERIAL
% =============================================================================

\section{V141: faithful-quotient spin bordism, vanishing eta holonomy,
and the remaining boundary comparison card}
\label{sec:V141}

\subsection{Purpose and source boundary}
\label{subsec:V141-purpose}

V140 proved from the field representations that the maximal faithful
Standard Model gauge group is
\begin{equation}
 G_6:=
 \frac{SU(3)\times SU(2)\times U(1)_{\widehat Y}}{\mathbb Z_6}.
 \label{eq:V141-G6}
\end{equation}
V141 now computes the bordism group which carries the residual closed
five-dimensional eta character on this branch.  The low-degree input and
the useful Puppe fibration are taken from the primary computation of
J.~Davighi, B.~Gripaios, and N.~Lohitsiri, \emph{Global anomalies in the
Standard Model(s) and Beyond}, JHEP 07 (2020) 232,
\url{https://arxiv.org/abs/1910.11277}, Section 4.5.  The total-degree-five
diagonal is reproduced below, so the vanishing conclusion does not rest on
an unseen differential.

Combined with the V138 local anomaly cancellation, the bordism result kills
the complete closed mapping-torus eta character for the faithful branch.
For the boundary programme, it also makes the Dai--Freed vector obtained
from a filling independent of the chosen filling whenever the boundary
datum is null-bordant.  Actual renewal mapping-torus boundaries have this
property by construction.  What remains is to build comparison maps across
cutoffs and prove their cocycle and convergence, not to search for another
faithful-branch global anomaly.

\subsection{The Puppe fibration}
\label{subsec:V141-fibration}

The quotient by the order-two subgroup of the common kernel is
\begin{equation}
 G_2\cong U(2)\times SU(3).
 \label{eq:V141-G2}
\end{equation}
The residual quotient from \(G_2\) to \(G_6\) has kernel \(\mathbb Z_3\),
so there is a group fibration
\begin{equation}
 \mathbb Z_3\longrightarrow U(2)\times SU(3)
 \longrightarrow G_6.
 \label{eq:V141-group-fibration}
\end{equation}

\begin{lemma}[Classifying-space Puppe fibration]
\label{lem:V141-Puppe}
The fibration \eqref{eq:V141-group-fibration} induces
\begin{equation}
 B(U(2)\times SU(3))
 \longrightarrow BG_6
 \longrightarrow K(\mathbb Z_3,2).
 \label{eq:V141-Puppe-fibration}
\end{equation}
\end{lemma}

\begin{proof}
Classifying \eqref{eq:V141-group-fibration} gives
\(B\mathbb Z_3\to B(U(2)\times SU(3))\to BG_6\).  The next map in its
Puppe sequence has target
\(B(B\mathbb Z_3)\simeq K(\mathbb Z_3,2)\), and its homotopy fibre is
\(B(U(2)\times SU(3))\).  This is exactly
\eqref{eq:V141-Puppe-fibration}.
\end{proof}

\subsection{Low-degree input}
\label{subsec:V141-input}

For the fibre in \eqref{eq:V141-Puppe-fibration}, the required spin-bordism
groups are
\begin{equation}
\begin{array}{c|cccccc}
q&0&1&2&3&4&5\\ \hline
\Omega_q^{\rm Spin}(B(U(2)\times SU(3)))
&\mathbb Z&\mathbb Z_2&\mathbb Z\oplus\mathbb Z_2&0&\mathbb Z^4&0.
\end{array}
\label{eq:V141-fibre-bordism}
\end{equation}
The base homology needed through degree five is
\begin{equation}
\begin{array}{c|cccccc}
p&0&1&2&3&4&5\\ \hline
H_p(K(\mathbb Z_3,2);\mathbb Z)
&\mathbb Z&0&\mathbb Z_3&0&\mathbb Z_3&0\\
H_p(K(\mathbb Z_3,2);\mathbb Z_2)
&\mathbb Z_2&0&0&0&0&0.
\end{array}
\label{eq:V141-base-homology}
\end{equation}

\begin{remark}[Why these inputs suffice]
\label{rem:V141-input-suffices}
The base \(K(\mathbb Z_3,2)\) is simply connected, so no twisted
\(\pi_1\)-action occurs in the coefficient system.  Every coefficient in
\eqref{eq:V141-fibre-bordism} is a direct sum of free abelian and
\(\mathbb Z_2\) factors.  Thus the two rows of
\eqref{eq:V141-base-homology}, together with additivity in the coefficient
group, determine every entry used below.
\end{remark}

\subsection{The total-degree-five diagonal}
\label{subsec:V141-diagonal}

The generalized homology spectral sequence associated with
\eqref{eq:V141-Puppe-fibration} has second page
\begin{equation}
 E^2_{p,q}
 =H_p\!\left(K(\mathbb Z_3,2);
 \Omega_q^{\rm Spin}(B(U(2)\times SU(3)))\right)
 \Longrightarrow\Omega_{p+q}^{\rm Spin}(BG_6).
 \label{eq:V141-AHSS}
\end{equation}

\begin{theorem}[Every total-degree-five entry vanishes]
\label{thm:V141-diagonal-zero}
On the diagonal \(p+q=5\),
\begin{equation}
 \boxed{E^2_{p,5-p}=0\qquad(0\le p\le5).}
 \label{eq:V141-diagonal-zero}
\end{equation}
More explicitly,
\begin{align}
 E^2_{0,5}&=H_0(K;0)=0,
 &E^2_{1,4}&=H_1(K;\mathbb Z^4)=0,
 \notag\\
 E^2_{2,3}&=H_2(K;0)=0,
 &E^2_{3,2}&=H_3(K;\mathbb Z\oplus\mathbb Z_2)=0,
 \notag\\
 E^2_{4,1}&=H_4(K;\mathbb Z_2)=0,
 &E^2_{5,0}&=H_5(K;\mathbb Z)=0,
 \label{eq:V141-six-zeroes}
\end{align}
where \(K=K(\mathbb Z_3,2)\).
\end{theorem}

\begin{proof}
Insert the coefficient group in degree \(q=5-p\) from
\eqref{eq:V141-fibre-bordism}.  The first and third entries vanish because
their coefficient groups are zero.  The second uses
\(H_1(K;\mathbb Z)=0\).  The fourth is the direct sum of
\(H_3(K;\mathbb Z)=0\) and
\(H_3(K;\mathbb Z_2)=0\).  The fifth and sixth are the corresponding
entries of \eqref{eq:V141-base-homology}.  This gives all six equalities.
\end{proof}

\begin{corollary}[Faithful-branch spin bordism vanishes]
\label{cor:V141-Omega5-zero}
\begin{equation}
 \boxed{\Omega_5^{\rm Spin}(BG_6)=0.}
 \label{eq:V141-Omega5-zero}
\end{equation}
\end{corollary}

\begin{proof}
Every later page is a subquotient of the second page.  Hence
\eqref{eq:V141-diagonal-zero} implies
\(E^\infty_{p,5-p}=0\) for all \(p\).  These entries are the associated
graded pieces of the finite filtration of
\(\Omega_5^{\rm Spin}(BG_6)\).  A filtered abelian group whose every graded
piece is zero is itself zero.  No differential or extension problem remains
on this diagonal.
\end{proof}

\begin{firewall}[What was imported]
\label{fire:V141-imported-input}
V141 reproduces the entire total-degree-five spectral-sequence deduction.
It invokes the published low-degree groups
\eqref{eq:V141-fibre-bordism} and homology table
\eqref{eq:V141-base-homology}; it does not rederive those two independent
computations from cellular chains.  The vanishing result uses no unknown
differential because its relevant second-page entries are already zero.
\end{firewall}

\subsection{The closed eta character is trivial}
\label{subsec:V141-closed-eta}

Let \(\tau_{\rm SM}(Y)\) be the product of exponentiated reduced eta
invariants of all physical Standard Model Weyl representations on a closed
spin five-manifold \(Y\) with \(G_6\)-bundle.

\begin{theorem}[No faithful-branch closed global anomaly]
\label{thm:V141-closed-anomaly-zero}
For the Standard Model representation of V138,
\begin{equation}
 \boxed{\tau_{\rm SM}(Y)=1}
 \label{eq:V141-closed-eta-one}
\end{equation}
for every closed spin \(G_6\) five-manifold \(Y\).  Equivalently, the
faithful-branch character
\begin{equation}
 \chi_{\rm anom}^{6}:
 \Omega_5^{\rm Spin}(BG_6)\longrightarrow U(1)
 \label{eq:V141-faithful-character}
\end{equation}
is trivial.
\end{theorem}

\begin{proof}
Corollary~\ref{cor:V138-local-cancellation} gives
\(I_6^{\rm SM}=0\).  Therefore the exponentiated eta invariant is unchanged
under spin \(G_6\)-bordism and defines the homomorphism
\eqref{eq:V141-faithful-character}.  Its domain is the zero group by
Corollary~\ref{cor:V141-Omega5-zero}; the only homomorphism from the zero
group sends its sole element to the identity of \(U(1)\).
\end{proof}

\begin{corollary}[Every closed mapping-torus holonomy is one]
\label{cor:V141-closed-holonomy}
For a closed four-manifold and every loop \(\gamma\) of faithful-branch
Standard Model backgrounds,
\begin{equation}
 \operatorname{Hol}_\gamma(\mathcal L_{\rm ch})=1.
 \label{eq:V141-closed-holonomy}
\end{equation}
Thus the chiral determinant line has a nonzero parallel section on each
connected parameter component, unique up to one constant phase.
\end{corollary}

\begin{proof}
The Bismut--Freed formula
\eqref{eq:V139-BF-holonomy} identifies the holonomy with the adiabatic eta
invariant of the closed five-dimensional mapping torus.  Apply
Theorem~\ref{thm:V141-closed-anomaly-zero}, then use
Theorem~\ref{thm:V139-flat-trivialization}.
\end{proof}

\begin{remark}[The mod-two count is absorbed differently]
\label{rem:V141-interplay}
For the direct-product and \(\mathbb Z_3\)-quotient branches, a
\(\mathbb Z_2\) bordism class retains the traditional \(SU(2)\) test.  For
the \(\mathbb Z_6\) faithful branch, the bordism group itself vanishes.
The even-doublet restriction is nevertheless already encoded in the
allowed hypercharges and the local \(SU(2)^2U(1)\) anomaly equation.  This
is consistent with, and sharper in global form than, the explicit Standard
Model parity count in Theorem~\ref{thm:V138-Witten-parity}.
\end{remark}

\subsection{Boundary filling independence}
\label{subsec:V141-boundary-filling}

Let \(W\) be a closed spin four-manifold with faithful \(G_6\)-bundle and
suppose \([W]=0\) in \(\Omega_4^{\rm Spin}(BG_6)\).  Choose a spin
\(G_6\) five-manifold \(X\) with \(\partial X=-W\).  By the Dai--Freed
theorem, its eta element has type
\begin{equation}
 \ell_X:=\tau_{\rm SM}(X)
 \in\operatorname{Det}(D_W^+),
 \label{eq:V141-filling-vector}
\end{equation}
using the orientation-duality convention of V139.

\begin{theorem}[Canonical vector on a null-bordant boundary datum]
\label{thm:V141-filling-independent}
The vector \(\ell_X\) in \eqref{eq:V141-filling-vector} is independent of
the chosen filling \(X\).  Denote it by \(\ell_W\).  If
\(Y\) is any five-manifold with \(\partial Y=W\), then
\begin{equation}
 \boxed{
 \left\langle\tau_{\rm SM}(Y),\ell_W\right\rangle=1.}
 \label{eq:V141-boundary-pairing-one}
\end{equation}
\end{theorem}

\begin{proof}
Let \(X_0\) and \(X_1\) both fill \(-W\).  Glue \(-X_1\) to \(X_0\) along
their common boundary.  The ratio of the two line vectors is, by the
Dai--Freed gluing law, the scalar eta invariant of the resulting closed
five-manifold.  It is one by
Theorem~\ref{thm:V141-closed-anomaly-zero}; hence
\(\ell_{X_0}=\ell_{X_1}\).

For the second assertion, glue \(Y\) to any filling \(X\) of \(-W\).
The pairing is the eta invariant of the closed glued five-manifold, again
one by Theorem~\ref{thm:V141-closed-anomaly-zero}.
\end{proof}

\begin{corollary}[Renewal mapping-torus phase]
\label{cor:V141-renewal-mapping-phase}
For every loop \(\gamma\) of four-dimensional renewal backgrounds with
boundary, the induced four-dimensional boundary mapping torus
\(W_\gamma=\partial Y_\gamma\) is null-bordant.  The canonical vector
\(\ell_{W_\gamma}\) exists and the paired eta holonomy of
\eqref{eq:V139-paired-holonomy} is one.
\end{corollary}

\begin{proof}
The five-dimensional mapping cylinder closed in the parameter direction is
\(Y_\gamma\), and
\(\partial Y_\gamma=W_\gamma\) by
\eqref{eq:V139-boundary-mapping-torus}.  Thus \(W_\gamma\) is null-bordant.
Apply Theorem~\ref{thm:V141-filling-independent}.
\end{proof}

\begin{firewall}[Arbitrary boundary components]
\label{fire:V141-arbitrary-boundary}
The group \(\Omega_4^{\rm Spin}(BG_6)\) is not zero.  V141 therefore does
not assign a filling vector to an arbitrary closed four-dimensional
spin--gauge datum.  The canonical conclusion applies to null-bordant data,
including the boundary mapping tori that actually arise from a renewal
loop.  A functor defined on every bordism component still needs one coherent
reference vector per nonzero component.
\end{firewall}

\subsection{Updated phase ledger}
\label{subsec:V141-phase-ledger}

\begin{theorem}[Finite-cutoff faithful phase card]
\label{thm:V141-finite-phase-card}
At a fixed renewal cutoff on the faithful \(G_6\) branch, assume a smooth
spectral section as in Definition~\ref{def:V139-spectral-section}.  Then:
\begin{enumerate}[label=\textup{(\roman*)},leftmargin=2.8em]
\item the determinant line exists;
\item its local anomaly curvature vanishes;
\item every closed and renewal-boundary mapping-torus holonomy is one;
\item consequently it admits a nonzero parallel phase section on each
      connected component, unique up to a constant phase.
\end{enumerate}
\end{theorem}

\begin{proof}
Item \textup{(i)} is Theorem~\ref{thm:V139-determinant-line}.
Item \textup{(ii)} is Theorem~\ref{thm:V138-local-flatness}.
The two cases in \textup{(iii)} are Corollaries
\ref{cor:V141-closed-holonomy} and
\ref{cor:V141-renewal-mapping-phase}.  The flat-line criterion
Theorem~\ref{thm:V139-flat-trivialization} proves \textup{(iv)}.
\end{proof}

\begin{firewall}[Finite-cutoff phase is not the cofinal phase]
\label{fire:V141-finite-not-cofinal}
Theorem~\ref{thm:V141-finite-phase-card} removes the faithful-branch anomaly
obstruction.  It does not construct the renewal maps
\(J_{q'q}\), prove their exact triangle identity, or pay the phase first
moment in Theorem~\ref{thm:V140-l1-phase}.  Independent constant choices at
each cutoff can still fail to define one continuum phase.
\end{firewall}

\subsection{V141 ledger and next acquisition}
\label{subsec:V141-ledger}

\paragraph{New exact conclusions.}
\begin{enumerate}[label=\textup{(\arabic*)},leftmargin=2.8em]
\item Lemma~\ref{lem:V141-Puppe} places the faithful quotient in the useful
      \(K(\mathbb Z_3,2)\) fibration.
\item Theorem~\ref{thm:V141-diagonal-zero} checks all six entries of the
      relevant spectral-sequence diagonal explicitly.
\item Corollary~\ref{cor:V141-Omega5-zero} proves
      \(\Omega_5^{\rm Spin}(BG_6)=0\) with no differential ambiguity.
\item Theorem~\ref{thm:V141-closed-anomaly-zero} closes the faithful-branch
      global eta character after the V138 local cancellation.
\item Theorem~\ref{thm:V141-filling-independent} constructs the canonical
      Dai--Freed vector on every null-bordant boundary datum.
\item Corollary~\ref{cor:V141-renewal-mapping-phase} applies that vector to
      the actual boundary mapping tori of the programme.
\item Theorem~\ref{thm:V141-finite-phase-card} closes the finite-cutoff
      faithful chiral phase, conditional only on its spectral-section datum.
\end{enumerate}

\paragraph{Next forced step.}
V142 should construct \(J_{q'q}\) from the actual renewal collar and
interior comparison operators.  It must split the comparison into a
finite-dimensional spectral-flow line and an invertible high-energy
determinant, prove the exact three-cutoff cocycle before estimates, and then
seek either an independently summable eta increment or an exact signed
telescope.  If no uniform boundary spectral section exists, the analysis
must first identify its family-index obstruction rather than choose
cutoff-wise projectors independently.

\begin{assessment}[V141 continuum status]
\label{ass:V141-continuum-status}
The full Standard--Model--Einstein continuum remains unproved.  V141 closes
the local and global chiral anomaly and finite-cutoff phase cards on the
faithful \(\mathbb Z_6\) branch.  Renewal comparison maps and their cofinal
phase control, the extended BRST boundary complex, microscopic gravitational
normalization, nonlinear existence, Einstein stress closure, and the
cofinal interacting limit remain open.
\end{assessment}

\section*{V141 exact checks and append-only record}
\addcontentsline{toc}{section}{V141 exact checks and append-only record}

The faithful-phase checksum is
\begin{equation}
 \boxed{
 E^2_{p,5-p}=0\ (0\le p\le5),
 \qquad
 \Omega_5^{\rm Spin}(BG_6)=0,
 \qquad
 \tau_{\rm SM}(Y_{\rm closed})=1.}
 \label{eq:V141-final-check}
\end{equation}
The append operation preserves the complete V140 prefix byte for byte before
its former live terminator.  The assembled V141 file is checked for one live
final terminator, balanced environments and braces, duplicate labels, newly
introduced unresolved references, display math delimiters, and control
characters before the exact V140 predecessor is removed.  The companion
files are not modified.

% =============================================================================
% END APPENDED VERSION V141 MATERIAL
% =============================================================================

% =============================================================================
% BEGIN APPENDED VERSION V142 MATERIAL
% =============================================================================

\section{V142: exact determinant-functor comparison maps, shell fusion,
and the spectral-section obstruction}
\label{sec:V142}

\subsection{Purpose and the first-order object}
\label{subsec:V142-purpose}

V141 removes the faithful-branch local and global anomaly, but a vanishing
anomaly does not manufacture maps between different cutoff lines.  This
note constructs such maps at a common finite regulator.  The construction
uses the first-order chiral complex and not merely the positive operator
appearing in the Gaussian modulus.  Its hypotheses are printed as exact
algebraic gates, so that a small cutoff defect cannot be mistaken for a
determinant-line morphism.

At cutoff (q), let
\begin{equation}
 C_q:=\left[,H_q^+\xrightarrow{D_q}H_q^-,\right]
 \label{eq:V142-complex}
\end{equation}
be the regulated Standard Model chiral complex with the spectral-section
boundary condition of V139.  It is placed in degrees zero and one.  Write
\begin{equation}
 \lambda(C_q)
 :=\det H^0(C_q)\otimes\det H^1(C_q)^{-1}
 =\operatorname{Det}(D_q).
 \label{eq:V142-determinant-line}
\end{equation}
All finite spaces are expressed in the transported physical and boundary
frames used by the collar and Calderon analysis of V119--V124.

\begin{definition}[Admissible renewal comparison]
\label{def:V142-admissible-comparison}
For (q\preceq q'), an admissible regulated comparison consists of
injective maps
\begin{equation}
 \iota_{q'q}^{\pm}:H_q^\pm\longrightarrow H_{q'}^\pm
 \label{eq:V142-inclusions}
\end{equation}
which obey the exact chain identity
\begin{equation}
 \boxed{D_{q'}\iota_{q'q}^+
 =\iota_{q'q}^-D_q.}
 \label{eq:V142-chain-gate}
\end{equation}
The shell complex is the quotient
\begin{equation}
 C_{q'/q}:=C_{q'}/\iota_{q'q}C_q.
 \label{eq:V142-shell-complex}
\end{equation}
For three cutoffs the inclusions are required to satisfy
\begin{equation}
 \iota_{q''q'}^\pm\iota_{q'q}^\pm=\iota_{q''q}^\pm.
 \label{eq:V142-inclusion-cocycle}
\end{equation}
\end{definition}

\begin{proposition}[The exact chain gate is necessary]
\label{prop:V142-defect-gate}
For arbitrary transported injections define
\begin{equation}
 \mathcal E_{q'q}
 :=D_{q'}\iota_{q'q}^+-\iota_{q'q}^-D_q.
 \label{eq:V142-chain-defect}
\end{equation}
If (\mathcal E_{q'q}\ne0), the injections do not define a map on
cohomology and hence do not induce the determinant-functor map used below.
The estimate (\|\mathcal E_{q'q}\|\ll1) does not change this conclusion.
One must either prove \eqref{eq:V142-chain-gate} or enlarge the datum by a
specified mapping-cone complex carrying the defect.
\end{proposition}

\begin{proof}
If (x\in\ker D_q), then
\[
 D_{q'}\iota_{q'q}^+x=\mathcal E_{q'q}x.
\]
Thus a closed vector need not be sent to a closed vector.  Similarly, the
class of (\iota_{q'q}^-y) in the target cokernel can depend on the chosen
representative (y) modulo (\operatorname{im}D_q).  These are algebraic
failures and persist for every nonzero defect, regardless of its norm.
\end{proof}

\begin{remark}[A regulator which preserves the gate]
\label{rem:V142-paired-regulator}
At a fixed geometry, complete singular pairs of (D_q) give a regulator
with zero chain defect: if (D_q^*D_qu=\lambda u) with (\lambda>0), retain
(u) and (\lambda^{-1/2}D_qu) together.  Across two renewal geometries,
this observation does not prove \eqref{eq:V142-chain-gate}; the transported
collar/interior comparison must still be checked.  This is the chiral
analogue of the common Hodge-cutoff gate
\eqref{eq:V120-regulator-chain-gate}.
\end{remark}

\subsection{The determinant line of one shell}
\label{subsec:V142-one-shell}

An admissible comparison gives a short exact sequence of two-term
complexes
\begin{equation}
 0\longrightarrow C_q\xrightarrow{\iota_{q'q}}C_{q'}
 \longrightarrow C_{q'/q}\longrightarrow0.
 \label{eq:V142-short-exact}
\end{equation}

\begin{theorem}[Canonical finite determinant isomorphism]
\label{thm:V142-determinant-isomorphism}
The exact sequence \eqref{eq:V142-short-exact} has a canonical isomorphism
\begin{equation}
 \mu_{q'q}:\lambda(C_q)\otimes\lambda(C_{q'/q})
 \xrightarrow{\ \cong\ }\lambda(C_{q'}).
 \label{eq:V142-mu}
\end{equation}
For a filtration by three cutoffs, these isomorphisms are associative when
the factors are kept in increasing-shell order, including the determinant
functor's Koszul sign.
\end{theorem}

\begin{proof}
Choose degreewise splittings of \eqref{eq:V142-short-exact}.  Wedge an
ordered basis of (C_q^j) with lifted ordered quotient bases in each degree
(j=0,1).  This gives
\[
 \det C_{q'}^j\cong\det C_q^j\otimes\det C_{q'/q}^j.
\]
Taking the degree-zero line tensored with the inverse degree-one line and
then applying the canonical torsion isomorphism from chain spaces to
cohomology gives \eqref{eq:V142-mu}.  Changing a lift adds an upper
triangular matrix with identity diagonal, so the wedge and the resulting
map do not change.  For a three-step filtration, both parenthesizations
wedge the same ordered basis.  The sign incurred when degree-one factors
are moved into the declared order is precisely the Koszul sign in the
determinant functor, proving strict associativity in that convention.
\end{proof}

Choose a common spectral band in the shell complex and write
\begin{equation}
 C_{q'/q}=F_{q'/q}\oplus G_{q'/q},
 \qquad D_{q'/q}|_{G_{q'/q}}=:A_{q'/q},
 \label{eq:V142-low-high-split}
\end{equation}
where (F_{q'/q}) contains every crossing mode and
(A_{q'/q}) is invertible.

\begin{lemma}[Low spectral-flow line and high acyclic torsion]
\label{lem:V142-low-high-line}
There is a canonical factorization
\begin{equation}
 \lambda(C_{q'/q})\cong
 \lambda(F_{q'/q})\otimes\lambda(G_{q'/q}).
 \label{eq:V142-low-high-line}
\end{equation}
The invertible complex (G_{q'/q}) has a canonical nonzero torsion vector
(\tau(A_{q'/q})).  All kernel crossings are retained in the finite line
(\lambda(F_{q'/q})); none is divided out of the high determinant.
\end{lemma}

\begin{proof}
Apply Theorem~\ref{thm:V142-determinant-isomorphism} to the direct-sum
short exact sequence.  Since (A_{q'/q}) is an isomorphism, the high
complex is acyclic.  If (e) is a basis of its degree-zero space, the
basis (A_{q'/q}e) in degree one gives the torsion vector; a change
(e\mapsto eM) multiplies both wedges by (\det M), which cancels.  The
finite complement is exactly the relative spectral-section line of
Proposition~\ref{prop:V139-section-change}.
\end{proof}

On the faithful branch, the shell is a renewal mapping cylinder.  Its
low-mode line has the null-bordant Dai--Freed vector supplied by
Theorem~\ref{thm:V141-filling-independent}.  Denote it by
(\ell_{q'/q}), and put
\begin{equation}
 \sigma_{q'/q}
 :=\ell_{q'/q}\otimes\tau(A_{q'/q})
 \in\lambda(C_{q'/q})\setminus\{0\}.
 \label{eq:V142-shell-vector}
\end{equation}
For the phase line use its Quillen-unit normalization
(\widehat\sigma_{q'/q}:=\sigma_{q'/q}/\|\sigma_{q'/q}\|_Q).

\begin{definition}[Renewal determinant-line comparison]
\label{def:V142-J}
For an admissible comparison define
\begin{equation}
 \boxed{
 J_{q'q}(v):=\mu_{q'q}
 \bigl(v\otimes\widehat\sigma_{q'/q}\bigr),
 \qquad v\in\lambda(C_q).}
 \label{eq:V142-J}
\end{equation}
\end{definition}

\subsection{Exact shell fusion and the comparison cocycle}
\label{subsec:V142-fusion}

For (q\preceq q'\preceq q''), the quotient filtration gives
\begin{equation}
 0\longrightarrow C_{q'/q}\longrightarrow C_{q''/q}
 \longrightarrow C_{q''/q'}\longrightarrow0.
 \label{eq:V142-shell-filtration}
\end{equation}
Let (\mu^{\rm sh}_{q''q'q}) be its determinant isomorphism, with the
same increasing-shell convention as in
Theorem~\ref{thm:V142-determinant-isomorphism}.

\begin{theorem}[Exact shell-vector fusion]
\label{thm:V142-shell-fusion}
If the spectral bands are chosen from a common refinement, then
\begin{equation}
 \boxed{
 \mu^{\rm sh}_{q''q'q}
 \bigl(\widehat\sigma_{q'/q}\otimes
       \widehat\sigma_{q''/q'}\bigr)
 =\widehat\sigma_{q''/q}.}
 \label{eq:V142-shell-fusion}
\end{equation}
\end{theorem}

\begin{proof}
On the finite crossing spaces, exact-sequence multiplicativity is
Proposition~\ref{prop:V139-section-change}.  On the high spaces, eliminate
successive shell blocks in filtration order.  The eliminators are block
triangular with identity diagonal, exactly as in
Theorem~\ref{thm:V120-diagonal-equivalence}; hence their determinants are
one and the two high torsions multiply to the direct-shell torsion.  This
is also the finite first-order counterpart of the Schur determinant
identity \eqref{eq:V123-gluing-determinant}.  Finally, the two
Dai--Freed filling vectors pair on the intermediate boundary.  The eta
gluing law of Theorem~\ref{thm:V139-DF-element} gives the direct filling
vector.  V141 makes the possible closed glued correction equal to one.
Thus the unnormalized shell vectors fuse.  The Quillen gluing norm is
multiplicative for this same finite exact sequence, so normalization
preserves the equality.
\end{proof}

\begin{corollary}[The comparison triangle is exactly flat]
\label{cor:V142-exact-cocycle}
For admissible nested renewal comparisons on the faithful branch,
\begin{equation}
 \boxed{J_{q''q'}J_{q'q}=J_{q''q}.}
 \label{eq:V142-exact-cocycle}
\end{equation}
Consequently the triangle defect of \eqref{eq:V140-triangle-defect} is
identically one before any cofinal estimate is made.
\end{corollary}

\begin{proof}
Apply the two sides to (v\in\lambda(C_q)).  Associativity of the
determinant isomorphisms and
Theorem~\ref{thm:V142-shell-fusion} turn both expressions into
\[
 \mu\bigl(v\otimes\widehat\sigma_{q''/q}\bigr).
\]
This is (J_{q''q}v).
\end{proof}

\begin{firewall}[The positive Schur determinant supplies the norm only]
\label{fire:V142-modulus-not-phase}
For an invertible chiral block,
\begin{equation}
 |\det D|^2=\det(D^*D).
 \label{eq:V142-modulus-square}
\end{equation}
Thus the normalized positive matching operators of V123--V124 determine
the Quillen norm of the high shell vector.  They cannot determine
(\widehat\sigma_{q'/q}), its eta phase, or the finite spectral-flow line.
Using (\det\mathcal W>0) as (J_{q'q}) would erase exactly the chiral
datum that V139 retained.
\end{firewall}

\subsection{Uniform spectral sections and their obstruction}
\label{subsec:V142-index-obstruction}

Let (A_{q,b}) be the self-adjoint tangential chiral family over a compact
parameter space (\mathcal B).  Its analytic family index is
\begin{equation}
 \operatorname{Ind}(A_q)\in K^1(\mathcal B).
 \label{eq:V142-K1-index}
\end{equation}

\begin{theorem}[Spectral-section existence gate]
\label{thm:V142-spectral-section-gate}
A global spectral section for (A_{q,b}) exists if and only if
\begin{equation}
 \boxed{\operatorname{Ind}(A_q)=0\quad\hbox{in }K^1(\mathcal B).}
 \label{eq:V142-index-zero}
\end{equation}
For a compatible cofinal family, one additionally needs spectral sections
whose relative (K^0(\mathcal B)) classes obey the shell-fusion law.
\end{theorem}

\begin{proof}
This is the Melrose--Piazza spectral-section theorem.  Necessity follows
because a spectral section is a finite-rank perturbation of the positive
polarization and supplies a global trivialization of the odd Fredholm
class.  Conversely, vanishing of that class permits a finite-rank
self-adjoint perturbation making the family invertible; its positive
spectral projection is a spectral section.  Differences of two choices
are finite-dimensional virtual bundles, hence classes in
(K^0(\mathcal B)), and add under composition as in
\eqref{eq:V139-relative-cocycle}.  A detailed primary proof is given by
Melrose and Piazza, \emph{Families of Dirac operators, boundaries and the
b-calculus}.
\end{proof}

\begin{assessment}[What V142 closes and what it exposes]
\label{ass:V142-comparison-status}
V142 closes the algebraic comparison triangle, conditional on two
checkable inputs: the transported renewal maps must satisfy the exact chain
gate \eqref{eq:V142-chain-gate}, and the tangential family must satisfy the
(K^1)-vanishing gate \eqref{eq:V142-index-zero}.  Neither gate follows
from a norm-small defect.  Once they hold, no separate triangle estimate
and no phase first moment are needed to construct the abstract directed
system of lines.  A first moment or a direct-tail theorem is still needed
to realize that system as a continuous scalar phase in a fixed continuum
frame.
\end{assessment}

\subsection{V142 ledger and next acquisition}
\label{subsec:V142-ledger}

\paragraph{New exact conclusions.}
\begin{enumerate}[label=\textup{(\arabic*)},leftmargin=2.8em]
\item Proposition~\ref{prop:V142-defect-gate} proves that a merely small
      cutoff intertwining defect does not induce a determinant-line map.
\item Theorem~\ref{thm:V142-determinant-isomorphism} constructs the
      comparison from the exact sequence of the physical and shell
      complexes.
\item Lemma~\ref{lem:V142-low-high-line} separates the finite spectral-flow
      line from the invertible high-energy torsion.
\item Definition~\ref{def:V142-J} gives the actual phase comparison map
      after the faithful Dai--Freed shell vector is inserted.
\item Theorem~\ref{thm:V142-shell-fusion} and
      Corollary~\ref{cor:V142-exact-cocycle} prove the exact three-cutoff
      cocycle before estimates.
\item Theorem~\ref{thm:V142-spectral-section-gate} identifies the precise
      (K^1) obstruction to a uniform spectral section.
\end{enumerate}

\paragraph{Next forced step.}
V143 should distinguish the algebraic direct-limit line from its analytic
realization in one continuum frame.  It should derive a direct-tail
Fredholm-determinant criterion which implies phase convergence without
summing adjacent absolute increments, and then identify the Schatten class
actually furnished by the three-dimensional boundary comparison.  If raw
order is too weak, V144 must go upstream to the symbol subtraction rather
than assume the missing trace norm.

\begin{assessment}[V142 continuum status]
\label{ass:V142-continuum-status}
The full Standard--Model--Einstein continuum remains unproved.  The
faithful chiral phase now has an exact finite comparison construction and
an exact conditional cocycle.  The chain gate, family-index gate, analytic
continuum realization, extended BRST boundary complex, microscopic
gravitational normalization, nonlinear existence, Einstein stress closure,
and interacting cofinal limit remain open.
\end{assessment}

\section*{V142 exact checks and append-only record}
\addcontentsline{toc}{section}{V142 exact checks and append-only record}

The determinant-line comparison checksum is
\begin{equation}
 \boxed{
 D_{q'}\iota_{q'q}^+=\iota_{q'q}^-D_q,
 \qquad
 J_{q''q'}J_{q'q}=J_{q''q},
 \qquad
 \operatorname{Ind}(A_q)=0\in K^1(\mathcal B).}
 \label{eq:V142-final-check}
\end{equation}
The append operation preserves the complete V141 prefix byte for byte before
its former live terminator.  The assembled V142 file is checked for one live
final terminator, balanced environments and braces, duplicate labels, newly
introduced unresolved references, display math delimiters, and control
characters before the exact V141 predecessor is removed.  The companion
files are not modified.

% =============================================================================
% END APPENDED VERSION V142 MATERIAL
% =============================================================================

% =============================================================================
% BEGIN APPENDED VERSION V143 MATERIAL
% =============================================================================

\section{V143: direct-tail phase convergence, the boundary Schatten
threshold, and the trace-class deficit}
\label{sec:V143}

\subsection{Purpose}
\label{subsec:V143-purpose}

V142 constructs an exact directed system of finite determinant lines when
the chain and spectral-section gates hold.  Such a directed system exists
algebraically without an infinite product.  To obtain a scalar continuum
phase in one analytic frame, however, the comparison operators must also
converge in a topology on which a determinant is continuous.

V143 first proves a direct-tail criterion.  It is stronger in structure
than estimating independently chosen adjacent phases and can succeed even
when their absolute variations are not summable.  The second part computes
the Schatten threshold forced by a three-dimensional boundary.  It shows
that a generic order-(-1) Calderon comparison is in
(\mathfrak S_p) only for (p>3), whereas the ordinary Fredholm
determinant requires (\mathfrak S_1).  This is the precise analytic gap
which the next symbol subtraction must close.

\subsection{The algebraic limit line needs no infinite product}
\label{subsec:V143-algebraic-limit}

Let (q_0\preceq q_1\preceq\cdots) be cofinal and suppose the maps of
V142 satisfy \eqref{eq:V142-exact-cocycle}.  Define an equivalence relation
on pairs ((q,v)), (v\in\lambda(C_q)), by
\begin{equation}
 (q,v)\sim(q',v')
 \quad\Longleftrightarrow\quad
 J_{r q}v=J_{r q'}v'
 \text{ for some }r\succeq q,q'.
 \label{eq:V143-equivalence}
\end{equation}

\begin{proposition}[Algebraic directed-limit line]
\label{prop:V143-algebraic-line}
The quotient
\begin{equation}
 \lambda_{\rm alg}:=\varinjlim_q(\lambda(C_q),J_{q'q})
 \label{eq:V143-algebraic-line}
\end{equation}
is a one-dimensional complex vector space, and every canonical map
(\lambda(C_q)\to\lambda_{\rm alg}) is an isomorphism.  No convergence of
scalar representatives is required for this statement.
\end{proposition}

\begin{proof}
The exact cocycle makes \eqref{eq:V143-equivalence} transitive and makes
addition and scalar multiplication independent of a common upper cutoff.
Every (J_{q'q}) is an isomorphism of one-dimensional spaces, so every
class has a unique representative at any fixed later cutoff.  Hence the
canonical maps are bijective and the quotient is one dimensional.
\end{proof}

\begin{firewall}[Algebraic existence is not analytic convergence]
\label{fire:V143-algebraic-not-analytic}
The space \eqref{eq:V143-algebraic-line} has no continuum Quillen norm or
preferred embedding into a fixed scalar line until a compatible analytic
topology is supplied.  Proposition~\ref{prop:V143-algebraic-line} closes
the categorical cocycle problem; it does not prove convergence of the
regularized determinant, stress tensor, or effective action.
\end{firewall}

\subsection{A direct-tail Fredholm criterion}
\label{subsec:V143-direct-tail}

Fix a polarized boundary Hilbert space (\mathcal H_\Sigma) and transport
all high-energy comparisons to it.  Suppose the phase comparison at cutoff
(q) is represented by an invertible operator
\begin{equation}
 T_q=I+K_q,
 \qquad K_q\in\mathfrak S_1(\mathcal H_\Sigma).
 \label{eq:V143-Tq}
\end{equation}
For (q\preceq q'), put
\begin{equation}
 r_{q'q}
 :=\operatorname{ph}\det_{\!F}(T_{q'}T_q^{-1}),
 \qquad
 \operatorname{ph}z:=z/|z|.
 \label{eq:V143-relative-phase}
\end{equation}

\begin{lemma}[Fredholm determinant continuity]
\label{lem:V143-det-continuity}
For (A,B\in\mathfrak S_1),
\begin{equation}
 \left|\det_{\!F}(I+A)-\det_{\!F}(I+B)\right|
 \le
 \|A-B\|_1
 \exp\bigl(1+\|A\|_1+\|B\|_1\bigr).
 \label{eq:V143-det-continuity}
\end{equation}
\end{lemma}

\begin{proof}
The exterior-power expansion
\(
\det_{\!F}(I+A)=\sum_{k\ge0}\operatorname{tr}(\wedge^kA)
\)
is absolutely convergent in trace norm.  Multilinearity of the (k)-fold
exterior product replaces one factor at a time and gives
\[
 \|\wedge^kA-\wedge^kB\|_1
 \le
 \frac{\|A-B\|_1}{(k-1)!}
 (\|A\|_1+\|B\|_1)^{k-1}.
\]
Summing and using the harmless standard factor (e) at (k=0,1) yields
\eqref{eq:V143-det-continuity}.
\end{proof}

\begin{theorem}[Direct-tail phase theorem]
\label{thm:V143-direct-tail}
Assume
\begin{equation}
 K_q\longrightarrow K_\infty
 \quad\text{in }\mathfrak S_1,
 \qquad
 T_\infty:=I+K_\infty\text{ is invertible}.
 \label{eq:V143-trace-limit}
\end{equation}
Then, for every base cutoff (q_0),
\begin{equation}
 \prod_{n=0}^{N-1}r_{q_{n+1}q_n}
 =\operatorname{ph}\det_{\!F}(T_{q_N}T_{q_0}^{-1})
 \longrightarrow
 \operatorname{ph}\det_{\!F}(T_\infty T_{q_0}^{-1}).
 \label{eq:V143-direct-telescope}
\end{equation}
In particular, convergence follows from one direct trace-norm tail and
does not require \(
\sum_n|1-r_{q_{n+1}q_n}|<\infty
\).
\end{theorem}

\begin{proof}
The group (I+\mathfrak S_1) is closed under products and inverses.
Multiplicativity of the Fredholm determinant gives exact cancellation of
all intermediate factors:
\[
 \prod_{n=0}^{N-1}
 \det_{\!F}(T_{q_{n+1}}T_{q_n}^{-1})
 =\det_{\!F}(T_{q_N}T_{q_0}^{-1}).
\]
Taking unit phases gives the equality in
\eqref{eq:V143-direct-telescope}.  Trace-norm convergence of (T_{q_N}),
boundedness of the fixed inverse (T_{q_0}^{-1}), and
Lemma~\ref{lem:V143-det-continuity} give convergence of the nonzero
determinants.  The phase map is continuous away from zero, and
invertibility of (T_\infty T_{q_0}^{-1}) makes its determinant nonzero.
\end{proof}

\begin{remark}[Relation to the V140 first moment]
\label{rem:V143-first-moment}
Theorem~\ref{thm:V140-l1-phase} controls a product through absolute
adjacent variation.  Theorem~\ref{thm:V143-direct-tail} instead computes
that product as one direct comparison.  It is an exact signed telescope of
operators.  The new burden is to prove \eqref{eq:V143-trace-limit} from
renewal estimates already available; assuming it as a definition would be
as circular as the norm warned against in
Firewall~\ref{fire:V140-first-moment-circular}.
\end{remark}

\subsection{Schatten membership from boundary order}
\label{subsec:V143-Schatten}

Let (Y) be a closed (d)-dimensional manifold and
(K\in\Psi^{-m}(Y;E)), (m>0), a classical pseudodifferential operator.

\begin{theorem}[Order-to-Schatten bound]
\label{thm:V143-order-Schatten}
One has
\begin{equation}
 K\in\mathfrak S_p
 \qquad\text{for every }p>\frac d m.
 \label{eq:V143-Schatten-threshold}
\end{equation}
The endpoint is false in general.  In particular, on a
three-dimensional boundary,
\begin{equation}
 \Psi^{-1}\subset\bigcap_{p>3}\mathfrak S_p,
 \qquad
 \Psi^{-4}\subset\mathfrak S_1,
 \label{eq:V143-three-dimensional-threshold}
\end{equation}
while an order-(-1) operator need not be trace class.
\end{theorem}

\begin{proof}
Let (\Delta_Y) be a positive Laplace-type operator.  The operator
\[
 B:=K(I+\Delta_Y)^{m/2}
\]
has order zero and is bounded on (L^2).  Hence the ideal inequality for
singular values gives
\[
 s_j(K)\le\|B\|,s_j((I+\Delta_Y)^{-m/2}).
\]
Weyl's law says that the latter singular values are
(O(j^{-m/d})).  Thus (\sum_js_j(K)^p<\infty) when (mp/d>1), proving
\eqref{eq:V143-Schatten-threshold}.  Sharpness is seen on the flat torus:
the Fourier multiplier ((I+\Delta_Y)^{-m/2}) has singular values
asymptotic to a nonzero multiple of (j^{-m/d}).  At (p=d/m) its
(p)-power sum is harmonic and diverges.  Taking (d=3), (m=1) and
(m=4) gives \eqref{eq:V143-three-dimensional-threshold}.
\end{proof}

\subsection{Application to the renewal boundary comparison}
\label{subsec:V143-renewal-application}

After the common principal polarization and universal local factor have
been removed, two Calderon projectors for first-order Dirac operators with
the same principal symbol differ generically by an operator of order
(-1).  Accordingly write the transported high-energy comparison as
\begin{equation}
 T_q=I+K_q,
 \qquad K_q\in\Psi^{-1}(\Sigma;E_\Sigma).
 \label{eq:V143-raw-comparison-order}
\end{equation}
This is consistent with the explicit order-(-1) mean-curvature row of the
normalized interface logarithm found in V125.  The exponentially small
finite-length frozen remainder of
Theorem~\ref{thm:V124-trace-class} is smoothing; it does not remove the
geometric order-(-1) row.

\begin{corollary}[The raw comparison misses trace class by three rows]
\label{cor:V143-trace-deficit}
For (\dim\Sigma=3), \eqref{eq:V143-raw-comparison-order} guarantees
(K_q\in\mathfrak S_4), but does not guarantee
(K_q\in\mathfrak S_1).  Therefore the ordinary Fredholm determinant in
Theorem~\ref{thm:V143-direct-tail} is not yet defined by the available
symbol order.  Subtracting the complete homogeneous orders
\begin{equation}
 -1,\quad-2,\quad-3
 \label{eq:V143-three-local-rows}
\end{equation}
leaves order at most (-4), which is trace class.
\end{corollary}

\begin{proof}
Apply Theorem~\ref{thm:V143-order-Schatten} with (d=3).  Removing the
three displayed classical symbol components lowers the remainder from
order (-1) to order (-4).  The second inclusion in
\eqref{eq:V143-three-dimensional-threshold} then applies.
\end{proof}

\begin{firewall}[A modified determinant has a cocycle anomaly]
\label{fire:V143-det4-anomaly}
Membership in (\mathfrak S_4) permits the modified determinant
\begin{equation}
 \det{}_4(I+K)
 :=\det_{\!F}\!\left((I+K)
 \exp\left[-K+\frac12K^2-\frac13K^3\right]\right).
 \label{eq:V143-det4}
\end{equation}
The operator inside the Fredholm determinant differs from (I) by a
trace-class operator.  However, (\det{}_4) is not multiplicative without
a finite trace-polynomial correction.  Using it alone would reintroduce a
nontrivial triangle defect into the exact V142 system.  The local rows in
\eqref{eq:V143-three-local-rows} must be transgressed with the same
three-cutoff convention before a renormalized phase comparison is declared.
\end{firewall}

\begin{assessment}[Analytic bottleneck after V143]
\label{ass:V143-analytic-bottleneck}
The desired signed telescope now has an exact sufficient form:
trace-norm convergence of one direct tail.  The existing geometric
comparison supplies only order (-1), hence (\mathfrak S_4), on the
three-dimensional boundary.  The obstruction is ultraviolet and local,
not a remaining global anomaly.  The next step is to compute and subtract
the complete orders (-1,-2,-3), and to prove that their counterterm
transgression cancels the modified-determinant multiplicative anomaly.
\end{assessment}

\subsection{V143 ledger and next acquisition}
\label{subsec:V143-ledger}

\paragraph{New exact conclusions.}
\begin{enumerate}[label=\textup{(\arabic*)},leftmargin=2.8em]
\item Proposition~\ref{prop:V143-algebraic-line} constructs the algebraic
      direct-limit phase line with no convergence assumption.
\item Theorem~\ref{thm:V143-direct-tail} turns all adjacent phases into one
      exact trace-class determinant telescope.
\item Theorem~\ref{thm:V143-order-Schatten} proves the precise
      order-to-Schatten threshold from Weyl's law.
\item Corollary~\ref{cor:V143-trace-deficit} proves that the raw
      three-dimensional comparison is only (\mathfrak S_4) and identifies
      the three local rows which must be removed to reach trace class.
\item Firewall~\ref{fire:V143-det4-anomaly} prevents a bare modified
      determinant from breaking the exact comparison cocycle.
\end{enumerate}

\paragraph{Next forced step.}
V144 should construct the renormalized relative phase from
(\det{}_4) plus an explicit finite trace-polynomial transgression.  It
must prove the three-cutoff multiplication law at a finite regulator before
removing that regulator.  It should then identify which parts of the
orders (-1,-2,-3) were already computed in V125--V128 and isolate the
missing chiral symbol coefficients rather than reuse bosonic ones.

\begin{assessment}[V143 continuum status]
\label{ass:V143-continuum-status}
The full Standard--Model--Einstein continuum remains unproved.  The
finite chiral phase system and its algebraic limit are exact, and a
non-adjacent analytic convergence criterion is now available.  Its
trace-class hypothesis, the chain and family-index gates, the extended
BRST boundary complex, microscopic gravitational normalization, nonlinear
existence, Einstein stress closure, and interacting cofinal limit remain
open.
\end{assessment}

\section*{V143 exact checks and append-only record}
\addcontentsline{toc}{section}{V143 exact checks and append-only record}

The direct-tail and Schatten checksum is
\begin{equation}
 \boxed{
 \prod_{n<N}r_{q_{n+1}q_n}
 =\operatorname{ph}\det_{\!F}(T_{q_N}T_{q_0}^{-1}),
 \qquad
 \Psi^{-1}(\Sigma^3)\subset\mathfrak S_p\ (p>3),
 \qquad
 \Psi^{-4}(\Sigma^3)\subset\mathfrak S_1.}
 \label{eq:V143-final-check}
\end{equation}
The append operation preserves the complete V142 prefix byte for byte before
its former live terminator.  The assembled V143 file is checked for one live
final terminator, balanced environments and braces, duplicate labels, newly
introduced unresolved references, display math delimiters, and control
characters before the exact V142 predecessor is removed.  The companion
files are not modified.

% =============================================================================
% END APPENDED VERSION V143 MATERIAL
% =============================================================================

% =============================================================================
% BEGIN APPENDED VERSION V144 MATERIAL
% =============================================================================

\section{V144: the fourth modified determinant, local transgression, and
an exact renormalized comparison cocycle}
\label{sec:V144}

\subsection{Purpose}
\label{subsec:V144-purpose}

V143 proves that the raw chiral boundary comparison is naturally an
(I+\mathfrak S_4) operator, not an (I+\mathfrak S_1) operator.  The
fourth modified determinant is therefore the first determinant defined by
the available order.  It is not multiplicative by itself.  V144 derives
its correction from an exact finite-regulator identity and proves that the
finite-part correction restores the V142 three-cutoff cocycle.

The construction is deliberately made before computing any new symbol
coefficient.  It shows exactly which local data are required: the complete
orders (-1,-2,-3) of the first-order chiral comparison.  The bosonic
gauge--ghost symbols of V125--V128 cannot be substituted for them.

\subsection{The exact finite identity}
\label{subsec:V144-finite-identity}

For a finite matrix (K) with (I+K) invertible, define
\begin{align}
 c_3(K)
 &:=\operatorname{tr}
 \left(K-\frac12K^2+\frac13K^3\right),
 \label{eq:V144-c3}\
 \det{}_4(I+K)
 &:=\det\left((I+K)
 \exp\left[-K+\frac12K^2-\frac13K^3\right]\right).
 \label{eq:V144-finite-det4}
\end{align}

\begin{lemma}[Exact determinant decomposition]
\label{lem:V144-exact-decomposition}
At every finite regulator,
\begin{equation}
 \boxed{
 \det(I+K)=\det{}_4(I+K)\exp c_3(K).}
 \label{eq:V144-exact-decomposition}
\end{equation}
\end{lemma}

\begin{proof}
Take the determinant of the exponential in
\eqref{eq:V144-finite-det4} and use
(\det(e^A)=e^{\operatorname{tr}A}).  Its exponent is
(-c_3(K)), which gives \eqref{eq:V144-exact-decomposition}.
\end{proof}

Let (T_{10,N}=I+K_{10,N}) and
(T_{21,N}=I+K_{21,N}) be two composable finite comparison matrices in
one common regulator, and require the direct matrix to be
\begin{equation}
 T_{20,N}:=T_{21,N}T_{10,N},
 \qquad
 K_{20,N}=K_{21,N}+K_{10,N}+K_{21,N}K_{10,N}.
 \label{eq:V144-composed-K}
\end{equation}

\begin{theorem}[Finite modified-determinant anomaly identity]
\label{thm:V144-finite-anomaly}
The failure of (\det{}_4) to multiply is exactly
\begin{equation}
 \boxed{
 \frac{\det{}_4(T_{20,N})}
 {\det{}_4(T_{21,N})\det{}_4(T_{10,N})}
 =
 \exp\left[
 c_3(K_{21,N})+c_3(K_{10,N})-c_3(K_{20,N})
 \right].}
 \label{eq:V144-finite-anomaly}
\end{equation}
Thus the corrected finite determinant
\begin{equation}
 \mathfrak D_N(T):=\det{}_4(T)\exp c_3(T-I)
 \label{eq:V144-corrected-finite}
\end{equation}
is exactly multiplicative.
\end{theorem}

\begin{proof}
Apply Lemma~\ref{lem:V144-exact-decomposition} to the three matrices and
use the ordinary identity
(\det(T_{21,N}T_{10,N})=\det T_{21,N}\det T_{10,N}).
Cancel the three nonzero ordinary determinants.  This proves
\eqref{eq:V144-finite-anomaly}, and substitution into
\eqref{eq:V144-corrected-finite} proves multiplicativity.
\end{proof}

\begin{firewall}[The common-regulator product is essential]
\label{fire:V144-compression-product}
If (P_N) is a projection on an infinite boundary space, then generally
\[
 P_NT_{21}T_{10}P_N
 \ne(P_NT_{21}P_N)(P_NT_{10}P_N).
\]
Theorem~\ref{thm:V144-finite-anomaly} applies to the right-hand side,
which is the direct map in the declared finite category.  Replacing it by
an independently compressed direct operator creates a corner leakage term
and can create a spurious triangle defect.
\end{firewall}

\subsection{Why precisely three local rows occur}
\label{subsec:V144-three-rows}

Let (K\in\Psi^{-1}(\Sigma^3;E)) be classical and let (P_N) be a common
covariant spectral cutoff.  The three terms in (c_3(P_NKP_N)) have
orders (-1,-2,-3), up to the lower symbol terms created by composition.

\begin{proposition}[Local divergence ledger]
\label{prop:V144-local-ledger}
The regulated trace (c_3(K_N)), (K_N=P_NKP_N), has a polyhomogeneous
asymptotic expansion of the form
\begin{equation}
 c_3(K_N)
 =a_2(K)\Lambda_N^2+a_1(K)\Lambda_N
 +a_{\log}(K)\log\Lambda_N+c_0(K)+o(1),
 \label{eq:V144-c3-expansion}
\end{equation}
with the understood omission of coefficients forbidden by parity or
symbol cancellation.  The coefficients (a_2,a_1,a_{\log}) are integrals
of local densities determined by the complete symbol through orders
(-1,-2,-3).  Every contribution of order at most (-4) is trace class
and affects only (c_0) and the vanishing tail.
\end{proposition}

\begin{proof}
In local boundary covariables, the trace of a homogeneous symbol of order
(-j) has radial factor
\[
 \int_1^{\Lambda_N}\rho^{2-j}\,d\rho.
\]
It is quadratic for (j=1), linear for (j=2), logarithmic for (j=3),
and convergent for (j>3).  Pseudodifferential composition expresses the
symbols of (K^2) and (K^3) as universal finite sums of products and
derivatives of the symbol of (K).  Their angular integrals and fibre
traces are local densities.  A partition of unity gives the global
statement.  The trace-class assertion is
Theorem~\ref{thm:V143-order-Schatten} with (m=4).
\end{proof}

Define the common-cutoff finite part by
\begin{equation}
 c_3^{\rm fp}(K)
 :=\operatorname*{FP}_{N\to\infty}c_3(K_N):=c_0(K).
 \label{eq:V144-c3-fp}
\end{equation}
This definition includes the regulator convention; changing the covariant
cutoff changes (c_3^{\rm fp}) by a local functional of the same three
symbol rows.

\subsection{Removal of the multiplicative anomaly}
\label{subsec:V144-anomaly-removal}

Suppose (K_{ij,N}\to K_{ij}) in (\mathfrak S_4), the finite matrices
obey \eqref{eq:V144-composed-K}, and the expansions
\eqref{eq:V144-c3-expansion} exist for the three comparisons in the same
cutoff convention.  Define
\begin{equation}
 \boxed{
 \det{}_{\rm ren}T_{ij}
 :=\det{}_4(I+K_{ij})\exp c_3^{\rm fp}(K_{ij}).}
 \label{eq:V144-ren-det}
\end{equation}

\begin{theorem}[Exact renormalized multiplication law]
\label{thm:V144-ren-multiplication}
Under the preceding hypotheses,
\begin{equation}
 \boxed{
 \det{}_{\rm ren}T_{20}
 =\det{}_{\rm ren}T_{21}\det{}_{\rm ren}T_{10}.}
 \label{eq:V144-ren-multiplication}
\end{equation}
Consequently the unit phases of these determinants have triangle defect
one and are compatible with the determinant-line cocycle
\eqref{eq:V142-exact-cocycle}.
\end{theorem}

\begin{proof}
The map
\begin{equation}
 \mathcal R_4(K)
 :=(I+K)\exp\left[-K+\frac12K^2-\frac13K^3\right]-I
 \label{eq:V144-R4}
\end{equation}
is continuous from (\mathfrak S_4) to (\mathfrak S_1).  Indeed, its
power series starts in total degree four, and Holder's inequality for
Schatten ideals places every product of four (\mathfrak S_4) factors in
(\mathfrak S_1).  Hence the three finite modified determinants converge.

Take a continuous logarithm along the finite comparison paths, or
equivalently retain the exponential identity
\eqref{eq:V144-finite-anomaly}.  Its left side has a finite nonzero limit.
Therefore the divergent coefficients on the right cancel between
(c_3(K_{21,N})+c_3(K_{10,N})-c_3(K_{20,N})).  Taking the common finite
part in the logarithmic identity gives
\begin{align*}
 &\log\det{}_4(T_{20})-log\det{}_4(T_{21})
 -\log\det{}_4(T_{10})\\
 &\qquad=
 c_3^{\rm fp}(K_{21})+c_3^{\rm fp}(K_{10})
 -c_3^{\rm fp}(K_{20})
 \quad\pmod{2\pi i\mathbb Z}.
\end{align*}
Exponentiating and rearranging proves
\eqref{eq:V144-ren-multiplication}.  Taking unit phases preserves the
identity.
\end{proof}

\begin{corollary}[Renormalized direct-tail criterion]
\label{cor:V144-ren-direct-tail}
Let (T_q=I+K_q) be direct comparisons in one fixed continuum frame.  If
\begin{equation}
 K_q\longrightarrow K_\infty\quad\text{in }\mathfrak S_4,
 \qquad
 c_3^{\rm fp}(K_q)\longrightarrow c_\infty,
 \label{eq:V144-ren-tail-hypotheses}
\end{equation}
and (I+K_\infty) is invertible, then
\begin{equation}
 \operatorname{ph}\det{}_{\rm ren}T_q
 \longrightarrow
 \operatorname{ph}\left[
 \det{}_4(I+K_\infty)e^{c_\infty}\right].
 \label{eq:V144-ren-tail-limit}
\end{equation}
By Theorem~\ref{thm:V144-ren-multiplication}, the adjacent product equals
the direct renormalized comparison at every finite stage.
\end{corollary}

\begin{proof}
Continuity of \(
\mathcal R_4:\mathfrak S_4\to\mathfrak S_1
\) and Fredholm determinant continuity imply
\(
\det{}_4(I+K_q)\to\det{}_4(I+K_\infty)
\).
The second convergence in \eqref{eq:V144-ren-tail-hypotheses} handles the
local factor.  Invertibility keeps the limit away from zero.  The product
statement is repeated application of
Theorem~\ref{thm:V144-ren-multiplication}.
\end{proof}

\subsection{Separation from the existing bosonic symbol calculation}
\label{subsec:V144-sector-separation}

The V125--V128 calculations concern
\begin{equation}
 \log\mathcal W_{0,R}
 -\frac12\operatorname{tr}\log\mathcal W_{1,R},
 \label{eq:V144-bosonic-object}
\end{equation}
the second-order gauge--ghost modulus obtained from Dirichlet-to-Neumann
Schur forms.  The present (K) is the first-order chiral polarization
comparison entering (\widehat\sigma_{q'/q}).  Its symbol contains spin,
chirality, the Standard Model representation, and the selected spectral
section.

\begin{proposition}[No transfer of the three coefficients by squaring]
\label{prop:V144-no-bosonic-transfer}
Equation \eqref{eq:V142-modulus-square} determines only the real part of
the logarithm of the chiral determinant.  It does not determine the
imaginary finite parts of
\begin{equation}
 \operatorname{tr}K_N,
 \qquad
 \operatorname{tr}K_N^2,
 \qquad
 \operatorname{tr}K_N^3,
 \label{eq:V144-chiral-three-traces}
\end{equation}
and therefore does not determine the local phase transgression in
\eqref{eq:V144-ren-det}.
\end{proposition}

\begin{proof}
Replacing a chiral determinant (z) by the determinant of its positive
square replaces (z) by (|z|^2).  The latter is unchanged under
(z\mapsto e^{i\theta}z).  The imaginary parts of
\eqref{eq:V144-chiral-three-traces} generate precisely such phase changes
in \eqref{eq:V144-exact-decomposition}.  Hence no calculation of the
positive square can recover them.
\end{proof}

\begin{assessment}[Upstream symbol acquisition]
\label{ass:V144-symbol-bottleneck}
The modified-determinant cocycle is now exact once the common finite parts
exist.  The missing analytic input is sharply localized: compute the
orders (-1,-2,-3) of the first-order chiral Calderon comparison in the
same spin, gauge, Higgs, and spectral-section frame, and prove convergence
of their finite parts together with (\mathfrak S_4) convergence of the
remainder.  The bosonic interface rows already in the manuscript do not
pay this chiral phase debit.
\end{assessment}

\subsection{V144 ledger and next acquisition}
\label{subsec:V144-ledger}

\paragraph{New exact conclusions.}
\begin{enumerate}[label=\textup{(\arabic*)},leftmargin=2.8em]
\item Lemma~\ref{lem:V144-exact-decomposition} splits every finite
      determinant into its fourth modified determinant and three explicit
      traces.
\item Theorem~\ref{thm:V144-finite-anomaly} computes the entire
      multiplicative anomaly at the common finite regulator.
\item Proposition~\ref{prop:V144-local-ledger} proves that exactly the
      orders (-1,-2,-3) contribute ultraviolet counterterms on
      (\Sigma^3).
\item Theorem~\ref{thm:V144-ren-multiplication} proves that their common
      finite-part transgression restores the exact comparison cocycle.
\item Corollary~\ref{cor:V144-ren-direct-tail} lowers the direct analytic
      convergence target from trace norm to (\mathfrak S_4) plus
      convergence of one local finite part.
\item Proposition~\ref{prop:V144-no-bosonic-transfer} proves why the
      existing positive gauge--ghost symbol does not supply the chiral
      phase counterterm.
\end{enumerate}

\paragraph{Next forced step.}
V145 is the next compulsory NS/YM audit.  After that audit it should build
the first-order chiral Calderon comparison in a fixed product-collar model
and compute at least its principal cancellation and order-(-1) phase
row.  If the companion programmes reveal a stronger exact difference
identity, it should replace separate coefficient estimates by that direct
identity.

\begin{assessment}[V144 continuum status]
\label{ass:V144-continuum-status}
The full Standard--Model--Einstein continuum remains unproved.  V144
constructs a multiplicative renormalized phase determinant at the exact
finite and conditional infinite levels.  The chiral symbol inputs, chain
and family-index gates, local finite-part convergence, extended BRST
boundary complex, microscopic gravitational normalization, nonlinear
existence, Einstein stress closure, and interacting cofinal limit remain
open.
\end{assessment}

\section*{V144 exact checks and append-only record}
\addcontentsline{toc}{section}{V144 exact checks and append-only record}

The renormalized-cocycle checksum is
\begin{equation}
 \boxed{
 \det T=\det{}_4(T)e^{c_3(T-I)},
 \qquad
 \det{}_{\rm ren}(T_{21}T_{10})
 =\det{}_{\rm ren}(T_{21})\det{}_{\rm ren}(T_{10}).}
 \label{eq:V144-final-check}
\end{equation}
The append operation preserves the complete V143 prefix byte for byte before
its former live terminator.  The assembled V144 file is checked for one live
final terminator, balanced environments and braces, duplicate labels, newly
introduced unresolved references, display math delimiters, and control
characters before the exact V143 predecessor is removed.  The companion
files are not modified.

% =============================================================================
% END APPENDED VERSION V144 MATERIAL
% =============================================================================

% =============================================================================
% BEGIN APPENDED VERSION V145 MATERIAL
% =============================================================================

\section{V145: compulsory companion audit, boundary-family cobordism,
and intrinsic cofinal phase closure}
\label{sec:V145}

\subsection{Purpose}
\label{subsec:V145-purpose}

V144 formulated the phase in a fixed pseudodifferential frame.  The
compulsory companion audit suggests going one level upstream: retain the
complete determinant line over the joint background--cutoff family before
compressing it to eigenchannels or scalar representatives.  At that level,
the tangential operator is a genuine boundary family.  Family cobordism
forces its odd (K)-theory index to vanish, so spectral sections exist on
every compact cutoff slab.  The Bismut--Freed connection is flat by V138
and has trivial holonomy by V141.  Its parallel transport therefore gives
the renewal comparison maps and their exact cocycle directly.

This closes the intrinsic faithful-branch chiral phase system.  V143--V144
remain the correct framework if a particular nonparallel scalar symbol
frame is required, but their local finite parts are not a new obstruction
to the abstract determinant-line phase.

\subsection{Mandatory V145 companion audit}
\label{subsec:V145-audit}

The stable read-only shared-folder snapshot at (09{:}59{:}23) Pacific
time on 5 September 2026 was
\begin{center}
\begin{tabular}{|p{0.49\textwidth}|r|p{0.31\textwidth}|}
\hline
Companion file & Bytes & SHA--256\\
\hline
\texttt{Renewal\_NS\_Stability\_Research\_Notes\_V31.tex}
& (887537) &
\texttt{F1121B62238AD5491BB7CF03635EA9934942EDF573ED8FAAA9EE79E1D38A6696}\\
\hline
\texttt{Renewal\_Yang\_Mills\_Mass\_Gap\_Research\_Notes\_V125.tex}
& (2906323) &
\texttt{B6E14D0813553508E323603D05BE1F44199D5E55557ABA9708231F607E256C28}\\
\hline
\end{tabular}
\end{center}

The NS note remains byte-identical to the V140 audit.  Its V31 conclusion
still forbids replacing the restricted formation correlation by the full
continuation-class stock.  The relevant instruction here is to avoid
inventing a phase norm whose value is the phase convergence one is trying
to prove.

YM has advanced from V122 to V125.  V125 writes the difference of two
Wilson scalars on an arbitrary tensor-product output as the compression of
one complete operator difference.  A single averaged square then gives a
square-root debit uniformly in the output label, lowering depth, copy
multiplicity, and number of coordinates.  Orthogonal outer projections are
aggregated before a norm is taken.  This closes its bounded-(a)-heat
off-cut branch; the hot-(a) scale mismatch remains open.

\begin{assessment}[V145 audit adjustment]
\label{ass:V145-audit-adjustment}
For the SM phase, the analogue of the complete YM compression is the full
determinant line with its connection.  Kernel crossings, spectral-section
copies, and boundary modes should remain inside that line until parallel
transport is taken.  Estimating phases mode by mode would introduce a
fictitious multiplicity problem.  The correct upstream questions are
whether the boundary-family (K^1) index vanishes and whether the complete
connection has nontrivial holonomy.
\end{assessment}

\subsection{The joint background--cutoff family}
\label{subsec:V145-joint-family}

Let (\mathcal B_0) be a compact connected chart of faithful
(G_6)-backgrounds.  For (Q>q_0), set
\begin{equation}
 \mathcal S_Q:=\mathcal B_0\times[q_0,Q].
 \label{eq:V145-cutoff-slab}
\end{equation}
Assume the renewal collar, metric, gauge connection, and Higgs--Yukawa data
vary smoothly on this slab.  Let (D_{b,q}^+) be the bulk chiral family and
(A_{b,q}=A_{b,q}^*) its induced tangential family on the closed
three-dimensional boundary.

\begin{theorem}[Boundary-family index vanishing]
\label{thm:V145-boundary-index-zero}
On every compact cutoff slab,
\begin{equation}
 \boxed{
 \operatorname{Ind}(A) =0
 \quad\text{in }K^1(\mathcal S_Q).}
 \label{eq:V145-boundary-index-zero}
\end{equation}
Consequently the tangential family admits a spectral section on
(\mathcal S_Q).
\end{theorem}

\begin{proof}
The family (A_{b,q}) is not an arbitrary family of odd Dirac operators:
it is the oriented boundary family of the even-dimensional vertical Dirac
symbol of (D_{b,q}^+).  Cobordism invariance of the family index says that
the analytic (K^1) index of such a boundary family is zero.  One way to
see the mechanism is to place the vertical symbol in the relative
(K)-group of the fibrewise cotangent ball and its boundary.  The symbol
of (A) is its boundary class.  The analytic-index map annihilates this
boundary because the two boundary faces of the doubled vertical family
occur with opposite orientations.  Hence
\eqref{eq:V145-boundary-index-zero}.

The spectral-section existence theorem
\ref{thm:V142-spectral-section-gate} now applies.  These two statements are
the cobordism and spectral-section theorems of Melrose--Piazza,
\emph{Families of Dirac operators, boundaries and the b-calculus}.
\end{proof}

\begin{remark}[Compact slabs rather than a hidden uniform gap]
\label{rem:V145-compact-slabs}
Theorem~\ref{thm:V145-boundary-index-zero} does not assert a uniform
positive spectral gap for all (q\to\infty).  It produces a smooth
spectral section on each compact slab.  On overlapping slabs, two choices
differ by the finite relative line of
Proposition~\ref{prop:V139-section-change}; its exact cocycle patches the
determinant lines over the increasing union.
\end{remark}

\subsection{The complete determinant connection}
\label{subsec:V145-connection}

Use the spectral sections of
Theorem~\ref{thm:V145-boundary-index-zero} and their relative-line
identifications to form the determinant line
\begin{equation}
 \mathcal L\longrightarrow
 \mathcal S:=\mathcal B_0\times[q_0,\infty).
 \label{eq:V145-joint-line}
\end{equation}
At a renewal boundary, its Bismut--Freed connection is paired with the
canonical Dai--Freed vectors of V141.

\begin{theorem}[Flatness and trivial holonomy on the joint family]
\label{thm:V145-joint-flatness}
For the complete three-generation Standard Model representation on the
faithful (G_6) branch,
\begin{equation}
 F_{\nabla^{\rm BF}}=0
 \label{eq:V145-curvature-zero}
\end{equation}
on (\mathcal S), and the paired holonomy is one around every loop in
(\mathcal S).
\end{theorem}

\begin{proof}
The curvature of the Bismut--Freed determinant connection is the degree-two
component on the parameter space of the fibre integral of the six-form
anomaly polynomial.  V138 proves that this polynomial vanishes for the
complete Standard Model representation.  The proof is algebraic in the
representation traces, so it also applies when one of the parameter
directions is the cutoff direction.  This gives
\eqref{eq:V145-curvature-zero}.

A loop in (\mathcal S) produces a five-dimensional faithful
spin--(G_6) mapping torus.  If the four-dimensional fibre is closed,
Theorem~\ref{thm:V141-closed-anomaly-zero} makes its eta holonomy one.  If
the fibre has the renewal boundary, the induced boundary mapping torus is
null-bordant and Corollary~\ref{cor:V141-renewal-mapping-phase} makes the
paired holonomy one.  These exhaust the loops relevant to the joint family.
\end{proof}

\begin{corollary}[A global parallel phase section]
\label{cor:V145-global-parallel-section}
On each connected component of (\mathcal S), the Hermitian line
(\mathcal L) has a nowhere-zero unit parallel section, unique up to one
constant element of (U(1)).
\end{corollary}

\begin{proof}
Fix a unit vector at one point and parallel transport it.  Flatness makes
transport invariant under endpoint-fixed homotopy, and trivial holonomy
makes it independent of the path.  This constructs the section.  The ratio
of two unit parallel sections is a constant unit scalar on a connected
base.
\end{proof}

\subsection{Geometric renewal comparisons}
\label{subsec:V145-geometric-comparisons}

For fixed (b\in\mathcal B_0) and (q\preceq q'), let
(\gamma_{q'q}(t)=(b,(1-t)q+tq')) be the canonical cutoff path.  Define
\begin{equation}
 \boxed{
 J^{\rm geom}_{q'q}
 :=\operatorname{PT}_{\gamma_{q'q}}^{\nabla^{\rm BF}}:
 \mathcal L_{b,q}\longrightarrow\mathcal L_{b,q'}.}
 \label{eq:V145-geometric-J}
\end{equation}

\begin{theorem}[Exact geometric comparison cocycle]
\label{thm:V145-geometric-cocycle}
For all (q\preceq q'\preceq q''),
\begin{equation}
 \boxed{
 J^{\rm geom}_{q''q'}J^{\rm geom}_{q'q}
 =J^{\rm geom}_{q''q}.}
 \label{eq:V145-geometric-cocycle}
\end{equation}
The map is unchanged if the canonical interval is replaced by any other
path in the same connected component with the same endpoints.
\end{theorem}

\begin{proof}
Parallel transport along a concatenated path is the composition of the two
parallel transports.  The concatenation of the \(q\)-to-\(q'\) and
\(q'\)-to-\(q''\) intervals and the direct \(q\)-to-\(q''\) interval differ
by a loop.  Theorem~\ref{thm:V145-joint-flatness} gives unit holonomy around
that loop, proving \eqref{eq:V145-geometric-cocycle}.  The same loop
argument proves path independence.
\end{proof}

\begin{proposition}[Agreement with the finite shell construction]
\label{prop:V145-shell-agreement}
Whenever the algebraic chain gate \eqref{eq:V142-chain-gate} holds for a
finite renewal comparison, the geometric map
(J^{\rm geom}_{q'q}) agrees with the shell map (J_{q'q}) of
\eqref{eq:V142-J}, provided both use the same spectral-section relative
line and Dai--Freed filling vector.
\end{proposition}

\begin{proof}
The determinant-line parallel transport theorem identifies transport
across a mapping cylinder with its exponentiated eta element.  Cutting the
cylinder at the intermediate interface gives the determinant exact
sequence \eqref{eq:V142-short-exact}; its high acyclic torsion and finite
crossing line are exactly the two factors in
\eqref{eq:V142-shell-vector}.  The Dai--Freed gluing theorem pairs the same
boundary vector.  Thus both constructions are the determinant-functor map
of the same bordism and agree.
\end{proof}

\begin{lemma}[Product-collar principal phase check]
\label{lem:V145-product-collar-phase}
On an exact product collar with
\begin{equation}
 \mathbb D=c(n)(\partial_r+A),
 \qquad A=A^*,
 \label{eq:V145-product-Dirac}
\end{equation}
the finite-regulator collar evolution is (E_R=e^{-RA}), and
\begin{equation}
 \operatorname{ph}\det E_R=1.
 \label{eq:V145-product-phase-one}
\end{equation}
Thus the principal product-collar comparison carries modulus but no local
phase.
\end{lemma}

\begin{proof}
The equation (\mathbb D\psi=0) is
(\partial_r\psi=-A\psi), so (E_R=e^{-RA}).  At a finite regulator its
eigenvalues are the positive numbers (e^{-R\lambda_j(A)}).  Therefore
\[
 \det E_R=\exp[-R\operatorname{tr}A]>0,
\]
which proves \eqref{eq:V145-product-phase-one}.
\end{proof}

\subsection{Intrinsic cofinal phase theorem}
\label{subsec:V145-phase-theorem}

\begin{theorem}[Faithful-branch intrinsic phase closure]
\label{thm:V145-intrinsic-phase-closure}
Let (s) be the unit parallel section of
Corollary~\ref{cor:V145-global-parallel-section} and write
(s_q:=s|_{\mathcal L_{b,q}}).  Then
\begin{equation}
 J^{\rm geom}_{q'q}s_q=s_{q'},
 \qquad
 r_{q'q}=1
 \label{eq:V145-ratio-one}
\end{equation}
for every pair of finite cutoffs.  The compatible family defines a unit
vector in the algebraic cofinal line
\eqref{eq:V143-algebraic-line}.  No adjacent phase first moment is required
in this parallel frame.
\end{theorem}

\begin{proof}
The defining property of a parallel section gives the first equality.
Using (s_q) as the local unit vectors in
\eqref{eq:V140-transition-ratio} gives (r_{q'q}=1).  Compatibility and
Proposition~\ref{prop:V143-algebraic-line} give the unit vector in the
direct-limit line.
\end{proof}

\begin{firewall}[What the phase theorem does not supply]
\label{fire:V145-phase-not-full-determinant}
Theorem~\ref{thm:V145-intrinsic-phase-closure} constructs the intrinsic
phase line and its compatible unit section.  It does not prove that the
cutoff Dirac operators converge to a Fredholm continuum operator, that
their Quillen moduli converge after renormalization, or that variations of
the full interacting effective action converge.  If one insists on a
fixed nonparallel pseudodifferential scalar frame, V144 gives the required
(\det{}_4) transgression.  Those are analytic modulus and realization
questions, not residual gauge anomalies.
\end{firewall}

\begin{assessment}[Phase-sector disposition]
\label{ass:V145-phase-disposition}
The family-index gate left open in V142 is closed on every compact cutoff
slab by boundary cobordism.  The comparison-map and phase first-moment
gates are bypassed intrinsically by complete determinant-line parallel
transport, whose curvature and holonomy both vanish on the faithful
branch.  Subject to the smooth-family hypothesis printed before
Theorem~\ref{thm:V145-boundary-index-zero}, the chiral phase sector is no
longer a bottleneck for the cofinal programme.
\end{assessment}

\subsection{V145 ledger and next acquisition}
\label{subsec:V145-ledger}

\paragraph{New exact conclusions.}
\begin{enumerate}[label=\textup{(\arabic*)},leftmargin=2.8em]
\item The mandatory audit records stable NS V31 and YM V125 snapshots and
      transfers the complete-operator strategy without importing an
      estimate.
\item Theorem~\ref{thm:V145-boundary-index-zero} closes the spectral-section
      (K^1) obstruction on every compact cutoff slab.
\item Theorem~\ref{thm:V145-joint-flatness} proves flatness and trivial
      paired holonomy on the joint background--cutoff family.
\item Theorem~\ref{thm:V145-geometric-cocycle} constructs the actual
      geometric renewal comparison maps and proves their exact cocycle.
\item Proposition~\ref{prop:V145-shell-agreement} identifies those maps
      with the V142 finite shell maps whenever the algebraic chain model is
      available.
\item Lemma~\ref{lem:V145-product-collar-phase} checks directly that the
      product-collar principal comparison has no phase.
\item Theorem~\ref{thm:V145-intrinsic-phase-closure} constructs the
      compatible cofinal unit phase without an adjacent first-moment sum.
\end{enumerate}

\paragraph{Next forced step.}
With the faithful chiral phase disposed of, V146 should return to the V137
off-shell BRST boundary problem.  It should include the
Nakanishi--Lautrup auxiliary field, derive the complete normal boundary
complex under the nilpotent BRST differential, and determine whether the
coupled normal-gauge/Higgs Robin block of V137 is preserved without using
bulk equations of motion.  If it fails, the missing boundary multiplier or
edge field must be introduced explicitly.

\begin{assessment}[V145 continuum status]
\label{ass:V145-continuum-status}
The full Standard--Model--Einstein continuum remains unproved.  The
faithful chiral phase and its spectral-section obstruction are now closed
intrinsically for smooth finite-cutoff families.  Convergence of the
renormalized modulus, the off-shell BRST boundary complex, microscopic
gravitational normalization, nonlinear existence, Einstein stress closure,
and the interacting cofinal limit remain open.
\end{assessment}

\section*{V145 exact checks and append-only record}
\addcontentsline{toc}{section}{V145 exact checks and append-only record}

The intrinsic-phase checksum is
\begin{equation}
 \boxed{
 \operatorname{Ind}(A)=0\in K^1(\mathcal S_Q),
 \qquad
 \operatorname{Hol}_{\gamma}(\mathcal L)=1,
 \qquad
 J^{\rm geom}_{q'q}s_q=s_{q'}.}
 \label{eq:V145-final-check}
\end{equation}
The append operation preserves the complete V144 prefix byte for byte before
its former live terminator.  The assembled V145 file is checked for one live
final terminator, balanced environments and braces, duplicate labels, newly
introduced unresolved references, display math delimiters, and control
characters before the exact V144 predecessor is removed.  The companion
files are not modified.

% =============================================================================
% END APPENDED VERSION V145 MATERIAL
% =============================================================================

% =============================================================================
% BEGIN APPENDED VERSION V146 MATERIAL
% =============================================================================

\section{V146: the off-shell BRST form domain, the natural Robin operator
domain, and the exact background-equation defect}
\label{sec:V146}

\subsection{Purpose}
\label{subsec:V146-purpose}

V137 found that the BRST variation of the normal gauge-fixing Robin row
contains a second normal ghost jet and left the extended off-shell boundary
complex open.  The obstruction comes from asking an off-shell differential
to preserve a natural Euler boundary equation pointwise.  These are two
different domains.  The BRST differential acts on the quadratic-form field
space.  The matrix Robin row is part of the operator domain associated with
that form after the Nakanishi--Lautrup field has been eliminated.

V146 constructs the form-domain complex and proves its off-shell closure.
It then differentiates the exact gauge Noether identity.  This yields the
precise obstruction to an operator-level ghost--longitudinal intertwiner on
a nonstationary background: a zeroth-order term proportional to the
classical background equations.  On a stationary background the defect
vanishes, and the relative ghost spectral core maps into the repaired V137
matrix Robin domain without imposing a second ghost boundary condition.

\subsection{The nilpotent quartet}
\label{subsec:V146-quartet}

Let
\begin{equation}
 R_\Phi c:=\binom{Dc}{-\tau_\Phi c}
 \label{eq:V146-gauge-generator}
\end{equation}
be the linearized gauge generator on the gauge--Higgs fluctuation
(u=(a,\varphi)).  In the linearized BRST complex set
\begin{equation}
 su=R_\Phi c,
 \qquad sc=0,
 \qquad s\bar c=b,
 \qquad sb=0,
 \label{eq:V146-BRST-quartet}
\end{equation}
where (\bar c) is the antighost and (b) is the bosonic
Nakanishi--Lautrup multiplier.  Background fields are held fixed in this
linearized complex.

\begin{lemma}[Off-shell nilpotency]
\label{lem:V146-nilpotency}
The differential \eqref{eq:V146-BRST-quartet} obeys (s^2=0) without any
field equation.
\end{lemma}

\begin{proof}
The fixed-background generator (R_\Phi) is linear and (sc=0), so
(s^2u=R_\Phi sc=0).  The remaining two identities are
(s^2\bar c=sb=0) and (s^2b=0).
\end{proof}

The background gauge functional of V137 is exactly the formal adjoint of
the gauge generator:
\begin{equation}
 F_\Phi u:=D^*a-\tau_\Phi^*\varphi
 =R_\Phi^*u.
 \label{eq:V146-F-adjoint}
\end{equation}
Thus the ghost operator is
\begin{equation}
 P_c:=F_\Phi R_\Phi,
 \label{eq:V146-Pc}
\end{equation}
including the lower-order background terms described in
Proposition~\ref{prop:V137-ghost-operator}.

For gauge parameter (\alpha>0), introduce the gauge fermion
\begin{equation}
 \Psi_\alpha
 :=\left\langle\bar c,
 F_\Phi u+\frac\alpha2b\right\rangle_{L^2(M)}.
 \label{eq:V146-gauge-fermion}
\end{equation}

\begin{proposition}[Exact off-shell gauge-fixing action]
\label{prop:V146-BRST-exact-action}
Up to the fixed Grassmann sign convention,
\begin{equation}
 s\Psi_\alpha
 =\langle b,F_\Phi u\rangle
 +\frac\alpha2\|b\|_2^2
 -\langle\bar c,P_cc\rangle.
 \label{eq:V146-gf-action}
\end{equation}
Moreover (s(s\Psi_\alpha)=0) identically.
\end{proposition}

\begin{proof}
Apply the graded Leibniz rule, use (s\bar c=b), (sb=0), and
(s(F_\Phi u)=F_\Phi R_\Phi c=P_cc).  Nilpotency follows from
Lemma~\ref{lem:V146-nilpotency}; no integration by parts and no ghost
equation is used.
\end{proof}

\subsection{The BRST form domain}
\label{subsec:V146-form-domain}

On the Robin Higgs branch define the smooth core
\begin{equation}
 \mathscr V_{\rm BRST}:=
 \left\{
 (a,\varphi,c,\bar c,b):
 \gamma_\top a=0,
 \ \gamma c=0
 \right\},
 \label{eq:V146-form-core}
\end{equation}
where (\varphi) has free boundary trace and (\bar c,b) are target
variables in the bulk (L^2) pairing.  The completed gauge--Higgs form
domain is
\begin{equation}
 \mathcal V_{AH}
 :=\{a\in H^1(T^*M\otimes\mathfrak g):\gamma_\top a=0\}
 \oplus H^1(\mathcal H_{\mathbb R}),
 \label{eq:V146-AH-form-domain}
\end{equation}
and the ghost source domain is (H^1_0(M;\mathfrak g)), restricted to a
smooth (H^2) core when (R_\Phi c) is required to lie in
\(mathcal V_{AH}).

\begin{theorem}[Off-shell boundary closure of the form domain]
\label{thm:V146-form-BRST-closure}
The differential \eqref{eq:V146-BRST-quartet} preserves the boundary
conditions in \eqref{eq:V146-form-core}.  In particular,
\begin{equation}
 \gamma c=0
 \quad\Longrightarrow\quad
 \gamma_\top(Dc)=D_\Sigma(\gamma c)=0.
 \label{eq:V146-tangential-closure}
\end{equation}
No condition on (D_nc) or (P_cc) at the boundary is required.
\end{theorem}

\begin{proof}
The only essential bosonic boundary row is
(\gamma_\top a=0).  Its BRST variation is
(D_\Sigma(\gamma c)), which vanishes for a Dirichlet ghost.  The Higgs
trace is unrestricted on the Robin form domain, so
(-\tau_\Phi c) is an allowed variation; in fact its trace vanishes.
The ghost trace is fixed and (sc=0).  The antighost and multiplier appear
as an (L^2) target doublet in \eqref{eq:V146-gf-action}; an (L^2)
field has no Sobolev boundary trace to constrain.  Finally
(s\bar c=b) and (sb=0) preserve that target space.
\end{proof}

\begin{remark}[Why no Dirichlet row is imposed on (b)]
\label{rem:V146-no-b-trace}
The multiplier has no derivative term in
\eqref{eq:V146-gf-action}; its natural space is (L^2(M)).  A boundary
condition (b|_\Sigma=0) is therefore not defined at this level.  Giving
(b) a Dirichlet trace would silently add regularity or a kinetic term and
would define a different gauge-fixing system.
\end{remark}

\begin{remark}[The antighost is a dual variable]
\label{rem:V146-antighost-dual}
The determinant (\det P_c) integrates (c\) over the Dirichlet operator
domain and (\bar c) over its (L^2) dual.  One may use self-adjointness to
identify source and target after the Gaussian is evaluated.  Imposing the
same Sobolev operator domain on (\bar c) before applying
(s\bar c=b) would confuse this Hilbert-space identification with the
off-shell BRST field space.
\end{remark}

\subsection{How the V137 Robin block is generated}
\label{subsec:V146-natural-domain}

The algebraic (b)-equation from \eqref{eq:V146-gf-action} is
\begin{equation}
 b=-\alpha^{-1}F_\Phi u.
 \label{eq:V146-b-elimination}
\end{equation}
Eliminating (b) produces the gauge-fixing form proportional to
(\|F_\Phi u\|^2).  In Feynman gauge, combining it with the invariant
Yang--Mills--Higgs quadratic form gives the Laplace-type matrix operator of
\eqref{eq:V137-matrix-Laplace}.

At (\gamma_\top a=0), the boundary trace of the gauge functional is
\begin{equation}
 \gamma F_\Phi u
 =-left[(D_n+H)a_n+\tau_\Phi^*\varphi\right].
 \label{eq:V146-F-boundary}
\end{equation}
The two natural Euler rows of the closed quadratic form are therefore
\begin{align}
 G(u)&:=(D_n+H)a_n+\tau_\Phi^*\varphi=0,
 \label{eq:V146-upper-natural}\
 R_H(u)&:=(D_n+S_H)\varphi+\tau_\Phi a_n=0.
 \label{eq:V146-lower-natural}
\end{align}
Together they are exactly the matrix Robin condition
\eqref{eq:V137-matrix-Robin}.

\begin{proposition}[Form domain versus operator domain]
\label{prop:V146-two-domains}
The essential rows in \eqref{eq:V146-form-core} define the off-shell
quadratic-form domain.  Equations
\eqref{eq:V146-upper-natural}--\eqref{eq:V146-lower-natural} define the
operator domain selected by the first variation of that closed form.  BRST
invariance requires preservation of the former.  It does not require
(sG(u)=0) for every off-shell Dirichlet ghost.
\end{proposition}

\begin{proof}
For a closed semibounded form, essential boundary data are imposed in its
form domain.  The associated operator consists of those form-domain
vectors for which the weak Euler functional is represented by an (L^2)
vector.  Integration by parts then forces the remaining boundary Green
form to vanish against arbitrary form-domain variations; this gives the
natural rows \eqref{eq:V146-upper-natural}--
\eqref{eq:V146-lower-natural}.  They are consequences of the operator
representation of the form, not additional restrictions on every
off-shell configuration.  Theorem~\ref{thm:V146-form-BRST-closure} proves
the required form-domain invariance.
\end{proof}

The calculation in V137 now has the precise interpretation
\begin{equation}
 sR_H(u)|_\Sigma=0,
 \qquad
 sG(u)|_\Sigma=-\gamma(P_cc),
 \label{eq:V146-natural-variations}
\end{equation}
where the sign follows from \eqref{eq:V146-F-boundary} and
(sF_\Phi u=P_cc).

\begin{corollary}[No overdetermined ghost row]
\label{cor:V146-no-second-ghost-row}
Equation \eqref{eq:V146-natural-variations} does not require imposing both
(c|_\Sigma=0) and (P_cc|_\Sigma=0) on the off-shell ghost.  The first is
the source boundary condition; the second appears only on the smooth
spectral core of the Dirichlet ghost Laplacian, where it follows from the
spectral equation.
\end{corollary}

\begin{proof}
The first statement is Proposition~\ref{prop:V146-two-domains}.  If
(P_cc=\lambda c) and (c|_\Sigma=0), then
((P_cc)|_\Sigma=\lambda c|_\Sigma=0), so the upper natural row is
preserved on each smooth ghost eigenspace without a second independent
boundary condition.
\end{proof}

\subsection{The exact nonstationary-background defect}
\label{subsec:V146-Noether-defect}

Let (S_0(X)) be the gauge-invariant classical Yang--Mills--Higgs action,
(X=(A,\Phi)), and let (R_X\epsilon) be its infinitesimal gauge action.
Gauge invariance is the Noether identity
\begin{equation}
 S_0'(X)[R_X\epsilon]=0.
 \label{eq:V146-Noether}
\end{equation}

\begin{theorem}[Differentiated Noether identity]
\label{thm:V146-differentiated-Noether}
Let (H_X=S_0''(X)) be the invariant Hessian and
(\mathcal E_X:=S_0'(X)) the background Euler covector.  Then
\begin{equation}
 \boxed{
 H_XR_X=\mathcal N_{\mathcal E_X},}
 \label{eq:V146-HR-defect}
\end{equation}
where (\mathcal N_{\mathcal E_X}) is the operator defined weakly by
\begin{equation}
 \langle v,\mathcal N_{\mathcal E_X}\epsilon\rangle
 :=-\mathcal E_X\!left[(D R)_X[v]\epsilon\right].
 \label{eq:V146-N-def}
\end{equation}
In particular, (\mathcal N_{\mathcal E_X}=0) on every classical
background.
\end{theorem}

\begin{proof}
Differentiate \eqref{eq:V146-Noether} in the arbitrary field direction
(v).  The derivative falling on (S_0') is
(S_0''(X)[v,R_X\epsilon]).  The derivative falling on the field-dependent
gauge generator is
(S_0'(X)[(D R)_X[v]\epsilon]).  Their sum is zero.  Move the second term
to the other side and use the (L^2) pairing to obtain
\eqref{eq:V146-HR-defect}--\eqref{eq:V146-N-def}.  If
(\mathcal E_X=0), the defining pairing is zero for every (v,epsilon).
\end{proof}

Let the Feynman-gauge Hessian be
\begin{equation}
 P_{AH}:=H_X+F_X^*F_X,
 \qquad P_c:=F_XR_X,
 \qquad F_X=R_X^*.
 \label{eq:V146-gauge-fixed-operators}
\end{equation}

\begin{corollary}[Exact ghost--longitudinal intertwining defect]
\label{cor:V146-intertwining-defect}
The complete quadratic operators satisfy
\begin{equation}
 \boxed{
 P_{AH}R_X-R_XP_c=\mathcal N_{\mathcal E_X}.}
 \label{eq:V146-intertwining-defect}
\end{equation}
Hence exact heat and determinant pairing holds on the longitudinal spectral
core of a stationary background, while a general off-shell background
carries the explicit Noether defect
(\mathcal N_{\mathcal E_X}).
\end{corollary}

\begin{proof}
Using \eqref{eq:V146-gauge-fixed-operators},
\[
 P_{AH}R_X
 =H_XR_X+F_X^*F_XR_X
 =\mathcal N_{\mathcal E_X}+R_XP_c.
\]
This proves the identity.  On a stationary background the defect vanishes.
Corollary~\ref{cor:V146-no-second-ghost-row} and
Lemma~\ref{lem:V137-ghost-jet-cancel} verify the two natural boundary rows
on the smooth ghost spectral core, so functional calculus preserves the
pairing there.
\end{proof}

\begin{firewall}[Quadratic off-shell determinants need the BV/source row]
\label{fire:V146-BV-needed}
The full nonlinear action (S_0+s\Psi_\alpha) is BRST invariant off shell.
Its quadratic truncation around a nonstationary background is not a closed
BRST complex by itself: the exact defect is
\(\mathcal N_{\mathcal E_X}\).  Deleting it by writing the background
equations would silently restrict an off-shell effective action to classical
solutions.  A genuinely off-shell determinant and stress variation must
retain the tadpole together with the BV antifield/source term whose
linearized Slavnov identity contains
\eqref{eq:V146-intertwining-defect}.
\end{firewall}

\begin{assessment}[Disposition of the V137 second-jet obstruction]
\label{ass:V146-second-jet-disposition}
The standard Nakanishi--Lautrup quartet is sufficient for nilpotent
off-shell BRST closure of the relative gauge--Higgs form domain.  No
ill-defined boundary trace of (b), no second Dirichlet ghost row, and no
new edge field is needed for that statement.  Operator-domain spectral
pairing is exact on stationary backgrounds.  On nonstationary backgrounds
the remaining obstruction is not the boundary jet; it is the explicit
bulk Noether defect \(
\mathcal N_{\mathcal E_X}
\), which belongs to the BV/source completion.
\end{assessment}

\subsection{V146 ledger and next acquisition}
\label{subsec:V146-ledger}

\paragraph{New exact conclusions.}
\begin{enumerate}[label=\textup{(\arabic*)},leftmargin=2.8em]
\item Proposition~\ref{prop:V146-BRST-exact-action} constructs the
      nilpotent gauge-fixing action with the Nakanishi--Lautrup field.
\item Theorem~\ref{thm:V146-form-BRST-closure} proves off-shell closure of
      the essential relative gauge--Higgs boundary data.
\item Proposition~\ref{prop:V146-two-domains} separates that form domain
      from the natural matrix Robin operator domain of V137.
\item Corollary~\ref{cor:V146-no-second-ghost-row} proves why no
      overdetermining second ghost boundary condition is needed.
\item Theorem~\ref{thm:V146-differentiated-Noether} derives the exact
      background-equation term from gauge invariance.
\item Corollary~\ref{cor:V146-intertwining-defect} identifies the precise
      obstruction to off-shell quadratic heat/determinant pairing and
      closes it on stationary backgrounds.
\end{enumerate}

\paragraph{Next forced step.}
V147 should build the minimal BV quadratic complex around a general
background.  It must add the antifield/source coupling dictated by the
differentiated Noether identity, prove the linearized Slavnov operator is
nilpotent on the boundary form domain, and show how the tadpole and
\(\mathcal N_{\mathcal E_X}\) cancel in the one-loop Ward identity.  Only
then may the off-shell Einstein stress variation use the V129--V137 heat
coefficients without a stationarity assumption.

\begin{assessment}[V146 continuum status]
\label{ass:V146-continuum-status}
The full Standard--Model--Einstein continuum remains unproved.  V146 closes
the off-shell BRST form-domain boundary complex and the stationary-background
operator pairing.  The nonstationary BV/source completion, renormalized
modulus, microscopic gravitational normalization, nonlinear existence,
Einstein stress closure, and interacting cofinal limit remain open.
\end{assessment}

\section*{V146 exact checks and append-only record}
\addcontentsline{toc}{section}{V146 exact checks and append-only record}

The BRST-domain and Noether-defect checksum is
\begin{equation}
 \boxed{
 s^2=0,
 \qquad
 \gamma c=0\Longrightarrow\gamma_\top(Dc)=0,
 \qquad
 P_{AH}R_X-R_XP_c=\mathcal N_{\mathcal E_X}.}
 \label{eq:V146-final-check}
\end{equation}
The append operation preserves the complete V145 prefix byte for byte before
its former live terminator.  The assembled V146 file is checked for one live
final terminator, balanced environments and braces, duplicate labels, newly
introduced unresolved references, display math delimiters, and control
characters before the exact V145 predecessor is removed.  The companion
files are not modified.

% =============================================================================
% END APPENDED VERSION V146 MATERIAL
% =============================================================================

% =============================================================================
% BEGIN APPENDED VERSION V147 MATERIAL
% =============================================================================

\section{V147: the BV source completion, the off-shell Slavnov identity,
and master-equation-preserving cutoff pushforward}
\label{sec:V147}

\subsection{Purpose}
\label{subsec:V147-purpose}

V146 proves that the relative gauge--Higgs form domain is BRST invariant
off shell and that the gauge-fixed quadratic operators obey
\[
 P_{AH}R_X-R_XP_c=\mathcal N_{\mathcal E_X}.
\]
The right side is not a boundary-condition failure.  It is the derivative
of the field-dependent gauge generator paired with the background Euler
covector.  The Batalin--Vilkovisky source term is the structure which
retains it in the Ward identity.

V147 constructs that source completion on the boundary form domain and
proves the exact off-shell Slavnov identity.  It also identifies the proper
finite-cutoff operation.  A naive Galerkin compression need not preserve a
nonlinear master equation; a BV fibre integral over a symplectic high-mode
pair does preserve it.  This gives a rigorous algebraic target for the
interacting renewal comparison maps.

\subsection{The minimal and nonminimal master action}
\label{subsec:V147-master-action}

Collect the classical gauge and Higgs fields into (X), let (c) be the
ghost, and introduce antifields (X^*,c^*,\bar c^*) of opposite parity.
The nonminimal fields are the antighost (\bar c) and multiplier (b).
For a closed gauge algebra, define
\begin{equation}
 S_{\rm BV}
 :=S_0(X)
 +\langle X^*,R_Xc\rangle
 -\frac12\langle c^*,[c,c]\rangle
 +\langle\bar c^*,b\rangle.
 \label{eq:V147-master-action}
\end{equation}
Antifields are placed in the continuous duals of the form domains from
\eqref{eq:V146-AH-form-domain}; the ghost has zero boundary trace.  The
canonical odd antibracket is
\begin{equation}
 (F,G)
 :=\sum_A\left(
 \frac{\delta_rF}{\delta\Phi^A}
 \frac{\delta_lG}{\delta\Phi_A^*}
 -
 \frac{\delta_rF}{\delta\Phi_A^*}
 \frac{\delta_lG}{\delta\Phi^A}
 \right),
 \label{eq:V147-antibracket}
\end{equation}
with the standard graded signs absorbed into the right/left derivative
notation.

\begin{theorem}[Classical master equation with the renewal boundary]
\label{thm:V147-CME}
On the V146 form-domain core,
\begin{equation}
 \boxed{(S_{\rm BV},S_{\rm BV})=0.}
 \label{eq:V147-CME}
\end{equation}
No boundary charge appears because (c|_\Sigma=0).
\end{theorem}

\begin{proof}
The term in \((S_{\rm BV},S_{\rm BV})/2\) with no antifield is
\(S_0'(X)[R_Xc]\), which is zero by the gauge Noether identity
\eqref{eq:V146-Noether}.  The coefficient of (X^*) states that the
commutator of two infinitesimal gauge transformations is the transformation
with parameter ([c,c]/2); it cancels the bracket term involving (c^*).
The Jacobi identity cancels the coefficient of (c^*).  Finally
(\langle\bar c^*,b\rangle) is a contractible pair and has zero
self-antibracket and zero bracket with the minimal sector.

In deriving the Noether identity, any boundary contribution is the gauge
charge paired with the boundary value of (c).  That value is zero on the
relative ghost domain.  The tangential gauge trace is preserved by
\eqref{eq:V146-tangential-closure}, so no further boundary term remains.
\end{proof}

Define the BV differential and the linearized Slavnov operator by
\begin{equation}
 \mathcal B_{S_{\rm BV}}F:=(S_{\rm BV},F).
 \label{eq:V147-linearized-Slavnov}
\end{equation}

\begin{corollary}[Off-shell nilpotent Slavnov operator]
\label{cor:V147-Slavnov-nilpotent}
For every functional (F),
\begin{equation}
 \boxed{\mathcal B_{S_{\rm BV}}^2F=0}
 \label{eq:V147-Slavnov-square}
\end{equation}
without imposing (S_0'(X)=0).
\end{corollary}

\begin{proof}
The graded Jacobi identity gives
\[
 (S_{\rm BV},(S_{\rm BV},F))
 =\frac12((S_{\rm BV},S_{\rm BV}),F).
\]
Apply Theorem~\ref{thm:V147-CME}.
\end{proof}

\begin{remark}[Why a Hessian alone was not nilpotent]
\label{rem:V147-Hessian-not-master}
At (c=X^*=0), the BV vector field has antifield component
(sX^*=-S_0'(X)).  A nonstationary background is therefore not a fixed
point of that vector field.  Dropping antifields before linearization
deletes this component and leaves precisely
\(\mathcal N_{\mathcal E_X}\).  Nilpotency belongs to the full functional
operator \eqref{eq:V147-linearized-Slavnov}, not to the isolated physical
Hessian at an arbitrary point.
\end{remark}

\subsection{Gauge fixing as a canonical BV operation}
\label{subsec:V147-gauge-fixing}

Let (\Psi_\alpha) be the gauge fermion
\eqref{eq:V146-gauge-fermion}.  The gauge-fixed Lagrangian subspace is
\begin{equation}
 \Phi_A^*=\frac{\delta\Psi_\alpha}{\delta\Phi^A}.
 \label{eq:V147-gauge-fixing-Lagrangian}
\end{equation}

\begin{proposition}[Boundary-compatible canonical gauge fixing]
\label{prop:V147-canonical-gauge-fixing}
Equation \eqref{eq:V147-gauge-fixing-Lagrangian} is a canonical
transformation of the BV phase space and preserves the classical master
equation.  Its restriction to the V146 form domain gives
\begin{equation}
 S_0+s\Psi_\alpha,
 \label{eq:V147-gauge-fixed-action}
\end{equation}
with (s\Psi_\alpha) as in \eqref{eq:V146-gf-action}.  No integration by
parts is required to prove the master identity.
\end{proposition}

\begin{proof}
The graph of the exact one-form (d\Psi_\alpha) is Lagrangian for the
canonical odd symplectic form, so translation of antifields by
(d\Psi_\alpha) preserves the antibracket.  Therefore it sends a solution
of \eqref{eq:V147-CME} to another solution.  Substitution of the graph
equations into \eqref{eq:V147-master-action} produces
\eqref{eq:V147-gauge-fixed-action}.  The form-domain boundary closure was
proved in Theorem~\ref{thm:V146-form-BRST-closure}.
\end{proof}

\subsection{The source row is the Noether defect}
\label{subsec:V147-source-row}

Differentiate the antifield vertex in \eqref{eq:V147-master-action}:
\begin{align}
 \frac{\delta^2S_{\rm BV}}{\delta X^*\delta c}
 \bigg|_{c=X^*=0}
 &=R_X,
 \label{eq:V147-source-R}\
 \frac{\delta^3S_{\rm BV}}
 {\delta X^*\delta X\delta c}
 \bigg|_{c=X^*=0}
 &=(D R)_X.
 \label{eq:V147-source-DR}
\end{align}

\begin{theorem}[BV form of the differentiated Noether identity]
\label{thm:V147-BV-Noether}
The coefficient linear in one field variation and one ghost in the master
equation is
\begin{equation}
 \boxed{
 H_XR_X+\mathcal E_X(DR)_X=0}
 \label{eq:V147-BV-Noether}
\end{equation}
in weak pairing notation.  After gauge fixing with (F_X=R_X^*), this is
equivalent to
\begin{equation}
 P_{AH}R_X-R_XP_c=\mathcal N_{\mathcal E_X}.
 \label{eq:V147-BV-intertwiner}
\end{equation}
Thus the antifield source vertex \eqref{eq:V147-source-DR} carries exactly
the term missing from the isolated quadratic complex.
\end{theorem}

\begin{proof}
Differentiate \eqref{eq:V147-CME} once with respect to (X) and once with
respect to (c), and then set ghosts and antifields to zero.  The derivative
which hits (S_0'(X)) gives (H_XR_X); the derivative which hits the
source derivative \eqref{eq:V147-source-R} gives
(\mathcal E_X(DR)_X).  This proves
\eqref{eq:V147-BV-Noether}.  Interpreting the second term by
\eqref{eq:V146-N-def} and adding (F_X^*F_XR_X=R_XP_c) proves
\eqref{eq:V147-BV-intertwiner}.
\end{proof}

\subsection{The one-loop off-shell Slavnov identity}
\label{subsec:V147-one-loop}

Let
\begin{equation}
 \Gamma=S_{\rm BV}+\hbar\Gamma_1+O(\hbar^2)
 \label{eq:V147-effective-expansion}
\end{equation}
be the renormalized effective action with antifield sources retained.  The
quantum master equation is
\begin{equation}
 \frac12(\Gamma,\Gamma)-\hbar\mathcal A=0,
 \label{eq:V147-QME-anomaly}
\end{equation}
where (\mathcal A) is the regulated measure anomaly.

\begin{theorem}[Anomaly-free one-loop Slavnov equation]
\label{thm:V147-one-loop-Slavnov}
On the faithful three-generation Standard Model branch, a
BV-compatible regulator and local counterterm scheme may be chosen so that
the nontrivial anomaly class vanishes and
\begin{equation}
 \boxed{
 \mathcal B_{S_{\rm BV}}\Gamma_1=0.}
 \label{eq:V147-one-loop-Slavnov}
\end{equation}
At zero ghosts and antifields, the coefficient linear in (c) is
\begin{equation}
 \boxed{
 \Gamma_{1,X}R_X
 +\mathcal E_X\Gamma_{1,X^*c}=0.}
 \label{eq:V147-offshell-Ward}
\end{equation}
\end{theorem}

\begin{proof}
The order-(\hbar) part of
\eqref{eq:V147-QME-anomaly} is
\[
 \mathcal B_{S_{\rm BV}}\Gamma_1=\mathcal A_1.
\]
The possible nontrivial local gauge-anomaly class is represented by the
six-form polynomial computed in V138, and it vanishes for the complete
Standard Model representation.  V141 removes the faithful-branch global
eta obstruction, including the renewal boundary pairing.  Regulator
breakings in a trivial local cohomology class are removed by local
counterterms, leaving \eqref{eq:V147-one-loop-Slavnov}.  Differentiate that
identity once with respect to (c) and then set ghosts and antifields to
zero.  The derivative of the classical antifield source gives (R_X), and
the derivative of the one-loop antifield source is
(\Gamma_{1,X^*c}).  The two surviving terms are exactly
\eqref{eq:V147-offshell-Ward}.
\end{proof}

\begin{firewall}[The determinant card alone is not off-shell gauge
independent]
\label{fire:V147-determinant-alone}
On a nonstationary background, the gauge-fixed determinant contribution to
(\Gamma_1) need not satisfy
(\Gamma_{1,X}R_X=0) by itself.  Equation
\eqref{eq:V147-offshell-Ward} shows the missing term explicitly:
it is the background Euler covector paired with the one-loop antifield
source.  Gauge independence of the source-free effective action follows on
its quantum stationary locus, or after the source insertion is retained;
it cannot be proved by deleting
\(\mathcal N_{\mathcal E_X}\) inside a heat trace.
\end{firewall}

\subsection{A master-equation-preserving finite cutoff}
\label{subsec:V147-BV-pushforward}

At a common finite regulator, let the odd symplectic BV space split as
\begin{equation}
 \mathcal E=\mathcal E_{\rm low}\oplus\mathcal E_{\rm high},
 \qquad
 \Delta=\Delta_{\rm low}+\Delta_{\rm high},
 \label{eq:V147-BV-splitting}
\end{equation}
with fields and antifields projected in dual pairs.  Let
(\mathcal L_{\rm high}\subset\mathcal E_{\rm high}) be a Lagrangian
gauge-fixing subspace.  Define the effective low-mode density by
\begin{equation}
 e^{-S_{\rm eff}/\hbar}
 :=\int_{\mathcal L_{\rm high}}
 e^{-S/\hbar}.
 \label{eq:V147-BV-pushforward}
\end{equation}

\begin{theorem}[Finite BV pushforward preserves the quantum master
equation]
\label{thm:V147-pushforward-QME}
If
\begin{equation}
 \Delta e^{-S/\hbar}=0
 \label{eq:V147-density-QME}
\end{equation}
and the finite-dimensional BV Stokes boundary term on
(\mathcal L_{\rm high}) is zero, then
\begin{equation}
 \boxed{
 \Delta_{\rm low}e^{-S_{\rm eff}/\hbar}=0.}
 \label{eq:V147-effective-QME}
\end{equation}
The result is unchanged under a deformation of
(\mathcal L_{\rm high}) through Lagrangian gauge-fixing subspaces, up to
a (\Delta_{\rm low})-exact canonical transformation.
\end{theorem}

\begin{proof}
Using \eqref{eq:V147-BV-splitting} and
\eqref{eq:V147-density-QME},
\begin{align*}
 \Delta_{\rm low}e^{-S_{\rm eff}/\hbar}
 &=\int_{\mathcal L_{\rm high}}
 \Delta_{\rm low}e^{-S/\hbar}\\
 &=-\int_{\mathcal L_{\rm high}}
 \Delta_{\rm high}e^{-S/\hbar}=0.
\end{align*}
The last equality is the finite-dimensional BV Stokes theorem.  For a
Lagrangian deformation, differentiate the fibre integral.  The variation
is the integral of a (\Delta_{\rm high})-exact Hamiltonian deformation;
moving the total BV Laplacian through the integral leaves a
(\Delta_{\rm low})-exact term.  This is precisely an infinitesimal
canonical transformation of the low effective density.
\end{proof}

\begin{corollary}[Exact composition of renewal BV integrations]
\label{cor:V147-pushforward-composition}
For a symplectic splitting into low, middle, and high sectors, integrating
high and then middle gives the same effective density as integrating their
product Lagrangian at once:
\begin{equation}
 \boxed{
 (\pi_{20})_*=(\pi_{10})_*(\pi_{21})_* .}
 \label{eq:V147-pushforward-composition}
\end{equation}
Both sides satisfy the low-mode quantum master equation.
\end{corollary}

\begin{proof}
At a finite regulator this is Fubini's theorem for the product Lagrangian
measure, with the fixed BV orientation.  Apply
Theorem~\ref{thm:V147-pushforward-QME} at either stage.
\end{proof}

\begin{firewall}[Naive mode compression is not BV pushforward]
\label{fire:V147-no-naive-compression}
Replacing (S) by its restriction to low fields discards vertices with
high internal legs and does not implement
\eqref{eq:V147-BV-pushforward}.  Nonlinear gauge transformations also need
not preserve a linear spectral subspace.  Therefore a projected classical
action can violate the master equation even when the full action satisfies
it.  The renewal cutoff map must be a dual-pair BV fibre integral, or it
must print its exact master-equation defect.
\end{firewall}

\begin{assessment}[BV completion status]
\label{ass:V147-BV-status}
The missing nonstationary source row is now identified and incorporated in
an exact nilpotent functional complex.  At finite dimension, symplectic BV
pushforward gives associative cutoff maps which preserve the quantum master
equation.  The remaining analytic task is to construct these paired
splittings for the actual renewal regulators and prove that the
renormalized infinite-dimensional fibre integrals converge.  V147 does not
replace that analytic task by a formal low-mode restriction.
\end{assessment}

\subsection{V147 ledger and next acquisition}
\label{subsec:V147-ledger}

\paragraph{New exact conclusions.}
\begin{enumerate}[label=\textup{(\arabic*)},leftmargin=2.8em]
\item Theorem~\ref{thm:V147-CME} constructs the boundary-compatible
      classical BV master action.
\item Corollary~\ref{cor:V147-Slavnov-nilpotent} proves off-shell
      nilpotency with the antifield component retained.
\item Theorem~\ref{thm:V147-BV-Noether} identifies the V146
      background-equation defect with the derivative of the antifield
      source vertex.
\item Theorem~\ref{thm:V147-one-loop-Slavnov} states the anomaly-free
      one-loop off-shell Ward identity with its indispensable source term.
\item Theorem~\ref{thm:V147-pushforward-QME} proves that a finite
      symplectic BV fibre integral preserves the quantum master equation.
\item Corollary~\ref{cor:V147-pushforward-composition} gives the exact
      three-cutoff composition law for those interacting integrations.
\end{enumerate}

\paragraph{Next forced step.}
V148 should construct an explicit renewal spectral splitting which is
symplectic in field--antifield pairs and compatible with the V120 Hodge and
V142 chiral regulators.  It must calculate the master-equation defect of a
naive cutoff and show how BV pushforward restores the omitted high-leg
contractions.  If this requires an unproved uniform propagator bound, the
programme should go upstream to the cofinal quadratic resolvent rather than
assume convergence of the interacting fibre integral.

\begin{assessment}[V147 continuum status]
\label{ass:V147-continuum-status}
The full Standard--Model--Einstein continuum remains unproved.  V147 closes
the algebraic nonstationary BV/source identity and the exact finite
pushforward composition law.  Actual cofinal propagator estimates,
renormalized interacting BV integration, microscopic gravitational
normalization, nonlinear existence, Einstein stress closure, and the
interacting continuum limit remain open.
\end{assessment}

\section*{V147 exact checks and append-only record}
\addcontentsline{toc}{section}{V147 exact checks and append-only record}

The BV/source checksum is
\begin{equation}
 \boxed{
 (S_{\rm BV},S_{\rm BV})=0,
 \qquad
 \mathcal B_{S_{\rm BV}}^2=0,
 \qquad
 \Delta_{\rm low}(\pi_*e^{-S/\hbar})=0.}
 \label{eq:V147-final-check}
\end{equation}
The append operation preserves the complete V146 prefix byte for byte before
its former live terminator.  The assembled V147 file is checked for one live
final terminator, balanced environments and braces, duplicate labels, newly
introduced unresolved references, display math delimiters, and control
characters before the exact V146 predecessor is removed.  The companion
files are not modified.

% =============================================================================
% END APPENDED VERSION V147 MATERIAL
% =============================================================================

% =============================================================================
% BEGIN APPENDED VERSION V148 MATERIAL
% =============================================================================

\section{V148: symplectic renewal spectral pairs, the exact naive-cutoff
master defect, and the high-leg restoration formula}
\label{sec:V148}

\subsection{Purpose}
\label{subsec:V148-purpose}

V147 proves that finite BV pushforward is the correct cutoff operation but
leaves the renewal splitting abstract.  V148 constructs a field--antifield
spectral splitting from dual projectors and combines the established Hodge,
gauge--Higgs, and chiral regulators.  It then computes exactly what goes
wrong if one merely sets high modes to zero.  The defect is the high-sector
part of the original antibracket, containing precisely the gauge variations
and equations which leave the low spectral space.

The final formula writes the genuine effective action as a connected
high-mode cumulant.  It displays the contractions omitted by naive
restriction and isolates the next analytic input: a uniformly controlled
high propagator together with local renormalization of its coincident
contractions.

\subsection{Dual projectors and the odd symplectic form}
\label{subsec:V148-dual-projectors}

Let (V) be a finite regulated field space and (V^*[-1]) its shifted
antifield space.  The canonical odd symplectic form is
\begin{equation}
 \omega\bigl((v,\lambda),(w,\mu)\bigr)
 :=\mu(v)-(-1)^{(|v|+1)(|w|+1)}\lambda(w).
 \label{eq:V148-odd-symplectic}
\end{equation}
Let (P:V\to V) be a projector and (Q=I-P).  Define the dual projectors
by
\begin{equation}
 (P^*\lambda)(v):=\lambda(Pv),
 \qquad Q^*:=I-P^*.
 \label{eq:V148-dual-projectors}
\end{equation}

\begin{lemma}[Symplectic low--high decomposition]
\label{lem:V148-symplectic-split}
The decomposition
\begin{align}
 \mathcal E_{\rm low}&:=PV\oplus P^*V^*[-1],
 &
 \mathcal E_{\rm high}&:=QV\oplus Q^*V^*[-1]
 \label{eq:V148-low-high}
\end{align}
is symplectic and orthogonal:
\begin{equation}
 \omega(\mathcal E_{\rm low},\mathcal E_{\rm high})=0.
 \label{eq:V148-symplectic-orthogonal}
\end{equation}
In Darboux coordinates adapted to this splitting, the finite BV Laplacian
satisfies
\begin{equation}
 \Delta=\Delta_{\rm low}+\Delta_{\rm high}.
 \label{eq:V148-Delta-split}
\end{equation}
\end{lemma}

\begin{proof}
If (v=Pv) and (\mu=Q^*\mu), then
(\mu(v)=\mu(QPv)=0).  The other mixed pairing vanishes identically in the
same way.  The restrictions of \(
\omega
\) to the two summands are nondegenerate because (P^*) and (Q^*) are
the full algebraic duals of (P) and (Q).  Choose dual bases separately
on the two summands.  The coordinate definition of (\Delta) then splits
into the two disjoint sums in \eqref{eq:V148-Delta-split}.
\end{proof}

\begin{firewall}[Equal-looking projectors are not enough]
\label{fire:V148-dual-gate}
Projecting fields and antifields by two independently chosen orthogonal
matrices need not preserve their evaluation pairing.  The antifield
projector must be the exact dual (P^*).  Otherwise the mixed symplectic
block is nonzero, \eqref{eq:V148-Delta-split} fails, and the BV Stokes proof
of Theorem~\ref{thm:V147-pushforward-QME} cannot be applied.
\end{firewall}

\subsection{The Standard Model renewal projector}
\label{subsec:V148-SM-projector}

At a stationary smooth background and a common finite threshold
(\Lambda_q), choose the field projector (P_q) blockwise as follows.
\begin{enumerate}[label=\textup{(\roman*)},leftmargin=2.8em]
\item On the scalar ghost and longitudinal gauge complex, retain complete
      common Hodge eigenspaces as in
      Proposition~\ref{prop:V120-common-Hodge-cutoff}, including the exact
      finite Schur sector of V120.
\item On the coupled gauge--Higgs block, use complete spectral subspaces of
      the self-adjoint matrix realization
      \eqref{eq:V137-matrix-Robin}; include a common infrared shift justified
      by Proposition~\ref{prop:V137-cofinal-form}.
\item On the chiral sector, use the spectral-section finite band and its
      invertible complement from \eqref{eq:V142-low-high-split}.
\item Project every antighost--multiplier target doublet with the dual of
      the corresponding ghost source projector.
\item Project all antifields by the exact dual (P_q^*).
\end{enumerate}

\begin{theorem}[Linear symplectic and BRST compatibility]
\label{thm:V148-linear-compatible-projector}
On a stationary background, the projector above can be chosen so that
\begin{equation}
 P_{AH,q}R_X=R_XP_{c,q}
 \label{eq:V148-projected-chain}
\end{equation}
on the regulated longitudinal spectral core.  Together with the dual
antifield projectors it defines the symplectic splitting
\eqref{eq:V148-low-high} for the full quadratic Standard Model BV complex.
\end{theorem}

\begin{proof}
Stationarity removes the defect in
\eqref{eq:V146-intertwining-defect}.  Hence functional calculus gives
\[
 f(P_{AH})R_X=R_Xf(P_c)
\]
for every bounded spectral cutoff function (f), after kernels are treated
by the exact V120 finite Schur complex.  Complete common Hodge eigenspaces
give the same conclusion in the de Rham rows.  The gauge--Higgs matrix
projector is orthogonal because its realization is self-adjoint; the chiral
spectral section retains its finite relative line rather than deleting
crossings.  Applying Lemma~\ref{lem:V148-symplectic-split} to the direct sum
of these field projectors and their exact duals proves the symplectic claim.
\end{proof}

\begin{remark}[Nonstationary backgrounds]
\label{rem:V148-nonstationary-projector}
When (\mathcal E_X\ne0), the commutator in
\eqref{eq:V148-projected-chain} contains the projected Noether defect
\(\mathcal N_{\mathcal E_X}\).  One must then retain the antifield source
and perform the full BV pushforward.  There is no exact field-only spectral
subcomplex to project.
\end{remark}

\subsection{Exact master defect of naive restriction}
\label{subsec:V148-naive-defect}

Let (\iota:\mathcal E_{\rm low}\hookrightarrow\mathcal E) be the
inclusion and define the naive restricted action
\begin{equation}
 S_{\rm res}:=\iota^*S_{\rm BV},
 \label{eq:V148-restricted-action}
\end{equation}
which simply sets every high field and antifield to zero.  Split the full
antibracket as
\begin{equation}
 (F,G)=(F,G)_{\rm low}+(F,G)_{\rm high}.
 \label{eq:V148-bracket-split}
\end{equation}

\begin{theorem}[Exact naive-cutoff master defect]
\label{thm:V148-naive-master-defect}
If the full action obeys the classical master equation, then
\begin{equation}
 \boxed{
 \frac12(S_{\rm res},S_{\rm res})_{\rm low}
 =-\frac12\,\iota^*(S_{\rm BV},S_{\rm BV})_{\rm high}.}
 \label{eq:V148-naive-master-defect}
\end{equation}
In dual high coordinates \((h^a,h_a^*)\), the right side is
\begin{equation}
 -\sum_{a\in\rm high}(-1)^{|h^a|}
 \left.
 \frac{\delta_rS_{\rm BV}}{\delta h^a}
 \frac{\delta_lS_{\rm BV}}{\delta h_a^*}
 \right|_{h=h^*=0}.
 \label{eq:V148-high-defect-coordinates}
\end{equation}
Thus restriction satisfies the master equation if and only if this complete
high-sector contraction vanishes.
\end{theorem}

\begin{proof}
Pull back the full master equation and use
\eqref{eq:V148-bracket-split}:
\[
 0=\frac12\iota^*(S_{\rm BV},S_{\rm BV})
 =\frac12(S_{\rm res},S_{\rm res})_{\rm low}
 +\frac12\iota^*(S_{\rm BV},S_{\rm BV})_{\rm high}.
\]
Rearrange to obtain \eqref{eq:V148-naive-master-defect}.  Expanding the
high antibracket in adapted Darboux coordinates gives
\eqref{eq:V148-high-defect-coordinates}.
\end{proof}

\begin{corollary}[Gauge-theory content of the defect]
\label{cor:V148-gauge-content}
The term in \eqref{eq:V148-high-defect-coordinates} with one classical
field equation and one gauge-source derivative contains
\begin{equation}
 -\left\langle
 Q_q\mathcal E_X(P_qX),
 Q_qR_{P_qX}(P_{c,q}c)
 \right\rangle.
 \label{eq:V148-high-Noether-pair}
\end{equation}
Further terms contain the high component of the ghost bracket and the
field dependence of the gauge generator.  Therefore closure of the linear
quadratic projector in Theorem~\ref{thm:V148-linear-compatible-projector}
does not imply closure of the interacting master action.
\end{corollary}

\begin{proof}
At zero high variables, differentiation of (S_0) in a high field
direction gives \(Q_q\mathcal E_X(P_qX)\).  Differentiation of the antifield
vertex with respect to its high antifield gives
\(Q_qR_{P_qX}(P_{c,q}c)\).  Their product is
\eqref{eq:V148-high-Noether-pair}.  Applying the same coordinate formula to
(-\langle c^*,[c,c]\rangle/2) and to the (X)-dependence of (R_X)
gives the other stated terms.
\end{proof}

\begin{firewall}[Small high fields do not make the defect zero]
\label{fire:V148-small-not-zero}
Equation \eqref{eq:V148-naive-master-defect} is an algebraic identity.
Even if \eqref{eq:V148-high-Noether-pair} is small in norm at large cutoff,
the naively restricted action does not satisfy the master equation at that
finite cutoff.  A controlled defect estimate can be useful for convergence,
but it cannot replace the exact BV pushforward when deriving a Ward
identity.
\end{firewall}

\subsection{The omitted high-leg contractions}
\label{subsec:V148-high-leg}

On a gauge-fixed high Lagrangian, write the finite action as
\begin{equation}
 S(\ell,h)
 =S_{\rm res}(\ell)
 +\frac12\langle h,A_qh\rangle
 +I_q(\ell,h),
 \qquad I_q(\ell,0)=0.
 \label{eq:V148-action-split}
\end{equation}
Let (d\mu_{C_q}) be the normalized graded Gaussian with covariance
(C_q=A_q^{-1}).  The exact BV pushforward is
\begin{equation}
 \boxed{
 S_{{\rm eff},q}(\ell)
 =S_{\rm res}(\ell)
 -\hbar\log
 \left\langle e^{-I_q(\ell,h)/\hbar}\right\rangle_{C_q}.}
 \label{eq:V148-effective-cumulant}
\end{equation}

\begin{proposition}[Connected high-leg expansion]
\label{prop:V148-connected-expansion}
At a finite regulator, \eqref{eq:V148-effective-cumulant} has the exact
formal cumulant expansion
\begin{equation}
 S_{{\rm eff},q}
 =S_{\rm res}
 +\sum_{n\ge1}
 \frac{(-1)^{n+1}}{n!\,\hbar^{,n-1}}
 \left\langle I_q^n\right\rangle_{C_q}^{\rm conn}.
 \label{eq:V148-cumulant-expansion}
\end{equation}
Every term after (S_{\rm res}) contracts at least one high leg.  These
are exactly the terms erased by setting (h=0).
\end{proposition}

\begin{proof}
For any finite normalized measure, the logarithm of the moment-generating
function is the exponential generating series of connected cumulants:
\[
 \log\langle e^{tI_q}\rangle
 =\sum_{n\ge1}\frac{t^n}{n!}
 \langle I_q^n\rangle^{\rm conn}.
\]
Set (t=-1/\hbar) and multiply by (-\hbar).  Since
(I_q(\ell,0)=0), each monomial contains a high variable and each nonzero
Gaussian cumulant pairs or otherwise contracts high legs.
\end{proof}

\begin{corollary}[Restoration of the master equation is collective]
\label{cor:V148-collective-restoration}
Although (S_{\rm res}) has defect
\eqref{eq:V148-naive-master-defect}, the complete sum
\eqref{eq:V148-effective-cumulant} satisfies the low quantum master equation
by Theorem~\ref{thm:V147-pushforward-QME}.  No single modewise contraction
is required to cancel the defect separately.
\end{corollary}

\begin{proof}
Equation \eqref{eq:V148-effective-cumulant} is merely the logarithm of the
BV fibre integral \eqref{eq:V147-BV-pushforward}.  Apply
Theorem~\ref{thm:V147-pushforward-QME}.  The equality concerns the complete
connected sum, so decomposing it into separately bounded channels is not
part of the proof.
\end{proof}

\subsection{The first uniform propagator gate}
\label{subsec:V148-propagator}

Suppose the shifted bosonic high operator has spectrum at least
(\Lambda_q^2) and the absolute chiral high operator has spectrum at least
(\Lambda_q).  Then their covariances obey
\begin{equation}
 \|C_{B,q}\|_{\rm op}\le\Lambda_q^{-2},
 \qquad
 \|C_{F,q}\|_{\rm op}\le\Lambda_q^{-1}.
 \label{eq:V148-propagator-op-bounds}
\end{equation}

\begin{lemma}[Spectral high-mode bound]
\label{lem:V148-propagator-bound}
Equation \eqref{eq:V148-propagator-op-bounds} follows directly from the
spectral theorem.  It is uniform in (q) if the shifts and spectral
thresholds are chosen uniformly as in
Proposition~\ref{prop:V137-cofinal-form}.
\end{lemma}

\begin{proof}
On the high bosonic subspace every shifted eigenvalue (lambda) satisfies
(\lambda\ge\Lambda_q^2), so every inverse eigenvalue is at most
(\Lambda_q^{-2}).  The Dirac statement is identical with absolute
eigenvalues at least (\Lambda_q).
\end{proof}

\begin{firewall}[Operator norm does not renormalize coincident
contractions]
\label{fire:V148-opnorm-insufficient}
The estimates \eqref{eq:V148-propagator-op-bounds} control a fixed high
line, but the number of high lines grows by Weyl's law and coincident
propagators have local divergences.  Summing diagrams using only operator
norm can therefore recreate a power of the cutoff.  The next step must
combine the exact cumulant before norms with heat-kernel/local counterterm
subtraction and a genuine trace-ideal or kernel estimate.
\end{firewall}

\begin{assessment}[Cutoff status after V148]
\label{ass:V148-cutoff-status}
The full quadratic Standard Model regulator now has an explicit symplectic
field--antifield pairing, and naive restriction has an exact, computable
master defect.  Finite BV integration restores the master equation through
the complete high-leg cumulant.  The remaining bottleneck is analytic:
uniformly renormalize those connected contractions and prove their cofinal
convergence without counting spectral outputs separately.
\end{assessment}

\subsection{V148 ledger and next acquisition}
\label{subsec:V148-ledger}

\paragraph{New exact conclusions.}
\begin{enumerate}[label=\textup{(\arabic*)},leftmargin=2.8em]
\item Lemma~\ref{lem:V148-symplectic-split} proves that dual projectors give
      an exact symplectic low--high BV decomposition.
\item Theorem~\ref{thm:V148-linear-compatible-projector} combines the
      existing Hodge, gauge--Higgs, and chiral spectral regulators on a
      stationary background.
\item Theorem~\ref{thm:V148-naive-master-defect} computes the full master
      defect of low-mode restriction before estimates.
\item Corollary~\ref{cor:V148-gauge-content} identifies its omitted
      high-equation/high-gauge-transformation pairing.
\item Proposition~\ref{prop:V148-connected-expansion} writes every omitted
      high-leg contraction as one connected cumulant.
\item Lemma~\ref{lem:V148-propagator-bound} supplies the first uniform high
      operator-norm decay and Firewall~\ref{fire:V148-opnorm-insufficient}
      identifies why it is not yet a renormalized continuum estimate.
\end{enumerate}

\paragraph{Next forced step.}
V149 should return to the microscopic gravitational normalization while the
interacting trace-ideal estimates are developed.  It should search the
attached renewal source metric and grand-tensor rules for an independently
fixed scalar curvature coupling (\xi_H) and left/right boundary weight
(g_L-2g_R).  If those data are absent, it must prove the normalization is
a free renormalized coupling rather than select a value from the one-loop
Einstein--GHY ratio alone.

\begin{assessment}[V148 continuum status]
\label{ass:V148-continuum-status}
The full Standard--Model--Einstein continuum remains unproved.  V148
constructs the finite symplectic renewal regulator and exact interacting
master-defect/restoration formulas.  Uniform renormalized cumulant bounds,
microscopic gravitational normalization, nonlinear existence, Einstein
stress closure, and the interacting cofinal limit remain open.
\end{assessment}

\section*{V148 exact checks and append-only record}
\addcontentsline{toc}{section}{V148 exact checks and append-only record}

The symplectic-cutoff checksum is
\begin{equation}
 \boxed{
 \omega(\mathcal E_{\rm low},\mathcal E_{\rm high})=0,
 \qquad
 \frac12(S_{\rm res},S_{\rm res})_{\rm low}
 =-\frac12\iota^*(S,S)_{\rm high},
 \qquad
 S_{\rm eff}=S_{\rm res}-\hbar\log\langle e^{-I/\hbar}\rangle.}
 \label{eq:V148-final-check}
\end{equation}
The append operation preserves the complete V147 prefix byte for byte before
its former live terminator.  The assembled V148 file is checked for one live
final terminator, balanced environments and braces, duplicate labels, newly
introduced unresolved references, display math delimiters, and control
characters before the exact V147 predecessor is removed.  The companion
files are not modified.

% =============================================================================
% END APPENDED VERSION V148 MATERIAL
% =============================================================================

% =============================================================================
% BEGIN APPENDED VERSION V149 MATERIAL
% =============================================================================

\section{V149: microscopic source audit, scalar--curvature counterterm
freedom, and the gravitational normalization no-selection theorem}
\label{sec:V149}

\subsection{Purpose and source status}
\label{subsec:V149-purpose}

V136 reduced the weight-two Einstein--GHY problem to the affine equation
\eqref{eq:V136-microscopic-equation}, but deliberately did not infer a
Higgs curvature coupling or Newton coefficient from that equation.  V149
now checks the attached microscopic source documents for such an independent
normalization.

The read-only source snapshot used here is
\begin{center}
\begin{tabular}{|p{0.50\textwidth}|r|p{0.29\textwidth}|}
\hline
Source file & Bytes & SHA--256\\
\hline
\texttt{Renewal\_Grand\_Tensor\_Monograph\_V067.tex}
& (2804854) &
\texttt{FACD058EC8BB210AD8B296C9D7F4D78A229D08D7A2AE378808700655CB8BC0FE}\\
\hline
\texttt{Renewal\_Einstein\_Standard\_Model\_Physical\_Source\_Metric\_}
\texttt{Duhamel\_Ancestry\_and\_Quantum\_Continuum\_Programme\_V35.tex}
& (1410598) &
\texttt{AE1FBC0A8125C37465E38E7649E2992725F4579A86F0BB17483B9AE829CFB17C}\\
\hline
\end{tabular}
\end{center}
The originally supplied V30 file was no longer present in Downloads; the
current V35 continuation of the same programme was inspected instead.  The
manuscripts are used only as mathematical source data.

\subsection{What the microscopic source fixes}
\label{subsec:V149-source-findings}

The Grand Tensor's displayed regulated Standard Model action (CA.12)
contains the gauge plaquette term, Higgs covariant edge term, Higgs quartic,
and Dirac--Yukawa term.  It contains no term proportional to
(R\,H^\dagger H).  The theorem immediately following CA.12 states that
the coefficients remain physical parameters.  Elsewhere the gravitational
classification gives
\begin{equation}
 \alpha L_{\rm Holst}+\beta L_{\rm Palatini}+\lambda L_{\rm vol},
 \label{eq:V149-gravity-class}
\end{equation}
and its non-null boundary theorem gives the canonical Dirichlet completion
\begin{equation}
 \chi\int_M\sqrt{|g|}\,R
 +2\chi\int_\Sigma\sqrt{|\gamma|}\,K
 \label{eq:V149-EH-GHY}
\end{equation}
up to orientation signs and corner terms.  The source requires the
Palatini coefficient to be nonzero on the active branch, but does not assign
its numerical value.  The inspected V35 source programme likewise contains
no explicit (R\,H^\dagger H) coupling and no numerical Newton
normalization.

\begin{proposition}[Exact source conclusion]
\label{prop:V149-source-conclusion}
The inspected microscopic sources fix the geometric bulk-to-face ratio in
\eqref{eq:V149-EH-GHY} and supply a finite minimal Standard Model action.
They do not independently fix
\begin{equation}
 \xi_H,
 \qquad g_R,
 \qquad g_L-2g_R.
 \label{eq:V149-unfixed-data}
\end{equation}
In particular, absence of (R\,H^\dagger H) from CA.12 specifies a
minimal finite classical ansatz; it is not a renormalization condition for
the continuum theory.
\end{proposition}

\begin{proof}
The terms printed in CA.12 exhaust its displayed action and do not include
the mixed curvature invariant.  Its theorem explicitly classifies the
coefficients as physical parameters.  Equation
\eqref{eq:V149-EH-GHY} fixes the face coefficient once (\chi) is given,
but leaves (\chi) itself arbitrary.  No inspected equation relates that
coefficient to (\xi_H) or supplies a numerical value for either.  The last
statement follows because a continuum renormalization condition concerns
the finite part after all allowed local counterterms have been included,
whereas CA.12 is a finite regulated classical action.
\end{proof}

\subsection{The allowed scalar--curvature coupling}
\label{subsec:V149-scalar-curvature}

Let (H) denote the complex Higgs doublet.  In four dimensions its
gauge-covariant engineering dimension is one.  Therefore
\begin{equation}
 I_\xi[H,g]
 :=\xi_H\int_M\sqrt{|g|}\,R\,H^\dagger H
 \label{eq:V149-Ixi}
\end{equation}
is a local, gauge-invariant, diffeomorphism-invariant operator of dimension
four.

\begin{theorem}[Independent renormalized-coupling theorem]
\label{thm:V149-independent-xi}
The coefficient (\xi_H) is an independent renormalized coupling.  Gauge
symmetry, diffeomorphism symmetry, the faithful (G_6) quotient, and the
Einstein--GHY variational ratio do not determine it.
\end{theorem}

\begin{proof}
The integrand in \eqref{eq:V149-Ixi} is neutral under (G_6), since
(H^\dagger H) is gauge invariant, and is a scalar density under
diffeomorphisms.  Its dimension is (2+1+1=4), so it belongs to the local
counterterm space of a four-dimensional renormalizable matter theory on a
curved background.  It is linearly independent of the Einstein--Hilbert
term because it varies with (H), and independent of the Higgs mass and
quartic terms because it varies with the curvature.  No listed symmetry
forbids it.  Hence a regulator can generate its counterterm even if the
finite classical coefficient is initially zero.  Subtraction removes the
divergence but leaves one finite coefficient which must be fixed by a
renormalization condition.  A bulk-to-boundary variational ratio constrains
the partner of this coefficient; it does not assign its finite value.
\end{proof}

For fixed induced metric and fixed Higgs boundary value, the differentiable
scalar--tensor completion is
\begin{equation}
 I_\xi^{\rm glob}
 =\xi_H\int_M\sqrt{|g|}\,R\,H^\dagger H
 +2\xi_H\int_\Sigma\sqrt{|\gamma|}\,K\,H^\dagger H.
 \label{eq:V149-Ixi-boundary}
\end{equation}

\begin{lemma}[Scalar--tensor GHY partner]
\label{lem:V149-scalar-GHY}
Under Dirichlet variations
(\delta\gamma|_\Sigma=0) and
(\delta H|_\Sigma=0), the second term in
\eqref{eq:V149-Ixi-boundary} cancels the normal derivative of
(\delta g) produced by the first.  The relative coefficient is uniquely
two among local terms linear in (K) and containing no additional normal
derivative.
\end{lemma}

\begin{proof}
Put (F=\xi_HH^\dagger H).  In Gaussian normal coordinates, variation of
(\int_M\sqrt{|g|}FR) contains the boundary derivative row
\[
 \int_\Sigma\sqrt{|\gamma|}\,F
 \left(n^\mu\nabla^\nu\delta g_{\mu\nu}
       -n^\mu\nabla_\mu(g^{\rho\sigma}\delta g_{\rho\sigma})\right).
\]
Variation of (2\int_\Sigma\sqrt{|\gamma|}FK) gives the negative of this
row.  Terms proportional to (\delta\gamma) or (\delta F) vanish under
the stated Dirichlet data.  A term (a\int FK) cancels the row only when
(a=2), proving uniqueness in the declared local class.
\end{proof}

\begin{remark}[Relation to the V136 Higgs Robin coefficient]
\label{rem:V149-Robin-relation}
Metric differentiability of \eqref{eq:V149-Ixi-boundary} and stationarity
of the Higgs quadratic action are distinct variations.  The latter produced
(\sigma_H=-2\xi_H) in
\eqref{eq:V136-natural-Robin}.  Neither relation selects (\xi_H).
\end{remark}

\subsection{Renormalized Einstein--GHY matching}
\label{subsec:V149-ren-matching}

Use the V136 normalized weight-two convention and write the total local
coefficients as
\begin{equation}
 \chi_{\rm tot}:=g_R+C_R,
 \qquad
 \zeta_{\rm tot}:=g_L+C_L,
 \label{eq:V149-total-coefficients}
\end{equation}
where (C_R,C_L) are the matter/interface coefficients and (g_R,g_L)
include bare gravity and local counterterms in the same regulator units.

\begin{theorem}[Matching is one affine renormalization condition]
\label{thm:V149-affine-ren-condition}
The differentiable Dirichlet gravitational action requires
\begin{equation}
 \zeta_{\rm tot}=2\chi_{\rm tot}.
 \label{eq:V149-total-ratio}
\end{equation}
Equivalently,
\begin{equation}
 g_L-2g_R=
 \begin{cases}
 12+4\xi_H,&\mathrm{(H_D)},\\
 12+4(\xi_H+\sigma_H),&\mathrm{(H_R)}.
 \end{cases}
 \label{eq:V149-affine-family}
\end{equation}
This fixes one affine combination and leaves at least two independent
finite renormalized parameters.  It cannot select (\xi_H).
\end{theorem}

\begin{proof}
Equation \eqref{eq:V149-total-ratio} is the ratio proved by the microscopic
GHY theorem \eqref{eq:V149-EH-GHY}.  Substitute
\eqref{eq:V149-total-coefficients} and rearrange:
\[
 g_L-2g_R=-(C_L-2C_R).
\]
The two defects (C_L-2C_R) were computed exactly in
\eqref{eq:V136-defects}, giving \eqref{eq:V149-affine-family}.  For each
choice of (g_R) and (\xi_H) there is a corresponding (g_L), so the
equation has a family of solutions rather than a selected point.
\end{proof}

\begin{corollary}[The two zero-subtraction values are scheme choices]
\label{cor:V149-zero-subtraction-values}
If one additionally sets (g_L-2g_R=0), then the Dirichlet Higgs branch
gives (\xi_H=-3).  On the action-compatible Robin branch
(\sigma_H=-2\xi_H), it gives (\xi_H=3).  These values follow from the
extra zero-defect convention and are not selected by the microscopic
source.
\end{corollary}

\begin{proof}
Set the left side of \eqref{eq:V149-affine-family} to zero.  The Dirichlet
equation is (12+4\xi_H=0).  The action-compatible Robin equation is
(12+4(\xi_H-2\xi_H)=0).  Solve the two scalar equations.  Proposition
\ref{prop:V149-source-conclusion} shows that the source does not impose the
extra convention on the matter counterterm split.
\end{proof}

\begin{firewall}[Canonical total GHY does not mean zero matter
counterterm]
\label{fire:V149-total-versus-sector}
The source theorem fixes the ratio of the \emph{total renormalized}
gravitational bulk and face terms.  It does not require the pure microscopic
gravity row and every matter determinant row to satisfy that ratio
separately.  Imposing (g_L=2g_R) before the matter contribution is added
and then forbidding the boundary counterterm would be an additional
subtraction scheme.  Corollary~\ref{cor:V149-zero-subtraction-values} shows
how that scheme manufactures an apparent value of (\xi_H).
\end{firewall}

\subsection{Running does not supply the missing boundary value}
\label{subsec:V149-running}

Let the renormalized couplings obey
\begin{equation}
 \mu\frac{d\xi_H}{d\mu}=\beta_\xi(\xi_H,g_i,Y_i,\lambda_H),
 \qquad
 \mu\frac{d\chi}{d\mu}=\beta_\chi.
 \label{eq:V149-RG}
\end{equation}

\begin{proposition}[An RG equation is not a normalization]
\label{prop:V149-RG-not-normalization}
Even if the beta functions in \eqref{eq:V149-RG} are computed exactly,
their solutions require finite data at one scale.  Therefore running cannot
replace the missing microscopic or experimental normalization of
(\xi_H) and (\chi).
\end{proposition}

\begin{proof}
Equation \eqref{eq:V149-RG} is a first-order system of ordinary differential
equations in (log\mu).  Local uniqueness determines a trajectory only
after its value at one scale is specified.  The beta functions determine
tangent vectors on coupling space, not the initial point.
\end{proof}

\begin{assessment}[Normalization disposition]
\label{ass:V149-normalization-disposition}
The microscopic normalization question is now decided rather than merely
uncomputed.  The attached source fixes the canonical total GHY ratio but
leaves the overall gravitational coefficient and Standard Model action
coefficients physical.  Curved-space renormalization introduces the
independent allowed coupling (\xi_H).  The continuum programme must
therefore be stated as a family indexed by finite renormalization conditions,
or be supplemented by measured/microscopic values.  No rigorous argument
may infer those values from the one-loop ratio alone.
\end{assessment}

\subsection{V149 ledger and next acquisition}
\label{subsec:V149-ledger}

\paragraph{New exact conclusions.}
\begin{enumerate}[label=\textup{(\arabic*)},leftmargin=2.8em]
\item Proposition~\ref{prop:V149-source-conclusion} records that the source
      fixes the total GHY ratio but not (\xi_H), (g_R), or the finite
      defect split.
\item Theorem~\ref{thm:V149-independent-xi} proves from locality, dimension,
      and symmetry that (R\,H^\dagger H) carries an independent
      renormalized coupling.
\item Lemma~\ref{lem:V149-scalar-GHY} proves the corresponding factor-two
      scalar--tensor boundary partner.
\item Theorem~\ref{thm:V149-affine-ren-condition} identifies the V136
      equation as one affine renormalization condition.
\item Corollary~\ref{cor:V149-zero-subtraction-values} proves that the two
      apparent special values arise only after an extra scheme convention.
\item Proposition~\ref{prop:V149-RG-not-normalization} proves that beta
      functions do not select the missing finite values.
\end{enumerate}

\paragraph{Next forced step.}
V150 is the next compulsory NS/YM audit.  It should then formulate the
renormalized Einstein stress identity for the parameter family identified
here, including the variation of (R\,H^\dagger H), its face momentum,
the BV source term from V147, and the condition under which the Bianchi
identity yields covariant conservation on the quantum stationary locus.

\begin{assessment}[V149 continuum status]
\label{ass:V149-continuum-status}
The full Standard--Model--Einstein continuum remains unproved.  V149
classifies the gravitational normalization as independent finite physical
data and closes the false selection route.  Uniform renormalized BV
cumulants, the full improved Einstein stress and boundary momentum,
nonlinear existence, and the interacting cofinal limit remain open.
\end{assessment}

\section*{V149 exact checks and append-only record}
\addcontentsline{toc}{section}{V149 exact checks and append-only record}

The normalization checksum is
\begin{equation}
 \boxed{
 \zeta_{\rm tot}=2\chi_{\rm tot},
 \qquad
 g_L-2g_R=-(C_L-2C_R),
 \qquad
 \xi_H\text{ is independent finite data}.}
 \label{eq:V149-final-check}
\end{equation}
The append operation preserves the complete V148 prefix byte for byte before
its former live terminator.  The assembled V149 file is checked for one live
final terminator, balanced environments and braces, duplicate labels, newly
introduced unresolved references, display math delimiters, and control
characters before the exact V148 predecessor is removed.  The companion
files are not modified.

% =============================================================================
% END APPENDED VERSION V149 MATERIAL
% =============================================================================
\clearpage
% =============================================================================
% BEGIN APPENDED VERSION V150 MATERIAL
% =============================================================================

\section{V150: compulsory companion audit, the Higgs-improved Einstein
source, and the complete BV--diffeomorphism Ward row}
\label{sec:V150}

\subsection{Purpose and proof order}
\label{subsec:V150-purpose}

V149 proved that the curvature coupling
\(\xi_H R H^\dagger H\) is independent finite data and that its
differentiable Dirichlet completion contains the scalar--tensor GHY term.
It did not compute the resulting Einstein source or face momentum.  V90
already proved the common minimal Standard Model Ward identity, and V147
already proved the off-shell BV source identity.  The present version joins
those earlier results without repeating them.

There are two reasons to do the joining at the level of one first variation.
First, the curvature term contributes simultaneously to the metric equation,
the Higgs equation, and the face momentum.  Keeping only one of those
contributions violates diffeomorphism invariance.  Second, the conversion of
an ordinary diffeomorphism to a gauge-covariant Lie derivative uses a gauge
transformation.  On a nonstationary background its one-loop Ward identity
contains the antifield source of V147.  The determinant stress by itself is
therefore not the complete off-shell source.

All calculations in this version are first made for smooth fields and test
variations supported away from corners.  Their distributional meaning is the
one induced by the displayed first-variation identities.  No product of two
uncontrolled distributions is introduced.

\subsection{The V150 NS/YM audit}
\label{subsec:V150-audit}

The shared folder was inspected read-only at
\(2026\text{-}09\text{-}05\ 10{:}31{:}25\ {-}07{:}00\).  The newest complete
companion snapshots then available were
\begin{center}
\begin{tabular}{@{}p{0.49\linewidth}r p{0.31\linewidth}@{}}
\toprule
file & bytes & SHA--256\\
\midrule
\texttt{Renewal\_NS\_Stability\_Research\_Notes\_V31.tex}
&887537&\texttt{f1121b62238ad5491bb7cf03635ea9934942edf573ed8faaa9ee79e1d38a6696}\\
\texttt{Renewal\_Yang\_Mills\_Mass\_Gap\_Research\_Notes\_V131.tex}
&3016528&\texttt{17822f375b9d01e521e45c56a42c90b416161c4d41b20065434576dd19a4059c}\\
\bottomrule
\end{tabular}
\end{center}

The NS snapshot is unchanged from V145.  Its transferable instruction remains
that the restricted signed formation quantity must be retained until the
correlated identity is formed; replacing it by a global positive stock loses
the only possible cancellation.

YM V130 replaced a coordinate-subset proliferation by the finite
fifteen-letter quartic partition alphabet and separated genuine orthogonal
output blocks from nonorthogonal algebraic partition words.  YM V131 goes
further upstream: it introduces edge interpolation before taking a norm,
projects onto row connectivity by partition-lattice Moebius inversion, and
places derivatives on actual connected links.  It does not yet prove its
global collision estimate.

\begin{proposition}[Transfer of the V131 proof order]
\label{prop:V150-audit-transfer}
At a finite Standard Model cutoff, form the complete metric and
diffeomorphism variation of the BV-pushed-forward action before splitting it
into determinants, counterterms, Higgs improvement, antifield sources, or
field sectors.  Ward cancellation is an algebraic cancellation in this
assembled functional.  A norm estimate applied to the summands before their
signed sum cannot prove the assembled Ward estimate.
\end{proposition}

\begin{proof}
Let \(\Gamma_q=\sum_{r=1}^N\Gamma_{q,r}\) be any finite decomposition and
let \(\mathcal W_\zeta\) denote its linear diffeomorphism Ward operator.  By
linearity,
\[
 \mathcal W_\zeta\Gamma_q
 =\sum_{r=1}^N\mathcal W_\zeta\Gamma_{q,r}.
\]
The vanishing of the left side does not imply vanishing of any summand.  In
the one-dimensional example
\(\mathcal W_\zeta\Gamma_{q,1}=1\) and
\(\mathcal W_\zeta\Gamma_{q,2}=-1\), the assembled norm is zero whereas the
sum of the termwise norms is two.  This is the same elementary obstruction
used by YM V130 for nonorthogonal partition words.  A decomposition into
mutually orthogonal physical output blocks permits an exact trace-norm sum;
the determinant, counterterm, improvement, and source labels do not define
such blocks.  Hence the first valid estimate is an estimate of the assembled
Ward residual.
\end{proof}

\subsection{One convention-fixed scalar--tensor first variation}
\label{subsec:V150-first-variation}

Use the V36/V90 outward spacelike normal and put
\begin{equation}
 f:=H^\dagger H,
 \qquad
 \Pi^{ab}:=K^{ab}-K\gamma^{ab}.
 \label{eq:V150-f-Pi}
\end{equation}
The nonminimal action and its face partner are those of V149,
\begin{equation}
 I_\xi^{\rm glob}[g,H]
 :=\xi_H\int_M\sqrt{|g|}\,fR
   +2\xi_H\int_\Sigma\sqrt{|\gamma|}\,fK.
 \label{eq:V150-Ixi}
\end{equation}
For a complex Higgs variation, Euler pairings below use
\(2\operatorname{Re}\langle\mathcal E_H,\delta H\rangle\).

\begin{theorem}[Complete first variation of the nonminimal pair]
\label{thm:V150-Ixi-variation}
For variations preserving the location of the face and supported away from
its corners,
\begin{align}
 \delta I_\xi^{\rm glob}
 ={}&-\xi_H\int_M\sqrt{|g|}\,
 \Bigl[fG^{\mu\nu}
 +(g^{\mu\nu}\Box-\nabla^\mu\nabla^\nu)f\Bigr]
 \delta g_{\mu\nu}
 \notag\\
 &+2\operatorname{Re}\int_M\sqrt{|g|}\,
       \langle \xi_HRH,\delta H\rangle
 \notag\\
 &+\xi_H\int_\Sigma\sqrt{|\gamma|}\,
 \Bigl[f\Pi^{ab}-\gamma^{ab}\nabla_n f\Bigr]
 \delta\gamma_{ab}
 \notag\\
 &+2\operatorname{Re}\int_\Sigma\sqrt{|\gamma|}\,
       \langle2\xi_HKH,\delta H\rangle .
 \label{eq:V150-Ixi-variation}
\end{align}
In particular, the scalar--tensor face momentum is not obtained by the
replacement \(\Pi\mapsto f\Pi\) alone; the normal-derivative term is forced.
\end{theorem}

\begin{proof}
The standard weighted-curvature identity, written for variation of the
covariant metric, is
\begin{align*}
 \delta\int_M\sqrt{|g|}\,fR
 ={}&-\int_M\sqrt{|g|}\,
 \bigl[fG^{\mu\nu}
 +(g^{\mu\nu}\Box-\nabla^\mu\nabla^\nu)f\bigr]
 \delta g_{\mu\nu}\\
 &+\int_M\sqrt{|g|}\,R\,\delta f
 +\hbox{the metric derivative face row}.
\end{align*}
Lemma~\ref{lem:V149-scalar-GHY} proves that the last row is cancelled by
the metric-derivative part of \(2\int_\Sigma\sqrt{|\gamma|}fK\).

It remains to identify the coefficient of an unfixed induced metric.  This
can be done without a coordinate expansion.  Where \(f>0\), set
\(\widetilde g=f g\).  In four dimensions,
\begin{align*}
 \widetilde\gamma_{ab}&=f\gamma_{ab},\\
 \widetilde K_{ab}
 &=f^{1/2}\left(K_{ab}+\frac{1}{2f}\gamma_{ab}\nabla_nf\right),\\
 \widetilde\Pi^{ab}
 &=f^{-3/2}\left(\Pi^{ab}
                  -\frac1f\gamma^{ab}\nabla_nf\right).
\end{align*}
The Einstein-frame GHY variation is the V36 expression
\(\int_\Sigma\sqrt{|\widetilde\gamma|}\,
\widetilde\Pi^{ab}\delta\widetilde\gamma_{ab}\).  Holding \(f\) fixed in
the metric variation and substituting
\(\delta\widetilde\gamma_{ab}=f\delta\gamma_{ab}\) gives
\[
 \int_\Sigma\sqrt{|\gamma|}\,
 \bigl(f\Pi^{ab}-\gamma^{ab}\nabla_nf\bigr)
 \delta\gamma_{ab}.
\]
The identity is polynomial in \(f\) and its first derivative, so localization
and continuity remove the temporary condition \(f>0\).

Finally,
\(\delta f=(\delta H)^\dagger H+H^\dagger\delta H\).
The bulk and face terms proportional to \(\delta f\) are respectively
\(R\delta f\) and \(2K\delta f\), which give the two Higgs Euler pairings in
\eqref{eq:V150-Ixi-variation}.  Multiplication by \(\xi_H\) completes the
proof.
\end{proof}

\begin{remark}[External consistency check]
\label{rem:V150-scalar-tensor-reference}
The bulk operator and the factor-two face partner agree with the general
scalar--tensor variation in
\url{https://arxiv.org/abs/0809.4033}.  The normal-derivative term in
\eqref{eq:V150-Ixi-variation} is fixed here in the manuscript's V36
covariant-metric convention rather than imported with an unspecified
Brown--York sign.
\end{remark}

\subsection{Improved Hilbert stress and its exact divergence}
\label{subsec:V150-improved-stress}

The V90 convention is
\begin{equation}
 \delta S_{\rm matter}
 =\frac12\int_M\sqrt{|g|}\,
 T^{\mu\nu}\delta g_{\mu\nu}+\text{Euler and face terms}.
 \label{eq:V150-stress-convention}
\end{equation}

\begin{corollary}[Higgs improvement tensor]
\label{cor:V150-improvement-tensor}
The contribution of \(I_\xi^{\rm glob}\) to the bulk Hilbert stress is
\begin{equation}
 \boxed{
 T_\xi^{\mu\nu}
 =-2\xi_H\Bigl[
 fG^{\mu\nu}
 +(g^{\mu\nu}\Box-\nabla^\mu\nabla^\nu)f
 \Bigr].}
 \label{eq:V150-improvement-tensor}
\end{equation}
Its bulk Higgs Euler covector is
\begin{equation}
 \mathcal E_H^\xi=\xi_HRH.
 \label{eq:V150-Higgs-Euler-xi}
\end{equation}
\end{corollary}

\begin{proof}
Compare the bulk coefficients of \(\delta g_{\mu\nu}\) and \(\delta H\) in
\eqref{eq:V150-Ixi-variation} with
\eqref{eq:V150-stress-convention} and the declared complex Euler pairing.
\end{proof}

\begin{theorem}[The improvement is one complete Noether pair]
\label{thm:V150-improvement-divergence}
For smooth \(H\),
\begin{equation}
 \boxed{
 \nabla_\mu T_\xi^{\mu}{}_{\nu}
 =\xi_HR\nabla_\nu f
 =2\operatorname{Re}
   \langle\mathcal E_H^\xi,D_\nu H\rangle .}
 \label{eq:V150-improvement-divergence}
\end{equation}
Consequently \(T_\xi\) is not, in general, conserved separately.  It
supplies exactly the curvature part of the Higgs Euler term in the complete
Standard Model Ward identity.
\end{theorem}

\begin{proof}
Metric compatibility and the contracted Bianchi identity give
\begin{align*}
 \nabla_\mu\bigl[fG^\mu{}_{\nu}
 +(\delta^\mu_{\nu}\Box-\nabla^\mu\nabla_\nu)f\bigr]
 &=(\nabla_\mu f)G^\mu{}_{\nu}
   +\nabla_\nu\Box f-\Box\nabla_\nu f\\
 &=-\frac12R\nabla_\nu f.
\end{align*}
The last equality is the curvature commutator on the one-form \(df\).
Multiplication by \(-2\xi_H\) proves the first equality in
\eqref{eq:V150-improvement-divergence}.  Gauge invariance of \(f\) gives
\[
 \nabla_\nu f
 =2\operatorname{Re}\langle H,D_\nu H\rangle.
\]
Insert \eqref{eq:V150-Higgs-Euler-xi} to obtain the second equality.
\end{proof}

\begin{corollary}[The complete classical Standard Model source]
\label{cor:V150-total-SM-conservation}
Let \(T_{\min}\) and \(\mathcal E_H^{\min}\) denote the complete minimal
gauge--Higgs--fermion quantities of V90, and set
\begin{equation}
 T_{\rm SM}:=T_{\min}+T_\xi,
 \qquad
 \mathcal E_H:=\mathcal E_H^{\min}+\mathcal E_H^\xi.
 \label{eq:V150-total-SM-source}
\end{equation}
Then the V90 off-shell Ward identity remains valid after these replacements.
On the simultaneous gauge, Higgs, and fermion Euler shell,
\begin{equation}
 \boxed{\nabla_\mu T_{\rm SM}^{\mu}{}_{\nu}=0.}
 \label{eq:V150-total-SM-conservation}
\end{equation}
\end{corollary}

\begin{proof}
The minimal terms obey Theorem~\ref{thm:V90-classical-SM-Ward}.  Equation
\eqref{eq:V150-improvement-divergence} is precisely the additional Higgs
Euler pairing produced by the additional invariant action.  Adding the two
identities proves the off-shell statement.  Vanishing of the complete Euler
bank gives \eqref{eq:V150-total-SM-conservation}.
\end{proof}

\subsection{Einstein equation and the effective Planck coefficient}
\label{subsec:V150-Einstein-equation}

Write the renormalized Einstein--Hilbert coefficient as \(\chi_R\) and use
the action
\begin{align}
 S_{gH}:={}&\chi_R\int_M\sqrt{|g|}(R-2\Lambda_R)
 +2\chi_R\int_\Sigma\sqrt{|\gamma|}K
 \notag\\
 &+S_{\min}[g,A,H,\Psi]+I_\xi^{\rm glob}[g,H].
 \label{eq:V150-total-action}
\end{align}

\begin{theorem}[Equivalent Einstein and scalar--tensor forms]
\label{thm:V150-Einstein-forms}
The metric Euler equation of \eqref{eq:V150-total-action} is equivalently
\begin{align}
 2\chi_R(G_{\mu\nu}+\Lambda_Rg_{\mu\nu})
 &=T_{\min,\mu\nu}+T_{\xi,\mu\nu},
 \label{eq:V150-Einstein-improved}\\
 2(\chi_R+\xi_Hf)G_{\mu\nu}
 +2\chi_R\Lambda_Rg_{\mu\nu}
 &=T_{\min,\mu\nu}
 -2\xi_H(g_{\mu\nu}\Box-\nabla_\mu\nabla_\nu)f.
 \label{eq:V150-Einstein-Jordan}
\end{align}
No division by \(\chi_R+\xi_Hf\) is used in this equivalence.
\end{theorem}

\begin{proof}
The V36 covariant-metric variation of the gravitational bulk term is
\(-\chi_R(G^{\mu\nu}+\Lambda_Rg^{\mu\nu})\delta g_{\mu\nu}\).
Equation~\eqref{eq:V150-stress-convention} supplies
\(\tfrac12(T_{\min}^{\mu\nu}+T_\xi^{\mu\nu})\delta g_{\mu\nu}\).
Stationarity gives \eqref{eq:V150-Einstein-improved}.  Substitute
\eqref{eq:V150-improvement-tensor} and move its \(fG_{\mu\nu}\) term to the
left to obtain \eqref{eq:V150-Einstein-Jordan}.
\end{proof}

\begin{proposition}[Bianchi compatibility is the Higgs Ward identity]
\label{prop:V150-Bianchi-compatibility}
The divergence of \eqref{eq:V150-Einstein-improved} vanishes if and only if
the complete Standard Model Euler pairings in its diffeomorphism Ward row
vanish.  Equivalently, the apparent derivative terms generated by applying
\(\nabla^\mu\) to \eqref{eq:V150-Einstein-Jordan} cancel the curvature part
of the Higgs equation through
\eqref{eq:V150-improvement-divergence}.
\end{proposition}

\begin{proof}
The contracted Bianchi identity and metric compatibility annihilate the
left side of \eqref{eq:V150-Einstein-improved}.  Its right side has the
off-shell divergence given by the complete V90 Ward identity with the
additional pair in Theorem~\ref{thm:V150-improvement-divergence}.  This
proves the first statement.  Since
\eqref{eq:V150-Einstein-Jordan} is obtained by moving one algebraic summand
of the same equation, the second statement is the same equality with its
two sides regrouped.
\end{proof}

\begin{firewall}[The effective coefficient is a new analytic gate]
\label{fire:V150-effective-Planck-gate}
Equation~\eqref{eq:V150-Einstein-Jordan} is meaningful without a sign
assumption.  Solving it as an Einstein evolution system, or applying the
Einstein-frame argument uniformly along a cofinal sequence, requires at
least a nondegeneracy bound such as
\begin{equation}
 \inf_{q,x}\bigl|\chi_{R,q}+\xi_{H,q}H_q^\dagger H_q(x)\bigr|>0.
 \label{eq:V150-Planck-nondegeneracy}
\end{equation}
A positive-energy Einstein-frame construction normally needs the stronger
positive lower bound.  V149's finite renormalization classification does not
imply either bound.
\end{firewall}

\subsection{The scalar--tensor face momentum and boundary Ward package}
\label{subsec:V150-face-momentum}

\begin{corollary}[Renormalized scalar--tensor face momentum]
\label{cor:V150-face-momentum}
The coefficient of \(\delta\gamma_{ab}\) supplied by the Einstein--Hilbert,
GHY, and nonminimal pair is
\begin{equation}
 \boxed{
 \mathcal P_{gH}^{ab}
 =(\chi_R+\xi_Hf)\Pi^{ab}
   -\xi_H\gamma^{ab}\nabla_nf.}
 \label{eq:V150-face-momentum}
\end{equation}
The corresponding Higgs face Euler covector contributed by the nonminimal
pair is
\begin{equation}
 e_{H,\Sigma}^{\xi}=2\xi_HKH.
 \label{eq:V150-Higgs-face-Euler}
\end{equation}
\end{corollary}

\begin{proof}
Add \(\chi_R\Pi^{ab}\), from
\eqref{eq:V36-GHY-variation}, to the face metric coefficient in
\eqref{eq:V150-Ixi-variation}.  The last line of that variation gives
\eqref{eq:V150-Higgs-face-Euler}.
\end{proof}

\begin{theorem}[Whole-pair boundary Ward identity]
\label{thm:V150-boundary-Ward}
Let \(\tau_{\min}^{ab}\) be the remaining minimal/boundary Standard Model
stress in the V90 convention, and let \(\mathcal E_{\partial}\) contain the
kinetic Higgs conormal together with
\eqref{eq:V150-Higgs-face-Euler} and all other boundary Euler rows.  For a
tangential test vector field \(\zeta\), diffeomorphism invariance of the
complete first variation gives
\begin{equation}
 \boxed{
 D_a\tau_{\rm SM}^{a}{}_{b}
 =n^\mu\gamma_b{}^\nu T^{\rm SM}_{\mu\nu}
  +\mathcal W^{\rm SM}_{\partial,b},}
 \label{eq:V150-boundary-Ward}
\end{equation}
where \(T^{\rm SM}=T_{\min}+T_\xi\), and
\begin{equation}
 \int_\Sigma\sqrt{|\gamma|}\,
 \mathcal W^{\rm SM}_{\partial,b}\zeta^b
 :=\langle\mathcal E_{\partial},
             \mathbb L_\zeta^{A,K}\phi_{\partial}\rangle.
 \label{eq:V150-boundary-source}
\end{equation}
Here \(\tau_{\rm SM}\) may either retain the face coefficient of
\(I_\xi^{\rm glob}\), or that coefficient may be moved in its entirety into
\(\mathcal P_{gH}\) in \eqref{eq:V150-face-momentum}.  The resulting
interface equation is identical.
\end{theorem}

\begin{proof}
Apply the proof of Theorem~\ref{thm:V90-boundary-SM-Ward} to the enlarged
invariant action.  Theorem~\ref{thm:V150-Ixi-variation} supplies all four new
first-variation rows, so no unmatched metric, Higgs, or face term remains.
The V90 integration by parts therefore gives
\eqref{eq:V150-boundary-Ward}.  Moving the complete face metric coefficient
from one side of the interface equation to the other is an algebraic
regrouping.  Because its bulk stress and Higgs Euler partners remain in the
same total variation, the Ward identity is unchanged.
\end{proof}

\begin{firewall}[Forbidden partial regroupings]
\label{fire:V150-no-partial-regrouping}
It is inconsistent to use \((\chi_R+\xi_Hf)\Pi^{ab}\) while dropping
\(-\xi_H\gamma^{ab}\nabla_nf\), or to retain \(T_\xi\) while omitting
\(\mathcal E_H^\xi\) or \(e_{H,\Sigma}^\xi\).  These are not independent
counterterm channels.  They are coefficients of one diffeomorphism-invariant
first variation.
\end{firewall}

\subsection{The BV source inside the one-loop diffeomorphism row}
\label{subsec:V150-BV-diffeomorphism}

Let \(X\) collect the gauge and Higgs fields, and retain the antifield source
\(X^*\) and ghost \(c\) of V147.  Write
\(\mathbb L_\zeta^A X\) for the gauge-covariant Lie derivative.  With the
V90 transformation conventions, the ordinary and covariant variations are
related by
\begin{equation}
 \mathcal L_\zeta X
 =\mathbb L_\zeta^A X+R_X(\iota_\zeta A).
 \label{eq:V150-Lie-split}
\end{equation}

\begin{theorem}[One-loop BV-completed diffeomorphism Ward identity]
\label{thm:V150-BV-diffeomorphism}
Assume the anomaly-free BV-compatible one-loop scheme of V147 and a
diffeomorphism-covariant regulator/counterterm action.  At zero ghosts and
antifields, the complete one-loop first variation obeys
\begin{align}
 0={}&T_1(\nabla_{(\mu}\zeta_{\nu)})
 +\langle\Gamma_{1,X},\mathbb L_\zeta^A X\rangle
 \notag\\
 &-\left\langle\mathcal E_X,
 \Gamma_{1,X^*c}[\iota_\zeta A]\right\rangle
 +\mathcal W_{1,\partial}(\zeta).
 \label{eq:V150-BV-diffeomorphism}
\end{align}
The third term is the V147 antifield source.  It vanishes on a classical
stationary background, but it cannot be deleted in an off-shell determinant
calculation.
\end{theorem}

\begin{proof}
Diffeomorphism invariance of the complete one-loop effective action gives
the ordinary-Lie identity
\[
 0=T_1(\nabla_{(\mu}\zeta_{\nu)})
   +\langle\Gamma_{1,X},\mathcal L_\zeta X\rangle
   +\mathcal W_{1,\partial}(\zeta).
\]
Insert \eqref{eq:V150-Lie-split}.  Apply
\eqref{eq:V147-offshell-Ward} with gauge parameter
\(\iota_\zeta A\):
\[
 \langle\Gamma_{1,X},R_X(\iota_\zeta A)\rangle
 =-\left\langle\mathcal E_X,
 \Gamma_{1,X^*c}[\iota_\zeta A]\right\rangle.
\]
Substitution proves \eqref{eq:V150-BV-diffeomorphism}.  The last assertion
follows because the source is multiplied by the classical Euler covector.
\end{proof}

\begin{corollary}[Quantum stationary conservation criterion]
\label{cor:V150-quantum-conservation}
Suppose a cofinal family of complete renormalized BV effective actions
satisfies the common-domain convergence hypotheses of
Theorem~\ref{thm:V90-renormalized-Ward-gate}, including its boundary first
variation, and suppose the assembled diffeomorphism Ward defect tends to
zero.  If the limiting quantum Euler covectors and boundary Euler covectors
vanish, then
\begin{equation}
 \boxed{
 \nabla_\mu T_{\rm ren}^{\mu}{}_{\nu}=0,
 \qquad
 D_a\tau_{\rm ren}^{a}{}_{b}
 =n^\mu\gamma_b{}^\nu T^{\rm ren}_{\mu\nu}}
 \label{eq:V150-quantum-conservation}
\end{equation}
as operator-valued distribution identities on the common rigging.
\end{corollary}

\begin{proof}
At every finite regulator, use the assembled identity from
Theorem~\ref{thm:V150-BV-diffeomorphism}; do not take norms of its component
rows.  The common-form Cauchy convergence and vanishing assembled defect
permit passage to the limit exactly as in the proof of
Theorem~\ref{thm:V90-renormalized-Ward-gate}.  The quantum stationary
conditions remove the Euler and BV-source pairings.  The boundary statement
then follows from the limiting form of
\eqref{eq:V150-boundary-Ward}.
\end{proof}

\begin{assessment}[What V150 proves and what remains analytic]
\label{ass:V150-status}
V150 closes the classical algebraic source package for arbitrary finite
\(\xi_H\): the bulk improvement, Higgs Euler term, scalar--tensor face
momentum, Einstein equation, Bianchi identity, and boundary Ward identity
are coefficients of one action variation.  It also identifies the V147
antifield insertion in the one-loop gauge-covariant diffeomorphism row.

The corollary is conditional because the programme has not proved uniform
renormalized BV cumulant convergence, the effective Planck lower bound
\eqref{eq:V150-Planck-nondegeneracy}, the trace needed for the normal quantum
flux, or nonlinear cofinal compactness.  No continuum Einstein solution is
claimed here.
\end{assessment}

\subsection{V150 ledger and next acquisition}
\label{subsec:V150-ledger}

\paragraph{New exact conclusions.}
\begin{enumerate}[label=\textup{(\arabic*)},leftmargin=2.em]
\item The compulsory audit records unchanged NS V31 and current YM V131,
      and imports only their proof-order lesson.
\item Theorem~\ref{thm:V150-Ixi-variation} computes all bulk and face rows of
      the nonminimal Higgs pair in the V36 sign convention.
\item Corollary~\ref{cor:V150-improvement-tensor} gives the exact improved
      Hilbert stress.
\item Theorem~\ref{thm:V150-improvement-divergence} proves that its divergence
      is exactly the curvature part of the Higgs Euler pairing.
\item Theorem~\ref{thm:V150-Einstein-forms} derives both equivalent forms of
      the Einstein equation without dividing by the effective coefficient.
\item Corollary~\ref{cor:V150-face-momentum} proves the normal-derivative
      correction to the scalar--tensor Brown--York momentum.
\item Theorem~\ref{thm:V150-boundary-Ward} inserts the whole nonminimal pair
      into the corrected V90 boundary flux identity.
\item Theorem~\ref{thm:V150-BV-diffeomorphism} proves the one-loop
      gauge-covariant diffeomorphism identity with the V147 antifield source.
\end{enumerate}

\paragraph{Next forced step.}
V151 should attack the first nonlinear cofinal passage which these identities
make well typed.  It must state sharp strong/weak hypotheses under which
\(f_qG(g_q)\),
\((g_q\Box_{g_q}-\nabla^{g_q}\nabla^{g_q})f_q\), and the face momentum
\((\chi_{R,q}+\xi_{H,q}f_q)\Pi_q-\xi_{H,q}\gamma_q\nabla_{n_q}f_q\)
converge as distributions.  It should identify any product for which weak
convergence alone is insufficient and provide a counterexample rather than
hide the missing compactness in the Bianchi identity.

\begin{assessment}[V150 continuum status]
\label{ass:V150-continuum-status}
The full Standard--Model--Einstein continuum remains unproved.  V150 closes
the exact improved classical source and boundary momentum and locates the BV
source in the one-loop diffeomorphism identity.  Uniform interacting BV
estimates, strong/weak product compactness, the effective Planck lower bound,
nonlinear existence, and the cofinal quantum limit remain open.
\end{assessment}

\section*{V150 exact checks and append-only record}
\addcontentsline{toc}{section}{V150 exact checks and append-only record}

The convention-fixed checksum is
\begin{equation}
 \boxed{
 \begin{gathered}
 T_\xi^{\mu\nu}
 =-2\xi_H\{fG^{\mu\nu}
 +(g^{\mu\nu}\Box-\nabla^\mu\nabla^\nu)f\},\\
 \nabla_\mu T_\xi^\mu{}_{\nu}=\xi_HR\nabla_\nu f,\\
 \mathcal P_{gH}^{ab}
 =(\chi_R+\xi_Hf)\Pi^{ab}
 -\xi_H\gamma^{ab}\nabla_nf.
 \end{gathered}}
 \label{eq:V150-final-check}
\end{equation}
The append operation preserves the complete V149 prefix byte for byte before
its former live terminator.  The assembled V150 file is checked for one live
final terminator, balanced environments and braces, duplicate labels, newly
introduced unresolved references, display math delimiters, and control
characters before the exact V149 predecessor is removed.  The audit is
read-only and neither companion programme is modified.

% =============================================================================
% END APPENDED VERSION V150 MATERIAL
% =============================================================================
\clearpage
% =============================================================================
% BEGIN APPENDED VERSION V151 MATERIAL
% =============================================================================

\section{V151: strong--weak passage for the improved Einstein source,
the four-dimensional critical threshold, and the face-trace gate}
\label{sec:V151}

\subsection{Purpose}
\label{subsec:V151-purpose}

V150 put the renormalized nonminimal Higgs contribution in the form
\[
 f_qG(g_q)
 +(g_q\Box_{g_q}-\nabla^{g_q}\nabla^{g_q})f_q,
 \qquad f_q=H_q^\dagger H_q.
\]
This expression is algebraically complete, but a cofinal limit still
contains a nonlinear weak product.  The present version gives a sufficient
compactness theorem, proves why the exponent \(p>2\) is the natural
four-dimensional subcritical range for an energy-bounded Higgs field, and
shows that the endpoint \(p=2\) cannot be passed by weak convergence alone.

The boundary term is treated separately.  Bulk energy compactness gives a
Higgs trace, but it does not give a normal derivative trace.  V151 therefore
states two honest routes for the face momentum: a strong conormal regime or
direct convergence of the variational boundary functional.  This is the same
distinction already forced for the quantum flux in V90.

\subsection{A fixed-manifold cofinal chart}
\label{subsec:V151-fixed-chart}

Let \(M\) be a compact four-manifold with smooth timelike boundary
\(\Sigma\).  Cofinal discrete geometries are assumed to have been pulled
back to this fixed manifold by the reconstruction maps of the programme.
This assumption is necessary: Sobolev convergence of fields living on
different manifolds has no meaning until such identifications are supplied.

Fix \(p>1\), write \(p'=p/(p-1)\), and let \(g_q\) be Lorentzian metrics of
one fixed signature.  In every fixed finite atlas assume
\begin{align}
 g_q&\longrightarrow g
      &&\text{strongly in }W^{1,\infty},
 \label{eq:V151-metric-strong}\\
 g_q&\rightharpoonup g
      &&\text{weakly in }W^{2,p},
 \label{eq:V151-metric-weak}\\
 \sup_q\bigl(\|g_q\|_{W^{2,p}}+|g_q^{-1}\|_{L^\infty}\bigr)&<\infty.
 \label{eq:V151-metric-uniform}
\end{align}
The strong convergence is understood also for the inverses, which follows
from the uniform nondegeneracy and
\eqref{eq:V151-metric-strong}.

\begin{lemma}[Weak curvature continuity in the strong \(C^1\) chart]
\label{lem:V151-curvature-continuity}
Under \eqref{eq:V151-metric-strong}--\eqref{eq:V151-metric-uniform},
\begin{equation}
 G(g_q)\rightharpoonup G(g)
 \quad\text{weakly in }L^p(M).
 \label{eq:V151-Einstein-weak}
\end{equation}
The same conclusion holds for \(R(g_q)\).
\end{lemma}

\begin{proof}
In a fixed chart, curvature is a finite sum of terms of the form
\[
 a(g_q^{-1})\,\partial^2g_q
 \quad\hbox{and}\quad
 b(g_q^{-1})\,\partial g_q\,\partial g_q.
\]
The coefficients in the first family converge strongly in \(L^\infty\),
while \(\partial^2g_q\rightharpoonup\partial^2g\) in \(L^p\).  Hence their
products converge weakly in \(L^p\).  Every factor in the second family
converges strongly in \(L^\infty\), so those products converge strongly in
\(L^p\).  Contraction with the strongly convergent inverse metric proves the
assertions for the Ricci, scalar, and Einstein tensors.  A partition of unity
gives the global statement.
\end{proof}

\subsection{Strong Higgs density and the nonlinear curvature product}
\label{subsec:V151-bulk-product}

Let the gauge bundle be identified in the same finite trivializing cover.
Assume its background connections are bounded so that their covariant
\(H^1\) norms are uniformly equivalent to a fixed Sobolev norm.

\begin{lemma}[Strong square map]
\label{lem:V151-square-map}
If
\begin{equation}
 H_q\longrightarrow H
 \quad\text{strongly in }L^{2p'}(M),
 \label{eq:V151-Higgs-strong}
\end{equation}
then
\begin{equation}
 f_q:=H_q^\dagger H_q
 \longrightarrow f:=H^\dagger H
 \quad\text{strongly in }L^{p'}(M).
 \label{eq:V151-f-strong}
\end{equation}
\end{lemma}

\begin{proof}
The pointwise identity
\[
 |H_q^\dagger H_q-H^\dagger H|
 \le(|H_q|+|H|)|H_q-H|
\]
and Holder's inequality give
\[
 \|f_q-f\|_{p'}
 \le(\|H_q\|_{2p'}+\|H\|_{2p'})
      \|H_q-H\|_{2p'}.
\]
The right side tends to zero.
\end{proof}

\begin{lemma}[Strong--weak curvature product]
\label{lem:V151-strong-weak-product}
Under Lemmas~\ref{lem:V151-curvature-continuity} and
\ref{lem:V151-square-map},
\begin{equation}
 f_qG(g_q)\rightharpoonup fG(g)
 \quad\text{weakly in }L^1(M).
 \label{eq:V151-product-limit}
\end{equation}
The analogous scalar product \(f_qR(g_q)\) has the same limit.
\end{lemma}

\begin{proof}
For a bounded test tensor \(\varphi\), split
\begin{align*}
 \int\varphi(f_qG_q-fG)
 &=\int\varphi(f_q-f)G_q
   +\int(\varphi f)(G_q-G).
\end{align*}
The first term tends to zero by Holder's inequality, the uniform
\(L^p\) curvature bound, and \eqref{eq:V151-f-strong}.  The second tends to
zero by weak \(L^p\) convergence because \(\varphi f\in L^{p'}\).  The scalar
case is identical.
\end{proof}

\begin{lemma}[Distributional continuity of the Hessian improvement]
\label{lem:V151-Hessian-limit}
Under \eqref{eq:V151-metric-strong}--\eqref{eq:V151-metric-uniform} and
\eqref{eq:V151-f-strong},
\begin{equation}
 (g_{q,\mu\nu}\Box_{g_q}
  -\nabla^{g_q}_\mu\nabla^{g_q}_\nu)f_q
 \longrightarrow
 (g_{\mu\nu}\Box_g-\nabla^g_\mu\nabla^g_\nu)f
 \quad\text{in }\mathcal D'(M).
 \label{eq:V151-Hessian-limit}
\end{equation}
\end{lemma}

\begin{proof}
Work in one fixed chart and pair with a compactly supported smooth test
tensor density.  Expand the left side as a second-order scalar operator
\[
 a_q^{\alpha\beta}\partial_\alpha\partial_\beta f_q
 +b_q^\alpha\partial_\alpha f_q,
\]
where \(a_q\) is a smooth rational expression in \(g_q\) and \(g_q^{-1}\),
and \(b_q\) is linear in \(\partial g_q\) with such coefficients.  Integrate
the first term twice and the second once.  After the integrations, every
term is a pairing of \(f_q\) with one of the following:
\begin{enumerate}[label=\textup{(\roman*)},leftmargin=2.8em]
\item a strongly convergent coefficient times a derivative of the test;
\item a strongly convergent first metric derivative times a derivative of
      the test;
\item a coefficient times \(\partial^2g_q\) times the test;
\item a product of strongly convergent first metric derivatives times the
      test.
\end{enumerate}
The first, second, and fourth families pass by strong convergence.  In the
third family, \(\partial^2g_q\rightharpoonup\partial^2g\) in \(L^p\), while
\(f_q\) converges strongly in \(L^{p'}\); Lemma
\ref{lem:V151-strong-weak-product} applies.  Summing the finite chart
expansion proves \eqref{eq:V151-Hessian-limit}.
\end{proof}

\begin{theorem}[Cofinal continuity of the improved bulk source]
\label{thm:V151-improved-source-limit}
Assume \eqref{eq:V151-metric-strong}--\eqref{eq:V151-metric-uniform},
\eqref{eq:V151-Higgs-strong}, and
\begin{equation}
 \xi_{H,q}\longrightarrow\xi_H.
 \label{eq:V151-xi-limit}
\end{equation}
Then the V150 improvement tensors satisfy
\begin{equation}
 \boxed{
 T_{\xi,q}\longrightarrow T_\xi
 \quad\text{in }\mathcal D'(M;\operatorname{Sym}^2T^*M).}
 \label{eq:V151-improved-source-limit}
\end{equation}
If in addition
\begin{align}
 \chi_{R,q}&\longrightarrow\chi_R,
 &\Lambda_{R,q}&\longrightarrow\Lambda_R,
 \label{eq:V151-gravity-couplings}\\
 T_{\min,q}&\longrightarrow T_{\min}
 &&\text{in }\mathcal D'(M),
 \label{eq:V151-minimal-stress-limit}
\end{align}
then every distributional limit of the finite-cutoff Einstein residual is
the residual of the limiting fields.  In particular,
\begin{equation}
 2\chi_{R,q}(G(g_q)+\Lambda_{R,q}g_q)
 -T_{\min,q}-T_{\xi,q}\longrightarrow0
 \label{eq:V151-residual-limit}
\end{equation}
in \(\mathcal D'\) implies the limiting improved Einstein equation
\eqref{eq:V150-Einstein-improved}.
\end{theorem}

\begin{proof}
Use Lemma~\ref{lem:V151-strong-weak-product} for \(f_qG(g_q)\), Lemma
\ref{lem:V151-Hessian-limit} for the derivative improvement, and
\eqref{eq:V151-xi-limit} for its scalar coefficient.  This proves
\eqref{eq:V151-improved-source-limit}.  Lemma
\ref{lem:V151-curvature-continuity}, the strong metric convergence, and
\eqref{eq:V151-gravity-couplings} pass the gravitational side to the limit.
Combine with \eqref{eq:V151-minimal-stress-limit} and test
\eqref{eq:V151-residual-limit} against an arbitrary smooth compactly
supported tensor.
\end{proof}

\subsection{Why \(p>2\) is the energy-subcritical range}
\label{subsec:V151-threshold}

\begin{theorem}[Four-dimensional Higgs compactness threshold]
\label{thm:V151-p-threshold}
Suppose \(H_q\rightharpoonup H\) weakly in \(H^1(M)\) and the sequence is
bounded there.  If \(p>2\), then after passage to a subsequence
\begin{equation}
 H_q\longrightarrow H
 \quad\text{strongly in }L^{2p'}(M),
 \label{eq:V151-Rellich-exponent}
\end{equation}
so the Higgs hypothesis of
Theorem~\ref{thm:V151-improved-source-limit} follows from the energy bound.
At \(p=2\), the required exponent is the critical exponent four and the
same conclusion does not follow from a bounded \(H^1\) sequence.
\end{theorem}

\begin{proof}
On a compact four-manifold, the Sobolev embedding is
\(H^1\hookrightarrow L^4\), while Rellich compactness holds into every
\(L^s\) with \(s<4\).  The inequality
\[
 2p'<4
\]
is equivalent to \(p>2\).  Hence Rellich gives
\eqref{eq:V151-Rellich-exponent} in that range.  When \(p=2\), one has
\(2p'=4\), where the Sobolev embedding is continuous but not compact.
\end{proof}

\begin{proposition}[Endpoint concentration counterexample]
\label{prop:V151-endpoint-counterexample}
There is a smooth sequence bounded in \(H^1(B_1(0)\subset\mathbb R^4)\)
which converges weakly to zero but not strongly in \(L^4\).  Consequently
its squared Higgs densities are bounded and weakly convergent in \(L^2\)
but need not converge strongly there.
\end{proposition}

\begin{proof}
Choose nonzero \(\phi\in C_c^\infty(B_1)\) and set
\begin{equation}
 H_q(x):=q\phi(qx).
 \label{eq:V151-bubble}
\end{equation}
A change of variables gives
\begin{align*}
 \|H_q\|_2^2&=q^{-2}\|\phi\|_2^2,\\
 \|\nabla H_q\|_2^2&=\|\nabla\phi\|_2^2,\\
 \|H_q\|_4^4&=\|\phi\|_4^4.
\end{align*}
Thus \(H_q\rightharpoonup0\) in \(H^1\), while its \(L^4\) norm does not
tend to zero.  The densities \(f_q=|H_q|^2\) have constant nonzero
\(L^2\) norm and concentrate at the origin, so they do not converge strongly
in \(L^2\).
\end{proof}

\begin{proposition}[Weak times weak does not identify the source]
\label{prop:V151-weak-weak-failure}
Even without concentration, separate weak convergence of \(f_q\) in
\(L^2\) and a curvature component \(C_q\) in \(L^2\) does not determine the
limit of \(f_qC_q\).
\end{proposition}

\begin{proof}
On a flat periodic chart let
\[
 f_q(x)=1+\sin(qx_1),
 \qquad C_q(x)=\sin(qx_1).
\]
Then \(f_q\rightharpoonup1\) and \(C_q\rightharpoonup0\) in \(L^2\), but
\[
 f_qC_q=\sin(qx_1)+\sin^2(qx_1)
 \rightharpoonup\frac12
\]
as a distribution.  This differs from the product of the weak limits.  The
example is a statement about the functional-analytic inference needed for
\(f_qG_q\); it does not assert that the displayed scalar \(C_q\) by itself
is an Einstein tensor.
\end{proof}

\begin{firewall}[What can repair the endpoint]
\label{fire:V151-endpoint-repairs}
At \(p=2\), any one of the following is additional information rather than
a consequence of the energy bound: strong \(L^4\) convergence of \(H_q\),
equi-integrability of \(|H_q|^4\) together with convergence in measure, a
uniform \(H^{1+\varepsilon}\) bound, or a compensated-compactness theorem
whose hypotheses couple the actual Einstein and Higgs equations.  The
Bianchi identity alone does not supply such compactness.
\end{firewall}

\subsection{Persistence of the effective Planck lower bound}
\label{subsec:V151-Planck-limit}

\begin{lemma}[Closedness of uniform nondegeneracy]
\label{lem:V151-Planck-closed}
Assume \(\chi_{R,q}\to\chi_R\), \(\xi_{H,q}\to\xi_H\),
\(f_q\to f\) in \(L^{p'}\), and, for one \(c_*>0\),
\begin{equation}
 \chi_{R,q}+\xi_{H,q}f_q\ge c_*
 \quad\text{almost everywhere for every }q.
 \label{eq:V151-Planck-uniform}
\end{equation}
Then
\begin{equation}
 \chi_R+\xi_Hf\ge c_*
 \quad\text{almost everywhere}.
 \label{eq:V151-Planck-limit}
\end{equation}
\end{lemma}

\begin{proof}
Strong \(L^{p'}\) convergence has an almost-everywhere convergent
subsequence.  Along it,
\(\chi_{R,q}+\xi_{H,q}f_q\to\chi_R+\xi_Hf\) pointwise almost everywhere.
Pass to the limit in \eqref{eq:V151-Planck-uniform}.  Since every subsequence
has the same \(L^{p'}\) limit, the asserted inequality belongs to the full
limit.
\end{proof}

\subsection{Face momentum: strong trace and variational routes}
\label{subsec:V151-face-limit}

Write
\[
 \mathcal P_q^{ab}
 =(\chi_{R,q}+\xi_{H,q}f_q)\Pi_q^{ab}
 -\xi_{H,q}\gamma_q^{ab}\nabla_{n_q}f_q.
\]

\begin{theorem}[Sufficient strong-conormal face passage]
\label{thm:V151-face-passage}
Let \(r>1\).  Suppose on \(\Sigma\)
\begin{align}
 \gamma_q&\longrightarrow\gamma
 &&\text{strongly in }W^{1,\infty},
 \label{eq:V151-gamma-strong}\\
 \Pi_q&\rightharpoonup\Pi
 &&\text{weakly in }L^r,
 \label{eq:V151-Pi-weak}\\
 f_q|_\Sigma&\longrightarrow f|_\Sigma
 &&\text{strongly in }L^{r'},
 \label{eq:V151-face-f-strong}\\
 \nabla_{n_q}f_q&\rightharpoonup\nabla_nf
 &&\text{weakly in }L^1.
 \label{eq:V151-normal-f-weak}
\end{align}
Together with convergence of \(\chi_{R,q}\) and \(\xi_{H,q}\), these imply
\begin{equation}
 \boxed{\mathcal P_q\longrightarrow\mathcal P
 \quad\text{in }\mathcal D'(\Sigma).}
 \label{eq:V151-face-passage}
\end{equation}
\end{theorem}

\begin{proof}
The product \(f_q\Pi_q\) passes by the strong--weak argument of
Lemma~\ref{lem:V151-strong-weak-product}, now on \(\Sigma\).  The inverse
induced metrics converge strongly in \(L^\infty\), so multiplication of
\eqref{eq:V151-normal-f-weak} by \(\gamma_q^{-1}\) preserves weak
distributional convergence.  The scalar couplings converge, and adding the
two rows proves the result.
\end{proof}

\begin{proposition}[Energy traces do not supply the conormal]
\label{prop:V151-no-energy-conormal}
A uniform \(H^1(M)\) bound on \(H_q\) supplies a bounded trace in
\(H^{1/2}(\Sigma)\), but it does not define
\(\nabla_{n_q}H_q|_\Sigma\), hence does not establish
\eqref{eq:V151-normal-f-weak}.  That row requires additional equation
regularity or a variational definition.
\end{proposition}

\begin{proof}
The trace map \(H^1(M)\to H^{1/2}(\Sigma)\) is continuous.  A classical
normal derivative trace is continuous from \(H^s(M)\) only when
\(s>3/2\), while a weak conormal at energy regularity is defined through a
Green identity using the equation.  Therefore the energy norm alone has no
continuous map to the asserted normal derivative trace.
\end{proof}

\begin{theorem}[Variational face passage]
\label{thm:V151-variational-face}
Let \(\Theta_{q,\Sigma}\) be the complete boundary first-variation
functional of the finite-cutoff action, including the two terms of
\eqref{eq:V150-face-momentum}.  Suppose
\begin{equation}
 \Theta_{q,\Sigma}\longrightarrow\Theta_\Sigma
 \quad\text{in the continuous dual of the common V88/V90 test space}.
 \label{eq:V151-boundary-functional-limit}
\end{equation}
Then the limiting face momentum is well defined by transposition and the
V150 boundary Ward identity passes to the limit whenever its assembled Ward
defect tends to zero.  No separate pointwise trace
\(\nabla_nf|_\Sigma\) is required.
\end{theorem}

\begin{proof}
Equation~\eqref{eq:V151-boundary-functional-limit} defines a unique
continuous covector on boundary metric and field variations.  Its metric
component is, by definition, the distributional face momentum.  Test the
finite-cutoff Ward identity with one tangential V88 test vector.  Every term
is a component of the same continuous first variation; pass to the dual
limit and remove the assumed assembled defect.  This is precisely route
\textup{(b)} of Proposition~\ref{prop:V90-boundary-flux-routes}, now applied
to the scalar--tensor pair.
\end{proof}

\subsection{Weak stability of conservation}
\label{subsec:V151-conservation-limit}

\begin{proposition}[Conservation passes only with the assembled residual]
\label{prop:V151-conservation-limit}
Assume the hypotheses of
Theorem~\ref{thm:V151-improved-source-limit}, let
\(T_q=T_{\min,q}+T_{\xi,q}\to T\) in distributions, and suppose the
finite-cutoff complete Ward rows satisfy
\begin{equation}
 \nabla^{g_q}_\mu T_q^\mu{}_{\nu}=\mathcal R_{q,\nu},
 \qquad
 \mathcal R_q\longrightarrow0
 \quad\text{in }\mathcal D'(M).
 \label{eq:V151-Ward-residual}
\end{equation}
If the convergence of \(T_q\) is strong enough to pair with the strongly
convergent Christoffel coefficients in
\eqref{eq:V151-metric-strong}, then
\begin{equation}
 \nabla^g_\mu T^\mu{}_{\nu}=0
 \quad\text{in }\mathcal D'(M).
 \label{eq:V151-limit-conservation}
\end{equation}
\end{proposition}

\begin{proof}
Pair \eqref{eq:V151-Ward-residual} with a compactly supported smooth vector
field and integrate the divergence once by parts.  The derivative falls on
the test; the connection terms are products of \(T_q\) with coefficients
converging strongly by \eqref{eq:V151-metric-strong}.  Pass to the
distributional limit, then use \(\mathcal R_q\to0\).  The result is the weak
form of \eqref{eq:V151-limit-conservation}.
\end{proof}

\begin{firewall}[No sectorwise substitute]
\label{fire:V151-no-sectorwise-limit}
The residual in \eqref{eq:V151-Ward-residual} is the complete signed Ward
residual from V150.  Bounds on the absolute values of separate determinant,
counterterm, improvement, or BV-source residuals may be sufficient if each
tends to zero, but their boundedness alone does not preserve the cancellation
required for conservation.
\end{firewall}

\subsection{V151 ledger and next acquisition}
\label{subsec:V151-ledger}

\paragraph{New exact conclusions.}
\begin{enumerate}[label=\textup{(\arabic*)},leftmargin=2.8em]
\item Strong \(W^{1,\infty}\) and weak \(W^{2,p}\) metric convergence gives
      weak \(L^p\) convergence of the Einstein tensor.
\item Strong \(L^{2p'}\) Higgs convergence gives the exact strong factor
      needed to multiply weak curvature.
\item The derivative improvement is continuous as a distribution under the
      same hypotheses.
\item For \(p>2\), a bounded four-dimensional Higgs energy supplies the
      required strong convergence by Rellich compactness.
\item At \(p=2\), the explicit concentrating bubble proves that energy
      compactness alone does not close the source.
\item A uniform positive effective Planck bound is closed under the strong
      Higgs-density limit but is not created by it.
\item The scalar--tensor face momentum has separate strong-conormal and
      variational-dual convergence routes.
\item Conservation passes through the cofinal limit only for the assembled
      Ward residual.
\end{enumerate}

\paragraph{Next forced step.}
V152 should address the endpoint rather than assume it away.  The upstream
question is whether the actual finite-cutoff Higgs measure supplies uniform
integrability of \(|H_q|^4\), or an equivalent concentration exclusion,
after the gauge--Higgs BV pushforward.  The note should derive a precise
criterion from a coercive renormalized quartic action, identify the sign and
metric-volume hypotheses it needs, and test whether that criterion survives
normalization of the interacting partition function.  If it fails, the
programme must remain in the curvature range \(p>2\) or add a stronger Higgs
regularity estimate.

\begin{assessment}[V151 continuum status]
\label{ass:V151-continuum-status}
The full Standard--Model--Einstein continuum remains unproved.  V151 proves
a sufficient strong--weak passage for the improved source and face momentum
and isolates the critical four-dimensional Higgs concentration obstruction.
It does not prove the required metric bounds, uniform integrability,
variational boundary convergence, minimal interacting stress convergence,
or existence of a limiting solution.
\end{assessment}

\section*{V151 exact checks and append-only record}
\addcontentsline{toc}{section}{V151 exact checks and append-only record}

The compactness checksum is
\begin{equation}
 \boxed{
 \begin{gathered}
 g_q\to g\text{ in }W^{1,\infty},\quad
 g_q\rightharpoonup g\text{ in }W^{2,p},\quad
 H_q\to H\text{ in }L^{2p'}\\
 \Longrightarrow\quad
 T_{\xi,q}\to T_\xi\text{ in }\mathcal D',\\
 H^1(M^4)\Subset L^{2p'}(M^4)
 \quad\Longleftrightarrow\quad p>2
 \text{ at the required exponent.}
 \end{gathered}}
 \label{eq:V151-final-check}
\end{equation}
The append operation preserves the complete V150 prefix byte for byte before
its former live terminator.  The assembled V151 file is checked for one live
final terminator, balanced environments and braces, duplicate labels, newly
introduced unresolved references, display math delimiters, and control
characters before the exact V150 predecessor is removed.  No companion file
is modified.

% =============================================================================
% END APPENDED VERSION V151 MATERIAL
% =============================================================================
\clearpage
% =============================================================================
% BEGIN APPENDED VERSION V152 MATERIAL
% =============================================================================

\section{V152: what quartic Gibbs coercivity controls, why it does not
exclude spacetime concentration, and a sharp Orlicz tightness gate}
\label{sec:V152}

\subsection{Purpose}
\label{subsec:V152-purpose}

At the curvature endpoint \(p=2\), V151 requires strong \(L^4\) compactness
of the Higgs field.  The Standard Model potential contains a positive
quartic term on its stable branch, so the first possible repair is a uniform
quartic Gibbs estimate.  This version separates two notions which that
argument can easily conflate:
\begin{enumerate}[label=\textup{(\roman*)},leftmargin=2.8em]
\item integrability of the single random variable
      \(Y_q=\int_M|H_q|^4\), and
\item uniform integrability in the spacetime variable of the densities
      \(\rho_q(x)=|H_q(x)|^4\).
\end{enumerate}
The first controls large field-amplitude tails.  The second excludes the
V151 concentrating bubble and is what strong \(L^4\) compactness needs.
Plain quartic coercivity gives the first under a normalization hypothesis,
but it does not give the second.  A local superlinear Orlicz estimate does.

\subsection{Finite-cutoff Gibbs normalization}
\label{subsec:V152-Gibbs}

Let \((\mathcal X_q,\nu_q)\) be the normalized finite-dimensional reference
probability space at cutoff \(q\).  Let \(S_q\) be the complete Euclidean
gauge-fixed BV action after the high-mode pushforward of V147, evaluated on
the chosen real integration cycle, and put
\begin{equation}
 Z_q:=\int_{\mathcal X_q}e^{-S_q/\hbar}\,d\nu_q,
 \qquad
 d\mu_q:=Z_q^{-1}e^{-S_q/\hbar}\,d\nu_q.
 \label{eq:V152-Gibbs-law}
\end{equation}
This is a probability statement only on a cycle for which the density is
real, nonnegative, and integrable.  A formal oscillatory Lorentzian integral
does not meet that hypothesis.

Define
\begin{equation}
 Y_q(H):=\int_M|H(x)|^4\,dV_{g_q}(x).
 \label{eq:V152-global-quartic}
\end{equation}

\begin{theorem}[Global exponential moment from normalized coercivity]
\label{thm:V152-global-exponential}
Assume constants \(a>0\), \(C<\infty\), and \(z_*>0\), independent of
\(q\), such that
\begin{equation}
 S_q(H,\text{other fields})\ge aY_q(H)-C,
 \qquad Z_q\ge z_*.
 \label{eq:V152-coercivity-normalization}
\end{equation}
Then for every \(0<\theta<a/\hbar\),
\begin{equation}
 \boxed{
 \sup_q\mathbb E_{\mu_q}e^{\theta Y_q}
 \le z_*^{-1}e^{C/\hbar}.}
 \label{eq:V152-global-exponential-bound}
\end{equation}
Consequently the random variables \(Y_q\) are uniformly integrable with
respect to the cutoff probability laws.
\end{theorem}

\begin{proof}
By \eqref{eq:V152-Gibbs-law} and
\eqref{eq:V152-coercivity-normalization},
\begin{align*}
 \mathbb E_{\mu_q}e^{\theta Y_q}
 &=Z_q^{-1}\int e^{\theta Y_q-S_q/\hbar}\,d\nu_q\\
 &\le z_*^{-1}e^{C/\hbar}
      \int e^{-(a/\hbar-\theta)Y_q}\,d\nu_q\\
 &\le z_*^{-1}e^{C/\hbar},
\end{align*}
because \(Y_q\ge0\) and \(\nu_q\) is a probability measure.  The
de la Vallee Poussin criterion, applied to the scalar random variables with
the superlinear function \(t\mapsto e^{\theta t}\), proves uniform
integrability.
\end{proof}

\begin{remark}[Why the partition function hypothesis is separate]
\label{rem:V152-Z-separate}
The lower action bound controls the numerator in
\eqref{eq:V152-Gibbs-law}; it does not bound the normalizer from below.
Moreover an additive cutoff-dependent vacuum constant changes \(Z_q\) but
not \(\mu_q\).  A useful normalization statement must therefore be phrased
in an additive-constant-invariant way, for example as a uniform relative
free-energy or small-sublevel probability estimate.  V152 records
\(Z_q\ge z_*\) only after a reference normalization has been fixed.
\end{remark}

\begin{proposition}[A sufficient normalized sublevel condition]
\label{prop:V152-sublevel-Z}
If there are \(m_*>0\), \(C_0<\infty\), and measurable sets
\(E_q\subset\mathcal X_q\) such that
\begin{equation}
 \nu_q(E_q)\ge m_*,
 \qquad
 \sup_{E_q}S_q\le C_0,
 \label{eq:V152-sublevel-condition}
\end{equation}
then \(Z_q\ge m_*e^{-C_0/\hbar}\).
\end{proposition}

\begin{proof}
Restrict the integral defining \(Z_q\) to \(E_q\) and use
\(e^{-S_q/\hbar}\ge e^{-C_0/\hbar}\) there.
\end{proof}

\begin{firewall}[The full pushed-forward action is required]
\label{fire:V152-full-action-coercivity}
Positivity of the classical polynomial \(\lambda_H|H|^4\) does not prove
\eqref{eq:V152-coercivity-normalization} for the full BV-pushed-forward
action.  Gauge fixing, fermion determinants, counterterms, curvature
coupling, and the chosen integration cycle all enter \(S_q\).  Coercivity
must be proved for their assembled real part, with constants independent of
the cutoff.
\end{firewall}

\subsection{Random-variable tails are not spatial equi-integrability}
\label{subsec:V152-two-integrabilities}

For a deterministic family \(H_q\), spatial uniform integrability means
\begin{equation}
 \lim_{\delta\downarrow0}
 \sup_q\sup_{\substack{E\subset M\\ |E|\le\delta}}
 \int_E|H_q|^4\,dV=0.
 \label{eq:V152-spatial-UI}
\end{equation}

\begin{proposition}[Quartic energy does not imply spatial uniform
integrability]
\label{prop:V152-quartic-not-spatial}
A uniform bound on \(Y_q\), even a deterministic bound, does not imply
\eqref{eq:V152-spatial-UI}.  Therefore the conclusion of
Theorem~\ref{thm:V152-global-exponential} does not by itself close the
\(p=2\) source limit of V151.
\end{proposition}

\begin{proof}
Use the V151 bubble \(H_q(x)=q\phi(qx)\) in a four-dimensional chart.  It
satisfies
\[
 Y_q=\int|H_q|^4=\int|\phi|^4
\]
for every \(q\).  Thus every exponential moment of the scalar \(Y_q\) is
bounded if this deterministic sequence is regarded as a point-mass law.
Nevertheless, with
\(E_q=\operatorname{supp}\phi(q\,\cdot)\), one has
\(|E_q|\to0\) and
\[
 \int_{E_q}|H_q|^4=\int|\phi|^4>0.
\]
This contradicts \eqref{eq:V152-spatial-UI}.  Hence global amplitude-tail
control and spacetime concentration control are logically independent.
\end{proof}

\begin{corollary}[Plain quartic stability is critical, not compact]
\label{cor:V152-quartic-critical}
The positive Standard Model quartic potential is compatible with the
critical scaling of the V151 bubble.  It may bound the \(L^4\) norm but does
not break the four-dimensional concentration symmetry needed for compact
embedding at that exponent.
\end{corollary}

\begin{proof}
Both the \(H^1\) gradient energy and the \(L^4\) potential energy of the
bubble are invariant under \(H(x)\mapsto qH(qx)\) in dimension four.  The
sequence therefore stays in bounded sublevels of their sum while losing
strong \(L^4\) compactness.
\end{proof}

\subsection{The local Orlicz repair}
\label{subsec:V152-Orlicz}

Let \(\Phi:[0,\infty)\to[0,\infty)\) be convex, increasing, and satisfy
\begin{equation}
 \frac{\Phi(t)}t\longrightarrow\infty
 \qquad(t\to\infty).
 \label{eq:V152-superlinear-Phi}
\end{equation}

\begin{theorem}[Deterministic Orlicz compactness at the critical exponent]
\label{thm:V152-Orlicz-compactness}
Let \(H_q\rightharpoonup H\) weakly in \(H^1(M^4)\) and assume
\begin{equation}
 \sup_q\int_M\Phi(|H_q|^4)\,dV<\infty.
 \label{eq:V152-local-Orlicz}
\end{equation}
Then, after passage to a subsequence,
\begin{equation}
 \boxed{H_q\longrightarrow H\quad\text{strongly in }L^4(M).}
 \label{eq:V152-strong-L4}
\end{equation}
Thus \eqref{eq:V152-local-Orlicz} closes the \(p=2\) hypothesis in
Theorem~\ref{thm:V151-improved-source-limit}.
\end{theorem}

\begin{proof}
Rellich compactness gives a subsequence converging strongly in \(L^2\), and
hence in measure and almost everywhere after a further subsequence.  By
\eqref{eq:V152-superlinear-Phi} and the de la Vallee Poussin criterion,
the densities \(|H_q|^4\) are uniformly integrable in the spacetime
variable.  Their pointwise limit is \(|H|^4\).  Vitali's theorem gives
\[
 \int_M\bigl||H_q|^4-|H|^4\bigr|\,dV\longrightarrow0,
\]
so \(\|H_q\|_4\to\|H\|_4\).  The sequence is weakly convergent in
\(L^4\), because the \(H^1\) bound gives an \(L^4\) bound and its weak limit
is fixed by the convergence in measure.  Uniform convexity of \(L^4\) now
turns weak convergence plus convergence of norms into strong convergence.
\end{proof}

\begin{corollary}[A simple higher-moment criterion]
\label{cor:V152-higher-moment}
For any \(\varepsilon>0\),
\begin{equation}
 \sup_q\|H_q\|_{L^{4+\varepsilon}}<\infty
 \label{eq:V152-L4epsilon}
\end{equation}
is sufficient for \eqref{eq:V152-strong-L4}.
\end{corollary}

\begin{proof}
Take \(\Phi(t)=t^{1+\varepsilon/4}\).  Then
\(\Phi(|H_q|^4)=|H_q|^{4+\varepsilon}\), so
Theorem~\ref{thm:V152-Orlicz-compactness} applies.
\end{proof}

\begin{remark}[A logarithmic alternative]
\label{rem:V152-log-Orlicz}
The strictly weaker estimate
\[
 \sup_q\int_M|H_q|^4\log(1+|H_q|)\,dV<\infty
\]
also suffices, by taking
\(\Phi(t)=t\log(1+t^{1/4})\).  The continuum programme therefore needs a
superlinear local density estimate, not necessarily a full positive power
beyond four.
\end{remark}

\subsection{Probability-law tightness in strong \(L^4\)}
\label{subsec:V152-law-tightness}

\begin{theorem}[Uniform local Orlicz expectation implies strong-law
tightness]
\label{thm:V152-law-tightness}
Suppose the finite-cutoff Higgs laws, transported to the fixed manifold of
V151, satisfy
\begin{equation}
 \sup_q\mathbb E_{\mu_q}\left[
 \|H\|_{H^1}^2+int_M\Phi(|H|^4)\,dV\right]<\infty
 \label{eq:V152-expected-Orlicz}
\end{equation}
for one \(\Phi\) satisfying \eqref{eq:V152-superlinear-Phi}.  Then their
pushforward laws are tight on \(L^4(M)\).
\end{theorem}

\begin{proof}
For \(R>0\), define
\begin{equation}
 K_R:=\left\{H:
 \|H\|_{H^1}^2+\int_M\Phi(|H|^4)\,dV\le R\right\}.
 \label{eq:V152-compact-sublevel}
\end{equation}
Every sequence in \(K_R\) has an \(H^1\)-weakly convergent subsequence.
The proof of Theorem~\ref{thm:V152-Orlicz-compactness} then makes that
subsequence strongly convergent in \(L^4\).  Thus the closure of \(K_R\) is
compact in \(L^4\).  Markov's inequality and
\eqref{eq:V152-expected-Orlicz} give
\[
 \sup_q\mu_q(K_R^c)\le C/R.
\]
Choosing \(R\) large proves tightness.
\end{proof}

\begin{corollary}[Subsequential quantum Higgs limit]
\label{cor:V152-Prokhorov}
Because \(L^4(M)\) is a separable Banach space, the tight family in
Theorem~\ref{thm:V152-law-tightness} has weakly convergent subsequences of
probability laws.  Any such subsequence admits no spatial \(L^4\)
concentration defect.
\end{corollary}

\begin{proof}
Prokhorov's theorem gives a weakly convergent subsequence of laws on
\(L^4(M)\).  Compact containment in the strong \(L^4\) topology excludes
loss of quartic mass through the V151 bubble mechanism.
\end{proof}

\subsection{The defect-measure formulation of the missing estimate}
\label{subsec:V152-defect-measure}

\begin{proposition}[Critical Higgs concentration measure]
\label{prop:V152-defect-measure}
Let \(H_q\rightharpoonup H\) in \(H^1(M^4)\).  After a subsequence there is
a nonnegative Radon measure \(\nu_H\) such that
\begin{equation}
 |H_q|^4dV\overset{*}{\rightharpoonup}
 |H|^4dV+\nu_H.
 \label{eq:V152-defect-measure}
\end{equation}
Moreover,
\begin{equation}
 \nu_H=0
 \quad\Longleftrightarrow\quad
 H_q\longrightarrow H\text{ strongly in }L^4(M).
 \label{eq:V152-defect-equivalence}
\end{equation}
\end{proposition}

\begin{proof}
The bounded \(H^1\) sequence is bounded in \(L^4\), so the measures
\(|H_q|^4dV\) have uniformly bounded mass and a weak-star convergent
subsequence.  Rellich compactness gives almost-everywhere convergence after
refinement.  Fatou's lemma against every nonnegative continuous test
function shows that the limiting measure dominates \(|H|^4dV\); their
difference is \(\nu_H\ge0\).

If convergence is strong in \(L^4\), the fourth-power densities converge in
\(L^1\), so \(\nu_H=0\).  Conversely, \(\nu_H=0\) gives convergence of the
total fourth powers and hence of the \(L^4\) norms.  Weak \(L^4\) convergence
and uniform convexity then give strong \(L^4\) convergence.
\end{proof}

\begin{corollary}[Exact endpoint acquisition target]
\label{cor:V152-endpoint-target}
At \(p=2\), the missing compactness condition in V151 is exactly
\(\nu_H=0\).  The local Orlicz estimate
\eqref{eq:V152-local-Orlicz} is one sufficient way to prove it; the global
Gibbs estimate \eqref{eq:V152-global-exponential-bound} is not.
\end{corollary}

\subsection{V152 ledger and next acquisition}
\label{subsec:V152-ledger}

\paragraph{New exact conclusions.}
\begin{enumerate}[label=\textup{(\arabic*)},leftmargin=2.8em]
\item A coercive complete Euclidean action plus a normalized partition
      lower bound gives an exponential moment of the total quartic energy.
\item A uniform sublevel mass is one explicit sufficient normalization
      condition.
\item Total quartic-energy control does not exclude spatial concentration;
      the V151 bubble proves the separation exactly.
\item A superlinear local Orlicz estimate upgrades weak \(H^1\) convergence
      to strong \(L^4\) convergence.
\item The expected local Orlicz estimate makes the transported Higgs laws
      tight in strong \(L^4\).
\item The endpoint obstruction is represented exactly by the nonnegative
      concentration measure \(\nu_H\).
\end{enumerate}

\paragraph{Next forced step.}
V153 should not assume the new Orlicz estimate.  It should pass the full
finite-cutoff Higgs Euler and stress identities with a possibly nonzero
critical defect measure and determine which terms of the limiting Einstein
and Ward equations acquire defects.  This will show whether the equations
themselves force \(\nu_H=0\), or whether an additional analytic estimate is
indispensable.  The complete signed first variation must again be passed
before individual defect terms are named.

\begin{assessment}[V152 continuum status]
\label{ass:V152-continuum-status}
The full Standard--Model--Einstein continuum remains unproved.  V152
eliminates a false compactness route: positive quartic Gibbs coercivity alone
does not close the critical improved source.  It gives a sharp sufficient
local Orlicz/tightness gate and an exact concentration measure, but the
programme has not yet derived the required superlinear local estimate from
the pushed-forward Standard Model action.
\end{assessment}

\section*{V152 exact checks and append-only record}
\addcontentsline{toc}{section}{V152 exact checks and append-only record}

The endpoint checksum is
\begin{equation}
 \boxed{
 \sup_q\mathbb E e^{\theta\int|H_q|^4}<\infty
 \ \not\Longrightarrow\ \nu_H=0,
 \qquad
 \sup_q\int\Phi(|H_q|^4)<\infty
 \ \Longrightarrow\ H_q\to H\text{ in }L^4.}
 \label{eq:V152-final-check}
\end{equation}
The append operation preserves the complete V151 prefix byte for byte before
its former live terminator.  The assembled V152 file is checked for one live
final terminator, balanced environments and braces, duplicate labels, newly
introduced unresolved references, display math delimiters, and control
characters before the exact V151 predecessor is removed.  No companion file
is modified.

% =============================================================================
% END APPENDED VERSION V152 MATERIAL
% =============================================================================
\clearpage
% =============================================================================
% BEGIN APPENDED VERSION V153 MATERIAL
% =============================================================================

\section{V153: critical Higgs defect stress, survival of the Euler equation,
and a monotonicity criterion which removes concentration}
\label{sec:V153}

\subsection{Purpose and model row}
\label{subsec:V153-purpose}

V152 identified the nonnegative measure \(\nu_H\) carried by a loss of
strong \(L^4\) Higgs compactness.  The next question is whether the coupled
field equations automatically force that measure to vanish.  The answer at
the level of weak equations is negative: the cubic Higgs Euler term can pass
weakly at the critical exponent even when the quartic energy in the stress
develops a measure defect.

This version proves that separation for a convention-fixed Higgs row,
then returns to the complete Standard Model identity.  Write
\begin{equation}
 S_H^{\min}
 :=-\int_M\sqrt{|g|}\left(
 |D_AH|_g^2+m_H^2H^\dagger H+\lambda_H(H^\dagger H)^2
 \right),
 \qquad \lambda_H>0,
 \label{eq:V153-Higgs-action}
\end{equation}
and add \(I_\xi^{\rm glob}\) from V150.  The Lorentzian sign in
\eqref{eq:V153-Higgs-action} is used only to identify the Hilbert stress;
the compactness arguments are local Banach-space statements.  The later
monotonicity theorem is explicitly Euclidean.

With \(f=H^\dagger H\), the minimal Hilbert stress is
\begin{equation}
 T^{H,\min}_{\mu\nu}
 =2\operatorname{Re}\langle D_\mu H,D_\nu H\rangle
 -g_{\mu\nu}\left(|D_AH|_g^2+m_H^2f+\lambda_Hf^2\right).
 \label{eq:V153-minimal-Higgs-stress}
\end{equation}
The complete Higgs stress is
\begin{equation}
 T^H:=T^{H,\min}+T_\xi,
 \label{eq:V153-complete-Higgs-stress}
\end{equation}
with \(T_\xi\) as in \eqref{eq:V150-improvement-tensor}.

\subsection{Critical compactness assumptions}
\label{subsec:V153-critical-assumptions}

Work on the fixed compact four-manifold of V151.  Assume
\begin{align}
 g_q&\to g &&\text{strongly in }W^{1,\infty},
 \label{eq:V153-g-strong}\\
 g_q&\rightharpoonup g &&\text{weakly in }W^{2,2},
 \label{eq:V153-g-weak}\\
 H_q&\rightharpoonup H &&\text{weakly in }H^1,
 \label{eq:V153-H-weak}\\
 A_q&\to A &&\text{strongly in }L^4,
 \qquad \sup_q\|A_q\|_{L^4}<\infty.
 \label{eq:V153-A-strong}
\end{align}
Uniform metric nondegeneracy and the corresponding inverse bounds are
retained.  By Rellich compactness, after a subsequence,
\begin{equation}
 H_q\to H\text{ strongly in }L^s\ (s<4)
 \quad\text{and almost everywhere}.
 \label{eq:V153-Rellich}
\end{equation}

\begin{lemma}[Covariant derivative weak limit]
\label{lem:V153-DH-limit}
Under \eqref{eq:V153-H-weak}--\eqref{eq:V153-A-strong},
\begin{equation}
 D_{A_q}H_q\rightharpoonup D_AH
 \quad\text{weakly in }L^2.
 \label{eq:V153-DH-limit}
\end{equation}
\end{lemma}

\begin{proof}
The ordinary derivatives converge weakly in \(L^2\).  Split
\[
 A_qH_q-AH=(A_q-A)H_q+A(H_q-H).
\]
The first term tends strongly to zero in \(L^2\) by Holder, since
\(A_q-A\to0\) in \(L^4\) and \(H_q\) is bounded in \(L^4\).  The second
converges strongly in \(L^r\) for every \(r<2\), by
\eqref{eq:V153-Rellich}, and weakly in \(L^2\) by boundedness.  Its weak
limit is zero.  Adding the derivative part proves the claim.
\end{proof}

\subsection{The Euler nonlinearity passes without quartic norm convergence}
\label{subsec:V153-Euler-passage}

\begin{lemma}[Critical cubic weak continuity along almost-everywhere
convergence]
\label{lem:V153-cubic-limit}
Under \eqref{eq:V153-H-weak}--\eqref{eq:V153-Rellich},
\begin{equation}
 |H_q|^2H_q\rightharpoonup |H|^2H
 \quad\text{weakly in }L^{4/3}(M).
 \label{eq:V153-cubic-limit}
\end{equation}
This conclusion does not require \(H_q\to H\) strongly in \(L^4\).
\end{lemma}

\begin{proof}
The sequence converges almost everywhere to \(|H|^2H\), and
\[
 \bigl\||H_q|^2H_q\bigr\|_{4/3}^{4/3}
 =\|H_q\|_4^4
\]
is uniformly bounded by the \(H^1\) bound.  In a reflexive \(L^r\) space
with \(r>1\), boundedness and almost-everywhere convergence identify every
weakly convergent subsequence with the almost-everywhere limit.  Hence the
whole sequence has the weak limit in \eqref{eq:V153-cubic-limit}.
\end{proof}

\begin{lemma}[Curvature--Higgs Euler product]
\label{lem:V153-RH-limit}
If \(R(g_q)\rightharpoonup R(g)\) in \(L^2\), then
\begin{equation}
 R(g_q)H_q\longrightarrow R(g)H
 \quad\text{in }\mathcal D'(M).
 \label{eq:V153-RH-limit}
\end{equation}
\end{lemma}

\begin{proof}
For a smooth bounded test section \(\varphi\), write
\begin{align*}
 \int\langle R_qH_q-RH,\varphi\rangle
 &=\int R_q\langle H_q-H,\varphi\rangle
   +\int(R_q-R)\langle H,\varphi\rangle.
\end{align*}
The first term tends to zero because \(H_q\to H\) strongly in \(L^2\)
and \(R_q\) is bounded in \(L^2\).  The second tends to zero by weak
\(L^2\) convergence of \(R_q\).
\end{proof}

Define the Higgs Euler operator in weak notation by
\begin{equation}
 \mathcal E_{H,q}
 :=D_{A_q}^{*_{g_q}}D_{A_q}H_q
 +m_{H,q}^2H_q
 +2\lambda_{H,q}|H_q|^2H_q
 +\xi_{H,q}R(g_q)H_q+J_q,
 \label{eq:V153-Higgs-Euler}
\end{equation}
where \(J_q\) collects the assembled Yukawa and gauge-fixing Euler
insertions, not their absolute values.

\begin{theorem}[Critical weak passage of the Higgs Euler equation]
\label{thm:V153-Higgs-Euler-passage}
Assume \eqref{eq:V153-g-strong}--\eqref{eq:V153-A-strong}, convergence of
the scalar couplings, and
\begin{equation}
 J_q\to J,
 \qquad \mathcal E_{H,q}\to0
 \quad\text{in }H^{-1}_{\rm loc}+L^1_{\rm loc}.
 \label{eq:V153-Euler-residual}
\end{equation}
Then the limiting Higgs field satisfies
\begin{equation}
 D_A^{*_{g}}D_AH+m_H^2H+2\lambda_H|H|^2H
 +\xi_HR(g)H+J=0
 \label{eq:V153-limit-Higgs-Euler}
\end{equation}
as a distribution.  No hypothesis \(\nu_H=0\) is used.
\end{theorem}

\begin{proof}
Test \eqref{eq:V153-Higgs-Euler} against a smooth compactly supported
section.  The kinetic weak form contains one factor
\(D_{A_q}H_q\), which passes by Lemma~\ref{lem:V153-DH-limit}, multiplied by
strongly convergent metric coefficients.  The mass term passes by strong
\(L^2\) convergence.  Lemma~\ref{lem:V153-cubic-limit} passes the critical
cubic term, and Lemma~\ref{lem:V153-RH-limit} passes the curvature term.
The assumed source and residual limits finish the proof.
\end{proof}

\begin{assessment}[Euler passage does not settle stress passage]
\label{ass:V153-Euler-not-stress}
The Higgs Euler equation tests the quartic potential through the cubic map
\(|H|^2H\), which is weakly compact in \(L^{4/3}\).  The metric equation
tests the same potential through the scalar energy density \(|H|^4\), which
is only bounded in \(L^1\) and can converge to a measure.  Thus the Euler
equation may have the expected limit while the Einstein source does not.
\end{assessment}

\subsection{The complete critical stress defect}
\label{subsec:V153-stress-defect}

By the uniform \(L^2\) bound in Lemma~\ref{lem:V153-DH-limit}, after a
subsequence the quadratic kinetic stresses have weak-star limits as finite
tensor-valued Radon measures.  Define \(\mathfrak K_{\mu\nu}\) by
\begin{align}
 &2\operatorname{Re}\langle D_{A_q,\mu}H_q,D_{A_q,\nu}H_q\rangle
 -g_{q,\mu\nu}|D_{A_q}H_q|_{g_q}^2
 \notag\\
 &\hspace{4em}\overset{*}{\rightharpoonup}
 2\operatorname{Re}\langle D_{A,\mu}H,D_{A,\nu}H\rangle
 -g_{\mu\nu}|D_AH|_g^2+\mathfrak K_{\mu\nu}.
 \label{eq:V153-kinetic-defect}
\end{align}
Let \(\nu_H\) be the V152 quartic concentration measure.

The V150 improvement must be kept assembled.  Whenever its finite-cutoff
sequence is bounded in one fixed distribution space and has a subsequential
limit, define
\begin{align}
 \mathfrak I_{\mu\nu}
 :={}&\mathcal D'\!\operatorname{-lim}_q
 \left[f_qG_{\mu\nu}(g_q)
 +(g_{q,\mu\nu}\Box_{g_q}
   -\nabla^{g_q}_\mu\nabla^{g_q}_\nu)f_q\right]
 \notag\\
 &-\left[fG_{\mu\nu}(g)
 +(g_{\mu\nu}\Box_g-\nabla^g_\mu\nabla^g_\nu)f\right].
 \label{eq:V153-improvement-defect}
\end{align}
This definition deliberately does not split the two improvement terms.

\begin{theorem}[Defect decomposition of the limiting Higgs stress]
\label{thm:V153-stress-defect}
Assume the critical hypotheses above, convergence of the scalar couplings,
and existence of the limit in
\eqref{eq:V153-improvement-defect}.  Then
\begin{equation}
 T^H_q\longrightarrow T^H+\mathfrak S_H
 \quad\text{in }\mathcal D'(M),
 \label{eq:V153-stress-defect-limit}
\end{equation}
where
\begin{equation}
 \boxed{
 \mathfrak S_{H,\mu\nu}
 =\mathfrak K_{\mu\nu}
 -\lambda_Hg_{\mu\nu}\nu_H
 -2\xi_H\mathfrak I_{\mu\nu}.}
 \label{eq:V153-total-defect}
\end{equation}
The mass term has no defect.
\end{theorem}

\begin{proof}
The kinetic row is the definition
\eqref{eq:V153-kinetic-defect}.  Strong \(L^2\) convergence of \(H_q\)
gives \(f_q\to f\) strongly in \(L^1\), so the mass energy converges in
\(L^1\).  By Proposition~\ref{prop:V152-defect-measure},
\[
 f_q^2dV\overset{*}{\rightharpoonup}f^2dV+\nu_H.
\]
Strong metric convergence therefore makes the quartic stress converge to
its ordinary limit minus \(\lambda_Hg\nu_H\).  Finally the improvement
stress is \(-2\xi_H\) times the assembled bracket in
\eqref{eq:V153-improvement-defect}.  Adding the rows gives
\eqref{eq:V153-total-defect}.
\end{proof}

\begin{firewall}[The improvement defect is correlated]
\label{fire:V153-improvement-correlated}
At \(p=2\), neither a separate defect for \(f_qG(g_q)\) nor one for the
Hessian term should be estimated before their sum.  Both come from the same
metric variation of \(\int f_qR(g_q)\), and integrations by parts can move
concentration between them.  Only \(\mathfrak I\) in
\eqref{eq:V153-improvement-defect} is invariant under that regrouping.
\end{firewall}

\subsection{Einstein and Ward equations with defect stress}
\label{subsec:V153-defect-Einstein}

\begin{corollary}[The naive limit solves a defect Einstein equation]
\label{cor:V153-defect-Einstein}
Suppose the finite-cutoff improved Einstein equations hold with residual
tending to zero, the minimal non-Higgs Standard Model stress has its expected
distributional limit, and the hypotheses of
Theorem~\ref{thm:V153-stress-defect} hold.  Then
\begin{equation}
 2\chi_R(G_{\mu\nu}+\Lambda_Rg_{\mu\nu})
 =T^{\rm expected}_{\mu\nu}+\mathfrak S_{H,\mu\nu}.
 \label{eq:V153-defect-Einstein}
\end{equation}
The desired Standard--Model--Einstein equation follows only after proving
\(\mathfrak S_H=0\), not merely the limiting Higgs Euler equation.
\end{corollary}

\begin{proof}
Pass the finite-cutoff equation against a smooth test tensor.  All regular
rows converge by hypothesis, while the Higgs row converges by
\eqref{eq:V153-stress-defect-limit}.  This gives
\eqref{eq:V153-defect-Einstein}.
\end{proof}

\begin{theorem}[The complete defect is conserved]
\label{thm:V153-defect-conservation}
Assume each finite-cutoff complete stress obeys the exact V150 Ward identity
on its Euler shell, the metric converges as in
\eqref{eq:V153-g-strong}, and the total stresses converge weak-star as
measures.  Then
\begin{equation}
 \nabla_g^\mu(T^{\rm expected}_{\mu\nu}
              +\mathfrak S_{H,\mu\nu})=0
 \label{eq:V153-total-defect-conservation}
\end{equation}
in distributions.  If the limiting regular fields satisfy their complete
Euler equations so that
\(\nabla_g^\mu T^{\rm expected}_{\mu\nu}=0\), then
\begin{equation}
 \boxed{\nabla_g^\mu\mathfrak S_{H,\mu\nu}=0.}
 \label{eq:V153-defect-conserved}
\end{equation}
This condition does not force \(\mathfrak S_H=0\).
\end{theorem}

\begin{proof}
Test the finite-cutoff conservation law with a smooth compactly supported
vector field.  One integration by parts pairs the stress measures with the
covariant derivative of that test field.  The latter converges uniformly by
\eqref{eq:V153-g-strong}; hence weak-star convergence passes the identity to
the limit and proves \eqref{eq:V153-total-defect-conservation}.  Subtract the
regular limiting Ward identity to obtain
\eqref{eq:V153-defect-conserved}.

Nonzero divergence-free tensor measures exist.  For example, on a flat
periodic chart every constant symmetric tensor density is divergence free.
Therefore the last equation is a constraint on the defect, not a vanishing
theorem.
\end{proof}

\begin{corollary}[Bianchi does not remove critical concentration]
\label{cor:V153-Bianchi-not-enough}
The contracted Bianchi identity propagates the complete stress balance but
cannot distinguish the expected stress from a conserved concentration
stress.  Hence V150's exact Bianchi compatibility and V153's limiting Higgs
Euler equation do not imply \(\nu_H=0\).
\end{corollary}

\subsection{An upstream condition which does kill the defect}
\label{subsec:V153-monotonicity}

The preceding obstruction concerns weak Lorentzian/cofinal passage.  On a
Euclidean stationary branch, a stronger operator property can exclude it.
Let \(V\) be a real Hilbert space continuously embedded in \(L^4(M)\), and
let \(\mathcal A_q:V\to V'\) be the complete gauge-fixed Higgs Euler
operator after the other background fields have been fixed.

\begin{theorem}[Uniform monotonicity removes the Higgs defect]
\label{thm:V153-monotonicity}
Assume constants \(c_1,c_2>0\), independent of \(q\), such that
\begin{equation}
 \langle\mathcal A_q(u)-\mathcal A_q(v),u-v\rangle
 \ge c_1\|u-v\|_V^2+c_2\|u-v\|_4^4
 \label{eq:V153-uniform-monotonicity}
\end{equation}
for all \(u,v\in V\).  Suppose
\begin{align}
 H_q&\rightharpoonup H\quad\text{in }V,
 \label{eq:V153-V-weak}\\
 \mathcal A_q(H_q)&=J_q,
 \qquad J_q\to J\quad\text{strongly in }V',
 \label{eq:V153-source-strong}\\
 \mathcal A_q(H)&\to\mathcal A(H)=J
 \quad&\text{strongly in }V'.
 \label{eq:V153-operator-consistency}
\end{align}
Then
\begin{equation}
 \boxed{H_q\to H\text{ strongly in }V\cap L^4(M),
 \qquad \nu_H=0.}
 \label{eq:V153-monotonicity-conclusion}
\end{equation}
If the remaining metric and gauge coefficients converge strongly enough,
the kinetic and improvement defects vanish as well.
\end{theorem}

\begin{proof}
Apply \eqref{eq:V153-uniform-monotonicity} to \((u,v)=(H_q,H)\):
\begin{align*}
 c_1\|H_q-H\|_V^2+c_2\|H_q-H\|_4^4
 &\le\langle J_q-\mathcal A_q(H),H_q-H\rangle.
\end{align*}
By \eqref{eq:V153-source-strong}--
\eqref{eq:V153-operator-consistency}, the first factor on the right tends
strongly to zero in \(V'\), while \(H_q-H\) is bounded in \(V\).  The right
side tends to zero, proving strong convergence in both norms.  Proposition
\ref{prop:V152-defect-measure} gives \(\nu_H=0\).  Strong convergence of the
covariant derivatives removes \(\mathfrak K\); Theorem
\ref{thm:V151-improved-source-limit} removes \(\mathfrak I\) under its
coefficient hypotheses.
\end{proof}

\begin{lemma}[The positive quartic map supplies the \(L^4\) monotonicity]
\label{lem:V153-quartic-monotonicity}
For vectors in a real or complex finite-dimensional unitary Higgs fibre,
\begin{equation}
 \operatorname{Re}\langle |u|^2u-|v|^2v,u-v\rangle
 \ge\frac14|u-v|^4.
 \label{eq:V153-quartic-monotonicity}
\end{equation}
Thus a term \(2\lambda_H|H|^2H\) with \(\lambda_H>0\) contributes
\((\lambda_H/2)\|u-v\|_4^4\) to
\eqref{eq:V153-uniform-monotonicity}.
\end{lemma}

\begin{proof}
The standard uniform monotonicity inequality for
\(z\mapsto|z|^{p-2}z\) in a Hilbert fibre is
\[
 \operatorname{Re}\langle |u|^{p-2}u-|v|^{p-2}v,u-v\rangle
 \ge2^{2-p}|u-v|^p,
 \qquad p\ge2.
\]
Set \(p=4\) and integrate the pointwise inequality.
\end{proof}

\begin{firewall}[Why monotonicity is not yet the Lorentzian continuum]
\label{fire:V153-monotonicity-scope}
The covariant Lorentzian wave operator is not coercive on spacetime
\(H^1\), and the symmetry-breaking quadratic Higgs potential need not make
the full stationary operator globally monotone.  Gauge zero modes must also
be quotiented or fixed.  Therefore
Theorem~\ref{thm:V153-monotonicity} is an exact sufficient Euclidean branch,
not a proof that the actual coupled Lorentzian BV measure satisfies its
hypotheses.
\end{firewall}

\subsection{V153 ledger and next acquisition}
\label{subsec:V153-ledger}

\paragraph{New exact conclusions.}
\begin{enumerate}[label=\textup{(\arabic*)},leftmargin=2.8em]
\item The critical cubic Higgs Euler nonlinearity converges weakly in
      \(L^{4/3}\) without strong \(L^4\) convergence.
\item The curvature--Higgs Euler product passes because Rellich supplies
      strong \(L^2\) Higgs convergence against weak \(L^2\) curvature.
\item The limiting Higgs Euler equation can therefore be correct while its
      Einstein stress carries a defect.
\item The complete defect splits into kinetic, quartic-measure, and
      assembled improvement contributions as in
      \eqref{eq:V153-total-defect}.
\item The finite-cutoff Ward identities make the total defect stress
      divergence free, but do not make it zero.
\item A uniform strong-monotonicity estimate with strong dual sources does
      remove the defect on a coercive Euclidean branch.
\end{enumerate}

\paragraph{Next forced step.}
V154 should attack the other factor in the critical product: the metric.
In background-harmonic gauge, determine a sufficient local estimate on the
reduced Einstein operator which upgrades weak \(W^{2,2}\) curvature control
to strong convergence below the principal order.  The argument must retain
the complete Einstein residual and gauge constraint, identify the Lorentzian
compactness loss, and state a continuation criterion rather than importing
elliptic compactness into a hyperbolic problem.

\begin{assessment}[V153 continuum status]
\label{ass:V153-continuum-status}
The full Standard--Model--Einstein continuum remains unproved.  V153 proves
that weak Higgs Euler convergence and Bianchi conservation do not eliminate
critical stress concentration.  It supplies an exact defect equation and a
Euclidean monotonicity route to strong convergence, but the required
Lorentzian or probabilistic defect-removal estimate is still open.
\end{assessment}

\section*{V153 exact checks and append-only record}
\addcontentsline{toc}{section}{V153 exact checks and append-only record}

The defect checksum is
\begin{equation}
 \boxed{
 \mathcal E_{H,q}\to\mathcal E_H,
 \qquad
 T^H_q\to T^H+\mathfrak K-\lambda_Hg\nu_H-2\xi_H\mathfrak I,
 \qquad
 \nabla^\mu\mathfrak S_{H,\mu\nu}=0\ \not\Rightarrow\
 \mathfrak S_H=0.}
 \label{eq:V153-final-check}
\end{equation}
The append operation preserves the complete V152 prefix byte for byte before
its former live terminator.  The assembled V153 file is checked for one live
final terminator, balanced environments and braces, duplicate labels, newly
introduced unresolved references, display math delimiters, and control
characters before the exact V152 predecessor is removed.  No companion file
is modified.

% =============================================================================
% END APPENDED VERSION V153 MATERIAL
% =============================================================================
\clearpage
% =============================================================================
% BEGIN APPENDED VERSION V154 MATERIAL
% =============================================================================

\section{V154: Einstein-frame principal diagonalization, curvature-level
hyperbolic stability, and harmonic-constraint propagation}
\label{sec:V154}

\subsection{Purpose}
\label{subsec:V154-purpose}

V153 showed that critical weak compactness can leave a conserved stress
defect.  One possible response is to strengthen the metric side.  The
Jordan-frame equation \eqref{eq:V150-Einstein-Jordan}, however, mixes second
metric derivatives with second Higgs derivatives through
\((g\Box-\nabla\nabla)(H^\dagger H)\).  Estimating those rows separately
would return to the forbidden split of V150.

This version instead goes upstream to the action.  Where the effective
curvature coefficient
\begin{equation}
 F(H):=\chi_R+\xi_HH^\dagger H
 \label{eq:V154-F}
\end{equation}
has a uniform positive lower bound, one conformal change converts the whole
scalar--tensor pair into a constant Einstein--Hilbert coefficient and a
modified positive Higgs target metric.  The metric and Higgs principal
operators are then diagonal in the Einstein frame after harmonic and gauge
fixing.

Diagonalization does not create compactness.  A homogeneous high-frequency
wave can converge strongly at one derivative while its second derivatives
retain order-one oscillation.  V154 proves this obstruction and then states
a curvature-level continuous-dependence theorem whose hypotheses are a
concrete continuation gate for the nonlinear programme.

\subsection{Exact conformal transformation of the scalar--tensor pair}
\label{subsec:V154-conformal}

Write the realified Higgs components as \(h^A\), with positive fibre metric
\(G^0_{AB}\), and normalize the kinetic term as
\begin{equation}
 -\frac12\int_M\sqrt{|g|}\,
 G^0_{AB}D_\mu h^A D^\mu h^B.
 \label{eq:V154-real-Higgs-kinetic}
\end{equation}
Fix a positive constant \(\chi_*>0\) and, on the region \(F>0\), define
\begin{equation}
 \Omega^2:=\frac{F}{\chi_*},
 \qquad
 \widetilde g_{\mu\nu}:=\Omega^2g_{\mu\nu}.
 \label{eq:V154-Einstein-frame}
\end{equation}

\begin{theorem}[Einstein-frame action identity including the face]
\label{thm:V154-frame-identity}
Up to corner terms and the usual transformation of noncurvature potentials,
\begin{align}
 &\int_M\sqrt{|g|}\,F R(g)
 +2\int_\Sigma\sqrt{|\gamma|}\,F K(g)
 -\frac12\int_M\sqrt{|g|}\,
  G^0_{AB}D_\mu h^A D^\mu h^B
 \notag\\
 &\quad=
 \chi_*\int_M\sqrt{|\widetilde g|}\,R(\widetilde g)
 +2\chi_*\int_\Sigma\sqrt{|\widetilde\gamma|}\,
            \widetilde K
 \notag\\
 &\qquadquad
 -\frac12\int_M\sqrt{|\widetilde g|}\,
 \mathcal G_{AB}(h)
 \widetilde g^{\mu\nu}D_\mu h^A D_\nu h^B,
 \label{eq:V154-frame-action}
\end{align}
where
\begin{equation}
 \boxed{
 \mathcal G_{AB}
 =\frac{\chi_*}{F}G^0_{AB}
  +\frac{3\chi_*}{F^2}
    (\partial_AF)(\partial_BF).}
 \label{eq:V154-target-metric}
\end{equation}
The coefficient three in the second term corresponds to the
\(-\tfrac12\mathcal G_{AB}\partial h^A\partial h^B\) normalization.
\end{theorem}

\begin{proof}
In four dimensions, for \(\widetilde g=\Omega^2g\),
\begin{align*}
 \sqrt{|g|}&=\Omega^{-4}\sqrt{|\widetilde g|},\\
 R(g)&=\Omega^2\left(
 R(\widetilde g)+6\widetilde\Box\log\Omega
 -6|\widetilde\nabla\log\Omega|_{\widetilde g}^2
 \right).
\end{align*}
Since \(F=\chi_*\Omega^2\), the weighted curvature density becomes
\[
 \chi_*\sqrt{|\widetilde g|}\left(
 R(\widetilde g)+6\widetilde\Box\log\Omega
 -6|\widetilde\nabla\log\Omega|^2
 \right).
\]
The induced metric and extrinsic curvature transform as
\begin{align*}
 \widetilde\gamma_{ab}&=\Omega^2\gamma_{ab},\\
 \widetilde K&=\Omega^{-1}
 (K+3\nabla_n\log\Omega).
\end{align*}
Consequently the difference between
\(2\int\sqrt{|\gamma|}FK\) and
\(2\chi_*\int\sqrt{|\widetilde\gamma|}\widetilde K\)
is exactly the boundary term produced by integrating
\(6\chi_*\widetilde\Box\log\Omega\).  They cancel with the fixed outward
normal convention.

The remaining conformal derivative is
\[
 -6\chi_*|\widetilde\nabla\log\Omega|^2
 =-\frac{3\chi_*}{2F^2}|\widetilde\nabla F|^2.
\]
Writing \(\partial_\mu F=(\partial_AF)D_\mu h^A\), this is
\(-\tfrac12(3\chi_*/F^2)\partial_AF\partial_BF
D h^AD h^B\).  Finally,
\[
 \sqrt{|g|}g^{\mu\nu}
 =\frac{\chi_*}{F}sqrt{|\widetilde g|},
   \widetilde g^{\mu\nu},
\]
which supplies the first term of \(\mathcal G\).  Adding the two kinetic
rows proves \eqref{eq:V154-frame-action}.
\end{proof}

\begin{corollary}[Positive scalar principal metric]
\label{cor:V154-target-positive}
If
\begin{equation}
 F\ge c_*>0,
 \qquad \chi_*>0,
 \label{eq:V154-F-positive}
\end{equation}
then \(\mathcal G\) is positive definite and
\begin{equation}
 \mathcal G_{AB}v^Av^B
 \ge\frac{\chi_*}{\sup F}
 G^0_{AB}v^Av^B
 \label{eq:V154-target-lower}
\end{equation}
whenever \(F\) also has a finite upper bound.  Thus the conformal change
does not introduce a scalar ghost on this branch.
\end{corollary}

\begin{proof}
The first term in \eqref{eq:V154-target-metric} is a positive multiple of
\(G^0\), and the second is a nonnegative rank-one quadratic form.  The
displayed lower bound follows directly.
\end{proof}

\begin{proposition}[Principal-symbol diagonalization]
\label{prop:V154-principal-diagonal}
On a smooth background satisfying \eqref{eq:V154-F-positive}, impose
background-harmonic gauge for \(\widetilde g\) and a standard covariant gauge
for the internal connection.  The metric--Higgs second-order principal
symbol of \eqref{eq:V154-frame-action} is block diagonal:
\begin{equation}
 \sigma_2(\mathcal P)(x,k)
 =-\widetilde g^{\alpha\beta}k_\alpha k_\beta
 \begin{pmatrix}
  \chi_*\mathbb G_{\rm DW}&0\\
  0&\mathcal G_{AB}
 \end{pmatrix}.
 \label{eq:V154-principal-symbol}
\end{equation}
Here \(\mathbb G_{\rm DW}\) is the nondegenerate DeWitt fibre form.  In
particular, no mixed second derivative of \(\widetilde g\) and \(h\) remains.
\end{proposition}

\begin{proof}
The first line on the right of \eqref{eq:V154-frame-action} is the ordinary
constant-coefficient Einstein--Hilbert plus GHY action.  Its harmonic-gauge
principal Hessian is the DeWitt wave block derived in V36.  The last line is
a first-derivative sigma-model action.  Varying it in \(h\) gives
\(\mathcal G_{AB}\widetilde\Box h^B\) at principal order; varying it in the
metric gives only first derivatives of \(h\) in the metric source and no
second derivative of \(h\) in the metric principal operator.  The curvature
coefficient is constant, so the metric equation contains no Hessian of
\(F\).  This proves the block symbol.
\end{proof}

\begin{firewall}[Frame equivalence needs uniform control]
\label{fire:V154-frame-uniformity}
Pointwise positivity of each \(F_q\) is insufficient for a cofinal change of
variables.  Uniform equivalence of Jordan- and Einstein-frame Sobolev norms
requires lower and upper bounds for \(F_q\) and derivative bounds for
\(F_q\) at the order being estimated.  If \(\inf F_q\downarrow0\), the
conformal map degenerates even though every finite-cutoff action can still be
rewritten formally.
\end{firewall}

\subsection{Hyperbolic evolution does not smooth curvature}
\label{subsec:V154-no-smoothing}

\begin{proposition}[Exact high-frequency wave obstruction]
\label{prop:V154-wave-obstruction}
On \([0,T]\times\mathbb T^3\), let
\begin{equation}
 u_q(t,x):=q^{-2}\sin(qt)\sin(qx_1).
 \label{eq:V154-wave-packet}
\end{equation}
Then
\begin{equation}
 (\partial_t^2-\Delta)u_q=0,
 \qquad
 u_q\to0\text{ strongly in }W^{1,\infty},
 \label{eq:V154-wave-low-convergence}
\end{equation}
but \(\partial_t^2u_q\) and \(\partial_1^2u_q\) have nonzero constant
\(L^2\) size and do not converge strongly to zero.  Thus zero reduced source
and strong one-derivative convergence do not imply strong curvature
convergence.
\end{proposition}

\begin{proof}
Both second derivatives in the wave operator equal
\(-\sin(qt)\sin(qx_1)\), so their difference is zero.  The function and its
first derivatives have sup norms bounded by \(q^{-2}\) and \(q^{-1}\),
respectively.  The second derivatives have amplitude one.  Their squared
\(L^2\) norms converge to the positive product of the average values of two
\(\sin^2\) factors, while the functions converge weakly to zero by
oscillation.  Strong \(L^2\) convergence is therefore impossible.
\end{proof}

\begin{corollary}[Required initial topology]
\label{cor:V154-curvature-data}
A hyperbolic compactness theorem intended to control curvature must control
the corresponding curvature-level initial energy.  Convergence of metric
data only in \(H^1\times L^2\), even with a vanishing source, cannot provide
that conclusion.
\end{corollary}

\subsection{A curvature-level reduced-Einstein stability theorem}
\label{subsec:V154-stability}

Let \(s\ge4\) be an integer and let
\(\widetilde g_q\) solve a background-harmonic reduced system on a fixed
slab \([0,T]\times\Omega\subset\mathbb R^{1+3}\):
\begin{equation}
 \widetilde g_q^{\alpha\beta}
 \partial_\alpha\partial_\beta\widetilde g_q
 =\mathcal N(\widetilde g_q,\partial\widetilde g_q)
  +\mathcal S_q.
 \label{eq:V154-reduced-system}
\end{equation}
The tensor \(\mathcal S_q\) is the complete assembled Einstein-frame
source, including gauge fixing and all Standard Model/BV contributions.

\begin{theorem}[Strong curvature stability on a common hyperbolic slab]
\label{thm:V154-strong-stability}
Assume:
\begin{enumerate}[label=\textup{(\roman*)},leftmargin=2.8em]
\item the metrics are uniformly Lorentzian with a common time function and
      \[
       \sup_q\bigl(
       \|\widetilde g_q\|_{L^\infty_tH^s_x}
       +\|\partial_t\widetilde g_q\|_{L^\infty_tH^{s-1}_x}
       \bigr)<\infty;
      \]
\item their initial data converge strongly in
      \(H^s(\Omega)\times H^{s-1}(\Omega)\);
\item \(\mathcal S_q\to\mathcal S\) strongly in
      \(L^1([0,T];H^{s-1}(\Omega))\);
\item the boundary conditions have a common constraint-preserving,
      maximally dissipative realization, and the difference boundary flux is
      bounded by its \(H^s\) energy plus a remainder tending to zero.
\end{enumerate}
Then the reduced solutions converge strongly,
\begin{align}
 \widetilde g_q&\to\widetilde g
 &&\text{in }C([0,T];H^s),
 \label{eq:V154-g-strong-Hs}\\
 \partial_t\widetilde g_q&\to\partial_t\widetilde g
 &&\text{in }C([0,T];H^{s-1}),
 \label{eq:V154-dtg-strong}\\
 \operatorname{Riem}(\widetilde g_q)&\to
 \operatorname{Riem}(\widetilde g)
 &&\text{in }C([0,T];H^{s-2}).
 \label{eq:V154-curvature-strong}
\end{align}
\end{theorem}

\begin{proof}
Put \(w_q=\widetilde g_q-\widetilde g\), subtract the two reduced systems,
and place the principal operator with coefficient \(\widetilde g_q\) on the
left.  The right side is the sum of
\(\mathcal S_q-\mathcal S\), a coefficient difference multiplying fixed
second derivatives of \(\widetilde g\), and the difference of
\(\mathcal N\).  Since \(s\ge4>5/2\), the three-dimensional Sobolev--Moser
estimates make the coefficient and nonlinear maps locally Lipschitz from
\(H^s\) to \(H^{s-1}\) on the uniformly hyperbolic bounded set.

The standard differentiated wave energy \(E_s(w_q;t)\) therefore satisfies
\begin{align}
 \frac{d}{dt}E_s(w_q;t)
 \le C E_s(w_q;t)
 +C\|\mathcal S_q-\mathcal S\|_{H^{s-1}}^2
 +\varepsilon_q(t),
 \label{eq:V154-energy-difference}
\end{align}
where the common boundary flux hypothesis gives
\(\int_0^T|\varepsilon_q(t)|dt\to0\).  The initial energy tends to zero.
Gronwall's inequality and the strong \(L^1_tH^{s-1}_x\) source convergence
(using its boundedness to pass to the squared form, or the equivalent
integrated energy inequality) give
\(\sup_tE_s(w_q;t)\to0\).  This proves
\eqref{eq:V154-g-strong-Hs}--\eqref{eq:V154-dtg-strong}.

Curvature is a smooth expression in
\((\widetilde g^{-1},\partial\widetilde g,partial^2\widetilde g)\).
Sobolev--Moser continuity at \(s\ge4\) then gives
\eqref{eq:V154-curvature-strong}.
\end{proof}

\begin{remark}[The strong source is the unresolved input]
\label{rem:V154-source-gate}
Theorem~\ref{thm:V154-strong-stability} is a continuation theorem, not a
construction of \(\mathcal S_q\).  V148--V153 show why strong convergence of
the complete interacting BV source is difficult: sectorwise determinant
limits and weak Higgs energy bounds do not imply hypothesis \textup{(iii)}.
\end{remark}

\subsection{Harmonic constraints are forced by the Ward residual}
\label{subsec:V154-constraint}

Let \(C_{q,\mu}\) be the background-harmonic gauge covector and define the
reduced Einstein tensor by
\begin{equation}
 \widehat G_{q,\mu\nu}
 :=G_{q,\mu\nu}-\nabla^{q}_{(\mu}C_{q,\nu)}
 +\frac12g_{q,\mu\nu}\nabla_q^\alpha C_{q,\alpha}.
 \label{eq:V154-reduced-Einstein}
\end{equation}
Suppose the reduced equation is
\begin{equation}
 \widehat G_{q,\mu\nu}=\kappa_qT_{q,\mu\nu}
 \label{eq:V154-reduced-equation}
\end{equation}
and define the complete assembled Ward residual by
\begin{equation}
 \mathcal R_{q,\nu}:=\nabla_q^\mu T_{q,\mu\nu}.
 \label{eq:V154-Ward-residual}
\end{equation}

\begin{theorem}[Exact subsidiary equation]
\label{thm:V154-subsidiary-equation}
With the curvature commutator convention of V150, equations
\eqref{eq:V154-reduced-Einstein}--\eqref{eq:V154-Ward-residual} imply
\begin{equation}
 \boxed{
 \Box_{g_q}C_{q,\nu}
 +R_{q,\nu}{}^\mu C_{q,\mu}
 =-2\kappa_q\mathcal R_{q,\nu}.}
 \label{eq:V154-subsidiary-equation}
\end{equation}
Thus the source of the harmonic constraint is the complete Ward residual,
including the V150 BV antifield row on a nonstationary background.
\end{theorem}

\begin{proof}
Take the covariant divergence of
\eqref{eq:V154-reduced-equation}.  The contracted Bianchi identity removes
\(\nabla^\mu G_{\mu\nu}\).  For a one-form \(C\), the curvature commutator
gives
\[
 2\nabla^\mu\nabla_{(\mu}C_{\nu)}
 -\nabla_\nu\nabla^\alpha C_\alpha
 =\Box C_\nu+R_\nu{}^\mu C_\mu.
\]
Remembering the two signs in
\eqref{eq:V154-reduced-Einstein}, the divergence equation is precisely
\eqref{eq:V154-subsidiary-equation}.
\end{proof}

\begin{corollary}[Constraint propagation and cofinal stability]
\label{cor:V154-constraint-propagation}
If \(\mathcal R_q=0\), the initial harmonic constraint and its normal
derivative vanish, and the boundary realization has homogeneous
constraint-preserving data, then \(C_q=0\) throughout the slab.  More
generally, under uniform subsidiary energy estimates,
\begin{equation}
 \|(C_q,\nabla C_q)(t)\|_{\mathcal H}
 \le C\left(
 \|(C_q,\nabla C_q)(0)\|_{\mathcal H}
 +\|\mathcal R_q\|_{L^1_tL^2_x}
 +\|b_q\|_{L^1_t\mathcal B}
 \right),
 \label{eq:V154-constraint-estimate}
\end{equation}
where \(b_q\) is the boundary constraint defect.  Vanishing of the three
right-hand terms implies \(C_q\to0\) in the subsidiary energy norm.
\end{corollary}

\begin{proof}
Apply the homogeneous uniqueness theorem, or the inhomogeneous wave energy
estimate, to \eqref{eq:V154-subsidiary-equation}.  Uniform bounds on the
curvature potential and boundary flux give the constant in
\eqref{eq:V154-constraint-estimate}.
\end{proof}

\begin{firewall}[Reduced solutions are not automatically Einstein
solutions]
\label{fire:V154-reduced-not-full}
Solving \eqref{eq:V154-reduced-system} while estimating determinant and BV
source rows separately does not prove \(\mathcal R_q\to0\).  Without the
assembled V150 Ward estimate, equation
\eqref{eq:V154-subsidiary-equation} permits a nonzero harmonic-constraint
source, so a limit of reduced solutions need not solve the full Einstein
equation.
\end{firewall}

\subsection{V154 ledger and next acquisition}
\label{subsec:V154-ledger}

\paragraph{New exact conclusions.}
\begin{enumerate}[label=\textup{(\arabic*)},leftmargin=2.8em]
\item The complete \(FR+2FK\) pair transforms to constant
      Einstein--Hilbert plus GHY form, with the derivative of \(F\) absorbed
      into the Higgs target metric.
\item The Einstein-frame target metric is positive wherever the effective
      Planck coefficient has a positive lower bound.
\item Harmonic and internal gauge fixing give a block-diagonal
      metric--Higgs principal symbol.
\item An exact homogeneous wave sequence proves that hyperbolic evolution
      supplies no curvature compactification from lower-order convergence.
\item Strong curvature convergence follows from curvature-level initial
      data, a strong complete source, uniform hyperbolicity, and a common
      dissipative boundary estimate.
\item The harmonic constraint obeys an exact wave equation forced by the
      assembled Ward residual.
\end{enumerate}

\paragraph{Next forced step and scheduled audit.}
V155 is the next compulsory NS/YM audit.  After that read-only inspection it
should formulate the complete coupled Einstein-frame continuation norm for
the gauge, Higgs, fermion, ghost, antifield, metric, and boundary port
variables.  It must distinguish classical Sobolev solution norms from
operator-valued quantum form norms and identify the precise bridge still
needed between the BV pushforward of V147/V148 and the strong source
hypothesis of Theorem~\ref{thm:V154-strong-stability}.

\begin{assessment}[V154 continuum status]
\label{ass:V154-continuum-status}
The full Standard--Model--Einstein continuum remains unproved.  V154 removes
the Jordan-frame principal mixing on the uniformly positive branch and
proves exact metric stability and constraint-propagation criteria.  The
programme still lacks the uniform effective Planck bound, strong complete BV
source convergence, compatible nonlinear boundary estimate, and cofinal
construction of solutions satisfying those criteria.
\end{assessment}

\section*{V154 exact checks and append-only record}
\addcontentsline{toc}{section}{V154 exact checks and append-only record}

The principal and subsidiary checksum is
\begin{equation}
 \boxed{
 FR+2FK
 \xrightarrow{\ \widetilde g=(F/\chi_*)g\ }
 \chi_*\widetilde R+2\chi_*\widetilde K
 -\frac12\frac{3\chi_*}{F^2}|dF|^2,
 \qquad
 \Box_gC_\nu+R_\nu{}^\mu C_\mu=-2\kappa\mathcal R_\nu.}
 \label{eq:V154-final-check}
\end{equation}
The append operation preserves the complete V153 prefix byte for byte before
its former live terminator.  The assembled V154 file is checked for one live
final terminator, balanced environments and braces, duplicate labels, newly
introduced unresolved references, display math delimiters, and control
characters before the exact V153 predecessor is removed.  No companion file
is modified.

% =============================================================================
% END APPENDED VERSION V154 MATERIAL
% =============================================================================
\clearpage
% =============================================================================
% BEGIN APPENDED VERSION V155 MATERIAL
% =============================================================================

\section{V155: compulsory companion audit, the coupled continuation norm,
and a shell-summability bridge from BV forms to a semiclassical source}
\label{sec:V155}

\subsection{Purpose}
\label{subsec:V155-purpose}

V154 gives a strong Einstein-frame continuation theorem if the complete
source converges in a curvature-level Sobolev norm.  V90 and V147, by
contrast, place the renormalized quantum stress in a common rigged form
space and construct finite BV pushforwards.  These are different analytic
objects.  An operator-valued distribution estimate does not become a
classical \(H^{s-1}\) source estimate merely by changing notation.

This version first performs the scheduled companion audit.  It then defines
the classical and quantum continuation spaces separately and proves a
bridge theorem.  The bridge has two additional inputs: a compatible family
of states, and an absolutely summable estimate for the complete metric
derivative of each BV shell pushforward in a state-evaluated Sobolev norm.
Those hypotheses are strong, but they are typed correctly and telescope
without a cutoff multiplicity.

\subsection{The V155 NS/YM audit}
\label{subsec:V155-audit}

The shared folder was inspected read-only at
\(2026\text{-}09\text{-}05\ 10{:}47{:}15\ {-}07{:}00\).  The newest
companion snapshots were
\begin{center}
\begin{tabular}{@{}p{0.49\linewidth}r p{0.31\linewidth}@{}}
\toprule
file & bytes & SHA--256\\
\midrule
\texttt{Renewal\_NS\_Stability\_Research\_Notes\_V31.tex}
&887537&\texttt{f1121b62238ad5491bb7cf03635ea9934942edf573ed8faaa9ee79e1d38a6696}\\
\texttt{Renewal\_Yang\_Mills\_Mass\_Gap\_Research\_Notes\_V133.tex}
&3051393&\texttt{c7cbbf6894a94fd856df7e71150e31473ef39df9587fd9dd06af7af418c85d5a}\\
\bottomrule
\end{tabular}
\end{center}

The NS note is unchanged.  Its relevant conclusion remains that a signed,
restricted formation correlation must be estimated before it is enlarged to
a global positive stock.

YM V132 constructed an exact connected forest/Moebius representation with
contractive unmarked factors.  YM V133 then found a decisive obstruction to
the proposed analytic estimate: the number of differentiated graph edges
can exceed the merger rank.  A parallel edge may supply no new fluctuation
factor, and an individual forest integrand can contain lower-rank moments.
Those terms cancel only after the complete forest/Moebius sum.  YM V133
therefore retracts the termwise three-link estimate while retaining the exact
connected identity.

\begin{proposition}[Transferable derivative-fidelity rule]
\label{prop:V155-derivative-fidelity}
A formal derivative, diagram edge, determinant factor, or field-sector label
does not by itself count a small shell fluctuation in the Standard Model
source.  A cofinal estimate may charge a decay factor only to an actual
difference insertion in the complete BV shell first variation, or to a
previously proved connected cumulant estimate.  All lower-rank terms must be
assembled before a norm is taken.
\end{proposition}

\begin{proof}
YM V133 supplies an explicit algebraic counterexample in which two
differentiated parallel bridges collapse to one covariance, so the left side
has order \(s\) while the proposed product of two bridge weights has order
\(s^2\).  The same logical implication would be used if one assigned an
independent small factor to a metric derivative, a gauge block, and a Higgs
block before forming the logarithmic shell derivative.  Since those labels
can act on the same carrier and share one fluctuation, no orthogonal range or
independence theorem justifies their product.  The only invariant object is
the derivative of the complete logarithmic fibre integral.
\end{proof}

\subsection{The classical Einstein-frame continuation space}
\label{subsec:V155-classical-space}

Fix an integer \(s\ge4\) and a slab \(M_T=[0,T]\times\Omega\).  Let
\(U\) denote the gauge-fixed Einstein-frame classical variables
\begin{equation}
 U=(\widetilde g,A,H,\Psi;\ c,\bar c,b),
 \label{eq:V155-classical-variables}
\end{equation}
where the last three entries are graded auxiliary variables when the
classical BV complex is retained.  Antifields are sources in the continuous
dual and are not counted as additional physical Cauchy data.

Define the bulk evolution norm
\begin{align}
 \|U\|_{\mathcal X_T^s}
 :={}&\sup_{0\le t\le T}\Bigl[
 \|(\widetilde g,A,H,c,\bar c,b)(t)\|_{H^s_x}
 \notag\\
 &\qquad
 +\|\partial_t(\widetilde g,A,H,c,\bar c,b)(t)\|_{H^{s-1}_x}
 +\|\Psi(t)\|_{H^s_x}\Bigr]
 \notag\\
 &+\|\mathcal S(U)\|_{L^1_tH^{s-1}_x}
 +\|\mathcal R_{\rm Ward}(U)\|_{L^1_tH^{s-1}_x}
 +\|\mathcal B(U)\|_{\mathcal Y^s_\Sigma}.
 \label{eq:V155-classical-norm}
\end{align}
Here \(\mathcal S\) is the complete reduced Einstein-frame source,
\(\mathcal R_{\rm Ward}\) is the assembled gauge/diffeomorphism Ward
residual, and \(\mathcal B(U)\) is the boundary trace/conormal row in the
common constraint-preserving port space \(\mathcal Y^s_\Sigma\).  That
boundary space is required to control the maximally dissipative flux in
Theorem~\ref{thm:V154-strong-stability}; formula
\eqref{eq:V155-classical-norm} does not assert that the nonlinear renewal
boundary already has such a realization.

\begin{definition}[Admissible classical continuation set]
\label{def:V155-admissible-classical}
For constants \(C,c_*,C_F>0\), let \(\mathfrak C_T^s(C,c_*,C_F)\) be the
set of configurations satisfying
\begin{equation}
 \|U\|_{\mathcal X_T^s}\le C,
 \qquad
 c_*\le F(H)\le C_F,
 \qquad
 \|F(H)\|_{W^{s,\infty}_{\rm adapted}}\le C_F,
 \label{eq:V155-continuation-set}
\end{equation}
together with uniform Lorentzian hyperbolicity, internal gauge fixing, and
the corner compatibility conditions through order \(s-1\).  The adapted
last norm means precisely the derivatives needed for uniform equivalence of
the Jordan- and Einstein-frame Sobolev norms; it may be replaced by any
weaker multiplier norm that proves those product estimates.
\end{definition}

\begin{proposition}[What the classical norm controls]
\label{prop:V155-classical-controls}
On an admissible continuation set, the Einstein-frame metric--Higgs
principal symbol is uniformly hyperbolic, the target metric
\(\mathcal G_{AB}\) is uniformly positive, the nonlinear Moser constants in
V154 are bounded, and the harmonic subsidiary equation has a uniform energy
estimate.
\end{proposition}

\begin{proof}
The bounds on \(F\) and its multiplier derivatives make the conformal change
\eqref{eq:V154-Einstein-frame} uniformly invertible at the stated Sobolev
order.  Corollary~\ref{cor:V154-target-positive} gives positivity of
\(\mathcal G\).  The \(H^s\) bound with \(s\ge4\) controls the coefficient
algebras and Lipschitz norms used in
\eqref{eq:V154-energy-difference}.  Uniform hyperbolicity and the boundary
port bound supply the metric and subsidiary energy inequalities.
\end{proof}

\begin{firewall}[Grassmann variables and physical data]
\label{fire:V155-graded-classical}
The Sobolev entries for \(c,\bar c\), and antifield sources are norms of
coefficients in a graded perturbative algebra.  They do not turn ghosts into
classical observable Cauchy fields.  Physical solutions are obtained on the
appropriate BV cohomology/gauge-fixed Lagrangian, while these rows retain the
identities needed for the source estimate.
\end{firewall}

\subsection{The quantum rigged-form space is different}
\label{subsec:V155-quantum-space}

At cutoff \(q\), let
\begin{equation}
 \mathscr D_q\subset\mathscr H_q\subset\mathscr D_q'
 \label{eq:V155-rigging}
\end{equation}
be the V90 common form rigging, enlarged by the dual BV pairs of V147/V148.
For a Banach test space \(Z\) on spacetime, define the form-valued
distribution norm
\begin{equation}
 \|\mathcal T_q\|_{Z'(\mathcal B_q)}
 :=\sup_{\|\varphi\|_Z\le1}
 \|\mathcal T_q(\varphi)\|_{\mathcal B(\mathscr D_q,\mathscr D_q')}.
 \label{eq:V155-form-distribution-norm}
\end{equation}
For the source topology in V154, the formal choice is
\(Z=H^{1-s}_0(M_T)\), so that \(Z'=H^{s-1}\).  Whether the renormalized
stress has this positive regularity is a substantive state-dependent
hypothesis.

\begin{proposition}[No automatic Sobolev upgrade]
\label{prop:V155-no-automatic-upgrade}
Convergence of \(\mathcal T_q(\varphi)\) in
\(\mathcal B(\mathscr D_q,\mathscr D_q')\) for each smooth test
\(\varphi\) does not imply boundedness, Cauchy convergence, or even
membership in
\(H^{s-1}(M_T;\mathcal B(\mathscr D,\mathscr D'))\).
\end{proposition}

\begin{proof}
Scalar distributions already disprove the implication.  The sequence
\(e^{iqx_1}\) tends to zero against each fixed smooth test after integration
by parts, but its positive Sobolev norms grow like \(q^{s-1}\).  Multiplying
that sequence by one fixed nonzero bounded form gives the same counterexample
in the form-valued setting.
\end{proof}

\subsection{State evaluation and common transport}
\label{subsec:V155-state-evaluation}

Let \(J_{q'q}:\mathscr D_q\to\mathscr D_{q'}\) be the exact low-sector
transport furnished by the determinant/BV comparison system.  It is assumed
to respect the form duality and to satisfy
\begin{equation}
 J_{q''q'}J_{q'q}=J_{q''q}.
 \label{eq:V155-J-cocycle}
\end{equation}
Transport every form back to one reference rigging before subtracting it.

Let \(\omega_q\) be a state functional continuous on the transported form
space, with
\begin{equation}
 |\omega_q(B)|\le C_\omega\|B\|_{\mathcal B(\mathscr D,\mathscr D')}
 \label{eq:V155-state-bound}
\end{equation}
uniformly.  Define the state-evaluated renormalized source by
\begin{equation}
 t_q(\varphi):=\omega_q(\mathcal T_q^{\rm ren}(\varphi)).
 \label{eq:V155-state-source}
\end{equation}

\begin{lemma}[Form convergence plus state convergence]
\label{lem:V155-form-state-convergence}
Suppose, after common transport,
\begin{align}
 \|\mathcal T_{q'}-\mathcal T_q\|_{Z'(\mathcal B)}&\to0,
 \label{eq:V155-form-Cauchy}\\
 \|\omega_{q'}-\omega_q\|_{\mathcal B^*}&\to0,
 \label{eq:V155-state-Cauchy}
\end{align}
and the forms are uniformly bounded in \(Z'(\mathcal B)\).  Then
\(t_q\) is Cauchy in \(Z'\).  In particular, for
\(Z=H^{1-s}_0\), it converges strongly in \(H^{s-1}\).
\end{lemma}

\begin{proof}
For \(\|\varphi\|_Z\le1\), add and subtract
\(\omega_{q'}(\mathcal T_q(\varphi))\):
\begin{align*}
 |t_{q'}(\varphi)-t_q(\varphi)|
 &\le C_\omega
 \|\mathcal T_{q'}-\mathcal T_q\|_{Z'(\mathcal B)}\\
 &\quad+|\omega_{q'}-\omega_q\|_{\mathcal B^*}
       \|\mathcal T_q\|_{Z'(\mathcal B)}.
\end{align*}
Take the supremum over the unit ball of \(Z\).  Both terms tend to zero by
hypothesis.
\end{proof}

\begin{firewall}[The bridge is semiclassical state data]
\label{fire:V155-semiclassical-bridge}
Equation~\eqref{eq:V155-state-source} selects a c-number expectation source
for a semiclassical Einstein equation.  It is not an operator equation for a
quantized metric.  A full quantum-gravity construction would require a
different common state space and is not supplied by the Grand Tensor or the
present notes.
\end{firewall}

\subsection{Complete BV shell derivatives}
\label{subsec:V155-shell-derivative}

Let \(q_0<q_1<\cdots\) be cofinal cutoffs.  Write the exact BV pushforward
over the shell \((q_n,q_{n+1}]\) as
\begin{equation}
 e^{-\Gamma_{n+1}/\hbar}
 =\int_{\mathcal L_{n+1,n}}
   e^{-S_{n+1}/\hbar},
 \label{eq:V155-shell-pushforward}
\end{equation}
after transport to the common low variables.  Define the complete shell
metric derivative
\begin{equation}
 \mathfrak s_n
 :=\delta_{\widetilde g}\Gamma_{n+1}
   -\delta_{\widetilde g}\Gamma_n.
 \label{eq:V155-shell-source}
\end{equation}
It contains every determinant, counterterm, Higgs-improvement, ghost,
antifield-source, and interaction contribution.

\begin{lemma}[First derivative of one logarithmic shell integral]
\label{lem:V155-log-derivative}
Where differentiation under the finite-dimensional integral is valid,
\begin{equation}
 \delta_{\widetilde g}\Gamma_{n+1}
 =\left\langle
   \delta_{\widetilde g}S_{n+1}
  \right\rangle_{n+1,n}.
 \label{eq:V155-log-derivative}
\end{equation}
For two variations,
\begin{align}
 \delta_1\delta_2\Gamma_{n+1}
 ={}&\langle\delta_1\delta_2S_{n+1}\rangle
 \notag\\
 &-\hbar^{-1}\left(
 \langle\delta_1S_{n+1}\,\delta_2S_{n+1}\rangle
 -\langle\delta_1S_{n+1}\rangle
  \langle\delta_2S_{n+1}\rangle\right).
 \label{eq:V155-second-log-derivative}
\end{align}
Thus connected subtraction is part of the logarithmic derivative before any
diagram norm is taken.
\end{lemma}

\begin{proof}
Differentiate \(-\hbar\log Z\), where
\(Z=\int e^{-S/\hbar}\).  The first derivative is
\(-\hbar Z^{-1}\delta Z=\langle\delta S\rangle\).  Differentiate that
normalized expectation once more.  The derivative of its numerator gives
the first and first-product terms in
\eqref{eq:V155-second-log-derivative}; the derivative of \(Z^{-1}\) gives
the final product of expectations with the opposite sign.
\end{proof}

\begin{assessment}[Why the YM V133 warning applies exactly]
\label{ass:V155-YM-transfer}
Expanding the covariance in \eqref{eq:V155-second-log-derivative} into Wick
or renewal diagrams can create several formal derivative edges which depend
on the same shell fluctuation.  YM V133 proves that edge count need not equal
fluctuation rank.  V155 therefore assigns no decay exponent diagram by
diagram; the acquisition target is a norm of \(\mathfrak s_n\) after the
complete connected subtraction.
\end{assessment}

\subsection{The shell-summability bridge}
\label{subsec:V155-bridge}

\begin{theorem}[Absolute shell variation gives a strong source limit]
\label{thm:V155-shell-bridge}
Let \(Z=H^{1-s}_0(M_T)\).  Assume the exact transport cocycle
\eqref{eq:V155-J-cocycle}, the BV pushforward composition of V147, and
\begin{equation}
 \sum_{n=0}^\infty
 \|\mathfrak s_n\|_{Z'(\mathcal B)}<\infty.
 \label{eq:V155-shell-summability}
\end{equation}
Then the transported complete metric derivatives
\(\delta_{\widetilde g}\Gamma_n\) converge in
\(H^{s-1}(M_T;\mathcal B(\mathscr D,\mathscr D'))\).

If, in addition, the transported states obey
\begin{equation}
 \sum_{n=0}^\infty
 \|\omega_{n+1}-\omega_n\|_{\mathcal B^*}<\infty
 \label{eq:V155-state-summability}
\end{equation}
and are uniformly bounded as in \eqref{eq:V155-state-bound}, then their
state-evaluated sources converge strongly in \(H^{s-1}(M_T)\).
\end{theorem}

\begin{proof}
For \(N>M\), the exact definition
\eqref{eq:V155-shell-source} telescopes:
\begin{equation}
 \delta_{\widetilde g}\Gamma_N
 -\delta_{\widetilde g}\Gamma_M
 =\sum_{n=M}^{N-1}\mathfrak s_n.
 \label{eq:V155-source-telescope}
\end{equation}
The triangle inequality applied after the complete shell derivative has
been formed gives
\[
 \|\delta\Gamma_N-\delta\Gamma_M\|_{Z'(\mathcal B)}
 \le\sum_{n=M}^{N-1}\|\mathfrak s_n\|_{Z'(\mathcal B)}.
\]
The tail tends to zero by
\eqref{eq:V155-shell-summability}; completeness gives the form-valued
source limit.  The same telescoping argument gives the state Cauchy property
from \eqref{eq:V155-state-summability}.  Apply
Lemma~\ref{lem:V155-form-state-convergence}.
\end{proof}

\begin{corollary}[No cutoff-count factor]
\label{cor:V155-no-cutoff-factor}
The estimate in Theorem~\ref{thm:V155-shell-bridge} depends on the tail
\(\sum_{n\ge M}\|\mathfrak s_n\|\), not on the number of intervening
cutoffs.  Exact BV composition and the direct telescope remove any separate
ancestry or refinement multiplicity.
\end{corollary}

\begin{proof}
Equation~\eqref{eq:V155-source-telescope} contains each adjacent shell once.
There is no second sum over paths because
\eqref{eq:V155-J-cocycle} identifies every composite transport with the
direct one.
\end{proof}

\begin{theorem}[Conditional coupled continuation]
\label{thm:V155-conditional-continuation}
Suppose a cofinal family of finite-cutoff solutions lies in one admissible
set \(\mathfrak C_T^s(C,c_*,C_F)\), its initial and boundary data are Cauchy
in the V154 topologies, the shell hypotheses
\eqref{eq:V155-shell-summability}--\eqref{eq:V155-state-summability} hold for
the complete source, and its assembled Ward and boundary defects tend to
zero in the corresponding strong norms.  Then:
\begin{enumerate}[label=\textup{(\roman*)},leftmargin=2.8em]
\item the state-evaluated Einstein-frame sources converge strongly in
      \(L^1_tH^{s-1}_x\);
\item the metrics converge strongly at curvature level by
      Theorem~\ref{thm:V154-strong-stability};
\item the harmonic constraints tend to zero by
      Corollary~\ref{cor:V154-constraint-propagation}; and
\item every strongly convergent matter field solves the limiting
      gauge-fixed Euler equation, while the exact limiting Ward identity
      identifies the full Einstein rather than merely reduced system.
\end{enumerate}
\end{theorem}

\begin{proof}
Theorem~\ref{thm:V155-shell-bridge} gives item \textup{(i)}.  The admissible
set supplies the uniform hyperbolicity, coefficient, and boundary constants
required in V154, so its stability theorem gives item \textup{(ii)}.  The
assembled Ward residual is the source in
\eqref{eq:V154-subsidiary-equation}; the assumed initial and boundary
compatibility gives item \textup{(iii)}.  Strong convergence permits passage
through the classical nonlinearities by Sobolev--Moser continuity.  The
limiting master/Ward identity removes gauge-reduction defects, proving item
\textup{(iv)}.
\end{proof}

\begin{firewall}[Conditional theorem, not acquired estimate]
\label{fire:V155-bridge-open}
Neither the Grand Tensor source nor V1--V154 proves
\eqref{eq:V155-shell-summability} in a positive Sobolev source norm.  Exact
BV composition supplies the telescope, but it does not bound its summands.
Likewise, compatible state transport is extra data and cannot be inferred
from determinant-line phase transport alone.
\end{firewall}

\subsection{V155 ledger and next acquisition}
\label{subsec:V155-ledger}

\paragraph{New exact conclusions.}
\begin{enumerate}[label=\textup{(\arabic*)},leftmargin=2.8em]
\item The compulsory audit records unchanged NS V31 and current YM V133.
\item YM V133's derivative-rank counterexample forbids charging formal
      diagram edges as independent small shell factors.
\item The classical Sobolev continuation space and the quantum rigged-form
      space are now explicitly separated.
\item State evaluation is identified as the additional semiclassical map
      from an operator-valued stress to a c-number Einstein source.
\item Form convergence plus compatible state convergence gives strong
      convergence in the declared source topology.
\item The metric derivative of a logarithmic shell integral automatically
      retains its connected subtraction.
\item Absolute summability of complete BV shell derivatives yields an exact
      cutoff-multiplicity-free source telescope.
\item Combining that bridge with V154 gives a conditional coupled
      continuation theorem.
\end{enumerate}

\paragraph{Next forced step.}
V156 should attack \eqref{eq:V155-shell-summability} rather than restate it.
At the Gaussian/one-loop level, subtract the complete local heat coefficients
allowed by the V129--V149 counterterm classification and estimate the
large-eigenvalue remainder of the metric derivative.  The proof must
determine the derivative order needed for an \(H^{s-1}\) source and decide
whether the resulting dyadic shell exponent is summable.  If positive
Sobolev convergence is too strong for the unsmeared stress, the note must
identify the weakest negative Sobolev topology actually obtained and return
that loss to the V154 stability theorem.

\begin{assessment}[V155 continuum status]
\label{ass:V155-continuum-status}
The full Standard--Model--Einstein continuum remains unproved.  V155 gives
a correctly typed conditional bridge from BV shell pushforwards to a
semiclassical strong source and incorporates the new YM derivative-fidelity
warning.  The summable complete shell estimate, compatible state family,
positive effective Planck bound, nonlinear boundary realization, and
interacting cofinal construction remain open.
\end{assessment}

\section*{V155 exact checks and append-only record}
\addcontentsline{toc}{section}{V155 exact checks and append-only record}

The bridge checksum is
\begin{equation}
 \boxed{
 \sum_n\|\delta_{\widetilde g}\Gamma_{n+1}
            -\delta_{\widetilde g}\Gamma_n\|_{H^{s-1}(\mathcal B)}<\infty
 \quad\Longrightarrow\quad
 \delta_{\widetilde g}\Gamma_n\text{ is Cauchy},}
 \label{eq:V155-final-check}
\end{equation}
where every summand is formed after the complete connected BV shell
variation.  The append operation preserves the complete V154 prefix byte for
byte before its former live terminator.  The assembled V155 file is checked
for one live final terminator, balanced environments and braces, duplicate
labels, newly introduced unresolved references, display math delimiters,
and control characters before the exact V154 predecessor is removed.  The
audit is read-only and no companion file is modified.

% =============================================================================
% END APPENDED VERSION V155 MATERIAL
% =============================================================================
\clearpage
% =============================================================================
% BEGIN APPENDED VERSION V156 MATERIAL
% =============================================================================

\section{V156: one-loop proper-time shell summability after complete local
subtraction, and the topology of the resulting stress limit}
\label{sec:V156}

\subsection{Purpose and scope}
\label{subsec:V156-purpose}

V155 reduced the strong semiclassical source bridge to summability of the
complete adjacent-shell metric derivatives.  This version tests that
condition for a smooth Gaussian one-loop regulator.  The result is positive
under a uniform differentiable heat-parametrix hypothesis: in four
dimensions, subtraction of every bulk and boundary heat coefficient through
weight four leaves a boundary remainder of order \(t^{1/2}\).  A dyadic
proper-time shell then has size \(O(\Lambda^{-1})\), which is summable.

The conclusion is deliberately limited.  It applies to the assembled
one-loop supertrace and its metric derivative, not to the interacting BV
cumulants of V148.  It also applies directly to a smooth proper-time cutoff;
a sharp spectral projector requires a separate local Weyl/Tauberian
remainder.  Finally, positive Sobolev convergence of the stress follows only
when the heat remainder is uniform in that positive Sobolev norm.

\subsection{The complete one-loop heat functional}
\label{subsec:V156-heat-functional}

Let \(P_{q,j}\) be the finite family of positive Laplace-type modulus
operators obtained from the common gauge--Higgs--ghost, doubled-fermion, and
auxiliary realization at cutoff \(q\).  Let \(\sigma_j\) include the exact
Gaussian, ghost, and multiplicity weights proved in V129--V136.  Define the
complete heat supertrace
\begin{equation}
 K_q(t;g):=\sum_j\sigma_j\operatorname{Tr}e^{-tP_{q,j}(g)}.
 \label{eq:V156-supertrace}
\end{equation}
The chiral determinant-line phase is not represented by
\eqref{eq:V156-supertrace}; on the anomaly-free admissible slab it is
transported by the flat parallel trivialization of V145.

For a metric variation \(h\) compactly supported away from corners, write
\begin{equation}
 \delta_gK_q(t)[h]
 =\int_M\langle\mathcal K_{q,M}(t),h\rangle,dV
 +\int_\Sigma\langle\mathcal K_{q,\Sigma}(t),h_\Sigma\rangle,dA.
 \label{eq:V156-heat-stress-density}
\end{equation}
All field sectors and their signs are summed before the densities
\(\mathcal K_{q,M},\mathcal K_{q,\Sigma}\) are named.

\begin{lemma}[Duhamel first variation]
\label{lem:V156-Duhamel-variation}
At finite cutoff, or for a trace-class heat semigroup with a differentiable
operator family,
\begin{equation}
 \delta_gK_q(t)[h]
 =-t\sum_j\sigma_j
 \operatorname{Tr}\bigl((\delta_gP_{q,j}[h])e^{-tP_{q,j}}\bigr).
 \label{eq:V156-Duhamel-variation}
\end{equation}
No commutativity of \(P\) and \(\delta P\) is assumed.
\end{lemma}

\begin{proof}
Duhamel's formula gives
\[
 \delta e^{-tP}
 =-\int_0^t e^{-(t-u)P}(\delta P)e^{-uP}\,du.
\]
Take the trace and use cyclicity to make its integrand independent of
\(u\).  Integration over \([0,t]\) gives the factor \(t\).  Sum with the
fixed weights \(\sigma_j\).
\end{proof}

In dimension four with smooth local boundary conditions, the differentiated
heat trace has the expansion
\begin{equation}
 \delta_gK_q(t)[h]
 \sim\sum_{k=0}^\infty
 t^{(k-4)/2}\,\delta_gA_{k,q}[h],
 \qquad t\downarrow0,
 \label{eq:V156-heat-expansion}
\end{equation}
where odd \(k\) are boundary-supported and the even coefficients can have
both bulk and boundary parts.  The functional \(A_{k,q}\) is the complete
weighted local coefficient; cancellations between sectors have already
occurred in it.

\begin{definition}[Complete weight-four subtraction]
\label{def:V156-complete-subtraction}
Set
\begin{equation}
 \delta_gK_q^{>4}(t)[h]
 :=\delta_gK_q(t)[h]
 -\sum_{k=0}^4t^{(k-4)/2}\delta_gA_{k,q}[h].
 \label{eq:V156-subtracted-heat}
\end{equation}
The subtraction is admissible only when the five coefficients are first
variations of one common local gauge-, BV-, and diffeomorphism-compatible
counterterm action, including all required face partners.
\end{definition}

\subsection{Uniform differentiated heat remainder}
\label{subsec:V156-remainder}

\begin{hypothesis}[Uniform heat-parametrix card]
\label{hyp:V156-parametrix}
Fix an integer \(r\ge0\).  Assume the following uniformly in \(q\):
\begin{enumerate}[label=\textup{(H\arabic*)},leftmargin=3.4em]
\item the coefficients of every \(P_{q,j}\), the metrics, and the boundary
      projectors/Robin endomorphisms lie in a bounded set of the finite
      \(C^N\) norms required for an \(H^r\) differentiated parametrix;
\item the realizations are strongly elliptic with one uniform parameter and
      one uniform boundary complementing constant;
\item zero modes are projected out or controlled by a uniform positive
      infrared shift; and
\item the common field multiplicities and supertrace weights are those of
      the fixed faithful Standard Model representation.
\end{enumerate}
\end{hypothesis}

\begin{theorem}[Uniform remainder through weight four]
\label{thm:V156-uniform-remainder}
Under Hypothesis~\ref{hyp:V156-parametrix}, there are \(t_0,C_r>0\),
independent of \(q\), such that
\begin{equation}
 \boxed{
 \|\delta_gK_q^{>4}(t)\|_{H^r(M)\oplus H^r(\Sigma)}
 \le C_r t^{1/2},
 \qquad0<t\le t_0.}
 \label{eq:V156-remainder-bound}
\end{equation}
If \(\Sigma=\varnothing\), the first omitted coefficient is \(k=6\) and
the right side improves to \(C_rt\).
\end{theorem}

\begin{proof}
The local heat-parametrix recursion constructs, for every \(N\), a finite
sum through \(k=N\) whose error solves an inhomogeneous heat equation with a
source vanishing to the corresponding order in \(t^{1/2}\).  Uniform bounds
on the coefficients and the strong ellipticity/complementing constants make
the parametrix and heat-estimate constants uniform on the bounded family in
Hypothesis~\ref{hyp:V156-parametrix}.  Differentiating the recursion in the
background metric and in the spacetime variables through the finite order
needed for \(H^r\) gives the same asymptotic powers, with a larger uniform
constant.

After removal of \(k=0,1,2,3,4\), the first possible boundary coefficient
is \(k=5\), whose power in four dimensions is
\(t^{(5-4)/2}=t^{1/2}\).  On a closed manifold the odd coefficients vanish,
so the first possible term is \(k=6\), of order \(t\).  The parametrix
remainder after one additional finite truncation is smaller and can be
absorbed into the same bounds for \(0<t\le t_0\).
\end{proof}

\begin{remark}[What is assumed rather than hidden]
\label{rem:V156-parametrix-scope}
The theorem proves the decay exponent from a standard uniform parametrix
card.  It does not prove that the cofinal renewal backgrounds satisfy the
uniform \(C^N\) and complementing bounds.  V137 supplies the finite-cutoff
self-adjoint mixed gauge--Higgs domain; uniform high-regularity control of
that domain is an additional analytic input.
\end{remark}

\subsection{Dyadic proper-time shell estimate}
\label{subsec:V156-dyadic}

Let \(\Lambda_n=2^n\Lambda_0\) and
\(\varepsilon_n=\Lambda_n^{-2}\).  Define the renormalized ultraviolet
shell derivative by
\begin{equation}
 \mathfrak s_n^{(1)}[h]
 :=-\frac{\hbar}{2}
 \int_{\varepsilon_{n+1}}^{\varepsilon_n}
 \frac{dt}{t}\,\delta_gK_q^{>4}(t)[h].
 \label{eq:V156-one-loop-shell}
\end{equation}
The overall sign and factor can be changed to the declared boson/fermion
effective-action convention; the summability statement is invariant under
that fixed scalar change.

\begin{theorem}[One-loop shell summability]
\label{thm:V156-shell-summability}
Under Hypothesis~\ref{hyp:V156-parametrix},
\begin{equation}
 \boxed{
 \|\mathfrak s_n^{(1)}\|_{H^r(M)\oplus H^r(\Sigma)}
 \le \hbar C_r'\Lambda_n^{-1},
 \qquad
 \sum_{n=0}^\infty\|\mathfrak s_n^{(1)}\|<\infty.}
 \label{eq:V156-shell-summability}
\end{equation}
On a closed manifold the bound improves to
\(C_r'\hbar\Lambda_n^{-2}\).
\end{theorem}

\begin{proof}
Use \eqref{eq:V156-remainder-bound} in
\eqref{eq:V156-one-loop-shell}:
\begin{align*}
 \|\mathfrak s_n^{(1)}\|
 &\le\frac{\hbar C_r}{2}
 \int_{\varepsilon_{n+1}}^{\varepsilon_n}t^{-1/2}\,dt\\
 &=\hbar C_r
 (\varepsilon_n^{1/2}-\varepsilon_{n+1}^{1/2})
 \le\hbar C_r\Lambda_n^{-1}.
\end{align*}
The geometric series \(\sum_n\Lambda_n^{-1}\) converges.  In the closed
case, integrate the bound \(C_rt\) against \(dt/t\) to obtain a constant
times \(\varepsilon_n=\Lambda_n^{-2}\).
\end{proof}

\begin{corollary}[Gaussian realization of the V155 telescope]
\label{cor:V156-V155-bridge}
If \(r=s-1\), the background and boundary families satisfy the uniform
parametrix card at that order, and the compatible state functionals obey
\eqref{eq:V155-state-summability}, then the Gaussian modulus contribution
to the state-evaluated Einstein source satisfies the shell hypothesis
\eqref{eq:V155-shell-summability}.
\end{corollary}

\begin{proof}
Theorem~\ref{thm:V156-shell-summability} is absolute summability in the
required \(H^{s-1}\) norm.  The complete Gaussian supertrace is already
formed before that norm is taken.  Apply
Theorem~\ref{thm:V155-shell-bridge} and its state part.
\end{proof}

\subsection{Necessity of the logarithmic subtraction}
\label{subsec:V156-necessity}

\begin{proposition}[An unsubtracted weight-four row is nonsummable]
\label{prop:V156-weight-four-necessary}
Suppose the complete metric variation \(\delta_gA_{4,q}\) has a nonzero
cutoff-stable limit and is not subtracted.  Then its contribution to every
dyadic shell is
\begin{equation}
 -\frac{\hbar}{2}\delta_gA_{4,q}
 \int_{\varepsilon_{n+1}}^{\varepsilon_n}\frac{dt}{t}
 =-\hbar(\log2)\delta_gA_{4,q},
 \label{eq:V156-log-shell}
\end{equation}
which does not tend to zero and cannot be absolutely summable.
\end{proposition}

\begin{proof}
The \(k=4\) power in \eqref{eq:V156-heat-expansion} is \(t^0\).
Since \(\varepsilon_n/\varepsilon_{n+1}=4\), the integral of \(dt/t\) over
the shell equals \(\log4=2\log2\).  Multiplication by
\(-\hbar/2\) gives \eqref{eq:V156-log-shell}.  A sequence bounded away from
zero cannot have a convergent sum of norms.
\end{proof}

\begin{corollary}[All divergent face rows belong to the common subtraction]
\label{cor:V156-face-subtractions}
Omitting a nonzero boundary coefficient with \(k<4\) is worse: its shell
size grows as a positive power of \(\Lambda_n\).  Thus area, GHY/scalar--GHY,
and all other allowed boundary divergences through weight four must be
included in the same invariant counterterm action before the one-loop
telescope is estimated.
\end{corollary}

\begin{proof}
The \(k\)-th heat term contributes
\[
 \int_{\varepsilon_{n+1}}^{\varepsilon_n}
 t^{(k-4)/2}\frac{dt}{t},
\]
whose magnitude has order \(\Lambda_n^{4-k}\) for \(k<4\).  Such a term
does not tend to zero.
\end{proof}

\begin{firewall}[Cancellation must precede subtraction and norm]
\label{fire:V156-complete-supertrace}
If the complete Standard Model supertrace makes some \(A_k\) vanish, that
is an exact cancellation and no counterterm is needed for that combined
coefficient.  Taking absolute values of field-sector heat traces first can
destroy the cancellation and manufacture a false divergence.  Conversely,
separate finite subtractions are admissible only if their sum is the first
variation of one invariant local action; otherwise the V150 Ward identity
need not survive.
\end{firewall}

\subsection{Which Sobolev topology has actually been acquired?}
\label{subsec:V156-topology}

\begin{proposition}[Remainder regularity determines source regularity]
\label{prop:V156-topology}
Let \(X\) be any Banach space of stress densities for which the uniform
remainder estimate
\begin{equation}
 \|\delta_gK_q^{>4}(t)\|_X\le C_Xt^{1/2}
 \label{eq:V156-X-remainder}
\end{equation}
holds.  Then the one-loop shells are absolutely summable in \(X\).  No
stronger topology follows from the decay exponent alone.
\end{proposition}

\begin{proof}
The proof of Theorem~\ref{thm:V156-shell-summability} uses only the norm in
\eqref{eq:V156-X-remainder}, so it applies verbatim to \(X\).  To disprove an
automatic upgrade, take any scalar sequence which is Cauchy in a weak or
negative Sobolev norm but unbounded in a stronger positive norm, and multiply
it by shell coefficients with summable scalar magnitudes.  Summability in
the weaker norm does not control the stronger one.
\end{proof}

\begin{assessment}[One-loop bridge acquired under a uniform parametrix]
\label{ass:V156-one-loop-status}
For smooth proper-time regularization, complete local subtraction through
weight four gives exactly the positive dyadic decay needed by the V155
telescope.  This is a genuine one-loop advance.  Its positive
\(H^{s-1}\) version requires uniform differentiated heat-kernel control at
that order.  The result does not estimate interacting connected BV
cumulants, construct compatible states, or prove uniform regularity of the
cofinal backgrounds.
\end{assessment}

\subsection{V156 ledger and next acquisition}
\label{subsec:V156-ledger}

\paragraph{New exact conclusions.}
\begin{enumerate}[label=\textup{(\arabic*)},leftmargin=2.8em]
\item Duhamel's formula identifies the complete one-loop metric derivative
      without assuming commutativity.
\item A uniform differentiated heat parametrix leaves an
      \(O(t^{1/2})\) boundary remainder after weights zero through four are
      removed.
\item Dyadic proper-time shells are consequently \(O(\Lambda^{-1})\) and
      absolutely summable; the closed-manifold rate is
      \(O(\Lambda^{-2})\).
\item A nonzero unsubtracted weight-four coefficient contributes the same
      logarithmic amount to every shell and is necessarily nonsummable.
\item The topology of the source telescope is exactly the topology in which
      the differentiated heat remainder is uniform.
\item Under positive-order parametrix bounds and compatible states, the
      Gaussian modulus realizes the V155 strong-source bridge.
\end{enumerate}

\paragraph{Next forced step.}
V157 should pass from the Gaussian identity to the first interacting BV
cumulant.  It should differentiate the complete logarithmic shell integral,
retain the connected covariance in
\eqref{eq:V155-second-log-derivative}, and seek an \(O(\Lambda^{-1-\epsilon})\)
or otherwise summable bound in the actual source topology.  Following the
V155/YM V133 audit, it must count genuine shell difference insertions rather
than diagram edges.  Failure at the first cumulant should be recorded as the
precise interacting obstruction, with no all-order claim.

\begin{assessment}[V156 continuum status]
\label{ass:V156-continuum-status}
The full Standard--Model--Einstein continuum remains unproved.  V156 closes
the smooth proper-time Gaussian shell summability gate conditionally on a
uniform differentiated parametrix and complete invariant subtraction.  The
interacting cumulant estimate, compatible state transport, nonlinear
background bounds, boundary realization, and coupled cofinal solution remain
open.
\end{assessment}

\section*{V156 exact checks and append-only record}
\addcontentsline{toc}{section}{V156 exact checks and append-only record}

The one-loop shell checksum is
\begin{equation}
 \boxed{
 \delta K^{>4}(t)=O_{H^r}(t^{1/2})
 \quad\Longrightarrow\quad
 \|\mathfrak s_n^{(1)}\|_{H^r}=O(\Lambda_n^{-1}),
 \qquad
 \sum_n\|\mathfrak s_n^{(1)}\|_{H^r}<\infty.}
 \label{eq:V156-final-check}
\end{equation}
The append operation preserves the complete V155 prefix byte for byte before
its former live terminator.  The assembled V156 file is checked for one live
final terminator, balanced environments and braces, duplicate labels, newly
introduced unresolved references, display math delimiters, and control
characters before the exact V155 predecessor is removed.  No companion file
is modified.

% =============================================================================
% END APPENDED VERSION V156 MATERIAL
% =============================================================================
\clearpage
% =============================================================================
% BEGIN APPENDED VERSION V157 MATERIAL
% =============================================================================

\section{V157: the first interacting shell derivative, invariant local RG
projection, and summability of the irrelevant remainder}
\label{sec:V157}

\subsection{Purpose}
\label{subsec:V157-purpose}

V156 proves dyadic summability after the complete Gaussian heat coefficients
through weight four have been subtracted.  Interactions add a familiar but
essential obstruction: a relevant or marginal local coupling can receive an
order-one increment on every logarithmic shell.  The raw shell derivative is
then not summable even though a renormalized continuum observable can exist
after the coupling is run.

This version corrects the V155 bridge accordingly.  The complete connected
shell derivative is first formed.  Its finite-dimensional invariant local
jet is absorbed into the running renormalized couplings.  Only the
irrelevant remainder is required to be absolutely summable.  The procedure
is algebraically exact and preserves the Ward identity when the local
projection commutes with the linearized Slavnov/diffeomorphism operator.
The analytic bound on the actual interacting remainder remains a named
gate.

\subsection{Exact first-cumulant metric derivative}
\label{subsec:V157-first-cumulant}

On one finite BV shell, split the action as
\begin{equation}
 S_\lambda=S_0+\lambda I,
 \label{eq:V157-action-split}
\end{equation}
where \(S_0\) includes the complete quadratic gauge-fixed action and
\(I\) is the assembled interaction, including every vertex required by the
BV master action.  Define
\begin{equation}
 \Gamma_\lambda:=-\hbar\log
 \int_{\mathcal L_{\rm shell}}e^{-S_\lambda/\hbar}.
 \label{eq:V157-effective-shell}
\end{equation}

\begin{lemma}[First interacting coefficient]
\label{lem:V157-first-coefficient}
At \(\lambda=0\),
\begin{equation}
 \partial_\lambda\Gamma_\lambda\big|_{\lambda=0}
 =\langle I\rangle_0.
 \label{eq:V157-first-coefficient}
\end{equation}
Its metric derivative is
\begin{equation}
 \boxed{
 \delta_g\langle I\rangle_0
 =\langle\delta_g I\rangle_0
 -\hbar^{-1}\operatorname{Cov}_0(I,\delta_gS_0),}
 \label{eq:V157-first-metric-derivative}
\end{equation}
where
\begin{equation}
 \operatorname{Cov}_0(A,B)
 :=\langle AB\rangle_0-\langle A\rangle_0\langle B\rangle_0.
 \label{eq:V157-covariance}
\end{equation}
\end{lemma}

\begin{proof}
Differentiate \eqref{eq:V157-effective-shell} in \(\lambda\) as in
Lemma~\ref{lem:V155-log-derivative}.  This gives
\eqref{eq:V157-first-coefficient}.  For a normalized expectation,
differentiate its numerator and denominator in the metric.  The explicit
variation of \(I\) gives the first term in
\eqref{eq:V157-first-metric-derivative}; variation of the Gaussian weight
and normalization gives minus \(\hbar^{-1}\) times
\eqref{eq:V157-covariance}.
\end{proof}

\begin{firewall}[One covariance, not two formal factors]
\label{fire:V157-one-covariance}
The two insertions in
\(\operatorname{Cov}_0(I,\delta_gS_0)\) can share the same shell
fluctuation.  Expanding them into diagrams can create several formal edges
without creating independent difference variables.  Following the V155
audit of YM V133, no product decay is charged until the connected covariance
has been rewritten in a derivative-faithful shell-difference form.
\end{firewall}

\subsection{Why raw marginal shells are not summable}
\label{subsec:V157-marginal}

Let \(\mathcal L_{\le4}\) be the finite-dimensional space of local bulk and
face first variations of gauge-, diffeomorphism-, and BV-compatible
operators of engineering dimension at most four.  It includes the
cosmological, Einstein--GHY, scalar--GHY, gauge kinetic, Higgs kinetic/mass/
quartic, Yukawa, gauge-fixing/BV-exact, and allowed boundary counterterm
rows, modulo identities and total derivatives.  Let
\begin{equation}
 \mathbf P_{\rm loc}:\mathcal X\to\mathcal L_{\le4}
 \label{eq:V157-local-projector}
\end{equation}
be a bounded projection in a chosen shell-source space \(\mathcal X\).

\begin{proposition}[Constant logarithmic shell obstruction]
\label{prop:V157-marginal-obstruction}
If for some nonzero \(\ell\in\mathcal L_{\le4}\),
\begin{equation}
 \mathbf P_{\rm loc}\mathfrak s_n\longrightarrow\ell,
 \label{eq:V157-marginal-limit}
\end{equation}
then \(\sum_n\|\mathfrak s_n\|_{\mathcal X}\) diverges.  Thus absolute
summability of the raw shell sources would incorrectly exclude a nonzero
beta function for a marginal coupling.
\end{proposition}

\begin{proof}
Boundedness of \(\mathbf P_{\rm loc}\) gives
\[
 \|\mathfrak s_n\|_{\mathcal X}
 \ge\|\mathbf P_{\rm loc}\mathfrak s_n\|_{\mathcal X}/
      \|\mathbf P_{\rm loc}\|.
\]
The right side is bounded below by a positive constant for all sufficiently
large \(n\), because of \eqref{eq:V157-marginal-limit}.  The series cannot
converge.
\end{proof}

\begin{assessment}[Interpretation]
\label{ass:V157-beta-interpretation}
The obstruction is not a failure of renormalization.  It says that adjacent
bare or fixed-coordinate actions need not approach one another in a
marginal local direction.  Their coordinates must be transported by the RG
recurrence before the irrelevant tail is compared.  V149 already established
that the finite coordinates \(\chi_R,\xi_H\), and the Standard Model
couplings are renormalization data rather than numbers selected by symmetry.
\end{assessment}

\subsection{Exact local-jet absorption}
\label{subsec:V157-local-absorption}

Choose a basis \(\{\mathcal O_a\}_{a=1}^N\) of invariant local action
functionals whose first variations span \(\mathcal L_{\le4}\).  Write the
running local action as
\begin{equation}
 C_n:=\sum_{a=1}^Nc_n^a\mathcal O_a.
 \label{eq:V157-running-local-action}
\end{equation}

\begin{theorem}[Algebraic RG telescope]
\label{thm:V157-RG-telescope}
Let \(\mathfrak s_n\) be the complete adjacent-shell first variation before
the new local counterterm is added.  Choose \(C_{n+1}-C_n\) so that
\begin{equation}
 \delta_g(C_{n+1}-C_n)
 =-\mathbf P_{\rm loc}\mathfrak s_n.
 \label{eq:V157-RG-recurrence}
\end{equation}
Then the renormalized adjacent-shell source is exactly
\begin{equation}
 \boxed{
 \mathfrak s_n^{\rm ren}
 =(I-\mathbf P_{\rm loc})\mathfrak s_n.}
 \label{eq:V157-ren-shell}
\end{equation}
If
\begin{equation}
 \sum_n\|(I-\mathbf P_{\rm loc})\mathfrak s_n\|_{\mathcal X}<\infty,
 \label{eq:V157-irrelevant-summability}
\end{equation}
then the renormalized effective metric derivatives are Cauchy in
\(\mathcal X\), irrespective of absolute summability of the local running
increments.
\end{theorem}

\begin{proof}
Adding \(C_n\) to the effective action adds \(\delta_gC_n\) to its source.
The difference of two adjacent renormalized sources is therefore
\[
 \mathfrak s_n+\delta_g(C_{n+1}-C_n).
\]
Insert \eqref{eq:V157-RG-recurrence} to obtain
\eqref{eq:V157-ren-shell}.  The difference from level \(M\) to level \(N\)
is the direct sum
\(\sum_{n=M}^{N-1}\mathfrak s_n^{\rm ren}\).  Applying the norm after each
complete projected shell has been formed and using
\eqref{eq:V157-irrelevant-summability} proves the Cauchy property.
\end{proof}

\begin{corollary}[Corrected V155 bridge hypothesis]
\label{cor:V157-corrected-bridge}
In Theorem~\ref{thm:V155-shell-bridge}, condition
\eqref{eq:V155-shell-summability} may be replaced by
\eqref{eq:V157-irrelevant-summability}, provided the comparison transports
also apply the running local-coordinate change
\eqref{eq:V157-RG-recurrence}.  The state and Ward hypotheses are otherwise
unchanged.
\end{corollary}

\begin{proof}
The proof of V155 uses only a Cauchy estimate for adjacent renormalized
sources.  Equation~\eqref{eq:V157-ren-shell} and
\eqref{eq:V157-irrelevant-summability} supply exactly that estimate.
\end{proof}

\subsection{Ward-compatible local projection}
\label{subsec:V157-Ward-projector}

Let \(\mathcal W\) denote the linearized combined Slavnov and
diffeomorphism Ward operator acting on first variations.  Absence of the
local and global anomaly classes was established in V138 and V141 on the
faithful Standard Model branch.

\begin{theorem}[Invariant local subtraction preserves the Ward row]
\label{thm:V157-Ward-preservation}
Assume
\begin{equation}
 \mathcal W\mathbf P_{\rm loc}
 =\mathbf P_{\rm loc}\mathcal W
 \label{eq:V157-Ward-commutation}
\end{equation}
on the complete shell first-variation space, and suppose
\begin{equation}
 \mathcal W\mathfrak s_n=\mathfrak a_n
 \label{eq:V157-shell-Ward-defect}
\end{equation}
with the complete anomaly/regulator defect \(\mathfrak a_n\).  Then
\begin{equation}
 \mathcal W\mathfrak s_n^{\rm ren}
 =(I-\mathbf P_{\rm loc})\mathfrak a_n.
 \label{eq:V157-ren-Ward-defect}
\end{equation}
In particular, an exact finite-shell Ward identity is preserved by the
local subtraction.
\end{theorem}

\begin{proof}
Apply \(\mathcal W\) to
\eqref{eq:V157-ren-shell}, commute it through the projection by
\eqref{eq:V157-Ward-commutation}, and use
\eqref{eq:V157-shell-Ward-defect}.
\end{proof}

\begin{firewall}[Component Taylor subtraction is not enough]
\label{fire:V157-component-Taylor}
A Taylor projector applied independently to gauge, Higgs, fermion,
improvement, or face components need not satisfy
\eqref{eq:V157-Ward-commutation}.  The local basis must consist of complete
invariant action functionals, with BRST-exact terms treated in the same BV
complex and bulk curvature terms paired with their required face terms.
Otherwise an apparently finite shell source can create a harmonic-constraint
source through \eqref{eq:V154-subsidiary-equation}.
\end{firewall}

\subsection{Taylor gain for the irrelevant remainder}
\label{subsec:V157-Taylor-gain}

The following elementary estimate records the analytic gain which a local
projection must realize in an external-momentum chart.

\begin{lemma}[Low-momentum Taylor remainder]
\label{lem:V157-Taylor-remainder}
Let \(k_n(p)\) be a Banach-valued shell kernel on \(|p|\le\mu\), and let
\(T_dk_n\) be its Taylor polynomial at \(p=0\) through degree \(d\).  If
\begin{equation}
 \sup_{|p|\le\mu}
 \|\partial_p^{d+1}k_n(p)\|
 \le C\Lambda_n^{-\varepsilon}
 \label{eq:V157-derivative-bound}
\end{equation}
for some \(\varepsilon\ge0\), then
\begin{equation}
 \sup_{|p|\le\mu}\|k_n(p)-T_dk_n(p)\|
 \le\frac{C}{(d+1)!}\mu^{d+1}Lambda_n^{-\varepsilon}.
 \label{eq:V157-Taylor-gain}
\end{equation}
If scaling supplies the sharper bound
\(C\Lambda_n^{-(d+1+\varepsilon)}\), the right side contains the summable
factor \((\mu/\Lambda_n)^{d+1}\Lambda_n^{-\varepsilon}\).
\end{lemma}

\begin{proof}
Apply the integral Taylor formula along the segment from zero to \(p\) and
take the Banach norm.  The length factor is at most \(\mu^{d+1}\), and the
simplex volume is \(1/(d+1)!\).
\end{proof}

\begin{assessment}[What remains to prove for the first cumulant]
\label{ass:V157-first-cumulant-gate}
The algebraic renormalization step is now complete.  The unresolved estimate
is derivative fidelity and scale decay for the assembled expression
\eqref{eq:V157-first-metric-derivative} after projection by
\(I-\mathbf P_{\rm loc}\).  A proof must show that actual shell difference
insertions yield the derivative bound in
\eqref{eq:V157-derivative-bound}, uniformly in the common source topology.
Diagram count or formal edge count does not prove it.
\end{assessment}

\subsection{V157 ledger and next acquisition}
\label{subsec:V157-ledger}

\paragraph{New exact conclusions.}
\begin{enumerate}[label=\textup{(\arabic*)},leftmargin=2.8em]
\item The first interacting metric derivative is an explicit insertion
      minus one connected covariance.
\item A nonzero marginal local shell increment makes raw absolute shell
      summability impossible.
\item Running invariant local counterterms absorb that increment exactly.
\item The corrected continuum criterion is summability of the irrelevant
      projected shell source.
\item A local projection commuting with the combined Ward operator preserves
      the complete finite-shell identity.
\item Taylor's theorem identifies the external-momentum gain needed from an
      actual derivative-faithful shell estimate.
\end{enumerate}

\paragraph{Next forced step.}
V158 should construct a concrete resampling or finite-difference
representation for the covariance in
\eqref{eq:V157-first-metric-derivative}.  It should prove how many distinct
shell variables are changed by each difference, test coincident insertions,
and derive the first projected momentum gain without splitting the
covariance into diagram norms.  If a single difference is all that the
identity contains, the resulting exponent must be accepted and compared
honestly with the dyadic summability threshold.

\begin{assessment}[V157 continuum status]
\label{ass:V157-continuum-status}
The full Standard--Model--Einstein continuum remains unproved.  V157 repairs
the V155 criterion in the presence of running marginal couplings and closes
the invariant algebraic RG telescope.  It does not prove the interacting
irrelevant-shell bound, compatible state transport, nonlinear background
estimates, or the coupled cofinal solution.
\end{assessment}

\section*{V157 exact checks and append-only record}
\addcontentsline{toc}{section}{V157 exact checks and append-only record}

The interacting-shell checksum is
\begin{equation}
 \boxed{
 \delta_g\langle I\rangle_0
 =\langle\delta_gI\rangle_0
 -\hbar^{-1}\operatorname{Cov}_0(I,\delta_gS_0),
 \qquad
 \mathfrak s_n^{\rm ren}
 =(I-\mathbf P_{\rm loc})\mathfrak s_n.}
 \label{eq:V157-final-check}
\end{equation}
The append operation preserves the complete V156 prefix byte for byte before
its former live terminator.  The assembled V157 file is checked for one live
final terminator, balanced environments and braces, duplicate labels, newly
introduced unresolved references, display math delimiters, and control
characters before the exact V156 predecessor is removed.  No companion file
is modified.

% =============================================================================
% END APPENDED VERSION V157 MATERIAL
% =============================================================================
\end{document}
